% !TeX program = xelatex

% BEGIN preamble.tex
% !TeX program = xelatex
\documentclass[UTF8,a4paper]{ctexbook}
% The class determines paper geometry, type size, Chinese fonts, heading layout,
% double-sided pagination and chapter opening. No geometry, fancyhdr or font override.
\usepackage{amsmath,amssymb,amsthm,mathtools,mathrsfs}
\usepackage{tikz-cd}
\usepackage{enumitem,longtable,booktabs,array,needspace}
\usepackage[perpage,bottom]{footmisc}
\usepackage{graphicx,adjustbox}
\usepackage[unicode,hidelinks,bookmarksnumbered,hypertexnames=false]{hyperref}
\usepackage{bookmark}
% Retain EGA's historical chapter and paragraph numbering.
\ctexset{part={name={第,卷},number=\Roman{part}},chapter={number={\ifnum\value{chapter}=0 0\else\Roman{chapter}\fi}},section={name={\texorpdfstring{\S\,}{§},}}}
\renewcommand{\thesection}{\arabic{section}}
\renewcommand{\thesubsection}{\thesection.\arabic{subsection}}
\setcounter{secnumdepth}{2}
\setcounter{tocdepth}{2}
\setlength{\emergencystretch}{2em}
\allowdisplaybreaks[2]
\newcommand{\egapar}[1]{\par\addvspace{0.55\baselineskip}\Needspace{3\baselineskip}\noindent\hypertarget{para:\arabic{chapter}:#1}{}\hypertarget{entry:\arabic{chapter}:#1}{}\label{entry:\arabic{chapter}:#1}\textbf{（#1）}\quad\ignorespaces}
\newcommand{\egastatement}[2]{\par\addvspace{0.55\baselineskip}\Needspace{3\baselineskip}\noindent\hypertarget{para:\arabic{chapter}:#1}{}\hypertarget{entry:\arabic{chapter}:#1}{}\label{entry:\arabic{chapter}:#1}\textbf{#2（#1）}\quad\ignorespaces}
\newcommand{\eref}[1]{\hyperlink{entry:\arabic{chapter}:#1}{（#1）}}
\newcommand{\erritem}[2]{\par\addvspace{0.6\baselineskip}\Needspace{3\baselineskip}\noindent\hypertarget{errata:I:#1}{}\textbf{#2（#1）}\quad\ignorespaces}
\newcommand{\unnumberedchapter}[1]{\chapter*{#1}\addcontentsline{toc}{chapter}{#1}\markboth{#1}{#1}}
\newcommand{\src}[1]{\par\nobreak\noindent{\small 原文印刷页码：#1。}\par\nobreak\vspace{0.25\baselineskip}}
\newcommand{\origfootnote}[1]{\ifhmode\unskip\nobreak\fi\footnote{\textbf{原注：}#1}}
\newcommand{\transnote}[1]{\ifhmode\unskip\nobreak\fi\footnote{\textbf{译注：}#1}}
\newcommand{\separation}{\par\medskip\centerline{$*\qquad *\qquad *$}\medskip}
\newcommand{\F}{\mathscr F}
\newcommand{\G}{\mathscr G}
\newcommand{\HH}{\mathscr H}
\newcommand{\B}{\mathfrak B}
\newcommand{\OO}{\mathscr O}
\newcommand{\Kcat}{\boldsymbol K}
\newcommand{\restrict}[2]{#1\vert_{#2}}
\newcommand{\restmod}[2]{#1_{[#2]}}
\DeclareMathOperator{\Hom}{Hom}
\DeclareMathOperator{\Supp}{Supp}
\DeclareMathOperator{\Ker}{Ker}
\DeclareMathOperator{\Ann}{Ann}
\DeclareMathOperator{\Spec}{Spec}
\newcommand{\ilim}{\varinjlim}
\newcommand{\plim}{\varprojlim}
\newcommand{\A}{\mathscr A}
\newcommand{\BB}{\mathscr B}
\newcommand{\C}{\mathscr C}
\newcommand{\E}{\mathscr E}
\newcommand{\LL}{\mathscr L}
\newcommand{\M}{\mathscr M}
\newcommand{\N}{\mathscr N}
\newcommand{\PP}{\mathscr P}
\newcommand{\J}{\mathscr J}
\newcommand{\shHom}{\mathscr{H}\!om}
\newcommand{\dual}[1]{\check{#1}}
\DeclareMathOperator{\Coker}{Coker}
\let\Im\relax
\DeclareMathOperator{\Im}{Im}
\DeclareMathOperator{\Tor}{Tor}
\newcommand{\frJ}{\mathfrak J}
\newcommand{\frK}{\mathfrak K}
\newcommand{\frQ}{\mathfrak Q}
\newcommand{\frT}{\mathfrak T}
\newcommand{\frR}{\mathfrak R}
\newcommand{\frS}{\mathfrak S}
\newcommand{\comp}[1]{\widehat{#1}}
\newcommand{\ctensor}{\mathbin{\widehat{\otimes}}}
\DeclareMathOperator{\grad}{grad}
\newcommand{\wt}[1]{\widetilde{#1}}
\newcommand{\jj}{\mathfrak j}
\newcommand{\rr}{\mathfrak r}
\newcommand{\mm}{\mathfrak m}
\newcommand{\nrad}{\mathfrak N}
\newcommand{\aphi}{{}^{a}\varphi}
\DeclareMathOperator{\Spf}{Spf}
\newcommand{\frX}{\mathfrak X}
\newcommand{\frY}{\mathfrak Y}
\newcommand{\frZ}{\mathfrak Z}
\newcommand{\frU}{\mathfrak U}
\newcommand{\frV}{\mathfrak V}
\newcommand{\frW}{\mathfrak W}
\newcommand{\frD}{\mathfrak D}
\newcommand{\SSS}{\mathscr S}
\newcommand{\nn}{\mathfrak n}
\newcommand{\pp}{\mathfrak p}
\newcommand{\qq}{\mathfrak q}
\newcommand{\frj}{\mathfrak j}
\newcommand{\frk}{\mathfrak k}
\DeclareMathOperator{\Proj}{Proj}
\DeclareMathOperator{\Tr}{Tr}
\DeclareMathOperator{\Sym}{S}
\DeclareMathOperator{\gr}{gr}
\newcommand{\red}{\mathrm{red}}
\newcommand{\Gstar}{\Gamma_*}
\newcommand{\Tstar}{\Gamma_*}
\newcommand{\R}{\mathscr R}
\newcommand{\TT}{\mathscr T}
\newcommand{\I}{\mathscr I}
\newcommand{\K}{\mathscr K}
\newcommand{\V}{\mathbf V}
\newcommand{\Pb}{\mathbf P}
\newcommand{\Z}{\mathbf Z}
\newcommand{\Gr}{\operatorname{Gr}}
\newcommand{\id}{\mathrm{id}}
\DeclareMathOperator{\codim}{codim}
\DeclareMathOperator{\rk}{rg}
\DeclareMathOperator{\Pic}{Pic}
\newcommand{\sProj}{\mathscr{P}\!roj}
\newcommand{\catHom}{\mathscr{H}\!om}
\newcommand{\Ens}{\mathbf{Ens}}
\newcommand{\OI}{0_{\mathrm I}}
\newcommand{\OIII}{0_{\mathrm{III}}}
\newcommand{\Ext}{\operatorname{Ext}}
\newcommand{\shExt}{\mathscr{E}\!xt}
\newcommand{\shTor}{\mathscr{T}\!or}
\newcommand{\Hh}{\mathbb H}
\newcommand{\len}{\operatorname{long}}
\newcommand{\cplx}[1]{#1^{\bullet}}
\newcommand{\Cat}{\mathbf C}
\newcommand{\ra}{\longrightarrow}
\newcommand{\la}{\longleftarrow}
\newcommand{\iso}{\xrightarrow{\sim}}
\newcommand{\eproof}{\par\noindent\textbf{证明。}\ignorespaces}
\newcommand{\finishproof}{\hfill$\square$\par}
\newcommand{\srcpage}[1]{\par\noindent{\small\itshape〔底本第#1页〕}\par\nobreak}
\newcommand{\cocomp}[1]{\complement #1}
\DeclareMathOperator{\Aut}{Aut}
\DeclareMathOperator{\GL}{GL}
\DeclareMathOperator{\trdeg}{trdeg}

% Macros for the reviewed continuation. Existing mathematical macros are unchanged.
\newcommand{\kk}{\mathfrak k}
\newcommand{\cH}{\mathscr H}
\newcommand{\hTor}{\mathbb{T}\!or}
\newcommand{\lTor}{\mathscr{T}\!or}
\newcommand{\spE}{\mathscr E}
\DeclareMathOperator{\coprof}{coprof}
\newcommand{\source}[1]{\par\smallskip\noindent\textit{底本位置：#1。}\par\smallskip}
\newcommand{\newscope}{new}
\newcommand{\egap}[1]{\par\addvspace{.6\baselineskip}\Needspace{3\baselineskip}\noindent\hypertarget{new:\newscope:#1}{}\textbf{（#1）}\quad\ignorespaces}
\newcommand{\egas}[2]{\par\addvspace{.6\baselineskip}\Needspace{3\baselineskip}\noindent\hypertarget{new:\newscope:#1}{}\textbf{#2（#1）}\quad\ignorespaces}
\newcommand{\cmt}[2]{\footnote{\textbf{译校注〔#1〕：}#2}}
\newcommand{\ednote}[1]{\footnote{\textbf{英译编者注：}#1}}
\newcommand{\orig}[1]{\footnote{\textbf{原注：}#1}}
\newcommand{\head}[1]{\chapter*{#1}\addcontentsline{toc}{chapter}{#1}\markboth{#1}{#1}}
\newcommand{\egasection}[1]{\section*{#1}\addcontentsline{toc}{section}{#1}}
\newcommand{\subsec}[1]{\subsection*{#1}\addcontentsline{toc}{subsection}{#1}}
\newcommand{\audit}[3]{\Needspace{5\baselineskip}\hypertarget{audit:#1}{}\section*{#1\quad #2}\addcontentsline{toc}{section}{#1\quad #2}\noindent\textbf{判定：#3。}\par\smallskip}
\pdfstringdefDisableCommands{\def\quad{ }\def\qquad{ }\def\S{§}}

\newcommand{\bTor}{\operatorname{\mathbf{Tor}}}
\newcommand{\termTor}{\mathop{\mathscr{T}\!or^{\mathrm{term}}}\nolimits}
\newcommand{\NewNote}[1]{\par\smallskip\noindent\textbf{编排说明。}#1\par\smallskip}
\hypersetup{pdftitle={代数几何基础：按内容顺序编排的中文译校合订本（第七版）},pdfauthor={原著 A. Grothendieck；与 J. Dieudonné 合作；预备笔记英译编者见各篇},pdfsubject={第0章归并；正文连续；已译范围和未译范围分开列明}}
\title{代数几何基础\\中文译校合订本\\[1em]\large 按内容顺序编排·第七版}
\author{A. Grothendieck\\与 J. Dieudonné 合作撰写\\[.6em]\normalsize 预备笔记英译及编辑：Piotr Blass、Joseph Blass}
\date{}

% Additions for the fourth continuation; no font or page-geometry overrides.
\DeclareMathOperator{\Ass}{Ass}
\DeclareMathOperator{\prof}{prof}
\DeclareMathOperator{\htp}{ht}
\newcommand{\pd}{\operatorname{dim.proj}}
\newcommand{\injdim}{\operatorname{dim.inj}}
\newcommand{\gldim}{\operatorname{dim.coh}}
\newcommand{\Salg}{\mathrm S}

\DeclareMathOperator{\End}{End}

% Additions for chapter 0, section 18. Keep class defaults.
\DeclareMathOperator{\Exan}{Exan}
\DeclareMathOperator{\Exal}{Exal}
\DeclareMathOperator{\Exalcom}{Exalcom}
\DeclareMathOperator{\Exantop}{Exantop}
\DeclareMathOperator{\Exaltop}{Exaltop}
\DeclareMathOperator{\Exalcotop}{Exalcotop}
\newcommand{\HomCont}{\operatorname{Hom.cont}}

% END preamble.tex

\begin{document}
\frontmatter

% BEGIN frontmatter.tex
\maketitle
\head{本版编排与收录说明}
\noindent\textbf{这是一份按内容顺序整理的中文工作译稿，不是全部已上传材料的完结全译本。}第I、II章及第III章的两个已出版分册正文已经连续收入；第0章当前收入§§1--18及§19的19.0--19.6；第IV章主体、第0章19.7--19.10及§§20--23尚未连续译入；第V卷英文预备笔记只收入已译出的部分。未译内容没有用原文扫描页、提要或空白标题代替译文。

\section*{正文次序}
本版不再按上传时间、译稿批次或独立PDF页序排列，而按原著内容编号组织：导言；第0章预备知识；第I章概形的语言；第II章若干类态射的初等整体研究；第III章凝聚层的上同调研究；后续待译范围说明；独立编排的第V卷预备笔记已译部分；最后为各原分册的书末资料。第III章第一部分的§§1--5与第二部分的§§6--7直接衔接，不再被第一部分的索引隔开。

第0章原来分载于第一卷、第三卷第一部分、第四卷第一部分。本版将已译的§§1--7、§§8--13、§§14--18和§19已译部分归并在一起，但保留 $0_{\mathrm I}$、$0_{\mathrm{III}}$、$0_{\mathrm{IV}}$ 等出处记号。它们标示原出版分册，并非重复的三章。这样的安排依内容编号，而不声称每一节的全部前提均已在前文出现；原著跨章引用仍按原号保留。

\section*{实际收录范围}
\begin{longtable}{@{}p{.25\linewidth}p{.67\linewidth}@{}}
\toprule 材料 & 已收入的连续译文及状态\\\midrule\endhead
第一卷（1960） & 导言；第0章§§1--7；第I章§§1--10，止于（10.15.7）；原书书末资料。承接既有V2校订源码。\\
第二卷（1961） & 第II章§§1--8，止于（8.14.14）；原书参考文献、索引、目录、勘误清单一。其书末资料沿用上一合订版依据底本恢复的译排来源。\\
第三卷第一部分（1961） & 第0章§§8--13；第III章§§1--5，止于（5.5.2）；原书书末资料。\\
第三卷第二部分（1963） & 第III章§§6--7，止于（7.9.14），连同其证明、记号索引、术语索引、原目录及勘误清单二。第三版已补齐（6.7.1）至本分册末尾；本版保留。\\
第四卷第一部分（1964） & 第0章补篇章首说明、§§14--18及§19的19.0--19.6，止于（19.6.7）末尾。本版新增（19.5.1）--（19.6.7）；19.7以后尚未译入。\\
第V卷英文预备笔记 & \texttt{EGAV1-2.ps} 的六页可读材料；\texttt{EGAVaiello.pdf} 的卷首说明及“超平面截面与锥投影”§§1--6。余下内容不计入已译范围。\\
\bottomrule
\end{longtable}

\section*{译文、勘误与核查的关系}
正文采用本轮已经核实的正确陈述、公式和证明；原文错误写法、修改理由及依据记入脚注。作者正式勘误与本次译者补证分开，较长反例与验证收入独立核查报告。尚未确认的疑点不擅自改成确定结论；原著本来省略的证明、预备笔记中未写完或无法辨认的内容，不凭通常知识补造成“原文”。本轮新增注号以G7开头，旧注号保留以便追踪。

这一编辑原则已落实在本轮新增的19.5--19.6，并回溯修改了（19.3.5，v）、（19.3.10.1）、（19.3.12）和（19.4.1）。\textbf{它不表示此前全部旧译文中的历史核查意见均已逐项回写正文。}其他既有部分继续保存原稿；其校订状态须参看各条脚注和历史报告，后续仍需逐项核定。

对作者已撤销的（19.3.12），正文仅保留编号和撤销说明，原陈述及问题写入脚注；完整历史陈述与原证明保存在报告附录，不再占据有效论述的正文位置。作者勘误清单二、清单三相关选译仍在书末收录。未经核实的替代结论没有补造。

第V卷上传材料的标题明确标为预备笔记的翻译与编辑，英译编者为Piotr Blass与Joseph Blass。它们包含重组建议、未完成句子及待补证明，不按正式刊行的第I--IV章同等状态标示。与英文合并稿重合的PS载体不重复计作新译文。

\section*{第七版新增范围与回溯修订}
本版从第六版实际接续点（19.5.1）起，译入19.5“形式光滑性与相伴分次环”以及19.6“域上的代数的情形”，止于（19.6.7）末尾。对应第四卷第一部分分册印刷第90--104页（上传PDF第89--103页）。本轮包括一般定理的全部分段证明、Cohen定理、系数域和几何正则性判据；19.7章首及（19.7.1）仍为下一接续点，未算作译完。

（19.6.2）在正文补明相容的子域嵌入 $k\subset A$；原文仅给 $k\subset K$ 的问题及 $\mathbf Z_p$ 反例放在脚注和报告中。（19.5.6.7）的核比较步骤换成有限滤过逐级消去论证；（19.6.6）以Galois作用在局部分量上的传递性及忠实平坦下降修补原证明。后两项是证明修补，不宣称原命题结论有反例。

第0章其余部分、第I--III章及预备笔记的旧正文，除上列回溯项外没有进行一轮新的逐句复审。本次新核查范围只包括新增内容及上述直接复核的相关条目。累计核查报告中“第六版”“第五版”等章节是历史档案，其中旧的正文状态须以本版实际源码为准。

\section*{版式与使用}
采用 \texttt{\string\documentclass[UTF8,a4paper]\{ctexbook\}}，由文档类决定默认中文字体、字号、版心、标题与双面分页。不另设页边距，不用自定义页眉。为保留原著结构，仅保留原章号、节号与公式号；默认双面排版所需的空白偶数页不删除。前置部分使用罗马数字页码，其后连续编号；总目录指向本合订本页面，各原著目录的页码仍指原版。

PDF可用于对照研读；可独立编译的单文件TeX与分文件工程同时提供。旧稿内容的编排归并与技术校验，不等于本次将旧稿的每一句译文或每一个证明重新独立审定。累计核查报告单独交付，历史核查意见带有批次说明；对同一事项，本版直接核得的作者勘误和明确论证须与历史意见对照阅读，不能将历史“待核”标签当作最新状态。
\tableofcontents

% END frontmatter.tex

\mainmatter

% BEGIN text/introduction.tex
% BEGIN chapters/introduction.tex
\chapter*{导言}
\addcontentsline{toc}{chapter}{导言}
\markboth{导言}{导言}
\src{5—9}
\begin{flushright}
献给 Oscar Zariski 与 André Weil。
\end{flushright}

本篇论文以及计划陆续发表的许多后续论文，旨在构成一部关于代数几何基础的专著。
原则上，我们不预设读者对这一学科有任何专门知识。事实甚至表明，虽然这类知识显然有其益处，
但有时也可能妨碍读者熟悉本书所阐述的观点与技术，因为它容易使人养成过于偏重双有理观点的习惯。
另一方面，我们将假定读者对下列主题具有良好的知识基础：
\begin{enumerate}[label=\alph*)]
\item 交换代数。例如 N. Bourbaki 的《数学原本》中正在编写的有关各卷所阐述的内容；
在这些卷册出版以前，可参阅 Samuel—Zariski\cite{b13} 以及 Samuel\cite{b11,b12}。
\item 同调代数。我们参照 Cartan—Eilenberg\cite{b2}（记作 (M)）、Godement\cite{b4}（记作 (G)），
以及 A. Grothendieck 最近的论文\cite{b6}（记作 (T)）。
\item 层论。主要参考文献仍为 (G) 与 (T)。这一理论提供了不可或缺的语言，使我们能够用“几何”术语
解释交换代数的基本概念，并将这些概念“整体化”。
\item 最后，读者若对函子语言有所熟悉，将会很有帮助；本专著将不断使用这种语言。
有关内容可参阅 (M)、(G)，尤其是 (T)。这一语言的原理以及函子一般理论的主要结果，
将在本专著两位作者正在编写的另一部著作中作更详细的阐述。
\end{enumerate}

\separation

在这篇导言中，我们不打算对代数几何中的“概形”观点作或详或略的描述，也不打算列出使我们不得不采用
这一观点的种种理由，特别是为何有必要系统地允许所研究的“簇”的局部环中存在幂零元。
这必然使有理映射的概念退居次位，而以正则映射或“态射”的概念居于首位。
本专著的目的恰在于系统地发展“概形”的语言，并且——我们希望——显示出它的必要性。
虽然这样做并不困难，我们在这里也不尝试对第一章所发展的概念作“直观的”引导。
希望预先了解本专著内容的读者，可以参阅 A. Grothendieck 在 1958 年 Edinburgh 国际数学家大会上的报告\cite{b7}，
以及同一作者的报告\cite{b8}。
J.-P. Serre 的论文\cite{b14}（记作 (FAC)）也可以看作介于代数几何的经典观点与概形观点之间的一种阐述；
就此而言，阅读该文可为阅读我们的《基础》作很好的准备。

\separation

为供参考，下面列出本专著预定的总体计划。这一计划以后仍可能修改，最后几章尤其如此：
\begin{enumerate}[label=第\Roman*章,align=left,leftmargin=5.3em,labelwidth=4.8em]
\item 概形的语言。
\item 若干类态射的初等整体研究。
\item 凝聚代数层的上同调及其应用。
\item 态射的局部研究。
\item 构造概形的初等方法。
\item 下降技术；构造概形的一般方法。
\item 群概形与主纤维空间。
\item 纤维空间的微分研究。
\item 基本群。
\item 留数与对偶性。
\item 相交理论、陈类与 Riemann—Roch 定理。
\item 阿贝尔概形与 Picard 概形。
\item Weil 上同调。
\end{enumerate}

原则上，各章都被视为开放的，随时可以在以后增加补充的节。
为了减少所采用的出版方式带来的不便，这些补充节将以单独的分册刊行。
当某章出版时，如果某个补充节已经列入计划或正在准备，即使出于轻重缓急的考虑，
其实际出版时间可能明显较晚，我们仍将在该章的目录中列出它。
为方便读者，我们在“第 0 章”中汇集一些关于交换代数、同调代数与层论的补充知识。
这些知识将在本专著各章中使用，虽然它们在不同程度上已为人所知，我们却未能找到便于引用的文献。
建议读者只在阅读本专著正文的过程中，并且只在对所引用的结果不够熟悉时，才查阅第 0 章。
我们还认为，这种读法会使本专著成为初学者熟悉交换代数和同调代数的一种良好途径：
倘若不伴随具体可感的应用，相当多的人会觉得这两门学科的学习枯燥，甚至令人沮丧。

\separation

在这篇导言中，即使只是对所阐述的概念和结果作简要的历史概述，也超出了我们的能力范围。
正文只列出我们认为对理解内容特别有用的文献，且只说明最重要结果的来源。
至少从形式上看，本书所论述的主题还相当新；这也解释了为什么我们很少引用十九世纪及二十世纪初
代数几何奠基者的工作——我们对这些工作的了解，仅来自转述。
不过，这里仍应略述那些最直接影响作者，并促成概形观点发展的著作。

首先必须提到 J.-P. Serre 的奠基性论文 (FAC)。对于不少被 A. Weil 经典著作
\textit{Foundations}\cite{b18} 的艰涩所吓退的年轻学习者而言，它是进入代数几何的门径；
本专著的一位作者也在其中。
该文第一次证明：“抽象”代数簇的“Zariski 拓扑”完全适合应用代数拓扑的某些技术，
尤其能够由此建立上同调理论。
此外，该文所给的代数簇定义，是最自然地适合扩展为我们在这里发展的概念的定义。\origfootnote{
据 J.-P. Serre 告知，应当指出：通过给定一个环层来定义簇的结构这一思想源于 H. Cartan，
他把这一思想作为其解析空间理论的出发点。当然，与代数几何中一样，在“解析几何”中，
也应当允许解析空间的局部环中出现幂零元。H. Grauert 最近已着手研究 H. Cartan 与 J.-P. Serre
定义的这种扩展\cite{b5}；我们有理由希望，不久将会出现一部在这一一般框架下系统阐述解析几何的著作。
显然，本专著所发展的概念和技术在解析几何中仍有意义，尽管应当预料到，后一理论中的技术困难将更为显著。
可以预见，代数几何由于其方法较为简单，能够为解析空间理论今后的发展提供某种形式上的模型。}

Serre 自己也曾指出：只要把域上的仿射代数换成任意交换环，仿射代数簇的上同调理论就可以毫无困难地转写过去。
因此，本专著的第一章、第二章以及第三章的头两节，就其主要内容而言，可以看作 (FAC) 以及同一作者
稍后的一篇论文\cite{b15} 中主要结果在这一扩大了的框架中的直接转写。
我们还从 C. Chevalley 的代数几何讨论班\cite{b1} 中获益良多；特别是系统地使用他引入的“可构造集”，
在概形理论中显示出很大的效用（参阅第四章）。
我们也借用了他从维数观点对态射所作的研究（第四章）；这一研究几乎无须作显著改动便可以转写到概形框架中。

另一方面，应当指出：Chevalley 引入的“局部环概形”概念，自然地适合用来扩展代数几何，
虽然它尚不具有我们在此希望赋予这一学科的全部灵活性和一般性。
关于这一概念与本书理论的关系，参见第一章第 8 节。
M. Nagata 已经在一系列论文\cite{b9} 中发展了这样的扩展，其中包含关于 Dedekind 环上的代数几何的许多特殊结果。\origfootnote{
在接近我们代数几何观点的工作中，应当提及 E. Kähler 的重要论文\cite{b22}，
以及 Chow 与 Igusa 最近的一篇短文\cite{b3}；后者在 Nagata—Chevalley 理论的框架内重新建立了 (FAC) 的若干结果，
并给出了一个 Künneth 公式。}

\separation

最后，不言而喻，一部关于代数几何的著作，尤其是一部关于其基础的著作，必然受到 O. Zariski 与 A. Weil
这样一些数学家的影响，哪怕这种影响是经由他人传递的。
特别是 Zariski 的全纯函数理论\cite{b20}，在借助上同调方法作适当调整，并补入一个存在定理
（第三章第 4、5 节）之后，与第六章阐述的下降技术一起，成为本专著所用的主要工具之一；
在我们看来，它也是代数几何现有的最强有力的工具之一。

这一理论所处的一般技术框架可以概述如下。第九章对基本群的研究将提供一个典型例子。
设有一个代数簇到另一个代数簇的固有态射（第二章）$f\colon X\to Y$，更一般地，也可以是一个概形到另一个概形的固有态射。
我们希望在一点 $y\in Y$ 的附近研究它，以解决一个关于 $y$ 的某邻域的问题 $P$。
研究按下列步骤进行：
\begin{enumerate}[label=\arabic*$^\circ$]
\item 可以假定 $Y$ 是仿射的，于是 $X$ 成为定义在 $Y$ 的仿射环 $A$ 上的概形；
甚至可以用 $y$ 的局部环来代替 $A$。在实际应用中，这一化归总是容易做到的（第五章），
它把问题归结为 $A$ 是局部环的情形。
\item 在 $A$ 是阿廷局部环时研究所考虑的问题。
为了使问题在不假定 $A$ 为整环时仍然确有意义，有时需要重新表述问题 $P$；
这样做往往反而能够使我们更好地理解问题，在这一阶段，这种理解具有“无穷小”的性质。
\item 形式概形理论（第三章第 3、4、5 节）使我们能够从阿廷环的情形过渡到完备局部环的情形。
\item 最后，若 $A$ 是任意局部环，则在 $X$ 上适当的概形中，考虑逼近某个给定“形式”截面的“多值截面”
（第四章），常常可以把对 $X$ 向 $A$ 的完备化作标量扩张后所得概形已经知道的结果，
转化为关于 $A$ 的一个足够简单的有限扩张——例如非分歧扩张——的类似结果。
\end{enumerate}

这个概述表明了系统研究阿廷环 $A$ 上的概形的重要性。
Serre 在表述局部类域论时所采用的观点，以及 Greenberg 最近的工作，似乎提示：
可以通过以函子方式给这样的概形 $X$ 配备一个定义在 $A$ 的剩余域 $k$ 上的概形 $X'$，来开展这种研究。
这里假定 $k$ 是完美域；在有利的情形下，$X'$ 的维数等于 $n\dim X$，其中 $n$ 为 $A$ 的长度。

至于 A. Weil 的影响，我们只须说明：写作本专著的主要动机之一，是需要发展必要的工具，
以便在所要求的一般性下表述“Weil 上同调”的定义，并着手证明\origfootnote{
为避免任何误解，须明确说明：在写这篇导言时，这项工作才刚刚开始，因此尚未得到 Weil 猜想的证明。}
为建立其著名的丢番图几何猜想\cite{b19} 所需的一切形式性质。
与此同样重要的动机，是希望为代数几何中的常用概念和方法找到自然的框架，
并使作者有机会理解这些概念与技术。

\separation

在结束这篇导言时，我们认为有必要预先提醒读者：与作者自己一样，读者大概也会经历一些困难，
才能习惯概形的语言，并确信几何直观所提示的通常构造，基本上都可以用唯一合理的方式转写成这种语言。
如同现代数学的许多领域一样，初始的直观与能够以所需的全部精确性和一般性来表达它的语言，
表面上正变得越来越远。
在这里，心理上的困难在于：我们必须把对集合而言熟悉的概念——笛卡儿积、群运算、环运算、模运算、
纤维空间、主齐性纤维空间等等——搬到一个已经颇不同于集合范畴的范畴的对象上，
也就是预概形范畴，或者给定预概形上的预概形范畴。
今后的数学家大概很难回避这种新的抽象努力；不过，归根结底，与我们的前辈为熟悉集合论而付出的努力相比，
它也许还相当轻微。

\separation

本书采用十进式层级编号给出引用。
例如在 III, 4.9.3 中，III 表示章，4 表示节，9 表示该节内的小节。
在同一章内部，省略章号。


% END chapters/introduction.tex
% END text/introduction.tex


% BEGIN text/ch0_01_07.tex
\renewcommand{\B}{\mathfrak B}
\setcounter{chapter}{-1}
\chapter{预备知识}

% BEGIN chapters/ch0_s1.tex
\section{分式环}\hypertarget{heading:0:1}{}
\setcounter{subsection}{-1}
\subsection{环与代数}\hypertarget{heading:0:1.0}{}
\src{11—12}

\egapar{1.0.1}
本专著所考虑的一切环都具有单位元；这样的环上的一切模都假定为酉模；
环同态总是假定将单位元映为单位元。
除非明确作出相反说明，环 $A$ 的子环总是假定包含 $A$ 的单位元。
我们主要考虑交换环；当不加限定地说到环时，即默认为交换环。
若 $A$ 是一个不一定交换的环，则除非明确作出相反说明，$A$-模总是指左模。

\egapar{1.0.2}
设 $A,B$ 为两个不一定交换的环，$\varphi\colon A\to B$ 为同态。
任一左 $B$-模（分别地，右 $B$-模）$M$ 都可以通过规定
\[
 a\cdot m=\varphi(a)\cdot m
 \qquad\text{（分别地， }m\cdot a=m\cdot\varphi(a)\text{）}
\]
而赋予左 $A$-模（分别地，右 $A$-模）结构。
需要区别 $M$ 上的 $A$-模结构与 $B$-模结构时，将如此定义的左（分别地，右）$A$-模记为 $M_{[\varphi]}$。
因此，若 $L$ 为 $A$-模，则同态 $u\colon L\to M_{[\varphi]}$ 就是一个交换群同态，满足
\[
 u(a\cdot x)=\varphi(a)\cdot u(x)\qquad(a\in A,\ x\in L).
\]
也称它为一个 $\varphi$-同态 $L\to M$，并称二元组 $(\varphi,u)$
（或略不严格地说，$u$）为从 $(A,L)$ 到 $(B,M)$ 的\emph{二重同态}
（\textit{di-homomorphisme}）。
由环 $A$ 及 $A$-模 $L$ 组成的二元组 $(A,L)$，因而构成一个以二重同态为态射的范畴。

\egapar{1.0.3}
在（\hyperlink{para:0:1.0.2}{1.0.2}）的假设下，若 $\mathfrak I$ 为 $A$ 的左理想（分别地，右理想），
则以 $B\mathfrak I$（分别地，$\mathfrak I B$）表示由 $\varphi(\mathfrak I)$ 生成的
$B$ 的左理想 $B\varphi(\mathfrak I)$（分别地，右理想 $\varphi(\mathfrak I)B$）。
它也是左 $B$-模（分别地，右 $B$-模）的典范同态
\[
 B\otimes_A\mathfrak I\longrightarrow B
 \qquad\text{（分别地， }\mathfrak I\otimes_A B\longrightarrow B\text{）}
\]
的像。

\egapar{1.0.4}
若 $A$ 为交换环，而 $B$ 为不一定交换的环，则在 $B$ 上给定一个 $A$-代数结构，
等价于给定一个环同态 $\varphi\colon A\to B$，使得 $\varphi(A)$ 包含在 $B$ 的中心内。
这时，对 $A$ 的任一理想 $\mathfrak I$，$\mathfrak I B=B\mathfrak I$ 是 $B$ 的双边理想；
对任一 $B$-模 $M$，$\mathfrak I M$ 是一个 $B$-模，并且等于 $(B\mathfrak I)M$。

\egapar{1.0.5}
我们不再重述有限型模及有限型交换代数的概念。
说 $A$-模 $M$ 是有限型的，就是说存在正合列
\[
 A^p\longrightarrow M\longrightarrow0.
\]
若一个 $A$-模 $M$ 同构于某个同态 $A^p\to A^q$ 的余核，换言之，若存在正合列
\[
 A^p\longrightarrow A^q\longrightarrow M\longrightarrow0,
\]
则称 $M$ 具有\emph{有限表示}。
注意，在诺特环 $A$ 上，每个有限型 $A$-模都具有有限表示。

回忆一下，若 $A$-代数 $B$ 的每个元素都是某个以 $A$ 中元素为系数的首一多项式在 $B$ 中的根，
则称 $B$ 在 $A$ 上是\emph{整的}。
这也等价于说，$B$ 的每个元素都包含在 $B$ 的某个子代数中，而这个子代数是有限型 $A$-模。
当此条件成立且 $B$ 交换时，由 $B$ 的任一有限子集生成的子代数都是有限型 $A$-模。
因此，交换代数 $B$ 在 $A$ 上既整又有限型，当且仅当 $B$ 是有限型 $A$-模；
此时也称 $B$ 为\emph{有限整 $A$-代数}，或在不致混淆时简称为\emph{有限 $A$-代数}。
须注意，这些定义并不要求定义 $A$-代数结构的同态 $A\to B$ 是单射。

\egapar{1.0.6}
\emph{整环}是指满足下述条件的环：任一有限族非零元素的乘积都非零。
等价地说，在这样的环中，$0\ne1$，且两个非零元素的乘积非零。
环 $A$ 的\emph{素理想}是使 $A/\mathfrak p$ 成为整环的理想 $\mathfrak p$；
这特别蕴含 $\mathfrak p\ne A$。
环 $A$ 至少具有一个素理想，当且仅当 $A\ne\{0\}$。

\egapar{1.0.7}
\emph{局部环}是恰有一个极大理想的环 $A$；这个极大理想就是可逆元集合的补集，
并包含所有不等于 $A$ 的理想。
若 $A,B$ 为两个局部环，$\mathfrak m,\mathfrak n$ 分别为其极大理想，
则称同态 $\varphi\colon A\to B$ 是\emph{局部的}，若
$\varphi(\mathfrak m)\subseteq\mathfrak n$，或等价地，
$\varphi^{-1}(\mathfrak n)=\mathfrak m$。
通过取商，这样的同态定义了一个从剩余域 $A/\mathfrak m$ 到剩余域 $B/\mathfrak n$ 的单同态。
两个局部同态的复合仍为局部同态。

\subsection{理想的根；环的幂零根与根}\hypertarget{heading:0:1.1}{}
\src{12}

\egapar{1.1.1}
设 $\mathfrak a$ 为环 $A$ 的理想。
$\mathfrak a$ 的\emph{根}记为 $\mathfrak r(\mathfrak a)$，是所有满足
$x^n\in\mathfrak a$（对至少一个整数 $n>0$）的元素 $x\in A$ 组成的集合；
它是包含 $\mathfrak a$ 的理想。
有 $\mathfrak r(\mathfrak r(\mathfrak a))=\mathfrak r(\mathfrak a)$；
$\mathfrak a\subseteq\mathfrak b$ 蕴含
$\mathfrak r(\mathfrak a)\subseteq\mathfrak r(\mathfrak b)$；
有限个理想之交的根等于这些理想的根之交。
若 $\varphi$ 为从环 $A'$ 到 $A$ 的同态，则对任一理想 $\mathfrak a\subseteq A$，有
\[
 \mathfrak r\bigl(\varphi^{-1}(\mathfrak a)\bigr)
 =\varphi^{-1}\bigl(\mathfrak r(\mathfrak a)\bigr).
\]
一个理想是某个理想的根，当且仅当它是素理想的交。
理想 $\mathfrak a$ 的根等于包含 $\mathfrak a$ 的素理想中所有极小者之交；
若 $A$ 是诺特环，则这些极小素理想只有有限个。

理想 $(0)$ 的根也称为 $A$ 的\emph{幂零根}；它是 $A$ 的所有幂零元组成的集合 $\mathfrak N$。
若 $\mathfrak N=(0)$，则称环 $A$ 是\emph{约化的}。
对任意环 $A$，它对其幂零根的商环 $A/\mathfrak N$ 都是约化环。

\egapar{1.1.2}
回忆一下，环 $A$（不一定交换）的\emph{根} $\mathfrak R(A)$，是 $A$ 的所有极大左理想的交，
也等于所有极大右理想的交。
$A/\mathfrak R(A)$ 的根为 $(0)$。

\subsection{分式模与分式环}\hypertarget{heading:0:1.2}{}
\src{13—14}

\egapar{1.2.1}
若环 $A$ 的子集 $S$ 满足 $1\in S$，且 $S$ 中任意两个元素的乘积仍在 $S$ 中，
则称 $S$ 是\emph{乘法子集}。
以下最重要的例子是：
$1^\circ$，某个元素 $f\in A$ 的所有幂 $f^n$（$n\ge0$）组成的集合 $S_f$；
$2^\circ$，$A$ 的一个素理想 $\mathfrak p$ 的补集 $A\setminus\mathfrak p$。

\egapar{1.2.2}
设 $S$ 为环 $A$ 的乘法子集，$M$ 为 $A$-模。
在集合 $M\times S$ 中，规定二元组 $(m_1,s_1)$ 与 $(m_2,s_2)$ 之间的关系为
\[
 \text{“存在 }s\in S\text{，使得 }s(s_2m_1-s_1m_2)=0\text{”}.
\]
这是一个等价关系。
以 $S^{-1}M$ 表示 $M\times S$ 对此关系的商集，以 $m/s$ 表示二元组 $(m,s)$ 在 $S^{-1}M$ 中的典范像；
称映射 $i_M^S\colon m\mapsto m/1$（也记为 $i^S$）为从 $M$ 到 $S^{-1}M$ 的\emph{典范映射}。
这个映射通常既不是单射也不是满射；其核由那些满足 $sm=0$（对某个 $s\in S$）的 $m\in M$ 组成。

在 $S^{-1}M$ 中，通过规定
\[
 \frac{m_1}{s_1}+\frac{m_2}{s_2}
 =\frac{s_2m_1+s_1m_2}{s_1s_2}
\]
定义加法群运算；可以验证，这与所考虑的 $S^{-1}M$ 中元素的表达方式无关。
在 $S^{-1}A$ 上，再通过
\[
 \frac{a_1}{s_1}\frac{a_2}{s_2}=\frac{a_1a_2}{s_1s_2}
\]
定义乘法。最后，在 $S^{-1}M$ 上以 $S^{-1}A$ 为算子集合，通过
\[
 \frac a s\frac m{s'}=\frac{am}{ss'}
\]
定义外部运算。
于是可以验证，$S^{-1}A$ 被赋予一个环结构，称为 $A$ 的\emph{以 $S$ 中元素为分母的分式环}；
$S^{-1}M$ 被赋予一个 $S^{-1}A$-模结构，称为 $M$ 的\emph{以 $S$ 中元素为分母的分式模}。
对每个 $s\in S$，$s/1$ 在 $S^{-1}A$ 中可逆，其逆元为 $1/s$。
典范映射 $i_A^S$ 是环同态，$i_M^S$ 是 $A$-模同态；
这里借助同态 $i_A^S\colon A\to S^{-1}A$ 将 $S^{-1}M$ 看作 $A$-模。

\egapar{1.2.3}
若对某个 $f\in A$ 有 $S=S_f=\{f^n\}_{n\ge0}$，则将 $S^{-1}A,S^{-1}M$ 分别记为 $A_f,M_f$。
把 $A_f$ 看作 $A$ 上的代数时，也可以写成 $A_f=A[1/f]$。
$A_f$ 同构于商代数 $A[T]/(fT-1)A[T]$。
当 $f=1$ 时，$A_f,M_f$ 分别典范地等同于 $A,M$；若 $f$ 是幂零元，则 $A_f,M_f$ 均为零。

当 $S=A\setminus\mathfrak p$，其中 $\mathfrak p$ 为 $A$ 的素理想时，
将 $S^{-1}A,S^{-1}M$ 分别记为 $A_{\mathfrak p},M_{\mathfrak p}$。
$A_{\mathfrak p}$ 是局部环，其极大理想 $\mathfrak q$ 由 $i_A^S(\mathfrak p)$ 生成，且
$(i_A^S)^{-1}(\mathfrak q)=\mathfrak p$。
通过取商，$i_A^S$ 给出从整环 $A/\mathfrak p$ 到域 $A_{\mathfrak p}/\mathfrak q$ 的单同态；
后者等同于 $A/\mathfrak p$ 的分式域。

\egapar{1.2.4}
分式环 $S^{-1}A$ 连同典范同态 $i_A^S$ 是一个\emph{泛映射问题}的解：
凡是从 $A$ 到环 $B$ 的同态 $u$，只要 $u(S)$ 中的元素在 $B$ 中均可逆，就唯一地分解为
\[
 u\colon A\xrightarrow{i_A^S}S^{-1}A\xrightarrow{u^*}B,
\]
其中 $u^*$ 为环同态。
在相同假设下，设 $M$ 为 $A$-模，$N$ 为 $B$-模，$v\colon M\to N$ 为 $A$-模同态；
这里使用 $u\colon A\to B$ 在 $N$ 上诱导的 $A$-模结构。\transnote{
原版第 14 页括号内印作“由 $u\colon A\to B$ 定义的 $N$ 上的 $B$-模结构”。
结合本句的“$A$-模同态”以及（\hyperlink{para:0:1.0.2}{1.0.2}），此处应为限制标量所得的 $A$-模结构，译文据此作了明示的记号校订。}
则 $v$ 唯一地分解为
\[
 v\colon M\xrightarrow{i_M^S}S^{-1}M\xrightarrow{v^*}N,
\]
其中 $v^*$ 是 $S^{-1}A$-模同态；$N$ 上的 $S^{-1}A$-模结构由 $u^*$ 给定。

\egapar{1.2.5}
令元素 $(a/s)\otimes m$ 对应于 $(am)/s$，就定义了一个 $S^{-1}A$-模的典范同构
\[
 S^{-1}A\otimes_A M\xrightarrow{\sim}S^{-1}M.
\]
其逆同构将 $m/s$ 映为 $(1/s)\otimes m$。

\egapar{1.2.6}
对 $S^{-1}A$ 的每个理想 $\mathfrak a'$，$\mathfrak a=(i_A^S)^{-1}(\mathfrak a')$ 是 $A$ 的理想，
且 $\mathfrak a'$ 是 $S^{-1}A$ 中由 $i_A^S(\mathfrak a)$ 生成的理想，它等同于 $S^{-1}\mathfrak a$（\hyperlink{para:0:1.3.2}{1.3.2}）。
映射 $\mathfrak p'\mapsto(i_A^S)^{-1}(\mathfrak p')$，是从 $S^{-1}A$ 的素理想集合到
$A$ 中满足 $\mathfrak p\cap S=\varnothing$ 的素理想集合的序同构。
此外，这时局部环 $A_{\mathfrak p}$ 与 $(S^{-1}A)_{S^{-1}\mathfrak p}$ 典范同构（\hyperlink{para:0:1.5.1}{1.5.1}）。

\egapar{1.2.7}
若 $A$ 为整环，以 $K$ 表示其分式域，则对每个不含 $0$ 的乘法子集 $S$，
典范映射 $i_A^S\colon A\to S^{-1}A$ 都是单射；
于是 $S^{-1}A$ 典范地等同于 $K$ 的一个包含 $A$ 的子环。
特别地，对 $A$ 的任一素理想 $\mathfrak p$，$A_{\mathfrak p}$ 是包含 $A$ 的局部环，
其极大理想为 $\mathfrak pA_{\mathfrak p}$，并有 $\mathfrak pA_{\mathfrak p}\cap A=\mathfrak p$。

\egapar{1.2.8}
若 $A$ 是约化环（\hyperlink{para:0:1.1.1}{1.1.1}），则 $S^{-1}A$ 也约化。
事实上，若 $x\in A,s\in S$ 且 $(x/s)^n=0$，则存在 $s'\in S$ 使 $s'x^n=0$，
从而 $(s'x)^n=0$。由假设得 $s'x=0$，所以 $x/s=0$。

\subsection{函子性质}\hypertarget{heading:0:1.3}{}
\src{14—15}

\egapar{1.3.1}
设 $M,N$ 为两个 $A$-模，$u\colon M\to N$ 为 $A$-同态。
若 $S$ 为 $A$ 的乘法子集，则由
\[
 (S^{-1}u)(m/s)=u(m)/s
\]
定义一个 $S^{-1}A$-同态 $S^{-1}M\to S^{-1}N$，记为 $S^{-1}u$。
若将 $S^{-1}M,S^{-1}N$ 分别典范地等同于 $S^{-1}A\otimes_A M,S^{-1}A\otimes_A N$（\hyperlink{para:0:1.2.5}{1.2.5}），
则 $S^{-1}u$ 等同于 $1\otimes u$。
若 $P$ 为第三个 $A$-模，$v\colon N\to P$ 为 $A$-同态，则
\[
 S^{-1}(v\circ u)=(S^{-1}v)\circ(S^{-1}u).
\]
换言之，在固定 $A,S$ 时，$S^{-1}M$ 是关于 $M$ 的协变函子，
从 $A$-模范畴取值于 $S^{-1}A$-模范畴。

\egapar{1.3.2}
函子 $S^{-1}M$ 是正合的；就是说，若列
\[
 M\xrightarrow{u}N\xrightarrow{v}P
\]
正合，则列
\[
 S^{-1}M\xrightarrow{S^{-1}u}S^{-1}N\xrightarrow{S^{-1}v}S^{-1}P
\]
也正合。
特别地，若 $u\colon M\to N$ 为单射（分别地，满射），则 $S^{-1}u$ 也为单射（分别地，满射）。
若 $N,P$ 是 $M$ 的两个子模，则 $S^{-1}N,S^{-1}P$ 典范地等同于 $S^{-1}M$ 的子模，且有
\[
 \begin{aligned}
 S^{-1}(N+P)&=S^{-1}N+S^{-1}P,\\
 S^{-1}(N\cap P)&=(S^{-1}N)\cap(S^{-1}P).
 \end{aligned}
\]

\egapar{1.3.3}
设 $(M_\alpha,\varphi_{\beta\alpha})$ 为 $A$-模的一个归纳系，
则 $(S^{-1}M_\alpha,S^{-1}\varphi_{\beta\alpha})$ 是 $S^{-1}A$-模的归纳系。
将 $S^{-1}M_\alpha$ 及 $S^{-1}\varphi_{\beta\alpha}$ 写为张量积（\hyperlink{para:0:1.2.5}{1.2.5}、\hyperlink{para:0:1.3.1}{1.3.1}），
由张量积运算与归纳极限可交换，得到典范同构
\[
 S^{-1}\bigl(\ilim M_\alpha\bigr)\xrightarrow{\sim}\ilim S^{-1}M_\alpha.
\]
也就是说，关于 $M$ 的函子 $S^{-1}M$ 与归纳极限可交换。

\egapar{1.3.4}
设 $M,N$ 为两个 $A$-模。存在一个关于 $M,N$ 的典范函子性同构
\[
 (S^{-1}M)\otimes_{S^{-1}A}(S^{-1}N)
 \xrightarrow{\sim}S^{-1}(M\otimes_A N),
\]
它将 $(m/s)\otimes(n/t)$ 映为 $(m\otimes n)/(st)$。

\egapar{1.3.5}
同样，有一个关于 $M,N$ 的函子性同态
\[
 S^{-1}\Hom_A(M,N)\longrightarrow
 \Hom_{S^{-1}A}(S^{-1}M,S^{-1}N),
\]
将 $u/s$ 映为同态 $m/t\mapsto u(m)/(st)$。
当 $M$ 具有有限表示时，这个同态是同构。
$M$ 具有 $A^r$ 的形式时，这一点是显然的；
从正合列 $A^p\to A^q\to M\to0$ 出发，利用函子 $S^{-1}M$ 的正合性以及
$\Hom_A(M,N)$ 关于 $M$ 的左正合性，就过渡到一般情形。
注意，当 $A$ 是诺特环而 $M$ 为有限型 $A$-模时，总会出现这种情形。

\subsection{乘法子集的变换}\hypertarget{heading:0:1.4}{}
\src{15—17}

\egapar{1.4.1}
设 $S,T$ 为环 $A$ 的两个乘法子集，且 $S\subseteq T$。
存在一个典范同态 $\rho_A^{T,S}$（或简记为 $\rho^{T,S}$）从 $S^{-1}A$ 到 $T^{-1}A$，
将 $S^{-1}A$ 中记为 $a/s$ 的元素映为 $T^{-1}A$ 中记为 $a/s$ 的元素；有
$i_A^T=\rho_A^{T,S}\circ i_A^S$。
对任一 $A$-模 $M$，同样存在一个 $S^{-1}A$-线性映射 $S^{-1}M\to T^{-1}M$，
将 $S^{-1}M$ 中的 $m/s$ 映为 $T^{-1}M$ 中的 $m/s$；
后者借助 $\rho_A^{T,S}$ 看作 $S^{-1}A$-模。
这个映射记为 $\rho_M^{T,S}$，或简记为 $\rho^{T,S}$，并且仍有
$i_M^T=\rho_M^{T,S}\circ i_M^S$。
在（\hyperlink{para:0:1.2.5}{1.2.5}）的典范等同下，$\rho_M^{T,S}$ 等同于 $\rho_A^{T,S}\otimes1$。
同态 $\rho^{T,S}$ 是从函子 $S^{-1}M$ 到函子 $T^{-1}M$ 的\emph{函子态射}，或\emph{自然变换}；
就是说，对任一同态 $u\colon M\to N$，交换图
\[
\begin{tikzcd}[column sep=large,row sep=large]
S^{-1}M \arrow[r,"S^{-1}u"] \arrow[d,"\rho_M^{T,S}"'] & S^{-1}N \arrow[d,"\rho_N^{T,S}"]\\
T^{-1}M \arrow[r,"T^{-1}u"'] & T^{-1}N
\end{tikzcd}
\]
成立。
还须注意，$T^{-1}u$ 完全由 $S^{-1}u$ 决定，因为对 $m\in M,t\in T$，有
\[
 (T^{-1}u)(m/t)=(t/1)^{-1}\rho^{T,S}\bigl((S^{-1}u)(m/1)\bigr).
\]

\egapar{1.4.2}
沿用上述记号。对两个 $A$-模 $M,N$，下列两个图交换（参阅（\hyperlink{para:0:1.3.4}{1.3.4}）、（\hyperlink{para:0:1.3.5}{1.3.5}））：
\[
\begin{tikzcd}[column sep=large,row sep=large]
(S^{-1}M)\otimes_{S^{-1}A}(S^{-1}N) \arrow[r,"\sim"] \arrow[d] & S^{-1}(M\otimes_A N) \arrow[d]\\
(T^{-1}M)\otimes_{T^{-1}A}(T^{-1}N) \arrow[r,"\sim"'] & T^{-1}(M\otimes_A N)
\end{tikzcd}
\]
\[
\begin{tikzcd}[column sep=large,row sep=large]
S^{-1}\Hom_A(M,N) \arrow[r] \arrow[d] & \Hom_{S^{-1}A}(S^{-1}M,S^{-1}N) \arrow[d]\\
T^{-1}\Hom_A(M,N) \arrow[r] & \Hom_{T^{-1}A}(T^{-1}M,T^{-1}N).
\end{tikzcd}
\]

\egapar{1.4.3}
在一个重要情形中，同态 $\rho^{T,S}$ 是双射：即 $T$ 中每个元素都是 $S$ 中某个元素的因子。
这时通过 $\rho^{T,S}$ 将模 $S^{-1}M$ 与 $T^{-1}M$ 等同。
若 $S$ 中元素在 $A$ 中的任一因子都属于 $S$，则称 $S$ 是\emph{饱和的}。
将 $S$ 换成其所有元素的所有因子组成的集合 $T$，这个集合既是乘法子集又是饱和的，
便知总可以只考虑 $S$ 饱和时的分式模 $S^{-1}M$。

\egapar{1.4.4}
若 $S,T,U$ 为 $A$ 的三个乘法子集，且 $S\subseteq T\subseteq U$，则
\[
 \rho^{U,S}=\rho^{U,T}\circ\rho^{T,S}.
\]

\egapar{1.4.5}
考虑 $A$ 的一族滤过递增的乘法子集 $(S_\alpha)$；以 $\alpha\le\beta$ 表示
$S_\alpha\subseteq S_\beta$。令 $S=\bigcup_\alpha S_\alpha$，这是一个乘法子集。
当 $\alpha\le\beta$ 时，置 $\rho_{\beta\alpha}=\rho^{S_\beta,S_\alpha}$。
由（\hyperlink{para:0:1.4.4}{1.4.4}），同态 $\rho_{\beta\alpha}$ 定义了归纳系
$(S_\alpha^{-1}A,\rho_{\beta\alpha})$ 的归纳极限环 $A'$。
以 $\rho_\alpha$ 表示典范映射 $S_\alpha^{-1}A\to A'$，并置
$\varphi_\alpha=\rho^{S,S_\alpha}$。
由（\hyperlink{para:0:1.4.4}{1.4.4}），对 $\alpha\le\beta$ 有
$\varphi_\alpha=\varphi_\beta\circ\rho_{\beta\alpha}$，故可以唯一定义同态
$\varphi\colon A'\to S^{-1}A$，使下图交换：
\[
\begin{tikzcd}[column sep=large,row sep=large]
& S_\alpha^{-1}A \arrow[d,"\rho_{\beta\alpha}"] \arrow[ddl,"\rho_\alpha"',bend right=12] \arrow[ddr,"\varphi_\alpha",bend left=12] &\\
& S_\beta^{-1}A \arrow[dl,"\rho_\beta"] \arrow[dr,"\varphi_\beta"'] &\\
A' \arrow[rr,"\varphi"'] && S^{-1}A
\end{tikzcd}
\qquad(\alpha\le\beta).
\]
事实上，$\varphi$ 是同构。由构造立即可知它是满射。
另一方面，若 $\rho_\alpha(a/s_\alpha)\in A'$ 满足
$\varphi(\rho_\alpha(a/s_\alpha))=0$，则 $a/s_\alpha=0$ 在 $S^{-1}A$ 中成立，
即存在 $s\in S$ 使 $sa=0$。
但存在 $\beta\ge\alpha$ 使 $s\in S_\beta$，于是
\[
 \rho_\alpha(a/s_\alpha)=\rho_\beta(a/s_\alpha)=0,
\]
故 $\varphi$ 也是单射。
$A$-模 $M$ 的情形同样处理，由此得到典范同构
\[
 \begin{aligned}
 \ilim S_\alpha^{-1}A&\xrightarrow{\sim}(\ilim S_\alpha)^{-1}A,\\
 \ilim S_\alpha^{-1}M&\xrightarrow{\sim}(\ilim S_\alpha)^{-1}M,
 \end{aligned}
\]
第二个同构关于 $M$ 是函子性的。

\egapar{1.4.6}
设 $S_1,S_2$ 为 $A$ 的两个乘法子集，则 $S_1S_2$ 也是 $A$ 的乘法子集。
以 $S'_2$ 表示 $S_2$ 在环 $S_1^{-1}A$ 中的典范像，它是这个环的乘法子集。
对任一 $A$-模 $M$，存在函子性同构
\[
 (S'_2)^{-1}(S_1^{-1}M)\xrightarrow{\sim}(S_1S_2)^{-1}M,
\]
将 $(m/s_1)/(s_2/1)$ 映为 $m/(s_1s_2)$。

\subsection{环的变换}\hypertarget{heading:0:1.5}{}
\src{17—19}

\egapar{1.5.1}
设 $A,A'$ 为两个环，$\varphi\colon A'\to A$ 为同态，
$S,S'$ 分别为 $A,A'$ 的乘法子集，且 $\varphi(S')\subseteq S$。
由（\hyperlink{para:0:1.2.4}{1.2.4}），复合同态 $A'\xrightarrow{\varphi}A\to S^{-1}A$ 分解为
\[
 A'\longrightarrow S'^{-1}A'\xrightarrow{\varphi^{S'}}S^{-1}A;
\]
有 $\varphi^{S'}(a'/s')=\varphi(a')/\varphi(s')$。
若 $A=\varphi(A')$ 且 $S=\varphi(S')$，则 $\varphi^{S'}$ 是满射。
若 $A'=A$ 且 $\varphi$ 为恒等映射，则 $\varphi^{S'}$ 就是（\hyperlink{para:0:1.4.1}{1.4.1}）中定义的同态 $\rho^{S,S'}$。

\egapar{1.5.2}
在（\hyperlink{para:0:1.5.1}{1.5.1}）的假设下，设 $M$ 为 $A$-模。
存在 $S'^{-1}A'$-模的典范函子性同态
\[
 \sigma\colon S'^{-1}(M_{[\varphi]})\longrightarrow
 (S^{-1}M)_{[\varphi^{S'}]},
\]
它将 $S'^{-1}(M_{[\varphi]})$ 的每个元素 $m/s'$ 映为
$(S^{-1}M)_{[\varphi^{S'}]}$ 中的 $m/\varphi(s')$。
立即可以验证，这一定义与所考虑元素的表达式 $m/s'$ 无关。
当 $S=\varphi(S')$ 时，$\sigma$ 是双射。
当 $A'=A$ 且 $\varphi$ 为恒等映射时，$\sigma$ 就是（\hyperlink{para:0:1.4.1}{1.4.1}）中定义的同态 $\rho^{S,S'}$。

特别取 $M=A$ 时，同态 $\varphi$ 在 $A$ 上定义一个 $A'$-代数结构。
这时 $S'^{-1}(A_{[\varphi]})$ 带有环结构，作为环，它等同于 $(\varphi(S'))^{-1}A$；
同态 $\sigma\colon S'^{-1}(A_{[\varphi]})\to S^{-1}A$ 是 $S'^{-1}A'$-代数同态。

\egapar{1.5.3}
设 $M,N$ 为两个 $A$-模。
复合（\hyperlink{para:0:1.3.4}{1.3.4}）与（\hyperlink{para:0:1.5.2}{1.5.2}）中定义的同态，得到同态
\[
 \bigl((S^{-1}M)\otimes_{S^{-1}A}(S^{-1}N)\bigr)_{[\varphi^{S'}]}
 \longleftarrow S'^{-1}\bigl((M\otimes_A N)_{[\varphi]}\bigr),
\]
当 $\varphi(S')=S$ 时它是同构。
同样，复合（\hyperlink{para:0:1.3.5}{1.3.5}）与（\hyperlink{para:0:1.5.2}{1.5.2}）中的同态，得到
\[
 \begin{split}
 S'^{-1}\bigl((\Hom_A(M,N))_{[\varphi]}\bigr)
 \longrightarrow{}&\bigl(\Hom_{S^{-1}A}(S^{-1}M,S^{-1}N)\bigr)_{[\varphi^{S'}]},
 \end{split}
\]
当 $\varphi(S')=S$ 且 $M$ 具有有限表示时，它是同构。

\egapar{1.5.4}
现在考虑一个 $A'$-模 $N'$，并构造张量积 $N'\otimes_{A'}A_{[\varphi]}$。
通过规定 $a\cdot(n'\otimes b)=n'\otimes(ab)$，可将它看作 $A$-模。
存在 $S^{-1}A$-模的函子性同构
\[
 \tau\colon (S'^{-1}N')\otimes_{S'^{-1}A'}(S^{-1}A)_{[\varphi^{S'}]}
 \xrightarrow{\sim} S^{-1}(N'\otimes_{A'}A_{[\varphi]}),
\]
将 $(n'/s')\otimes(a/s)$ 映为 $(n'\otimes a)/(\varphi(s')s)$。
事实上，分别将 $n'/s'$ 或 $a/s$ 换成同一元素的另一表达式时，
可以验证 $(n'\otimes a)/(\varphi(s')s)$ 不变。
另一方面，将 $(n'\otimes a)/s$ 对应于 $(n'/1)\otimes(a/s)$，就可以定义 $\tau$ 的逆同态。
这里使用了如下事实：
$S^{-1}(N'\otimes_{A'}A_{[\varphi]})$ 典范同构于
$(N'\otimes_{A'}A_{[\varphi]})\otimes_A S^{-1}A$（\hyperlink{para:0:1.2.5}{1.2.5}），
因而也典范同构于 $N'\otimes_{A'}(S^{-1}A)_{[\psi]}$，
其中 $\psi$ 表示从 $A'$ 到 $S^{-1}A$ 的复合同态 $a'\mapsto\varphi(a')/1$。

\egapar{1.5.5}
若 $M',N'$ 为两个 $A'$-模，则复合（\hyperlink{para:0:1.3.4}{1.3.4}）与（\hyperlink{para:0:1.5.4}{1.5.4}）的同构，得到同构
\[
 \begin{split}
 S'^{-1}M'\otimes_{S'^{-1}A'}S'^{-1}N'\otimes_{S'^{-1}A'}S^{-1}A
 \xrightarrow{\sim}{}&
 S^{-1}(M'\otimes_{A'}N'\otimes_{A'}A).
 \end{split}
\]
同样，若 $M'$ 具有有限表示，则由（\hyperlink{para:0:1.3.5}{1.3.5}）及（\hyperlink{para:0:1.5.4}{1.5.4}）得到同构
\[
 \begin{split}
 \Hom_{S'^{-1}A'}(S'^{-1}M',S'^{-1}N')\otimes_{S'^{-1}A'}S^{-1}A
 \xrightarrow{\sim}{}& S^{-1}\bigl(\Hom_{A'}(M',N')\otimes_{A'}A\bigr).
 \end{split}
\]

\egapar{1.5.6}
在（\hyperlink{para:0:1.5.1}{1.5.1}）的假设下，设 $T,T'$ 分别为 $A,A'$ 的另一个乘法子集，
且 $S\subseteq T$、$S'\subseteq T'$、$\varphi(T')\subseteq T$。
于是下图交换：
\[
\begin{tikzcd}[column sep=large,row sep=large]
 S'^{-1}A' \arrow[r,"\varphi^{S'}"] \arrow[d,"\rho^{T',S'}"'] & S^{-1}A \arrow[d,"\rho^{T,S}"]\\
 T'^{-1}A' \arrow[r,"\varphi^{T'}"'] & T^{-1}A.
\end{tikzcd}
\]
若 $M$ 为 $A$-模，则下图交换：
\[
\begin{tikzcd}[column sep=large,row sep=large]
 S'^{-1}(M_{[\varphi]}) \arrow[r,"\sigma"] \arrow[d,"\rho^{T',S'}"'] & (S^{-1}M)_{[\varphi^{S'}]} \arrow[d,"\rho^{T,S}"]\\
 T'^{-1}(M_{[\varphi]}) \arrow[r,"\sigma"'] & (T^{-1}M)_{[\varphi^{T'}]}.
\end{tikzcd}
\]
最后，若 $N'$ 为 $A'$-模，则下图交换：
\[
\begin{tikzcd}[column sep=large,row sep=large]
 (S'^{-1}N')\otimes_{S'^{-1}A'}(S^{-1}A)_{[\varphi^{S'}]}
 \arrow[r,"\tau"] \arrow[d]
 & S^{-1}(N'\otimes_{A'}A_{[\varphi]}) \arrow[d,"\rho^{T,S}"]\\
 (T'^{-1}N')\otimes_{T'^{-1}A'}(T^{-1}A)_{[\varphi^{T'}]}
 \arrow[r,"\tau"'] & T^{-1}(N'\otimes_{A'}A_{[\varphi]}).
\end{tikzcd}
\]
其中左侧竖箭头由在 $S'^{-1}N'$ 上应用 $\rho_{N'}^{T',S'}$，
并在 $S^{-1}A$ 上应用 $\rho_A^{T,S}$ 得到。

\egapar{1.5.7}
设 $A''$ 为第三个环，$\varphi'\colon A''\to A'$ 为环同态，
$S''$ 为 $A''$ 的乘法子集，且 $\varphi'(S'')\subseteq S'$。
置 $\varphi''=\varphi\circ\varphi'$，则有
\[
 (\varphi'')^{S''}=\varphi^{S'}\circ(\varphi')^{S''}.
\]
设 $M$ 为 $A$-模，显然有 $M_{[\varphi'']}=(M_{[\varphi]})_{[\varphi']}$。
若 $\sigma',\sigma''$ 分别由 $\varphi',\varphi''$ 定义，其定义方式与（\hyperlink{para:0:1.5.2}{1.5.2}）中从 $\varphi$ 定义 $\sigma$ 相同，
则有传递公式
\[
 \sigma''=\sigma\circ\sigma'.
\]
最后，设 $N''$ 为 $A''$-模。
$A$-模 $N''\otimes_{A''}A_{[\varphi'']}$ 典范地等同于
$(N''\otimes_{A''}A'_{[\varphi']})\otimes_{A'}A_{[\varphi]}$；
同样，$S^{-1}A$-模
\[
 (S''^{-1}N'')\otimes_{S''^{-1}A''}(S^{-1}A)_{[(\varphi'')^{S''}]}
\]
典范地等同于
\[
 \bigl((S''^{-1}N'')\otimes_{S''^{-1}A''}
       (S'^{-1}A')_{[(\varphi')^{S''}]}\bigr)
       \otimes_{S'^{-1}A'}(S^{-1}A)_{[\varphi^{S'}]}.
\]
在这些等同之下，若 $\tau',\tau''$ 分别由 $\varphi',\varphi''$ 定义，其方式与（\hyperlink{para:0:1.5.4}{1.5.4}）中
从 $\varphi$ 定义 $\tau$ 相同，则有传递公式
\[
 \tau''=\tau\circ(\tau'\otimes1).
\]

\egapar{1.5.8}
设 $A$ 为环 $B$ 的子环。
对 $A$ 的每个极小素理想 $\mathfrak p$，都存在 $B$ 的一个极小素理想 $\mathfrak q$，
使得 $\mathfrak p=A\cap\mathfrak q$。
事实上，$A_{\mathfrak p}$ 是 $B_{\mathfrak p}$ 的子环（\hyperlink{para:0:1.3.2}{1.3.2}），并且只有一个素理想
$\mathfrak p'$（\hyperlink{para:0:1.2.6}{1.2.6}）。
由于 $B_{\mathfrak p}$ 非零，它至少有一个素理想 $\mathfrak q'$，
且必有 $\mathfrak q'\cap A_{\mathfrak p}=\mathfrak p'$。
于是 $\mathfrak q'$ 在 $B$ 中的逆像素理想 $\mathfrak q_1$ 满足
$\mathfrak q_1\cap A=\mathfrak p$；
故对每个包含在 $\mathfrak q_1$ 中的 $B$ 的极小素理想 $\mathfrak q$，更有
$\mathfrak q\cap A=\mathfrak p$。

\subsection{将模 \texorpdfstring{$M_f$}{Mf} 等同于一个归纳极限}\hypertarget{heading:0:1.6}{}
\src{19—20}

\egapar{1.6.1}
设 $M$ 为 $A$-模，$f$ 为 $A$ 的元素。
考虑一个 $A$-模序列 $(M_n)$，其中每个模都等于 $M$；
对每对整数 $m\le n$，令 $\varphi_{nm}\colon M_m\to M_n$ 为同态
$z\mapsto f^{n-m}z$。
立即可见，$((M_n),(\varphi_{nm}))$ 是一个 $A$-模归纳系。
以 $N=\ilim M_n$ 表示其归纳极限。
我们将定义一个从 $N$ 到 $M_f$ 的典范函子性 $A$-同构。

为此，注意，对每个 $n$，$\theta_n\colon z\mapsto z/f^n$ 是从 $M=M_n$ 到 $M_f$ 的 $A$-同态，
并且由定义，当 $m\le n$ 时有 $\theta_n\circ\varphi_{nm}=\theta_m$。
因此存在 $A$-同态 $\theta\colon N\to M_f$，使得对每个 $n$ 有
$\theta_n=\theta\circ\varphi_n$，其中 $\varphi_n$ 为典范同态 $M_n\to N$。
按假设，$M_f$ 的每个元素都至少可以写成某个 $n$ 对应的 $z/f^n$，故 $\theta$ 显然是满射。
另一方面，若 $\theta(\varphi_n(z))=0$，即 $z/f^n=0$，
则存在整数 $k>0$ 使 $f^kz=0$，从而 $\varphi_{n+k,n}(z)=0$，故 $\varphi_n(z)=0$。
因此可以通过 $\theta$ 将 $M_f$ 与 $\ilim M_n$ 等同。

\egapar{1.6.2}
现在将 $M_n,\varphi_{nm},\varphi_n$ 分别记为 $M_{f,n},\varphi^f_{nm},\varphi^f_n$。
设 $g$ 为 $A$ 的另一个元素。
由于 $f^n$ 整除 $f^ng^n$，有函子性同态
\[
 \rho_{fg,f}\colon M_f\longrightarrow M_{fg}
 \qquad\text{（\hyperlink{para:0:1.4.1}{1.4.1}、\hyperlink{para:0:1.4.3}{1.4.3}）}.
\]
若分别将 $M_f,M_{fg}$ 等同于 $\ilim M_{f,n},\ilim M_{fg,n}$，
则 $\rho_{fg,f}$ 等同于由
\[
 \rho^n_{fg,f}\colon M_{f,n}\longrightarrow M_{fg,n},
 \qquad \rho^n_{fg,f}(z)=g^nz
\]
定义的各映射的归纳极限。
事实上，这立即由下图的交换性得到：
\[
\begin{tikzcd}[column sep=large,row sep=large]
 M_{f,n} \arrow[r,"\rho^n_{fg,f}"] \arrow[d,"\varphi^f_n"']
 & M_{fg,n} \arrow[d,"\varphi^{fg}_n"]\\
 M_f \arrow[r,"\rho_{fg,f}"'] & M_{fg}.
\end{tikzcd}
\]

\subsection{模的支集}\hypertarget{heading:0:1.7}{}
\src{20}

\egapar{1.7.1}
给定 $A$-模 $M$，称 $A$ 中所有满足 $M_{\mathfrak p}\ne0$ 的素理想 $\mathfrak p$ 组成的集合为
$M$ 的\emph{支集}，记为 $\Supp(M)$。
$M=0$ 当且仅当 $\Supp(M)=\varnothing$。
这是因为，若对所有 $\mathfrak p$ 都有 $M_{\mathfrak p}=0$，
则每个元素 $x\in M$ 的零化理想都不能包含在 $A$ 的任何素理想中，因而等于整个 $A$。

\egapar{1.7.2}
若 $0\to N\to M\to P\to0$ 为 $A$-模的正合列，则
\[
 \Supp(M)=\Supp(N)\cup\Supp(P).
\]
因为对 $A$ 的每个素理想 $\mathfrak p$，列
$0\to N_{\mathfrak p}\to M_{\mathfrak p}\to P_{\mathfrak p}\to0$ 正合（\hyperlink{para:0:1.3.2}{1.3.2}），
而 $M_{\mathfrak p}=0$ 当且仅当 $N_{\mathfrak p}=P_{\mathfrak p}=0$。

\egapar{1.7.3}
若 $M$ 是一族子模 $(M_\lambda)$ 的和，则对每个素理想 $\mathfrak p$，
$M_{\mathfrak p}$ 是各 $(M_\lambda)_{\mathfrak p}$ 的和（\hyperlink{para:0:1.3.3}{1.3.3}、\hyperlink{para:0:1.3.2}{1.3.2}），故
\[
 \Supp(M)=\bigcup_\lambda\Supp(M_\lambda).
\]

\egapar{1.7.4}
若 $M$ 是有限型 $A$-模，则 $\Supp(M)$ 是包含 $M$ 的零化理想的所有素理想组成的集合。
事实上，若 $M$ 由单个元素 $x$ 生成，则 $M_{\mathfrak p}=0$ 的意思是存在
$s\notin\mathfrak p$ 使 $s\cdot x=0$，也就是 $\mathfrak p$ 不包含 $x$ 的零化理想。
现在若 $M$ 有有限生成元系 $(x_i)_{1\le i\le n}$，并以 $\mathfrak a_i$ 表示 $x_i$ 的零化理想，
则由（\hyperlink{para:0:1.7.3}{1.7.3}），$\Supp(M)$ 是包含某个 $\mathfrak a_i$ 的素理想的集合，
等价地，也是包含 $\mathfrak a=\bigcap_i\mathfrak a_i$ 的素理想的集合；
而 $\mathfrak a$ 恰为 $M$ 的零化理想。

\egapar{1.7.5}
若 $M,N$ 为两个有限型 $A$-模，则
\[
 \Supp(M\otimes_A N)=\Supp(M)\cap\Supp(N).
\]
结合（\hyperlink{para:0:1.3.4}{1.3.4}），所需证明的是：对 $A$ 的素理想 $\mathfrak p$，
条件 $M_{\mathfrak p}\otimes_{A_{\mathfrak p}}N_{\mathfrak p}\ne0$，
等价于“$M_{\mathfrak p}\ne0$ 且 $N_{\mathfrak p}\ne0$”。
换言之，只需证明：若 $P,Q$ 为局部环 $B$ 上的两个非零有限型模，则 $P\otimes_B Q\ne0$。
设 $\mathfrak m$ 为 $B$ 的极大理想。
由 Nakayama 引理，向量空间 $P/\mathfrak mP$ 和 $Q/\mathfrak mQ$ 均非零，故它们的张量积
\[
 (P/\mathfrak mP)\otimes_{B/\mathfrak m}(Q/\mathfrak mQ)
 =(P\otimes_B Q)\otimes_B(B/\mathfrak m)
\]
也非零，由此得到结论。

特别地，若 $M$ 是有限型 $A$-模，$\mathfrak a$ 是 $A$ 的理想，则
$\Supp(M/\mathfrak aM)$ 是同时包含 $\mathfrak a$ 与 $M$ 的零化理想 $\mathfrak n$ 的素理想的集合（\hyperlink{para:0:1.7.4}{1.7.4}），
也就是包含 $\mathfrak a+\mathfrak n$ 的素理想的集合。


% END chapters/ch0_s1.tex

\par

% BEGIN chapters/ch0_s2.tex
\section{不可约空间；诺特空间}\hypertarget{heading:0:2}{}
\subsection{不可约空间}\hypertarget{heading:0:2.1}{}
\src{21—22}

\egapar{2.1.1}
若拓扑空间 $X$ 非空，且不能表示成两个均不等于 $X$ 的闭子空间的并，则称 $X$ 是\emph{不可约的}。
等价地说，$X\ne\varnothing$，且 $X$ 中任意两个非空开集的交非空，因而任意有限个非空开集的交也非空。
在 $X$ 非空的条件下，这也等价于以下任一说法：每个非空开集都处处稠密；
每个不等于 $X$ 的闭子集都无处稠密；或者，$X$ 的每个开集都是连通的。

\egapar{2.1.2}
拓扑空间 $X$ 的子空间 $Y$ 不可约，当且仅当其闭包 $\overline Y$ 不可约。
特别地，单点子空间的闭包 $\overline{\{x\}}$ 总是不可约的。
对于关系 $y\in\overline{\{x\}}$，等价地，
$\overline{\{y\}}\subseteq\overline{\{x\}}$，
我们称 $y$ 是 $x$ 的一个\emph{特殊化}，或称 $x$ 是 $y$ 的一个\emph{一般化}。
若不可约空间 $X$ 中存在一点 $x$，使 $X=\overline{\{x\}}$，则称 $x$ 是 $X$ 的\emph{泛点}。
这时 $X$ 的每个非空开集都包含 $x$，而每个包含 $x$ 的子空间都以 $x$ 为泛点。

\egapar{2.1.3}
回忆一下，\emph{Kolmogoroff 空间}是指满足下列分离公理的拓扑空间 $X$：
\begin{quote}
$(\mathrm T_0)$ 若 $x\ne y$ 为 $X$ 中的任意两点，则存在一个开集包含其中一点而不包含另一点。
\end{quote}
若不可约的 Kolmogoroff 空间具有泛点，则泛点唯一，因为每个非空开集都包含所有泛点。

又回忆，若拓扑空间 $X$ 的任一开覆盖都可以抽取有限子覆盖，则称 $X$ 是\emph{拟紧的}；
等价地，任一滤过递减的非空闭集族都具有非空交。
若 $X$ 拟紧，则 $X$ 的每个非空闭子集 $A$ 都包含一个极小非空闭集 $M$，
因为 $A$ 的非空闭子集所成的集合对于关系 $\supseteq$ 是归纳有序的。
若 $X$ 还是 Kolmogoroff 空间，则 $M$ 必由一个点组成；略不严格地说，$M$ 是一个\emph{闭点}。

\egapar{2.1.4}
在不可约空间 $X$ 中，每个非空开子空间 $U$ 都不可约；若 $X$ 有泛点 $x$，则 $x$ 也是 $U$ 的泛点。

设 $(U_\alpha)$ 为拓扑空间 $X$ 的一个覆盖，其指标集非空，且每个 $U_\alpha$ 都是非空开集。
$X$ 不可约，当且仅当每个 $U_\alpha$ 都不可约，并且对任意 $\alpha,\beta$ 有
$U_\alpha\cap U_\beta\ne\varnothing$。
这一条件显然必要。
为证明充分性，只须证明：若 $V$ 是 $X$ 的非空开集，则对每个 $\alpha$ 都有
$V\cap U_\alpha\ne\varnothing$。
因为这样 $V\cap U_\alpha$ 就在每个 $U_\alpha$ 中稠密，从而 $V$ 在 $X$ 中稠密。
现在至少存在一个指标 $\gamma$ 使 $V\cap U_\gamma\ne\varnothing$，
于是 $V\cap U_\gamma$ 在 $U_\gamma$ 中稠密。
又由于对每个 $\alpha$ 有 $U_\alpha\cap U_\gamma\ne\varnothing$，
故 $V\cap U_\alpha\cap U_\gamma\ne\varnothing$。

\egapar{2.1.5}
设 $X$ 为不可约空间，$f$ 为从 $X$ 到拓扑空间 $Y$ 的连续映射。
则 $f(X)$ 不可约；若 $x$ 是 $X$ 的泛点，则 $f(x)$ 是 $f(X)$ 的泛点，
因而也是 $\overline{f(X)}$ 的泛点。
特别地，若 $Y$ 也不可约且只有一个泛点 $y$，则在 $x$ 为 $X$ 的泛点的上述情形下，
$f(X)$ 处处稠密，当且仅当 $f(x)=y$。

\egapar{2.1.6}
拓扑空间 $X$ 的每个不可约子空间都包含在一个极大不可约子空间中，后者必为闭集。
$X$ 的极大不可约子空间称为 $X$ 的\emph{不可约分支}。
若 $Z_1,Z_2$ 为 $X$ 的两个不同不可约分支，则 $Z_1\cap Z_2$ 是在各子空间 $Z_1,Z_2$ 中
均无处稠密的闭集。
特别地，若 $X$ 的一个不可约分支具有泛点（\hyperlink{para:0:2.1.2}{2.1.2}），则此泛点不能属于任何其他不可约分支。
若 $X$ 只有有限个不可约分支 $Z_i$（$1\le i\le n$），并令
\[
 U_i=Z_i\cap\left(X\setminus\bigcup_{j\ne i}Z_j\right),
\]
则各 $U_i$ 是开集、不可约、两两不交，且它们的并在 $X$ 中稠密。

设 $U$ 为拓扑空间 $X$ 的开子集。
若 $Z$ 为 $X$ 中与 $U$ 相交的不可约子集，则 $Z\cap U$ 在 $Z$ 中开且稠密，因而不可约。
反过来，对 $U$ 的任一不可约闭子集 $Y$，它在 $X$ 中的闭包 $\overline Y$ 不可约，
并且 $\overline Y\cap U=Y$。
因此，$U$ 的不可约分支与 $X$ 中那些与 $U$ 相交的不可约分支之间存在一一对应。

\egapar{2.1.7}
若拓扑空间 $X$ 是有限个不可约闭子空间 $Y_i$ 的并，则 $X$ 的不可约分支就是集合
$\{Y_i\}$ 中的极大元素。
因为对 $X$ 的任一不可约闭子集 $Z$，$Z$ 是各 $Z\cap Y_i$ 的并，
由此立刻推出 $Z$ 必包含在某个 $Y_i$ 中。
设 $Y$ 为拓扑空间 $X$ 的子空间，并假定 $Y$ 只有有限个不可约分支
$Y_i$（$1\le i\le n$），则它们在 $X$ 中的闭包 $\overline{Y_i}$ 是 $\overline Y$ 的不可约分支。

\egapar{2.1.8}
设 $Y$ 为不可约空间，且具有唯一泛点 $y$。
设 $X$ 为拓扑空间，$f\colon X\to Y$ 为连续映射。
则对 $X$ 中每个与 $f^{-1}(y)$ 相交的不可约分支 $Z$，$f(Z)$ 都在 $Y$ 中稠密。
反命题未必成立；但是，若 $Z$ 有泛点 $z$，且 $f(Z)$ 在 $Y$ 中稠密，则必有 $f(z)=y$（\hyperlink{para:0:2.1.5}{2.1.5}）。
此外，这时 $Z\cap f^{-1}(y)$ 是 $\{z\}$ 在 $f^{-1}(y)$ 中的闭包，因而不可约；
而 $f^{-1}(y)$ 中每个包含 $z$ 的不可约子集都必包含在 $Z$ 中（\hyperlink{para:0:2.1.6}{2.1.6}），
故 $z$ 是 $Z\cap f^{-1}(y)$ 的泛点。
由于 $f^{-1}(y)$ 的每个不可约分支都包含在 $X$ 的某个不可约分支中，得到：
若 $X$ 中每个与 $f^{-1}(y)$ 相交的不可约分支 $Z$ 都有泛点，
则这些分支的集合与 $f^{-1}(y)$ 的不可约分支 $Z\cap f^{-1}(y)$ 的集合之间存在一一对应，
并且 $Z$ 的泛点与 $Z\cap f^{-1}(y)$ 的泛点相同。

% Begin the second subsection with its opening definition.
\clearpage
\subsection{诺特空间}\hypertarget{heading:0:2.2}{}
\src{23}

\egapar{2.2.1}
若拓扑空间 $X$ 的开集集合满足极大条件，等价地，其闭集集合满足极小条件，
则称 $X$ 是\emph{诺特的}。
若每个 $x\in X$ 都有一个作为子空间是诺特空间的邻域，则称 $X$ 是\emph{局部诺特的}。

\egapar{2.2.2}
设 $E$ 为满足极小条件的有序集，$P$ 为 $E$ 的元素的一种性质，且满足下述条件：
若 $a\in E$ 使得对每个 $x<a$ 都有 $P(x)$ 成立，则 $P(a)$ 成立。
在这些条件下，$P(x)$ 对每个 $x\in E$ 都成立，这就是\emph{诺特归纳原理}。
事实上，设 $F$ 为所有不满足 $P(x)$ 的元素 $x\in E$ 组成的集合；
若 $F$ 非空，则它有极小元素 $a$。
由于此时对每个 $x<a$ 都有 $P(x)$ 成立，故 $P(a)$ 也成立，矛盾。
我们将特别在 $E$ 是诺特空间的闭子集的某个集合时应用此原理。

\egapar{2.2.3}
诺特空间的每个子空间都是诺特的。
反过来，一个拓扑空间若是有限个诺特子空间的并，则它也是诺特的。

\egapar{2.2.4}
每个诺特空间都拟紧。
反过来，一个拓扑空间若每个开集都拟紧，则它是诺特的。

\egapar{2.2.5}
诺特空间只有有限个不可约分支；这可以用诺特归纳法看出。


% END chapters/ch0_s2.tex

\par

% BEGIN chapters/ch0_s3.tex
\section{层的补充知识}\hypertarget{heading:0:3}{}
\subsection{取值于一个范畴的层}\hypertarget{heading:0:3.1}{}
\src{23—25}

\egapar{3.1.1}
设 $\Kcat$ 为范畴，$(A_\alpha)_{\alpha\in I}$、
$(A_{\alpha\beta})_{(\alpha,\beta)\in I\times I}$ 为 $\Kcat$ 的两族对象，
满足 $A_{\alpha\beta}=A_{\beta\alpha}$；
$(\rho_{\alpha\beta})_{(\alpha,\beta)\in I\times I}$ 为一族态射
$\rho_{\alpha\beta}\colon A_\alpha\to A_{\alpha\beta}$。
设有一个对象 $A$ 及一族态射 $\rho_\alpha\colon A\to A_\alpha$。
若对 $\Kcat$ 的每个对象 $B$，映射
\[
 \Hom(B,A)\longrightarrow\prod_\alpha\Hom(B,A_\alpha),
 \qquad f\longmapsto(\rho_\alpha\circ f)_\alpha
\]
将 $\Hom(B,A)$ 双射到由那些对任意指标对 $(\alpha,\beta)$ 满足
\[
 \rho_{\alpha\beta}\circ f_\alpha
 =\rho_{\beta\alpha}\circ f_\beta
\]
的族 $(f_\alpha)$ 组成的集合，则称二元组 $(A,(\rho_\alpha))$ 是由
$(A_\alpha)$、$(A_{\alpha\beta})$ 及 $(\rho_{\alpha\beta})$ 定义的\emph{泛问题的解}。
立即可见，这样的解若存在，则在相差一个同构的意义下唯一。

\egapar{3.1.2}
我们不再重述拓扑空间 $X$ 上取值于范畴 $\Kcat$ 的预层 $U\mapsto\F(U)$ 的定义
(G, I, 1.9)。若这样的预层满足下列公理，则称其为\emph{取值于 $\Kcat$ 的层}：
\begin{quote}
\textbf{(F)} 对 $X$ 的任一开集 $U$ 及其由包含在 $U$ 中的开集 $U_\alpha$ 组成的任一覆盖，
以 $\rho_\alpha$、$\rho_{\alpha\beta}$ 分别表示限制态射
\[
 \F(U)\longrightarrow\F(U_\alpha),
 \qquad
 \F(U_\alpha)\longrightarrow\F(U_\alpha\cap U_\beta).
\]
则由 $\F(U)$ 及族 $(\rho_\alpha)$ 组成的二元组，是关于
$(\F(U_\alpha))$、$(\F(U_\alpha\cap U_\beta))$ 及 $(\rho_{\alpha\beta})$
的泛问题的解（\hyperlink{para:0:3.1.1}{3.1.1}）。\origfootnote{
这是一般射影极限概念的一个特例，此处的射影系不一定滤过；参阅 (T, I, 1.8)，
以及导言中提到的正在编写的著作。}
\end{quote}
这等价于说，对 $\Kcat$ 的每个对象 $T$，$U\mapsto\Hom(T,\F(U))$ 是集合层。

\egapar{3.1.3}
设 $\Kcat$ 是由某种“附有态射的结构类型”$\Sigma$ 定义的范畴；
也就是说，$\Kcat$ 的对象是带有 $\Sigma$ 型结构的集合，其态射就是 $\Sigma$ 的态射。
再假定 $\Kcat$ 满足下列条件：
\begin{quote}
\textbf{(E)} 若 $(A,(\rho_\alpha))$ 是范畴 $\Kcat$ 中由族
$(A_\alpha)$、$(A_{\alpha\beta})$、$(\rho_{\alpha\beta})$ 定义的泛映射问题的解，
则它在集合范畴中也是这些族所定义的泛映射问题的解；
也就是说，把 $A$、$A_\alpha$、$A_{\alpha\beta}$ 看作集合，
把 $\rho_\alpha$、$\rho_{\alpha\beta}$ 看作映射时，这一结论仍成立。\origfootnote{
可以证明，这也等价于说典范函子 $\Kcat\to(\mathrm{Ens})$ 与射影极限可交换，
这里的射影系不一定滤过。}
\end{quote}
在这些条件下，(F) 蕴含：$U\mapsto\F(U)$ 作为集合预层是一个层。
此外，由 (F)，映射 $u\colon T\to\F(U)$ 是 $\Kcat$ 的态射，当且仅当每个复合映射
$\rho_\alpha\circ u\colon T\to\F(U_\alpha)$ 都是态射。
这就是说，$\F(U)$ 上的 $\Sigma$ 型结构是相对于各态射 $\rho_\alpha$ 的\emph{初始结构}。
反过来，假定 $X$ 上取值于 $\Kcat$ 的预层 $U\mapsto\F(U)$ 是一个集合层，
并满足上述条件，则它显然满足 (F)，故是取值于 $\Kcat$ 的层。

\egapar{3.1.4}
当 $\Sigma$ 是群结构或环结构这一类型时，取值于 $\Kcat$ 的预层 $U\mapsto\F(U)$
一旦是集合层，就自动成为取值于 $\Kcat$ 的层，也就是 (G) 意义下的群层或环层。\origfootnote{
这是因为，在范畴 $\Kcat$ 中，每个作为集合映射是双射的态射都是同构。
例如，当 $\Kcat$ 是拓扑空间范畴时，这一点不再成立。}
但是，例如当 $\Kcat$ 是拓扑环范畴，而其态射为连续同态时，就不再如此。
这时，一个取值于 $\Kcat$ 的层，就是满足以下条件的环层 $U\mapsto\F(U)$：
对每个开集 $U$ 及其由开集 $U_\alpha\subseteq U$ 组成的任一覆盖，
环 $\F(U)$ 的拓扑是使所有同态 $\F(U)\to\F(U_\alpha)$ 连续的最粗拓扑。
在此情形下，把不带拓扑的环层 $U\mapsto\F(U)$ 称为这一拓扑环层的\emph{底层环层}。

因此，拓扑环层的态射 $u_V\colon\F(V)\to\G(V)$，其中 $V$ 遍历 $X$ 的所有开集，
是其底层环层的同态，并且要求对每个开集 $V\subseteq X$，$u_V$ 都连续。
为区别于底层环层的任意同态，称这些态射为\emph{拓扑环层的连续同态}。
对拓扑空间层及拓扑群层，有类似的定义和约定。

\egapar{3.1.5}
显然，对任意范畴 $\Kcat$，若 $\F$ 为 $X$ 上取值于 $\Kcat$ 的预层（分别地，层），
且 $U$ 为 $X$ 的开集，则开集 $V\subseteq U$ 对应的各 $\F(V)$ 构成取值于 $\Kcat$ 的预层
（分别地，层）。称其为 $\F$ 在 $U$ 上\emph{诱导的预层}（分别地，\emph{诱导的层}），记为 $\F|U$。
对 $X$ 上取值于 $\Kcat$ 的预层态射 $u\colon\F\to\G$，
以 $u|U$ 表示由所有 $V\subseteq U$ 对应的 $u_V$ 组成的态射 $\F|U\to\G|U$。

\egapar{3.1.6}
现在假定范畴 $\Kcat$ 具有归纳极限 (T, 1.8)。
于是，对 $X$ 上取值于 $\Kcat$ 的每个预层，特别是每个层 $\F$，以及每个 $x\in X$，
都可以将\emph{茎} $\F_x$ 定义为 $\Kcat$ 中如下的归纳极限对象：
令 $U$ 遍历 $x$ 在 $X$ 中的所有开邻域，这些邻域关于 $\supseteq$ 构成滤过集，
并采用态射 $\rho_U^V\colon\F(V)\to\F(U)$，取各 $\F(U)$ 的归纳极限。
若 $u\colon\F\to\G$ 为取值于 $\Kcat$ 的预层态射，则对每个 $x\in X$，
将态射 $u_x\colon\F_x\to\G_x$ 定义为各 $u_U\colon\F(U)\to\G(U)$
沿 $x$ 的开邻域集合的归纳极限。
如此，对每个 $x\in X$，$\F_x$ 就成为关于 $\F$ 的、取值于 $\Kcat$ 的协变函子。

当 $\Kcat$ 还由某种附有态射的结构类型 $\Sigma$ 定义时，
仍将 $\F(U)$ 的元素称为层 $\F$ 在 $U$ 上的\emph{截面}，并把 $\F(U)$ 写作 $\Gamma(U,\F)$。
对 $s\in\Gamma(U,\F)$ 及包含在 $U$ 中的开集 $V$，把 $\rho_V^U(s)$ 写作 $s|V$。
对每个 $x\in U$，$s$ 在 $\F_x$ 中的典范像是 $s$ 在点 $x$ 的\emph{芽}，记为 $s_x$。
我们绝不在这个意义下使用记号 $s(x)$；该记号保留给本专著中将要考虑的特殊层的另一个概念（\hyperlink{para:0:5.5.1}{5.5.1}）。
若 $u\colon\F\to\G$ 为取值于 $\Kcat$ 的层态射，则对任意 $s\in\Gamma(U,\F)$，
将 $u_U(s)$ 写作 $u(s)$。

若 $\F$ 是交换群层、环层或模层，则将所有满足 $\F_x\ne\{0\}$ 的点 $x\in X$ 组成的集合
称为 $\F$ 的\emph{支集}，记为 $\Supp(\F)$；这个集合在 $X$ 中未必是闭的。

当 $\Kcat$ 由附有态射的结构类型定义时，对取值于 $\Kcat$ 的层，
我们将一贯避免引入“展布空间”（\textit{espaces étalés}）的观点。
换言之，我们绝不把层看成一个拓扑空间，甚至不把它看成其各茎的并所成的集合；
也不把 $X$ 上这种层的态射 $u\colon\F\to\G$ 看成拓扑空间之间的连续映射。

\subsection{开集基上的预层}\hypertarget{heading:0:3.2}{}
\src{25—28}

\egapar{3.2.1}
以下我们只考虑具有广义射影极限的范畴 $\Kcat$；
所谓广义，是说对应的预序集不一定滤过，参阅 (T, 1.8)。
设 $X$ 为拓扑空间，$\B$ 为 $X$ 的拓扑的一个开集基。
我们把下列数据称为\emph{$\B$ 上取值于 $\Kcat$ 的预层}：
给每个 $U\in\B$ 配备一个对象 $\F(U)\in\Kcat$，
并给每对满足 $U\subseteq V$ 的 $\B$ 中元素配备一个态射
\[
 \rho_U^V\colon\F(V)\longrightarrow\F(U),
\]
使得 $\rho_U^U$ 为恒等态射，且当 $U,V,W\in\B$、$U\subseteq V\subseteq W$ 时，有
$\rho_U^W=\rho_U^V\circ\rho_V^W$。

可以给它配备一个通常意义下取值于 $\Kcat$ 的预层 $U\mapsto\F'(U)$：
对任一开集 $U$，令
\[
 \F'(U)=\plim_{\substack{V\in\B\\V\subseteq U}}\F(V).
\]
这里 $V$ 遍历包含在 $U$ 中的基元素所成的有序集，其次序为 $\subseteq$，一般并不滤过。
当 $V\subseteq W\subseteq U$ 且 $V,W\in\B$ 时，各 $\F(V)$ 关于 $\rho_V^W$ 构成射影系。
事实上，若 $U,U'$ 为 $X$ 的开集且 $U\subseteq U'$，
就把 $\rho'{}^{U'}_U$ 定义为各典范态射 $\F'(U')\to\F(V)$（$V\subseteq U$）的射影极限；
也就是说，它是唯一的态射 $\F'(U')\to\F'(U)$，与各典范态射 $\F'(U)\to\F(V)$ 复合后，
得到相应的典范态射 $\F'(U')\to\F(V)$。
立即可以验证各 $\rho'{}^{U'}_U$ 的传递性。
此外，若 $U\in\B$，则典范态射 $\F'(U)\to\F(U)$ 是同构，从而可以将这两个对象等同。\origfootnote{
若 $X$ 为诺特空间，即使只假定 $\Kcat$ 对有限射影系具有射影极限，也仍可定义 $\F'(U)$，
并证明它是通常意义下的预层。
事实上，若 $U$ 为 $X$ 的任意开集，则存在由 $\B$ 的元素组成的有限覆盖 $(V_i)$；
对每对指标 $(i,j)$，取 $V_i\cap V_j$ 的一个由 $\B$ 的元素组成的有限覆盖 $(V_{ijk})$。
令 $I$ 由指标 $i$ 与三元组 $(i,j,k)$ 组成，并只规定次序关系
$i>(i,j,k)$、$j>(i,j,k)$。
取 $\F'(U)$ 为各 $\F(V_i)$ 与 $\F(V_{ijk})$ 所成系统的射影极限；
容易验证这与覆盖 $(V_i)$ 和 $(V_{ijk})$ 的选择无关，且 $U\mapsto\F'(U)$ 为预层。}

\egapar{3.2.2}
如此定义的预层 $\F'$ 是层，当且仅当 $\B$ 上的预层 $\F$ 满足以下条件：
\begin{quote}
\textbf{($\mathrm F_0$)} 对任一 $U\in\B$，任一由包含在 $U$ 中的基元素 $U_\alpha\in\B$
组成的覆盖 $(U_\alpha)$，以及 $\Kcat$ 的任一对象 $T$，映射
\[
 \Hom(T,\F(U))\longrightarrow\prod_\alpha\Hom(T,\F(U_\alpha)),
 \qquad f\longmapsto(\rho_{U_\alpha}^U\circ f)_\alpha
\]
是从 $\Hom(T,\F(U))$ 到下述族 $(f_\alpha)$ 的集合的双射：
对每对指标 $(\alpha,\beta)$ 和每个满足 $V\in\B$、$V\subseteq U_\alpha\cap U_\beta$ 的 $V$，都有
\[
 \rho_V^{U_\alpha}\circ f_\alpha
 =\rho_V^{U_\beta}\circ f_\beta.
\]
\end{quote}
\origfootnote{
这也就是说，由 $\F(U)$ 及 $\rho_\alpha=\rho_{U_\alpha}^U$ 组成的二元组，是（\hyperlink{para:0:3.1.1}{3.1.1}）的泛问题的解。
这里取 $A_\alpha=\F(U_\alpha)$、$A_{\alpha\beta}=\prod\F(V)$，
其中 $V$ 遍历包含在 $U_\alpha\cap U_\beta$ 中的基元素；
$\rho_{\alpha\beta}=(\rho''_V)\colon\F(U_\alpha)\to\prod\F(V)$，
其组成态射满足：当 $V,V',W\in\B$、$V\cup V'\subseteq U_\alpha\cap U_\beta$、$W\subseteq V\cap V'$ 时，
有 $\rho_W^V\circ\rho''_V=\rho_W^{V'}\circ\rho''_{V'}$。}

该条件显然必要。
为证明充分性，先考虑 $X$ 的拓扑的另一个基 $\B'\subseteq\B$。
设 $\F''$ 为从子族 $(\F(V))_{V\in\B'}$ 导出的预层；我们证明 $\F''$ 与 $\F'$ 典范同构。
首先，对任一开集 $U$，典范态射 $\F'(U)\to\F(V)$ 对 $V\in\B'$、$V\subseteq U$ 的射影极限，
给出态射 $\F'(U)\to\F''(U)$。
若 $U\in\B$，则此态射为同构：因为根据假设，各典范态射 $\F''(U)\to\F(V)$，
其中 $V\in\B'$、$V\subseteq U$，分解为
$\F''(U)\to\F(U)\to\F(V)$；
立即可以看出，如此定义的态射 $\F(U)\to\F''(U)$ 与 $\F''(U)\to\F(U)$ 的两个复合均为恒等态射。
因此，对任一开集 $U$，各态射 $\F''(U)\to\F''(W)=\F(W)$，其中 $W\in\B$、$W\subseteq U$，
满足刻画各 $\F(W)$ 的射影极限的条件。
由射影极限在相差同构意义下的唯一性，就证明了我们的断言。

现在，设 $U$ 为 $X$ 的任一开集，$(U_\alpha)$ 为其开覆盖，
并以 $\B'$ 表示 $\B$ 中那些至少包含在某个 $U_\alpha$ 内的基元素组成的子族。
显然，$\B'$ 仍是 $X$ 的拓扑的一个基。\transnote{
原文此处写作“$X$ 的拓扑的一个基”。当 $U\ne X$ 时，所定义的 $\B'$ 实际上是开子空间 $U$ 的拓扑基；
以下只对 $U$ 及各 $U_\alpha$ 取射影极限，应在这个开子空间内理解此句。这里保留原句并标明其适用范围。}
故 $\F'(U)$ 是各 $\F(V)$ 对 $V\in\B'$、$V\subseteq U$ 的射影极限，
而 $\F'(U_\alpha)$ 是对 $V\in\B'$、$V\subseteq U_\alpha$ 的射影极限。
由射影极限的定义立即可验证公理 (F)。

当 ($\mathrm F_0$) 成立时，我们也略不严格地将基 $\B$ 上的预层 $\F$ 称为一个\emph{层}。

\egapar{3.2.3}
设 $\F,\G$ 为基 $\B$ 上取值于 $\Kcat$ 的两个预层。
态射 $u\colon\F\to\G$ 定义为一族态射
$u_V\colon\F(V)\to\G(V)$（$V\in\B$），满足与限制态射 $\rho_U^V$ 的通常相容条件。
采用（\hyperlink{para:0:3.2.1}{3.2.1}）的记号，对每个开集 $U$，取 $u'_U$ 为各 $u_V$ 对 $V\in\B$、$V\subseteq U$ 的射影极限，
就得到相应通常预层之间的态射 $u'\colon\F'\to\G'$。
与限制态射的相容性由射影极限的函子性得出。

\egapar{3.2.4}
若范畴 $\Kcat$ 具有归纳极限，而 $\F$ 为基 $\B$ 上取值于 $\Kcat$ 的预层，
则对每个 $x\in X$，属于 $\B$ 的 $x$ 的邻域在 $x$ 的邻域集合中，关于 $\supseteq$ 构成共尾子集。
因此，若 $\F'$ 为与 $\F$ 对应的通常预层，则茎 $\F'_x$ 等于
$\ilim_{V\in\B,\,x\in V}\F(V)$。
若 $u\colon\F\to\G$ 为 $\B$ 上取值于 $\Kcat$ 的预层态射，
$u'\colon\F'\to\G'$ 为相应通常预层之间的态射，
则 $u'_x$ 同样是各态射 $u_V\colon\F(V)\to\G(V)$ 对 $V\in\B$、$x\in V$ 的归纳极限。

\egapar{3.2.5}
回到（\hyperlink{para:0:3.2.1}{3.2.1}）的一般条件。
若 $\F$ 为取值于 $\Kcat$ 的通常层，$\F_0$ 为将 $\F$ 限制到 $\B$ 所得的基上的层，
则按（\hyperlink{para:0:3.2.1}{3.2.1}）的方法由 $\F_0$ 构造的通常层 $\F_1$，由条件 (F) 及射影极限的唯一性，
与 $\F$ 典范同构。
通常将 $\F$ 与 $\F_1$ 等同。

若 $\G$ 为 $X$ 上取值于 $\Kcat$ 的另一个通常层，$u\colon\F\to\G$ 为态射，
则上述说明表明：只给出 $V\in\B$ 时的 $u_V\colon\F(V)\to\G(V)$，就完全决定了 $u$。
反过来，若已给出每个 $V\in\B$ 对应的 $u_V$，
则只须对 $V,W\in\B$、$V\subseteq W$ 验证它们与限制态射 $\rho_V^W$ 的交换图，
便存在唯一态射 $u'\colon\F\to\G$，使对每个 $V\in\B$ 有 $u'_V=u_V$（\hyperlink{para:0:3.2.3}{3.2.3}）。

\egapar{3.2.6}
仍假定 $\Kcat$ 具有射影极限。
则 $X$ 上取值于 $\Kcat$ 的层范畴也具有射影极限。
事实上，若 $(\F_\lambda)$ 为 $X$ 上取值于 $\Kcat$ 的层的一个射影系，则
\[
 \F(U)=\plim_\lambda\F_\lambda(U)
\]
定义了取值于 $\Kcat$ 的预层；
由射影极限的传递性可验证公理 (F)，而 $\F$ 是各 $\F_\lambda$ 的射影极限这一事实也是显然的。

当 $\Kcat$ 是集合范畴时，对任一射影系 $(\G_\lambda)$，
若对每个 $\lambda$，$\G_\lambda$ 是 $\F_\lambda$ 的子层，
则 $\plim\G_\lambda$ 典范地等同于 $\plim\F_\lambda$ 的子层。
若 $\Kcat$ 是交换群范畴，则协变函子 $\plim\F_\lambda$ 是可加且左正合的。

\subsection{层的粘合}\hypertarget{heading:0:3.3}{}
\src{28—29}

\egapar{3.3.1}
仍假定范畴 $\Kcat$ 具有广义射影极限。
设 $X$ 为拓扑空间，$\mathfrak U=(U_\lambda)_{\lambda\in L}$ 为 $X$ 的开覆盖。
对每个 $\lambda\in L$，设 $\F_\lambda$ 为 $U_\lambda$ 上取值于 $\Kcat$ 的层。
假定对每对指标 $(\lambda,\mu)$，给定同构
\[
 \theta_{\lambda\mu}\colon
 \F_\mu|(U_\lambda\cap U_\mu)
 \xrightarrow{\sim}
 \F_\lambda|(U_\lambda\cap U_\mu).
\]
再假定对每个三元组 $(\lambda,\mu,\nu)$，
以 $\theta'_{\lambda\mu}$、$\theta'_{\lambda\nu}$、$\theta'_{\mu\nu}$
表示相应同构在 $U_\lambda\cap U_\mu\cap U_\nu$ 上的限制，则有
\[
 \theta'_{\lambda\nu}=\theta'_{\lambda\mu}\circ\theta'_{\mu\nu};
\]
这是各 $\theta_{\lambda\mu}$ 的\emph{粘合条件}。

于是存在 $X$ 上取值于 $\Kcat$ 的层 $\F$，并且对每个 $\lambda$ 存在同构
$\eta_\lambda\colon\F|U_\lambda\xrightarrow{\sim}\F_\lambda$，
使得对每对 $(\lambda,\mu)$，若以 $\eta'_\lambda,\eta'_\mu$ 表示它们在
$U_\lambda\cap U_\mu$ 上的限制，则有
\[
 \theta_{\lambda\mu}=\eta'_\lambda\circ(\eta'_\mu)^{-1}.
\]
此外，这些条件将 $\F$ 及各 $\eta_\lambda$ 确定到唯一同构。
事实上，唯一性立即由（\hyperlink{para:0:3.2.5}{3.2.5}）得出。

为证明 $\F$ 的存在性，以 $\B$ 表示那些至少包含在某个 $U_\lambda$ 中的开集组成的基。
对每个 $U\in\B$，利用 Hilbert 的 $\tau$ 函数，
从所有满足 $U\subseteq U_\lambda$ 的指标 $\lambda$ 所对应的对象 $\F_\lambda(U)$ 中选取一个。
将此对象记为 $\F(U)$。
当 $U\subseteq V$、$U,V\in\B$ 时，借助各 $\theta_{\lambda\mu}$，限制态射 $\rho_U^V$ 的定义是显然的，
而传递条件是粘合条件的推论。
此外，($\mathrm F_0$) 的验证是直接的，故如此定义的基上的预层确实是层。
由一般方法（\hyperlink{para:0:3.2.1}{3.2.1}）得到一个通常层，仍记为 $\F$，它满足所要求的条件。
称 $\F$ 是\emph{通过同构 $\theta_{\lambda\mu}$ 将各层 $\F_\lambda$ 粘合所得的层}；
通常借助 $\eta_\lambda$ 将 $\F_\lambda$ 与 $\F|U_\lambda$ 等同。

显然，$X$ 上取值于 $\Kcat$ 的任一层 $\F$ 都可以看作将各
$\F_\lambda=\F|U_\lambda$ 通过恒等同构 $\theta_{\lambda\mu}$ 粘合所得的层，
其中 $(U_\lambda)$ 为 $X$ 的任意开覆盖。

\egapar{3.3.2}
沿用上述记号。
对每个 $\lambda\in L$，设 $\G_\lambda$ 为 $U_\lambda$ 上取值于 $\Kcat$ 的另一个层。
假定对每对 $(\lambda,\mu)$ 给定同构
\[
 \omega_{\lambda\mu}\colon
 \G_\mu|(U_\lambda\cap U_\mu)
 \xrightarrow{\sim}\G_\lambda|(U_\lambda\cap U_\mu),
\]
且这些同构满足粘合条件。
最后假定对每个 $\lambda$ 给定态射 $u_\lambda\colon\F_\lambda\to\G_\lambda$，并且下图交换：
\begin{equation}
\begin{tikzcd}[column sep=large,row sep=large]
\F_\mu|(U_\lambda\cap U_\mu)
 \arrow[r,"u_\mu"] \arrow[d,"\theta_{\lambda\mu}"']
 &\G_\mu|(U_\lambda\cap U_\mu) \arrow[d,"\omega_{\lambda\mu}"]\\
\F_\lambda|(U_\lambda\cap U_\mu)
 \arrow[r,"u_\lambda"'] &\G_\lambda|(U_\lambda\cap U_\mu).
\end{tikzcd}
\tag{3.3.2.1}\label{eq:0:3.3.2.1}
\end{equation}
若 $\G$ 是通过各 $\omega_{\lambda\mu}$ 粘合 $\G_\lambda$ 所得的层，
则存在唯一态射 $u\colon\F\to\G$，使下图交换：
\[
\begin{tikzcd}[column sep=large,row sep=large]
 \F|U_\lambda \arrow[r,"u|U_\lambda"] \arrow[d,"\sim"']
 & \G|U_\lambda \arrow[d,"\sim"]\\
 \F_\lambda \arrow[r,"u_\lambda"'] & \G_\lambda.
\end{tikzcd}
\]
这立即由（\hyperlink{para:0:3.2.3}{3.2.3}）得出。
族 $(u_\lambda)$ 与 $u$ 之间的对应，是从
$\prod_\lambda\Hom(\F_\lambda,\G_\lambda)$ 中满足（\hyperref[eq:0:3.3.2.1]{3.3.2.1}）的部分
到 $\Hom(\F,\G)$ 的函子性双射。

\egapar{3.3.3}
采用（\hyperlink{para:0:3.3.1}{3.3.1}）的记号，设 $V$ 为 $X$ 的开集。
立即可知，各 $\theta_{\lambda\mu}$ 在 $V\cap U_\lambda\cap U_\mu$ 上的限制，
对诱导层 $\F_\lambda|(V\cap U_\lambda)$ 满足粘合条件；
将这些诱导层粘合得到的 $V$ 上的层，典范地等同于 $\F|V$。

\subsection{预层的直像}\hypertarget{heading:0:3.4}{}
\src{29—30}

\egapar{3.4.1}
设 $X,Y$ 为两个拓扑空间，$\psi\colon X\to Y$ 为连续映射。
设 $\F$ 为 $X$ 上取值于范畴 $\Kcat$ 的预层。
对每个开集 $U\subseteq Y$，置 $\G(U)=\F(\psi^{-1}(U))$；
若 $U,V$ 为 $Y$ 的开集且 $U\subseteq V$，
以 $\rho_U^V$ 表示态射 $\F(\psi^{-1}(V))\to\F(\psi^{-1}(U))$。
立即可见，各 $\G(U)$ 及 $\rho_U^V$ 定义了 $Y$ 上取值于 $\Kcat$ 的预层，
称为 $\F$ 在 $\psi$ 下的\emph{直像}，记为 $\psi_*(\F)$。
若 $\F$ 是层，则立即可以对预层 $\G(U)$ 验证公理 (F)，故 $\psi_*(\F)$ 是层。

\egapar{3.4.2}
设 $\F_1,\F_2$ 为 $X$ 上取值于 $\Kcat$ 的两个预层，$u\colon\F_1\to\F_2$ 为态射。
当 $U$ 遍历 $Y$ 的开子集时，态射族
\[
 u_{\psi^{-1}(U)}\colon\F_1(\psi^{-1}(U))\longrightarrow\F_2(\psi^{-1}(U))
\]
满足与限制态射的相容条件，因而定义态射
$\psi_*(u)\colon\psi_*(\F_1)\to\psi_*(\F_2)$。
若 $v\colon\F_2\to\F_3$ 为到第三个这种预层的态射，则
$\psi_*(v\circ u)=\psi_*(v)\circ\psi_*(u)$。
换言之，$\psi_*(\F)$ 是关于 $\F$ 的协变函子，
从 $X$ 上取值于 $\Kcat$ 的预层范畴（分别地，层范畴），
取值于 $Y$ 上取值于 $\Kcat$ 的预层范畴（分别地，层范畴）。

\egapar{3.4.3}
设 $Z$ 为第三个拓扑空间，$\psi'\colon Y\to Z$ 为连续映射，并令 $\psi''=\psi'\circ\psi$。
显然，对 $X$ 上每个取值于 $\Kcat$ 的预层 $\F$，有
$\psi''_*(\F)=\psi'_*(\psi_*(\F))$；
对每个这样的预层态射 $u\colon\F\to\G$，还有
$\psi''_*(u)=\psi'_*(\psi_*(u))$。
因此，$\psi''_*$ 是函子 $\psi'_*$ 与 $\psi_*$ 的复合，可写为
\[
 (\psi'\circ\psi)_*=\psi'_*\circ\psi_*.
\]
此外，对 $Y$ 的任一开集 $U$，诱导预层 $\F|\psi^{-1}(U)$
在限制映射 $\psi|\psi^{-1}(U)$ 下的直像，就是诱导预层 $\psi_*(\F)|U$。

\egapar{3.4.4}
假定范畴 $\Kcat$ 具有归纳极限，$\F$ 为 $X$ 上取值于 $\Kcat$ 的预层。
对每个 $x\in X$，各态射
$\Gamma(\psi^{-1}(U),\F)\to\F_x$，其中 $U$ 为 $\psi(x)$ 在 $Y$ 中的开邻域，
构成一个归纳系；取其归纳极限得到茎之间的态射
\[
 \psi_x\colon(\psi_*(\F))_{\psi(x)}\longrightarrow\F_x.
\]
通常这个态射既不是单射也不是满射。
它是函子性的：事实上，若 $u\colon\F\to\G$ 为 $X$ 上取值于 $\Kcat$ 的预层态射，则下图交换：
\[
\begin{tikzcd}[column sep=large,row sep=large]
 (\psi_*(\F))_{\psi(x)} \arrow[r,"\psi_x"] \arrow[d,"(\psi_*(u))_{\psi(x)}"']
 & \F_x \arrow[d,"u_x"]\\
 (\psi_*(\G))_{\psi(x)} \arrow[r,"\psi_x"'] & \G_x.
\end{tikzcd}
\]
若 $Z$ 为第三个拓扑空间，$\psi'\colon Y\to Z$ 为连续映射，$\psi''=\psi'\circ\psi$，
则对每个 $x\in X$，有 $\psi''_x=\psi_x\circ\psi'_{\psi(x)}$。

\egapar{3.4.5}
在（\hyperlink{para:0:3.4.4}{3.4.4}）的假设下，再假定 $\psi$ 是从 $X$ 到 $Y$ 的子空间 $\psi(X)$ 的同胚。
则对每个 $x\in X$，$\psi_x$ 是同构。
这特别适用于 $Y$ 的子集 $X$ 到 $Y$ 的典范单射 $j$。

\egapar{3.4.6}
假定 $\Kcat$ 是群范畴、环范畴等等。
若 $\F$ 为 $X$ 上取值于 $\Kcat$ 的层，支集为 $S$，且 $y\notin\overline{\psi(S)}$，
则由 $\psi_*(\F)$ 的定义，$(\psi_*(\F))_y=\{0\}$。
也就是说，$\psi_*(\F)$ 的支集包含在 $\overline{\psi(S)}$ 中，
但未必包含在 $\psi(S)$ 中。
在相同假设下，若 $j$ 为子集 $X$ 到 $Y$ 的典范单射，则层 $j_*(\F)$ 在 $X$ 上诱导 $\F$。
若 $X$ 还是 $Y$ 的闭子集，则 $j_*(\F)$ 是 $Y$ 上在 $X$ 上诱导 $\F$、
在 $Y\setminus X$ 上诱导零层的层 (G, II, 2.9.2)；
但是，若只假定 $X$ 局部闭而不闭，则 $j_*(\F)$ 通常不等于这一层。

\subsection{预层的逆像}\hypertarget{heading:0:3.5}{}
\src{30—33}

\egapar{3.5.1}
在（\hyperlink{para:0:3.4.1}{3.4.1}）的假设下，若 $\F,\G$ 分别为 $X,Y$ 上取值于 $\Kcat$ 的预层，
则 $Y$ 上的预层态射 $u\colon\G\to\psi_*(\F)$ 也称为从 $\G$ 到 $\F$ 的一个
\emph{$\psi$-态射}，也记作 $\G\to\F$。
还以 $\Hom_\psi(\G,\F)$ 表示 $\psi$-态射的集合 $\Hom_Y(\G,\psi_*(\F))$。

对每一对 $(U,V)$，其中 $U$ 为 $X$ 的开集，$V$ 为 $Y$ 的开集，且 $\psi(U)\subseteq V$，
将限制态射 $\F(\psi^{-1}(V))\to\F(U)$ 与态射
$u_V\colon\G(V)\to\psi_*(\F)(V)=\F(\psi^{-1}(V))$ 复合，
得到态射 $u_{U,V}\colon\G(V)\to\F(U)$。
立即可见，这些态射使下图交换：
\begin{equation}
\begin{tikzcd}[column sep=large,row sep=large]
 \G(V) \arrow[r,"u_{U,V}"] \arrow[d] & \F(U) \arrow[d]\\
 \G(V') \arrow[r,"u_{U',V'}"'] & \F(U').
\end{tikzcd}
\tag{3.5.1.1}\label{eq:0:3.5.1.1}
\end{equation}
这里 $U'\subseteq U$、$V'\subseteq V$、$\psi(U')\subseteq V'$。
反过来，给定一族使图（\hyperref[eq:0:3.5.1.1]{3.5.1.1}）交换的态射 $(u_{U,V})$，就定义了一个 $\psi$-态射 $u$；
只须取 $u_V=u_{\psi^{-1}(V),V}$。

若 $\Kcat$ 具有广义射影极限，而 $\B,\B'$ 分别为 $X,Y$ 的拓扑基，
则为了定义\emph{层}的 $\psi$-态射 $u$，
只须对 $U\in\B$、$V\in\B'$、$\psi(U)\subseteq V$ 给出 $u_{U,V}$，
并对 $U,U'\in\B$、$V,V'\in\B'$ 验证相容条件（\hyperref[eq:0:3.5.1.1]{3.5.1.1}）。
事实上，对 $Y$ 的任一开集 $W$，只须将 $u_W$ 定义为
各 $u_{U,V}$ 对 $V\in\B'$、$V\subseteq W$、$U\in\B$、$\psi(U)\subseteq V$ 的射影极限。

若 $\Kcat$ 具有归纳极限，则对每个 $x\in X$ 及 $\psi(x)$ 在 $Y$ 中的任一开邻域 $V$，有态射
\[
 \G(V)\longrightarrow\F(\psi^{-1}(V))\longrightarrow\F_x.
\]
这些态射构成归纳系，取归纳极限得到态射 $\G_{\psi(x)}\to\F_x$。

\egapar{3.5.2}
在（\hyperlink{para:0:3.4.3}{3.4.3}）的假设下，设 $\F,\G,\HH$ 分别为 $X,Y,Z$ 上取值于 $\Kcat$ 的预层，
$u\colon\G\to\psi_*(\F)$、$v\colon\HH\to\psi'_*(\G)$
分别为 $\psi$-态射和 $\psi'$-态射。
由此得到一个 $\psi''$-态射
\[
 w\colon\HH\xrightarrow{v}\psi'_*(\G)
 \xrightarrow{\psi'_*(u)}\psi'_*(\psi_*(\F))=\psi''_*(\F),
\]
根据定义，称它为 $u$ 与 $v$ 的\emph{复合}。
因此，由拓扑空间 $X$ 及其上一个取值于 $\Kcat$ 的预层 $\F$ 组成的二元组 $(X,\F)$，可以看作一个范畴的对象。
其中的态射 $(\psi,\theta)\colon(X,\F)\to(Y,\G)$，
由连续映射 $\psi\colon X\to Y$ 及 $\psi$-态射 $\theta\colon\G\to\F$ 组成。

\egapar{3.5.3}
设 $\psi\colon X\to Y$ 为连续映射，$\G$ 为 $Y$ 上取值于 $\Kcat$ 的预层。
称二元组 $(\G',\rho)$ 为 $\G$ 在 $\psi$ 下的\emph{逆像}，若 $\G'$ 为 $X$ 上取值于 $\Kcat$ 的\emph{层}，
$\rho\colon\G\to\G'$ 为 $\psi$-态射，即同态 $\G\to\psi_*(\G')$，
并且对 $X$ 上取值于 $\Kcat$ 的每个层 $\F$，映射
\begin{equation}
 \Hom_X(\G',\F)\longrightarrow\Hom_\psi(\G,\F)
 =\Hom_Y(\G,\psi_*(\F)),
 \tag{3.5.3.1}\label{eq:0:3.5.3.1}
\end{equation}
将 $v$ 映为 $\psi_*(v)\circ\rho$，是一个双射。
由于此映射关于 $\F$ 是函子性的，它定义了关于 $\F$ 的函子同构。
二元组 $(\G',\rho)$ 是一个泛问题的解，故只要存在，就由上述条件确定到唯一同构。
这时写作 $\G'=\psi^*(\G)$、$\rho=\rho_\G$。
我们也略不严格地称 $\psi^*(\G)$ 为 $\G$ 在 $\psi$ 下的\emph{逆像层}，
其中总将 $\psi^*(\G)$ 看作配备了典范 $\psi$-态射
$\rho_\G\colon\G\to\psi^*(\G)$，也就是 $Y$ 上预层的典范同态
\begin{equation}
 \rho_\G\colon\G\longrightarrow\psi_*(\psi^*(\G)).
 \tag{3.5.3.2}\label{eq:0:3.5.3.2}
\end{equation}

对任一同态 $v\colon\psi^*(\G)\to\F$，其中 $\F$ 是 $X$ 上取值于 $\Kcat$ 的层，置
\[
 v^\flat=\psi_*(v)\circ\rho_\G\colon\G\longrightarrow\psi_*(\F).
\]
由定义，每个预层态射 $u\colon\G\to\psi_*(\F)$ 都唯一地写成 $v^\flat$，
将相应的 $v$ 记为 $u^\sharp$。
换言之，任一预层态射 $u\colon\G\to\psi_*(\F)$ 都唯一地分解为
\begin{equation}
 u\colon\G\xrightarrow{\rho_\G}\psi_*(\psi^*(\G))
 \xrightarrow{\psi_*(u^\sharp)}\psi_*(\F).
 \tag{3.5.3.3}\label{eq:0:3.5.3.3}
\end{equation}

\egapar{3.5.4}
现在假定范畴 $\Kcat$ 使得\origfootnote{
在导言中提到的著作里，我们将给出关于范畴 $\Kcat$ 的非常一般的条件，
以保证取值于 $\Kcat$ 的预层的逆像存在。}
$Y$ 上取值于 $\Kcat$ 的任一预层 $\G$ 都具有在 $\psi$ 下的逆像，记为 $\psi^*(\G)$。
我们将证明，可以把 $\psi^*(\G)$ 定义为关于 $\G$ 的协变函子，
从 $Y$ 上取值于 $\Kcat$ 的预层范畴取值于 $X$ 上取值于 $\Kcat$ 的层范畴，
使对应 $v\mapsto v^\flat$ 成为关于 $\G,\F$ 的双函子同构
\begin{equation}
 \Hom_X(\psi^*(\G),\F)\xrightarrow{\sim}\Hom_Y(\G,\psi_*(\F)).
 \tag{3.5.4.1}\label{eq:0:3.5.4.1}
\end{equation}

事实上，对 $Y$ 上取值于 $\Kcat$ 的任一预层态射 $w\colon\G_1\to\G_2$，考虑复合态射
\[
 \G_1\xrightarrow{w}\G_2\xrightarrow{\rho_{\G_2}}\psi_*(\psi^*(\G_2)).
\]
与它对应的态射 $(\rho_{\G_2}\circ w)^\sharp\colon\psi^*(\G_1)\to\psi^*(\G_2)$，记为 $\psi^*(w)$。
由（\hyperref[eq:0:3.5.3.3]{3.5.3.3}），有
\begin{equation}
 \psi_*(\psi^*(w))\circ\rho_{\G_1}=\rho_{\G_2}\circ w.
 \tag{3.5.4.2}\label{eq:0:3.5.4.2}
\end{equation}
对每个态射 $u\colon\G_2\to\psi_*(\F)$，其中 $\F$ 为 $X$ 上取值于 $\Kcat$ 的层，
由（\hyperref[eq:0:3.5.3.3]{3.5.3.3}）、（\hyperref[eq:0:3.5.4.2]{3.5.4.2}）及 $v^\flat$ 的定义，得到
\[
 \begin{aligned}
 (u^\sharp\circ\psi^*(w))^\flat
 &=\psi_*(u^\sharp)\circ\psi_*(\psi^*(w))\circ\rho_{\G_1}\\
 &=\psi_*(u^\sharp)\circ\rho_{\G_2}\circ w
 =u\circ w,
 \end{aligned}
\]
即
\begin{equation}
 (u\circ w)^\sharp=u^\sharp\circ\psi^*(w).
 \tag{3.5.4.3}\label{eq:0:3.5.4.3}
\end{equation}
特别取 $u$ 为态射 $\G_2\xrightarrow{w'}\G_3\xrightarrow{\rho_{\G_3}}\psi_*(\psi^*(\G_3))$，则有
\[
 \begin{aligned}
 \psi^*(w'\circ w)
 &=(\rho_{\G_3}\circ w'\circ w)^\sharp\\
 &=(\rho_{\G_3}\circ w')^\sharp\circ\psi^*(w)
 =\psi^*(w')\circ\psi^*(w),
 \end{aligned}
\]
由此得到我们的断言。

最后，对 $X$ 上取值于 $\Kcat$ 的任一层 $\F$，令 $i_\F$ 为 $\psi_*(\F)$ 的恒等态射，并以
\[
 \sigma_\F\colon\psi^*(\psi_*(\F))\longrightarrow\F
\]
表示态射 $(i_\F)^\sharp$。
公式（\hyperref[eq:0:3.5.4.3]{3.5.4.3}）特别给出下列分解：
\begin{equation}
 u^\sharp\colon\psi^*(\G)\xrightarrow{\psi^*(u)}\psi^*(\psi_*(\F))
 \xrightarrow{\sigma_\F}\F,
 \tag{3.5.4.4}\label{eq:0:3.5.4.4}
\end{equation}
其中 $u\colon\G\to\psi_*(\F)$ 为任一态射。
称 $\sigma_\F$ 为\emph{典范态射}。

\egapar{3.5.5}
设 $\psi'\colon Y\to Z$ 为连续映射，并假定 $Z$ 上取值于 $\Kcat$ 的任一预层 $\HH$
都有在 $\psi'$ 下的逆像 $\psi'^*(\HH)$。
于是，在（\hyperlink{para:0:3.5.4}{3.5.4}）的假设下，$Z$ 上取值于 $\Kcat$ 的每个预层 $\HH$，
都有在 $\psi''=\psi'\circ\psi$ 下的逆像，并有典范函子性同构
\begin{equation}
 \psi''{}^*(\HH)\xrightarrow{\sim}\psi^*(\psi'^*(\HH)).
 \tag{3.5.5.1}\label{eq:0:3.5.5.1}
\end{equation}
事实上，考虑到 $\psi''_*=\psi'_*\circ\psi_*$，这立即由定义得出。
此外，若 $u\colon\G\to\psi_*(\F)$ 为 $\psi$-态射，
$v\colon\HH\to\psi'_*(\G)$ 为 $\psi'$-态射，
$w=\psi'_*(u)\circ v$ 为其复合（\hyperlink{para:0:3.5.2}{3.5.2}），则立即可验证 $w^\sharp$ 为复合态射
\[
 w^\sharp\colon\psi^*(\psi'^*(\HH))
 \xrightarrow{\psi^*(v^\sharp)}\psi^*(\G)\xrightarrow{u^\sharp}\F.
\]

\egapar{3.5.6}
特别取 $\psi$ 为恒等映射 $i_X\colon X\to X$。
若 $X$ 上取值于 $\Kcat$ 的预层 $\F$ 在 $\psi$ 下的逆像存在，
则称这个逆像为\emph{与预层 $\F$ 相伴的层}。
因此，从 $\F$ 到一个取值于 $\Kcat$ 的层 $\F'$ 的任一态射 $u\colon\F\to\F'$，
都唯一地分解为
\[
 \F\longrightarrow i_X^*(\F)\longrightarrow\F'.
\]

\subsection{常值层与局部常值层}\hypertarget{heading:0:3.6}{}
\src{33}

\egapar{3.6.1}
若 $X$ 上取值于 $\Kcat$ 的预层 $\F$，对每个非空开集 $U\subseteq X$ 都使典范态射
$\F(X)\to\F(U)$ 为同构，则称 $\F$ 为\emph{常值预层}。
注意，$\F$ 未必是层。
若一个层与某个常值预层相伴（\hyperlink{para:0:3.5.6}{3.5.6}），则称其为\emph{常值层}
（原文：\textit{faisceau simple}）。
若每个 $x\in X$ 都有一个开邻域 $U$，使 $\F|U$ 为常值层，
则称层 $\F$ 为\emph{局部常值层}（原文：\textit{faisceau localement simple}）。

\egapar{3.6.2}
假定 $X$ 不可约（\hyperlink{para:0:2.1.1}{2.1.1}）。则以下性质等价：
\begin{enumerate}[label=\alph*)]
\item $\F$ 是 $X$ 上的常值预层；
\item $\F$ 是 $X$ 上的常值层；
\item $\F$ 是 $X$ 上的局部常值层。
\end{enumerate}
事实上，设 $\F$ 为 $X$ 上的常值预层。
若 $U,V$ 为 $X$ 的两个非空开集，则 $U\cap V$ 非空；
由于 $\F(X)\to\F(U)\to\F(U\cap V)$ 的复合以及 $\F(X)\to\F(U)$ 都是同构，
故 $\F(U)\to\F(U\cap V)$ 也是同构。
同样，$\F(V)\to\F(U\cap V)$ 也是同构。
由此立即得到（\hyperlink{para:0:3.1.2}{3.1.2}）的公理 (F)，故 $\F$ 同构于其相伴层，因而 a) 蕴含 b)。

现在，设 $(U_\alpha)$ 为 $X$ 的一个由非空开集组成的开覆盖，
而 $\F$ 为 $X$ 上的一个层，使得每个 $\F|U_\alpha$ 都是常值层。
由于 $U_\alpha$ 不可约，由前述讨论，$\F|U_\alpha$ 为常值预层。
由于 $U_\alpha\cap U_\beta$ 非空，两个态射
$\F(U_\alpha)\to\F(U_\alpha\cap U_\beta)$ 与
$\F(U_\beta)\to\F(U_\alpha\cap U_\beta)$ 均为同构，
从而对每对指标得到典范同构 $\theta_{\alpha\beta}\colon\F(U_\alpha)\to\F(U_\beta)$。
这时对 $U=X$ 应用条件 (F)，可见对每个指标 $\alpha_0$，
$\F(U_{\alpha_0})$ 及各 $\theta_{\alpha_0\alpha}$ 也是相应泛问题的解。
由唯一性可知，$\F(X)\to\F(U_{\alpha_0})$ 为同构，故 c) 蕴含 a)。

\subsection{群预层或环预层的逆像}\hypertarget{heading:0:3.7}{}
\src{34}

\egapar{3.7.1}
我们将证明，当 $\Kcat$ 取为集合范畴时，任一取值于 $\Kcat$ 的预层 $\G$ 在 $\psi$ 下的逆像总是存在。
关于 $X,Y,\psi$ 的记号及假设沿用（\hyperlink{para:0:3.5.3}{3.5.3}）。
事实上，对任一开集 $U\subseteq X$，如下定义 $\G'(U)$：
其元素 $s'$ 是一个族 $(s'_x)_{x\in U}$，其中对每个 $x\in U$ 有 $s'_x\in\G_{\psi(x)}$；
且对每个 $x\in U$，满足下述条件：
存在 $\psi(x)$ 在 $Y$ 中的一个开邻域 $V$，
$x$ 的一个邻域 $W\subseteq\psi^{-1}(V)\cap U$，以及元素 $s\in\G(V)$，
使得对每个 $z\in W$ 有 $s'_z=s_{\psi(z)}$。
立即可以验证，$U\mapsto\G'(U)$ 满足层的公理。

现在设 $\F$ 为 $X$ 上的集合层，
$u\colon\G\to\psi_*(\F)$ 与 $v\colon\G'\to\F$ 为态射。
如下定义 $u^\sharp$ 和 $v^\flat$：
若 $s'$ 是 $\G'$ 在点 $x\in X$ 的一个邻域 $U$ 上的截面，
而 $V$ 为 $\psi(x)$ 的开邻域，$s\in\G(V)$ 使得对 $x$ 的某个包含在
$\psi^{-1}(V)\cap U$ 中的邻域中的每个 $z$，都有 $s'_z=s_{\psi(z)}$，
则规定 $u^\sharp(s')$ 在 $x$ 的芽为 $s_{\psi(x)}$ 经（\hyperlink{para:0:3.5.1}{3.5.1}）所定义的茎映射后的像\transnote{
原版此处的公式右端写为 $u_{\psi(x)}(s_{\psi(x)})$。
若 $u_{\psi(x)}$ 按 $Y$ 上预层同态 $u\colon\G\to\psi_*(\F)$ 的茎映射理解，
其值位于 $(\psi_*(\F))_{\psi(x)}$，还须接上（\hyperlink{para:0:3.4.4}{3.4.4}）的比较映射 $\psi_x$ 才位于 $\F_x$。
译文将这个复合显式写出；它正是（\hyperlink{para:0:3.5.1}{3.5.1}）末段使用的映射。}，即
\[
 (u^\sharp(s'))_x=\psi_x\bigl(u_{\psi(x)}(s_{\psi(x)})\bigr).
\]

同样，若 $s\in\G(V)$，其中 $V$ 为 $Y$ 的开集，则 $v^\flat(s)$ 是 $\F$ 在 $\psi^{-1}(V)$ 上的截面，
它是 $\G'$ 的截面 $s'$ 在 $v$ 下的像，而 $s'$ 由对每个 $x\in\psi^{-1}(V)$ 规定
$s'_x=s_{\psi(x)}$ 给出。

此外，典范同态（\hyperlink{para:0:3.5.3}{3.5.3}）$\rho\colon\G\to\psi_*(\psi^*(\G))$ 如下定义：
对每个开集 $V\subseteq Y$ 及每个截面 $s\in\Gamma(V,\G)$，
$\rho(s)$ 是 $\psi^*(\G)$ 在 $\psi^{-1}(V)$ 上的截面
$(s_{\psi(x)})_{x\in\psi^{-1}(V)}$。
关系
\[
 (u^\sharp)^\flat=u,\qquad (v^\flat)^\sharp=v,
 \qquad v^\flat=\psi_*(v)\circ\rho
\]
的验证是直接的；由此证明了我们的断言。

若 $w\colon\G_1\to\G_2$ 为 $Y$ 上的集合预层同态，
可以验证，$\psi^*(w)$ 的显式描述如下：
若 $s'=(s'_x)_{x\in U}$ 为 $\psi^*(\G_1)$ 在 $X$ 的开集 $U$ 上的截面，
则 $(\psi^*(w))(s')$ 就是族 $(w_{\psi(x)}(s'_x))_{x\in U}$。
最后，立即可知，对 $Y$ 的任一开集 $V$，$\G|V$ 在 $\psi$ 限制于 $\psi^{-1}(V)$ 的映射下的逆像，
与诱导层 $\psi^*(\G)|\psi^{-1}(V)$ 相同。

当 $\psi$ 为恒等映射 $i_X$ 时，就重新得到与预层相伴的集合层的定义 (G, II, 1.2)。
当 $\Kcat$ 为群范畴或不一定交换的环的范畴时，上述讨论无须改动仍然适用。

若 $X$ 为拓扑空间 $Y$ 的任意子集，$j\colon X\to Y$ 为典范单射，
则对 $Y$ 上取值于范畴 $\Kcat$ 的任一层 $\G$，
称逆像 $j^*(\G)$（若存在）为\emph{$\G$ 在 $X$ 上诱导的层}。
对集合层、群层或环层，这重新给出了通常的定义 (G, II, 1.5)。

\egapar{3.7.2}
保留（\hyperlink{para:0:3.5.3}{3.5.3}）的记号与假设，并假定 $\G$ 为 $Y$ 上的群层（分别地，环层）。
由 $\psi^*(\G)$ 的截面定义（\hyperlink{para:0:3.7.1}{3.7.1}），结合（\hyperlink{para:0:3.4.4}{3.4.4}）可知，茎的同态
\[
 \psi_x\circ\rho_{\psi(x)}\colon
 \G_{\psi(x)}\longrightarrow(\psi^*(\G))_x
\]
是关于 $\G$ 的函子性同构，从而可以将这两个茎等同。
在此等同下，$u_x^\sharp$ 就是（\hyperlink{para:0:3.5.1}{3.5.1}）中定义的同态。
特别地，有
\[
 \Supp(\psi^*(\G))=\psi^{-1}(\Supp(\G)).
\]
这一结果的直接推论是：在交换群层的阿贝尔范畴中，函子 $\psi^*(\G)$ 关于 $\G$ 是正合的。

\subsection{伪离散空间层}\hypertarget{heading:0:3.8}{}
\src{35；止于（\hyperlink{para:0:3.8.3}{3.8.3}）}

\egapar{3.8.1}
设 $X$ 为拓扑空间，其拓扑具有由拟紧开集组成的基 $\B$。
设 $\F$ 为 $X$ 上的集合层。
若给每个 $\F(U)$ 赋予离散拓扑，则 $U\mapsto\F(U)$ 为拓扑空间预层。
我们将证明，存在与它相伴的拓扑空间层 $\F'$（\hyperlink{para:0:3.5.6}{3.5.6}），
使得对每个拟紧开集 $U$，$\Gamma(U,\F')$ 就是离散空间 $\F(U)$。
为此，只须证明 $\B$ 上的离散拓扑空间预层 $U\mapsto\F(U)$ 满足（\hyperlink{para:0:3.2.2}{3.2.2}）的条件 ($\mathrm F_0$)。
更一般地，只须证明：若 $U$ 是拟紧开集，$(U_\alpha)$ 是由 $\B$ 的元素组成的 $U$ 的覆盖，
则 $\Gamma(U,\F)$ 上使各映射 $\Gamma(U,\F)\to\Gamma(U_\alpha,\F)$ 连续的最粗拓扑
$\mathscr T$ 是离散拓扑。

现在存在有限个指标 $\alpha_i$ 使 $U=\bigcup_i U_{\alpha_i}$。
设 $s\in\Gamma(U,\F)$，以 $s_i$ 表示它在 $\Gamma(U_{\alpha_i},\F)$ 中的像。
根据定义，单点集 $\{s_i\}$ 的各逆像的交是 $s$ 在拓扑 $\mathscr T$ 下的一个邻域。
但由于 $\F$ 为集合层，且各 $U_{\alpha_i}$ 覆盖 $U$，这个交只有 $s$ 一个元素；
由此得到我们的断言。

注意，若 $U$ 是 $X$ 的非拟紧开集，则拓扑空间 $\Gamma(U,\F')$ 的底层集合仍为 $\Gamma(U,\F)$，
但其拓扑一般不离散：它是使各映射
\[
 \Gamma(U,\F)\longrightarrow\Gamma(V,\F)
 \qquad(V\in\B,\ V\subseteq U)
\]
连续的最粗拓扑，其中各 $\Gamma(V,\F)$ 带离散拓扑。

上述讨论无须改动就适用于群层或不一定交换的环层，并分别为它们配备拓扑群层或拓扑环层。
为简便起见，称层 $\F'$ 为与集合层（分别地，群层、环层）$\F$ 相伴的
\emph{伪离散空间层}（分别地，\emph{伪离散群层}、\emph{伪离散环层}）。

\egapar{3.8.2}
设 $\F,\G$ 为 $X$ 上的两个集合层（分别地，群层、环层），$u\colon\F\to\G$ 为同态。
以 $\F',\G'$ 分别表示与 $\F,\G$ 相伴的伪离散层，则 $u$ 也是一个连续同态 $\F'\to\G'$。
事实上，这由（\hyperlink{para:0:3.2.5}{3.2.5}）得出。

\egapar{3.8.3}
设 $\F$ 为集合层，$\HH$ 为 $\F$ 的子层，$\F',\HH'$ 分别为与 $\F,\HH$ 相伴的伪离散层。
则对每个开集 $U\subseteq X$，$\Gamma(U,\HH')$ 在 $\Gamma(U,\F')$ 中是闭的。
事实上，它是各 $\Gamma(V,\HH)$（$V\in\B$、$V\subseteq U$）在连续映射
\[
 \Gamma(U,\F')\longrightarrow\Gamma(V,\F)
\]
下的逆像之交，而 $\Gamma(V,\HH)$ 在离散空间 $\Gamma(V,\F)$ 中是闭的。


% END chapters/ch0_s3.tex

\par

% BEGIN chapters/ch0_s4.tex
\section{赋环空间}\hypertarget{heading:0:4}{}
\subsection{赋环空间、\texorpdfstring{$\A$}{A}-模层与\texorpdfstring{$\A$}{A}-代数层}\hypertarget{heading:0:4.1}{}
\src{35—39}

\egapar{4.1.1}
一个\emph{赋环空间}（分别地，\emph{赋拓扑环空间}）是一个二元组 $(X,\A)$，
其中 $X$ 是拓扑空间，$\A$ 是 $X$ 上的环层（环不必交换；分别地，$\A$ 是拓扑环层）。
称 $X$ 为赋环空间 $(X,\A)$ 的\emph{底拓扑空间}，称 $\A$ 为其\emph{结构层}。
结构层记作 $\OO_X$；它在点 $x\in X$ 的茎记作 $\OO_{X,x}$，在不会引起混淆时也简记为 $\OO_x$。

用 $1$ 或 $e$ 表示 $\OO_X$ 在 $X$ 上的\emph{单位截面}，即环 $\Gamma(X,\OO_X)$ 的单位元。

本专著主要考察交换环层；因此，未加说明地谈到赋环空间 $(X,\A)$ 时，均默认 $\A$ 是交换环层。

结构层不必交换的赋环空间（分别地，赋拓扑环空间）构成一个范畴：
把态射 $(X,\A)\to(Y,\BB)$ 定义为二元组 $\Psi=(\psi,\theta)$，
其中 $\psi\colon X\to Y$ 是连续映射，$\theta\colon\BB\to\A$ 是环层的一个 $\psi$-态射（\hyperlink{para:0:3.5.1}{3.5.1}）
（分别地，是拓扑环层的一个 $\psi$-态射）。
另一个态射 $\Psi'=(\psi',\theta')\colon(Y,\BB)\to(Z,\C)$ 与 $\Psi$ 的复合，
记为 $\Psi''=\Psi'\circ\Psi$，是态射 $(\psi'',\theta'')$，其中
\[
 \psi''=\psi'\circ\psi,\qquad
 \theta''=\psi'_*(\theta)\circ\theta'
\]
（参见（\hyperlink{para:0:3.5.2}{3.5.2}））。对于赋环空间，回顾（\hyperlink{para:0:3.5.5}{3.5.5}），有
\[
 \theta''{}^\sharp=\theta^\sharp\circ\psi^*(\theta'{}^\sharp).
\]
因此，若 $\theta^\sharp$ 与 $\theta'{}^\sharp$ 都是单射同态（分别地，满射同态），
则 $\theta''{}^\sharp$ 也如此；这里用到了对每个 $x\in X$，
$\psi_x\circ\rho_{\psi(x)}$ 都是同构这一事实（\hyperlink{para:0:3.7.2}{3.7.2}）。
由以上事实立即验证：当 $\psi$ 是单射连续映射，且 $\theta^\sharp$ 是环层的满射同态时，
$(\psi,\theta)$ 是赋环空间范畴中的\emph{单态射}（T，I.1）。

在不会引起混淆时，我们往往略去记号上的区别，以 $\Psi$ 代替 $\psi$；
例如，对于 $Y$ 的子集 $U$，写 $\Psi^{-1}(U)$ 而不写 $\psi^{-1}(U)$。

\egapar{4.1.2}
对于 $X$ 的任意子集 $M$，二元组 $(M,\A|M)$ 显然是赋环空间，称为由 $(X,\A)$
\emph{在 $M$ 上诱导的赋环空间}，也称为 $(X,\A)$ \emph{在 $M$ 上的限制}。
若 $j\colon M\to X$ 是典范包含映射，$\omega$ 是 $\A|M$ 的恒等映射，则
\[
 (j,\omega^\flat)\colon(M,\A|M)\longrightarrow(X,\A)
\]
是赋环空间的单态射，称为\emph{典范包含态射}。
态射 $\Psi\colon(X,\A)\to(Y,\BB)$ 与这一包含态射的复合，称为 $\Psi$ 在 $M$ 上的限制。

\egapar{4.1.3}
我们不再重述赋环空间 $(X,\A)$ 上的 $\A$-\emph{模层}，或\emph{代数层}的定义（G，II，2.2）。
当 $\A$ 是不必交换的环层时，除非明确作相反说明，$\A$-模层始终指\emph{左 $\A$-模层}。
$\A$ 的子模层称为 $\A$ 中的\emph{理想层}（左、右或双边理想层），也称为 $\A$-\emph{理想层}。

当 $\A$ 是交换环层时，在 $\A$-模层的定义中处处以代数结构代替模结构，
就得到 $X$ 上的 $\A$-\emph{代数层}的定义。
等价地，一个不必交换的 $\A$-代数层是一个 $\A$-模层 $\C$，连同一个 $\A$-模层同态
$\varphi\colon\C\otimes_{\A}\C\to\C$ 和一个在 $X$ 上的截面 $e$，满足：
第一，图
\[
\begin{tikzcd}[column sep=large]
\C\otimes_{\A}\C\otimes_{\A}\C
 \arrow[r,"\varphi\otimes1"]\arrow[d,"1\otimes\varphi"'] &
\C\otimes_{\A}\C\arrow[d,"\varphi"]\\
\C\otimes_{\A}\C\arrow[r,"\varphi"'] & \C
\end{tikzcd}
\]
交换；第二，对每个开集 $U\subseteq X$ 及每个截面 $s\in\Gamma(U,\C)$，有
\[
 \varphi\bigl((e|U)\otimes s\bigr)
 =\varphi\bigl(s\otimes(e|U)\bigr)=s.
\]
称 $\C$ 是\emph{交换的} $\A$-代数层，还要求图
\[
\begin{tikzcd}[column sep=large]
 \C\otimes_{\A}\C\arrow[rr,"\sigma"]\arrow[dr,"\varphi"']&&
 \C\otimes_{\A}\C\arrow[dl,"\varphi"]\\
 &\C&
\end{tikzcd}
\]
交换，其中 $\sigma$ 是张量积 $\C\otimes_{\A}\C$ 的典范对称同构。

$\A$-代数层同态也像（G，II，2.2）中的 $\A$-模层同态那样定义，
但它们当然不再构成一个阿贝尔群。

若 $\F$ 是 $\A$-代数层 $\C$ 的一个子 $\A$-模层，
则\emph{由 $\F$ 生成的 $\C$ 的子 $\A$-代数层}是各同态
\[
 \bigotimes_{\A}^{n}\F\longrightarrow\C\qquad(n\geq0)
\]
的像之和。它也是下述代数预层 $U\mapsto\BB(U)$ 的相伴层：
$\BB(U)$ 是由子模 $\Gamma(U,\F)$ 生成的 $\Gamma(U,\C)$ 的子代数。

\egapar{4.1.4}
称拓扑空间 $X$ 上的环层 $\A$ \emph{在点 $x$ 约化}，是指茎 $\A_x$ 是约化环（\hyperlink{para:0:1.1.1}{1.1.1}）；
称 $\A$ \emph{约化}，是指它在 $X$ 的每一点都约化。
回顾：一个环 $A$ 称为\emph{正则环}，是指对 $A$ 的每个素理想 $\mathfrak p$，
局部环 $A_{\mathfrak p}$ 都是正则局部环。
称 $X$ 上的环层 $\A$ \emph{在点 $x$ 正则}（分别地，\emph{正则}），
是指茎 $\A_x$ 是正则环（分别地，$\A$ 在每一点都正则）。
最后，称 $\A$ \emph{在点 $x$ 正规}（分别地，\emph{正规}），是指 $\A_x$ 是整闭整环
（分别地，$\A$ 在每一点都正规）。
称赋环空间 $(X,\A)$ 具有上述某一性质，是指其环层 $\A$ 具有该性质。

一个\emph{分次环层} $\A$，按定义是作为阿贝尔群层等于一族阿贝尔群层
$(\A_n)_{n\in\mathbf Z}$ 的直和（G，II，2.7）的环层，并满足
$\A_m\A_n\subseteq\A_{m+n}$。
一个\emph{分次 $\A$-模层}是一个 $\A$-模层 $\F$，它等于一族阿贝尔群层
$(\F_n)_{n\in\mathbf Z}$ 的直和，并满足 $\A_m\F_n\subseteq\F_{m+n}$。
这些条件也分别等价于在每一点 $x$ 都有
\[
 (\A_m)_x(\A_n)_x\subseteq(\A_{m+n})_x,
 \qquad
 (\A_m)_x(\F_n)_x\subseteq(\F_{m+n})_x.
\]

\egapar{4.1.5}
给定一个结构层不必交换的赋环空间 $(X,\A)$，我们不再重述双函子
\[
 \F\otimes_{\A}\G,\qquad
 \shHom_{\A}(\F,\G),\qquad
 \Hom_{\A}(\F,\G)
\]
的定义（G，II，2.8 与 2.2）。视具体情形，它们定义在左或右 $\A$-模层的范畴上，
取值于阿贝尔群层的范畴（更一般地，若 $\C$ 是 $\A$ 的中心，则取值于 $\C$-模层的范畴）。
\transnote{本句依原文并列三种记号；其中前两种为层，而 $\Hom_{\A}(\F,\G)$ 表示模层的全体同态所成的阿贝尔群。后者与同态层的关系由下文的截面公式给出。}
对于每个 $x\in X$，茎 $(\F\otimes_{\A}\G)_x$ 典范地等同于
$\F_x\otimes_{\A_x}\G_x$。
还可以定义典范函子性同态
\[
 \bigl(\shHom_{\A}(\F,\G)\bigr)_x
 \longrightarrow\Hom_{\A_x}(\F_x,\G_x),
\]
它一般既不是单射，也不是满射。
上述双函子都是加性的，特别地，与有限直和可交换。
$\F\otimes_{\A}\G$ 关于 $\F$ 和 $\G$ 均右正合，并与归纳极限可交换；
$\A\otimes_{\A}\G$ 与 $\F\otimes_{\A}\A$ 分别典范地等同于 $\G$ 与 $\F$。
函子 $\shHom_{\A}(\F,\G)$ 与 $\Hom_{\A}(\F,\G)$ 关于 $\F$ 和 $\G$ 均左正合。
更准确地说，若有正合列 $0\to\G'\to\G\to\G''$，则
\[
 0\longrightarrow\shHom_{\A}(\F,\G')
 \longrightarrow\shHom_{\A}(\F,\G)
 \longrightarrow\shHom_{\A}(\F,\G'')
\]
正合；若有正合列 $\F'\to\F\to\F''\to0$，则
\[
 0\longrightarrow\shHom_{\A}(\F'',\G)
 \longrightarrow\shHom_{\A}(\F,\G)
 \longrightarrow\shHom_{\A}(\F',\G)
\]
正合。函子 $\Hom$ 也具有相应的性质。
此外，$\shHom_{\A}(\A,\G)$ 典范地等同于 $\G$；对每个开集 $U\subseteq X$，有
\[
 \Gamma\bigl(U,\shHom_{\A}(\F,\G)\bigr)
 =\Hom_{\A|U}(\F|U,\G|U).
\]

对于每个左 $\A$-模层（分别地，右 $\A$-模层）$\F$，
称右 $\A$-模层（分别地，左 $\A$-模层）$\shHom_{\A}(\F,\A)$ 为 $\F$ 的\emph{对偶}，
记为 $\dual{\F}$。

最后，若 $\A$ 是交换环层，$\F$ 是 $\A$-模层，则
$U\mapsto\bigwedge^p\Gamma(U,\F)$ 是一个预层，
其相伴层是一个 $\A$-模层，记为 $\bigwedge^p\F$，称为 $\F$ 的\emph{第 $p$ 次外幂}。
容易验证，从预层 $U\mapsto\bigwedge^p\Gamma(U,\F)$ 到相伴层 $\bigwedge^p\F$ 的典范映射是单射，
并且对每个 $x\in X$，有
\[
 \left(\bigwedge^p\F\right)_x=\bigwedge^p(\F_x).
\]
显然，$\bigwedge^p\F$ 是关于 $\F$ 的协变函子。
\transnote{本句关于典范映射为单射的断言，在所列的一般性下不成立，而不仅是对任意预层不能自动成立。
在（I，\hyperlink{para:1:2.3.2}{2.3.2}）由 $s=t^{-1}$ 粘合的两张仿射图上，取按 $e_s=t e_t$ 粘合的可逆层 $\LL$。
则 $\Gamma(X,\OO_X)=k$，而 $\LL$ 的全局截面对应于
$b(t)=t\,a(1/t)$ 的多项式对，故 $\Gamma(X,\LL)\simeq k^2$。
另一方面，$\bigwedge^2\LL=0$，所以非零的
$\bigwedge_k^2\Gamma(X,\LL)\simeq k$ 到 $\Gamma(X,\bigwedge^2\LL)=0$ 的典范映射不是单射。
正文保留原著断言；上述反例为译者补充。各外幂分别在相应的截面环或茎环上取。}

\egapar{4.1.6}
设 $\A$ 是不必交换的环层，$\J$ 是 $\A$ 的左理想层，$\F$ 是左 $\A$-模层。
记 $\J\F$ 为典范映射 $\J\otimes_{\mathbf Z}\F\to\F$ 的像所给出的 $\F$ 的子 $\A$-模层，
这里 $\mathbf Z$ 表示常值预层 $U\mapsto\mathbf Z$ 的相伴层。
显然，对于每个 $x\in X$，有 $(\J\F)_x=\J_x\F_x$。
当 $\A$ 交换时，$\J\F$ 也是 $\J\otimes_{\A}\F\to\F$ 的典范像。
立即可见，$\J\F$ 也是预层
\[
 U\longmapsto\Gamma(U,\J)\Gamma(U,\F)
\]
的相伴 $\A$-模层。
若 $\J_1,\J_2$ 是 $\A$ 的两个左理想层，则
$\J_1(\J_2\F)=(\J_1\J_2)\F$。

\egapar{4.1.7}
设 $(X_\lambda,\A_\lambda)_{\lambda\in L}$ 是一族赋环空间。
对于每对指标 $(\lambda,\mu)$，假定给定了 $X_\lambda$ 的一个开子集 $V_{\lambda\mu}$，以及赋环空间的同构
\[
 \varphi_{\mu\lambda}\colon
 (V_{\lambda\mu},\A_\lambda|V_{\lambda\mu})
 \xrightarrow{\sim}
 (V_{\mu\lambda},\A_\mu|V_{\mu\lambda}),
\]
其中 $V_{\lambda\lambda}=X_\lambda$，$\varphi_{\lambda\lambda}$ 为恒等态射。
再假定，对于每个三元组 $(\lambda,\mu,\nu)$，若以 $\varphi'_{\mu\lambda}$ 表示
$\varphi_{\mu\lambda}$ 在 $V_{\lambda\mu}\cap V_{\lambda\nu}$ 上的限制，则它是
\[
 (V_{\lambda\mu}\cap V_{\lambda\nu},
  \A_\lambda|(V_{\lambda\mu}\cap V_{\lambda\nu}))
\]
到
\[
 (V_{\mu\lambda}\cap V_{\mu\nu},
  \A_\mu|(V_{\mu\lambda}\cap V_{\mu\nu}))
\]
的同构，并且有
\[
 \varphi'_{\lambda\nu}
 =\varphi'_{\lambda\mu}\circ\varphi'_{\mu\nu}
\]
（各 $\varphi_{\mu\lambda}$ 的\emph{粘合条件}）。
首先，沿各 $V_{\lambda\mu}$，利用 $\varphi_{\mu\lambda}$ 把各 $X_\lambda$ 粘合，得到拓扑空间 $X$。
把 $X_\lambda$ 与 $X$ 中相应的开子集 $X'_\lambda$ 等同，则由假设，三个集合
\[
 V_{\lambda\mu}\cap V_{\lambda\nu},\quad
 V_{\mu\lambda}\cap V_{\mu\nu},\quad
 V_{\nu\lambda}\cap V_{\nu\mu}
\]
都与 $X'_\lambda\cap X'_\mu\cap X'_\nu$ 等同。
于是可以把 $X_\lambda$ 的赋环空间结构搬移到 $X'_\lambda$ 上。
若 $\A'_\lambda$ 是由 $\A_\lambda$ 搬移所得的环层，则各 $\A'_\lambda$ 满足粘合条件（\hyperlink{para:0:3.3.1}{3.3.1}），
因而定义了 $X$ 上的环层 $\A$。
称 $(X,\A)$ 是\emph{沿各 $V_{\lambda\mu}$、利用各 $\varphi_{\mu\lambda}$ 粘合各
$(X_\lambda,\A_\lambda)$ 所得的赋环空间}。

\subsection{\texorpdfstring{$\A$}{A}-模层的直像}\hypertarget{heading:0:4.2}{}
\src{39—40}

\egapar{4.2.1}
设 $(X,\A),(Y,\BB)$ 为赋环空间，
$\Psi=(\psi,\theta)\colon(X,\A)\to(Y,\BB)$ 为态射。
于是 $\psi_*(\A)$ 是 $Y$ 上的环层，$\theta\colon\BB\to\psi_*(\A)$ 是环层同态。
设 $\F$ 为 $\A$-模层；它的直像 $\psi_*(\F)$ 是 $Y$ 上的阿贝尔群层。
对每个开集 $U\subseteq Y$，
\[
 \Gamma(U,\psi_*(\F))=\Gamma(\psi^{-1}(U),\F)
\]
是环 $\Gamma(U,\psi_*(\A))=\Gamma(\psi^{-1}(U),\A)$ 上的模。
定义这些模结构的双线性映射与限制运算相容，因此在 $\psi_*(\F)$ 上定义了
$\psi_*(\A)$-模层结构。
再利用同态 $\theta\colon\BB\to\psi_*(\A)$，就得到 $\psi_*(\F)$ 上的 $\BB$-模层结构；
称这个 $\BB$-模层为\emph{$\F$ 在态射 $\Psi$ 下的直像}，记作 $\Psi_*(\F)$。

若 $\F_1,\F_2$ 是 $X$ 上的两个 $\A$-模层，$u\colon\F_1\to\F_2$ 为 $\A$-同态，
则考察 $Y$ 上各开集的截面，立即可见 $\psi_*(u)$ 是
$\psi_*(\F_1)\to\psi_*(\F_2)$ 的 $\psi_*(\A)$-同态，因而也是
$\Psi_*(\F_1)\to\Psi_*(\F_2)$ 的 $\BB$-同态。
作为 $\BB$-同态，把它记作 $\Psi_*(u)$。
这样，$\Psi_*$ 是从 $\A$-模层范畴到 $\BB$-模层范畴的协变函子。
而且，立即可见该函子左正合（G，II，2.12）。

在 $\psi_*(\A)$ 上，$\BB$-模层结构与环层结构共同定义了一个 $\BB$-代数层结构；
将这一 $\BB$-代数层记作 $\Psi_*(\A)$。

\egapar{4.2.2}
设 $\M,\N$ 是两个 $\A$-模层。对 $Y$ 的每个开集 $U$，有典范映射
\[
 \Gamma(\psi^{-1}(U),\M)\times\Gamma(\psi^{-1}(U),\N)
 \longrightarrow\Gamma(\psi^{-1}(U),\M\otimes_{\A}\N).
\]
它在环 $\Gamma(\psi^{-1}(U),\A)=\Gamma(U,\psi_*(\A))$ 上双线性，
因而在 $\Gamma(U,\BB)$ 上也双线性。于是得到同态
\[
 \Gamma(U,\Psi_*(\M))\otimes_{\Gamma(U,\BB)}\Gamma(U,\Psi_*(\N))
 \longrightarrow\Gamma(U,\Psi_*(\M\otimes_{\A}\N)).
\]
容易验证，这些同态与限制运算相容，因而给出一个典范函子性 $\BB$-模层同态
\begin{equation}
 \Psi_*(\M)\otimes_{\BB}\Psi_*(\N)
 \longrightarrow\Psi_*(\M\otimes_{\A}\N).
 \tag{4.2.2.1}\label{eq:0:4.2.2.1}
\end{equation}
这个同态一般既不是单射，也不是满射。
若 $\PP$ 是第三个 $\A$-模层，则立即验证图
\begin{equation}
\begin{tikzcd}[column sep=small,row sep=large]
\Psi_*(\M)\otimes_{\BB}\Psi_*(\N)\otimes_{\BB}\Psi_*(\PP)
 \arrow[r]\arrow[d] &
\Psi_*(\M\otimes_{\A}\N)\otimes_{\BB}\Psi_*(\PP)\arrow[d]\\
\Psi_*(\M)\otimes_{\BB}\Psi_*(\N\otimes_{\A}\PP)\arrow[r]&
\Psi_*(\M\otimes_{\A}\N\otimes_{\A}\PP)
\end{tikzcd}
\tag{4.2.2.2}\label{eq:0:4.2.2.2}
\end{equation}
交换。

\egapar{4.2.3}
设 $\M,\N$ 为两个 $\A$-模层。
对每个开集 $U\subseteq Y$，为简化记号，置 $V=\psi^{-1}(U)$；依定义，有
\[
 \Gamma(\psi^{-1}(U),\shHom_{\A}(\M,\N))
 =\Hom_{\A|V}(\M|V,\N|V).
\]
映射 $u\mapsto\Psi_*(u)$ 是同态
\[
 \Hom_{\A|V}(\M|V,\N|V)
 \longrightarrow
 \Hom_{\BB|U}(\Psi_*(\M)|U,\Psi_*(\N)|U),
\]
这里同态是就 $\Gamma(U,\BB)$-模结构而言的。
这些同态与限制运算相容，因而定义了典范函子性 $\BB$-模层同态
\begin{equation}
 \Psi_*(\shHom_{\A}(\M,\N))
 \longrightarrow\shHom_{\BB}(\Psi_*(\M),\Psi_*(\N)).
 \tag{4.2.3.1}\label{eq:0:4.2.3.1}
\end{equation}

\egapar{4.2.4}
若 $\C$ 是 $\A$-代数层，则复合同态
\[
 \Psi_*(\C)\otimes_{\BB}\Psi_*(\C)
 \longrightarrow\Psi_*(\C\otimes_{\A}\C)
 \longrightarrow\Psi_*(\C)
\]
在 $\Psi_*(\C)$ 上定义了 $\BB$-代数层结构；这由（\hyperref[eq:0:4.2.2.2]{4.2.2.2}）得到。
同样，若 $\M$ 是 $\C$-模层，则 $\Psi_*(\M)$ 典范地具有 $\Psi_*(\C)$-模层结构。

\egapar{4.2.5}
特别考察 $X$ 是 $Y$ 的闭子空间、$\psi$ 是典范包含映射 $j\colon X\to Y$ 的情形。
设 $\BB'=\BB|X=j^*(\BB)$ 是环层 $\BB$ 在 $X$ 上的限制。
利用同态 $\theta^\sharp\colon\BB'\to\A$，可把 $\A$-模层 $\M$ 看作 $\BB'$-模层；
于是 $\Psi_*(\M)$ 是在 $X$ 上诱导 $\M$、在其余部分诱导零层的 $\BB$-模层。
若 $\N$ 是另一个 $\A$-模层，则
\[
 \Psi_*(\M)\otimes_{\BB}\Psi_*(\N)
 \simeq\Psi_*(\M\otimes_{\BB'}\N),
\]
并且
\[
 \shHom_{\BB}(\Psi_*(\M),\Psi_*(\N))
 \simeq\Psi_*(\shHom_{\BB'}(\M,\N)),
\]
这些等同都是典范的。

\egapar{4.2.6}
设 $(Z,\C)$ 是第三个赋环空间，
$\Psi'=(\psi',\theta')\colon(Y,\BB)\to(Z,\C)$ 是态射。
若 $\Psi''=\Psi'\circ\Psi$，则显然
\[
 \Psi''_*=\Psi'_*\circ\Psi_*.
\]

\separation
\subsection{\texorpdfstring{$\BB$}{B}-模层的逆像}\hypertarget{heading:0:4.3}{}
\src{40—42}

\egapar{4.3.1}
假设及记号仍如（\hyperlink{para:0:4.2.1}{4.2.1}）。设 $\G$ 是 $\BB$-模层，$\psi^*(\G)$ 为其逆像（\hyperlink{para:0:3.7.1}{3.7.1}），
它是 $X$ 上的阿贝尔群层。
由 $\psi^*(\G)$ 与 $\psi^*(\BB)$ 的截面定义（\hyperlink{para:0:3.7.1}{3.7.1}），
可见 $\psi^*(\G)$ 典范地具有 $\psi^*(\BB)$-模层结构。
另一方面，同态 $\theta^\sharp\colon\psi^*(\BB)\to\A$ 赋予 $\A$ 以
$\psi^*(\BB)$-模层结构；为避免混淆，必要时记此结构为 $\A_{[\theta]}$。
于是张量积
\[
 \psi^*(\G)\otimes_{\psi^*(\BB)}\A_{[\theta]}
\]
具有 $\A$-模层结构。
称这个 $\A$-模层为\emph{$\G$ 在态射 $\Psi$ 下的逆像}，记作 $\Psi^*(\G)$。

若 $\G_1,\G_2$ 是 $Y$ 上的两个 $\BB$-模层，$v\colon\G_1\to\G_2$ 为 $\BB$-同态，
则立即验证 $\psi^*(v)$ 是 $\psi^*(\G_1)\to\psi^*(\G_2)$ 的 $\psi^*(\BB)$-同态；
故 $\psi^*(v)\otimes1$ 是 $\Psi^*(\G_1)\to\Psi^*(\G_2)$ 的 $\A$-同态，
记为 $\Psi^*(v)$。
由此把 $\Psi^*$ 定义成从 $\BB$-模层范畴到 $\A$-模层范畴的协变函子。
与 $\psi^*$ 不同，这个函子一般不正合，而只是右正合；
这是因为在 $\psi^*(\BB)$-模层范畴中，与 $\A$ 作张量积是右正合函子。

依（\hyperlink{para:0:3.7.2}{3.7.2}），对每个 $x\in X$，有
\[
 (\Psi^*(\G))_x
 =\G_{\psi(x)}\otimes_{\BB_{\psi(x)}}\A_x.
\]
因此，$\Psi^*(\G)$ 的支集包含于 $\psi^{-1}(\Supp(\G))$。

\egapar{4.3.2}
设 $(\G_\lambda)$ 是 $\BB$-模层的归纳系，$\G=\ilim\G_\lambda$ 是其归纳极限。
典范同态 $\G_\lambda\to\G$ 定义了 $\psi^*(\BB)$-同态
$\psi^*(\G_\lambda)\to\psi^*(\G)$，从而得到典范同态
\[
 \ilim\psi^*(\G_\lambda)\longrightarrow\psi^*(\G).
\]
层的归纳极限在一点的茎等于各层在该点的茎的归纳极限（G，II，1.11），
故以上典范同态是双射（\hyperlink{para:0:3.7.2}{3.7.2}）。
又因张量积与层的归纳极限可交换，故有典范函子性 $\A$-模层同构
\[
 \ilim\Psi^*(\G_\lambda)
 \xrightarrow{\sim}\Psi^*(\ilim\G_\lambda).
\]
另一方面，对于 $\BB$-模层的有限直和 $\bigoplus_i\G_i$，显然
$\psi^*(\bigoplus_i\G_i)=\bigoplus_i\psi^*(\G_i)$。
与 $\A$ 作张量积，得到
\begin{equation}
 \Psi^*\!\left(\bigoplus_i\G_i\right)
 =\bigoplus_i\Psi^*(\G_i).
 \tag{4.3.2.1}\label{eq:0:4.3.2.1}
\end{equation}
通过取归纳极限，并利用以上结果，得到该等式对于任意直和也成立。

\egapar{4.3.3}
设 $\G_1,\G_2$ 是两个 $\BB$-模层。
由阿贝尔群层逆像的定义（\hyperlink{para:0:3.7.1}{3.7.1}），立即得到典范 $\psi^*(\BB)$-模层同态
\[
 \psi^*(\G_1)\otimes_{\psi^*(\BB)}\psi^*(\G_2)
 \longrightarrow\psi^*(\G_1\otimes_{\BB}\G_2).
\]
因为层的张量积在一点的茎就是各层在该点的茎的张量积（G，II，2.8），
由（\hyperlink{para:0:3.7.2}{3.7.2}）可知，以上同态实际上是同构。
与 $\A$ 作张量积，得到典范函子性同构
\begin{equation}
 \Psi^*(\G_1)\otimes_{\A}\Psi^*(\G_2)
 \xrightarrow{\sim}\Psi^*(\G_1\otimes_{\BB}\G_2).
 \tag{4.3.3.1}\label{eq:0:4.3.3.1}
\end{equation}

\egapar{4.3.4}
设 $\C$ 是 $\BB$-代数层。
给定 $\C$ 的代数层结构，等价于给定一个满足结合性与交换性条件的 $\BB$-同态
$\C\otimes_{\BB}\C\to\C$；这些条件可在各茎上验证。
利用以上同构，可把它搬移成满足相同条件的 $\A$-模层同态
\[
 \Psi^*(\C)\otimes_{\A}\Psi^*(\C)\longrightarrow\Psi^*(\C),
\]
于是 $\Psi^*(\C)$ 具有 $\A$-代数层结构。
特别地，由定义立即可见，$\A$-代数层 $\Psi^*(\BB)$ 等于 $\A$，这里的相等取到典范同构为止。

同样，若 $\M$ 是 $\C$-模层，则给定该模层结构，等价于给定满足结合性条件的
$\BB$-同态 $\C\otimes_{\BB}\M\to\M$；
通过搬移，得到 $\Psi^*(\M)$ 上的 $\Psi^*(\C)$-模层结构。

\egapar{4.3.5}
设 $\J$ 是 $\BB$ 的理想层。
因函子 $\psi^*$ 正合，$\psi^*(\BB)$-模层 $\psi^*(\J)$ 典范地等同于 $\psi^*(\BB)$ 的一个理想层。
典范包含同态 $\psi^*(\J)\to\psi^*(\BB)$ 给出 $\A$-模层同态
\[
 \Psi^*(\J)=\psi^*(\J)\otimes_{\psi^*(\BB)}\A_{[\theta]}
 \longrightarrow\A.
\]
把 $\Psi^*(\J)$ 在此同态下的像记为 $\Psi^*(\J)\A$；在不致混淆时，也记为 $\J\A$。
因此，依定义有
\[
 \J\A=\theta^\sharp(\psi^*(\J))\A.
\]
特别地，利用 $\psi^*(\J)$ 与 $\J$ 的茎的典范等同（\hyperlink{para:0:3.7.2}{3.7.2}），对每个 $x\in X$，有
\[
 (\J\A)_x=\theta_x(\J_{\psi(x)})\A_x.
\]
若 $\J_1,\J_2$ 是 $\BB$ 的两个理想层，则
\[
 (\J_1\J_2)\A=\J_1(\J_2\A)=(\J_1\A)(\J_2\A).
\]
若 $\F$ 是 $\A$-模层，则置 $\J\F=(\J\A)\F$。

\egapar{4.3.6}
设 $(Z,\C)$ 是第三个赋环空间，
$\Psi'=(\psi',\theta')\colon(Y,\BB)\to(Z,\C)$ 是态射，$\Psi''=\Psi'\circ\Psi$。
由定义（\hyperlink{para:0:4.3.1}{4.3.1}）及（\hyperref[eq:0:4.3.3.1]{4.3.3.1}），有
\[
 \Psi''{}^*=\Psi^*\circ\Psi'{}^*.
\]

\subsection{直像与逆像的关系}\hypertarget{heading:0:4.4}{}
\src{42—44}

\egapar{4.4.1}
假设及记号仍如（\hyperlink{para:0:4.2.1}{4.2.1}），设 $\G$ 为 $\BB$-模层。
按定义，$\BB$-模层同态 $u\colon\G\to\Psi_*(\F)$ 也称为从 $\G$ 到 $\F$ 的
\emph{$\Psi$-态射}，或简称从 $\G$ 到 $\F$ 的\emph{同态}；
在不会引起混淆时，写作 $u\colon\G\to\F$。
给定这样的同态，等价于对每对开集 $U\subseteq X,V\subseteq Y$，只要 $\psi(U)\subseteq V$，
给定一个 $\Gamma(V,\BB)$-模同态
\[
 u_{U,V}\colon\Gamma(V,\G)\longrightarrow\Gamma(U,\F),
\]
其中利用环同态 $\theta_{U,V}\colon\Gamma(V,\BB)\to\Gamma(U,\A)$，
把 $\Gamma(U,\F)$ 看作 $\Gamma(V,\BB)$-模。
各 $u_{U,V}$ 还须使图（\hyperref[eq:0:3.5.1.1]{3.5.1.1}）交换。
事实上，要定义 $u$，只需对 $X,Y$ 的拓扑基 $\B,\B'$ 中的 $U,V$ 给定这些 $u_{U,V}$，
并对于这些限制验证（\hyperref[eq:0:3.5.1.1]{3.5.1.1}）的交换性。

\egapar{4.4.2}
在（\hyperlink{para:0:4.2.1}{4.2.1}）及（\hyperlink{para:0:4.2.6}{4.2.6}）的假设下，设 $\HH$ 为 $\C$-模层，
$v\colon\HH\to\Psi'_*(\G)$ 为 $\Psi'$-态射。则
\[
 w\colon\HH\xrightarrow{v}\Psi'_*(\G)
 \xrightarrow{\Psi'_*(u)}\Psi'_*(\Psi_*(\F))
\]
是一个 $\Psi''$-态射，称为 $u$ 与 $v$ 的复合。

\egapar{4.4.3}
现在证明，可以定义关于 $\G,\F$ 的双函子的典范同构
\begin{equation}
 \Hom_{\A}(\Psi^*(\G),\F)
 \xrightarrow{\sim}\Hom_{\BB}(\G,\Psi_*(\F)).
 \tag{4.4.3.1}\label{eq:0:4.4.3.1}
\end{equation}
将它记为 $v\mapsto v^\flat_{\Psi}$，在不致混淆时简记为 $v\mapsto v^\flat$；
其逆同构记为 $u\mapsto u^\sharp_{\Psi}$ 或 $u\mapsto u^\sharp$。
定义如下：把 $v\colon\Psi^*(\G)\to\F$ 与典范映射 $\psi^*(\G)\to\Psi^*(\G)$ 复合，
得到群层同态 $v'\colon\psi^*(\G)\to\F$；它也是 $\psi^*(\BB)$-模层同态。
由（\hyperlink{para:0:3.7.1}{3.7.1}），得到同态
\[
 v'{}^\flat\colon\G\longrightarrow\psi_*(\F)=\Psi_*(\F).
\]
不难验证，它也是 $\BB$-模层同态；置 $v^\flat=v'{}^\flat$。
同样，从 $\BB$-模层同态 $u\colon\G\to\Psi_*(\F)$，由（\hyperlink{para:0:3.7.1}{3.7.1}）得到
$\psi^*(\BB)$-模层同态 $u'{}^\sharp\colon\psi^*(\G)\to\F$；
与 $\A$ 作张量积，得到 $\A$-模层同态 $\Psi^*(\G)\to\F$，记为 $u^\sharp$。
立即验证
\[
 (u^\sharp)^\flat=u,\qquad(v^\flat)^\sharp=v,
\]
以及同构 $v\mapsto v^\flat$ 关于 $\F$ 的函子性。
再像（\hyperlink{para:0:3.5.4}{3.5.4}）那样，形式地推得 $u\mapsto u^\sharp$ 关于 $\G$ 的函子性；
这一论证也能证明（\hyperlink{para:0:4.3.1}{4.3.1}）中已直接证明的 $\Psi^*$ 的函子性。

若 $v$ 取为 $\Psi^*(\G)$ 的恒等同态，则 $v^\flat$ 为同态
\begin{equation}
 \rho_{\G}\colon\G\longrightarrow\Psi_*(\Psi^*(\G));
 \tag{4.4.3.2}\label{eq:0:4.4.3.2}
\end{equation}
若 $u$ 取为 $\Psi_*(\F)$ 的恒等同态，则 $u^\sharp$ 为同态
\begin{equation}
 \sigma_{\F}\colon\Psi^*(\Psi_*(\F))\longrightarrow\F.
 \tag{4.4.3.3}\label{eq:0:4.4.3.3}
\end{equation}
这些同态称为\emph{典范同态}。它们一般既不是单射，也不是满射。
这里也有与（\hyperref[eq:0:3.5.3.3]{3.5.3.3}）及（\hyperref[eq:0:3.5.4.4]{3.5.4.4}）类似的典范分解。

注意，若 $s$ 是 $\G$ 在 $Y$ 的开集 $V$ 上的截面，则 $\rho_{\G}(s)$ 是
$\Psi^*(\G)$ 在 $\psi^{-1}(V)$ 上的截面 $s'\otimes1$，其中对每个
$x\in\psi^{-1}(V)$，有 $s'_x=s_{\psi(x)}$。
还要注意，若 $u\colon\G\to\Psi_*(\F)$ 为同态，则对每个 $x\in X$，它在茎上定义同态
\[
 u_x\colon\G_{\psi(x)}\longrightarrow\F_x.
\]
这个同态由典范同态 $\xi\mapsto\xi\otimes1$ 与 $(u^\sharp)_x$ 复合而成：
\[
 \G_{\psi(x)}\longrightarrow
 \G_{\psi(x)}\otimes_{\BB_{\psi(x)}}\A_x
 =(\Psi^*(\G))_x
 \xrightarrow{(u^\sharp)_x}\F_x.
\]
$u_x$ 也可由同态
\[
 \Gamma(V,\G)\longrightarrow\Gamma(\psi^{-1}(V),\F)
 \longrightarrow\F_x
\]
取归纳极限得到，其中 $V$ 遍历 $\psi(x)$ 的开邻域。

\egapar{4.4.4}
设 $\F_1,\F_2$ 为 $\A$-模层，$\G_1,\G_2$ 为 $\BB$-模层，
$u_i$ 是从 $\G_i$ 到 $\F_i$ 的同态（$i=1,2$）。
记 $u_1\otimes u_2$ 为满足
\[
 u^\sharp=(u_1)^\sharp\otimes(u_2)^\sharp
\]
的同态 $u\colon\G_1\otimes_{\BB}\G_2\to\F_1\otimes_{\A}\F_2$，
这里使用了（\hyperref[eq:0:4.3.3.1]{4.3.3.1}）的等同。
可以验证，$u$ 也是复合同态
\[
 \G_1\otimes_{\BB}\G_2
 \longrightarrow\Psi_*(\F_1)\otimes_{\BB}\Psi_*(\F_2)
 \longrightarrow\Psi_*(\F_1\otimes_{\A}\F_2),
\]
其中第一箭头是普通的张量积 $u_1\otimes u_2$，第二箭头是典范同态（\hyperref[eq:0:4.2.2.1]{4.2.2.1}）。

\egapar{4.4.5}
设 $(\G_\lambda)_{\lambda\in L}$ 是 $\BB$-模层的归纳系；对每个 $\lambda\in L$，
设 $u_\lambda\colon\G_\lambda\to\Psi_*(\F)$ 为同态，且它们组成一个归纳系。
置 $\G=\ilim\G_\lambda$，$u=\ilim u_\lambda$。
则各 $(u_\lambda)^\sharp$ 组成同态 $\Psi^*(\G_\lambda)\to\F$ 的归纳系，
其归纳极限恰为 $u^\sharp$。

\egapar{4.4.6}
设 $\M,\N$ 为两个 $\BB$-模层，$V$ 是 $Y$ 的开集，$U=\psi^{-1}(V)$。
映射 $v\mapsto\Psi^*(v)$ 是同态
\[
 \Hom_{\BB|V}(\M|V,\N|V)
 \longrightarrow
 \Hom_{\A|U}(\Psi^*(\M)|U,\Psi^*(\N)|U),
\]
这里按 $\Gamma(V,\BB)$-模结构理解：右端自然具有
$\Gamma(U,\psi^*(\BB))$-模结构，再借助典范同态（\hyperlink{para:0:3.7.2}{3.7.2}）
\[
 \Gamma(V,\BB)\longrightarrow\Gamma(U,\psi^*(\BB)),
\]
也成为 $\Gamma(V,\BB)$-模。
立即可见，这些同态与限制相容，因而定义典范函子性同态
\[
 \gamma\colon\shHom_{\BB}(\M,\N)
 \longrightarrow\Psi_*\bigl(\shHom_{\A}(\Psi^*(\M),\Psi^*(\N))\bigr).
\]
与之对应的同态为
\[
 \gamma^\sharp\colon\Psi^*(\shHom_{\BB}(\M,\N))
 \longrightarrow\shHom_{\A}(\Psi^*(\M),\Psi^*(\N)).
\]
这些典范同态关于 $\M,\N$ 均具有函子性。

\egapar{4.4.7}
设 $\F$ 是 $\A$-代数层，$\G$ 是 $\BB$-代数层。
若 $u\colon\G\to\Psi_*(\F)$ 是 $\BB$-代数层同态，
则 $u^\sharp\colon\Psi^*(\G)\to\F$ 是 $\A$-代数层同态。
这由图
\[
\begin{tikzcd}[column sep=large]
 \G\otimes_{\BB}\G\arrow[r]\arrow[d]&\G\arrow[d,"u"]\\
 \Psi_*(\F\otimes_{\A}\F)\arrow[r]&\Psi_*(\F)
\end{tikzcd}
\]
的交换性及（\hyperlink{para:0:4.4.4}{4.4.4}）得到。
同样，若 $v\colon\Psi^*(\G)\to\F$ 是 $\A$-代数层同态，
则 $v^\flat\colon\G\to\Psi_*(\F)$ 是 $\BB$-代数层同态。

\egapar{4.4.8}
设 $(Z,\C)$ 是第三个赋环空间，
$\Psi'=(\psi',\theta')\colon(Y,\BB)\to(Z,\C)$ 为态射，
$\Psi''\colon(X,\A)\to(Z,\C)$ 为复合态射 $\Psi'\circ\Psi$。
设 $\HH$ 为 $\C$-模层，$v'$ 是从 $\HH$ 到 $\G$ 的同态。
复合 $v''=v\circ v'$ 按定义是由
\[
 \HH\xrightarrow{v'}\Psi'_*(\G)
 \xrightarrow{\Psi'_*(v)}\Psi'_*(\Psi_*(\F))
\]
给出的从 $\HH$ 到 $\F$ 的同态。
可以验证，$v''{}^\sharp$ 就是同态
\[
 \Psi^*(\Psi'{}^*(\HH))
 \xrightarrow{\Psi^*(v'{}^\sharp)}\Psi^*(\G)
 \xrightarrow{v^\sharp}\F.
\]
\transnote{原版本条使用 $v''=v\circ v'$，而末式最后一箭头标为 $u^\sharp$。
本译稿统一记为 $v^\sharp$；这里的 $v$ 是 $\G\to\Psi_*(\F)$ 的同态，即前文同型的 $u$。}


% END chapters/ch0_s4.tex

\par

% BEGIN chapters/ch0_s5.tex
\section{拟凝聚层与凝聚层}\hypertarget{heading:0:5}{}
\subsection{拟凝聚层}\hypertarget{heading:0:5.1}{}
\src{44—45}

\egapar{5.1.1}
设 $(X,\OO_X)$ 是赋环空间，$\F$ 是 $\OO_X$-模层。
给定一个 $\OO_X$-模层同态 $u\colon\OO_X\to\F$，等价于给定截面
$s=u(1)\in\Gamma(X,\F)$。
事实上，给定 $s$ 后，对每个截面 $t\in\Gamma(U,\OO_X)$，必有
$u(t)=t\cdot(s|U)$；称 $u$ 是\emph{由截面 $s$ 定义的}。
现在设 $I$ 为任意指标集，考察直和层 $\OO_X^{(I)}$。
对每个 $i\in I$，令 $h_i$ 为第 $i$ 个因子到 $\OO_X^{(I)}$ 的典范包含同态。
已知 $u\mapsto(u\circ h_i)_{i\in I}$ 给出从
$\Hom_{\OO_X}(\OO_X^{(I)},\F)$ 到直积
$\bigl(\Hom_{\OO_X}(\OO_X,\F)\bigr)^I$ 的同构。
因此，$u\colon\OO_X^{(I)}\to\F$ 这些同态与 $\F$ 在 $X$ 上的截面族
$(s_i)_{i\in I}$ 之间存在典范的一一对应。
与 $(s_i)$ 对应的同态 $u$ 把
$(a_i)\in\Gamma(U,\OO_X)^{(I)}$ 映到 $\sum_{i\in I}a_i\cdot(s_i|U)$。

若由 $(s_i)$ 定义的同态 $\OO_X^{(I)}\to\F$ 是满射，则称 $\F$ \emph{由截面族 $(s_i)$ 生成}；
等价地，对每个 $x\in X$，$\OO_x$-模 $\F_x$ 由各 $(s_i)_x$ 生成。
若 $\F$ 由其在 $X$ 上的全部截面组成的族，或其中某个子族生成，
则称 $\F$ \emph{由其在 $X$ 上的截面生成}。
换言之，存在某个适当的 $I$，使得有满射同态 $\OO_X^{(I)}\to\F$。

注意，一个 $\OO_X$-模层 $\F$ 可能具有如下性质：存在 $x_0\in X$，
使得无论取 $x_0$ 的哪个开邻域 $U$，$\F|U$ 都不由其在 $U$ 上的截面生成。
只需取 $X=\mathbf R$，$\OO_X$ 为常值层 $\mathbf Z$，
$\F$ 为 $\OO_X$ 中满足 $\F_0=\{0\}$ 而 $\F_x=\mathbf Z$（$x\ne0$）的子模层，
并取 $x_0=0$：对于 $0$ 的一个邻域 $U$，$\F|U$ 在 $U$ 上的截面只有零截面。
\transnote{上述“只有零截面”的说法应在 $0$ 的连通开邻域（例如开区间）上理解。
若开邻域另有不含 $0$ 的连通分支，该分支上可以有非零截面；这不影响原文所述“不由截面生成”的反例。}

\egapar{5.1.2}
设 $f\colon X\to Y$ 为赋环空间态射。
若 $\OO_X$-模层 $\F$ 由其在 $X$ 上的截面生成，则典范同态
\[
 f^*(f_*(\F))\longrightarrow\F
\]
（\hyperref[eq:0:4.4.3.3]{4.4.3.3}）是满射。
事实上，沿用（\hyperlink{para:0:5.1.1}{5.1.1}）的记号，$s_i\otimes1$ 是 $f^*(f_*(\F))$ 在 $X$ 上的截面，
它在 $\F$ 中的像是 $s_i$。
在（\hyperlink{para:0:5.1.1}{5.1.1}）的例子中取 $f$ 为恒等态射，可见此命题的逆命题一般不成立。

\egapar{5.1.3}
称 $\OO_X$-模层 $\F$ 是\emph{拟凝聚的}，是指对每个 $x\in X$，
存在 $x$ 的一个开邻域 $U$，使 $\F|U$ 同构于某个形如
\[
 \OO_X^{(I)}|U\longrightarrow\OO_X^{(J)}|U
\]
的同态的余核，其中 $I,J$ 是任意指标集。
显然，$\OO_X$ 自身是拟凝聚 $\OO_X$-模层，
任意一族拟凝聚 $\OO_X$-模层的直和也是拟凝聚 $\OO_X$-模层。
称 $\OO_X$-代数层 $\A$ \emph{拟凝聚}，是指它作为 $\OO_X$-模层拟凝聚。

\egapar{5.1.4}
设 $f\colon X\to Y$ 为赋环空间态射。
若 $\G$ 是拟凝聚 $\OO_Y$-模层，则 $f^*(\G)$ 是拟凝聚 $\OO_X$-模层。
事实上，对每个 $x\in X$，在 $Y$ 中存在 $f(x)$ 的开邻域 $V$，
使 $\G|V$ 为同态 $\OO_Y^{(I)}|V\to\OO_Y^{(J)}|V$ 的余核。
若 $U=f^{-1}(V)$，$f_U$ 是 $f$ 在 $U$ 上的限制，则
\[
 f^*(\G)|U=f_U^*(\G|V).
\]
由于 $f_U^*$ 右正合且与直和可交换，$f_U^*(\G|V)$ 是同态
$\OO_X^{(I)}|U\to\OO_X^{(J)}|U$ 的余核。

\subsection{有限型层}\hypertarget{heading:0:5.2}{}
\src{45—47}

\egapar{5.2.1}
称 $\OO_X$-模层 $\F$ 是\emph{有限型的}，是指对每个 $x\in X$，
存在开邻域 $U$，使 $\F|U$ 由有限个在 $U$ 上的截面生成；
也就是说，$\F|U$ 同构于某个形如 $(\OO_X|U)^p$ 的层的商层，其中 $p$ 有限。
有限型层的每个商层仍是有限型层；有限个有限型层的直和、张量积也都是有限型层。
有限型 $\OO_X$-模层未必拟凝聚，例如 $\OO_X/\F$，其中 $\F$ 是（\hyperlink{para:0:5.1.1}{5.1.1}）中的例子。
若 $\F$ 有限型，则对每个 $x\in X$，$\F_x$ 都是有限型 $\OO_x$-模；
但（\hyperlink{para:0:5.1.1}{5.1.1}）的例子表明，这个必要条件一般并不充分。

\egapar{5.2.2}
设 $\F$ 是有限型 $\OO_X$-模层。
若 $s_i$（$1\leq i\leq n$）是 $\F$ 在点 $x\in X$ 的一个开邻域 $U$ 上的截面，
且各 $(s_i)_x$ 生成 $\F_x$，则存在开邻域 $V\subseteq U$，
使得对每个 $y\in V$，各 $(s_i)_y$ 都生成 $\F_y$（FAC，I，2，12，命题 1）。
特别地，由此得到 $\F$ 的支集是闭集。

同样，若同态 $u\colon\F\to\G$ 满足 $u_x=0$，
则存在 $x$ 的邻域 $U$，使得对每个 $y\in U$，都有 $u_y=0$。

\egapar{5.2.3}
假设 $X$ 拟紧，$\F,\G$ 是两个 $\OO_X$-模层，其中 $\G$ 有限型，
$u\colon\F\to\G$ 是满射同态。
再假设 $\F$ 是某个 $\OO_X$-模层归纳系 $(\F_\lambda)$ 的归纳极限。
则存在指标 $\mu$，使同态 $\F_\mu\to\G$ 是满射。

事实上，对每个 $x\in X$，存在 $x$ 的开邻域 $U(x)$ 及 $\G$ 在其上的有限个截面 $s_i$，
使得对每个 $y\in U(x)$，各 $(s_i)_y$ 都生成 $\G_y$。
于是存在开邻域 $V(x)\subseteq U(x)$ 及 $\F$ 在 $V(x)$ 上的有限个截面 $t_i$，
使每个 $i$ 都满足 $s_i|V(x)=u(t_i)$。
还可假设这些 $t_i$ 是同一个层 $\F_{\lambda(x)}$ 在 $V(x)$ 上的截面的典范像。
取有限个这样的邻域 $V(x_j)$ 覆盖 $X$，再取一个不小于所有 $\lambda(x_j)$ 的指标 $\mu$。
显然，这个指标满足要求。

仍假设 $X$ 拟紧。若有限型 $\OO_X$-模层 $\F$ 由其在 $X$ 上的截面生成（\hyperlink{para:0:5.1.1}{5.1.1}），
则 $\F$ 由这些截面的某个有限子族生成。
事实上，只需用有限个开邻域 $U_k$ 覆盖 $X$，使得对每个 $k$，
存在 $\F$ 在 $X$ 上的有限个截面 $s_{k,j}$，其在 $U_k$ 上的限制生成 $\F|U_k$。
显然，这些 $s_{k,j}$ 就生成 $\F$。

\egapar{5.2.4}
设 $f\colon X\to Y$ 为赋环空间态射。
若 $\G$ 是有限型 $\OO_Y$-模层，则 $f^*(\G)$ 是有限型 $\OO_X$-模层。
事实上，对每个 $x\in X$，存在 $f(x)$ 的开邻域 $V$ 以及满射同态
$v\colon\OO_Y^p|V\to\G|V$。
若 $U=f^{-1}(V)$，$f_U$ 为 $f$ 在 $U$ 上的限制，则
$f^*(\G)|U=f_U^*(\G|V)$。
因为 $f_U^*$ 右正合（\hyperlink{para:0:4.3.1}{4.3.1}）且与直和可交换（\hyperlink{para:0:4.3.2}{4.3.2}），
$f_U^*(v)$ 是满射同态 $\OO_X^p|U\to f^*(\G)|U$。

\egapar{5.2.5}
称 $\OO_X$-模层 $\F$ \emph{具有有限表示}，是指对每个 $x\in X$，
存在开邻域 $U$，使 $\F|U$ 同构于某个 $(\OO_X|U)$-同态
\[
 \OO_X^p|U\longrightarrow\OO_X^q|U
\]
的余核，其中 $p,q$ 是两个正整数。
这样的 $\OO_X$-模层因而是有限型且拟凝聚的。
若 $f\colon X\to Y$ 为赋环空间态射，$\G$ 为具有有限表示的 $\OO_Y$-模层，
则（\hyperlink{para:0:5.1.4}{5.1.4}）的论证表明，$f^*(\G)$ 也具有有限表示。

\egapar{5.2.6}
设 $\F$ 是具有有限表示的 $\OO_X$-模层（\hyperlink{para:0:5.2.5}{5.2.5}）。
则对每个 $\OO_X$-模层 $\HH$，典范函子性同态
\[
 \bigl(\shHom_{\OO_X}(\F,\HH)\bigr)_x
 \longrightarrow\Hom_{\OO_x}(\F_x,\HH_x)
\]
是双射（T，4.1.1）。

\egapar{5.2.7}
设 $\F,\G$ 是两个具有有限表示的 $\OO_X$-模层。
若对于某个 $x\in X$，$\F_x$ 与 $\G_x$ 是同构的 $\OO_x$-模，
则存在 $x$ 的开邻域 $U$，使 $\F|U$ 与 $\G|U$ 同构。
事实上，设 $\varphi\colon\F_x\to\G_x$ 与 $\psi\colon\G_x\to\F_x$ 为互逆同构。
依（\hyperlink{para:0:5.2.6}{5.2.6}），存在开邻域 $V$，以及
$\shHom_{\OO_X}(\F,\G)$、$\shHom_{\OO_X}(\G,\F)$ 在 $V$ 上的截面 $u,v$，
分别满足 $u_x=\varphi,v_x=\psi$。
由于 $(u\circ v)_x$ 与 $(v\circ u)_x$ 是恒等自同构，
存在开邻域 $U\subseteq V$，使 $(u\circ v)|U$ 与 $(v\circ u)|U$ 都是恒等自同构。
命题得证。

\subsection{凝聚层}\hypertarget{heading:0:5.3}{}
\src{47—48}

\egapar{5.3.1}
称 $\OO_X$-模层 $\F$ 是\emph{凝聚的}，是指它满足以下两个条件：
\begin{enumerate}[label=\alph*)]
\item $\F$ 是有限型的。
\item 对每个开集 $U\subseteq X$、每个正整数 $n$，以及每个同态
$u\colon\OO_X^n|U\to\F|U$，$u$ 的核都是有限型的。
\end{enumerate}
注意，这两个条件都是局部性的。
以下回顾的凝聚层性质，其大部分证明可参阅（FAC），I，2。

\egapar{5.3.2}
每个凝聚 $\OO_X$-模层都具有有限表示（\hyperlink{para:0:5.2.5}{5.2.5}）；
逆命题未必成立，因为 $\OO_X$ 本身未必是凝聚 $\OO_X$-模层。

凝聚 $\OO_X$-模层的每个有限型子 $\OO_X$-模层都是凝聚的；
有限个凝聚 $\OO_X$-模层的直和也是凝聚 $\OO_X$-模层。

\egapar{5.3.3}
若 $0\to\F\to\G\to\HH\to0$ 是 $\OO_X$-模层的正合列，
且其中两个模层凝聚，则第三个也凝聚。

\egapar{5.3.4}
若 $\F,\G$ 是两个凝聚 $\OO_X$-模层，$u\colon\F\to\G$ 为同态，
则 $\Im(u),\Ker(u),\Coker(u)$ 都是凝聚 $\OO_X$-模层。
特别地，若 $\F,\G$ 是某个凝聚 $\OO_X$-模层的凝聚子模层，
则 $\F+\G$ 与 $\F\cap\G$ 都凝聚。

\egapar{5.3.5}
若 $\F,\G$ 是两个凝聚 $\OO_X$-模层，则
$\F\otimes_{\OO_X}\G$ 与 $\shHom_{\OO_X}(\F,\G)$ 也凝聚。

\egapar{5.3.6}
设 $\F$ 是凝聚 $\OO_X$-模层，$\J$ 是 $\OO_X$ 的凝聚理想层。
则 $\OO_X$-模层 $\J\F$ 凝聚，因为它是典范同态
$\J\otimes_{\OO_X}\F\to\F$ 的像（\hyperlink{para:0:5.3.4}{5.3.4} 与 \hyperlink{para:0:5.3.5}{5.3.5}）。

\egapar{5.3.7}
称 $\OO_X$-代数层 $\A$ \emph{凝聚}，是指 $\A$ 作为 $\OO_X$-模层凝聚。
特别地，$\OO_X$ 是\emph{凝聚环层}，当且仅当对每个开集 $U\subseteq X$，
以及每个形如 $u\colon\OO_X^n|U\to\OO_X|U$ 的同态，
其核都是有限型 $(\OO_X|U)$-模层。

若 $\OO_X$ 是凝聚环层，则由（\hyperlink{para:0:5.3.4}{5.3.4}），每个具有有限表示的 $\OO_X$-模层 $\F$ 都凝聚。

$\OO_X$-模层 $\F$ 的\emph{零化理想层}是典范同态
\[
 \OO_X\longrightarrow\shHom_{\OO_X}(\F,\F)
\]
的核 $\J$；该同态把每个截面 $s\in\Gamma(U,\OO_X)$ 映到
$\Hom_{\OO_X|U}(\F|U,\F|U)$ 中由 $s$ 给出的乘法同态。
若 $\OO_X$ 与 $\F$ 都凝聚，则 $\J$ 凝聚（\hyperlink{para:0:5.3.4}{5.3.4} 与 \hyperlink{para:0:5.3.5}{5.3.5}），
且对每个 $x\in X$，$\J_x$ 是 $\F_x$ 的零化理想（\hyperlink{para:0:5.2.6}{5.2.6}）。

\egapar{5.3.8}
假设 $\OO_X$ 凝聚。设 $\F$ 是凝聚 $\OO_X$-模层，$x\in X$，
$M$ 是 $\F_x$ 的一个有限型子模。
则存在 $x$ 的开邻域 $U$，以及 $\F|U$ 的一个凝聚子 $(\OO_X|U)$-模层 $\G$，
使 $\G_x=M$（T，4.1，引理 1）。

这一结果，连同凝聚 $\OO_X$-模层的子模层的性质，给出了 $\OO_X$ 凝聚时各环 $\OO_x$ 必须满足的条件。
例如，依（\hyperlink{para:0:5.3.4}{5.3.4}），$\OO_x$ 中两个有限型理想的交仍须是有限型理想。

\egapar{5.3.9}
假设 $\OO_X$ 凝聚。设 $M$ 为具有有限表示的 $\OO_x$-模，
因此它同构于某个同态 $\varphi\colon\OO_x^p\to\OO_x^q$ 的余核。
则存在 $x$ 的开邻域 $U$，以及一个凝聚 $(\OO_X|U)$-模层 $\F$，使 $\F_x$ 同构于 $M$。
事实上，依（\hyperlink{para:0:5.2.6}{5.2.6}），存在 $\shHom_{\OO_X}(\OO_X^p,\OO_X^q)|U$ 的截面 $u$，
使 $u_x=\varphi$；同态 $u\colon\OO_X^p|U\to\OO_X^q|U$ 的余核 $\F$ 就满足要求（\hyperlink{para:0:5.3.4}{5.3.4}）。
\transnote{原版先写 $\varphi\colon\OO_x^p\to\OO_x^q$，而末句写为
$u\colon\OO_X^q|U\to\OO_X^p|U$；译文使 $u$ 的源、靶与前面的 $\varphi$ 一致。}

\egapar{5.3.10}
假设 $\OO_X$ 凝聚，$\J$ 是 $\OO_X$ 的凝聚理想层。
一个 $(\OO_X/\J)$-模层 $\F$ 凝聚，当且仅当它作为 $\OO_X$-模层凝聚。
特别地，$\OO_X/\J$ 是凝聚环层。

\egapar{5.3.11}
设 $f\colon X\to Y$ 为赋环空间态射，并假设 $\OO_X$ 凝聚。
则对每个凝聚 $\OO_Y$-模层 $\G$，$f^*(\G)$ 都是凝聚 $\OO_X$-模层。
事实上，沿用（\hyperlink{para:0:5.2.4}{5.2.4}）的记号，可假设 $\G|V$ 是同态
$v\colon\OO_Y^q|V\to\OO_Y^p|V$ 的余核。
由于 $f_U^*$ 右正合，$f^*(\G)|U=f_U^*(\G|V)$ 是同态
$f_U^*(v)\colon\OO_X^q|U\to\OO_X^p|U$ 的余核，由此得到结论。

\egapar{5.3.12}
设 $Y$ 是 $X$ 的闭子集，$j\colon Y\to X$ 为典范包含映射，
$\OO_Y$ 是 $Y$ 上的环层，并置 $\OO_X=j_*(\OO_Y)$。
一个 $\OO_Y$-模层 $\G$ 是有限型的（分别地，拟凝聚的、凝聚的），
当且仅当 $j_*(\G)$ 是有限型 $\OO_X$-模层（分别地，拟凝聚 $\OO_X$-模层、凝聚 $\OO_X$-模层）。

\subsection{局部自由层}\hypertarget{heading:0:5.4}{}
\src{48—53}

\egapar{5.4.1}
设 $X$ 是赋环空间。称 $\OO_X$-模层 $\F$ \emph{局部自由}，是指对每个 $x\in X$，
存在 $x$ 的开邻域 $U$，使 $\F|U$ 同构于形如 $\OO_X^{(I)}|U$ 的 $(\OO_X|U)$-模层，
其中 $I$ 可以依赖于 $U$。
若在这些邻域上都可取 $I$ 为有限集，则称 $\F$ 具有\emph{有限秩}；
若 $I$ 的元素个数在这些邻域上都等于同一个有限数 $n$，则称 $\F$ 的\emph{秩为 $n$}。
秩为 $1$ 的局部自由 $\OO_X$-模层也称为\emph{可逆层}（参见（\hyperlink{para:0:5.4.3}{5.4.3}））。
若 $\F$ 是有限秩局部自由 $\OO_X$-模层，则对每个 $x\in X$，
$\F_x$ 是有限秩 $n(x)$ 的自由 $\OO_x$-模，且存在 $x$ 的邻域 $U$，
使 $\F|U$ 的秩为 $n(x)$；因此，若 $X$ 连通，$n(x)$ 就是常数。

显然，每个局部自由层都是拟凝聚的；若 $\OO_X$ 是凝聚环层，
则每个有限秩局部自由 $\OO_X$-模层都凝聚。

若 $\LL$ 局部自由，则 $\LL\otimes_{\OO_X}\F$ 是 $\OO_X$-模层范畴中关于 $\F$ 的正合函子。

我们主要考察有限秩局部自由 $\OO_X$-模层；以后未加说明地谈到局部自由层时，均默认其秩有限。

\egapar{5.4.2}
若 $\LL,\F$ 是两个 $\OO_X$-模层，则有典范函子性同态
\begin{equation}
 \dual{\LL}\otimes_{\OO_X}\F
 =\shHom_{\OO_X}(\LL,\OO_X)\otimes_{\OO_X}\F
 \longrightarrow\shHom_{\OO_X}(\LL,\F).
 \tag{5.4.2.1}\label{eq:0:5.4.2.1}
\end{equation}
它的定义如下。
对每个开集 $U$，以及每对 $(u,t)$，其中
\[
 u\in\Gamma(U,\shHom_{\OO_X}(\LL,\OO_X))
 =\Hom_{\OO_X|U}(\LL|U,\OO_X|U),
 \qquad t\in\Gamma(U,\F),
\]
令其对应于 $\Hom_{\OO_X|U}(\LL|U,\F|U)$ 中的同态；
这个同态在每个 $x\in U$ 把 $s_x\in\LL_x$ 映到 $u_x(s_x)t_x\in\F_x$。
若 $\LL$ 局部自由且秩有限，则此同态是双射。
因为这是局部性质，只需考虑 $\LL=\OO_X^n$ 的情形。
又因对每个 $\OO_X$-模层 $\F$，$\shHom_{\OO_X}(\OO_X^n,\F)$ 都典范同构于 $\F^n$，
所以进一步归结到 $\LL=\OO_X$ 的情形，而此时结论显然成立。

\egapar{5.4.3}
若 $\LL$ 可逆，则其对偶 $\dual{\LL}=\shHom_{\OO_X}(\LL,\OO_X)$ 也可逆；
因为问题是局部的，立即归结到 $\LL=\OO_X$ 的情形。
此外，有典范同构
\begin{equation}
 \shHom_{\OO_X}(\LL,\OO_X)\otimes_{\OO_X}\LL
 \xrightarrow{\sim}\OO_X.
 \tag{5.4.3.1}\label{eq:0:5.4.3.1}
\end{equation}
事实上，依（\hyperlink{para:0:5.4.2}{5.4.2}），只需定义典范同构
$\shHom_{\OO_X}(\LL,\LL)\simeq\OO_X$。
而对每个 $\OO_X$-模层 $\F$，已有典范同态
$\OO_X\to\shHom_{\OO_X}(\F,\F)$（\hyperlink{para:0:5.3.7}{5.3.7}）。
只须证明当 $\F=\LL$ 可逆时此同态为双射；这是局部问题，仍归结到 $\LL=\OO_X$ 的显然情形。

由此，置 $\LL^{-1}=\shHom_{\OO_X}(\LL,\OO_X)$，并称 $\LL^{-1}$ 为 $\LL$ 的\emph{逆层}。
当 $X$ 只有一个点、$\OO_X$ 为极大理想是 $\mathfrak m$ 的局部环 $A$ 时，
“可逆层”这一术语可作如下解释。
设 $M,M'$ 是两个 $A$-模，其中 $M$ 有限型，并且 $M\otimes_A M'$ 同构于 $A$。
因为
\[
 (A/\mathfrak m)\otimes_A(M\otimes_A M')
 \simeq(M/\mathfrak mM)\otimes_{A/\mathfrak m}(M'/\mathfrak mM'),
\]
右端这两个 $A/\mathfrak m$-向量空间的张量积同构于 $A/\mathfrak m$，
故 $M/\mathfrak mM$ 与 $M'/\mathfrak mM'$ 的维数都必须为 $1$。
于是，对每个不属于 $\mathfrak mM$ 的元素 $\xi\in M$，有
$M=A\xi+\mathfrak mM$；因 $M$ 有限型，由 Nakayama 引理，$M=A\xi$。
再者，$\xi$ 的零化理想零化 $M\otimes_A M'\simeq A$，故该理想为 $\{0\}$。
所以 $M$ 同构于 $A$。

在一般情形，这表明：若 $\LL$ 是有限型 $\OO_X$-模层，
存在 $\OO_X$-模层 $\F$ 使 $\LL\otimes_{\OO_X}\F\simeq\OO_X$，
并且各环 $\OO_x$ 都是局部环，则每个 $\LL_x$ 都是同构于 $\OO_x$ 的 $\OO_x$-模。
若还假设 $\OO_X$ 与 $\LL$ 凝聚，则由（\hyperlink{para:0:5.2.7}{5.2.7}）可得 $\LL$ 可逆。

\egapar{5.4.4}
若 $\LL,\LL'$ 是两个可逆 $\OO_X$-模层，则 $\LL\otimes_{\OO_X}\LL'$ 也可逆。
因为这是局部问题，可假设 $\LL=\OO_X$，此时结果显然。
对每个整数 $n\geq1$，以 $\LL^{\otimes n}$ 表示 $n$ 个 $\LL$ 的张量积。
约定 $\LL^{\otimes0}=\OO_X$；对 $n\geq1$，置
$\LL^{\otimes(-n)}=(\LL^{-1})^{\otimes n}$。
采用这些记号，则对任意整数 $m,n$，存在典范函子性同构
\begin{equation}
 \LL^{\otimes m}\otimes_{\OO_X}\LL^{\otimes n}
 \xrightarrow{\sim}\LL^{\otimes(m+n)}.
 \tag{5.4.4.1}\label{eq:0:5.4.4.1}
\end{equation}
事实上，依定义，立即归结到 $m=-1,n=1$ 的情形，而此时该同构已在（\hyperlink{para:0:5.4.3}{5.4.3}）中定义。

\egapar{5.4.5}
设 $f\colon Y\to X$ 为赋环空间态射。
若 $\LL$ 为局部自由 $\OO_X$-模层（分别地，可逆 $\OO_X$-模层），
则 $f^*(\LL)$ 为局部自由 $\OO_Y$-模层（分别地，可逆 $\OO_Y$-模层）。
这是以下事实的直接结果：局部同构的两个 $\OO_X$-模层的逆像仍局部同构；
$f^*$ 与有限直和可交换；并且 $f^*(\OO_X)=\OO_Y$（\hyperlink{para:0:4.3.4}{4.3.4}）。
此外，已知存在典范函子性同态
\[
 f^*(\dual{\LL})\longrightarrow\dual{f^*(\LL)}
\]
（\hyperlink{para:0:4.4.6}{4.4.6}）。若 $\LL$ 局部自由，则此同态为双射；
事实上，仍可归结到 $\LL=\OO_X$ 的显然情形。
由此得到，若 $\LL$ 可逆，则对每个整数 $n$，
$f^*(\LL^{\otimes n})$ 典范地等同于 $(f^*(\LL))^{\otimes n}$。

\egapar{5.4.6}
设 $\LL$ 为可逆 $\OO_X$-模层。
以 $\Gamma_*(X,\LL)$，或简记 $\Gamma_*(\LL)$，表示阿贝尔群直和
\[
 \bigoplus_{n\in\mathbf Z}\Gamma(X,\LL^{\otimes n}).
\]
赋予它一个分次环结构如下：
若 $s_1\in\Gamma(X,\LL^{\otimes m})$、$s_2\in\Gamma(X,\LL^{\otimes n})$，
令 $(s_1,s_2)$ 对应于 $\LL^{\otimes(m+n)}$ 在 $X$ 上的截面，
即在典范同构（\hyperref[eq:0:5.4.4.1]{5.4.4.1}）下与 $\LL^{\otimes m}\otimes_{\OO_X}\LL^{\otimes n}$
的截面 $s_1\otimes s_2$ 对应的截面。
立即可验证此乘法的结合性。
显然，$\Gamma_*(X,\LL)$ 是关于 $\LL$ 的协变函子，取值于分次环范畴。
\transnote{本条的分次包括所有整数次数。涉及负张量幂时，对一般同态
$\LL\to\LL'$，对偶给出的箭头方向相反；因此，关于 $\LL$ 的上述协变函子性
须在以同构为态射的可逆层范畴中理解。本条后文关于 $\LL$ 的同类表述亦同。
关于 $\F$ 的函子性以及赋环空间态射下的逆像构造不受此问题影响。}

若 $\F$ 是任意 $\OO_X$-模层，置
\[
 \Gamma_*(\LL,\F)
 =\bigoplus_{n\in\mathbf Z}
   \Gamma(X,\F\otimes_{\OO_X}\LL^{\otimes n}).
\]
把这个阿贝尔群赋成分次环 $\Gamma_*(\LL)$ 上的分次模如下：
对于 $s\in\Gamma(X,\LL^{\otimes m})$、
$t\in\Gamma(X,\F\otimes_{\OO_X}\LL^{\otimes n})$，
令 $(s,t)$ 对应于 $\F\otimes_{\OO_X}\LL^{\otimes(m+n)}$ 的截面，
即利用（\hyperref[eq:0:5.4.4.1]{5.4.4.1}）与 $s\otimes t$ 典范对应的截面。
模公理的验证是直接的。
固定 $X,\LL$ 时，$\Gamma_*(\LL,\F)$ 是关于 $\F$ 的协变函子，
取值于分次 $\Gamma_*(\LL)$-模范畴；固定 $X,\F$ 时，它是关于 $\LL$ 的协变函子，取值于阿贝尔群范畴。

若 $f\colon Y\to X$ 为赋环空间态射，则典范同态（\hyperref[eq:0:4.4.3.2]{4.4.3.2}）
\[
 \rho\colon\LL^{\otimes n}\longrightarrow f_*(f^*(\LL^{\otimes n}))
\]
定义阿贝尔群同态
\[
 \Gamma(X,\LL^{\otimes n})\longrightarrow\Gamma(Y,f^*(\LL^{\otimes n})).
\]
因为 $f^*(\LL^{\otimes n})=(f^*(\LL))^{\otimes n}$，
由典范同态（\hyperref[eq:0:4.4.3.2]{4.4.3.2}）与（\hyperref[eq:0:5.4.4.1]{5.4.4.1}）的定义可见，以上同态定义了函子性分次环同态
\[
 \Gamma_*(\LL)\longrightarrow\Gamma_*(f^*(\LL)).
\]
同一个典范同态（\hyperlink{para:0:4.4.3}{4.4.3}）还定义阿贝尔群同态
\[
 \Gamma(X,\F\otimes_{\OO_X}\LL^{\otimes n})
 \longrightarrow\Gamma(Y,f^*(\F\otimes_{\OO_X}\LL^{\otimes n})).
\]
又因
\[
 f^*(\F\otimes_{\OO_X}\LL^{\otimes n})
 =f^*(\F)\otimes_{\OO_Y}(f^*(\LL))^{\otimes n}
 \qquad\text{（\hyperref[eq:0:4.3.3.1]{4.3.3.1}），}
\]
让 $n$ 变化，这些同态定义了分次模的二重同态
\[
 \Gamma_*(\LL,\F)\longrightarrow\Gamma_*(f^*(\LL),f^*(\F)).
\]

\egapar{5.4.7}
可以证明，存在一个由可逆 $\OO_X$-模层组成的\emph{集合} $\mathfrak M$，
也记为 $\mathfrak M(X)$，使每个可逆 $\OO_X$-模层恰好与 $\mathfrak M$ 中的一个元素同构。
\origfootnote{参见导言中提到的那部正在准备的著作。}
在 $\mathfrak M$ 中定义合成法则：令 $\mathfrak M$ 中两个元素 $\LL,\LL'$ 的乘积为
$\mathfrak M$ 中唯一与 $\LL\otimes\LL'$ 同构的元素。
在此合成法则下，$\mathfrak M$ 是一个群，且同构于上同调群
$H^1(X,\OO_X^*)$。
这里 $\OO_X^*$ 是 $\OO_X$ 的子层，满足：对每个开集 $U\subseteq X$，
$\Gamma(U,\OO_X^*)$ 是环 $\Gamma(U,\OO_X)$ 的可逆元群；
因此 $\OO_X^*$ 是乘法记法的阿贝尔群层。

为证明这一点，注意对于每个开集 $U\subseteq X$，
截面群 $\Gamma(U,\OO_X^*)$ 典范地等同于 $(\OO_X|U)$-模层 $\OO_X|U$ 的自同构群。
这一等同把 $\OO_X^*$ 在 $U$ 上的截面 $\varepsilon$ 对应到 $\OO_X|U$ 的自同构 $u$，
它在每个 $x\in U$ 满足 $u_x(s_x)=\varepsilon_xs_x$。
设 $\mathfrak U=(U_\lambda)$ 是 $X$ 的开覆盖。
对每对指标 $(\lambda,\mu)$ 给定 $\OO_X|(U_\lambda\cap U_\mu)$ 的自同构 $\theta_{\lambda\mu}$，
等价于给定覆盖 $\mathfrak U$ 的一个取值于 $\OO_X^*$ 的 $1$-余链。
各 $\theta_{\lambda\mu}$ 满足粘合条件（\hyperlink{para:0:3.3.1}{3.3.1}），等价于相应余链是余闭链。
同样，对每个 $\lambda$ 给定 $\OO_X|U_\lambda$ 的自同构 $\omega_\lambda$，
等价于给定覆盖 $\mathfrak U$ 的一个取值于 $\OO_X^*$ 的 $0$-余链；
它的余边界对应于自同构族
\[
 (\omega_\lambda|(U_\lambda\cap U_\mu))
 \circ(\omega_\mu|(U_\lambda\cap U_\mu))^{-1}.
\]

对于 $\mathfrak U$ 的每个取值于 $\OO_X^*$ 的 $1$-余闭链，
利用与之对应的自同构族 $(\theta_{\lambda\mu})$ 粘合，得到一个可逆 $\OO_X$-模层；
令该余闭链对应到 $\mathfrak M$ 中与所得模层同构的元素。
上同调等价的两个余闭链，对应于 $\mathfrak M$ 中相同的元素（\hyperlink{para:0:3.3.2}{3.3.2}）。
也就是说，我们定义了映射
$\varphi_{\mathfrak U}\colon H^1(\mathfrak U,\OO_X^*)\to\mathfrak M$。
若 $\mathfrak V$ 是 $X$ 的另一个开覆盖，且比 $\mathfrak U$ 更细，则图
\[
\begin{tikzcd}[column sep=large,row sep=large]
 H^1(\mathfrak U,\OO_X^*)\arrow[d]\arrow[dr,"\varphi_{\mathfrak U}"]&\\
 H^1(\mathfrak V,\OO_X^*)\arrow[r,"\varphi_{\mathfrak V}"']&\mathfrak M
\end{tikzcd}
\]
交换，其中竖直箭头为典范同态（G，II，5.7）；这由（\hyperlink{para:0:3.3.3}{3.3.3}）得到。
取归纳极限，便得到映射
\[
 H^1(X,\OO_X^*)\longrightarrow\mathfrak M,
\]
这里用到 Čech 上同调群 $\check H^1(X,\OO_X^*)$ 与第一上同调群
$H^1(X,\OO_X^*)$ 的已知等同（G，II，5.9，定理 5.9.1 的推论）。
此映射为满射：依定义，对于每个可逆 $\OO_X$-模层 $\LL$，
存在 $X$ 的开覆盖 $(U_\lambda)$，使 $\LL$ 由各 $\OO_X|U_\lambda$ 粘合而成（\hyperlink{para:0:3.3.1}{3.3.1}）。
此映射也为单射：只须对映射 $H^1(\mathfrak U,\OO_X^*)\to\mathfrak M$ 证明这一点，
而它由（\hyperlink{para:0:3.3.2}{3.3.2}）得到。

剩下只需证明所得双射是群同态。
给定两个可逆 $\OO_X$-模层 $\LL,\LL'$，可以取开覆盖 $(U_\lambda)$，
使 $\LL|U_\lambda$ 与 $\LL'|U_\lambda$ 对每个 $\lambda$ 都同构于 $\OO_X|U_\lambda$。
于是，对每个指标 $\lambda$，存在
$a_\lambda\in\Gamma(U_\lambda,\LL)$ 与 $a'_\lambda\in\Gamma(U_\lambda,\LL')$，
使这两个截面模的元素分别为 $s_\lambda a_\lambda$ 与 $s'_\lambda a'_\lambda$，
其中 $s_\lambda,s'_\lambda$ 遍历 $\Gamma(U_\lambda,\OO_X)$。
相应余闭链 $(\varepsilon_{\lambda\mu})$ 与 $(\varepsilon'_{\lambda\mu})$ 满足：在 $U_\lambda\cap U_\mu$ 上，
\[
 s_\lambda a_\lambda=s_\mu a_\mu
 \quad\Longleftrightarrow\quad
 s_\mu=\varepsilon_{\lambda\mu}s_\lambda,
\]
相应地，
\[
 s'_\lambda a'_\lambda=s'_\mu a'_\mu
 \quad\Longleftrightarrow\quad
 s'_\mu=\varepsilon'_{\lambda\mu}s'_\lambda.
\]
由于 $\LL\otimes_{\OO_X}\LL'$ 在 $U_\lambda$ 上的截面都是
$s_\lambda s'_\lambda(a_\lambda\otimes a'_\lambda)$ 的有限和，
显然，余闭链 $(\varepsilon_{\lambda\mu}\varepsilon'_{\lambda\mu})$
对应于 $\LL\otimes_{\OO_X}\LL'$。
证明完成。
\origfootnote{关于本结果的一般形式，参见原书第 51 页脚注所引著作。}

\egapar{5.4.8}
设 $f=(\psi,\omega)\colon Y\to X$ 为赋环空间态射。
可逆 $\OO_X$-模层范畴中的函子 $f^*(\LL)$ 定义了从集合 $\mathfrak M(X)$
到集合 $\mathfrak M(Y)$ 的映射，略去记号上的区别，也把它记为 $f^*$。
\transnote{原版此处印作“自由 $\OO_X$-模层的范畴”（\textit{Modules libres}）；
按前后文的 $\mathfrak M(X)$ 定义，此处所指是可逆模层，译文据此注明其对象范围。}
另一方面，有典范同态（T，3.2.2）
\begin{equation}
 H^1(X,\OO_X^*)\longrightarrow H^1(Y,\OO_Y^*).
 \tag{5.4.8.1}\label{eq:0:5.4.8.1}
\end{equation}
在（\hyperlink{para:0:5.4.7}{5.4.7}）的典范等同
$\mathfrak M(X)=H^1(X,\OO_X^*)$、$\mathfrak M(Y)=H^1(Y,\OO_Y^*)$ 下，
同态（\hyperref[eq:0:5.4.8.1]{5.4.8.1}）与映射 $f^*$ 等同。

事实上，若 $\LL$ 来自与 $X$ 的开覆盖 $(U_\lambda)$ 对应的余闭链
$(\varepsilon_{\lambda\mu})$，则只需证明 $f^*(\LL)$ 来自一个余闭链，
其上同调类是 $(\varepsilon_{\lambda\mu})$ 的上同调类在（\hyperref[eq:0:5.4.8.1]{5.4.8.1}）下的像。
设 $\theta_{\lambda\mu}$ 是与 $\varepsilon_{\lambda\mu}$ 对应的
$\OO_X|(U_\lambda\cap U_\mu)$ 的自同构。
显然，$f^*(\LL)$ 是利用自同构 $f^*(\theta_{\lambda\mu})$，
粘合各 $\OO_Y|\psi^{-1}(U_\lambda)$ 所得的层。
所以只需验证这些自同构对应于余闭链 $(\omega^\sharp(\varepsilon_{\lambda\mu}))$，
而这由定义立即得到。
这里把 $\varepsilon_{\lambda\mu}$ 与其在 $\rho$ 下的典范像等同（\hyperlink{para:0:3.7.2}{3.7.2}），
即把它看作 $\psi^*(\OO_X)$ 在 $\psi^{-1}(U_\lambda\cap U_\mu)$ 上的截面。

\egapar{5.4.9}
设 $\E,\F$ 是两个 $\OO_X$-模层，且 $\F$ 局部自由。
设 $\G$ 是 $\F$ 以 $\E$ 为核的一个扩张，即存在正合列
\[
 0\longrightarrow\E\longrightarrow\G\xrightarrow{p}\F\longrightarrow0.
\]
则对每个 $x\in X$，存在 $x$ 的开邻域 $U$，
使 $\G|U$ 同构于直和 $\E|U\oplus\F|U$。
事实上，只需考虑 $\F=\OO_X^n$ 的情形。
设 $e_i$（$1\leq i\leq n$）是 $\OO_X^n$ 的典范截面（\hyperlink{para:0:5.5.5}{5.5.5}）。
存在 $x$ 的开邻域 $U$，以及 $\G$ 在 $U$ 上的 $n$ 个截面 $s_i$，
使 $p(s_i)=e_i|U$（$1\leq i\leq n$）。
令 $f\colon\F|U\to\G|U$ 为这些截面所定义的同态（\hyperlink{para:0:5.1.1}{5.1.1}）。
立即可见，对于每个开集 $V\subseteq U$ 及每个截面 $s\in\Gamma(V,\G)$，
都有 $s-f(p(s))\in\Gamma(V,\E)$，从而得到结论。

\egapar{5.4.10}
设 $f\colon X\to Y$ 为赋环空间态射，$\F$ 是 $\OO_X$-模层，
$\LL$ 是有限秩局部自由 $\OO_Y$-模层。则存在典范同构
\begin{equation}
 f_*(\F)\otimes_{\OO_Y}\LL
 \xrightarrow{\sim}f_*(\F\otimes_{\OO_X}f^*(\LL)).
 \tag{5.4.10.1}\label{eq:0:5.4.10.1}
\end{equation}
事实上，对于任意 $\OO_Y$-模层 $\LL$，都有典范同态
\[
 \begin{aligned}
 f_*(\F)\otimes_{\OO_Y}\LL
 &\xrightarrow{1\otimes\rho}
 f_*(\F)\otimes_{\OO_Y}f_*(f^*(\LL))\\
 &\xrightarrow{\alpha}f_*(\F\otimes_{\OO_X}f^*(\LL)),
 \end{aligned}
\]
其中 $\rho$ 是同态（\hyperref[eq:0:4.4.3.2]{4.4.3.2}），$\alpha$ 是同态（\hyperref[eq:0:4.2.2.1]{4.2.2.1}）。
要证明当 $\LL$ 局部自由时此同态为双射，只需注意问题在 $Y$ 上是局部的，
从而可设 $\LL=\OO_Y^n$。
又因 $f_*$ 与 $f^*$ 都与有限直和可交换，可进一步设 $n=1$；
这时，命题由定义及等式 $f^*(\OO_Y)=\OO_X$ 立即得到。

\subsection{茎为局部环的赋环空间上的层}\hypertarget{heading:0:5.5}{}
\src{53—54}

\egapar{5.5.1}
称赋环空间 $(X,\OO_X)$ 为\emph{茎为局部环的赋环空间}，
是指对每个 $x\in X$，$\OO_x$ 都是局部环；
在本书考察的赋环空间中，这类空间将占绝大多数。
此时，用 $\mathfrak m_x$ 表示 $\OO_x$ 的极大理想，
用 $k(x)$ 表示剩余域 $\OO_x/\mathfrak m_x$。
对于任意 $\OO_X$-模层 $\F$、任意开集 $U\subseteq X$、任意 $x\in U$，
以及任意截面 $f\in\Gamma(U,\F)$，
以 $f(x)$ 表示芽 $f_x\in\F_x$ 模 $\mathfrak m_x\F_x$ 的类，
称它为\emph{$f$ 在点 $x$ 的值}。
因此，关系 $f(x)=0$ 意味着 $f_x\in\mathfrak m_x\F_x$；
当它成立时，略用术语，称 $f$ \emph{在 $x$ 处为零}。
须注意，不要把这个关系与 $f_x=0$ 混淆。

\egapar{5.5.2}
设 $X$ 是茎为局部环的赋环空间，$\LL$ 是可逆 $\OO_X$-模层，$f$ 是 $\LL$ 在 $X$ 上的截面。
对于一点 $x\in X$，以下三个性质等价：
\begin{enumerate}[label=\alph*)]
\item $f_x$ 是 $\LL_x$ 的一个生成元。
\item $f_x\notin\mathfrak m_x\LL_x$，也就是说，$f(x)\ne0$。
\item 存在 $x$ 的开邻域 $V$ 和 $\LL^{-1}$ 在 $V$ 上的截面 $g$，
使 $f\otimes g$ 在 $\Gamma(V,\OO_X)$ 中的典范像（\hyperlink{para:0:5.4.3}{5.4.3}）为单位截面。
\end{enumerate}
事实上，问题是局部的，可以只考虑 $\LL=\OO_X$ 的情形。
此时 a) 与 b) 的等价性显然，c) 推出 b) 也显然。
另一方面，若 $f_x\notin\mathfrak m_x$，则 $f_x$ 在 $\OO_x$ 中可逆；设 $f_xg_x=1_x$。
依截面芽的定义，这表示存在 $x$ 的邻域 $V$ 和 $\OO_X$ 在 $V$ 上的截面 $g$，
使得在 $V$ 上有 $fg=1$，从而得到 c)。

由条件 c) 立即可知，满足等价条件 a)、b)、c) 的各点组成的集合 $X_f$ 是 $X$ 的开集。
按（\hyperlink{para:0:5.5.1}{5.5.1}）的术语，它就是 $f$ 不为零的点的集合。

\egapar{5.5.3}
在（\hyperlink{para:0:5.5.2}{5.5.2}）的假设下，再设 $\LL'$ 为另一个可逆 $\OO_X$-模层。
若 $f\in\Gamma(X,\LL)$，$g\in\Gamma(X,\LL')$，则
\[
 X_f\cap X_g=X_{f\otimes g}.
\]
事实上，问题是局部的，立即可归结到 $\LL=\LL'=\OO_X$ 的情形；
此时 $f\otimes g$ 典范地等同于乘积 $fg$，命题显然成立。

\egapar{5.5.4}
设 $\F$ 是秩为 $n$ 的局部自由 $\OO_X$-模层。
立即可见，当 $p\leq n$ 时，$\bigwedge^p\F$ 是秩为 $\binom np$ 的局部自由 $\OO_X$-模层；
当 $p>n$ 时，它为零层。
这是因为问题是局部的，可归结到 $\F=\OO_X^n$ 的情形。
此外，对每个 $x\in X$，
\[
 (\bigwedge^p\F)_x\big/\mathfrak m_x(\bigwedge^p\F)_x
\]
是 $k(x)$ 上维数为 $\binom np$ 的向量空间，
并且典范地等同于 $\bigwedge^p(\F_x/\mathfrak m_x\F_x)$。
设 $s_1,\ldots,s_p$ 为 $\F$ 在开集 $U$ 上的截面，
置 $s=s_1\wedge\cdots\wedge s_p$，它是 $\bigwedge^p\F$ 在 $U$ 上的截面（\hyperlink{para:0:4.1.5}{4.1.5}）。
有
\[
 s(x)=s_1(x)\wedge\cdots\wedge s_p(x).
\]
因此，$s_1(x),\ldots,s_p(x)$ 线性相关，等价于 $s(x)=0$。
由此得到，$s_1(x),\ldots,s_p(x)$ 线性无关的点 $x$ 组成一个开集。
事实上，归结到 $\F=\OO_X^n$ 后，只需对 $s$ 在
$\bigwedge^p\F=\OO_X^{\binom np}$ 的各因子上的投影应用（\hyperlink{para:0:5.5.2}{5.5.2}）。

特别地，若 $s_1,\ldots,s_n$ 是 $\F$ 在 $U$ 上的 $n$ 个截面，
且对每个 $x\in U$，$s_1(x),\ldots,s_n(x)$ 都线性无关，
则由这些截面定义的同态 $u\colon\OO_X^n|U\to\F|U$（\hyperlink{para:0:5.1.1}{5.1.1}）是同构。
事实上，可以只考虑 $\F=\OO_X^n$ 并把 $\bigwedge^n\F$ 与 $\OO_X$ 典范等同的情形。
这时，$s=s_1\wedge\cdots\wedge s_n$ 是 $\OO_X$ 在 $U$ 上的可逆截面，
利用 Cramer 公式即可定义 $u$ 的逆同态。

\egapar{5.5.5}
设 $\E,\F$ 是两个有限秩局部自由 $\OO_X$-模层，$u\colon\E\to\F$ 为同态。
存在 $x\in X$ 的邻域 $U$，使 $u|U$ 为单射，且 $\F|U$ 是 $u(\E)|U$
与某个局部自由子 $(\OO_X|U)$-模层 $\G$ 的直和，
其充要条件是 $u_x\colon\E_x\to\F_x$ 取商后给出的向量空间同态
\[
 \E_x/\mathfrak m_x\E_x\longrightarrow\F_x/\mathfrak m_x\F_x
\]
为单射。

条件的必要性如下：此时 $\F_x$ 是自由 $\OO_x$-模 $u_x(\E_x)$ 与 $\G_x$ 的直和，
因此 $\F_x/\mathfrak m_x\F_x$ 是
$u_x(\E_x)/\mathfrak m_xu_x(\E_x)$ 与 $\G_x/\mathfrak m_x\G_x$ 的直和。

为证充分性，可以只考虑 $\E=\OO_X^m$ 的情形。
设 $e_i$ 为 $\OO_X^m$ 的截面，使对于每个 $y\in X$，
$(e_i)_y$ 都等于 $\OO_y^m$ 的典范基的第 $i$ 个元素；这些称为 $\OO_X^m$ 的\emph{典范截面}。
令 $s_1,\ldots,s_m$ 为这些截面在 $u$ 下的像。
按假设，$s_1(x),\ldots,s_m(x)$ 线性无关。
若 $\F$ 的秩为 $n$，则存在 $x$ 的邻域 $V$，以及 $\F$ 在 $V$ 上的 $n-m$ 个截面
$s_{m+1},\ldots,s_n$，使各 $s_i(x)$（$1\leq i\leq n$）构成
$\F_x/\mathfrak m_x\F_x$ 的一组基。
依（\hyperlink{para:0:5.5.4}{5.5.4}），存在邻域 $U\subseteq V$，使得对每个 $y\in U$，
各 $s_i(y)$（$1\leq i\leq n$）都构成 $\F_y/\mathfrak m_y\F_y$ 的一组基。
再依（\hyperlink{para:0:5.5.4}{5.5.4}），得到从 $\F|U$ 到 $\OO_X^n|U$ 的同构，
把 $s_i|U$（$1\leq i\leq m$）分别映到 $e_i|U$。
证明完成。


% END chapters/ch0_s5.tex

\par

% BEGIN chapters/ch0_s6.tex
\section{平坦性}\hypertarget{heading:0:6}{}
\src{54—55（引言）}

\egapar{6.0}
平坦性的概念归于 J.-P. Serre\cite{b16}。
以下对于已经在 N. Bourbaki 的《交换代数》中阐述的结果，省略其证明，读者可参阅该书。
我们假设所有环都交换。
\origfootnote{关于这些结果中的大部分如何推广到非交换情形，参见上述 N. Bourbaki 的论述。}

若 $M,N$ 为两个 $A$-模，$M'$ 是 $M$ 的子模，$N'$ 是 $N$ 的子模，
则以 $\Im(M'\otimes_A N')$ 表示典范映射
$M'\otimes_A N'\to M\otimes_A N$ 的像，即 $M\otimes_A N$ 的相应子模。

\subsection{平坦模}\hypertarget{heading:0:6.1}{}
\src{55}

\egapar{6.1.1}
设 $M$ 为 $A$-模。以下性质等价：
\begin{enumerate}[label=\alph*)]
\item $M\otimes_A N$ 是 $A$-模范畴中关于 $N$ 的正合函子。
\item 对每个 $i>0$ 和每个 $A$-模 $N$，都有 $\Tor_i^A(M,N)=0$。
\item 对每个 $A$-模 $N$，都有 $\Tor_1^A(M,N)=0$。
\end{enumerate}
当 $M$ 满足这些条件时，称 $M$ 是\emph{平坦 $A$-模}。
显然，每个自由 $A$-模都是平坦的。

要使 $M$ 为平坦 $A$-模，只需对于 $A$ 的每个有限型理想 $\mathfrak I$，
典范映射 $M\otimes_A\mathfrak I\to M\otimes_A A=M$ 为单射。

\egapar{6.1.2}
平坦 $A$-模的归纳极限仍为平坦 $A$-模。
$A$-模的直和 $\bigoplus_{\lambda\in L}M_\lambda$ 平坦，当且仅当每个 $M_\lambda$ 都平坦。
特别地，每个射影 $A$-模都平坦。

设 $0\to M'\to M\to M''\to0$ 为 $A$-模的正合列，且 $M''$ 平坦。
则对每个 $A$-模 $N$，序列
\[
 0\longrightarrow M'\otimes_A N\longrightarrow M\otimes_A N
 \longrightarrow M''\otimes_A N\longrightarrow0
\]
正合。此外，此时 $M$ 平坦，当且仅当 $M'$ 平坦。
但有可能 $M,M'$ 都平坦，而 $M''=M/M'$ 不平坦。

\egapar{6.1.3}
设 $M$ 为平坦 $A$-模，$N$ 为任意 $A$-模。对于 $N$ 的两个子模 $N',N''$，有
\[
 \begin{aligned}
 \Im\bigl(M\otimes_A(N'+N'')\bigr)
 &=\Im(M\otimes_A N')+\Im(M\otimes_A N''),\\
 \Im\bigl(M\otimes_A(N'\cap N'')\bigr)
 &=\Im(M\otimes_A N')\cap\Im(M\otimes_A N''),
 \end{aligned}
\]
这里的像都取在 $M\otimes_A N$ 中。

\egapar{6.1.4}
设 $M,N$ 为两个 $A$-模，$M'\subseteq M,N'\subseteq N$ 为子模，
并假设 $M/M'$ 与 $N/N'$ 中至少一个平坦。则
\[
 \Im(M'\otimes_A N')
 =\Im(M'\otimes_A N)\cap\Im(M\otimes_A N'),
\]
这里的像取在 $M\otimes_A N$ 中。
特别地，若 $\mathfrak I$ 是 $A$ 的理想，且 $M/M'$ 平坦，则
$\mathfrak I M'=M'\cap\mathfrak I M$。

\subsection{环的变更}\hypertarget{heading:0:6.2}{}
\src{55—56}

当一个加法群 $M$ 具有相对于环 $A,B,\ldots$ 的若干模结构时，
我们有时说 $M$ 是 $A$-平坦的、$B$-平坦的，等等，
来代替“$M$ 作为 $A$-模、$B$-模等是平坦的”这一说法。

\egapar{6.2.1}
设 $A,B$ 为两个环，$M$ 为 $A$-模，$N$ 为 $(A,B)$-双模。
若 $M$ 平坦，且 $N$ 是 $B$-平坦的，则 $M\otimes_A N$ 是 $B$-平坦的。
特别地，若 $M,N$ 都是平坦 $A$-模，则 $M\otimes_A N$ 是平坦 $A$-模。
若 $B$ 为 $A$-代数，$M$ 为平坦 $A$-模，则 $B$-模 $M_{(B)}=M\otimes_A B$ 平坦。
最后，若 $B$ 为作为 $A$-模平坦的 $A$-代数，$N$ 为平坦 $B$-模，则 $N$ 也是 $A$-平坦的。

\egapar{6.2.2}
设 $A$ 为环，$B$ 为作为 $A$-模平坦的 $A$-代数。
设 $M,N$ 为两个 $A$-模，且 $M$ 具有有限表示。
则典范同态
\begin{equation}
 \Hom_A(M,N)\otimes_A B
 \longrightarrow\Hom_B(M\otimes_A B,N\otimes_A B)
 \tag{6.2.2.1}\label{eq:0:6.2.2.1}
\end{equation}
是同构；此同态把 $u\otimes b$ 映到同态
$m\otimes b'\mapsto u(m)\otimes bb'$。

\egapar{6.2.3}
设 $(A_\lambda,\varphi_{\mu\lambda})$ 为环的滤过归纳系，$A=\ilim A_\lambda$。
另设对每个 $\lambda$，$M_\lambda$ 是 $A_\lambda$-模；
对 $\lambda\leq\mu$，$\theta_{\mu\lambda}\colon M_\lambda\to M_\mu$ 是
$\varphi_{\mu\lambda}$-同态，且 $(M_\lambda,\theta_{\mu\lambda})$ 构成归纳系。
于是 $M=\ilim M_\lambda$ 是 $A$-模。
若对每个 $\lambda$，$M_\lambda$ 都是平坦 $A_\lambda$-模，则 $M$ 是平坦 $A$-模。

事实上，设 $\mathfrak I$ 是 $A$ 的有限型理想。
依归纳极限的定义，存在指标 $\lambda$ 及 $A_\lambda$ 的理想 $\mathfrak I_\lambda$，
使 $\mathfrak I=\mathfrak I_\lambda A$。
对 $\mu\geq\lambda$，置 $\mathfrak I_\mu=\mathfrak I_\lambda A_\mu$，
则还有 $\mathfrak I=\ilim_{\mu\geq\lambda}\mathfrak I_\mu$。
因归纳极限函子正合且与张量积可交换，有
\[
 M\otimes_A\mathfrak I
 =\ilim_{\mu\geq\lambda}(M_\mu\otimes_{A_\mu}\mathfrak I_\mu)
 =\ilim_{\mu\geq\lambda}\mathfrak I_\mu M_\mu
 =\mathfrak I M.
\]

\subsection{平坦性的局部化}\hypertarget{heading:0:6.3}{}
\src{56—57}

\egapar{6.3.1}
若 $A$ 为环，$S$ 为 $A$ 的乘法子集，则 $S^{-1}A$ 是平坦 $A$-模。
事实上，对每个 $A$-模 $N$，$N\otimes_A S^{-1}A$ 与 $S^{-1}N$ 等同（\hyperlink{para:0:1.2.5}{1.2.5}），
而已知 $S^{-1}N$ 是关于 $N$ 的正合函子（\hyperlink{para:0:1.3.2}{1.3.2}）。

若 $M$ 为平坦 $A$-模，则 $S^{-1}M=M\otimes_A S^{-1}A$ 为平坦 $S^{-1}A$-模（\hyperlink{para:0:6.2.1}{6.2.1}），
由以上结果及（\hyperlink{para:0:6.2.1}{6.2.1}），它也是 $A$-平坦的。
特别地，若 $P$ 为 $S^{-1}A$-模，则可把它看作与 $S^{-1}P$ 同构的 $A$-模；
$P$ 是 $A$-平坦的，当且仅当它是 $S^{-1}A$-平坦的。

\egapar{6.3.2}
设 $A$ 为环，$B$ 为 $A$-代数，$T$ 为 $B$ 的乘法子集。
若 $P$ 为 $A$-平坦的 $B$-模，则 $T^{-1}P$ 是 $A$-平坦的。
事实上，对于任意 $A$-模 $N$，有
\[
 \begin{aligned}
 (T^{-1}P)\otimes_A N
 &=(T^{-1}B\otimes_B P)\otimes_A N\\
 &=T^{-1}B\otimes_B(P\otimes_A N)
 =T^{-1}(P\otimes_A N).
 \end{aligned}
\]
而 $T^{-1}(P\otimes_A N)$ 是关于 $N$ 的正合函子，
因为它是关于 $N$ 的正合函子 $P\otimes_A N$ 与关于 $Q$ 的正合函子 $T^{-1}Q$ 的复合。
若 $S$ 是 $A$ 的乘法子集，且其在 $B$ 中的像包含于 $T$，
则 $T^{-1}P=S^{-1}(T^{-1}P)$，因而依（\hyperlink{para:0:6.3.1}{6.3.1}）也是 $S^{-1}A$-平坦的。

\egapar{6.3.3}
设 $\varphi\colon A\to B$ 为环同态，$M$ 为 $B$-模。以下性质等价：
\begin{enumerate}[label=\alph*)]
\item $M$ 是平坦 $A$-模。
\item 对 $B$ 的每个极大理想 $\mathfrak n$，$M_{\mathfrak n}$ 是平坦 $A$-模。
\item 对 $B$ 的每个极大理想 $\mathfrak n$，置 $\mathfrak m=\varphi^{-1}(\mathfrak n)$，
则 $M_{\mathfrak n}$ 是平坦 $A_{\mathfrak m}$-模。
\end{enumerate}
由于 $M_{\mathfrak n}=(M_{\mathfrak n})_{\mathfrak m}$，
b) 与 c) 的等价性由（\hyperlink{para:0:6.3.1}{6.3.1}）得到；a) 推出 b) 是（\hyperlink{para:0:6.3.2}{6.3.2}）的特殊情形。
只须证明 b) 推出 a)，即对每个 $A$-模的单射同态 $u\colon N'\to N$，
同态 $v=1\otimes u\colon M\otimes_A N'\to M\otimes_A N$ 为单射。
$v$ 也是 $B$-模同态；已知要使它为单射，只需对 $B$ 的每个极大理想 $\mathfrak n$，
同态 $v_{\mathfrak n}\colon(M\otimes_A N')_{\mathfrak n}\to(M\otimes_A N)_{\mathfrak n}$ 为单射。
而
\[
 (M\otimes_A N)_{\mathfrak n}
 =B_{\mathfrak n}\otimes_B(M\otimes_A N)
 =M_{\mathfrak n}\otimes_A N,
\]
所以 $v_{\mathfrak n}$ 正是
$1\otimes u\colon M_{\mathfrak n}\otimes_A N'\to M_{\mathfrak n}\otimes_A N$；
因 $M_{\mathfrak n}$ 为 $A$-平坦的，该同态为单射。

特别地，取 $B=A$ 可知：$A$-模 $M$ 平坦，当且仅当对 $A$ 的每个极大理想 $\mathfrak m$，
$M_{\mathfrak m}$ 都是 $A_{\mathfrak m}$-平坦的。

\egapar{6.3.4}
设 $M$ 为 $A$-模。若 $M$ 平坦，且 $f\in A$ 不是 $A$ 中的零因子，
则 $f$ 不零化 $M$ 的任何非零元素。
因为把 $M$ 与 $M\otimes_A A$ 等同后，同态 $m\mapsto fm$ 写成 $1\otimes u$，
其中 $u$ 是 $A$ 中的乘法 $a\mapsto fa$；
若 $u$ 为单射，则因 $M$ 平坦，$1\otimes u$ 也为单射。
特别地，若 $A$ 是整环，则 $M$ 无挠。

反之，假设 $A$ 为整环，$M$ 无挠，且对 $A$ 的每个极大理想 $\mathfrak m$，
$A_{\mathfrak m}$ 都是离散赋值环，则 $M$ 是 $A$-平坦的。
事实上，只需证明 $M_{\mathfrak m}$ 是 $A_{\mathfrak m}$-平坦的（\hyperlink{para:0:6.3.3}{6.3.3}），
因此可以假设 $A$ 本身就是离散赋值环。
又因 $M$ 是其有限型子模的归纳极限，且这些子模无挠，
还可只考虑 $M$ 有限型的情形（\hyperlink{para:0:6.1.2}{6.1.2}）。
此时，命题由 $M$ 是自由 $A$-模这一事实得到。

特别地，若 $A$ 为整环，环同态 $\varphi\colon A\to B$ 使 $B$ 成为非零的平坦 $A$-模，
则 $\varphi$ 必为单射。
反之，若 $B$ 为整环，$A$ 为 $B$ 的子环，且对 $A$ 的每个极大理想 $\mathfrak m$，
$A_{\mathfrak m}$ 都是离散赋值环，则 $B$ 为 $A$-平坦的。

\subsection{忠实平坦模}\hypertarget{heading:0:6.4}{}
\src{57—58}

\egapar{6.4.1}
对于 $A$-模 $M$，以下四个性质等价：
\begin{enumerate}[label=\alph*)]
\item $A$-模序列 $N'\to N\to N''$ 正合，当且仅当序列
$M\otimes_A N'\to M\otimes_A N\to M\otimes_A N''$ 正合。
\item $M$ 平坦，且对每个 $A$-模 $N$，关系 $M\otimes_A N=0$ 蕴含 $N=0$。
\item $M$ 平坦，且对每个 $A$-模同态 $v\colon N\to N'$，
关系 $1_M\otimes v=0$ 蕴含 $v=0$，其中 $1_M$ 为 $M$ 的恒等自同构。
\item $M$ 平坦，且对 $A$ 的每个极大理想 $\mathfrak m$，有 $\mathfrak mM\ne M$。
\end{enumerate}
当 $M$ 满足这些条件时，称 $M$ 为\emph{忠实平坦 $A$-模}；此时 $M$ 必为忠实模。
此外，若 $u\colon N\to N'$ 为 $A$-模同态，则 $u$ 是单射（分别地，满射、双射），
当且仅当 $1\otimes u\colon M\otimes_A N\to M\otimes_A N'$ 是单射（分别地，满射、双射）。

\egapar{6.4.2}
非零自由模是忠实平坦的；一个平坦模与一个忠实平坦模的直和也是忠实平坦的。
若 $S$ 为 $A$ 的乘法子集，则 $S^{-1}A$ 只有在 $S$ 全由可逆元组成时才是忠实平坦 $A$-模，
此时 $S^{-1}A\simeq A$。

\egapar{6.4.3}
设 $0\to M'\to M\to M''\to0$ 为 $A$-模的正合列。
若 $M',M''$ 都平坦，且其中一个忠实平坦，则 $M$ 忠实平坦。

\egapar{6.4.4}
设 $A,B$ 为两个环，$M$ 为 $A$-模，$N$ 为 $(A,B)$-双模。
若 $M$ 忠实平坦，且 $N$ 为忠实平坦 $B$-模，则 $M\otimes_A N$ 为忠实平坦 $B$-模。
特别地，若 $M,N$ 都为忠实平坦 $A$-模，则 $M\otimes_A N$ 也为忠实平坦 $A$-模。
若 $B$ 为 $A$-代数，$M$ 为忠实平坦 $A$-模，则 $B$-模 $M_{(B)}$ 忠实平坦。

\egapar{6.4.5}
若 $M$ 为忠实平坦 $A$-模，$S$ 为 $A$ 的乘法子集，
则 $S^{-1}M$ 为忠实平坦 $S^{-1}A$-模，
因为 $S^{-1}M=M\otimes_A S^{-1}A$（\hyperlink{para:0:6.4.4}{6.4.4}）。
反之，若对 $A$ 的每个极大理想 $\mathfrak m$，
$M_{\mathfrak m}$ 都为忠实平坦 $A_{\mathfrak m}$-模，则 $M$ 为忠实平坦 $A$-模。
因为 $M$ 是 $A$-平坦的（\hyperlink{para:0:6.3.3}{6.3.3}），并且
\[
 \begin{aligned}
 M_{\mathfrak m}/\mathfrak mM_{\mathfrak m}
 &=(M\otimes_A A_{\mathfrak m})
   \otimes_{A_{\mathfrak m}}(A_{\mathfrak m}/\mathfrak mA_{\mathfrak m})\\
 &=M\otimes_A(A/\mathfrak m)=M/\mathfrak mM.
 \end{aligned}
\]
所以假设蕴含对 $A$ 的每个极大理想 $\mathfrak m$，$M/\mathfrak mM\ne0$，
由（\hyperlink{para:0:6.4.1}{6.4.1}）即得结论。

\subsection{标量限制}\hypertarget{heading:0:6.5}{}
\src{58}

\egapar{6.5.1}
设 $A$ 为环，$\varphi\colon A\to B$ 为把 $B$ 赋成 $A$-代数的环同态。
假设存在一个 $B$-模 $N$，它作为 $A$-模是忠实平坦的。
则对每个 $A$-模 $M$，从 $M$ 到 $B\otimes_A M=M_{(B)}$ 的同态 $x\mapsto1\otimes x$ 为单射。
特别地，$\varphi$ 为单射；
对 $A$ 的每个理想 $\mathfrak a$，有 $\varphi^{-1}(\mathfrak aB)=\mathfrak a$；
对 $A$ 的每个极大理想（分别地，素理想）$\mathfrak m$，
存在 $B$ 的一个极大理想（分别地，素理想）$\mathfrak n$，
使 $\varphi^{-1}(\mathfrak n)=\mathfrak m$。

\egapar{6.5.2}
当（\hyperlink{para:0:6.5.1}{6.5.1}）的条件满足时，通过 $\varphi$ 把 $A$ 与 $B$ 的子环等同；
更一般地，对每个 $A$-模 $M$，把 $M$ 与 $M_{(B)}$ 的子 $A$-模等同。
注意，若此时 $B$ 为诺特环，则 $A$ 也为诺特环。
因为映射 $\mathfrak a\mapsto\mathfrak aB$ 是从 $A$ 的理想集到 $B$ 的理想集的保序单射；
若 $A$ 中存在无限严格递增的理想序列，就会导致 $B$ 中也存在这样的序列。

\subsection{忠实平坦环}\hypertarget{heading:0:6.6}{}
\src{58—59}

\egapar{6.6.1}
设 $\varphi\colon A\to B$ 为把 $B$ 赋成 $A$-代数的环同态。以下五个性质等价：
\begin{enumerate}[label=\alph*)]
\item $B$ 为忠实平坦 $A$-模；换言之，$M_{(B)}$ 是关于 $M$ 的正合且忠实的函子。
\item $\varphi$ 为单射，且 $A$-模 $B/\varphi(A)$ 平坦。
\item $A$-模 $B$ 平坦，即函子 $M_{(B)}$ 正合；并且对每个 $A$-模 $M$，
从 $M$ 到 $M_{(B)}$ 的同态 $x\mapsto1\otimes x$ 为单射。
\item $A$-模 $B$ 平坦，且对 $A$ 的每个理想 $\mathfrak a$，有
$\varphi^{-1}(\mathfrak aB)=\mathfrak a$。
\item $A$-模 $B$ 平坦，且对 $A$ 的每个极大理想 $\mathfrak m$，
存在 $B$ 的极大理想 $\mathfrak n$，使 $\varphi^{-1}(\mathfrak n)=\mathfrak m$。
\end{enumerate}
当这些条件成立时，把 $A$ 与 $B$ 的子环等同。

\egapar{6.6.2}
设 $A$ 为局部环，$\mathfrak m$ 为其极大理想，$B$ 为满足 $\mathfrak mB\ne B$ 的 $A$-代数。
例如，当 $B$ 为局部环且 $A\to B$ 是局部同态时，这一条件就成立。
若 $B$ 为平坦 $A$-模，则 $B$ 为忠实平坦 $A$-模。
事实上，因 $\mathfrak mB\ne B$，存在 $B$ 的极大理想 $\mathfrak n$ 包含 $\mathfrak mB$。
由于 $\varphi^{-1}(\mathfrak n)$ 包含 $\mathfrak m$ 而不包含 $1$，
有 $\varphi^{-1}(\mathfrak n)=\mathfrak m$，故可应用（\hyperlink{para:0:6.6.1}{6.6.1}）的判据 e)。

在以上条件下，若 $B$ 为诺特环，则 $A$ 也为诺特环（\hyperlink{para:0:6.5.2}{6.5.2}）。

\egapar{6.6.3}
设 $B$ 为忠实平坦的 $A$-代数。
对于任意 $A$-模 $M$ 及其子 $A$-模 $M'$，有
\[
 M'=M\cap M'_{(B)},
\]
这里把 $M$ 与 $M_{(B)}$ 的子 $A$-模等同。
$M$ 为平坦 $A$-模（分别地，忠实平坦 $A$-模），
当且仅当 $M_{(B)}$ 为平坦 $B$-模（分别地，忠实平坦 $B$-模）。

\egapar{6.6.4}
设 $B$ 为 $A$-代数，$N$ 为忠实平坦 $B$-模。
$B$ 为平坦 $A$-模（分别地，忠实平坦 $A$-模），
当且仅当 $N$ 为平坦 $A$-模（分别地，忠实平坦 $A$-模）。

特别地，设 $C$ 为 $B$-代数。
若环 $C$ 在 $B$ 上忠实平坦，且 $B$ 在 $A$ 上忠实平坦，则 $C$ 在 $A$ 上忠实平坦；
若 $C$ 在 $B$ 上和在 $A$ 上都忠实平坦，则 $B$ 在 $A$ 上忠实平坦。

\subsection{赋环空间的平坦态射}\hypertarget{heading:0:6.7}{}
\src{59—60}

\egapar{6.7.1}
设 $f\colon X\to Y$ 为赋环空间态射，$\F$ 为 $\OO_X$-模层。
若 $\F_x$ 是平坦 $\OO_{f(x)}$-模，称 $\F$ \emph{在点 $x\in X$ 是 $f$-平坦的}；
当 $f$ 不致混淆时，也称为\emph{在 $x$ 是 $Y$-平坦的}。
若 $\F$ 在每个 $x\in f^{-1}(y)$ 都是 $f$-平坦的，称 $\F$ \emph{在 $y\in Y$ 之上是 $f$-平坦的}。
若 $\F$ 在 $X$ 的每一点都 $f$-平坦，称 $\F$ \emph{是 $f$-平坦的}。
若 $\OO_X$ 在 $x$ 是 $f$-平坦的（分别地，在 $y$ 之上是 $f$-平坦的、是 $f$-平坦的），
则称态射 $f$ \emph{在 $x$ 平坦}（分别地，\emph{在 $y$ 之上平坦}、\emph{平坦}）。

\egapar{6.7.2}
沿用（\hyperlink{para:0:6.7.1}{6.7.1}）的记号。
若 $\F$ 在 $x$ 是 $f$-平坦的，则对 $y=f(x)$ 的每个开邻域 $U$，
\[
 (f^*(\G)\otimes_{\OO_X}\F)_x
\]
是 $(\OO_Y|U)$-模层范畴中关于 $\G$ 的正合函子。
事实上，这个茎典范地等同于 $\G_y\otimes_{\OO_y}\F_x$，结论由定义得到。
特别地，若 $f$ 是平坦态射，则 $f^*$ 在 $\OO_Y$-模层范畴中为正合函子。

\egapar{6.7.3}
反之，假设环层 $\OO_Y$ 凝聚，并假设对 $y$ 的每个开邻域 $U$，
函子 $(f^*(\G)\otimes_{\OO_X}\F)_x$ 在凝聚 $(\OO_Y|U)$-模层范畴中关于 $\G$ 正合。
则 $\F$ 在 $x$ 是 $f$-平坦的。
事实上，只需证明对于 $\OO_y$ 的每个有限型理想 $\mathfrak I$，
典范同态 $\mathfrak I\otimes_{\OO_y}\F_x\to\F_x$ 为单射（\hyperlink{para:0:6.1.1}{6.1.1}）。
而由（\hyperlink{para:0:5.3.8}{5.3.8}），存在 $y$ 的开邻域 $U$，以及 $\OO_Y|U$ 的凝聚理想层 $\J$，
使 $\J_y=\mathfrak I$；由此即得结论。

\egapar{6.7.4}
（\hyperlink{heading:0:6.1}{6.1}）中关于平坦模的结果，立即转写为关于在一点 $f$-平坦的层的命题。

若 $0\to\F'\to\F\to\F''\to0$ 是 $\OO_X$-模层的正合列，
且 $\F''$ 在点 $x\in X$ 是 $f$-平坦的，
则对于 $y=f(x)$ 的每个开邻域 $U$ 及每个 $(\OO_Y|U)$-模层 $\G$，序列
\[
 \begin{aligned}
 0\longrightarrow(f^*(\G)\otimes_{\OO_X}\F')_x
 &\longrightarrow(f^*(\G)\otimes_{\OO_X}\F)_x\\
 &\longrightarrow(f^*(\G)\otimes_{\OO_X}\F'')_x
 \longrightarrow0
 \end{aligned}
\]
正合。此时，$\F$ 在 $x$ 是 $f$-平坦的，当且仅当 $\F'$ 在 $x$ 是 $f$-平坦的。
对于在 $y\in Y$ 之上 $f$-平坦的 $\OO_X$-模层，以及 $f$-平坦的 $\OO_X$-模层，
都有相应的类似命题。

\egapar{6.7.5}
设 $f\colon X\to Y,g\colon Y\to Z$ 为两个赋环空间态射，
$x\in X,y=f(x)$，$\F$ 为 $\OO_X$-模层。
若 $\F$ 在 $x$ 是 $f$-平坦的，且态射 $g$ 在 $y$ 平坦，
则 $\F$ 在 $x$ 是 $(g\circ f)$-平坦的（\hyperlink{para:0:6.2.1}{6.2.1}）。
特别地，若 $f,g$ 都是平坦态射，则 $g\circ f$ 平坦。

\egapar{6.7.6}
设 $X,Y$ 为赋环空间，$f\colon X\to Y$ 为平坦态射。
则双函子的典范同态（\hyperlink{para:0:4.4.6}{4.4.6}）
\begin{equation}
 f^*(\shHom_{\OO_Y}(\F,\G))
 \longrightarrow\shHom_{\OO_X}(f^*(\F),f^*(\G))
 \tag{6.7.6.1}\label{eq:0:6.7.6.1}
\end{equation}
在 $\F$ 具有有限表示时为同构（\hyperlink{para:0:5.2.5}{5.2.5}）。

事实上，问题是局部的，可以假设存在正合列
$\OO_Y^p\to\OO_Y^q\to\F\to0$。
由关于 $f$ 的假设，（\hyperref[eq:0:6.7.6.1]{6.7.6.1}）的两边都是关于 $\F$ 的左正合函子。
因此只需在 $\F=\OO_Y$ 时证明命题，而此时结论显然。

\egapar{6.7.8}
称赋环空间态射 $f\colon X\to Y$ \emph{忠实平坦}，
是指 $f$ 为满射，且对每个 $x\in X$，$\OO_x$ 都是忠实平坦 $\OO_{f(x)}$-模。
当 $X,Y$ 是茎为局部环的赋环空间时（\hyperlink{para:0:5.5.1}{5.5.1}），
这等价于说态射 $f$ 为满射且平坦（\hyperlink{para:0:6.6.2}{6.6.2}）。
\transnote{原版在（\hyperlink{para:0:6.7.6}{6.7.6}）之后直接编号为（\hyperlink{para:0:6.7.8}{6.7.8}），此处照录。
本句引用（\hyperlink{para:0:6.6.2}{6.6.2}）从平坦推出忠实平坦时，还须保证每个茎同态
$\OO_{Y,f(x)}\to\OO_{X,x}$ 是局部同态，或至少满足
$\mathfrak m_{f(x)}\OO_{X,x}\ne\OO_{X,x}$。
仅知道两边的茎环都是局部环，并不提供这个条件。
例如，在两个单点赋环空间上，取结构环同态 $\mathbb Z_{(p)}\hookrightarrow\mathbb Q$。
对应态射的底空间映射满射，并且它平坦（\hyperlink{para:0:6.3.1}{6.3.1}），但
$\mathbb Q\otimes_{\mathbb Z_{(p)}}\mathbb F_p=0$，而 $\mathbb F_p\ne0$，所以它不忠实平坦。
这里该茎同态不是局部同态。正文保留原版措辞；反例为译者补充。}
当 $f$ 忠实平坦时，$f^*$ 是 $\OO_Y$-模层范畴中的正合且忠实的函子（\hyperlink{para:0:6.6.1}{6.6.1}，a)）。
一个 $\OO_Y$-模层 $\G$ 是 $Y$-平坦的，当且仅当 $f^*(\G)$ 是 $Y$-平坦的（\hyperlink{para:0:6.6.3}{6.6.3}）。


% END chapters/ch0_s6.tex

\par

% BEGIN chapters/ch0_s7_1.tex
\section{进制环}\hypertarget{heading:0:7}{}
\subsection{可容许环}\hypertarget{heading:0:7.1}{}
\src{60—62}

\egapar{7.1.1}
回顾一下：在拓扑环 $A$（不必是分离的）中，若 $0$ 是序列 $(x^n)_{n\geq 0}$ 的一个极限，
则称元素 $x$ 是\emph{拓扑幂零的}。若 $A$ 中存在一个由理想（它们必然是开的）组成的
$0$ 的基本邻域系，则称拓扑环 $A$ 是\emph{线性拓扑化的}。

\egastatement{7.1.2}{定义}
在线性拓扑化环 $A$ 中，若理想 $\frJ$ 是开的，并且对于 $0$ 的每个邻域 $V$，
都存在整数 $n>0$ 使得 $\frJ^n\subset V$，则称 $\frJ$ 是一个\emph{定义理想}
（通常也略称序列 $(\frJ^n)$ 趋于 $0$）。若线性拓扑化环 $A$ 中存在定义理想，
则称 $A$ 是\emph{预可容许的}；若 $A$ 是预可容许的，并且还是分离且完备的，
则称 $A$ 是\emph{可容许的}。

显然，若 $\frJ$ 是定义理想，$\frQ$ 是 $A$ 的开理想，则 $\frJ\cap\frQ$ 仍是定义理想；
因而，预可容许环 $A$ 的定义理想组成 $0$ 的一个基本邻域系。

\egastatement{7.1.3}{引理}
设 $A$ 是线性拓扑化环。
\begin{enumerate}[label=(\roman*)]
\item 元素 $x\in A$ 为拓扑幂零元，当且仅当对于 $A$ 的每个开理想 $\frJ$，
$x$ 在 $A/\frJ$ 中的典范像是幂零元。$A$ 中拓扑幂零元的集合 $\frT$ 是一个理想。
\item 再设 $A$ 是预可容许的，且 $\frJ$ 是 $A$ 的一个定义理想。
元素 $x\in A$ 为拓扑幂零元，当且仅当它在 $A/\frJ$ 中的典范像是幂零元；
理想 $\frT$ 是 $A/\frJ$ 的幂零根的逆像，因而是开的。
\end{enumerate}

(i) 直接由定义得到。为证明 (ii)，只需注意，对于 $A$ 中 $0$ 的每个邻域 $V$，
存在 $n>0$ 使得 $\frJ^n\subset V$；若 $x\in A$ 满足 $x^m\in\frJ$，
则对 $q\geq n$ 有 $x^{mq}\in V$，故 $x$ 是拓扑幂零元。

\egastatement{7.1.4}{命题}
设 $A$ 是预可容许环，$\frJ$ 是 $A$ 的一个定义理想。
\begin{enumerate}[label=(\roman*)]
\item $A$ 的理想 $\frJ'$ 包含于某个定义理想，当且仅当存在整数 $n>0$ 使得
$\frJ'{}^n\subset\frJ$。
\item 元素 $x\in A$ 属于某个定义理想，当且仅当它是拓扑幂零元。
\end{enumerate}

(i) 若 $\frJ'{}^n\subset\frJ$，则对于 $A$ 中 $0$ 的每个开邻域 $V$，
存在 $m$ 使得 $\frJ^m\subset V$，从而 $\frJ'{}^{mn}\subset V$。

(ii) 必要性显然。反之，若条件成立，则存在 $n$ 使得 $x^n\in\frJ$；
因此 $\frJ'=\frJ+Ax$ 是定义理想，因为它是开的，而且 $\frJ'{}^n\subset\frJ$。

\egastatement{7.1.5}{推论}
在预可容许环 $A$ 中，每个开素理想都包含所有定义理想。

\egastatement{7.1.6}{推论}
沿用（\hyperlink{para:0:7.1.4}{7.1.4}）的记号和假设。对于 $A$ 的理想 $\frJ_0$，以下性质等价：
\begin{enumerate}[label=\alph*)]
\item $\frJ_0$ 是 $A$ 的最大定义理想；
\item $\frJ_0$ 是一个极大定义理想；
\item $\frJ_0$ 是定义理想，而且环 $A/\frJ_0$ 是约化的。
\end{enumerate}
具有这些性质的理想 $\frJ_0$ 存在，当且仅当 $A/\frJ$ 的幂零根是幂零理想；
此时 $\frJ_0$ 就等于 $A$ 中拓扑幂零元所成的理想 $\frT$。

显然 a) 蕴含 b)；由（\hyperlink{para:0:7.1.4}{7.1.4}，(ii)）和（\hyperlink{para:0:7.1.3}{7.1.3}，(ii)），b) 蕴含 c)；
出于同样的理由，c) 蕴含 a)。最后的断言由（\hyperlink{para:0:7.1.4}{7.1.4}，(i)）和（\hyperlink{para:0:7.1.3}{7.1.3}，(ii)）得到。

当 $A/\frJ$ 的幂零根 $\frT/\frJ$ 是幂零理想时，以 $A_{\mathrm{red}}$ 记商环
$A/\frT$；这是一个约化环。

\egastatement{7.1.7}{推论}
预可容许的诺特环具有最大定义理想。

\egastatement{7.1.8}{推论}
设预可容许环 $A$ 的某个定义理想 $\frJ$ 的各次幂 $\frJ^n$（$n>0$）
组成 $0$ 的一个基本邻域系。则 $A$ 的任意定义理想 $\frJ'$ 的各次幂
$\frJ'{}^n$ 也具有这一性质。

\egastatement{7.1.9}{定义}
若预可容许环 $A$ 具有一个定义理想 $\frJ$，使得 $\frJ^n$ 组成 $A$ 中 $0$ 的
一个基本邻域系（等价地，各 $\frJ^n$ 都是开的），则称 $A$ 是\emph{预进制环}。
分离且完备的预进制环称为\emph{进制环}。

若 $\frJ$ 是预进制环（分别地，进制环）$A$ 的一个定义理想，也称 $A$ 为
\emph{$\frJ$-预进制环}（分别地，\emph{$\frJ$-进制环}），并称其拓扑为
\emph{$\frJ$-预进拓扑}（分别地，\emph{$\frJ$-进拓扑}）。更一般地，若 $M$ 是
$A$-模，则以子模 $\frJ^nM$ 为 $0$ 的基本邻域系的拓扑称为
$\frJ$-预进拓扑（分别地，$\frJ$-进拓扑）。由（\hyperlink{para:0:7.1.8}{7.1.8}），这些拓扑与定义理想
$\frJ$ 的选择无关。\transnote{本册保留原著对 \textit{préadique} 与 \textit{adique} 的术语区分；
不要仅因出现“进”字便略去原文中的分离性、完备性条件。对模的陈述还应逐项核对其所列假设。}

\egastatement{7.1.10}{命题}
设 $A$ 是可容许环，$\frJ$ 是 $A$ 的一个定义理想。则 $\frJ$ 包含于 $A$ 的根基中。
\transnote{这里的“根基”是（\hyperlink{para:0:1.1.2}{1.1.2}）中的 Jacobson 根基，而不是幂零根。}

这一命题与下列任意一个推论等价。

\egastatement{7.1.11}{推论}
对每个 $x\in\frJ$，元素 $1+x$ 在 $A$ 中可逆。

\egastatement{7.1.12}{推论}
元素 $f\in A$ 在 $A$ 中可逆，当且仅当它在 $A/\frJ$ 中的典范像可逆。

\egastatement{7.1.13}{推论}
对于每个有限型 $A$-模 $M$，关系 $M=\frJ M$
（等价于 $M\otimes_A(A/\frJ)=0$）蕴含 $M=0$。

\egastatement{7.1.14}{推论}
设 $u:M\to N$ 是 $A$-模同态，且 $N$ 是有限型的。则 $u$ 为满射，当且仅当
\[
 u\otimes 1:M\otimes_A(A/\frJ)\longrightarrow N\otimes_A(A/\frJ)
\]
为满射。

事实上，（\hyperlink{para:0:7.1.10}{7.1.10}）与（\hyperlink{para:0:7.1.11}{7.1.11}）的等价性由 Bourbaki，\textit{Alg.}，第 VIII 章，
§6，第 3 节，定理 1 得到；（\hyperlink{para:0:7.1.10}{7.1.10}）与（\hyperlink{para:0:7.1.13}{7.1.13}）的等价性由同处定理 2 得到。
（\hyperlink{para:0:7.1.10}{7.1.10}）蕴含（\hyperlink{para:0:7.1.14}{7.1.14}），由同处命题 6 的推论 4 得到；另一方面，
将（\hyperlink{para:0:7.1.14}{7.1.14}）用于零同态，便得到（\hyperlink{para:0:7.1.13}{7.1.13}）。最后，由（\hyperlink{para:0:7.1.10}{7.1.10}），若 $f$ 在
$A/\frJ$ 中的像可逆，则 $f$ 不属于 $A$ 的任何极大理想，所以 $f$ 在 $A$ 中可逆；
也就是说，（\hyperlink{para:0:7.1.10}{7.1.10}）蕴含（\hyperlink{para:0:7.1.12}{7.1.12}）。反之，（\hyperlink{para:0:7.1.12}{7.1.12}）蕴含（\hyperlink{para:0:7.1.11}{7.1.11}）。

因而只须证明（\hyperlink{para:0:7.1.11}{7.1.11}）。由于 $A$ 分离且完备，又由于序列 $(\frJ^n)$ 趋于 $0$，
级数
\[
 \sum_{n=0}^{\infty}(-1)^n x^n
\]
显然在 $A$ 中收敛；若以 $y$ 记其和，则有 $y(1+x)=1$。


% END chapters/ch0_s7_1.tex

\par

% BEGIN chapters/ch0_s7_2.tex
\subsection{进制环与射影极限}\hypertarget{heading:0:7.2}{}
\src{62—66}

\egapar{7.2.1}
离散环的任意射影极限显然都是线性拓扑化、分离且完备的环。
反之，设 $A$ 是线性拓扑化环，$(\frJ_\lambda)$ 是 $A$ 中由理想组成的
$0$ 的一个开基本邻域系。典范映射 $\varphi_\lambda:A\to A/\frJ_\lambda$
组成连续同态的射影系，因而定义连续同态
\[
 \varphi:A\longrightarrow\plim_\lambda A/\frJ_\lambda.
\]
若 $A$ 是分离的，则 $\varphi$ 是 $A$ 到 $\plim A/\frJ_\lambda$ 的一个稠密子环的
拓扑同构；若 $A$ 还完备，则 $\varphi$ 是 $A$ 到 $\plim A/\frJ_\lambda$ 的拓扑同构。

\egastatement{7.2.2}{引理}
线性拓扑化环为可容许环，当且仅当它同构于某个射影极限 $A=\plim A_\lambda$，
其中 $(A_\lambda,u_{\lambda\mu})$ 是离散环的射影系，其指标集是一个关于 $\leq$
滤过的有序集 $L$，具有记为 $0$ 的最小元，而且满足下列条件：
\begin{enumerate}[label=\arabic*°]
\item 各典范同态 $u_\lambda:A\to A_\lambda$ 都是满射；
\item 同态 $u_{0\lambda}:A_\lambda\to A_0$ 的核 $\frJ_\lambda$ 是幂零理想。
\end{enumerate}
此时，$u_0:A\to A_0$ 的核 $\frJ$ 等于 $\plim\frJ_\lambda$。

必要性由（\hyperlink{para:0:7.2.1}{7.2.1}）得到：取一个由定义理想组成的 $0$ 的基本邻域系
$(\frJ_\lambda)$，使这些理想都包含于其中某个 $\frJ_0$，再应用（\hyperlink{para:0:7.1.4}{7.1.4}，(i)）。
充分性由射影极限的定义和（\hyperlink{para:0:7.1.2}{7.1.2}）得到；最后的断言是直接的。

\egapar{7.2.3}
设 $A$ 是可容许拓扑环，$\frJ$ 是 $A$ 的一个包含于某个定义理想的理想
（依（\hyperlink{para:0:7.1.4}{7.1.4}），这等价于 $(\frJ^n)$ 趋于 $0$）。可以在 $A$ 上考虑以各次幂
$\frJ^n$（$n>0$）为 $0$ 的基本邻域系的环拓扑；仍称之为 $\frJ$-预进拓扑。
$A$ 可容许的假设蕴含 $\bigcap_n\frJ^n=0$，所以 $A$ 上的 $\frJ$-预进拓扑是分离的。
设
\[
 \comp A=\plim_n A/\frJ^n
\]
是 $A$ 关于这一拓扑的完备化，其中各 $A/\frJ^n$ 均赋以离散拓扑；以 $u$ 记环同态
$A\to\comp A$（它不一定连续），即同态序列 $u_n:A\to A/\frJ^n$ 的射影极限。
另一方面，$A$ 上的 $\frJ$-预进拓扑比给定的拓扑 $\mathscr T$ 更细；
由于 $A$ 关于 $\mathscr T$ 是分离且完备的，可以将从赋有 $\frJ$-预进拓扑的 $A$
到赋有 $\mathscr T$ 的 $A$ 的恒等映射作连续延拓，得到连续同态
$v:\comp A\to A$。

\NewNote{\textbf{作者勘误提示：}本命题已有作者正式勘误：须增加各理想幂在原拓扑中闭的条件，原证明不正确。此前译注中的待核标记应结合该勘误阅读。 参见\hyperlink{errIII:3}{《勘误与增补（清单二）》第3条}。}
\egastatement{7.2.4}{命题}
若 $A$ 是可容许环，且 $\frJ$ 包含于 $A$ 的某个定义理想，则 $A$ 关于
$\frJ$-预进拓扑是分离且完备的。

事实上，沿用（\hyperlink{para:0:7.2.3}{7.2.3}）的记号，显然 $v\circ u$ 是 $A$ 的恒等映射。
另一方面，$u_n\circ v:\comp A\to A/\frJ^n$ 是典范映射 $u_n$ 的连续延拓
（这里 $A$ 赋以 $\frJ$-预进拓扑，$A/\frJ^n$ 赋以离散拓扑）；
也就是说，它是从 $\comp A=\plim_k A/\frJ^k$ 到 $A/\frJ^n$ 的典范映射。
因此，$u\circ v$ 是这一映射序列的射影极限，按定义就是 $\comp A$ 的恒等映射；命题得证。
\transnote{上述陈述和论证依底本保留。原文在此将 $u_n\circ v$ 直接称为“连续延拓”，
但此前只说明了 $v$ 对 $A$ 的原拓扑连续，并未另证这一复合映射到离散商环的连续性；
不能仅凭较粗拓扑下的连续性略去这一步。本译稿不在此补写原文没有给出的论证。}

\NewNote{\textbf{作者勘误提示：}作者要求本推论同样增加理想幂闭的条件。 参见\hyperlink{errIII:3}{《勘误与增补（清单二）》第3条}。}
\egastatement{7.2.5}{推论}
在（\hyperlink{para:0:7.2.3}{7.2.3}）的假设下，下列条件等价：
\begin{enumerate}[label=\alph*)]
\item 同态 $u$ 连续；
\item 同态 $v$ 双连续；
\item $A$ 是 $\frJ$-进制环。
\end{enumerate}

\NewNote{\textbf{作者勘误提示：}作者要求本推论同样增加理想幂闭的条件。 参见\hyperlink{errIII:3}{《勘误与增补（清单二）》第3条}。}
\egastatement{7.2.6}{推论}
设 $A$ 是可容许环，$\frJ$ 是 $A$ 的定义理想。则 $A$ 为诺特环，当且仅当
$A/\frJ$ 是诺特环，而且 $\frJ/\frJ^2$ 是有限型 $(A/\frJ)$-模。

这些条件显然必要。反之，设它们成立。由（\hyperlink{para:0:7.2.4}{7.2.4}），$A$ 关于 $\frJ$-预进拓扑完备，
所以 $A$ 为诺特环，当且仅当关于滤过 $(\frJ^n)$ 的相伴分次环 $\grad(A)$ 是诺特环
（[1]，18-07 页，定理 4）。设 $a_1,\ldots,a_n$ 是 $\frJ$ 的元素，
其模 $\frJ^2$ 的类生成 $(A/\frJ)$-模 $\frJ/\frJ^2$。
用归纳法立即可见：关于 $a_i$（$1\leq i\leq n$）的总次数为 $m$ 的单项式，
其模 $\frJ^{m+1}$ 的类生成 $(A/\frJ)$-模 $\frJ^m/\frJ^{m+1}$。
于是 $\grad(A)$ 同构于 $(A/\frJ)[T_1,\ldots,T_n]$ 的一个商环，其中 $T_i$ 是不定元。
证明完成。

\egastatement{7.2.7}{命题}
设 $(A_i,u_{ij})$ 是离散环的射影系，$i\in\mathbb N$；对每个整数 $i$，
以 $\frJ_i$ 记同态 $u_{0i}:A_i\to A_0$ 在 $A_i$ 中的核。假设：
\begin{enumerate}[label=\alph*)]
\item 当 $i\leq j$ 时，$u_{ij}$ 是满射，且其核为 $\frJ_j^{i+1}$
（所以 $A_i$ 同构于 $A_j/\frJ_j^{i+1}$）；
\item $\frJ_1/\frJ_1^2\;(=\frJ_1)$ 是环 $A_0=A_1/\frJ_1$ 上的有限型模。
\end{enumerate}
令 $A=\plim_i A_i$；对每个整数 $n\geq0$，以 $u_n$ 记典范同态
$A\to A_n$，并以 $\frJ^{(n+1)}\subset A$ 记其核。则：
\begin{enumerate}[label=(\roman*)]
\item $A$ 是进制环，且 $\frJ=\frJ^{(1)}$ 是它的定义理想；
\item 对每个 $n\geq1$，有 $\frJ^{(n)}=\frJ^n$；
\item $\frJ/\frJ^2$ 同构于 $\frJ_1=\frJ_1/\frJ_1^2$，因而是环
$A_0=A/\frJ$ 上的有限型模。
\end{enumerate}

各 $u_{ij}$ 的满射性蕴含 $u_n$ 是满射；此外，假设 a) 蕴含 $\frJ_j^{j+1}=0$，
所以 $A$ 是可容许环（\hyperlink{para:0:7.2.2}{7.2.2}）。依定义，各 $\frJ^{(n)}$ 组成 $A$ 中 $0$ 的一个
基本邻域系，故 (ii) 蕴含 (i)。又有 $\frJ=\plim_i\frJ_i$，且映射
$\frJ\to\frJ_i$ 都是满射，因此 (ii) 还蕴含 (iii)。只须证明 (ii)。

依定义，$\frJ^{(n)}$ 由 $A$ 中满足 $x_k=0$（$k<n$）的元素
$(x_k)_{k\geq0}$ 组成，所以
\[
 \frJ^{(n)}\frJ^{(m)}\subset\frJ^{(n+m)};
\]
也就是说，各 $\frJ^{(n)}$ 构成 $A$ 的一个滤过。另一方面，
$\frJ^{(n)}/\frJ^{(n+1)}$ 同构于 $\frJ^{(n)}$ 在 $A_n$ 中的投影；
由 $\frJ^{(n)}=\plim_{i\geq n}\frJ_i^n$，这一投影就是 $\frJ_n^n$，
它是 $A_0=A_n/\frJ_n$ 上的模。

取 $\frJ=\frJ^{(1)}$ 的 $r$ 个元素 $a_i=(a_{ik})_{k\geq0}$，使
$a_{11},\ldots,a_{r1}$ 生成 $A_0$-模 $\frJ_1$。
我们将证明：关于 $a_i$ 的总次数为 $n$ 的单项式所成的集合 $S_n$，
生成 $A$ 的理想 $\frJ^{(n)}$。
由于 $\frJ_i^{i+1}=0$，首先显然有 $S_n\subset\frJ^{(n)}$；
又因 $A$ 关于滤过 $(\frJ^{(m)})$ 完备，只需证明 $S_n$ 中元素模
$\frJ^{(n+1)}$ 的类所成的集合 $\overline S_n$，在上述滤过下，
生成 $\grad(A)$ 上的分次模 $\grad(\frJ^{(n)})$
（[1]，18-06 页，引理）。
根据 $\grad(A)$ 的乘法定义，只须证明：对每个 $m$，$\overline S_m$
生成 $A_0$-模 $\frJ^{(m)}/\frJ^{(m+1)}$；等价地，$\frJ_m^m$
由关于 $a_{jm}$（$1\leq j\leq r$）的次数为 $m$ 的单项式生成。

为此，只须证明 $\frJ_m$ 作为 $A_m$-模，由关于 $a_{jm}$ 的次数不超过 $m$
的单项式生成。$m=1$ 时，命题依定义显然；现对 $m$ 作归纳，并令
$\frJ'_m$ 为这些单项式生成的 $\frJ_m$ 的 $A_m$-子模。
关系 $\frJ_{m-1}=\frJ_m/\frJ_m^m$ 和归纳假设给出
\[
 \frJ_m=\frJ'_m+\frJ_m^m.
\]
由于 $\frJ_m^{m+1}=0$，由此得到 $\frJ_m^m=\frJ_m'{}^m$，最终得到
$\frJ_m=\frJ'_m$。

\egastatement{7.2.8}{推论}
在（\hyperlink{para:0:7.2.7}{7.2.7}）的条件下，$A$ 是诺特环，当且仅当 $A_0$ 是诺特环。

这直接由（\hyperlink{para:0:7.2.6}{7.2.6}）得到。

\egastatement{7.2.9}{命题}
假设（\hyperlink{para:0:7.2.7}{7.2.7}）的条件成立。对每个整数 $i$，设 $M_i$ 是 $A_i$-模；
当 $i\leq j$ 时，设 $v_{ij}:M_j\to M_i$ 是 $u_{ij}$-同态，
使得 $(M_i,v_{ij})$ 构成射影系。再假设 $M_0$ 是有限型 $A_0$-模，
各 $v_{ij}$ 都是满射，而且其核为 $\frJ_j^{i+1}M_j$。
则 $M=\plim_i M_i$ 是有限型 $A$-模，且满射 $u_n$-同态
$v_n:M\to M_n$ 的核为 $\frJ^{n+1}M$
（因而 $M_n$ 可等同于 $M/\frJ^{n+1}M=M\otimes_A(A/\frJ^{n+1})$）。

取 $M$ 的 $s$ 个元素 $z_h=(z_{hk})_{k\geq0}$，使得 $z_{h0}$（$1\leq h\leq s$）
生成 $M_0$；我们将证明这些 $z_h$ 生成 $A$-模 $M$。
$A$-模 $M$ 关于由 $M^{(n)}$ 组成的滤过分离且完备，其中 $M^{(n)}$ 是
$M$ 中满足 $y_k=0$（$k<n$）的元素 $y=(y_k)_{k\geq0}$ 所成的集合。
显然有
\[
 \frJ^{(n)}M\subset M^{(n)},\qquad
 M^{(n)}/M^{(n+1)}=\frJ_n^nM_n.
\]
因此，只须证明 $z_h$ 模 $M^{(1)}$ 的类，
\transnote{底本印刷第 65 页此处作 $M^{(0)}$。按本段刚给出的定义，$M^{(0)}=M$；
相伴分次模的零次部分是 $M^{(0)}/M^{(1)}$。本译文据此将此处模数写为 $M^{(1)}$。}
在环 $\grad(A)$ 上生成关于上述滤过的分次模 $\grad(M)$
（[1]，18-06 页，引理）。为此，容易看出，只须再证明各 $z_{hn}$
（$1\leq h\leq s$）生成 $A_n$-模 $M_n$。

仍对 $n$ 作归纳；$n=0$ 时，命题依定义显然。关系
$M_{n-1}=M_n/\frJ_n^nM_n$ 和归纳假设表明：若 $M'_n$ 是 $z_{hn}$
生成的 $M_n$ 的子模，则
\[
 M_n=M'_n+\frJ_n^nM_n.
\]
由于 $\frJ_n$ 是幂零理想，这蕴含 $M_n=M'_n$。
同样转到相伴分次模的论证表明，$\frJ^{(n)}M$ 到 $M^{(n)}$ 的典范映射为满射
（因而为双射）；也就是说，$\frJ^{(n)}M=\frJ^nM$ 是
$M\to M_{n-1}$ 的核（$n\geq1$）。
\transnote{底本本段末句印作“$M\to M_n$ 的核”。按 $M^{(n)}$ 的定义，
应为 $M\to M_{n-1}$；相应地，$M\to M_n$ 的核为 $\frJ^{n+1}M$，与命题正文一致。}

\egastatement{7.2.10}{推论}
设 $(N_i,w_{ij})$ 是满足（\hyperlink{para:0:7.2.9}{7.2.9}）条件的另一个 $A_i$-模射影系，并令
$N=\plim N_i$。则 $A_i$-同态 $h_i:M_i\to N_i$ 的射影系 $(h_i)$
与 $A$-模同态 $h:M\to N$ 之间存在一一对应；后者对 $\frJ$-进拓扑必然连续。

显然，若 $h:M\to N$ 是 $A$-同态，则 $h(\frJ^nM)\subset\frJ^nN$，
由此得到 $h$ 的连续性。因而，通过取商，$h$ 给出一个 $A_i$-同态
$h_i:M_i\to N_i$ 的射影系，且 $h$ 正是其射影极限。推论得证。

\egastatement{7.2.11}{注记}
设 $A$ 是进制环，且有定义理想 $\frJ$，使得 $\frJ/\frJ^2$ 是有限型
$(A/\frJ)$-模。显然，环 $A_i=A/\frJ^{i+1}$ 满足（\hyperlink{para:0:7.2.7}{7.2.7}）的条件；
由于 $A$ 是各 $A_i$ 的射影极限，命题（\hyperlink{para:0:7.2.7}{7.2.7}）描述了所考虑类型的所有进制环
（特别是所有诺特进制环）。

\egastatement{7.2.12}{例}
设 $B$ 是环，$\frJ$ 是 $B$ 的一个理想，使得 $\frJ/\frJ^2$ 是有限型
$(B/\frJ)$-模（等价地，是有限型 $B$-模）；令
\[
 A=\plim_n B/\frJ^{n+1}.
\]
则 $A$ 是赋有 $\frJ$-预进拓扑的 $B$ 的分离完备化。
若 $A_n=B/\frJ^{n+1}$，则各 $A_n$ 显然满足（\hyperlink{para:0:7.2.7}{7.2.7}）的条件；
所以 $A$ 是进制环。以 $\overline{\frJ}$ 记 $\frJ$ 的典范像在 $A$ 中的闭包，
则 $\overline{\frJ}$ 是 $A$ 的定义理想，$\overline{\frJ}{}^n$ 是
$\frJ^n$ 的典范像的闭包，$A/\overline{\frJ}{}^n$ 可等同于 $B/\frJ^n$，
而 $\overline{\frJ}/\overline{\frJ}{}^2$ 与 $\frJ/\frJ^2$
作为 $(A/\overline{\frJ})$-模同构。

同样，若 $N$ 满足 $N/\frJ N$ 是有限型 $B$-模，并令
$M_i=N/\frJ^{i+1}N$，则 $M=\plim M_i$ 是有限型 $A$-模，
同构于 $N$ 关于 $\frJ$-预进拓扑的分离完备化；
$\overline{\frJ}{}^nM$ 可等同于 $\frJ^nN$ 的典范像的闭包，
而 $M/\overline{\frJ}{}^nM$ 可等同于 $N/\frJ^nN$。


% END chapters/ch0_s7_2.tex

\par

% BEGIN chapters/ch0_s7_3.tex
\subsection{诺特预进制环}\hypertarget{heading:0:7.3}{}
\src{66—68}

\egapar{7.3.1}
设 $A$ 是环，$\frJ$ 是 $A$ 的理想，$M$ 是 $A$-模。以
\[
 \comp A=\plim_n A/\frJ^n,\qquad
 \comp M=\plim_n M/\frJ^nM
\]
分别记 $A$ 和 $M$ 关于 $\frJ$-预进拓扑的分离完备化。
设 $M'\xrightarrow{u}M\xrightarrow{v}M''\to0$ 是 $A$-模正合列；
由 $M/\frJ^nM=M\otimes_A(A/\frJ^n)$，对每个 $n$，序列
\[
 M'/\frJ^nM'\xrightarrow{u_n}M/\frJ^nM
 \xrightarrow{v_n}M''/\frJ^nM''\longrightarrow0
\]
都是正合的。又由于
$v(\frJ^nM)=\frJ^nv(M)=\frJ^nM''$，
$\comp v=\plim v_n$ 为满射（Bourbaki，\textit{Top. gén.}，第 IX 章，
第 2 版，60 页，推论 2）。
另一方面，若 $z=(z_k)$ 属于 $\comp v$ 的核，则对每个整数 $k$，
存在 $z'_k\in M'/\frJ^kM'$ 使得 $u_k(z'_k)=z_k$。
由此可知存在 $z'=(z'_n)\in\comp{M'}$，使得 $\comp u(z')$ 的前 $k$ 个分量
与 $z$ 的相应分量相同；换言之，$\comp{M'}$ 在 $\comp u$ 下的像在
$\Ker\comp v$ 中稠密。

若再设 $A$ 是诺特环，则由（\hyperlink{para:0:7.2.12}{7.2.12}），$\comp A$ 也是诺特环，
因为此时 $\frJ/\frJ^2$ 是有限型 $A$-模。此外还有下述定理。

\egastatement{7.3.2}{Krull 定理}
设 $A$ 是诺特环，$\frJ$ 是 $A$ 的理想，$M$ 是有限型 $A$-模，
$M'$ 是 $M$ 的子模。则 $M$ 的 $\frJ$-预进拓扑在 $M'$ 上诱导的拓扑，
就是 $M'$ 的 $\frJ$-预进拓扑。

这直接由下述引理得到。

\egastatement{7.3.2.1}{Artin--Rees 引理}
在（\hyperlink{para:0:7.3.2}{7.3.2}）的假设下，存在整数 $p$，使得当 $n\geq p$ 时，有
\[
 M'\cap\frJ^nM=\frJ^{n-p}(M'\cap\frJ^pM).
\]
证明见 [1]，2-04 页。

\egastatement{7.3.3}{推论}
在（\hyperlink{para:0:7.3.2}{7.3.2}）的假设下，典范映射 $M\otimes_A\comp A\to\comp M$ 是双射，
且在有限型 $A$-模范畴中，$M\otimes_A\comp A$ 作为 $M$ 的函子是正合的；
因而，$\frJ$-进分离完备化 $\comp A$ 是平坦 $A$-模（\hyperlink{para:0:6.1.1}{6.1.1}）。

首先注意，在有限型 $A$-模范畴中，$\comp M$ 是关于 $M$ 的正合函子。
事实上，设
$0\to M'\xrightarrow{u}M\xrightarrow{v}M''\to0$ 是正合列。
已知 $\comp v:\comp M\to\comp{M''}$ 为满射（\hyperlink{para:0:7.3.1}{7.3.1}）；
另一方面，若 $i:M\to\comp M$ 是典范同态，则由 Krull 定理，
$i(u(M'))$ 在 $\comp M$ 中的闭包可等同于 $M'$ 关于 $\frJ$-预进拓扑的
分离完备化。因此 $\comp u$ 为单射；再由（\hyperlink{para:0:7.3.1}{7.3.1}），$\comp u$ 的像就是
$\comp v$ 的核。

典范映射 $M\otimes_A\comp A\to\comp M$ 可由映射
\[
 M\otimes_A\comp A\longrightarrow
 M\otimes_A(A/\frJ^n)=M/\frJ^nM
\]
取射影极限而得到。当 $M$ 具有 $A^p$ 的形式时，这个映射显然是双射。
若 $M$ 是有限型 $A$-模，则存在正合列 $A^p\to A^q\to M\to0$。
利用有限型 $A$-模范畴中函子 $M\otimes_A\comp A$ 和 $\comp M$ 关于 $M$
的右正合性，得到交换图
\[
\begin{tikzcd}[column sep=large,row sep=large]
 A^p\otimes_A\comp A \arrow[r] \arrow[d] &
 A^q\otimes_A\comp A \arrow[r] \arrow[d] &
 M\otimes_A\comp A \arrow[r] \arrow[d] &0\\
 \comp A^p \arrow[r] &\comp A^q \arrow[r] &\comp M\arrow[r]&0.
\end{tikzcd}
\]
两行均正合，且前两个竖直箭头均为同构，立即得到结论。

\egastatement{7.3.4}{推论}
设 $A$ 是诺特环，$\frJ$ 是 $A$ 的理想，$M,N$ 是两个有限型 $A$-模。
存在典范的、函子性的同构
\[
 \comp{M\otimes_A N}\xrightarrow{\sim}
 \comp M\otimes_{\comp A}\comp N,
 \qquad
 \comp{\Hom_A(M,N)}\xrightarrow{\sim}
 \Hom_{\comp A}(\comp M,\comp N).
\]
这由（\hyperlink{para:0:7.3.3}{7.3.3}）、（\hyperlink{para:0:6.2.1}{6.2.1}）和（\hyperlink{para:0:6.2.2}{6.2.2}）得到。

\egastatement{7.3.5}{推论}
设 $A$ 是诺特环，$\frJ$ 是 $A$ 的理想。下列条件等价：
\begin{enumerate}[label=\alph*)]
\item $\frJ$ 包含于 $A$ 的根基中；
\item $\comp A$ 是忠实平坦 $A$-模（\hyperlink{para:0:6.4.1}{6.4.1}）；
\item 每个有限型 $A$-模关于 $\frJ$-预进拓扑都是分离的；
\item 每个有限型 $A$-模的每个子模关于 $\frJ$-预进拓扑都是闭的。
\end{enumerate}

由于 $\comp A$ 是平坦 $A$-模，条件 b) 与 c) 等价：b) 等价于，对于每个有限型
$A$-模 $M$，典范映射 $M\to\comp M=M\otimes_A\comp A$ 是单射（\hyperlink{para:0:6.6.1}{6.6.1}，c)）。
c) 显然蕴含 d)，因为若 $N$ 是有限型 $A$-模 $M$ 的子模，则 $M/N$ 关于
$\frJ$-预进拓扑分离，故 $N$ 在 $M$ 中闭。

下面证明 d) 蕴含 a)。若 $\mathfrak m$ 是 $A$ 的极大理想，则
$\mathfrak m$ 关于 $\frJ$-预进拓扑在 $A$ 中闭，因此
\[
 \mathfrak m=\bigcap_{p\geq0}(\mathfrak m+\frJ^p).
\]
由于 $\mathfrak m+\frJ^p$ 必等于 $A$ 或 $\mathfrak m$，对充分大的 $p$，
有 $\mathfrak m+\frJ^p=\mathfrak m$，因而 $\frJ^p\subset\mathfrak m$；
由 $\mathfrak m$ 为素理想，得 $\frJ\subset\mathfrak m$。

最后，a) 蕴含 b)。事实上，设 $P$ 是有限型 $A$-模 $M$ 中的 $\{0\}$ 关于
$\frJ$-预进拓扑的闭包。由 Krull 定理（\hyperlink{para:0:7.3.2}{7.3.2}），$M$ 的 $\frJ$-预进拓扑在
$P$ 上诱导的拓扑就是 $P$ 的 $\frJ$-预进拓扑，所以 $\frJ P=P$。
$P$ 是有限型的，且 $\frJ$ 包含于 $A$ 的根基中；由 Nakayama 引理，得 $P=0$。

注意，当 $A$ 是诺特局部环、$\frJ\ne A$ 是 $A$ 的任意理想时，
（\hyperlink{para:0:7.3.5}{7.3.5}）的各个条件均成立。

\egastatement{7.3.6}{推论}
若 $A$ 是诺特 $\frJ$-进制环，则每个有限型 $A$-模关于 $\frJ$-预进拓扑
都分离且完备。

由于此时 $\comp A=A$，结论立即由（\hyperlink{para:0:7.3.3}{7.3.3}）得到。
因此，命题（\hyperlink{para:0:7.2.9}{7.2.9}）描述了诺特进制环上的所有有限型模。

\egastatement{7.3.7}{推论}
在（\hyperlink{para:0:7.3.2}{7.3.2}）的假设下，典范映射 $M\to\comp M=M\otimes_A\comp A$ 的核，
是被 $1+\frJ$ 中某个元素零化的所有 $x\in M$ 所成的集合。

事实上，$x\in M$ 属于这个核，当且仅当子模 $Ax$ 的分离完备化为零
（由（\hyperlink{para:0:7.3.2}{7.3.2}））；换言之，当且仅当 $x\in\frJ x$。


% END chapters/ch0_s7_3.tex

\par

% BEGIN chapters/ch0_s7_4.tex
\subsection{局部环上的拟有限模}\hypertarget{heading:0:7.4}{}
\src{68—69}

\egastatement{7.4.1}{定义}
给定局部环 $A$，其极大理想为 $\mathfrak m$。若 $M/\mathfrak mM$ 在剩余域
$k=A/\mathfrak m$ 上具有有限秩，则称 $A$-模 $M$ 是（在 $A$ 上）\emph{拟有限的}。

若 $A$ 是诺特环，则 $M$ 关于 $\mathfrak m$-预进拓扑的分离完备化
$\comp M$ 是有限型 $\comp A$-模。事实上，此时 $\mathfrak m/\mathfrak m^2$
是有限型 $A$-模，故结论由（\hyperlink{para:0:7.2.12}{7.2.12}）以及关于 $M/\mathfrak mM$ 的假设得到。

特别地，若还假设 $A$ 完备，且 $M$ 关于 $\mathfrak m$-预进拓扑分离
（也就是说，$\bigcap_n\mathfrak m^nM=0$），则 $M$ 本身就是有限型 $A$-模。
事实上，此时 $\comp M$ 是有限型 $A$-模，而 $M$ 可等同于 $\comp M$ 的一个子模，
故 $M$ 也是有限型的（而且由（\hyperlink{para:0:7.3.6}{7.3.6}），它与其完备化相同）。

\egastatement{7.4.2}{命题}
设 $A,B$ 是两个局部环，$\mathfrak m,\mathfrak n$ 分别为其极大理想，
并假设 $B$ 是诺特环。设 $\varphi:A\to B$ 是局部同态，$M$ 是有限型 $B$-模。
若 $M$ 是拟有限 $A$-模，则 $M$ 上的 $\mathfrak m$-预进拓扑与
$\mathfrak n$-预进拓扑相同，因而都是分离的。

注意，依假设，$M/\mathfrak mM$ 作为 $A$-模具有有限长度，
所以作为 $B$-模也具有有限长度。由此可知，$\mathfrak n$ 是 $B$ 中唯一包含
$M/\mathfrak mM$ 的零化理想的素理想：事实上，由（\hyperlink{para:0:1.7.4}{1.7.4}）和（\hyperlink{para:0:1.7.2}{1.7.2}），
立即归结到 $M/\mathfrak mM$ 为单模的情形；此时它必同构于 $B/\mathfrak n$，
断言显然成立。\transnote{这一证明段落按非零情形叙述；若 $M=0$，命题关于两种拓扑的结论直接成立，
但“唯一包含零化理想的素理想”这一中间说法不适用于零模。}
另一方面，由于 $M$ 是有限型 $B$-模，包含 $M/\mathfrak mM$ 的零化理想的素理想，
恰是包含 $\mathfrak mB+\mathfrak b$ 的素理想，其中 $\mathfrak b$ 是
$B$-模 $M$ 的零化理想（\hyperlink{para:0:1.7.5}{1.7.5}）。
$B$ 为诺特环，故由 [11]，127 页，推论 4，可知 $\mathfrak mB+\mathfrak b$
是 $B$ 的定义理想；也就是说，存在 $k>0$ 使得
\[
 \mathfrak n^k\subset\mathfrak mB+\mathfrak b\subset\mathfrak n.
\]
从而，对每个 $h>0$，有
\[
 \mathfrak n^{kh}M\subset(\mathfrak mB+\mathfrak b)^hM
 =\mathfrak m^hM\subset\mathfrak n^hM.
\]
这证明了 $M$ 上的 $\mathfrak m$-预进拓扑与 $\mathfrak n$-预进拓扑相同；
又由（\hyperlink{para:0:7.3.5}{7.3.5}），后者是分离的。

\egastatement{7.4.3}{推论}
在（\hyperlink{para:0:7.4.2}{7.4.2}）的假设下，若 $A$ 还是诺特环，并且关于 $\mathfrak m$-预进拓扑完备，
则 $M$ 是有限型 $A$-模。

事实上，此时 $M$ 关于 $\mathfrak m$-预进拓扑分离，结论由（\hyperlink{para:0:7.4.1}{7.4.1}）中的说明得到。

\egapar{7.4.4}
（\hyperlink{para:0:7.4.2}{7.4.2}）最重要的应用情形是 $B$ 本身为拟有限 $A$-模。
这等价于说，$B/\mathfrak mB$ 是域 $k=A/\mathfrak m$ 上的有限秩代数；
根据上文，这个条件又可以分解为以下两个条件的合取：
\begin{enumerate}[label=(\roman*)]
\item $\mathfrak mB$ 是 $B$ 的定义理想；
\item $B/\mathfrak n$ 是域 $A/\mathfrak m$ 的有限秩扩张。
\end{enumerate}
在这种情形下，每个有限型 $B$-模显然都是拟有限 $A$-模。

\egastatement{7.4.5}{推论}
在（\hyperlink{para:0:7.4.2}{7.4.2}）的假设下，若 $\mathfrak b$ 是 $B$-模 $M$ 的零化理想，
则 $B/\mathfrak b$ 是拟有限 $A$-模。

设 $M\ne0$（否则推论显然）。可以将 $M$ 看成诺特局部环 $B/\mathfrak b$
上的模；它的零化理想此时为零，所以（\hyperlink{para:0:7.4.2}{7.4.2}）的证明表明
$\mathfrak m(B/\mathfrak b)$ 是 $B/\mathfrak b$ 的定义理想。
此外，$M/\mathfrak nM$ 是域 $A/\mathfrak m$ 上的有限秩向量空间，因为它是
$M/\mathfrak mM$ 的商，而后者依假设在 $A/\mathfrak m$ 上具有有限秩。
由 $M\ne0$ 及 Nakayama 引理，有 $M\ne\mathfrak nM$。
由于 $M/\mathfrak nM$ 是 $B/\mathfrak n$ 上的非零向量空间，
它在 $A/\mathfrak m$ 上具有有限秩，便蕴含 $B/\mathfrak n$ 在
$A/\mathfrak m$ 上也具有有限秩。对环 $B/\mathfrak b$ 应用（\hyperlink{para:0:7.4.4}{7.4.4}），即得结论。


% END chapters/ch0_s7_4.tex

\par

% BEGIN chapters/ch0_s7_5.tex
\subsection{限制形式幂级数环}\hypertarget{heading:0:7.5}{}
\src{69—72}

\egapar{7.5.1}
设 $A$ 是线性拓扑化、分离且完备的拓扑环；设 $(\frJ_\lambda)$ 是 $A$ 中由
（开）理想组成的 $0$ 的一个基本邻域系，从而 $A$ 可典范地等同于
$\plim A/\frJ_\lambda$（\hyperlink{para:0:7.2.1}{7.2.1}）。对每个 $\lambda$，令
$B_\lambda=(A/\frJ_\lambda)[T_1,\ldots,T_r]$，其中 $T_j$ 是不定元；
各 $B_\lambda$ 显然组成离散环的射影系。记
\[
 \plim_\lambda B_\lambda=A\{T_1,\ldots,T_r\}.
\]
我们将证明，这个拓扑环与所选的理想基本系 $(\frJ_\lambda)$ 无关。
更确切地说，设 $A'$ 是形式幂级数环 $A[[T_1,\ldots,T_r]]$ 中由下列级数组成的子环：
\[
 \sum_{\alpha}c_\alpha T^\alpha,
 \qquad \alpha=(\alpha_1,\ldots,\alpha_r)\in\mathbb N^r,
 \qquad \lim c_\alpha=0,
\]
这里的极限是沿 $\mathbb N^r$ 的有限子集之补集所组成的滤子取的。
这样的级数称为以 $T_j$ 为变量、以 $A$ 中元素为系数的\emph{限制形式幂级数}。

对 $A$ 中 $0$ 的每个邻域 $V$，令 $V'$ 为所有满足每个 $c_\alpha\in V$ 的
$x=\sum_\alpha c_\alpha T^\alpha\in A'$ 所成的集合。
立即可验证，各 $V'$ 组成 $0$ 的一个基本邻域系，在 $A'$ 上定义一个分离的环拓扑。
下面典范地定义一个从环 $A\{T_1,\ldots,T_r\}$ 到 $A'$ 的拓扑同构。

对每个 $\alpha\in\mathbb N^r$ 及每个 $\lambda$，设 $\varphi_{\lambda,\alpha}$ 是从
$(A/\frJ_\lambda)[T_1,\ldots,T_r]$ 到 $A/\frJ_\lambda$ 的取系数映射，
即将多项式映为其中 $T^\alpha$ 的系数。显然，各 $\varphi_{\lambda,\alpha}$
组成 $(A/\frJ_\lambda)$-模同态的射影系，其射影极限为连续同态
\[
 \varphi_\alpha:A\{T_1,\ldots,T_r\}\longrightarrow A.
\]
我们将证明，对于每个 $y\in A\{T_1,\ldots,T_r\}$，形式幂级数
$\sum_\alpha\varphi_\alpha(y)T^\alpha$ 是限制的。
事实上，设 $y_\lambda$ 是 $y$ 在 $B_\lambda$ 中的分量，$H_\lambda$ 是
多项式 $y_\lambda$ 的非零系数所对应的 $\alpha\in\mathbb N^r$ 的有限集合。
若 $\frJ_\mu\subset\frJ_\lambda$ 且 $\alpha\notin H_\lambda$，则
$y_\mu$ 的 $T^\alpha$ 系数在 $A/\frJ_\lambda$ 中的像为零；
取极限，得到 $\varphi_\alpha(y)\in\frJ_\lambda$（$\alpha\notin H_\lambda$）。
因此，可以令
\[
 \varphi(y)=\sum_\alpha\varphi_\alpha(y)T^\alpha
\]
来定义环同态 $\varphi:A\{T_1,\ldots,T_r\}\to A'$，且 $\varphi$ 显然连续。

反之，设 $\theta_\lambda:A\to A/\frJ_\lambda$ 是典范同态。
对每个 $z=\sum_\alpha c_\alpha T^\alpha\in A'$ 和每个 $\lambda$，
只有有限多个指标 $\alpha$ 满足 $\theta_\lambda(c_\alpha)\ne0$，因此
\[
 \psi_\lambda(z)=\sum_\alpha\theta_\lambda(c_\alpha)T^\alpha
\]
属于 $B_\lambda$。各 $\psi_\lambda$ 连续，并组成同态的射影系；其射影极限为连续同态
$\psi:A'\to A\{T_1,\ldots,T_r\}$。最后，直接验证
$\varphi\circ\psi$ 和 $\psi\circ\varphi$ 均为恒等自同构，即告完成。

\egapar{7.5.2}
通过（\hyperlink{para:0:7.5.1}{7.5.1}）中定义的同构，将 $A\{T_1,\ldots,T_r\}$ 与限制形式幂级数环
$A'$ 等同。典范同构
\[
 \bigl((A/\frJ_\lambda)[T_1,\ldots,T_r]\bigr)[T_{r+1},\ldots,T_s]
 \xrightarrow{\sim}(A/\frJ_\lambda)[T_1,\ldots,T_s]
\]
通过取射影极限，定义典范同构
\[
 \bigl(A\{T_1,\ldots,T_r\}\bigr)\{T_{r+1},\ldots,T_s\}
 \xrightarrow{\sim} A\{T_1,\ldots,T_s\}.
\]

\egapar{7.5.3}
对每个从 $A$ 到线性拓扑化、分离且完备的环 $B$ 的连续同态 $u:A\to B$，
以及 $B$ 的任意 $r$ 个元素组成的系 $(b_1,\ldots,b_r)$，
存在唯一连续同态
\[
 \overline u:A\{T_1,\ldots,T_r\}\longrightarrow B,
\]
使得对每个 $a\in A$ 有 $\overline u(a)=u(a)$，且对 $1\leq j\leq r$
有 $\overline u(T_j)=b_j$。事实上，只须令
\[
 \overline u\left(\sum_\alpha c_\alpha T^\alpha\right)
 =\sum_\alpha u(c_\alpha)b_1^{\alpha_1}\cdots b_r^{\alpha_r}.
\]
族 $(u(c_\alpha)b_1^{\alpha_1}\cdots b_r^{\alpha_r})$ 在 $B$ 中可求和，
以及 $\overline u$ 的连续性，都可直接验证，留给读者。
注意，这一性质（对于任意 $B$ 及任意 $b_j$）在唯一同构的意义下刻画了
拓扑环 $A\{T_1,\ldots,T_r\}$。

\egastatement{7.5.4}{命题}
\begin{enumerate}[label=(\roman*)]
\item 若 $A$ 是可容许环，则 $A'=A\{T_1,\ldots,T_r\}$ 也是可容许环。
\item 设 $A$ 是进制环，$\frJ$ 是 $A$ 的定义理想，且 $\frJ/\frJ^2$ 是有限型
$(A/\frJ)$-模。令 $\frJ'=\frJ A'$，则 $A'$ 是 $\frJ'$-进制环，
而且 $\frJ'/\frJ'{}^2$ 是有限型 $(A'/\frJ')$-模。
若 $A$ 还是诺特环，则 $A'$ 也是诺特环。
\end{enumerate}

(i) 若 $\frJ$ 是 $A$ 的理想，而 $\frJ'$ 是 $A'$ 中所有满足
$c_\alpha\in\frJ$ 的 $\sum_\alpha c_\alpha T^\alpha$ 所成的理想，
则 $(\frJ')^n\subset(\frJ^n)'$。因此，若 $\frJ$ 是 $A$ 的定义理想，
则 $\frJ'$ 是 $A'$ 的定义理想。

(ii) 令 $A_i=A/\frJ^{i+1}$；当 $i\leq j$ 时，以 $u_{ij}$ 记典范同态
$A/\frJ^{j+1}\to A/\frJ^{i+1}$。
令 $A'_i=A_i[T_1,\ldots,T_r]$，并以 $u'_{ij}:A'_j\to A'_i$ 记将
$u_{ij}$ 作用于 $A'_j$ 的多项式系数所得到的同态。
我们证明射影系 $(A'_i,u'_{ij})$ 满足（\hyperlink{para:0:7.2.7}{7.2.7}）的条件；由于 $\frJ'$ 是
$A'\to A'_0$ 的核，这将证明 (ii) 的第一个断言。

显然，各 $u'_{ij}$ 都是满射。$u'_{0i}$ 的核 $\frJ'_i$ 由
$A_i[T_1,\ldots,T_r]$ 中所有系数属于 $\frJ/\frJ^{i+1}$ 的多项式组成；
特别地，$\frJ'_1$ 由 $A_1[T_1,\ldots,T_r]$ 中所有系数属于
$\frJ/\frJ^2$ 的多项式组成。由于 $\frJ/\frJ^2$ 是有限型
$A_1=A/\frJ^2$-模，可见 $\frJ'_1/\frJ_1'{}^2$ 是有限型 $A'_1$-模
（等价地，是有限型 $A'_0=A'_1/\frJ'_1$-模）。

下面证明 $u'_{ij}$ 的核为 $\frJ_j'{}^{i+1}$。
显然 $\frJ_j'{}^{i+1}$ 包含于该核。另一方面，取 $\frJ$ 的元素
$a_1,\ldots,a_m$，使其模 $\frJ^2$ 的类生成 $\frJ/\frJ^2$。
立即可验证：关于 $a_k$（$1\leq k\leq m$）的次数不超过 $j$ 的单项式，
其模 $\frJ^{j+1}$ 的类生成 $\frJ/\frJ^{j+1}$；
次数大于 $i$ 且不超过 $j$ 的单项式的类，因而生成
$\frJ^{i+1}/\frJ^{j+1}$。
以这样的元素为系数的、关于 $T_k$ 的单项式，是
$\frJ'_j$ 中 $i+1$ 个元素的乘积，由此得到断言。
最后，若 $A$ 是诺特环，则
$A'/\frJ'=(A/\frJ)[T_1,\ldots,T_r]$ 是诺特环；
所以 $A'$ 是诺特环（\hyperlink{para:0:7.2.8}{7.2.8}）。

\egastatement{7.5.5}{命题}
设 $A$ 是诺特 $\frJ$-进制环，$B$ 是可容许拓扑环，
$\varphi:A\to B$ 是连续同态，赋予 $B$ 一个 $A$-代数结构。下列条件等价：
\begin{enumerate}[label=\alph*)]
\item $B$ 是诺特 $\frJ B$-进制环，而且 $B/\frJ B$ 是有限型 $(A/\frJ)$-代数。
\item $B$ 在拓扑 $A$-代数意义下同构于 $\plim B_n$，其中当 $m\geq n$ 时，
有 $B_n=B_m/\frJ^{n+1}B_m$，而 $B_1$ 是有限型 $A_1=A/\frJ^2$-代数。
\item $B$ 在拓扑 $A$-代数意义下同构于某个形如 $A\{T_1,\ldots,T_r\}$ 的代数
关于一个理想的商（由（\hyperlink{para:0:7.3.6}{7.3.6}）及（\hyperlink{para:0:7.5.4}{7.5.4}，(ii)），这个理想必然闭）。
\end{enumerate}

由于 $A$ 是诺特环，$A'=A\{T_1,\ldots,T_r\}$ 也是诺特环（\hyperlink{para:0:7.5.4}{7.5.4}），
所以 c) 蕴含 $B$ 为诺特环。由于 $\frJ'=\frJ A'$ 是 $A'$ 中 $0$ 的开邻域，
且各 $\frJ'{}^n$ 组成 $0$ 的一个基本邻域系，它们的像 $\frJ^nB$
组成 $B$ 中 $0$ 的一个基本邻域系。已知 $B$ 分离且完备，所以 $B$ 是
$\frJ B$-进制环。最后，$B/\frJ B$ 作为 $(A/\frJ)$-代数，是
$A'/\frJ A'=(A/\frJ)[T_1,\ldots,T_r]$ 的商，因而是有限型的。
这就证明了 c) 蕴含 a)。

若 $B$ 是诺特 $\frJ B$-进制环，则 $B$ 同构于 $\plim B_n$，其中
$B_n=B/\frJ^{n+1}B$（\hyperlink{para:0:7.2.11}{7.2.11}），而 $\frJ B/\frJ^2B$ 是有限型
$(B/\frJ B)$-模。设 $(a_j)_{1\leq j\leq s}$ 是这个模的一个生成系，
$(c_i)_{1\leq i\leq r}$ 是 $B/\frJ^2B$ 的一组元素，
其模 $\frJ B/\frJ^2B$ 的类生成 $(A/\frJ)$-代数 $B/\frJ B$。
立即可见，各 $c_i$ 生成 $(A/\frJ^2)$-代数 $B/\frJ^2B$，
所以 a) 蕴含 b)。

只须再证明 b) 蕴含 c)。由假设，$B_1$ 是诺特环；
又因 $B_1/\frJ B_1=B_0$，有 $\frJ^2B_1=0$，故
$\frJ_1=\frJ B_1=\frJ B/\frJ^2B$ 是有限型 $B_1$-模。
因此，射影系 $(B_n)$ 满足（\hyperlink{para:0:7.2.7}{7.2.7}）的条件，且 $B$ 是 $\frJ B$-进制环。

取 $B$ 的有限个元素 $(c_i)_{1\leq i\leq r}$，使其模 $\frJ B$ 的类生成
$(A/\frJ)$-代数 $B/\frJ B$；并且，使这些元素以 $\frJ$ 中元素为系数的
线性组合，其模 $\frJ^2B$ 的类生成 $B_1$-模 $\frJ B/\frJ^2B$。
存在连续 $A$-同态 $u:A'=A\{T_1,\ldots,T_r\}\to B$，
它在 $A$ 上限制为 $\varphi$，且对 $1\leq i\leq r$ 有 $u(T_i)=c_i$（\hyperlink{para:0:7.5.3}{7.5.3}）。
若证明 $u$ 为满射，即得到 c)：因为由 $u(A')=B$ 得到
$u(\frJ^nA')=\frJ^nB$，也就是说，$u$ 是拓扑环的严格态射，
所以 $B$ 同构于 $A'$ 关于某个闭理想的商。
\transnote{底本印刷第 72 页在这两处写作 $u(A)=B$ 和 $u(\frJ^nA)=\frJ^nB$。
按本段同态的定义域，应为 $A'$；本译文补回这两个撇号。}

由于 $B$ 关于 $\frJ B$-进拓扑完备，只须证明（[1]，18-07 页）：
由 $u$ 关于 $A'$ 与 $B$ 上的 $\frJ$-进滤过典范诱导的同态
$\grad(A')\to\grad(B)$ 是满射。
依定义，由 $u$ 诱导的同态
\[
 A'/\frJ A'\longrightarrow B/\frJ B,
 \qquad
 \frJ A'/\frJ^2A'\longrightarrow\frJ B/\frJ^2B
\]
都是满射。对 $n$ 作归纳，立即得到
\[
 \frJ A'/\frJ^nA'\longrightarrow\frJ B/\frJ^nB
\]
也是满射，进而
\[
 \frJ^nA'/\frJ^{n+1}A'\longrightarrow
 \frJ^nB/\frJ^{n+1}B
\]
也是满射。证明完成。


% END chapters/ch0_s7_5.tex

\par

% BEGIN chapters/ch0_s7_6.tex
\subsection{完备分式环}\hypertarget{heading:0:7.6}{}
\src{72—75}

\egapar{7.6.1}
设 $A$ 是线性拓扑化环，$(\frJ_\lambda)$ 是 $A$ 中由理想组成的 $0$ 的一个基本邻域系，
$S$ 是 $A$ 的乘法子集。设 $u_\lambda:A\to A_\lambda=A/\frJ_\lambda$ 是典范同态；
当 $\frJ_\mu\subset\frJ_\lambda$ 时，设 $u_{\lambda\mu}:A_\mu\to A_\lambda$
是典范同态。令 $S_\lambda=u_\lambda(S)$，从而
$u_{\lambda\mu}(S_\mu)=S_\lambda$。
各 $u_{\lambda\mu}$ 典范地给出满同态
$S_\mu^{-1}A_\mu\to S_\lambda^{-1}A_\lambda$，这些环关于上述同态构成射影系；
以 $A\{S^{-1}\}$ 记该系的射影极限。
这个定义与所选的基本邻域系 $(\frJ_\lambda)$ 无关；事实上，有下述命题。

\egastatement{7.6.2}{命题}
环 $A\{S^{-1}\}$ 在拓扑意义下同构于环 $S^{-1}A$ 的分离完备化，
这里 $S^{-1}A$ 所赋的拓扑以各 $S^{-1}\frJ_\lambda$ 为 $0$ 的一个基本邻域系。

事实上，若 $v_\lambda:S^{-1}A\to S_\lambda^{-1}A_\lambda$ 是由
$u_\lambda$ 诱导的典范同态，则 $v_\lambda$ 的核为 $S^{-1}\frJ_\lambda$，
且 $v_\lambda$ 是满射。由（\hyperlink{para:0:7.2.1}{7.2.1}）即得命题。

\egastatement{7.6.3}{推论}
若 $S'$ 是 $S$ 在 $A$ 的分离完备化 $\comp A$ 中的典范像，
则 $A\{S^{-1}\}$ 可典范地等同于 $\comp A\{S'{}^{-1}\}$。

注意，即使 $A$ 分离且完备，$S^{-1}A$ 关于各 $S^{-1}\frJ_\lambda$ 所定义的拓扑
也不一定分离且完备。例如，取 $S$ 为各 $f^n$（$n\geq0$）组成的集合，
其中 $f$ 是拓扑幂零元，但不是幂零元。此时 $S^{-1}A$ 非零；
另一方面，对每个 $\lambda$，存在 $n$ 使得 $f^n\in\frJ_\lambda$，
从而 $1=f^n/f^n\in S^{-1}\frJ_\lambda$，故
$S^{-1}\frJ_\lambda=S^{-1}A$。

\egastatement{7.6.4}{推论}
若在 $A$ 中有 $0\notin\overline S$，则环 $A\{S^{-1}\}$ 非零。

事实上，在环 $S^{-1}A$ 中，$0$ 不属于 $\{1\}$ 的闭包；
否则，对 $A$ 的每个开理想 $\frJ_\lambda$，都有 $1\in S^{-1}\frJ_\lambda$，
从而对每个 $\lambda$ 都有 $\frJ_\lambda\cap S\ne\varnothing$，与假设矛盾。

\egapar{7.6.5}
称 $A\{S^{-1}\}$ 为 $A$ 的\emph{以 $S$ 中元素为分母的完备分式环}。
沿用上述记号，显然 $S^{-1}\frJ_\lambda$ 在 $A$ 中的逆像包含
$\frJ_\lambda$，所以典范映射 $A\to S^{-1}A$ 连续。
将它与典范映射 $S^{-1}A\to A\{S^{-1}\}$ 复合，得到典范连续同态
\[
 A\longrightarrow A\{S^{-1}\},
\]
它是各同态 $A\to S_\lambda^{-1}A_\lambda$ 的射影极限。

\egapar{7.6.6}
由 $A\{S^{-1}\}$ 与典范映射 $A\to A\{S^{-1}\}$ 组成的对，
由下列泛性质刻画：对于每个从 $A$ 到线性拓扑化、分离且完备的环 $B$ 的连续同态
$u$，若 $u(S)$ 由 $B$ 中的可逆元组成，则 $u$ 唯一分解为
\[
 A\longrightarrow A\{S^{-1}\}\xrightarrow{u'}B,
\]
其中 $u'$ 连续。

事实上，$u$ 唯一分解为 $A\to S^{-1}A\xrightarrow{v'}B$；
对于 $B$ 的每个开理想 $\frK$，$u^{-1}(\frK)$ 包含某个 $\frJ_\lambda$，
所以 $v'{}^{-1}(\frK)$ 必包含 $S^{-1}\frJ_\lambda$，从而 $v'$ 连续。
由于 $B$ 分离且完备，$v'$ 唯一分解为
$S^{-1}A\to A\{S^{-1}\}\xrightarrow{u'}B$，其中 $u'$ 连续；断言得证。

\egapar{7.6.7}
设 $B$ 是另一个线性拓扑化环，$T$ 是 $B$ 的乘法子集，
$\varphi:A\to B$ 是连续同态，且 $\varphi(S)\subset T$。
由上文，连续同态 $A\to B\to B\{T^{-1}\}$ 唯一分解为
\[
 A\longrightarrow A\{S^{-1}\}\xrightarrow{\varphi'}B\{T^{-1}\},
\]
其中 $\varphi'$ 连续。特别地，若 $B=A$ 且 $\varphi$ 是恒等映射，
则当 $S\subset T$ 时，存在连续同态
\[
 \rho^{T,S}:A\{S^{-1}\}\longrightarrow A\{T^{-1}\},
\]
它由 $S^{-1}A\to T^{-1}A$ 通过取分离完备化而得到。
若 $U$ 是 $A$ 的第三个乘法子集，且 $S\subset T\subset U$，则
\[
 \rho^{U,S}=\rho^{U,T}\circ\rho^{T,S}.
\]

\egapar{7.6.8}
设 $S_1,S_2$ 是 $A$ 的两个乘法子集，$S'_2$ 是 $S_2$ 在
$A\{S_1^{-1}\}$ 中的典范像。存在典范拓扑同构
\[
 A\{S_1^{-1}\}\{S_2'{}^{-1}\}
 \xrightarrow{\sim}A\{S_1^{-1}S_2^{-1}\}.
\]
这可以从典范同构
$(S_1S_2)^{-1}A\simeq S_2''{}^{-1}(S_1^{-1}A)$ 得到，
其中 $S''_2$ 是 $S_2$ 在 $S_1^{-1}A$ 中的典范像；这一同构是双连续的。

\egapar{7.6.9}
设 $\mathfrak a$ 是 $A$ 的开理想。可以假设对每个 $\lambda$ 都有
$\frJ_\lambda\subset\mathfrak a$，于是环 $S^{-1}A$ 中有
$S^{-1}\frJ_\lambda\subset S^{-1}\mathfrak a$；
也就是说，$S^{-1}\mathfrak a$ 是 $S^{-1}A$ 的开理想。
以 $\mathfrak a\{S^{-1}\}$ 记它的分离完备化，等于
\[
 \plim_\lambda
 \bigl(S^{-1}\mathfrak a/S^{-1}\frJ_\lambda\bigr).
\]
这是 $A\{S^{-1}\}$ 的一个开理想，同构于 $S^{-1}\mathfrak a$ 的典范像的闭包。
此外，存在离散环的典范同构
\[
 A\{S^{-1}\}/\mathfrak a\{S^{-1}\}
 \simeq S^{-1}A/S^{-1}\mathfrak a=S^{-1}(A/\mathfrak a).
\]
反之，若 $\mathfrak a'$ 是 $A\{S^{-1}\}$ 的开理想，
则 $\mathfrak a'$ 包含某个形如 $\frJ_\lambda\{S^{-1}\}$ 的理想，
因而是 $S^{-1}A/S^{-1}\frJ_\lambda$ 的某个理想的逆像；
由（\hyperlink{para:0:1.2.6}{1.2.6}），该理想必来自某个 $S^{-1}\mathfrak a$，其中
$\mathfrak a\supset\frJ_\lambda$。于是有
$\mathfrak a'=\mathfrak a\{S^{-1}\}$。特别地，由（\hyperlink{para:0:1.2.6}{1.2.6}），得到下述命题。

\egastatement{7.6.10}{命题}
映射 $\mathfrak p\mapsto\mathfrak p\{S^{-1}\}$ 是从 $A$ 中满足
$\mathfrak p\cap S=\varnothing$ 的开素理想所成的集合，
到 $A\{S^{-1}\}$ 的开素理想所成的集合的保序双射。
此外，$A\{S^{-1}\}/\mathfrak p\{S^{-1}\}$ 的分式域
与 $A/\mathfrak p$ 的分式域典范同构。

\egastatement{7.6.11}{命题}
\begin{enumerate}[label=(\roman*)]
\item 若 $A$ 是可容许环，则 $A'=A\{S^{-1}\}$ 也是可容许环；
而且对 $A$ 的每个定义理想 $\frJ$，$\frJ'=\frJ\{S^{-1}\}$ 是 $A'$ 的定义理想。
\item 设 $A$ 是进制环，$\frJ$ 是 $A$ 的定义理想，且
$\frJ/\frJ^2$ 是有限型 $(A/\frJ)$-模。则 $A'$ 是 $\frJ'$-进制环，
$\frJ'/\frJ'{}^2$ 是有限型 $(A'/\frJ')$-模。
若 $A$ 还是诺特环，则 $A'$ 也是诺特环。
\end{enumerate}

(i) 若 $\frJ$ 是 $A$ 的定义理想，则 $S^{-1}\frJ$ 显然是拓扑环
$S^{-1}A$ 的定义理想，因为 $(S^{-1}\frJ)^n=S^{-1}\frJ^n$。
设 $A''$ 是与 $S^{-1}A$ 相伴的分离环，$\frJ''$ 是 $S^{-1}\frJ$ 在 $A''$
中的像。$S^{-1}\frJ^n$ 的像为 $\frJ''{}^n$，所以 $\frJ''{}^n$ 在 $A''$ 中趋于 $0$。
由于 $\frJ'$ 是 $\frJ''$ 在 $A'$ 中的闭包，$\frJ'{}^n$ 包含于
$\frJ''{}^n$ 的闭包，因而在 $A'$ 中趋于 $0$。

(ii) 令 $A_i=A/\frJ^{i+1}$；当 $i\leq j$ 时，以 $u_{ij}$ 记典范同态
$A/\frJ^{j+1}\to A/\frJ^{i+1}$。设 $S_i$ 是 $S$ 在 $A_i$ 中的典范像，
并令 $A'_i=S_i^{-1}A_i$；再以 $u'_{ij}:A'_j\to A'_i$ 记由
$u_{ij}$ 典范诱导的同态。
我们证明射影系 $(A'_i,u'_{ij})$ 满足命题（\hyperlink{para:0:7.2.7}{7.2.7}）的条件。
显然，各 $u'_{ij}$ 都是满射；另一方面，由（\hyperlink{para:0:1.3.2}{1.3.2}），$u'_{ij}$ 的核为
\[
 S_j^{-1}(\frJ^{i+1}/\frJ^{j+1})=\frJ_j'{}^{i+1},
 \qquad
 \frJ'_j=S_j^{-1}(\frJ/\frJ^{j+1}).
\]
最后，$\frJ'_1/\frJ_1'{}^2=S_1^{-1}(\frJ/\frJ^2)$；
由于 $\frJ/\frJ^2$ 是有限型 $(A/\frJ^2)$-模，
$\frJ'_1/\frJ_1'{}^2$ 是有限型 $A'_1$-模。
若 $A$ 是诺特环，则 $A'_0=S_0^{-1}(A/\frJ)$ 也是诺特环。
由（\hyperlink{para:0:7.2.8}{7.2.8}），命题的证明完成。

\egastatement{7.6.12}{推论}
在（\hyperlink{para:0:7.6.11}{7.6.11}，(ii)）的假设下，有
\[
 \bigl(\frJ\{S^{-1}\}\bigr)^n=\frJ^n\{S^{-1}\}.
\]
这由（\hyperlink{para:0:7.2.7}{7.2.7}）和（\hyperlink{para:0:7.6.11}{7.6.11}）的证明得到。

\egastatement{7.6.13}{命题}
设 $A$ 是诺特进制环，$S$ 是 $A$ 的乘法子集。则 $A\{S^{-1}\}$ 是平坦 $A$-模。

事实上，若 $\frJ$ 是 $A$ 的定义理想，则 $A\{S^{-1}\}$ 是诺特环
$S^{-1}A$ 关于 $S^{-1}\frJ$-预进拓扑的分离完备化；
因而由（\hyperlink{para:0:7.3.3}{7.3.3}），$A\{S^{-1}\}$ 是平坦 $S^{-1}A$-模。
又由于 $S^{-1}A$ 是平坦 $A$-模（\hyperlink{para:0:6.3.1}{6.3.1}），
由平坦性的传递性（\hyperlink{para:0:6.2.1}{6.2.1}）即得命题。

\egastatement{7.6.14}{推论}
在（\hyperlink{para:0:7.6.13}{7.6.13}）的假设下，设 $S'\subset S$ 是 $A$ 的另一个乘法子集。
则 $A\{S^{-1}\}$ 是平坦 $A\{S'{}^{-1}\}$-模。

事实上，由（\hyperlink{para:0:7.6.8}{7.6.8}），$A\{S^{-1}\}$ 可典范地等同于
$A\{S'{}^{-1}\}\{S_0^{-1}\}$，其中 $S_0$ 是 $S$ 在
$A\{S'{}^{-1}\}$ 中的典范像；而 $A\{S'{}^{-1}\}$ 是诺特环（\hyperlink{para:0:7.6.11}{7.6.11}）。

\egapar{7.6.15}
对线性拓扑化环 $A$ 的每个元素 $f$，以 $A_{\{f\}}$ 记完备分式环
$A\{S_f^{-1}\}$，其中 $S_f$ 是各 $f^n$（$n\geq0$）所成的乘法集合；
对 $A$ 的每个开理想 $\mathfrak a$，将 $\mathfrak a\{S_f^{-1}\}$ 记为
$\mathfrak a_{\{f\}}$。若 $g$ 是 $A$ 的另一个元素，
则存在典范连续同态 $A_{\{f\}}\to A_{\{fg\}}$（\hyperlink{para:0:7.6.7}{7.6.7}）。
因而，当 $f$ 遍历 $A$ 的乘法子集 $S$ 时，各 $A_{\{f\}}$ 关于上述同态
构成滤过的环归纳系；令
\[
 A_{\{S\}}=\ilim_{f\in S}A_{\{f\}}.
\]
对每个 $f\in S$，存在同态 $A_{\{f\}}\to A\{S^{-1}\}$（\hyperlink{para:0:7.6.7}{7.6.7}），
且这些同态组成归纳系；取归纳极限，便定义典范同态
$A_{\{S\}}\to A\{S^{-1}\}$。

\egastatement{7.6.16}{命题}
若 $A$ 是诺特环，则 $A\{S^{-1}\}$ 是平坦 $A_{\{S\}}$-模。
\transnote{底本此处只写“诺特环”。紧接的证明引用（\hyperlink{para:0:7.6.14}{7.6.14}），而该推论承接
（\hyperlink{para:0:7.6.13}{7.6.13}）的诺特进制环假设。本译文保留底本陈述，不将缺少的假设静默补入；
在后文（\hyperlink{para:0:7.6.18}{7.6.18}）的应用中，$A$ 明确被假设为诺特进制环。}

事实上，由（\hyperlink{para:0:7.6.14}{7.6.14}），对每个 $f\in S$，$A\{S^{-1}\}$ 都在环
$A_{\{f\}}$ 上平坦；由（\hyperlink{para:0:6.2.3}{6.2.3}）即得结论。

\egastatement{7.6.17}{命题}
设 $\mathfrak p$ 是可容许环 $A$ 的开素理想，令 $S=A-\mathfrak p$。
则 $A\{S^{-1}\}$ 与 $A_{\{S\}}$ 都是局部环，
典范同态 $A_{\{S\}}\to A\{S^{-1}\}$ 是局部同态；
而且这两个环的剩余域都与 $A/\mathfrak p$ 的分式域典范同构。

事实上，取 $A$ 的定义理想 $\frJ\subset\mathfrak p$；有
$S^{-1}\frJ\subset S^{-1}\mathfrak p=\mathfrak pA_{\mathfrak p}$，
故 $A_{\mathfrak p}/S^{-1}\frJ$ 是局部环。
由（\hyperlink{para:0:7.1.12}{7.1.12}）、（\hyperlink{para:0:7.6.9}{7.6.9}）和（\hyperlink{para:0:7.6.11}{7.6.11}，(i)）可知，$A\{S^{-1}\}$ 是局部环。
令 $\mathfrak m=\ilim_{f\in S}\mathfrak p_{\{f\}}$，这是 $A_{\{S\}}$ 的一个理想；
我们将证明 $A_{\{S\}}$ 中每个不属于 $\mathfrak m$ 的元素都可逆。

这样的元素是某个 $z\in A_{\{f\}}$ 在 $A_{\{S\}}$ 中的像，
其中 $f\in S$ 且 $z\notin\mathfrak p_{\{f\}}$。
其在
$A_{\{f\}}/\frJ_{\{f\}}=S_f^{-1}(A/\frJ)$ 中的典范像 $z_0$
因而不属于 $S_f^{-1}(\mathfrak p/\frJ)$（\hyperlink{para:0:7.6.9}{7.6.9}）；
也就是说，$z_0=\overline x/\overline f^{\,k}$，其中 $x\notin\mathfrak p$，
而 $\overline x,\overline f$ 分别是 $x,f$ 模 $\frJ$ 的类。
由 $x\in S$，有 $g=xf\in S$；在 $S_g^{-1}A$ 中，
$x/f^k\in S_f^{-1}A$ 的典范像
\[
 y_0=\frac{x^{k+1}}{g^k}
\]
具有逆元 $x^{k-1}f^{2k}/g^k$。
因而，$y_0$ 在 $S_g^{-1}A/S_g^{-1}\frJ$ 中的像可逆；
由（\hyperlink{para:0:7.6.9}{7.6.9}）和（\hyperlink{para:0:7.1.12}{7.1.12}），$z$ 在 $A_{\{g\}}$ 中的典范像 $y$ 可逆。
所以 $z$ 在 $A_{\{S\}}$ 中的像（等于 $y$ 的像）可逆。

因此，$A_{\{S\}}$ 是局部环，其极大理想为 $\mathfrak m$。
此外，$\mathfrak p_{\{f\}}$ 在 $A\{S^{-1}\}$ 中的像包含于该环的极大理想
$\mathfrak p\{S^{-1}\}$；所以 $\mathfrak m$ 的像也包含于
$\mathfrak p\{S^{-1}\}$，从而典范同态
$A_{\{S\}}\to A\{S^{-1}\}$ 是局部同态。
最后，$A\{S^{-1}\}/\mathfrak p\{S^{-1}\}$ 的每个元素，都是某个适当
$f\in S$ 所对应的环 $S_f^{-1}A$ 中某个元素的像，
故同态 $A_{\{S\}}\to A\{S^{-1}\}/\mathfrak p\{S^{-1}\}$ 为满射；
取商后，即得到两个剩余域之间的同构。

\egastatement{7.6.18}{推论}
在（\hyperlink{para:0:7.6.17}{7.6.17}）的假设下，若 $A$ 还是诺特进制环，则局部环 $A\{S^{-1}\}$
和 $A_{\{S\}}$ 都是诺特环，而且 $A\{S^{-1}\}$ 是忠实平坦 $A_{\{S\}}$-模。

已知 $A\{S^{-1}\}$ 是诺特环（\hyperlink{para:0:7.6.11}{7.6.11}，(ii)），
且在 $A_{\{S\}}$ 上平坦（\hyperlink{para:0:7.6.16}{7.6.16}）。
由于同态 $A_{\{S\}}\to A\{S^{-1}\}$ 是局部同态，
由（\hyperlink{para:0:6.6.2}{6.6.2}），$A\{S^{-1}\}$ 是忠实平坦 $A_{\{S\}}$-模；
进而由（\hyperlink{para:0:6.5.2}{6.5.2}），$A_{\{S\}}$ 是诺特环。


% END chapters/ch0_s7_6.tex

\par

% BEGIN chapters/ch0_s7_7.tex
\subsection{完备张量积}\hypertarget{heading:0:7.7}{}
\src{75—77}

\egapar{7.7.1}
设 $A$ 是线性拓扑化环，$M,N$ 是两个线性拓扑化 $A$-模。
设 $\frJ,V,W$ 分别是 $A,M,N$ 中 $0$ 的开邻域，它们都是 $A$-模，
且满足 $\frJ M\subset V$、$\frJ N\subset W$，从而 $M/V$ 与 $N/W$
都可看成 $(A/\frJ)$-模。
当 $\frJ,V,W$ 遍历满足上述条件的开邻域系时，模
$(M/V)\otimes_{A/\frJ}(N/W)$ 显然构成一个定义在环射影系 $A/\frJ$
上的模射影系。因此，取射影极限，便得到 $A$ 的分离完备化 $\comp A$ 上的一个模，
称为 $M$ 与 $N$ 的\emph{完备张量积}，记为 $(M\otimes_A N)^\wedge$。

注意，$M/V$ 与 $\comp M/\overline V$ 典范同构，其中 $\comp M$ 是 $M$ 的
分离完备化，$\overline V$ 是 $V$ 的像在 $\comp M$ 中的闭包。
由此可见，完备张量积 $(M\otimes_A N)^\wedge$ 可典范地等同于
$(\comp M\otimes_{\comp A}\comp N)^\wedge$；后者也记为
$\comp M\ctensor_{\comp A}\comp N$。

\egapar{7.7.2}
沿用上述记号，张量积 $(M/V)\otimes_A(N/W)$ 与
$(M/V)\otimes_{A/\frJ}(N/W)$ 可典范地等同；
因而也可等同于
\[
 (M\otimes_A N)\big/
 \bigl(\Im(V\otimes_A N)+\Im(M\otimes_A W)\bigr).
\]
所以，$(M\otimes_A N)^\wedge$ 是 $A$-模 $M\otimes_A N$ 的分离完备化，
其中所赋的拓扑以各子模
\[
 \Im(V\otimes_A N)+\Im(M\otimes_A W)
\]
为 $0$ 的一个基本邻域系，$V,W$ 分别遍历 $M,N$ 的开子模。
为简便起见，称这一拓扑为 $M,N$ 上给定拓扑的\emph{张量积拓扑}。

\egapar{7.7.3}
设 $M',N'$ 是两个线性拓扑化 $A$-模，$u:M\to M'$、$v:N\to N'$ 是连续同态。
显然，$u\otimes v$ 关于 $M\otimes N$ 和 $M'\otimes N'$ 上的张量积拓扑连续；
取分离完备化，得到连续同态
\[
 (M\otimes N)^\wedge\longrightarrow(M'\otimes N')^\wedge,
\]
记为 $u\ctensor v$。
因此，在以线性拓扑化 $A$-模为对象的范畴中，
$(M\otimes_A N)^\wedge$ 是关于 $M,N$ 的双函子。

\egapar{7.7.4}
同样可以定义任意有限个线性拓扑化 $A$-模的完备张量积；
显然，这个积具有通常的结合性与交换性。

\egapar{7.7.5}
若 $B,C$ 是两个线性拓扑化的拓扑 $A$-代数，则 $B\otimes_A C$ 上的张量积拓扑，
以该代数中的理想
\[
 \Im(\frK\otimes_A C)+\Im(B\otimes_A\mathfrak L)
\]
为 $0$ 的一个基本邻域系，其中 $\frK$、$\mathfrak L$ 分别遍历 $B,C$ 的开理想。
因此，$(B\otimes_A C)^\wedge$ 具有拓扑 $\comp A$-代数结构，
它是 $(A/\frJ)$-代数射影系
\[
 (B/\frK)\otimes_{A/\frJ}(C/\mathfrak L)
\]
的射影极限，其中 $\frJ$ 是满足
$\frJ B\subset\frK$、$\frJ C\subset\mathfrak L$ 的 $A$ 的开理想
（这样的理想总存在）。称这一代数为代数 $B,C$ 的\emph{完备张量积}。

\egapar{7.7.6}
当 $B\otimes_A C$ 赋以张量积拓扑时，$A$-代数同态
$b\mapsto b\otimes1$、$c\mapsto1\otimes c$，分别从 $B,C$ 到 $B\otimes_A C$，
都是连续的。将它们与 $B\otimes_A C$ 到其分离完备化的典范同态复合，
便得到典范同态
\[
 \rho:B\longrightarrow(B\otimes_A C)^\wedge,
 \qquad
 \sigma:C\longrightarrow(B\otimes_A C)^\wedge.
\]
代数 $(B\otimes_A C)^\wedge$ 及同态 $\rho,\sigma$ 还具有下列泛性质：
对于每个线性拓扑化、分离且完备的 $A$-代数 $D$，以及任意一对连续
$A$-同态 $u:B\to D$、$v:C\to D$，存在唯一连续 $A$-同态
\[
 w:(B\otimes_A C)^\wedge\longrightarrow D,
\]
使得 $u=w\circ\rho$、$v=w\circ\sigma$。

事实上，已经存在唯一 $A$-同态 $w_0:B\otimes_A C\to D$，
使得 $u(b)=w_0(b\otimes1)$、$v(c)=w_0(1\otimes c)$。
只须证明 $w_0$ 连续，因为此后取分离完备化便得到连续同态 $w$。
设 $\mathfrak M$ 是 $D$ 的开理想；由假设，存在开理想
$\frK\subset B$、$\mathfrak L\subset C$，满足
$u(\frK)\subset\mathfrak M$、$v(\mathfrak L)\subset\mathfrak M$。
于是 $\Im(\frK\otimes C)+\Im(B\otimes\mathfrak L)$ 在 $w_0$ 下的像
仍包含于 $\mathfrak M$，断言得证。

\egastatement{7.7.7}{命题}
若 $B,C$ 是两个预可容许的 $A$-代数，则 $(B\otimes_A C)^\wedge$ 是可容许的。
若 $\frK$、$\mathfrak L$ 分别是 $B,C$ 的定义理想，则
\[
 \mathfrak H=\Im(\frK\otimes C)+\Im(B\otimes\mathfrak L)
\]
的典范像在 $(B\otimes_A C)^\wedge$ 中的闭包是一个定义理想。

只须证明 $\mathfrak H^n$ 关于张量积拓扑趋于 $0$；
这直接由下列包含关系得到：
\[
 \mathfrak H^{2n}\subset
 \Im(\frK^n\otimes C)+\Im(B\otimes\mathfrak L^n).
\]

\egastatement{7.7.8}{命题}
设 $A$ 是预进制环，$\frJ$ 是 $A$ 的定义理想，$M$ 是赋有 $\frJ$-预进拓扑的
有限型 $A$-模。对于每个诺特进制的拓扑 $A$-代数 $B$，
$B\otimes_A M$ 可等同于完备张量积 $(B\otimes_A M)^\wedge$。

若 $\frK$ 是 $B$ 的定义理想，则由假设，存在整数 $m$ 使得
$\frJ^mB\subset\frK$，所以
\begin{align*}
 \Im(B\otimes_A\frJ^{nm}M)
 &=\Im(\frJ^{nm}B\otimes_A M)\\
 &\subset\Im(\frK^nB\otimes_A M)
 =\frK^n(B\otimes_A M).
\end{align*}
由此可知，$B\otimes_A M$ 上由 $B$ 和 $M$ 的拓扑构成的张量积拓扑，
就是 $\frK$-预进拓扑。
由于 $B\otimes_A M$ 是有限型 $B$-模，命题立即由（\hyperlink{para:0:7.3.6}{7.3.6}）得到。


% END chapters/ch0_s7_7.tex

\par

% BEGIN chapters/ch0_s7_8.tex
\subsection{同态模上的拓扑}\hypertarget{heading:0:7.8}{}
\src{77—78}

\egapar{7.8.1}
设 $A$ 是诺特 $\frJ$-进制环，$M,N$ 是两个赋有 $\frJ$-预进拓扑的有限型
$A$-模。已知它们分离且完备（\hyperlink{para:0:7.3.6}{7.3.6}）；此外，每个 $A$-同态 $M\to N$
都自动连续，且 $A$-模 $\Hom_A(M,N)$ 是有限型的。
对每个整数 $i\geq0$，令
\[
 A_i=A/\frJ^{i+1},\qquad
 M_i=M/\frJ^{i+1}M,\qquad
 N_i=N/\frJ^{i+1}N.
\]
当 $i\leq j$ 时，每个同态 $u_j:M_j\to N_j$ 都将 $\frJ^{i+1}M_j$
映入 $\frJ^{i+1}N_j$，所以通过取商，得到同态 $u_i:M_i\to N_i$。
由此定义典范同态
\[
 \Hom_{A_j}(M_j,N_j)\longrightarrow\Hom_{A_i}(M_i,N_i).
\]
各 $\Hom_{A_i}(M_i,N_i)$ 关于这些同态组成射影系；
由（\hyperlink{para:0:7.2.10}{7.2.10}），有典范同构
\[
 \varphi:\Hom_A(M,N)\longrightarrow
 \plim_i\Hom_{A_i}(M_i,N_i).
\]
此外，还有下述命题。

\egastatement{7.8.2}{命题}
若 $M,N$ 是诺特 $\frJ$-进制环 $A$ 上的有限型模，则子模
$\Hom_A(M,\frJ^{i+1}N)$ 组成 $\Hom_A(M,N)$ 中关于 $\frJ$-进拓扑的
$0$ 的一个基本邻域系；而且典范同态
\[
 \varphi:\Hom_A(M,N)\longrightarrow
 \plim_i\Hom_{A_i}(M_i,N_i)
\]
是拓扑同构。

事实上，可以将 $M$ 看成有限型自由 $A$-模 $L$ 的商，
从而将 $\Hom_A(M,N)$ 等同于 $\Hom_A(L,N)$ 的子模。
在这个等同下，$\Hom_A(M,\frJ^{i+1}N)$ 是
$\Hom_A(M,N)$ 与 $\Hom_A(L,\frJ^{i+1}N)$ 的交。
由（\hyperlink{para:0:7.3.2}{7.3.2}），$\Hom_A(L,N)$ 的 $\frJ$-进拓扑在
$\Hom_A(M,N)$ 上诱导的拓扑就是 $\frJ$-进拓扑。
因而，第一个断言归结到 $M=L=A^m$ 的情形；此时
\[
 \Hom_A(L,N)=N^m,\qquad
 \Hom_A(L,\frJ^{i+1}N)=(\frJ^{i+1}N)^m=\frJ^{i+1}N^m,
\]
结论显然。

\begin{sloppypar}
为证明第二个断言，注意当 $j\leq i$ 时，$\Hom_A(M,\frJ^{i+1}N)$ 在
$\Hom_{A_j}(M_j,N_j)$ 中的像为零，故 $\varphi$ 连续。
反之，$\Hom_{A_i}(M_i,N_i)$ 中的 $0$ 在 $\Hom_A(M,N)$ 中的逆像为
$\Hom_A(M,\frJ^{i+1}N)$，所以 $\varphi$ 双连续。
\end{sloppypar}


若只假设 $A$ 是诺特 $\frJ$-预进制环，$M,N$ 是两个关于
$\frJ$-预进拓扑分离的有限型 $A$-模，则上述证明表明，
（\hyperlink{para:0:7.8.2}{7.8.2}）的第一个断言仍成立，而且 $\varphi$ 是
$\Hom_A(M,N)$ 到 $\plim_i\Hom_{A_i}(M_i,N_i)$ 的一个子模的拓扑同构。

\egastatement{7.8.3}{命题}
在（\hyperlink{para:0:7.8.2}{7.8.2}）的假设下，从 $M$ 到 $N$ 的单射同态
（分别地，满射同态、双射同态）所成的集合，是 $\Hom_A(M,N)$ 的开子集。

事实上，由（\hyperlink{para:0:7.3.5}{7.3.5}）和（\hyperlink{para:0:7.1.14}{7.1.14}），$u$ 为满射，当且仅当相应同态
$u_0:M/\frJ M\to N/\frJ N$ 为满射。因此，从 $M$ 到 $N$ 的满射同态所成的集合，
是某个离散空间的子集在连续映射
$\Hom_A(M,N)\to\Hom_{A_0}(M_0,N_0)$ 下的逆像。

现证明单射同态所成的集合是开的。设 $v$ 是这样的同态，令 $M'=v(M)$。
由 Artin--Rees 引理（\hyperlink{para:0:7.3.2.1}{7.3.2.1}），存在整数 $k\geq0$，使得对每个 $m>0$，有
\[
 M'\cap\frJ^{m+k}N\subset\frJ^mM'.
\]
我们将证明，对于 $w\in\frJ^{k+1}\Hom_A(M,N)$，$u=v+w$ 为单射；
这便完成证明。

设 $x\in M$ 满足 $u(x)=0$。对每个 $i\geq0$，我们证明
$x\in\frJ^iM$ 蕴含 $x\in\frJ^{i+1}M$；由此得到
$x\in\bigcap_{i\geq0}\frJ^iM=(0)$。
事实上，此时 $w(x)\in\frJ^{i+k+1}N$；另一方面，$w(x)=-v(x)\in M'$，所以
\[
 v(x)\in M'\cap\frJ^{i+k+1}N\subset\frJ^{i+1}M'.
\]
由于 $v$ 是 $M$ 到 $M'$ 的同构，得 $x\in\frJ^{i+1}M$。证毕。

\begin{flushright}
（待续。）
\end{flushright}


% END chapters/ch0_s7_8.tex

\par




% END text/ch0_01_07.tex


% BEGIN text/ch0_08_13.tex
\renewcommand{\B}{\mathscr B}
\NewNote{以下§§8--13原载第三卷第一部分，现接在本章§7之后。原编号及原注保留。}
\origfootnote{为便于查找引文，今后引用与第一章一同发表的第0章各节时，采用记号 $0_{\mathrm I}$。}
\setcounter{section}{7}
\section{可表函子}
\subsection{可表函子}
\src{5--8}
\egapar{8.1.1}
以 $\Ens$ 表示集合范畴。设 $C$ 为一个范畴；对 $C$ 中的两个对象 $X,Y$，令 $h_X(Y)=\Hom(Y,X)$；对 $C$ 中任一态射 $u:Y\to Y'$，以 $h_X(u)$ 表示从 $\Hom(Y',X)$ 到 $\Hom(Y,X)$ 的映射 $v\mapsto vu$。由这些定义立即可见，$h_X:C\to\Ens$ 是反变函子，即由 $C$ 的对偶范畴 $C^\circ$ 到 $\Ens$ 的协变函子所成范畴 $\catHom(C^\circ,\Ens)$ 中的一个对象（T，1.7，d）及 [29]）。

\egapar{8.1.2}
现在设 $w:X\to X'$ 是 $C$ 中的一个态射；对每个 $Y\in C$ 及每个 $v\in\Hom(Y,X)=h_X(Y)$，有 $wv\in\Hom(Y,X')=h_{X'}(Y)$。以 $h_w(Y)$ 表示从 $h_X(Y)$ 到 $h_{X'}(Y)$ 的映射 $v\mapsto wv$。立即可见，对每个态射 $u:Y\to Y'$，图
\[
\begin{tikzcd}[column sep=large]
h_X(Y')\arrow[r,"h_X(u)"]\arrow[d,"h_w(Y')"']&h_X(Y)\arrow[d,"h_w(Y)"]\\
h_{X'}(Y')\arrow[r,"h_{X'}(u)"']&h_{X'}(Y)
\end{tikzcd}
\]
交换；换言之，$h_w$ 是函子态射 $h_X\to h_{X'}$（T，1.2），也就是范畴 $\catHom(C^\circ,\Ens)$ 中的态射（T，1.7，d））。因此，$h_X$ 和 $h_w$ 的定义给出一个典范协变函子
\begin{equation}\tag{8.1.2.1}
h:C\longrightarrow\catHom(C^\circ,\Ens).
\end{equation}

\egapar{8.1.3}
设 $X$ 是 $C$ 中的对象，$F$ 是从 $C$ 到 $\Ens$ 的反变函子（即 $\catHom(C^\circ,\Ens)$ 中的对象）。设 $g:h_X\to F$ 为函子态射：于是对每个 $Y\in C$，$g(Y)$ 是映射 $h_X(Y)\to F(Y)$，而且对 $C$ 中每个态射 $u:Y\to Y'$，图
\begin{equation}\tag{8.1.3.1}
\begin{tikzcd}[column sep=large]
h_X(Y')\arrow[r,"h_X(u)"]\arrow[d,"g(Y')"']&h_X(Y)\arrow[d,"g(Y)"]\\
F(Y')\arrow[r,"F(u)"']&F(Y)
\end{tikzcd}
\end{equation}
交换。特别地，有映射 $g(X):h_X(X)=\Hom(X,X)\to F(X)$，从而有元素
\begin{equation}\tag{8.1.3.2}
\alpha(g)=(g(X))(1_X)\in F(X),
\end{equation}
进而有典范映射
\begin{equation}\tag{8.1.3.3}
\alpha:\Hom(h_X,F)\longrightarrow F(X).
\end{equation}
反过来，取一个元素 $\xi\in F(X)$；对 $C$ 中每个态射 $v:Y\to X$，$F(v)$ 是从 $F(X)$ 到 $F(Y)$ 的映射。考虑从 $h_X(Y)$ 到 $F(Y)$ 的映射
\begin{equation}\tag{8.1.3.4}
v\longmapsto(F(v))(\xi).
\end{equation}
若以 $(\beta(\xi))(Y)$ 表示此映射，则
\begin{equation}\tag{8.1.3.5}
\beta(\xi):h_X\longrightarrow F
\end{equation}
是函子态射，因为对 $C$ 中每个态射 $u:Y\to Y'$，有
\[
(F(vu))(\xi)=(F(v)\circ F(u))(\xi),
\]
这就验证了在 $g=\beta(\xi)$ 时（8.1.3.1）的交换性。\transnote{此处底本右端确实印作 $F(v)\circ F(u)$。按本段的反变函子约定，当 $v:Y'\to X$ 时，应读作 $F(u)\circ F(v)$；这里保留底本的式样，并明确标出复合次序疑误。}
这样就定义了典范映射
\begin{equation}\tag{8.1.3.6}
\beta:F(X)\longrightarrow\Hom(h_X,F).
\end{equation}

\egastatement{8.1.4}{命题}
映射 $\alpha$ 和 $\beta$ 是互为逆映射的双射。
\begin{proof}
对 $\xi\in F(X)$ 计算 $\alpha(\beta(\xi))$。对每个 $Y\in C$，$(\beta(\xi))(Y)$ 是从 $h_X(Y)$ 到 $F(Y)$ 的映射 $g_1(Y):v\mapsto(F(v))(\xi)$。因此
\[
\alpha(\beta(\xi))=(g_1(X))(1_X)=(F(1_X))(\xi)=1_{F(X)}(\xi)=\xi.
\]
现在对 $g\in\Hom(h_X,F)$ 计算 $\beta(\alpha(g))$。对每个 $Y\in C$，$(\beta(\alpha(g)))(Y)$ 是映射
\[
v\longmapsto(F(v))\bigl((g(X))(1_X)\bigr).
\]
由（8.1.3.1）的交换性，此映射就是
\[
v\longmapsto(g(Y))\bigl((h_X(v))(1_X)\bigr)=(g(Y))(v),
\]
最后的等式来自 $h_X(v)$ 的定义。也就是说，此映射等于 $g(Y)$，命题得证。
\end{proof}

\egapar{8.1.5}
回顾一下：范畴 $C$ 的一个\textbf{子范畴} $C'$，由以下条件定义：其对象是 $C$ 的对象；若 $X',Y'$ 是 $C'$ 中的两个对象，则 $C'$ 中从 $X'$ 到 $Y'$ 的态射集合 $\Hom_{C'}(X',Y')$ 是 $C$ 中态射集合 $\Hom_C(X',Y')$ 的一个子集，而且“态射复合”的典范映射
\[
\Hom_{C'}(X',Y')\times\Hom_{C'}(Y',Z')\longrightarrow\Hom_{C'}(X',Z')
\]
是典范映射
\[
\Hom_C(X',Y')\times\Hom_C(Y',Z')\longrightarrow\Hom_C(X',Z')
\]
的限制。

若对 $C'$ 中每一对对象都有 $\Hom_{C'}(X',Y')=\Hom_C(X',Y')$，则称 $C'$ 是 $C$ 的\textbf{满子范畴}。由 $C$ 中所有同构于 $C'$ 的对象的对象组成的子范畴 $C''$，这时也为 $C$ 的满子范畴，并且容易验证，它与 $C'$ 等价（T，1.2）。

协变函子 $F:C_1\to C_2$ 称为\textbf{全忠实}的，若对 $C_1$ 中每一对对象 $X_1,Y_1$，映射
\[
u\longmapsto F(u):\Hom(X_1,Y_1)\longrightarrow\Hom(F(X_1),F(Y_1))
\]
为双射。这蕴含 $C_2$ 的子范畴 $F(C_1)$ 是满的。此外，若两个对象 $X_1,X'_1$ 有同一个像 $X_2$，则存在唯一同构 $u:X_1\to X'_1$，使 $F(u)=1_{X_2}$。对 $F(C_1)$ 中每个对象 $X_2$，取满足 $F(X_1)=X_2$ 的对象 $X_1$ 之一，记为 $G(X_2)$（这里用选择公理定义 $G$）；对 $F(C_1)$ 中任一态射 $v:X_2\to Y_2$，以 $G(v)$ 表示满足 $F(u)=v$ 的唯一态射 $u:G(X_2)\to G(Y_2)$。于是 $G$ 是从 $F(C_1)$ 到 $C_1$ 的函子，$FG$ 是 $F(C_1)$ 上的恒等函子。上述论述又表明，存在函子同构 $\varphi:1_{C_1}\to GF$，使 $F,G,\varphi$ 与恒等同构 $1_{F(C_1)}\to FG$ 给出范畴 $C_1$ 与 $C_2$ 的满子范畴 $F(C_1)$ 之间的等价（T，1.2）。

\egapar{8.1.6}
将命题（8.1.4）用于 $F=h_{X'}$ 的情形，其中 $X'$ 是 $C$ 的任意对象。映射
\[
\beta:\Hom(X,X')\longrightarrow\Hom(h_X,h_{X'})
\]
在这里正是（8.1.2）定义的 $w\mapsto h_w$。此映射为双射，故依（8.1.5）的术语，有：

\egastatement{8.1.7}{命题}
典范函子 $h:C\to\catHom(C^\circ,\Ens)$ 是全忠实的。

\egapar{8.1.8}
设 $F$ 为从 $C$ 到 $\Ens$ 的反变函子；若存在对象 $X\in C$，使 $F$ 同构于 $h_X$，就称 $F$ 是\textbf{可表}的。由（8.1.7）可知，给定 $X\in C$ 和函子同构 $g:h_X\to F$，便在唯一同构意义下确定了 $X$。命题（8.1.7）还意味着，$h$ 给出 $C$ 与 $\catHom(C^\circ,\Ens)$ 中由所有反变可表函子组成的满子范畴之间的等价。另一方面，由（8.1.4），给定函子态射 $g:h_X\to F$ 等价于给定一个元素 $\xi\in F(X)$；$g$ 为同构等价于 $\xi$ 满足以下条件：\emph{对 $C$ 的每个对象 $Y$，从 $\Hom(Y,X)$ 到 $F(Y)$ 的映射 $v\mapsto(F(v))(\xi)$ 为双射。} 当 $\xi$ 满足此条件时，就说偶对 $(X,\xi)$ \textbf{表示}可表函子 $F$。若存在这样的 $\xi\in F(X)$，也常略称对象 $X\in C$ 表示 $F$，即 $h_X$ 同构于 $F$。

设 $F,F'$ 是从 $C$ 到 $\Ens$ 的两个反变可表函子，且给定函子同构 $h_X\to F$ 与 $h_{X'}\to F'$。由（8.1.6），$\Hom(X,X')$ 与函子态射 $F\to F'$ 的集合 $\Hom(F,F')$ 之间有典范的一一对应。

\egapar{8.1.9}
\textbf{例 I：射影极限。} 反变可表函子的概念，特别包含通常所谓“泛问题的解”的对偶概念。更一般地，我们将看到，射影极限是可表函子的一个特例。回顾一下：在范畴 $C$ 中，给定一个预序集 $I$、一族对象 $(A_\alpha)_{\alpha\in I}$，以及对每对满足 $\alpha\leq\beta$ 的指标给定态射 $u_{\alpha\beta}:A_\beta\to A_\alpha$，便定义一个射影系。此系在 $C$ 中的一个射影极限，由对象 $B$（记作 $\plim A_\alpha$）及对每个 $\alpha\in I$ 给定的态射 $u_\alpha:B\to A_\alpha$ 组成，满足：1° 当 $\alpha\leq\beta$ 时，$u_\alpha=u_{\alpha\beta}u_\beta$；2° 对 $C$ 中每个对象 $X$ 及每族满足 $v_\alpha=u_{\alpha\beta}v_\beta$（$\alpha\leq\beta$）的态射 $v_\alpha:X\to A_\alpha$，存在唯一态射 $v:X\to B$（记作 $\plim v_\alpha$），使每个 $\alpha$ 都有 $v_\alpha=u_\alpha v$（T，1.8）。可以作如下解释：$u_{\alpha\beta}$ 典范地定义映射
\[
\overline u_{\alpha\beta}:\Hom(X,A_\beta)\longrightarrow\Hom(X,A_\alpha),
\]
这些映射给出集合的射影系 $(\Hom(X,A_\alpha),\overline u_{\alpha\beta})$；按定义，$(v_\alpha)$ 是集合 $\plim_\alpha\Hom(X,A_\alpha)$ 的一个元素。显然 $X\mapsto\plim_\alpha\Hom(X,A_\alpha)$ 是从 $C$ 到 $\Ens$ 的反变函子；射影极限 $B$ 的存在，等价于 $(v_\alpha)\mapsto\plim v_\alpha$ 是关于 $X$ 的函子同构
\begin{equation}\tag{8.1.9.1}
\plim_\alpha\Hom(X,A_\alpha)\simeq\Hom(X,B),
\end{equation}
即函子 $X\mapsto\plim_\alpha\Hom(X,A_\alpha)$ 可表。

\egapar{8.1.10}
\textbf{例 II：终对象。} 设 $C$ 为一个范畴，$\{a\}$ 为单元素集合。考虑反变函子 $F:C\to\Ens$：它把每个对象 $X$ 对应于集合 $\{a\}$，把每个态射 $X\to X'$ 对应于唯一映射 $\{a\}\to\{a\}$。此函子可表，意味着存在对象 $e\in C$，使每个 $Y\in C$ 都满足 $\Hom(Y,e)=h_e(Y)$ 为单元素集合；称 $e$ 为 $C$ 的一个\textbf{终对象}。显然，$C$ 的任意两个终对象同构（这就容许通常借助选择公理，选定 $C$ 的一个终对象并记为 $C_e$）。例如，集合范畴 $\Ens$ 中的终对象是单元素集合；在域 $K$ 上的增广代数范畴中（态射为与增广相容的代数同态），$K$ 是一个终对象；在 $S$-预概形范畴中（I，2.5.1），$S$ 是一个终对象。

\egapar{8.1.11}
对范畴 $C$ 中两个对象 $X,Y$，令 $h^X(Y)=\Hom(X,Y)$；对每个态射 $u:Y\to Y'$，以 $h^X(u)$ 表示从 $\Hom(X,Y)$ 到 $\Hom(X,Y')$ 的映射 $v\mapsto uv$。于是 $h^X$ 是协变函子 $C\to\Ens$。如同（8.1.2），可由此定义一个典范协变函子 $h':C^\circ\to\catHom(C,\Ens)$。从 $C$ 到 $\Ens$ 的协变函子 $F$，即 $\catHom(C,\Ens)$ 中的对象，若存在对象 $X\in C$（必在唯一同构意义下唯一），使 $F$ 同构于 $h^X$，便称为可表函子。与上述论述“对偶”的考虑留给读者展开；这一概念所涵盖的，此时是归纳极限，尤其是通常“泛问题的解”的概念。\transnote{这里及下文“在唯一同构意义下唯一”均须连同表示同构（或泛元素）理解；不带这一结构的对象本身未必没有非平凡自同构。}

\subsection{范畴中的代数结构}
\src{9--12}
\egapar{8.2.1}
给定从 $C$ 到 $\Ens$ 的两个反变函子 $F,F'$，回顾其乘积的定义：对每个 $Y\in C$，令 $(F\times F')(Y)=F(Y)\times F'(Y)$；对每个态射 $u:Y\to Y'$，令 $(F\times F')(u)=F(u)\times F'(u)$，即从 $F(Y')\times F'(Y')$ 到 $F(Y)\times F'(Y)$ 的映射
\[
(\xi,\xi')\longmapsto\bigl((F(u))(\xi),(F'(u))(\xi')\bigr).
\]
于是 $F\times F':C\to\Ens$ 是反变函子（它也正是范畴 $\catHom(C^\circ,\Ens)$ 中对象 $F,F'$ 的乘积）。给定对象 $X\in C$，把函子态射
\begin{equation}\tag{8.2.1.1}
\gamma_X:h_X\times h_X\longrightarrow h_X
\end{equation}
称为 $X$ 上的\textbf{内合成律}。

换言之（T，1.2），对每个 $Y\in C$，$\gamma_X(Y)$ 是映射 $h_X(Y)\times h_X(Y)\to h_X(Y)$，因此按定义是集合 $h_X(Y)$ 上的一个内合成律；它满足：对 $C$ 中每个态射 $u:Y\to Y'$，图
\[
\begin{tikzcd}[column sep=huge]
h_X(Y')\times h_X(Y')\arrow[r,"h_X(u)\times h_X(u)"]\arrow[d,"\gamma_X(Y')"']&h_X(Y)\times h_X(Y)\arrow[d,"\gamma_X(Y)"]\\
h_X(Y')\arrow[r,"h_X(u)"']&h_X(Y)
\end{tikzcd}
\]
交换。这意味着，关于合成律 $\gamma_X(Y')$ 与 $\gamma_X(Y)$，$h_X(u)$ 是从 $h_X(Y')$ 到 $h_X(Y)$ 的同态。

同样，给定 $C$ 中两个对象 $Z,X$，将函子态射
\begin{equation}\tag{8.2.1.2}
\omega_{X,Z}:h_Z\times h_X\longrightarrow h_X
\end{equation}
称为 $X$ 上\textbf{以 $Z$ 为算子域的外合成律}。如上可见，对每个 $Y\in C$，$\omega_{X,Z}(Y)$ 是 $h_X(Y)$ 上以 $h_Z(Y)$ 为算子域的外合成律；而且对每个态射 $u:Y\to Y'$，$h_Z(u)$ 与 $h_X(u)$ 构成从 $(h_Z(Y'),h_X(Y'))$ 到 $(h_Z(Y),h_X(Y))$ 的双同态。

\egapar{8.2.2}
设 $X'$ 是 $C$ 的另一个对象，并在 $X'$ 上给定内合成律 $\gamma_{X'}$。若对每个 $Y\in C$，$h_w(Y):h_X(Y)\to h_{X'}(Y)$ 都是关于合成律 $\gamma_X(Y)$ 与 $\gamma_{X'}(Y)$ 的同态，则称 $C$ 中态射 $w:X\to X'$ 是关于这些合成律的\textbf{同态}。若第三个对象 $X''$ 带有内合成律 $\gamma_{X''}$，且态射 $w':X'\to X''$ 为关于 $\gamma_{X'}$ 与 $\gamma_{X''}$ 的同态，则显然 $w'w:X\to X''$ 是关于 $\gamma_X$ 与 $\gamma_{X''}$ 的同态。若 $C$ 中同构 $w:X\iso X'$ 及其逆态射 $w^{-1}$ 分别为关于两方向合成律的同态，就称 $w$ 为关于合成律 $\gamma_X,\gamma_{X'}$ 的\textbf{同构}。

对于带外合成律的 $C$ 中对象偶对，同样定义双同态。

\egapar{8.2.3}
若对象 $X\in C$ 上的内合成律 $\gamma_X$ 满足：对每个 $Y\in C$，$\gamma_X(Y)$ 都是 $h_X(Y)$ 上的群律，则称带此合成律的 $X$ 为一个 \textbf{$C$-群}，或一个 \textbf{$C$-群对象}。同样定义 $C$-环、$C$-模等等。

\egapar{8.2.4}
假设对象 $X\in C$ 与自身的乘积 $X\times X$ 在 $C$ 中存在。按定义，在典范同构意义下有 $h_{X\times X}=h_X\times h_X$，因为乘积是射影极限的特例（8.1.9）。因此，$X$ 上的内合成律可视为函子态射 $\gamma_X:h_{X\times X}\to h_X$，从而典范地确定（8.1.6）一个满足 $h_{c_X}=\gamma_X$ 的元素 $c_X\in\Hom(X\times X,X)$。所以在此情形下，给定 $X$ 上的内合成律等价于给定态射 $X\times X\to X$；当 $C=\Ens$ 时，就得到集合上内合成律的通常概念。对于外合成律，当乘积 $Z\times X$ 在 $C$ 中存在时，也有类似结果。

\egapar{8.2.5}
在上述记号下，进一步假设 $X\times X\times X$ 在 $C$ 中存在。由乘积作为表示函子的对象的刻画（8.1.9），存在典范同构
\[
(X\times X)\times X\simeq X\times X\times X\simeq X\times(X\times X).
\]
若将 $X\times X\times X$ 典范地等同于 $(X\times X)\times X$，则对每个 $Y\in C$，映射 $\gamma_X(Y)\times1_{h_X(Y)}$ 等同于 $h_{c_X\times1_X}(Y)$。因此，下面的说法彼此等价：对每个 $Y\in C$，内合成律 $\gamma_X(Y)$ 是结合的；映射图
\[
\begin{tikzcd}[column sep=large]
h_X(Y)\times h_X(Y)\times h_X(Y)\arrow[r,"\gamma_X(Y)\times1"]\arrow[d,"1\times\gamma_X(Y)"']&h_X(Y)\times h_X(Y)\arrow[d,"\gamma_X(Y)"]\\
h_X(Y)\times h_X(Y)\arrow[r,"\gamma_X(Y)"']&h_X(Y)
\end{tikzcd}
\]
交换；态射图
\[
\begin{tikzcd}[column sep=large]
X\times X\times X\arrow[r,"c_X\times1_X"]\arrow[d,"1_X\times c_X"']&X\times X\arrow[d,"c_X"]\\
X\times X\arrow[r,"c_X"']&X
\end{tikzcd}
\]
交换。

\egapar{8.2.6}
在（8.2.5）的假设下，要表达每个 $Y\in C$ 上的内合成律 $\gamma_X(Y)$ 是群律，一方面须表达其结合性，另一方面须表达存在具有群中取逆性质的映射 $\alpha_X(Y):h_X(Y)\to h_X(Y)$。由于对每个态射 $u:Y\to Y'$，已知 $h_X(u)$ 必须是群同态 $h_X(Y')\to h_X(Y)$，首先可见 $\alpha_X:h_X\to h_X$ 必须是函子态射。另一方面，群 $G$ 中取逆映射 $s\mapsto s^{-1}$ 的特征性质可以不借助单位元表达，只需写出下列两个复合映射
\[
\begin{aligned}
(s,t)&\longmapsto(s,s^{-1},t)\longmapsto(s,s^{-1}t)\longmapsto s(s^{-1}t),\\
(s,t)&\longmapsto(s,s^{-1},t)\longmapsto(s,ts^{-1})\longmapsto(ts^{-1})s
\end{aligned}
\]
均等于第二投影 $(s,t)\mapsto t$。由（8.1.3），有 $\alpha_X=h_{a_X}$，其中 $a_X\in\Hom(X,X)$。于是上述第一个条件表示复合态射
\[
X\times X\xrightarrow{(1_X,a_X)\times1_X}X\times X\times X
\xrightarrow{1_X\times c_X}X\times X\xrightarrow{c_X}X
\]
为 $C$ 中的第二投影 $X\times X\to X$；第二个条件也有同样的翻译。\transnote{本段不借助单位元的描述按底本保留。若把其中的等式单独作为集合上的公理，尚须排除空集情形；这里不把这一补充静默加入原文。}

\egapar{8.2.7}
现在假设 $C$ 中存在终对象 $e$（8.1.10）。仍假设每个 $Y\in C$ 上的 $\gamma_X(Y)$ 是 $h_X(Y)$ 的群律，以 $\eta_X(Y)$ 表示该群律的单位元。因为对每个态射 $u:Y\to Y'$，$h_X(u)$ 都为群同态，故
\[
\eta_X(Y)=(h_X(u))(\eta_X(Y')).
\]
特别取 $Y'=e$；这时 $u$ 为 $\Hom(Y,e)$ 中唯一元素 $\varepsilon$，可见 $\eta_X(e)$ 完全确定了所有的 $\eta_X(Y)$。令 $e_X=\eta_X(X)$，即群 $h_X(X)=\Hom(X,X)$ 的单位元。图
\[
\begin{tikzcd}[column sep=large]
h_X(e)\arrow[r,"h_X(\varepsilon)"]\arrow[d,"h_{e_X}(e)"']&h_X(Y)\arrow[d,"h_{e_X}(Y)"]\\
h_X(e)\arrow[r,"h_X(\varepsilon)"']&h_X(Y)
\end{tikzcd}
\]
的交换性（参见8.1.2）表明，在集合 $h_X(Y)$ 上，$h_{e_X}(Y)$ 正是把每个元素送到单位元的映射 $s\mapsto\eta_X(Y)$。于是可验证：对每个 $Y\in C$，$\eta_X(Y)$ 是 $\gamma_X(Y)$ 的单位元，等价于复合态射
\[
X\xrightarrow{(1_X,1_X)}X\times X
\xrightarrow{1_X\times e_X}X\times X\xrightarrow{c_X}X
\]
以及交换其中 $1_X$ 和 $e_X$ 所得的类似复合，都等于 $1_X$。

\egapar{8.2.8}
当然，范畴中代数结构的例子还很容易列举许多。这里对群的例子作了相当详细的处理；今后遇到其他代数结构的例子时，通常将类似考虑留给读者自行展开。\par

% ---- END c0_s08.tex ----
% ---- BEGIN c0_s09.tex ----
\section{可构造集}
\subsection{可构造集}
\src{12--14}
\egastatement{9.1.1}{定义}
连续映射 $f:X\to Y$ 称为\textbf{拟紧}的，若对 $Y$ 的每个拟紧开集 $U$，$f^{-1}(U)$ 都是拟紧的。拓扑空间 $X$ 的子集 $Z$ 称为在 $X$ 中\textbf{逆紧}的（rétrocompact），若典范嵌入 $Z\to X$ 为拟紧映射；换言之，对 $X$ 的每个拟紧开集 $U$，$U\cap Z$ 都拟紧。

$X$ 的闭子集在 $X$ 中逆紧；但 $X$ 的拟紧子集不一定在 $X$ 中逆紧。若 $X$ 拟紧，则在 $X$ 中逆紧的每个开子集都拟紧。显然，有限个在 $X$ 中逆紧的集合之并仍在 $X$ 中逆紧，因为有限个拟紧集之并拟紧。有限个在 $X$ 中逆紧的开集之交，仍是在 $X$ 中逆紧的开集。在局部诺特空间 $X$ 中，每个拟紧集都是诺特子空间，从而 $X$ 的每个子集都在 $X$ 中逆紧。

\egastatement{9.1.2}{定义}
给定拓扑空间 $X$。若 $X$ 的一个子集属于满足下述条件的最小的子集族 $\mathfrak F$，便称它为\textbf{可构造}的：$\mathfrak F$ 包含 $X$ 中所有逆紧开子集，且对有限交与取补集封闭（这也蕴含 $\mathfrak F$ 对有限并封闭）。

\egastatement{9.1.3}{命题}
$X$ 的一个子集可构造，当且仅当它是有限个形如 $U\cap\complement V$ 的集合之并，其中 $U,V$ 为在 $X$ 中逆紧的开集。
\begin{proof}
充分性显然。为证明必要性，考虑所有上述形状的有限并组成的集合族 $\mathfrak G$；只须证明 $\mathfrak G$ 中每个集合的补集仍属于 $\mathfrak G$。设
\[
Z=\bigcup_{i\in I}(U_i\cap\complement V_i),
\]
其中 $I$ 有限，$U_i,V_i$ 都是在 $X$ 中逆紧的开集。则
\[
\complement Z=\bigcap_{i\in I}(V_i\cup\complement U_i).
\]
因此 $\complement Z$ 是一些集合的有限并，而其中每个集合，是若干 $V_i$ 与若干 $\complement U_i$ 的交，故可写为 $V\cap\complement U$，其中 $U$ 是若干 $U_i$ 的并，$V$ 是若干 $V_i$ 的交。前面已经指出，在 $X$ 中逆紧的开集的有限并与有限交，仍是这样的开集，故结论成立。
\end{proof}

\egastatement{9.1.4}{推论}
$X$ 的每个可构造子集都在 $X$ 中逆紧。
\begin{proof}
只须证明：若 $U,V$ 是在 $X$ 中逆紧的开集，则 $U\cap\complement V$ 在 $X$ 中逆紧。事实上，若 $W$ 为 $X$ 的拟紧开集，则 $W\cap U\cap\complement V$ 在拟紧空间 $W\cap U$ 中为闭集，因而拟紧。
\end{proof}
特别地：

\egastatement{9.1.5}{推论}
$X$ 的开子集 $U$ 可构造，当且仅当它在 $X$ 中逆紧。$X$ 的闭子集 $F$ 可构造，当且仅当开集 $\complement F$ 逆紧。

\egapar{9.1.6}
一个重要情形是：$X$ 的每个拟紧开子集都逆紧；换言之，$X$ 的任意两个拟紧开子集之交都拟紧（参见 I，5.5.6）。当 $X$ 自身拟紧时，这意味着 $X$ 中的逆紧开子集恰好就是拟紧开子集，而 $X$ 的可构造子集恰好是有限个形如 $U\cap\complement V$ 的集合之并，其中 $U,V$ 为拟紧开集。

\egastatement{9.1.7}{推论}
诺特空间 $X$ 的一个子集可构造，当且仅当它是 $X$ 的有限个局部闭子集之并。

\egastatement{9.1.8}{命题}
设 $X$ 为拓扑空间，$U$ 是其开子集。
\begin{enumerate}[label=(\roman*)]
\item 若 $T$ 为 $X$ 的可构造子集，则 $T\cap U$ 是 $U$ 的可构造子集。
\item 进一步假设 $U$ 在 $X$ 中逆紧。则 $U$ 的子集 $Z$ 在 $X$ 中可构造，当且仅当它在 $U$ 中可构造。
\end{enumerate}
\begin{proof}
（i）利用（9.1.3），归结为证明：若 $T$ 为在 $X$ 中逆紧的开集，则 $T\cap U$ 为在 $U$ 中逆紧的开集。也就是说，对每个拟紧开集 $W\subset U$，$T\cap U\cap W=T\cap W$ 拟紧；这由假设立即得出。

（ii）必要性由（i）得出，故只须证明充分性。由（9.1.3），只需考虑 $Z$ 为在 $U$ 中逆紧的开集的情形：因为这样便知 $U-Z$ 在 $X$ 中可构造；若 $Z,Z'$ 都是在 $U$ 中逆紧的开集，$Z\cap(U-Z')$ 就也在 $X$ 中可构造。若 $W$ 为 $X$ 的拟紧开集，$Z$ 为在 $U$ 中逆紧的开集，则 $Z\cap W=Z\cap(W\cap U)$；由假设，$W\cap U$ 为 $U$ 的拟紧开集。因此 $W\cap Z$ 拟紧，故 $Z$ 在 $X$ 中为逆紧开集，特别地在 $X$ 中可构造。
\end{proof}

\egastatement{9.1.9}{推论}
设 $X$ 为拓扑空间，$(U_i)_{i\in I}$ 是由在 $X$ 中逆紧的开集组成的有限覆盖。$X$ 的子集 $Z$ 在 $X$ 中可构造，当且仅当对每个 $i\in I$，$Z\cap U_i$ 在 $U_i$ 中可构造。

\egapar{9.1.10}
特别假设 $X$ 拟紧，并且 $X$ 的每一点都有由在 $X$ 中逆紧的开邻域组成的基本邻域系（这些邻域也拟紧）。那么，子集 $Z$ 在 $X$ 中可构造是一个局部性质；也就是说，这等价于：对每个 $x\in X$，存在 $x$ 的开邻域 $V$，使 $V\cap Z$ 在 $V$ 中可构造。事实上，若满足此条件，则由对 $X$ 的假设以及（9.1.8，i），对每个 $x\in X$ 可取一个在 $X$ 中逆紧的开邻域 $V$，使 $V\cap Z$ 在 $V$ 中可构造。再用有限个这样的邻域覆盖 $X$，并应用（9.1.9）即可。

\egastatement{9.1.11}{定义}
设 $X$ 为拓扑空间。若对每个 $x\in X$，存在 $x$ 的开邻域 $V$，使 $T\cap V$ 在 $V$ 中可构造，则称 $X$ 的子集 $T$ 在 $X$ 中\textbf{局部可构造}。

由（9.1.8，i）立即可知：若 $V\cap T$ 在 $V$ 中可构造，则对每个开集 $W\subset V$，$W\cap T$ 在 $W$ 中可构造。若 $T$ 在 $X$ 中局部可构造，则对 $X$ 的每个开集 $U$，$T\cap U$ 在 $U$ 中局部可构造，这也是上述论述的推论。同样的论述表明，$X$ 中局部可构造子集的全体对有限并和有限交封闭；显然，对取补集也封闭。

\egastatement{9.1.12}{命题}
设 $X$ 为拓扑空间。在 $X$ 中可构造的每个集合都局部可构造。若 $X$ 拟紧且其拓扑具有由在 $X$ 中逆紧的集合组成的基，则逆命题也成立。
\begin{proof}
第一断言由定义（9.1.11）得出，第二断言由（9.1.10）得出。
\end{proof}

\egastatement{9.1.13}{推论}
设拓扑空间 $X$ 的拓扑具有由在 $X$ 中逆紧的集合组成的基。则 $X$ 中每个局部可构造子集 $T$ 都在 $X$ 中逆紧。
\begin{proof}
设 $U$ 为 $X$ 的拟紧开集；$T\cap U$ 在 $U$ 中局部可构造，故由（9.1.12），它在 $U$ 中可构造，因而由（9.1.4）拟紧。
\end{proof}

\subsection{诺特空间中的可构造集}
\src{14--16}
\egapar{9.2.1}
已经看到（9.1.7），在诺特空间 $X$ 中，可构造子集就是局部闭子集的有限并。

若 $X'$ 也是诺特空间，则 $X$ 中可构造集在连续映射 $X'\to X$ 下的逆像在 $X'$ 中可构造。若 $Y$ 为诺特空间 $X$ 的可构造子集，则 $Y$ 的子集中，那些作为 $Y$ 的子空间可构造的，恰好就是作为 $X$ 的子空间可构造的。

\egastatement{9.2.2}{命题}
设 $X$ 为不可约诺特空间，$E$ 是 $X$ 的可构造子集。$E$ 在 $X$ 中稠密，当且仅当 $E$ 包含 $X$ 的一个非空开子集。
\begin{proof}
因为每个非空开集都在 $X$ 中稠密，充分性显然。反过来，设
\[
E=\bigcup_{i=1}^{n}(U_i\cap F_i),
\]
其中 $U_i$ 非空开、$F_i$ 闭。则 $\overline E\subset\bigcup_iF_i$。所以若 $\overline E=X$，$X$ 就等于某个 $F_i$，从而 $E\supset U_i$，证明完成。
\end{proof}
当 $X$ 有泛点 $x$（$0_{\mathrm I}$，2.1.2）时，（9.2.2）的条件也等价于 $x\in E$。

\egastatement{9.2.3}{命题}
设 $X$ 为诺特空间。$X$ 的子集 $E$ 可构造，当且仅当对 $X$ 的每个不可约闭子集 $Y$，$E\cap Y$ 或者在 $Y$ 中稀疏，或者包含 $Y$ 的一个非空开子集。
\begin{proof}
必要性由 $E\cap Y$ 须为 $Y$ 的可构造子集及（9.2.2）得出，因为不可约空间 $Y$ 的非稠密子集必为稀疏集（$0_{\mathrm I}$，2.1.1）。为证明充分性，对由满足 $Y\cap E$ 可构造（相对于 $Y$ 或相对于 $X$，二者等价）的闭子集 $Y$ 所成集合 $\mathfrak M$ 应用诺特归纳原理（$0_{\mathrm I}$，2.2.2）。于是可假设，对 $X$ 的每个真闭子集 $Y$，$E\cap Y$ 均可构造。

首先，若 $X$ 不是不可约的，设 $X_i$（$1\leq i\leq m$）为其不可约分支；它们只有有限个（$0_{\mathrm I}$，2.2.5）。由假设，各 $E\cap X_i$ 可构造，故其并 $E$ 也可构造。其次，若 $X$ 不可约，则由假设，或是 $E$ 稀疏，因而 $\overline E\ne X$，于是 $E=E\cap\overline E$ 可构造；或是 $E$ 包含一个非空开集 $U$，此时 $E$ 为 $U$ 与 $E\cap(X-U)$ 之并。$X-U$ 是不同于 $X$ 的闭集，故 $E\cap(X-U)$ 可构造，因此 $E$ 也可构造，证明完成。
\end{proof}

\egastatement{9.2.4}{推论}
设 $X$ 为诺特空间，$(E_\alpha)$ 是 $X$ 的可构造子集的滤过递增族，满足：
\begin{enumerate}[label=\arabic*°]
\item $X$ 为此族 $(E_\alpha)$ 的并。
\item $X$ 的每个不可约闭子集都包含于某个 $E_\alpha$ 的闭包中。
\end{enumerate}
那么存在指标 $\alpha$，使 $X=E_\alpha$。

当每个不可约闭子集都具有泛点时，可以去掉假设2°。
\begin{proof}
对由至少包含于某个 $E_\alpha$ 的闭子集组成的集合 $\mathfrak M$ 应用诺特归纳原理（$0_{\mathrm I}$，2.2.2）。于是可假设，$X$ 的每个真闭子集 $Y$ 都包含于某个 $E_\alpha$。若 $X$ 不是不可约的，结论显然：因为每个不可约分支 $X_i$（$1\leq i\leq m$）都包含于某个 $E_{\alpha_i}$，而存在一个 $E_\alpha$ 包含所有这些 $E_{\alpha_i}$。

现设 $X$ 不可约。由假设，存在 $\beta$ 使 $X=\overline{E_\beta}$；所以由（9.2.2），$E_\beta$ 包含 $X$ 的一个非空开集 $U$。闭集 $X-U$ 包含于某个 $E_\gamma$，只须取同时包含 $E_\beta,E_\gamma$ 的 $E_\alpha$ 即可。若每个不可约闭子集 $Y$ 都有泛点 $y$，则存在 $\alpha$ 使 $y\in E_\alpha$，故 $Y=\overline{\{y\}}\subset\overline{E_\alpha}$，即条件2°是条件1°的推论。
\end{proof}

\egastatement{9.2.5}{命题}
设 $X$ 为诺特空间，$x\in X$，$E$ 为 $X$ 的可构造子集。$E$ 是 $x$ 的邻域，当且仅当对每个包含 $x$ 的不可约闭子集 $Y$，$E\cap Y$ 在 $Y$ 中稠密（若 $Y$ 有泛点 $y$，由9.2.2，这也表示 $y\in E$）。
\begin{proof}
必要性显然，证明充分性。对由满足 $x\in Y$ 且 $E\cap Y$ 是 $x$ 在 $Y$ 中的邻域的闭子集 $Y$ 所成集合 $\mathfrak M$，应用诺特归纳原理。于是可假设：每个包含 $x$ 的真闭子集 $Y\ne X$ 都属于 $\mathfrak M$。

若 $X$ 不是不可约的，则每个包含 $x$ 的不可约分支 $X_i$ 均不同于 $X$，故 $E\cap X_i$ 是 $x$ 相对于 $X_i$ 的邻域。因此 $E$ 是 $x$ 在所有包含 $x$ 的不可约分支之并中的邻域；而此并是 $x$ 在 $X$ 中的邻域，故 $E$ 也是。

若 $X$ 不可约，则由假设，$E$ 在 $X$ 中稠密，故包含一个非空开子集 $U$（9.2.2）。当 $x\in U$ 时结论显然；否则，由归纳假设，$x$ 相对于 $X-U$ 是 $E\cap(X-U)$ 的内点。因此 $X-E$ 在 $X$ 中的闭包不包含 $x$；该闭包的补集就是包含于 $E$ 的 $x$ 的邻域，证明完成。
\end{proof}

\egastatement{9.2.6}{推论}
设 $X$ 为诺特空间，$E$ 为其子集。$E$ 在 $X$ 中开，当且仅当对每个与 $E$ 相交的不可约闭子集 $Y$，$E\cap Y$ 包含 $Y$ 的一个非空开子集。
\begin{proof}
必要性显然。反过来，若此条件成立，由（9.2.3）便知 $E$ 可构造；而（9.2.5）表明 $E$ 是其每一点的邻域，故结论成立。
\end{proof}

\subsection{可构造函数}
\src{16--17}
\egastatement{9.3.1}{定义}
设 $h$ 为从拓扑空间 $X$ 到集合 $T$ 的映射。若每个 $t\in T$ 的逆像 $h^{-1}(t)$ 可构造，且除了有限个 $t$ 外这些逆像均为空集，则称 $h$ 为\textbf{可构造}映射。这时对 $T$ 的任一子集 $S$，$h^{-1}(S)$ 都可构造。若每个 $x\in X$ 都有开邻域 $V$，使 $h|V$ 可构造，则称 $h$ 为\textbf{局部可构造}映射。

每个可构造函数都局部可构造。当 $X$ 拟紧且有由在 $X$ 中逆紧的开集组成的基时（特别当 $X$ 诺特时），逆命题也成立。

\egastatement{9.3.2}{命题}
设 $h$ 为从诺特空间 $X$ 到集合 $T$ 的映射。$h$ 可构造，当且仅当对 $X$ 的每个不可约闭子集 $Y$，存在一个相对于 $Y$ 开的非空子集 $U\subset Y$，使 $h$ 在 $U$ 上为常值。
\begin{proof}
此条件必要：由假设，$h$ 在 $Y$ 上只取有限多个值 $t_i$，且每个 $h^{-1}(t_i)\cap Y$ 在 $Y$ 中可构造（9.2.1）。这些集合不可能全为 $Y$ 中的稀疏集，故至少一个包含非空开集（9.2.3）。

为证明充分性，对由满足 $h|Y$ 可构造的闭子集 $Y$ 所成集合 $\mathfrak M$ 应用诺特归纳原理。于是可假设，对每个真闭子集 $Y\ne X$，$h|Y$ 都可构造。若 $X$ 不是不可约的，则 $h$ 在每个不可约分支 $X_i$ 上的限制都可构造；分支只有有限个，故由定义（9.3.1）立即得到 $h$ 可构造。若 $X$ 不可约，则由假设存在非空开子集 $U$，使 $h$ 在其上为常值；另一方面，$h|X-U$ 由假设可构造，所以 $h$ 可构造。
\end{proof}

\egastatement{9.3.3}{推论}
设 $X$ 是每个不可约闭子集均有泛点的诺特空间。若映射 $h:X\to T$ 满足对每个 $t\in T$，$h^{-1}(t)$ 都可构造，则 $h$ 可构造。
\begin{proof}
若 $Y$ 为 $X$ 的不可约闭子集，$y$ 是其泛点，则 $Y\cap h^{-1}(h(y))$ 可构造且包含 $y$；由（9.2.2），此集包含 $Y$ 的一个非空开子集。应用（9.3.2）即可。
\end{proof}

\egastatement{9.3.4}{命题}
设 $X$ 为每个不可约闭子集均有泛点的诺特空间，$h$ 是从 $X$ 到一个有序集的可构造映射。$h$ 在 $X$ 上上半连续，当且仅当对每个 $x\in X$ 及 $x$ 的每个泛化 $x'$（$0_{\mathrm I}$，2.1.2），有 $h(x')\leq h(x)$。
\begin{proof}
函数 $h$ 只取有限多个值。因此，上半连续表示：对每个 $x\in X$，由满足 $h(y)\leq h(x)$ 的 $y\in X$ 组成的集合 $E$ 是 $x$ 的邻域。由假设，$E$ 是 $X$ 的可构造子集；另一方面，不可约闭子集 $Y$ 包含 $x$，意味着 $Y$ 的泛点 $y$ 是 $x$ 的泛化。结论由（9.2.5）得出。
\end{proof}\par

% ---- END c0_s09.tex ----
% ---- BEGIN c0_s10.tex ----
\section{平坦模的补充知识}
第10.1和10.2小节中只陈述而未证明的性质，其证明请参见 Bourbaki，\textit{Algèbre commutative}，第二章和第三章。
\subsection{平坦模与自由模之间的关系}
\src{17--18}
\egapar{10.1.1}
设 $A$ 为环，$\frj$ 为 $A$ 的理想，$M$ 为 $A$-模。对每个整数 $p\geq0$，有典范的 $(A/\frj)$-模同态
\begin{equation}\tag{10.1.1.1}
\varphi_p:(M/\frj M)\otimes_{A/\frj}(\frj^p/\frj^{p+1})\longrightarrow\frj^pM/\frj^{p+1}M,
\end{equation}
它显然满射。以 $\gr(A)=\bigoplus_{p\geq0}\frj^p/\frj^{p+1}$ 表示 $A$ 对滤过 $(\frj^p)$ 的相伴分次环，以 $\gr(M)=\bigoplus_{p\geq0}\frj^pM/\frj^{p+1}M$ 表示 $M$ 对滤过 $(\frj^pM)$ 的相伴分次 $\gr(A)$-模。于是 $\gr^p(A)=\frj^p/\frj^{p+1}$，$\gr^0(M)=M/\frj M$；各 $\varphi_p$ 给出分次 $\gr(A)$-模的满同态
\begin{equation}\tag{10.1.1.2}
\varphi:\gr^0(M)\otimes_{A/\frj}\gr(A)\longrightarrow\gr(M).
\end{equation}

\egapar{10.1.2}
假设满足下列两种假设之一：
\begin{enumerate}[label=(\roman*)]
\item $\frj$ 幂零；
\item $A$ 诺特，$\frj$ 包含于 $A$ 的根基，且 $M$ 为有限型模。
\end{enumerate}
那么，下列性质等价：
\begin{enumerate}[label=\alph*)]
\item $M$ 是自由 $A$-模。
\item $M/\frj M=M\otimes_A(A/\frj)$ 是自由 $(A/\frj)$-模，且 $\Tor_1^A(M,A/\frj)=0$。
\item $M/\frj M$ 是自由 $(A/\frj)$-模，而且典范同态（10.1.1.2）单射（因而双射）。
\end{enumerate}

\egapar{10.1.3}
假设 $A/\frj$ 为域（即 $\frj$ 为极大理想），且满足（10.1.2）的假设（i）或（ii）之一。那么，下列性质等价：
\begin{enumerate}[label=\alph*)]
\item $M$ 是自由 $A$-模。
\item $M$ 是投射 $A$-模。
\item $M$ 是平坦 $A$-模。
\item $\Tor_1^A(M,A/\frj)=0$。
\item 典范同态（10.1.1.2）为双射。
\end{enumerate}
此结果特别适用于下列两种情形：
\begin{enumerate}[label=(\roman*)]
\item $M$ 是局部环 $A$ 上的任意模，而 $A$ 的极大理想 $\frj$ 幂零（例如 $A$ 是阿廷局部环）。
\item $M$ 是诺特局部环上的有限型模。
\end{enumerate}

\subsection{平坦性的局部判别准则}
\src{18--20}
\egapar{10.2.1}
假设与记号同（10.1.1），考虑以下条件：
\begin{enumerate}[label=\alph*)]
\item $M$ 是平坦 $A$-模。
\item $M/\frj M$ 是平坦 $(A/\frj)$-模，且 $\Tor_1^A(M,A/\frj)=0$。
\item $M/\frj M$ 是平坦 $(A/\frj)$-模，而且典范同态（10.1.1.2）为双射。
\item 对每个 $n\geq1$，$M/\frj^nM$ 是平坦 $(A/\frj^n)$-模。
\end{enumerate}
有蕴含关系
\[
\mathrm{a)}\Longrightarrow\mathrm{b)}\Longrightarrow\mathrm{c)}\Longrightarrow\mathrm{d)}.
\]
若 $\frj$ 幂零，则 a）、b）、c）、d）四个条件等价。如果 $A$ 诺特且 $M$ \textbf{理想分离}，也有相同结论；这里理想分离是指：对 $A$ 的每个理想 $\mathfrak a$，$A$-模 $\mathfrak a\otimes_A M$ 对 $\frj$-预进拓扑分离。

\egapar{10.2.2}
设 $A$ 为诺特环，$B$ 为诺特交换 $A$-代数，$\frj$ 为 $A$ 的理想且 $\frj B$ 包含于 $B$ 的根基，$M$ 为有限型 $B$-模。那么，当把 $M$ 看成 $A$-模时，（10.2.1）的条件 a）、b）、c）、d）等价。

此结果尤其适用于 $A,B$ 均为诺特局部环而 $A\to B$ 为局部同态的情形。更特别地，若 $\frj$ 为 $A$ 的极大理想，则在条件 b）与 c）中，可以去掉 $M/\frj M$ 平坦的假设，因为它自动满足；而条件 d）表示各 $M/\frj^nM$ 都是 $A/\frj^n$ 上的自由模。

\egapar{10.2.3}
对 $A,B,\frj,M$ 的假设同（10.2.2）开头。设 $\widehat A$ 为 $A$ 对 $\frj$-预进拓扑的分离完备化，$\widehat M$ 为 $M$ 对 $\frj B$-预进拓扑的分离完备化。那么，$M$ 为平坦 $A$-模，当且仅当 $\widehat M$ 为平坦 $\widehat A$-模。

\egapar{10.2.4}
设 $\rho:A\to B$ 为诺特局部环的局部同态，$k$ 为 $A$ 的剩余域，$M,N$ 为有限型 $B$-模，并假设 $N$ 是 $A$-平坦的。设 $u:M\to N$ 为 $B$-同态。那么以下条件等价：
\begin{enumerate}[label=\alph*)]
\item $u$ 单射，且 $\Coker(u)$ 为平坦 $A$-模。
\item $u\otimes1:M\otimes_Ak\to N\otimes_Ak$ 单射。
\end{enumerate}

\egapar{10.2.5}
设 $\rho:A\to B$、$\sigma:B\to C$ 为诺特局部环的局部同态，$k$ 为 $A$ 的剩余域，$M$ 为有限型 $C$-模。假设 $B$ 为平坦 $A$-模。那么以下条件等价：
\begin{enumerate}[label=\alph*)]
\item $M$ 为平坦 $B$-模。
\item $M$ 为平坦 $A$-模，且 $M\otimes_Ak$ 为平坦 $(B\otimes_Ak)$-模。
\end{enumerate}

\egastatement{10.2.6}{命题}
设 $A,B$ 为诺特局部环，$\rho:A\to B$ 为局部同态，$\frj$ 为包含于 $B$ 的极大理想的一个理想，$M$ 为有限型 $B$-模。假设对每个 $n\geq0$，$M_n=M/\frj^{n+1}M$ 都为平坦 $A$-模。那么 $M$ 为平坦 $A$-模。
\begin{proof}
须证明：对有限型 $A$-模之间的每个单同态 $u:N'\to N$，$v=1\otimes u:M\otimes_AN'\to M\otimes_AN$ 单射。$M\otimes_AN'$ 与 $M\otimes_AN$ 都为有限型 $B$-模，故对 $\frj$-预进拓扑分离（$0_{\mathrm I}$，7.3.5）。所以只须证明分离完备化之间的同态
\[
\widehat v:(M\otimes_AN')^{\wedge}\longrightarrow(M\otimes_AN)^{\wedge}
\]
单射。这里 $\widehat v=\plim v_n$，其中 $v_n$ 为同态 $1\otimes u:M_n\otimes_AN'\to M_n\otimes_AN$。由假设，$M_n$ 是 $A$-平坦的，故每个 $v_n$ 都单射；函子 $\plim$ 左正合，因此 $\widehat v$ 也单射。
\end{proof}

\egastatement{10.2.7}{推论}
设 $A$ 为诺特环，$B$ 为诺特局部环，$\rho:A\to B$ 为同态，$f$ 为 $B$ 的极大理想中的元素，$M$ 为有限型 $B$-模。假设 $M$ 的位似 $f_M:x\mapsto fx$ 单射，且 $M/fM$ 为平坦 $A$-模。那么 $M$ 为平坦 $A$-模。
\begin{proof}
对 $i\geq0$，令 $M_i=f^iM$。因为 $f_M$ 单射，$M_i/M_{i+1}$ 同构于 $M/fM$，故对每个 $i\geq0$ 都是 $A$-平坦的。从正合列
\[
0\longrightarrow M_i/M_{i+1}\longrightarrow M/M_{i+1}\longrightarrow M/M_i\longrightarrow0
\]
出发，对 $i$ 归纳，得到每个 $M/M_i$ 都为 $A$-平坦的（$0_{\mathrm I}$，6.1.2）。于是可以应用（10.2.6）。\transnote{底本此处直接引用（10.2.6），而该命题假设 $A$ 局部且同态局部；本推论只假设 $A$ 诺特。若采用这一路证明，需要先在 $B$ 的极大理想的逆像处局部化 $A$。下文是原著另给的直接证明，不依赖这一步省略。}
也可直接作如下论证。对每个有限型 $A$-模 $N$，$M\otimes_AN$ 为有限型 $B$-模。因为 $f$ 属于 $B$ 的根基 $\nn$，$M\otimes_AN$ 上的 $(f)$-进拓扑比 $\nn$-进拓扑细，而后者已知是分离的（$0_{\mathrm I}$，7.3.5）。又因 $M/M_i$ 是 $A$-平坦的，有
\[
\begin{aligned}
f^i(M\otimes_AN)
&=\Im(M_i\otimes_AN\longrightarrow M\otimes_AN)\\
&=\Ker(M\otimes_AN\longrightarrow(M/M_i)\otimes_AN)
\end{aligned}
\]
（$0_{\mathrm I}$，6.1.2）。取有限型 $A$-模 $N$、其子模 $N'$，以 $j:N'\to N$ 为典范嵌入。在交换图
\[
\begin{tikzcd}[column sep=large]
M\otimes_AN'\arrow[r]\arrow[d,"M\otimes j"']&(M/M_i)\otimes_AN'\arrow[d,"(M/M_i)\otimes j"]\\
M\otimes_AN\arrow[r]&(M/M_i)\otimes_AN
\end{tikzcd}
\]
中，因为 $M/M_i$ 为 $A$-平坦，右侧箭头单射。于是对任意 $i$ 都有
\[
\Ker(M\otimes_AN'\to M\otimes_AN)
\subset\Ker(M\otimes_AN'\to(M/M_i)\otimes_AN').
\]
如上所见，右端各核的交为零，左端也只能为零；所以 $M$ 为 $A$-平坦的。
\end{proof}

\egastatement{10.2.8}{命题}
设 $A$ 为约化诺特环，$M$ 为有限型 $A$-模。假设对每个同时为离散赋值环的 $A$-代数 $B$，$M\otimes_AB$ 都为平坦 $B$-模（因而为自由模，10.1.3）。那么 $M$ 为平坦 $A$-模。
\begin{proof}
已知 $M$ 平坦，当且仅当对 $A$ 的每个极大理想 $\mm$，$M_\mm$ 为平坦 $A_\mm$-模（$0_{\mathrm I}$，6.3.3）；所以可以归结到 $A$ 局部的情形（$0_{\mathrm I}$，1.2.8）。设 $\mm$ 为 $A$ 的极大理想，$\pp_i$（$1\leq i\leq r$）为其极小素理想，$k=A/\mm$ 为剩余域。已知（II，7.1.7），对每个 $i$，存在一个离散赋值环 $B_i$，它支配整环 $A/\pp_i$，并与后者有相同的分式域 $K_i$。令 $M_i=M\otimes_AB_i$。由假设，$M_i$ 在 $B_i$ 上自由；因此，以 $k_i$ 为 $B_i$ 的剩余域，有
\begin{equation}\tag{10.2.8.1}
\rk_{k_i}(M_i\otimes_{B_i}k_i)=\rk_{K_i}(M_i\otimes_{B_i}K_i).
\end{equation}
显然复合同态 $A\to A/\pp_i\to B_i$ 为局部同态，故 $k_i$ 是 $k$ 的扩张，并且
\[
M_i\otimes_{B_i}k_i=M\otimes_Ak_i=(M\otimes_Ak)\otimes_kk_i,
\qquad M_i\otimes_{B_i}K_i=M\otimes_AK_i.
\]
等式（10.2.8.1）于是给出
\[
\rk_k(M\otimes_Ak)=\rk_{K_i}(M\otimes_AK_i)\qquad(1\leq i\leq r).
\]
因为 $A$ 约化，已知此条件推出 $M$ 为自由 $A$-模（Bourbaki，\textit{Algèbre commutative}，第二章，§3，第2小节，命题7）。
\end{proof}

\subsection{局部环的平坦扩张的存在性}
\src{20--23}
\egastatement{10.3.1}{命题}
设 $A$ 为诺特局部环，$\frj$ 为其极大理想，$k=A/\frj$ 为其剩余域。设 $K$ 为 $k$ 的扩张。则存在从 $A$ 到某个诺特局部环 $B$ 的局部同态，使 $B/\frj B$ 与 $K$ 作为 $k$-代数同构，且 $B$ 为平坦 $A$-模。

我们分几个步骤证明此命题。

\egapar{10.3.1.1}
首先假设 $K=k(T)$，其中 $T$ 为不定元。在多项式环 $A'=A[T]$ 中，考虑由所有系数均属于 $\frj$ 的多项式组成的素理想 $\frj'=\frj A'$；显然 $A'/\frj'$ 典范地同构于 $k[T]$。证明分式环 $B=A'_{\frj'}$ 符合要求：它显然为诺特局部环，极大理想为 $\frk=\frj B$。此外，
\[
B/\frk=(A'/\frj')_{\frj'}=(k[T])_{(0)}
\]
就是 $k[T]$ 的分式域 $K$。最后，$B$ 为平坦 $A'$-模，$A'$ 为自由 $A$-模，故 $B$ 为平坦 $A$-模（$0_{\mathrm I}$，6.2.1）。

\egapar{10.3.1.2}
其次假设 $K=k(t)=k[t]$，其中 $t$ 在 $k$ 上代数；设 $f\in k[T]$ 为 $t$ 的极小多项式。存在首一多项式 $F\in A[T]$，其在 $k[T]$ 中的典范像为 $f$。仍令 $A'=A[T]$，并令 $\frj'=\frj A'+(F)$。这次我们证明商环 $B=A'/(F)$ 符合要求。首先，显然 $B$ 为自由 $A$-模，因而平坦。环 $A'/\frj'$ 同构于
\[
(A'/\frj A')\big/\bigl((\frj A'+(F))/\frj A'\bigr)=k[T]/(f)=K.
\]
所以 $\frj'$ 在 $B$ 中的像 $\frk$ 为极大理想，且显然 $\frk=\frj B$。最后，$B$ 为半局部环，因为它是有限型 $A$-模（Bourbaki，\textit{Algèbre commutative}，第四章，§2，第5小节，命题9的推论3）；其极大理想与 $B/\frj B$ 的极大理想一一对应（[13]，第一卷，第259页）。因此上述论述表明 $B$ 为局部环。

\egastatement{10.3.1.3}{引理}
设 $(A_\lambda,f_{\mu\lambda})$ 为局部环的滤过归纳系，且各 $f_{\mu\lambda}$ 均为局部同态。令 $\mm_\lambda$ 为 $A_\lambda$ 的极大理想，$K_\lambda=A_\lambda/\mm_\lambda$。则 $A'=\ilim A_\lambda$ 为局部环，其极大理想为 $\mm'=\ilim\mm_\lambda$，剩余域为 $K=\ilim K_\lambda$。此外，若 $\lambda<\mu$ 时有 $\mm_\mu=\mm_\lambda A_\mu$，则对每个 $\lambda$，有 $\mm'=\mm_\lambda A'$。若再假设当 $\lambda<\mu$ 时 $A_\mu$ 为平坦 $A_\lambda$-模，且各 $A_\lambda$ 均诺特，则 $A'$ 诺特，且对每个 $\lambda$ 都为平坦 $A_\lambda$-模。
\begin{proof}
由假设，当 $\lambda<\mu$ 时，$f_{\mu\lambda}(\mm_\lambda)\subset\mm_\mu$；故各 $\mm_\lambda$ 组成归纳系，其极限 $\mm'$ 显然为 $A'$ 的理想。若 $x'\notin\mm'$，存在 $\lambda$ 及 $x_\lambda\in A_\lambda$，使 $x'=f_\lambda(x_\lambda)$，其中 $f_\lambda:A_\lambda\to A'$ 是典范同态。因为 $x'\notin\mm'$，必有 $x_\lambda\notin\mm_\lambda$，故 $x_\lambda$ 在 $A_\lambda$ 中有逆元 $y_\lambda$，而 $y'=f_\lambda(y_\lambda)$ 为 $x'$ 在 $A'$ 中的逆元。这就证明了 $A'$ 为以 $\mm'$ 为极大理想的局部环。关于 $K$ 的断言由 $\ilim$ 的正合性立即得到。

假设 $\mm_\mu=\mm_\lambda A_\mu$ 表示典范映射 $\mm_\lambda\otimes_{A_\lambda}A_\mu\to\mm_\mu$ 满射。因此，由 $\ilim$ 的正合性及其与张量积的可交换性，得到 $\mm'=\mm_\lambda A'$。

现在假设当 $\lambda<\mu$ 时有 $\mm_\mu=\mm_\lambda A_\mu$，并且 $A_\mu$ 为平坦 $A_\lambda$-模。由（$0_{\mathrm I}$，6.2.3），$A'$ 对每个 $\lambda$ 都为平坦 $A_\lambda$-模。因为 $A',A_\lambda$ 均为局部环且 $\mm'=\mm_\lambda A'$，$A'$ 甚至是忠实平坦 $A_\lambda$-模（$0_{\mathrm I}$，6.6.2）。最后再假设所有 $A_\lambda$ 诺特；这时各 $\mm_\lambda$-预进拓扑分离（$0_{\mathrm I}$，7.3.5）。先证明这推出 $A'$ 上的 $\mm'$-进拓扑分离。若 $x'\in A'$ 属于每个 $\mm'^n$（$n>0$），它是某个 $x_\mu\in A_\mu$ 的像。$\mm'^n=\mm_\mu^nA'$ 在 $A_\mu$ 中的逆像为 $\mm_\mu^n$（$0_{\mathrm I}$，6.6.1），故 $x_\mu$ 属于所有 $\mm_\mu^n$；于是 $x_\mu=0$，从而 $x'=0$。

以 $\widehat{A'}$ 表示 $A'$ 对 $\mm'$-进拓扑的完备化。由上可知 $A'\subset\widehat{A'}$。我们将证明 $\widehat{A'}$ 诺特，且对每个 $\lambda$ 都为 $A_\lambda$-平坦。由此将推出 $\widehat{A'}$ 为 $A'$-平坦（$0_{\mathrm I}$，6.2.3）；又因 $\mm'\widehat{A'}\ne\widehat{A'}$，它为忠实平坦 $A'$-模（$0_{\mathrm I}$，6.6.2），最终便可推出 $A'$ 诺特（$0_{\mathrm I}$，6.5.2），从而完成引理的证明。

有 $\widehat{A'}=\plim_n A'/\mm'^n$。因为 $A'$ 为 $A_\lambda$-平坦，
\[
\begin{aligned}
\mm'^n/\mm'^{n+1}
&=(\mm_\lambda^n/\mm_\lambda^{n+1})\otimes_{A_\lambda}A'\\
&=(\mm_\lambda^n/\mm_\lambda^{n+1})\otimes_{K_\lambda}(A'/\mm')\\
&=(\mm_\lambda^n/\mm_\lambda^{n+1})\otimes_{K_\lambda}K.
\end{aligned}
\]
$\mm_\lambda^n/\mm_\lambda^{n+1}$ 是有限维 $K_\lambda$-向量空间，故每个 $\mm'^n/\mm'^{n+1}$ 都为有限维 $K$-向量空间。由（$0_{\mathrm I}$，7.2.12）及（$0_{\mathrm I}$，7.2.8），$\widehat{A'}$ 诺特。又已知 $\widehat{A'}$ 的极大理想为 $\mm'\widehat{A'}$，而且 $\widehat{A'}/\mm'^n\widehat{A'}$ 同构于 $A'/\mm'^n$。因
\[
A'/\mm'^n=(A_\lambda/\mm_\lambda^n)\otimes_{A_\lambda}A',
\]
$A'/\mm'^n$ 为平坦 $(A_\lambda/\mm_\lambda^n)$-模（$0_{\mathrm I}$，6.2.1）。因此可以对诺特 $A_\lambda$-代数 $\widehat{A'}$ 应用判别准则（10.2.2），从而得到 $\widehat{A'}$ 为 $A_\lambda$-平坦。证毕。
\end{proof}

\egapar{10.3.1.4}
现在处理一般情形。存在一个序数 $\gamma$，并对每个序数 $\lambda\leq\gamma$ 存在 $K$ 的含 $k$ 的子域 $k_\lambda$，满足：1° 对每个 $\lambda<\gamma$，$k_{\lambda+1}$ 是由一个元素生成的 $k_\lambda$ 的扩张；2° 对每个无前驱的序数 $\mu$，$k_\mu=\bigcup_{\lambda<\mu}k_\lambda$；3° $K=k_\gamma$。事实上，只须取从所有 $\xi\leq\beta$ 的序数组成的集合到 $K$ 的一个双射 $\xi\mapsto t_\xi$（$\beta$ 适当），然后用超限归纳定义各 $k_\lambda$（$\lambda\leq\beta$）：若 $\lambda$ 无前驱，取为所有 $\mu<\lambda$ 的 $k_\mu$ 之并；若 $\lambda=\nu+1$，则取 $k_\lambda=k_\nu(t_\xi)$，其中 $\xi$ 为使 $t_\xi\notin k_\nu$ 的最小序数。按定义，$\gamma$ 就是满足 $k_\gamma=K$ 的最小序数 $\gamma\leq\beta$。

现在以超限归纳定义诺特局部环族 $A_\lambda$（$\lambda\leq\gamma$），以及对 $\lambda\leq\mu$ 的局部同态 $f_{\mu\lambda}:A_\lambda\to A_\mu$，使满足：
\begin{enumerate}[label=(\roman*)]
\item $(A_\lambda,f_{\mu\lambda})$ 为归纳系，且 $A_0=A$。
\item 对每个 $\lambda$，有 $k$-同构 $A_\lambda/\frj A_\lambda\simeq k_\lambda$。
\item 对 $\lambda\leq\mu$，$A_\mu$ 是平坦 $A_\lambda$-模。
\end{enumerate}
假定 $\lambda<\mu<\zeta$ 时的 $A_\lambda$ 与 $f_{\mu\lambda}$ 已定义。首先设 $\zeta=\xi+1$，故 $k_\zeta=k_\xi(t)$。若 $t$ 在 $k_\xi$ 上超越，就按（10.3.1.1）的方法令
\[
A_\zeta=(A_\xi[T])_{\frj A_\xi[T]}.
\]
以典范同态为 $f_{\zeta\xi}$；当 $\lambda<\xi$ 时，令 $f_{\zeta\lambda}=f_{\zeta\xi}\circ f_{\xi\lambda}$。由（10.3.1.1）的证明，立即可验证条件（i）至（iii）。若 $t$ 在 $k_\xi$ 上代数，设 $h$ 为其在 $k_\xi[T]$ 中的极小多项式，$H$ 为 $A_\xi[T]$ 中在 $k_\xi[T]$ 内的像为 $h$ 的首一多项式；取 $A_\zeta=A_\xi[T]/(H)$，各 $f_{\zeta\lambda}$ 同前定义。条件（i）至（iii）的验证由（10.3.1.2）得出。

现设 $\zeta$ 无前驱。以局部环归纳系 $(A_\lambda,f_{\mu\lambda})_{\lambda<\mu<\zeta}$ 的归纳极限为 $A_\zeta$；当 $\lambda<\zeta$ 时，$f_{\zeta\lambda}$ 定义为典范同态。$A_\zeta$ 为诺特局部环、各 $f_{\zeta\lambda}$ 为局部同态，以及对 $\lambda\leq\zeta$ 的条件（i）至（iii），都由归纳假设及引理（10.3.1.3）得出。完成此构造后，显然 $B=A_\gamma$ 满足（10.3.1）的断言。

注意，由（10.2.1，c）），有典范同构
\begin{equation}\tag{10.3.1.5}
\gr(A)\otimes_kK\simeq\gr(B).
\end{equation}
另一方面，可以把 $B$ 换成它的 $\frj B$-进完备化 $\widehat B$ 而不改变（10.3.1）的结论，因为 $\widehat B$ 为平坦 $B$-模（$0_{\mathrm I}$，7.3.3），从而为平坦 $A$-模（$0_{\mathrm I}$，6.2.1）。

此外，还证明了：
\egastatement{10.3.2}{推论}
若 $K$ 是有限次扩张，可以假设 $B$ 为有限 $A$-代数。\par

% ---- END c0_s10.tex ----
% ---- BEGIN c0_s11a.tex ----
\section{同调代数的补充知识}
\subsection{谱序列回顾}
\src{23--27}
\egapar{11.1.1}
今后要使用比（T，2.4）所定义的更一般的谱序列概念。沿用（T，2.4）的记号，把由下列数据组成的系 $E$ 称为阿贝尔范畴 $C$ 中的一个\textbf{谱序列}：
\begin{enumerate}[label=\alph*)]
\item 一族 $C$ 的对象 $(E_r^{pq})$，对 $p,q\in\mathbf Z$ 及 $r\geq2$ 定义。
\item 一族态射 $d_r^{pq}:E_r^{pq}\to E_r^{p+r,q-r+1}$，满足
\[
d_r^{p+r,q-r+1}d_r^{pq}=0.
\]
令 $Z_{r+1}(E_r^{pq})=\Ker(d_r^{pq})$，$B_{r+1}(E_r^{pq})=\Im(d_r^{p+r,q-r+1})$，于是
\[
B_{r+1}(E_r^{pq})\subset Z_{r+1}(E_r^{pq})\subset E_r^{pq}.
\]
\transnote{底本在定义 $B_{r+1}(E_r^{pq})$ 时确实将微分上标印作 $p+r,q-r+1$。若要使其像成为 $E_r^{pq}$ 的子对象，上标应为 $p-r,q+r-1$。正文保留底本印式，以本注标明疑误。}
\item 一族同构
\[
\alpha_r^{pq}:Z_{r+1}(E_r^{pq})/B_{r+1}(E_r^{pq})\iso E_{r+1}^{pq}.
\]
\end{enumerate}
当 $k\geq r+1$ 时，再对 $k$ 归纳定义子对象 $B_k(E_r^{pq})$ 与 $Z_k(E_r^{pq})$：它们分别是在典范态射 $E_r^{pq}\to E_r^{pq}/B_{r+1}(E_r^{pq})$ 下，对应于商对象中由 $\alpha_r^{pq}$ 等同为 $B_k(E_{r+1}^{pq})$ 与 $Z_k(E_{r+1}^{pq})$ 的子对象的逆像。显然，在同构意义下，有
\begin{equation}\tag{11.1.1.1}
Z_k(E_r^{pq})/B_k(E_r^{pq})=E_k^{pq}\qquad(k\geq r+1).
\end{equation}
若再令 $B_r(E_r^{pq})=0$，$Z_r(E_r^{pq})=E_r^{pq}$，则有包含关系
\begin{equation}\tag{11.1.1.2}
\begin{aligned}
0=B_r(E_r^{pq})&\subset B_{r+1}(E_r^{pq})\subset B_{r+2}(E_r^{pq})\subset\cdots\\
&\cdots\subset Z_{r+2}(E_r^{pq})\subset Z_{r+1}(E_r^{pq})\subset Z_r(E_r^{pq})=E_r^{pq}.
\end{aligned}
\end{equation}
系 $E$ 的其余数据如下：
\begin{enumerate}[label=\alph*),start=4]
\item $E_2^{pq}$ 的两个子对象 $B_\infty(E_2^{pq})$、$Z_\infty(E_2^{pq})$，使
\[
B_\infty(E_2^{pq})\subset Z_\infty(E_2^{pq}),
\]
而且对每个 $k\geq2$，
\[
B_k(E_2^{pq})\subset B_\infty(E_2^{pq}),\qquad
Z_\infty(E_2^{pq})\subset Z_k(E_2^{pq}).
\]
令
\begin{equation}\tag{11.1.1.3}
E_\infty^{pq}=Z_\infty(E_2^{pq})/B_\infty(E_2^{pq}).
\end{equation}
\item 一族 $C$ 的对象 $(E^n)$，每个都带有递降滤过 $(F^p(E^n))_{p\in\mathbf Z}$。如常，以 $\gr(E^n)$ 表示滤过对象 $E^n$ 的相伴分次对象，即各 $\gr_p(E^n)=F^p(E^n)/F^{p+1}(E^n)$ 的直和。
\item 对每对 $(p,q)\in\mathbf Z\times\mathbf Z$，给定同构 $\beta^{pq}:E_\infty^{pq}\iso\gr_p(E^{p+q})$。
\end{enumerate}
不计各滤过时，对象族 $(E^n)$ 称为谱序列 $E$ 的\textbf{终项}。

假设范畴 $C$ 具有无限直和，或者对每个 $r\geq2$ 和 $n\in\mathbf Z$，满足 $p+q=n$ 且 $E_r^{pq}\ne0$ 的偶对 $(p,q)$ 只有有限个（只须对 $r=2$ 如此）。那么 $E_r^n=\bigoplus_{p+q=n}E_r^{pq}$ 有定义。以 $d_r^n$ 表示从 $E_r^n$ 到 $E_r^{n+1}$ 的态射，其在每个满足 $p+q=n$ 的分量 $E_r^{pq}$ 上的限制为 $d_r^{pq}$。则 $d_r^{n+1}d_r^n=0$，即 $(E_r^n)_{n\in\mathbf Z}$ 构成 $C$ 中微分算子次数为 $+1$ 的复形 $E_r^\bullet$；由 c），$\mathrm H^n(E_r^\bullet)$ 对每个 $r\geq2$ 都同构于 $E_{r+1}^n$。

\egapar{11.1.2}
从谱序列 $E$ 到谱序列 $E'=(E_r'{}^{pq},E'^n)$ 的\textbf{态射} $u:E\to E'$，由一族态射
\[
u_r^{pq}:E_r^{pq}\to E_r'{}^{pq},\qquad u^n:E^n\to E'^n
\]
组成，其中各 $u^n$ 与 $E^n,E'^n$ 的滤过相容，并且下图交换：
\[
\begin{tikzcd}[column sep=large]
E_r^{pq}\arrow[r,"d_r^{pq}"]\arrow[d,"u_r^{pq}"']&E_r^{p+r,q-r+1}\arrow[d,"u_r^{p+r,q-r+1}"]\\
E_r'{}^{pq}\arrow[r,"d_r'{}^{pq}"']&E_r'{}^{p+r,q-r+1}.
\end{tikzcd}
\]
此外，经由商对象，$u_r^{pq}$ 给出态射
\[
\overline u_r^{pq}:Z_{r+1}(E_r^{pq})/B_{r+1}(E_r^{pq})\longrightarrow
Z_{r+1}(E_r'{}^{pq})/B_{r+1}(E_r'{}^{pq}),
\]
须有 $\alpha_r'{}^{pq}\circ\overline u_r^{pq}=u_{r+1}^{pq}\circ\alpha_r^{pq}$。最后须有
\[
u_2^{pq}(B_\infty(E_2^{pq}))\subset B_\infty(E_2'{}^{pq}),\qquad
u_2^{pq}(Z_\infty(E_2^{pq}))\subset Z_\infty(E_2'{}^{pq}).
\]
经由商对象，$u_2^{pq}$ 给出态射 $u_\infty^{pq}:E_\infty^{pq}\to E_\infty'{}^{pq}$，且下图必须交换：
\[
\begin{tikzcd}[column sep=large]
E_\infty^{pq}\arrow[r,"u_\infty^{pq}"]\arrow[d,"\beta^{pq}"']&E_\infty'{}^{pq}\arrow[d,"\beta'{}^{pq}"]\\
\gr_p(E^{p+q})\arrow[r,"\gr_p(u^{p+q})"']&\gr_p(E'^{p+q}).
\end{tikzcd}
\]
由上述定义，对 $r$ 归纳可知：若各 $u_2^{pq}$ 为同构，则对 $r\geq2$，各 $u_r^{pq}$ 都为同构。若进一步知道
\[
u_2^{pq}(B_\infty(E_2^{pq}))=B_\infty(E_2'{}^{pq}),\qquad
u_2^{pq}(Z_\infty(E_2^{pq}))=Z_\infty(E_2'{}^{pq}),
\]
且各 $u^n$ 为同构，就可以断定 $u$ 为同构。

\egapar{11.1.3}
回顾一下，对对象 $X\in C$ 的递降滤过 $(F^p(X))_{p\in\mathbf Z}$，若 $\inf_pF^p(X)=0$，称滤过\textbf{分离}；若存在 $p$ 使 $F^p(X)=0$，称\textbf{离散}；若 $\sup_pF^p(X)=X$，称\textbf{穷尽}（或\textbf{余分离}）；若存在 $p$ 使 $F^p(X)=X$，称\textbf{余离散}。

若
\[
B_\infty(E_2^{pq})=\sup_k B_k(E_2^{pq}),\qquad
Z_\infty(E_2^{pq})=\inf_k Z_k(E_2^{pq}),
\]
称谱序列 $E=(E_r^{pq},E^n)$ \textbf{弱收敛}（也就是说，$B_\infty$ 和 $Z_\infty$ 由 $E$ 的数据 a）至 c）决定）。若 $E$ 弱收敛，并满足下列条件，就称其\textbf{正则}：
\begin{enumerate}[label=\arabic*°]
\item 对每个偶对 $(p,q)$，递降列 $(Z_k(E_2^{pq}))_{k\geq2}$ 稳定。弱收敛性这时蕴含，当 $k$ 足够大时（依赖于 $p,q$），$Z_k(E_2^{pq})=Z_\infty(E_2^{pq})$。
\item 对每个 $n$，$E^n$ 的滤过 $(F^p(E^n))$ 离散且穷尽。
\end{enumerate}
若 $E$ 弱收敛并满足以下条件，则称其\textbf{余正则}：
\begin{enumerate}[label=\arabic*°,start=3]
\item 对每个 $(p,q)$，递增列 $(B_k(E_2^{pq}))$ 稳定，这蕴含 $B_\infty(E_2^{pq})=B_k(E_2^{pq})$，从而 $E_\infty^{pq}=\inf_k E_k^{pq}$。
\item 对每个 $n$，$E^n$ 的滤过余离散。
\end{enumerate}
最后，若 $E$ 同时正则且余正则，称其\textbf{双正则}。换言之，满足以下条件：
\begin{enumerate}[label=\alph*)]
\item 对每个 $(p,q)$，列 $(B_r(E_2^{pq}))_{r\geq2}$ 与 $(Z_r(E_2^{pq}))_{r\geq2}$ 均稳定，且当 $r$ 足够大时，$B_r(E_2^{pq})=B_\infty(E_2^{pq})$、$Z_r(E_2^{pq})=Z_\infty(E_2^{pq})$（这蕴含 $E_r^{pq}=E_\infty^{pq}$）。
\item 对每个 $n$，滤过 $(F^p(E^n))$ 离散且余离散（也称此滤过\textbf{有限}）。
\end{enumerate}
因此，（T，2.4）定义的谱序列就是双正则谱序列。

\egapar{11.1.4}
假设范畴 $C$ 中存在滤过归纳极限，而且函子 $\ilim$ 正合（这等价于（T，1.5）的公理 AB5）成立，参见 T，1.8）。对象 $X\in C$ 的滤过 $(F^p(X))$ 穷尽这件事，也可写为 $\ilim_{p\to-\infty}F^p(X)=X$。若谱序列 $E$ 弱收敛，则
\[
B_\infty(E_2^{pq})=\ilim_{k\to\infty}B_k(E_2^{pq}).
\]
若再有从 $E$ 到 $C$ 中另一个弱收敛谱序列 $E'$ 的态射 $u:E\to E'$，则由 $\ilim$ 的正合性，有
\[
u_2^{pq}(B_\infty(E_2^{pq}))=B_\infty(E_2'{}^{pq}).
\]
\transnote{此等号按底本保留。仅假设 $u$ 为态射时，一般只能得到包含关系；下一命题另有各 $u_2^{pq}$ 为同构的假设，在该假设下可得到这里需要的等号。}
此外：

\egastatement{11.1.5}{命题}
设 $C$ 为滤过归纳极限正合的阿贝尔范畴，$E,E'$ 为 $C$ 中两个正则谱序列，$u:E\to E'$ 为谱序列态射。若各 $u_2^{pq}$ 为同构，则 $u$ 也为同构。
\begin{proof}
已知（11.1.2）各 $u_r^{pq}$ 为同构，且
\[
u_2^{pq}(B_\infty(E_2^{pq}))=B_\infty(E_2'{}^{pq}).
\]
$E,E'$ 正则的假设又推出
\[
u_2^{pq}(Z_\infty(E_2^{pq}))=Z_\infty(E_2'{}^{pq}).
\]
因 $u_2^{pq}$ 为同构，$u_\infty^{pq}$ 也为同构，故 $\gr_p(u^{p+q})$ 为同构。$E^n,E'^n$ 的滤过均离散且穷尽，所以各 $u^n$ 也为同构（Bourbaki，\textit{Algèbre commutative}，第三章，§2，第8小节，定理1）。
\end{proof}

\egapar{11.1.6}
由（11.1.1.2）与定义（11.1.1.3），在谱序列 $E$ 中，若 $E_r^{pq}=0$，则 $k\geq r$ 时 $E_k^{pq}=0$，并且 $E_\infty^{pq}=0$。若存在整数 $r\geq2$，并对每个整数 $n\in\mathbf Z$ 存在整数 $q(n)$，使每个 $q\ne q(n)$ 都有 $E_r^{n-q,q}=0$，则称谱序列\textbf{退化}。由上述论述首先可知，当 $k\geq r$（包括 $k=\infty$）且 $q\ne q(n)$ 时，也有 $E_k^{n-q,q}=0$。此外，由 $E_{r+1}^{pq}$ 的定义，有
\[
E_{r+1}^{n-q(n),q(n)}=E_r^{n-q(n),q(n)}.
\]
若 $E$ 弱收敛，便还有 $E_\infty^{n-q(n),q(n)}=E_r^{n-q(n),q(n)}$；换言之，对每个 $n\in\mathbf Z$，$p\ne q(n)$ 时 $\gr_p(E^n)=0$，且 $\gr_{q(n)}(E^n)=E_r^{n-q(n),q(n)}$。若 $E^n$ 的滤过又离散且穷尽，则 $E$ 正则，且在同构意义下 $E^n=E_r^{n-q(n),q(n)}$。\transnote{本条的退化定义、稳定性断言及 $\gr_{q(n)}$ 下标均按底本保留，没有以通常的“某一页后所有微分为零”的定义替换。按（11.1.1，f）的约定，$E_\infty^{n-q(n),q(n)}$ 对应的滤过次数应为 $n-q(n)$；原文下标应谨慎核对。本译稿在此只标明文本疑点，不把额外条件补写成原著内容。}

\egapar{11.1.7}
假设范畴 $C$ 中滤过归纳极限存在且正合，设 $(E_\lambda,u_{\mu\lambda})$ 为 $C$ 中谱序列的归纳系，其指标集滤过。那么，此归纳系的归纳极限在 $C$ 中对象的谱序列所成加法范畴中存在。只须分别把
\[
E_r^{pq},d_r^{pq},\alpha_r^{pq},B_\infty(E_2^{pq}),Z_\infty(E_2^{pq}),E^n,F^p(E^n),\beta^{pq}
\]
定义为相应各
\[
E_{r,\lambda}^{pq},d_{r,\lambda}^{pq},\alpha_{r,\lambda}^{pq},
B_\infty(E_{2,\lambda}^{pq}),Z_\infty(E_{2,\lambda}^{pq}),E_\lambda^n,F^p(E_\lambda^n),\beta_\lambda^{pq}
\]
的归纳极限。条件（11.1.1）的验证由 $\ilim$ 在 $C$ 中的正合性得出。

\egastatement{11.1.8}{注}
假设 $C$ 为诺特环 $A$ 上的 $A$-模范畴（分别地，为环 $A$ 上的模范畴）。定义（11.1.1）表明：若对某个给定的 $r$，各 $E_r^{pq}$ 为有限型 $A$-模（分别地，有限长度模），则对每个 $s\geq r$，各 $E_s^{pq}$ 以及 $E_\infty^{pq}$ 都有同样性质。若进一步对每个 $n$，终项 $E^n$ 的滤过离散且余离散，则每个 $E^n$ 也为有限型 $A$-模（分别地，有限长度模）。

\egapar{11.1.9}
我们将要考虑保证谱序列 $E$ 关于 $p+q=n$“一致”双正则的条件，为此使用下列引理。

\egastatement{11.1.10}{引理}
设 $(E_r^{pq})$ 是由（11.1.1）的数据 a）、b）、c）联系起来的一族 $C$ 的对象。对一个固定整数 $n$，下列性质等价：
\begin{enumerate}[label=\alph*)]
\item 存在整数 $r(n)$，使当 $r\geq r(n)$ 且 $p+q=n$ 或 $p+q=n-1$ 时，各态射 $d_r^{pq}$ 均为零。
\item 存在整数 $r(n)$，使当 $p+q=n$ 或 $p+q=n+1$ 且 $s\geq r\geq r(n)$ 时，有 $B_r(E_2^{pq})=B_s(E_2^{pq})$。
\item 存在整数 $r(n)$，使当 $p+q=n$ 或 $p+q=n-1$ 且 $s\geq r\geq r(n)$ 时，有 $Z_r(E_2^{pq})=Z_s(E_2^{pq})$。
\item 存在整数 $r(n)$，使当 $p+q=n$ 且 $s\geq r\geq r(n)$ 时，有 $B_r(E_2^{pq})=B_s(E_2^{pq})$ 和 $Z_r(E_2^{pq})=Z_s(E_2^{pq})$。
\end{enumerate}
\begin{proof}
由（11.1.1）的条件 a）、b）、c），$Z_{r+1}(E_2^{pq})=Z_r(E_2^{pq})$ 等价于 $d_r^{pq}=0$，而
\[
B_r(E_2^{p+r,q-r+1})=B_{r+1}(E_2^{p+r,q-r+1})
\]
也等价于 $d_r^{pq}=0$。引理由此立即得出。
\end{proof}

\subsection{滤过复形的谱序列}
\src{27--29}
\egapar{11.2.1}
给定阿贝尔范畴 $C$，约定用 $K^\bullet$ 这样的记号表示微分算子次数为 $+1$ 的 $C$ 中对象的复形 $(K^i)_{i\in\mathbf Z}$，用 $K_\bullet$ 这样的记号表示微分算子次数为 $-1$ 的复形 $(K_i)_{i\in\mathbf Z}$。对微分 $d$ 次数为 $+1$ 的复形 $K^\bullet=(K^i)$，令 $K'_i=K^{-i}$，并以 $d:K^{-i}\to K^{-i+1}$ 为微分 $K'_i\to K'_{i-1}$，便得到复形 $K'_\bullet$；反之亦然。这样可根据情形选用其中一种复形，并把关于一种类型的结果转述为另一种类型的结果。

同样，用 $K^{\bullet\bullet}=(K^{ij})$（分别地，$K_{\bullet\bullet}=(K_{ij})$）表示两个微分算子次数均为 $+1$（分别地，$-1$）的双复形；改变指标的符号就可在两种类型之间转换。对任意多重复形，也采用相应记号并有类似说明。$K^\bullet$ 或 $K_\bullet$ 还用来表示未必为复形的、以 $\mathbf Z$ 分次的 $C$ 中对象（也可令其微分为零而把它看作复形）。例如，以
\[
\mathrm H^\bullet(K^\bullet)=(\mathrm H^i(K^\bullet))_{i\in\mathbf Z}
\]
表示微分次数为 $+1$ 的复形的上同调；以 $\mathrm H_\bullet(K_\bullet)=(\mathrm H_i(K_\bullet))_{i\in\mathbf Z}$ 表示微分次数为 $-1$ 的复形的同调。按上述方法从 $K^\bullet$ 变成 $K'_\bullet$ 时，有 $\mathrm H_i(K'_\bullet)=\mathrm H^{-i}(K^\bullet)$。

还要指出，对复形 $K^\bullet$，通常写
\[
Z^i(K^\bullet)=\Ker(K^i\to K^{i+1}),\qquad
B^i(K^\bullet)=\Im(K^{i-1}\to K^i),
\]
分别称为\textbf{上闭链对象}与\textbf{上边界对象}；对 $K_\bullet$，写
\[
Z_i(K_\bullet)=\Ker(K_i\to K_{i-1}),\qquad
B_i(K_\bullet)=\Im(K_{i+1}\to K_i),
\]
分别称为\textbf{闭链对象}与\textbf{边界对象}。于是
\[
\mathrm H^i(K^\bullet)=Z^i(K^\bullet)/B^i(K^\bullet),\qquad
\mathrm H_i(K_\bullet)=Z_i(K_\bullet)/B_i(K_\bullet).
\]
若 $K^\bullet=(K^i)$（分别地，$K_\bullet=(K_i)$）是 $C$ 中复形，$T:C\to C'$ 为到阿贝尔范畴 $C'$ 的函子，则以 $T(K^\bullet)$（分别地，$T(K_\bullet)$）表示 $C'$ 中的复形 $(T(K^i))$（分别地，$(T(K_i))$）。

这里不再重复 $\partial$-函子的定义（T，2.1），只指出：当态射 $\partial$ 将次数降低一时，也称“$\partial$-函子”，而不另称“$\partial^*$-函子”；若有疑问，每次均由上下文明确这一点。
最后，若存在 $i_0$，使 $i<i_0$（分别地，$i>i_0$）时 $A^i=0$，便称分次对象 $(A^i)_{i\in\mathbf Z}$ \textbf{有下界}（分别地，\textbf{有上界}）。

\egapar{11.2.2}
设 $K^\bullet$ 为 $C$ 中微分 $d$ 的次数为 $+1$ 的复形，带有由分次子对象组成的滤过 $F(K^\bullet)=(F^p(K^\bullet))_{p\in\mathbf Z}$，即
\[
F^p(K^\bullet)=(K^i\cap F^p(K^\bullet))_{i\in\mathbf Z};
\]
并假设对每个 $p\in\mathbf Z$，有 $d(F^p(K^\bullet))\subset F^p(K^\bullet)$。现在简要回顾如何由 $K^\bullet$ 函子地定义一个谱序列 $E(K^\bullet)$（M，XV，4；G，I，4.3）。对 $r\geq2$，典范态射
\[
F^p(K^\bullet)/F^{p+r}(K^\bullet)\longrightarrow F^p(K^\bullet)/F^{p+1}(K^\bullet)
\]
给出上同调态射
\[
\mathrm H^{p+q}(F^p(K^\bullet)/F^{p+r}(K^\bullet))\longrightarrow
\mathrm H^{p+q}(F^p(K^\bullet)/F^{p+1}(K^\bullet)).
\]
以 $Z_r^{pq}(K^\bullet)$ 表示其像。同样，从正合列
\[
\begin{aligned}
0\longrightarrow F^p(K^\bullet)/F^{p+1}(K^\bullet)
&\longrightarrow F^{p-r+1}(K^\bullet)/F^{p+1}(K^\bullet)\\
&\longrightarrow F^{p-r+1}(K^\bullet)/F^p(K^\bullet)\longrightarrow0
\end{aligned}
\]
由上同调正合列得到态射
\[
\mathrm H^{p+q-1}(F^{p-r+1}(K^\bullet)/F^p(K^\bullet))\longrightarrow
\mathrm H^{p+q}(F^p(K^\bullet)/F^{p+1}(K^\bullet)),
\]
以 $B_r^{pq}(K^\bullet)$ 表示其像。可以证明 $B_r^{pq}(K^\bullet)\subset Z_r^{pq}(K^\bullet)$；令
\[
E_r^{pq}(K^\bullet)=Z_r^{pq}(K^\bullet)/B_r^{pq}(K^\bullet).
\]
这里不具体写出 $d_r^{pq}$ 和 $\alpha_r^{pq}$ 的定义。

注意，对固定 $p,q$，各 $Z_r^{pq}(K^\bullet)$ 与 $B_r^{pq}(K^\bullet)$ 都是同一个对象
\[
\mathrm H^{p+q}(F^p(K^\bullet)/F^{p+1}(K^\bullet))
\]
的子对象；把该对象记为 $Z_1^{pq}(K^\bullet)$，并令 $B_1^{pq}(K^\bullet)=0$。这样，前述定义也适用于 $r=1$；再令 $E_1^{pq}(K^\bullet)=Z_1^{pq}(K^\bullet)$。也定义 $d_1^{pq}$、$\alpha_1^{pq}$，使（11.1.1）的条件对 $r=1$ 成立。另一方面，把态射
\[
\mathrm H^{p+q}(F^p(K^\bullet))\longrightarrow
\mathrm H^{p+q}(F^p(K^\bullet)/F^{p+1}(K^\bullet))=E_1^{pq}(K^\bullet)
\]
的像记为 $Z_\infty^{pq}(K^\bullet)$；把如前由上同调正合列所得态射
\[
\mathrm H^{p+q-1}(K^\bullet/F^p(K^\bullet))\longrightarrow
\mathrm H^{p+q}(F^p(K^\bullet)/F^{p+1}(K^\bullet))=E_1^{pq}(K^\bullet)
\]
的像记为 $B_\infty^{pq}(K^\bullet)$。以这两个子对象在 $E_2^{pq}(K^\bullet)$ 中的典范像，分别作为 $Z_\infty(E_2^{pq}(K^\bullet))$ 与 $B_\infty(E_2^{pq}(K^\bullet))$。

最后，以 $F^p(\mathrm H^n(K^\bullet))$ 表示由典范嵌入 $F^p(K^\bullet)\to K^\bullet$ 所得态射
\[
\mathrm H^n(F^p(K^\bullet))\longrightarrow\mathrm H^n(K^\bullet)
\]
的像。由上同调正合列，这也等于态射 $\mathrm H^n(K^\bullet)\to\mathrm H^n(K^\bullet/F^p(K^\bullet))$ 的核。这样就在 $E^n(K^\bullet)=\mathrm H^n(K^\bullet)$ 上定义了滤过；这里也不写出各同构 $\beta^{pq}$ 的定义。

\egapar{11.2.3}
$E(K^\bullet)$ 的函子性应作如下理解：给定 $C$ 中两个滤过复形 $K^\bullet,K'^\bullet$ 和一个与滤过相容的复形态射 $u:K^\bullet\to K'^\bullet$，便显然得到各 $u_r^{pq}$（$r\geq1$）和 $u^n$；可证明这些态射在（11.1.2）的意义下与 $d_r^{pq}$、$\alpha_r^{pq}$ 及 $\beta^{pq}$ 相容，因此给出谱序列态射 $E(u):E(K^\bullet)\to E(K'^\bullet)$。此外，可证明：若 $u,v:K^\bullet\to K'^\bullet$ 为上述类型的态射，并以阶数 $\leq k$ 同伦，则 $r>k$ 时 $u_r^{pq}=v_r^{pq}$，且对每个 $n$，$u^n=v^n$（M，XV，3.1）。

\egapar{11.2.4}
假设 $C$ 中滤过归纳极限正合。若 $K^\bullet$ 的滤过 $(F^p(K^\bullet))$ 穷尽，则对每个 $n$，滤过 $(F^p(\mathrm H^n(K^\bullet)))$ 也穷尽；因为由假设，$K^\bullet=\ilim_{p\to-\infty}F^p(K^\bullet)$，而对 $C$ 的假设使上同调与归纳极限交换。同理，有
\[
B_\infty(E_2^{pq}(K^\bullet))=\sup_k B_k(E_2^{pq}(K^\bullet)).
\]
若对每个 $n$ 存在整数 $u(n)$，使 $p>u(n)$ 时 $\mathrm H^n(F^p(K^\bullet))=0$，则称 $K^\bullet$ 的滤过\textbf{正则}。特别地，离散滤过具有此性质。可证明（M，XV，4）：当 $K^\bullet$ 的滤过正则且穷尽，且 $C$ 中滤过归纳极限正合时，谱序列 $E(K^\bullet)$ 正则。

\subsection{双复形的谱序列}
\src{29--32}
\egapar{11.3.1}
关于双复形的约定，这里采用（T，2.4）而不是（M）的约定；因此，对双复形 $K^{\bullet\bullet}=(K^{ij})$，假设两个次数为 $+1$ 的微分 $d',d''$ 可交换。假设以下条件之一成立：1° $C$ 中存在无限直和；2° 对每个 $n\in\mathbf Z$，满足 $p+q=n$ 且 $K^{pq}\ne0$ 的偶对 $(p,q)$ 只有有限个。那么双复形给出一个（单）复形 $(K^n)_{n\in\mathbf Z}$，其中 $K^n=\bigoplus_{i+j=n}K^{ij}$，次数为 $+1$ 的微分由
\[
dx=d'x+(-1)^i d''x\qquad(x\in K^{ij})
\]
定义。今后提到双复形 $K^{\bullet\bullet}$ 定义的（单）复形时，总是默认满足上述条件之一。对多重复形采用类似约定。

以 $K^{i,\bullet}$（分别地，$K^{\bullet,j}$）表示单复形 $(K^{ij})_{j\in\mathbf Z}$（分别地，$(K^{ij})_{i\in\mathbf Z}$），并以
\[
Z_{\mathrm{II}}^p(K^{i,\bullet}),\ B_{\mathrm{II}}^p(K^{i,\bullet}),\ \mathrm H_{\mathrm{II}}^p(K^{i,\bullet})
\]
（分别地，$Z_{\mathrm I}^p(K^{\bullet,j}), B_{\mathrm I}^p(K^{\bullet,j}),\mathrm H_{\mathrm I}^p(K^{\bullet,j})$）表示其第 $p$ 个上闭链、上边界和上同调对象。微分 $d':K^{i,\bullet}\to K^{i+1,\bullet}$ 为复形态射，故在上闭链、上边界和上同调上给出算子
\[
\begin{aligned}
d'&:Z_{\mathrm{II}}^p(K^{i,\bullet})\to Z_{\mathrm{II}}^p(K^{i+1,\bullet}),\\
d'&:B_{\mathrm{II}}^p(K^{i,\bullet})\to B_{\mathrm{II}}^p(K^{i+1,\bullet}),\\
d'&:\mathrm H_{\mathrm{II}}^p(K^{i,\bullet})\to\mathrm H_{\mathrm{II}}^p(K^{i+1,\bullet}).
\end{aligned}
\]
显然，以这些算子为微分，各相应的 $i$-指标对象族都是复形。将复形 $(\mathrm H_{\mathrm{II}}^p(K^{i,\bullet}))_{i\in\mathbf Z}$ 记为 $\mathrm H_{\mathrm{II}}^p(K^{\bullet\bullet})$；其第 $q$ 个上闭链、上边界和上同调对象分别记为 $Z_{\mathrm I}^q(\mathrm H_{\mathrm{II}}^p(K^{\bullet\bullet}))$、$B_{\mathrm I}^q(\mathrm H_{\mathrm{II}}^p(K^{\bullet\bullet}))$、$\mathrm H_{\mathrm I}^q(\mathrm H_{\mathrm{II}}^p(K^{\bullet\bullet}))$。同样定义复形 $\mathrm H_{\mathrm I}^p(K^{\bullet\bullet})$ 及其上同调对象 $\mathrm H_{\mathrm{II}}^q(\mathrm H_{\mathrm I}^p(K^{\bullet\bullet}))$。回顾 $\mathrm H^n(K^{\bullet\bullet})$ 表示由该双复形定义的单复形的第 $n$ 个上同调对象。

\egapar{11.3.2}
在由双复形定义的复形上，可以考虑两个典范滤过 $(F_{\mathrm I}^p)$、$(F_{\mathrm{II}}^p)$，定义为
\begin{equation}\tag{11.3.2.1}
F_{\mathrm I}^p(K^{\bullet\bullet})=
\left(\bigoplus_{i+j=n,\ i\geq p}K^{ij}\right)_n,
\qquad
F_{\mathrm{II}}^p(K^{\bullet\bullet})=
\left(\bigoplus_{i+j=n,\ j\geq p}K^{ij}\right)_n.
\end{equation}
按定义，它们确为相应单复形的分次子对象，从而使之成为滤过复形；显然这些滤过穷尽且分离。

每个滤过对应一个谱序列（11.2.2）；分别以 $'E(K^{\bullet\bullet})$、$''E(K^{\bullet\bullet})$ 表示对应于 $(F_{\mathrm I}^p)$ 和 $(F_{\mathrm{II}}^p)$ 的谱序列，称为双复形 $K^{\bullet\bullet}$ 的\textbf{谱序列}；二者均以上同调 $(\mathrm H^n(K^{\bullet\bullet}))$ 为终项。还可证明（M，XV，6）：
\begin{equation}\tag{11.3.2.2}
\begin{aligned}
'E_2^{pq}(K^{\bullet\bullet})&=\mathrm H_{\mathrm I}^p(\mathrm H_{\mathrm{II}}^q(K^{\bullet\bullet})),\\
''E_2^{pq}(K^{\bullet\bullet})&=\mathrm H_{\mathrm{II}}^p(\mathrm H_{\mathrm I}^q(K^{\bullet\bullet})).
\end{aligned}
\end{equation}
每个双复形态射 $u:K^{\bullet\bullet}\to K'^{\bullet\bullet}$ 自动与同类型的滤过相容，因此对两个谱序列分别给出态射。此外，两个同伦态射在相应滤过单复形之间给出阶数 $\leq1$ 的同伦，因此在每个谱序列上诱导同一态射（M，XV，6.1）。

\egastatement{11.3.3}{命题}
设 $K^{\bullet\bullet}=(K^{ij})$ 为阿贝尔范畴 $C$ 中的双复形。
\begin{enumerate}[label=(\roman*)]
\item 若存在 $i_0,j_0$，使 $i<i_0$ 或 $j<j_0$（分别地，$i>i_0$ 或 $j>j_0$）时 $K^{ij}=0$，则两个谱序列 $'E$ 与 $''E$ 均双正则。
\item 若存在 $i_0,i_1$，使 $i<i_0$ 或 $i>i_1$ 时 $K^{ij}=0$（分别地，存在 $j_0,j_1$，使 $j<j_0$ 或 $j>j_1$ 时 $K^{ij}=0$），则两个谱序列均双正则。
\end{enumerate}
再假设 $C$ 中滤过归纳极限存在且正合。那么：
\begin{enumerate}[label=(\roman*),start=3]
\item 若存在 $i_0$，使 $i>i_0$ 时 $K^{ij}=0$（分别地，存在 $j_0$，使 $j<j_0$ 时 $K^{ij}=0$），则 $'E$ 正则。
\item 若存在 $i_0$，使 $i<i_0$ 时 $K^{ij}=0$（分别地，存在 $j_0$，使 $j>j_0$ 时 $K^{ij}=0$），则 $''E$ 正则。
\end{enumerate}
\begin{proof}
命题由定义（11.1.3）、（11.2.4）以及下述关于滤过 $F_{\mathrm I}$ 的观察立即得到；交换 $K^{\bullet\bullet}$ 的两指标，就得到关于 $F_{\mathrm{II}}$ 的类似观察。

1° 若 $i>i_0$ 时 $K^{ij}=0$，则滤过 $F_{\mathrm I}$ 离散。

2° 若 $i<i_0$ 时 $K^{ij}=0$，则 $F_{\mathrm I}$ 余离散。于是对每个 $n$，相应的上同调滤过也余离散；而由相应的 $B_r^{pq}$ 的定义（11.2.2），每个固定 $(p,q)$ 上的列 $(B_r^{pq})_{r\geq2}$ 稳定。

3° 若 $j<j_0$ 时 $K^{ij}=0$，则
\[
F_{\mathrm I}^{p+r}(K^{\bullet\bullet})\cap\left(\bigoplus_{i+j=n}K^{ij}\right)=0
\]
只要 $p+r+j_0>n$。因此当 $r\geq q-j_0+1$ 时，$Z_r^{pq}=Z_\infty(E_2^{pq})$；另一方面，$p>n-j_0+1$ 时 $\mathrm H^n(F_{\mathrm I}^p(K^{\bullet\bullet}))=0$。

4° 若 $j>j_0$ 时 $K^{ij}=0$，则
\[
F_{\mathrm I}^{p-r+1}(K^{\bullet\bullet})\cap\left(\bigoplus_{i+j=n}K^{ij}\right)
=\bigoplus_{i+j=n}K^{ij}
\]
只要 $p-r+1+j_0<n$。于是，底本写作当 $r<j_0-q+1$ 时 $B_r^{pq}=B_\infty(E_2^{pq})$；另一方面，$p+j_0<n-1$ 时，$\mathrm H^n(F_{\mathrm I}^p(K^{\bullet\bullet}))=\mathrm H^n(K^{\bullet\bullet})$。\transnote{这里关于 $r$ 的不等号按底本第31页保留。由上一句 $p-r+1+j_0<n$ 解得的是 $r>p+j_0+1-n$；令 $n=p+q$，应是 $r>j_0-q+1$。}
\end{proof}

\egapar{11.3.4}
假设双复形 $K^{\bullet\bullet}=(K^{ij})$ 满足 $i<0$ 或 $j<0$ 时 $K^{ij}=0$。已知此时对每个 $p\in\mathbf Z$ 可以定义典范的\textbf{边缘同态}
\begin{equation}\tag{11.3.4.1}
'E_2^{p0}(K^{\bullet\bullet})\longrightarrow\mathrm H^p(K^{\bullet\bullet})
\end{equation}
（M，XV，6）。简要回顾其理由：一方面，在谱序列 $'E$ 中，当 $2\leq r\leq+\infty$ 时有 $Z_r^{p0}=Z_2^{p0}=Z^p(Z_{\mathrm{II}}^0(K^{\bullet\bullet}))$；另一方面，$\mathrm H^p(F_{
\mathrm I}^{p+1}(K^{\bullet\bullet}))=0$，因此同构
\[
\beta^{p0}:'E_\infty^{p0}\iso\mathrm H^p(F_{\mathrm I}^p)/\mathrm H^p(F_{\mathrm I}^{p+1})
\]
给出同态 $'E_\infty^{p0}\to\mathrm H^p(F_{\mathrm I}^p(K^{\bullet\bullet}))\to\mathrm H^p(K^{\bullet\bullet})$。所有 $Z_r^{p0}$ 相等，又给出当 $r\leq s$ 时的典范同态 $'E_r^{p0}\to'E_s^{p0}$，特别给出 $'E_2^{p0}\to'E_\infty^{p0}$；复合即得边缘同态（11.3.4.1）。此外，立即可验证：若 $z\in Z_{\mathrm{II}}^0(K^{p,\bullet})\subset K^{p0}$ 且 $d'z=0$，则此边缘同态先把 $z$ 模 $B_2^{p0}$ 的类送到 $'E_\infty^{p0}$ 中模 $B_\infty^{p0}$ 的类，再送到 $\mathrm H^p(K^{\bullet\bullet})$ 中 $z$ 的上同调类。

所以最终可见，（11.3.4.1）来自典范嵌入 $Z_{\mathrm{II}}^0(K^{\bullet\bullet})\to K^{\bullet\bullet}$（后者视为单复形）在上同调上诱导的映射。同样，边缘同态
\begin{equation}\tag{11.3.4.2}
''E_2^{p0}(K^{\bullet\bullet})\longrightarrow\mathrm H^p(K^{\bullet\bullet})
\end{equation}
来自典范嵌入 $Z_{\mathrm I}^0(K^{\bullet\bullet})\to K^{\bullet\bullet}$。

\egapar{11.3.5}
现设 $K_{\bullet\bullet}=(K_{ij})$ 为 $C$ 中两个微分次数均为 $-1$ 的双复形。以 $K_{i,\bullet}$（分别地，$K_{\bullet,j}$）表示单复形 $(K_{ij})_{j\in\mathbf Z}$（分别地，$(K_{ij})_{i\in\mathbf Z}$）；以 $\mathrm H_p^{\mathrm{II}}(K_{i,\bullet})$（分别地，$\mathrm H_p^{\mathrm I}(K_{\bullet,j})$）表示其第 $p$ 个同调对象，以 $\mathrm H_p^{\mathrm{II}}(K_{\bullet\bullet})$（分别地，$\mathrm H_p^{\mathrm I}(K_{\bullet\bullet})$）表示复形 $(\mathrm H_p^{\mathrm{II}}(K_{i,\bullet}))_i$（分别地，$(\mathrm H_p^{\mathrm I}(K_{\bullet,j}))_j$），并以 $\mathrm H_q^{\mathrm I}(\mathrm H_p^{\mathrm{II}}(K_{\bullet\bullet}))$（分别地，$\mathrm H_q^{\mathrm{II}}(\mathrm H_p^{\mathrm I}(K_{\bullet\bullet}))$）表示其第 $q$ 个同调对象。对闭链和边界对象采用类似记号；最后，$\mathrm H_n(K_{\bullet\bullet})$ 表示由此双复形定义的微分次数为 $-1$ 的单复形的第 $n$ 个同调对象（若存在）。

令 $K'^{ij}=K_{-i,-j}$，得到与之相伴的微分次数均为 $+1$ 的双复形 $K'^{\bullet\bullet}$。按定义，$K_{\bullet\bullet}$ 的谱序列就是 $K'^{\bullet\bullet}$ 的谱序列，记为 $'E(K_{\bullet\bullet})$ 和 $''E(K_{\bullet\bullet})$，但改变记号，令
\[
'E^r_{pq}(K_{\bullet\bullet})='E_r^{-p,-q}(K'^{\bullet\bullet}),\qquad
''E^r_{pq}(K_{\bullet\bullet})=''E_r^{-p,-q}(K'^{\bullet\bullet})
\]
（$2\leq r\leq\infty$）。于是
\[
'E^2_{pq}=\mathrm H_p^{\mathrm I}(\mathrm H_q^{\mathrm{II}}(K_{\bullet\bullet})),\qquad
''E^2_{pq}=\mathrm H_p^{\mathrm{II}}(\mathrm H_q^{\mathrm I}(K_{\bullet\bullet})).
\]
为避免符号错误，在讨论这些谱序列与终项的关系时，一般宜回到复形 $K'^{\bullet\bullet}$。不过仍记下与（11.3.3）对应的准则。

\egapar{11.3.6}
在下列情形中，$'E(K_{\bullet\bullet})$ 与 $''E(K_{\bullet\bullet})$ 都双正则：a）存在 $i_0,j_0$，使 $i>i_0$ 或 $j>j_0$（分别地，$i<i_0$ 或 $j<j_0$）时 $K_{ij}=0$；b）存在 $i_0,i_1$，使 $i<i_0$ 与 $i>i_1$ 的各情形中均有 $K_{ij}=0$；c）存在 $j_0,j_1$，使 $j<j_0$ 与 $j>j_1$ 的各情形中均有 $K_{ij}=0$。

若存在 $i_0$，使 $i<i_0$ 时 $K_{ij}=0$，或存在 $j_0$，使 $j>j_0$ 时 $K_{ij}=0$，则 $'E(K_{\bullet\bullet})$ 正则。

若存在 $i_0$，使 $i>i_0$ 时 $K_{ij}=0$，或存在 $j_0$，使 $j<j_0$ 时 $K_{ij}=0$，则 $''E(K_{\bullet\bullet})$ 正则。\par

% ---- END c0_s11a.tex ----
% ---- BEGIN c0_s11b.tex ----
\subsection{函子相对于复形的超上同调}
\src{32--35}
\egapar{11.4.1}
设 $C$ 为阿贝尔范畴。回顾 $C$ 中对象 $A$ 的\textbf{右解消}（或\textbf{上同调解消}）：它是微分次数为 $+1$ 的 $C$ 中复形
\[
0\longrightarrow L^0\longrightarrow L^1\longrightarrow L^2\longrightarrow\cdots,
\]
配有一个称为解消的\textbf{增广}的态射 $\varepsilon:A\to L^0$（可把它看作复形态射
\[
\begin{tikzcd}[column sep=small]
0\arrow[r]&A\arrow[r]\arrow[d,"\varepsilon"]&0\arrow[r]\arrow[d]&0\arrow[r]\arrow[d]&\cdots\\
0\arrow[r]&L^0\arrow[r]&L^1\arrow[r]&L^2\arrow[r]&\cdots
\end{tikzcd}
\]
），使序列 $0\to A\to L^0\to L^1\to\cdots$ 正合。同样，$A$ 的\textbf{左解消}（或\textbf{同调解消}）是微分次数为 $-1$ 的复形 $0\leftarrow L_0\leftarrow L_1\leftarrow\cdots$，配有增广 $\varepsilon:L_0\to A$，使序列 $0\leftarrow A\leftarrow L_0\leftarrow L_1\leftarrow\cdots$ 正合。

若对象 $A$ 的右解消 $(L^i)_{i\geq0}$ 满足 $i\geq n+1$ 时 $L^i=0$，就称其\textbf{长度 $\leq n$}。同样定义长度 $\leq n$ 的左解消。若对某个整数 $n$，解消长度 $\leq n$，就称其\textbf{有限}。

若组成解消的除 $A$ 外的对象全为投射对象（分别地，内射对象），就称该解消为\textbf{投射解消}（分别地，\textbf{内射解消}）。当 $C$ 是某环上例如左模的范畴时，若解消中除 $A$ 外的各模全为平坦模（分别地，自由模），则相应称该解消为\textbf{平坦解消}（分别地，\textbf{自由解消}）。

\egapar{11.4.2}
设 $K^\bullet=(K^i)_{i\in\mathbf Z}$ 为 $C$ 中微分次数为 $+1$ 的复形。

$K^\bullet$ 的\textbf{右 Cartan--Eilenberg 解消}是由双复形 $L^{\bullet\bullet}=(L^{ij})$ 和单复形态射 $\varepsilon:K^\bullet\to L^{\bullet,0}$ 组成的偶对；其中双复形的两个微分次数均为 $+1$，$j<0$ 时 $L^{ij}=0$，并满足下列条件：
\begin{enumerate}[label=(\roman*)]
\item 对每个指标 $i$，下列序列正合：
\[
\begin{aligned}
0&\longrightarrow K^i\longrightarrow L^{i0}\longrightarrow L^{i1}\longrightarrow\cdots,\\
0&\longrightarrow B^i(K^\bullet)\longrightarrow B_{\mathrm I}^i(L^{\bullet,0})\longrightarrow B_{\mathrm I}^i(L^{\bullet,1})\longrightarrow\cdots,\\
0&\longrightarrow Z^i(K^\bullet)\longrightarrow Z_{\mathrm I}^i(L^{\bullet,0})\longrightarrow Z_{\mathrm I}^i(L^{\bullet,1})\longrightarrow\cdots,\\
0&\longrightarrow\mathrm H^i(K^\bullet)\longrightarrow\mathrm H_{\mathrm I}^i(L^{\bullet,0})\longrightarrow\mathrm H_{\mathrm I}^i(L^{\bullet,1})\longrightarrow\cdots.
\end{aligned}
\]
即 $L^{i,\bullet}$、$B_{\mathrm I}^i(L^{\bullet\bullet})$、$Z_{\mathrm I}^i(L^{\bullet\bullet})$、$\mathrm H_{\mathrm I}^i(L^{\bullet\bullet})$ 分别是 $K^i,B^i(K^\bullet),Z^i(K^\bullet),\mathrm H^i(K^\bullet)$ 的解消。
\item 对每个 $j$，单复形 $L^{\bullet,j}$ 分裂；也就是说，正合列
\begin{equation}\tag{11.4.2.1}
0\longrightarrow B_{\mathrm I}^i(L^{\bullet,j})\longrightarrow Z_{\mathrm I}^i(L^{\bullet,j})\longrightarrow\mathrm H_{\mathrm I}^i(L^{\bullet,j})\longrightarrow0
\end{equation}
和
\begin{equation}\tag{11.4.2.2}
0\longrightarrow Z_{\mathrm I}^i(L^{\bullet,j})\longrightarrow L^{ij}\longrightarrow B_{\mathrm I}^{i+1}(L^{\bullet,j})\longrightarrow0
\end{equation}
均分裂。
\end{enumerate}
可证明（M，XVII，1.2）：若 $C$ 的每个对象都是某个内射对象的子对象，则 $C$ 中每个复形 $K^\bullet$ 都有内射 Cartan--Eilenberg 解消，即由内射对象 $L^{ij}$ 组成的解消（上述条件（ii）这时又推出 $B_{\mathrm I}^i(L^{\bullet,j}),Z_{\mathrm I}^i(L^{\bullet,j}),\mathrm H_{\mathrm I}^i(L^{\bullet,j})$ 也内射）。此外，对 $C$ 中每个复形态射 $f:K^\bullet\to K'^\bullet$、$K^\bullet$ 的任意 Cartan--Eilenberg 解消 $L^{\bullet\bullet}$ 及 $K'^\bullet$ 的任意内射 Cartan--Eilenberg 解消 $L'^{\bullet\bullet}$，存在与 $f$ 和增广相容的双复形态射 $F:L^{\bullet\bullet}\to L'^{\bullet\bullet}$。若 $f,g:K^\bullet\to K'^\bullet$ 同伦，则相应的双复形态射也同伦（同上）。

当 $K^\bullet$ 有下界（分别地，有上界）时，若 $i<i_0$（分别地，$i>i_0$）时 $K^i=0$，则可取 $L^{\bullet\bullet}$ 使相同范围内 $L^{ij}=0$（M，XVII，1.3）。另一方面，若存在整数 $n$，使 $C$ 的每个对象都有长度 $\leq n$ 的内射解消，则可假设 $j>n$ 时 $L^{ij}=0$（M，XVII，1.4）。

\egapar{11.4.3}
现在设 $T$ 为从 $C$ 到阿贝尔范畴 $C'$ 的协变加法函子。给定 $C$ 中复形 $K^\bullet$ 及其内射 Cartan--Eilenberg 解消 $L^{\bullet\bullet}$，假设双复形 $T(L^{\bullet\bullet})$ 定义的单复形存在（参见11.3.1）。则此双复形的两个谱序列 $'E(T(L^{\bullet\bullet}))$ 和 $''E(T(L^{\bullet\bullet}))$，称为 $T$ 相对于复形 $K^\bullet$ 的\textbf{超上同调谱序列}。由（11.4.2）与（11.3.2），它们确实只依赖于 $K^\bullet$，不依赖所选的内射 Cartan--Eilenberg 解消；并且它们函子地依赖于 $K^\bullet$。二者具有共同终项 $\mathrm H^\bullet(T(L^{\bullet\bullet}))$，称为 $T$ 相对于 $K^\bullet$ 的\textbf{超上同调}，记为 $\mathrm R^\bullet T(K^\bullet)$。可以证明，这两个谱序列的 $E_2$ 项分别为
\begin{equation}\tag{11.4.3.1}
'E_2^{pq}=\mathrm H^p(\mathrm R^qT(K^\bullet)),
\end{equation}
\begin{equation}\tag{11.4.3.2}
''E_2^{pq}=\mathrm R^pT(\mathrm H^q(K^\bullet)).
\end{equation}
如常，$\mathrm R^pT$ 表示 $T$ 的第 $p$ 个导出函子（$p\in\mathbf Z$）；$\mathrm R^qT(K^\bullet)$ 表示复形 $(\mathrm R^qT(K^i))_{i\in\mathbf Z}$。除非明确声明例外，今后假设 $C$ 的每个对象都是 $C$ 中某个内射对象的子对象，因此每个复形都有内射 Cartan--Eilenberg 解消。由于 $j<0$ 时 $L^{ij}=0$，由（11.3.3）的准则可知，在下列任一情形中，两个超上同调谱序列都存在且双正则：1° $K^\bullet$ 有下界；2° 存在与对象无关的整数 $n$，使 $C$ 的每个对象都有长度 $\leq n$ 的内射解消。事实上，在第一种情形，可以假设（11.4.2）存在 $i_0$ 使 $i<i_0$ 时 $L^{ij}=0$；在第二种情形，可以假设存在 $j_1$ 使 $j>j_1$ 时 $L^{ij}=0$。在两种情形中，对固定 $n$，满足 $L^{ij}\ne0$ 及 $i+j=n$ 的偶对 $(i,j)$ 显然都只有有限个，这就证明了上述断言。

若假设 $C'$ 中滤过归纳极限存在且正合（这尤其蕴含 $C'$ 中无限直和存在），则双复形 $T(L^{\bullet\bullet})$ 定义的复形存在；准则（11.3.3）表明，第一个谱序列 $'E(T(L^{\bullet\bullet}))$ 总是正则的。

若 $K^\bullet$ 除唯一一项 $K^{i_0}$ 外全为零，则 $\mathrm R^nT(K^\bullet)$ 同构于 $\mathrm R^{n-i_0}T(K^{i_0})$；取在 $i\ne i_0$ 时 $L^{ij}=0$ 的 Cartan--Eilenberg 解消，便从定义立即得到此结论。

若 $K^\bullet,K'^\bullet$ 是 $C$ 中两个复形，$f,g:K^\bullet\to K'^\bullet$ 为两个同伦态射，则由 $f,g$ 得到的态射 $\mathrm R^\bullet T(K^\bullet)\to\mathrm R^\bullet T(K'^\bullet)$ 相同，相应上同调谱序列的态射也相同。\transnote{底本此后直接列出（11.4.5），未另印（11.4.4）；后文引用（11.4.4）时，指向这里关于正则性等性质的说明。本译稿不补造该编号。}

\egastatement{11.4.5}{命题}
假设 $C'$ 中滤过归纳极限存在且正合。若对每个 $n>0$ 和 $i\in\mathbf Z$，有 $\mathrm R^nT(K^i)=0$，则有函子同构
\begin{equation}\tag{11.4.5.1}
\mathrm R^iT(K^\bullet)\simeq\mathrm H^i(T(K^\bullet))\qquad(i\in\mathbf Z).
\end{equation}
\begin{proof}
此时第一个谱序列（11.4.3.1）中仅有可能非零的 $E_2$ 项为 $'E_2^{p0}=\mathrm H^p(T(K^\bullet))$；换言之，该谱序列退化。又因其正则（11.4.4），结论由（11.1.6）得出。
\end{proof}

\egapar{11.4.6}
现在例如考虑从 $C\times C'$ 到 $C''$ 的协变双函子 $(M,N)\mapsto T(M,N)$，其中 $C,C',C''$ 均为阿贝尔范畴。为简便起见，假设 $T$ 关于每个变元都加法；再假设 $C,C'$ 的每个对象都是某个内射对象的子对象，而且 $C''$ 中滤过归纳极限存在且正合。给定分别属于 $C,C'$、微分次数为 $+1$ 的复形 $K^\bullet,K'^\bullet$，取其内射 Cartan--Eilenberg 解消 $L^{\bullet\bullet},L'^{\bullet\bullet}$。则 $T(L^{\bullet\bullet},L'^{\bullet\bullet})$ 为 $C''$ 中的四重复形；令 $T(L^{ij},L'^{hk})$ 的两个次数为 $i+h$ 与 $j+k$，将它视为双复形。按定义，$T$ 相对于 $K^\bullet,K'^\bullet$ 的超上同调，就是此双复形（也就是相应单复形）的上同调
\[
\mathrm H^\bullet(T(L^{\bullet\bullet},L'^{\bullet\bullet})),
\]
记作 $\mathrm R^\bullet T(K^\bullet,K'^\bullet)$。它为两个谱序列的共同终项，其 $E_2$ 项如下（参见 M，XVII，2）：
\begin{equation}\tag{11.4.6.1}
'E_2^{pq}=\mathrm H^p(\mathrm R^qT(K^\bullet,K'^\bullet)),
\end{equation}
\begin{equation}\tag{11.4.6.2}
''E_2^{pq}=\bigoplus_{q'+q''=q}\mathrm R^pT(\mathrm H^{q'}(K^\bullet),\mathrm H^{q''}(K'^\bullet)).
\end{equation}
这里 $\mathrm R^qT(K^\bullet,K'^\bullet)$ 是双复形 $(\mathrm R^qT(K^i,K'^j))_{i,j\in\mathbf Z}$，（11.4.6.1）右端表示将其视为单复形时的上同调。

此外，第一个谱序列总正则；若存在 $n$，使 $C,C'$ 的每个对象都有长度 $\leq n$ 的内射解消，或若 $K^\bullet,K'^\bullet$ 均有下界，则两个谱序列均双正则。在后一情形，还可省去底本所说 $C$ 与 $C''$ 中归纳极限存在的假设。

若 $K_1^\bullet,K_1'^\bullet$ 分别为 $C,C'$ 中的另外两个复形，$f,g:K^\bullet\to K_1^\bullet$ 同伦，$f',g':K'^\bullet\to K_1'^\bullet$ 同伦，则分别由 $(f,f')$ 与 $(g,g')$ 得到的态射
\[
\mathrm R^\bullet T(K^\bullet,K'^\bullet)\longrightarrow\mathrm R^\bullet T(K_1^\bullet,K_1'^\bullet)
\]
相同；两个超上同调谱序列的相应态射也相同。容易推广到任意协变加法多函子。

\egastatement{11.4.7}{命题}
假设对 $C$ 的每个内射对象 $I$（分别地，$C'$ 的每个内射对象 $I'$），$A'\mapsto T(I,A')$（分别地，$A\mapsto T(A,I')$）为正合函子。那么，在（11.4.6）的记号下，有典范同构
\begin{equation}\tag{11.4.7.1}
\begin{aligned}
\mathrm R^\bullet T(K^\bullet,K'^\bullet)
&\simeq\mathrm H^\bullet(T(L^{\bullet\bullet},K'^\bullet))\\
&\simeq\mathrm H^\bullet(T(K^\bullet,L'^{\bullet\bullet})),
\end{aligned}
\end{equation}
其中后两项分别为三重复形 $T(L^{\bullet\bullet},K'^\bullet)$ 与 $T(K^\bullet,L'^{\bullet\bullet})$ 定义的单复形的上同调。
\begin{proof}
例如定义第一个同构。可以把四重复形 $T(L^{\bullet\bullet},L'^{\bullet\bullet})$ 视为双复形，令 $T(L^{ij},L'^{hk})$ 的次数为 $i+j+h$ 与 $k$。因每个 $h$ 所对应的 $L'^{h,\bullet}$ 是 $K'^h$ 的解消，由对 $T$ 的假设，对此双复形有 $\mathrm H_{\mathrm{II}}^q=0$（$q\ne0$），并且
\[
\mathrm H_{\mathrm{II}}^0(T(L^{\bullet\bullet},L'^{\bullet\bullet}))=T(L^{\bullet\bullet},K'^\bullet).
\]
所以其第一个谱序列退化；又因 $k<0$ 时 $L'^{hk}=0$，此谱序列正则（11.3.3）。结论由（11.1.6）得出。
\end{proof}
对于任意 $n$ 个变元的协变多函子，也有类似结果：计算超上同调时，不必把所有复形都换成 Cartan--Eilenberg 解消，只须替换其中 $n-1$ 个；条件是，当把任意 $n-1$ 个变元固定为内射对象时，关于余下变元的协变函子正合。

\subsection{超上同调中的归纳极限}
\src{35--39}
\egastatement{11.5.1}{引理}
设 $K^\bullet=(K^i)_{i\in\mathbf Z}$ 为 $C$ 中的复形，对每个整数 $r\in\mathbf Z$，令 $K_{(r)}^\bullet$ 为满足 $i<r$ 时 $K_{(r)}^i=0$、$i\geq r$ 时 $K_{(r)}^i=K^i$ 的复形。设 $T:C\to C'$ 为与归纳极限交换的协变加法函子（假设 $C,C'$ 中滤过归纳极限存在且正合）。则 $\mathrm R^\bullet T(K^\bullet)$ 同构于当 $r\to-\infty$ 时的归纳极限 $\ilim_r\mathrm R^\bullet T(K_{(r)}^\bullet)$。
\begin{proof}
构造 $K^\bullet$ 的内射 Cartan--Eilenberg 解消时，可对每个 $i$ 任意选择 $B^i(K^\bullet)$ 的内射解消 $(X_B^{ij})_{j\geq0}$ 及 $\mathrm H^i(K^\bullet)$ 的内射解消 $(X_H^{ij})_{j\geq0}$。构造方法表明，$K^i$ 的内射解消 $(L^{ij})_{j\geq0}$ 及微分 $L^{ij}\to L^{i+1,j}$，只依赖于解消 $(X_B^{i,\bullet})$、$(X_H^{i,\bullet})$ 与 $(X_B^{i+1,\bullet})$（M，XVII，1.2）。显然，$i\geq r+1$ 时有
\[
B^i(K_{(r)}^\bullet)=B^i(K^\bullet),\qquad
\mathrm H^i(K_{(r)}^\bullet)=\mathrm H^i(K^\bullet).
\]
另一方面，每个 $r$ 都有典范嵌入 $\varphi_{r-1,r}^\bullet:K_{(r)}^\bullet\to K_{(r-1)}^\bullet$，其在 $i\ne r-1$ 的次数上为恒等映射。因此，若 $L^{\bullet\bullet}=(L^{ij})$ 是 $K^\bullet$ 的内射 Cartan--Eilenberg 解消，则对每个 $r$，可以定义 $K_{(r)}^\bullet$ 的内射 Cartan--Eilenberg 解消 $L_{(r)}^{\bullet\bullet}$，使 $i<r$ 时 $L_{(r)}^{ij}=0$，$i\geq r+1$ 时 $L_{(r)}^{ij}=L^{ij}$。

又可定义对应于 $\varphi_{r-1,r}^\bullet$ 的双复形态射 $\Phi_{r-1,r}:L_{(r)}^{\bullet\bullet}\to L_{(r-1)}^{\bullet\bullet}$；此态射的构造方法（同上）表明，可使 $i\ne r$ 且 $i\ne r-1$ 时的 $\Phi_{r-1,r}^{ij}$ 为恒等映射。这样得到双复形的归纳系 $(L_{(r)}^{\bullet\bullet})$，当 $r\to-\infty$ 时，其归纳极限显然为 $L^{\bullet\bullet}$。因为直和与归纳极限可交换，$L^{\bullet\bullet}$ 的相伴单复形也是各 $L_{(r)}^{\bullet\bullet}$ 的相伴单复形的归纳极限。由假设，$T$ 与归纳极限交换；上同调也如此（由于 $\ilim$ 正合）。于是，在同构意义下，确有
\[
\mathrm H^\bullet(T(L^{\bullet\bullet}))=\ilim_r\mathrm H^\bullet(T(L_{(r)}^{\bullet\bullet})).
\]
\end{proof}
引理（11.5.1）允许通过归纳极限，把对有下界复形成立的性质推广到任意复形。作为第一个例子，证明下列命题。

\egastatement{11.5.2}{命题}
关于 $C,C',T$ 的假设同（11.5.1）。则 $\mathrm R^\bullet T(K^\bullet)$ 是 $C$ 中复形所成阿贝尔范畴上的上同调函子。
\eproof
先说明可归结到有下界复形。若有复形的正合列 $0\to K'^\bullet\to K^\bullet\to K''^\bullet\to0$，则对每个 $r$ 显然得到正合列 $0\to K'_{(r)}{}^\bullet\to K_{(r)}^\bullet\to K''_{(r)}{}^\bullet\to0$，因此按所假定的有下界情形，有正合列
\[
\cdots\to\mathrm R^nT(K'_{(r)}{}^\bullet)\to\mathrm R^nT(K_{(r)}^\bullet)
\to\mathrm R^nT(K''_{(r)}{}^\bullet)\xrightarrow{\partial}\mathrm R^{n+1}T(K'_{(r)}{}^\bullet)\to\cdots.
\]
这些正合列组成归纳系。引理（11.5.1）与 $\ilim$ 的正合性表明，得到正合列
\[
\cdots\to\mathrm R^nT(K'^\bullet)\to\mathrm R^nT(K^\bullet)
\to\mathrm R^nT(K''^\bullet)\xrightarrow{\partial}\mathrm R^{n+1}T(K'^\bullet)\to\cdots.
\]
处理有下界复形时，可只考虑 $i<0$ 时 $K^i=0$ 的那些复形；它们显然组成阿贝尔范畴 $\mathbf K$。

\egastatement{11.5.2.1}{引理}
在 $\mathbf K$ 中，令 $\mathfrak I$ 为具有以下性质的复形 $Q^\bullet=(Q^i)_{i\geq0}$ 的全体：1° 每个 $Q^i$ 为 $C$ 的内射对象；2° 对每个 $i\geq0$，$Z^i(Q^\bullet)=B^i(Q^\bullet)$，且 $Z^i(Q^\bullet)$ 为 $Q^i$ 的直因子。那么：
\begin{enumerate}[label=(\roman*)]
\item 每个 $Q^\bullet\in\mathfrak I$ 都是 $\mathbf K$ 中的内射对象。
\item $\mathbf K$ 的每个对象都同构于某个属于 $\mathfrak I$ 的复形的子复形。
\end{enumerate}
\begin{proof}
（i）设 $A^\bullet=(A^i)$ 为 $\mathbf K$ 的对象，$A'^\bullet=(A'^i)$ 为其子对象，$Q^\bullet=(Q^i)$ 为 $\mathfrak I$ 的对象，且给定态射 $f=(f^i):A'^\bullet\to Q^\bullet$。须将其延拓为态射 $g=(g^i):A^\bullet\to Q^\bullet$。为简化叙述，使用模范畴的语言（参见 [27]）。

将 $Q^i$ 等同于 $B^i(Q^\bullet)\oplus B^{i+1}(Q^\bullet)$。对 $i$ 归纳，假设各 $g^j$ 已在 $j<i$ 时定义，并对 $j<i-1$ 与微分 $d^j:A^j\to A^{j+1}$、$d'^j:Q^j\to Q^{j+1}$ 相容；又满足：1° $g^{i-1}(Z^{i-1}(A^\bullet))\subset Z^{i-1}(Q^\bullet)$；2° 若对每个 $j$ 令 $C^j=(d^j)^{-1}(A'^{j+1})$，则 $d'^{i-1}\circ g^{i-1}$ 在 $C^{i-1}$ 上与 $f^i\circ d^{i-1}$ 一致。

态射 $f^i:A'^i\to Q^i$ 与两个投影复合，给出 $f'^i:A'^i\to B^i(Q^\bullet)$ 和 $f''^i:A'^i\to B^{i+1}(Q^\bullet)$。因为 $d'^{i-1}\circ g^{i-1}$ 把 $A^{i-1}$ 送入 $B^i(Q^\bullet)$，并在 $Z^{i-1}(A^\bullet)$ 上为零，便给出态射 $h^i:B^i(A^\bullet)\to B^i(Q^\bullet)$；由其在 $C^{i-1}$ 上与 $f^i\circ d^{i-1}$ 一致，$h^i$ 与 $f'^i$ 在 $B^i(A^\bullet)\cap A'^i$ 上一致。$B^i(Q^\bullet)$ 为 $Q^i$ 的直因子，因而内射；故存在态射 $g'^i:A^i\to B^i(Q^\bullet)$，分别在 $B^i(A^\bullet)$ 和 $A'^i$ 上与 $h^i$ 和 $f'^i$ 一致。

另一方面，考虑态射 $f''^{i+1}\circ d^i:C^i\to B^{i+1}(Q^\bullet)$，它在 $Z^i(A^\bullet)$ 上为零。因为 $B^{i+1}(Q^\bullet)$ 内射，存在态射 $g''^i:A^i\to B^{i+1}(Q^\bullet)$，在 $C^i$ 上与 $f''^{i+1}\circ d^i$ 一致，在 $Z^i(A^\bullet)$ 上为零。取 $g^i=g'^i+g''^i$，即可继续归纳。\transnote{底本此处在 $f'$、$f''$ 的上标与撇号上存在源靶不一致的印式，本文照录其段落构造；尤其按本段定义，$f''^{i+1}$ 的靶是 $B^{i+2}(Q^\bullet)$，与这里所写 $B^{i+1}(Q^\bullet)$ 不同。应把这一印式与前述投影分解对照使用。}

（ii）为了把 $A^\bullet=(A^i)$ 嵌入一个属于 $\mathfrak I$ 的复形，对每个 $i\geq1$，取 $C$ 的内射对象 $Q'^i$ 及嵌入 $f'^i:A^i\to Q'^i$。令 $i<0$ 时 $Q^i=0$，令 $Q^0=Q'^1$，并令 $i\geq1$ 时 $Q^i=Q'^i\oplus Q'^{i+1}$，采用显然的微分算子。对每个 $i$ 令
\[
f^i=f'^i+(f'^{i+1}\circ d^i),
\]
其中 $i\leq0$ 时取 $f'^i=0$。立即可见，每个 $f^i$ 单射，且这些 $f^i$ 与微分算子相容。\transnote{本引理的零次条件及 $Q^0$ 的构造按底本第37页保留。例如一般复形的 $d^0$ 未必单射，故末句在第0次不能仅由所写公式推出。这里未用另一种内射复形构造替换原文。}
\end{proof}

\egastatement{11.5.2.2}{推论}
$\mathbf K$ 的每个对象 $K^\bullet$ 都有由 $\mathfrak I$ 的对象组成的右解消。若 $L^{\bullet\bullet}$ 是这样的解消，则对 $K^\bullet$ 的任意由 $\mathbf K$ 的对象组成的解消 $L'^{\bullet\bullet}$，存在与增广相容的双复形态射 $F:L'^{\bullet\bullet}\to L^{\bullet\bullet}$；任意两个这样的态射 $F,F'$ 同伦。
\begin{proof}
这正是（M，V，1.1 a））应用于阿贝尔范畴 $\mathbf K$。
\end{proof}

\egapar{11.5.2.3}
有了这些预备结果，考虑 $K^\bullet$ 的内射 Cartan--Eilenberg 解消 $L'^{\bullet\bullet}$ 及由 $\mathfrak I$ 的对象组成的解消 $L^{\bullet\bullet}$，并证明
\[
\mathrm H^\bullet(T(L'^{\bullet\bullet}))\simeq\mathrm H^\bullet(T(L^{\bullet\bullet})).
\]
由（11.5.2.2），得到双复形态射 $T(L'^{\bullet\bullet})\to T(L^{\bullet\bullet})$，从而得到其第一个谱序列之间的态射 $'E(T(L'^{\bullet\bullet}))\to'E(T(L^{\bullet\bullet}))$。这些谱序列由（11.3.3）正则，所以由（11.1.5），只须证明它在 $E_2$ 项上为同构；也就是证明 $\mathrm H_{\mathrm{II}}^q(T(L^{i,\bullet}))=\mathrm R^qT(K^i)$。由于 $L^{i,\bullet}$ 为 $K^i$ 的右解消，归结为下面的引理。

\egastatement{11.5.2.4}{引理}
若 $L^\bullet=(L^i)_{i\in\mathbf Z}$ 是 $C$ 中对象 $A$ 的右解消，且对每个 $i\in\mathbf Z$、$n>0$ 有 $\mathrm R^nT(L^i)=0$，则
\[
\mathrm H^\bullet(T(L^\bullet))=\mathrm R^\bullet T(A).
\]
这是（T，2.5.2）的特例。

\egapar{11.5.2.5}
现在，（11.5.2）的证明立即完成：因为 $L^{\bullet\bullet}$ 为 $K^\bullet$ 在阿贝尔范畴 $\mathbf K$ 中的内射解消，换言之，$K^\bullet\mapsto\mathrm H^\bullet(T(L^{\bullet\bullet}))$ 就是 $T$ 在范畴 $\mathbf K$ 中的右导出上同调函子（T，2.3）。\finishproof

\egastatement{11.5.3}{命题}
关于 $C,C',T$ 的假设同（11.5.1）。设 $L^{\bullet\bullet}=(L^{ij})$ 为双复形，满足 $j<0$ 时 $L^{ij}=0$，并且对每个 $i$，$L^{i,\bullet}$ 为 $K^i$ 的解消；再假设对每对 $(i,j)$ 及每个 $n>0$，有 $\mathrm R^nT(L^{ij})=0$。那么存在函子同构
\begin{equation}\tag{11.5.3.1}
\mathrm R^\bullet T(K^\bullet)\simeq\mathrm H^\bullet(T(L^{\bullet\bullet})).
\end{equation}
\begin{proof}
令 $L_{(r)}^{\bullet\bullet}$ 为满足 $i<r$ 时 $L_{(r)}^{ij}=0$、$i\geq r$ 时 $L_{(r)}^{ij}=L^{ij}$ 的双复形。显然，当 $r\to-\infty$ 时，$L^{\bullet\bullet}$ 是这些双复形的归纳极限。由对 $T$ 的假设及（11.5.1），只须在 $K^\bullet$ 有下界的情形证明命题；例如可设 $i<0$ 时 $K^i=0$，并有 $L^{ij}=0$。取一个由 $\mathfrak I$ 的对象组成的 $K^\bullet$ 的右解消 $L'^{\bullet\bullet}$（11.5.2.2）。存在与增广相容的双复形态射 $L^{\bullet\bullet}\to L'^{\bullet\bullet}$，从而得到第一个谱序列之间的态射 $'E(T(L^{\bullet\bullet}))\to'E(T(L'^{\bullet\bullet}))$。如同（11.5.2.3），引理（11.5.2.4）表明此态射为同构，故结论成立。
\end{proof}

\egastatement{11.5.4}{注}
上述论证表明，在有下界复形的范畴中，（11.5.2）和（11.5.3）的结论成立而不必假设 $T$ 与滤过归纳极限交换。此外，当仅考虑满足 $i<0$ 时 $K^i=0$ 的复形所成范畴 $\mathbf K$ 时，把 $\mathrm R^\bullet T(K^\bullet)$ 刻画为 $T$ 在 $\mathbf K$ 中的右导出函子系，便说明此上同调函子是泛的（T，2.3）。

还有一种无需对 $T$ 另加假设而（11.5.2）仍成立的情形：存在整数 $m>0$，使 $C$ 中每个对象都有长度 $\leq m$ 的内射解消。事实上，在（11.5.1）的证明中，所出现的各对象的内射解消都可取为长度 $\leq m$。于是立即可见，当 $r$ 足够大时，$T(L_{(r)}^{\bullet\bullet})$ 中总次数为 $n$ 的项与 $T(L^{\bullet\bullet})$ 的相应项相等，且只有有限个；这推出对每个 $n$，当 $r$ 足够大时，$\mathrm H^n(T(L^{\bullet\bullet}))=\mathrm H^n(T(L_{(r)}^{\bullet\bullet}))$。按（11.5.2）的记号，于是当 $r$ 足够大（依赖于 $n$）时，$\mathrm R^nT(K_{(r)}^\bullet)=\mathrm R^nT(K^\bullet)$，对 $K'^\bullet,K''^\bullet$ 亦然，结论随之得到。\transnote{本段底本写“$r$ 足够大”；按照（11.5.1）对 $K_{(r)}$ 的定义和 $r\to-\infty$ 的方向，实际须理解为 $r$ 充分小（或 $-r$ 充分大）。这里保留原文用语并说明方向。}
同理，当 $C$ 满足上述假设且各解消 $L^{i,\bullet}$ 的长度 $\leq m$ 时，（11.5.3）也无需对 $T$ 另加条件。

\egapar{11.5.5}
（11.5.2）的结果可推广到协变多函子。例如考虑（11.4.6）的情形，再假设 $C,C',C''$ 中滤过归纳极限存在且正合，且 $T$ 与归纳极限交换。那么 $\mathrm R^\bullet T(K^\bullet,K'^\bullet)$ 关于每个复形 $K^\bullet,K'^\bullet$ 都为上同调双函子。证明时，像（11.5.2）那样先归结到两个复形都有下界的情形，再为二者取（11.5.2.2）所述类型的内射解消，便归结到（M，V，4.1）的普遍性质。

\egapar{11.5.6}
同样，（11.4.7）和（11.5.3）可作如下推广。在（11.5.5）的假设下，设有两个双复形 $L^{\bullet\bullet}=(L^{ij})$、$L'^{\bullet\bullet}=(L'^{ij})$，当 $j<0$ 时二者的项均为零，对每个 $i$，$L^{i,\bullet}$ 为 $K^i$ 的解消，$L'^{i,\bullet}$ 为 $K'^i$ 的解消；并假设对任意 $(i,j,h,k)$ 及 $n>0$，有 $\mathrm R^nT(L^{ij},L'^{hk})=0$。那么存在关于 $K^\bullet,K'^\bullet$ 的函子同构
\begin{equation}\tag{11.5.6.1}
\mathrm R^\bullet T(K^\bullet,K'^\bullet)\simeq\mathrm H^\bullet(T(L^{\bullet\bullet},L'^{\bullet\bullet})).
\end{equation}
像（11.5.3）那样，归结到两个复形都有下界的情形即可证明。

进一步假设对每对 $(i,j)$ 及 $(h,k)$，函子 $A\mapsto T(A,L'^{hk})$ 与 $A'\mapsto T(L^{ij},A')$ 分别在 $C,C'$ 中正合。那么还有函子同构
\begin{equation}\tag{11.5.6.2}
\begin{aligned}
\mathrm R^\bullet T(K^\bullet,K'^\bullet)
&\simeq\mathrm H^\bullet(T(L^{\bullet\bullet},K'^\bullet))\\
&\simeq\mathrm H^\bullet(T(K^\bullet,L'^{\bullet\bullet})).
\end{aligned}
\end{equation}
证明与（11.4.7）相同。

\egapar{11.5.7}
还要指出，只要限制于有下界复形 $K^\bullet,K'^\bullet$，或假设 $C$（分别地，$C'$）的各对象具有长度一致有界的内射解消，且在（11.5.6）中两个双复形的第二次数有上界，则（11.5.5）与（11.5.6）的结果无需假设 $T$ 与归纳极限交换也成立。\par

% ---- END c0_s11b.tex ----
% ---- BEGIN c0_s11c.tex ----
\subsection{函子相对于复形的超同调}
\src{39--41}
\egapar{11.6.1}
设 $C$ 为阿贝尔范畴，$K_\bullet=(K_i)_{i\in\Z}$ 为 $C$ 中微分算子次数为 $-1$ 的复形。$K_\bullet$ 的一个\textbf{左 Cartan--Eilenberg 解消}由一个微分算子次数均为 $-1$ 的双复形 $L_{\bullet\bullet}=(L_{ij})$ 和一个单复形态射 $\varepsilon:L_{\bullet0}\to K_\bullet$ 组成，其中 $j<0$ 时 $L_{ij}=0$，并且满足将（11.4.2）的条件中的“箭头反向”所得的条件。如果 $C$ 的每个对象都是一个投射对象的商，则 $C$ 的每个复形 $K_\bullet$ 都有一个\textbf{投射 Cartan--Eilenberg 解消}，即由投射对象 $L_{ij}$ 组成的解消；它具有类似于（11.4.2）的函子性质。此外，若 $K_\bullet$ 有下界（分别地，有上界），可以在 $i<i_0$（分别地，$i>i_0$）时 $K_i=0$ 的条件下，令相应的 $L_{ij}=0$。若 $C$ 的每个对象都有长度 $\leq n$ 的投射解消，则可以令 $j>n$ 时 $L_{ij}=0$。

\egapar{11.6.2}
设 $T$ 为从 $C$ 到阿贝尔范畴 $C'$ 的加性协变函子。相对于 $C$ 的复形 $K_\bullet$ 的\textbf{超同调} $\mathrm L_\bullet T(K_\bullet)$ 及\textbf{超同调谱序列}（在它们存在时），仍由（11.4.3）经“箭头反向”来定义；这样得到的两个谱序列的 $E^2$ 项为
\begin{equation}\tag{11.6.2.1}
{}'E^2_{pq}=\mathrm H_p(\mathrm L_qT(K_\bullet)),
\end{equation}
\begin{equation}\tag{11.6.2.2}
{}''E^2_{pq}=\mathrm L_pT(\mathrm H_q(K_\bullet)),
\end{equation}
其中 $\mathrm L_pT$ 在 $p\geq0$ 时表示 $T$ 的第 $p$ 个导出函子，在 $p<0$ 时为零；$\mathrm L_qT(K_\bullet)$ 表示复形 $(\mathrm L_qT(K_i))_{i\in\Z}$。

超同调的性质并不全能由超上同调的性质简单地经“箭头反向”得到（除非对范畴 $C'$ 作（T，1.5）中 AB~5* 型的附加假设），这是由于上述两个谱序列的正则性条件不同；这次必须对它们应用（11.3.4）的判别准则。\transnote{此处原文引作（11.3.4）；关于同调型双复形正则性的对应判别见（11.3.6）。保留原引文编号。}这些判别说明，若 $C'$ 中滤过归纳极限存在且正合，则由双复形 $T(L_{\bullet\bullet})$ 定义的复形存在，而这次是\emph{第二个}谱序列 ${}''E(T(L_{\bullet\bullet}))$ 正则。若假设 $K_\bullet$ 有下界，或存在整数 $n$ 使 $C$ 的每个对象都有长度 $\leq n$ 的投射解消，则两个超同调谱序列都存在（无需对 $C'$ 作假设）且双正则。

\egastatement{11.6.3}{命题}
设 $C,C'$ 为两个阿贝尔范畴，$T:C\to C'$ 为加性协变函子。则：
\begin{enumerate}[label=(\roman*)]
\item 超同调 $\mathrm L_\bullet T(K_\bullet)$ 是 $C$ 的有下界复形所成阿贝尔范畴上的同调函子。
\item 设 $K_\bullet$ 为 $C$ 的有下界复形。若对每个 $n>0$ 及 $i\in\Z$ 都有 $\mathrm L_nT(K_i)=0$，则存在函子同构
\begin{equation}\tag{11.6.3.1}
\mathrm L_iT(K_\bullet)\simeq\mathrm H_i(T(K_\bullet))\qquad(i\in\Z).
\end{equation}
\item 设 $K_\bullet$ 为 $C$ 的有下界复形，$L_{\bullet\bullet}=(L_{ij})$ 为双复形，满足 $j<0$ 时 $L_{ij}=0$，且对每个 $i$，$L_{i\bullet}$ 是 $K_i$ 的一个解消；再假设对每对 $(i,j)$ 及每个 $n>0$ 都有 $\mathrm L_nT(L_{ij})=0$。则存在函子同构
\begin{equation}\tag{11.6.3.2}
\mathrm L_\bullet T(K_\bullet)\simeq\mathrm H_\bullet(T(L_{\bullet\bullet})).
\end{equation}
\end{enumerate}
\begin{proof}
证明与（11.5.2）、（11.4.5）及（11.5.3）在有下界复形情形下的证明相同。这些论证的细节留给读者。
\end{proof}

\egapar{11.6.4}
对于在每个变元上加性的协变多函子，有完全类似的结果。例如，对双函子 $T$，两个超同调谱序列的 $E^2$ 项为
\begin{equation}\tag{11.6.4.1}
{}'E^2_{pq}=\mathrm H_p(\mathrm L_qT(K_\bullet,K'_\bullet)),
\end{equation}
\begin{equation}\tag{11.6.4.2}
{}''E^2_{pq}=\bigoplus_{q'+q''=q}\mathrm L_pT\bigl(\mathrm H_{q'}(K_\bullet),\mathrm H_{q''}(K'_\bullet)\bigr).
\end{equation}
这里仍是第二个谱序列正则；当复形 $K_\bullet,K'_\bullet$ 有下界，或所考虑的阿贝尔范畴的对象具有长度有固定上界的投射解消时，两个谱序列均双正则。此外：

\egastatement{11.6.5}{命题}
设 $C,C',C''$ 为三个阿贝尔范畴，$T:C\times C'\to C''$ 为双加性协变双函子。
\begin{enumerate}[label=(\roman*)]
\item $\mathrm L_\bullet T(K_\bullet,K'_\bullet)$ 对分别由 $C,C'$ 的对象组成的有下界复形 $K_\bullet,K'_\bullet$ 中的每一个而言都是同调函子。
\item 设 $K_\bullet,K'_\bullet$ 有下界。设 $L_{\bullet\bullet}=(L_{ij})$、$L'_{\bullet\bullet}=(L'_{ij})$ 为两个双复形，满足 $j<0$ 时 $L_{ij}=L'_{ij}=0$，对每个 $i$，$L_{i\bullet}$ 是 $K_i$ 的解消而 $L'_{i\bullet}$ 是 $K'_i$ 的解消；并且对 $n>0$ 及每组 $(i,j,h,k)$ 都有 $\mathrm L_nT(L_{ij},L'_{hk})=0$。则有函子同构
\begin{equation}\tag{11.6.5.1}
\mathrm L_\bullet T(K_\bullet,K'_\bullet)\simeq\mathrm H_\bullet(T(L_{\bullet\bullet},L'_{\bullet\bullet})).
\end{equation}
\item 再假设对每对 $(i,j)$ 及 $(h,k)$，函子 $A\mapsto T(A,L'_{hk})$、$A'\mapsto T(L_{ij},A')$ 分别在 $C,C'$ 中正合。则存在函子同构
\begin{equation}\tag{11.6.5.2}
\begin{aligned}
\mathrm L_\bullet T(K_\bullet,K'_\bullet)
&\simeq\mathrm H_\bullet(T(L_{\bullet\bullet},K'_\bullet))\\
&\simeq\mathrm H_\bullet(T(K_\bullet,L'_{\bullet\bullet})).
\end{aligned}
\end{equation}
\end{enumerate}
\begin{proof}
证明与（11.5.5）及（11.5.6）的证明类似。
\end{proof}

\subsection{函子相对于双复形的超同调}
\src{41--43}
\egapar{11.7.1}
设 $C$ 为每个对象都是某个投射对象之商的阿贝尔范畴。考虑由 $C$ 的对象组成、两个次数都有下界的双复形 $K_{\bullet\bullet}=(K_{ij})$；总可以归结到 $i<0$ 或 $j<0$ 时 $K_{ij}=0$ 的情形，以下就这样假设。可以把 $K_{\bullet\bullet}$ 看成由阿贝尔范畴 $\mathbf K$ 的对象 $K_{i\bullet}=(K_{ij})_{j\geq0}$ 组成的单复形，其中 $\mathbf K$ 是 $C$ 的非负次数复形范畴。

由引理（11.5.2.1）（更准确地说，由将其“箭头反向”所得的对偶引理）及（M，V，2.2），$K_{\bullet\bullet}$ 在范畴 $\mathbf K$ 中有一个投射 Cartan--Eilenberg 解消。这样的解消是一个由投射对象组成、各次数均非负的三重复形 $M_{\bullet\bullet\bullet}=(M_{ijk})$，使对每个 $i$，$M_{i\bullet\bullet}$、$\mathrm B_i^{\mathrm I}(M_{\bullet\bullet\bullet})$、$\mathrm Z_i^{\mathrm I}(M_{\bullet\bullet\bullet})$、$\mathrm H_i^{\mathrm I}(M_{\bullet\bullet\bullet})$ 分别在 $\mathbf K$ 中构成 $K_{i\bullet}$、$\mathrm B_i^{\mathrm I}(K_{\bullet\bullet})$、$\mathrm Z_i^{\mathrm I}(K_{\bullet\bullet})$、$\mathrm H_i^{\mathrm I}(K_{\bullet\bullet})$ 的投射解消；特别地，对每对 $(i,j)$，$M_{ij\bullet}$ 是 $K_{ij}$ 在 $C$ 中的投射解消。

\egastatement{11.7.2}{命题}
设 $T$ 为从 $C$ 到阿贝尔范畴 $C'$ 的加性协变函子。采用（11.7.1）的记号，则三重复形 $T(M_{\bullet\bullet\bullet})$ 所对应单复形的同调 $\mathrm H_\bullet(T(M_{\bullet\bullet\bullet}))$，典范地同构于 $K_{\bullet\bullet}$ 所对应单复形的超同调 $\mathrm L_\bullet T(K_{\bullet\bullet})$（11.6.2）；它也是六个双正则谱序列 ${}^{(t)}E$（$t=a,b,a',b',c,d$）的共同终项，其 $E^2$ 项由下列公式给出：
\[
\begin{aligned}
{}^{(a)}E^2_{pq}&=\mathrm L_pT(\mathrm H_q(K_{\bullet\bullet})),\\
{}^{(b)}E^2_{pq}&=\mathrm H_p(\mathrm L_q^{\mathrm{II}}T(K_{\bullet\bullet})),\\
{}^{(a')}E^2_{pq}&=\mathrm L_pT(\mathrm H_q^{\mathrm I}(K_{\bullet\bullet})),\\
{}^{(b')}E^2_{pq}&=\mathrm H_p(\mathrm L_qT(K_{\bullet\bullet})),\\
{}^{(c)}E^2_{pq}&=\mathrm L_pT(\mathrm H_q^{\mathrm{II}}(K_{\bullet\bullet})),\\
{}^{(d)}E^2_{pq}&=\mathrm H_p(\mathrm L_q^{\mathrm I}T(K_{\bullet\bullet})).
\end{aligned}
\]
这里再次提醒：对任意复形 $A_\bullet=(A_i)$，用 $F(A_\bullet)$ 表示由 $F(A_i)$ 组成的复形。例如，$\mathrm L_q^{\mathrm{II}}T(K_{\bullet\bullet})$ 表示复形 $(\mathrm L_q^{\mathrm{II}}T(K_{i\bullet}))_{i\geq0}$，其中 $\mathrm L_q^{\mathrm{II}}T(K_{i\bullet})$ 是函子 $T$ 相对于单复形 $K_{i\bullet}$ 的第 $q$ 个超同调。
\begin{proof}
以 $L_\bullet$ 表示与 $K_{\bullet\bullet}$ 相伴的单复形，从而 $L_i=\bigoplus_{r+s=i}K_{rs}$；令 $N_{ij}=\bigoplus_{r+s=i}M_{rsj}$。显然每个 $N_{i\bullet}$ 都是 $L_i$ 在 $C$ 中的投射解消；由（11.6.3）及（11.6.4），有函子同构
\[
\mathrm L_\bullet T(L_\bullet)\simeq\mathrm H_\bullet(T(N_{\bullet\bullet})).
\]
由于与双复形 $T(N_{\bullet\bullet})$ 相伴的单复形也与三重复形 $T(M_{\bullet\bullet\bullet})$ 相伴，第一断言得证。

此外，$\mathrm L_\bullet T(L_\bullet)$ 是函子 $T$ 相对于单复形 $L_\bullet$ 的两个超同调谱序列（11.6.2）的终项；这两个谱序列分别就是 ${}^{(b')}E$ 与 ${}^{(a)}E$。

再将 $M_{\bullet\bullet\bullet}$ 视为双复形 $U_{\bullet\bullet}$，其中 $U_{ij}=\bigoplus_{r+s=j}M_{irs}$。$\mathrm H_\bullet(T(M_{\bullet\bullet\bullet}))$ 仍是双复形 $T(U_{\bullet\bullet})$ 的两个谱序列的终项。对每个 $i$，$M_{i\bullet\bullet}$ 相对于单复形 $K_{i\bullet}$ 满足（11.6.3）的条件，故
\[
\mathrm H_q^{\mathrm{II}}(T(U_{i\bullet}))=\mathrm L_q^{\mathrm{II}}T(K_{i\bullet});
\]
所以 $T(U_{\bullet\bullet})$ 的第一个谱序列正是 ${}^{(b)}E$。另一方面，对每个 $r$，$M_{\bullet r\bullet}$ 是单复形 $K_{\bullet r}$ 的一个 Cartan--Eilenberg 解消；（M，XV，2）的计算给出
\[
\mathrm H_q^{\mathrm I}(T(M_{\bullet rs}))=T(\mathrm H_q^{\mathrm I}(M_{\bullet rs})),
\]
从而
\[
\mathrm H_q^{\mathrm I}(T(U_{\bullet j}))
=\bigoplus_{r+s=j}T(\mathrm H_q^{\mathrm I}(M_{\bullet rs})).
\]
换言之，单复形 $\mathrm H_q^{\mathrm I}(T(U_{\bullet\bullet}))$ 就是双复形 $T(\mathrm H_q^{\mathrm I}(M_{\bullet\bullet\bullet}))$ 所对应的单复形。对每个 $q$，$\mathrm H_q^{\mathrm I}(M_{\bullet\bullet\bullet})$ 在范畴 $\mathbf K$ 中是单复形 $\mathrm H_q^{\mathrm I}(K_{\bullet\bullet})$ 的投射解消；应用（11.6.3）可得
\[
\mathrm H_p^{\mathrm{II}}\bigl(T(\mathrm H_q^{\mathrm I}(M_{\bullet\bullet\bullet}))\bigr)
=\mathrm L_pT(\mathrm H_q^{\mathrm I}(K_{\bullet\bullet})),
\]
因而得到谱序列 ${}^{(a')}E$。最后，在三重复形 $M_{\bullet\bullet\bullet}$ 的定义中交换下标 $i,j$，并对所得三重复形应用上述论证，就得到 ${}^{(c)}E$ 与 ${}^{(d)}E$。
\end{proof}
称 $\mathrm L_\bullet T(K_{\bullet\bullet})$ 为 $T$ 相对于双复形 $K_{\bullet\bullet}$ 的\textbf{超同调}。

\egastatement{11.7.3}{注}
（i）由（11.6.3），$\mathrm L_\bullet T(K_{\bullet\bullet})$ 是 $C$ 中各次数都有下界的双复形范畴上的同调函子。

（ii）设 $M_{\bullet\bullet\bullet}$ 为 $C$ 的三重复形，对每对 $(r,s)$，$M_{rs\bullet}$ 是 $K_{rs}$ 的解消，且对每个三元组 $(i,j,k)$ 及 $n>0$ 都有 $\mathrm L_nT(M_{ijk})=0$。则有同构
\[
\mathrm L_\bullet T(K_{\bullet\bullet})\simeq\mathrm H_\bullet(T(M_{\bullet\bullet\bullet})).
\]
事实上，用（11.7.2）证明中的记号，$N_{i\bullet}$ 是 $L_i$ 的解消，且对 $n>0$ 及每对 $(i,j)$ 有 $\mathrm L_nT(N_{ij})=0$，故只需应用（11.6.3，iii）。

（iii）（11.7.2）的结果立即可推广到协变多函子。例如，设 $C'$ 为第二个每个对象都是投射对象之商的阿贝尔范畴，$K'_{\bullet\bullet}$ 为 $C'$ 中两个次数都有下界的双复形，$T:C\times C'\to C''$ 为到阿贝尔范畴 $C''$ 的加性协变双函子。若 $L_\bullet,L'_\bullet$ 分别为与 $K_{\bullet\bullet},K'_{\bullet\bullet}$ 相伴的单复形，则把 $T$ 相对于两个双复形的超同调定义为 $\mathrm L_\bullet T(L_\bullet,L'_\bullet)$；应用（11.6.4）、（11.6.5），像（11.7.2）一样可得六个以此超同调为终项的谱序列，其具体写法留给读者。\par

% ---- END c0_s11c.tex ----
% ---- BEGIN c0_s11d.tex ----
\subsection{单纯复形上同调的补充知识}
\src{43--46}
\egapar{11.8.1}
设 $A$ 为有限集，$\Sigma(A)$ 为由 $A$ 的元素组成的有限序列 $\sigma=(\alpha_0,\ldots,\alpha_h)$ 的集合（称为 $A$ 的“单形”）；置 $|\sigma|=\{\alpha_0,\ldots,\alpha_h\}$。回顾链复形 $C_\bullet(A)$：它是以 $\Sigma(A)$ 的元素为生成元的分次自由阿贝尔群，$(\alpha_0,\ldots,\alpha_h)$ 的次数为 $h$，微分定义为
\[
d(\alpha_0,\ldots,\alpha_h)
=\sum_{j=0}^h(-1)^j(\alpha_0,\ldots,\widehat{\alpha_j},\ldots,\alpha_h).
\]
在 $C_\bullet(A)$ 中，以两项相等的链 $\sigma=(\alpha_0,\ldots,\alpha_h)$ 及链 $\pi(\sigma)-\varepsilon_\pi\sigma$ 生成的子群 $D_\bullet(A)$ 是子复形；这里 $\pi$ 为任意置换，$\pi(\sigma)=(\alpha_{\pi(0)},\ldots,\alpha_{\pi(h)})$，$\varepsilon_\pi$ 是 $\pi$ 的符号。$D_\bullet(A)$ 的元素称为退化链。置 $L_\bullet(A)=C_\bullet(A)/D_\bullet(A)$，有自然的复形同态 $p:C_\bullet(A)\to L_\bullet(A)$。

另一方面，按如下方法定义复形同态 $j:L_\bullet(A)\to C_\bullet(A)$：在 $A$ 上选定全序；对于单形 $\sigma=(\alpha_0,\ldots,\alpha_h)$ 模 $D_\bullet(A)$ 的类，若有两项相等，则对应于零，否则对应于将各 $\alpha_i$ 按递增次序排列所得的序列 $(\beta_0,\ldots,\beta_h)$。显然 $p\circ j$ 是 $L_\bullet(A)$ 的恒等映射。\transnote{原文此处未写出排序置换的符号。若将其理解为对任意代表单形的规则，须计入该符号；等价地，可以只对严格递增单形所给的基规定 $j$，再线性延拓。这里保留原文表述。}

\egapar{11.8.2}
设 $B$ 为另一个有限集。若 $d',d''$ 分别为 $C_\bullet(A),C_\bullet(B)$ 的微分，则张量积复形 $C_\bullet(A)\otimes C_\bullet(B)$ 可以看成以 $\Sigma(A)\times\Sigma(B)$ 的元素为生成元的自由阿贝尔群，其微分为
\[
d(\sigma,\tau)=(d'\sigma,\tau)+(-1)^{h+1}(\sigma,d''\tau)
\quad\text{若 }\operatorname{Card}|\sigma|=h+1.
\]
\transnote{上式及其条件按原版排印保留。以（11.8.1）的次数约定写通常张量积微分时，符号应依单形的次数而定；含重复顶点时，$\operatorname{Card}|\sigma|$ 也不能代替序列长度。}
自然同态 $C_\bullet(A)\to L_\bullet(A)$ 与 $C_\bullet(B)\to L_\bullet(B)$ 给出同态
\[
p:C_\bullet(A)\otimes C_\bullet(B)\longrightarrow L_\bullet(A)\otimes L_\bullet(B),
\]
其右端同构于
\[
\frac{C_\bullet(A)\otimes C_\bullet(B)}
{C_\bullet(A)\otimes D_\bullet(B)+D_\bullet(A)\otimes C_\bullet(B)}.
\]
同样，利用（11.8.1）中通过 $A,B$ 上的全序定义的同态 $L_\bullet(A)\to C_\bullet(A)$、$L_\bullet(B)\to C_\bullet(B)$，可定义同态
\[
j:L_\bullet(A)\otimes L_\bullet(B)\longrightarrow C_\bullet(A)\otimes C_\bullet(B),
\]
使 $p\circ j$ 为恒等映射。采用这些记号：

\egastatement{11.8.3}{命题}
存在同伦
\[
h:C_\bullet(A)\otimes C_\bullet(B)\longrightarrow C_\bullet(A)\otimes C_\bullet(B),
\]
使 $h(\sigma,\tau)$ 是满足 $|\sigma_i|\subset|\sigma|$、$|\tau_i|\subset|\tau|$ 的单形对 $(\sigma_i,\tau_i)$ 的线性组合，且令 $f=j\circ p$，则有
\begin{equation}\tag{11.8.3.1}
f-1=h\circ d+d\circ h.
\end{equation}
\begin{proof}
只需按 $\sigma,\tau$ 次数之和归纳，定义每个单形对上的 $h$；次数之和为零时可取 $h=0$。令
\[
\omega=f(\sigma,\tau)-(\sigma,\tau)-h(d(\sigma,\tau)).
\]
由归纳假设及 $d$ 的定义，$\omega\in C_\bullet(|\sigma|)\otimes C_\bullet(|\tau|)$。又由（11.8.3.1）及归纳假设，
\[
d\omega=f(d(\sigma,\tau))-d(\sigma,\tau)-dh(d(\sigma,\tau))
=h(d(d(\sigma,\tau)))=0.
\]
对于 $q>0$，有 $\mathrm H_q(C_\bullet(A))=0$（G，I，3.7.4）；故由 Künneth 公式（G，I，5.5.2），也有 $\mathrm H_q(C_\bullet(A)\otimes C_\bullet(B))=0$。将这一事实中的 $A,B$ 换成 $|\sigma|,|\tau|$，可知存在 $\omega'\in C_\bullet(|\sigma|)\otimes C_\bullet(|\tau|)$，使 $\omega=d\omega'$。取 $h(\sigma,\tau)=\omega'$，即对该单形对满足（11.8.3.1），从而归纳可以继续。
\end{proof}

\egapar{11.8.4}
采用（11.8.2）的记号，若 $|\sigma|\subset|\sigma'|$、$|\tau|\subset|\tau'|$，就置 $(\sigma,\tau)\leq(\sigma',\tau')$。$\Sigma(A)\times\Sigma(B)$ 上的一个\textbf{系数系统} $\SSS$ 由阿贝尔群族 $(\Gamma_{\sigma,\tau})$ 及对 $(\sigma,\tau)\leq(\sigma',\tau')$ 给定的同态 $\Gamma_{\sigma,\tau}\to\Gamma_{\sigma',\tau'}$ 组成；其中 $\Gamma_{\sigma,\tau}$ 只依赖集合 $|\sigma|,|\tau|$，且这些同态对于上述预序构成归纳系统。

定义上链复形 $C^\bullet(A,B;\SSS)$ 为各族 $\lambda=(\lambda(\sigma,\tau))$ 所成的复形，其中 $(\sigma,\tau)$ 遍历 $\Sigma(A)\times\Sigma(B)$，且 $\lambda(\sigma,\tau)\in\Gamma_{\sigma,\tau}$。微分如下给出：若 $d(\sigma,\tau)=\sum_i\pm(\sigma_i,\tau_i)$，则对每个 $i$，有 $|\sigma_i|\subset|\sigma|$、$|\tau_i|\subset|\tau|$，并令
\[
d\lambda(\sigma,\tau)=\sum_i\pm\lambda(\sigma_i,\tau_i),
\]
右端的各项表示它们在 $\Gamma_{\sigma,\tau}$ 中的典范像。

若上链 $\lambda\in C^\bullet(A,B;\SSS)$ 满足：当 $\sigma,\tau$ 中任一单形有两项相等时 $\lambda(\sigma,\tau)=0$，并且对任意下标置换 $\pi,\pi'$，有
\[
\lambda(\pi(\sigma),\tau)=\varepsilon_\pi\lambda(\sigma,\tau),\qquad
\lambda(\sigma,\pi'(\tau))=\varepsilon_{\pi'}\lambda(\sigma,\tau),
\]
\begin{sloppypar}
则称它为\textbf{双交错上链}。显然这些上链构成 $C^\bullet(A,B;\SSS)$ 的子复形 $L^\bullet(A,B;\SSS)$。
\end{sloppypar}


\egastatement{11.8.5}{命题}
典范嵌入 $L^\bullet(A,B;\SSS)\to C^\bullet(A,B;\SSS)$ 在这两个复形的上同调之间给出同构。
\begin{proof}
若 $p,j$ 取（11.8.2）中的意义，则映射 ${}^tp:\lambda\mapsto\lambda\circ p$、${}^tj:\lambda\mapsto\lambda\circ j$ 分别在 $L^\bullet(A,B;\SSS)$ 与 $C^\bullet(A,B;\SSS)$ 上有定义；前一个就是典范嵌入。由于 ${}^tj\circ{}^tp$ 为恒等映射，只需证明 ${}^tp\circ{}^tj$ 与恒等映射同伦。由（11.8.3），${}^th:\lambda\mapsto\lambda\circ h$ 在 $C^\bullet(A,B;\SSS)$ 中有定义，故可转置恒等式（11.8.3.1），得到所需结果。
\end{proof}

\egapar{11.8.6}
命题（11.8.5）将 $L^\bullet(A,B;\SSS)$ 的上同调计算归结为 $C^\bullet(A,B;\SSS)$ 的上同调计算。另一方面，回顾 Eilenberg--Zilber 定理（G，I，3.10.2）：后一个上同调典范地同构于如下定义的上链复形的上同调。构造链复形 $P_\bullet(A,B)$，其元素是满足 $\sigma,\tau$ 次数相同的 $(\sigma,\tau)\in\Sigma(A)\times\Sigma(B)$ 的线性组合。其微分给为
\[
\begin{gathered}
d:(\sigma,\tau)\longmapsto\sum_{j,k}(-1)^{j+k}(\sigma_j,\tau_k),\\
\text{其中 }d\sigma=\sum_j(-1)^j\sigma_j,\quad d\tau=\sum_k(-1)^k\tau_k.
\end{gathered}
\]
\transnote{此处双重求和公式照录原版；按通常的对角单纯复形构造，应同步取两因子的第 $i$ 个面并以 $(-1)^i$ 求和。原版后面的上链公式亦采用上述双重求和，译文同样保留而不静默改写。}
于是有两个典范复形同态
\[
f:P_\bullet(A,B)\to C_\bullet(A)\otimes C_\bullet(B),\qquad
g:C_\bullet(A)\otimes C_\bullet(B)\to P_\bullet(A,B),
\]
并可证明（同上）存在同伦 $h,h'$，使
\[
f\circ g-1=d\circ h+h\circ d,\qquad
g\circ f-1=d\circ h'+h'\circ d.
\]
此外，$f(\sigma,\tau)\in C_\bullet(|\sigma|)\otimes C_\bullet(|\tau|)$、$g(\sigma,\tau)\in P_\bullet(|\sigma|,|\tau|)$；同伦还可选为满足
\[
h(\sigma,\tau)\in C_\bullet(|\sigma|)\otimes C_\bullet(|\tau|),\qquad
h'(\sigma,\tau)\in P_\bullet(|\sigma|,|\tau|).
\]
这是因为 $h(\sigma,\tau)$、$h'(\sigma,\tau)$ 可以按次数之和归纳定义，并且 $q>0$ 时 $\mathrm H_q(C_\bullet(|\sigma|)\otimes C_\bullet(|\tau|))$ 与 $\mathrm H_q(P_\bullet(|\sigma|,|\tau|))$ 都为零（同上）；像（11.8.3）那样论证即可。

定义 $P^\bullet(A,B;\SSS)$ 为各族 $\lambda=(\lambda(\sigma,\tau))$ 所成的复形，其中 $(\sigma,\tau)$ 遍历次数相同的单形对，且 $\lambda(\sigma,\tau)\in\Gamma_{\sigma,\tau}$。由于
\[
d\sigma=\sum_j(-1)^j\sigma_j\in C_\bullet(|\sigma|),\qquad
d\tau=\sum_k(-1)^k\tau_k\in C_\bullet(|\tau|),
\]
故
\[
d\lambda(\sigma,\tau)=\sum_{j,k}(-1)^{j+k}\lambda(\sigma_j,\tau_k)
\]
有定义，并给出 $P^\bullet(A,B;\SSS)$ 的微分。由上述说明，${}^tf:\lambda\mapsto\lambda\circ f$、${}^tg:\lambda\mapsto\lambda\circ g$、${}^th:\lambda\mapsto\lambda\circ h$、${}^th':\lambda\mapsto\lambda\circ h'$ 均有定义；因此 ${}^tf\circ{}^tg$、${}^tg\circ{}^tf$ 都与恒等映射同伦。这就得到 $C^\bullet(A,B;\SSS)$ 与 $P^\bullet(A,B;\SSS)$ 的上同调之间所需的同构。

\egastatement{11.8.7}{注}
将（11.8.3）的同一论证应用于 $C_\bullet(A),L_\bullet(A)$，可知：若 $j,p$ 按（11.8.1）定义，则 $f=j\circ p$ 仍满足（11.8.3.1）型关系，其中 $|h(\sigma)|\subset|\sigma|$。像（11.8.5）那样，可由此得到 $L^\bullet(A;\SSS)$ 到 $C^\bullet(A;\SSS)$ 的上同调同构；这两个复形以显然的方式定义。这就是（G，I，3.8.1）中概述证明的结果。

\egapar{11.8.8}
再采用（11.8.2）的记号与假设，考虑 $\Sigma(A)\times\Sigma(B)$ 上系数系统的复形 $\SSS^\bullet=(\SSS^k)$。对每个 $(\sigma,\tau)$，$\Gamma^k_{\sigma,\tau}$（$k\in\Z$）构成阿贝尔群复形；当 $(\sigma,\tau)\leq(\sigma',\tau')$ 时，有交换图
\[
\begin{tikzcd}
\Gamma^k_{\sigma,\tau}\arrow[r]\arrow[d]&\Gamma^{k+1}_{\sigma,\tau}\arrow[d]\\
\Gamma^k_{\sigma',\tau'}\arrow[r]&\Gamma^{k+1}_{\sigma',\tau'}.
\end{tikzcd}
\]
立即可验证
\[
C^\bullet(A,B;\SSS^\bullet)=(C^h(A,B;\SSS^k))_{(h,k)\in\Z^2}
\]
是阿贝尔群双复形，且 $L^\bullet(A,B;\SSS^\bullet)=(L^h(A,B;\SSS^k))$ 是其子双复形。

\egastatement{11.8.9}{命题}
典范嵌入 $L^\bullet(A,B;\SSS^\bullet)\to C^\bullet(A,B;\SSS^\bullet)$ 在两个双复形的上同调之间给出同构。
\begin{proof}
简记 $C^{\bullet\bullet}=C^\bullet(A,B;\SSS^\bullet)$、$L^{\bullet\bullet}=L^\bullet(A,B;\SSS^\bullet)$。由于 $h<0$ 时 $C^{hk}=L^{hk}=0$，两个双复形的第二谱序列正则（11.3.3）；同态 $L^{\bullet\bullet}\to C^{\bullet\bullet}$ 因而给出谱序列态射 ${}''E(L^{\bullet\bullet})\to{}''E(C^{\bullet\bullet})$，其在 $E_2$ 项上为
\begin{equation}\tag{11.8.9.1}
\mathrm H_{\mathrm{II}}^p(\mathrm H_{\mathrm I}^q(L^{\bullet\bullet}))
\longrightarrow\mathrm H_{\mathrm{II}}^p(\mathrm H_{\mathrm I}^q(C^{\bullet\bullet})).
\end{equation}
对每个 $k\in\Z$，（11.8.3）说明同态 $\mathrm H^q(L^{\bullet k})\to\mathrm H^q(C^{\bullet k})$ 为双射；故结论由（11.1.5）得出。
\end{proof}

\egapar{11.8.10}
同样，采用（11.8.6）的记号，有典范双复形同态 $C^\bullet(A,B;\SSS^\bullet)\to P^\bullet(A,B;\SSS^\bullet)$（记号的含义是显然的）；将（11.8.9）的同一论证改为以（11.8.6）为基础，便知此同态仍在上同调上给出同构。

\subsection{关于有限型复形的一个引理}
\src{46--48}
\egastatement{11.9.1}{命题}
设 $C$ 为阿贝尔范畴，$K',K''$ 为其对象集合的子集，$K''\subset K'$，并满足：
\begin{enumerate}[label=(\roman*)]
\item 对每个对象 $A'\in K'$ 及每个满态射 $u:A\to A'$，存在对象 $B\in K''$ 和态射 $v:B\to A$，使 $uv$ 为满态射。
\item 对每对对象 $A\in K''$、$B\in K'$ 及每个满态射 $u:A\to B$，$\Ker u$ 属于 $K'$。
\item 两个 $K''$ 中对象的积仍属于 $K''$。
\end{enumerate}
设 $P_\bullet=(P_i)_{i\in\Z}$ 为 $C$ 的复形，对每个 $i$ 都有 $\mathrm H_i(P_\bullet)\in K'$，并存在 $d$ 使 $i<d$ 时 $\mathrm H_i(P_\bullet)=0$。则存在 $C$ 的复形 $Q_\bullet=(Q_i)$ 及复形态射 $u:Q_\bullet\to P_\bullet$，使每个 $Q_i\in K''$，$i<d$ 时 $Q_i=0$，且相应的态射 $\mathrm H_\bullet(Q_\bullet)\to\mathrm H_\bullet(P_\bullet)$ 为同构。
\begin{proof}
先证明性质（i）的下述推论。

（i bis）设 $u:C\to B$ 为 $C$ 中的满态射，$A\in K'$，$v:A\to B$ 为态射。则存在对象 $D\in K''$、满态射 $u':D\to A$ 及态射 $v':D\to C$，使图
\begin{equation}\tag{11.9.1.1}
\begin{tikzcd}
D\arrow[r,"u'"]\arrow[d,"v'"']&A\arrow[d,"v"]\\
C\arrow[r,"u"']&B
\end{tikzcd}
\end{equation}
交换。

事实上，考虑 $C$ 中的纤维积 $C\times_B A$ 及典范投影 $p:C\times_B A\to C$、$q:C\times_B A\to A$，给出交换图
\begin{equation}\tag{11.9.1.2}
\begin{tikzcd}
C\times_B A\arrow[r,"q"]\arrow[d,"p"']&A\arrow[d,"v"]\\
C\arrow[r,"u"']&B.
\end{tikzcd}
\end{equation}
已知 $q$ 的余核为 $A/v^{-1}(u(C))$（[27]，第1--12页）。由于 $u$ 满，$u(C)=B$、$v^{-1}(u(C))=A$，故 $q$ 为满态射。将（i）应用于 $q$，可得 $D\in K''$ 及 $w:D\to C\times_B A$，使 $qw$ 满；取 $u'=qw$、$v'=pw$ 即可。

现在归纳证明命题。对某个 $i\geq d-1$，设已对 $j\leq i$ 构造了对象 $Q_j$、态射 $d_j:Q_j\to Q_{j-1}$ 及 $u_j:Q_j\to P_j$，满足 $j<d$ 时 $Q_j=0$，且对 $j\leq i$，有 $d_{j-1}d_j=0$、$d'_ju_j=u_{j-1}d_j$。另假设以下条件成立：

$(\mathrm I_i)$：$j\leq i$ 时 $Q_j\in K''$，$j<i$ 时 $\mathrm B_j(Q_\bullet)\in K'$。

$(\mathrm{II}_i)$：$j<i$ 时，族 $(u_k)_{k\leq i}$ 给出的同态 $\mathrm H_j(Q_\bullet)\to\mathrm H_j(P_\bullet)$ 为同构。

$(\mathrm{III}_i)$：复合态射
\[
v_i:\mathrm Z_i(Q_\bullet)\longrightarrow\mathrm Z_i(P_\bullet)
\longrightarrow\mathrm H_i(P_\bullet)
\]
为满态射；这里第一箭头是 $u_i$ 的限制，第二箭头是典范态射。

由（ii）及 $(\mathrm I_i)$，满态射 $Q_i\to\mathrm B_{i-1}(Q_\bullet)$ 的核 $\mathrm Z_i(Q_\bullet)$ 属于 $K'$。再由（ii）及 $(\mathrm{III}_i)$，$N_i=\Ker v_i$ 也属于 $K'$。\transnote{此句按原文保留；（ii）只对来源对象属于 $K''$ 的满态射作出保证，而上句只得到 $\mathrm Z_i(Q_\bullet)\in K'$。一般表述下这一引用需要补充论证或相应的核稳定性假设。}
由（i bis），存在 $Q'_{i+1}\in K''$、满态射 $d'_{i+1}:Q'_{i+1}\to N_i$ 和态射 $u'_{i+1}:Q'_{i+1}\to P_{i+1}$，使图
\begin{equation}\tag{11.9.1.3}
\begin{tikzcd}
Q'_{i+1}\arrow[r,"d'_{i+1}"]\arrow[d,"u'_{i+1}"']&N_i\arrow[d,"u_i"]\\
P_{i+1}\arrow[r,"d'_{i+1}"']&\mathrm B_i(P_\bullet)
\end{tikzcd}
\end{equation}
交换。\transnote{原文在微分及新构造映射上重复使用带撇号的 $d$；此图保留原记号，箭头的定义域和值域可区分它们。}

典范态射 $\mathrm Z_{i+1}(P_\bullet)\to\mathrm H_{i+1}(P_\bullet)$ 为满态射，而 $\mathrm H_{i+1}(P_\bullet)\in K'$，故由（i），存在 $Q''_{i+1}\in K''$ 及态射 $u''_{i+1}:Q''_{i+1}\to\mathrm Z_{i+1}(P_\bullet)$，使复合态射到 $\mathrm H_{i+1}(P_\bullet)$ 为满态射。取 $d''_{i+1}:Q''_{i+1}\to N_i$ 为零，则图
\begin{equation}\tag{11.9.1.4}
\begin{tikzcd}
Q''_{i+1}\arrow[r,"d''_{i+1}"]\arrow[d,"u''_{i+1}"']&N_i\arrow[d,"u_i"]\\
\mathrm Z_{i+1}(P_\bullet)\arrow[r,"d'_{i+1}"']&P_i
\end{tikzcd}
\end{equation}
交换，其中下面的水平箭头为零。

取 $Q_{i+1}=Q'_{i+1}\times Q''_{i+1}$，由（iii）它属于 $K''$；置 $d_{i+1}=d'_{i+1}+d''_{i+1}$、$u_{i+1}=u'_{i+1}+u''_{i+1}$。由于 $d_{i+1}(Q_{i+1})=d'_{i+1}(Q'_{i+1})=N_i\subset\mathrm Z_i(Q_\bullet)$，有 $d_id_{i+1}=0$，并按通常记号有 $\mathrm B_i(Q_\bullet)=N_i$，这就验证了 $(\mathrm I_{i+1})$。两个图的交换性给出 $d'_{i+1}u_{i+1}=u_id_{i+1}$。

由 $N_i$ 的定义，族 $(u_k)_{k\leq i+1}$ 给出的态射
\[
\mathrm H_i(Q_\bullet)=\mathrm Z_i(Q_\bullet)/N_i\longrightarrow\mathrm H_i(P_\bullet)
\]
正是 $v_i$ 经商得到的态射；由于 $v_i$ 满，它为同构，故 $(\mathrm{II}_{i+1})$ 成立。最后，按定义 $Q''_{i+1}\subset\mathrm Z_{i+1}(Q_\bullet)$；$u''_{i+1}$ 的选择说明态射
\[
v_{i+1}:\mathrm Z_{i+1}(Q_\bullet)\to\mathrm Z_{i+1}(P_\bullet)\to\mathrm H_{i+1}(P_\bullet)
\]
为满态射，因为其在 $Q''_{i+1}$ 上的限制已是满态射。这验证了 $(\mathrm{III}_{i+1})$；归纳可以继续，命题得证。
\end{proof}

\egastatement{11.9.2}{推论}
设 $A$ 为不必交换的诺特环，$P_\bullet=(P_i)_{i\in\Z}$ 为右 $A$-模复形。假设各 $\mathrm H_i(P_\bullet)$ 是有限型 $A$-模，并存在 $d$ 使 $i<d$ 时 $\mathrm H_i(P_\bullet)=0$。则存在由有限秩自由右 $A$-模组成的复形 $Q_\bullet=(Q_i)$ 及复形同态 $u:Q_\bullet\to P_\bullet$，使 $i<d$ 时 $Q_i=0$，且相应的同调同态 $\mathrm H_\bullet(Q_\bullet)\to\mathrm H_\bullet(P_\bullet)$ 为双射。
\begin{proof}
在（11.9.1）中取 $C$ 为右 $A$-模范畴，$K'$ 为有限型 $A$-模的集合，$K''$ 为有限秩自由 $A$-模的集合。由 $A$ 诺特，（i）、（ii）、（iii）的验证是直接的。
\end{proof}

\egastatement{11.9.3}{注}
（i）在（11.9.2）的条件下，再假设各 $P_i$ 为平坦右 $A$-模。则对每个左 $A$-模 $M$，复形同态 $u\otimes1:Q_\bullet\otimes_A M\to P_\bullet\otimes_A M$ 仍给出同调同构
\[
\mathrm H_\bullet(Q_\bullet\otimes_A M)\simeq\mathrm H_\bullet(P_\bullet\otimes_A M),
\]
这一点将在第 III 章中证明。

（ii）若不假设 $A$ 诺特，（11.9.2）的结论未必正确。事实上，将它应用于除一个项外其余均为零的复形，就会推出每个有限型左 $A$-模都有由有限型自由模组成的解消；一般并非如此（见 Bourbaki，《交换代数》，第 I 章，§2，习题6）。不过，可以把 $A$ 诺特的假设换成较弱的假设：各 $\mathrm H_i(P_\bullet)$ 都具有有限 $\infty$-表现（见第 IV 章）。

\subsection{有限长度模复形的 Euler--Poincaré 示性数}
\src{48--49}
\egapar{11.10.1}
设 $A$ 为不必交换的环，
\begin{equation}\tag{11.10.1.1}
M^\bullet:0\longrightarrow M^0\longrightarrow M^1\longrightarrow\cdots\longrightarrow M^n\longrightarrow0
\end{equation}
为有限长度左 $A$-模组成的复形。此复形的\textbf{Euler--Poincaré 示性数}定义为数
\begin{equation}\tag{11.10.1.2}
\chi(M^\bullet)=\sum_{i=0}^n(-1)^i\len M^i.
\end{equation}

\egastatement{11.10.2}{命题}
对每个由有限长度左 $A$-模组成的有限复形 $M^\bullet$，有
\[
\chi(M^\bullet)=\chi(\mathrm H^\bullet(M^\bullet)),
\]
其中把 $\mathrm H^\bullet(M^\bullet)$ 看成微分为零的复形。特别地，若（11.10.1.1）正合，则 $\chi(M^\bullet)=0$。
\begin{proof}
简记 $B^i=\mathrm B^i(M^\bullet)$、$Z^i=\mathrm Z^i(M^\bullet)$、$H^i=\mathrm H^i(M^\bullet)=Z^i/B^i$；它们都是有限长度模。由正合列
\[
0\to B^i\to Z^i\to H^i\to0,\qquad
0\to Z^i\to M^i\to B^{i+1}\to0
\]
可得
\[
\len Z^i=\len H^i+\len B^i,\qquad
\len M^i=\len Z^i+\len B^{i+1},
\]
故
\[
\len M^i-\len H^i=\len B^{i+1}+\len B^i.
\]
将这一关系乘以 $(-1)^i$，再对 $0\leq i\leq n$ 求和，注意 $B^0=B^{n+1}=0$，即得所需等式。
\end{proof}

\egastatement{11.10.3}{推论}
设 $E=(E_r^{pq})$ 为环 $A$ 上模范畴中的谱序列。假设各 $E_2^{pq}$ 为有限长度 $A$-模，且只有有限多个 $(p,q)$ 满足 $E_2^{pq}\ne0$。则所有复形 $E_r^\bullet=(E_r^n)_{n\in\Z}$（11.1.1）的 Euler--Poincaré 示性数 $\chi(E_r^\bullet)$ 都相等。若再假设 $E$ 弱收敛，并对每个 $n\in\Z$ 置 $E_\infty^n=\bigoplus_{p+q=n}E_\infty^{pq}$，则还有 $\chi(E_r^\bullet)=\chi(E_\infty^\bullet)$，其中 $E_\infty^\bullet=(E_\infty^n)_{n\in\Z}$ 看成微分为零的复形。
\begin{proof}
若 $E_2^{pq}=0$，则 $2\leq r<+\infty$ 时 $E_r^{pq}=0$；因此所有 $E_r^\bullet$ 都是由有限长度 $A$-模组成的有限复形。对于有限的 $r$，关系 $\chi(E_r^\bullet)=\chi(E_{r+1}^\bullet)$ 由（11.10.2）及 $\mathrm H^\bullet(E_r^\bullet)$ 与 $E_{r+1}^\bullet$ 的同构（看成微分为零的复形）得出。各 $E_2^{pq}$ 的有限长度性保证对每对 $(p,q)$，序列 $(\mathrm B_k(E_2^{pq}))_k$ 与 $(\mathrm Z_k(E_2^{pq}))_k$ 都稳定；再由弱收敛性及只有有限多个非零 $E_2^{pq}$，可知存在 $r\geq2$，使对每对 $(p,q)$ 均有 $E_r^{pq}=E_\infty^{pq}$。这就证明了关于 $\chi(E_\infty^\bullet)$ 的断言。
\end{proof}\par

% ---- END c0_s11d.tex ----
% ---- BEGIN c0_s12a.tex ----
\section{层上同调的补充知识}
\subsection{赋环空间上模层的上同调}
\src{49--57}
\egapar{12.1.1}
设 $(X,\OO_X)$ 为赋环空间。回顾：对每个 $\OO_X$-模 $\F$，定义了上同调 $\mathrm H^\bullet(X,\F)$；这是从 $\OO_X$-模范畴 $C(X)$ 到阿贝尔群范畴的泛上同调函子（T，2.2），亦即左正合函子 $\F\mapsto\Gamma(X,\F)$ 的导出函子。函子 $\mathrm H^\bullet(X,\F)$ 同构于在 $X$ 上阿贝尔群层范畴中以同样方式定义的上同调函子在 $C(X)$ 上的限制（G，II，7.2.1）。

\egapar{12.1.2}
置 $A=\Gamma(X,\OO_X)$。由于 $A$ 的每个元素都定义阿贝尔群 $\Gamma(X,\F)$ 的一个自同态，由函子性，它也定义 $\mathrm H^\bullet(X,\F)$ 的一个 $\partial$-函子自同态；这些自同态在每个 $\mathrm H^p(X,\F)$ 上定义 $A$-模结构，且算子 $\partial$ 是 $A$-线性的。此外，对任意两个非负整数 $p,q$ 及两个 $\OO_X$-模 $\F,\G$，有称为\textbf{杯积}的 $A$-模同态
\begin{equation}\tag{12.1.2.1}
\mathrm H^p(X,\F)\otimes_A\mathrm H^q(X,\G)
\longrightarrow\mathrm H^{p+q}(X,\F\otimes_{\OO_X}\G)
\end{equation}
（G，II，6.6）。这些同态使 $S=\bigoplus_{p\geq0}\mathrm H^p(X,\OO_X)$ 成为反交换分次 $A$-代数，并使 $\bigoplus_{p\geq0}\mathrm H^p(X,\F)$ 成为分次 $S$-模。

对 $X$ 的任意开覆盖 $\mathfrak U$，我们总用 $C^\bullet(\mathfrak U,\F)$ 表示其神经在系数系统 $\Gamma(U_\sigma,\F)$ 中取值的\emph{交错上链}复形（这一约定与（G，II，5.1）不同）。显然它是分次 $A$-模，因而它的上同调群 $\mathrm H^p(\mathfrak U,\F)$ 都具有 $A$-模结构；此外，典范映射 $\mathrm H^p(\mathfrak U,\F)\to\mathrm H^p(X,\F)$（G，II，5.4）都是 $A$-线性的。

\egapar{12.1.3}
设 $(X',\OO_{X'})$ 为第二个赋环空间，$f=(\psi,\theta):X'\to X$ 为态射。置 $A'=\Gamma(X',\OO_{X'})$；$\psi$-态射 $\theta$ 典范地定义环同态 $A\to A'$。设 $\F$ 为 $\OO_X$-模，$\F'$ 为 $\OO_{X'}$-模。对每个 $f$-态射 $u:\F\to\F'$（$\OI$，4.4.1），将证明：对每个 $p\geq0$，可定义双同态
\begin{equation}\tag{12.1.3.1}
u_p:\mathrm H^p(X,\F)\longrightarrow\mathrm H^p(X',\F').
\end{equation}
事实上，由于 $\psi^*$ 在 $X$ 上阿贝尔群层范畴中正合，$\F\mapsto\mathrm H^\bullet(X',\psi^*\F)$ 是该范畴上的 $\partial$-函子；已知有典范的 $\partial$-函子同态
\begin{equation}\tag{12.1.3.2}
\mathrm H^\bullet(X,\F)\longrightarrow\mathrm H^\bullet(X',\psi^*\F),
\end{equation}
它由在零次上等于典范同态 $\Gamma(X,\F)\to\Gamma(X',\psi^*\F)$ 的条件唯一确定（T，3.2.2）。此外，$A$ 的每个元素都确定 $\Gamma(X,\F)$ 的自同态 $\mu$ 及 $\Gamma(X',\psi^*\F)$ 的自同态 $\mu'$，使图
\begin{equation}\tag{12.1.3.3}
\begin{tikzcd}
\Gamma(X,\F)\arrow[r]\arrow[d,"\mu"']&\Gamma(X',\psi^*\F)\arrow[d,"\mu'"]\\
\Gamma(X,\F)\arrow[r]&\Gamma(X',\psi^*\F)
\end{tikzcd}
\end{equation}
交换。由泛上同调函子态射延拓的唯一性（T，2.2），$\mu,\mu'$ 唯一延拓到上同调，并使类似于（12.1.3.3）的图交换；这意味着（12.1.3.2）是 $A$-模同态。

再注意 $f^*\F=\psi^*\F\otimes_{\psi^*\OO_X}\OO_{X'}$，故有从 $\psi^*\OO_X$-模 $\psi^*\F$ 到 $\OO_{X'}$-模 $f^*\F$ 的典范双同态。由函子性，可得函子性双同态
\begin{equation}\tag{12.1.3.4}
\mathrm H^p(X',\psi^*\F)\longrightarrow\mathrm H^p(X',f^*\F),
\end{equation}
相应的环为 $A,A'$。将它与（12.1.3.2）复合，得到对 $\F$ 具有函子性的典范双同态
\begin{equation}\tag{12.1.3.5}
\theta_p:\mathrm H^p(X,\F)\longrightarrow\mathrm H^p(X',f^*\F).
\end{equation}
最后，由 $u^\sharp:f^*\F\to\F'$ 的函子性，得到 $A'$-模同态 $\mathrm H^p(X',f^*\F)\to\mathrm H^p(X',\F')$；与（12.1.3.5）复合即得（12.1.3.1）。

设 $f'=(\psi',\theta'):X''\to X'$ 为第二个赋环空间态射，$f''=f\circ f'$。利用 $\psi'^*$ 与张量积的可交换性（$\OI$，4.3.3），立即可验证：将双同态 $\mathrm H^p(X',f^*\F)\to\mathrm H^p(X'',f'^*(f^*\F))$ 与（12.1.3.5）复合，所得就是相应的双同态 $\mathrm H^p(X,\F)\to\mathrm H^p(X'',f''^*\F)$。

\egapar{12.1.4}
（12.1.3.2）还可以直接定义如下。取 $\F$ 在 $X$ 上阿贝尔群层范畴中的内射解消 $\LL^\bullet=(\LL^i)$；由于 $\psi^*$ 正合，$\psi^*(\LL^\bullet)$ 是 $\psi^*\F$ 的一个层解消。若 $\LL'^\bullet=(\LL'^i)$ 是 $\psi^*\F$ 在 $X'$ 上阿贝尔群层范畴中的内射解消，则存在与增广相容的阿贝尔群层复形态射 $\psi^*(\LL^\bullet)\to\LL'^\bullet$（M，V，1.1 a）），它在同伦意义下唯一确定。于是有阿贝尔群复形同态
\[
\Gamma(X,\LL^\bullet)\longrightarrow\Gamma(X',\psi^*\LL^\bullet)
\longrightarrow\Gamma(X',\LL'^\bullet).
\]
其复合在上同调上给出 $\partial$-函子态射 $\mathrm H^\bullet(X,\F)\to\mathrm H^\bullet(X',\psi^*\F)$；由于在零次上与（12.1.3.2）一致，它就是该同态（T，2.2）。

考虑 $X$ 的开覆盖 $\mathfrak U=(U_\alpha)$，及其逆像覆盖 $\mathfrak U'=(f^{-1}(U_\alpha))$。对 $X$ 的各开集 $V$，典范同态 $\Gamma(V,\F)\to\Gamma(f^{-1}(V),f^*\F)$ 立即定义复形同态（G，II，5.1）
\[
C^\bullet(\mathfrak U,\F)\longrightarrow C^\bullet(\mathfrak U',f^*\F),
\]
故得到典范同态
\begin{equation}\tag{12.1.4.1}
\theta'_p:\mathrm H^p(\mathfrak U,\F)\longrightarrow\mathrm H^p(\mathfrak U',f^*\F).
\end{equation}
并且有交换图
\begin{equation}\tag{12.1.4.2}
\begin{tikzcd}
\mathrm H^p(\mathfrak U,\F)\arrow[r,"\theta'_p"]\arrow[d]&\mathrm H^p(\mathfrak U',f^*\F)\arrow[d]\\
\mathrm H^p(X,\F)\arrow[r,"\theta_p"']&\mathrm H^p(X',f^*\F),
\end{tikzcd}
\end{equation}
竖直箭头为（G，II，5.2）的典范同态。

为证明交换性，考虑 $\F$ 相对于 $\mathfrak U$ 的交错上链层复形 $\C^\bullet(\mathfrak U,\F)$，满足
\[
\Gamma(X,\C^\bullet(\mathfrak U,\F))=C^\bullet(\mathfrak U,\F)
\]
（G，II，5.2）。典范同态 $\Gamma(V,\F)\to\Gamma(f^{-1}(V),\psi^*\F)$ 定义一个 $\psi$-态射
\[
\C^\bullet(\mathfrak U,\F)\longrightarrow\C^\bullet(\mathfrak U',\psi^*\F).
\]
采用上述记号，有交换图
\[
\begin{tikzcd}[column sep=small]
\Gamma(X,\C^\bullet(\mathfrak U,\F))\arrow[r]\arrow[d]&\Gamma(X',\C^\bullet(\mathfrak U',\psi^*\F))\arrow[d]&{}\\
\Gamma(X,\LL^\bullet)\arrow[r]&\Gamma(X',\psi^*\LL^\bullet)\arrow[r]&\Gamma(X',\LL'^\bullet).
\end{tikzcd}
\]
在上同调上得到交换图
\[
\begin{tikzcd}
\mathrm H^p(\mathfrak U,\F)\arrow[r]\arrow[d]&\mathrm H^p(\mathfrak U',\psi^*\F)\arrow[d]\\
\mathrm H^p(X,\F)\arrow[r]&\mathrm H^p(X',\psi^*\F),
\end{tikzcd}
\]
其中竖直箭头仍为（G，II，5.2）的典范同态。将它与交换图
\[
\begin{tikzcd}
\mathrm H^p(\mathfrak U',\psi^*\F)\arrow[r]\arrow[d]&\mathrm H^p(\mathfrak U',f^*\F)\arrow[d]\\
\mathrm H^p(X',\psi^*\F)\arrow[r]&\mathrm H^p(X',f^*\F)
\end{tikzcd}
\]
组合即可；后者由同态 $\psi^*\F\to f^*\F$ 及（G，II，5.2）的典范同态的函子性得出。

注意，若 $A=\Gamma(X,\OO_X)$、$A'=\Gamma(X',\OO_{X'})$，则同态（12.1.4.3）是与环 $A,A'$ 相应的模的双同态。\transnote{此处原文引为（12.1.4.3），但本条列出的编号只有（12.1.4.1）、（12.1.4.2）；所指同态依上下文应为（12.1.4.1）。}对两个态射的复合，（12.1.4.1）也有类似于（12.1.3.5）的传递性。最后，在上述定义中，也可以不用内射解消 $\LL^\bullet$，而从满足对每个 $i$ 及 $p>0$ 都有 $\mathrm H^p(X,\LL^i)=0$ 的解消出发（G，II，4.7.1）。

\egapar{12.1.5}
设 $\F,\G,\K$ 为三个 $\OO_X$-模，并考虑 $\OO_X$-同态 $u:\F\otimes_{\OO_X}\G\to\K$。由杯积（12.1.2.1），它在上同调上给出同态
\begin{equation}\tag{12.1.5.1}
\mathrm H^p(X,\F)\otimes_A\mathrm H^q(X,\G)
\longrightarrow\mathrm H^{p+q}(X,\K).
\end{equation}
将证明，在（12.1.3）的假设及记号下，有交换图
\begin{equation}\tag{12.1.5.2}
\begin{tikzcd}[column sep=small]
\mathrm H^p(X,\F)\otimes_A\mathrm H^q(X,\G)\arrow[r]\arrow[d]&\mathrm H^{p+q}(X,\K)\arrow[d]\\
\mathrm H^p(X',f^*\F)\otimes_{A'}\mathrm H^q(X',f^*\G)\arrow[r]&\mathrm H^{p+q}(X',f^*\K),
\end{tikzcd}
\end{equation}
其中竖直箭头来自典范同态（12.1.3.5）。

为此，回顾（12.1.5.1）的构造：取 $\F,\G,\K$ 的典范解消（G，II，4.3）$\LL^\bullet,\M^\bullet,\N^\bullet$（它们由 $\OO_X$-模组成）；与 $u$ 相应的 $\OO_X$-模复形的线性映射 $\LL^\bullet\otimes_{\OO_X}\M^\bullet\to\N^\bullet$ 给出 $A$-模复形同态
\[
\Gamma(X,\LL^\bullet)\otimes_A\Gamma(X,\M^\bullet)\longrightarrow\Gamma(X,\N^\bullet).
\]
再经由上同调同态
\[
\mathrm H^p(\Gamma(X,\LL^\bullet))\otimes_A\mathrm H^q(\Gamma(X,\M^\bullet))
\longrightarrow\mathrm H^{p+q}(\Gamma(X,\N^\bullet))
\]
即可（G，II，6.6）。显然有交换图
\begin{equation}\tag{12.1.5.3}
\begin{tikzcd}[column sep=small]
\Gamma(X,\LL^\bullet)\otimes_A\Gamma(X,\M^\bullet)\arrow[r]\arrow[d]&\Gamma(X,\N^\bullet)\arrow[d]\\
\Gamma(X',\psi^*\LL^\bullet)\otimes_{\Gamma(X',\psi^*\OO_X)}\Gamma(X',\psi^*\M^\bullet)\arrow[r]&\Gamma(X',\psi^*\N^\bullet).
\end{tikzcd}
\end{equation}
它在上同调上给出交换图
\begin{equation}\tag{12.1.5.4}
\adjustbox{max width=\dimexpr\linewidth-5em\relax}{\begin{tikzcd}[column sep=small]
\mathrm H^p(X,\F)\otimes_A\mathrm H^q(X,\G)\arrow[r]\arrow[d]&\mathrm H^{p+q}(X,\K)\arrow[d]\\
\mathrm H^p(\Gamma(X',\psi^*\LL^\bullet))\otimes_{\Gamma(X',\psi^*\OO_X)}\mathrm H^q(\Gamma(X',\psi^*\M^\bullet))\arrow[r]&\mathrm H^{p+q}(\Gamma(X',\psi^*\N^\bullet)).
\end{tikzcd}}
\end{equation}
由于 $\psi^*\LL^\bullet,\psi^*\M^\bullet,\psi^*\N^\bullet$ 分别是 $\psi^*\F,\psi^*\G,\psi^*\K$ 的解消，故有交换图（G，II，6.6.1）
\begin{equation}\tag{12.1.5.5}
\adjustbox{max width=\dimexpr\linewidth-5em\relax}{\begin{tikzcd}[column sep=small]
\mathrm H^p(\Gamma(X',\psi^*\LL^\bullet))\otimes_{\Gamma(X',\psi^*\OO_X)}\mathrm H^q(\Gamma(X',\psi^*\M^\bullet))\arrow[r]\arrow[d]&\mathrm H^{p+q}(\Gamma(X',\psi^*\N^\bullet))\arrow[d]\\
\mathrm H^p(X',\psi^*\F)\otimes_{\Gamma(X',\psi^*\OO_X)}\mathrm H^q(X',\psi^*\G)\arrow[r]&\mathrm H^{p+q}(X',\psi^*\K).
\end{tikzcd}}
\end{equation}
最后，由函子性，有交换图
\begin{equation}\tag{12.1.5.6}
\begin{tikzcd}[column sep=small]
\mathrm H^p(X',\psi^*\F)\otimes_{\Gamma(X',\psi^*\OO_X)}\mathrm H^q(X',\psi^*\G)\arrow[r]\arrow[d]&\mathrm H^{p+q}(X',\psi^*\K)\arrow[d]\\
\mathrm H^p(X',f^*\F)\otimes_{A'}\mathrm H^q(X',f^*\G)\arrow[r]&\mathrm H^{p+q}(X',f^*\K).
\end{tikzcd}
\end{equation}
组合（12.1.5.4）、（12.1.5.5）、（12.1.5.6）三个图，即得到所需交换图（12.1.5.2）。

\egastatement{12.1.6}{注}
采用（12.1.3）的记号，设有交换图
\begin{equation}\tag{12.1.6.1}
\begin{tikzcd}
0\arrow[r]&\F\arrow[r,"r"]\arrow[d,"u"]&\G\arrow[r,"s"]\arrow[d,"v"]&\K\arrow[r]\arrow[d,"w"]&0\\
0\arrow[r]&\F'\arrow[r,"r'"']&\G'\arrow[r,"s'"']&\K'\arrow[r]&0,
\end{tikzcd}
\end{equation}
其中 $r,s$ 为 $\OO_X$-模同态，$r',s'$ 为 $\OO_{X'}$-模同态，$u,v,w$ 为 $f$-态射，各行正合。于是得到交换图
\begin{equation}\tag{12.1.6.2}
\begin{tikzcd}[column sep=tiny]
\cdots\arrow[r]&\mathrm H^p(X,\F)\arrow[r]\arrow[d,"u_p"]&\mathrm H^p(X,\G)\arrow[r]\arrow[d,"v_p"]&\mathrm H^p(X,\K)\arrow[r,"\partial"]\arrow[d,"w_p"]&\mathrm H^{p+1}(X,\F)\arrow[r]\arrow[d,"u_{p+1}"]&\cdots\\
\cdots\arrow[r]&\mathrm H^p(X',\F')\arrow[r]&\mathrm H^p(X',\G')\arrow[r]&\mathrm H^p(X',\K')\arrow[r]&\mathrm H^{p+1}(X',\F')\arrow[r]&\cdots.
\end{tikzcd}
\end{equation}
事实上，（12.1.6.1）可分解为
\[
\begin{tikzcd}[column sep=small]
0\arrow[r]&\F\arrow[r]\arrow[d]&\G\arrow[r]\arrow[d]&\K\arrow[r]\arrow[d]&0\\
0\arrow[r]&\psi^*\F\arrow[r]\arrow[d]&\psi^*\G\arrow[r]\arrow[d]&\psi^*\K\arrow[r]\arrow[d]&0\\
0\arrow[r]&\F'\arrow[r]&\G'\arrow[r]&\K'\arrow[r]&0,
\end{tikzcd}
\]
其中中间一行正合（$\OI$，3.7.2）。只需利用（12.1.3.2）是 $\partial$-函子同态，以及各 $\mathrm H^p(X',\F')$ 关于 $\F'$ 构成 $\partial$-函子这两个事实即可。

\egapar{12.1.7}
采用（12.1.3）的假设及记号，现在考虑 $\F=f_*\F'=\psi_*\F'$ 的情形。将证明，（12.1.3）中定义的双同态
\begin{equation}\tag{12.1.7.1}
\mathrm H^p(X,f_*\F')\longrightarrow\mathrm H^p(X',\F')
\end{equation}
可作为复合函子 $\F'\mapsto\Gamma(X,\psi_*\F')$ 的谱序列的边缘同态得到\transnote{本段开始处原著把直像层的截面所取的空间写成 $X'$；直像 $\psi_*\F'$ 是 $X$ 上的层，故此处明确写作 $\Gamma(X,\psi_*\F')$。}（允许在 $\mathrm H^p(X',\F')$ 上复合一个自同构）（T，2.4）。

按（12.1.4）的描述，可取 $\psi_*\F'$ 与 $\F'$ 的内射解消 $\LL^\bullet,\LL'^\bullet$，然后取位于典范同态 $\psi^*(\psi_*\F')\to\F'$ 之上的复形同态 $v:\psi^*\LL^\bullet\to\LL'^\bullet$。注意 $\Gamma(X',\LL'^\bullet)=\Gamma(X,\psi_*\LL'^\bullet)$，并且复合同态
\[
\Gamma(X,\LL^\bullet)\to\Gamma(X',\psi^*\LL^\bullet)\to\Gamma(X',\LL'^\bullet)
\]
就是
\begin{equation}\tag{12.1.7.2}
\Gamma(v^\flat):\Gamma(X,\LL^\bullet)\longrightarrow\Gamma(X,\psi_*\LL'^\bullet)
\end{equation}
（$\OI$，3.7.1）；（12.1.7.1）由它取上同调而得。

另一方面，取复形 $\psi_*\LL'^\bullet$ 在 $X$ 上阿贝尔群层范畴中的内射 Cartan--Eilenberg 解消 $\M^{\bullet\bullet}=(\M^{ij})$。复合函子的谱序列就是双复形 $\Gamma(X,\M^{\bullet\bullet})$ 的谱序列；由于 $i<0$ 或 $j<0$ 时 $\M^{ij}=0$，它们双正则。其中第一个谱序列退化：因为各 $\psi_*\LL'^i$ 是松弛层（G，II，3.1.1），故
\[
\mathrm H^q(\Gamma(X,\M^{i\bullet}))
\simeq\mathrm H^q(X,\psi_*\LL'^i)=0\qquad(q>0)
\]
（G，II，4.4.3）。于是有双射的边缘同态（11.1.6）
\begin{equation}\tag{12.1.7.3}
{}'E_2^{p0}=\mathrm H_{\mathrm I}^p(\mathrm H_{\mathrm{II}}^0(\Gamma(X,\M^{\bullet\bullet})))
\longrightarrow\mathrm H^p(\Gamma(X,\M^{\bullet\bullet})).
\end{equation}
已知（11.3.4）它由增广
\begin{equation}\tag{12.1.7.4}
\Gamma(X,\psi_*\LL'^\bullet)\longrightarrow\Gamma(X,\M^{\bullet\bullet})
\end{equation}
取上同调得到；后者本身来自增广 $\eta:\psi_*\LL'^\bullet\to\M^{\bullet\bullet}$。

对于第二个谱序列，则有边缘同态
\begin{equation}\tag{12.1.7.5}
{}''E_2^{p0}=\mathrm H_{\mathrm{II}}^p(\mathrm H_{\mathrm I}^0(\Gamma(X,\M^{\bullet\bullet})))
\longrightarrow\mathrm H^p(\Gamma(X,\M^{\bullet\bullet})),
\end{equation}
由复形同态 $\mathrm Z_{\mathrm I}^0(\Gamma(X,\M^{\bullet\bullet}))\to\Gamma(X,\M^{\bullet\bullet})$ 取上同调而得（11.3.4）。由于 $\psi_*$ 左正合，列
\[
0\longrightarrow\psi_*\F'\longrightarrow\psi_*\LL'^0\longrightarrow\psi_*\LL'^1
\]
正合；依 Cartan--Eilenberg 解消的定义（11.4.2），可取 $\mathrm B_{\mathrm I}^0(\M^{\bullet\bullet})=0$、$\mathrm Z_{\mathrm I}^0(\M^{\bullet\bullet})=\LL^\bullet$。由于图
\[
\begin{tikzcd}
\LL^0\arrow[r,"i_0"]&\M^{00}\\
\psi_*\F'\arrow[u,"\varepsilon"]\arrow[r,"\varepsilon'"']&\psi_*\LL'^0\arrow[u,"\eta_0"']
\end{tikzcd}
\]
交换，复形嵌入 $i:\LL^\bullet\to\M^{\bullet\bullet}$ 与增广 $\varepsilon,\varepsilon'$ 相容。因而有两个与增广相容的复形同态：
\[
\begin{tikzcd}
\LL^\bullet\arrow[rr,"i"]\arrow[dr,"v^\flat"']&&\M^{\bullet\bullet}\\
&\psi_*\LL'^\bullet\arrow[ur,"\eta"']&
\end{tikzcd}
\]
由于 $\LL^\bullet$ 是内射解消而 $\M^{\bullet\bullet}$ 由内射层组成，由（M，V，1.1 a）），这两个同态同伦。因此相应的两个同态 $\Gamma(X,\LL^\bullet)\to\Gamma(X,\M^{\bullet\bullet})$ 也同伦，在上同调上给出同一同态。换言之，已证明边缘同态（12.1.7.5）
\[
\mathrm H^p(X,\psi_*\F')\to\mathrm H^p(\Gamma(X,\M^{\bullet\bullet}))
\]
\begin{sloppypar}
等于（12.1.7.1）与（12.1.7.3）之复合；而后者即 $\mathrm H^p(X',\F')\to\mathrm H^p(\Gamma(X,\M^{\bullet\bullet}))$，已经证明是同构。这就得到所述断言。\par
\end{sloppypar}


% ---- END c0_s12a.tex ----
% ---- BEGIN c0_s12b.tex ----
\subsection{高阶直像}
\src{57--60}
\egapar{12.2.1}
设 $(X,\OO_X),(Y,\OO_Y)$ 为赋环空间，$f=(\psi,\theta):X\to Y$ 为态射；它定义直像函子 $f_*:C(X)\to C(Y)$，其底层阿贝尔群层函子与 $\psi_*$ 在 $C(X)$ 上的限制相同。$\psi_*$ 加性且左正合；又由于每个阿贝尔群层都同构于某个内射阿贝尔群层的子层，可以定义 $\psi_*$ 的右导出函子 $\F\mapsto\mathrm R^p\psi_*\F$。这些是 $Y$ 上的阿贝尔群层，且各 $\mathrm R^p\psi_*$ 构成泛上同调函子（T，2.3）。

此外，$\mathrm R^p\psi_*\F$ 是预层 $V\mapsto\mathrm H^p(f^{-1}(V),\F)$ 的相伴层（T，3.7.2）。若 $\F$ 是 $\OO_X$-模，则 $\mathrm H^p(f^{-1}(V),\F)$ 自然具有 $\Gamma(f^{-1}(V),\OO_X)$-模结构，即 $\Gamma(V,\psi_*\OO_X)$-模结构；再利用 $\theta:\OO_Y\to\psi_*\OO_X$，得到 $\Gamma(V,\OO_Y)$-模结构。对于这些结构，从开集 $V$ 限制到 $V'\subset V$ 显然给出双同态，所以它们在各 $\mathrm R^p\psi_*\F$ 上定义 $\OO_Y$-模结构。我们以 $\mathrm R^pf_*\F$ 表示这个 $\OO_Y$-模，从而得到加性函子 $\mathrm R^pf_*:C(X)\to C(Y)$。

这些函子构成 $\partial$-函子。事实上，若 $0\to\F'\to\F\to\F''\to0$ 是 $\OO_X$-模正合列，由上述高阶直像及模结构的描述立即可知，同态 $\partial:\mathrm R^p\psi_*\F''\to\mathrm R^{p+1}\psi_*\F'$ 是 $\OO_Y$-模同态。最后，各 $\mathrm R^pf_*$ 就是 $f_*$ 的右导出函子：每个 $\OO_X$-模都有由 $\OO_X$-模组成的内射解消，而这样的解消由松弛阿贝尔群层组成（G，II，7.1），所以能用来计算 $\mathrm R^p\psi_*\F$，因为对每个松弛层 $\G$ 及 $n\geq1$，都有 $\mathrm R^n\psi_*\G=0$（T，2.4.1，注3，以及命题3.3.2的推论）。因此，各 $\mathrm R^pf_*$ 构成从 $C(X)$ 到 $C(Y)$ 的泛上同调函子（T，2.3）。

\egapar{12.2.2}
设 $\F,\G$ 为两个 $\OO_X$-模。采用（12.2.1）的记号，对 $Y$ 的每个开集 $V$，有杯积同态（12.1.2.1）
\[
\mathrm H^p(f^{-1}(V),\F)\otimes_{\Gamma(f^{-1}(V),\OO_X)}
\mathrm H^q(f^{-1}(V),\G)
\longrightarrow\mathrm H^{p+q}(f^{-1}(V),\F\otimes_{\OO_X}\G).
\]
由杯积的定义（G，II，6.6），这些同态与从 $V$ 到开子空间 $V'$ 的限制交换。另一方面，$\theta$ 给出环同态
\[
\Gamma(V,\OO_Y)\to\Gamma(V,\psi_*\OO_X)=\Gamma(f^{-1}(V),\OO_X),
\]
故给出典范张量积同态
\[
\begin{aligned}
&\mathrm H^p(f^{-1}(V),\F)\otimes_{\Gamma(V,\OO_Y)}\mathrm H^q(f^{-1}(V),\G)\\
&\hspace{2em}\longrightarrow
\mathrm H^p(f^{-1}(V),\F)\otimes_{\Gamma(f^{-1}(V),\OO_X)}\mathrm H^q(f^{-1}(V),\G),
\end{aligned}
\]
它同样与限制相容。复合后得到 $\Gamma(V,\OO_Y)$-模同态；在相伴层上，它定义关于 $\F,\G$ 的典范函子同态
\begin{equation}\tag{12.2.2.1}
\mathrm R^pf_*\F\otimes_{\OO_Y}\mathrm R^qf_*\G
\longrightarrow\mathrm R^{p+q}f_*(\F\otimes_{\OO_X}\G).
\end{equation}
当 $p=q=0$ 时，此同态即为（$\OI$，4.2.2.1）。

\egastatement{12.2.3}{命题}
对每个 $\OO_X$-模 $\F$ 及每个有限秩局部自由 $\OO_Y$-模 $\LL$，有典范的函子同构
\begin{equation}\tag{12.2.3.1}
\mathrm R^pf_*\F\otimes_{\OO_Y}\LL
\iso\mathrm R^pf_*(\F\otimes_{\OO_X}f^*\LL).
\end{equation}
\begin{proof}
（12.2.2.1）的一个特例为
\begin{equation}\tag{12.2.3.2}
\mathrm R^pf_*\F\otimes_{\OO_Y}f_*(f^*\LL)
\longrightarrow\mathrm R^pf_*(\F\otimes_{\OO_X}f^*\LL).
\end{equation}
将它与由典范同态（$\OI$，4.4.3.2）导出的、从（12.2.3.1）左端到（12.2.3.2）左端的同态复合，即得（12.2.3.1）。为验证局部自由情形下它为同构，由于问题在 $Y$ 上局部，且所涉及函子对 $\LL$ 加性，可立即归结到 $\LL=\OO_Y$。由（12.2.2.1）的定义，只需验证相应的预层同态为双射；而由 $f^*\OO_Y=\OO_X$，这一点显然成立。
\end{proof}

\egapar{12.2.4}
设 $(Z,\OO_Z)$ 为第三个赋环空间，$g:Y\to Z$ 为赋环空间态射。已知（G，II，7.1及3.1.1），对于每个内射 $\OO_X$-模 $\G$，$f_*\G$ 是松弛阿贝尔群层，故 $p>0$ 时 $\mathrm R^pg_*(f_*\G)=0$（12.2.1）。因此可对复合函子 $g_*f_*$ 应用 Leray 复合函子谱序列（T，2.4.1）：令 $h=g\circ f$，则存在以函子 $\mathrm R^\bullet h_*$ 为终项的双正则谱序列，其 $E_2$ 项为
\begin{equation}\tag{12.2.4.1}
E_2^{pq}=\mathrm R^pg_*(\mathrm R^qf_*\F).
\end{equation}

\egapar{12.2.5}
在（12.2.4）的条件下，将直接定义典范 $\OO_Z$-模同态
\begin{equation}\tag{12.2.5.1}
\mathrm R^ng_*(f_*\F)\longrightarrow\mathrm R^nh_*\F,
\end{equation}
\begin{equation}\tag{12.2.5.2}
\mathrm R^nh_*\F\longrightarrow g_*(\mathrm R^nf_*\F).
\end{equation}
它们可以与 Leray 谱序列的边缘同态相识别（见12.1.7）。只需在高阶直像所对应的预层上进行构造（12.2.1）。对 $Z$ 的任意开集 $W$，及其在 $Y$ 中的逆像 $g^{-1}(W)$，有典范双同态
\begin{equation}\tag{12.2.5.3}
\mathrm H^n(g^{-1}(W),f_*\F)
\longrightarrow\mathrm H^n(f^{-1}(g^{-1}(W)),f^*(f_*\F)),
\end{equation}
相应的环为 $\Gamma(g^{-1}(W),\OO_Y)$ 与 $\Gamma(h^{-1}(W),\OO_X)$。另一方面，典范同态（$\OI$，4.4.3.3）由函子性给出
\begin{equation}\tag{12.2.5.4}
\mathrm H^n(h^{-1}(W),f^*(f_*\F))\longrightarrow\mathrm H^n(h^{-1}(W),\F),
\end{equation}
这是 $\Gamma(h^{-1}(W),\OO_X)$-模同态。考虑环同态 $\Gamma(W,\OO_Z)\to\Gamma(h^{-1}(W),\OO_X)$，可知复合（12.2.5.3）与（12.2.5.4）得到预层同态，进而得到（12.2.5.1）。

（12.2.5.2）的定义更简单。按定义，$\mathrm R^nh_*\F$ 对应于预层 $W\mapsto\mathrm H^n(f^{-1}(g^{-1}(W)),\F)$，而 $\mathrm R^nf_*\F$ 对应于 $V\mapsto\mathrm H^n(f^{-1}(V),\F)$；故有典范同态
\[
\mathrm H^n(f^{-1}(g^{-1}(W)),\F)
\longrightarrow\Gamma(g^{-1}(W),\mathrm R^nf_*\F).
\]
它们显然定义预层同态，后者又定义（12.2.5.2）。

\egapar{12.2.6}
在（12.2.4）的假设下，设 $\F,\G,\K$ 为三个 $\OO_X$-模，$u:\F\otimes_{\OO_X}\G\to\K$ 为 $\OO_X$-同态。则有交换图
\begin{equation}\tag{12.2.6.1}
\begin{tikzcd}[column sep=small]
\mathrm R^pg_*(f_*\F)\otimes_{\OO_Z}\mathrm R^qg_*(f_*\G)\arrow[r]\arrow[d]&\mathrm R^{p+q}g_*(f_*\K)\arrow[d]\\
\mathrm R^ph_*\F\otimes_{\OO_Z}\mathrm R^qh_*\G\arrow[r]&\mathrm R^{p+q}h_*\K
\end{tikzcd}
\end{equation}
及
\begin{equation}\tag{12.2.6.2}
\begin{tikzcd}[column sep=small]
\mathrm R^ph_*\F\otimes_{\OO_Z}\mathrm R^qh_*\G\arrow[r]\arrow[d]&\mathrm R^{p+q}h_*\K\arrow[d]\\
g_*(\mathrm R^pf_*\F)\otimes_{\OO_Z}g_*(\mathrm R^qf_*\G)\arrow[r]&g_*(\mathrm R^{p+q}f_*\K).
\end{tikzcd}
\end{equation}
水平箭头来自（12.2.2.1）（最后一个还结合（$\OI$，4.2.2.1）），两个图的竖直箭头分别来自（12.2.5.1）、（12.2.5.2）。只需对相应的预层同态验证。回到（12.2.2）及（12.2.5）的定义，（12.2.6.1）立即归结为交换图（12.1.5.2）；（12.2.6.2）的验证更为简单。

\subsection{层的 \texorpdfstring{$\Ext$}{Ext} 函子的补充知识}
\src{60--62}
\egapar{12.3.1}
考虑赋环空间 $(X,\OO_X)$。这里不再重复双函子 $\Ext^p_{\OO_X}(X;\F,\G)$ 与 $\shExt^p_{\OO_X}(\F,\G)$ 的定义和主要性质；前者从 $\OO_X$-模范畴取值于 $\Gamma(X,\OO_X)$-模范畴，后者取值于 $\OO_X$-模范畴本身。也不再重复将它们联系起来的双正则谱序列 $E(\F,\G)$ 的定义（T，4.2及G，II，7.3）。

\egapar{12.3.2}
与（M，XIV，1）中一样，可以定义 $\OO_X$-模 $\F$ 被 $\OO_X$-模 $\G$ 扩张的概念，以及等价扩张类之间的合成律：对于模所作的论证，显然适用于任意阿贝尔范畴。（M，XIV，1.1）的第二个证明只用到嵌入内射对象的存在性，故也适用于 $\OO_X$-模范畴。因此 $\Ext^p_{\OO_X}(X;\F,\G)$ 典范地识别为 $\F$ 被 $\G$ 扩张的等价类所成阿贝尔群。

\egastatement{12.3.3}{命题}
设 $(X,\OO_X)$ 为赋环空间，且环层 $\OO_X$ 凝聚。则对每对凝聚 $\OO_X$-模 $\F,\G$ 及每个 $p\geq0$，$\shExt^p_{\OO_X}(\F,\G)$ 都是凝聚 $\OO_X$-模。
\begin{proof}
注意各 $\shExt^p_{\OO_X}(\F,\G)$ 关于 $\F$ 构成反变上同调函子。由于 $\F$ 凝聚，对每个 $p$ 及每点 $x\in X$，存在 $x$ 的开邻域 $U$ 及 $\OO_X|U$-模正合列
\[
0\longrightarrow\R\longrightarrow\LL_{p-1}\longrightarrow\cdots
\longrightarrow\LL_0\longrightarrow\F|U\longrightarrow0,
\]
其中每个 $\LL_i$（$0\leq i\leq p-1$）同构于某个 $\OO_X^{n_i}|U$，且 $\R$ 凝聚。这由（$\OI$，5.3.2、5.3.4）对 $p$ 归纳得出，使用了 $\OO_X$ 凝聚的假设。

对于 $p\geq1$，当 $\LL|U$ 同构于 $\OO_X^n|U$ 时，\transnote{原文紧接的消失式用直立的 $\Ext$，但本证明讨论的是层 $\shExt$；这里依上下文使用层的记号，不能将其解释为任意开集上全局 $\Ext$ 的消失。}$\shExt^p_{\OO_X|U}(\LL|U,\G|U)=0$（T，4.2.3）。因此（M，V，7.2）的论证可用于反变上同调函子，从而得到正合列
\[
\begin{gathered}
\shHom_{\OO_X|U}(\LL_{p-1},\G|U)
\longrightarrow\shHom_{\OO_X|U}(\R,\G|U)\\
\longrightarrow\shExt^p_{\OO_X|U}(\F|U,\G|U)\longrightarrow0.
\end{gathered}
\]
其前两项凝聚（$\OI$，5.3.5），故第三项也凝聚（$\OI$，5.3.4）。
\end{proof}

\egastatement{12.3.4}{命题}
设 $f:X\to Y$ 为赋环空间的平坦态射，$\F,\G$ 为两个 $\OO_Y$-模。
\begin{enumerate}[label=(\roman*)]
\item 存在上同调双函子的同态
\begin{equation}\tag{12.3.4.1}
f^*(\shExt^p_{\OO_Y}(\F,\G))
\longrightarrow\shExt^p_{\OO_X}(f^*\F,f^*\G),
\end{equation}
在零次上等于典范同态（$\OI$，4.4.6）。
\item 存在典范谱序列态射
\begin{equation}\tag{12.3.4.2}
E(\F,\G)\longrightarrow E(f^*\F,f^*\G),
\end{equation}
其在 $E_2$ 项上为由（12.3.4.1）及（12.1.3.1）导出的同态
\begin{equation}\tag{12.3.4.3}
\mathrm H^p(Y,\shExt^q_{\OO_Y}(\F,\G))
\longrightarrow\mathrm H^p(X,\shExt^q_{\OO_X}(f^*\F,f^*\G)).
\end{equation}
\end{enumerate}
\begin{proof}
（i）由于 $f^*$ 在 $\OO_Y$-模范畴中正合（$\OI$，6.7.2），函子 $\G\mapsto f^*(\shHom_{\OO_Y}(\F,\G))$ 与 $\G\mapsto\shHom_{\OO_X}(f^*\F,f^*\G)$ 左正合；从（$\OI$，4.4.6）典范地得到它们的导出函子之间的同态。取 $\G$ 的内射解消 $\LL^\bullet=(\LL^i)$ 来计算，得到层复形上同调的态射
\[
\HH^p(f^*(\shHom_{\OO_Y}(\F,\LL^\bullet)))
\longrightarrow\HH^p(\shHom_{\OO_X}(f^*\F,f^*\LL^\bullet)).
\]
由 $f^*$ 的正合性，左端按定义等于
\[
f^*(\HH^p(\shHom_{\OO_Y}(\F,\LL^\bullet)))
=f^*(\shExt^p_{\OO_Y}(\F,\G)).
\]
另一方面，$f^*\LL^\bullet$ 是 $f^*\G$ 的解消。若 $\LL'^\bullet=(\LL'^i)$ 为 $f^*\G$ 在 $\OO_X$-模范畴中的内射解消，则有在同伦意义下确定的复形同态 $f^*\LL^\bullet\to\LL'^\bullet$，因此在上同调上给出确定的同态。与上面构造的同态复合，即得到（12.3.4.1）。

（ii）采用上述记号，有 $\OO_X$-模层复形同态
\[
f^*(\shHom_{\OO_Y}(\F,\LL^\bullet))
\longrightarrow\shHom_{\OO_X}(f^*\F,\LL'^\bullet).
\]
设 $\M^{\bullet\bullet}$ 为复形 $\shHom_{\OO_Y}(\F,\LL^\bullet)$ 在 $\OO_Y$-模范畴中的内射 Cartan--Eilenberg 解消。由 $f^*$ 正合，$f^*\M^{\bullet\bullet}$ 是复形 $f^*(\shHom_{\OO_Y}(\F,\LL^\bullet))$ 的 Cartan--Eilenberg 解消。若 $\M'^{\bullet\bullet}$ 是 $\shHom_{\OO_X}(f^*\F,\LL'^\bullet)$ 的内射 Cartan--Eilenberg 解消，则有与上述同态相容、在同伦意义下确定的同态 $f^*\M^{\bullet\bullet}\to\M'^{\bullet\bullet}$（11.4.2），即层双复形的 $f$-态射 $\M^{\bullet\bullet}\to\M'^{\bullet\bullet}$。

\begin{sloppypar}
由此得到在同伦意义下确定的模双复形双同态 $\Gamma(Y,\M^{\bullet\bullet})\to\Gamma(X,\M'^{\bullet\bullet})$，因而得到确定的谱序列态射（11.3.2）；它就是所需的（12.3.4.2）。（12.3.4.3）的描述立即由定义得出。
\end{sloppypar}
\end{proof}

\egastatement{12.3.5}{命题}
在（12.3.4）的假设下，再假设环层 $\OO_Y$ 凝聚；则对每个凝聚 $\OO_Y$-模 $\F$，典范同态（12.3.4.1）均为双射。
\begin{proof}
问题在 $Y$ 上局部，可设存在正合列 $0\to\R\to\OO_Y^n\to\F\to0$，此时 $\R$ 也是凝聚 $\OO_Y$-模（$\OI$，5.3.4）。对 $p$ 归纳证明同态
\[
f^*(\shExt^p_{\OO_Y}(\F,\G))\longrightarrow\shExt^p_{\OO_X}(f^*\F,f^*\G)
\]
为双射；$p=0$ 时由（$\OI$，6.7.6.1）成立。有交换图
\[
\adjustbox{max width=\linewidth}{\begin{tikzcd}[column sep=tiny]
f^*\shExt^{p-1}_{\OO_Y}(\OO_Y^n,\G)\arrow[r]\arrow[d]&f^*\shExt^{p-1}_{\OO_Y}(\R,\G)\arrow[r]\arrow[d]&f^*\shExt^p_{\OO_Y}(\F,\G)\arrow[r]\arrow[d]&f^*\shExt^p_{\OO_Y}(\OO_Y^n,\G)\arrow[d]\\
\shExt^{p-1}_{\OO_X}(\OO_X^n,f^*\G)\arrow[r]&\shExt^{p-1}_{\OO_X}(f^*\R,f^*\G)\arrow[r]&\shExt^p_{\OO_X}(f^*\F,f^*\G)\arrow[r]&\shExt^p_{\OO_X}(\OO_X^n,f^*\G).
\end{tikzcd}}
\]
这里使用 $f^*\OO_Y=\OO_X$；由于 $f^*$ 正合，两行正合。又对 $p>0$，有 $\shExt^p_{\OO_Y}(\OO_Y,\G)=0$，同样 $\shExt^p_{\OO_X}(\OO_X,f^*\G)=0$（T，4.2.3）。由归纳假设，前两个竖直箭头为同构，最右两项为零，因而第三个竖直箭头也为同构。
\end{proof}

\subsection{直像函子的超上同调}
\src{62--64}
\egapar{12.4.1}
设 $(X,\OO_X),(Y,\OO_Y)$ 为赋环空间，$f:X\to Y$ 为态射。可以取 $f_*$ 相对于任意 $\OO_X$-模复形 $\K^\bullet=(\K^i)_{i\in\Z}$ 的超上同调（11.4.4），因为在阿贝尔范畴 $C(Y)$ 中滤过归纳极限存在且正合（T，3.1.1）。超上同调 $\OO_Y$-模 $\mathrm R^pf_*(\K^\bullet)$ 也记作 $\HH^p(f,\K^\bullet)$ 或 $\HH^p(\K^\bullet)$。

回顾 $\HH^\bullet(f,\K^\bullet)$ 是 $\OO_Y$-模双复形 $f_*\LL^{\bullet\bullet}$ 的上同调，其中 $\LL^{\bullet\bullet}$ 为 $\K^\bullet$ 在 $\OO_X$-模范畴中的内射 Cartan--Eilenberg 解消。它是两个谱序列 ${}'\E(f,\K^\bullet)$、${}''\E(f,\K^\bullet)$ 的终项；其 $E_2$ 项为
\begin{equation}\tag{12.4.1.1}
{}'\E_2^{pq}=\HH^p(\HH^q(f,\K^\bullet)),
\end{equation}
\begin{equation}\tag{12.4.1.2}
{}''\E_2^{pq}=\HH^p(f,\HH^q(\K^\bullet))
=\mathrm R^pf_*(\HH^q(\K^\bullet)).
\end{equation}
这里采用函子作用于复形的一般记号 $T(A^\bullet)$（11.2.1），并对单个 $\OO_X$-模 $\F$，用 $\HH^p(f,\F)$ 表示 $\mathrm R^pf_*\F$。还要回顾，第一个谱序列总正则；若 $\K^\bullet$ 有下界，或存在整数 $m$ 使每个 $\OO_X$-模都有长度 $\leq m$ 的松弛解消，则两个谱序列均双正则（11.4.4）。

\egapar{12.4.2}
同样，以 $\mathrm H^\bullet(X,\K^\bullet)$ 表示截面函子 $\Gamma$ 相对于 $\OO_X$-模复形 $\K^\bullet$ 的超上同调；因此各 $\mathrm H^p(X,\K^\bullet)$ 是 $\Gamma(X,\OO_X)$-模。它也可视为 $\HH^\bullet(f,\K^\bullet)$ 的特例，其中 $f$ 从 $(X,\OO_X)$ 映向底空间只有一点、环为 $\Gamma(X,\OO_X)$ 的赋环空间。对 $X$ 的开集 $V$，简记 $\mathrm H^\bullet(V,\K^\bullet)=\mathrm H^\bullet(V,\K^\bullet|V)$。

\egastatement{12.4.3}{命题}
对每个整数 $p\in\Z$，$\OO_Y$-模 $\HH^p(f,\K^\bullet)$ 典范地同构于 $Y$ 上预层 $U\mapsto\mathrm H^p(f^{-1}(U),\K^\bullet)$ 的相伴层。
\begin{proof}
用（12.4.1）的记号，上同调层 $\HH^p(f_*\LL^{\bullet\bullet})$ 对应于预层
\[
U\longmapsto\mathrm H^p(\Gamma(U,f_*\LL^{\bullet\bullet}))
=\mathrm H^p(\Gamma(f^{-1}(U),\LL^{\bullet\bullet})).
\]
$\LL^{\bullet\bullet}|f^{-1}(U)$ 显然是 $\K^\bullet|f^{-1}(U)$ 的内射 Cartan--Eilenberg 解消（T，3.1.3），所以按定义，右端等于 $\mathrm H^p(f^{-1}(U),\K^\bullet)$。
\end{proof}

\egastatement{12.4.4}{命题}
在以下每一种情形下，超上同调 $\HH^\bullet(f,\K^\bullet)$ 都是关于 $\K^\bullet$ 的上同调函子：
\begin{enumerate}[label=\alph*)]
\item $\K^\bullet$ 在有下界复形范畴中变化；
\item 存在整数 $m$，使每个 $\OO_X$-模都有长度 $\leq m$ 的松弛解消；
\item $X$ 是诺特空间。
\end{enumerate}
\begin{proof}
情形 a）、b）是（11.5.4）的特例。情形 c）由（11.5.2）得出，因为已知此时 $f_*$ 与归纳极限交换（G，II，3.10.1）。
\end{proof}

\egapar{12.4.5}
考虑 $X$ 的开覆盖 $\mathfrak U=(U_\alpha)$。对 $X$ 上每个预层复形 $\K^\bullet=(\K^j)$，考虑双复形 $C^\bullet(\mathfrak U,\K^\bullet)$，其 $(i,j)$ 分量为 $C^i(\mathfrak U,\K^j)$，即覆盖神经在 $\K^j$ 中取值的交错 $i$-上链群（G，II，5.1）。称此双复形的上同调为覆盖 $\mathfrak U$ 在 $\K^\bullet$ 中取系数的\textbf{超上同调}，记为
\[
\mathrm H^\bullet(\mathfrak U,\K^\bullet)
=\mathrm H^\bullet(C^\bullet(\mathfrak U,\K^\bullet)).
\]
覆盖的 Leray 谱序列（T，3.8.1及G，II，5.9.1）推广到超上同调如下。

\egastatement{12.4.6}{命题}
设 $\K^\bullet=(\K^j)$ 为 $\OO_X$-模复形。存在关于 $\K^\bullet$ 的正则谱函子，以超上同调 $\mathrm H^\bullet(X,\K^\bullet)$ 为终项，其 $E_2$ 项为
\begin{equation}\tag{12.4.6.1}
E_2^{pq}=\mathrm H^p(\mathfrak U,h^q(\K^\bullet)),
\end{equation}
其中 $h^q(\K^\bullet)$ 表示 $X$ 上由预层 $V\mapsto\mathrm H^q(V,\K^j)$ 组成的复形。\transnote{原文在这一定义中写 $V\mapsto H^q(V,\K^\bullet)$，易与已定义的超上同调混淆；由随后证明中的 $h^q(\K^j)$ 可知，此处指逐项取第 $q$ 次普通上同调所成的预层复形，本译文将指标 $j$ 明写。}当 $\K^\bullet$ 有下界时，此谱序列双正则。
\begin{proof}
\begin{sloppypar}
取 $\K^\bullet$ 的内射 Cartan--Eilenberg 解消 $\LL^{\bullet\bullet}$，及三重复形 $C^\bullet(\mathfrak U,\LL^{\bullet\bullet})=(C^i(\mathfrak U,\LL^{jk}))$。先将它按次数 $i,j+k$ 视为双复形。由于 $i$ 只取非负值，其第二谱序列正则（11.3.3）且退化：各 $\LL^{jk}$ 为松弛层，故对 $q>0$ 有 $\mathrm H^q(\mathfrak U,\LL^{jk})=0$（G，II，5.2.3）。因此有典范同构（11.1.6）
\end{sloppypar}
\[
\mathrm H^\bullet(C^\bullet(\mathfrak U,\LL^{\bullet\bullet}))
\simeq\mathrm H^\bullet(\Gamma(X,\LL^{\bullet\bullet}))
\]
（G，II，5.2.2），亦即按（12.4.2）有
\[
\mathrm H^\bullet(C^\bullet(\mathfrak U,\LL^{\bullet\bullet}))
\simeq\mathrm H^\bullet(X,\K^\bullet).
\]
再将三重复形按次数 $i+j,k$ 视为双复形。由于 $k$ 只取非负值，其第一谱序列总正则；若 $j<j_0$ 时 $\LL^{jk}=0$，即 $\K^\bullet$ 有下界时，它双正则（11.3.3）。这就是所求谱序列：对每个 $j$，$\LL^{j\bullet}$ 是 $\K^j$ 的内射解消，所以 $\mathrm H^q(C^\bullet(\mathfrak U,\LL^{j\bullet}))$ 正是上链复形 $C^\bullet(\mathfrak U,h^q(\K^j))$。命题得证。
\end{proof}

\egastatement{12.4.7}{推论}
若对覆盖神经的每个单形 $\sigma$ 及每个整数 $i$，均有 $q>0$ 时 $\mathrm H^q(U_\sigma,\K^i)=0$，则有典范同构
\begin{equation}\tag{12.4.7.1}
\mathrm H^\bullet(\mathfrak U,\K^\bullet)\simeq\mathrm H^\bullet(X,\K^\bullet).
\end{equation}
\begin{proof}
假设说明 $q>0$ 时 $C^\bullet(\mathfrak U,h^q(\K^\bullet))=0$，故 $E_2^{pq}=0$。谱序列（12.4.6.1）退化且正则，所以结论由定义（12.4.5）及（11.1.6）得出。
\end{proof}

\egapar{12.4.8}
设 $(X',\OO_{X'})$ 为第二个赋环空间，$f=(\psi,\theta):X'\to X$ 为态射。用（12.1.3）、（12.1.4）的同一方法，对 $\OO_X$-模复形 $\K^\bullet$ 的超上同调定义双同态
\begin{equation}\tag{12.4.8.1}
\mathrm H^p(X,\K^\bullet)\longrightarrow\mathrm H^p(X',f^*\K^\bullet).
\end{equation}
从 $\K^\bullet$ 的内射 Cartan--Eilenberg 解消 $\LL^{\bullet\bullet}$ 出发。由于 $\psi^*$ 正合，$\psi^*\LL^{\bullet\bullet}$ 是 $\psi^*\K^\bullet$ 在 $\psi^*\OO_X$-模范畴中的 Cartan--Eilenberg 解消；于是有态射 $\psi^*\LL^{\bullet\bullet}\to\LL'^{\bullet\bullet}$，其中 $\LL'^{\bullet\bullet}$ 为 $\psi^*\K^\bullet$ 的内射 Cartan--Eilenberg 解消，因而得到上同调态射
\[
\mathrm H^\bullet(X,\K^\bullet)\longrightarrow\mathrm H^\bullet(X',\psi^*\K^\bullet).
\]
与由 $\psi^*\K^\bullet\to f^*\K^\bullet$ 的函子性导出的态射复合，即得（12.4.8.1）。

由（12.4.8.1）、（12.4.3）出发，像（12.2.5）那样论证，对赋环空间态射 $f:X\to Y$、$g:Y\to Z$，可定义 $\OO_X$-模复形 $\K^\bullet$ 的超上同调同态
\begin{equation}\tag{12.4.8.2}
\HH^p(g,f_*\K^\bullet)\longrightarrow\HH^p(g\circ f,\K^\bullet),
\end{equation}
\begin{equation}\tag{12.4.8.3}
\HH^p(g\circ f,\K^\bullet)\longrightarrow g_*\HH^p(f,\K^\bullet).
\end{equation}
定义的细节留给读者。\par

% ---- END c0_s12b.tex ----
% ---- BEGIN c0_s13a.tex ----
\section{同调代数中的射影极限}
\subsection{Mittag--Leffler 条件}
\src{64--65}
\egapar{13.1.1}
设 $C$ 为存在无限积的阿贝尔范畴（T，1.5的公理 AB~3*）。则 $C$ 的任一对象的一族子对象有下确界；每个 $C$-对象射影系统都有射影极限，并且射影极限是关于所考虑射影系统的左正合函子（T，1.8）。

设 $(A_\alpha,f_{\alpha\beta})$ 为 $C$ 的对象射影系统，其指标集 $I$ 向右滤过。令 $A=\plim A_\alpha$，并令 $f_\alpha:A\to A_\alpha$ 为典范态射。对每个 $\alpha$，子对象 $f_{\alpha\beta}(A_\beta)$（$\beta\geq\alpha$）构成滤过递减族；子对象
\[
A'_\alpha=\inf_{\beta\geq\alpha}f_{\alpha\beta}(A_\beta)
\]
称为 $A_\alpha$ 中的“普遍像”子对象。显然 $f_\alpha(A)\subset A'_\alpha$，且 $\alpha\leq\beta$ 时 $f_{\alpha\beta}(A'_\beta)\subset A'_\alpha$；因此 $(A'_\alpha,f_{\alpha\beta}|A'_\beta)$ 是射影系统，且 $A=\plim A'_\alpha$。

\egapar{13.1.2}
给定 $C$ 中的射影系统 $(A_\alpha,f_{\alpha\beta})$，称以下条件为\textbf{Mittag--Leffler 条件}：

\noindent（ML）对每个指标 $\alpha$，存在 $\beta\geq\alpha$，使对每个 $\gamma\geq\beta$ 都有
\[
f_{\alpha\gamma}(A_\gamma)=f_{\alpha\beta}(A_\beta).
\]
显然，若各 $f_{\alpha\beta}$ 均为满态射，则（ML）成立。反过来，若（ML）成立，且 $A'_\alpha$ 为普遍像子对象，则 $\alpha\leq\beta$ 时，$f_{\alpha\beta}$ 在 $A'_\beta$ 上的限制是满态射 $A'_\beta\to A'_\alpha$。事实上，若 $\gamma\geq\beta$ 满足 $\delta\geq\gamma$ 时 $f_{\beta\delta}(A_\delta)=f_{\beta\gamma}(A_\gamma)$，则 $A'_\beta=f_{\beta\gamma}(A_\gamma)$；而这又说明 $\delta\geq\gamma$ 时 $f_{\alpha\delta}(A_\delta)=f_{\alpha\gamma}(A_\gamma)$，故
\[
A'_\alpha=f_{\alpha\gamma}(A_\gamma)=f_{\alpha\beta}(A'_\beta).
\]
还要指出，当各 $A_\alpha$ 在 $C$ 中是阿廷对象，即其任意子对象族均有极小元时，（ML）成立。因为滤过递减子对象族 $(f_{\alpha\beta}(A_\beta))$ 的一个极小元，必为该族的最小子对象。

\egastatement{13.1.3}{注}
例如当 $C$ 是集合范畴时，也可以表述（ML）条件；这时仍可定义 $A_\alpha$ 的“普遍像”子集，而（13.1.1）、（13.1.2）中的有关说明仍成立。

\subsection{阿贝尔群的 Mittag--Leffler 条件}
\src{65--67}
\egastatement{13.2.1}{命题}
设
\[
0\longrightarrow A_\alpha\xrightarrow{u_\alpha}B_\alpha
\xrightarrow{v_\alpha}C_\alpha\longrightarrow0
\]
是阿贝尔群射影系统的正合列，具有同一滤过指标集 $I$。
\begin{enumerate}[label=(\roman*)]
\item 若 $(B_\alpha)$ 满足（ML），则 $(C_\alpha)$ 也满足。
\item 若 $(A_\alpha)$ 与 $(C_\alpha)$ 满足（ML），则 $(B_\alpha)$ 也满足。
\end{enumerate}
\begin{proof}
以 $(f_{\alpha\beta})$、$(g_{\alpha\beta})$、$(h_{\alpha\beta})$ 分别表示定义三个射影系统的同态。

（i）设 $\lambda\geq\beta$ 时 $g_{\alpha\beta}(B_\beta)=g_{\alpha\lambda}(B_\lambda)$。由 $v_\beta,v_\lambda$ 满射，
\[
h_{\alpha\beta}(C_\beta)
=v_\alpha(g_{\alpha\beta}(B_\beta))
=v_\alpha(g_{\alpha\lambda}(B_\lambda))
=h_{\alpha\lambda}(C_\lambda).
\]
（ii）给定 $\alpha$，取 $\beta\geq\alpha$，使 $\lambda\geq\beta$ 时 $f_{\alpha\beta}(A_\beta)=f_{\alpha\lambda}(A_\lambda)$；再取 $\gamma\geq\beta$，使 $\lambda\geq\gamma$ 时 $h_{\beta\lambda}(C_\lambda)=h_{\beta\gamma}(C_\gamma)$。设 $y_\alpha\in g_{\alpha\gamma}(B_\gamma)$，即 $y_\alpha=g_{\alpha\gamma}(y_\gamma)$，其中 $y_\gamma\in B_\gamma$。置 $y_\beta=g_{\beta\gamma}(y_\gamma)$，则 $v_\beta(y_\beta)=h_{\beta\gamma}(v_\gamma(y_\gamma))$。

对任意 $\lambda\geq\gamma$，由假设存在 $y_\lambda\in B_\lambda$，使
\[
h_{\beta\gamma}(v_\gamma(y_\gamma))
=h_{\beta\lambda}(v_\lambda(y_\lambda))
=v_\beta(g_{\beta\lambda}(y_\lambda)).
\]
故 $v_\beta(y_\beta-g_{\beta\lambda}(y_\lambda))=0$，存在 $x_\beta\in A_\beta$，使 $y_\beta=g_{\beta\lambda}(y_\lambda)+u_\beta(x_\beta)$。于是
\[
y_\alpha=g_{\alpha\lambda}(y_\lambda)+u_\alpha(f_{\alpha\beta}(x_\beta)).
\]
因 $\lambda\geq\beta$，存在 $x_\lambda\in A_\lambda$，使 $f_{\alpha\beta}(x_\beta)=f_{\alpha\lambda}(x_\lambda)$，所以
\[
y_\alpha=g_{\alpha\lambda}(y_\lambda+u_\lambda(x_\lambda))
\in g_{\alpha\lambda}(B_\lambda).
\]
证明完成。
\end{proof}

\egastatement{13.2.2}{命题}
设 $I$ 为具有可数共尾子集的滤过有序集，且
\[
0\longrightarrow A_\alpha\xrightarrow{u_\alpha}B_\alpha
\xrightarrow{v_\alpha}C_\alpha\longrightarrow0
\]
是以 $I$ 为指标集的阿贝尔群射影系统正合列。若 $(A_\alpha)$ 满足（ML），则列
\[
0\longrightarrow\plim A_\alpha\longrightarrow\plim B_\alpha
\longrightarrow\plim C_\alpha\longrightarrow0
\]
正合。
\begin{proof}
只需证明 $v=\plim v_\alpha:\plim B_\alpha\to\plim C_\alpha$ 满射。取 $z=(z_\alpha)\in\plim C_\alpha$，置 $E_\alpha=v_\alpha^{-1}(z_\alpha)$。各 $E_\alpha$ 在 $g_{\alpha\beta}$ 的限制下构成非空集合射影系统。证明此系统满足（ML）：通过 $u_\alpha$ 将 $A_\alpha$ 识别为 $B_\alpha$ 的子群，取 $\beta\geq\alpha$，使 $\lambda\geq\beta$ 时 $g_{\alpha\beta}(A_\beta)=g_{\alpha\lambda}(A_\lambda)$。将证明同时有 $g_{\alpha\beta}(E_\beta)=g_{\alpha\lambda}(E_\lambda)$。

事实上，取 $y_\lambda\in E_\lambda$，置 $y_\beta=g_{\beta\lambda}(y_\lambda)$、$y_\alpha=g_{\alpha\lambda}(y_\lambda)$。若 $y'_\alpha\in g_{\alpha\beta}(E_\beta)$，写成 $y'_\alpha=g_{\alpha\beta}(y'_\beta)$，则 $x_\beta=y'_\beta-y_\beta\in A_\beta$；存在 $x_\lambda\in A_\lambda$，使 $g_{\alpha\beta}(x_\beta)=g_{\alpha\lambda}(x_\lambda)$。因此
\[
\begin{aligned}
y'_\alpha&=g_{\alpha\beta}(y_\beta)+g_{\alpha\beta}(x_\beta)\\
&=g_{\alpha\lambda}(y_\lambda)+g_{\alpha\lambda}(x_\lambda)
=g_{\alpha\lambda}(y_\lambda+x_\lambda)\in g_{\alpha\lambda}(E_\lambda).
\end{aligned}
\]
这证明了（ML）。已知在上述对 $I$ 的假设下，满足（ML）的非空集合射影系统有非空射影极限（Bourbaki，《一般拓扑》，第 II 章，第3版，§3，定理1）。故存在 $y=(y_\alpha)\in\plim E_\alpha$；按定义每个 $v_\alpha(y_\alpha)=z_\alpha$，所以 $v(y)=z$。
\end{proof}

\egastatement{13.2.3}{命题}
对 $I$ 作（13.2.2）的假设。设 $(K_\alpha^\bullet)_{\alpha\in I}$ 为阿贝尔群复形的射影系统，各复形 $K_\alpha^\bullet=(K_\alpha^n)_{n\in\Z}$ 的微分次数为 $+1$。对每个 $n$，存在典范的函子同态
\begin{equation}\tag{13.2.3.1}
h_n:\mathrm H^n(\plim K_\alpha^\bullet)\longrightarrow\plim\mathrm H^n(K_\alpha^\bullet).
\end{equation}
若对每个次数 $n$，$(K_\alpha^n)_{\alpha\in I}$ 满足（ML），则所有 $h_n$ 均满射。若再对某个 $n$，$(\mathrm H^{n-1}(K_\alpha^\bullet))_{\alpha\in I}$ 满足（ML），则 $h_n$ 为双射。
\begin{proof}
对每个 $n$ 置 $K^n=\plim K_\alpha^n$。由交换图
\[
\begin{tikzcd}[column sep=small]
\cdots\arrow[r]&K^{n-1}\arrow[r]\arrow[d]&K^n\arrow[r]\arrow[d]&K^{n+1}\arrow[r]\arrow[d]&\cdots\\
\cdots\arrow[r]&K_\alpha^{n-1}\arrow[r]&K_\alpha^n\arrow[r]&K_\alpha^{n+1}\arrow[r]&\cdots
\end{tikzcd}
\]
即可定义 $h_n$，其中 $K^\bullet$ 的微分是各 $K_\alpha^\bullet$ 相应微分的射影极限。考虑正合列
\[
(*_n)\qquad0\to\mathrm B^n(K_\alpha^\bullet)\to\mathrm Z^n(K_\alpha^\bullet)
\to\mathrm H^n(K_\alpha^\bullet)\to0,
\]
\[
(**_n)\qquad0\to\mathrm Z^{n-1}(K_\alpha^\bullet)\to K_\alpha^{n-1}
\to\mathrm B^n(K_\alpha^\bullet)\to0.
\]
由假设及（13.2.1，i），每个 $(\mathrm B^n(K_\alpha^\bullet))_\alpha$ 满足（ML）；于是由（13.2.2），列
\[
(***_n)\quad0\to\plim\mathrm B^n(K_\alpha^\bullet)
\to\plim\mathrm Z^n(K_\alpha^\bullet)
\to\plim\mathrm H^n(K_\alpha^\bullet)\to0
\]
正合。显然 $\plim\mathrm B^n(K_\alpha^\bullet)$ 可识别为 $K^n$ 中包含 $\mathrm B^n(K^\bullet)$ 的子群，而 $\plim\mathrm Z^n(K_\alpha^\bullet)$ 可识别为 $\mathrm Z^n(K^\bullet)$ 的子群，所以 $h_n$ 满射。\transnote{原版此处将包含 $\plim\mathrm B^n(K_\alpha^\bullet)$ 的群印作 $K^{n+1}$；按本段定义应为 $K^n$，这里据其定义明示订正。}

再假设 $(\mathrm H^{n-1}(K_\alpha^\bullet))_\alpha$ 满足（ML）。由 $(*_{n-1})$ 及（13.2.1，ii），$(\mathrm Z^{n-1}(K_\alpha^\bullet))_\alpha$ 满足（ML）；将（13.2.2）应用于 $(**_n)$，得到正合列
\[
0\to\plim\mathrm Z^{n-1}(K_\alpha^\bullet)\to K^{n-1}
\xrightarrow{u}\plim\mathrm B^n(K_\alpha^\bullet)\to0.
\]
由于将 $u$ 与嵌入 $\plim\mathrm B^n(K_\alpha^\bullet)\to K^n$ 复合就是微分 $K^{n-1}\to K^n$，$u$ 的满射性说明 $\plim\mathrm B^n(K_\alpha^\bullet)=\mathrm B^n(K^\bullet)$，故 $h_n$ 单射。
\end{proof}

\egastatement{13.2.4}{注}
（i）（13.2.2）的论证（见 Bourbaki，同上）说明，其结论在以下较弱的假设下仍成立：各 $A_\alpha$ 能赋予完备可度量空间结构，使平移为同胚，定义射影系统的 $f_{\alpha\beta}:A_\beta\to A_\alpha$ 对所取距离一致连续，并且满足

\noindent（ML$'$）对每个 $\alpha$，存在 $\beta\geq\alpha$，使每个 $\gamma\geq\beta$ 的 $f_{\alpha\gamma}(A_\gamma)$ 在 $f_{\alpha\beta}(A_\beta)$ 中稠密。

这也可给（13.2.3）作相应补充：假设对每个 $\alpha$，$n<0$ 时 $K_\alpha^n=0$；又假设对 $n\geq0$，$(K_\alpha^n)_\alpha$ 满足（ML），并且 $A_\alpha=\mathrm H^0(K_\alpha^\bullet)$ 能赋予满足上述性质的度量空间结构。那么（13.2.3）对于 $n\geq2$ 的结论不变，且 $h_1$ 为双射。因为（13.2.2）的论证仍说明 $(\mathrm B^1(K_\alpha^\bullet))_\alpha$ 满足（ML），$(***_1)$ 正合，并由以上说明有 $\plim\mathrm B^1(K_\alpha^\bullet)=\mathrm B^1(K^\bullet)$。这样特别证明了（T，3.10.2）中的断言。

（ii）可以引入函子 $\plim$ 的右导出函子 $\plim{}^{(p)}$，从而得到比上述更完整的陈述 [28]。

\subsection{应用：层的射影极限的上同调}
\src{68--69}
\egastatement{13.3.1}{命题}
设 $X$ 为拓扑空间，$(\F_k)_{k\in\mathbf N}$ 为 $X$ 上阿贝尔群层的射影系统，$\F=\plim\F_k$。假设：
\begin{enumerate}[label=(\roman*)]
\item $X$ 的拓扑有基 $\mathfrak B$，使对每个 $U\in\mathfrak B$ 及每个 $i\geq0$，射影系统 $(\mathrm H^i(U,\F_k))_k$ 满足（ML）。
\item 对每个 $x\in X$ 及 $i>0$，有
\[
\ilim_U\bigl(\plim_k\mathrm H^i(U,\F_k)\bigr)=0,
\]
其中 $U$ 遍历 $\mathfrak B$ 中的 $x$ 的邻域。
\item 定义射影系统的同态 $u_{hk}:\F_k\to\F_h$（$h\leq k$）均满射。
\end{enumerate}
则对每个 $i>0$，典范同态 $h_i:\mathrm H^i(X,\F)\to\plim_k\mathrm H^i(X,\F_k)$ 满射。若再对某个 $i$，$(\mathrm H^{i-1}(X,\F_k))_k$ 满足（ML），则 $h_i$ 为双射。
\begin{proof}
a）先假设各 $\F_k$ 以及 $u_{hk}$ 的核 $\J_{hk}$ 都为松弛层。将证明此时（iii）推出 $i>0$ 时 $\mathrm H^i(X,\F)=0$。只需证明对 $X$ 的每个开集 $U$ 及 $U$ 的任意开覆盖 $\mathfrak U$，$i>0$ 时都有 $\mathrm H^i(\mathfrak U,\F)=0$。这首先说明 Čech 上同调 $\check{\mathrm H}^i(U,\F)=0$，再将（G，II，5.9.2）应用于 $X$ 的所有开集，即得 $\mathrm H^i(X,\F)=0$。

由于各 $\F_k$ 松弛，$i>0$ 时 $\mathrm H^i(\mathfrak U,\F_k)=0$（G，II，5.2.3）。对每个 $k$，取覆盖神经的交错上链复形 $C^\bullet(\mathfrak U,\F_k)$（G，II，5.1），它们显然构成阿贝尔群复形射影系统。将证明所有映射 $C^\bullet(\mathfrak U,\F_k)\to C^\bullet(\mathfrak U,\F_h)$（$h\leq k$）均满射。依定义，只需证明对每个开集 $V$，$\Gamma(V,\F_k)\to\Gamma(V,\F_h)$ 满射；但由正合列 $0\to\J_{hk}\to\F_k\to\F_h\to0$，得到
\[
\Gamma(V,\F_k)\to\Gamma(V,\F_h)\to\mathrm H^1(V,\J_{hk})=0,
\]
因为 $\J_{hk}$ 松弛。故 $(C^i(\mathfrak U,\F_k))_k$ 满足（ML）。各 $(\mathrm H^i(\mathfrak U,\F_k))_k$ 也满足（ML）：$i>0$ 时显然，而 $i=0$ 时，因 $\mathrm H^0(\mathfrak U,\F_k)=\Gamma(U,\F_k)$（G，II，5.2.2），由刚才的结论成立。因此可用（13.2.3），得到 $i>0$ 时
\[
\mathrm H^i(\mathfrak U,\F)=\plim_k\mathrm H^i(\mathfrak U,\F_k)=0.
\]

b）转到一般情形。对每个 $k$，取 $\F_k$ 的典范松弛解消 $\C^\bullet(X,\F_k)=(\C^i(X,\F_k))_{i\geq0}$（G，II，4.3）。对每个 $i\geq0$，$(\C^i(X,\F_k))_k$ 显然是松弛层射影系统，且满足 a）的条件：若 $\J_{hk}=\Ker u_{hk}$，由（iii），列 $0\to\J_{hk}\to\F_k\to\F_h\to0$ 正合；而函子 $\F\mapsto\C^i(X,\F)$ 正合（G，II，4.3）。

置 $\G^i=\plim_k\C^i(X,\F_k)$；由 a），$j>0$、$i\geq0$ 时 $\mathrm H^j(X,\G^i)=0$。将证明 $\G^\bullet=(\G^i)_{i\geq0}$ 是 $\F$ 的解消；于是由上述消失性，$\mathrm H^\bullet(X,\F)=\mathrm H^\bullet(\Gamma(X,\G^\bullet))$（G，II，4.7.1）。

对正合列 $0\to\F_k\to\C^0(X,\F_k)\to\C^1(X,\F_k)\to\cdots$ 取射影极限，显然得到层复形 $0\to\F\to\G^0\to\G^1\to\cdots$。只需证明 $i>0$ 时 $\HH^i(\G^\bullet)=0$；此层由预层 $U\mapsto\mathrm H^i(\Gamma(U,\G^\bullet))$ 生成（G，II，4.1）。而 $\Gamma(U,\G^\bullet)$ 是复形系统 $(\Gamma(U,\C^\bullet(X,\F_k)))_k$ 的射影极限（$\OI$，3.2.6）。由 a），各项间映射都满射；又因典范解消的限制仍为典范解消，
\[
\mathrm H^i(U,\F_k)=\mathrm H^i(\Gamma(U,\C^\bullet(X,\F_k))).
\]
由（i），对每个 $U\in\mathfrak B$ 可应用（13.2.3），故 $i\geq0$ 时
\[
\mathrm H^i(\Gamma(U,\G^\bullet))=\plim_k\mathrm H^i(U,\F_k).
\]
（ii）于是按定义说明 $i>0$ 时 $\HH^i(\G^\bullet)=0$。

现在 $\mathrm H^i(X,\F)=\mathrm H^i(\Gamma(X,\G^\bullet))$，且
\[
\Gamma(X,\G^\bullet)=\plim_k\Gamma(X,\C^\bullet(X,\F_k)).
\]
刚才已经说明此系统在每个次数上的映射都满射；再次应用（13.2.3）即得结论。
\end{proof}

\egastatement{13.3.2}{注}
（i）（13.3.1）只有在 $i>0$ 时才有实质内容，因为 $i=0$ 时，不作假设，$h_0$ 也总为同构（$\OI$，3.2.6）。

（ii）特别地，若对每个 $k$、$i>0$、$U\in\mathfrak B$ 都有 $\mathrm H^i(U,\F_k)=0$，且各 $\Gamma(U,\F_k)\to\Gamma(U,\F_h)$ 满射，则（13.3.1）的（i）、（ii）成立。这将是（13.3.1）最常用的应用情形。\par

% ---- END c0_s13a.tex ----
% ---- BEGIN c0_s13b.tex ----
\subsection{Mittag--Leffler 条件与射影系统的相伴分次对象}
\src{69--71}
\egapar{13.4.1}
设 $\mathbf A=(A_k,u_{hk})_{k\in\Z}$ 为阿贝尔范畴 $C$ 中的射影系统；若存在 $k_0$ 使 $k<k_0$ 时 $A_k=0$，则称它\textbf{有下界}。在每个 $A_k$ 上以下列公式定义滤过 $(F^p(A_k))_{p\in\Z}$：
\begin{equation}\tag{13.4.1.1}
\begin{cases}
F^p(A_k)=\Ker(A_k\to A_{p-1}),&p\leq k+1,\\
F^p(A_k)=0,&p\geq k+1.
\end{cases}
\end{equation}
依假设 $F^{k_0}(A_k)=A_k$ 且 $F^{k+1}(A_k)=0$，所以这是有限滤过（11.1.3）。相伴分次对象为
\[
\gr^p(A_k)=\Ker(A_k\to A_{p-1})/\Ker(A_k\to A_p),
\]
因而同构于 $\Ker(A_k\to A_{p-1})$ 经 $A_k\to A_p$ 的像。由射影系统态射的传递性，
\begin{equation}\tag{13.4.1.2}
\gr^p(A_k)=\Ker(A_p\to A_{p-1})\cap\Im(A_k\to A_p).
\end{equation}
由（13.4.1.1），$\Ker(A_p\to A_{p-1})=\gr^p(A_p)$，所以也有
\begin{equation}\tag{13.4.1.3}
\gr^p(A_k)=\gr^p(A_p)\cap\Im(A_k\to A_p).
\end{equation}
此外，这些定义说明，当 $k\geq h$ 时 $u_{hk}(F^p(A_k))\subset F^p(A_h)$；故对每个 $p\in\Z$，各 $\gr^p(u_{hk})$ 定义射影系统 $(\gr^p(A_k))_k$。

\egapar{13.4.2}
若 $k$ 充分大时态射 $A_{k+1}\to A_k$ 都为同构，则称射影系统 $\mathbf A$ \textbf{最终常值}。若各 $A_j\to A_i$（$j\geq i$）均为满态射，则称 $\mathbf A$ \textbf{严格}。

当 $\mathbf A$ 严格时，由（13.4.1.3），对 $p\leq k\leq h$，典范态射 $\gr^p(A_h)\to\gr^p(A_k)$ 为同构；换言之，$(\gr^p(A_k))_k$ 最终常值。这时，对每个 $p$，把 $\gr^p(A_p)$ 与 $\plim_k\gr^p(A_k)$ 相识别，所成的对象列记为 $\gr^\bullet(\mathbf A)$，称为严格射影系统 $\mathbf A$ 的\textbf{相伴分次对象}。

现在设有下界射影系统 $\mathbf A$ 满足（ML）。已知其普遍像对象构成的系统 $\mathbf A'=(A'_k)$ 严格（13.1.2），且仍有下界。将 $\mathbf A'$ 的相伴分次对象 $\gr^\bullet(\mathbf A')$ 也称为 $\mathbf A$ 的相伴分次对象，并记作 $\gr^\bullet(\mathbf A)$。

\egastatement{13.4.3}{命题}
设 $\mathbf A=(A_k)_{k\in\Z}$ 为阿贝尔范畴 $C$ 中有下界的射影系统。下列条件等价：

\noindent a）$\mathbf A$ 满足（ML）。\qquad
b）对每个 $p\in\Z$，$(\gr^p(A_k))_{k\in\Z}$ 最终常值。

这些条件成立时，对每个 $p\in\Z$ 有典范同构
\begin{equation}\tag{13.4.3.1}
\gr^p(\mathbf A)\simeq\plim_k\gr^p(A_k).
\end{equation}
\begin{proof}
由（13.4.1.2），a）立即推出 b）；将同一公式应用于 $\mathbf A'$（记号见13.4.2），按定义即得（13.4.3.1）。对 $k\leq h$，置 $A_{kh}=\Im(A_h\to A_k)$；若 $k\leq h\leq j$，则 $A_{kj}\subset A_{kh}\subset A_k$。

在 $A_{kh}$ 上取由 $(F^p(A_k))$ 诱导的滤过。由定义系统的态射的传递性，立即可验证它也是 $(F^p(A_h))$ 的商滤过；故
\begin{equation}\tag{13.4.3.2}
\gr^p(A_{kh})=\Im(\gr^p(A_h)\to\gr^p(A_k)).
\end{equation}
设 b）成立。对每个 $p\in\Z$ 及 $k\geq p$，存在整数 $L(p,k)$，使 $h\geq L(p,k)$ 时（13.4.3.2）右端恒定。由于 $p<k_0$ 时 $\gr^p(A_{kh})=0$，对于给定 $k$，只有有限多个 $p\leq k$ 的 $L(p,k)$ 非零。以 $L(k)=m$ 表示这些整数的最大值。则 $h\geq m$ 时 $A_{kh}\subset A_{km}$，且由 $m$ 的定义，典范嵌入 $A_{kh}\to A_{km}$ 给出同构 $\gr^\bullet(A_{kh})\simeq\gr^\bullet(A_{km})$。由于滤过有限，嵌入本身为双射（Bourbaki，《交换代数》，第 III 章，§2，第8小节，定理1），所以 $\mathbf A$ 满足（ML）。
\end{proof}

\egapar{13.4.4}
设 $C$ 中存在射影极限 $A=\plim A_k$。在（13.4.1）的定义中可以把 $A_k$ 换成 $A$，如此得到的滤过仍满足
\begin{equation}\tag{13.4.4.1}
\gr^p(A)=\gr^p(A_p)\cap\Im(A\to A_p).
\end{equation}

\egastatement{13.4.5}{推论}
设 $C$ 为阿贝尔群范畴。若 $\mathbf A$ 满足（ML），且 $A=\plim A_k$，则对每个 $p\in\Z$ 有典范同构
\begin{equation}\tag{13.4.5.1}
\gr^p(A)\simeq\plim_k\gr^p(A_k).
\end{equation}
\begin{proof}
当 $k$ 充分大时，$\Im(A_k\to A_p)=\Im(A\to A_p)$（Bourbaki，《一般拓扑》，第 II 章，第3版，§3，第5小节，定理1）。因此结论由（13.4.1.3）、（13.4.4.1）得出。
\end{proof}

\subsection{滤过复形谱序列的射影极限}
\src{71--73}
\egapar{13.5.1}
设 $C$ 为阿贝尔范畴，$X^\bullet$ 为 $C$ 的复形，赋有滤过 $(F^p(X^\bullet))_p$，且存在 $p_0$ 使 $F^{p_0}(X^\bullet)=X^\bullet$。对每个 $k\in\Z$，考虑 $X_k^\bullet=X^\bullet/F^{k+1}(X^\bullet)$。它典范地具有由下列各项组成的滤过：$p\leq k$ 时
\[
F^p(X_k^\bullet)=F^p(X^\bullet)/F^{k+1}(X^\bullet),
\]
$p\geq k+1$ 时 $F^p(X_k^\bullet)=0$。此外，有典范态射 $X_{k+1}^\bullet\to X_k^\bullet$，使 $\mathbf X^\bullet=(X_k^\bullet)_{k\in\Z}$ 成为 $C$ 的滤过复形射影系统。注意它严格，且 $k<p_0$ 时 $X_k^\bullet=0$。

\egapar{13.5.2}
更一般地，考虑 $C$ 中有下界的严格复形射影系统 $\mathbf X^\bullet=(X_k^\bullet)_{k\in\Z}$；在每个 $X_k^\bullet$ 上取（13.4.1）定义的滤过（在 $C$ 的有下界复形所成阿贝尔范畴中进行）。各 $X_k^\bullet\to X_p^\bullet$（$p\leq k$）成为有限滤过复形的态射。滤过复形谱序列的函子性（11.2.3）说明，定义 $\mathbf X^\bullet$ 的各态射也给出谱序列之间的态射，从而 $E(\mathbf X^\bullet)=(E(X_k^\bullet))$ 成为谱序列的射影系统。

\egastatement{13.5.3}{引理}
设滤过复形射影系统 $\mathbf X^\bullet=(X_k^\bullet)_{k\in\Z}$ 按（13.5.2）得到。则：
\begin{enumerate}[label=\alph*)]
\item 当 $r>p-p_0$ 时，对每个 $k\in\Z$，有 $B_r^{pq}(X_k^\bullet)=B_\infty^{pq}(X_k^\bullet)$。
\item 当 $k+1\leq p+r$ 时，有 $Z_r^{pq}(X_k^\bullet)=Z_\infty^{pq}(X_k^\bullet)$。
\item 当 $k+1\geq p+r$ 时，对每个 $h\geq k$，态射
\[
Z_r^{pq}(X_h^\bullet)\to Z_r^{pq}(X_k^\bullet),\qquad
B_r^{pq}(X_h^\bullet)\to B_r^{pq}(X_k^\bullet)
\]
均为同构。
\end{enumerate}
\begin{proof}
这三个性质立即由（11.2.2）的定义得出，并使用 $p-r<p_0$ 时 $F^{p-r+1}(X_k^\bullet)=X_k^\bullet$。
\end{proof}

\egapar{13.5.4}
设（13.5.3）的假设成立。对固定 $p,q,r$（$r$ 有限），射影系统 $(Z_r^{pq}(X_k^\bullet))_k$、$(B_r^{pq}(X_k^\bullet))_k$、$(E_r^{pq}(X_k^\bullet))_k$ 都最终常值；分别以 $Z_r^{pq}(\mathbf X^\bullet)$、$B_r^{pq}(\mathbf X^\bullet)$ 及
\[
E_r^{pq}(\mathbf X^\bullet)=Z_r^{pq}(\mathbf X^\bullet)/B_r^{pq}(\mathbf X^\bullet)
\]
表示它们的射影极限。前两个对象典范地识别为 $E_1^{pq}(\mathbf X^\bullet)$ 的子对象。$d_r^{pq}$ 的定义（M，XV，1）说明，相对于 $X_k^\bullet$ 的这些态射也最终常值，因而定义态射
\begin{equation}\tag{13.5.4.1}
d_r^{pq}:E_r^{pq}(\mathbf X^\bullet)\to E_r^{p+r,q-r+1}(\mathbf X^\bullet),
\end{equation}
满足 $d_r^{p+r,q-r+1}\circ d_r^{pq}=0$。此外，$\Ker d_r^{pq}$ 典范地同构于 $Z_{r+1}^{pq}(\mathbf X^\bullet)/B_r^{pq}(\mathbf X^\bullet)$，而 $\Im d_r^{pq}$ 典范地同构于
\[
B_{r+1}^{p+r,q-r+1}(\mathbf X^\bullet)/B_r^{p+r,q-r+1}(\mathbf X^\bullet).
\]

\egastatement{13.5.5}{引理}
在（13.5.3）的假设下，对 $s\geq r>p-p_0$，有典范单态射
\begin{equation}\tag{13.5.5.1}
i:E_s^{pq}(\mathbf X^\bullet)\to E_r^{pq}(\mathbf X^\bullet)
\end{equation}
及典范同构
\begin{equation}\tag{13.5.5.2}
j_r:E_r^{pq}(\mathbf X^\bullet)\iso E_\infty^{pq}(X_{p+r-1}^\bullet),
\end{equation}
使图
\begin{equation}\tag{13.5.5.3}
\begin{tikzcd}
E_s^{pq}(\mathbf X^\bullet)\arrow[r,"j_s"]\arrow[d,"i"']&E_\infty^{pq}(X_{p+s-1}^\bullet)\arrow[d]\\
E_r^{pq}(\mathbf X^\bullet)\arrow[r,"j_r"']&E_\infty^{pq}(X_{p+r-1}^\bullet)
\end{tikzcd}
\end{equation}
交换；右边的竖直箭头来自 $X_{p+s-1}^\bullet\to X_{p+r-1}^\bullet$。
\begin{proof}
$i$ 的存在性来自 $r>p-p_0$ 时 $B_r^{pq}(X_k^\bullet)=B_\infty^{pq}(X_k^\bullet)$（13.5.3，a）。当 $k+1\leq p+r$ 时，$Z_r^{pq}(X_k^\bullet)=Z_\infty^{pq}(X_k^\bullet)$（13.5.3，b）；特别地，对 $k=p+r-1$ 成立。另一方面，由（13.5.3，c），$Z_r^{pq}(X_{p+r-1}^\bullet)$、$B_r^{pq}(X_{p+r-1}^\bullet)$ 分别典范地识别为 $Z_r^{pq}(\mathbf X^\bullet)$、$B_r^{pq}(\mathbf X^\bullet)$。这就得到 $j_r$ 的存在性及图的交换性。
\end{proof}

\egastatement{13.5.6}{推论}
在（13.5.3）的假设下，若 $\plim_rE_r^{pq}(\mathbf X^\bullet)$、$\plim_kE_\infty^{pq}(X_k^\bullet)$ 中一个存在，则另一个也存在，并有典范同构
\begin{equation}\tag{13.5.6.1}
j_\infty:\plim_rE_r^{pq}(\mathbf X^\bullet)\iso\plim_kE_\infty^{pq}(X_k^\bullet).
\end{equation}
此外，第一个射影系统最终常值，当且仅当第二个射影系统最终常值。

\egapar{13.5.7}
以 $B_\infty^{pq}(\mathbf X^\bullet)$、$Z_\infty^{pq}(\mathbf X^\bullet)$ 表示 $E_1^{pq}(\mathbf X^\bullet)$ 中分别等于 $r>p-p_0$ 时 $B_r^{pq}(\mathbf X^\bullet)$，及 $\inf_r Z_r^{pq}(\mathbf X^\bullet)$（在其存在时）的子对象。于是 $\plim_rE_r^{pq}(\mathbf X^\bullet)$ 典范地识别为
\[
E_\infty^{pq}(\mathbf X^\bullet)
=Z_\infty^{pq}(\mathbf X^\bullet)/B_\infty^{pq}(\mathbf X^\bullet).
\]
注意各 $Z_r^{pq}(\mathbf X^\bullet)$、$B_r^{pq}(\mathbf X^\bullet)$、$E_r^{pq}(\mathbf X^\bullet)$（$1\leq r\leq+\infty$）及 $d_r^{pq}$，都以函子方式依赖于满足（13.5.5）限制条件的系统 $\mathbf X^\bullet$；而（13.5.5）、（13.5.6）所定义的态射也具有函子性。

\subsection{函子相对于有限滤过对象的谱序列}
\src{73--75}
\egapar{13.6.1}
设 $C,C'$ 为阿贝尔范畴，$T:C\to C'$ 为加性协变函子。假设 $C$ 的每个对象都同构于某个内射对象的子对象，故右导出函子 $\mathrm R^pT$（$p\geq0$）存在。

\egastatement{13.6.2}{引理}
设对象 $A\in C$ 具有有限滤过 $(F^i(A))_{i\in\Z}$。则存在 $A$ 的内射解消 $X^\bullet=(X^n)_{n\geq0}$，并具有有限滤过 $(F^i(X^\bullet))_{i\in\Z}$，使 $F^i(A)=A$（分别地，$F^i(A)=0$）推出 $F^i(X^\bullet)=X^\bullet$（分别地，$F^i(X^\bullet)=0$），且对每个 $i$，$F^i(X^\bullet)$ 是 $F^i(A)$ 的内射解消。
\begin{proof}
设 $p$ 为满足 $F^p(A)=A$ 的最大指标，$q>p$ 为满足 $F^q(A)=0$ 的最小指标。对 $q-p$ 归纳；$q-p=1$ 时显然。已构造好 $A/F^{q-1}(A)$ 的满足所需性质的内射解消 $X''^\bullet$ 后，考虑正合列
\[
0\to F^{q-1}(A)\to A\to A/F^{q-1}(A)\to0.
\]
取 $F^{q-1}(A)$ 的内射解消 $X'^\bullet$，再确定 $A$ 的内射解消 $X^\bullet$，使有与上述列相容的正合列 $0\to X'^\bullet\to X^\bullet\to X''^\bullet\to0$（M，V，2.2）；显然所得 $X^\bullet$ 满足要求。
\end{proof}

\egastatement{13.6.3}{推论}
设 $B$ 为 $C$ 的另一个有限滤过对象，$s$ 为整数，$u:A\to B$ 满足对每个 $i\in\Z$ 有 $u(F^i(A))\subset F^{i+s}(B)$。若 $Y^\bullet=(Y^n)_{n\geq0}$ 是 $B$ 的内射解消，赋有满足（13.6.2）性质的滤过，则存在与 $u$ 相容的态射 $v:X^\bullet\to Y^\bullet$，使 $v(F^i(X^\bullet))\subset F^{i+s}(Y^\bullet)$。此外，任意两个这样的态射 $v,v'$ 都同伦。
\begin{proof}
由上述构造及（M，V，2.3），对 $q-p$ 归纳即得。
\end{proof}

\egapar{13.6.4}
在（13.6.2）的假设下，考虑 $C'$ 中的复形 $T(X^\bullet)$。由于各 $F^p(X^\bullet)$ 逐项是 $X^\bullet$ 的直和因子，$T(F^p(X^\bullet))$ 显然给它一个滤过。由（13.6.3），该滤过复形的谱序列在同构意义下只依赖于滤过对象 $A$。其终项为 $\mathrm R^\bullet T(A)$，所带滤过为
\begin{equation}\tag{13.6.4.1}
\begin{aligned}
F^p(\mathrm R^nT(A))
&=\Im(\mathrm R^nT(F^p(A))\to\mathrm R^nT(A))\\
&=\Ker(\mathrm R^nT(A)\to\mathrm R^nT(A/F^p(A)))
\end{aligned}
\end{equation}
（11.2.2），其 $E_1$ 项为
\begin{equation}\tag{13.6.4.2}
E_1^{pq}=\mathrm R^{p+q}T(\gr^p(A)),
\end{equation}
这里照例 $\gr^p(A)=F^p(A)/F^{p+1}(A)$。由（11.2.2），终项的滤过有限；对给定 $p,q$，$B_r^{pq}(A)=B_r^{pq}(T(X^\bullet))$ 与 $Z_r^{pq}(A)=Z_r^{pq}(T(X^\bullet))$ 的序列都稳定，因此谱序列双正则（11.1.3）。将它记为 $E(A)=(E_r^{pq}(A))$，称作\textbf{函子 $T$ 相对于滤过对象 $A$ 的谱序列}。

\egapar{13.6.5}
现在设（13.6.3）的假设成立，并沿用其记号。各 $F^i(X^\bullet)$、$F^i(Y^\bullet)$ 分别逐项为 $X^\bullet$、$Y^\bullet$ 的直和因子，故对每个 $i$，
\[
(Tv)(T(F^i(X^\bullet)))\subset T(F^{i+s}(Y^\bullet)).
\]
由（11.2.2）的定义，对 $1\leq r\leq+\infty$，$Tv$ 给出态射
\[
B_r^{pq}(T(X^\bullet))\to B_r^{p+s,q-s}(T(Y^\bullet)),\qquad
Z_r^{pq}(T(X^\bullet))\to Z_r^{p+s,q-s}(T(Y^\bullet)),
\]
从而得到 $w_r:E_r^{pq}(A)\to E_r^{p+s,q-s}(B)$。同样在终项上得到 $u_n:\mathrm R^nT(A)\to\mathrm R^nT(B)$，满足
\[
u_n(F^p(\mathrm R^nT(A)))\subset F^{p+s}(\mathrm R^nT(B)).
\]
$d_r^{pq}$ 的定义（M，XV，1）还说明以下图交换：
\[
\begin{tikzcd}
E_r^{pq}(A)\arrow[r,"d_r^{pq}"]\arrow[d,"w_r"']&E_r^{p+r,q-r+1}(A)\arrow[d,"w_r"]\\
E_r^{p+s,q-s}(B)\arrow[r,"d_r^{p+s,q-s}"']&E_r^{p+r+s,q-r-s+1}(B).
\end{tikzcd}
\]
由此得到涉及 $\alpha_r^{pq}$ 的类似交换图，具体写法留给读者。最后，对终项也有交换图（同上）：
\[
\begin{tikzcd}
E_\infty^{pq}(A)\arrow[r,"\beta^{pq}"]\arrow[d]&\gr^p(\mathrm R^{p+q}T(A))\arrow[d,"u_{p+q}"]\\
E_\infty^{p+s,q-s}(B)\arrow[r,"\beta^{p+s,q-s}"']&\gr^{p+s}(\mathrm R^{p+q}T(B)).
\end{tikzcd}
\]

\egapar{13.6.6}
特别地，设存在赋有滤过 $(F^i(S))_{i\in\Z}$ 的环 $S$ 及环同态
\begin{equation}\tag{13.6.6.1}
h:S\longrightarrow\Hom_C(A,A),
\end{equation}
使对 $t\in F^j(S)$ 及任意 $i,j$，均有 $h_t(F^i(A))\subset F^{i+j}(A)$。简言之，称 $A$ 具有滤过环 $S$ 上的\textbf{滤过 $S$-$C$-模}结构。

取相伴分次对象，$t\in F^j(S)$ 给出的每个 $h_t$ 都在 $\gr^\bullet(A)$ 上定义次数为 $j$ 的齐次分次自同态 $\bar h_t$，且它只依赖于 $t$ 在 $\gr^j(S)$ 中的类。这样得到分次环同态
\[
\bar h:\gr^\bullet(S)\longrightarrow\Hom_C(\gr^\bullet(A),\gr^\bullet(A)),
\]
右端为分次自同态环。称 $\gr^\bullet(A)$ 具有\textbf{分次 $\gr^\bullet(S)$-$C$-模}结构。

由（13.6.5），对 $1\leq r\leq+\infty$，每个 $t\in\gr^s(S)$ 在双分次对象 $(B_r^{pq}(A))_{p,q}$、$(Z_r^{pq}(A))_{p,q}$、$E_r(A)=(E_r^{pq}(A))_{p,q}$ 上典范地定义双次数为 $(s,-s)$ 的自同态；当 $r$ 有限时，$E_r(A)$ 上此自同态与由 $d_r^{pq}$ 定义的双分次自同态交换。由于这些自同态满足通常的结合性及关于环和对象加法的分配性，简言之，称这些对象为\textbf{双分次 $\gr^\bullet(S)$-$C'$-模}；各 $\alpha_r^{pq}$ 显然给出这种结构的同构。

对每个整数 $n$，用 $B_r^n(A)$（分别地，$Z_r^n(A),E_r^n(A)$）表示只取 $p+q=n$ 的各 $B_r^{pq}(A)$（分别地，$Z_r^{pq}(A),E_r^{pq}(A)$）所成的分次子对象，$1\leq r<+\infty$。它们显然为分次 $\gr^\bullet(S)$-$C'$-模。最后，每个 $t\in\gr^j(S)$ 也在 $\gr^\bullet(\mathrm R^nT(A))$ 上定义次数为 $j$ 的分次自同态，使它成为分次 $\gr^\bullet(S)$-$C'$-模；$p+q=n$ 的各 $\beta^{pq}$ 给出 $E_\infty^n(A)$ 到 $\gr^\bullet(\mathrm R^nT(A))$ 的这种结构的同构。

当 $C'$ 为阿贝尔群范畴时，$S$-$C'$-模、分次或双分次 $\gr^\bullet(S)$-$C'$-模结构，就是通常的 $S$-模、分次或双分次 $\gr^\bullet(S)$-模结构。\par

% ---- END c0_s13b.tex ----
% ---- BEGIN c0_s13c.tex ----
\subsection{变元的射影极限与导出函子}
\src{75--79}
\egapar{13.7.1}
设 $C,C'$ 为阿贝尔范畴，且 $C$ 的每个对象都是某个内射对象的子对象；设 $T:C\to C'$ 为加性协变函子。考虑 $C$ 中有下界的严格射影系统 $\mathbf A=(A_k)_{k\in\Z}$；具体假设 $k<k_0$ 时 $A_k=0$。依公式（13.4.1.1），在每个 $A_k$ 上典范地赋予滤过 $(F^i(A_k))_{i\in\Z}$。由于系统严格，对 $i\leq j\leq k+1$，典范态射
\begin{equation}\tag{13.7.1.1}
F^i(A_h)/F^j(A_h)\longrightarrow F^i(A_k)/F^j(A_k)\qquad(h\geq k)
\end{equation}
均为同构。还要回顾，对每个 $k$，$F^{k_0}(A_k)=A_k$、$F^{k+1}(A_k)=0$。

\egapar{13.7.2}
对每个 $k$，构造 $A_k$ 的具有（13.6.2）性质的内射解消 $X_k^\bullet=(X_k^n)_{n\geq0}$。典范态射 $A_{k+1}\to A_k$ 使得可以为每个 $k$ 定义与滤过相容的复形态射 $X_{k+1}^\bullet\to X_k^\bullet$（13.6.3），从而 $\mathbf X^\bullet=(X_k^\bullet)_{k\in\Z}$ 成为复形射影系统。

还可以假设此系统严格。事实上，由同构（13.7.1.1），$A_k\simeq A_{k+1}/F^{k+1}(A_{k+1})$；因此构造 $X_{k+1}^\bullet$ 时，可把该商对象的内射解消取为 $X_k^\bullet$。由（M，V，2.3），可以将复形态射构造为经商给出保持滤过的同构
\[
X_{k+1}^\bullet/F^{k+1}(X_{k+1}^\bullet)\iso X_k^\bullet,
\]
这就是（13.5.1）的条件。

\egapar{13.7.3}
按构造，复形 $T(X_k^\bullet)$ 的射影系统 $T(\mathbf X^\bullet)$ 满足（13.5.3）的假设。因此（13.5.4）、（13.5.5）、（13.5.6）可用于谱序列 $E(T(X_k^\bullet))=E(A_k)$。对 $1\leq r\leq+\infty$，用 $E_r^{pq}(\mathbf A)$ 表示 $E_r^{pq}(T(\mathbf X^\bullet))$（$r=+\infty$ 时见13.5.7），其他类似记号亦然。特别地，由（13.6.4.2）及 $(\gr^p(A_k))_k$ 最终常值，
\begin{equation}\tag{13.7.3.1}
E_1^{pq}(\mathbf A)=\mathrm R^{p+q}T(\gr^p(\mathbf A)).
\end{equation}
这些结果与（13.4.3）给出下述命题，它首先由施维术（Shih Weishu）用另一种未发表的方法证明。

\egastatement{13.7.4}{命题（Shih）}
设 $n$ 为整数。下列两个条件等价：

\noindent a）对每对满足 $p+q=n$ 的 $(p,q)$，射影系统 $(E_r^{pq}(\mathbf A))_{r\geq2}$ 最终常值。

\noindent b）射影系统 $\mathrm R^nT(\mathbf A)=(\mathrm R^nT(A_k))_{k\in\Z}$ 满足（ML）。

当这些条件成立时，对每个 $p\in\Z$ 有典范同构
\begin{equation}\tag{13.7.4.1}
\gr^p(\mathrm R^nT(\mathbf A))\simeq E_\infty^{p,n-p}(\mathbf A).
\end{equation}
\begin{proof}
由（13.5.6），a）等价于 $p+q=n$ 时 $(E_\infty^{pq}(A_k))_k$ 最终常值。另一方面，$\gr^p(\mathrm R^nT(A_k))$ 典范地同构于 $E_\infty^{p,n-p}(A_k)$，所以由（13.4.3），a）、b）等价。同构（13.7.4.1）就是将（13.5.6.1）应用于当前情形所得。
\end{proof}

\egastatement{13.7.5}{推论}
设 $\mathbf F=(\F_k)_{k\in\mathbf N}$ 为满足（13.3.1）的（i）、（ii）、（iii）的阿贝尔群层射影系统，$\F=\plim\F_k$。对于函子 $\G\mapsto\Gamma(X,\G)$，假设每对满足 $p+q=n$ 或 $p+q=n+1$ 的 $(p,q)$，其射影系统 $(E_r^{pq}(\mathbf F))_r$ 都最终常值。在 $\mathrm H^{n+1}(X,\F)$ 上取滤过
\[
F^p(\mathrm H^{n+1}(X,\F))
=\Ker(\mathrm H^{n+1}(X,\F)\to\mathrm H^{n+1}(X,\F_{p-1})).
\]
则对每个 $p\in\Z$ 有典范同构
\begin{equation}\tag{13.7.5.1}
\gr^p(\mathrm H^{n+1}(X,\F))\simeq E_\infty^{p,n-p+1}(\mathbf F).
\end{equation}
\begin{proof}
在（13.7.4）中取 $C$ 为 $X$ 上阿贝尔群层范畴，$C'$ 为阿贝尔群范畴，$T=\Gamma$，得到对每个 $p$ 的典范同构
\[
\gr^p(\mathrm R^{n+1}T(\mathbf F))\simeq E_\infty^{p,n-p+1}(\mathbf F).
\]
另一方面，由（13.7.4），$(\mathrm H^n(X,\F_k))_k$ 满足（ML），故由（13.3.1）有典范同构
\begin{equation}\tag{13.7.5.1}
\mathrm H^{n+1}(X,\F)\iso\plim_k\mathrm H^{n+1}(X,\F_k).
\end{equation}
\transnote{原版将这一中间公式也编号为（13.7.5.1），与推论陈述中的公式重号；这里保留重号。}
又由（13.7.4），$\mathrm R^{n+1}\Gamma(\mathbf F)$ 满足（ML），故有典范同构
\[
\gr^p(\mathrm R^{n+1}\Gamma(\mathbf F))
\simeq\plim_k\gr^p(\mathrm H^{n+1}(X,\F_k))
\]
（13.4.3），以及
\[
\plim_k\gr^p(\mathrm H^{n+1}(X,\F_k))
\simeq\gr^p(\plim_k\mathrm H^{n+1}(X,\F_k))
\]
（13.4.5）。因此只需证明中间的同构（13.7.5.1）与两端的滤过相容；这立即由定义及对每个 $p$ 的交换图得出：
\[
\begin{tikzcd}
\mathrm H^{n+1}(X,\F)\arrow[rr,"\sim"]\arrow[dr]&&\plim_k\mathrm H^{n+1}(X,\F_k)\arrow[dl]\\
&\mathrm H^{n+1}(X,\F_{p-1})&
\end{tikzcd}
\]
\end{proof}

\egapar{13.7.6}
设环 $S$ 具有滤过 $(F^i(S))_{i\in\Z}$，且 $F^0(S)=S$（故 $i<0$ 时 $\gr^i(S)=0$）。假设每个 $A_k$ 连同（13.7.1）的滤过都是滤过 $S$-$C$-模（13.6.6），且 $A_k\to A_h$（$k\geq h$）均为这种结构的态射；简言之，称 $\mathbf A$ 为\textbf{滤过 $S$-$C$-模的射影系统}。

立即可知，对 $k\geq h$、$1\leq r\leq+\infty$，态射 $B_r^{pq}(A_k)\to B_r^{pq}(A_h)$、$Z_r^{pq}(A_k)\to Z_r^{pq}(A_h)$ 都是双分次 $\gr^\bullet(S)$-$C'$-模的态射（13.6.5）；而对有限的 $r$，族 $(Z_r^{pq}(\mathbf A))$、$(B_r^{pq}(\mathbf A))$、$(E_r^{pq}(\mathbf A))$ 也为双分次 $\gr^\bullet(S)$-$C'$-模，前两个是 $(E_1^{pq}(\mathbf A))$ 的子模。仍以 $Z_r^n(\mathbf A)$、$B_r^n(\mathbf A)$、$E_r^n(\mathbf A)$ 表示只取 $p+q=n$ 各项所得的子对象；它们是分次 $\gr^\bullet(S)$-$C'$-模。

当 $(E_r^{pq}(\mathbf A))_r$ 最终常值时，$(E_\infty^{pq}(\mathbf A))$ 也是双分次 $\gr^\bullet(S)$-$C'$-模，每个 $E_\infty^n(\mathbf A)$ 是分次模。此外，对每个 $k$，各
\[
\beta^{p,n-p}:E_\infty^{p,n-p}(A_k)\iso\gr^p(\mathrm R^nT(A_k))
\]
构成 $E_\infty^n(A_k)\to\gr^\bullet(\mathrm R^nT(A_k))$ 的分次 $\gr^\bullet(S)$-$C'$-模同构。在上述条件下，相应的
\[
E_\infty^n(\mathbf A)\iso\plim_k\gr^\bullet(\mathrm R^nT(A_k))
\]
也为这种结构的同构；而典范同构
\[
\gr^\bullet(\mathrm R^nT(\mathbf A))\iso\plim_k\gr^\bullet(\mathrm R^nT(A_k))
\]
显然亦然。因此，（13.7.4.1）的各同构合成一个分次 $\gr^\bullet(S)$-$C'$-模同构。

\NewNote{\textbf{作者勘误提示：}本命题的有限性假设与若干次数范围已有作者正式修正。 参见\hyperlink{errIII:24}{《勘误与增补（清单二）》第24条}。}
\egastatement{13.7.7}{命题}
设 $S$ 为诺特 $\frj$-进环。设 $C$ 为阿贝尔范畴，每个对象都同构于某个内射对象的子对象；$T$ 为从 $C$ 到阿贝尔群范畴的加性协变函子。设 $\mathbf A=(A_k)_{k\in\Z}$ 为有下界的严格滤过 $S$-$C$-模射影系统（$S$ 上取 $\frj$-进滤过）。对于给定整数 $n$，假设：

\noindent（$\mathrm F_n$）当 $m=n$ 及 $m=n+1$ 时，分次 $\gr^\bullet(S)$-模
\[
E_1^{(m)}(\mathbf A)
=(\mathrm R^mT(\gr^p(\mathbf A)))_{p\in\Z}
\]
为有限型（见13.7.3.1）。则：
\begin{enumerate}[label=(\roman*)]
\item 射影系统 $(\mathrm R^nT(A_k))_k$ 与 $(\mathrm R^{n+1}T(A_k))_k$ 满足（ML）。
\item 若置 $\mathrm R'^nT(\mathbf A)=\plim_k\mathrm R^nT(A_k)$，则 $\mathrm R'^nT(\mathbf A)$ 是有限型 $S$-模。
\item 在 $\mathrm R'^nT(\mathbf A)$ 上以
\[
F^p(\mathrm R'^nT(\mathbf A))
=\Ker(\mathrm R'^nT(\mathbf A)\to\mathrm R^nT(A_{p-1}))
\]
定义的滤过是 $\frj$-好滤过，即对每个 $p$，
\[
\frj F^p(\mathrm R'^nT(\mathbf A))\subset F^{p+1}(\mathrm R'^nT(\mathbf A)),
\]
且当 $p$ 充分大时等号成立。特别地，此滤过所定义的拓扑就是 $\frj$-进拓扑。
\item $p+q=n$ 或 $p+q=n+1$ 时，射影系统 $(E_r^{pq}(\mathbf A))_r$ 最终常值。因此 $E_\infty^{pq}(\mathbf A)$ 有定义（13.5.7），并有典范的分次 $\gr^\bullet(S)$-模同构
\begin{equation}\tag{13.7.7.1}
\gr^p(\mathrm R'^nT(\mathbf A))\simeq E_\infty^{p,n-p}(\mathbf A)\qquad(p\in\Z).
\end{equation}
\end{enumerate}
注意，利用（13.7.4.1），同构（13.7.7.1）允许我们略滥用记号，将射影系统 $\mathrm R^nT(\mathbf A)$ 的极限 $\mathrm R'^nT(\mathbf A)$ 也写成 $\mathrm R^nT(\mathbf A)$。

\eproof
由于分次环 $\gr^\bullet(S)$ 诺特（Bourbaki，《交换代数》，第 III 章，§2，第9小节，定理2的推论5），当 $m=n,n+1$ 时，$E_1^{(m)}(\mathbf A)$ 的分次子模递增列 $B_r^{(m)}(\mathbf A)$（13.6.6）稳定，所以（11.1.10）的条件 b）成立。因而（13.7.4）的条件 a）对 $n,n+1$ 成立，这已证明（i）。

此外，由（13.7.4.1）及（13.7.6）的说明，$\gr^\bullet(\mathrm R^nT(\mathbf A))$ 是同构于
\[
E_\infty^{(n)}(\mathbf A)=Z_\infty^{(n)}(\mathbf A)/B_\infty^{(n)}(\mathbf A)
\]
的分次 $\gr^\bullet(S)$-模。由于 $Z_\infty^{(n)}(\mathbf A)$ 为 $E_1^{(n)}(\mathbf A)$ 的子模，它有限型，故 $E_\infty^{(n)}(\mathbf A)$ 也有限型。又对于滤过 $(F^p(\mathrm R'^nT(\mathbf A)))$，（13.4.5）给出 $\gr^\bullet(\mathrm R'^nT(\mathbf A))$ 与 $\gr^\bullet(\mathrm R^nT(\mathbf A))$ 作为 $\gr^\bullet(S)$-模的同构，因此（iv）得证。（ii）、（iii）将由这些结果及下述引理得出。

\egastatement{13.7.7.2}{引理}
设 $S$ 为诺特 $\frj$-进环，$M$ 为赋有余离散滤过 $(F^p(M))_{p\in\Z}$ 的 $S$-模，并满足 $\frj F^p(M)\subset F^{p+1}(M)$（即 $M$ 是带有 $\frj$-进滤过的 $S$ 上的滤过模）。再假设 $M$ 对此滤过所定义的拓扑分离。则下列条件等价：
\begin{enumerate}[label=\alph*)]
\item $M$ 为有限型 $S$-模，且 $(F^p(M))$ 为 $\frj$-好滤过；
\item $\gr^\bullet(M)$ 为有限型 $\gr^\bullet(S)$-模；
\item 各 $\gr^p(M)$ 为有限型 $S$-模，且 $p$ 充分大时典范同态
\begin{equation}\tag{13.7.7.3}
\frj\otimes_S\gr^p(M)\longrightarrow\gr^{p+1}(M)
\end{equation}
均满射。
\end{enumerate}
这些典范同态由 $\frj\otimes_S F^p(M)\to F^{p+1}(M)$ 导出，因为复合同态
\[
\frj\otimes_S F^{p+1}(M)\to\frj\otimes_S F^p(M)\to F^{p+1}(M)
\]
的像是 $\frj F^{p+1}(M)\subset F^{p+2}(M)$。

证明见 Bourbaki，《交换代数》，第 III 章，§3，第1小节，命题3。

\egapar{13.7.7.4}
为应用引理（13.7.7.2），只需注意：上述滤过在 $\mathrm R'^nT(\mathbf A)$ 上定义的拓扑使它成为分离完备 $S$-模，因为这正是离散群 $\mathrm R^nT(A_k)$ 的射影极限拓扑。另一方面，若 $k<k_0$ 时 $A_k=0$，则此时 $\mathrm R^nT(A_k)=0$，故
\[
F^{k_0}(\mathrm R'^nT(\mathbf A))=\mathrm R'^nT(\mathbf A).
\]
因此引理的条件都满足。
\finishproof

\egastatement{13.7.8}{推论}
若假设（$\mathrm F_n$）成立，则对每个 $f\in S$，有典范同构
\begin{equation}\tag{13.7.8.1}
\plim_k((\mathrm R^nT(A_k))_f)
\simeq\mathrm R^nT(\mathbf A)\otimes_S S_{\{f\}}.
\end{equation}
\begin{proof}
$\mathrm R^nT(\mathbf A)$ 为有限型 $S$-模；$S_{\{f\}}$ 为诺特进 $S$-代数（$\OI$，7.6.11），是 $S_f$ 对 $\frj$-预进拓扑的分离完备化（$\OI$，7.6.2）。由（$\OI$，7.7.8、7.7.1），$\mathrm R^nT(\mathbf A)\otimes_S S_{\{f\}}$ 同构于 $\mathrm R^nT(\mathbf A)\otimes_S S_f$ 对 $\frj$-预进拓扑的分离完备化。

该拓扑的零邻域基为 $(\frj^p\mathrm R^nT(\mathbf A))\otimes_S S_f$，所以 $F^p(\mathrm R^nT(\mathbf A))\otimes_S S_f$ 也是零邻域基；后者是典范映射
\[
(\mathrm R^nT(\mathbf A))_f\longrightarrow(\mathrm R^nT(A_{p-1}))_f
\]
的核。因此，与 $\mathrm R^nT(\mathbf A)\otimes_S S_f$ 相伴的分离群可识别为 $\plim_k((\mathrm R^nT(A_k))_f)$ 的子群 $G$。

系统 $((\mathrm R^nT(A_k))_f)_k$ 显然满足（ML）；而 $(\mathrm R^nT(\mathbf A))_f$ 在每个 $(\mathrm R^nT(A_k))_f$ 中的像，等于 $h\geq k$ 充分大时各 $(\mathrm R^nT(A_h))_f$ 的共同像。所以 $G$ 在此射影极限中稠密；又该极限分离且完备，推论得证。
\end{proof}
\begin{flushright}（待续。）\end{flushright}\par

% ---- END c0_s13c.tex ----
% ---- BEGIN c3_s01a.tex ----

% END text/ch0_08_13.tex

\renewcommand{\newscope}{0-IV}

% BEGIN text/ch0_14_15a.tex
\NewNote{以下为原载第四卷第一部分的第0章补篇。原章首预告§§14--23；当前只有已译出的§§14--15收入正文。}
\source{\texttt{EGAIV.1.pdf}，PDF第4--12页，原书第5--13页}
\subsection*{原补篇内容提要与章首说明}
\noindent §14\quad 拓扑空间的组合维数。\\
§15\quad $M$-正则序列与 $\F$-正则序列。\\
§16\quad 诺特局部环中的维数与深度。\\
§17\quad 正则环。\\
§18\quad 代数扩张的补充知识。\\
§19\quad 形式光滑代数与 Cohen 环。\\
§20\quad 导子与微分。\\
§21\quad 特征 $p$ 环中的微分。\\
§22\quad 形式光滑性与正则性的微分判据。\\
§23\quad 日本环。

上述各节几乎全部用于阐述第IV章中将使用的交换代数概念。尽管这些概念中的相当一部分已见于若干部著作（[1]、[12]、[13]、[17]、[30]），我们仍认为，为读者提供一个连贯、且大致独立的阐述会更加方便。与使用概形语言的第IV章§§5、6、7合在一起，这些节在本论著内部构成一部小型专论，基本独立于第I至III章；它旨在连贯地表述那些在完备化、整闭包等运算下“表现良好”的环的性质，并把这些性质系统地联系到少数一般观念上。\orig{这里回顾：符号 $0_N$（分别地，$\mathrm{Err}_N$）指随第 $N$ 章刊行的第0章部分（分别地，勘误表）。}\orig{这里讨论的大多数性质由 Chevalley、Zariski、Nagata 和 Serre 发现。此处采用的方法最初由 A. Grothendieck 于1961年秋在 Harvard University 的课程中发展。}

\setcounter{section}{13}\section{拓扑空间的组合维数}
\subsection{拓扑空间的组合维数}
\egap{14.1.1}
设 $I$ 为有序集。按定义，$I$ 中元素的一条链是一个有限严格递增序列
$i_0<i_1<\cdots<i_n$（$n\geq0$），这条链的长度定义为 $n$。若 $X$ 为拓扑空间，它的不可约闭子集按包含关系构成有序集，从而得到 $X$ 的不可约闭子集链的概念。

\egas{14.1.2}{定义}
拓扑空间 $X$ 的\textbf{组合维数}（不致混淆时简称\textbf{维数}），记为 $\dim_c(X)$ 或 $\dim(X)$，是 $X$ 的不可约闭子集链长度的上确界。对于 $x\in X$，$X$ 在 $x$ 处的\textbf{组合维数}（简称\textbf{维数}），记为 $\dim_x(X)$，定义为
\[
 \dim_x(X)=\inf_U\dim(U),
\]
其中 $U$ 遍历 $x$ 在 $X$ 中的所有开邻域。
由定义，$\dim(\varnothing)=-\infty$（实数空集的上确界约定为 $-\infty$）。若 $(X_\alpha)$ 为 $X$ 的不可约分支族，则
\begin{equation}
 \dim(X)=\sup_\alpha\dim(X_\alpha).\tag{14.1.2.1}
\end{equation}
这是因为每条不可约闭子集链都包含在某个不可约分支中；反过来，各分支在 $X$ 中闭，所以分支内的不可约闭子集也是 $X$ 的不可约闭子集。

\egas{14.1.3}{定义}
若拓扑空间 $X$ 的所有不可约分支具有同一维数，则称 $X$ \textbf{等维}。由（14.1.2.1），这个共同维数等于 $\dim(X)$。

\egas{14.1.4}{命题}
（i）对于拓扑空间 $X$ 的任意子集 $Y$，有 $\dim(Y)\leq\dim(X)$。

（ii）若 $X$ 是有限个闭子集 $X_i$ 的并，则 $\dim(X)=\sup_i\dim(X_i)$。
\begin{proof}
对于 $Y$ 的每个不可约闭子集 $Z$，它在 $X$ 中的闭包 $\overline Z$ 不可约（$0_{\mathrm I}$，2.1.2），且 $\overline Z\cap Y=Z$，故得（i）。另一方面，若 $X=\bigcup_{i=1}^nX_i$ 且各 $X_i$ 闭，则 $X$ 的每个不可约闭子集都包含于某个 $X_i$（$0_{\mathrm I}$，2.1.1）；于是每条不可约闭子集链都包含于某个 $X_i$，得（ii）。
\end{proof}
由（i），对于每个 $x\in X$，还可写
\begin{equation}
 \dim_x(X)=\lim_U\dim(U),\tag{14.1.4.1}
\end{equation}
极限沿 $x$ 的开邻域的递减滤过系统取。

\egas{14.1.5}{推论}
设 $X$ 为拓扑空间，$x\in X$，$U$ 为 $x$ 的邻域；设 $Y_i$（$1\leq i\leq n$）为 $U$ 的闭子集，都包含 $x$，并且 $U=\bigcup_iY_i$。则
\begin{equation}
 \dim_x(X)=\sup_i\dim_x(Y_i).\tag{14.1.5.1}
\end{equation}
\begin{proof}
由（14.1.4，（ii）），$\dim_x(X)=\inf_V\sup_i\dim(Y_i\cap V)$，其中 $V$ 遍历包含于 $U$ 的 $x$ 的开邻域；同样，$\dim_x(Y_i)=\inf_V\dim(Y_i\cap V)$。若 $\sup_i\dim_x(Y_i)=+\infty$，结论显然；否则存在开邻域 $V_0\subset U$，使对所有 $i$ 及所有开邻域 $V\subset V_0$，有 $\dim(Y_i\cap V)=\dim_x(Y_i)$，于是结论成立。
\end{proof}

\egas{14.1.6}{命题}
对于每个拓扑空间 $X$，有 $\dim(X)=\sup_{x\in X}\dim_x(X)$。
\begin{proof}
定义及（14.1.4）说明对所有 $x$ 都有 $\dim_x(X)\leq\dim(X)$。反过来，取不可约闭子集链 $Z_0\subsetneq Z_1\subsetneq\cdots\subsetneq Z_n$，并取 $x\in Z_0$。对于包含 $x$ 的每个开集 $U$，各 $U\cap Z_i$ 都不可约（$0_{\mathrm I}$，2.1.6）且在 $U$ 中闭；它们在 $X$ 中的闭包分别为 $Z_i$，所以彼此不同。因此 $\dim(U)\geq n$，证明完成。
\end{proof}

\egas{14.1.7}{推论}
若 $(X_\alpha)$ 是 $X$ 的开覆盖，或局部有限闭覆盖，则
$\dim(X)=\sup_\alpha\dim(X_\alpha)$。
\begin{proof}
若 $X_\alpha$ 是 $x$ 的邻域，则 $\dim_x(X)\leq\dim(X_\alpha)$，得第一断言。若各 $X_\alpha$ 闭，且 $x$ 的邻域 $U$ 只与有限多个 $X_\alpha$ 相交，则由（14.1.4）
\[
 \dim_x(X)\leq\dim(U)=\sup_\alpha\dim(U\cap X_\alpha)
 \leq\sup_\alpha\dim(X_\alpha),
\]
得到第二断言。
\end{proof}

\egas{14.1.8}{推论}
设 $X$ 为诺特 Kolmogoroff 空间（$0_{\mathrm I}$，2.1.3），$F$ 为其闭点集。则
$\dim(X)=\sup_{x\in F}\dim_x(X)$。
\begin{proof}
在（14.1.6）的证明中，只须注意 $Z_0$ 中存在闭点（$0_{\mathrm I}$，2.1.3）。
\end{proof}

\egas{14.1.9}{命题}
设 $X$ 为非空诺特 Kolmogoroff 空间。则 $\dim(X)=0$ 当且仅当 $X$ 有限且离散。
\begin{proof}
若空间 $X$ 是 Hausdorff 的，特别地若它离散，则其不可约闭子集只能是单点，因而 $\dim(X)=0$。反之，若 $X$ 诺特、为 Kolmogoroff 空间且维数为0，则每个不可约分支都包含闭点（$0_{\mathrm I}$，2.1.3），故这个分支只能等于该闭点。$X$ 只有有限多个不可约分支，因此有限且离散。
\end{proof}

\egas{14.1.10}{推论}
设 $X$ 为诺特 Kolmogoroff 空间。点 $x\in X$ 孤立，当且仅当 $\dim_x(X)=0$。
\begin{proof}
必要性显然，且不需对 $X$ 作额外假设。充分性如下：存在 $x$ 的开邻域 $U$ 使 $\dim(U)=0$；$U$ 仍是诺特 Kolmogoroff 空间，故有限且离散。
\end{proof}

\egas{14.1.11}{命题}
函数 $x\mapsto\dim_x(X)$ 在 $X$ 上上半连续。
\begin{proof}
在其值为 $+\infty$ 的点处，上半连续性是显然的。设 $\dim_x(X)=n<+\infty$。由（14.1.4.1），存在 $x$ 的开邻域 $U_0$，使每个包含于 $U_0$ 的 $x$ 的开邻域 $U$ 都满足 $\dim(U)=n$。于是，对于每个 $y\in U_0$ 及其开邻域 $V\subset U_0$，有
$\dim(V)\leq\dim(U_0)=n$（14.1.4）；再由（14.1.4.1），$\dim_y(X)\leq n$。
\end{proof}

\egas{14.1.12}{备注}
若 $f:X\to Y$ 是拓扑空间之间的连续映射，可能有 $\dim(f(X))>\dim(X)$。例如，取 $X$ 为两点 $a,b$ 的离散空间；$Y$ 为同一个集合，闭集规定为 $\varnothing,\{a\},\{a,b\}$。恒等集合映射 $f:X\to Y$ 连续，而 $\dim(Y)=1$、$\dim(X)=0$。注意，$Y$ 是一个离散赋值环 $A$ 的谱，其中 $a$ 是唯一闭点，$b$ 是泛点。若 $K,k$ 分别为 $A$ 的分式域和剩余域，则 $X$ 是环 $k\times K$ 的谱，$f$ 对应于典范同态组成的环同态
\[
 (\varphi,\psi):A\longrightarrow k\times K,\qquad
 \varphi:A\to k,\quad\psi:A\to K
\]
（参见IV，5.4.3）。

\subsection{闭子集的余维数}
\egas{14.2.1}{定义}
给定拓扑空间 $X$ 的不可约闭子集 $Y$，称以 $Y$ 为最小项的 $X$ 的不可约闭子集链长度的上确界为 $Y$ 在 $X$ 中的\textbf{组合余维数}，简称\textbf{余维数}，记为 $\codim(Y,X)$。若 $Y$ 为任意闭子集，则其余维数定义为它的不可约分支在 $X$ 中余维数的下确界。若 $X$ 的所有极小不可约闭子集在 $X$ 中具有相同余维数，则称 $X$ \textbf{等余维}。

由此，$\codim(\varnothing,X)=+\infty$，因为实数空集的下确界约定为 $+\infty$。设 $(X_\alpha)$、$(Y_\beta)$ 分别为 $X,Y$ 的不可约分支族。每个 $Y_\beta$ 包含于某个 $X_\alpha$；更一般地，以 $Y_\beta$ 为最小项的链完全包含于某个 $X_\alpha$。于是
\begin{equation}
 \codim(Y,X)=\inf_\beta\sup_{\alpha:\,Y_\beta\subset X_\alpha}
                      \codim(Y_\beta,X_\alpha).\tag{14.2.1.1}
\end{equation}

\egas{14.2.2}{命题}
设 $X$ 为拓扑空间。

（i）若 $\Phi$ 为 $X$ 的不可约闭子集族，则
\begin{equation}\dim(X)=\sup_{Y\in\Phi}\codim(Y,X).\tag{14.2.2.1}\end{equation}
（ii）对于每个非空闭子集 $Y\subset X$，
\begin{equation}\dim(Y)+\codim(Y,X)\leq\dim(X).\tag{14.2.2.2}\end{equation}
（iii）若 $Y\subset Z\subset T$ 为 $X$ 的闭子集，则
\begin{equation}\codim(Y,Z)+\codim(Z,T)\leq\codim(Y,T).\tag{14.2.2.3}\end{equation}
（iv）闭子集 $Y$ 满足 $\codim(Y,X)=0$ 当且仅当它包含 $X$ 的某个不可约分支。
\begin{proof}
（i）、（iv）直接来自定义。对于（ii），可归结为 $Y$ 不可约，此时公式由定义（14.1.1）、（14.2.1）得出。对于（iii），按定义先可设 $Y$ 不可约；这时
$\codim(Y,Z)=\sup_\alpha\codim(Y,Z_\alpha)$，$Z_\alpha$ 遍历包含 $Y$ 的 $Z$ 的不可约分支。显然 $\codim(Y,T)\geq\codim(Y,Z)$，所以当 $\codim(Y,Z)=+\infty$ 时不等式成立。否则，存在 $\alpha$ 使该上确界在 $Z_\alpha$ 处达到；由（14.2.1），可再归结为 $Z$ 也不可约，而此时（14.2.2.3）直接来自定义。
\end{proof}

\egas{14.2.3}{命题}
设 $Y$ 为拓扑空间 $X$ 的闭子集。对于每个开集 $U\subset X$，有
\begin{equation}
 \codim(Y\cap U,U)\geq\codim(Y,X).\tag{14.2.3.1}
\end{equation}
此外，两边相等，当且仅当对与 $U$ 相交的 $Y$ 的不可约分支族 $(Y_\alpha)$，有
$\codim(Y,X)=\inf_\alpha\codim(Y_\alpha,X)$。
\begin{proof}
已知（$0_{\mathrm I}$，2.1.6），取闭包 $Z\mapsto\overline Z$ 给出 $U$ 的不可约闭子集与 $X$ 中和 $U$ 相交的不可约闭子集之间的双射；特别地，对应 $Y\cap U$ 的不可约分支与和 $U$ 相交的 $Y$ 的不可约分支。对于这样的 $Y_\alpha$，
$\codim(Y_\alpha,X)=\codim(Y_\alpha\cap U,U)$；结论由定义得出。
\end{proof}

\egas{14.2.4}{定义}
设 $Y$ 为 $X$ 的闭子集，$x\in X$。$Y$ 在 $X$ 中于 $x$ 处的余维数定义为
\[
 \codim_x(Y,X)=\sup_U\codim(Y\cap U,U),
\]
其中 $U$ 遍历 $x$ 的开邻域。由（14.2.3），也可写为
\begin{equation}
 \codim_x(Y,X)=\lim_U\codim(Y\cap U,U).\tag{14.2.4.1}
\end{equation}
此极限沿开邻域的递减滤过系统取。若 $x\in X\setminus Y$，则 $\codim_x(Y,X)=+\infty$。

\egas{14.2.5}{命题}
若 $(Y_i)_{1\leq i\leq n}$ 是 $X$ 的有限闭子集族，$Y=\bigcup_iY_i$，则
\begin{equation}
 \codim(Y,X)=\inf_i\codim(Y_i,X).\tag{14.2.5.1}
\end{equation}
\begin{proof}
每个 $Y_i$ 的不可约分支都包含于 $Y$ 的某个不可约分支；反过来，$Y$ 的每个不可约分支也都是某个 $Y_i$ 的不可约分支（$0_{\mathrm I}$，2.1.1）。所以结论来自余维数的定义与（14.2.2.3）。
\end{proof}

\egas{14.2.6}{推论}
设 $Y$ 为 $X$ 的局部诺特闭子空间。

（i）每个 $x\in X$ 只包含于 $Y$ 的有限多个不可约分支 $Y_i$（$1\leq i\leq n$）中，并且
$\codim_x(Y,X)=\inf_i\codim(Y_i,X)$。

（ii）函数 $x\mapsto\codim_x(Y,X)$ 在 $X$ 上下半连续。
\begin{proof}
按假设，有 $x$ 的开邻域 $U_0$，使 $Y\cap U_0$ 诺特，因而只有有限多个不可约分支；它们是 $Y$ 的分支在 $U_0$ 上的迹。所以包含 $x$ 的 $Y$ 的分支只有有限多个。缩小 $U_0$，可避开不包含 $x$ 的那些分支，从而设 $Y_i\cap U$ 正是 $Y\cap U$ 的所有不可约分支。对每个开邻域 $V\subset U$，$Y_i\cap V$ 仍是 $Y\cap V$ 的分支，且由（14.2.3），
$\codim(Y_i,X)=\codim(Y_i\cap V,V)$，得到（i）。对于每个 $x'\in V$，包含 $x'$ 的分支只可能是这些 $Y_i$ 中的一部分，因此
$\codim_{x'}(Y,X)\geq\codim_x(Y,X)$，得到（ii）。
\end{proof}

\subsection{链条件}
\egap{14.3.1}
在拓扑空间 $X$ 中，称不可约闭子集链
$Z_0\subsetneq Z_1\subsetneq\cdots\subsetneq Z_n$ 为\textbf{饱和链}，如果对于每个相邻的两项，不存在严格介于它们之间的不可约闭子集。

\egas{14.3.2}{命题}
设 $X$ 为拓扑空间，且每当不可约闭子集 $Y\subset Z$ 时都有 $\codim(Y,Z)<+\infty$。以下条件等价：

（a）具有相同端点的任意两条饱和不可约闭子集链长度相同。

（b）若 $Y\subset Z\subset T$ 是三个不可约闭子集，则
\begin{equation}
 \codim(Y,T)=\codim(Y,Z)+\codim(Z,T).\tag{14.3.2.1}
\end{equation}
\begin{proof}
（a）推出（b）是直接的。反过来，设（b）成立。证明两条具有相同端点、长度为 $m,n$（取 $n\leq m$）的饱和链必有 $m=n$。对 $n$ 归纳，$n=1$ 时由定义成立。设 $n>1$，并反设 $n<m$；取饱和链
$Z_0\subsetneq Z_1\subsetneq\cdots\subsetneq Z_n$，且存在端点 $Z_0,Z_n$、长度 $m$ 的另一饱和链。因为 $\codim(Z_0,Z_n)\geq m>n$，而 $\codim(Z_0,Z_1)=1$，由（b）得到
\[
 \codim(Z_1,Z_n)=\codim(Z_0,Z_n)-1>n-1,
\]
与归纳假设矛盾。
\end{proof}
满足上述等价条件时，称 $X$ 满足\textbf{链条件}，或称 $X$ 为\textbf{悬链空间}。悬链空间的每个闭子空间仍是悬链空间。

\egas{14.3.3}{命题〔原文；存在实质性错误〕}
设 $X$ 为有限维诺特 Kolmogoroff 空间。原文声称以下条件等价：\cmt{E01}{这一等价性有误。正确关系为（a）$\Longleftrightarrow$（c）$\Longleftrightarrow$（d）$\Longrightarrow$（b），而（b）不推出其余条件。反例、修正及文献核对见E01。}

（a）$X$ 的任意两条极大不可约闭子集链具有相同长度。

（b）$X$ 等维、等余维且悬链。

（c）$X$ 等维，且每当不可约闭子集 $Y\subset Z$ 时，
\begin{equation}\dim(Z)=\dim(Y)+\codim(Y,Z).\tag{14.3.3.1}\end{equation}
（d）$X$ 等余维，且每当不可约闭子集 $Y\subset Z$ 时，
\begin{equation}\codim(Y,X)=\codim(Y,Z)+\codim(Z,X).\tag{14.3.3.2}\end{equation}
\begin{proof}[原文证明]
对 $X$ 的假设说明，每条极大不可约闭子集链的两端必为闭点和不可约分支（$0_{\mathrm I}$，2.1.3）；而端点为 $Y,Z$（$Y\subset Z$）的每条饱和链都包含于一条极大链中，极大链新增的项不是包含于 $Y$，就是包含 $Z$。原文据此断言（a）与（b）的等价性立即成立。\cmt{E01}{问题正出在这里：不同闭点在不同不可约分支中的余维数不由“整个空间等余维”统一控制。}
这些说明还表明，（a）成立时，对于任意不可约闭子集 $Y$，
\begin{equation}\dim(Y)+\codim(Y,X)=\dim(X).\tag{14.3.3.3}\end{equation}
结合（14.3.2.1），即得（14.3.3.1）、（14.3.3.2）。反之，（14.3.3.1）推出（14.3.2.1），所以由（14.3.2）得到链条件；再取 $Y$ 为闭点 $x$、$Z$ 为包含它的不可约分支，得到 $\codim(\{x\},X)=\dim(Z)$，所以（c）推出（b）。同样，（14.3.3.2）推出链条件；取同样的 $Y,Z$，也得 $\codim(\{x\},X)=\dim(Z)$，从而（d）推出（b）（每个分支都包含闭点）。
\end{proof}
原文把有限维且满足“（14.3.3）的等价条件”的诺特 Kolmogoroff 空间称为\textbf{双等维}。\cmt{E01}{因上述条件并不等价，此定义必须明确采用（a）这一较强条件。只满足（b）应称“弱双等维”，参见E01。}

\egas{14.3.4}{推论}
设 $X$ 为诺特 Kolmogoroff 双等维空间。则对于任意闭点 $x$ 和任意不可约分支 $Z$，
\begin{equation}
 \dim(X)=\dim(Z)=\codim(\{x\},X)=\dim_x(X).\tag{14.3.4.1}
\end{equation}
\begin{proof}
最后一个等号来自以下事实：若
$Y_0=\{x\}\subsetneq Y_1\subsetneq\cdots\subsetneq Y_n$
为极大链，且 $U$ 为 $x$ 的邻域，则各 $U\cap Y_i$ 在 $U$ 中闭且不可约，因其在 $X$ 中的闭包为 $Y_i$ 而彼此不同；于是由（14.1.4），$\dim(U)=\dim(X)$。
\end{proof}

\egas{14.3.5}{推论}
设 $X$ 为诺特 Kolmogoroff 空间。若 $X$ 双等维，则它的任意不可约分支之并，以及任意不可约闭子集，也双等维。另对每个闭子集 $Y$，有\cmt{边界约定}{对非空 $Y$ 按有限维公式理解。若 $Y=\varnothing$，本节既有约定给出 $-\infty+(+\infty)$，此和未定义；因此这一等式还需明确排除空集，或另行制定不同的扩展算术约定。}
\begin{equation}
 \dim(Y)+\codim(Y,X)=\dim(X).\tag{14.3.5.1}
\end{equation}

\begin{proof}
每条不可约闭子集链都包含于某个不可约分支，故第一断言来自（14.3.3）。若 $X'$ 是不可约闭子集，它满足（14.3.3，（c）），得第二断言。最后，在（14.3.3）的证明中已经验证了 $Y$ 不可约时的（14.3.5.1）；若 $Y_i$（$1\leq i\leq m$）是 $Y$ 的不可约分支，则维数最大的分支也是在 $X$ 中余维数最小的分支，故该公式来自维数和余维数的定义。
\end{proof}

\egas{14.3.6}{备注}
读者会注意到，（14.3.2）的证明适用于任意有序集，其证明并未使用这些元素是拓扑空间不可约闭子集这一事实。（14.3.3）的证明也同样适用于满足以下条件的有序集 $E$：每个 $x\in E$ 之下都有一个极小元素 $y\leq x$，且元素链的长度有界。\cmt{E01}{关于（14.3.3）的这一概括仍受同一错误影响；本册的反例本来就是有限有序集。}

\section{正则序列}
\subsection{\texorpdfstring{$M$}{M}-正则序列与 \texorpdfstring{$M$}{M}-拟正则序列}
\egap{15.1.1}
设 $A$ 为环，$N$ 为 $A$-模。以下将分次模 $N\otimes_AA[T_1,\ldots,T_n]$ 简记为 $N[T_1,\ldots,T_n]$，其中 $N$ 取平凡分次，多项式环按总次数分次。设 $M$ 为 $A$-模，$f_i\in A$（$1\leq i\leq n$），$\jj=\sum_{i=1}^nf_iA$，赋予 $M$ 以 $\jj$-预进制滤过。定义零次分次满同态
\begin{equation}
 \varphi:(\operatorname{gr}_0(M))[T_1,\ldots,T_n]
 \longrightarrow\operatorname{gr}_\bullet(M)\tag{15.1.1.1}
\end{equation}
如下：若 $\xi$ 是 $x\in M$ 模 $\jj M$ 的类，$P[T_1,\ldots,T_n]$ 为次数 $k$ 的齐次多项式，则把 $\xi,P$ 对应到 $P(f_1,\ldots,f_n)x$ 模 $\jj^{k+1}M$ 的类。这个类显然只依赖 $\xi,P$。若 $S_k$ 为齐次次数 $k$ 的多项式组成的模，得到 $A$-线性映射
\[
 \varphi_k:\operatorname{gr}_0(M)\otimes_AS_k\longrightarrow\operatorname{gr}_k(M),
\]
它由 $\jj^kM$ 的定义而满射。以下称 $f\in A$ 在 $M$ 中为非零因子，如果数乘 $f_M:x\mapsto fx$ 单射。

\egas{15.1.2}{引理}
设 $M$ 为 $A$-模，$f\in A$，赋予 $M$ 以 $(fA)$-预进制滤过。若 $f$ 在 $M$ 中为非零因子，则典范同态
\begin{equation}
 \varphi:(M/fM)[T]\longrightarrow\operatorname{gr}_\bullet(M)\tag{15.1.2.1}
\end{equation}
双射。如果 $M$ 对 $(fA)$-预进制拓扑分离，则逆命题也成立。
\begin{proof}
这里 $\operatorname{gr}_k(M)=f^kM/f^{k+1}M$。$\varphi$ 单射意味着，对每个 $x,k$，$f^kx\in f^{k+1}M$ 推出 $x\in fM$；当 $f$ 是非零因子时显然如此。反过来，数乘 $f_M$ 在商上诱导 $\operatorname{gr}_k(M)\to\operatorname{gr}_{k+1}(M)$，从而诱导一个次数1的分次同态 $g$。$\varphi$ 单射亦可表述为 $g$ 单射；滤过分离时，由此得到 $f_M$ 单射（Bourbaki，\textit{Alg. comm.}，III，§2，8，定理1推论1）。
\end{proof}

\egas{15.1.3}{推论}
设 $A$ 为诺特环，$M$ 为有限型 $A$-模。则 $f$ 在 $M$ 中为非零因子，当且仅当 $\varphi$ 双射。
\begin{proof}
只需证明充分性。要证明 $f_M$ 单射，只须对每个素理想 $\mathfrak p$，证明 $M'=M_\mathfrak p$ 上由 $f'=f/1\in A'=A_\mathfrak p$ 定义的数乘单射（Bourbaki，\textit{Alg. comm.}，II，§3，4，定理1）。有
\[
 f'^kM'/f'^{k+1}M'=(f^kM/f^{k+1}M)\otimes_AA',
\]
且相应 $\varphi'$ 就是 $\varphi\otimes1_{A'}$；局部化正合，所以 $\varphi$ 单射时 $\varphi'$ 也单射（$0_{\mathrm I}$，1.3.2）。于是归结到 $A$ 为诺特局部环，极大理想记为 $\mm$。若 $f\notin\mm$，它可逆，无须再证。若 $f\in\mm$，则 $M$ 对 $(fA)$-预进制拓扑分离（$0_{\mathrm I}$，7.3.5），应用（15.1.2）即可。
\end{proof}

\egas{15.1.4}{定义}
若 $f$ 在 $M$ 中为非零因子，则称 $f$ \textbf{$M$-正则}；若（15.1.2.1）双射，则称 $f$ \textbf{$M$-拟正则}。

我们已经证明，每个 $M$-正则元素都 $M$-拟正则；若 $A$ 诺特且 $M$ 有限型，逆命题成立。不再假设 $M$ 有限型时，逆命题可能失败。例如，设 $A$ 为不是域的主理想整环，$K$ 为其分式域，$M=K/A$。对于任意 $f\ne0$，有 $fM=M$，所以（15.1.2.1）的两边都为0；原文接着说，“但是 $f$ 此时是 $M$ 中的零因子”。\cmt{E02}{最后一句必须限定 $f$ 为非零非单位。取 $f=1$ 时，数乘是恒等映射，当然单射；而若 $f$ 非单位，则 $1/f+A\ne0$，但被 $f$ 乘后为0。见E02。}

% END text/ch0_14_15a.tex


% BEGIN text/ch0_15b.tex
% 本次续译：0_IV 15.1.5--15.2.5。原文错误保留并以译校注说明。
\source{\texttt{EGAIV.1.pdf}，PDF第13--21页，原书第14--22页的相应部分}
\egap{15.1.5}
现在设 $M$ 带有滤过 $(M_k)$，且 $M_0=M$，并以
\[
 \gr_\bullet(M)=\bigoplus_{k\geq0}M_k/M_{k+1}
\]
表示相伴分次 $A$-模。沿用（15.1.1）的记号，对每个 $k\geq0$，置
\begin{equation}
 M'_k=M_k+\jj M_{k-1}+\cdots+\jj^{k-1}M_1+\jj^kM_0.
 \tag{15.1.5.1}
\end{equation}
显然 $(M'_k)$ 是 $M$ 的一个滤过。若 $M_k=\mathfrak q^kM$，其中 $\mathfrak q$ 为 $A$ 的理想，则 $(M'_k)$ 正是 $(\jj+\mathfrak q)$-预进制滤过。以 $\gr'_\bullet(M)$ 表示它的相伴分次模。仍用
$(\gr_\bullet(M)\otimes_A(A/\jj))[T_1,\ldots,T_n]$ 表示分次模的张量积
\[
 (\gr_\bullet(M)\otimes_A(A/\jj))\otimes_A A[T_1,\ldots,T_n]
\]
（II，2.1.2）。如下定义一个次数为0的分次满同态：
\begin{equation}
 \psi:(\gr_\bullet(M)\otimes_A(A/\jj))[T_1,\ldots,T_n]
 \longrightarrow\gr'_\bullet(M).
 \tag{15.1.5.2}
\end{equation}
若 $\alpha\in A/\jj$ 是 $a\in A$ 的类，$\xi\in M_k/M_{k+1}$ 是 $x\in M_k$ 的类，$P[T_1,\ldots,T_n]$ 为 $h$ 次齐次多项式，则令三元组 $(\alpha,\xi,P)$ 对应于 $P(f_1,\ldots,f_n)ax$ 模 $M'_{h+k+1}$ 的类。仍可立即验证，此元素只依赖于 $\xi,\alpha,P$，从而定义出（15.1.5.2）；由（15.1.5.1），它满射。取满足 $M_1=0$ 的滤过，就回到（15.1.1.1）的同态。

\egas{15.1.6}{引理}
沿用（15.1.5）的记号，设 $n=1$，记 $f=f_1$。若 $f$ 是 $\gr_\bullet(M)$-正则的，则（15.1.5.2）双射。
\begin{proof}
以 $Q_k,Q'_k$ 分别表示（15.1.5.2）两端的 $k$ 次部分。在 $Q_k$ 上取滤过
\begin{equation}
 (Q_k)_i=\sum_{j\leq k-i}
 \bigl(\gr_{k-j}(M)\otimes_A(A/fA)\bigr)T^j.
 \tag{15.1.6.1}
\end{equation}
于是 $(Q_k)_0=Q_k$、$(Q_k)_{k+1}=0$，且
\[
 \gr_i(Q_k)=\bigl(\gr_i(M)\otimes_A(A/fA)\bigr)T^{k-i}.
\]
给 $Q'_k$ 取滤过 $(Q'_k)_i=\psi((Q_k)_i)$。这些滤过有限，故只须证明由 $\psi_k$ 诱导的
$\psi_{ki}:\gr_i(Q_k)\to\gr_i(Q'_k)$ 单射（Bourbaki，\textit{交换代数}，第III章，§3，第8小节，定理1的推论1）。\cmt{B2-T05}{本页此处印为§3；本页末及其他相关引用指向第III章§2第8小节。这里只记录内部引文不一致，不以此判定数学结论错误。}
现在
\[
 \gr_i(Q_k)
 =\bigl((M_i/M_{i+1})\otimes_A(A/fA)\bigr)T^{k-i}
 =\bigl(M_i/(fM_i+M_{i+1})\bigr)T^{k-i},
\]
而 $(Q'_k)_{i+1}$ 是
$M_k+fM_{k-1}+\cdots+f^{k-i-1}M_{i+1}$ 在 $M'_k/M'_{k+1}$ 中的像。因此归结为证明：对 $x\in M_i$，关系
\begin{equation}
 f^{k-i}x\in M_k+fM_{k-1}+\cdots+f^{k-i-1}M_{i+1}+M'_{k+1}
 \tag{15.1.6.2}
\end{equation}
推出 $x\in fM_i+M_{i+1}$。

（15.1.6.2）的右端包含于 $M_{i+1}+M'_{k+1}$；由（15.1.5.1），又有
$M'_{k+1}\subset M_{i+1}+f^{k+1-i}M$。关于 $f$ 的假设说明，它对每个 $M/M_h$ 都是非零因子（Bourbaki，同处）。由
$f^{k-i}x\in f^{k+1-i}M+M_{i+1}$，在 $M/M_{i+1}$ 中考虑，可得某个 $y\in M$ 使
$x-fy\in M_{i+1}$。由于 $x\in M_i$，有 $fy\in M_i$；再在 $M/M_i$ 中考虑，得到 $y\in M_i$。所以 $x\in fM_i+M_{i+1}$，证毕。
\end{proof}

\egas{15.1.7}{定义}
设 $(f_i)_{1\leq i\leq n}$ 是 $A$ 中的一个序列。若对每个 $1\leq i\leq n$，$f_i$ 都是
\[
 M\Big/\Bigl(\sum_{j=1}^{i-1}f_jM\Bigr)
\]
的正则元素，则称此序列\textbf{$M$-正则}。若（15.1.1.1）的典范同态 $\varphi$ 双射，则称它\textbf{$M$-拟正则}。当 $M=A$ 时，简称“正则序列”和“拟正则序列”。\cmt{B2-H03}{原文的定义没有要求最终商模非零；因此后面（15.1.12）说 $(1,0)$ 正则，并非错误。不能静默加入其他文献中常见的非零商条件。}

\egas{15.1.8}{命题}
采用（15.1.5）的记号。若 $(f_i)_{1\leq i\leq n}$ 是 $\gr_\bullet(M)$-正则序列，则（15.1.5.2）的典范同态双射。
\begin{proof}
当 $n=1$ 时，这就是（15.1.6）。对 $n$ 归纳。置 $\mathfrak k=\sum_{i=1}^{n-1}f_iA$，故 $\jj=\mathfrak k+f_nA$，并令
\[
 M''_k=M_k+\mathfrak kM_{k-1}+\cdots+\mathfrak k^{k-1}M_1+\mathfrak k^kM_0.
\]
于是
\[
 M'_k=M''_k+f_nM''_{k-1}+\cdots+f_n^{k-1}M''_1+f_n^kM''_0.
\]
以 $\gr''_\bullet(M)$ 表示 $(M''_k)$ 的相伴分次模。在典范同构意义下，可以写
\[
\begin{split}
 &(\gr_\bullet(M)\otimes_A(A/\jj))[T_1,\ldots,T_n]\\
 &\quad\simeq
 (\gr_\bullet(M)\otimes_A(A/\mathfrak k))[T_1,\ldots,T_{n-1}]
 \otimes_A(A/f_nA)\otimes_A A[T_n].
\end{split}
\]
故（15.1.5.2）分解为
\begin{equation}
 (\gr''_\bullet(M)\otimes_A(A/f_nA))[T_n]
 \longrightarrow\gr'_\bullet(M)
 \tag{15.1.8.1}
\end{equation}
与下列同态的复合：
\[
\begin{split}
 &(\gr_\bullet(M)\otimes_A(A/\mathfrak k))[T_1,\ldots,T_{n-1}]
       \otimes_A(A/f_nA)\otimes_A A[T_n]\\
 &\xrightarrow{\ \psi'\otimes1\ }
 \gr''_\bullet(M)\otimes_A(A/f_nA)\otimes_A A[T_n].
\end{split}
\]
这里 $\psi'$ 是把（15.1.5.2）中的 $n,\jj,\gr'_\bullet(M)$ 分别换为 $n-1,\mathfrak k,\gr''_\bullet(M)$ 所得的映射。由归纳假设，$\psi'$ 双射，所以只须证明（15.1.8.1）双射。由（15.1.6），只需证明 $f_n$ 是 $\gr''_\bullet(M)$-正则的。

但 $\gr''_\bullet(M)$ 与 $(\gr_\bullet(M)\otimes_A(A/\mathfrak k))[T_1,\ldots,T_{n-1}]$ 等同。原文写道：由假设，乘以 $f_n$ 在 $\gr_\bullet(M)\otimes_A(A/\mathfrak k)$ 上“双射”；由于 $A[T_1,\ldots,T_{n-1}]$ 是自由 $A$-模，乘以 $f_n$ 在 $\gr''_\bullet(M)$ 上也“双射”。\cmt{B2-T01}{这两处“双射”均应为“单射”。正则性只给单射，而单射已足以应用（15.1.6）。例如 $A=k[x,y]$，序列 $(x,y)$，则乘 $y$ 在 $A/(x)=k[y]$ 上单射而不满射。}
\end{proof}

\egas{15.1.9}{命题}
设 $M$ 为 $A$-模，$(f_i)_{1\leq i\leq n}$ 为 $A$ 中的序列。若它 $M$-正则，则它 $M$-拟正则。若还假设
\[
 M,\quad M/f_1M,\quad\ldots,\quad
 M\Big/\Bigl(\sum_{j=1}^{n-1}f_jM\Bigr)
\]
对于 $\jj$-预进制拓扑都分离，其中 $\jj=\sum_{i=1}^n f_iA$，则逆命题也成立。
\begin{proof}[原文证明]
在（15.1.8）中取 $M_1=0$，则（15.1.5.2）回到（15.1.1.1），从而得到第一断言。反过来，假定 $\varphi$ 双射。置 $\mathfrak k=\sum_{i=1}^{n-1}f_iA$，以 $\gr'_\bullet(M)$ 表示 $\mathfrak k$-预进制滤过的相伴分次模。$\varphi$ 分解为
\begin{equation}
 (\gr'_\bullet(M)\otimes_A(A/f_nA))[T_n]
 \longrightarrow\gr_\bullet(M)
 \tag{15.1.9.1}
\end{equation}
以及 $\varphi'\otimes1$ 的复合，其中
\[
 \varphi':(\gr'_0(M))[T_1,\ldots,T_{n-1}]
       \longrightarrow\gr'_\bullet(M)
\]
是典范同态。原文接着说：由于这两个同态满射，若 $\varphi$ 双射，它们就都双射；“所以 $\varphi'$ 必然单射，从而双射，因为它满射”。\cmt{B2-G01}{从 $\varphi'\otimes1$ 单射到 $\varphi'$ 单射，不能只靠非忠实的张量化消去。报告保留这一证明缺步，并另给利用分离滤过和首项的完整补证；不据此否定命题本身。}

于是可对 $n$ 归纳，$n=1$ 的情形由（15.1.2）得到。在各模
\[
 M,\quad M/f_1M,\quad\ldots,\quad M/\Bigl(\sum_{i=1}^{n-2}f_iM\Bigr)
\]
上，$\mathfrak k$-预进制拓扑比 $\jj$-预进制拓扑更细，所以也分离。故 $(f_i)_{1\leq i\leq n-1}$ 是 $M$-正则序列。

另一方面，证明假设推出：若 $y\in M/\mathfrak kM$ 满足
$f_n^k y\in f_n^{k+1}(M/\mathfrak kM)$，则 $y\in f_n(M/\mathfrak kM)$。取 $y$ 的代表元 $x$，先证明 $f_n^k x\in\jj^{k+1}M$，再由（15.1.9.1）的单射性得到所需结论。事实上，假设给出
\[
 f_n^k x\in f_n^{k+1}M+\mathfrak kM\subset\jj M;
\]
由于 $\varphi$ 单射，得到 $f_n^{k-1}x\in\jj M$。如此递推，得到 $x\in\jj M$，最终得到 $f_n^kx\in\jj^{k+1}M$。

像（15.1.2）一样结束论证：在 $M/\mathfrak kM$ 上，$\jj$-预进制滤过与 $(f_nA)$-预进制滤过相同，所以按假设分离；因此乘以 $f_n$ 的同态单射。这就表明整个序列 $(f_i)_{1\leq i\leq n}$ 是 $M$-正则的。
\end{proof}

\egas{15.1.10}{推论}
采用（15.1.9）的记号。假设对每个有限型 $A$-模 $N$，$M\otimes_A N$ 对于 $\jj$-预进制拓扑都分离，则 $(f_i)$ 是 $M$-正则序列，当且仅当它是 $M$-拟正则序列。

注意，这个假设还推出：若 $(f_i)_{1\leq i\leq n}$ 是 $M$-正则序列，则对区间 $[1,n]$ 的任意排列 $\sigma$，$(f_{\sigma(i)})$ 也 $M$-正则。

\egas{15.1.11}{推论}
设 $A$ 是诺特环，$M$ 是有限型 $A$-模，$(f_i)_{1\leq i\leq n}$ 是由 $A$ 的根基中的元素组成的有限序列。则它 $M$-正则，当且仅当它 $M$-拟正则。
\begin{proof}
这是（15.1.10）的特例，因为此时每个有限型 $A$-模对于 $\jj$-预进制拓扑都分离（$0_{\mathrm I}$，7.3.5）。
\end{proof}

\egas{15.1.12}{备注}
（i）当 $A$ 是诺特局部环、$M$ 为有限型 $A$-模且所有 $f_i$ 属于 $A$ 的极大理想时，J.-P. Serre 曾以“$M$-序列”之名引入 $M$-正则序列的概念［17］。

（ii）即使 $A$ 诺特且 $M=A$，两个元素组成的拟正则序列 $(f,g)$ 也未必正则。例如，$(0,1)$ 是拟正则的，因为 $\jj=0A+1A=A$，故 $\gr_\bullet(A)=0$；若 $A\ne0$，它不正则。不过 $(1,0)$ 正则，因为 $1$ 在 $A$ 中不是零因子，而 $A/A=0$，故 $0$ 在这个零模中不是零因子（参见15.1.1）。

（iii）在（15.1.6）中，一般不能把“$f$ 是 $\gr_\bullet(M)$-正则的”换成“$f$ 是 $\gr_\bullet(M)$-拟正则的”。为此，回到（15.1.4）的例子：$A$ 为主理想整环，其分式域为 $K\ne A$。取 $M=K$，并令滤过为
\[
 M_0=M=K,\qquad M_1=A,\qquad M_k=0\quad(k\geq2).
\]
于是 $\gr_0(M)=K/A$、$\gr_1(M)=A$、$\gr_k(M)=0$（$k\geq2$）。原文取 $0\ne f\in A$，说已知 $f$ 并非 $\gr_\bullet(M)$-正则；\cmt{B2-T02}{与前批E02相同，这里须取非零非单位 $f$。若 $f$ 为单位，它对任何模都正则，而且后文（15.1.5.2）的源也为零，不能构成所说反例。这是同一遗漏的再次出现。}
但因为
\[
 f^k\gr_\bullet(M)/f^{k+1}\gr_\bullet(M)=f^kA/f^{k+1}A,
\]
立即可见（15.1.2.1）（将其中的 $M$ 换为 $\gr_\bullet(M)$）单射，所以 $f$ 是拟正则的。另一方面，采用（15.1.6）的记号，有 $M'_k=K$（$k\geq0$），故 $\gr'_\bullet(M)=0$，而（15.1.5.2）左端不为零。

\egas{15.1.13}{命题}
设 $A$ 为环，$M,N$ 为两个 $A$-模，$(f_i)_{1\leq i\leq n}$ 为 $A$ 中的序列。若 $(f_i)$ 是 $M$-正则序列且 $N$ 平坦，则它是 $M\otimes_A N$-正则序列。如果 $N$ 忠实平坦，则逆命题也成立。
\begin{proof}
无需对 $N$ 作假设，就有
\[
 (M\otimes_A N)\Big/\sum_{j=1}^{i-1}f_j(M\otimes_A N)
 \simeq\Bigl(M/\sum_{j=1}^{i-1}f_jM\Bigr)\otimes_A N.
\]
对 $M/(\sum_{j=1}^{i-1}f_jM)$ 上乘以 $f_i$ 的同态应用（$0_{\mathrm I}$，6.1.1、6.4.1），即得结论。
\end{proof}

\egas{15.1.14}{推论}
设 $A$ 为环，$M$ 为 $A$-模，$(f_i)$ 为 $A$ 中的有限序列，$\rho:A\to B$ 为环同态。置
$M'=M\otimes_A B$、$f'_i=\rho(f_i)$。若 $B$ 是平坦 $A$-模且 $(f_i)$ 是 $M$-正则序列，则 $(f'_i)$ 是 $M'$-正则序列；若 $B$ 忠实平坦，逆命题也成立。
\begin{proof}
$M'/(\sum_{j=1}^{i-1}f'_jM')$ 与 $(M/(\sum_{j=1}^{i-1}f_jM))\otimes_A B$ 等同，前一 $B$-模上乘以 $f'_i$ 的同态与后一 $A$-模上乘以 $f_i$ 的同态等同。所以结论由（15.1.13）得出。
\end{proof}

\egas{15.1.15}{命题〔按底本保留〕}
设 $A$ 为环，$B$ 为交换 $A$-代数，$(f_i)_{1\leq i\leq n}$ 为 $B$ 中的有限序列，$M$ 为 $B$-模。假设序列 $M$-正则，且各 $A$-模
\[
 M\Big/\Bigl(\sum_{j=1}^{i-1}f_jM\Bigr)\qquad(1\leq i\leq n)
\]
平坦。那么对任意环同态 $\rho:A\to A'$，置
\[
 B'=B\otimes_A A',\quad M'=M\otimes_A A',\quad
 f'_i=f_i\otimes1\quad(1\leq i\leq n),
\]
序列 $(f'_i)$ 是 $M'$-正则的，且各 $A'$-模
$M'/(\sum_{j=1}^{i-1}f'_jM')$（$1\leq i\leq n$）平坦。\cmt{B2-E01}{底本两处求和上限确为 $i-1$，范围确为 $1\leq i\leq n$，遗漏了最终商的平坦性。取 $A=B=M=k[t]$、$n=1$、$f_1=t$、$A'=k$，满足所列假设，但 $f'_1=0$ 不是 $M'=k$ 的正则元素。报告给出保留全部前缀商（包括最终商）平坦的修订版本及证明。}
\begin{proof}[原文证明]
考虑按假设正合的列
\[
 0\longrightarrow M\xrightarrow{f_1}M\longrightarrow M/f_1M\longrightarrow0,
\]
其中中间箭头是乘以 $f_1$ 的同态。原文使用“$M/f_1M$ 是平坦 $A$-模”的假设，推出
\[
 0\longrightarrow M\otimes_A A'\xrightarrow{f_1\otimes1}
 M\otimes_A A'\longrightarrow(M/f_1M)\otimes_A A'\longrightarrow0
\]
正合（$0_{\mathrm I}$，6.1.2），从而乘以 $f'_1$ 在 $M'$ 上单射。

再令 $M_i=M/(\sum_{j=1}^i f_jM)$、$M'_i=M'/(\sum_{j=1}^i f'_jM')$（$1\leq i\leq n-1$）。有
\[
 M'_i=M_i\otimes_A A',\qquad
 M_i/f_{i+1}M_i=M_{i+1},\qquad
 M'_i/f'_{i+1}M'_i=M'_{i+1}.
\]
把以上论证中的 $M,f_1$ 换为 $M_i,f_{i+1}$，由 $M_i/f_{i+1}M_i$ 的平坦性，得到乘以 $f'_{i+1}$ 在 $M'_i$ 上单射。最后关于平坦性的断言由（$0_{\mathrm I}$，6.2.1）得出。
\end{proof}

\egas{15.1.16}{命题}
设 $A,B$ 是诺特局部环，$k$ 为 $A$ 的剩余域，$\varphi:A\to B$ 为局部同态，$M$ 为有限型 $B$-模。设 $(f_i)_{1\leq i\leq n}$ 是 $B$ 的极大理想中的元素序列，$g_i$ 为 $f_i$ 在 $B\otimes_A k$ 中的像。以下条件等价：
\begin{enumerate}[label=\alph*)]
\item $(f_i)$ 是 $M$-正则序列，且 $M_i=M/(\sum_{j=1}^i f_jM)$ 对 $1\leq i\leq n$ 均为平坦 $A$-模；
\item $(f_i)$ 是 $M$-正则序列，且 $M_n$ 是平坦 $A$-模；
\item $M$ 是平坦 $A$-模，且 $(g_i)$ 是 $M\otimes_A k$-正则序列；
\item $M$ 是平坦 $A$-模，且对任意环同态 $\rho:A\to A'$，若置 $B'=B\otimes_A A'$、$M'=M\otimes_A A'$、$f'_i=f_i\otimes1$，则 $(f'_i)$ 是 $M'$-正则序列。
\end{enumerate}
\begin{proof}
显然a）推出b）。为证明b）推出d），由（15.1.15），只须证明 $M_i$ 对 $0\leq i\leq n-1$ 也平坦，其中 $M_0=M$。\cmt{B2-E01}{这里按报告给出的（15.1.15）修订版使用：以下递降论证连同b）确实给出所有 $M_0,\ldots,M_n$ 的平坦性。因此前一条的印误不构成本命题的反例。}
由 $M_{i+1}=M_i/f_{i+1}M_i$，作递降归纳，归结为：若 $h$ 属于 $B$ 的极大理想，乘以 $h$ 在 $M$ 上单射，且 $M/hM$ 对 $A$ 平坦，则 $M$ 对 $A$ 平坦。这正是（$0_{\mathrm{III}}$，10.2.7）。

c）是d）取 $A'=k$ 的特例。最后证明c）推出a）。$M$ 是平坦 $A$-模及有限型 $B$-模，由c）及（$0_{\mathrm{III}}$，10.2.4），乘以 $f_1$ 的同态单射，且余核 $M_1=M/f_1M$ 对 $A$ 平坦。对 $i$ 归纳，假定 $M_i$ 对 $A$ 平坦；按假设，乘以 $g_{i+1}$ 在 $M_i\otimes_A k$ 上单射，再用同一结果，乘以 $f_{i+1}$ 在 $M_i$ 上单射，且余核 $M_{i+1}$ 对 $A$ 平坦，证明完成。注意c）推出a）的证明并未使用 $f_i$ 属于 $B$ 的极大理想这一条件。
\end{proof}

\egas{15.1.17}{推论}
$A,B$ 满足（15.1.16）的假设。设 $M$ 是对 $A$ 平坦的有限型 $B$-模，$R\subset M$ 为 $B$-子模，且 $N=M/R$ 对 $A$ 平坦。设 $(f_i)_{1\leq i\leq n}$ 是 $B$ 的极大理想中的元素，满足：
\begin{enumerate}[label=(\roman*)]
\item $f_iM\subset R$；
\item 若 $g_i$ 为 $f_i$ 在 $B\otimes_A k$ 中的像，则 $(g_i)$ 是 $M\otimes_A k$-正则序列；
\item 子模 $g_i(M\otimes_A k)$ 的和，等于 $R\otimes_A k$ 在 $M\otimes_A k$ 中的典范像，即典范同态 $M\otimes_A k\to N\otimes_A k$ 的核。
\end{enumerate}
那么 $(f_i)$ 是 $M$-正则序列，且 $R=\sum_{i=1}^n f_iM$。
\begin{proof}
由于 $N$ 对 $A$ 平坦，由 $0\to R\to M\to N\to0$ 得到正合列
\[
 0\longrightarrow R\otimes_A k\longrightarrow M\otimes_A k
 \longrightarrow N\otimes_A k\longrightarrow0
\]
（$0_{\mathrm I}$，6.1.2），故 $R\otimes_A k=\sum_i g_i(M\otimes_A k)$。$M$ 有限型且 $B$ 诺特，所以 $R$ 也是有限型 $B$-模。原文由 Nakayama 引理以及 $f_i$ 属于 $B$ 的极大理想，推出 $R=\sum_i f_iM$。\cmt{B2-H04}{更精确地，对有限型 $B$-模 $R/(\sum_i f_iM)$ 应用Nakayama，所用理想为 $\mm_A B\subset\mm_B$，包含关系来自局部同态。正则性则由（15.1.16）的c）推出a）。这补明所用对象，不是否定推论。}
\end{proof}

\egas{15.1.18}{引理}
设 $M$ 为 $A$-模，带有分离且穷尽的滤过 $(M_k)_{k\geq0}$；设 $f\in A$ 是 $\gr_\bullet(M)$-正则元素。则：
\begin{enumerate}[label=(\roman*)]
\item $f$ 是 $M$-正则的；
\item 对每个 $k$，有 $M_k\cap fM=fM_k$。因此 $M_k$ 在 $M'=M/fM$ 中的典范像 $M'_k$ 典范同构于 $M_k/fM_k$；
\item 列
\[
 0\longrightarrow M_{k+1}/fM_{k+1}\longrightarrow M_k/fM_k
 \longrightarrow\gr_k(M)/f\gr_k(M)\longrightarrow0
\]
正合。给 $M'$ 取商滤过 $(M'_k)$ 后，$\gr_\bullet(M')$ 典范同构于 $\gr_\bullet(M)/f\gr_\bullet(M)$。
\end{enumerate}
\begin{proof}
（i）由假设及Bourbaki，\textit{交换代数}第III章§2第8小节定理1的推论1得出。由于 $\gr_\bullet(M/M_k)$ 是 $\gr_\bullet(M)$ 的直和因子，$f$ 对每个 $\gr_\bullet(M/M_k)$ 都正则，由（i），它对 $M/M_k$ 正则，故得到（ii）。为证（iii），只须验证 $M_{k+1}/fM_{k+1}\to M_k/fM_k$ 单射：若 $x\in M_{k+1}$ 且 $x\in fM_k$，由（ii）有 $x\in fM_{k+1}$，引理得证。
\end{proof}

\egas{15.1.19}{命题}
设 $M$ 为 $A$-模，带有满足 $M_0=M$ 的滤过 $(M_k)_{k\geq0}$，并设 $(f_i)_{1\leq i\leq n}$ 为 $\gr_\bullet(M)$-正则序列。另假设对 $0\leq i\leq n-1$，$(M_k)$ 在
$M/(\sum_{j\leq i}f_jM)$ 上诱导的商滤过分离。那么 $(f_i)$ 是 $M$-正则序列；若置 $M'=M/(\sum_{i\leq n}f_iM)$ 并赋予商滤过，则
\[
 \gr_\bullet(M')\simeq
 \gr_\bullet(M)\Big/\Bigl(\sum_{i\leq n}f_i\gr_\bullet(M)\Bigr)
\]
典范同构。
\begin{proof}
$n=1$ 时由（15.1.18）得到。对 $n$ 归纳，令
$N=M/(\sum_{i\leq n-1}f_iM)$，$(N_k)$ 为商滤过；它按假设分离，且
\[
 \gr_\bullet(N)\simeq
 \gr_\bullet(M)\Big/\Bigl(\sum_{i\leq n-1}f_i\gr_\bullet(M)\Bigr).
\]
按假设，$f_n$ 对 $\gr_\bullet(N)$ 正则。将（15.1.18）用于 $N$，得到 $f_n$ 对 $N$ 正则，且 $\gr_\bullet(M')\simeq\gr_\bullet(N)/f_n\gr_\bullet(N)$，从而证明命题。
\end{proof}
特别地，以下任一条件下可应用（15.1.19）：一、滤过 $(M_k)$ 有限且分离，因为这时 $k$ 足够大就有 $M_k=0$；二、$A$ 诺特，$\jj$ 是包含于 $A$ 根基的理想，$M$ 有限型，而 $(M_k)$ 是 $\jj$-预进制滤过（$0_{\mathrm I}$，7.3.5）。

\egas{15.1.20}{推论}
设 $A$ 为环，$M$ 为 $A$-模，$(f_i)_{1\leq i\leq n}$ 为 $M$-正则序列。则对任意整数 $\nu_i\geq1$，$(f_i^{\nu_i})$ 也是 $M$-正则序列，且
$M/(\sum_{i\leq n}f_i^{\nu_i}M)$ 有一个子模列，其逐次商都同构于 $M/(\sum_{i\leq n}f_iM)$。\cmt{B2-H05}{法文此处为 \textit{suite de composition}，在本句中是有限子模滤过，不要求逐次商为单模；不可按“单模商组成的合成列”的狭义定义误读。}
\begin{proof}
对 $n$ 归纳。$n=1$ 时，由乘以 $f_1$ 单射，有 $f_1^kM/f_1^{k+1}M\simeq M/f_1M$，断言立即成立。令
\[
 N=M\Big/\Bigl(\sum_{i\leq n-1}f_i^{\nu_i}M\Bigr).
\]
由归纳假设，$N$ 有一个有限滤过 $(N_k)$，各非零逐次商 $\gr_k(N)$ 同构于 $M/(\sum_{i\leq n-1}f_iM)$。按假设 $f_n$ 对 $\gr_\bullet(N)$ 正则，故由（15.1.19）它对 $N$ 正则。再用已证的单元素情形，$f_n^{\nu_n}$ 对 $N$ 正则，且 $N/f_n^{\nu_n}N$ 有逐次商同构于 $N/f_nN$ 的子模列。对于由 $(N_k)$ 诱导的 $N/f_nN$ 的滤过，由（15.1.19），各分次项为
\[
 \gr_k(N/f_nN)\simeq\gr_k(N)/f_n\gr_k(N)
 \simeq M\Big/\Bigl(\sum_{i\leq n}f_iM\Bigr),
\]
推论得证。
\end{proof}

\egas{15.1.21}{命题}
设 $A,B$ 为诺特局部环，$\varphi:A\to B$ 为局部同态，$M$ 为有限型 $B$-模，$(f_i)_{1\leq i\leq n}$ 为 $A$ 的极大理想中的一个 $A$-正则序列。以下条件等价：
\begin{enumerate}[label=\alph*)]
\item $M$ 是平坦 $A$-模；
\item 若 $\jj=\sum_{i\leq n}f_iA$，则 $M/\jj M$ 是平坦 $A/\jj$-模，且 $(f_i)$ 为 $M$-正则序列。
\end{enumerate}
\begin{proof}
a）推出b），由（15.1.13）及（$0_{\mathrm{III}}$，10.2.2）得到。反过来，$M$-正则性推出 $M$-拟正则性（15.1.9），所以（15.1.1.1）双射。序列还假定为 $A$-正则，因此它也 $A$-拟正则，且相应典范同态
$(A/\jj)[T_1,\ldots,T_n]\to\gr_\bullet(A)$ 双射。这表明典范同态（$0_{\mathrm{III}}$，10.1.1.2）
\[
 \gr_0(M)\otimes_A\gr_\bullet(A)\longrightarrow\gr_\bullet(M)
\]
双射。因 $f_i$ 属于 $A$ 的极大理想且 $\varphi$ 局部，由（$0_{\mathrm{III}}$，10.2.2）得到结论。
\end{proof}

\subsection{\texorpdfstring{$\F$}{F}-正则序列}
\egap{15.2.1}
设 $X$ 为赋环空间，$\F$ 为 $\OO_X$-模，$f_i$（$1\leq i\leq n$）为 $\OO_X$ 在 $X$ 上的截面，$\J$ 为它们生成的理想层，即它们所典范定义的同态 $\OO_X^n\to\OO_X$ 的像（$0_{\mathrm I}$，5.1.1）。子模 $\J^k\F$ 给出滤过，仍令
$\gr_k(\F)=\J^k\F/\J^{k+1}\F$。设 $S_k$ 为 $\mathbf Z[T_1,\ldots,T_n]$ 中的 $k$ 次齐次多项式群，把它看作 $X$ 上的常层。如（15.1.1）一样，定义典范同态
\begin{equation}
 \varphi_k:\gr_0(\F)\otimes_{\mathbf Z}S_k\longrightarrow\gr_k(\F).
 \tag{15.2.1.1}
\end{equation}
具体地，对每个开集 $U$（或拓扑的一组基中的开集），考虑由截面 $f_i|_U$ 定义的（15.1.1.1）同态
\[
 \varphi_{k,U}:\gr_0(\Gamma(U,\F))\otimes_{\mathbf Z}S_k
       \longrightarrow\gr_k(\Gamma(U,\F)),
\]
其中 $\Gamma(U,\F)$ 按子模 $(\Gamma(U,\J))^k\Gamma(U,\F)$ 滤过。
每个 $\gr_k(\F)$ 是预层 $U\mapsto\gr_k(\Gamma(U,\F))$ 的相伴层，因而各 $\varphi_{k,U}$ 给出（15.2.1.1）。由模的滤过归纳极限的正合性及其与张量积的交换性，对每个 $x\in X$，该同态在 $x$ 处的茎同态与由 $((f_i)_x)_{1\leq i\leq n}$ 定义的同态
\begin{equation}
 \gr_0(\F_x)\otimes_{\mathbf Z}S_k\longrightarrow\gr_k(\F_x)
 \tag{15.2.1.2}
\end{equation}
等同。

\egap{15.2.2}
采用（15.2.1）的记号，以 $f_i\F$ 表示由 $f_i$ 定义的 $\F$ 的自同态的像。若由 $f_i$ 定义的
$\F/(\sum_{j=1}^{i-1}f_j\F)$ 的自同态对所有 $1\leq i\leq n$ 都单射，则称 $(f_i)$ 为\textbf{$\F$-正则序列}。若（15.2.1.1）的各同态单射，则称它为\textbf{$\F$-拟正则序列}。

由定义及（15.2.1），序列 $\F$-正则（分别地，$\F$-拟正则），当且仅当对每个 $x\in X$，其茎序列 $((f_i)_x)$ 是 $\F_x$-正则（分别地，$\F_x$-拟正则）的。于是每个正则序列均拟正则（15.1.9）；若 $\F$ 有限型、每个 $\OO_{X,x}$ 都诺特，且所有 $(f_i)_x$ 都属于 $\OO_{X,x}$ 的根基，则逆命题也成立（15.1.11）。

若 $X=\Spec(A)$ 仿射且 $\F=\widetilde M$，则由（I，1.3），$(f_i)$ 为 $\F$-正则（分别地，$\F$-拟正则）序列，等价于它为 $M$-正则（分别地，$M$-拟正则）序列。当 $\F=\OO_X$ 时，省略定义中的 $\F$。

\egas{15.2.3}{备注}
当 $X=\Spec(A)$、$\jj=\sum_{i=1}^n f_iA$ 时，在不与 $V(\jj)$ 相交的每个开集 $U$ 上，对于任意 $A$-模 $M$，该序列都拟正则。因为对每个 $x\in U$，有 $\J_x=\OO_{X,x}$，所以（15.2.1.2）两端都为零。因此，原文说，模层的正则序列概念只在 $V(\jj)$ 的点的邻域中才有意义。\cmt{B2-H06}{前一句直接证明的是“拟正则”在此开集上自动成立。若把末句读成“正则性在该开集上也自动成立”，则不对：原文（15.1.12）的 $(0,1)$ 已是反例。末句应按研究兴趣的说明阅读。}

\egas{15.2.4}{命题}
设 $X$ 为赋环空间，$\F$ 为凝聚 $\OO_X$-模，$x\in X$，$f_i\in\Gamma(X,\OO_X)$（$1\leq i\leq n$）。若 $((f_i)_x)$ 为 $\F_x$-正则序列，则存在 $x$ 的开邻域 $U$，使 $(f_i|_U)$ 为 $(\F|_U)$-正则序列。
\begin{proof}
对每个 $i$，$\F/(\sum_{j=1}^{i-1}f_j\F)$ 凝聚（$0_{\mathrm I}$，5.3.4），故由 $f_i$ 定义的自同态的核也凝聚。该核在 $x$ 处的茎为零，因而在 $x$ 的某个邻域上的限制为零（$0_{\mathrm I}$，5.2.2），结论得证。
\end{proof}

\egas{15.2.5}{命题}
设 $u=(\psi,\theta):X\to Y$ 为赋环空间的平坦态射（$0_{\mathrm I}$，6.7.1），$\F$ 为 $\OO_Y$-模，$\G=u^*\F$；设 $f_i\in\Gamma(Y,\OO_Y)$、$g_i=\theta(f_i)\in\Gamma(Y,\psi_*\OO_X)=\Gamma(X,\OO_X)$。如果 $(f_i)$ 是 $\F$-正则序列，则 $(g_i)$ 是 $\G$-正则序列；若 $u$ 忠实平坦，则逆命题也成立（$0_{\mathrm I}$，6.7.8）。
\begin{proof}
令 $x\in X$、$y=u(x)=\psi(x)$，则
\[
 \G_x=\F_{u(x)}\otimes_{\OO_{Y,u(x)}}\OO_{X,x},\qquad
 (g_i)_x=\theta_x^\sharp((f_i)_{u(x)}).
\]
按假设，$\theta_x^\sharp$ 使 $\OO_{X,x}$ 成为平坦 $\OO_{Y,u(x)}$-模，且取茎函子正合；第一断言由（15.1.14）的相应结论得出。第二断言同理，因为 $u$ 忠实平坦表示 $\psi$ 满射，且各 $\theta_x^\sharp$ 使 $\OO_{X,x}$ 成为忠实平坦的 $\OO_{Y,u(x)}$-模。
\end{proof}



% END text/ch0_15b.tex


% BEGIN continuations/iv0_16_12.tex
% Continuation from the uploaded EGAIV.1.pdf, PDF pp.21--27 (printed pp.22--28).
\setcounter{section}{15}
\section{诺特局部环中的维数与深度}
\renewcommand{\newscope}{0-IV}
\subsection{环的维数}
\source{EGA IV 第一部分，第0章，印刷第22--25页；上传PDF第21--24页}

\egap{16.1.1}
环 $A$ 的\textbf{维数}（或\textbf{Krull维数}），记作 $\dim A$，是其谱 $\Spec(A)$ 的（组合）维数（14.2.1）。\cmt{D4-T01}{这里所引原号为（14.2.1），该处定义的是余维数；组合维数的定义在（14.1.2）。正文保留原引号。}
由于 $\Spec(A)$ 的不可约闭集为 $V(\pp)$，其中 $\pp$ 为 $A$ 的素理想（Bourbaki，\textit{交换代数}，第二章，§4，第3目，命题14的推论1），也就是说，$\dim A$ 是 $A$ 的素理想链
\[
 \pp_0\subsetneq\pp_1\subsetneq\cdots\subsetneq\pp_m
\]
的长度（14.1.1）的上确界。若 $A\ne0$，则 $\dim A\geq0$。非零Artin环（特别地，域）的维数为零（Bourbaki，\textit{交换代数}，第四章，§2，第5目，命题7）；不是域的Dedekind环（特别地，$\mathbf Z$ 或域 $k$ 上的单变元多项式环 $k[T]$）的维数为1（同书，第七章，§2）。诺特环不一定具有有限维数[30]。
若 $(\pp_\alpha)$ 为 $A$ 的极小素理想族，则由（14.1.2.1），
\begin{equation}
 \dim A=\sup_\alpha\dim(A/\pp_\alpha).
 \tag{16.1.1.1}
\end{equation}

\egap{16.1.2}
对于 $A$ 的任意理想 $\mathfrak a$，$A/\mathfrak a$ 的素理想与 $A$ 中包含 $\mathfrak a$ 的素理想一一对应，故
\begin{equation}
 \dim(A/\mathfrak a)\leq\dim A.
 \tag{16.1.2.1}
\end{equation}
若进一步假定 $\mathfrak a$ 不包含于 $A$ 的任何极小素理想 $\pp_\alpha$，则任意一条由包含 $\mathfrak a$ 的素理想组成的链，都可添入某个 $\pp_\alpha$ 而加长，因而
\begin{equation}
 1+\dim(A/\mathfrak a)\leq\dim A.
 \tag{16.1.2.2}
\end{equation}

\egap{16.1.3}
若 $S$ 是 $A$ 的乘法子集，$S^{-1}A$ 的素理想与 $A$ 中不交于 $S$ 的素理想一一对应，故
\begin{equation}
 \dim(S^{-1}A)\leq\dim A.
 \tag{16.1.3.1}
\end{equation}
对 $A$ 的每个素理想 $\pp$，按定义（14.2.1），
\begin{equation}
 \dim A_{\pp}=\codim(V(\pp),\Spec A).
 \tag{16.1.3.2}
\end{equation}
这个数也称为素理想 $\pp$ 的\textbf{高度}，记作 $\htp(\pp)$。更一般地，对于 $A$ 的理想 $\mathfrak a$，称
\begin{equation}
 \htp(\mathfrak a)=\codim(V(\mathfrak a),\Spec A)
 \tag{16.1.3.3}
\end{equation}
为 $\mathfrak a$ 的高度。若 $(\pp_\lambda)$ 为包含 $\mathfrak a$ 的素理想中所有极小者组成的族，则由（14.2.1），
\begin{equation}
 \htp(\mathfrak a)=\inf_\lambda\htp(\pp_\lambda).
 \tag{16.1.3.4}
\end{equation}
并且有
\begin{equation}
 \dim A=\sup_{\mm}\dim A_{\mm},
 \tag{16.1.3.5}
\end{equation}
其中 $\mm$ 遍历 $A$ 的极大理想。

\egap{16.1.4}
对于 $A$ 的每个素理想 $\pp$，有（14.2.2.2）
\begin{equation}
 \dim A_{\pp}+\dim(A/\pp)\leq\dim A.
 \tag{16.1.4.1}
\end{equation}
当 $\Spec A$ 是悬链的（分别地，等维的、等余维的、双等维的）时，称 $A$ 具有同一性质。当 $A$ 诺特且双等维时，对每个素理想，式（16.1.4.1）的两端相等（14.3.5）。\cmt{D4-H01}{本句中的“双等维”须结合前文（14.3.3）的译校注阅读；不能只采用该条中有问题的较弱条件b）来代替等长极大链条件。}

对任意理想 $\mathfrak a$，$A/\mathfrak a$ 的素理想对应于 $A$ 中包含 $\mathfrak a$ 的素理想，因此若 $A$ 悬链，则 $A/\mathfrak a$ 也悬链。同样，对任意乘法子集 $S$，由素理想的对应关系可知，$A$ 悬链蕴含 $S^{-1}A$ 悬链。

此外，$A$ 中包含于素理想 $\pp$ 的素理想，与 $A_{\pp}$ 的素理想一一对应。因此，$A$ 悬链的充要条件是每个 $A_{\pp}$ 都悬链。于是，说 $A$ 悬链，就是说（14.3.2）、（16.1.3.2）对于任意两个素理想 $\qq\subset\pp$，有
\cmt{D4-A01}{作者正式勘误 \hyperlink{ErrIV:5}{Err$_{\mathrm{IV}}$，5}（EGA IV第四部分，印刷第347页）要求将本句的“说 $A$ 悬链”改为“说一个整环 $A$ 悬链”。没有整环条件，上述等价表述一般不成立。这里保留初版原句，同时明示作者修订；详见本版核查报告。}
\begin{equation}
 \dim A_{\qq}+\dim(A_{\pp}/\qq A_{\pp})=\dim A_{\pp}.
 \tag{16.1.4.2}
\end{equation}
若 $A$ 双等维，则对每个素理想 $\pp$，$A_{\pp}$ 也双等维。

\egas{16.1.5}{命题}
设 $\rho:A\to B$ 为环同态，使 $B$ 成为整的 $A$-代数。则 $\dim B\leq\dim A$；若 $\rho$ 还是单射，则 $\dim B=\dim A$。
\begin{proof}
若 $\mathfrak n=\Ker\rho$，则 $\rho(A)\simeq A/\mathfrak n$，所以由（16.1.2.1），$\dim\rho(A)\leq\dim A$。由于 $B$ 整于 $\rho(A)$，只须在 $A\subset B$ 时证明第二断言。

若 $\pp_0\subsetneq\cdots\subsetneq\pp_m$ 是 $B$ 的素理想链，则 $\pp_i\cap A$ 是 $A$ 中两两不同的素理想（Bourbaki，\textit{交换代数}，第五章，§2，第1目，命题1的推论1），因而构成素理想链。反过来，对 $A$ 的每条素理想链 $\pp_0\subsetneq\cdots\subsetneq\pp_m$，可对每个 $i$ 找到 $B$ 的素理想 $\pp'_i$，满足 $\pp'_i\cap A=\pp_i$，且 $\pp'_i\subset\pp'_{i+1}$（$0\leq i<m$）（同处，定理1的推论2）。由此得到结论。
\end{proof}

\egas{16.1.6}{命题}
设 $B$ 是整环，$A$ 是 $B$ 的局部子环，$\mm$ 为 $A$ 的极大理想。假定 $A$ 整闭且 $B$ 整于 $A$。则对 $B$ 的每个极大理想 $\mathfrak n$，有 $\dim B_{\mathfrak n}=\dim A$。
\begin{proof}
（16.1.5）证明的第一部分表明 $\dim B_{\mathfrak n}\leq\dim A$。反过来，考虑 $A$ 的素理想链 $\pp_0\subsetneq\cdots\subsetneq\pp_m=\mm$。已知 $\mathfrak n\cap A=\mm$（Bourbaki，\textit{交换代数}，第五章，§2，第1目，命题1）；根据假设，可用第二个Cohen--Seidenberg定理（同书，第五章，§2，第4目，定理3），得到 $B$ 的素理想链
$\pp'_0\subsetneq\cdots\subsetneq\pp'_m=\mathfrak n$，其中 $\pp'_i\cap A=\pp_i$。于是 $\dim B_{\mathfrak n}\geq\dim A$，命题得证。
\end{proof}

\egap{16.1.7}
设 $M$ 为 $A$-模。称环 $A/\Ann M$ 的维数为 $M$ 的\textbf{维数}，记作 $\dim_A M$ 或 $\dim M$；这个环同构于 $M$ 的数乘同态环 $A_M$。因此，对 $k>0$，有 $\dim(M^k)=\dim M$。若 $M$ 有限型，则包含 $\Ann M$ 的素理想恰为 $\Supp M$ 的点（$0_{\mathrm I}$，1.7.4），且这个支集在 $\Spec A$ 中闭，因而
\begin{equation}
 \dim M=\dim(\Supp M)\leq\dim A.
 \tag{16.1.7.1}
\end{equation}
若 $N\subset M$ 是两个 $A$-模，则
\begin{equation}
 \dim N\leq\dim M,\qquad \dim(M/N)\leq\dim M.
 \tag{16.1.7.2}
\end{equation}
若 $S$ 为乘法子集且 $M$ 有限型，则
\begin{equation}
 \dim_{S^{-1}A}(S^{-1}M)\leq\dim_A M.
 \tag{16.1.7.3}
\end{equation}
事实上，$\Ann(S^{-1}M)=S^{-1}(\Ann M)$（Bourbaki，\textit{交换代数}，第二章，§2，第4目），且
$S^{-1}A/S^{-1}(\Ann M)=S^{-1}(A/\Ann M)$。

另一方面，每条包含 $\Ann M$ 的素理想极大链都以某个极大理想 $\mm$ 结束；其长度等于相应的 $A_{\mm}$ 中素理想链的长度，而这些素理想包含
$\Ann M_{\mm}=(\Ann M)_{\mm}$。由此及（16.1.7.3），得到
\begin{equation}
 \dim_A M=\sup_{\mm}\dim_{A_{\mm}}M_{\mm},
 \tag{16.1.7.4}
\end{equation}
其中 $\mm$ 遍历 $A$ 的极大理想。当 $A/\Ann M$ 悬链（分别地，等维、等余维、双等维）时，称 $M$ 具有同一性质。

\egas{16.1.8}{命题}
设 $A$ 为环，$M$ 为有限型 $A$-模，$\pp$ 为 $A$ 的素理想。则
\begin{equation}
 \dim_{A_{\pp}}M_{\pp}+\dim_A(M/\pp M)\leq\dim_A M.
 \tag{16.1.8.1}
\end{equation}
\begin{proof}
有 $\Ann M_{\pp}=(\Ann M)_{\pp}$；所以 $A_{\pp}$ 中包含 $\Ann M_{\pp}$ 的素理想，与 $A$ 中包含 $\Ann M$ 而包含于 $\pp$ 的素理想一一对应。另一方面，若 $\mathfrak n=\Ann M$，则
$V(\Ann(M/\pp M))=V(\pp+\mathfrak n)$（$0_{\mathrm I}$，1.7.5）；所以包含 $\Ann(M/\pp M)$ 的素理想包含 $\pp+\mathfrak n$。由此得到结论。最后这个观察还表明
\begin{equation}
 \dim_{A/\pp}(M/\pp M)=\dim_A(M/\pp M).
 \tag{16.1.8.2}
\end{equation}
\end{proof}

\egas{16.1.9}{命题}
设 $A$ 为诺特环，$\rho:A\to B$ 为环同态，$M$ 为 $B$-模，且 $M_{[\rho]}$ 是有限型 $A$-模。则
\begin{equation}
 \dim_A M_{[\rho]}=\dim_B M.
 \tag{16.1.9.1}
\end{equation}
\begin{proof}
$B$-模 $M$ 的数乘同态环 $B_M$ 可等同于 $C=\End_A(M_{[\rho]})$ 的子环。$C$ 是有限型 $A$-模（Bourbaki，\textit{交换代数}，第三章，§3，第1目，引理2），故 $B_M$ 也是。又 $A$-模 $M_{[\rho]}$ 的数乘同态环 $A_M$ 是 $B_M$ 的子环，等于 $A$ 在 $C$ 中的典范像；所以 $B_M$ 有限于 $A_M$。结论由（16.1.5）得到。
\cmt{D4-A02}{作者正式勘误 \hyperlink{ErrIV:6}{Err$_{\mathrm{IV}}$，6}删去 $A$ 诺特的假设。理由是对 $u\in B_M$，$A[u]$ 是 $C$ 的子 $A$-代数；其在有限型 $A$-模 $M_{[\rho]}$ 上的作用给出整性，从而 $u$ 整于 $A$，也整于 $A_M$，仍可用（16.1.5）。这是作者的推广，不是初版证明已在无诺特假设下成立。}
\end{proof}

\egap{16.1.10}
设 $A$ 为诺特环。若 $M$ 是有限长度 $A$-模，则 $\Supp M$ 是有限离散空间（Bourbaki，\textit{交换代数}，第四章，§2，第5目，命题7及其推论1），所以当 $M\ne0$ 时 $\dim M=0$。反过来，若 $M$ 是满足 $\dim M=0$ 的有限型 $A$-模，则 $\Supp M$ 的每一点都是闭点，即 $A$ 的极大理想，因此 $M$ 有限长度（同处，命题7）。

\subsection{诺特半局部环的维数}
\source{EGA IV 第一部分，印刷第25--28页；上传PDF第24--27页\orig{本小节重述Serre于1958年在法兰西学院课程中所作的讲述。}}

\egap{16.2.1}
设 $A$ 为诺特半局部环，$\rr$ 为其根基。回忆 $A$ 的\textbf{定义理想} $\qq$ 是满足 $\qq\subset\rr$ 且包含 $\rr$ 的某个幂的理想；这等价于说，包含 $\qq$ 的素理想都是极大理想，或说 $A/\qq$ 为Artin环。\cmt{D4-H02}{在半局部情形，后两种说法还须保留前面的 $\qq\subset\rr$ 限制。仅 $A/\qq$ 为Artin环，并不能推出 $\qq$ 包含于每个极大理想。例子及（16.3.6）中相关表述的处理见核查报告。}
设 $M$ 为有限型 $A$-模。对每个整数 $n>0$，$M/\qq^nM$ 的长度有限，等于
\[
 \sum_{i=0}^{n-1}\len(\gr_i M),
\]
其中 $\gr_\bullet M$ 是 $M$ 的 $\qq$-预进制滤过的相伴分次 $A/\qq$-模。已知有有理系数多项式 $Q_{\qq}(n)$，使得对充分大的 $n$，$\len(\gr_nM)=Q_{\qq}(n)$（III，2.5.4）。于是存在关于 $n$ 的有理系数多项式 $P_{\qq}(M,n)$，满足
\[
 \len(M/\qq^nM)=P_{\qq}(M,n)\qquad(n\text{ 充分大}).
\]
这个多项式唯一；当 $M\ne0$ 时，其首项系数为正。称它为 $M$ 关于 $\qq$-预进制滤过的\textbf{Hilbert--Samuel多项式}。

\egas{16.2.2.1}{引理}
设 $(M_n)$ 是 $M$ 的一个 $\qq$-良好滤过（Bourbaki，\textit{交换代数}，第三章，§3，第1目）。则当 $n$ 充分大时，$\len(M/M_n)$ 等于关于 $n$ 的一个多项式，其次数和首项系数都与 $P_{\qq}(M,n)$ 的相同。\cmt{D4-T02}{初版标题误排为（16.12.2.1）；本译稿按后续引用统一排为（16.2.2.1），在此记录该调整。}
\begin{proof}
存在 $n_0$，使 $n\geq n_0$ 时 $M_{n+1}=\qq M_n$。因此，对充分大的 $n$，有
\[
 \qq^{n+n_0}M\subset M_{n+n_0}=\qq^nM_{n_0}
 \subset\qq^nM\subset M_n,
\]
从而
\[
 \begin{aligned}
 P_{\qq}(M,n+n_0)&\geq\len(M/M_{n+n_0})\\
 &\geq P_{\qq}(M,n)\geq\len(M/M_n),
 \end{aligned}
\]
由此得到结论。
\end{proof}
若 $\qq'$ 是另一个定义理想，则对某个整数 $m$ 有 $\qq'\supset\qq^m$，因而
$P_{\qq'}(M,n)\leq P_{\qq}(M,mn)$。所以 $P_{\qq}(M,n)$ 的次数与所选定义理想无关；记作 $d(M)$。

\egas{16.2.2.2}{引理}
设 $0\to M'\to M\to M''\to0$ 是有限型 $A$-模的正合列。则多项式
\[
 P_{\qq}(M,n)-P_{\qq}(M',n)-P_{\qq}(M'',n)
\]
的次数不超过 $d(M')-1\leq d(M)-1$。
\begin{proof}
滤过 $M'_n=M'\cap\qq^nM$ 是 $\qq$-良好的（Bourbaki，\textit{交换代数}，第三章，§3，第1目，命题1）。\cmt{D4-A03}{初版此处印为 $M'\cap\qq^nM'$；作者在勘误清单三A部明确改为 $M'\cap\qq^nM$。本处采用已发表的改正，并保留其来源。}
又有
\[
 \len(M/\qq^nM)=\len(M''/\qq^nM'')+\len(M'/M'_n),
\]
所以把（16.2.2.1）用于 $M'$ 就得到本引理。
\end{proof}

\egas{16.2.3}{定理（Krull--Chevalley--Samuel）}
设 $A$ 为诺特半局部环，$\rr$ 为其根基，$M\ne0$ 为有限型 $A$-模，$d(M)$ 为 $M$ 关于任意定义理想的Hilbert--Samuel多项式的次数。另令 $s(M)$ 为下述整数 $n$ 的下确界：存在 $\rr$ 中的元素 $x_1,\ldots,x_n$，使
$M/(x_1M+\cdots+x_nM)$ 的长度有限。则
\[
 d(M)=s(M)=\dim M.
\]

\egas{16.2.3.1}{引理}
对 $x\in\rr$，令 ${}_xM$ 为 $M$ 中被 $x$ 零化的元素组成的子模。则：
\begin{enumerate}[label=(\roman*)]
 \item $s(M)\leq s(M/xM)+1$。
 \item 设 $\pp_i$（$1\leq i\leq m$）为 $\Supp M$ 中满足 $\dim(A/\pp_i)=\dim M$ 的素理想。若对每个 $i$ 都有 $x\notin\pp_i$，则
 $\dim(M/xM)\leq\dim M-1$。
 \item 若 $\qq$ 是 $A$ 的定义理想，则多项式
 \[
   P_{\qq}({}_xM,n)-P_{\qq}(M/xM,n)
 \]
 的次数不超过 $d(M)-1$。
\end{enumerate}
\begin{proof}
（i）中，若令 $N=M/xM$，且 $N/(x_1N+\cdots+x_mN)$ 长度有限，只须注意这个模同构于 $M/(xM+x_1M+\cdots+x_mM)$。

为证明（ii），令 $\mathfrak a=\Ann M$。包含 $\Ann(M/xM)$ 的素理想，正是包含 $\mathfrak a+Ax$ 的素理想（$0_{\mathrm I}$，1.7.5）。因此没有一个 $\pp_i$ 包含 $\mathfrak a+Ax$。按 $\dim M$ 的定义（16.1.7），包含 $\mathfrak a+Ax$ 的素理想链长度至多为 $\dim M-1$。

最后，为证明（iii），使用两个正合列
\[
 0\longrightarrow{}_xM\longrightarrow M\xrightarrow{x}xM\longrightarrow0,
 \qquad
 0\longrightarrow xM\longrightarrow M\longrightarrow M/xM\longrightarrow0.
\]
断言由（16.2.2.2）得到。\cmt{D4-H03}{左下标 $x$ 属于 ${}_xM=\Ker(x:M\to M)$，不能漏掉。若误读为 $P_{\qq}(M,n)-P_{\qq}(M/xM,n)$，即使 $A=k[[t]]$、$x=t$、$M=A$，所述次数界也不成立。}
\end{proof}
定理（16.2.3）的证明分三步进行。

\egap{16.2.3.2}
有 $\dim M\leq d(M)$。

对 $d(M)$ 作归纳；$d(M)=0$ 的情形是直接的。设 $d(M)\geq1$，取 $\pp_0\in\Supp M$，使 $\dim(A/\pp_0)=\dim M$。则 $\pp_0$ 是 $\Supp M$ 的极小元，因而属于 $\Ass M$（Bourbaki，\textit{交换代数}，第四章，§1，第4目，定理2）；于是 $M$ 含有同构于 $A/\pp_0$ 的子模 $N$（同处，第1目）。由（16.2.2.2），有 $d(M)\geq d(N)$，故归结为证明 $\dim N\leq d(N)$。

设 $\pp_0\subsetneq\pp_1\subsetneq\cdots\subsetneq\pp_n$ 为 $A$ 的一条素理想链，来证明 $n\leq d(N)$。$n=0$ 时无须证明。否则，存在
$x\in\pp_1\cap\rr$ 且 $x\notin\pp_0$：若 $\pp_0\supset\pp_1\cap\rr$，由于 $\pp_0$ 不包含 $\pp_1$，便有 $\pp_0\supset\rr$，从而 $\pp_0$ 极大，矛盾。对于 $i\geq1$，有 $\pp_i\in\Supp(N/xN)$，故 $n-1\leq\dim(N/xN)$（16.1.7）。又由 $x$ 的选择，${}_xN=0$，所以（16.2.3.1（iii））给出 $d(N/xN)\leq d(N)-1$。原文接着由归纳假设写出
\[
 d(N/xN)=\dim(N/xN)\geq n-1,
\]
从而 $n-1\leq d(N)-1$。\cmt{D4-G01}{此步骤只在归纳证明 $\dim\leq d$；此时可直接得到的是 $\dim(N/xN)\leq d(N/xN)$，不必预先写等号。将这一等号替换为所需的不等式，即足以完成该步，不是否定定理。}

\egap{16.2.3.3}
有 $d(M)\leq s(M)$。

设 $x_1,\ldots,x_n\in\rr$，且令 $\mathfrak a=\sum_i x_iA$ 后，$M/\mathfrak aM$ 有限长度。其伴随素理想于是都是极大理想（Bourbaki，\textit{交换代数}，第四章，§2，第5目，命题7）。由于包含
$\qq=\mathfrak a+(\rr\cap\Ann M)$ 的素理想都极大，$A/\qq$ 为Artin环（同处，命题9），而 $\qq$ 是 $A$ 的定义理想。

有 $\qq^mM/\qq^{m+1}M=\mathfrak a^mM/\mathfrak a^{m+1}M$。若 $M$ 由 $r$ 个元素 $\xi_j$ 生成，则这个 $A/\qq$-模由
\[
 x_1^{h_1}\cdots x_n^{h_n}\xi_j,
 \qquad h_1+\cdots+h_n=m,
\]
的像生成，因此可由 $r\binom{m+n-1}{n-1}$ 个元素生成。它的长度至多为
$a r\binom{m+n-1}{n-1}$，其中 $a$ 为常数。这就推出关于 $m$ 的多项式 $P_{\qq}(M,m)$ 的次数至多为参数个数 $n$；故 $d(M)\leq s(M)$。

\egap{16.2.3.4}
有 $s(M)\leq\dim M$。

首先，（16.2.3.2）表明 $n=\dim M$ 有限。对 $n$ 作归纳。$n=0$ 时，$M$ 有限长度（16.1.10），结论成立。设 $n\geq1$，令 $\pp_i$（$1\leq i\leq m$）为 $\Supp M$ 中满足 $\dim(A/\pp_i)=n$ 的素理想。这些是 $\Supp M$ 的极小元，因此只有有限个（Bourbaki，\textit{交换代数}，第四章，§1，第4目，定理2）。由于 $n\geq1$，它们都不是极大理想，所以可取 $x\in\rr$，使 $x$ 不属于任何 $\pp_i$（同书，第二章，§1，第1目，命题2）。由（16.2.3.1），
\[
 s(M)\leq s(M/xM)+1,\qquad
 \dim M\geq\dim(M/xM)+1.
\]
归纳假设给出 $s(M/xM)\leq\dim(M/xM)$，因而得到结论。

\egas{16.2.4}{推论}
在（16.2.3）的假设下，
\begin{equation}
 \dim_A M=\dim_{\widehat A}\widehat M.
 \tag{16.2.4.1}
\end{equation}
事实上，$\widehat M/\widehat{\rr}^{\,n}\widehat M$ 与 $M/\rr^nM$ 同构，故 $d(\widehat M)=d(M)$。

\egas{16.2.5}{推论}
诺特半局部环 $A$ 的维数有限，且等于从 $A$ 的根基中取元素、生成一个定义理想所需的最少元素数。

这就是 $M=A$ 时的等式 $s(M)=\dim M$。特别地：

\egas{16.2.6}{推论}
设 $A$ 为诺特局部环，$\mm$ 为其极大理想，$k=A/\mm$ 为其剩余域。则
\begin{equation}
 \dim A\leq\rk_k(\mm/\mm^2).
 \tag{16.2.6.1}
\end{equation}
事实上，若 $x_1,\ldots,x_n\in\mm$ 的模 $\mm^2$ 类组成 $k$-向量空间 $\mm/\mm^2$ 的一组基，则 $x_i$ 生成 $\mm$（Bourbaki，\textit{交换代数}，第二章，§3，第2目，命题4的推论2）。

\egas{16.2.7}{命题}
设 $k$ 为域，$B$ 为非负分次有限型 $k$-代数，由一次齐次元生成且 $B_0=k$。设 $E$ 为有限型分次 $B$-模。对于充分大的 $n$，$\rk_kE_n=r(n)$ 可写成 $\Phi(n)-\Phi(n-1)$，其中 $\Phi$ 为次数 $d$ 的多项式（III，2.5.4）。则存在 $B$ 中 $d$ 个次数至少为1的齐次元 $t_1,\ldots,t_d$，使 $E/(\sum_i t_iE)$ 在 $k$ 上有限维。
\begin{proof}
令 $\qq=\bigoplus_{n\geq1}B_n$，这是 $B$ 的极大理想；对每个 $n$，令 $E'_n=\bigoplus_{h\geq n+1}E_h$。每个 $B-\qq$ 中的元素，在Artin $B$-模 $E/E'_n$ 上给出的数乘是单射，因而是双射（Bourbaki，\textit{代数}，第八章，§2，第2目，引理3）。所以 $(E/E'_n)_{\qq}$ 可等同于 $E/E'_n$。由于 $B_{\qq}$ 的剩余域是 $k$，由假设可知 $\len_{B_{\qq}}((E/E'_n)_{\qq})$ 是次数 $d$ 的关于 $n$ 的多项式。

$E$ 有限型且 $B$ 由一次齐次元生成，故 $(E'_n)$ 是 $E$ 上的 $\qq$-良好滤过（Bourbaki，\textit{交换代数}，第三章，§1，第3目，引理1）；所以 $((E'_n)_{\qq})$ 是 $E_{\qq}$ 上的 $\qq B_{\qq}$-良好滤过。由（16.2.3）、（16.2.2.1）及 $B_{\qq}$ 诺特局部，得到 $\dim E_{\qq}=d$。

对 $d$ 作归纳，$d=0$ 时结论直接成立。设 $d\geq1$。$\Ass E$ 的极小元 $\pp_i$（$1\leq i\leq r$）是分次理想（同书，第四章，§3，第1目，命题1），且 $\qq$ 不等于这些理想的并。因此存在齐次元 $t_1\in\qq$，不属于任何 $\pp_i$（同书，第三章，§1，第4目，命题8）。$(\pp_i)_{\qq}$ 是 $\Ass E_{\qq}$ 的极小元（同书，第四章，§1，第2目，命题5），所以
$\dim(E_{\qq}/t_1E_{\qq})=d-1$（16.3.4）。将归纳假设用于分次模 $E/t_1E$，即得结论。
\end{proof}

% END continuations/iv0_16_12.tex


% BEGIN continuations/iv0_16_3.tex
% EGAIV.1.pdf, PDF pp.27--31, printed pp.28--32.
\subsection{诺特局部环中的参数系}
\source{EGA IV 第一部分，印刷第28--32页；上传PDF第27--31页}

\egas{16.3.1}{命题}
设 $A$ 为诺特环，$\pp$ 为 $A$ 的素理想，$n$ 为整数。以下条件等价：
\begin{enumerate}[label=\alph*)]
 \item $\htp(\pp)=\dim A_{\pp}\leq n$；
 \item 存在 $A$ 中 $n$ 个元素 $x_i$（$1\leq i\leq n$），使 $\pp$ 在包含它们所生成的理想 $\mathfrak a$ 的素理想集合中极小（也就是说，$V(\pp)$ 是 $V(\mathfrak a)$ 的一个不可约分支）。
\end{enumerate}
\begin{proof}
条件b）表明 $\mathfrak a A_{\pp}$ 为 $A_{\pp}$ 的定义理想，故由（16.2.5）得到a）。反过来，若a）成立，则 $A_{\pp}$ 有一个由 $n$ 个元素 $x_i/s$（$s\in A-\pp$）生成的定义理想 $\mathfrak b$。由 $A_{\pp}$ 的素理想与 $A$ 中包含于 $\pp$ 的素理想之间的一一对应，$A$ 中由 $x_i$ 生成的理想 $\mathfrak a$ 满足b）。
\end{proof}
取 $n=1$，得到特别情形：

\egas{16.3.2}{推论（主理想定理，Hauptidealsatz）}
诺特环中的一个素理想高度至多为1，当且仅当它在包含某个适当主理想的素理想集合中极小。

\egas{16.3.3}{推论（Artin--Tate）}
设 $A$ 为诺特整环。以下条件等价：
\begin{enumerate}[label=\alph*)]
 \item $A$ 为维数至多为1的半局部环；
 \item 存在 $A$ 中的非零元素 $f$，使 $A_f$ 为域。
\end{enumerate}
\begin{proof}
条件a）等价于说，$A$ 只有有限个非零素理想 $\mm_i$，且它们都是极大理想，因为 $(0)$ 是素理想。由于 $A$ 为整环，$\mm_i$ 的乘积非零；取这个乘积中的非零元素 $f$，则 $f$ 属于每个 $\mm_i$，所以 $(0)$ 是整环 $A_f$ 的唯一素理想，即 $A_f$ 是域。（这一部分证明没有用到 $A$ 诺特。）

现在证明b）推出a）。条件b）表示 $A$ 的每个非零素理想都包含 $f$。令 $\pp_i$（$1\leq i\leq n$）为包含 $f$ 的素理想中的极小者；由于 $A/fA$ 诺特，这些理想只有有限个。只须证明所有 $\pp_i$ 都极大：这样每个非零素理想必包含某个 $\pp_i$，因而等于它。

假设存在极大理想 $\mm$，包含某个 $\pp_i$ 而与所有 $\pp_i$ 不同。则 $\mm$ 不包含于这些 $\pp_i$ 的并（Bourbaki，\textit{交换代数}，第二章，§1，第1目，命题2），故可取 $g\in\mm$，使 $g$ 不属于任何 $\pp_i$。令 $\qq$ 为包含 $g$ 的素理想中的一个极小者。主理想定理（16.3.2）表明 $\qq$ 高度为1，因而非零；所以它包含 $f$，也包含某个 $\pp_i$。但是 $g\in\qq$ 而 $g\notin\pp_i$，故 $\qq\ne\pp_i$，这使 $\qq$ 的高度至少为2，矛盾。
\end{proof}

\egas{16.3.4}{命题}
设 $A$ 为诺特半局部环，$\rr$ 为其根基，$M$ 为有限型 $A$-模，$\pp_i$ 为 $\Supp M$ 中满足 $\dim(A/\pp_i)=\dim M$ 的素理想。则 $\pp_i$ 是 $\Ass M$ 的极小元，且
$\dim(M/\pp_iM)=\dim M$。对于每个 $x\in\rr$，有
\begin{equation}
 \dim(M/xM)\geq\dim M-1.
 \tag{16.3.4.1}
\end{equation}
此外，两端相等的充要条件是 $x$ 不属于任何 $\pp_i$。
\begin{proof}
第一断言来自 $\pp_i$ 是 $\Supp M$ 的极小元（原文引（16.1.6）），而这些极小元也正是 $\Ass M$ 的极小元（Bourbaki，\textit{交换代数}，第四章，§1，第4目，定理2）。又 $\Supp(M/\pp_iM)$ 等于包含 $\pp_i$ 的素理想集合（$0_{\mathrm I}$，1.7.5），故
$\dim(M/\pp_iM)=\dim(A/\pp_i)=\dim M$。

令 $N=M/xM$。若 $x_1,\ldots,x_n\in\rr$，且 $N/(x_1N+\cdots+x_nN)$ 长度有限，则这表示 $M/(xM+x_1M+\cdots+x_nM)$ 长度有限；由（16.2.3）得到（16.3.4.1）。

当 $x$ 不属于任何 $\pp_i$ 时，由（16.3.4.1）和（16.2.3.1）得到等号。反过来，若 $\dim(M/xM)=\dim M-1$，则任何 $\pp_i$ 都不能属于 $\Supp(M/xM)$。后一个支集由包含 $\Ann M+Ax$ 的素理想组成；由于 $\pp_i$ 都包含 $\Ann M$，它们不能包含 $x$。
\end{proof}

\egas{16.3.5}{推论}
采用（16.3.4）的记号。对每个包含于某个 $\pp_i$ 的理想 $\mathfrak a$，有 $\dim(M/\mathfrak aM)=\dim M$；若令 $N=M/\mathfrak aM$，则
$\dim(M/xM)=\dim M-1$ 蕴含 $\dim(N/xN)=\dim N-1$。
\begin{proof}
$\Supp N$ 等于包含 $\mathfrak a+\Ann M$ 的素理想集合，故 $\pp_i\in\Supp N$。由
$\dim N\leq\dim M=\dim(A/\pp_i)$，得到 $\dim(A/\pp_i)=\dim N$。此外，$\Supp N$ 中满足 $\dim(A/\qq)=\dim N$ 的素理想 $\qq$，都是某些 $\pp_j$。若 $x$ 不属于任何 $\pp_j$，由（16.3.4）即得 $\dim(N/xN)=\dim N-1$。
\end{proof}

\egas{16.3.6}{定义}
设 $A$ 为诺特半局部环，$\rr$ 为其根基，$M$ 为有限型 $A$-模，令 $n=\dim M$。称 $\rr$ 中由 $n$ 个元素组成的一个族 $(x_i)_{1\leq i\leq n}$ 为 $M$ 的\textbf{参数系}，若 $M/(x_1M+\cdots+x_nM)$ 的长度有限。

由（16.2.3）可知，总存在这样的参数系。已在（16.2.3.3）的证明中看到，若
$\jj=\sum_{i=1}^n x_iA$ 且 $M/\jj M$ 长度有限，则 $\qq=\jj+\Ann M$ 为 $A$ 的定义理想；反过来，若如此，$M/\jj M$ 是Artin环 $A/\qq$ 上的有限型模，因此长度有限。\cmt{D4-H04}{本句在半局部情形的“$A$ 的定义理想”说法过强：$\Ann M$ 可能不包含于 $\rr$。（16.2.3.3）实际取的是 $\jj+(\rr\cap\Ann M)$。准确的商环说法是 $\jj$ 在 $A/\Ann M$ 中的像为定义理想；报告给出半局部反例。另，参数个数 $n$ 的定义及随后参数系论证按非零有限型模理解；零模在本书约定下维数为 $-\infty$，不能将其直接当作参数个数。}
所以，说 $(x_i)$ 是 $M$ 的参数系，等价于说它是 $A/\Ann M$ 的参数系，也就是其在 $A/\Ann M$ 中的像组成该环的参数系。

若 $(x_i)_{1\leq i\leq n}$ 是 $M$ 的参数系，则对任意整数 $k\geq1$，$(x_i^k)$ 也如此。根据前述说明，可归结到 $M=A$；若 $\jj$ 为 $A$ 的定义理想，则 $\sum_{i=1}^n x_i^kA$ 包含 $\jj^{nk}$，因而仍为定义理想。

\egas{16.3.7}{命题}
设 $A$ 为诺特半局部环，$\rr$ 为其根基，$x_1,\ldots,x_k\in\rr$，$M$ 为有限型 $A$-模。则
\begin{equation}
 \dim\bigl(M/(x_1M+\cdots+x_kM)\bigr)\geq\dim M-k.
 \tag{16.3.7.1}
\end{equation}
此外，两端相等，当且仅当能在 $\rr$ 中选取 $x_i$（$k+1\leq i\leq n=\dim M$），使 $(x_i)_{1\leq i\leq n}$ 为 $M$ 的参数系。
\begin{proof}
对 $k$ 归纳，应用（16.3.4.1），即得不等式。若两端相等，取
$M/(x_1M+\cdots+x_kM)$ 的参数系 $x_{k+1},\ldots,x_n$。则
$M/(x_1M+\cdots+x_kM+x_{k+1}M+\cdots+x_nM)$ 长度有限，故 $(x_i)_{1\leq i\leq n}$ 为 $M$ 的参数系。

反过来，若存在这样的 $x_i$，令 $N=M/(x_1M+\cdots+x_kM)$，则 $N/(x_{k+1}N+\cdots+x_nN)$ 长度有限，所以 $\dim N\leq n-k$。结合（16.3.7.1），得到等号。
\end{proof}

\egas{16.3.8}{命题}
设 $A$ 为诺特半局部环，$X=\Spec A$，$\mathfrak a\ne A$ 为理想，且
$\codim(V(\mathfrak a),X)=r>0$。则存在 $\mathfrak a$ 中的 $r$ 个元素 $x_1,\ldots,x_r$，为 $A$ 的一个参数系的一部分，并满足
\[
 \codim\bigl(V(\mathfrak a),V(Ax_1+\cdots+Ax_r)\bigr)=0.
\]
\begin{proof}
令 $\pp_i$（$1\leq i\leq m$）为 $A$ 的极小素理想，$\qq_j$（$1\leq j\leq n$）为包含 $\mathfrak a$ 的素理想中的极小者。由（14.2.1）、（16.1.3.2），
\[
 \codim(V(\mathfrak a),X)
 =\inf_j\codim(V(\qq_j),X)=\inf_j\dim A_{\qq_j}.
\]
假设 $r>0$ 表明 $\mathfrak a$ 不包含于任何 $\pp_i$，所以存在 $x\in\mathfrak a$，不属于任何 $\pp_i$（Bourbaki，\textit{交换代数}，第二章，§1，第1目，命题2）。同样有
\[
 \begin{aligned}
 \codim(V(\mathfrak a),V(Ax))
 &=\inf_j\dim(A/Ax)_{\qq_j}\\
 &=\inf_j\dim(A_{\qq_j}/A_{\qq_j}x).
 \end{aligned}
\]
因为 $x/1$ 不属于 $A_{\qq_j}$ 的任何极小素理想，由（16.3.4），
$\dim(A_{\qq_j}/A_{\qq_j}x)=\dim A_{\qq_j}-1$。故
$\codim(V(\mathfrak a),V(Ax))=r-1$。原文接着说，同理 $\dim(A/Ax)=\dim A-1$。\cmt{D4-P01}{本段仅写 $x\in\mathfrak a$、避开极小素理想，尚未保证 $x\in\rr$，而（16.3.4）及参数系定义要求这一点。这里记录为半局部情形的证明适用条件缺口；未据此宣布命题结论错误。报告给出能使这一具体取元步骤失效的例子，并区分局部环情形。}

现在对 $r$ 归纳。若 $r>1$，在 $A'=A/Ax$ 中存在 $r-1$ 个元素 $x'_i$（$2\leq i\leq r$），属于 $\mathfrak a/Ax$，为 $A'$ 的参数系的一部分，并满足
\[
 \codim\bigl(V(\mathfrak a/Ax),V(A'x'_2+\cdots+A'x'_r)\bigr)=0.
\]
取元素 $x'_i$（$r+1\leq i\leq s$），使 $(x'_i)_{2\leq i\leq s}$ 成为 $A'$ 的一个参数系。对 $2\leq i\leq s$，在类 $x'_i$ 中选取 $A$ 的元素 $x_i$，并在 $2\leq i\leq r$ 时令 $x_i\in\mathfrak a$。于是 $A/(Ax+Ax_2+\cdots+Ax_s)$ 长度有限；又 $\dim A=s$，故 $x_1=x$ 及其余 $x_i$ 满足命题要求。
\end{proof}

\egas{16.3.9}{命题}
设 $A,B$ 为诺特局部环，$\mm$ 为 $A$ 的极大理想，$k$ 为其剩余域，$\varphi:A\to B$ 为局部同态。
\begin{enumerate}[label=(\roman*)]
 \item 有
 \begin{equation}
  \dim B\leq\dim A+\dim(B\otimes_A k).
  \tag{16.3.9.1}
 \end{equation}
 \item 若 $M\ne0$ 为有限型 $A$-模，$N\ne0$ 为有限型 $B$-模，则
 \begin{equation}
  \dim_B(M\otimes_A N)\leq\dim_A M+
  \dim_{B\otimes_A k}(N\otimes_A k).
  \tag{16.3.9.2}
 \end{equation}
\end{enumerate}
\begin{proof}
令 $m=\dim A$，取 $A$ 的参数系 $(s_i)_{1\leq i\leq m}$（16.3.6），并置 $\mathfrak a=\sum_i s_iA$。则 $A/\mathfrak a$ 长度有限，因而是Artin环；它的极大理想 $\mm/\mathfrak a$ 幂零。因此 $B/\mathfrak a B$ 的理想 $\mm B/\mathfrak a B$ 也幂零，因为它是 $(\mm/\mathfrak a)\otimes_A B$ 的像。于是
$\dim(B/\mathfrak a B)=\dim(B/\mm B)$，这两个环的谱具有相同的底空间。

又 $B/\mathfrak a B=B/(s_1B+\cdots+s_mB)$，且 $s_i$ 在 $B$ 中的像属于其极大理想。由（16.3.7），得到
$\dim B\leq m+\dim(B\otimes_A k)$，即（16.3.9.1）。

为证明（16.3.9.2），令 $\mathfrak r=\Ann_A M$、$\mathfrak s=\Ann_B N$。按定义（16.1.7），有
$\dim_A M=\dim(A/\mathfrak r)$、$\dim_B N=\dim(B/\mathfrak s)$。另一方面，$\Supp(M\otimes_A N)$ 是 $\Spec B$ 的闭子集，可等同于（I，9.1.13）
\[
 \begin{aligned}
 \Supp M\times_{\Spec A}\Supp N
 &=\Spec\bigl((B/\mathfrak s)\otimes_A(A/\mathfrak r)\bigr)\\
 &=\Spec\bigl(B/(\mathfrak s+\mathfrak rB)\bigr).
 \end{aligned}
\]
由于 $N\otimes_A k$ 是有限型 $B$-模，有
$\dim_{B\otimes_Ak}(N\otimes_Ak)=\dim_B(N\otimes_Ak)$（16.1.9）；同样的论证表明，它的支集为 $\Spec(B/(\mathfrak s+\mm B))$。将（16.3.9.1）用于 $A/\mathfrak r$ 与 $B/(\mathfrak s+\mathfrak rB)$，得到
\[
 \dim B/(\mathfrak s+\mathfrak rB)
 \leq\dim(A/\mathfrak r)+\dim B/(\mathfrak s+\mm B),
\]
这就是（16.3.9.2）。
\end{proof}

\egas{16.3.10}{推论}
在（16.3.9）的假设下，若 $\mm B$ 是 $B$ 的定义理想，则 $\dim B\leq\dim A$。若进一步 $A$ 为整环且 $\dim A=\dim B$，则 $\varphi$ 单射。
\begin{proof}
此时 $\dim(B/\mm B)=0$，故第一断言由（16.3.9）得到。还可将 $A$ 替换为 $\varphi(A)$，于是 $\dim B\leq\dim\varphi(A)$。若 $\mathfrak a=\Ker\varphi\ne0$，且 $A$ 为整环，则由（16.1.2.2），
$\dim\varphi(A)=\dim(A/\mathfrak a)<\dim A$；因此 $\dim A=\dim B$ 只能在 $\mathfrak a=0$ 时发生。
\end{proof}

% END continuations/iv0_16_3.tex


% BEGIN continuations/iv0_16_45.tex
% EGAIV.1.pdf, PDF pp.31--38, printed pp.32--39.
\subsection{深度与余深度}
\source{EGA IV 第一部分，印刷第32--36页；上传PDF第31--35页\orig{在第III章第三部分中，我们将从上同调的观点，并在更一般的框架中，发展深度的概念。}}

\egas{16.4.1}{命题}
设 $A$ 为诺特局部环，$\mm$ 为其极大理想，$M$ 为有限型 $A$-模。则每个由 $\mm$ 的元素组成的 $M$-正则序列 $(x_i)_{1\leq i\leq r}$（15.1.7），都为 $M$ 的一个参数系的一部分。
\begin{proof}
对 $r$ 作归纳。考虑 $A$-模 $N=M/(x_1M+\cdots+x_{r-1}M)$。由归纳假设，$x_1,\ldots,x_{r-1}$ 为 $M$ 的参数系的一部分，所以 $\dim N=\dim M-(r-1)$（16.3.7）。由于乘以 $x_r$ 在 $N$ 上单射，$x_r$ 不属于 $N$ 的任何伴随素理想。由（16.3.4），
\[
 \dim(N/x_rN)=\dim N-1,
\]
也就是
\[
 \dim\bigl(M/(x_1M+\cdots+x_rM)\bigr)=\dim M-r.
\]
结论因而由（16.3.7）得到。
\end{proof}
因此，当 $M\ne0$ 时，由 $\mm$ 的元素组成的 $M$-正则序列，至多有 $\dim M$ 个元素。\cmt{D4-H05}{与（16.3.6）相同，前面的“为参数系的一部分”须按 $M\ne0$ 阅读。对零模，任意有限序列在本书的正则性定义下都正则，而其维数为 $-\infty$。原文紧接着明确在长度界中加入了 $M\ne0$。}

\egas{16.4.2}{引理}
设 $A$ 为诺特局部环，$\mm$ 为其极大理想，$k=A/\mm$ 为其剩余域，$M$ 为一个 $A$-模。由 $\mm$ 的元素组成的 $M$-正则序列 $(x_i)_{1\leq i\leq n}$ 为极大序列，当且仅当
\cmt{D4-E01}{此处原文没有假定 $M$ 有限型。按这样的字面范围，引理不成立：在 $A=k[x,y]_{(x,y)}$ 上取所有高度1素理想商模的直和，空正则序列已极大，却没有非零 $\Hom_A(k,M)$。错误在于对可能无限的伴随素理想族使用有限素理想回避。加入 $M$ 有限型可修补；本版报告给出完整反例及修补证明。}
\[
 \Hom_A\bigl(k,M/(x_1M+\cdots+x_nM)\bigr)\ne0.
\]
\begin{proof}
说 $(x_i)$ 极大，表示对任何 $x\in\mm$，乘以 $x$ 在
$N=M/(x_1M+\cdots+x_nM)$ 上都不单射；等价地，$\mm$ 包含于 $N$ 的伴随素理想之并（Bourbaki，\textit{交换代数}，第四章，§1，第1目，命题2的推论2）。原文随后由 $\mm$ 极大，断言它必等于其中一个伴随素理想（同书，第二章，§1，第1目，命题2）。于是存在 $\xi\in N$，其零化理想为 $\mm$；故 $N$ 的子模 $A\xi$ 同构于 $A/\mm=k$，由此得到引理。
\end{proof}

\egas{16.4.3}{引理}
设 $A$ 为诺特局部环，$\mm$ 为其极大理想，$k=A/\mm$ 为其剩余域，$M$ 为 $A$-模，$(x_i)_{1\leq i\leq n}$ 为由 $\mm$ 的元素组成的 $M$-正则序列。则 $A$-模
\[
 \Hom_A\bigl(k,M/(x_1M+\cdots+x_nM)\bigr)
 \quad\text{与}\quad \Ext_A^n(k,M)
\]
同构；当 $n\geq1$ 时，它们还同构于 $\Ext_A^{n-1}(k,M/x_1M)$。
\begin{proof}
对 $n$ 作归纳；$n=0$ 时结论直接成立。令 $N=M/x_1M$。序列 $(x_i)_{2\leq i\leq n}$ 为 $N$-正则序列，故
\[
 \Hom_A\bigl(k,M/(x_1M+\cdots+x_nM)\bigr)
 =\Hom_A\bigl(k,N/(x_2N+\cdots+x_nN)\bigr)
\]
由归纳假设同构于 $\Ext_A^{n-1}(k,N)$。正合列
$0\to M\xrightarrow{x_1}M\to N\to0$ 给出Ext正合列
\begin{equation}
 \begin{aligned}
 \Ext_A^{n-1}(k,M)&\longrightarrow\Ext_A^{n-1}(k,N)\\
 &\longrightarrow\Ext_A^n(k,M)\xrightarrow{x_1}\Ext_A^n(k,M).
 \end{aligned}
 \tag{16.4.3.1}
\end{equation}
归纳假设表明，$\Ext_A^{n-1}(k,M)$ 同构于
$\Hom_A(k,M/(x_1M+\cdots+x_{n-1}M))$；后者为零，因为 $(x_i)_{1\leq i\leq n-1}$ 不是极大的 $M$-正则序列（16.4.2）。\cmt{D4-H06}{这里仅用到（16.4.2）中正确的一个方向：若存在非零 $\Hom_A(k,N)$，则 $\mm$ 中每个元素都在 $N$ 上为零因子。因而上一引理在任意模情形的反例，不直接反驳本引理；报告作了区分。}

另一方面，$x_1\in\mm$，故对任意 $A$-模 $T$，乘以 $x_1$ 在 $\Hom_A(k,T)$ 上为零，从而在所有 $\Ext_A^i(k,T)$ 上也为零。由（16.4.3.1）得到引理的各断言。
\end{proof}

\egas{16.4.4}{推论}
设 $A$ 为诺特局部环，$\mm$ 为其极大理想，$k=A/\mm$ 为其剩余域，$M\ne0$ 为有限型 $A$-模。所有由 $\mm$ 的元素组成的极大 $M$-正则序列都有相同的元素个数；这个数等于使 $\Ext_A^n(k,M)\ne0$ 的最小整数 $n$。

\egas{16.4.5}{定义}
设 $A$ 为诺特局部环，$\mm$ 为其极大理想，$M$ 为有限型 $A$-模。称由 $\mm$ 中元素组成的 $M$-正则序列的元素个数的上确界为 $M$ 的\textbf{深度}，记作 $\prof_A M$ 或 $\prof M$。

所以 $\prof M=+\infty$ 当且仅当 $M=0$，而由（16.4.1），
\begin{equation}
 \prof M\leq\dim M\qquad(M\ne0).
 \tag{16.4.5.1}
\end{equation}
对任意 $k>0$，有 $\prof(M^k)=\prof M$。

\egas{16.4.6}{命题}
设 $A$ 为诺特局部环，$\mm$ 为其极大理想，$M$ 为有限型 $A$-模。
\begin{enumerate}[label=(\roman*)]
 \item $\prof M=0$ 的充要条件是 $\mm\in\Ass M$，也就是 $M$ 含有同构于剩余域 $k=A/\mm$ 的子模。
 \item 对每个属于 $\mm$ 的 $M$-正则元素 $x$，有
 \begin{equation}
  \prof(M/xM)=\prof M-1.
  \tag{16.4.6.1}
 \end{equation}
 \item 若 $M\ne0$，则
 \begin{equation}
  \prof M\leq\inf_{\pp\in\Ass M}\dim(A/\pp)\leq\dim M.
  \tag{16.4.6.2}
 \end{equation}
 \item 有 $\prof_A M=\prof_{\widehat A}\widehat M$。
\end{enumerate}
\begin{proof}
（i）$\prof M=0$ 表示不存在 $x\in\mm$ 使数乘在 $M$ 上单射，即每个 $x\in\mm$ 都属于 $M$ 的某个伴随素理想（Bourbaki，\textit{交换代数}，第四章，§1，第1目，命题2的推论2）。而 $\mm$ 只有在等于某个伴随素理想时，才能被这些理想的并覆盖（同书，第二章，§1，第1目，命题2）。

（ii）若 $(x_i)_{1\leq i\leq m}$ 是由 $\mm$ 中元素组成的 $(M/xM)$-正则序列，则将 $x$ 放在这些元素之前，得到 $M$-正则序列。根据定义，前一序列极大时，所得序列也极大。由此得到结论。

（iii）对 $n=\prof M$ 作归纳，$n=0$ 时结论直接成立。使用下面的引理。
\renewcommand{\qedsymbol}{}
\end{proof}

\egas{16.4.6.3}{引理}
设 $A$ 为诺特环，$M$ 为有限型 $A$-模，$t\in A$ 为 $M$-正则元素。则对每个 $\pp\in\Ass M$，包含 $\pp+At$ 的素理想中的每个极小者都属于 $\Ass(M/tM)$。
\begin{proof}
已知有正合列 $0\to M'\to M\to M''\to0$，其中
$\Ass M'=\{\pp\}$、$\Ass M''=\Ass M-\{\pp\}$（Bourbaki，\textit{交换代数}，第四章，§1，第1目，命题4）。由于 $t$ 为 $M$-正则，它也是 $M'$-正则和 $M''$-正则的（同处，命题2的推论2）。由此，（15.1.18）给出正合列
\[
 0\longrightarrow M'/tM'\longrightarrow M/tM\longrightarrow M''/tM''\longrightarrow0.
\]
因此 $\Ass(M'/tM')\subset\Ass(M/tM)$。又包含 $\pp$ 的素理想正是 $\Supp M'$ 的点（同处，第3目，命题7），所以包含 $\pp+At$ 的素理想正是
$M'/tM'=M'\otimes_A(A/tA)$ 的支集的点（$0_{\mathrm I}$，1.7.5）。这些素理想中的极小元属于 $\Ass(M'/tM')$（Bourbaki，同处，第3目，命题7的推论1），从而也属于 $\Ass(M/tM)$。
\end{proof}
回到（iii）的证明。取 $M$-正则元素 $x\in\mm$，由（ii），有 $\prof(M/xM)=n-1$。对任意 $\pp\in\Ass M$，引理给出属于 $\Ass(M/xM)$ 且包含 $\pp+xA$ 的素理想 $\pp'$。因 $x$ 为 $M$-正则，$x\notin\pp$，故 $\pp'\ne\pp$，从而
$\dim(A/\pp')\leq\dim(A/\pp)-1$。归纳假设给出 $n-1\leq\dim(A/\pp')$，所以 $n\leq\dim(A/\pp)$。

（iv）可限于 $M\ne0$。回忆 $\widehat M\simeq M\otimes_A\widehat A$（$0_{\mathrm I}$，7.3.3）；由于 $\widehat A$ 在 $A$ 上平坦，有典范同构
\[
 \Ext_A^i(k,M)\otimes_A\widehat A
 \simeq\Ext_{\widehat A}^i(k,M\otimes_A\widehat A),
\]
因为 $k\simeq k\otimes_A\widehat A$（$0_{\mathrm{III}}$，12.3.5）。由（16.4.4）及 $\widehat A$ 在 $A$ 上忠实平坦（$0_{\mathrm I}$，7.3.5），得到（iv）。

\egap{16.4.6.4}
若 $\pp$ 是 $A$ 的素理想，既可能有
$\prof_{A_{\pp}}M_{\pp}<\prof_A M$，也可能有反向的严格不等式。
第一种情形的例子：取维数至少为1的局部整环 $A$。则 $\prof A\geq1$，但其分式域 $K$ 满足 $\prof K=0$。

第二种情形：取维数至少为2的诺特局部整环 $A_0$，极大理想为 $\mm_0$，剩余域为 $k=A_0/\mm_0$。令 $A=A_0\oplus\jj$ 为 $A_0$ 由 $A_0$-模 $k$ 作出的平凡扩张（18.2.3），其中理想 $\jj$ 在 $A_0$-模意义下同构于 $k$。乘法为
\[
 (x,\xi)(x',\xi')=(xx',x\xi'+x'\xi).
\]
$\jj$ 是 $A$ 的幂零根，且 $A/\jj\simeq A_0$。因此 $A$ 的每个素理想都形如 $\pp_0\oplus\jj$，其中 $\pp_0$ 为 $A_0$ 的素理想；特别地，$\mm=\mm_0\oplus\jj$ 是唯一极大理想。$\mm$ 的每个元素都零化 $\jj$，所以 $\prof A=\prof A_{\mm}=0$。但若 $\pp_0\ne\mm_0$，$\jj$ 中的元素被 $\mm_0-\pp_0$ 中的元素零化，因此 $A_{\pp}$ 典范同构于 $(A_0)_{\pp_0}$；由于 $A_0$ 为整环，在 $\pp_0\ne0$ 时 $\prof A_{\pp}=\prof(A_0)_{\pp_0}\geq1$。

\egas{16.4.7}{推论}
诺特约化局部环 $A$ 的深度为零，当且仅当 $A$ 是域。
\begin{proof}
由于 $A$ 的所有极小素理想之交为零，$A$ 没有嵌入伴随素理想。由（16.4.6（i）），若 $\prof A=0$，则极大理想 $\mm$ 也是极小素理想，因而是唯一素理想。由于 $A$ 约化，有 $\mm=0$。
\end{proof}

\egas{16.4.8}{命题}
设 $A,B$ 为诺特局部环，$\rho:A\to B$ 为局部同态，$M$ 为 $B$-模，且 $M_{[\rho]}$ 为有限型 $A$-模。则
\begin{equation}
 \prof_A M_{[\rho]}=\prof_B M.
 \tag{16.4.8.1}
\end{equation}
\begin{proof}
可限于 $M\ne0$。设 $\prof_A M_{[\rho]}=n$，取由 $A$ 的极大理想中元素组成的极大 $M_{[\rho]}$-正则序列 $(x_i)_{1\leq i\leq n}$。则 $\rho(x_i)$ 组成由 $B$ 的极大理想中元素构成的 $M$-正则序列。令
$N=M/(\rho(x_1)M+\cdots+\rho(x_n)M)$，由（16.4.6（ii）），
$\prof_B N=\prof_B M-n$。同样，
$N_{[\rho]}=M_{[\rho]}/(x_1M_{[\rho]}+\cdots+x_nM_{[\rho]})$，且 $\prof_A N_{[\rho]}=0$。故只须证明 $n=0$ 的情形。

令 $k$ 为 $A$ 的剩余域。由于 $k$ 是 $A$ 的商模，
$P=\Hom_A(k,M_{[\rho]})$ 是 $\Hom_A(A,M_{[\rho]})=M_{[\rho]}$ 的子模，且由假设与（16.4.6（i）），$P\ne0$。$P$ 有限型，并被 $A$ 的极大理想零化，故它是有限维 $k$-向量空间，因而是有限长度 $A$-模。另一方面，$P$ 也是 $B$-模 $M$ 的子模；作为 $A$-模长度有限，作为 $B$-模更是如此。因为非零，它含有一个简单 $B$-子模，即同构于 $B$ 的剩余域的子模。于是由（16.4.6（i）），$\prof_B M=0$。
\end{proof}

\egas{16.4.9}{定义}
设 $A$ 为诺特局部环，$M$ 为有限型 $A$-模。若 $M\ne0$，称有限非负整数
$\dim_A M-\prof_A M$（16.4.5.1）为 $M$ 的\textbf{余深度}，记作 $\coprof_A M$ 或 $\coprof M$。若 $M=0$，置 $\coprof M=0$。

对每个 $k>0$，有 $\coprof(M^k)=\coprof M$。

\egas{16.4.10}{命题}
设 $A$ 为诺特局部环，$M$ 为有限型 $A$-模。
\begin{enumerate}[label=(\roman*)]
 \item 对每个属于 $A$ 的极大理想的 $M$-正则元素 $x$，有
 $\coprof(M/xM)=\coprof M$。
 \item $\coprof_A M=\coprof_{\widehat A}\widehat M$。
\end{enumerate}
事实上，（i）由（16.3.4）、（16.4.6（ii））得到；（ii）由（16.2.4）、（16.4.6（iv））得到。
以后将证明（IV，6.11.5），对 $A$ 的每个素理想 $\pp$，有
$\coprof_{A_{\pp}}M_{\pp}\leq\coprof_A M$。

\egas{16.4.11}{命题}
在（16.4.8）的假设下，
\begin{equation}
 \coprof_A M_{[\rho]}=\coprof_B M.
 \tag{16.4.11.1}
\end{equation}
这由（16.1.9）、（16.4.8）得到。

\subsection{Cohen--Macaulay模}
\source{EGA IV 第一部分，印刷第36--39页；上传PDF第35--38页}

\egas{16.5.1}{定义}
设 $A$ 为诺特局部环。若有限型 $A$-模 $M$ 满足 $\coprof M=0$，则称其为\textbf{Cohen--Macaulay模}。也就是说，或者 $M=0$，或者 $M\ne0$ 且 $\dim M=\prof M$。若 $A$-模 $A$ 为Cohen--Macaulay模，则称 $A$ 为\textbf{Cohen--Macaulay环}。

有限长度 $A$-模是Cohen--Macaulay模，因为其维数为零（16.1.10）。维数1的约化局部环 $A$ 是Cohen--Macaulay环：它的极大理想 $\mm$ 不可能属于 $\Ass A$（Bourbaki，\textit{交换代数}，第四章，§2，第5目，命题10），因而由（16.4.6（i））有 $\prof A=1$。

若 $A$ 是维数至少为2的整闭局部整环，则 $\prof A\geq2$。事实上，若 $0\ne x\in\mm$，已知 $A/xA$ 的伴随素理想高度为1（同书，第七章，§1，第4目，命题8），因而都不等于 $\mm$；它们的并因此也不等于 $\mm$（同书，第二章，§1，第1目，命题2）。所以存在由 $\mm$ 中元素组成的 $A$-正则序列 $(x,y)$，这证明了上述断言。由此，维数2的整闭局部环是Cohen--Macaulay环。

\egas{16.5.2}{命题}
设 $A$ 为诺特局部环，$M$ 为有限型 $A$-模。$M$ 是Cohen--Macaulay模，当且仅当 $\widehat M$ 是Cohen--Macaulay $\widehat A$-模。

这直接由（16.4.10（ii））得到。

\egas{16.5.3}{命题}
在（16.4.8）的假设下，$M$ 是Cohen--Macaulay $B$-模，当且仅当 $M_{[\rho]}$ 是Cohen--Macaulay $A$-模。

这由（16.4.11）得到。

\egas{16.5.4}{命题}
设 $A$ 为诺特局部环，$M\ne0$ 为有限型 $A$-模。若 $M$ 是Cohen--Macaulay模，则对每个 $\pp\in\Ass M$，有 $\dim(A/\pp)=\dim M$；特别地，$M$ 没有嵌入伴随素理想。

这由（16.4.6.2）直接得到。也可以说，$M$（或 $\Supp M$）等维（16.1.7）。

\egas{16.5.5}{命题}
设 $A$ 为诺特局部环，$\mm$ 为其极大理想，$M\ne0$ 为有限型 $A$-模，$x\in\mm$。假定 $M$ 是Cohen--Macaulay模。以下条件等价：
\begin{enumerate}[label=\alph*)]
 \item $x$ 是 $M$-正则元素；
 \item $\dim(M/xM)=\dim M-1$；
 \item $x$ 不属于 $\Supp M$ 的任何极小元。
\end{enumerate}
此外，在这些条件下，$M/xM$ 也是Cohen--Macaulay模。
\begin{proof}
由（16.5.4），$\Ass M$ 的所有元素都是 $\Supp M$ 的极小元，且满足 $\dim(A/\pp)=\dim M$。因此a）与c）的等价来自Bourbaki，\textit{交换代数}，第四章，§1，第1目，命题2的推论2；b）与c）的等价来自（16.3.4）。最后，$M/xM$ 为Cohen--Macaulay模，由（16.4.10（i））得到。
\end{proof}

\egas{16.5.6}{推论}
在（16.5.5）的假设下，设 $(x_i)_{1\leq i\leq r}$ 为 $\mm$ 中元素的序列。以下条件等价：
\begin{enumerate}[label=\alph*)]
 \item $(x_i)$ 是 $M$-正则序列；
 \item $\dim(M/(x_1M+\cdots+x_rM))=\dim M-r$；
 \item $x_i$ 为 $M$ 的一个参数系的一部分。
\end{enumerate}
此外，当这些条件成立时，$N=M/(x_1M+\cdots+x_rM)$ 是Cohen--Macaulay模；故对每个 $\pp\in\Ass N$，有 $\dim(A/\pp)=\dim M-r$。
\begin{proof}
不需对 $M$ 加Cohen--Macaulay假设，就已知道b）与c）等价（原文引（16.3.6）），且a）推出c）（16.4.1）。为证明c）推出a）并证明最后的断言，对 $r$ 归纳，$r=0$ 时结论直接成立。

由于 $x_1$ 为参数系的一部分，有 $\dim(M/x_1M)=\dim M-1$（原文引（16.3.6））；由（16.5.5），$x_1$ 是 $M$-正则元素，且 $P=M/x_1M$ 为Cohen--Macaulay模。$(x_i)_{2\leq i\leq r}$ 为 $P$ 的参数系的一部分，所以归纳假设表明它是 $P$-正则序列，而且 $P/(x_2P+\cdots+x_rP)=N$ 是Cohen--Macaulay模；于是整个 $(x_i)_{1\leq i\leq r}$ 是 $M$-正则序列。
\end{proof}

\egas{16.5.7}{备注}
由（16.5.6），当 $M$ 为Cohen--Macaulay模时，$M$ 的参数系恰好就是由 $\mm$ 中元素组成的极大 $M$-正则序列。反过来，由（16.4.1），若某个非零有限型 $A$-模有一个 $M$-正则序列，同时又是其参数系，则 $M$ 为Cohen--Macaulay模。（16.5.6）中最后的性质也刻画Cohen--Macaulay模：

\egas{16.5.8}{命题}
设 $A$ 为诺特局部环，$M$ 为有限型 $A$-模。假定对于每个为 $M$ 的一个参数系的一部分的序列 $(x_i)_{1\leq i\leq r}$，以及每个属于
$\Ass N$ 的素理想 $\pp$，其中 $N=M/(x_1M+\cdots+x_rM)$，都有 $\dim(A/\pp)=\dim N$。则 $M$ 是Cohen--Macaulay模。
\begin{proof}
取 $M$ 的参数系 $(y_i)_{1\leq i\leq n}$；只须证明它为 $M$-正则序列。对 $n=\dim M$ 作归纳，$n=0$ 时直接成立。将假设用于 $r=0$，可知 $\Ass M$ 的所有元素 $\pp_j$ 满足 $\dim(A/\pp_j)=\dim M$。由于 $y_1$ 为参数系的一部分，由（16.3.7）、（16.3.4），$y_1$ 不属于任何 $\pp_j$，故为 $M$-正则元素。令 $P=M/y_1M$；则 $(y_i)_{2\leq i\leq n}$ 为 $P$ 的参数系，并且 $P$ 满足命题的假设。应用归纳假设，得到 $(y_i)_{2\leq i\leq n}$ 为 $P$-正则序列，命题得证。
\end{proof}

\egas{16.5.9}{命题}
设 $A$ 为诺特局部环，$M$ 为有限型Cohen--Macaulay $A$-模。取 $\pp\in\Supp M$，并置 $r=\dim M-\dim(M/\pp M)$。则存在由 $\pp$ 的元素组成的长度 $r$ 的 $M$-正则序列 $(x_i)_{1\leq i\leq r}$。对每个这样的序列，均有
\[
 \dim\bigl(M/(x_1M+\cdots+x_rM)\bigr)
 =\dim(M/\pp M)=\dim(A/\pp),
\]
且 $\pp$ 是 $\Ass(M/(x_1M+\cdots+x_rM))$ 的极小元。
\begin{proof}
先证明 $x_i$ 的存在性。$r=0$ 时无须证明；$r>0$ 时对 $r$ 归纳。对每个 $\qq_j\in\Ass M$，有 $\dim(A/\qq_j)=\dim M$（16.5.4）。假设 $\dim(M/\pp M)<\dim M$ 表明 $\pp$ 不能包含于任何 $\qq_j$，这些 $\qq_j$ 正是包含 $\Ann M$ 的极小素理想（16.1.2.2）、（16.1.7）。由素理想回避，$\pp$ 也不包含于它们的并（Bourbaki，\textit{交换代数}，第二章，§1，第1目，命题2）。所以存在 $x_1\in\pp$，为 $M$-正则元素（同书，第四章，§1，第1目，命题2的推论2）。

由（16.5.5），$N=M/x_1M$ 是维数 $\dim M-1$ 的Cohen--Macaulay模。因 $x_1\in\pp$，有 $N/\pp N=M/\pp M$，于是
$\dim N-\dim(N/\pp N)=r-1$。将归纳假设用于 $N$；若 $(x_i)_{2\leq i\leq r}$ 为由 $\pp$ 的元素组成的 $N$-正则序列，则 $(x_i)_{1\leq i\leq r}$ 即为所需 $M$-正则序列。

令 $P=M/(x_1M+\cdots+x_rM)$。因 $x_i\in\pp$，有 $P/\pp P=M/\pp M$；由（16.5.6），$\dim P=\dim(P/\pp P)$。又 $P$ 为Cohen--Macaulay模，故每个 $\pp'\in\Ass P$ 都满足 $\dim P=\dim(A/\pp')$（16.5.4）。由于 $\pp\in\Supp P$（$0_{\mathrm I}$，1.7.5），$\pp$ 包含某个 $\pp'\in\Ass P$（Bourbaki，同书，第四章，§1，第4目，定理2）。而
\[
 \dim P=\dim(P/\pp P)\leq\dim(A/\pp)
 \leq\dim(A/\pp')=\dim P,
\]
故必有 $\pp=\pp'$（16.1.2.2），完成证明。
\end{proof}

\egas{16.5.10}{推论}
在（16.5.9）的假设下：
\begin{enumerate}[label=(\roman*)]
 \item $M_{\pp}$ 是Cohen--Macaulay $A_{\pp}$-模；
 \item 有
 \begin{equation}
  \dim M=\dim_A(M/\pp M)+\dim_{A_{\pp}}M_{\pp}.
  \tag{16.5.10.1}
 \end{equation}
\end{enumerate}
\begin{proof}
沿用（16.5.9）的记号，令 $x'_i$ 为 $x_i$ 在 $A_{\pp}$ 中的像。$x'_i$ 属于 $A_{\pp}$ 的极大理想，并由平坦性组成 $M_{\pp}$-正则序列（15.1.14）。因此
\[
 \prof_{A_{\pp}}M_{\pp}\geq r
 =\dim M-\dim(M/\pp M)\geq\dim_{A_{\pp}}M_{\pp},
\]
最后一个不等式来自（16.1.8.1）。结合（16.4.5.1），这些数必全相等，推论得证。
\end{proof}

\egas{16.5.11}{推论}
设 $A$ 为诺特局部环。假定存在有限型Cohen--Macaulay $A$-模 $M$，其支集为整个 $\Spec A$。则对 $A$ 的每个素理想 $\pp$，有
\begin{equation}
 \dim A=\dim(A/\pp)+\dim A_{\pp}.
 \tag{16.5.11.1}
\end{equation}
事实上，$\Supp(M/\pp M)=\Spec(A/\pp)$（$0_{\mathrm I}$，1.7.5），且 $\Supp M_{\pp}=\Spec A_{\pp}$，所以（16.5.11.1）是（16.5.10.1）的特例。

\egas{16.5.12}{推论}
满足（16.5.11）条件的诺特局部环 $A$ 的每个商环，都是悬链环。特别地，局部Cohen--Macaulay环的每个商环都悬链。
\begin{proof}
只须证明 $A$ 本身悬链，即对于两个素理想 $\pp\subset\qq$，有（16.1.4.2）
\[
 \dim A_{\qq}=\dim A_{\pp}+\dim(A_{\qq}/\pp A_{\qq}).
\]
由（16.5.10（i）），$A_{\qq}$ 满足与 $A$ 相同的假设，所以这一等式就是将（16.5.11.1）用于 $A_{\qq}$ 及其素理想 $\pp A_{\qq}$。
\cmt{D4-A04}{作者正式勘误 \hyperlink{ErrIV:7}{Err$_{\mathrm{IV}}$，7}补充：式（16.5.11.1）等价于将 $A$ 换为 $A/\mathfrak r$ 后的同一公式，其中 $\mathfrak r$ 遍历 $A$ 的极小素理想。因此本证明可先归结到整环情形，再应用修正后的（16.1.4.2）。不应将（16.1.4）的初版缺条件误判为本推论结论错误。}
\end{proof}

\egap{16.5.13}
设 $A$ 为诺特环，$M$ 为有限型 $A$-模。若对 $A$ 的每个素理想 $\pp$，$M_{\pp}$ 都是Cohen--Macaulay $A_{\pp}$-模，则称 $M$ 为Cohen--Macaulay $A$-模。由（16.5.10），$A$ 局部时，这一定义与（16.5.1）一致。若每个 $A_{\pp}$ 都为Cohen--Macaulay环，则称 $A$ 为Cohen--Macaulay环。

若 $M$ 是Cohen--Macaulay $A$-模（分别地，$A$ 为Cohen--Macaulay环），则对任意乘法子集 $S$，$S^{-1}M$ 是Cohen--Macaulay $S^{-1}A$-模（分别地，$S^{-1}A$ 为Cohen--Macaulay环）。

% END continuations/iv0_16_45.tex


% BEGIN continuations/iv0_17_1.tex
% EGAIV.1.pdf, PDF pp.38--41, printed pp.39--42.
\section{正则环}
\subsection{正则环的定义}
\source{EGA IV 第一部分，印刷第39--42页；上传PDF第38--41页}

\egas{17.1.1}{命题}
设 $A$ 为维数 $n$ 的诺特局部环，$\mm$ 为其极大理想，$k=A/\mm$ 为其剩余域。以下条件等价：
\begin{enumerate}[label=\alph*)]
 \item 典范满的分次 $k$-模同态（15.1.1.1）
 \[
  \varphi:\Salg^\bullet(\mm/\mm^2)\longrightarrow\gr_\bullet A
 \]
 是双射；这里右端是 $A$ 关于 $\mm$-预进制滤过的相伴分次 $k$-模。
 \item $\rk_k(\mm/\mm^2)=n$。
 \item $\mm$ 有一个由 $n$ 个元素组成的生成系。
 \item $\mm$ 有一个生成系，它是 $A$-正则序列。
\end{enumerate}
\begin{proof}
由Nakayama引理，$\rk_k(\mm/\mm^2)$ 等于生成 $\mm$ 所需的最少元素数，故b）与c）等价。另一方面，若 $(x_i)_{1\leq i\leq m}$ 生成 $\mm$，且它们的模 $\mm^2$ 类 $\overline x_i$ 组成 $\mm/\mm^2$ 的 $k$-基，则
$\Salg^\bullet(\mm/\mm^2)\simeq k[T_1,\ldots,T_m]$。又任意有限型 $A$-模的 $\mm$-预进制滤过都是分离的（$0_{\mathrm I}$，7.3.5），所以由（15.1.9），a）与d）等价。由 $\mm$ 中元素组成的 $A$-正则序列长度至多为 $n$（16.4.1），故d）推出c）。

只须证明c）推出a）。简记
$S=\Salg^\bullet(\mm/\mm^2)=k[T_1,\ldots,T_n]$、$G=\gr_\bullet A$。考虑正合列
\[
 0\longrightarrow\jj\longrightarrow S\xrightarrow{\varphi}G\longrightarrow0,
\]
其中 $\jj=\Ker\varphi$ 为 $S$ 的分次理想。对每个整数 $s>0$，
\begin{equation}
 \binom{s+n-1}{n-1}=\len(S_s)=\len(G_s)+\len(\jj_s).
 \tag{17.1.1.1}
\end{equation}
若 $\jj\ne0$，则有 $\jj$ 中的齐次元 $u$，次数为 $h>0$。当 $s\geq h$ 时，$\jj_s$ 包含 $uS_{s-h}$；因为 $u\ne0$ 且 $S$ 为整环，$uS_{s-h}$ 与 $S_{s-h}$ 作为 $k$-向量空间同构。因此
$\len\jj_s\geq\len S_{s-h}=\binom{s-h+n-1}{n-1}$，从而
\begin{equation}
 \len G_s\leq\binom{s+n-1}{n-1}-\binom{s-h+n-1}{n-1}
 \qquad(s\geq h).
 \tag{17.1.1.2}
\end{equation}
右端是次数至多为 $n-2$ 的关于 $s$ 的多项式。可是
$G_s=\mm^s/\mm^{s+1}$；对于充分大的 $s$，$\len G_s$ 是次数恰为 $n-1$、首项系数为正的多项式（16.2.1）。这与（17.1.1.2）矛盾。
\end{proof}

\egas{17.1.2}{定义}
满足（17.1.1）的等价条件的诺特局部环称为\textbf{正则局部环}。

\egas{17.1.3}{推论}
正则局部环是整环、整闭环，也是Cohen--Macaulay环。
\begin{proof}
最后一个断言由（17.1.1 d））得到。另一方面，若 $A$ 正则，则 $\gr_\bullet A\simeq k[T_1,\ldots,T_n]$，故为完全整闭的整环。于是 $A$ 也具有这两种性质（Bourbaki，\textit{交换代数}，第三章，§2，第3目，命题1的推论；第五章，§1，第5目，命题15）。
\end{proof}

\egas{17.1.4}{例}
（i）维数零的正则局部环由（17.1.3）为整环，故必为域；反过来也成立。

（ii）维数1的诺特局部环正则，当且仅当它是离散赋值环。事实上，维数1的诺特局部环正则，表示其极大理想是主理想；结论由Bourbaki，\textit{交换代数}，第六章，§3，第6目，命题9得到。

（iii）设 $k$ 为域。形式幂级数环 $A=k[[T_1,\ldots,T_n]]$ 是维数 $n$ 的正则环。事实上，$T_i$（$1\leq i\leq n$）生成 $A$ 的极大理想 $\mm$，并组成 $A$-正则序列，因为
\[
 A/(T_1A+\cdots+T_iA)\simeq k[[T_{i+1},\ldots,T_n]].
\]

\egas{17.1.5}{命题}
诺特局部环 $A$ 正则，当且仅当其完备化 $\widehat A$ 正则。
\begin{proof}
$\widehat A$ 的极大理想为 $\mm\widehat A$，且已知
$\mm^h/\mm^{h+1}$ 与 $(\mm\widehat A)^h/(\mm\widehat A)^{h+1}$ 为同构的 $k$-向量空间。因此（17.1.1 a））对 $A$ 和 $\widehat A$ 是同一个条件。
\end{proof}

\egas{17.1.6}{定义}
给定极大理想为 $\mm$ 的诺特局部环 $A$，称一个生成 $\mm$ 的 $A$ 的参数系（原文引（16.3.5））为 $A$ 的\textbf{正则参数系}。

这样的参数系存在，当且仅当 $A$ 正则（17.1.1）。若 $(x_i)_{1\leq i\leq n}$ 为正则参数系，则它是 $\mm$ 的极小生成系，且其模 $\mm^2$ 类组成 $\mm/\mm^2$ 的 $k$-基。由（17.1.1 a））及（15.1.9），$(x_i)$ 为极大 $A$-正则序列，这也解释了术语的来源。

但应注意：在正则环中，一个同时为 $A$-正则序列的参数系，不一定是正则参数系；正则参数系各元素的 $k$ 次幂给出例子（15.1.20）。\cmt{D4-H07}{“正则序列构成的参数系”与“生成极大理想的正则参数系”须区分。例如 $k[[t]]$ 中 $(t^2)$ 为正则序列和参数系，却不生成极大理想 $(t)$。这里的幂次反例应取 $k\geq2$；$k=1$ 时当然仍为原正则参数系。}

\egas{17.1.7}{命题}
设 $A$ 为诺特局部环，$\mm$ 为其极大理想，$k=A/\mm$ 为其剩余域，$(x_i)_{1\leq i\leq r}$ 为 $\mm$ 中元素的序列，$\jj=x_1A+\cdots+x_rA$。以下条件等价：
\begin{enumerate}[label=\alph*)]
 \item $A$ 正则，且 $x_i$ 为 $A$ 的正则参数系的一部分；
 \item[$\mathrm a')$] $A$ 正则，且 $x_i$ 的模 $\mm^2$ 类 $\overline x_i$ 在 $k$ 上线性无关；
 \item[\(\mathrm b)\)] $x_i$ 为 $A$ 的参数系的一部分，且 $A/\jj$ 为正则环。
\end{enumerate}
此外，在这些条件下，$\jj$ 是素理想。
\begin{proof}
最后的断言由 $A/\jj$ 正则从而为整环（17.1.3）得到。令 $n=\dim A$。a）与a'）的等价是直接的：正则参数系的模 $\mm^2$ 类构成 $\mm/\mm^2$ 的一组 $k$-基。反过来，若 $\overline x_1,\ldots,\overline x_r$ 线性无关，可在 $\mm$ 中取另外 $n-r$ 个元素 $x_i$（$r+1\leq i\leq n$），使全部 $\overline x_i$ 构成 $\mm/\mm^2$ 的基；则 $(x_i)_{1\leq i\leq n}$ 为正则参数系。

为证明a）推出b），注意 $x_i$ 为参数系的一部分，故 $\dim(A/\jj)=n-r$（原文引（16.3.6））。令 $\mathfrak n=\mm/\jj$ 为 $A/\jj$ 的极大理想。因为 $\mathfrak n^2=(\mm^2+\jj)/\jj$，有正合列
\begin{equation}
 0\longrightarrow(\mm^2+\jj)/\mm^2\longrightarrow\mm/\mm^2
 \longrightarrow\mathfrak n/\mathfrak n^2\longrightarrow0.
 \tag{17.1.7.1}
\end{equation}
a）表明 $\rk_k((\mm^2+\jj)/\mm^2)=r$，于是
$\rk_k(\mathfrak n/\mathfrak n^2)=n-r=\dim(A/\jj)$。由（17.1.1 b））得到b）。

反过来，假定b）。仍有 $\dim(A/\jj)=n-r$。$A/\jj$ 正则，故 $\rk_k(\mathfrak n/\mathfrak n^2)=n-r$；又
$\rk_k((\mm^2+\jj)/\mm^2)\leq r$。由（17.1.7.1），
$\rk_k(\mm/\mm^2)\leq r+(n-r)=n$。结合（16.2.6）和（17.1.1），$A$ 正则。此外，$\rk_k(\mm/\mm^2)=n$ 迫使
$\rk_k((\mm^2+\jj)/\mm^2)=r$，所以 $\overline x_i$ 在 $k$ 上线性无关。
\end{proof}

\egas{17.1.8}{推论}
设 $A$ 为诺特局部环，$\mm$ 为其极大理想，$t\in\mm$。以下条件等价：
\begin{enumerate}[label=\alph*)]
 \item $A/tA$ 正则，且 $t$ 在 $A$ 中不是零因子；
 \item $A$ 正则且 $t\notin\mm^2$。
\end{enumerate}
这是（17.1.7）在 $r=1$ 时的情形，并用到 $\mm$ 中的非零因子为参数系的一部分（16.4.1）。

\egas{17.1.9}{推论}
设 $A$ 为正则局部环，$\mm$ 为其极大理想，$\jj\subset\mm$ 为理想。以下条件等价：
\begin{enumerate}[label=\alph*)]
 \item $A/\jj$ 正则；
 \item $\jj$ 由为 $A$ 的某个正则参数系的一部分的序列 $(x_i)_{1\leq i\leq r}$ 生成。
\end{enumerate}
\begin{proof}
b）推出a）已由（17.1.7）证明。反过来，假定 $A/\jj$ 正则，所以 $\jj$ 是素理想。令 $n=\dim A$、$n-r=\dim(A/\jj)$。沿用（17.1.7）的记号，由（17.1.7.1）得到
$\rk_k((\mm^2+\jj)/\mm^2)=r$。故可选取 $\jj$ 中元素 $x_1,\ldots,x_r$，为 $A$ 的正则参数系的一部分。

令 $\jj'=x_1A+\cdots+x_rA$。由（17.1.7），$\jj'$ 为素理想，且 $A/\jj'$ 是维数 $n-r$ 的正则环。但是
$A/\jj=(A/\jj')/(\jj/\jj')$，而 $\jj/\jj'$ 为整环 $A/\jj'$ 的素理想。因此由（16.1.2.2），$A/\jj$ 与 $A/\jj'$ 维数相等只能在 $\jj=\jj'$ 时发生。
\end{proof}

% END continuations/iv0_17_1.tex


% BEGIN continuations/iv0_17_2.tex
% EGAIV.1.pdf, PDF pp.41--45, printed pp.42--46.
\subsection{模的射影维数与内射维数回顾}
\source{EGA IV 第一部分，印刷第42--46页；上传PDF第41--45页}

\egap{17.2.1}
设 $A$ 为环，$M$ 为 $A$-模。回忆（M，第六章，§2），$M$ 的\textbf{射影维数}（分别地，\textbf{内射维数}），记作 $\pd M$ 或 $\pd_A M$（分别地，$\injdim M$ 或 $\injdim_A M$），是存在长度为 $n$ 的射影左解消（分别地，内射右解消）的最小 $n$，取值为非负整数或 $+\infty$。

等价地（同处），$n$ 是使任意 $A$-模 $N$ 均满足
$\Ext_A^i(M,N)=0$（分别地，$\Ext_A^i(N,M)=0$）对所有 $i>n$ 成立的最小数；只要求 $i=n+1$ 时消失也等价。

对 $M$ 的每个直和因子 $M'$，因 $\Ext_A^i(M',N)$ 为 $\Ext_A^i(M,N)$ 的直和因子，有
$\pd M'\leq\pd M$；同样有 $\injdim M'\leq\injdim M$。

还可如下刻画 $\pd M=n$（分别地，$\injdim M=n$）：$n$ 是使每个正合列
\[
 0\longrightarrow R\longrightarrow P_{n-1}\longrightarrow\cdots
 \longrightarrow P_0\longrightarrow M\longrightarrow0
\]
中，只要 $P_i$（$0\leq i<n$）都射影，便有 $R$ 射影的最小数；分别地，对于正合列
\[
 0\longrightarrow M\longrightarrow Q_0\longrightarrow\cdots
 \longrightarrow Q_{n-1}\longrightarrow C\longrightarrow0,
\]
只要各 $Q_i$ 都内射，便有 $C$ 内射。

\egas{17.2.2}{备注}
（i）$M$ 为射影模（分别地，内射模），等价于 $\pd M=0$（分别地，$\injdim M=0$）。

（ii）设 $T$ 为从 $A$-模范畴到一个阿贝尔范畴的协变加性函子。由（17.2.1）及导出函子的定义，若 $\pd M\leq n$（分别地，$\injdim M\leq n$），则对所有 $i>n$，有 $L_iT(M)=0$（分别地，$R^iT(M)=0$）。

（iii）若 $A$ 诺特且 $M$ 有限型，则（17.2.1）中最后的解释表明，当 $\pd M=n$ 时，$M$ 有一个长度 $n$ 的左解消，由有限型射影模构成。若 $A$ 还局部，则此解消可取为有限型自由模构成，因为有限型射影 $A$-模是自由的（Bourbaki，\textit{交换代数}，第二章，§3，第2目，命题5的推论2）。

\egas{17.2.3}{引理}
设 $A$ 为环，$M$ 为 $A$-模。$\injdim M\leq n$ 的充要条件是对每个单生成 $A$-模 $N$，有 $\Ext_A^{n+1}(N,M)=0$。
\begin{proof}
沿用（17.2.1）的记号，只须证明 $C$ 内射。对任意 $A$-模 $N$，$\Ext_A^{n+1}(N,M)\simeq\Ext_A^1(N,C)$（M，第五章，§7），所以对所有形如 $N=A/\jj$ 的模，均有 $\Ext_A^1(N,C)=0$。由Ext正合列，对 $A$ 的每个理想 $\jj$，典范同态
$\Hom_A(A,C)\to\Hom_A(\jj,C)$ 都满射。这表明 $C$ 内射（M，第一章，3.2）。
\end{proof}

\egas{17.2.4}{引理}
设 $A$ 为诺特环，$M$ 为有限型 $A$-模。$\pd M\leq n$ 的充要条件是对每个单生成 $A$-模 $N$，有 $\Ext_A^{n+1}(M,N)=0$。
\begin{proof}
已知 $\pd M\leq n$ 等价于对每个有限型 $A$-模 $N$，有
\[
 \Ext_A^{n+1}(M,N)=0
\]
（M，第六章，2.5）。为证明命题中的条件也充分，对 $N$ 的生成元个数 $m$ 归纳。取一个由 $m-1$ 个元素生成的子模 $N_1$，使 $N_2=N/N_1$ 单生成。正合列 $0\to N_1\to N\to N_2\to0$ 给出
\[
 \Ext_A^{n+1}(M,N_1)\longrightarrow\Ext_A^{n+1}(M,N)
 \longrightarrow\Ext_A^{n+1}(M,N_2).
\]
归纳假设表明，中间项为零。
\end{proof}

\egas{17.2.5}{推论}
设 $A$ 为诺特环，$M$ 为 $A$-模。则
\begin{equation}
 \injdim_A M=\sup_{\mm}\injdim_{A_{\mm}}M_{\mm}.
 \tag{17.2.5.1}
\end{equation}
若 $M$ 有限型，还有
\begin{equation}
 \pd_A M=\sup_{\mm}\pd_{A_{\mm}}M_{\mm},
 \tag{17.2.5.2}
\end{equation}
其中 $\mm$ 遍历 $A$ 的所有素理想，或只遍历所有极大理想。
\begin{proof}
若 $N$ 为有限型 $A$-模，因而有有限表示，则对任意乘法子集 $S$ 及 $i\geq0$，有
\[
 S^{-1}\Ext_A^i(N,M)
 \simeq\Ext_{S^{-1}A}^i(S^{-1}N,S^{-1}M).
\]
原文说，由平坦性，考虑 $M$ 的一个自由解消，并利用 $i=0$ 时上式成立（Bourbaki，\textit{交换代数}，第二章，§2，第7目，命题19），即可得到这一同构。\cmt{D4-T03}{原文这里写“$M$ 的自由解消”。要以自由解消计算 $\Ext_A^i(N,M)$，应解消第一变量 $N$；因 $A$ 诺特、$N$ 有限型，可取逐项有限自由的解消，从而局部化与各项Hom交换。报告给出具体核对。}
特别地，对任意素理想 $\mm$ 及理想 $\jj$，
\[
 \Ext_{A_{\mm}}^i(A_{\mm}/\jj A_{\mm},M_{\mm})
 =(\Ext_A^i(A/\jj,M))_{\mm}.
\]
结合（17.2.3）、$A_{\mm}$ 的每个理想都形如 $\jj A_{\mm}$，以及Bourbaki，同书，第二章，§3，第3目，定理1，得到（17.2.5.1）。用（17.2.4），并交换 $M,N$ 的地位，可同样证明（17.2.5.2）。
\end{proof}
因此，对于诺特环及有限型模，射影维数或内射维数的研究，由（17.2.5）归结到局部环。此时有：

\egas{17.2.6}{引理}
设 $A$ 为诺特局部环，$k$ 为其剩余域，$M$ 为有限型 $A$-模。$\pd M\leq n$ 的必要条件是 $\Tor_i^A(M,k)=0$ 对所有 $i>n$ 成立；充分条件是 $\Tor_{n+1}^A(M,k)=0$。
\begin{proof}
必要性是（17.2.2（ii））用于协变函子 $M\mapsto k\otimes_A M$ 的特例。为证明充分性，沿用（17.2.1）的记号，令 $P_i$ 有限型，只须证明 $R$ 射影。又 $\Tor_{n+1}^A(M,k)\simeq\Tor_1^A(R,k)$（M，第五章，§7）；由于 $R$ 有限型，$\Tor_1^A(R,k)=0$ 蕴含 $R$ 自由（Bourbaki，\textit{交换代数}，第二章，§3，第2目，命题5的推论2）。
\end{proof}

\egas{17.2.7}{推论}
在（17.2.6）的假设下，设 $x$ 是属于 $A$ 的极大理想的 $M$-正则元素。则
\cmt{D4-E02}{按（17.2.1）、（17.2.2）的约定，零模也是射影模，故其射影维数为0。于是本式须加 $M\ne0$；否则 $M=0$ 时成为 $0=1$。本处保留初版陈述并注明边界条件。非零模且射影维数无限的情形，也可由下面的短正合列验证。}
\begin{equation}
 \pd(M/xM)=\pd M+1.
 \tag{17.2.7.1}
\end{equation}
\begin{proof}
乘以 $x$ 给出的正合列 $0\to M\xrightarrow{x}M\to M/xM\to0$ 导出Tor正合列：对 $i\geq1$，
\[
 \begin{aligned}
 \Tor_i^A(M,k)&\xrightarrow{x}\Tor_i^A(M,k)
 \longrightarrow\Tor_i^A(M/xM,k)\\
 &\longrightarrow\Tor_{i-1}^A(M,k)
 \xrightarrow{x}\Tor_{i-1}^A(M,k).
 \end{aligned}
\]
对任意 $A$-模 $N$，乘以 $x$ 在 $M\otimes_A N$ 上的同态，既可由 $M$ 上的数乘产生，也可由 $N$ 上的数乘产生。由于 $x$ 在 $k$ 上的数乘为零，所有 $\Tor_i^A(M,k)$ 上由 $x$ 诱导的同态也为零。因此有短正合列
\[
 0\longrightarrow\Tor_i^A(M,k)\longrightarrow\Tor_i^A(M/xM,k)
 \longrightarrow\Tor_{i-1}^A(M,k)\longrightarrow0.
\]
若 $\pd M=n$，则由（17.2.6），$\Tor_n^A(M,k)\ne0$ 且
$\Tor_{n+1}^A(M,k)=\Tor_{n+2}^A(M,k)=0$。所以
$\Tor_{n+1}^A(M/xM,k)\ne0$、$\Tor_{n+2}^A(M/xM,k)=0$；再用（17.2.6），得到 $\pd(M/xM)=n+1$。
\end{proof}

\egas{17.2.8}{命题（M. Auslander）}
设 $A$ 为环。以下条件等价：
\begin{enumerate}[label=\alph*)]
 \item 每个 $A$-模的射影维数都至多为 $n$；
 \item[$\mathrm a')$] 每个有限型 $A$-模的射影维数都至多为 $n$；
 \item[\(\mathrm b)\)] 每个 $A$-模的内射维数都至多为 $n$；
 \item[\(\mathrm c)\)] 对每对 $A$-模 $M,N$，有 $\Ext_A^{n+1}(M,N)=0$；
 \item[$\mathrm c')$] 对每对 $A$-模 $M,N$，只要 $M$ 有限型，甚至仅单生成，就有
 \[\Ext_A^{n+1}(M,N)=0.\]
\end{enumerate}
这由（17.2.1）、（17.2.3）直接得到。

使这些等价条件成立的最小数 $n$（整数或 $+\infty$），称为 $A$ 的\textbf{整体上同调维数}，或简称\textbf{上同调维数}，记作 $\gldim A$。

\egas{17.2.9}{命题}
设 $A$ 为诺特环。以下条件等价：
\begin{enumerate}[label=\alph*)]
 \item $\gldim A\leq n$；
 \item 每个有限型 $A$-模的内射维数都至多为 $n$；
 \item 对任意两个有限型 $A$-模 $M,N$，有 $\Ext_A^{n+1}(M,N)=0$。
\end{enumerate}
这由定义和（17.2.4）直接得到。

\egas{17.2.10}{推论}
若 $A$ 为诺特环，则
\begin{equation}
 \gldim A=\sup_{\mm}\gldim A_{\mm},
 \tag{17.2.10.1}
\end{equation}
其中 $\mm$ 遍历 $A$ 的谱，或遍历其所有极大理想。

这由（17.2.9）、（17.2.5）得到。

\egas{17.2.11}{命题}
设 $A$ 为诺特局部环，$k$ 为其剩余域。$\gldim A\leq n$ 的必要条件是对所有 $i>n$，有 $\Tor_i^A(k,k)=0$；充分条件是 $\Tor_{n+1}^A(k,k)=0$。
\begin{proof}
结合（17.2.6），只须证明 $\gldim A\leq n$ 与 $\pd_A k\leq n$ 等价。第一条件由定义蕴含第二条件。反过来，若 $\pd_A k\leq n$，则由（17.2.2（ii）），对每个 $A$-模 $M$ 及 $i>n$，$\Tor_i^A(M,k)=0$，从而 $\pd M\leq n$，并由（17.2.8）得到结论。\cmt{D4-H08}{使用（17.2.6）作最后一步时，应先令 $M$ 遍历有限型模；这已足以应用（17.2.8 a'））。不能把“剩余域检测射影维数”不加条件地直接推广为任意模的判据。}
\end{proof}

\egas{17.2.12}{推论}
设 $A$ 为诺特环。$\gldim A\leq n$ 的必要条件是对 $A$ 的每个极大理想 $\mm$ 及 $i>n$，有
$\Tor_i^{A_{\mm}}(A/\mm,A/\mm)=0$；这些等式只在 $i=n+1$ 成立，便已充分。

这由（17.2.11）、（17.2.10）得到。

\egas{17.2.13}{命题}
设 $A,B$ 为诺特局部环，$\varphi:A\to B$ 为局部同态，使 $B$ 成为平坦 $A$-模。则
\begin{equation}
 \gldim A\leq\gldim B.
 \tag{17.2.13.1}
\end{equation}
\begin{proof}
设 $n=\gldim B$ 有限。只须证明，对任意两个有限型 $A$-模 $M,N$，$i>n$ 时有 $\Tor_i^A(M,N)=0$（17.2.11）。$B$ 是忠实平坦 $A$-模（$0_{\mathrm I}$，6.6.2），故这等价于证明（$0_{\mathrm I}$，6.4.1）
\begin{equation}
 \Tor_i^A(M,N)\otimes_A B=0.
 \tag{17.2.13.2}
\end{equation}
取 $M$ 的自由左解消 $L_\bullet=(L_j)$。由于 $B$ 平坦，$L_\bullet\otimes_A B$ 是 $B$-模 $M\otimes_A B$ 的自由左解消。\cmt{D4-T04}{初版在本段两次写“右解消”（résolution droite）。这里按Tor的构造排为“左解消”，并明示所作改正；这不是新的定理条件。}
又
\[
 (L_\bullet\otimes_A B)\otimes_B(N\otimes_A B)
 =(L_\bullet\otimes_A N)\otimes_A B.
\]
所以（17.2.13.2）的左端等于 $\Tor_i^B(M\otimes_A B,N\otimes_A B)$。由 $\gldim B=n$，此模在 $i>n$ 时为零（17.2.2（ii）），结论得证。
\end{proof}

\egap{17.2.14}
设 $(X,\OO_X)$ 为赋环空间，$\F$ 为 $\OO_X$-模。称
\[
 \sup_{x\in X}\pd_{\OO_{X,x}}\F_x
 \quad\text{（分别地，}\sup_{x\in X}\injdim_{\OO_{X,x}}\F_x\text{）}
\]
为 $\F$ 的\textbf{逐点射影维数}（分别地，\textbf{逐点内射维数}），记作 $\pd\F$（分别地，$\injdim\F$）。称
$\sup_{x\in X}\gldim\OO_{X,x}$ 为 $X$ 的\textbf{逐点上同调维数}，记作 $\gldim X$。

由（17.2.5）、（17.2.10），当 $A$ 为诺特环、$M$ 为有限型 $A$-模、$X=\Spec A$ 时，有
\[
 \begin{aligned}
 \pd\widetilde M&=\pd M,\\
 \injdim\widetilde M&=\injdim M,\\
 \gldim X&=\gldim A.
 \end{aligned}
\]
称 $\F_x$ 的射影维数（分别地，内射维数）为 $\F$ 在点 $x$ 的射影维数（分别地，内射维数）；称 $\OO_{X,x}$ 的上同调维数为 $X$ 在点 $x$ 的上同调维数。

\egas{17.2.15}{命题}
设 $X,Y$ 为茎环均是诺特局部环的赋环空间，$f:X\to Y$ 为平坦态射。若 $\gldim X\leq n$，则 $Y$ 在 $f(X)$ 的每一点的上同调维数都至多为 $n$。

这由（17.2.13）直接得到。\cmt{D4-H09}{应用（17.2.13）需各茎同态为局部同态。本句应在局部赋环空间及其局部态射的语境中使用；仅说茎环各自局部而允许任意赋环空间态射，不足以推出结论。此项列作适用条件提醒，不因省略重申约定而直接判为作者定理错误。}

% END continuations/iv0_17_2.tex


% BEGIN continuations/iv0_17_3.tex
% EGAIV.1.pdf, PDF pp.45--50, printed pp.46--51.
\subsection{正则环的上同调理论}
\source{EGA IV 第一部分，印刷第46--51页；上传PDF第45--50页}

\egas{17.3.1}{定理（Hilbert--Serre）}
设 $A$ 为诺特局部环。$A$ 的上同调维数有限，当且仅当 $A$ 正则；此时
\begin{equation}
 \gldim A=\dim A.
 \tag{17.3.1.1}
\end{equation}
\begin{proof}
先设 $A$ 正则。令 $\mm$ 为其极大理想，取正则参数系
$x=(x_i)_{1\leq i\leq n}$（17.1.6），并考虑外代数复形 $K_\bullet(x)$（III，1.1.1）。它由自由 $A$-模构成，且 $i>n$ 时 $K_i(x)=0$。因为 $(x_i)$ 为正则序列，$i>0$ 时
$H_i(K_\bullet(x))=0$（III，1.1.4及1.1.3.3），并且
\[
 H_0(K_\bullet(x))=A/(x_1A+\cdots+x_nA)=A/\mm=k.
\]
因此 $K_i(x)$ 构成 $k$ 的一个长度 $n$ 的自由解消。由于 $x_i\in\mm$，复形
$K_\bullet(x,k)=K_\bullet(x)\otimes_A k$ 的边缘算子在每个次数上均为零。所以按定义，
\[
 \Tor_i^A(k,k)\simeq\bigwedge\nolimits^i(k^n).
\]
式（17.3.1.1）即由（17.2.11）得到；这个结果实质上就是Hilbert的“合冲定理”。

现在证明：若 $A$ 为诺特局部环，极大理想为 $\mm$，并且 $\pd_A\mm$ 有限，则 $A$ 正则；这将完成（17.3.1）的证明。对 $n=\rk_k(\mm/\mm^2)$ 作归纳。$n=0$ 时 $\mm=0$，断言直接成立。需要以下引理。
\renewcommand{\qedsymbol}{}
\end{proof}

\egas{17.3.1.2}{引理（Nagata）}
若 $\mm-\mm^2$ 的每个元素都是 $A$ 中的零因子，则存在 $0\ne c\in A$，使 $c\mm=0$；也就是说，$\mm\in\Ass A$。
\begin{proof}
可限于 $\mm\ne0$，因而 $\mm\ne\mm^2$。假设表明，$\mm-\mm^2$ 包含于 $\Ass A$ 中各素理想 $\pp_i$ 的并（Bourbaki，\textit{交换代数}，第四章，§1，第1目，命题2的推论3）；故 $\mm$ 包含于 $\mm^2$ 与各 $\pp_i$ 的并。因 $\mm\not\subset\mm^2$，由素理想回避（同书，第二章，§1，第1目，命题2），$\mm$ 包含于某个 $\pp_i$。$\mm$ 极大，因而引理成立。
\end{proof}

\egas{17.3.1.3}{引理}
若 $a\in\mm-\mm^2$，则 $\mm/Aa$ 同构于 $\mm/a\mm$ 的一个直和因子。
\begin{proof}
可取一个包含 $a$ 的 $\mm$ 的极小生成系，其模 $\mm^2$ 类为 $\mm/\mm^2$ 的基。令 $\mathfrak b$ 为该生成系中除 $a$ 以外的元素所生成的理想。若 $xa\in\mathfrak b$，考察模 $\mm^2$ 的像可知 $x\in\mm$，故
$\mathfrak b\cap Aa\subset a\mm$。由包含 $\mathfrak b\to\mm$，取商得到单射
\[
 \mathfrak b/(\mathfrak b\cap Aa)\longrightarrow\mm/a\mm.
\]
又 $\mathfrak b+Aa=\mm$，复合同态
\[
 \mm/Aa=(\mathfrak b+Aa)/Aa
 \simeq\mathfrak b/(\mathfrak b\cap Aa)
 \longrightarrow\mm/a\mm\longrightarrow\mm/Aa
\]
为恒等，故引理成立。
\end{proof}

\egas{17.3.1.4}{引理}
设 $A$ 为诺特局部环，$\mm$ 为其极大理想，$E$ 为有限型且射影维数有限的 $A$-模。若 $a\in\mm$ 同时为 $A$-正则及 $E$-正则元素，则 $E/aE$ 作为 $A/aA$-模的射影维数有限，且至多为 $\pd_AE$。
\begin{proof}
对 $h=\pd_AE$ 归纳。$h=0$ 时，$E$ 射影，所以 $E/aE$ 为射影 $A/aA$-模，结论成立。存在正合列
$0\to N\to L\to E\to0$，其中 $L$ 自由、$N$ 有限型，且 $\pd_AN=h-1$（17.2.2（iii））。此外，
\[
 0\longrightarrow N/aN\longrightarrow L/aL\longrightarrow E/aE\longrightarrow0
\]
正合（15.1.18）。由于 $a$ 是 $A$-正则元素，$L$ 自由，所以 $a$ 也是 $N$-正则元素。归纳假设表明 $N/aN$ 的 $A/aA$-射影维数至多为 $h-1$；而 $L/aL$ 是自由 $A/aA$-模，所以 $E/aE$ 的射影维数至多为 $h$。
\end{proof}
现分两种情形，继续定理的证明。

\noindent\textbf{I.}
假定 $\mm-\mm^2$ 的每个元素都是 $A$ 中的零因子。由（17.3.1.2），存在 $0\ne c\in A$，使 $c\mm=0$。证明此时 $\mm=0$。

若不然，$\mm$ 不能是射影 $A$-模，否则它便自由（Bourbaki，\textit{交换代数}，第二章，§3，第2目，命题5的推论2），与 $c\mm=0$ 矛盾。因此 $\gldim A\geq1$。令 $d=\gldim A$。\cmt{D4-H10}{初版在此将同一个字母 $n$ 改用于 $\gldim A$，而前文 $n$ 表示嵌入维数 $\rk_k(\mm/\mm^2)$；为免混读，此处仅将前者改记为 $d$。其有限性来自 $\pd_A\mm<\infty$、正合列 $0\to\mm\to A\to k\to0$ 及（17.2.11），不依赖正在证明的正则性。}
由于 $\mm\in\Ass A$，存在正合列
\[
 0\longrightarrow k\longrightarrow A\longrightarrow E\longrightarrow0.
\]
因 $d\geq1$，Tor正合列给出
\[
 0\longrightarrow\Tor_{d+1}^A(E,k)\longrightarrow\Tor_d^A(k,k)\longrightarrow0.
\]
但 $\Tor_{d+1}^A(E,k)=0$（17.2.2（ii）），而 $\Tor_d^A(k,k)\ne0$（17.2.11），矛盾。

\noindent\textbf{II.}
因此可限于存在 $a\in\mm-\mm^2$ 为 $A$-正则元素的情形。它也为 $\mm$-正则元素。令 $A'=A/aA$，其极大理想为 $\mm'=\mm/aA$，且
$\rk_k(\mm'/\mm'^2)=n-1$。由（17.3.1.4），$\mm/a\mm$ 是射影维数有限的 $A'$-模；$\mm'=\mm/aA$ 是它的直和因子（17.3.1.3），所以也有有限射影维数（17.2.1）。由归纳假设，$A'$ 正则；再由（17.1.8），$A$ 正则。\hfill$\square$

\egas{17.3.2}{推论}
若 $A$ 为正则局部环，则对 $A$ 的每个素理想 $\pp$，$A_{\pp}$ 正则。

由（17.2.10），$\gldim A_{\pp}\leq\gldim A$，所以结论由（17.3.1）得到。

\egas{17.3.3}{命题}
设 $A,B$ 为诺特局部环，$\rho:A\to B$ 为局部同态；$\mm,\mathfrak n$ 分别为其极大理想，$k=A/\mm$、$k'=B/\mathfrak n$ 分别为剩余域。由此有典范 $k'$-线性同态
\begin{equation}
 \psi:(\mm/\mm^2)\otimes_k k'\longrightarrow\mathfrak n/\mathfrak n^2.
 \tag{17.3.3.1}
\end{equation}
\begin{enumerate}[label=(\roman*)]
 \item 若 $B$ 正则且为平坦 $A$-模，则 $A$ 正则。
 \item 以下条件等价：
 \begin{enumerate}[label=\alph*)]
  \item $B$ 正则且（17.3.3.1）中的 $\psi$ 单射；
  \item $A,B$ 都正则，且对 $A$ 的一个正则参数系 $(x_i)_{1\leq i\leq m}$，$\rho(x_i)$ 为 $B$ 的正则参数系的一部分；此时对 $A$ 的每个正则参数系都如此；
  \item $B$ 与 $B\otimes_A k$ 正则，且 $B$ 为平坦 $A$-模；
  \item $A$ 与 $B\otimes_A k$ 正则，且 $B$ 为平坦 $A$-模；
  \item $A$ 与 $B\otimes_A k$ 正则，且
  \begin{equation}
   \dim B=\dim A+\dim(B\otimes_A k).
   \tag{17.3.3.2}
  \end{equation}
 \end{enumerate}
\end{enumerate}
\begin{proof}
（i）由（17.2.13），$\gldim A\leq\gldim B$，故只须用（17.3.1）。

（ii）当 $A$ 正则且 $(x_i)$ 为其正则参数系时，由（15.1.21），$B$ 为平坦 $A$-模等价于 $y_i=\rho(x_i)$ 组成 $B$-正则序列，因为 $B\otimes_A k$ 是平坦 $k$-模。又 $\dim A=m$，$y_i$ 生成 $\mm B$，所以
$\dim(B/\mm B)=\dim B-m$，即（17.3.3.2），等价于 $(y_i)_{1\leq i\leq m}$ 为 $B$ 的参数系的一部分（16.3.7）。因此b）、d）、e）分别等价于：
\begin{enumerate}
 \item[b')] $A,B$ 正则，且 $(y_i)$ 为 $B$ 的正则参数系的一部分；
 \item[d')] $A$ 正则，$B/(\sum_i y_iB)$ 正则，且 $(y_i)$ 为 $B$-正则序列；
 \item[e')] $A$ 正则，$B/(\sum_i y_iB)$ 正则，且 $(y_i)$ 为 $B$ 的参数系的一部分。
\end{enumerate}
由（17.1.7），b'）与e'）等价；d'）推出e'）（16.4.1），且由b'）推出（17.1.7），所以三者等价。b）与d）的合取显然推出c），而c）由（i）推出d）；因此b）、c）、d）、e）等价。

b）也推出a）：$x_i$ 的模 $\mm^2$ 类为 $\mm/\mm^2$ 的 $k$-基，而 $y_i$ 的模 $\mathfrak n^2$ 类在 $k'$ 上线性无关。只须证明a）推出e）。令
$V=\mm/\mm^2$、$W=\mathfrak n/\mathfrak n^2$。有
$\psi(V\otimes_k k')=(\mathfrak n^2+\mm B)/\mathfrak n^2$。由（16.3.9），
\begin{equation}
 \dim B\leq\dim A+\dim(B\otimes_A k).
 \tag{17.3.3.3}
\end{equation}
由（16.2.6），
\begin{equation}
 \dim A\leq\rk_{k'}(V\otimes_k k').
 \tag{17.3.3.4}
\end{equation}
在局部环 $B\otimes_A k=B/\mm B$ 中，极大理想为
$\mathfrak n'=\mathfrak n/\mm B$，剩余域为 $k'$；且
\[
 \mathfrak n'/\mathfrak n'^2
 \simeq\mathfrak n/(\mathfrak n^2+\mm B)
 \simeq W/\psi(V\otimes_k k').
\]
再次用（16.2.6），
\begin{equation}
 \dim(B\otimes_A k)\leq\rk_{k'}W-\rk_{k'}\psi(V\otimes_k k').
 \tag{17.3.3.5}
\end{equation}
又 $B$ 正则，故 $\dim B=\rk_{k'}W$（17.1.1）。结合以上三式，
\[
 \begin{aligned}
 \rk_{k'}W
 &\leq\dim A+\dim(B\otimes_A k)\\
 &\leq\rk_{k'}W+\rk_{k'}(V\otimes_k k')
       -\rk_{k'}\psi(V\otimes_k k').
 \end{aligned}
\]
两端相等，等价于 $\psi$ 单射。故a）迫使（17.3.3.3）、（17.3.3.4）、（17.3.3.5）中的各不等式均为等式。而（17.3.3.4）与（17.3.3.5）取等号，分别表示 $A$ 与 $B\otimes_A k$ 正则（17.1.1）。因此a）推出e）。
\end{proof}

\egas{17.3.4}{命题}
设 $A$ 为维数 $n$ 的正则局部环，$\mm$ 为其极大理想。对每个非零有限型 $A$-模 $M$，有
\begin{equation}
 \prof M+\pd M=n.
 \tag{17.3.4.1}
\end{equation}
\begin{proof}
对 $m=\prof M$ 作归纳。若 $m=0$，则存在 $M$ 的子模 $N\simeq k=A/\mm$（16.4.6（i））。将Tor正合列用于
$0\to N\to M\to M/N\to0$，得到
\[
 \Tor_{n+1}^A(M/N,k)\longrightarrow\Tor_n^A(N,k)
 \longrightarrow\Tor_n^A(M,k).
\]
$A$ 正则，故 $\Tor_{n+1}^A(M/N,k)=0$（17.3.1.1）、（17.2.6）。所以 $\Tor_n^A(k,k)$ 同构于 $\Tor_n^A(M,k)$ 的一个子模。再次用这两条结果，得 $\Tor_n^A(M,k)\ne0$，故 $\pd M\geq n$（17.2.6）。而 $\pd M\leq n$（17.3.1.1），于是 $\pd M=n$。

若 $m>0$，取 $M$-正则元素 $x\in\mm$。有 $\prof(M/xM)=m-1$（原文引（16.4.6（i））），且 $\pd(M/xM)=\pd M+1$（17.2.7）。\cmt{D4-T05}{深度下降1的公式位于（16.4.6（ii）），不是原文所引的（i）。保留原引并在此指出。}
由归纳假设即得（17.3.4.1）。
\end{proof}

\egas{17.3.5}{推论}
（i）设 $A$ 为维数 $n$ 的正则局部环。非零有限型 $A$-模 $M$ 自由，当且仅当它是维数 $n$ 的Cohen--Macaulay $A$-模。

（ii）设 $A$ 为正则局部环，$B$ 为局部环，$\rho:A\to B$ 为局部同态，使 $B$ 成为有限型 $A$-模。$B$ 是自由 $A$-模，当且仅当 $B$ 是Cohen--Macaulay环且 $\rho$ 单射；也可将后一个条件换成 $\dim B=\dim A$（16.1.5）。
\begin{proof}
（i）由（17.3.4.1），并用到有限型 $A$-模射影与自由等价（Bourbaki，\textit{交换代数}，第二章，§3，第2目，命题5的推论2）。有限型自由模的特征正是 $\pd M=0$（17.2.2（i））。

（ii）$B$ 为Cohen--Macaulay环，等价于它为Cohen--Macaulay $A$-模（16.5.3）；因此在 $\dim B=\dim A$（16.1.5）的条件下，应用（i）即可。
\end{proof}

\egap{17.3.6}
若诺特环 $A$ 满足：对每个素理想 $\pp$，局部环 $A_{\pp}$ 都正则，则称 $A$ 为\textbf{正则环}。当 $A$ 本身局部时，由（17.3.2），这一定义等价于（17.1.2）。仍由（17.3.2），只须对每个极大理想 $\mm$ 验证 $A_{\mm}$ 正则。由定义，对任意乘法子集 $S$，$A$ 正则蕴含 $S^{-1}A$ 正则。

\egas{17.3.7}{命题}
若 $A$ 为诺特正则环，则每个多项式环 $A[T_1,\ldots,T_n]$ 都正则。
\begin{proof}
只须证明 $B=A[T]$ 正则。$B$ 为自由 $A$-模，所以对 $B$ 的每个素理想 $\qq$，$B_{\qq}$ 为平坦 $A_{\pp}$-模，其中 $\pp=\qq\cap A$（$0_{\mathrm I}$，6.3.1）。由（17.3.3），只须证明 $B_{\qq}/\pp B_{\qq}$ 正则。\cmt{D4-A05}{初版这里引（17.1.10），而本节并无此条；作者勘误清单三A部将它明确改为（17.3.3）。本处采用该正式改正。}
这个环是 $A_{\pp}[T]/\pp A_{\pp}[T]=k[T]$ 在一个素理想处的局部化，其中 $k=A_{\pp}/\pp A_{\pp}$。所以只须证明 $C=k[T]$ 正则。$C$ 为主理想环，其在素理想处的局部环，或是离散赋值环，或是域（在零理想处），故都正则（17.1.4）。
\end{proof}

\egas{17.3.8}{推论}
若 $A$ 为正则环，则每个形式幂级数环 $A[[T_1,\ldots,T_n]]$ 都正则。

令 $\jj$ 为多项式环 $A[T_1,\ldots,T_n]$ 中由 $T_i$ 生成的理想。多项式环由（17.3.7）正则，而形式幂级数环是它关于 $\jj$-预进制拓扑的完备化（Bourbaki，\textit{交换代数}，第三章，§2，第12目，例1）。结论由下述引理得到。

\egas{17.3.8.1}{引理}
设 $A$ 为正则环，$\jj$ 为 $A$ 的理想，$\widehat A$ 为 $A$ 关于 $\jj$-预进制拓扑的分离完备化。则 $\widehat A$ 正则。
\begin{proof}
由（17.3.6），只须对 $\widehat A$ 的每个极大理想 $\mathfrak n$，证明 $\widehat A_{\mathfrak n}$ 正则。已知（Bourbaki，\textit{交换代数}，第三章，§3，第4目，命题8），$\mathfrak n=\mm\widehat A$，其中 $\mm$ 为 $A$ 中包含 $\jj$ 的极大理想。典范同态 $A\to\widehat A$ 诱导单射
$A_{\mm}\to\widehat A_{\mathfrak n}$；后者的 $\mathfrak n\widehat A_{\mathfrak n}$-预进制拓扑在 $A_{\mm}$ 上诱导 $\mm A_{\mm}$-预进制拓扑，且 $A_{\mm}$ 对该拓扑在 $\widehat A_{\mathfrak n}$ 中稠密。

因此两个诺特局部环 $A_{\mm}$ 与 $\widehat A_{\mathfrak n}$ 的完备化典范同构。按假设，$A_{\mm}$ 正则，所以其完备化正则（17.1.5）；于是 $\widehat A_{\mathfrak n}$ 的完备化正则，再由（17.1.5），$\widehat A_{\mathfrak n}$ 正则。
\end{proof}

\egas{17.3.9}{推论}
设 $A$ 为诺特环，且为某个诺特正则环 $B$ 的商环。若 $C$ 为有限型 $A$-代数，则 $C$ 的每个分式环 $S^{-1}C$ 都是正则环的商环。
\begin{proof}
$C$ 是多项式环 $A[T_1,\ldots,T_n]$ 的商环。若 $A=B/\jj$，则
\[
 A[T_1,\ldots,T_n]
 =B[T_1,\ldots,T_n]/\jj B[T_1,\ldots,T_n].
\]
由（17.3.7），可归结到 $C$ 是正则环 $B'$ 的商环的情形。若 $S'$ 是 $S$ 在 $B'$ 中的逆像，则 $S'$ 为乘法子集，且 $S^{-1}C$ 为 $S'^{-1}B'$ 的商环。\cmt{D4-T06}{初版此句写成“$C$ 为 $S'^{-1}B'$ 的商环”；这里应为所讨论的 $S^{-1}C$。例如 $B'=C=\mathbf Z$、$S=\mathbf Z-\{0\}$ 时，$\mathbf Z$ 不是 $\mathbf Q$ 的商环，而 $S^{-1}C=\mathbf Q$ 是。正文采用所需对象，并在此记录原式。}
故最后只须知道：正则环的分式环仍正则；这已在（17.3.6）指出。
\end{proof}
正则环商环的重要性，除上述性质外，还在于它们是悬链环（16.5.12）。

在代数几何的应用中，一切重要的环都是正则环的商环。\cmt{D4-H11}{这是原著对应用范围的概括性叙述，不应读成“每一个环都是诺特正则环的商环”的无条件数学断言。原句作为原著叙述保留。}

% END continuations/iv0_17_3.tex


% BEGIN continuations/iv0_18_1.tex
% Source: EGAIV.1.pdf, PDF pp.50--53; printed pp.51--54.
\setcounter{section}{17}
\section{代数扩张的补充知识}
\source{EGA IV第一部分，印刷第51--68页；上传PDF第50--67页}
本节汇集若干关于环的函子性构造，这些构造将在§§19、20中反复使用；本节不包含非平凡的结果。

\subsection{增广环的逆像}
\egap{18.1.1}
给定一个环 $A$（不必交换），$A$-环范畴的对象是由环 $B$ 和环同态 $\rho:A\to B$ 组成的偶对 $(B,\rho)$。其态射（也称为 $A$-同态）是环同态 $u:B\to C$，满足：若 $\rho:A\to B$ 和 $\sigma:A\to C$ 分别是定义 $B,C$ 的 $A$-环结构的同态（称为结构同态），则下图交换：
\begin{equation}
\begin{tikzcd}[column sep=large]
 B\arrow[rr,"u"] && C\\
 & A\arrow[ul,"\rho"]\arrow[ur,"\sigma"'] &
\end{tikzcd}
\tag{18.1.1.1}
\end{equation}
$u$ 的核 $\mathfrak j$ 是一个双边理想，并经由 $\rho$ 成为 $A$-双模。

若 $f:A'\to A$ 是环同态，则对每个 $A$-环 $(B,\rho)$，$(B,\rho\circ f)$ 是一个 $A'$-环。若 $u:B\to C$ 是从 $A$-环 $(B,\rho)$ 到 $A$-环 $(C,\sigma)$ 的 $A$-同态，它也是从 $A'$-环 $(B,\rho\circ f)$ 到 $A'$-环 $(C,\sigma\circ f)$ 的 $A'$-同态。这样就定义了从 $A$-环范畴到 $A'$-环范畴的一个典范函子。

\egap{18.1.2}
设 $B,E,F$ 是三个 $A$-环，$f:E\to B$、$g:F\to B$ 是两个 $A$-同态。回顾：$E$ 与 $F$ 在 $B$ 上的\emph{纤维积}（相对于 $f,g$）记为 $E\times_BF$，是乘积环 $E\times F$ 中由满足 $f(x)=g(y)$ 的偶对 $(x,y)$ 组成的子环 $G$。投影 $\operatorname{pr}_1,\operatorname{pr}_2$ 在 $G$ 上的限制
\[
 p_1:G\longrightarrow E,\qquad p_2:G\longrightarrow F
\]
仍称为\emph{典范投影}。$G$ 的 $A$-环结构由同态 $a\mapsto(\rho(a),\sigma(a))$ 定义，其中 $\rho:A\to E$、$\sigma:A\to F$ 是结构同态；由（18.1.1），该同态确实取值于 $G$。

$E\times_BF$ 的特征性质如下：对每一对满足 $f\circ u=g\circ v$ 的 $A$-同态 $u:C\to E$、$v:C\to F$，恰有一个 $A$-同态 $w:C\to G$，使 $u=p_1\circ w$、$v=p_2\circ w$。也可以说，$E\times_BF$ 是由 $B,E,F$ 以及 $f,g$ 组成的射影系在 $A$-环范畴中的射影极限（$0_{\mathrm{III}}$，8.1.9）。

\egap{18.1.3}
令 $\mathfrak j=\Ker f$ 为 $E$ 的双边理想。$p_2:G\to F$ 的核 $\mathfrak j'$ 由 $(x,0)$（$x\in\mathfrak j$）组成；因此 $p_1$ 的限制 $i_1:\mathfrak j'\to\mathfrak j$ 是不要求含幺的环之间的同构，同时也是 $A$-双模同构。同样，若 $\mathfrak k=\Ker g$，则 $p_1:G\to E$ 的核 $\mathfrak k'$ 由 $(0,y)$（$y\in\mathfrak k$）组成；$p_2$ 的限制 $i_2:\mathfrak k'\to\mathfrak k$ 是同样意义下的同构。最后，$A$-同态 $\varphi=f\circ p_1=g\circ p_2:G\to B$ 的核为
\[
 \mathfrak j\times\mathfrak k=\mathfrak j'\oplus\mathfrak k',
\]
因而有交换图
\begin{equation}
\begin{tikzcd}[column sep=1.8em,row sep=1.5em]
 & 0\arrow[d] & 0\arrow[d] & 0\arrow[d]\\
 & \mathfrak j'\arrow[r]\arrow[d,equal] &
 \mathfrak j'\oplus\mathfrak k'\arrow[r]\arrow[d] &
 \mathfrak k\arrow[d]\\
 0\arrow[r] & \mathfrak j'\arrow[r]\arrow[d,"i_1"'] &
 G\arrow[r,"p_2"]\arrow[d,"p_1"]\arrow[dr,"\varphi",description] &
 F\arrow[d,"g"]\\
 0\arrow[r] & \mathfrak j\arrow[r] & E\arrow[r,"f"] & B
\end{tikzcd}
\tag{18.1.3.1}
\end{equation}
\cmt{E5-A01}{图中的 $p_2$ 已按作者正式勘误移至 $G\to F$ 旁，见第四卷第四部分勘误清单三A部，印刷第345页；本书书末亦收入该条。图中未假定 $f,g$ 满射，不能在各行右端擅加 $\to0$。}
\cmt{E5-D01}{本译图按本段的核计算重排：上行使用 $0\to\mathfrak j\,{}^{\prime}\to\mathfrak j\,{}^{\prime}\oplus\mathfrak k\,{}^{\prime}\to\mathfrak k\to0$，其中通过 $i_2$ 识别 $\mathfrak k\,{}^{\prime}$ 与 $\mathfrak k$。这是对初版左上角未标名箭头组合的明示重绘，不是初版图形的逐箭头复制；核查报告另作说明。}（18.1.2）、（18.1.3）的定义和结果立即推广到由一族 $A$-同态 $f_\alpha:E_\alpha\to B$ 所定义的任意族 $(E_\alpha)_{\alpha\in I}$ 的纤维积。这些结果的表述留给读者。

\egap{18.1.4}
依照（M，VIII）的术语，我们称配备了一个满的 $A$-同态 $f:E\to B$ 的 $A$-环 $E$ 为\emph{在 $B$ 上增广的 $A$-环}；$f$ 称为 $E$ 的\emph{增广}，其核 $\mathfrak j$ 称为\emph{增广理想}。若存在 $A$-环同态 $s:B\to E$，它是 $f$ 的右逆，即 $f\circ s=1_B$，则称这个增广 $A$-环\emph{平凡}。这时，$A$-双模的正合列
\[
 0\longrightarrow\mathfrak j\longrightarrow E\longrightarrow B\longrightarrow0
\]
分裂；也就是说，可以将 $A$-双模 $E$ 等同于 $B\times\mathfrak j$，且 $E$ 中的乘法写为
\[
 (b,z)(b',z')=(bb',\,bz'+zb'+zz'),
\]
其中经由 $s:B\to E$ 把 $\mathfrak j$ 看作 $B$-双模。

\egap{18.1.5}
沿用（18.1.2）的记号，并设 $f$ 满射，即 $f$ 使 $E$ 成为在 $B$ 上增广的 $A$-环。这时 $p_2:G\to F$ 也满射，即它在 $G$ 上定义了一个在 $F$ 上增广的 $A$-环结构；称其为增广环 $E$ 沿 $g:F\to B$ 的\emph{逆像}。

\egas{18.1.6}{命题}
增广 $A$-环 $E$ 沿 $g:F\to B$ 的逆像 $G=E\times_BF$ 平凡，当且仅当存在一个 $A$-同态 $u:F\to E$，使下图交换：
\[
\begin{tikzcd}[column sep=large]
 & F\arrow[dl,"u"']\arrow[dr,"g"] &\\
 E\arrow[rr,"f"'] && B.
\end{tikzcd}
\]
\begin{proof}
必要性由 $u=p_1\circ s$ 得到，其中 $s:F\to G$ 是 $p_2$ 的一个 $A$-同态右逆。反之，若存在满足命题条件的 $u$，则将纤维积的泛性质（18.1.2）应用于 $u:F\to E$ 与 $1_F:F\to F$，便得到 $p_2$ 的右逆 $s$。
\end{proof}
特别地，若 $E$ 是平凡增广环，则它的每个逆像也平凡。

\egap{18.1.7}
回到（18.1.2）、（18.1.3）的情形。$\mathfrak k$ 通过其 $F$-双模结构和环同态 $p_2:G\to F$，自然具有 $G$-双模结构。由于 $i_2$ 双射，还有一个单射 $\theta=j\circ i_2^{-1}:\mathfrak k\to G$，它是 $G$-双模同态，其中 $j:\mathfrak k'\hookrightarrow G$ 是包含映射。我们将反过来说明：若另假设 $f,g$ 都满射，即 $E,F$ 都是在 $B$ 上增广的 $A$-环，那么给出这样的同态 $\theta$，便可由在 $F$ 上增广的 $G$ 重建在 $B$ 上增广的 $E$。

确切地说，给定在 $B$ 上增广的 $A$-环 $F$，以及在 $F$ 上增广的 $A$-环 $G$；分别将其增广理想记为 $\mathfrak k$ 和 $\mathfrak j'$：
\begin{equation}
\begin{tikzcd}[column sep=1.9em,row sep=1.5em]
 &&&0\arrow[d] &\\
 &&&\mathfrak k\arrow[d,hook]\arrow[dl,"\theta"'] &\\
 0\arrow[r]&\mathfrak j'\arrow[r]&G\arrow[r,"h"]&F\arrow[r]\arrow[d,"g"]&0\\
 &&&B\arrow[d]&\\
 &&&0&
\end{tikzcd}
\tag{18.1.7.1}
\end{equation}
同态 $h$ 在 $F$ 的每个理想上，特别是在 $\mathfrak k$ 上，定义一个 $G$-双模结构。设 $\theta:\mathfrak k\to G$ 是使图（18.1.7.1）交换的 $G$-双模同态。于是 $\theta$ 单射，且 $\mathfrak j'\cap\theta(\mathfrak k)=0$。由于 $\theta(\mathfrak k)$ 是 $G$ 的双边理想，并且 $h(\theta(\mathfrak k))=\mathfrak k$，复合同态 $g\circ h:G\to B$ 分解为
\[
 G\xrightarrow{q}G/\theta(\mathfrak k)\xrightarrow{f}B,
\]
其中 $f$ 满射。再令 $E=G/\theta(\mathfrak k)$，则 $\mathfrak j'$ 在 $E$ 中的像 $\mathfrak j$ 为 $f$ 的核，因为 $g\circ h$ 的核为 $\mathfrak j'\oplus\theta(\mathfrak k)$；且典范同态 $q:G\to E$ 在 $\mathfrak j'$ 上的限制单射。

于是可以形成纤维积 $G'=E\times_BF$。两个 $A$-同态 $q:G\to E$ 和 $h:G\to F$ 满足 $f\circ q=g\circ h$，故由 $G'$ 的泛性质（18.1.2），它们定义唯一的 $A$-同态
\[
 u:G\longrightarrow G'.
\]
\cmt{E5-T01}{原版此处把泛性质产生的箭头印成 $u:G'\to G$，后面的相伴分次映射方向也随之反写。本译文明确改为 $u(g)=(q(g),h(g))$，因而下一句中的一级分次映射为 $\mathfrak j'\to\mathfrak j''$；报告给出不依赖滤过记号的直接双射验证。这里修订的是构造方向，不是否定重建结论。}
我们来说明 $u$ 双射。只须将其作为 $A$-双模同态证明这一点。注意 $u$ 与 $G$ 上由 $G,\mathfrak j'$ 组成的有限滤过，以及 $G'$ 上由 $G',\mathfrak j''=\Ker(G'\to F)$ 组成的有限滤过相容。由（18.1.3），$\gr_0u:F\to F$ 为恒等映射，$\gr_1u:\mathfrak j'\to\mathfrak j''$ 双射，因此 $u$ 本身双射（Bourbaki，\textit{交换代数}，第三章，§2，第8目，定理1的推论3）。

% END continuations/iv0_18_1.tex


% BEGIN continuations/iv0_18_2.tex
% Source: EGAIV.1.pdf, PDF pp.53--58; printed pp.54--59.
\subsection{以双模为核的环扩张}
\source{EGA IV第一部分，印刷第54--59页；上传PDF第53--58页}
\egap{18.2.1}
设 $E$ 是在 $B$ 上增广的 $A$-环，$f:E\to B$ 为增广，$\mathfrak j=\Ker f$ 为增广理想。若 $\mathfrak j^2=0$，则 $\mathfrak j$ 不仅是 $E$-双模，也是 $B$-双模，因为 $B\simeq E/\mathfrak j$。具体地，每个 $b\in B$ 可写成 $f(x)$，其中 $x\in E$。若 $z,z'\in\mathfrak j$，则
\[
 (x+z)z'=xz',\qquad z'(x+z)=z'x.
\]
因此 $xz'$（分别地，$z'x$）的值不依赖于所选的 $x\in f^{-1}(b)$，可以写作 $bz'$（分别地，$z'b$）；这就定义了所说的 $B$-双模结构。反过来，若 $\mathfrak j$ 配备了一个 $B$-双模结构，且对 $x\in E,z'\in\mathfrak j$ 有 $xz'=f(x)z'$、$z'x=z'f(x)$，则 $\mathfrak j^2=0$。

\egap{18.2.2}
称 $A$-双模同态的正合列
\[
 0\longrightarrow L\xrightarrow{j}E\xrightarrow{f}B\longrightarrow0
\]
为 $A$-环 $B$ 的一个\emph{以 $B$-双模 $L$ 为核的 $A$-扩张}，若 $E$ 是 $A$-环，$f$ 是 $A$-环同态，并且对 $x\in E,z\in L$ 有
\[
 j(f(x)z)=xj(z),\qquad j(zf(x))=j(z)x.
\]
由（18.2.1），$j(L)$ 是 $E$ 中平方为零的双边理想。略用术语时，也说 $E$ 是 $B$ 的一个以 $L$ 为核的扩张。若存在 $A$-环同构 $u:E\xrightarrow{\sim}E'$，使下图交换，则称 $B$ 的两个以 $L$ 为核的 $A$-扩张 $E,E'$ 为\emph{$A$-等价}；$u$ 也称为 $A$-扩张的一个\emph{$A$-等价}：
\begin{equation}
\begin{tikzcd}[column sep=1.7em]
0\arrow[r]&L\arrow[r,"j"]\arrow[d,equal]&E\arrow[r,"f"]\arrow[d,"u","\sim"']&B\arrow[r]\arrow[d,equal]&0\\
0\arrow[r]&L\arrow[r,"j'"']&E'\arrow[r,"f'"']&B\arrow[r]&0.
\end{tikzcd}
\tag{18.2.2.1}
\end{equation}

\egap{18.2.3}
若以 $B$-双模 $L$ 为核的 $A$-扩张 $E$ 作为在 $B$ 上增广的 $A$-环是平凡的（18.1.4），则称它为\emph{$A$-平凡}。在乘积 $A$-双模 $B\times L$ 上置
\[
 (x,s)(y,t)=(xy,xt+sy),
\]
即可定义一个 $A$-环结构。典范映射 $j:L\to B\times L$ 和 $f:B\times L\to B$ 定义了一个以 $L$ 为核的 $B$ 的 $A$-平凡扩张：典范映射 $g:B\to B\times L$ 就是 $f$ 的一个 $A$-同态右逆。这个扩张称为 $B$ 的以 $L$ 为核的\emph{标准平凡扩张}，记为 $D_B(L)$。$B$ 的每个以 $L$ 为核的 $A$-平凡扩张均与 $D_B(L)$ $A$-等价。

还应注意：$A$-环 $A$ 自身的每个以 $A$-双模为核的扩张，都必为 $A$-平凡。

\egap{18.2.4}
给定两个 $A$-扩张 $L\xrightarrow{j}E\xrightarrow{f}B$、$L'\xrightarrow{j'}E'\xrightarrow{f'}B'$。从前者到后者的\emph{态射}，按定义是一个由 $A$-双模同态组成的三元组 $(u,v,w)$，使下图交换：
\begin{equation}
\begin{tikzcd}[column sep=1.7em]
0\arrow[r]&L\arrow[r,"j"]\arrow[d,"w"']&E\arrow[r,"f"]\arrow[d,"u"]&B\arrow[r]\arrow[d,"v"]&0\\
0\arrow[r]&L'\arrow[r,"j'"']&E'\arrow[r,"f'"']&B'\arrow[r]&0,
\end{tikzcd}
\tag{18.2.4.1}
\end{equation}
其中 $u,v$ 是 $A$-环同态，且对 $z\in L,b\in B$ 有
\[
 w(bz)=v(b)w(z),\qquad w(zb)=w(z)v(b).
\]
换言之，$(v,w)$ 是从 $B$-双模 $L$ 到 $B'$-双模 $L'$ 的双同态。若 $(u',v',w')$ 是从 $L'\to E'\to B'$ 到另一 $A$-扩张 $L''\to E''\to B''$ 的态射，那么 $(u'\circ u,v'\circ v,w'\circ w)$ 仍是态射，这说明上述术语是适当的。图（18.2.4.1）中两个交换正方形的考察，将导出 $A$-环扩张上的两种运算。

\egap{18.2.5}
先考虑 $B'$ 的一个以 $L'$ 为核的 $A$-扩张
\[
 0\longrightarrow L'\xrightarrow{j'}E'\xrightarrow{f'}B'\longrightarrow0
\]
以及一个 $A$-环同态 $v:B\to B'$。令 $F=E'\times_{B'}B$ 为增广 $A$-环 $E'$ 沿 $v$ 的逆像（18.1.5），于是有交换图
\[
\begin{tikzcd}[column sep=1.7em]
0\arrow[r]&L_0\arrow[r]\arrow[d,"i"',"\sim"]&F\arrow[r,"p_2"]\arrow[d,"p_1"]&B\arrow[r]\arrow[d,"v"]&0\\
0\arrow[r]&L'\arrow[r,"j'"']&E'\arrow[r,"f'"']&B'\arrow[r]&0,
\end{tikzcd}
\]
其两行正合，$p_1,p_2$ 是典范同态。由（18.1.3），$i$ 双射；由定义（18.1.2），$L_0^2=0$。因此可将 $F$ 看作 $B$ 的以 $L'$ 为核的 $A$-扩张，称为扩张 $E'$ 沿 $v$ 的\emph{逆像}；这里借助环同态 $v$ 将 $L'$ 看作 $B$-双模。

纤维积对每个因子的函子性表明，若有扩张态射
\[
\begin{tikzcd}[column sep=1.7em]
0\arrow[r]&L'_1\arrow[r]\arrow[d,"h"]&E'_1\arrow[r]\arrow[d,"g"]&B'\arrow[r]\arrow[d,equal]&0\\
0\arrow[r]&L'_2\arrow[r]&E'_2\arrow[r]&B'\arrow[r]&0,
\end{tikzcd}
\]
则可导出逆像态射
\[
\begin{tikzcd}[column sep=1.2em]
0\arrow[r]&L'_1\arrow[r]\arrow[d,"h"]&E'_1\times_{B'}B\arrow[r]\arrow[d,"g\times_{B'}1_B"]&B\arrow[r]\arrow[d,equal]&0\\
0\arrow[r]&L'_2\arrow[r]&E'_2\times_{B'}B\arrow[r]&B\arrow[r]&0.
\end{tikzcd}
\]
特别地，若 $E'_1,E'_2$ 是 $B'$ 的两个以 $L'$ 为核的 $A$-等价扩张，则其沿 $v$ 的逆像也是 $A$-等价的。

由纤维积的定义，（18.2.4.1）中的任意 $A$-扩张态射都经由 $E'$ 沿 $v$ 的逆像分解。具体地，恰有一个 $A$-同态 $u_0:E\to F=E'\times_{B'}B$，使下图交换：
\[
\begin{tikzcd}[column sep=1.8em,row sep=1.9em]
0\arrow[r]&L\arrow[r,"j"]\arrow[d,"w_0"']&E\arrow[r,"f"]\arrow[d,"u_0"]&B\arrow[r]\arrow[d,equal]&0\\
0\arrow[r]&L_0\arrow[r]\arrow[d,"i"',"\sim"]&F\arrow[r,"p_2"]\arrow[d,"p_1"]&B\arrow[r]\arrow[d,"v"]&0\\
0\arrow[r]&L'\arrow[r,"j'"']&E'\arrow[r,"f'"']&B'\arrow[r]&0,
\end{tikzcd}
\]
其中 $w_0$ 是 $u_0$ 在 $L$ 上的限制，且 $p_1\circ u_0=u$、$i\circ w_0=w$。

\egap{18.2.6}
特别研究沿满的 $A$-同态的扩张逆像。设 $v:B\to B'$ 是满的 $A$-同态，$F$ 是 $B$ 的以 $B$-双模 $L$ 为核的一个 $A$-扩张，$\mathfrak k=\Ker v$ 是 $B$ 的双边理想；经由增广 $F\to B$，它可以看作 $F$-双模。由（18.1.7），每个使下图交换的 $F$-双模同态 $\theta:\mathfrak k\to F$
\[
\begin{tikzcd}[column sep=large]
 &\mathfrak k\arrow[dl,"\theta"']\arrow[dr,hook]&\\
 F\arrow[rr]&&B
\end{tikzcd}
\]
都确定 $B'=B/\mathfrak k$ 的一个以 $L$ 为核的扩张 $E'$，其沿 $v$ 的逆像与 $F$ 等价；而具有后一性质的每个 $A$-扩张 $E'$，在 $A$-等价意义下都由这种方式得到。

此外，两个从 $\mathfrak k$ 到 $F$ 的 $F$-双模同态 $\theta_1,\theta_2$ 给出 $B'$ 的以 $L$ 为核的 $A$-等价扩张，当且仅当存在 $A$-扩张 $F$ 到自身的一个 $A$-等价 $u$，使 $\theta_2=u\circ\theta_1$。
\cmt{E5-T02}{原版在“两个同态给出等价扩张”一句把商环写为 $B$；这里按该句的构造对象明确改为 $B'$。由 $F/\theta(\mathfrak k)$ 得到的是 $B/\mathfrak k=B'$ 的扩张，并非 $B$ 的扩张。若这样的 $\theta$ 存在，$\mathfrak k$ 对 $L$ 的左右作用均为零，故 $L$ 可下降为 $B'$-双模；不能在没有这一步时预先作该识别。}
这立即由（18.2.5）和 $F$ 到纤维积 $E'\times_{B'}B$ 的典范双射的定义（18.1.7）得出。

\egap{18.2.7}
现在考察（18.2.4.1）左侧的正方形。先回顾 $A$-双模范畴中的\emph{并和}：给定三个 $A$-双模 $X,Y,Z$ 以及两个 $A$-同态 $f:X\to Y$、$g:X\to Z$，并和 $Y\oplus_X Z$ 是由 $X,Y,Z$ 和 $f,g$ 组成的归纳系在 $A$-双模范畴中的归纳极限（$0_{\mathrm{III}}$，8.1.11）。它定义为乘积 $Y\times Z$ 对 $A$-子双模 $M$ 的商，其中 $M$ 是同态 $x\mapsto(f(x),-g(x))$ 的像。

其特征性质是：对每一对满足 $u\circ f=v\circ g$ 的 $A$-双模同态 $u:Y\to T$、$v:Z\to T$，恰有一个同态 $w:Y\oplus_X Z\to T$，满足 $u=w\circ j_1$、$v=w\circ j_2$，其中 $j_1:Y\to Y\oplus_X Z$、$j_2:Z\to Y\oplus_X Z$ 是典范映射。
\cmt{E5-T03}{原版第二个等式印作 $u=w\circ j_2$；按定义，$j_2$ 的定义域为 $Z$，故左端应为 $v$。本处已据泛性质作此明确订正。}

\egap{18.2.8}
现给定 $B$ 的一个以 $B$-双模 $L$ 为核的 $A$-扩张
\[
 0\longrightarrow L\xrightarrow{j}E\xrightarrow{f}B\longrightarrow0,
\]
以及一个 $B$-双模同态 $w:L\to L'$。令 $H$ 为并和 $A$-双模 $E\oplus_L L'$。说明如何在 $H$ 上赋予 $A$-环结构，并定义一个 $A$-扩张
\[
 0\longrightarrow L'\longrightarrow H\longrightarrow B\longrightarrow0.
\]
为此注意：经由增广同态 $E\to B$，$L'$ 具有 $E$-双模结构，因而可构造标准平凡 $A$-扩张 $G=D_E(L')$（18.2.3）。考虑映射
\[
 \theta:L\longrightarrow G,\qquad z\longmapsto(j(z),-w(z)).
\]
这是 $G$-双模同态，这里经由典范同态 $p:G\to E$ 将 $L$ 看作 $G$-双模。事实上，对 $(x,z')\in G,z\in L$，有
\[
 (j(z),-w(z))(x,z')=(j(z)x,\,j(z)z'-w(z)x).
\]
按照 $L'$ 的 $E$-双模结构的定义，$j(z)z'=f(j(z))z'=0$；又有 $j(z)x=j(zf(x))$、$w(z)x=w(zf(x))$。同样验证 $\theta$ 是左 $G$-模同态。于是可对交换图
\[
\begin{tikzcd}[column sep=1.7em,row sep=1.3em]
 &&&0\arrow[d]&\\
 &&&L\arrow[d,"j"]\arrow[dl,"\theta"']&\\
0\arrow[r]&L'\arrow[r]&G\arrow[r,"p"]&E\arrow[r]\arrow[d,"f"]&0\\
 &&&B\arrow[d]&\\
 &&&0&
\end{tikzcd}
\]
应用（18.1.7）。由于按定义 $H=G/\theta(L)$，所需断言立即由（18.1.7）得出。

称 $B$ 的以 $L'$ 为核的 $A$-扩张 $H$ 为\emph{由 $E$ 借助同态 $w:L\to L'$ 导出的扩张}。并和对各个加项的函子性还表明，若有扩张态射
\[
\begin{tikzcd}[column sep=1.7em]
0\arrow[r]&L\arrow[r]\arrow[d,equal]&E_1\arrow[r]\arrow[d,"g"]&B_1\arrow[r]\arrow[d,"h"]&0\\
0\arrow[r]&L\arrow[r]&E_2\arrow[r]&B_2\arrow[r]&0,
\end{tikzcd}
\]
则可典范地导出扩张态射
\[
\begin{tikzcd}[column sep=1.2em]
0\arrow[r]&L'\arrow[r]\arrow[d,equal]&E_1\oplus_LL'\arrow[r]\arrow[d,"g\oplus_L1_{L'}"]&B_1\arrow[r]\arrow[d,"h"]&0\\
0\arrow[r]&L'\arrow[r]&E_2\oplus_LL'\arrow[r]&B_2\arrow[r]&0.
\end{tikzcd}
\]
\cmt{E5-H01}{这两个图在改变商环时，须将 $L'$ 配备相容的 $B_2$-双模结构，并经 $h:B_1\to B_2$ 限制为 $B_1$-双模，且要求 $w$ 相容。否则仅有 $A$-双模的并和，还不足以调用上面的环扩张构造。商环固定为 $B$ 的应用没有这一歧义。}
特别地，若 $E_1,E_2$ 是 $B$ 的以 $L$ 为核的两个 $A$-等价扩张，则借助 $w$ 导出的两个以 $L'$ 为核的扩张也 $A$-等价。

若给定一个 $A$-扩张态射（18.2.4.1），它就经由上述由 $E$ 借助 $w:L\to L'$ 导出的扩张 $H$ 分解；此时通过 $v:B\to B'$ 将 $L'$ 看作 $B$-双模。事实上，并和的定义表明，存在唯一的 $A$-双模同态 $u_0:H\to E'$，使下图交换：
\[
\begin{tikzcd}[column sep=1.8em,row sep=1.8em]
0\arrow[r]&L\arrow[r,"j"]\arrow[d,"w_0"']&E\arrow[r,"f"]\arrow[d,"j_1"]&B\arrow[r]\arrow[d,equal]&0\\
0\arrow[r]&L'_0\arrow[r,"j_2"']\arrow[d,"j_0"',"\sim"]&H\arrow[r,"f_0"]\arrow[d,"u_0"]&B\arrow[r]\arrow[d,"v"]&0\\
0\arrow[r]&L'\arrow[r,"j'"']&E'\arrow[r,"f'"']&B'\arrow[r]&0,
\end{tikzcd}
\]
且 $u=u_0\circ j_1$、$w=j_0\circ w_0$。这里 $j_1,j_2$ 为典范同态。直接验证可知，$u_0$ 也是环同态。

\egas{18.2.9}{命题}
设 $B,B'$ 是两个 $A$-环，$L$ 是 $B$-双模，$L'$ 是 $B'$-双模，$v:B\to B'$ 是 $A$-环同态，$w:L\to L'$ 是 $A$-双模同态，且 $(v,w)$ 为双模的双同态。则存在唯一的 $A$-环同态 $u:D_B(L)\to D_{B'}(L')$，使下列两图交换：
\[
\begin{tikzcd}[column sep=large]
D_B(L)\arrow[r,"u"]&D_{B'}(L')\\
B\arrow[u,hook]\arrow[r,"v"']&B'\arrow[u,hook]
\end{tikzcd}
\qquad
\begin{tikzcd}[column sep=large]
D_B(L)\arrow[r,"u"]&D_{B'}(L')\\
L\arrow[u,hook]\arrow[r,"w"']&L'\arrow[u,hook]
\end{tikzcd}
\]
其中竖直箭头均为典范单射。
\begin{proof}
$u$ 只能是映射 $(x,s)\mapsto(v(x),w(s))$。其为 $A$-环同态直接由定义（18.2.3）验证。还应注意，$u$ 同样使下图交换：
\[
\begin{tikzcd}[column sep=large]
D_B(L)\arrow[r,"u"]\arrow[d]&D_{B'}(L')\arrow[d]\\
B\arrow[r,"v"']&B',
\end{tikzcd}
\]
这次竖直箭头为增广。
\end{proof}

\egas{18.2.10}{命题}
设 $B$ 是 $A$-环，$L$ 是 $B$-双模，$E$ 是 $B$ 的一个以 $L$ 为核的 $A$-扩张。令 $G$ 为从 $B$ 到 $E$ 的 $B$-环同态的集合（也就是说，增广 $E\to B$ 的 $A$-同态右逆的集合）。令 $G'$ 为从 $D_B(L)$ 到 $E$ 的 $A$-等价的集合。将每个 $g\in G$ 对应于
\[
 g':(x,s)\longmapsto g(x)+s
\]
定义了从 $G$ 到 $G'$ 的双射；其逆映射将 $g'\in G'$ 对应于 $B$-同态 $x\mapsto g'(x,0)$。
\cmt{E5-T04}{这里保留原版两处“$B$-同态”的措辞，但其精确定义以原文括号内说明为准：$g:B\to E$ 必须是满足 $f\circ g=1_B$ 的 $A$-环同态。选择 $g$ 以前并未给 $E$ 指定一个 $B$-环结构；不能把此处读成预先固定该结构的下范畴中的同态。式中将 $L$ 通过 $j$ 与其像等同。}
这由定义立即得到。

% END continuations/iv0_18_2.tex


% BEGIN continuations/iv0_18_3.tex
% Source: EGAIV.1.pdf, PDF pp.58--63; printed pp.59--64.
\subsection{\texorpdfstring{$A$}{A}-扩张类的群}
\source{EGA IV第一部分，印刷第59--64页；上传PDF第58--63页}
\egap{18.3.1}
固定一个 $A$-环 $B$ 和一个 $B$-双模 $L$。在 $B$ 的以 $L$ 为核的 $A$-扩张之间，“$E$ 与 $E'$ $A$-等价”是一个等价关系；对于这一关系，可以谈论 $B$ 的以 $L$ 为核的 $A$-扩张的等价类\emph{集合}。

只须注意：对每个 $x\in B$，若选取 $E$ 中映到 $x$ 的一个元素 $c_x$，则每个 $z\in E$ 都恰有一种写法 $c_x+t$，其中 $t\in L$。而且可写
\[
 c_x+c_y=c_{x+y}+\varphi(x,y),\qquad
 c_{xy}=c_xc_y+\psi(x,y),
\]
其中 $\varphi(x,y),\psi(x,y)\in L$。映射 $\varphi,\psi:B\times B\to L$ 必须满足一些表达 $E$ 为 $A$-环的条件，这里无须把这些条件写出。因此，每个这样的 $A$-扩张都与一个以 $B\times L$ 为底层集合的 $A$-扩张 $A$-等价，所说的结论随即得到。

\egap{18.3.2}
固定 $A$-环 $B$，暂将 $B$ 的以 $L$ 为核的 $A$-扩张类集合记为 $T(L)$。对每个 $B$-双模同态 $w:L\to L'$，将扩张 $E$ 的类对应于（18.2.8）中借助 $w$ 导出的扩张 $E\oplus_LL'$ 的类，便典范地定义了映射
\[
 T(w):T(L)\longrightarrow T(L').
\]
若 $w':L'\to L''$ 是另一个 $B$-双模同态，则还有
\begin{equation}
 T(w'\circ w)=T(w')\circ T(w).
 \tag{18.3.2.1}
\end{equation}
事实上，由归纳极限的一般性质（例如参见I，3.3.9），有 $A$-双模的典范同构
\[
 E\oplus_LL''\simeq (E\oplus_LL')\oplus_{L'}L''.
\]
直接验证可知，这确实是 $A$-扩张的 $A$-等价，故得到（18.3.2.1）。若 $\mathbf C_B$ 表示 $B$-双模范畴，则这样就定义了一个协变函子 $T:\mathbf C_B\to\mathbf{Ens}$。

\egap{18.3.3}
考虑一族 $B$-双模 $(L_\alpha)_{\alpha\in I}$ 及其乘积 $L=\prod_\alpha L_\alpha$。投影 $\operatorname{pr}_\alpha:L\to L_\alpha$ 定义映射 $T(\operatorname{pr}_\alpha):T(L)\to T(L_\alpha)$，因而得到典范映射
\begin{equation}
 \prod_\alpha T(\operatorname{pr}_\alpha):
 T(L)\longrightarrow\prod_{\alpha\in I}T(L_\alpha).
 \tag{18.3.3.1}
\end{equation}
我们说明它是双射。对每个 $\alpha\in I$，设 $E_\alpha$ 是 $B$ 的以 $L_\alpha$ 为核的一个 $A$-扩张，并令 $F$ 是各 $E_\alpha$ 在 $B$ 上的纤维积（18.1.3）。则 $F$ 是 $B$ 的以 $L=\prod_\alpha L_\alpha$ 为核的一个 $A$-扩张。若将每个 $E_\alpha$ 换成与之 $A$-等价的 $E'_\alpha$，则所得纤维积 $F'$ 与 $F$ $A$-等价。这样定义了一个映射 $\prod_\alpha T(L_\alpha)\to T(L)$；它与（18.3.3.1）互逆，故断言成立。

以后通过（18.3.3.1）典范地将 $T(L)$ 与 $\prod_
\alpha T(L_\alpha)$ 等同。若 $(L'_\alpha)_{\alpha\in I}$ 是另一族 $B$-双模，对每个 $\alpha$ 给定 $B$-双模同态 $w_\alpha:L_\alpha\to L'_\alpha$，并置 $w=\prod_\alpha w_\alpha:L\to L'=\prod_\alpha L'_\alpha$，则在上述识别下，$T(w)$ 与 $\prod_\alpha T(w_\alpha)$ 等同。

\egap{18.3.4}
对于一个 $B$-双模 $L$，加法是一个 $B$-双模同态 $s:L\times L\to L$；加法群的取负映射 $t:L\to L$ 也是如此。由（18.3.3），得到 $T(L)$ 上的合成律
\[
 T(s):T(L)\times T(L)\longrightarrow T(L).
\]
此合成律是交换群的群律，取逆映射为 $T(t)$；这由借助交换图定义群对象的方式得到（$0_{\mathrm{III}}$，8.2.5、8.2.6）。将所得交换群记为 $\Exan_A(B,L)$，称为\emph{$B$ 的以 $L$ 为核的 $A$-扩张类群}。

\egap{18.3.5}
令 $\mathbf K$ 为如下范畴。对象是三元组 $(A,B,L)$，其中 $A$ 是环，$B$ 是 $A$-环，$L$ 是 $B$-双模。态射是三元组 $(u,v,w)$，其中 $u:A'\to A$、$v:B'\to B$ 是环同态，使下图左侧的正方形交换：
\[
\begin{tikzcd}[column sep=large]
 A\arrow[r]&B&L\arrow[d,"w"]\\
 A'\arrow[u,"u"]\arrow[r]&B'\arrow[u,"v"]&L'.
\end{tikzcd}
\]
其中横箭头为结构同态；$w:L\to L'$ 是交换群同态，满足对所有 $z\in L,b'\in B'$ 有
\[
 w(v(b')z)=b'w(z),\qquad w(zv(b'))=w(z)b'.
\]
换言之，通过 $v$ 在 $L$ 上赋予 $B'$-双模结构后，$w$ 是 $B'$-双模同态。态射的合成定义为
\[
 (u',v',w')\circ(u,v,w)=(u\circ u',v\circ v',w'\circ w),
\]
其合理性可立即验证。我们的目的是证明
\begin{equation}
 (A,B,L)\longmapsto\Exan_A(B,L)
 \tag{18.3.5.1}
\end{equation}
是从 $\mathbf K$ 到交换群范畴 $\mathbf{Ab}$ 的一个协变函子。因此，对每个上述三元组 $(u,v,w)$，需要定义一个交换群同态
\[
 (u,v,w)_*:\Exan_A(B,L)\longrightarrow\Exan_{A'}(B',L').
\]
由 $\mathbf K$ 中态射的定义，可写
\[
 (u,v,w)=(1_{A'},1_{B'},w)\circ(1_{A'},v,1_L)\circ(u,1_B,1_L),
\]
其中第一个因子里的 $L$ 具有由 $v$ 定义的 $B'$-双模结构。我们先在 $u,v,w$ 中有两个为恒等同态时定义 $(u,v,w)_*$。

\egap{18.3.6}
首先，取 $(1_A,1_B,w)_*$ 为（18.3.2）中记作 $T(w)$ 的映射
\begin{equation}
 w_*:\Exan_A(B,L)\longrightarrow\Exan_A(B,L').
 \tag{18.3.6.1}
\end{equation}
它是群同态：这一性质可以用交换图表示，而这些图正是 $L,L'$ 上相应交换图经 $T$ 所得的像。

取 $(1_A,v,1_L)_*$ 为如下映射
\begin{equation}
 v^*:\Exan_A(B,L)\longrightarrow\Exan_A(B',L).
 \tag{18.3.6.2}
\end{equation}
若 $E$ 是 $B$ 的以 $L$ 为核的一个 $A$-扩张，由（18.2.5），$E\times_BB'$ 是 $B'$ 的以 $L$ 为核的一个 $A$-扩张；若以 $A$-等价的 $E'$ 替换 $E$，则 $E'\times_BB'$ 与 $E\times_BB'$ $A$-等价。定义 $v^*$ 将 $E$ 的类映到 $E\times_BB'$ 的类。对每个 $B$-双模同态 $w:L\to L'$，直接验证下图交换：
\begin{equation}
\begin{tikzcd}[column sep=large,row sep=large]
\Exan_A(B,L)\arrow[r,"v^*"]\arrow[d,"w_*"']&
\Exan_A(B',L)\arrow[d,"w_*"]\\
\Exan_A(B,L')\arrow[r,"v^*"']&\Exan_A(B',L').
\end{tikzcd}
\tag{18.3.6.3}
\end{equation}
在右列中，通过 $v$ 把 $L,L'$ 看作 $B'$-双模。分别以 $L\times L,L$ 代替 $L,L'$，以加法 $s$ 代替 $w$，便得到 $v^*$ 确实是群同态。

最后，取 $(u,1_B,1_L)_*$ 为映射
\begin{equation}
 u^*:\Exan_A(B,L)\longrightarrow\Exan_{A'}(B,L).
 \tag{18.3.6.4}
\end{equation}
它将 $B$ 的一个以 $L$ 为核的 $A$-扩张 $E$ 对应于经 $u$ 看作 $A'$-环的同一个环 $E$（18.1.1）；这显然是 $B$ 的以 $L$ 为核的一个 $A'$-扩张，其中 $B$ 也经 $u$ 看作 $A'$-环。$A$-等价也是 $A'$-等价，故得到（18.3.6.4）。对每个 $B$-双模同态 $w:L\to L'$，它仍使下图交换：
\begin{equation}
\begin{tikzcd}[column sep=large,row sep=large]
\Exan_A(B,L)\arrow[r,"u^*"]\arrow[d,"w_*"']&
\Exan_{A'}(B,L)\arrow[d,"w_*"]\\
\Exan_A(B,L')\arrow[r,"u^*"']&\Exan_{A'}(B,L').
\end{tikzcd}
\tag{18.3.6.5}
\end{equation}
因此仍可如上推出 $u^*$ 为群同态。

现在置
\[
 (u,v,w)_*=(1_{A'},1_{B'},w)_*\circ(1_{A'},v,1_L)_*\circ(u,1_B,1_L)_*.
\]
利用（18.3.6.3）、（18.3.6.5）的交换性，容易验证
\[
 (u\circ u',v\circ v',w'\circ w)_*=(u',v',w')_*\circ(u,v,w)_*,
\]
这就证明（18.3.5.1）是函子。
\cmt{E5-T05}{原版此处回指（18.3.4.1），但该编号未出现；所指的函子式为（18.3.5.1），本处明确据此前式订正。}

群同态（18.3.6.1）的存在特别表明：若 $E$ 是 $B$ 的以 $L$ 为核的平凡 $A$-扩张，那么（18.2.8）中所定义的 $E\oplus_LL'$ 也平凡；当然，这一点也易于直接验证。

另一方面，对于 $B$ 的中心 $Z$ 中每个元素 $z$，数乘
\[
 h_z:y\longmapsto yz=zy
\]
是 $B$-双模 $L$ 的自同态，故 $(h_z)_*$ 是交换群 $\Exan_A(B,L)$ 的自同态；由函子性，这些自同态在 $\Exan_A(B,L)$ 上定义一个典范的 $Z$-模结构。
\cmt{E5-H02}{这里忠实保留原文的 $yz=zy$。对通常意义下的任意 $B$-双模，这个等式不由 $z\in Z(B)$ 推出：左右中心作用仍可能不同。若将“双模”约定为中心作用对称，则该句成立；否则须区分左、右作用，再另行论证它们在扩张类群上的关系。报告给出扭曲双模的直接检验例，但不据此否定可另行构造的中心作用。}

\egap{18.3.7}
设 $A,A'$ 为两个环，$u:A'\to A$ 为同态，$B$ 为 $A$-环，$L$ 为 $B$-双模。群同态
\[
 u^*:\Exan_A(B,L)\longrightarrow\Exan_{A'}(B,L)
\]
的核，按定义由那些经 $u$ 看作 $A'$-扩张后为 $A'$-平凡的 $A$-扩张类组成。不致混淆时，将这个核记为 $\Exan_{A/A'}(B,L)$。

若 $\Lambda$ 是一个环，并给定 $\Lambda$-环的 $\Lambda$-同态交换图
\begin{equation}
\begin{tikzcd}[column sep=large]
 B'\arrow[r]&B\\
 A'\arrow[r]\arrow[u]&A\arrow[u],
\end{tikzcd}
\tag{18.3.7.1}
\end{equation}
则可典范地导出同态
\begin{equation}
 \Exan_{A/\Lambda}(B,L)\longrightarrow
 \Exan_{A'/\Lambda}(B,L)\longrightarrow
 \Exan_{A'/\Lambda}(B',L).
 \tag{18.3.7.2}
\end{equation}
它们来自下图的交换性：
\[
\begin{tikzcd}[column sep=1.3em,row sep=large]
\Exan_A(B,L)\arrow[r]\arrow[d]&\Exan_{A'}(B,L)\arrow[r]\arrow[d]&
\Exan_{A'}(B',L)\arrow[d]\\
\Exan_\Lambda(B,L)\arrow[r,equal]&\Exan_\Lambda(B,L)\arrow[r]&
\Exan_\Lambda(B',L),
\end{tikzcd}
\]
其中各箭头通过函子性由（18.3.7.1）中的箭头导出。

\egas{18.3.8}{命题}
设 $B$ 是一个环，$\mathfrak j$ 是其双边理想，$C=B/\mathfrak j$ 为商环；这时 $\mathfrak j/\mathfrak j^2$ 典范地具有 $C$-双模结构。对每个 $C$-双模 $L$，令 $\Hom_C(\mathfrak j/\mathfrak j^2,L)$ 表示 $C$-双模同态的加法群。将 $C$-双模同态 $w:\mathfrak j/\mathfrak j^2\to L$（它当然也是 $B$-同态）对应于由扩张 $B/\mathfrak j^2$ 借助 $w$ 导出的扩张
\[
 (B/\mathfrak j^2)\oplus_{\mathfrak j/\mathfrak j^2}L
\]
的类，便定义了交换群的一个典范同构
\begin{equation}
 \eta_L:\Hom_C(\mathfrak j/\mathfrak j^2,L)
 \xrightarrow{\sim}\Exan_B(C,L).
 \tag{18.3.8.1}
\end{equation}
其逆同构将 $C$ 的以 $L$ 为核的 $B$-扩张 $E$ 的类对应于同态 $w$；这里 $w$ 由如下条件确定：复合 $\mathfrak j\to\mathfrak j/\mathfrak j^2\xrightarrow{w}L$ 是结构同态 $B\to E$ 在 $\mathfrak j$ 上的限制。
\begin{proof}
令 $F=B/\mathfrak j^2$，它是 $C$ 的以 $\mathfrak j/\mathfrak j^2$ 为核的一个 $B$-扩张。对任意 $B$-扩张 $0\to L\xrightarrow{j}E\xrightarrow{p}C\to0$，结构同态 $f:B\to E$ 满足 $p\circ f:B\to C$ 为典范商同态。因为 $p\circ f$ 将 $\mathfrak j$ 映为零，$f(\mathfrak j)$ 包含在 $\Ker p=j(L)$ 中；而 $j(L)$ 平方为零，所以 $f(\mathfrak j^2)=0$。故 $f$ 分解为
\[
 B\longrightarrow B/\mathfrak j^2=F\xrightarrow{u}E.
\]
若 $w$ 是 $u$ 在 $\mathfrak j/\mathfrak j^2$ 上的限制，则有交换图
\[
\begin{tikzcd}[column sep=1.7em]
0\arrow[r]&\mathfrak j/\mathfrak j^2\arrow[r]\arrow[d,"w"']&F\arrow[r]\arrow[d,"u"]&C\arrow[r]\arrow[d,equal]&0\\
0\arrow[r]&L\arrow[r,"j"']&E\arrow[r,"p"']&C\arrow[r]&0.
\end{tikzcd}
\]
换言之，$(u,1_C,w)$ 是扩张态射（18.2.4）。由此导出扩张态射
\[
\begin{tikzcd}[column sep=1.8em]
0\arrow[r]&L\arrow[r]\arrow[d,equal]&E'\arrow[r]\arrow[d,"u'"]&C\arrow[r]\arrow[d,equal]&0\\
0\arrow[r]&L\arrow[r,"j"']&E\arrow[r,"p"']&C\arrow[r]&0,
\end{tikzcd}
\]
其中 $E'=F\oplus_{\mathfrak j/\mathfrak j^2}L$（18.2.8）。须证明 $u'$ 是 $B$-等价，这就证明 $\eta_L$ 双射。由（18.2.8）的构造，并考虑（18.1.7），只须证明扩张 $E$ 沿典范同态 $g:F\to C$ 的逆像 $G$ 是 $F$-平凡的。因为 $g$ 分解为 $F\xrightarrow{u}E\xrightarrow{p}C$，这由（18.1.6）得到。

还须证明 $\eta_L$ 为群同态。对每个 $B$-同态 $h:L\to L'$，下图直接由定义交换：
\[
\begin{tikzcd}[column sep=large,row sep=large]
\Hom_C(\mathfrak j/\mathfrak j^2,L)\arrow[r,"\eta_L"]\arrow[d,"{\Hom(1,h)}"']&
\Exan_B(C,L)\arrow[d,"h_*"]\\
\Hom_C(\mathfrak j/\mathfrak j^2,L')\arrow[r,"\eta_{L'}"']&\Exan_B(C,L').
\end{tikzcd}
\]
将此用于定义加法的同态 $L\times L\to L$ 即得结论。
\end{proof}

% END continuations/iv0_18_3.tex


% BEGIN continuations/iv0_18_4.tex
% Source: EGAIV.1.pdf, PDF pp.63--65; printed pp.64--66.
\subsection{代数的扩张}
\source{EGA IV第一部分，印刷第64--66页；上传PDF第63--65页}
\egap{18.4.1}
设 $A$ 为交换环。$A$-代数范畴是 $A$-环范畴的满子范畴，其特征是结构同态 $\rho:A\to B$ 的像包含于 $B$ 的中心。

若 $B$ 是 $A$-代数，$L$ 是 $B$-双模，则标准平凡 $A$-扩张 $D_B(L)$ 是 $A$-代数。
\cmt{E5-H03}{按允许左右作用不同的一般双模定义，此断言须增加 $\rho(a)\ell=\ell\rho(a)$（$a\in A,\ell\in L$），即 $L$ 对 $A$ 的作用对称。若“$B$-双模”在此已隐含“在 $A$ 上的双模”，则是须明说的约定；若按前文任意环的双模理解，字面断言有反例。报告取 $A=B=\mathbb Q[t]$、右作用经 $t\mapsto t+1$ 扭曲，直接验证 $D_B(L)$ 不是该结构下的 $A$-代数。不把这一约定问题说成交换代数扩张理论失效。}
同样，在（18.2.5）的构造中，若 $B,B',E'$ 都是 $A$-代数，则 $E'\times_{B'}B$ 也是 $A$-代数；在（18.2.8）的构造中，若 $B,E$ 都是 $A$-代数，则 $E\oplus_LL'$ 也是 $A$-代数。
最后，若 $E$ 是 $B$ 的以 $B$-双模 $L$ 为核的一个 $A$-扩张，并且是 $A$-代数，则与 $E$ $A$-等价的每个 $A$-扩张也是 $A$-代数。

由定义（18.3.4）便得：那些作为环是 $A$-代数的 $A$-扩张，其等价类在 $\Exan_A(B,L)$ 中形成一个子群，记为 $\Exal_A(B,L)$。令 $\mathbf K'$ 是（18.3.5）中范畴 $\mathbf K$ 的满子范畴，其对象 $(A,B,L)$ 满足 $A$ 交换、$B$ 是 $A$-代数。上述讨论说明
\[
 (A,B,L)\longmapsto\Exal_A(B,L)
\]
是从 $\mathbf K'$ 到 $\mathbf{Ab}$ 的协变函子。将（18.3.7）中的 $\Exan$ 全部换成 $\Exal$，结果保持不变。\cmt{E5-H04}{本段关于子群和函子的表述也须连同E5-H03的相容条件理解；特别是（18.2.8）的 $L'$ 也须对 $A$ 的作用对称。本译文保留原段组织，并在此明示这一适用界限。}

\egap{18.4.2}
仍设 $A$ 为交换环。在（18.4.1）中将“代数”换成“交换代数”，将“双模”换成“模”，上述讨论仍成立。若 $B$ 是交换 $A$-代数，$L$ 是 $B$-模，则那些是交换 $A$-代数的 $A$-扩张（等价地，是交换 $A$-环的扩张），其等价类在 $\Exal_A(B,L)$ 中形成子群，记为 $\Exalcom_A(B,L)$。若 $\mathbf K''$ 是 $\mathbf K'$ 中由 $A,B$ 交换且 $L$ 为 $B$-模的三元组组成的满子范畴，则
\[
 (A,B,L)\longmapsto\Exalcom_A(B,L)
\]
仍是从 $\mathbf K''$ 到 $\mathbf{Ab}$ 的协变函子。也可在（18.3.7）中将 $\Exan$ 全部换成 $\Exalcom$。最后，若在（18.3.8）中假设 $B$ 交换且 $L$ 为 $C$-模，则同一论证给出典范同构
\begin{equation}
 \Hom_C(\mathfrak j/\mathfrak j^2,L)\xrightarrow{\sim}\Exalcom_B(C,L),
 \tag{18.4.2.1}
\end{equation}
其中左端是 $C$-模同态群。

\egap{18.4.3}
设 $A$ 为交换环。一类重要的 $A$-代数扩张如下：$E$ 是 $A$-代数 $B$ 的以 $B$-双模 $L$ 为核的 $A$-代数扩张，并且 $E$ 作为 $A$-模，是 $A$-模 $B$ 的以 $A$-模 $L$ 为核的平凡扩张。也就是说，$A$-模正合列
\[
 0\longrightarrow L\longrightarrow E\longrightarrow B\longrightarrow0
\]
分裂。\cmt{E5-T06}{原版此处把正合列印作 $0\to E\to L\to B\to0$。按本节扩张的定义及紧接着的 $E\simeq B\times L$，前两个对象应交换；本处明确订正。}
这时称 $E$ 为 $B$ 的以 $L$ 为核的\emph{Hochschild扩张}。当 $B$ 是射影 $A$-模时总是如此，特别地，$A$ 是交换域时总是如此。

作为 $A$-模，可以把 $E$ 与 $B\times L$ 等同；$E$ 中的乘法则写为
\[
 (x,s)(y,t)=(xy,xt+sy+f(x,y)),\qquad f(x,y)\in L.
\]
写出此乘法在 $B\times L$ 上定义 $A$-代数结构的条件，可得（M，XIV，2）：$f$ 必须是从 $B\times B$ 到 $L$ 的 $A$-双线性映射，且
\begin{equation}
 f(xy,z)+f(x,y)z=xf(y,z)+f(x,yz).
 \tag{18.4.3.1}
\end{equation}
换言之，$f$ 是 $B$ 上取值于 $L$ 的一个Hochschild意义下的 $2$-余循环。扩张 $E$ 为 $A$-平凡，当且仅当
\begin{equation}
 f(x,y)=xg(y)-g(xy)+g(x)y,
 \tag{18.4.3.2}
\end{equation}
其中 $g:B\to L$ 是 $A$-线性映射；也就是说，$f$ 必须是Hochschild意义下的一个 $2$-余边界。于是，$B$ 的以 $L$ 为核的Hochschild扩张类形成 $\Exal_A(B,L)$ 的一个子群，同构于Hochschild上同调群 $H_A^2(B,L)$。
\cmt{E5-H05}{这里的分裂只是 $A$-模分裂，不先验保持单位。在未经归一化的坐标中，单位可为 $(1,-f(1,1))$，不必是 $(1,0)$。这不反驳原段：归一化余循环或选取保持单位的分裂后可写成常用形式；报告说明两个写法的关系，不向原文添入未给出的归一化假设。}

若 $B$ 是交换 $A$-代数，$L$ 是 $B$-模，则 $B$ 的以 $L$ 为核的交换Hochschild扩张对应于\emph{对称}的 $2$-余循环，即满足 $f(y,x)=f(x,y)$ 的那些余循环。这些交换Hochschild扩张的类在 $\Exalcom_A(B,L)$ 中形成子群，同构于对称 $2$-余循环群在 $H_A^2(B,L)$ 中的像；将这一子群记为 $H_A^2(B,L)^s$。

\egap{18.4.4}
现在限于 $A,B$ 交换且 $L$ 为 $B$-模的情形，并回顾此时Hochschild群 $H_A^i(B,L)$ 的一个等价定义（M，IX，6）。考虑 $B$-模复形 $P_\bullet=(P_n)_{n\geq0}$，其中
\[
 P_n=B^{\otimes_A(n+1)}=B\otimes_AB\otimes_A\cdots\otimes_AB
 \quad\text{（共 $n+1$ 个因子）},
\]
其 $B$-模结构由
\[
 x(y_1\otimes\cdots\otimes y_{n+1})=(xy_1)\otimes y_2\otimes\cdots\otimes y_{n+1}
\]
定义。次数为 $-1$ 的微分 $d_n:P_n\to P_{n-1}$ 定义为
\[
\begin{aligned}
 d_n(x_1\otimes\cdots\otimes x_{n+1})
 ={}&(x_1x_2)\otimes x_3\otimes\cdots\otimes x_{n+1}\\
 &-x_1\otimes(x_2x_3)\otimes\cdots\otimes x_{n+1}+\cdots\\
 &+(-1)^{n-1}x_1\otimes\cdots\otimes x_{n-1}\otimes(x_nx_{n+1})\\
 &+(-1)^n(x_{n+1}x_1)\otimes x_2\otimes\cdots\otimes x_n.
\end{aligned}
\]
因为 $B$ 交换，这确实是 $B$-线性的。

这个复形的取值于 $L$ 的 $2$-余循环，是满足 $h(d_3(x\otimes y\otimes z\otimes t))=0$ 的 $B$-线性映射 $h:B\otimes_AB\otimes_AB\to L$。由于
\[
 h(x\otimes y\otimes z)=xh(1\otimes y\otimes z),
\]
余循环 $h$ 由 $A$-双线性映射
\[
 (y,z)\longmapsto f(y,z)=h(1\otimes y\otimes z)
\]
确定；将前述关于 $h$ 的条件写出，就对 $f$ 重新得到（18.4.3.1）。同样，$2$-余边界是形如
\[
 x\otimes y\otimes z\longmapsto h'(d_2(x\otimes y\otimes z))
\]
的映射，其中 $h':B\otimes_AB\to L$ 是线性映射。这里 $h'$ 仍由 $A$-线性映射 $g:y\mapsto h'(1\otimes y)$ 确定；并有
\[
 h'(d_2(1\otimes x\otimes y))=xg(y)-g(xy)+yg(x),
\]
这就重新得到（18.4.3.2）。\cmt{E5-T07}{原版称 $g$ 为“$B$-线性”，本处改为“$A$-线性”。$h'$ 对第一个张量因子 $B$-线性，并不使 $y\mapsto h'(1\otimes y)$ 对 $B$-线性；例如在 $A=\mathbb Q,B=L=\mathbb Q[t]$ 时取 $h'(x\otimes y)=x\,dy/dt$，便得到仅 $A$-线性的导数映射。此修订与前段（18.4.3.2）的条件一致。}
对每个 $i$ 作同样处理，可见
\begin{equation}
 H_A^i(B,L)=H^i\bigl(\Hom_B(P_\bullet,L)\bigr).
 \tag{18.4.4.1}
\end{equation}

\egap{18.4.5}
在（18.4.4）的条件下，也可用同样方式解释群 $H_A^2(B,L)^s$（18.4.3）。为此，将 $P_\bullet$ 在次数 $3$ 处修改，考虑新复形
\[
 P'_\bullet:\quad P'_3\xrightarrow{d'_3}P_2\xrightarrow{d_2}P_1.
\]
取 $P'_3=P_3\oplus(B\otimes_AB\otimes_AB)$，令 $d'_3$ 在 $P_3$ 上与 $d_3$ 一致，并在 $B\otimes_AB\otimes_AB$ 上由
\[
 d'_3(x\otimes y\otimes z)=x\otimes y\otimes z-x\otimes z\otimes y
\]
给出。关系 $d_2d'_3=0$ 由 $B$ 的交换性得到。用前面的记号，$P'_\bullet$ 的一个 $2$-余循环现在对应于一个对称且满足（18.4.3.1）的 $A$-双线性映射 $f:B\times B\to L$；因此
\[
 H_A^2(B,L)^s=H^2\bigl(\Hom_B(P'_\bullet,L)\bigr).
\]

\egap{18.4.6}
在特殊情形下，若给定一个交换域 $k$ 及其扩域 $K$，则对任意 $K$-向量空间 $L,M$ 都有 $\Ext_K^1(M,L)=0$；所以（M，VI，3.3 a））有典范同构
\begin{equation}
 H_k^i(K,L)\simeq\Hom_K(H_i(P_\bullet),L),
 \tag{18.4.6.1}
\end{equation}
以及
\begin{equation}
 H_k^2(K,L)^s\simeq\Hom_K(H_2(P'_\bullet),L).
 \tag{18.4.6.2}
\end{equation}

% END continuations/iv0_18_4.tex


% BEGIN continuations/iv0_18_5.tex
% Source: EGAIV.1.pdf, PDF pp.65--67; printed pp.66--68.
\subsection{拓扑环的情形}
\source{EGA IV第一部分，印刷第66--68页；上传PDF第65--67页}
\egap{18.5.1}
设 $A,B$ 是两个具有线性拓扑的拓扑环，$\rho:A\to B$ 是连续同态，$L$ 是拓扑 $B$-双模，并假设 $B$ 中存在开双边理想 $\mathfrak k_0$，使 $\mathfrak k_0L=L\mathfrak k_0=0$。于是，对每个开理想 $\mathfrak k\subset\mathfrak k_0$，都可将 $L$ 看作 $(B/\mathfrak k)$-双模。

令 $\mathfrak k\subset\mathfrak k_0$ 为 $B$ 的开双边理想，$\mathfrak j$ 为 $A$ 的开双边理想，且 $\rho(\mathfrak j)\subset\mathfrak k$；于是 $B/\mathfrak k$ 可看作离散的 $(A/\mathfrak j)$-环。由上述讨论，群 $\Exan_{A/\mathfrak j}(B/\mathfrak k,L)$ 有定义。再设 $\mathfrak k'\subset\mathfrak k$、$\mathfrak j'\subset\mathfrak j$ 分别是 $B,A$ 的开双边理想，满足 $\rho(\mathfrak j')\subset\mathfrak k'$，则由（18.3.5.1）有典范同态
\begin{equation}
 \Exan_{A/\mathfrak j}(B/\mathfrak k,L)
 \longrightarrow\Exan_{A/\mathfrak j'}(B/\mathfrak k',L).
 \tag{18.5.1.1}
\end{equation}
\cmt{E5-T08}{原版在本句“分别属于”后写成 $A,B$，但 $\mathfrak k'$ 属于 $B$，$\mathfrak j'$ 属于 $A$；本处按前后定义明确换回 $B,A$。}
满足 $\rho(\mathfrak j)\subset\mathfrak k\subset\mathfrak k_0$ 的开理想偶对 $(\mathfrak j,\mathfrak k)$，对于关系“$\mathfrak j\supset\mathfrak j'$ 且 $\mathfrak k\supset\mathfrak k'$”显然构成向右滤过的有序集；（18.5.1.1）在这一指标集上定义了加法群的归纳系。略用记号，置
\begin{equation}
 \Exantop_A(B,L)=
 \varinjlim_{(\mathfrak j,\mathfrak k)}
 \Exan_{A/\mathfrak j}(B/\mathfrak k,L).
 \tag{18.5.1.2}
\end{equation}
这里称为略用记号，是因为不再有一个扩张集合与此群自然地一一对应。

说（18.5.1.2）右端为零，意味着：对满足 $\rho(\mathfrak j)\subset\mathfrak k\subset\mathfrak k_0$ 的每一对开理想 $\mathfrak j\subset A$、$\mathfrak k\subset B$，以及 $B/\mathfrak k$ 的每个以 $L$ 为核的 $(A/\mathfrak j)$-扩张 $E$，都存在开理想 $\mathfrak j'\subset\mathfrak j$、$\mathfrak k'\subset\mathfrak k$，满足 $\rho(\mathfrak j')\subset\mathfrak k'$，使 $E$ 沿 $B/\mathfrak k'\to B/\mathfrak k$ 的逆像平凡。

在 $A$ 交换、$B$ 为拓扑 $A$-代数（分别地，交换拓扑 $A$-代数）的情形，从 $\Exal_{A/\mathfrak j}(B/\mathfrak k,L)$、$\Exalcom_{A/\mathfrak j}(B/\mathfrak k,L)$ 类似定义归纳极限 $\Exaltop_A(B,L)$、$\Exalcotop_A(B,L)$；其细节留给读者。

\egap{18.5.2}
若给定连续环同态的交换图
\[
\begin{tikzcd}[column sep=large]
B'\arrow[r]&B\\
A'\arrow[r]\arrow[u]&A\arrow[u],
\end{tikzcd}
\]
则由（18.3.5.1）取归纳极限，典范地得到两个加法群同态
\[
 \Exantop_A(B,L)\longrightarrow\Exantop_{A'}(B,L)
 \longrightarrow\Exantop_{A'}(B',L).
\]
由于归纳极限函子在交换群范畴中正合，同态
\[
 \Exantop_A(B,L)\longrightarrow\Exantop_{A'}(B,L)
\]
的核是下列各同态之核的归纳极限：
\[
 \Exan_{A/\mathfrak j}(B/\mathfrak k,L)
 \longrightarrow\Exan_{A'/\mathfrak j'}(B/\mathfrak k,L),
\]
其中 $\mathfrak j'$ 取为 $\mathfrak j$ 在 $A'$ 中的逆像。将此核记为 $\Exantop_{A/A'}(B,L)$；同样定义 $\Exaltop_{A/A'}(B,L)$ 和 $\Exalcotop_{A/A'}(B,L)$。最后，若有连续的 $B$-双模同态 $L\to L'$，则由（18.3.6.1）取归纳极限，典范地得到加法群同态
\[
 \Exantop_A(B,L)\longrightarrow\Exantop_A(B,L').
\]

\egap{18.5.3}
给定拓扑环 $C$ 及两个拓扑 $C$-双模 $M,N$，将从 $M$ 到 $N$ 的连续 $C$-同态所成的加法群记为 $\HomCont_C(M,N)$。

\egas{18.5.3.1}{引理}
设 $C$ 是拓扑环，$E,L$ 是两个拓扑 $C$-双模。假设所考虑的拓扑均为线性拓扑，$L$ 离散，且被 $C$ 的某个开双边理想零化。则有典范同构
\begin{equation}
 \varinjlim_{(\mathfrak k,V)}\Hom_{C/\mathfrak k}(E/V,L)
 \xrightarrow{\sim}\HomCont_C(E,L),
 \tag{18.5.3.2}
\end{equation}
其中左端沿如下偶对 $(\mathfrak k,V)$ 所成的滤过有序集取归纳极限：$\mathfrak k$ 为 $C$ 的开双边理想，$V$ 为 $E$ 的开 $C$-子双模，并且
\[
 \mathfrak kL=L\mathfrak k=0,\qquad
 \mathfrak kE\subset V,\qquad E\mathfrak k\subset V.
\]
\begin{proof}
由于 $C/\mathfrak k$ 与 $E/V$ 都离散，存在典范同态
\[
 u_{\mathfrak k,V}:\Hom_{C/\mathfrak k}(E/V,L)
 \longrightarrow\HomCont_C(E,L).
\]
这些同态形成归纳系，因此取归纳极限得到（18.5.3.2）。当 $V'\subset V$ 时，$E/V'\to E/V$ 满射，故由定义，同态
\[
 \Hom_{C/\mathfrak k}(E/V,L)\longrightarrow
 \Hom_{C/\mathfrak k}(E/V',L)
\]
为单射；这里假设 $\mathfrak kL=L\mathfrak k=0$，并有 $\mathfrak kE\subset V'$、$E\mathfrak k\subset V'$。当 $\mathfrak k'\subset\mathfrak k$ 时，同态
\[
 \Hom_{C/\mathfrak k}(E/V,L)\longrightarrow
 \Hom_{C/\mathfrak k'}(E/V,L)
\]
也单射。因此（18.5.3.2）为单射。

另一方面，若 $u:E\to L$ 是连续 $C$-同态，则其核是 $E$ 的开子双模 $V_0$。若 $\mathfrak k_0$ 是 $C$ 的开双边理想，满足
\[
 \mathfrak k_0L=L\mathfrak k_0=0,\qquad
 \mathfrak k_0E\subset V_0,\qquad E\mathfrak k_0\subset V_0,
\]
则 $u$ 显然是从 $E/V_0$ 到 $L$ 的某个 $(C/\mathfrak k_0)$-同态的典范像，故（18.5.3.2）满射。
\end{proof}

于是（18.3.8）可以如下推广到拓扑环。

\egas{18.5.4}{命题}
设 $B$ 是线性拓扑环，$\mathfrak j$ 是其双边理想，$C=B/\mathfrak j$ 为商拓扑环。则 $\mathfrak j/\mathfrak j^2$ 典范地具有拓扑 $C$-双模结构；这里 $\mathfrak j$ 取 $B$ 诱导的拓扑，$\mathfrak j/\mathfrak j^2$ 取 $\mathfrak j$ 的商拓扑。对每个被 $C$ 的某个开理想零化的离散 $C$-双模 $L$，都有典范同构
\begin{equation}
 \HomCont_C(\mathfrak j/\mathfrak j^2,L)
 \xrightarrow{\sim}\Exantop_B(C,L).
 \tag{18.5.4.1}
\end{equation}
\begin{proof}
对 $B$ 的每个开理想 $\mathfrak k$，若 $(\mathfrak j+\mathfrak k)/\mathfrak j$ 零化 $L$，则由（18.3.8）有典范同构
\[
 \Hom_{B/(\mathfrak j+\mathfrak k)}
 \bigl((\mathfrak j+\mathfrak k)/(\mathfrak j^2+\mathfrak k),L\bigr)
 \simeq
 \Exan_{B/\mathfrak k}\bigl(B/(\mathfrak j+\mathfrak k),L\bigr).
\]
再取归纳极限，并利用（18.5.3.1）即可。
\cmt{E5-T09}{原版将这里调用的命题编号印为（18.3.9）；本节18.3止于（18.3.8），所需商环扩张同构也正是（18.3.8.1），故此处明确订正。关于上述单指标子系在（18.5.1.2）双指标系中的共尾性，见核查报告的补充验证；该补充不属于原文。}
\end{proof}

% END continuations/iv0_18_5.tex


% BEGIN continuations/iv0_19_01.tex
% Source: EGAIV.1.pdf, PDF pp.68--77; printed pp.69--78.
\setcounter{section}{18}
\section{形式光滑代数与Cohen环}
\source{EGA IV第一部分，本节原文为印刷第69--115页；本轮连续译入第69--89页，至（19.4.7）}
\setcounter{subsection}{-1}
\subsection{引言}
\egap{19.0.1}
在第IV章中，我们将引入并研究一类重要的预概形态射，即\emph{光滑态射}\orig{某些近期著作也称之为“简单态射”，这一用语取自经典的“简单点”术语；但是它容易引起混淆，尤其是在代数群理论中。}。它们的一项基本性质——这一性质与一个有限性条件合起来便刻画了它们——是态射的\emph{提升性质}：若 $f:X\to Y$ 是光滑态射，$g:Y'\to Y$ 是态射，则对于使 $Y''$ 成为与 $Y'$“相差甚微”的预概形的任意态射 $h:Y''\to Y'$，每个 $Y$-态射 $Y''\to X$ 都分解为 $Y''\xrightarrow{h}Y'\to X$。更精确地，先限于 $Y=\Spec(A)$、$X=\Spec(B)$ 均仿射的情形：若对每个 $A$-代数 $C$ 及 $C$ 的每个幂零理想 $\mathfrak j$，任意 $A$-同态 $B\to C/\mathfrak j$ 都可以提升为 $B\to C\to C/\mathfrak j$，则称 $B$ 为\emph{形式光滑的 $A$-代数}。换言之，映射
\[
 \Hom_A(B,C)\longrightarrow\Hom_A(B,C/\mathfrak j)
\]
是满射。在许多应用中，$X$ 将作为某个反变函子 $Y'\mapsto F(Y')$ 的代表对象出现；该函子从 $Y$-预概形范畴取值于集合范畴，因此有 $F(Y')\simeq\Hom_Y(Y',X)$（$0_{\mathrm{III}}$，8.1.8）。在仿射情形，令 $F^\circ(C)=F(\Spec(C))$，则验证上述条件下 $F^\circ(C)\to F^\circ(C/\mathfrak j)$ 为满射——即使尚不知道 $F$ 可表，也能进行这一验证——便可在附加的有限性条件满足时证明 $f$ 光滑。\cmt{F6-H01}{本段开头是引导性的描述，不应脱离后面的仿射条件，读作任意非仿射测试预概形都存在全局提升。正文在（19.3.1）给出拓扑代数的精确定义；几何情形须区分局部提升与全局提升。}

\egap{19.0.2}
对于拓扑环 $A$ 上的拓扑代数，也有相应的形式光滑代数概念；这里暂不作精确定义（参见19.3.1）。§19首先从初等观点研究这些概念；随后借助将在§§20、21中发展的微分性质，在§22证明更深入的定理。这里概述关于形式光滑代数的主要结果。

\noindent I．设 $A,B$ 为两个诺特局部环，$\varphi:A\to B$ 为局部同态，$k$ 为 $A$ 的剩余域，并令 $B_0=B\otimes_Ak$。给 $A,B,B_0$ 配备各自的预进制拓扑，给 $k$ 配备离散拓扑。$B$ 为形式光滑 $A$-代数的充要条件是：$B$ 为平坦 $A$-模，且 $B_0$ 为形式光滑 $k$-代数（19.7.1）。因此，对于诺特局部环，该定理把形式光滑性归结为域上的诺特局部代数的同一问题。

\noindent II．设 $k$ 为域，$A$ 为一个诺特局部 $k$-代数。$A$ 形式光滑的充要条件是它在 $k$ 上\emph{几何正则}，即对于 $k$ 的每个有限扩张 $k'$，半局部环 $A\otimes_kk'$ 都正则（22.5.8）；这时 $A$ 的完备化 $\widehat A$ 同构于形式幂级数环 $K[[T_1,\ldots,T_n]]$（19.6.5）。此外，当 $A$ 完备且在 $k$ 上形式光滑时，它的 $k$-代数结构完全由其剩余域 $K$ 及维数确定；只要满足不等式
\[
 \dim(A)\geq\operatorname{rg}_{K}(\Upsilon_{K/k}),
\]
该维数可以任意，其中 $\Upsilon_{K/k}$ 是 $K$ 的“不完美模”（22.2.6）。\cmt{F6-T01}{初版本处秩的下标印作 $k$；本轮已直接核对同卷（22.2.6）（印刷第188页），该定理明确以 $K$-向量空间的秩表述同一维数界。因此这里明示统一为 $K$。核对该引文不计作§22已经译入。}特别地，$k$ 的一个扩张 $K$ 形式光滑的充要条件是 $K/k$ 可分（19.6.1）。

\noindent III．设 $A$ 为诺特局部环，$\mathfrak j\ne A$ 为其理想，$A_0=A/\mathfrak j$；设 $B_0$ 为完备诺特局部环，$A_0\to B_0$ 为使 $B_0$ 成为形式光滑 $A_0$-代数的局部同态。那么存在完备诺特局部环 $B$、使 $B$ 成为平坦 $A$-模的局部同态 $A\to B$，以及 $A_0$-同构 $B\otimes_AA_0\simeq B_0$（于是由I，$B$ 是形式光滑的 $A$-代数）；此外，这些性质将 $B$ 确定到同构（19.7.2）。该定理特别包含关于完备诺特局部环结构的Cohen定理（19.8），它们将在第IV章§§6、7中起重要作用。

\noindent IV．研究一个诺特局部环在另一个诺特局部环上的形式光滑性，其意义来自光滑态射的如下“逐点”刻画：若 $X,Y$ 为局部诺特预概形，$f:X\to Y$ 为局部有限型态射，则 $f$ 光滑的充要条件是，对每个 $x\in X$，局部环 $\OO_x$ 在 $\OO_{f(x)}$ 上形式光滑。特别地，若 $Y=\Spec(k)$ 且 $k$ 为完美域，则 $f$ 光滑等价于 $X$ 正则。

\noindent V．最后，我们将在§§20、21、22中看到，形式光滑代数概念自然地出现在Kähler微分理论中，这两种理论彼此阐明。

\egap{19.0.3}
本节及以后各节中，拓扑环与拓扑模均假定为\emph{线性拓扑化}的（$0_{\mathrm I}$，7.1.1）；除非另有明确说明，所考虑的拓扑环均假定交换。回顾：若 $A,B$ 为拓扑环，环同态 $\rho:A\to B$ 定义了 $B$ 的 $A$-代数结构，则称 $B$ 为\emph{拓扑 $A$-代数}，若 $\rho$ 对所考虑的拓扑连续。

为简洁起见，在拓扑环 $A$（或拓扑 $A$-模 $M$）中，将“由理想（或子模）构成的零点基本邻域系”简称为“开理想（或开子模）基本系”。给定拓扑环 $A$ 和一个 $A$-模 $M$，当 $\mathfrak j$ 遍历一个开理想基本系时，子模 $\mathfrak jM$ 构成 $M$ 上某一拓扑的开子模基本系；该拓扑使 $M$ 成为拓扑 $A$-模，称为\emph{由 $A$ 的拓扑诱导的拓扑}。

设拓扑 $A$-模 $M$ 的拓扑比由 $A$ 诱导的拓扑\emph{更粗}。若 $N$ 为 $M$ 的开子模，则离散 $A$-模 $M/N$ 被 $A$ 的某个开理想零化，因为按假设存在开理想 $\mathfrak k$ 使 $\mathfrak kM\subset N$。
若 $M,N$ 的拓扑都由 $A$ 诱导，则每个 $A$-同态 $u:M\to N$ 连续：对 $N$ 中任意零点邻域 $V$，存在开理想 $\mathfrak j$ 使 $\mathfrak jN\subset V$，于是 $u(\mathfrak jM)\subset\mathfrak jN\subset V$。
当 $B$ 为拓扑 $A$-代数时，$B$ 上由 $A$ 诱导的拓扑比其给定拓扑\emph{更细}：对 $B$ 的每个开理想 $\mathfrak k$，按假设存在 $A$ 的开理想 $\mathfrak j$ 使 $\mathfrak jB\subset\mathfrak k$。

\subsection{形式满态射与形式单态射}
\source{印刷第71--78页；上传PDF第70--77页}
\egas{19.1.1}{命题}
设 $A$ 为拓扑环，$M,N$ 为拓扑 $A$-模，$(W_\lambda)$ 为 $N$ 的开子模基本系，$u:M\to N$ 为连续 $A$-同态。以 $\widehat M,\widehat N$ 表示其分离完备化，以 $\widehat u:\widehat M\to\widehat N$ 表示 $u$ 在分离完备化上的连续延拓。

（i）下列条件等价：
\begin{enumerate}[label=\alph*)]
\item $u(M)$ 在 $N$ 中稠密。
\item[$a'$）] $\widehat u(\widehat M)$ 在 $\widehat N$ 中稠密。
\item[$a''$）] 对每个 $\lambda$，复合同态 $M\xrightarrow{u}N\to N/W_\lambda$ 满射。
\end{enumerate}
（ii）下列条件等价：
\begin{enumerate}[label=\alph*),start=2]
\item $N$ 的拓扑经 $u$ 的逆像等于 $M$ 的拓扑。
\item[$b'$）] $\widehat u$ 将拓扑 $\widehat A$-模 $\widehat M$ 同构地映到 $\widehat N$ 的拓扑子 $\widehat A$-模 $\widehat u(\widehat M)$ 上（该子模必为闭子模）。
\item[$b''$）] $u^{-1}(W_\lambda)$ 构成 $M$ 的零点基本邻域系。
\end{enumerate}
这些断言直接来自分离完备化的定义及 $N/W_\lambda$ 的离散性。

\egap{19.1.2}
当（19.1.1）中（i）（或（ii））的等价条件满足时，称 $u$ 为\emph{形式满态射}（或\emph{形式单态射}）。若二者同时满足，则称 $u$ 为\emph{形式双态射}；由（19.1.1），这等价于 $\widehat u$ 是拓扑 $\widehat A$-模 $\widehat M\to\widehat N$ 的同构。\cmt{F6-H02}{这里的“形式”不可删去。例如环 $k[t]$ 配备 $(t)$-预进制拓扑后，到 $k[[t]]$ 的典范映射稠密且诱导原拓扑，但一般并不满射；它是形式双态射而非原模之间的双射。}

\egas{19.1.3}{命题}
设 $A$ 为拓扑环，$M,N$ 为拓扑 $A$-模，$u:M\to N$ 为连续 $A$-同态。假定 $M,N$ 分别有同指标集的开子模基本系 $(V_\lambda)$、$(W_\lambda)$，且 $u(V_\lambda)\subset W_\lambda$。令 $u_\lambda:M/V_\lambda\to N/W_\lambda$ 为商同态。则：

（i）$u$ 为形式满态射，当且仅当每个 $u_\lambda$ 满射。

（ii）$u$ 为形式单态射，当且仅当对每个 $\lambda$，存在 $\mu$ 使 $V_\mu\subset V_\lambda$，且 $\Ker u_\mu$ 包含于典范映射 $M/V_\mu\to M/V_\lambda$ 的核 $V_\lambda/V_\mu$。

（iii）$u$ 为形式双态射，当且仅当对每个 $\lambda$，$u_\lambda$ 满射，且存在 $\mu$ 使 $V_\mu\subset V_\lambda$，并使典范同态 $M/V_\mu\to M/V_\lambda$ 分解为
\[
 M/V_\mu\xrightarrow{u_\mu}N/W_\mu\longrightarrow M/V_\lambda
\]
（其中后一个同态必唯一）。

（i）是直接的。另一方面，$\Ker u_\mu=u^{-1}(W_\mu)/V_\mu$，而典范映射 $M/V_\mu\to M/V_\lambda$ 的核为 $V_\lambda/V_\mu$；所以（ii）的条件等价于 $u^{-1}(W_\mu)\subset V_\lambda$，从而立即得到（ii）。最后，若 $u_\mu$ 满射，则 $N/W_\mu$ 与 $(M/V_\mu)/\Ker u_\mu$ 相同构，因此 $\Ker u_\mu\subset V_\lambda/V_\mu$ 等价于 $M/V_\mu\to M/V_\lambda$ 分解为 $v\circ u_\mu$，其中 $v:N/W_\mu\to M/V_\lambda$ 为同态。

\egas{19.1.4}{推论}
设 $M,N$ 的拓扑均由拓扑环 $A$ 诱导，$u:M\to N$ 为形式满态射，$(\mathfrak j_\lambda)$ 为 $A$ 的开理想基本系。$u$ 是形式双态射的充要条件是，每个商同态 $u_\lambda:M/\mathfrak j_\lambda M\to N/\mathfrak j_\lambda N$ 双射。

事实上，$u(\mathfrak j_\lambda M)\subset\mathfrak j_\lambda N$，故可用（19.1.3，iii）。若有分解
\[
 M/\mathfrak j_\mu M\xrightarrow{u_\mu}N/\mathfrak j_\mu N\xrightarrow{v}M/\mathfrak j_\lambda M,
\]
则 $\mathfrak j_\lambda M/\mathfrak j_\mu M$ 经 $u_\mu$ 的像为 $\mathfrak j_\lambda N/\mathfrak j_\mu N$，所以后者经 $v$ 的像为零，且 $v$ 分解为
\[
 N/\mathfrak j_\mu N\to N/\mathfrak j_\lambda N\xrightarrow{w}M/\mathfrak j_\lambda M.
\]
于是 $M/\mathfrak j_\lambda M$ 的恒等自同构分解为
\[
 M/\mathfrak j_\lambda M\xrightarrow{u_\lambda}N/\mathfrak j_\lambda N\xrightarrow{w}M/\mathfrak j_\lambda M.
\]
这表明 $u_\lambda$ 单射；它已经是满射，推论得证。

\egas{19.1.5}{命题}
设 $A$ 为拓扑环，$M,N$ 为拓扑 $A$-模，$u:M\to N$ 为连续 $A$-同态。下列条件等价：

$a)$ 对每个离散拓扑 $A$-模 $E$ 和每个连续 $A$-同态 $v:M\to E$，存在连续 $A$-同态 $w:N\to E$ 使 $v=w\circ u$。

$b)$ 对 $M$ 的每个开子模 $V'$，存在 $N$ 的开子模 $W''$、$M$ 的开子模 $V''\subset V'\cap u^{-1}(W'')$，以及同态 $h:N/W''\to M/V'$，使典范同态 $M/V''\to M/V'$ 分解为
\[
 M/V''\xrightarrow{u''}N/W''\xrightarrow{h}M/V',
\]
其中 $u''$ 由 $u$ 取商得到。

要由 $a)$ 推出 $b)$，只须将 $a)$ 用于 $E=M/V'$ 和典范同态 $v:M\to M/V'$。这时 $W''=w^{-1}(0)$ 是 $N$ 的开子模，$V''=u^{-1}(W'')$ 是包含于 $V'$ 的开子模；这些子模和由 $w$ 取商所得的 $h:N/W''\to M/V'$ 满足要求。
反之，将 $E$ 换成 $v(M)$，可在证明 $b)\Rightarrow a)$ 时假定 $v$ 满射。于是 $V'=\Ker v$ 是开子模，$E$ 同构于离散 $A$-模 $M/V'$。应用 $b)$，取复合 $N\to N/W''\xrightarrow{h}M/V'$ 即得所求连续同态 $w$，因为下图交换：
\[
\begin{tikzcd}
 M\arrow[r,"u"]\arrow[d] & N\arrow[d]\\
 M/V''\arrow[r,"u''"] & N/W''.
\end{tikzcd}
\]
当这些条件满足时，称 $u$ \emph{形式可左逆}。由于 $b)$ 蕴含 $u^{-1}(W'')\subset V'$，$u$ 此时为形式单态射。以下推论也说明该术语的合理性。

\egas{19.1.6}{推论}
若存在连续 $\widehat A$-同态 $s:\widehat N\to\widehat M$，满足 $s\circ\widehat u=1_{\widehat M}$，则 $u$ 形式可左逆。

注意，此时 $\widehat u$ 将 $\widehat M$ 拓扑同构到 $\widehat N$ 的一个拓扑直和因子上。沿用（19.1.5，$a$）的记号，因 $E$ 离散，$v:M\to E$ 可延拓为连续 $\widehat A$-同态 $\widehat v:\widehat M\to E$。设 $j:M\to\widehat M$、$j':N\to\widehat N$ 为典范同态；令 $w=\widehat v\circ s\circ j'$，则
\[
 w\circ u=\widehat v\circ s\circ\widehat u\circ j=\widehat v\circ j=v.
\]

\egas{19.1.7}{推论}
设 $M,N$ 的拓扑均由 $A$ 诱导，$(\mathfrak j_\lambda)$ 为 $A$ 的开理想基本系。$u$ 形式可左逆，当且仅当每个商同态
$u_\lambda:M/\mathfrak j_\lambda M\to N/\mathfrak j_\lambda N$ 可左逆，即它将 $M/\mathfrak j_\lambda M$ 同构到 $N/\mathfrak j_\lambda N$ 的一个直和因子上。

充分性：在（19.1.5，$b$）中取 $V'=\mathfrak j_\lambda M$、$W''=\mathfrak j_\lambda N$、$V''=V'$，以及满足 $h\circ u_\lambda=1$ 的 $h$。反之，若 $u$ 形式可左逆，则对每个 $\lambda$，由（19.1.5，$b$）存在 $\mathfrak j_\mu\subset\mathfrak j_\lambda$ 及 $h:N/\mathfrak j_\mu N\to M/\mathfrak j_\lambda M$，使典范同态分解为
\[
 M/\mathfrak j_\mu M\xrightarrow{u_\mu}N/\mathfrak j_\mu N\xrightarrow{h}M/\mathfrak j_\lambda M.
\]
由于
\[
 h\bigl(\mathfrak j_\lambda(N/\mathfrak j_\mu N)\bigr)
 =\mathfrak j_\lambda h(N/\mathfrak j_\mu N)
 \subset\mathfrak j_\lambda(M/\mathfrak j_\lambda M)=0,
\]
$h$ 典范地分解为 $N/\mathfrak j_\mu N\to N/\mathfrak j_\lambda N\xrightarrow{s}M/\mathfrak j_\lambda M$，且 $s$ 是 $u_\lambda$ 的左逆。

\egas{19.1.8}{命题}
设拓扑环 $A$ 有可数递降开理想基本系 $(\mathfrak j_n)$；$M,N$ 为拓扑由 $A$ 诱导的拓扑 $A$-模，$u:M\to N$ 为 $A$-同态。令 $M_n=M/\mathfrak j_nM$、$N_n=N/\mathfrak j_nN$，$u_n:M_n\to N_n$ 为商同态。假定每个 $P_n=\Coker u_n$ 为投射 $(A/\mathfrak j_n)$-模，且 $M$ 分离完备。则 $u$ 形式可左逆的充要条件是 $u$ 可左逆；此时 $u$ 将 $M$ 拓扑同构到 $N$ 的一个拓扑直和因子上。

由（19.1.7），有交换图
\[
\begin{tikzcd}[column sep=1.8em]
 0\arrow[r]&M_n\arrow[r,"u_n"]&N_n\arrow[r,"p_n"]&P_n\arrow[r]&0\\
 0\arrow[r]&M_{n+1}\arrow[u,"f"]\arrow[r,"u_{n+1}"]&N_{n+1}\arrow[u,"g"]\arrow[r,"p_{n+1}"]&P_{n+1}\arrow[u,"h"]\arrow[r]&0,
\end{tikzcd}
\]
各行是分裂正合列，亦即对每个 $n$ 存在 $s'_n:P_n\to N_n$ 使 $p_n\circ s'_n=1_{P_n}$。我们将对 $n$ 递归构造 $s_n:P_n\to N_n$，满足 $p_n\circ s_n=1_{P_n}$，且
\[
\begin{tikzcd}
 P_n\arrow[r,"s_n"]&N_n\\
 P_{n+1}\arrow[u,"h"]\arrow[r,"s_{n+1}"]&N_{n+1}\arrow[u,"g"]
\end{tikzcd}
\]
交换。事实上，$g\circ s'_{n+1}-s_n\circ h$ 是从 $P_{n+1}$ 到
\[
 u_n(M_n)=u_{n+1}(M_{n+1})/\mathfrak j_nu_{n+1}(M_{n+1})
\]
的同态；由于 $P_{n+1}$ 为投射 $(A/\mathfrak j_{n+1})$-模，该同态分解为
\[
 P_{n+1}\xrightarrow{t_{n+1}}u_{n+1}(M_{n+1})
 \longrightarrow u_{n+1}(M_{n+1})/\mathfrak j_nu_{n+1}(M_{n+1}).
\]
因此 $s_{n+1}=s'_{n+1}-t_{n+1}$ 满足要求。由 $N_n=u_n(M_n)\oplus s_n(P_n)$，得到 $u_n$ 的左逆 $w_n:N_n\to M_n$，而且下图交换：
\[
\begin{tikzcd}
 N_n\arrow[r,"w_n"]&M_n\\
 N_{n+1}\arrow[u,"g"]\arrow[r,"w_{n+1}"]&M_{n+1}\arrow[u,"f"].
\end{tikzcd}
\]
射影系 $(w_n)$ 有射影极限 $\widehat w:\widehat N\to\widehat M=M$。再与 $N\to\widehat N$ 复合，得 $w:N\to M$，使 $w\circ u$ 在每个 $M_n$ 上诱导恒等映射。$M$ 分离完备，故 $w\circ u=1_M$。

\egas{19.1.9}{命题}
设 $A$ 为预可容许拓扑环（$0_{\mathrm I}$，7.1.2），$\mathfrak a$ 为定义理想，$(\mathfrak j_\lambda)$ 为开理想基本系。设 $M,N$ 的拓扑由 $A$ 诱导，且每个 $N/\mathfrak j_\lambda N$ 为投射 $(A/\mathfrak j_\lambda)$-模（参见19.2.3）。若 $u:M\to N$ 为 $A$-同态，则下列条件等价：

$a)$ $u$ 形式可左逆。

$b)$ 商同态 $u_0:M/\mathfrak aM\to N/\mathfrak aN$ 可左逆。

由（19.1.7），$a)$ 等价于每个 $u_\lambda$ 可左逆。因 $\mathfrak a$ 开，故包含某个 $\mathfrak j_\lambda$，从而 $u_0$ 可左逆。反过来，为证 $b)\Rightarrow a)$，注意对每个 $\lambda$，$\mathfrak a/\mathfrak j_\lambda$ 按假设为 $A/\mathfrak j_\lambda$ 的幂零理想。\cmt{F6-T02}{此处须将开理想基本系限于包含在 $\mathfrak a$ 中的共尾子系；否则 $\mathfrak a/\mathfrak j_\lambda$ 未必有定义。也可统一写为 $(\mathfrak a+\mathfrak j_\lambda)/\mathfrak j_\lambda$。这是证明中可直接实行的基本系收缩，不是增加命题的实质假设。}所需断言由下命题推出。

\egas{19.1.10}{命题}
设 $A$ 为环，$M,N$ 为 $A$-模，$N$ 投射，$u:M\to N$ 为 $A$-同态。设理想 $\mathfrak j$ 满足下列条件之一：

（i）$\mathfrak j$ 幂零；

（ii）$\mathfrak j$ 包含于 $A$ 的根中，且 $M$ 有限型。

则 $u$ 可左逆，当且仅当商同态 $u_0:M/\mathfrak jM\to N/\mathfrak jN$ 可左逆。

必要性是直接的。为证充分性，设 $v_0$ 为 $u_0$ 的左逆。因 $N$ 投射，复合 $N\to N/\mathfrak jN\xrightarrow{v_0}M/\mathfrak jM$ 分解为 $N\xrightarrow{v}M\to M/\mathfrak jM$。于是 $w=v\circ u$ 在 $M/\mathfrak jM$ 上诱导恒等自同构 $w_0$。只须证明 $w$ 双射，因为此时 $w^{-1}\circ v$ 即为 $u$ 的左逆。分两种情况。

（i）对每个 $n$ 有交换图
\[
\begin{tikzcd}[column sep=large]
 (\mathfrak j^n/\mathfrak j^{n+1})\otimes_{A/\mathfrak j}(M/\mathfrak jM)
 \arrow[r]\arrow[d,"1\otimes\gr^0(w)"']&\mathfrak j^nM/\mathfrak j^{n+1}M\arrow[d,"\gr^n(w)"]\\
 (\mathfrak j^n/\mathfrak j^{n+1})\otimes_{A/\mathfrak j}(M/\mathfrak jM)
 \arrow[r]&\mathfrak j^nM/\mathfrak j^{n+1}M.
\end{tikzcd}
\]
横箭头满射，而 $\gr^0(w)=w_0$ 为恒等映射，故 $\gr^n(w)$ 也为恒等映射，特别是双射。因 $\mathfrak j$ 幂零，$M$ 的 $\mathfrak j$-预进制滤过有限；所以 $w$ 双射（Bourbaki，\textit{交换代数}，第三章，§2，第8目，定理1的推论3）。

（ii）只须对每个极大理想 $\mathfrak m$ 证明 $w_{\mathfrak m}$ 在 $M_{\mathfrak m}$ 上双射（同书，第二章，§3，第3目，定理1）。由于 $\mathfrak jA_{\mathfrak m}\subset\mathfrak mA_{\mathfrak m}$，且 $A_{\mathfrak m}/\mathfrak jA_{\mathfrak m}=(A/\mathfrak j)_{\mathfrak m}$，可归结为局部环。此外可假定 $\mathfrak j$ 为极大理想 $\mathfrak j_0$：若 $u_0$ 可左逆，则它与 $1_{A/\mathfrak j_0}$ 张量所得的 $u_{00}:M/\mathfrak j_0M\to N/\mathfrak j_0N$ 也可左逆，因为 $M/\mathfrak j_0M=(M/\mathfrak jM)\otimes_{A/\mathfrak j}(A/\mathfrak j_0)$。

故设 $\mathfrak j$ 极大，$A/\mathfrak j$ 为域。只须证明 $M$ 在给定条件下是自由 $A$-模：这时 $\det w_0$ 是 $\det w$ 在 $A/\mathfrak j$ 中的像，故 $\det w\notin\mathfrak j$，从而可逆。因 $M/\mathfrak jM$ 为有限维向量空间，可取有限型自由 $A$-模 $L$ 及 $f:L\to M$，使商同态 $f_0:L/\mathfrak jL\to M/\mathfrak jM$ 双射。\cmt{F6-T03}{初版此处仅印“有限型 $A$-模 $L$”，紧接着却使用“这里 $L$ 是自由的”。应在选择 $L$ 时明确取有限型自由模：提升 $M/\mathfrak jM$ 的一组基即可。此处在正文补出“自由”，不改变原证明的方法。}由Nakayama引理，$f$ 满射（同书，第二章，§3，第2目，命题4的推论1）。令 $g=u\circ f$，则 $g_0:L/\mathfrak jL\to N/\mathfrak jN$ 可左逆。由于 $L$ 自由，开头的论证证明 $g$ 本身可左逆；这迫使 $f$ 单射，证明结束。

顺便记录（19.1.10）的如下推论。

\egas{19.1.11}{推论}
设 $A$ 为局部环，$k$ 为其剩余域，$M$ 为有限型 $A$-模，$N$ 为投射 $A$-模，$u:M\to N$ 为同态。$u$ 可左逆，当且仅当 $M$ 有一个生成系 $(x_i)_{1\leq i\leq m}$，使 $x_i\otimes1$ 经 $u\otimes1:M\otimes_Ak\to N\otimes_Ak$ 的像线性无关；此时 $x_i$ 是 $M$ 的一组基。

若 $u$ 可左逆，则 $M$ 为有限型投射模，因而自由，必要性成立。反之，条件满足时，$x_i\otimes1$ 是 $M\otimes_Ak$ 的一组基，且 $u\otimes1$ 可左逆。对极大理想 $\mathfrak j$ 应用（19.1.10）即可。

\egas{19.1.12}{推论}
设 $A$ 为环，$M$ 为有限型 $A$-模，$N$ 为投射 $A$-模，$u:M\to N$ 为同态。对每个 $\mathfrak p\in\Spec(A)$，下列条件等价：

$a)$ $u_{\mathfrak p}:M_{\mathfrak p}\to N_{\mathfrak p}$ 可左逆。

$b)$ $u\otimes1:M\otimes_Ak(\mathfrak p)\to N\otimes_Ak(\mathfrak p)$ 单射，其中 $k(\mathfrak p)$ 为 $A$ 在 $\mathfrak p$ 处的剩余域。

$c)$ 存在有限个 $x_i\in M$（$1\leq i\leq m$），其在 $M_{\mathfrak p}$ 中的像生成该模，并存在 $N$ 上的 $m$ 个线性形式 $v_i$，使
$\det(\langle v_i,u(x_j)\rangle)\notin\mathfrak p$。

$d)$ 存在 $f\in A\setminus\mathfrak p$，使 $u_f:M_f\to N_f$ 可左逆。

此外，满足这些条件的素理想构成 $\Spec(A)$ 的开集。

最后一项断言直接由 $d)$ 得出。因为 $N$ 投射，它是某个自由模 $A^{(I)}$ 的直和因子；又因 $M$ 有限型，$u(M)$ 包含于一个有限秩自由坐标子模 $A^n$。将 $N$ 换成 $N\oplus P$（$u$ 的像仍在 $N$ 内），$a)$、$b)$、$c)$、$d)$ 各自均保持等价，故可限于 $N$ 有限型自由的情形。
$d)\Rightarrow a)$ 是直接的；$a)$ 与 $b)$ 由（19.1.11）等价。$a)$ 还蕴含 $M_{\mathfrak p}$ 自由，因此可取 $x_i$ 的像为其一组基；注意 $N$ 有限型自由，$N_{\mathfrak p}$ 上的任意线性形式可写成 $v/f$，其中 $f\notin\mathfrak p$、$v$ 为 $N$ 上的线性形式，便得 $c)$。$c)\Rightarrow b)$ 也是直接的。

只剩 $a)\Rightarrow d)$。因 $N$ 有限表示，有典范同构
\[
 (\Hom_A(N,M))_{\mathfrak p}\simeq
 \Hom_{A_{\mathfrak p}}(N_{\mathfrak p},M_{\mathfrak p})
\]
（Bourbaki，\textit{交换代数}，第二章，§2，第7目，命题19）。若 $w_{\mathfrak p}$ 为 $u_{\mathfrak p}$ 的左逆，则存在 $w:N\to M$ 和 $f\notin\mathfrak p$，使 $w_{\mathfrak p}$ 是 $w$ 在局部化后除以 $f$ 所得的同态。\cmt{F6-T04}{初版此处将 $w$ 与 $w_{\mathfrak p}$ 的等式方向排混；本句按随后的关系 $(w\circ u)_{\mathfrak p}=f\,1_{M_{\mathfrak p}}$ 明示记号，以免把局部同态与全局同态写成无类型的等式。}于是 $(w\circ u)_{\mathfrak p}=f\,1_{M_{\mathfrak p}}$。由 $M$ 有限型，存在 $g\notin\mathfrak p$，使 $g(w\circ u-f1_M)$ 在所有生成元上为零，故为零。令 $h=fg$，则在 $A_h$ 上 $g$、$f$ 均可逆，并有 $w_h\circ u_h=(f/1)1_{M_h}$。所以 $(f/1)^{-1}w_h$ 是 $u_h$ 的左逆。

\egas{19.1.13}{命题}
在（19.1.9）的假设外，再假定开理想基本系为递降序列 $(\mathfrak j_n)$，且 $M$ 分离完备。那么（19.1.9）的条件 $a)$、$b)$ 还等价于：

$c)$ $u$ 可左逆。

已知 $c)\Rightarrow a)$（19.1.6）。反过来，若 $a)$ 成立，沿用（19.1.8）的记号，$M_n$ 是投射 $(A/\mathfrak j_n)$-模 $N_n$ 的直和因子，故 $P_n$ 也同构于 $N_n$ 的直和因子，因而投射；应用（19.1.8）即可。

最后记录以下命题。

\egas{19.1.14}{命题}
设 $A$ 为环，$M$ 为有限型 $A$-模，$N$ 为投射 $A$-模，$u:M\to N$ 为同态。

（i）$u$ 可左逆，当且仅当对 $A$ 的每个极大理想 $\mathfrak m$，$u_{\mathfrak m}:M_{\mathfrak m}\to N_{\mathfrak m}$ 可左逆。

（ii）设 $A'$ 为忠实平坦的 $A$-代数。$u$ 可左逆，当且仅当 $u\otimes1:M\otimes_AA'\to N\otimes_AA'$ 可左逆。

如（19.1.12），可设 $N$ 有限型自由。于是说 $u$ 可左逆，等价于 $u$ 单射且商模 $P=N/u(M)$ 投射，因为这时 $u(M)$ 为 $N$ 的直和因子。注意 $M$ 有限型，所以 $P$ 有限表示。

（i）必要性是直接的。反之，条件成立时，$u$ 单射（Bourbaki，\textit{交换代数}，第二章，§3，第3目，定理1）；而每个 $P_{\mathfrak m}=N_{\mathfrak m}/u_{\mathfrak m}(M_{\mathfrak m})$ 投射，故 $P$ 投射（同书，§5，第2目，定理1）。

（ii）必要性仍是直接的。反之，$u$ 单射（$0_{\mathrm I}$，6.4.1），且 $P\otimes_AA'=\Coker(u\otimes1)$ 投射，因而平坦。故 $P$ 为平坦 $A$-模（$0_{\mathrm I}$，6.6.3）；它又有限表示，所以投射（Bourbaki，\textit{交换代数}，第二章，§5，第2目，定理1的推论2）。

\egas{19.1.15}{注记}
形式可左逆同态的“对偶”概念是\emph{形式可右逆}同态。按定义，连续 $A$-同态 $u:M\to N$ 满足如下条件：对 $N$ 的每个开子模 $W'$，存在 $M$ 的开子模 $V'\subset u^{-1}(W')$、$N$ 的开子模 $W''\subset W'$，以及 $h:N/W''\to M/V'$，使典范同态 $N/W''\to N/W'$ 分解为
\[
 N/W''\xrightarrow{h}M/V'\xrightarrow{u'}N/W',
\]
其中 $u'$ 由 $u$ 取商得到。这蕴含 $u$ 为形式满态射。若存在连续 $\widehat A$-同态 $r:\widehat N\to\widehat M$，使 $\widehat u\circ r=1_{\widehat N}$，则可直接验证 $u$ 形式可右逆。
若（19.1.7）的条件成立，$u$ 形式可右逆当且仅当每个 $u_\lambda$ 可右逆，即 $\Ker u_\lambda$ 为 $M/\mathfrak j_\lambda M$ 的直和因子，且 $u_\lambda$ 将该核的某个补模同构到 $N/\mathfrak j_\lambda N$。
与（19.1.8）、（19.1.9）对应的命题留给读者表述并证明：在（19.1.8）的对应命题中，须假定 $M$ 分离完备且 $M_n$ 为投射 $(A/\mathfrak j_n)$-模；在（19.1.9）的对应命题中，应删去关于 $N/\mathfrak j_\lambda N$ 的假设，改为假定每个 $M/\mathfrak j_\lambda M$ 为投射 $(A/\mathfrak j_\lambda)$-模。

% END continuations/iv0_19_01.tex


% BEGIN continuations/iv0_19_23.tex
% Source: EGAIV.1.pdf, PDF pp.77--85; printed pp.78--86.
\Needspace{7\baselineskip}
\subsection{形式投射模}
\source{印刷第78--79页；上传PDF第77--78页}
\egas{19.2.1}{定义}
设 $A$ 为拓扑环。称拓扑 $A$-模 $M$ 为\emph{形式投射模}，若它满足下列条件：对 $A$ 的每个开理想 $\mathfrak j$、每对离散 $(A/\mathfrak j)$-模 $P,Q$、每个满 $A$-同态 $u:P\to Q$ 及每个连续 $A$-同态 $v:M\to Q$，存在连续 $A$-同态 $w:M\to P$，使 $v=u\circ w$。

\egap{19.2.2}
在验证（19.2.1）时，将 $Q$ 换成 $v(M)$、将 $P$ 换成 $u^{-1}(v(M))$，即可限于 $v$ 本身满射的情形。于是 $Q\simeq M/V$，其中 $V$ 为满足 $\mathfrak jM\subset V$ 的开子模。条件因而等价于：对每个离散 $(A/\mathfrak j)$-模 $P$ 及每个满 $A$-同态 $u:P\to M/V$，存在 $M$ 的开子模 $V'\subset V$ 及 $w:M/V'\to P$，使典范同态分解为
\[
 M/V'\xrightarrow{w}P\xrightarrow{u}M/V.
\]
注意，只须让 $\mathfrak j$ 遍历 $A$ 的一个由理想组成的零点基本邻域系来验证此条件。

\egap{19.2.3}
假设 $A$ 有开理想基本系 $(\mathfrak j_\lambda)$，$M$ 有同指标集的开子模基本系 $(V_\lambda)$，并且每个 $M/V_\lambda$ 为投射 $(A/\mathfrak j_\lambda)$-模。那么 $M$ 形式投射。事实上，沿用（19.2.2）的记号，取 $\lambda$ 使 $\mathfrak j_\lambda\subset\mathfrak j$ 且 $V_\lambda\subset V$。$P$ 也是 $(A/\mathfrak j_\lambda)$-模，故典范 $(A/\mathfrak j_\lambda)$-同态 $M/V_\lambda\to M/V$ 因 $M/V_\lambda$ 投射而分解为 $M/V_\lambda\to P\to M/V$。
当本段的这一更强条件满足时，称 $M$ \emph{严格形式投射}。

\egas{19.2.4}{命题}
设 $A$ 为拓扑环，拓扑 $A$-模 $M$ 的拓扑由 $A$ 诱导。那么 $M$ 形式投射，当且仅当它严格形式投射。

充分性已证。反之，设 $M$ 形式投射，$(\mathfrak j_\lambda)$ 为开理想基本系。对每个 $\lambda$，取自由 $(A/\mathfrak j_\lambda)$-模 $P$ 及满同态 $u:P\to M/\mathfrak j_\lambda M$。存在 $\mathfrak j_\mu\subset\mathfrak j_\lambda$ 使典范同态分解为
\[
 M/\mathfrak j_\mu M\xrightarrow{w}P\xrightarrow{u}M/\mathfrak j_\lambda M.
\]
因 $w(\mathfrak j_\lambda(M/\mathfrak j_\mu M))\subset\mathfrak j_\lambda P=0$，$w$ 分解为
$M/\mathfrak j_\mu M\to M/\mathfrak j_\lambda M\xrightarrow{v}P$。
于是 $u\circ v$ 为恒等映射，证明了 $M/\mathfrak j_\lambda M$ 为投射 $(A/\mathfrak j_\lambda)$-模。

\egas{19.2.5}{命题}
设 $A$ 为拓扑环，$M$ 为拓扑 $A$-模。

（i）$M$ 形式投射（或严格形式投射），当且仅当拓扑 $\widehat A$-模 $\widehat M$ 具有同一性质。

（ii）设 $A'$ 为拓扑 $A$-代数。若 $M$ 形式投射（或严格形式投射），则配备张量积拓扑的 $M\otimes_AA'$ 是形式投射（或严格形式投射）的拓扑 $A'$-模。

（i）只须注意：当 $\mathfrak j$（或 $V$）遍历 $A$ 的开理想（或 $M$ 的开子模）时，它的分离完备化 $\widehat{\mathfrak j}$（或 $\widehat V$）遍历 $\widehat A$ 的开理想（或 $\widehat M$ 的开子模），且有典范同构
\[
 \widehat A/\widehat{\mathfrak j}\simeq A/\mathfrak j,
 \qquad \widehat M/\widehat V\simeq M/V.
\]
形式投射性和严格形式投射性只涉及这些商，因此得到（i）。

（ii）先设 $M$ 形式投射，令 $M'=M\otimes_AA'$。设 $\mathfrak j'$ 为 $A'$ 的开理想，$P',Q'$ 为离散 $(A'/\mathfrak j')$-模，$u':P'\to Q'$ 为满 $A'$-同态，$v':M'\to Q'$ 为连续 $A'$-同态。存在 $A$ 的开理想 $\mathfrak j$ 使 $\mathfrak jA'\subset\mathfrak j'$，因此 $P',Q'$ 也是离散 $(A/\mathfrak j)$-模。考虑连续复合同态 $v:M\xrightarrow{j}M'\xrightarrow{v'}Q'$；由假设，有连续 $A$-同态 $w:M\to P'$ 使 $v=u'\circ w$。因 $P'$ 是拓扑 $A'$-模，$w$ 分解为 $M\xrightarrow{j}M'\xrightarrow{w'}P'$，其中 $w'$ 为连续 $A'$-同态。由 $v'\circ j=(u'\circ w')\circ j$ 得 $v'=u'\circ w'$。

再设 $M$ 严格形式投射，取 $A$ 的开理想基本系 $(\mathfrak j_\lambda)$ 和 $M$ 的开子模基本系 $(V_\lambda)$，使 $M/V_\lambda$ 为投射 $(A/\mathfrak j_\lambda)$-模；取 $A'$ 的开理想基本系 $(\mathfrak j'_{\mu})$。$M'$ 的零点基本邻域可取为
\[
 W_{\lambda\mu}=
 \Im(M\otimes_A\mathfrak j'_\mu)+\Im(V_\lambda\otimes_AA'),
\]
其中 $\lambda,\mu$ 满足 $\mathfrak j_\lambda A'\subset\mathfrak j'_\mu$。\cmt{F6-T05}{初版此段的 $A'$、$\mathfrak j'_\mu$ 漏印了若干撇号，并直接写了不同环中理想的包含。本处明确写出 $\mathfrak j_\lambda A'\subset\mathfrak j'_\mu$，使标量扩张与商模的基环一致。}由于
\[
 M'/W_{\lambda\mu}\simeq(M/V_\lambda)\otimes_A(A'/\mathfrak j'_\mu),
\]
且 $M/V_\lambda$ 为投射 $(A/\mathfrak j_\lambda)$-模，$M'/W_{\lambda\mu}$ 为投射 $(A'/\mathfrak j'_\mu)$-模。

\subsection{形式光滑代数}
\source{印刷第79--86页；上传PDF第78--85页}
\egas{19.3.1}{定义}
设 $A$ 为拓扑环，$B$ 为拓扑 $A$-代数。称 $B$ 为\emph{形式光滑的 $A$-代数}，若对每个离散拓扑 $A$-代数 $C$ 及 $C$ 的每个幂零理想 $\mathfrak j$，每个连续 $A$-同态 $u:B\to C/\mathfrak j$ 都分解为
\[
 B\xrightarrow{v}C\xrightarrow{\varphi}C/\mathfrak j,
\]
其中 $v$ 连续，$\varphi$ 为典范同态。

这一定义等价于下列性质：

\noindent（P）对 $B$ 的每个开理想 $\mathfrak k$ 和每个 $A$-同态 $u':B/\mathfrak k\to C/\mathfrak j$，存在开理想 $\mathfrak k'\subset\mathfrak k$，使复合同态 $B\to B/\mathfrak k\xrightarrow{u'}C/\mathfrak j$ 分解为
\[
 B\to B/\mathfrak k'\xrightarrow{v'}C\xrightarrow{\varphi}C/\mathfrak j,
\]
其中 $v'$ 为 $A$-同态。

事实上，若 $u:B\to C/\mathfrak j$ 连续，其核 $\mathfrak k$ 为开理想，故 $u$ 分解为 $B\to B/\mathfrak k\xrightarrow{u'}C/\mathfrak j$。若（P）满足，对 $u'$ 应用它，取 $v$ 为复合 $B\to B/\mathfrak k'\xrightarrow{v'}C$ 即满足定义。反之，设 $B$ 形式光滑。给定开理想 $\mathfrak k$ 和 $u':B/\mathfrak k\to C/\mathfrak j$，对 $u:B\to B/\mathfrak k\xrightarrow{u'}C/\mathfrak j$ 应用定义。若 $u=\varphi\circ v$ 且 $v$ 连续，则 $\mathfrak k'=\Ker v\cap\mathfrak k$ 为包含于 $\mathfrak k$ 的开理想；$v$ 分解为 $B\to B/\mathfrak k'\xrightarrow{v'}C$，从而得到（P）。

\egas{19.3.2}{命题}
设 $A$ 为离散环，$V$ 为投射 $A$-模。配备离散拓扑的对称代数 $B=\mathrm S_A(V)$ 是形式光滑的 $A$-代数。

令 $C,\mathfrak j$ 如（19.3.1）。将 $A$-代数同态 $u:B\to C/\mathfrak j$ 限制到 $V=\mathrm S_A^1(V)$，得 $A$-模同态 $u_1:V\to C/\mathfrak j$。因 $V$ 投射，$u_1$ 分解为 $V\xrightarrow{v_1}C\to C/\mathfrak j$。$v_1$ 延拓成 $A$-代数同态 $v:\mathrm S_A(V)\to C$，其与 $\varphi$ 的复合在 $V$ 上等于 $u$，故 $\varphi\circ v=u$。

\egas{19.3.3}{推论}
若 $A$ 为离散环，则任意多项式代数 $B=A[T_\alpha]_{\alpha\in I}$ 配备离散拓扑后，都是形式光滑的 $A$-代数。

\egas{19.3.4}{命题}
设 $A$ 为拓扑环，$B=A[[T_\alpha]]_{\alpha\in I}$ 为形式幂级数环，元素写作 $\sum_{\lambda\in\mathbf N^{(I)}}c_\lambda T^\lambda$（即幺半群 $\mathbf N^{(I)}$ 在 $A$ 上的大代数，作为 $A$-模与乘积 $A^{\mathbf N^{(I)}}$ 识别）。给 $B$ 配备乘积拓扑，则它是形式光滑的 $A$-代数。

设 $(\mathfrak a_\mu)_{\mu\in L}$ 为 $A$ 的开理想基本系。对每个有限子集 $H\subset I$、$\mu\in L$ 和整数 $n$，令 $\mathfrak k_{H,\mu,n}$ 为下述零点邻域：它由满足如下条件的 $(c_\lambda)$ 构成——若 $\lambda=(\lambda_\alpha)$ 满足 $\lambda_\alpha=0$（$\alpha\notin H$）且 $\sum_{\alpha\in H}\lambda_\alpha\leq n$，则 $c_\lambda\in\mathfrak a_\mu$。直接验证可知，这些 $\mathfrak k_{H,\mu,n}$ 是 $B$ 的理想，且构成零点基本邻域系；因此乘积拓扑与 $A$-代数结构相容。
先记录以下引理。

\egas{19.3.4.1}{引理}
设 $E$ 为离散拓扑 $A$-代数。

（i）若 $f:B\to E$ 为连续 $A$-同态，则存在有限子集 $H\subset I$，使 $f(T_\alpha)=0$（$\alpha\notin H$），且对每个 $\alpha\in H$，$f(T_\alpha)$ 在 $E$ 中幂零。

（ii）反之，给定有限子集 $H\subset I$ 和 $E$ 中一族幂零元 $(\xi_\alpha)_{\alpha\in H}$，存在唯一连续 $A$-同态 $g:B\to E$，使 $g(T_\alpha)=\xi_\alpha$（$\alpha\in H$）、$g(T_\alpha)=0$（$\alpha\notin H$）。

（i）直接来自 $f^{-1}(0)$ 是乘积 $A$-模 $A^{\mathbf N^{(I)}}$ 的零点邻域，因此除有限多个 $\lambda$ 外均有 $f(T^\lambda)=0$。
为证（ii），注意多项式环 $B'=A[T_\alpha]_{\alpha\in I}$ 在 $B$ 中稠密。$g|_{B'}$ 的存在唯一性是直接的；其连续性由 $g(T_\alpha)$（$\alpha\in H$）幂零得到：若 $(g(T_\alpha))^n=0$，则 $g(T^\lambda)$ 只可能在满足 $\lambda_\alpha=0$（$\alpha\notin H$）且 $\lambda_\alpha\leq n$（$\alpha\in H$）的有限多个 $\lambda$ 处非零。\cmt{F6-H03}{本引理中的“离散 $A$-代数”仍须有连续结构同态 $A\to E$。延拓还使用离散 $E$ 分离完备；不能把任意不连续的基环作用代入。}

引理既证，仍令 $C,\mathfrak j$ 如（19.3.1）。存在有限集 $H$，使 $u(T_\alpha)=0$（$\alpha\notin H$），且 $\eta_\alpha=u(T_\alpha)$（$\alpha\in H$）在 $C/\mathfrak j$ 中幂零。由于 $\mathfrak j$ 幂零，可取幂零元 $\xi_\alpha\in C$ 映到 $\eta_\alpha$。令连续 $A$-同态 $v:B\to C$ 满足 $v(T_\alpha)=0$（$\alpha\notin H$）、$v(T_\alpha)=\xi_\alpha$（$\alpha\in H$），则 $u$ 分解为 $B\xrightarrow{v}C\to C/\mathfrak j$。

\egas{19.3.5}{命题}
设 $A$ 为拓扑环。

（i）$A$ 是形式光滑的 $A$-代数。

（ii）若 $B$ 是形式光滑的 $A$-代数，$C$ 是形式光滑的 $B$-代数，则 $C$ 是形式光滑的 $A$-代数。

（iii）若 $B$ 是形式光滑的 $A$-代数，$A'$ 是拓扑 $A$-代数，则拓扑 $A'$-代数 $B\otimes_AA'$（$0_{\mathrm I}$，7.7.5、7.7.6）形式光滑。

（iv）设 $B$ 为拓扑 $A$-代数，$S\subset A$、$T\subset B$ 为乘性子集，$S$ 在 $B$ 中的典范像包含于 $T$。若 $B$ 是形式光滑的 $A$-代数，则 $T^{-1}B$ 是形式光滑的 $S^{-1}A$-代数。

（v）设 $B_i$（$1\leq i\leq n$）为拓扑 $A$-代数。$\prod_{i=1}^nB_i$ 形式光滑，当且仅当每个 $B_i$ 形式光滑。

（i）设 $C$ 为离散 $A$-代数，$\varphi:C\to C/\mathfrak j$ 为任意商代数的典范映射。从 $A$ 到 $C/\mathfrak j$ 的唯一 $A$-同态是 $A\xrightarrow{\psi}C\xrightarrow{\varphi}C/\mathfrak j$，其中 $\psi$ 定义 $C$ 的 $A$-代数结构；$\psi$ 连续，故满足定义。

（ii）设 $\alpha:A\to B$、$\beta:B\to C$ 为连续结构同态。取离散 $A$-代数 $E$、其幂零理想 $\mathfrak a$ 和连续 $A$-同态 $u:C\to E/\mathfrak a$。于是 $u\circ\beta\circ\alpha$ 定义了 $E/\mathfrak a$ 的 $A$-代数结构。由 $B$ 的形式光滑性，$u\circ\beta:B\to E/\mathfrak a$ 分解为 $B\xrightarrow{v}E\to E/\mathfrak a$，其中 $v$ 连续。$v$、$u\circ\beta$ 分别给 $E,E/\mathfrak a$ 定义拓扑 $B$-代数结构，后者仍为前者的离散商代数。而 $u$ 是连续 $B$-同态，故分解为 $C\xrightarrow{w}E\to E/\mathfrak a$，其中 $w$ 是连续 $B$-同态。因 $v\circ\alpha$ 正是 $E$ 的 $A$-结构同态，$w$ 也是连续 $A$-同态。

（iii）设 $C$ 为离散拓扑 $A'$-代数，$\mathfrak j$ 为其幂零理想，$u:B\otimes_AA'\to C/\mathfrak j$ 为连续 $A'$-同态。把 $u$ 与典范 $\rho:B\to B\otimes_AA'$ 复合，得连续 $A$-同态，因为 $A\to B\xrightarrow{\rho}B\otimes_AA'$ 等于 $A\to A'\xrightarrow{\sigma}B\otimes_AA'$。由假设，$u\circ\rho$ 分解为 $B\xrightarrow{v}C\to C/\mathfrak j$，其中 $v$ 连续；这里 $C$ 的 $A$-结构由 $A\to A'\to C$ 给定。由 $A\to B\xrightarrow{v}C$ 与 $A\to A'\xrightarrow{w}C$ 相等，存在连续环同态 $f:B\otimes_AA'\to C$，满足 $v=f\circ\rho$、$w=f\circ\sigma$（$0_{\mathrm I}$，7.7.6）。两个经由 $B$ 的复合及两个经由 $A'$ 的复合分别相等，故 $u$ 确实分解为 $B\otimes_AA'\xrightarrow{f}C\to C/\mathfrak j$。

（iv）$S^{-1}A$、$T^{-1}B$ 上的拓扑取为以 $S^{-1}\mathfrak j$、$T^{-1}\mathfrak k$ 为零点基本邻域的拓扑，其中 $\mathfrak j,\mathfrak k$ 分别遍历 $A,B$ 的开理想基本系（$0_{\mathrm I}$，7.6.1）。若 $\alpha:A\to B$ 为结构同态，则由 $\alpha(S)\subset T$ 所得的典范 $\alpha':S^{-1}A\to T^{-1}B$ 连续（$0_{\mathrm I}$，7.6.6）。设 $C$ 为离散拓扑 $S^{-1}A$-代数，$\mathfrak j$ 为其幂零理想，$u:T^{-1}B\to C/\mathfrak j$ 为连续 $S^{-1}A$-同态。复合 $B\to T^{-1}B\xrightarrow{u}C/\mathfrak j$ 由假设分解为 $B\xrightarrow{v}C\to C/\mathfrak j$，$v$ 为连续 $A$-同态。对 $t\in T$，$u(t/1)$ 可逆；因为 $\mathfrak j$ 幂零，该剩余类的每个提升，特别是 $v(t)$，均在 $C$ 中可逆（$0_{\mathrm I}$，7.1.12）。于是 $v$ 分解为 $B\to T^{-1}B\xrightarrow{w}C$。由 $v$ 连续，$w$ 连续（$0_{\mathrm I}$，7.6.6）；而且复合 $A\to B\to T^{-1}B\xrightarrow{w}C$ 等于 $A\to S^{-1}A\to T^{-1}B\xrightarrow{w}C$，所以 $w$ 是 $S^{-1}A$-同态。最后，由经由 $B$ 的两个复合相同，同样的论证给出 $u$ 等于 $T^{-1}B\xrightarrow{w}C\to C/\mathfrak j$。

最后证明（v）。设 $B=\prod_i B_i$，并以 $e_i\in B$ 表示各直因子的单位幂等元。若每个 $B_i$ 形式光滑，给定离散 $A$-代数 $C$、幂零理想 $\mathfrak j$ 和连续 $u:B\to C/\mathfrak j$，则 $\bar e_i=u(e_i)$ 是一组正交幂等元，且和为1。交换环中的幂等元可沿幂零理想唯一提升，故它们提升为 $C$ 中和为1的正交幂等元 $\epsilon_i$。于是
\[
 C\simeq\prod_i C_i,\qquad C_i=\epsilon_i C,
 \qquad C/\mathfrak j\simeq\prod_i C_i/(\epsilon_i\mathfrak j).
\]
$u$ 对应于各连续 $A$-同态 $u_i:B_i\to C_i/(\epsilon_i\mathfrak j)$；由 $B_i$ 形式光滑，它们可连续提升到 $C_i$，取有限乘积便得到所求 $B\to C$。

反之，设 $B$ 形式光滑，给定 $u_i:B_i\to C/\mathfrak j$，与投影 $B\to B_i$ 复合后连续提升为 $v:B\to C$。$v(e_i)$ 模 $\mathfrak j$ 为1，而 $v(e_j)$（$j\ne i$）模 $\mathfrak j$ 为零；由幂等元提升的唯一性，它们在 $C$ 中就分别为1和零。因此 $v$ 通过 $B_i$ 分解，给出 $u_i$ 的连续提升，证明 $B_i$ 形式光滑。\cmt{F6-G01}{初版把“从有限乘积到 $C$ 的同态”误说成“分别从每个因子到同一个 $C$ 的同态”，漏掉了目标中的幂等元分解。第六版在报告中补证；本版将正确证明移入正文。幂等元提升可逐次归结为平方零核：对 $e^2-e\in I$、$I^2=0$，用 $e-(2e-1)^{-1}(e^2-e)$ 修正；唯一性及正交性再由幂零理想不含非零幂等元得到。}

\egas{19.3.6}{命题}
设 $A$ 为拓扑环，$B$ 为拓扑 $A$-代数，$\widehat A,\widehat B$ 为各自的分离完备化，故 $\widehat B$ 为拓扑 $\widehat A$-代数。以下条件等价：

$a)$ $B$ 是形式光滑的 $A$-代数；

$b)$ $\widehat B$ 是形式光滑的 $A$-代数；

$c)$ $\widehat B$ 是形式光滑的 $\widehat A$-代数。

当然，若 $\varphi:A\to B$ 定义原结构，则 $\widehat\varphi$ 定义 $\widehat B$ 的 $\widehat A$-代数结构。每个离散 $A$-代数 $C$ 分离完备，因而结构同态可典范延拓，使其成为 $\widehat A$-代数；从 $B$ 到 $C$ 的连续 $A$-同态也典范延拓为从 $\widehat B$ 到 $C$ 的 $\widehat A$-同态，当然也是 $A$-同态。换言之，每个这样的同态唯一分解为 $B\to\widehat B\to C$。由此及（19.3.1）立即得到命题。

\egas{19.3.7}{推论}
在（19.3.5，iv）的假设下，拓扑 $A\{S^{-1}\}$-代数 $B\{T^{-1}\}$ 形式光滑。
这由定义（$0_{\mathrm I}$，7.6.1、7.6.7）、（19.3.5，iv）及（19.3.6）得出。

\egas{19.3.8}{命题}
设 $A$ 为拓扑环，$B$ 为拓扑 $A$-代数，并假定 $B$ 的每个开理想 $\mathfrak k$ 的平方 $\mathfrak k^2$ 也开。给 $A,B$ 配备比原拓扑更细的拓扑，得到拓扑环 $A',B'$，且假定原典范同态 $\varphi:A\to B$ 仍是从 $A'$ 到 $B'$ 的连续同态。若 $B'$ 为形式光滑的 $A'$-代数，则 $B$ 为形式光滑的 $A$-代数。
只须应用下列引理。

\egas{19.3.8.1}{引理}
设 $C$ 为离散 $A$-代数，$\mathfrak j$ 为其幂零理想。假定 $B$ 的每个开理想的平方仍开。如果复合同态 $v:B\xrightarrow{u}C\to C/\mathfrak j$ 连续，则 $u$ 连续。

事实上，$\mathfrak k=u^{-1}(\mathfrak j)$ 为 $v$ 的核，按假设开。若 $\mathfrak j^n=0$，则 $\mathfrak k^n\subset\Ker u$。由假设对 $h$ 递推，$\mathfrak k^{2^h}$ 开，因而每个 $\mathfrak k^n$ 都开。所以 $\Ker u$ 开，即 $u$ 连续。

注意，对 $B$ 的假设等价于它的拓扑为各个 $\mathfrak k$-预进制拓扑的上确界，其中 $\mathfrak k$ 遍历所有开理想。若 $B$ 是预进制环（$0_{\mathrm I}$，7.1.9），这一条件总满足。

\egap{19.3.9}
现在证明，形式光滑性还蕴含比（19.3.1）的定义更一般的同态“提升”性质。

\egas{19.3.10}{命题}
设 $A$ 为拓扑环，$B$ 为形式光滑 $A$-代数，$C$ 为拓扑 $A$-代数，$\mathfrak j$ 为其理想，满足：

$1^\circ$ $C$ 可度量且完备；

$2^\circ$ $\mathfrak j$ 闭，且序列 $(\mathfrak j^n)$ 趋于零。

则每个连续 $A$-同态 $u:B\to C/\mathfrak j$ 都分解为 $B\xrightarrow{v}C\to C/\mathfrak j$，其中 $v$ 为连续 $A$-同态。

设 $(\mathfrak a_n)$ 为 $C$ 的递降开理想基本系。对每个 $n$，考虑离散代数 $C/\mathfrak a_n$ 及其理想 $(\mathfrak j+\mathfrak a_n)/\mathfrak a_n$；存在 $k$ 使 $\mathfrak j^k\subset\mathfrak a_n$，故该理想幂零。复合
\[
 u_n:B\xrightarrow{u}C/\mathfrak j\longrightarrow
 C/(\mathfrak j+\mathfrak a_n)
 =(C/\mathfrak a_n)/((\mathfrak j+\mathfrak a_n)/\mathfrak a_n)
\]
由假设分解为 $B\xrightarrow{v_n}C/\mathfrak a_n\xrightarrow{\varphi_n}C/(\mathfrak j+\mathfrak a_n)$，其中 $v_n$ 连续。我们将逐次选择 $v_n$，使 $v_n$ 对每个 $n$ 均分解为 $B\xrightarrow{v_{n+1}}C/\mathfrak a_{n+1}\to C/\mathfrak a_n$；亦即 $(v_n)$ 为射影系。这由下列引理得到。

\egas{19.3.10.1}{引理〔校订〕}
设 $B$ 为形式光滑 $A$-代数，$E,E'$ 为离散拓扑 $A$-代数，$\mathfrak k,\mathfrak k'$ 分别为它们的幂零理想；设 $f:E\to E'$ 为满足
\[
 f(\mathfrak k)=\mathfrak k'
\]
的满 $A$-同态，$f':E/\mathfrak k\to E'/\mathfrak k'$ 为商同态。设 $u:B\to E/\mathfrak k$ 连续，$u'=f'\circ u$；若连续 $A$-同态 $v':B\to E'$ 提升 $u'$，则存在连续 $A$-同态 $v:B\to E$，同时满足
\[
 f\circ v=v',\qquad (E\to E/\mathfrak k)\circ v=u.
\]
\cmt{F6-E01}{初版仅要求 $f(\mathfrak k)\subseteq\mathfrak k'$，但按这一条件，引理有反例：取 $A=k$、$B=k[x,x^{-1}]$、$E=k[t]$、$E'=k[t]/(t^2)$、$\mathfrak k=0$、$\mathfrak k'=(\bar t)$，令 $u(x)=1$、$v'(x)=1+\bar t$，则 $v'$ 不能提升，因为 $k[t]$ 的单位都是非零常数。正文采用使原纤维积证明成立的充分修订条件“相等”，并明确同时提升 $u$；不声称这是所有个别情形的必要条件。}

在交换图
\[
\begin{tikzcd}
 E\arrow[r,"f"]\arrow[d]&E'\arrow[d]\\
 E/\mathfrak k\arrow[r,"f'"]&E'/\mathfrak k'
\end{tikzcd}
\]
中，各同态均满。若 $\mathfrak a=\Ker f$，则 $E'=E/\mathfrak a$，而由修订后的条件，$\mathfrak k'$ 与 $(\mathfrak k+\mathfrak a)/\mathfrak a$ 识别。\cmt{F6-G02}{这一识别在初版只有包含条件时不能成立；本版已将所需相等条件列入引理。主命题（19.3.10）及（19.5.6.3）的实际应用均满足它。}现在使用以下初等引理。

\egas{19.3.10.2}{引理}
设 $F$ 为环（不必交换），$\mathfrak a,\mathfrak b$ 为双边理想。则纤维积
\[
 (F/\mathfrak a)\times_{F/(\mathfrak a+\mathfrak b)}(F/\mathfrak b)
\]
与 $F/(\mathfrak a\cap\mathfrak b)$ 典范识别。

这是（18.1.7）的特殊应用：在其图（18.1.7.1）中，将 $G$ 换成 $F/(\mathfrak a\cap\mathfrak b)$，$\mathfrak j'$ 换成 $\mathfrak b/(\mathfrak a\cap\mathfrak b)$，$F$ 换成 $F/\mathfrak b$，$\mathfrak k$ 换成 $(\mathfrak a+\mathfrak b)/\mathfrak b$，$B$ 换成 $F/(\mathfrak a+\mathfrak b)$，$\theta$ 换成典范双射 $F$-同态 $(\mathfrak a+\mathfrak b)/\mathfrak b\simeq\mathfrak a/(\mathfrak a\cap\mathfrak b)$。

将该引理应用于（19.3.10.1），图
\[
\begin{tikzcd}
 B\arrow[r,"v'"]\arrow[d,"u"']&E'=E/\mathfrak a\arrow[d]\\
 E/\mathfrak k\arrow[r,"f'"]&E'/\mathfrak k'
\end{tikzcd}
\]
交换，故存在 $w:B\to E/(\mathfrak k\cap\mathfrak a)$，使 $v'$ 和 $u$ 分别分解为经此商环再到 $E/\mathfrak a$ 和 $E/\mathfrak k$ 的复合。$\Ker w$ 包含 $\Ker v'\cap\Ker u$，故为开理想，$w$ 连续。$\mathfrak k\cap\mathfrak a$ 幂零，故 $w$ 分解为 $B\xrightarrow{v}E\to E/(\mathfrak k\cap\mathfrak a)$，其中 $v$ 为满足要求的连续 $A$-同态。

回到（19.3.10）的证明：在（19.3.10.1）中依次以
\[
 C/\mathfrak a_{n+1},\quad C/\mathfrak a_n,\quad
 (\mathfrak j+\mathfrak a_{n+1})/\mathfrak a_{n+1},\quad
 (\mathfrak j+\mathfrak a_n)/\mathfrak a_n
\]
替换 $E,E',\mathfrak k,\mathfrak k'$，并以 $u_{n+1},u_n,v_n$ 替换 $u,u',v'$，便得 $v_{n+1}$。由于 $C$ 分离完备，$C=\varprojlim_n C/\mathfrak a_n$，且 $v=\varprojlim_n v_n$ 为连续 $A$-同态。$\mathfrak j$ 闭，所以 $\mathfrak j=\bigcap_n(\mathfrak j+\mathfrak a_n)$；$C/\mathfrak j$ 可度量且完备，而 $(\mathfrak j+\mathfrak a_n)/\mathfrak j$ 构成其零点基本邻域系，因此
\[
 C/\mathfrak j=\varprojlim_n C/(\mathfrak j+\mathfrak a_n),
 \qquad \varprojlim_n\varphi_n=\varphi:C\to C/\mathfrak j.
\]
由 $\varprojlim_nu_n=u$，得 $u=\varphi\circ v$，证毕。\cmt{F6-T06}{初版末行印为 $u=v\circ\varphi$，与定义域、陪域不合；此处明确更正为 $u=\varphi\circ v$。}

\egas{19.3.11}{推论}
设 $A$ 为拓扑环，$B$ 为形式光滑 $A$-代数，$C$ 为完备诺特局部环，$\mathfrak j\ne C$ 为其理想；设连续 $\varphi:A\to C$ 使 $C$ 成为拓扑 $A$-代数。则每个连续 $u_0:B\to C/\mathfrak j$ 均分解为 $B\xrightarrow{u}C\to C/\mathfrak j$，其中 $u$ 为连续 $A$-同态。
因为（19.3.10）的所有条件均满足（$0_{\mathrm I}$，7.3.5）。

\egas{19.3.12}{撤销条目}
本编号保留供原著引文定位。作者已在Err$_{\mathrm{IV}}$，8中要求删除本条；本版不以未经核实的替代定理填补此编号。\cmt{F6-A01}{初版断言：设 $B$ 为形式光滑拓扑 $A$-代数，$C\simeq\varprojlim C_\lambda$，其中 $C_\lambda$ 为离散Artin商环；若 $\mathfrak j\subset C$ 闭且拓扑幂零，且 $B$ 的每个开理想的平方仍开，则每个连续 $B\to C/\mathfrak j$ 均可连续提升到 $C$。作者正式勘误（IV第四部分印刷第347页）要求删除，因为原证明使用的 $W_\lambda,U_\lambda,V_\lambda$ 不是线性流形。原证明把代数同态的提升集合当作线性紧模中的仿射子空间，此步骤不成立；这不等于本版已构造该断言的反例。原陈述和原证明全文移存本版核查报告的历史附录。另见书末\hyperlink{ErrIV:8}{作者勘误}。}

% END continuations/iv0_19_23.tex


% BEGIN continuations/iv0_19_4.tex
% Source: EGAIV.1.pdf, PDF pp.85--88; printed pp.86--89.
\subsection{形式光滑性的初步判据}
\source{印刷第86--89页；上传PDF第85--88页}
\egas{19.4.1}{命题}
设 $A$ 为拓扑环，$B$ 为拓扑 $A$-代数。假定分别有 $A,B$ 中的两个递降滤过族 $(\mathfrak j_\alpha)_{\alpha\in I}$、$(\mathfrak k_\alpha)_{\alpha\in I}$，满足：

$1^\circ$ $(\mathfrak j_\alpha)$ 在 $A$ 中趋于零，$(\mathfrak k_\alpha)$ 在 $B$ 中趋于零；

$2^\circ$ 对每个 $\alpha$，$\mathfrak j_\alpha B\subset\mathfrak k_\alpha$，因而 $B/\mathfrak k_\alpha$ 为拓扑 $(A/\mathfrak j_\alpha)$-代数；

$3^\circ$ 每个 $B/\mathfrak k_\alpha$ 均为形式光滑的 $(A/\mathfrak j_\alpha)$-代数。

则 $B$ 为形式光滑的 $A$-代数。

事实上，设 $C$ 为离散拓扑 $A$-代数，$\mathfrak a$ 为其幂零理想，$\mathfrak k$ 为 $B$ 的开理想。可同时取足够后的 $\alpha$，使
\[
 \mathfrak k_\alpha\subset\mathfrak k,
 \qquad\mathfrak j_\alpha C=0.
\]
第一项由 $\mathfrak k_\alpha\to0$ 得到；第二项由连续结构同态 $A\to C$ 的核开，以及 $\mathfrak j_\alpha\to0$ 得到。于是 $C$ 是 $(A/\mathfrak j_\alpha)$-代数，$B/\mathfrak k$ 是 $B/\mathfrak k_\alpha$ 对一个开理想的商；每个 $A$-同态 $u':B/\mathfrak k\to C/\mathfrak a$ 也为 $(A/\mathfrak j_\alpha)$-同态。\cmt{F6-G04}{初版只明写 $\mathfrak k_\alpha\subset\mathfrak k$，随后直接应用形式光滑性；还须使测试代数 $C$ 本身具有上述商基环结构。本版把同一组原假设所保证的同时选取步骤补入正文，不增加命题假设。}因此存在 $B$ 的开理想 $\mathfrak k'$，满足 $\mathfrak k_\alpha\subset\mathfrak k'\subset\mathfrak k$，并使
\[
 B/\mathfrak k_\alpha\to B/\mathfrak k\xrightarrow{u'}C/\mathfrak a
\]
分解为 $B/\mathfrak k_\alpha\to B/\mathfrak k'\xrightarrow{v'}C\to C/\mathfrak a$，其中 $v'$ 为 $(A/\mathfrak j_\alpha)$-同态。由（19.3.1）后的说明即得结论。

\egas{19.4.2}{推论}
设 $A$ 为拓扑环，$B$ 为拓扑 $A$-代数，$(\mathfrak j_\alpha)$ 为 $A$ 中趋于零的递降滤过理想族。令 $A_\alpha=A/\mathfrak j_\alpha$、$B_\alpha=B/\mathfrak j_\alpha B$。则 $B$ 是形式光滑 $A$-代数，当且仅当每个 $B_\alpha$ 是形式光滑 $A_\alpha$-代数。

充分性来自（19.4.1），必要性来自（19.3.5，iii）。

\egas{19.4.3}{命题}
设 $A$ 为拓扑环，$B$ 为拓扑 $A$-代数。假定对每个离散 $A$-代数 $C$ 及满足 $\mathfrak j^2=0$ 的理想 $\mathfrak j$，每个连续 $A$-同态 $u:B\to C/\mathfrak j$ 都分解为 $B\xrightarrow{v}C\to C/\mathfrak j$，其中 $v$ 连续。则 $B$ 为形式光滑 $A$-代数。

事实上，设 $\mathfrak k$ 为 $C$ 的任意幂零理想，$u':B\to C/\mathfrak k$ 为连续 $A$-同态。若 $\mathfrak k^m=0$，令 $C_i=C/\mathfrak k^i$（$1\leq i\leq m$）。则 $C_m=C$；对 $2\leq i\leq m$，$\mathfrak k^{i-1}/\mathfrak k^i$ 为 $C_i$ 的平方零理想 $\mathfrak j_i$，且 $C_i=C_{i+1}/\mathfrak j_{i+1}$。由假设对 $i$ 递归，得到连续 $u_i:B\to C_i$，使 $u_1=u'$，且 $u_i$ 分解为 $B\xrightarrow{u_{i+1}}C_{i+1}\to C_i$。从而得出结论。

\egas{19.4.4}{命题}
设 $A$ 为拓扑环，$B$ 为交换拓扑 $A$-代数。$B$ 形式光滑，当且仅当对每个被 $B$ 的某个开理想零化的离散拓扑 $B$-模 $L$，均有（参见18.5.1）
\[
 \Exalcotop_A(B,L)=0.
\]

设 $(\mathfrak k_\lambda)$ 为 $B$ 的递降开理想基本系，令 $B_\lambda=B/\mathfrak k_\lambda$。取离散拓扑 $A$-代数 $C$ 和平方零理想 $\mathfrak j$，于是 $C$ 是 $C/\mathfrak j$ 以 $\mathfrak j$ 为核的 $A$-扩张。给定 $u:B_\lambda\to C/\mathfrak j$，沿 $u$ 拉回该扩张，得到 $B_\lambda$ 以 $\mathfrak j$ 为核的 $A$-扩张 $E_\lambda$（18.2.5），并给它配备离散拓扑。若所述扩张类群为零，则按定义（18.5.1）存在 $\mu$，使 $\mathfrak k_\mu\subset\mathfrak k_\lambda$ 且拉回扩张 $E_\mu$ 为 $A$-平凡扩张。这意味着（18.1.6）有连续 $A$-同态 $v:B_\mu\to E_\lambda$，使典范 $B_\mu\to B_\lambda$ 分解为 $B_\mu\xrightarrow{v}E_\lambda\to B_\lambda$。于是 $B_\mu\to B_\lambda\to C/\mathfrak j$ 分解为 $B_\mu\to E_\lambda\to C\to C/\mathfrak j$；由（19.3.1）及（19.4.3），$B$ 形式光滑。
反向结论是直接的：将（19.3.1）用于以 $L$ 为核的 $B_\lambda$ 的拓扑 $A$-扩张 $C$，以及恒等同态 $B_\lambda\to B_\lambda=C/L$ 即可。

若 $A,B$ 均离散，判据化为
\begin{equation}
 \Exalcom_A(B,L)=0\qquad\text{对每个 $B$-模 $L$}
 \tag{19.4.4.1}
\end{equation}
亦即，$B$ 以一个 $B$-模为核的每个交换 $A$-扩张都须为 $A$-平凡扩张。

\egas{19.4.5}{推论}
设 $A$ 为拓扑环，$B$ 为交换拓扑 $A$-代数。

（i）若 $B$ 形式光滑，则对 $B$ 的每个开理想 $\mathfrak k$、每个 $(B/\mathfrak k)$-模 $L$，以及每个从 $(B/\mathfrak k)\times(B/\mathfrak k)$ 到 $L$ 的 $A$-双线性对称Hochschild二余循环 $f$（18.4.3），存在开理想 $\mathfrak k'\subset\mathfrak k$，使对于典范 $\varphi:B/\mathfrak k'\to B/\mathfrak k$，$f\circ(\varphi\times\varphi)$ 是从 $(B/\mathfrak k')\times(B/\mathfrak k')$ 到 $L$ 的 $A$-双线性Hochschild二余边界。\cmt{F6-T07}{初版最后一处定义域仍印为 $(B/\mathfrak k)\times(B/\mathfrak k)$。复合 $f\circ(\varphi\times\varphi)$ 的定义域应为带撇号的商环乘积，此处按映射类型明示更正。}

（ii）若 $B$ 作为 $A$-模形式投射（19.2.1），且（i）的条件满足，则 $B$ 为形式光滑 $A$-代数。

（i）二余循环 $f$ 定义一个 $B/\mathfrak k$ 以 $L$ 为核的Hochschild $A$-扩张 $C$（18.4.3）。把（19.3.1）用于 $C$、其平方零理想 $L$ 及恒等同态 $B/\mathfrak k\to B/\mathfrak k=C/L$，再用（18.4.3），即得（i）。

（ii）应用（19.4.3）：考虑 $B$ 的开理想 $\mathfrak k$、$A$ 的开理想 $\mathfrak j$ 满足 $\mathfrak jB\subset\mathfrak k$，以及 $B/\mathfrak k$ 以 $L$ 为核的一个 $(A/\mathfrak j)$-扩张 $E$。因 $B$ 作为 $A$-模形式投射，连续 $A$-线性典范映射 $p:B\to B/\mathfrak k$ 分解为 $B\xrightarrow{w}E\to B/\mathfrak k$，其中 $w$ 连续且 $A$-线性（19.2.1）。故 $w$ 又分解为 $B\to B/\mathfrak k'\xrightarrow{w'}E$，其中 $\mathfrak k'\subset\mathfrak k$ 为开理想。必要时缩小 $\mathfrak j$，可假定 $\mathfrak jB\subset\mathfrak k'$。将 $E$ 沿 $B/\mathfrak k'\to B/\mathfrak k$ 拉回，所得 $E'$ 是 $B/\mathfrak k'$ 以 $L$ 为核的Hochschild $(A/\mathfrak j)$-扩张。对定义它的余循环应用（i）（18.4.3），得到开理想 $\mathfrak k''\subset\mathfrak k'$，使 $E$ 沿 $B/\mathfrak k''\to B/\mathfrak k$ 的拉回 $E''$ 为 $A$-平凡扩张，证毕。

\egas{19.4.6}{推论}
设 $A$ 为拓扑环，$B$ 为作为 $A$-模形式投射的拓扑 $A$-代数。设 $A'$ 为配备由 $A$ 诱导拓扑的 $A$-代数，且 $A'$ 作为 $A$-模忠实平坦。再假定满足以下条件之一：

$1^\circ$ $A$ 有开理想基本系 $(\mathfrak j_\lambda)$，$B$ 有同指标集的开理想基本系 $(\mathfrak m_\lambda)$，使 $\mathfrak j_\lambda B\subset\mathfrak m_\lambda$，且每个 $B/\mathfrak m_\lambda$ 为有限型投射 $(A/\mathfrak j_\lambda)$-模；

$2^\circ$ $A'$ 为有限型投射 $A$-模。

那么，配备张量积拓扑的 $B'=B\otimes_AA'$ 是形式光滑的 $A'$-代数，当且仅当 $B$ 是形式光滑的 $A$-代数。

由（19.3.5，iii），充分性不需关于 $B,A'$ 的任何额外假设。为证必要性，应用（19.4.5）。在假设 $2^\circ$ 下，仍以 $(\mathfrak m_\lambda)$ 表示 $B$ 的开理想基本系，并对每个 $\lambda$ 取 $A$ 的开理想 $\mathfrak j_\lambda$ 使 $\mathfrak j_\lambda B\subset\mathfrak m_\lambda$。两种情形都令
\[
 A_\lambda=A/\mathfrak j_\lambda,\quad
 B_\lambda=B/\mathfrak m_\lambda,\quad
 A'_\lambda=A_\lambda\otimes_AA',\quad
 B'_\lambda=B_\lambda\otimes_AA'.
\]
设 $f$ 为从 $B_\lambda\times B_\lambda$ 到某个 $B_\lambda$-模 $L$ 的 $A_\lambda$-双线性对称Hochschild二余循环。标量扩张给出 $A'_\lambda$-双线性对称二余循环 $f'=f\otimes1:B'_\lambda\times B'_\lambda\to L'=L\otimes_AA'$。由 $B'$ 形式光滑及（19.4.5，i），存在 $\mu$ 使 $\mathfrak m_\mu\subset\mathfrak m_\lambda$，并使对于典范 $\varphi':B'_\mu\to B'_\lambda$，$f'\circ(\varphi'\times\varphi')$ 为二余边界。亦即，其在Hochschild群 $H^2_{A'_\mu}(B'_\mu,L')^s$ 中的类 $c'$ 为零。

设 $\varphi:B_\mu\to B_\lambda$ 为典范同态，$c$ 为 $f\circ(\varphi\times\varphi)$ 在 $H^2_{A_\mu}(B_\mu,L)^s$ 中的类，则 $c'$ 为 $c$ 的典范像。令 $P'_\bullet$ 为（18.4.5）中用于计算该对称Hochschild群的、对应 $A_\mu,B_\mu$ 的复形；对应 $A'_\mu,B'_\mu$ 的复形为 $P'_\bullet\otimes_AA'$。在假设 $1^\circ$ 下，由构造，其各项为有限型投射 $A_\mu$-模。由Bourbaki，\textit{代数}，第二章，第3版，§5，第3目，命题7，在两种假设下都有典范同构
\[
 \Hom_{A'_\mu}(P'_\bullet\otimes_AA',L\otimes_AA')
 \simeq\Hom_{A_\mu}(P'_\bullet,L)\otimes_AA'.
\]
又因 $A'$ 平坦，由（18.4.5）得到
\[
 H^2_{A'_\mu}(B'_\mu,L')^s
 \simeq H^2_{A_\mu}(B_\mu,L)^s\otimes_AA'.
\]
故可写 $c'=c\otimes1$；由 $A'$ 忠实平坦，$c'=0$ 蕴含 $c=0$，证明完成。

\egas{19.4.7}{命题}
设 $A$ 为预可容许拓扑环，$\mathfrak j$ 为其定义理想（$0_{\mathrm I}$，7.1.2），$B$ 为作为 $A$-模形式投射的拓扑 $A$-代数。给商环 $A_0=A/\mathfrak j$、$B_0=B/\mathfrak jB$ 配备商拓扑。则 $B$ 为形式光滑 $A$-代数，当且仅当 $B_0$ 为形式光滑 $A_0$-代数。

必要性由（19.4.2）给出。为证充分性，先取包含于 $\mathfrak j$ 的开理想构成 $A$ 的零点基本邻域系 $(\mathfrak j_\alpha)$。由（19.4.2）及 $B/\mathfrak j_\alpha B$ 作为 $(A/\mathfrak j_\alpha)$-模形式投射（19.2.5），可归结为 $A$ 离散；由预可容许性，$\mathfrak j$ 此时幂零。只须再证 $\mathfrak j^2=0$ 的情形：若 $\mathfrak j^m=0$，则可对 $k=2,3,\ldots,m$，依次将该结果应用于
$A_k=A/\mathfrak j^k$、$B_k=B/\mathfrak j^kB$ 和 $A_k$ 的平方零理想 $\mathfrak j^{k-1}/\mathfrak j^k$；注意 $B_k=B\otimes_AA_k$ 为形式投射的 $A_k$-模（19.2.5）。

现在应用（19.4.5，ii）。设 $\mathfrak k$ 为 $B$ 的开理想，$f:(B/\mathfrak k)\times(B/\mathfrak k)\to L$ 为取值于 $(B/\mathfrak k)$-模 $L$ 的 $A$-双线性对称Hochschild二余循环。先考虑 $\mathfrak jL=0$ 的特殊情形。此时 $L$ 也是 $(B_0/\mathfrak kB_0)$-模，而且 $f$ 分解为
\[
 (B/\mathfrak k)\times(B/\mathfrak k)
 \longrightarrow(B_0/\mathfrak kB_0)\times(B_0/\mathfrak kB_0)
 \xrightarrow{f_0}L,
\]
其中 $f_0$ 为双线性对称二余循环。由假设，有开理想 $\mathfrak k'\subset\mathfrak k$ 和 $A_0$-线性映射 $g_0:B_0/\mathfrak k'B_0\to L$，使
\[
 f_0(\varphi_0(x),\varphi_0(y))=xg_0(y)-g_0(xy)+g_0(x)y
\]
对 $x,y\in B_0/\mathfrak k'B_0$ 成立，其中 $\varphi_0:B_0/\mathfrak k'B_0\to B_0/\mathfrak kB_0$ 为典范同态。于是复合 $A$-线性映射 $g:B/\mathfrak k'\to B_0/\mathfrak k'B_0\xrightarrow{g_0}L$ 满足 $dg=f\circ(\varphi\times\varphi)$，$\varphi:B/\mathfrak k'\to B/\mathfrak k$ 为典范同态。

再考虑一般情形。先取 $L'=L/\mathfrak jL$，则 $\mathfrak jL'=0$。将 $f$ 与 $L\to L'$ 复合，得对称二余循环 $f'$。由前一情形，存在开理想 $\mathfrak k'\subset\mathfrak k$ 和 $A$-线性 $g':B/\mathfrak k'\to L'$，满足 $dg'=f'\circ(\varphi'\times\varphi')$，其中 $\varphi':B/\mathfrak k'\to B/\mathfrak k$ 为典范同态。因 $B$ 是形式投射 $A$-模，存在开理想 $\mathfrak k_1\subset\mathfrak k'$ 和 $A$-线性 $g_1:B/\mathfrak k_1\to L$，使 $B/\mathfrak k_1\to B/\mathfrak k'\xrightarrow{g'}L'$ 分解为 $B/\mathfrak k_1\xrightarrow{g_1}L\to L'$。

若 $\varphi_1:B/\mathfrak k_1\to B/\mathfrak k$ 为典范同态，考虑
\[
 f_1(x,y)=f(\varphi_1(x),\varphi_1(y))-xg_1(y)+g_1(xy)-g_1(x)y.
\]
它是 $(B/\mathfrak k_1)\times(B/\mathfrak k_1)\to L$ 的 $A$-双线性对称二余循环。由 $g_1$ 的选择，$f_1$ 取值于 $\mathfrak jL$。因 $\mathfrak j(\mathfrak jL)=0$，再用第一种情形，得到开理想 $\mathfrak k''\subset\mathfrak k_1$ 和 $A$-线性 $g_2:B/\mathfrak k''\to\mathfrak jL$，满足
\[
 f_1(\varphi_2(x),\varphi_2(y))=xg_2(y)-g_2(xy)+g_2(x)y,
\]
其中 $\varphi_2:B/\mathfrak k''\to B/\mathfrak k_1$ 为典范同态。因此 $A$-线性映射
\[
 g=g_2+g_1\circ\varphi_2:B/\mathfrak k''\longrightarrow L
\]
满足 $dg=f\circ(\varphi\times\varphi)$，$\varphi:B/\mathfrak k''\to B/\mathfrak k$ 为典范同态。证毕。

% END continuations/iv0_19_4.tex


% BEGIN continuations/iv0_19_5a.tex
% 底本 EGAIV.1.pdf，PDF第89--93页，分册印刷第90--94页。
\subsection{形式光滑性与相伴分次环}
\source{第四卷第一部分，印刷第90--100页；上传PDF第89--99页}
\egap{19.5.1}
设 $C$ 为交换拓扑环，$V$ 为拓扑 $C$-模。考虑对称代数
\[
 \Salg_C(V)=\bigoplus_{n\geq0}\Salg_C^n(V).
\]
我们将典范地在它上面赋予一个与其 $C$-代数结构相容的线性拓扑。为此，取 $V$ 的开子模 $U$，以及它在 $\Salg_C(V)$ 中生成的分次理想 $U\Salg_C(V)$；令
\[
 \mathfrak a\Salg_C(V)+U\Salg_C(V)
\]
构成 $\Salg_C(V)$ 的零点基本邻域系，其中 $\mathfrak a$ 与 $U$ 分别遍历 $C$ 的开理想基本系和 $V$ 的开子模基本系。

若 $V$ 的拓扑不比由 $C$ 的拓扑诱导的模拓扑更细，就可以只取满足 $\mathfrak aV\subset U$ 的偶对 $(\mathfrak a,U)$。此时
\[
 \mathfrak a\Salg_C(V)+U\Salg_C(V)
   =\mathfrak a\,1_C+U\Salg_C(V).
\]
因此，对每个 $n\geq1$，在 $\Salg_C^n(V)$ 上的诱导拓扑以 $U\Salg_C^{n-1}(V)$ 为零点基本邻域，其中 $U$ 遍历 $V$ 的开子模。特别地，它在 $V=\Salg_C^1(V)$ 上诱导的拓扑就是原给定拓扑；在一般情形中，诱导拓扑可能比原拓扑更粗。

在所有情形中，这样定义的拓扑代数 $\Salg_C(V)$，在拓扑 $C$-模与拓扑 $C$-代数的范畴之间，解决与不带拓扑时相同的泛问题。事实上，设 $E$ 为拓扑 $C$-代数，$u:V\to E$ 为 $C$-模同态，$\Salg(u)=v$ 为它在 $\Salg_C(V)$ 上的典范延拓。若 $u$ 连续，对 $E$ 的任意开理想 $\mathfrak b$，其逆像 $U=u^{-1}(\mathfrak b)$ 是 $V$ 的开子模，且
\[
 v\bigl(U\Salg_C(V)\bigr)\subset\mathfrak b.
\]
另一方面，存在 $C$ 的开理想 $\mathfrak a$，使 $\mathfrak aE\subset\mathfrak b$，从而
$v(\mathfrak a\Salg_C(V))\subset\mathfrak b$。这就证明了 $v$ 连续。
反之，若 $v$ 连续，对 $E$ 的任意开理想 $\mathfrak b$，存在 $V$ 的开子模 $U$，使
$v(U\Salg_C(V))\subset\mathfrak b$；特别有 $v(U\cdot1_C)\subset\mathfrak b$，即 $u(U)\subset\mathfrak b$，故 $u$ 连续。

还应记住，对 $n\geq1$，有典范的 $C$-模同构
\begin{equation}
 \Salg_C^n(V)\big/\bigl(U\Salg_C^{n-1}(V)\bigr)
       \simeq\Salg_C^n(V/U).
 \tag{19.5.1.1}
\end{equation}
在上述 $\mathfrak aV\subset U$ 的情形中，两端带相应商拓扑时都是离散拓扑模。\cmt{G7-H01}{原文不加限制地在此称两端为“离散”拓扑模。代数同构总成立，但按本段刚定义的拓扑，离散性并非自动；存在开理想 $\mathfrak a$ 满足 $\mathfrak aV\subset U$ 足以保证它，前面所讨论的模拓扑条件正提供这一点。本版将该条件写入正文。一般反例见核查报告G7-H01；后文用于 $V=\mathfrak j/\mathfrak j^2$ 的情形满足此条件。}

\egap{19.5.2}
设 $A$ 为拓扑环，$B$ 为交换拓扑 $A$-代数，$\mathfrak j$ 为 $B$ 的理想，不要求开或闭。在以下各段中，$\mathfrak j$ 及其幂 $\mathfrak j^n$（$n\geq2$）均取 $B$ 的诱导拓扑，而
\[
 C=B/\mathfrak j,\qquad
 \mathfrak j^n/\mathfrak j^{n+1}=\gr_{\mathfrak j}^n(B)
\]
取商拓扑。于是各 $\mathfrak j^n/\mathfrak j^{n+1}$ 都是拓扑 $C$-模。典范包含
$\mathfrak j/\mathfrak j^2\to\gr_{\mathfrak j}^{\bullet}(B)$ 延拓为一个先不计拓扑的 $C$-代数同态
\[
 \varphi:\Salg_C(\mathfrak j/\mathfrak j^2)
        \longrightarrow\gr_{\mathfrak j}^{\bullet}(B),
\]
并在每个次数 $n$ 给出满 $C$-模同态
\begin{equation}
 \varphi_n:\Salg_C^n(\mathfrak j/\mathfrak j^2)
         \longrightarrow\gr_{\mathfrak j}^n(B).
 \tag{19.5.2.1}
\end{equation}
给对称代数赋予（19.5.1）中的拓扑后，各 $\varphi_n$ 连续。这可以将该段的泛性质应用于拓扑 $A$-代数
\[
 E=\gr_{\mathfrak j}^{\bullet}(B/\mathfrak j^{n+1}),
\]
其中有限个分次部分 $\mathfrak j^h/\mathfrak j^{h+1}$ 配备乘积拓扑，再应用于典范包含 $\mathfrak j/\mathfrak j^2\to E$ 而得到。这里 $\mathfrak j/\mathfrak j^2$ 的拓扑不比由 $C$ 诱导的模拓扑更细；后者也就是由 $B$ 诱导的模拓扑。

\egas{19.5.3}{定理}
设 $A$ 为拓扑环，$B$ 为拓扑 $A$-代数，$\mathfrak j$ 为 $B$ 的理想，$C=B/\mathfrak j$ 为商拓扑 $A$-代数。假设 $C$ 为形式光滑 $A$-代数。则：

（i）若 $B$ 为形式光滑 $A$-代数，则 $\mathfrak j/\mathfrak j^2$ 是形式投射的拓扑 $C$-模（19.2.1）。

（ii）若 $B$ 为形式光滑 $A$-代数，且为预可容许环（$0_{\mathrm I}$，7.1.2），则（19.5.2）中的各 $\varphi_n$ 是形式双态射（19.1.2）。

（iii）反之，若 $B$ 预可容许，$\mathfrak j^n\to0$，$\mathfrak j/\mathfrak j^2$ 为形式投射 $C$-模，且各 $\varphi_n$ 为形式双态射，则 $B$ 为形式光滑 $A$-代数。

此定理的证明较长，并牵涉许多技术细节。因此，我们先证明一个更简单的推论，使主要思想显现出来；这个推论也是以后最常使用的（19.5.3）的特殊情形。

\egas{19.5.4}{推论}
设 $A$ 为拓扑环，$B$ 为拓扑 $A$-代数，$\mathfrak j$ 为 $B$ 的理想，且 $B$ 的拓扑就是 $\mathfrak j$-预进制拓扑。假设离散 $A$-代数 $C=B/\mathfrak j$ 形式光滑。则以下三条件等价：

$a)$ $B$ 为形式光滑 $A$-代数；

$b)$ $\mathfrak j/\mathfrak j^2$ 是投射 $C$-模，且典范同态
\[
 \varphi:\Salg_C(\mathfrak j/\mathfrak j^2)
           \longrightarrow\gr_{\mathfrak j}^{\bullet}(B)
\]
为双射；

$c)$ $B$ 的分离完备化 $\widehat B$，作为拓扑 $A$-代数，同构于某个 $\widehat{B'}$。这里 $B'=\Salg_C(V)$，$V$ 为投射 $C$-模，而 $B'$ 取由其增广理想 $B'^+$ 定义的预进制拓扑。

\egap{19.5.4.1}
先证明 $a)$ 蕴含 $\mathfrak j/\mathfrak j^2$ 是投射 $C$-模。设 $P,Q$ 为 $C$-模，$u:P\to Q$ 为满 $C$-同态，$v:\mathfrak j/\mathfrak j^2\to Q$ 为 $C$-同态。将 $B/\mathfrak j^2$ 看作 $C$ 以 $\mathfrak j/\mathfrak j^2$ 为核的 $B$-扩张，由 $v$ 得到 $C$ 以 $Q$ 为核的 $B$-扩张
\[
 E=(B/\mathfrak j^2)\boxplus_{\mathfrak j/\mathfrak j^2}Q
\]
（18.2.8）。因 $C$ 是形式光滑的离散 $A$-代数，$E$ 为 $A$-平凡扩张（19.4.4），因而可识别为 $D_C(Q)$（18.2.3）。由 $u$ 典范得到满 $A$-同态
$\theta:D_C(P)\to D_C(Q)$（18.2.9），其核 $\mathfrak a$ 包含于 $P$，故平方为零。令 $E'=D_C(P)$，并令 $f:B\to E=E'/\mathfrak a$ 为定义 $E$ 的 $B$-代数结构的同态。$f$ 连续；由 $B$ 形式光滑，它分解为
\[
 B\xrightarrow{g}E'\xrightarrow{\theta}E,
\]
其中 $g$ 是连续 $A$-同态。若 $v':B/\mathfrak j^2\to E$ 是由 $v$ 得到的同态，则
\[
\begin{tikzcd}
 E'\arrow[r,"\theta"]&E\\
 B\arrow[u,"g"]\arrow[r]&B/\mathfrak j^2\arrow[u,"v'"']
\end{tikzcd}
\]
交换。因 $(\theta g)(\mathfrak j)\subset Q$，有 $g(\mathfrak j)\subset P+\mathfrak a=P$，故 $g(\mathfrak j^2)=0$。限制到 $\mathfrak j$，得到交换图
\[
\begin{tikzcd}
 P\arrow[r,"u"]&Q\\
 &\mathfrak j/\mathfrak j^2\arrow[ul,"w"]\arrow[u,"v"']
\end{tikzcd}
\]
这证明了 $\mathfrak j/\mathfrak j^2$ 的投射性。

\egap{19.5.4.2}
再证明 $a)$ 蕴含 $\varphi$ 双射，从而证明 $a)\Rightarrow b)$。令
\[
 E_n=B/\mathfrak j^{n+1},\qquad
 F_n=\Salg_C(\mathfrak j/\mathfrak j^2)/(\Salg_C^+(\mathfrak j/\mathfrak j^2))^{n+1}.
\]
$\mathfrak j/\mathfrak j^{n+1}$ 是 $E_n$ 的幂零理想，而离散 $A$-代数 $C$ 形式光滑，因此 $C$ 的恒等同态分解为
\[
 C\xrightarrow{f}E_n\longrightarrow C=E_n/(\mathfrak j/\mathfrak j^{n+1}),
\]
其中 $f$ 为 $A$-同态。再由（19.5.4.1），$\mathfrak j/\mathfrak j^2$ 的恒等映射可分解为
\[
 \mathfrak j/\mathfrak j^2\xrightarrow{g}\mathfrak j/\mathfrak j^{n+1}
          \longrightarrow\mathfrak j/\mathfrak j^2,
\]
其中 $g$ 为 $C$-线性映射。$f,g$ 典范地给出 $C$-代数同态
$\Salg_C(\mathfrak j/\mathfrak j^2)\to E_n$，它将增广理想的 $(n+1)$ 次幂送到零，故诱导满 $A$-代数同态
\[
 v:F_n\longrightarrow E_n.
\]
这里 $\gr^0(v),\gr^1(v)$ 分别为 $C$、$\mathfrak j/\mathfrak j^2$ 的恒等同态，且由 $\varphi$ 的定义，对 $j\leq n$ 有 $\gr^j(v)=\varphi_j$。

$\mathfrak N=\Ker v$ 是 $F_n$ 的幂零理想，故 $E_n=F_n/\mathfrak N$。由 $B$ 形式光滑，连续典范同态 $p_n:B\to E_n$ 分解为
\[
 B\xrightarrow{w}F_n\xrightarrow{v}E_n,
\]
$w$ 连续。由于 $\gr^0(v)$ 为恒等映射，$w(\mathfrak j)\subset F_n^+$，所以 $w(\mathfrak j^{n+1})=0$。于是 $w=w' p_n$，其中 $w':E_n\to F_n$，且 $vw'=1_{E_n}$。$\gr^0(w'),\gr^1(w')$ 也为恒等映射。
因 $\gr^\bullet(F_n)$ 作为 $C$-代数由一次部分生成，复合
\[
 \gr^j(F_n)\xrightarrow{\gr^j(v)}\gr^j(E_n)
       \xrightarrow{\gr^j(w')}\gr^j(F_n)
\]
对每个 $j\leq n$ 都为恒等映射：它在零次、一次上为恒等映射，而乘法决定其余次数。\cmt{G7-T01}{原文在这一步写“$\gr^\bullet(E_n)$ 由一次部分生成”，而所要判断的是 $\gr^\bullet(F_n)$ 上的复合自同态。正文将对象改为 $F_n$。两种分次代数确实都由一次部分生成，但此处的源、靶必须与复合一致。}
取 $j=n$ 得 $\varphi_n$ 单射；它本已满射，因此 $a)\Rightarrow b)$ 得证。

\egap{19.5.4.3}
接着证明 $b)\Rightarrow a)$。如（19.5.4.2）开头所构造的满 $A$-同态 $v:F_n\to E_n$ 满足 $\gr^j(v)=\varphi_j$。由 $b)$，这些同态均双射；两个滤过都有限，所以 $v$ 双射（Bourbaki，\textit{交换代数}，第三章，§2，第8目，定理1的推论3）。

设 $G$ 为离散拓扑 $A$-代数，$\mathfrak N^2=0$，且 $f:B\to G/\mathfrak N$ 为连续 $A$-同态。存在 $m\geq1$ 使 $f(\mathfrak j^m)=0$；取 $n=2m$，则 $f$ 分解为 $B\to E_n\xrightarrow{f_n}G/\mathfrak N$。复合
\[
 r:C\longrightarrow F_n\xrightarrow{v}E_n\xrightarrow{f_n}G/\mathfrak N
\]
是连续 $A$-同态。由 $C$ 形式光滑，$r$ 可提升为 $r':C\to G$，使 $G$ 成为 $C$-代数。
将 $f_n v$ 限制到 $\mathfrak j/\mathfrak j^2$ 得 $C$-线性映射 $g:\mathfrak j/\mathfrak j^2\to G/\mathfrak N$；由投射性，它提升为 $C$-线性 $h:\mathfrak j/\mathfrak j^2\to G$。延拓后得 $C$-代数同态
\[
 w:\Salg_C(\mathfrak j/\mathfrak j^2)\to G.
\]
由构造，每个 $m$ 次齐次元素在 $w$ 下的像属于 $\mathfrak N$，所以每个 $2m$ 次齐次元素的像为零。因而 $w$ 通过 $F_n$ 分解为 $w':F_n\to G$。复合 $F_n\xrightarrow{w'}G\to G/\mathfrak N$ 与 $f_n v$ 在 $C$ 及 $\mathfrak j/\mathfrak j^2$ 上相同，故相等。于是
\[
 B\longrightarrow E_n\xrightarrow{v^{-1}}F_n\xrightarrow{w'}G\longrightarrow G/\mathfrak N
\]
等于 $f$；中间的提升连续，因为它通过离散商 $E_n$。由（19.4.3），$B$ 形式光滑。

\egap{19.5.4.4}
最后证明 $a)$ 与 $c)$ 等价。先设 $c)$ 成立。$B'=\Salg_C(V)$ 在离散拓扑下是形式光滑 $C$-代数（19.3.2），故在 $B'^+$-预进制拓扑下也形式光滑（19.3.8）。$C$ 作为 $A$-代数形式光滑，故 $B'$ 作为 $A$-代数形式光滑（19.3.5，ii）；再由（19.3.6），$\widehat{B'}$、从而 $\widehat B$ 和 $B$ 均形式光滑。

反之，设 $a)$ 成立。由 $C$ 形式光滑，（19.3.10）给出分解
\[
 C\longrightarrow\widehat B\longrightarrow B/\mathfrak j=C
\]
为恒等同态，其中第一箭头为 $A$-同态。\cmt{G7-H02}{此处按原版页面保留 $\widehat B$ 上的帽号；不是断言 $C\to B/\mathfrak j$ 可直接提升到尚未完备的 $B$。解析文本多次漏识该帽号，已与原图核对，不列为原作者错误。}
因此 $\widehat B$ 以及所有 $B/\mathfrak j^{n+1}$ 都有相容的 $C$-代数结构。
因 $\mathfrak j/\mathfrak j^2$ 投射，可以从其到 $B/\mathfrak j^2$ 的典范包含出发，逐次选取相容的 $C$-同态
\[
 u_n:\mathfrak j/\mathfrak j^2\longrightarrow B/\mathfrak j^{n+1}
 \qquad(n\geq1),
\]
并使像落在 $\mathfrak j/\mathfrak j^{n+1}$ 中。由对称代数的泛性质，得到相容的 $C$-代数同态
\[
 v_n:B'=\Salg_C(\mathfrak j/\mathfrak j^2)\longrightarrow B/\mathfrak j^{n+1}.
\]
$v_n$ 在 $(B'^+)^{n+1}$ 上为零，且 $\gr^j(v_n)=\varphi_j$（$0\leq j\leq n$）。由 $b)$ 和有限滤过判据，诱导同态
\[
 B'/(B'^+)^{n+1}\longrightarrow B/\mathfrak j^{n+1}
\]
对每个 $n$ 都是同构。取射影极限，便得到 $c)$。

\egas{19.5.5}{注}
在（19.5.4）的一般假设下，进一步假设 $\mathfrak j$ 是 $B$ 的极大理想，因而 $C=B/\mathfrak j$ 是域，并假设 $\mathfrak j/\mathfrak j^2$ 为有限维 $C$-向量空间。则该推论的三个条件还等价于：

$d)$ 对每个多项式代数 $E=C[T_1,\ldots,T_n]$，令 $\mathfrak n=(T_1,\ldots,T_n)$；对每个包含某个 $\mathfrak n^k$ 的理想 $\mathfrak b$，每个保持 $C$-增广的 $A$-代数同态 $v:B\to E/\mathfrak b$ 均分解为
\[
 B\xrightarrow{w}E/\mathfrak n^k\longrightarrow E/\mathfrak b,
\]
其中 $w$ 为 $A$-同态。这里 $E/\mathfrak b$ 的 $C$-增广要求 $\mathfrak b\subseteq\mathfrak n$。\cmt{G7-H03}{“$C$-增广代数”本身要求增广能通过商环；这里将其包含条件写明。由保持增广及 $\mathfrak n^k\subset\mathfrak b$ 可得 $v(\mathfrak j^k)=0$，因此所用同态对预进制拓扑连续；这不是省略连续性后任意环同态的提升断言。}

必要性来自形式光滑性；商同态 $E/\mathfrak n^k\to E/\mathfrak b$ 的核幂零。为证充分性，由 $\mathfrak j/\mathfrak j^2$ 是自由 $C$-模，只须证明典范 $\varphi$ 双射。此时
\[
 \Salg_C(\mathfrak j/\mathfrak j^2)\simeq C[T_1,\ldots,T_n],
 \qquad n=\rk_C(\mathfrak j/\mathfrak j^2).
\]
对每个 $k\geq0$，如（19.5.4.2），$C$ 形式光滑给出满 $A$-同态
\[
 f:E/\mathfrak n^{k+1}\longrightarrow B/\mathfrak j^{k+1},
 \qquad \gr^j(f)=\varphi_j\quad(j\leq k).
\]
若 $\mathfrak b/\mathfrak n^{k+1}=\Ker f$，则 $B/\mathfrak j^{k+1}\simeq E/\mathfrak b$；由 $d)$，典范同态 $B\to B/\mathfrak j^{k+1}$ 分解为
\[
 B\xrightarrow{w}E/\mathfrak n^{k+1}\xrightarrow{f}B/\mathfrak j^{k+1}.
\]
因 $\gr^0(f)$ 为恒等映射，$w(\mathfrak j)\subset\mathfrak n/\mathfrak n^{k+1}$，故 $w(\mathfrak j^{k+1})=0$。如（19.5.4.2），得到 $\varphi_k$ 对每个 $k$ 都单射，从而得到断言。

% END continuations/iv0_19_5a.tex


% BEGIN continuations/iv0_19_5b.tex
% 一般拓扑情形的完整证明。原证明的核比较缺步在19.5.6.7明示修补。
\egap{19.5.6}
\textbf{定理（19.5.3）的证明。}
设 $(\mathfrak b_\alpha)$ 是 $B$ 的一个递降开理想基本系，令
\[
 B_\alpha=B/\mathfrak b_\alpha,\qquad
 C_\alpha=B/(\mathfrak b_\alpha+\mathfrak j),\qquad
 \mathfrak j_\alpha=(\mathfrak b_\alpha+\mathfrak j)/\mathfrak b_\alpha.
\]
于是
\[
 C_\alpha=B_\alpha/\mathfrak j_\alpha,\qquad
 \mathfrak j_\alpha^n/\mathfrak j_\alpha^{n+1}
  =(\mathfrak b_\alpha+\mathfrak j^n)/(\mathfrak b_\alpha+\mathfrak j^{n+1}).
\]
$C$-模 $\mathfrak j^n/\mathfrak j^{n+1}$ 的开子模
\[
 \bigl((\mathfrak b_\alpha\cap\mathfrak j^n)+\mathfrak j^{n+1}\bigr)
         /\mathfrak j^{n+1}
\]
构成其零点基本邻域系。由于
\[
 (\mathfrak b_\alpha\cap\mathfrak j^n)+\mathfrak j^{n+1}
   =\mathfrak j^n\cap(\mathfrak b_\alpha+\mathfrak j^{n+1}),
\]
有典范识别
\[
 \frac{\mathfrak j^n}{(\mathfrak b_\alpha\cap\mathfrak j^n)+\mathfrak j^{n+1}}
 \simeq\frac{\mathfrak b_\alpha+\mathfrak j^n}
                  {\mathfrak b_\alpha+\mathfrak j^{n+1}}
 \simeq\frac{\mathfrak j_\alpha^n}{\mathfrak j_\alpha^{n+1}}.
\]

\egap{19.5.6.1}
\textbf{（19.5.3，i）的证明。}
设 $P,Q$ 为离散 $C_\alpha$-模，$u:P\to Q$ 为满 $C_\alpha$-同态。离散环 $B/(\mathfrak b_\alpha+\mathfrak j^2)$ 是 $C_\alpha$ 以平方零理想 $\mathfrak j_\alpha/\mathfrak j_\alpha^2$ 为核的 $B$-扩张。
给定连续 $C$-同态 $v:\mathfrak j/\mathfrak j^2\to Q$。必要时将 $\mathfrak b_\alpha$ 换成更小的开理想，可假设
\[
 \bigl((\mathfrak b_\alpha\cap\mathfrak j)+\mathfrak j^2\bigr)/\mathfrak j^2
       \subseteq\Ker v.
\]
因此 $v$ 诱导离散模同态 $v':\mathfrak j_\alpha/\mathfrak j_\alpha^2\to Q$。
沿 $v'$ 变换上述扩张的核，得到 $C_\alpha$ 以 $Q$ 为核的 $B_\alpha$-扩张 $E_\alpha$；令
\[
 \psi:B/(\mathfrak b_\alpha+\mathfrak j^2)\longrightarrow E_\alpha
\]
为对应同态（18.2.8）。$E_\alpha$ 是离散拓扑 $A$-代数，而典范 $C_\alpha\simeq E_\alpha/Q$ 与 $C\to C_\alpha$ 复合，给出连续 $A$-同态 $f:C\to E_\alpha/Q$。

由 $C$ 形式光滑及 $Q^2=0$，存在连续 $A$-同态 $g:C\to E_\alpha$ 提升 $f$。因 $E_\alpha$ 离散，有 $\beta\geq\alpha$ 使 $g$ 分解为
\[
 C\longrightarrow C_\beta\xrightarrow{g'}E_\alpha.
\]
沿 $C_\beta\to C_\alpha$ 拉回 $E_\alpha$，得到 $C_\beta$ 以 $Q$ 为核的 $A$-扩张 $E_\beta$。$g'$ 的存在表明此扩张 $A$-平凡，因此可识别为 $D_{C_\beta}(Q)$。满射 $u:P\to Q$ 给出
\[
 E'_\beta=D_{C_\beta}(P)\longrightarrow D_{C_\beta}(Q)=E_\beta,
\]
其核 $\mathfrak a\subset P$ 平方为零。由拉回构造有连续 $A$-同态 $f':B\to E_\beta$；由 $B$ 形式光滑，它提升为连续 $h':B\to E'_\beta$。

令 $\theta:E'_\beta\to E_\beta\to E_\alpha$ 为复合。图
\[
\begin{tikzcd}
 E'_\beta\arrow[r,"\theta"]&E_\alpha\\
 B\arrow[u,"h'"]\arrow[r]&B/(\mathfrak b_\alpha+\mathfrak j^2)\arrow[u,"\psi"']
\end{tikzcd}
\]
交换。$h'(\mathfrak j)\subset P$，故 $h'(\mathfrak j^2)=0$。限制到 $\mathfrak j$，得到
\[
\begin{tikzcd}
 P\arrow[r,"u"]&Q\\
 &\mathfrak j/\mathfrak j^2\arrow[ul,"w"]\arrow[u,"v"']
\end{tikzcd}
\]
其中 $w$ 连续。于是 $\mathfrak j/\mathfrak j^2$ 是形式投射 $C$-模。

\egap{19.5.6.2}
对每个 $n\geq0$，令 $E_n=B/\mathfrak j^{n+1}$，特别 $E_0=C$。其开理想基本系为
$(\mathfrak b_\alpha+\mathfrak j^{n+1})/\mathfrak j^{n+1}$，相应离散商记为
\[
 E_{\alpha,n}=B/(\mathfrak b_\alpha+\mathfrak j^{n+1})
           =B_\alpha/\mathfrak j_\alpha^{n+1}.
\]
如上所见，$\gr_{\mathfrak j}^n(B)$ 的相应离散商为
$\mathfrak j_\alpha^n/\mathfrak j_\alpha^{n+1}$。再令
\[
 F_{\alpha,n}=
 \frac{\Salg_{C_\alpha}(\mathfrak j_\alpha/\mathfrak j_\alpha^2)}
 {\bigl(\Salg_{C_\alpha}^+(\mathfrak j_\alpha/\mathfrak j_\alpha^2)\bigr)^{n+1}}.
\]
对固定的 $n\geq1$，由（19.5.1.1），各
$\Salg_{C_\alpha}^n(\mathfrak j_\alpha/\mathfrak j_\alpha^2)$
是拓扑 $C$-模 $\Salg_C^n(\mathfrak j/\mathfrak j^2)$ 对一个开子模基本系的商。
当 $\alpha\leq\beta$ 时，$B_\beta\to B_\alpha$、$C_\beta\to C_\alpha$、
$\mathfrak j_\beta/\mathfrak j_\beta^2\to\mathfrak j_\alpha/\mathfrak j_\alpha^2$、
$E_{\beta,n}\to E_{\alpha,n}$、$F_{\beta,n}\to F_{\alpha,n}$ 等典范同态，统称为\emph{过渡同态}。

\egas{19.5.6.3}{引理}
假设 $C$ 为形式光滑 $A$-代数，$B$ 预可容许，且 $\mathfrak j/\mathfrak j^2$ 为形式投射 $C$-模。则：

（i）对每个 $\alpha$，存在 $\beta\geq\alpha$ 和满 $A$-代数同态
\[
 v_{\alpha\beta}:F_{\beta,n}\longrightarrow E_{\alpha,n},
\]
使 $\gr^0(v_{\alpha\beta}):C_\beta\to C_\alpha$ 及
$\gr^1(v_{\alpha\beta}):\mathfrak j_\beta/\mathfrak j_\beta^2\to
\mathfrak j_\alpha/\mathfrak j_\alpha^2$ 均为过渡同态。

（ii）若 $\beta,v_{\alpha\beta}$ 满足（i），则对每个 $\gamma\geq\beta$，存在 $\delta\geq\gamma$ 和满足（i）的满 $A$-代数同态
$v_{\gamma\delta}:F_{\delta,n}\to E_{\gamma,n}$，使下图交换：
\begin{equation}
\begin{tikzcd}
 F_{\beta,n}\arrow[r,"v_{\alpha\beta}"]&E_{\alpha,n}\\
 F_{\delta,n}\arrow[r,"v_{\gamma\delta}"]\arrow[u]&E_{\gamma,n}\arrow[u]
\end{tikzcd}
\tag{19.5.6.4}
\end{equation}
其中竖直箭头为过渡同态。

\eproof
（i）在离散拓扑 $A$-代数 $E_{\alpha,n}$ 中，理想 $\mathfrak j_\alpha/\mathfrak j_\alpha^{n+1}$ 幂零，且其商为 $C_\alpha$。由 $C$ 形式光滑，典范同态 $C\to C_\alpha$ 提升为连续 $A$-同态
$f_\alpha:C\to E_{\alpha,n}$。
以 $f_\alpha$ 定义 $C$ 的作用，则 $\mathfrak j_\alpha/\mathfrak j_\alpha^{n+1}$ 是被 $C$ 的某个开理想零化的离散拓扑 $C$-模。形式投射性给出连续 $C$-线性映射
\[
 g_\alpha:\mathfrak j/\mathfrak j^2\longrightarrow
                   \mathfrak j_\alpha/\mathfrak j_\alpha^{n+1}
\]
使
\[
\begin{tikzcd}
 \mathfrak j_\alpha/\mathfrak j_\alpha^{n+1}\arrow[r]&\mathfrak j_\alpha/\mathfrak j_\alpha^2\\
 &\mathfrak j/\mathfrak j^2\arrow[ul,"g_\alpha"]\arrow[u]
\end{tikzcd}
\]
交换。由连续性，可以取 $\beta\geq\alpha$，使 $f_\alpha,g_\alpha$ 分别通过 $C_\beta$ 和 $\mathfrak j_\beta/\mathfrak j_\beta^2$ 分解。它们给出
\[
 \Salg_{C_\beta}(\mathfrak j_\beta/\mathfrak j_\beta^2)\longrightarrow E_{\alpha,n},
\]
并将增广理想的 $(n+1)$ 次幂送到零，因而诱导所求 $v_{\alpha\beta}$。其零次和一次相伴分次同态均满；由于目标相伴分次代数由零次和一次部分生成，$\gr(v_{\alpha\beta})$ 满射。两个滤过均有限，故 $v_{\alpha\beta}$ 满射（Bourbaki，\textit{交换代数}，第三章，§2，第8目，定理1的推论3）。

（ii）由预可容许性，可以将基本系限制在同一个定义理想 $\mathfrak b_{\alpha_0}$ 内，其幂趋于零。因此，对 $\gamma\geq\alpha$，过渡同态 $E_{\gamma,n}\to E_{\alpha,n}$ 的核
\[
 (\mathfrak b_\alpha+\mathfrak j^{n+1})/
       (\mathfrak b_\gamma+\mathfrak j^{n+1})
\]
幂零。应用已校订的引理（19.3.10.1），可取连续 $f_\gamma:C\to E_{\gamma,n}$，使它提升 $C\to C_\gamma$ 并满足
\[
\begin{tikzcd}
 C\arrow[r,"f_\gamma"]\arrow[dr,"f_\alpha"']&E_{\gamma,n}\arrow[d]\\
 &E_{\alpha,n}.
\end{tikzcd}
\]
\cmt{G7-H04}{这里确实满足校订引理要求的“理想像相等”：过渡同态将 $\mathfrak j_\gamma/\mathfrak j_\gamma^{n+1}$ 满射到 $\mathfrak j_\alpha/\mathfrak j_\alpha^{n+1}$。因此补入该条件不破坏本处应用。旧引理按包含条件书写的问题已在第六版报告F6-E01/G02论证，本版正文已经修正。}

再由 $\mathfrak j/\mathfrak j^2$ 形式投射，可选择连续 $C$-同态
$g_\gamma:\mathfrak j/\mathfrak j^2\to\mathfrak j_\gamma/\mathfrak j_\gamma^{n+1}$，同时满足其到 $\mathfrak j_\gamma/\mathfrak j_\gamma^2$ 的复合为典范映射 $\rho_\gamma$，而到 $\mathfrak j_\alpha/\mathfrak j_\alpha^{n+1}$ 的复合为 $g_\alpha$。亦即，下图交换：
\[
\begin{tikzcd}[column sep=small]
 &\mathfrak j_\gamma/\mathfrak j_\gamma^{n+1}
       \arrow[dl]\arrow[dr] &\\
 \mathfrak j_\alpha/\mathfrak j_\alpha^{n+1}\arrow[dr]
 & &\mathfrak j_\gamma/\mathfrak j_\gamma^2\arrow[dl]\\
 &\mathfrak j_\alpha/\mathfrak j_\alpha^2&
\end{tikzcd}
\qquad
\begin{gathered}
 \mathfrak j/\mathfrak j^2\xrightarrow{g_\gamma}
              \mathfrak j_\gamma/\mathfrak j_\gamma^{n+1},\\
 g_\alpha=\operatorname{pr}_1 g_\gamma,\quad
 \rho_\gamma=\operatorname{pr}_2 g_\gamma.
\end{gathered}
\]
具体说，$(g_\alpha,\rho_\gamma)$ 的像包含于上述两个目标之积中，由
$P=\mathfrak j_\gamma/\mathfrak j_\gamma^{n+1}$ 的典范映射所得的像 $Q$。这是因为二者在 $\mathfrak j_\alpha/\mathfrak j_\alpha^2$ 中相容，而
$\mathfrak j_\gamma^2\to\mathfrak j_\alpha^2$ 满射。对 $P\twoheadrightarrow Q$ 应用形式投射性的定义，就得到 $g_\gamma$。再如（i）使 $f_\gamma,g_\gamma$ 通过足够后面的同一个指标 $\delta$ 分解，所构造的 $v_{\gamma\delta}$ 就满足（19.5.6.4）。
\finishproof

\egas{19.5.6.5}{引理}
假设 $B,C$ 都是形式光滑 $A$-代数，且 $B$ 预可容许。由（19.5.6.1），（19.5.6.3）的条件满足。对每个满足该引理（i）的 $\alpha\leq\beta$ 和 $v_{\alpha\beta}$，存在 $\lambda\geq\beta$ 及 $A$-代数同态
\[
 w_{\beta\lambda}:E_{\lambda,n}\longrightarrow F_{\beta,n},
\]
使：$1^\circ$ $\gr^0(w_{\beta\lambda})$ 和 $\gr^1(w_{\beta\lambda})$ 是对应的过渡同态；$2^\circ$ 复合
\[
 E_{\lambda,n}\xrightarrow{w_{\beta\lambda}}F_{\beta,n}
             \xrightarrow{v_{\alpha\beta}}E_{\alpha,n}
\]
是过渡同态。

\eproof
将（19.5.6.3，ii）用于 $\gamma=\beta$，得到 $\delta\geq\beta$ 及满同态
$v_{\beta\delta}:F_{\delta,n}\to E_{\beta,n}$，并使其与 $v_{\alpha\beta}$ 相容。
其核 $\mathfrak N$ 幂零：$F_{\delta,n}$ 的增广理想本来幂零，而零次同态 $C_\delta\to C_\beta$ 的核因 $B$ 预可容许也幂零。因此 $E_{\beta,n}=F_{\delta,n}/\mathfrak N$。
由 $B$ 形式光滑，典范连续同态 $B\to E_{\beta,n}$ 提升为连续同态 $w:B\to F_{\delta,n}$。因目标离散，可取 $\lambda\geq\delta$ 使 $w(\mathfrak b_\lambda)=0$，得到
$w_\lambda:B/\mathfrak b_\lambda\to F_{\delta,n}$。

令 $q_{\beta\delta}:F_{\delta,n}\to F_{\beta,n}$ 为过渡同态，并令
$w'_\lambda=q_{\beta\delta}w_\lambda$。由于 $\gr^0(q_{\beta\delta})$ 与 $\gr^0(v_{\beta\delta})$ 都是 $C_\delta\to C_\beta$，而 $v_{\beta\delta}w_\lambda$ 是典范同态，$w'_\lambda$ 将 $\mathfrak j_\lambda$ 送入 $F_{\beta,n}$ 的增广理想，故将 $\mathfrak j_\lambda^{n+1}$ 送到零。于是
\[
 B/\mathfrak b_\lambda\longrightarrow E_{\lambda,n}
        \xrightarrow{w_{\beta\lambda}}F_{\beta,n}
\]
给出 $w'_\lambda$ 的分解。\cmt{G7-T02}{原证明在此将 $w'_\lambda$ 的目标增广理想写成 $F_{\delta,n}^{+}$，而按定义它的目标是 $F_{\beta,n}$。下文最后一处 $w$ 的下标也应与 $\beta\lambda$ 一致；正文按实际源靶统一。}
由（19.5.6.4）的交换性，$v_{\alpha\beta}w_{\beta\lambda}$ 是过渡同态。
上述零次同态的比较证明 $\gr^0(w_{\beta\lambda})$ 为过渡同态；又因
$\gr^1(q_{\beta\delta})=\gr^1(v_{\beta\delta})$，同样的比较证明一次同态也为过渡同态。
\finishproof

\egap{19.5.6.6}
\textbf{（19.5.3，ii）的证明。}
我们使用（19.1.3，iii）的判据。对每个 $\alpha$，按（19.5.6.5）选取 $v_{\alpha\beta}$ 和 $w_{\beta\lambda}$。同态
\[
 \gr^n(v_{\alpha\beta}):\gr^n(F_{\beta,n})\to\gr^n(E_{\alpha,n})
\]
就是由（19.5.2.1）的 $\varphi_n$ 取商得到的
\[
 \varphi_{\alpha\beta,n}:
 \Salg_{C_\beta}^n(\mathfrak j_\beta/\mathfrak j_\beta^2)
       \to\mathfrak j_\alpha^n/\mathfrak j_\alpha^{n+1}.
\]
事实上，这在零次和一次由构造成立，而相伴分次同态与乘法相容，所以对 $j\leq n$ 均有
$\gr^j(v_{\alpha\beta})=\varphi_{\alpha\beta,j}$。

$\varphi_n$ 满射，故为形式满态射。对 $j\leq n$，复合
\[
 \gr^j(F_{\lambda,n})\xrightarrow{\varphi_{\lambda\lambda,j}}
 \gr^j(E_{\lambda,n})\xrightarrow{\gr^j(w_{\beta\lambda})}
 \gr^j(F_{\beta,n})
\]
为过渡同态：它在零次和一次成立，而源分次代数 $\gr^\bullet(F_{\lambda,n})$ 由一次部分在零次环上生成。再与 $F_{\beta,n}\to F_{\alpha,n}$ 的相伴分次同态复合，并取 $j=n$，便把过渡同态
$\gr^n(F_{\lambda,n})\to\gr^n(F_{\alpha,n})$ 分解为
\[
 \gr^n(F_{\lambda,n})\xrightarrow{\varphi_{\lambda\lambda,n}}
 \gr^n(E_{\lambda,n})\longrightarrow\gr^n(F_{\alpha,n}).
\]
这正是（19.1.3）的形式单态射判据；所以 $\varphi_n$ 为形式双态射。\cmt{G7-T03}{原文再次以 $\gr^\bullet(E_{\lambda,n})$ 的生成性来判断源为 $\gr^\bullet(F_{\lambda,n})$ 的复合。正文明确使用源的生成性，与（19.5.4.2）的处理一致。}

\egas{19.5.6.7}{引理}
假设 $B$ 预可容许，$\mathfrak j/\mathfrak j^2$ 为形式投射 $C$-模，$C$ 为形式光滑 $A$-代数，且各 $\varphi_j$ 为形式双态射。对满足（19.5.6.3，i）的 $\alpha\leq\beta$ 和 $v_{\alpha\beta}$，存在 $\gamma\geq\beta$，使对每个 $\mu\geq\gamma$ 均有交换图
\[
\begin{tikzcd}
 F_{\beta,n}\arrow[r,"v_{\alpha\beta}"]&E_{\alpha,n}\\
 &E_{\mu,n}\arrow[ul,"u_{\beta\mu}"]\arrow[u,"p_{\alpha\mu}"']
\end{tikzcd}
\]
其中 $p_{\alpha\mu}$ 为过渡同态。

\eproof
对 $0\leq j\leq n$，由（19.1.3）及 $\varphi_j$ 为形式双态射，可以将指标取得足够后，使相应过渡同态通过
$\varphi_{\gamma\gamma,j}:\gr^j(F_{\gamma,n})\to\gr^j(E_{\gamma,n})$ 分解。次数只有有限个，故可对所有次数选同一个指标。由该满同态的满射性，所得分解唯一；其各次分量与乘法相容，因此组成满的分次代数同态
\[
 w_{\beta\gamma}:\gr^\bullet(E_{\gamma,n})
           \longrightarrow\gr^\bullet(F_{\beta,n})=F_{\beta,n},
 \qquad
 w_{\beta\gamma}\varphi_{\gamma\gamma}=q_{\beta\gamma}.
\]
其中 $q_{\beta\gamma}$ 为 $F$ 系的过渡同态。

以下用有限滤过逐次消去核，补足从相伴分次信息返回原环同态的步骤。\cmt{G7-G01}{原证明在印刷第98页直接使用 $\Ker h\subseteq\Ker(\gr h)$，把滤过源识别为其分次直和。对一般保持滤过的同态，此包含并不成立：只能得到最低非零齐次项落在相应分次核中。本版用下述有限次逐级消去替换原证明的这一步；反例、原论证的定位及补证细节见核查报告G7-G01。这是证明修补，不是对引理结论的反例。}
令 $\alpha_0=\alpha$、$\beta_0=\beta$、$v_0=v_{\alpha\beta}$。递归进行 $N=n+1$ 次以下选取。
已有 $\alpha_i,\beta_i,v_i:F_{\beta_i,n}\to E_{\alpha_i,n}$ 时，先选
$\alpha_{i+1}\geq\beta_i$，使对每个 $0\leq r\leq n$ 都有分解
\[
 q_{\beta_i\alpha_{i+1}}^{\,r}
   =w_i^r\,\varphi_{\alpha_{i+1}\alpha_{i+1},r}.
\]
再由（19.5.6.3，ii）选 $\beta_{i+1}\geq\alpha_{i+1}$ 和满同态
$v_{i+1}:F_{\beta_{i+1},n}\to E_{\alpha_{i+1},n}$，使
\[
 p_{\alpha_i\alpha_{i+1}}v_{i+1}
     =v_iq_{\beta_i\beta_{i+1}}.
\]

以 $F_{\beta_i,n}^{\geq r}$ 表示增广理想的 $r$ 次幂。若
$x\in\Ker(v_{i+1})\cap F_{\beta_{i+1},n}^{\geq r}$，则其 $r$ 次分量 $x_r$ 满足
\[
 \varphi_{\alpha_{i+1}\alpha_{i+1},r}
       q_{\alpha_{i+1}\beta_{i+1}}^{\,r}(x_r)=0,
\]
因为 $v_{i+1}(x)=0$。应用 $w_i^r$，得
$q_{\beta_i\beta_{i+1}}^{\,r}(x_r)=0$；故
\[
 q_{\beta_i\beta_{i+1}}(x)
   \in \Ker(v_i)\cap F_{\beta_i,n}^{\geq r+1}.
\]
这里属于 $\Ker(v_i)$ 使用上面的交换关系。连续应用 $N=n+1$ 次后，因
$F_{\beta_0,n}^{\geq n+1}=0$，得到
\[
 q_{\beta_0\beta_N}(\Ker v_N)=0.
\]
由于 $v_N$ 满射，存在唯一 $A$-代数同态
\[
 u:E_{\alpha_N,n}\longrightarrow F_{\beta_0,n},
 \qquad uv_N=q_{\beta_0\beta_N}.
\]
又有
$v_0uv_N=v_0q_{\beta_0\beta_N}=p_{\alpha_0\alpha_N}v_N$，由满射性可消去 $v_N$，从而
$v_0u=p_{\alpha_0\alpha_N}$。令 $\gamma=\alpha_N$，对任意 $\mu\geq\gamma$，取
$u_{\beta\mu}=u\,p_{\gamma\mu}$，即得所求交换图。
\finishproof

\egap{19.5.6.8}
\textbf{（19.5.3，iii）的证明。}
设 $G$ 为离散拓扑 $A$-代数，$\mathfrak N^2=0$，$f:B\to G/\mathfrak N$ 为连续 $A$-同态。存在 $\alpha$ 使 $f(\mathfrak b_\alpha)=0$；由假设 $\mathfrak j^m\to0$，又有 $m\geq1$ 使 $\mathfrak j^m\subset\mathfrak b_\alpha$。取 $n=2m$，则
\[
 f:B\xrightarrow{p_\alpha}E_{\alpha,n}
                    \xrightarrow{f_{\alpha,n}}G/\mathfrak N.
\]
由（19.5.6.3），先取 $\beta\geq\alpha$ 及复合同态
\begin{equation}
 f_{\alpha,n}v_{\alpha\beta}:
 F_{\beta,n}\xrightarrow{v_{\alpha\beta}}E_{\alpha,n}
                  \xrightarrow{f_{\alpha,n}}G/\mathfrak N.
 \tag{19.5.6.9}
\end{equation}
$F_{\beta,n}$ 为 $C_\beta$-代数，因而是离散拓扑 $C$-代数；限制（19.5.6.9）到系数，得到连续 $A$-同态 $r:C\to G/\mathfrak N$。由 $C$ 形式光滑，取连续提升 $r':C\to G$，使 $G$ 成为离散拓扑 $C$-代数。

将（19.5.6.9）与
$\mathfrak j/\mathfrak j^2\to\mathfrak j_\beta/\mathfrak j_\beta^2\to F_{\beta,n}$
复合，得到连续 $C$-同态 $g:\mathfrak j/\mathfrak j^2\to G/\mathfrak N$。由形式投射性，取连续提升 $h:\mathfrak j/\mathfrak j^2\to G$。可取 $\gamma\geq\beta$，使 $h$ 在
$((\mathfrak b_\gamma\cap\mathfrak j)+\mathfrak j^2)/\mathfrak j^2$ 上为零，同时 $r'$ 在
$(\mathfrak b_\gamma+\mathfrak j)/\mathfrak j$ 上为零。于是 $h$ 通过 $C_\gamma$-线性映射
$h':\mathfrak j_\gamma/\mathfrak j_\gamma^2\to G$ 分解，并延拓为
\[
 w:\Salg_{C_\gamma}(\mathfrak j_\gamma/\mathfrak j_\gamma^2)\to G.
\]
由构造，$m$ 次齐次元素的像属于 $\mathfrak N$，$n=2m$ 次齐次元素的像因此为零。所以 $w$ 诱导 $w'_\gamma:F_{\gamma,n}\to G$。

对 $v_{\alpha\gamma}=v_{\alpha\beta}q_{\beta\gamma}$ 应用（19.5.6.7），取得 $\delta\geq\gamma$ 及
$u_{\gamma\delta}:E_{\delta,n}\to F_{\gamma,n}$，满足
$v_{\alpha\gamma}u_{\gamma\delta}=p_{\alpha\delta}$。于是有交换图
\[
\begin{tikzcd}[column sep=large,row sep=large]
 &E_{\alpha,n}\arrow[r,"f_{\alpha,n}"]&G/\mathfrak N\\
 &F_{\beta,n}\arrow[u,"v_{\alpha\beta}"]&\\
 B\arrow[r,"p_\delta"]&E_{\delta,n}\arrow[uu,"p_{\alpha\delta}",bend left=40]
   \arrow[r,"u_{\gamma\delta}"']&F_{\gamma,n}\arrow[ul,"q_{\beta\gamma}"']
   \arrow[r,"w'_\gamma"']&G\arrow[uul,bend right=35]
\end{tikzcd}
\]
下方的复合 $f':B\to G$ 连续，且 $G\to G/\mathfrak N$ 与它复合为 $f$。因此 $B$ 形式光滑；定理（19.5.3）的证明完成。

\egas{19.5.7}{推论}
设 $A$ 为拓扑环，$B$ 为拓扑 $A$-代数，$(\mathfrak b_\lambda)$ 为其开理想基本系，$\mathfrak j$ 为 $B$ 的理想，$C=B/\mathfrak j$ 取商拓扑，并令 $C_\lambda=B/(\mathfrak b_\lambda+\mathfrak j)$。假设：

$1^\circ$ 对每个 $n$，$B$ 在 $\mathfrak j^n$ 上诱导的子空间拓扑，也等于由 $B$ 的拓扑给出的 $B$-模拓扑（19.0.2）。特别地，当 $B$ 诺特且拓扑为预进制拓扑时，此条件满足（$0_{\mathrm I}$，7.3.2）。

$2^\circ$ $C$ 为形式光滑 $A$-代数。

则：

（i）若 $B$ 形式光滑，则对每个 $\lambda$，$(\mathfrak j/\mathfrak j^2)\otimes_CC_\lambda$ 为投射 $C_\lambda$-模。

（ii）若 $B$ 形式光滑且预可容许，则对每个 $\lambda$，典范同态
\begin{equation}
 \varphi_\lambda=\varphi\otimes1_{C_\lambda}:
 \Salg_{C_\lambda}\bigl((\mathfrak j/\mathfrak j^2)\otimes_CC_\lambda\bigr)
      \longrightarrow\gr_{\mathfrak j}^{\bullet}(B)\otimes_CC_\lambda
 \tag{19.5.7.1}
\end{equation}
双射。

（iii）反之，若 $B$ 预可容许，$\mathfrak j^n\to0$，并且对每个 $\lambda$ 上述模为投射 $C_\lambda$-模、（19.5.7.1）双射，则 $B$ 为形式光滑 $A$-代数。\cmt{G7-T04}{初版（iii）将公式引作（19.5.4.1），应为本推论的（19.5.7.1）。正文修正该引文；原页中模的基环确为 $C_\lambda$，不得把提取文本遗漏下标误判为原版错误。}

事实上，在此条件下 $\mathfrak j/\mathfrak j^2$、各 $\Salg_C^n(\mathfrak j/\mathfrak j^2)$ 以及各相伴分次部分的拓扑都由 $B$ 诱导（19.5.1）；结论因此由（19.5.3）及适用于此类拓扑的（19.1.4）、（19.2.4）直接得到。

\egas{19.5.8}{注}
在（19.5.7）的假设下，进一步假设 $\mathfrak j/\mathfrak j^2$ 为有限型 $C$-模，而 $C$ 配备商拓扑后是Zariski环。取 $C$ 的定义理想 $\mathfrak r$，其幂构成零点基本邻域系。则（19.5.7）中的（i）、（ii）可分别替换为：

（i$'$）$\mathfrak j/\mathfrak j^2$ 是投射 $C$-模；

（ii$'$）$\varphi:\Salg_C(\mathfrak j/\mathfrak j^2)\to\gr_{\mathfrak j}^{\bullet}(B)$ 双射。

（i$'$）显然蕴含原（i）。反之，若各 $(\mathfrak j/\mathfrak j^2)\otimes_CC_\lambda$ 投射，则对每个 $n$，
$(\mathfrak j/\mathfrak j^2)\otimes_C(C/\mathfrak r^n)$
为投射、从而平坦的 $(C/\mathfrak r^n)$-模。由（$0_{\mathrm{III}}$，10.2.2），$\mathfrak j/\mathfrak j^2$ 平坦；它又是有限表现模，故投射（Bourbaki，\textit{交换代数}，第二章，§5，第2目，定理1的推论2）。
另一方面，$\Salg_C^n(\mathfrak j/\mathfrak j^2)$ 与 $\gr_{\mathfrak j}^n(B)$ 都是有限型 $C$-模。对Zariski环上的有限型模，同态双射等价于其完备化双射（同书第三章，§3，第5目，命题9）。故（ii）与（ii$'$）等价。\cmt{G7-H05}{这里沿用原著所引用的Zariski环约定及（19.5.7）的拓扑相容假设；不能仅保留“某个线性拓扑环”和“有限型模”便推出有限表现性或上述完备化判据。}

% END continuations/iv0_19_5b.tex


% BEGIN continuations/iv0_19_6.tex
% 底本 EGAIV.1.pdf，PDF第99--103页，分册印刷第100--104页。
\subsection{域上的代数的情形}
\source{第四卷第一部分，印刷第100--104页；上传PDF第99--103页}
\egas{19.6.1}{定理（Cohen）}
设 $k$ 为域，$K$ 为 $k$ 的扩域，二者均取离散拓扑。则 $K$ 是形式光滑的 $k$-代数，当且仅当扩张 $K/k$ 可分。

必要性将在（19.6.5.1）中证明，在此之前不使用它。这里仅证明充分性，分两种情形。

\noindent\textbf{I．$K/k$ 为有限生成可分扩张。}
由域论（Bourbaki，\textit{代数}，第五章，§9，第3目，定理2），存在纯超越子扩张
$K'=k(T_1,\ldots,T_n)$，使 $K/K'$ 为有限可分代数扩张。由（19.3.5，ii），只须处理纯超越扩张及有限可分代数扩张两种情形。
前一种情形中，$k[T_1,\ldots,T_n]$ 为形式光滑 $k$-代数（19.3.3），故其分式域也形式光滑（19.3.5，iv）。
后一种情形中，对任意与基域作用相容的 $(K,K)$-双模 $L$，均有
\[
 H_k^i(K,L)=0\qquad(i>0).
\]
\cmt{G7-T05}{原文这里说“所有Hochschild上同调群”均为零，须理解为正次数；例如 $K=L=k$ 时 $H_k^0(k,k)=k$。本处证明仅需二次上同调消失，正文写明 $i>0$，不改变后续结论。}
事实上，令 $C=K\otimes_kK$，则 $L$ 是左 $C$-模，而
$H_k^i(K,L)=\Ext_C^i(K,L)$；这里 $K$ 通过乘法同态成为 $C$-模（(M)，IX，4）。因 $K/k$ 有限可分，$C$ 是有限个域的直积，其中通过乘法同态得到的一个直因子就是 $K$（Bourbaki，\textit{代数}，第八章，§8，命题3）。因此 $K$ 是投射 $C$-模，得到上述消失。以 $L$ 为核的每个 $k$-扩张都为 $k$-平凡扩张，尤其交换扩张如此；由（19.4.4）即得结论。

\noindent\textbf{II．一般情形。}
沿用（18.4.5）的记号，只须证明对每个 $K$-向量空间 $L$，有
\[
 \Hom_K\bigl(H_2(P'_\bullet),L\bigr)=0,
\]
亦即 $H_2(P'_\bullet)=0$。若 $K$ 是其子扩张 $K_\alpha$ 的滤过并，则 $P'_\bullet$ 是对应复形 $P'^\alpha_\bullet$ 的归纳极限，因为模的滤过归纳极限与张量积可交换。又因该归纳极限函子正合，
\[
 H_2(P'_\bullet)=\varinjlim_\alpha H_2(P'^\alpha_\bullet).
\]
$K/k$ 可分蕴含每个子扩张可分；而 $K$ 是其有限生成子扩张的滤过并。由第一部分，各右端同调模为零，故 $K$ 形式光滑。\finishproof

\egas{19.6.2}{推论}
设 $A$ 为分离完备的局部环，$K$ 为其剩余域，并\emph{给定一个子域 $k\subset A$}，使通过剩余映射识别后的扩张 $K/k$ 可分。则存在包含 $k$ 的子域 $K'\subset A$，使典范同态 $A\to K$ 在 $K'$ 上的限制为 $K'\xrightarrow{\sim}K$。\cmt{G7-E01}{初版只写“$k$ 为 $K$ 的子域”，未给出相容的 $k\hookrightarrow A$，但结论要求 $K'\subset A$ 包含 $k$，证明也使用 $A$ 的 $k$-代数结构。按初版假设取 $A=\mathbf Z_p$、$K=k=\mathbf F_p$ 即不成立：$A$ 的特征为零，不能包含特征为 $p$ 的子域。正文补明所需嵌入；它在下一推论的应用中本来已经具备。}

设 $\mathfrak m$ 为 $A$ 的极大理想。由（19.6.1），$K$ 为形式光滑的离散 $k$-代数。对完备环 $A$ 及理想 $\mathfrak m$ 应用（19.3.10），恒等同态 $K\to A/\mathfrak m=K$ 提升为 $k$-同态 $f:K\to A$。域到非零环的保单位同态必单射；因此 $f$ 将 $K$ 同构地映到包含 $k$ 的子域 $K'=f(K)$，而剩余映射限制到 $K'$ 是它的逆。

\egas{19.6.3}{推论}
设 $A$ 为完备诺特局部环，$\mathfrak m$ 为其极大理想，$k=A/\mathfrak m$ 为剩余域。以下条件等价：

$a)$ $A$ 包含一个子域；

$b)$ 若 $p=\operatorname{char}k$，则 $pA=0$；

$c)$ 存在域 $K$，使 $A$ 同构于形式幂级数环 $B=K[[T_1,\ldots,T_n]]$ 的一个商环。

这些条件满足时，可以使 $A\simeq B/\mathfrak b$，其中 $\mathfrak b$ 包含于 $B$ 的极大理想的平方。

$c)\Rightarrow a)$：局部环 $A$ 非零，复合同态 $K\to B\to A$ 非零，故单射。
$a)\Rightarrow b)$：若 $K\subset A$ 为子域，其到 $k$ 的剩余同态单射，所以两域特征相同；故 $p\cdot1_A=0$。
$b)\Rightarrow a)$：设 $f:\mathbf Z\to A$ 和 $g:A\to k$ 为典范同态。$\Ker(gf)=p\mathbf Z$，而 $f(p\mathbf Z)=0$，故 $\Ker f=p\mathbf Z$。于是 $A$ 含有同构于 $\mathbf Z/p\mathbf Z$ 的子环 $A'$，且 $A'\cap\mathfrak m=0$。若 $p>0$，$A'$ 已是域；若 $p=0$，则 $A'$ 的每个\emph{非零}元素在 $A$ 中可逆，故 $A$ 还含有它的分式域 $\mathbf Q$。\cmt{G7-T06}{原文称 $A'$ 的“每个元素”可逆，须排除零元。前面的“不与 $\mathfrak m$ 相交”也应按 $A'\cap\mathfrak m=\{0\}$ 理解，不能把零元从两个子环中删除。}

最后证明 $a)\Rightarrow c)$。$A$ 含有与 $k$ 的素子域同构的素子域 $k_0$；$k/k_0$ 可分。由已校订的（19.6.2），存在同构 $f:k\xrightarrow{\sim}K\subset A$，且与剩余映射相容。取 $x_1,\ldots,x_n\in\mathfrak m$，使它们模 $\mathfrak m^2$ 的类为 $\mathfrak m/\mathfrak m^2$ 的 $k$-基。完备性给出唯一连续同态
\[
 u:k[[T_1,\ldots,T_n]]\longrightarrow A,
 \qquad u|_k=f,\quad u(T_i)=x_i.
\]
它由 $x_i$ 的选取而满射（Bourbaki，\textit{交换代数}，第三章，§2，第9目，命题11）。它在常数项和一次相伴分次部分均为同构，故其核包含于 $(T_1,\ldots,T_n)^2$，也得到最后的断言。

\egas{19.6.4}{定理}
设 $k$ 为域，$A$ 为诺特局部 $k$-代数，$\mathfrak m$ 为其极大理想，$K=A/\mathfrak m$，$A$ 取 $\mathfrak m$-预进制拓扑。假设 $K/k$ 可分。则以下条件等价：

$a)$ $A$ 为形式光滑 $k$-代数；

$b)$ $A$ 的完备化 $\widehat A$ 作为拓扑 $k$-代数，同构于 $K[[T_1,\ldots,T_n]]$，其 $k$-代数结构由
$k\to K\to K[[T_1,\ldots,T_n]]$ 给出；

$b')$ $\widehat A$ 作为环同构于某个 $K[[T_1,\ldots,T_n]]$；

$c)$ $A$ 为正则局部环。

$a)\Rightarrow b)$ 是（19.5.4）中 $a)$ 与 $c)$ 等价的特殊情形：将其中的基环换成 $k$，$B$ 换成 $A$，$\mathfrak j$ 换成 $\mathfrak m$，并用（19.6.1）。此时 $\mathfrak m/\mathfrak m^2$ 为有限维 $K$-向量空间，其对称代数是有限变量多项式环。
$b)\Rightarrow b')$ 是直接的；$b')\Rightarrow c)$ 由（17.1.1）得到。
若 $c)$ 成立，由（17.1.1，d），$\gr_{\mathfrak m}^{\bullet}(A)$ 为 $\mathfrak m/\mathfrak m^2$ 的对称代数；再将（19.5.4）中 $b)\Rightarrow a)$ 用于 $C=K$，就得到 $A$ 形式光滑。

\egas{19.6.5}{推论}
设 $k$ 为域，$A$ 为诺特局部 $k$-代数，$\mathfrak m$ 为其极大理想，$A$ 取 $\mathfrak m$-预进制拓扑。若 $A$ 为形式光滑 $k$-代数，则 $A$ 在 $k$ 上几何正则：即对每个有限扩张 $k'/k$，半局部环 $A'=A\otimes_kk'$ 为正则环（17.3.6；另见IV，6.7.6）。此外，若 $K$ 为 $A$ 的剩余域，则 $\widehat A$ 作为环同构于某个 $K[[T_1,\ldots,T_n]]$。

由于 $\widehat{A'}\simeq\widehat A\otimes_kk'$，由（19.3.6）及用于 $A'$ 各极大理想处局部环的（17.1.5），可假设 $A$ 完备。$A'$ 为形式光滑 $k'$-代数（19.3.5，iii），并且是若干完备局部 $k'$-代数的有限直积；每个直因子也形式光滑（19.3.5，v）。所以只须证明一个这样的形式光滑完备局部代数正则。

将其基域仍记为 $k$，令 $k_0$ 为 $k$ 的素子域。$k/k_0$ 可分，故为形式光滑扩张（19.6.1）；再用（19.3.5，ii），$A$ 为形式光滑 $k_0$-代数。其剩余域 $K$ 也在 $k_0$ 上可分，因此对 $A,k_0$ 应用（19.6.4），同时得到正则性和所述形式幂级数环结构。

\egas{19.6.5.1}{推论}
设 $A$ 为诺特局部环，$k\subset A$ 为子域，并假设 $A$ 对其预进制拓扑为形式光滑 $k$-代数。则每个满足 $k\subset K\subset A$ 的子域 $K$ 都是 $k$ 的可分扩张。

事实上，对每个有限扩张 $k'/k$，$K\otimes_kk'$ 单射到 $A'=A\otimes_kk'$ 中。由（19.6.5），$A'$ 正则，因而约化；其子环 $K\otimes_kk'$ 也约化。这就证明 $K/k$ 可分（Bourbaki，\textit{代数}，第八章，§7，第3目，定理1）。\cmt{G7-T07}{原证明末句把由约化性得到的对象写成 $K$；实际须得到 $K\otimes_kk'$ 约化。域 $K$ 本身约化并不能判定扩张可分。正文写出正确对象。}

特别地，对 $A=K$ 应用此推论，便证明了（19.6.1）中的必要性。

\egas{19.6.5.2}{注}
若 $K/k$ 不可分，则配备通常 $k$-代数结构的形式幂级数环 $K[[T_1,\ldots,T_n]]$（其结构同态如（19.6.4））不是形式光滑 $k$-代数。
另一方面，存在形式光滑的完备诺特局部 $k$-代数 $A$，其剩余域 $K$ 在 $k$ 上不可分。例如，取 $B=k[T_1,\ldots,T_n]$，在一个使 $B/\mathfrak n$ 为 $k$ 的有限不可分扩张的极大理想 $\mathfrak n$ 处局部化，然后完备化。这样的 $A$ 不能与通常结构的 $K[[T_1,\ldots,T_n]]$ 为 $k$-代数同构；但由（19.6.5），二者仍可以作为环同构。

某些情形下，可以将（19.6.5）加强如下。

\egas{19.6.6}{命题}
设 $k$ 为域，$A$ 为诺特局部 $k$-代数，$\mathfrak m$ 为其极大理想，并给 $A$ 赋予 $\mathfrak m$-预进制拓扑。假设剩余域 $K=A/\mathfrak m$ 满足：存在有限根式扩张（即有限纯不可分扩张）$k_0/k$，使局部环 $K\otimes_kk_0$ 的剩余域为 $k_0$ 的可分扩张。以后把这一条件称为 $K$ 在 $k$ 上具有有限根式重数（IV，4.7.4）；还将证明，$K/k$ 有限生成时总满足此条件。则以下条件等价：

$a)$ $A$ 为形式光滑 $k$-代数；

$b)$ 存在有限根式扩张 $k'/k$，使 $\widehat A\otimes_kk'$ 作为 $k'$-代数，同构于某个 $K'[[T_1,\ldots,T_n]]$，其中 $K'/k'$ 可分；

$b')$ 存在有限扩张 $k'/k$，使完备半局部环 $\widehat A\otimes_kk'$ 的至少一个局部分量，作为 $k'$-代数同构于这样的 $K'[[T_1,\ldots,T_n]]$，其中 $K'/k'$ 可分；

$c)$ $A$ 在 $k$ 上几何正则（19.6.5）。

先作几点说明。若 $k'/k$ 为有限根式扩张，在 $A'=A\otimes_kk'$ 中只有一个素理想位于 $\mathfrak m$ 之上：它由某个 $p^h$ 次幂属于 $\mathfrak m A'$ 的元素构成，其中 $p$ 为 $k$ 的特征指数（Bourbaki，\textit{交换代数}，第五章，§2，第3目，引理4）。因此 $A'$ 为局部环，$K\otimes_kk'=A'/\mathfrak m A'$ 也为局部环，二者剩余域相同。

若 $K/k$ 可分，则对每个有限扩张 $k''/k$，$K\otimes_kk''$ 为有限个域的直积（Bourbaki，\textit{代数}，第八章，§7，第3目，定理1的推论1）；于是 $\mathfrak m\otimes_kk''$ 是 $A\otimes_kk''$ 的根，而上述各域就是其极大理想处的剩余域。它们都在 $k''$ 上可分，因为对任意扩张 $k'''/k''$，
\[
 (K\otimes_kk'')\otimes_{k''}k'''=K\otimes_kk'''
\]
约化（同处定理1）。将此用于有限根式重数条件，可知对每个有限扩张 $k''/k_0$，$(K\otimes_kk'')_{\mathrm{red}}$ 在 $k''$ 上可分。

现在证明各条件的等价。$b)\Rightarrow b')$ 是直接的。下面证明 $b')\Rightarrow a)$，并明确补足其中的下降步骤。\cmt{G7-G02}{原证明称扩域后的各局部分量与原来的 $\widehat A\otimes_k k'_0$ 同构，这一般错误；例如 $A=k=\mathbf R$、$k'=\mathbf C$ 时，$\mathbf C$ 不与 $\mathbf R$ 为 $\mathbf R$-代数同构。正确结论是各分量彼此由Galois作用联系。正文用这一结论以及（19.4.6）的忠实平坦下降给出修补；初版在此引（19.4.7），与所需判据不符。详细验证见报告G7-G02。}

将给定 $k'$ 放入有限准Galois扩张 $L/k$，并令 $l$ 为 $L$ 中 $k$ 的纯不可分闭包，使 $l/k$ 为有限根式扩张、$L/l$ 为有限Galois扩张。给定的良好局部分量经 $k'\to L$ 扩域后，是有限个形式幂级数环的直积：这是因为 $K'/k'$ 可分，$K'\otimes_{k'}L$ 是在 $L$ 上可分的域的有限直积，且有限扩域与形式幂级数环可交换。因此
\[
 B_0=\widehat A\otimes_kl,\qquad D=B_0\otimes_lL
\]
中的 $D$ 至少有一个这样的局部分量，而 $B_0$ 为完备局部环。

群 $\operatorname{Gal}(L/l)$ 作用在 $D$ 上，并传递地置换 $D$ 的局部分量。确实，$D$ 为完备半局部环，若分量分成不止一个轨道，某个真轨道对应的单位幂等元之和就给出 $D^{\operatorname{Gal}(L/l)}=B_0$ 的非平凡幂等元，矛盾于 $B_0$ 局部。各分量因而作为 $B_0$-代数彼此同构；这些同构对 $L$ 的作用可以是半线性的，但仍保持环的正则性以及剩余域在 $L$ 上的可分性。由（19.6.4），每个分量都为形式光滑的 $L$-代数，故其有限直积 $D$ 也如此（19.3.5，v）。

拓扑 $k$-向量空间 $\widehat A$ 形式投射：一个到离散向量空间的连续线性映射通过某个开子空间商分解，而该商作为 $k$-向量空间投射，因此可连续提升。又因 $L/k$ 为有限自由且忠实平坦的扩张，可应用（19.4.6）的第二种条件，将 $D=\widehat A\otimes_kL$ 的形式光滑性下降到 $\widehat A$。最后由（19.3.6），$A$ 形式光滑。

$a)\Rightarrow c)$ 已由（19.6.5）证明。若 $c)$ 成立，取假设所给的 $k_0$，或包含它的任意有限根式扩张 $k'/k$；则 $A\otimes_kk'$ 为正则局部环，且其剩余域 $K'$ 在 $k'$ 上可分。由（19.6.4），它的完备化
\[
 \widehat{A\otimes_kk'}=\widehat A\otimes_kk'
\]
作为 $k'$-代数同构于 $K'[[T_1,\ldots,T_n]]$，即得 $b)$。
\finishproof

\egas{19.6.7}{注}
（i）在上述证明中，剩余域 $K$ 具有有限根式重数的假设用于 $c)\Rightarrow b)$：它保证能通过有限根式扩域使剩余域可分。修补后的 $b')\Rightarrow a)$ 不需要这一假设。\cmt{G7-T08}{初版此处写成“该假设只用于 $b')\Rightarrow a)$”。这与前一段 $c)\Rightarrow b)$ 中选取 $k_0$ 的实际依赖不符；正文按依赖关系改正，不将该备注误作无条件存在有限根式扩域的保证。}

（ii）以后将在（22.5.8）证明：无须对扩张 $K/k$ 增加条件，$a)$ 与 $c)$ 仍等价。

% END continuations/iv0_19_6.tex

\NewNote{本章本版连续译至§19的（19.6.7）末尾。19.7--19.10、第0章§§20--23尚未译入；下一接续点为19.7“局部同态的情形：存在性与唯一性定理”的章首及（19.7.1）。下接第I章不表示第0章余文已译完。}

% BEGIN text/ch1.tex
\renewcommand{\B}{\mathfrak B}
\setcounter{chapter}{0}
\chapter{概形的语言}
\src{79}
\begin{center}\textbf{本章提要}\end{center}
\begin{enumerate}[label=\S\,\arabic*.,itemsep=0.1em]
\item 仿射概形。
\item 预概形与预概形态射。
\item 预概形的积。
\item 子预概形与浸入态射。
\item 约化预概形；分离条件。
\item 有限性条件。
\item 有理映射。
\item Chevalley 概形。
\item 拟凝聚层的补充知识。
\item 形式概形。
\end{enumerate}

第 1—8 节主要是展开一种将在后文始终使用的语言。不过应当指出，依照本书的总体精神，
第 7、8 节的使用将少于其他各节，其作用也较不根本；事实上，讨论 Chevalley 概形，
只是为了与 Chevalley\cite{b1} 和 Nagata\cite{b9} 的语言建立联系。
第 9 节给出拟凝聚层的一些定义与结果，其中某些内容已不再只是把交换代数中已知的概念
翻译成“几何”语言，而具有整体性质；从以后的各章起，它们将成为态射整体研究中不可缺少的工具。
最后，第 10 节引入概形概念的一种推广；在第 III 章中，我们将借助它来方便地表述并证明
固有态射的上同调研究中的基本结果。另外还应指出，形式概形的概念似乎是表述
“模空间理论”（代数簇的分类问题）中某些事实所不可缺少的。\transnote{
此处原文为 \textit{théorie des modules}，作者紧接着用“代数簇的分类问题”作解释。
这里的 \textit{modules} 指分类参数意义上的模问题，而不是环上模的代数理论。}
第 10 节的结果在第 III 章 §3 之前不会使用，建议读者在读到那里以前暂且跳过这一节。
\par


% END chapters/ch1_intro.tex

\par

% BEGIN chapters/ch1_s1_1.tex
\section{仿射概形}\hypertarget{heading:1:1}{}
\subsection{环的素谱}\hypertarget{heading:1:1.1}{}
\src{80—83}

\egapar{1.1.1}
\textbf{记号。}设 $A$ 为一个（交换）环，$M$ 为一个 $A$-模。在本章及以后各章中，
我们将经常使用下列记号：

$\Spec(A)$ 是 $A$ 的素理想的集合，也称为 $A$ 的\emph{素谱}；
对于 $x\in X=\Spec(A)$，常常用 $\jj_x$ 代替 $x$ 更为方便。
$\Spec(A)$ 为空，当且仅当环 $A$ 为零环。

$A_x=A_{\jj_x}$ 是分母取自 $S=A\setminus\jj_x$ 的（局部）分式环 $S^{-1}A$。

$\mm_x=\jj_xA_{\jj_x}$ 是 $A_x$ 的极大理想。

$k(x)=A_x/\mm_x$ 是 $A_x$ 的剩余域，它与整环 $A/\jj_x$ 的分式域典范同构，
我们将两者等同。

对于 $f\in A$ 及 $x\in X$，$f(x)$ 是 $f$ 模 $\jj_x$ 的类，
视为 $A/\jj_x\subset k(x)$ 的元素。也称 $f(x)$ 为 $f$ 在点 $x\in\Spec(A)$ 的\emph{值}；
关系 $f(x)=0$ 与 $f\in\jj_x$ 等价。

$M_x=M\otimes_A A_x$ 是分母取自 $A\setminus\jj_x$ 的分式模。

$\rr(E)$ 是由 $A$ 的子集 $E$ 生成的理想的根。

对于 $E\subset A$，$V(E)$ 是满足 $E\subset\jj_x$ 的点 $x\in X$ 的集合，
也就是使每个 $f\in E$ 都满足 $f(x)=0$ 的点的集合。因此有
\begin{equation}
\rr(E)=\bigcap_{x\in V(E)}\jj_x.
\tag{1.1.1.1}\label{eq:1:1.1.1.1}
\end{equation}
对于 $f\in A$，记 $V(f)=V(\{f\})$。

$D(f)=X\setminus V(f)$ 是满足 $f(x)\ne0$ 的点 $x\in X$ 的集合。

\egastatement{1.1.2}{命题}
有下列性质：
\begin{enumerate}[label=(\roman*)]
\item $V(0)=X$，$V(1)=\varnothing$。
\item $E\subset E'$ 蕴含 $V(E)\supset V(E')$。
\item 对于 $A$ 的任意子集族 $(E_\lambda)$，有
\[
V\Bigl(\bigcup_\lambda E_\lambda\Bigr)
=V\Bigl(\sum_\lambda E_\lambda\Bigr)
=\bigcap_\lambda V(E_\lambda).
\]
\item $V(EE')=V(E)\cup V(E')$。
\item $V(E)=V(\rr(E))$。
\end{enumerate}

性质 (i)、(ii)、(iii) 是显然的，(v) 由 (ii) 及公式（\hyperref[eq:1:1.1.1.1]{1.1.1.1}）得到。
显然 $V(EE')\supset V(E)\cup V(E')$；反之，若 $x\notin V(E)$ 且 $x\notin V(E')$，
则存在 $f\in E$ 和 $f'\in E'$，使 $f(x)$ 和 $f'(x)$ 在 $k(x)$ 中均非零。
于是 $f(x)f'(x)\ne0$，即 $x\notin V(EE')$，这就证明了 (iv)。

命题（\hyperlink{para:1:1.1.2}{1.1.2}）特别表明，当 $E$ 遍历 $A$ 的所有子集时，形如 $V(E)$ 的集合
恰为 $X$ 上某个拓扑的闭集。我们称此拓扑为\emph{谱拓扑}。\origfootnote{
将这一拓扑引入代数几何的是 Zariski。因此它也常被称为 $X$ 的“Zariski 拓扑”。}
除非另作明确说明，始终假定 $X=\Spec(A)$ 赋予谱拓扑。

\egapar{1.1.3}
对于 $X$ 的任意子集 $Y$，以 $\jj(Y)$ 表示满足 $f(y)=0$（对所有 $y\in Y$）的
元素 $f\in A$ 所成的集合；等价地，$\jj(Y)$ 是 $y\in Y$ 对应的素理想 $\jj_y$ 的交。
显然，$Y\subset Y'$ 蕴含 $\jj(Y)\supset\jj(Y')$，且对于 $X$ 的任意子集族 $(Y_\lambda)$，有
\begin{equation}
\jj\Bigl(\bigcup_\lambda Y_\lambda\Bigr)
=\bigcap_\lambda\jj(Y_\lambda).
\tag{1.1.3.1}\label{eq:1:1.1.3.1}
\end{equation}
最后，有
\begin{equation}
\jj(\{x\})=\jj_x.
\tag{1.1.3.2}\label{eq:1:1.1.3.2}
\end{equation}

\egastatement{1.1.4}{命题}
(i) 对于 $A$ 的任意子集 $E$，有 $\jj(V(E))=\rr(E)$。

(ii) 对于 $X$ 的任意子集 $Y$，有 $V(\jj(Y))=\overline Y$，其中 $\overline Y$ 是 $Y$ 在 $X$ 中的闭包。

(i) 是定义和（\hyperref[eq:1:1.1.1.1]{1.1.1.1}）的直接推论。另一方面，$V(\jj(Y))$ 是包含 $Y$ 的闭集；
反之，若 $Y\subset V(E)$，则对于每个 $f\in E$ 及每个 $y\in Y$ 都有 $f(y)=0$，
因而 $E\subset\jj(Y)$，$V(E)\supset V(\jj(Y))$，这就证明了 (ii)。

\egastatement{1.1.5}{推论}
$X=\Spec(A)$ 的闭子集与 $A$ 中等于自身的根的理想（即素理想的交）之间，
通过两个反序映射 $Y\mapsto\jj(Y)$、$\mathfrak a\mapsto V(\mathfrak a)$ 建立一一对应。
两个闭子集之并 $Y_1\cup Y_2$ 对应于 $\jj(Y_1)\cap\jj(Y_2)$；
任意闭子集族 $(Y_\lambda)$ 的交对应于各 $\jj(Y_\lambda)$ 之和的根。

\egastatement{1.1.6}{推论}
若 $A$ 是诺特环，则 $X=\Spec(A)$ 是诺特空间。

应当注意，这一推论的逆命题不成立。例如，一个非诺特整环可以恰好只有一个非零素理想，
秩为 $1$ 的非离散赋值环就是这样的例子。

作为素谱不是诺特空间的环 $A$ 的例子，可以取无限紧空间 $Y$ 上的数值连续函数环
$\mathscr C(Y)$。我们知道，作为集合，$Y$ 可与 $A$ 的极大理想集合等同；
而且容易看出，$X=\Spec(A)$ 的拓扑在 $Y$ 上诱导的拓扑就是 $Y$ 原有的拓扑。
由于 $Y$ 不是诺特空间，$X$ 也不是诺特空间。

\egastatement{1.1.7}{推论}
对于每个 $x\in X$，$\{x\}$ 的闭包是满足 $\jj_x\subset\jj_y$ 的点 $y\in X$ 的集合。
$\{x\}$ 是闭集，当且仅当 $\jj_x$ 是极大理想。

\egastatement{1.1.8}{推论}
空间 $X=\Spec(A)$ 是 Kolmogoroff 空间。

事实上，若 $x,y$ 是 $X$ 的两个不同点，则或者 $\jj_x\not\subset\jj_y$，
或者 $\jj_y\not\subset\jj_x$，因而其中一个点不属于另一个点的闭包。

\egapar{1.1.9}
由命题（\hyperlink{para:1:1.1.2}{1.1.2}，(iv)），对于 $A$ 的两个元素 $f,g$，有
\begin{equation}
D(fg)=D(f)\cap D(g).
\tag{1.1.9.1}\label{eq:1:1.1.9.1}
\end{equation}
还应指出，由命题（\hyperlink{para:1:1.1.4}{1.1.4}，(i)）及命题（\hyperlink{para:1:1.1.2}{1.1.2}，(v)），
关系 $D(f)=D(g)$ 表示 $\rr(f)=\rr(g)$，也等价于包含 $(f)$ 与包含 $(g)$ 的极小素理想相同。
特别地，当 $f=ug$ 且 $u$ 可逆时，此关系成立。

\egastatement{1.1.10}{命题}
(i) 当 $f$ 遍历 $A$ 时，集合 $D(f)$ 构成 $X$ 的拓扑的一组基。

(ii) 对于每个 $f\in A$，$D(f)$ 都是拟紧的。特别地，$X=D(1)$ 是拟紧的。

(i) 设 $U$ 是 $X$ 的开集。依定义，$U=X\setminus V(E)$，其中 $E$ 是 $A$ 的某个子集。
而 $V(E)=\bigcap_{f\in E}V(f)$，故 $U=\bigcup_{f\in E}D(f)$。

(ii) 由 (i)，只需证明：若 $A$ 的元素族 $(f_\lambda)_{\lambda\in L}$ 满足
$D(f)\subset\bigcup_{\lambda\in L}D(f_\lambda)$，
则存在 $L$ 的有限子集 $J$，使 $D(f)\subset\bigcup_{\lambda\in J}D(f_\lambda)$。
设 $\mathfrak a$ 是由各 $f_\lambda$ 生成的理想；由假设有
$V(f)\supset V(\mathfrak a)$，从而 $\rr(f)\subset\rr(\mathfrak a)$。
由于 $f\in\rr(f)$，存在整数 $n>0$ 使 $f^n\in\mathfrak a$。
于是 $f^n$ 已属于某个有限子族 $(f_\lambda)_{\lambda\in J}$ 生成的理想 $\mathfrak b$，且
\[
V(f)=V(f^n)\supset V(\mathfrak b)
=\bigcap_{\lambda\in J}V(f_\lambda),
\]
即 $D(f)\subset\bigcup_{\lambda\in J}D(f_\lambda)$。

\egastatement{1.1.11}{命题}
对于 $A$ 的每个理想 $\mathfrak a$，$\Spec(A/\mathfrak a)$ 典范地等同于
$\Spec(A)$ 的闭子空间 $V(\mathfrak a)$。

事实上，我们知道，$A/\mathfrak a$ 的理想（分别地，素理想）与 $A$ 中包含 $\mathfrak a$ 的理想
（分别地，素理想）之间，存在保持包含序的典范一一对应。

回忆一下，$A$ 的幂零元集合 $\nrad$（即 $A$ 的幂零根）是等于 $\rr(0)$ 的理想，
也就是 $A$ 的所有素理想之交（0，\hyperlink{para:0:1.1.1}{1.1.1}）。

\egastatement{1.1.12}{推论}
拓扑空间 $\Spec(A)$ 与 $\Spec(A/\nrad)$ 典范同构。

\egastatement{1.1.13}{命题}
$X=\Spec(A)$ 不可约（0，\hyperlink{para:0:2.1.1}{2.1.1}），当且仅当环 $A/\nrad$ 是整环；
等价地，当且仅当理想 $\nrad$ 是素理想。

由推论（\hyperlink{para:1:1.1.12}{1.1.12}），可以只考虑 $\nrad=0$ 的情形。
若 $X$ 可约，则存在两个均不等于 $X$ 的闭子集 $Y_1,Y_2$，使 $X=Y_1\cup Y_2$，从而
$\jj(X)=\jj(Y_1)\cap\jj(Y_2)=0$，且 $\jj(Y_1)$ 与 $\jj(Y_2)$ 均不等于 $(0)$（\hyperlink{para:1:1.1.5}{1.1.5}）；
因此 $A$ 不是整环。反之，若 $A$ 中有两个非零元素 $f,g$，使 $fg=0$，则
$V(f)\ne X$，$V(g)\ne X$（因为 $A$ 的所有素理想之交为 $(0)$），而
$X=V(fg)=V(f)\cup V(g)$。

\egastatement{1.1.14}{推论}
(i) 在 $X=\Spec(A)$ 的闭子集与 $A$ 的根理想之间的一一对应中，
$X$ 的不可约闭子集对应于 $A$ 的素理想。特别地，$X$ 的不可约分支对应于 $A$ 的极小素理想。

(ii) 映射 $x\mapsto\overline{\{x\}}$ 在 $X$ 与 $X$ 的不可约闭子集的集合之间建立一一对应；
换言之，$X$ 的每个不可约闭子集恰有一个泛点。

(i) 直接由（\hyperlink{para:1:1.1.13}{1.1.13}）及（\hyperlink{para:1:1.1.11}{1.1.11}）得到。为证明 (ii)，由（\hyperlink{para:1:1.1.11}{1.1.11}），可以只考虑 $X$ 不可约的情形。
这时，由命题（\hyperlink{para:1:1.1.13}{1.1.13}），$A$ 中存在一个最小素理想 $\nrad$，它对应于 $X$ 的一个泛点。
此外，由于 $X$ 是 Kolmogoroff 空间，其泛点唯一（（\hyperlink{para:1:1.1.8}{1.1.8}）及（0，\hyperlink{para:0:2.1.3}{2.1.3}））。

\egastatement{1.1.15}{命题}
若 $\frJ$ 是 $A$ 的一个包含于根 $\frR(A)$ 中的理想，
则 $V(\frJ)$ 在 $X=\Spec(A)$ 中的唯一邻域是整个空间 $X$。

事实上，依定义，$A$ 的每个极大理想都属于 $V(\frJ)$。
又因为 $A$ 的每个真理想 $\mathfrak a$ 都包含于一个极大理想，
故 $V(\mathfrak a)\cap V(\frJ)\ne\varnothing$，从而得到本命题。\transnote{
原文此处写作“每个理想”；这里只涉及非空闭集 $V(\mathfrak a)$，
因而 $\mathfrak a\ne A$。译文将这一限制明确写出。
$\frR(A)$ 是第 0 章（\hyperlink{para:1:1.1.2}{1.1.2}）定义的根，不是幂零根 $\nrad$。}
\par


% END chapters/ch1_s1_1.tex

\par

% BEGIN chapters/ch1_s1_2.tex
\subsection{环的素谱的函子性质}\hypertarget{heading:1:1.2}{}
\src{83—84}

\egapar{1.2.1}
设 $A,A'$ 是两个环，$\varphi\colon A'\to A$ 是环同态。
对于每个素理想 $x=\jj_x\in\Spec(A)=X$，环 $A'/\varphi^{-1}(\jj_x)$
与 $A/\jj_x$ 的一个子环典范同构，因而是整环；也就是说，
$\varphi^{-1}(\jj_x)$ 是 $A'$ 的素理想。以 $\aphi(x)$ 表示这个素理想，便定义了映射
\[
\aphi\colon X=\Spec(A)\longrightarrow X'=\Spec(A'),
\]
也记作 $\Spec(\varphi)$，称为\emph{与同态 $\varphi$ 相伴的映射}。
以 $\varphi^x$ 表示由 $\varphi$ 取商得到的单同态
$A'/\varphi^{-1}(\jj_x)\to A/\jj_x$，也以同一记号表示它到分式域的典范延拓
\[
\varphi^x\colon k(\aphi(x))\longrightarrow k(x).
\]
因此，依定义，对于每个 $f'\in A'$，有
\begin{equation}
\varphi^x\bigl(f'(\aphi(x))\bigr)=\bigl(\varphi(f')\bigr)(x)
\qquad(x\in X).
\tag{1.2.1.1}\label{eq:1:1.2.1.1}
\end{equation}

\egastatement{1.2.2}{命题}
(i) 对于 $A'$ 的每个子集 $E'$，有
\begin{equation}
(\aphi)^{-1}(V(E'))=V(\varphi(E'));
\tag{1.2.2.1}\label{eq:1:1.2.2.1}
\end{equation}
特别地，对于每个 $f'\in A'$，有
\begin{equation}
(\aphi)^{-1}(D(f'))=D(\varphi(f')).
\tag{1.2.2.2}\label{eq:1:1.2.2.2}
\end{equation}
(ii) 对于 $A$ 的每个理想 $\mathfrak a$，有
\begin{equation}
\overline{\aphi(V(\mathfrak a))}=V(\varphi^{-1}(\mathfrak a)).
\tag{1.2.2.3}\label{eq:1:1.2.2.3}
\end{equation}

事实上，关系 $\aphi(x)\in V(E')$ 依定义等价于
$E'\subset\varphi^{-1}(\jj_x)$，也就等价于 $\varphi(E')\subset\jj_x$，
最终等价于 $x\in V(\varphi(E'))$，从而得到 (i)。
为证明 (ii)，可以假定 $\mathfrak a$ 等于自身的根，因为
$V(\rr(\mathfrak a))=V(\mathfrak a)$（\hyperlink{para:1:1.1.2}{1.1.2}，(v)），且
$\varphi^{-1}(\rr(\mathfrak a))=\rr(\varphi^{-1}(\mathfrak a))$。
令 $Y=V(\mathfrak a)$，$\mathfrak a'=\jj(\aphi(Y))$，则
$\overline{\aphi(Y)}=V(\mathfrak a')$（命题（\hyperlink{para:1:1.1.4}{1.1.4}，(ii)））。
关系 $f'\in\mathfrak a'$ 依定义等价于：对于每个 $x'\in\aphi(Y)$，有 $f'(x')=0$。
由公式（\hyperref[eq:1:1.2.1.1]{1.2.1.1}），这又等价于：对于每个 $x\in Y$，有 $\varphi(f')(x)=0$；
也等价于 $\varphi(f')\in\jj(Y)=\mathfrak a$，这里用了 $\mathfrak a$ 为根理想。
由此得到 (ii)。

\egastatement{1.2.3}{推论}
映射 $\aphi$ 连续。

注意，若 $A''$ 是第三个环，$\varphi'\colon A''\to A'$ 是同态，则有
\[
{}^a(\varphi\circ\varphi')=({}^a\varphi')\circ({}^a\varphi).
\]
\noindent 此结果与推论（\hyperlink{para:1:1.2.3}{1.2.3}）说明：$\Spec(A)$ 是关于 $A$ 的反变函子，
从环范畴取值于拓扑空间范畴。\transnote{
原版在此印作 ${}^a(\varphi'\circ\varphi)=({}^a\varphi)\circ({}^a\varphi')$。
按本段给定的 $\varphi\colon A'\to A$、$\varphi'\colon A''\to A'$，
两边的复合次序均应交换；译文采用与这些源、靶一致的式子。
同一关系在原文（\hyperlink{para:1:1.6.10}{1.6.10}）中按正确次序写出。}

\egastatement{1.2.4}{推论}
假设 $\varphi$ 满足：每个 $f\in A$ 都可写成 $f=h\varphi(f')$，其中 $h$ 在 $A$ 中可逆
（特别地，$\varphi$ 满射时满足此条件）。则 $\aphi$ 是 $X$ 到 $\aphi(X)$ 的同胚。

先证明，对于每个 $E\subset A$，存在 $A'$ 的子集 $E'$，使 $V(E)=V(\varphi(E'))$。
由公理 $(\mathrm T_0)$（\hyperlink{para:1:1.1.8}{1.1.8}）及公式（\hyperref[eq:1:1.2.2.1]{1.2.2.1}），这首先蕴含 $\aphi$ 单射；
再由同一公式可知 $\aphi$ 是同胚。
而对于每个 $f\in E$，只需取一个 $f'\in A'$，使 $h\varphi(f')=f$，其中 $h$ 在 $A$ 中可逆；
这些 $f'$ 所成的集合 $E'$ 就满足要求。

\egapar{1.2.5}
特别地，若 $\varphi$ 是由 $A$ 到商环 $A/\mathfrak a$ 的典范同态，
便重新得到（\hyperlink{para:1:1.1.12}{1.1.12}），\transnote{
此处引文编号照录原版。一般理想 $\mathfrak a$ 的相应结论在（\hyperlink{para:1:1.1.11}{1.1.11}）；
（\hyperlink{para:1:1.1.12}{1.1.12}）是 $\mathfrak a$ 为幂零根时的特例。}
而 $\aphi$ 就是把与 $\Spec(A/\mathfrak a)$ 等同的 $V(\mathfrak a)$
嵌入 $X=\Spec(A)$ 的典范包含映射。

（\hyperlink{para:1:1.2.4}{1.2.4}）的另一个特例给出：

\egastatement{1.2.6}{推论}
若 $S$ 是 $A$ 的一个乘法子集，则素谱 $\Spec(S^{-1}A)$ 连同其拓扑，
典范地等同于 $X=\Spec(A)$ 中由满足 $\jj_x\cap S=\varnothing$ 的点 $x$ 所成的子空间。

事实上，我们知道（0，\hyperlink{para:0:1.2.6}{1.2.6}），$S^{-1}A$ 的素理想恰是满足
$\jj_x\cap S=\varnothing$ 的理想 $S^{-1}\jj_x$，且有
$\jj_x=(i_A^S)^{-1}(S^{-1}\jj_x)$。
因此，只需将推论（\hyperlink{para:1:1.2.4}{1.2.4}）应用于 $i_A^S$。

\egastatement{1.2.7}{推论}
$\aphi(X)$ 在 $X'$ 中稠密，当且仅当核 $\Ker\varphi$ 的每个元素均为幂零元。

事实上，将公式（\hyperref[eq:1:1.2.2.3]{1.2.2.3}）应用于理想 $\mathfrak a=(0)$，得到
$\overline{\aphi(X)}=V(\Ker\varphi)$。
而 $V(\Ker\varphi)=X'$ 当且仅当 $\Ker\varphi$ 包含于 $A'$ 的每个素理想，
即包含于 $A'$ 的幂零根 $\nrad$。\transnote{
原版这一证明的末句将 $X'$、$A'$ 印作 $X$、$A$。
由于 $\Ker\varphi$ 是 $\varphi\colon A'\to A$ 的定义域 $A'$ 中的理想，
这里补出撇号，以与命题及（\hyperref[eq:1:1.2.2.3]{1.2.2.3}）一致。}
\par


% END chapters/ch1_s1_2.tex

\par

% BEGIN chapters/ch1_s1_3.tex
\subsection{与模相伴的层}\hypertarget{heading:1:1.3}{}
\src{84—90}

\egapar{1.3.1}
设 $A$ 是交换环，$M$ 是 $A$-模，$f\in A$，$S_f$ 是由各 $f^n$（$n\geq0$）组成的乘法子集。
回忆一下，我们记 $A_f=S_f^{-1}A$，$M_f=S_f^{-1}M$。
若 $S'_f$ 是由能够整除 $S_f$ 的某个元素的所有 $g\in A$ 所成的饱和乘法子集，
则我们知道，$A_f$ 与 $M_f$ 分别典范地等同于 $(S'_f)^{-1}A$ 与 $(S'_f)^{-1}M$（0，\hyperlink{para:0:1.4.3}{1.4.3}）。

\egastatement{1.3.2}{引理}
下列条件等价：
\begin{enumerate}[label=\alph*)]
\item $g\in S'_f$；
\item $S'_g\subset S'_f$；
\item $f\in\rr(g)$；
\item $\rr(f)\subset\rr(g)$；
\item $V(g)\subset V(f)$；
\item $D(f)\subset D(g)$。
\end{enumerate}
这直接由定义及（\hyperlink{para:1:1.1.5}{1.1.5}）得到。

\egapar{1.3.3}
若 $D(f)=D(g)$，则由引理（\hyperlink{para:1:1.3.2}{1.3.2}，b)）可知 $M_f=M_g$。
更一般地，若 $D(f)\supset D(g)$，即 $S'_f\subset S'_g$，
则由（0，\hyperlink{para:0:1.4.1}{1.4.1}）知存在典范的函子性同态
\[
\rho_{g,f}\colon M_f\longrightarrow M_g.
\]
若 $D(f)\supset D(g)\supset D(h)$，则有（0，\hyperlink{para:0:1.4.4}{1.4.4}）
\begin{equation}
\rho_{h,g}\circ\rho_{g,f}=\rho_{h,f}.
\tag{1.3.3.1}\label{eq:1:1.3.3.1}
\end{equation}

对于给定的 $x\in X=\Spec(A)$，当 $f$ 遍历 $A\setminus\jj_x$ 时，
各 $S'_f$ 构成 $A\setminus\jj_x$ 的一个递增滤过子集族：
因为对 $A\setminus\jj_x$ 的任意两个元素 $f,g$，$S'_f$ 与 $S'_g$ 都包含于 $S'_{fg}$。
又因为这些 $S'_f$ 的并为 $A\setminus\jj_x$，由（0，\hyperlink{para:0:1.4.5}{1.4.5}）可知，
$A_x$-模 $M_x$ 与相对于同态族 $(\rho_{g,f})$ 的归纳极限 $\ilim M_f$ 典范地等同。
以
\[
\rho_{x,f}\colon M_f\longrightarrow M_x
\]
表示 $f\in A\setminus\jj_x$（等价地，$x\in D(f)$）时的典范同态。

\egastatement{1.3.4}{定义}
在 $X=\Spec(A)$ 的由各 $D(f)$（$f\in A$）组成的开集基 $\B$ 上，
考虑预层 $D(f)\mapsto A_f$（分别地，$D(f)\mapsto M_f$）。
与此预层相伴的环层（分别地，$\wt A$-模层）称为素谱 $X$ 的\emph{结构层}
（分别地，\emph{与 $A$-模 $M$ 相伴的层}），
记作 $\wt A$ 或 $\OO_X$（分别地，$\wt M$）
（（\hyperlink{para:1:1.1.10}{1.1.10}），（0，\hyperlink{para:0:3.2.1}{3.2.1} 及 \hyperlink{para:0:3.5.6}{3.5.6}））。

由（0，\hyperlink{para:0:3.2.4}{3.2.4}）可知，茎 $\wt A_x$（分别地，$\wt M_x$）
与环 $A_x$（分别地，$A_x$-模 $M_x$）等同。以
\[
\theta_f\colon A_f\longrightarrow\Gamma(D(f),\wt A),
\qquad
\theta_f\colon M_f\longrightarrow\Gamma(D(f),\wt M)
\]
分别表示典范映射。因而，对于每个 $x\in D(f)$ 及每个 $z\in M_f$，有
\begin{equation}
\bigl(\theta_f(z)\bigr)_x=\rho_{x,f}(z).
\tag{1.3.4.1}\label{eq:1:1.3.4.1}
\end{equation}

\egastatement{1.3.5}{命题}
$\wt M$ 是关于 $M$ 的正合协变函子，从 $A$-模范畴取值于 $\wt A$-模层范畴。

事实上，设 $M,N$ 是两个 $A$-模，$u\colon M\to N$ 是同态。
对于每个 $f\in A$，$u$ 典范地给出一个 $A_f$-模同态 $u_f\colon M_f\to N_f$，
且当 $D(g)\subset D(f)$ 时，下图交换（0，\hyperlink{para:0:1.4.1}{1.4.1}）：
\[
\begin{tikzcd}[column sep=large,row sep=large]
M_f\arrow[r,"u_f"]\arrow[d,"\rho_{g,f}"'] & N_f\arrow[d,"\rho_{g,f}"]\\
M_g\arrow[r,"u_g"'] & N_g.
\end{tikzcd}
\]
这些同态因而定义一个 $\wt A$-模层同态 $\wt u\colon\wt M\to\wt N$（0，\hyperlink{para:0:3.2.3}{3.2.3}）。
此外，对于每个 $x\in X$，$\wt u_x$ 是各 $u_f$ 沿满足 $x\in D(f)$ 的 $f\in A$ 的归纳极限。
因而，由（0，\hyperlink{para:0:1.4.5}{1.4.5}），当将 $\wt M_x$、$\wt N_x$ 分别典范地等同于 $M_x$、$N_x$ 时，
$\wt u_x$ 就与由 $u$ 典范得到的同态 $u_x$ 等同。
若 $P$ 是第三个 $A$-模，$v\colon N\to P$ 为同态，且 $w=v\circ u$，
则显然 $\wt w_x=\wt v_x\circ\wt u_x$，故 $\wt w=\wt v\circ\wt u$。
因此，确实定义了所述协变函子。
这一函子正合，因为对于每个 $x\in X$，$M_x$ 是关于 $M$ 的正合函子（0，\hyperlink{para:0:1.3.2}{1.3.2}）。
另外，由两边的定义（0，\hyperlink{para:0:1.7.1}{1.7.1} 及 \hyperlink{para:0:3.1.6}{3.1.6}），有 $\Supp(\wt M)=\Supp(M)$。

\egastatement{1.3.6}{命题}
对于每个 $f\in A$，开集 $D(f)\subset X$ 典范地等同于素谱 $\Spec(A_f)$；
与 $A_f$-模 $M_f$ 相伴的层 $\wt{M_f}$ 典范地等同于限制 $\wt M|D(f)$。

第一项断言是（\hyperlink{para:1:1.2.6}{1.2.6}）的特例。又若 $g\in A$ 满足 $D(g)\subset D(f)$，
则 $M_g$ 典范地等同于 $M_f$ 的分式模，其分母是 $g$ 在 $A_f$ 中的典范像的各次幂（0，\hyperlink{para:0:1.4.6}{1.4.6}）。
由定义即得 $\wt{M_f}$ 与 $\wt M|D(f)$ 的典范等同。

\egastatement{1.3.7}{定理}
对于每个 $A$-模 $M$ 及每个 $f\in A$，同态
\[
\theta_f\colon M_f\longrightarrow\Gamma(D(f),\wt M)
\]
是双射；换言之，预层 $D(f)\mapsto M_f$ 是一个层。
特别地，$M$ 通过 $\theta_1$ 与 $\Gamma(X,\wt M)$ 等同。

注意，当 $M=A$ 时，$\theta_f$ 是环结构的同态。
因此，定理（\hyperlink{para:1:1.3.7}{1.3.7}）将蕴含：当借助 $\theta_f$ 将环 $A_f$ 与
$\Gamma(D(f),\wt A)$ 等同时，同态 $\theta_f\colon M_f\to\Gamma(D(f),\wt M)$ 是模同构。

先证明 $\theta_f$ 单射。若 $z\in M_f$ 满足 $\theta_f(z)=0$，
这表示对于 $A_f$ 的每个素理想 $\mathfrak p$，存在 $h\notin\mathfrak p$，使 $hz=0$。
于是 $z$ 的零化理想不包含于 $A_f$ 的任何素理想，因而是整个 $A_f$，故 $z=0$。

还需证明 $\theta_f$ 满射。可以归结到 $f=1$ 的情形，
一般情形借助（\hyperlink{para:1:1.3.6}{1.3.6}）作“局部化”即可得到。
于是设 $s$ 是 $\wt M$ 在 $X$ 上的截面；由（\hyperlink{para:1:1.3.4}{1.3.4}）及（\hyperlink{para:1:1.1.10}{1.1.10}，(ii)），
存在 $X$ 的一个有限覆盖 $(D(f_i))_{i\in I}$（$f_i\in A$），
使每个限制 $s_i=s|D(f_i)$ 都可写成 $\theta_{f_i}(t_i)$，其中 $t_i\in M_{f_i}$。
若 $i,j\in I$，将 $s_i$ 与 $s_j$ 在 $D(f_i)\cap D(f_j)=D(f_if_j)$ 上的限制相等写出，
由 $\wt M$ 的定义得到
\begin{equation}
\rho_{f_if_j,f_i}(t_i)=\rho_{f_if_j,f_j}(t_j).
\tag{1.3.7.1}\label{eq:1:1.3.7.1}
\end{equation}
依定义，对于每个 $i\in I$ 可写 $t_i=z_i/f_i^{n_i}$，其中 $z_i\in M$。
由于 $I$ 有限，在将各 $z_i$ 乘以 $f_i$ 的适当次幂之后，可以假定所有 $n_i$ 都等于同一个正整数 $n$。
这时，（\hyperref[eq:1:1.3.7.1]{1.3.7.1}）表示存在整数 $m_{ij}\geq0$，使
\[
(f_if_j)^{m_{ij}}(f_j^nz_i-f_i^nz_j)=0.
\]
还可假定所有 $m_{ij}$ 都等于同一个整数 $m$。
再将 $z_i$ 换成 $f_i^mz_i$，便归结到 $m=0$ 的情形，\transnote{
这里分母指数也相应从 $n$ 改为 $n+m$，之后仍以 $n$ 表示新的共同指数。
这是保持 $t_i$ 不变的同一分数表示的调整。}
也就是对于任意 $i,j$ 都有
\begin{equation}
f_j^nz_i=f_i^nz_j.
\tag{1.3.7.2}\label{eq:1:1.3.7.2}
\end{equation}
这里 $D(f_i^n)=D(f_i)$，且各 $D(f_i)$ 覆盖 $X$，
故各 $f_i^n$ 生成的理想为 $A$；换言之，存在 $g_i\in A$，使 $\sum_i g_if_i^n=1$。
考虑 $M$ 中的元素 $z=\sum_i g_iz_i$。由（\hyperref[eq:1:1.3.7.2]{1.3.7.2}），有
\[
f_i^nz=\sum_j g_jf_i^nz_j
=\Bigl(\sum_j g_jf_j^n\Bigr)z_i=z_i,
\]
因而依定义，$t_i=z/1$ 在 $M_{f_i}$ 中成立。
所以 $s_i$ 是 $\theta_1(z)$ 在 $D(f_i)$ 上的限制，从而 $s=\theta_1(z)$，证明完毕。

\egastatement{1.3.8}{推论}
设 $M,N$ 是两个 $A$-模；典范同态
\[
u\longmapsto\wt u\colon
\Hom_A(M,N)\longrightarrow\Hom_{\wt A}(\wt M,\wt N)
\]
是双射。特别地，$\wt M=0$ 与 $M=0$ 等价。

\begin{sloppypar}
事实上，考虑典范同态 $v\mapsto\Gamma(v)$，它将
$\Hom_{\wt A}(\wt M,\wt N)$ 映到
$\Hom_{\Gamma(\wt A)}(\Gamma(\wt M),\Gamma(\wt N))$；
由定理（\hyperlink{para:1:1.3.7}{1.3.7}），后一个模典范地等同于 $\Hom_A(M,N)$。
只需验证 $u\mapsto\wt u$ 与 $v\mapsto\Gamma(v)$ 互逆。
由 $\wt u$ 的定义，显然 $\Gamma(\wt u)=u$。
另一方面，对于 $v\in\Hom_{\wt A}(\wt M,\wt N)$，令 $u=\Gamma(v)$；
由 $v$ 典范得到的映射
$w\colon\Gamma(D(f),\wt M)\to\Gamma(D(f),\wt N)$ 使下图交换：
\end{sloppypar}
\[
\begin{tikzcd}[column sep=large,row sep=large]
M\arrow[r,"u"]\arrow[d,"\rho_{f,1}"'] & N\arrow[d,"\rho_{f,1}"]\\
M_f\arrow[r,"w"'] & N_f.
\end{tikzcd}
\]
因此，对于每个 $f\in A$，必有 $w=u_f$（0，\hyperlink{para:0:1.2.4}{1.2.4}），这就说明 $\wt{\Gamma(v)}=v$。

\egastatement{1.3.9}{推论}
(i) 设 $u\colon M\to N$ 是 $A$-模同态；
则与 $\Ker u$、$\Im u$、$\Coker u$ 相伴的层分别为
$\Ker\wt u$、$\Im\wt u$、$\Coker\wt u$。
特别地，$\wt u$ 单射（分别地，满射、双射），当且仅当 $u$ 具有相应性质。

(ii) 若 $M$ 是某个 $A$-模族 $(M_\lambda)$ 的归纳极限（分别地，直和），
则 $\wt M$ 在一个典范同构的意义下，是层族 $(\wt{M_\lambda})$ 的归纳极限（分别地，直和）。

(i) 只需将 $M\mapsto\wt M$ 的正合性（\hyperlink{para:1:1.3.5}{1.3.5}）应用于两个 $A$-模正合列
\[
0\longrightarrow\Ker u\longrightarrow M\longrightarrow\Im u\longrightarrow0,
\]
\[
0\longrightarrow\Im u\longrightarrow N\longrightarrow\Coker u\longrightarrow0.
\]
第二项断言随后由定理（\hyperlink{para:1:1.3.7}{1.3.7}）得到。

(ii) 设 $(M_\lambda,g_{\mu\lambda})$ 是以 $M$ 为归纳极限的 $A$-模归纳系，
$g_\lambda\colon M_\lambda\to M$ 为典范同态。
由于对 $\lambda\leq\mu\leq\nu$ 有
$\wt g_\mu\circ\wt g_{\mu\lambda}=\wt g_\lambda$ 及
$\wt g_{\nu\lambda}=\wt g_{\nu\mu}\circ\wt g_{\mu\lambda}$，
故 $(\wt{M_\lambda},\wt g_{\mu\lambda})$ 是 $X$ 上的层的归纳系。
若以 $h_\lambda\colon\wt{M_\lambda}\to\ilim\wt{M_\lambda}$ 表示典范同态，
则存在唯一同态 $v\colon\ilim\wt{M_\lambda}\to\wt M$，使 $v\circ h_\lambda=\wt g_\lambda$。
为证明 $v$ 双射，只需验证对每个 $x\in X$，$v_x$ 是
$(\ilim\wt{M_\lambda})_x$ 到 $\wt M_x$ 的双射。
而 $\wt M_x=M_x$，并且
\[
\Bigl(\ilim\wt{M_\lambda}\Bigr)_x
=\ilim(\wt{M_\lambda})_x
=\ilim(M_\lambda)_x=M_x
\qquad\text{（0，\hyperlink{para:0:1.3.3}{1.3.3}）}.
\]
另一方面，依定义，$(\wt g_\lambda)_x$ 与 $(h_\lambda)_x$ 都等于
由 $(M_\lambda)_x$ 到 $M_x$ 的典范映射。
由 $(\wt g_\lambda)_x=v_x\circ(h_\lambda)_x$ 可知 $v_x$ 是恒等映射。

最后，若 $M$ 是两个 $A$-模 $N,P$ 的直和，则显然 $\wt M=\wt N\oplus\wt P$。
任意直和都是有限直和的归纳极限，故 (ii) 的断言均已证明。

注意，（\hyperlink{para:1:1.3.8}{1.3.8}）证明了：与各 $A$-模的相伴层同构的层组成一个阿贝尔范畴（T，I，1.4）。
还应指出，由（\hyperlink{para:1:1.3.9}{1.3.9}），若 $M$ 是有限型 $A$-模，即存在满同态 $A^n\to M$，
则存在满同态 $\wt A^n\to\wt M$；也就是说，$\wt A$-模层 $\wt M$
由其在 $X$ 上的某个有限截面族生成（0，\hyperlink{para:0:5.1.1}{5.1.1}）；反之亦然。

\egapar{1.3.10}
若 $N$ 是 $A$-模 $M$ 的子模，则由（\hyperlink{para:1:1.3.9}{1.3.9}），典范包含 $j\colon N\to M$
给出单同态 $\wt N\to\wt M$，从而可以将 $\wt N$ 典范地等同于 $\wt M$ 的一个
$\wt A$-子模层；我们始终作此等同。
若 $N,P$ 是 $M$ 的两个子模，则
\begin{equation}
\wt{N+P}=\wt N+\wt P,
\tag{1.3.10.1}\label{eq:1:1.3.10.1}
\end{equation}
\begin{equation}
\wt{N\cap P}=\wt N\cap\wt P.
\tag{1.3.10.2}\label{eq:1:1.3.10.2}
\end{equation}
因为 $N+P$ 是典范同态 $N\oplus P\to M$ 的像，
$N\cap P$ 是典范同态 $M\to(M/N)\oplus(M/P)$ 的核，所以只需应用（\hyperlink{para:1:1.3.9}{1.3.9}）。
由（\hyperref[eq:1:1.3.10.1]{1.3.10.1}）与（\hyperref[eq:1:1.3.10.2]{1.3.10.2}）可知，若 $\wt N=\wt P$，则 $N=P$。

\egastatement{1.3.11}{推论}
在与 $A$-模的相伴层同构的层所组成的范畴中，函子 $\Gamma$ 正合。

事实上，设 $\wt M\xrightarrow{\wt u}\wt N\xrightarrow{\wt v}\wt P$ 是正合列，
对应于两个 $A$-模同态 $u\colon M\to N$、$v\colon N\to P$。
若 $Q=\Im u$，$R=\Ker v$，则
$\wt Q=\Im\wt u=\Ker\wt v=\wt R$（推论（\hyperlink{para:1:1.3.9}{1.3.9}）），故 $Q=R$。

\egastatement{1.3.12}{推论}
设 $M,N$ 是两个 $A$-模。

(i) 与 $M\otimes_A N$ 相伴的层典范地等同于 $\wt M\otimes_{\wt A}\wt N$。

(ii) 若 $M$ 还有有限表示，则与 $\Hom_A(M,N)$ 相伴的层和
$\shHom_{\wt A}(\wt M,\wt N)$ 典范同构。

(i) 层 $\G=\wt M\otimes_{\wt A}\wt N$ 与预层
\[
U\longmapsto\G(U)=\Gamma(U,\wt M)\otimes_{\Gamma(U,\wt A)}\Gamma(U,\wt N)
\]
相伴，其中 $U$ 遍历由各 $D(f)$（$f\in A$）组成的 $X$ 的基（\hyperlink{para:1:1.1.10}{1.1.10}，(i)）。
由（\hyperlink{para:1:1.3.7}{1.3.7}）和（\hyperlink{para:1:1.3.6}{1.3.6}），$\G(D(f))$ 典范地等同于 $M_f\otimes_{A_f}N_f$。
另一方面，此 $A_f$-模与 $(M\otimes_A N)_f$ 典范同构（0，\hyperlink{para:0:1.3.4}{1.3.4}），
而后者又典范地等同于 $\Gamma(D(f),\wt{M\otimes_A N})$（\hyperlink{para:1:1.3.7}{1.3.7} 及 \hyperlink{para:1:1.3.6}{1.3.6}）。
此外，容易验证所得典范同构
\[
\G(D(f))\xrightarrow{\sim}\Gamma(D(f),\wt{M\otimes_A N})
\]
满足与限制算子的相容条件（0，\hyperlink{para:0:1.4.2}{1.4.2}），因而定义一个典范的函子性同构
\[
\wt M\otimes_{\wt A}\wt N\xrightarrow{\sim}\wt{M\otimes_A N}.
\]

(ii) 层 $\G=\shHom_{\wt A}(\wt M,\wt N)$ 与预层
\[
U\longmapsto\G(U)=\Hom_{\wt A|U}(\wt M|U,\wt N|U)
\]
\begin{sloppypar}
相伴，其中 $U$ 遍历由各 $D(f)$ 组成的 $X$ 的基。
而 $\G(D(f))$ 与 $\Hom_{A_f}(M_f,N_f)$ 典范地等同（\hyperlink{para:1:1.3.6}{1.3.6} 及 \hyperlink{para:1:1.3.8}{1.3.8}）；
由关于 $M$ 的假设，后者又典范地等同于 $(\Hom_A(M,N))_f$（0，\hyperlink{para:0:1.3.5}{1.3.5}）。
最后，$(\Hom_A(M,N))_f$ 典范地等同于
$\Gamma(D(f),\wt{\Hom_A(M,N)})$（\hyperlink{para:1:1.3.6}{1.3.6} 及 \hyperlink{para:1:1.3.7}{1.3.7}）。
所得典范同构与限制算子相容（0，\hyperlink{para:0:1.4.2}{1.4.2}），因此定义了典范同构
\end{sloppypar}
\[
\shHom_{\wt A}(\wt M,\wt N)\xrightarrow{\sim}\wt{\Hom_A(M,N)}.
\]

\egapar{1.3.13}
现在设 $B$ 是一个（交换）$A$-代数。
这可解释为：$B$ 是一个 $A$-模，并且给定了元素 $e\in B$ 及
$A$-同态 $\varphi\colon B\otimes_A B\to B$，使下列条件成立：
第一，下面两图交换，其中 $\sigma$ 是典范对称同构：
\[
\begin{tikzcd}[column sep=large,row sep=large]
B\otimes_A B\otimes_A B\arrow[r,"\varphi\otimes1"]\arrow[d,"1\otimes\varphi"']
 & B\otimes_A B\arrow[d,"\varphi"]\\
B\otimes_A B\arrow[r,"\varphi"'] & B;
\end{tikzcd}
\qquad
\begin{tikzcd}[column sep=small,row sep=large]
B\otimes_A B\arrow[rr,"\sigma"]\arrow[dr,"\varphi"'] && B\otimes_A B\arrow[dl,"\varphi"]\\
& B &
\end{tikzcd}
\]
第二，$\varphi(e\otimes x)=\varphi(x\otimes e)=x$。
由（\hyperlink{para:1:1.3.12}{1.3.12}），$\wt A$-模层同态
$\wt\varphi\colon\wt B\otimes_{\wt A}\wt B\to\wt B$ 满足相应条件，
因而在 $\wt B$ 上定义一个 $\wt A$-代数层结构。

同样，给定一个 $B$-模 $N$，等价于给定一个 $A$-模 $N$ 及一个
$A$-同态 $\psi\colon B\otimes_A N\to N$，使下图交换，且 $\psi(e\otimes n)=n$：
\[
\begin{tikzcd}[column sep=large,row sep=large]
B\otimes_A B\otimes_A N\arrow[r,"\varphi\otimes1"]\arrow[d,"1\otimes\psi"']
 & B\otimes_A N\arrow[d,"\psi"]\\
B\otimes_A N\arrow[r,"\psi"'] & N.
\end{tikzcd}
\]
同态 $\wt\psi\colon\wt B\otimes_{\wt A}\wt N\to\wt N$ 满足相应条件，
因而在 $\wt N$ 上定义一个 $\wt B$-模层结构。

同理，若 $u\colon B\to B'$（分别地，$v\colon N\to N'$）是
$A$-代数同态（分别地，$B$-模同态），则 $\wt u$（分别地，$\wt v$）是
$\wt A$-代数层同态（分别地，$\wt B$-模层同态），
$\Ker\wt u$ 是 $\wt B$-理想层（分别地，$\Ker\wt v$、$\Coker\wt v$ 和 $\Im\wt v$
都是 $\wt B$-模层）。若 $N$ 是 $B$-模，则 $\wt N$ 是有限型 $\wt B$-模层，
当且仅当 $N$ 是有限型 $B$-模（0，\hyperlink{para:0:5.2.3}{5.2.3}）。

若 $M,N$ 是两个 $B$-模，则 $\wt B$-模层 $\wt M\otimes_{\wt B}\wt N$
典范地等同于 $\wt{M\otimes_B N}$；同样，当 $M$ 有有限表示时，
$\shHom_{\wt B}(\wt M,\wt N)$ 典范地等同于 $\wt{\Hom_B(M,N)}$。
证明与（\hyperlink{para:1:1.3.12}{1.3.12}）相同。

若 $\frJ$ 是 $B$ 的理想，$N$ 是 $B$-模，则 $\wt{\frJ N}=\wt{\frJ}\,\wt N$。

最后，若 $B$ 是由 $A$-子模 $B_n$（$n\in\mathbb Z$）给出分次的 $A$-代数，
则 $\wt A$-代数层 $\wt B$ 是各 $\wt A$-模层 $\wt{B_n}$ 的直和（\hyperlink{para:1:1.3.9}{1.3.9}），
并由这些 $\wt A$-子模层给出分次；分次公理表示同态
$B_m\otimes_A B_n\to B$ 的像包含于 $B_{m+n}$。
同样，若 $M$ 是由子模 $M_n$ 给出分次的 $B$-模，
则 $\wt M$ 是由各 $\wt{M_n}$ 给出分次的 $\wt B$-模层。

\egapar{1.3.14}
若 $B$ 是 $A$-代数，$M$ 是 $B$ 的子模，则由 $\wt M$ 生成的
$\wt B$ 的 $\wt A$-子代数层（0，\hyperlink{para:0:4.1.3}{4.1.3}）就是 $\wt C$，
其中 $C$ 是由 $M$ 生成的 $B$ 的子代数。
事实上，$C$ 是同态 $\bigotimes_A^n M\to B$（$n\geq0$）的像所成的各子模之和，
因此只需应用（\hyperlink{para:1:1.3.9}{1.3.9}）和（\hyperlink{para:1:1.3.12}{1.3.12}）。
\par


% END chapters/ch1_s1_3.tex

\par

% BEGIN chapters/ch1_s1_4.tex
\subsection{素谱上的拟凝聚层}\hypertarget{heading:1:1.4}{}
\src{90—92}

\egastatement{1.4.1}{定理}
设 $X$ 是环 $A$ 的素谱，$V$ 是 $X$ 的拟紧开子集，$\F$ 是一个 $(\OO_X|V)$-模层。
下列四个条件等价：
\begin{enumerate}[label=\alph*)]
\item 存在 $A$-模 $M$，使 $\F$ 同构于 $\wt M|V$。
\item 存在 $V$ 的有限开覆盖 $(V_i)$，其中各 $V_i$ 均为包含于 $V$ 的
$D(f_i)$（$f_i\in A$），并且对每个 $i$，$\F|V_i$ 同构于形如 $\wt{M_{f_i}}$ 的层，
其中 $M_{f_i}$ 是一个 $A_{f_i}$-模。
\item 层 $\F$ 拟凝聚（0，\hyperlink{para:0:5.1.3}{5.1.3}）。
\item 下列两个性质成立：
\begin{enumerate}[label={$d\arabic*)$},leftmargin=2.7em]
\item 对于每个满足 $D(f)\subset V$ 的 $f\in A$，及每个截面
$s\in\Gamma(D(f),\F)$，存在整数 $n\geq0$，使 $f^ns$ 可延拓为 $\F$ 在 $V$ 上的一个截面。
\item 对于每个满足 $D(f)\subset V$ 的 $f\in A$，及每个在 $D(f)$ 上的限制为 $0$ 的截面
$t\in\Gamma(V,\F)$，存在整数 $n\geq0$，使 $f^nt=0$。
\end{enumerate}
\end{enumerate}

在条件 $d1)$ 和 $d2)$ 的写法中，已依（\hyperlink{para:1:1.3.7}{1.3.7}）默认将 $A$ 与 $\Gamma(\wt A)$ 等同。

$a)\Rightarrow b)$ 直接由（\hyperlink{para:1:1.3.6}{1.3.6}）及各 $D(f_i)$ 构成 $X$ 的拓扑的一组基（\hyperlink{para:1:1.1.10}{1.1.10}）得到。
由于每个 $A$-模都同构于某个同态 $A^{(I)}\to A^{(J)}$ 的余核，
（\hyperlink{para:1:1.3.9}{1.3.9}）说明与每个 $A$-模相伴的层都是拟凝聚的；故 $b)\Rightarrow c)$。
反之，若 $\F$ 拟凝聚，则每个 $x\in V$ 都有一个形如 $D(f)\subset V$ 的邻域，
使 $\F|D(f)$ 同构于某个同态
$\wt{A_f}^{(I)}\to\wt{A_f}^{(J)}$ 的余核，
因而同构于 $A_f$-模 $N$ 的相伴层，其中 $N$ 是对应同态
$A_f^{(I)}\to A_f^{(J)}$ 的余核（\hyperlink{para:1:1.3.8}{1.3.8} 及 \hyperlink{para:1:1.3.9}{1.3.9}）。
因为 $V$ 拟紧，显然 $c)\Rightarrow b)$。

为证明 $b)$ 蕴含 $d1)$ 和 $d2)$，先假定 $V=D(g)$，其中 $g\in A$，
并且 $\F$ 同构于某个 $A_g$-模 $N$ 的相伴层 $\wt N$。
将 $X$ 换成 $V$、$A$ 换成 $A_g$（\hyperlink{para:1:1.3.6}{1.3.6}），便可归结到 $g=1$ 的情形。
这时 $\Gamma(D(f),\wt N)$ 与 $N_f$ 典范地等同（\hyperlink{para:1:1.3.6}{1.3.6} 及 \hyperlink{para:1:1.3.7}{1.3.7}），
故截面 $s\in\Gamma(D(f),\wt N)$ 与形如 $z/f^n$ 的元素等同，其中 $z\in N$。
截面 $f^ns$ 与 $N_f$ 中的元素 $z/1$ 等同，因而是与元素 $z\in N$ 等同的
$\wt N$ 在 $X$ 上的截面在 $D(f)$ 上的限制。这就证明了此情形中的 $d1)$。
同样，$t\in\Gamma(X,\wt N)$ 与某个元素 $z\in N$ 等同，
$t$ 在 $D(f)$ 上的限制与 $z$ 在 $N_f$ 中的像 $z/1$ 等同；
这一像为零，表示存在整数 $n\geq0$，使 $f^nz=0$ 在 $N$ 中成立，
也等价于 $f^nt=0$。

为完成 $b)$ 蕴含 $d1)$、$d2)$ 的证明，只需建立下面的引理。

\egastatement{1.4.1.1}{引理}
假定 $V$ 是有限多个形如 $D(g_i)$ 的集合之并，且各层
$\F|D(g_i)$、$\F|(D(g_i)\cap D(g_j))=\F|D(g_ig_j)$ 都满足 $d1)$ 和 $d2)$。
则 $\F$ 具有以下两个性质：
\begin{enumerate}[label={$d'\arabic*)$},leftmargin=2.7em]
\item 对于每个 $f\in A$ 及每个截面 $s\in\Gamma(D(f)\cap V,\F)$，
存在整数 $n\geq0$，使 $f^ns$ 可延拓为 $\F$ 在 $V$ 上的一个截面。
\item 对于每个 $f\in A$ 及每个在 $D(f)\cap V$ 上的限制为 $0$ 的截面
$t\in\Gamma(V,\F)$，存在整数 $n\geq0$，使 $f^nt=0$。
\end{enumerate}

先证明 $d'2)$。因为 $D(f)\cap D(g_i)=D(fg_i)$，对于每个 $i$，
存在整数 $n_i$，使 $(fg_i)^{n_i}t$ 在 $D(g_i)$ 上的限制为零。
由于 $g_i$ 在 $A_{g_i}$ 中的像可逆，$f^{n_i}t$ 在 $D(g_i)$ 上的限制也为零。
取 $n$ 为各 $n_i$ 中的最大者，即得 $d'2)$。

为证明 $d'1)$，将 $d1)$ 应用于层 $\F|D(g_i)$。
存在整数 $n_i\geq0$ 及 $\F$ 在 $D(g_i)$ 上的截面 $s'_i$，
使 $s'_i$ 延拓了 $(fg_i)^{n_i}s$ 在 $D(fg_i)$ 上的限制。
由于 $g_i$ 在 $A_{g_i}$ 中的像可逆，存在 $\F$ 在 $D(g_i)$ 上的截面 $s_i$，
使 $s'_i=g_i^{n_i}s_i$；于是 $s_i$ 延拓了 $f^{n_i}s$ 在 $D(fg_i)$ 上的限制。
还可以假定所有 $n_i$ 都等于同一个整数 $n$。
由构造，$s_i-s_j$ 在
$D(f)\cap D(g_i)\cap D(g_j)=D(fg_ig_j)$ 上的限制为零。
将 $d2)$ 应用于层 $\F|D(g_ig_j)$，可知存在整数 $m_{ij}\geq0$，
使 $(fg_ig_j)^{m_{ij}}(s_i-s_j)$ 在 $D(g_ig_j)$ 上的限制为零。
由于 $g_ig_j$ 在 $A_{g_ig_j}$ 中的像可逆，
$f^{m_{ij}}(s_i-s_j)$ 在 $D(g_ig_j)$ 上的限制为零。
于是可假定所有 $m_{ij}$ 都等于同一个整数 $m$，
从而存在截面 $s'\in\Gamma(V,\F)$，延拓各 $f^ms_i$。
此截面因而延拓了 $f^{n+m}s$，这就证明了 $d'1)$。

现在还需证明，$d1)$ 和 $d2)$ 蕴含 $a)$。
先证明：若 $\F$ 满足 $d1)$ 和 $d2)$，则对每个满足 $D(g)\subset V$ 的 $g\in A$，
层 $\F|D(g)$ 也满足这两个条件。对于 $d1)$，这是显然的。
另一方面，若 $t\in\Gamma(D(g),\F)$ 在 $D(f)\subset D(g)$ 上的限制为零，
则由 $d1)$，存在整数 $m\geq0$，使 $g^mt$ 延拓为 $\F$ 在 $V$ 上的截面 $s$。
应用 $d2)$ 可知，存在整数 $n\geq0$，使 $f^ns=0$；
由于 $g$ 在 $A_g$ 中的像可逆，故 $f^nt=0$。

由此，因为 $V$ 拟紧，引理（\hyperlink{para:1:1.4.1.1}{1.4.1.1}）说明条件 $d'1)$ 与 $d'2)$ 成立。
考虑 $A$-模 $M=\Gamma(V,\F)$，并定义一个 $\wt A$-模层同态
\[
u\colon\wt M\longrightarrow j_*(\F),
\]
其中 $j\colon V\to X$ 是典范包含映射。
由于各 $D(f)$ 构成 $X$ 的拓扑的一组基，只需对每个 $f\in A$ 定义同态
\[
u_f\colon M_f\longrightarrow\Gamma(D(f),j_*(\F))
=\Gamma(D(f)\cap V,\F),
\]
并验证通常的相容条件（0，\hyperlink{para:0:3.2.5}{3.2.5}）。
因为 $f$ 在 $A_f$ 中的典范像可逆，限制同态
$M=\Gamma(V,\F)\to\Gamma(D(f)\cap V,\F)$ 分解为
\[
M\longrightarrow M_f\xrightarrow{u_f}\Gamma(D(f)\cap V,\F)
\qquad\text{（0，\hyperlink{para:0:1.2.4}{1.2.4}）};
\]
当 $D(g)\subset D(f)$ 时，相容条件的验证是直接的。

下面证明，条件 $d'1)$（分别地，$d'2)$）蕴含每个 $u_f$ 满射（分别地，单射）。
这将说明 $u$ 双射，从而 $\F$ 是一个同构于 $\wt M$ 的 $\wt A$-模层在 $V$ 上的限制。
事实上，若 $s\in\Gamma(D(f)\cap V,\F)$，则由 $d'1)$，存在整数 $n\geq0$，
使 $f^ns$ 延拓为某个截面 $z\in M$。于是 $u_f(z/f^n)=s$，故 $u_f$ 满射。
同样，若 $z\in M$ 满足 $u_f(z/1)=0$，这表示截面 $z$ 在 $D(f)\cap V$ 上的限制为零。
由 $d'2)$，存在整数 $n\geq0$，使 $f^nz=0$，于是 $z/1=0$ 在 $M_f$ 中成立，
故 $u_f$ 单射。证毕。

\egastatement{1.4.2}{推论}
$X$ 的一个拟紧开集上的每个拟凝聚层，都由 $X$ 上的某个拟凝聚层诱导。

\egastatement{1.4.3}{推论}
$X=\Spec(A)$ 上的每个拟凝聚 $\OO_X$-代数层，都同构于形如 $\wt B$ 的
$\OO_X$-代数层，其中 $B$ 是一个 $A$-代数；
每个拟凝聚 $\wt B$-模层都同构于形如 $\wt N$ 的 $\wt B$-模层，其中 $N$ 是一个 $B$-模。

事实上，一个拟凝聚 $\OO_X$-代数层是拟凝聚 $\OO_X$-模层，
因而形如 $\wt B$，其中 $B$ 是一个 $A$-模。
由用 $\wt A$-模层同态 $\wt B\otimes_{\wt A}\wt B\to\wt B$
对 $\OO_X$-代数层结构的刻画，以及（\hyperlink{para:1:1.3.12}{1.3.12}），可知 $B$ 是 $A$-代数。
若 $\F$ 是拟凝聚 $\wt B$-模层，只需证明 $\F$ 也是拟凝聚 $\wt A$-模层，即可同样得出结论。
由于问题是局部的，限制到 $X$ 的某个形如 $D(f)$ 的开集上，便可假定 $\F$ 是某个
$\wt B$-模层同态 $\wt B^{(I)}\to\wt B^{(J)}$ 的余核；
此同态当然也是 $\wt A$-模层同态。
于是由（\hyperlink{para:1:1.3.8}{1.3.8}）及（\hyperlink{para:1:1.3.9}{1.3.9}）得到本命题。
\par


% END chapters/ch1_s1_4.tex

\par

% BEGIN chapters/ch1_s1_5.tex
\subsection{素谱上的凝聚层}\hypertarget{heading:1:1.5}{}
\src{92—93}

\egastatement{1.5.1}{定理}
设 $A$ 是诺特环，$X=\Spec(A)$ 是其素谱，$V$ 是 $X$ 的开子集，
$\F$ 是一个 $(\OO_X|V)$-模层。下列条件等价：
\begin{enumerate}[label=\alph*)]
\item $\F$ 凝聚。
\item $\F$ 有限型且拟凝聚。
\item 存在有限型 $A$-模 $M$，使 $\F$ 同构于层 $\wt M|V$。
\end{enumerate}

$a)$ 显然蕴含 $b)$。为证明 $b)$ 蕴含 $c)$，先注意，由于 $V$ 拟紧（0，\hyperlink{para:0:2.2.3}{2.2.3}），
$\F$ 同构于一个形如 $\wt N|V$ 的层，其中 $N$ 是 $A$-模（\hyperlink{para:1:1.4.1}{1.4.1}）。
而 $N=\ilim M_\lambda$，其中 $M_\lambda$ 遍历 $N$ 的有限型 $A$-子模，
因而由（\hyperlink{para:1:1.3.9}{1.3.9}）有
$\F=\wt N|V=\ilim(\wt{M_\lambda}|V)$。
因为 $\F$ 有限型且 $V$ 拟紧，存在指标 $\lambda$，使
$\F=\wt{M_\lambda}|V$（0，\hyperlink{para:0:5.2.3}{5.2.3}）。

最后证明 $c)$ 蕴含 $a)$。显然，这时 $\F$ 有限型（\hyperlink{para:1:1.3.6}{1.3.6} 及 \hyperlink{para:1:1.3.9}{1.3.9}）；
又因为问题是局部的，可以只考虑 $V=D(f)$、$f\in A$ 的情形。
由于 $A_f$ 是诺特环，最终只需证明：形如 $\wt A^n\to\wt M$ 的同态的核有限型，
其中 $M$ 是一个 $A$-模。
而这种同态形如 $\wt u$，其中 $u$ 是同态 $A^n\to M$（\hyperlink{para:1:1.3.8}{1.3.8}）；
若 $P=\Ker u$，则 $\wt P=\Ker\wt u$（\hyperlink{para:1:1.3.9}{1.3.9}）。
因为 $A$ 是诺特环，$P$ 有限型，证明完毕。

\egastatement{1.5.2}{推论}
在（\hyperlink{para:1:1.5.1}{1.5.1}）的假设下，层 $\OO_X$ 是凝聚环层。

\egastatement{1.5.3}{推论}
在（\hyperlink{para:1:1.5.1}{1.5.1}）的假设下，$X$ 的任意开集上的每个凝聚层，
都由 $X$ 上的某个凝聚层诱导。

\egastatement{1.5.4}{推论}
在（\hyperlink{para:1:1.5.1}{1.5.1}）的假设下，每个拟凝聚 $\OO_X$-模层 $\F$
都是其凝聚 $\OO_X$-子模层的归纳极限。

事实上，$\F=\wt M$，其中 $M$ 是某个 $A$-模，
而 $M$ 是其有限型子模的归纳极限；由（\hyperlink{para:1:1.3.9}{1.3.9}）及（\hyperlink{para:1:1.5.1}{1.5.1}）即得结论。
\par


% END chapters/ch1_s1_5.tex

\par

% BEGIN chapters/ch1_s1_6.tex
\subsection{素谱上拟凝聚层的函子性质}\hypertarget{heading:1:1.6}{}
\src{93—96}

\egapar{1.6.1}
设 $A,A'$ 是两个环，$\varphi\colon A'\to A$ 是同态，
\[
\aphi\colon X=\Spec(A)\longrightarrow X'=\Spec(A')
\]
是与 $\varphi$ 相伴的连续映射（\hyperlink{para:1:1.2.1}{1.2.1}）。我们将定义一个典范的环层同态
\[
\wt\varphi\colon\OO_{X'}\longrightarrow(\aphi)_*(\OO_X).
\]
对于每个 $f'\in A'$，令 $f=\varphi(f')$；由（\hyperref[eq:1:1.2.2.2]{1.2.2.2}），有
$(\aphi)^{-1}(D(f'))=D(f)$。
环 $\Gamma(D(f'),\wt{A'})$ 与 $\Gamma(D(f),\wt A)$
分别等同于 $A'_{f'}$ 与 $A_f$（\hyperlink{para:1:1.3.6}{1.3.6} 及 \hyperlink{para:1:1.3.7}{1.3.7}）。
而同态 $\varphi$ 典范地定义同态
$\varphi_{f'}\colon A'_{f'}\to A_f$（0，\hyperlink{para:0:1.5.1}{1.5.1}）；也就是说，有环同态
\[
\Gamma(D(f'),\wt{A'})
\longrightarrow\Gamma((\aphi)^{-1}(D(f')),\wt A)
=\Gamma(D(f'),(\aphi)_*(\wt A)).
\]
此外，这些同态满足通常的相容条件：当 $D(f')\supset D(g')$ 时，下图交换（0，\hyperlink{para:0:1.5.1}{1.5.1}）：
\[
\begin{tikzcd}[column sep=large,row sep=large]
\Gamma(D(f'),\wt{A'})\arrow[r]\arrow[d]
 & \Gamma(D(f'),(\aphi)_*(\wt A))\arrow[d]\\
\Gamma(D(g'),\wt{A'})\arrow[r]
 & \Gamma(D(g'),(\aphi)_*(\wt A)).
\end{tikzcd}
\]
由于各 $D(f')$ 构成 $X'$ 的拓扑的一组基（0，\hyperlink{para:0:3.2.3}{3.2.3}），
便确实定义了一个 $\OO_{X'}$-代数层同态。
因此，二元组 $\Phi=(\aphi,\wt\varphi)$ 是赋环空间态射
\[
\Phi\colon(X,\OO_X)\longrightarrow(X',\OO_{X'})
\qquad\text{（0，\hyperlink{para:0:4.1.1}{4.1.1}）}.
\]

还应指出，若令 $x'=\aphi(x)$，则同态 $\wt\varphi_x^\sharp$（0，\hyperlink{para:0:3.7.1}{3.7.1}）
正是由 $\varphi\colon A'\to A$ 典范得到的同态
\[
\varphi_x\colon A'_{x'}\longrightarrow A_x
\qquad\text{（0，\hyperlink{para:0:1.5.1}{1.5.1}）}.
\]
事实上，每个 $z'\in A'_{x'}$ 均可写成 $g'/f'$，其中 $f',g'\in A'$ 且 $f'\notin\jj_{x'}$。
所以 $D(f')$ 是 $x'$ 在 $X'$ 中的邻域，而由 $\wt\varphi$ 得到的同态
$\Gamma(D(f'),\wt{A'})\to\Gamma((\aphi)^{-1}(D(f')),\wt A)$ 就是 $\varphi_{f'}$。
考虑对应于 $g'/f'\in A'_{f'}$ 的截面 $s'\in\Gamma(D(f'),\wt{A'})$，
便得到 $\wt\varphi_x^\sharp(z')=\varphi(g')/\varphi(f')$ 在 $A_x$ 中成立。

\egapar{1.6.2}
\textbf{例。}设 $S$ 是 $A$ 的乘法子集，$\varphi$ 是典范同态 $A\to S^{-1}A$。
由（\hyperlink{para:1:1.2.6}{1.2.6}）已知，$\aphi$ 将 $Y=\Spec(S^{-1}A)$ 同胚地映到
$X=\Spec(A)$ 中由满足 $\jj_x\cap S=\varnothing$ 的点组成的子空间。
此外，对于属于这个子空间的每个 $x$，写作 $x=\aphi(y)$（$y\in Y$），
同态 $\wt\varphi_y^\sharp\colon\OO_x\to\OO_y$ 是双射（0，\hyperlink{para:0:1.2.6}{1.2.6}）。
换言之，$\OO_Y$ 与 $\OO_X$ 在 $Y$ 上诱导的层等同。

\egastatement{1.6.3}{命题}
对于每个 $A$-模 $M$，存在从 $\OO_{X'}$-模层 $\wt{\restmod M\varphi}$
到直像 $\Phi_*(\wt M)$ 的典范函子性同构。

为简便，令 $M'=\restmod M\varphi$；对于每个 $f'\in A'$，仍令 $f=\varphi(f')$。
截面模 $\Gamma(D(f'),\wt{M'})$ 与 $\Gamma(D(f),\wt M)$ 分别等同于
$M'_{f'}$ 与 $M_f$，它们分别是 $A'_{f'}$-模与 $A_f$-模。
此外，$A'_{f'}$-模 $\restmod{(M_f)}{\varphi_{f'}}$ 与 $M'_{f'}$ 典范同构（0，\hyperlink{para:0:1.5.2}{1.5.2}）。
因此有 $\Gamma(D(f'),\wt{A'})$-模的函子性同构
\[
\Gamma(D(f'),\wt{M'})
\xrightarrow{\sim}
\restmod{\Gamma((\aphi)^{-1}(D(f')),\wt M)}{\varphi_{f'}}.
\]
这些同构满足通常的与限制相容的条件（0，\hyperlink{para:0:1.5.6}{1.5.6}），因而定义了所述函子性同构。
更具体地，若 $u\colon M_1\to M_2$ 是 $A$-模同态，
则也可将它视为 $A'$-模同态
$\restmod{(M_1)}\varphi\to\restmod{(M_2)}\varphi$；
若将后者记为 $u_{[\varphi]}$，则 $\Phi_*(\wt u)$ 与 $\wt{u_{[\varphi]}}$ 等同。

此证明还说明，对于每个 $A$-代数 $B$，典范的函子性同构
$\wt{\restmod B\varphi}\to\Phi_*(\wt B)$ 是 $\OO_{X'}$-代数层同构。
若 $M$ 是 $B$-模，则典范的函子性同构
$\wt{\restmod M\varphi}\to\Phi_*(\wt M)$ 是 $\Phi_*(\wt B)$-模层同构。

\egastatement{1.6.4}{推论}
直像函子 $\Phi_*$ 在拟凝聚 $\OO_X$-模层范畴上正合。

事实上，显然 $\restmod M\varphi$ 是关于 $M$ 的正合函子，
而 $\wt{M'}$ 是关于 $M'$ 的正合函子（\hyperlink{para:1:1.3.5}{1.3.5}）。

\egastatement{1.6.5}{命题}
设 $N'$ 是 $A'$-模，$N$ 是 $A$-模 $N'\otimes_{A'}\restmod A\varphi$。
存在从 $\OO_X$-模层 $\Phi^*(\wt{N'})$ 到 $\wt N$ 的典范函子性同构。

先注意，映射 $j\colon z'\mapsto z'\otimes1$ 是从 $N'$ 到
$\restmod N\varphi$ 的 $A'$-同态。事实上，依定义，对于 $f'\in A'$，有
\[
(f'z')\otimes1=z'\otimes\varphi(f')=\varphi(f')(z'\otimes1).
\]
由（\hyperlink{para:1:1.3.8}{1.3.8}）得到 $\OO_{X'}$-模层同态
$\wt j\colon\wt{N'}\to\wt{\restmod N\varphi}$；
再由（\hyperlink{para:1:1.6.3}{1.6.3}），可以将 $\wt j$ 看成从 $\wt{N'}$ 到 $\Phi_*(\wt N)$ 的同态。
与此同态 $\wt j$ 典范对应的是同态
$h=\wt j^\sharp\colon\Phi^*(\wt{N'})\to\wt N$（0，\hyperlink{para:0:4.4.3}{4.4.3}）。
我们将证明每个茎上的 $h_x$ 都是双射。

令 $x'=\aphi(x)$，取 $f'\in A'$ 使 $x'\in D(f')$，并令 $f=\varphi(f')$。
环 $\Gamma(D(f),\wt A)$ 与 $A_f$ 等同；
模 $\Gamma(D(f),\wt N)$、$\Gamma(D(f'),\wt{N'})$ 分别与 $N_f$、$N'_{f'}$ 等同。
设截面 $s'\in\Gamma(D(f'),\wt{N'})$ 对应于 $n'/f'^p$（$n'\in N'$），
$s$ 是它经 $\wt j$ 在 $\Gamma(D(f),\wt N)$ 中的像；
则 $s$ 对应于 $(n'\otimes1)/f^p$。
另设 $t\in\Gamma(D(f),\wt A)$ 对应于 $g/f^q$（$g\in A$）；依定义，有
\[
h_x(s'_{x'}\otimes t_x)=t_xs_x
\qquad\text{（0，\hyperlink{para:0:4.4.3}{4.4.3}）}.
\]
而 $N_f$ 可以典范地等同于
$N'_{f'}\otimes_{A'_{f'}}\restmod{(A_f)}{\varphi_{f'}}$（0，\hyperlink{para:0:1.5.4}{1.5.4}）；
这时 $s$ 对应于元素 $(n'/f'^p)\otimes1$，而截面 $y\mapsto t_ys_y$
对应于 $(n'/f'^p)\otimes(g/f^q)$。
（0，\hyperlink{para:0:1.5.6}{1.5.6}）中的相容交换图说明，$h_x$ 就是典范同构
\begin{equation}
N'_{x'}\otimes_{A'_{x'}}\restmod{(A_x)}{\varphi_x}
\xrightarrow{\sim}
N_x=(N'\otimes_{A'}\restmod A\varphi)_x.
\tag{1.6.5.1}\label{eq:1:1.6.5.1}
\end{equation}

此外，设 $v\colon N'_1\to N'_2$ 是 $A'$-模同态。
由于每个 $x'\in X'$ 都满足 $\wt v_{x'}=v_{x'}$，
从上述论证立即得到：$\Phi^*(\wt v)$ 典范地等同于 $\wt{v\otimes1}$。
这就完成了（\hyperlink{para:1:1.6.5}{1.6.5}）的证明。

若 $B'$ 是 $A'$-代数，则从 $\Phi^*(\wt{B'})$ 到
$\wt{B'\otimes_{A'}\restmod A\varphi}$ 的典范同构是 $\OO_X$-代数层同构。
若再有一个 $B'$-模 $N'$，则从 $\Phi^*(\wt{N'})$ 到
$\wt{N'\otimes_{A'}\restmod A\varphi}$ 的典范同构是 $\Phi^*(\wt{B'})$-模层同构。

\egastatement{1.6.6}{推论}
当 $s'$ 遍历 $A'$-模 $\Gamma(\wt{N'})$ 时，
这些截面 $s'$ 在 $\Phi^*(\wt{N'})$ 中的典范像生成
$A$-模 $\Gamma(\Phi^*(\wt{N'}))$。\transnote{
原版此处将 $\Gamma(\wt{N'})$ 称为“$A$-模”；按（\hyperlink{para:1:1.6.5}{1.6.5}），$N'$ 是 $A'$-模，
故这里将该处写作“$A'$-模”。结论中的 $\Gamma(\Phi^*(\wt{N'}))$ 则确为 $A$-模。}

事实上，当将 $N'$、$N$ 分别与 $\Gamma(\wt{N'})$、$\Gamma(\wt N)$ 等同（\hyperlink{para:1:1.3.7}{1.3.7}）时，
这些像就是 $N$ 中的元素 $z'\otimes1$，其中 $z'$ 遍历 $N'$。

\egapar{1.6.7}
在（\hyperlink{para:1:1.6.5}{1.6.5}）的证明中，已顺便证明典范映射（0，\hyperref[eq:0:4.4.3.2]{4.4.3.2}）
\[
\rho\colon\wt{N'}\longrightarrow\Phi_*(\Phi^*(\wt{N'}))
\]
就是同态 $\wt j$，其中
$j\colon N'\to N'\otimes_{A'}\restmod A\varphi$ 是映射 $z'\mapsto z'\otimes1$。
同样，典范映射（0，\hyperref[eq:0:4.4.3.3]{4.4.3.3}）
\[
\sigma\colon\Phi^*(\Phi_*(\wt M))\longrightarrow\wt M
\]
就是 $\wt p$，其中
$p\colon\restmod M\varphi\otimes_{A'}\restmod A\varphi\to M$
是将每个纯张量 $z\otimes a$（$z\in M$，$a\in A$）映为 $az$ 的典范同态。
这直接由定义得到（0，\hyperlink{para:0:3.7.1}{3.7.1}；0，\hyperlink{para:0:4.4.3}{4.4.3}；I，\hyperlink{para:1:1.3.7}{1.3.7}）。

由此可知（0，\hyperlink{para:0:4.4.3}{4.4.3} 及 \hyperref[eq:0:3.5.4.4]{3.5.4.4}），若 $v\colon N'\to\restmod M\varphi$
是 $A'$-同态，则有 $\wt v^\sharp=\wt{v\otimes1}$。\transnote{
此式照录原文的简写。若 $v\otimes1$ 按通常的张量积同态理解，
它的靶是 $\restmod M\varphi\otimes_{A'}A$，还须与上一段的 $p$ 复合。
将靶明确写为 $M$ 时，完整的同态为
$z'\otimes a\mapsto a\,v(z')$，亦即
$\wt v^\sharp=\wt{p\circ(v\otimes1)}$。}

\egapar{1.6.8}
设 $N'_1,N'_2$ 是两个 $A'$-模，且 $N'_1$ 有有限表示。
由（\hyperlink{para:1:1.6.7}{1.6.7}）及（\hyperlink{para:1:1.3.12}{1.3.12}，(ii)）可知，典范同态（0，\hyperlink{para:0:4.4.6}{4.4.6}）
\[
\Phi^*\bigl(\shHom_{\wt{A'}}(\wt{N'_1},\wt{N'_2})\bigr)
\longrightarrow
\shHom_{\wt A}\bigl(\Phi^*(\wt{N'_1}),\Phi^*(\wt{N'_2})\bigr)
\]
就是 $\wt\gamma$，其中 $\gamma$ 是典范的 $A$-模同态
\[
\Hom_{A'}(N'_1,N'_2)\otimes_{A'}A
\longrightarrow
\Hom_A(N'_1\otimes_{A'}A,N'_2\otimes_{A'}A).
\]

\egapar{1.6.9}
设 $\frJ'$ 是 $A'$ 的理想，$M$ 是 $A$-模。
依定义，$\wt{\frJ'}\,\wt M$ 是典范同态
\[
\Phi^*(\wt{\frJ'})\otimes_{\wt A}\wt M\longrightarrow\wt M
\]
的像，故由（\hyperlink{para:1:1.6.5}{1.6.5}）及（\hyperlink{para:1:1.3.12}{1.3.12}，(i)），
$\wt{\frJ'}\,\wt M$ 典范地等同于 $\wt{\frJ'M}$。
特别地，$\Phi^*(\wt{\frJ'})\wt A$ 与 $\wt{\frJ'A}$ 等同。
再考虑函子 $\Phi^*$ 的右正合性，$\wt A$-代数层
$\Phi^*(\wt{A'/\frJ'})$ 便与 $\wt{A/\frJ'A}$ 等同。

\egapar{1.6.10}
设 $A''$ 是第三个环，$\varphi'\colon A''\to A'$ 是同态，令 $\varphi''=\varphi\circ\varphi'$。
直接由定义，有
\[
{}^a\varphi''=({}^a\varphi')\circ({}^a\varphi),
\qquad
\wt{\varphi''}=\wt\varphi\circ\wt{\varphi'}
\qquad\text{（0，\hyperlink{para:0:1.5.7}{1.5.7}）}.
\]
因此 $\Phi''=\Phi'\circ\Phi$；换言之，$(\Spec(A),\wt A)$ 是关于 $A$ 的函子，
从环范畴取值于赋环空间范畴。\transnote{
这里 $\wt\varphi\circ\wt{\varphi'}$ 按不同底空间上层态射的复合来理解（0，\hyperlink{para:0:3.5.2}{3.5.2}）；
若均写成直像层之间的同态，则为
$({}^a\varphi')_*(\wt\varphi)\circ\wt{\varphi'}$。
由 $\Phi''=\Phi'\circ\Phi$ 可见，这个环到赋环空间的函子是反变的。}
\par


% END chapters/ch1_s1_6.tex

\par

% BEGIN chapters/ch1_s1_7.tex
\subsection{仿射概形态射的刻画}\hypertarget{heading:1:1.7}{}
\src{96—97}

\egastatement{1.7.1}{定义}
若赋环空间 $(X,\OO_X)$ 同构于某个形如 $(\Spec(A),\wt A)$ 的赋环空间，其中 $A$ 是环，
则称 $(X,\OO_X)$ 为\emph{仿射概形}。
这时，与环 $A$ 典范地等同的环 $\Gamma(X,\OO_X)$（\hyperlink{para:1:1.3.7}{1.3.7}）
称为仿射概形 $(X,\OO_X)$ 的\emph{环}；在不致混淆时，将它记作 $A(X)$。

为简便起见，当说到仿射概形 $\Spec(A)$ 时，始终是指赋环空间 $(\Spec(A),\wt A)$。

\egapar{1.7.2}
设 $A,B$ 是两个环，$(X,\OO_X)$、$(Y,\OO_Y)$ 是对应于素谱
$X=\Spec(A)$、$Y=\Spec(B)$ 的仿射概形。
由（\hyperlink{para:1:1.6.1}{1.6.1}）已知，每个环同态 $\varphi\colon B\to A$ 对应于态射
\[
\Phi=(\aphi,\wt\varphi)=\Spec(\varphi)
\colon(X,\OO_X)\longrightarrow(Y,\OO_Y).
\]
注意，$\varphi$ 由 $\Phi$ 完全确定，因为依定义有
\[
\varphi=\Gamma(\wt\varphi)\colon
\Gamma(\wt B)\longrightarrow\Gamma((\aphi)_*(\wt A))=\Gamma(\wt A).
\]

\egastatement{1.7.3}{定理}
设 $(X,\OO_X)$、$(Y,\OO_Y)$ 是两个仿射概形。
赋环空间态射 $(\psi,\theta)\colon(X,\OO_X)\to(Y,\OO_Y)$
可写成 $(\aphi,\wt\varphi)$，其中 $\varphi$ 是环同态 $A(Y)\to A(X)$，
当且仅当对于每个 $x\in X$，$\theta_x^\sharp$ 是局部同态
$\OO_{\psi(x)}\to\OO_x$。

令 $A=A(X)$，$B=A(Y)$。此条件是必要的，因为由（\hyperlink{para:1:1.6.1}{1.6.1}）已知，
$\wt\varphi_x^\sharp$ 是由 $\varphi$ 典范得到的同态 $B_{\aphi(x)}\to A_x$；
按 $\aphi(x)=\varphi^{-1}(\jj_x)$ 的定义，此同态是局部的。

证明此条件充分。依定义，$\theta$ 是同态 $\OO_Y\to\psi_*(\OO_X)$，
由此典范地得到环同态
\[
\varphi=\Gamma(\theta)\colon
B=\Gamma(Y,\OO_Y)\longrightarrow
\Gamma(Y,\psi_*(\OO_X))=\Gamma(X,\OO_X)=A.
\]
关于 $\theta_x^\sharp$ 的假设使我们可以取商，得到剩余域之间的单同态
$\theta^x\colon k(\psi(x))\to k(x)$，并且对于每个截面 $f\in\Gamma(Y,\OO_Y)=B$，有
\[
\theta^x(f(\psi(x)))=\varphi(f)(x).
\]
因此，关系 $f(\psi(x))=0$ 等价于 $\varphi(f)(x)=0$。
这表示 $\jj_{\psi(x)}=\jj_{\aphi(x)}$，也可写成
$\psi(x)=\aphi(x)$ 对所有 $x\in X$ 成立，即 $\psi=\aphi$。
我们还知道下图交换（0，\hyperlink{para:0:3.7.1}{3.7.1}）：
\[
\begin{tikzcd}[column sep=large,row sep=large]
B=\Gamma(Y,\OO_Y)\arrow[r,"\varphi"]\arrow[d]
 & \Gamma(X,\OO_X)=A\arrow[d]\\
B_{\psi(x)}\arrow[r,"\theta_x^\sharp"'] & A_x.
\end{tikzcd}
\]
这表示 $\theta_x^\sharp$ 等于由 $\varphi$ 典范得到的同态
$\varphi_x\colon B_{\aphi(x)}\to A_x$（0，\hyperlink{para:0:1.5.1}{1.5.1}）。
由于各 $\theta_x^\sharp$ 完全刻画 $\theta^\sharp$，从而也完全刻画 $\theta$（0，\hyperlink{para:0:3.7.1}{3.7.1}），
按 $\wt\varphi$ 的定义（\hyperlink{para:1:1.6.1}{1.6.1}），便得到 $\theta=\wt\varphi$。

满足（\hyperlink{para:1:1.7.3}{1.7.3}）的条件的赋环空间态射 $(\psi,\theta)$，
称为\emph{仿射概形态射}。

\egastatement{1.7.4}{推论}
设 $(X,\OO_X)$、$(Y,\OO_Y)$ 是两个仿射概形，
$A=\Gamma(\OO_X)$，$B=\Gamma(\OO_Y)$。
则仿射概形态射集合 $\Hom((X,\OO_X),(Y,\OO_Y))$
与从 $B$ 到 $A$ 的环同态集合之间存在典范同构。

也可以说，关于 $A$ 的函子 $(\Spec(A),\wt A)$ 与关于 $(X,\OO_X)$ 的函子
$\Gamma(X,\OO_X)$，给出交换环范畴与仿射概形范畴的对偶范畴之间的等价（T，I，1.2）。

\egastatement{1.7.5}{推论}
若 $\varphi\colon B\to A$ 满射，则对应态射 $(\aphi,\wt\varphi)$
是赋环空间的单态射（参见（\hyperlink{para:1:4.1.7}{4.1.7}））。

事实上，我们知道 $\aphi$ 单射（\hyperlink{para:1:1.2.5}{1.2.5}）。又因为 $\varphi$ 满射，
对于每个 $x\in X$，由 $\varphi$ 取分式环得到的同态
$\wt\varphi_x^\sharp\colon B_{\aphi(x)}\to A_x$ 也满射（0，\hyperlink{para:0:1.5.1}{1.5.1}）。
由此得到结论（0，\hyperlink{para:0:4.1.1}{4.1.1}）。
\par


% END chapters/ch1_s1_7.tex

\par

% BEGIN chapters/ch1_s2_1.tex
\section{预概形与预概形态射}\hypertarget{heading:1:2}{}
\subsection{预概形的定义}\hypertarget{heading:1:2.1}{}
\src{97—98}

\egapar{2.1.1}
给定一个赋环空间 $(X,\OO_X)$。若 $X$ 的开子集 $V$ 所对应的赋环空间
$(V,\OO_X|V)$ 是仿射概形（\hyperlink{para:1:1.7.1}{1.7.1}），则称 $V$ 为一个\emph{仿射开集}。

\egastatement{2.1.2}{定义}
所谓\emph{预概形}，是指一个赋环空间 $(X,\OO_X)$，其中 $X$ 的每一点都具有仿射开邻域。

\egastatement{2.1.3}{命题}
若 $(X,\OO_X)$ 是预概形，则仿射开集构成 $X$ 的拓扑的一组基。

事实上，设 $V$ 为 $x\in X$ 的任意开邻域。依假设，$x$ 具有一个开邻域 $W$，
使得 $(W,\OO_X|W)$ 是仿射概形；以 $A$ 表示它的环。
在空间 $W$ 中，$V\cap W$ 是 $x$ 的开邻域，故存在 $f\in A$，使得
$D(f)$ 是包含在 $V\cap W$ 中的 $x$ 的开邻域（\hyperlink{para:1:1.1.10}{1.1.10}，(i)）。
赋环空间 $(D(f),\OO_X|D(f))$ 是环同构于 $A_f$ 的仿射概形（\hyperlink{para:1:1.3.6}{1.3.6}），命题得证。

\egastatement{2.1.4}{命题}
预概形的底空间是 Kolmogoroff 空间。

事实上，若 $x,y$ 是预概形 $X$ 的两个不同的点，且它们不属于同一个仿射开集，
则显然存在其中一点的一个开邻域不含另一点；若它们属于同一个仿射开集，
则由（\hyperlink{para:1:1.1.8}{1.1.8}）得到同一结论。

\egastatement{2.1.5}{命题}
若 $(X,\OO_X)$ 是预概形，则 $X$ 的每个不可约闭子集都有且只有一个泛点；
因此，映射 $x\mapsto\overline{\{x\}}$ 是 $X$ 到其不可约闭子集的集合的双射。

事实上，设 $Y$ 为 $X$ 的不可约闭子集，$y\in Y$，$U$ 为 $y$ 在 $X$ 中的一个仿射开邻域。
则 $U\cap Y$ 在 $Y$ 中稠密且不可约（0，\hyperlink{para:0:2.1.1}{2.1.1} 及 \hyperlink{para:0:2.1.4}{2.1.4}）。
由（\hyperlink{para:1:1.1.14}{1.1.14}），$U\cap Y$ 是某一点 $x$ 在 $U$ 中的闭包，因而 $Y$ 是 $x$ 在 $X$ 中的闭包。
泛点的唯一性由（\hyperlink{para:1:2.1.4}{2.1.4}）及（0，\hyperlink{para:0:2.1.3}{2.1.3}）得到。

\egapar{2.1.6}
若 $Y$ 是 $X$ 的不可约闭子集，$y$ 是它的泛点，则局部环 $\OO_y$ 也记为
$\OO_{X/Y}$，称为\emph{$X$ 沿 $Y$ 的局部环}，或\emph{$Y$ 在 $X$ 中的局部环}。

若 $X$ 本身不可约，且 $x$ 为其泛点，则也称 $\OO_x$ 为\emph{$X$ 上的有理函数环}（参见 §7）。

\egastatement{2.1.7}{命题}
若 $(X,\OO_X)$ 是预概形，则对于 $X$ 的任意开子集 $U$，
赋环空间 $(U,\OO_X|U)$ 是预概形。

这直接由定义（\hyperlink{para:1:2.1.2}{2.1.2}）及命题（\hyperlink{para:1:2.1.3}{2.1.3}）得到。

称 $(U,\OO_X|U)$ 为 $(X,\OO_X)$ \emph{在 $U$ 上诱导的预概形}，
或 $(X,\OO_X)$ \emph{在 $U$ 上的限制}。

\egapar{2.1.8}
若预概形 $(X,\OO_X)$ 的底空间 $X$ 不可约（分别为连通），则称该预概形
\emph{不可约}（分别为\emph{连通}）。称一个预概形为\emph{整的}，
若它不可约且约化（参见（\hyperlink{para:1:5.1.4}{5.1.4}））。称预概形 $(X,\OO_X)$ \emph{局部整}，
若每个 $x\in X$ 都有一个开邻域 $U$，使得 $(X,\OO_X)$ 在 $U$ 上诱导的预概形是整的。


% END chapters/ch1_s2_1.tex

\par

% BEGIN chapters/ch1_s2_2.tex
\subsection{预概形态射}\hypertarget{heading:1:2.2}{}
\src{98—100}

\egastatement{2.2.1}{定义}
给定两个预概形 $(X,\OO_X)$、$(Y,\OO_Y)$。
所谓从 $(X,\OO_X)$ 到 $(Y,\OO_Y)$ 的\emph{（预概形）态射}，
是指一个赋环空间态射 $(\psi,\theta)$，使得对于每个 $x\in X$，
$\theta_x^\sharp$ 都是局部同态 $\OO_{\psi(x)}\to\OO_x$。

取商以后，同态 $\theta_x^\sharp\colon\OO_{\psi(x)}\to\OO_x$
给出单同态 $\theta^x\colon k(\psi(x))\to k(x)$；由此可将 $k(x)$ 看成域 $k(\psi(x))$ 的扩张。

\egapar{2.2.2}
两个预概形态射 $(\psi,\theta)$、$(\psi',\theta')$ 的复合 $(\psi'',\theta'')$
仍为预概形态射，这由公式
$\theta^{\prime\prime\sharp}=\theta^\sharp\circ\psi^*(\theta^{\prime\sharp})$
（0，\hyperlink{para:0:3.5.5}{3.5.5}）得到。因此，预概形构成一个范畴。依照一般的记号，
以 $\Hom(X,Y)$ 表示从预概形 $X$ 到预概形 $Y$ 的态射集合。

\egastatement{2.2.3}{例}
若 $U$ 是 $X$ 的开子集，则诱导预概形 $(U,\OO_X|U)$ 到 $(X,\OO_X)$ 的
典范嵌入（0，\hyperlink{para:0:4.1.2}{4.1.2}）是预概形态射。而且它是赋环空间的单态射，
因而也是预概形的单态射；这直接由（0，\hyperlink{para:0:4.1.1}{4.1.1}）得到。

\egastatement{2.2.4}{命题}
设 $(X,\OO_X)$ 是预概形，$(S,\OO_S)$ 是环为 $A$ 的仿射概形。
则从预概形 $(X,\OO_X)$ 到预概形 $(S,\OO_S)$ 的态射，
与从 $A$ 到环 $\Gamma(X,\OO_X)$ 的同态之间存在典范的一一对应。

首先指出，若 $(X,\OO_X)$、$(Y,\OO_Y)$ 是任意两个赋环空间，
则从前者到后者的态射 $(\psi,\theta)$ 典范地确定环同态
\[
\Gamma(\theta)\colon\Gamma(Y,\OO_Y)
\longrightarrow\Gamma(Y,\psi_*(\OO_X))=\Gamma(X,\OO_X).
\]
在所考虑的情形中，只须证明：任意同态 $\varphi\colon A\to\Gamma(X,\OO_X)$
都能写成 $\Gamma(\theta)$，且相应的态射唯一。
依假设，$X$ 有一个由仿射开集构成的覆盖 $(V_\alpha)$。
将 $\varphi$ 与限制同态
$\Gamma(X,\OO_X)\to\Gamma(V_\alpha,\OO_X|V_\alpha)$ 复合，得到同态
\[
\varphi_\alpha\colon A\longrightarrow\Gamma(V_\alpha,\OO_X|V_\alpha).
\]
由（\hyperlink{para:1:1.7.3}{1.7.3}），它对应于唯一的态射
$(\psi_\alpha,\theta_\alpha)\colon(V_\alpha,\OO_X|V_\alpha)\to(S,\OO_S)$。
此外，对于任意一对指标 $(\alpha,\beta)$，$V_\alpha\cap V_\beta$ 的每一点都有
一个包含在 $V_\alpha\cap V_\beta$ 中的仿射开邻域 $W$（\hyperlink{para:1:2.1.3}{2.1.3}）。
显然，将 $\varphi_\alpha$、$\varphi_\beta$ 分别与到 $W$ 的限制同态复合，
得到的是同一个同态 $\Gamma(S,\OO_S)\to\Gamma(W,\OO_X|W)$。
因而，由对任意 $\alpha$ 及 $x\in V_\alpha$ 都成立的关系
$(\theta_\alpha^\sharp)_x=(\varphi_\alpha)_x$（\hyperlink{para:1:1.6.1}{1.6.1}），
态射 $(\psi_\alpha,\theta_\alpha)$ 与 $(\psi_\beta,\theta_\beta)$ 在 $W$ 上的限制相同。
于是，存在唯一的赋环空间态射
$(\psi,\theta)\colon(X,\OO_X)\to(S,\OO_S)$，
其在每个 $V_\alpha$ 上的限制为 $(\psi_\alpha,\theta_\alpha)$。
显然，它是预概形态射，且 $\Gamma(\theta)=\varphi$。

设 $u\colon A\to\Gamma(X,\OO_X)$ 是环同态，
$\varphi=(\psi,\theta)\colon(X,\OO_X)\to(S,\OO_S)$ 是对应态射。
则对于每个 $f\in A$，有
\begin{equation}
\varphi^{-1}(D(f))=X_{u(f)}.\tag{2.2.4.1}\label{eq:1:2.2.4.1}
\end{equation}
这里使用的是（0，\hyperlink{para:0:5.5.2}{5.5.2}）中关于局部自由层 $\OO_X$ 的记号。
事实上，只须在 $X$ 本身为仿射概形时验证此式，而这时它就是（\hyperref[eq:1:1.2.2.2]{1.2.2.2}）。

\egastatement{2.2.5}{命题}
在（\hyperlink{para:1:2.2.4}{2.2.4}）的假设下，设 $\varphi\colon A\to\Gamma(X,\OO_X)$ 是环同态，
$f\colon(X,\OO_X)\to(S,\OO_S)$ 是对应的预概形态射，
$\F$、$\G$ 分别是拟凝聚的 $\OO_S$-模层与 $\OO_X$-模层，并令 $M=\Gamma(S,\F)$。
则 $f$-态射 $\F\to\G$（0，\hyperlink{para:0:4.4.1}{4.4.1}）与 $A$-同态
$M\to\restmod{\Gamma(X,\G)}{\varphi}$ 之间存在典范的一一对应。

事实上，仿照（\hyperlink{para:1:2.2.4}{2.2.4}）中的论证，即可归结到 $X$ 为仿射概形的情形；
这时命题由（\hyperlink{para:1:1.6.3}{1.6.3}）及（\hyperlink{para:1:1.3.8}{1.3.8}）得到。

\egapar{2.2.6}
称预概形态射 $(\psi,\theta)\colon(X,\OO_X)\to(Y,\OO_Y)$ 为\emph{开的}
（分别为\emph{闭的}），若对于 $X$ 的每个开子集 $U$（分别为闭子集 $F$），
$\psi(U)$ 在 $Y$ 中开（分别为 $\psi(F)$ 在 $Y$ 中闭）。
若 $\psi(X)$ 在 $Y$ 中稠密，则称 $(\psi,\theta)$ 为\emph{支配态射}；
若 $\psi$ 满射，则称它为\emph{满射态射}。注意，这些条件只涉及连续映射 $\psi$。

\egastatement{2.2.7}{命题}
设
\[
f=(\psi,\theta)\colon(X,\OO_X)\to(Y,\OO_Y),\qquad
 g=(\psi',\theta')\colon(Y,\OO_Y)\to(Z,\OO_Z)
\]
是两个预概形态射。
\begin{enumerate}[label=(\roman*)]
\item 若 $f,g$ 都是开的（分别为闭的、支配的、满射的），则 $g\circ f$ 也具有相应性质。
\item 若 $f$ 满射，且 $g\circ f$ 是闭的，则 $g$ 是闭的。
\item 若 $g\circ f$ 满射，则 $g$ 满射。
\end{enumerate}
断言 (i)、(iii) 是显然的。令 $g\circ f=(\psi'',\theta'')$。
若 $F$ 在 $Y$ 中闭，则 $\psi^{-1}(F)$ 在 $X$ 中闭，故 $\psi''(\psi^{-1}(F))$ 在 $Z$ 中闭。
但由于 $\psi$ 满射，有 $\psi(\psi^{-1}(F))=F$，从而
$\psi''(\psi^{-1}(F))=\psi'(F)$，这就证明了 (ii)。

\egastatement{2.2.8}{命题}
设 $f=(\psi,\theta)$ 为态射 $(X,\OO_X)\to(Y,\OO_Y)$，$(U_\alpha)$ 是 $Y$ 的开覆盖。
则 $f$ 是开的（分别为闭的、满射的、支配的），当且仅当它在每个诱导预概形
$(\psi^{-1}(U_\alpha),\OO_X|\psi^{-1}(U_\alpha))$ 上的限制，
视为从此诱导预概形到诱导预概形 $(U_\alpha,\OO_Y|U_\alpha)$ 的态射时，
是开的（分别为闭的、满射的、支配的）。

此命题直接由定义得到，只须注意：$Y$ 的子集 $F$ 在 $Y$ 中闭（分别为开、稠密），
当且仅当每个 $F\cap U_\alpha$ 在 $U_\alpha$ 中闭（分别为开、稠密）。

\egapar{2.2.9}
设 $(X,\OO_X)$、$(Y,\OO_Y)$ 为两个预概形，并假设 $X,Y$ 具有相同的有限个数的
不可约分支，分别记为 $X_i,Y_i$（$1\leq i\leq n$）；
以 $\xi_i,\eta_i$ 分别表示 $X_i,Y_i$ 的泛点（\hyperlink{para:1:2.1.5}{2.1.5}）。
若态射
\[
f=(\psi,\theta)\colon(X,\OO_X)\longrightarrow(Y,\OO_Y)
\]
满足：对每个 $i$，有 $\psi(\xi_i)=\eta_i$，且
$\theta_{\xi_i}^\sharp\colon\OO_{\eta_i}\to\OO_{\xi_i}$ 是同构，
则称 $f$ 为\emph{双有理态射}。显然，双有理态射是支配的（0，\hyperlink{para:0:2.1.8}{2.1.8}），
因而，若它还是闭的，则它满射。

\egastatement{2.2.10}{记号约定}
在本书以下各处，只要不致混淆，我们将在预概形（分别为态射）的记号中省略
结构层（分别为结构层之间的态射）。若 $U$ 是预概形 $X$ 的底空间的开子集，
则当把 $U$ 作为预概形来谈论时，始终是指在 $U$ 上诱导的预概形。


% END chapters/ch1_s2_2.tex

\par

% BEGIN chapters/ch1_s2_3.tex
\subsection{预概形的粘合}\hypertarget{heading:1:2.3}{}
\src{101}

\egapar{2.3.1}
由定义（\hyperlink{para:1:2.1.2}{2.1.2}）可知，由预概形粘合而得的任意赋环空间（0，\hyperlink{para:0:4.1.6}{4.1.6}）仍然是预概形。
\transnote{此处原文引作（0，\hyperlink{para:0:4.1.6}{4.1.6}）；本版关于赋环空间粘合的条目实际为（0，\hyperlink{para:0:4.1.7}{4.1.7}）。}
特别地，由于每个预概形依定义都有由仿射开集构成的覆盖，
可见每个预概形都能由仿射概形粘合而得。

\egastatement{2.3.2}{例}
设 $K$ 是域，$B=K[s]$、$C=K[t]$ 是 $K$ 上的两个一元多项式环；
令 $X_1=\Spec(B)$、$X_2=\Spec(C)$，它们是两个同构的仿射概形。
在 $X_1$（分别在 $X_2$）中，令 $U_{12}$（分别令 $U_{21}$）为仿射开集
$D(s)$（分别为 $D(t)$）。其环 $B_s$（分别为 $C_t$）由形如
$f(s)/s^n$（分别为 $g(t)/t^n$）的有理分式组成，其中 $f\in B$（分别为 $g\in C$）。
设 $u_{12}\colon U_{21}\to U_{12}$ 为与环同构 $B_s\to C_t$ 对应的预概形同构（\hyperlink{para:1:2.2.4}{2.2.4}），
此环同构把 $f(s)/s^n$ 映到有理分式 $f(1/t)/(1/t)^n$。
可以借助 $u_{12}$ 沿 $U_{12}$、$U_{21}$ 将 $X_1$ 与 $X_2$ 粘合，
因为这里显然没有另需满足的粘合条件。

以后，我们将把如此得到的预概形 $X$ 看成一种一般构造方法的特殊情形（II，2.4.3）。
这里仅证明 $X$ 不是仿射概形；这将由环 $\Gamma(X,\OO_X)$ 同构于 $K$、
因而其素谱仅有一点而得到。事实上，$\OO_X$ 的一个全局截面在 $X_1$
（分别在 $X_2$）上的限制，是一个多项式 $f(s)$（分别为 $g(t)$）；
这里已将 $X_1,X_2$ 等同于 $X$ 中相应的仿射开集。
由定义可知，必须有 $g(t)=f(1/t)$，而这只有在 $f=g\in K$ 时才可能。


% END chapters/ch1_s2_3.tex

\par

% BEGIN chapters/ch1_s2_4.tex
\subsection{局部概形}\hypertarget{heading:1:2.4}{}
\src{101—103}

\egapar{2.4.1}
所谓\emph{局部概形}，是指其环 $A$ 为局部环的仿射概形；
这时 $X=\Spec(A)$ 中有且只有一个闭点 $a$，且对于任意其他点 $b\in X$，
有 $a\in\overline{\{b\}}$（\hyperlink{para:1:1.1.7}{1.1.7}）。

对于任意预概形 $Y$ 及点 $y\in Y$，称局部概形 $\Spec(\OO_y)$ 为
\emph{$Y$ 在点 $y$ 的局部概形}。
设 $V$ 为 $Y$ 中含 $y$ 的仿射开集，$B$ 为仿射概形 $V$ 的环。
$\OO_y$ 与 $B_y$ 典范地等同（\hyperlink{para:1:1.3.4}{1.3.4}），故典范同态 $B\to B_y$ 对应于一个
预概形态射 $\Spec(\OO_y)\to V$（\hyperlink{para:1:1.6.1}{1.6.1}）。
将此态射与典范嵌入 $V\to Y$ 复合，便得到态射 $\Spec(\OO_y)\to Y$；
它不依赖于所选的含 $y$ 的仿射开集 $V$。
事实上，若 $V'$ 为第二个含 $y$ 的仿射开集，则存在第三个含 $y$ 的仿射开集 $W$，
满足 $W\subset V\cap V'$（\hyperlink{para:1:2.1.3}{2.1.3}）；因而可归结到 $V\subset V'$ 的情形。
若 $B'$ 为 $V'$ 的环，则只须注意下图交换（0，\hyperlink{para:0:1.5.1}{1.5.1}）：
\[
\begin{tikzcd}[column sep=large,row sep=large]
 B'\arrow[r]\arrow[dr] & B\arrow[d]\\
 & \OO_y.
\end{tikzcd}
\]
如此定义的态射 $\Spec(\OO_y)\to Y$ 称为\emph{典范态射}。

\egastatement{2.4.2}{命题}
设 $(Y,\OO_Y)$ 为预概形。对于每个 $y\in Y$，令 $(\psi,\theta)$ 为典范态射
\[
(\Spec(\OO_y),\wt{\OO_y})\longrightarrow(Y,\OO_Y).
\]
则 $\psi$ 是从 $\Spec(\OO_y)$ 到 $Y$ 的子空间 $S_y$ 的同胚，
其中 $S_y$ 由满足 $y\in\overline{\{z\}}$ 的点 $z$ 组成，
即由 $y$ 的一般化点组成（0，\hyperlink{para:0:2.1.2}{2.1.2}）。
此外，若 $z=\psi(\mathfrak p)$，则相应的茎同态
$\OO_z\to(\OO_y)_{\mathfrak p}$ 是同构；因此，$(\psi,\theta)$ 是赋环空间的单态射。

由于 $\Spec(\OO_y)$ 的唯一闭点 $a$ 属于此空间每一点的闭包，且 $\psi(a)=y$，
连续映射 $\psi$ 下 $\Spec(\OO_y)$ 的像包含于 $S_y$。
由于 $S_y$ 包含于每个含 $y$ 的仿射开集，可以归结到 $Y$ 为仿射概形的情形；
而这时命题由（\hyperlink{para:1:1.6.2}{1.6.2}）得到。

因此可见（\hyperlink{para:1:2.1.5}{2.1.5}），$\Spec(\OO_y)$ 与 $Y$ 中含 $y$ 的不可约闭子集的集合之间
存在一一对应。

\egastatement{2.4.3}{推论}
点 $y\in Y$ 是 $Y$ 的某个不可约分支的泛点，当且仅当局部环 $\OO_y$ 的
唯一素理想是它的极大理想，换言之，当且仅当 $\OO_y$ 的维数为零。

\egastatement{2.4.4}{命题}
设 $(X,\OO_X)$ 是环为 $A$ 的局部概形，$a$ 是它的唯一闭点，$(Y,\OO_Y)$ 为预概形。
每个态射 $u=(\psi,\theta)\colon(X,\OO_X)\to(Y,\OO_Y)$ 都唯一地分解为
\[
 X\longrightarrow\Spec(\OO_{\psi(a)})\longrightarrow Y,
\]
其中第二个箭头为典范态射，第一个箭头对应于一个局部同态 $\OO_{\psi(a)}\to A$。
这就在态射 $(X,\OO_X)\to(Y,\OO_Y)$ 的集合与局部同态
$\OO_y\to A$（$y\in Y$）的集合之间建立了典范的一一对应。

事实上，对于每个 $x\in X$，有 $a\in\overline{\{x\}}$，
因而 $\psi(a)\in\overline{\{\psi(x)\}}$。
这说明 $\psi(X)$ 包含于每个含 $\psi(a)$ 的仿射开集。
因此可以归结到 $(Y,\OO_Y)$ 为环是 $B$ 的仿射概形的情形；
这时 $u=(\aphi,\wt\varphi)$，其中 $\varphi\in\Hom(B,A)$（\hyperlink{para:1:1.7.3}{1.7.3}）。
此外，$\varphi^{-1}(\jj_a)=\jj_{\psi(a)}$，故 $B\setminus\jj_{\psi(a)}$ 中
每个元素在 $\varphi$ 下的像都在局部环 $A$ 中可逆。
命题所述的分解于是由分式环的泛性质（0，\hyperlink{para:0:1.2.4}{1.2.4}）得到。
反过来，每个局部同态 $\OO_y\to A$ 都对应于唯一的态射
$(\psi',\theta')\colon X\to\Spec(\OO_y)$，使得 $\psi'(a)$ 为对应于 $y$ 的闭点（\hyperlink{para:1:1.7.3}{1.7.3}）；
将它与典范态射 $\Spec(\OO_y)\to Y$ 复合，便得到态射 $X\to Y$，命题得证。

\egapar{2.4.5}
环为域 $K$ 的仿射概形，其底空间只有一点。
若 $A$ 为极大理想是 $\mathfrak m$ 的局部环，则每个局部同态 $A\to K$ 的核都等于
$\mathfrak m$，因而分解为 $A\to A/\mathfrak m\to K$，其中第二个箭头是单同态。
因此，态射 $\Spec(K)\to\Spec(A)$ 与域单同态 $A/\mathfrak m\to K$ 一一对应。
\transnote{上一句按原文保留，但其适用范围须与前一句的“局部同态”一致：
它描述的是将 $\Spec(K)$ 的唯一点映到 $\Spec(A)$ 的闭点的态射。
对于所有态射 $\Spec(K)\to\Spec(A)$，还必须允许该点映到 $\Spec(A)$ 的其他点；
完整的对应关系是下述（\hyperlink{para:1:2.4.6}{2.4.6}）中的 $k(y)\to K$，其中 $y$ 遍历目标的各点。
不能仅由 $A$ 是局部环就省去这一选择。}

设 $(Y,\OO_Y)$ 是预概形。对于每个 $y\in Y$ 及 $\OO_y$ 的理想 $\mathfrak a_y$，
典范同态 $\OO_y\to\OO_y/\mathfrak a_y$ 定义了一个态射
$\Spec(\OO_y/\mathfrak a_y)\to\Spec(\OO_y)$；
将其与典范态射 $\Spec(\OO_y)\to Y$ 复合，得到态射
$\Spec(\OO_y/\mathfrak a_y)\to Y$，也称为\emph{典范态射}。
取 $\mathfrak a_y=\mathfrak m_y$，即 $\OO_y/\mathfrak a_y=k(y)$，
命题（\hyperlink{para:1:2.4.4}{2.4.4}）遂给出：

\egastatement{2.4.6}{推论}
设 $(X,\OO_X)$ 是环为域 $K$ 的局部概形，$\xi$ 是 $X$ 的唯一点，$(Y,\OO_Y)$ 是预概形。
每个态射 $u\colon(X,\OO_X)\to(Y,\OO_Y)$ 都唯一地分解为
\[
 X\longrightarrow\Spec(k(u(\xi)))\longrightarrow Y,
\]
其中第二个箭头是典范态射，第一个箭头对应于单同态 $k(u(\xi))\to K$。
这在态射 $(X,\OO_X)\to(Y,\OO_Y)$ 的集合与单同态 $k(y)\to K$（$y\in Y$）
的集合之间建立了典范的一一对应。

\egastatement{2.4.7}{推论}
对于每个 $y\in Y$，典范态射 $\Spec(\OO_y/\mathfrak a_y)\to Y$ 是赋环空间的单态射。

在 $\mathfrak a_y=0$ 时，我们已知此结论（\hyperlink{para:1:2.4.2}{2.4.2}）；再应用（\hyperlink{para:1:1.7.5}{1.7.5}）即可。

\egastatement{2.4.8}{注}
设 $X$ 是局部概形，$a$ 是其唯一闭点。
由于每个含 $a$ 的仿射开集必为整个 $X$，每个可逆的 $\OO_X$-模层（0，\hyperlink{para:0:5.4.1}{5.4.1}）
必同构于 $\OO_X$，也称它为\emph{平凡的}。
对于一般的仿射概形 $\Spec(A)$，此性质并不一定成立；
我们将在第 V 章看到，若 $A$ 是正规环，则当 $A$ 为析因环时，此性质成立。


% END chapters/ch1_s2_4.tex

\par

% BEGIN chapters/ch1_s2_5.tex
\subsection{一个预概形上的预概形}\hypertarget{heading:1:2.5}{}
\src{103—104}

\egastatement{2.5.1}{定义}
给定一个预概形 $S$。给出一个预概形 $X$ 及一个预概形态射
$\varphi\colon X\to S$，就称为给出了一个\emph{预概形 $S$ 上的预概形 $X$}，
或一个\emph{$S$-预概形}；称 $S$ 为 $S$-预概形 $X$ 的\emph{基预概形}，
称态射 $\varphi$ 为 $S$-预概形 $X$ 的\emph{结构态射}。
当 $S$ 是环为 $A$ 的仿射概形时，也称赋予 $\varphi$ 的 $X$ 为
\emph{环 $A$ 上的预概形}，或\emph{$A$-预概形}。

由（\hyperlink{para:1:2.2.4}{2.2.4}）可知，给出一个环 $A$ 上的预概形，等价于给出一个预概形
$(X,\OO_X)$，其结构层 $\OO_X$ 为一个 $A$-代数层。
因此，\emph{任意预概形都能以唯一的方式看成一个 $\mathbf Z$-预概形}。

若 $\varphi\colon X\to S$ 是 $S$-预概形 $X$ 的结构态射，
则当 $\varphi(x)=s$ 时，称点 $x\in X$ \emph{位于点 $s\in S$ 之上}。
若 $\varphi$ 是支配态射（\hyperlink{para:1:2.2.6}{2.2.6}），则称 $X$ \emph{支配} $S$。

\egapar{2.5.2}
设 $X,Y$ 是两个 $S$-预概形。
若预概形态射 $u\colon X\to Y$ 使下图交换：
\[
\begin{tikzcd}[column sep=large,row sep=large]
X\arrow[r,"u"]\arrow[dr] & Y\arrow[d]\\
& S,
\end{tikzcd}
\]
其中斜向及向下的箭头为结构态射，则称 $u$ 为一个\emph{$S$ 上的预概形态射}，
或一个\emph{$S$-态射}。这蕴含：对于每个 $s\in S$ 及位于 $s$ 之上的 $x\in X$，
$u(x)$ 也位于 $s$ 之上。

由此定义立即可见，两个 $S$-态射的复合仍为 $S$-态射，故 $S$-预概形构成一个范畴。
以 $\Hom_S(X,Y)$ 表示从 $S$-预概形 $X$ 到 $S$-预概形 $Y$ 的 $S$-态射集合；
$S$-预概形 $X$ 的恒等态射记为 $1_X$。
当 $S$ 是环为 $A$ 的仿射概形时，也称 $S$-态射为\emph{$A$-态射}。

\egapar{2.5.3}
若 $X$ 是 $S$-预概形，$v\colon X'\to X$ 是预概形态射，
则复合态射 $X'\to X\to S$ 使 $X'$ 成为 $S$-预概形。
特别地，借助典范嵌入，可以把在 $X$ 的任意开集 $U$ 上诱导的预概形看成 $S$-预概形。

因此，若 $u\colon X\to Y$ 是 $S$-预概形之间的 $S$-态射，
则 $u$ 在 $X$ 的任意开集 $U$ 上诱导的预概形上的限制，是一个 $S$-态射 $U\to Y$。
反过来，设 $(U_\alpha)$ 为 $X$ 的开覆盖，并对每个 $\alpha$ 给出一个
$S$-态射 $u_\alpha\colon U_\alpha\to Y$。
若对于每对指标 $(\alpha,\beta)$，$u_\alpha$ 与 $u_\beta$ 在
$U_\alpha\cap U_\beta$ 上的限制相同，则存在唯一的 $S$-态射 $X\to Y$，
其在每个 $U_\alpha$ 上的限制为 $u_\alpha$。

若 $S$-态射 $u\colon X\to Y$ 满足 $u(X)\subset V$，其中 $V$ 是 $Y$ 的开子集，
则把 $u$ 看成从 $X$ 到 $V$ 的态射时，它仍为 $S$-态射。

\egapar{2.5.4}
设 $S'\to S$ 是预概形态射。
对于每个 $S'$-预概形 $X$，复合态射 $X\to S'\to S$ 使 $X$ 成为 $S$-预概形。
反过来，设 $S'$ 是在 $S$ 的一个开集上诱导的预概形；
若 $X$ 是 $S$-预概形，且其结构态射 $f\colon X\to S$ 满足 $f(X)\subset S'$，
则可以把 $X$ 看成 $S'$-预概形。
在后一情形中，若 $Y$ 是一个 $S$-预概形，其结构态射也将底空间映入 $S'$，
则从 $X$ 到 $Y$ 的每个 $S$-态射也都是 $S'$-态射。

\egapar{2.5.5}
若 $X$ 是 $S$-预概形，$\varphi\colon X\to S$ 是结构态射，
则称从 $S$ 到 $X$ 的一个 $S$-态射为 $X$ 的一个\emph{$S$-截面}，
即一个满足 $\varphi\circ\psi=1_S$ 的预概形态射 $\psi\colon S\to X$。
以 $\Gamma(X/S)$ 表示 $X$ 的 $S$-截面的集合。
\transnote{本条首句原文将 $X$ 排作“un $S$-morphisme”；依（\hyperlink{para:1:2.5.1}{2.5.1}）及本句的结构态射，
应为“un $S$-préschéma”，此处明确按后者译出。}


% END chapters/ch1_s2_5.tex

\par

% BEGIN chapters/ch1_s3_1.tex
\section{预概形的积}\hypertarget{heading:1:3}{}
\subsection{预概形的和}\hypertarget{heading:1:3.1}{}
\src{104}

设 $(X_\alpha)$ 是任意一族预概形，$X$ 是其底空间 $X_\alpha$ 的拓扑和。
于是 $X$ 是一族两两不交的开子空间 $X'_\alpha$ 的并，且对于每个 $\alpha$，
存在从 $X_\alpha$ 到 $X'_\alpha$ 的同胚 $\varphi_\alpha$。
在每个 $X'_\alpha$ 上赋予层 $(\varphi_\alpha)_*(\OO_{X_\alpha})$，
显然便使 $X$ 成为预概形；称其为预概形族 $(X_\alpha)$ 的\emph{和}，记为
$\coprod_\alpha X_\alpha$。
若 $Y$ 为预概形，则映射 $f\mapsto(f\circ\varphi_\alpha)_\alpha$ 是从
$\Hom(X,Y)$ 到积集合 $\prod_\alpha\Hom(X_\alpha,Y)$ 的双射。
特别地，若各 $X_\alpha$ 是结构态射为 $\psi_\alpha$ 的 $S$-预概形，
则唯一的态射 $\psi\colon X\to S$，满足对每个 $\alpha$ 有
$\psi\circ\varphi_\alpha=\psi_\alpha$，使 $X$ 成为 $S$-预概形。
两个预概形 $X,Y$ 的和记为 $X\amalg Y$。
若 $X=\Spec(A)$、$Y=\Spec(B)$，则 $X\amalg Y$ 与 $\Spec(A\times B)$ 典范地等同。


% END chapters/ch1_s3_1.tex

\par

% BEGIN chapters/ch1_s3_2.tex
\subsection{预概形的积}\hypertarget{heading:1:3.2}{}
\src{104—108}

\egastatement{3.2.1}{定义}
给定两个 $S$-预概形 $X,Y$。设三元组 $(Z,p_1,p_2)$ 由一个 $S$-预概形 $Z$
及两个 $S$-态射 $p_1\colon Z\to X$、$p_2\colon Z\to Y$ 组成。
若对于每个 $S$-预概形 $T$，映射 $f\mapsto(p_1\circ f,p_2\circ f)$ 是从
$T$ 到 $Z$ 的 $S$-态射集合，到由一个 $S$-态射 $T\to X$ 和一个 $S$-态射 $T\to Y$
组成的有序对的集合的双射，换言之，给出双射
\[
\Hom_S(T,Z)\xrightarrow{\sim}\Hom_S(T,X)\times\Hom_S(T,Y),
\]
则称 $(Z,p_1,p_2)$ 是 $S$-预概形 $X,Y$ 的一个\emph{积}。

这就是将范畴中两个对象的积的一般概念用于 $S$-预概形范畴（T，I，1.1）。
特别地，两个 $S$-预概形的积在唯一的 $S$-同构意义下唯一。
由于这一唯一性，通常用 $X\times_S Y$（不致混淆时简记为 $X\times Y$）
表示两个 $S$-预概形 $X,Y$ 的积，而在记号中省略态射 $p_1,p_2$；
它们分别称为从 $X\times_S Y$ 到 $X,Y$ 的\emph{典范投影}。
若 $g\colon T\to X$、$h\colon T\to Y$ 是两个 $S$-态射，
则以 $(g,h)_S$ 或简记 $(g,h)$ 表示满足 $p_1\circ f=g$、$p_2\circ f=h$ 的
$S$-态射 $f\colon T\to X\times_S Y$。
若 $X',Y'$ 是两个 $S$-预概形，$p'_1,p'_2$ 是假设存在的积 $X'\times_S Y'$ 的典范投影，
$u\colon X'\to X$、$v\colon Y'\to Y$ 为两个 $S$-态射，
则以 $u\times_S v$（或简记 $u\times v$）表示 $S$-态射
$(u\circ p'_1,v\circ p'_2)_S\colon X'\times_S Y'\to X\times_S Y$。

当 $S$ 是环为 $A$ 的仿射概形时，上述记号中的 $S$ 常被换成 $A$。

\egastatement{3.2.2}{命题}
设 $X,Y,S$ 是环分别为 $B,C,A$ 的三个仿射概形。
令 $Z=\Spec(B\otimes_A C)$，并令 $p_1,p_2$ 为与典范 $A$-同态
\[
u\colon B\longrightarrow B\otimes_A C,\quad b\longmapsto b\otimes1,
\qquad
v\colon C\longrightarrow B\otimes_A C,\quad c\longmapsto1\otimes c
\]
对应的 $S$-态射（\hyperlink{para:1:2.2.4}{2.2.4}）。则 $(Z,p_1,p_2)$ 是 $X,Y$ 的一个积。

由（\hyperlink{para:1:2.2.4}{2.2.4}），只须验证：若将每个 $A$-同态 $f\colon B\otimes_A C\to L$
（$L$ 为 $A$-代数）映到有序对 $(f\circ u,f\circ v)$，便定义了双射
\[
\Hom_A(B\otimes_A C,L)\xrightarrow{\sim}\Hom_A(B,L)\times\Hom_A(C,L).
\]
\origfootnote{这里的记号 $\Hom_A$ 表示 $A$-代数同态的集合。}
这直接由定义及关系 $b\otimes c=(b\otimes1)(1\otimes c)$ 得到。

\egastatement{3.2.3}{推论}
设 $T$ 是环为 $D$ 的仿射概形，
$\alpha=({}^a\rho,\wt\rho)$、$\beta=({}^a\sigma,\wt\sigma)$ 分别是
$S$-态射 $T\to X$、$T\to Y$，其中 $\rho\colon B\to D$、$\sigma\colon C\to D$
是 $A$-同态。则 $(\alpha,\beta)_S=({}^a\tau,\wt\tau)$，其中
$\tau\colon B\otimes_A C\to D$ 是满足 $\tau(b\otimes c)=\rho(b)\sigma(c)$ 的同态。

\egastatement{3.2.4}{命题}
设 $f\colon S'\to S$ 是预概形的单态射（T，I，1.1），
$X,Y$ 是两个 $S'$-预概形，并借助 $f$ 同时看成 $S$-预概形。
则 $S$-预概形 $X,Y$ 的每个积也是 $S'$-预概形 $X,Y$ 的积，反之亦然。

设 $\varphi\colon X\to S'$、$\psi\colon Y\to S'$ 为结构态射。
若 $T$ 是 $S$-预概形，$u\colon T\to X$、$v\colon T\to Y$ 是两个 $S$-态射，
则依定义有 $f\circ\varphi\circ u=f\circ\psi\circ v=\theta$，其中 $\theta$ 为 $T$ 的结构态射。
关于 $f$ 的假设蕴含 $\varphi\circ u=\psi\circ v=\theta'$。
可见，可以将 $T$ 看成结构态射为 $\theta'$ 的 $S'$-预概形，
并将 $u,v$ 看成 $S'$-态射。结合（\hyperlink{para:1:3.2.1}{3.2.1}），命题的结论立即得到。

\egastatement{3.2.5}{推论}
设 $X,Y$ 是两个 $S$-预概形，$\varphi\colon X\to S$、$\psi\colon Y\to S$ 是其结构态射。
若 $S'$ 是 $S$ 的开子集，且 $\varphi(X)\subset S'$、$\psi(Y)\subset S'$，
则 $S$-预概形 $X,Y$ 的每个积也是 $S'$-预概形 $X,Y$ 的积，反之亦然。

只须将（\hyperlink{para:1:3.2.4}{3.2.4}）用于典范嵌入 $S'\to S$。

\egastatement{3.2.6}{定理}
给定两个 $S$-预概形 $X,Y$，其积 $X\times_S Y$ 存在。

我们分几个步骤证明。

\egastatement{3.2.6.1}{引理}
设 $(Z,p,q)$ 是 $X,Y$ 的积，$U,V$ 分别为 $X,Y$ 的开子集。
令 $W=p^{-1}(U)\cap q^{-1}(V)$。
则由 $W$ 及 $p,q$ 在 $W$ 上的限制（分别视为态射 $W\to U$、$W\to V$）
组成的三元组，是 $U,V$ 的一个积。

事实上，若 $T$ 为 $S$-预概形，则可以将 $S$-态射 $T\to W$ 与将 $T$ 映入 $W$ 的
$S$-态射 $T\to Z$ 等同。
设 $g\colon T\to U$、$h\colon T\to V$ 是任意两个 $S$-态射；
可分别把它们看成从 $T$ 到 $X,Y$ 的 $S$-态射。
依假设，存在唯一的 $S$-态射 $f\colon T\to Z$，使得 $g=p\circ f$、$h=q\circ f$。
由于 $p(f(T))\subset U$、$q(f(T))\subset V$，有
$f(T)\subset p^{-1}(U)\cap q^{-1}(V)$，断言得证。

\egastatement{3.2.6.2}{引理}
设 $Z$ 为 $S$-预概形，$p\colon Z\to X$、$q\colon Z\to Y$ 是两个 $S$-态射，
$(U_\alpha)$、$(V_\lambda)$ 分别是 $X,Y$ 的开覆盖。
假设对于每对指标 $(\alpha,\lambda)$，$S$-预概形
$W_{\alpha\lambda}=p^{-1}(U_\alpha)\cap q^{-1}(V_\lambda)$
与 $p,q$ 在其上的限制构成 $U_\alpha,V_\lambda$ 的一个积。
则 $(Z,p,q)$ 是 $X,Y$ 的一个积。

先证明：若 $f_1,f_2$ 是两个 $S$-态射 $T\to Z$，则关系
$p\circ f_1=p\circ f_2$、$q\circ f_1=q\circ f_2$ 蕴含 $f_1=f_2$。
事实上，$Z$ 是各 $W_{\alpha\lambda}$ 的并，故各 $f_1^{-1}(W_{\alpha\lambda})$
构成 $T$ 的开覆盖，各 $f_2^{-1}(W_{\alpha\lambda})$ 也如此。
此外，依假设有
\[
\begin{aligned}
f_1^{-1}(W_{\alpha\lambda})
 &=f_1^{-1}(p^{-1}(U_\alpha))\cap f_1^{-1}(q^{-1}(V_\lambda))\\
 &=f_2^{-1}(p^{-1}(U_\alpha))\cap f_2^{-1}(q^{-1}(V_\lambda))
 =f_2^{-1}(W_{\alpha\lambda}).
\end{aligned}
\]
因此，只须证明对于每对指标，$f_1,f_2$ 在此共同开集上的限制相同。
这些限制可看成从 $f_1^{-1}(W_{\alpha\lambda})$ 到 $W_{\alpha\lambda}$ 的 $S$-态射，
故断言由假设与定义（\hyperlink{para:1:3.2.1}{3.2.1}）得到。

现在设已给出两个 $S$-态射 $g\colon T\to X$、$h\colon T\to Y$。
令 $T_{\alpha\lambda}=g^{-1}(U_\alpha)\cap h^{-1}(V_\lambda)$，
它们构成 $T$ 的开覆盖。依假设，存在 $S$-态射 $f_{\alpha\lambda}$，
使 $p\circ f_{\alpha\lambda}$、$q\circ f_{\alpha\lambda}$ 分别为
$g,h$ 在 $T_{\alpha\lambda}$ 上的限制。
再证明 $f_{\alpha\lambda}$、$f_{\beta\mu}$ 在
$T_{\alpha\lambda}\cap T_{\beta\mu}$ 上的限制相同，即可完成证明。
由定义，此交集在这两个态射下的像均包含于
$W_{\alpha\lambda}\cap W_{\beta\mu}$。由于
\[
W_{\alpha\lambda}\cap W_{\beta\mu}
=p^{-1}(U_\alpha\cap U_\beta)\cap q^{-1}(V_\lambda\cap V_\mu),
\]
由（\hyperlink{para:1:3.2.6.1}{3.2.6.1}），此预概形与 $p,q$ 的限制构成
$U_\alpha\cap U_\beta$、$V_\lambda\cap V_\mu$ 的积。
而 $p\circ f_{\alpha\lambda}$ 与 $p\circ f_{\beta\mu}$ 在
$T_{\alpha\lambda}\cap T_{\beta\mu}$ 上相同，对 $q$ 亦然，
故 $f_{\alpha\lambda}$ 与 $f_{\beta\mu}$ 在此交集上相同。证毕。

\egapar{3.2.6.3}
\emph{设 $(U_\alpha)$、$(V_\lambda)$ 分别为 $X,Y$ 的开覆盖，
并假设每对 $(\alpha,\lambda)$ 所对应的积 $U_\alpha\times_S V_\lambda$ 都存在。
则 $X,Y$ 的积存在。}

将引理（\hyperlink{para:1:3.2.6.1}{3.2.6.1}）用于开集 $U_\alpha\cap U_\beta$、$V_\lambda\cap V_\mu$，
可知分别由 $X,Y$ 在它们上面诱导的 $S$-预概形的积存在。
此外，令 $i=(\alpha,\lambda)$、$j=(\beta,\mu)$。
由积的唯一性，此积到 $U_\alpha\times_S V_\lambda$
（分别到 $U_\beta\times_S V_\mu$）的某个开子集上诱导的 $S$-预概形
$W_{ij}$（分别为 $W_{ji}$）存在典范同构 $h_{ij}$（分别为 $h_{ji}$）。
因而 $f_{ij}=h_{ij}\circ h_{ji}^{-1}$ 是 $W_{ji}$ 到 $W_{ij}$ 的同构。
对于第三对指标 $k=(\gamma,\nu)$，在 $W_{ki}\cap W_{kj}$ 上有
$f_{ik}=f_{ij}\circ f_{jk}$；这是将（\hyperlink{para:1:3.2.6.1}{3.2.6.1}）用于 $U_\beta$、$V_\mu$ 中的开集
$U_\alpha\cap U_\beta\cap U_\gamma$、$V_\lambda\cap V_\mu\cap V_\nu$ 得到的。

因此，存在预概形 $Z$、其底空间的开覆盖 $(Z_i)$，及对每个 $i$ 的一个同构
$g_i\colon Z_i\to U_\alpha\times_S V_\lambda$，
使得对每对 $(i,j)$，在相应的交集上有 $f_{ij}=g_i\circ g_j^{-1}$（\hyperlink{para:1:2.3.1}{2.3.1}）；
并且 $g_i(Z_i\cap Z_j)=W_{ij}$。
若 $p_i,q_i,\theta_i$ 分别为 $S$-预概形 $U_\alpha\times_S V_\lambda$ 的
两个投影及结构态射，则立即可见在 $Z_i\cap Z_j$ 上有
$p_i\circ g_i=p_j\circ g_j$，另外两个态射也满足同样的关系。
故可以由在每个 $Z_i$ 上分别等于 $p_i\circ g_i$、$q_i\circ g_i$、$\theta_i\circ g_i$
来定义预概形态射 $p\colon Z\to X$、$q\colon Z\to Y$、$\theta\colon Z\to S$。
赋予 $\theta$ 的 $Z$ 于是成为 $S$-预概形。

现证明 $Z'_i=p^{-1}(U_\alpha)\cap q^{-1}(V_\lambda)$ 等于 $Z_i$。
对于每个 $j=(\beta,\mu)$，有
\[
Z_j\cap Z'_i=g_j^{-1}\bigl(p_j^{-1}(U_\alpha)\cap q_j^{-1}(V_\lambda)\bigr).
\]
又
\[
p_j^{-1}(U_\alpha)\cap q_j^{-1}(V_\lambda)
=p_j^{-1}(U_\alpha\cap U_\beta)\cap q_j^{-1}(V_\lambda\cap V_\mu).
\]
由（\hyperlink{para:1:3.2.6.1}{3.2.6.1}），$p_j,q_j$ 在该预概形上的限制使其成为
$U_\alpha\cap U_\beta$、$V_\lambda\cap V_\mu$ 的积。
积的唯一性于是给出
$p_j^{-1}(U_\alpha)\cap q_j^{-1}(V_\lambda)=W_{ji}$。
故对于每个 $j$，有 $Z_j\cap Z'_i=Z_j\cap Z_i$，从而 $Z'_i=Z_i$。
由（\hyperlink{para:1:3.2.6.2}{3.2.6.2}），$(Z,p,q)$ 是 $X,Y$ 的积。

\egapar{3.2.6.4}
\emph{设 $\varphi\colon X\to S$、$\psi\colon Y\to S$ 为 $X,Y$ 的结构态射，
$(S_i)$ 为 $S$ 的开覆盖，并令 $X_i=\varphi^{-1}(S_i)$、$Y_i=\psi^{-1}(S_i)$。
若各积 $X_i\times_S Y_i$ 都存在，则 $X\times_S Y$ 存在。}

由（\hyperlink{para:1:3.2.6.3}{3.2.6.3}），只须证明对任意 $i,j$，积 $X_i\times_S Y_j$ 都存在。
令
\[
X_{ij}=X_i\cap X_j=\varphi^{-1}(S_i\cap S_j),\qquad
Y_{ij}=Y_i\cap Y_j=\psi^{-1}(S_i\cap S_j).
\]
由（\hyperlink{para:1:3.2.6.1}{3.2.6.1}），积 $Z_{ij}=X_{ij}\times_S Y_{ij}$ 存在。
若 $T$ 是 $S$-预概形，$g\colon T\to X_i$、$h\colon T\to Y_j$ 为 $S$-态射，
则依 $S$-态射的定义，必有
$\varphi(g(T))=\psi(h(T))\subset S_i\cap S_j$，
故 $g(T)\subset X_{ij}$、$h(T)\subset Y_{ij}$。
由此立即可见，$Z_{ij}$ 是 $X_i,Y_j$ 的积。

\egapar{3.2.6.5}
现在可以完成定理（\hyperlink{para:1:3.2.6}{3.2.6}）的证明。
若 $S$ 为仿射概形，则 $X,Y$ 分别有由仿射开集构成的覆盖 $(U_\alpha)$、$(V_\lambda)$。
由（\hyperlink{para:1:3.2.2}{3.2.2}），$U_\alpha\times_S V_\lambda$ 存在，
故由（\hyperlink{para:1:3.2.6.3}{3.2.6.3}），$X\times_S Y$ 也存在。
若 $S$ 为任意预概形，则 $S$ 有由仿射开集构成的覆盖 $(S_i)$。
设 $\varphi\colon X\to S$、$\psi\colon Y\to S$ 为结构态射，并令
$X_i=\varphi^{-1}(S_i)$、$Y_i=\psi^{-1}(S_i)$。
依上述结果，各积 $X_i\times_{S_i}Y_i$ 存在，
故各积 $X_i\times_S Y_i$ 也存在（\hyperlink{para:1:3.2.5}{3.2.5}）。
由（\hyperlink{para:1:3.2.6.4}{3.2.6.4}），$X\times_S Y$ 存在。

\egastatement{3.2.7}{推论}
设 $Z=X\times_S Y$ 是两个 $S$-预概形的积，$p,q$ 分别为到 $X,Y$ 的投影，
$\varphi,\psi$ 分别为 $X,Y$ 的结构态射。
设 $S'$ 为 $S$ 的开子集，$U,V$ 分别为 $X,Y$ 的开子集，满足
$U\subset\varphi^{-1}(S')$、$V\subset\psi^{-1}(S')$。
则积 $U\times_{S'}V$ 与 $Z$ 在 $p^{-1}(U)\cap q^{-1}(V)$ 上诱导的预概形
（视为 $S'$-预概形）典范地等同。
此外，若 $f\colon T\to X$、$g\colon T\to Y$ 是满足 $f(T)\subset U$、$g(T)\subset V$
的 $S$-态射，则 $S'$-态射 $(f,g)_{S'}$ 与将 $(f,g)_S$ 的目标限制为
$p^{-1}(U)\cap q^{-1}(V)$ 所得的态射等同。

这由（\hyperlink{para:1:3.2.5}{3.2.5}）及（\hyperlink{para:1:3.2.6.1}{3.2.6.1}）得到。

\egapar{3.2.8}
设 $(X_\alpha)$、$(Y_\lambda)$ 为两族 $S$-预概形，
$X,Y$ 分别为这两族的和（\hyperlink{heading:1:3.1}{3.1}）。
则 $X\times_S Y$ 与预概形族 $(X_\alpha\times_S Y_\lambda)_{\alpha,\lambda}$ 的和等同；
这直接由（\hyperlink{para:1:3.2.6.3}{3.2.6.3}）得到。


% END chapters/ch1_s3_2.tex

\par

% BEGIN chapters/ch1_s3_3.tex
\subsection{积的形式性质；基预概形的变换}\hypertarget{heading:1:3.3}{}
\src{108—111}

\egapar{3.3.1}
读者可以注意到，除了（\hyperlink{para:1:3.3.13}{3.3.13}）及（\hyperlink{para:1:3.3.15}{3.3.15}），本小节中所述的一切性质，在任意范畴中，
只要陈述中涉及的积存在，都不加改变地成立。
事实上，显然对于范畴中的任意对象 $S$，都可以完全仿照（\hyperlink{heading:1:2.5}{2.5}）定义
$S$-对象与 $S$-态射的概念。

\egapar{3.3.2}
首先，$X\times_S Y$ 是 $S$-预概形范畴中关于 $X,Y$ 的协变双函子。
只须注意下图交换：
\[
\begin{tikzcd}[column sep=large,row sep=large]
X\times Y\arrow[r,"f\times g"]\arrow[d]
 &X'\times Y'\arrow[r,"f'\times g'"]\arrow[d]
 &X''\times Y''\arrow[d]\\
X\arrow[r,"f"']&X'\arrow[r,"f'"']&X''.
\end{tikzcd}
\]

\egastatement{3.3.3}{命题}
对于任意 $S$-预概形 $X$，$X\times_S S$ 的第一投影
（分别为 $S\times_S X$ 的第二投影）是从该积到 $X$ 的函子性同构，
其逆同构是 $(1_X,\varphi)_S$（分别为 $(\varphi,1_X)_S$），
其中 $\varphi\colon X\to S$ 为结构态射。
因此，在典范同构意义下，可以写成
\[
 X\times_S S=S\times_S X=X.
\]

只须证明三元组 $(X,1_X,\varphi)$ 是 $X,S$ 的积。
若 $T$ 为 $S$-预概形，则从 $T$ 到 $S$ 的唯一 $S$-态射必为结构态射 $\psi\colon T\to S$。
若 $f$ 为从 $T$ 到 $X$ 的 $S$-态射，则必有 $\psi=\varphi\circ f$，断言得证。

\egastatement{3.3.4}{推论}
设 $X,Y$ 为两个 $S$-预概形，$\varphi\colon X\to S$、$\psi\colon Y\to S$ 是其结构态射。
若将 $X$ 与 $X\times_S S$、$Y$ 与 $S\times_S Y$ 典范地等同，
则投影 $X\times_S Y\to X$、$X\times_S Y\to Y$ 分别等同于
$1_X\times\psi$、$\varphi\times1_Y$。

验证是直接的，留给读者。

\egapar{3.3.5}
可以如（\hyperlink{heading:1:3.2}{3.2}）那样定义任意有限个 $n$ 个 $S$-预概形的积。
注意到 $(X_1\times_S X_2\times\cdots\times_S X_{n-1})\times_S X_n$
满足积的定义，对 $n$ 归纳，即可由（\hyperlink{para:1:3.2.6}{3.2.6}）得到这些积的存在性。
与任意范畴中的情形一样，积的唯一性蕴含其交换性与结合性。
例如，若 $p_1,p_2,p_3$ 为 $X_1\times X_2\times X_3$ 的投影，
并将此预概形与 $(X_1\times X_2)\times X_3$ 等同，
则到 $X_1\times X_2$ 的投影与 $(p_1,p_2)_S$ 等同。

\egapar{3.3.6}
设 $S,S'$ 为两个预概形，$\varphi\colon S'\to S$ 是使 $S'$ 成为 $S$-预概形的态射。
对于每个 $S$-预概形 $X$，考虑积 $X\times_S S'$，
并以 $p,\pi'$ 分别表示它到 $X,S'$ 的投影。
赋予 $\pi'$ 后，此积是一个 $S'$-预概形；当将其如此看待时，
记为 $X_{(S')}$ 或 $X_{(\varphi)}$，称其为借助态射 $\varphi$
\emph{将基预概形从 $S$ 扩张到 $S'$ 所得的预概形}，
或 $X$ 在 $\varphi$ 下的\emph{逆像}。
注意，若 $\pi$ 为 $X$ 的结构态射，$\theta$ 为将 $X\times_S S'$ 看成 $S$-预概形时的结构态射，
则下图交换：
\[
\begin{tikzcd}[column sep=large,row sep=large]
X\arrow[d,"\pi"'] & X_{(S')}\arrow[l,"p"']\arrow[d,"\pi'"]\arrow[dl,"\theta" description]\\
S & S'\arrow[l,"\varphi"].
\end{tikzcd}
\]

\egapar{3.3.7}
沿用（\hyperlink{para:1:3.3.6}{3.3.6}）的记号。对于任意 $S$-态射 $f\colon X\to Y$，
以 $f_{(S')}$ 表示 $S'$-态射
$f\times_S1\colon X_{(S')}\to Y_{(S')}$，
称 $f_{(S')}$ 为态射 $f$ 在 $\varphi$ 下的\emph{逆像}。
因此，$X_{(S')}$ 是关于 $X$ 的协变函子，从 $S$-预概形范畴取值于 $S'$-预概形范畴。

\egapar{3.3.8}
预概形 $X_{(S')}$ 也可看成一个泛映射问题的解：
每个 $S'$-预概形 $T$ 都可借助 $\varphi$ 看成 $S$-预概形；
每个 $S$-态射 $g\colon T\to X$ 都唯一地写成 $g=p\circ f$，
其中 $f$ 为一个 $S'$-态射 $T\to X_{(S')}$。
这是将积的定义用于 $S$-态射 $g$ 及 $\psi\colon T\to S'$（$T$ 的结构态射）得到的。
\transnote{本条最后一句原文把与 $\psi$ 配对的态射写为 $f$；
按该句所应用的积的定义及前文的命名，应为 $g\colon T\to X$，此处据此译出。}

\egastatement{3.3.9}{命题（基预概形扩张的传递性）}
设 $S''$ 为预概形，$\varphi'\colon S''\to S'$ 为态射。
则对于每个 $S$-预概形 $X$，存在从 $S''$-预概形
$(X_{(\varphi)})_{(\varphi')}$ 到 $S''$-预概形 $X_{(\varphi\circ\varphi')}$ 的典范函子性同构。

事实上，设 $T$ 为 $S''$-预概形，$\psi$ 为其结构态射，$g$ 为从 $T$ 到 $X$ 的 $S$-态射；
这里将 $T$ 看成结构态射为 $\varphi\circ\varphi'\circ\psi$ 的 $S$-预概形。
由于 $T$ 也可看成结构态射为 $\varphi'\circ\psi$ 的 $S'$-预概形，
可以写 $g=p\circ g'$，其中 $g'$ 为 $S'$-态射 $T\to X_{(\varphi)}$；
再写 $g'=p'\circ g''$，其中 $g''$ 为 $S''$-态射
$T\to(X_{(\varphi)})_{(\varphi')}$：
\[
\begin{tikzcd}[column sep=large,row sep=large]
X\arrow[d,"\pi"'] & X_{(\varphi)}\arrow[l,"p"']\arrow[d,"\pi'"]
 &(X_{(\varphi)})_{(\varphi')}\arrow[l,"p'"']\arrow[d,"\pi''"]\\
S&S'\arrow[l,"\varphi"]&S''\arrow[l,"\varphi'"].
\end{tikzcd}
\]
由泛映射问题的解的唯一性，即得命题。

\begin{sloppypar}
在不致混淆时，此结果也可表述为在典范同构意义下的等式
$(X_{(S')})_{(S'')}=X_{(S'')}$，或
\end{sloppypar}
\begin{equation}
(X\times_S S')\times_{S'}S''=X\times_S S''.\tag{3.3.9.1}\label{eq:1:3.3.9.1}
\end{equation}
（\hyperlink{para:1:3.3.9}{3.3.9}）中所定义的同构的函子性，同样可用态射逆像的传递公式表述为
\begin{equation}
(f_{(S')})_{(S'')}=f_{(S'')}\tag{3.3.9.2}\label{eq:1:3.3.9.2}
\end{equation}
对每个 $S$-态射 $f\colon X\to Y$ 成立。

\egastatement{3.3.10}{推论}
若 $X,Y$ 是两个 $S$-预概形，则从 $S'$-预概形
$X_{(S')}\times_{S'}Y_{(S')}$ 到 $S'$-预概形 $(X\times_S Y)_{(S')}$
存在典范函子性同构。

事实上，在典范同构意义下，由（\hyperref[eq:1:3.3.9.1]{3.3.9.1}）及 $S$-预概形的积的结合性，有
\[
(X\times_S S')\times_{S'}(Y\times_S S')
 =X\times_S(Y\times_S S')=(X\times_S Y)\times_S S'.
\]
（\hyperlink{para:1:3.3.10}{3.3.10}）中同构的函子性表述为公式
\begin{equation}
(u_{(S')},v_{(S')})_{S'}=((u,v)_S)_{(S')}\tag{3.3.10.1}\label{eq:1:3.3.10.1}
\end{equation}
对任意一对 $S$-态射 $u\colon T\to X$、$v\colon T\to Y$ 成立。

换言之，逆像函子 $X_{(S')}$ 与取积运算交换；它也与取和运算交换（\hyperlink{para:1:3.2.8}{3.2.8}）。

\egastatement{3.3.11}{推论}
设 $Y$ 为 $S$-预概形，$f\colon X\to Y$ 为使 $X$ 成为 $Y$-预概形
（从而也成为 $S$-预概形）的态射。
则预概形 $X_{(S')}$ 与积 $X\times_Y Y_{(S')}$ 等同，
且投影 $X\times_Y Y_{(S')}\to Y_{(S')}$ 与 $f_{(S')}$ 等同。

设 $\psi\colon Y\to S$ 为 $Y$ 的结构态射。有交换图
\[
\begin{tikzcd}[column sep=large,row sep=large]
S'\arrow[d] &Y_{(S')}\arrow[l]\arrow[d]&X_{(S')}\arrow[l,"f_{(S')}"']\arrow[d]\\
S&Y\arrow[l,"\psi"]&X\arrow[l,"f"].
\end{tikzcd}
\]
而 $Y_{(S')}$ 与 $S'_{(\psi)}$ 等同，$X_{(S')}$ 与 $S'_{(\psi\circ f)}$ 等同；
结合（\hyperlink{para:1:3.3.9}{3.3.9}）及（\hyperlink{para:1:3.3.4}{3.3.4}），即得推论。

\egapar{3.3.12}
设 $f\colon X\to X'$、$g\colon Y\to Y'$ 是两个 $S$-态射，且都是预概形的单态射
（T，I，1.1）。则 $f\times_S g$ 是单态射。
事实上，设 $p,q$ 为 $X\times_S Y$ 的投影，$p',q'$ 为 $X'\times_S Y'$ 的投影，
$u,v$ 为两个 $S$-态射 $T\to X\times_S Y$。
关系 $(f\times_S g)\circ u=(f\times_S g)\circ v$ 蕴含
\[
p'\circ(f\times_S g)\circ u=p'\circ(f\times_S g)\circ v,
\]
即 $f\circ p\circ u=f\circ p\circ v$。
由于 $f$ 是单态射，有 $p\circ u=p\circ v$；
同样利用 $g$ 是单态射，得 $q\circ u=q\circ v$，故 $u=v$。

因此，对于基预概形的每个扩张 $S'\to S$，
$f_{(S')}\colon X_{(S')}\to Y_{(S')}$ 是单态射。

\egapar{3.3.13}
设 $S,S'$ 是环分别为 $A,A'$ 的两个仿射概形；
态射 $S'\to S$ 对应于环同态 $A\to A'$。
若 $X$ 为 $S$-预概形，则也以 $X_{(A')}$ 或 $X\otimes_A A'$ 表示 $S'$-预概形 $X_{(S')}$。
当 $X$ 本身是环为 $B$ 的仿射概形时，$X_{(A')}$ 为仿射概形，其环为
$B_{(A')}=B\otimes_A A'$，即将 $A$-代数 $B$ 的标量环扩张到 $A'$ 所得的环。

\egapar{3.3.14}
沿用（\hyperlink{para:1:3.3.6}{3.3.6}）的记号。对于每个 $S$-态射 $f\colon S'\to X$，
$f'=(f,1_{S'})_S$ 是一个 $S'$-态射 $S'\to X'=X_{(S')}$，
满足 $p\circ f'=f$、$\pi'\circ f'=1_{S'}$，即 $X'$ 的一个 $S'$-截面。
反过来，若 $f'$ 是这样的 $S'$-截面，则 $f=p\circ f'$ 是 $S$-态射 $S'\to X$。
由此定义了典范的一一对应
\[
\Hom_S(S',X)\xrightarrow{\sim}\Hom_{S'}(S',X').
\]
称 $f'$ 为 $f$ 的\emph{图态射}，记为 $\Gamma_f$。

\egapar{3.3.15}
给定预概形 $X$，它总可看成 $\mathbf Z$-预概形。
由（\hyperlink{para:1:3.3.14}{3.3.14}）特别可知，$X\otimes_{\mathbf Z}\mathbf Z[T]$
（$T$ 为不定元）的 $X$-截面与态射 $X\to\Spec(\mathbf Z[T])$ 一一对应。
\transnote{原文此处写作“morphismes $\mathbf Z[T]\to X$”。
按照所引用的（\hyperlink{para:1:3.3.14}{3.3.14}），应为 $X\to\Spec(\mathbf Z[T])$；
正文明确采用这一方向，而不是把环同态与预概形态射的方向混同。}
我们证明：这些 $X$-截面也与结构层 $\OO_X$ 在 $X$ 上的截面一一对应。
事实上，设 $(U_\alpha)$ 为 $X$ 的仿射开覆盖，
$u\colon X\to X\otimes_{\mathbf Z}\mathbf Z[T]$ 为 $X$-态射，
$u_\alpha$ 为其在 $U_\alpha$ 上的限制。
若 $A_\alpha$ 为仿射概形 $U_\alpha$ 的环，则
$U_\alpha\otimes_{\mathbf Z}\mathbf Z[T]$ 是环为 $A_\alpha[T]$ 的仿射概形（\hyperlink{para:1:3.2.2}{3.2.2}），
而 $u_\alpha$ 典范地对应于一个 $A_\alpha$-同态
$A_\alpha[T]\to A_\alpha$（\hyperlink{para:1:1.7.3}{1.7.3}）。
这样的同态完全由 $T$ 在 $A_\alpha$ 中的像确定，记此像为
$s_\alpha\in A_\alpha=\Gamma(U_\alpha,\OO_X)$。
写出 $u_\alpha,u_\beta$ 在仿射开集 $V\subset U_\alpha\cap U_\beta$ 上的限制相同这一条件，
立即得到 $s_\alpha,s_\beta$ 在 $V$ 上相同。
因此，族 $(s_\alpha)$ 是 $\OO_X$ 的某个全局截面 $s$ 在各 $U_\alpha$ 上的限制。
反过来，显然这样的截面定义了一族态射 $(u_\alpha)$，
它们是某个 $X$-态射 $X\to X\otimes_{\mathbf Z}\mathbf Z[T]$ 在各 $U_\alpha$ 上的限制。
此结果将在 II，1.7.12 中推广。


% END chapters/ch1_s3_3.tex

\par

% BEGIN chapters/ch1_s3_4.tex
\subsection{在预概形中取值的点；几何点}\hypertarget{heading:1:3.4}{}
\src{111—114}

\egapar{3.4.1}
设 $X$ 为预概形。对于任意预概形 $T$，也以 $X(T)$ 表示态射 $T\to X$ 的集合
$\Hom(T,X)$；此集合的元素又称为 $X$ 的\emph{在 $T$ 中取值的点}。
将每个态射 $f\colon T\to T'$ 对应到映射
$X(T')\to X(T)$、$u'\mapsto u'\circ f$，
便可见，对于固定的 $X$，$X(T)$ 是关于 $T$ 的反变函子，
从预概形范畴取值于集合范畴。
此外，每个预概形态射 $g\colon X\to Y$ 都定义一个函子性同态
$X(T)\to Y(T)$，将 $v\in X(T)$ 映到 $g\circ v$。

\egapar{3.4.2}
给定三个集合 $P,Q,R$ 及两个映射 $\varphi\colon P\to R$、$\psi\colon Q\to R$。
积集合 $P\times Q$ 中由满足 $\varphi(p)=\psi(q)$ 的有序对 $(p,q)$ 组成的子集，
称为 $P,Q$ \emph{在 $R$ 上的纤维积}（相对于 $\varphi,\psi$），记为 $P\times_R Q$。
使用（\hyperlink{para:1:3.4.1}{3.4.1}）的记号，$S$-预概形的积的定义（\hyperlink{para:1:3.2.1}{3.2.1}）也可解释为公式
\begin{equation}
(X\times_S Y)(T)=X(T)\times_{S(T)}Y(T),\tag{3.4.2.1}\label{eq:1:3.4.2.1}
\end{equation}
其中映射 $X(T)\to S(T)$、$Y(T)\to S(T)$ 分别对应于结构态射 $X\to S$、$Y\to S$。

\egapar{3.4.3}
若给定预概形 $S$，并只考虑 $S$-预概形与 $S$-态射，则也以 $X(T)_S$ 表示
$S$-态射 $T\to X$ 的集合 $\Hom_S(T,X)$；在不致混淆时，允许省略下标 $S$。
$X(T)_S$ 的元素又称为 $S$-预概形 $X$ 的\emph{在 $S$-预概形 $T$ 中取值的点}
（可能混淆时称为\emph{$S$-点}）。
特别地，$X$ 的一个 $S$-截面就是 $X$ 的一个在 $S$ 中取值的点。
这时（\hyperref[eq:1:3.4.2.1]{3.4.2.1}）也写成
\begin{equation}
(X\times_S Y)(T)_S=X(T)_S\times Y(T)_S.\tag{3.4.3.1}\label{eq:1:3.4.3.1}
\end{equation}
更一般地，若 $Z$ 是 $S$-预概形，$X,Y,T$ 是 $Z$-预概形
（从而也自然是 $S$-预概形），则
\begin{equation}
(X\times_Z Y)(T)_S=X(T)_S\times_{Z(T)}Y(T)_S.\tag{3.4.3.2}\label{eq:1:3.4.3.2}
\end{equation}
注意，若要证明由 $S$-预概形 $W$ 及两个 $S$-态射
$r\colon W\to X$、$s\colon W\to Y$ 组成的三元组 $(W,r,s)$
是 $X,Y$ 在 $Z$ 上的积，则依定义只须验证：对于每个 $S$-预概形 $T$，下图
\[
\begin{tikzcd}[column sep=large,row sep=large]
W(T)_S\arrow[r,"r'"]\arrow[d,"s'"'] &X(T)_S\arrow[d,"\varphi'"]\\
Y(T)_S\arrow[r,"\psi'"']&Z(T)_S
\end{tikzcd}
\]
使 $W(T)_S$ 成为 $X(T)_S,Y(T)_S$ 在 $Z(T)_S$ 上的纤维积。
其中 $r',s'$ 分别对应于 $r,s$，$\varphi',\psi'$ 分别对应于结构态射
$\varphi\colon X\to Z$、$\psi\colon Y\to Z$。

\egapar{3.4.4}
当上述 $T$（分别为 $S$）是环为 $B$（分别为 $A$）的仿射概形时，
在上述记号中将 $T$（分别将 $S$）换成 $B$（分别换成 $A$）。
于是，对于 $X(B)$（分别为 $X(B)_A$）的元素，称之为
$X$ 的\emph{在环 $B$ 中取值的点}（分别为 $A$-预概形 $X$ 的
\emph{在 $A$-代数 $B$ 中取值的点}）。
注意，现在 $X(B)$、$X(B)_A$ 都是关于 $B$ 的协变函子。
同样，以 $X(T)_A$ 表示 $A$-预概形 $X$ 的在 $A$-预概形 $T$ 中取值的点的集合。

\egapar{3.4.5}
特别考虑 $T=\Spec(A)$、$A$ 为局部环的情形。
这时，$X(A)$ 的元素与局部同态 $\OO_x\to A$（$x\in X$）一一对应（\hyperlink{para:1:2.4.4}{2.4.4}）；
称底空间 $X$ 的点 $x$ 为它所对应的 $A$-值点的\emph{落点}。

更特别地，称预概形 $X$ 的在一个域 $K$ 中取值的点为 $X$ 的\emph{几何点}。
因此，给出这样的点，就是给出它在底空间 $X$ 中的落点 $x$，以及 $k(x)$ 的一个扩张域 $K$。
称 $K$ 为相应几何点的\emph{取值域}，也说此几何点\emph{位于} $x$。
由此定义了映射 $X(K)\to X$，将一个 $K$-值几何点映到其落点。
\transnote{这里忠实保留本版的定义：取值域 $K$ 可以是任意域，未附加代数闭等条件。
“落点”对应原文 \textit{localité}，区别于一个给定取值域中的点本身。}

若 $S'=\Spec(K)$ 是 $S$-预概形（即把 $K$ 看成 $S$ 的某个剩余域 $k(s)$ 的扩张，$s\in S$），
且 $X$ 为 $S$-预概形，则 $X(K)_S$ 的一个元素，也称为
\emph{$X$ 的位于 $s$ 之上、在 $K$ 中取值的几何点}，
就是给出一个 $k(s)$-单同态 $k(x)\to K$，
其中 $x$ 是 $X$ 中位于 $s$ 之上的一点，故 $k(x)$ 是 $k(s)$ 的扩张。

特别地，若 $S=\Spec(K)=\{\xi\}$，则 $X$ 的 $K$-值几何点与满足 $k(x)=K$ 的
点 $x\in X$ 等同；这样的点又称为 $K$-预概形 $X$ 的\emph{在 $K$ 上有理的点}。
若 $K'$ 为 $K$ 的扩张，则 $X$ 的 $K'$-值几何点与 $X_{(K')}$ 的
在 $K'$ 上有理的点一一对应（\hyperlink{para:1:3.3.14}{3.3.14}）。

\egastatement{3.4.6}{引理}
设 $X_i$（$1\leq i\leq n$）为 $S$-预概形，$s\in S$，
且对每个 $i$，$x_i$ 是 $X_i$ 中位于 $s$ 之上的一点。
则存在 $k(s)$ 的一个扩张域 $K$，以及积
$Y=X_1\times_S X_2\times_S\cdots\times_S X_n$ 的一个 $K$-值几何点，
使它的各个投影分别位于 $x_i$。

事实上，各域 $k(x_i)$ 存在到同一个 $k(s)$ 的扩张域 $K$ 的
$k(s)$-单同态（Bourbaki，\textit{Alg.}，第 V 章，§4，命题 2）。
各复合 $k(s)\to k(x_i)\to K$ 都相同，故与 $k(x_i)\to K$ 对应的态射
$\Spec(K)\to X_i$ 都是 $S$-态射，从而确定唯一的态射 $\Spec(K)\to Y$。
若 $y$ 为 $Y$ 中相应的点，则它在每个 $X_i$ 中的投影显然为 $x_i$。

\egastatement{3.4.7}{命题}
设 $X_i$（$1\leq i\leq n$）为 $S$-预概形，对每个 $i$ 给定点 $x_i\in X_i$。
则存在 $Y=X_1\times_S X_2\times_S\cdots\times_S X_n$ 的一点 $y$，
使其第 $i$ 个投影为 $x_i$（$1\leq i\leq n$），
当且仅当各 $x_i$ 都位于 $S$ 的同一点 $s$ 之上。

此条件显然必要；引理（\hyperlink{para:1:3.4.6}{3.4.6}）说明它也充分。

换言之，以 $(X)$ 表示 $X$ 的底集合，便有典范满射
$(X\times_S Y)\to(X)\times_{(S)}(Y)$。
必须注意，此映射一般不是单射；也就是说，$X\times_S Y$ 中可能存在多个不同的点，
具有同样的投影 $x\in X$、$y\in Y$。
即使 $S,X,Y$ 是域 $k,K,K'$ 的素谱，也已如此，
因为张量积 $K\otimes_k K'$ 一般具有多个不同的素理想（参见（\hyperlink{para:1:3.4.9}{3.4.9}））。

\egastatement{3.4.8}{推论}
设 $f\colon X\to Y$ 为 $S$-态射，
$f_{(S')}\colon X_{(S')}\to Y_{(S')}$ 为由基预概形扩张 $S'\to S$ 得到的 $S'$-态射。
令 $p\colon X_{(S')}\to X$、$q\colon Y_{(S')}\to Y$ 为投影。
则对于 $X$ 的任意子集 $M$，有
\[
 q^{-1}(f(M))=f_{(S')}(p^{-1}(M)).
\]

事实上，由（\hyperlink{para:1:3.3.11}{3.3.11}），借助交换图
\[
\begin{tikzcd}[column sep=large,row sep=large]
X\arrow[d,"f"']&X_{(S')}\arrow[l,"p"']\arrow[d,"f_{(S')}"]\\
Y&Y_{(S')}\arrow[l,"q"]
\end{tikzcd}
\]
可将 $X_{(S')}$ 与积 $X\times_Y Y_{(S')}$ 等同。
由（\hyperlink{para:1:3.4.7}{3.4.7}），对于 $x\in M$、$y'\in Y_{(S')}$，关系 $q(y')=f(x)$
等价于存在 $x'\in X_{(S')}$，使 $p(x')=x$、$f_{(S')}(x')=y'$。
推论由此得到。

引理（\hyperlink{para:1:3.4.6}{3.4.6}）可进一步精确化如下。

\egastatement{3.4.9}{命题}
设 $X,Y$ 为两个 $S$-预概形，$x\in X$、$y\in Y$ 位于同一点 $s\in S$ 之上。
则 $X\times_S Y$ 中投影分别为 $x,y$ 的点的集合，与把 $k(x),k(y)$
看成 $k(s)$ 的扩张时，二者的复合扩张的类型的集合之间，存在典范的一一对应
（Bourbaki，\textit{Alg.}，第 VIII 章，§8，命题 2）。

令 $p,q$ 分别为 $X\times_S Y$ 到 $X,Y$ 的投影，
$E$ 为底空间中的子空间 $p^{-1}(x)\cap q^{-1}(y)$。
首先注意，态射 $\Spec(k(x))\to S$、$\Spec(k(y))\to S$ 分别分解为
\[
\Spec(k(x))\longrightarrow\Spec(k(s))\longrightarrow S,\qquad
\Spec(k(y))\longrightarrow\Spec(k(s))\longrightarrow S.
\]
由于 $\Spec(k(s))\to S$ 是单态射（\hyperlink{para:1:2.4.7}{2.4.7}），由（\hyperlink{para:1:3.2.4}{3.2.4}）有
\[
\begin{aligned}
P&=\Spec(k(x))\times_S\Spec(k(y))\\
 &=\Spec(k(x))\times_{\Spec(k(s))}\Spec(k(y))
 =\Spec\bigl(k(x)\otimes_{k(s)}k(y)\bigr).
\end{aligned}
\]
以 $P_0$ 表示预概形 $P$ 的底集合。
我们定义两个互逆的映射 $\alpha\colon P_0\to E$、$\beta\colon E\to P_0$。
若 $i\colon\Spec(k(x))\to X$、$j\colon\Spec(k(y))\to Y$ 是典范态射（\hyperlink{para:1:2.4.5}{2.4.5}），
则取 $\alpha$ 为与态射 $i\times_S j$ 对应的底空间映射。
另一方面，依假设，每个 $z\in E$ 都确定两个 $k(s)$-单同态
$k(x)\to k(z)$、$k(y)\to k(z)$，从而确定张量积上的 $k(s)$-代数同态
\[
 k(x)\otimes_{k(s)}k(y)\longrightarrow k(z),
\]
\transnote{此处原文称所得映射为“$k(s)$-monomorphisme produit tensoriel”。
由两个域嵌入所确定的张量积到共同域的映射，一般只是代数同态，并不必单射；
本段构造 $\Spec(k(z))\to P$ 只需要它是同态。正文在此明确译作“代数同态”。}
因而确定态射 $\Spec(k(z))\to P$。
令 $\beta(z)$ 为该态射将 $\Spec(k(z))$ 的唯一点映到 $P_0$ 所得的点。
验证 $\alpha\circ\beta$ 与 $\beta\circ\alpha$ 为恒等映射，
只须使用（\hyperlink{para:1:2.4.5}{2.4.5}）及积的定义（\hyperlink{para:1:3.2.1}{3.2.1}）。
最后，已知 $P_0$ 与 $k(x),k(y)$ 的复合扩张的类型的集合一一对应
（Bourbaki，\textit{Alg.}，第 VIII 章，§8，命题 1）。


% END chapters/ch1_s3_4.tex

\par

% BEGIN chapters/ch1_s3_5.tex
\subsection{满射与单射}\hypertarget{heading:1:3.5}{}
\src{114—117}

\egapar{3.5.1}
一般地，考虑预概形态射的一个性质 $P$，以及以下两个命题：
\begin{enumerate}[label=(\roman*)]
\item 若 $f\colon X\to X'$、$g\colon Y\to Y'$ 是具有性质 $P$ 的两个 $S$-态射，
则 $f\times_S g$ 具有性质 $P$。
\item 若 $f\colon X\to Y$ 是具有性质 $P$ 的 $S$-态射，
则由基预概形扩张得到的每个 $S'$-态射
$f_{(S')}\colon X_{(S')}\to Y_{(S')}$ 都具有性质 $P$。
\end{enumerate}
由于 $f_{(S')}=f\times_S1_{S'}$，可见：若对于每个预概形 $X$，恒等态射 $1_X$
都具有性质 $P$，则 (i) 蕴含 (ii)。
另一方面，$f\times_S g$ 是复合态射
\[
X\times_S Y\xrightarrow{f\times1_Y}X'\times_S Y
\xrightarrow{1_{X'}\times g}X'\times_S Y',
\]
故可见：若两个具有性质 $P$ 的态射的复合仍具有此性质，则 (ii) 蕴含 (i)。
此观察的第一个应用如下。

\egastatement{3.5.2}{命题}
\begin{enumerate}[label=(\roman*)]
\item 若 $f\colon X\to X'$、$g\colon Y\to Y'$ 是满射的 $S$-态射，
则 $f\times_S g$ 满射。
\item 若 $f\colon X\to Y$ 是满射的 $S$-态射，
则对于基预概形的任意扩张 $S'$，$f_{(S')}$ 满射。
\end{enumerate}
两个满射的复合仍为满射，故只须证明 (ii)；
而将（\hyperlink{para:1:3.4.8}{3.4.8}）用于 $M=X$，立即得到此结论。

\egastatement{3.5.3}{命题}
态射 $f\colon X\to Y$ 满射，当且仅当对于每个域 $K$ 及每个态射 $\Spec(K)\to Y$，
存在 $K$ 的一个扩张域 $K'$ 及一个态射 $\Spec(K')\to X$，使下图交换：
\[
\begin{tikzcd}[column sep=large,row sep=large]
X\arrow[d,"f"']&\Spec(K')\arrow[l]\arrow[d]\\
Y&\Spec(K)\arrow[l].
\end{tikzcd}
\]

条件是充分的，因为对于每个 $y\in Y$，只须将它用于与单同态 $k(y)\to K$
对应的态射 $\Spec(K)\to Y$，其中 $K$ 是 $k(y)$ 的一个扩张（\hyperlink{para:1:2.4.6}{2.4.6}）。
反过来，设 $f$ 满射，并令 $y\in Y$ 为 $\Spec(K)$ 的唯一点的像。
存在 $x\in X$，使 $f(x)=y$；考虑对应的单同态 $k(y)\to k(x)$（\hyperlink{para:1:2.2.1}{2.2.1}）。
只须取 $k(y)$ 的一个扩张域 $K'$，使存在 $k(y)$-单同态 $k(x)\to K'$ 与 $K\to K'$
（Bourbaki，\textit{Alg.}，第 V 章，§4，命题 2）。
与 $k(x)\to K'$ 对应的态射 $\Spec(K')\to X$ 即符合要求。

用（\hyperlink{para:1:3.4.5}{3.4.5}）引入的语言，也可说：$Y$ 的每个 $K$-值几何点，
都来自 $X$ 的一个在 $K$ 的某个扩张中取值的几何点。

\egastatement{3.5.4}{定义}
称预概形态射 $f\colon X\to Y$ 为\emph{普遍单射}，或一个\emph{根式态射}，
若对于每个域 $K$，相应的映射 $X(K)\to Y(K)$ 都是单射。

由定义立即得到，每个预概形的单态射（T，1.1）都是根式态射。

\egapar{3.5.5}
为了使态射 $f\colon X\to Y$ 为根式态射，只须定义（\hyperlink{para:1:3.5.4}{3.5.4}）的条件
对每个代数闭域成立。事实上，若 $K$ 是任意域，$K'$ 是 $K$ 的一个代数闭扩张，
则下图交换：
\[
\begin{tikzcd}[column sep=large,row sep=large]
X(K)\arrow[r,"\alpha"]\arrow[d,"\varphi"']&Y(K)\arrow[d,"\varphi'"]\\
X(K')\arrow[r,"\alpha'"']&Y(K').
\end{tikzcd}
\]
其中 $\varphi,\varphi'$ 来自态射 $\Spec(K')\to\Spec(K)$，
$\alpha,\alpha'$ 对应于 $f$。
而 $\varphi$ 单射，$\alpha'$ 依假设也单射，故 $\alpha$ 必单射。

\egastatement{3.5.6}{命题}
设 $f\colon X\to Y$、$g\colon Y\to Z$ 是两个预概形态射。
\begin{enumerate}[label=(\roman*)]
\item 若 $f,g$ 都是根式态射，则 $g\circ f$ 也是根式态射。
\item 反过来，若 $g\circ f$ 是根式态射，则 $f$ 也是根式态射。
\end{enumerate}
结合定义（\hyperlink{para:1:3.5.4}{3.5.4}），此命题归结为对映射
$X(K)\to Y(K)\to Z(K)$ 的相应断言，后者是显然的。

\egastatement{3.5.7}{命题}
\begin{enumerate}[label=(\roman*)]
\item 若 $S$-态射 $f\colon X\to X'$、$g\colon Y\to Y'$ 是根式态射，
则 $f\times_S g$ 也是根式态射。
\item 若 $S$-态射 $f\colon X\to Y$ 是根式态射，
则对于基预概形的每个扩张 $S'\to S$，
$f_{(S')}\colon X_{(S')}\to Y_{(S')}$ 也是根式态射。
\end{enumerate}
由（\hyperlink{para:1:3.5.1}{3.5.1}），只须证明 (i)。我们已知（\hyperref[eq:1:3.4.2.1]{3.4.2.1}）
\[
\begin{gathered}
(X\times_S Y)(K)=X(K)\times_{S(K)}Y(K),\\
(X'\times_S Y')(K)=X'(K)\times_{S(K)}Y'(K).
\end{gathered}
\]
与 $f\times_S g$ 对应的映射
$(X\times_S Y)(K)\to(X'\times_S Y')(K)$
与 $(u,v)\mapsto(f\circ u,g\circ v)$ 等同，命题由此立即得到。

\egastatement{3.5.8}{命题}
态射 $f=(\psi,\theta)\colon X\to Y$ 是根式态射，当且仅当 $\psi$ 单射，
且对于每个 $x\in X$，单同态
$\theta^x\colon k(\psi(x))\to k(x)$ 使 $k(x)$ 成为 $k(\psi(x))$ 的纯不可分扩张。

设 $f$ 为根式态射，先证明 $\psi(x_1)=\psi(x_2)=y$ 必蕴含 $x_1=x_2$。
事实上，存在 $k(y)$ 的扩张域 $K$ 及 $k(y)$-单同态
$k(x_1)\to K$、$k(x_2)\to K$
（Bourbaki，\textit{Alg.}，第 V 章，§4，命题 2）。
对应的态射 $u_1,u_2\colon\Spec(K)\to X$ 满足 $f\circ u_1=f\circ u_2$，
故依假设有 $u_1=u_2$；特别地，$x_1=x_2$。
现在借助 $\theta^x$ 将 $k(x)$ 看成 $k(\psi(x))$ 的扩张。
若它不是纯不可分扩张，则存在到 $k(\psi(x))$ 的某个代数闭扩张 $K$ 的
两个不同的 $k(\psi(x))$-单同态 $k(x)\to K$，
与之对应的两个态射 $\Spec(K)\to X$ 将违背假设。
反过来，结合（\hyperlink{para:1:2.4.6}{2.4.6}），命题中的条件显然足以保证 $f$ 为根式态射。

\egastatement{3.5.9}{推论}
若 $A$ 是环，$S$ 是 $A$ 的乘法子集，则典范态射
$\Spec(S^{-1}A)\to\Spec(A)$ 是根式态射。

事实上，此态射是单态射（\hyperlink{para:1:1.6.2}{1.6.2}）。

\egastatement{3.5.10}{推论}
设 $f\colon X\to Y$ 是根式态射，$g\colon Y'\to Y$ 为态射，
并令 $X'=X_{(Y')}=X\times_Y Y'$。
则根式态射 $f_{(Y')}$（\hyperlink{para:1:3.5.7}{3.5.7}，(ii)）将底空间 $X'$ 双射到 $g^{-1}(f(X))$。
此外，对于每个域 $K$，集合 $X'(K)$ 与以下子集等同：
$Y'(K)$ 在与 $g$ 对应的映射 $Y'(K)\to Y(K)$ 下，对子集 $X(K)\subset Y(K)$ 的逆像。

第一项断言直接由（\hyperlink{para:1:3.5.8}{3.5.8}）及（\hyperlink{para:1:3.4.8}{3.4.8}）得到；第二项由交换图
\[
\begin{tikzcd}[column sep=large,row sep=large]
X'(K)\arrow[r,hook]\arrow[d]&Y'(K)\arrow[d]\\
X(K)\arrow[r,hook]&Y(K)
\end{tikzcd}
\]
及前述纤维积的描述得到。

\egastatement{3.5.11}{注}
称预概形态射 $f=(\psi,\theta)$ 为\emph{单射}，若映射 $\psi$ 是单射。
态射 $f=(\psi,\theta)\colon X\to Y$ 为根式态射，当且仅当对于每个态射 $Y'\to Y$，
$f_{(Y')}\colon X_{(Y')}\to Y'$ 是单射；这就说明了“普遍单射态射”这一术语的缘由。
事实上，必要性由（\hyperlink{para:1:3.5.7}{3.5.7}，(ii)）及（\hyperlink{para:1:3.5.8}{3.5.8}）得到。
反过来，此条件首先蕴含 $\psi$ 单射。
若对于某个 $x\in X$，单同态 $\theta^x\colon k(\psi(x))\to k(x)$ 不是纯不可分的，
则存在 $k(\psi(x))$ 的一个扩张域 $K$，以及对应于同一个态射 $\Spec(K)\to Y$ 的
两个不同的态射 $\Spec(K)\to X$（\hyperlink{para:1:3.5.8}{3.5.8}）。
令 $Y'=\Spec(K)$，便得到 $X_{(Y')}$ 的两个不同的 $Y'$-截面（\hyperlink{para:1:3.3.14}{3.3.14}），
这与 $f_{(Y')}$ 单射的假设矛盾。


% END chapters/ch1_s3_5.tex

\par

% BEGIN chapters/ch1_s3_6.tex
\subsection{纤维}\hypertarget{heading:1:3.6}{}
\src{117—118}

\egastatement{3.6.1}{命题}
设 $f\colon X\to Y$ 为态射，$y\in Y$，
$\mathfrak a_y$ 是 $\OO_y$ 关于 $\mathfrak m_y$-预进制拓扑的一个定义理想。
则投影 $p\colon X\times_Y\Spec(\OO_y/\mathfrak a_y)\to X$
是从预概形 $X\times_Y\Spec(\OO_y/\mathfrak a_y)$ 的底空间
到纤维 $f^{-1}(y)$ 的同胚，后者赋予由 $X$ 的底空间拓扑诱导的拓扑。

由于 $\Spec(\OO_y/\mathfrak a_y)\to Y$ 是根式态射（\hyperlink{para:1:3.5.4}{3.5.4} 及 \hyperlink{para:1:2.4.7}{2.4.7}），
且 $\mathfrak m_y/\mathfrak a_y$ 依假设幂零，故
$\Spec(\OO_y/\mathfrak a_y)$ 只有一点（\hyperlink{para:1:1.1.12}{1.1.12}）。
于是已知（\hyperlink{para:1:3.5.10}{3.5.10} 及 \hyperlink{para:1:3.3.4}{3.3.4}），$p$ 将底集合与 $f^{-1}(y)$ 等同；
只须再证明 $p$ 为同胚。
由（\hyperlink{para:1:3.2.7}{3.2.7}），这是 $X,Y$ 上的局部问题，因而可假设
$X=\Spec(B)$、$Y=\Spec(A)$，其中 $B$ 为 $A$-代数。
态射 $p$ 对应于同态 $1\otimes\varphi\colon B\to B\otimes_A A'$，
这里 $A'=A_y/\mathfrak a_y$，$\varphi$ 为 $A$ 到 $A'$ 的典范映射。
而 $B\otimes_A A'$ 的每个元素都写成
\[
\sum_i b_i\otimes\frac{\varphi(a_i)}{\varphi(s)}
 =\left(\sum_i b_i a_i\otimes1\right)(1\otimes\varphi(s))^{-1},
\qquad s\notin\jj_y,
\]
因此可应用命题（\hyperlink{para:1:1.2.4}{1.2.4}）。

\egapar{3.6.2}
在本专著以下各处，当把态射的一条纤维 $f^{-1}(y)$ 看成赋有 $k(y)$-预概形结构时，
始终是指借助到 $X$ 的投影，将 $X\times_Y\Spec(k(y))$ 的结构搬运到此纤维上所得的预概形。
后一个积也写作 $X\otimes_Y k(y)$ 或 $X\otimes_{\OO_y}k(y)$。
更一般地，若 $B$ 为 $\OO_y$-代数，则以 $X\otimes_Y B$ 或
$X\otimes_{\OO_y}B$ 表示积 $X\times_Y\Spec(B)$。

按上述约定，由（\hyperlink{para:1:3.5.10}{3.5.10}），$X$ 的位于 $y$ 之上、在 $k(y)$ 的一个扩张域 $K$
中取值的点，与 $f^{-1}(y)$ 的 $K$-值点等同。

\egapar{3.6.3}
设 $f\colon X\to Y$、$g\colon Y\to Z$ 是两个态射，$h=g\circ f$ 为其复合。
对于每个 $z\in Z$，纤维 $h^{-1}(z)$ 是一个同构于下式的预概形：
\[
X\times_Z\Spec(k(z))
 =(X\times_Y Y)\times_Z\Spec(k(z))=X\times_Y g^{-1}(z).
\]
特别地，若 $U$ 为 $X$ 的开子集，则由预概形 $f^{-1}(y)$ 在 $U\cap f^{-1}(y)$ 上
诱导的预概形，同构于 $f_U^{-1}(y)$，其中 $f_U$ 是 $f$ 在 $U$ 上的限制。

\egastatement{3.6.4}{命题（纤维的传递性）}
设 $f\colon X\to Y$、$g\colon Y'\to Y$ 为两个态射；
令 $X'=X\times_Y Y'=X_{(Y')}$，$f'=f_{(Y')}\colon X'\to Y'$。
对于每个 $y'\in Y'$，若令 $y=g(y')$，则预概形 $f'^{-1}(y')$
同构于 $f^{-1}(y)\otimes_{k(y)}k(y')$。

事实上，只须注意，两个预概形
\[
(X\otimes_Y k(y))\otimes_{k(y)}k(y'),\qquad
(X\times_Y Y')\otimes_{Y'}k(y')
\]
都由（\hyperref[eq:1:3.3.9.1]{3.3.9.1}）与 $X\times_Y\Spec(k(y'))$ 典范同构。

特别地，若 $V$ 为 $y$ 在 $Y$ 中的开邻域，$f_V$ 表示 $f$ 在 $f^{-1}(V)$
上诱导的预概形上的限制，则预概形 $f^{-1}(y)$ 与 $f_V^{-1}(y)$ 典范地等同。

\egastatement{3.6.5}{命题}
设 $f\colon X\to Y$ 为态射，$y\in Y$，$Z=\Spec(\OO_y)$ 为局部预概形，
$p=(\psi,\theta)\colon X\times_Y Z\to X$ 为投影。
将 $Z$ 的底空间与 $Y$ 的一个子空间等同（\hyperlink{para:1:2.4.2}{2.4.2}）以后，
$p$ 是从 $X\times_Y Z$ 的底空间到 $X$ 的子空间 $f^{-1}(Z)$ 的同胚。
此外，对于每个 $t\in X\times_Y Z$，若令 $z=\psi(t)$，
则 $\theta_t^\sharp\colon\OO_z\to\OO_t$ 是同构。

由于 $Z$（视为 $Y$ 的子空间）包含于每个含 $y$ 的仿射开集（\hyperlink{para:1:2.4.2}{2.4.2}），
可以如（\hyperlink{para:1:3.6.1}{3.6.1}）那样归结到 $X=\Spec(A)$、$Y=\Spec(B)$ 为仿射概形的情形，
其中 $A$ 为 $B$-代数。
这时 $X\times_Y Z$ 是 $A\otimes_B B_y$ 的素谱，
而此环与 $S^{-1}A$ 典范地等同，其中 $S$ 是 $B\setminus\jj_y$ 在 $A$ 中的像（0，\hyperlink{para:0:1.5.2}{1.5.2}）。
由于 $p$ 对应于典范同态 $A\to S^{-1}A$，命题由（\hyperlink{para:1:1.6.2}{1.6.2}）得到。


% END chapters/ch1_s3_6.tex

\par

% BEGIN chapters/ch1_s3_7.tex
\subsection[应用：预概形的模理想约化]{应用：预概形的模 \texorpdfstring{$\mathfrak J$}{J} 约化}
\src{118—119}
\origfootnote{本小节使用第 I 章后文及第 II 章中的概念和结果，
此后不会用到它；本小节只面向熟悉经典代数几何的读者。}

\egapar{3.7.1}
设 $A$ 为环，$X$ 为 $A$-预概形，$\mathfrak J$ 为 $A$ 的理想。
则 $X_0=X\otimes_A(A/\mathfrak J)$ 是一个 $(A/\mathfrak J)$-预概形，
有时称它为由 $X$ 作\emph{模 $\mathfrak J$ 约化}而得。

\egapar{3.7.2}
上述术语尤其用于 $A$ 为局部环、$\mathfrak J$ 为其极大理想的情形。
这时 $X_0$ 是 $A$ 的剩余域 $k=A/\mathfrak J$ 上的预概形。
若 $A$ 还是整环，分式域为 $K$，也可考虑 $K$-预概形 $X'=X\otimes_A K$。
迄今为止，人们通常不甚严格地说 $X_0$ 是由 $X'$ 作模 $\mathfrak J$ 约化而得；
我们将不采用这种说法。
使用此种说法时，$A$ 是一个一维局部环（通常为离散赋值环），
而且或多或少地暗含：给定的 $K$-预概形 $X'$ 是某个 $K$-预概形 $P'$
的闭子预概形；实际上 $P'$ 是典型的射影空间 $\mathbf P_K^r$（参见 II，4.1.1）。
此外，$P'$ 本身具有形式 $P'=P\otimes_A K$，其中 $P$ 是给定的 $A$-预概形；
实际上，使用 II，4.1.1 的记号，$P$ 是 $A$-概形 $\mathbf P_A^r$。
用我们的语言，从 $X'$ 定义 $X_0$ 的过程如下。

考虑仿射概形 $Y=\Spec(A)$。它由两个点组成：唯一闭点 $y=\mathfrak J$，以及泛点 $(0)$。
因此，由泛点组成的单点集 $U$ 是 $Y$ 中的开集 $U=\Spec(K)$。
若 $X$ 是 $A$-预概形，即 $Y$-预概形，且 $\psi\colon X\to Y$ 为结构态射，
则 $X'=X\otimes_A K$ 就是 $X$ 在 $\psi^{-1}(U)$ 上诱导的预概形。
特别地，若 $\varphi\colon P\to Y$ 是结构态射，
则 $P'=\varphi^{-1}(U)$ 的一个闭子预概形 $X'$ 是 $P$ 的一个局部闭子预概形。
若 $P$ 为诺特预概形，例如 $A$ 诺特且 $P$ 为有限型 $A$-预概形，
则存在 $P$ 的一个不小于 $X'$ 的最小闭子预概形 $X=\overline{X'}$（\hyperlink{para:1:9.5.10}{9.5.10}）。
$X'$ 是 $X$ 在开集 $\varphi^{-1}(U)\cap X$ 上诱导的预概形，
故同构于 $X\otimes_A K$（\hyperlink{para:1:9.5.10}{9.5.10}）。
所以，$X'$ 到 $P'=P\otimes_A K$ 的浸入，使我们能典范地把 $X'$ 看成
形如 $X'=X\otimes_A K$ 的预概形，其中 $X$ 为 $A$-预概形。

于是可以考虑模 $\mathfrak J$ 约化的预概形 $X_0=X\otimes_A k$；
它也就是闭点 $y$ 的纤维 $\psi^{-1}(y)$。
迄今由于缺乏适当的术语，人们一直避免明确引入 $A$-预概形 $X$。
但应当注意，通常关于“模 $\mathfrak J$ 约化”预概形 $X_0$ 所作的一切断言，
都应当看成关于 $X$ 本身的更完整断言的推论；
只有按照这种方式解释，才能令人满意地表述和理解这些断言。
而且，所作的假设似乎总归结为关于 $X$ 本身的假设，
与事先给出 $X'$ 到某个 $P'$ 的浸入无关，因而可以得到更加内在的陈述。

\egapar{3.7.3}
最后指出一个很特殊的事实，它无疑曾使这里所讨论情形的概念澄清受到延缓。
若 $A$ 是离散赋值环，且 $X$ 在 $A$ 上固有
（当 $X$ 是某个 $\mathbf P_A^r$ 的闭子预概形时确实如此，参见 II，5.5.4），
则 $X$ 的 $A$-值点与 $X'$ 的 $K$-值点一一对应（II，7.3.8）。
正因为如此，人们常以为自己证明了关于 $X'$ 的结果，
实际却是在证明关于 $X$ 的陈述；
这些陈述以此形式表达时，在不再假设基局部环为一维的情况下仍然成立。


% END chapters/ch1_s3_7.tex

\par

% BEGIN chapters/ch1_s4_1.tex
\section{子预概形与浸入态射}\hypertarget{heading:1:4}{}
\subsection{子预概形}\hypertarget{heading:1:4.1}{}
\src{119—122}

\egapar{4.1.1}
由于拟凝聚层的概念（0，\hyperlink{para:0:5.1.3}{5.1.3}）具有局部性，预概形 $X$ 上的拟凝聚 $\OO_X$-模 $\F$
可以由下述条件刻画：对 $X$ 的每个仿射开集 $V$，$\F|V$ 同构于与某个
$\Gamma(V,\OO_X)$-模相伴的层（\hyperlink{para:1:1.4.1}{1.4.1}）。显然，在预概形 $X$ 上，结构层 $\OO_X$
是拟凝聚的；拟凝聚 $\OO_X$-模之间同态的核、余核、像，以及拟凝聚
$\OO_X$-模的归纳极限和直和，也都是拟凝聚的（\hyperlink{para:1:1.3.7}{1.3.7} 和 \hyperlink{para:1:1.3.9}{1.3.9}）。

\egastatement{4.1.2}{命题}
设 $X$ 为预概形，$\J$ 为 $\OO_X$ 中的拟凝聚理想层。
则层 $\OO_X/\J$ 的支集 $Y$ 是闭集；若以 $\OO_Y$ 表示 $\OO_X/\J$ 在 $Y$ 上的限制，
则 $(Y,\OO_Y)$ 是预概形。

显然，只需考虑 $X$ 为仿射概形的情形（\hyperlink{para:1:2.1.3}{2.1.3}），并证明此时 $Y$ 是 $X$ 的闭子集且为仿射概形。
事实上，若 $X=\Spec(A)$，则 $\OO_X=\wt A$，$\J=\wt{\mathfrak J}$，
其中 $\mathfrak J$ 为 $A$ 的理想（\hyperlink{para:1:1.4.1}{1.4.1}）。此时 $Y=V(\mathfrak J)$，
是 $X$ 的闭子集，并与环 $B=A/\mathfrak J$ 的素谱相认同（\hyperlink{para:1:1.1.11}{1.1.11}）。
此外，若 $\varphi:A\to B=A/\mathfrak J$ 为典范同态，
则直像 ${}^a\varphi_*(\wt B)$ 典范地认同于层
$\wt A/\wt{\mathfrak J}=\OO_X/\J$（\hyperlink{para:1:1.6.3}{1.6.3} 和 \hyperlink{para:1:1.3.9}{1.3.9}），证明完成。

称 $(Y,\OO_Y)$ 为由理想层 $\J$ 定义的 $(X,\OO_X)$ 的\emph{子预概形}。
这是下述一般子预概形概念的一个特例。

\egastatement{4.1.3}{定义}
称赋环空间 $(Y,\OO_Y)$ 为预概形 $(X,\OO_X)$ 的\emph{子预概形}，如果：
\begin{enumerate}[label=\arabic*°]
\item $Y$ 是 $X$ 的局部闭子空间；
\item 以 $U$ 表示 $X$ 中包含 $Y$ 且使 $Y$ 在其中为闭集的最大开集
（即 $X$ 中 $\overline Y\setminus Y$ 的补集），
则 $(Y,\OO_Y)$ 是由 $\OO_X|U$ 的一个拟凝聚理想层定义的 $(U,\OO_X|U)$ 的子预概形。
\end{enumerate}
若 $Y$ 在 $X$ 中为闭集，则称子预概形 $(Y,\OO_Y)$ 是\emph{闭的}；此时 $U=X$。

由此定义及（\hyperlink{para:1:4.1.2}{4.1.2}）立即可知，$X$ 的闭子预概形与 $\OO_X$ 的拟凝聚理想层 $\J$
之间有典范的一一对应。事实上，若两个这样的理想层 $\J,\J'$
所定义的商层具有同一个闭支集 $Y$，且 $\OO_X/\J$ 与 $\OO_X/\J'$ 在 $Y$ 上的限制相同，
则立即得到 $\J'=\J$。\transnote{这里的“相同”保留商结构中的认同，
不是只要求两个商层作为抽象环层同构。}

\egapar{4.1.4}
设 $(Y,\OO_Y)$ 是 $X$ 的子预概形，$U$ 为包含 $Y$ 且使 $Y$ 在其中闭的最大开集，
$V$ 为 $X$ 中包含于 $U$ 的开集。则 $V\cap Y$ 在 $V$ 中闭。
此外，若 $Y$ 由 $\OO_X|U$ 的拟凝聚理想层 $\J$ 定义，
则 $\J|V$ 是 $\OO_X|V$ 的拟凝聚理想层；立即可见，
$Y$ 在 $Y\cap V$ 上诱导的预概形正是由理想层 $\J|V$ 定义的 $V$ 的闭子预概形。反之：

\egastatement{4.1.5}{命题}
设 $(Y,\OO_Y)$ 是赋环空间，$Y$ 为 $X$ 的子空间。
假设存在由 $X$ 的开集组成的 $Y$ 的覆盖 $(V_\alpha)$，使得对每个 $\alpha$，
$Y\cap V_\alpha$ 在 $V_\alpha$ 中闭，且赋环空间
$(Y\cap V_\alpha,\OO_Y|(Y\cap V_\alpha))$ 是 $X$ 在 $V_\alpha$ 上诱导的预概形的闭子预概形。
则 $(Y,\OO_Y)$ 是 $X$ 的子预概形。

假设蕴含 $Y$ 在 $X$ 中局部闭，且包含 $Y$ 并使 $Y$ 在其中闭的最大开集 $U$
包含所有 $V_\alpha$。因而可以归结到 $U=X$、$Y$ 在 $X$ 中闭的情形。
定义 $\OO_X$ 的拟凝聚理想层 $\J$，使 $\J|V_\alpha$ 为定义闭子预概形
$(Y\cap V_\alpha,\OO_Y|(Y\cap V_\alpha))$ 的 $\OO_X|V_\alpha$-理想层，
并对每个不与 $Y$ 相交的开集 $W\subset X$ 令 $\J|W=\OO_X|W$。
由（\hyperlink{para:1:4.1.3}{4.1.3}）和（\hyperlink{para:1:4.1.4}{4.1.4}）立即验证，满足这些条件的理想层 $\J$ 存在且唯一，
并且它定义闭子预概形 $(Y,\OO_Y)$。

特别地，$X$ 在其一个开集上诱导的预概形是 $X$ 的子预概形。

\egastatement{4.1.6}{命题}
$X$ 的一个子预概形（分别地，闭子预概形）的子预概形（分别地，闭子预概形），
典范地认同于 $X$ 的一个子预概形（分别地，闭子预概形）。

$X$ 的局部闭子空间的局部闭子集仍是 $X$ 的局部闭子空间。
因而问题显然具有局部性（\hyperlink{para:1:4.1.5}{4.1.5}），可以假设 $X$ 仿射。
命题遂由下述典范认同得到：若 $\mathfrak J\subset\mathfrak J'$ 是环 $A$ 的两个理想，则
\[
 A/\mathfrak J'\simeq (A/\mathfrak J)/(\mathfrak J'/\mathfrak J).
\]
以下总作上述认同。

\egapar{4.1.7}
设 $Y$ 是预概形 $X$ 的子预概形，以 $\psi$ 表示底空间的典范嵌入 $Y\to X$。
已知逆像 $\psi^*(\OO_X)$ 为限制 $\OO_X|Y$（0，\hyperlink{para:0:3.7.1}{3.7.1}）。
若对每个 $y\in Y$，以 $\omega_y$ 表示典范同态
$(\OO_X)_y\to(\OO_Y)_y$，则这些同态是某个满射环层同态
$\omega:\OO_X|Y\to\OO_Y$ 在各茎上的同态。
只需在 $Y$ 上局部验证此事，因而可假设 $X$ 仿射且 $Y$ 为闭子预概形；
若此时 $\J$ 是定义 $Y$ 的理想层，则 $\omega_y$ 正是
$\OO_X|Y\to(\OO_X/\J)|Y$ 在各茎上的同态。

于是定义了赋环空间的单态射（0，\hyperlink{para:0:4.1.1}{4.1.1}）$j=(\psi,\omega^\flat)$。
它显然是预概形态射 $Y\to X$（\hyperlink{para:1:2.2.1}{2.2.1}），称为\emph{典范嵌入态射}。
若 $f:X\to Z$ 为态射，则也称复合态射 $Y\to X\to Z$ 为 $f$ 在子预概形 $Y$ 上的\emph{限制}。

\egapar{4.1.8}
按照一般定义（T，I，1.1），若预概形态射 $f:Z\to X$ 可分解为
$Z\xrightarrow{g}Y\xrightarrow{j}X$，其中 $Y$ 是 $X$ 的子预概形，
$j$ 为嵌入态射，$g$ 为预概形态射，则称 $f$ \emph{不大于} $j$。
由于 $j$ 是单态射，$g$ 必然唯一。

\egastatement{4.1.9}{命题}
态射 $f:Z\to X$ 不大于嵌入态射 $j:Y\to X$ 的充要条件是：
$f(Z)\subset Y$，且对每个 $z\in Z$，令 $y=f(z)$，由 $f$ 给出的同态
$(\OO_X)_y\to(\OO_Z)_z$ 可分解为
$(\OO_X)_y\to(\OO_Y)_y\to(\OO_Z)_z$；
等价地，前一同态的核包含 $(\OO_X)_y\to(\OO_Y)_y$ 的核。

条件显然必要。为证明充分性，必要时以一个使 $Y$ 在其中闭的开集 $U$ 代替 $X$（\hyperlink{para:1:4.1.3}{4.1.3}），
即可假设 $Y$ 是 $X$ 的闭子预概形。它由 $\OO_X$ 的拟凝聚理想层 $\J$ 定义。
写 $f=(\psi,\theta)$，并以 $\mathscr I$ 表示
$\theta^\sharp:\psi^*(\OO_X)\to\OO_Z$ 的核理想层。
由函子 $\psi^*$ 的性质（0，\hyperlink{para:0:3.7.2}{3.7.2}），假设蕴含
$(\psi^*(\J))_z\subset\mathscr I_z$ 对每个 $z\in Z$ 成立，故 $\psi^*(\J)\subset\mathscr I$。
因此 $\theta^\sharp$ 分解为
\[
 \psi^*(\OO_X)\longrightarrow
 \psi^*(\OO_X)/\psi^*(\J)=\psi^*(\OO_X/\J)
 \xrightarrow{\omega}\OO_Z,
\]
其中第一箭头为典范同态。以 $\psi':Z\to Y$ 表示与 $\psi$ 一致的连续映射，
显然 $\psi'^*(\OO_Y)=\psi^*(\OO_X/\J)$。
另一方面，$\omega$ 显然在茎上是局部同态，故
$g=(\psi',\omega^\flat)$ 是预概形态射 $Z\to Y$（\hyperlink{para:1:2.2.1}{2.2.1}）；
由以上所述 $f=j\circ g$，命题得证。

\egastatement{4.1.10}{推论}
嵌入态射 $Z\to X$ 不大于嵌入态射 $Y\to X$ 的充要条件是 $Z$ 为 $Y$ 的子预概形。

此时记作 $Z\leq Y$；这显然是 $X$ 的子预概形集合上的一个序关系。


% END chapters/ch1_s4_1.tex

\par

% BEGIN chapters/ch1_s4_2.tex
\subsection{浸入态射}\hypertarget{heading:1:4.2}{}
\src{122—124}

\egastatement{4.2.1}{定义}
若态射 $f:Y\to X$ 可分解为 $Y\xrightarrow{g}Z\xrightarrow{j}X$，
其中 $g$ 为同构，$Z$ 为 $X$ 的子预概形（分别地，闭子预概形、在一个开集上诱导的子预概形），
$j$ 为嵌入态射，则称 $f$ 为\emph{浸入}（分别地，\emph{闭浸入}、\emph{开浸入}）。

子预概形 $Z$ 和同构 $g$ 此时被唯一确定。
事实上，若 $Z'$ 为另一个子预概形，$j':Z'\to X$ 为嵌入，$g':Y\to Z'$ 为同构，
且 $j\circ g=j'\circ g'$，则 $j'=j\circ g\circ g'^{-1}$，从而 $Z'\leq Z$（\hyperlink{para:1:4.1.10}{4.1.10}）。
同样得到 $Z\leq Z'$，故 $Z'=Z$；又因 $j$ 是单态射，$g'=g$。
称 $f=j\circ g$ 为浸入 $f$ 的\emph{典范分解}，称 $Z$ 和 $g$ 为与 $f$ \emph{相伴}的子预概形和同构。

显然，浸入是预概形的单态射（\hyperlink{para:1:4.1.7}{4.1.7}），从而是根式态射（\hyperlink{para:1:3.5.4}{3.5.4}）。

\egastatement{4.2.2}{命题}
\begin{enumerate}[label=\alph*)]
\item 态射 $f=(\psi,\theta):Y\to X$ 为开浸入的充要条件是：
$\psi$ 是 $Y$ 到 $X$ 的一个开子集的同胚，且对每个 $y\in Y$，
同态 $\theta_y^\sharp:\OO_{\psi(y)}\to\OO_y$ 为双射。
\item 态射 $f=(\psi,\theta):Y\to X$ 为浸入（分别地，闭浸入）的充要条件是：
$\psi$ 是 $Y$ 到 $X$ 的一个局部闭子集（分别地，闭子集）的同胚，
且对每个 $y\in Y$，同态 $\theta_y^\sharp:\OO_{\psi(y)}\to\OO_y$ 为满射。
\end{enumerate}

 a) 条件显然必要。反之，若条件成立，$\theta^\sharp$ 给出 $\psi^*(\OO_X)$ 与 $\OO_Y$ 的同构。
而 $\psi^*(\OO_X)$ 是从 $\OO_X|\psi(Y)$ 经 $\psi^{-1}$ 搬运结构所得的层，故结论成立。

 b) 条件显然必要，现证明充分性。先设 $X$ 仿射且 $Z=\psi(Y)$ 在 $X$ 中闭。
已知 $\psi_*(\OO_Y)$ 的支集为 $Z$（0，\hyperlink{para:0:3.4.6}{3.4.6}）；若以 $\OO'_Z$ 表示它在 $Z$ 上的限制，
则赋环空间 $(Z,\OO'_Z)$ 是从 $(Y,\OO_Y)$ 经同胚 $\psi:Y\to Z$ 搬运结构所得的。
先证明 $f_*(\OO_Y)=\psi_*(\OO_Y)$ 是拟凝聚 $\OO_X$-模。
对 $x\notin Z$，$\psi_*(\OO_Y)$ 在 $x$ 的适当邻域上为零。
若 $x\in Z$，则有唯一 $y\in Y$ 使 $x=\psi(y)$。
取 $y$ 的仿射开邻域 $V$；$\psi(V)$ 在 $Z$ 中开，故为某个开集 $U\subset X$ 在 $Z$ 上的迹。
$\psi_*(\OO_Y)$ 在 $U$ 上的限制等于直像
$(\psi_V)_*(\OO_Y|V)$ 在 $U$ 上的限制，其中 $\psi_V=\psi|V$。
态射 $(\psi,\theta)$ 在 $(V,\OO_Y|V)$ 上的限制是到 $(X,\OO_X)$ 的态射，
因而形如 $({}^a\varphi,\wt\varphi)$，其中
$\varphi:A=\Gamma(X,\OO_X)\to\Gamma(V,\OO_Y)$ 为环同态（\hyperlink{para:1:1.7.3}{1.7.3}）。
所以 $(\psi_V)_*(\OO_Y|V)$ 是拟凝聚 $\OO_X$-模（\hyperlink{para:1:1.6.3}{1.6.3}）；
拟凝聚性具有局部性，故所断言的拟凝聚性得证。

此外，由 $\psi$ 的同胚性质，对每个 $y\in Y$，典范映射
$\widetilde\psi_y:(\psi_*(\OO_Y))_{\psi(y)}\to\OO_y$ 为同构（0，\hyperlink{para:0:3.4.5}{3.4.5}）。
交换图
\[
\begin{tikzcd}[column sep=large,row sep=large]
\OO_{\psi(y)} \arrow[r,"\theta_{\psi(y)}"] \arrow[d,"\widetilde\psi_y\circ\alpha_{\psi(y)}"'] &
(\psi_*(\OO_Y))_{\psi(y)} \arrow[d,"\widetilde\psi_y"]\\
(\psi^*(\OO_X))_y \arrow[r,"\theta_y^\sharp"'] & \OO_y
\end{tikzcd}
\]
的两个竖箭头都是同构（0，\hyperlink{para:0:3.7.2}{3.7.2}），所以 $\theta_y^\sharp$ 的满射性蕴含
$\theta_{\psi(y)}$ 的满射性。
由于 $\psi_*(\OO_Y)$ 的支集为 $Z=\psi(Y)$，$\theta$ 是从
$\OO_X=\wt A$ 到拟凝聚 $\OO_X$-模 $f_*(\OO_Y)$ 的满同态。
于是存在唯一同构 $\omega:\wt A/\wt{\mathfrak J}\to f_*(\OO_Y)$，
其中 $\mathfrak J$ 为 $A$ 的一个理想，使它与典范同态
$\wt A\to\wt A/\wt{\mathfrak J}$ 的复合等于 $\theta$（\hyperlink{para:1:1.3.8}{1.3.8}）。
若 $\OO_Z$ 是 $\wt A/\wt{\mathfrak J}$ 在 $Z$ 上的限制，则 $(Z,\OO_Z)$ 是 $(X,\OO_X)$ 的子预概形。
$f$ 分解为此子预概形到 $X$ 的典范嵌入与同构 $(\psi_0,\omega_0)$ 的复合，
其中 $\psi_0$ 是把 $\psi$ 看成 $Y\to Z$ 的映射，$\omega_0$ 为 $\omega$ 在 $\OO_Z$ 上的限制。

再考虑一般情形。设 $U$ 为 $X$ 的仿射开集，且 $U\cap\psi(Y)$ 非空并在 $U$ 中闭。
把 $f$ 限制到 $Y$ 在 $\psi^{-1}(U)$ 上诱导的预概形，并看作到 $U$ 的态射，即归结为第一种情形。
故限制是闭浸入 $\psi^{-1}(U)\to U$，其典范分解为 $j_U\circ g_U$，
其中 $g_U$ 是 $\psi^{-1}(U)$ 到 $U$ 的一个子预概形 $Z_U$ 的同构，$j_U:Z_U\to U$ 为嵌入。
若 $V\subset U$ 是另一个仿射开集，则 $Z_U$ 在 $V$ 上的限制 $Z'_V$ 是 $V$ 的子预概形；
$f|\psi^{-1}(V)$ 因而分解为 $j'_V\circ g'_V$。
由典范分解的唯一性（\hyperlink{para:1:4.2.1}{4.2.1}），$Z'_V=Z_V$，$g'_V=g_V$。
所以存在 $X$ 的子预概形 $Z$（\hyperlink{para:1:4.1.5}{4.1.5}），其底空间为 $\psi(Y)$，
在每个 $U\cap\psi(Y)$ 上的限制为 $Z_U$。
各 $g_U$ 是一个同构 $g:Y\to Z$ 在 $\psi^{-1}(U)$ 上的限制，且 $f=j\circ g$，
其中 $j:Z\to X$ 为典范嵌入。

\egastatement{4.2.3}{推论}
设 $X$ 为仿射概形。态射 $f=(\psi,\theta):Y\to X$ 为闭浸入的充要条件是：
$Y$ 为仿射概形，且同态 $\Gamma(\theta):\Gamma(\OO_X)\to\Gamma(\OO_Y)$ 为满射。
\transnote{原版此处写作 $\Gamma(\psi)$；这里按结构层同态 $\theta$ 写作 $\Gamma(\theta)$。}

\egastatement{4.2.4}{推论}
\begin{enumerate}[label=\alph*)]
\item 设 $f:Y\to X$ 为态射，$(V_\lambda)$ 是由 $X$ 的开集组成的 $f(Y)$ 的覆盖。
则 $f$ 为浸入（分别地，开浸入）的充要条件是：
它在每个诱导预概形 $f^{-1}(V_\lambda)$ 上的限制是到 $V_\lambda$ 的浸入（分别地，开浸入）。
\item 设 $f:Y\to X$ 为态射，$(V_\lambda)$ 是 $X$ 的开覆盖。
则 $f$ 为闭浸入的充要条件是：它在每个 $f^{-1}(V_\lambda)$ 上的限制是到 $V_\lambda$ 的闭浸入。
\end{enumerate}

写 $f=(\psi,\theta)$。在 a) 中，各 $\theta_y^\sharp$ 为满射（分别地，双射）；
在 b) 中，它们为满射。故只须分别验证 $\psi$ 是到局部闭子集（分别地，开子集）或到闭子集的同胚。
由假设，$\psi$ 显然为单射，且对每个 $y\in Y$，它把 $y$ 的任一邻域映为
$\psi(y)$ 在 $\psi(Y)$ 中的邻域。
在 a) 中，$\psi(Y)\cap V_\lambda$ 在 $V_\lambda$ 中局部闭（分别地，开），
故 $\psi(Y)$ 在 $\bigcup V_\lambda$ 中局部闭（分别地，开），从而在 $X$ 中也如此。
在 b) 中，$\psi(Y)\cap V_\lambda$ 在 $V_\lambda$ 中闭，而 $X=\bigcup V_\lambda$，
故 $\psi(Y)$ 在 $X$ 中闭。

\egastatement{4.2.5}{命题}
两个浸入（分别地，两个开浸入、两个闭浸入）的复合仍是浸入（分别地，开浸入、闭浸入）。

这由（\hyperlink{para:1:4.1.6}{4.1.6}）直接得到。


% END chapters/ch1_s4_2.tex

\par

% BEGIN chapters/ch1_s4_3.tex
\subsection{浸入的积}\hypertarget{heading:1:4.3}{}
\src{124—125}

\egastatement{4.3.1}{命题}
设 $\alpha:X'\to X$、$\beta:Y'\to Y$ 为两个 $S$-态射。
若 $\alpha,\beta$ 是浸入（分别地，开浸入、闭浸入），则 $\alpha\times_S\beta$
也是浸入（分别地，开浸入、闭浸入）。
此外，若 $\alpha$、$\beta$ 分别把 $X'$、$Y'$ 认同于 $X$、$Y$ 的子预概形 $X''$、$Y''$，
则 $\alpha\times_S\beta$ 把 $X'\times_S Y'$ 的底空间认同于 $X\times_S Y$ 的底空间中的子空间
$p^{-1}(X'')\cap q^{-1}(Y'')$，其中 $p,q$ 为到 $X,Y$ 的投影。

由定义（\hyperlink{para:1:4.2.1}{4.2.1}），可限于 $X',Y'$ 是子预概形而 $\alpha,\beta$ 是嵌入态射的情形。
对于在开集上诱导的子预概形，命题已在（\hyperlink{para:1:3.2.7}{3.2.7}）中证明。
每个子预概形都是某个开集上诱导的预概形的闭子预概形（\hyperlink{para:1:4.1.3}{4.1.3}），
故归结到 $X',Y'$ 为闭子预概形的情形。

首先证明可以假设 $S$ 仿射。取 $S$ 的仿射开覆盖 $(S_\lambda)$，
令 $X_\lambda=\varphi^{-1}(S_\lambda)$、$Y_\lambda=\psi^{-1}(S_\lambda)$，
其中 $\varphi,\psi$ 为结构态射。
$X',Y'$ 分别在 $X_\lambda\cap X'$、$Y_\lambda\cap Y'$ 上的限制 $X'_\lambda,Y'_\lambda$
是 $X_\lambda,Y_\lambda$ 的闭子预概形。
这些预概形都可看作 $S_\lambda$-预概形；相应在 $S$ 和 $S_\lambda$ 上的积相同（\hyperlink{para:1:3.2.5}{3.2.5}）。
若命题在 $S$ 仿射时成立，则 $\alpha\times_S\beta$ 在每个
$X'_\lambda\times_S Y'_\lambda$ 上的限制都是浸入（\hyperlink{para:1:3.2.7}{3.2.7}）。
又由于混合积 $X'_\lambda\times_S Y'_\mu$ 和 $X_\lambda\times_S Y_\mu$
分别认同于
\[
 (X'_\lambda\cap X'_\mu)\times_S(Y'_\lambda\cap Y'_\mu),\qquad
 (X_\lambda\cap X_\mu)\times_S(Y_\lambda\cap Y_\mu)
\]
（\hyperlink{para:1:3.2.6.4}{3.2.6.4}），在每个 $X'_\lambda\times_S Y'_\mu$ 上的限制仍为浸入。
因此由（\hyperlink{para:1:4.2.4}{4.2.4}），$\alpha\times_S\beta$ 也为浸入。

其次证明还可以假设 $X,Y$ 仿射。取它们的仿射开覆盖 $(U_i)$、$(V_j)$，
并以 $X'_i,Y'_j$ 表示 $X',Y'$ 在 $X'\cap U_i,Y'\cap V_j$ 上的限制。
它们分别是 $U_i,V_j$ 的闭子预概形。
$U_i\times_S V_j$ 认同于 $X\times_S Y$ 在 $p^{-1}(U_i)\cap q^{-1}(V_j)$ 上的限制（\hyperlink{para:1:3.2.7}{3.2.7}）；
同样，$X'_i\times_S Y'_j$ 认同于 $X'\times_S Y'$ 在
$p'^{-1}(X'_i)\cap q'^{-1}(Y'_j)$ 上的限制。
置 $\gamma=\alpha\times_S\beta$；由定义，$p\circ\gamma=\alpha\circ p'$、
$q\circ\gamma=\beta\circ q'$。因 $X'_i=\alpha^{-1}(U_i)$、$Y'_j=\beta^{-1}(V_j)$，
\[
 p'^{-1}(X'_i)\cap q'^{-1}(Y'_j)
 =\gamma^{-1}\bigl(p^{-1}(U_i)\cap q^{-1}(V_j)\bigr)
 =\gamma^{-1}(U_i\times_S V_j).
\]
遂可如第一部分那样得出结论。

现设 $X,Y,S$ 仿射，其环分别为 $B,C,A$。
$B,C$ 为 $A$-代数，$X',Y'$ 是环分别为 $B,C$ 的商代数 $B',C'$ 的仿射概形。
此外 $\alpha=({}^a\rho,\wt\rho)$、$\beta=({}^a\sigma,\wt\sigma)$，
其中 $\rho:B\to B'$、$\sigma:C\to C'$ 为典范同态（\hyperlink{para:1:1.7.3}{1.7.3}）。
已知 $X\times_S Y$、$X'\times_S Y'$ 分别为环 $B\otimes_A C$、$B'\otimes_A C'$ 的仿射概形，
且 $\alpha\times_S\beta=({}^a\tau,\wt\tau)$，其中 $\tau=\rho\otimes\sigma$（\hyperlink{para:1:3.2.2}{3.2.2} 和 \hyperlink{para:1:3.2.3}{3.2.3}）。
该同态为满射，故 $\alpha\times_S\beta$ 确为浸入。
若 $\mathfrak b=\Ker\rho$、$\mathfrak c=\Ker\sigma$，且
$u(b)=b\otimes1$、$v(c)=1\otimes c$，则 $\Ker\tau$ 为
$u(\mathfrak b)$ 和 $v(\mathfrak c)$ 所生成的理想之和。
\transnote{原文将此核简写为 $u(\mathfrak b)+v(\mathfrak c)$；
这里按生成理想理解，即
$\Im(\mathfrak b\otimes_A C)+\Im(B\otimes_A\mathfrak c)$，而不是两个集合像的简单和。}
由于 $p=({}^a u,\wt u)$、$q=({}^a v,\wt v)$，
此核在 $\Spec(B\otimes_A C)$ 中对应闭集 $p^{-1}(X')\cap q^{-1}(Y')$
（\hyperref[eq:1:1.2.2.1]{1.2.2.1} 和 \hyperlink{para:1:1.1.2}{1.1.2}，(iii)），证明完成。

\egastatement{4.3.2}{推论}
若 $f:X\to Y$ 是浸入（分别地，开浸入、闭浸入）且为 $S$-态射，
则对任意基预概形扩张 $S'\to S$，$f_{(S')}$ 也是浸入（分别地，开浸入、闭浸入）。


% END chapters/ch1_s4_3.tex

\par

% BEGIN chapters/ch1_s4_4.tex
\subsection{子预概形的逆像}\hypertarget{heading:1:4.4}{}
\src{125—126}

\egastatement{4.4.1}{命题}
设 $f:X\to Y$ 为态射，$Y'$ 为 $Y$ 的子预概形
（分别地，闭子预概形、在一个开集上诱导的预概形），$j:Y'\to Y$ 为嵌入态射。
则投影 $p:X\times_Y Y'\to X$ 是浸入（分别地，闭浸入、开浸入）；
与 $p$ 相伴的 $X$ 的子预概形的底空间为 $f^{-1}(Y')$。
此外，若 $j'$ 是此子预概形到 $X$ 的嵌入态射，
则态射 $h:Z\to X$ 满足 $f\circ h$ 不大于 $j$ 的充要条件是 $h$ 不大于 $j'$。

由于 $p=1_X\times_Y j$（\hyperlink{para:1:3.3.4}{3.3.4}），第一断言由（\hyperlink{para:1:4.3.1}{4.3.1}）得到；
第二断言是（\hyperlink{para:1:3.5.10}{3.5.10}）的特例，交换其中 $X$ 与 $Y'$ 的角色即可。
最后，若 $f\circ h=j\circ h'$，其中 $h':Z\to Y'$，
则由积的定义，$h=p\circ u$，其中 $u:Z\to X\times_Y Y'$，从而得到最后的断言。

称如此定义的 $X$ 的子预概形为 $Y$ 的子预概形 $Y'$ 在态射 $f$ 下的\emph{逆像}，
与（\hyperlink{para:1:3.3.6}{3.3.6}）中更一般的术语一致。
以后把 $f^{-1}(Y')$ 看作 $X$ 的子预概形时，总指此子预概形。

若预概形 $f^{-1}(Y')$ 与 $X$ 相同，则 $j'$ 为恒等态射，故任意 $h:Z\to X$ 均不大于 $j'$；
因此态射 $f:X\to Y$ 分解为 $X\xrightarrow{g}Y'\xrightarrow{j}Y$。

若 $y$ 为 $Y$ 的闭点，且 $Y'=\Spec(k(y))$ 为以 $\{y\}$ 为底空间的最小闭子预概形（\hyperlink{para:1:4.1.9}{4.1.9}），
则闭子预概形 $f^{-1}(Y')$ 典范地同构于（\hyperlink{para:1:3.6.2}{3.6.2}）定义的纤维 $f^{-1}(y)$，以后将它们认同。

\egastatement{4.4.2}{推论}
设 $f:X\to Y$、$g:Y\to Z$ 为态射，$h=g\circ f$。
对 $Z$ 的每个子预概形 $Z'$，$X$ 的子预概形 $f^{-1}(g^{-1}(Z'))$ 与 $h^{-1}(Z')$ 相同。

这由典范同构 $X\times_Y(Y\times_Z Z')\simeq X\times_Z Z'$ 的存在得到（\hyperref[eq:1:3.3.9.1]{3.3.9.1}）。

\egastatement{4.4.3}{推论}
设 $X',X''$ 是 $X$ 的两个子预概形，$j':X'\to X$、$j'':X''\to X$ 为嵌入态射。
则 $j'^{-1}(X'')$ 与 $j''^{-1}(X')$ 均等于子预概形序关系下 $X',X''$ 的下确界
$\inf(X',X'')$，并典范地同构于 $X'\times_X X''$。

这由（\hyperlink{para:1:4.4.1}{4.4.1}）及（\hyperlink{para:1:4.1.10}{4.1.10}）立即得到。

\egastatement{4.4.4}{推论}
设 $f:X\to Y$ 为态射，$Y',Y''$ 是 $Y$ 的子预概形。则
\[
 f^{-1}(\inf(Y',Y''))=\inf(f^{-1}(Y'),f^{-1}(Y'')).
\]
这是因为 $(X\times_Y Y')\times_X(X\times_Y Y'')$ 与
$X\times_Y(Y'\times_Y Y'')$ 之间存在典范同构（\hyperref[eq:1:3.3.9.1]{3.3.9.1}）。

\egastatement{4.4.5}{命题}
设 $f:X\to Y$ 为态射，$Y'$ 是由 $\OO_Y$ 的拟凝聚理想层 $\mathscr K$
定义的闭子预概形（\hyperlink{para:1:4.1.3}{4.1.3}）。则 $X$ 的闭子预概形 $f^{-1}(Y')$
由 $\OO_X$ 的拟凝聚理想层 $f^*(\mathscr K)\OO_X$ 定义。

问题显然在 $X,Y$ 上具有局部性。
只须注意：若 $A$ 为 $B$-代数，$\mathfrak K$ 为 $B$ 的理想，
则 $A\otimes_B(B/\mathfrak K)=A/\mathfrak K A$，再应用（\hyperlink{para:1:1.6.9}{1.6.9}）即可。
\transnote{原文此处写“$B$ 是 $A$-代数”，与紧接的 $A\otimes_B(B/\mathfrak K)$ 不合；
译文按该公式所需的结构写为“$A$ 是 $B$-代数”。}

\egastatement{4.4.6}{推论}
设 $X'$ 是由 $\OO_X$ 的拟凝聚理想层 $\J$ 定义的闭子预概形，$i:X'\to X$ 为嵌入。
$f$ 在 $X'$ 上的限制 $f\circ i$ 不大于嵌入 $j:Y'\to Y$
（即分解为 $j\circ g$，其中 $g:X'\to Y'$）的充要条件是
$f^*(\mathscr K)\OO_X\subset\J$。

只须将（\hyperlink{para:1:4.4.1}{4.4.1}）应用于 $i$，并利用（\hyperlink{para:1:4.4.5}{4.4.5}）。


% END chapters/ch1_s4_4.tex

\par

% BEGIN chapters/ch1_s4_5.tex
\subsection{局部浸入与局部同构}\hypertarget{heading:1:4.5}{}
\src{126—127}

\egastatement{4.5.1}{定义}
设 $f:X\to Y$ 为预概形态射。
若存在 $x\in X$ 的开邻域 $U$ 及 $f(x)$ 在 $Y$ 中的开邻域 $V$，
使 $f$ 在诱导预概形 $U$ 上的限制是到诱导预概形 $V$ 的闭浸入，
则称 $f$ 在 $x$ 处为\emph{局部浸入}。
若 $f$ 在 $X$ 的每一点处均为局部浸入，则称 $f$ 为\emph{局部浸入}。

\egastatement{4.5.2}{定义}
若存在 $x\in X$ 的开邻域 $U$，使态射 $f:X\to Y$ 在诱导预概形 $U$ 上的限制
为到 $Y$ 的开浸入，则称 $f$ 在 $x$ 处为\emph{局部同构}。
若 $f$ 在 $X$ 的每一点处均为局部同构，则称 $f$ 为\emph{局部同构}。

\egapar{4.5.3}
因此，浸入（分别地，闭浸入）$f:X\to Y$ 可以刻画为这样一个局部浸入：
其底空间映射是 $X$ 到 $Y$ 的一个子集（分别地，闭子集）的同胚。
开浸入则可以刻画为单射的局部同构。

\egastatement{4.5.4}{命题}
设 $X$ 为不可约预概形，$f:X\to Y$ 为单射且支配的态射。
若 $f$ 是局部浸入，则 $f$ 是浸入，且 $f(X)$ 在 $Y$ 中开。

取 $x\in X$，并取 $x$ 的开邻域 $U$、$f(x)$ 的开邻域 $V$，
使 $f|U$ 为到 $V$ 的闭浸入。
由于 $U$ 在 $X$ 中稠密，由假设 $f(U)$ 在 $Y$ 中稠密，故 $f(U)=V$，
且 $f$ 是 $U$ 到 $V$ 的同胚。
$f$ 的单射性蕴含 $f^{-1}(V)=U$，命题遂立即得到。

\egastatement{4.5.5}{命题}
\begin{enumerate}[label=(\roman*)]
\item 两个局部浸入（分别地，局部同构）的复合仍是局部浸入（分别地，局部同构）。
\item 设 $f:X\to X'$、$g:Y\to Y'$ 为两个 $S$-态射。
若它们是局部浸入（分别地，局部同构），则 $f\times_S g$ 也如此。
\item 若 $S$-态射 $f$ 是局部浸入（分别地，局部同构），
则任意基预概形扩张 $S'\to S$ 所得的 $f_{(S')}$ 也如此。
\end{enumerate}

根据（\hyperlink{para:1:3.5.1}{3.5.1}），只须证明 (i) 和 (ii)。
(i) 由闭浸入（分别地，开浸入）的传递性，以及下述事实立即得到：
若 $f$ 是 $X$ 到 $Y$ 的一个闭子集的同胚，则对每个开集 $U\subset X$，
$f(U)$ 在 $f(X)$ 中开，故存在开集 $V\subset Y$ 使
$f(U)=V\cap f(X)$，于是 $f(U)$ 在 $V$ 中闭。
\transnote{原文将浸入的传递性引用为（\hyperlink{para:1:4.2.4}{4.2.4}）；直接陈述该性质的是（\hyperlink{para:1:4.2.5}{4.2.5}）。}

为证明 (ii)，以 $p,q$ 表示 $X\times_S Y$ 的投影，$p',q'$ 表示 $X'\times_S Y'$ 的投影。
取 $z\in X\times_S Y$，令 $z'=(f\times_S g)(z)$。
由假设，可分别取 $x=p(z)$、$x'=p'(z')$、$y=q(z)$、$y'=q'(z')$ 的开邻域
$U,U',V,V'$，使 $f|U$、$g|V$ 分别为到 $U',V'$ 的闭浸入（分别地，开浸入）。
$U\times_S V$ 与 $U'\times_S V'$ 的底空间分别认同于 $z,z'$ 的开邻域
$p^{-1}(U)\cap q^{-1}(V)$ 与 $p'^{-1}(U')\cap q'^{-1}(V')$（\hyperlink{para:1:3.2.7}{3.2.7}），
因而命题由（\hyperlink{para:1:4.3.1}{4.3.1}）得到。


% END chapters/ch1_s4_5.tex

\par

% BEGIN chapters/ch1_s5_1.tex
\section{约化预概形；分离条件}\hypertarget{heading:1:5}{}
\subsection{约化预概形}\hypertarget{heading:1:5.1}{}
\src{127—131}

\egastatement{5.1.1}{命题}
设 $(X,\OO_X)$ 为预概形，$\BB$ 为拟凝聚 $\OO_X$-代数。
存在唯一拟凝聚 $\OO_X$-子模 $\N$，使得对每个 $x\in X$，$\N_x$ 为环 $\BB_x$ 的幂零根。
若 $X$ 仿射，因而 $\BB=\wt B$，其中 $B$ 是 $A(X)$-代数，则
$\N=\wt{\mathfrak N}$，其中 $\mathfrak N$ 为 $B$ 的幂零根。

问题具有局部性，故归结到证明最后的断言。
已知 $\wt{\mathfrak N}$ 是拟凝聚 $\OO_X$-模（\hyperlink{para:1:1.4.1}{1.4.1}），
其在 $x$ 处的茎为分式环 $B_x$ 的理想 $\mathfrak N_x$。
只须证明 $B_x$ 的幂零根包含于 $\mathfrak N_x$，反向包含显然。
设 $z/s$ 属于 $B_x$ 的幂零根，其中 $z\in B$，$s\notin\jj_x$。
由假设存在整数 $k$ 使 $(z/s)^k=0$，即存在 $t\notin\jj_x$ 使 $tz^k=0$。
于是 $(tz)^k=0$，故 $z/s=(tz)/(ts)$ 属于 $\mathfrak N_x$。

称如此定义的拟凝聚 $\OO_X$-模 $\N$ 为 $\OO_X$-代数 $\BB$ 的\emph{幂零根}。
特别地，以 $\N_X$ 表示 $\OO_X$ 的幂零根。

\egastatement{5.1.2}{推论}
设 $X$ 为预概形。由理想层 $\N_X$ 定义的 $X$ 的闭子预概形，
是以 $X$ 为底空间的唯一约化子预概形（0，\hyperlink{para:0:4.1.4}{4.1.4}）；
它也是以 $X$ 为底空间的最小子预概形。

由 $\N_X$ 定义的闭子预概形 $Y$ 的结构层为 $\OO_X/\N_X$。
立即可见它是约化的；又因对每个 $x$，$\N_{X,x}\subset\mm_x$，其底空间为 $X$。
为证明其余断言，注意以 $X$ 为底空间的子预概形 $Z$ 由理想层 $\J$
定义（\hyperlink{para:1:4.1.3}{4.1.3}），且每个 $x\in X$ 处 $\J_x\ne\OO_x$。
可以限于仿射情形 $X=\Spec(A)$、$\J=\wt{\mathfrak J}$，其中 $\mathfrak J$ 为 $A$ 的理想。
此时对每个 $x$，$\mathfrak J_x\subset\mm_x$，所以 $\mathfrak J$ 包含于 $A$ 的每个素理想，
从而包含于它们的交，即 $A$ 的幂零根 $\mathfrak N$。
因此 $Y$ 是以 $X$ 为底空间的最小子预概形（\hyperlink{para:1:4.1.9}{4.1.9}）。
若 $Z\ne Y$，则至少存在一个 $x$ 使 $\J_x\ne\N_{X,x}$，
于是由（\hyperlink{para:1:5.1.1}{5.1.1}），$Z$ 不是约化的。

\egastatement{5.1.3}{定义}
称以 $X$ 为底空间的唯一约化子预概形为与 $X$ \emph{相伴的约化预概形}，记为 $X_{\mathrm{red}}$。

所以，预概形 $X$ 约化意味着 $X=X_{\mathrm{red}}$。

\egastatement{5.1.4}{命题}
环 $A$ 的素谱为约化（分别地，整）预概形的充要条件是
$A$ 为约化环（分别地，整环）。

由（\hyperlink{para:1:5.1.1}{5.1.1}）立即可知，$\mathfrak N=(0)$ 是 $X=\Spec(A)$ 约化的充要条件；
关于整环的断言遂由（\hyperlink{para:1:1.1.13}{1.1.13}）得到。
由于整环的每个非零分式环仍为整环，由（\hyperlink{para:1:5.1.4}{5.1.4}）可知，
对每个局部整预概形 $X$ 和每个 $x\in X$，$\OO_x$ 都是整环。
当 $X$ 的底空间局部诺特时，逆命题成立。
事实上，此时 $X$ 约化。若 $U$ 是 $X$ 中作为拓扑空间为诺特的仿射开集，
则 $U$ 只有有限个不可约分支，所以它的环 $A$ 只有有限个极小素理想（\hyperlink{para:1:1.1.14}{1.1.14}）。
若两个这样的分支有公共点 $x$，则 $\OO_x$ 至少有两个不同的极小素理想，因而不是整环。
所以各分支 $U_i$ 是两两不交的开集，并且都是整的。

\egapar{5.1.5}
设 $f=(\psi,\theta):X\to Y$ 为预概形态射。
同态 $\theta_x^\sharp:\OO_{\psi(x)}\to\OO_x$ 将幂零元映为幂零元；
因此取商后得到同态
\[
 \omega:\psi^*(\OO_Y/\N_Y)\longrightarrow\OO_X/\N_X.
\]
显然，各茎上的同态
$\omega_x:\OO_{Y_{\mathrm{red}},\psi(x)}\to\OO_{X_{\mathrm{red}},x}$ 都是局部同态。
故 $(\psi,\omega^\flat)$ 是预概形态射
$X_{\mathrm{red}}\to Y_{\mathrm{red}}$，记为 $f_{\mathrm{red}}$，
称为与 $f$ \emph{相伴的约化态射}。
对两个态射 $f:X\to Y$、$g:Y\to Z$，立即有
$(g\circ f)_{\mathrm{red}}=g_{\mathrm{red}}\circ f_{\mathrm{red}}$。
因此 $X_{\mathrm{red}}$ 是 $X$ 的协变函子。

上述定义说明交换图
\[
\begin{tikzcd}
X_{\mathrm{red}}\arrow[r,"f_{\mathrm{red}}"]\arrow[d] &Y_{\mathrm{red}}\arrow[d]\\
X\arrow[r,"f"']&Y
\end{tikzcd}
\]
成立，竖箭头为嵌入态射。换言之，$X_{\mathrm{red}}\to X$ 为函子态射。
特别地，若 $X$ 约化，则每个态射 $f:X\to Y$ 都分解为
$X\xrightarrow{f_{\mathrm{red}}}Y_{\mathrm{red}}\to Y$；
即 $f$ 不大于嵌入态射 $Y_{\mathrm{red}}\to Y$。

\egastatement{5.1.6}{命题}
若态射 $f:X\to Y$ 为满射、根式态射、浸入、闭浸入、开浸入、局部浸入或局部同构，
则 $f_{\mathrm{red}}$ 具有相应性质。
反之，若 $f_{\mathrm{red}}$ 为满射（分别地，根式态射），则 $f$ 也如此。

对于满射，命题显然。对于根式态射，由 $X$ 和 $X_{\mathrm{red}}$ 在每个点 $x$
处的剩余域 $k(x)$ 相同及（\hyperlink{para:1:3.5.8}{3.5.8}）即可得到。
最后，若 $f=(\psi,\theta)$ 为浸入、闭浸入或局部浸入
（分别地，开浸入或局部同构），则命题由下述事实得到：
若 $\theta_x^\sharp$ 为满射（分别地，双射），
则对两端取幂零根的商后所得的同态也如此（\hyperlink{para:1:5.1.2}{5.1.2} 和 \hyperlink{para:1:4.2.2}{4.2.2}；另参见 \hyperlink{para:1:5.5.12}{5.5.12}）。

\egastatement{5.1.7}{命题}
若 $X,Y$ 为两个 $S$-预概形，则
$X_{\mathrm{red}}\times_{S_{\mathrm{red}}}Y_{\mathrm{red}}$
与 $X_{\mathrm{red}}\times_S Y_{\mathrm{red}}$ 相同，
并典范地认同于 $X\times_S Y$ 的一个具有相同底空间的子预概形。

$X_{\mathrm{red}}\times_S Y_{\mathrm{red}}$ 与 $X\times_S Y$ 的一个具有相同底空间的子预概形的典范认同，
由（\hyperlink{para:1:4.3.1}{4.3.1}）得到。
另一方面，结构态射 $X_{\mathrm{red}}\to S$、$Y_{\mathrm{red}}\to S$
均通过 $S_{\mathrm{red}}$ 分解（\hyperlink{para:1:5.1.5}{5.1.5}）；由于 $S_{\mathrm{red}}\to S$ 为单态射，
第一断言由（\hyperlink{para:1:3.2.4}{3.2.4}）得到。

\egastatement{5.1.8}{推论}
预概形 $(X\times_S Y)_{\mathrm{red}}$ 与
$(X_{\mathrm{red}}\times_{S_{\mathrm{red}}}Y_{\mathrm{red}})_{\mathrm{red}}$ 典范地认同。

这由（\hyperlink{para:1:5.1.2}{5.1.2}）和（\hyperlink{para:1:5.1.7}{5.1.7}）得到。
应当注意，即使 $X,Y$ 约化，$X\times_S Y$ 也未必约化，
因为两个约化代数的张量积可能具有幂零元。

\egastatement{5.1.9}{命题}
设 $X$ 为预概形，$\J$ 为 $\OO_X$ 的拟凝聚理想层，且对某个整数 $n>0$ 有 $\J^n=0$。
以 $X_0$ 表示闭子预概形 $(X,\OO_X/\J)$。
则 $X$ 为仿射概形的充要条件是 $X_0$ 为仿射概形。

条件显然必要，现证明充分性。置 $X_k=(X,\OO_X/\J^{k+1})$，
只须对 $k$ 归纳证明各 $X_k$ 仿射，因而归结到 $\J^2=0$ 的情形。
置
\[
 A=\Gamma(X,\OO_X),\qquad
 A_0=\Gamma(X_0,\OO_{X_0})=\Gamma(X,\OO_X/\J).
\]
典范同态 $\OO_X\to\OO_X/\J$ 诱导环同态 $\varphi:A\to A_0$。
下面将证明 $\varphi$ 为满射，即序列
\begin{equation}
0\longrightarrow\Gamma(X,\J)\longrightarrow\Gamma(X,\OO_X)
\longrightarrow\Gamma(X,\OO_X/\J)\longrightarrow0
\tag{5.1.9.1}\label{eq:1:5.1.9.1}
\end{equation}
正合。先假设这一点成立，并证明它蕴含命题。

注意 $\mathfrak K=\Gamma(X,\J)$ 是 $A$ 中平方为零的理想，
所以是 $A_0=A/\mathfrak K$-模。
由假设 $X_0=\Spec(A_0)$；$X_0$ 与 $X$ 的底空间相同，故
$\mathfrak K=\Gamma(X_0,\J)$。
又因 $\J^2=0$，$\J$ 是拟凝聚 $(\OO_X/\J)$-模，
因而 $\J\simeq\wt{\mathfrak K}$，且对每个 $x\in X_0$ 有 $\J_x=\mathfrak K_x$（\hyperlink{para:1:1.4.1}{1.4.1}）。

置 $X'=\Spec(A)$，考虑与恒等映射 $A\to\Gamma(X,\OO_X)$
相对应的预概形态射 $f=(\psi,\theta):X\to X'$（\hyperlink{para:1:2.2.4}{2.2.4}）。
对 $X$ 的每个仿射开集 $V$，图
\[
\begin{tikzcd}[column sep=large]
A\arrow[r]\arrow[d]&\Gamma(V,\OO_X|V)\arrow[d]\\
A_0=A/\mathfrak K\arrow[r]&\Gamma(V,\OO_{X_0}|V)
\end{tikzcd}
\]
交换，因此图
\[
\begin{tikzcd}
X'&X\arrow[l,"f"']\\
X'_0\arrow[u,"j'"]&X_0\arrow[l,"f_0"]\arrow[u,"j"']
\end{tikzcd}
\]
交换，其中 $X'_0$ 是由拟凝聚理想层 $\wt{\mathfrak K}$ 定义的 $X'$ 的闭子预概形，
$j,j'$ 为典范嵌入。
由于 $X_0$ 仿射，$f_0$ 是同构；又因 $j,j'$ 的底空间映射为恒等映射，
首先得到 $\psi:X\to X'$ 为同胚。
此外，$\J_x=\mathfrak K_x$ 表明
$\theta^\sharp:\psi^*(\OO_{X'})\to\OO_X$ 的限制给出
$\psi^*(\wt{\mathfrak K})\to\J$ 的同构。
另一方面，由 $f_0$ 为同构，取商后 $\theta^\sharp$ 诱导同构
$\psi^*(\OO_{X'}/\wt{\mathfrak K})\to\OO_X/\J$。
由五引理（M，I，1.1）立即得到 $\theta^\sharp$ 本身为同构，故 $f$ 为同构，$X$ 仿射。

所以只剩证明（\hyperref[eq:1:5.1.9.1]{5.1.9.1}）的正合性；它可由 $H^1(X,\J)=0$ 得到。
而 $H^1(X,\J)=H^1(X_0,\J)$，且已知 $\J$ 为拟凝聚 $\OO_{X_0}$-模。
因此所需结论由下述引理得到。

\egastatement{5.1.9.2}{引理}
若 $Y$ 为仿射概形，$\F$ 为拟凝聚 $\OO_Y$-模，则 $H^1(Y,\F)=0$。

此引理将在第 III 章 §1 中，作为更一般的定理
“对每个 $i>0$，$H^i(Y,\F)=0$”的推论予以证明。
为给出独立的证明，注意 $H^1(Y,\F)$ 认同于以 $\F$ 为核的 $\OO_Y$ 的扩张类所构成的模
$\operatorname{Ext}^1_{\OO_Y}(Y;\OO_Y,\F)$（T，4.2.3）。
所以只须证明每个这样的扩张 $\G$ 平凡。
对每个 $y\in Y$，存在其邻域 $V$ 使
$\G|V\simeq\F|V\oplus\OO_Y|V$（0，\hyperlink{para:0:5.4.9}{5.4.9}），故 $\G$ 为拟凝聚 $\OO_Y$-模。
若 $A$ 为 $Y$ 的环，则 $\F=\wt M$、$\G=\wt N$，其中 $M,N$ 为 $A$-模。
由假设，$N$ 是以 $M$ 为核的 $A$-模 $A$ 的扩张（\hyperlink{para:1:1.3.11}{1.3.11}）。
这样的扩张必然平凡，引理以及（\hyperlink{para:1:5.1.9}{5.1.9}）遂得证。

\egastatement{5.1.10}{推论}
设预概形 $X$ 的幂零根 $\N_X$ 为幂零理想层。
则 $X$ 为仿射概形的充要条件是 $X_{\mathrm{red}}$ 为仿射概形。


% END chapters/ch1_s5_1.tex

\par

% BEGIN chapters/ch1_s5_2.tex
\subsection{具有给定底空间的子预概形的存在性}\hypertarget{heading:1:5.2}{}
\src{131—132}

\egastatement{5.2.1}{命题}
对于预概形 $X$ 的底空间中的每个局部闭子空间 $Y$，
存在唯一以 $Y$ 为底空间的约化子预概形。

唯一性由（\hyperlink{para:1:5.1.2}{5.1.2}）得到，故只须证明存在性。
若 $X$ 仿射、环为 $A$，且 $Y$ 在 $X$ 中闭，则命题立即成立：
$\jj(Y)$ 是满足 $V(\mathfrak a)=Y$ 的最大理想 $\mathfrak a\subset A$，
且它等于自身的根（\hyperlink{para:1:1.1.4}{1.1.4}，(i)），所以 $A/\jj(Y)$ 是约化环。

一般情形下，对每个使 $U\cap Y$ 在 $U$ 中闭的仿射开集 $U\subset X$，
考虑由 $A(U)$ 的理想 $\jj(U\cap Y)$ 所相伴的理想层定义的闭子预概形 $Y_U$；它是约化的。
若 $V\subset U$ 是仿射开集，则 $Y_V$ 是 $Y_U$ 在 $V\cap Y$ 上诱导的预概形。
事实上，该诱导预概形是 $V$ 的约化闭子预概形，底空间为 $V\cap Y$，
所以 $Y_V$ 的唯一性给出此断言。

\egastatement{5.2.2}{命题}
设 $X$ 为约化预概形，$f:X\to Y$ 为态射，$Z$ 为 $Y$ 的闭子预概形，且 $f(X)\subset Z$。
则 $f$ 分解为 $X\to Z\xrightarrow{j}Y$，其中 $j$ 为嵌入态射。

由假设，$X$ 的闭子预概形 $f^{-1}(Z)$ 的底空间为整个 $X$（\hyperlink{para:1:4.4.1}{4.4.1}）。
由于 $X$ 约化，该闭子预概形等于 $X$（\hyperlink{para:1:5.1.2}{5.1.2}），故命题由（\hyperlink{para:1:4.4.1}{4.4.1}）得到。

\egastatement{5.2.3}{推论}
设 $X$ 为预概形 $Y$ 的约化子预概形，$Z$ 为以 $\overline X$ 为底空间的 $Y$ 的约化闭子预概形。
则 $X$ 是 $Z$ 在一个开集上诱导的子预概形。

事实上，存在 $Y$ 的开集 $U$ 使 $X=U\cap\overline X$。
由（\hyperlink{para:1:5.2.2}{5.2.2}），$X$ 是 $Z$ 的约化子预概形；
因而由唯一性（\hyperlink{para:1:5.2.1}{5.2.1}），它正是 $Z$ 在开子空间 $X$ 上诱导的子预概形。

\egastatement{5.2.4}{推论}
设 $f:X\to Y$ 为态射，$X'$、$Y'$ 分别为由拟凝聚理想层
$\J\subset\OO_X$、$\mathscr K\subset\OO_Y$ 定义的闭子预概形。
若 $X'$ 约化且 $f(X')\subset Y'$，则 $f^*(\mathscr K)\OO_X\subset\J$。

根据（\hyperlink{para:1:5.2.2}{5.2.2}），$f|X'$ 分解为 $X'\to Y'\to Y$，再应用（\hyperlink{para:1:4.4.6}{4.4.6}）即可。


% END chapters/ch1_s5_2.tex

\par

% BEGIN chapters/ch1_s5_3.tex
\subsection{对角线；态射的图}\hypertarget{heading:1:5.3}{}
\src{132—135}

\egapar{5.3.1}
设 $X$ 为 $S$-预概形。
称 $S$-态射 $(1_X,1_X)_S$ 为 $X$ 到 $X\times_S X$ 的\emph{对角态射}，
记为 $\Delta_{X|S}$、$\Delta_X$，或在无混淆时记为 $\Delta$。
它是满足
\begin{equation}
 p_1\circ\Delta_X=p_2\circ\Delta_X=1_X
 \tag{5.3.1.1}\label{eq:1:5.3.1.1}
\end{equation}
的唯一 $S$-态射，其中 $p_1,p_2$ 为 $X\times_S X$ 的投影（\hyperlink{para:1:3.2.1}{3.2.1}）。
若 $f:T\to X$、$g:T\to Y$ 为两个 $S$-态射，立即可验证
\begin{equation}
 (f,g)_S=(f\times_S g)\circ\Delta_{T|S}.
 \tag{5.3.1.2}\label{eq:1:5.3.1.2}
\end{equation}
读者应当注意，上述定义及（\hyperlink{para:1:5.3.1}{5.3.1}）至（\hyperlink{para:1:5.3.8}{5.3.8}）的结果在任意范畴中都成立，
只要其中出现的积在该范畴中存在。

\egastatement{5.3.2}{命题}
设 $X,Y$ 为两个 $S$-预概形。
若将 $(X\times Y)\times(X\times Y)$ 典范地认同于
$(X\times X)\times(Y\times Y)$，则 $\Delta_{X\times Y}$ 认同于 $\Delta_X\times\Delta_Y$。

事实上，若 $p_1,q_1$ 分别为 $X\times X\to X$、$Y\times Y\to Y$ 的第一投影，
则 $(X\times Y)\times(X\times Y)\to X\times Y$ 的第一投影认同于 $p_1\times q_1$，且
\[
(p_1\times q_1)\circ(\Delta_X\times\Delta_Y)
=(p_1\circ\Delta_X)\times(q_1\circ\Delta_Y)=1_{X\times Y}.
\]
第二投影的论证相同。

\egastatement{5.3.4}{推论}
对任意基预概形扩张 $S'\to S$，$\Delta_{X_{(S')}}$ 典范地认同于 $(\Delta_X)_{(S')}$。
\transnote{原版在（\hyperlink{para:1:5.3.2}{5.3.2}）之后直接编号为（\hyperlink{para:1:5.3.4}{5.3.4}）；此处保留原版编号，未补造（5.3.3）。}

只须注意 $(X\times_S X)_{(S')}$ 典范地认同于
$X_{(S')}\times_{S'}X_{(S')}$（\hyperlink{para:1:3.3.10}{3.3.10}）。

\egastatement{5.3.5}{命题}
设 $X,Y$ 为两个 $S$-预概形，$\varphi:S\to T$ 为态射，因而使每个 $S$-预概形成为 $T$-预概形。
设 $f:X\to S$、$g:Y\to S$ 为结构态射，$p,q$ 为 $X\times_S Y$ 的投影，
$\pi=f\circ p=g\circ q$ 为其到 $S$ 的结构态射。
则图
\begin{equation}
\begin{tikzcd}[column sep=large]
X\times_S Y\arrow[r,"{(p,q)_T}"]\arrow[d,"\pi"']&X\times_T Y\arrow[d,"f\times_T g"]\\
S\arrow[r,"\Delta_{S|T}"']&S\times_T S
\end{tikzcd}
\tag{5.3.5.1}\label{eq:1:5.3.5.1}
\end{equation}
交换，并把 $X\times_S Y$ 认同于两个 $(S\times_T S)$-预概形 $S$ 与 $X\times_T Y$ 的积，
其投影分别为 $\pi$ 与 $(p,q)_T$。

\begin{sloppypar}
根据（\hyperlink{para:1:3.4.3}{3.4.3}），取任意 $T$-预概形 $Z$，将 $X,Y,S$ 分别替换为
$X(Z)_T,Y(Z)_T,S(Z)_T$，即归结到证明集合范畴中对应的命题。
在集合范畴中验证立即可作，留给读者。
\end{sloppypar}


\egastatement{5.3.6}{推论}
令 $P=S\times_T S$，则态射 $(p,q)_T$ 认同于
$1_{X\times_T Y}\times_P\Delta_S$。

这由（\hyperlink{para:1:5.3.5}{5.3.5}）和（\hyperlink{para:1:3.3.4}{3.3.4}）得到。

\egastatement{5.3.7}{推论}
若 $f:X\to Y$ 为 $S$-态射，则图
\[
\begin{tikzcd}[column sep=large]
X\arrow[r,"{(1_X,f)_S}"]\arrow[d,"f"']&X\times_S Y\arrow[d,"f\times_S1_Y"]\\
Y\arrow[r,"\Delta_Y"']&Y\times_S Y
\end{tikzcd}
\]
交换，并把 $X$ 认同于两个 $(Y\times_S Y)$-预概形 $Y$ 与 $X\times_S Y$ 的积。

只须在（\hyperlink{para:1:5.3.5}{5.3.5}）中以 $Y$ 代替 $S$、以 $S$ 代替 $T$，
并注意 $X\times_Y Y=X$（\hyperlink{para:1:3.3.3}{3.3.3}）。

\egastatement{5.3.8}{命题}
态射 $f:X\to Y$ 为预概形单态射的充要条件是
$\Delta_{X|Y}$ 为 $X$ 到 $X\times_Y X$ 的同构。

$f$ 为单态射意味着：对每个 $Y$-预概形 $Z$，对应映射
$f':X(Z)_Y\to Y(Z)_Y$ 为单射。
由于 $Y(Z)_Y$ 只有一个元素，这意味着 $X(Z)_Y$ 至多有一个元素。
\transnote{原文此处将 $X(Z)_Y$ 也说成“只有一个元素”；应允许空集，
因此这里明确写作“至多有一个元素”。}
这又等价于说 $X(Z)_Y\times X(Z)_Y$ 经对角映射典范地同构于 $X(Z)_Y$。
前一集合等于 $(X\times_Y X)(Z)_Y$（\hyperref[eq:1:3.4.3.1]{3.4.3.1}），故此条件正表示 $\Delta_{X|Y}$ 为同构。

\NewNote{\textbf{作者勘误提示：}作者确认此处原证明没有验证对角像局部闭；书末保留原作者的补证。结论与证明缺步应分开判断。 参见\hyperlink{errIII:10}{《勘误与增补（清单二）》第10条}。}
\egastatement{5.3.9}{命题}
对角态射 $\Delta_X$ 是 $X$ 到 $X\times_S X$ 的浸入。

底空间上的连续映射 $p_1$、$\Delta_X$ 满足 $p_1\circ\Delta_X=1_X$，
所以 $\Delta_X$ 是 $X$ 到 $\Delta_X(X)$ 的同胚。
同样，对应于 $p_1$ 与 $\Delta_X$ 的同态的复合
$\OO_x\to\OO_{\Delta_X(x)}\to\OO_x$ 为恒等同态，
所以对应于 $\Delta_X$ 的同态为满射；命题遂由（\hyperlink{para:1:4.2.2}{4.2.2}）得到。
\transnote{原文略去（\hyperlink{para:1:4.2.2}{4.2.2}）所需的像的局部闭性；这一步可用本卷已有结果补足。
令 $f:X\to S$ 为结构态射。对每个 $x\in X$，取 $f(x)$ 的仿射开邻域 $V$，
以及包含于 $f^{-1}(V)$ 的 $x$ 的仿射开邻域 $U$。
在开子预概形 $U\times_S U=U\times_V U$ 内，$\Delta_X(U)$ 是闭子空间，
因为乘法同态 $A(U)\otimes_{A(V)}A(U)\to A(U)$ 为满射（\hyperlink{para:1:3.2.2}{3.2.2} 及 \hyperlink{para:1:4.2.3}{4.2.3}）。
各 $U\times_S U$ 的并是 $\Delta_X(X)$ 的开邻域；对每个这样的开集，
它与 $\Delta_X(X)$ 的交恰为 $\Delta_X(U)$，所以 $\Delta_X(X)$ 在该并内闭，因而局部闭。
本段是译者补证，不属于原著正文。}

称与浸入 $\Delta_X$ 相伴的 $X\times_S X$ 的子预概形（\hyperlink{para:1:4.2.1}{4.2.1}）为 $X\times_S X$ 的\emph{对角线}。

\egastatement{5.3.10}{推论}
在（\hyperlink{para:1:5.3.5}{5.3.5}）的假设下，$(p,q)_T$ 为浸入。

这由（\hyperlink{para:1:5.3.6}{5.3.6}）和（\hyperlink{para:1:4.3.1}{4.3.1}）得到。
称它为 $X\times_S Y$ 到 $X\times_T Y$ 的\emph{典范浸入}。

\egastatement{5.3.11}{推论}
设 $X,Y$ 为两个 $S$-预概形，$f:X\to Y$ 为 $S$-态射。
则 $f$ 的图态射 $\Gamma_f=(1_X,f)_S$（\hyperlink{para:1:3.3.14}{3.3.14}）为 $X$ 到 $X\times_S Y$ 的浸入。

这是（\hyperlink{para:1:5.3.10}{5.3.10}）中以 $Y$ 代替 $S$、以 $S$ 代替 $T$ 的特例（参见 \hyperlink{para:1:5.3.7}{5.3.7}）。
与浸入 $\Gamma_f$ 相伴的子预概形（\hyperlink{para:1:4.2.1}{4.2.1}）称为态射 $f$ 的\emph{图}。
$X\times_S Y$ 中作为态射 $X\to Y$ 的图的子预概形 $G$，由下述性质刻画：
投影 $p_1:X\times_S Y\to X$ 在 $G$ 上的限制是同构 $g:G\to X$。
此时 $G$ 为态射 $p_2\circ g^{-1}$ 的图，其中 $p_2$ 为到 $Y$ 的投影。

特别取 $X=S$，则 $S$-态射 $S\to Y$，即 $Y$ 的 $S$-截面（\hyperlink{para:1:2.5.5}{2.5.5}），等于它们的图态射。
作为 $S$-截面的图的 $Y$ 的子预概形
（即结构态射 $Y\to S$ 的限制将其同构地映到 $S$ 的子预概形），
也称为这些截面的\emph{像}，或略称为 $Y$ 的 $S$-截面。

\egastatement{5.3.12}{推论}
沿用（\hyperlink{para:1:5.3.11}{5.3.11}）的假设和记号。
对态射 $g:S'\to S$，设 $f'$ 为 $f$ 在 $g$ 下的逆像（\hyperlink{para:1:3.3.7}{3.3.7}），
则 $\Gamma_{f'}$ 为 $\Gamma_f$ 在 $g$ 下的逆像。

这是（\hyperref[eq:1:3.3.10.1]{3.3.10.1}）的特例。

\egastatement{5.3.13}{推论}
设 $f:X\to Y$、$g:Y\to Z$ 为两个态射。
若 $g\circ f$ 是浸入（分别地，局部浸入），则 $f$ 也如此。

事实上，$f$ 分解为
$X\xrightarrow{\Gamma_f}X\times_Z Y\xrightarrow{p_2}Y$，
且 $p_2=(g\circ f)\times_Z1_Y$（\hyperlink{para:1:3.3.4}{3.3.4}）。
若 $g\circ f$ 为浸入（分别地，局部浸入），则 $p_2$ 也如此（\hyperlink{para:1:4.3.1}{4.3.1} 和 \hyperlink{para:1:4.5.5}{4.5.5}）；
又因 $\Gamma_f$ 为浸入（\hyperlink{para:1:5.3.11}{5.3.11}），由浸入的传递性
（\hyperlink{para:1:4.2.4}{4.2.4}），分别由（\hyperlink{para:1:4.5.5}{4.5.5}），得到结论。\transnote{原文引用（\hyperlink{para:1:4.2.4}{4.2.4}）；浸入的传递性直接见（\hyperlink{para:1:4.2.5}{4.2.5}）。}

\egastatement{5.3.14}{推论}
设 $j:X\to Y$、$g:X\to Z$ 为两个 $S$-态射。
若 $j$ 是浸入（分别地，局部浸入），则 $(j,g)_S$ 也如此。

若 $p:Y\times_S Z\to Y$ 为第一投影，则 $j=p\circ(j,g)_S$，应用（\hyperlink{para:1:5.3.13}{5.3.13}）即可。

\egastatement{5.3.15}{命题}
若 $f:X\to Y$ 为 $S$-态射，则图
\begin{equation}
\begin{tikzcd}
X\arrow[r,"\Delta_X"]\arrow[d,"f"']&X\times_S X\arrow[d,"f\times_S f"]\\
Y\arrow[r,"\Delta_Y"']&Y\times_S Y
\end{tikzcd}
\tag{5.3.15.1}\label{eq:1:5.3.15.1}
\end{equation}
交换。换言之，$\Delta_X$ 是预概形范畴中的函子态射。

验证立即可作，留给读者。

\egastatement{5.3.16}{推论}
若 $X$ 为 $Y$ 的子预概形，则对角线 $\Delta_X(X)$ 认同于
$\Delta_Y(Y)$ 的一个子预概形，其底空间认同于
\[
 \Delta_Y(Y)\cap p_1^{-1}(X)=\Delta_Y(Y)\cap p_2^{-1}(X),
\]
其中 $p_1,p_2$ 是 $Y\times_S Y$ 的投影。

对嵌入态射 $f:X\to Y$ 应用（\hyperlink{para:1:5.3.15}{5.3.15}）。
已知 $f\times_S f$ 为浸入，并把 $X\times_S X$ 的底空间认同于
$p_1^{-1}(X)\cap p_2^{-1}(X)$（\hyperlink{para:1:4.3.1}{4.3.1}）。
若 $z\in\Delta_Y(Y)\cap p_1^{-1}(X)$，则 $z=\Delta_Y(y)$，
且 $y=p_1(z)\in X$，所以可写 $y=f(y)$。
由图（\hyperref[eq:1:5.3.15.1]{5.3.15.1}）的交换性，$z=\Delta_Y(f(y))$ 属于 $\Delta_X(X)$。

\egastatement{5.3.17}{推论}
设 $f_1,f_2:Y\to X$ 为两个 $S$-态射，$y\in Y$ 满足
$f_1(y)=f_2(y)=x$，且对应的同态 $k(x)\to k(y)$ 相同。
若 $f=(f_1,f_2)_S$，则 $f(y)$ 属于对角线 $\Delta_{X|S}(X)$。

对应于 $f_i$ 的同态 $k(x)\to k(y)$ 定义 $S$-态射
$g_i:\Spec(k(y))\to\Spec(k(x))$，使图
\[
\begin{tikzcd}
\Spec(k(y))\arrow[r,"g_i"]\arrow[d]&\Spec(k(x))\arrow[d]\\
Y\arrow[r,"f_i"']&X
\end{tikzcd}
\]
交换。于是图
\[
\begin{tikzcd}[column sep=large]
\Spec(k(y))\arrow[r,"{(g_1,g_2)_S}"]\arrow[d]&
\Spec(k(x))\times_S\Spec(k(x))\arrow[d]\\
Y\arrow[r,"{(f_1,f_2)_S}"']&X\times_S X
\end{tikzcd}
\]
也交换。由 $g_1=g_2$，$\Spec(k(y))$ 的唯一点在 $(g_1,g_2)_S$ 下的像属于
$\Spec(k(x))\times_S\Spec(k(x))$ 的对角线。
因此结论由（\hyperlink{para:1:5.3.15}{5.3.15}）得到。


% END chapters/ch1_s5_3.tex

\par

% BEGIN chapters/ch1_s5_4.tex
\subsection{分离态射与分离预概形}\hypertarget{heading:1:5.4}{}
\src{135—136}

\egastatement{5.4.1}{定义}
若预概形态射 $f:X\to Y$ 的对角态射 $X\to X\times_Y X$ 是闭浸入，
则称 $f$ 是\emph{分离的}；此时也称 $X$ 是 $Y$ 上的\emph{分离预概形}，或称 $X$ 为 $Y$-\emph{概形}。
若预概形 $X$ 在 $\Spec(\mathbf Z)$ 上分离，则称 $X$ \emph{分离}；
此时也称 $X$ 为\emph{概形}（参见 \hyperlink{para:1:5.5.7}{5.5.7}）。

根据（\hyperlink{para:1:5.3.9}{5.3.9}），$X$ 在 $Y$ 上分离的充要条件是
$\Delta_{X|Y}(X)$ 在 $X\times_Y X$ 的底空间中闭。

\egastatement{5.4.2}{命题}
设 $S\to T$ 为分离态射。若 $X,Y$ 为两个 $S$-预概形，
则典范浸入 $X\times_S Y\to X\times_T Y$（\hyperlink{para:1:5.3.10}{5.3.10}）是闭的。

参照图（\hyperref[eq:1:5.3.5.1]{5.3.5.1}），$(p,q)_T$ 可看作由 $\Delta_{S|T}$ 经基预概形 $S\times_T S$ 的扩张
$f\times_T g:X\times_T Y\to S\times_T S$ 所得的态射。
所以命题由（\hyperlink{para:1:4.3.2}{4.3.2}）得到。

\egastatement{5.4.3}{推论}
设 $Y$ 为 $S$-概形，$f:X\to Y$ 为 $S$-态射。
则图态射 $\Gamma_f:X\to X\times_S Y$（\hyperlink{para:1:5.3.11}{5.3.11}）为闭浸入。

这是（\hyperlink{para:1:5.4.2}{5.4.2}）中以 $Y$ 代替 $S$、以 $S$ 代替 $T$ 的特例。

\egastatement{5.4.4}{推论}
设 $f:X\to Y$、$g:Y\to Z$ 为两个态射，$g$ 分离。
若 $g\circ f$ 为闭浸入，则 $f$ 也为闭浸入。

从（\hyperlink{para:1:5.4.3}{5.4.3}）出发的证明，与从（\hyperlink{para:1:5.3.11}{5.3.11}）出发证明（\hyperlink{para:1:5.3.13}{5.3.13}）相同。

\egastatement{5.4.5}{推论}
设 $Z$ 为 $S$-概形，$j:X\to Y$、$g:X\to Z$ 为两个 $S$-态射。
若 $j$ 为闭浸入，则 $(j,g)_S:X\to Y\times_S Z$ 也为闭浸入。

从（\hyperlink{para:1:5.4.4}{5.4.4}）出发的证明，与从（\hyperlink{para:1:5.3.13}{5.3.13}）出发证明（\hyperlink{para:1:5.3.14}{5.3.14}）相同。

\egastatement{5.4.6}{推论}
若 $X$ 为 $S$-概形，则 $X$ 的每个 $S$-截面（\hyperlink{para:1:2.5.5}{2.5.5}）都是闭浸入。

若 $\varphi:X\to S$ 为结构态射，$\psi:S\to X$ 为 $S$-截面，
只须对 $\varphi\circ\psi=1_S$ 应用（\hyperlink{para:1:5.4.5}{5.4.5}）即可。

\egastatement{5.4.7}{推论}
设 $S$ 为整预概形，$s$ 为其泛点，$X$ 为 $S$-概形。
若 $X$ 的两个 $S$-截面 $f,g$ 满足 $f(s)=g(s)$，则 $f=g$。

事实上，若 $x=f(s)=g(s)$，则对应于 $f,g$ 的同态 $k(x)\to k(s)$ 必然相同。
置 $h=(f,g)_S$，由（\hyperlink{para:1:5.3.17}{5.3.17}）可知 $h(s)$ 属于对角线 $Z=\Delta_X(X)$。
因 $S=\overline{\{s\}}$ 且 $Z$ 由假设为闭，$h(S)\subset Z$。
根据（\hyperlink{para:1:5.2.2}{5.2.2}），$h$ 分解为 $S\to Z\to X\times_S X$，
再由对角线的定义得到 $f=g$。

\egastatement{5.4.8}{注}
反过来，若（\hyperlink{para:1:5.4.3}{5.4.3}）的结论在 $f=1_Y$ 时成立，便得到 $Y$ 在 $S$ 上分离。
同样，若（\hyperlink{para:1:5.4.5}{5.4.5}）的结论可用于两个态射
$Y\xrightarrow{\Delta_Y}Y\times_Z Y\xrightarrow{p_1}Y$，
则得到 $\Delta_Y$ 为闭浸入，因而 $Y$ 在 $Z$ 上分离。
\transnote{原文在后一处引用（\hyperlink{para:1:5.4.5}{5.4.5}）；这里所用的复合态射判据直接见（\hyperlink{para:1:5.4.4}{5.4.4}）。}
最后，若（\hyperlink{para:1:5.4.6}{5.4.6}）的结论对 $Y$-预概形 $Y\times_S Y\to Y$ 的 $Y$-截面 $\Delta_Y$ 成立，
也蕴含 $Y$ 在 $S$ 上分离。


% END chapters/ch1_s5_4.tex

\par

% BEGIN chapters/ch1_s5_5.tex
\subsection{分离性判据}\hypertarget{heading:1:5.5}{}
\src{136—140}

\egastatement{5.5.1}{命题}
\begin{enumerate}[label=(\roman*)]
\item 每个预概形单态射，特别是每个浸入，都是分离态射。
\item 两个分离态射的复合仍分离。
\item 若 $f:X\to X'$、$g:Y\to Y'$ 为两个分离的 $S$-态射，则 $f\times_S g$ 分离。
\item 若 $f:X\to Y$ 为分离的 $S$-态射，则对任意基预概形扩张 $S'\to S$，
$S'$-态射 $f_{(S')}$ 分离。
\item 若两个态射的复合 $g\circ f$ 分离，则 $f$ 分离。
\item $f$ 分离的充要条件是 $f_{\mathrm{red}}$（\hyperlink{para:1:5.1.5}{5.1.5}）分离。
\end{enumerate}

(i) 由（\hyperlink{para:1:5.3.8}{5.3.8}）立即得到。若 $f:X\to Y$、$g:Y\to Z$ 为态射，图
\begin{equation}
\begin{tikzcd}[column sep=large]
X\arrow[r,"\Delta_{X|Z}"]\arrow[dr,"\Delta_{X|Y}"']&X\times_Z X\\
&X\times_Y X\arrow[u,"j"']
\end{tikzcd}
\tag{5.5.1.1}\label{eq:1:5.5.1.1}
\end{equation}
交换，验证立即可作，其中 $j$ 为典范浸入（\hyperlink{para:1:5.3.10}{5.3.10}）。
若 $f,g$ 分离，则 $\Delta_{X|Y}$ 由定义为闭浸入，$j$ 由（\hyperlink{para:1:5.4.2}{5.4.2}）为闭浸入，
所以 $\Delta_{X|Z}$ 也为闭浸入，这证明 (ii)。
\transnote{原文在此引用（\hyperlink{para:1:4.2.4}{4.2.4}）；闭浸入复合的陈述见（\hyperlink{para:1:4.2.5}{4.2.5}）。}
由 (i)、(ii)，(iii)、(iv) 等价（\hyperlink{para:1:3.5.1}{3.5.1}），故只须证明 (iv)。
由（\hyperlink{para:1:3.3.11}{3.3.11}）和（\hyperref[eq:1:3.3.9.1]{3.3.9.1}），有典范认同
\[
 X_{(S')}\times_{Y_{(S')}}X_{(S')}
 \simeq (X\times_Y X)\times_Y Y_{(S')}.
\]
立即可验证，此时对角态射认同于 $\Delta_X\times_Y1_{Y_{(S')}}$，
故命题由（\hyperlink{para:1:4.3.1}{4.3.1}）得到。

为证明 (v)，如（\hyperlink{para:1:5.3.13}{5.3.13}）中那样考虑分解
$X\xrightarrow{\Gamma_f}X\times_Z Y\xrightarrow{p_2}Y$。
由 $p_2=(g\circ f)\times_Z1_Y$，假设及 (iii)、(i) 蕴含 $p_2$ 分离。
$\Gamma_f$ 为浸入，故由 (i) 分离；再由 (ii)，$f$ 分离。
最后，为证明 (vi)，回忆
$X_{\mathrm{red}}\times_{Y_{\mathrm{red}}}X_{\mathrm{red}}$
与 $X_{\mathrm{red}}\times_Y X_{\mathrm{red}}$ 典范地认同（\hyperlink{para:1:5.1.7}{5.1.7}）。
若 $j:X_{\mathrm{red}}\to X$ 为嵌入，则图
\[
\begin{tikzcd}[column sep=large]
X_{\mathrm{red}}\arrow[r,"\Delta_{X_{\mathrm{red}}}"]\arrow[d,"j"']&
X_{\mathrm{red}}\times_Y X_{\mathrm{red}}\arrow[d,"j\times_Y j"]\\
X\arrow[r,"\Delta_{X|Y}"']&X\times_Y X
\end{tikzcd}
\]
交换（\hyperlink{para:1:5.3.15}{5.3.15}）。竖箭头均为底空间的同胚（\hyperlink{para:1:4.3.1}{4.3.1}），故命题成立。

\egastatement{5.5.2}{推论}
若 $f:X\to Y$ 分离，则它在 $X$ 的任意子预概形上的限制分离。

这由（\hyperlink{para:1:5.5.1}{5.5.1}，(i)、(ii)）得到。

\egastatement{5.5.3}{推论}
若 $X,Y$ 为两个 $S$-预概形，且 $Y$ 在 $S$ 上分离，则 $X\times_S Y$ 在 $X$ 上分离。

这是（\hyperlink{para:1:5.5.1}{5.5.1}，(iv)）的特例。

\egastatement{5.5.4}{命题}
设预概形 $X$ 的底空间为有限个闭子集 $X_k$（$1\leq k\leq n$）的并。
仍以 $X_k$ 表示具有该底空间的约化子预概形（\hyperlink{para:1:5.2.1}{5.2.1}）。
设 $f:X\to Y$ 为态射，并为每个 $k$ 取闭子集 $Y_k\subset Y$，使 $f(X_k)\subset Y_k$；
同样以 $Y_k$ 表示相应的约化子预概形。
于是 $f|X_k$ 分解为 $X_k\xrightarrow{f_k}Y_k\to Y$（\hyperlink{para:1:5.2.2}{5.2.2}）。
则 $f$ 分离的充要条件是每个 $f_k$ 分离。

必要性由（\hyperlink{para:1:5.5.1}{5.5.1}，(i)、(ii)、(v)）得到。
反之，若每个 $f_k$ 分离，则 $f$ 的每个限制 $X_k\to Y$ 分离（\hyperlink{para:1:5.5.1}{5.5.1}，(i)、(ii)）。
若 $p_1,p_2$ 为 $X\times_Y X$ 的投影，则子空间 $\Delta_{X_k}(X_k)$
认同于 $\Delta_X(X)\cap p_1^{-1}(X_k)$（\hyperlink{para:1:5.3.16}{5.3.16}）。
这些子空间在 $X\times_Y X$ 中闭，所以它们的并 $\Delta_X(X)$ 也闭。

特别地，若各 $X_k$ 为 $X$ 的不可约分支，则可取 $Y_k$ 为 $Y$ 的不可约分支（0，\hyperlink{para:0:2.1.5}{2.1.5}）。
因而在此情形，（\hyperlink{para:1:5.5.4}{5.5.4}）将分离性的研究归结到整预概形的情形。

\egastatement{5.5.5}{命题}
设 $(Y_\lambda)$ 为预概形 $Y$ 的开覆盖。
态射 $f:X\to Y$ 分离的充要条件是每个限制
$f^{-1}(Y_\lambda)\to Y_\lambda$ 分离。

置 $X_\lambda=f^{-1}(Y_\lambda)$。
由（\hyperlink{para:1:4.2.4}{4.2.4}，b)）及积 $X_\lambda\times_Y X_\lambda$
与 $X_\lambda\times_{Y_\lambda}X_\lambda$ 相同（\hyperlink{para:1:3.2.5}{3.2.5}），
只须证明 $X_\lambda\times_Y X_\lambda$ 覆盖 $X\times_Y X$。
置 $Y_{\lambda\mu}=Y_\lambda\cap Y_\mu$、
$X_{\lambda\mu}=X_\lambda\cap X_\mu=f^{-1}(Y_{\lambda\mu})$。
则 $X_\lambda\times_Y X_\mu$ 认同于
$X_{\lambda\mu}\times_Y X_{\lambda\mu}$（\hyperlink{para:1:3.2.6.4}{3.2.6.4}），
也即 $X_{\lambda\mu}\times_{Y_\lambda}X_{\lambda\mu}$（\hyperlink{para:1:3.2.5}{3.2.5}），
从而认同于 $X_\lambda\times_Y X_\lambda$ 的一个开集（\hyperlink{para:1:3.2.7}{3.2.7}），所断言的覆盖成立。

因此，取 $Y$ 的仿射开覆盖，可将分离态射的研究归结到以仿射概形为靶的分离态射。
\transnote{原文在这一句引用（\hyperlink{para:1:5.5.4}{5.5.4}）；对应的开覆盖判据是刚证明的（\hyperlink{para:1:5.5.5}{5.5.5}）。}

\egastatement{5.5.6}{命题}
设 $Y$ 为仿射概形，$X$ 为预概形，$(U_\alpha)$ 为 $X$ 的仿射开覆盖。
态射 $f:X\to Y$ 分离的充要条件是：对每对指标 $(\alpha,\beta)$，
$U_\alpha\cap U_\beta$ 为仿射开集，且环
$\Gamma(U_\alpha\cap U_\beta,\OO_X)$
由两个环 $\Gamma(U_\alpha,\OO_X)$、$\Gamma(U_\beta,\OO_X)$ 的典范像之并生成。

各 $U_\alpha\times_Y U_\beta$ 构成 $X\times_Y X$ 的开覆盖（\hyperlink{para:1:3.2.7}{3.2.7}）。
以 $p,q$ 表示投影，有
\[
\begin{aligned}
\Delta_X^{-1}(U_\alpha\times_Y U_\beta)
 &=\Delta_X^{-1}\bigl(p^{-1}(U_\alpha)\cap q^{-1}(U_\beta)\bigr)\\
 &=\Delta_X^{-1}(p^{-1}(U_\alpha))\cap
   \Delta_X^{-1}(q^{-1}(U_\beta))=U_\alpha\cap U_\beta.
\end{aligned}
\]
故只须表述 $\Delta_X$ 在 $U_\alpha\cap U_\beta$ 上的限制为到
$U_\alpha\times_Y U_\beta$ 的闭浸入的条件。
由定义，该限制为 $(j_\alpha,j_\beta)_Y$，
其中 $j_\alpha,j_\beta$ 分别是交集到 $U_\alpha,U_\beta$ 的嵌入。
而 $U_\alpha\times_Y U_\beta$ 仿射，其环典范地同构于
\[
 \Gamma(U_\alpha,\OO_X)\otimes_{\Gamma(Y,\OO_Y)}\Gamma(U_\beta,\OO_X)
\]
（\hyperlink{para:1:3.2.2}{3.2.2}）。所以 $U_\alpha\cap U_\beta$ 必须仿射，且将
$h_\alpha\otimes h_\beta$ 映为 $h_\alpha h_\beta$ 的环同态必须为满射（\hyperlink{para:1:4.2.3}{4.2.3}）。证明完成。

\egastatement{5.5.7}{推论}
仿射概形是分离的，因而确实是概形，这说明（\hyperlink{para:1:5.4.1}{5.4.1}）的术语与此前一致。

\egastatement{5.5.8}{推论}
设 $Y$ 为仿射概形。态射 $f:X\to Y$ 分离的充要条件是 $X$ 分离，即 $X$ 为概形。

事实上，（\hyperlink{para:1:5.5.6}{5.5.6}）的判据不依赖于 $f$。

\egastatement{5.5.9}{推论}
态射 $f:X\to Y$ 分离的必要条件是：
对每个使 $Y$ 在其上诱导的预概形分离的开集 $U$，预概形 $f^{-1}(U)$ 分离；
充分条件是：对每个仿射开集 $U\subset Y$，$f^{-1}(U)$ 分离。

必要性由开覆盖判据及（\hyperlink{para:1:5.5.1}{5.5.1}，(ii)）得到；
充分性由开覆盖判据及（\hyperlink{para:1:5.5.8}{5.5.8}）得到，并利用 $Y$ 的仿射开覆盖的存在性。
\transnote{原文在这两处将开覆盖判据引为（\hyperlink{para:1:5.5.4}{5.5.4}）；这里的对应条目是（\hyperlink{para:1:5.5.5}{5.5.5}）。}
特别地，若 $X,Y$ 都是仿射概形，则每个态射 $X\to Y$ 都分离。

\egastatement{5.5.10}{命题}
设 $Y$ 为概形，$f:X\to Y$ 为态射。
对 $X$ 的每个仿射开集 $U$ 及 $Y$ 的每个仿射开集 $V$，$U\cap f^{-1}(V)$ 仿射。

以 $p_1,p_2$ 表示 $X\times_{\mathbf Z}Y$ 的投影。
子空间 $U\cap f^{-1}(V)$ 是
$\Gamma_f(X)\cap p_1^{-1}(U)\cap p_2^{-1}(V)$ 在 $p_1$ 下的像。
而 $p_1^{-1}(U)\cap p_2^{-1}(V)$ 认同于 $U\times_{\mathbf Z}V$ 的底空间（\hyperlink{para:1:3.2.7}{3.2.7}），
故为仿射概形（\hyperlink{para:1:3.2.2}{3.2.2}）。
由于 $\Gamma_f(X)$ 在 $X\times_{\mathbf Z}Y$ 中闭（\hyperlink{para:1:5.4.3}{5.4.3}），
上述交集在 $U\times_{\mathbf Z}V$ 中闭。
于是与 $\Gamma_f$ 相伴的子预概形在该开子空间上诱导的预概形，
是一个仿射概形的闭子预概形，故仿射（\hyperlink{para:1:4.2.3}{4.2.3}）。
最后，由 $\Gamma_f$ 是浸入得到命题。

\egastatement{5.5.11}{例}
例（\hyperlink{para:1:2.3.2}{2.3.2}）的预概形，即“域 $K$ 上的射影直线”，是分离的。
事实上，对于仿射开覆盖 $(X_1,X_2)$，$X_1\cap X_2=U_{12}$ 仿射；
$\Gamma(U_{12},\OO_X)$ 为形如 $f(s)/s^m$、$f\in K[s]$ 的有理分式所成的环，
它由 $K[s]$ 和 $1/s$ 生成，所以（\hyperlink{para:1:5.5.6}{5.5.6}）的条件成立。

仍如例（\hyperlink{para:1:2.3.2}{2.3.2}）选取 $X_1,X_2,U_{12},U_{21}$，但这次取将 $f(s)$ 映为 $f(t)$ 的同构为 $u_{12}$。
粘合后得到一个整的、非分离的预概形 $X$：
（\hyperlink{para:1:5.5.6}{5.5.6}）的第一条件成立，而第二条件不成立。
这里立即可见 $\Gamma(X,\OO_X)\to\Gamma(X_1,\OO_X)=K[s]$ 为同构。
逆同构定义满射态射 $f:X\to\Spec(K[s])$。
若 $y\in\Spec(K[s])$ 满足 $\jj_y\ne(s)$，则 $f^{-1}(y)$ 只有一个点；
若 $\jj_y=(s)$，则 $f^{-1}(y)$ 由两个不同的点组成。
称 $X$ 为“域 $K$ 上将点 $0$ 加倍的仿射直线”。
\transnote{原版在上述两处写作理想 $(0)$，但在本例中加倍的是原点，其理想为 $(s)$；
零理想对应的泛点并未加倍。此处已明确改写为 $(s)$。}

还可给出（\hyperlink{para:1:5.5.6}{5.5.6}）的两个条件都不成立的例子。
首先注意，在域 $K$ 上二元多项式环 $A=K[s,t]$ 的素谱 $Y$ 中，
开集 $U=D(s)\cup D(t)$ 不是仿射开集。
事实上，若 $z$ 为 $\OO_Y$ 在 $U$ 上的截面，则存在整数 $m,n\geq0$，
使 $s^m z$、$t^n z$ 均为关于 $s,t$ 的多项式在 $U$ 上的限制（\hyperlink{para:1:1.4.1}{1.4.1}）。
这显然只有在 $z$ 能延拓到整个 $Y$ 上、并与一个关于 $s,t$ 的多项式相认同时才可能。
若 $U$ 仿射，则嵌入态射 $U\to Y$ 将是同构（\hyperlink{para:1:1.7.3}{1.7.3}），与 $U\ne Y$ 矛盾。

现取仿射概形 $Y_1,Y_2$，分别为环
$A_1=K[s_1,t_1]$、$A_2=K[s_2,t_2]$ 的素谱。
置 $U_{12}=D(s_1)\cup D(t_1)$、$U_{21}=D(s_2)\cup D(t_2)$，
并取 $u_{12}$ 为同构 $Y_2\to Y_1$ 在 $U_{21}$ 上的限制，
该同构对应于将 $f(s_1,t_1)$ 映为 $f(s_2,t_2)$ 的环同构。
于是得到一个（\hyperlink{para:1:5.5.6}{5.5.6}）的两个条件都不满足的例子。
如此得到的整预概形称为“域 $K$ 上将点 $0$ 加倍的仿射平面”。

\egastatement{5.5.12}{注}
给定预概形态射的一个性质 $P$，考虑下述命题：
\begin{enumerate}[label=(\roman*)]
\item 每个闭浸入具有性质 $P$。
\item 两个具有性质 $P$ 的态射的复合具有性质 $P$。
\item 若 $f:X\to X'$、$g:Y\to Y'$ 为具有性质 $P$ 的 $S$-态射，
则 $f\times_S g$ 具有性质 $P$。
\item 若 $f:X\to Y$ 为具有性质 $P$ 的 $S$-态射，
则任意基预概形扩张 $S'\to S$ 所得的 $f_{(S')}$ 具有性质 $P$。
\item 若 $f:X\to Y$、$g:Y\to Z$ 的复合 $g\circ f$ 具有性质 $P$，
且 $g$ 分离，则 $f$ 具有性质 $P$。
\item 若 $f:X\to Y$ 具有性质 $P$，则 $f_{\mathrm{red}}$（\hyperlink{para:1:5.1.5}{5.1.5}）具有性质 $P$。
\end{enumerate}
若 (i)、(ii) 成立，则 (iii) 与 (iv) 等价，且 (v)、(vi) 都是 (i)、(ii)、(iii) 的推论。

第一断言已在（\hyperlink{para:1:3.5.1}{3.5.1}）中证明。
考虑分解（\hyperlink{para:1:5.3.13}{5.3.13}）$X\xrightarrow{\Gamma_f}X\times_Z Y\xrightarrow{p_2}Y$。
等式 $p_2=(g\circ f)\times_Z1_Y$ 表明：
若 $g\circ f$ 具有性质 $P$，则 $p_2$ 由 (iii) 具有性质 $P$。
若 $g$ 分离，则 $\Gamma_f$ 为闭浸入（\hyperlink{para:1:5.4.3}{5.4.3}），因而由 (i) 具有性质 $P$；
最后由 (ii)，$f$ 具有性质 $P$。

再考虑交换图
\[
\begin{tikzcd}
X_{\mathrm{red}}\arrow[r,"f_{\mathrm{red}}"]\arrow[d]&Y_{\mathrm{red}}\arrow[d]\\
X\arrow[r,"f"']&Y
\end{tikzcd}
\]
其中竖箭头为闭浸入（\hyperlink{para:1:5.1.5}{5.1.5}），故由 (i) 具有性质 $P$。
若 $f$ 具有性质 $P$，则由 (ii)，复合态射
$X_{\mathrm{red}}\to Y_{\mathrm{red}}\to Y$ 具有性质 $P$。
由于闭浸入分离（\hyperlink{para:1:5.5.1}{5.5.1}，(i)），由 (v) 得 $f_{\mathrm{red}}$ 具有性质 $P$。

若考虑下述两个命题：
\begin{enumerate}[label=\textup{(\roman*')},start=1]
\item 每个浸入具有性质 $P$；
\end{enumerate}
\noindent(v') 若 $g\circ f$ 具有性质 $P$，则 $f$ 具有性质 $P$；
则以上论证表明，(v') 是 (i')、(ii)、(iii) 的推论。

\egapar{5.5.13}
还应注意，(v)、(vi) 也是 (i)、(iii) 以及下述命题的推论：

\noindent(ii') 若 $j:X\to Y$ 为闭浸入，$g:Y\to Z$ 为具有性质 $P$ 的态射，
则 $g\circ j$ 具有性质 $P$。

同样，(v') 是 (i')、(iii) 以及下述命题的推论：

\noindent(ii'') 若 $j:X\to Y$ 为浸入，$g:Y\to Z$ 为具有性质 $P$ 的态射，
则 $g\circ j$ 具有性质 $P$。

这些断言由（\hyperlink{para:1:5.5.12}{5.5.12}）的论证立即得到。


% END chapters/ch1_s5_5.tex

\par

% BEGIN chapters/ch1_s6_1.tex
\section{有限性条件}\hypertarget{heading:1:6}{}
\subsection{诺特预概形与局部诺特预概形}\hypertarget{heading:1:6.1}{}
\src{140—143}

\egastatement{6.1.1}{定义}
若预概形 $X$ 是有限个仿射开集 $V_i$ 的并（分别地，任意族仿射开集的并），
且每个 $V_i$ 上诱导的仿射概形的环都是诺特环，
则称 $X$ 为\emph{诺特预概形}（分别地，\emph{局部诺特预概形}）。

由（\hyperlink{para:1:1.5.2}{1.5.2}）及问题的局部性立即可知：若 $X$ 局部诺特，
则结构层 $\OO_X$ 是凝聚环层。
此时凝聚 $\OO_X$-模 $\F$ 的每个拟凝聚子模（分别地，拟凝聚商模）都是凝聚的。
问题仍具有局部性；只须应用（\hyperlink{para:1:1.5.1}{1.5.1}）、（\hyperlink{para:1:1.4.1}{1.4.1}）、（\hyperlink{para:1:1.3.10}{1.3.10}），
并利用诺特环上有限型模的子模、商模仍为有限型模的事实。
特别地，$\OO_X$ 的每个拟凝聚理想层都是凝聚的。

若 $X$ 是有限个开集 $W_i$ 的并（分别地，任意族开集的并），
且每个 $W_i$ 上诱导的预概形为诺特（分别地，局部诺特）的，
则显然 $X$ 为诺特（分别地，局部诺特）的。

\egastatement{6.1.2}{命题}
预概形 $X$ 为诺特的充要条件是它局部诺特且底空间拟紧。
此时 $X$ 的底空间是诺特空间。

第一断言由定义及（\hyperlink{para:1:1.1.10}{1.1.10}，(ii)）立即得到。
第二断言由（\hyperlink{para:1:1.1.6}{1.1.6}）及有限个诺特子空间之并为诺特空间的事实（0，\hyperlink{para:0:2.2.3}{2.2.3}）得到。

\egastatement{6.1.3}{命题}
设 $X$ 为环为 $A$ 的仿射概形。
下述条件等价：a) $X$ 诺特；b) $X$ 局部诺特；c) $A$ 诺特。

a)、b) 的等价性由（\hyperlink{para:1:6.1.2}{6.1.2}）及仿射概形的底空间拟紧（\hyperlink{para:1:1.1.10}{1.1.10}）得到；
c) 显然蕴含 a)。
为证明 a) 蕴含 c)，取 $X$ 的有限仿射开覆盖 $(V_i)$，使每个 $A_i=A(V_i)$ 诺特。
设 $(\mathfrak a_n)$ 为 $A$ 的递增理想序列。
它典范地一一对应于 $\wt A=\OO_X$ 中的递增理想层序列
$(\wt{\mathfrak a_n})$（\hyperlink{para:1:1.3.7}{1.3.7}）。
只须证明后一序列稳定。
限制 $\wt{\mathfrak a_n}|V_i$ 是 $\OO_X|V_i$ 的拟凝聚理想层，
因为它是沿典范嵌入 $V_i\to X$ 的逆像（0，\hyperlink{para:0:5.1.4}{5.1.4}）。
故它形如 $\wt{\mathfrak a_{ni}}$，其中 $\mathfrak a_{ni}$ 为 $A_i$ 的理想（\hyperlink{para:1:1.3.7}{1.3.7}）。
由于 $A_i$ 诺特，每个序列 $(\mathfrak a_{ni})_n$ 稳定，从而命题成立。

还应注意，这个论证也证明：若 $X$ 为诺特预概形，
则 $\OO_X$ 的每个递增凝聚理想层序列均稳定。

\egastatement{6.1.4}{命题}
诺特（分别地，局部诺特）预概形的每个子预概形仍为诺特（分别地，局部诺特）的。

只须对诺特预概形 $X$ 证明；又由定义（\hyperlink{para:1:6.1.1}{6.1.1}）立即可归结到 $X$ 仿射的情形。
每个子预概形都是某个开集上诱导的预概形的闭子预概形（\hyperlink{para:1:4.1.3}{4.1.3}），
所以只须分别考虑闭子预概形 $Y$ 与在开集上诱导的子预概形 $Y$。
闭的情形立即成立：若 $X$ 的环为 $A$，则 $Y$ 仿射且其环为 $A/\mathfrak J$（\hyperlink{para:1:4.2.3}{4.2.3}）；
$A$ 诺特（\hyperlink{para:1:6.1.3}{6.1.3}），故 $A/\mathfrak J$ 诺特。
若 $Y$ 在 $X$ 中开，则底空间 $Y$ 诺特（\hyperlink{para:1:6.1.2}{6.1.2}），所以拟紧，
可表示为有限个 $D(f_i)$（$f_i\in A$）的并。
因而归结到 $Y=D(f)$ 的情形。此时 $Y$ 仿射，环同构于 $A_f$（\hyperlink{para:1:1.3.6}{1.3.6}）；
$A$ 诺特，故 $A_f$ 诺特。

\egapar{6.1.5}
两个诺特 $S$-预概形的积未必诺特，即使它们都是仿射概形也是如此；
因为两个诺特代数的张量积未必是诺特环（参见 \hyperlink{para:1:6.3.8}{6.3.8}）。

\egastatement{6.1.6}{命题}
若 $X$ 为诺特预概形，则 $\OO_X$ 的幂零根 $\N_X$ 是幂零的。

可用有限个仿射开集 $U_i$ 覆盖 $X$。
只须证明存在整数 $n_i$ 使 $(\N_X|U_i)^{n_i}=0$；
若 $n$ 为这些 $n_i$ 的最大值，便有 $\N_X^n=0$。
所以归结到 $X=\Spec(A)$、$A$ 诺特的情形。
根据（\hyperlink{para:1:5.1.1}{5.1.1}）和（\hyperlink{para:1:1.3.13}{1.3.13}），只须注意诺特环的幂零根是幂零的（文献 \cite{b11}，127 页，推论 4）。

\egastatement{6.1.7}{推论}
设 $X$ 为诺特预概形。$X$ 为仿射概形的充要条件是 $X_{\mathrm{red}}$ 为仿射概形。

这由（\hyperlink{para:1:6.1.6}{6.1.6}）和（\hyperlink{para:1:5.1.10}{5.1.10}）得到。

\egastatement{6.1.8}{引理}
设 $X$ 为拓扑空间，$x\in X$，$U$ 为 $x$ 的一个只有有限个不可约分支的开邻域。
则存在 $x$ 的邻域 $V$，使包含于 $V$ 的每个 $x$ 的开邻域都连通。

设 $U_i$（$1\leq i\leq m$）为 $U$ 中不包含 $x$ 的不可约分支。
它们的并在 $U$ 中的补集 $V$ 是 $x$ 在 $U$ 中、从而在 $X$ 中的开邻域。
它也是 $U$ 与 $X$ 中所有不包含 $x$ 的不可约分支之并的补集的交（0，\hyperlink{para:0:2.1.6}{2.1.6}）。
\transnote{原文将 $V$ 直接称为上述并在 $X$ 中的补集；
一般还须与 $U$ 相交。例如 $X$ 不可约且 $U\ne X$ 时，原来的直接等式不成立。}
若 $W\subset V$ 为 $x$ 的开邻域，则 $W$ 的不可约分支是 $U$ 中与 $W$ 相交的不可约分支在 $W$ 上的迹（0，\hyperlink{para:0:2.1.6}{2.1.6}）。
所以它们都包含 $x$；由于它们连通，$W$ 也连通。

\egastatement{6.1.9}{推论}
局部诺特拓扑空间局部连通；特别地，其连通分支是开集。

\egastatement{6.1.10}{命题}
设 $X$ 为局部诺特拓扑空间。下述条件等价：
\begin{enumerate}[label=\alph*)]
\item $X$ 的不可约分支是开集；
\item $X$ 的不可约分支与连通分支相同；
\item $X$ 的连通分支不可约；
\item $X$ 的两个不同不可约分支不相交。
\end{enumerate}
若 $X$ 为预概形，则它们还等价于：

\noindent e) 对每个 $x\in X$，$\Spec(\OO_x)$ 不可约，即 $\OO_x$ 的幂零根是素理想。

a) 蕴含 b) 立即成立：不可约空间连通，且 a) 使各不可约分支既开又闭。
b) 显然蕴含 c)。反之，包含一个连通分支 $C$ 且不同于 $C$ 的闭集 $F$ 不可能不可约：
$F$ 不连通，故为两个非空、不交、在 $F$ 中既开又闭、因而在 $X$ 中闭的子集之并。
所以 c) 蕴含 b)。由此立即得到 c) 蕴含 d)，因为不同连通分支不相交。

以上尚未用到局部诺特假设。现利用此假设证明 d) 蕴含 a)。
由（0，\hyperlink{para:0:2.1.6}{2.1.6}），可限于 $X$ 诺特的情形。
此时不可约分支只有有限个，它们闭且两两不交，故也开。

最后，d) 与 e) 的等价性无需假设底空间局部诺特。
由（0，\hyperlink{para:0:2.1.6}{2.1.6}），可归结到 $X=\Spec(A)$ 仿射的情形。
$x$ 只属于一个不可约分支，意味着 $\jj_x$ 只包含一个极小素理想（\hyperlink{para:1:1.1.14}{1.1.14}）；
这又等价于说 $\jj_x\OO_x$ 只包含 $\OO_x$ 的一个极小素理想，从而得到结论。

\egastatement{6.1.11}{推论}
设 $X$ 为局部诺特空间。
$X$ 不可约的充要条件是：$X$ 非空且连通，并且两个不同的不可约分支不相交。
若 $X$ 为预概形，最后一个条件等价于每个 $\Spec(\OO_x)$ 不可约。

最后部分已在（\hyperlink{para:1:6.1.10}{6.1.10}）中证明。
只须证明第一部分条件的充分性：由（\hyperlink{para:1:6.1.10}{6.1.10}），这些条件使不可约分支与连通分支相同；
$X$ 非空且连通，故不可约。

\egastatement{6.1.12}{推论}
设 $X$ 为局部诺特预概形。
$X$ 为整的充要条件是 $X$ 非空且连通，并且对每个 $x\in X$，$\OO_x$ 为整环。
\transnote{原文在本条只写“连通”，未另写“非空”。这里明确补出非空条件，
以与（0，\hyperlink{para:0:2.1.1}{2.1.1}）及（\hyperlink{para:1:6.1.11}{6.1.11}）中的不可约性约定一致。}

\egastatement{6.1.13}{命题}
设 $X$ 为局部诺特预概形，$x\in X$，且 $\OO_x$ 的幂零根 $\N_{X,x}$ 为素理想
（分别地，$\OO_x$ 约化；$\OO_x$ 为整环）。
则存在 $x$ 的一个不可约（分别地，约化；整）的开邻域 $U$。

只须考虑幂零根为素理想和幂零根为零两种情形，第三种假设是前两种的合取。
若 $\N_{X,x}$ 为素理想，则 $x$ 只属于 $X$ 的一个不可约分支 $Y$（\hyperlink{para:1:6.1.10}{6.1.10}）。
不包含 $x$ 的不可约分支之并是闭的，因为这些分支构成局部有限族。
其补集 $U$ 因而开且包含于 $Y$，所以不可约（0，\hyperlink{para:0:2.1.6}{2.1.6}）。
若 $\N_{X,x}=0$，则在 $x$ 的某个邻域内每个 $y$ 都有 $\N_{X,y}=0$：
$\N_X$ 拟凝聚（\hyperlink{para:1:5.1.1}{5.1.1}），因 $X$ 局部诺特而凝聚，故结论由（0，\hyperlink{para:0:5.2.2}{5.2.2}）得到。


% END chapters/ch1_s6_1.tex

\par

% BEGIN chapters/ch1_s6_2.tex
\subsection{阿廷预概形}\hypertarget{heading:1:6.2}{}
\src{143—144}

\egastatement{6.2.1}{定义}
若一个预概形仿射且其环为阿廷环，则称它为\emph{阿廷预概形}。

\egastatement{6.2.2}{命题}
对于预概形 $X$，下述条件等价：
\begin{enumerate}[label=\alph*)]
\item $X$ 为阿廷概形；
\item $X$ 诺特且其底空间离散；
\item $X$ 诺特且其底空间的各点都是闭点，即满足条件 $T_1$。
\end{enumerate}
此时 $X$ 的底空间有限，且 $X$ 的环 $A$ 为各点的局部环（均为阿廷环）的直积。

已知 a) 蕴含最后的断言（文献 \cite{b13}，205 页，定理 3）。
此时 $A$ 的每个素理想都是极大理想，并且是一个局部直积分量的极大理想的逆像，
故 $X$ 有限且离散。于是 a) 蕴含 b)，而 b) 显然蕴含 c)。
为证明 c) 蕴含 a)，先证明 $X$ 有限。
可归结到 $X$ 仿射的情形；一个各素理想均极大的诺特环为阿廷环
（文献 \cite{b13}，203 页），故所断言的有限性成立。
此时 $X$ 的底空间离散，为有限个点 $x_i$ 的拓扑和，局部环 $\OO_{x_i}=A_i$ 为阿廷环。
显然 $X$ 同构于环 $A=\prod_i A_i$ 的素谱所定义的仿射概形（\hyperlink{para:1:1.7.3}{1.7.3}）。


% END chapters/ch1_s6_2.tex

\par

% BEGIN chapters/ch1_s6_3.tex
\subsection{有限型态射}\hypertarget{heading:1:6.3}{}
\src{144—147}

\egastatement{6.3.1}{定义}
若预概形态射 $f:X\to Y$ 满足：$Y$ 是一族具有下述性质的仿射开集 $(V_\alpha)$ 的并，
则称 $f$ 为\emph{有限型态射}：

\noindent(P) $f^{-1}(V_\alpha)$ 是有限个仿射开集 $U_{\alpha\beta}$ 的并，
且每个环 $A(U_{\alpha\beta})$ 都是 $A(V_\alpha)$ 上的有限型代数。

此时也称 $X$ 是 $Y$ 上的\emph{有限型预概形}，或称 $X$ 为\emph{有限型 $Y$-预概形}。

\egastatement{6.3.2}{命题}
若 $f:X\to Y$ 为有限型态射，则 $Y$ 的每个仿射开集 $W$
都具有定义（\hyperlink{para:1:6.3.1}{6.3.1}）中的性质 (P)。

先证明下述引理。

\egastatement{6.3.2.1}{引理}
若仿射开集 $T\subset Y$ 具有性质 (P)，则对每个 $g\in A(T)$，$D(g)$ 也具有性质 (P)。

由假设，$f^{-1}(T)$ 是有限个仿射开集 $Z_j$ 的并，且每个 $A(Z_j)$ 都是有限型 $A(T)$-代数。
设 $\varphi_j:A(T)\to A(Z_j)$ 为对应于 $f|Z_j$ 的环同态（\hyperlink{para:1:2.2.4}{2.2.4}），置 $g_j=\varphi_j(g)$。
则 $f^{-1}(D(g))\cap Z_j=D(g_j)$（\hyperref[eq:1:1.2.2.2]{1.2.2.2}）。
而 $A(D(g_j))=A(Z_j)_{g_j}=A(Z_j)[1/g_j]$ 在 $A(Z_j)$ 上有限型，
由假设在 $A(T)$ 上也有限型，从而在 $A(D(g))=A(T)[1/g]$ 上有限型。引理得证。

由于 $W$ 拟紧（\hyperlink{para:1:1.1.10}{1.1.10}），存在 $W$ 的有限覆盖，成员形如 $D(g_i)$，
每个 $g_i$ 属于某个 $A(V_{\alpha(i)})$。
每个 $D(g_i)$ 拟紧，故为有限个 $D(h_{ik})$ 的并，其中 $h_{ik}\in A(W)$。
若 $\varphi_i:A(W)\to A(D(g_i))$ 为典范映射，则
$D(h_{ik})=D(\varphi_i(h_{ik}))$（\hyperref[eq:1:1.2.2.2]{1.2.2.2}）。
由引理（\hyperlink{para:1:6.3.2.1}{6.3.2.1}），每个 $f^{-1}(D(h_{ik}))$ 都有有限仿射开覆盖 $(U_{ikj})$，
使 $A(U_{ikj})$ 在 $A(D(h_{ik}))=A(W)[1/h_{ik}]$ 上有限型，命题遂得证。

所以可以说，$Y$ 上有限型预概形的概念在 $Y$ 上具有局部性。

\egastatement{6.3.3}{命题}
设 $X,Y$ 为两个仿射概形。
$X$ 在 $Y$ 上有限型的充要条件是 $A(X)$ 为有限型 $A(Y)$-代数。

充分性显然，现证明必要性。置 $A=A(Y)$、$B=A(X)$。
由（\hyperlink{para:1:6.3.2}{6.3.2}），$X$ 有有限仿射开覆盖 $(V_i)$，使各 $A(V_i)$ 为有限型 $A$-代数。
由于 $V_i$ 拟紧，可用有限个 $D(g_{ij})\subset V_i$、$g_{ij}\in B$ 覆盖它。
若 $\varphi_i:B\to A(V_i)$ 对应于典范嵌入 $V_i\to X$，则
\[
 B_{g_{ij}}=A(V_i)_{\varphi_i(g_{ij})}=A(V_i)[1/\varphi_i(g_{ij})],
\]
故 $B_{g_{ij}}$ 为有限型 $A$-代数。
于是可归结到 $V_i=D(g_i)$、$g_i\in B$ 的情形。
由假设存在有限子集 $F_i\subset B$ 及整数 $n_i\geq0$，
使 $B_{g_i}$ 由元素 $b_i/g_i^{n_i}$（$b_i\in F_i$）作为 $A$-代数生成。
因 $g_i$ 只有有限个，可假设各 $n_i$ 相等。
又因 $D(g_i)$ 覆盖 $X$，各 $g_i$ 在 $B$ 中生成单位理想，
故存在 $h_i\in B$ 使 $\sum_i h_i g_i=1$。

令 $F$ 为各 $F_i$ 以及所有 $g_i,h_i$ 的并。证明 $B'=A[F]$ 等于 $B$。
对任意 $b\in B$ 及每个 $i$，$b$ 在 $B_{g_i}$ 中的典范像形如
$b'_i/g_i^{m_i}$，其中 $b'_i\in B'$。
乘以适当的 $g_i$ 的幂后可使各 $m_i$ 同为整数 $m$。
由分式环的定义，存在依赖于 $b$ 的整数 $N>m$，使 $g_i^N b\in B'$ 对每个 $i$ 成立。
在环 $B'$ 中，各 $g_i^N$ 生成单位理想，因为各 $g_i$ 已经如此，且 $h_i\in B'$。
故有 $c_i\in B'$ 使 $\sum_i c_i g_i^N=1$，从而
$b=\sum_i c_i g_i^N b\in B'$。证毕。

\egastatement{6.3.4}{命题}
\begin{enumerate}[label=(\roman*)]
\item 每个闭浸入为有限型态射。
\item 两个有限型态射的复合为有限型态射。
\item 两个有限型 $S$-态射 $f:X\to X'$、$g:Y\to Y'$ 的积 $f\times_S g$ 为有限型态射。
\item 若 $f:X\to Y$ 为有限型 $S$-态射，则任意基预概形扩张 $S'\to S$ 所得的 $f_{(S')}$ 为有限型态射。
\item 若 $g\circ f$ 为有限型态射且 $g$ 分离，则 $f$ 为有限型态射。
\item 若 $f$ 为有限型态射，则 $f_{\mathrm{red}}$ 为有限型态射。
\end{enumerate}

根据（\hyperlink{para:1:5.5.12}{5.5.12}），只须证明 (i)、(ii)、(iv)。
为证明 (i)，可限于闭子预概形的典范嵌入 $X\to Y$，
并可假设 $Y$ 仿射（\hyperlink{para:1:6.3.2}{6.3.2}）。
此时 $X$ 仿射（\hyperlink{para:1:4.2.3}{4.2.3}），其环为 $A/\mathfrak J$，其中 $A=A(Y)$。
$A/\mathfrak J$ 在 $A$ 上有限型，故结论成立。

为证明 (ii)，设 $f:X\to Y$、$g:Y\to Z$ 为有限型态射，$U$ 为 $Z$ 的仿射开集。
$g^{-1}(U)$ 有有限仿射开覆盖 $(V_i)$，使 $A(V_i)$ 在 $A(U)$ 上有限型（\hyperlink{para:1:6.3.2}{6.3.2}）。
同样，每个 $f^{-1}(V_i)$ 有有限仿射开覆盖 $(W_{ij})$，
使 $A(W_{ij})$ 在 $A(V_i)$ 上有限型，故在 $A(U)$ 上有限型。结论成立。

最后证明 (iv)。可限于 $S=Y$ 的情形：
$f_{(S')}$ 也等于把 $f$ 看作 $Y$-态射、沿 $Y_{(S')}\to Y$ 扩张所得的 $f_{(Y_{(S')})}$（\hyperlink{para:1:3.3.9}{3.3.9}）。
以 $p:X_{(S')}\to X$、$q:X_{(S')}\to S'$ 表示投影。
若 $V\subset S$ 为仿射开集，则 $f^{-1}(V)$ 是有限个仿射开集 $W_i$ 的并，
且各 $A(W_i)$ 在 $A(V)$ 上有限型（\hyperlink{para:1:6.3.2}{6.3.2}）。
设 $V'\subset S'$ 为包含于 $g^{-1}(V)$ 的仿射开集。
由 $f\circ p=g\circ q$，$q^{-1}(V')$ 包含于各 $p^{-1}(W_i)$ 的并。
而 $p^{-1}(W_i)\cap q^{-1}(V')$ 认同于 $W_i\times_V V'$（\hyperlink{para:1:3.2.7}{3.2.7}），
它仿射，其环同构于 $A(W_i)\otimes_{A(V)}A(V')$（\hyperlink{para:1:3.2.2}{3.2.2}）。
此环在 $A(V')$ 上有限型，命题得证。

\egastatement{6.3.5}{推论}
设 $f:X\to Y$ 为浸入。
若 $Y$ 的底空间局部诺特，或 $X$ 的底空间诺特，则 $f$ 为有限型态射。

总可假设 $Y$ 仿射（\hyperlink{para:1:6.3.2}{6.3.2}）。
若 $Y$ 的底空间局部诺特，还可假设它诺特；此时子空间 $X$ 也诺特。
所以可以假设 $Y$ 仿射、$X$ 的底空间诺特。
由 $X$ 在 $Y$ 中局部闭（\hyperlink{para:1:4.1.3}{4.1.3}），存在由有限个 $Y$ 的仿射开集
$D(g_i)$（$g_i\in A(Y)$）组成的 $X$ 的覆盖，
使 $X\cap D(g_i)$ 在 $D(g_i)$ 中闭，因而仿射（\hyperlink{para:1:4.2.3}{4.2.3}）。
由（\hyperlink{para:1:6.3.4}{6.3.4}，(i)）和（\hyperlink{para:1:6.3.3}{6.3.3}），$A(X\cap D(g_i))$ 在 $A(D(g_i))$ 上有限型；
而 $A(D(g_i))=A(Y)[1/g_i]$ 在 $A(Y)$ 上有限型，证明完成。

\egastatement{6.3.6}{推论}
设 $f:X\to Y$、$g:Y\to Z$ 为态射。
若 $g\circ f$ 为有限型态射，且 $X$ 诺特或 $X\times_Z Y$ 局部诺特，则 $f$ 为有限型态射。

这由（\hyperlink{para:1:5.5.12}{5.5.12}）的证明以及将（\hyperlink{para:1:6.3.5}{6.3.5}）应用于图浸入 $\Gamma_f$ 立即得到。

\egastatement{6.3.7}{命题}
设 $f:X\to Y$ 为有限型态射。
若 $Y$ 诺特（分别地，局部诺特），则 $X$ 诺特（分别地，局部诺特）。

只须在 $Y$ 诺特时证明。
此时 $Y$ 是有限个仿射开集 $V_i$ 的并，且各 $A(V_i)$ 诺特。
由（\hyperlink{para:1:6.3.2}{6.3.2}），每个 $f^{-1}(V_i)$ 为有限个仿射开集 $W_{ij}$ 的并，
各 $A(W_{ij})$ 为诺特环 $A(V_i)$ 上的有限型代数，故为诺特环。
所以 $X$ 诺特。

\egastatement{6.3.8}{推论}
设 $X$ 为有限型 $S$-预概形。
对任意使 $S'$ 诺特（分别地，局部诺特）的基扩张 $S'\to S$，
$X_{(S')}$ 诺特（分别地，局部诺特）。

由（\hyperlink{para:1:6.3.4}{6.3.4}，(iv)），$X_{(S')}$ 在 $S'$ 上有限型，再用（\hyperlink{para:1:6.3.7}{6.3.7}）即可。
也可表述为：在 $S$-预概形的积 $X\times_S Y$ 中，
若一因子在 $S$ 上有限型，另一因子诺特（分别地，局部诺特），
则该积诺特（分别地，局部诺特）。

\egastatement{6.3.9}{推论}
设 $S$ 为局部诺特预概形，$X$ 为有限型 $S$-预概形。
则每个 $S$-态射 $f:X\to Y$ 都为有限型态射。

可假设 $S$ 诺特。若 $\varphi:X\to S$、$\psi:Y\to S$ 为结构态射，
则 $\varphi=\psi\circ f$；由（\hyperlink{para:1:6.3.7}{6.3.7}），$X$ 诺特，
故 $f$ 由（\hyperlink{para:1:6.3.6}{6.3.6}）为有限型态射。

\egastatement{6.3.10}{命题}
设 $f:X\to Y$ 为有限型态射。
$f$ 为满射的充要条件是：对每个代数闭域 $\Omega$，
由 $f$ 对应的映射 $X(\Omega)\to Y(\Omega)$（\hyperlink{para:1:3.4.1}{3.4.1}）为满射。

条件充分：对每个 $y\in Y$，考虑 $k(y)$ 的代数闭扩张 $\Omega$，
及交换图
\[
\begin{tikzcd}
X\arrow[d,"f"']&\Spec(\Omega)\arrow[l]\arrow[dl]\\
Y&
\end{tikzcd}
\]
即可看出（参见 \hyperlink{para:1:3.5.3}{3.5.3}）。
反之，设 $f$ 满射，$g:\Spec(\Omega)\to Y$ 为态射，$\Omega$ 代数闭。
考虑图
\[
\begin{tikzcd}
X\arrow[d,"f"']&X_{(\Omega)}\arrow[l]\arrow[d,"f_{(\Omega)}"]\\
Y&\Spec(\Omega)\arrow[l,"g"]
\end{tikzcd}
\]
只须证明 $X_{(\Omega)}$ 有一个 $\Omega$-有理点（\hyperlink{para:1:3.3.14}{3.3.14}、\hyperlink{para:1:3.4.3}{3.4.3}、\hyperlink{para:1:3.4.4}{3.4.4}）。
$f$ 满射，故 $X_{(\Omega)}$ 非空；$f$ 有限型，故 $f_{(\Omega)}$ 有限型（\hyperlink{para:1:6.3.4}{6.3.4}，(iv)）。
所以 $X_{(\Omega)}$ 有非空仿射开集 $Z$，其环 $A(Z)$ 为非零有限型 $\Omega$-代数。
由希尔伯特零点定理（文献 \cite{b21}），存在 $\Omega$-同态 $A(Z)\to\Omega$，
从而存在 $X_{(\Omega)}$ 在 $\Spec(\Omega)$ 上的截面，命题得证。


% END chapters/ch1_s6_3.tex

\par

% BEGIN chapters/ch1_s6_4.tex
\subsection{代数预概形}\hypertarget{heading:1:6.4}{}
\src{147—150}

\egastatement{6.4.1}{定义}
给定域 $K$，称 $K$ 上的有限型预概形 $X$ 为\emph{代数 $K$-预概形}；
$K$ 称为 $X$ 的\emph{基域}。
若 $X$ 还是概形，等价地是 $K$-概形（\hyperlink{para:1:5.5.8}{5.5.8}），则称它为\emph{代数 $K$-概形}。

每个代数 $K$-预概形都是诺特的（\hyperlink{para:1:6.3.7}{6.3.7}）。

\egastatement{6.4.2}{命题}
设 $X$ 为代数 $K$-预概形。
点 $x\in X$ 为闭点的充要条件是 $k(x)$ 为 $K$ 的有限次代数扩张。

可假设 $X$ 仿射，其环 $A$ 为有限型 $K$-代数。
事实上，满足 $A(U)$ 为有限型 $K$-代数的仿射开集 $U$ 构成 $X$ 的覆盖（\hyperlink{para:1:6.3.1}{6.3.1}）。
此时 $X$ 的闭点正是对应理想 $\jj_x$ 极大，亦即 $A/\jj_x$ 为域的点，
且该域必然为 $k(x)$。
$A/\jj_x$ 是有限型 $K$-代数，所以若 $x$ 闭，$k(x)$ 是作为 $K$-代数有限型的域，
故必为有限秩 $K$-代数（文献 \cite{b21}）。
反之，若 $k(x)$ 在 $K$ 上秩有限，则 $A/\jj_x\subset k(x)$ 在 $K$ 上也秩有限。
每个作为有限秩 $K$-代数的整环都是域，故 $A/\jj_x=k(x)$，所以 $x$ 闭。

\egastatement{6.4.3}{推论}
设 $K$ 为代数闭域，$X$ 为代数 $K$-预概形。
则 $X$ 的闭点正是 $K$-有理点（\hyperlink{para:1:3.4.4}{3.4.4}），并典范地认同于 $X$ 的 $K$-值点。

\egastatement{6.4.4}{命题}
设 $X$ 为域 $K$ 上的代数预概形。下述性质等价：
\begin{enumerate}[label=\alph*)]
\item $X$ 为阿廷的；
\item $X$ 的底空间离散；
\item $X$ 的底空间只有有限个闭点；
\end{enumerate}
\noindent c') $X$ 的底空间有限；
\begin{enumerate}[label=\alph*),start=4]
\item $X$ 的所有点都是闭点；
\item $X$ 同构于 $\Spec(A)$，其中 $A$ 为有限秩 $K$-代数。
\end{enumerate}

由于 $X$ 诺特，由（\hyperlink{para:1:6.2.2}{6.2.2}），a)、b)、d) 等价，并蕴含 c)、c')；
e) 显然蕴含 a)。剩下证明 c) 蕴含 d)、e)；可限于 $X$ 仿射的情形。
此时 $A(X)$ 是有限型 $K$-代数（\hyperlink{para:1:6.3.3}{6.3.3}），所以是雅各布森环
（文献 \cite{b1}，3-11 至 3-12 页），并且由假设只有有限个极大理想。
有限个素理想的交若仍为素理想，必等于其中一个，
所以 $A(X)$ 的每个素理想都极大，d) 得证。
此外，已知此时 $A(X)$ 是有限型阿廷 $K$-代数（\hyperlink{para:1:6.2.2}{6.2.2}），
所以必为有限秩 $K$-代数（文献 \cite{b21}）。

\egapar{6.4.5}
当（\hyperlink{para:1:6.4.4}{6.4.4}）的条件满足时，称 $X$ 为 $K$ 上的\emph{有限概形}
（参见 II，6.1.1），或\emph{有限 $K$-概形}。
其秩为 $[A:K]$，亦记为 $\operatorname{rg}_K(X)$。
若 $X,Y$ 为两个有限 $K$-概形，则
\begin{equation}
 \operatorname{rg}_K(X\amalg Y)=\operatorname{rg}_K(X)+\operatorname{rg}_K(Y),
 \tag{6.4.5.1}\label{eq:1:6.4.5.1}
\end{equation}
\begin{equation}
 \operatorname{rg}_K(X\times_KY)=\operatorname{rg}_K(X)\operatorname{rg}_K(Y),
 \tag{6.4.5.2}\label{eq:1:6.4.5.2}
\end{equation}
这由（\hyperlink{para:1:3.2.2}{3.2.2}）得到。

\egastatement{6.4.6}{推论}
设 $X$ 为有限 $K$-概形。
对 $K$ 的任意扩张 $K'$，$X\otimes_K K'$ 为有限 $K'$-概形，
其在 $K'$ 上的秩等于 $X$ 在 $K$ 上的秩。

事实上，若 $A=A(X)$，则 $[A\otimes_K K':K']=[A:K]$。

\egastatement{6.4.7}{推论}
设 $X$ 为有限 $K$-概形，置
\[
 n=\sum_{x\in X}[k(x):K]_s.
\]
这里回忆：$[K':K]_s$ 表示扩张 $K'/K$ 的可分次数，
即包含于 $K'$ 的 $K$ 的最大代数可分扩张的次数。
则对 $K$ 的任意代数闭扩张 $\Omega$，$X\otimes_K\Omega$ 的底空间恰有 $n$ 个点，
这些点与 $X$ 的 $\Omega$-值点相认同。

显然可限于 $A=A(X)$ 为局部环的情形（\hyperlink{para:1:6.2.2}{6.2.2}）。
设 $\mathfrak m$ 为其极大理想，$L=A/\mathfrak m$ 为其剩余域，即 $K$ 的一个代数扩张。
$X$ 的 $\Omega$-值点与 $X\otimes_K\Omega$ 的 $\Omega$-截面一一对应（\hyperlink{para:1:3.4.1}{3.4.1} 和 \hyperlink{para:1:3.3.14}{3.3.14}），
也与 $L\to\Omega$ 的 $K$-同态一一对应（\hyperlink{para:1:1.7.3}{1.7.3}）。
因此，由 Bourbaki《代数》第 V 章 §7 第 5 节命题 8，并利用（\hyperlink{para:1:6.4.3}{6.4.3}），得到本推论。

\egapar{6.4.8}
（\hyperlink{para:1:6.4.7}{6.4.7}）中定义的数 $n$ 称为 $A$ 或 $X$ 在 $K$ 上的\emph{可分秩}，
也称为 $X$ 的\emph{几何点数}；它等于 $X(\Omega)_K$ 的元素数。
由此定义立即可知，对 $K$ 的任意扩张 $K'$，$X\otimes_K K'$ 与 $X$ 的几何点数相同。
以 $n(X)$ 表示此数。若 $X,Y$ 为两个有限 $K$-概形，显然
\begin{equation}
 n(X\amalg Y)=n(X)+n(Y).
 \tag{6.4.8.1}\label{eq:1:6.4.8.1}
\end{equation}
同样有
\begin{equation}
 n(X\times_KY)=n(X)n(Y),
 \tag{6.4.8.2}\label{eq:1:6.4.8.2}
\end{equation}
这由将 $n(X)$ 解释为 $X(\Omega)_K$ 的元素数及公式（\hyperref[eq:1:3.4.3.1]{3.4.3.1}）立即得到。

\egastatement{6.4.9}{命题}
设 $K$ 为域，$X,Y$ 为两个代数 $K$-预概形，$f:X\to Y$ 为 $K$-态射，
$\Omega$ 为 $K$ 的一个代数闭扩张，且在 $K$ 上的超越次数无限。
则 $f$ 为满射的充要条件是对应映射
$X(\Omega)_K\to Y(\Omega)_K$（\hyperlink{para:1:3.4.1}{3.4.1}）为满射。

必要性由（\hyperlink{para:1:6.3.10}{6.3.10}）得到，因为 $f$ 必为有限型态射（\hyperlink{para:1:6.3.9}{6.3.9}）。
充分性的论证与（\hyperlink{para:1:6.3.10}{6.3.10}）相同；只须注意，对每个 $y\in Y$，
$k(y)$ 是 $K$ 的有限型扩张，所以 $K$-同构于 $\Omega$ 的一个子域。

\egastatement{6.4.10}{注}
将在第 IV 章看到，（\hyperlink{para:1:6.4.9}{6.4.9}）的结论不加关于 $\Omega/K$ 的超越次数的假设也成立。

\egastatement{6.4.11}{命题}
若 $f:X\to Y$ 为有限型态射，则对每个 $y\in Y$，纤维 $f^{-1}(y)$
是剩余域 $k(y)$ 上的代数预概形；且对每个 $x\in f^{-1}(y)$，
$k(x)$ 是 $k(y)$ 的有限型扩张。

由于 $f^{-1}(y)=X\otimes_Y k(y)$，命题由（\hyperlink{para:1:6.3.4}{6.3.4}，(iv)）和（\hyperlink{para:1:6.3.3}{6.3.3}）得到。

\egastatement{6.4.12}{命题}
设 $f:X\to Y$、$g:Y'\to Y$ 为两个态射，置 $X'=X\times_Y Y'$、$f'=f_{(Y')}$。
取 $y'\in Y'$，令 $y=g(y')$。
若纤维 $f^{-1}(y)$ 为有限 $k(y)$-概形，则纤维 $f'^{-1}(y')$ 为有限 $k(y')$-概形，
且其秩和几何点数分别等于 $f^{-1}(y)$ 的秩和几何点数。

利用纤维的传递性，结论由（\hyperlink{para:1:6.4.6}{6.4.6}）和（\hyperlink{para:1:6.4.8}{6.4.8}）立即得到。
\transnote{原文在此将纤维的传递性引为（\hyperlink{para:1:3.6.5}{3.6.5}）；对应的陈述为（\hyperlink{para:1:3.6.4}{3.6.4}）。}

\egapar{6.4.13}
命题（\hyperlink{para:1:6.4.11}{6.4.11}）表明，有限型态射直观上对应于“代数簇的代数族”，
其中 $Y$ 的点扮演“参数”的角色；这赋予此类态射以“几何”意义。
非有限型态射在下文中主要出现于“基预概形的变换”问题，例如局部化或完备化。


% END chapters/ch1_s6_4.tex

\par

% BEGIN chapters/ch1_s6_5.tex
\subsection{态射的局部确定}\hypertarget{heading:1:6.5}{}
\src{150—152}

\egastatement{6.5.1}{命题}
设 $X,Y$ 为两个 $S$-预概形，$Y$ 在 $S$ 上有限型；
设 $x\in X$、$y\in Y$ 位于同一点 $s\in S$ 之上。
\begin{enumerate}[label=(\roman*)]
\item 若两个 $S$-态射 $f=(\psi,\theta)$、$f'=(\psi',\theta')$ 从 $X$ 到 $Y$，
满足 $\psi(x)=\psi'(x)=y$，且局部 $\OO_s$-同态
$\theta_x^\sharp,\theta_x'^\sharp:\OO_y\to\OO_x$ 相同，
则 $f,f'$ 在 $x$ 的一个开邻域内相同。
\item 再设 $S$ 局部诺特。则对每个局部 $\OO_s$-同态
$\varphi:\OO_y\to\OO_x$，存在 $x$ 的开邻域 $U$
及 $S$-态射 $f=(\psi,\theta):U\to Y$，使
$\psi(x)=y$、$\theta_x^\sharp=\varphi$。
\end{enumerate}

(i) 问题在 $S,X,Y$ 上具有局部性，故可假设它们仿射，环分别为 $A,B,C$。
$f,f'$ 分别形如 $({}^a\varphi,\wt\varphi)$、$({}^a\varphi',\wt\varphi')$，
其中 $\varphi,\varphi':C\to B$ 为 $A$-同态，
$\varphi^{-1}(\jj_x)=\varphi'^{-1}(\jj_x)=\jj_y$，
且由它们诱导的同态 $\varphi_x,\varphi'_x:C_y\to B_x$ 相同。
还可假设 $C$ 为有限型 $A$-代数。
取生成元 $c_i$（$1\leq i\leq n$），置
$b_i=\varphi(c_i)$、$b'_i=\varphi'(c_i)$。
由假设在 $B_x$ 中 $b_i/1=b'_i/1$，
所以有 $s_i\in B\setminus\jj_x$ 使 $s_i(b_i-b'_i)=0$。
可取所有 $s_i$ 相同，记为 $g\in B\setminus\jj_x$。
于是 $B_g$ 中各 $b_i/1=b'_i/1$。
若 $i_g:B\to B_g$ 为典范同态，则 $i_g\circ\varphi=i_g\circ\varphi'$，
故 $f,f'$ 在 $D(g)$ 上的限制相同。

(ii) 可归结到与 (i) 相同的情形，并再假设 $A$ 诺特。
取 $A$-代数 $C$ 的生成元 $c_i$（$1\leq i\leq n$），
设 $\alpha:A[X_1,\ldots,X_n]\to C$ 为将 $X_i$ 映为 $c_i$ 的满同态。
设 $i_y:C\to C_y$ 为典范同态，考虑复合同态
\[
 \beta:A[X_1,\ldots,X_n]\xrightarrow{\alpha}C
 \xrightarrow{i_y}C_y\xrightarrow{\varphi}B_x.
\]
记 $\mathfrak a=\Ker\beta$。
由于 $A$ 诺特，$A[X_1,\ldots,X_n]$ 诺特，故 $\mathfrak a$ 有有限生成系
$Q_j(X_1,\ldots,X_n)$（$1\leq j\leq m$）。
每个 $\varphi(i_y(c_i))$ 可写成 $b_i/s_i$，其中 $b_i\in B$、$s_i\notin\jj_x$；
并可取各分母均为同一个 $g\notin\jj_x$。
由假设，在 $B_x$ 中有 $Q_j(b_1/g,\ldots,b_n/g)=0$。
写
\[
 Q_j(X_1/T,\ldots,X_n/T)=P_j(X_1,\ldots,X_n,T)/T^{k_j},
\]
其中 $P_j$ 为 $k_j$ 次齐次多项式。
置 $z_j=P_j(b_1,\ldots,b_n,g)\in B$。
由假设存在 $h_j\in B\setminus\jj_x$ 使 $h_jz_j=0$，并可使各 $h_j$ 同为 $h$。
于是 $P_j(hb_1,\ldots,hb_n,hg)=0$。
\transnote{此处可在齐次化时取 $k_j\geq1$；若直接取零次多项式的次数为零，
从 $hz_j=0$ 到所写等式的步骤需改为在 $B_{hg}$ 中验证。原文未说明这一选择。}

考虑将 $X_i$ 映为 $hb_i/(hg)$ 的同态
$\rho:A[X_1,\ldots,X_n]\to B_{hg}$。
$\mathfrak a$ 在该同态下的像为零，故 $\Ker\alpha\subset\mathfrak a$ 的像也为零。
所以 $\rho$ 分解为
$A[X_1,\ldots,X_n]\xrightarrow{\alpha}C\xrightarrow{\gamma}B_{hg}$，
且 $\gamma(c_i)=hb_i/(hg)$。
若 $i_x:B_{hg}\to B_x$ 为典范同态，则图
\begin{equation}
\begin{tikzcd}
C\arrow[r,"\gamma"]\arrow[d,"i_y"']&B_{hg}\arrow[d,"i_x"]\\
C_y\arrow[r,"\varphi"']&B_x
\end{tikzcd}
\tag{6.5.1.1}\label{eq:1:6.5.1.1}
\end{equation}
交换。因此 $\varphi=\gamma_x$；因 $\varphi$ 是局部同态，${}^a\gamma(x)=y$。
所以 $f=({}^a\gamma,\wt\gamma)$ 是从邻域 $D(hg)$ 到 $Y$ 的满足要求的 $S$-态射。

\egastatement{6.5.2}{推论}
在（\hyperlink{para:1:6.5.1}{6.5.1}，(ii)）的假设下，若 $X$ 在 $S$ 上也有限型，则可取 $f$ 为有限型态射。

这由（\hyperlink{para:1:6.3.6}{6.3.6}）得到。

\egastatement{6.5.3}{推论}
在（\hyperlink{para:1:6.5.1}{6.5.1}，(ii)）的假设下，再设 $Y$ 为整预概形且 $\varphi$ 为单同态。
则可取 $f=({}^a\gamma,\wt\gamma)$，其中 $\gamma$ 为单同态。

可假设 $C$ 为整环（\hyperlink{para:1:5.1.4}{5.1.4}），于是 $i_y$ 为单同态；
由图（\hyperref[eq:1:6.5.1.1]{6.5.1.1}）得到 $\gamma$ 为单同态。

\egastatement{6.5.4}{命题}
设 $f=(\psi,\theta):X\to Y$ 为有限型态射，$x\in X$、$y=\psi(x)$。
\begin{enumerate}[label=(\roman*)]
\item $f$ 在 $x$ 处为局部浸入（\hyperlink{para:1:4.5.1}{4.5.1}）的充要条件是
$\theta_x^\sharp:\OO_y\to\OO_x$ 为满射。
\item 再设 $Y$ 局部诺特。$f$ 在 $x$ 处为局部同构（\hyperlink{para:1:4.5.2}{4.5.2}）的充要条件是
$\theta_x^\sharp$ 为同构。
\end{enumerate}

(ii) 由（\hyperlink{para:1:6.5.1}{6.5.1}），存在 $y$ 的开邻域 $V$ 及态射 $g:V\to X$，
使 $g\circ f$ 和 $f\circ g$ 分别在 $x$ 和 $y$ 的某个邻域内有定义且等于恒等态射。
由此容易得到 $f$ 为局部同构。

(i) 问题在 $X,Y$ 上局部，可假设它们仿射，环分别为 $A,B$。
$f=({}^a\varphi,\wt\varphi)$，其中 $\varphi:B\to A$ 使 $A$ 成为有限型 $B$-代数；
$\varphi^{-1}(\jj_x)=\jj_y$，且诱导的同态 $\varphi_x:B_y\to A_x$ 为满射。
设 $t_i$（$1\leq i\leq n$）为 $B$-代数 $A$ 的生成元。
由满射性，存在 $b_i\in B$ 及 $c\in B\setminus\jj_y$，
使 $A_x$ 中有 $t_i/1=\varphi(b_i)/\varphi(c)$。
于是存在 $a\in A\setminus\jj_x$，令 $g=a\varphi(c)$，
则在 $A_g$ 中也有 $t_i/1=a\varphi(b_i)/g$（\hyperlink{para:1:1.3.3}{1.3.3}）。

由假设存在系数在 $\varphi(B)$ 中的多项式 $Q$ 使 $a=Q(t_1,\ldots,t_n)$。
写
\[
 Q(X_1/T,\ldots,X_n/T)=P(X_1,\ldots,X_n,T)/T^m,
\]
其中 $P$ 为 $m$ 次齐次多项式。
在 $A_g$ 中有
\[
 a/1=a^mP(\varphi(b_1),\ldots,\varphi(b_n),\varphi(c))/g^m
     =a^m\varphi(d)/g^m
\]
其中 $d\in B$。
由定义 $g/1=(a/1)(\varphi(c)/1)$ 可逆，所以 $a/1$、$\varphi(c)/1$ 均可逆；
从而可写 $a/1=(\varphi(d)/1)(\varphi(c)/1)^{-m}$。
因此 $\varphi(d)/1$ 在 $A_g$ 中也可逆。
置 $h=cd$，则 $\varphi(h)/1$ 在 $A_g$ 中可逆，
所以 $B\xrightarrow{\varphi}A\to A_g$ 分解为
$B\to B_h\xrightarrow{\gamma}A_g$（0，\hyperlink{para:0:1.2.4}{1.2.4}）。

证明 $\gamma$ 为满射，只须证明其像包含各 $t_i/1$ 及 $(g/1)^{-1}$。
事实上，
\[
 (g/1)^{-1}=(\varphi(c)/1)^{m-1}(\varphi(d)/1)^{-1}
          =\gamma(c^m/h),
\]
且 $a/1=\gamma(d^{m+1}/h^m)$，所以
$(a\varphi(b_i))/1=\gamma(b_i d^{m+1}/h^m)$。
又 $t_i/1=(a\varphi(b_i)/1)(g/1)^{-1}$，故满射性得证。
$h$ 的选择使 $\psi(D(g))\subset D(h)$，
且 $f|D(g)=({}^a\gamma,\wt\gamma)$。
因为 $\gamma$ 为满射，该限制是 $D(g)$ 到 $D(h)$ 的闭浸入（\hyperlink{para:1:4.2.3}{4.2.3}）。

\egastatement{6.5.5}{推论}
设 $f=(\psi,\theta):X\to Y$ 为有限型态射，$X$ 不可约，$x$ 为其泛点，$y=\psi(x)$。
\begin{enumerate}[label=(\roman*)]
\item $f$ 在 $X$ 的某一点处为局部浸入的充要条件是
$\theta_x^\sharp:\OO_y\to\OO_x$ 为满射。
\item 再设 $Y$ 不可约且局部诺特。
$f$ 在 $X$ 的某一点处为局部同构的充要条件是：
$y$ 为 $Y$ 的泛点，等价地 $f$ 支配，且 $\theta_x^\sharp$ 为同构；
亦即 $f$ 为双有理态射（\hyperlink{para:1:2.2.9}{2.2.9}）。
\end{enumerate}

由于 $X$ 的每个非空开集都包含 $x$，(i) 显然由（\hyperlink{para:1:6.5.4}{6.5.4}，(i)）得到；
(ii) 同样由（\hyperlink{para:1:6.5.4}{6.5.4}，(ii)）得到。


% END chapters/ch1_s6_5.tex

\par

% BEGIN chapters/ch1_s6_6.tex
\subsection{拟紧态射与局部有限型态射}\hypertarget{heading:1:6.6}{}
\src{152—155}

\egastatement{6.6.1}{定义}
若 $Y$ 的每个拟紧开集在态射 $f:X\to Y$ 下的逆像都是拟紧的，
则称 $f$ 为\emph{拟紧态射}。

设 $\mathfrak B$ 为由拟紧开集组成的 $Y$ 的拓扑基，例如仿射开集基。
$f$ 拟紧的充要条件是：$\mathfrak B$ 中每个成员在 $f$ 下的逆像拟紧，
等价地是有限个仿射开集的并；因为 $Y$ 的每个拟紧开集都是有限个基元素的并。
例如，若 $X$ 拟紧而 $Y$ 仿射，则任意态射 $f:X\to Y$ 均拟紧。
事实上，$X$ 是有限个仿射开集 $U_i$ 的并，
且对 $Y$ 的每个仿射开集 $V$，$U_i\cap f^{-1}(V)$ 仿射（\hyperlink{para:1:5.5.10}{5.5.10}），从而拟紧。

若 $f:X\to Y$ 拟紧，则对每个开集 $V\subset Y$，
限制 $f^{-1}(V)\to V$ 显然拟紧。
反之，若 $(U_\alpha)$ 为 $Y$ 的开覆盖，且各限制
$f^{-1}(U_\alpha)\to U_\alpha$ 拟紧，则 $f$ 拟紧。

\egastatement{6.6.2}{定义}
若对每个 $x\in X$，存在其开邻域 $U$ 及 $y=f(x)$ 的开邻域 $V\supset f(U)$，
使 $f|U:U\to V$ 为有限型态射，
则称态射 $f:X\to Y$ 为\emph{局部有限型态射}。
此时也称 $X$ 为 $Y$ 上的\emph{局部有限型预概形}，或\emph{局部有限型 $Y$-预概形}。

由（\hyperlink{para:1:6.3.2}{6.3.2}）立即可知：若 $f$ 局部有限型，则对每个开集 $W\subset Y$，
限制 $f^{-1}(W)\to W$ 局部有限型。
若 $Y$ 局部诺特，且 $X$ 在 $Y$ 上局部有限型，则 $X$ 局部诺特（\hyperlink{para:1:6.3.7}{6.3.7}）。

\egastatement{6.6.3}{命题}
态射 $f:X\to Y$ 为有限型的充要条件是：它拟紧且局部有限型。

必要性由（\hyperlink{para:1:6.3.1}{6.3.1}）及（\hyperlink{para:1:6.6.1}{6.6.1}）之后的说明立即得到。
反之，设这些条件成立，$U$ 为 $Y$ 的仿射开集，环为 $A$。
对每个 $x\in f^{-1}(U)$，由假设存在开邻域 $V(x)\subset f^{-1}(U)$
及 $y=f(x)$ 的开邻域 $W(x)\subset U$，
使 $f(V(x))\subset W(x)$ 且 $V(x)\to W(x)$ 为有限型态射。
用形如 $D(g)$、$g\in A$ 的较小邻域 $W_1(x)\subset W(x)$ 代替 $W(x)$，
并以 $V(x)\cap f^{-1}(W_1(x))$ 代替 $V(x)$，
可假设 $W(x)=D(g)$。
此时 $W(x)$ 的环为 $A[1/g]$，所以它在 $U$ 上有限型，
从而 $V(x)$ 在 $U$ 上有限型。
又因 $f^{-1}(U)$ 拟紧，它是有限个 $V(x_i)$ 的并，证明完成。

\egastatement{6.6.4}{命题}
\begin{enumerate}[label=(\roman*)]
\item 浸入 $X\to Y$ 在下列任一种情形下拟紧：它是闭浸入；
$Y$ 的底空间局部诺特；$X$ 的底空间诺特。
\item 两个拟紧态射的复合拟紧。
\item 若 $f:X\to Y$ 为拟紧 $S$-态射，则任意基扩张 $g:S'\to S$
所得的 $f_{(S')}:X_{(S')}\to Y_{(S')}$ 拟紧。
\item 两个拟紧 $S$-态射 $f:X\to X'$、$g:Y\to Y'$ 的积 $f\times_S g$ 拟紧。
\item 若 $f:X\to Y$、$g:Y\to Z$ 的复合拟紧，
且 $g$ 分离，或 $X$ 的底空间局部诺特，则 $f$ 拟紧。
\item $f$ 拟紧的充要条件是 $f_{\mathrm{red}}$ 拟紧。
\end{enumerate}

(vi) 显然，因为拟紧性只依赖于底空间的连续映射。
先证明 (v) 中 $X$ 的底空间局部诺特的情形。
置 $h=g\circ f$，设 $U$ 为 $Y$ 的拟紧开集。
$g(U)$ 是 $Z$ 中拟紧但未必开的子空间，故包含于有限个拟紧开集 $V_j$ 的并（\hyperlink{para:1:2.1.3}{2.1.3}）。
所以 $f^{-1}(U)$ 包含于各 $h^{-1}(V_j)$ 的并。
这些是 $X$ 中的拟紧子空间，从而是诺特子空间。
由（0，\hyperlink{para:0:2.2.3}{2.2.3}），$f^{-1}(U)$ 为诺特空间，当然拟紧。

其余断言只须证明 (i)、(ii)、(iii)（\hyperlink{para:1:5.5.12}{5.5.12}）。
(ii) 显然；(i) 在 $Y$ 的底空间局部诺特或 $X$ 的底空间诺特时由（\hyperlink{para:1:6.3.5}{6.3.5}）得到，
对闭浸入则显然。
为证明 (iii)，可限于 $S=Y$（\hyperlink{para:1:3.3.11}{3.3.11}）。
置 $f'=f_{(S')}$，设 $U'$ 为 $S'$ 的拟紧开集。
对每个 $s'\in U'$，取 $g(s')$ 在 $S$ 中的仿射开邻域 $T$，
及 $s'$ 的仿射开邻域 $W\subset U'\cap g^{-1}(T)$。
只须证明 $f'^{-1}(W)$ 拟紧；即归结到 $S,S'$ 都仿射时证明 $X\times_S S'$ 的底空间拟紧。
此时由假设 $X$ 为有限个仿射开集 $V_j$ 的并，
故 $X\times_S S'$ 为仿射概形 $V_j\times_S S'$ 的底空间之并（\hyperlink{para:1:3.2.2}{3.2.2} 和 \hyperlink{para:1:3.2.7}{3.2.7}）。命题得证。

另外，若 $X=X'\amalg X''$ 是两个预概形的和，
则态射 $f:X\to Y$ 拟紧的充要条件是它在 $X',X''$ 上的限制都拟紧。

\egastatement{6.6.5}{命题}
设 $f:X\to Y$ 为拟紧态射。
$f$ 支配的充要条件是：对 $Y$ 的每个不可约分支的泛点 $y$，
$f^{-1}(y)$ 包含 $X$ 的某个不可约分支的泛点。

充分性立即成立，且不必假设 $f$ 拟紧。
为证明必要性，取 $y$ 的仿射开邻域 $U$。
$f^{-1}(U)$ 拟紧，故为有限个仿射开集 $V_i$ 的并。
$f$ 支配的假设蕴含 $y$ 属于某个 $f(V_i)$ 在 $U$ 中的闭包。
显然可假设 $X,Y$ 约化。
由于 $V_i$ 的不可约分支在 $X$ 中的闭包是 $X$ 的不可约分支（0，\hyperlink{para:0:2.1.6}{2.1.6}），
可用 $V_i$ 代替 $X$，用以 $\overline{f(V_i)}\cap U$ 为底空间的 $U$ 的约化闭子预概形代替 $Y$（\hyperlink{para:1:5.2.1}{5.2.1}）。
于是归结到 $X=\Spec(A)$、$Y=\Spec(B)$ 均仿射且约化的情形。
$f$ 支配，所以 $B$ 为 $A$ 的子环（\hyperlink{para:1:1.2.7}{1.2.7}）；
再由 $B$ 的每个极小素理想都是 $B$ 与 $A$ 的某个极小素理想的交（0，\hyperlink{para:0:1.5.8}{1.5.8}），得到命题。

\egastatement{6.6.6}{命题}
\begin{enumerate}[label=(\roman*)]
\item 每个局部浸入都局部有限型。
\item 两个局部有限型态射的复合局部有限型。
\item 局部有限型 $S$-态射沿任意 $S'\to S$ 作基扩张后仍局部有限型。
\item 两个局部有限型 $S$-态射的积局部有限型。
\item 若复合态射 $g\circ f$ 局部有限型，则 $f$ 局部有限型。
\item 若 $f$ 局部有限型，则 $f_{\mathrm{red}}$ 局部有限型。
\end{enumerate}

由（\hyperlink{para:1:5.5.12}{5.5.12}），只须证明 (i)、(ii)、(iii)。
若 $j:X\to Y$ 为局部浸入，则对每个 $x\in X$，
存在 $j(x)$ 的开邻域 $V$ 和 $x$ 的开邻域 $U$，
使 $j|U:U\to V$ 为闭浸入（\hyperlink{para:1:4.5.1}{4.5.1}），所以为有限型态射。

为证明 (ii)，设 $f:X\to Y$、$g:Y\to Z$ 局部有限型，取 $x\in X$。
由假设存在 $g(f(x))$ 的开邻域 $W$ 和 $f(x)$ 的开邻域 $V$，
使 $g(V)\subset W$ 且 $V$ 在 $W$ 上有限型。
$f^{-1}(V)$ 在 $V$ 上局部有限型（\hyperlink{para:1:6.6.2}{6.6.2}），
因而可取 $x$ 的开邻域 $U\subset f^{-1}(V)$，使它在 $V$ 上有限型。
于是 $g(f(U))\subset W$，且 $U$ 在 $W$ 上有限型（\hyperlink{para:1:6.3.4}{6.3.4}，(ii)）。

最后证明 (iii)，可限于 $Y=S$（\hyperlink{para:1:3.3.11}{3.3.11}）。
对每个 $x'\in X'=X_{(S')}$，令 $x$ 为其在 $X$ 中的像，$s$ 为 $x$ 在 $S$ 中的像。
取 $s$ 的开邻域 $T$，以 $T'$ 表示它在 $S'$ 中的逆像；
取 $x$ 的开邻域 $U$，使它的像包含于 $T$ 且在 $T$ 上有限型。
则 $U\times_S T'=U\times_T T'$ 是 $x'$ 的开邻域（\hyperlink{para:1:3.2.7}{3.2.7}），
且在 $T'$ 上有限型（\hyperlink{para:1:6.3.4}{6.3.4}，(iv)）。

\egastatement{6.6.7}{推论}
设 $X,Y$ 为两个在 $S$ 上局部有限型的预概形。
若 $S$ 局部诺特，则 $X\times_S Y$ 局部诺特。

$X$ 在 $S$ 上局部有限型，故局部诺特。
$X\times_S Y$ 在 $X$ 上局部有限型，所以也局部诺特。

\egastatement{6.6.8}{注}
命题（\hyperlink{para:1:6.3.10}{6.3.10}）及其证明立即推广到仅假设 $f$ 局部有限型的情形。
同样，命题（\hyperlink{para:1:6.4.2}{6.4.2}）、（\hyperlink{para:1:6.4.9}{6.4.9}）在其中的 $X,Y$ 仅假设为域 $K$ 上局部有限型预概形时仍成立。


% END chapters/ch1_s6_6.tex

\par

% BEGIN chapters/ch1_s7_1.tex
\section{有理映射}\hypertarget{heading:1:7}{}
\subsection{有理映射与有理函数}\hypertarget{heading:1:7.1}{}
\src{155—158}
\egapar{7.1.1} 设 $X,Y$ 为两个预概形，$U,V$ 为 $X$ 中的两个稠密开集，$f$（分别为 $g$）为 $U$（分别为 $V$）到 $Y$ 的态射。若 $f,g$ 在 $U\cap V$ 的一个稠密开集中重合，就称它们\emph{等价}。由于 $X$ 中有限个稠密开集之交仍是 $X$ 中的稠密开集，这显然是一个等价关系。
\egastatement{7.1.2}{定义} 给定两个预概形 $X,Y$，称从 $X$ 到 $Y$ 的\emph{有理映射}为：从 $X$ 的稠密开子集到 $Y$ 的态射，关于（\hyperlink{para:1:7.1.1}{7.1.1}）所定义关系的一个等价类。若 $X,Y$ 为 $S$-预概形，且该类中存在一个 $S$-态射，则称该类为一个 $S$-有理映射。称从 $S$ 到 $X$ 的任一 $S$-有理映射为 $X$ 的一个\emph{$S$-有理截面}。称 $X$-预概形 $X\otimes_{\mathbf Z}\mathbf Z[T]$（$T$ 为不定元）的任一 $X$-有理截面为预概形 $X$ 上的一个\emph{有理函数}。

当只讨论 $S$-预概形且不致混淆时，将滥用术语，以“有理映射”代替“$S$-有理映射”。

设 $f$ 为从 $X$ 到 $Y$ 的有理映射，$U$ 为 $X$ 的开集。若 $f_1,f_2$ 为属于类 $f$ 的态射，分别定义于 $X$ 的稠密开集 $V,W$，则限制 $f_1|(U\cap V)$ 与 $f_2|(U\cap W)$ 在 $U\cap V\cap W$ 的一个稠密开集中重合；该开集也在 $U$ 中稠密。\transnote{原文此处直接写“在 $U\cap V\cap W$ 中重合”。依（\hyperlink{para:1:7.1.1}{7.1.1}），一般只能断言在其中的一个稠密开集中重合；这里明确写出这一点，不预先使用（\hyperlink{para:1:7.2.2}{7.2.2}）所需的约化性和分离性假设。}因此类 $f$ 定义了从 $U$ 到 $Y$ 的一个有理映射，称为 $f$ 在 $U$ 上的\emph{限制}，记作 $f|U$。

将每个 $S$-态射 $f:X\to Y$ 映到它所属的 $S$-有理映射，就定义了从 $\Hom_S(X,Y)$ 到从 $X$ 到 $Y$ 的 $S$-有理映射集合的典范映射。将 $X$ 的 $Y$-有理截面集合记作 $\Gamma_{\mathrm{rat}}(X/Y)$，于是有典范映射
\[
\Gamma(X/Y)\longrightarrow\Gamma_{\mathrm{rat}}(X/Y).
\]
另外，若 $X,Y$ 为两个 $S$-预概形，从 $X$ 到 $Y$ 的 $S$-有理映射集合显然典范地等同于 $\Gamma_{\mathrm{rat}}((X\times_S Y)/X)$（\hyperlink{para:1:3.3.14}{3.3.14}）。
\egapar{7.1.3} 由（\hyperlink{para:1:7.1.2}{7.1.2}）与（\hyperlink{para:1:3.3.14}{3.3.14}）立即得知，$X$ 上的有理函数典范地等同于结构层 $\OO_X$ 在 $X$ 的稠密开集上的截面的等价类；两个这样的截面等价，是指它们在包含于各自定义集交集中的一个稠密开集上重合。特别地，$X$ 上的有理函数组成一个环 $R(X)$。
\egapar{7.1.4} 当 $X$ 为不可约预概形时，每个非空开集都在 $X$ 中稠密；也可以说，$X$ 的非空开集正是其泛点 $x$ 的开邻域。两个从 $X$ 的非空开子集到 $Y$ 的态射等价，此时就意味着它们在 $x$ 处有相同的芽。换言之，从 $X$ 到 $Y$ 的有理映射（分别为 $S$-有理映射），等同于从 $X$ 的非空开子集到 $Y$ 的态射（分别为 $S$-态射）在 $X$ 的泛点 $x$ 处的芽。特别地：
\egastatement{7.1.5}{命题} 若 $X$ 为不可约预概形，则 $X$ 上的有理函数环 $R(X)$ 典范地等同于 $X$ 的泛点 $x$ 处的局部环 $\OO_x$。这是一个零维局部环，因而在 $X$ 为诺特预概形时是阿廷局部环；在 $X$ 为整预概形时是域，而当 $X$ 还为仿射概形时，它等同于 $A(X)$ 的分式域。

由上述讨论，以及有理函数与 $\OO_X$ 在稠密开集上的截面的等同，第一项断言就是层在一点处的茎的定义。对其余断言，可归结到 $X$ 为环 $A$ 的仿射概形的情形。此时 $\jj_x$ 为 $A$ 的幂零根，因此 $\OO_x$ 为零维；若 $A$ 为整环，则 $\jj_x=(0)$，从而 $\OO_x$ 为 $A$ 的分式域。最后，若 $A$ 为诺特环，则已知 $\jj_x$ 是幂零的（\cite{b11}，第 127 页，推论 4），且 $\OO_x=A_{\jj_x}$ 为阿廷环。

若 $X$ 为整预概形，则对任意 $z\in X$，环 $\OO_z$ 都是整环。任一包含 $z$ 的仿射开集 $U$ 也包含 $x$，而 $R(U)$ 等于 $A(U)$ 的分式域，典范地等同于 $R(X)$；因此 $R(X)$ 也等同于 $\OO_z$ 的分式域。将 $\OO_z$ 典范地等同于 $R(X)$ 的子环，是将每个截面芽 $s\in\OO_z$ 映到唯一的有理函数：该函数为一个在 $z$ 处具有芽 $s$ 的 $\OO_X$-截面的类；此截面必定义于一个稠密开集上。
\egapar{7.1.6} 现在设 $X$ 只有有限个不可约分支 $X_i$（$1\leq i\leq n$）；当 $X$ 的底空间为诺特空间时即属此种情形。令
\[
X_i'=X_i\setminus\bigcup_{j\ne i}(X_i\cap X_j).
\]
这是 $X$ 的一个开集，是不可约的，其泛点 $x_i$ 也是 $X_i$ 的泛点；各 $X_i'$ 两两不交，其并在 $X$ 中稠密（0，\hyperlink{para:0:2.1.6}{2.1.6}）。对 $X$ 的任一稠密开集 $U$，$U_i=U\cap X_i'$ 为 $X_i'$ 的非空稠密开集，各 $U_i$ 两两不交，故 $U'=\bigcup_iU_i$ 在 $X$ 中稠密。给定从 $U'$ 到 $Y$ 的态射，等价于任意给定各 $U_i$ 到 $Y$ 的态射。因此：
\egastatement{7.1.7}{命题} 设 $X,Y$ 为两个预概形（分别为 $S$-预概形），且 $X$ 有有限个不可约分支 $X_i$，泛点为 $x_i$（$1\leq i\leq n$）。若 $R_i$ 为从 $X$ 的开子集到 $Y$ 的态射（分别为 $S$-态射）在 $x_i$ 处的芽的集合，则从 $X$ 到 $Y$ 的有理映射（分别为 $S$-有理映射）集合等同于 $\prod_{i=1}^nR_i$。
\egastatement{7.1.8}{推论} 设 $X$ 为诺特预概形。$X$ 上的有理函数环为阿廷环，其局部分量为 $X$ 的不可约分支的泛点 $x_i$ 处的局部环 $\OO_{x_i}$。
\egastatement{7.1.9}{推论} 设 $A$ 为诺特环，$X=\Spec(A)$。若 $Q$ 为 $A$ 的所有极小素理想之并的补集，则 $X$ 上的有理函数环典范地等同于分式环 $Q^{-1}A$。

这将由下述引理得出。
\egastatement{7.1.9.1}{引理} 元素 $f\in A$ 满足 $D(f)$ 在 $X$ 中稠密，当且仅当 $f\in Q$；$X$ 的每个稠密开集都包含一个形如 $D(f)$ 的开集，其中 $f\in Q$。

先假定此引理已经证明。由于截面环 $\Gamma(D(f),\OO_X)$ 等同于 $A_f$（\hyperlink{para:1:1.3.6}{1.3.6} 与 \hyperlink{para:1:1.3.7}{1.3.7}），且 $f\in Q$ 时的 $D(f)$ 在按 $\supset$ 排序的稠密开集集合中构成一个共尾子集，故由定义（\hyperlink{para:1:7.1.1}{7.1.1}），$X$ 的有理函数环等同于各 $A_f$（$f\in Q$）的归纳极限，其中采用“$g$ 是 $f$ 的倍数”这一预序关系；此极限即 $Q^{-1}A$（0，\hyperlink{para:0:1.4.5}{1.4.5}）。

为证明（\hyperlink{para:1:7.1.9.1}{7.1.9.1}），仍将 $X$ 的不可约分支记作 $X_i$（$1\leq i\leq n$）。$D(f)$ 在 $X$ 中稠密，当且仅当 $D(f)\cap X_i\ne\varnothing$ 对一切 $i$ 成立。令 $\mathfrak p_i=\jj(X_i)$，上述条件就是 $f\notin\mathfrak p_i$ 对所有 $i$ 成立；而各 $\mathfrak p_i$ 正是 $A$ 的极小素理想（\hyperlink{para:1:1.1.14}{1.1.14}），故这些条件等价于 $f\in Q$，得到引理的第一项断言。另一方面，若 $U$ 为 $X$ 的稠密开集，其补集形如 $V(\mathfrak a)$，其中理想 $\mathfrak a$ 不包含于任何一个 $\mathfrak p_i$；因而它也不包含于这些理想之并（\cite{b10}，第 13 页）。故可取 $f\in\mathfrak a\cap Q$，于是 $D(f)\subset U$。证明完毕。
\egapar{7.1.10} 再设 $X$ 不可约，泛点为 $x$。每个非空开集 $U$ 都包含 $x$，因而也包含所有满足 $x\in\overline{\{z\}}$ 的 $z\in X$，所以可将任一态射 $U\to Y$ 与典范态射 $\Spec(\OO_x)\to X$（\hyperlink{para:1:2.4.1}{2.4.1}）复合。两个分别从 $X$ 的非空开子集到 $Y$ 的态射，若在一个非空开子集中重合，经此复合即给出同一态射 $\Spec(\OO_x)\to Y$。换言之，每个从 $X$ 到 $Y$ 的有理映射因此对应一个确定的态射 $\Spec(\OO_x)\to Y$。
\egastatement{7.1.11}{命题} 设 $X,Y$ 为两个 $S$-预概形，$X$ 不可约，泛点为 $x$，且 $Y$ 在 $S$ 上为有限型。若两个从 $X$ 到 $Y$ 的 $S$-有理映射对应同一 $S$-态射 $\Spec(\OO_x)\to Y$，则这两个有理映射相同。若还假设 $S$ 局部诺特，则每个从 $\Spec(\OO_x)$ 到 $Y$ 的 $S$-态射都对应唯一一个从 $X$ 到 $Y$ 的 $S$-有理映射。

考虑到 $X$ 的每个非空开集都是稠密的，这立即由（\hyperlink{para:1:6.5.1}{6.5.1}）得出。
\egastatement{7.1.12}{推论} 假设 $S$ 局部诺特，并满足（\hyperlink{para:1:7.1.11}{7.1.11}）的其余假设。那么，从 $X$ 到 $Y$ 的 $S$-有理映射，等同于 $S$-预概形 $Y$ 的取值于 $S$-预概形 $\Spec(\OO_x)$ 的点。

这就是用（\hyperlink{para:1:3.4.1}{3.4.1}）的术语重述（\hyperlink{para:1:7.1.11}{7.1.11}）。
\egastatement{7.1.13}{推论} 假设满足（\hyperlink{para:1:7.1.12}{7.1.12}）的条件。令 $s$ 为 $x$ 在 $S$ 中的像。给定从 $X$ 到 $Y$ 的一个 $S$-有理映射，等价于给定 $Y$ 中位于 $s$ 之上的一点 $y$，以及一个局部 $\OO_s$-同态 $\OO_y\to\OO_x=R(X)$。

这由（\hyperlink{para:1:7.1.11}{7.1.11}）与（\hyperlink{para:1:2.4.4}{2.4.4}）得出。特别地：
\egastatement{7.1.14}{推论} 在（\hyperlink{para:1:7.1.12}{7.1.12}）的条件下，对给定的 $Y$，从 $X$ 到 $Y$ 的 $S$-有理映射只依赖于 $S$-预概形 $\Spec(\OO_x)$；特别地，对任一 $z\in X$，将 $X$ 换成 $\Spec(\OO_z)$ 后，它们仍保持不变。

事实上，因 $z\in\overline{\{x\}}$，$x$ 是 $Z=\Spec(\OO_z)$ 的泛点，且 $\OO_{X,x}=\OO_{Z,x}$。

当 $X$ 为整预概形时，$R(X)=\OO_x=k(x)$ 是域（\hyperlink{para:1:7.1.5}{7.1.5}），以上推论于是给出：
\egastatement{7.1.15}{推论} 假设满足（\hyperlink{para:1:7.1.12}{7.1.12}）的条件，且 $X$ 为整预概形。令 $s$ 为 $x$ 在 $S$ 中的像。那么，从 $X$ 到 $Y$ 的 $S$-有理映射等同于 $Y\otimes_S k(s)$ 的取值于扩域 $R(X)/k(s)$ 的几何点。换言之，每个这样的有理映射等价于给定 $Y$ 中位于 $s$ 之上的一点 $y$，以及一个 $k(s)$-单同态 $k(y)\to k(x)=R(X)$。

事实上，$Y$ 中位于 $s$ 之上的点等同于 $Y\otimes_S k(s)$ 的点（\hyperlink{para:1:3.6.3}{3.6.3}），而局部 $\OO_s$-同态 $\OO_y\to R(X)$ 等同于 $k(s)$-单同态 $k(y)\to R(X)$。更特别地：
\egastatement{7.1.16}{推论} 设 $k$ 为域，$X,Y$ 为两个 $k$ 上的代数预概形（\hyperlink{para:1:6.4.1}{6.4.1}），且 $X$ 为整预概形。那么，从 $X$ 到 $Y$ 的 $k$-有理映射等同于 $Y$ 的取值于扩域 $R(X)/k$ 的几何点（\hyperlink{para:1:3.4.4}{3.4.4}）。


% END chapters/ch1_s7_1.tex

\par

% BEGIN chapters/ch1_s7_2.tex
\subsection{有理映射的定义域}\hypertarget{heading:1:7.2}{}
\src{158—161}
\egapar{7.2.1} 设 $X,Y$ 为两个预概形，$f$ 为从 $X$ 到 $Y$ 的有理映射。若存在一个包含 $x$ 的稠密开集 $U$，以及属于等价类 $f$ 的态射 $U\to Y$，就称 $f$ \emph{在点 $x\in X$ 处有定义}。所有使 $f$ 有定义的点组成的集合称为 $f$ 的\emph{定义域}；这显然是 $X$ 中的一个稠密开集。
\egastatement{7.2.2}{命题} 设 $X,Y$ 为两个 $S$-预概形，其中 $X$ 约化，$Y$ 在 $S$ 上分离。设 $f$ 为从 $X$ 到 $Y$ 的一个 $S$-有理映射，$U_0$ 为其定义域。那么，存在唯一一个属于类 $f$ 的 $S$-态射 $U_0\to Y$。

由于对属于类 $f$ 的任一态射 $U\to Y$ 都必有 $U\subset U_0$，本命题显然由以下引理得出。
\egastatement{7.2.2.1}{引理} 在（\hyperlink{para:1:7.2.2}{7.2.2}）的假设下，设 $U_1,U_2$ 为 $X$ 中的两个稠密开集，$f_i:U_i\to Y$（$i=1,2$）为两个 $S$-态射。若存在 $V\subset U_1\cap U_2$，它在 $X$ 中开且稠密，并且 $f_1,f_2$ 在 $V$ 中重合，则它们在整个 $U_1\cap U_2$ 中重合。

显然可归结到 $X=U_1=U_2$ 的情形。由于 $X$（因而 $V$）约化，$X$ 是 $X$ 的不小于 $V$ 的最小闭子预概形（\hyperlink{para:1:5.2.2}{5.2.2}）。令 $g=(f_1,f_2)_S:X\to Y\times_S Y$。按假设，对角线 $T=\Delta_Y(Y)$ 为 $Y\times_S Y$ 的闭子预概形，故 $Z=g^{-1}(T)$ 为 $X$ 的闭子预概形（\hyperlink{para:1:4.4.1}{4.4.1}）。若 $h:V\to Y$ 为 $f_1,f_2$ 在 $V$ 上的共同限制，则 $g$ 在 $V$ 上的限制为 $g'=(h,h)_S=\Delta_Y\circ h$。由 $\Delta_Y^{-1}(T)=Y$，得 $g'^{-1}(T)=V$，故 $Z$ 在 $V$ 上诱导 $V$，因而不小于 $V$，于是 $Z=X$。由 $g^{-1}(T)=X$，可知 $g$ 分解为 $\Delta_Y\circ f$，其中 $f:X\to Y$ 为态射（\hyperlink{para:1:4.4.1}{4.4.1}）；按对角态射的定义，$f_1=f_2=f$。

（\hyperlink{para:1:7.2.2}{7.2.2}）所定义的态射 $U_0\to Y$，显然是类 $f$ 中唯一不能延拓到 $X$ 的严格包含 $U_0$ 的开子集上的态射。因此，在（\hyperlink{para:1:7.2.2}{7.2.2}）的假设下，可将从 $X$ 到 $Y$ 的有理映射，等同于从 $X$ 的稠密开子集到 $Y$ 的不能继续延拓至更大开集的态射。在此等同下，（\hyperlink{para:1:7.2.2}{7.2.2}）给出：
\egastatement{7.2.3}{推论} 对 $X,Y$ 作（\hyperlink{para:1:7.2.2}{7.2.2}）的假设，并设 $U$ 为 $X$ 的一个稠密开集。则从 $U$ 到 $Y$ 的 $S$-态射，与从 $X$ 到 $Y$ 的在 $U$ 的每一点都有定义的 $S$-有理映射之间，存在典范双射。

事实上，由（\hyperlink{para:1:7.2.2}{7.2.2}），每个 $S$-态射 $f:U\to Y$ 都有唯一一个延拓它的 $S$-有理映射 $\bar f:X\dashrightarrow Y$。
\egastatement{7.2.4}{推论} 设 $S$ 为概形，$X$ 为约化 $S$-预概形，$Y$ 为 $S$-概形，$f:U\to Y$ 为从 $X$ 的稠密开集 $U$ 到 $Y$ 的 $S$-态射。若 $\bar f$ 为延拓 $f$ 的从 $X$ 到 $Y$ 的 $\mathbf Z$-有理映射，则 $\bar f$ 是一个 $S$-态射（因此也是延拓 $f$ 的 $S$-有理映射）。

事实上，令 $\varphi:X\to S$、$\psi:Y\to S$ 为结构态射，$U_0$ 为 $\bar f$ 的定义域，$j:U_0\to X$ 为嵌入。只须证明 $\psi\circ\bar f=\varphi\circ j$；由于 $f$ 为 $S$-态射，这立即由（\hyperlink{para:1:7.2.2.1}{7.2.2.1}）得出。
\egastatement{7.2.5}{推论} 设 $X,Y$ 为两个 $S$-预概形，$X$ 约化，且 $X,Y$ 均在 $S$ 上分离。设 $p:Y\to X$ 为 $S$-态射，使 $Y$ 成为 $X$-预概形；设 $U$ 为 $X$ 的稠密开集，$f$ 为 $Y$ 的一个 $U$-截面。那么，延拓 $f$ 的从 $X$ 到 $Y$ 的有理映射 $\bar f$ 是 $Y$ 的一个 $X$-有理截面。

须证明 $p\circ\bar f$ 在 $\bar f$ 的定义域上为恒等态射；因 $X$ 在 $S$ 上分离，这仍由（\hyperlink{para:1:7.2.2.1}{7.2.2.1}）得出。
\egastatement{7.2.6}{推论} 设 $X$ 为约化预概形，$U$ 为 $X$ 的一个稠密开集。$\OO_X$ 在 $U$ 上的截面，与 $X$ 上在 $U$ 的每一点都有定义的有理函数 $f$ 之间，存在典范双射。

由（\hyperlink{para:1:7.2.3}{7.2.3}）、（\hyperlink{para:1:7.1.2}{7.1.2}）与（\hyperlink{para:1:7.1.3}{7.1.3}），只须注意到 $X$-预概形 $X\otimes_{\mathbf Z}\mathbf Z[T]$ 在 $X$ 上分离（\hyperlink{para:1:5.5.1}{5.5.1}，(iv)）。
\egastatement{7.2.7}{推论} 设 $Y$ 为约化预概形，$f:X\to Y$ 为分离态射，$U$ 为 $Y$ 的稠密开集，$g:U\to f^{-1}(U)$ 为 $f^{-1}(U)$ 的一个 $U$-截面。令 $Z$ 为 $X$ 的以 $\overline{g(U)}$ 为底空间的约化子预概形（\hyperlink{para:1:5.2.1}{5.2.1}）。则 $g$ 是 $X$ 的一个 $Y$-截面的限制，当且仅当 $f$ 在 $Z$ 上的限制为 $Z$ 到 $Y$ 的同构。换言之，延拓 $g$ 的从 $Y$ 到 $X$ 的有理映射处处有定义，当且仅当上述条件成立（\hyperlink{para:1:7.2.5}{7.2.5}）。

$f$ 在 $f^{-1}(U)$ 上的限制是分离态射（\hyperlink{para:1:5.5.1}{5.5.1}，(i)），故 $g$ 为闭浸入（\hyperlink{para:1:5.4.6}{5.4.6}）。于是 $g(U)=Z\cap f^{-1}(U)$，且 $Z$ 在其开集 $g(U)$ 上诱导的子预概形，与 $f^{-1}(U)$ 中和 $g$ 相伴的闭子预概形相同（\hyperlink{para:1:5.2.1}{5.2.1}）。陈述中的条件显然充分：若成立，令 $f_Z:Z\to Y$ 为 $f$ 的限制，$\bar g:Y\to Z$ 为其逆同构，则 $\bar g$ 延拓 $g$。反之，若 $g$ 为 $X$ 的一个 $Y$-截面 $h$ 在 $U$ 上的限制，则 $h$ 为闭浸入（\hyperlink{para:1:5.4.6}{5.4.6}），故 $h(Y)$ 为闭集；由于它包含于 $Z$，它就等于 $Z$。由（\hyperlink{para:1:5.2.1}{5.2.1}），$h$ 必为 $Y$ 到 $X$ 的闭子预概形 $Z$ 的同构。
\egapar{7.2.8} 设 $X,Y$ 为两个 $S$-预概形，其中 $X$ 约化，$Y$ 在 $S$ 上分离。设 $f$ 为从 $X$ 到 $Y$ 的 $S$-有理映射，$x\in X$。只要 $f$ 的定义域在 $\Spec(\OO_x)$ 上的迹稠密于 $\Spec(\OO_x)$，就可将 $f$ 与典范 $S$-态射 $\Spec(\OO_x)\to X$（\hyperlink{para:1:2.4.1}{2.4.1}）复合；这里将 $\Spec(\OO_x)$ 等同于所有满足 $x\in\overline{\{z\}}$ 的 $z\in X$ 的集合（\hyperlink{para:1:2.4.2}{2.4.2}）。以下情形满足此条件。

$1^\circ$ $X$ 不可约，因而是整预概形。此时 $X$ 的泛点 $\xi$ 也是 $\Spec(\OO_x)$ 的泛点；由于 $f$ 的定义域 $U$ 包含 $\xi$，$U\cap\Spec(\OO_x)$ 也包含 $\xi$，故在 $\Spec(\OO_x)$ 中稠密。

$2^\circ$ $X$ 局部诺特。此时所需断言由以下引理得出。
\egastatement{7.2.8.1}{引理} 设 $X$ 为底空间局部诺特的预概形，$x\in X$。$\Spec(\OO_x)$ 的不可约分支，正是 $X$ 的包含 $x$ 的不可约分支在 $\Spec(\OO_x)$ 上的迹。开集 $U\subset X$ 满足 $U\cap\Spec(\OO_x)$ 稠密于 $\Spec(\OO_x)$，当且仅当它与 $X$ 的每个包含 $x$ 的不可约分支相交；特别地，当 $U$ 在 $X$ 中稠密时，此条件成立。

第二项断言显然由第一项得出，故只须证明后者。$\Spec(\OO_x)$ 包含于每个包含 $x$ 的仿射开集 $U$，而 $U$ 的包含 $x$ 的不可约分支，是 $X$ 的包含 $x$ 的不可约分支在 $U$ 上的迹（0，\hyperlink{para:0:2.1.6}{2.1.6}）。所以可假设 $X$ 为环 $A$ 的仿射概形。$A_x$ 的素理想与 $A$ 中包含于 $\jj_x$ 的素理想一一对应（0，\hyperlink{para:0:1.2.6}{1.2.6}）；因此 $A_x$ 的极小素理想与 $A$ 中包含于 $\jj_x$ 的极小素理想一一对应，从而得到引理（\hyperlink{para:1:1.1.14}{1.1.14}）。

现假设处于以上两种情形之一。若 $U$ 为 $S$-有理映射 $f$ 的定义域，记 $f'$ 为从 $\Spec(\OO_x)$ 到 $Y$ 的、在 $U\cap\Spec(\OO_x)$ 中与 $f$ 重合的有理映射（利用（\hyperlink{para:1:2.4.2}{2.4.2}）的等同）。称 $f'$ 为由 $f$ \emph{诱导}的有理映射。
\egastatement{7.2.9}{命题} 设 $S$ 为局部诺特预概形，$X$ 为约化 $S$-预概形，$Y$ 为有限型 $S$-概形。再假设 $X$ 不可约或局部诺特。设 $f$ 为从 $X$ 到 $Y$ 的 $S$-有理映射，$x\in X$。则 $f$ 在 $x$ 处有定义，当且仅当由 $f$ 诱导的有理映射 $f':\Spec(\OO_x)\dashrightarrow Y$（\hyperlink{para:1:7.2.8}{7.2.8}）是一个态射。

必要性显然，因为 $\Spec(\OO_x)$ 包含于每个包含 $x$ 的开集。证明充分性。由（\hyperlink{para:1:6.5.1}{6.5.1}），存在 $x$ 在 $X$ 中的一个开邻域 $U$，以及一个 $S$-态射 $g:U\to Y$，在 $\Spec(\OO_x)$ 上诱导 $f'$。若 $X$ 不可约，$U$ 在 $X$ 中稠密，故由（\hyperlink{para:1:7.2.3}{7.2.3}），可将 $g$ 看作 $S$-有理映射。此外，$X$ 的泛点属于 $\Spec(\OO_x)$ 及 $f$ 的定义域，所以 $f,g$ 在此点处重合，因而在 $X$ 的一个非空开集中重合（\hyperlink{para:1:6.5.1}{6.5.1}）。由于 $f,g$ 为 $S$-有理映射，它们相同（\hyperlink{para:1:7.2.3}{7.2.3}），故 $f$ 在 $x$ 处有定义。

若 $X$ 局部诺特，可假设 $U$ 诺特。此时 $X$ 的包含 $x$ 的不可约分支 $X_i$ 只有有限个（\hyperlink{para:1:7.2.8.1}{7.2.8.1}），且必要时缩小 $U$，可使它们成为仅有的与 $U$ 相交的不可约分支；这是因为与诺特开集 $U$ 相交的不可约分支只有有限个。如上可知，$f,g$ 在每个 $X_i$ 的一个非空开集中重合。考虑到各 $X_i$ 包含于 $\overline U$，取一个定义于 $U\cup(X\setminus\overline U)$ 的稠密开子集上的态射 $f_1$：它在 $U$ 中等于 $g$，在 $X\setminus\overline U$ 与 $f$ 的定义域之交上等于 $f$。由于 $U\cup(X\setminus\overline U)$ 在 $X$ 中稠密，$f_1,f$ 在 $X$ 的一个稠密开集中重合；又因 $f$ 是有理映射，$f$ 就是 $f_1$ 的延拓（\hyperlink{para:1:7.2.3}{7.2.3}），故在 $x$ 处有定义。


% END chapters/ch1_s7_2.tex

\par

% BEGIN chapters/ch1_s7_3.tex
\subsection{有理函数层}\hypertarget{heading:1:7.3}{}
\src{161—163}
\egapar{7.3.1} 设 $X$ 为预概形。对每个开集 $U\subset X$，记 $R(U)$ 为 $U$ 上的有理函数环（\hyperlink{para:1:7.1.3}{7.1.3}）；这是一个 $\Gamma(U,\OO_X)$-代数。若 $V\subset U$ 为另一开集，则 $\OO_X$ 在 $U$ 的稠密开子集上的任一截面，限制到 $V$ 后，给出 $V$ 的一个稠密开子集上的截面；若两个截面在 $U$ 的一个稠密开子集上重合，它们在 $V$ 上的限制也在 $V$ 的一个稠密开子集上重合。于是定义了代数的双同态 $\rho_{V,U}:R(U)\to R(V)$。显然，对三个开集 $U\supset V\supset W$，有 $\rho_{W,U}=\rho_{W,V}\circ\rho_{V,U}$；因此这些 $R(U)$ 定义了 $X$ 上的一个代数预层。
\egastatement{7.3.2}{定义} 称由这些 $R(U)$ 构成的预层所伴随的 $\OO_X$-代数层为预概形 $X$ 上的\emph{有理函数层}，记作 $\mathscr R(X)$。

对任意预概形 $X$ 及任意开集 $U\subset X$，显然诱导层 $\mathscr R(X)|U$ 就是 $\mathscr R(U)$。
\egastatement{7.3.3}{命题} 设预概形 $X$ 的不可约分支族 $(X_\lambda)$ 局部有限；特别地，当 $X$ 的底空间局部诺特时即如此。那么，$\OO_X$-模 $\mathscr R(X)$ 为拟凝聚层。对 $X$ 的每个只与有限个分支 $X_\lambda$ 相交的开集 $U$，$R(U)=\Gamma(U,\mathscr R(X))$，并典范地等同于那些满足 $U\cap X_\lambda\ne\varnothing$ 的 $X_\lambda$ 的泛点处的局部环的直积。

显然可限于 $X$ 只有有限个不可约分支 $X_i$，其泛点为 $x_i$（$1\leq i\leq n$）的情形。由（\hyperlink{para:1:7.1.7}{7.1.7}），$R(U)$ 典范地等同于满足 $U\cap X_i\ne\varnothing$ 的各 $\OO_{x_i}$ 的直积；原文将这些局部环记为 $R(X_i)$。\transnote{这里的 $R(X_i)$ 旨在表示对应泛点处的有理函数环 $\OO_{X,x_i}$。在尚未给不可约分支指定子预概形结构时，可用（\hyperlink{para:1:7.1.6}{7.1.6}）的不可约开集 $X_i'$ 理解这一记号；不能在非约化情形中擅自换成约化闭分支的函数域。}
再证明预层 $U\mapsto R(U)$ 满足层的公理，于是就有 $R(U)=\Gamma(U,\mathscr R(X))$。由前述表示，它满足（F1）。为验证（F2），设 $(V_\alpha)$ 为开集 $U$ 的开覆盖，$V_\alpha\subset U$，且 $s_\alpha\in R(V_\alpha)$ 对每对指标都在 $V_\alpha\cap V_\beta$ 上有相同限制。于是，对满足 $U\cap X_i\ne\varnothing$ 的每个 $i$，所有满足 $V_\alpha\cap X_i\ne\varnothing$ 的 $s_\alpha$ 在 $R(X_i)$ 中的分量都相同。将此分量记作 $t_i$，则以各 $t_i$ 为分量的 $R(U)$ 中的元素，显然在每个 $V_\alpha$ 上限制为 $s_\alpha$。

最后，为证明 $\mathscr R(X)$ 拟凝聚，可限于 $X=\Spec(A)$ 为仿射概形的情形。取形如 $D(f)$（$f\in A$）的仿射开集作为 $U$，由前面的讨论及定义（\hyperlink{para:1:1.3.4}{1.3.4}）得 $\mathscr R(X)=\wt M$，其中 $M$ 为各 $A$-模 $A_{x_i}$ 的直和。
\egastatement{7.3.4}{推论} 设 $X$ 为只有有限个不可约分支的约化预概形，并令 $X_i$（$1\leq i\leq n$）为以这些不可约分支为底空间的约化闭子预概形（\hyperlink{para:1:5.2.1}{5.2.1}）。若 $h_i:X_i\to X$ 为典范嵌入，则 $\mathscr R(X)$ 为各 $\OO_X$-代数 $(h_i)_*(\mathscr R(X_i))$ 的直积。
\egastatement{7.3.5}{推论} 若 $X$ 不可约，则每个拟凝聚 $\mathscr R(X)$-模 $\F$ 都是常层。

只须证明每个 $x\in X$ 都有一个邻域 $U$，使 $\F|U$ 为常层（0，\hyperlink{para:0:3.6.2}{3.6.2}），故可归结到 $X$ 仿射的情形。还可假设 $\F$ 为某个同态
\[
(\mathscr R(X))^{(I)}\longrightarrow(\mathscr R(X))^{(J)}
\]
的余核（0，\hyperlink{para:0:5.1.3}{5.1.3}），因而只须验证 $\mathscr R(X)$ 为常层。而这显然成立，因为每个非空开集 $U$ 都包含 $X$ 的泛点，故 $\Gamma(U,\mathscr R(X))=R(X)$。
\egastatement{7.3.6}{推论} 若 $X$ 不可约，则对每个拟凝聚 $\OO_X$-模 $\F$，$\F\otimes_{\OO_X}\mathscr R(X)$ 为常层。若 $X$ 还约化，因而是整预概形，则 $\F\otimes_{\OO_X}\mathscr R(X)$ 同构于一个形如 $(\mathscr R(X))^{(I)}$ 的层。

第二项断言来自此时 $R(X)$ 为域这一事实。
\egastatement{7.3.7}{命题} 设预概形 $X$ 局部整或局部诺特。那么 $\mathscr R(X)$ 是拟凝聚 $\OO_X$-代数。若 $X$ 还约化（当 $X$ 局部整时即如此），则典范同态 $\OO_X\to\mathscr R(X)$ 为单射。

问题是局部的，故第一项断言由（\hyperlink{para:1:7.3.3}{7.3.3}）得出；第二项断言立即由（\hyperlink{para:1:7.2.3}{7.2.3}）得出。
\egapar{7.3.8} 设 $X,Y$ 为各有有限个不可约分支的两个预概形，$f:X\to Y$ 为态射，其在 $X$ 的不可约分支的泛点集合上的限制，是到 $Y$ 的不可约分支的泛点集合的满射。原文断言，此时有
\begin{equation}
 f^*(\mathscr R(Y))=\mathscr R(X).
 \tag{7.3.8.1}\label{eq:1:7.3.8.1}
\end{equation}
\transnote{此式及下述论证按原文保留，但在所列假设下，该式一般不成立。譬如 $X=\Spec(k[t])$、$Y=\Spec(k)$，结构态射满足所述泛点条件，却有 $f^*\mathscr R(Y)=\OO_X$，而 $\mathscr R(X)$ 是函数域 $k(t)$ 所给出的常层。仅在泛点处计算茎，不能推出全层同构；（\hyperlink{para:1:7.3.5}{7.3.5}）也不能直接用于这里的任意 $\OO_X$-模。}
事实上，由（\hyperlink{para:1:7.3.3}{7.3.3}），原文将问题归结到 $X,Y$ 不可约、泛点分别为 $x,y$ 且 $f(x)=y$ 的情形。于是
\[
 (f^*(\mathscr R(Y)))_x
 =\OO_y\otimes_{\OO_y}\OO_x=\OO_x
 \qquad\text{（0，\hyperlink{para:0:4.3.1}{4.3.1}）},
\]
并据（\hyperlink{para:1:7.3.5}{7.3.5}）得出（\hyperref[eq:1:7.3.8.1]{7.3.8.1}）。


% END chapters/ch1_s7_3.tex

\par

% BEGIN chapters/ch1_s7_4.tex
\subsection{挠层与无挠层}\hypertarget{heading:1:7.4}{}
\src{163—164}
\egapar{7.4.1} 设 $X$ 为整预概形。对每个 $\OO_X$-模 $\F$，典范同态 $\OO_X\to\mathscr R(X)$ 经张量积定义一个同态（仍称典范同态）
\[
 \F\longrightarrow\F\otimes_{\OO_X}\mathscr R(X).
\]
在每个茎上，它就是 $\F_x\to\F_x\otimes_{\OO_x}R(X)$ 的同态 $z\mapsto z\otimes1$。此同态的核 $\mathscr T$ 是 $\F$ 的子 $\OO_X$-模，称为 $\F$ 的\emph{挠层}；若 $\F$ 拟凝聚，则 $\mathscr T$ 也拟凝聚（\hyperlink{para:1:4.1.1}{4.1.1} 与 \hyperlink{para:1:7.3.6}{7.3.6}）。若 $\mathscr T=0$，称 $\F$ \emph{无挠}；若 $\F=\mathscr T$，称 $\F$ 是\emph{挠层}。对任意 $\OO_X$-模 $\F$，$\F/\mathscr T$ 都无挠。由（\hyperlink{para:1:7.3.5}{7.3.5}）得：
\egastatement{7.4.2}{命题} 若 $X$ 为整预概形，则每个无挠拟凝聚 $\OO_X$-模 $\F$，都同构于一个形如 $(\mathscr R(X))^{(I)}$ 的常层的子层 $\G$，且该常层由 $\G$ 作为 $\mathscr R(X)$-模生成。

$I$ 的基数称为 $\F$ 的\emph{秩}。对 $X$ 的每个非空仿射开集 $U$，$\F$ 的秩等于 $\Gamma(U,\F)$ 作为 $\Gamma(U,\OO_X)$-模的秩；考察包含于 $U$ 的 $X$ 的泛点，即立即得到这一点。特别地：
\egastatement{7.4.3}{推论} 在整预概形 $X$ 上，每个无挠、秩为 $1$ 的拟凝聚 $\OO_X$-模，特别是每个可逆 $\OO_X$-模，都同构于 $\mathscr R(X)$ 的一个子 $\OO_X$-模；反之亦然。\transnote{逆向陈述应就非零拟凝聚子 $\OO_X$-模理解；零子模秩为 $0$，而任意子层也未必拟凝聚。原文逆向措辞未重复这些限制。}
\egastatement{7.4.4}{推论} 设 $X$ 为整预概形，$\LL,\LL'$ 为两个无挠 $\OO_X$-模，$f$（分别为 $f'$）为 $\LL$（分别为 $\LL'$）在 $X$ 上的截面。则 $f\otimes f'=0$，当且仅当 $f,f'$ 中至少一个为零。\transnote{对于所列的任意无挠模层，结论不成立。
取 $X=\Spec(k[t])$、$K=k(t)$，在闭点 $(t)$、$(t-1)$ 分别取从该单点直像得到的层
$\LL=i_{0*}K$、$\LL'=i_{1*}K$，其 $\OO_X$-作用由 $\OO_X\to\mathscr R(X)$ 诱导。
它们的非零茎均为 $K$，按（\hyperlink{para:1:7.4.1}{7.4.1}）都无挠；各自的全局截面 $1$ 非零。
但两者的支集不交，故 $\LL\otimes_{\OO_X}\LL'=0$。
这两个层都不是拟凝聚层。若另外假设两层拟凝聚，则可由（\hyperlink{para:1:7.4.2}{7.4.2}）嵌入函数域上的向量空间所给出的常层，
从而证明结论；这一证明不要求泛点处的茎只有一维。
正文保留原版陈述与证明；以上反例和附加条件分析为译者说明。}

令 $x$ 为 $X$ 的泛点。按假设有 $(f\otimes f')_x=f_x\otimes f'_x=0$。原文指出，$\LL_x,\LL'_x$ 等同于域 $\OO_x$ 的子 $\OO_x$-模，因此该关系推出 $f_x=0$ 或 $f'_x=0$；再由 $\LL,\LL'$ 无挠，得 $f=0$ 或 $f'=0$（\hyperlink{para:1:7.3.5}{7.3.5}）。
\egastatement{7.4.5}{命题} 设 $X,Y$ 为两个整预概形，$f:X\to Y$ 为支配态射。对每个无挠拟凝聚 $\OO_X$-模 $\F$，$f_*(\F)$ 都是无挠 $\OO_Y$-模。

由于 $f_*$ 左正合（0，\hyperlink{para:0:4.2.1}{4.2.1}），由（\hyperlink{para:1:7.4.2}{7.4.2}），只须在 $\F=(\mathscr R(X))^{(I)}$ 时证明命题。$Y$ 的每个非空开集 $U$ 都包含 $Y$ 的泛点，故 $f^{-1}(U)$ 包含 $X$ 的泛点（0，\hyperlink{para:0:2.1.5}{2.1.5}）。因此，此时有
\[
 \Gamma(U,f_*(\F))=\Gamma(f^{-1}(U),\F)=(R(X))^{(I)}.
\]
换言之，$f_*(\F)$ 是茎为 $(R(X))^{(I)}$ 的常层，将其视为 $\mathscr R(Y)$-模，显然无挠。
\egastatement{7.4.6}{命题} 设 $X$ 为整预概形，$x$ 为其泛点。对每个有限型拟凝聚 $\OO_X$-模 $\F$，下列条件等价：
\begin{enumerate}[label=\alph*)]
\item $\F$ 为挠层；
\item $\F_x=0$；
\item $\Supp(\F)\ne X$。
\end{enumerate}
由（\hyperlink{para:1:7.3.5}{7.3.5}）与（\hyperlink{para:1:7.4.1}{7.4.1}），$\F_x=0$ 和 $\F\otimes_{\OO_X}\mathscr R(X)=0$ 等价，故 a) 与 b) 等价。另一方面，$\Supp(\F)$ 在 $X$ 中闭（0，\hyperlink{para:0:5.2.2}{5.2.2}），且 $X$ 的每个非空开集都包含 $x$，故 b) 与 c) 等价。
\egapar{7.4.7} 将（\hyperlink{para:1:7.4.1}{7.4.1}）的定义以滥用术语的方式推广到 $X$ 为只有有限个不可约分支的约化预概形的情形。原文据（\hyperlink{para:1:7.3.4}{7.3.4}）指出，（\hyperlink{para:1:7.4.6}{7.4.6}）中 a) 与 c) 的等价性对这样的预概形仍成立。\transnote{这里原文对 c) 的推广不足以成立：例如两个整仿射概形的不交并上，仅在其中一个分支取结构层、在另一个分支取零的拟凝聚层，其支集不等于全空间，但不是挠层。与 a) 对应的条件应要求支集不包含任何不可约分支，亦即所有分支泛点处的茎都为零；不能只保留 $\Supp(\F)\ne X$。}


% END chapters/ch1_s7_4.tex

\par

% BEGIN chapters/ch1_s8_1.tex
\section{Chevalley 概形}\hypertarget{heading:1:8}{}
\subsection{相关的局部环}\hypertarget{heading:1:8.1}{}
\src{164—165}
对每个局部环 $A$，将其极大理想记作 $\mathfrak m(A)$。
\egastatement{8.1.1}{引理} 设 $A,B$ 为两个局部环，$A\subset B$。下列条件等价：
\begin{enumerate}[label=(\roman*)]
\item $\mathfrak m(B)\cap A=\mathfrak m(A)$；
\item $\mathfrak m(A)\subset\mathfrak m(B)$；
\item $1$ 不属于由 $\mathfrak m(A)$ 在 $B$ 中生成的理想。
\end{enumerate}
显然 (i) 推出 (ii)，而 (ii) 推出 (iii)。最后，若 (iii) 成立，则 $\mathfrak m(B)\cap A$ 包含 $\mathfrak m(A)$ 且不包含 $1$，故等于 $\mathfrak m(A)$。

当（\hyperlink{para:1:8.1.1}{8.1.1}）的等价条件成立时，称 $B$ \emph{支配} $A$；这等价于说嵌入 $A\to B$ 是局部同态。显然，在一个环 $R$ 的所有局部子环的集合中，支配关系是一个序关系。
\egapar{8.1.2} 现在考虑一个域 $R$。对 $R$ 的每个子环 $A$，以 $L(A)$ 表示各局部环 $A_{\mathfrak p}$ 组成的集合，其中 $\mathfrak p$ 遍历 $A$ 的素谱；将这些环等同于 $R$ 的包含 $A$ 的子环。由于 $\mathfrak p=(\mathfrak p A_{\mathfrak p})\cap A$，映射 $\mathfrak p\mapsto A_{\mathfrak p}$ 是 $\Spec(A)$ 到 $L(A)$ 的双射。
\egastatement{8.1.3}{引理} 设 $R$ 为域，$A$ 为其子环。$R$ 的一个局部子环 $M$ 支配某个 $A_{\mathfrak p}\in L(A)$，当且仅当 $A\subset M$。此时，被 $M$ 支配的局部环 $A_{\mathfrak p}$ 唯一，并对应于 $\mathfrak p=\mathfrak m(M)\cap A$。

若 $M$ 支配 $A_{\mathfrak p}$，由（\hyperlink{para:1:8.1.1}{8.1.1}），有 $\mathfrak m(M)\cap A_{\mathfrak p}=\mathfrak p A_{\mathfrak p}$，故 $\mathfrak p$ 唯一。另一方面，若 $A\subset M$，则 $\mathfrak m(M)\cap A=\mathfrak p$ 为 $A$ 的素理想，且 $A\setminus\mathfrak p$ 的元素在 $M$ 中可逆，故 $A_{\mathfrak p}\subset M$ 且 $\mathfrak p A_{\mathfrak p}\subset\mathfrak m(M)$。因此 $M$ 支配 $A_{\mathfrak p}$。
\egastatement{8.1.4}{引理} 设 $R$ 为域，$M,N$ 为 $R$ 的两个局部子环，$P$ 为由 $M\cup N$ 生成的 $R$ 的子环。下列条件等价：
\begin{enumerate}[label=(\roman*)]
\item 存在 $P$ 的素理想 $\mathfrak p$，使 $\mathfrak m(M)=\mathfrak p\cap M$，$\mathfrak m(N)=\mathfrak p\cap N$；
\item 由 $\mathfrak m(M)\cup\mathfrak m(N)$ 在 $P$ 中生成的理想 $\mathfrak a$ 不等于 $P$；
\item 存在 $R$ 的局部子环 $Q$，同时支配 $M,N$。
\end{enumerate}
显然 (i) 推出 (ii)。反之，若 $\mathfrak a\ne P$，则 $\mathfrak a$ 包含于 $P$ 的某个极大理想 $\mathfrak n$；由于 $1\notin\mathfrak n$，$\mathfrak n\cap M$ 包含 $\mathfrak m(M)$ 且不等于 $M$，故 $\mathfrak n\cap M=\mathfrak m(M)$；同理 $\mathfrak n\cap N=\mathfrak m(N)$。若 $Q$ 同时支配 $M,N$，则显然 $P\subset Q$，且
\[
\mathfrak m(M)=(\mathfrak m(Q)\cap P)\cap M,
\qquad
\mathfrak m(N)=(\mathfrak m(Q)\cap P)\cap N,
\]
故 (iii) 推出 (i)；取 $Q=P_{\mathfrak p}$ 即可得到逆向蕴含。

当（\hyperlink{para:1:8.1.4}{8.1.4}）的条件成立时，遵循 C. Chevalley 的术语，称局部环 $M,N$ \emph{相关}。
\egastatement{8.1.5}{命题} 设 $A,B$ 为域 $R$ 的两个子环，$C$ 为由 $A\cup B$ 生成的子环。下列条件等价：
\begin{enumerate}[label=(\roman*)]
\item 对每个包含 $A,B$ 的局部环 $Q$，令 $\mathfrak p=\mathfrak m(Q)\cap A$、$\mathfrak q=\mathfrak m(Q)\cap B$，则有 $A_{\mathfrak p}=B_{\mathfrak q}$；
\item 对 $C$ 的每个素理想 $\mathfrak r$，令 $\mathfrak p=\mathfrak r\cap A$、$\mathfrak q=\mathfrak r\cap B$，则有 $A_{\mathfrak p}=B_{\mathfrak q}$；
\item 若 $M\in L(A)$ 与 $N\in L(B)$ 相关，则它们相同；
\item $L(A)\cap L(B)=L(C)$。
\end{enumerate}
引理（\hyperlink{para:1:8.1.3}{8.1.3}）与（\hyperlink{para:1:8.1.4}{8.1.4}）证明了 (i) 与 (iii) 等价。将 (i) 用于 $Q=C_{\mathfrak r}$，显然得到 (ii)。反之，(ii) 推出 (i)：若 $Q$ 包含 $A\cup B$，则包含 $C$；令 $\mathfrak r=\mathfrak m(Q)\cap C$，由（\hyperlink{para:1:8.1.3}{8.1.3}）有 $\mathfrak p=\mathfrak r\cap A$、$\mathfrak q=\mathfrak r\cap B$。

(iv) 显然推出 (i)。事实上，若 $Q$ 包含 $A\cup B$，由（\hyperlink{para:1:8.1.3}{8.1.3}），它支配一个 $C_{\mathfrak r}\in L(C)$。按假设 $C_{\mathfrak r}\in L(A)\cap L(B)$，而（\hyperlink{para:1:8.1.1}{8.1.1}）与（\hyperlink{para:1:8.1.3}{8.1.3}）给出 $C_{\mathfrak r}=A_{\mathfrak p}=B_{\mathfrak q}$。

最后证明 (iii) 推出 (iv)。取 $Q\in L(C)$。由（\hyperlink{para:1:8.1.3}{8.1.3}），$Q$ 支配某个 $M\in L(A)$ 和某个 $N\in L(B)$；它们相关，故按假设相同。于是 $C\subset M$，所以 $M$ 支配某个 $Q'\in L(C)$（\hyperlink{para:1:8.1.3}{8.1.3}）。因此 $Q$ 支配 $Q'$，由（\hyperlink{para:1:8.1.3}{8.1.3}）必有 $Q=Q'=M$，故 $Q\in L(A)\cap L(B)$。反之，若 $Q\in L(A)\cap L(B)$，则 $C\subset Q$；因而（\hyperlink{para:1:8.1.3}{8.1.3}）$Q$ 支配某个 $Q'\in L(C)\subset L(A)\cap L(B)$。$Q,Q'$ 相关，故相同，于是 $Q=Q'\in L(C)$。证明完毕。


% END chapters/ch1_s8_1.tex

\par

% BEGIN chapters/ch1_s8_2.tex
\subsection{整概形的局部环}\hypertarget{heading:1:8.2}{}
\src{165—168}
\egapar{8.2.1} 设 $X$ 为整预概形，$R$ 为其有理函数域，即 $X$ 的泛点 $a$ 处的局部环。对任意 $x\in X$，已知 $\OO_x$ 典范地等同于 $R$ 的子环（\hyperlink{para:1:7.1.5}{7.1.5}）；对任意有理函数 $f\in R$，其定义域 $\delta(f)$ 为满足 $f\in\OO_x$ 的点 $x\in X$ 组成的开集。由（\hyperlink{para:1:7.2.6}{7.2.6}），对每个开集 $U\subset X$，有
\begin{equation}
\Gamma(U,\OO_X)=\bigcap_{x\in U}\OO_x.
\tag{8.2.1.1}\label{eq:1:8.2.1.1}
\end{equation}
\egastatement{8.2.2}{命题} 设 $X$ 为整预概形，$R$ 为其有理函数域。$X$ 为概形，当且仅当对 $X$ 的任意两点 $x,y$，关系“$\OO_x$ 与 $\OO_y$ 相关”（\hyperlink{para:1:8.1.4}{8.1.4}）蕴含 $x=y$。

假设此条件成立，证明 $X$ 分离。设 $U,V$ 为 $X$ 中的两个不同仿射开集，其环分别为 $A,B$，并等同于 $R$ 的子环。于是 $U$（分别为 $V$）等同于 $L(A)$（分别为 $L(B)$）（\hyperlink{para:1:8.1.2}{8.1.2}）。由假设与（\hyperlink{para:1:8.1.5}{8.1.5}），若 $C$ 为由 $A\cup B$ 生成的子环，则 $W=U\cap V$ 等同于 $L(A)\cap L(B)=L(C)$。已知 $R$ 的任一子环 $E$ 都等于 $L(E)$ 中各局部环的交（\cite{b1}，第 5-03 页，命题 4 bis），所以 $C$ 等同于所有 $z\in W$ 处的 $\OO_z$ 的交，即 $\Gamma(W,\OO_X)$（\hyperref[eq:1:8.2.1.1]{8.2.1.1}）。

考虑 $X$ 在 $W$ 上诱导的预概形。恒等同态 $\varphi:C\to\Gamma(W,\OO_X)$ 对应一个态射 $\Phi=(\psi,\theta):W\to\Spec(C)$（\hyperlink{para:1:2.2.4}{2.2.4}）。我们证明 $\Phi$ 为预概形同构，于是 $W$ 为仿射开集。由 $W$ 与 $L(C)=\Spec(C)$ 的等同，$\psi$ 为双射。另一方面，对每个 $x\in W$，令 $\mathfrak r=\mathfrak m_x\cap C$，则 $\theta_x^\sharp$ 为嵌入 $C_{\mathfrak r}\to\OO_x$；按定义 $C_{\mathfrak r}$ 等同于 $\OO_x$，故该同态为双射。

只须再证明 $\psi$ 为同胚，即 $W$ 的每个闭子集 $F$ 的像 $\psi(F)$ 在 $\Spec(C)$ 中闭。$F$ 是 $U$ 的某个形如 $V(\mathfrak a)$ 的闭集在 $W$ 上的迹，其中 $\mathfrak a$ 为 $A$ 的理想。证明 $\psi(F)=V(\mathfrak a C)$ 即足够。事实上，$C$ 中包含 $\mathfrak a C$ 的素理想就是包含 $\mathfrak a$ 的素理想，亦即形如 $\psi(x)=\mathfrak m_x\cap C$ 且 $\mathfrak a\subset\mathfrak m_x$、$x\in W$ 的理想；而在 $U$ 中，$\mathfrak a\subset\mathfrak m_x$ 等价于 $x\in V(\mathfrak a)$，故有 $\psi(F)=V(\mathfrak a C)$。

因此 $X$ 分离，因为 $U\cap V$ 仿射且其环 $C$ 由 $U,V$ 的环的并 $A\cup B$ 生成（\hyperlink{para:1:5.5.6}{5.5.6}）。

反之，假设 $X$ 分离，并取两点 $x,y\in X$，使 $\OO_x,\OO_y$ 相关。取包含 $x$（分别为 $y$）的仿射开集 $U$（分别为 $V$），环为 $A$（分别为 $B$）。已知 $U\cap V$ 仿射，其环 $C$ 由 $A\cup B$ 生成（\hyperlink{para:1:5.5.6}{5.5.6}）。令 $\mathfrak p=\mathfrak m_x\cap A$、$\mathfrak q=\mathfrak m_y\cap B$，则 $A_{\mathfrak p}=\OO_x$、$B_{\mathfrak q}=\OO_y$。因它们相关，存在 $C$ 的素理想 $\mathfrak r$，使 $\mathfrak p=\mathfrak r\cap A$、$\mathfrak q=\mathfrak r\cap B$（\hyperlink{para:1:8.1.4}{8.1.4}）。由 $U\cap V$ 仿射，存在 $z\in U\cap V$ 使 $\mathfrak r=\mathfrak m_z\cap C$，显然 $x=z$ 且 $y=z$，故 $x=y$。
\egastatement{8.2.3}{推论} 设 $X$ 为整概形，$x,y\in X$。则 $x\in\overline{\{y\}}$，当且仅当 $\OO_x\subset\OO_y$；换言之，当且仅当每个在 $x$ 处有定义的有理函数也在 $y$ 处有定义。

必要性显然，因为有理函数 $f\in R$ 的定义域 $\delta(f)$ 是开集。证明充分性。若 $\OO_x\subset\OO_y$，则存在 $\OO_x$ 的素理想 $\mathfrak p$，使 $\OO_y$ 支配 $(\OO_x)_{\mathfrak p}$（\hyperlink{para:1:8.1.3}{8.1.3}）。又由（\hyperlink{para:1:2.4.2}{2.4.2}），存在 $z\in X$ 使 $x\in\overline{\{z\}}$，且 $\OO_z=(\OO_x)_{\mathfrak p}$。因 $\OO_z,\OO_y$ 相关，由（\hyperlink{para:1:8.2.2}{8.2.2}）得 $z=y$，从而得到推论。
\egastatement{8.2.4}{推论} 若 $X$ 为整概形，则映射 $x\mapsto\OO_x$ 为单射。换言之，若 $x,y$ 为 $X$ 的两个不同点，则存在一个有理函数，在其中一点有定义而在另一点无定义。

这由（\hyperlink{para:1:8.2.3}{8.2.3}）与公理（$T_0$）（\hyperlink{para:1:2.1.4}{2.1.4}）得出。
\egastatement{8.2.5}{推论} 设 $X$ 为底空间诺特的整概形。当 $f$ 遍历 $X$ 的有理函数域 $R$ 时，各集合 $\delta(f)$ 生成 $X$ 的拓扑。

此时 $X$ 的每个闭子集都是有限个不可约闭集的并，也就是有限个形如 $\overline{\{y\}}$ 的集合之并（\hyperlink{para:1:2.1.5}{2.1.5}）。若 $x\notin\overline{\{y\}}$，则存在在 $x$ 处有定义而在 $y$ 处无定义的有理函数 $f$（\hyperlink{para:1:8.2.3}{8.2.3}），亦即 $x\in\delta(f)$ 且 $\delta(f)$ 不与 $\overline{\{y\}}$ 相交。因此 $\overline{\{y\}}$ 的补集为形如 $\delta(f)$ 的集合之并；由第一句话，$X$ 的每个开集都是这类开集的有限交的并。
\egapar{8.2.6} 推论（\hyperlink{para:1:8.2.5}{8.2.5}）表明，给定以 $R$ 为分式域的局部环族 $(\OO_x)_{x\in X}$，就完全确定了 $X$ 的拓扑。也可以等价地如下定义 $X$ 的闭集：给定 $X$ 的有限子集 $\{x_1,\ldots,x_n\}$，考虑所有满足 $\OO_y\subset\OO_{x_i}$ 对至少一个 $i$ 成立的 $y\in X$；遍历所有有限子集，所得集合正是 $X$ 的闭集。此外，一旦知道 $X$ 的拓扑，结构层 $\OO_X$ 也被各 $\OO_x$ 唯一确定，因为（\hyperref[eq:1:8.2.1.1]{8.2.1.1}）给出 $\Gamma(U,\OO_X)=\bigcap_{x\in U}\OO_x$。所以，当 $X$ 为底空间诺特的整概形时，族 $(\OO_x)_{x\in X}$ 完全确定预概形 $X$。
\egastatement{8.2.7}{命题} 设 $X,Y$ 为两个整概形，$f:X\to Y$ 为支配态射（\hyperlink{para:1:2.2.6}{2.2.6}），$K$（分别为 $L$）为 $X$（分别为 $Y$）的有理函数域。则 $L$ 等同于 $K$ 的子域，且对每个 $x\in X$，$\OO_{f(x)}$ 是 $Y$ 的唯一被 $\OO_x$ 支配的局部环。

若 $f=(\psi,\theta)$，$a$ 为 $X$ 的泛点，则 $\psi(a)$ 是 $Y$ 的泛点（0，\hyperlink{para:0:2.1.5}{2.1.5}）。因此 $\theta_a^\sharp$ 为域 $L=\OO_{\psi(a)}$ 到域 $K=\OO_a$ 的单同态。由于 $Y$ 的每个非空仿射开集 $U$ 都包含 $\psi(a)$，由（\hyperlink{para:1:2.2.4}{2.2.4}），与 $f$ 对应的同态
\[
 \Gamma(U,\OO_Y)\longrightarrow\Gamma(\psi^{-1}(U),\OO_X)
\]
就是 $\theta_a^\sharp$ 在 $\Gamma(U,\OO_Y)$ 上的限制。因此，对每个 $x\in X$，$\theta_x^\sharp$ 为 $\theta_a^\sharp$ 在 $\OO_{\psi(x)}$ 上的限制，故也是单同态。又已知 $\theta_x^\sharp$ 为局部同态，所以通过 $\theta_a^\sharp$ 将 $L$ 等同于 $K$ 的子域后，$\OO_{\psi(x)}$ 被 $\OO_x$ 支配（\hyperlink{para:1:8.1.1}{8.1.1}）。它也是 $Y$ 的唯一被 $\OO_x$ 支配的局部环，因为 $Y$ 的两个相关局部环必相同（\hyperlink{para:1:8.2.2}{8.2.2}）。
\egastatement{8.2.8}{命题} 设 $X$ 为不可约预概形，$f:X\to Y$ 为局部浸入（分别为局部同构），并假设态射 $f$ 分离。则 $f$ 为浸入（分别为开浸入）。

令 $f=(\psi,\theta)$。两种情形都只须证明 $\psi$ 为 $X$ 到 $\psi(X)$ 的同胚（\hyperlink{para:1:4.5.3}{4.5.3}）。以 $f_{\mathrm{red}}$ 代替 $f$（\hyperlink{para:1:5.1.6}{5.1.6} 与 \hyperlink{para:1:5.5.1}{5.5.1}，(vi)），可假设 $X,Y$ 约化。令 $Y'$ 为 $Y$ 中以 $\overline{\psi(X)}$ 为底空间的约化闭子预概形，则 $f$ 分解为 $X\xrightarrow{f'}Y'\xrightarrow{j}Y$，其中 $j$ 为典范嵌入（\hyperlink{para:1:5.2.2}{5.2.2}）。由（\hyperlink{para:1:5.5.1}{5.5.1}，(v)），$f'$ 仍分离；它也仍是局部浸入（分别为局部同构）：问题在 $X,Y$ 上都是局部的，故只须考虑 $f$ 为闭浸入（分别为开浸入）的情形，此时断言立即由（\hyperlink{para:1:4.2.2}{4.2.2}）得出。

所以可假设 $f$ 支配，这又使 $Y$ 不可约（0，\hyperlink{para:0:2.1.5}{2.1.5}），因而 $X,Y$ 均为整预概形。问题在 $Y$ 上是局部的，故可假设 $Y$ 仿射；因 $f$ 分离，$X$ 为概形（\hyperlink{para:1:5.5.1}{5.5.1}，(ii)）。于是已处于（\hyperlink{para:1:8.2.7}{8.2.7}）的假设之下。因此对每个 $x\in X$，$\theta_x^\sharp$ 为单射；而 $f$ 为局部浸入，又使 $\theta_x^\sharp$ 满射（\hyperlink{para:1:4.2.2}{4.2.2}），故为双射，即在（\hyperlink{para:1:8.2.7}{8.2.7}）的等同下有 $\OO_{\psi(x)}=\OO_x$。由（\hyperlink{para:1:8.2.4}{8.2.4}），$\psi$ 为单射；当 $f$ 为局部同构时，这就已证明命题（\hyperlink{para:1:4.5.3}{4.5.3}）。

若仅假设 $f$ 为局部浸入，对每个 $x\in X$，存在 $x$ 在 $X$ 中的开邻域 $U$，以及 $\psi(x)$ 在 $Y$ 中的开邻域 $V$，使 $\psi$ 在 $U$ 上的限制为 $U$ 到 $V$ 的一个闭子集的同胚。由于 $U$ 在 $X$ 中稠密，$\psi(U)$ 在 $Y$ 中稠密，因而在 $V$ 中稠密，故 $\psi(U)=V$。由 $\psi$ 单射，$\psi^{-1}(V)=U$，这就证明了 $\psi$ 为 $X$ 到 $\psi(X)$ 的同胚。


% END chapters/ch1_s8_2.tex

\par

% BEGIN chapters/ch1_s8_3.tex
\subsection{Chevalley 概形}\hypertarget{heading:1:8.3}{}
\src{168—169}
\egapar{8.3.1} 设 $X$ 为诺特整概形，$R$ 为其有理函数域；记 $X'$ 为所有局部子环 $\OO_x\subset R$（$x$ 遍历 $X$）组成的集合。$X'$ 满足以下三个条件。

\noindent\textbf{（Sch.1）} 对每个 $M\in X'$，$R$ 是 $M$ 的分式域。

\noindent\textbf{（Sch.2）} 存在 $R$ 的有限个诺特子环 $A_i$，使 $X'=\bigcup_iL(A_i)$，并且对每对指标 $i,j$，由 $A_i\cup A_j$ 生成的 $R$ 的子环 $A_{ij}$ 是有限型 $A_i$-代数。

\noindent\textbf{（Sch.3）} 若 $X'$ 的两个元素 $M,N$ 相关，则它们相同。

事实上，（Sch.1）已见于（\hyperlink{para:1:8.2.1}{8.2.1}），而（Sch.3）由（\hyperlink{para:1:8.2.2}{8.2.2}）得出。为证明（Sch.2），只须用有限个环为诺特环的仿射开集 $U_i$ 覆盖 $X$，并取 $A_i=\Gamma(U_i,\OO_X)$。$X$ 为概形的假设蕴含 $U_i\cap U_j$ 仿射，且 $\Gamma(U_i\cap U_j,\OO_X)=A_{ij}$（\hyperlink{para:1:5.5.6}{5.5.6}）。此外，由于空间 $U_i$ 诺特，浸入 $U_i\cap U_j\to U_i$ 为有限型（\hyperlink{para:1:6.3.5}{6.3.5}），故 $A_{ij}$ 是有限型 $A_i$-代数（\hyperlink{para:1:6.3.3}{6.3.3}）。
\egapar{8.3.2} 以（Sch.1）、（Sch.2）、（Sch.3）为公理的结构，推广了 C. Chevalley 意义下的“概形”。Chevalley 还假设 $R$ 是某个域 $K$ 的有限型扩张，且各 $A_i$ 是有限型 $K$-代数；这使（Sch.2）中的一部分成为多余的\cite{b1}。反之，若在集合 $X'$ 上给定这样的结构，便可利用（\hyperlink{para:1:8.2.6}{8.2.6}）的说明，为它构造一个相伴的整概形 $X$：取 $X$ 的底空间为配备（\hyperlink{para:1:8.2.6}{8.2.6}）所定义拓扑的 $X'$，并取层 $\OO_X$ 使得对每个开集 $U\subset X$，都有
\[
 \Gamma(U,\OO_X)=\bigcap_{x\in U}\OO_x,
\]
限制同态按显然的方式定义。这样确实得到一个整概形，其局部环正是 $X'$ 的元素；我们将此验证留给读者。以后不使用这个结果。


% END chapters/ch1_s8_3.tex

\par

% BEGIN chapters/ch1_s9_1.tex
\section{拟凝聚层的补充知识}\hypertarget{heading:1:9}{}
\subsection{拟凝聚层的张量积}\hypertarget{heading:1:9.1}{}
\src{169—171}
\egastatement{9.1.1}{命题} 设 $X$ 为预概形（分别为局部诺特预概形），$\F,\G$ 为两个拟凝聚（分别为凝聚）$\OO_X$-模。则 $\F\otimes_{\OO_X}\G$ 拟凝聚（分别为凝聚）；若 $\F,\G$ 为有限型，则该张量积也为有限型。若 $\F$ 有有限表示，且 $\G$ 拟凝聚（分别为凝聚），则 $\shHom_{\OO_X}(\F,\G)$ 拟凝聚（分别为凝聚）。

问题是局部的，可假设 $X$ 仿射（分别为诺特仿射）。此外，若 $\F$ 凝聚，可假设它是某个同态 $\OO_X^p\to\OO_X^q$ 的余核。拟凝聚层的断言由（\hyperlink{para:1:1.3.12}{1.3.12}）与（\hyperlink{para:1:1.3.9}{1.3.9}）得出；凝聚层的断言由（\hyperlink{para:1:1.5.1}{1.5.1}）及以下事实得出：若 $M,N$ 为诺特环 $A$ 上的有限型模，则 $M\otimes_A N$ 与 $\Hom_A(M,N)$ 都是有限型 $A$-模。
\egastatement{9.1.2}{定义} 设 $X,Y$ 为两个 $S$-预概形，$p,q$ 为 $X\times_S Y$ 的投影，$\F$（分别为 $\G$）为拟凝聚 $\OO_X$-模（分别为 $\OO_Y$-模）。称预概形 $X\times_S Y$ 上的张量积
\[
 p^*(\F)\otimes_{\OO_{X\times_S Y}}q^*(\G)
\]
为 $\F,\G$ 在 $\OO_S$ 上（或 $S$ 上）的\emph{张量积}，记作 $\F\otimes_{\OO_S}\G$ 或 $\F\otimes_S\G$。

若 $X_i$（$1\leq i\leq n$）为 $S$-预概形，$\F_i$ 为拟凝聚 $\OO_{X_i}$-模，则同样可在 $Z=X_1\times_S\cdots\times_S X_n$ 上定义张量积 $\F_1\otimes_S\cdots\otimes_S\F_n$。由（\hyperlink{para:1:9.1.1}{9.1.1}）与（0，\hyperlink{para:0:5.1.4}{5.1.4}），它为拟凝聚 $\OO_Z$-模；若各 $\F_i$ 凝聚且 $Z$ 局部诺特，则由（\hyperlink{para:1:9.1.1}{9.1.1}）、（0，\hyperlink{para:0:5.3.11}{5.3.11}）与（\hyperlink{para:1:6.1.1}{6.1.1}），它凝聚。

取 $X=Y=S$ 时，定义（\hyperlink{para:1:9.1.2}{9.1.2}）给出通常的 $\OO_S$-模的张量积。此外，由 $q^*(\OO_Y)=\OO_{X\times_S Y}$（0，\hyperlink{para:0:4.3.4}{4.3.4}），张量积 $\F\otimes_S\OO_Y$ 典范地等同于 $p^*(\F)$；同理，$\OO_X\otimes_S\G$ 典范地等同于 $q^*(\G)$。更特别地，取 $Y=S$，并令 $f:X\to Y$ 为结构态射，就有 $\OO_X\otimes_Y\G=f^*(\G)$。因此通常的张量积与逆像都是此一般张量积的特例。

由定义（\hyperlink{para:1:9.1.2}{9.1.2}）立即得知，对固定的 $X,Y$，$\F\otimes_S\G$ 是关于 $\F,\G$ 的协变、可加、右正合双函子。
\egastatement{9.1.3}{命题} 设 $S,X,Y$ 为三个仿射概形，环分别为 $A,B,C$，其中 $B,C$ 为 $A$-代数。设 $M$（分别为 $N$）为 $B$-模（分别为 $C$-模），$\F=\wt M$（分别为 $\G=\wt N$）为相伴的拟凝聚层。则 $\F\otimes_S\G$ 典范地同构于 $(B\otimes_A C)$-模 $M\otimes_A N$ 所伴随的层。

事实上，由（\hyperlink{para:1:1.6.5}{1.6.5}），$\F\otimes_S\G$ 典范地同构于以下 $(B\otimes_A C)$-模所伴随的层：
\[
 \bigl(M\otimes_B(B\otimes_A C)\bigr)
 \otimes_{B\otimes_A C}
 \bigl((B\otimes_A C)\otimes_C N\bigr).
\]
由张量积间的典范同构，该模同构于
\[
 M\otimes_B(B\otimes_A C)\otimes_C N
 =(M\otimes_B B)\otimes_A(C\otimes_C N)=M\otimes_A N.
\]
\egastatement{9.1.4}{命题} 设 $f:T\to X$、$g:T\to Y$ 为两个 $S$-态射，$\F$（分别为 $\G$）为拟凝聚 $\OO_X$-模（分别为 $\OO_Y$-模）。则
\[
 (f,g)_S^*(\F\otimes_S\G)
 =f^*(\F)\otimes_{\OO_T}g^*(\G).
\]
若 $p,q$ 为 $X\times_S Y$ 的投影，则此式由 $(f,g)_S^*\circ p^*=f^*$ 与 $(f,g)_S^*\circ q^*=g^*$（0，\hyperlink{para:0:3.5.5}{3.5.5}），以及代数层的张量积之逆像等于其逆像的张量积（0，\hyperlink{para:0:4.3.3}{4.3.3}）得出。
\egastatement{9.1.5}{推论} 设 $f:X\to X'$、$g:Y\to Y'$ 为两个 $S$-态射，$\F'$（分别为 $\G'$）为拟凝聚 $\OO_{X'}$-模（分别为 $\OO_{Y'}$-模）。则
\[
 (f\times_S g)^*(\F'\otimes_S\G')
 =f^*(\F')\otimes_Sg^*(\G').
\]
这由（\hyperlink{para:1:9.1.4}{9.1.4}）及 $f\times_Sg=(f\circ p,g\circ q)_S$ 得出，其中 $p,q$ 为 $X\times_S Y$ 的投影。
\egastatement{9.1.6}{推论} 设 $X,Y,Z$ 为三个 $S$-预概形，$\F,\G,\HH$ 分别为其上的拟凝聚模。则层 $\F\otimes_S\G\otimes_S\HH$，是 $(\F\otimes_S\G)\otimes_S\HH$ 经典范同构
\[
 X\times_S Y\times_S Z\longrightarrow(X\times_S Y)\times_S Z
\]
的逆像。

事实上，记 $p_1,p_2,p_3$ 为三重积的投影，则此同构写为 $((p_1,p_2)_S,p_3)_S$。

同样，$\G\otimes_S\F$ 经 $X\times_S Y\to Y\times_S X$ 的典范同构的逆像，是 $\F\otimes_S\G$。
\egastatement{9.1.7}{推论} 若 $X$ 为 $S$-预概形，则每个拟凝聚 $\OO_X$-模 $\F$ 都是 $\F\otimes_S\OO_S$ 经典范同构 $X\to X\times_S S$（\hyperlink{para:1:3.3.3}{3.3.3}）的逆像。

此同构为 $(1_X,\varphi)_S$，其中 $\varphi:X\to S$ 为结构态射；推论由（\hyperlink{para:1:9.1.4}{9.1.4}）及 $\varphi^*(\OO_S)=\OO_X$ 得出。
\egapar{9.1.8} 设 $X$ 为 $S$-预概形，$\F$ 为拟凝聚 $\OO_X$-模，$\varphi:S'\to S$ 为态射。将 $X\times_S S'=X_{(\varphi)}=X_{(S')}$ 上的拟凝聚层 $\F\otimes_S\OO_{S'}$ 记作 $\F_{(\varphi)}$ 或 $\F_{(S')}$。因此 $\F_{(S')}=p^*(\F)$，其中 $p:X_{(S')}\to X$ 为投影。
\egastatement{9.1.9}{命题} 设 $\varphi':S''\to S'$ 为态射。对 $S$-预概形 $X$ 上的每个拟凝聚 $\OO_X$-模 $\F$，$(\F_{(\varphi)})_{(\varphi')}$ 是 $\F_{(\varphi\circ\varphi')}$ 经典范同构 $(X_{(\varphi)})_{(\varphi')}\to X_{(\varphi\circ\varphi')}$（\hyperlink{para:1:3.3.9}{3.3.9}）的逆像。

这立即由定义与（\hyperlink{para:1:3.3.9}{3.3.9}）得出，也可写为
\begin{equation}
 (\F\otimes_S\OO_{S'})\otimes_{S'}\OO_{S''}
 =\F\otimes_S\OO_{S''}.
\tag{9.1.9.1}\label{eq:1:9.1.9.1}
\end{equation}
\egastatement{9.1.10}{命题} 设 $Y$ 为 $S$-预概形，$f:X\to Y$ 为 $S$-态射。对每个拟凝聚 $\OO_Y$-模 $\G$ 及每个态射 $S'\to S$，有
\[
 (f_{(S')})^*(\G_{(S')})=(f^*(\G))_{(S')}.
\]
这立即由下图的交换性得出：
\[
\begin{tikzcd}
 X_{(S')}\arrow[r,"f_{(S')}"]\arrow[d]&Y_{(S')}\arrow[d]\\
 X\arrow[r,"f"]&Y.
\end{tikzcd}
\]
\egastatement{9.1.11}{推论} 设 $X,Y$ 为两个 $S$-预概形，$\F$（分别为 $\G$）为拟凝聚 $\OO_X$-模（分别为 $\OO_Y$-模）。层 $\F_{(S')}\otimes_{S'}\G_{(S')}$ 经典范同构
\[
 (X\times_S Y)_{(S')}\longrightarrow X_{(S')}\times_{S'}Y_{(S')}
 \qquad\text{（\hyperlink{para:1:3.3.10}{3.3.10}）}
\]
的逆像，等于 $(\F\otimes_S\G)_{(S')}$。

若 $p,q$ 为 $X\times_S Y$ 的投影，此同构就是 $(p_{(S')},q_{(S')})_{S'}$；推论由（\hyperlink{para:1:9.1.4}{9.1.4}）与（\hyperlink{para:1:9.1.10}{9.1.10}）得出。
\egastatement{9.1.12}{命题} 沿用（\hyperlink{para:1:9.1.2}{9.1.2}）的记号，设 $z\in X\times_S Y$，$x=p(z)$，$y=q(z)$。则茎 $(\F\otimes_S\G)_z$ 同构于
\[
 (\F_x\otimes_{\OO_x}\OO_z)\otimes_{\OO_z}
 (\G_y\otimes_{\OO_y}\OO_z)
 =\F_x\otimes_{\OO_x}\OO_z\otimes_{\OO_y}\G_y.
\]
可归结到仿射情形，所以命题由（\hyperlink{para:1:1.6.5}{1.6.5}）中的茎公式得出。\transnote{原文引用“（\hyperref[eq:1:1.6.5.1]{1.6.5.1}）”，但该处的相应公式在底本中未单独编号。}
\egastatement{9.1.13}{推论} 若 $\F,\G$ 为有限型，则
\[
 \Supp(\F\otimes_S\G)
 =p^{-1}(\Supp(\F))\cap q^{-1}(\Supp(\G)).
\]
由于 $p^*(\F)$ 与 $q^*(\G)$ 都是有限型 $\OO_{X\times_S Y}$-模，利用（\hyperlink{para:1:9.1.12}{9.1.12}）与（0，\hyperlink{para:0:1.7.5}{1.7.5}），可归结到 $\G=\OO_Y$ 的特例，即证明
\begin{equation}
 \Supp(p^*(\F))=p^{-1}(\Supp(\F)).
\tag{9.1.13.1}\label{eq:1:9.1.13.1}
\end{equation}
\transnote{原文左端印作 $\Supp(p^{-1}(\F))$；上下文涉及的是模的拉回 $p^*(\F)$，此处按该含义书写，并保留原编号。}
与（0，\hyperlink{para:0:1.7.5}{1.7.5}）相同的论证，将问题归结到验证：对每个 $z\in X\times_S Y$，令 $x=p(z)$，有 $\OO_z/\mathfrak m_x\OO_z\ne0$。这来自同态 $\OO_x\to\OO_z$ 按假设为局部同态。

我们将本小节中两个因子情形的结果推广到任意多个因子的工作留给读者。


% END chapters/ch1_s9_1.tex

\par

% BEGIN chapters/ch1_s9_2.tex
\subsection{拟凝聚层的直像}\hypertarget{heading:1:9.2}{}
\src{171—172}
\egastatement{9.2.1}{命题} 设 $f:X\to Y$ 为预概形态射。假设 $Y$ 存在一个由仿射开集组成的覆盖 $(Y_\alpha)$，具有下述性质：每个 $f^{-1}(Y_\alpha)$ 都有由包含于它的仿射开集组成的有限覆盖 $(X_{\alpha i})$，并且每个交集 $X_{\alpha i}\cap X_{\alpha j}$ 本身也是有限个仿射开集之并。那么，对每个拟凝聚 $\OO_X$-模 $\F$，$f_*(\F)$ 为拟凝聚 $\OO_Y$-模。

问题在 $Y$ 上是局部的，可假设 $Y$ 等于某个 $Y_\alpha$，并省略下标 $\alpha$。

a) 先设各 $X_i\cap X_j$ 本身就是仿射开集。令 $\F_i=\F|X_i$、$\F_{ij}=\F|(X_i\cap X_j)$，并将它们分别经 $f$ 在 $X_i$ 及 $X_i\cap X_j$ 上的限制所得的直像记作 $\F_i'$ 与 $\F_{ij}'$；已知这些直像拟凝聚（\hyperlink{para:1:1.6.3}{1.6.3}）。令 $\G=\bigoplus_i\F_i'$、$\HH=\bigoplus_{i,j}\F_{ij}'$；它们为拟凝聚 $\OO_Y$-模。定义同态 $u:\G\to\HH$，使 $f_*(\F)$ 为其核，即可推出 $f_*(\F)$ 拟凝聚（\hyperlink{para:1:1.3.9}{1.3.9}）。只须将 $u$ 定义为预层同态。因此，对每个开集 $W\subset Y$，只须定义同态
\[
 u_W:\bigoplus_i\Gamma(f^{-1}(W)\cap X_i,\F)
 \longrightarrow\bigoplus_{i,j}\Gamma(f^{-1}(W)\cap X_i\cap X_j,\F),
\]
使其满足关于 $W$ 的通常相容条件。对截面 $s_i\in\Gamma(f^{-1}(W)\cap X_i,\F)$，以 $s_{ij}$ 表示其在 $f^{-1}(W)\cap X_i\cap X_j$ 上的限制，取
\[
 u_W((s_i)_i)=(s_{ij}-s_{ji})_{i,j}.
\]
相容条件显然成立。令 $\mathscr K=\Ker(u)$，定义 $f_*(\F)\to\mathscr K$：将每个 $s\in\Gamma(f^{-1}(W),\F)$ 映到其在各 $f^{-1}(W)\cap X_i$ 上的限制组成的族 $(s_i)$。层公理（F1）与（F2）（G，II，1.1）说明此同态为双射，从而证明了此种情形。

b) 一般情形，只要证明各 $\F_{ij}'$ 拟凝聚，上面的论证就仍适用。按假设，$X_i\cap X_j$ 为有限个仿射开集 $X_{ijk}$ 之并；这些 $X_{ijk}$ 是一个概形中的仿射开集，故其中任意两个之交仍是仿射开集（\hyperlink{para:1:5.5.6}{5.5.6}）。于是归结到第一种情形，（\hyperlink{para:1:9.2.1}{9.2.1}）得证。
\egastatement{9.2.2}{推论} （\hyperlink{para:1:9.2.1}{9.2.1}）的结论在下列各情形成立：
\begin{enumerate}[label=\alph*)]
\item $f$ 分离且拟紧；
\item $f$ 分离且为有限型；
\item $f$ 拟紧，且 $X$ 的底空间局部诺特。
\end{enumerate}
在 a) 中，各 $X_i\cap X_j$ 仿射（\hyperlink{para:1:5.5.6}{5.5.6}）。b) 是 a) 的特例（\hyperlink{para:1:6.6.3}{6.6.3}）。最后，在 c) 中，可归结到 $Y$ 仿射且 $X$ 的底空间诺特的情形；此时 $X$ 有有限仿射开覆盖 $(X_i)$，而各 $X_i\cap X_j$ 拟紧，故为有限个仿射开集之并（\hyperlink{para:1:2.1.3}{2.1.3}）。


% END chapters/ch1_s9_2.tex

\par

% BEGIN chapters/ch1_s9_3.tex
\subsection{拟凝聚层截面的延拓}\hypertarget{heading:1:9.3}{}
\src{172—173}
\egastatement{9.3.1}{定理} 设 $X$ 为底空间诺特的预概形，或为底空间拟紧的概形。设 $\LL$ 为可逆 $\OO_X$-模（0，\hyperlink{para:0:5.4.1}{5.4.1}），$f$ 为 $\LL$ 在 $X$ 上的截面，$X_f$ 为所有满足 $f(x)\ne0$ 的 $x\in X$ 组成的开集（0，\hyperlink{para:0:5.5.1}{5.5.1}），$\F$ 为拟凝聚 $\OO_X$-模。
\begin{enumerate}[label=(\roman*)]
\item 若 $s\in\Gamma(X,\F)$ 且 $s|X_f=0$，则存在整数 $n>0$ 使 $s\otimes f^{\otimes n}=0$。
\item 对每个截面 $s\in\Gamma(X_f,\F)$，存在整数 $n>0$，使 $s\otimes f^{\otimes n}$ 可延拓为 $\F\otimes\LL^{\otimes n}$ 在 $X$ 上的截面。
\end{enumerate}
(i) 由于 $X$ 的底空间拟紧，它是有限个仿射开集 $U_i$ 之并，可使每个 $\LL|U_i$ 同构于 $\OO_X|U_i$。故归结到 $X$ 仿射且 $\LL=\OO_X$ 的情形。此时 $f$ 等同于 $A(X)$ 的一个元素，且 $X_f=D(f)$；$s$ 等同于某个 $A(X)$-模 $M$ 中的元素，而 $s|X_f$ 等同于 $M_f$ 中的相应元素。由分式模的定义，结论显然。

(ii) 同样可将 $X$ 表为有限个仿射开集 $U_i$（$1\leq i\leq r$）之并，使 $\LL|U_i\simeq\OO_X|U_i$。对每个 $i$，经此同构，$(s\otimes f^{\otimes n})|(U_i\cap X_f)$ 等同于 $(f|(U_i\cap X_f))^n(s|(U_i\cap X_f))$。由（\hyperlink{para:1:1.4.1}{1.4.1}），存在整数 $n>0$，使对每个 $i$，此截面可延拓为 $\F\otimes\LL^{\otimes n}$ 在 $U_i$ 上的截面 $s_i$。记 $s_{ij}$ 为 $s_i$ 在 $U_i\cap U_j$ 上的限制；按定义，$s_{ij}-s_{ji}$ 在 $X_f\cap U_i\cap U_j$ 中为零。

若 $X$ 的底空间诺特，则 $U_i\cap U_j$ 拟紧；若 $X$ 为概形，则 $U_i\cap U_j$ 为仿射开集（\hyperlink{para:1:5.5.6}{5.5.6}），故仍拟紧。由 (i)，存在与 $i,j$ 无关的整数 $m$，使 $(s_{ij}-s_{ji})\otimes f^{\otimes m}=0$。于是立即得到 $\F\otimes\LL^{\otimes(n+m)}$ 在 $X$ 上的一个截面 $s'$，在各 $U_i$ 上诱导 $s_i\otimes f^{\otimes m}$，从而在 $X_f$ 上诱导 $s\otimes f^{\otimes(n+m)}$。

以下推论用更代数化的语言解释定理（\hyperlink{para:1:9.3.1}{9.3.1}）。
\egastatement{9.3.2}{推论} 在（\hyperlink{para:1:9.3.1}{9.3.1}）的假设下，考虑分次环 $A_* =\Gamma_*(\LL)$ 及分次 $A_*$-模 $M_* =\Gamma_*(\LL,\F)$（0，\hyperlink{para:0:5.4.6}{5.4.6}）。若 $f\in A_n$，$n\in\mathbf Z$，则有典范同构
\[
 \Gamma(X_f,\F)\simeq((M_*)_f)_0,
\]
其中右端为分式模 $(M_*)_f$ 的零次齐次元素组成的子群。
\egastatement{9.3.3}{推论} 假设满足（\hyperlink{para:1:9.3.1}{9.3.1}）的假设，且 $\LL=\OO_X$。令 $A=\Gamma(X,\OO_X)$、$M=\Gamma(X,\F)$，则 $A$-模 $\Gamma(X_f,\F)$ 典范地同构于 $M_f$。
\egastatement{9.3.4}{命题} 设 $X$ 为诺特预概形，$\F$ 为凝聚 $\OO_X$-模，$\J$ 为 $\OO_X$ 中的凝聚理想层，且 $\F$ 的支集包含于 $\OO_X/\J$ 的支集中。那么，存在整数 $n>0$，使 $\J^n\F=0$。

由于 $X$ 为有限个环诺特的仿射开集之并，可假设 $X$ 为诺特环 $A$ 的仿射概形。此时 $\F=\wt M$，其中 $M=\Gamma(X,\F)$ 为有限型 $A$-模；$\J=\wt J$，其中 $J=\Gamma(X,\J)$ 为 $A$ 的理想（\hyperlink{para:1:1.4.1}{1.4.1} 与 \hyperlink{para:1:1.5.1}{1.5.1}）。因 $A$ 诺特，$J$ 有有限生成系 $f_i$（$1\leq i\leq m$）。按假设，$\F$ 在 $X$ 上的任一截面在每个 $D(f_i)$ 上都为零。若 $s_j$（$1\leq j\leq q$）为生成 $M$ 的截面，则存在与 $i,j$ 无关的整数 $h$，使 $f_i^h s_j=0$（\hyperlink{para:1:1.4.1}{1.4.1}）；故 $f_i^hs=0$ 对每个 $s\in M$ 成立。取 $n=mh$，即有 $J^nM=0$；相应的 $\OO_X$-模 $\J^n\F=\widetilde{J^nM}$（\hyperlink{para:1:1.3.13}{1.3.13}）也为零。
\egastatement{9.3.5}{推论} 在（\hyperlink{para:1:9.3.4}{9.3.4}）的假设下，存在 $X$ 的闭子预概形 $Y$，其底空间为 $\OO_X/\J$ 的支集，并且若 $j:Y\to X$ 为典范嵌入，则 $\F=j_*(j^*(\F))$。

首先，$\OO_X/\J$ 与 $\OO_X/\J^n$ 的支集相同：若 $\J_x=\OO_x$，则 $\J_x^n=\OO_x$，而另一方面对每个 $x\in X$ 都有 $\J_x^n\subset\J_x$。故利用（\hyperlink{para:1:9.3.4}{9.3.4}），可假设 $\J\F=0$。取 $Y$ 为由 $\J$ 定义的闭子预概形；由于 $\F$ 此时为 $(\OO_X/\J)$-模，结论立即成立。


% END chapters/ch1_s9_3.tex

\par

% BEGIN chapters/ch1_s9_4.tex
\subsection{拟凝聚层的延拓}\hypertarget{heading:1:9.4}{}
\src{174—176}
\egapar{9.4.1} 设 $X$ 为拓扑空间，$\F$ 为 $X$ 上的集合层（分别为群层、环层），$U$ 为 $X$ 的开子集，$\psi:U\to X$ 为典范嵌入，$\G$ 为 $\F|U=\psi^*(\F)$ 的子层。由于 $\psi_*$ 左正合，$\psi_*(\G)$ 为 $\psi_*(\psi^*(\F))$ 的子层。考虑典范同态 $\rho:\F\to\psi_*(\psi^*(\F))$（0，\hyperlink{para:0:3.5.3}{3.5.3}），将 $\F$ 的子层 $\rho^{-1}(\psi_*(\G))$ 记作 $\overline\G$。由定义立即得知，对 $X$ 的任一开集 $V$，$\Gamma(V,\overline\G)$ 由那些在 $V\cap U$ 上的限制属于 $\Gamma(V\cap U,\G)$ 的截面 $s\in\Gamma(V,\F)$ 组成。因此 $\overline\G|U=\psi^*(\overline\G)=\G$；$\overline\G$ 是 $\F$ 的在 $U$ 上诱导 $\G$ 的最大子层。称它为 $\F|U$ 的子层 $\G$ 在 $\F$ 中的\emph{典范延拓}。
\egastatement{9.4.2}{命题} 设 $X$ 为预概形，$U$ 为 $X$ 的开子集，且典范嵌入 $j:U\to X$ 为拟紧态射；当 $X$ 的底空间局部诺特时，对每个 $U$ 都如此（\hyperlink{para:1:6.6.4}{6.6.4}，(i)）。则：
\begin{enumerate}[label=(\roman*)]
\item 对每个拟凝聚 $(\OO_X|U)$-模 $\G$，$j_*(\G)$ 为拟凝聚 $\OO_X$-模，且 $j_*(\G)|U=j^*(j_*(\G))=\G$。
\item 对每个拟凝聚 $\OO_X$-模 $\F$，以及 $\F|U$ 的每个拟凝聚子 $(\OO_X|U)$-模 $\G$，其典范延拓 $\overline\G$（\hyperlink{para:1:9.4.1}{9.4.1}）为 $\F$ 的拟凝聚子 $\OO_X$-模。
\end{enumerate}
若 $j=(\psi,\theta)$，其中 $\psi$ 为底空间的嵌入 $U\to X$，则对每个 $(\OO_X|U)$-模 $\G$，按定义 $j_*(\G)=\psi_*(\G)$。此外，对每个 $\OO_X$-模 $\HH$，由开子集上诱导预概形的定义，有 $j^*(\HH)=\psi^*(\HH)=\HH|U$。因此 (i) 为（\hyperlink{para:1:9.2.2}{9.2.2}，a)）的特例。出于同一理由，$j_*(j^*(\F))$ 拟凝聚；而 $\overline\G$ 为 $j_*(\G)$ 在同态 $\rho:\F\to j_*(j^*(\F))$ 下的原像，故 (ii) 由（\hyperlink{para:1:4.1.1}{4.1.1}）得出。

当 $U$ 拟紧且 $X$ 为概形时，嵌入 $j:U\to X$ 拟紧的假设也成立：此时 $U$ 为有限个仿射开集 $U_i$ 之并，而对 $X$ 的每个仿射开集 $V$，$V\cap U_i$ 仿射（\hyperlink{para:1:5.5.6}{5.5.6}），因而拟紧。
\egastatement{9.4.3}{推论} 设 $X$ 为预概形，$U$ 为 $X$ 的拟紧开集，且嵌入态射 $j:U\to X$ 拟紧。再假设每个拟凝聚 $\OO_X$-模都是其有限型拟凝聚子 $\OO_X$-模的归纳极限；当 $X$ 仿射时，此条件成立。设 $\F$ 为拟凝聚 $\OO_X$-模，$\G$ 为 $\F|U$ 的有限型拟凝聚子 $(\OO_X|U)$-模。那么，存在 $\F$ 的有限型拟凝聚子 $\OO_X$-模 $\G'$，使 $\G'|U=\G$。

事实上，$\G=\overline\G|U$，而由（\hyperlink{para:1:9.4.2}{9.4.2}），$\overline\G$ 拟凝聚，故为其有限型拟凝聚子 $\OO_X$-模 $\HH_\lambda$ 的归纳极限。于是 $\G$ 为各 $\HH_\lambda|U$ 的归纳极限；因其为有限型，$U$ 又拟紧，故它等于其中某个 $\HH_\lambda|U$（0，\hyperlink{para:0:5.2.3}{5.2.3}）。
\egastatement{9.4.4}{注} 假设对每个仿射开集 $U\subset X$，嵌入态射 $U\to X$ 都拟紧。若（\hyperlink{para:1:9.4.3}{9.4.3}）的结论对每个仿射开集 $U$ 及 $\F|U$ 的每个有限型拟凝聚子 $(\OO_X|U)$-模 $\G$ 都成立，则 $\F$ 是其有限型拟凝聚子 $\OO_X$-模的归纳极限。

事实上，对每个仿射开集 $U\subset X$，有 $\F|U=\wt M$，其中 $M$ 为 $A(U)$-模。$M$ 为其有限型子模的归纳极限，因此 $\F|U$ 为其有限型拟凝聚子 $(\OO_X|U)$-模的归纳极限（\hyperlink{para:1:1.3.9}{1.3.9}）。按假设，其中每个子模层都由 $\F$ 的某个有限型拟凝聚子 $\OO_X$-模 $\G_{\lambda,U}$ 在 $U$ 上诱导。这些 $\G_{\lambda,U}$ 的有限和仍为有限型拟凝聚 $\OO_X$-模：问题是局部的，仿射情形已在（\hyperlink{para:1:1.3.10}{1.3.10}）中处理。显然 $\F$ 是这些有限和的归纳极限，断言得证。
\egastatement{9.4.5}{推论} 在（\hyperlink{para:1:9.4.3}{9.4.3}）的假设下，对每个有限型拟凝聚 $(\OO_X|U)$-模 $\G$，存在有限型拟凝聚 $\OO_X$-模 $\G'$，使 $\G'|U=\G$。

因 $\F=j_*(\G)$ 拟凝聚（\hyperlink{para:1:9.4.2}{9.4.2}），且 $\F|U=\G$，只须将（\hyperlink{para:1:9.4.3}{9.4.3}）用于 $\F$。
\egastatement{9.4.6}{引理} 设 $X$ 为预概形，$L$ 为良序集合，$(V_\lambda)_{\lambda\in L}$ 为 $X$ 的仿射开覆盖，$U$ 为 $X$ 的开集；对每个 $\lambda\in L$，令 $W_\lambda=\bigcup_{\mu<\lambda}V_\mu$。假设：$1^\circ$ 对每个 $\lambda$，$V_\lambda\cap W_\lambda$ 拟紧；$2^\circ$ 浸入态射 $U\to X$ 拟紧。那么，对每个拟凝聚 $\OO_X$-模 $\F$，以及 $\F|U$ 的每个有限型拟凝聚子 $(\OO_X|U)$-模 $\G$，存在 $\F$ 的有限型拟凝聚子 $\OO_X$-模 $\G'$，使 $\G'|U=\G$。

令 $U_\lambda=U\cup W_\lambda$。用超限归纳定义族 $(\G_\lambda')$：$\G_\lambda'$ 是 $\F|U_\lambda$ 的有限型拟凝聚子 $(\OO_X|U_\lambda)$-模，满足 $\G_\lambda'|U_\mu=\G_\mu'$（$\mu<\lambda$）及 $\G_\lambda'|U=\G$。将这些层粘合所得的唯一子 $\OO_X$-模 $\G'$（0，\hyperlink{para:0:3.3.1}{3.3.1}），即满足要求。\transnote{在归纳的初始步取 $U_0=U$、$\G_0'=\G$；若原良序指标集有末项，还须做完该末项对应的后继步骤，再取覆盖全体 $X$ 的粘合。原文省略了这两项端点约定。}

设对 $\mu<\lambda$ 已定义 $\G_\mu'$ 并满足上述性质。若 $\lambda$ 为无前驱的非初始指标，取 $\G_\lambda'$ 为 $\F|U_\lambda$ 的唯一子 $(\OO_X|U_\lambda)$-模，使其在每个 $U_\mu$ 上限制为 $\G_\mu'$；这是允许的，因为此时各 $U_\mu$（$\mu<\lambda$）覆盖 $U_\lambda$。若 $\lambda=\mu+1$，则 $U_\lambda=U_\mu\cup V_\mu$，只须定义 $\F|V_\mu$ 的有限型拟凝聚子 $(\OO_X|V_\mu)$-模 $\G_\mu''$，使
\[
 \G_\mu''|(U_\mu\cap V_\mu)=\G_\mu'|(U_\mu\cap V_\mu).
\]
再取 $\G_\lambda'$ 为在 $U_\mu$ 上限制为 $\G_\mu'$、在 $V_\mu$ 上限制为 $\G_\mu''$ 的子模层（0，\hyperlink{para:0:3.3.1}{3.3.1}）。由于 $V_\mu$ 仿射，只要证明 $U_\mu\cap V_\mu$ 拟紧，$\G_\mu''$ 的存在性就由（\hyperlink{para:1:9.4.3}{9.4.3}）保证。而 $U_\mu\cap V_\mu$ 为 $U\cap V_\mu$ 与 $W_\mu\cap V_\mu$ 之并，按假设二者都拟紧。
\egastatement{9.4.7}{定理} 设 $X$ 为预概形，$U$ 为 $X$ 的开集，并假设下列条件之一成立：
\begin{enumerate}[label=\alph*)]
\item $X$ 的底空间局部诺特；
\item $X$ 为拟紧概形，且 $U$ 为拟紧开集。
\end{enumerate}
对每个拟凝聚 $\OO_X$-模 $\F$，以及 $\F|U$ 的每个有限型拟凝聚子 $(\OO_X|U)$-模 $\G$，存在 $\F$ 的有限型拟凝聚子 $\OO_X$-模 $\G'$，使 $\G'|U=\G$。

取 $X$ 的仿射开覆盖 $(V_\lambda)_{\lambda\in L}$，在 b) 中令 $L$ 有限。赋予 $L$ 一个良序结构，只须验证引理（\hyperlink{para:1:9.4.6}{9.4.6}）的条件。在 a) 中这显然成立，因为各 $V_\lambda$ 为诺特空间。在 b) 中，各 $V_\lambda\cap V_\mu$ 仿射（\hyperlink{para:1:5.5.6}{5.5.6}），因而拟紧；由于 $L$ 有限，$V_\lambda\cap W_\lambda$ 拟紧。定理得证。
\egastatement{9.4.8}{推论} 在（\hyperlink{para:1:9.4.7}{9.4.7}）的假设下，对每个有限型拟凝聚 $(\OO_X|U)$-模 $\G$，存在有限型拟凝聚 $\OO_X$-模 $\G'$，使 $\G'|U=\G$。

将（\hyperlink{para:1:9.4.7}{9.4.7}）用于 $\F=j_*(\G)$ 即可；由（\hyperlink{para:1:9.4.2}{9.4.2}），它拟凝聚，且 $\F|U=\G$。
\egastatement{9.4.9}{推论} 设 $X$ 为底空间局部诺特的预概形，或为拟紧概形。那么，每个拟凝聚 $\OO_X$-模都是其有限型拟凝聚子 $\OO_X$-模的归纳极限。

这由（\hyperlink{para:1:9.4.7}{9.4.7}）与注（\hyperlink{para:1:9.4.4}{9.4.4}）得出。
\egastatement{9.4.10}{推论} 在（\hyperlink{para:1:9.4.9}{9.4.9}）的假设下，若拟凝聚 $\OO_X$-模 $\F$ 的每个有限型拟凝聚子 $\OO_X$-模都由其在 $X$ 上的截面生成，则 $\F$ 也由其在 $X$ 上的截面生成。

取点 $x\in X$ 的仿射开邻域 $U$，并取 $\F$ 在 $U$ 上的一个截面 $s$。由 $s$ 生成的 $\F|U$ 的子模层 $\G$ 拟凝聚且为有限型，因此存在 $\F$ 的有限型拟凝聚子 $\OO_X$-模 $\G'$，使 $\G'|U=\G$（\hyperlink{para:1:9.4.7}{9.4.7}）。按假设，存在 $\G'$ 在 $X$ 上的有限个截面 $t_i$，以及 $\OO_X$ 在 $x$ 的某个邻域 $V\subset U$ 上的截面 $a_i$，使 $s|V=\sum_i a_i(t_i|V)$。推论得证。


% END chapters/ch1_s9_4.tex

\par

% BEGIN chapters/ch1_s9_5.tex
\subsection{预概形的闭像；子预概形的闭包}\hypertarget{heading:1:9.5}{}
\src{176—179}
\egastatement{9.5.1}{命题} 设 $f:X\to Y$ 为预概形态射，使 $f_*(\OO_X)$ 为拟凝聚 $\OO_Y$-模；当 $f$ 拟紧，且 $f$ 分离或 $X$ 局部诺特时即如此（\hyperlink{para:1:9.2.2}{9.2.2}）。那么，存在 $Y$ 的最小闭子预概形 $Y'$，使典范嵌入 $j:Y'\to Y$ 不小于 $f$；等价地，使 $X$ 的子预概形 $f^{-1}(Y')$ 与 $X$ 相同（\hyperlink{para:1:4.4.1}{4.4.1}）。\transnote{原文陈述漏写“闭”（\textit{fermé}）；以下证明、推论（\hyperlink{para:1:9.5.2}{9.5.2}）及定义（\hyperlink{para:1:9.5.3}{9.5.3}）均在闭子预概形中讨论最小性。这里补明“闭”，而不把结论解释为对所有子预概形的最小性。}

更精确地：
\egastatement{9.5.2}{推论} 在（\hyperlink{para:1:9.5.1}{9.5.1}）的条件下，设 $f=(\psi,\theta)$，$\J$ 为同态 $\theta:\OO_Y\to f_*(\OO_X)$ 的拟凝聚核。那么，由 $\J$ 定义的 $Y$ 的闭子预概形 $Y'$ 满足（\hyperlink{para:1:9.5.1}{9.5.1}）的条件。

由于函子 $\psi^*$ 正合，典范分解
\[
 \theta:\OO_Y\longrightarrow\OO_Y/\J
 \xrightarrow{\theta'}\psi_*(\OO_X)
\]
给出分解（0，\hyperref[eq:0:3.5.4.3]{3.5.4.3}）
\[
 \theta^\sharp:\psi^*(\OO_Y)
 \longrightarrow\psi^*(\OO_Y)/\psi^*(\J)
 \xrightarrow{\theta'^\sharp}\OO_X.
\]
对每个 $x\in X$，$\theta_x^\sharp$ 为局部同态，故 $\theta_x'^\sharp$ 也为局部同态。将 $\psi$ 视作从 $X$ 到 $Y'$ 的连续映射，记作 $\psi_0$；将限制
\[
 \theta_0=\theta'|Y':(\OO_Y/\J)|Y'
 \longrightarrow\psi_*(\OO_X)|Y'=(\psi_0)_*(\OO_X)
\]
与之配对，便得到预概形态射 $f_0=(\psi_0,\theta_0):X\to Y'$（\hyperlink{para:1:2.2.1}{2.2.1}），满足 $f=j\circ f_0$。\transnote{原文在这段证明中混用了 $X'$、$Y'$，并在下一段使用 $X''$。译文统一用 $Y'$、$Y''$ 表示 $Y$ 的子预概形；不改变它们的角色。}

若 $Y''$ 为 $Y$ 的另一闭子预概形，由拟凝聚理想层 $\J'$ 定义，且嵌入 $j':Y''\to Y$ 不小于 $f$，则首先 $Y''\supset\psi(X)$，因 $Y''$ 闭，故包含 $Y'$ 的底空间。\transnote{此处底空间包含关系使用了 $|Y'|=\overline{\psi(X)}$，该事实紧接着在（\hyperlink{para:1:9.5.4}{9.5.4}）独立证明；也可先用同态分解所得的 $\J'\subset\J$ 直接推出包含关系。}
此外，对每个 $y\in Y''$，$\theta_y$ 必分解为
\[
 \OO_y\longrightarrow\OO_y/\J'_y
 \longrightarrow(\psi_*(\OO_X))_y.
\]
按定义这蕴含 $\J'_y\subset\J_y$，因此 $Y'$ 是 $Y''$ 的闭子预概形（\hyperlink{para:1:4.1.10}{4.1.10}）。
\egastatement{9.5.3}{定义} 若存在 $Y$ 的最小闭子预概形 $Y'$，使典范嵌入 $j:Y'\to Y$ 不小于 $f$，就称 $Y'$ 为 $X$ 经态射 $f$ 所得的\emph{闭像预概形}。
\egastatement{9.5.4}{命题} 若 $f_*(\OO_X)$ 为拟凝聚 $\OO_Y$-模，则 $X$ 经 $f$ 所得闭像的底空间，是 $f(X)$ 在 $Y$ 中的闭包 $\overline{f(X)}$。

由于 $f_*(\OO_X)$ 的支集包含于 $\overline{f(X)}$，沿用（\hyperlink{para:1:9.5.2}{9.5.2}）的记号，对 $y\notin\overline{f(X)}$，有 $\J_y=\OO_y$；因此 $\OO_Y/\J$ 的支集包含于 $\overline{f(X)}$。此外，该支集为闭集，且包含 $f(X)$。事实上，若 $y\in f(X)$，则环 $(\psi_*(\OO_X))_y$ 的单位元不为零，它是截面
\[
 1\in\Gamma(X,\OO_X)=\Gamma(Y,\psi_*(\OO_X))
\]
在 $y$ 处的芽。它是 $\OO_y$ 的单位元经 $\theta_y$ 所得的像，故后者不属于 $\J_y$，从而 $\OO_y/\J_y\ne0$。证明完毕。
\egastatement{9.5.5}{命题（闭像的传递性）} 设 $f:X\to Y$、$g:Y\to Z$ 为两个预概形态射。假设 $X$ 经 $f$ 的闭像 $Y'$ 存在；令 $g'$ 为 $g$ 在 $Y'$ 上的限制，并假设 $Y'$ 经 $g'$ 的闭像 $Z'$ 存在。那么，$X$ 经 $g\circ f$ 的闭像存在，且等于 $Z'$。

只须证明（\hyperlink{para:1:9.5.1}{9.5.1}），$Z'$ 是 $Z$ 的满足以下条件的最小闭子预概形 $Z_1$：$X$ 的闭子预概形 $(g\circ f)^{-1}(Z_1)=f^{-1}(g^{-1}(Z_1))$（\hyperlink{para:1:4.4.2}{4.4.2}）等于 $X$。等价地，须证明 $Z'$ 是使嵌入 $g^{-1}(Z_1)\to Y$ 不小于 $f$ 的最小闭子预概形（\hyperlink{para:1:4.4.1}{4.4.1}）。由闭像 $Y'$ 的存在性，$Z_1$ 具有此性质，当且仅当 $g^{-1}(Z_1)$ 不小于 $Y'$，亦即令 $j:Y'\to Y$ 为嵌入后，有
\[
 j^{-1}(g^{-1}(Z_1))=g'^{-1}(Z_1)=Y'.
\]
按 $Z'$ 的定义，$Z'$ 正是满足此条件的最小闭子预概形。
\egastatement{9.5.6}{推论} 设 $f:X\to Y$ 为 $S$-态射，且 $Y$ 为 $X$ 经 $f$ 的闭像。设 $Z$ 为 $S$-概形。若两个 $S$-态射 $g_1,g_2:Y\to Z$ 满足 $g_1\circ f=g_2\circ f$，则 $g_1=g_2$。

令 $h=(g_1,g_2)_S:Y\to Z\times_S Z$。因对角线 $T=\Delta_Z(Z)$ 为 $Z\times_S Z$ 的闭子预概形，$Y'=h^{-1}(T)$ 为 $Y$ 的闭子预概形（\hyperlink{para:1:4.4.1}{4.4.1}）。令 $u=g_1\circ f=g_2\circ f$，则按积的定义 $h'=h\circ f=(u,u)_S=\Delta_Z\circ u$。因 $\Delta_Z^{-1}(T)=Z$，有 $h'^{-1}(T)=u^{-1}(Z)=X$，故 $f^{-1}(Y')=X$。由（\hyperlink{para:1:4.4.1}{4.4.1}），典范嵌入 $Y'\to Y$ 不小于 $f$，所以按假设 $Y'=Y$。因此 $h$ 分解为 $\Delta_Z\circ v$，其中 $v:Y\to Z$ 为态射（\hyperlink{para:1:4.4.1}{4.4.1}），从而 $g_1=g_2=v$。
\egastatement{9.5.7}{注} 若 $X,Y$ 为 $S$-概形，（\hyperlink{para:1:9.5.6}{9.5.6}）的含义是：当 $Y$ 为 $X$ 经 $f$ 的闭像时，$f$ 是 $S$-概形范畴中的满态射（T，1.1）。我们将在第 V 章证明其逆命题：若闭像 $Y'$ 存在，且 $f$ 为 $S$-概形的满态射，则必有 $Y'=Y$。
\egastatement{9.5.8}{命题} 假设满足（\hyperlink{para:1:9.5.1}{9.5.1}）的假设，$Y'$ 为 $X$ 经 $f$ 的闭像。对 $Y$ 的每个开集 $V$，令 $f_V:f^{-1}(V)\to V$ 为 $f$ 的限制。那么，$f^{-1}(V)$ 经 $f_V$ 在 $V$ 中的闭像存在，且等于 $Y'$ 在其开集 $V\cap Y'$ 上诱导的预概形；换言之，等于 $Y$ 的子预概形 $\inf(V,Y')$（\hyperlink{para:1:4.4.3}{4.4.3}）。

令 $X'=f^{-1}(V)$。由于 $\OO_{X'}$ 经 $f_V$ 的直像就是 $f_*(\OO_X)$ 在 $V$ 上的限制，同态 $\OO_V\to(f_V)_*(\OO_{X'})$ 的核 $\J'$ 显然是 $\J$ 在 $V$ 上的限制，命题随即得出。

此结果可解释为：闭像的形成与作为开浸入的底预概形扩张 $Y_1\to Y$ 可交换。第 IV 章将说明，若 $f$ 分离且拟紧，则对作为平坦态射的扩张 $Y_1\to Y$ 也有相同结论。
\egastatement{9.5.9}{命题} 设态射 $f:X\to Y$ 的闭像 $Y'$ 存在。
\begin{enumerate}[label=(\roman*)]
\item 若 $X$ 约化，则 $Y'$ 也约化。
\item 若满足（\hyperlink{para:1:9.5.1}{9.5.1}）的假设，且 $X$ 不可约（分别为整），则 $Y'$ 也不可约（分别为整）。
\end{enumerate}
按假设，$f$ 分解为 $X\xrightarrow{g}Y'\xrightarrow{j}Y$，其中 $j$ 为典范嵌入。若 $X$ 约化，则 $g$ 分解为 $X\to Y'_{\mathrm{red}}\xrightarrow{j'}Y'$，其中 $j'$ 为典范嵌入（\hyperlink{para:1:5.2.2}{5.2.2}）；由 $Y'$ 的定义，得到 $Y'_{\mathrm{red}}=Y'$。若还满足（\hyperlink{para:1:9.5.1}{9.5.1}）的条件，则由（\hyperlink{para:1:9.5.4}{9.5.4}），$f(X)$ 在 $Y'$ 中稠密；若 $X$ 不可约，则 $Y'$ 也不可约（0，\hyperlink{para:0:2.1.5}{2.1.5}）。关于整预概形的断言是前两项的合取。
\egastatement{9.5.10}{命题} 设 $Y$ 为预概形 $X$ 的子预概形，且典范嵌入 $j:Y\to X$ 拟紧。那么，存在 $X$ 的不小于 $Y$ 的最小闭子预概形 $\overline Y$；其底空间为 $Y$ 的底空间的闭包。$Y$ 的底空间在该闭包中开，且预概形 $Y$ 由 $\overline Y$ 在此开集上诱导。

将（\hyperlink{para:1:9.5.1}{9.5.1}）用于嵌入 $j$ 即可：它分离（\hyperlink{para:1:5.5.1}{5.5.1}），并按假设拟紧。因此（\hyperlink{para:1:9.5.1}{9.5.1}）给出 $\overline Y$ 的存在性，（\hyperlink{para:1:9.5.4}{9.5.4}）说明其底空间为 $Y$ 在 $X$ 中的闭包。因 $Y$ 在 $X$ 中局部闭，它在 $\overline Y$ 中开；最后一项断言由（\hyperlink{para:1:9.5.8}{9.5.8}）用于 $X$ 的一个使 $Y$ 在其中闭的开集 $V$ 得出。

沿用这些记号，若嵌入 $V\to X$ 拟紧，$\J$ 为 $\OO_X|V$ 中定义 $V$ 的闭子预概形 $Y$ 的拟凝聚理想层，则由（\hyperlink{para:1:9.5.1}{9.5.1}），定义 $\overline Y$ 的 $\OO_X$ 的拟凝聚理想层为 $\J$ 的典范延拓 $\overline\J$（\hyperlink{para:1:9.4.1}{9.4.1}），因为它显然是 $\OO_X$ 的在 $V$ 上诱导 $\J$ 的最大拟凝聚理想子层。
\egastatement{9.5.11}{推论} 在（\hyperlink{para:1:9.5.10}{9.5.10}）的假设下，$\OO_{\overline Y}$ 在 $\overline Y$ 的开集 $V$ 上的任一截面，若在 $V\cap Y$ 中为零，则该截面为零。

由（\hyperlink{para:1:9.5.8}{9.5.8}），可归结到 $V=\overline Y$ 的情形。$\OO_{\overline Y}$ 在 $\overline Y$ 上的截面，典范地对应于 $\overline Y\otimes_{\mathbf Z}\mathbf Z[T]$ 的 $\overline Y$-截面（\hyperlink{para:1:3.3.15}{3.3.15}），而后一个预概形在 $\overline Y$ 上分离，故本推论是（\hyperlink{para:1:9.5.6}{9.5.6}）的特例。

若存在 $X$ 的不小于子预概形 $Y$ 的最小闭子预概形 $\overline Y$，则在不致混淆时，称 $\overline Y$ 为 $Y$ 在 $X$ 中的\emph{闭包}。


% END chapters/ch1_s9_5.tex

\par

% BEGIN chapters/ch1_s9_6.tex
\subsection{拟凝聚代数层；结构层的变换}\hypertarget{heading:1:9.6}{}
\src{179—180}
\egastatement{9.6.1}{命题} 设 $X$ 为预概形，$\BB$ 为拟凝聚 $\OO_X$-代数（0，\hyperlink{para:0:5.1.3}{5.1.3}）。一个 $\BB$-模 $\F$ 在赋环空间 $(X,\BB)$ 上拟凝聚，当且仅当 $\F$ 作为 $\OO_X$-模拟凝聚。

问题是局部的，可假设 $X$ 为环 $A$ 的仿射概形；此时 $\BB=\wt B$，其中 $B$ 为 $A$-代数（\hyperlink{para:1:1.4.3}{1.4.3}）。若 $\F$ 在 $(X,\BB)$ 上拟凝聚，还可假设它是某个 $\BB$-同态 $\BB^{(I)}\to\BB^{(J)}$ 的余核。此同态也是 $\OO_X$-模间的 $\OO_X$-同态，而 $\BB^{(I)},\BB^{(J)}$ 为拟凝聚 $\OO_X$-模（\hyperlink{para:1:1.3.9}{1.3.9}，(ii)），故 $\F$ 也是拟凝聚 $\OO_X$-模（\hyperlink{para:1:1.3.9}{1.3.9}，(i)）。

反之，若 $\F$ 为拟凝聚 $\OO_X$-模，则 $\F=\wt M$，其中 $M$ 为 $B$-模（\hyperlink{para:1:1.4.3}{1.4.3}）。$M$ 同构于某个 $B$-同态 $B^{(I)}\to B^{(J)}$ 的余核，故 $\F$ 作为 $\BB$-模，同构于相应同态 $\BB^{(I)}\to\BB^{(J)}$ 的余核（\hyperlink{para:1:1.3.13}{1.3.13}）。证明完毕。

特别地，若 $\F,\G$ 为两个拟凝聚 $\BB$-模，则 $\F\otimes_{\BB}\G$ 为拟凝聚 $\BB$-模；若 $\F$ 还有有限表示，则 $\shHom_{\BB}(\F,\G)$ 也拟凝聚（\hyperlink{para:1:1.3.13}{1.3.13}）。
\egapar{9.6.2} 给定预概形 $X$，若对每个 $x\in X$，存在一个仿射开邻域 $U$，使 $B=\Gamma(U,\BB)$ 是 $A=\Gamma(U,\OO_X)$ 上的有限型代数，就称拟凝聚 $\OO_X$-代数 $\BB$ 为\emph{有限型}。此时 $\BB|U=\wt B$，而对每个 $f\in A$，诱导的 $(\OO_X|D(f))$-代数 $\BB|D(f)$ 为有限型：它同构于 $\widetilde{B_f}$，而 $B_f=B\otimes_A A_f$ 显然是有限型 $A_f$-代数。因各 $D(f)$ 构成 $U$ 的拓扑基，若 $\BB$ 为有限型拟凝聚 $\OO_X$-代数，则对 $X$ 的任一开集 $V$，$\BB|V$ 为有限型拟凝聚 $(\OO_X|V)$-代数。
\egastatement{9.6.3}{命题} 设 $X$ 为局部诺特预概形。则每个有限型拟凝聚 $\OO_X$-代数 $\BB$ 都是凝聚环层（0，\hyperlink{para:0:5.3.7}{5.3.7}）。

仍可限于 $X$ 为诺特环 $A$ 的仿射概形、$\BB=\wt B$ 且 $B$ 为有限型 $A$-代数的情形；于是 $B$ 是诺特环。须证明某个 $\BB$-同态 $\BB^m\to\BB$ 的核 $\N$ 为有限型 $\BB$-模。$\N$ 作为 $\BB$-模同构于 $\wt N$，其中 $N$ 为相应的 $B$-模同态 $B^m\to B$ 的核（\hyperlink{para:1:1.3.13}{1.3.13}）。由于 $B$ 诺特，$B^m$ 的子模 $N$ 为有限型 $B$-模，故存在像为 $N$ 的同态 $B^p\to B^m$。序列 $B^p\to B^m\to B$ 正合，因此相应序列 $\BB^p\to\BB^m\to\BB$ 正合（\hyperlink{para:1:1.3.5}{1.3.5}）。$\N$ 为 $\BB^p\to\BB^m$ 的像（\hyperlink{para:1:1.3.9}{1.3.9}，(i)），命题得证。
\egastatement{9.6.4}{推论} 在（\hyperlink{para:1:9.6.3}{9.6.3}）的假设下，$\BB$-模 $\F$ 凝聚，当且仅当它作为 $\OO_X$-模拟凝聚，且作为 $\BB$-模为有限型。若这些条件成立，且 $\G$ 为 $\F$ 的子 $\BB$-模或商 $\BB$-模，则 $\G$ 为凝聚 $\BB$-模，当且仅当 $\G$ 为拟凝聚 $\OO_X$-模。

由（\hyperlink{para:1:9.6.1}{9.6.1}），对 $\F$ 的条件显然必要。证明充分性。可限于以下情形：$X$ 为诺特环 $A$ 的仿射概形，$\BB=\wt B$，$B$ 为有限型 $A$-代数，$\F=\wt M$，$M$ 为 $B$-模，且存在满射 $\BB^m\to\F\to0$。此时有相应正合列 $B^m\to M\to0$，故 $M$ 为有限型 $B$-模；又因 $B$ 诺特，此同态的核 $P$ 为有限型 $B$-模。因此 $\F$ 为某个 $\BB$-同态 $\BB^p\to\BB^m$ 的余核（\hyperlink{para:1:1.3.13}{1.3.13}），因 $\BB$ 为凝聚环层，$\F$ 凝聚（0，\hyperlink{para:0:5.3.4}{5.3.4}）。同样的论证说明，$\F$ 的拟凝聚子 $\BB$-模（分别为商 $\BB$-模）是有限型，从而得到推论的第二部分。
\egastatement{9.6.5}{命题} 设 $X$ 为拟紧概形，或为底空间诺特的预概形。对每个有限型拟凝聚 $\OO_X$-代数 $\BB$，存在 $\BB$ 的有限型拟凝聚子 $\OO_X$-模 $\E$，使 $\E$ 生成 $\OO_X$-代数 $\BB$。

按假设，存在 $X$ 的有限仿射开覆盖 $(U_i)$，使 $B_i=\Gamma(U_i,\BB)$ 为 $A_i=\Gamma(U_i,\OO_X)$ 上的有限型代数。取生成 $A_i$-代数 $B_i$ 的有限型子 $A_i$-模 $E_i$。由（\hyperlink{para:1:9.4.7}{9.4.7}），存在 $\BB$ 的有限型拟凝聚子 $\OO_X$-模 $\E_i$，使 $\E_i|U_i=\wt E_i$。各 $\E_i$ 的和 $\E$ 显然满足要求。
\egastatement{9.6.6}{命题} 设 $X$ 为底空间局部诺特的预概形，或为拟紧概形。那么，每个拟凝聚 $\OO_X$-代数 $\BB$ 都是其有限型拟凝聚子 $\OO_X$-代数的归纳极限。

由（\hyperlink{para:1:9.4.9}{9.4.9}），$\BB$ 作为 $\OO_X$-模，是其有限型拟凝聚子 $\OO_X$-模的归纳极限。这些子模生成 $\BB$ 的有限型拟凝聚子 $\OO_X$-代数（\hyperlink{para:1:1.3.14}{1.3.14}），而 $\BB$ 当然也是这些子代数的归纳极限。


% END chapters/ch1_s9_6.tex

\par

% BEGIN chapters/ch1_s10_1.tex
\section{形式概形}\hypertarget{heading:1:10}{}
\subsection{仿射形式概形}\hypertarget{heading:1:10.1}{}
\src{180—182}
\egapar{10.1.1} 设 $A$ 为可容许拓扑环（0，\hyperlink{para:0:7.1.2}{7.1.2}）。对 $A$ 的每个定义理想 $J$，$\Spec(A/J)$ 等同于 $\Spec(A)$ 的闭子空间 $V(J)$（\hyperlink{para:1:1.1.11}{1.1.11}），即 $A$ 的开素理想组成的集合。这个拓扑空间不依赖于所取的定义理想 $J$，将它记作 $\frX$。取由定义理想组成的 $A$ 中零点的一个基本邻域系 $(J_\lambda)$。对每个 $\lambda$，令 $\OO_\lambda$ 为 $\Spec(A/J_\lambda)$ 的结构层；它由 $\widetilde{A/J_\lambda}$ 在 $\frX$ 上诱导，后者在 $\frX$ 之外为零。对 $J_\lambda\subset J_\mu$，典范同态 $A/J_\lambda\to A/J_\mu$ 定义环层同态 $u_{\mu\lambda}:\OO_\lambda\to\OO_\mu$（\hyperlink{para:1:1.6.1}{1.6.1}）；各 $\OO_\lambda$ 对这些同态构成一个射影系。

由于 $\frX$ 的拓扑有由拟紧开集组成的基，可为每个 $\OO_\lambda$ 配备相伴的伪离散拓扑环层结构（0，\hyperlink{para:0:3.8.1}{3.8.1}），其底层环层仍为 $\OO_\lambda$，并仍用 $\OO_\lambda$ 表示。它们仍构成拓扑环层的射影系（0，\hyperlink{para:0:3.8.2}{3.8.2}）。将该射影系的极限记作 $\OO_{\frX}$；这是 $\frX$ 上的拓扑环层。对每个拟紧开集 $U\subset\frX$，$\Gamma(U,\OO_{\frX})$ 因而为离散环系 $\Gamma(U,\OO_\lambda)$ 的拓扑射影极限（0，\hyperlink{para:0:3.2.6}{3.2.6}）。
\egastatement{10.1.2}{定义} 给定可容许拓扑环 $A$，称 $\Spec(A)$ 的由 $A$ 的开素理想组成的闭子空间 $\frX$ 为 $A$ 的\emph{形式谱}，记作 $\Spf(A)$。若一个拓扑赋环空间同构于形式谱 $\frX=\Spf(A)$，其拓扑环层 $\OO_{\frX}$ 为伪离散环层 $(\widetilde{A/J_\lambda})|\frX$ 的射影极限，其中 $J_\lambda$ 遍历 $A$ 的定义理想构成的有向集合，就称该拓扑赋环空间为\emph{仿射形式概形}。

谈到形式谱 $\frX=\Spf(A)$ 作为仿射形式概形时，总是指按上述方式定义的拓扑赋环空间 $(\frX,\OO_{\frX})$。

每个仿射概形 $X=\Spec(A)$ 都可通过将 $A$ 赋予离散拓扑，以唯一的方式视作仿射形式概形：此时对拟紧开集 $U$，拓扑环 $\Gamma(U,\OO_X)$ 是离散的；但对 $X$ 的一般开集，通常并非如此。
\egastatement{10.1.3}{命题} 若 $\frX=\Spf(A)$，其中 $A$ 为可容许环，则 $\Gamma(\frX,\OO_{\frX})$ 与 $A$ 拓扑同构。

因 $\frX$ 在 $\Spec(A)$ 中闭，它拟紧，所以 $\Gamma(\frX,\OO_{\frX})$ 拓扑同构于离散环 $\Gamma(\frX,\OO_\lambda)$ 的射影极限。而 $\Gamma(\frX,\OO_\lambda)$ 同构于 $A/J_\lambda$（\hyperlink{para:1:1.3.7}{1.3.7}）；由于 $A$ 分离且完备，它拓扑同构于 $\plim A/J_\lambda$（0，\hyperlink{para:0:7.2.1}{7.2.1}），命题得证。
\egastatement{10.1.4}{命题} 设 $A$ 为可容许环，$\frX=\Spf(A)$。对每个 $f\in A$，令 $\frD(f)=D(f)\cap\frX$；则拓扑赋环空间 $(\frD(f),\OO_{\frX}|\frD(f))$ 同构于仿射形式谱 $\Spf(A_{\{f\}})$（0，\hyperlink{para:0:7.6.15}{7.6.15}）。

对 $A$ 的每个定义理想 $J$，离散环 $S_f^{-1}A/S_f^{-1}J$ 典范地等同于 $A_{\{f\}}/J_{\{f\}}$（0，\hyperlink{para:0:7.6.9}{7.6.9}），所以拓扑空间 $\Spf(A_{\{f\}})$ 典范地等同于 $\frD(f)$（\hyperlink{para:1:1.2.5}{1.2.5} 与 \hyperlink{para:1:1.2.6}{1.2.6}）。此外，对 $\frX$ 的每个包含于 $\frD(f)$ 的拟紧开集 $U$，$\Gamma(U,\OO_\lambda)$ 等同于 $\Spec(S_f^{-1}A/S_f^{-1}J_\lambda)$ 的结构层在 $U$ 上的截面模（\hyperlink{para:1:1.3.6}{1.3.6}）。因此，令 $\frY=\Spf(A_{\{f\}})$，便有 $\Gamma(U,\OO_{\frX})$ 与 $\Gamma(U,\OO_{\frY})$ 的典范等同，命题得证。
\egapar{10.1.5} 将 $\Spf(A)$ 的结构层 $\OO_{\frX}$ 只看作不带拓扑的环层时，它在每个 $x\in\frX$ 处的茎，由（\hyperlink{para:1:10.1.4}{10.1.4}）等同于各 $A_{\{f\}}$（$f\notin\jj_x$）的归纳极限。因此，由（0，\hyperlink{para:0:7.6.17}{7.6.17} 与 \hyperlink{para:0:7.6.18}{7.6.18}）得：
\egastatement{10.1.6}{命题} 对每个 $x\in\frX=\Spf(A)$，茎 $\OO_x$ 为局部环，其剩余域同构于 $k(x)=A_x/\jj_xA_x$。若 $A$ 还为诺特进制环，则 $\OO_x$ 为诺特环。

由于 $k(x)$ 不为零，由此得出环层 $\OO_{\frX}$ 的支集就是 $\frX$。


% END chapters/ch1_s10_1.tex

\par

% BEGIN chapters/ch1_s10_2.tex
\subsection{仿射形式概形的态射}\hypertarget{heading:1:10.2}{}
\src{182—183}
\egapar{10.2.1} 设 $A,B$ 为两个可容许环，$\varphi:B\to A$ 为连续同态。连续映射 ${}^a\varphi:\Spec(A)\to\Spec(B)$（\hyperlink{para:1:1.2.1}{1.2.1}）将 $\frX=\Spf(A)$ 映入 $\frY=\Spf(B)$，因为 $A$ 的开素理想在 $\varphi$ 下的原像是 $B$ 的开素理想。另一方面，对每个 $g\in B$，$\varphi$ 定义一个连续同态
\[
 \Gamma(\frD(g),\OO_{\frY})
 \longrightarrow\Gamma(\frD(\varphi(g)),\OO_{\frX})
\]
（\hyperlink{para:1:10.1.4}{10.1.4}、\hyperlink{para:1:10.1.3}{10.1.3} 与 0，\hyperlink{para:0:7.6.7}{7.6.7}）。这些同态满足由 $g$ 到其倍数所对应的限制相容条件，且 $\frD(\varphi(g))={}^a\varphi^{-1}(\frD(g))$，故它们定义了拓扑环层的连续同态 $\OO_{\frY}\to{}^a\varphi_*(\OO_{\frX})$（0，\hyperlink{para:0:3.2.5}{3.2.5}），仍记作 $\wt\varphi$。于是得到拓扑赋环空间态射 $\Phi=({}^a\varphi,\wt\varphi):\frX\to\frY$。

作为不带拓扑的环层同态，$\wt\varphi$ 对每个 $x\in\frX$ 都定义茎上的同态 $\wt\varphi_x^\sharp:\OO_{{}^a\varphi(x)}\to\OO_x$。
\egastatement{10.2.2}{命题} 设 $A,B$ 为两个可容许拓扑环，$\frX=\Spf(A)$、$\frY=\Spf(B)$。拓扑赋环空间态射 $u=(\psi,\theta):\frX\to\frY$ 形如 $({}^a\varphi,\wt\varphi)$，其中 $\varphi:B\to A$ 为连续环同态，当且仅当对每个 $x\in\frX$，$\theta_x^\sharp:\OO_{\psi(x)}\to\OO_x$ 为局部同态。

必要性：令 $\mathfrak p=\jj_x\in\Spf(A)$，$\mathfrak q=\varphi^{-1}(\mathfrak p)$。若 $g\notin\mathfrak q$，则 $\varphi(g)\notin\mathfrak p$；由 $\varphi$ 导出的同态 $B_{\{g\}}\to A_{\{\varphi(g)\}}$（0，\hyperlink{para:0:7.6.7}{7.6.7}）显然将 $\mathfrak q_{\{g\}}$ 映入 $\mathfrak p_{\{\varphi(g)\}}$。取归纳极限，利用（\hyperlink{para:1:10.1.5}{10.1.5}）与（0，\hyperlink{para:0:7.6.17}{7.6.17}），得 $\wt\varphi_x^\sharp$ 为局部同态。

反之，设 $(\psi,\theta)$ 满足条件。由（\hyperlink{para:1:10.1.3}{10.1.3}），$\theta$ 定义连续环同态
\[
 \varphi=\Gamma(\theta):B=\Gamma(\frY,\OO_{\frY})
 \longrightarrow\Gamma(\frX,\OO_{\frX})=A.
\]
由 $\theta$ 的假设，截面 $\varphi(g)$ 在 $x$ 处的芽可逆，当且仅当 $g$ 在 $\psi(x)$ 处的芽可逆。而由（0，\hyperlink{para:0:7.6.17}{7.6.17}），$\OO_{\frX}$（分别为 $\OO_{\frY}$）的整体截面中，在 $x$（分别为 $\psi(x)$）处具有不可逆芽的，正是 $\jj_x$（分别为 $\jj_{\psi(x)}$）的元素。因此 ${}^a\varphi=\psi$。最后，对每个 $g\in B$，下图交换：
\[
\begin{tikzcd}[column sep=large]
 B=\Gamma(\frY,\OO_{\frY})\arrow[r,"\varphi"]\arrow[d]&\Gamma(\frX,\OO_{\frX})=A\arrow[d]\\
 B_{\{g\}}\arrow[r,"\Gamma(\theta_{\frD(g)})"]&A_{\{\varphi(g)\}}.
\end{tikzcd}
\]
下行的环分别是 $\Gamma(\frD(g),\OO_{\frY})$ 与 $\Gamma(\frD(\varphi(g)),\OO_{\frX})$。由完备分式环的泛性质（0，\hyperlink{para:0:7.6.6}{7.6.6}），对每个 $g$ 有 $\theta_{\frD(g)}=\wt\varphi_{\frD(g)}$，故 $\theta=\wt\varphi$（0，\hyperlink{para:0:3.2.5}{3.2.5}）。

称满足（\hyperlink{para:1:10.2.2}{10.2.2}）条件的拓扑赋环空间态射 $(\psi,\theta)$ 为\emph{仿射形式概形态射}。也可以说，函子 $A\mapsto\Spf(A)$ 与 $\frX\mapsto\Gamma(\frX,\OO_{\frX})$ 给出了可容许环范畴与仿射形式概形范畴的对偶范畴之间的等价（T，I，1.2）。
\egapar{10.2.3} 作为（\hyperlink{para:1:10.2.2}{10.2.2}）的特例，对 $f\in A$，$\frX$ 在 $\frD(f)$ 上诱导的仿射形式概形的典范嵌入，对应于典范连续同态 $A\to A_{\{f\}}$。在（\hyperlink{para:1:10.2.2}{10.2.2}）的假设下，取 $h\in B$，并取 $g\in A$ 为 $\varphi(h)$ 的倍数，则 $\psi(\frD(g))\subset\frD(h)$；$u$ 在 $\frD(g)$ 上的限制，作为 $\frD(g)$ 到 $\frD(h)$ 的态射，是使下图交换的唯一态射 $v$：
\[
\begin{tikzcd}
 \frD(g)\arrow[r,"v"]\arrow[d]&\frD(h)\arrow[d]\\
 \frX\arrow[r,"u"]&\frY.
\end{tikzcd}
\]
它对应于使下图交换的唯一连续同态 $\varphi':B_{\{h\}}\to A_{\{g\}}$（0，\hyperlink{para:0:7.6.7}{7.6.7}）：
\[
\begin{tikzcd}
 A\arrow[d]&B\arrow[l,"\varphi"']\arrow[d]\\
 A_{\{g\}}&B_{\{h\}}\arrow[l,"\varphi'"].
\end{tikzcd}
\]


% END chapters/ch1_s10_2.tex

\par

% BEGIN chapters/ch1_s10_3.tex
\subsection{仿射形式概形的定义理想}\hypertarget{heading:1:10.3}{}
\src{183—185}
\egapar{10.3.1} 设 $A$ 为可容许环，$J$ 为 $A$ 的开理想，$\frX=\Spf(A)$。以 $(J_\lambda)$ 表示 $A$ 中包含于 $J$ 的定义理想所组成的集合。则 $\wt J/\wt J_\lambda$ 是 $\wt A/\wt J_\lambda$ 的理想层。将这些层在 $\frX$ 上的限制的射影极限记作 $J^\Delta$；它等同于 $\OO_{\frX}$ 的理想层（0，\hyperlink{para:0:3.2.6}{3.2.6}）。对每个 $f\in A$，$\Gamma(\frD(f),J^\Delta)$ 为各 $S_f^{-1}J/S_f^{-1}J_\lambda$ 的射影极限，即等同于环 $A_{\{f\}}$ 的开理想 $J_{\{f\}}$（0，\hyperlink{para:0:7.6.9}{7.6.9}）；特别地，$\Gamma(\frX,J^\Delta)=J$。由于各 $\frD(f)$ 构成 $\frX$ 的拓扑基，得
\begin{equation}
 J^\Delta|\frD(f)=(J_{\{f\}})^\Delta.
\tag{10.3.1.1}\label{eq:1:10.3.1.1}
\end{equation}
\egapar{10.3.2} 沿用（\hyperlink{para:1:10.3.1}{10.3.1}）的记号。对每个 $f\in A$，从 $A_{\{f\}}=\Gamma(\frD(f),\OO_{\frX})$ 到
\[
 \Gamma(\frD(f),(\widetilde{A/J})|\frX)
 =S_f^{-1}A/S_f^{-1}J
\]
的典范映射为满射，其核为 $\Gamma(\frD(f),J^\Delta)=J_{\{f\}}$（0，\hyperlink{para:0:7.6.9}{7.6.9}）。因此，这些映射定义了从拓扑环层 $\OO_{\frX}$ 到伪离散环层 $(\widetilde{A/J})|\frX$ 的连续满同态，其核为 $J^\Delta$，称为\emph{典范同态}。这个同态也就是（\hyperlink{para:1:10.2.1}{10.2.1}）中的 $\wt\varphi$，其中 $\varphi:A\to A/J$ 是连续商同态。仿射形式概形态射 $({}^a\varphi,\wt\varphi):\Spec(A/J)\to\frX$ 也称为\emph{典范态射}。原文指出，${}^a\varphi$ 是 $\frX$ 到自身的恒等同胚。\transnote{本小节起初只假设 $J$ 是开理想；在此一般性下，$\Spec(A/J)$ 是 $\frX$ 的闭子空间，不一定等于 $\frX$。这里的“恒等同胚”适用于 $J$ 为定义理想的情形。另，原文此处称目标为“离散环层”，应结合（\hyperlink{para:1:10.1.1}{10.1.1}）的伪离散层约定理解，不能断言任意开集的截面环都离散。}
因此有典范同构
\begin{equation}
 \OO_{\frX}/J^\Delta\simeq(\widetilde{A/J})|\frX.
\tag{10.3.2.1}\label{eq:1:10.3.2.1}
\end{equation}
由 $\Gamma(\frX,J^\Delta)=J$，映射 $J\mapsto J^\Delta$ 显然严格保序。根据上述讨论，对 $J\subset J'$，层 $(J')^\Delta/J^\Delta$ 典范地同构于 $(\widetilde{J'/J})|\frX$，也记作 $(J'/J)^\Delta$。
\egapar{10.3.3} 继续沿用（\hyperlink{para:1:10.3.1}{10.3.1}）的假设与记号。若 $\OO_{\frX}$ 的一个理想层 $\J$ 满足：对每个 $x\in\frX$，存在形如 $\frD(f)$（$f\in A$）的开邻域，使 $\J|\frD(f)=K^\Delta$，其中 $K$ 为 $A_{\{f\}}$ 的定义理想，则称 $\J$ 为 $\frX$ 的\emph{定义理想层}，或简称 $\frX$ 的\emph{定义理想}。
\egastatement{10.3.4}{命题} 对每个 $f\in A$，$\frX$ 的每个定义理想都在 $\frD(f)$ 上诱导一个定义理想。

这由（\hyperref[eq:1:10.3.1.1]{10.3.1.1}）得出。
\egastatement{10.3.5}{命题} 若 $A$ 为可容许环，则 $\frX=\Spf(A)$ 的每个定义理想都形如 $J^\Delta$，其中 $J$ 为 $A$ 的唯一确定的定义理想。

设 $\J$ 为 $\frX$ 的定义理想。按假设及 $\frX$ 的拟紧性，可取有限个 $f_i\in A$，使各 $\frD(f_i)$ 覆盖 $\frX$，并且 $\J|\frD(f_i)=K_i^\Delta$，其中 $K_i$ 为 $A_{\{f_i\}}$ 的定义理想。对每个 $i$，存在 $A$ 的开理想 $H_i$，使 $(H_i)_{\{f_i\}}=K_i$（0，\hyperlink{para:0:7.6.9}{7.6.9}）。取 $A$ 的一个包含于所有 $H_i$ 的定义理想 $H$。

$\J/H^\Delta$ 在 $\Spec(A/H)$ 的结构层 $\widetilde{A/H}$ 中的典范像（\hyperlink{para:1:10.3.2}{10.3.2}），限制到 $\frD(f_i)$ 后与 $\widetilde{H_i/H}$ 的限制相同。因此，此典范像为 $\Spec(A/H)$ 上的拟凝聚层，故形如 $\widetilde{J/H}$，其中 $J$ 为 $A$ 中包含 $H$ 的理想（\hyperlink{para:1:1.4.1}{1.4.1}）。于是 $\J=J^\Delta$（\hyperlink{para:1:10.3.2}{10.3.2}）。此外，对每个 $i$，存在整数 $n_i$ 使 $K_i^{n_i}\subset H_{\{f_i\}}$。令 $n$ 为所有 $n_i$ 的最大值，则 $(\J/H^\Delta)^n=0$，从而（\hyperlink{para:1:10.3.2}{10.3.2}）$\widetilde{(J/H)^n}=0$，最终得到 $(J/H)^n=0$（\hyperlink{para:1:1.3.13}{1.3.13}）。因此 $J$ 为 $A$ 的定义理想（0，\hyperlink{para:0:7.1.4}{7.1.4}）。唯一性来自 $\Gamma(\frX,J^\Delta)=J$。
\egastatement{10.3.6}{命题} 设 $A$ 为进制环，$J$ 为 $A$ 的定义理想，且 $J/J^2$ 为有限型 $A/J$-模。那么，对每个整数 $n>0$，有 $(J^\Delta)^n=(J^n)^\Delta$。

对每个 $f\in A$，由于 $J^n$ 为开理想，利用（\hyperref[eq:1:10.3.1.1]{10.3.1.1}）与（0，\hyperlink{para:0:7.6.12}{7.6.12}），有
\[
 (\Gamma(\frD(f),J^\Delta))^n
 =(J_{\{f\}})^n=(J^n)_{\{f\}}
 =\Gamma(\frD(f),(J^n)^\Delta).
\]
而 $(J^\Delta)^n$ 是预层 $U\mapsto(\Gamma(U,J^\Delta))^n$ 所伴随的层（0，\hyperlink{para:0:4.1.6}{4.1.6}）。因各 $\frD(f)$ 构成拓扑基，命题得证。
\egapar{10.3.7} 若定义理想族 $(\J_\lambda)$ 满足 $\frX$ 的每个定义理想都包含其中某个 $\J_\lambda$，就称该族为\emph{定义理想基本系}。由于 $\J_\lambda=J_\lambda^\Delta$，这等价于各 $J_\lambda$ 构成 $A$ 中零点的基本邻域系。取一族 $f_\alpha\in A$，使各 $\frD(f_\alpha)$ 覆盖 $\frX$。若 $\OO_{\frX}$ 的一个向下有向的理想层族 $(\J_\lambda)$ 满足：对每个 $\alpha$，各 $\J_\lambda|\frD(f_\alpha)$ 构成 $\frD(f_\alpha)$ 的定义理想基本系，那么 $(\J_\lambda)$ 就是 $\frX$ 的定义理想基本系。

事实上，对 $\frX$ 的每个定义理想 $\J$，存在由有限个 $\frD(f_i)$ 组成的覆盖，使对每个 $i$，某个 $\J_{\lambda_i}|\frD(f_i)$ 是包含于 $\J|\frD(f_i)$ 的定义理想。取指标 $\mu$，使 $\J_\mu\subset\J_{\lambda_i}$ 对所有 $i$ 成立。由（\hyperlink{para:1:10.3.3}{10.3.3}），$\J_\mu$ 是 $\frX$ 的定义理想，并显然包含于 $\J$，得到断言。


% END chapters/ch1_s10_3.tex

\par

% BEGIN chapters/ch1_s10_4.tex
\subsection{形式预概形与形式预概形态射}\hypertarget{heading:1:10.4}{}
\src{185—186}
\egapar{10.4.1} 给定拓扑赋环空间 $\frX$，若 $\frX$ 在开集 $U\subset\frX$ 上诱导的拓扑赋环空间为仿射形式概形，就称 $U$ 为\emph{仿射形式开集}。若其环还为进制环（分别为诺特进制环），则称 $U$ 为\emph{进制仿射形式开集}（分别为\emph{诺特仿射形式开集}）。
\egastatement{10.4.2}{定义} 若拓扑赋环空间 $\frX$ 的每一点都具有仿射形式开邻域，就称它为\emph{形式预概形}。若其每一点都具有进制仿射形式开邻域（分别为诺特仿射形式开邻域），则称 $\frX$ 为\emph{进制的}（分别为\emph{局部诺特的}）。若 $\frX$ 局部诺特，且其底空间拟紧，因而诺特，则称 $\frX$ \emph{诺特}。
\egastatement{10.4.3}{命题} 若 $\frX$ 为形式预概形（分别为局部诺特形式预概形），则仿射形式开集（分别为诺特仿射形式开集）构成 $\frX$ 的拓扑基。

这由（\hyperlink{para:1:10.4.2}{10.4.2}）与（\hyperlink{para:1:10.1.4}{10.1.4}）得出，同时用到：若 $A$ 为诺特进制环，则对每个 $f\in A$，$A_{\{f\}}$ 也为诺特进制环（0，\hyperlink{para:0:7.6.11}{7.6.11}）。
\egastatement{10.4.4}{推论} 若 $\frX$ 为形式预概形（分别为局部诺特形式预概形、诺特形式预概形），则它在任意开集上诱导的拓扑赋环空间仍为形式预概形（分别为局部诺特形式预概形、诺特形式预概形）。
\egastatement{10.4.5}{定义} 给定两个形式预概形 $\frX,\frY$，称拓扑赋环空间态射 $(\psi,\theta):\frX\to\frY$ 为一个\emph{形式预概形态射}，若对每个 $x\in\frX$，$\theta_x^\sharp:\OO_{\psi(x)}\to\OO_x$ 为局部同态。

两个形式预概形态射的复合显然仍为这样的态射，因此形式预概形构成一个范畴。将从 $\frX$ 到 $\frY$ 的态射集合记作 $\Hom(\frX,\frY)$。

若 $U$ 为 $\frX$ 的开子集，则在 $U$ 上诱导的形式预概形到 $\frX$ 的典范嵌入，为形式预概形态射，甚至是拓扑赋环空间的单态射（0，\hyperlink{para:0:4.1.1}{4.1.1}）。
\egastatement{10.4.6}{命题} 设 $\frX$ 为形式预概形，$\frS=\Spf(A)$ 为仿射形式概形。从 $\frX$ 到 $\frS$ 的形式预概形态射，与从环 $A$ 到拓扑环 $\Gamma(\frX,\OO_{\frX})$ 的连续同态之间，存在典范双射。

证明与（\hyperlink{para:1:2.2.4}{2.2.4}）相同，只须将“同态”改为“连续同态”，“仿射开集”改为“仿射形式开集”，并使用（\hyperlink{para:1:10.2.2}{10.2.2}）代替（\hyperlink{para:1:1.7.3}{1.7.3}）；细节留给读者。
\egapar{10.4.7} 给定形式预概形 $\frS$，若给定形式预概形 $\frX$ 及态射 $\varphi:\frX\to\frS$，就说它们定义了\emph{$\frS$ 上的形式预概形}，或\emph{$\frS$-形式预概形} $\frX$，并称 $\varphi$ 为其\emph{结构态射}。若 $\frS=\Spf(A)$，其中 $A$ 可容许，则也称 $\frX$ 为\emph{$A$-形式预概形}，或\emph{$A$ 上的形式预概形}。任意形式预概形总可视作离散环 $\mathbf Z$ 上的形式预概形。

若 $\frX,\frY$ 为两个 $\frS$-形式预概形，且态射 $u:\frX\to\frY$ 使下图交换，则称 $u$ 为 $\frS$-态射：
\[
\begin{tikzcd}
 \frX\arrow[rr,"u"]\arrow[dr]&&\frY\arrow[dl]\\
 &\frS&
\end{tikzcd}
\]
其中斜箭头为结构态射。按此定义，对固定的 $\frS$，$\frS$-形式预概形构成范畴。将从 $\frX$ 到 $\frY$ 的 $\frS$-态射集合记作 $\Hom_{\frS}(\frX,\frY)$。当 $\frS=\Spf(A)$ 时，也称 $A$-态射。
\egapar{10.4.8} 由于每个仿射概形可视作仿射形式概形（\hyperlink{para:1:10.1.2}{10.1.2}），每个通常的预概形也可视作形式预概形。此外，由（\hyperlink{para:1:10.4.5}{10.4.5}），对通常的预概形，形式预概形的态射（分别为 $S$-态射）与 §2 定义的态射（分别为 $S$-态射）相同。


% END chapters/ch1_s10_4.tex

\par

% BEGIN chapters/ch1_s10_5.tex
\subsection{形式预概形的定义理想}\hypertarget{heading:1:10.5}{}
\src{186—188}
\egapar{10.5.1}设 $\frX$ 为形式预概形。称 $\OO_{\frX}$ 的理想层 $\J$ 为 $\frX$ 的\emph{定义理想层}（或\emph{定义理想}），如果每个 $x\in\frX$ 都有一个仿射形式开邻域 $U$，使 $\J|U$ 是 $\frX$ 在 $U$ 上诱导的仿射形式概形的定义理想（\hyperlink{para:1:10.3.3}{10.3.3}）。由（\hyperref[eq:1:10.3.1.1]{10.3.1.1}）及（\hyperlink{para:1:10.4.3}{10.4.3}），对于每个开集 $V\subset\frX$，$\J|V$ 都是 $\frX$ 在 $V$ 上诱导的形式预概形的定义理想。

称 $\frX$ 的一族定义理想 $(\J_\lambda)$ 为\emph{定义理想基本系}，如果存在由仿射形式开集组成的覆盖 $(U_\alpha)$，使得对每个 $\alpha$，族 $(\J_\lambda|U_\alpha)$ 是 $\frX$ 在 $U_\alpha$ 上诱导的仿射形式概形的定义理想基本系（\hyperlink{para:1:10.3.7}{10.3.7}）。由（\hyperlink{para:1:10.3.7}{10.3.7}）末尾的说明，当 $\frX$ 是仿射形式概形时，这一定义与该处的定义一致。对于每个开集 $V\subset\frX$，族 $(\J_\lambda|V)$ 也是 $V$ 上诱导的形式预概形的定义理想基本系，见（\hyperref[eq:1:10.3.1.1]{10.3.1.1}）。若 $\frX$ 局部诺特，而 $\J$ 是其定义理想，则由（\hyperlink{para:1:10.3.6}{10.3.6}），各次幂 $\J^n$ 构成 $\frX$ 的定义理想基本系。

\egapar{10.5.2}设 $\frX$ 为形式预概形，$\J$ 为其定义理想。赋环空间 $(\frX,\OO_{\frX}/\J)$ 是一个通常意义的预概形；当 $\frX$ 为仿射形式概形（分别为局部诺特、诺特形式预概形）时，它是仿射的（分别为局部诺特、诺特的）。事实上，立即可归结到仿射情形，而这一情形已在（\hyperlink{para:1:10.3.2}{10.3.2}）中证明。此外，若 $\theta:\OO_{\frX}\to\OO_{\frX}/\J$ 为典范同态，则
\[
u=(1_{\frX},\theta):(\frX,\OO_{\frX}/\J)\longrightarrow(\frX,\OO_{\frX})
\]
是一个形式预概形态射，称为\emph{典范态射}。这也已在仿射情形中证明（\hyperlink{para:1:10.3.2}{10.3.2}），而一般情形立即归结于此。

\egastatement{10.5.3}{命题}设 $\frX$ 为形式预概形，$(\J_\lambda)$ 为其定义理想基本系。则拓扑环层 $\OO_{\frX}$ 是各伪离散环层 $\OO_{\frX}/\J_\lambda$ 的射影极限（0，\hyperlink{para:0:3.8.1}{3.8.1}）。

由于 $\frX$ 的拓扑具有由拟紧仿射形式开集组成的基（\hyperlink{para:1:10.4.3}{10.4.3}），可归结到仿射情形；此时命题是（\hyperlink{para:1:10.3.5}{10.3.5}）、（\hyperlink{para:1:10.3.2}{10.3.2}）和定义（\hyperlink{para:1:10.1.1}{10.1.1}）的结果。

不能确定每个形式预概形都具有定义理想。然而：

\egastatement{10.5.4}{命题}设 $\frX$ 为局部诺特形式预概形。存在 $\frX$ 的一个最大定义理想 $\mathscr T$；它是使预概形 $(\frX,\OO_{\frX}/\J)$ 约化的唯一一个定义理想 $\J$。若 $\J$ 是 $\frX$ 的任一定义理想，则 $\mathscr T$ 是 $\OO_{\frX}/\J$ 的幂零根在 $\OO_{\frX}\to\OO_{\frX}/\J$ 下的逆像。

先设 $\frX=\Spf(A)$，其中 $A$ 为诺特进制环。由 $A$ 的最大定义理想的存在性及其性质（0，\hyperlink{para:0:7.1.6}{7.1.6}及\hyperlink{para:0:7.1.7}{7.1.7}），结合（\hyperlink{para:1:10.3.4}{10.3.4}）与（\hyperlink{para:1:5.1.1}{5.1.1}），立即得到 $\mathscr T$ 的存在性及上述性质。

为证明一般情形，只须证明：若 $U\supset V$ 是 $\frX$ 的两个诺特仿射形式开集，则 $U$ 的最大定义理想 $\mathscr T_U$ 在 $V$ 上诱导 $V$ 的最大定义理想 $\mathscr T_V$。但
\[
\bigl(V,(\OO_{\frX}|V)/(\mathscr T_U|V)\bigr)
\]
是约化的，故此结论由前面的结果得出。

以 $\frX_{\mathrm{red}}$ 表示通常意义的约化预概形 $(\frX,\OO_{\frX}/\mathscr T)$。

\egastatement{10.5.5}{推论}设 $\frX$ 为局部诺特形式预概形，$\mathscr T$ 为其最大定义理想。对于 $\frX$ 的每个开集 $V$，$\mathscr T|V$ 是 $\frX$ 在 $V$ 上诱导的形式预概形的最大定义理想。

\egastatement{10.5.6}{命题}设 $\frX,\frY$ 为两个形式预概形，$\J$（分别为 $\mathscr K$）为 $\frX$（分别为 $\frY$）的定义理想，$f:\frX\to\frY$ 为形式预概形态射。
\begin{enumerate}[label=(\roman*)]
\item 若 $f^*(\mathscr K)\OO_{\frX}\subset\J$，则存在唯一一个通常预概形的态射
\[
f':(\frX,\OO_{\frX}/\J)\longrightarrow(\frY,\OO_{\frY}/\mathscr K),
\]
使下图交换，其中竖直箭头是典范态射：
\begin{equation}\tag{10.5.6.1}\label{eq:1:10.5.6.1}
\begin{tikzcd}
(\frX,\OO_{\frX})\arrow[r,"f"] &(\frY,\OO_{\frY})\\
(\frX,\OO_{\frX}/\J)\arrow[u]\arrow[r,"f'"'] &(\frY,\OO_{\frY}/\mathscr K)\arrow[u]
\end{tikzcd}
\end{equation}
\item 设 $\frX=\Spf(A)$、$\frY=\Spf(B)$ 是仿射形式概形，$\J=J^\Delta$、$\mathscr K=K^\Delta$，其中 $J$（分别为 $K$）为 $A$（分别为 $B$）的定义理想，并设 $f=({}^{a}\varphi,\wt\varphi)$，其中 $\varphi:B\to A$ 为连续同态。则 $f^*(\mathscr K)\OO_{\frX}\subset\J$ 的充要条件是 $\varphi(K)\subset J$；此时 $f'=({}^{a}\varphi',\wt{\varphi'})$，其中 $\varphi':B/K\to A/J$ 是由 $\varphi$ 取商所得的同态。
\end{enumerate}

(i) 写 $f=(\psi,\theta)$。假设说明，$\psi^*(\OO_{\frY})$ 的理想层 $\psi^*(\mathscr K)$ 在 $\theta^\sharp:\psi^*(\OO_{\frY})\to\OO_{\frX}$ 下的像包含于 $\J$（0，\hyperlink{para:0:4.3.5}{4.3.5}）。取商，得到环层同态
\[
\omega:\psi^*(\OO_{\frY}/\mathscr K)
=\psi^*(\OO_{\frY})/\psi^*(\mathscr K)
\longrightarrow\OO_{\frX}/\J.
\]
又因每个 $\theta_x^\sharp$ 是局部同态，$\omega_x$ 也是局部同态。因此赋环空间态射 $(\psi,\omega^\flat)$ 是所求的唯一态射 $f'$（\hyperlink{para:1:2.2.1}{2.2.1}）。

(ii) 仿射形式预概形态射与连续环同态之间的典范函子对应（\hyperlink{para:1:10.2.2}{10.2.2}）表明，所考虑的情形中，关系 $f^*(\mathscr K)\OO_{\frX}\subset\J$ 蕴含 $f'=({}^{a}\varphi',\wt{\varphi'})$，其中 $\varphi'$ 是使下图交换的唯一同态：
\begin{equation}\tag{10.5.6.2}\label{eq:1:10.5.6.2}
\begin{tikzcd}
B\arrow[r,"\varphi"]\arrow[d]&A\arrow[d]\\
B/K\arrow[r,"\varphi'"']&A/J
\end{tikzcd}
\end{equation}
$\varphi'$ 的存在性蕴含 $\varphi(K)\subset J$。反之，若此条件成立，以 $\varphi'$ 表示使（\hyperref[eq:1:10.5.6.2]{10.5.6.2}）交换的唯一同态，并置 $f'=({}^{a}\varphi',\wt{\varphi'})$，则（\hyperref[eq:1:10.5.6.1]{10.5.6.1}）显然交换。考察分别对应于 $f,f'$ 的同态 $\psi^*(\OO_{\frY})\to\OO_{\frX}$ 及 $\psi^*(\OO_{\frY}/\mathscr K)\to\OO_{\frX}/\J$，即知这蕴含 $f^*(\mathscr K)\OO_{\frX}\subset\J$。

上述对应 $f\mapsto f'$ 显然具有函子性。


% END chapters/ch1_s10_5.tex

\par

% BEGIN chapters/ch1_s10_6.tex
\subsection{形式预概形作为预概形的归纳极限}\hypertarget{heading:1:10.6}{}
\src{188—193}
\egapar{10.6.1}设 $\frX$ 为形式预概形，$(\J_\lambda)$ 为其定义理想基本系；对每个 $\lambda$，以 $f_\lambda$ 表示典范态射 $(\frX,\OO_{\frX}/\J_\lambda)\to\frX$（\hyperlink{para:1:10.5.2}{10.5.2}）。对于 $\J_\mu\subset\J_\lambda$，典范同态 $\OO_{\frX}/\J_\mu\to\OO_{\frX}/\J_\lambda$ 定义通常预概形的典范态射
\[
f_{\mu\lambda}:(\frX,\OO_{\frX}/\J_\lambda)\longrightarrow(\frX,\OO_{\frX}/\J_\mu),
\qquad f_\lambda=f_\mu\circ f_{\mu\lambda}.
\]
因此，预概形 $X_\lambda=(\frX,\OO_{\frX}/\J_\lambda)$ 及态射 $f_{\mu\lambda}$ 构成形式预概形范畴中的一个归纳系（\hyperlink{para:1:10.4.8}{10.4.8}）。

\egastatement{10.6.2}{命题}沿用（\hyperlink{para:1:10.6.1}{10.6.1}）的记号，形式预概形 $\frX$ 与各态射 $f_\lambda$ 构成归纳系 $(X_\lambda,f_{\mu\lambda})$ 在形式预概形范畴中的归纳极限（T，I，1.8）。

设 $\frY$ 为形式预概形，对每个 $\lambda$ 给定态射 $g_\lambda=(\psi_\lambda,\theta_\lambda):X_\lambda\to\frY$，且对于 $\J_\mu\subset\J_\lambda$ 有 $g_\lambda=g_\mu\circ f_{\mu\lambda}$。后一条件及 $X_\lambda$ 的定义首先说明，所有 $\psi_\lambda$ 都是底空间之间同一个连续映射 $\psi:\frX\to\frY$。此外，各环层同态
\[
\theta_\lambda^\sharp:\psi^*(\OO_{\frY})\longrightarrow
\OO_{X_\lambda}=\OO_{\frX}/\J_\lambda
\]
构成射影系。取射影极限，得到同态 $\omega:\psi^*(\OO_{\frY})\to\plim\OO_{\frX}/\J_\lambda=\OO_{\frX}$，而赋环空间态射 $g=(\psi,\omega^\flat)$ 显然是使下图交换的唯一态射：
\begin{equation}\tag{10.6.2.1}\label{eq:1:10.6.2.1}
\begin{tikzcd}
X_\lambda\arrow[rr,"g_\lambda"]\arrow[dr,"f_\lambda"']&&\frY\\
&\frX\arrow[ur,"g"']&
\end{tikzcd}
\end{equation}
尚须证明 $g$ 是形式预概形态射。这一问题在 $\frX,\frY$ 上是局部的，因而可设 $\frX=\Spf(A)$、$\frY=\Spf(B)$，其中 $A,B$ 为可容许环，且 $\J_\lambda=J_\lambda^\Delta$，$(J_\lambda)$ 是 $A$ 的定义理想基本系（\hyperlink{para:1:10.3.5}{10.3.5}）。由于 $A=\plim A/J_\lambda$，根据（\hyperlink{para:1:10.2.2}{10.2.2}）的双射对应及射影极限的定义，存在使（\hyperref[eq:1:10.6.2.1]{10.6.2.1}）交换的仿射形式概形态射 $g$。但赋环空间态射的唯一性表明，它与证明开头定义的 $g$ 相同。

下一个命题在某些附加条件下，给出一个给定的通常预概形归纳系在形式预概形范畴中归纳极限的存在性。

\egastatement{10.6.3}{命题}设 $X$ 为拓扑空间，$(\OO_i,u_{ji})$ 为以 $\mathbb N$ 为指标集的环层射影系，其中 $u_{ji}:\OO_i\to\OO_j$（$i\ge j$）。令 $\J_i$ 为 $u_{0i}:\OO_i\to\OO_0$ 的核。假设：
\begin{enumerate}[label=\alph*)]
\item 赋环空间 $(X,\OO_i)$ 是一个预概形 $X_i$；
\item 对每个 $x\in X$ 及每个 $i$，存在 $x$ 的开邻域 $U_i$，使 $\J_i|U_i$ 幂零；
\item 各同态 $u_{ji}$ 都是满射。
\end{enumerate}
令 $\OO_{\frX}$ 为各伪离散环层 $\OO_i$ 的拓扑环层射影极限，$u_i:\OO_{\frX}\to\OO_i$ 为典范同态。则拓扑赋环空间 $\frX=(X,\OO_{\frX})$ 是形式预概形；各 $u_i$ 是满射，其核 $\J^{(i)}$ 构成 $\frX$ 的定义理想基本系，且 $\J^{(0)}$ 是各理想层 $\J_i$ 的射影极限。

首先，在每个茎上 $u_{ji}$ 是满同态，因而是局部同态；故 $v_{ij}=(1_X,u_{ji}):X_j\to X_i$（$i\ge j$）是预概形态射（\hyperlink{para:1:2.2.1}{2.2.1}）。先设每个 $X_i$ 都是环为 $A_i$ 的仿射概形。存在环同态 $\varphi_{ji}:A_i\to A_j$，使 $u_{ji}=\wt{\varphi_{ji}}$（\hyperlink{para:1:1.7.3}{1.7.3}）。由（\hyperlink{para:1:1.6.3}{1.6.3}），$\OO_j$ 是 $X_i$ 上的拟凝聚 $\OO_i$-模层，其外乘由 $u_{ji}$ 定义，它与借助 $\varphi_{ji}$ 视为 $A_i$-模的 $A_j$ 相伴。

对 $f\in A_i$，置 $f'=\varphi_{ji}(f)$。由假设，$D(f)$ 与 $D(f')$ 在 $X$ 中是同一个开集，而对应于 $u_{ji}$ 的截面同态
\[
\Gamma(D(f),\OO_i)=(A_i)_f\longrightarrow
\Gamma(D(f'),\OO_j)=(A_j)_{f'}
\]
就是 $(\varphi_{ji})_f$（\hyperlink{para:1:1.6.1}{1.6.1}）。但将 $A_j$ 视为 $A_i$-模时，$(A_j)_f$ 就是 $(A_i)_f$-模 $(A_j)_{f'}$；所以，将 $\varphi_{ji}$ 视为 $A_i$-模同态时，仍有 $u_{ji}=\wt{\varphi_{ji}}$。由 $u_{ji}$ 满射即得 $\varphi_{ji}$ 满射（\hyperlink{para:1:1.3.9}{1.3.9}）；若 $J_{ji}$ 是其核，则 $u_{ji}$ 的核是拟凝聚 $\OO_i$-模层 $\wt{J_{ji}}$。特别地，$\J_i=\wt{J_i}$，其中 $J_i=\Ker(\varphi_{0i}:A_i\to A_0)$。

假设 b) 蕴含 $\J_i$ 幂零：由于 $X$ 拟紧，可用有限个开集 $U_k$ 覆盖 $X$，且 $(\J_i|U_k)^{n_k}=0$；取 $n=\max n_k$，即有 $\J_i^n=0$。于是 $J_i$ 幂零（\hyperlink{para:1:1.3.13}{1.3.13}）。环 $A=\plim A_i$ 因而可容许（0，\hyperlink{para:0:7.2.2}{7.2.2}），典范同态 $\varphi_i:A\to A_i$ 是满射，其核 $J^{(i)}$ 是 $J_{ik}$（$k\ge i$）的射影极限；各 $J^{(i)}$ 构成 $A$ 的零点邻域基本系。此时（\hyperlink{para:1:10.6.3}{10.6.3}）的各项结论由（\hyperlink{para:1:10.1.1}{10.1.1}）及（\hyperlink{para:1:10.3.2}{10.3.2}）得出，$(X,\OO_{\frX})$ 正是 $\Spf(A)$。

仍在这一特殊情形中，若 $f=(f_i)\in A=\plim A_i$，则所有 $D(f_i)$（$X_i$ 中的仿射开集）都与 $\frX$ 的开集 $\frD(f)$ 相同；$X_i$ 在其上诱导的预概形因而与 $\Spec((A_i)_{f_i})$ 相同。

在一般情形，先注意到，对 $X$ 的每个拟紧开集 $U$，每个 $\J_i|U$ 都幂零，证明与上述相同。现证明每个 $x\in X$ 都有一个同时为所有 $X_i$ 的仿射开集的开邻域 $U$。取 $U$ 为 $X_0$ 的仿射开邻域，并注意 $\OO_{X_0}=\OO_{X_i}/\J_i$。由于 $\J_i|U$ 幂零，由（\hyperlink{para:1:5.1.9}{5.1.9}）知 $U$ 也是每个 $X_i$ 的仿射开集。对于这样的 $U$，上面的仿射情形表明 $(U,\OO_{\frX}|U)$ 为形式预概形，各 $\J^{(i)}|U$ 构成其定义理想基本系，而 $\J^{(0)}|U$ 是各 $\J_i|U$ 的射影极限。结论得证。

\egastatement{10.6.4}{推论}再假设对于 $i\ge j$，$u_{ji}$ 的核是 $\J_i^{j+1}$，且 $\J_1/\J_1^2$ 是 $\OO_0=\OO_1/\J_1$ 上的有限型模层。则 $\frX$ 为进制形式预概形；若 $\J=\J^{(0)}$，$\J^{(n)}$ 为 $\OO_{\frX}\to\OO_n$ 的核，则 $\J^{(n)}=\J^{n+1}$，且 $\J/\J^2$ 与 $\J_1$ 同构。若 $X_0$ 还是局部诺特的（分别为诺特的），则 $\frX$ 也是局部诺特的（分别为诺特的）。

由于 $\frX$ 与 $X_0$ 底空间相同，问题是局部的，可设所有 $X_i$ 仿射。利用 $\J_i=\wt{J_i}$，立即归结到（0，\hyperlink{para:0:7.2.7}{7.2.7}及\hyperlink{para:0:7.2.8}{7.2.8}）的相应断言；此时 $J_1/J_1^2$ 是有限型 $A_0$-模（\hyperlink{para:1:1.3.9}{1.3.9}）。

特别地，\emph{每个局部诺特形式预概形} $\frX$ 都是满足（\hyperlink{para:1:10.6.3}{10.6.3}）与（\hyperlink{para:1:10.6.4}{10.6.4}）条件的一列通常的局部诺特预概形 $(X_n)$ 的归纳极限：只须取 $\frX$ 的一个定义理想 $\J$（\hyperlink{para:1:10.5.4}{10.5.4}），并置 $X_n=(\frX,\OO_{\frX}/\J^{n+1})$，见（\hyperlink{para:1:10.5.1}{10.5.1}）与（\hyperlink{para:1:10.6.2}{10.6.2}）。

\egastatement{10.6.5}{推论}设 $A$ 为可容许环。仿射形式概形 $\frX=\Spf(A)$ 为诺特的充要条件是 $A$ 为诺特进制环。

充分性显然。反之，设 $\frX$ 诺特，$J$ 为 $A$ 的定义理想，$\J=J^\Delta$ 为相应的定义理想。通常预概形 $X_n=(\frX,\OO_{\frX}/\J^{n+1})$ 仿射且诺特，故环 $A_n=A/J^{n+1}$ 诺特（\hyperlink{para:1:6.1.3}{6.1.3}），从而 $J/J^2$ 是有限型 $A/J$-模。由于各 $\J^n$ 构成定义理想基本系（\hyperlink{para:1:10.5.1}{10.5.1}），有 $\OO_{\frX}=\plim\OO_{\frX}/\J^n$（\hyperlink{para:1:10.5.3}{10.5.3}）。据此由（\hyperlink{para:1:10.1.3}{10.1.3}），$A$ 与 $\plim A/J^n$ 拓扑同构，因此是诺特进制环（0，\hyperlink{para:0:7.2.8}{7.2.8}）。

\egastatement{10.6.6}{注}沿用（\hyperlink{para:1:10.6.3}{10.6.3}）的记号，设 $\F_i$ 为 $\OO_i$-模层，并对 $i\ge j$ 给定 $v_{ij}$-态射 $\theta_{ji}:\F_i\to\F_j$，使 $\theta_{kj}\circ\theta_{ji}=\theta_{ki}$（$k\le j\le i$）。由于 $v_{ij}$ 的底连续映射为恒等映射，$\theta_{ji}$ 是空间 $X$ 上的阿贝尔群层同态。若 $\F$ 为阿贝尔群层射影系 $(\F_i)$ 的射影极限，则各 $\theta_{ji}$ 为 $v_{ij}$-态射这一事实，使得通过射影极限在 $\F$ 上定义 $\OO_{\frX}$-模层结构。赋予这一结构后，称 $\F$ 为 $\OO_i$-模层系 $(\F_i)$ 关于 $\theta_{ji}$ 的\emph{射影极限}。在 $v_{ij}^*(\F_i)=\F_j$ 且 $\theta_{ji}$ 为恒等态射的特殊情形中，简言之，称 $\F$ 为满足 $v_{ij}^*(\F_i)=\F_j$（$j\le i$）的系 $(\F_i)$ 的射影极限，而不再提及 $\theta_{ji}$。

\egapar{10.6.7}设 $\frX,\frY$ 为形式预概形，$\J,\mathscr K$ 分别为其定义理想，$f:\frX\to\frY$ 为满足 $f^*(\mathscr K)\OO_{\frX}\subset\J$ 的态射。则对于每个 $n>0$，
\[
f^*(\mathscr K^n)\OO_{\frX}
=\bigl(f^*(\mathscr K)\OO_{\frX}\bigr)^n\subset\J^n.
\]
故由 $f$ 可得通常预概形的态射 $f_n:X_n\to Y_n$（\hyperlink{para:1:10.5.6}{10.5.6}），其中
\[
X_n=(\frX,\OO_{\frX}/\J^{n+1}),\qquad
Y_n=(\frY,\OO_{\frY}/\mathscr K^{n+1}).
\]
由定义立即可知，对 $m\le n$，下图交换；换言之，$(f_n)$ 为态射的归纳系：
\begin{equation}\tag{10.6.7.1}\label{eq:1:10.6.7.1}
\begin{tikzcd}
X_m\arrow[r,"f_m"]\arrow[d]&Y_m\arrow[d]\\
X_n\arrow[r,"f_n"']&Y_n
\end{tikzcd}
\end{equation}

\egapar{10.6.8}反之，设 $(X_n)$（分别为 $(Y_n)$）为满足（\hyperlink{para:1:10.6.3}{10.6.3}）的 b)、c) 条件的通常预概形归纳系，$\frX$（分别为 $\frY$）为其归纳极限。由归纳极限的定义，任一构成归纳系的态射序列 $f_n:X_n\to Y_n$ 都有归纳极限 $f:\frX\to\frY$；它是使下图交换的唯一形式预概形态射：
\[
\begin{tikzcd}
X_n\arrow[r,"f_n"]\arrow[d]&Y_n\arrow[d]\\
\frX\arrow[r,"f"']&\frY
\end{tikzcd}
\]

\egastatement{10.6.9}{命题}设 $\frX,\frY$ 为两个局部诺特形式预概形，$\J,\mathscr K$ 分别为其定义理想。（\hyperlink{para:1:10.6.7}{10.6.7}）所定义的映射 $f\mapsto(f_n)$ 是从满足 $f^*(\mathscr K)\OO_{\frX}\subset\J$ 的态射 $f:\frX\to\frY$ 的集合，到使（\hyperref[eq:1:10.6.7.1]{10.6.7.1}）交换的态射序列 $(f_n)$ 的集合的双射。

若 $f$ 是这样的一个序列的归纳极限，须证明 $f^*(\mathscr K)\OO_{\frX}\subset\J$。问题在 $\frX,\frY$ 上是局部的，故可设 $\frX=\Spf(A)$、$\frY=\Spf(B)$ 仿射，$A,B$ 为诺特进制环，$\J=J^\Delta$、$\mathscr K=K^\Delta$，$J,K$ 分别为 $A,B$ 的定义理想。由（\hyperlink{para:1:10.3.6}{10.3.6}）及（\hyperlink{para:1:10.3.2}{10.3.2}），$X_n=\Spec(A_n)$、$Y_n=\Spec(B_n)$，其中 $A_n=A/J^{n+1}$、$B_n=B/K^{n+1}$。有 $f_n=({}^{a}\varphi_n,\wt{\varphi_n})$，各 $\varphi_n:B_n\to A_n$ 构成射影系，所以 $f=({}^{a}\varphi,\wt\varphi)$，其中 $\varphi=\plim\varphi_n$。（\hyperref[eq:1:10.6.7.1]{10.6.7.1}）在 $m=0$ 时的交换性给出
\[
\varphi_n(K/K^{n+1})\subset J/J^{n+1}
\]
对每个 $n$ 成立；取射影极限得 $\varphi(K)\subset J$，再由（\hyperlink{para:1:10.5.6}{10.5.6}，(ii)）得到所求包含关系。

\egastatement{10.6.10}{推论}设 $\frX,\frY$ 为两个局部诺特形式预概形，$\mathscr T$ 为 $\frX$ 的最大定义理想（\hyperlink{para:1:10.5.4}{10.5.4}）。
\begin{enumerate}[label=(\roman*)]
\item 对 $\frY$ 的每个定义理想 $\mathscr K$ 及每个态射 $f:\frX\to\frY$，都有 $f^*(\mathscr K)\OO_{\frX}\subset\mathscr T$。
\item $\Hom(\frX,\frY)$ 与使（\hyperref[eq:1:10.6.7.1]{10.6.7.1}）交换的态射序列 $(f_n)$ 的集合之间有典范双射，其中
\[
X_n=(\frX,\OO_{\frX}/\mathscr T^{n+1}),\qquad
Y_n=(\frY,\OO_{\frY}/\mathscr K^{n+1}).
\]
\end{enumerate}
(ii) 立即由 (i) 及（\hyperlink{para:1:10.6.9}{10.6.9}）得出。证明 (i) 可归结到 $\frX=\Spf(A)$、$\frY=\Spf(B)$，其中 $A,B$ 为诺特进制环，$\mathscr T=T^\Delta$、$\mathscr K=K^\Delta$，$T$ 为 $A$ 的最大定义理想，$K$ 为 $B$ 的定义理想。写 $f=({}^{a}\varphi,\wt\varphi)$，其中 $\varphi:B\to A$ 连续。$K$ 的元素拓扑幂零（0，\hyperlink{para:0:7.1.4}{7.1.4}，(ii)），故 $\varphi(K)$ 的元素也拓扑幂零；由于 $T$ 是 $A$ 中拓扑幂零元素的集合（0，\hyperlink{para:0:7.1.6}{7.1.6}），有 $\varphi(K)\subset T$。据（\hyperlink{para:1:10.5.6}{10.5.6}，(ii)）即得结论。

\egastatement{10.6.11}{推论}设 $\frS,\frX,\frY$ 为三个局部诺特形式预概形，态射 $f:\frX\to\frS$、$g:\frY\to\frS$ 使 $\frX,\frY$ 成为 $\frS$-形式预概形。设 $\J$、$\mathscr K$、$\LL$ 分别为 $\frS,\frX,\frY$ 的定义理想，并假设
\[
f^*(\J)\OO_{\frX}\subset\mathscr K,\qquad
g^*(\J)\OO_{\frY}=\LL.
\]
置 $S_n=(\frS,\OO_{\frS}/\J^{n+1})$、$X_n=(\frX,\OO_{\frX}/\mathscr K^{n+1})$、$Y_n=(\frY,\OO_{\frY}/\LL^{n+1})$。则 $\Hom_{\frS}(\frX,\frY)$ 与使（\hyperref[eq:1:10.6.7.1]{10.6.7.1}）交换的 $S_n$-态射序列 $u_n:X_n\to Y_n$ 的集合之间有典范双射。

对每个 $\frS$-态射 $u:\frX\to\frY$，由定义 $f=g\circ u$，故
\[
u^*(\LL)\OO_{\frX}
=u^*\bigl(g^*(\J)\OO_{\frY}\bigr)\OO_{\frX}
=f^*(\J)\OO_{\frX}\subset\mathscr K.
\]
推论遂由（\hyperlink{para:1:10.6.9}{10.6.9}）得出。

注意，当 $m\le n$ 时，给定态射 $f_n:X_n\to Y_n$ 就决定使（\hyperref[eq:1:10.6.7.1]{10.6.7.1}）交换的唯一态射 $f_m:X_m\to Y_m$，归结到仿射情形立即可见。因此定义了映射
\[
\varphi_{mn}:\Hom_{S_n}(X_n,Y_n)\longrightarrow\Hom_{S_m}(X_m,Y_m),
\]
各集合 $\Hom_{S_n}(X_n,Y_n)$ 关于这些映射构成射影系。（\hyperlink{para:1:10.6.11}{10.6.11}）也就是说，存在典范双射
\[
\Hom_{\frS}(\frX,\frY)\simeq\plim_n\Hom_{S_n}(X_n,Y_n).
\]


% END chapters/ch1_s10_6.tex

\par

% BEGIN chapters/ch1_s10_7.tex
\subsection{形式预概形的积}\hypertarget{heading:1:10.7}{}
\src{193—194}
\egapar{10.7.1}设 $\frS$ 为形式预概形。由于各 $\frS$-形式预概形构成一个范畴，可以定义 $\frS$-形式预概形的积。

\egastatement{10.7.2}{命题}设 $\frX=\Spf(B)$、$\frY=\Spf(C)$ 为仿射形式概形 $\frS=\Spf(A)$ 上的两个仿射形式概形。令 $\frZ=\Spf(B\ctensor_A C)$，$p_1,p_2$ 为分别对应于典范连续 $A$-同态 $\rho:B\to B\ctensor_A C$、$\sigma:C\to B\ctensor_A C$ 的 $\frS$-态射（\hyperlink{para:1:10.2.2}{10.2.2}）。则 $(\frZ,p_1,p_2)$ 是 $\frS$-仿射形式概形 $\frX,\frY$ 的一个积。

根据（\hyperlink{para:1:10.4.6}{10.4.6}），只须验证：对每个连续 $A$-同态 $\varphi:B\ctensor_A C\to D$，其中 $D$ 为同时是拓扑 $A$-代数的可容许环，使它对应于 $(\varphi\circ\rho,\varphi\circ\sigma)$，便定义了一个双射
\[
\Hom_A(B\ctensor_A C,D)\simeq\Hom_A(B,D)\times\Hom_A(C,D).
\]
这正是完备张量积的泛性质（0，\hyperlink{para:0:7.7.6}{7.7.6}）；这里的同态均为连续同态。

\egastatement{10.7.3}{命题}给定两个 $\frS$-形式预概形 $\frX,\frY$，积 $\frX\times_{\frS}\frY$ 存在。

证明与（\hyperlink{para:1:3.2.6}{3.2.6}）相同，只须将仿射概形（分别为仿射开集）换为仿射形式概形（分别为仿射形式开集），并将命题（\hyperlink{para:1:3.2.2}{3.2.2}）换为（\hyperlink{para:1:10.7.2}{10.7.2}）。预概形积的全部形式性质（\hyperlink{para:1:3.2.7}{3.2.7}、\hyperlink{para:1:3.2.8}{3.2.8}及\hyperlink{para:1:3.3.1}{3.3.1}—\hyperlink{para:1:3.3.12}{3.3.12}），对形式预概形的积均不加修改地成立。

\egapar{10.7.4}设 $\frS,\frX,\frY$ 为三个形式预概形，$f:\frX\to\frS$、$g:\frY\to\frS$ 为两个态射。假设它们分别具有以同一个集合 $I$ 为指标集的定义理想基本系 $(\J_\lambda)$、$(\mathscr K_\lambda)$、$(\LL_\lambda)$，且对每个 $\lambda$ 有
\[
f^*(\J_\lambda)\OO_{\frX}\subset\mathscr K_\lambda,\qquad
g^*(\J_\lambda)\OO_{\frY}\subset\LL_\lambda.
\]
置
\[
S_\lambda=(\frS,\OO_{\frS}/\J_\lambda),\quad
X_\lambda=(\frX,\OO_{\frX}/\mathscr K_\lambda),\quad
Y_\lambda=(\frY,\OO_{\frY}/\LL_\lambda).
\]
当 $\J_\mu\subset\J_\lambda$、$\mathscr K_\mu\subset\mathscr K_\lambda$、$\LL_\mu\subset\LL_\lambda$ 时，$S_\lambda$（分别为 $X_\lambda,Y_\lambda$）是 $S_\mu$（分别为 $X_\mu,Y_\mu$）的具有相同底空间的闭子预概形（\hyperlink{para:1:10.6.1}{10.6.1}）。由于 $S_\lambda\to S_\mu$ 为单态射，积 $X_\lambda\times_{S_\lambda}Y_\lambda$ 与 $X_\lambda\times_{S_\mu}Y_\lambda$ 首先相同（\hyperlink{para:1:3.2.4}{3.2.4}）；后者又与 $X_\mu\times_{S_\mu}Y_\mu$ 的具有相同底空间的一个闭子预概形相同（\hyperlink{para:1:4.3.1}{4.3.1}）。

此时，$\frX\times_{\frS}\frY$ 是各通常预概形 $X_\lambda\times_{S_\lambda}Y_\lambda$ 的归纳极限。事实上，如（\hyperlink{para:1:10.6.2}{10.6.2}）中一样，可归结到 $\frS,\frX,\frY$ 都为仿射形式概形的情形。根据（\hyperlink{para:1:10.5.6}{10.5.6}，(ii)）及关于定义理想基本系的假设，该断言立即由两个代数的完备张量积的定义得出（0，\hyperlink{para:0:7.7.1}{7.7.1}）。

再设 $\frZ$ 为 $\frS$-形式预概形，$(\M_\lambda)$ 为以 $I$ 为指标集的定义理想基本系，$u:\frZ\to\frX$、$v:\frZ\to\frY$ 为两个 $\frS$-态射，满足
\[
u^*(\mathscr K_\lambda)\OO_{\frZ}\subset\M_\lambda,\qquad
v^*(\LL_\lambda)\OO_{\frZ}\subset\M_\lambda.
\]
置 $Z_\lambda=(\frZ,\OO_{\frZ}/\M_\lambda)$，并以 $u_\lambda:Z_\lambda\to X_\lambda$、$v_\lambda:Z_\lambda\to Y_\lambda$ 表示由 $u,v$ 得到的 $S_\lambda$-态射（\hyperlink{para:1:10.5.6}{10.5.6}）。立即可验，$(u,v)_{\frS}$ 为 $(u_\lambda,v_\lambda)_{S_\lambda}$ 的归纳极限。

特别地，当 $\frS,\frX,\frY$ 都局部诺特时，可取某个定义理想的各次幂所组成的基本系（\hyperlink{para:1:10.5.1}{10.5.1}），因而适用上述讨论。但须注意，$\frX\times_{\frS}\frY$ 不一定局部诺特，另见（\hyperlink{para:1:10.13.5}{10.13.5}）。


% END chapters/ch1_s10_7.tex

\par

% BEGIN chapters/ch1_s10_8.tex
\subsection{预概形沿闭子集的形式完备化}\hypertarget{heading:1:10.8}{}
\src{194—198}
\egapar{10.8.1}设 $X$ 为通常意义的局部诺特预概形，$X'$ 为其底空间的闭子集。以 $\Phi$ 表示 $\OO_X$ 中使 $\Supp(\OO_X/\J)=X'$ 的凝聚理想层 $\J$ 的集合。$\Phi$ 非空（\hyperlink{para:1:5.2.1}{5.2.1}、\hyperlink{para:1:4.1.4}{4.1.4}及\hyperlink{para:1:6.1.1}{6.1.1}）；我们用关系 $\supset$ 将它排序。

\egastatement{10.8.2}{引理}有序集 $\Phi$ 是滤过的。若 $X$ 诺特，则对于每个 $\J_0\in\Phi$，其各次幂 $\J_0^n$（$n>0$）的集合在 $\Phi$ 中共尾。

事实上，若 $\J_1,\J_2\in\Phi$，置 $\J=\J_1\cap\J_2$。由于 $\OO_X$ 凝聚，$\J$ 也凝聚（\hyperlink{para:1:6.1.1}{6.1.1}及0，\hyperlink{para:1:5.3.4}{5.3.4}）；对于每个 $x\in X$，$\J_x=(\J_1)_x\cap(\J_2)_x$，故当 $x\notin X'$ 时 $\J_x=\OO_x$，而当 $x\in X'$ 时 $\J_x\ne\OO_x$，即 $\J\in\Phi$。另一方面，若 $X$ 诺特，$\J_0,\J\in\Phi$，则存在 $n>0$ 使 $\J_0^n(\OO_X/\J)=0$（\hyperlink{para:1:9.3.4}{9.3.4}），即 $\J_0^n\subset\J$。

\egapar{10.8.3}设 $\F$ 为凝聚 $\OO_X$-模层。对于每个 $\J\in\Phi$，$\F\otimes_{\OO_X}(\OO_X/\J)$ 为凝聚 $\OO_X$-模层（\hyperlink{para:1:9.1.1}{9.1.1}），其支集包含于 $X'$；通常将它与其在 $X'$ 上的限制等同。当 $\J$ 遍历 $\Phi$ 时，这些层构成阿贝尔群层的射影系。

\egastatement{10.8.4}{定义}给定局部诺特预概形 $X$ 的闭子集 $X'$ 及凝聚 $\OO_X$-模层 $\F$，称
\[
\plim_{\J\in\Phi}\bigl(\F\otimes_{\OO_X}(\OO_X/\J)\bigr)
\]
在 $X'$ 上的限制为 $\F$ \emph{沿 $X'$ 的完备化}，记为 $\F_{/X'}$，无混淆时也记为 $\widehat\F$。它在 $X'$ 上的截面称为 $\F$ 沿 $X'$ 的\emph{形式截面}。

显然，对每个开集 $U\subset X$，有 $(\F|U)_{/(U\cap X')}=\F_{/X'}|(U\cap X')$。通过射影极限，$(\OO_X)_{/X'}$ 成为环层，$\F_{/X'}$ 可视为其上的模层。又因 $X'$ 的拓扑有拟紧开集基，可将 $(\OO_X)_{/X'}$（分别为 $\F_{/X'}$）视为各伪离散环层 $\OO_X/\J$（分别为伪离散群层 $\F\otimes_{\OO_X}(\OO_X/\J)=\F/\J\F$）的射影极限，从而成为拓扑环层（分别为拓扑群层）。通过射影极限，$\F_{/X'}$ 因而成为拓扑 $(\OO_X)_{/X'}$-模层（0，\hyperlink{para:0:3.8.1}{3.8.1}及\hyperlink{para:0:3.8.2}{3.8.2}）。回忆，对每个拟紧开集 $U\subset X$，$\Gamma(U\cap X',(\OO_X)_{/X'})$（分别为 $\Gamma(U\cap X',\F_{/X'})$）是离散环 $\Gamma(U,\OO_X/\J)$（分别为离散群 $\Gamma(U,\F/\J\F)$）的射影极限。

若 $u:\F\to\G$ 为 $\OO_X$-模层同态，则典范地得到
\[
u_\J:\F\otimes_{\OO_X}(\OO_X/\J)\longrightarrow\G\otimes_{\OO_X}(\OO_X/\J)
\]
的射影系。取射影极限并限制到 $X'$，得到连续的 $(\OO_X)_{/X'}$-同态 $u_{/X'}:\F_{/X'}\to\G_{/X'}$，也记为 $\widehat u$，称为 $u$ \emph{沿 $X'$ 的完备化}。若 $v:\G\to\HH$ 为另一个同态，则显然 $(v\circ u)_{/X'}=v_{/X'}\circ u_{/X'}$。因此 $\F_{/X'}$ 是从凝聚 $\OO_X$-模层范畴到拓扑 $(\OO_X)_{/X'}$-模层范畴的加性协变函子。

\egastatement{10.8.5}{命题}$(\OO_X)_{/X'}$ 的支集为 $X'$。拓扑赋环空间 $(X',(\OO_X)_{/X'})$ 为局部诺特形式预概形；若 $\J\in\Phi$，则 $\J_{/X'}$ 为此形式预概形的定义理想。当 $X=\Spec(A)$ 仿射、$A$ 诺特、$\J=\wt J$ 且 $X'=V(J)$ 时，$(X',(\OO_X)_{/X'})$ 典范地与 $\Spf(\widehat A)$ 等同，其中 $\widehat A$ 是 $A$ 关于 $J$-预进制拓扑的分离完备化。

显然只须证明最后一项。由（0，\hyperlink{para:0:7.3.3}{7.3.3}），$J$ 的 $J$-预进制分离完备化 $\widehat J$ 与 $\widehat A$ 中的理想 $J\widehat A$ 等同；$\widehat A$ 为诺特 $\widehat J$-进制环，且 $\widehat A/\widehat J^{\,n}=A/J^n$。后一关系说明，$\widehat A$ 的开素理想是 $\widehat{\mathfrak p}$，其中 $\mathfrak p$ 为 $A$ 中包含 $J$ 的素理想，且其在 $A$ 中的逆像为 $\mathfrak p$，故 $\Spf(\widehat A)=X'$。又有 $\OO_X/\J^n=(A/J^n)^\sim$，所以命题立即由定义得出。

称这样定义的形式预概形为 $X$ \emph{沿 $X'$ 的完备化}，记为 $X_{/X'}$，无混淆时记为 $\widehat X$。取 $X'=X$ 时可以取 $\J=0$，故 $X_{/X}=X$。若 $U$ 为 $X$ 在一个开集上诱导的子预概形，则 $U_{/(U\cap X')}$ 显然典范地与 $X_{/X'}$ 在开集 $U\cap X'\subset X'$ 上诱导的形式子预概形等同。

\egastatement{10.8.6}{推论}通常预概形 $\widehat X_{\mathrm{red}}$ 是 $X$ 的底空间为 $X'$ 的唯一约化子预概形（\hyperlink{para:1:5.2.1}{5.2.1}）。$\widehat X$ 诺特的充要条件是 $\widehat X_{\mathrm{red}}$ 诺特；$X$ 诺特则足以保证这一点。

$\widehat X_{\mathrm{red}}$ 的确定是局部问题（\hyperlink{para:1:10.5.4}{10.5.4}），可再设 $X$ 仿射且环 $A$ 诺特。沿用（\hyperlink{para:1:10.8.5}{10.8.5}）的记号，$\widehat A$ 中拓扑幂零元素的理想 $T$ 是 $A/J$ 的幂零根在 $\widehat A\to\widehat A/\widehat J=A/J$ 下的逆像（0，\hyperlink{para:0:7.1.3}{7.1.3}），故 $\widehat A/T$ 同构于 $A/J$ 对其幂零根的商。第一项断言遂由（\hyperlink{para:1:10.5.4}{10.5.4}）与（\hyperlink{para:1:5.1.1}{5.1.1}）得出。若 $\widehat X_{\mathrm{red}}$ 诺特，则底空间 $X'$ 诺特，故 $X_n=(X',\OO_X/\J^{n+1})$ 诺特（\hyperlink{para:1:6.1.2}{6.1.2}），从而 $\widehat X$ 诺特（\hyperlink{para:1:10.6.4}{10.6.4}）；反向由（\hyperlink{para:1:6.1.2}{6.1.2}）立即得出。

\egapar{10.8.7}典范同态 $\OO_X\to\OO_X/\J$（$\J\in\Phi$）构成射影系；取射影极限，得到环层同态
\[
\theta:\OO_X\longrightarrow\psi_*((\OO_X)_{/X'})
=\plim_{\J\in\Phi}\OO_X/\J,
\]
其中 $\psi:X'\hookrightarrow X$ 为底空间的典范嵌入。以 $i$ 或 $i_X$ 表示典范赋环空间态射 $(\psi,\theta):X_{/X'}\to X$。

通过张量积，对每个凝聚 $\OO_X$-模层 $\F$，典范同态 $\OO_X\to\OO_X/\J$ 给出 $\OO_X$-模层同态 $\F\to\F\otimes_{\OO_X}(\OO_X/\J)$ 的射影系。取射影极限，便得到典范函子同态 $\gamma:\F\to\psi_*(\F_{/X'})$。

\egastatement{10.8.8}{命题}(i) 函子 $\F\mapsto\F_{/X'}$ 是正合的。

(ii) 函子同态 $\gamma^\sharp:i^*(\F)\to\F_{/X'}$ 是 $(\OO_X)_{/X'}$-模层的同构。

(i) 只须证明：若 $0\to\F'\to\F\to\F''\to0$ 为凝聚 $\OO_X$-模层的正合列，$U$ 为 $X$ 的一个仿射开集，环 $A$ 诺特，则
\[
0\to\Gamma(U\cap X',\F'_{/X'})\to\Gamma(U\cap X',\F_{/X'})
\to\Gamma(U\cap X',\F''_{/X'})\to0
\]
正合。此时 $\F|U=\wt M$、$\F'|U=\wt{M'}$、$\F''|U=\wt{M''}$，其中 $M,M',M''$ 为有限型 $A$-模，且 $0\to M'\to M\to M''\to0$ 正合（\hyperlink{para:1:1.5.1}{1.5.1}及\hyperlink{para:1:1.3.11}{1.3.11}）。取 $\J\in\Phi$，并令 $J\subset A$ 使 $\J|U=\wt J$。有
\[
\Gamma\bigl(U\cap X',\F\otimes_{\OO_X}(\OO_X/\J^n)\bigr)
=M\otimes_A(A/J^n)
\]
（\hyperlink{para:1:1.3.12}{1.3.12}），故由射影极限的定义，$\Gamma(U\cap X',\F_{/X'})=\plim_n M\otimes_A(A/J^n)=\widehat M$，即 $M$ 的 $J$-预进制分离完备化。同理，另两项为 $\widehat{M'}$、$\widehat{M''}$。所求断言由诺特环上有限型模的完备化函子的正合性得到（0，\hyperlink{para:0:7.3.3}{7.3.3}）。

(ii) 问题是局部的，可设有正合列 $\OO_X^p\to\OO_X^q\to\F\to0$（0，\hyperlink{para:0:5.3.2}{5.3.2}）。由于 $\gamma^\sharp$ 具有函子性，而 $i^*(\F)$ 与 $\F_{/X'}$ 都为右正合函子，得到具有正合行的交换图
\begin{equation}\tag{10.8.8.1}\label{eq:1:10.8.8.1}
\begin{tikzcd}[column sep=small]
i^*(\OO_X^p)\arrow[r]\arrow[d,"\gamma^\sharp"]&i^*(\OO_X^q)\arrow[r]\arrow[d,"\gamma^\sharp"]&i^*(\F)\arrow[r]\arrow[d,"\gamma^\sharp"]&0\\
(\OO_X^p)_{/X'}\arrow[r]&(\OO_X^q)_{/X'}\arrow[r]&\F_{/X'}\arrow[r]&0
\end{tikzcd}
\end{equation}
又因为这两个函子与有限直和可交换（0，\hyperlink{para:0:3.2.6}{3.2.6}及\hyperlink{para:0:4.3.2}{4.3.2}），可归结到 $\F=\OO_X$。此时 $i^*(\OO_X)=(\OO_X)_{/X'}=\OO_{\widehat X}$（0，\hyperlink{para:0:4.3.4}{4.3.4}），$\gamma^\sharp$ 为 $\OO_{\widehat X}$-模层同态。只须验证它将 $X'$ 的开集上的单位截面映为自身；这是显然的，并说明此时 $\gamma^\sharp$ 为恒等同态。

\egastatement{10.8.9}{推论}赋环空间态射 $i:X_{/X'}\to X$ 是平坦的。

这由（0，\hyperlink{para:0:6.7.3}{6.7.3}）及（\hyperlink{para:1:10.8.8}{10.8.8}，(i)）得出。

\egastatement{10.8.10}{推论}若 $\F,\G$ 为凝聚 $\OO_X$-模层，则存在关于 $\F,\G$ 的典范函子同构
\begin{align}
\F_{/X'}\otimes_{(\OO_X)_{/X'}}\G_{/X'}&\simeq(\F\otimes_{\OO_X}\G)_{/X'},\tag{10.8.10.1}\label{eq:1:10.8.10.1}\\
\bigl(\shHom_{\OO_X}(\F,\G)\bigr)_{/X'}&\simeq
\shHom_{(\OO_X)_{/X'}}(\F_{/X'},\G_{/X'}).\tag{10.8.10.2}\label{eq:1:10.8.10.2}
\end{align}
这由 $i^*(\F)$ 与 $\F_{/X'}$ 的典范等同得出。第一个同构的存在是对任意赋环空间态射成立的结果（0，\hyperref[eq:0:4.3.3.1]{4.3.3.1}）；第二个是对任意平坦态射成立的结果（0，\hyperlink{para:0:6.7.6}{6.7.6}），故由（\hyperlink{para:1:10.8.9}{10.8.9}）得出。

\egastatement{10.8.11}{命题}对于每个凝聚 $\OO_X$-模层 $\F$，由 $\F\to\F_{/X'}$ 得到的典范同态 $\Gamma(X,\F)\to\Gamma(X',\F_{/X'})$ 的核，由在 $X'$ 的某个邻域中为零的截面组成。

由 $\F_{/X'}$ 的定义，这种截面的典范像为零。反之，若 $s\in\Gamma(X,\F)$ 的典范像为零，只须证明每个 $x\in X'$ 在 $X$ 中都有一个使 $s$ 为零的邻域。可因此归结到 $X=\Spec(A)$、$A$ 诺特、$X'=V(J)$、$\F=\wt M$，其中 $M$ 为有限型 $A$-模。此时 $\Gamma(X',\F_{/X'})$ 为 $M$ 的 $J$-预进制分离完备化，截面同态即 $M\to\widehat M$。其核是被 $1+J$ 中的某个元素消去的元素的集合（0，\hyperlink{para:0:7.3.7}{7.3.7}），故存在 $f\in J$ 使 $(1+f)s=0$。对每个 $x\in X'$，$(1_x+f_x)s_x=0$；由于 $J_x\OO_x$ 包含于 $\OO_x$ 的极大理想，$1_x+f_x$ 可逆，所以 $s_x=0$。命题得证。

\egastatement{10.8.12}{推论}$\Supp(\F_{/X'})=\Supp(\F)\cap X'$。

$\F_{/X'}$ 显然为有限型 $(\OO_X)_{/X'}$-模层，见（\hyperlink{para:1:10.8.8}{10.8.8}，(ii)）及（0，\hyperlink{para:0:5.2.4}{5.2.4}）；故其支集闭（0，\hyperlink{para:0:5.2.2}{5.2.2}），并显然包含于 $\Supp(\F)\cap X'$。要证相等，立即归结到证明 $\Gamma(X',\F_{/X'})=0$ 蕴含 $\Supp(\F)\cap X'=\varnothing$；这由（\hyperlink{para:1:10.8.11}{10.8.11}）及（\hyperlink{para:1:1.4.1}{1.4.1}）得出。

\egastatement{10.8.13}{推论}设 $u:\F\to\G$ 为凝聚 $\OO_X$-模层的同态。$u_{/X'}:\F_{/X'}\to\G_{/X'}$ 为零的充要条件是 $u$ 在 $X'$ 的某个邻域中为零。

由（\hyperlink{para:1:10.8.8}{10.8.8}，(ii)），$u_{/X'}$ 与 $i^*(u)$ 等同。将 $u$ 视为 $\HH=\shHom_{\OO_X}(\F,\G)$ 在 $X$ 上的截面，则 $u_{/X'}$ 为 $i^*(\HH)=\HH_{/X'}$ 在 $X'$ 上与之典范对应的截面，见（\hyperref[eq:1:10.8.10.2]{10.8.10.2}）及（0，\hyperlink{para:0:4.4.6}{4.4.6}）。只须对凝聚 $\OO_X$-模层 $\HH$ 应用（\hyperlink{para:1:10.8.11}{10.8.11}）。

\egastatement{10.8.14}{推论}设 $u:\F\to\G$ 为凝聚 $\OO_X$-模层同态。$u_{/X'}$ 为单态射（分别为满态射）的充要条件是 $u$ 在 $X'$ 的某个邻域中为单态射（分别为满态射）。

设 $\mathscr T,\mathscr S$ 分别为 $u$ 的核与余核，故有正合列 $0\to\mathscr T\xrightarrow{v}\F\xrightarrow{u}\G\xrightarrow{w}\mathscr S\to0$。由（\hyperlink{para:1:10.8.8}{10.8.8}，(i)）得到正合列
\[
0\to\mathscr T_{/X'}\xrightarrow{v_{/X'}}\F_{/X'}
\xrightarrow{u_{/X'}}\G_{/X'}\xrightarrow{w_{/X'}}\mathscr S_{/X'}\to0.
\]
若 $u_{/X'}$ 为单态射（分别为满态射），则 $v_{/X'}=0$（分别为 $w_{/X'}=0$），故由（\hyperlink{para:1:10.8.13}{10.8.13}），在 $X'$ 的某个邻域中 $v=0$（分别为 $w=0$）。


% END chapters/ch1_s10_8.tex

\par

% BEGIN chapters/ch1_s10_9.tex
\subsection{将态射延拓到完备化}\hypertarget{heading:1:10.9}{}
\src{198—201}
\egapar{10.9.1}设 $X,Y$ 为两个通常意义的局部诺特预概形，$f:X\to Y$ 为态射，$X',Y'$ 分别为其底空间的闭子集，且 $f(X')\subset Y'$。设 $\J,\mathscr K$ 分别为 $\OO_X,\OO_Y$ 的理想层，使 $\Supp(\OO_X/\J)=X'$、$\Supp(\OO_Y/\mathscr K)=Y'$，且 $f^*(\mathscr K)\OO_X\subset\J$。这样的理想层总存在：例如可取 $\J$ 为定义底空间为 $X'$ 的子预概形的最大理想层（\hyperlink{para:1:5.2.1}{5.2.1}），此时由 $f(X')\subset Y'$ 即得该包含关系（\hyperlink{para:1:5.2.4}{5.2.4}）。

对于每个 $n>0$，有 $f^*(\mathscr K^n)\OO_X\subset\J^n$（0，\hyperlink{para:0:4.3.5}{4.3.5}）。因而置
\[
X_n=(X',(\OO_X/\J^{n+1})|X'),\qquad
Y_n=(Y',(\OO_Y/\mathscr K^{n+1})|Y'),
\]
可由 $f$ 得到态射 $f_n:X_n\to Y_n$；显然各 $f_n$ 构成归纳系。以 $\widehat f:X_{/X'}\to Y_{/Y'}$ 表示其归纳极限（\hyperlink{para:1:10.6.8}{10.6.8}），并略滥用术语，称它为 $f$ 到 $X,Y$ 分别沿 $X',Y'$ 的完备化的\emph{延拓}。

容易验证，这一态射与满足上述条件的 $\J,\mathscr K$ 的选择无关。只须在 $X,Y$ 为环分别是 $A,B$ 的诺特仿射概形时验证。此时 $\J=\wt J$、$\mathscr K=\wt K$，其中 $J\subset A$、$K\subset B$ 为理想，$f$ 对应于满足 $\varphi(K)\subset J$ 的环同态 $\varphi:B\to A$（\hyperlink{para:1:4.4.6}{4.4.6}及\hyperlink{para:1:1.7.4}{1.7.4}）。于是 $\widehat f$ 对应于连续同态 $\widehat\varphi:\widehat B\to\widehat A$（\hyperlink{para:1:10.2.2}{10.2.2}），其中 $\widehat A,\widehat B$ 分别为 $J$-与 $K$-预进制分离完备化（\hyperlink{para:1:10.6.8}{10.6.8}）。若以另一个理想层 $\J'=\wt{J'}$ 代替 $\J$，而仍使 $\Supp(\OO_X/\J')=X'$，则 $A$ 的 $J$-与 $J'$-预进制拓扑相同（\hyperlink{para:1:10.8.2}{10.8.2}）。

注意，按这一定义，对应于 $\widehat f$ 的底空间连续映射 $X'\to Y'$，就是 $f$ 在 $X'$ 上的限制。

\egapar{10.9.2}由上一定义立即可知，赋环空间态射的下图交换，其中竖直箭头为典范态射（\hyperlink{para:1:10.8.7}{10.8.7}）：
\[
\begin{tikzcd}
\widehat X\arrow[r,"\widehat f"]\arrow[d,"i_X"']&\widehat Y\arrow[d,"i_Y"]\\
X\arrow[r,"f"']&Y
\end{tikzcd}
\]

\egapar{10.9.3}设 $Z$ 为第三个预概形，$g:Y\to Z$ 为态射，$Z'\subset Z$ 为满足 $g(Y')\subset Z'$ 的闭子集。若 $\widehat g$ 表示 $g$ 沿 $Y',Z'$ 的完备化，则由（\hyperlink{para:1:10.9.1}{10.9.1}）立即得到 $\widehat{g\circ f}=\widehat g\circ\widehat f$。

\egastatement{10.9.4}{命题}设 $X,Y$ 为两个局部诺特 $S$-预概形，$Y$ 在 $S$ 上为有限型。设 $f,g:X\to Y$ 为满足 $f(X')\subset Y'$、$g(X')\subset Y'$ 的两个 $S$-态射。则 $\widehat f=\widehat g$ 的充要条件是 $f,g$ 在 $X'$ 的某个邻域中相同。

充分性显然，且不需要关于 $Y$ 的有限性假设。为证必要性，先注意到 $\widehat f=\widehat g$ 蕴含对每个 $x\in X'$ 有 $f(x)=g(x)$。问题是局部的，故可设 $X,Y$ 分别为 $x$ 与 $y=f(x)=g(x)$ 的仿射开邻域，环均诺特；又可设 $S$ 仿射，$\Gamma(Y,\OO_Y)$ 为有限型 $\Gamma(S,\OO_S)$-代数（\hyperlink{para:1:6.3.3}{6.3.3}）。于是 $f,g$ 对应于两个 $\Gamma(S,\OO_S)$-同态 $\rho,\sigma:\Gamma(Y,\OO_Y)\to\Gamma(X,\OO_X)$（\hyperlink{para:1:1.7.3}{1.7.3}）；由假设，它们到 $\Gamma(Y,\OO_Y)$ 的分离完备化的连续延拓相同。据（\hyperlink{para:1:10.8.11}{10.8.11}），对每个 $s\in\Gamma(Y,\OO_Y)$，$\rho(s),\sigma(s)$ 在 $X'$ 的某个依赖于 $s$ 的邻域中相同。由有限型假设，可得 $X'$ 的一个邻域 $V$，使得所有这些截面对在 $V$ 中同时相同。若 $h\in\Gamma(X,\OO_X)$ 使 $D(h)$ 为包含于 $V$ 的 $x$ 的邻域，则由前述事实及（\hyperlink{para:1:1.4.1}{1.4.1}，d)）知 $f,g$ 在 $D(h)$ 中相同。\transnote{本句原版将 $h$ 所属的截面环印为 $\Gamma(Y,\OO_X)$；这里按其用途——定义 $X$ 中的主开邻域——写为 $\Gamma(X,\OO_X)$。}

\egastatement{10.9.5}{命题}在（\hyperlink{para:1:10.9.1}{10.9.1}）的假设下，对于每个凝聚 $\OO_Y$-模层 $\G$，存在典范函子同构
\[
(f^*(\G))_{/X'}\simeq\widehat f^{\,*}(\G_{/Y'})
\]
（作为 $(\OO_X)_{/X'}$-模层）。

由（\hyperlink{para:1:10.8.8}{10.8.8}），将左端典范地与 $i_X^*(f^*(\G))$ 等同，右端与 $\widehat f^{\,*}(i_Y^*(\G))$ 等同，命题立即由（\hyperlink{para:1:10.9.2}{10.9.2}）的交换图得出。

\egapar{10.9.6}现在设 $\F$ 为凝聚 $\OO_X$-模层，$\G$ 为凝聚 $\OO_Y$-模层。若 $u:\G\to\F$ 为 $f$-态射，则相应有 $\OO_X$-同态 $u^\sharp:f^*(\G)\to\F$，进而通过完备化得到连续同态
\[
(u^\sharp)_{/X'}:(f^*(\G))_{/X'}\longrightarrow\F_{/X'}.
\]
由（\hyperlink{para:1:10.9.5}{10.9.5}），存在唯一的 $\widehat f$-态射 $v:\G_{/Y'}\to\F_{/X'}$，使 $v^\sharp=(u^\sharp)_{/X'}$。

将三元组 $(\F,X,X')$（$\F$ 为凝聚 $\OO_X$-模层，$X'$ 为 $X$ 的闭子集）视为一个范畴的对象；态射 $(\F,X,X')\to(\G,Y,Y')$ 由满足 $f(X')\subset Y'$ 的预概形态射 $f:X\to Y$ 和 $f$-态射 $u:\G\to\F$ 组成。于是可说，$(X_{/X'},\F_{/X'})$ 是该范畴上的函子，其值为由局部诺特形式预概形 $\frZ$ 与 $\OO_{\frZ}$-模层 $\HH$ 所成的二元组 $(\frZ,\HH)$；后一范畴的态射由形式预概形态射 $g$ 及一个 $g$-态射组成。

\egastatement{10.9.7}{命题}设 $S,X,Y$ 为三个局部诺特预概形，$g:X\to S$、$h:Y\to S$ 为态射，$S'\subset S$ 为闭子集，$X'\subset X$、$Y'\subset Y$ 为满足 $g(X')\subset S'$、$h(Y')\subset S'$ 的闭子集。令 $Z=X\times_S Y$，并假设 $Z$ 局部诺特；令 $Z'=p^{-1}(X')\cap q^{-1}(Y')$，其中 $p,q$ 为积的投影。则 $Z_{/Z'}$ 与 $S_{/S'}$-形式预概形的积
\[
X_{/X'}\times_{S_{/S'}}Y_{/Y'}
\]
相同，其中结构态射为 $\widehat g,\widehat h$，投影为 $\widehat p,\widehat q$。

问题在 $S,X,Y$ 上是局部的，故归结到 $S=\Spec(A)$、$X=\Spec(B)$、$Y=\Spec(C)$，$S'=V(J)$、$X'=V(K)$、$Y'=V(L)$，其中理想满足 $\varphi(J)\subset K$、$\psi(J)\subset L$；$\varphi:A\to B$、$\psi:A\to C$ 分别对应于 $g,h$。于是 $Z=\Spec(B\otimes_A C)$，而 $Z'=V(M)$，其中
\[
M=\Im(K\otimes_A C)+\Im(B\otimes_A L).
\]
若 $\widehat A,\widehat B,\widehat C$ 分别为 $J$-、$K$-、$L$-预进制分离完备化，则 $\widehat B\ctensor_{\widehat A}\widehat C$ 是 $B\otimes_A C$ 关于 $M$-预进制拓扑的分离完备化（0，\hyperlink{para:0:7.7.2}{7.7.2}），故由（\hyperlink{para:1:10.7.2}{10.7.2}）得到结论。

还须注意：若 $T$ 为局部诺特 $S$-预概形，$u:T\to X$、$v:T\to Y$ 为 $S$-态射，$T'\subset T$ 为满足 $u(T')\subset X'$、$v(T')\subset Y'$ 的闭子集，则延拓 $\widehat{(u,v)_S}$ 与 $(\widehat u,\widehat v)_{S_{/S'}}$ 等同。

\egastatement{10.9.8}{推论}设 $X,Y$ 为局部诺特 $S$-预概形，且 $X\times_S Y$ 局部诺特。设 $S'$ 为 $S$ 的闭子集，$X',Y'$ 分别为 $X,Y$ 中其在 $S$ 中的像包含于 $S'$ 的闭子集。对于每个满足 $f(X')\subset Y'$ 的 $S$-态射 $f:X\to Y$，图态射 $\Gamma_{\widehat f}$ 与 $f$ 的图态射的延拓 $\widehat{\Gamma_f}$ 等同。

\egastatement{10.9.9}{推论}设 $X,Y$ 为两个局部诺特预概形，$f:X\to Y$ 为态射，$Y'\subset Y$ 为闭子集，$X'=f^{-1}(Y')$。则 $X_{/X'}$ 借助下列交换图，与形式预概形的积 $X\times_Y Y_{/Y'}$ 等同：
\[
\begin{tikzcd}
X&X_{/X'}\arrow[l,"i_X"']\arrow[d,"\widehat f"]\\
Y\arrow[from=u,"f"']&Y_{/Y'}\arrow[l,"i_Y"]
\end{tikzcd}
\]
只须在（\hyperlink{para:1:10.9.7}{10.9.7}）中将 $S,S'$ 换为 $Y$，并将其中第一因子及其闭子集都换为 $X$。

\egastatement{10.9.10}{注}若 $X=X_1\amalg X_2$ 为和（\hyperlink{heading:1:3.1}{3.1}），$X'=X'_1\cup X'_2$，其中 $X'_i$ 为 $X_i$ 的闭子集（$i=1,2$），则立即可见
\[
X_{/X'}=(X_1)_{/X'_1}\amalg(X_2)_{/X'_2}.
\]


% END chapters/ch1_s10_9.tex

\par

% BEGIN chapters/ch1_s10_10.tex
\subsection{对仿射形式概形上凝聚层的应用}\hypertarget{heading:1:10.10}{}
\src{201—203}
\egapar{10.10.1}本小节中，$A$ 始终表示诺特进制环，$J$ 为 $A$ 的一个定义理想。令 $X=\Spec(A)$、$\frX=\Spf(A)$，后者与 $X$ 的闭子集 $V(J)$ 等同（\hyperlink{para:1:10.1.2}{10.1.2}）。由（\hyperlink{para:1:10.1.2}{10.1.2}）及（\hyperlink{para:1:10.8.4}{10.8.4}），仿射形式概形 $\frX$ 就是仿射概形 $X$ 沿其底空间的闭子集 $\frX$ 的完备化 $X_{/\frX}$。因此，对每个凝聚 $\OO_X$-模层 $\F$，有相应的有限型 $\OO_{\frX}$-模层 $\F_{/\frX}$，它还是拓扑环层 $\OO_{\frX}$ 上的拓扑模层。

每个凝聚 $\OO_X$-模层 $\F$ 都形如 $\wt M$，其中 $M$ 为有限型 $A$-模（\hyperlink{para:1:1.5.1}{1.5.1}）。置 $(\wt M)_{/\frX}=M^\Delta$。若 $u:M\to N$ 为有限型 $A$-模的同态，则相应有 $\wt u:\wt M\to\wt N$，进而有连续同态 $(\wt u)_{/\frX}:(\wt M)_{/\frX}\to(\wt N)_{/\frX}$，记为 $u^\Delta$。显然 $(v\circ u)^\Delta=v^\Delta\circ u^\Delta$。这样定义了从有限型 $A$-模范畴到有限型 $\OO_{\frX}$-模层范畴的加性协变函子 $M\mapsto M^\Delta$。当 $A$ 为离散环时，$M^\Delta=\wt M$。

\egastatement{10.10.2}{命题}(i) $M\mapsto M^\Delta$ 是正合函子，且存在 $A$-模的典范函子同构 $\Gamma(\frX,M^\Delta)\simeq M$。

(ii) 若 $M,N$ 为两个有限型 $A$-模，则存在典范函子同构
\begin{align}
(M\otimes_A N)^\Delta&\simeq M^\Delta\otimes_{\OO_{\frX}}N^\Delta,\tag{10.10.2.1}\label{eq:1:10.10.2.1}\\
(\Hom_A(M,N))^\Delta&\simeq\shHom_{\OO_{\frX}}(M^\Delta,N^\Delta).\tag{10.10.2.2}\label{eq:1:10.10.2.2}
\end{align}
(iii) 映射 $u\mapsto u^\Delta$ 是函子同构
\begin{equation}\tag{10.10.2.3}\label{eq:1:10.10.2.3}
\Hom_A(M,N)\simeq\Hom_{\OO_{\frX}}(M^\Delta,N^\Delta).
\end{equation}

$M^\Delta$ 的正合性来自 $\wt M$（\hyperlink{para:1:1.3.5}{1.3.5}）及 $\F_{/\frX}$（\hyperlink{para:1:10.8.8}{10.8.8}）的正合性。由定义，$\Gamma(\frX,M^\Delta)$ 为 $\Gamma(X,\wt M)=M$ 的 $J$-预进制分离完备化；但由于 $A$ 完备且 $M$ 有限型，$M$ 已分离且完备（0，\hyperlink{para:0:7.3.6}{7.3.6}），从而 (i) 得证。（\hyperref[eq:1:10.10.2.1]{10.10.2.1}）来自（\hyperlink{para:1:1.3.12}{1.3.12}，(i)）与（\hyperref[eq:1:10.8.10.1]{10.8.10.1}）的复合；（\hyperref[eq:1:10.10.2.2]{10.10.2.2}）来自（\hyperlink{para:1:1.3.12}{1.3.12}，(ii)）与（\hyperref[eq:1:10.8.10.2]{10.8.10.2}）的复合。最后，$\Hom_A(M,N)$ 为有限型 $A$-模，可对它应用 (i)，将 $\Gamma(\frX,(\Hom_A(M,N))^\Delta)$ 与 $\Hom_A(M,N)$ 等同，再利用（\hyperref[eq:1:10.10.2.2]{10.10.2.2}），便知（\hyperref[eq:1:10.10.2.3]{10.10.2.3}）为同构。

由（\hyperlink{para:1:10.10.2}{10.10.2}）可得一系列与由（\hyperlink{para:1:1.3.7}{1.3.7}）、（\hyperlink{para:1:1.3.12}{1.3.12}）所得结果类似的推论，我们留给读者表述。还须指出，对正合列 $0\to J\to A\to A/J\to0$ 应用正合函子 $M^\Delta$，结合（\hyperlink{para:1:10.3.2}{10.3.2}），可知此处记作 $J^\Delta$ 的 $\OO_{\frX}$-理想层，与（\hyperlink{para:1:10.3.1}{10.3.1}）中同样记号表示的理想层一致。

\egastatement{10.10.3}{命题}在（\hyperlink{para:1:10.10.1}{10.10.1}）的假设下，$\OO_{\frX}$ 为凝聚环层。

对 $f\in A$，$A_{\{f\}}$ 为诺特进制环（0，\hyperlink{para:0:7.6.11}{7.6.11}）。问题是局部的，故由（\hyperlink{para:1:10.1.4}{10.1.4}），只须证明同态 $v:\OO_{\frX}^n\to\OO_{\frX}$ 的核为有限型 $\OO_{\frX}$-模层。此时 $v=u^\Delta$，其中 $u:A^n\to A$ 为 $A$-同态（\hyperlink{para:1:10.10.2}{10.10.2}）。由于 $A$ 诺特，$\Ker u$ 有限型，故存在同态 $w:A^m\to A^n$ 使 $A^m\xrightarrow{w}A^n\xrightarrow{u}A$ 正合。由（\hyperlink{para:1:10.10.2}{10.10.2}），列 $\OO_{\frX}^m\xrightarrow{w^\Delta}\OO_{\frX}^n\xrightarrow{v}\OO_{\frX}$ 正合，因而 $\Ker v$ 有限型。

\egapar{10.10.4}沿用上述记号，置 $A_n=A/J^{n+1}$，并令
\[
X_n=\Spec(A_n)=(\frX,\OO_{\frX}/\J^{n+1}),\qquad\J=J^\Delta.
\]
对于 $m\le n$，以 $u_{mn}:X_m\to X_n$ 表示对应于典范同态 $A_n\to A_m$ 的预概形态射。形式概形 $\frX$ 是各 $X_n$ 关于 $u_{mn}$ 的归纳极限（\hyperlink{para:1:10.6.3}{10.6.3}）。

\egastatement{10.10.5}{命题}在（\hyperlink{para:1:10.10.1}{10.10.1}）的假设下，设 $\F$ 为 $\OO_{\frX}$-模层。下列条件等价：
\begin{enumerate}[label=\alph*)]
\item $\F$ 凝聚；
\item $\F$ 同构于一列凝聚 $\OO_{X_n}$-模层 $\F_n$ 的射影极限（\hyperlink{para:1:10.6.6}{10.6.6}），且 $u_{mn}^*(\F_n)=\F_m$；
\item 存在有限型 $A$-模 $M$，使 $\F\simeq M^\Delta$；$M$ 由（\hyperlink{para:1:10.10.2}{10.10.2}，(i)）确定到典范同构。
\end{enumerate}

先证 b) 蕴含 c)。有 $\F_n=\wt{M_n}$，其中 $M_n$ 为有限型 $A_n$-模。假设给出 $M_m=M_n\otimes_{A_n}A_m$（$m\le n$）（\hyperlink{para:1:1.6.5}{1.6.5}）。于是各 $M_n$ 关于典范双同态 $M_n\to M_m$ 构成射影系；由 $A_n$ 的定义，它满足（0，\hyperlink{para:0:7.2.9}{7.2.9}）的条件。因此 $M=\plim M_n$ 为有限型 $A$-模，且对每个 $n$，$M_n=M\otimes_A A_n$。故 $\F_n$ 由 $\wt M\otimes_{\OO_X}(\OO_X/\wt{J^{n+1}})$ 在 $X_n$ 上诱导，从而按定义（\hyperlink{para:1:10.8.4}{10.8.4}）有 $\F=M^\Delta$。

反之，c) 蕴含 b)。若 $u_n:X_n\to X$ 为浸入态射，则 $u_n^*(\wt M)=(M\otimes_A A_n)^\sim$ 由 $\wt M\otimes_{\OO_X}(\OO_X/\wt{J^{n+1}})$ 在 $X_n$ 上诱导，并由定义 $M^\Delta=\plim u_n^*(\wt M)$。由于 $u_m=u_n\circ u_{mn}$，层 $\F_n=u_n^*(\wt M)$ 满足 b)。

再证 c) 蕴含 a)。由定义 $\OO_{\frX}=A^\Delta$；由于 $M$ 为某个同态 $A^m\to A^n$ 的余核，据（\hyperlink{para:1:10.10.2}{10.10.2}），$M^\Delta$ 为相应同态 $\OO_{\frX}^m\to\OO_{\frX}^n$ 的余核。环层 $\OO_{\frX}$ 凝聚（\hyperlink{para:1:10.10.3}{10.10.3}），所以 $M^\Delta$ 也凝聚（0，\hyperlink{para:0:5.3.4}{5.3.4}）。

最后证 a) 蕴含 b)。视为 $\OO_{\frX}$-模层，有 $\OO_{X_n}=\OO_{\frX}/\J^{n+1}=A_n^\Delta$。令 $\F_n=\F\otimes_{\OO_{\frX}}\OO_{X_n}$，则它作为 $\OO_{\frX}$-模层凝聚（0，\hyperlink{para:0:5.3.5}{5.3.5}）；又由于 $\J^{n+1}$ 凝聚，它作为 $\OO_{X_n}$-模层也凝聚（0，\hyperlink{para:0:5.3.10}{5.3.10}）。记住 $X_m\to X_n$ 的底连续映射为恒等映射，立即可验 $u_{mn}^*(\F_n)=\F_m$。因此 $\G=\plim\F_n$ 为凝聚 $\OO_{\frX}$-模层，因为已知 b) 蕴含 a)。各典范同态 $\F\to\F_n$ 组成射影系，给出典范同态 $w:\F\to\G$。只须证明 $w$ 双射。

这一问题现在是局部的，可设 $\F$ 为同态 $\OO_{\frX}^m\to\OO_{\frX}^n$ 的余核。由（\hyperlink{para:1:10.10.2}{10.10.2}），此同态形如 $v^\Delta$，其中 $v:A^m\to A^n$，故 $\F\simeq M^\Delta$，$M=\Coker v$。于是
\[
\F_n=M^\Delta\otimes_{\OO_{\frX}}A_n^\Delta=(M\otimes_A A_n)^\Delta.
\]
$M\otimes_A A_n$ 的 $J$-进制拓扑为离散拓扑，所以作为 $\OO_{X_n}$-模层，$(M\otimes_A A_n)^\Delta=(M\otimes_A A_n)^\sim$。前面已知 $M^\Delta=\plim\F_n$，此时 $w$ 就是恒等同态。证毕。

\egastatement{10.10.6}{推论}若 $\F$ 满足（\hyperlink{para:1:10.10.5}{10.10.5}）的条件 b)，则射影系 $(\F_n)$ 同构于系 $(\F\otimes_{\OO_{\frX}}\OO_{X_n})$。

\egapar{10.10.7}现设 $A,B$ 为两个诺特进制环，$\varphi:B\to A$ 为连续同态。以 $J,K$ 分别表示 $A,B$ 的定义理想，且 $\varphi(K)\subset J$。置 $X=\Spec(A)$、$Y=\Spec(B)$、$\frX=\Spf(A)$、$\frY=\Spf(B)$。设 $f:X\to Y$ 为对应于 $\varphi$ 的预概形态射（\hyperlink{para:1:1.6.1}{1.6.1}），$\widehat f:\frX\to\frY$ 为其到完备化的延拓（\hyperlink{para:1:10.9.1}{10.9.1}）；后者也是对应于 $\varphi$ 的形式预概形态射（\hyperlink{para:1:10.2.2}{10.2.2}）。

\egastatement{10.10.8}{命题}对于每个有限型 $B$-模 $N$，存在 $\OO_{\frX}$-模层的典范函子同构
\[
\widehat f^{\,*}(N^\Delta)\simeq(N\otimes_B A)^\Delta.
\]

以 $i_X:\frX\to X$、$i_Y:\frY\to Y$ 表示典范态射，则在典范函子同构意义下，由（\hyperlink{para:1:10.8.8}{10.8.8}）及（\hyperlink{para:1:1.6.5}{1.6.5}）有
\[
N^\Delta=i_Y^*(\wt N),\qquad
(N\otimes_B A)^\Delta=i_X^*((N\otimes_B A)^\sim)=i_X^*(f^*(\wt N)).
\]
命题因而由（\hyperlink{para:1:10.9.2}{10.9.2}）的交换图得出。

\egastatement{10.10.9}{推论}对 $B$ 的每个理想 $\mathfrak b$，有
\[
\widehat f^{\,*}(\mathfrak b^\Delta)\OO_{\frX}=(\mathfrak b A)^\Delta.
\]

设 $j:\mathfrak b\hookrightarrow B$ 为典范嵌入，它对应于 $\OO_{\frY}$-模层的典范嵌入 $j^\Delta:\mathfrak b^\Delta\to\OO_{\frY}$。由定义，左端为同态 $\widehat f^{\,*}(j^\Delta):\widehat f^{\,*}(\mathfrak b^\Delta)\to\OO_{\frX}=\widehat f^{\,*}(\OO_{\frY})$ 的像。由（\hyperlink{para:1:10.10.8}{10.10.8}），这个同态与 $(j\otimes1)^\Delta:(\mathfrak b\otimes_B A)^\Delta\to(B\otimes_B A)^\Delta$ 等同。由于 $j\otimes1$ 的像是理想 $\mathfrak bA$，据（\hyperlink{para:1:10.10.2}{10.10.2}），其相伴层同态的像为 $(\mathfrak bA)^\Delta$，结论得证。


% END chapters/ch1_s10_10.tex

\par

% BEGIN chapters/ch1_s10_11.tex
\subsection{形式预概形上的凝聚层}\hypertarget{heading:1:10.11}{}
\src{204—205}
\egastatement{10.11.1}{命题}若 $\frX$ 为局部诺特形式预概形，则环层 $\OO_{\frX}$ 凝聚，且 $\frX$ 的每个定义理想层都凝聚。

问题是局部的，故归结到诺特仿射形式概形的情形；此时命题由（\hyperlink{para:1:10.10.3}{10.10.3}）及（\hyperlink{para:1:10.10.5}{10.10.5}）得出。

\egapar{10.11.2}设 $\frX$ 为局部诺特形式预概形，$\J$ 为其定义理想层，$X_n=(\frX,\OO_{\frX}/\J^{n+1})$ 为相应的通常局部诺特预概形。因此 $\frX$ 为序列 $(X_n)$ 关于典范态射 $u_{mn}:X_m\to X_n$ 的归纳极限（\hyperlink{para:1:10.6.3}{10.6.3}）。沿用这些记号，有：

\egastatement{10.11.3}{定理}$\OO_{\frX}$-模层 $\F$ 凝聚的充要条件是，它同构于一列凝聚 $\OO_{X_n}$-模层 $\F_n$ 的射影极限，且对 $m\le n$ 有 $u_{mn}^*(\F_n)=\F_m$（\hyperlink{para:1:10.6.6}{10.6.6}）。此时射影系 $(\F_n)$ 同构于系
\[
u_n^*(\F)=\F\otimes_{\OO_{\frX}}\OO_{X_n},
\]
其中 $u_n:X_n\to\frX$ 为典范态射。

问题是局部的，故归结到 $\frX$ 为诺特仿射形式概形的情形，此时定理是（\hyperlink{para:1:10.10.5}{10.10.5}）及（\hyperlink{para:1:10.10.6}{10.10.6}）的结果。因此可以说，\emph{给定一个凝聚 $\OO_{\frX}$-模层，等价于给定一个满足 $u_{mn}^*(\F_n)=\F_m$（$m\le n$）的凝聚 $\OO_{X_n}$-模层射影系。}

\egastatement{10.11.4}{推论}若 $\F,\G$ 为两个凝聚 $\OO_{\frX}$-模层，则沿用（\hyperlink{para:1:10.11.3}{10.11.3}）的记号，可以定义典范函子同构
\begin{equation}\tag{10.11.4.1}\label{eq:1:10.11.4.1}
\Hom_{\OO_{\frX}}(\F,\G)\simeq\plim_n\Hom_{\OO_{X_n}}(\F_n,\G_n).
\end{equation}
右端的射影极限是关于映射 $\theta_n\mapsto u_{mn}^*(\theta_n)$（$m\le n$）而言。同态（\hyperref[eq:1:10.11.4.1]{10.11.4.1}）使 $\theta\in\Hom_{\OO_{\frX}}(\F,\G)$ 对应于序列 $(u_n^*(\theta))$。由（\hyperlink{para:1:10.11.3}{10.11.3}），将一个相容系 $(\theta_n)$ 映至它在 $\Hom_{\OO_{\frX}}(\F,\G)$ 中的射影极限，立即定义了逆同态。

\egastatement{10.11.5}{推论}同态 $\theta:\F\to\G$ 为满射的充要条件是相应的 $\theta_0=u_0^*(\theta):\F_0\to\G_0$ 为满射。

问题是局部的，可设 $\frX=\Spf(A)$、$A$ 为诺特进制环，$\F=M^\Delta$、$\G=N^\Delta$、$\theta=u^\Delta$，其中 $M,N$ 为有限型 $A$-模，$u:M\to N$ 为同态。此时 $\theta_0=\wt{u_0}$，其中 $u_0=u\otimes1:M\otimes_A A/J\to N\otimes_A A/J$。由（\hyperlink{para:1:1.3.9}{1.3.9}）及（\hyperlink{para:1:10.10.2}{10.10.2}），$\theta,u$ 同时满射，$\theta_0,u_0$ 同时满射；而 $u,u_0$ 同时满射（0，\hyperlink{para:0:7.1.14}{7.1.14}），故得结论。

\egapar{10.11.6}定理（\hyperlink{para:1:10.11.3}{10.11.3}）说明，将每个凝聚 $\OO_{\frX}$-模层 $\F$ 视为各伪离散群层 $\F_n$ 的射影极限（0，\hyperlink{para:0:3.8.1}{3.8.1}），可将它视为拓扑 $\OO_{\frX}$-模层。于是由（\hyperlink{para:1:10.11.4}{10.11.4}），任意凝聚 $\OO_{\frX}$-模层的同态 $u:\F\to\G$ 都自动连续（0，\hyperlink{para:0:3.8.2}{3.8.2}）。此外，若 $\HH$ 为凝聚 $\OO_{\frX}$-模层 $\F$ 的凝聚子模层，则对于每个开集 $U\subset\frX$，$\Gamma(U,\HH)$ 是拓扑群 $\Gamma(U,\F)$ 的闭子群。事实上，由 $\Gamma$ 的左正合性，它是 $\Gamma(U,\F)\to\Gamma(U,\F/\HH)$ 的核；$\F/\HH$ 凝聚（0，\hyperlink{para:0:5.3.4}{5.3.4}），所以上述同态连续，而 $\Gamma(U,\F/\HH)$ 为分离拓扑群，即得断言。

\egastatement{10.11.7}{命题}设 $\F,\G$ 为两个凝聚 $\OO_{\frX}$-模层。沿用（\hyperlink{para:1:10.11.3}{10.11.3}）的记号，可以定义拓扑 $\OO_{\frX}$-模层的典范函子同构
\begin{align}
\F\otimes_{\OO_{\frX}}\G&\simeq\plim_n(\F_n\otimes_{\OO_{X_n}}\G_n),\tag{10.11.7.1}\label{eq:1:10.11.7.1}\\
\shHom_{\OO_{\frX}}(\F,\G)&\simeq\plim_n\shHom_{\OO_{X_n}}(\F_n,\G_n).\tag{10.11.7.2}\label{eq:1:10.11.7.2}
\end{align}
这里拓扑如（\hyperlink{para:1:10.11.6}{10.11.6}）所定义。

（\hyperref[eq:1:10.11.7.1]{10.11.7.1}）来自公式
\[
\F_n\otimes_{\OO_{X_n}}\G_n
=(\F\otimes_{\OO_{\frX}}\OO_{X_n})\otimes_{\OO_{X_n}}
(\G\otimes_{\OO_{\frX}}\OO_{X_n})
=(\F\otimes_{\OO_{\frX}}\G)\otimes_{\OO_{\frX}}\OO_{X_n}
\]
及（\hyperlink{para:1:10.11.3}{10.11.3}）。将（\hyperref[eq:1:10.11.7.2]{10.11.7.2}）的两端暂视为不带拓扑的模层，则其同构来自相应同态层的截面的定义，以及对任意诺特仿射形式开集应用（\hyperref[eq:1:10.11.4.1]{10.11.4.1}）。尚须证明它在拟紧开集上的截面同态双连续，因而可归结到 $\frX=\Spf(A)$、$A$ 为诺特进制环。此时 $\F=M^\Delta$、$\G=N^\Delta$，$M,N$ 为有限型 $A$-模（\hyperlink{para:1:10.10.5}{10.10.5}）。结合（\hyperlink{para:1:10.10.2}{10.10.2}）及（\hyperlink{para:1:1.3.12}{1.3.12}，(ii)），只须证明典范同构
\[
\begin{gathered}
\Hom_A(M,N)\simeq\plim_n\Hom_{A_n}(M_n,N_n),\\
M_n=M\otimes_A A_n,\quad N_n=N\otimes_A A_n,
\end{gathered}
\]
是拓扑同构；这已在（0，\hyperlink{para:0:7.8.2}{7.8.2}）中证明。

\egapar{10.11.8}由于 $\Hom_{\OO_{\frX}}(\F,\G)$ 是拓扑群层 $\shHom_{\OO_{\frX}}(\F,\G)$ 的截面群，它自然带有群拓扑。若 $\frX$ 诺特，由（\hyperref[eq:1:10.11.7.2]{10.11.7.2}），各子群 $\Hom_{\OO_{\frX}}(\F,\J^n\G)$ 构成该群的零点邻域基本系。

\egastatement{10.11.9}{命题}设 $\frX$ 为诺特形式预概形，$\F,\G$ 为两个凝聚 $\OO_{\frX}$-模层。在拓扑群 $\Hom_{\OO_{\frX}}(\F,\G)$ 中，满同态（分别为单同态、双射同态）组成一个开子集。

由（\hyperlink{para:1:10.11.5}{10.11.5}），满同态的集合是离散群 $\Hom_{\OO_{X_0}}(\F_0,\G_0)$ 的一个子集在连续映射 $\Hom_{\OO_{\frX}}(\F,\G)\to\Hom_{\OO_{X_0}}(\F_0,\G_0)$ 下的逆像，故第一项成立。为证第二项，用有限个诺特仿射形式开集 $U_i$ 覆盖 $\frX$。$\theta\in\Hom_{\OO_{\frX}}(\F,\G)$ 为单射，当且仅当它在各连续限制映射
\[
\Hom_{\OO_{\frX}}(\F,\G)\longrightarrow
\Hom_{\OO_{\frX}|U_i}(\F|U_i,\G|U_i)
\]
下的像均为单射。于是归结到仿射情形，而这一情形已在（0，\hyperlink{para:0:7.8.3}{7.8.3}）中证明。


% END chapters/ch1_s10_11.tex

\par

% BEGIN chapters/ch1_s10_12.tex
\subsection{形式预概形的进制态射}\hypertarget{heading:1:10.12}{}
\src{206—207}
\egapar{10.12.1}设 $\frX,\frS$ 为两个局部诺特形式预概形。称态射 $f:\frX\to\frS$ 为\emph{进制态射}，如果存在 $\frS$ 的一个定义理想 $\J$，使 $\mathscr K=f^*(\J)\OO_{\frX}$ 为 $\frX$ 的定义理想；也称 $\frX$ 为一个\emph{进制 $\frS$-预概形}（关于 $f$）。此时，对 $\frS$ 的每个定义理想 $\J_1$，$\mathscr K_1=f^*(\J_1)\OO_{\frX}$ 都是 $\frX$ 的定义理想。

事实上，问题是局部的，可设 $\frX,\frS$ 仿射且诺特；于是存在 $n$ 使 $\J^n\subset\J_1$、$\J_1^n\subset\J$（\hyperlink{para:1:10.3.6}{10.3.6}及0，\hyperlink{para:1:7.1.4}{7.1.4}），故 $\mathscr K^n\subset\mathscr K_1$、$\mathscr K_1^n\subset\mathscr K$。第一个关系说明 $\mathscr K_1=K_1^\Delta$，其中 $K_1$ 为 $A=\Gamma(\frX,\OO_{\frX})$ 的开理想；第二个关系说明 $K_1$ 为 $A$ 的定义理想（0，\hyperlink{para:0:7.1.4}{7.1.4}），断言得证。

由此立即可知，若 $\frX,\frY$ 为两个进制 $\frS$-预概形，则每个 $\frS$-态射 $u:\frX\to\frY$ 都是进制的。事实上，若 $f:\frX\to\frS$、$g:\frY\to\frS$ 为结构态射，$\J$ 为 $\frS$ 的定义理想，则 $f=g\circ u$，所以
\[
u^*\bigl(g^*(\J)\OO_{\frY}\bigr)\OO_{\frX}=f^*(\J)\OO_{\frX}
\]
是 $\frX$ 的定义理想，而由假设 $g^*(\J)\OO_{\frY}$ 是 $\frY$ 的定义理想。

\egapar{10.12.2}以下固定一个局部诺特形式预概形 $\frS$ 及其定义理想 $\J$，并置 $S_n=(\frS,\OO_{\frS}/\J^{n+1})$。各局部诺特进制 $\frS$-预概形显然构成范畴。

称通常的局部诺特 $S_n$-预概形的归纳系 $(X_n)$ 为\emph{$(S_n)$-进制归纳系}，如果其结构态射 $f_n:X_n\to S_n$ 使得对 $m\le n$，下图交换，并将 $X_m$ 与 $X_n\times_{S_n}S_m=(X_n)_{(S_m)}$ 等同：
\begin{equation}\tag{10.12.2.1}\label{eq:1:10.12.2.1}
\begin{tikzcd}
X_n\arrow[d,"f_n"']&X_m\arrow[l]\arrow[d,"f_m"]\\
S_n&S_m\arrow[l]
\end{tikzcd}
\end{equation}
进制归纳系构成范畴：只须将态射 $(X_n)\to(Y_n)$ 定义为一个 $S_n$-态射 $u_n:X_n\to Y_n$ 的归纳系，且对 $m\le n$，$u_m$ 与 $(u_n)_{(S_m)}$ 等同。

\egastatement{10.12.3}{定理}进制 $\frS$-预概形的范畴与 $(S_n)$-进制归纳系的范畴典范等价。

这一等价如下得到。若 $\frX$ 为进制 $\frS$-预概形，$f:\frX\to\frS$ 为结构态射，则 $\mathscr K=f^*(\J)\OO_{\frX}$ 为 $\frX$ 的定义理想。使 $\frX$ 对应于 $X_n=(\frX,\OO_{\frX}/\mathscr K^{n+1})$ 的归纳系，结构态射 $f_n:X_n\to S_n$ 由 $f$ 得到（\hyperlink{para:1:10.5.6}{10.5.6}）。

先证 $(X_n)$ 是进制归纳系。若 $f=(\psi,\theta)$，则 $\psi^*(\J)\OO_{\frX}=\mathscr K$，所以对每个 $n$ 有 $\psi^*(\J^n)\OO_{\frX}=\mathscr K^n$；并由 $\psi^*$ 的正合性，对 $m\le n$ 有
\[
\mathscr K^{m+1}/\mathscr K^{n+1}
=\psi^*(\J^{m+1}/\J^{n+1})\,(\OO_{\frX}/\mathscr K^{n+1}).
\]
于是由（\hyperlink{para:1:4.4.5}{4.4.5}）得到结论。此外，立即可验，对进制 $\frS$-预概形的 $\frS$-态射 $u:\frX\to\frY$，相应有 $S_n$-态射 $u_n:X_n\to Y_n$ 的归纳系，且对 $m\le n$，$u_m=(u_n)_{(S_m)}$。

由此确实得到范畴等价，将由下一个更精确的命题得出。

\egastatement{10.12.3.1}{命题}设 $(X_n)$ 为 $S_n$-预概形归纳系，结构态射 $f_n:X_n\to S_n$ 使（\hyperref[eq:1:10.12.2.1]{10.12.2.1}）交换，并对 $m\le n$ 将 $X_m$ 与 $X_n\times_{S_n}S_m$ 等同。则 $(X_n)$ 满足（\hyperlink{para:1:10.6.3}{10.6.3}）的 b)、c) 条件。令 $\frX$ 为其归纳极限，$f:\frX\to\frS$ 为 $(f_n)$ 的归纳极限。若 $X_0$ 局部诺特，则 $\frX$ 局部诺特，且 $f$ 为进制态射。

由于在 $S_n$ 中定义 $S_m$ 的理想层幂零，由（\hyperlink{para:1:4.4.5}{4.4.5}），在 $X_n$ 中定义 $X_m$ 的理想层也幂零，故（\hyperlink{para:1:10.6.3}{10.6.3}）的条件成立。问题因而在 $\frX,\frS$ 上是局部的，可设 $\frS=\Spf(A)$、$\J=J^\Delta$，$A$ 为诺特 $J$-进制环，且 $X_n=\Spec(B_n)$。令 $A_n=A/J^{n+1}$、$J_n=J/J^{n+1}$。假设说明 $B_0$ 诺特，且 $B_m=B_n/J_n^{m+1}B_n$。因此 $B_n\to B_0$ 的核为 $K_n=J_nB_n$，而 $B_n\to B_m$ 的核为 $K_n^{m+1}$。由于 $A_1$ 诺特，$J_1$ 为有限型 $A_1$-模，故 $K_1=K_1/K_1^2$ 为有限型 $B_1$-模，因而也是有限型 $B_0=B_1/K_1$-模。由（\hyperlink{para:1:10.6.4}{10.6.4}），$\frX$ 诺特。

若 $B=\plim B_n$，则 $\frX=\Spf(B)$；令 $K=\Ker(B\to B_0)$，有 $B_n=B/K^{n+1}$。若 $\rho_n:A/J^{n+1}\to B/K^{n+1}$ 为对应于 $f_n$ 的同态，则
\[
K/K^{n+1}=(B/K^{n+1})\rho_n(J/J^{n+1}).
\]
由于对应于 $f$ 的 $\rho:A\to B$ 等于 $\plim\rho_n$，理想 $JB$ 在 $K$ 中稠密；$B$ 的每个理想都闭（0，\hyperlink{para:0:7.3.5}{7.3.5}），故 $K=JB$。若 $\mathscr K=K^\Delta$，则由（\hyperlink{para:1:10.10.9}{10.10.9}）有 $f^*(\J)\OO_{\frX}=\mathscr K$，证明完成。

\egapar{10.12.3.2}对于两个进制 $\frS$-预概形 $\frX,\frY$，上述等价给出典范双射
\[
\Hom_{\frS}(\frX,\frY)\simeq\plim_n\Hom_{S_n}(X_n,Y_n),
\]
其中射影极限关于映射 $u_n\mapsto(u_n)_{(S_m)}$（$m\le n$）而言。


% END chapters/ch1_s10_12.tex

\par

% BEGIN chapters/ch1_s10_13.tex
\subsection{有限型态射}\hypertarget{heading:1:10.13}{}
\src{207—209}
\egastatement{10.13.1}{命题}设 $\frY$ 为局部诺特形式预概形，$\mathscr K$ 为其定义理想，$f:\frX\to\frY$ 为形式预概形态射。下列条件等价：
\begin{enumerate}[label=\alph*)]
\item $\frX$ 局部诺特，$f$ 为进制态射（\hyperlink{para:1:10.12.1}{10.12.1}），且置 $\J=f^*(\mathscr K)\OO_{\frX}$ 后，由 $f$ 得到的态射 $f_0:(\frX,\OO_{\frX}/\J)\to(\frY,\OO_{\frY}/\mathscr K)$ 为有限型；
\item $\frX$ 局部诺特，且为一个 $(Y_n)$-进制归纳系 $(X_n)$ 的归纳极限，其中 $X_0\to Y_0$ 为有限型；
\item $\frY$ 的每一点都有满足下列性质的诺特仿射形式开邻域 $V$：

(Q) $f^{-1}(V)$ 为有限个诺特仿射形式开集 $U_i$ 的并，每个诺特进制环 $\Gamma(U_i,\OO_{\frX})$ 都拓扑同构于 $\Gamma(V,\OO_{\frY})$ 上某个限制形式幂级数代数（0，\hyperlink{para:0:7.5.1}{7.5.1}）对一个必为闭的理想的商。
\end{enumerate}

由（\hyperlink{para:1:10.12.3}{10.12.3}），a) 蕴含 b) 是立即的。为证 b) 蕴含 c)，由于问题在 $\frY$ 上是局部的，可设 $\frY=\Spf(B)$，$B$ 为诺特进制环，$\mathscr K=K^\Delta$。由假设 $X_0$ 在 $Y_0$ 上为有限型，故 $X_0$ 为有限个仿射开集 $U_i$ 的并，且 $X_0$ 在 $U_i$ 上诱导的仿射概形的环 $A_{i0}$ 为 $B/K$ 上的有限型代数（\hyperlink{para:1:6.3.2}{6.3.2}）。由（\hyperlink{para:1:5.1.9}{5.1.9}），$U_i$ 也是每个诺特预概形 $X_n$ 的仿射开集。若 $A_{in}$ 为相应的环，则 b) 蕴含 $A_{im}=A_{in}/K^{m+1}A_{in}$（$m\le n$）。因此 $\frX$ 在 $U_i$ 上诱导的形式预概形同构于 $\Spf(A_i)$，其中 $A_i=\plim_n A_{in}$（\hyperlink{para:1:10.6.4}{10.6.4}）。$A_i$ 为 $KA_i$-进制环，且 $A_i/KA_i\simeq A_{i0}$ 为有限型 $B/K$-代数。由（0，\hyperlink{para:0:7.5.5}{7.5.5}），$A_i$ 拓扑同构于 $B$ 上某个限制形式幂级数代数的商；这样的代数诺特（0，\hyperlink{para:0:7.5.4}{7.5.4}），所取理想必闭。

为证 c) 蕴含 a)，可限于 $\frX=\Spf(A)$ 也仿射，其中 $A$ 是诺特进制环，且拓扑同构于 $B$ 上某个限制形式幂级数代数对闭理想的商。由（0，\hyperlink{para:0:7.5.5}{7.5.5}），$A/KA$ 为有限型 $B/K$-代数，且 $J=KA$ 为 $A$ 的定义理想；所以据（\hyperlink{para:1:10.10.9}{10.10.9}），a) 的条件成立。

注意，当命题（\hyperlink{para:1:10.13.1}{10.13.1}）的条件成立时，由 c) 可知，性质 a) 对 $\frY$ 的任一定义理想 $\mathscr K$ 都成立；因而在 b) 中，所有 $f_n$ 都为有限型态射。

\egastatement{10.13.2}{推论}若（\hyperlink{para:1:10.13.1}{10.13.1}）的条件成立，则 $\frY$ 的每个诺特仿射形式开集 $V$ 都具有性质 (Q)；若 $\frY$ 诺特，则 $\frX$ 也诺特。

这立即由（\hyperlink{para:1:10.13.1}{10.13.1}）及（\hyperlink{para:1:6.3.2}{6.3.2}）得出。

\egastatement{10.13.3}{定义}当（\hyperlink{para:1:10.13.1}{10.13.1}）的等价性质 a)、b)、c) 成立时，称态射 $f$ 为\emph{有限型态射}，或称 $\frX$ 为\emph{有限型 $\frY$-形式预概形}、\emph{$\frY$ 上的有限型形式预概形}。

\egastatement{10.13.4}{推论}设 $\frX=\Spf(A)$、$\frY=\Spf(B)$ 为两个诺特仿射形式概形。$\frX$ 在 $\frY$ 上为有限型的充要条件是，诺特进制环 $A$ 同构于 $B$ 上一个限制形式幂级数代数对闭理想的商。

沿用（\hyperlink{para:1:10.13.1}{10.13.1}）的记号，若 $\frX$ 在 $\frY$ 上为有限型，则 $A/KA$ 为有限型 $B/K$-代数（\hyperlink{para:1:6.3.3}{6.3.3}），而 $KA$ 为 $A$ 的定义理想（\hyperlink{para:1:10.10.9}{10.10.9}）。所以由（0，\hyperlink{para:0:7.5.5}{7.5.5}）得出结论。

\egastatement{10.13.5}{命题}(i) 两个有限型形式预概形态射的复合仍为有限型。

(ii) 设 $\frX,\frS,\frS'$ 为三个局部诺特（分别为诺特）形式预概形，$f:\frX\to\frS$、$g:\frS'\to\frS$ 为两个态射。若 $f$ 为有限型，则 $\frX\times_{\frS}\frS'$ 局部诺特（分别为诺特），并在 $\frS'$ 上为有限型。

(iii) 设 $\frS$ 为局部诺特形式预概形，$\frX',\frY'$ 为两个局部诺特 $\frS$-形式预概形，且 $\frX'\times_{\frS}\frY'$ 局部诺特。若 $\frX,\frY$ 为局部诺特 $\frS$-形式预概形，$f:\frX\to\frX'$、$g:\frY\to\frY'$ 为有限型 $\frS$-态射，则 $\frX\times_{\frS}\frY$ 局部诺特，$f\times_{\frS}g$ 为有限型 $\frS$-态射。

(iii) 由 (i)、(ii) 通过（\hyperlink{para:1:3.5.1}{3.5.1}）的形式论证得出，故只须证明 (i)、(ii)。

设 $f:\frX\to\frY$、$g:\frY\to\frZ$ 为局部诺特形式预概形的有限型态射。若 $\LL$ 为 $\frZ$ 的定义理想，则 $\mathscr K=g^*(\LL)\OO_{\frY}$ 为 $\frY$ 的定义理想，$\J=f^*(g^*(\LL))\OO_{\frX}$ 为 $\frX$ 的定义理想。置 $X_0=(\frX,\OO_{\frX}/\J)$、$Y_0=(\frY,\OO_{\frY}/\mathscr K)$、$Z_0=(\frZ,\OO_{\frZ}/\LL)$，并令 $f_0:X_0\to Y_0$、$g_0:Y_0\to Z_0$ 对应于 $f,g$。由假设 $f_0,g_0$ 为有限型，故对应于 $g\circ f$ 的 $g_0\circ f_0$ 也为有限型（\hyperlink{para:1:6.3.4}{6.3.4}）；据（\hyperlink{para:1:10.13.1}{10.13.1}），$g\circ f$ 为有限型。

在 (ii) 的条件下，$\frS,\frX,\frS'$ 分别为局部诺特预概形序列 $(S_n),(X_n),(S_n')$ 的归纳极限，而且可设 $X_m=X_n\times_{S_n}S_m$（$m\le n$）（\hyperlink{para:1:10.13.1}{10.13.1}）。于是 $\frX\times_{\frS}\frS'$ 为各 $X_n\times_{S_n}S_n'$ 的归纳极限（\hyperlink{para:1:10.7.4}{10.7.4}），并有
\[
X_m\times_{S_m}S_m'
=(X_n\times_{S_n}S_m)\times_{S_m}S_m'
=(X_n\times_{S_n}S_n')\times_{S_n'}S_m'.
\]
由于 $X_0$ 在 $S_0$ 上为有限型，$X_0\times_{S_0}S_0'$ 局部诺特（\hyperlink{para:1:6.3.8}{6.3.8}）。由（\hyperlink{para:1:10.12.3.1}{10.12.3.1}），先知 $\frX\times_{\frS}\frS'$ 局部诺特；又因 $X_0\times_{S_0}S_0'$ 在 $S_0'$ 上为有限型（\hyperlink{para:1:6.3.8}{6.3.8}），由（\hyperlink{para:1:10.12.3.1}{10.12.3.1}）及（\hyperlink{para:1:10.13.1}{10.13.1}），此形式预概形在 $\frS'$ 上为有限型。关于诺特预概形的断言是（\hyperlink{para:1:6.3.8}{6.3.8}）的直接结果，(ii) 得证。

\egastatement{10.13.6}{推论}在（\hyperlink{para:1:10.9.9}{10.9.9}）的假设下，若 $f$ 为有限型态射，则它到完备化的延拓 $\widehat f$ 也是有限型态射。


% END chapters/ch1_s10_13.tex

\par

% BEGIN chapters/ch1_s10_14.tex
\subsection{形式预概形的闭子预概形}\hypertarget{heading:1:10.14}{}
\src{209—212}
\egastatement{10.14.1}{命题}设 $\frX$ 为局部诺特形式预概形，$\A$ 为 $\OO_{\frX}$ 的凝聚理想层。若 $\frY$ 为 $\OO_{\frX}/\A$ 的支集（它是闭集），则拓扑赋环空间 $(\frY,(\OO_{\frX}/\A)|\frY)$ 是局部诺特形式预概形；若 $\frX$ 诺特，则它也诺特。

由（\hyperlink{para:1:10.10.3}{10.10.3}）及（0，\hyperlink{para:0:5.3.4}{5.3.4}），$\OO_{\frX}/\A$ 凝聚，故其支集 $\frY$ 闭（0，\hyperlink{para:0:5.2.2}{5.2.2}）。取 $\frX$ 的定义理想 $\J$，置 $X_n=(\frX,\OO_{\frX}/\J^{n+1})$。环层 $\OO_{\frX}/\A$ 为各
\[
\OO_{\frX}/(\A+\J^{n+1})
=(\OO_{\frX}/\A)\otimes_{\OO_{\frX}}(\OO_{\frX}/\J^{n+1})
\]
的射影极限（\hyperlink{para:1:10.11.3}{10.11.3}），且这些层的支集都为 $\frY$。理想层 $(\A+\J^{n+1})/\J^{n+1}$ 是凝聚 $\OO_{\frX}$-模层，因而也是凝聚 $(\OO_{\frX}/\J^{n+1})$-模层（0，\hyperlink{para:0:5.3.10}{5.3.10}）。令 $Y_n$ 为它在 $X_n$ 中定义的闭子预概形，则 $(\frY,(\OO_{\frX}/\A)|\frY)$ 显然是各 $Y_n$ 的归纳极限。由于（\hyperlink{para:1:10.6.4}{10.6.4}）的条件成立，它局部诺特；若 $\frX$ 诺特，则 $Y_0$ 诺特（\hyperlink{para:1:6.1.4}{6.1.4}），故它也诺特。

\egastatement{10.14.2}{定义}称形如 $(\frY,(\OO_{\frX}/\A)|\frY)$ 的形式预概形为形式预概形 $\frX$ 的\emph{闭子预概形}，其中 $\A$ 为凝聚理想层；称它为\emph{由 $\A$ 定义的闭子预概形}。\transnote{原版此处及下一句写作“凝聚 $\OO_{\frX}$-模层”。根据（\hyperlink{para:1:10.14.1}{10.14.1}）及商环层 $\OO_{\frX}/\A$ 的定义，这里须指 $\OO_{\frX}$ 的凝聚理想层，而非任意凝聚模层；正文据此明确写出。}

凝聚理想层与 $\frX$ 的闭子预概形之间这样定义的对应显然是双射。拓扑赋环空间态射 $j=(\psi,\theta):\frY\to\frX$，其中 $\psi$ 为底空间嵌入、$\theta^\sharp$ 为典范商同态的相应限制，显然是形式预概形态射（\hyperlink{para:1:10.4.5}{10.4.5}），称为 $\frY$ 到 $\frX$ 的\emph{典范嵌入}。

若 $\frX=\Spf(A)$，$A$ 为诺特进制环，则 $\A=\mathfrak a^\Delta$，其中 $\mathfrak a$ 为 $A$ 的理想（\hyperlink{para:1:10.10.5}{10.10.5}）。于是从上面的讨论立即得知，在同构意义下 $\frY=\Spf(A/\mathfrak a)$，且 $j$ 对应于典范同态 $A\to A/\mathfrak a$（\hyperlink{para:1:10.2.2}{10.2.2}）。

称局部诺特形式预概形的态射 $f:\frZ\to\frX$ 为\emph{闭浸入}，如果它可分解为 $\frZ\xrightarrow{g}\frY\xrightarrow{j}\frX$，其中 $g$ 为 $\frZ$ 到 $\frX$ 的一个闭子预概形 $\frY$ 的同构，$j$ 为典范嵌入。由于 $j$ 为赋环空间的单态射，$g,\frY$ 必然唯一。

\egastatement{10.14.3}{命题}闭浸入是有限型态射。

立即归结到 $\frX=\Spf(A)$、$\frY=\Spf(A/\mathfrak a)$ 的情形，再用（\hyperlink{para:1:10.13.1}{10.13.1}，c)）即得命题。

\egastatement{10.14.4}{引理}设 $f:\frY\to\frX$ 为局部诺特形式预概形的态射，$(U_\alpha)$ 为由 $\frX$ 的诺特仿射形式开集组成的 $f(\frY)$ 的覆盖，且各 $f^{-1}(U_\alpha)$ 都是 $\frY$ 的诺特仿射形式开集。$f$ 为闭浸入的充要条件是，$f(\frY)$ 在 $\frX$ 中闭，且对每个 $\alpha$，$f$ 在 $f^{-1}(U_\alpha)$ 上的限制对应于满同态
\[
\Gamma(U_\alpha,\OO_{\frX})\longrightarrow
\Gamma(f^{-1}(U_\alpha),\OO_{\frY})
\]
（\hyperlink{para:1:10.4.6}{10.4.6}）。

必要性显然。反之，若这些条件成立，以 $\mathfrak a_\alpha$ 表示上述截面同态的核。取 $\A|U_\alpha=\mathfrak a_\alpha^\Delta$，并在 $\frX\setminus f(\frY)$ 上取 $\A=\OO_{\frX}$，便可定义凝聚理想层 $\A$。\transnote{原版写成在各 $U_\alpha$ 的并的补集中令 $\A$ 为零，并写 $\mathfrak a_\alpha^\Delta$ 的支集为 $U_\alpha\cap f(\frY)$。对于定义闭子预概形的理想层，零支集应当属于商层 $\OO_{U_\alpha}/\mathfrak a_\alpha^\Delta$；在闭像之外理想层应等于整个结构层。正文明确采用这一与局部商结构一致的表述。}
事实上，由于 $f(\frY)$ 闭，且 $\OO_{U_\alpha}/\mathfrak a_\alpha^\Delta$ 的支集为 $U_\alpha\cap f(\frY)$，只须验证 $\mathfrak a_\alpha^\Delta,\mathfrak a_\beta^\Delta$ 在任意诺特仿射形式开集 $V\subset U_\alpha\cap U_\beta$ 上诱导相同的层。

$f$ 在 $f^{-1}(U_\alpha)$ 上的限制是到 $U_\alpha$ 的闭浸入，因此 $f^{-1}(V)$ 是 $f^{-1}(U_\alpha)$ 的诺特仿射形式开集，且 $f$ 在其上的限制仍为闭浸入。若 $\mathfrak b$ 为相应满同态 $\Gamma(V,\OO_{\frX})\to\Gamma(f^{-1}(V),\OO_{\frY})$ 的核，则由（\hyperlink{para:1:10.10.2}{10.10.2}）立即可知，$\mathfrak a_\alpha^\Delta$ 在 $V$ 上诱导 $\mathfrak b^\Delta$。定义 $\A$ 后，显然 $f=j\circ g$，其中 $j:\frZ\to\frX$ 为 $\A$ 定义的闭子预概形的典范嵌入，$g:\frY\to\frZ$ 为同构。\transnote{原版此处将复合印为 $g\circ j$；按给定的定义域、值域，应为 $j\circ g$。}

\egastatement{10.14.5}{命题}(i) 若 $f:\frZ\to\frY$、$g:\frY\to\frX$ 为局部诺特形式预概形的闭浸入，则 $g\circ f$ 为闭浸入。

(ii) 设 $\frX,\frY,\frS$ 为三个局部诺特形式预概形，$f:\frX\to\frS$ 为闭浸入，$g:\frY\to\frS$ 为态射。则 $\frX\times_{\frS}\frY\to\frY$ 为闭浸入。

(iii) 设 $\frS$ 为局部诺特形式预概形，$\frX',\frY'$ 为局部诺特 $\frS$-形式预概形，且 $\frX'\times_{\frS}\frY'$ 局部诺特。若 $\frX,\frY$ 为局部诺特 $\frS$-形式预概形，$f:\frX\to\frX'$、$g:\frY\to\frY'$ 为闭浸入的 $\frS$-态射，则 $f\times_{\frS}g$ 为闭浸入。

由（\hyperlink{para:1:3.5.1}{3.5.1}），仍只须证明 (i)、(ii)。为证 (i)，可设 $\frY$ 为 $\frX$ 的闭子预概形，由凝聚理想层 $\J\subset\OO_{\frX}$ 定义；$\frZ$ 为 $\frY$ 的闭子预概形，由凝聚理想层 $\mathscr K\subset\OO_{\frY}$ 定义。若 $\psi:\frY\hookrightarrow\frX$ 为底空间嵌入，则 $\psi_*(\mathscr K)$ 为 $\psi_*(\OO_{\frY})=\OO_{\frX}/\J$ 的凝聚理想层（0，\hyperlink{para:0:5.3.12}{5.3.12}），故也是凝聚 $\OO_{\frX}$-模层（0，\hyperlink{para:0:5.3.10}{5.3.10}）。于是 $\OO_{\frX}\to(\OO_{\frX}/\J)/\psi_*(\mathscr K)$ 的核 $\mathscr K_1$ 凝聚（0，\hyperlink{para:0:5.3.4}{5.3.4}），且 $\OO_{\frX}/\mathscr K_1\simeq\psi_*(\OO_{\frY}/\mathscr K)$，这说明 $\frZ$ 同构于 $\frX$ 的闭子预概形。

为证 (ii)，立即可限于 $\frS=\Spf(A)$、$\frX=\Spf(B)$、$\frY=\Spf(C)$，其中 $A$ 为诺特 $J$-进制环，$B=A/\mathfrak a$，$C$ 为拓扑 $A$-代数且为诺特进制环。只须证明 $C\to C\ctensor_A(A/\mathfrak a)$ 为满射。由于 $A/\mathfrak a$ 为有限型 $A$-模，且带 $J$-进制拓扑，根据（0，\hyperlink{para:0:7.7.8}{7.7.8}），$C\ctensor_A(A/\mathfrak a)=C\otimes_A(A/\mathfrak a)=C/\mathfrak aC$，从而得证。

\egastatement{10.14.6}{推论}在（\hyperlink{para:1:10.14.5}{10.14.5}，(ii)）的假设下，设 $p,q$ 为投影，使下图交换：
\[
\begin{tikzcd}
\frX\arrow[d,"f"']&\frX\times_{\frS}\frY\arrow[l,"p"']\arrow[d,"q"]\\
\frS&\frY\arrow[l,"g"]
\end{tikzcd}
\]
则对每个凝聚 $\OO_{\frX}$-模层 $\F$，存在 $\OO_{\frY}$-模层的典范同构
\begin{equation}\tag{10.14.6.1}\label{eq:1:10.14.6.1}
u:g^*(f_*(\F))\simeq q_*(p^*(\F)).
\end{equation}

为定义这样的同态 $u$，由（0，\hyperlink{para:0:4.4.3}{4.4.3}），等价地只须定义
\[
f_*(\F)\longrightarrow g_*(q_*(p^*(\F)))=f_*(p_*(p^*(\F))).
\]
取此同态为 $u^\flat=f_*(\rho)$，其中 $\rho:\F\to p_*(p^*(\F))$ 为典范同态（0，\hyperlink{para:0:4.4.3}{4.4.3}）。为证 $u$ 为同构，立即归结到 $\frS,\frX,\frY$ 分别为诺特进制环 $A,B,C$ 的形式谱，且满足（\hyperlink{para:1:10.14.5}{10.14.5}，(ii)）中上述条件。此时 $\F=M^\Delta$，$M$ 为有限型 $(A/\mathfrak a)$-模（\hyperlink{para:1:10.10.5}{10.10.5}）。由（\hyperlink{para:1:10.10.8}{10.10.8}），（\hyperref[eq:1:10.14.6.1]{10.14.6.1}）的两端分别与 $(C\otimes_A M)^\Delta$ 及 $((C/\mathfrak aC)\otimes_{A/\mathfrak a}M)^\Delta$ 等同。由于
\[
(C/\mathfrak aC)\otimes_{A/\mathfrak a}M
=(C\otimes_A(A/\mathfrak a))\otimes_{A/\mathfrak a}M
\simeq C\otimes_A M,
\]
推论得证。

\egastatement{10.14.7}{推论}设 $X$ 为通常的局部诺特预概形，$Y$ 为其闭子预概形，$j:Y\to X$ 为典范嵌入，$X'$ 为 $X$ 的闭子集，$Y'=Y\cap X'$。则 $\widehat j:Y_{/Y'}\to X_{/X'}$ 为闭浸入；对每个凝聚 $\OO_Y$-模层 $\F$，有
\[
\widehat j_*(\F_{/Y'})=(j_*(\F))_{/X'}.
\]
由于 $Y'=j^{-1}(X')$，只须利用（\hyperlink{para:1:10.9.9}{10.9.9}）并应用（\hyperlink{para:1:10.14.5}{10.14.5}）、（\hyperlink{para:1:10.14.6}{10.14.6}）。


% END chapters/ch1_s10_14.tex

\par

% BEGIN chapters/ch1_s10_15.tex
\subsection{分离的形式预概形}\hypertarget{heading:1:10.15}{}
\src{212—213}
\egastatement{10.15.1}{定义}设 $\frS$ 为形式预概形，$\frX$ 为 $\frS$-形式预概形，$f:\frX\to\frS$ 为结构态射。称态射 $(1_{\frX},1_{\frX})_{\frS}:\frX\to\frX\times_{\frS}\frX$ 为\emph{对角态射}，记为 $\Delta_{\frX/\frS}$，也记为 $\Delta_f$。若底空间 $\frX$ 在 $\Delta_f$ 下的像为 $\frX\times_{\frS}\frX$ 的底空间的闭子集，则称 $\frX$ \emph{在 $\frS$ 上分离}，或称它为一个\emph{$\frS$-形式概形}，也称 $f$ 为\emph{分离态射}。若形式预概形 $\frX$ 在 $\mathbb Z$ 上分离，则称它\emph{分离}，或称它为一个\emph{形式概形}。

\egastatement{10.15.2}{命题}假设形式预概形 $\frS,\frX$ 分别为通常预概形序列 $(S_n),(X_n)$ 的归纳极限，且态射 $f:\frX\to\frS$ 为态射序列 $f_n:X_n\to S_n$ 的归纳极限。则 $f$ 分离的充要条件是 $f_0:X_0\to S_0$ 分离。

此时 $\Delta_f$ 为各 $\Delta_{X_n/S_n}$ 的归纳极限（\hyperlink{para:1:10.7.4}{10.7.4}），而 $\frX$ 的底空间在 $\Delta_{\frX/\frS}$ 下的像（分别为 $\frX\times_{\frS}\frX$ 的底空间）与 $X_0$ 在 $\Delta_{X_0/S_0}$ 下的像（分别为 $X_0\times_{S_0}X_0$ 的底空间）相同，故得结论。

\egastatement{10.15.3}{命题}以下假定所有考虑的形式预概形（分别为形式预概形态射），都是通常预概形序列（分别为通常预概形态射序列）的归纳极限。
\begin{enumerate}[label=(\roman*)]
\item 两个分离态射的复合分离；
\item 若 $f:\frX\to\frX'$、$g:\frY\to\frY'$ 为两个分离的 $\frS$-态射，则 $f\times_{\frS}g$ 分离；
\item 若 $f:\frX\to\frY$ 为分离的 $\frS$-态射，则对于任意基形式预概形扩张 $\frS'\to\frS$，$\frS'$-态射 $f_{(\frS')}$ 分离；
\item 若两个态射的复合 $g\circ f$ 分离，则 $f$ 分离。
\end{enumerate}

这里约定：若同一个形式预概形 $\frZ$ 在同一命题中出现多次，则在各处都将它视为同一个通常预概形序列 $(Z_n)$ 的归纳极限；从 $\frZ$ 出发的态射（分别为到 $\frZ$ 的态射）均视为从 $Z_n$ 出发（分别为到 $Z_n$）的通常预概形态射的归纳极限。

沿用（\hyperlink{para:1:10.15.2}{10.15.2}）的记号，有 $(g\circ f)_0=g_0\circ f_0$，以及 $(f\times_{\frS}g)_0=f_0\times_{S_0}g_0$。因此（\hyperlink{para:1:10.15.3}{10.15.3}）的各项断言立即由（\hyperlink{para:1:10.15.2}{10.15.2}）和通常预概形的相应断言（\hyperlink{para:1:5.5.1}{5.5.1}）得出。

对于同样一类形式预概形及态射，我们将与（\hyperlink{para:1:5.5.5}{5.5.5}）、（\hyperlink{para:1:5.5.9}{5.5.9}）、（\hyperlink{para:1:5.5.10}{5.5.10}）对应的命题留给读者表述，其中应将“仿射开集”换为“满足（\hyperlink{para:1:10.6.3}{10.6.3}）的条件 b) 的仿射形式开集”。类似论证也表明，每个诺特仿射形式概形都分离，这使所用术语得到说明。

\egastatement{10.15.4}{命题}设 $\frS$ 为局部诺特形式预概形，$\frX,\frY$ 为两个局部诺特 $\frS$-形式预概形，且 $\frX$ 或 $\frY$ 在 $\frS$ 上为有限型（故 $\frX\times_{\frS}\frY$ 局部诺特），并且 $\frY$ 在 $\frS$ 上分离。若 $f:\frX\to\frY$ 为 $\frS$-态射，则图态射
\[
\Gamma_f=(1_{\frX},f)_{\frS}:\frX\longrightarrow\frX\times_{\frS}\frY
\]
为闭浸入。

可设 $\frS$ 为局部诺特预概形序列 $(S_n)$ 的归纳极限，$\frX,\frY$ 分别为 $S_n$-预概形序列 $(X_n),(Y_n)$ 的归纳极限，$f$ 为 $S_n$-态射序列 $f_n:X_n\to Y_n$ 的归纳极限。于是 $\frX\times_{\frS}\frY$ 为 $(X_n\times_{S_n}Y_n)$ 的归纳极限，$\Gamma_f$ 为 $(\Gamma_{f_n})$ 的归纳极限（\hyperlink{para:1:10.7.4}{10.7.4}）。由假设，$Y_0$ 在 $S_0$ 上分离（\hyperlink{para:1:10.15.2}{10.15.2}），故 $\Gamma_{f_0}(X_0)$ 为 $X_0\times_{S_0}Y_0$ 的闭子空间。由于这些空间分别与 $\Gamma_f(\frX)$ 及 $\frX\times_{\frS}\frY$ 的底空间相同，首先得到 $\Gamma_f(\frX)$ 闭。

当 $(U,V)$ 遍历满足 $f(U)\subset V$ 的 $\frX,\frY$ 的诺特仿射形式开集对时，各 $U\times_{\frS}V$ 在 $\frX\times_{\frS}\frY$ 中覆盖 $\Gamma_f(\frX)$；若 $f_U:U\to V$ 为限制，则 $\Gamma_{f_U}:U\to U\times_{\frS}V$ 为 $\Gamma_f$ 的相应限制。若证明 $\Gamma_{f_U}$ 为闭浸入，则 $\Gamma_f$ 也为闭浸入（\hyperlink{para:1:10.14.4}{10.14.4}）。换言之，归结到 $\frS=\Spf(A)$、$\frX=\Spf(B)$、$\frY=\Spf(C)$，$A,B,C$ 为诺特进制环，且 $f$ 对应于连续 $A$-同态 $\varphi:C\to B$。此时 $\Gamma_f$ 对应于唯一连续同态 $\omega:B\ctensor_A C\to B$，它与典范同态 $B\to B\ctensor_A C$、$C\to B\ctensor_A C$ 的复合分别为恒等同态及 $\varphi$。显然 $\omega$ 为满射，故断言得证。

\egastatement{10.15.5}{推论}设 $\frS$ 为局部诺特形式预概形，$\frX$ 为有限型 $\frS$-形式预概形。$\frX$ 在 $\frS$ 上分离的充要条件是，对角态射 $\frX\to\frX\times_{\frS}\frX$ 为闭浸入。

\egastatement{10.15.6}{命题}局部诺特形式预概形的闭浸入 $j:\frY\to\frX$ 是分离态射。

沿用（\hyperlink{para:1:10.14.2}{10.14.2}）的记号，$j_0:Y_0\to X_0$ 为闭浸入，故为分离态射；只须应用（\hyperlink{para:1:10.15.2}{10.15.2}）。

\egastatement{10.15.7}{命题}设 $X$ 为通常的局部诺特预概形，$X'$ 为 $X$ 的闭子集，$\widehat X=X_{/X'}$。$\widehat X$ 分离的充要条件是 $\widehat X_{\mathrm{red}}$ 分离；$X$ 分离则足以保证这一点。

沿用（\hyperlink{para:1:10.8.5}{10.8.5}）的记号，$\widehat X$ 分离当且仅当 $X_0$ 分离（\hyperlink{para:1:10.15.2}{10.15.2}）；又因 $\widehat X_{\mathrm{red}}=(X_0)_{\mathrm{red}}$，后者等价于 $\widehat X_{\mathrm{red}}$ 分离（\hyperlink{para:1:5.5.1}{5.5.1}，(vi)）。


% END chapters/ch1_s10_15.tex

\par

% END text/ch1.tex


% BEGIN text/ch2.tex
\renewcommand{\B}{\mathscr B}
\setcounter{chapter}{1}

\chapter{若干类态射的初等整体研究}
\src{5}
\subsection*{内容提要}
\begin{enumerate}[label=\S\,\arabic*.]
\item 仿射态射。
\item 齐次素谱。
\item 分次代数层的齐次谱。
\item 射影丛。丰沛层。
\item 拟仿射态射；拟射影态射；固有态射；射影态射。
\item 整态射与有限态射。
\item 赋值判别准则。
\item 爆破概形；投影锥；射影闭包。
\end{enumerate}

本章研究各种类的态射时，并不大量使用上同调方法；利用这些方法作更深入的研究，将在第三章进行，在那里主要使用第二章的第2、4和5节。初读时可以略去第8节：它对第1至第3节所发展的形式体系作一些补充，这些补充归结为该形式体系的简单应用；我们使用它们的频率，将低于使用本章其他结果的频率。

\section{仿射态射}
本节的大多数结果，是第一章第1节中相应结果的“整体”版本；因此，它们在本质上并非新的结果，只是为后面的讨论提供一种方便的语言。

\subsection{\texorpdfstring{$S$-预概形与 $\OO_S$-代数}{S-预概形与结构层代数}}
\src{5--6}
\egapar{1.1.1}
设 $S$ 是一个预概形，$X$ 是一个 $S$-预概形，$f:X\to S$ 是其结构态射。我们知道（0，4.2.4），直像 $f_*(\OO_X)$ 是一个 $\OO_S$-代数；在不会引起混淆时，将它记作 $\A(X)$。若 $U$ 是 $S$ 的一个开集，则有
\[
\A\bigl(f^{-1}(U)\bigr)=\A(X)|U.
\]
同样，对任意 $\OO_X$-模 $\F$（分别地，任意 $\OO_X$-代数 $\B$），我们将直像 $f_*(\F)$（分别地，$f_*(\B)$）写作 $\A(\F)$（分别地，$\A(\B)$）；它不仅是一个 $\OO_S$-模（分别地，$\OO_S$-代数），而且是一个 $\A(X)$-模（分别地，$\A(X)$-代数）。

\egapar{1.1.2}
设 $Y$ 是另一个 $S$-预概形，$g:Y\to S$ 是其结构态射，$h:X\to Y$ 是一个 $S$-态射；于是有交换图
\begin{equation}\tag{1.1.2.1}
\begin{tikzcd}[column sep=large,row sep=large]
X \arrow[rr,"h"] \arrow[dr,"f"'] && Y \arrow[dl,"g"]\\
& S &
\end{tikzcd}
\end{equation}
按定义，$h=(\psi,\theta)$，其中
$\theta:\OO_Y\to\psi_*(\OO_X)=h_*(\OO_X)$ 是环层同态。由此得到（0，4.2.2）一个 $\OO_S$-代数同态
\[
g_*(\theta):g_*(\OO_Y)\longrightarrow g_*\bigl(h_*(\OO_X)\bigr)=f_*(\OO_X),
\]
也就是一个 $\OO_S$-代数同态 $\A(Y)\to\A(X)$，我们仍将它记作 $\A(h)$。若 $h':Y\to Z$ 是另一个 $S$-态射，立即有
\[
\A(h'\circ h)=\A(h)\circ\A(h').
\]
因此，我们在 $S$-预概形范畴上定义了一个取值于 $\OO_S$-代数范畴的反变函子 $\A(X)$。

现设 $\F$ 是一个 $\OO_X$-模，$\G$ 是一个 $\OO_Y$-模，$u:\G\to\F$ 是一个 $h$-态射，即一个 $\OO_Y$-模同态 $\G\to h_*(\F)$（0，4.4.1）。则
\[
g_*(u):g_*(\G)\longrightarrow g_*\bigl(h_*(\F)\bigr)=f_*(\F)
\]
是一个 $\OO_S$-模同态 $\A(\G)\to\A(\F)$，记作 $\A(u)$。此外，二元组 $(\A(h),\A(u))$ 构成从 $\A(Y)$-模 $\A(\G)$ 到 $\A(X)$-模 $\A(\F)$ 的双同态。

\egapar{1.1.3}
固定预概形 $S$。我们可以将二元组 $(X,\F)$（其中 $X$ 是一个 $S$-预概形，$\F$ 是一个 $\OO_X$-模）看作一个范畴的对象：从 $(X,\F)$ 到 $(Y,\G)$ 的态射定义为二元组 $(h,u)$，其中 $h:X\to Y$ 是 $S$-态射，而 $u:\G\to\F$ 是 $h$-态射。于是可以说，$(\A(X),\A(\F))$ 是一个反变函子，取值于如下范畴：其对象是由一个 $\OO_S$-代数以及该代数上的一个模组成的二元组，其态射是双同态。

\subsection{一个预概形之上的仿射预概形}
\src{6--8}
\egastatement{1.2.1}{定义}
设 $X$ 是一个 $S$-预概形，$f:X\to S$ 是其结构态射。若 $S$ 存在一个由仿射开集组成的覆盖 $(S_\alpha)$，使得对每个 $\alpha$，$X$ 在开集 $f^{-1}(S_\alpha)$ 上诱导的预概形是仿射的，则称 $X$ 在 $S$ 之上是\textbf{仿射的}。

\egastatement{1.2.2}{例}
$S$ 的每个闭子预概形，都是一个在 $S$ 之上仿射的 $S$-预概形（I，4.2.3及4.2.4）。

\egastatement{1.2.3}{注}
在 $S$ 之上仿射的预概形 $X$，未必是仿射概形，例 $X=S$ 即说明这一点（1.2.2）。另一方面，若仿射概形 $X$ 是一个 $S$-预概形，$X$ 也未必在 $S$ 之上仿射（见1.3.3）。不过，回顾一下：若 $S$ 是一个概形，则每个本身为仿射概形的 $S$-预概形，都是在 $S$ 之上仿射的（I，5.5.10）。

\egastatement{1.2.4}{命题}
每个在 $S$ 之上仿射的 $S$-预概形，都在 $S$ 之上分离；换言之，它是一个 $S$-概形。
\begin{proof}
这立即由（I，5.5.5）及（I，5.5.8）得出。
\end{proof}

\egastatement{1.2.5}{命题}
设 $X$ 是一个在 $S$ 之上仿射的 $S$-概形，$f:X\to S$ 是结构态射。对任意开集 $U\subset S$，$f^{-1}(U)$ 在 $U$ 之上仿射。
\begin{proof}
根据定义（1.2.1），可归结到 $S=\Spec(A)$ 与 $X=\Spec(B)$ 均仿射的情形；这时
$f=({}^{a}\varphi,\wt\varphi)$，其中 $\varphi:A\to B$ 是同态。由于 $D(g)$（$g\in A$）构成 $S$ 的一组拓扑基，可归结到 $U=D(g)$ 的情形。但这时已知
$f^{-1}(U)=D(\varphi(g))$（I，1.2.2.2），命题由此得证。
\end{proof}

\egastatement{1.2.6}{命题}
设 $X$ 是一个在 $S$ 之上仿射的 $S$-概形，$f:X\to S$ 是结构态射。对任意拟凝聚 $\OO_X$-模 $\F$，$f_*(\F)$ 是一个拟凝聚 $\OO_S$-模。
\begin{proof}
结合（1.2.4），这由（I，9.2.2，a）得出。
\end{proof}
特别地，$\OO_S$-代数 $\A(X)=f_*(\OO_X)$ 是拟凝聚的。

\egastatement{1.2.7}{命题}
设 $X$ 是一个在 $S$ 之上仿射的 $S$-概形。对任意 $S$-预概形 $Y$，从集合 $\Hom_S(Y,X)$ 到集合 $\Hom(\A(X),\A(Y))$ 的映射 $h\mapsto\A(h)$（1.1.2）是双射。
\begin{proof}
设 $f:X\to S$、$g:Y\to S$ 是结构态射。先设 $S=\Spec(A)$ 与 $X=\Spec(B)$ 仿射。需要证明：对每个 $\OO_S$-代数同态
$\omega:f_*(\OO_X)\to g_*(\OO_Y)$，存在唯一的 $S$-态射 $h:Y\to X$，使 $\A(h)=\omega$。

按定义，对每个开集 $U\subset S$，$\omega$ 定义一个 $\Gamma(U,\OO_S)$-代数同态
\[
\omega_U=\Gamma(U,\omega):
\Gamma(f^{-1}(U),\OO_X)\longrightarrow\Gamma(g^{-1}(U),\OO_Y).
\]
特别取 $U=S$，便得到 $\Gamma(S,\OO_S)$-代数同态
\[
\varphi:\Gamma(X,\OO_X)\longrightarrow\Gamma(Y,\OO_Y).
\]
由于 $X$ 仿射，它对应于一个完全确定的 $S$-态射 $h:Y\to X$（I，2.2.4）。还需证明 $\A(h)=\omega$，也就是说，对 $S$ 的一组拓扑基中的每个开集 $U$，$\omega_U$ 与 $h$ 的限制
$g^{-1}(U)\to f^{-1}(U)$ 所对应的代数同态 $\varphi_U$ 相同。

可只考虑 $U=D(\lambda)$，其中 $\lambda\in A$。若
$f=({}^{a}\rho,\wt\rho)$，$\rho:A\to B$ 是环同态，则
$f^{-1}(U)=D(\mu)$，其中 $\mu=\rho(\lambda)$，而
$\Gamma(f^{-1}(U),\OO_X)$ 是分式环 $B_\mu$。
\transnote{此处“$\lambda\in A$”依照 $S=\Spec(A)$ 的记号书写；印刷原文写作“$\lambda\in S$”。}
图
\[
\begin{tikzcd}[column sep=large]
B \arrow[r,"\varphi"] \arrow[d] & \Gamma(Y,\OO_Y) \arrow[d]\\
B_\mu \arrow[r,"\varphi_U"'] & \Gamma(g^{-1}(U),\OO_Y)
\end{tikzcd}
\]
交换；将 $\varphi_U$ 换成 $\omega_U$ 后，相应的图也交换。于是由分式环的泛性质（0，1.2.4）得到 $\varphi_U=\omega_U$。

现在转到一般情形。取 $S$ 的仿射开覆盖 $(S_\alpha)$，使各 $f^{-1}(S_\alpha)$ 仿射。任意 $\OO_S$-代数同态
$\omega:\A(X)\to\A(Y)$，经限制给出一族 $\OO_{S_\alpha}$-代数同态
\[
\omega_\alpha:
\A\bigl(f^{-1}(S_\alpha)\bigr)\longrightarrow
\A\bigl(g^{-1}(S_\alpha)\bigr).
\]
由前面的结果，这对应于一族 $S_\alpha$-态射
$h_\alpha:g^{-1}(S_\alpha)\to f^{-1}(S_\alpha)$。
只需证明：对 $S_\alpha\cap S_\beta$ 的一组拓扑基中的每个仿射开集 $U$，$h_\alpha$ 与 $h_\beta$ 在 $g^{-1}(U)$ 上的限制相同。这是显然的，因为根据上面的证明，这两个限制都对应于 $\omega$ 的限制
$\A(X)|U\to\A(Y)|U$。
\end{proof}

\egastatement{1.2.8}{推论}
设 $X,Y$ 是两个在 $S$ 之上仿射的 $S$-概形。$S$-态射 $h:Y\to X$ 是同构，当且仅当
$\A(h):\A(X)\to\A(Y)$ 是同构。
\begin{proof}
这立即由（1.2.7）及 $\A(X)$ 的函子性得出。
\end{proof}

\subsection{与一个 \texorpdfstring{$\OO_S$}{结构层}-代数相伴的、在 \texorpdfstring{$S$}{S} 之上仿射的预概形}
\src{8--9}
\egastatement{1.3.1}{命题}
设 $S$ 是一个预概形。对任意拟凝聚 $\OO_S$-代数 $\B$，存在一个在 $S$ 之上仿射的预概形 $X$，使 $\A(X)=\B$；$X$ 在唯一的 $S$-同构意义下确定。
\begin{proof}
唯一性由（1.2.8）得出。下面证明 $X$ 的存在性。对每个仿射开集 $U\subset S$，令
\[
X_U=\Spec\bigl(\Gamma(U,\B)\bigr).
\]
由于 $\Gamma(U,\B)$ 是 $\Gamma(U,\OO_S)$-代数，$X_U$ 是一个 $S$-预概形（I，1.6.1）。又因为 $\B$ 拟凝聚，$\A(X_U)$ 典范地等同于 $\B|U$（I，1.3.7、1.3.13及1.6.3）。

设 $V$ 是 $S$ 的另一个仿射开集，$f_U:X_U\to S$ 是结构态射，并令 $X_{UV}$ 是 $X_U$ 在 $f_U^{-1}(U\cap V)$ 上诱导的预概形。$X_{UV}$ 与 $X_{VU}$ 都在 $U\cap V$ 之上仿射（1.2.5），且按定义，$\A(X_{UV})$ 与 $\A(X_{VU})$ 都典范地等同于 $\B|(U\cap V)$。于是存在一个典范 $S$-同构（1.2.8）
\[
\theta_{UV}:X_{VU}\longrightarrow X_{UV}.
\]
此外，若 $W$ 是第三个仿射开集，将 $\theta_{UV}$、$\theta_{VW}$、$\theta_{UW}$ 分别限制到 $U\cap V\cap W$ 在 $X_V,X_W,X_W$ 中的逆像，并记作 $\theta'_{UV}$、$\theta'_{VW}$、$\theta'_{UW}$，则有
$\theta'_{UV}\circ\theta'_{VW}=\theta'_{UW}$。

因此存在预概形 $X$、它的一个仿射开覆盖 $(T_U)$，以及对每个 $U$ 的同构 $\varphi_U:X_U\to T_U$，使 $\varphi_U$ 将 $f_U^{-1}(U\cap V)$ 映到 $T_U\cap T_V$，并满足
$\varphi_U\circ\theta_{UV}=\varphi_V$（I，2.3.1）。态射
$g_U=f_U\circ\varphi_U^{-1}$ 使 $T_U$ 成为一个 $S$-预概形；$g_U$ 与 $g_V$ 在 $T_U\cap T_V$ 上相同，因而 $X$ 是一个 $S$-预概形。由定义显然可见，$X$ 在 $S$ 之上仿射，且 $\A(T_U)=\B|U$，所以 $\A(X)=\B$。
\end{proof}
称如此定义的 $S$-概形 $X$ \textbf{与 $\OO_S$-代数 $\B$ 相伴}，或称为 $\B$ 的\textbf{谱}，记作 $\Spec(\B)$。

\egastatement{1.3.2}{推论}
设 $X$ 是一个在 $S$ 之上仿射的预概形，$f:X\to S$ 是结构态射。对任意仿射开集 $U\subset S$，$f^{-1}(U)$ 上诱导的预概形是仿射概形，其环为 $\Gamma(U,\A(X))$。
\begin{proof}
由（1.2.6）及（1.3.1），可设 $X$ 与一个 $\OO_S$-代数相伴；推论于是由（1.3.1）中 $X$ 的构造得出。
\end{proof}

\egastatement{1.3.3}{例}
设 $S$ 是一个域 $K$ 上将原点 $0$ 加倍所得的仿射平面（I，5.5.11）。采用（I，5.5.11）的记号，$S$ 是两个仿射开集 $Y_1,Y_2$ 的并。若 $f$ 是开浸入 $Y_1\to S$，则
$f^{-1}(Y_2)=Y_1\cap Y_2$ 不是 $Y_1$ 中的仿射开集（同上）。因此，这给出了一个本身为仿射概形、但在 $S$ 之上并不仿射的例子。

\egastatement{1.3.4}{推论}
设 $S$ 是一个仿射概形。一个 $S$-预概形 $X$ 在 $S$ 之上仿射，当且仅当 $X$ 是仿射概形。

\egastatement{1.3.5}{推论}
设 $X$ 是一个在预概形 $S$ 之上仿射的预概形，$Y$ 是一个 $X$-预概形。则 $Y$ 在 $S$ 之上仿射，当且仅当 $Y$ 在 $X$ 之上仿射。
\begin{proof}
立即归结到 $S$ 仿射的情形；这时 $X$ 也仿射（1.3.4）。命题中的两个条件便都表示 $Y$ 是一个仿射概形（1.3.4）。
\end{proof}

\egapar{1.3.6}
设 $X$ 是一个在 $S$ 之上仿射的预概形。根据（1.3.5），定义一个在 $X$ 之上仿射的预概形 $Y$，等价于给定一个在 $S$ 之上仿射的预概形 $Y$，以及一个 $S$-态射 $g:Y\to X$。换言之（1.3.1及1.2.7），这等价于给定一个拟凝聚 $\OO_S$-代数 $\B$，以及一个 $\OO_S$-代数同态 $\A(X)\to\B$；后者可看作在 $\B$ 上定义了一个 $\A(X)$-代数结构。若 $f:X\to S$ 是结构态射，则有
\[
\B=f_*\bigl(g_*(\OO_Y)\bigr).
\]

\egastatement{1.3.7}{推论}
设 $X$ 是一个在 $S$ 之上仿射的预概形。则 $X$ 在 $S$ 上为有限型，当且仅当拟凝聚 $\OO_S$-代数 $\A(X)$ 为有限型（I，9.6.2）。
\begin{proof}
由定义（I，9.6.2），可归结到 $S$ 仿射的情形。这时 $X$ 是仿射概形（1.3.4）；若 $S=\Spec(A)$、$X=\Spec(B)$，则 $\A(X)$ 是 $\OO_S$-代数 $\wt B$。由于 $\Gamma(S,\wt B)=B$，推论由（I，9.6.2）及（I，6.3.3）得出。
\transnote{本句印刷原文写作 $\Gamma(U,\wt B)=B$，但前文的仿射底空间记作 $S$；此处将该字母统一为 $S$。}
\end{proof}

\egastatement{1.3.8}{推论}
设 $X$ 是一个在 $S$ 之上仿射的预概形。则 $X$ 约化，当且仅当拟凝聚 $\OO_S$-代数 $\A(X)$ 约化（0，4.1.4）。
\begin{proof}
问题显然在 $S$ 上是局部的；根据（1.3.2），推论由（I，5.1.1）及（I，5.1.4）得出。
\end{proof}

\subsection{在 \texorpdfstring{$S$}{S} 之上仿射的预概形上的拟凝聚层}
\src{9--12}
\egastatement{1.4.1}{命题}
设 $X$ 是一个在 $S$ 之上仿射的预概形，$Y$ 是一个 $S$-预概形，$\F$（分别地，$\G$）是一个拟凝聚 $\OO_X$-模（分别地，$\OO_Y$-模）。则从态射 $(Y,\G)\to(X,\F)$ 的集合到双同态
\[
(\A(X),\A(\F))\longrightarrow(\A(Y),\A(\G))
\]
的集合的映射 $(h,u)\mapsto(\A(h),\A(u))$（1.1.2及1.1.3）是双射。
\begin{proof}
证明与（1.2.7）完全按同样步骤进行，使用（I，2.2.5）及（I，2.2.4）；细节留给读者。
\end{proof}

\egastatement{1.4.2}{推论}
在（1.4.1）的假设之外，再设 $Y$ 在 $S$ 之上仿射，则 $(h,u)$ 是同构，当且仅当 $(\A(h),\A(u))$ 是双同构。

\egastatement{1.4.3}{命题}
对任意二元组 $(\B,\M)$，其中 $\B$ 是一个拟凝聚 $\OO_S$-代数，$\M$ 是一个拟凝聚 $\B$-模（作为 $\OO_S$-模或作为 $\B$-模拟凝聚是等价的，见I，9.6.1），存在二元组 $(X,\F)$，其中 $X$ 是一个在 $S$ 之上仿射的预概形，$\F$ 是拟凝聚 $\OO_X$-模，并满足
$\A(X)=\B$ 及 $\A(\F)=\M$。此外，这个二元组在唯一同构意义下确定。
\begin{proof}
唯一性由（1.4.1）及（1.4.2）得出；存在性与（1.3.1）同样证明，细节仍留给读者。
\end{proof}
将 $\OO_X$-模 $\F$ 记作 $\wt\M$，称它\textbf{与拟凝聚 $\B$-模 $\M$ 相伴}。对每个仿射开集 $U\subset S$，$\wt\M|p^{-1}(U)$（$p$ 是结构态射 $X\to S$）典范地同构于 $\Gamma(U,\M)^\sim$。

\egastatement{1.4.4}{推论}
在拟凝聚 $\B$-模范畴中，$\wt\M$ 是关于 $\M$ 的协变加性正合函子，并与归纳极限及直和交换。
\begin{proof}
立即归结到 $S$ 仿射的情形；这时结论由（I，1.3.5、1.3.9及1.3.11）得出。
\end{proof}

\egastatement{1.4.5}{推论}
在（1.4.3）的假设下，$\wt\M$ 是有限型 $\OO_X$-模，当且仅当 $\M$ 是有限型 $\B$-模。
\begin{proof}
立即归结到 $S=\Spec(A)$ 仿射的情形。于是 $\B=\wt B$，其中 $B$ 是一个有限型 $A$-代数（I，9.6.2），而 $\M=\wt M$，其中 $M$ 是一个 $B$-模（I，1.3.13）。\transnote{原文此处称 $B$ 为“有限型 $A$-代数”，但推论所引用的（1.4.3）未作这一假设；这里保留原文的表述。}
在预概形 $X$ 上，$\OO_X$ 与环 $B$ 相伴，$\wt\M$ 与 $B$-模 $M$ 相伴。故 $\wt\M$ 为有限型，当且仅当 $M$ 为有限型（I，1.3.13），结论由此得出。
\end{proof}

\egastatement{1.4.6}{命题}
设 $Y$ 是一个在 $S$ 之上仿射的预概形，$X,X'$ 是两个在 $Y$ 之上仿射的预概形，因而也在 $S$ 之上仿射（1.3.5）。令
\[
\B=\A(Y),\qquad\A=\A(X),\qquad\A'=\A(X').
\]
则 $X\times_YX'$ 在 $Y$ 之上仿射，因而在 $S$ 之上仿射；$\A(X\times_YX')$ 典范地等同于 $\A\otimes_\B\A'$。
\begin{proof}
$\A\otimes_\B\A'$ 是一个拟凝聚 $\B$-代数，因而也是拟凝聚 $\OO_S$-代数（I，9.6.1）。令 $Z$ 为其谱（1.3.1）。典范 $\B$-同态
$\A\to\A\otimes_\B\A'$ 与 $\A'\to\A\otimes_\B\A'$ 对应于 $Y$-态射
$p:Z\to X$、$p':Z\to X'$（1.2.7）。

为证明 $(Z,p,p')$ 是积 $X\times_YX'$，可只考虑 $S$ 为环 $C$ 的仿射概形的情形（I，3.2.6.4）。这时 $Y,X,X'$ 都是仿射概形（1.3.4），其环 $B,A,A'$ 是 $C$-代数，且
$\B=\wt B$、$\A=\wt A$、$\A'=\wt{A'}$。已知 $\A\otimes_\B\A'$ 典范地等同于 $\OO_S$-代数 $(A\otimes_BA')^\sim$（I，1.3.13）；故 $A(Z)$ 等同于 $A\otimes_BA'$，$p,p'$ 对应于典范同态
$A\to A\otimes_BA'$、$A'\to A\otimes_BA'$。命题于是由（I，3.2.2）得出。
\end{proof}

\egastatement{1.4.7}{推论}
设 $\F$（分别地，$\F'$）是一个拟凝聚 $\OO_X$-模（分别地，$\OO_{X'}$-模）。则
$\A(\F\otimes_Y\F')$ 典范地等同于 $\A(\F)\otimes_{\A(Y)}\A(\F')$。
\begin{proof}
已知 $\F\otimes_Y\F'$ 在 $X\times_YX'$ 上拟凝聚（I，9.1.2）。设
$g:Y\to S$、$f:X\to Y$、$f':X'\to Y$ 是结构态射，并令 $Z=X\times_YX'$。则结构态射 $h:Z\to S$ 同时等于 $g\circ f\circ p$ 和 $g\circ f'\circ p'$。

如下定义典范同态
\[
\A(\F)\otimes_{\A(Y)}\A(\F')\longrightarrow\A(\F\otimes_Y\F').
\]
对任意开集 $U\subset S$，有典范同态（0，4.4.3）
\begin{align*}
\Gamma(f^{-1}(g^{-1}(U)),\F)&\longrightarrow\Gamma(h^{-1}(U),p^*\F),\\
\Gamma(f'^{-1}(g^{-1}(U)),\F')&\longrightarrow\Gamma(h^{-1}(U),p'^*\F').
\end{align*}
由此得到典范同态
\begin{multline*}
\Gamma(f^{-1}(g^{-1}(U)),\F)
\otimes_{\Gamma(g^{-1}(U),\OO_Y)}
\Gamma(f'^{-1}(g^{-1}(U)),\F')\\
\longrightarrow
\Gamma(h^{-1}(U),p^*\F)\otimes_{\Gamma(h^{-1}(U),\OO_Z)}
\Gamma(h^{-1}(U),p'^*\F').
\end{multline*}

为证明所得同态确实是 $\A(Z)$-模同构，可设 $S$，因而 $X,X',Y,X\times_YX'$ 都仿射。采用（1.4.6）的记号，设 $\F=\wt M$、$\F'=\wt{M'}$，其中 $M$ 为 $A$-模，$M'$ 为 $A'$-模。这时 $\F\otimes_Y\F'$ 等同于 $X\times_YX'$ 上与 $(A\otimes_BA')$-模 $M\otimes_BM'$ 相伴的层（I，9.1.3）。将 $M,M',B$ 看作 $C$-模，利用 $\OO_S$-模
$(M\otimes_BM')^\sim$ 与 $\wt M\otimes_{\wt B}\wt{M'}$ 的典范等同（I，1.3.13及1.6.3），便得到推论。
\end{proof}
特别将（1.4.7）用于 $X=Y$、$\F'=\OO_{X'}$，可见 $\A'$-模 $\A(f'^*\F)$ 等同于 $\A(\F)\otimes_\B\A'$。

\egapar{1.4.8}
特别地，令 $X=X'=Y$，其中 $X$ 在 $S$ 之上仿射。若 $\F,\G$ 是两个拟凝聚 $\OO_X$-模，则有典范函子同构
\begin{equation}\tag{1.4.8.1}
\A(\F\otimes_{\OO_X}\G)=\A(\F)\otimes_{\A(X)}\A(\G).
\end{equation}
若再设 $\F$ 有有限表示，则由（I，1.6.3及1.3.12）得到典范同构
\begin{equation}\tag{1.4.8.2}
\A\bigl(\shHom_{\OO_X}(\F,\G)\bigr)
=\shHom_{\A(X)}\bigl(\A(\F),\A(\G)\bigr).
\end{equation}

\egastatement{1.4.9}{注}
若 $X,X'$ 是两个在 $S$ 之上仿射的预概形，则它们的和 $X\amalg X'$ 也在 $S$ 之上仿射，因为两个仿射概形的和仍是仿射概形。

\egastatement{1.4.10}{命题}
设 $S$ 是预概形，$\B$ 是拟凝聚 $\OO_S$-代数，$X=\Spec(\B)$。对 $\B$ 的任意拟凝聚理想层 $\J$，$\wt\J$ 是 $\OO_X$ 的一个拟凝聚理想层；由 $\wt\J$ 定义的 $X$ 的闭子预概形 $Y$，典范地同构于 $\Spec(\B/\J)$。
\begin{proof}
由（I，4.2.3）立即知道 $Y$ 在 $S$ 之上仿射。根据（1.3.1），可归结到 $S$ 仿射的情形；命题于是立即由（I，4.1.2）得出。
\end{proof}
（1.4.10）也可以这样表述：若 $h:\B\to\B'$ 是拟凝聚 $\OO_S$-代数的满同态，则相应的态射
$\Spec(\B')\to\Spec(\B)$ 是闭浸入。

\egastatement{1.4.11}{命题}
设 $S$ 是预概形，$\B$ 是拟凝聚 $\OO_S$-代数，$X=\Spec(\B)$。对 $\OO_S$ 的任意拟凝聚理想 $\J$，若 $f:X\to S$ 是结构态射，则在典范同构意义下有
\[
f^*(\J)\OO_X=(\J\B)^\sim.
\]
\begin{proof}
问题在 $S$ 上是局部的，可归结到 $S=\Spec(A)$ 仿射的情形；这时命题就是（I，1.6.9）。
\end{proof}

\subsection{底预概形的变换}
\src{12--14}
\egastatement{1.5.1}{命题}
设 $X$ 是一个在 $S$ 之上仿射的预概形。对任意底预概形的扩张 $g:S'\to S$，
$X'=X_{(S')}=X\times_SS'$ 在 $S'$ 之上仿射。
\begin{proof}
设 $f'$ 是投影 $X'\to S'$。只需证明：对每个满足 $g(U')$ 包含于 $S$ 的某个仿射开集 $U$ 的仿射开集 $U'\subset S'$，$f'^{-1}(U')$ 仿射（1.2.1）。因此可只考虑 $S,S'$ 均仿射的情形，并只需证明 $X'$ 是仿射概形（1.3.4）。这时 $X$ 也仿射（1.3.4）；若 $A,A',B$ 分别是 $S,S',X$ 的环，则已知 $X'$ 为环 $A'\otimes_AB$ 的仿射概形（I，3.2.2）。
\end{proof}

\egastatement{1.5.2}{推论}
在（1.5.1）的假设下，设 $f:X\to S$ 是结构态射，$f':X'\to S'$、$g':X'\to X$ 是投影，因而图
\[
\begin{tikzcd}
X' \arrow[r,"g'"] \arrow[d,"f'"'] & X \arrow[d,"f"]\\
S' \arrow[r,"g"'] & S
\end{tikzcd}
\]
交换。对任意拟凝聚 $\OO_X$-模 $\F$，存在典范 $\OO_{S'}$-模同构
\begin{equation}\tag{1.5.2.1}
u:g^*\bigl(f_*(\F)\bigr)\xrightarrow{\sim}f'_*\bigl(g'^*(\F)\bigr).
\end{equation}
特别地，$\A(X')$ 与 $g^*(\A(X))$ 之间存在典范同构。
\begin{proof}
为定义 $u$，只需定义同态
\[
v:f_*(\F)\longrightarrow g_*\bigl(f'_*\bigl(g'^*\F\bigr)\bigr)
=f_*\bigl(g'_*\bigl(g'^*\F\bigr)\bigr),
\]
并取 $u=v^\sharp$（0，4.4.3）。我们取 $v=f_*(\rho)$，其中
$\rho:\F\to g'_*(g'^*\F)$ 为典范同态（0，4.4.3）。
为证明 $u$ 是同构，可设 $S,S'$，从而 $X,X'$ 均仿射。采用（1.5.1）的记号，设 $\F=\wt M$，其中 $M$ 为 $B$-模。立即可见，$g^*(f_*\F)$ 和 $f'_*(g'^*\F)$ 都等于与 $A'$-模 $A'\otimes_AM$ 相伴的 $\OO_{S'}$-模，这里将 $M$ 看作 $A$-模；而 $u$ 正是与恒等映射相伴的同态（I，1.6.3、1.6.5及1.6.7）。
\end{proof}

\egastatement{1.5.3}{注}
不能认为：不再假设 $X$ 在 $S$ 之上仿射后，（1.5.2）仍然成立；即使 $S'=\Spec(k(s))$（$s\in S$），且 $S'\to S$ 为典范态射（I，2.4.5），也不能如此。这时 $X'$ 就是纤维 $f^{-1}(s)$（I，3.6.2）。换言之，当 $X$ 不在 $S$ 之上仿射时，“取拟凝聚层的直像”与“取纤维”这两个运算并不交换。不过，我们将在第三章（III，4.2.4）中看到一个这方面的“渐近”结果：当 $f$ 固有（5.4）且 $S$ 诺特时，它对 $X$ 上的凝聚层成立。

\egastatement{1.5.4}{推论}
对任意在 $S$ 之上仿射的预概形 $X$ 及任意 $s\in S$，纤维 $f^{-1}(s)$ 是仿射概形；这里 $f$ 是结构态射 $X\to S$。
\begin{proof}
只需将（1.5.1）用于 $S'=\Spec(k(s))$，再使用（1.3.4）。
\end{proof}

\egastatement{1.5.5}{推论}
设 $X$ 是一个 $S$-预概形，$S'$ 是一个在 $S$ 之上仿射的预概形。则 $X'=X_{(S')}$ 在 $X$ 之上仿射。此外，若 $f:X\to S$ 是结构态射，则有典范 $\OO_X$-代数同构
\[
\A(X')\simeq f^*\bigl(\A(S')\bigr).
\]
对任意拟凝聚 $\A(S')$-模 $\M$，还有典范双同构
\[
f^*(\M)\simeq\A\bigl(f'^*(\wt\M)\bigr),
\]
其中 $f'=f_{(S')}$ 是结构态射 $X'\to S'$。
\begin{proof}
在（1.5.1）及（1.5.2）中交换 $X$ 与 $S'$ 的角色即可。
\end{proof}

\egapar{1.5.6}
现设 $S,S'$ 是两个预概形，$q:S'\to S$ 是态射，$\B$（分别地，$\B'$）是拟凝聚 $\OO_S$-代数（分别地，$\OO_{S'}$-代数），而 $u:\B\to\B'$ 是一个 $q$-态射，即 $\OO_S$-代数同态 $\B\to q_*(\B')$。若
$X=\Spec(\B)$、$X'=\Spec(\B')$，则由此典范地得到态射
\[
v=\Spec(u):X'\longrightarrow X,
\]
使图
\begin{equation}\tag{1.5.6.1}
\begin{tikzcd}
X' \arrow[r,"v"] \arrow[d] & X \arrow[d]\\
S' \arrow[r,"q"'] & S
\end{tikzcd}
\end{equation}
交换，其中竖直箭头为结构态射。

事实上，给定 $u$ 等价于给定拟凝聚 $\OO_{S'}$-代数同态
$u^\sharp:q^*(\B)\to\B'$（0，4.4.3）。它典范地定义一个 $S'$-态射
\[
w:\Spec(\B')\longrightarrow\Spec\bigl(q^*(\B)\bigr),
\]
满足 $\A(w)=u^\sharp$（1.2.7）。另一方面，由（1.5.2），$\Spec(q^*\B)$ 典范地等同于 $X\times_SS'$；$v$ 就是 $w$ 与第一投影的复合
$X'\to X\times_SS'\to X$，而（1.5.6.1）的交换性由定义得出。

设 $U$（分别地，$U'$）是 $S$（分别地，$S'$）的仿射开集，满足 $q(U')\subset U$。记
\[
A=\Gamma(U,\OO_S),\quad A'=\Gamma(U',\OO_{S'}),\quad
B=\Gamma(U,\B),\quad B'=\Gamma(U',\B').
\]
$u$ 的限制作为 $(q|U')$-态射 $\B|U\to\B'|U'$，对应于代数的双同态 $B\to B'$。若 $V,V'$ 分别是 $U,U'$ 在 $X,X'$ 中沿结构态射的逆像，则 $v$ 的限制 $V'\to V$ 对应于上述双同态（I，1.7.3）。

\egapar{1.5.7}
仍采用（1.5.6）的假设，设 $\M$ 是拟凝聚 $\B$-模。则有典范 $\OO_{X'}$-模同构
\begin{equation}\tag{1.5.7.1}
v^*(\wt\M)\simeq\bigl(q^*(\M)\otimes_{q^*(\B)}\B'\bigr)^\sim.
\end{equation}
事实上，典范同构（1.5.2.1）给出 $p^*(\wt\M)$ 与 $\Spec(q^*\B)$ 上同 $q^*\B$-模 $q^*\M$ 相伴的层之间的典范同构；随后应用（1.4.7）即可。

\subsection{仿射态射}
\src{14}
\egapar{1.6.1}
若预概形态射 $f:X\to Y$ 使 $X$ 成为一个在 $Y$ 之上仿射的预概形，则称 $f$ 是\textbf{仿射态射}。一个预概形在另一个预概形之上仿射的诸性质，可用这种语言表述如下。

\egastatement{1.6.2}{命题}
\begin{enumerate}[label=(\roman*)]
\item 闭浸入是仿射的。
\item 两个仿射态射的复合是仿射的。
\item 若 $f:X\to Y$ 是仿射 $S$-态射，则对任意底扩张 $S'\to S$，$f_{(S')}:X_{(S')}\to Y_{(S')}$ 仿射。
\item 若 $f:X\to Y$、$f':X'\to Y'$ 是两个仿射 $S$-态射，则
$f\times_Sf':X\times_SX'\to Y\times_SY'$ 仿射。
\item 若 $f:X\to Y$、$g:Y\to Z$ 是两个态射，$g\circ f$ 仿射且 $g$ 分离，则 $f$ 仿射。
\item 若 $f$ 仿射，则 $f_{\red}$ 也仿射。
\end{enumerate}
\begin{proof}
由（I，5.5.12），只需证明（i）、（ii）及（iii）。（i）就是例（1.2.2），（ii）就是（1.3.5）；最后，（iii）由（1.5.1）得出，因为 $X_{(S')}$ 等同于积 $X\times_YY_{(S')}$（I，3.3.11）。
\end{proof}

\egastatement{1.6.3}{推论}
若 $X$ 是仿射概形而 $Y$ 是概形，则任意态射 $X\to Y$ 都仿射。

\egastatement{1.6.4}{命题}
设 $Y$ 是局部诺特预概形，$f:X\to Y$ 是有限型态射。则 $f$ 仿射，当且仅当 $f_{\red}$ 仿射。
\begin{proof}
根据（1.6.2，vi），只需证明充分性。只需证明：当 $Y$ 仿射且诺特时，$X$ 仿射。这时 $Y_{\red}$ 仿射，故由假设 $X_{\red}$ 也仿射。又 $X$ 诺特，结论便由（I，6.1.7）得出。
\end{proof}

\subsection{与一个模层相伴的向量丛}
\src{14--19}
\egapar{1.7.1}
设 $A$ 是一个环，$E$ 是一个 $A$-模。回顾一下：$E$ 上的\textbf{对称代数}，记作 $\Sym(E)$ 或 $\Sym_A(E)$，是张量代数 $T(E)$ 对由所有元素 $x\otimes y-y\otimes x$（$x,y\in E$）生成的双边理想所取的商代数。

代数 $\Sym(E)$ 由如下泛性质刻画：设 $\sigma:E\to\Sym(E)$ 是典范映射，即 $E\to T(E)$ 与典范映射 $T(E)\to\Sym(E)$ 的复合。对任意交换 $A$-代数 $B$，每个 $A$-线性映射 $E\to B$ 都唯一地分解为
\[
E\xrightarrow{\sigma}\Sym(E)\xrightarrow{g}B,
\]
其中 $g$ 是 $A$-代数同态。由此立即得出：对两个 $A$-模 $E,F$，有典范同构
\[
\Sym(E\oplus F)=\Sym(E)\otimes_A\Sym(F).
\]
此外，$\Sym(E)$ 是从 $A$-模范畴到交换 $A$-代数范畴的协变函子。上述刻画还说明：若 $E=\ilim E_\alpha$，则在典范同构意义下有
$\Sym(E)=\ilim\Sym(E_\alpha)$。

在不会引起混淆时，常将乘积 $\sigma(x_1)\sigma(x_2)\cdots\sigma(x_n)$ 简记为 $x_1x_2\cdots x_n$（$x_i\in E$）。代数 $\Sym(E)$ 是分次的：$\Sym^n(E)$ 由 $E$ 中 $n$ 个元素的乘积的线性组合组成（$n\geq0$）。代数 $\Sym(A)$ 典范地同构于一个不定元的多项式代数 $A[T]$，而 $\Sym(A^n)$ 典范地同构于 $n$ 个不定元的多项式代数 $A[T_1,\ldots,T_n]$。

\egapar{1.7.2}
设 $\varphi:A\to B$ 是环同态。若 $F$ 是 $B$-模，典范映射 $F\to\Sym(F)$ 给出典范映射
$F_{[\varphi]}\to\Sym(F)_{[\varphi]}$，所以它分解为
\[
F_{[\varphi]}\longrightarrow\Sym(F_{[\varphi]})
\longrightarrow\Sym(F)_{[\varphi]}.
\]
典范同态 $\Sym(F_{[\varphi]})\to\Sym(F)_{[\varphi]}$ 是满射，但未必为双射。\transnote{上一句按印刷原文译出；原文未在此处附加关于环同态 $\varphi$ 的满射假设。此项表述留待单独校订，不在译文中暗中补入条件。}
若 $E$ 是 $A$-模，则任意双同态 $E\to F$，即任意 $A$-同态 $E\to F_{[\varphi]}$，典范地给出 $A$-代数同态
\[
\Sym(E)\longrightarrow\Sym(F_{[\varphi]})\longrightarrow\Sym(F)_{[\varphi]},
\]
也就是一个代数的双同态 $\Sym(E)\to\Sym(F)$。

仍用这些记号，对任意 $A$-模 $E$，$\Sym(E\otimes_AB)$ 典范地等同于 $\Sym(E)\otimes_AB$；这立即由 $\Sym(E)$ 的泛性质（1.7.1）得出。

\egapar{1.7.3}
设 $R$ 是环 $A$ 的乘法子集。将（1.7.2）用于 $B=R^{-1}A$，并记住
$R^{-1}E=E\otimes_AR^{-1}A$，可见有典范同构
\[
\Sym(R^{-1}E)=R^{-1}\Sym(E).
\]
若 $R'\supset R$ 是另一个乘法子集，则图
\[
\begin{tikzcd}[column sep=large]
R^{-1}E \arrow[r] \arrow[d] & R'^{-1}E \arrow[d]\\
\Sym(R^{-1}E) \arrow[r] & \Sym(R'^{-1}E)
\end{tikzcd}
\]
交换。

\egapar{1.7.4}
现在设 $(S,\A)$ 是赋环空间，$\E$ 是 $S$ 上的一个 $\A$-模。将每个开集 $U\subset S$ 对应于 $\Gamma(U,\A)$-代数
$\Sym(\Gamma(U,\E))$，由 $\Sym(E)$ 的函子性（1.7.2）定义了一个代数预层。称其相伴层为 $\A$-模 $\E$ 的\textbf{对称 $\A$-代数}，记作 $\Sym(\E)$ 或 $\Sym_\A(\E)$。

由（1.7.1）立即得到，$\Sym(\E)$ 是一个泛问题的解：若 $\B$ 是 $\A$-代数，则任意 $\A$-模同态 $\E\to\B$ 唯一地分解为
$\E\to\Sym(\E)\to\B$，其中第二个箭头是 $\A$-代数同态。因此，$\A$-模同态 $\E\to\B$ 与 $\A$-代数同态 $\Sym(\E)\to\B$ 之间存在一一对应。特别地，每个 $\A$-模同态 $u:\E\to\F$ 定义一个 $\A$-代数同态 $\Sym(u):\Sym(\E)\to\Sym(\F)$，所以 $\Sym(\E)$ 是关于 $\E$ 的协变函子。

根据（1.7.2）以及 $\Sym$ 与归纳极限的交换性，对任意 $s\in S$，有
\[
\bigl(\Sym(\E)\bigr)_s=\Sym(\E_s).
\]
若 $\E,\F$ 是两个 $\A$-模，则 $\Sym(\E\oplus\F)$ 典范地等同于 $\Sym(\E)\otimes_\A\Sym(\F)$，这可从相应的预层看出。

还应注意，$\Sym(\E)$ 是一个分次 $\A$-代数，是各 $\Sym^n(\E)$ 的无限直和；$\A$-模 $\Sym^n(\E)$ 是预层 $U\mapsto\Sym^n(\Gamma(U,\E))$ 的相伴层。特别取 $\E=\A$，可见 $\Sym_\A(\A)$ 等同于
$\A[T]=\A\otimes_{\mathbf Z}\mathbf Z[T]$，其中 $T$ 为不定元，$\mathbf Z$ 看作常层。

\egapar{1.7.5}
设 $(T,\B)$ 是另一个赋环空间，$f:(S,\A)\to(T,\B)$ 是态射。若 $\F$ 是 $\B$-模，则 $\Sym(f^*\F)$ 典范地等同于 $f^*(\Sym(\F))$。事实上，若 $f=(\psi,\theta)$，由定义（0，4.3.1）及（1.7.2），有
\[
\Sym(f^*\F)
=\Sym\bigl(\psi^*\F\otimes_{\psi^*\B}\A\bigr)
=\Sym(\psi^*\F)\otimes_{\psi^*\B}\A.
\]
对 $S$ 的任意开集 $U$ 及 $\Sym(\psi^*\F)$ 在 $U$ 上的任意截面 $h$，在每个点 $s\in U$ 的某个邻域 $V$ 上，$h$ 与 $\Sym(\Gamma(V,\psi^*\F))$ 中的一个元素相同。考察 $\psi^*\F$ 的定义（0，3.7.1），并注意到任意模 $E$ 的 $\Sym(E)$ 中的元素，都是有限多个由 $E$ 的元素构成的乘积的线性组合，可见存在 $T$ 中 $\psi(s)$ 的邻域 $W$、$\Sym(\F)$ 在 $W$ 上的截面 $h'$，以及 $s$ 的邻域 $V'\subset V\cap\psi^{-1}(W)$，使 $h$ 在 $V'$ 上与 $t\mapsto h'_{\psi(t)}$ 相同；我们的断言由此得出。

\egastatement{1.7.6}{命题}
设 $A$ 是环，$S=\Spec(A)$ 是其素谱，$\E=\wt M$ 是与 $A$-模 $M$ 相伴的 $\OO_S$-模。则 $\OO_S$-代数 $\Sym(\E)$ 与 $A$-代数 $\Sym(M)$ 相伴。
\begin{proof}
对每个 $f\in A$，有 $\Sym(M_f)=\Sym(M)_f$（1.7.3）；命题因此由定义（I，1.3.4）得出。
\end{proof}

\egastatement{1.7.7}{推论}
若 $S$ 是预概形，$\E$ 是拟凝聚 $\OO_S$-模，则 $\OO_S$-代数 $\Sym(\E)$ 拟凝聚。若 $\E$ 还是有限型，则各个 $\OO_S$-模 $\Sym^n(\E)$ 都是有限型。
\begin{proof}
第一断言立即由（1.7.6）及（I，1.4.1）得出。第二断言来自如下事实：若 $E$ 是有限型 $A$-模，则 $\Sym^n(E)$ 是有限型 $A$-模；再应用（I，1.3.13）即可。
\end{proof}

\egastatement{1.7.8}{定义}
设 $\E$ 是拟凝聚 $\OO_S$-模。称拟凝聚 $\OO_S$-代数 $\Sym(\E)$ 的谱（1.3.1）为 $\E$ 定义的 $S$ 上的\textbf{向量丛}，记作 $\mathbf V(\E)$。

由（1.2.7），对每个 $S$-预概形 $X$，$S$-态射 $X\to\mathbf V(\E)$ 与 $\OO_S$-代数同态 $\Sym(\E)\to\A(X)$ 之间存在典范一一对应，因而也与 $\OO_S$-模同态
$\E\to\A(X)=f_*(\OO_X)$ 之间存在典范一一对应，其中 $f$ 是结构态射 $X\to S$。特别地，有以下各点。

\egapar{1.7.9}
取 $X$ 为 $S$ 在开集 $U\subset S$ 上诱导的子预概形。则 $S$-态射 $U\to\mathbf V(\E)$，就是 $\mathbf V(\E)$ 在开集 $p^{-1}(U)$ 上诱导的 $U$-预概形的 $U$-截面（I，2.5.5），其中 $p:\mathbf V(\E)\to S$ 是结构态射。由刚才的结果，这些 $U$-截面与 $\OO_S$-模同态
$\E\to j_*(\OO_S|U)$ 一一对应，其中 $j:U\to S$ 是典范嵌入；等价地，它们与 $(\OO_S|U)$-同态
$j^*\E=\E|U\to\OO_S|U$ 一一对应（0，4.4.3）。

此外，将 $S$-态射 $U\to\mathbf V(\E)$ 限制到开集 $U'\subset U$，显然对应于将同态 $\E|U\to\OO_S|U$ 限制到 $U'$。由此得出：$\mathbf V(\E)$ 的 $S$-截面芽所成的层，典范地等同于 $\E$ 的对偶 $\dual\E$。

特别取 $X=U=S$，零同态 $\E\to\OO_S$ 对应于 $\mathbf V(\E)$ 的一个典范 $S$-截面，称为\textbf{零 $S$-截面}（参见8.3.3）。

\egapar{1.7.10}
现在取 $X$ 为域 $K$ 的谱 $\{\xi\}$。结构态射 $f:X\to S$ 对应于一个单同态 $k(s)\to K$，其中 $s=f(\xi)$（I，2.4.6）。于是 $S$-态射 $\{\xi\}\to\mathbf V(\E)$，就是 $\mathbf V(\E)$ 的取值于 $k(s)$ 的扩张 $K$ 的几何点（I，3.4.5），这些点位于 $p^{-1}(s)$ 的点之上。这个集合，也可称为 $\mathbf V(\E)$ 在 $s$ 之上的、在 $K$ 上有理的几何纤维。

根据（1.7.8），它等同于 $\OO_S$-模同态 $\E\to f_*(\OO_X)$ 的集合；等价地，等同于 $\OO_X$-模同态 $f^*\E\to\OO_X=K$ 的集合（0，4.4.3）。而按定义（0，4.3.1），
\[
f^*\E=\E_s\otimes_{\OO_{S,s}}K
=\E^s\otimes_{k(s)}K,\qquad \E^s=\E_s/\mm_s\E_s.
\]
因此，上述在 $K$ 上有理的几何纤维，等同于 $K$-向量空间 $\E^s\otimes_{k(s)}K$ 的对偶。若 $\E^s$ 或 $K$ 在 $k(s)$ 上为有限维，则此对偶也等同于
$\dual{(\E^s)}\otimes_{k(s)}K$，其中 $\dual{(\E^s)}$ 表示 $k(s)$-向量空间 $\E^s$ 的对偶。

\egastatement{1.7.11}{命题}
\begin{enumerate}[label=(\roman*)]
\item $\mathbf V(\E)$ 是关于 $\E$ 的反变函子，从拟凝聚 $\OO_S$-模范畴到仿射 $S$-概形范畴。
\item 若 $\E$ 是有限型 $\OO_S$-模，则 $\mathbf V(\E)$ 在 $S$ 上为有限型。
\item 对两个拟凝聚 $\OO_S$-模 $\E,\F$，$\mathbf V(\E\oplus\F)$ 典范地等同于 $\mathbf V(\E)\times_S\mathbf V(\F)$。
\item 设 $g:S'\to S$ 是态射。对任意拟凝聚 $\OO_S$-模 $\E$，$\mathbf V(g^*\E)$ 典范地等同于
$\mathbf V(\E)_{(S')}=\mathbf V(\E)\times_SS'$。
\item 拟凝聚 $\OO_S$-模的满同态 $\E\to\F$ 对应于闭浸入 $\mathbf V(\F)\to\mathbf V(\E)$。
\end{enumerate}
\begin{proof}
（i）立即由（1.2.7）得出，并使用如下事实：每个 $\OO_S$-模同态 $\E\to\F$ 都典范地定义一个 $\OO_S$-代数同态 $\Sym(\E)\to\Sym(\F)$。
（ii）立即来自定义（I，6.3.1）及如下事实：若 $E$ 为有限型 $A$-模，则 $\Sym(E)$ 为有限型 $A$-代数。
为证明（iii），从典范同构
$\Sym(\E\oplus\F)\simeq\Sym(\E)\otimes_{\OO_S}\Sym(\F)$（1.7.4）出发，应用（1.4.6）即可。
同样，为证明（iv），从典范同构 $\Sym(g^*\E)\simeq g^*(\Sym(\E))$（1.7.5）出发，应用（1.5.2）即可。
最后，为证明（v），只需注意：若 $\E\to\F$ 满，则相应的 $\OO_S$-代数同态 $\Sym(\E)\to\Sym(\F)$ 也满；结论便由（1.4.10）得出。
\end{proof}

\egapar{1.7.12}
特别取 $\E=\OO_S$。预概形 $\mathbf V(\OO_S)$ 是仿射 $S$-概形，它是 $\OO_S$-代数 $\Sym(\OO_S)$ 的谱，而后者等同于
$\OO_S[T]=\OO_S\otimes_{\mathbf Z}\mathbf Z[T]$，其中 $T$ 为不定元。
当 $S=\Spec(\mathbf Z)$ 时，这由（1.7.6）显然成立；考虑结构态射
$S\to\Spec(\mathbf Z)$ 并使用（1.7.11，iv），就从此情形得到一般情形。
根据这一结果，我们还记 $\mathbf V(\OO_S)=S[T]$，并得到公式
\begin{equation}\tag{1.7.12.1}
S[T]=S\otimes_{\mathbf Z}\mathbf Z[T].
\end{equation}
在（I，3.3.15）中已经见到的、将 $S[T]$ 的 $S$-截面芽层与 $\OO_S$ 等同的结果，现在在更一般的背景中重新出现，成为（1.7.9）的特殊情形。

\egapar{1.7.13}
对任意 $S$-预概形 $X$，由（1.7.8）知道，$\Hom_S(X,S[T])$ 典范地等同于
$\Hom_{\OO_S}(\OO_S,\A(X))$，后者典范地同构于 $\Gamma(S,\A(X))$，所以带有环结构。此外，每个 $S$-态射 $h:X\to Y$ 对应于环结构的态射
\[
\Gamma(\A(h)):\Gamma(S,\A(Y))\longrightarrow\Gamma(S,\A(X))
\]
（1.1.2）。当赋予 $\Hom_S(X,S[T])$ 上述环结构时，便可将它看作关于 $X$ 的反变函子，从 $S$-预概形范畴到环范畴。

另一方面，$\Hom_S(X,\mathbf V(\E))$ 同样等同于
$\Hom_{\OO_S}(\E,\A(X))$，这里将 $\A(X)$ 看作 $\OO_S$-模。因此，可典范地赋予它一个 $\Hom_S(X,S[T])$-模结构。与上面同样可见，二元组
\[
\bigl(\Hom_S(X,S[T]),\Hom_S(X,\mathbf V(\E))\bigr)
\]
是关于 $X$ 的反变函子，取值于以二元组 $(A,M)$ 为对象、以双同态为态射的范畴，其中 $A$ 是环，$M$ 是 $A$-模。

我们将上述事实表述为：$S[T]$ 是一个\textbf{环 $S$-概形}，而 $\mathbf V(\E)$ 是环 $S$-概形 $S[T]$ 上的一个\textbf{模 $S$-概形}（参见第0章第8节）。

\egapar{1.7.14}
我们将说明：在 $S$-概形 $\mathbf V(\E)$ 上定义的模 $S$-概形结构，使我们能够在唯一同构意义下重构 $\OO_S$-模 $\E$。为此，将证明 $\E$ 典范地同构于 $\Sym(\E)=\A(\mathbf V(\E))$ 的一个 $\OO_S$-子模，而且这个子模可通过该结构来定义。

事实上（1.7.4），$\OO_S$-代数同态的集合
$\Hom_{\OO_S}(\Sym(\E),\A(X))$ 典范地等同于 $\OO_S$-模同态的集合
$\Hom_{\OO_S}(\E,\A(X))$。若 $h,h'$ 是后一个集合的两个元素，$s_i$（$1\leq i\leq n$）是 $\E$ 在开集 $U\subset S$ 上的截面，$t$ 是 $\A(X)$ 在 $U$ 上的截面，则按定义
\[
(h+h')(s_1\cdots s_n)=\prod_{i=1}^n\bigl(h(s_i)+h'(s_i)\bigr),
\qquad
(th)(s_1\cdots s_n)=t^n\prod_{i=1}^nh(s_i).
\]
于是，若 $z$ 是 $\Sym(\E)$ 在 $U$ 上的截面，则 $h\mapsto h(z)$ 是从
$\Hom_S(X,\mathbf V(\E))=\Hom_{\OO_S}(\Sym(\E),\A(X))$
到 $\Gamma(U,\A(X))$ 的映射。

我们将证明：$\E$ 等同于 $\Sym(\E)$ 中如下子模；对任意开集 $U\subset S$、该子模在 $U$ 上的任意截面 $z$，以及任意 $S$-预概形 $X$，映射
\[
h\longmapsto h(z):
\Hom_{\OO_S|U}\bigl(\Sym(\E)|U,\A(X)|U\bigr)
\longrightarrow\Gamma(U,\A(X))
\]
都是 $\Gamma(U,\A(X))$-模同态。

$\E$ 显然具有这一性质。为证明逆命题，只需证明：当 $S=\Spec(A)$、$\E=\wt M$ 时，$\Sym(\E)$ 在 $S$ 上具有上述性质（取 $U=S$）的截面 $z$ 必是 $\E$ 的截面。此时可写
$z=\sum_{n\geq0}z_n$，其中 $z_n\in\Sym^n(M)$；所需证明的是，当 $n\ne1$ 时 $z_n=0$。

令 $B=\Sym(M)$，取 $X=\Spec(B[T])$，其中 $T$ 为不定元。集合
$\Hom_{\OO_S}(\Sym(\E),\A(X))$ 等同于环同态 $h:B\to B[T]$ 的集合（I，1.3.13）。由上面的讨论，
\[
(Th)(z)=\sum_{n\geq0}T^nh(z_n).
\]
关于 $z$ 的假设说明，对任意同态 $h$，均有
\[
\sum_{n\geq0}T^nh(z_n)=T\sum_{n\geq0}h(z_n).
\]
特别取 $h$ 为典范嵌入，就得到
$\sum_{n\geq0}T^nz_n=T\sum_{n\geq0}z_n$；这确实推出 $z_n=0$（$n\ne1$）。

\egastatement{1.7.15}{命题}
设 $Y$ 是底空间诺特的预概形，或是拟紧概形。任意在 $Y$ 上有限型的仿射 $Y$-概形 $X$，都 $Y$-同构于某个形如 $\mathbf V(\E)$ 的 $Y$-概形的闭 $Y$-子概形，其中 $\E$ 是有限型拟凝聚 $\OO_Y$-模。
\begin{proof}
拟凝聚 $\OO_Y$-代数 $\A(X)$ 为有限型（1.3.7）。所作假设推出 $\A(X)$ 由一个有限型拟凝聚 $\OO_Y$-子模 $\E$ 生成（I，9.6.5）。按定义，这说明将嵌入 $\E\to\A(X)$ 典范地延拓所得的典范同态
$\Sym(\E)\to\A(X)$ 是满射。结论于是由（1.4.10）得出。
\end{proof}

\section{齐次素谱}
\subsection{分次环与分次模的一般性质}
\src{19--23}
\egastatement{2.1.1}{记号}
给定一个非负分次环 $S$，我们用 $S_n$ 表示由次数为 $n$ 的齐次元素组成的部分（$n\geq0$），用 $S_+$ 表示所有 $n>0$ 的 $S_n$ 的直和。于是 $1\in S_0$，$S_0$ 是 $S$ 的子环，$S_+$ 是 $S$ 的分次理想，而 $S$ 是 $S_0$ 与 $S_+$ 的直和。若 $M$ 是 $S$ 上的分次模，允许正次与负次，则同样以 $M_n$ 表示由 $M$ 的 $n$ 次齐次元素组成的 $S_0$-模，其中 $n\in\mathbf Z$。

对每个整数 $d>0$，以 $S^{(d)}$ 表示各 $S_{nd}$ 的直和。将 $S_{nd}$ 的元素看作 $n$ 次齐次元素，这些 $S_{nd}$ 就在 $S^{(d)}$ 上定义了分次环结构。

对每个满足 $0\leq k\leq d-1$ 的整数 $k$，以 $M^{(d,k)}$ 表示各 $M_{nd+k}$（$n\in\mathbf Z$）的直和。将 $M_{nd+k}$ 的元素看作 $n$ 次齐次元素，它就成为分次 $S^{(d)}$-模。当 $k=0$ 时，将 $M^{(d,0)}$ 简记为 $M^{(d)}$。

对每个整数 $n$，无论正负，以 $M(n)$ 表示由
\[
(M(n))_k=M_{n+k}\qquad(k\in\mathbf Z)
\]
定义的分次 $S$-模。特别地，$S(n)$ 是满足 $(S(n))_k=S_{n+k}$ 的分次 $S$-模，其中约定 $S_n=0$（$n<0$）。若分次 $S$-模 $M$ 作为分次模同构于若干形如 $S(n)$ 的模的直和，则称 $M$ \textbf{自由}。由于 $S(n)$ 是单生成 $S$-模，由看作 $-n$ 次元素的 $1\in S$ 生成，这等价于说，$M$ 有一组由齐次元素组成的 $S$-基。

若存在正合列 $P\to Q\to M\to0$，其中 $P,Q$ 都是有限多个形如 $S(n)$ 的模的直和，且各同态均为零次同态（见2.1.2），则称分次 $S$-模 $M$ 有\textbf{有限表示}。

\egapar{2.1.2}
设 $M,N$ 是两个分次 $S$-模。如下在 $M\otimes_SN$ 上定义分次 $S$-模结构。先在 $M\otimes_{\mathbf Z}N$ 上定义分次 $\mathbf Z$-模结构，这里 $\mathbf Z_0=\mathbf Z$，$\mathbf Z_n=0$（$n\ne0$）：令
\[
(M\otimes_{\mathbf Z}N)_q=\bigoplus_{m+n=q}M_m\otimes_{\mathbf Z}N_n.
\]
因为 $M,N$ 分别是各 $M_m,N_n$ 的直和，已知可将 $M\otimes_{\mathbf Z}N$ 典范地等同于所有 $M_m\otimes_{\mathbf Z}N_n$ 的直和。又有
\[
M\otimes_SN=(M\otimes_{\mathbf Z}N)/P,
\]
其中 $P$ 是由 $(xs)\otimes y-x\otimes(sy)$（$x\in M,y\in N,s\in S$）生成的 $\mathbf Z$-子模。显然 $P$ 是分次 $\mathbf Z$-子模；立即可见，取商后确实在 $M\otimes_SN$ 上得到分次 $S$-模结构。

对两个分次 $S$-模 $M,N$，回顾如下定义：若 $S$-模同态 $u:M\to N$ 满足
$u(M_j)\subset N_{j+k}$ 对所有 $j\in\mathbf Z$ 成立，则称 $u$ 的\textbf{次数为 $k$}。令 $H_n$ 为从 $M$ 到 $N$ 的所有 $n$ 次同态所成的集合。我们用 $\Hom_S(M,N)$ 表示各 $H_n$（$n\in\mathbf Z$）在所有 $S$-模同态 $M\to N$ 所成的 $S$-模 $H$ 中的直和。一般而言，$\Hom_S(M,N)$ 不等于 $H$。

不过，当 $M$ 为有限型时，有 $H=\Hom_S(M,N)$。事实上，这时可设 $M$ 由有限多个齐次元素 $x_i$（$1\leq i\leq n$）生成。每个同态 $u\in H$ 唯一地写为 $\sum_{k\in\mathbf Z}u_k$，其中对每个 $k$，$u_k(x_i)$ 是 $u(x_i)$ 的次数为 $k+\deg(x_i)$ 的齐次分量。这推出除有限多个指标外都有 $u_k=0$。按定义，$u_k\in H_k$，所述结论由此得出。

$\Hom_S(M,N)$ 的零次元素称为\textbf{分次 $S$-模同态}。显然有 $S_mH_n\subset H_{m+n}$，所以这些 $H_n$ 在 $\Hom_S(M,N)$ 上定义分次 $S$-模结构。

由这些定义，对任意两个分次 $S$-模 $M,N$，立即有
\begin{align}
M(m)\otimes_SN(n)&=(M\otimes_SN)(m+n),\tag{2.1.2.1}\\
\Hom_S(M(m),N(n))&=(\Hom_S(M,N))(n-m).\tag{2.1.2.2}
\end{align}

设 $S,S'$ 是两个分次环。\textbf{分次环同态} $\varphi:S\to S'$ 是满足
$\varphi(S_n)\subset S'_n$（$n\in\mathbf Z$）的环同态；换言之，$\varphi$ 必须是分次 $\mathbf Z$-模的零次同态。给定这样的同态，就在 $S'$ 上定义了分次 $S$-模结构。连同此结构及其分次环结构，称 $S'$ 为\textbf{分次 $S$-代数}。

若 $M$ 是分次 $S$-模，则分次 $S$-模的张量积 $M\otimes_SS'$ 自然带有分次 $S'$-模结构，分次就是上面定义的分次。

\egastatement{2.1.3}{引理}
设 $S$ 是一个非负分次环。由齐次元素组成的子集 $E\subset S_+$ 作为 $S$-模生成 $S_+$，当且仅当 $E$ 作为 $S_0$-代数生成 $S$。
\begin{proof}
充分性显然。下面证明必要性。令 $E_n$（分别地，$E'_n$）为 $E$ 中次数等于 $n$（分别地，不大于 $n$）的元素的集合。只需对 $n>0$ 归纳证明：$S_n$ 是由 $E'_n$ 中元素的乘积所成的 $n$ 次元素生成的 $S_0$-模。根据假设，当 $n=1$ 时这显然成立。假设还给出
\[
S_n=\sum_{p=0}^{n-1}S_pE_{n-p},
\]
归纳论证于是立即完成。
\end{proof}

\egastatement{2.1.4}{推论}
$S_+$ 为有限型理想，当且仅当 $S$ 为有限型 $S_0$-代数。
\begin{proof}
无论是 $S_0$-代数 $S$ 的有限生成系，还是 $S$-理想 $S_+$ 的有限生成系，都可将每个生成元换成它的齐次分量，从而总可假设生成系由齐次元素组成。
\end{proof}

\egastatement{2.1.5}{推论}
$S$ 诺特，当且仅当 $S_0$ 诺特且 $S$ 为有限型 $S_0$-代数。
\begin{proof}
充分性显然。必要性来自 $S_0\simeq S/S_+$，以及 $S_+$ 必为有限型理想这一事实（2.1.4）。
\end{proof}

\egastatement{2.1.6}{引理}
设 $S$ 是非负分次环，且为有限型 $S_0$-代数。设 $M$ 是有限型分次 $S$-模，则：
\begin{enumerate}[label=(\roman*)]
\item 各 $M_n$ 是有限型 $S_0$-模，且存在整数 $n_0$，使得 $n\leq n_0$ 时有 $M_n=0$。
\item 存在整数 $n_1$ 及整数 $h>0$，使对所有 $n\geq n_1$ 都有 $M_{n+h}=S_hM_n$。
\item 对每对满足 $d>0$、$0\leq k\leq d-1$ 的整数 $(d,k)$，$M^{(d,k)}$ 是有限型 $S^{(d)}$-模。
\item 对每个整数 $d>0$，$S^{(d)}$ 是有限型 $S_0$-代数。
\item 存在整数 $h>0$，使对每个 $m>0$ 都有 $S_{mh}=(S_h)^m$。
\item 对每个整数 $n>0$，存在整数 $m_0$，使 $m\geq m_0$ 时有 $S_m\subset S_+^n$。
\end{enumerate}
\begin{proof}
可设 $S$ 作为 $S_0$-代数，由齐次元素 $f_i$ 生成，其中 $\deg(f_i)=h_i$（$1\leq i\leq r$）；$M$ 作为 $S$-模，由齐次元素 $x_j$ 生成，其中 $\deg(x_j)=k_j$（$1\leq j\leq s$）。显然 $M_n$ 由元素
$f_1^{\alpha_1}\cdots f_r^{\alpha_r}x_j$ 的 $S_0$-线性组合构成，其中各 $\alpha_i\geq0$，并满足
\[
k_j+\sum_i\alpha_ih_i=n.
\]
对每个 $j$，由于各 $h_i>0$，满足此方程的数组 $(\alpha_i)$ 只有有限多个；这证明（i）的第一断言，第二断言显然。

令 $h$ 为各 $h_i$ 的最小公倍数，并令 $g_i=f_i^{h/h_i}$（$1\leq i\leq r$），于是各 $g_i$ 的次数都为 $h$。令 $z_\mu$ 为所有形如
$f_1^{\alpha_1}\cdots f_r^{\alpha_r}x_j$ 且 $0\leq\alpha_i<h/h_i$ 的元素。这些元素只有有限多个；令 $n_1$ 为其次数的最大值。显然，当 $n\geq n_1$ 时，$M_{n+h}$ 中的每个元素都是各 $z_\mu$ 的线性组合，系数是关于各 $g_i$ 的正次单项式。因此
$M_{n+h}=S_hM_n$，证明了（ii）。

同样可见，对任意 $d>0$，$M^{(d,k)}$ 的元素是形如
$g^df_1^{\alpha_1}\cdots f_r^{\alpha_r}x_j$ 的元素的 $S_0$-线性组合，其中 $0\leq\alpha_i<d$，$g$ 是 $S$ 的齐次元素；这给出（iii）。在（iii）及引理（2.1.3）中取 $M=S_+$，并注意 $(S_+)^{(d)}=(S^{(d)})_+$，便得到（iv）。在（ii）中取 $M=S$，可得（v）。

最后，对给定的 $n$，满足 $\alpha_i\geq0$ 及 $\sum_i\alpha_i<n$ 的数组 $(\alpha_i)$ 只有有限多个。若 $m_0$ 是这些数组对应的 $\sum_i\alpha_ih_i$ 的最大值，则 $m>m_0$ 时有 $S_m\subset S_+^n$，这证明（vi）。
\end{proof}

\egastatement{2.1.7}{推论}
若 $S$ 诺特，则对每个整数 $d>0$，$S^{(d)}$ 也诺特。
\begin{proof}
这由（2.1.5）及（2.1.6，iv）得出。
\end{proof}

\egapar{2.1.8}
设 $\pp$ 是分次环 $S$ 的一个分次素理想；于是它是子群
$\pp_n=\pp\cap S_n$ 的直和。假设 $\pp$ 不包含 $S_+$。若 $f\in S_+$ 不属于 $\pp$，则 $f^nx\in\pp$ 等价于 $x\in\pp$。特别地，若 $f\in S_d$（$d>0$），则对每个 $x\in S_{m-nd}$，$f^nx\in\pp_m$ 等价于 $x\in\pp_{m-nd}$。

\egastatement{2.1.9}{命题}
设 $n_0>0$ 为整数，并对每个 $n\geq n_0$ 给定 $S_n$ 的子群 $\pp_n$。存在一个不包含 $S_+$ 的分次素理想 $\pp\subset S$，使所有 $n\geq n_0$ 都满足 $\pp\cap S_n=\pp_n$，当且仅当下列条件成立：
\begin{enumerate}[label=\arabic*${}^{\circ}$]
\item 对每个 $m\geq0$、$n\geq n_0$，$S_m\pp_n\subset\pp_{m+n}$。
\item 若 $m,n\geq n_0$，$f\in S_m$、$g\in S_n$，则 $fg\in\pp_{m+n}$ 推出 $f\in\pp_m$ 或 $g\in\pp_n$。
\item 至少存在一个 $n\geq n_0$，使 $\pp_n\ne S_n$。
\end{enumerate}
此外，这样的分次素理想 $\pp$ 是唯一的。
\begin{proof}
条件1和2的必要性显然。此外，由 $\pp\not\supset S_+$，至少存在一个 $k>0$ 使 $\pp\cap S_k\ne S_k$。若 $f\in S_k$ 不属于 $\pp$，则根据（2.1.8），$\pp\cap S_n=S_n$ 推出
$\pp\cap S_{n-mk}=S_{n-mk}$。因此，若从某个 $n$ 开始恒有 $\pp\cap S_n=S_n$，就会有 $\pp\supset S_+$，与假设矛盾；这证明条件3的必要性。

反过来，假设条件1、2、3成立。注意：若某个 $d\geq n_0$ 及 $f\in S_d$ 满足 $f\notin\pp_d$，则若 $\pp$ 存在，对 $m<n_0$，$\pp_m$ 必须是所有满足如下条件的 $x\in S_m$ 所成的集合：除有限多个 $r$ 外，均有
$f^rx\in\pp_{m+rd}$。这已证明 $\pp$ 若存在就是唯一的。

还需证明：用上述条件定义 $m<n_0$ 时的 $\pp_m$，则
$\pp=\sum_{n\geq0}\pp_n$ 是素理想。先注意，由条件2，对 $m\geq n_0$，$\pp_m$ 也可定义为满足“除有限多个 $r$ 外，$f^rx\in\pp_{m+rd}$”的所有 $x\in S_m$ 的集合。

于是，若 $x\in S_n$、$y\in\pp_m$，则除有限多个 $r$ 外有
$f^rxy\in\pp_{m+n+rd}$，所以 $xy\in\pp_{m+n}$；这证明 $\pp$ 是 $S$ 的理想。为证明它是素理想，即由各 $S_n/\pp_n$ 分次的环 $S/\pp$ 为整环，只需考察两个元素的最高次分量，证明：若 $x\in S_m$、$y\in S_n$ 且 $x\notin\pp_m$、$y\notin\pp_n$，则 $xy\notin\pp_{m+n}$。

否则，对充分大的 $r$，有 $f^{2r}xy\in\pp_{m+n+2rd}$。但对每个 $r>0$，均有 $f^ry\notin\pp_{n+rd}$。由条件2可得，除有限多个 $r$ 外，$f^rx\in\pp_{m+rd}$，进而推出 $x\in\pp_m$，与假设矛盾。
\end{proof}

\egapar{2.1.10}
若 $S_+$ 的子集 $\frj$ 是 $S$ 的理想，则称它为 $S_+$ 的一个\textbf{理想}。若 $\frj$ 是 $S_+$ 与 $S$ 的某个不包含 $S_+$ 的分次素理想的交，则称它为 $S_+$ 的一个\textbf{分次素理想}；根据（2.1.9），这个 $S$ 的素理想是唯一的。

若 $\frj$ 是 $S_+$ 的理想，则 $\frj$ \textbf{在 $S_+$ 中的根}，是 $S_+$ 中某个幂属于 $\frj$ 的元素所成的集合，即
\[
\mathfrak r_+(\frj)=\mathfrak r(\frj)\cap S_+.
\]
特别地，零理想在 $S_+$ 中的根也称为 $S_+$ 的\textbf{幂零根}，记作 $\nrad_+$；它就是 $S_+$ 中所有幂零元素的集合。

若 $\frj$ 是 $S_+$ 的分次理想，则其根 $\mathfrak r_+(\frj)$ 也是分次理想。取商环 $S/\frj$，可归结到 $\frj=0$ 的情形；只需证明：若
$x=x_h+x_{h+1}+\cdots+x_k$ 幂零，则各 $x_i\in S_i$（$h\leq i\leq k$）也幂零。可设 $x_k\ne0$；这时 $x^n$ 的最高次分量为 $x_k^n$，所以 $x_k$ 幂零，再对 $k$ 归纳即可。

若 $\nrad_+=0$，即 $S_+$ 没有非零幂零元素，则称分次环 $S$ \textbf{本质约化}。

\egapar{2.1.11}
注意：若分次环 $S$ 中的元素 $x$ 是零因子，则它的最高次分量也是零因子。若无单位元的环 $S_+$ 非零且没有零因子，则称 $S$ \textbf{本质整}；为此，只需 $S_+$ 中的非零齐次元素在该环中都不是零因子。显然，若 $\pp$ 是 $S_+$ 的分次素理想，则 $S/\pp$ 本质整。

设 $S$ 本质整，$x_0\in S_0$。若存在非零齐次元素 $f\in S_+$ 使 $x_0f=0$，则有 $x_0S_+=0$。事实上，对每个 $g\in S_+$，
$(x_0g)f=(x_0f)g=0$，假设便推出 $x_0g=0$。因此，$S$ 为整环，当且仅当 $S_0$ 为整环且 $S_+$ 在 $S_0$ 中的零化子为零。

\subsection{分次环的分式环}
\src{23--25}
\egapar{2.2.1}
设 $S$ 是非负分次环，$f$ 是 $S$ 中次数为 $d>0$ 的齐次元素。分式环 $S'=S_f$ 也是分次环：令 $S'_n$ 为所有 $x/f^k$ 所成的集合，其中
$x\in S_{n+kd}$、$k\geq0$；这里应注意，$n$ 可取任意负值。以 $S_{(f)}$ 表示由零次元素组成的子环 $S'_0=(S_f)_0$。

若 $f\in S_d$，则 $S_f$ 中的单项式 $(f/1)^k$（$k$ 为任意整数）在环 $S_{(f)}$ 上线性无关，其线性组合所成的集合就是环 $(S_f)^{(d)}$；因此该环同构于
\[
S_{(f)}[T,T^{-1}]=S_{(f)}\otimes_{\mathbf Z}\mathbf Z[T,T^{-1}],
\]
其中 $T$ 为不定元。事实上，若有关系
$\sum_{h=-a}^{b}\xi_h(f/1)^h=0$，其中 $\xi_h=x_h/f^n$、$x_h\in S_{nd}$，则按定义存在 $k\geq a$，使
$\sum_{h=-a}^{b}f^{h+k}x_h=0$。由于此和中各项次数不同，对每个 $h$ 都有 $f^{h+k}x_h=0$，从而 $\xi_h=0$。

若 $M$ 为分次 $S$-模，则 $M'=M_f$ 是分次 $S_f$-模：$M'_n$ 由所有 $x/f^k$ 组成，其中 $x\in M_{n+kd}$、$k\geq0$。以 $M_{(f)}$ 表示 $M'$ 的零次齐次元素所成的集合。立即可见，$M_{(f)}$ 是 $S_{(f)}$-模，并且有
\[
(M_f)^{(d)}=M_{(f)}\otimes_{S_{(f)}}(S_f)^{(d)}.
\]

\egastatement{2.2.2}{引理}
设 $d,e$ 为正整数，$f\in S_d$、$g\in S_e$。存在典范环同构
\[
S_{(fg)}\simeq(S_{(f)})_{g^d/f^e}.
\]
经此同构将两个环典范地等同后，存在典范模同构
\[
M_{(fg)}\simeq(M_{(f)})_{g^d/f^e}.
\]
\begin{proof}
$fg$ 整除 $f^eg^d$，而后者整除 $(fg)^{d+e}$；所以分次环 $S_{fg}$ 与 $S_{f^eg^d}$ 典范地等同。另一方面，$S_{f^eg^d}$ 也等同于 $(S_{f^e})_{g^d/1}$（0，1.4.6）。由于 $f^e/1$ 在 $S_{f^e}$ 中可逆，$S_{f^eg^d}$ 还等同于 $(S_{f^e})_{g^d/f^e}$。元素 $g^d/f^e$ 在 $S_{f^e}$ 中的次数为零，所以后一个环的零次部分就是
$(S_{(f^e)})_{g^d/f^e}$。又显然 $S_{(f^e)}=S_{(f)}$，这证明了第一部分。第二部分同样证明。
\end{proof}

\egapar{2.2.3}
在（2.2.2）的假设下，典范同态 $S_f\to S_{fg}$（0，1.4.1）将 $x/f^k$ 映到 $g^kx/(fg)^k$，显然是零次同态。因此限制到零次部分，得到典范同态 $S_{(f)}\to S_{(fg)}$，使图
\[
\begin{tikzcd}[column sep=large]
& S_{(f)} \arrow[dl] \arrow[dr] &\\
S_{(fg)} \arrow[rr,"\sim"'] && (S_{(f)})_{g^d/f^e}
\end{tikzcd}
\]
交换。同样定义典范同态 $M_{(f)}\to M_{(fg)}$。

\egastatement{2.2.4}{引理}
若 $f,g$ 是 $S_+$ 的两个齐次元素，则环 $S_{(fg)}$ 由 $S_{(f)}$ 与 $S_{(g)}$ 的典范像的并生成。
\begin{proof}
由（2.2.2），只需注意
\[
\frac{1}{g^d/f^e}=\frac{f^{d+e}}{(fg)^d}
\]
属于 $S_{(g)}$ 在 $S_{(fg)}$ 中的典范像；这从定义显然可见。
\end{proof}

\egastatement{2.2.5}{命题}
设 $d>0$ 为整数，$f\in S_d$。存在典范环同构
\[
S_{(f)}\simeq S^{(d)}/(f-1)S^{(d)}.
\]
通过这个同构将两个环等同后，存在典范模同构
\[
M_{(f)}\simeq M^{(d)}/(f-1)M^{(d)}.
\]
\begin{proof}
将 $x/f^n$（$x\in S_{nd}$）对应于 $x$ 模 $(f-1)S^{(d)}$ 的剩余类 $\bar x$，便定义第一个同构。这一映射良定义，因为对每个 $x\in S^{(d)}$，有
$f^hx\equiv x\pmod{(f-1)S^{(d)}}$；所以若某个 $h>0$ 满足 $f^hx=0$，则 $\bar x=0$。

另一方面，若 $x\in S_{nd}$ 满足 $x=(f-1)y$，其中
$y=y_{hd}+y_{(h+1)d}+\cdots+y_{kd}$、$y_{jd}\in S_{jd}$ 且 $y_{hd}\ne0$，则必有 $h=n$、$x=-y_{nd}$，以及关系
$y_{(j+1)d}=fy_{jd}$（$n\leq j\leq k-1$）和 $fy_{kd}=0$。于是某个 $f$ 的幂消去 $x$。
\transnote{按上述递推，直接得到的是 $f^{k-n+1}x=0$；印刷原文在此写作 $f^{k-n}x=0$。本句仅保留其所用的结论“某个幂消去 $x$”，并在此列明指数差异。}
因此，可以将齐次元素 $x\in S_{nd}$ 的模 $(f-1)S^{(d)}$ 的剩余类对应于 $x/f^n\in S_{(f)}$，上面的观察保证它良定义。立即可见，如此得到的两个映射是互逆的环同态。对 $M$ 完全同样处理。
\end{proof}

\egastatement{2.2.6}{推论}
若 $S$ 诺特，则对每个正次齐次元素 $f$，$S_{(f)}$ 也诺特。
\begin{proof}
这立即由（2.1.7）及（2.2.5）得出。
\end{proof}

\egapar{2.2.7}
设 $T$ 是 $S_+$ 中由齐次元素组成的乘法子集，则 $T_0=T\cup\{1\}$ 是 $S$ 的乘法子集。由于 $T_0$ 的元素齐次，环 $T_0^{-1}S$ 以显然的方式分次。以 $S_{(T)}$ 表示其零次子环，即所有形如 $x/h$ 的元素，其中 $h\in T$，$x$ 为与 $h$ 同次的齐次元素。

已知 $T_0^{-1}S$ 典范地等同于各环 $S_f$（$f\in T$）相对于典范同态 $S_f\to S_{fg}$ 的归纳极限（0，1.4.5）。由于这一等同保持次数，它将 $S_{(T)}$ 等同于 $S_{(f)}$（$f\in T$）的归纳极限。对任意分次 $S$-模 $M$，同样定义 $S_{(T)}$-模 $M_{(T)}$，它由 $T_0^{-1}M$ 的零次元素组成；同样可见，它是 $M_{(f)}$（$f\in T$）的归纳极限。

若 $\pp$ 是 $S_+$ 的分次素理想，则分别以 $S_{(\pp)}$、$M_{(\pp)}$ 表示上述 $S_{(T)}$、$M_{(T)}$，其中 $T$ 为 $S_+$ 中不属于 $\pp$ 的齐次元素的集合。

\subsection{分次环的齐次素谱}
\src{25--28}
\egapar{2.3.1}
给定一个非负分次环 $S$，称 $S_+$ 的分次素理想所成的集合（2.1.10）为 $S$ 的\textbf{齐次素谱}，记作 $\Proj(S)$。等价地，它是 $S$ 中不包含 $S_+$ 的分次素理想所成的集合。我们将定义一个以 $\Proj(S)$ 为底集合的概形结构。

\egapar{2.3.2}
对 $S$ 的任意子集 $E$，以 $V_+(E)$ 表示所有包含 $E$ 而不包含 $S_+$ 的分次素理想所成的集合；它就是 $\Spec(S)$ 的子集 $V(E)\cap\Proj(S)$。由（I，1.1.2）得到
\begin{align}
V_+(0)&=\Proj(S),\qquad V_+(S)=V_+(S_+)=\varnothing,\tag{2.3.2.1}\\
V_+\Bigl(\bigcup_\lambda E_\lambda\Bigr)&=\bigcap_\lambda V_+(E_\lambda),\tag{2.3.2.2}\\
V_+(EE')&=V_+(E)\cup V_+(E').\tag{2.3.2.3}
\end{align}
将 $E$ 换成由 $E$ 生成的分次理想，不改变 $V_+(E)$。此外，若 $\frj$ 是 $S$ 的分次理想，则对每个 $n>0$，有
\begin{equation}\tag{2.3.2.4}
V_+(\frj)=V_+\Bigl(\bigcup_{q\geq n}(\frj\cap S_q)\Bigr).
\end{equation}
事实上，若 $\pp\in\Proj(S)$ 包含 $\frj$ 中所有次数不小于 $n$ 的齐次元素，由于存在齐次元素 $f\in S_+$ 不属于 $\pp$，对任意 $m\geq0$ 及 $x\in S_m\cap\frj$，除有限多个 $r$ 外，均有
$f^rx\in\frj\cap S_{m+rd}$，所以 $f^rx\in\pp\cap S_{m+rd}$；这推出 $x\in\pp\cap S_m$（2.1.9）。
最后，对 $S$ 的任意分次理想 $\frj$，有
\begin{equation}\tag{2.3.2.5}
V_+(\frj)=V_+\bigl(\mathfrak r_+(\frj)\bigr).
\end{equation}

\egapar{2.3.3}
按定义，$V_+(E)$ 是 $X=\Proj(S)$ 在由 $\Spec(S)$ 的谱拓扑诱导的拓扑下的闭子集；这个拓扑仍称为 $X$ 上的\textbf{谱拓扑}。对每个 $f\in S$，令
\begin{equation}\tag{2.3.3.1}
D_+(f)=D(f)\cap\Proj(S)=\Proj(S)-V_+(f).
\end{equation}
于是，对任意 $f,g\in S$，有（I，1.1.9.1）
\begin{equation}\tag{2.3.3.2}
D_+(fg)=D_+(f)\cap D_+(g).
\end{equation}

\egastatement{2.3.4}{命题}
当 $f$ 遍历 $S_+$ 中的齐次元素时，各 $D_+(f)$ 构成 $X=\Proj(S)$ 的一组拓扑基。
\begin{proof}
由（2.3.2.2）及（2.3.2.4），每个 $X$ 的闭子集都是形如 $V_+(f)$ 的集合的交，其中 $f$ 为正次齐次元素。
\end{proof}

\egapar{2.3.5}
设 $f$ 是 $S_+$ 中次数为 $d>0$ 的齐次元素。对 $S$ 的每个不包含 $f$ 的分次素理想 $\pp$，已知所有形如 $x/f^n$（$x\in\pp,n\geq0$）的元素所成的集合，是分式环 $S_f$ 的素理想（0，1.2.6）。它与 $S_{(f)}$ 的交因而是 $S_{(f)}$ 的素理想，记作 $\psi_f(\pp)$；它由 $x/f^n$（$n\geq0,x\in\pp\cap S_{nd}$）组成。这样便定义了映射
\[
\psi_f:D_+(f)\longrightarrow\Spec(S_{(f)}).
\]
若 $g\in S_e$ 是 $S_+$ 的第二个齐次元素，则图
\begin{equation}\tag{2.3.5.1}
\begin{tikzcd}[column sep=large]
D_+(f) \arrow[r,"\psi_f"] & \Spec(S_{(f)})\\
D_+(fg) \arrow[r,"\psi_{fg}"'] \arrow[u,hook] & \Spec(S_{(fg)}) \arrow[u,"{}^{a}\omega"']
\end{tikzcd}
\end{equation}
交换，其中左边的竖直箭头为包含映射，右边的竖直箭头是由典范同态
$\omega=\omega_{fg,f}:S_{(f)}\to S_{(fg)}$ 诱导的映射（I，1.2.1）。事实上，若 $x/f^n\in\omega^{-1}(\psi_{fg}(\pp))$ 且 $f\notin\pp$，则按定义
$xg^n/(fg)^n\in\psi_{fg}(\pp)$，所以 $xg^n\in\pp$，进而 $x\in\pp$；逆向推理显然。

\egastatement{2.3.6}{命题}
映射 $\psi_f$ 是从 $D_+(f)$ 到 $\Spec(S_{(f)})$ 的同胚。
\begin{proof}
首先，$\psi_f$ 连续。若 $h\in S_{nd}$，则按定义，$h/f^n\in\psi_f(\pp)$ 当且仅当 $h\in\pp$，所以
\[
\psi_f^{-1}\bigl(D(h/f^n)\bigr)=D_+(hf).
\]
连续性由（2.3.3.2）得出。此外，当 $h$ 遍历各 $S_{nd}$ 时，$D_+(hf)$ 构成 $D_+(f)$ 的一组拓扑基，这是（2.3.4）及（2.3.3.2）的结果。结合 $D_+(f)$ 与 $\Spec(S_{(f)})$ 都满足公理 $(T_0)$，以上说明 $\psi_f$ 单射，且其逆映射
$\psi_f(D_+(f))\to D_+(f)$ 连续。

最后证明满射。给定 $S_{(f)}$ 的素理想 $\qq_0$，对每个 $n\geq0$，令 $\pp_n$ 为满足 $x^d/f^n\in\qq_0$ 的所有 $x\in S_n$。这些 $\pp_n$ 满足（2.1.9）的条件。事实上，若 $x,y\in S_n$ 且
$x^d/f^n,y^d/f^n\in\qq_0$，则 $(x+y)^{2d}/f^{2n}\in\qq_0$；由于 $\qq_0$ 为素理想，得到 $(x+y)^d/f^n\in\qq_0$。这证明各 $\pp_n$ 是 $S_n$ 的子群。利用 $\qq_0$ 为素理想，其余条件的验证立即完成。

若 $\pp$ 是如此定义的分次素理想，则确有 $\psi_f(\pp)=\qq_0$：对 $x\in S_{nd}$，由于 $\qq_0$ 为素理想，$x/f^n\in\qq_0$ 与 $x^d/f^{nd}\in\qq_0$ 等价。
\end{proof}

\egastatement{2.3.7}{推论}
$D_+(f)=\varnothing$，当且仅当 $f$ 幂零。
\begin{proof}
$\Spec(S_{(f)})=\varnothing$ 等价于 $S_{(f)}=0$，也等价于 $S_f$ 中 $1=0$；按定义，这恰好表示 $f$ 幂零。
\end{proof}

\egastatement{2.3.8}{推论}
设 $E\subset S_+$。下列条件等价：
\begin{enumerate}[label=\alph*)]
\item $V_+(E)=X=\Proj(S)$；
\item $E$ 的每个元素都幂零；
\item $E$ 的每个元素的所有齐次分量都幂零。
\end{enumerate}
\begin{proof}
显然 c）推出 b），b）推出 a）。若 $\frj$ 是 $E$ 生成的 $S$ 的分次理想，条件 a）等价于 $V_+(\frj)=X$。因此，对 $\frj$ 中的每个齐次元素 $f$，有 $V_+(f)=X$，由（2.3.7）知 $f$ 幂零。
\end{proof}

\egastatement{2.3.9}{推论}
若 $\frj$ 是 $S_+$ 的分次理想，则 $\mathfrak r_+(\frj)$ 是所有包含 $\frj$ 的 $S_+$ 的分次素理想的交。
\begin{proof}
考虑分次环 $S/\frj$，可归结到 $\frj=0$ 的情形。需要证明：若 $f\in S_+$ 不幂零，则 $S$ 中存在不包含 $f$ 的分次素理想。由于 $f$ 至少有一个齐次分量不幂零，可设 $f$ 齐次；断言于是由（2.3.7）得出。
\end{proof}

\egapar{2.3.10}
对 $X=\Proj(S)$ 的任意子集 $Y$，以 $\frj_+(Y)$ 表示满足 $Y\subset V_+(f)$ 的所有 $f\in S_+$ 所成的集合；等价地，
$\frj_+(Y)=\frj(Y)\cap S_+$。所以 $\frj_+(Y)$ 是 $S_+$ 的一个理想，且等于它在 $S_+$ 中的根。

\egastatement{2.3.11}{命题}
\begin{enumerate}[label=(\roman*)]
\item 对 $S_+$ 的任意子集 $E$，$\frj_+(V_+(E))$ 是由 $E$ 生成的分次理想在 $S_+$ 中的根。
\item 对 $X$ 的任意子集 $Y$，$V_+(\frj_+(Y))=\overline Y$，其中 $\overline Y$ 是 $Y$ 在 $X$ 中的闭包。
\end{enumerate}
\begin{proof}
（i）若 $\frj$ 是 $E$ 生成的 $S_+$ 的分次理想，则 $V_+(E)=V_+(\frj)$，结论由（2.3.9）得出。

（ii）由于 $V_+(\frj)=\bigcap_{f\in\frj}V_+(f)$，$Y\subset V_+(\frj)$ 推出每个 $f\in\frj$ 都满足 $Y\subset V_+(f)$，所以 $\frj_+(Y)\supset\frj$，从而
$V_+(\frj_+(Y))\subset V_+(\frj)$。由闭集的定义，这证明（ii）。
\end{proof}

\egastatement{2.3.12}{推论}
$X=\Proj(S)$ 的闭子集 $Y$，与 $S_+$ 中等于其自身根的分次理想 $\frj$，经反序映射
$Y\mapsto\frj_+(Y)$、$\frj\mapsto V_+(\frj)$ 一一对应。
两个闭子集的并 $Y_1\cup Y_2$ 对应于 $\frj_+(Y_1)\cap\frj_+(Y_2)$；任意一族闭子集 $(Y_\lambda)$ 的交，对应于各 $\frj_+(Y_\lambda)$ 的和在 $S_+$ 中的根。

\egastatement{2.3.13}{推论}
设 $\frj$ 为 $S_+$ 的分次理想。则 $V_+(\frj)=\varnothing$，当且仅当 $S_+$ 的每个元素都有某个幂属于 $\frj$。

这个推论也可用以下等价形式表述。

\egastatement{2.3.14}{推论}
设 $(f_\alpha)$ 是 $S_+$ 的一族齐次元素。各 $D_+(f_\alpha)$ 构成 $X=\Proj(S)$ 的覆盖，当且仅当 $S_+$ 的每个元素都有某个幂属于各 $f_\alpha$ 生成的理想。

\egastatement{2.3.15}{推论}
设 $(f_\alpha)$ 是 $S_+$ 的一族齐次元素，$f\in S_+$。下列关系等价：
\begin{enumerate}[label=\alph*)]
\item $D_+(f)\subset\bigcup_\alpha D_+(f_\alpha)$；
\item $V_+(f)\supset\bigcap_\alpha V_+(f_\alpha)$；
\item $f$ 的某个幂属于各 $f_\alpha$ 生成的理想。
\end{enumerate}

\egastatement{2.3.16}{推论}
$X=\Proj(S)$ 为空，当且仅当 $S_+$ 的每个元素都幂零。

\egastatement{2.3.17}{推论}
在（2.3.12）所述的一一对应中，不可约闭子集对应于 $S_+$ 的分次素理想。
\begin{proof}
若 $Y=Y_1\cup Y_2$，其中 $Y_1,Y_2$ 闭且都不等于 $Y$，则
$\frj_+(Y)=\frj_+(Y_1)\cap\frj_+(Y_2)$，而两个理想 $\frj_+(Y_1),\frj_+(Y_2)$ 都不等于 $\frj_+(Y)$，故 $\frj_+(Y)$ 不是素理想。

反过来，若 $\frj$ 是 $S_+$ 的非素分次理想，则存在 $f,g\in S_+$，满足
$f,g\notin\frj$ 而 $fg\in\frj$。于是 $V_+(f),V_+(g)$ 均不包含 $V_+(\frj)$，但根据（2.3.2.3），
$V_+(\frj)\subset V_+(f)\cup V_+(g)$。由此得出，$V_+(\frj)$ 是两个闭集
$V_+(f)\cap V_+(\frj)$ 与 $V_+(g)\cap V_+(\frj)$ 的并，而且这两个闭集都不等于 $V_+(\frj)$。
\end{proof}

\subsection{\texorpdfstring{$\Proj(S)$}{Proj(S)} 上的概形结构}
\src{28--30}
\egapar{2.4.1}
设 $f,g$ 为 $S_+$ 的两个齐次元素。考虑仿射概形
\[
Y_f=\Spec(S_{(f)}),\qquad Y_g=\Spec(S_{(g)}),\qquad
Y_{fg}=\Spec(S_{(fg)}).
\]
由（2.2.2），对应于典范同态
$\omega_{fg,f}:S_{(f)}\to S_{(fg)}$ 的态射
$u_{f,fg}=({}^{a}\omega_{fg,f},\wt\omega_{fg,f}):Y_{fg}\to Y_f$ 是开浸入（I，1.3.6）。利用同胚
$\psi_f:D_+(f)\to Y_f$（2.3.6）的逆映射，可将 $Y_f$ 的仿射概形结构搬到 $D_+(f)$ 上。由图（2.3.5.1）的交换性，仿射概形 $D_+(fg)$ 便等同于仿射概形 $D_+(f)$ 在其底空间的开集 $D_+(fg)$ 上诱导的概形。

结合（2.3.4），显然 $X=\Proj(S)$ 带有唯一的预概形结构，使得在每个 $D_+(f)$ 上的限制，都是刚刚定义的仿射概形。进一步，有如下结果。

\egastatement{2.4.2}{命题}
预概形 $\Proj(S)$ 是一个概形。
\begin{proof}
只需证明（I，5.5.6）：对任意齐次元素 $f,g\in S_+$，
$D_+(f)\cap D_+(g)=D_+(fg)$ 仿射，且它的环由 $D_+(f)$ 与 $D_+(g)$ 的环的典范像生成。第一点由定义显然成立，第二点由（2.2.4）得出。
\end{proof}
当把齐次素谱 $\Proj(S)$ 称作一个概形时，总是指刚才定义的这个结构。

\egastatement{2.4.3}{例}
取 $S=K[T_1,T_2]$，其中 $K$ 为域，$T_1,T_2$ 为两个不定元，$S$ 按总次数分次。由（2.3.14），$\Proj(S)$ 是 $D_+(T_1)$ 与 $D_+(T_2)$ 的并。立即可见，这两个仿射概形都典范地同构于 $K[T]$，而 $\Proj(S)$ 是按（I，2.3.2）所述方式将这两个仿射概形粘合得到的（参见7.4.14）。

\egastatement{2.4.4}{命题}
设 $S$ 为非负分次环，$X=\Proj(S)$。
\begin{enumerate}[label=(\roman*)]
\item 若 $\nrad_+$ 是 $S_+$ 的幂零根（2.1.10），则概形 $X_{\red}$ 典范地同构于 $\Proj(S/\nrad_+)$。特别地，若 $S$ 本质约化，则 $\Proj(S)$ 约化。
\item 假设 $S$ 本质约化。则 $X$ 为整概形，当且仅当 $S$ 本质整。
\end{enumerate}
\begin{proof}
（i）令 $\overline S=S/\nrad_+$，以 $x\mapsto\bar x$ 表示典范零次同态 $S\to\overline S$。对每个 $f\in S_d$（$d>0$），典范同态 $S_f\to\overline S_{\bar f}$（0，1.5.1）满且为零次，故限制后给出满同态
$S_{(f)}\to\overline S_{(\bar f)}$。若 $f\notin\nrad_+$，立即可验证 $\overline S_{\bar f}$ 约化，而这个同态的核为 $S_{(f)}$ 的幂零根，也就是说
$\overline S_{(\bar f)}=(S_{(f)})_{\red}$。因此，对应的闭浸入
$D_+(\bar f)\to D_+(f)$ 将 $D_+(\bar f)$ 等同于 $(D_+(f))_{\red}$（I，5.1.2），特别地，它是底空间之间的同胚。

若 $g\notin\nrad_+$ 是 $S_+$ 的另一个齐次元素，则图
\[
\begin{tikzcd}[column sep=large]
S_{(f)} \arrow[r] \arrow[d] & \overline S_{(\bar f)} \arrow[d]\\
S_{(fg)} \arrow[r] & \overline S_{(\overline{fg})}
\end{tikzcd}
\]
交换。又各 $D_+(f)$，其中 $f$ 为不属于 $\nrad_+$ 的正次齐次元素，覆盖 $X=\Proj(S)$（2.3.7）。所以态射
$D_+(\bar f)\to D_+(f)$ 是某个闭浸入 $\Proj(\overline S)\to\Proj(S)$ 的限制，且后者是底空间的同胚；结论由（I，5.1.2）得出。

（ii）若 $S$ 本质整，即 $(0)$ 是 $S_+$ 的一个与 $S_+$ 不同的分次素理想，则由（i），$X$ 约化；由（2.3.17），$X$ 不可约。
反过来，设 $S$ 本质约化且 $X$ 为整概形。则对每个非零齐次元素 $f\in S_+$，$D_+(f)\ne\varnothing$（2.3.7）。$X$ 的不可约性推出，对非零齐次元素 $f,g\in S_+$，有
$D_+(f)\cap D_+(g)\ne\varnothing$。由（2.3.3.2）知 $fg\ne0$，因而 $S_+$ 没有零因子，结论得证。
\end{proof}

\egapar{2.4.5}
给定交换环 $A$，回顾如下定义：若分次环 $S$ 带有 $A$-代数结构，使每个子群 $S_n$ 都是 $A$-模，则称 $S$ 是\textbf{分次 $A$-代数}。事实上，只需 $S_0$ 是 $A$-代数即可；换言之，只要在 $S_0$ 上给定 $A$-代数结构，并对 $a\in A,x\in S_n$ 令
$a\cdot x=(a\cdot1)x$，就定义了 $S$ 上的分次 $A$-代数结构。

\egastatement{2.4.6}{命题}
假设 $S$ 为分次 $A$-代数。则在 $X=\Proj(S)$ 上，结构层 $\OO_X$ 是一个 $A$-代数，这里将 $A$ 看作 $X$ 上的常层；换言之，$X$ 是 $\Spec(A)$ 之上的概形。
\begin{proof}
只需注意：对 $S_+$ 中的每个齐次元素 $f$，$S_{(f)}$ 是 $A$-代数；而对齐次元素 $f,g\in S_+$，图
\[
\begin{tikzcd}[column sep=large]
S_{(f)} \arrow[rr] && S_{(fg)}\\
& A \arrow[ul] \arrow[ur] &
\end{tikzcd}
\]
交换。
\end{proof}

\egastatement{2.4.7}{命题}
设 $S$ 为非负分次环。
\begin{enumerate}[label=(\roman*)]
\item 对每个整数 $d>0$，概形 $\Proj(S)$ 到概形 $\Proj(S^{(d)})$ 有一个典范同构。
\item 令 $S'$ 为满足 $S'_0=\mathbf Z$，以及 $S'_n=S_n$（$n>0$，看作 $\mathbf Z$-模）的分次环。则 $\Proj(S)$ 到 $\Proj(S')$ 有一个典范同构。
\end{enumerate}
\begin{proof}
（i）先证明映射 $\pp\mapsto\pp\cap S^{(d)}$ 是从集合 $\Proj(S)$ 到 $\Proj(S^{(d)})$ 的双射。给定分次素理想 $\pp'\in\Proj(S^{(d)})$，令
$\pp_{nd}=\pp'\cap S_{nd}$（$n\geq0$）。对不是 $d$ 的倍数的 $n>0$，将 $\pp_n$ 定义为满足 $x^d\in\pp'\cap S_{nd}$ 的所有 $x\in S_n$。
若 $x^d,y^d\in\pp'\cap S_{nd}$，则 $(x+y)^{2d}\in\pp'\cap S_{2nd}$；由于 $\pp'$ 为素理想，得到 $(x+y)^d\in\pp'\cap S_{nd}$。立即可见，如此定义的各 $\pp_n$ 满足（2.1.9）的条件，故存在唯一的 $\pp\in\Proj(S)$，使 $\pp\cap S^{(d)}=\pp'$。

对 $S_+$ 中任意齐次元素 $f$，有 $V_+(f)=V_+(f^d)$（2.3.2.3），所以这个双射是拓扑空间的同胚。最后，采用同样记号，$S_{(f^d)}$ 与 $S^{(d)}_{(f^d)}$ 典范地等同（2.2.2）；故 $\Proj(S)$ 与 $\Proj(S^{(d)})$ 作为概形典范地等同。

（ii）将每个 $\pp\in\Proj(S)$ 对应于唯一的 $\pp'\in\Proj(S')$，使所有 $n>0$ 都满足 $\pp'\cap S'_n=\pp\cap S_n$（2.1.9）。显然，这定义一个底空间的典范同胚
$\Proj(S)\simeq\Proj(S')$，因为当 $f$ 是 $S_+$ 中齐次元素时，对 $S,S'$ 而言，$V_+(f)$ 是同一个集合。又有 $S'_{(f)}=S_{(f)}$，因此两者作为概形等同。
\end{proof}

\egastatement{2.4.8}{推论}
若 $S$ 是分次 $A$-代数，$S'_A$ 是满足
$(S'_A)_0=A$、$(S'_A)_n=S_n$（$n>0$）的分次 $A$-代数，则存在典范同构
$\Proj(S)\simeq\Proj(S'_A)$。
\begin{proof}
采用（2.4.7，ii）的记号，这两个概形都典范地同构于 $\Proj(S')$。
\end{proof}

\subsection{与一个分次模相伴的层}
\src{30--36}
\egapar{2.5.1}
设 $M$ 是一个分次 $S$-模。对 $S_+$ 中的每个齐次元素 $f$，$M_{(f)}$ 是 $S_{(f)}$-模，因而在仿射概形 $\Spec(S_{(f)})$ 上有相伴的拟凝聚层 $(M_{(f)})^\sim$（I，1.3.4）；这个仿射概形已等同于 $D_+(f)$。

\egastatement{2.5.2}{命题}
在 $X=\Proj(S)$ 上，存在唯一的拟凝聚 $\OO_X$-模 $\wt M$，使对每个齐次元素 $f\in S_+$，有
\[
\Gamma(D_+(f),\wt M)=M_{(f)},
\]
而对于齐次元素 $f,g\in S_+$，限制同态
$\Gamma(D_+(f),\wt M)\to\Gamma(D_+(fg),\wt M)$ 是典范同态
$M_{(f)}\to M_{(fg)}$（2.2.3）。
\begin{proof}
设 $f\in S_d,g\in S_e$。由（2.2.2），$D_+(fg)$ 等同于
$(S_{(f)})_{g^d/f^e}$ 的素谱。因此，$D_+(f)$ 上的层 $(M_{(f)})^\sim$ 限制到 $D_+(fg)$，典范地等同于与模
$(M_{(f)})_{g^d/f^e}$ 相伴的层（I，1.3.6），也就等同于 $(M_{(fg)})^\sim$（2.2.2）。由此得到典范同构
\[
\theta_{fg}:(M_{(g)})^\sim|D_+(fg)\xrightarrow{\sim}
(M_{(f)})^\sim|D_+(fg).
\]
若 $h$ 是 $S_+$ 的第三个齐次元素，则在 $D_+(fgh)$ 上有
$\theta_{fh}=\theta_{fg}\circ\theta_{gh}$。

于是（0，3.3.1），在 $X$ 上存在拟凝聚 $\OO_X$-模 $\F$，以及对每个齐次元素 $f\in S_+$ 的同构
$\eta_f:\F|D_+(f)\to(M_{(f)})^\sim$，使 $\theta_{fg}=\eta_f\circ\eta_g^{-1}$。

现在考虑 $X$ 的拓扑基 $(D_+(f))$ 上由
$D_+(f)\mapsto M_{(f)}$ 及典范限制同态 $M_{(f)}\to M_{(fg)}$ 所定义的预层，并令 $\G$ 为其相伴层。上述结果结合（I，1.3.7），证明 $\F$ 与 $\G$ 同构。将 $\G$ 记作 $\wt M$，它确实满足命题中的条件。特别地，$\wt S=\OO_X$。
\end{proof}

\egastatement{2.5.3}{定义}
称（2.5.2）中定义的拟凝聚 $\OO_X$-模 $\wt M$ \textbf{与分次 $S$-模 $M$ 相伴}。

回顾一下：将分次模同态限于零次同态时，分次 $S$-模构成一个范畴。采用这个约定，有如下结果。

\egastatement{2.5.4}{命题}
关于 $M$ 的函子 $M\mapsto\wt M$ 是从分次 $S$-模范畴到拟凝聚 $\OO_X$-模范畴的协变加性正合函子，并与归纳极限及直和交换。
\begin{proof}
这些性质都是局部的，只需在层
$\wt M|D_+(f)=(M_{(f)})^\sim$ 上验证。函子 $M\mapsto M_f$、分次 $S_f$-模范畴中的函子 $N\mapsto N_0$，以及 $S_{(f)}$-模范畴中的函子 $P\mapsto\wt P$，都正合并与归纳极限及直和交换（I，1.3.5及1.3.9）；命题由此得出。
\end{proof}
将对应于零次同态 $u:M\to N$ 的同态 $\wt M\to\wt N$ 记为 $\wt u$。由（2.5.4）立即得到，（I，1.3.9及1.3.10）的结果对分次 $S$-模及零次同态仍成立，这里 $\wt M$ 取本节所赋予的意义；因为那些证明都是纯形式的。

\egastatement{2.5.5}{命题}
对每个 $\pp\in X=\Proj(S)$，有
\[
(\wt M)_\pp=M_{(\pp)}.
\]
\begin{proof}
按定义，$(\wt M)_\pp=\ilim\Gamma(D_+(f),\wt M)$，其中 $f$ 遍历 $S_+$ 中不属于 $\pp$ 的齐次元素。由于 $\Gamma(D_+(f),\wt M)=M_{(f)}$，命题由 $M_{(\pp)}$ 的定义（2.2.7）得出。
\end{proof}
特别地，局部环 $(\OO_X)_\pp$ 就是 $S_{(\pp)}$，即由分式 $x/f$ 构成的环，其中 $f$ 是 $S_+$ 中不属于 $\pp$ 的齐次元素，而 $x$ 是与 $f$ 同次的齐次元素。

更特别地，若 $S$ 本质整，则 $X=\Proj(S)$ 为整概形（2.4.4）。若 $\xi=(0)$ 是 $X$ 的泛点，则有理函数域
$R(X)=\OO_{X,\xi}$ 是由 $fg^{-1}$ 构成的域，其中 $f,g$ 是 $S_+$ 中同次的齐次元素，且 $g\ne0$。

\egastatement{2.5.6}{命题}
若对每个 $z\in M$ 及 $S_+$ 中每个齐次元素 $f$，都存在 $f$ 的某个幂消去 $z$，则 $\wt M=0$。当 $S_0$-代数 $S$ 由一次齐次元素的集合 $S_1$ 生成时，这个充分条件也是必要条件。
\begin{proof}
条件 $\wt M=0$ 等价于对 $S_+$ 中每个齐次元素 $f$ 都有 $M_{(f)}=0$。另一方面，若 $f\in S_d$，则 $M_{(f)}=0$ 表示：对每个次数为 $d$ 的倍数的齐次元素 $z\in M$，存在 $n$ 使 $f^nz=0$。因此，若对每个 $f\in S_1$ 都有 $M_{(f)}=0$，命题的条件就对每个 $f\in S_1$ 成立。若 $S_1$ 生成 $S$，则条件对 $S_+$ 中任意齐次元素也成立，因为这些元素都是 $S_1$ 中元素的乘积的线性组合。
\end{proof}

\egastatement{2.5.7}{命题}
设 $d>0$ 为整数，$f\in S_d$。则对每个 $n\in\mathbf Z$，
$(\OO_X|D_+(f))$-模 $(S(nd))^\sim|D_+(f)$ 典范地同构于 $\OO_X|D_+(f)$。
\begin{proof}
乘以 $S_f$ 中的可逆元素 $(f/1)^n$，给出双射
\[
S_{(f)}=(S_f)_0\xrightarrow{\sim}(S_f)_{nd}
=(S_f(nd))_0=(S(nd)_f)_0=S(nd)_{(f)}.
\]
换言之，$S_{(f)}$-模 $S_{(f)}$ 与 $S(nd)_{(f)}$ 典范地同构，命题由此得出。
\end{proof}

\egastatement{2.5.8}{推论}
在开集 $U=\bigcup_{f\in S_d}D_+(f)$ 上，$\OO_X$-模 $(S(nd))^\sim$ 的限制是一个可逆 $(\OO_X|U)$-模（0，5.4.1）。

\egastatement{2.5.9}{推论}
若理想 $S_+$ 由一次齐次元素的集合 $S_1$ 生成，则对每个 $n\in\mathbf Z$，$\OO_X$-模 $(S(n))^\sim$ 可逆。
\begin{proof}
由假设，$X=\bigcup_{f\in S_1}D_+(f)$（2.3.14）；在（2.5.8）中取 $U=X$ 即可。
\end{proof}

\egapar{2.5.10}
在本节其余部分，对每个 $n\in\mathbf Z$，令
\begin{equation}\tag{2.5.10.1}
\OO_X(n)=(S(n))^\sim.
\end{equation}
对 $X$ 的每个开集 $U$ 及每个 $(\OO_X|U)$-模 $\F$，令
\begin{equation}\tag{2.5.10.2}
\F(n)=\F\otimes_{\OO_X}(\OO_X(n)|U)\qquad(n\in\mathbf Z).
\end{equation}
若 $S_+$ 由 $S_1$ 生成，则对每个 $n\in\mathbf Z$，关于 $\F$ 的函子 $\F(n)$ 正合，因为这时 $\OO_X(n)$ 是可逆 $\OO_X$-模。

\egapar{2.5.11}
设 $M,N$ 为两个分次 $S$-模。对每个 $f\in S_d$（$d>0$），如下定义典范函子性 $S_{(f)}$-模同态
\begin{equation}\tag{2.5.11.1}
\lambda_f:M_{(f)}\otimes_{S_{(f)}}N_{(f)}\longrightarrow(M\otimes_SN)_{(f)}.
\end{equation}
将由嵌入 $M_{(f)}\to M_f$、$N_{(f)}\to N_f$ 及 $S_{(f)}\to S_f$ 得到的同态
$M_{(f)}\otimes_{S_{(f)}}N_{(f)}\to M_f\otimes_{S_f}N_f$，与典范同构
$M_f\otimes_{S_f}N_f\simeq(M\otimes_SN)_f$（0，1.3.4）复合。按两个分次模的张量积的定义，后一个同构保持次数。因此，对 $x\in M_{md}$、$y\in N_{nd}$，$m,n\geq0$，有
\[
\lambda_f\bigl((x/f^m)\otimes(y/f^n)\bigr)
=(x\otimes y)/f^{m+n}.
\]
由此定义立即得到：对 $g\in S_e$（$e>0$），图
\[
\begin{tikzcd}[column sep=large]
M_{(f)}\otimes_{S_{(f)}}N_{(f)} \arrow[r,"\lambda_f"] \arrow[d] & (M\otimes_SN)_{(f)} \arrow[d]\\
M_{(fg)}\otimes_{S_{(fg)}}N_{(fg)} \arrow[r,"\lambda_{fg}"'] & (M\otimes_SN)_{(fg)}
\end{tikzcd}
\]
交换，其中右侧竖直箭头为典范同态，左侧竖直箭头来自典范同态。因此，各 $\lambda_f$ 给出一个典范函子性 $\OO_X$-模同态
\begin{equation}\tag{2.5.11.2}
\lambda:\wt M\otimes_{\OO_X}\wt N\longrightarrow(M\otimes_SN)^\sim.
\end{equation}

特别考虑 $S$ 的两个分次理想 $\frj,\frk$。由于 $\wt\frj,\wt\frk$ 是 $\OO_X$ 的理想层，有典范同态
$\wt\frj\otimes_{\OO_X}\wt\frk\to\OO_X$。这时图
\begin{equation}\tag{2.5.11.3}
\begin{tikzcd}[column sep=large]
\wt\frj\otimes_{\OO_X}\wt\frk \arrow[rr,"\lambda"] \arrow[dr] && (\frj\otimes_S\frk)^\sim \arrow[dl]\\
& \OO_X &
\end{tikzcd}
\end{equation}
交换。归结到每个 $D_+(f)$（$f$ 为 $S_+$ 中齐次元素）上验证，即立即由 $\lambda_f$ 的定义（2.5.11.1）及（I，1.3.13）得出。

最后，若 $M,N,P$ 是三个分次 $S$-模，则图
\begin{equation}\tag{2.5.11.4}
\begin{tikzcd}[column sep=large]
\wt M\otimes\wt N\otimes\wt P \arrow[r,"\lambda\otimes1"] \arrow[d,"1\otimes\lambda"'] & (M\otimes_SN)^\sim\otimes\wt P \arrow[d,"\lambda"]\\
\wt M\otimes(N\otimes_SP)^\sim \arrow[r,"\lambda"'] & (M\otimes_SN\otimes_SP)^\sim
\end{tikzcd}
\end{equation}
交换，图中未标底环的层张量积均取在 $\OO_X$ 上。仍只需在每个 $D_+(f)$ 上验证，便立即由定义及（I，1.3.13）得出。

\egapar{2.5.12}
在（2.5.11）的假设下，如下定义典范函子性 $S_{(f)}$-模同态
\begin{equation}\tag{2.5.12.1}
\mu_f:(\Hom_S(M,N))_{(f)}\longrightarrow
\Hom_{S_{(f)}}(M_{(f)},N_{(f)}).
\end{equation}
将 $u/f^n$，其中 $u$ 是 $nd$ 次同态，对应于把 $x/f^m$（$x\in M_{md}$）映到
$u(x)/f^{m+n}$ 的同态 $M_{(f)}\to N_{(f)}$。对 $g\in S_e$（$e>0$），仍有交换图
\[
\begin{tikzcd}[column sep=large]
(\Hom_S(M,N))_{(f)} \arrow[r,"\mu_f"] \arrow[d] & \Hom_{S_{(f)}}(M_{(f)},N_{(f)}) \arrow[d]\\
(\Hom_S(M,N))_{(fg)} \arrow[r,"\mu_{fg}"'] & \Hom_{S_{(fg)}}(M_{(fg)},N_{(fg)})
\end{tikzcd}
\]
其中左侧箭头为典范同态，右侧箭头来自典范同态。结合（I，1.3.8），这些 $\mu_f$ 定义典范函子性 $\OO_X$-模同态
\begin{equation}\tag{2.5.12.2}
\mu:(\Hom_S(M,N))^\sim\longrightarrow\shHom_{\OO_X}(\wt M,\wt N).
\end{equation}

\egastatement{2.5.13}{命题}
假设理想 $S_+$ 由 $S_1$ 生成。则同态 $\lambda$（2.5.11.2）是同构。若分次 $S$-模 $M$ 有有限表示（2.1.1），同态 $\mu$（2.5.12.2）也为同构。
\begin{proof}
由于 $X$ 是 $D_+(f)$（$f\in S_1$）的并（2.3.14），可归结到在 $f$ 齐次且次数为1时，证明在所述假设下 $\lambda_f,\mu_f$ 是同构。

先定义 $\mathbf Z$-双线性映射
\[
M_m\times N_n\longrightarrow M_{(f)}\otimes_{S_{(f)}}N_{(f)},
\quad(x,y)\longmapsto(x/f^m)\otimes(y/f^n),
\]
其中当 $m<0$ 时，用 $x/f^m$ 表示 $f^{-m}x/1$。这些映射定义 $\mathbf Z$-线性映射
$M\otimes_{\mathbf Z}N\to M_{(f)}\otimes_{S_{(f)}}N_{(f)}$。
若 $s\in S_q$，则 $(sx)\otimes y$ 映到
\[
(s/f^q)\bigl((x/f^m)\otimes(y/f^n)\bigr)
\qquad(x\in M_m,y\in N_n).
\]
于是得到模的双同态
$\gamma_f:M\otimes_SN\to M_{(f)}\otimes_{S_{(f)}}N_{(f)}$，它相对于将 $s\in S_q$ 映到 $s/f^q$ 的典范环同态 $S\to S_{(f)}$。

若 $\sum_i(x_i\otimes y_i)\in M\otimes_SN$，其中 $x_i,y_i$ 分别齐次、次数为 $m_i,n_i$，且某个 $r$ 满足
$f^r\sum_i(x_i\otimes y_i)=0$，也就是说 $\sum_i(f^rx_i\otimes y_i)=0$，则由（0，1.3.4），
\[
\sum_i\left(\frac{f^rx_i}{f^{m_i+r}}\otimes\frac{y_i}{f^{n_i}}\right)=0.
\]
这就是 $\gamma_f(\sum_i x_i\otimes y_i)=0$。因此 $\gamma_f$ 分解为
\[
M\otimes_SN\longrightarrow(M\otimes_SN)_f
\xrightarrow{\gamma'_f}M_{(f)}\otimes_{S_{(f)}}N_{(f)}.
\]
令 $\lambda'_f$ 为 $\gamma'_f$ 在 $(M\otimes_SN)_{(f)}$ 上的限制，立即可验证 $\lambda_f$ 与 $\lambda'_f$ 是互逆的 $S_{(f)}$-同态，证明了第一部分。

为证明第二部分，设 $M$ 是分次 $S$-模同态 $P\to Q$ 的余核，其中 $P,Q$ 是有限多个形如 $S(n)$ 的模的直和。利用 $\Hom_S(L,N)$ 关于 $L$ 的左正合性，以及 $M_{(f)}$ 关于 $M$ 的正合性，立即归结到 $M=S(n)$ 时证明 $\mu_f$ 为同构。

对每个齐次元素 $z\in N$，令 $u_z:S(n)\to N$ 为满足 $u_z(1)=z$ 的同态。立即可见，$\eta:z\mapsto u_z$ 是从 $N(-n)$ 到 $\Hom_S(S(n),N)$ 的零次同构。它对应于同构
\[
\eta_f:(N(-n))_{(f)}\longrightarrow(\Hom_S(S(n),N))_{(f)}.
\]
另一方面，令
\[
\eta'_f:N_{(f)}\longrightarrow\Hom_{S_{(f)}}(S(n)_{(f)},N_{(f)})
\]
为将 $z'\in N_{(f)}$ 对应于同态 $v_{z'}$ 的同构，其中
$v_{z'}(s/f^k)=sz'/f^{n+k}$，$s\in S_{n+k}=(S(n))_k$。容易验证，复合
\begin{multline*}
(N(-n))_{(f)}\xrightarrow{\eta_f}(\Hom_S(S(n),N))_{(f)}
\xrightarrow{\mu_f}\Hom_{S_{(f)}}(S(n)_{(f)},N_{(f)})\\
\xrightarrow{\eta_f'^{-1}}N_{(f)}
\end{multline*}
是同构 $z/f^h\mapsto z/f^{h-n}$。所以 $\mu_f$ 是同构。
\end{proof}

若 $S_+$ 由 $S_1$ 生成，由（2.5.13），对 $S$ 的任意分次理想 $\frj$ 及任意分次 $S$-模 $M$，有典范同构
\begin{equation}\tag{2.5.13.1}
\wt\frj\cdot\wt M=(\frj M)^\sim.
\end{equation}
事实上，这由图
\[
\begin{tikzcd}[column sep=large]
\wt\frj\otimes_{\OO_X}\wt M \arrow[rr,"\lambda"] \arrow[dr] && (\frj\otimes_SM)^\sim \arrow[dl]\\
& \wt M &
\end{tikzcd}
\]
的交换性得出，其验证与（2.5.11.3）相同。

\egastatement{2.5.14}{推论}
假设 $S$ 由 $S_1$ 生成。对任意 $m,n\in\mathbf Z$，有典范同构
\begin{align}
\OO_X(m)\otimes_{\OO_X}\OO_X(n)&=\OO_X(m+n),\tag{2.5.14.1}\\
\OO_X(n)&=\OO_X(1)^{\otimes n}.\tag{2.5.14.2}
\end{align}
\begin{proof}
第一个公式由（2.5.13）及典范零次同构
$S(m)\otimes_SS(n)\to S(m+n)$ 得出：该同构将 $1\otimes1$ 映到
$1\in(S(m+n))_{-m-n}$，其中两个因子 $1$ 分别位于 $(S(m))_{-m}$ 与 $(S(n))_{-n}$。
随后，只需在 $n=-1$ 时证明第二个公式。根据（2.5.13），这归结为证明
$\Hom_S(S(1),S)$ 典范地同构于 $S(-1)$；回到定义（2.1.2），并记住 $S(1)$ 是单生成 $S$-模，即可立即验证。
\end{proof}

\egastatement{2.5.15}{推论}
假设 $S$ 由 $S_1$ 生成。对任意分次 $S$-模 $M$ 及任意 $n\in\mathbf Z$，有典范同构
\begin{equation}\tag{2.5.15.1}
(M(n))^\sim=\wt M(n).
\end{equation}
\begin{proof}
这来自定义（2.5.10.2）、（2.5.10.1）、命题（2.5.13），以及典范零次同构
$M(n)\simeq M\otimes_SS(n)$；后者将
$z\in(M(n))_h=M_{n+h}$ 映到
\[
z\otimes1\in M_{n+h}\otimes(S(n))_{-n}\subset(M\otimes_SS(n))_h.
\]
\end{proof}

\egapar{2.5.16}
令 $S'$ 为满足 $S'_0=\mathbf Z$、$S'_n=S_n$（$n>0$）的分次环。对 $f\in S_d$（$d>0$）及任意 $n\in\mathbf Z$，有
$(S(n))_{(f)}=(S'(n))_{(f)}$。事实上，后者的元素形如 $x/f^k$，其中 $x\in S_{n+kd}$、$k>0$，而总能选取 $k$ 使 $n+kd\ne0$。由于
$X=\Proj(S)$ 与 $X'=\Proj(S')$ 典范地等同（2.4.7，ii），可见在此等同下，$\OO_X(n)$ 与 $\OO_{X'}(n)$ 相互对应。

另一方面，对每个 $d>0$、$n\in\mathbf Z$，有
\[
(S^{(d)}(n))_k=S_{(n+k)d}=(S(nd))_{kd}.
\]
所以对 $f\in S_d$，有
$(S^{(d)}(n))_{(f)}=(S(nd))_{(f)}$。已知概形 $X=\Proj(S)$ 与
$X^{(d)}=\Proj(S^{(d)})$ 典范地等同（2.4.7，i）；上述结果说明：若
$S_0$-代数 $S^{(d)}$ 由 $S_d$ 生成，则在此等同下，对任意 $n\in\mathbf Z$，$\OO_X(nd)$ 与 $\OO_{X^{(d)}}(n)$ 相互对应。

\egastatement{2.5.17}{命题}
设 $d>0$ 为整数，$U=\bigcup_{f\in S_d}D_+(f)$。对每个整数 $n$，典范同态
\[
\OO_X(nd)\otimes_{\OO_X}\OO_X(-nd)\longrightarrow\OO_X
\]
在 $U$ 上的限制是同构。
\begin{proof}
由（2.5.16），可只考虑 $d=1$；这时结论由（2.5.13）的证明得出。
\end{proof}

\subsection{与 \texorpdfstring{$\Proj(S)$}{Proj(S)} 上的层相伴的分次 \texorpdfstring{$S$}{S}-模}
\src{36--38}
本小节始终假设理想 $S_+$ 由一次齐次元素的集合 $S_1$ 生成。

\egapar{2.6.1}
这时 $\OO_X$-模 $\OO_X(1)$ 可逆（2.5.9）。对任意 $\OO_X$-模 $\F$，结合（2.5.14.2），令（0，5.4.6）
\begin{equation}\tag{2.6.1.1}
\Gstar(\F)=\Gamma_*(\OO_X(1),\F)
=\bigoplus_{n\in\mathbf Z}\Gamma(X,\F(n)).
\end{equation}
回顾（0，5.4.6），$\Gstar(\OO_X)$ 带有分次环结构，而 $\Gstar(\F)$ 带有分次 $\Gstar(\OO_X)$-模结构。

由于 $\OO_X(n)$ 局部自由，$\Gstar(\F)$ 是关于 $\F$ 的协变加性左正合函子。特别地，若 $\J$ 是 $\OO_X$ 的理想层，则 $\Gstar(\J)$ 典范地等同于 $\Gstar(\OO_X)$ 的一个分次理想。

\egapar{2.6.2}
设 $M$ 为分次 $S$-模。对每个 $f\in S_d$（$d>0$），$x\mapsto x/1$ 是阿贝尔群同态
$M_0\to M_{(f)}$。由于 $M_{(f)}$ 典范地等同于 $\Gamma(D_+(f),\wt M)$，由此得到阿贝尔群同态
$\alpha_0^f:M_0\to\Gamma(D_+(f),\wt M)$。显然，对每个 $g\in S_e$（$e>0$），图
\[
\begin{tikzcd}[column sep=large]
& \Gamma(D_+(f),\wt M) \arrow[dd]\\
M_0 \arrow[ur,"\alpha_0^f"] \arrow[dr,"\alpha_0^{fg}"'] &\\
& \Gamma(D_+(fg),\wt M)
\end{tikzcd}
\]
交换。这表示对每个 $x\in M_0$，截面 $\alpha_0^f(x)$ 与 $\alpha_0^g(x)$ 在
$D_+(f)\cap D_+(g)$ 上相同。因此存在唯一截面
$\alpha_0(x)\in\Gamma(X,\wt M)$，使它在每个 $D_+(f)$ 上的限制为 $\alpha_0^f(x)$。这样便定义了阿贝尔群同态
\begin{equation}\tag{2.6.2.1}
\alpha_0:M_0\longrightarrow\Gamma(X,\wt M),
\end{equation}
其定义并未使用 $S$ 由 $S_1$ 生成的假设。

对每个 $n\in\mathbf Z$，将此结果用于分次模 $M(n)$，结合（2.5.15），得到阿贝尔群同态
\begin{equation}\tag{2.6.2.2}
\alpha_n:M_n=(M(n))_0\longrightarrow\Gamma(X,\wt M(n)).
\end{equation}
由此得到分次阿贝尔群的函子性零次同态
\begin{equation}\tag{2.6.2.3}
\alpha:M\longrightarrow\Gstar(\wt M),
\end{equation}
也记作 $\alpha_M$，它在每个 $M_n$ 上等于 $\alpha_n$。

特别取 $M=S$，结合 $\Gstar(\OO_X)$ 中乘法的定义（0，5.4.6），立即可验证
$\alpha:S\to\Gstar(\OO_X)$ 为分次环同态；对每个分次 $S$-模 $M$，
（2.6.2.3）是分次模的双同态。

\egastatement{2.6.3}{命题}
对每个 $f\in S_d$（$d>0$），$D_+(f)$ 恰好是 $\OO_X(d)$ 的截面 $\alpha_d(f)$ 不消失的点 $\pp\in X$ 的集合（0，5.5.2）。
\begin{proof}
由假设 $X=\bigcup_{g\in S_1}D_+(g)$，只需对每个 $g\in S_1$ 验证：
$D_+(g)$ 中 $\alpha_d(f)$ 不消失的点集恰好为 $D_+(fg)$。
按定义，$\alpha_d(f)$ 在 $D_+(g)$ 上的限制，是对应于 $(S(d))_{(g)}$ 中元素 $f/1$ 的截面。经过典范同构
$(S(d))_{(g)}\simeq S_{(g)}$（2.5.7），它等同于 $\OO_X$ 在 $D_+(g)$ 上对应于
$f/g^d\in S_{(g)}$ 的截面。

说这个截面在 $\pp\in D_+(g)$ 处消失，表示 $f/g^d\in\qq$，其中 $\qq$ 是与 $\pp$ 对应的 $S_{(g)}$ 的素理想（2.3.6）。按定义，这等价于 $f\in\pp$，命题得证。
\end{proof}

\egapar{2.6.4}
现在设 $\F$ 是 $\OO_X$-模，令 $M=\Gstar(\F)$。利用分次环同态
$\alpha:S\to\Gstar(\OO_X)$，可将 $M$ 看作分次 $S$-模。对每个 $f\in S_d$（$d>0$），由（2.6.3），$\OO_X(d)$ 的截面 $\alpha_d(f)$ 限制到 $D_+(f)$ 后可逆。于是对每个 $n>0$，$\OO_X(nd)$ 的截面 $\alpha_{nd}(f^n)$ 限制到 $D_+(f)$ 后也可逆。

设 $z\in M_{nd}=\Gamma(X,\F(nd))$，$n\geq0$。若存在整数 $k\geq0$，使
$f^kz$ 在 $D_+(f)$ 上的限制为零，即 $\F((n+k)d)$ 的截面
\[
(z|D_+(f))\otimes(\alpha_{kd}(f^k)|D_+(f))
\]
为零，则由上述可逆性，$z|D_+(f)=0$。因此，将 $z/f^n$ 对应于 $\F$ 在 $D_+(f)$ 上的截面
\[
(z|D_+(f))\bigl(\alpha_{nd}(f^n)|D_+(f)\bigr)^{-1},
\]
就定义了 $S_{(f)}$-同态 $\beta_f:M_{(f)}\to\Gamma(D_+(f),\F)$。
立即可验证，对 $g\in S_e$（$e>0$），图
\begin{equation}\tag{2.6.4.1}
\begin{tikzcd}[column sep=large]
M_{(f)} \arrow[r,"\beta_f"] \arrow[d] & \Gamma(D_+(f),\F) \arrow[d]\\
M_{(fg)} \arrow[r,"\beta_{fg}"'] & \Gamma(D_+(fg),\F)
\end{tikzcd}
\end{equation}
交换。

记住 $M_{(f)}$ 典范地等同于 $\Gamma(D_+(f),\wt M)$，以及各 $D_+(f)$ 构成 $X$ 的拓扑基（2.3.4），可见这些 $\beta_f$ 来自唯一的典范 $\OO_X$-模同态
\begin{equation}\tag{2.6.4.2}
\beta:(\Gstar(\F))^\sim\longrightarrow\F,
\end{equation}
也记作 $\beta_\F$；它显然具有函子性。

\egastatement{2.6.5}{命题}
设 $M$ 为分次 $S$-模，$\F$ 为 $\OO_X$-模。则以下复合都是恒等同构：
\begin{align}
\wt M\xrightarrow{\wt\alpha_M}(\Gstar(\wt M))^\sim
\xrightarrow{\beta_{\wt M}}\wt M,\tag{2.6.5.1}\\
\Gstar(\F)\xrightarrow{\alpha_{\Gstar(\F)}}
\Gstar\bigl((\Gstar(\F))^\sim\bigr)
\xrightarrow{\Gstar(\beta_\F)}\Gstar(\F).\tag{2.6.5.2}
\end{align}
\begin{proof}
（2.6.5.1）的验证是局部的。在开集 $D_+(f)$ 中，它立即来自定义，以及如下事实：作用于拟凝聚层时，$\beta$ 由它对 $D_+(f)$ 上截面的作用确定（I，1.3.8）。

（2.6.5.2）逐次验证。令 $M=\Gstar(\F)$，则
$M_n=\Gamma(X,\F(n))$，并且
\[
(\Gstar(\wt M))_n=\Gamma(X,\wt M(n))
=\Gamma(X,(M(n))^\sim).
\]
若 $f\in S_1$、$z\in M_n$，则 $\alpha_n^f(z)$ 为 $(M(n))_{(f)}$ 中的元素 $z/1$，它等于 $(f/1)^n(z/f^n)$。在 $\beta_f$ 下，它对应于 $D_+(f)$ 上的截面
\[
\bigl(\alpha_1(f)^n|D_+(f)\bigr)
\Bigl((z|D_+(f))\bigl(\alpha_1(f)^n|D_+(f)\bigr)^{-1}\Bigr),
\]
也就是 $z$ 在 $D_+(f)$ 上的限制；这便验证了（2.6.5.2）。
\end{proof}

\subsection{有限性条件}
\src{38--41}
\egastatement{2.7.1}{命题}
\begin{enumerate}[label=(\roman*)]
\item 若 $S$ 为诺特分次环，则 $X=\Proj(S)$ 为诺特概形。
\item 若 $S$ 为有限型分次 $A$-代数，则 $X=\Proj(S)$ 是 $Y=\Spec(A)$ 上的有限型概形。
\end{enumerate}
\begin{proof}
（i）若 $S$ 诺特，则理想 $S_+$ 有有限齐次生成系 $f_i$（$1\leq i\leq p$）。由（2.3.14），$X$ 是各
$D_+(f_i)=\Spec(S_{(f_i)})$ 的并。因此，只需各 $S_{(f_i)}$ 诺特；这由（2.2.6）得出。

（ii）假设推出 $S_0$ 为有限型 $A$-代数，$S$ 为有限型 $S_0$-代数，所以 $S_+$ 为有限型理想（2.1.4）。与（i）同样，归结到对每个 $f\in S_d$，证明 $S_{(f)}$ 为有限型 $A$-代数。由（2.2.5），只需证明 $S^{(d)}$ 为有限型 $A$-代数，而这由（2.1.6）得出。
\end{proof}

\egapar{2.7.2}
以下对分次 $S$-模 $M$ 考察这两个有限性条件：
\begin{description}[leftmargin=3.5em,style=nextline]
\item[(TF)] 存在整数 $n$，使子模 $\bigoplus_{k\geq n}M_k$ 为有限型 $S$-模。
\item[(TN)] 存在整数 $n$，使 $k\geq n$ 时有 $M_k=0$。
\end{description}
若 $M$ 满足（TN），则对 $S_+$ 中每个齐次元素 $f$，都有 $M_{(f)}=0$，所以 $\wt M=0$。

设 $M,N$ 为两个分次 $S$-模。若零次同态 $u:M\to N$ 满足：存在整数 $n$，使
$u_k:M_k\to N_k$ 在 $k\geq n$ 时单射（分别地，满射、双射），则称 $u$ 是\textbf{（TN）-单射}（分别地，\textbf{（TN）-满射}、\textbf{（TN）-双射}）。所以，$u$ 为（TN）-单射（分别地，（TN）-满射），等价于 $\Ker u$（分别地，$\Coker u$）满足（TN）。由（2.5.4），若 $u$ 为（TN）-单射（分别地，满射、双射），则 $\wt u$ 单射（分别地，满射、双射）。当 $u$ 为（TN）-双射时，也称 $u$ 为\textbf{（TN）-同构}。

\egastatement{2.7.3}{命题}
设 $S$ 为分次环，理想 $S_+$ 为有限型，$M$ 为分次 $S$-模。
\begin{enumerate}[label=(\roman*)]
\item 若 $M$ 满足（TF），则 $\OO_X$-模 $\wt M$ 为有限型。
\item 假设 $M$ 满足（TF）。则 $\wt M=0$，当且仅当 $M$ 满足（TN）。
\end{enumerate}
\transnote{本命题按原文所列假设翻译。第（ii）项的证明中使用了由 $M_{(f_j)}=0$ 推出 $f_j$ 的某个幂消去模生成元的步骤；（2.5.6）已对这类逆向推理特别说明一次生成假设。本译稿在此不暗中增加原文未列出的条件。}
\begin{proof}
已经说明（TN）推出 $\wt M=0$。若 $M$ 满足（TF），令
$M'=\bigoplus_{k\geq n}M_k$ 为假设中的有限型分次子模。则 $M/M'$ 满足（TN），所以
$(M/M')^\sim=0$，由函子 $M\mapsto\wt M$ 的正合性（2.5.4）得到
$\wt M=\wt{M'}$。因此，为证明 $\wt M$ 为有限型，可归结到 $M$ 本身为有限型。
问题是局部的，只需证明 $M_{(f)}$ 是有限型 $S_{(f)}$-模（I，1.3.9）。而 $M^{(d)}$ 是有限型 $S^{(d)}$-模（2.1.6，iii），断言由（2.2.5）得出。

现设 $M$ 满足（TF）且 $\wt M=0$。仍用上述记号，则 $\wt{M'}=0$；$M'$ 满足（TN）与 $M$ 满足（TN）等价。因此，为证明 $\wt M=0$ 推出（TN），仍可设 $M$ 由有限多个齐次元素 $x_i$（$1\leq i\leq p$）生成。再令
$(f_j)_{1\leq j\leq q}$ 是理想 $S_+$ 的齐次生成系。

由假设，每个 $j$ 都满足 $M_{(f_j)}=0$，故存在整数 $n$，使任意 $i,j$ 都有 $f_j^nx_i=0$。
令 $n_j=\deg f_j$；在满足 $\sum_jr_j\leq nq$ 的有限多个非负整数数组 $(r_j)$ 中，取
$\sum_jr_jn_j$ 的最大值为 $m$。显然，当 $k>m$ 时，每个 $i$ 都有 $S_kx_i=0$。若 $h$ 是各 $x_i$ 的次数的最大值，便得 $M_k=0$（$k>h+m$），证明完成。
\end{proof}

\egastatement{2.7.4}{推论}
设 $S$ 为分次环，理想 $S_+$ 为有限型。则 $X=\Proj(S)=\varnothing$，当且仅当存在整数 $n$，使 $k\geq n$ 时有 $S_k=0$。
\begin{proof}
条件 $X=\varnothing$ 等价于 $\OO_X=\wt S=0$，而 $S$ 是单生成 $S$-模。
\end{proof}

\egastatement{2.7.5}{定理}
假设理想 $S_+$ 由有限多个一次齐次元素生成，令 $X=\Proj(S)$。则对每个拟凝聚 $\OO_X$-模 $\F$，典范同态
$\beta:(\Gstar(\F))^\sim\to\F$（2.6.4）是同构。
\begin{proof}
若 $S_+$ 由有限多个 $f_i\in S_1$ 生成，则 $X$ 是拟紧子空间
$\Spec(S_{(f_i)})$（2.3.6）的并，故 $X$ 拟紧。此外，$X$ 是概形（2.4.2）。由（I，9.3.2）、（2.5.14.2）及（2.6.3），对每个 $f\in S_d$（$d>0$），有典范同构
\[
(\Gstar(\F))_{(\alpha_d(f))}\simeq\Gamma(D_+(f),\F).
\]
按定义，将 $\Gstar(\F)$ 看作 $\Gstar(\OO_X)$-模时的
$(\Gstar(\F))_{(\alpha_d(f))}$，就是将它看作 $S$-模时的 $(\Gstar(\F))_{(f)}$。
回到（I，9.3.1）中上述典范同构的定义，可见它与 $\beta_f$ 相同，定理得证。
\end{proof}

\egastatement{2.7.6}{注}
若假设分次环 $S$ 诺特，则只要 $S_+$ 由一次齐次元素的集合 $S_1$ 生成，（2.7.5）的条件就自动满足。

\egastatement{2.7.7}{推论}
在（2.7.5）的假设下，每个拟凝聚 $\OO_X$-模 $\F$ 都同构于形如 $\wt M$ 的 $\OO_X$-模，其中 $M$ 是分次 $S$-模。

\egastatement{2.7.8}{推论}
在（2.7.5）的假设下，每个有限型拟凝聚 $\OO_X$-模 $\F$ 都同构于形如 $\wt N$ 的 $\OO_X$-模，其中 $N$ 是有限型分次 $S$-模。
\begin{proof}
可设 $\F=\wt M$，$M$ 为分次 $S$-模（2.7.7）。令 $(z_\lambda)_{\lambda\in L}$ 为 $M$ 的齐次生成系。对每个有限子集 $H\subset L$，令 $M_H$ 为由 $\lambda\in H$ 对应的 $z_\lambda$ 生成的分次子模。显然，$M$ 为各 $M_H$ 的归纳极限，故 $\F$ 为各 $\wt{M_H}$ 的归纳极限（2.5.4）。但 $\F$ 有限型，底空间 $X$ 拟紧，由（0，5.2.3），对某个有限子集 $H$，有 $\F=\wt{M_H}$。
\end{proof}

\egastatement{2.7.9}{推论}
在（2.7.5）的假设下，设 $\F$ 为有限型拟凝聚 $\OO_X$-模。则存在整数 $n_0$，使对每个 $n\geq n_0$，$\F(n)$ 同构于形如 $\OO_X^k$ 的模的商，其中 $k>0$ 可依赖于 $n$；因此，$\F(n)$ 由它在 $X$ 上的有限多个截面生成（0，5.1.1）。
\begin{proof}
由（2.7.8），可设 $\F=\wt M$，其中 $M$ 是有限多个形如 $S(m_i)$ 的模的直和的商。根据（2.5.4），可只考虑 $M=S(m)$，于是
$\F(n)=(S(m+n))^\sim=\OO_X(m+n)$（2.5.15）。所以只需证明以下引理。
\end{proof}

\egastatement{2.7.9.1}{引理}
在（2.7.5）的假设下，对每个 $n\geq0$，存在依赖于 $n$ 的整数 $k$，以及满同态
$\OO_X^k\to\OO_X(n)$。
\begin{proof}
只需证明（2.7.2）：对适当的 $k$，存在从分次 $S$-模 $S^k$ 到 $S(n)$ 的零次（TN）-满同态 $u$。
有 $(S(n))_0=S_n$，并由假设，对每个 $h>0$ 有 $S_h=S_1^h$，所以
$SS_n=S_n+S_{n+1}+\cdots$。由于 $S_n$ 是有限型 $S_0$-模（2.1.5及2.1.6，i），取其生成系 $(z_i)_{1\leq i\leq k}$。令 $u$ 将 $S^k$ 的第 $i$ 个标准基元素映到 $z_i$。\transnote{此处按照同态 $u:S^k\to S(n)$ 的方向写出基元素的像；印刷原文的对应方向写反。}
这时 $\Coker u$ 等同于
$(S(n))_{-n}+\cdots+(S(n))_{-1}$，所以 $u$ 满足要求。
\end{proof}

\egastatement{2.7.10}{推论}
在（2.7.5）的假设下，设 $\F$ 为有限型拟凝聚 $\OO_X$-模。则存在整数 $n_0$，使对每个 $n\geq n_0$，$\F$ 同构于形如 $(\OO_X(-n))^k$ 的模的商，其中 $k$ 可依赖于 $n$。

\egastatement{2.7.11}{命题}
假设（2.7.5）的条件成立，设 $M$ 为分次 $S$-模。则：
\begin{enumerate}[label=(\roman*)]
\item 典范同态 $\wt\alpha:\wt M\to(\Gstar(\wt M))^\sim$ 是同构。
\item 设 $\G$ 为 $\wt M$ 的拟凝聚 $\OO_X$-子模，$N$ 为 $\Gstar(\G)$ 在 $\alpha$ 下的逆像所成的分次 $S$-子模。则 $\wt N=\G$；这里根据（2.5.4），将 $\wt N$ 等同于 $\wt M$ 的子模。
\end{enumerate}
\begin{proof}
由（2.7.5），$\beta:(\Gstar(\wt M))^\sim\to\wt M$ 是同构。由（2.6.5.1），$\wt\alpha$ 是它的逆同构，证明了（i）。

令 $P=\alpha(M)$ 为 $\Gstar(\wt M)$ 的分次子模。由于 $M\mapsto\wt M$ 正合（2.5.4），$\wt M$ 在 $\wt\alpha$ 下的像等于 $\wt P$；故由（i），
$\wt P=(\Gstar(\wt M))^\sim$。令 $Q=\Gstar(\G)\cap P$。由上述结果及（2.5.4），有
$\wt Q=(\Gstar(\G))^\sim$，故由（2.7.5），$\beta$ 在 $\wt Q$ 上的限制是到 $\G$ 的同构。

按 $N$ 的定义及（2.5.4），同构 $\wt\alpha$ 在 $\wt N$ 上的限制是到 $\wt Q$ 的同构。于是结合（2.6.5.1）得到所需结论。
\end{proof}

\subsection{函子性}
\src{41--48}
\egapar{2.8.1}
设 $S,S'$ 是两个非负分次环，$\varphi:S'\to S$ 是分次环同态。以 $G(\varphi)$ 表示
$X=\Proj(S)$ 中 $V_+(\varphi(S'_+))$ 的补集；等价地，它是所有
$D_+(\varphi(f'))$ 的并，其中 $f'$ 遍历 $S'_+$ 中的齐次元素。因此，从 $\Spec(S)$ 到 $\Spec(S')$ 的连续映射 ${}^{a}\varphi$（I，1.2.1）在 $G(\varphi)$ 上的限制，是从 $G(\varphi)$ 到 $\Proj(S')$ 的连续映射；仍将它简记为 ${}^{a}\varphi$。若 $f'$ 是 $S'_+$ 中齐次元素，则结合（I，1.2.2.2）及 ${}^{a}\varphi(G(\varphi))\subset\Proj(S')$，有
\begin{equation}\tag{2.8.1.1}
({}^{a}\varphi)^{-1}(D_+(f'))=D_+(\varphi(f')).
\end{equation}

另一方面，令 $f=\varphi(f')$。同态 $\varphi$ 典范地定义分次环同态
$S'_{f'}\to S_f$，限制到零次元素，得到环同态
$\varphi_{(f')}:S'_{(f')}\to S_{(f)}$。它对应于仿射概形态射（I，1.6.1）
\[
({}^{a}\varphi_{(f')},\wt\varphi_{(f')}):
\Spec(S_{(f)})\longrightarrow\Spec(S'_{(f')}).
\]
将 $\Spec(S_{(f)})$ 典范地等同于 $\Proj(S)$ 在 $D_+(f)$ 上诱导的概形（2.3.6），就定义了态射
$\Phi_f:D_+(f)\to D_+(f')$；其底映射等同于 ${}^{a}\varphi$ 在 $D_+(f)$ 上的限制。

若 $g'$ 是 $S'_+$ 的另一个齐次元素，$g=\varphi(g')$，则图
\[
\begin{tikzcd}
D_+(f) \arrow[r,"\Phi_f"] & D_+(f')\\
D_+(fg) \arrow[r,"\Phi_{fg}"'] \arrow[u,hook] & D_+(f'g') \arrow[u,hook]
\end{tikzcd}
\]
立即交换，因为相应的环同态图
\[
\begin{tikzcd}[column sep=large]
S'_{(f')} \arrow[r,"\varphi_{(f')}"] \arrow[d] & S_{(f)} \arrow[d]\\
S'_{(f'g')} \arrow[r,"\varphi_{(f'g')}"'] & S_{(fg)}
\end{tikzcd}
\]
交换。结合 $G(\varphi)$ 的定义及（2.3.3.2），得到以下命题。

\egastatement{2.8.2}{命题}
给定分次环同态 $\varphi:S'\to S$，存在唯一的态射
$({}^{a}\varphi,\wt\varphi):G(\varphi)\to\Proj(S')$，使对 $S'_+$ 中每个齐次元素 $f'$，它在 $D_+(\varphi(f'))$ 上的限制，都等于与 $\varphi$ 对应的同态
$S'_{(f')}\to S_{(\varphi(f'))}$ 所相伴的态射。称此态射\textbf{与 $\varphi$ 相伴}，记作 $\Proj(\varphi)$。

注意，采用上述记号，若 $f'\in S'_d$，则图
\begin{equation}\tag{2.8.2.1}
\begin{tikzcd}[column sep=large]
S'_{(f')} \arrow[r,"\varphi_{(f')}"] \arrow[d,"\sim"'] & S_{(f)} \arrow[d,"\sim"]\\
S'^{(d)}/(f'-1)S'^{(d)} \arrow[r] & S^{(d)}/(f-1)S^{(d)}
\end{tikzcd}
\end{equation}
交换，其中竖直箭头是（2.2.5）中的同构。

\egastatement{2.8.3}{推论}
\begin{enumerate}[label=(\roman*)]
\item 态射 $\Proj(\varphi)$ 仿射。
\item 若 $\Ker\varphi$ 幂零，特别地，若 $\varphi$ 单射，则 $\Proj(\varphi)$ 是支配态射。
\end{enumerate}
\begin{proof}
（i）立即由（2.8.2）及（2.8.1.1）得出。若 $\Ker\varphi$ 幂零，则对每个齐次元素 $f'\in S'_+$，令 $f=\varphi(f')$，立即可验证
$\Ker(\varphi_{f'})$ 幂零，因而 $\Ker(\varphi_{(f')})$ 也幂零。结论由（2.8.2）及（I，1.2.7）得出。
\end{proof}

应注意：一般存在从 $\Proj(S)$ 到 $\Proj(S')$ 的非仿射态射，因而它们并不来自分次环同态 $S'\to S$。一个例子是 $A$ 为域时的结构态射
$\Proj(S)\to\Spec(A)$，其中将 $\Spec(A)$ 等同于 $\Proj(A[T])$（见3.1.7）；这由（I，2.3.2）得出。

\egapar{2.8.4}
设 $S''$ 是第三个非负分次环，$\varphi':S''\to S'$ 是分次环同态，令
$\varphi''=\varphi\circ\varphi'$。由（2.8.1.1）及
${}^{a}\varphi''=({}^{a}\varphi')\circ({}^{a}\varphi)$，立即可验证
\[
G(\varphi'')\subset G(\varphi).
\]
若 $\Phi,\Phi',\Phi''$ 分别是与 $\varphi,\varphi',\varphi''$ 相伴的态射，则
\[
\Phi''=\Phi'\circ\bigl(\Phi|G(\varphi'')\bigr).
\]

\egapar{2.8.5}
假设 $S$（分别地，$S'$）是分次 $A$-代数（分别地，分次 $A'$-代数），并设环同态
$\psi:A'\to A$ 使图
\begin{equation}\tag{2.8.5.1}
\begin{tikzcd}
A' \arrow[r,"\psi"] \arrow[d] & A \arrow[d]\\
S' \arrow[r,"\varphi"'] & S
\end{tikzcd}
\end{equation}
交换。于是可分别把 $G(\varphi)$、$\Proj(S')$ 看作 $\Spec(A)$、$\Spec(A')$ 上的概形。若 $\Phi,\Psi$ 分别是与 $\varphi,\psi$ 相伴的态射，则图
\[
\begin{tikzcd}[column sep=large]
G(\varphi) \arrow[r,"\Phi"] \arrow[d] & \Proj(S') \arrow[d]\\
\Spec(A) \arrow[r,"\Psi"'] & \Spec(A')
\end{tikzcd}
\]
交换。事实上，只需对 $\Phi$ 在 $D_+(f)$ 上的限制证明这一点，其中 $f=\varphi(f')$、$f'$ 是 $S'_+$ 中齐次元素；这时结论来自交换图
\[
\begin{tikzcd}
A' \arrow[r,"\psi"] \arrow[d] & A \arrow[d]\\
S'_{(f')} \arrow[r,"\varphi_{(f')}"'] & S_{(f)}.
\end{tikzcd}
\]

\egapar{2.8.6}
现在设 $M$ 是分次 $S$-模，并考虑显然分次的 $S'$-模 $M_{[\varphi]}$。设
$f'$ 是 $S'_+$ 中齐次元素，$f=\varphi(f')$。已知存在典范同构（0，1.5.2）
\[
(M_{[\varphi]})_{f'}\simeq(M_f)_{[\varphi_{f'}]}.
\]
显然这个同构保持次数，所以还有典范同构
\[
(M_{[\varphi]})_{(f')}\simeq(M_{(f)})_{[\varphi_{(f')}]}.
\]
它典范地对应于层同构（2.5.2及I，1.6.3）
\[
(M_{[\varphi]})^\sim|D_+(f')\simeq(\Phi_f)_*\bigl(\wt M|D_+(f)\bigr).
\]
若 $g'$ 是 $S'_+$ 的另一个齐次元素，$g=\varphi(g')$，则图
\[
\begin{tikzcd}[column sep=large]
(M_{[\varphi]})_{(f')} \arrow[r,"\sim"] \arrow[d] & (M_{(f)})_{[\varphi_{(f')}]} \arrow[d]\\
(M_{[\varphi]})_{(f'g')} \arrow[r,"\sim"'] & (M_{(fg)})_{[\varphi_{(f'g')}]} 
\end{tikzcd}
\]
交换。立即得出，$D_+(f'g')$ 上对应的层同构，就是 $D_+(f')$ 上上述层同构的限制。由于 $\Phi_f$ 为 $\Phi$ 在 $D_+(f)$ 上的限制，结合（2.8.1.1），并令 $X'=\Proj(S')$，有如下命题。

\egastatement{2.8.7}{命题}
从 $\OO_{X'}$-模 $(M_{[\varphi]})^\sim$ 到 $\OO_{X'}$-模
$\Phi_*(\wt M|G(\varphi))$ 存在典范函子性同构。

由此立即得到：从一个分次 $S'$-模 $M'$ 到分次 $S$-模 $M$ 的 $\varphi$-态射的集合，到
$\Phi$-态射 $\wt{M'}\to\wt M|G(\varphi)$ 的集合，存在一个典范函子性映射。
采用（2.8.4）的记号，若 $M''$ 是分次 $S''$-模，则一个 $\varphi$-态射 $M'\to M$ 与一个 $\varphi'$-态射 $M''\to M'$ 的复合，典范地对应于
\[
\wt{M''}\longrightarrow\wt{M'}|G(\varphi')
\longrightarrow\wt M|G(\varphi'')
\]
所给出的复合态射。

\egastatement{2.8.8}{命题}
在（2.8.1）的假设下，设 $M'$ 为分次 $S'$-模。从 $(\OO_X|G(\varphi))$-模
$\Phi^*(\wt{M'})$ 到 $(\OO_X|G(\varphi))$-模
$(M'\otimes_{S'}S)^\sim|G(\varphi)$ 存在典范函子性同态 $\nu$。若理想 $S'_+$ 由 $S'_1$ 生成，则 $\nu$ 是同构。
\begin{proof}
对 $f'\in S'_d$（$d>0$），令 $f=\varphi(f')$，定义典范函子性 $S_{(f)}$-模同态
\begin{equation}\tag{2.8.8.1}
\nu_f:M'_{(f')}\otimes_{S'_{(f')}}S_{(f)}
\longrightarrow(M'\otimes_{S'}S)_{(f)}.
\end{equation}
其定义是：将同态
$M'_{(f')}\otimes_{S'_{(f')}}S_{(f)}\to M'_{f'}\otimes_{S'_{f'}}S_f$
与典范同构
$M'_{f'}\otimes_{S'_{f'}}S_f\simeq(M'\otimes_{S'}S)_f$（0，1.5.4）复合，并注意后者保持次数。对另一个 $g'\in S'_+$，令 $g=\varphi(g')$，立即可验证各 $\nu_f$ 与从 $D_+(f)$ 到 $D_+(fg)$ 的限制运算相容；结合（I，1.6.5），这定义了所需同态
\[
\nu:\Phi^*(\wt{M'})\longrightarrow(M'\otimes_{S'}S)^\sim|G(\varphi).
\]

为证明第二断言，只需对每个 $f'\in S'_1$ 证明 $\nu_f$ 为同构，因为这时 $G(\varphi)$ 是各 $D_+(\varphi(f'))$ 的并。先定义 $\mathbf Z$-双线性映射
\[
M'_m\times S_n\longrightarrow M'_{(f')}\otimes_{S'_{(f')}}S_{(f)},
\quad(x',s)\longmapsto(x'/f'^m)\otimes(s/f^n),
\]
其中对 $m<0$，约定 $x'/f'^m=f'^{-m}x'/1$。如（2.5.13）的证明中一样，可见它给出模的双同态
\[
\gamma_f:M'\otimes_{S'}S\longrightarrow
M'_{(f')}\otimes_{S'_{(f')}}S_{(f)}.
\]
若 $r>0$ 且 $f'^r\sum_i(x'_i\otimes s_i)=0$，也就是
$\sum_i(f'^rx'_i\otimes s_i)=0$，则由（0，1.5.4），
\[
\sum_i\left(\frac{f'^rx'_i}{f'^{m_i+r}}\otimes\frac{s_i}{f^{n_i}}\right)=0,
\]
即 $\gamma_f(\sum_i x'_i\otimes s_i)=0$。这说明 $\gamma_f$ 经由局部化
$(M'\otimes_{S'}S)_f$ 分解。将所得同态限制到零次部分，得到
$\nu'_f:(M'\otimes_{S'}S)_{(f)}\to M'_{(f')}\otimes_{S'_{(f')}}S_{(f)}$。
最后验证 $\nu_f$ 与 $\nu'_f$ 互逆，即完成证明。
\end{proof}

特别地，由（2.1.2.1），对每个 $n\in\mathbf Z$，有典范同态
\begin{equation}\tag{2.8.8.2}
\Phi^*(\OO_{X'}(n))\longrightarrow\OO_X(n)|G(\varphi).
\end{equation}

\egapar{2.8.9}
设 $A,A'$ 是两个环，$\psi:A'\to A$ 是环同态，定义态射
$\Psi:\Spec(A)\to\Spec(A')$。设 $S'$ 是非负分次 $A'$-代数，令
$S=S'\otimes_{A'}A$；它显然由各 $S'_n\otimes_{A'}A$ 分次，成为分次 $A$-代数。映射
$\varphi:s'\mapsto s'\otimes1$ 是使图（2.8.5.1）交换的分次环同态。
由于这时 $S_+$ 由 $\varphi(S'_+)$ 作为 $A$-模生成，有
$G(\varphi)=\Proj(S)=X$。令 $X'=\Proj(S')$，得到交换图
\begin{equation}\tag{2.8.9.1}
\begin{tikzcd}
X \arrow[r,"\Phi"] \arrow[d,"p"'] & X' \arrow[d,"p'"]\\
Y \arrow[r,"\Psi"'] & Y'
\end{tikzcd}
\end{equation}
其中 $Y=\Spec(A)$、$Y'=\Spec(A')$。
现设 $M'$ 是分次 $S'$-模，令
$M=M'\otimes_{A'}A=M'\otimes_{S'}S$。在这些条件下，有以下结果。

\egastatement{2.8.10}{命题}
图（2.8.9.1）将概形 $X$ 等同于积 $X'\times_{Y'}Y$。此外，典范同态
$\nu:\Phi^*(\wt{M'})\to\wt M$（2.8.8）是同构。
\begin{proof}
为证明第一断言，只需对 $S'_+$ 中每个齐次元素 $f'$，令 $f=\varphi(f')$，证明
$\Phi$ 和 $p$ 在 $D_+(f)$ 上的限制，将此概形等同于积
$D_+(f')\times_{Y'}Y$（I，3.2.6.2）。换言之，只需证明（I，3.2.2）
$S_{(f)}$ 典范地等同于 $S'_{(f')}\otimes_{A'}A$。而这立即来自保持次数的典范同构
$S_f\simeq S'_{f'}\otimes_{A'}A$（0，1.5.4）。

由此，$M'_{(f')}\otimes_{S'_{(f')}}S_{(f)}$ 等同于
$M'_{(f')}\otimes_{A'}A$；又因保持次数的典范同构
$M_f\simeq M'_{f'}\otimes_{A'}A$，后者等同于 $M_{(f)}$。这证明第二断言。
\end{proof}

\egastatement{2.8.11}{推论}
对每个 $n\in\mathbf Z$，$\wt M(n)$ 等同于
$\Phi^*(\wt{M'}(n))=\wt{M'}(n)\otimes_{Y'}\OO_Y$。特别地，
\[
\OO_X(n)=\Phi^*(\OO_{X'}(n))=\OO_{X'}(n)\otimes_{Y'}\OO_Y.
\]
\begin{proof}
这由（2.8.10）及（2.5.15）得出。
\end{proof}

\egapar{2.8.12}
在（2.8.9）的假设下，对 $f'\in S'_d$（$d>0$）及 $f=\varphi(f')$，图
\begin{equation}\tag{2.8.12.1}
\begin{tikzcd}[column sep=large]
M'_{(f')} \arrow[r,"\sim"] \arrow[d] & M'^{(d)}/(f'-1)M'^{(d)} \arrow[d]\\
M_{(f)} \arrow[r,"\sim"'] & M^{(d)}/(f-1)M^{(d)}
\end{tikzcd}
\end{equation}
交换（参见2.2.5）。

\egapar{2.8.13}
保持（2.8.9）的记号与假设。设 $\F'$ 是 $\OO_{X'}$-模，令 $\F=\Phi^*\F'$。由（2.8.11）及（0，4.3.3），对每个 $n\in\mathbf Z$，有
$\F(n)=\Phi^*(\F'(n))$。因此有典范同态（0，3.7.1）
\[
\Gamma(\rho):\Gamma(X',\F'(n))\longrightarrow\Gamma(X,\F(n)),
\]
给出分次模的典范双同态 $\Gstar(\F')\to\Gstar(\F)$。

假设 $S'_+$ 由 $S'_1$ 生成，且 $\F'=\wt{M'}$，于是 $\F=\wt M$，其中 $M=M'\otimes_{A'}A$。若 $f'$ 为 $S'_+$ 中齐次元素，$f=\varphi(f')$，已知
$M_{(f)}=M'_{(f')}\otimes_{A'}A$。因此图
\[
\begin{tikzcd}[column sep=large]
M'_0 \arrow[r] \arrow[d] & M'_{(f')}=\Gamma(D_+(f'),\wt{M'}) \arrow[d]\\
M_0 \arrow[r] & M_{(f)}=\Gamma(D_+(f),\wt M)
\end{tikzcd}
\]
交换。由此观察及同态 $\alpha$ 的定义（2.6.2），立即得出图
\begin{equation}\tag{2.8.13.1}
\begin{tikzcd}[column sep=large]
M' \arrow[r,"\alpha_{M'}"] \arrow[d] & \Gstar(\wt{M'}) \arrow[d]\\
M \arrow[r,"\alpha_M"'] & \Gstar(\wt M)
\end{tikzcd}
\end{equation}
交换。同样，图
\begin{equation}\tag{2.8.13.2}
\begin{tikzcd}[column sep=large]
(\Gstar(\F'))^\sim \arrow[r,"\beta_{\F'}"] \arrow[d] & \F' \arrow[d]\\
(\Gstar(\F))^\sim \arrow[r,"\beta_\F"'] & \F
\end{tikzcd}
\end{equation}
交换，其中右侧竖直箭头是典范 $\Phi$-态射 $\F'\to\Phi^*\F'=\F$。

\egapar{2.8.14}
继续保持（2.8.9）的假设与记号。设 $N'$ 是第二个分次 $S'$-模，$N=N'\otimes_{A'}A$。典范双同态 $M'\to M$、$N'\to N$ 显然给出双同态
$M'\otimes_{S'}N'\to M\otimes_SN$，相对于典范环同态 $S'\to S$；因而也给出零次 $S$-同态
\[
(M'\otimes_{S'}N')\otimes_{A'}A\longrightarrow M\otimes_SN.
\]
结合（2.8.10），它对应于 $\OO_X$-模同态
\begin{equation}\tag{2.8.14.1}
\Phi^*\bigl((M'\otimes_{S'}N')^\sim\bigr)
\longrightarrow(M\otimes_SN)^\sim.
\end{equation}
此外，立即可验证图
\begin{equation}\tag{2.8.14.2}
\begin{tikzcd}[column sep=large]
\Phi^*(\wt{M'}\otimes_{\OO_{X'}}\wt{N'}) \arrow[r,"\sim"] \arrow[d,"\Phi^*(\lambda)"'] & \wt M\otimes_{\OO_X}\wt N \arrow[d,"\lambda"]\\
\Phi^*((M'\otimes_{S'}N')^\sim) \arrow[r] & (M\otimes_SN)^\sim
\end{tikzcd}
\end{equation}
交换，其中第一行是典范同构（0，4.3.3），右上角也等于
$\Phi^*(\wt{M'})\otimes_{\OO_X}\Phi^*(\wt{N'})$。
若 $S'_+$ 由 $S'_1$ 生成，显然 $S_+$ 由 $S_1$ 生成；这时图中的两个竖直箭头均为同构（2.5.13），所以（2.8.14.1）也为同构。

同样，将一个 $k$ 次同态 $u'$ 对应于同为 $k$ 次的同态 $u'\otimes1$，得到典范双同态
$\Hom_{S'}(M',N')\to\Hom_S(M,N)$。由此又得到零次 $S$-同态
\[
\Hom_{S'}(M',N')\otimes_{A'}A\longrightarrow\Hom_S(M,N),
\]
进而得到 $\OO_X$-模同态
\begin{equation}\tag{2.8.14.3}
\Phi^*\bigl((\Hom_{S'}(M',N'))^\sim\bigr)
\longrightarrow(\Hom_S(M,N))^\sim.
\end{equation}
此外，图
\[
\begin{tikzcd}[column sep=large]
\Phi^*((\Hom_{S'}(M',N'))^\sim) \arrow[r] \arrow[d,"\Phi^*(\mu)"'] & (\Hom_S(M,N))^\sim \arrow[d,"\mu"]\\
\Phi^*(\shHom_{\OO_{X'}}(\wt{M'},\wt{N'})) \arrow[r] & \shHom_{\OO_X}(\wt M,\wt N)
\end{tikzcd}
\]
交换，其中第二行是典范同态（0，4.4.6）。

\egapar{2.8.15}
采用（2.8.1）的记号与假设。由（2.4.7），先将 $S_0,S'_0$ 换成 $\mathbf Z$、将 $\varphi_0$ 换成恒等映射，再将 $S,S'$ 分别换成 $S^{(d)},S'^{(d)}$（$d>0$），将 $\varphi$ 换成其在 $S'^{(d)}$ 上的限制 $\varphi^{(d)}$，在同构意义下，并不改变态射 $\Phi$。

\subsection{概形 \texorpdfstring{$\Proj(S)$}{Proj(S)} 的闭子概形}
\src{48--49}
\egapar{2.9.1}
若 $\varphi:S\to S'$ 是分次环同态，并存在整数 $n$，使 $k\geq n$ 时
$\varphi_k:S_k\to S'_k$ 满射（分别地，单射、双射），则称 $\varphi$ 为\textbf{（TN）-满射}（分别地，\textbf{（TN）-单射}、\textbf{（TN）-双射}）。由（2.8.15），相应态射 $\Phi$ 的研究，可以归结到 $\varphi$ 满射（分别地，单射、双射）的情形。（TN）-双射也称为\textbf{（TN）-同构}。

\egastatement{2.9.2}{命题}
设 $S$ 为非负分次环，$X=\Proj(S)$。
\begin{enumerate}[label=(\roman*)]
\item 若 $\varphi:S\to S'$ 是分次环的（TN）-满同态，则相应态射 $\Phi$（2.8.1）在整个 $\Proj(S')$ 上有定义，而且是从 $\Proj(S')$ 到 $X$ 的闭浸入。若 $\frj=\Ker\varphi$，则与 $\Phi$ 相伴的闭子概形由 $\OO_X$ 的拟凝聚理想层 $\wt\frj$ 定义。
\item 再假设 $S_+$ 由有限多个一次齐次元素生成。设 $X'$ 是 $X$ 的闭子概形，由 $\OO_X$ 的拟凝聚理想层 $\J$ 定义。令 $\frj$ 为 $\Gstar(\J)$ 在典范同态
$\alpha:S\to\Gstar(\OO_X)$（2.6.2）下的逆像所成的分次理想，并令 $S'=S/\frj$。则 $X'$ 是与典范分次环同态 $S\to S'$ 所对应的闭浸入 $\Proj(S')\to X$ 相伴的子概形。
\end{enumerate}
\begin{proof}
（i）可设 $\varphi$ 满射（2.9.1）。由假设，$\varphi(S_+)$ 生成 $S'_+$，所以
$G(\varphi)=\Proj(S')$。第二断言在 $X$ 上局部验证。令 $f$ 为 $S_+$ 的齐次元素，$f'=\varphi(f)$。
由于 $\varphi$ 为分次环满同态，立即可验证
$\varphi_{(f)}:S_{(f)}\to S'_{(f')}$ 满射，且其核为 $\frj_{(f)}$。由（I，4.2.3），（i）得证。

（ii）由（i），只需验证：由典范嵌入 $j:\frj\to S$ 诱导的同态
$\wt j:\wt\frj\to\OO_X$，是从 $\wt\frj$ 到 $\J$ 的同构；这由（2.7.11）得出。
\end{proof}
注意，$\frj$ 是所有满足 $\wt j(\wt{\frj'})=\J$ 的分次理想 $\frj'\subset S$ 中最大的一个。因为回到定义（2.6.2），立即可验证上述关系推出
$\alpha(\frj')\subset\Gstar(\J)$。

\egastatement{2.9.3}{推论}
假设（2.9.2，i）的条件成立，且 $S_+$ 由 $S_1$ 生成。则对每个 $n\in\mathbf Z$，
$\Phi^*((S(n))^\sim)$ 典范地同构于 $(S'(n))^\sim$。所以对每个 $\OO_X$-模 $\F$，$\Phi^*(\F(n))$ 典范地同构于 $(\Phi^*\F)(n)$。
\begin{proof}
这是（2.8.8）的特殊情形，并使用（2.5.10.2）。
\end{proof}

\egastatement{2.9.4}{推论}
假设（2.9.2，ii）的条件成立。则 $X$ 的闭子预概形 $X'$ 为整预概形，当且仅当分次理想 $\frj$ 是 $S$ 的素理想。
\begin{proof}
因为 $X'\simeq\Proj(S/\frj)$，充分性由（2.4.4）得出。为证明必要性，考虑正合列
$0\to\J\to\OO_X\to\OO_X/\J$，它给出正合列
\[
0\longrightarrow\Gstar(\J)\longrightarrow\Gstar(\OO_X)
\longrightarrow\Gstar(\OO_X/\J).
\]
只需证明：若 $f\in S_m,g\in S_n$，且 $\alpha_{m+n}(fg)$ 在 $\Gstar(\OO_X/\J)$ 中的像为零，则 $\alpha_m(f),\alpha_n(g)$ 的像之一为零。

按定义，这两个像是整概形 $X'$ 上可逆 $(\OO_X/\J)$-模
$\LL=(\OO_X/\J)(m)$ 与 $\LL'=(\OO_X/\J)(n)$ 的截面。假设表示它们的积在 $\LL\otimes\LL'$ 中为零（2.9.3及2.5.14.1），所以其中一个截面为零（I，7.4.4）。
\end{proof}

\egastatement{2.9.5}{推论}
设 $A$ 是环，$M$ 是 $A$-模，$S$ 是由一次齐次元素的集合 $S_1$ 生成的分次 $A$-代数，
$u:M\to S_1$ 是 $A$-模满同态，$\bar u:\Sym(M)\to S$ 是从 $M$ 的对称代数到 $S$、延拓 $u$ 的 $A$-代数同态。则与 $\bar u$ 对应的态射是从 $\Proj(S)$ 到 $\Proj(\Sym(M))$ 的闭浸入。
\begin{proof}
由假设，$\bar u$ 满射；应用（2.9.2）即可。
\end{proof}


\section{分次代数层的齐次谱}
\subsection{拟凝聚分次 \texorpdfstring{$\OO_Y$}{OY}-代数的齐次谱}
\src{49--54}
\egapar{3.1.1}
设 $Y$ 是预概形，$\SSS$ 是分次 $\OO_Y$-代数，$\M$ 是分次 $\SSS$-模。若 $\SSS$ 拟凝聚，则它的每个齐次分量 $\SSS_n$ 都是拟凝聚 $\OO_Y$-模，因为它是 $\SSS$ 的一个自同态的像（I，1.3.8及1.3.9）。同样，若 $\M$ 作为 $\OO_Y$-模拟凝聚，则它的各个齐次分量 $\M_n$ 也拟凝聚；逆命题也成立。对整数 $d>0$，用 $\SSS^{(d)}$ 表示 $\OO_Y$-模 $\SSS_{nd}$（$n\in\Z$）的直和；若 $\SSS$ 拟凝聚，则它也拟凝聚（I，1.3.9）。对 $0\leq k\leq d-1$，用 $\M^{(d,k)}$ 表示各个 $\M_{nd+k}$（$n\in\Z$）的直和；$k=0$ 时也记作 $\M^{(d)}$。它是一个分次 $\SSS^{(d)}$-模；当 $\SSS$ 与 $\M$ 拟凝聚时，它也拟凝聚（I，9.6.1）。用 $\M(n)$ 表示满足 $(\M(n))_k=\M_{n+k}$（$k\in\Z$）的分次 $\SSS$-模；若 $\SSS$ 与 $\M$ 拟凝聚，则 $\M(n)$ 也为拟凝聚分次 $\SSS$-模（I，9.6.1）。

若每个 $y\in Y$ 都有一个开邻域 $U$，并存在有限多个整数 $n_j$，使有零次满同态
\[
\bigoplus_{j=1}^{q}\bigl(\SSS(n_j)|U\bigr)\longrightarrow\M|U,
\]
则称 $\M$ 为\textbf{有限型分次 $\SSS$-模}。若每个 $y$ 都有开邻域 $U$，以及有限多个整数 $m_i,n_j$，使 $\M|U$ 同构于某个零次同态
\[
\bigoplus_{i=1}^{p}\bigl(\SSS(m_i)|U\bigr)\longrightarrow
\bigoplus_{j=1}^{q}\bigl(\SSS(n_j)|U\bigr)
\]
的余核，则称 $\M$ \textbf{具有有限表示}。

设 $U$ 是 $Y$ 的仿射开集，$A=\Gamma(U,\OO_Y)$ 是其环。由假设，分次 $(\OO_Y|U)$-代数 $\SSS|U$ 同构于 $\wt S$，其中 $S=\Gamma(U,\SSS)$ 是分次 $A$-代数（I，1.4.3）。令 $X_U=\Proj(\Gamma(U,\SSS))$。设 $U'\subset U$ 是另一个仿射开集，$A'$ 是它的环，$j:U'\to U$ 为典范嵌入，对应于限制同态 $A\to A'$。有 $\SSS|U'=j^*(\SSS|U)$，因而 $S'=\Gamma(U',\SSS)$ 典范地等同于 $S\otimes_A A'$（I，1.6.4）。由（2.8.10），$X_{U'}$ 典范地等同于 $X_U\times_U U'$，从而也等同于 $f_U^{-1}(U')$，其中 $f_U:X_U\to U$ 是结构态射（I，4.4.1）。把这样得到的典范同构 $f_U^{-1}(U')\xrightarrow{\sim}X_{U'}$ 记作 $\sigma_{U',U}$，把 $\sigma_{U',U}^{-1}$ 与典范嵌入 $f_U^{-1}(U')\to X_U$ 的复合所给出的开浸入记作 $\rho_{U',U}:X_{U'}\to X_U$。若 $U''\subset U'$ 是第三个仿射开集，则立即有
\[
\rho_{U',U}\circ\rho_{U'',U'}=\rho_{U'',U}.
\]

\egastatement{3.1.2}{命题}
设 $Y$ 是预概形。对每个非负分次拟凝聚 $\OO_Y$-代数 $\SSS$，存在一个 $Y$ 上的预概形 $X$，在 $Y$-同构意义下唯一，具有如下性质：若 $f:X\to Y$ 为结构态射，则对每个仿射开集 $U\subset Y$，存在从诱导预概形 $f^{-1}(U)$ 到 $X_U=\Proj(\Gamma(U,\SSS))$ 的同构 $\eta_U$；且对每个包含于 $U$ 的仿射开集 $V$，图
\begin{equation}\tag{3.1.2.1}
\begin{tikzcd}
f^{-1}(V) \arrow[r,"\eta_V"] \arrow[d,hook] & X_V \arrow[d,"\rho_{V,U}"]\\
f^{-1}(U) \arrow[r,"\eta_U"'] & X_U
\end{tikzcd}
\end{equation}
交换。
\begin{proof}
对 $Y$ 的两个仿射开集 $U,V$，令 $X_{U,V}$ 为 $X_U$ 在 $f_U^{-1}(U\cap V)$ 上诱导的预概形。下面定义 $Y$-同构 $\theta_{V,U}:X_{U,V}\to X_{V,U}$。对仿射开集 $W\subset U\cap V$，复合同构
\[
f_U^{-1}(W)\xrightarrow{\sigma_{W,U}}X_W
\xrightarrow{\sigma_{W,V}^{-1}}f_V^{-1}(W),
\]
得到同构 $\tau_W$。立即可验证：若 $W'\subset W$ 是仿射开集，则 $\tau_{W'}$ 是 $\tau_W$ 在 $f_U^{-1}(W')$ 上的限制。因此这些 $\tau_W$ 确实是某个 $Y$-同构 $\theta_{V,U}$ 的限制。又若 $U,V,W$ 是 $Y$ 的三个仿射开集，把 $\theta_{U,V},\theta_{V,W},\theta_{U,W}$ 分别限制到 $U\cap V\cap W$ 在 $X_V,X_W$ 中的相应逆像上，并记作 $\theta'_{U,V},\theta'_{V,W},\theta'_{U,W}$，则由定义有
$\theta'_{U,V}\circ\theta'_{V,W}=\theta'_{U,W}$。
故满足所述性质的 $X$ 的存在性由（I，2.3.1）得出；结合（3.1.2.1），其在 $Y$-同构意义下的唯一性是显然的。
\end{proof}

\egapar{3.1.3}
称（3.1.2）定义的预概形 $X$ 为拟凝聚分次 $\OO_Y$-代数 $\SSS$ 的\textbf{齐次谱}，记作 $\Proj(\SSS)$。立即可知 $\Proj(\SSS)$ 在 $Y$ 上分离（2.4.2及I，5.5.5）；若 $\SSS$ 是有限型 $\OO_Y$-代数（I，9.6.2），则 $\Proj(\SSS)$ 在 $Y$ 上为有限型（2.7.1，ii及I，6.3.1）。若 $f:X\to Y$ 是结构态射，则对 $Y$ 在任一开集 $U$ 上诱导的预概形，$f^{-1}(U)$ 立即等同于齐次谱 $\Proj(\SSS|U)$。

\egastatement{3.1.4}{命题}
设 $f\in\Gamma(Y,\SSS_d)$，$d>0$。则 $X=\Proj(\SSS)$ 的底空间中存在开子集 $X_f$，具有如下性质：对每个仿射开集 $U\subset Y$，在等同于 $X_U=\Proj(\Gamma(U,\SSS))$ 的 $\varphi^{-1}(U)$ 中有
\[
X_f\cap\varphi^{-1}(U)=D_+(f|U),
\]
其中 $\varphi:X\to Y$ 是结构态射。此外，$X_f$ 上诱导的预概形在 $Y$ 上仿射，且典范地同构于
$\Spec\bigl(\SSS^{(d)}/(f-1)\SSS^{(d)}\bigr)$（1.3.1）。
\begin{proof}
有 $f|U\in\Gamma(U,\SSS_d)=(\Gamma(U,\SSS))_d$。若 $U'\subset U$ 是 $Y$ 的仿射开集，则 $f|U'$ 是 $f|U$ 在限制同态 $\Gamma(U,\SSS)\to\Gamma(U',\SSS)$ 下的像。因此，采用（3.1.1）的记号，$D_+(f|U')$ 是 $X_{U'}$ 在 $\rho_{U',U}^{-1}(D_+(f|U))$ 上诱导的预概形（2.8.1）。第一断言由此得出。又 $X_U$ 在 $D_+(f|U)$ 上诱导的预概形典范地等同于
$\Spec((\Gamma(U,\SSS))_{(f|U)})$，这些等同与限制同态相容（2.8.1）。第二断言由（2.2.5）和图（2.8.2.1）的交换性得出。
\end{proof}
仍称 $X_f$（作为底空间 $X$ 的开集）为 $f$ \textbf{不为零的点} $x\in X$ 所成的集合。

\egastatement{3.1.5}{推论}
若 $f\in\Gamma(Y,\SSS_m)$、$g\in\Gamma(Y,\SSS_n)$，则
\begin{equation}\tag{3.1.5.1}X_{fg}=X_f\cap X_g.\end{equation}
\begin{proof}
只须把两边分别与 $\varphi^{-1}(U)$ 相交，其中 $U$ 是 $Y$ 的仿射开集，再应用（2.3.3.2）。
\end{proof}

\egastatement{3.1.6}{推论}
设 $(f_\alpha)$ 是 $\SSS$ 在 $Y$ 上的一族截面，且 $f_\alpha\in\Gamma(Y,\SSS_{n_\alpha})$。若这族截面生成的 $\SSS$-理想层（0，5.1.1）包含一切充分大次数的 $\SSS_n$，则底空间 $X$ 是各 $X_{f_\alpha}$ 的并。
\begin{proof}
对每个仿射开集 $U\subset Y$，$\varphi^{-1}(U)$ 是各 $X_{f_\alpha}\cap\varphi^{-1}(U)$ 的并（2.3.14）。
\end{proof}

\egastatement{3.1.7}{推论}
设 $\A$ 是拟凝聚 $\OO_Y$-代数，令
\[
\SSS=\A[T]=\A\otimes_{\Z}\Z[T],
\]
其中 $T$ 为不定元，$\Z$ 和 $\Z[T]$ 看成 $Y$ 上的常值层。则 $X=\Proj(\SSS)$ 典范地等同于 $\Spec(\A)$。特别地，$\Proj(\OO_Y[T])$ 等同于 $Y$。
\begin{proof}
对在每一点等于 $T$ 的唯一截面 $f\in\Gamma(Y,\SSS)$ 应用（3.1.6），得 $X_f=X$。此外，此时 $d=1$，且
$\SSS^{(1)}/(f-1)\SSS^{(1)}=\SSS/(f-1)\SSS$ 典范地同构于 $\A$；故结论成立（1.2.2）。
\end{proof}
设 $g\in\Gamma(Y,\OO_Y)$，并取 $\SSS=\OO_Y[T]$；于是 $g\in\Gamma(Y,\SSS_0)$。令
$h=gT\in\Gamma(Y,\SSS_1)$。若 $X=\Proj(\SSS)$，则（3.1.7）的典范等同把 $X_h$ 等同于 $Y$ 的开子集 $Y_g$（按0，5.5.2的意义）。事实上，可设 $Y=\Spec(A)$ 仿射；结合（2.2.5），只须注意分式环 $A_g$ 典范地等同于 $A[T]/(gT-1)A[T]$（0，1.2.3）。

\egastatement{3.1.8}{命题}
设 $\SSS$ 是非负分次拟凝聚 $\OO_Y$-代数。
\begin{enumerate}[label=(\roman*)]
\item 对每个 $d>0$，存在从 $\Proj(\SSS)$ 到 $\Proj(\SSS^{(d)})$ 的典范 $Y$-同构。
\item 令 $\SSS'$ 为 $\OO_Y$ 与各 $\SSS_n$（$n>0$）的直和所成的分次 $\OO_Y$-代数，则 $\Proj(\SSS')$ 与 $\Proj(\SSS)$ 典范地 $Y$-同构。
\item 设 $\LL$ 是可逆 $\OO_Y$-模（0，5.4.1），$\SSS_{(\LL)}$ 是各 $\SSS_d\otimes\LL^{\otimes d}$（$d\geq0$）的直和所成的分次 $\OO_Y$-代数。则 $\Proj(\SSS_{(\LL)})$ 与 $\Proj(\SSS)$ 典范地 $Y$-同构。
\end{enumerate}
\begin{proof}
三种情形中，都只须在 $Y$ 上局部地定义同构；与从开集到更小开集的限制操作的相容性容易验证。因此可设 $Y$ 仿射；此时（i）由（2.4.7，i）得出，（ii）由（2.4.8）得出。

对于（iii），还可设 $\LL\simeq\OO_Y$，因为问题在 $Y$ 上是局部的；这时两齐次谱的同构性显然。为定义典范同构，设 $Y=\Spec(A)$、$\SSS=\wt S$，其中 $S$ 是分次 $A$-代数；取自由 $A$-模 $L$ 的一个生成元 $c$，使 $\LL=\wt L$。对每个 $d\geq0$，映射 $x\mapsto x\otimes c^{\otimes d}$ 是从 $S_d$ 到 $S_d\otimes L^{\otimes d}$ 的 $A$-同构；这些同构定义分次代数的 $A$-同构
\[
\varphi_c:S\longrightarrow S_{(L)}=\bigoplus_{d\geq0}S_d\otimes L^{\otimes d}.
\]
设 $f\in S_d$ 为 $d$ 次齐次元。对 $x\in S_{nd}$ 和每个单位 $\varepsilon\in A$，有
\[
\frac{x\otimes c^{\otimes nd}}{(f\otimes c^{\otimes d})^n}
=\frac{x\otimes(\varepsilon c)^{\otimes nd}}{(f\otimes(\varepsilon c)^{\otimes d})^n}.
\]
所以由 $\varphi_c$ 诱导的同构
$S_{(f)}\to(S_{(L)})_{(f\otimes c^{\otimes d})}$ 与 $L$ 的生成元 $c$ 的选择无关。证明完毕。
\end{proof}

\egapar{3.1.9}
回顾（0，4.1.3及I，1.3.14），拟凝聚分次 $\OO_Y$-代数 $\SSS$ 由 $\OO_Y$-模 $\SSS_1$ 生成，当且仅当存在由仿射开集组成的覆盖 $(U_\alpha)$，使分次 $\Gamma(U_\alpha,\OO_Y)$-代数 $\Gamma(U_\alpha,\SSS)$ 由其一次齐次元的集合 $\Gamma(U_\alpha,\SSS_1)$ 生成。此时，对 $Y$ 的每个开集 $V$，$\SSS|V$ 由 $(\OO_Y|V)$-模 $\SSS_1|V$ 生成。

\egastatement{3.1.10}{命题}
假设存在 $Y$ 的有限仿射开覆盖 $(U_i)$，使每个分次代数 $\Gamma(U_i,\SSS)$ 在 $\Gamma(U_i,\OO_Y)$ 上为有限型。则存在 $d>0$，使 $\SSS^{(d)}$ 由 $\SSS_d$ 生成，且 $\SSS_d$ 是有限型 $\OO_Y$-模。
\begin{proof}
由（2.1.6，v），对每个 $i$，存在整数 $m_i$，使对每个 $n>0$ 有
$\Gamma(U_i,\SSS_{nm_i})=(\Gamma(U_i,\SSS_{m_i}))^n$。取 $d$ 为各 $m_i$ 的一个公倍数即可，并使用（2.1.6，i）。
\end{proof}

\egastatement{3.1.11}{推论}
在（3.1.10）的假设下，$\Proj(\SSS)$ 与某个齐次谱 $\Proj(\SSS')$ 是 $Y$-同构的，其中分次 $\OO_Y$-代数 $\SSS'$ 由 $\SSS'_1$ 生成，且 $\SSS'_1$ 是有限型 $\OO_Y$-模。
\begin{proof}
取 $\SSS'=\SSS^{(d)}$，其中 $d$ 满足（3.1.10）的性质，再应用（3.1.8，i）。
\end{proof}

\egapar{3.1.12}
若 $\SSS$ 是非负分次拟凝聚 $\OO_Y$-代数，则其幂零根 $\N$ 是拟凝聚 $\OO_Y$-模（I，5.1.1）。称 $\N_+=\N\cap\SSS_+$ 为 $\SSS_+$ 的幂零根；它是拟凝聚分次 $\SSS$-模，因为可立即归结到 $Y$ 仿射的情形，再用（2.1.10）。对每个 $y\in Y$，$(\N_+)_y$ 是 $(\SSS_+)_y=(\SSS_y)_+$ 的幂零根（I，5.1.1）。

若 $\N_+=0$，则称分次 $\OO_Y$-代数 $\SSS$ \textbf{本质既约}；这等价于对每个 $y\in Y$，$\SSS_y$ 是本质既约分次 $\OO_y$-代数。对任意分次 $\OO_Y$-代数 $\SSS$，$\SSS/\N_+$ 都本质既约。

若对每个 $y\in Y$，$\SSS_y$ 都是整环，并且 $(\SSS_y)_+=(\SSS_+)_y\ne0$，则称 $\SSS$ \textbf{为整的}。

\egastatement{3.1.13}{命题}
设 $\SSS$ 是非负分次 $\OO_Y$-代数。若 $X=\Proj(\SSS)$，则 $Y$-概形 $X_{\red}$ 典范地同构于 $\Proj(\SSS/\N_+)$。特别地，若 $\SSS$ 本质既约，则 $X$ 既约。
\begin{proof}
$X'=\Proj(\SSS/\N_+)$ 既约，由（2.4.4，i）立即得出，因为此性质是局部的。又对每个仿射开集 $U\subset Y$，有 $\varphi'^{-1}(U)=(\varphi^{-1}(U))_{\red}$，其中 $\varphi,\varphi'$ 是结构态射。立即可验证这些典范 $U$-态射 $\varphi'^{-1}(U)\to\varphi^{-1}(U)$ 与限制操作相容，故定义闭浸入 $X'\to X$，它在底空间上是同胚。结论由（I，5.1.2）得出。
\end{proof}

\egastatement{3.1.14}{命题}
设 $Y$ 是整预概形，$\SSS$ 是满足 $\SSS_0=\OO_Y$ 的拟凝聚分次 $\OO_Y$-代数。
\begin{enumerate}[label=(\roman*)]
\item 若 $\SSS$ 为整的（3.1.12），则 $X=\Proj(\SSS)$ 为整预概形，且结构态射 $\varphi:X\to Y$ 占优。
\item 再假设 $\SSS$ 本质既约。反过来，若 $X$ 为整预概形且 $\varphi$ 占优，则 $\SSS$ 为整的。
\end{enumerate}
\begin{proof}
（i）设 $(U_\alpha)$ 为 $Y$ 的非空仿射开集组成的一个基。只须在把 $Y$ 换成任一 $U_\alpha$、把 $\SSS$ 换成 $\SSS|U_\alpha$ 后证明命题。事实上，这将说明各 $\varphi^{-1}(U_\alpha)$ 是 $X$ 的不可约开集，因而非空；而任意两个这样的开集交非空，因为 $U_\alpha\cap U_\beta$ 包含某个 $U_\gamma$。故 $X$ 不可约（0，2.1.4）。另一方面，既约性为局部性质，故 $X$ 既约，从而为整预概形，且 $\varphi(X)$ 在 $Y$ 中稠密。

于是可设 $Y=\Spec(A)$，其中 $A$ 为整环（I，5.1.4），并设 $\SSS=\wt S$，其中 $S$ 是分次 $A$-代数。假设表示：对每个 $y\in Y$，$S_y=\wt S_y$ 是整分次环，且 $(S_y)_+\ne0$。只须证明 $S$ 是整环；因为此时 $S_+\ne0$，可应用（2.4.4，ii）。设 $f,g\in S$ 均非零而 $fg=0$。在每个 $S_y$ 中有 $(f/1)(g/1)=0$，故由假设 $f/1=0$ 或 $g/1=0$。例如若在某个 $S_y$ 中 $f/1=0$，则存在 $a\in A\setminus\mathfrak j_y$，使 $af=0$。对每个 $z\in Y$，在整环 $S_z$ 中有 $(a/1)(f/1)=0$；由于 $A$ 是整环，$a/1\ne0$，故 $f/1=0$。这推出 $f=0$，矛盾。

（ii）问题在 $Y$ 上是局部的，仍可设 $Y=\Spec(A)$，$A$ 为整环，$\SSS=\wt S$。由假设，对每个 $y$，$(S_y)_+$ 不含非零幂零元；而 $(S_0)_y=A_y$ 也不含非零幂零元。因此每个 $S_y$ 都既约，从而 $S$ 既约（I，5.1.1）。$X$ 为整预概形又推出 $S$ 本质为整（2.4.4，ii），故只须证明 $S_+$ 在 $A=S_0$ 中的零化子 $\mathfrak j$ 为零（2.1.11）。否则，对某个非零 $h\in\mathfrak j$ 有 $(S_h)_+=0$，所以 $\varphi^{-1}(D(h))=\varnothing$（3.1.1）。由于 $h$ 非幂零，$D(h)\ne\varnothing$，这与 $\varphi(X)$ 在 $Y$ 中稠密矛盾。
\end{proof}

\subsection{与分次 \texorpdfstring{$\SSS$}{S}-模相伴的 \texorpdfstring{$\Proj(\SSS)$}{Proj(S)} 上的层}
\src{54--56}
\egapar{3.2.1}
设 $Y$ 是预概形，$\SSS$ 是非负分次拟凝聚 $\OO_Y$-代数，$\M$ 是拟凝聚分次 $\SSS$-模。这里在 $(Y,\OO_Y)$ 上拟凝聚与在赋环空间 $(Y,\SSS)$ 上拟凝聚等价（I，9.6.1）。采用（3.1.1）的记号，把拟凝聚 $\OO_{X_U}$-模 $(\Gamma(U,\M))^\sim$ 记作 $\wt\M_U$。当 $U'\subset U$ 时，$\Gamma(U',\M)$ 典范地等同于 $\Gamma(U,\M)\otimes_A A'$（I，1.6.4），故
$\wt\M_{U'}=\rho_{U',U}^*(\wt\M_U)$（2.8.11）。

\egastatement{3.2.2}{命题}
在 $X=\Proj(\SSS)$ 上，存在唯一的拟凝聚 $\OO_X$-模 $\wt\M$，使对 $Y$ 的每个仿射开集 $U$，有
\[
\eta_U^*((\Gamma(U,\M))^\sim)=\wt\M|f^{-1}(U),
\]
其中 $f:X\to Y$ 为结构态射，$\eta_U:f^{-1}(U)\xrightarrow{\sim}\Proj(\Gamma(U,\SSS))$ 为前述同构。
\begin{proof}
由（3.1.2.1），$\rho_{U',U}$ 等同于嵌入态射 $f^{-1}(U')\to f^{-1}(U)$。故命题由关系
$\wt\M_{U'}=\rho_{U',U}^*(\wt\M_U)$ 和层的粘合原理（0，3.3.1）立即得出。
\end{proof}
称 $\wt\M$ 为与拟凝聚分次 $\SSS$-模 $\M$ \textbf{相伴的 $\OO_X$-模}。

\egastatement{3.2.3}{命题}
设 $\M$ 是拟凝聚分次 $\SSS$-模，$f\in\Gamma(Y,\SSS_d)$，$d>0$。若 $\gamma_f$ 是从 $X_f$ 到 $Y$-预概形
$Z_f=\Spec(\SSS^{(d)}/(f-1)\SSS^{(d)})$ 的典范同构（3.1.4），则
$\gamma_{f*}(\wt\M|X_f)$ 是与
$\M^{(d)}/(f-1)\M^{(d)}$ 相伴的 $\OO_{Z_f}$-模（1.4.3）。
\begin{proof}
问题在 $Y$ 上是局部的；结合图（2.8.12.1）的交换性，立即归结到（2.2.5）。
\end{proof}

\egastatement{3.2.4}{命题}
$\OO_X$-模 $\wt\M$ 是关于 $\M$ 的协变可加正合函子，从拟凝聚分次 $\SSS$-模范畴到拟凝聚 $\OO_X$-模范畴；它与归纳极限和直和可交换。
\begin{proof}
问题在 $Y$ 上是局部的，故归结到（I，1.3.11及1.3.9）和（2.5.4）。
\end{proof}
特别地，若 $\N$ 是 $\M$ 的拟凝聚分次 $\SSS$-子模，则 $\wt\N$ 典范地等同于 $\wt\M$ 的拟凝聚 $\OO_X$-子模。尤其，对 $\SSS$ 的每个拟凝聚分次理想层 $\J$，$\wt\J$ 是 $\OO_X$ 的拟凝聚理想层。

若 $\M$ 是拟凝聚分次 $\SSS$-模，$\J$ 是 $\OO_Y$ 的拟凝聚理想层，则 $\J\M$ 是 $\M$ 的拟凝聚分次 $\SSS$-子模，并有
\begin{equation}\tag{3.2.4.1}(\J\M)^\sim=\J\wt\M,\end{equation}
其中右边按（0，4.3.5）的意义解释。事实上，只须在 $Y=\Spec(A)$ 仿射时验证：设 $\SSS=\wt S$，其中 $S$ 是分次 $A$-代数；$\M=\wt M$，其中 $M$ 是分次 $S$-模；$\J=\wt{\mathfrak j}$，其中 $\mathfrak j$ 是 $A$ 的理想。对 $S_+$ 的每个齐次元 $f$，（3.2.4.1）左边在 $D_+(f)=\Spec(S_{(f)})$ 上的限制与
$(\mathfrak jM)_{(f)}=\mathfrak jM_{(f)}$ 相伴；由（I，1.3.13及1.6.9），右边的限制也如此。

\egastatement{3.2.5}{命题}
设 $f\in\Gamma(Y,\SSS_d)$，$d>0$。在开集 $X_f$ 上，对每个 $n\in\Z$，$(\OO_X|X_f)$-模
$(\SSS(nd))^\sim|X_f$ 典范地同构于 $\OO_X|X_f$。特别地，若 $\OO_Y$-代数 $\SSS$ 由 $\SSS_1$ 生成（3.1.9），则对每个 $n\in\Z$，$(\SSS(n))^\sim$ 是可逆 $\OO_X$-模。
\begin{proof}
设 $\varphi:X\to Y$ 为结构态射。对每个仿射开集 $U\subset Y$，结合（3.1.4），已在（2.5.7）中定义从
$(\SSS(nd))^\sim|(X_f\cap\varphi^{-1}(U))$ 到
$\OO_X|(X_f\cap\varphi^{-1}(U))$ 的典范同构。立即可验证这些同构与从 $U$ 到仿射开集 $U'\subset U$ 的限制相容，第一断言由此得出。若 $\SSS$ 由 $\SSS_1$ 生成，则存在 $Y$ 的仿射开覆盖 $(U_\alpha)$，使 $\Gamma(U_\alpha,\SSS)$ 由其一次分量 $\Gamma(U_\alpha,\SSS)_1=\Gamma(U_\alpha,\SSS_1)$ 生成。可逆性是局部性质，故第二断言归结到（2.5.9）。
\end{proof}
仍对每个 $n\in\Z$ 置
\begin{equation}\tag{3.2.5.1}\OO_X(n)=(\SSS(n))^\sim,\end{equation}
并对每个 $\OO_X$-模 $\F$ 置
\begin{equation}\tag{3.2.5.2}\F(n)=\F\otimes_{\OO_X}\OO_X(n).\end{equation}
立即由定义可得：对 $Y$ 的每个开集 $U$，
$((\SSS|U)(n))^\sim=\OO_X(n)|f^{-1}(U)$，其中 $f:X\to Y$ 为结构态射。

\egastatement{3.2.6}{命题}
设 $\M,\N$ 是两个拟凝聚分次 $\SSS$-模。存在关于 $\M,\N$ 的典范函子性同态
\begin{equation}\tag{3.2.6.1}
\lambda:\wt\M\otimes_{\OO_X}\wt\N\longrightarrow(\M\otimes_{\SSS}\N)^\sim,
\end{equation}
以及典范函子性同态
\begin{equation}\tag{3.2.6.2}
\mu:(\shHom_{\SSS}(\M,\N))^\sim\longrightarrow\shHom_{\OO_X}(\wt\M,\wt\N).
\end{equation}
此外，若 $\SSS$ 由 $\SSS_1$ 生成，则 $\lambda$ 是同构；若再有 $\M$ 具有有限表示（3.1.1），则 $\mu$ 也是同构。
\begin{proof}
当 $Y$ 仿射时，$\lambda,\mu$ 已在（2.5.11.2）和（2.5.12.2）中定义。这些定义是局部的，结合（2.8.14），立即推广到这里的一般情形。
\end{proof}

\egastatement{3.2.7}{推论}
若 $\SSS$ 由 $\SSS_1$ 生成，则对任意 $m,n\in\Z$，在典范同构意义下有
\begin{equation}\tag{3.2.7.1}\OO_X(m)\otimes_{\OO_X}\OO_X(n)=\OO_X(m+n),\end{equation}
\begin{equation}\tag{3.2.7.2}\OO_X(n)=(\OO_X(1))^{\otimes n}.\end{equation}

\egastatement{3.2.8}{推论}
假设 $\SSS$ 由 $\SSS_1$ 生成。对每个分次 $\SSS$-模 $\M$ 和每个 $n\in\Z$，在典范同构意义下有
\begin{equation}\tag{3.2.8.1}(\M(n))^\sim=\wt\M(n).\end{equation}
\begin{proof}
这由 $Y$ 仿射时的相应性质（2.5.14及2.5.15）及（2.8.11）得出。
\end{proof}

\egastatement{3.2.9}{注记}
（i）若 $\SSS=\A[T]$，其中 $\A$ 是拟凝聚 $\OO_Y$-代数（3.1.7），则立即可验证所有可逆 $\OO_X$-模 $\OO_X(n)$ 都典范地同构于 $\OO_X$。又设 $\N$ 是拟凝聚 $\A$-模，置 $\M=\N\otimes_{\A}\A[T]$。由（3.2.3）和（3.1.7），在 $X=\Proj(\A[T])$ 与 $X'=\Spec(\A)$ 的典范等同下，$\OO_X$-模 $\wt\M$ 等同于按（1.4.3）的意义与 $\N$ 相伴的 $\OO_{X'}$-模 $\wt\N$。

（ii）设 $\SSS$ 是任意分次 $\OO_Y$-代数，$\SSS'$ 为满足 $\SSS'_0=\OO_Y$ 且 $\SSS'_n=\SSS_n$（$n>0$）的分次 $\OO_Y$-代数。$X=\Proj(\SSS)$ 到 $X'=\Proj(\SSS')$ 的典范同构（3.1.8，ii）把 $\OO_X(n)$ 与 $\OO_{X'}(n)$ 等同。这由仿射情形的相应命题（2.5.16）及这些局部等同与限制操作的相容性得出。同样，令 $X^{(d)}=\Proj(\SSS^{(d)})$；$X$ 到 $X^{(d)}$ 的典范同构（3.1.8，i）把 $\OO_X(nd)$ 与 $\OO_{X^{(d)}}(n)$ 等同，对每个 $n\in\Z$ 皆如此。

\egastatement{3.2.10}{命题}
设 $\LL$ 是可逆 $\OO_Y$-模，$g:X_{(\LL)}=\Proj(\SSS_{(\LL)})\to X=\Proj(\SSS)$ 是典范同构（3.1.8，iii）。则对每个 $n\in\Z$，$g_*(\OO_{X_{(\LL)}}(n))$ 典范地同构于 $\OO_X(n)\otimes_Y\LL^{\otimes n}$。
\begin{proof}
先设 $Y$ 仿射，环为 $A$，且 $\LL=\wt L$，其中 $L$ 是秩为一的自由 $A$-模。采用（3.1.8，iii）的证明中的记号，对 $f\in S_d$，定义从
$(S(n))_{(f)}\otimes_A L^{\otimes n}$ 到 $(S_{(L)}(n))_{(f\otimes c^{\otimes d})}$ 的同构，把 $(x/f^k)\otimes c^{\otimes n}$（$x\in S_{kd+n}$）映为
\[
\frac{x\otimes c^{\otimes(n+kd)}}{(f\otimes c^{\otimes d})^k}.
\]
立即可见该同构不依赖 $L$ 的生成元 $c$ 的选择，并且对各个 $f\in S_+$ 如此定义的同构，与从 $D_+(f)$ 到 $D_+(fg)$ 的限制相容。一般情形下，回到（3.1.1）的定义，即可验证对 $Y$ 的各仿射开集 $U$ 所定义的同构与从 $U$ 到仿射开集 $U'\subset U$ 的限制相容。
\end{proof}

\subsection{与 \texorpdfstring{$\Proj(\SSS)$}{Proj(S)} 上的层相伴的分次 \texorpdfstring{$\SSS$}{S}-模}
\src{56--58}
本小节始终假设分次 $\OO_Y$-代数 $\SSS$ 由 $\SSS_1$ 生成（3.1.9）。回顾（3.1.8，i），在有限性条件（3.1.10）下，这一限制并非本质限制。

\egapar{3.3.1}
设 $p:X=\Proj(\SSS)\to Y$ 为结构态射。对每个 $\OO_X$-模 $\F$，置
\begin{equation}\tag{3.3.1.1}\Gstar(\F)=\bigoplus_{n\in\Z}p_*(\F(n)),\end{equation}
特别地，
\begin{equation}\tag{3.3.1.2}\Gstar(\OO_X)=\bigoplus_{n\in\Z}p_*(\OO_X(n)).\end{equation}
已知对两个 $\OO_X$-模 $\F,\G$ 有典范同态
$p_*(\F)\otimes_{\OO_Y}p_*(\G)\to p_*(\F\otimes_{\OO_X}\G)$（0，4.2.2）。因此由（3.2.7.1），$\Gstar(\OO_X)$ 具有分次 $\OO_Y$-代数结构；同样，（3.2.5.2）在 $\Gstar(\F)$ 上定义一个分次 $\Gstar(\OO_X)$-模结构。

由（3.2.5）及函子 $p_*$ 的左正合性（0，4.2.1），$\Gstar(\F)$ 是关于 $\F$ 的协变可加左正合函子，从 $\OO_X$-模范畴到分次 $\OO_Y$-模范畴，后者的态射为零次同态。特别地，若 $\J$ 是 $\OO_X$ 的理想层，则 $\Gstar(\J)$ 等同于 $\Gstar(\OO_X)$ 的分次理想层。

\egapar{3.3.2}
设 $\M$ 为拟凝聚分次 $\SSS$-模。对 $Y$ 的每个仿射开集 $U$，已在（2.6.2）定义阿贝尔群同态
\[
\alpha_{0,U}:\Gamma(U,\M_0)\longrightarrow\Gamma(p^{-1}(U),\wt\M).
\]
这些同态立即与限制操作可交换（2.8.13.1），因而定义阿贝尔群层同态
\begin{equation}\tag{3.3.2.1}\alpha_0:\M_0\longrightarrow p_*(\wt\M).\end{equation}
这里无需使用 $\SSS$ 由 $\SSS_1$ 生成的假设。对每个 $\M_n=(\M(n))_0$ 应用此结果，并使用（3.2.8.1），得到阿贝尔群层同态
\begin{equation}\tag{3.3.2.2}\alpha_n:\M_n\longrightarrow p_*(\wt\M(n))\qquad(n\in\Z),\end{equation}
从而得到分次阿贝尔群层的零次函子性同态
\begin{equation}\tag{3.3.2.3}\alpha:\M\longrightarrow\Gstar(\wt\M),\end{equation}
亦记作 $\alpha_{\M}$。特别取 $\M=\SSS$，可验证 $\alpha:\SSS\to\Gstar(\OO_X)$ 是分次 $\OO_Y$-代数同态，而（3.3.2.3）是相对于这个分次代数同态的分次模双同态。

还应注意，每个 $\alpha_n$ 对应一个典范 $\OO_X$-模同态（0，4.4.3）
\begin{equation}\tag{3.3.2.4}\alpha_n^\sharp:p^*(\M_n)\longrightarrow\wt\M(n).\end{equation}
不难验证：它正是按函子（3.2.4）对应于下列分次 $\OO_Y$-模的典范零次同态的同态：
\begin{equation}\tag{3.3.2.5}\M_n\otimes_{\OO_Y}\SSS\longrightarrow\M(n),\end{equation}
左边的分次自然来自 $\SSS$。事实上，可设 $Y=\Spec(A)$、$\M=\wt M$、$\SSS=\wt S$，且分次 $A$-代数 $S$ 由 $S_1$ 生成，故当 $f$ 遍历 $S_1$ 时，$D_+(f)$ 覆盖 $X$。回到（2.6.2）的定义并使用（I，1.6.7），可见（3.3.2.4）在 $D_+(f)$ 上的限制对应于（I，1.3.8）$S_{(f)}$-模同态
\[
M_n\otimes_A S_{(f)}\longrightarrow(S(n))_{(f)},\qquad x\otimes1\longmapsto x/1\quad(x\in M_n),
\]
由此即证得上述断言。\transnote{底本印刷第58页此处的目标写作 $(S(n))_{(f)}$；按照前面的（3.3.2.4）及 $x\in M_n$，目标应为 $(M(n))_{(f)}$。上式保留原印记号，特此标明。}

\egastatement{3.3.3}{命题}
对每个截面 $f\in\Gamma(Y,\SSS_d)$（$d>0$），$X_f$ 等于把 $\alpha_d(f)$ 看成 $\OO_X(d)$ 的截面时它不为零的点集（0，5.5.2）。
\begin{proof}
$\alpha_d(f)$ 是 $p_*(\OO_X(d))$ 在 $Y$ 上的截面；按定义，它也是 $\OO_X(d)$ 在 $X$ 上的截面（0，4.2.1）。由 $X_f$ 的定义（3.1.4），问题归结到 $Y$ 仿射的情形；这已在（2.6.3）处理。
\end{proof}

\egapar{3.3.4}
从现在起，除本小节开头的假设外，再假设：对每个拟凝聚 $\OO_X$-模 $\F$，各个 $p_*(\F(n))$ 都在 $Y$ 上拟凝聚，因而 $\Gstar(\F)=\bigoplus_{n\in\Z}p_*(\F(n))$ 也是拟凝聚 $\OO_Y$-模（I，1.4.1及1.3.9）。特别地，$X$ 在 $Y$ 上为有限型时，总有此性质（I，9.2.2）。因此 $(\Gstar(\F))^\sim$ 有定义，并为拟凝聚 $\OO_X$-模。对每个仿射开集 $U\subset Y$，由（I，1.3.9）和（2.5.4）有
\[
\begin{aligned}
\left(\Gamma\left(U,\bigoplus_{n\in\Z}p_*(\F(n))\right)\right)^\sim
&=\bigoplus_{n\in\Z}(\Gamma(U,p_*(\F(n))))^\sim\\
&=\bigoplus_{n\in\Z}(\Gamma(p^{-1}(U),\F(n)))^\sim\\
&=\left(\bigoplus_{n\in\Z}\Gamma(p^{-1}(U),\F(n))\right)^\sim\\
&=(\Gstar(\F|p^{-1}(U)))^\sim.
\end{aligned}
\]
故有典范同态（2.6.4）
\[
\beta_U:\left(\Gamma\left(U,\bigoplus_{n\in\Z}p_*(\F(n))\right)\right)^\sim
\longrightarrow\F|p^{-1}(U).
\]
图（2.8.13.2）的交换性说明这些同态与 $Y$ 上的限制操作可交换，故对拟凝聚 $\OO_X$-模得到典范函子性同态
\begin{equation}\tag{3.3.4.1}\beta:(\Gstar(\F))^\sim\longrightarrow\F,\end{equation}
也记作 $\beta_{\F}$。

\egastatement{3.3.5}{命题}
设 $\M$ 是拟凝聚分次 $\SSS$-模，$\F$ 是拟凝聚 $\OO_X$-模。则复合同态
\begin{equation}\tag{3.3.5.1}
\wt\M\xrightarrow{\wt\alpha}(\Gstar(\wt\M))^\sim\xrightarrow{\beta}\wt\M
\end{equation}
及
\begin{equation}\tag{3.3.5.2}
\Gstar(\F)\xrightarrow{\alpha}\Gstar((\Gstar(\F))^\sim)
\xrightarrow{\Gstar(\beta)}\Gstar(\F)
\end{equation}
均为恒等同构。
\begin{proof}
问题在 $Y$ 上是局部的，故归结到（2.6.5）。
\end{proof}

\subsection{有限性条件}
\src{59--61}
\egastatement{3.4.1}{命题}
设 $Y$ 为预概形，$\SSS$ 为由 $\SSS_1$ 生成的拟凝聚 $\OO_Y$-代数（3.1.9），并假设 $\SSS_1$ 为有限型。则 $X=\Proj(\SSS)$ 在 $Y$ 上为有限型。
\begin{proof}
问题在 $Y$ 上是局部的，可设 $Y$ 仿射，环为 $A$。于是 $\SSS=\wt S$，其中 $S=\Gamma(Y,\SSS)$。由假设，$S$ 由 $S_1=\Gamma(Y,\SSS_1)$ 生成，并且可设 $S_1$ 是有限型 $A$-模（I，1.3.9及1.3.12）。因此 $S$ 是有限型分次 $A$-代数，命题归结到（2.7.1，ii）。
\end{proof}

\egapar{3.4.2}
设 $\SSS$ 为拟凝聚分次 $\OO_Y$-代数。对拟凝聚分次 $\SSS$-模 $\M$，考虑下列有限性条件：
\begin{description}[leftmargin=3.8em,style=nextline]
\item[（TF）] 存在整数 $n$，使 $\SSS$-模 $\bigoplus_{k\geq n}\M_k$ 为有限型。
\item[（TN）] 存在整数 $n$，使 $k\geq n$ 时 $\M_k=0$。
\end{description}
若 $\M$ 满足（TN），则 $\wt\M=0$，因为问题在 $Y$ 上是局部的（2.7.2）。

设 $\M,\N$ 为拟凝聚分次 $\SSS$-模。若存在整数 $n$，使 $k\geq n$ 时 $u_k:\M_k\to\N_k$ 单射（分别地，满射、双射），则称零次同态 $u:\M\to\N$ 为\textbf{（TN）-单射}（分别地，\textbf{（TN）-满射}、\textbf{（TN）-双射}）。此时 $\wt u:\wt\M\to\wt\N$ 单射（分别地，满射、双射）；由（2.7.2），结合（I，1.3.9）及问题在 $Y$ 上的局部性，即得此结论。（TN）-双射也称为\textbf{（TN）-同构}。

\egastatement{3.4.3}{命题}
设 $Y$ 是预概形，$\SSS$ 是由 $\SSS_1$ 生成的拟凝聚分次 $\OO_Y$-代数，并假设 $\SSS_1$ 为有限型。设 $\M$ 是拟凝聚分次 $\SSS$-模。
\begin{enumerate}[label=(\roman*)]
\item 若 $\M$ 满足（TF），则 $\wt\M$ 为有限型。
\item 假设 $\M$ 满足（TF）。则 $\wt\M=0$，当且仅当 $\M$ 满足（TN）。
\end{enumerate}
\begin{proof}
这些问题在 $Y$ 上是局部的，故归结到 $Y$ 仿射的情形。设其环为 $A$，$\SSS=\wt S$，其中分次 $A$-代数 $S$ 的理想 $S_+$ 为有限型；设 $\M=\wt M$，其中 $M$ 是分次 $S$-模。命题由（2.7.3）得出。
\end{proof}

\egastatement{3.4.4}{定理}
设 $Y$ 为预概形，$\SSS$ 为由有限型 $\SSS_1$ 生成的拟凝聚分次 $\OO_Y$-代数，并置 $X=\Proj(\SSS)$。则对每个拟凝聚 $\OO_X$-模 $\F$，典范同态 $\beta$（3.3.4）是同构。
\begin{proof}
首先，（3.4.1）保证 $\beta$ 有定义。验证它是同构，可归结到 $Y$ 仿射，环为 $A$，$\SSS=\wt S$ 的情形，其中 $S$ 是由 $S_1$ 生成的分次 $A$-代数，且 $S_1$ 是有限型 $A$-模。应用（2.7.5）即可。
\end{proof}

\egastatement{3.4.5}{推论}
在（3.4.4）的假设下，每个拟凝聚 $\OO_X$-模 $\F$ 都同构于某个形如 $\wt\M$ 的 $\OO_X$-模，其中 $\M$ 为拟凝聚分次 $\SSS$-模。若再有 $\F$ 为有限型，且 $Y$ 为拟紧概形或其底空间为诺特空间，则可取 $\M$ 为有限型。
\begin{proof}
取 $\M=\Gstar(\F)$，第一断言立即由（3.4.4）得出。要证明第二断言，只须证明 $\M$ 是它的有限型分次 $\SSS$-子模 $\M_\lambda$ 的归纳极限。因为此时 $\wt\M$ 是各 $\wt\M_\lambda$ 的归纳极限（3.2.4），从而 $\F$ 是各 $\beta(\wt\M_\lambda)$ 的归纳极限；而 $X$ 拟紧（3.4.1及I，6.3.1），$\F$ 为有限型，故 $\F$ 必等于某个 $\beta(\wt\M_\lambda)$（0，5.2.3）。

为构造这样的 $\M_\lambda$，对每个 $n\in\Z$ 考虑拟凝聚 $\OO_Y$-模 $\M_n$。由关于 $Y$ 的假设，它是有限型 $\OO_Y$-子模 $\M_{n,\nu}$ 的归纳极限（I，9.4.9）。显然 $\N_{n,\nu}=\SSS\M_{n,\nu}$ 是有限型分次 $\SSS$-模；取这些 $\N_{n,\nu}$ 的有限和作为 $\M_\lambda$，立即可验证满足要求。
\end{proof}

\egastatement{3.4.6}{推论}
假设（3.4.4）的条件成立，且 $Y$ 的底空间拟紧。设 $\F$ 为有限型拟凝聚 $\OO_X$-模。则存在整数 $n_0$，使对每个 $n\geq n_0$，典范同态
$\sigma:p^*(p_*(\F(n)))\to\F(n)$（0，4.4.3）满射。
\begin{proof}
对每个 $y\in Y$，取其仿射开邻域 $U$。存在整数 $n_0(U)$，使 $n\geq n_0(U)$ 时，$\F(n)|p^{-1}(U)$ 由有限多个在 $p^{-1}(U)$ 上的截面生成（2.7.9）。这些截面是 $p^*(p_*(\F(n)))$ 在 $p^{-1}(U)$ 上的截面的典范像（0，3.7.1及4.4.3），故 $\F(n)|p^{-1}(U)$ 等于 $p^*(p_*(\F(n)))|p^{-1}(U)$ 的典范像。由于 $Y$ 拟紧，取其有限仿射开覆盖 $(U_i)$，再令 $n_0$ 为各 $n_0(U_i)$ 的最大值，即完成证明。
\end{proof}

\egastatement{3.4.7}{注记}
设 $p=(\psi,\theta):X\to Y$ 为赋环空间态射，$\F$ 为 $\OO_X$-模。典范同态 $\sigma:p^*(p_*(\F))\to\F$ 满射，可以具体表述如下（0，4.4.1）：对每个 $x\in X$ 和 $\F$ 在 $x$ 的某个开邻域 $V$ 上的每个截面 $s$，存在 $p(x)$ 在 $Y$ 中的开邻域 $U$、$\F$ 在 $p^{-1}(U)$ 上的有限多个截面 $t_i$（$1\leq i\leq m$）、$x$ 的开邻域 $W\subset V\cap p^{-1}(U)$，以及 $\OO_X$ 在 $W$ 上的截面 $a_i$，使
\[
s|W=\sum_i a_i(t_i|W).
\]
若 $Y$ 是仿射概形，且 $p_*(\F)$ 拟凝聚，则这一条件等价于 $\F$ 由它在 $X$ 上的截面生成（0，5.5.1）。事实上，若 $Y=\Spec(A)$，可设 $U=D(f)$，$f\in A$。对 $p_*(\F)$ 应用（I，1.4.1），可得整数 $n>0$ 和 $\F$ 在 $X$ 上的截面 $u_i$，使 $g^nt_i$ 为 $u_i$ 在 $p^{-1}(U)$ 上的限制，其中 $g=\theta(f)$。由于 $g$ 在 $p^{-1}(U)$ 上可逆，遂有
\[
s|W=\sum_i b_i(u_i|W),\qquad b_i=a_i(g^n|W)^{-1}.
\]
断言由此得证。当 $Y$ 仿射时，推论（3.4.6）因此重新给出（2.7.9）。

由此推出：若 $Y$ 为任意预概形，$\F$ 为拟凝聚 $\OO_X$-模且 $p_*(\F)$ 拟凝聚，则下列条件等价：
\begin{enumerate}[label=\alph*)]
\item 典范同态 $\sigma:p^*(p_*(\F))\to\F$ 满射。
\item 存在拟凝聚 $\OO_Y$-模 $\G$ 和满同态 $p^*(\G)\to\F$。
\item 对 $Y$ 的每个仿射开集 $U$，$\F|p^{-1}(U)$ 由其在 $p^{-1}(U)$ 上的截面生成。
\end{enumerate}
刚才已证得 a）与 c）等价。由于 $p_*(\F)$ 按假设拟凝聚，a）显然推出 b）。反之，每个同态 $u:p^*(\G)\to\F$ 都分解为
\[
p^*(\G)\longrightarrow p^*(p_*(\F))\xrightarrow{\sigma}\F
\]
（0，3.5.4.4）。故 $u$ 满射时 $\sigma$ 也满射，b）推出 a）。

\egastatement{3.4.8}{推论}
假设（3.4.4）的条件成立，且 $Y$ 是拟紧概形，或其底空间为诺特空间。设 $\F$ 为有限型拟凝聚 $\OO_X$-模。则存在整数 $n_0$，使对每个 $n\geq n_0$，$\F$ 同构于某个形如 $(p^*(\G))(-n)$ 的 $\OO_X$-模的商，其中 $\G$ 为有限型拟凝聚 $\OO_Y$-模，可依赖于 $n$。
\begin{proof}
结构态射 $p:X\to Y$ 分离且为有限型，故 $p_*(\F(n))$ 拟凝聚（I，9.2.2，b）。由 $Y$ 的假设，它是其有限型拟凝聚 $\OO_Y$-子模的归纳极限（I，9.4.9）。结合（3.4.6）、（0，4.3.2）及（0，5.2.3），$\F(n)$ 是某个 $p^*(\G)$ 的典范像，其中 $\G$ 是 $p_*(\F(n))$ 的有限型拟凝聚 $\OO_Y$-子模。推论由（3.2.5.2）和（3.2.7.1）得出。
\end{proof}

\subsection{函子性质}
\src{61--64}
\egapar{3.5.1}
设 $Y$ 是预概形，$\SSS,\SSS'$ 是两个非负分次拟凝聚 $\OO_Y$-代数。置 $X=\Proj(\SSS)$、$X'=\Proj(\SSS')$，并以 $p,p'$ 分别表示它们到 $Y$ 的结构态射。设 $\varphi:\SSS'\to\SSS$ 为分次 $\OO_Y$-代数同态。对 $Y$ 的每个仿射开集 $U$，置 $S_U=\Gamma(U,\SSS)$、$S'_U=\Gamma(U,\SSS')$、$A_U=\Gamma(U,\OO_Y)$。则 $\varphi$ 定义分次 $A_U$-代数同态 $\varphi_U:S'_U\to S_U$，它对应于 $p^{-1}(U)$ 中的开集 $G(\varphi_U)$ 和态射
$\Phi_U:G(\varphi_U)\to p'^{-1}(U)$（2.8.1）。若 $V\subset U$ 是仿射开集，则图
\begin{equation}\tag{3.5.1.1}
\begin{tikzcd}
S'_U \arrow[r,"\varphi_U"] \arrow[d] & S_U \arrow[d]\\
S'_V \arrow[r,"\varphi_V"'] & S_V
\end{tikzcd}
\end{equation}
交换。由定义（2.8.1）立即可验证
$G(\varphi_V)=G(\varphi_U)\cap p^{-1}(V)$，且 $\Phi_V$ 为 $\Phi_U$ 在 $G(\varphi_V)$ 上的限制。因此定义了 $X$ 的开子集 $G(\varphi)$，满足
$G(\varphi)\cap p^{-1}(U)=G(\varphi_U)$，以及仿射 $Y$-态射
$\Phi:G(\varphi)\to X'$；称它\textbf{与 $\varphi$ 相伴}，记作 $\Proj(\varphi)$。

若对每个 $y\in Y$，存在仿射邻域 $U$，使 $\Gamma(U,\SSS_+)$ 作为 $\Gamma(U,\OO_Y)$-模由 $\varphi(\Gamma(U,\SSS'_+))$ 生成，则 $G(\varphi_U)=p^{-1}(U)$，故 $G(\varphi)=X$。

\egastatement{3.5.2}{命题}
\begin{enumerate}[label=(\roman*)]
\item 若 $\M$ 是拟凝聚分次 $\SSS$-模，则 $\OO_{X'}$-模 $(\M_{[\varphi]})^\sim$ 与 $\Phi_*(\wt\M|G(\varphi))$ 之间有典范函子性同构。
\item 若 $\M'$ 是拟凝聚分次 $\SSS'$-模，则存在从 $(\OO_X|G(\varphi))$-模 $\Phi^*(\wt{\M'})$ 到
$(\M'\otimes_{\SSS'}\SSS)^\sim|G(\varphi)$ 的典范函子性同态 $\nu$。若 $\SSS'$ 由 $\SSS'_1$ 生成，则 $\nu$ 是同构。
\end{enumerate}
\begin{proof}
在 $Y$ 仿射的情形，所涉及的同态已在（2.8.7）和（2.8.8）定义。一般情形只须验证它们与从 $Y$ 的仿射开集到更小仿射开集的限制相容；这立即由（3.5.1.1）的交换性得出。
\end{proof}
特别地，对每个 $n\in\Z$ 有典范同态
\begin{equation}\tag{3.5.2.1}
\Phi^*(\OO_{X'}(n))\longrightarrow\OO_X(n)|G(\varphi).
\end{equation}

\egastatement{3.5.3}{命题}
设 $Y,Y'$ 是两个预概形，$\psi:Y'\to Y$ 为态射，$\SSS$ 是拟凝聚分次 $\OO_Y$-代数，置 $\SSS'=\psi^*(\SSS)$。则 $Y'$-概形 $X'=\Proj(\SSS')$ 典范地等同于 $\Proj(\SSS)\times_Y Y'$。此外，若 $\M$ 是拟凝聚分次 $\SSS$-模，则 $\OO_{X'}$-模 $(\psi^*(\M))^\sim$ 等同于 $\wt\M\otimes_Y\OO_{Y'}$。
\begin{proof}
首先，$\psi^*(\SSS),\psi^*(\M)$ 及其齐次分量均为拟凝聚 $\OO_{Y'}$-模（0，5.1.4）。设 $U$ 是 $Y$ 的仿射开集，$U'\subset\psi^{-1}(U)$ 是 $Y'$ 的仿射开集，$A,A'$ 分别为二者的环。于是 $\SSS|U=\wt S$，其中 $S$ 是分次 $A$-代数，而 $\SSS'|U'$ 等同于 $(S\otimes_A A')^\sim$（I，1.6.5）。第一断言由（2.8.10）及（I，3.2.6.2）得出，因为由这个等同所定义的投影
$\Proj(\SSS'|U')\to\Proj(\SSS|U)$ 立即与 $U,U'$ 上的限制相容，故确实定义态射 $\Proj(\SSS')\to\Proj(\SSS)$。

设 $q:\Proj(\SSS)\to Y$、$q':\Proj(\SSS')\to Y'$ 为结构态射。于是 $q'^{-1}(U')$ 等同于 $q^{-1}(U)\times_U U'$。由（2.8.10）和（I，1.6.5），层
$(\psi^*(\M))^\sim|q'^{-1}(U')$ 与
$(\wt\M\otimes_Y\OO_{Y'})|q'^{-1}(U')$ 都典范地等同于 $(M\otimes_A A')^\sim$，其中 $M=\Gamma(U,\M)$。这些等同与限制操作的相容性仍可立即验证，第二断言因此成立。
\end{proof}

\egastatement{3.5.4}{推论}
采用（3.5.3）的记号，并置 $X=\Proj(\SSS)$。则对每个 $n\in\Z$，$\OO_{X'}(n)$ 典范地等同于 $\OO_X(n)\otimes_Y\OO_{Y'}$。
\begin{proof}
由（3.5.3）的记号，显然对每个 $n\in\Z$ 有 $\psi^*(\SSS(n))=\SSS'(n)$。
\end{proof}

\egapar{3.5.5}
保留前述记号，以 $\Psi$ 表示典范投影 $X'\to X$，并置 $\M'=\psi^*(\M)$。再假设 $\SSS$ 由 $\SSS_1$ 生成，且 $X$ 在 $Y$ 上为有限型。由此 $\SSS'$ 由 $\SSS'_1$ 生成，可在 $Y,Y'$ 均仿射时立即验证；且 $X'$ 在 $Y'$ 上为有限型（I，6.3.4）。设 $\F$ 是 $\OO_X$-模，置 $\F'=\Psi^*(\F)$。由（3.5.4）及（0，4.3.3），对每个 $n\in\Z$ 有 $\F'(n)=\Psi^*(\F(n))$。

按如下方式定义典范 $\psi$-同态
$\theta_n:q_*(\F(n))\to q'_*(\F'(n))$。由于图
\[
\begin{tikzcd}
X \arrow[d,"q"'] & X' \arrow[l,"\Psi"'] \arrow[d,"q'"]\\
Y & Y' \arrow[l,"\psi"]
\end{tikzcd}
\]
交换，只须定义同态
\[
q_*(\F(n))\longrightarrow\psi_*(q'_*(\Psi^*(\F(n))))
=q_*(\Psi_*(\Psi^*(\F(n)))).
\]
取 $\theta_n=q_*(\rho_n)$ 即可，其中 $\rho_n:\F(n)\to\Psi_*(\Psi^*(\F(n)))$ 是典范同态（0，4.4.3）。对仿射开集 $U\subset Y$ 和仿射开集 $U'\subset\psi^{-1}(U)$，$\theta_n$ 在截面上给出的正是典范同态（0，3.7.2）
\[
\Gamma(q^{-1}(U),\F(n))\longrightarrow\Gamma(q'^{-1}(U'),\F'(n)).
\]
由（2.8.13.2）的交换性，当 $\F$ 拟凝聚时，图
\[
\begin{tikzcd}[column sep=large]
\F \arrow[r,"\rho"] & \F'\\
(\Gstar(\F))^\sim \arrow[u,"\beta_{\F}"] \arrow[r] & (\Gstar(\F'))^\sim \arrow[u,"\beta_{\F'}"']
\end{tikzcd}
\]
交换，其中上方横箭头为典范 $\Psi$-态射 $\F\to\Psi^*(\F)$。
同样，由（2.8.13.1）的交换性，图
\[
\begin{tikzcd}[column sep=large]
\Gstar(\wt\M) \arrow[r] & \Gstar(\wt{\M'})\\
\M \arrow[u,"\alpha_{\M}"] \arrow[r] & \M' \arrow[u,"\alpha_{\M'}"']
\end{tikzcd}
\]
交换，其中下方横箭头为典范 $\psi$-态射 $\M\to\psi^*(\M)$。

\egapar{3.5.6}
现考虑两个预概形 $Y,Y'$、态射 $g:Y'\to Y$、拟凝聚分次 $\OO_Y$-代数 $\SSS$ 与拟凝聚分次 $\OO_{Y'}$-代数 $\SSS'$，以及分次代数的 $g$-态射 $u:\SSS\to\SSS'$，即分次 $\OO_Y$-代数同态 $\SSS\to g_*(\SSS')$。已知这等价于给定分次 $\OO_{Y'}$-代数同态 $u^\sharp:g^*(\SSS)\to\SSS'$。由 $u^\sharp$ 典范地得到 $Y'$-态射
\[
w=\Proj(u^\sharp):G(u^\sharp)\longrightarrow\Proj(g^*(\SSS)),
\]
其中 $G(u^\sharp)$ 是 $X'=\Proj(\SSS')$ 的开集（3.5.1）。另一方面，置 $X=\Proj(\SSS)$，则 $X''=\Proj(g^*(\SSS))$ 典范地等同于 $X\times_Y Y'$（3.5.3）。将 $w$ 与第一投影 $p:X\times_Y Y'\to X$ 复合，得到态射
$v:G(u^\sharp)\to X$，记作 $\Proj(u)$，且图
\[
\begin{tikzcd}
G(u^\sharp) \arrow[r,"v"] \arrow[d] & X \arrow[d]\\
Y' \arrow[r,"g"'] & Y
\end{tikzcd}
\]
交换。

此外，对每个拟凝聚分次 $\SSS$-模 $\M$，有典范 $v$-态射
\begin{equation}\tag{3.5.6.1}
u:\wt\M\longrightarrow
(g^*(\M)\otimes_{g^*(\SSS)}\SSS')^\sim|G(u^\sharp).
\end{equation}
事实上，$u^\sharp$ 由下列同态复合而成：
\[
\begin{aligned}
v^*(\wt\M)=w^*(p^*(\wt\M))
&\longrightarrow w^*((g^*(\M))^\sim)\\
&\longrightarrow(g^*(\M)\otimes_{g^*(\SSS)}\SSS')^\sim|G(u^\sharp).
\end{aligned}
\]
第一箭头来自（3.5.3）的同构，第二箭头是（3.5.2，i）的同态。\transnote{底本这里引用（3.5.2，i）；所用的逆像同态实际上列在（3.5.2，ii）。保留原引文并指出其位置差异。}
当 $\SSS$ 由 $\SSS_1$ 生成时，由（3.5.2），$u^\sharp$ 是同构。
特别地，对每个 $n\in\Z$，有典范 $v$-态射
\begin{equation}\tag{3.5.6.2}u:\OO_X(n)\longrightarrow\OO_{X'}(n)|G(u^\sharp).\end{equation}

\subsection{预概形 \texorpdfstring{$\Proj(\SSS)$}{Proj(S)} 的闭子预概形}
\src{64--65}
\egapar{3.6.1}
设 $Y$ 为预概形，$\varphi:\SSS\to\SSS'$ 为拟凝聚分次 $\OO_Y$-代数的零次同态。若存在整数 $n$，使 $k\geq n$ 时 $\varphi_k:\SSS_k\to\SSS'_k$ 满射（分别地，单射、双射），则称 $\varphi$ 为\textbf{（TN）-满射}（分别地，\textbf{（TN）-单射}、\textbf{（TN）-双射}）。此时，相应态射 $\Phi:\Proj(\SSS')\to\Proj(\SSS)$ 的研究可以归结到 $\varphi$ 满射（分别地，单射、双射）的情形。证明与（2.9.1）相同，使用（3.1.8）；（2.9.1）就是 $Y$ 仿射的特殊情形。（TN）-双射也称为\textbf{（TN）-同构}。

\egastatement{3.6.2}{命题}
设 $Y$ 是预概形，$\SSS$ 是拟凝聚分次 $\OO_Y$-代数，$X=\Proj(\SSS)$。
\begin{enumerate}[label=(\roman*)]
\item 若 $\varphi:\SSS\to\SSS'$ 是分次 $\OO_Y$-代数的（TN）-满同态，则相应态射 $\Phi=\Proj(\varphi)$（3.5.1）在整个 $\Proj(\SSS')$ 上有定义，并且是到 $X$ 的闭浸入。若 $\J=\Ker\varphi$，则与 $\Phi$ 相伴的 $X$ 的闭子预概形由 $\OO_X$ 的拟凝聚理想层 $\wt\J$ 定义。
\item 再假设 $\SSS_0=\OO_Y$，$\SSS$ 由 $\SSS_1$ 生成，且 $\SSS_1$ 为有限型。设 $X'$ 是 $X$ 的闭子预概形，由 $\OO_X$ 的拟凝聚理想层 $\K$ 定义。令 $\J$ 为 $\Gstar(\K)$ 在典范同态 $\alpha:\SSS\to\Gstar(\OO_X)$（3.3.2）下的逆像所成的拟凝聚分次理想层，并置 $\SSS'=\SSS/\J$。则 $X'$ 是与典范分次 $\OO_Y$-代数同态 $\SSS\to\SSS'$ 所对应的闭浸入 $\Proj(\SSS')\to X$ 相伴的子预概形（I，4.2.1）。
\end{enumerate}
\begin{proof}
（i）可设 $\varphi$ 满射（3.6.1）。于是，对 $Y$ 的每个仿射开集 $U$，$\Gamma(U,\SSS)\to\Gamma(U,\SSS')$ 满射（I，1.3.9），故由（3.5.1），$G(\varphi)$ 为整个源齐次谱。立即可把问题归结到 $Y$ 仿射的情形，再应用（2.9.2，i）。

（ii）只须证明：由典范嵌入 $\J\to\SSS$ 所诱导的同态 $\wt\J\to\OO_X$ 是从 $\wt\J$ 到 $\K$ 的同构。问题在 $Y$ 上是局部的，可设 $Y$ 仿射，环为 $A$；于是 $\SSS=\wt S$，其中分次 $A$-代数 $S$ 由有限型 $A$-模 $S_1$ 生成。应用（2.9.2，ii）即可。
\end{proof}

\egastatement{3.6.3}{推论}
在（3.6.2，i）的条件下，再假设 $\SSS$ 由 $\SSS_1$ 生成。则对每个 $n\in\Z$，$\Phi^*(\OO_X(n))$ 典范地等同于 $\OO_{X'}(n)$。
\begin{proof}
当 $Y$ 仿射时，已经定义这样的典范同构（2.9.3）。一般情形只须验证，对 $Y$ 的各仿射开集 $U$ 如此定义的同构，与从 $U$ 到仿射开集 $U'\subset U$ 的限制相容；这是直接的。
\end{proof}

\egastatement{3.6.4}{推论}
设 $Y$ 为预概形，$\SSS$ 为由 $\SSS_1$ 生成的拟凝聚分次 $\OO_Y$-代数，$\M$ 为拟凝聚 $\OO_Y$-模，$u:\M\to\SSS_1$ 为 $\OO_Y$-模满同态。设 $\bar u:\Sym_{\OO_Y}(\M)\to\SSS$ 为延拓 $u$ 的分次 $\OO_Y$-代数同态（1.7.4）。则与 $\bar u$ 相应的态射是从 $\Proj(\SSS)$ 到 $\Proj(\Sym_{\OO_Y}(\M))$ 的闭浸入。
\begin{proof}
由假设 $\bar u$ 满射，应用（3.6.1，i）即可。\transnote{原文引用（3.6.1，i）；该处并无分项，所用结果为（3.6.2，i）。}
\end{proof}

\subsection{从预概形到齐次谱的态射}
\src{65--68}
\egapar{3.7.1}
设 $q:X\to Y$ 是预概形态射，$\LL$ 为可逆 $\OO_X$-模，$\SSS$ 为非负分次拟凝聚 $\OO_Y$-代数；于是 $q^*(\SSS)$ 是非负分次拟凝聚 $\OO_X$-代数。考虑拟凝聚分次 $\OO_X$-代数
$\SSS'=\bigoplus_{n\geq0}\LL^{\otimes n}$，并假设给定分次 $\OO_X$-代数同态
\[
\psi:q^*(\SSS)\longrightarrow\SSS'=\bigoplus_{n\geq0}\LL^{\otimes n}.
\]
这也等价于给定分次代数的 $q$-态射 $\psi^\flat:\SSS\to q_*(\SSS')$。已知 $\Proj(\SSS')$ 典范地等同于 $X$（3.1.7及3.1.8，iii）。由 $\psi$ 典范地得到 $X$ 的开集 $G(\psi)$ 和 $Y$-态射
\begin{equation}\tag{3.7.1.1}r_{\LL,\psi}:G(\psi)\longrightarrow\Proj(\SSS)=P,\end{equation}
称它为\textbf{与 $\psi$ 及 $\LL$ 相伴的态射}。回顾（3.5.6），此态射是 $Y$-态射
$r=\Proj(\psi):G(\psi)\to\Proj(q^*(\SSS))$ 与第一投影
$\pi:\Proj(q^*(\SSS))\simeq P\times_Y X\to P$ 的复合。

\egapar{3.7.2}
在 $Y=\Spec(A)$ 仿射时具体写出 $r=r_{\LL,\psi}$。此时 $\SSS=\wt S$，其中 $S$ 是非负分次 $A$-代数。先设 $X=\Spec(B)$ 也仿射，且 $\LL=\wt L$，其中 $L$ 为秩一自由 $B$-模。于是 $q^*(\SSS)=(S\otimes_A B)^\sim$（I，1.6.5）。若 $c$ 为 $L$ 的生成元，则
$\psi_n:q^*(\SSS_n)\to\LL^{\otimes n}$ 对应于同态
\[
w_n:S_n\otimes_A B\longrightarrow L^{\otimes n},\qquad
s\otimes b\longmapsto bv_n(s)c^{\otimes n},
\]
其中 $v_n:S_n\to B$ 为 $A$-模同态，诸 $v_n$ 合成代数同态 $S\to B$。设 $f\in S_d$，$d>0$，并置 $g=v_d(f)$。由（2.8.10）及 $D_+(f)$ 与 $\Spec(S_{(f)})$ 的等同（2.3.6），有
$\pi^{-1}(D_+(f))=D_+(f\otimes1)$。另一方面，由（2.8.1.1）并使用 $X$ 与 $\Proj(\SSS')$ 的典范等同，得
$r^{-1}(D_+(f\otimes1))=D(g)$。因此
\begin{equation}\tag{3.7.2.1}r_{\LL,\psi}^{-1}(D_+(f))=D(g).\end{equation}
态射 $\Proj(\psi)$ 在 $D(g)$ 上的限制对应于将
$(s\otimes1)/(f\otimes1)^n$（$s\in S_{nd}$）映为 $v_{nd}(s)/g^n$ 的同态（2.8.1），而 $\pi$ 在 $D_+(f\otimes1)$ 上的限制对应于
$s/f^n\mapsto(s\otimes1)/(f\otimes1)^n$。由此，$r_{\LL,\psi}$ 在 $D(g)$ 上的限制对应于 $A$-代数同态
$\omega:S_{(f)}\to B_g$，满足
$\omega(s/f^n)=v_{nd}(s)/g^n$（$s\in S_{nd},n>0$）。令 $X$ 为任意预概形而仍保持 $Y$ 仿射，结合（2.8.1），遂得下述命题。

\egastatement{3.7.3}{命题}
若 $Y=\Spec(A)$ 仿射，$\SSS=\wt S$，其中 $S$ 是分次 $A$-代数，则对每个 $f\in S_d=\Gamma(Y,\SSS_d)$，有
\begin{equation}\tag{3.7.3.1}r_{\LL,\psi}^{-1}(D_+(f))=X_{\psi^\flat(f)},\end{equation}
其中 $\psi^\flat(f)\in\Gamma(X,\LL^{\otimes d})$。而 $r_{\LL,\psi}$ 的限制
$X_{\psi^\flat(f)}\to D_+(f)=\Spec(S_{(f)})$ 对应于代数同态（I，2.2.4）
\begin{equation}\tag{3.7.3.2}
\psi^\flat_{(f)}:S_{(f)}\longrightarrow\Gamma(X_{\psi^\flat(f)},\OO_X),
\end{equation}
其对 $s\in S_{nd}=\Gamma(Y,\SSS_{nd})$ 的取值为
\begin{equation}\tag{3.7.3.3}
\psi^\flat_{(f)}(s/f^n)
=(\psi^\flat(s)|X_{\psi^\flat(f)})
(\psi^\flat(f)|X_{\psi^\flat(f)})^{-n}.
\end{equation}

若 $G(\psi)=X$，则称 $r_{\LL,\psi}$ \textbf{处处有定义}。显然，这当且仅当对每个仿射开集 $U\subset Y$，有
$G(\psi)\cap q^{-1}(U)=q^{-1}(U)$；换言之，这个问题在 $Y$ 上是局部的。若 $Y$ 仿射，则 $G(\psi)$ 是 $f$ 遍历 $S_+$ 中齐次元时诸 $r_{\LL,\psi}^{-1}(D_+(f))$ 的并（2.8.1）。由（3.7.3.1），诸 $X_{\psi^\flat(f)}$ 须覆盖 $X$，亦即如下所述。

\egastatement{3.7.4}{推论}
在（3.7.3）的假设下，$r_{\LL,\psi}$ 处处有定义，当且仅当对每个 $x\in X$，存在整数 $n>0$ 和 $\SSS_n$ 在 $Y$ 上的截面 $s$，使 $t=\psi^\flat(s)\in\Gamma(X,\LL^{\otimes n})$ 满足 $t(x)\ne0$。

当 $\psi$ 为（TN）-满射时，这个条件总满足。同样，$r_{\LL,\psi}$ 是否占优也是 $Y$ 上的局部问题，并有：

\egastatement{3.7.5}{推论}
在（3.7.3）的假设下，$r_{\LL,\psi}$ 占优，当且仅当对每个整数 $n>0$，凡截面 $s\in S_n$ 满足 $\psi^\flat(s)\in\Gamma(X,\LL^{\otimes n})$ 局部幂零，则 $s$ 本身幂零。
\begin{proof}
所须表达的是：$r_{\LL,\psi}^{-1}(D_+(s))$ 为空，仅当 $D_+(s)$ 为空。故推论由（3.7.3.1）和（2.3.7）得出。
\end{proof}

\egastatement{3.7.6}{命题}
设 $q:X\to Y$ 是态射，$\LL$ 是可逆 $\OO_X$-模，$\SSS,\SSS'$ 是拟凝聚分次 $\OO_Y$-代数，$u:\SSS'\to\SSS$ 是分次代数同态，$\psi:q^*(\SSS)\to\bigoplus_{n\geq0}\LL^{\otimes n}$ 是分次代数同态，并置 $\psi'=\psi\circ q^*(u)$。若 $r_{\LL,\psi'}$ 处处有定义，则 $r_{\LL,\psi}$ 也处处有定义。若 $u$ 为（TN）-满射，且 $r_{\LL,\psi'}$ 占优，则 $r_{\LL,\psi}$ 也占优。反之，若 $u$ 为（TN）-单射，且 $r_{\LL,\psi}$ 占优，则 $r_{\LL,\psi'}$ 占优。
\begin{proof}
由 $G(\psi')\subset G(\psi)$（2.8.4）即得第一断言。若 $u$ 为（TN）-满射，则 $\Proj(u):\Proj(\SSS)\to\Proj(\SSS')$ 处处有定义并为闭浸入。由于 $r_{\LL,\psi'}$ 是 $r_{\LL,\psi}|G(\psi')$ 与 $\Proj(u)$ 的复合，$r_{\LL,\psi'}$ 占优就推出 $r_{\LL,\psi}$ 占优。最后，若 $u$ 为（TN）-单射，则 $\Proj(u):G(u)\to\Proj(\SSS')$ 占优（2.8.3）。由于 $G(\psi')$ 是 $G(u)$ 在 $r_{\LL,\psi}$ 下的逆像，$r_{\LL,\psi}$ 占优便推出 $r_{\LL,\psi'}$ 占优。
\end{proof}

\egastatement{3.7.7}{命题}
设 $Y$ 是拟紧预概形，$q:X\to Y$ 为拟紧态射，$\LL$ 为可逆 $\OO_X$-模，$\SSS$ 为拟凝聚分次 $\OO_Y$-代数，它是拟凝聚分次 $\OO_Y$-代数归纳系 $(\SSS_\lambda)$ 的滤过归纳极限。令 $\varphi_\lambda:\SSS_\lambda\to\SSS$ 为典范同态，设
$\psi:q^*(\SSS)\to\bigoplus_{n\geq0}\LL^{\otimes n}$ 为分次代数同态，并置 $\psi_\lambda=\psi\circ q^*(\varphi_\lambda)$。则 $r_{\LL,\psi}$ 处处有定义，当且仅当存在 $\lambda$，使 $r_{\LL,\psi_\lambda}$ 处处有定义；此时，对每个 $\mu\geq\lambda$，$r_{\LL,\psi_\mu}$ 都处处有定义。
\begin{proof}
充分性由（3.7.6）得出。反过来，假设 $r_{\LL,\psi}$ 处处有定义。可以归结到 $Y$ 仿射的情形：若对每个仿射开集 $U\subset Y$，存在 $\lambda(U)$ 使 $r_{\LL,\psi_{\lambda(U)}}$ 在 $q^{-1}(U)$ 上处处有定义，则因 $Y$ 拟紧，可取其有限仿射开覆盖 $(U_i)$，再取对所有 $i$ 均有 $\lambda\geq\lambda(U_i)$ 的指标，应用（3.7.6）即可。

若 $Y$ 仿射，假设推出对每个 $x\in X$，存在某个 $S_n$ 中的截面 $s^{(x)}$，使 $t^{(x)}=\psi^\flat(s^{(x)})$ 满足 $t^{(x)}(x)\ne0$；这里把 $t^{(x)}$ 视为 $\LL^{\otimes n}$ 在 $X$ 上的截面。因此存在 $x$ 的邻域 $V(x)$，使其中每个 $z$ 都满足 $t^{(x)}(z)\ne0$。取有限多个 $V(x_i)$ 覆盖 $X$，令 $s^{(i)}$ 为相应的截面。则存在 $\lambda$，使全部 $s^{(i)}$ 都可写成 $\varphi_\lambda(s_\lambda^{(i)})$，$s_\lambda^{(i)}\in S^\lambda$。由（3.7.4），$r_{\LL,\psi_\lambda}$ 处处有定义。最后断言是（3.7.6）的直接结果。
\end{proof}

\egastatement{3.7.8}{推论}
在（3.7.7）的假设下，若诸 $r_{\LL,\psi_\lambda}$ 占优，则 $r_{\LL,\psi}$ 占优；若各 $\varphi_\lambda$ 单射，逆命题也成立。
\begin{proof}
第二断言是（3.7.6）的特例。为证明第一断言，可设 $Y$ 仿射。若 $s\in S$ 满足 $\psi^\flat(s)$ 局部幂零，因存在某个 $\lambda$ 使 $s=\varphi_\lambda(s_\lambda)$，由假设和（3.7.5）可得 $s_\lambda$ 幂零，从而 $s$ 幂零。应用（3.7.5）的准则即可。
\end{proof}

\egastatement{3.7.9}{注记}
（i）采用（3.7.1）的记号，并使用（3.2.10），对每个 $n\in\Z$ 有典范同态
\begin{equation}\tag{3.7.9.1}r_{\LL,\psi}^*(\OO_P(n))\longrightarrow\LL^{\otimes n},\end{equation}
其一般定义见（3.5.6.2）。立即可见，在（3.7.3）的条件下，它在 $X_{\psi^\flat(f)}$ 上的限制具体把
$s/f^k$（$s\in S_{n+kd}$）映到
\[
(\psi^\flat(s)|X_{\psi^\flat(f)})(\psi^\flat(f)|X_{\psi^\flat(f)})^{-k}.
\]

（ii）设 $\F$ 是拟凝聚 $\OO_X$-模，并假设 $q$ 拟紧且分离，使对每个 $n\geq0$，$q_*(\F\otimes\LL^{\otimes n})$ 为拟凝聚 $\OO_Y$-模（I，9.2.2）。令
$\M'=\bigoplus_{n\geq0}\F\otimes\LL^{\otimes n}$，它是拟凝聚分次 $\SSS'$-模。考虑其直像
\[
\M=q_*(\M')=\bigoplus_{n\geq0}q_*(\F\otimes\LL^{\otimes n}),
\]
通过 $\psi$，它成为拟凝聚 $\SSS$-模。下面构造典范 $\OO_X$-模同态
\begin{equation}\tag{3.7.9.2}\zeta:r_{\LL,\psi}^*(\wt\M)\longrightarrow\F|G(\psi).\end{equation}
事实上，已在（3.5.6.1）定义典范同态
\begin{equation}\tag{3.7.9.3}
r_{\LL,\psi}^*(\wt\M)\longrightarrow
(q^*(\M)\otimes_{q^*(\SSS)}\SSS')^\sim|G(\psi),
\end{equation}
其中右边看成 $\Proj(\SSS')$ 上的拟凝聚层。另一方面，对每个拟凝聚分次 $\SSS'$-模 $\M'$，有典范同态
\begin{equation}\tag{3.7.9.4}
q^*(q_*(\M'))\otimes_{q^*(\SSS)}\SSS'\longrightarrow\M'.
\end{equation}
其定义为：对 $X$ 的每个开集 $U$、$q^*(q_*(\M'))$ 在 $U$ 上的截面 $t'$ 和 $\SSS'$ 在 $U$ 上的截面 $b'$，把 $t'\otimes b'$ 映为 $\M'$ 的截面 $b'\sigma(t')$；其中 $\sigma(t')$ 是与 $t'$ 典范相应的截面（0，4.4.3）。因此得到典范同态
\begin{equation}\tag{3.7.9.5}
(q^*(q_*(\M'))\otimes_{q^*(\SSS)}\SSS')^\sim|G(\psi)
\longrightarrow\wt{\M'}|G(\psi).
\end{equation}
最后，$\wt{\M'}$ 典范地等同于 $\F$（3.2.9，i），故得到所求典范同态。

在（3.7.3）的条件下，它在 $X_{\psi^\flat(f)}$ 上的限制可如下具体说明。给定 $\F\otimes\LL^{\otimes nd}$ 在 $X$ 上的截面 $z$，令 $z'$ 为把 $z$ 看成 $q_*(\F\otimes\LL^{\otimes nd})$ 在 $Y$ 上的截面，则 $z'/f^n$ 对应于 $\F$ 在 $X_{\psi^\flat(f)}$ 上的截面
\[
(z|X_{\psi^\flat(f)})(\psi^\flat(f)|X_{\psi^\flat(f)})^{-n}.
\]

\subsection{浸入齐次谱的判别准则}
\src{69--71}
\egapar{3.8.1}
采用（3.7.1）的记号。$r_{\LL,\psi}$ 是否为浸入（分别地，开浸入、闭浸入），显然是 $Y$ 上的局部问题。

\egastatement{3.8.2}{命题}
在（3.7.3）的假设下，$r_{\LL,\psi}$ 处处有定义且为浸入，当且仅当存在一族截面 $s_\alpha\in S_{n_\alpha}$（$n_\alpha>0$），使置 $f_\alpha=\psi^\flat(s_\alpha)$ 后，满足下列条件：
\begin{enumerate}[label=(\roman*)]
\item 各 $X_{f_\alpha}$ 覆盖 $X$。
\item 各 $X_{f_\alpha}$ 为仿射开集。
\item 对每个 $\alpha$ 及每个 $t\in\Gamma(X_{f_\alpha},\OO_X)$，存在整数 $m>0$ 和 $s\in S_{mn_\alpha}$，使
$t=(\psi^\flat(s)|X_{f_\alpha})(f_\alpha|X_{f_\alpha})^{-m}$。
\end{enumerate}
$r_{\LL,\psi}$ 处处有定义且为开浸入，当且仅当存在一族 $(s_\alpha)$ 满足（i）、（ii）、（iii）及：
\begin{enumerate}[label=(\roman*),start=4]
\item 对每个 $n>0$ 和每个 $s\in S_n$，若 $\psi^\flat(s)|X_{f_\alpha}=0$，则存在整数 $k>0$，使 $s_\alpha^ks=0$。
\end{enumerate}
$r_{\LL,\psi}$ 处处有定义且为闭浸入，当且仅当存在一族 $(s_\alpha)$ 满足（i）、（ii）、（iii）及：
\begin{enumerate}[label=(\roman*),start=5]
\item 各 $D_+(s_\alpha)$ 覆盖 $P=\Proj(S)$。
\end{enumerate}
\begin{proof}
$r=r_{\LL,\psi}$ 为浸入（分别地，闭浸入），当且仅当存在由 $D_+(s_\alpha)$ 组成的 $r(G(\psi))$（分别地，$P$）的覆盖，使置 $U_\alpha=r^{-1}(D_+(s_\alpha))$ 后，$r|U_\alpha$ 是从 $U_\alpha$ 到 $D_+(s_\alpha)$ 的闭浸入（I，4.2.4）。由（3.7.3.1），条件（i）同时表达 $r$ 处处有定义和诸 $D_+(s_\alpha)$ 覆盖 $r(X)$。因为 $D_+(s_\alpha)$ 仿射，（ii）与（iii）表达 $r|X_{f_\alpha}$ 是到 $D_+(s_\alpha)$ 的闭浸入（I，4.2.3）。最后，（iii）与（iv）表达 $\psi^\flat_{(s_\alpha)}$ 为双射（记号见3.7.3.2）；故（ii）、（iii）、（iv）表达每个 $r|X_{f_\alpha}$ 是到 $D_+(s_\alpha)$ 的同构，而（i）至（iv）一起表达 $r$ 是开浸入。
\end{proof}

\egastatement{3.8.3}{推论}
在（3.7.6）的假设下，若 $r_{\LL,\psi'}$ 处处有定义且为浸入，则 $r_{\LL,\psi}$ 也如此。若再假设 $u$ 为（TN）-满射，且 $r_{\LL,\psi'}$ 为开浸入（分别地，闭浸入），则 $r_{\LL,\psi}$ 也为开浸入（分别地，闭浸入）。
\begin{proof}
由（3.8.2），存在一族 $s'_\alpha\in S'_{n_\alpha}$，使置 $f_\alpha=\psi'^\flat(s'_\alpha)$ 后，（i）、（ii）、（iii）成立。置 $s_\alpha=u(s'_\alpha)$，仍有 $f_\alpha=\psi^\flat(s_\alpha)$；若
$t=(\psi'^\flat(s')|X_{f_\alpha})(f_\alpha|X_{f_\alpha})^{-m}$，则置 $s=u(s')$ 后，也有
$t=(\psi^\flat(s)|X_{f_\alpha})(f_\alpha|X_{f_\alpha})^{-m}$。第一断言由此得证。第二断言立即由 $\Proj(u)$ 是闭浸入得出。
\end{proof}

\egastatement{3.8.4}{命题}
假设（3.7.7）的条件成立，且 $q:X\to Y$ 为有限型态射。则 $r_{\LL,\psi}$ 处处有定义且为浸入，当且仅当存在 $\lambda$，使 $r_{\LL,\psi_\lambda}$ 处处有定义且为浸入；此时对每个 $\mu\geq\lambda$，$r_{\LL,\psi_\mu}$ 都处处有定义且为浸入。
\begin{proof}
结合（3.8.3），只须证明：若 $r_{\LL,\psi}$ 处处有定义且为浸入，则至少对一个 $\lambda$，$r_{\LL,\psi_\lambda}$ 也如此。与（3.7.7）中相同，使用 $Y$ 的拟紧性，先归结到 $Y$ 仿射的情形。由于 $X$ 拟紧，由（3.8.2）存在 $S$ 中的有限族元素 $s_i\in S_{n_i}$，具有（i）、（ii）、（iii）的性质。令 $f_i=\psi^\flat(s_i)$。态射 $X_{f_i}\to Y$ 为有限型：它是仿射概形间的态射，因而拟紧（I，6.6.1）；又因 $q$ 为有限型，它局部为有限型（I，6.3.2），故由（I，6.6.3）得出结论。因此 $X_{f_i}$ 的环 $B_i$ 是有限型 $A$-代数（I，6.3.3）。取该代数的有限生成元族 $(t_{ij})$。由假设，存在 $s_{ij}\in S_{m_{ij}n_i}$，使
\[
t_{ij}=(\psi^\flat(s_{ij})|X_{f_i})(\psi^\flat(s_i)|X_{f_i})^{-m_{ij}}.
\]
存在 $\lambda$ 以及 $s'_i\in S^\lambda_{n_i}$、$s'_{ij}\in S^\lambda_{m_{ij}n_i}$，其在 $\varphi_\lambda$ 下的像分别为 $s_i,s_{ij}$。于是 $r_{\LL,\psi_\lambda}$ 显然满足（3.8.2）的（i）、（ii）、（iii）。
\end{proof}

\egastatement{3.8.5}{命题}
设 $Y$ 是拟紧概形，或底空间为诺特空间的预概形，$q:X\to Y$ 是有限型态射，$\LL$ 是可逆 $\OO_X$-模，$\SSS$ 是拟凝聚分次 $\OO_Y$-代数，$\psi:q^*(\SSS)\to\bigoplus_{n\geq0}\LL^{\otimes n}$ 为分次代数同态。则 $r_{\LL,\psi}$ 处处有定义且为浸入，当且仅当存在整数 $n>0$ 及 $\SSS_n$ 的有限型拟凝聚 $\OO_Y$-子模 $\E$，使：
\begin{enumerate}[label=\alph*)]
\item 同态 $\psi_n\circ q^*(j_n):q^*(\E)\to\LL^{\otimes n}$ 满射，其中 $j_n:\E\to\SSS_n$ 为典范嵌入。
\item 以 $\SSS'$ 表示 $\E$ 生成的分次 $\OO_Y$-子代数，$j':\SSS'\to\SSS$ 表示嵌入，并置 $\psi'=\psi\circ q^*(j')$，则 $r_{\LL,\psi'}$ 处处有定义且为浸入。
\end{enumerate}
此时，每个包含 $\E$ 的 $\SSS_n$ 的拟凝聚 $\OO_Y$-子模 $\E'$ 具有同样性质；对每个 $k>0$，$\E^{\otimes k}$ 在 $\SSS_{nk}$ 中的像也具有同样性质。
\begin{proof}
条件的充分性和最后两个断言都是（3.8.3）的特殊情形，并使用 $\Proj(\SSS)$ 与 $\Proj(\SSS^{(d)})$ 的典范同构（3.1.8）。

设 $(U_i)$ 是 $Y$ 的有限仿射开覆盖，$A_i=A(U_i)$。由于 $q^{-1}(U_i)$ 拟紧，$r_{\LL,\psi}$ 为处处有定义的浸入这一假设与（3.8.2）一起，推出存在 $S^{(i)}=\Gamma(U_i,\SSS)$ 中的有限族元素 $s_{ij}\in S^{(i)}_{n_{ij}}$，具有（i）、（ii）、（iii）的性质。令 $f_{ij}=\psi^\flat(s_{ij})$。与（3.8.4）的证明相同，态射 $X_{f_{ij}}\to U_i$ 为有限型，故 $X_{f_{ij}}$ 的环 $B_{ij}$ 是有限型 $A_i$-代数，且有形如
\[
(\psi^\flat(t_{ijk})|X_{f_{ij}})
(\psi^\flat(s_{ij})|X_{f_{ij}})^{-m_{ijk}}
\]
的生成元系，其中 $t_{ijk}\in S^{(i)}_{m_{ijk}n_{ij}}$。令 $n$ 为诸 $m_{ijk}n_{ij}$ 的一个公倍数；对每个 $(i,j,k)$，把 $s_{ij}$ 换成满足 $p m_{ijk}n_{ij}=n$ 的适当幂 $s_{ij}^{p}$，并把 $t_{ijk}$ 乘以 $s_{ij}^{(p-1)m_{ijk}}$，可设对每个 $i$，全部 $s_{ij},t_{ijk}$ 都属于 $S^{(i)}_n$，且 $m_{ijk}=1$。\transnote{这一句依原文保留。原文在替换分母生成元与统一 $m_{ijk}=1$ 的叙述之间省略了重新选取次数为 $n$ 的分母元的步骤；上述等次数化是原文所要使用的结论。}

固定 $i$，以 $E_i$ 表示这些元素生成的 $S^{(i)}_n$ 的 $A_i$-子模。存在 $\SSS_n$ 的有限型凝聚 $\OO_Y$-子模 $\E_i$，使 $\E_i|U_i=\wt E_i$（I，9.4.7）。\transnote{底本这里写“凝聚”（cohérent），而命题及所引延拓结果所需的是“拟凝聚”（quasi-cohérent）；这里保留原文字样，不另增诺特假设。}
显然，诸 $\E_i$ 的和 $\E$ 满足要求：每个截面 $f_{ij}$ 都有如下性质，对每个 $x\in X_{f_{ij}}$，存在仿射邻域 $V\subset X_{f_{ij}}$，使 $f_{ij}|V$ 是 $\Gamma(V,\LL^{\otimes n})$ 的一组基。
\end{proof}

\section{射影丛。丰沛层}
\subsection{射影丛的定义}
\src{71--72}
\egastatement{4.1.1}{定义}
设 $Y$ 是预概形，$\E$ 是拟凝聚 $\OO_Y$-模，$\Sym_{\OO_Y}(\E)$ 是 $\E$ 的对称 $\OO_Y$-代数（1.7.4）；它是拟凝聚的（1.7.7）。称 $Y$-概形
$P=\Proj(\Sym_{\OO_Y}(\E))$ 为\textbf{由 $\E$ 定义的 $Y$ 上的射影丛}，记作 $\Pb(\E)$。$\OO_P$-模 $\OO_P(1)$ 称为 $P$ 上的\textbf{基本层}。

若 $Y$ 仿射，环为 $A$，则 $\E=\wt E$，其中 $E$ 是 $A$-模；此时把 $\Pb(\wt E)$ 写作 $\Pb(E)$。当 $\E=\OO_Y^n$ 时，把 $\Pb(\E)$ 写作 $\Pb_Y^{n-1}$；若 $Y$ 仿射，环为 $A$，也写作 $\Pb_A^{n-1}$。由于 $\Sym_{\OO_Y}(\OO_Y)$ 典范地等同于 $\OO_Y[T]$，$\Pb_Y^0$ 典范地等同于 $Y$（3.1.7）。（2.4.3）的例子正是 $\Pb_K^1$。

\egapar{4.1.2}
设 $\E,\F$ 为拟凝聚 $\OO_Y$-模，$u:\E\to\F$ 为 $\OO_Y$-模同态；它典范地对应于分次 $\OO_Y$-代数同态
$\Sym(u):\Sym_{\OO_Y}(\E)\to\Sym_{\OO_Y}(\F)$（1.7.4）。若 $u$ 满射，则 $\Sym(u)$ 也满射，故 $\Proj(\Sym(u))$ 为闭浸入 $\Pb(\F)\to\Pb(\E)$（3.6.2，i），记作 $\Pb(u)$。因此，若把拟凝聚 $\OO_Y$-模之间的态射限制为满同态，则 $\Pb(\E)$ 是关于 $\E$ 的反变函子。

仍假设 $u$ 满射，并置 $P=\Pb(\E)$、$Q=\Pb(\F)$、$j=\Pb(u)$；由（3.6.3），在同构意义下有
\begin{equation}\tag{4.1.2.1}j^*(\OO_P(n))=\OO_Q(n)\qquad(n\in\Z).\end{equation}

\egapar{4.1.3}
设 $\psi:Y'\to Y$ 为态射，置 $\E'=\psi^*(\E)$。则
$\Sym_{\OO_{Y'}}(\E')=\psi^*(\Sym_{\OO_Y}(\E))$（1.7.5），故由（3.5.3），在典范同构意义下有
\begin{equation}\tag{4.1.3.1}\Pb(\psi^*(\E))=\Pb(\E)\times_Y Y'.\end{equation}
而且，对每个 $n\in\Z$，显然有
$\psi^*((\Sym_{\OO_Y}(\E))(n))=(\Sym_{\OO_{Y'}}(\E'))(n)$。
因此，置 $P=\Pb(\E)$、$P'=\Pb(\E')$，由（3.5.4），在同构意义下有
\begin{equation}\tag{4.1.3.2}\OO_{P'}(n)=\OO_P(n)\otimes_Y\OO_{Y'}\qquad(n\in\Z).\end{equation}

\egastatement{4.1.4}{命题}
设 $\LL$ 是可逆 $\OO_Y$-模。对每个拟凝聚 $\OO_Y$-模 $\E$，存在典范 $Y$-同构
$i_{\LL}:\Pb(\E)\xrightarrow{\sim}\Pb(\E\otimes\LL)$。若置 $P=\Pb(\E)$、$Q=\Pb(\E\otimes\LL)$，则对每个 $n\in\Z$，$i_{\LL}^*(\OO_Q(n))$ 典范地同构于 $\OO_P(n)\otimes_Y\LL^{\otimes n}$。
\begin{proof}
先注意：若 $A$ 是环，$E$ 是 $A$-模，$L$ 是秩一自由 $A$-模，则有典范 $A$-模同态
\[
\Sym_A^n(E\otimes L)\longrightarrow\Sym_A^n(E)\otimes L^{\otimes n},
\]
把 $(x_1\otimes y_1)\cdots(x_n\otimes y_n)$ 映为
$(x_1\cdots x_n)\otimes(y_1\otimes\cdots\otimes y_n)$，其中 $x_i\in E,y_i\in L$。归结到 $L=A$ 的情形，立即可验证它是同构。因此有分次 $A$-代数的典范同构
\[
\Sym_A(E\otimes L)\xrightarrow{\sim}\bigoplus_{n\geq0}\Sym_A^n(E)\otimes L^{\otimes n}.
\]
回到命题的条件，前述说明给出分次 $\OO_Y$-代数的典范同构
\begin{equation}\tag{4.1.4.1}
\Sym_{\OO_Y}(\E\otimes\LL)\xrightarrow{\sim}
\bigoplus_{n\geq0}\Sym_{\OO_Y}^n(\E)\otimes\LL^{\otimes n}.
\end{equation}
在预层上定义此同构，并使用（1.7.4）、（I，1.3.9）和（I，1.3.12）即可。命题由（3.1.8，iii）及（3.2.10）得出。
\end{proof}

\egapar{4.1.5}
采用（4.1.1）的假设，置 $P=\Pb(\E)$，并以 $p:P\to Y$ 表示结构态射。由于按定义 $\E=(\Sym_{\OO_Y}(\E))_1$，存在典范同态 $\alpha_1:\E\to p_*(\OO_P(1))$（3.3.2.2），从而也存在典范同态（0，4.4.3）
\begin{equation}\tag{4.1.5.1}\alpha_1^\sharp:p^*(\E)\longrightarrow\OO_P(1).\end{equation}

\egastatement{4.1.6}{命题}
典范同态（4.1.5.1）满射。
\begin{proof}
已在（3.3.2）看到，$\alpha_1^\sharp$ 按函子性对应于典范同态
$\E\otimes_{\OO_Y}\Sym_{\OO_Y}(\E)\to(\Sym_{\OO_Y}(\E))(1)$。
由定义，$\E$ 生成 $\Sym_{\OO_Y}(\E)$，所以该同态在所需的正次数部分满射；结论由（3.2.4）得出。\transnote{原文直接称此分次同态“满射”。若包括移位模的负次数，准确地说它是（TN）-满射；相伴层同态的满射性不受影响，参见（3.4.2）。}
\end{proof}

\subsection{从预概形到射影丛的态射}
\src{72--76}
\egapar{4.2.1}
保留（4.1.5）的记号。设 $X$ 是 $Y$-预概形，$q:X\to Y$ 是结构态射，$r:X\to P$ 是 $Y$-态射，即图
\[
\begin{tikzcd}
P \arrow[d,"p"'] & X \arrow[l,"r"'] \arrow[dl,"q"]\\
Y &
\end{tikzcd}
\]
交换。由于函子 $r^*$ 右正合（0，4.3.1），由满同态（4.1.5.1）得到满同态
\[
r^*(\alpha_1^\sharp):r^*(p^*(\E))\longrightarrow r^*(\OO_P(1)).
\]
这里 $r^*(p^*(\E))=q^*(\E)$，而 $r^*(\OO_P(1))$ 局部同构于 $r^*(\OO_P)=\OO_X$，故是可逆 $\OO_X$-模 $\LL_r$。这样，由 $r$ 定义了典范的满 $\OO_X$-模同态
\begin{equation}\tag{4.2.1.1}\varphi_r:q^*(\E)\longrightarrow\LL_r.\end{equation}

若 $Y=\Spec(A)$ 仿射，$\E=\wt E$，则此同态还可具体表述如下。给定 $f\in E$，先由（2.6.3）得
\begin{equation}\tag{4.2.1.2}r^{-1}(D_+(f))=X_{\varphi_r^\flat(f)}.\end{equation}
设 $V$ 是 $X$ 中包含于 $r^{-1}(D_+(f))$ 的仿射开集，其环为 $A$-代数 $B$。置 $S=\Sym_A(E)$。$r|V$ 对应于 $A$-同态 $\omega:S_{(f)}\to B$；有 $q^*(\E)|V=(E\otimes_A B)^\sim$，且 $\LL_r|V=\wt L_V$，其中
$L_V=(S(1))_{(f)}\otimes_{S_{(f)}}B$（I，1.6.5）。$\varphi_r$ 在 $q^*(\E)|V$ 上的限制对应于 $B$-同态
\[
u:E\otimes_A B\longrightarrow L_V,\qquad
x\otimes1\longmapsto(x/1)\otimes1=(f/1)\otimes\omega(x/f).
\]
$\varphi_r$ 到 $\OO_X$-代数同态
\[
\psi_r:q^*(\Sym(\E))=\Sym(q^*(\E))\longrightarrow
\Sym(\LL_r)=\bigoplus_{n\geq0}\LL_r^{\otimes n}
\]
的典范延拓具有如下性质：其在 $q^*(\Sym_n(\E))|V$ 上的限制，对应于把 $x\otimes1$ 映为 $(f/1)^{\otimes n}\otimes\omega(x/f^n)$ 的同态
$\Sym_A^n(E)\otimes_A B\to L_V^{\otimes n}$。

\egapar{4.2.2}
反过来，设给定态射 $q:X\to Y$、可逆 $\OO_X$-模 $\LL$ 和拟凝聚 $\OO_Y$-模 $\E$。每个同态 $\varphi:q^*(\E)\to\LL$ 典范地对应于拟凝聚 $\OO_X$-代数同态
\[
\psi:\Sym(q^*(\E))=q^*(\Sym(\E))\longrightarrow\bigoplus_{n\geq0}\LL^{\otimes n},
\]
从而对应于 $Y$-态射（3.7.1）
$r_{\LL,\psi}:G(\psi)\to\Proj(\Sym(\E))=\Pb(\E)$，记作 $r_{\LL,\varphi}$。若 $\varphi$ 满射，则 $\psi$ 也满射，故 $r_{\LL,\varphi}$ 处处有定义（3.7.4）。此外，有如下命题。

\egastatement{4.2.3}{命题}
给定态射 $q:X\to Y$ 与拟凝聚 $\OO_Y$-模 $\E$。映射
$r\mapsto(\LL_r,\varphi_r)$ 和 $(\LL,\varphi)\mapsto r_{\LL,\varphi}$，在 $Y$-态射 $r:X\to\Pb(\E)$ 的集合与二元组 $(\LL,\varphi)$ 的等价类集合之间建立一一对应，其中 $\LL$ 是可逆 $\OO_X$-模，$\varphi:q^*(\E)\to\LL$ 为满同态。这里二元组 $(\LL,\varphi)$ 与 $(\LL',\varphi')$ 等价，是指存在 $\OO_X$-模同构 $\tau:\LL\xrightarrow{\sim}\LL'$，使 $\varphi'=\tau\circ\varphi$。
\begin{proof}
先从 $Y$-态射 $r:X\to\Pb(\E)$ 出发，构造 $\LL_r,\varphi_r$（4.2.1），再构造 $r'=r_{\LL_r,\varphi_r}$。由（4.2.1）及（3.7.2），立即可知 $r=r'$；定义（3.7.2）中的 $v_n$ 时，取 $(f/1)\otimes1$ 为 $\LL_r$ 的局部生成元即可。

反之，从 $(\LL,\varphi)$ 出发，构造 $r''=r_{\LL,\varphi}$，再构造 $\LL_{r''},\varphi_{r''}$。下面证明存在典范同构 $\tau:\LL_{r''}\xrightarrow{\sim}\LL$，满足 $\varphi=\tau\circ\varphi_{r''}$。为定义它，可设 $Y=\Spec(A)$、$X=\Spec(B)$ 均仿射。采用（4.2.1）及（3.7.2）的记号，把 $L_{r''}$ 中的每个 $(x/1)\otimes1$（$x\in E$）映为 $L$ 中的 $v_1(x)c$。立即可验证 $\tau$ 不依赖 $L$ 的生成元 $c$ 的选择。由假设 $\varphi$ 满射，为证明 $\tau$ 是同构，只须验证：若 $x/1=0$ 在 $(S(1))_{(f)}$ 中成立，则 $v_1(x)/1=0$ 在 $B_g$ 中成立。前一关系表示存在 $n$ 使 $f^nx=0$ 在 $S_{n+1}=\Sym_A^{n+1}(E)$ 中成立，故
$v_{n+1}(f^nx)=g^nv_1(x)=0$ 在 $B$ 中成立，结论由此得出。最后，两个等价二元组 $(\LL,\varphi)$、$(\LL',\varphi')$ 显然给出相同态射 $r_{\LL,\varphi}=r_{\LL',\varphi'}$。
\end{proof}

特别地，取 $X=Y$，得到：
\egastatement{4.2.4}{定理}
$\Pb(\E)$ 的 $Y$-截面集合，与 $\E$ 的满足 $\E/\F$ 可逆的拟凝聚 $\OO_Y$-子模 $\F$ 所成的集合之间，存在典范一一对应。

注意，这个性质对应于把“射影空间”定义为向量空间的超平面集合的经典定义。经典情形为 $Y=\Spec(K)$，$K$ 为域，$\E=\wt E$，其中 $E$ 是有限维 $K$-向量空间；此时具有（4.2.4）所述性质的 $\F$ 对应于 $E$ 的超平面，而 $\Pb(\E)$ 的 $Y$-截面就是 $\Pb(\E)$ 的 $K$-有理点（I，3.4.5）。

\egastatement{4.2.5}{注记}
从 $X$ 到 $P$ 的 $Y$-态射与它们的图态射之间有典范一一对应；后者为 $P\times_Y X$ 的 $X$-截面（I，3.3.14）。因此反过来，（4.2.3）也可由（4.2.4）推出。以 $\operatorname{Hyp}_Y(X,\E)$ 表示 $\E\otimes_Y\OO_X=q^*(\E)$ 的拟凝聚 $\OO_X$-子模 $\F$ 所成的集合，要求 $q^*(\E)/\F$ 为可逆 $\OO_X$-模。若 $g:X'\to X$ 为 $Y$-态射，则由 $g^*$ 右正合，有
\[
g^*(q^*(\E)/\F)=g^*(q^*(\E))/\Im(g^*(\F)),
\]
所以这个商层可逆。\transnote{原文把分母简写为 $g^*(\F)$；此处显式写出其在 $g^*(q^*(\E))$ 中的典范像，未断言逆像函子保持子模的单射。}
因此 $\operatorname{Hyp}_Y(X,\E)$ 是 $Y$-预概形范畴上的反变函子。定理（4.2.4）可解释为：当 $X$ 遍历 $Y$-预概形时，反变函子 $\Hom_Y(X,\Pb(\E))$ 与 $\operatorname{Hyp}_Y(X,\E)$ 之间有典范同构。

这也通过下述泛性质刻画射影丛 $P=\Pb(\E)$，比第1至第3节的构造更接近几何直观：对每个态射 $q:X\to Y$ 及每个作为 $\E\otimes_Y\OO_X$ 的商的可逆 $\OO_X$-模 $\LL$，存在唯一 $Y$-态射 $r:X\to P$，使 $\LL=r^*(\OO_P(1))$。\transnote{这里“作为商”包括给定商同态，等同也与相应商同态相容，正如（4.2.3）；只给出可逆层的同构类不足以唯一确定 $r$。}
以后将看到如何同样定义“格拉斯曼”概形等对象。

\egastatement{4.2.6}{推论}
假设每个可逆 $\OO_Y$-模都平凡（I，2.4.8）。设 $V=\Hom_{\OO_Y}(\E,\OO_Y)$，把它看成环 $A=\Gamma(Y,\OO_Y)$ 上的模；设 $V^*$ 是其中满同态组成的子集。则 $\Pb(\E)$ 的 $Y$-截面集合典范地等同于 $V^*/A^*$，其中 $A^*$ 是 $A$ 的单位群。

特别地：

1° 当 $Y$ 是局部概形时，（4.2.6）适用（I，2.4.8）。设 $Y$ 是任意预概形，$y\in Y$，$Y'=\Spec(k(y))$。由（4.1.3.1），$\Pb(\E)$ 的纤维 $p^{-1}(y)$ 等同于 $\Pb(\E^y)$，其中
$\E^y=\E_y\otimes_{\OO_y}k(y)=\E_y/\mm_y\E_y$ 看成 $k(y)$-向量空间。更一般地，若 $K$ 为 $k(y)$ 的扩域，则 $p^{-1}(y)\otimes_{k(y)}K$ 等同于 $\Pb(\E^y\otimes_{k(y)}K)$。因此，（4.2.6）说明：$\Pb(\E)$ 在 $k(y)$ 的扩域 $K$ 中取值的几何点集合（I，3.4.5），也可称为 $\Pb(\E)$ 在 $y$ 上方的 $K$-有理几何纤维，等同于 $K$-向量空间 $\E^y\otimes_{k(y)}K$ 的对偶空间所对应的射影空间。

2° 假设 $Y$ 仿射，环为 $A$，且每个可逆 $\OO_Y$-模都平凡；再取 $\E=\OO_Y^n$。则（4.2.6）中的 $V$ 等同于 $A^n$（I，1.3.8），$V^*$ 等同于生成单位理想 $A$ 的元素组 $(f_i)_{1\leq i\leq n}$ 所成的集合。两个这样的元素组定义 $\Pb_Y^{n-1}=\Pb_A^{n-1}$ 的同一 $Y$-截面，也就是说定义同一 $A$-值点，当且仅当其中一个可由另一个乘以 $A$ 的一个单位得到。

这些性质说明把 $\Pb(\E)$ 称为“射影丛”是适当的。须注意，如此得到的“射影空间”定义实际上与经典定义对偶；这是因为我们必须能够对任意拟凝聚 $\OO_Y$-模 $\E$ 定义 $\Pb(\E)$，并不要求 $\E$ 局部自由。

\egastatement{4.2.7}{注记}
将在第五章看到：若 $Y$ 局部诺特且连通，$\E$ 局部自由，则置 $P=\Pb(\E)$ 后，每个可逆 $\OO_P$-模 $\LL$ 都同构于形如 $\LL'\otimes_{\OO_Y}\OO_P(m)$ 的 $\OO_P$-模，其中 $\LL'$ 是在同构意义下确定的可逆 $\OO_Y$-模，$m$ 是确定的整数。换言之，$H^1(P,\OO_P^*)$ 同构于 $\Z\times H^1(Y,\OO_Y^*)$（0，5.4.7）。\transnote{这里按底本保留原述；秩为零或一等退化情形并不满足上述唯一整数的断言。例如 $\E=\OO_Y$ 时 $P=Y$ 且 $\OO_P(1)$ 平凡。}

还将看到（III，2.1.14，并使用0，5.4.10）：当 $m<0$ 时 $p_*(\LL^{\otimes m})=0$，当 $m\geq0$ 时 $p_*(\LL^{\otimes m})$ 同构于 $\LL'\otimes_{\OO_Y}(\Sym_{\OO_Y}(\E))_m$。\transnote{底本在这一句确实写作 $p_*(\LL^{\otimes m})$；但上一句已把 $\LL$ 写成 $\LL'\otimes\OO_P(m)$。通常所需的直像公式涉及 $p_*(\LL)$，而非再次取 $m$ 次张量幂。本译保留原式并标出这个记号问题。}
因此，对拟凝聚 $\OO_Y$-模 $\F$，每个 $Y$-态射 $\Pb(\E)\to\Pb(\F)$ 都由一个可逆 $\OO_Y$-模 $\LL'$、整数 $m\geq0$ 和 $\OO_Y$-同态
\[
\psi:\F\longrightarrow\LL'\otimes_{\OO_Y}(\Sym_{\OO_Y}(\E))_m
\]
确定，要求相应的层同态 $\psi^\sharp$ 满射。还将看到：若所考虑的 $Y$-态射为同构，则 $m=1$，且 $\F\simeq\E\otimes_{\OO_Y}\LL'$，这就是（4.1.4）的逆命题。由此可以把 $\Pb(\E)$ 的自同构芽层确定为群层 $\mathscr{A}ut(\E)$ 模去 $\OO_Y^*$ 的商；若 $\E$ 的秩为 $n$，则 $\mathscr{A}ut(\E)$ 局部同构于 $\operatorname{GL}(n,\OO_Y)$。

\egapar{4.2.8}
保留（4.2.1）的记号，设 $u:X'\to X$ 为态射。若 $Y$-态射 $r:X\to P$ 对应于同态 $\varphi:q^*(\E)\to\LL$，则 $Y$-态射 $r\circ u$ 对应于
$u^*(\varphi):u^*(q^*(\E))\to u^*(\LL)$；这立即由定义得出。

\egapar{4.2.9}
设 $\E,\F$ 是拟凝聚 $\OO_Y$-模，$v:\E\to\F$ 为满同态，$j=\Pb(v):\Pb(\F)\to\Pb(\E)$ 为相应闭浸入（4.1.2）。若 $Y$-态射 $r:X\to\Pb(\F)$ 对应于同态 $\varphi:q^*(\F)\to\LL$，则 $j\circ r$ 对应于
$q^*(\E)\xrightarrow{q^*(v)}q^*(\F)\xrightarrow{\varphi}\LL$；这仍由（4.2.1）的定义得出。

\egapar{4.2.10}
设 $\psi:Y'\to Y$ 为态射，并置 $\E'=\psi^*(\E)$。若 $Y$-态射 $r:X\to P$ 对应于同态 $\varphi:q^*(\E)\to\LL$，则 $Y'$-态射
$r_{(Y')}:X_{(Y')}\to P'=\Pb(\E')$ 对应于
\[
\varphi_{(Y')}:q_{(Y')}^*(\E')
\simeq q^*(\E)\otimes_Y\OO_{Y'}\longrightarrow\LL\otimes_Y\OO_{Y'}.
\]
事实上，由（4.1.3.1），有交换图
\[
\begin{tikzcd}[column sep=large]
Y' \arrow[d,"\psi"'] & P' \arrow[l,"p_{(Y')}"'] \arrow[d,"\psi_{(P)}"] & X_{(Y')} \arrow[l,"r_{(Y')}"'] \arrow[d,"\psi_{(X)}"]\\
Y & P \arrow[l,"p"] & X \arrow[l,"r"]
\end{tikzcd}
\]
由（4.1.3.2），有
\[
\begin{aligned}
r_{(Y')}^*(\OO_{P'}(1))
&=r_{(Y')}^*(\psi_{(P)}^*(\OO_P(1)))\\
&=\psi_{(X)}^*(r^*(\OO_P(1)))
=\psi_{(X)}^*(\LL)=\LL\otimes_Y\OO_{Y'}.
\end{aligned}
\]
另一方面，$\psi_{(P)}^*(\alpha_1^\sharp)$ 正是典范同态
$\alpha_1'^\sharp:p_{(Y')}^*(\E')\to\OO_{P'}(1)$；像（4.1.6）那样具体写出 $P,P'$ 中的典范同态即可看出。所述结论由此成立。

\subsection{Segre 态射}
\src{76--78}
\egapar{4.3.1}
设 $Y$ 为预概形，$\E,\F$ 为拟凝聚 $\OO_Y$-模，置 $P_1=\Pb(\E)$、$P_2=\Pb(\F)$，并以 $p_1:P_1\to Y$、$p_2:P_2\to Y$ 表示结构态射。令 $Q=P_1\times_Y P_2$，$q_1:Q\to P_1$、$q_2:Q\to P_2$ 为典范投影。则 $\OO_Q$-模
\[
\LL=\OO_{P_1}(1)\otimes_Y\OO_{P_2}(1)
=q_1^*(\OO_{P_1}(1))\otimes_{\OO_Q}q_2^*(\OO_{P_2}(1))
\]
是可逆的，因为它是两个可逆 $\OO_Q$-模的张量积（0，5.4.4）。若 $r=p_1\circ q_1=p_2\circ q_2:Q\to Y$ 为结构态射，则
$r^*(\E\otimes_{\OO_Y}\F)=q_1^*(p_1^*(\E))\otimes_{\OO_Q}q_2^*(p_2^*(\F))$（0，4.3.3）。因此把典范满同态
$p_1^*(\E)\to\OO_{P_1}(1)$、$p_2^*(\F)\to\OO_{P_2}(1)$（4.1.5.1）作张量积，得到典范满同态
\begin{equation}\tag{4.3.1.1}\varepsilon:r^*(\E\otimes_{\OO_Y}\F)\longrightarrow\LL.\end{equation}
从而得到典范态射（4.2.2），称为\textbf{Segre 态射}：
\begin{equation}\tag{4.3.1.2}\zeta:\Pb(\E)\times_Y\Pb(\F)\longrightarrow\Pb(\E\otimes_{\OO_Y}\F).\end{equation}

在 $Y=\Spec(A)$ 仿射时，把 $\zeta$ 具体写出。设 $\E=\wt E$、$\F=\wt F$，其中 $E,F$ 是 $A$-模，因此
$\E\otimes_{\OO_Y}\F=(E\otimes_A F)^\sim$（I，1.3.12）。置
$R=\Sym_A(E)$、$S=\Sym_A(F)$、$T=\Sym_A(E\otimes_A F)$。
取 $f\in E,g\in F$，考虑 $Q$ 的仿射开集
\[
D_+(f)\times_Y D_+(g)=\Spec(B),\qquad B=R_{(f)}\otimes_A S_{(g)}.
\]
$\LL$ 在此开集上的限制为 $\wt L$，其中
$L=(R(1))_{(f)}\otimes_A(S(1))_{(g)}$。把 $L$ 看成自由 $B$-模时，$c=(f/1)\otimes(g/1)$ 是其生成元（2.5.7）。同态（4.3.1.1）对应于从 $(E\otimes_A F)\otimes_A B$ 到 $L$ 的同态
\[
(x\otimes y)\otimes b\longmapsto b((x/1)\otimes(y/1)).
\]
采用（3.7.2）的记号，即有 $v_1(x\otimes y)=(x/f)\otimes(y/g)$。因此 $\zeta$ 在 $D_+(f)\times_YD_+(g)$ 上的限制是到 $D_+(f\otimes g)$ 的态射，对应于环同态
$\omega:T_{(f\otimes g)}\to R_{(f)}\otimes_A S_{(g)}$，满足
\begin{equation}\tag{4.3.1.3}
\omega\left(\frac{x\otimes y}{f\otimes g}\right)=\frac{x}{f}\otimes\frac{y}{g}
\qquad(x\in E,\ y\in F).
\end{equation}

\egapar{4.3.2}
置 $P=\Pb(\E\otimes_{\OO_Y}\F)$。由（4.2.3），存在典范同构
\begin{equation}\tag{4.3.2.1}
\tau:\zeta^*(\OO_P(1))\xrightarrow{\sim}\OO_{P_1}(1)\otimes_Y\OO_{P_2}(1).
\end{equation}
下面证明：对 $x\in\Gamma(Y,\E)$、$y\in\Gamma(Y,\F)$，有
\begin{equation}\tag{4.3.2.2}\tau(\alpha_1(x\otimes y))=\alpha_1(x)\otimes\alpha_1(y).\end{equation}
事实上，可归结到 $Y$ 仿射的情形。采用（4.3.1）和（2.6.2）的记号，$\alpha_1(x\otimes y)$ 是 $(T(1))_{(f\otimes g)}$ 中的 $(x\otimes y)/1$，$\alpha_1(x)$ 是 $(R(1))_{(f)}$ 中的 $x/1$，$\alpha_1(y)$ 是 $(S(1))_{(g)}$ 中的 $y/1$。（4.2.3）中 $\tau$ 的定义及（4.3.1）对 $\omega$ 的计算立即给出（4.3.2.2）。

由此推出公式
\begin{equation}\tag{4.3.2.3}\zeta^{-1}(P_{x\otimes y})=(P_1)_x\times_Y(P_2)_y,\end{equation}
记号如（3.1.4）。事实上，结合（3.3.3），公式（4.3.2.2）把问题归结到下面的引理；通过（I，3.2.7及3.2.3）回到仿射情形即可。

\egastatement{4.3.2.4}{引理}
设 $B,B'$ 是 $A$-代数，$Y=\Spec(A)$、$Z=\Spec(B)$、$Z'=\Spec(B')$。对任意 $t\in B,t'\in B'$，有
$D(t\otimes t')=D(t)\times_Y D(t')$。
\begin{proof}
若 $p:Z\times_Y Z'\to Z$、$p':Z\times_Y Z'\to Z'$ 为典范投影，则由（I，1.2.2.2）有
$p^{-1}(D(t))=D(t\otimes1)$、$p'^{-1}(D(t'))=D(1\otimes t')$。
因为 $(t\otimes1)(1\otimes t')=t\otimes t'$，结论由（I，3.2.7）及（I，1.1.9.1）得出。
\end{proof}

\egastatement{4.3.3}{命题}
Segre 态射是闭浸入。
\begin{proof}
问题在 $Y$ 上是局部的，可设 $Y$ 仿射。采用（4.3.1）及（4.3.2）的记号，诸 $D_+(f\otimes g)$ 构成 $P$ 的拓扑基，因为 $f,g$ 分别遍历 $E,F$ 时，诸 $f\otimes g$ 生成 $T$。\transnote{底本在这里说“拓扑基”；由这些一次生成元直接得到的是开覆盖，后面的局部验证只需要开覆盖。}
另一方面，由（4.3.2.3），$\zeta^{-1}(D_+(f\otimes g))=D_+(f)\times_Y D_+(g)$。故只须证明 $\zeta$ 在该逆像上的限制是到 $D_+(f\otimes g)$ 的闭浸入（I，4.2.4）。公式（4.3.1.3）说明相应的 $\omega$ 满射，证明完毕。
\end{proof}

\egapar{4.3.4}
当把拟凝聚 $\OO_Y$-模的同态限制为满同态时，Segre 态射关于 $\E,\F$ 具有函子性。事实上，须证明：若 $\E\to\E'$ 为满 $\OO_Y$-同态，则图
\[
\begin{tikzcd}[column sep=large]
\Pb(\E')\times_Y\Pb(\F) \arrow[r,"j\times1"] \arrow[d,"\zeta"'] & \Pb(\E)\times_Y\Pb(\F) \arrow[d,"\zeta"]\\
\Pb(\E'\otimes\F) \arrow[r] & \Pb(\E\otimes\F)
\end{tikzcd}
\]
交换，其中 $j:\Pb(\E')\to\Pb(\E)$ 为典范闭浸入。置 $P'_1=\Pb(\E')$，其余记号仍如（4.3.1）。$j\times1$ 为闭浸入（I，4.3.1）；由（4.1.2.1）和（I，9.1.5），在同构意义下有
\[
\begin{aligned}
(j\times1)^*(\OO_{P_1}(1)\otimes_Y\OO_{P_2}(1))
&=j^*(\OO_{P_1}(1))\otimes_Y\OO_{P_2}(1)\\
&=\OO_{P'_1}(1)\otimes_Y\OO_{P_2}(1).
\end{aligned}
\]
因此由（4.2.8）和（4.2.9）得到所述断言。

\egapar{4.3.5}
采用（4.3.1）的记号，设 $\psi:Y'\to Y$ 为态射，并置 $\E'=\psi^*(\E)$、$\F'=\psi^*(\F)$。则 Segre 态射
$\Pb(\E')\times_{Y'}\Pb(\F')\to\Pb(\E'\otimes\F')$ 等同于 $\zeta_{(Y')}$。

事实上，另置 $P'_1=\Pb(\E')$、$P'_2=\Pb(\F')$。由（4.1.3.1），$P'_i$ 等同于 $(P_i)_{(Y')}$（$i=1,2$），故结构态射 $P'_1\times_{Y'}P'_2\to Y'$ 等同于 $r_{(Y')}$。又 $\E'\otimes\F'$ 等同于 $\psi^*(\E\otimes\F)$，故 $\Pb(\E'\otimes\F')$ 等同于 $(\Pb(\E\otimes\F))_{(Y')}$（4.1.3.1）。最后，由（4.1.3.2）及（I，9.1.11），$\LL'=\OO_{P'_1}(1)\otimes_{Y'}\OO_{P'_2}(1)$ 等同于 $\LL\otimes_Y\OO_{Y'}$。典范同态
$r_{(Y')}^*(\E'\otimes\F')\to\LL'$ 于是等同于 $\varepsilon_{(Y')}$；所述断言由（4.2.10）得出。

\egastatement{4.3.6}{注记}
$\Pb(\E)$ 与 $\Pb(\F)$ 的和预概形同样典范地同构于 $\Pb(\E\oplus\F)$ 的一个闭子预概形。事实上，满同态 $\E\oplus\F\to\E$、$\E\oplus\F\to\F$ 对应于闭浸入
$\Pb(\E)\to\Pb(\E\oplus\F)$、$\Pb(\F)\to\Pb(\E\oplus\F)$。
只须说明如此得到的两个闭子预概形的底空间不交。问题在 $Y$ 上是局部的，可采用（4.3.1）的记号。$\Sym_A^n(E)$ 和 $\Sym_A^n(F)$ 等同于 $\Sym_A^n(E\oplus F)$ 中交为零的两个子模；若 $\mathfrak p$ 是 $\Sym(E)$ 的分次素理想，对每个 $n>0$ 有 $\mathfrak p\cap\Sym_n(E)\ne\Sym_n(E)$，则它在 $\Sym(E\oplus F)$ 中对应的分次素理想，在 $\Sym_n(E)$ 上的迹为 $\mathfrak p\cap\Sym_n(E)$，同时包含 $\Sym_n(F)$。因此分别来自 $\Proj(\Sym(E))$ 和 $\Proj(\Sym(F))$ 的两点，不可能在 $\Proj(\Sym(E\oplus F))$ 中有相同的像。

\subsection{到射影丛的浸入。甚丰沛层}
\src{78--83}
\egastatement{4.4.1}{命题}
设 $Y$ 是拟紧概形，或底空间为诺特空间的预概形；$q:X\to Y$ 是有限型态射，$\LL$ 是可逆 $\OO_X$-模。
\begin{enumerate}[label=(\roman*)]
\item 设 $\SSS$ 是非负分次拟凝聚 $\OO_Y$-代数，$\psi:q^*(\SSS)\to\bigoplus_{n\geq0}\LL^{\otimes n}$ 为分次代数同态。则 $r_{\LL,\psi}$ 处处有定义且为浸入，当且仅当存在整数 $n>0$ 和 $\SSS_n$ 的有限型拟凝聚 $\OO_Y$-子模 $\E$，使
$\psi'=\psi_n\circ q^*(j):q^*(\E)\to\LL^{\otimes n}=\LL'$ 满射，且态射 $r_{\LL',\psi'}:X\to\Pb(\E)$ 为浸入；这里 $j:\E\to\SSS_n$ 是嵌入。
\item 设 $\F$ 是拟凝聚 $\OO_Y$-模，$\varphi:q^*(\F)\to\LL$ 为满同态。则 $r_{\LL,\varphi}:X\to\Pb(\F)$ 为浸入，当且仅当存在 $\F$ 的有限型拟凝聚 $\OO_Y$-子模 $\E$，使 $\varphi'=\varphi\circ q^*(j):q^*(\E)\to\LL$ 满射，且 $r_{\LL,\varphi'}:X\to\Pb(\E)$ 为浸入；这里 $j:\E\to\F$ 为典范嵌入。
\end{enumerate}
\begin{proof}
（i）由（3.8.5），$r_{\LL,\psi}$ 处处有定义且为浸入，等价于存在 $n>0$ 和 $\E$，使以 $\SSS'$ 表示 $\E$ 生成的子代数时，同态 $q^*(\E)\to\LL^{\otimes n}$ 满射，且相应的 $r_{\LL,\psi'}:X\to\Proj(\SSS')$ 处处有定义并为浸入。另一方面，有典范满同态 $\Sym_{\OO_Y}(\E)\to\SSS'$，它对应于闭浸入 $\Proj(\SSS')\to\Pb(\E)$（3.6.2），故结论成立。

（ii）$\F$ 是其有限型拟凝聚子模 $\E_\lambda$ 的归纳极限（I，9.4.9），所以 $\Sym(\F)$ 是诸 $\Sym(\E_\lambda)$ 的归纳极限。结论由（3.8.4）得出，并注意在该命题的证明中可以把全部 $n_i$ 取为一：事实上，设 $Y$ 仿射，且 $r=r_{\LL,\varphi}$ 为浸入，则 $r(X)$ 是 $\Pb(\F)$ 的拟紧子空间，可用有限多个形如 $D_+(f)$、$f\in F$ 的开集覆盖，使各 $D_+(f)\cap r(X)$ 为闭集。
\end{proof}

\egastatement{4.4.2}{定义}
设 $Y$ 是预概形，$q:X\to Y$ 是态射。若存在拟凝聚 $\OO_Y$-模 $\E$ 和 $Y$-浸入 $i:X\to P=\Pb(\E)$，使可逆 $\OO_X$-模 $\LL$ 同构于 $i^*(\OO_P(1))$，则称 $\LL$ \textbf{相对于 $q$ 甚丰沛}，或\textbf{相对于 $Y$ 甚丰沛}；不会混淆时简称\textbf{甚丰沛}（très ample）。

由（4.2.3），这等价于存在拟凝聚 $\OO_Y$-模 $\E$ 和满同态 $\varphi:q^*(\E)\to\LL$，使 $r_{\LL,\varphi}:X\to\Pb(\E)$ 为浸入。注意，存在相对于 $Y$ 甚丰沛的 $\OO_X$-模，蕴含 $q$ 分离（3.1.3及I，5.5.1，i、ii）。

\egastatement{4.4.3}{推论}
假设存在由 $\SSS_1$ 生成的拟凝聚分次 $\OO_Y$-代数 $\SSS$，以及 $Y$-浸入 $i:X\to P=\Proj(\SSS)$，使 $\LL\simeq i^*(\OO_P(1))$。则 $\LL$ 相对于 $q$ 甚丰沛。
\begin{proof}
令 $\F=\SSS_1$，则典范同态 $\Sym(\F)\to\SSS$ 满射。将相应闭浸入 $\Proj(\SSS)\to\Pb(\F)$（3.6.2）与 $i$ 复合，得到浸入 $j:X\to\Pb(\F)=P'$，使 $\LL\simeq j^*(\OO_{P'}(1))$（3.6.3）。
\end{proof}

\egastatement{4.4.4}{命题}
设 $q:X\to Y$ 是拟紧态射，$\LL$ 是可逆 $\OO_X$-模。下列性质等价：
\begin{enumerate}[label=\alph*)]
\item $\LL$ 相对于 $q$ 甚丰沛。
\item $q_*(\LL)$ 拟凝聚，典范同态 $\sigma:q^*(q_*(\LL))\to\LL$ 满射，且 $r_{\LL,\sigma}:X\to\Pb(q_*(\LL))$ 是浸入。
\end{enumerate}
\begin{proof}
由于 $q$ 拟紧，当 $q$ 分离时，$q_*(\LL)$ 拟凝聚（I，9.2.2）。已知存在满同态 $\varphi:q^*(\E)\to\LL$、其中 $\E$ 拟凝聚，便推出 $\sigma$ 满射（3.4.7）。而 $\varphi$ 的分解
$q^*(\E)\to q^*(q_*(\LL))\to\LL$ 典范地对应于分解
\[
q^*(\Sym(\E))\longrightarrow q^*(\Sym(q_*(\LL)))
\longrightarrow\bigoplus_{n\geq0}\LL^{\otimes n}.
\]
因此由（3.8.3），若 $r_{\LL,\varphi}$ 是浸入，则 $j=r_{\LL,\sigma}$ 也是浸入。置 $P'=\Pb(q_*(\LL))$，又有 $\LL\simeq j^*(\OO_{P'}(1))$（4.2.4）。由此 a）、b）等价。
\end{proof}

\egastatement{4.4.5}{推论}
假设 $q$ 拟紧。则 $\LL$ 相对于 $Y$ 甚丰沛，当且仅当存在 $Y$ 的开覆盖 $(U_\alpha)$，使每个 $\LL|q^{-1}(U_\alpha)$ 相对于 $U_\alpha$ 甚丰沛。
\begin{proof}
（4.4.4）的条件 b）是 $Y$ 上的局部条件。
\end{proof}

\egastatement{4.4.6}{命题}
设 $Y$ 是拟紧概形，或底空间为诺特空间的预概形；$q:X\to Y$ 是有限型态射，$\LL$ 是可逆 $\OO_X$-模。则（4.4.4）的条件 a）、b）还等价于下列条件：
\begin{enumerate}[label=\alph*'),start=1]
\item 存在有限型拟凝聚 $\OO_Y$-模 $\E$ 及满同态 $\varphi:q^*(\E)\to\LL$，使 $r_{\LL,\varphi}$ 为浸入。
\item $q_*(\LL)$ 中存在具有 a'）所述性质的有限型凝聚 $\OO_Y$-子模 $\E$。\transnote{条件 b'）在底本中用的是“凝聚”；在本命题所允许的非诺特情形，应注意这比“有限型拟凝聚”更强。这里依原文保留。}
\end{enumerate}
\begin{proof}
a'）或 b'）显然推出 a）。另一方面，a）由（4.4.1）推出 a'），同理 b）推出 b'）。
\end{proof}

\egastatement{4.4.7}{推论}
假设 $Y$ 是拟紧概形或诺特预概形。若 $\LL$ 相对于 $q$ 甚丰沛，则存在拟凝聚分次 $\OO_Y$-代数 $\SSS$，使 $\SSS_1$ 为有限型并生成 $\SSS$，以及占优的开 $Y$-浸入 $i:X\to P=\Proj(\SSS)$，使 $\LL\simeq i^*(\OO_P(1))$。
\begin{proof}
（4.4.6）的条件 b'）成立；结构态射 $p:P'=\Pb(\E)\to Y$ 分离且为有限型（3.1.3）。所以当 $Y$ 为拟紧概形（分别地，诺特预概形）时，$P'$ 也是拟紧概形（分别地，诺特预概形）。设 $Z$ 是与浸入 $j=r_{\LL,\varphi}:X\to P'$ 相伴的子预概形 $X'$ 在 $P'$ 中的闭包（I，9.5.11）。显然，$j$ 分解为占优开浸入 $i:X\to Z$ 与典范嵌入 $Z\to P'$ 的复合。由（3.6.2），$Z$ 等同于某个 $\Proj(\SSS)$，其中 $\SSS$ 是 $\Sym(\E)$ 模去一个拟凝聚分次理想层的商。显然 $\SSS_1$ 为有限型并生成 $\SSS$；又 $\OO_Z(1)$ 是 $\OO_{P'}(1)$ 在典范嵌入下的逆像（3.6.3），故 $\LL=i^*(\OO_Z(1))$。
\end{proof}

\egastatement{4.4.8}{命题}
设 $q:X\to Y$ 是态射，$\LL$ 是相对于 $q$ 甚丰沛的 $\OO_X$-模，$\LL'$ 为可逆 $\OO_X$-模，并假设存在拟凝聚 $\OO_Y$-模 $\E'$ 及满同态 $q^*(\E')\to\LL'$。则 $\LL\otimes_{\OO_X}\LL'$ 相对于 $q$ 甚丰沛。
\begin{proof}
由假设，存在 $Y$-态射 $r':X\to P'=\Pb(\E')$，使 $\LL'=r'^*(\OO_{P'}(1))$（4.2.1）。另由假设，存在拟凝聚 $\OO_Y$-模 $\E$ 和 $Y$-浸入 $r:X\to P=\Pb(\E)$，使 $\LL=r^*(\OO_P(1))$。置 $Q=\Pb(\E\otimes\E')$，考虑 Segre 态射 $\zeta:P\times_Y P'\to Q$（4.3.1）。由于 $r$ 为浸入，$(r,r')_Y:X\to P\times_Y P'$ 也为浸入（I，5.3.14）；而 $\zeta$ 是浸入（4.3.3），故复合 $r'':X\to P\times_Y P'\to Q$ 为浸入。由（4.3.2.1），$\zeta^*(\OO_Q(1))\simeq\OO_P(1)\otimes_Y\OO_{P'}(1)$，因此由（I，9.1.4），$r''^*(\OO_Q(1))\simeq\LL\otimes\LL'$，命题得证。
\end{proof}

\egastatement{4.4.9}{推论}
设 $q:X\to Y$ 是态射。
\begin{enumerate}[label=(\roman*)]
\item 设 $\K$ 是可逆 $\OO_Y$-模，$\LL$ 是可逆 $\OO_X$-模。则 $\LL$ 相对于 $q$ 甚丰沛，当且仅当 $\LL\otimes q^*(\K)$ 相对于 $q$ 甚丰沛。
\item 若 $\LL,\LL'$ 都相对于 $q$ 甚丰沛，则 $\LL\otimes\LL'$ 也如此。特别地，对每个 $n>0$，$\LL^{\otimes n}$ 相对于 $q$ 甚丰沛。
\end{enumerate}
\begin{proof}
（ii）及（i）的必要性立即由（4.4.8）得出。反之，若 $\LL\otimes q^*(\K)$ 甚丰沛，则由已证部分，$(\LL\otimes q^*(\K))\otimes q^*(\K^{-1})$ 也甚丰沛；而它同构于 $\LL$（0，5.4.3及5.4.5）。
\end{proof}

\egastatement{4.4.10}{命题}
\begin{enumerate}[label=(\roman*)]
\item 对每个预概形 $Y$，每个可逆 $\OO_Y$-模 $\LL$ 都相对于恒等态射 $1_Y$ 甚丰沛。
\end{enumerate}
\noindent（i bis）设 $f:X\to Y$ 是态射，$j:X'\to X$ 是浸入。若 $\LL$ 是相对于 $f$ 甚丰沛的 $\OO_X$-模，则 $j^*(\LL)$ 相对于 $f\circ j$ 甚丰沛。
\begin{enumerate}[label=(\roman*),start=2]
\item 设 $Z$ 是拟紧预概形，$f:X\to Y$ 是有限型态射，$g:Y\to Z$ 是拟紧态射；$\LL$ 为相对于 $f$ 甚丰沛的 $\OO_X$-模，$\K$ 为相对于 $g$ 甚丰沛的 $\OO_Y$-模。则存在整数 $n_0>0$，使对每个 $n\geq n_0$，$\LL\otimes f^*(\K^{\otimes n})$ 相对于 $g\circ f$ 甚丰沛。
\item 设 $f:X\to Y$、$g:Y'\to Y$ 是态射，$X'=X_{(Y')}$。若 $\LL$ 相对于 $f$ 甚丰沛，则 $\LL'=\LL\otimes_Y\OO_{Y'}$ 是相对于 $f_{(Y')}$ 甚丰沛的 $\OO_{X'}$-模。
\item 设 $f_i:X_i\to Y_i$（$i=1,2$）为 $S$-态射。若 $\LL_i$ 为相对于 $f_i$ 甚丰沛的 $\OO_{X_i}$-模，则 $\LL_1\otimes_S\LL_2$ 相对于 $f_1\times_S f_2$ 甚丰沛。
\item 设 $f:X\to Y$、$g:Y\to Z$ 是态射。若 $\OO_X$-模 $\LL$ 相对于 $g\circ f$ 甚丰沛，则 $\LL$ 相对于 $f$ 甚丰沛。
\item 设 $f:X\to Y$ 是态射，$j:X_{\red}\to X$ 为典范嵌入。若 $\LL$ 相对于 $f$ 甚丰沛，则 $j^*(\LL)$ 相对于 $f_{\red}$ 甚丰沛。
\end{enumerate}
\begin{proof}
（i bis）立即由定义（4.4.2）得出。像（I，5.5.12）那样推理，即可由（i bis）和（v）形式地推出（vi），细节留给读者。为证明（v），与（I，5.5.12）相同，考虑分解
$X\xrightarrow{\Gamma_f}X\times_Z Y\xrightarrow{p_2}Y$，注意 $p_2=(g\circ f)\times1_Y$。由假设及（i）、（iv），$\LL\otimes_Z\OO_Y$ 相对于 $p_2$ 甚丰沛；又 $\LL=\Gamma_f^*(\LL\otimes_Z\OO_Y)$（I，9.1.4），$\Gamma_f$ 为浸入（I，5.3.11），故可应用（i bis）。

为证明（i），在定义（4.4.2）中取 $\E=\LL$，注意此时 $\Pb(\E)$ 等同于 $Y$（4.1.4）。

证明（iii）。存在拟凝聚 $\OO_Y$-模 $\E$ 和 $Y$-浸入 $i:X\to P=\Pb(\E)$，使 $\LL=i^*(\OO_P(1))$。置 $\E'=g^*(\E)$，它是拟凝聚 $\OO_{Y'}$-模，且 $P'=\Pb(\E')=P_{(Y')}$（4.1.3.1）。$i_{(Y')}$ 为从 $X_{(Y')}$ 到 $P'$ 的浸入（I，4.3.2），而 $\LL'\simeq i_{(Y')}^*(\OO_{P'}(1))$（4.2.10）。

为证明（iv），由假设存在 $Y_i$-浸入 $r_i:X_i\to P_i=\Pb(\E_i)$，其中 $\E_i$ 为拟凝聚 $\OO_{Y_i}$-模，且 $\LL_i=r_i^*(\OO_{P_i}(1))$。$r_1\times_S r_2$ 是从 $X_1\times_S X_2$ 到 $P_1\times_S P_2$ 的 $S$-浸入（I，4.3.1）；$\OO_{P_1}(1)\otimes_S\OO_{P_2}(1)$ 在它下面的逆像为 $\LL_1\otimes_S\LL_2$。置 $T=Y_1\times_S Y_2$，以 $p_1,p_2$ 表示到 $Y_1,Y_2$ 的投影。令 $P'_i=\Pb(p_i^*(\E_i))$；由（4.1.3.1），$P'_i=P_i\times_{Y_i}T$，故在典范同构意义下有
\[
\begin{aligned}
P'_1\times_TP'_2
&=(P_1\times_{Y_1}T)\times_T(T\times_{Y_2}P_2)\\
&=P_1\times_{Y_1}(T\times_{Y_2}P_2)\\
&=P_1\times_{Y_1}(Y_1\times_S P_2)=P_1\times_S P_2.
\end{aligned}
\]
同样，$\OO_{P'_i}(1)=\OO_{P_i}(1)\otimes_{Y_i}\OO_T$；主要依据（I，9.1.9.1及9.1.2）的类似计算说明，上述等同把
$\OO_{P'_1}(1)\otimes_T\OO_{P'_2}(1)$ 与
$\OO_{P_1}(1)\otimes_S\OO_{P_2}(1)$ 等同。因此可把 $r_1\times_Sr_2$ 看成到 $P'_1\times_TP'_2$ 的 $T$-浸入，上述可逆层在它下面的逆像为 $\LL_1\otimes_S\LL_2$。最后与（4.4.8）相同，使用 Segre 态射即完成证明。

只剩（ii）。先可归结到 $Z$ 仿射的情形。因为一般地有 $Z$ 的有限仿射开覆盖 $(U_i)$；若对
$\K|g^{-1}(U_i)$、$\LL|f^{-1}(g^{-1}(U_i))$ 及整数 $n_i$ 已证得命题，则取 $n_0=\max_i n_i$，由（4.4.5）即得一般情形。假设推出 $f,g$ 分离，故在 $Z$ 仿射时 $X,Y$ 都是拟紧概形。

由（4.4.6），存在浸入 $r:X\to P=\Pb(\E)$，其中 $\E$ 是有限型拟凝聚 $\OO_Y$-模，且 $\LL=r^*(\OO_P(1))$。下面证明存在相对于复合 $P\to Y\to Z$ 甚丰沛的 $\OO_P$-模 $\M$，使 $\OO_P(1)\simeq\M\otimes_Y\K^{\otimes(-m)}$，其中 $m$ 为某整数。于是对 $n\geq m+1$，由假设及应用于 $h:P\to Y$ 和 $1_Y$ 的（iv），$\OO_P(1)\otimes_Y\K^{\otimes n}$ 相对于 $Z$ 甚丰沛。由于 $r$ 为浸入，且
$\LL\otimes f^*(\K^{\otimes n})=r^*(\OO_P(1)\otimes_Y\K^{\otimes n})$，结论将由（i bis）得出。为证明关于 $\OO_P(1)$ 的断言，使用下列引理。\renewcommand{\qedsymbol}{}
\end{proof}

\egastatement{4.4.10.1}{引理}
设 $Z$ 是拟紧概形，或底空间为诺特空间的预概形；$g:Y\to Z$ 是拟紧态射，$\K$ 是相对于 $g$ 甚丰沛的可逆 $\OO_Y$-模，$\E$ 是有限型拟凝聚 $\OO_Y$-模。则存在整数 $m_0$，使对每个 $m\geq m_0$，$\E$ 同构于某个形如 $g^*(\F)\otimes\K^{\otimes(-m)}$ 的 $\OO_Y$-模的商，其中 $\F$ 是有限型拟凝聚 $\OO_Z$-模，可以依赖于 $m$。

此引理将在（4.5.10.1）证明；读者可验证，第4.5小节没有任何地方使用（4.4.10）。

\noindent\textbf{（4.4.10）的证明续。}
承认引理后，存在闭浸入
$j_1:P\to P_1=\Pb(g^*(\F)\otimes\K^{\otimes(-m)})$，使
$\OO_P(1)\simeq j_1^*(\OO_{P_1}(1))$（4.1.2）。另一方面，$P_1$ 与 $P_2=\Pb(g^*(\F))$ 同构，且把 $\OO_{P_1}(1)$ 与 $\OO_{P_2}(1)\otimes_Y\K^{\otimes(-m)}$ 等同（4.1.4）。因此存在闭浸入 $j_2:P\to P_2$，使
$\OO_P(1)\simeq j_2^*(\OO_{P_2}(1))\otimes_Y\K^{\otimes(-m)}$。

最后，$P_2$ 等同于 $P_3\times_Z Y$，其中 $P_3=\Pb(\F)$；而 $\OO_{P_2}(1)$ 等同于 $\OO_{P_3}(1)\otimes_Z\OO_Y$（4.1.3）。按定义，$\OO_{P_3}(1)$ 相对于 $Z$ 甚丰沛；$\K$ 也如此。因此由（iv），$\OO_{P_2}(1)\otimes_Y\K$ 相对于 $Z$ 甚丰沛；由（i bis），
$\M=j_2^*(\OO_{P_2}(1)\otimes_Y\K)$ 也如此。于是
$\OO_P(1)\simeq\M\otimes_Y\K^{\otimes(-m-1)}$，证明完成。\hfill$\square$

\egastatement{4.4.11}{命题}
设 $f:X\to Y$、$f':X'\to Y$ 是态射，$X''$ 是和预概形 $X\amalg X'$，$f'':X''\to Y$ 在 $X$（分别地，$X'$）上的限制为 $f$（分别地，$f'$）。设 $\LL$ 为可逆 $\OO_X$-模，$\LL'$ 为可逆 $\OO_{X'}$-模，$\LL''$ 为在 $X$ 上等于 $\LL$、在 $X'$ 上等于 $\LL'$ 的可逆 $\OO_{X''}$-模。则 $\LL''$ 相对于 $f''$ 甚丰沛，当且仅当 $\LL$ 相对于 $f$ 甚丰沛且 $\LL'$ 相对于 $f'$ 甚丰沛。
\begin{proof}
立即可归结到 $Y$ 仿射的情形。若 $\LL''$ 甚丰沛，则由（4.4.10，i bis），$\LL,\LL'$ 也甚丰沛。反之，若二者甚丰沛，则由定义（4.4.2）和（4.3.6），$\LL''$ 也甚丰沛。
\end{proof}

\subsection{丰沛层}
\src{83--89}

\egapar{4.5.1}
给定预概形 $X$ 和可逆 $\OO_X$-模 $\LL$，对每个 $\OO_X$-模 $\F$，在关于 $\LL$ 不会发生混淆时，记
$\F(n)=\F\otimes_{\OO_X}\LL^{\otimes n}$（$n\in\mathbf Z$）。又记
$S=\bigoplus_{n\geq0}\Gamma(X,\LL^{\otimes n})$；这是（0，5.4.6）所定义的环 $\Gamma_*(\LL)$ 的分次子环。把 $X$ 看成 $\mathbf Z$-预概形，记结构态射为 $p:X\to\Spec(\mathbf Z)$，则分次 $\OO_X$-代数同态
$p^*(\wt S)\to\bigoplus_{n\geq0}\LL^{\otimes n}$ 与分次环 $S$ 的自同态一一对应（I，2.2.5）。与 $S$ 的恒等自同态相对应的同态
$\varepsilon:p^*(\wt S)\to\bigoplus_{n\geq0}\LL^{\otimes n}$ 称为\emph{典范同态}。与它相应的态射 $G(\varepsilon)\to\Proj(S)$（3.7.1）也称为典范态射。

\egastatement{4.5.2}{定理}
设 $X$ 是拟紧概形，或底空间为诺特空间的预概形；$\LL$ 是可逆 $\OO_X$-模，$S=\bigoplus_{n\geq0}\Gamma(X,\LL^{\otimes n})$。下列条件等价：
\begin{enumerate}[label=\alph*)]
\item 当 $f$ 遍历 $S_+$ 的齐次元素时，$X_f$ 构成 $X$ 的拓扑的一组基。
\item[\textup{a')}] 当 $f$ 遍历 $S_+$ 的齐次元素时，其中为仿射的那些 $X_f$ 覆盖 $X$。
\item[\textup{b)}] 典范态射 $G(\varepsilon)\to\Proj(S)$（4.5.1）处处有定义，且是占优开浸入。
\item[\textup{b')}] 典范态射 $G(\varepsilon)\to\Proj(S)$ 处处有定义，且在底空间上给出 $X$ 到 $\Proj(S)$ 的一个子空间的同胚。
\item[\textup{c)}] 对每个拟凝聚 $\OO_X$-模 $\F$，若 $\G_n$ 是 $\F(n)$ 中由其在 $X$ 上的截面生成的子 $\OO_X$-模，则 $\F$ 是所有 $n>0$ 对应的子 $\OO_X$-模 $\G_n(-n)$ 之和。
\item[\textup{c')}] 条件 \textup{c)} 对 $\OO_X$ 的每个拟凝聚理想层成立。
\end{enumerate}
此外，若 $(f_\alpha)$ 是 $S_+$ 的一族齐次元素，且各 $X_{f_\alpha}$ 仿射，则典范态射在 $\bigcup_\alpha X_{f_\alpha}$ 上的限制是
$\bigcup_\alpha X_{f_\alpha}\simeq\bigcup_\alpha(\Proj(S))_{f_\alpha}$ 的同构。
\begin{proof}
显然 b) 推出 b')；由公式（3.7.3.1），并注意 $\varepsilon^\flat$ 为恒等同态，b') 推出 a)。条件 a) 推出 a')：每个 $x\in X$ 都有仿射邻域 $U$，使 $\LL|U\simeq\OO_X|U$；若 $f\in\Gamma(X,\LL^{\otimes n})$ 满足 $x\in X_f\subset U$，则 $X_f$ 也是 $U$ 中满足 $(f|U)(x')\ne0$ 的点 $x'$ 的集合，因而是仿射开集（I，1.3.6）。要证明 a') 推出 b)，只须证明定理的最后一项断言，并验证：若 $X=\bigcup X_{f_\alpha}$，则（3.8.2）的条件（iv）成立。后一点直接由（I，9.3.1，i）得出。至于最后一项断言，由于 $X_f$ 是 $G(\varepsilon)\to\Proj(S)$ 下 $(\Proj(S))_f$ 的逆像，应用（I，9.3.2）即可。因此 a)、a')、b)、b') 等价。

为证明 a') 推出 c)，注意：若 $f\in S_m$ 且 $X_f$ 仿射，则 $\F|X_f$ 由它在 $X_f$ 上的截面生成（I，1.3.9）；另一方面，这样的截面 $s$ 可写成
$(t|X_f)(f|X_f)^{-k}$，其中 $t\in\Gamma(X,\F(km))$（I，9.3.1，ii）。按定义，$t$ 也是 $\G_{km}$ 的截面，故 $s$ 确是 $\G_{km}(-km)$ 在 $X_f$ 上的截面，c) 得证。显然 c) 推出 c')。只剩证明 c') 推出 a)：设 $U$ 是 $x\in X$ 的开邻域，取 $\OO_X$ 的拟凝聚理想层 $\J$，使其定义的闭子预概形的底空间为 $X-U$（I，5.2.1）。由 c')，存在整数 $n>0$ 及 $\J(n)$ 在 $X$ 上的截面 $f$，使 $f(x)\ne0$。显然 $f\in S_n$ 且 $x\in X_f\subset U$，故 a) 成立。
\end{proof}
若 $X$ 的底空间为诺特空间，则（4.5.2）的等价条件推出 $X$ 是概形，因为由（4.5.2，b），它同构于概形 $\Proj(S)$ 的一个子预概形。

\egastatement{4.5.3}{定义}
称可逆 $\OO_X$-模 $\LL$ 是\emph{丰沛的}，若 $X$ 是拟紧概形，且（4.5.2）的等价条件成立。

由判别准则（4.5.2，a）显然可知：若 $\LL$ 是丰沛 $\OO_X$-模，则对 $X$ 的每个开集 $U$，$\LL|U$ 是丰沛的 $(\OO_X|U)$-模。\transnote{原文在这里写“每个开集”。按本条所采用的定义，称限制为“丰沛”还要求 $U$ 拟紧；一般拟紧概形的开集未必拟紧。原句关于拓扑基的性质仍可限制到开集，但这里没有在正文中把原定义替换为不要求拟紧的现代约定。}
（4.5.2）的证明还表明，其中仿射的 $X_f$ 已构成 $X$ 的拓扑的一组基。进一步有：

\egastatement{4.5.4}{推论}
设 $\LL$ 是丰沛 $\OO_X$-模。对每个有限子空间 $Z\subset X$ 及 $Z$ 的每个邻域 $U$，存在整数 $n$ 和 $f\in\Gamma(X,\LL^{\otimes n})$，使 $X_f$ 是包含在 $U$ 中的 $Z$ 的仿射邻域。
\begin{proof}
由（4.5.2，b），只须证明：对 $\Proj(S)$ 的每个有限子集 $Z'$ 及 $Z'$ 的每个开邻域 $U$，存在齐次元素 $f\in S_+$，使 $Z'\subset(\Proj(S))_f\subset U$（2.4.1）。按定义，$\Proj(S)$ 中 $U$ 的补集 $Y$ 为 $V_+(\frj)$，其中 $\frj$ 是不包含 $S_+$ 的分次理想（2.3.2）；$Z'$ 的各点则为不包含 $\frj$ 的分次素理想 $\pp_i$（2.3.1）。存在 $f\in\frj$，不属于任何 $\pp_i$（Bourbaki，\textit{交换代数}，第二章，§1，第1小节，命题2）；由于各 $\pp_i$ 分次，同处的论证甚至表明可令 $f$ 齐次。此元素即满足要求。
\end{proof}

\egastatement{4.5.5}{命题}
设 $X$ 是拟紧概形，或底空间为诺特空间的预概形。（4.5.2）的条件 \textup{a)} 至 \textup{c')} 也与下列各条件等价：
\begin{enumerate}[label=\textup{d)}]
\item 对每个有限型拟凝聚 $\OO_X$-模 $\F$，存在整数 $n_0$，使对所有 $n\geq n_0$，$\F(n)$ 由其在 $X$ 上的截面生成。
\item[\textup{d')}] 对每个有限型拟凝聚 $\OO_X$-模 $\F$，存在两个整数 $n>0$、$k>0$，使 $\F$ 同构于 $\OO_X$-模 $(\LL^{\otimes(-n)})^{\oplus k}$ 的一个商。
\item[\textup{d'')}] 条件 \textup{d')} 对 $\OO_X$ 的每个有限型拟凝聚理想层成立。
\end{enumerate}
\begin{proof}
由于 $X$ 拟紧，若有限型拟凝聚 $\OO_X$-模 $\F$ 的 $\F(n)$ 由全局截面生成，则有限型模 $\F(n)$ 由有限多个这样的截面生成（0，5.2.3）。故 d) 推出 d')，且显然 d') 推出 d'')。每个拟凝聚 $\OO_X$-模是其有限型子 $\OO_X$-模的归纳极限（I，9.4.9）；因而为验证（4.5.2）的 c')，只须对有限型拟凝聚理想层验证，故 d'') 推出 c')。

只须再证 $\LL$ 丰沛时 d) 成立。取 $X$ 的有限开覆盖 $X_{f_i}$，其中 $f_i\in S_{n_i}$，且可令各 $X_{f_i}$ 仿射。把 $f_i$ 替换为适当幂不会改变 $X_{f_i}$，故可令所有 $n_i$ 等于同一整数 $m$。由有限型假设，$\F|X_{f_i}$ 由它在 $X_{f_i}$ 上有限多个截面 $h_{ij}$ 生成（I，1.3.13）。存在整数 $k_0$，使对每个 $(i,j)$，截面 $h_{ij}f_i^{\otimes k_0}$ 可延拓为 $\F(k_0m)$ 在 $X$ 上的截面（I，9.3.1）。于是对每个 $k\geq k_0$，$h_{ij}f_i^{\otimes k}$ 也可延拓为 $\F(km)$ 的全局截面，故 $\F(km)$ 由其全局截面生成。

对每个 $0\leq p<m$，$\F(p)$ 仍为有限型，故存在整数 $k_p$，使 $k\geq k_p$ 时，$\F(p)(km)=\F(p+km)$ 由全局截面生成。取 $n_0$ 大于所有 $k_pm$，则每个 $n\geq n_0$ 都可写成 $n=km+p$，其中 $k\geq k_p$、$0\leq p<m$，于是 $\F(n)$ 由全局截面生成。
\end{proof}

\egastatement{4.5.6}{命题}
设 $X$ 是拟紧概形，$\LL$ 是可逆 $\OO_X$-模。
\begin{enumerate}[label=\textup{(\roman*)}]
\item 对每个整数 $n>0$，$\LL$ 丰沛，当且仅当 $\LL^{\otimes n}$ 丰沛。
\item 设 $\LL'$ 是可逆 $\OO_X$-模，且对每个 $x\in X$，存在整数 $n>0$ 及 $\LL'^{\otimes n}$ 的全局截面 $s'$，使 $s'(x)\ne0$。若 $\LL$ 丰沛，则 $\LL\otimes\LL'$ 也丰沛。
\end{enumerate}
\begin{proof}
（i）直接由（4.5.2，a）的准则得出，因为 $X_{f^{\otimes n}}=X_f$。
另一方面，若 $\LL$ 丰沛，则对每个 $x\in X$ 及其邻域 $U$，存在 $m>0$ 及 $f\in\Gamma(X,\LL^{\otimes m})$，使 $x\in X_f\subset U$（4.5.2，a）。若 $f'\in\Gamma(X,\LL'^{\otimes n})$ 满足 $f'(x)\ne0$，令
$s=f^{\otimes n}\otimes f'^{\otimes m}\in\Gamma(X,(\LL\otimes\LL')^{\otimes mn})$，则 $s(x)\ne0$，故 $x\in X_s\subset X_f\subset U$。由（4.5.2，a），$\LL\otimes\LL'$ 丰沛。
\end{proof}

\egastatement{4.5.7}{推论}
两个丰沛 $\OO_X$-模的张量积是丰沛的。

\egastatement{4.5.8}{推论}
设 $\LL$ 是丰沛 $\OO_X$-模，$\LL'$ 是可逆 $\OO_X$-模。则存在整数 $n_0>0$，使对每个 $n\geq n_0$，$\LL^{\otimes n}\otimes\LL'$ 丰沛，且由其全局截面生成。
\begin{proof}
由（4.5.5），存在整数 $m_0$，使 $m\geq m_0$ 时，$\LL^{\otimes m}\otimes\LL'$ 由全局截面生成；由（4.5.6），可取 $n_0=m_0+1$。
\end{proof}

\egastatement{4.5.9}{注记}
令 $P=H^1(X,\OO_X^*)$ 为可逆 $\OO_X$-模的同构类群（0，5.4.7），$P^+$ 为其中丰沛层的同构类构成的子集。设 $P^+\ne\varnothing$。由（4.5.7）和（4.5.8），有
\[
P^++P^+\subset P^+,
\qquad P^+-P^+=P.
\]
换言之，$P^+\cup\{0\}$ 是 $P$ 上某个与群结构相容的预序的正元素集合；由（4.5.8），这一预序甚至是阿基米德的。因此有时称丰沛层为“正层”，称其逆层为“负层”；我们不采用这种术语。

\egastatement{4.5.10}{命题}
设 $Y$ 是仿射概形，$q:X\to Y$ 是分离拟紧态射，$\LL$ 是可逆 $\OO_X$-模。
\begin{enumerate}[label=\textup{(\roman*)}]
\item 若 $\LL$ 相对于 $q$ 甚丰沛，则 $\LL$ 丰沛。
\item 再设 $q$ 为有限型。则 $\LL$ 丰沛，当且仅当它具有下列等价性质之一：
\begin{enumerate}[label=\textup{e)}]
\item 存在 $n_0>0$，使每个整数 $n\geq n_0$ 对应的 $\LL^{\otimes n}$ 相对于 $q$ 甚丰沛。
\item[\textup{e')}] 存在 $n>0$，使 $\LL^{\otimes n}$ 相对于 $q$ 甚丰沛。
\end{enumerate}
\end{enumerate}
\begin{proof}
第一项由甚丰沛模的定义（4.4.2）得出：若 $A$ 是 $Y$ 的环，则存在 $A$-模 $E$ 及满同态
$\psi:q^*((\Sym(E))^{\sim})\to\bigoplus_{n\geq0}\LL^{\otimes n}$，使 $i=r_{\LL,\psi}:X\to P=\Pb(E)$ 处处有定义且为浸入，并且 $\LL=i^*(\OO_P(1))$。当 $f$ 遍历 $(\Sym(E))_+$ 的齐次元素时，$D_+(f)$ 构成 $P$ 的拓扑的一组基；由（3.7.3.1），$i^{-1}(D_+(f))=X_{\psi^\flat(f)}$，故（4.5.2）的 a) 成立，$\LL$ 丰沛。

现设 $q$ 为有限型且 $\LL$ 丰沛，证明 e)。由（4.5.2，b）和（4.4.1，i），存在整数 $k_0>0$，使 $\LL^{\otimes k_0}$ 相对于 $q$ 甚丰沛。另一方面，由（4.5.5），存在整数 $m_0$，使每个 $m\geq m_0$ 对应的 $\LL^{\otimes m}$ 由全局截面生成。令 $n_0=k_0+m_0$。若 $n\geq n_0$，可写成 $n=k_0+m$，其中 $m\geq m_0$，于是
$\LL^{\otimes n}=\LL^{\otimes k_0}\otimes\LL^{\otimes m}$。由（4.4.8）的准则和（3.4.7），$\LL^{\otimes n}$ 相对于 $q$ 甚丰沛。

显然 e) 推出 e')；由（i）及（4.5.6，i），e') 推出 $\LL$ 丰沛。
\end{proof}

\egapar{4.5.10.1}
\textbf{引理（4.4.10.1）的证明。}
记 $\E(n)=\E\otimes\K^{\otimes n}$。由于 $g$ 分离（4.4.2），可应用（3.4.8）的论证，将问题归结为证明：当 $n$ 充分大时，典范同态
$g^*(g_*(\E(n)))\to\E(n)$ 满射。又因 $Z$ 拟紧，（3.4.6）的论证将命题归结到 $Z$ 仿射的情形。此时由（4.5.10，i），$\K$ 丰沛，结论由（4.5.5，d'）得出。\hfill$\square$

\egastatement{4.5.11}{推论}
设 $Y$ 是仿射概形，$q:X\to Y$ 是分离有限型态射，$\LL$ 是丰沛 $\OO_X$-模，$\LL'$ 是可逆 $\OO_X$-模。存在整数 $n_0$，使每个 $n\geq n_0$ 对应的 $\LL^{\otimes n}\otimes\LL'$ 相对于 $q$ 甚丰沛。
\begin{proof}
存在 $m_0$，使 $m\geq m_0$ 时，$\LL^{\otimes m}\otimes\LL'$ 由全局截面生成（4.5.8）；另有 $k_0$，使 $k\geq k_0$ 时，$\LL^{\otimes k}$ 相对于 $q$ 甚丰沛。因此当 $k\geq k_0$ 时，$\LL^{\otimes(k+m_0)}\otimes\LL'$ 相对于 $q$ 甚丰沛（4.4.8；3.4.7）。
\end{proof}

\egastatement{4.5.12}{注记}
尚不知道：对 $\OO_X$-模 $\LL$，假设 $\LL^{\otimes n}$ 相对于 $q$ 甚丰沛，是否推出 $\LL^{\otimes(n+1)}$ 也如此。\transnote{这里保留原著写作时的叙述，不将它作为有关当前研究状态的断言。}

\egastatement{4.5.13}{命题}
设 $X$ 是拟紧预概形，$Z$ 是由 $\OO_X$ 的拟凝聚幂零理想层 $\J$ 定义的闭子预概形，$j:Z\to X$ 为典范嵌入。可逆 $\OO_X$-模 $\LL$ 丰沛，当且仅当 $\LL'=j^*(\LL)$ 是丰沛 $\OO_Z$-模。
\begin{proof}
必要性：对 $\LL^{\otimes n}$ 的每个全局截面 $f$，记 $f'=f\otimes1$ 为其典范像；这是
$\LL'^{\otimes n}=\LL^{\otimes n}\otimes_{\OO_X}(\OO_X/\J)$ 在空间 $Z$（与 $X$ 相同）上的截面。显然 $X_f=Z_{f'}$，故由（4.5.2，a），$\LL'$ 丰沛。

为证充分性，先考虑有限列预概形 $X_k=(X,\OO_X/\J^{k+1})$；每一项都是下一项中由平方零理想层定义的闭子预概形。因此可归结到 $\J^2=0$ 的情形。并且，因假设 $X_0$ 为概形，$X$ 也是概形（4.5.3；I，5.5.1）。准则（4.5.2，a）表明，只须证明下面的引理。\renewcommand{\qedsymbol}{}
\end{proof}

\egastatement{4.5.13.1}{引理}
在（4.5.13）的假设下，再设 $\J^2=0$。设 $\LL$ 为可逆 $\OO_X$-模，$g$ 为 $\LL'^{\otimes n}$ 在 $Z$ 上的截面，且 $Z_g$ 仿射。则存在整数 $m>0$，使 $g^{\otimes m}$ 是 $\LL^{\otimes nm}$ 的一个全局截面 $f$ 的典范像。
\begin{proof}
由于 $\F\mapsto\F(n)$ 是正合函子，有 $\OO_X$-模的正合列
\[
0\longrightarrow\J(n)\longrightarrow\OO_X(n)=\LL^{\otimes n}
\longrightarrow\OO_Z(n)=\LL'^{\otimes n}\longrightarrow0.
\]
从而有上同调正合列
\[
\begin{split}
0\longrightarrow\Gamma(X,\J(n))\longrightarrow\Gamma(X,\LL^{\otimes n})
\longrightarrow\Gamma(X,\LL'^{\otimes n})
\longrightarrow H^1(X,\J(n)).
\end{split}
\]
它特别将 $g$ 对应于元素 $\delta g\in H^1(X,\J(n))$。

注意，由 $\J^2=0$，可将 $\J$ 看成拟凝聚 $\OO_Z$-模，且对每个 $k$，有
$\LL'^{\otimes k}\otimes_{\OO_Z}\J(n)=\J(n+k)$。因此，对任一截面
$s\in\Gamma(X,\LL'^{\otimes k})$，以 $s$ 作张量乘法给出 $\OO_Z$-模同态
$\J(n)\to\J(n+k)$，并由此给出上同调群同态
$H^1(X,\J(n))\to H^1(X,\J(n+k))$。

我们将证明，对充分大的某个 $m>0$，有
\begin{equation}
g^{\otimes m}\otimes\delta g=0.\tag{4.5.13.2}
\end{equation}
事实上，$Z_g$ 是 $Z$ 的仿射开集，因而把 $\J(n)$ 看成 $\OO_Z$-模时，有
$H^1(Z_g,\J(n))=0$（I，5.1.9.2）。特别地，记 $g'=g|Z_g$，考察它在
$\delta:\Gamma(Z_g,\LL'^{\otimes n})\to H^1(Z_g,\J(n))$ 下的像，则 $\delta g'=0$。

为明确这个关系，注意阿贝尔群层的一维上同调与其 Čech 上同调相同（G，II，5.9）。因此，为求出 $\delta g$，须取 $X$ 的一个足够细的开覆盖 $(U_\alpha)$；可设此覆盖有限，且由仿射开集组成。对每个 $\alpha$ 取
$g_\alpha\in\Gamma(U_\alpha,\LL^{\otimes n})$，使其在
$\Gamma(U_\alpha,\LL'^{\otimes n})$ 中的典范像为 $g|U_\alpha$，再取上闭链
$(g_{\alpha|\beta}-g_{\beta|\alpha})$ 的类；这里 $g_{\alpha|\beta}$ 为 $g_\alpha$ 在 $U_\alpha\cap U_\beta$ 上的限制，上闭链取值于 $\J(n)$。必要时加细覆盖，还可设 $\delta g'$ 是用覆盖 $U_\alpha\cap Z_g$ 和截面 $g_\alpha|(U_\alpha\cap Z_g)$ 以同样方式求得。于是关系 $\delta g'=0$ 表示：对每个 $\alpha$，存在
$h_\alpha\in\Gamma(U_\alpha\cap Z_g,\J(n))$，使
\[
(g_{\alpha|\beta}-g_{\beta|\alpha})|(U_\alpha\cap U_\beta\cap Z_g)
=h_{\alpha|\beta}-h_{\beta|\alpha},
\]
其中 $h_{\alpha|\beta}$ 表示相应限制（G，II，5.11）。由（I，9.3.1），存在整数 $m>0$，使对每个 $\alpha$，$g^{\otimes m}\otimes h_\alpha$ 是某个
$t_\alpha\in\Gamma(U_\alpha,\J(n+nm))$ 的限制。因此对每对指标，有
\[
g^{\otimes m}\otimes(g_{\alpha|\beta}-g_{\beta|\alpha})
=t_{\alpha|\beta}-t_{\beta|\alpha},
\]
这证明（4.5.13.2）。\transnote{原文从在 $Z_g$ 上的等式直接写出此全局重叠上的等式。这里还须按所引（I，9.3.1，i），在有限多个重叠上进一步增加同一个 $m$，以消去限制为零的差；原文把这一增加幂次的步骤略去了。}

另一方面，若 $s\in\Gamma(X,\OO_Z(p))$、$t\in\Gamma(X,\OO_Z(q))$，则在 $H^1(X,\J(p+q))$ 中有
\begin{equation}
\delta(s\otimes t)=(\delta s)\otimes t+s\otimes(\delta t).
\tag{4.5.13.3}
\end{equation}
计算两端时，仍可取 $X$ 的一个开覆盖 $(U_\alpha)$，并在每个 $U_\alpha$ 上取
$s_\alpha\in\Gamma(U_\alpha,\OO_X(p))$、$t_\alpha\in\Gamma(U_\alpha,\OO_X(q))$，使它们的典范像分别为 $s|U_\alpha$、$t|U_\alpha$。公式（4.5.13.3）即由下式得出：
\[
\begin{aligned}
&s_{\alpha|\beta}\otimes t_{\alpha|\beta}
-s_{\beta|\alpha}\otimes t_{\beta|\alpha}\\
&\quad=(s_{\alpha|\beta}-s_{\beta|\alpha})\otimes t_{\alpha|\beta}
+s_{\beta|\alpha}\otimes(t_{\alpha|\beta}-t_{\beta|\alpha}).
\end{aligned}
\]
由关于 $k$ 的归纳法，得到
\begin{equation}
\delta(g^{\otimes k})=k\,g^{\otimes(k-1)}\otimes\delta g.
\tag{4.5.13.4}
\end{equation}
由（4.5.13.2）和（4.5.13.4），$\delta(g^{\otimes(m+1)})=0$。因此 $g^{\otimes(m+1)}$ 是某个
$f\in\Gamma(X,\LL^{\otimes n(m+1)})$ 的典范像。这也完成了（4.5.13）的证明。
\end{proof}

\egastatement{4.5.14}{推论}
设 $X$ 是诺特概形，$j:X_{\red}\to X$ 是典范嵌入。可逆 $\OO_X$-模 $\LL$ 丰沛，当且仅当 $j^*(\LL)$ 是丰沛 $\OO_{X_{\red}}$-模。
\begin{proof}
由（I，6.1.6）得出。
\end{proof}

\subsection{相对丰沛层}
\src{89--94}

\egastatement{4.6.1}{定义}
设 $f:X\to Y$ 是拟紧态射，$\LL$ 是可逆 $\OO_X$-模。称 $\LL$ \emph{相对于 $f$ 丰沛}，或\emph{相对于 $Y$ 丰沛}，也称为 \emph{$f$-丰沛}、\emph{$Y$-丰沛}（若不致与（4.5.3）中的概念混淆，也简称\emph{丰沛}），若存在 $Y$ 的仿射开覆盖 $(U_\alpha)$，使令 $X_\alpha=f^{-1}(U_\alpha)$ 后，每个 $\LL|X_\alpha$ 都是丰沛 $\OO_{X_\alpha}$-模。

$f$-丰沛 $\OO_X$-模的存在，必然推出 $f$ 分离（4.5.3；I，5.5.5）。

\egastatement{4.6.2}{命题}
设 $f:X\to Y$ 是拟紧态射，$\LL$ 是可逆 $\OO_X$-模。若 $\LL$ 相对于 $f$ 甚丰沛，则它相对于 $f$ 丰沛。
\begin{proof}
由甚丰沛性的关于 $Y$ 的局部性（4.4.5）、定义（4.6.1）和准则（4.5.10，i）得出。
\end{proof}

\egastatement{4.6.3}{命题}
设 $f:X\to Y$ 是拟紧态射，$\LL$ 是可逆 $\OO_X$-模，令
$\SSS=\bigoplus_{n\geq0}f_*(\LL^{\otimes n})$ 为相应分次 $\OO_Y$-代数。下列条件等价：
\begin{enumerate}[label=\alph*)]
\item $\LL$ 是 $f$-丰沛的。
\item $\SSS$ 拟凝聚，且典范同态
$\sigma:f^*(\SSS)\to\bigoplus_{n\geq0}\LL^{\otimes n}$（0，4.4.3）使 $Y$-态射
$r_{\LL,\sigma}:G(\sigma)\to\Proj(\SSS)=P$ 处处有定义，且为占优开浸入。
\item[\textup{b')}] $f$ 分离，$Y$-态射 $r_{\LL,\sigma}$ 处处有定义，且在底空间上给出 $X$ 到 $\Proj(\SSS)$ 的一个子空间的同胚。
\end{enumerate}
此时，对每个 $n\in\mathbf Z$，在（3.7.9.1）中定义的典范同态
\begin{equation}
r_{\LL,\sigma}^*(\OO_P(n))\longrightarrow\LL^{\otimes n}
\tag{4.6.3.1}
\end{equation}
是同构。最后，对每个拟凝聚 $\OO_X$-模 $\F$，令
$\M=\bigoplus_{n\geq0}f_*(\F\otimes\LL^{\otimes n})$，则（3.7.9.2）中定义的典范同态
\begin{equation}
r_{\LL,\sigma}^*(\wt\M)\longrightarrow\F
\tag{4.6.3.2}
\end{equation}
也是同构。
\begin{proof}
已注意到 a) 推出 $f$ 分离，因而 $\SSS$ 拟凝聚（I，9.2.2，a）。由于处处有定义及为浸入都是关于 $Y$ 的局部性质，为证明 a) 推出 b)，可设 $Y$ 仿射且 $\LL$ 丰沛；于是由（4.5.2，b）得出。显然 b) 推出 b')；为证 b') 推出 a)，取 $Y$ 的仿射开覆盖 $(U_\alpha)$，再对各 $\LL|f^{-1}(U_\alpha)$ 应用准则（4.5.2，b'）即可。

\begin{sloppypar}
对最后两项断言，使用这里 $\sigma^\flat$ 为恒等同态这一事实，以及（3.7.9.1）、（3.7.9.2）的具体描述，立刻得出（4.6.3.1）是同构。至于（4.6.3.2），可归结到 $Y$ 仿射、从而 $\LL$ 丰沛的情形；此时该同态显然单射，而准则（4.5.2，c）表明它满射，故为同构。
\end{sloppypar}
\end{proof}

\egastatement{4.6.4}{推论}
设 $(U_\alpha)$ 为 $Y$ 的开覆盖。$\LL$ 相对于 $Y$ 丰沛，当且仅当每个 $\LL|f^{-1}(U_\alpha)$ 相对于 $U_\alpha$ 丰沛。
\begin{proof}
条件（4.6.3，b）是关于 $Y$ 的局部条件。
\end{proof}

\egastatement{4.6.5}{推论}
设 $\K$ 是可逆 $\OO_Y$-模。$\LL$ 是 $Y$-丰沛的，当且仅当 $\LL\otimes f^*(\K)$ 是 $Y$-丰沛的。
\begin{proof}
取开覆盖 $(U_\alpha)$ 使各 $\K|U_\alpha\simeq\OO_Y|U_\alpha$，由（4.6.4）即得。
\end{proof}

\egastatement{4.6.6}{推论}
设 $Y$ 仿射。$\LL$ 是 $Y$-丰沛的，当且仅当 $\LL$ 丰沛。
\begin{proof}
这直接由定义（4.6.1）及准则（4.6.3，b）、（4.5.2，b）得出，因为此时按定义有
$\Proj(\SSS)=\Proj(\Gamma(Y,\SSS))$。
\end{proof}

\egastatement{4.6.7}{推论}
设 $f:X\to Y$ 是拟紧态射。假设存在拟凝聚 $\OO_Y$-模 $\E$ 和 $Y$-态射 $g:X\to P=\Pb(\E)$，使 $g$ 在底空间上为 $X$ 到 $P$ 的一个子空间的同胚。则 $\LL=g^*(\OO_P(1))$ 是 $f$-丰沛的。
\begin{proof}
可设 $Y$ 仿射，此时由准则（4.5.2，a）、公式（3.7.3.1）和（4.2.3）得出。
\end{proof}

\egastatement{4.6.8}{命题}
设 $X$ 是拟紧概形，或底空间为诺特空间的预概形；$f:X\to Y$ 是分离拟紧态射。可逆 $\OO_X$-模 $\LL$ 是 $f$-丰沛的，当且仅当下列等价条件之一成立：
\begin{enumerate}[label=\textup{c)}]
\item 对每个有限型拟凝聚 $\OO_X$-模 $\F$，存在整数 $n_0>0$，使每个 $n\geq n_0$ 对应的典范同态
$\sigma:f^*(f_*(\F\otimes\LL^{\otimes n}))\to\F\otimes\LL^{\otimes n}$ 满射。
\item[\textup{c')}] 对 $\OO_X$ 的每个有限型拟凝聚理想层 $\J$，存在整数 $n_0>0$，使典范同态
$\sigma:f^*(f_*(\J\otimes\LL^{\otimes n_0}))\to\J\otimes\LL^{\otimes n_0}$ 满射。
\end{enumerate}
\begin{proof}
$X$ 拟紧，故 $f(X)$ 拟紧，因而可用 $Y$ 的有限多个仿射开集 $U_i$ 覆盖 $f(X)$。当 $\LL$ 为 $f$-丰沛时，为证 c)，可把 $Y$ 替换为各 $U_i$，把 $X$ 替换为 $f^{-1}(U_i)$：若对各 $i$ 得到整数 $n_i$，使 c) 对 $U_i$、$f^{-1}(U_i)$、$\LL|f^{-1}(U_i)$ 在 $n\geq n_i$ 时成立，则取 $n_0=\max_i n_i$ 即得一般情形。当 $Y$ 仿射时，考虑（4.6.6），c) 由（4.5.5，d）得出。显然 c) 推出 c')。

最后，为证 c') 推出 $\LL$ 为 $f$-丰沛，仍可归结到 $Y$ 仿射的情形：$\OO_X|f^{-1}(U_i)$ 的每个有限型拟凝聚理想层 $\J_i$，都是 $\OO_X$ 的某个有限型凝聚理想层的限制（I，9.4.7）。\transnote{这里“凝聚”对应原文 \textit{cohérent}；所引延拓结果在一般情形所给出的是“有限型拟凝聚”，不能在非诺特情形由有限型直接断定凝聚。此处沿原文保留措辞，并标明这一差别。}
由 c')，$\J_i\otimes\LL^{\otimes n_0}|f^{-1}(U_i)$ 由其截面生成（I，9.2.2；3.4.7）。应用准则（4.5.5，d''）即可。
\end{proof}

\egastatement{4.6.9}{命题}
设 $f:X\to Y$ 是拟紧态射，$\LL$ 是可逆 $\OO_X$-模。
\begin{enumerate}[label=\textup{(\roman*)}]
\item 对整数 $n>0$，$\LL$ 为 $f$-丰沛，当且仅当 $\LL^{\otimes n}$ 为 $f$-丰沛。
\item 设 $\LL'$ 是可逆 $\OO_X$-模，且存在整数 $n>0$，使典范同态
$\sigma:f^*(f_*(\LL'^{\otimes n}))\to\LL'^{\otimes n}$ 满射。若 $\LL$ 为 $f$-丰沛，则 $\LL\otimes\LL'$ 也为 $f$-丰沛。
\end{enumerate}
\begin{proof}
立即归结到 $Y$ 仿射的情形，再由（4.5.6）得出。
\end{proof}

\egastatement{4.6.10}{推论}
两个 $f$-丰沛 $\OO_X$-模的张量积是 $f$-丰沛的。

\egastatement{4.6.11}{命题}
设 $Y$ 是拟紧预概形，$f:X\to Y$ 是有限型态射，$\LL$ 是可逆 $\OO_X$-模。$\LL$ 为 $f$-丰沛，当且仅当它具有下列等价性质之一：
\begin{enumerate}[label=\textup{d)}]
\item 存在 $n_0>0$，使每个整数 $n\geq n_0$ 对应的 $\LL^{\otimes n}$ 相对于 $f$ 甚丰沛。
\item[\textup{d')}] 存在 $n>0$，使 $\LL^{\otimes n}$ 相对于 $f$ 甚丰沛。
\end{enumerate}
\begin{proof}
若 $\LL$ 相对于 $f$ 丰沛，则存在 $Y$ 的有限仿射开覆盖 $(U_i)$，使各 $\LL|f^{-1}(U_i)$ 丰沛。由（4.5.10），存在整数 $n_0$，使对每个 $i$ 及每个 $n\geq n_0$，$\LL^{\otimes n}|f^{-1}(U_i)$ 相对于 $f^{-1}(U_i)\to U_i$ 甚丰沛；于是由（4.4.5），$\LL^{\otimes n}$ 相对于 $f$ 甚丰沛。反之，d') 已推出 $\LL^{\otimes n}$ 为 $f$-丰沛（4.6.2），进而 $\LL$ 为 $f$-丰沛（4.6.9，i）。
\end{proof}

\egastatement{4.6.12}{推论}
设 $Y$ 是拟紧预概形，$f:X\to Y$ 是有限型态射，$\LL,\LL'$ 是可逆 $\OO_X$-模。若 $\LL$ 为 $f$-丰沛，则存在 $n_0$，使对每个 $n\geq n_0$，$\LL^{\otimes n}\otimes\LL'$ 相对于 $f$ 甚丰沛。
\begin{proof}
与（4.6.11）同样，使用 $Y$ 的有限仿射开覆盖以及（4.5.11）即可。
\end{proof}

\egastatement{4.6.13}{命题}
\begin{enumerate}[label=\textup{(\roman*)}]
\item 对每个预概形 $Y$，每个可逆 $\OO_Y$-模 $\LL$ 都相对于恒等态射 $1_Y$ 丰沛。
\item[\textup{(i bis)}] 设 $f:X\to Y$ 为拟紧态射，$j:X'\to X$ 为拟紧态射，且在底空间上是 $X'$ 到 $X$ 的一个子空间的同胚。若 $\LL$ 相对于 $f$ 丰沛，则 $j^*(\LL)$ 相对于 $f\circ j$ 丰沛。
\item 设 $Z$ 是拟紧预概形，$f:X\to Y$、$g:Y\to Z$ 是拟紧态射，$\LL$ 为 $f$-丰沛 $\OO_X$-模，$\K$ 为 $g$-丰沛 $\OO_Y$-模。则存在整数 $n_0>0$，使对每个 $n\geq n_0$，$\LL\otimes f^*(\K^{\otimes n})$ 相对于 $g\circ f$ 丰沛。
\item 设 $f:X\to Y$ 为拟紧态射，$g:Y'\to Y$ 为任意态射，$X'=X_{(Y')}$。若 $\LL$ 相对于 $f$ 丰沛，则 $\LL'=\LL\otimes_Y\OO_{Y'}$ 相对于 $f_{(Y')}$ 丰沛。
\item 设 $f_i:X_i\to Y_i$（$i=1,2$）为拟紧 $S$-态射。若 $\LL_i$ 相对于 $f_i$ 丰沛，则 $\LL_1\otimes_S\LL_2$ 相对于 $f_1\times_S f_2$ 丰沛。
\item 设 $f:X\to Y$、$g:Y\to Z$ 为态射，且 $g\circ f$ 拟紧。若 $\OO_X$-模 $\LL$ 相对于 $g\circ f$ 丰沛，并且 $g$ 分离或 $X$ 的底空间局部诺特，则 $\LL$ 相对于 $f$ 丰沛。
\item 设 $f:X\to Y$ 为拟紧态射，$j:X_{\red}\to X$ 为典范嵌入。若 $\LL$ 相对于 $f$ 丰沛，则 $j^*(\LL)$ 相对于 $f_{\red}$ 丰沛。
\end{enumerate}
\begin{proof}
先注意，（v）和（vi）可用（4.4.10）中同样的论证，由（i）、（i bis）、（iv）推出，只须将（4.4.5）改用（4.6.4）；细节留给读者。（i）显然由（4.4.10，i）及（4.6.2）得出。为证明（i bis）、（iii）、（iv），使用下面的引理。\renewcommand{\qedsymbol}{}
\end{proof}

\egastatement{4.6.13.1}{引理}
\begin{enumerate}[label=\textup{(\roman*)}]
\item 设 $u:Z\to S$ 是态射，$\LL$ 是可逆 $\OO_S$-模，$s$ 是其在 $S$ 上的截面，$s'$ 为与之典范对应的 $\LL'=u^*(\LL)$ 在 $Z$ 上的截面。则 $Z_{s'}=u^{-1}(S_s)$。
\item 设 $Z,Z'$ 为 $S$-预概形，$p,p'$ 是 $T=Z\times_S Z'$ 的投影。设 $\LL$（分别地，$\LL'$）是可逆 $\OO_Z$-模（分别地，$\OO_{Z'}$-模），$t,t'$ 是相应全局截面，$s,s'$ 是与之典范对应的 $p^*(\LL),p'^*(\LL')$ 在 $T$ 上的截面。则 $T_{s\otimes s'}=Z_t\times_S Z'_{t'}$。
\end{enumerate}
\begin{proof}
由定义，可归结到所考虑的预概形全部仿射的情形。在（i）中还可设 $\LL=\OO_S$，于是直接由（I，1.2.2.2）得出。同样，在（ii）中可设 $\LL=\OO_Z$、$\LL'=\OO_{Z'}$，此时断言即为引理（4.3.2.4）。
\end{proof}

\noindent\textbf{（4.6.13）的证明续。}
先证（i bis）。可设 $Y$ 仿射（4.6.4），故 $\LL$ 丰沛（4.6.6）。当 $s$ 遍历所有 $n>0$ 对应的 $\Gamma(X,\LL^{\otimes n})$ 时，$X_s$ 构成 $X$ 的拓扑的一组基（4.5.2，a）；由假设，$j^{-1}(X_s)$ 构成 $X'$ 的拓扑的一组基。由引理（4.6.13.1，i）及（4.5.2，a），$j^*(\LL)$ 丰沛。

再证（iii）。仍可设 $Y,Y'$ 仿射（4.6.4），从而投影 $h:X'\to X$ 仿射（1.5.5）。由于 $\LL$ 丰沛（4.6.6），当 $s$ 遍历各 $\LL^{\otimes n}$（$n>0$）的使 $X_s$ 仿射的全局截面时，这些 $X_s$ 覆盖 $X$（4.5.2，a'）；故 $h^{-1}(X_s)$ 仿射（1.2.5）且覆盖 $X'$。由引理（4.6.13.1，i）及（4.5.2，a'），$\LL'$ 丰沛；而 $f_{(Y')}$ 拟紧（I，6.6.4，iii）。

为证（iv），先注意 $f_1\times_S f_2$ 拟紧（I，6.6.4，iv）。又可设 $S,Y_1,Y_2$ 仿射（4.6.4；I，3.2.7），故各 $\LL_i$ 丰沛（4.6.6）。当 $s_i$ 遍历 $\LL_i^{\otimes n_i}$ 的使 $(X_i)_{s_i}$ 仿射的截面时，开集
$(X_1)_{s_1}\times_S(X_2)_{s_2}$ 覆盖 $X_1\times_S X_2$（4.5.2，a'）。把 $s_1,s_2$ 替换为适当幂不改变这些开集，故总可令 $n_1=n_2$。由（4.6.13.1，ii）及（4.5.2，a'），$\LL_1\otimes_S\LL_2$ 丰沛；又因 $Y_1\times_S Y_2$ 仿射，所求断言由（4.6.6）得出。

只剩（ii）。用（4.4.10）中的同样论证，这次使用（4.6.4），可归结到 $Z$ 仿射的情形。此时 $\K$ 丰沛而 $Y$ 拟紧，故存在有限多个截面
$s_i\in\Gamma(Y,\K^{\otimes k_i})$，使 $Y_{s_i}$ 仿射且覆盖 $Y$（4.5.2，a'）。把 $s_i$ 替换为适当幂，还可令各 $k_i$ 等于同一整数 $k$。记 $s'_i$ 为与 $s_i$ 典范对应的 $f^*(\K^{\otimes k})$ 的全局截面，则
$X_{s'_i}=f^{-1}(Y_{s_i})$ 覆盖 $X$（4.6.13.1，i）。由于 $\LL|X_{s'_i}$ 丰沛（4.6.4；4.6.6），对每个 $i$ 存在有限多个截面
$t_{ij}\in\Gamma(X,\LL^{\otimes n_{ij}})$，使 $X_{t_{ij}}$ 仿射、包含在 $X_{s'_i}$ 中并覆盖 $X_{s'_i}$（4.5.2，a'）；还可令各 $n_{ij}$ 等于同一整数 $n$。\transnote{原文将这里的截面空间写成 $\Gamma(X,\LL^{\otimes n_{ij}})$，但后文再对这些截面作延拓；从上下文看，这里所用的是 $\LL|X_{s'_i}$ 的幂在 $X_{s'_i}$ 上的截面，$X_{t_{ij}}$ 也相对于该开子预概形解释。}

$X$ 分离且拟紧，故存在整数 $m>0$，并对每个 $(i,j)$ 存在截面
\[
u_{ij}\in\Gamma\bigl(X,\LL^{\otimes n}\otimes_{\OO_X}f^*(\K^{\otimes mk})\bigr),
\]
使 $t_{ij}\otimes s_i'^{\otimes m}$ 为 $u_{ij}$ 在 $X_{s'_i}$ 上的限制（I，9.3.1）。又有 $X_{u_{ij}}=X_{t_{ij}}$，故各 $X_{u_{ij}}$ 仿射且覆盖 $X$。还可令 $m=nr$；置 $n_0=rk$，由（4.5.2，a'）可见 $\LL\otimes f^*(\K^{\otimes n_0})$ 丰沛。\transnote{为令延拓截面在 $X_{s'_i}$ 外消失，上述等式 $X_{u_{ij}}=X_{t_{ij}}$ 还使用了增加 $m$ 并再乘以 $s'_i$ 的步骤；原文将它包含在“存在 $m$”的选择中。}
另有 $h_0>0$，使 $h\geq h_0$ 时，$\K^{\otimes h}$ 由其全局截面生成（4.5.5）；由逆像的定义（0，3.7.1；4.4.1），$f^*(\K^{\otimes h})$ 也由其全局截面生成。故由（4.5.6），$h\geq h_0$ 时，$\LL\otimes f^*(\K^{\otimes(n_0+h)})$ 丰沛。证明完成。\hfill$\square$

\egastatement{4.6.14}{注记}
在（4.6.13，ii）的条件下，不可认为 $\LL\otimes f^*(\K)$ 就相对于 $g\circ f$ 丰沛。事实上，$\LL\otimes f^*(\K^{-1})$ 也相对于 $f$ 丰沛（4.6.5）；若作上述误判，便会推出 $\LL$ 相对于 $g\circ f$ 丰沛。特别取 $g$ 为恒等态射，就会得到每个可逆 $\OO_X$-模都相对于 $f$ 丰沛，而一般并非如此（参见（5.1.6）、（5.3.4，i）和（5.3.1））。\transnote{原文这里写“取 $g$ 为恒等态射”；若要从所述错误规则推出“每个可逆模都丰沛”的结论，应把两个复合因子的角色相应调整，即取前一因子为恒等态射。此注不改变前述命题。}

\egastatement{4.6.15}{命题}
设 $f:X\to Y$ 是拟紧态射，$\J$ 是 $\OO_X$ 的局部幂零拟凝聚理想层，$Z$ 是它定义的闭子预概形，$j:Z\to X$ 为典范嵌入。可逆 $\OO_X$-模 $\LL$ 相对于 $f$ 丰沛，当且仅当 $j^*(\LL)$ 相对于 $f\circ j$ 丰沛。
\begin{proof}
问题关于 $Y$ 局部（4.6.4），故可设 $Y$ 仿射。此时 $X$ 拟紧，因而可设 $\J$ 幂零。考虑（4.6.6），命题即为（4.5.13）。
\end{proof}

\egastatement{4.6.16}{推论}
设 $X$ 为局部诺特预概形，$f:X\to Y$ 为拟紧态射。可逆 $\OO_X$-模 $\LL$ 相对于 $f$ 丰沛，当且仅当它在典范嵌入 $X_{\red}\to X$ 下的逆像 $\LL'$ 相对于 $f_{\red}$ 丰沛。
\begin{proof}
必要性已在（4.6.13，vi）证明。反之，为证 $\LL$ 相对于 $f$ 丰沛，可设 $Y$ 仿射（4.6.4）；此时 $Y_{\red}$ 也仿射，故 $\LL'$ 丰沛（4.6.6）。此时 $X$ 诺特，且 $X_{\red}$ 是由拟凝聚幂零理想层定义的闭子预概形（I，6.1.6），所以由（4.5.13），$\LL$ 也丰沛。
\end{proof}

\egastatement{4.6.17}{命题}
沿用（4.4.11）的记号和假设。$\LL''$ 相对于 $f''$ 丰沛，当且仅当 $\LL$ 相对于 $f$ 丰沛且 $\LL'$ 相对于 $f'$ 丰沛。
\begin{proof}
必要性由（4.6.13，i bis）得出。为证充分性，可设 $Y$ 仿射；注意 $\LL$ 在 $X$ 上的截面可以在 $X'$ 上取零而延拓为 $\LL''$ 的全局截面，再分别对 $\LL,\LL',\LL''$ 应用（4.5.2，a），即得 $\LL''$ 丰沛。
\end{proof}

\egastatement{4.6.18}{命题}
设 $Y$ 是拟紧预概形，$\SSS$ 是有限型拟凝聚分次 $\OO_Y$-代数，$X=\Proj(\SSS)$，$f:X\to Y$ 为结构态射。则 $f$ 为有限型，且存在整数 $d>0$，使 $\OO_X(d)$ 可逆并且 $f$-丰沛。
\begin{proof}
由（3.1.10），存在整数 $d>0$，使 $\SSS^{(d)}$ 由 $\SSS_d$ 生成。在 $X$ 与 $X^{(d)}=\Proj(\SSS^{(d)})$ 的典范同构下，$\OO_X(d)$ 等同于 $\OO_{X^{(d)}}(1)$（3.2.9，ii）。故可归结到 $\SSS$ 由 $\SSS_1$ 生成的情形。由（4.4.3）及（4.6.2）即得结论，并注意 $f$ 为有限型（3.4.1）。
\end{proof}

\section{拟仿射态射；拟射影态射；固有态射；射影态射}
\subsection{拟仿射态射}
\src{94--96}

\egastatement{5.1.1}{定义}
称一个概形为\emph{拟仿射概形}，若它同构于某个仿射概形在其一个拟紧开集上诱导的子概形。称态射 $f:X\to Y$ 为\emph{拟仿射态射}，或称通过 $f$ 看成 $Y$-预概形的 $X$ 为\emph{拟仿射 $Y$-概形}，若存在 $Y$ 的仿射开覆盖 $(U_\alpha)$，使各 $f^{-1}(U_\alpha)$ 为拟仿射概形。

显然，拟仿射态射是分离的（I，5.5.5；5.5.8）和拟紧的（I，6.6.1）；每个仿射态射显然都是拟仿射态射。
回顾：对每个预概形 $X$，令 $A=\Gamma(X,\OO_X)$，恒等同态 $A\to A=\Gamma(X,\OO_X)$ 定义一个称为\emph{典范态射}的态射 $X\to\Spec(A)$（I，2.2.4）。回顾 $\Proj(A[T])$ 典范等同于 $\Spec(A)$（3.1.7），可见它正是（4.5.1）在 $\LL=\OO_X$ 时定义的典范态射。

\egastatement{5.1.2}{命题}
设 $X$ 是拟紧概形，或底空间为诺特空间的预概形；令 $A=\Gamma(X,\OO_X)$。下列条件等价：
\begin{enumerate}[label=\textup{a)}]
\item $X$ 是拟仿射概形。
\item[\textup{b)}] 典范态射 $u:X\to\Spec(A)$ 是开浸入。
\item[\textup{b')}] 典范态射 $u:X\to\Spec(A)$ 在底空间上是 $X$ 到 $\Spec(A)$ 的一个子空间的同胚。
\item[\textup{c)}] $\OO_X$-模 $\OO_X$ 相对于 $u$ 甚丰沛（4.4.2）。
\item[\textup{c')}] $\OO_X$-模 $\OO_X$ 丰沛（4.5.1）。
\item[\textup{d)}] 当 $f$ 遍历 $A$ 时，$X_f$ 构成 $X$ 的拓扑的一组基。
\item[\textup{d')}] 当 $f$ 遍历 $A$ 时，其中仿射的那些 $X_f$ 覆盖 $X$。
\item[\textup{e)}] 每个拟凝聚 $\OO_X$-模都由其在 $X$ 上的截面生成。
\item[\textup{e')}] $\OO_X$ 的每个有限型拟凝聚理想层都由其在 $X$ 上的截面生成。
\end{enumerate}
\begin{proof}
显然 b) 推出 a)。考虑命题前的说明，将（4.4.4）的准则 b) 应用于恒等态射，得 a) 推出 c)。另一方面，c) 推出 c')（4.5.10，i）；而 c')、b)、b') 由（4.5.2）的准则 b)、b') 等价。最后，由（4.5.2）的准则 a)、a')、c) 及（4.5.5）的准则 d'')，c') 分别等价于 d)、d')、e)、e')。
\end{proof}
还应注意，沿用上述记号，其中仿射的 $X_f$ 已构成 $X$ 的拓扑的一组基，且典范态射 $u$ 占优（4.5.2）。

\egastatement{5.1.3}{推论}
设 $X$ 是拟紧预概形。若存在从 $X$ 到仿射概形 $Y$ 的态射 $v:X\to Y$，使其在底空间上为 $X$ 到 $Y$ 的一个开子空间的同胚，则 $X$ 拟仿射。
\begin{proof}
存在 $\OO_Y$ 的一族全局截面 $(g_\alpha)$，使 $D(g_\alpha)$ 构成 $v(X)$ 的拓扑的一组基。若 $v=(\psi,\theta)$，令 $f_\alpha=\theta(g_\alpha)$，则
$X_{f_\alpha}=\psi^{-1}(D(g_\alpha))$（I，2.2.4.1）。故 $X_{f_\alpha}$ 构成 $X$ 的拓扑的一组基，（5.1.2）的准则 d) 成立。
\end{proof}

\egastatement{5.1.4}{推论}
若 $X$ 是拟仿射概形，则每个可逆 $\OO_X$-模都相对于典范态射甚丰沛，因而丰沛。
\begin{proof}
这样的模 $\LL$ 由其全局截面生成（5.1.2，e），故 $\LL\otimes\OO_X=\LL$ 甚丰沛（4.4.8）。
\end{proof}

\egastatement{5.1.5}{推论}
设 $X$ 是拟紧预概形。若存在可逆 $\OO_X$-模 $\LL$，使 $\LL$ 和 $\LL^{-1}$ 都丰沛，则 $X$ 是拟仿射概形。
\begin{proof}
此时 $\OO_X=\LL\otimes\LL^{-1}$ 丰沛（4.5.7）。
\end{proof}

\egastatement{5.1.6}{命题}
设 $f:X\to Y$ 是拟紧态射。下列条件等价：
\begin{enumerate}[label=\textup{a)}]
\item $f$ 拟仿射。
\item[\textup{b)}] $\OO_Y$-代数 $f_*(\OO_X)=\A(X)$ 拟凝聚，且与恒等同态 $\A(X)\to\A(X)$ 相对应的典范态射 $X\to\Spec(\A(X))$（1.2.7）为开浸入。
\item[\textup{b')}] $\A(X)$ 拟凝聚，且典范态射 $X\to\Spec(\A(X))$ 在底空间上为 $X$ 到 $\Spec(\A(X))$ 的一个子空间的同胚。
\item[\textup{c)}] $\OO_X$ 相对于 $f$ 甚丰沛。
\item[\textup{c')}] $\OO_X$ 相对于 $f$ 丰沛。
\item[\textup{d)}] $f$ 分离，且对每个拟凝聚 $\OO_X$-模 $\F$，典范同态
$\sigma:f^*(f_*(\F))\to\F$（0，4.4.3）满射。
\end{enumerate}
此外，当 $f$ 拟仿射时，每个可逆 $\OO_X$-模 $\LL$ 都相对于 $f$ 甚丰沛。
\begin{proof}
a) 与 c') 的等价性，由 $f$-丰沛性的关于 $Y$ 的局部性（4.6.4）、定义（5.1.1）及准则（5.1.2，c'）得出。其他性质也都关于 $Y$ 局部，故直接由（5.1.2）和（5.1.4）得出，并注意当 $f$ 分离时，$f_*(\F)$ 拟凝聚（I，9.2.2，a）。
\end{proof}

\egastatement{5.1.7}{推论}
设 $f:X\to Y$ 为拟仿射态射。对 $Y$ 的每个开集 $U$，$f$ 的限制 $f^{-1}(U)\to U$ 拟仿射。

\egastatement{5.1.8}{推论}
设 $Y$ 是仿射概形，$f:X\to Y$ 是拟紧态射。$f$ 拟仿射，当且仅当 $X$ 是拟仿射概形。
\begin{proof}
直接由（5.1.6）和（4.6.6）得出。
\end{proof}

\egastatement{5.1.9}{推论}
设 $Y$ 是拟紧概形，或底空间为诺特空间的预概形，$f:X\to Y$ 是有限型态射。若 $f$ 拟仿射，则 $\A(X)=f_*(\OO_X)$ 有一个有限型拟凝聚子 $\OO_Y$-代数 $\B$（I，9.6.2），使与典范嵌入 $\B\to\A(X)$ 相对应的态射 $X\to\Spec(\B)$（1.2.7）为浸入。此外，$\A(X)$ 的每个包含 $\B$ 的有限型拟凝聚子 $\OO_Y$-代数 $\B'$，都具有同一性质。
\begin{proof}
$\A(X)$ 是其有限型拟凝聚子 $\OO_Y$-代数的归纳极限（I，9.6.5）。考虑 $\Spec(\A(X))$ 与 $\Proj(\A(X)[T])$ 的等同（3.1.7），结论即为（3.8.4）的特例。
\end{proof}

\egastatement{5.1.10}{命题}
\begin{enumerate}[label=\textup{(\roman*)}]
\item 若拟紧态射 $X\to Y$ 在底空间上为 $X$ 到 $Y$ 的一个子空间的同胚，特别地若它为闭浸入，则它拟仿射。
\item 两个拟仿射态射的复合仍拟仿射。
\item 若 $f:X\to Y$ 是拟仿射 $S$-态射，则对基预概形的每个扩张 $S'\to S$，$f_{(S')}:X_{(S')}\to Y_{(S')}$ 拟仿射。
\item 若 $f:X\to Y$、$g:X'\to Y'$ 是拟仿射 $S$-态射，则 $f\times_Sg$ 拟仿射。
\item 若 $f:X\to Y$、$g:Y\to Z$ 的复合 $g\circ f$ 拟仿射，且 $g$ 分离或 $X$ 的底空间局部诺特，则 $f$ 拟仿射。
\item 若 $f$ 拟仿射，则 $f_{\red}$ 拟仿射。
\end{enumerate}
\begin{proof}
由准则（5.1.6，c'），（i）、（iii）、（iv）、（v）、（vi）分别直接由（4.6.13）的（i bis）、（iii）、（iv）、（v）、（vi）得出。为证（ii），可设 $Z$ 仿射，再将（4.6.13，ii）应用于 $\LL=\OO_X$、$\K=\OO_Y$ 即得。
\end{proof}

\egastatement{5.1.11}{注记}
设 $f:X\to Y$、$g:Y\to Z$ 为态射，且 $X\times_ZY$ 局部诺特。图浸入 $\Gamma_f:X\to X\times_ZY$ 拟紧（I，6.3.5），故拟仿射；由（I，5.5.12），在（5.1.10，v）中去掉 $g$ 分离的假设，结论仍成立。

\egastatement{5.1.12}{命题}
设 $f:X\to Y$ 是拟紧态射，$g:X'\to X$ 是拟仿射态射。若 $\LL$ 是 $f$-丰沛 $\OO_X$-模，则 $g^*(\LL)$ 是相对于 $f\circ g$ 丰沛的 $\OO_{X'}$-模。
\begin{proof}
$\OO_{X'}$ 相对于 $g$ 甚丰沛，且问题关于 $Y$ 局部（4.6.4）。因此由（4.6.13，ii），在 $Y$ 仿射时，存在整数 $n$，使
$g^*(\LL^{\otimes n})=(g^*(\LL))^{\otimes n}$ 相对于 $f\circ g$ 丰沛。由（4.6.9），$g^*(\LL)$ 也相对于 $f\circ g$ 丰沛。
\end{proof}

\subsection{Serre 判别准则}
\src{97--99}

\egastatement{5.2.1}{定理（Serre 判别准则）}
设 $X$ 是拟紧概形，或底空间为诺特空间的预概形。下列条件等价：
\begin{enumerate}[label=\textup{a)}]
\item $X$ 是仿射概形。
\item[\textup{b)}] 存在一族元素 $f_\alpha\in A=\Gamma(X,\OO_X)$，使各 $X_{f_\alpha}$ 仿射，且 $f_\alpha$ 在 $A$ 中生成的理想等于 $A$。
\item[\textup{c)}] 函子 $\Gamma(X,\F)$ 在拟凝聚 $\OO_X$-模范畴上关于 $\F$ 正合。换言之，若
\begin{equation}
0\longrightarrow\F'\longrightarrow\F\longrightarrow\F''\longrightarrow0
\tag{$*$}
\end{equation}
为拟凝聚 $\OO_X$-模的正合列，则
\[
0\longrightarrow\Gamma(X,\F')\longrightarrow\Gamma(X,\F)
\longrightarrow\Gamma(X,\F'')\longrightarrow0
\]
正合。
\item[\textup{c')}] 对每个满足 $\F$ 同构于某个有限积 $\OO_X^n$ 的子 $\OO_X$-模的拟凝聚模正合列 \textup{($*$)}，性质 \textup{c)} 成立。
\item[\textup{d)}] 对每个拟凝聚 $\OO_X$-模 $\F$，有 $H^1(X,\F)=0$。
\item[\textup{d')}] 对 $\OO_X$ 的每个拟凝聚理想层 $\J$，有 $H^1(X,\J)=0$。
\end{enumerate}
\begin{proof}
显然 a) 推出 b)。另一方面，b) 推出各 $X_{f_\alpha}$ 覆盖 $X$，因为按假设，截面 $1$ 是各 $f_\alpha$ 的线性组合，且各 $D(f_\alpha)$ 覆盖 $\Spec(A)$。由（4.5.2）的最后一项断言，$X\to\Spec(A)$ 为同构。

已知 a) 推出 c)（I，1.3.11），而 c) 显然推出 c')。现证 c') 推出 b)。首先，c') 推出：对每个闭点 $x\in X$ 及其每个开邻域 $U$，存在 $f\in A$，使 $x\in X_f\subset X-U$。\transnote{底本在这里确实印为 $X-U$；由于 $x\in U$，这一包含式不可能成立。后续构造实际给出的是 $x\in X_f\subset U$。此处保留原句并明确指出这一排印错误。}
事实上，令 $\J$（分别地，$\J'$）为定义底空间为 $X-U$（分别地，$(X-U)\cup\{x\}$）的约化闭子预概形的拟凝聚理想层（I，5.2.1）。显然 $\J'\subset\J$，且 $\J''=\J/\J'$ 为支集为 $\{x\}$ 的拟凝聚 $\OO_X$-模，并满足 $\J''_x=k(x)$。把 c') 应用于
$0\to\J'\to\J\to\J''\to0$，可知 $\Gamma(X,\J)\to\Gamma(X,\J'')$ 满射。故 $\J''$ 的在 $x$ 处的芽为 $1_x$ 的截面，是某个
$f\in\Gamma(X,\J)\subset\Gamma(X,\OO_X)$ 的像。按定义，$f(x)=1_x$，且在 $X-U$ 上 $f(y)=0$，从而得到所述断言。

若 $U$ 仿射，则 $X_f$ 也仿射（I，1.3.6）。因此所有仿射的 $X_f$（$f\in A$）之并是包含 $X$ 的全部闭点的开集 $Z$。由于 $X$ 为拟紧 Kolmogoroff 空间，必有 $Z=X$（0，2.1.3）。$X$ 拟紧，故存在有限多个元素 $f_i\in A$（$1\leq i\leq n$），使 $X_{f_i}$ 仿射且覆盖 $X$。考察由这些截面定义的同态 $\OO_X^n\to\OO_X$（0，5.1.1）。每个 $x\in X$ 处至少有一个 $(f_i)_x$ 可逆，故此同态满射，从而有正合列
$0\to\R\to\OO_X^n\to\OO_X\to0$，其中 $\R$ 为拟凝聚子模。由 c')，对应的同态
$\Gamma(X,\OO_X^n)\to\Gamma(X,\OO_X)$ 满射，b) 得证。

最后，a) 推出 d)（I，5.1.9.2），d) 显然推出 d')。只剩证明 d') 推出 c')。若 $\F$ 为 $\OO_X^n$ 的拟凝聚子模，则滤过
$0\subset\OO_X\subset\OO_X^2\subset\cdots\subset\OO_X^n$
在 $\F$ 上诱导滤过 $\F_k=\F\cap\OO_X^k$（$0\leq k\leq n$）。各 $\F_k$ 拟凝聚（I，4.1.1），且 $\F_{k+1}/\F_k$ 同构于
$\OO_X^{k+1}/\OO_X^k\simeq\OO_X$ 的一个拟凝聚子模，即 $\OO_X$ 的拟凝聚理想层。因此 d') 推出 $H^1(X,\F_{k+1}/\F_k)=0$。上同调正合列
\[
H^1(X,\F_k)\longrightarrow H^1(X,\F_{k+1})
\longrightarrow H^1(X,\F_{k+1}/\F_k)=0
\]
使我们可对 $k$ 归纳，得到每个 $H^1(X,\F_k)=0$，证毕。
\end{proof}

\egastatement{5.2.1.1}{注记}
当 $X$ 为诺特预概形时，在条件 c') 和 d') 中可将“拟凝聚”替换为“凝聚”。因为在 c') 推出 b) 的证明中，$\J,\J'$ 此时都是凝聚理想层；又因为凝聚模的每个拟凝聚子模都是凝聚的（I，6.1.1），故得结论。

\NewNote{\textbf{作者勘误提示：}原作者指出本陈述的局部性论证有缺陷，并在勘误中给出增加假设后的陈述。 参见\hyperlink{errIII:19}{《勘误与增补（清单二）》第19条}。}
\egastatement{5.2.2}{推论}
设 $f:X\to Y$ 为分离拟紧态射。下列条件等价：
\begin{enumerate}[label=\textup{a)}]
\item $f$ 为仿射态射。
\item[\textup{b)}] 函子 $f_*$ 在拟凝聚 $\OO_X$-模范畴上正合。
\item[\textup{c)}] 对每个拟凝聚 $\OO_X$-模 $\F$，有 $R^1f_*(\F)=0$。
\item[\textup{c')}] 对 $\OO_X$ 的每个拟凝聚理想层 $\J$，有 $R^1f_*(\J)=0$。
\end{enumerate}
\begin{proof}
由函子 $R^1f_*$ 的定义（T，3.7.3），所有这些条件都关于 $Y$ 局部，故可设 $Y$ 仿射。若 $f$ 仿射，则 $X$ 仿射，b) 就是（I，1.6.4）。反之，为证 b) 推出 a)，注意对每个拟凝聚 $\OO_X$-模 $\F$，$f_*(\F)$ 为拟凝聚 $\OO_Y$-模（I，9.2.2，a）。按假设，$f_*(\F)$ 关于 $\F$ 正合；由于 $Y$ 仿射，$\Gamma(Y,\G)$ 在拟凝聚 $\OO_Y$-模范畴上正合（I，1.3.11）。所以
$\Gamma(Y,f_*(\F))=\Gamma(X,\F)$ 关于 $\F$ 正合，由（5.2.1，c）即得 a)。

若 $f$ 仿射，则对每个仿射开集 $U\subset Y$，$f^{-1}(U)$ 仿射（1.3.2），从而
$H^1(f^{-1}(U),\F)=0$（5.2.1，d），这按定义推出 $R^1f_*(\F)=0$。最后假设 c') 成立。Leray 谱序列低次项给出的正合列（G，II，4.17.1；I，4.5.1）特别给出
\[
0\longrightarrow H^1(Y,f_*(\J))\longrightarrow H^1(X,\J)
\longrightarrow H^0(Y,R^1f_*(\J)).
\]
$Y$ 仿射且 $f_*(\J)$ 拟凝聚（I，9.2.2，a），故 $H^1(Y,f_*(\J))=0$（5.2.1）；由 c')，$H^1(X,\J)=0$。再由（5.2.1），$X$ 为仿射概形。
\end{proof}

\egastatement{5.2.3}{推论}
若 $f:X\to Y$ 为仿射态射，则对每个拟凝聚 $\OO_X$-模 $\F$，典范同态
$H^1(Y,f_*(\F))\to H^1(X,\F)$ 为双射。
\begin{proof}
Leray 谱序列低次项给出正合列
\[
0\longrightarrow H^1(Y,f_*(\F))\longrightarrow H^1(X,\F)
\longrightarrow H^0(Y,R^1f_*(\F)).
\]
由（5.2.2）即得。
\end{proof}

\egastatement{5.2.4}{注记}
在第三章 §1 中，我们将证明：若 $X$ 仿射，则对每个 $i>0$ 和每个拟凝聚 $\OO_X$-模 $\F$，有 $H^i(X,\F)=0$。

\subsection{拟射影态射}
\src{99--100}

\egastatement{5.3.1}{定义}
称态射 $f:X\to Y$ 为\emph{拟射影态射}，或称通过 $f$ 看成 $Y$-预概形的 $X$ \emph{在 $Y$ 上拟射影}，也称 $X$ 为\emph{拟射影 $Y$-概形}，若 $f$ 为有限型，且存在 $f$-丰沛的可逆 $\OO_X$-模。

应注意，这个概念\emph{不是关于 $Y$ 的局部概念}：Nagata 的反例 [26] 和 Hironaka 的反例表明，即使 $X,Y$ 为某个代数闭域上的非奇异代数概形，也可能 $Y$ 的每一点都有仿射邻域 $U$，使 $f^{-1}(U)$ 在 $U$ 上拟射影，而 $f$ 本身并不拟射影。

拟射影态射必然分离（4.6.1）。当 $Y$ 拟紧时，称 $f$ 拟射影，等价于称 $f$ 为有限型且存在相对于 $f$ 甚丰沛的 $\OO_X$-模（4.6.2；4.6.11）。此外：

\egastatement{5.3.2}{命题}
设 $Y$ 是拟紧概形，或底空间为诺特空间的预概形；$X$ 是 $Y$-预概形。下列条件等价：
\begin{enumerate}[label=\alph*)]
\item $X$ 是拟射影 $Y$-概形。
\item $X$ 在 $Y$ 上为有限型，且存在有限型拟凝聚 $\OO_Y$-模 $\E$，使 $X$ 与 $\Pb(\E)$ 的一个子预概形 $Y$-同构。
\item $X$ 在 $Y$ 上为有限型，且存在拟凝聚分次 $\OO_Y$-代数 $\SSS$，使 $\SSS_1$ 为有限型并生成 $\SSS$，而 $X$ 与 $\Proj(\SSS)$ 在某个处处稠密开集上诱导的子预概形 $Y$-同构。
\end{enumerate}
\begin{proof}
由前面的说明及（4.4.3）、（4.4.6）、（4.4.7）立即得出。
\end{proof}
当 $Y$ 为诺特预概形时，可在（5.3.2）的条件 b)、c) 中去掉 $X$ 在 $Y$ 上为有限型的假设，因为此时这一条件自动成立（I，6.3.5）。

\egastatement{5.3.3}{推论}
设 $Y$ 是存在丰沛 $\OO_Y$-模 $\LL$ 的拟紧概形（4.5.3）。$Y$-概形 $X$ 拟射影，当且仅当它在 $Y$ 上为有限型，且同构于某个形如 $\Pb_Y^r$ 的射影丛的子 $Y$-概形。
\begin{proof}
若 $\E$ 是有限型拟凝聚 $\OO_Y$-模，则 $\E$ 同构于某个
$(\LL^{\otimes(-n)})^{\oplus k}$ 的商（4.5.5），故 $\Pb(\E)$ 同构于 $\Pb_Y^{k-1}$ 的一个闭子概形（4.1.2；4.1.4）。
\end{proof}

\egastatement{5.3.4}{命题}
\begin{enumerate}[label=\textup{(\roman*)}]
\item 有限型拟仿射态射是拟射影的；特别地，拟紧浸入以及有限型仿射态射都是拟射影的。
\item 若 $f:X\to Y$、$g:Y\to Z$ 拟射影且 $Z$ 拟紧，则 $g\circ f$ 拟射影。
\item 若 $f:X\to Y$ 是拟射影 $S$-态射，则对基预概形的每个扩张 $S'\to S$，$f_{(S')}:X_{(S')}\to Y_{(S')}$ 拟射影。
\item 若 $f:X\to Y$、$g:X'\to Y'$ 是拟射影 $S$-态射，则 $f\times_Sg$ 拟射影。
\item 若 $f:X\to Y$、$g:Y\to Z$ 的复合 $g\circ f$ 拟射影，且 $g$ 分离或 $X$ 局部诺特，则 $f$ 拟射影。
\item 若 $f$ 拟射影，则 $f_{\red}$ 拟射影。
\end{enumerate}
\begin{proof}
（i）由（5.1.6）和（5.1.10，i）得出。其余断言都是定义（5.3.1）、有限型态射的性质（I，6.3.4）及（4.6.13）的直接推论。
\end{proof}

\egastatement{5.3.5}{注记}
可能出现 $f_{\red}$ 拟射影而 $f$ 并不拟射影的情形，即使假设 $Y$ 是一个有限秩 $\mathbf C$-代数的谱，且 $f$ 固有，也仍如此。

\egastatement{5.3.6}{推论}
若 $X,X'$ 是拟射影 $Y$-概形，则 $X\amalg X'$ 是拟射影 $Y$-概形。
\begin{proof}
由（4.6.18）得出。\transnote{原文引（4.6.18）；直接适用的关于有限和与相对丰沛性的结果为（4.6.17）。}
\end{proof}

\subsection{固有态射与普遍闭态射}
\src{100--103}

\egastatement{5.4.1}{定义}
称预概形态射 $f:X\to Y$ 为\emph{固有的}，若它满足以下两条件：
\begin{enumerate}[label=\alph*)]
\item $f$ 分离且为有限型。
\item 对每个预概形 $Y'$ 及每个态射 $Y'\to Y$，投影
$f_{(Y')}:X\times_Y Y'\to Y'$ 为闭态射（I，2.2.6）。
\end{enumerate}
此时，也称以 $f$ 为结构态射的 $Y$-预概形 $X$ \emph{在 $Y$ 上固有}。

条件 a)、b) 都关于 $Y$ 局部。要验证 $X\times_Y Y'$ 的闭子集 $Z$ 在投影
$q:X\times_Y Y'\to Y'$ 下的像为闭集，只须对 $Y'$ 的每个仿射开集 $U'$ 验证 $q(Z)\cap U'$ 在 $U'$ 中闭。由于
$q(Z)\cap U'=q(Z\cap q^{-1}(U'))$，而 $q^{-1}(U')$ 等同于 $X\times_Y U'$（I，4.4.1），可见验证 b) 时只须考虑 $Y'$ 为仿射概形的情形。后面将看到，若 $Y$ 局部诺特，甚至只须对在 $Y$ 上为有限型的 $Y'$ 验证 b)（5.6.3）。显然每个固有态射都是闭态射。

\egastatement{5.4.2}{命题}
\begin{enumerate}[label=\textup{(\roman*)}]
\item 闭浸入是固有态射。
\item 两个固有态射的复合是固有态射。
\item 若 $X,Y$ 是 $S$-预概形，$f:X\to Y$ 是固有 $S$-态射，则对基预概形的每个扩张 $S'\to S$，$f_{(S')}:X_{(S')}\to Y_{(S')}$ 固有。
\item 若 $f:X\to Y$、$g:X'\to Y'$ 是固有 $S$-态射，则
$f\times_Sg:X\times_S X'\to Y\times_S Y'$ 固有。\transnote{底本此处将乘积态射的源、靶印成 $X\times_S Y\to X'\times_S Y'$；按已给出的 $f,g$ 的类型，应为正文所列的乘积。}
\end{enumerate}
\begin{proof}
只须证明（i）、（ii）、（iii）（I，3.5.1）。在三种情形中，（5.4.1）的条件 a) 都由已有结果（I，5.5.1；6.3.4）得出。只须验证 b)。（i）中这是直接的：若 $X\to Y$ 为闭浸入，则
$X\times_Y Y'\to Y\times_Y Y'=Y'$ 也是闭浸入（I，4.3.2；3.3.3）。

为证（ii），考虑固有态射 $X\to Y$、$Y\to Z$ 及态射 $Z'\to Z$。可以写成
$X\times_Z Z'=X\times_Y(Y\times_Z Z')$（I，3.3.9.1），故投影分解为
\[
X\times_Y(Y\times_Z Z')\longrightarrow Y\times_Z Z'\longrightarrow Z'.
\]
考虑开头的说明，（ii）便由两个闭态射的复合仍闭这一事实得出。

最后，对每个 $S'\to S$，$X_{(S')}$ 等同于 $X\times_Y Y_{(S')}$（I，3.3.11）。对每个 $Z\to Y_{(S')}$，有
\[
X_{(S')}\times_{Y_{(S')}}Z
=(X\times_Y Y_{(S')})\times_{Y_{(S')}}Z
=X\times_Y Z.
\]
由假设，$X\times_Y Z\to Z$ 闭，故（iii）得证。
\end{proof}

\egastatement{5.4.3}{推论}
设 $f:X\to Y$、$g:Y\to Z$ 为态射，且 $g\circ f$ 固有。
\begin{enumerate}[label=\textup{(\roman*)}]
\item 若 $g$ 分离，则 $f$ 固有。
\item 若 $g$ 分离且为有限型，且 $f$ 满射，则 $g$ 固有。
\end{enumerate}
\begin{proof}
（i）由（5.4.2）及一般方法（I，5.5.12）得出。为证（ii），只须验证定义（5.4.1）的条件 b)。对每个态射 $Z'\to Z$，下图交换（I，3.2.1）：
\[
\begin{tikzcd}[column sep=small]
X\times_Z Z' \arrow[rr,"f\times1_{Z'}"] \arrow[dr,"p"'] &&Y\times_Z Z'\arrow[dl,"p'"]\\
&Z'&
\end{tikzcd}
\]
其中 $p,p'$ 为投影。由于 $f$ 满射，$f\times1_{Z'}$ 满射（I，3.5.2）；又按假设，$p$ 为闭态射。$Y\times_Z Z'$ 的每个闭子集 $F$ 都是某个闭子集 $E\subset X\times_Z Z'$ 在 $f\times1_{Z'}$ 下的像，于是
$p'(F)=p(E)$ 在 $Z'$ 中闭，结论得证。
\end{proof}

\egastatement{5.4.4}{推论}
若预概形 $X$ 在 $Y$ 上固有，$\SSS$ 是拟凝聚 $\OO_Y$-代数，则每个 $Y$-态射 $f:X\to\Proj(\SSS)$ 都固有，因而闭。
\begin{proof}
结构态射 $p:\Proj(\SSS)\to Y$ 分离，而 $p\circ f$ 按假设固有。
\end{proof}

\egastatement{5.4.5}{推论}
设 $f:X\to Y$ 是分离有限型态射。设 $(X_i)_{1\leq i\leq n}$、$(Y_i)_{1\leq i\leq n}$ 分别为 $X,Y$ 的有限族闭子预概形，$j_i:X_i\to X$、$h_i:Y_i\to Y$ 为典范嵌入。假设 $X$ 的底空间为各 $X_i$ 之并，且对每个 $i$ 存在态射 $f_i:X_i\to Y_i$ 使下图交换：
\[
\begin{tikzcd}
X_i\arrow[r,"f_i"]\arrow[d,"j_i"']&Y_i\arrow[d,"h_i"]\\
X\arrow[r,"f"']&Y
\end{tikzcd}
\]
则 $f$ 固有，当且仅当每个 $f_i$ 固有。
\begin{proof}
若 $f$ 固有，则 $j_i$ 为闭浸入，故 $f\circ j_i$ 固有（5.4.2）。由于 $h_i$ 为闭浸入，因而分离，由（5.4.3），$f_i$ 固有。
反之，设每个 $f_i$ 固有。取各 $X_i$ 的和预概形 $Z$，令 $u:Z\to X$ 在每个 $X_i$ 上为 $j_i$。$f\circ u$ 在 $X_i$ 上的限制为
$f\circ j_i=h_i\circ f_i$，故因 $h_i,f_i$ 固有而固有（5.4.2）；于是由定义（5.4.1），$u$ 固有。\transnote{此处原文写 $u$ 固有；为了随后应用（5.4.3）推出 $f$ 固有，前面的论证所需且实际得到的是 $f\circ u$ 固有。$u$ 本身也固有，但仅这一事实不足以推出所需结论。}
又按假设，$u$ 满射，故由（5.4.3），$f$ 固有。
\end{proof}

\egastatement{5.4.6}{推论}
设 $f:X\to Y$ 是分离有限型态射。$f$ 固有，当且仅当 $f_{\red}:X_{\red}\to Y_{\red}$ 固有。
\begin{proof}
这是（5.4.5）中取 $n=1$、$X_1=X_{\red}$、$Y_1=Y_{\red}$ 的特例（I，5.1.5）。
\end{proof}

\egapar{5.4.7}
若 $X,Y$ 为诺特预概形，$f:X\to Y$ 为分离有限型态射，则验证 $f$ 固有，可归结为验证整预概形间的占优态射固有。事实上，令 $X_i$（$1\leq i\leq n$）为 $X$ 的有限多个不可约分支，并仍以 $X_i$ 表示具有该底空间的唯一约化闭子预概形（I，5.2.1）。令 $Y_i$ 为底空间是 $\overline{f(X_i)}$ 的唯一约化闭子预概形。若 $\psi_i:X_i\to X$、$\varphi_i:Y_i\to Y$ 为嵌入，则
$f\circ\psi_i=\varphi_i\circ f_i$，其中 $f_i:X_i\to Y_i$ 占优（I，5.2.2）。于是满足（5.4.5）的条件，$f$ 固有当且仅当各 $f_i$ 固有。

\egastatement{5.4.8}{推论}
设 $X,Y$ 为在 $S$ 上分离且为有限型的 $S$-预概形，$f:X\to Y$ 是 $S$-态射。$f$ 固有，当且仅当对每个 $S$-预概形 $S'$，态射
$f\times_S1_{S'}:X\times_S S'\to Y\times_S S'$ 为闭态射。
\begin{proof}
若 $g:X\to S$、$h:Y\to S$ 为结构态射，则按定义 $g=h\circ f$，故 $f$ 分离且为有限型（I，5.5.1；6.3.4）。若 $f$ 固有，则 $f\times_S1_{S'}$ 固有（5.4.2），从而闭。

反之，假设所述条件成立，取 $Y$-预概形 $Y'$。它也是 $S$-预概形；由于 $Y\to S$ 分离，$X\times_Y Y'$ 等同于 $X\times_S Y'$ 的一个闭子预概形（I，5.4.2）。交换图
\[
\begin{tikzcd}[column sep=large]
X\times_Y Y'\arrow[r,"f\times1_{Y'}"]\arrow[d,hook]&Y\times_Y Y'=Y'\arrow[d,hook]\\
X\times_S Y'\arrow[r,"f\times_S1_{Y'}"']&Y\times_S Y'
\end{tikzcd}
\]
中的竖箭头为闭浸入。因此，若下面的横箭头为闭态射，上面的横箭头也为闭态射。
\end{proof}

\egastatement{5.4.9}{注记}
若态射 $f:X\to Y$ 满足定义（5.4.1）的条件 b)，称它为\emph{普遍闭态射}。读者可注意到：在（5.4.2）至（5.4.8）中，可以处处将“固有”替换为“普遍闭”，结论仍成立；而且在（5.4.3）、（5.4.5）、（5.4.6）、（5.4.8）的假设中，可以去掉有限性条件。

\egapar{5.4.10}
设 $f:X\to Y$ 是有限型态射。称 $X$ 的一个闭子集 $Z$ \emph{在 $Y$ 上固有}，或 \emph{$Y$-固有}，或\emph{相对于 $f$ 固有}，若 $f$ 在某个底空间为 $Z$ 的闭子预概形上的限制（I，5.2.1）固有。由于这样的限制分离，由（5.4.6）及（I，5.5.1，vi），此性质与具有底空间 $Z$ 的闭子预概形的选择无关。

若 $g:X'\to X$ 是固有态射，则 $g^{-1}(Z)$ 是 $X'$ 的固有子集。事实上，取底空间为 $Z$ 的子预概形 $T$，由（5.4.2，iii），$g$ 在闭子预概形 $g^{-1}(T)$ 上的限制 $g^{-1}(T)\to T$ 固有；再应用（5.4.2，ii）即可。

另一方面，若 $X''$ 是有限型 $Y$-概形，$u:X\to X''$ 为 $Y$-态射，则 $u(Z)$ 为 $X''$ 的固有子集。令 $T$ 为具有底空间 $Z$ 的约化闭子预概形；$f|T$ 固有，所以 $u|T$ 也固有（5.4.3，i），故 $u(Z)$ 在 $X''$ 中闭。取 $X''$ 中底空间为 $u(Z)$ 的闭子预概形 $T''$（I，5.2.1），使 $u|T$ 分解为
\[
T\xrightarrow{v}T''\xrightarrow{j}X'',
\]
其中 $j$ 为典范嵌入（I，5.2.2）。于是 $v$ 固有且满射（5.4.5）。若 $g$ 是结构态射 $X''\to Y$ 在 $T''$ 上的限制，则 $g$ 分离且为有限型，且 $f|T=g\circ v$；由（5.4.3，ii），$g$ 固有，从而得所求断言。

特别地，若 $Z$ 是 $X$ 的 $Y$-固有子集，则：
\begin{enumerate}[label=\arabic*°]
\item 对 $X$ 的每个闭子预概形 $X'$，$Z\cap X'$ 是 $X'$ 的 $Y$-固有子集。
\item 若 $X$ 是有限型 $Y$-概形 $X''$ 的子预概形，则 $Z$ 也是 $X''$ 的 $Y$-固有子集；特别地，它在 $X''$ 中闭。
\end{enumerate}

\subsection{射影态射}
\src{103--106}

\egastatement{5.5.1}{命题}
设 $X$ 为 $Y$-预概形。下列条件等价：
\begin{enumerate}[label=\alph*)]
\item $X$ 与某个射影丛 $\Pb(\E)$ 的闭子预概形 $Y$-同构，其中 $\E$ 为有限型拟凝聚 $\OO_Y$-模。
\item 存在拟凝聚分次 $\OO_Y$-代数 $\SSS$，使 $\SSS_1$ 为有限型并生成 $\SSS$，且 $X$ 与 $\Proj(\SSS)$ $Y$-同构。
\end{enumerate}
\begin{proof}
由（3.6.2，ii），a) 推出 b)：若 $\J$ 是 $\Sym(\E)$ 的拟凝聚分次理想层，则拟凝聚分次代数
$\SSS=\Sym(\E)/\J$ 由 $\SSS_1$ 生成，而 $\SSS_1$ 是 $\E$ 的典范像，故为有限型 $\OO_Y$-模。将（3.6.2）应用于 $\M\to\SSS_1$ 为恒等同态的情形，便得 b) 推出 a)。\transnote{后一处直接对应的“由对称代数满射得到闭浸入”的叙述为（3.6.4）；这里保留原文引号（3.6.2）。}
\end{proof}

\egastatement{5.5.2}{定义}
称 $Y$-预概形 $X$ \emph{在 $Y$ 上射影}，或称它为\emph{射影 $Y$-概形}，若它满足（5.5.1）的等价条件 a)、b)。若态射 $f:X\to Y$ 使 $X$ 成为射影 $Y$-概形，称 $f$ 为\emph{射影态射}。

显然，若 $f:X\to Y$ 射影，则存在相对于 $f$ 甚丰沛的 $\OO_X$-模（4.4.2）。

\egastatement{5.5.3}{定理}
\begin{enumerate}[label=\textup{(\roman*)}]
\item 每个射影态射都拟射影且固有。
\item 反之，若 $Y$ 是拟紧概形，或底空间为诺特空间的预概形，则每个拟射影且固有的态射 $f:X\to Y$ 都射影。
\end{enumerate}
\begin{proof}
（i）若 $f:X\to Y$ 射影，则显然为有限型且拟射影，特别地分离。另一方面，由（5.5.1，b）和（3.5.3），对每个态射 $Y'\to Y$，
$f\times_Y1_{Y'}:X\times_Y Y'\to Y'$ 也射影。因此，为证 $f$ 普遍闭，只须证明射影态射为闭态射。

问题关于 $Y$ 局部，可设 $Y=\Spec(A)$，从而 $X=\Proj(S)$（5.5.1），其中 $S$ 是由 $S_1$ 中有限多个元素生成的分次 $A$-代数。对每个 $y\in Y$，纤维 $f^{-1}(y)$ 等同于
$\Proj(S)\times_Y\Spec(k(y))$（I，3.6.1），因而等同于
$\Proj(S\otimes_A k(y))$（2.8.10）。所以 $f^{-1}(y)=\varnothing$ 当且仅当 $S\otimes_A k(y)$ 满足（TN）（2.7.4），即对充分大的 $n$，有
$S_n\otimes_A k(y)=0$。由于 $(S_n)^\sim$ 为有限型 $\OO_Y$-模，由 Nakayama 引理，此条件等价于对充分大的 $n$，有 $(S_n)^\sim_y=0$。若 $\mathfrak a_n$ 是 $A$-模 $S_n$ 在 $A$ 中的零化子，则这也等价于：对充分大的 $n$，有 $\mathfrak a_n\not\subset\mathfrak j_y$（0，1.7.4）。由假设 $S_nS_1=S_{n+1}$，有 $\mathfrak a_n\subset\mathfrak a_{n+1}$。令 $\mathfrak a=\bigcup_n\mathfrak a_n$，便得
$f(X)=V(\mathfrak a)$，故 $f(X)$ 在 $Y$ 中闭。

现取 $X$ 的任意闭子集 $X'$。存在底空间为 $X'$ 的闭子预概形（I，5.2.1）；显然其复合态射 $X'\to X\to Y$ 射影（5.5.1，a），故 $f(X')$ 在 $Y$ 中闭。

（ii）关于 $Y$ 的假设以及 $f$ 拟射影，推出存在有限型拟凝聚 $\OO_Y$-模 $\E$ 及 $Y$-浸入 $j:X\to\Pb(\E)$（5.3.2）。由于 $f$ 固有，由（5.4.4），$j$ 为闭浸入，故 $f$ 射影。
\end{proof}

\egastatement{5.5.4}{注记}
\begin{enumerate}[label=\textup{(\roman*)}]
\item 设态射 $f:X\to Y$ 满足：1° $f$ 固有；2° 存在相对于 $f$ 甚丰沛的 $\OO_X$-模 $\LL$；3° 拟凝聚 $\OO_Y$-模 $\E=f_*(\LL)$ 为有限型。则 $f$ 为射影态射。事实上，由（4.4.4），存在 $Y$-浸入 $r:X\to\Pb(\E)$；由于 $f$ 固有，$r$ 为闭浸入（5.4.4）。第三章 §3 将表明，当 $Y$ 局部诺特时，条件3°是前两条件的推论；故此时条件1°、2°刻画射影态射。若 $Y$ 还拟紧，则可将条件2°替换为存在 $f$-丰沛 $\OO_X$-模（4.6.11）。
\item 设 $Y$ 是存在丰沛 $\OO_Y$-模的拟紧概形。$Y$-概形 $X$ 射影，当且仅当它与某个形如 $\Pb_Y^r$ 的射影丛的\emph{闭}子 $Y$-概形 $Y$-同构。充分性显然。反之，若 $X$ 在 $Y$ 上射影，则它拟射影，故存在 $Y$-浸入 $j:X\to\Pb_Y^r$（5.3.3）；由（5.4.4）和（5.5.3），此浸入闭。
\item （5.5.3）的论证表明：对每个预概形 $Y$ 及每个整数 $r\geq0$，结构态射 $\Pb_Y^r\to Y$ 满射。因为令 $\SSS=\Sym_{\OO_Y}(\OO_Y^{r+1})$，显然有
$\SSS^y=\Sym_{k(y)}(k(y)^{r+1})$（1.7.3），故对每个 $y\in Y$ 和每个 $n\geq0$，$(\SSS^y)_n\ne0$。
\item Nagata 的例子 [26] 表明，存在不拟射影的固有态射。
\end{enumerate}

\egastatement{5.5.5}{命题}
\begin{enumerate}[label=\textup{(\roman*)}]
\item 闭浸入为射影态射。
\item 若 $f:X\to Y$、$g:Y\to Z$ 为射影态射，且 $Z$ 为拟紧概形或底空间为诺特空间的预概形，则 $g\circ f$ 射影。
\item 若 $f:X\to Y$ 为射影 $S$-态射，则对基预概形的每个扩张 $S'\to S$，$f_{(S')}:X_{(S')}\to Y_{(S')}$ 射影。
\item 若 $f:X\to Y$、$g:X'\to Y'$ 为射影 $S$-态射，则 $f\times_Sg$ 射影。
\item 若 $g\circ f$ 射影且 $g$ 分离，则 $f$ 射影。
\item 若 $f$ 射影，则 $f_{\red}$ 射影。
\end{enumerate}
\begin{proof}
（i）直接由（3.1.7）得出。由于（ii）中对 $Z$ 加了限制，须分别证明（iii）、（iv）（参见 I，3.5.1）。为证（iii），可归结到 $S=Y$（I，3.3.11），再由（5.5.1，b）和（3.5.3）即得。

为证（iv），立即可归结到 $X=\Pb(\E)$、$X'=\Pb(\E')$ 的情形，其中 $\E$（分别地，$\E'$）是有限型拟凝聚 $\OO_Y$-模（分别地，$\OO_{Y'}$-模）。令 $p,p'$ 为 $T=Y\times_S Y'$ 的投影。由（4.1.3.1），
$\Pb(p^*(\E))=\Pb(\E)\times_Y T$，$\Pb(p'^*(\E'))=\Pb(\E')\times_{Y'}T$；故
\[
\begin{aligned}
\Pb(p^*\E)\times_T\Pb(p'^*\E')
&=(\Pb(\E)\times_Y T)\times_T(T\times_{Y'}\Pb(\E'))\\
&=\Pb(\E)\times_Y(T\times_{Y'}\Pb(\E'))\\
&=\Pb(\E)\times_S\Pb(\E').
\end{aligned}
\]
最后一步把 $T$ 替换为 $Y\times_S Y'$ 并使用（I，3.3.9.1）。$p^*(\E),p'^*(\E')$ 在 $T$ 上为有限型（0，5.2.4），它们的张量积也如此。又由（4.3.3），
$\Pb(p^*\E)\times_T\Pb(p'^*\E')$ 等同于
$\Pb(p^*(\E)\otimes_{\OO_T}p'^*(\E'))$ 的闭子预概形，故（iv）得证。

为得（v）、（vi），可应用（I，5.5.13），因为由（5.5.1，a），射影 $Y$-概形的每个闭子预概形仍是射影 $Y$-概形。最后，（ii）在关于 $Z$ 的假设下，由（5.5.3）、（5.3.4，ii）和（5.4.2，ii）得出。
\end{proof}

\egastatement{5.5.6}{命题}
若 $X,X'$ 为射影 $Y$-概形，则 $X\amalg X'$ 为射影 $Y$-概形。
\begin{proof}
由（5.5.2）和（4.3.6）显然得出。
\end{proof}

\egastatement{5.5.7}{命题}
设 $X$ 是射影 $Y$-概形，$\LL$ 是 $Y$-丰沛 $\OO_X$-模。对 $\LL$ 的每个全局截面 $f$，$X_f$ 在 $Y$ 上仿射。
\begin{proof}
问题关于 $Y$ 局部，故可设 $Y=\Spec(A)$。又 $X_{f^{\otimes n}}=X_f$，因而把 $\LL$ 替换为适当的 $\LL^{\otimes n}$，可设 $\LL$ 相对于结构态射 $q:X\to Y$ 甚丰沛（4.6.11）。典范同态
$\sigma:q^*(q_*(\LL))\to\LL$ 满射，且对应的态射
$r=r_{\LL,\sigma}:X\to P=\Pb(q_*(\LL))$ 是浸入，满足 $\LL=r^*(\OO_P(1))$（4.4.4）。由于 $X$ 在 $Y$ 上固有，$r$ 还是闭浸入（5.4.4）。

按定义，$f\in\Gamma(Y,q_*(\LL))$，且 $\sigma^\flat$ 是 $q_*(\LL)$ 的恒等同态。由（3.7.3.1），$X_f=r^{-1}(D_+(f))$；故 $X_f$ 是仿射概形 $D_+(f)$ 的闭子预概形，因而仿射。
\end{proof}
特别取 $Y=X$，并考虑（4.6.13，i），得到如下推论；其直接证明也是立即的。

\egastatement{5.5.8}{推论}
设 $X$ 为预概形，$\LL$ 为可逆 $\OO_X$-模。对每个全局截面 $f$，$X_f$ 在 $X$ 上仿射；特别地，若 $X$ 仿射，则 $X_f$ 为仿射概形。

\subsection{Chow 引理}
\src{106--110}

\egastatement{5.6.1}{定理（Chow 引理）}
设 $S$ 是预概形，$X$ 是有限型 $S$-概形。假设下列条件之一成立：
\begin{enumerate}[label=\alph*)]
\item $S$ 诺特。
\item $S$ 是拟紧概形，且 $X$ 只有有限多个不可约分支。
\end{enumerate}
则：
\begin{enumerate}[label=\textup{(\roman*)}]
\item 存在在 $S$ 上拟射影的 $S$-概形 $X'$，以及射影且满射的 $S$-态射 $f:X'\to X$。
\item 可选择 $X',f$，使存在开集 $U\subset X$，满足 $U'=f^{-1}(U)$ 在 $X'$ 中稠密，且 $f|U'$ 为 $U'\simeq U$ 的同构。
\item 若 $X$ 约化（分别地，不可约、整），可令 $X'$ 也约化（分别地，不可约、整）。
\end{enumerate}
\begin{proof}
证明分若干步。

\textbf{A）} 首先可归结到 $X$ 不可约的情形。在假设 a) 下，$X$ 诺特，所以在两种假设下，$X$ 的不可约分支 $X_i$ 都只有有限多个。若定理对底空间为各 $X_i$ 的约化闭子预概形已证得，令 $X'_i$ 及 $f_i:X'_i\to X_i$ 为相应的预概形和态射。取 $X'$ 为各 $X'_i$ 的和，并令 $f:X'\to X$ 在 $X'_i$ 上的限制为 $j_i\circ f_i$，其中 $j_i:X_i\to X$ 是典范嵌入，则满足要求。\transnote{这一步在原文中取的是“约化”闭子预概形。若原来的 $X$ 非约化，直接这样替换不能保证在开集 $U$ 上与原 $X$ 的结构层同构，故关于（ii）的归结不能仅由这句话得到。这里忠实保留原著归结，不暗中把定理的 $X$ 加上约化假设。}
若各 $X'_i$ 约化，显然 $X'$ 约化。取
\[
U=\bigcup_i\left(U_i\cap\left(X-\bigcup_{j\ne i}X_j\right)\right),
\]
便满足（ii）。最后，各 $X'_i$ 在 $S$ 上拟射影，故 $X'$ 也拟射影（5.3.6）；由（5.5.5，i、ii），各 $X'_i\to X$ 射影，故 $f$ 射影（5.5.6），而按定义它显然满射。

\textbf{B）} 现设 $X$ 不可约。结构态射 $r:X\to S$ 为有限型，故存在 $S$ 的有限仿射开覆盖 $(S_i)$，并对每个 $i$ 存在 $r^{-1}(S_i)$ 的有限仿射开覆盖 $(T_{ij})$。态射 $T_{ij}\to S_i$ 仿射且为有限型，因而拟射影（5.3.4，i）。在 a)、b) 中，浸入 $S_i\to S$ 都拟紧，故为拟射影态射（5.3.4，i）；因此 $r|T_{ij}$ 为拟射影态射（5.3.4，ii）。将这些 $T_{ij}$ 记成 $U_k$（$1\leq k\leq n$）。对每个 $k$，存在开浸入 $\varphi_k:U_k\to P_k$，其中 $P_k$ 在 $S$ 上射影（5.3.2；5.5.2）。令 $U=\bigcap_kU_k$。由于 $X$ 不可约且各 $U_k$ 非空，$U$ 非空，因而在 $X$ 中处处稠密。各 $\varphi_k$ 在 $U$ 上的限制定义态射
\[
\varphi:U\longrightarrow P=P_1\times_S P_2\times_S\cdots\times_S P_n,
\]
使下图交换：
\begin{equation}
\begin{tikzcd}
U\arrow[r,"\varphi"]\arrow[d,"j_k"']&P\arrow[d,"p_k"]\\
U_k\arrow[r,"\varphi_k"']&P_k
\end{tikzcd}
\tag{5.6.1.1}
\end{equation}
这里 $j_k:U\to U_k$ 是典范嵌入，$p_k:P\to P_k$ 是典范投影。若 $j:U\to X$ 为典范嵌入，则
$\psi=(j,\varphi)_S:U\to X\times_S P$ 为浸入（I，5.3.14）。在 a) 下，$X\times_S P$ 局部诺特（3.4.1；I，6.3.7、6.3.8）；在 b) 下，$X\times_S P$ 为拟紧概形（I，5.5.1；6.6.4）。两种情形中，与 $\psi$ 相伴且底空间为 $\psi(U)$ 的子预概形 $Z$ 在 $X\times_S P$ 中的闭包 $X'$ 都存在，且 $\psi$ 分解为
\begin{equation}
\psi:U\xrightarrow{\psi'}X'\xrightarrow{h}X\times_S P,
\tag{5.6.1.2}
\end{equation}
其中 $\psi'$ 为开浸入，$h$ 为闭浸入（I，9.5.10）。记 $q_1:X\times_S P\to X$、$q_2:X\times_S P\to P$ 为典范投影，令
\begin{align}
f &:X'\xrightarrow{h}X\times_S P\xrightarrow{q_1}X,
\tag{5.6.1.3}\\
g &:X'\xrightarrow{h}X\times_S P\xrightarrow{q_2}P.
\tag{5.6.1.4}
\end{align}
下面证明 $X',f$ 满足要求。

\textbf{C）} 先证 $f$ 射影且满射，并且 $f|U'$ 为 $U'=f^{-1}(U)$ 到 $U$ 的同构。各 $P_k$ 在 $S$ 上射影，故 $P$ 也射影（5.5.5，iv）；所以 $X\times_S P$ 在 $X$ 上射影（5.5.5，iii），其闭子预概形 $X'$ 也在 $X$ 上射影。另一方面，
$f\circ\psi'=q_1\circ(h\circ\psi')=q_1\circ\psi=j$，故 $f(X')$ 包含 $X$ 的稠密开集 $U$。由于 $f$ 闭（5.5.3），有 $f(X')=X$。

$q_1^{-1}(U)=U\times_S P$ 是 $X\times_S P$ 在一个开集上诱导的预概形；按定义，$U'=h^{-1}(U\times_S P)$ 是 $X'$ 在开集 $U'$ 上的诱导预概形，故它是 $Z$ 相对于 $U\times_S P$ 的闭包（I，9.5.8）。而 $\psi$ 分解为
\[
U\xrightarrow{\Gamma_\varphi}U\times_S P
\xrightarrow{j\times1}X\times_S P.
\]
$P$ 在 $S$ 上分离，故图态射 $\Gamma_\varphi$ 为闭浸入（I，5.4.3），于是 $Z$ 为 $U\times_S P$ 的闭子预概形，因而 $U'=Z$。又 $\psi$ 为浸入，故 $f|U'$ 为到 $U$ 的同构，是 $\psi'$ 的逆同构；按 $X'$ 的定义，$U'$ 在 $X'$ 中稠密。

\textbf{D）} 现证 $g$ 为浸入。由于 $P$ 在 $S$ 上射影，这就说明 $X'$ 在 $S$ 上拟射影。令
\[
\begin{aligned}
V_k&=\varphi_k(U_k)&&\text{（$P_k$ 的开集）},\\
W_k&=p_k^{-1}(V_k)&&\text{（$P$ 的开集）},\\
U'_k&=f^{-1}(U_k),\qquad U''_k=g^{-1}(W_k)
&&\text{（$X'$ 的开集）}.
\end{aligned}
\]
显然各 $U'_k$ 覆盖 $X'$。先证明 $U'_k\subset U''_k$，从而各 $U''_k$ 也覆盖 $X'$。为此只须证明下图交换：
\begin{equation}
\begin{tikzcd}
U'_k\arrow[r,"g|U'_k"]\arrow[d,"f|U'_k"']&P\arrow[d,"p_k"]\\
U_k\arrow[r,"\varphi_k"']&P_k
\end{tikzcd}
\tag{5.6.1.5}
\end{equation}
预概形 $U'_k=h^{-1}(U_k\times_S P)$ 是 $X'$ 在开集 $U'_k$ 上的限制，因而是 $Z=U'\subset U'_k$ 相对于该开集的闭包（I，9.5.8）。由于 $P_k$ 为 $S$-概形，为证（5.6.1.5）交换，只须把图与典范嵌入 $U'\to U'_k$ 复合后验证交换；等价地，利用 $U\simeq U'$，与 $\psi$ 复合后验证（I，9.5.6）。按定义，所得的图正是（5.6.1.1），故断言成立。

各 $W_k$ 因而构成 $g(X')$ 的开覆盖。为证 $g$ 为浸入，只须证明各 $g|U''_k$ 是到 $W_k$ 的浸入（I，4.2.4）。考虑态射
\[
u_k:W_k\xrightarrow{p_k}V_k\xrightarrow{\varphi_k^{-1}}U_k\longrightarrow X.
\]
由于 $X$ 在 $S$ 上分离，图态射 $\Gamma_{u_k}:W_k\to X\times_S W_k$ 为闭浸入（I，5.4.3），故其图
$T_k=\Gamma_{u_k}(W_k)$ 为 $X\times_S W_k$ 的闭子预概形。若证明此子预概形在子预概形的序关系中包含 $U'$，则由（I，9.5.8），它也包含 $X'$ 在开集 $U''_k$ 上诱导的子预概形。\transnote{原文在此后两处把先前记作 $U''_k$ 的开集印作 $X''_k$；此处统一沿用已经定义的 $U''_k$。}
由于 $q_2|T_k$ 为到 $W_k$ 的同构，$g|U''_k$ 就是到 $W_k$ 的浸入，断言便得证。

令 $v_k:U'\to X\times_S W_k$ 为典范嵌入。须证明存在 $w_k:U'\to W_k$，使
$v_k=\Gamma_{u_k}\circ w_k$。按乘积的定义，只须证
$q_1\circ v_k=u_k\circ q_2\circ v_k$（I，3.2.1）；或在右边与同构 $\psi':U\to U'$ 复合，只须证
$q_1\circ\psi=u_k\circ q_2\circ\psi$。因 $q_1\circ\psi=j$、$q_2\circ\psi=\varphi$，此式由（5.6.1.1）的交换性及 $u_k$ 的定义得出。

\textbf{E）} 显然，$U$，因而 $U'$，不可约，所以以上构造中的 $X'$ 也不可约，态射 $f$ 因而双有理（I，2.2.9）。若 $X$ 还约化，则 $U'$ 约化，故 $X'$ 约化（I，9.5.9）。定理证毕。
\end{proof}

\egastatement{5.6.2}{推论}
假设（5.6.1）中的 a)、b) 之一成立。$X$ 在 $S$ 上固有，当且仅当存在在 $S$ 上射影的概形 $X'$ 以及满的 $S$-态射 $f:X'\to X$；此时 $f$ 由（5.5.5，v）为射影态射。并且可选择 $f$，使存在 $X$ 的稠密开集 $U$，令 $f|f^{-1}(U)$ 为 $f^{-1}(U)\simeq U$ 的同构，且 $f^{-1}(U)$ 在 $X'$ 中稠密。若 $X$ 不可约（分别地，约化），可令 $X'$ 也如此；当 $X,X'$ 不可约时，$f$ 双有理。
\begin{proof}
充分性由（5.5.3）及（5.4.3，ii）得出。必要性：沿用（5.6.1）的记号，若 $X$ 在 $S$ 上固有，则 $X'$ 在 $X$ 上射影，故在 $X$ 上固有（5.5.3）；由（5.4.2，ii），$X'$ 在 $S$ 上固有。它又在 $S$ 上拟射影，故由（5.5.3），$X'$ 在 $S$ 上射影。
\end{proof}

\egastatement{5.6.3}{推论}
设 $S$ 是局部诺特预概形，$X$ 是在 $S$ 上为有限型的 $S$-概形，$f_0:X\to S$ 为结构态射。$X$ 在 $S$ 上固有，当且仅当对每个有限型态射 $S'\to S$，
$(f_0)_{(S')}:X_{(S')}\to S'$ 为闭态射。甚至只须对所有形如
$S'=S\otimes_{\mathbf Z}\mathbf Z[T_1,\ldots,T_n]$ 的 $S$-预概形验证这一条件，其中各 $T_i$ 为不定元。
\begin{proof}
必要性显然，只证充分性。问题关于 $S,S'$ 局部（5.4.1），故可设 $S,S'$ 仿射且诺特。由 Chow 引理，存在射影 $S$-概形 $P$、浸入 $j:X'\to P$ 及射影满态射 $f:X'\to X$，使下图交换：
\[
\begin{tikzcd}
X'\arrow[r,"f"]\arrow[d,"j"']&X\arrow[d,"f_0"]\\
P\arrow[r,"r"']&S.
\end{tikzcd}
\]
$P$ 在 $S$ 上为有限型，故第一个假设推出投影 $q_2:X\times_S P\to P$ 为闭态射。浸入 $j$ 是 $q_2$ 与态射
$f\times1:X'\times_S P\to X\times_S P$ 的复合；\transnote{原文此句省略了源为 $X'$ 的图态射。完整的因子链是 $j=q_2\circ(f\times1)\circ\Gamma_j$。由于 $P$ 在 $S$ 上分离，$\Gamma_j$ 为闭浸入，因此所需闭性论证可完成。}
$f$ 射影，故固有（5.5.3），所以 $f\times1$ 闭。由此得 $j$ 为闭浸入，因而固有（5.4.2，i）。结构态射 $r:P\to S$ 射影，因而固有（5.5.3），故
$f_0\circ f=r\circ j$ 固有（5.4.2，ii）。最后，$f$ 满射，所以由（5.4.3），$f_0$ 固有。

为仅利用第二个假设证明结论，只须证第二个假设推出第一个。若 $S'$ 仿射且在 $S=\Spec(A)$ 上为有限型，则
$S'=\Spec(A[c_1,\ldots,c_n])$（I，6.3.3），所以它同构于
$S''=\Spec(A[T_1,\ldots,T_n])$ 的闭子预概形。交换图
\[
\begin{tikzcd}
X\times_S S'\arrow[r,"1_X\times j"]\arrow[d,"(f_0)_{(S')}"']&X\times_S S''\arrow[d,"(f_0)_{(S'')}"]\\
S'\arrow[r,"j"']&S''
\end{tikzcd}
\]
中的 $j$ 和 $1_X\times j$ 为闭浸入（I，4.3.1）。原文接着写道：按假设 $(f_0)_{(S')}$ 闭，故 $(f_0)_{(S'')}$ 也闭。\transnote{这里底本最后一句把 $S'$ 和 $S''$ 的角色颠倒了。所需论证是：第二个假设保证右边的 $(f_0)_{(S'')}$ 闭，由两个水平闭浸入，推出左边的 $(f_0)_{(S')}$ 闭。}
\end{proof}

\section{整态射与有限态射}
\subsection{在另一预概形上为整的预概形}
\src{110--114}

\egastatement{6.1.1}{定义}
设 $X$ 是 $S$-预概形，$f:X\to S$ 是其结构态射。称 $X$ \emph{在 $S$ 上为整}，或称 $f$ 为\emph{整态射}，若存在 $S$ 的仿射开覆盖 $(S_\alpha)$，使每个诱导预概形 $f^{-1}(S_\alpha)$ 为仿射概形，且其环 $B_\alpha$ 是 $S_\alpha$ 的环 $A_\alpha$ 上的整代数。若 $X$ 在 $S$ 上为整且为有限型，则称 $X$ \emph{在 $S$ 上有限}，或称 $f$ 为\emph{有限态射}。
当 $S$ 仿射、环为 $A$ 时，也称“在 $A$ 上为整（分别地，有限）”，以代替“在 $S$ 上为整（分别地，有限）”。

\egapar{6.1.2}
显然，若 $X$ 在 $S$ 上为整，则它在 $S$ 上仿射。在 $S$ 上仿射的预概形 $X$ 在 $S$ 上为整（分别地，有限），当且仅当相伴的拟凝聚 $\OO_S$-代数 $\A(X)$ 满足：存在 $S$ 的仿射开覆盖 $(S_\alpha)$，使每个
$\Gamma(S_\alpha,\A(X))$ 是 $\Gamma(S_\alpha,\OO_S)$ 上的整代数（分别地，整且有限型的代数）。具有此性质的拟凝聚 $\OO_S$-代数称为\emph{在 $\OO_S$ 上为整（分别地，有限）}。因此，给定一个在 $S$ 上为整（分别地，有限）的预概形，等价于给定一个在 $\OO_S$ 上为整（分别地，有限）的拟凝聚 $\OO_S$-代数（1.3.1）。应注意，拟凝聚 $\OO_S$-代数 $\B$ 有限，当且仅当它为有限型 $\OO_S$-模（I，1.3.9）；这也等价于说它为整且有限型的 $\OO_S$-代数，因为环 $A$ 上的整且有限型代数是有限型 $A$-模。

\egastatement{6.1.3}{命题}
设 $S$ 是局部诺特预概形。在 $S$ 上仿射的预概形 $X$ 在 $S$ 上有限，当且仅当 $\OO_S$-代数 $\A(X)$ 凝聚。
\begin{proof}
考虑前面的说明，只须注意：在 $S$ 局部诺特时，有限型拟凝聚 $\OO_S$-模恰为凝聚 $\OO_S$-模（I，1.5.1）。
\end{proof}

\egastatement{6.1.4}{命题}
设 $X$ 是在 $S$ 上为整（分别地，有限）的预概形，$f:X\to S$ 是结构态射。则对每个仿射开集 $U\subset S$，若其环为 $A$，则 $f^{-1}(U)$ 为仿射概形，且其环 $B$ 是 $A$ 上的整代数（分别地，有限代数）。

先证明下面的引理。

\egastatement{6.1.4.1}{引理}
设 $A$ 是环，$M$ 是 $A$-模，$(g_i)_{1\leq i\leq m}$ 是 $A$ 中的有限族元素，且各 $D(g_i)$ 覆盖 $\Spec(A)$。若每个 $M_{g_i}$ 都是有限型 $A_{g_i}$-模，则 $M$ 为有限型 $A$-模。
\begin{proof}
可以设 $M_{g_i}$ 有有限生成系 $(m_j/g_i^n)$，其中 $m_j\in M$，且对所有 $i$ 使用同一个 $n$。证明这些 $m_j$ 生成 $M$。令 $M'$ 为它们生成的子 $A$-模，取 $m\in M$。由假设，对每个 $i$，存在 $a_{ij}\in A$ 和不依赖于 $i$ 的整数 $p$，使在 $M_{g_i}$ 中
\[
\frac m1=\frac{\sum_j a_{ij}m_j}{g_i^p}.
\]
因此存在整数 $r\geq p$，使每个 $g_i^r m\in M'$。各 $D(g_i^r)=D(g_i)$ 覆盖 $\Spec(A)$，故各 $g_i^r$ 生成单位理想；即存在 $a_i\in A$ 使 $\sum_i a_i g_i^r=1$。所以
$m=(\sum_i a_i g_i^r)m\in M'$，引理得证。
\end{proof}

\noindent\textbf{（6.1.4）的证明。}
已知 $f^{-1}(U)$ 仿射（1.3.2）。令 $\varphi:A\to B$ 为与 $f$ 相应的同态。存在 $U$ 的有限主开覆盖 $D(g_i)$，$g_i\in A$，使令 $h_i=\varphi(g_i)$ 后，$B_{h_i}$ 在 $A_{g_i}$ 上为整（分别地，整且有限）。事实上，$U$ 有仿射开覆盖 $(V_\alpha)$，使令
$A_\alpha=A(V_\alpha)$、$B_\alpha=A(f^{-1}(V_\alpha))$ 后，$B_\alpha$ 在 $A_\alpha$ 上为整（分别地，有限）。每个 $x\in U$ 属于某个 $V_\alpha$，故存在 $g\in A$，使 $x\in D(g)\subset V_\alpha$。若 $g_\alpha$ 为 $g$ 在 $A_\alpha$ 中的像，则
$A(D(g))=A_g=(A_\alpha)_{g_\alpha}$。令 $h=\varphi(g)$，$h_\alpha$ 为 $h$ 在 $B_\alpha$ 中的像，则
$A(D(h))=B_h=(B_\alpha)_{h_\alpha}$。$B_\alpha$ 在 $A_\alpha$ 上为整（分别地，有限），故 $(B_\alpha)_{h_\alpha}$ 在 $(A_\alpha)_{g_\alpha}$ 上也如此。最后利用 $U$ 拟紧，即得到所求有限覆盖。

先设各 $B_{h_i}$ 在 $A_{g_i}$ 上整且有限。作为 $A_{g_i}$-模，$B_{h_i}$ 也就是 $B_{g_i}$，故引理（6.1.4.1）表明 $B$ 为有限型 $A$-模。
现只设各 $B_{h_i}$ 在 $A_{g_i}$ 上为整。取 $b\in B$，令 $C$ 为 $b$ 生成的子 $A$-代数。对每个 $i$，$C_{g_i}$ 是 $B_{h_i}$ 中由 $b/1$ 生成的 $A_{g_i}$-代数。由假设，每个 $C_{g_i}$ 为有限型 $A_{g_i}$-模，因此由（6.1.4.1），$C$ 为有限型 $A$-模。这证明 $B$ 在 $A$ 上为整。\hfill$\square$

\egastatement{6.1.5}{命题}
\begin{enumerate}[label=\textup{(\roman*)}]
\item 闭浸入有限，因而为整。
\item 两个有限态射（分别地，整态射）的复合有限（分别地，为整）。
\item 若 $f:X\to Y$ 为有限 $S$-态射（分别地，整 $S$-态射），则对每个基扩张 $S'\to S$，$f_{(S')}:X_{(S')}\to Y_{(S')}$ 有限（分别地，为整）。
\item 若 $f:X\to Y$、$g:X'\to Y'$ 为有限 $S$-态射（分别地，整 $S$-态射），则 $f\times_Sg:X\times_S X'\to Y\times_S Y'$ 也如此。\transnote{与（5.4.2，iv）相同，底本此处的乘积源、靶次序印错；这里按 $f,g$ 已给定的类型书写。}
\item 若 $f:X\to Y$、$g:Y\to Z$ 满足 $g\circ f$ 有限（分别地，为整）且 $g$ 分离，则 $f$ 有限（分别地，为整）。
\item 若 $f:X\to Y$ 有限（分别地，为整），则 $f_{\red}$ 也如此。
\end{enumerate}
\begin{proof}
由（I，5.5.12），只须证明（i）、（ii）、（iii）。为证闭浸入 $X\to S$ 有限，可设 $S=\Spec(A)$，再注意商环 $A/\frj$ 是单生成 $A$-模即可。

为证（ii），仍可设 $Z$，因而 $X,Y$，仿射（1.3.4）。此时断言等价于：若 $B$ 是有限 $A$-代数（分别地，整 $A$-代数），$C$ 是有限 $B$-代数（分别地，整 $B$-代数），则 $C$ 是有限 $A$-代数（分别地，整 $A$-代数）。这是立即的。

最后，为证（iii），由于 $X_{(S')}$ 等同于 $X\times_Y Y_{(S')}$（I，3.3.11），可设 $S=Y$；又可设 $S=\Spec(A)$、$S'=\Spec(A')$。此时 $X$ 仿射，环为 $B$（1.3.4）；$X_{(S')}$ 仿射，环为 $A'\otimes_A B$。只须注意：若 $B$ 是有限（分别地，整）$A$-代数，则 $A'\otimes_A B$ 是有限（分别地，整）$A'$-代数。
\end{proof}
还应注意：若 $S$-预概形 $X,Y$ 在 $S$ 上有限（分别地，为整），则和预概形 $X\amalg Y$ 也如此。因为这归结为：有限（分别地，整）$A$-代数 $B,C$ 的直积 $B\times C$ 仍是有限（分别地，整）$A$-代数。

\egastatement{6.1.6}{推论}
若 $X$ 在 $S$ 上为整（分别地，有限），则对每个开集 $U\subset S$，$f^{-1}(U)$ 在 $U$ 上为整（分别地，有限）。
\begin{proof}
这是（6.1.5，iii）的特例。
\end{proof}

\egastatement{6.1.7}{推论}
设 $f:X\to Y$ 是有限态射。则对每个 $y\in Y$，纤维 $f^{-1}(y)$ 是 $k(y)$ 上的有限代数概形；特别地，其底空间有限且离散。
\begin{proof}
作为 $k(y)$-预概形，$f^{-1}(y)$ 等同于 $X\times_Y\Spec(k(y))$（I，3.6.1），它在 $\Spec(k(y))$ 上有限（6.1.5，iii）。因此它为仿射概形，且其环为 $k(y)$ 上的有限秩代数（6.1.4）。由（I，6.4.4）即得结论。
\end{proof}

\egastatement{6.1.8}{推论}
设 $X,S$ 是整预概形，$f:X\to S$ 为占优态射。若 $f$ 为整（分别地，有限），则 $X$ 的有理函数域 $R(X)$ 在 $S$ 的有理函数域 $R(S)$ 上代数（分别地，为有限次代数扩张）。
\begin{proof}
令 $s$ 为 $S$ 的泛点。$k(s)$-预概形 $f^{-1}(s)$ 在 $\Spec(k(s))$ 上为整（分别地，有限）（6.1.5，iii），且按假设包含 $X$ 的泛点 $x$。$x$ 在 $f^{-1}(s)$ 中的局部环等于 $k(x)$（I，3.6.5），它是 $k(s)$ 上的某个整代数（分别地，有限代数）的局部环（6.1.4），故得推论。
\end{proof}

\egastatement{6.1.9}{注记}
为使（6.1.5，v）成立，$g$ 分离的假设是必要的。事实上，若 $Y$ 在 $Z$ 上不分离，则恒等态射 $1_Y$ 分解为
$Y\xrightarrow{\Delta_Y}Y\times_Z Y\xrightarrow{p_1}Y$，但由下面的（6.1.10），$\Delta_Y$ 不是整态射。

\egastatement{6.1.10}{命题}
每个整态射都是普遍闭态射。
\begin{proof}
设 $f:X\to Y$ 为整态射。由（6.1.5，iii），只须证明 $f$ 闭。取 $X$ 的闭子集 $Z$，存在具有该底空间的闭子预概形（I，5.2.1），所以由（6.1.5，i、ii），只须证明 $f(X)$ 在 $Y$ 中闭。由（6.1.5，vi），可设 $X,Y$ 约化。令 $T$ 为 $Y$ 中底空间为 $\overline{f(X)}$ 的约化闭子预概形（I，5.2.1）；$f$ 分解为
$X\xrightarrow{g}T\xrightarrow{j}Y$，其中 $j$ 为典范嵌入（I，5.2.2）。$j$ 分离（I，5.5.1，i），故由（6.1.5，v），$g$ 为整。因此可设 $f(X)$ 在 $Y$ 中稠密。

最后，问题关于 $Y$ 局部，可设 $Y=\Spec(A)$。于是 $X=\Spec(B)$，其中 $B$ 为整 $A$-代数（6.1.4）；$A$ 约化（I，5.1.4），而 $f(X)$ 在 $Y$ 中稠密，推出相应同态 $\varphi:A\to B$ 单射（I，1.2.7）。在此条件下，$f(X)=Y$ 就是说：$A$ 的每个素理想都是 $B$ 的某个素理想与 $A$ 的交。这正是 Cohen--Seidenberg 第一定理（[13]，第一卷，第257页，定理3）。
\end{proof}

\egastatement{6.1.11}{推论}
每个有限态射 $f:X\to Y$ 都是射影态射。
\begin{proof}
由于 $f$ 仿射，$\OO_X$ 相对于 $f$ 甚丰沛（5.1.2）；又 $f_*(\OO_X)$ 是有限型拟凝聚 $\OO_Y$-模（6.1.2）。最后，$f$ 分离、为有限型且普遍闭（6.1.10），因此满足准则（5.5.4，i）的条件。
\end{proof}

\egastatement{6.1.12}{命题}
设 $f:X'\to X$ 是有限态射，令 $\B=f_*(\OO_{X'})$，它为拟凝聚 $\OO_X$-代数及有限型 $\OO_X$-模。设 $\F$ 为拟凝聚 $\OO_{X'}$-模。$\F$ 局部自由且秩为 $r$，当且仅当 $f_*(\F)$ 为秩 $r$ 的局部自由 $\B$-模。

若对某个开集 $U\subset X$，$f_*(\F)|U\simeq\B^r|U$，则显然
$\F|f^{-1}(U)\simeq\OO_{X'}^r|f^{-1}(U)$（1.4.2）。反之，设 $\F$ 为秩 $r$ 的局部自由模，证明 $f_*(\F)$ 作为 $\B$-模局部同构于 $\B^r$。取 $x\in X$。由于 $f$ 闭（6.1.10），当 $U$ 遍历 $x$ 的仿射邻域基时，$f^{-1}(U)$ 遍历有限集 $f^{-1}(x)$ 的仿射邻域基（1.2.5）。命题即由下面的引理得出。

\egastatement{6.1.12.1}{引理}
设 $Y$ 为预概形，$\E$ 为秩 $r$ 的局部自由 $\OO_Y$-模，$Z$ 为包含在某个仿射开集 $V$ 中的有限子集。则存在 $Z$ 的邻域 $U\subset V$，使 $\E|U\simeq\OO_Y^r|U$。
\begin{proof}
显然可设 $Y$ 仿射。对每个 $z\in Z$，闭包 $\overline{\{z\}}$ 中至少有一个闭点 $z'$（0，2.1.3）。若 $Z'$ 为这些 $z'$ 的集合，则 $Z'$ 的每个邻域也是 $Z$ 的邻域，故可设 $Z$ 在 $Y$ 中闭且离散。取具有底空间 $Z$ 的约化闭子预概形（I，5.2.1），记其典范嵌入为 $j:Z\to Y$。则
$j^*(\E)=\E\otimes_Y\OO_Z$ 在离散概形 $Z$ 上局部自由且秩为 $r$，故同构于 $\OO_Z^r$。即存在 $\E\otimes_Y\OO_Z$ 的 $r$ 个全局截面 $s_i$，使它们定义的同态
$\OO_Z^r\to\E\otimes_Y\OO_Z$ 为双射。

设 $Y=\Spec(A)$，$Z$ 由理想 $\frj$ 定义，且 $\E=\wt M$。各 $s_i$ 是 $M\otimes_A(A/\frj)$ 的元素，故为 $r$ 个元素
$t_i\in M=\Gamma(Y,\E)$ 的像。对每个 $z\in Z$，存在邻域 $V_z$，使各 $t_i$ 的限制定义同构
$\OO_Y^r|V_z\to\E|V_z$（0，5.5.4）。于是 $U=\bigcup_{z\in Z}V_z$ 满足要求。
\end{proof}

\egastatement{6.1.13}{命题}
设 $g:X'\to X$ 是整态射，$Y$ 是正规且局部整的预概形，$f$ 是从 $Y$ 到 $X'$ 的有理映射，且 $g\circ f$ 为处处有定义的有理映射（I，7.2.1）。则 $f$ 处处有定义。

若 $f_1,f_2$ 是有理映射类 $f$ 中的两个态射，分别定义在 $Y$ 的稠密开集上并取值于 $X'$，则显然 $g\circ f_1$ 和 $g\circ f_2$ 等价。这说明用 $g\circ f$ 表示其等价类是合理的。还应回顾：若再假设 $Y$ 局部诺特，则 $Y$ 正规已经推出它局部整（I，6.1.13）。

为证（6.1.13），先注意问题关于 $Y$ 局部，故可设 $g\circ f$ 的类中有态射 $h:Y\to X$。考察逆像 $Y'=X'_{(Y)}$，则
$g'=g_{(Y)}:Y'\to Y$ 为整态射（6.1.5，iii）。利用从 $Y$ 到 $X'$ 的有理映射与 $Y'$ 的有理 $Y$-截面之间的对应（I，7.1.2），可归结为证明 $X=Y$ 的特例，即：

\egastatement{6.1.14}{推论}
设 $X$ 是正规且局部整的预概形，$g:X'\to X$ 是整态射，$f$ 是 $X'$ 的有理 $X$-截面。则 $f$ 处处有定义。

问题关于 $X$ 局部，可设 $X$ 为整预概形；此时 $f$ 等同于某个开集 $U\subset X$ 到 $X'$ 的态射（I，7.2.2），并且是 $g^{-1}(U)$ 的 $U$-截面。由于 $g$ 分离，$f$ 为 $U$ 到 $g^{-1}(U)$ 的闭浸入（I，5.4.6）。令 $Z$ 为与 $f$ 相伴的闭子预概形（I，4.2.1），它同构于 $U$，因而为整。令 $X_1$ 为 $X'$ 中底空间为 $Z$ 的闭包的约化子预概形（I，5.2.1）。则 $Z$ 是 $X_1$ 在一个开集上诱导的子预概形（I，5.2.3）；$Z$ 不可约，所以 $X_1$ 也不可约，进而为整。此时 $f$ 可看成 $X_1$ 的有理 $X$-截面，而 $g|X_1$ 为整态射（6.1.5，i、ii）。因此，最终归结为 $X'=X_1$ 的情形，即：

\egastatement{6.1.15}{推论}
设 $X$ 是整且正规的预概形，$X'$ 是整预概形，$g:X'\to X$ 是整态射。若 $X'$ 存在有理 $X$-截面 $f$，则 $g$ 为同构。
\begin{proof}
问题关于 $X$ 局部，可设 $X$ 仿射，环为整环 $A$。则 $X'$ 仿射，环 $A'$ 为整环且在 $A$ 上为整（6.1.4）。又由（6.1.14）的论证，$X$ 有一个稠密开集同构于 $X'$ 的一个稠密开集，故 $A,A'$ 的分式域相同。由（I，8.2.1.1）及各 $\OO_x$ 整闭的假设，$A$ 整闭，所以 $A'=A$。这也完成（6.1.13）的证明。
\end{proof}

\subsection{拟有限态射}
\src{114--116}

\egastatement{6.2.1}{命题}
设 $f:X\to Y$ 为局部有限型态射，$x\in X$。下列条件等价：
\begin{enumerate}[label=\alph*)]
\item 点 $x$ 在它的纤维 $f^{-1}(f(x))$ 中孤立。
\item 环 $\OO_x$ 是拟有限 $\OO_{f(x)}$-模（0，7.4.1）。
\end{enumerate}
\begin{proof}
问题显然关于 $X,Y$ 局部，故可设 $X=\Spec(A)$、$Y=\Spec(B)$ 仿射，$A$ 是有限型 $B$-代数（I，6.3.3）。还可将 $X$ 替换为
$X\times_Y\Spec(\OO_{f(x)})$ 而不改变纤维 $f^{-1}(f(x))$ 和局部环 $\OO_x$（I，3.6.5）。因此可设 $B$ 为局部环，等于 $\OO_{f(x)}$。若 $\nn$ 是 $B$ 的极大理想，则 $f^{-1}(f(x))$ 为环是 $A/\nn A$ 的仿射概形，并在 $k(f(x))=B/\nn$ 上为有限型（I，6.4.11）。

若 a) 成立，还可设 $f^{-1}(f(x))$ 只含点 $x$；于是 $A/\nn A$ 在 $B/\nn$ 上为有限秩（I，6.4.4），即 $A$ 为拟有限 $B$-模。反之，若如此，则 $f^{-1}(f(x))$ 为阿廷仿射概形，因而离散（I，6.4.4）；于是 $x$ 在纤维中孤立，故 b) 推出 a)。\transnote{原文证明在局部环 $\OO_x$ 与某个仿射邻域的坐标环 $A$ 之间省略了一步：由 $\OO_x/\nn\OO_x$ 的有限性回到纤维中 $x$ 的局部零维性，再缩小邻域。正文按原文保留这一证明的详略。}
\end{proof}

\egastatement{6.2.2}{推论}
设 $f:X\to Y$ 是有限型态射。下列条件等价：
\begin{enumerate}[label=\alph*)]
\item 每个 $x\in X$ 都在其纤维 $f^{-1}(f(x))$ 中孤立，即该纤维的底空间离散。
\item 对每个 $x\in X$，$f^{-1}(f(x))$ 为有限 $k(f(x))$-预概形。
\item 对每个 $x\in X$，环 $\OO_x$ 是拟有限 $\OO_{f(x)}$-模。
\end{enumerate}
\begin{proof}
a) 与 c) 的等价性由（6.2.1）得出。另一方面，$f^{-1}(f(x))$ 是 $k(f(x))$ 上的代数预概形（I，6.4.11），故 a) 与 b) 的等价性由（I，6.4.4）得出。
\end{proof}

\egastatement{6.2.3}{定义}
若有限型态射 $f:X\to Y$ 满足（6.2.2）的等价条件，称 $f$ 为\emph{拟有限态射}，或称 $X$ \emph{在 $Y$ 上拟有限}。
显然每个有限态射都拟有限（6.1.8）。\transnote{这里原文引（6.1.8）；直接给出有限纤维结论的是（6.1.7）。}

\egastatement{6.2.4}{命题}
\begin{enumerate}[label=\textup{(\roman*)}]
\item 闭浸入 $X\to Y$，以及源 $X$ 为诺特的浸入 $X\to Y$，都是拟有限态射。
\item 若 $f:X\to Y$、$g:Y\to Z$ 拟有限，则 $g\circ f$ 拟有限。
\item 若 $X,Y$ 为 $S$-预概形，$f:X\to Y$ 为拟有限 $S$-态射，则对每个基扩张 $g:S'\to S$，$f_{(S')}:X_{(S')}\to Y_{(S')}$ 拟有限。
\item 若 $f:X\to Y$、$g:X'\to Y'$ 为拟有限 $S$-态射，则
$f\times_Sg:X\times_S X'\to Y\times_S Y'$ 拟有限。\transnote{底本的乘积源、靶有与（6.1.5，iv）相同的错排；这里明确写出由 $f,g$ 决定的类型。}
\item 若 $f:X\to Y$、$g:Y\to Z$ 满足 $g\circ f$ 拟有限，并且 $g$ 分离，或 $X$ 诺特，或 $X\times_ZY$ 局部诺特，则 $f$ 拟有限。
\item 若 $f$ 拟有限，则 $f_{\red}$ 拟有限。
\end{enumerate}
\begin{proof}
若 $f$ 为浸入，其非空纤维只有一点，故（i）由（I，6.3.4，i；6.3.5）得出。为证（ii），先注意 $h=g\circ f$ 为有限型（I，6.3.4，ii）。若 $z=h(x)$、$y=f(x)$，则 $y$ 在 $g^{-1}(z)$ 中孤立；因此存在 $y$ 的开邻域 $V\subset Y$，不含 $g^{-1}(z)$ 中不同于 $y$ 的点。于是 $f^{-1}(V)$ 是 $x$ 的开邻域，不交任何 $f^{-1}(y')$，其中 $y'\in g^{-1}(z)$ 且 $y'\ne y$。由于 $x$ 在 $f^{-1}(y)$ 中孤立，它也在 $h^{-1}(z)=f^{-1}(g^{-1}(z))$ 中孤立。

为证（iii），可设 $Y=S$（I，3.3.11）。首先 $f'=f_{(S')}$ 为有限型（I，6.3.4，iii）。若 $x'\in X'=X_{(S')}$，令 $y'=f'(x')$、$y=g(y')$，则
$f'^{-1}(y')$ 等同于 $f^{-1}(y)\otimes_{k(y)}k(y')$（I，3.6.5）。按假设，$f^{-1}(y)$ 在 $k(y)$ 上为有限秩，故 $f'^{-1}(y')$ 在 $k(y')$ 上为有限秩，从而离散。

由一般方法（I，5.5.12），（iv）、（v）、（vi）由（i）、（ii）、（iii）得出；只有（v）中不采用 $g$ 分离假设的情形须另作说明。此时，若 $x$ 在
$f^{-1}(g^{-1}(g(f(x))))$ 中孤立，则它当然在 $f^{-1}(f(x))$ 中孤立；而 $f$ 为有限型由（I，6.3.6）得出。
\end{proof}

\egastatement{6.2.5}{命题}
设 $A$ 是完备诺特局部环，$X$ 是局部有限型 $A$-概形，$x\in X$ 位于 $Y=\Spec(A)$ 的闭点 $y$ 之上，且在纤维 $f^{-1}(y)$ 中孤立，其中 $f:X\to Y$ 为结构态射。则 $\OO_x$ 为有限型 $A$-模，且 $X$ 与
$X'=\Spec(\OO_x)$ 及某个 $A$-概形 $X''$ 的和（I，3.1）$Y$-同构；其中 $X'$ 为有限 $Y$-概形。
\begin{proof}
由（6.2.1），$\OO_x$ 为拟有限 $A$-模。由于 $\OO_x$ 诺特（I，6.3.7）且同态 $A\to\OO_x$ 为局部同态，$A$ 的完备性推出 $\OO_x$ 是有限型 $A$-模（0，7.4.3）。令 $X'=\Spec(\OO_x)$ 为 $X$ 在 $x$ 处的局部概形（I，2.4.1），$g:X'\to X$ 为典范态射。复合 $X'\to X\to Y$ 有限（6.1.1），而 $f$ 分离，故 $g$ 有限（6.1.5，v），于是 $g(X')$ 在 $X$ 中闭（6.1.10）。另一方面，$g$ 为有限型，由 $g$ 的定义及（I，6.5.4），它在 $X'$ 的闭点 $x'$ 处为局部同构。但 $X'$ 是 $x'$ 的唯一开邻域，因而 $g$ 为开浸入，故 $g(X')$ 在 $X$ 中也开，命题得证。
\end{proof}

\egastatement{6.2.6}{推论}
设 $A$ 是完备诺特局部环，$Y=\Spec(A)$，$f:X\to Y$ 是分离拟有限态射。则 $X$ 与一个和 $X'\amalg X''$ $Y$-同构，其中 $X'$ 为有限 $Y$-概形，$X''$ 为拟有限 $Y$-概形，且若 $y$ 是 $Y$ 的闭点，则 $X''\cap f^{-1}(y)=\varnothing$。
\begin{proof}
按假设，纤维 $f^{-1}(y)$ 有限且离散；对它的点数作归纳，由（6.2.5）即得推论。
\end{proof}

\egastatement{6.2.7}{注记}
在第五章中，我们将看到：若 $Y$ 局部诺特，则分离拟有限态射 $X\to Y$ 必然拟仿射。

\subsection{预概形的整闭包}
\src{116--120}

\egastatement{6.3.1}{命题}
设 $(X,\A)$ 为赋环空间，$\B$ 为一个交换 $\A$-代数，$f$ 为 $\B$ 在 $X$ 上的截面。下列性质等价：
\begin{enumerate}[label=\alph*)]
\item $\B$ 中由各 $f^n$（$n\geq0$）生成的子 $\A$-模（0，5.1.1）为有限型。
\item 存在 $\B$ 的一个子 $\A$-代数 $\C$，它作为 $\A$-模为有限型，且 $f\in\Gamma(X,\C)$。
\item 对每个 $x\in X$，$f_x$ 在茎环 $\A_x$ 上为整。
\end{enumerate}
\begin{proof}
由各 $f^n$ 生成的子 $\A$-模也是一个 $\A$-代数，故（a）推出（b）。另一方面，（b）说明对每个 $x\in X$，$\C_x$ 为有限型 $\A_x$-模，因而代数 $\C_x$ 中每个元素，特别是 $f_x$，都在 $\A_x$ 上为整。最后，若在某点 $x$ 有关系
\[
f_x^n+(a_1)_x f_x^{n-1}+\cdots+(a_n)_x=0,
\]
其中各 $a_i$ 是 $\A$ 在 $x$ 的某个开邻域 $U$ 上的截面，则截面
$f^n|U+a_1 f^{n-1}|U+\cdots+a_n$ 在 $x$ 的某个邻域 $V\subset U$ 上为零。因此，各 $f^k|V$（$k\geq0$）都是各 $f^j|V$（$0\leq j\leq n-1$）的以 $\Gamma(V,\A)$ 为系数的线性组合。这证明（c）推出（a）。
\end{proof}
当（6.3.1）的等价条件成立时，称截面 $f$ \emph{在 $\A$ 上为整}。

\egastatement{6.3.2}{推论}
在（6.3.1）的假设下，存在唯一的子 $\A$-模 $\A'\subset\B$，使对每个 $x\in X$，$\A'_x$ 为所有在 $\A_x$ 上为整的芽 $f_x\in\B_x$ 构成的集合。对每个开集 $U\subset X$，$\A'$ 在 $U$ 上的截面，就是在 $\A|U$ 上为整的截面 $f\in\Gamma(U,\B)$。
\begin{proof}
令 $\Gamma(U,\A')$ 为所有满足“对每个 $x\in U$，$f_x$ 在 $\A_x$ 上为整”的 $f\in\Gamma(U,\B)$ 所成的集合，便立即得到 $\A'$ 的存在性。第二个断言由（6.3.1）立即得出。
\end{proof}
显然，$\A'$ 是 $\B$ 的一个子 $\A$-代数；称它为 $\A$ 在 $\B$ 中的\emph{整闭包}。

\egapar{6.3.3}
设 $(X,\A)$、$(Y,\B)$ 为两个赋环空间，$g=(\psi,\theta):X\to Y$ 为态射。设 $\C$、$\mathscr D$ 分别为 $\A$-代数和 $\B$-代数，$u:\mathscr D\to\C$ 为一个 $g$-态射（0，4.4.1）。若 $\A'$、$\B'$ 分别为 $\A$ 在 $\C$ 中、$\B$ 在 $\mathscr D$ 中的整闭包，则 $u$ 在 $\B'$ 上的限制为一个 $g$-态射
\[
u':\B'\longrightarrow\A'.
\]
事实上，令 $j:\B'\to\mathscr D$ 为典范嵌入，只须证明
$u^\sharp\circ g^*(j):g^*(\B')\to\C$ 将 $g^*(\B')$ 映入 $\A'$。由 $\B'$ 的定义，
\[
(g^*(\B'))_x=\B'_{\psi(x)}\otimes_{\B_{\psi(x)}}\A_x
\]
中的元素在 $\A_x$ 上为整，所以它在 $u_x^\sharp$ 下的像也如此。这就证明了断言。

\egastatement{6.3.4}{命题}
设 $X$ 为预概形，$\A$ 为拟凝聚 $\OO_X$-代数。则 $\OO_X$ 在 $\A$ 中的整闭包 $\OO'_X$ 为拟凝聚 $\OO_X$-代数；而且对每个仿射开集 $U\subset X$，$\Gamma(U,\OO'_X)$ 是 $\Gamma(U,\OO_X)$ 在 $\Gamma(U,\A)$ 中的整闭包。
\begin{proof}
可设 $X=\Spec(B)$ 仿射，$\A=\wt A$，其中 $A$ 是 $B$-代数。令 $B'$ 为 $B$ 在 $A$ 中的整闭包。只须证明：对每个 $x\in X$，$A_x$ 中在 $B_x$ 上为整的元素必属于 $B'_x$。这来自如下事实：对交换环 $C$，在一个 $C$-代数中取整闭包，与相对于 $C$ 的一个乘法子集取分式环，这两种运算可交换（[13]，第一卷，第261、257页）。
\end{proof}
此时，$X$-概形 $X'=\Spec(\OO'_X)$ 称为 $X$ \emph{相对于 $\A$（或相对于 $\Spec(\A)$）的整闭包}。显然，$X'$ 在 $X$ 上为整（6.1.2）。
由（6.3.4）立即可知，若 $f:X'\to X$ 为结构态射，则对每个开集 $U\subset X$，$f^{-1}(U)$ 是 $X$ 在 $U$ 上诱导的预概形相对于 $\A|U$ 的整闭包。

\egapar{6.3.5}
设 $X,Y$ 为两个预概形，$f:X\to Y$ 为态射，$\A$、$\B$ 分别为拟凝聚 $\OO_X$-代数和拟凝聚 $\OO_Y$-代数，$u:\B\to\A$ 为 $f$-态射。已知（6.3.3）由此得到一个 $f$-态射 $u':\OO'_Y\to\OO'_X$，其中 $\OO'_X$、$\OO'_Y$ 分别为 $\OO_X$ 在 $\A$ 中、$\OO_Y$ 在 $\B$ 中的整闭包。因此，若 $X'$、$Y'$ 分别为 $X,Y$ 相对于 $\A,\B$ 的整闭包，则 $u$ 典范地给出态射 $f'=\Spec(u'):X'\to Y'$（1.5.6），使下图交换：
\begin{equation}\tag{6.3.5.1}
\begin{tikzcd}
X'\arrow[r,"f'"]\arrow[d]&Y'\arrow[d]\\
X\arrow[r,"f"']&Y.
\end{tikzcd}
\end{equation}

\egapar{6.3.6}
设 $X$ 只有有限个不可约分支 $X_i$（$1\leq i\leq r$），它们的泛点为 $\xi_i$。特别考虑 $X$ 相对于一个拟凝聚 $\R(X)$-代数 $\A$ 的整闭包；这里拟凝聚既可作为 $\OO_X$-代数理解，也可作为 $\R(X)$-代数理解，两者相同。由（I，7.3.5），$\A$ 是 $r$ 个拟凝聚 $\OO_X$-代数 $\A_i$ 的直积，$\A_i$ 的支集包含于 $X_i$，而它在 $X_i$ 上诱导的层为常层，其茎 $A_i$ 是 $\OO_{\xi_i}$ 上的代数。由（6.3.4），$\OO_X$ 在 $\A$ 中的整闭包 $\OO'_X$ 是 $\OO_X$ 在各 $\A_i$ 中的整闭包 $\OO'_{X_i}$ 的直积。因此，$X'=\Spec(\OO'_X)$ 为各 $X$-概形 $X'_i=\Spec(\OO'_{X_i})$（$1\leq i\leq r$）的和。

再设 $\A$ 约化，等价地，各代数 $A_i$ 约化；因而 $A_i$ 可看作域 $k(\xi_i)$ 上的代数，该域等于以 $X_i$ 为底空间的约化子预概形的有理函数域。于是各 $X'_i$ 都是约化 $X$-概形（1.3.8），而 $X'$ 也是 $X_{\red}$ 的整闭包。进一步，设每个代数 $A_i$ 是有限个域 $K_{ij}$（$1\leq j\leq s_i$）的直积。若 $\K_{ij}$ 为 $\A_i$ 中对应于 $K_{ij}$ 的子代数，则 $\OO'_{X_i}$ 显然是 $\OO_X$ 在各 $\K_{ij}$ 中的整闭包 $\OO'_{X_{ij}}$ 的直积。因此 $X'_i$ 为各 $X$-概形 $X'_{ij}=\Spec(\OO'_{X_{ij}})$（$1\leq j\leq s_i$）的和。此外，在上述假设和记号下，有：

\egastatement{6.3.7}{命题}
每个 $X'_{ij}$ 都是整的正规 $X_i$-预概形，其有理函数域 $R(X'_{ij})$ 典范地等同于 $k(\xi_i)$ 在 $K_{ij}$ 中的相对代数闭包 $K'_{ij}$。
\begin{proof}
由前面的论述，可设 $X$ 整，因而 $r=1$，且 $s_1=1$，于是唯一的代数 $A_1$ 是一个域 $K$。令 $\xi$ 为 $X$ 的泛点，$f:X'\to X$ 为结构态射。对每个非空仿射开集 $U\subset X$，$f^{-1}(U)$ 等同于整环 $B_U=\Gamma(U,\OO_X)$ 在域 $K$ 中的整闭包 $B'_U$ 的谱（6.3.4）。$B'_U$ 为整闭整环，故其谱的各点的局部环也如此。因此，按定义 $f^{-1}(U)$ 是整的正规概形（0，4.1.4；I，5.1.4）。
又因 $(0)$ 是 $B'_U$ 中位于 $B_U$ 的素理想 $(0)$ 之上的唯一素理想（[13]，第一卷，第259页），$f^{-1}(\xi)$ 只有一个点 $\xi'$；$k(\xi')$ 是 $B'_U$ 的分式域 $K'$，也就是 $k(\xi)$ 在 $K$ 中的相对代数闭包。最后，$U$ 遍历 $X$ 的非空仿射开集时，各 $f^{-1}(U)$ 构成 $X'$ 的不可约开覆盖；任意两个这样的开集的交 $f^{-1}(U\cap V)$ 都含 $\xi'$，故非空。由（0，2.1.4），$X'$ 不可约。
\end{proof}

\egastatement{6.3.8}{推论}
设 $X$ 是只有有限个不可约分支 $X_i$（$1\leq i\leq r$）的约化预概形，$\xi_i$ 为 $X_i$ 的泛点。则 $X$ 相对于 $\R(X)$ 的整闭包 $X'$ 是 $r$ 个整的正规 $X$-概形 $X'_i$ 的和。若 $f:X'\to X$ 是结构态射，则 $f^{-1}(\xi_i)$ 只含 $X'_i$ 的泛点 $\xi'_i$，且 $k(\xi_i)=k(\xi'_i)$；换言之，$f$ 为双有理态射。

此时称 $X'$ 为约化预概形 $X$ 的\emph{正规化}。应注意，$f$ 双有理且为整，故为满射（6.1.10）。有 $X'=X$，当且仅当 $X$ 正规。若 $X$ 为整预概形，则由（6.3.8），它的正规化 $X'$ 也为整。

\egapar{6.3.9}
设 $X,Y$ 为整预概形，$f:X\to Y$ 为占优态射，$L=R(X)$、$K=R(Y)$ 分别为其有理函数域。$f$ 典范地给出单射 $K\to L$；将 $K,L$ 分别等同于常层 $\R(Y),\R(X)$ 时，这个单射为一个 $f$-态射。设 $K_1,L_1$ 分别为 $K,L$ 的扩域，且给定使下图交换的单同态 $K_1\to L_1$：
\[
\begin{tikzcd}
K_1\arrow[r]&L_1\\
K\arrow[r]\arrow[u]&L.\arrow[u]
\end{tikzcd}
\]
将 $K_1,L_1$ 看作分别定义在 $Y,X$ 上的常层，即 $\R(Y)$-代数和 $\R(X)$-代数时，这就是说 $K_1\to L_1$ 是一个 $f$-态射。若 $X',Y'$ 分别为 $X,Y$ 相对于 $L_1,K_1$ 的整闭包，则 $X',Y'$ 是整的正规预概形（6.3.6），其有理函数域分别典范地等同于 $L$ 在 $L_1$ 中、$K$ 在 $K_1$ 中的相对代数闭包 $L',K'$；而且存在一个必为占优的典范态射 $f':X'\to Y'$，使图（6.3.5.1）交换。

最重要的情形是 $L_1=L$，故 $K_1$ 为包含在 $L$ 中的 $K$ 的扩域；并且设 $X$ 整且正规，因而 $X'=X$。以上说明：当 $X$ 正规，且 $Y'$ 是 $Y$ 相对于某个域 $K_1\subset L=R(X)$ 的整闭包时，每个占优态射 $f:X\to Y$ 都分解为
\[
f:X\xrightarrow{f'}Y'\longrightarrow Y,
\]
其中 $f'$ 占优。此外，当单同态 $K_1\to L$ 给定时，$f'$ 必唯一；化归 $X,Y$ 仿射即能看出。因此，对给定的 $Y,L$ 和一个 $K$-单同态 $K_1\to L$，$Y$ 相对于 $K_1$ 的整闭包 $Y'$ 是一个泛问题的解。

\egastatement{6.3.10}{注记}
回到（6.3.6）的假设，并进一步假设各代数 $A_i$ 在 $k(\xi_i)$ 上为有限秩（因此 $A_i$ 是有限个域的直积）。在某些情形，可以断言结构态射 $X'\to X$ 不仅为整，而且有限。只考虑 $X$ 约化的情形；问题关于 $X$ 局部，又可设 $X$ 仿射，环为 $C$，其极小素理想只有有限个 $\pp_i$（$1\leq i\leq r$），令 $C_i=C/\pp_i$，它们为整环。若每个 $C_i$ 在其分式域的任意有限次扩张中的整闭包均为有限型 $C$-模，则 $X'$ 在 $X$ 上有限（6.3.4）。已知以下情况下此条件总成立：$C$ 是一个域上的有限型代数（[13]，第一卷，第267页，定理9），或是 $\Z$ 上的有限型代数（[9，I]，第93页，定理3），或是某个完备诺特局部环上的有限型代数（[25]，第298页）。因此，当 $X$ 在一个域、$\Z$ 或一个完备诺特局部环上为有限型时，$X'\to X$ 为有限态射。

\subsection{模层自同态的行列式}
\src{120--125}

\egapar{6.4.1}
设 $A$ 为交换环，$E$ 为秩 $n$ 的自由 $A$-模，$u$ 为 $E$ 的自同态。回顾 $u$ 的\emph{特征多项式}的定义：考虑秩 $n$ 的自由 $A[T]$-模 $E\otimes_A A[T]$（$T$ 为不定元）上的自同态 $u\otimes1$，并令
\begin{equation}\tag{6.4.1.1}
P(u,T)=\det\bigl(T\cdot1-(u\otimes1)\bigr),
\end{equation}
其中 $1$ 表示 $E\otimes_A A[T]$ 的恒等自同构。有
\begin{equation}\tag{6.4.1.2}
P(u,T)=T^n-\sigma_1(u)T^{n-1}+\cdots+(-1)^n\sigma_n(u),
\end{equation}
其中 $\sigma_i(u)\in A$，是 $u$ 关于 $E$ 的任意一组基的矩阵元素的一个 $i$ 次齐次多项式，系数为整数。称各 $\sigma_i(u)$ 为 $u$ 的\emph{初等对称函数}；特别地，$\sigma_1(u)=\Tr u$，$\sigma_n(u)=\det u$。由 Hamilton--Cayley 定理，有
\begin{equation}\tag{6.4.1.3}
P(u,u)=u^n-\sigma_1(u)u^{n-1}+\cdots+(-1)^n\sigma_n(u)=0.
\end{equation}
也可写成
\begin{equation}\tag{6.4.1.4}
(\det u)\cdot1_E=uQ(u),
\end{equation}
其中 $1_E$ 为 $E$ 的恒等自同构，且
\begin{equation}\tag{6.4.1.5}
Q(u)=(-1)^{n+1}\bigl(u^{n-1}-\sigma_1(u)u^{n-2}+\cdots+(-1)^{n-1}\sigma_{n-1}(u)\bigr).
\end{equation}
设 $\varphi:A\to B$ 为环同态，使 $B$ 成为 $A$-代数。考虑秩 $n$ 的自由 $B$-模 $E_{(B)}=E\otimes_A B$，以及 $u$ 扩张得到的自同态 $u\otimes1$。立即可知，对每个 $i$，有 $\sigma_i(u\otimes1)=\varphi(\sigma_i(u))$。

\egapar{6.4.2}
现设 $A$ 为整环，分式域为 $K$，$E$ 为有限型 $A$-模，但不一定自由。令 $n$ 为 $E$ 的秩，即自由 $K$-模 $E\otimes_A K$ 的秩。$E$ 的每个自同态 $u$ 都典范地给出 $E\otimes_A K$ 的自同态 $u\otimes1$。稍滥用术语，仍称多项式 $P(u\otimes1,T)$ 为 $u$ 的\emph{特征多项式}，并记作 $P(u,T)$；其系数 $\sigma_i(u\otimes1)$ 属于 $K$，仍记为 $\sigma_i(u)$，并称为 $u$ 的初等对称函数。特别地，按定义 $\det u=\det(u\otimes1)$。采用这些记号后，公式（6.4.1.3）至（6.4.1.5）仍有意义且成立，只须把 $u^i$ 理解为同态 $E\to E\otimes_A K$：先施行典范同态 $x\mapsto x\otimes1$，再施行 $E\otimes_A K$ 的自同态 $u^i\otimes1=(u\otimes1)^i$。

令 $F$ 为 $E$ 的挠子模，$E_0=E/F$。有 $u(F)\subset F$，故取商后 $u$ 给出 $E_0$ 的自同态 $u_0$。又 $E\otimes_A K$ 等同于 $E_0\otimes_A K$，而 $u\otimes1$ 等同于 $u_0\otimes1$，所以对 $1\leq i\leq n$ 有 $\sigma_i(u)=\sigma_i(u_0)$。
当 $E$ 无挠时，$E$ 典范地等同于 $E\otimes_A K$ 的一个子 $A$-模，且 $u\otimes1=0$ 等价于 $u=0$。若 $E$ 自由，由前面的说明，（6.4.1）和（6.4.2）对 $\sigma_i(u)$ 给出的定义一致，从而上述记号是相容的。
还应注意，若 $E$ 为挠模，即 $E_0=\{0\}$，则 $E_0$ 的外代数归结为 $K$，而 $E_0$ 的唯一自同态 $u_0$ 的行列式等于 $1$。

\NewNote{\textbf{作者勘误提示：}本条证明已有作者正式更正；有限型子模中的元素不一定在底环上整。 参见\hyperlink{errIII:21}{《勘误与增补（清单二）》第21条}。}
\egastatement{6.4.3}{命题}
设 $A$ 为整环，$E$ 为有限型 $A$-模，$u$ 为 $E$ 的自同态。则 $u$ 的初等对称函数 $\sigma_i(u)$，特别是 $\det u$，是 $K$ 中在 $A$ 上为整的元素。
\begin{proof}
令 $A'$ 为 $A$ 的整闭包。因为 $A'[T]$ 为整闭环（[24]，第99页），它就是 $A[T]$ 在其分式域 $K(T)$ 中的整闭包。将 $u$ 换成 $T\cdot1-u\otimes1$，将 $A$ 换成 $A[T]$，便只须证明 $\det u$ 在 $A$ 上为整。若 $n$ 为 $E$ 的秩，则
\[
\det u=\det\bigl(\textstyle\bigwedge^n u\bigr),\qquad
\bigl(\textstyle\bigwedge^n u\bigr)\otimes1=\textstyle\bigwedge^n(u\otimes1),
\]
故可设 $n=1$。此时 $u\mapsto\det u$ 是从 $A$-模 $\Hom_A(E,E)$ 到 $K$ 的同态。由于 $E$ 为有限型，若它有 $n$ 个生成元，则 $\Hom_A(E,E)$ 同构于有限型 $A$-模 $E^n$ 的一个子模。因此，各 $\det u$ 属于 $K$ 的某个有限型子 $A$-模，从而在 $A$ 上为整。\transnote{末句不应理解为“有限型模的任意子模均为有限型”，因为这里未假设 $A$ 诺特。此处保持原文论证；秩一情形的整性也可由有限型模上的行列式技巧得到。}
\end{proof}

\egastatement{6.4.4}{推论}
在（6.4.3）的假设下，若进一步假设 $A$ 正规，则各 $\sigma_i(u)$，特别是 $\Tr u$ 和 $\det u$，都属于 $A$。

\egastatement{6.4.5}{命题}
设 $A$ 为整环，$E$ 为秩 $n$ 的有限型 $A$-模，$u$ 为 $E$ 的自同态，且各 $\sigma_i(u)$ 属于 $A$。若 $u$ 是 $E$ 的自同构，则 $\det u$ 必在 $A$ 中可逆；当 $E$ 无挠时，这个条件也是充分的。

条件充分，因为假设与公式（6.4.1.4）、（6.4.1.5）表明 $(\det u)^{-1}Q(u)$ 为 $u$ 的逆；由于 $E$ 无挠，这些公式在 $E$ 中成立，而不只在 $E\otimes_A K$ 中成立。
条件必要，因为若 $u$ 可逆，则由（6.4.3），$\det(u^{-1})$ 属于 $A$ 在 $K$ 中的整闭包 $A'$，并且是 $\det u$ 的逆。于是由下面的引理即得断言。

\egastatement{6.4.5.1}{引理}
设 $A$ 为环 $A'$ 的子环，且 $A'$ 在 $A$ 上为整。若 $x\in A$ 在 $A'$ 中可逆，则它在 $A$ 中也可逆。
\begin{proof}
否则，$x$ 属于 $A$ 的某个极大理想 $\mm$。由 Cohen--Seidenberg 第一定理（[13]，第一卷，第257页，定理3），存在 $A'$ 的极大理想 $\mm'$，使 $\mm=\mm'\cap A$。于是 $x\in\mm'$，与其可逆性矛盾。
\end{proof}

\egastatement{6.4.6}{推论}
设 $A$ 为整闭整环，$E$ 为有限型无挠 $A$-模，$u$ 为 $E$ 的自同态。则 $u$ 是 $E$ 的自同构，当且仅当 $\det u$ 在 $A$ 中可逆。
\begin{proof}
由（6.4.4）和（6.4.5）即得。
\end{proof}

\egastatement{6.4.7}{注记}
以后还要用到上述结果的一个推广。设 $A$ 为约化诺特环，$\pp_\alpha$（$1\leq\alpha\leq r$）为其极小素理想，$K_\alpha$ 为整环 $A/\pp_\alpha$ 的分式域，$K$ 为 $A$ 的全分式环，它是各域 $K_\alpha$ 的直积。设 $E$ 为有限型 $A$-模，并假设 $E\otimes_A K$ 是秩 $n$ 的自由 $K$-模；这里这一点不再是其余假设的推论。等价地，各 $K_\alpha$-向量空间 $E_\alpha=E\otimes_A K_\alpha$ 的维数均为 $n$。若 $u$ 为 $E$ 的自同态，仍令 $P(u,T)=P(u\otimes1,T)$、$\sigma_j(u)=\sigma_j(u\otimes1)$，特别地令 $\det u=\det(u\otimes1)$。各 $\sigma_j(u)$ 因而是 $K$ 的元素。
显然，$E\otimes_A K$ 是各 $E_\alpha$ 的直和，且各直和项在 $u\otimes1$ 下稳定；$u\otimes1$ 在 $E_\alpha$ 上的限制，就是 $u$ 扩张得到的自同态 $u_\alpha$。故 $\sigma_j(u)$ 在各 $K_\alpha$ 中的分量为 $\sigma_j(u_\alpha)$。因为 $A$ 在 $K$ 中的整闭包是它在各 $K_\alpha$ 中的整闭包的直积，各 $\sigma_j(u)$ 都在 $A$ 上为整。

\egastatement{6.4.7.1}{引理}
由所有 $\sigma_j(u)$（$1\leq j\leq n$，$u$ 遍历 $\Hom_A(E,E)$）生成的 $K$ 的子 $A$-代数，是有限型 $A$-模。
\begin{proof}
只须证明由各 $P(u,T)$ 生成的 $K[T]$ 的子 $A[T]$-代数为有限型 $A[T]$-模。事实上，若各多项式 $F_i(T)$（$1\leq i\leq m$）构成这个 $A[T]$-模的生成系，则各 $F_i(T)$ 的系数在 $A$ 上为整，因而它们生成一个作为 $A$-模为有限型的 $A$-代数（[13]，第一卷，第255页，定理1）。
所以可以将 $A$ 换成诺特环 $A[T]$，将 $E$ 换成 $E'=E\otimes_A A[T]$；有 $E'\otimes_{A[T]}K[T]=E\otimes_A K[T]$，这是秩 $n$ 的自由 $K[T]$-模。回到原来的记号，只须证明当 $u$ 遍历 $\Hom_A(E,E)$ 时，各 $\det u$ 生成的 $A$-模为有限型。因为有限型 $A$-模的子模仍为有限型，进一步只须证明：当 $v$ 遍历 $\bigwedge^nE$ 的自同态时，各 $\det v$ 生成的 $A$-模为有限型。也就是说，又化归 $n=1$。但此时 $\Hom_A(E,E)$ 为有限型 $A$-模，而 $v\mapsto\det v$ 是它到 $K$ 的一个同态，故得结论。
\end{proof}

令 $F$ 为典范同态 $E\to E\otimes_A K$ 的核，$E_0=E/F$。如前可知，$E\otimes_A K$ 等同于 $E_0\otimes_A K$，有 $u(F)\subset F$；若 $u_0$ 为 $u$ 取商后得到的自同态，则 $u\otimes1$ 等同于 $u_0\otimes1$，每个 $j$ 都满足 $\sigma_j(u)=\sigma_j(u_0)$。若 $F=0$，将 $E$ 等同于 $E\otimes_A K$ 的子模，将 $u^i$ 看作同态 $E\to E\otimes_A K$ 后，公式（6.4.1.3）和（6.4.1.5）仍有意义且成立。因此，（6.4.5）及其证明也推广到这个情形。

\egapar{6.4.8}
设 $(X,\A)$ 为赋环空间，$\E$ 为有限秩局部自由 $\A$-模，$u$ 为 $\E$ 的自同态。按假设，$X$ 有一个拓扑基 $\mathfrak B$，使对每个 $V\in\mathfrak B$，$\E|V$ 同构于 $\A^n|V$，其中 $n$ 可以随 $V$ 改变。对每个 $V\in\mathfrak B$，$u_V$ 为自由 $\Gamma(V,\A)$-模 $\Gamma(V,\E)$ 的自同态，所以 $\det u_V$ 有定义且属于 $\Gamma(V,\A)$。若 $e_1,\ldots,e_n$ 为 $\Gamma(V,\E)$ 的一组基，它们在每个开集 $W\subset V$ 上的限制构成 $\Gamma(W,\E)$ 在 $\Gamma(W,\A)$ 上的一组基。因此 $\det u_W$ 是 $\det u_V$ 在 $W$ 上的限制。
于是存在唯一的 $\A$ 在 $X$ 上的截面，记为 $\det u$，称为 $u$ 的\emph{行列式}，使其在每个 $V\in\mathfrak B$ 上的限制为 $\det u_V$。显然，对每个 $x\in X$，有 $(\det u)_x=\det u_x$。对 $\E$ 的两个自同态 $u,v$，有
\begin{equation}\tag{6.4.8.1}
\det(v\circ u)=(\det v)(\det u),
\end{equation}
以及
\begin{equation}\tag{6.4.8.2}
\det(1_\E)=1_\A.
\end{equation}
若 $\E$ 的秩恒等于 $n$，特别是 $X$ 连通时（0，5.4.1），则对每个 $s\in\Gamma(X,\A)$，有
\begin{equation}\tag{6.4.8.3}
\det(su)=s^n\det u.
\end{equation}
应注意，$n\geq1$ 时 $\det(0)=0_\A$，而 $n=0$ 时 $\det(0)=1_\A$。此外，$u$ 为 $\E$ 的自同构，当且仅当 $\det u$ 在 $\Gamma(X,\A)$ 中可逆。
若 $\E$ 的秩为常数，也可同样定义各初等对称函数 $\sigma_i(u)\in\Gamma(X,\A)$，关系（6.4.1.3）至（6.4.1.5）仍成立。

这样得到一个乘法幺半群同态
$\det:\Hom_\A(\E,\E)\to\Gamma(X,\A)$。由于按定义
$\Hom_\A(\E,\E)=\Gamma(X,\shHom_\A(\E,\E))$，在此定义中可把 $X$ 换成任意开集 $U\subset X$，从而立即得到乘法幺半群层的同态
$\det:\shHom_\A(\E,\E)\to\A$。当 $\E$ 的秩为常数时，同样定义集合层同态 $\sigma_i:\shHom_\A(\E,\E)\to\A$；在 $i=1$ 时，$\sigma_1=\Tr$ 是 $\A$-模同态。

设 $(Y,\B)$ 为第二个赋环空间，$f:(X,\A)\to(Y,\B)$ 为赋环空间态射。若 $\F$ 为局部自由 $\B$-模，则 $f^*(\F)$ 为局部自由 $\A$-模；$\F$ 秩恒定时二者的秩相同（0，5.4.5）。对 $\F$ 的每个自同态 $v$，$f^*(v)$ 是 $f^*(\F)$ 的自同态。由定义立即可知，$\det f^*(v)$ 是 $\A=f^*(\B)$ 在 $X$ 上的、典范地对应于 $\det v\in\Gamma(Y,\B)$ 的截面。也可说同态
$f^*(\det):f^*(\shHom_\B(\F,\F))\to f^*(\B)=\A$ 是复合
\[
f^*(\shHom_\B(\F,\F))
\xrightarrow{\gamma^\sharp}\shHom_\A(f^*(\F),f^*(\F))
\xrightarrow{\det}\A
\]
（0，4.4.6）。各 $\sigma_i$ 也有相应的结果。

\egapar{6.4.9}
现设 $X$ 为局部整预概形。其有理函数层 $\R(X)$ 为局部常值的域层（I，7.3.4），并且作为 $\OO_X$-模拟凝聚。若 $\E$ 为有限型拟凝聚 $\OO_X$-模，则 $\E'=\E\otimes_{\OO_X}\R(X)$ 为局部自由 $\R(X)$-模（I，7.3.6）。对 $\E$ 的每个自同态 $u$，$u\otimes1_{\R(X)}$ 是 $\E'$ 的自同态，而 $\det(u\otimes1)$ 是 $\R(X)$ 在 $X$ 上的截面，仍称为 $u$ 的行列式，记为 $\det u$。由（6.4.3），$\det u$ 是 $\OO_X$ 在 $\R(X)$ 中的整闭包（6.3.2）的一个截面。若 $X$ 进一步正规，则 $\det u$ 是 $\OO_X$ 在 $X$ 上的截面；再设 $\E$ 无挠，则由（6.4.6），$u$ 为 $\E$ 的自同构，当且仅当 $\det u$ 可逆。公式（6.4.8.1）至（6.4.8.3）仍成立。
将 $u\mapsto\det u$ 应用于 $\shHom_{\OO_X}(\E,\E)$ 的各截面模，得到层同态
$\det:\shHom_{\OO_X}(\E,\E)\to\R(X)$；当 $X$ 正规时，其值属于 $\OO_X$。当 $\E'$ 的秩为常数时，对其他初等对称函数 $\sigma_i(u)$ 也有类似的定义和结果；若 $X$ 进一步正规，则各 $\sigma_i(u)$ 为 $\OO_X$ 在 $X$ 上的截面。

最后，设 $X,Y$ 为两个整预概形，$f:X\to Y$ 为占优态射。已知存在典范同态 $f^*(\R(Y))\to\R(X)$（I，7.3.8）\origfootnote{见《勘误与增补：清单一》。}，因此对每个有限型拟凝聚 $\OO_Y$-模 $\F$，得到典范同态
\[
\theta:f^*(\F\otimes_{\OO_Y}\R(Y))\longrightarrow
f^*(\F)\otimes_{\OO_X}\R(X).
\]
若 $v$ 为 $\F$ 的自同态，则 $f^*(v\otimes1_{\R(Y)})$ 是 $f^*(\F\otimes_{\OO_Y}\R(Y))$ 的自同态，且下图交换：
\[
\begin{tikzcd}[column sep=large]
f^*(\F\otimes\R(Y))\arrow[r,"f^*(v\otimes1)"]\arrow[d,"\theta"']&f^*(\F\otimes\R(Y))\arrow[d,"\theta"]\\
f^*(\F)\otimes\R(X)\arrow[r,"f^*(v)\otimes1"']&f^*(\F)\otimes\R(X).
\end{tikzcd}
\]
容易推出，$\det f^*(v)$ 是截面 $\det v$ 经 $f^*(\R(Y))\to\R(X)$ 所得到的典范像。事实上，立即可化归 $X=\Spec(A)$、$Y=\Spec(B)$ 仿射，$A,B$ 为整环，分式域分别为 $K,L$。同态 $B\to A$ 单射，故延拓为单同态 $L\to K$。若 $\F=\wt M$，其中 $M$ 为有限型 $B$-模，则 $M\otimes_B L$ 在 $L$ 上的秩，等于 $(M\otimes_B A)\otimes_A K$ 在 $K$ 上的秩；对 $M$ 的每个自同态 $u$，$\det((u\otimes1)\otimes1)$ 是 $\det(u\otimes1)$ 在 $K$ 中的像。结论即得。

\egapar{6.4.10}
最后设 $X$ 为局部诺特约化预概形。此时 $\R(X)$ 仍为拟凝聚 $\OO_X$-模（I，7.3.4）。设 $\E$ 为凝聚 $\OO_X$-模，且 $\E'=\E\otimes_{\OO_X}\R(X)$ 局部自由且为有限秩。由（6.4.7），若 $\E'$ 秩恒定，则对 $\E$ 的每个自同态 $u$，可以定义 $\sigma_j(u)$，它们为 $\R(X)$ 在 $X$ 上的截面。即使不再假设 $\E'$ 秩恒定，仍可定义同态
$\det:\shHom_{\OO_X}(\E,\E)\to\R(X)$。

\subsection{可逆层的范数}
\src{125--130}

\egapar{6.5.1}
设 $(X,\A)$ 为赋环空间，$\B$ 为交换 $\A$-代数。$\A$-模 $\B$ 典范地等同于 $\shHom_\A(\B,\B)$ 的一个子 $\A$-模：$\B$ 在开集 $U\subset X$ 上的截面 $f$ 对应于乘以该截面的自同态。若 $(X,\A)$ 与 $\B$ 满足（6.4.8）、（6.4.9）或（6.4.10）中列出的某一条件，就可定义 $\det(f)$，在某些情形也可定义 $\sigma_i(f)$。它们为 $\A$ 或 $\R(X)$ 在 $U$ 上的截面，分别称为 $f$ 的\emph{范数}和\emph{初等对称函数}；范数记为 $\mathrm N_{\B/\A}(f)$。以下假设下列条件之一成立：
\begin{enumerate}[label=\textup{(\Roman*)}]
\item $\B$ 为有限秩局部自由 $\A$-模。
\item $(X,\A)$ 为局部诺特约化预概形，$\B$ 为凝聚 $\A$-模，$\B\otimes_\A\R(X)$ 为局部自由 $\R(X)$-模，而且对每个开集 $U\subset X$ 及每个 $f\in\Gamma(U,\B)$，$\mathrm N_{\B/\A}(f)$ 为 $\A$ 在 $U$ 上的截面。
\end{enumerate}
当局部诺特预概形 $X$ 正规时，最后的条件自动成立（6.4.9）。
另一方面，$\B\otimes_\A\R(X)$ 局部自由的假设，还可表述如下。令 $X_\alpha$ 为 $X$ 中以各不可约分支为底空间的约化闭子预概形（I，5.2.1）；它们为局部诺特整预概形。每个 $x\in X$ 只属于有限个 $X_\alpha$；又 $\B\otimes_\A\R(X_\alpha)$ 是秩恒为 $k_\alpha$ 的局部自由 $\R(X_\alpha)$-模（I，7.3.6）。$\B\otimes_\A\R(X)$ 为局部自由 $\R(X)$-模，就是说：对每个 $x\in X$，对应于所有满足 $x\in X_\alpha$ 的指标的秩 $k_\alpha$ 相等。这个问题是局部的，因而可化归 $X=\Spec(C)$，其中 $C$ 为约化诺特环，$\B=\wt D$，其中 $D$ 是作为 $C$-模为有限型的 $C$-代数。若 $\pp_i$（$1\leq i\leq m$）为 $C$ 的极小素理想，则 $C$ 的全分式环 $L$ 是各整环 $C/\pp_i$ 的分式域 $K_i$ 的直积，而 $D\otimes_C L$ 为各 $D\otimes_C K_i$ 的直和，故得结论。

在假设（I）或（II）下，这样定义了一个乘法幺半群层同态
$\mathrm N_{\B/\A}:\B\to\A$，不致混淆时也记为 $\mathrm N$，称为\emph{范数同态}。对 $\B$ 在同一开集 $U$ 上的两个截面 $f,g$，有
\begin{equation}\tag{6.5.1.1}
\mathrm N_{\B/\A}(fg)=\mathrm N_{\B/\A}(f)\mathrm N_{\B/\A}(g),
\end{equation}
这是相应的 $\A$ 在 $U$ 上的截面间的等式；并且
\begin{equation}\tag{6.5.1.2}
\mathrm N_{\B/\A}(1_\B)=1_\A.
\end{equation}
最后，若秩恒为 $n$，则对 $\A$ 在 $U$ 上的每个截面 $s$，有
\begin{equation}\tag{6.5.1.3}
\mathrm N_{\B/\A}(s\cdot1_\B)=s^n.
\end{equation}
令 $s=0_\A$，在 $n\geq1$ 时得 $\mathrm N(0_\B)=0_\A$；在 $n=0$ 时得 $\mathrm N(0_\B)=\mathrm N(1_\B)=1_\A$。

在假设（I）下，$f\in\Gamma(U,\B)$ 可逆，当且仅当 $\mathrm N(f)\in\Gamma(U,\A)$ 可逆。在假设（II）下，这个条件必要；若进一步假设 $\B\to\B\otimes_\A\R(X)$ 单射，且满足下面更强的假设，则由（6.4.7），该条件也充分：

\noindent\textbf{（II bis）} $(X,\A)$ 为局部诺特约化预概形，$\B$ 为凝聚 $\A$-模，且 $\B\otimes_\A\R(X)$ 为局部自由 $\R(X)$-模；对每个使 $\B\otimes_\A\R(X)|U$ 在 $\R(X)|U$ 上秩恒为 $n$ 的开集 $U$，以及每个 $f\in\Gamma(U,\B)$，各 $\sigma_j(f)$（$1\leq j\leq n$）都是 $\A$ 在 $U$ 上的截面。

同样应注意，当 $X$ 正规时，这个条件成立。

\egapar{6.5.2}
假设（6.5.1）中的（I）或（II）成立，且 $\LL'$ 为可逆 $\B$-模。以下将典范地给它配上一个可逆 $\A$-模 $\LL$，该对象在唯一同构意义下确定。令 $\A^*$、$\B^*$ 分别为 $\A,\B$ 的子层，使对每个开集 $U\subset X$，$\Gamma(U,\A^*)$、$\Gamma(U,\B^*)$ 分别为环 $\Gamma(U,\A)$、$\Gamma(U,\B)$ 的可逆元集合。它们为乘法群层，而 $\mathrm N_{\B/\A}$ 在 $\B^*$ 上的限制为群层同态 $\B^*\to\A^*$（6.5.1）。
令 $\Omega$ 为所有二元组 $(U_\lambda,\eta_\lambda)$ 构成的集合，其中 $U_\lambda$ 为 $X$ 的开集，$\eta_\lambda:\LL'|U_\lambda\xrightarrow{\sim}\B|U_\lambda$ 为 $(\B|U_\lambda)$-模同构。按假设，各 $U_\lambda$ 覆盖 $X$。对任意两个指标 $\lambda,\mu$，令
\[
\omega_{\lambda\mu}=(\eta_\lambda|U_\lambda\cap U_\mu)
\circ(\eta_\mu|U_\lambda\cap U_\mu)^{-1}.
\]
它是 $\B|U_\lambda\cap U_\mu$ 的自同构，典范地等同于 $\B^*$ 在 $U_\lambda\cap U_\mu$ 上的截面；族 $(\omega_{\lambda\mu})$ 是覆盖 $\mathfrak U=(U_\lambda)$ 的一个取值于 $\B^*$ 的 $1$-上闭链（0，5.4.7）。由于 $\mathrm N_{\B/\A}:\B^*\to\A^*$ 为同态，$(\mathrm N_{\B/\A}\omega_{\lambda\mu})$ 是 $\mathfrak U$ 的一个取值于 $\A^*$ 的 $1$-上闭链，因而它确定一个可逆 $\A$-模 $\LL$（0，5.4.7），在唯一同构意义下确定。将它记为 $\mathrm N_{\B/\A}(\LL')$，并称为可逆 $\B$-模 $\LL'$ 的\emph{范数}。

取 $\Omega$ 的一个子集，使相应的各 $U_\lambda$ 仍覆盖 $X$，记该子覆盖为 $\mathfrak V$。将 $(\omega_{\lambda\mu})$ 限制于 $\mathfrak V$，得到的 $(\mathrm N_{\B/\A}\omega_{\lambda\mu})$ 是上述 $\A^*$-值上闭链的限制。由它确定的可逆 $\A$-模，显然典范同构于 $\mathrm N_{\B/\A}(\LL')$。因此，可以使用 $\mathfrak U$ 的任意子覆盖来定义该模。这立即表明：若 $\LL'_1,\LL'_2$ 是可逆 $\B$-模，则由（6.5.1.1）、（6.5.1.2），有
\begin{equation}\tag{6.5.2.1}
\mathrm N(\LL'_1\otimes_\B\LL'_2)=\mathrm N(\LL'_1)\otimes_\A\mathrm N(\LL'_2),
\end{equation}
以及
\begin{equation}\tag{6.5.2.2}
\mathrm N_{\B/\A}(\B)=\A,
\end{equation}
和
\begin{equation}\tag{6.5.2.3}
\mathrm N(\LL'^{-1})=(\mathrm N(\LL'))^{-1},
\end{equation}
这些等式均在典范同构意义下成立。又由（6.5.1.3），若 $\LL$ 为可逆 $\A$-模，且在情形（I）中 $\B$ 在 $\A$ 上秩恒为 $n$，或在情形（II）中 $\B\otimes_\A\R(X)$ 在 $\R(X)$ 上秩恒为 $n$，则在典范同构意义下有
\begin{equation}\tag{6.5.2.4}
\mathrm N_{\B/\A}(\LL\otimes_\A\B)=\LL^{\otimes n}.
\end{equation}

\egapar{6.5.3}
现说明 $\mathrm N_{\B/\A}(\LL')$ 是可逆 $\B$-模范畴上的协变函子。设 $h':\LL'_1\to\LL'_2$ 为可逆 $\B$-模间的同态。取开覆盖 $\mathfrak V=(U_\lambda)$，使对每个 $\lambda$，存在 $(\B|U_\lambda)$-模同构
$\eta_\lambda^{(i)}:\LL'_i|U_\lambda\xrightarrow{\sim}\B|U_\lambda$，$i=1,2$。于是存在 $\B|U_\lambda$ 的自同态 $h'_\lambda$，使
$h'_\lambda\circ\eta_\lambda^{(1)}=\eta_\lambda^{(2)}\circ(h'|U_\lambda)$，并且可以将 $h'_\lambda$ 等同于 $\B$ 在 $U_\lambda$ 上的一个截面（0，5.1.1）。因此，对每个指标对 $\lambda,\mu$，
$(\eta_\lambda^{(2)})^{-1}\circ h'_\lambda\circ\eta_\lambda^{(1)}$
与
$(\eta_\mu^{(2)})^{-1}\circ h'_\mu\circ\eta_\mu^{(1)}$
在 $U_\lambda\cap U_\mu$ 上的限制相同。令 $(\omega_{\lambda\mu}^{(1)})$、$(\omega_{\lambda\mu}^{(2)})$ 分别为 $\LL'_1,\LL'_2$ 相应的 $\B^*$-值 $1$-上闭链，得到
\[
\omega_{\lambda\mu}^{(2)}h'_\mu=h'_\lambda\omega_{\lambda\mu}^{(1)}.
\]
令 $h_\lambda=\mathrm N(h'_\lambda)$，则有相应的关系
\[
\mathrm N(\omega_{\lambda\mu}^{(2)})h_\mu
=h_\lambda\mathrm N(\omega_{\lambda\mu}^{(1)}).
\]
所以各 $h_\lambda$ 定义同态 $\mathrm N(\LL'_1)\to\mathrm N(\LL'_2)$，记为 $\mathrm N_{\B/\A}(h')$ 或 $\mathrm N(h')$。在假设（I）下，由于问题是局部的，$h'$ 为同构，当且仅当 $\mathrm N_{\B/\A}(h')$ 为同构。在假设（II）下，该条件仍必要；若（II bis）成立，且 $\B\to\B\otimes_\A\R(X)$ 单射，则也充分。

特别取 $\LL'_1=\B$，此时同态 $\B\to\LL'$ 等同于 $\LL'$ 在 $X$ 上的截面（0，5.1.1），从而得到典范映射
\[
\mathrm N_{\B/\A}:\Gamma(X,\LL')\longrightarrow
\Gamma(X,\mathrm N_{\B/\A}(\LL')).
\]
再由（6.5.1.1），若 $f_1\in\Gamma(X,\LL'_1)$、$f_2\in\Gamma(X,\LL'_2)$，则
\begin{equation}\tag{6.5.3.1}
\mathrm N(f_1\otimes f_2)=\mathrm N(f_1)\otimes\mathrm N(f_2).
\end{equation}
对每个可逆 $\A$-模 $\LL$ 及每个 $f\in\Gamma(X,\LL)$，考虑（6.5.2.4），有
\begin{equation}\tag{6.5.3.2}
\mathrm N_{\B/\A}(f\otimes1_\B)=f^{\otimes n},
\end{equation}
其中在假设（I）下 $\B$ 秩恒为 $n$，或在假设（II）下 $\B\otimes_\A\R(X)$ 秩恒为 $n$。最后，与 $\LL'$ 在 $X$ 上的截面 $f'$ 对应的同态 $\B\to\LL'$ 为同构，当且仅当对每个 $x\in X$，$f'_x$ 是 $\LL'_x$ 的一组基。在假设（I）下，这又等价于对每个 $x$，$(\mathrm N(f'))_x$ 是 $(\mathrm N(\LL'))_x$ 的一组基。在假设（II）下，这个条件仍必要；当 $\B$ 满足（II bis）且 $\B\to\B\otimes_\A\R(X)$ 单射时，也充分。

\egapar{6.5.4}
设 $(X,\A)$、$(X',\A')$ 为两个赋环空间，$f:X'\to X$ 为态射，$\B$ 为 $\A$-代数，$\B'=f^*(\B)$ 为其逆像 $\A'$-代数。假设下列条件之一成立：
\begin{enumerate}[label=\arabic*\textsuperscript{o}]
\item $\B$ 满足（6.5.1）的假设（I）。
\item $(X,\A)$ 与 $\B$ 满足（6.5.1）的假设（II），$(X',\A')$ 为局部诺特约化预概形；分别令 $X_\alpha$、$X'_\beta$ 为以 $X,X'$ 各不可约分支为底空间的约化闭子预概形，则 $f$ 在每个 $X'_\beta$ 上的限制是它到某个 $X_\alpha$ 的占优态射。
\end{enumerate}
在这些条件下，$\B'$ 分别满足（6.5.1）的（I）或（II）。第一个断言是立即的。为证第二个，只须说明：对每个 $x'\in X'$，各 $\B'\otimes_{\A'}\R(X'_\beta)$ 的秩在所有满足 $x'\in X'_\beta$ 的指标 $\beta$ 上相等。若 $f|X'_\beta$ 占优于 $X_\alpha$，则 $\B'\otimes_{\A'}\R(X'_\beta)$ 的秩等于 $\B\otimes_\A\R(X_\alpha)$ 的秩；如（6.4.9），化归仿射情形即可立即看出。于是由假设（II）和（6.5.1）得到断言。
由（6.4.8）、（6.4.9）、（6.4.10），若 $s$ 为 $\B$ 在开集 $U\subset X$ 上的截面，$s'$ 为 $\B'$ 在 $f^{-1}(U)$ 上的相应截面，则 $\mathrm N_{\B'/\A'}(s')$ 是 $\A'$ 在 $f^{-1}(U)$ 上对应于 $\A$ 的截面 $\mathrm N_{\B/\A}(s)$ 的截面。
若 $\M$ 为可逆 $\B$-模，令 $\M'=f^*(\M)$，它为可逆 $\B'$-模，则由以上结论，在典范同构意义下有
\[
\mathrm N_{\B'/\A'}(\M')=f^*(\mathrm N_{\B/\A}(\M)).
\]

\egapar{6.5.5}
从现在起设 $(X,\A)$ 为预概形。已知，给定一个作为 $\A$-模为有限型的拟凝聚 $\A$-代数 $\B$，等价于给定一个有限态射 $g:X'\to X$，使 $g_*(\OO_{X'})=\B$，并在 $X$-同构意义下确定（6.1.2、1.3.1）。此外，给定一个拟凝聚 $\OO_{X'}$-模 $\F'$，等价于给定一个拟凝聚 $\B$-模 $\F$，使 $g_*(\F')=\F$（1.4.3）；而 $\F'$ 可逆，当且仅当 $\F$ 可逆（6.1.12）。因此，要用有限态射 $g$ 的语言表述前面的结果，须假设 $g_*(\OO_{X'})$ 为有限型局部自由 $\OO_X$-模，或假设 $(X,\OO_X)$ 与 $g_*(\OO_{X'})$ 满足（II）。对每个可逆 $\OO_{X'}$-模 $\LL'$，令
\begin{equation}\tag{6.5.5.1}
\mathrm N_{X'/X}(\LL')=
\mathrm N_{g_*(\OO_{X'})/\OO_X}\bigl(g_*(\LL')\bigr),
\end{equation}
称它为 $\LL'$ \emph{相对于 $g$ 的范数}。同样，若 $h':\LL'_1\to\LL'_2$ 为可逆 $\OO_{X'}$-模间的同态，令
\begin{equation}\tag{6.5.5.2}
\mathrm N_{X'/X}(h')=
\mathrm N_{g_*(\OO_{X'})/\OO_X}\bigl(g_*(h')\bigr):
\mathrm N_{X'/X}(\LL'_1)\longrightarrow\mathrm N_{X'/X}(\LL'_2).
\end{equation}
特别地，令 $\LL'_1=\OO_{X'}$，得到典范映射
\begin{equation}\tag{6.5.5.3}
\mathrm N_{X'/X}:\Gamma(X',\LL')\longrightarrow
\Gamma(X,\mathrm N_{X'/X}(\LL')).
\end{equation}
大部分语言转换留给读者；这里只明确写出以下结果。

\egastatement{6.5.6}{命题}
设 $g:X'\to X$ 为有限态射，并假设 $g_*(\OO_{X'})$ 为局部自由 $\OO_X$-模，或假设 $(X,\OO_X)$ 与 $g_*(\OO_{X'})$ 满足（II bis），特别是 $X$ 局部诺特且正规时。对于可逆 $\OO_{X'}$-模间的同态 $h':\LL'_1\to\LL'_2$，在第一种假设下，$h'$ 为同构，当且仅当 $\mathrm N_{X'/X}(h')$ 为同构。在第二种假设下，这个条件仍必要；若同态
$g_*(\OO_{X'})\to g_*(\OO_{X'})\otimes_{\OO_X}\R(X)$ 单射，则也充分。

这里使用了如下事实：$g_*(h')$ 为同构，当且仅当 $h'$ 为同构（1.4.2）。

\egastatement{6.5.7}{推论}
设 $g:X'\to X$ 为有限态射，并假设 $g_*(\OO_{X'})$ 为局部自由 $\OO_X$-模，或假设 $(X,\OO_X)$ 与 $g_*(\OO_{X'})$ 满足（II bis）且
$g_*(\OO_{X'})\to g_*(\OO_{X'})\otimes_{\OO_X}\R(X)$ 单射。
设 $\LL'$ 为可逆 $\OO_{X'}$-模，$f'$ 为它在 $X'$ 上的截面，$f=\mathrm N_{X'/X}(f')$ 为 $\LL=\mathrm N_{X'/X}(\LL')$ 在 $X$ 上的相应截面。则
\[
g(X'-X'_{f'})=X-X_f,
\]
而且 $X_f$ 是满足 $g^{-1}(U)\subset X'_{f'}$ 的最大开集 $U\subset X$。
\begin{proof}
$g(X'-X'_{f'})$ 在 $X$ 中闭（6.1.10），所以只须证明最后的断言。$U\subset X_f$ 等价于由 $f|U$ 定义的同态 $\OO_X|U\to\LL|U$ 为同构。由（6.5.6），这又等价于由 $f'|g^{-1}(U)$ 定义的同态
$\OO_{X'}|g^{-1}(U)\to\LL'|g^{-1}(U)$ 为同构，即 $g^{-1}(U)\subset X'_{f'}$。
\end{proof}

\egastatement{6.5.8}{命题}
设 $g:X'\to X$ 为有限态射，$f:Y\to X$ 为态射。令 $Y'=X'_{(Y)}$，$g'=g_{(Y)}$，$f'=f_{(X')}$，因而有交换图
\[
\begin{tikzcd}
X'\arrow[d,"g"']&Y'\arrow[l,"f'"']\arrow[d,"g'"]\\
X&Y.\arrow[l,"f"]
\end{tikzcd}
\]
假设 $g_*(\OO_{X'})$ 局部自由，或假设 $(X,\OO_X)$ 与 $g_*(\OO_{X'})$ 满足（II），$Y$ 为局部诺特约化预概形，且 $f$ 在 $Y$ 的每个不可约分支上的限制占优于 $X$ 的一个不可约分支。则对每个可逆 $\OO_{X'}$-模 $\LL'$，在典范同构意义下有
\[
\mathrm N_{Y'/Y}(f'^*(\LL'))=f^*(\mathrm N_{X'/X}(\LL')).
\]
\begin{proof}
由（1.5.2），有 $f^*(g_*(\LL'))=g'_*(f'^*(\LL'))$；特别地，$g'_*(\OO_{Y'})=f^*(g_*(\OO_{X'}))$。所以，若 $g_*(\OO_{X'})$ 局部自由，则 $g'_*(\OO_{Y'})$ 也如此。结论由定义和（6.5.4）得出。
\end{proof}

\egastatement{6.5.9}{注记}
以后将推广上述范数概念，并建立它与除子的直接像概念之间的联系。

\subsection{应用：丰沛性的判别准则}
\src{130--135}

\egastatement{6.6.1}{命题}
设 $Y$ 为预概形，$f:X\to Y$ 为拟紧态射，$g:X'\to X$ 为有限满态射。假设 $g_*(\OO_{X'})$ 为局部自由 $\OO_X$-模，或假设 $(X,\OO_X)$ 与 $g_*(\OO_{X'})$ 满足（II bis）。则对每个相对于 $f\circ g$ 丰沛的可逆 $\OO_{X'}$-模 $\LL'$，$\LL=\mathrm N_{X'/X}(\LL')$ 相对于 $f$ 丰沛。

可以设 $Y$ 仿射（4.6.4）；此时由（4.6.6），上述命题等价于下面的推论。

\egastatement{6.6.2}{推论}
设 $X$ 为拟紧预概形，$g:X'\to X$ 为有限满态射，且 $g_*(\OO_{X'})$ 为局部自由 $\OO_X$-模，或 $(X,\OO_X)$ 与 $g_*(\OO_{X'})$ 满足（II bis）。则对每个丰沛 $\OO_{X'}$-模 $\LL'$，$\LL=\mathrm N_{X'/X}(\LL')$ 丰沛。
\begin{proof}
在第二种假设下，可以进一步设典范同态
$g_*(\OO_{X'})\to g_*(\OO_{X'})\otimes_{\OO_X}\R(X)$ 单射。事实上，令 $\J$ 为该同态的核，它是 $\B=g_*(\OO_{X'})$ 的凝聚理想（I，6.1.1）；令 $X''=\Spec(\B/\J)$，得到交换图
\[
\begin{tikzcd}
X''\arrow[r,"h"]\arrow[dr,"g'"']&X'\arrow[d,"g"]\\
&X,
\end{tikzcd}
\]
其中 $h$ 为闭浸入（1.4.10）。$\J$ 的支集为闭集（0，5.2.2），且在 $X$ 中无处稠密（I，7.4.6）。所以，对 $X$ 的一个不可约分支的泛点 $x$，存在仿射开邻域 $U$，使 $\B|U=(\B/\J)|U$。由于 $g$ 满射，$x\in g'(X'')$，故 $g'$ 占优；而 $g'$ 有限，因而满射（6.1.10）。按定义有
\[
g'_*(\OO_{X''})\otimes_{\OO_X}\R(X)
=(\B/\J)\otimes_{\OO_X}\R(X)
=g_*(\OO_{X'})\otimes_{\OO_X}\R(X).
\]
因此 $(X,\OO_X)$ 与 $g'_*(\OO_{X''})$ 满足（II bis），且
$g'_*(\OO_{X''})\to g'_*(\OO_{X''})\otimes_{\OO_X}\R(X)$ 单射。最后，$\LL''=h^*(\LL')$ 为丰沛 $\OO_{X''}$-模（4.6.13，i bis），并且
\[
\mathrm N_{X''/X}(\LL'')=\mathrm N_{X'/X}(\LL').
\]
事实上，为定义这两个可逆 $\OO_X$-模，可使用同一个仿射开覆盖 $(U_\lambda)$，使 $g_*(\LL')$、$g'_*(\LL'')$ 在 $U_\lambda$ 上的限制分别同构于 $\B|U_\lambda$、$(\B/\J)|U_\lambda$。每个同构
$\eta_\lambda:g_*(\LL')|U_\lambda\xrightarrow{\sim}\B|U_\lambda$
都典范地给出同构
$\eta'_\lambda:g'_*(\LL'')|U_\lambda\xrightarrow{\sim}(\B/\J)|U_\lambda$。
若 $(\omega_{\lambda\mu})$、$(\omega'_{\lambda\mu})$ 为这两族同构对应的 $1$-上闭链（6.5.2），则 $\omega'_{\lambda\mu}$ 是 $\omega_{\lambda\mu}$ 在
$\Gamma(U_\lambda\cap U_\mu,\B/\J)$ 中的典范像。由 $X''$ 的定义，
\[
\mathrm N_{\B/\A}(\omega_{\lambda\mu})
=\mathrm N_{(\B/\J)/\A}(\omega'_{\lambda\mu}),\qquad\A=\OO_X,
\]
故得到上述等式。

于是，在（II bis）情形下，可以设该典范同态单射。只须证明：当 $f$ 遍历各 $\LL^{\otimes n}$（$n>0$）在 $X$ 上的截面时，各 $X_f$ 构成 $X$ 的拓扑基（4.5.2）。取 $x\in X$ 及其任意邻域 $U$。因为 $g^{-1}(x)$ 有限（6.1.7）且 $\LL'$ 丰沛，存在整数 $n>0$ 和 $\LL'^{\otimes n}$ 在 $X'$ 上的截面 $f'$，使 $X'_{f'}$ 为包含在 $g^{-1}(U)$ 中的 $g^{-1}(x)$ 的邻域（4.5.4）。由于
$\LL^{\otimes n}=\mathrm N_{X'/X}(\LL'^{\otimes n})$，取 $f=\mathrm N_{X'/X}(f')$ 即可；由（6.5.7），$X-X_f=g(X'-X'_{f'})$，所以 $x\in X_f\subset U$。
\end{proof}

\egastatement{6.6.3}{推论}
在（6.6.1）的假设下，可逆 $\OO_X$-模 $\LL$ 相对于 $f$ 丰沛，当且仅当 $\LL'=g^*(\LL)$ 相对于 $f\circ g$ 丰沛。
\begin{proof}
必要性由 $g$ 仿射和（5.1.12）得出。为证充分性，可以设 $Y$ 仿射（4.6.4），于是 $X,X'$ 拟紧，$\LL'$ 丰沛（4.6.6）；要证明 $\LL$ 丰沛。在第一种假设下，$g_*(\OO_{X'})$ 的秩在邻域内为某个给定整数 $n$ 的各点构成 $X$ 中的开闭集；在第二种假设下，将它换成 $g_*(\OO_{X'})\otimes_{\OO_X}\R(X)$，结论相同。故 $X$ 是有限个这样的开集的和，并可设 $X$ 就是其中一个（4.6.17）。此时 $\mathrm N_{X'/X}(\LL')=\LL^{\otimes n}$，故由（6.6.2），$\LL^{\otimes n}$ 丰沛，从而 $\LL$ 丰沛（4.5.6）。
\end{proof}

\egastatement{6.6.4}{推论}
假设（6.6.1）的条件成立，并进一步设 $f:X\to Y$ 为有限型。则 $f$ 拟射影，当且仅当 $f\circ g$ 拟射影。若再设 $Y$ 为拟紧概形或底空间诺特的预概形，则 $f$ 射影，当且仅当 $f\circ g$ 射影。
\begin{proof}
假设说明 $f\circ g$ 为有限型。考虑拟射影态射的定义（5.3.1），第一个断言由（6.6.1）和（6.6.3）得出。结合该结果与（5.5.3，ii），剩下只须证明：已知 $f$ 拟射影时，$f$ 固有，当且仅当 $f\circ g$ 固有。此时 $f$ 分离（5.3.1）且为有限型；由于 $g$ 满射，由（5.4.2，ii）和（5.4.3，ii）即得。
\end{proof}
特别地，有：

\egastatement{6.6.5}{推论}
设 $X$ 为域 $K$ 上的有限型预概形，$K'$ 为 $K$ 的有限次扩域。则 $X$ 在 $K$ 上射影（分别地，拟射影），当且仅当 $X'=X\otimes_K K'$ 在 $K'$ 上射影（分别地，拟射影）。
\begin{proof}
必要性来自（5.3.4，iii）和（5.5.5，iii）。反之，假设该条件成立，令 $g:X'\to X$ 为典范投影。显然 $g$ 有限（6.1.5，iii）且满射（I，3.5.2，ii）。又 $g_*(\OO_{X'})$ 同构于 $\OO_X\otimes_K K'$（1.5.2），故为局部自由 $\OO_X$-模。由假设、（6.1.11）及（5.5.5，ii），$X'$ 在 $K$ 上射影（分别地，拟射影）；再由（6.6.4）即知 $X$ 在 $K$ 上射影（分别地，拟射影）。
\end{proof}
第五章将证明，$K'$ 为 $K$ 的任意扩域时，（6.6.5）仍成立。
本小节余下部分用于证明准则（6.6.11），它是（6.6.1）的一个较技术性的加强；初读时可以略过。

\egastatement{6.6.6}{引理}
设 $X$ 为约化诺特预概形，$\E$ 为凝聚 $\OO_X$-模，且 $\E\otimes_{\OO_X}\R(X)$ 为秩恒为 $n$ 的局部自由 $\R(X)$-模。则存在约化诺特预概形 $Z$ 及有限双有理态射 $h:Z\to X$，具有如下性质：集合层态射
$\sigma_i:\shHom_{\OO_X}(\E,\E)\to\R(X)$（$1\leq i\leq n$；见6.4.10）将 $\shHom_{\OO_X}(\E,\E)$ 映入凝聚 $\OO_X$-代数 $h_*(\OO_Z)$。
\begin{proof}
取 $X$ 的仿射开集 $U$，环为 $A(U)=A$，令 $E=\Gamma(U,\E)$；令 $C_U$ 为由各 $\sigma_i(u)$（$u$ 遍历 $\Hom_A(E,E)$）生成的 $R(U)$ 的子代数。由（6.4.7.1），这个 $A$-代数为有限秩。又各代数 $C_U$ 的构造显然与仿射开集 $U$ 到仿射开集 $U'\subset U$ 的限制相容。因此得到 $\R(X)$ 的一个有限子 $\OO_X$-代数 $\C$，使对每个仿射开集 $U$，有 $\Gamma(U,\C)=C_U$。取 $Z=\Spec(\C)$，$h$ 为结构态射，它有限（6.1.2）；因为 $\C$ 约化，$Z$ 为约化诺特预概形（1.3.8）。最后，按定义 $C_U$ 的全分式环为 $R(U)$；$C_U$ 包含于 $A(U)$ 在 $R(U)$ 中的整闭包，故 $A(U)$ 的极小素理想与 $C_U$ 的极小素理想一一对应（[13]，第一卷，第259页）。这说明 $h$ 双有理，证明完毕。
\end{proof}

\egastatement{6.6.7}{推论}
在（6.6.6）的假设下，设 $W\subset X$ 为开集，且对每个 $x\in W$，$X$ 在 $x$ 处正规，或 $\E_x$ 为自由 $\OO_x$-模。则可将 $h$ 定义为使其限制给出同构 $h^{-1}(W)\xrightarrow{\sim}W$。
\begin{proof}
任一假设都说明：若 $U\subset W$ 为仿射开集，采用（6.6.6）的记号，则对每个 $x\in U$，有 $(\sigma_i(u))_x\in A_x$（6.4.3），故 $\sigma_i(u)\in A$。结论由（6.6.6）中 $h$ 的定义得出。
\end{proof}

\egapar{6.6.8}
设 $X$ 为约化诺特预概形，$g:X'\to X$ 为有限满态射，故 $\B=g_*(\OO_{X'})$ 为凝聚 $\OO_X$-代数。进一步设 $\B\otimes_{\OO_X}\R(X)$ 为秩恒为 $n$ 的局部自由 $\R(X)$-模。令 $\E=\B$，应用（6.6.6）；采用该引理的记号，得到乘法幺半群层同态
$\sigma_n:\shHom_{\OO_X}(\B,\B)\to h_*(\OO_Z)$。
将它与典范同态 $\B\to\shHom_{\OO_X}(\B,\B)$（6.5.1）复合，得到乘法幺半群层同态
\begin{equation}\tag{6.6.8.1}
\mathrm N':\B=g_*(\OO_{X'})\longrightarrow h_*(\OO_Z)=\C.
\end{equation}
对每个可逆 $\OO_{X'}$-模 $\LL'$，$g_*(\LL')$ 为可逆 $\B$-模（6.1.12）。用（6.5.2）的方法，便可函子性地给 $\LL'$ 配上一个可逆 $\C$-模，记作 $\mathrm N'(g_*(\LL'))$。

\egastatement{6.6.9}{引理}
设 $X$ 为约化诺特预概形，$g:X'\to X$ 为有限满态射，且 $g_*(\OO_{X'})\otimes_{\OO_X}\R(X)$ 为秩恒定的局部自由 $\R(X)$-模。则存在约化诺特预概形 $Z$ 及有限双有理态射 $h:Z\to X$，具有如下性质：对每个丰沛 $\OO_{X'}$-模 $\LL'$，满足 $h_*(\M)=\mathrm N'(g_*(\LL'))$ 的可逆 $\OO_Z$-模 $\M$ 丰沛；这里使用（6.6.8）的记号。
\begin{proof}
先设 $\B\to\B\otimes_{\OO_X}\R(X)$ 单射。按照（6.6.6），以 $\E=g_*(\OO_{X'})$ 定义 $Z,h$。取 $z\in Z$；要证明存在整数 $m>0$ 及 $\M^{\otimes m}$ 在 $Z$ 上的截面 $t$，使 $Z_t$ 为包含 $z$ 的仿射开集（4.5.2）。令 $x=h(z)$，取 $X$ 中包含 $x$ 的仿射开集 $U$；$h^{-1}(U)$ 为 $z$ 的仿射开邻域。只须找到 $t$ 使 $z\in Z_t\subset h^{-1}(U)$，因为此时 $Z_t$ 必仿射（5.5.8）。由假设与（4.5.4），存在整数 $n>0$ 及 $\LL'^{\otimes n}$ 在 $X'$ 上的截面 $s'$，使
\begin{equation}\tag{6.6.9.1}
g^{-1}(x)\subset X'_{s'}\subset g^{-1}(U).
\end{equation}
按定义，$s'$ 也是 $g_*(\LL'^{\otimes n})$ 在 $X$ 上的截面，像（6.5.2）那样，它对应于 $\mathrm N'(g_*(\LL'^{\otimes n}))$ 的截面 $s=\mathrm N'(s')$。\transnote{底本此句及紧接着的句子省略了张量幂 $n$；这里将其写出，以与前一句 $s'\in\Gamma(X',\LL'^{\otimes n})$ 及开头所求的 $\M$ 的正张量幂保持类型一致。}
把 $s$ 看作 $\M^{\otimes n}$ 在 $Z$ 上的截面 $t$，将证明它满足要求。令
\begin{equation}\tag{6.6.9.2}
V=X-g(X'-X'_{s'}).
\end{equation}
由（6.6.9.1）和（6.1.10），$V$ 是包含 $x$ 且包含于 $U$ 的 $X$ 的开集。将证明
\begin{equation}\tag{6.6.9.3}
h^{-1}(V)\subset Z_t\subset h^{-1}(U),
\end{equation}
从而完成证明。也就是说，由满足 $s_y$ 可逆的 $y\in X$ 构成的集合 $T$，包含 $V$ 且包含于 $U$。为此，先取 $V$ 中的仿射开集 $W$。$g^{-1}(W)$ 是 $X'$ 的仿射开集，而由（6.6.9.2），对每个 $y'\in g^{-1}(W)$，$s'_{y'}$ 可逆。由关于 $X,\B$ 的假设，可以应用（6.4.7）；令 $y=g(y')$，可知 $s_y$ 可逆，即 $V\subset T$。另一方面，由（6.4.7）也有逆向结论：若 $s_y$ 可逆，则 $s'_{y'}$ 可逆；由（6.6.9.1），$y'\in g^{-1}(U)$，故 $y\in U$。这说明 $T\subset U$。

一般情形使用（6.6.2）中的同一论证：将 $X'$ 换成使
$g'_*(\OO_{X''})\to g'_*(\OO_{X''})\otimes_{\OO_X}\R(X)$ 单射的 $X''$，并将 $\LL'$ 换成满足
$\mathrm N'(g_*(\LL'))=\mathrm N'(g'_*(\LL''))$ 的丰沛 $\OO_{X''}$-模 $\LL''$。于是，在适当选择 $h$ 后，（6.6.9）在所有情形下成立。
\end{proof}

\egastatement{6.6.10}{推论}
假设（6.6.9）的条件成立。对每个使 $g^*(\LL)$ 丰沛的可逆 $\OO_X$-模 $\LL$，$h^*(\LL)$ 丰沛。
\begin{proof}
令 $\LL'=g^*(\LL)$，则 $g_*(\LL')=\LL\otimes_{\OO_X}\B$（0，5.4.10），从而由与（6.5.2.4）相同的论证，
\[
\mathrm N'(g_*(\LL'))=\C\otimes_{\OO_X}\LL^{\otimes n}.
\]
故 $\M=(h^*(\LL))^{\otimes n}$。由于 $\M$ 丰沛，$h^*(\LL)$ 也丰沛（4.5.6）。
\end{proof}

\egastatement{6.6.11}{命题}
设 $Y$ 为仿射概形，$X$ 为约化诺特预概形，$f:X\to Y$ 为拟紧态射，$g:X'\to X$ 为有限满态射。设 $W$ 是 $X$ 的一个开子集，且对每个 $x\in W$，$X$ 在 $x$ 处正规，或存在开邻域 $T\subset W$，使 $g_*(\OO_{X'})|T$ 为局部自由 $(\OO_X|T)$-模。则存在约化 $Y$-预概形 $Z$ 和有限双有理 $Y$-态射 $h:Z\to X$，其限制给出同构 $h^{-1}(W)\xrightarrow{\sim}W$，并具有如下性质：对每个使 $g^*(\LL)$ 相对于 $f\circ g$ 丰沛的可逆 $\OO_X$-模 $\LL$，$h^*(\LL)$ 相对于 $f\circ h$ 丰沛。
\begin{proof}
$Y$ 仿射，故 $g^*(\LL)$ 丰沛；只须证明适当选择 $h$ 后，$h^*(\LL)$ 丰沛（4.6.6）。将说明可以把 $g$ 换成一个有限满态射 $g':X''\to X$，使 $g'^*(\LL)$ 丰沛，且 $g'_*(\OO_{X''})\otimes_{\OO_X}\R(X)$ 为秩恒定的局部自由 $\R(X)$-模。这样便化归（6.6.10）的情形。

令 $\B=g_*(\OO_{X'})$。令 $X_i$（$1\leq i\leq n$）为 $X$ 中以各不可约分支为底空间的约化闭子预概形（I，5.2.1），它们为整预概形。令 $X'_i=g^{-1}(X_i)$ 为 $X'$ 的闭子预概形，$g_i:X'_i\to X_i$ 为 $g$ 的限制，它有限（6.1.5，iii）且满射。令 $k_i$ 为 $\OO_{X_i}$-代数 $\B_i=(g_i)_*(\OO_{X'_i})$ 的秩。因为 $\B\otimes_{\OO_X}\R(X_i)$ 为常预层（I，7.3.5），$k_i$ 也是 $\B|U$ 作为 $(\OO_X|U)$-代数的秩，其中 $U$ 为只与分支 $X_i$ 相交的任意开集。如果 $T$ 为 $X$ 的开集，且 $\B|T\simeq\OO_X^m|T$，则对所有满足 $T\cap X_i\ne\varnothing$ 的指标 $i$，有 $k_i=m$。
令 $U$ 为 $X$ 中 $\B$ 在其邻域内为局部自由 $\OO_X$-模的各点构成的开集；令 $U_j$（$1\leq j\leq s$）为其连通分支。由于 $U$ 诺特，各 $U_j$ 在 $X$ 中开，且只有有限个。令 $V_j$ 为 $X'$ 中由开集 $g^{-1}(U_j)$ 上诱导的子预概形的闭包（I，9.5.11）。根据前面的说明，对所有满足 $X_i\cap U_j\ne\varnothing$ 的指标 $i$，秩 $k_i$ 等于同一个整数 $m_j$；而同一个 $X_i$ 不会与两个不同的 $U_j$ 相交。

用 $i_\lambda$ 表示满足 $X_i\cap U=\varnothing$ 的指标。令 $k$ 为所有 $k_i$ 的乘积，$n_i=k/k_i$。如下定义 $X''$：对每个 $j$，取 $k/m_j$ 个同构于 $V_j$ 的预概形；对每个 $\lambda$，取 $k/k_{i_\lambda}=n_{i_\lambda}$ 个同构于 $X'_{i_\lambda}$ 的预概形；$X''$ 为这些预概形的和。在每个和项上取典范嵌入，得到态射 $g'':X''\to X'$。显然，$g''$ 有限且占优，因而满射，因为有限态射闭（6.1.10）。\transnote{这里保留原文关于 $g''$ 的断言。若 $X'$ 含完全位于 $X-U$ 之上的额外分支，$g''$ 的占优性不能仅由这段构造立即推出；证明实际所需的是复合 $g\circ g''$ 在 $X$ 上的满射性。此处未据此增添原文假设。}
令 $g'=g\circ g''$，它为有限满态射 $X''\to X$。有
$g'^*(\LL)=g''^*(g^*(\LL))$，故 $g'^*(\LL)$ 为丰沛 $\OO_{X''}$-模（5.1.12）。对新预概形 $X''$，与原来各 $k_i$ 相应的秩均等于 $k$。结合（I，7.3.3），立即可知，对每个仿射开集 $T\subset X$，
\[
\bigl(g'_*(\OO_{X''})\otimes_{\OO_X}\R(X)\bigr)|T
\simeq (\R(X)|T)^k.
\]
\end{proof}

\egastatement{6.6.12}{推论}
在（6.6.11）中，若 $W=X$，则可逆 $\OO_X$-模 $\LL$ 相对于 $f$ 丰沛，当且仅当 $g^*(\LL)$ 相对于 $f\circ g$ 丰沛。

\egastatement{6.6.13}{注记}
第三章将证明：若 $Y$ 诺特，$f$ 为有限型，且 $f$ 在 $X$ 中以 $X-W$ 为底空间的约化闭子预概形上的限制固有，则（6.6.12）的结论仍成立。但第五章将给出域 $K$ 上的代数概形 $X$，它的结构态射 $X\to\Spec(K)$ 不固有，其正规化 $X'$ 拟仿射，而 $X$ 本身不拟仿射。因此 $\OO_X$ 不丰沛，尽管 $\OO_{X'}$ 丰沛（5.1.2），且 $X'\to X$ 为有限满态射（6.3.10）。下一小节将说明，将“拟仿射”换成“仿射”后，这种现象不会发生。

\subsection{Chevalley 定理}
\src{135--138}
下面借助 Serre 判别准则（5.2.1）证明一个定理；C. Chevalley 曾用其他方法，在代数概形的情形证明过它。

\egastatement{6.7.1}{定理}
设 $X$ 为仿射概形，$Y$ 为诺特预概形，$f:X\to Y$ 为有限满态射。则 $Y$ 为仿射概形。

显然，$f_{\red}:X_{\red}\to Y_{\red}$ 有限（6.1.5，vi）。$X_{\red}$ 仿射，而 $Y$ 诺特，故 $Y$ 仿射等价于 $Y_{\red}$ 仿射（I，6.1.7）。因此可设 $Y$ 约化。对 $Y$ 的每个闭子集 $Y'$，存在唯一的以它为底空间的约化子预概形（I，5.1.2）；其逆像 $f^{-1}(Y')$ 典范同构于 $X\times_Y Y'$（I，4.4.1），且作为 $X$ 的闭子预概形是仿射的。$f$ 在 $f^{-1}(Y')$ 上的限制等同于 $f\times_Y1_{Y'}$，因此为有限满态射（6.1.5，iii）。由诺特归纳原理（0，2.2.2），结合（I，6.1.7），可以在以下假设下证明定理：对每个真闭子集 $Y'\ne Y$，$Y$ 中以 $Y'$ 为底空间的每个闭子预概形都仿射。
于是，对每个支集 $Z$ 不等于 $Y$ 的凝聚 $\OO_Y$-模 $\F$，有 $H^1(Y,\F)=0$。事实上，存在 $Y$ 的闭子预概形 $Y'$，底空间为 $Z$，且若 $j:Y'\to Y$ 为典范嵌入，则 $\F=j_*(j^*(\F))$（I，9.3.5）。由（5.2.3）和（I，5.1.9.2），
\[
H^1(Y,\F)=H^1(Y',j^*(\F))=0.
\]

先设 $Y$ 不可约性不成立。取 $Y$ 的一个不可约分支 $Y'$，令 $Y''=Y-Y'$；仍以 $Y'$ 表示以该底空间确定的约化闭子预概形，以 $j:Y'\to Y$ 表示典范嵌入。取凝聚 $\OO_Y$-模 $\F$，考虑典范同态
\[
\rho:\F\longrightarrow\F'=j_*(j^*(\F))
\]
（0，4.4.3）。若 $\J$ 为定义 $Y'$ 的理想层，则 $j_*(\OO_{Y'})=\OO_Y/\J$，所以由（0，5.3.10）和（0，5.3.12），$\F'$ 为凝聚 $\OO_Y$-模。因此 $\G=\Ker\rho$ 和 $\K=\Im\rho$ 也凝聚（0，5.3.4）。令 $y$ 为 $Y'$ 的泛点；$y$ 是 $Y'$ 的内点，且 $Y$ 约化，故 $\J_y=0$。按定义，$\F'_y=\F_y$，因而 $y$ 不属于 $\G$ 的闭支集。另一方面，$\F'$ 的支集，因而 $\K$ 的支集，包含于 $Y'$。所以 $\G,\K$ 的支集都不等于 $Y$，从而
$H^1(Y,\G)=H^1(Y,\K)=0$。对正合列 $0\to\G\to\F\to\K\to0$ 应用上同调长正合列，得到 $H^1(Y,\F)=0$。由 Serre 判别准则（5.2.1）得结论。

于是可设 $Y$ 不可约，因而为整。也可设 $X$ 为整：若 $X_i$ 为 $X$ 中以各不可约分支为底空间的约化闭子预概形（I，5.2.1），$f_i$ 为 $f$ 在 $X_i$ 上的限制，则至少一个 $f_i$ 占优；它有限（6.1.5），故满射（6.1.10）。又 $X_i$ 仿射，因此在命题中可用这个 $X_i$ 代替 $X$。

\egastatement{6.7.1.1}{引理}
设 $X,Y$ 为两个诺特整预概形，$x,y$ 分别为它们的泛点，$f:X\to Y$ 为有限满态射。设 $\LL$ 为可逆 $\OO_X$-模，且存在 $y$ 的仿射开邻域 $U$ 和截面 $g\in\Gamma(X,\LL)$，使 $x\in X_g\subset f^{-1}(U)$。则存在两个整数 $m>0,n>0$、同态
$u:\OO_Y^m\to f_*(\LL^{\otimes n})$，以及 $y$ 的开邻域 $V$，使 $u|V$ 为 $\OO_Y^m|V$ 到 $f_*(\LL^{\otimes n})|V$ 的同构。
\begin{proof}
令 $C$ 为 $U$ 的整环，$k=\OO_y$ 为其分式域。$f$ 有限，故 $U'=f^{-1}(U)$ 仿射（1.3.2）；令其整环为 $D$，分式域为 $K=\OO_x$。按假设，$D$ 为有限型 $C$-模（6.1.4），故 $K$ 为 $k$ 的有限次扩张。纤维
$f^{-1}(y)=X\times_Y\Spec(k(y))=X\times_Y\Spec(\OO_y)$ 等同于 $\Spec(K)$（I，3.6.6）。取 $D$ 中的元素 $s_i$（$1\leq i\leq m$）构成 $K$ 在 $k$ 上的一组基。由（I，9.3.1），存在 $n>0$，使 $\LL^{\otimes n}$ 在 $X_g$ 上的各截面 $(s_i|X_g)g^{\otimes n}$ 延拓为其在 $X$ 上的截面 $b_i$。按定义，$b_i$ 也是 $f_*(\LL^{\otimes n})$ 在 $Y$ 上的截面，故定义同态
$u:\OO_Y^m\to f_*(\LL^{\otimes n})$（0，5.1.1）。将证明 $u$ 满足要求。
有 $\LL^{\otimes n}|U'=\wt M$，其中 $M$ 为有限型 $D$-模。若 $\varphi:C\to D$ 为与限制态射 $f^{-1}(U)\to U$ 对应的单射，则 $M_{[\varphi]}$ 为有限型 $C$-模，所以
\[
f_*(\LL^{\otimes n})|U=\wt{M_{[\varphi]}}
\]
（I，1.6.3）凝聚。因为 $U$ 可为 $Y$ 的任意仿射开集，$f_*(\LL^{\otimes n})$ 凝聚。此外，$u|U=\wt\theta$，其中 $\theta:C^m\to M_{[\varphi]}$ 为 $C$-同态；而 $u_y=\theta_y$ 为同态
\[
\theta\otimes1:k^m=C^m\otimes_C k\longrightarrow M_{[\varphi]}\otimes_C k.
\]
它按定义是同构，因为各 $(b_i)_x$ 构成右端在 $k$ 上的一组基；按假设，$g_x$ 生成 $\LL_x$，因而 $g_x^{\otimes n}$ 生成 $(\LL^{\otimes n})_x$。\transnote{底本此处在取 $Y$ 的泛点茎时将若干 $k$ 印为 $K$，并在最后一句省略了 $g$ 的张量幂。这里显式写为在 $C$ 的分式域 $k$ 上张量化，与 $u_y$ 的定义相符。}
所以 $\Ker u$、$\Coker u$ 的支集都不含 $y$。这些 $\OO_Y$-模凝聚（0，5.3.3），支集因而闭（0，5.2.2），引理得证。
\end{proof}

\noindent\textbf{（6.7.1）的证明续。}
因为 $X$ 仿射，令 $\LL=\OO_X$，引理（6.7.1.1）的假设便成立（I，1.1.10）。令 $\A=\OO_Y$、$\B=f_*(\OO_X)$。由 Serre 判别准则（5.2.1.1），只须证明：对每个凝聚 $\OO_Y$-模 $\F$，有 $H^1(Y,\F)=0$；甚至只须在 $\F\subset\OO_Y$ 时证明。因为 $Y$ 为整，此时 $\F$ 无挠；事实上，下面对任意凝聚无挠 $\OO_Y$-模 $\F$ 证明 $H^1(Y,\F)=0$。
同态 $u:\A^m\to\B$ 定义同态
\[
v:\G=\shHom_\A(\B,\F)\longrightarrow
\shHom_\A(\A^m,\F)=\F^m.
\]
先证明 $v$ 单射。按假设，$\TT=\Coker u$ 的支集不与 $V$ 相交，所以它为挠 $\OO_Y$-模（I，7.4.6）。对正合列 $\A^m\to\B\to\TT\to0$，利用 $\shHom_\A$ 的左正合性，得到
\[
0\longrightarrow\shHom_\A(\TT,\F)\longrightarrow\G
\xrightarrow{v}\F^m.
\]
由于 $\F$ 无挠，$\shHom_\A(\TT,\F)=0$（0，5.2.6），所以 $v$ 单射。于是有正合列
\[
0\longrightarrow\G\longrightarrow\F^m\longrightarrow\Coker v\longrightarrow0,
\]
其中 $\G$ 和 $\Coker v$ 为凝聚 $\OO_Y$-模（0，5.3.4、5.3.5）。由上同调长正合列，只须证明
$H^1(Y,\G)=H^1(Y,\Coker v)=0$，就得到
$H^1(Y,\F^m)=(H^1(Y,\F))^m=0$，进而 $H^1(Y,\F)=0$。
限制 $v|V$ 为同构，故 $\Coker v$ 的支集不等于 $Y$。根据证明开始的归纳假设，$H^1(Y,\Coker v)=0$。另一方面，$\G$ 为凝聚 $\B$-模（I，9.6.4）；由于 $X$ 在 $Y$ 上仿射，存在拟凝聚 $\OO_X$-模 $\K$，使 $\G\simeq f_*(\K)$（1.4.3）。又因 $X$ 仿射，$H^1(X,\K)=0$（I，5.1.9.2），从而由（5.2.3），$H^1(Y,\G)=0$。定理（6.7.1）得证。\hfill$\square$

\egastatement{6.7.2}{推论}
设 $X$ 为诺特预概形，$(X_i)_{1\leq i\leq n}$ 为其底空间的有限闭覆盖。则 $X$ 仿射，当且仅当对每个 $i$，存在 $X$ 的仿射闭子预概形，以 $X_i$ 为底空间。
\begin{proof}
若存在这些仿射闭子预概形，令 $X'$ 为它们的和；显然 $X'$ 仿射。在各 $X_i$ 上取典范嵌入，定义满态射 $f:X'\to X$。由（6.7.1），只须验证 $f$ 有限，而这已在（6.1.5）中说明。
\end{proof}

\egastatement{6.7.3}{推论}
诺特预概形 $X$ 仿射，当且仅当以 $X$ 各不可约分支为底空间的约化闭子预概形都仿射。

\section{赋值判别准则}
\src{138}
本节给出态射的分离性和固有性的赋值判别准则，也就是使用一个可变辅助概形 $\Spec(A)$ 的准则，其中 $A$ 为赋值环。在适当的“诺特性”假设下，可在这些准则中只考虑 $A$ 为离散赋值环的情形。这很可能是以后唯一需要应用的情形；在一般情形中引入任意赋值环，只是为了与经典论述建立联系。

\subsection{赋值环的回顾}
\src{138--141}

\egapar{7.1.1}
在刻画赋值环的各种等价性质中，下面使用这一种：称环 $A$ 为\emph{赋值环}，若它是不为域的整环，而且在所有包含于 $A$ 的分式域 $K$、又不等于 $K$ 的局部环中，$A$ 关于支配关系（I，8.1.1）为极大元。\transnote{本卷在此定义中明确排除域；译文保留这个约定，不把平凡赋值环自动加入定义。}
回顾，赋值环是整闭的。若 $A$ 为赋值环，则对 $A$ 的每个非零素理想 $\pp$，$A_\pp$ 也是赋值环。

\egapar{7.1.2}
设 $K$ 为域，$A$ 是 $K$ 的一个不为域的局部子环。则存在以 $K$ 为分式域且支配 $A$ 的赋值环（[1]，第1-07页，引理2）。
另一方面，设 $B$ 为赋值环，$k$ 为其剩余域，$K$ 为其分式域，$L$ 为 $k$ 的扩域。则存在支配 $B$、剩余域为 $L$ 的完备赋值环 $C$。事实上，$L$ 是 $k$ 的某个纯超越扩张 $L'=k(T_\lambda)_{\lambda\in\Lambda}$ 的代数扩张。已知可以将 $K$ 上对应于 $B$ 的赋值延拓到 $K'=K(T_\lambda)_{\lambda\in\Lambda}$ 上，使其剩余域为 $L'$（[24]，第98页）。将 $B$ 换成这个延拓赋值的赋值环的完备化，便可化归 $B$ 完备且 $L$ 是 $k$ 的代数闭包的情形。令 $\overline K$ 为 $K$ 的代数闭包；可以将定义 $B$ 的赋值延拓到 $\overline K$，相应的剩余域为 $k$ 的代数闭包，这只须把 $k[T]$ 中首一多项式的系数提升到 $\overline K$ 中即可看出。最后便化归 $L=k$ 的情形，取 $C$ 为 $B$ 的完备化即满足要求。

\egapar{7.1.3}
设 $K$ 为域，$A$ 为 $K$ 的子环。则 $A$ 在 $K$ 中的整闭包 $A'$ 是所有包含 $A$ 且以 $K$ 为分式域的赋值环的交（[11]，第51页，定理2）。（7.1.2）有以下等价的几何表述。

\egastatement{7.1.4}{命题}
设 $Y$ 为预概形，$p:X\to Y$ 为态射，$x\in X$，$y=p(x)$，$y'\ne y$ 为 $y$ 的一个特化（0，2.1.2）。则存在一个赋值环的谱 $Y'$，它为局部概形，以及一个分离态射 $f:Y'\to Y$，使若 $a$ 为 $Y'$ 的唯一闭点、$b$ 为其泛点，则 $f(a)=y'$、$f(b)=y$。还可要求下列附加性质之一成立：
\begin{enumerate}[label=\textup{(\roman*)}]
\item $Y'$ 为某个剩余域代数闭的完备赋值环的谱，且存在一个 $k(y)$-同态 $k(x)\to k(b)$。
\item 存在一个 $k(y)$-同构 $k(x)\xrightarrow{\sim}k(b)$。
\end{enumerate}
\begin{proof}
令 $Y_1$ 为 $Y$ 中以 $\overline{\{y\}}$ 为底空间的约化闭子预概形（I，5.2.1），令 $X_1=p^{-1}(Y_1)$ 为逆像闭子预概形。按假设 $y'\in\overline{\{y\}}$，而 $k(x)$ 在 $X$ 与 $X_1$ 中相同，故可设 $Y$ 为整，泛点为 $y$。$\OO_{y'}$ 此时为不为域的整局部环，其分式域为 $\OO_y=k(y)$，而 $k(x)$ 是 $k(y)$ 的扩域。为满足 $f(a)=y'$、$f(b)=y$ 和附加条件（i），取 $Y'=\Spec(A')$，其中 $A'$ 为支配 $\OO_{y'}$ 的完备赋值环，其剩余域为包含 $k(x)$ 的代数闭域；为满足（ii），取支配 $\OO_{y'}$ 且以 $k(x)$ 为分式域的赋值环 $A'$。$A'$ 的存在由（7.1.2）保证。\transnote{原文对（i）的构造写的是“剩余域为 $k(x)$ 的代数闭扩域”，而命题（i）要求的是 $k(x)$ 到分式域 $k(b)$ 的相容嵌入；这两种要求应予区分。这里保留原文叙述，不将剩余域与分式域混同。}
\end{proof}

\egapar{7.1.5}
回顾，称局部环 $A$ 的\emph{维数为 $1$}，若存在异于极大理想 $\mm$ 的素理想，且 $A$ 中每个异于 $\mm$ 的素理想都是极小素理想。若 $A$ 为整环，这等价于说 $\mm$ 与 $(0)$ 为仅有的两个素理想，且 $\mm\ne(0)$。换言之，$Y=\Spec(A)$ 只有两个点 $a,b$：$a$ 为唯一闭点，有 $\mathfrak j_a=\mm$，$k(a)=k=A/\mm$ 为剩余域；$b$ 为 $Y$ 的泛点，$\mathfrak j_b=(0)$，$\{b\}$ 为 $Y$ 中异于 $\varnothing,Y$ 的唯一开集，因而处处稠密；$k(b)=K$ 为 $A$ 的分式域。

\egapar{7.1.6}
对维数为 $1$ 的诺特局部环 $A$，已知以下条件等价（[1]，第2-08、17-01页）：（a）$A$ 正规；（b）$A$ 正则；（c）$A$ 为赋值环。此外，此时 $A$ 为离散赋值环。对离散赋值环，（7.1.2）和（7.1.3）有以下对应结果。

\egastatement{7.1.7}{命题}
设 $A$ 为不为域的诺特整局部环，$K$ 为其分式域，$L$ 为 $K$ 的有限型扩域。则存在以 $L$ 为分式域且支配 $A$ 的离散赋值环。
\begin{proof}
先设 $L=K$。令 $\mm$ 为 $A$ 的极大理想，$(x_1,\ldots,x_r)$ 为 $\mm$ 的一组非零生成元，令 $B=A[x_2/x_1,\ldots,x_r/x_1]\subset K$，它为诺特环。立即可知 $\mm B=x_1B$。若 $\pp$ 为 $x_1B$ 上的极小素理想，则已知 $\pp$ 的高度为 $1$（[13]，第一卷，第277页），也就是说 $B_\pp$ 是一维诺特局部环。\transnote{原文这里未说明分母生成元 $x_1$ 的选择。使用 $x_1B$ 上的极小素理想之前，还需要 $x_1B\ne B$；并非任意排列生成元都保证这一点。这里未将这项选择静默补作原文假设。}
显然，$\pp B_\pp\cap A$ 是包含 $\mm$ 而不含 $1$ 的 $A$ 的理想，所以它等于 $\mm$，从而 $B_\pp$ 支配 $A$（I，8.1.1）。由 Krull--Akizuki 定理（[25]，第293页），$B_\pp$ 的整闭包 $C$ 为诺特环，尽管 $C$ 不一定是有限型 $B_\pp$-模。若 $\nn$ 为 $C$ 的极大理想，则 $C_\nn$ 为一维正规诺特局部环（[25]，第295页），所以为离散赋值环；它支配 $B_\pp$，因而支配 $A$。

现设 $L$ 为 $K$ 的有限型扩域。由前面的情形，可以设 $A$ 已经是离散赋值环。令 $w$ 为 $K$ 上对应于 $A$ 的赋值；存在延拓 $w$ 的 $L$ 上的离散赋值 $w'$。事实上，按 $L$ 的生成元数目归纳，可化归 $L=K(a)$，此时结论是经典的（[24]，第106页）。
\end{proof}

\egastatement{7.1.8}{推论}
设 $A$ 为诺特整环，$K$ 为其分式域，$L$ 为 $K$ 的有限型扩域。则 $A$ 在 $L$ 中的整闭包，是所有以 $L$ 为分式域且包含 $A$ 的离散赋值环的交。
\begin{proof}
任意这样的离散赋值环正规，故包含 $L$ 中所有在 $A$ 上为整的元素。只须证明：若 $x\in L$ 不在 $A$ 上为整，则存在以 $L$ 为分式域、包含 $A$ 而不包含 $x$ 的离散赋值环 $C$。关于 $x$ 的假设等价于 $x\notin B=A[1/x]$，即 $1/x$ 在诺特环 $B$ 中不可逆。因此有包含 $1/x$ 的 $B$ 的素理想 $\pp$。整局部环 $B_\pp$ 诺特且包含于 $L$；$L$ 为其分式域的有限型扩张，因为该分式域包含 $K$。由（7.1.7），存在支配 $B_\pp$ 且以 $L$ 为分式域的离散赋值环 $C$。由于 $1/x\in\pp B_\pp$ 属于 $C$ 的极大理想，$x\notin C$。证明完毕。
\end{proof}
（7.1.7）的几何形式如下。

\egastatement{7.1.9}{命题}
设 $Y$ 为局部诺特预概形，$p:X\to Y$ 为局部有限型态射，$x\in X$，$y=p(x)$，$y'\ne y$ 为 $y$ 的特化。则存在一个离散赋值环的谱 $Y'$、一个分离态射 $f:Y'\to Y$，以及一个从 $Y'$ 到 $X$ 的 $Y$-有理映射 $g$，使若 $a$ 为 $Y'$ 的闭点、$b$ 为其泛点，则 $f(a)=y'$、$f(b)=y$、$g(b)=x$；而且在交换图
\begin{equation}\tag{7.1.9.1}
\begin{tikzcd}
&k(x)\arrow[dl,"\gamma"']\\
k(b)&k(y)\arrow[l,"\varphi"]\arrow[u,"\pi"']
\end{tikzcd}
\end{equation}
中，$\gamma$ 为双射；这里 $\pi,\varphi,\gamma$ 分别为对应于 $p,f,g$ 的同态。
\begin{proof}
如（7.1.4），结合（I，6.3.4，iv），可化归 $Y$ 整且泛点为 $y$。问题关于 $X,Y$ 局部，故可设 $p$ 为有限型。于是处于（7.1.4）的情形，并有额外性质：$k(x)$ 为 $k(y)$ 的有限型扩张（I，6.4.11），$\OO_{y'}$ 诺特。因此可以应用（7.1.7），取 $Y'=\Spec(A')$，其中 $A'$ 为支配 $\OO_{y'}$、以 $k(x)$ 为分式域的离散赋值环。这样得到交换图（7.1.9.1），其中 $\gamma$ 双射，$\pi,\varphi$ 分别对应于 $p,f$。又 $X,Y$ 局部诺特（I，6.6.2），$Y'$ 整，故由（I，7.1.15），存在唯一的从 $Y'$ 到 $X$ 的 $Y$-有理映射 $g$ 对应于同构 $\gamma$。证明完毕。
\end{proof}

\subsection{分离性的赋值判别准则}
\src{141--142}

\egastatement{7.2.1}{命题}
设 $X,Y$ 为预概形，$f:X\to Y$ 为拟紧态射。下列条件等价：
\begin{enumerate}[label=\alph*)]
\item $f$ 为闭态射。
\item 对每个 $x\in X$ 和 $y=f(x)$ 的每个异于 $y$ 的特化 $y'$，存在 $x$ 的一个特化 $x'$，使 $f(x')=y'$。
\end{enumerate}
\begin{proof}
条件（b）表示 $f(\overline{\{x\}})=\overline{\{y\}}$，因此是（a）的推论。\transnote{底本条件（b）的末尾印为 $f(x')=x$；根据所给点的所属空间及紧随其后的等式，这里明确校为 $f(x')=y'$。}
为证明（b）推出（a），取底空间 $X$ 的闭子集 $X'$，令 $Y'=\overline{f(X')}$，要证明 $Y'=f(X')$。分别考虑 $X,Y$ 中以 $X',Y'$ 为底空间的约化闭子预概形（I，5.2.1）。存在态射 $f':X'\to Y'$，使图
\[
\begin{tikzcd}
X'\arrow[r,"f'"]\arrow[d]&Y'\arrow[d]\\
X\arrow[r,"f"']&Y
\end{tikzcd}
\]
交换（I，5.2.2）；$f$ 拟紧，故 $f'$ 拟紧。因此只须证明：当 $f$ 拟紧且占优时，（b）推出 $f(X)=Y$。取 $y'\in Y$，令 $y$ 为包含 $y'$ 的一个不可约分支的泛点。由（b），只须证明 $f^{-1}(y)$ 非空；已知这由 $f$ 占优且拟紧推出（I，6.6.5）。
\end{proof}

\egastatement{7.2.2}{推论}
设 $f:X\to Y$ 为拟紧浸入。则底空间 $X$ 在 $Y$ 中闭，当且仅当它包含其每个点在 $Y$ 中的所有特化。

\egastatement{7.2.3}{命题}
设 $Y$ 为预概形（分别地，局部诺特预概形），$f:X\to Y$ 为态射（分别地，局部有限型态射）。下列条件等价：
\begin{enumerate}[label=\alph*)]
\item $f$ 分离。
\item 对角态射 $X\to X\times_Y X$ 拟紧；且对每个形如 $Y'=\Spec(A)$ 的 $Y$-预概形，其中 $A$ 为赋值环（分别地，离散赋值环），从 $Y'$ 到 $X$ 的两个 $Y$-态射若在 $Y'$ 的泛点处相同，则它们相等。
\item 对角态射 $X\to X\times_Y X$ 拟紧；且对每个形如 $Y'=\Spec(A)$ 的 $Y$-预概形，其中 $A$ 为赋值环（分别地，离散赋值环），$X'=X_{(Y')}$ 的两个 $Y'$-截面若在 $Y'$ 的泛点处相同，则它们相等。
\end{enumerate}
\begin{proof}
（b）与（c）的等价性，来自 $Y'$ 到 $X$ 的 $Y$-态射和 $X'$ 的 $Y'$-截面的一一对应（I，3.3.14）。若 $X$ 在 $Y$ 上分离，因为 $Y'$ 为整，由（I，7.2.2.1）知（b）成立。剩下只须证明（b）使对角态射 $\Delta:X\to X\times_Y X$ 闭，等价地，证明它满足（7.2.2）的准则。
取对角线 $\Delta(X)$ 上一点 $z$，及其在 $X\times_Y X$ 中异于 $z$ 的特化 $z'$。由（7.1.4），存在赋值环 $A$ 及态射 $h:Y'=\Spec(A)\to X\times_Y X$，将 $Y'$ 的闭点 $a$ 映至 $z'$，泛点 $b$ 映至 $z$。该态射使 $Y'$ 成为 $(X\times_Y X)$-预概形，因而成为 $Y$-预概形。将 $h$ 与 $X\times_Y X$ 的两个投影复合，得到两个 $Y$-态射 $g_1,g_2:Y'\to X$；它们在 $b$ 处相同，所以由假设等于同一个态射 $g$。这表示 $h=\Delta\circ g$（I，5.3.1），故 $z'\in\Delta(X)$。
若 $Y$ 局部诺特、$f$ 局部有限型，则 $X\times_Y X$ 局部诺特（I，6.6.7）；由（7.1.9），可重复上述论证，并取 $A$ 为离散赋值环。
\end{proof}

\egastatement{7.2.4}{注记}
（i）当 $Y$ 局部诺特、$f$ 局部有限型时，对角态射 $\Delta$ 拟紧的假设总成立，因为 $X\times_Y X$ 局部诺特（I，6.6.4，i）。在一般情形，这个假设也表示：对 $X$ 的任意仿射开覆盖 $(U_\alpha)$，各交 $U_\alpha\cap U_\beta$ 都拟紧。

（ii）为了使 $f$ 分离，只须对完备且剩余域代数闭的赋值环 $A$ 验证条件（b）或（c）；这由（7.2.3）的证明和（7.1.4）得出。

\subsection{固有性的赋值判别准则}
\src{143--147}

\egastatement{7.3.1}{命题}
设 $A$ 为赋值环，$Y=\Spec(A)$，$b$ 为 $Y$ 的泛点，$X$ 为整概形。设 $f:X\to Y$ 为闭态射，$f^{-1}(b)$ 只有一个点 $x$，且相应同态 $k(b)\to k(x)$ 双射。则 $f$ 为同构。
\begin{proof}
$f$ 闭且占优，故 $f(X)=Y$。只须证明（I，4.2.2）：对 $Y$ 中每个 $y\ne b$，存在唯一的点 $x'$ 使 $f(x')=y$，且相应同态 $\OO_y\to\OO_{x'}$ 双射；因为此时 $f$ 为同胚。
若 $f(x')=y$，则 $\OO_{x'}$ 是包含于 $K=k(x)=k(b)$ 且支配 $\OO_y$ 的局部环。后者为局部环 $A_{\mathfrak j_y}$，故为以 $K$ 为分式域的赋值环（7.1.1）。又 $x'$ 不是 $X$ 的泛点，所以 $\OO_{x'}\ne K$（0，2.1.3）。因此 $\OO_{x'}=\OO_y$。因为 $X$ 是整概形，关系 $\OO_{x'}=\OO_{x''}$ 推出 $x'=x''$（I，8.2.2），证明完毕。
\end{proof}

\egapar{7.3.2}
设 $A$ 为赋值环，$Y=\Spec(A)$，$b$ 为其泛点；$\OO_b=k(b)=K$ 为 $A$ 的分式域。设 $f:X\to Y$ 为态射。已知 $X$ 的 $Y$-有理截面，与定义在 $b$ 的邻域上的 $Y$-截面在 $b$ 处的芽一一对应（I，7.1.4），从而得到典范映射
\begin{equation}\tag{7.3.2.1}
\Gamma_{\mathrm{rat}}(X/Y)\longrightarrow
\Gamma(f^{-1}(b)/\Spec(K)).
\end{equation}
按定义（I，3.4.5），右端元素等同于 $f^{-1}(b)=X\otimes_A K$ 上的 $K$-有理点。若 $f$ 分离，因为 $Y$ 是整概形，由（I，5.4.7），映射（7.3.2.1）单射。将它与典范映射
$\Gamma(X/Y)\to\Gamma_{\mathrm{rat}}(X/Y)$（I，7.1.2）复合，得到典范映射
\begin{equation}\tag{7.3.2.2}
\Gamma(X/Y)\longrightarrow\Gamma(f^{-1}(b)/\Spec(K)).
\end{equation}
当 $f$ 分离时，这个映射仍单射（I，5.4.7）。

\egastatement{7.3.3}{命题}
设 $A$ 为分式域为 $K$ 的赋值环，$Y=\Spec(A)$，$b$ 为 $Y$ 的泛点，$f:X\to Y$ 为分离闭态射。则典范映射（7.3.2.2）双射；这等价于它为满射，并推出 $X$ 的各 $Y$-有理截面处处有定义。
\begin{proof}
取 $f^{-1}(b)$ 的一个 $K$-有理点 $x$。因为 $f$ 分离，相应态射 $f^{-1}(b)\to\Spec(K)$ 也分离（I，5.5.1，iv）；其每个截面都是闭浸入（I，5.4.6），故 $\{x\}$ 在 $f^{-1}(b)$ 中闭。考虑 $X$ 中以 $\overline{\{x\}}$ 为底空间的约化闭子预概形 $X'$。显然 $f|X'$ 满足（7.3.1）的条件，故为 $X'$ 到 $Y$ 的同构。其逆同构即为所求的 $X$ 的 $Y$-截面。
\end{proof}

\egapar{7.3.4}
为表述接下来的两个结果，采用一个将在第四章解释并讨论的术语：若 $F$ 是预概形 $Y$ 的子集，称整数 $\dim\OO_y$（$y$ 遍历 $F$）的下确界为 $F$ 在 $Y$ 中的\emph{余维数}，记为 $\codim_Y F$。

\egastatement{7.3.5}{推论}
设 $Y$ 为局部诺特约化预概形，$N$ 为 $Y$ 的非正则点的集合（0，4.1.4），并假设 $\codim_YN\geq2$。设 $f:X\to Y$ 为有限型、分离且闭的态射，$g$ 为 $X$ 的一个 $Y$-有理截面。令 $Y'$ 为 $g$ 没有定义的各点的集合，它为闭集（I，7.2.1），则 $\codim_YY'\geq2$。
\begin{proof}
只须证明：若 $y\in Y$ 且 $\dim\OO_y\leq1$，则 $g$ 在 $y$ 处有定义。若 $\dim\OO_y=0$，则 $y$ 为 $Y$ 的一个不可约分支的泛点（I，1.1.14），所以属于每个处处稠密的开集，特别属于 $g$ 的定义域。若 $\dim\OO_y=1$，由假设 $\OO_y$ 为正则诺特局部环，故为离散赋值环（7.1.6）。令 $Z=\Spec(\OO_y)$。因为 $U=Y-Y'$ 处处稠密，$U\cap Z$ 非空（I，2.4.2）；令 $g'$ 为 $g$ 诱导的从 $Z$ 到 $X$ 的有理映射（I，7.2.8）。只须证明 $g'$ 为态射（I，7.2.9）。可把 $g'$ 看作 $Z$-预概形 $f^{-1}(Z)=X\times_Y Z$ 的 $Z$-有理截面。与 $f$ 相应的态射 $f^{-1}(Z)\to Z$ 显然闭，并由（I，5.5.1，i）分离。由（7.3.3），$g'$ 处处有定义；$Z$ 约化且 $X$ 在 $Y$ 上分离，所以 $g'$ 为态射（I，7.2.2）。
\end{proof}

\egastatement{7.3.6}{推论}
设 $S$ 为局部诺特预概形，$X,Y$ 为两个 $S$-预概形。假设 $Y$ 约化，且 $Y$ 的非正则点集合 $N$ 满足 $\codim_YN\geq2$；再假设结构态射 $X\to S$ 固有。设 $f$ 为从 $Y$ 到 $X$ 的 $S$-有理映射，$Y'$ 为 $f$ 没有定义的各点的集合。则 $\codim_YY'\geq2$。
\begin{proof}
从 $Y$ 到 $X$ 的 $S$-有理映射等同于 $X\times_S Y$ 的 $Y$-有理截面（I，7.1.2）。结构态射 $X\times_S Y\to Y$ 闭（5.4.1），故应用（7.3.5）即得。
\end{proof}

\egastatement{7.3.7}{注记}
特别地，当 $Y$ 正规（0，4.1.4）时，由（7.1.6），关于 $Y$ 的上述假设成立。
可以利用（7.3.3）的逆命题，刻画普遍闭态射和固有态射。

\egastatement{7.3.8}{定理}
设 $Y$ 为预概形（分别地，局部诺特预概形），$f:X\to Y$ 为拟紧分离态射（分别地，有限型态射）。下列条件等价：
\begin{enumerate}[label=\alph*)]
\item $f$ 普遍闭（分别地，固有）。
\item 对每个形如 $Y'=\Spec(A)$ 的 $Y$-概形，其中 $A$ 为分式域为 $K$ 的赋值环（分别地，离散赋值环），与典范嵌入 $A\to K$ 对应的典范映射
\[
\Hom_Y(Y',X)\longrightarrow\Hom_Y(\Spec(K),X)
\]
为满射（分别地，双射）。
\item 对每个形如 $Y'=\Spec(A)$ 的 $Y$-概形，其中 $A$ 为赋值环（分别地，离散赋值环），与 $Y'$-预概形 $X_{(Y')}$ 有关的典范映射（7.3.2.2）为满射（分别地，双射）。
\end{enumerate}
\begin{proof}
（b）与（c）的等价性由（I，3.3.14）立即得到。（a）推出（c）：在两种假设下，（a）都说明 $f_{(Y')}$ 分离（I，5.5.1，iv）且闭，再应用（7.3.3）即可。
剩下证明（b）推出（a）。先考虑 $Y$ 任意、$f$ 分离且拟紧的情形。若 $f$ 满足（b），则对任意 $Y$-预概形 $Y''$，$f_{(Y'')}:X_{(Y'')}\to Y''$ 也满足（b）；这是因为（b）等价于（c），且对每个 $Y'\to Y''$，有
$X_{(Y'')}\times_{Y''}Y'=X\times_Y Y'$（I，3.3.9.1）。又 $f$ 分离且拟紧时，$f_{(Y'')}$ 也如此（I，5.5.1，iv；6.6.4，iii）。所以只须由（b）推出 $f$ 闭，即验证（7.2.1）的条件（b）。
取 $x\in X$，以及 $y=f(x)$ 的一个异于 $y$ 的特化 $y'$。由（7.1.4），存在赋值环的谱 $Y'$ 和分离态射 $g:Y'\to Y$，使其闭点 $a$、泛点 $b$ 满足 $g(a)=y'$、$g(b)=y$，且存在 $k(y)$-同态 $k(x)\to k(b)$。这个同态典范地对应于 $Y$-态射 $\Spec(k(b))\to X$（I，2.4.6）；由（b），它来自某个 $Y$-态射 $h:Y'\to X$。有 $h(b)=x$；令 $x'=h(a)$，则 $x'$ 为 $x$ 的特化，且
$f(x')=f(h(a))=g(a)=y'$。

现设 $Y$ 局部诺特，$f$ 为有限型。条件（b）首先由（7.2.3）推出 $f$ 分离，因为对角态射 $X\to X\times_Y X$ 拟紧（7.2.4）。此外，为验证 $f$ 固有，只须对每个有限型 $Y$-预概形 $Y''$ 证明 $f_{(Y'')}:X_{(Y'')}\to Y''$ 闭（5.6.3）。此时 $Y''$ 局部诺特，故可以重复第一种情形的论证，并取 $Y'$ 为离散赋值环的谱，以（7.1.9）代替（7.1.4）。
\end{proof}

\egastatement{7.3.9}{注记}
（i）当 $Y$ 为任意预概形、$f$ 为分离态射时，为使 $f$ 普遍闭，只须对完备且剩余域代数闭的赋值环 $A$ 验证（b）或（c）；这由证明和（7.1.4）得出。

（ii）从（7.3.8）的准则（c），可以得到“射影态射 $X\to Y$ 为闭态射”（5.5.3）的另一种证明，它更接近经典方法。可设 $Y$ 仿射，因而将 $X$ 等同于某个射影丛 $\Pb_Y^r$ 的闭子预概形（5.3.3）。为证 $X\to Y$ 闭，只须证明结构态射 $\Pb_Y^r\to Y$ 闭。结合（4.1.3.1），准则（7.3.8，c）将问题化为：若 $Y$ 为分式域为 $K$ 的赋值环 $A$ 的谱，则 $\Pb_Y^r$ 的每个 $K$-值点都由一个 $A$-值点限制到 $Y$ 的泛点而来。
每个可逆 $\OO_Y$-模都平凡（I，2.4.8），所以由（4.2.6），$\Pb_Y^r$ 的 $K$-值点等同于一类 $K$ 中的元素 $(\xi,\xi c_0,\ldots,\xi c_r)$，其中 $\xi\ne0$，各 $c_i$ 不全为零。将各 $c_i$ 乘以 $A$ 中具有适当赋值的一个元素，可设各 $c_i$ 都属于 $A$，且至少一个可逆。\transnote{底本写“$A$ 中的元素”；使最小赋值变为零的缩放因子一般应取自 $K^*$，不一定属于 $A$。例如所有坐标均有正赋值时，需要除以一个具有最小赋值的坐标。}
由（4.2.6），此时 $(c_0,\ldots,c_r)$ 也定义了 $\Pb_Y^r$ 的一个 $A$-值点，断言得证。

（iii）将给定 $Y$-预概形 $X$ 看作等价于给定关于 $Y$-预概形 $Y'$ 的函子
$X(Y')=\Hom_Y(Y',X)$ 时，（7.2.3）与（7.3.8）的准则特别方便。例如，利用它们将能在某些条件下证明“Picard 概形”固有。

\egastatement{7.3.10}{推论}
设 $Y$ 为整概形（分别地，局部诺特整概形），$X$ 为整概形，$f:X\to Y$ 为占优态射。
\begin{enumerate}[label=\textup{(\roman*)}]
\item 若 $f$ 拟紧且普遍闭，则每个以 $R(X)$ 为分式域、并支配 $Y$ 的某个局部环的赋值环，也支配 $X$ 的一个局部环。
\item 反之，设 $f$ 为有限型，且每个以 $R(X)$ 为分式域的赋值环（分别地，离散赋值环）都满足（i）中的性质。则 $f$ 固有。
\end{enumerate}
\begin{proof}
先注意，在所有情形下，假设都推出 $f$ 分离（I，5.5.9）。

（i）令 $K=R(Y)$、$L=R(X)$。取 $y\in Y$，设赋值环 $A$ 的分式域为 $L$，且支配 $\OO_y$。单射 $\OO_y\to A$ 定义态射 $h:Y'=\Spec(A)\to Y$（I，2.4.4），使 $h(a)=y$，其中 $a$ 为 $Y'$ 的闭点。又若 $\eta$ 为 $Y$ 的泛点，它也是 $\Spec(\OO_y)$ 的泛点，则 $h(b)=\eta$，其中 $b$ 为 $Y'$ 的泛点，因为 $K\subset L$。令 $\xi$ 为 $X$ 的泛点，有 $k(\xi)=k(b)=L$，从而得到 $Y$-态射 $g:\Spec(L)\to X$，使 $g(b)=\xi$。由（7.3.8），$g$ 来自某个 $Y$-态射 $g':Y'\to X$。令 $x=g'(a)$，则 $A$ 显然支配 $\OO_x$。

（ii）问题关于 $Y$ 局部，所以可设 $Y$ 仿射（分别地，仿射且诺特）。由于 $f$ 为有限型，两种情形都可以应用 Chow 引理（5.6.1）。因此有射影态射 $p:P\to Y$、浸入 $j:X'\to P$，以及射影、满射且双有理的态射 $g:X'\to X$，其中 $X'$ 为整，使下图交换：
\[
\begin{tikzcd}
P\arrow[d,"p"']&X'\arrow[l,"j"']\arrow[d,"g"]\\
Y&X.\arrow[l,"f"]
\end{tikzcd}
\]
只须证明 $j$ 为闭浸入；因为此时 $f\circ g=p\circ j$ 为射影态射，因而固有，且 $g$ 满射，所以 $f$ 固有（5.4.3）。令 $Z$ 为 $P$ 中以 $\overline{j(X')}$ 为底空间的约化闭子预概形（I，5.2.1）。由于 $X'$ 整，$j$ 分解为 $i\circ h$，其中 $i:Z\to P$ 为典范嵌入，$h:X'\to Z$ 为占优开浸入（I，5.2.3），且 $Z$ 整。此外，$Z$ 在 $Y$ 上射影。因此可设 $P$ 为整，$j$ 占优且双有理，只须证明 $j$ 满射。
取 $z\in P$。$\OO_z$ 是整局部环（分别地，诺特整局部环），其分式域为
$L=R(P)=R(X')=R(X)$。只须考虑 $z$ 不是 $P$ 的泛点的情形。由（7.1.2）、（7.1.7），存在以 $L$ 为分式域并支配 $\OO_z$ 的赋值环（分别地，离散赋值环）$A$。令 $y=p(z)$，则 $A$ 也支配 $\OO_y$。由假设，存在 $x\in X$ 使 $A$ 支配 $\OO_x$。$g$ 固有，故由证明的第一部分，$A$ 也支配 $X'$ 的某个局部环 $\OO_{x'}$。于是 $\OO_z$ 与 $\OO_{j(x')}=\OO_{x'}$ 为相关局部环（I，8.1.4）；而 $P$ 为概形，所以 $z=j(x')$（I，8.2.2）。证明完毕。
\end{proof}

\egastatement{7.3.11}{推论}
设 $X,Y$ 为整概形，$f:X\to Y$ 为占优、拟紧且普遍闭的态射。进一步设 $Y$ 仿射，环为整环 $B$。则 $\Gamma(X,\OO_X)$ 典范地同构于 $B$ 在 $R(X)$ 中的整闭包的一个子环。
\begin{proof}
由（I，8.2.1.1），$\Gamma(X,\OO_X)$ 等同于各 $\OO_x$（$x\in X$）的交。由（7.3.10）、（7.1.2）和（7.1.3），它包含于所有包含 $B$ 且以 $R(X)$ 为分式域的赋值环的交；再由（7.1.3）即得结论。
\end{proof}

\egastatement{7.3.12}{注记}
在（7.3.11）的假设下，若 $R(X)$ 为 $R(Y)$ 的有限型扩域，则在许多情形中可以推出：$\Gamma(X,\OO_X)$ 是 $B=\Gamma(Y,\OO_Y)$ 上的有限型模。例如，当 $B$ 是某个域上的有限型代数时就是如此，因为已知 $B$ 在其分式域的有限型扩张中的整闭包，为有限型 $B$-模（[13]，第一卷，第267页，定理9）；结合 $B$ 诺特与（7.3.11）即得。
特别地，在域 $K$ 上既固有又仿射的概形 $X$ 有限。由（1.6.4）、（5.4.6）及（I，6.4.4，c），可只考虑 $X$ 约化的情形。再只须对以 $X$ 的有限个不可约分支为底空间的各闭子预概形证明其在 $K$ 上有限，因而结合（5.4.5），最后化归 $X$ 为整的情形；此时结论由以上说明得出。
第三章将用其他方法重新得到这个命题，并将它置于更一般的结果之下：若 $f:X\to Y$ 固有，$Y$ 局部诺特，则对每个凝聚 $\OO_X$-模 $\F$，$f_*(\F)$ 凝聚（III，4.4.2）。
最后指出，在经典代数几何中，（7.3.10）的准则被用作固有态射的定义。这里只为这一原因提及它；就我们知道的应用而言，（7.3.8）的准则似乎更便于操作。

\subsection{代数曲线与一维函数域}
\src{148--152}
本小节的目的是说明如何用概形语言表述经典的代数曲线概念，例如 C. Chevalley 的著作[23]中所介绍的概念。\emph{本小节中，$k$ 表示一个域；所考虑的所有概形都是有限型 $k$-概形，所有态射都是 $k$-态射。}

\egastatement{7.4.1}{命题}
设 $X$ 为 $k$ 上的有限型预概形，因而诺特。令 $x_i$（$1\leq i\leq n$）为其不可约分支 $X_i$ 的泛点，$K_i=k(x_i)$。下列条件等价：
\begin{enumerate}[label=\alph*)]
\item 每个 $K_i$ 在 $k$ 上的超越次数为 $1$。
\item 对 $X$ 的每个闭点 $x$，局部环 $\OO_x$ 的维数为 $1$（7.1.5）。
\item $X$ 中异于各 $X_i$ 的不可约闭子集，恰为 $X$ 的闭点。
\end{enumerate}
\begin{proof}
因为 $X$ 拟紧，$X$ 的每个不可约闭子集 $F$ 包含一个闭点（0，2.1.3）。由（I，2.4.2），$\OO_x$ 的素理想与 $X$ 中包含 $x$ 的不可约闭子集一一对应（I，1.1.14），立即得到（b）、（c）的等价性。另一方面，若 $\pp_\alpha$（$1\leq\alpha\leq r$）为诺特局部环 $\OO_x$ 的极小素理想，则各 $\OO_x/\pp_\alpha$ 为整局部环，其分式域为那些满足 $x\in X_i$ 的 $K_i$。已知，有限型整局部 $k$-代数的维数等于其分式域在 $k$ 上的超越次数（[1]，第4-06页，定理2）。最后，$\dim\OO_x$ 为各 $\dim(\OO_x/\pp_\alpha)$ 的上确界。（a）说明这些维数都为 $1$，故推出（b）。反之，若 $\dim\OO_x=1$，则没有 $\pp_\alpha$ 能等于 $\OO_x$ 的极大理想，否则 $\OO_x$ 维数为零；因此每个 $\OO_x/\pp_\alpha$ 的维数都为 $1$，从而（b）推出（a）。
\end{proof}
由（I，6.4.4）立即可知，在（7.4.1）的条件下，集合 $X$ 或为空，或为无限集。

\egastatement{7.4.2}{定义}
称满足（7.4.1）条件的非空代数 $k$-概形为 $k$ 上的\emph{代数曲线}。
用将在第四章引入的维数语言来说，$k$ 上的代数曲线就是一个非空代数 $k$-概形，其每个不可约分支的维数都为 $1$。
应注意，若 $X$ 是 $k$ 上的代数曲线，则以其各不可约分支为底空间的约化闭子预概形 $X_i$（$1\leq i\leq n$）也都是 $k$ 上的代数曲线。

\egastatement{7.4.3}{推论}
设 $X$ 为不可约代数曲线。它的唯一非闭点是其泛点。$X$ 的真闭子集恰为由有限个闭点组成的集合；它们也是 $X$ 中仅有的非处处稠密的子集。
\begin{proof}
若 $x\in X$ 不是闭点，其闭包为 $X$ 的不可约闭子集，故由（7.4.1）必为 $X$ 本身，因此 $x$ 是泛点。真闭子集 $F\ne X$ 不含 $X$ 的泛点，所以其中每个点在 $X$ 中闭，因而在 $F$ 中也闭。考虑 $X$ 中以 $F$ 为底空间的约化闭子预概形（I，5.2.1），由（I，6.2.2）可知 $F$ 有限且离散。因此 $X$ 的每个无限子集的闭包必为 $X$。
\end{proof}
当 $X$ 为任意代数曲线时，对其各不可约分支应用（7.4.3），可知 $X$ 的非闭点恰为这些分支的泛点。

\egastatement{7.4.4}{推论}
设 $X,Y$ 为 $k$ 上的不可约代数曲线，$f:X\to Y$ 为 $k$-态射。则 $f$ 占优，当且仅当对每个 $y\in Y$，$f^{-1}(y)$ 有限。
\begin{proof}
若 $f$ 不占优，则由（7.4.3），$f(X)$ 必为 $Y$ 的有限子集，因而不可能每个 $f^{-1}(y)$ 都有限；否则 $X$ 有限，与（7.4.1）矛盾。反之，设 $f$ 占优。对每个异于 $Y$ 的泛点 $\eta$ 的点 $y$，由于 $\{y\}$ 闭，$f^{-1}(y)$ 在 $X$ 中闭（7.4.3）；它不含 $X$ 的泛点 $\xi$，故有限（7.4.3）。
最后，纤维 $f^{-1}(\eta)$ 为 $k(\eta)$ 上的有限型不可约概形，泛点为 $\xi$（I，6.3.9、6.4.11）。$k(\xi)$、$k(\eta)$ 都为 $k$ 的有限型扩域，超越次数同为 $1$，所以 $k(\xi)$ 为 $k(\eta)$ 的有限次扩张。因此 $\xi$ 在 $f^{-1}(\eta)$ 中闭（I，6.4.2），从而该纤维只有一点 $\xi$。
\end{proof}
第三章将证明：诺特预概形间的固有态射 $f:X\to Y$，若对每个 $y\in Y$，$f^{-1}(y)$ 有限，则 $f$ 必有限。因此，由（7.4.4）可知，不可约代数曲线到代数曲线的占优固有态射有限。

\egastatement{7.4.5}{推论}
设 $X$ 为 $k$ 上的代数曲线。则 $X$ 正则，当且仅当 $X$ 正规，也等价于其每个闭点的局部环都是离散赋值环。
\begin{proof}
由（7.4.1，b）和（7.1.6）立即得到。
\end{proof}

\egastatement{7.4.6}{推论}
设 $X$ 为约化代数曲线，$\A$ 为凝聚约化 $\R(X)$-代数。则 $X$ 相对于 $\A$ 的整闭包 $X'$（6.3.4）为正规代数曲线，且典范态射 $X'\to X$ 有限。
\begin{proof}
$X'\to X$ 有限来自（6.3.10），所以 $X'$ 为代数 $k$-概形。令 $x_i$（$1\leq i\leq n$）和 $x'_j$（$1\leq j\leq m$）分别为 $X,X'$ 各不可约分支的泛点；每个 $k(x'_j)$ 都是某个 $k(x_i)$ 的有限次代数扩张（6.3.6），故在 $k$ 上的超越次数为 $1$。因此 $X'$ 是代数曲线；而已知它是有限个整的正规概形的和（6.3.6、6.3.7）。
\end{proof}

\egapar{7.4.7}
称 $k$ 上的代数曲线 $X$ \emph{完备}，若它在 $k$ 上固有。

\egastatement{7.4.8}{推论}
约化代数曲线 $X$ 完备，当且仅当其正规化 $X'$ 完备。
\begin{proof}
典范态射 $f:X'\to X$ 有限（7.4.6），故固有（6.1.11）且满射（6.3.8）。令 $g:X\to\Spec(k)$ 为结构态射；由（5.4.2，ii）及（5.4.3，ii），并利用 $g$ 按假设分离，可知 $g$ 与 $g\circ f$ 同时固有。
\end{proof}

\egastatement{7.4.9}{命题}
设 $X$ 为 $k$ 上的正规代数曲线，$Y$ 为在 $k$ 上固有的代数 $k$-概形。则从 $X$ 到 $Y$ 的每个 $k$-有理映射都处处有定义，即为态射。
\begin{proof}
由（7.3.7），若这样的映射在 $x\in X$ 处没有定义，则 $\dim\OO_x\geq2$，所以这样的点不存在。最后的断言来自（I，7.2.3）。
\end{proof}

\egastatement{7.4.10}{推论}
$k$ 上的正规代数曲线在 $k$ 上拟射影。
\begin{proof}
$X$ 是有限个整的正规代数曲线的和（6.3.8），所以可设 $X$ 为整（5.3.6）。$X$ 拟紧，故有有限仿射开覆盖 $U_i$（$1\leq i\leq n$）。每个 $U_i$ 在 $k$ 上为有限型，故存在整数 $n_i$ 和 $k$-浸入 $f_i:U_i\to\Pb_k^{n_i}$（5.3.3；5.3.4，i）。$U_i$ 在 $X$ 中稠密，由（7.4.9），$f_i$ 延拓为 $k$-态射 $g_i:X\to\Pb_k^{n_i}$。于是得到 $k$-态射 $g=(g_1,\ldots,g_n)_k:X\to P$，其中 $P$ 为各 $\Pb_k^{n_i}$ 在 $k$ 上的乘积。对每个 $i$，$g_i|U_i$ 是浸入，故 $g|U_i$ 也是浸入（I，5.3.14）。各 $U_i$ 覆盖 $X$，而 $g$ 分离（I，5.5.13，v），所以 $g$ 是从 $X$ 到 $P$ 的浸入（I，8.2.8）。Segre 态射（4.3.3）给出 $P$ 到某个 $\Pb_k^r$ 的浸入，故 $X$ 拟射影。
\end{proof}

\egastatement{7.4.11}{推论}
正规代数曲线 $X$ 同构于某条正规完备代数曲线 $\overline X$ 在一个处处稠密开集上诱导的概形；$\overline X$ 在唯一同构意义下确定。
\begin{proof}
若 $X_1,X_2$ 为正规完备曲线，则由（7.4.9），$X_1$ 的一个稠密开集 $U_1$ 到 $X_2$ 的一个稠密开集 $U_2$ 的同构，唯一延拓为 $X_1$ 到 $X_2$ 的同构，故得到唯一性。
为证明存在性，将 $X$ 看作某个射影丛 $\Pb_k^r$ 的子概形（7.4.10）。令 $X_0$ 为 $X$ 在 $\Pb_k^r$ 中的闭包（I，9.5.11）。$X$ 是 $X_0$ 在一个稠密开集上诱导的概形（I，9.5.10），因此 $X$ 各不可约分支的泛点也是 $X_0$ 各不可约分支的泛点，且这些点的剩余域相同。由（7.4.1），$X_0$ 为代数曲线；它约化（I，9.5.9）且射影（5.5.1），因而完备（5.5.3）。取 $\overline X$ 为 $X_0$ 的正规化，它仍完备（7.4.8）。若 $h:\overline X\to X_0$ 为典范态射，则由于 $X$ 正规，$h|h^{-1}(X)$ 为到 $X$ 的同构（6.3.4）；$h^{-1}(X)$ 又包含 $\overline X$ 各不可约分支的泛点（6.3.8），故在 $\overline X$ 中稠密。证明完毕。
\end{proof}

\egastatement{7.4.12}{注记}
第五章将证明，不假设曲线正规，甚至不假设它约化，（7.4.10）的结论仍成立；还将证明，一条代数曲线，无论约化与否，仿射当且仅当它的各约化不可约分支都不完备。

\egastatement{7.4.13}{推论}
设 $X$ 为正规不可约曲线，$R(X)=K$；$Y$ 为整完备曲线，$R(Y)=L$。则占优 $k$-态射 $X\to Y$ 与 $k$-单同态 $L\to K$ 之间有典范一一对应。
\begin{proof}
由（7.4.9），从 $X$ 到 $Y$ 的 $k$-有理映射等同于 $k$-态射 $u:X\to Y$。分别令 $x,y$ 为 $X,Y$ 的泛点，则 $u$ 占优当且仅当 $u(x)=y$。由这些说明及（I，7.1.13）即得推论。
\end{proof}

\egapar{7.4.14}
当 $Y$ 取为射影直线 $\Pb_k^1=\Proj(k[T_0,T])$ 时，可以将（7.4.13）的结果说得更具体；其中 $T_0,T$ 是两个不定元。这是整概形（2.4.4），且其开集 $D_+(T_0)$ 上诱导的概形同构于 $\Spec(k[T])$（2.3.6）。所以 $Y$ 的泛点为 $k[T]$ 的理想 $(0)$，有理函数域为 $k(T)$，从而 $Y$ 为 $k$ 上的完备代数曲线。此外，$S=k[T_0,T]$ 中包含 $T_0$ 且异于 $S_+$ 的唯一齐次素理想为主理想 $(T_0)$。因此 $D_+(T_0)$ 在 $Y=\Pb_k^1$ 中的补集只有一个闭点，称为“无穷远点”，记为 $\infty$。向量丛与射影丛关系的一般讨论见8.4。采用这些记号，有：

\egastatement{7.4.15}{推论}
设 $X$ 为正规不可约曲线，$R(X)=K$。则从 $K$ 到所有从 $X$ 到 $\Pb_k^1$ 且异于取值为 $\infty$ 的常值态射的态射 $u$ 所成的集合，有一个典范双射。$u$ 占优，当且仅当其对应的 $K$ 中元素在 $k$ 上超越。

这个结论立即来自（7.4.9）以及下面的引理。

\egastatement{7.4.15.1}{引理}
设 $X$ 为 $k$ 上的整预概形，$K=R(X)$ 为其有理函数域。则 $K$ 与所有从 $X$ 到 $\Pb_k^1$ 且异于取值为 $\infty$ 的常值态射的有理映射 $u$ 所成的集合之间，有一个典范双射。这样的有理映射占优，当且仅当其对应的 $K$ 中元素在 $k$ 上超越。
\begin{proof}
从 $X$ 到 $\Pb_k^1$ 的有理映射，与 $\Pb_k^1$ 的 $K$-值点一一对应（I，7.1.12）。若这样的点定位于 $\Pb_k^1$ 的泛点（I，3.4.5），则对应的有理映射显然占优。否则，因为 $\Pb_k^1$ 中异于泛点的各点都闭（7.4.3），有理映射 $u$ 的类中唯一的态射 $U\to\Pb_k^1$（I，7.2.2）将定义域 $U$ 映到一个闭点 $y$。这个不一定在整个 $X$ 上有定义的态射因而不占优；稍滥用术语，称 $u$ 为“取值为 $y$ 的常值有理映射”。
剩下只须将定位点异于 $\infty$ 的 $\Pb_k^1$ 的 $K$-值点，与 $K$ 的元素一一对应，并验证：这种点的定位点为 $\Pb_k^1$ 的泛点，当且仅当对应元素在 $k$ 上超越。由（4.2.6，例1）立即可作此验证。
\end{proof}

\egastatement{7.4.16}{推论}
设 $X,Y$ 为 $k$ 上正规、完备且不可约的代数曲线，其有理函数域分别为 $K=R(X)$、$L=R(Y)$。则 $k$-同构 $X\xrightarrow{\sim}Y$ 的集合，与 $k$-同构 $L\xrightarrow{\sim}K$ 的集合之间有典范一一对应。
\begin{proof}
这是（7.4.13）的直接推论。
\end{proof}

\egapar{7.4.17}
推论（7.4.16）说明，$k$ 上正规、完备且不可约的代数曲线，在唯一同构意义下由其有理函数域 $K$ 确定。按定义，$K$ 是 $k$ 的超越次数为 $1$ 的有限型扩域；经典地称它为一个变量的代数函数域。此外，有：

\egastatement{7.4.18}{命题}
对 $k$ 的每个超越次数为 $1$ 的有限型扩域 $K$，存在正规、完备且不可约的代数曲线 $X$，使 $R(X)=K$；它在唯一同构意义下确定。$X$ 的局部环的集合（I，8.2.1），等同于由 $K$ 和所有包含 $k$ 且以 $K$ 为分式域的赋值环构成的集合。
\begin{proof}
$K$ 是 $k$ 的一个纯超越扩张 $k(T)$ 的有限次扩张；如前所见，$k(T)$ 等同于射影直线 $Y=\Pb_k^1$ 的有理函数域。令 $X$ 为 $Y$ 相对于 $K$ 的整闭包（6.3.4）。$X$ 为有理函数域为 $K$ 的正规代数曲线（6.3.7）；因为 $X\to Y$ 有限（7.4.6），它完备。$X$ 的局部环如下：$x$ 为泛点时，$\OO_x=K$；$x$ 异于泛点时，$\OO_x$ 为包含 $k$ 且以 $K$ 为分式域的离散赋值环（7.4.5）。
反之，设 $A$ 为这样的赋值环。$X\to\Spec(k)$ 固有，且 $A$ 支配 $k$，故它支配 $X$ 的一个局部环 $\OO_x$（7.3.10）。后者为以 $K$ 为分式域的赋值环，因而必等于 $A$。
\end{proof}

\egastatement{7.4.19}{注记}
由（7.4.16）和（7.4.18），给定 $k$ 上一条正规、完备且不可约的代数曲线，本质上等价于给定 $k$ 的一个超越次数为 $1$ 的有限型扩域 $K$。应注意，若 $k'$ 为基域 $k$ 的扩域，则 $X\otimes_k k'$ 仍为 $k'$ 上的完备代数曲线（5.4.2，iii），但一般既不约化，也不可约。然而，若 $K$ 在 $k$ 上可分，且 $k$ 在 $K$ 中相对代数闭，则它约化且不可约。经典术语把这种情形称为 $K$ 是 $k$ 的“正则扩张”，但本书不采用这一术语。即使如此，$X\otimes_k k'$ 仍可能不正规。细节见第四章。

\section[爆破概形、投影锥与射影闭包]{爆破概形；投影锥；射影闭包}
\subsection{爆破预概形}
\src{152--157}
\egapar{8.1.1}
设 $Y$ 为预概形，对每个整数 $n\geq0$，设 $\J_n$ 为 $\OO_Y$ 的拟凝聚理想。假设满足下列条件：
\begin{equation}
\J_0=\OO_Y,\qquad \J_{n+1}\subset\J_n\quad\text{对每个 }n;
\tag{8.1.1.1}
\end{equation}
\begin{equation}
\J_m\J_n\subset\J_{m+n}\quad\text{对任意 }m,n.
\tag{8.1.1.2}
\end{equation}
应注意，这些假设蕴含
\begin{equation}
\J_1^n\subset\J_n.
\tag{8.1.1.3}
\end{equation}
令
\begin{equation}
\SSS=\bigoplus_{n\geq0}\J_n.
\tag{8.1.1.4}
\end{equation}
由（8.1.1.1）和（8.1.1.2），$\SSS$ 为拟凝聚分次 $\OO_Y$-代数，因而定义了一个 $Y$-概形 $X=\Proj(\SSS)$。若 $\I$ 为 $\OO_Y$ 的可逆理想，则 $\J_n\otimes_{\OO_Y}\I^{\otimes n}$ 典范等同于 $\J_n\I^n$。因此，若把各 $\J_n$ 换成 $\J_n\I^n$，从而把 $\SSS$ 换成拟凝聚 $\OO_Y$-代数 $\SSS_{(\I)}$，则 $X_{(\I)}=\Proj(\SSS_{(\I)})$ 与 $X$ 典范同构（3.1.8）。

\egapar{8.1.2}
设 $Y$ 局部整，因而其有理函数层 $\R(Y)$ 为拟凝聚 $\OO_Y$-代数（I，7.3.7）。称 $\R(Y)$ 的一个子 $\OO_Y$-模 $\J$ 为 $\R(Y)$ 的\emph{分式理想}，若它为有限型（0，5.2.1）。假设对每个 $n\geq0$ 给定 $\R(Y)$ 的一个拟凝聚分式理想 $\J_n$，使 $\J_0=\OO_Y$，并满足条件（8.1.1.2），但不必满足（8.1.1.1）的第二个条件。仍可由公式（8.1.1.4）定义拟凝聚分次 $\OO_Y$-代数 $\SSS$，以及相应的 $Y$-概形 $X=\Proj(\SSS)$；对于 $\R(Y)$ 的每个可逆分式理想 $\I$，仍有 $X$ 到 $X_{(\I)}$ 的典范同构。

\egastatement{8.1.3}{定义}
设 $Y$ 为预概形（分别地，局部整预概形），$\J$ 为 $\OO_Y$ 的拟凝聚理想（分别地，$\R(Y)$ 的拟凝聚分式理想）。称 $Y$-概形
\[
X=\Proj\left(\bigoplus_{n\geq0}\J^n\right)
\]
是\emph{爆破理想 $\J$ 所得的概形}，或称为 $Y$ \emph{相对于 $\J$ 的爆破预概形}。当 $\J$ 为 $\OO_Y$ 的拟凝聚理想，$Y'$ 为它所定义的 $Y$ 的闭子预概形时，也称 $X$ 为\emph{爆破 $Y'$ 所得的 $Y$-概形}。

按定义，$\SSS=\bigoplus_{n\geq0}\J^n$ 由 $\SSS_1=\J$ 生成。若 $\J$ 为有限型 $\OO_Y$-模，则 $X$ 在 $Y$ 上射影（5.5.2）。不对 $\J$ 加上这个假设，$\OO_X$-模 $\OO_X(1)$ 仍可逆（3.2.5），且由（4.4.3），相对于结构态射 $X\to Y$ 甚丰沛。

应注意，若 $f:X\to Y$ 为结构态射，$\J$ 为 $\OO_Y$ 的理想，$Y'$ 为它所定义的闭子预概形，则 $f$ 在 $f^{-1}(Y-Y')$ 上的限制为到 $Y-Y'$ 的同构。事实上，这是 $Y$ 上的局部问题，所以只须考虑 $\J=\OO_Y$；此时结论来自（3.1.7）。

把 $\J$ 换成 $\J^d$（$d>0$）时，爆破 $Y$-概形 $X$ 换成一个与它典范同构的 $Y$-概形 $X'$（8.1.1）。同样，对每个可逆理想（分别地，可逆分式理想）$\I$，相对于 $\J\I$ 的爆破预概形 $X_{(\I)}$ 与 $X$ 典范同构（8.1.1）。特别地，若 $\J$ 是可逆理想或可逆分式理想，则爆破 $\J$ 所得的 $Y$-概形与 $Y$ 同构（3.1.7）。

\egastatement{8.1.4}{命题}
设 $Y$ 为整预概形。
\begin{enumerate}[label=(\roman*)]
\item 对 $\R(Y)$ 中满足（8.1.1.2）且 $\J_0=\OO_Y$ 的每个拟凝聚分式理想列 $(\J_n)$，$Y$-概形 $X=\Proj(\bigoplus_{n\geq0}\J_n)$ 为整，且结构态射 $f:X\to Y$ 占优。\transnote{原文（i）未排除所有正次分量均为零的情形。若 $\J_n=0$ 对所有 $n>0$ 成立，则 $X=\varnothing$，并非整概形。此处保留原断言；要按其后援引的（3.1.12）、（3.1.14）论证，须排除这一退化情形。（ii）已明确要求 $\J\ne0$。}
\item 设 $\J$ 为 $\R(Y)$ 的拟凝聚分式理想，$X$ 为 $Y$ 相对于 $\J$ 的爆破 $Y$-概形。若 $\J\ne0$，则结构态射 $f:X\to Y$ 双有理且满射。
\end{enumerate}
\begin{proof}
（i）来自 $\SSS=\bigoplus_{n\geq0}\J_n$ 为整的 $\OO_Y$-代数（3.1.12、3.1.14），因为对每个 $y\in Y$，$\OO_y$ 为整环（I，5.1.4）。

（ii）由（i），$X$ 为整。若 $x,y$ 分别为 $X,Y$ 的泛点，则 $f(x)=y$，只须证明 $k(x)$ 在 $k(y)$ 上的秩为 $1$。$x$ 也是纤维 $f^{-1}(y)$ 的泛点。令 $\psi:Z\to Y$ 为典范态射，其中 $Z=\Spec(k(y))$；则 $f^{-1}(y)$ 等同于 $\Proj(\SSS')$，其中 $\SSS'=\psi^*(\SSS)$（3.5.3）。显然 $\SSS'=\bigoplus_{n\geq0}(\J_y)^n$；由于 $\J$ 是 $\R(Y)$ 的非零拟凝聚分式理想，$\J_y\ne0$（I，7.3.6），故 $\J_y=k(y)$。于是 $\Proj(\SSS')$ 等同于 $\Spec(k(y))$（3.1.7），得到结论。
\end{proof}
将在（III，2.3.8）证明（8.1.4）的一个逆命题。

\egapar{8.1.5}
回到（8.1.1）的情形和记号。按定义，各单同态 $\J_{n+1}\to\J_n$（8.1.1.1）对每个 $k\in\Z$ 定义了分次 $\SSS$-模间的零次单同态
\begin{equation}
u_k:\SSS_+(k+1)\longrightarrow\SSS(k).
\tag{8.1.5.1}
\end{equation}
因为 $\SSS_+(k+1)$ 与 $\SSS(k+1)$ 典范地 $(\mathrm{TN})$-同构，$u_k$ 典范地对应于 $\OO_X$-模的单同态（3.4.2）
\begin{equation}
\wt u_k:\OO_X(k+1)\longrightarrow\OO_X(k).
\tag{8.1.5.2}
\end{equation}
另一方面，回忆（3.2.6）中定义的典范同态
\begin{equation}
\lambda:\OO_X(h)\otimes_{\OO_X}\OO_X(k)\longrightarrow\OO_X(h+k).
\tag{8.1.5.3}
\end{equation}
由于图表
\[
\begin{tikzcd}[column sep=large]
\SSS(h)\otimes_{\SSS}\SSS(k)\otimes_{\SSS}\SSS(l) \arrow[r] \arrow[d] &
\SSS(h+k)\otimes_{\SSS}\SSS(l) \arrow[d] \\
\SSS(h)\otimes_{\SSS}\SSS(k+l) \arrow[r] & \SSS(h+k+l)
\end{tikzcd}
\]
交换，由 $\lambda$ 的函子性（3.2.6），同态（8.1.5.3）在
\begin{equation}
\SSS_X=\bigoplus_{n\in\Z}\OO_X(n)
\tag{8.1.5.4}
\end{equation}
上定义了拟凝聚分次 $\OO_X$-代数结构。此外，图表
\[
\begin{tikzcd}[column sep=large]
\SSS(h)\otimes_{\SSS}\SSS_+(k+1) \arrow[r] \arrow[d,"1\otimes u_k"'] &
\SSS_+(h+k+1) \arrow[d,"u_{k+h}"] \\
\SSS(h)\otimes_{\SSS}\SSS(k) \arrow[r] & \SSS(h+k)
\end{tikzcd}
\]
交换，因而 $\lambda$ 的函子性给出交换图
\begin{equation}
\begin{tikzcd}[column sep=large]
\OO_X(h)\otimes_{\OO_X}\OO_X(k+1) \arrow[r,"\lambda"] \arrow[d,"1\otimes\wt u_k"'] &
\OO_X(h+k+1) \arrow[d,"\wt u_{k+h}"] \\
\OO_X(h)\otimes_{\OO_X}\OO_X(k) \arrow[r,"\lambda"] & \OO_X(h+k).
\end{tikzcd}
\tag{8.1.5.5}
\end{equation}
水平箭头为典范同态。因此，可以说各 $\wt u_k$ 定义了分次 $\SSS_X$-模间的零次单同态
\begin{equation}
\wt u:\SSS_X(1)\longrightarrow\SSS_X.
\tag{8.1.5.6}
\end{equation}

\egapar{8.1.6}
仍用（8.1.5）的记号。注意，对 $n\geq0$，复合同态
\[
\wt v_n=\wt u_{n-1}\circ\wt u_{n-2}\circ\cdots\circ\wt u_0
\]
是单同态 $\OO_X(n)\to\OO_X$。\transnote{上式按原版抄录；依（8.1.5.2）的源与靶，通常从右到左施行复合时，次序应为 $\wt u_0\circ\wt u_1\circ\cdots\circ\wt u_{n-1}$。$n=0$ 时取恒等同态。以下像层的构造按这条从 $\OO_X(n)$ 到 $\OO_X$ 的复合同态理解。}
以 $\J_{n,X}$ 表示它的像；这是 $\OO_X$ 的拟凝聚理想，与 $\OO_X(n)$ 同构。此外，对 $m\geq0,n\geq0$，图表
\[
\begin{tikzcd}[column sep=large]
\OO_X(m)\otimes_{\OO_X}\OO_X(n) \arrow[r,"\lambda"] \arrow[d,"\wt v_m\otimes\wt v_n"'] &
\OO_X(m+n) \arrow[d,"\wt v_{m+n}"] \\
\OO_X \arrow[r,"\id"'] & \OO_X
\end{tikzcd}
\]
交换。于是得到下列包含关系：
\begin{equation}
\J_{0,X}=\OO_X,\qquad \J_{n,X}\subset\J_{m,X}\quad(0\leq m\leq n),
\tag{8.1.6.1}
\end{equation}
\begin{equation}
\J_{m,X}\J_{n,X}\subset\J_{m+n,X}\quad(m\geq0,\ n\geq0).
\tag{8.1.6.2}
\end{equation}

\egastatement{8.1.7}{命题}
设 $Y$ 为预概形，$\J$ 为 $\OO_Y$ 的拟凝聚理想，$X=\Proj(\bigoplus_{n\geq0}\J^n)$ 为爆破 $\J$ 所得的 $Y$-概形。则对每个 $n>0$，有典范同构
\begin{equation}
\OO_X(n)\xrightarrow{\sim}\J^n\OO_X=\J_{n,X}
\tag{8.1.7.1}
\end{equation}
（参见0，4.3.5），因而 $n>0$ 时，$\J^n\OO_X$ 为可逆且甚丰沛的 $\OO_X$-模。
\begin{proof}
最后的断言是直接的，因为 $\OO_X(1)$ 可逆（3.2.5），并按定义相对于 $Y$ 甚丰沛（4.4.3、4.4.9）。另一方面，按定义，$v_n$ 的像就是 $\J^n\SSS$；因此（8.1.7.1）来自函子 $\wt\M$ 的正合性（3.2.4）及公式（3.2.4.1）。
\end{proof}

\egastatement{8.1.8}{推论}
在（8.1.7）的假设下，若 $f:X\to Y$ 为结构态射，$Y'$ 为 $Y$ 中由 $\J$ 定义的闭子预概形，则 $X$ 的闭子预概形 $X'=f^{-1}(Y')$ 由 $\J\OO_X$ 定义；这个理想典范同构于 $\OO_X(1)$，因而有典范正合列
\begin{equation}
0\longrightarrow\OO_X(1)\longrightarrow\OO_X\longrightarrow\OO_{X'}\longrightarrow0.
\tag{8.1.8.1}
\end{equation}
\begin{proof}
由（8.1.7.1）和（I，4.4.5）得到。
\end{proof}

\egapar{8.1.9}
在（8.1.7）的假设下，可以更精确地说明各 $\J_{n,X}$ 的结构。事实上，同态
\[
\wt u_{-1}:\OO_X\longrightarrow\OO_X(-1)
\]
典范地对应于 $\OO_X(-1)$ 在 $X$ 上的一个截面 $s$，称为相对于 $\J$ 的\emph{典范截面}（0，5.1.1）。另一方面，图（8.1.5.5）中的水平箭头是同构（3.2.7）；将其中 $h$ 换为 $k$，$k$ 换为 $-1$，得到 $\wt u_k=1_k\otimes\wt u_{-1}$，其中 $1_h$ 为 $\OO_X(h)$ 的恒等同态。\transnote{这里原版的下标也有错位：所述代换直接给出 $\wt u_{k-1}=1_k\otimes\wt u_{-1}$；等价地，$\wt u_k=1_{k+1}\otimes\wt u_{-1}$。其后“以 $s$ 作张量乘法”的解释不变。}
也就是说，对每个 $k\in\Z$，$\wt u_k$ 就是\emph{以典范截面 $s$ 作张量乘法}。同态 $\wt u$（8.1.5.6）也有相同解释。因此，对每个 $n\geq0$，同态 $\wt v_n:\OO_X(n)\to\OO_X$ 就是以 $s^{\otimes n}$ 作张量乘法。由此得到：

\egastatement{8.1.10}{推论}
采用（8.1.8）的记号，$X'$ 的底空间为满足 $s(x)=0$ 的点 $x\in X$ 所成的集合，其中 $s$ 为 $\OO_X(-1)$ 的典范截面。
\begin{proof}
若 $c_x$ 为 $(\OO_X(1))_x$ 的一个生成元，则 $s_x\otimes c_x$ 典范等同于 $\J_{1,X}$ 在 $x$ 处的茎的一个生成元；它可逆，当且仅当 $s_x\notin\mm_x(\OO_X(-1))_x$，即 $s(x)\ne0$。
\end{proof}

\egastatement{8.1.11}{命题}
设 $Y$ 为整预概形，$\J$ 为 $\R(Y)$ 的拟凝聚分式理想，$X$ 为爆破 $\J$ 所得的 $Y$-概形。则 $\J\OO_X$ 为相对于 $Y$ 甚丰沛的可逆 $\OO_X$-模。
\begin{proof}
这是 $Y$ 上的局部问题（4.4.5），所以可设 $Y=\Spec(A)$，其中 $A$ 为整环，分式域为 $K$，且 $\J=\wt{\frj}$，$\frj$ 为 $K$ 的分式理想。因此，存在 $A$ 的非零元素 $a$，使 $a\frj\subset A$。令 $S=\bigoplus_{n\geq0}\frj^n$。映射 $x\mapsto ax$ 给出 $\frj^{n+1}=(S(1))_n$ 到 $a\frj^{n+1}=a\frj S_n\subset\frj^n=S_n$ 的 $A$-同构，因而定义分次 $S$-模间的零次 $(\mathrm{TN})$-同构 $S_+(1)\to a\frj S$。另一方面，$x\mapsto a^{-1}x$ 为分次 $S$-模间的零次同构 $a\frj S\xrightarrow{\sim}\frj S$。于是复合后得到 $\OO_X$-模的同构 $\OO_X(1)\xrightarrow{\sim}\J\OO_X$（3.2.4）。由于 $S$ 由 $S_1=\frj$ 生成，$\OO_X(1)$ 可逆（3.2.5）且甚丰沛（4.4.3、4.4.9），得到结论。
\end{proof}

\subsection{分次环局部化的初步结果}
\src{157--162}
\egapar{8.2.1}
设 $S$ 为分次环，暂不假设其次数非负。令
\begin{equation}
S^{\geq}=\bigoplus_{n\geq0}S_n,\qquad S^{\leq}=\bigoplus_{n\leq0}S_n.
\tag{8.2.1.1}
\end{equation}
这是 $S$ 的两个分次子环，次数分别全非负、全非正。若 $f$ 为 $S$ 的次数为 $d$ 的齐次元素，$d$ 可正可负，则分式环 $S_f=S'$ 仍带有分次环结构：$S'_n$（$n\in\Z$）为各 $x/f^k$ 所成的集合，其中 $x\in S_{n+kd}$，$k\geq0$。仍记 $S_{(f)}=S'_0$，并分别以 $S_f^{\geq}$、$S_f^{\leq}$ 表示 $S'^{\geq}$、$S'^{\leq}$。若 $d>0$，则
\begin{equation}
(S^{\geq})_f=S_f.
\tag{8.2.1.2}
\end{equation}
事实上，若 $x\in S_{n+kd}$ 且 $n+kd<0$，则可写 $x/f^k=xf^h/f^{h+k}$，而当 $h>0$ 充分大时，$n+(h+k)d>0$。由定义得到
\begin{equation}
(S^{\geq})_{(f)}=(S_f^{\geq})_0=S_{(f)}.
\tag{8.2.1.3}
\end{equation}
若 $M$ 为分次 $S$-模，同样令
\begin{equation}
M^{\geq}=\bigoplus_{n\geq0}M_n,\qquad M^{\leq}=\bigoplus_{n\leq0}M_n.
\tag{8.2.1.4}
\end{equation}
它们分别为分次 $S^{\geq}$-模、分次 $S^{\leq}$-模，其交为 $S_0$-模 $M_0$。若 $f\in S_d$，仍将 $M_f$ 定义为分次 $S_f$-模，次数为 $n$ 的元素为各 $z/f^k$，其中 $z\in M_{n+kd}$、$k\geq0$。以 $M_{(f)}$ 表示 $M_f$ 中次数为零的元素所成的 $S_{(f)}$-模，分别简记 $(M_f)^{\geq}$、$(M_f)^{\leq}$ 为 $M_f^{\geq}$、$M_f^{\leq}$。若 $d>0$，同上可见
\begin{equation}
(M^{\geq})_f=M_f
\tag{8.2.1.5}
\end{equation}
以及
\begin{equation}
(M^{\geq})_{(f)}=(M_f^{\geq})_0=M_{(f)}.
\tag{8.2.1.6}
\end{equation}

\egapar{8.2.2}
设 $z$ 为一个不定元，称为\emph{齐次化变量}。若 $S$ 为分次环，次数可正可负，则多项式代数\origfootnote{这里不会与用 $\widehat S$ 表示一个环的分离完备化的记号混淆。}
\begin{equation}
\widehat S=S[z]
\tag{8.2.2.1}
\end{equation}
在对齐次 $f$ 和 $n\geq0$ 规定
\begin{equation}
\deg(fz^n)=n+\deg f
\tag{8.2.2.2}
\end{equation}
后，成为分次 $S$-代数。

\egastatement{8.2.3}{引理}
（i）有不分次环的典范同构
\begin{equation}
\widehat S_{(z)}\xrightarrow{\sim}\widehat S/(z-1)\widehat S\xrightarrow{\sim}S.
\tag{8.2.3.1}
\end{equation}
（ii）对每个 $f\in S_d$（$d>0$），有不分次环的典范同构
\begin{equation}
\widehat S_{(f)}\xrightarrow{\sim}S_f^{\leq}.
\tag{8.2.3.2}
\end{equation}
\begin{proof}
（8.2.3.1）中的第一个同构已在（2.2.5）中定义，第二个显然。因此，这样定义的同构 $\widehat S_{(z)}\xrightarrow{\sim}S$ 把 $xz^n/z^{n+k}$ 映到 $x$，其中 $\deg x=k$、$k\geq-n$。同态（8.2.3.2）将 $xz^n/f^k$，其中 $\deg x=kd-n$，映到 $S_f$ 中次数为 $-n$ 的元素 $x/f^k$；这显然也是同构。
\end{proof}

\egapar{8.2.4}
设 $M$ 为分次 $S$-模。显然，$S$-模
\begin{equation}
\widehat M=M\otimes_S\widehat S=M\otimes_S S[z]
\tag{8.2.4.1}
\end{equation}
为各 $S$-模 $M\otimes Sz^n$ 的直和，因而也是阿贝尔群 $M_k\otimes Sz^n$（$k\in\Z,n\geq0$）的直和。对 $M$ 的齐次元素 $x$，规定
\begin{equation}
\deg(x\otimes z^n)=n+\deg x,
\tag{8.2.4.2}
\end{equation}
即在 $\widehat M$ 上定义分次 $\widehat S$-模结构。下述（8.2.3）的类似结论留给读者证明。

\egastatement{8.2.5}{引理}
（i）有不分次模的典范双同构
\begin{equation}
\widehat M_{(z)}\xrightarrow{\sim}M.
\tag{8.2.5.1}
\end{equation}
（ii）对每个 $f\in S_d$（$d>0$），有不分次模的典范双同构
\begin{equation}
\widehat M_{(f)}\xrightarrow{\sim}M_f^{\leq}.
\tag{8.2.5.2}
\end{equation}

\egapar{8.2.6}
设 $S$ 为非负分次环，考虑其中的递减分次理想列
\begin{equation}
S_{[n]}=\bigoplus_{m\geq n}S_m\qquad(n\geq0).
\tag{8.2.6.1}
\end{equation}
特别地，$S_{[0]}=S$、$S_{[1]}=S_+$。显然 $S_{[m]}S_{[n]}\subset S_{[m+n]}$，所以可定义分次环 $S^{\natural}$：
\begin{equation}
S^{\natural}=\bigoplus_{n\geq0}S_n^{\natural},\qquad S_n^{\natural}=S_{[n]}.
\tag{8.2.6.2}
\end{equation}
因此 $S_0^{\natural}$ 为忘掉分次结构的环 $S$，$S^{\natural}$ 是一个 $S_0^{\natural}$-代数。对每个齐次元素 $f\in S_d$（$d>0$），以 $f^{\natural}$ 表示将 $f$ 看作 $S_{[d]}=S_d^{\natural}$ 的元素时所得的元素。采用这些记号，有：

\egastatement{8.2.7}{引理}
设 $S$ 为非负分次环，$f\in S_d$ 为齐次元素，$d>0$。有环的典范同构
\begin{equation}
S_f\xrightarrow{\sim}\bigoplus_{n\in\Z}S(n)_{(f)},
\tag{8.2.7.1}
\end{equation}
\begin{equation}
(S_f^{\geq})_{f/1}\xrightarrow{\sim}S_f,
\tag{8.2.7.2}
\end{equation}
\begin{equation}
S^{\natural}_{(f^{\natural})}\xrightarrow{\sim}S_f^{\geq},
\tag{8.2.7.3}
\end{equation}
其中前两个为分次环同构。
\begin{proof}
按定义，立即有 $(S_f)_n=(S(n)_f)_0$，从而得到（8.2.7.1）；它就是恒等映射。另一方面，$f/1$ 在 $S_f$ 中可逆，对 $S_f$ 应用（8.2.1.2），得到典范同构 $S_f\xrightarrow{\sim}(S_f^{\geq})_{f/1}=(S_f)_{f/1}$；其逆同构按定义就是（8.2.7.2）。最后，若 $x=\sum_{m\geq n}y_m\in S_{[n]}$，其中 $n=kd$，则把元素 $x/(f^{\natural})^k$ 映到 $S_f^{\geq}$ 中的 $\sum_m y_m/f^k$；立即可验证这定义了同构（8.2.7.3）。
\end{proof}

\egapar{8.2.8}
若 $M$ 为分次 $S$-模，同样对每个 $n\in\Z$ 令
\begin{equation}
M_{[n]}=\bigoplus_{m\geq n}M_m.
\tag{8.2.8.1}
\end{equation}
由于 $S_{[m]}M_{[n]}\subset M_{[m+n]}$（$m\geq0$），可定义分次 $S^{\natural}$-模 $M^{\natural}$：
\begin{equation}
M^{\natural}=\bigoplus_{n\in\Z}M_n^{\natural},\qquad M_n^{\natural}=M_{[n]}.
\tag{8.2.8.2}
\end{equation}
仍将下述引理的证明留给读者。

\egastatement{8.2.9}{引理}
采用（8.2.7）、（8.2.8）的记号，有模的典范双同构
\begin{equation}
M_f\xrightarrow{\sim}\bigoplus_{n\in\Z}M(n)_{(f)},
\tag{8.2.9.1}
\end{equation}
\begin{equation}
(M_f^{\geq})_{f/1}\xrightarrow{\sim}M_f,
\tag{8.2.9.2}
\end{equation}
\begin{equation}
M^{\natural}_{(f^{\natural})}\xrightarrow{\sim}M_f^{\geq},
\tag{8.2.9.3}
\end{equation}
其中前两个为分次模的双同构。

\egastatement{8.2.10}{引理}
设 $S$ 为非负分次环。
\begin{enumerate}[label=(\roman*)]
\item $S^{\natural}$ 为有限型（分别地，诺特）$S_0^{\natural}$-代数，当且仅当 $S$ 为有限型（分别地，诺特）$S_0$-代数。
\item 对 $n\geq n_0$ 有 $S_{n+1}^{\natural}=S_1^{\natural}S_n^{\natural}$，当且仅当对 $n\geq n_0$ 有 $S_{n+1}=S_1S_n$。
\item 对 $n\geq n_0$ 有 $S_n^{\natural}=(S_1^{\natural})^n$，当且仅当对 $n\geq n_0$ 有 $S_n=S_1^n$。
\item 若 $(f_\alpha)$ 为 $S_+$ 中的一族齐次元素，且 $S_+$ 是这些元素生成的 $S_+$-理想在 $S_+$ 中的根，则 $S_+^{\natural}$ 是各 $f_\alpha^{\natural}$ 生成的 $S_+^{\natural}$-理想在 $S_+^{\natural}$ 中的根。
\end{enumerate}
\begin{proof}
（i）若 $S^{\natural}$ 为有限型 $S_0^{\natural}$-代数，则由（2.1.6，i），$S_+=S_1^{\natural}$ 是有限型 $S=S_0^{\natural}$-模，故 $S$ 为有限型 $S_0$-代数（2.1.4）。若 $S^{\natural}$ 为诺特环，则 $S_0^{\natural}=S$ 也是诺特环（2.1.5）。反之，设 $S$ 为有限型 $S_0$-代数。由（2.1.6，ii），存在 $h>0$ 和 $m_0>0$，使 $n\geq m_0$ 时 $S_{n+h}=S_hS_n$；显然可设 $m_0\geq h$。此外，各 $S_m$ 为有限型 $S_0$-模（2.1.6，i）。于是，若 $n\geq m_0+h$，则
\[
S_n^{\natural}=S_hS_{n-h}^{\natural}=S_h^{\natural}S_{n-h}^{\natural}.
\]
若 $m<m_0+h$，令 $E=S_{m_0}+\cdots+S_{m_0+h-1}$，则
\[
S_m^{\natural}=S_m+\cdots+S_{m_0+h-1}+S_hE+S_h^2E+\cdots.
\]
对 $1\leq m\leq m_0$，令 $G_m$ 为各 $S_i$（$m\leq i\leq m_0+h-1$）作为 $S_0$-模的有限生成系的并，并将它看作 $S_{[m]}$ 的子集。对 $m_0+1\leq m\leq m_0+h-1$，同样令 $G_m$ 为各 $S_i$（$m\leq i\leq m_0+h-1$）及 $S_hE$ 的有限 $S_0$-生成系的并，看作 $S_{[m]}$ 的子集。显然，对 $1\leq m\leq m_0+h-1$，有 $S_m^{\natural}=S_0^{\natural}G_m$。因此，这些 $G_m$ 的并 $G$ 为 $S_0^{\natural}$-代数 $S^{\natural}$ 的一个生成系。由此还得到：若 $S=S_0^{\natural}$ 为诺特环，则 $S^{\natural}$ 也为诺特环。

（ii）若 $n\geq n_0$ 时 $S_{n+1}=S_1S_n$，则 $S_{n+1}^{\natural}=S_1S_n^{\natural}$，从而 $S_{n+1}^{\natural}=S_1^{\natural}S_n^{\natural}$。反之，后一个关系写成
\[
S_{n+1}+S_{n+2}+\cdots=(S_1+S_2+\cdots)(S_n+S_{n+1}+\cdots).
\]
比较两边在 $S$ 中次数为 $n+1$ 的分量，得到 $S_{n+1}=S_1S_n$。

（iii）若 $n\geq n_0$ 时 $S_n=S_1^n$，则 $S_n^{\natural}=S_1^n+S_1^{n+1}+\cdots$。$S_1^{\natural}$ 包含 $S_1+S_1^2+\cdots$，故 $S_n^{\natural}\subset(S_1^{\natural})^n$，因而 $n\geq n_0$ 时 $S_n^{\natural}=(S_1^{\natural})^n$。反之，$(S_1^{\natural})^n=(S_1+S_2+\cdots)^n$ 中，在 $S$ 中次数为 $n$ 的项恰为 $S_1^n$ 中的项；所以 $S_n^{\natural}=(S_1^{\natural})^n$ 蕴含 $S_n=S_1^n$。

（iv）只须证明：把 $g\in S_{k+h}$ 看作 $S_k^{\natural}$ 的元素，其中 $k>0,h\geq0$，则存在整数 $n>0$，使 $S_{kn}^{\natural}$ 中的 $g^n$ 为各 $f_\alpha^{\natural}$ 的线性组合，系数属于 $S_+^{\natural}$。按假设，存在 $m_0$，使对 $m\geq m_0$，在 $S$ 中有
\[
g^m=\sum_\alpha c_{\alpha m}f_\alpha,
\]
其中出现的指标 $\alpha$ 与 $m$ 无关。此外，显然可取 $c_{\alpha m}$ 为齐次的，且在 $S$ 中
\[
\deg(c_{\alpha m})=m(k+h)-\deg f_\alpha.
\]
将 $m_0$ 取到充分大，使对 $g^{m_0}$ 中出现的所有 $f_\alpha$，都有 $km_0>\deg f_\alpha$。对每个 $\alpha$，令 $c'_{\alpha m}$ 为将 $c_{\alpha m}$ 看作 $S^{\natural}$ 中次数为 $km-\deg f_\alpha$ 的元素所得的元素。于是，在 $S^{\natural}$ 中有 $g^m=\sum_\alpha c'_{\alpha m}f_\alpha^{\natural}$，证明完毕。
\end{proof}

\egapar{8.2.11}
考虑分次 $S_0$-代数
\begin{equation}
S^{\natural}\otimes_S S_0=S^{\natural}/S_+S^{\natural}
=\bigoplus_{n\geq0}S_{[n]}/S_+S_{[n]}.
\tag{8.2.11.1}
\end{equation}
由于 $S_n$ 是 $S_{[n]}/S_+S_{[n]}$ 的商 $S_0$-模，有分次 $S_0$-代数的典范同态
\begin{equation}
S^{\natural}\otimes_S S_0\longrightarrow S.
\tag{8.2.11.2}
\end{equation}
它显然满射，因而对应于典范闭浸入（2.9.2）
\begin{equation}
\Proj(S)\longrightarrow\Proj(S^{\natural}\otimes_S S_0).
\tag{8.2.11.3}
\end{equation}

\egastatement{8.2.12}{命题}
典范态射（8.2.11.3）为双射。同态（8.2.11.2）为 $(\mathrm{TN})$-双射，当且仅当存在 $n_0$，使 $n\geq n_0$ 时 $S_{n+1}=S_1S_n$。若满足后一条件，则（8.2.11.3）为同构；当 $S$ 诺特时，逆命题成立。
\begin{proof}
为证明第一个断言，只须证明同态（8.2.11.2）的核 $\frj$ 由幂零元素组成（2.8.3）。若 $f\in S_{[n]}$ 模 $S_+S_{[n]}$ 的类属于此核，则 $f\in S_{[n+1]}$。把 $f^{n+1}$ 看作 $S_{[n(n+1)]}$ 的元素，它属于 $S_+S_{[n(n+1)]}$，因为可写成 $f\cdot f^n$。所以它模 $S_+S_{[n(n+1)]}$ 的类为零，证得所述断言。

条件 $n\geq n_0$ 时 $S_{n+1}=S_1S_n$，等价于 $n\geq n_0$ 时 $S_{n+1}^{\natural}=S_1^{\natural}S_n^{\natural}$（8.2.10，ii）。按定义，这等价于（8.2.11.2）为 $(\mathrm{TN})$-单射，因而为 $(\mathrm{TN})$-双射；此时由（2.9.1），（8.2.11.3）为同构。

反之，若（8.2.11.3）为同构，则 $\Proj(S^{\natural}\otimes_S S_0)$ 上的层 $\wt{\frj}$ 为零（2.9.2，i）。在 $S$ 诺特的假设下，$S^{\natural}\otimes_S S_0$ 是诺特环 $S^{\natural}$ 的商，故诺特（8.2.10，i）。由（2.7.3），$\frj$ 满足 $(\mathrm{TN})$，所以对充分大的 $n$，$S_{n+1}^{\natural}=S_1^{\natural}S_n^{\natural}$。再用（8.2.10，ii）即完成证明。
\end{proof}

\egapar{8.2.13}
现考虑典范单同态 $(S_+)^n\to S_{[n]}$；它们定义了分次环间的零次单同态
\begin{equation}
\bigoplus_{n\geq0}(S_+)^n\longrightarrow S^{\natural}.
\tag{8.2.13.1}
\end{equation}

\egastatement{8.2.14}{命题}
同态（8.2.13.1）为 $(\mathrm{TN})$-同构，当且仅当存在 $n_0$，使对每个 $n\geq n_0$，有 $S_n=S_1^n$。此时，与（8.2.13.1）对应的态射处处有定义，且为同构
\[
\Proj(S^{\natural})\longrightarrow\Proj\left(\bigoplus_{n\geq0}(S_+)^n\right).
\]
当 $S$ 诺特时，逆命题成立。
\begin{proof}
由（8.2.10，iii）及（2.9.1），前两个断言显然。第三个断言由（8.2.10，i、iii）以及下述引理得到。
\end{proof}

\egastatement{8.2.14.1}{引理}
设 $T$ 为非负分次环，且为有限型 $T_0$-代数。若与单同态 $\bigoplus_{n\geq0}T_1^n\to T$ 对应的态射处处有定义，且为同构
\[
\Proj(T)\longrightarrow\Proj\left(\bigoplus_{n\geq0}T_1^n\right),
\]
则存在 $n_0$，使 $n\geq n_0$ 时 $T_n=T_1^n$。
\begin{proof}
令 $g_i$（$1\leq i\leq r$）为 $T_0$-模 $T_1$ 的生成元。假设首先蕴含，各 $D_+(g_i)$ 覆盖 $\Proj(T)$（2.8.1）。令 $h_j$（$1\leq j\leq s$）为 $T_+$ 中的一组齐次元素，$\deg h_j=n_j$，它们与各 $g_i$ 一起生成理想 $T_+$，等价地，生成 $T_0$-代数 $T$（2.1.3）。令 $T'=\bigoplus_{n\geq0}T_1^n$。环 $T_{(g_i)}$ 中的元素 $h_j/g_i^{n_j}$ 按假设属于子环 $T'_{(g_i)}$，所以存在整数 $k$，使对每个 $j$，有 $T_1^kh_j\subset T_1^{k+n_j}$。由归纳，对每个 $r\geq1$，$T_1^kh_j^r\subset T'$；再由各 $h_j$ 的定义，得到 $T_1^kT\subset T'$。

另一方面，对每个 $j$，存在整数 $m_j$，使 $h_j^{m_j}$ 属于各 $g_i$ 在 $T$ 中生成的理想（2.3.14），所以 $h_j^{m_j}\in T_1T$，进而 $h_j^{m_jk}\in T_1^kT\subset T'$。因此，存在整数 $m_0\geq k$，使 $m\geq m_0$ 时 $h_j^m\in T_1^{mn_j}$。若 $q$ 为各 $n_j$ 中的最大值，则 $n_0=qsm_0+k$ 满足要求。事实上，$n\geq n_0$ 时，$T_n$ 的元素为属于 $T_1^\alpha u$ 的单项式的和，其中 $u$ 为各 $h_j$ 的幂的乘积。\transnote{这一句原版写作“$S_n$ 的元素”；本引理只涉及 $T$，此处按上下文写作 $T_n$。}
若 $\alpha\geq k$，由前述结果有 $T_1^\alpha u\subset T_1^n$。否则，某个 $h_j$ 的指数至少为 $m_0$，所以 $u\in T_1^\beta v$，其中 $\beta\geq k$，$v$ 仍为各 $h_j$ 的幂的乘积，归结为前一种情形。故在所有情形中都有 $T_1^\alpha u\subset T_1^n$。
\end{proof}

\egastatement{8.2.15}{注记}
条件 $n\geq n_0$ 时 $S_n=S_1^n$ 显然蕴含 $n\geq n_0$ 时 $S_{n+1}=S_1S_n$，但逆命题不成立，即使假设 $S$ 诺特也如此。例如，令 $K$ 为域，$A=K[x]$，$B=K[y]/y^2K[y]$，其中 $x,y$ 为不定元，次数分别为 $1$、$2$，并令 $S=A\otimes_K B$。则 $S$ 为分次 $K$-代数，有一组基由 $1$、$x^n$（$n\geq1$）及 $x^ny$（$n\geq0$）组成。立即可见，$n\geq2$ 时 $S_{n+1}=S_1S_n$，但是 $S_1^n=Kx^n$，而 $n\geq2$ 时 $S_n=Kx^n+Kx^{n-2}y$。

\subsection{投影锥}
\src{162--168}
\egapar{8.3.1}
设 $Y$ 为预概形；本小节只讨论 $Y$-预概形及 $Y$-态射。设 $\SSS$ 为非负分次的拟凝聚 $\OO_Y$-代数，\emph{并假设 $\SSS_0=\OO_Y$}。依（8.2.2）的记号，令
\begin{equation}
\widehat{\SSS}=\SSS[z]=\SSS\otimes_{\OO_Y}\OO_Y[z],
\tag{8.3.1.1}
\end{equation}
按公式（8.2.2.2）定义次数，将它看作非负分次 $\OO_Y$-代数。这样，对 $Y$ 的每个仿射开集 $U$，有
\[
\Gamma(U,\widehat{\SSS})=(\Gamma(U,\SSS))[z].
\]
以下令
\begin{equation}
X=\Proj(\SSS),\qquad C=\Spec(\SSS),\qquad \widehat C=\Proj(\widehat{\SSS}),
\tag{8.3.1.2}
\end{equation}
其中定义 $C$ 时，把 $\SSS$ 看作不分次的 $\OO_Y$-代数。称 $C$、$\widehat C$ 分别为 $\SSS$ 定义的\emph{仿射锥}、\emph{射影锥}；也用“锥”简称“仿射锥”。稍滥用术语，也称 $C$、$\widehat C$ 分别为 $X$ 的\emph{仿射投影锥}、\emph{射影投影锥}，其中默认预概形 $X$ 是以 $\Proj(\SSS)$ 的形式给定的。最后，称 $\widehat C$ 为 $C$ 的\emph{射影闭包}，这里 $\SSS$ 的给定包含在 $C$ 的结构之中。

\egastatement{8.3.2}{命题}
存在典范 $Y$-态射
\begin{equation}
Y\xrightarrow{\varepsilon}C\xrightarrow{i}\widehat C,
\tag{8.3.2.1}
\end{equation}
\begin{equation}
X\xrightarrow{j}\widehat C,
\tag{8.3.2.2}
\end{equation}
使 $\varepsilon,j$ 为闭浸入，$i$ 为仿射态射，且是占优开浸入，满足
\begin{equation}
i(C)=\widehat C-j(X).
\tag{8.3.2.3}
\end{equation}
此外，$\widehat C$ 是 $\widehat C$ 中包含 $i(C)$ 的最小闭子预概形。
\begin{proof}
为定义 $i$，考虑 $\widehat C$ 的开集
\begin{equation}
\widehat C_z=\Spec\bigl(\widehat{\SSS}/(z-1)\widehat{\SSS}\bigr)
\tag{8.3.2.4}
\end{equation}
（3.1.4），其中 $z$ 典范等同于 $\widehat{\SSS}$ 在 $Y$ 上的一个截面。同构 $i:C\xrightarrow{\sim}\widehat C_z$ 对应于典范同构（8.2.3.1）
\[
\widehat{\SSS}/(z-1)\widehat{\SSS}\xrightarrow{\sim}\SSS.
\]
态射 $\varepsilon$ 对应于核为 $\SSS_+$ 的增广同态 $\SSS\to\SSS_0=\OO_Y$（1.2.7）；因为此同态满射，$\varepsilon$ 为闭浸入（1.4.10）。最后，$j$ 同样对应于零次满同态 $\widehat{\SSS}\to\SSS$（3.5.1）；它在 $\SSS$ 上为恒等，在其核 $z\widehat{\SSS}$ 上为零。由（3.6.2），$j$ 处处有定义且为闭浸入。

为证明其余断言，显然可设 $Y=\Spec(A)$ 为仿射，$\SSS=\wt S$，其中 $S$ 为分次 $A$-代数，从而 $\widehat{\SSS}=(\widehat S)^{\sim}$。此时，$S_+$ 的齐次元素 $f$ 等同于 $\widehat{\SSS}$ 在 $Y$ 上的截面；（2.3.3）中记作 $D_+(f)$ 的 $\widehat C$ 的开集也写作 $\widehat C_f$（3.1.4）。同样，（I，1.1.1）中记作 $D(f)$ 的 $C$ 的开集也写作 $C_f$（0，5.5.2）。由（2.3.14）及 $\widehat S$ 的定义，开集 $\widehat C_z=i(C)$ 与各 $\widehat C_f$，其中 $f$ 遍历 $S_+$ 中的齐次元素，覆盖 $\widehat C$。此外，有
\begin{equation}
i^{-1}(\widehat C_f)=C_f.
\tag{8.3.2.5}
\end{equation}
事实上，
\[
\widehat C_f\cap i(C)=\widehat C_f\cap\widehat C_z
=\widehat C_{fz}=\Spec(\widehat S_{(fz)}).
\]
若 $d=\deg f$，则 $\widehat S_{(fz)}$ 典范同构于 $(\widehat S_{(z)})_{f/z^d}$（2.2.2）。由同构（8.2.3.1）的定义，在所对应的分式环同构下，$(\widehat S_{(z)})_{f/z^d}$ 的像恰为 $S_f$。因为 $C_f=\Spec(S_f)$，这证明了（8.3.2.5），同时表明 $i$ 为仿射态射。此外，将 $i|C_f$ 看作到 $\widehat C_f$ 的态射，它对应于典范同态 $\widehat S_{(f)}\to\widehat S_{(fz)}$（I，1.7.3）。由此及（8.2.3.2），可得到下述结果。

\egapar{8.3.2.6}
\emph{若 $Y=\Spec(A)$ 为仿射，$\SSS=\wt S$，则对 $S_+$ 的每个齐次元素 $f$，$\widehat C_f$ 典范等同于 $\Spec(S_f^{\leq})$，而 $i$ 的限制 $C_f\to\widehat C_f$ 对应于典范单同态 $S_f^{\leq}\to S_f$。}

仍设 $Y$ 仿射。按定义，$\widehat C_z$ 在 $\widehat C=\Proj(\widehat S)$ 中的补集为 $\widehat S$ 中包含 $z$ 的齐次素理想所成的集合；由 $j$ 的定义，这就是 $j(X)$，所以（8.3.2.3）成立。

为证明最后的断言，仍可设 $Y$ 仿射。采用上述记号，注意 $z$ 在 $\widehat S$ 中不是零因子，而 $i(C)=\widehat C_z$；只须证明下面的引理。
\end{proof}

\egastatement{8.3.2.7}{引理}
设 $T$ 为非负分次环，$Z=\Proj(T)$，$g$ 为 $T$ 中次数为 $d>0$ 的齐次元素。若 $g$ 在 $T$ 中不是零因子，则 $Z$ 是 $Z$ 中包含 $Z_g=D_+(g)$ 的最小闭子预概形。
\begin{proof}
由（I，4.1.9），这是 $Z$ 上的局部问题。因此，对每个齐次元素 $h\in T_e$（$e>0$），只须证明 $Z_h$ 是 $Z_h$ 中包含 $Z_{gh}$ 的最小闭子预概形。由定义和（I，4.3.2），这等价于典范同态 $T_{(h)}\to T_{(gh)}$ 单射。这个同态等同于典范同态 $T_{(h)}\to(T_{(h)})_{g^e/h^d}$（2.2.3）。由于 $g^e$ 在 $T$ 中不是零因子，$g^e/h^d$ 在 $T_h$ 中也不是零因子，因而在 $T_{(h)}$ 中亦然。事实上，若 $t\in T,m>0$ 且 $(g^e/h^d)(t/h^m)=0$，则存在 $n>0$ 使 $h^ng^et=0$，从而 $h^nt=0$，故 $t/h^m=0$。由（0，1.2.2）完成证明。
\end{proof}

\egapar{8.3.3}
常通过开浸入 $i$，把仿射锥 $C$ 与射影锥 $\widehat C$ 在开集 $i(C)$ 上诱导的子预概形等同。与闭浸入 $\varepsilon$ 对应的 $C$ 的闭子预概形，称为 $C$ 的\emph{顶点预概形}。$\varepsilon$ 是 $C$ 的一个 $Y$-截面，也称为 $C$ 的\emph{顶点截面}或\emph{零截面}。通过 $\varepsilon$ 可把 $Y$ 与 $C$ 的顶点预概形等同。

此外，$i\circ\varepsilon$ 为 $\widehat C$ 的一个 $Y$-截面，因而也是闭浸入（I，5.4.6）；它对应于零次典范满同态 $\widehat{\SSS}=\SSS[z]\to\OO_Y[z]$（参见3.1.7），其核为 $\SSS_+[z]=\SSS_+\widehat{\SSS}$。这个闭浸入对应的 $\widehat C$ 的子预概形仍称为 $\widehat C$ 的\emph{顶点预概形}，$i\circ\varepsilon$ 称为 $\widehat C$ 的\emph{顶点截面}；通过 $i\circ\varepsilon$ 可将它与 $Y$ 等同。最后，与 $j$ 对应的 $\widehat C$ 的闭子预概形称为 $\widehat C$ 的\emph{无穷远子预概形}，可通过 $j$ 与 $X$ 等同。

\egapar{8.3.4}
分别在开集
\begin{equation}
E=C-\varepsilon(Y),\qquad \widehat E=\widehat C-i(\varepsilon(Y))
\tag{8.3.4.1}
\end{equation}
上由 $C$、$\widehat C$ 诱导的子预概形，分别称为 $\SSS$ 所定义的\emph{去顶点仿射锥}、\emph{去顶点射影锥}，这里稍滥用术语。应注意，尽管如此称呼，$E$ 不一定在 $Y$ 上仿射，$\widehat E$ 也不一定在 $Y$ 上射影（参见8.4.3）。把 $C$ 与 $i(C)$ 等同时，对底空间有
\begin{equation}
C\cup\widehat E=\widehat C,\qquad C\cap\widehat E=E.
\tag{8.3.4.2}
\end{equation}
因此，$\widehat C$ 可看作通过粘合开子预概形 $C$ 与 $\widehat E$ 得到。又由（8.3.2.3），
\begin{equation}
E=\widehat E-j(X).
\tag{8.3.4.3}
\end{equation}
若 $Y=\Spec(A)$ 仿射，采用（8.3.2）的记号，有
\begin{equation}
E=\bigcup_f C_f,\qquad \widehat E=\bigcup_f\widehat C_f,
\qquad C_f=C\cap\widehat C_f,
\tag{8.3.4.4}
\end{equation}
其中 $f$ 遍历 $S_+$ 的齐次元素；也可只遍历其中一个子集 $M$，使它生成的 $S_+$-理想在 $S_+$ 中的根为 $S_+$ 本身，等价地，使各 $X_f$（$f\in M$）覆盖 $X$（2.3.14）。沿 $C_f$ 粘合 $C$ 与 $\widehat C_f$，由浸入态射 $C_f\to C$、$C_f\to\widehat C_f$ 确定。如（8.3.2.6）所见，它们分别对应于典范同态 $S\to S_f$、$S_f^{\leq}\to S_f$。

\egastatement{8.3.5}{命题}
采用（8.3.1）、（8.3.4）的记号，与典范单同态 $\varphi:\SSS\to\widehat{\SSS}=\SSS[z]$ 对应的态射（3.5.1）为满射的仿射态射
\begin{equation}
p:\widehat E\longrightarrow X,
\tag{8.3.5.1}
\end{equation}
称为\emph{典范收缩}，并且
\begin{equation}
p\circ j=1_X.
\tag{8.3.5.2}
\end{equation}
\begin{proof}
可设 $Y$ 仿射。利用 $\widehat E$ 的表达式（8.3.4.4），$p$ 的定义域 $G(\varphi)$ 等于 $\widehat E$，将由以下第一个断言得到。

\egapar{8.3.5.3}
\emph{若 $Y=\Spec(A)$ 仿射，$\SSS=\wt S$，则对 $S_+$ 中每个齐次元素 $f$，有}
\begin{equation}
p^{-1}(X_f)=\widehat C_f;
\tag{8.3.5.4}
\end{equation}
\emph{将 $p|\widehat C_f$ 看作 $\widehat C_f=\Spec(S_f^{\leq})$ 到 $X_f$ 的态射，它对应于典范单同态 $S_{(f)}\to S_f^{\leq}$。若进一步 $f\in S_1$，则 $\widehat C_f$ 同构于 $X_f\otimes_{\Z}\Z[T]$，其中 $T$ 为不定元。}

公式（8.3.5.4）只是（2.8.1.1）的特殊情形，第二个断言就是 $Y$ 仿射时 $\Proj(\varphi)$ 的定义（2.8.1）。于是，公式（8.3.5.2）及 $p$ 的满射性，来自典范同态的复合 $\SSS\to\widehat{\SSS}\to\SSS$ 为 $\SSS$ 的恒等同态。最后，（8.3.5.3）的最后一个断言来自 $f\in S_1$ 时 $S_f^{\leq}\simeq S_{(f)}[T]$（2.2.1）。
\end{proof}

\egastatement{8.3.6}{推论}
$p$ 在 $E$ 上的限制
\begin{equation}
\pi:E\longrightarrow X
\tag{8.3.6.1}
\end{equation}
为满射的仿射态射。若 $Y$ 仿射，$f$ 为 $S_+$ 的齐次元素，则
\begin{equation}
\pi^{-1}(X_f)=C_f,
\tag{8.3.6.2}
\end{equation}
且 $\pi|C_f$ 对应于典范单同态 $S_{(f)}\to S_f$。若进一步 $f\in S_1$，则 $C_f$ 同构于 $X_f\otimes_{\Z}\Z[T,T^{-1}]$，其中 $T$ 为不定元。
\begin{proof}
（8.3.6.2）立即来自（8.3.5.3）和（8.3.2.5），并证明 $\pi$ 的满射性。又已知，浸入 $i$ 在 $C_f$ 上的限制对应于单同态 $S_f^{\leq}\to S_f$（8.3.2）。最后一个断言来自 $f\in S_1$ 时 $S_f\simeq S_{(f)}[T,T^{-1}]$（2.2.1）。
\end{proof}

\egastatement{8.3.7}{注记}
若 $Y$ 仿射，由浸入 $\varepsilon$ 的定义（8.3.2），$E$ 的底空间的元素就是 $S$ 中不包含 $S_+$ 的素理想 $\pp$，不要求它们齐次。对这样的 $\pp$，各 $\pp\cap S_n$ 显然满足（2.1.9）的条件，所以存在唯一的齐次素理想 $\qq$，使对每个 $n$，$\qq\cap S_n=\pp\cap S_n$。底空间之间的映射 $\pi:E\to X$ 可由下式解释：
\begin{equation}
\pi(\pp)=\qq.
\tag{8.3.7.1}
\end{equation}
事实上，为验证此式，只须取 $S_+$ 的一个齐次元素 $f$，使 $\pp\in D(f)$，并注意 $\qq_{(f)}$ 为 $\pp_f$ 在单同态 $S_{(f)}\to S_f$ 下的逆像。

\egastatement{8.3.8}{推论}
若 $\SSS$ 由 $\SSS_1$ 生成，则态射 $p,\pi$ 为有限型。对每个 $x\in X$，纤维 $p^{-1}(x)$ 同构于 $\Spec(k(x)[T])$，纤维 $\pi^{-1}(x)$ 同构于 $\Spec(k(x)[T,T^{-1}])$。
\begin{proof}
由（8.3.5）和（8.3.6）立即得到；只须注意，$Y$ 仿射且 $S$ 由 $S_1$ 生成时，各 $X_f$（$f\in S_1$）覆盖 $X$（2.3.14）。
\end{proof}

\egastatement{8.3.9}{注记}
由分次 $\OO_Y$-代数 $\OO_Y[T]$，其中 $T$ 为不定元，所定义的去顶点仿射锥，等同于
\[
G_m=\Spec(\OO_Y[T,T^{-1}]),
\]
因为它就是 $C_T$，如（8.3.2）所见；更一般的结果见（8.4.4）。此预概形典范地带有“交换 $Y$-群概形”结构。这个概念以后将详细展开，这里可简述如下。一个 $Y$-群概形是带有两个 $Y$-态射 $p:G\times_YG\to G$、$s:G\to G$ 的 $Y$-概形 $G$，它们满足与群的合成律及取逆映射的公理形式上类似的条件。图表
\[
\begin{tikzcd}[column sep=large]
G\times G\times G \arrow[r,"p\times1"] \arrow[d,"1\times p"'] & G\times G \arrow[d,"p"] \\
G\times G \arrow[r,"p"'] & G
\end{tikzcd}
\]
须交换，这是“结合性”；还须满足对应于下列群中事实的条件：映射
\[
\begin{aligned}
(x,y)&\longmapsto(x,x^{-1},y)\longmapsto(x,x^{-1}y)\longmapsto x(x^{-1}y),\\
(x,y)&\longmapsto(x,x^{-1},y)\longmapsto(x,yx^{-1})\longmapsto(yx^{-1})x
\end{aligned}
\]
都应等于 $(x,y)\mapsto y$。例如，对应于第一条复合映射的态射列为
\[
G\times G\xrightarrow{(1,s)\times1}G\times G\times G
\xrightarrow{1\times p}G\times G\xrightarrow{p}G;
\]
第二条态射列由读者同样写出。

由（I，3.4.3）立即可知，在 $Y$-概形 $G$ 上给定 $Y$-群概形结构，等价于对每个 $Y$-预概形 $Z$，在集合 $\Hom_Y(Z,G)$ 上给定群结构，并要求对每个 $Y$-态射 $Z\to Z'$，相应映射 $\Hom_Y(Z',G)\to\Hom_Y(Z,G)$ 为群同态。

在这里考虑的特殊情形 $G_m$ 中，$\Hom_Y(Z,G_m)$ 等同于 $Z\times_YG_m$ 的 $Z$-截面的集合（I，3.3.14），因而等同于 $\Spec(\OO_Z[T,T^{-1}])$ 的 $Z$-截面的集合。最终，与（I，3.3.15）相同的论证表明，这个集合典范等同于环 $\Gamma(Z,\OO_Z)$ 的可逆元素的集合，而它的群结构就是来自这个环的乘法的结构。读者可验证，上述态射 $p,s$ 按以下方式给出：由（1.2.7）、（1.4.6），它们对应于 $\OO_Y$-代数同态
\[
\begin{aligned}
\pi &: \OO_Y[T,T^{-1}]\longrightarrow\OO_Y[T,T^{-1},T',T'^{-1}],\\
\sigma &: \OO_Y[T,T^{-1}]\longrightarrow\OO_Y[T,T^{-1}],
\end{aligned}
\]
并完全由 $\pi(T)=TT'$、$\sigma(T)=T^{-1}$ 确定。

对于任意仿射锥 $C=\Spec(\SSS)$，其中 $\SSS$ 为非负分次拟凝聚 $\OO_Y$-代数，$G_m$ 可看作一个“泛算子域”。这意味着，可典范地定义 $Y$-态射 $G_m\times_Y C\to C$，它具有带群作用的集合的外部运算的形式性质。等价地，正如上面对群概形所作的说明，对每个 $Y$-预概形 $Z$，在 $\Hom_Y(Z,C)$ 上给定以群 $\Hom_Y(Z,G_m)$ 为算子集的外部运算，满足群作用于集合的通常公理，并与 $Y$-态射 $Z\to Z'$ 相容。在当前情形中，态射 $G_m\times_YC\to C$ 由 $\OO_Y$-代数同态
\[
\SSS\longrightarrow\SSS\otimes_{\OO_Y}\OO_Y[T,T^{-1}]=\SSS[T,T^{-1}]
\]
确定；它将截面 $s_n\in\Gamma(U,\SSS_n)$，其中 $U$ 为 $Y$ 的开集，映到截面
\[
s_nT^n\in\Gamma\bigl(U,\SSS\otimes_{\OO_Y}\OO_Y[T,T^{-1}]\bigr).
\]

反之，设给定一个事先没有分次结构的拟凝聚 $\OO_Y$-代数 $\SSS$，以及 $C=\Spec(\SSS)$ 上以 $Y$-群概形 $G_m$ 为算子域的“带群作用的集合的 $Y$-概形”结构，则可由此在 $\SSS$ 上典范地得到 $\OO_Y$-代数的分次结构。事实上，给定 $Y$-态射 $G_m\times_YC\to C$，等价于给定 $\OO_Y$-代数同态 $\psi:\SSS\to\SSS[T,T^{-1}]$。它可写为
\[
\psi=\sum_{n\in\Z}\psi_nT^n,
\]
其中 $\psi_n:\SSS\to\SSS$ 为 $\OO_Y$-模同态；对每个截面 $s\in\Gamma(U,\SSS)$，除有限个指标外均有 $\psi_n(s)=0$，这里 $U$ 为 $Y$ 的开集。可验证，群作用于集合的公理蕴含，各 $\psi_n(\SSS)=\SSS_n$ 在 $\SSS$ 上定义了次数可正可负的 $\OO_Y$-代数分次，且 $\psi_n$ 为相应的投影。由此得到任意仿射 $Y$-概形上的“仿射锥”结构的概念：它以“几何”方式定义，不必预先给定分次。

这里不再展开这一观点；与上述说明相应的精确定义及结果，留给读者表述。

\subsection{向量丛的射影闭包}
\src{168--169}
\egapar{8.4.1}
设 $Y$ 为预概形，$\E$ 为拟凝聚 $\OO_Y$-模。若取 $\SSS$ 为分次 $\OO_Y$-代数 $\Sym_{\OO_Y}(\E)$，则由定义（8.3.1.1），$\widehat{\SSS}$ 等同于 $\Sym_{\OO_Y}(\E\oplus\OO_Y)$。此时，$\SSS$ 定义的仿射锥 $\Spec(\SSS)$ 按定义就是 $\V(\E)$，而 $\Proj(\SSS)$ 按定义就是 $\Pb(\E)$。因此：

\egastatement{8.4.2}{命题}
$Y$ 上向量丛 $\V(\E)$ 的射影闭包，典范同构于 $\Pb(\E\oplus\OO_Y)$；后者的无穷远子预概形典范同构于 $\Pb(\E)$。

\egastatement{8.4.3}{注记}
例如，取 $\E=\OO_Y^r$，$r\geq2$。若 $Y\ne\varnothing$，则由 $\SSS$ 定义的去顶点锥 $E,\widehat E$ 在 $Y$ 上既不仿射，也不射影。第二个断言直接可见，因为 $\widehat C=\Pb(\OO_Y^{r+1})$ 在 $Y$ 上射影，而 $E,\widehat E$ 的底空间都是 $\widehat C$ 中不闭的开集。因此，典范浸入 $E\to\widehat C$、$\widehat E\to\widehat C$ 都不射影（5.5.3），再用（5.5.5，v）即可。另一方面，设 $Y=\Spec(A)$ 仿射，例如取 $r=2$，则 $C=\Spec(A[T_1,T_2])$，$E$ 为 $C$ 在开集 $D(T_1)\cup D(T_2)$ 上诱导的预概形。已见后者不仿射（I，5.5.11）；于是 $\widehat E$ 更不可能仿射，因为 $E$ 是 $\widehat E$ 上截面 $z$ 不为零的开集（8.3.2）。不过，有：

\egastatement{8.4.4}{命题}
若 $\LL$ 为可逆 $\OO_Y$-模，则对 $C=\V(\LL)$ 所对应的去顶点锥 $E,\widehat E$，有典范同构
\begin{equation}
\Spec\left(\bigoplus_{n\in\Z}\LL^{\otimes n}\right)\xrightarrow{\sim}E,
\tag{8.4.4.1}
\end{equation}
\begin{equation}
\V(\LL^{-1})\xrightarrow{\sim}\widehat E.
\tag{8.4.4.2}
\end{equation}
此外，存在 $\V(\LL)$ 的射影闭包到 $\V(\LL^{-1})$ 的射影闭包的典范同构，将前者的零截面、无穷远子预概形分别变成后者的无穷远子预概形、零截面。
\begin{proof}
此处 $\SSS=\bigoplus_{n\geq0}\LL^{\otimes n}$。典范单同态
\[
\SSS\longrightarrow\bigoplus_{n\in\Z}\LL^{\otimes n}
\]
定义典范占优态射
\begin{equation}
\Spec\left(\bigoplus_{n\in\Z}\LL^{\otimes n}\right)
\longrightarrow\V(\LL)=\Spec\left(\bigoplus_{n\geq0}\LL^{\otimes n}\right).
\tag{8.4.4.3}
\end{equation}
只须证明，它是左边的概形到 $E$ 的同构。这是 $Y$ 上的局部问题，可设 $Y=\Spec(A)$ 仿射且 $\LL=\OO_Y$。于是 $\SSS=(A[T])^{\sim}$，$\bigoplus_{n\in\Z}\LL^{\otimes n}=(A[T,T^{-1}])^{\sim}$。$A[T,T^{-1}]$ 是 $A[T]$ 的分式环 $A[T]_T$，所以（8.4.4.3）将左边与 $C=\V(\LL)$ 在 $D(T)$ 上诱导的预概形等同。此开集在 $C$ 中的补集 $V(T)$，是理想 $TA[T]$ 定义的闭子预概形的底空间，也就是 $C$ 的零截面。因此 $E=D(T)$。

同构（8.4.4.2）将由最后一个断言得到：$\V(\LL^{-1})$ 是其射影闭包中无穷远子预概形的补集，$\widehat E$ 是 $C=\V(\LL)$ 的射影闭包中零截面的补集。这两个射影闭包分别为 $\Pb(\LL^{-1}\oplus\OO_Y)$、$\Pb(\LL\oplus\OO_Y)$，而
\[
\LL\oplus\OO_Y=\LL\otimes_{\OO_Y}(\LL^{-1}\oplus\OO_Y).
\]
由（4.1.4），所求典范同构存在；剩下只须验证它交换零截面和无穷远子预概形。

可设 $Y=\Spec(A)$ 仿射，$L=Ac$、$L^{-1}=Ac'$，且典范同构 $L\otimes_A L^{-1}\to A$ 将 $c\otimes c'$ 映到 $1$。此时，$\Sym(L\oplus A)$ 为 $A[z]$ 与 $\bigoplus_{n\geq0}Ac^{\otimes n}$ 的张量积，$\Sym(L^{-1}\oplus A)$ 为 $A[z]$ 与 $\bigoplus_{n\geq0}Ac'^{\otimes n}$ 的张量积。（4.1.4）中定义的同构将 $z^h\otimes c'^{\otimes(n-h)}$ 对应于 $z^{n-h}\otimes c^{\otimes h}$。在 $\Pb(\LL^{-1}\oplus\OO_Y)$ 中，无穷远子预概形为截面 $z$ 的零点集合，零截面为截面 $c'$ 的零点集合；$\Pb(\LL\oplus\OO_Y)$ 有相同的描述。因此，由刚才的显式说明立即得到结论。
\end{proof}

\subsection{函子性质}
\src{169--171}
\egapar{8.5.1}
设 $Y,Y'$ 为两个预概形，$q:Y'\to Y$ 为态射；$\SSS,\SSS'$ 分别为非负分次的拟凝聚 $\OO_Y$-代数、$\OO_{Y'}$-代数。考虑分次代数的 $q$-态射
\begin{equation}
\varphi:\SSS\longrightarrow\SSS'.
\tag{8.5.1.1}
\end{equation}
已知（1.5.6），它典范地对应于态射 $\Phi=\Spec(\varphi):\Spec(\SSS')\to\Spec(\SSS)$，使图表
\begin{equation}
\begin{tikzcd}
C' \arrow[r,"\Phi"] \arrow[d] & C \arrow[d] \\
Y' \arrow[r,"q"'] & Y
\end{tikzcd}
\tag{8.5.1.2}
\end{equation}
交换，这里 $C=\Spec(\SSS)$、$C'=\Spec(\SSS')$。进一步假设 $\SSS_0=\OO_Y$、$\SSS'_0=\OO_{Y'}$；令 $\varepsilon:Y\to C$、$\varepsilon':Y'\to C'$ 为典范浸入（8.3.2），则有交换图
\begin{equation}
\begin{tikzcd}
Y' \arrow[r,"q"] \arrow[d,"\varepsilon'"'] & Y \arrow[d,"\varepsilon"] \\
C' \arrow[r,"\Phi"'] & C.
\end{tikzcd}
\tag{8.5.1.3}
\end{equation}
它对应于图表
\[
\begin{tikzcd}
\SSS \arrow[r,"\varphi"] \arrow[d] & \SSS' \arrow[d] \\
\OO_Y \arrow[r] & \OO_{Y'},
\end{tikzcd}
\]
其中竖直箭头为增广同态。此图的交换性来自 $\varphi$ 是分次代数同态这一假设。

\egastatement{8.5.2}{命题}
若 $E,E'$ 分别为 $\SSS,\SSS'$ 定义的去顶点仿射锥，则 $\Phi^{-1}(E)\subset E'$。进一步，$\Proj(\varphi):G(\varphi)\to\Proj(\SSS)$ 处处有定义，即 $G(\varphi)=\Proj(\SSS')$，当且仅当 $\Phi^{-1}(E)=E'$。
\begin{proof}
第一个断言来自（8.5.1.3）的交换性。为证明第二个断言，可设 $Y=\Spec(A)$、$Y'=\Spec(A')$ 都仿射，$\SSS=\wt S$、$\SSS'=\wt S'$。对 $S_+$ 的齐次元素 $f$，令 $f'=\varphi(f)$，则 $\Phi^{-1}(C_f)=C'_{f'}$（I，2.2.4.1）。$G(\varphi)=\Proj(S')$ 意味着：各 $f'=\varphi(f)$ 在 $S'_+$ 中生成的理想的根为 $S'_+$ 本身（2.8.1、2.3.14）。这又等价于各 $C'_{f'}$ 覆盖 $E'$（8.3.4.4）。
\end{proof}

\egapar{8.5.3}
令 $\widehat\varphi(z)=z$，即可将 $q$-态射 $\varphi$ 典范延拓为分次代数的 $q$-态射
\begin{equation}
\widehat\varphi:\widehat{\SSS}\longrightarrow\widehat{\SSS'}.
\tag{8.5.3.1}
\end{equation}
于是得到态射
\[
\widehat\Phi=\Proj(\widehat\varphi):G(\widehat\varphi)\longrightarrow
\widehat C=\Proj(\widehat{\SSS}),
\]
使图表
\[
\begin{tikzcd}
G(\widehat\varphi) \arrow[r,"\widehat\Phi"] \arrow[d] & \widehat C \arrow[d] \\
Y' \arrow[r,"q"'] & Y
\end{tikzcd}
\]
交换（3.5.6）。令 $i:C\to\widehat C$、$i':C'\to\widehat C'$ 为典范开浸入（8.3.2）。由定义立即得到 $i'(C')\subset G(\widehat\varphi)$，且图表
\begin{equation}
\begin{tikzcd}
C' \arrow[r,"\Phi"] \arrow[d,"i'"'] & C \arrow[d,"i"] \\
G(\widehat\varphi) \arrow[r,"\widehat\Phi"'] & \widehat C
\end{tikzcd}
\tag{8.5.3.2}
\end{equation}
交换。最后，令 $X=\Proj(\SSS)$、$X'=\Proj(\SSS')$，$j:X\to\widehat C$、$j':X'\to\widehat C'$ 为典范闭浸入（8.3.2）。由这些浸入的定义仍可得到 $j'(G(\varphi))\subset G(\widehat\varphi)$，且图表
\begin{equation}
\begin{tikzcd}
G(\varphi) \arrow[r,"\Proj(\varphi)"] \arrow[d,"j'"'] & X \arrow[d,"j"] \\
G(\widehat\varphi) \arrow[r,"\widehat\Phi"'] & \widehat C
\end{tikzcd}
\tag{8.5.3.3}
\end{equation}
交换。

\egastatement{8.5.4}{命题}
若 $\widehat E,\widehat E'$ 分别为 $\SSS,\SSS'$ 定义的去顶点射影锥，则 $\widehat\Phi^{-1}(\widehat E)\subset\widehat E'$。此外，若 $p:\widehat E\to X$、$p':\widehat E'\to X'$ 为典范收缩，则
\[
p'\bigl(\widehat\Phi^{-1}(\widehat E)\bigr)\subset G(\varphi),
\]
且图表
\begin{equation}
\begin{tikzcd}
\widehat\Phi^{-1}(\widehat E) \arrow[r,"\widehat\Phi"] \arrow[d,"p'"'] & \widehat E \arrow[d,"p"] \\
G(\varphi) \arrow[r,"\Proj(\varphi)"'] & X
\end{tikzcd}
\tag{8.5.4.1}
\end{equation}
交换。\transnote{原版在图表前的包含关系中写作 $G(\widehat\varphi)$；按 $p'$ 的靶 $X'$ 以及紧随其后的图表，应为 $G(\varphi)$。}
若 $\Proj(\varphi)$ 处处有定义，则 $\widehat\Phi$ 也处处有定义，并且 $\widehat\Phi^{-1}(\widehat E)=\widehat E'$。
\begin{proof}
第一个断言来自图（8.5.1.3）、（8.5.3.2）的交换性，接着两个断言来自典范收缩的定义（8.3.5）及 $\widehat\varphi$ 的定义。另一方面，为证明 $\Proj(\varphi)$ 处处有定义时 $\widehat\Phi$ 也处处有定义，可设 $Y=\Spec(A)$、$Y'=\Spec(A')$ 仿射，$\SSS=\wt S$、$\SSS'=\wt S'$。假设是：$f$ 遍历 $S_+$ 中的齐次元素时，各 $\varphi(f)$ 在 $S'_+$ 中生成的理想，其在 $S'_+$ 中的根为 $S'_+$ 本身。立即得到，$z$ 与各 $\varphi(f)$ 在 $(S'[z])_+$ 中生成的理想，其根为 $(S'[z])_+$ 本身，证得所述断言。同时这证明 $\widehat E'$ 为各 $\widehat C'_{\varphi(f)}$ 的并，因而等于 $\widehat\Phi^{-1}(\widehat E)$。
\end{proof}

\egastatement{8.5.5}{推论}
当 $\Proj(\varphi)$ 处处有定义时，$\widehat C$ 的无穷远子预概形（分别地，顶点预概形）的底空间在 $\widehat\Phi$ 下的逆像，为 $\widehat C'$ 的无穷远子预概形（分别地，顶点预概形）的底空间。\transnote{原版此处把 $\widehat C$ 和 $\widehat C'$ 的撇号写反；此处按已定义的箭头 $\widehat\Phi:\widehat C'\to\widehat C$ 标明源与靶。}
\begin{proof}
利用（8.3.4.1）、（8.3.4.2），由（8.5.4）和（8.5.2）立即得到。
\end{proof}

\subsection{去顶点锥的一个典范同构}
\src{171--173}
\egapar{8.6.1}
设 $Y$ 为预概形，$\SSS$ 为非负分次拟凝聚 $\OO_Y$-代数，且 $\SSS_0=\OO_Y$；令 $X$ 为 $Y$-概形 $\Proj(\SSS)$。现将8.5的结果应用于 $Y'=X$、$q:X\to Y$ 为结构态射的情形。令
\begin{equation}
\SSS_X=\bigoplus_{n\in\Z}\OO_X(n),
\tag{8.6.1.1}
\end{equation}
这是拟凝聚分次 $\OO_X$-代数，其乘法由典范同态（3.2.6.1）
\[
\OO_X(m)\otimes_{\OO_X}\OO_X(n)\longrightarrow\OO_X(m+n)
\]
定义；结合性由交换图（2.5.11.4）保证。取 $\SSS'$ 为 $\SSS_X$ 的非负分次拟凝聚子 $\OO_X$-代数
\[
\SSS_X^{\geq}=\bigoplus_{n\geq0}\OO_X(n).
\]
最后，考虑典范 $q$-态射
\begin{equation}
\alpha:\SSS\longrightarrow\SSS_X^{\geq}.
\tag{8.6.1.2}
\end{equation}
它在（3.3.2.3）中定义为同态 $\SSS\to q_*(\SSS_X)$，但显然把 $\SSS$ 映入 $q_*(\SSS_X^{\geq})$。令
\begin{equation}
C_X=\Spec(\SSS_X^{\geq}),\quad
\widehat C_X=\Proj(\SSS_X^{\geq}[z]),\quad
X'=\Proj(\SSS_X^{\geq}),
\tag{8.6.1.3}
\end{equation}
并以 $E_X,\widehat E_X$ 分别表示相应的去顶点仿射锥、去顶点射影锥。以
\[
\begin{gathered}
\varepsilon_X:X\to C_X,\qquad i_X:C_X\to\widehat C_X,\qquad
j_X:X'\to\widehat C_X,\\
p_X:\widehat E_X\to X',\qquad \pi_X:E_X\to X'
\end{gathered}
\]
表示8.3中定义的典范态射。\transnote{原版列出 $i_X$ 时将其源写成 $C$；此处按构造写成 $C_X$。}

\egastatement{8.6.2}{命题}
结构态射 $u:X'\to X$ 为同构，且态射 $\Proj(\alpha)$ 处处有定义并等于 $u$。态射 $\Proj(\widehat\alpha):\widehat C_X\to\widehat C$ 处处有定义，其在 $\widehat E_X$、$E_X$ 上的限制分别为到 $\widehat E$、$E$ 的同构。最后，通过 $u$ 将 $X'$ 与 $X$ 等同时，态射 $p_X,\pi_X$ 分别等同于 $X$-预概形 $\widehat E_X,E_X$ 的结构态射。
\begin{proof}
显然可设 $Y=\Spec(A)$ 仿射，$\SSS=\wt S$。于是 $X$ 为各仿射开集 $X_f$ 的并，其中 $f$ 遍历 $S_+$ 中的齐次元素，$X_f$ 的环为 $S_{(f)}$。又由（8.2.7.1），
\begin{equation}
\Gamma(X_f,\SSS_X^{\geq})=S_f^{\geq}.
\tag{8.6.2.1}
\end{equation}
因此 $u^{-1}(X_f)=\Proj(S_f^{\geq})$。若 $f\in S_d$（$d>0$），则 $\Proj(S_f^{\geq})$ 典范同构于 $\Proj((S_f^{\geq})^{(d)})$（2.4.7）；另一方面，$(S_f^{\geq})^{(d)}=(S_f^{(d)})^{\geq}$ 通过 $T\mapsto f/1$ 等同于 $S_{(f)}[T]$（2.2.1）。由（3.1.7），结构态射 $u^{-1}(X_f)\to X_f$ 为同构，得到第一个断言。

为证明第二个断言，注意 $\Proj(\alpha)$ 的限制 $u^{-1}(X_f)\cap G(\alpha)\to X=\Proj(S)$，对应于典范映射 $S\to S_f^{\geq}$，$x\mapsto x/1$（2.6.2）。首先得到 $G(\alpha)=X'$。又因 $u^{-1}(X_f)=(u^{-1}(X_f))_{f/1}$，由（2.8.1.1），$\Proj(\alpha)$ 把 $u^{-1}(X_f)$ 映入 $X_f$；将其限制看作到 $X_f=\Spec(S_{(f)})$ 的态射，它正是 $u$ 的限制。

最后，对 $p_X$ 应用（8.3.5.4），得到
\[
p_X^{-1}(u^{-1}(X_f))=\Spec\bigl((S_f^{\geq})_{f/1}^{\leq}\bigr).
\]
由（8.5.4.1），这个开集是在 $\Proj(\widehat\alpha)$ 下 $p^{-1}(X_f)=\Spec(S_f^{\leq})$ 的逆像（8.3.5.3）。利用（8.2.3.2），$\Proj(\widehat\alpha)$ 在此开集上的限制，对应于（8.2.7.2）的逆同构在 $S_f^{\leq}$ 上的限制，得到第三点。最后一个断言由定义显然。
\end{proof}
还应注意，由交换图（8.5.3.2），$\Proj(\widehat\alpha)$ 在 $C_X$ 上的限制恰为态射 $\Spec(\alpha)$。

\egastatement{8.6.3}{推论}
作为 $X$-概形，$\widehat E_X$ 典范同构于 $\Spec(\SSS_X^{\leq})$，$E_X$ 典范同构于 $\Spec(\SSS_X)$。
\begin{proof}
已知 $p_X,\pi_X$ 仿射（8.3.5、8.3.6），所以利用（1.3.1），只须在 $Y=\Spec(A)$ 仿射且 $\SSS=\wt S$ 时验证。第一个断言来自典范同构（8.2.7.2）
\[
(S_f^{\geq})_{f/1}^{\leq}\xrightarrow{\sim}S_f^{\leq},
\]
以及这些同构与从 $f$ 到 $fg$ 的过渡相容，其中 $f,g$ 为 $S_+$ 的齐次元素。同样，对 $\pi_X$ 应用（8.3.6.2），得到
\[
\pi_X^{-1}(u^{-1}(X_f))=\Spec\bigl((S_f^{\geq})_{f/1}\bigr).
\]
于是，第二个断言来自典范同构 $(S_f^{\geq})_{f/1}\xrightarrow{\sim}S_f$（8.2.7.2）。
\end{proof}
因此可以说，$X$-概形 $\widehat C_X$ 是通过粘合两个在 $X$ 上仿射的 $X$-概形 $C_X=\Spec(\SSS_X^{\geq})$、$\widehat E_X=\Spec(\SSS_X^{\leq})$ 得到的；它们的交为开集 $E_X=\Spec(\SSS_X)$。

\egastatement{8.6.4}{推论}
假设 $\OO_X(1)$ 为可逆 $\OO_X$-模，且 $\SSS_X$ 同构于 $\bigoplus_{n\in\Z}(\OO_X(1))^{\otimes n}$；特别地，当 $\SSS$ 由 $\SSS_1$ 生成时，满足此假设（3.2.5、3.2.7）。则去顶点射影锥 $\widehat E$ 等同于 $X$ 上的秩一向量丛 $\V(\OO_X(-1))$；去顶点仿射锥 $E$ 等同于该向量丛在零截面的补集上诱导的子预概形。在这个等同下，典范收缩 $\widehat E\to X$ 等同于 $X$-概形 $\V(\OO_X(-1))$ 的结构态射。最后，存在典范 $Y$-态射 $\V(\OO_X(1))\to C$，其在 $\V(\OO_X(1))$ 的零截面的补集上的限制，为该补集到去顶点仿射锥 $E$ 的同构。
\begin{proof}
令 $\LL=\OO_X(1)$。则 $\SSS_X^{\geq}=\Sym_{\OO_X}(\LL)$，所以由（8.6.3），$\widehat E_X$ 典范等同于 $\V(\LL^{-1})$，$C_X$ 等同于 $\V(\LL)$。态射 $\V(\LL)\to C$ 为 $\Proj(\widehat\alpha)$ 的限制；因此推论的各断言都是（8.6.2）的特殊情形。
\end{proof}
应注意，$C$ 的顶点预概形的底空间在态射 $\V(\OO_X(1))\to C$ 下的逆像，为 $\V(\OO_X(1))$ 的零截面的底空间（8.5.5）。但一般来说，$C$ 与 $\V(\OO_X(1))$ 中的相应子预概形并不同构。下面研究这个问题。

\subsection{投影锥的爆破}
\src{173--177}
\egapar{8.7.1}
在（8.6.1）的条件下，令 $r=\Proj(\widehat\alpha)$。由（8.5.1.3）和（8.5.3.2），有交换图
\begin{equation}
\begin{tikzcd}
X \arrow[r,"i_X\circ\varepsilon_X"] \arrow[d,"q"'] & \widehat C_X \arrow[d,"r"] \\
Y \arrow[r,"i\circ\varepsilon"'] & \widehat C.
\end{tikzcd}
\tag{8.7.1.1}
\end{equation}
此外，由（8.6.2），$r$ 在零截面的补集 $\widehat C_X-i_X(\varepsilon_X(X))$ 上的限制，为到零截面的补集 $\widehat C-i(\varepsilon(Y))$ 的同构。为简化起见，若设 $Y$ 仿射，$\SSS$ 为有限型且由 $\SSS_1$ 生成，则 $X$ 在 $Y$ 上射影，$\widehat C_X$ 在 $X$ 上射影（5.5.1），故 $\widehat C_X$ 在 $Y$ 上射影（5.5.5，ii），从而在 $\widehat C$ 上射影（5.5.5，v）。这样得到射影 $Y$-态射 $r:\widehat C_X\to\widehat C$；它在 $C_X$ 上的限制为射影 $Y$-态射 $C_X\to C$。它将 $X$ 收缩到 $Y$，并在 $X$ 与 $Y$ 的补集之间诱导同构。因此，$C_X$ 与 $C$ 的关系，类似于爆破预概形与原预概形之间的关系（8.1.3）。事实上，下面将证明 $C_X$ 可与某个分次 $\OO_C$-代数的齐次谱等同。

\egapar{8.7.2}
仍用（8.6.1）的记号。对每个 $n\geq0$，考虑分次 $\OO_Y$-代数 $\SSS$ 的拟凝聚理想
\begin{equation}
\SSS_{[n]}=\bigoplus_{m\geq n}\SSS_m.
\tag{8.7.2.1}
\end{equation}
显然有
\begin{equation}
\SSS_{[0]}=\SSS,\qquad \SSS_{[n]}\subset\SSS_{[m]}\quad(m\leq n),
\tag{8.7.2.2}
\end{equation}
\begin{equation}
\SSS_{[n]}\SSS_{[m]}\subset\SSS_{[m+n]}.
\tag{8.7.2.3}
\end{equation}
考虑与 $\SSS_{[n]}$ 相伴的 $\OO_C$-模；它是 $\OO_C=\wt{\SSS}$ 的拟凝聚理想（1.4.4），记为
\begin{equation}
\J_n=(\SSS_{[n]})^{\sim}.
\tag{8.7.2.4}
\end{equation}
由（8.7.2.2）、（8.7.2.3），利用（1.4.4）及（1.4.8.1），得到类似公式
\begin{equation}
\J_0=\OO_C,\qquad \J_n\subset\J_m\quad(m\leq n),
\tag{8.7.2.5}
\end{equation}
\begin{equation}
\J_n\J_m\subset\J_{n+m}.
\tag{8.7.2.6}
\end{equation}
因此满足（8.1.1）的条件，这使我们引入拟凝聚分次 $\OO_C$-代数
\begin{equation}
\SSS^{\natural}=\bigoplus_{n\geq0}\J_n
=\left(\bigoplus_{n\geq0}\SSS_{[n]}\right)^{\sim}.
\tag{8.7.2.7}
\end{equation}

\egastatement{8.7.3}{命题}
存在典范 $C$-同构
\begin{equation}
h:C_X\xrightarrow{\sim}\Proj(\SSS^{\natural}).
\tag{8.7.3.1}
\end{equation}
\begin{proof}
先设 $Y=\Spec(A)$ 仿射，则 $\SSS=\wt S$，其中 $S$ 为非负分次 $A$-代数，且 $C=\Spec(S)$。由定义（8.7.2.4），采用（8.2.6）的记号，有 $\SSS^{\natural}=(S^{\natural})^{\sim}$。\transnote{原版此处的引文写作（8.2.7.4）；相应定义实际为（8.7.2.4）。}
为定义（8.7.3.1），考虑齐次元素 $f\in S_d$（$d>0$）以及对应的 $f^{\natural}\in S_d^{\natural}$（8.2.6）。$S$-同构（8.2.7.3）定义了 $C$-同构
\begin{equation}
\Spec(S_f^{\geq})\xrightarrow{\sim}\Spec(S^{\natural}_{(f^{\natural})}).
\tag{8.7.3.2}
\end{equation}
采用（8.6.2）的记号，若 $v:C_X\to X$ 为结构态射，则由（8.6.2.1），$v^{-1}(X_f)=\Spec(S_f^{\geq})$。另一方面，$\Spec(S^{\natural}_{(f^{\natural})})=D_+(f^{\natural})$，所以（8.7.3.2）定义同构 $v^{-1}(X_f)\to D_+(f^{\natural})$。若 $g\in S_e$（$e>0$），由同构（8.2.7.3）的定义立即可见，图表
\[
\begin{tikzcd}[column sep=large]
v^{-1}(X_{fg}) \arrow[r,"\sim"] \arrow[d] & D_+(f^{\natural}g^{\natural}) \arrow[d] \\
v^{-1}(X_f) \arrow[r,"\sim"] & D_+(f^{\natural})
\end{tikzcd}
\]
交换。按定义，$S_+$ 由其齐次元素生成，因此由（8.2.10，iv）及（2.3.14），各 $D_+(f^{\natural})$ 覆盖 $\Proj(S^{\natural})$。各 $X_f$ 覆盖 $X$，故各 $v^{-1}(X_f)$ 覆盖 $C_X$。这就定义了当前情形中的同构（8.7.3.1）。

为在一般情形证明（8.7.3），只须验证：若 $U',U$ 为 $Y$ 的仿射开集，$U'\subset U$，其环分别为 $A',A$，且 $\SSS|U=\wt S$、$\SSS|U'=\wt S'$，则图表
\begin{equation}
\begin{tikzcd}[column sep=large]
C_{U'} \arrow[r] \arrow[d] & \Proj(S'^{\natural}) \arrow[d] \\
C_U \arrow[r] & \Proj(S^{\natural})
\end{tikzcd}
\tag{8.7.3.3}
\end{equation}
交换。$S'$ 典范等同于 $S\otimes_A A'$，因而 $S'^{\natural}$ 等同于
\[
S^{\natural}\otimes_S S'=S^{\natural}\otimes_A A'.
\]
所以 $\Proj(S'^{\natural})=\Proj(S^{\natural})\times_U U'$（2.8.10）。同样，令 $X=\Proj(S)$、$X'=\Proj(S')$，则 $X'=X\times_U U'$，且 $\SSS_{X'}=\SSS_X\otimes_{\OO_U}\OO_{U'}$（3.5.4）；换言之，$\SSS_{X'}=j^*(\SSS_X)$，其中 $j:X'\to X$ 为投影。因此由（1.5.2），$C_{U'}=C_U\times_X X'=C_U\times_U U'$，图（8.7.3.3）的交换性遂为直接的。
\end{proof}

\egastatement{8.7.4}{注记}
（i）（8.7.3）证明的最后部分立即推广如下。设 $g:Y'\to Y$ 为态射，$\SSS'=g^*(\SSS)$，$X'=\Proj(\SSS')$。则有交换图
\begin{equation}
\begin{tikzcd}[column sep=large]
C_{X'} \arrow[r] \arrow[d] & \Proj(\SSS'^{\natural}) \arrow[d] \\
C_X \arrow[r] & \Proj(\SSS^{\natural}).
\end{tikzcd}
\tag{8.7.4.1}
\end{equation}
另一方面，设 $\varphi:\SSS''\to\SSS$ 为分次 $\OO_Y$-代数同态，令 $X''=\Proj(\SSS'')$，并假设 $u=\Proj(\varphi):X\to X''$ 处处有定义。又有 $Y$-态射 $v:C\to C''$，其中 $C''=\Spec(\SSS'')$，使 $\A(v)=\varphi$。由于 $\varphi$ 为分次代数同态，由它得到分次代数的 $v$-态射 $\psi:\SSS''^{\natural}\to\SSS^{\natural}$（1.4.1）。由（8.2.10，iv）及关于 $\varphi$ 的假设，容易得到 $\Proj(\psi)$ 处处有定义。最后，由（3.5.6.1），有典范 $u$-态射 $\SSS_{X''}\to\SSS_X$，从而得到态射 $w:C_X\to C_{X''}$（1.5.6）。此时图表
\begin{equation}
\begin{tikzcd}[column sep=large]
C_X \arrow[r,"\sim"] \arrow[d,"w"'] & \Proj(\SSS^{\natural}) \arrow[d,"\Proj(\psi)"] \\
C_{X''} \arrow[r,"\sim"] & \Proj(\SSS''^{\natural})
\end{tikzcd}
\tag{8.7.4.2}
\end{equation}
交换，将 $Y$ 归结为仿射即可立即验证。\transnote{原版此处将 $w$ 写成反方向 $C_{X''}\to C_X$，图（8.7.4.2）的上下行撇号亦有错置。上述写法依已经给定的 $u:X\to X''$、$\psi:\SSS''^{\natural}\to\SSS^{\natural}$ 以及谱的反变方向，明确同一自然性图表的源和靶。}

（ii）由（8.7.2.5）、（8.7.2.6），对每个 $m>0$，有 $\J_1^m\subset\J_m\subset\J_1$。按定义，$\J_1=(\SSS_+)^{\sim}$，所以 $\J_1$ 在 $C$ 中定义闭子预概形 $\varepsilon(Y)$（1.4.10、8.3.2）。因此，对每个 $m>0$，$\OO_C/\J_m$ 的支集包含在顶点预概形 $\varepsilon(Y)$ 的底空间中。于是，在去顶点仿射锥 $E$ 的逆像上，结构态射 $\Proj(\SSS^{\natural})\to C$ 为同构；这也可由（8.7.3）和（8.7.1）得到。

把 $C$ 典范等同于 $\widehat C$ 的一个开集（8.3.3），显然可将 $\OO_C$ 的理想 $\J_m$ 延拓为 $\OO_{\widehat C}$ 的理想 $\I_m$，条件是在 $\widehat C$ 的开集 $\widehat E$ 上与 $\OO_{\widehat C}$ 一致。令 $\TT=\bigoplus_{m\geq0}\I_m$，这是拟凝聚分次 $\OO_{\widehat C}$-代数。同构（8.7.3.1）可延拓为 $\widehat C$-同构
\begin{equation}
\widehat C_X\xrightarrow{\sim}\Proj(\TT).
\tag{8.7.4.3}
\end{equation}
事实上，在 $\widehat E$ 上，由前述结果，$\Proj(\TT)$ 典范等同于 $\widehat E$；于是，要求（8.7.4.3）在 $\widehat E$ 上与典范同构 $\widehat E_X\to\widehat E$（8.6.2）一致，即定义了这一同构的相应部分。它与（8.7.3.1）显然在公共部分 $E$ 上一致。\transnote{原版最后一句写作“在 $\widehat E$ 上一致”；两部分的公共定义区域为 $C\cap\widehat E=E$，故此处明确写成 $E$。}

\egastatement{8.7.5}{推论}
假设存在 $n_0>0$，使
\begin{equation}
\SSS_{n+1}=\SSS_1\SSS_n\qquad(n\geq n_0).
\tag{8.7.5.1}
\end{equation}
则 $C_X$ 的顶点子预概形（与 $X$ 同构），是 $C$ 的顶点子预概形（与 $Y$ 同构）在典范态射 $r:C_X\to C$ 下的逆像。反之，若这一性质成立，并进一步假设 $Y$ 诺特、$\SSS$ 为有限型，则存在 $n_0>0$，使（8.7.5.1）成立。
\begin{proof}
第一个断言为 $Y$ 上的局部问题，可设 $Y=\Spec(A)$ 仿射，$\SSS=\wt S$，其中 $S$ 为非负分次 $A$-代数。此时由（8.2.12）得到结论，因为通过等同（8.7.3.1），
\[
\Proj(S^{\natural}\otimes_S S_0)=C_X\times_C\varepsilon(Y),
\]
即这个预概形是 $\varepsilon(Y)$ 在 $C_X$ 中的逆像（I，4.4.1）。当 $Y$ 仿射诺特且 $S$ 为有限型时，逆命题也来自（8.2.12）。

若 $Y$ 诺特而不必仿射，$\SSS$ 为有限型，则存在由诺特仿射开集 $U_i$ 构成的有限覆盖。由前述结果，对每个 $i$，存在整数 $n_i$，使 $n\geq n_i$ 时
\[
\SSS_{n+1}|U_i=(\SSS_1|U_i)(\SSS_n|U_i).
\]
取各 $n_i$ 中的最大值，即满足（8.7.5.1）。
\end{proof}

\egapar{8.7.6}
现考虑在仿射锥 $C$ 中爆破顶点子预概形 $\varepsilon(Y)$ 所得的 $C$-预概形 $Z$。按定义（8.1.3），它为 $\Proj(\bigoplus_{n\geq0}\J_1^n)$，即由各 $\SSS_+^n$ 在 $C$ 上的相伴层构成的分次代数的齐次谱。典范单同态
\begin{equation}
\iota:\bigoplus_{n\geq0}\SSS_+^n\longrightarrow\SSS^{\natural}
\tag{8.7.6.1}
\end{equation}
通过（8.7.3）的等同，定义典范占优 $C$-态射
\begin{equation}
G(\iota)\longrightarrow Z,
\tag{8.7.6.2}
\end{equation}
其中 $G(\iota)$ 为 $C_X$ 的开集（3.5.1）。应注意，可能有 $G(\iota)\ne C_X$。例如，取 $Y=\Spec(K)$，$K$ 为域，$\SSS=\wt S$，其中 $S=K[y]$，$y$ 为次数为 $2$ 的不定元。若 $R_n$ 表示把 $(S_+)^n$ 看作 $S_{[n]}=S_n^{\natural}$ 的子集，则 $S_+^{\natural}$ 不是各 $R_n$ 的并生成的理想在 $S_+^{\natural}$ 中的根（参见2.3.14）。

\egastatement{8.7.7}{推论}
假设存在 $n_0>0$，使
\begin{equation}
\SSS_n=\SSS_1^n\qquad(n\geq n_0).
\tag{8.7.7.1}
\end{equation}
则典范态射（8.7.6.2）处处有定义，并为同构 $C_X\xrightarrow{\sim}Z$。反之，若这一性质成立，并进一步假设 $Y$ 诺特、$\SSS$ 为有限型，则存在 $n_0$，使（8.7.7.1）成立。
\begin{proof}
第一个断言为 $Y$ 上的局部问题，所以来自（8.2.14）。逆命题同样成立，论证与（8.7.5）相同。
\end{proof}

\egastatement{8.7.8}{注记}
条件（8.7.7.1）蕴含（8.7.5.1）。所以，当它成立时，$C_X$ 不仅等同于爆破仿射锥 $C$ 的顶点（与 $Y$ 等同）所得的预概形，而且 $C_X$ 的顶点（与 $X$ 等同）等同于 $C$ 的顶点 $Y$ 的逆像闭子预概形。此外，假设（8.7.7.1）蕴含，$X=\Proj(\SSS)$ 上的 $\OO_X$-模 $\OO_X(n)$ 都可逆（3.2.5、3.2.9），并且若 $\LL=\OO_X(1)$，则 $\OO_X(n)=\LL^{\otimes n}$（3.2.7、3.2.9）。因此，按定义（8.6.1.1），$C_X$ 就是 $X$ 上的向量丛 $\V(\LL)$，其顶点就是此向量丛的零截面。

\subsection{丰沛层与收缩}
\src{177--182}

\egapar{8.8.1} 设 $Y$ 是预概形，$f:X\to Y$ 是分离拟紧态射，$\LL$ 是相对于 $f$ 丰沛的可逆 $\OO_X$-模。考虑非负分次 $\OO_Y$-代数
\begin{equation}\SSS=\OO_Y\oplus\bigoplus_{n\geq1}f_*(\LL^{\otimes n}),\tag{8.8.1.1}\end{equation}
它是拟凝聚的（I，9.2.2，a））。有典范的分次 $\OO_X$-代数同态
\begin{equation}\tau:f^*(\SSS)\longrightarrow\bigoplus_{n\geq0}\LL^{\otimes n},\tag{8.8.1.2}\end{equation}
在次数 $\geq1$ 的部分，它等于典范同态 $\sigma:f^*(f_*(\LL^{\otimes n}))\to\LL^{\otimes n}$（0，4.4.3），在次数 $0$ 的部分则为恒等同态。由 $\LL$ 的 $f$-丰沛性（4.6.3及3.6.1），相应的 $Y$-态射
\begin{equation}r=r_{\LL,\tau}:X\longrightarrow P=\Proj(\SSS)\tag{8.8.1.3}\end{equation}
处处有定义，并且是占优开浸入；而且，对每个 $n\in\Z$，有
\begin{equation}r^*(\OO_P(n))=\LL^{\otimes n}.\tag{8.8.1.4}\end{equation}

\egastatement{8.8.2}{命题} 设 $C=\Spec(\SSS)$ 是由 $\SSS$ 定义的仿射锥。若 $\LL$ 是 $f$-丰沛的，则有典范的 $Y$-态射
\begin{equation}g:V=\V(\LL)\longrightarrow C,\tag{8.8.2.1}\end{equation}
使图
\begin{equation}\begin{tikzcd}
X\arrow[r,"j"]\arrow[d,"f"']&\V(\LL)\arrow[r,"\pi"]\arrow[d,"g"]&X\arrow[d,"f"]\\
Y\arrow[r,"\varepsilon"']&C\arrow[r,"\psi"']&Y
\end{tikzcd}\tag{8.8.2.2}\end{equation}
交换；这里 $\psi,\pi$ 是结构态射，$j,\varepsilon$ 是分别把 $X,Y$ 映至 $\V(\LL)$ 的零截面和 $C$ 的顶点预概形的典范浸入。此外，$g$ 在 $\V(\LL)-j(X)$ 上的限制是开浸入
\begin{equation}\V(\LL)-j(X)\longrightarrow E=C-\varepsilon(Y),\tag{8.8.2.3}\end{equation}
其目标是与 $\SSS$ 对应的去顶点仿射锥 $E$。
\begin{proof}
沿用（8.8.1）的记号，令 $\SSS_P^{\geq}=\bigoplus_{n\geq0}\OO_P(n)$，$C_P=\Spec(\SSS_P^{\geq})$。由（8.6.2），有典范态射 $h=\Spec(\alpha):C_P\to C$，使图
\begin{equation}\begin{tikzcd}C_P\arrow[r]\arrow[d,"h"']&P\arrow[d,"p"]\\ C\arrow[r,"\psi"']&Y\end{tikzcd}\tag{8.8.2.4}\end{equation}
交换。此外，若 $\varepsilon_P:P\to C_P$ 是典范浸入，则图
\begin{equation}\begin{tikzcd}P\arrow[r,"\varepsilon_P"]\arrow[d,"p"']&C_P\arrow[d,"h"]\\Y\arrow[r,"\varepsilon"']&C\end{tikzcd}\tag{8.8.2.5}\end{equation}
交换（8.7.1.1）；最后，$h$ 在去顶点仿射锥 $E_P$ 上的限制是同构 $E_P\simeq E$（8.6.2）。另一方面，由（8.8.1.4），有 $r^*(\SSS_P^{\geq})=\Sym_{\OO_X}(\LL)$，因而有典范的 $P$-态射 $q:\V(\LL)\to C_P$，交换图
\begin{equation}\begin{tikzcd}\V(\LL)\arrow[r,"\pi"]\arrow[d,"q"']&X\arrow[d,"r"]\\C_P\arrow[r]&P\end{tikzcd}\tag{8.8.2.6}\end{equation}
把 $\V(\LL)$ 识别为积 $C_P\times_P X$（1.5.2）。因为 $r$ 是开浸入，$q$ 也是开浸入（I，4.3.2）。由（8.5.2），$q$ 在 $\V(\LL)-j(X)$ 上的限制将它映入 $E_P$，而图
\begin{equation}\begin{tikzcd}X\arrow[r,"j"]\arrow[d,"r"']&\V(\LL)\arrow[d,"q"]\\P\arrow[r,"\varepsilon_P"']&C_P\end{tikzcd}\tag{8.8.2.7}\end{equation}
交换（这是8.5.1.3的特例）。取 $g=h\circ q$，（8.8.2）的各项断言立即由这些事实得到。
\end{proof}

\egastatement{8.8.3}{注} 再假设 $Y$ 是诺特预概形，$f$ 是固有态射。此时 $r$ 是固有的（5.4.4），因而是闭的；又因为它是占优开浸入，所以必为同构 $X\simeq P$。第三章（III，2.3.5.1）还将证明，$\SSS$ 必为有限型 $\OO_Y$-代数。于是 $\SSS^{\natural}$ 是有限型的 $\SSS^{\natural}_0$-代数（8.2.10，(i)及8.7.2.7）。因为 $C_P$ 与 $\Proj(\SSS^{\natural})$ 在 $C$ 上同构（8.7.3），所以 $h:C_P\to C$ 是射影态射。$r$ 是同构，故 $q:\V(\LL)\to C_P$ 也是同构，于是 $g:\V(\LL)\to C$ 是射影态射。此外，$h$ 在 $E_P$ 上的限制是到 $E$ 的同构，且 $q$ 是同构，所以（8.8.2.3）是同构 $\V(\LL)-j(X)\simeq E$。

若再假设 $\LL$ 相对于 $f$ 甚丰沛，第三章（III，2.3.5.1）还将证明，存在整数 $n_0>0$，使 $n\geq n_0$ 时 $\SSS_n=\SSS_1^n$。由（8.7.7），$\V(\LL)$ 可识别为在仿射锥 $C$ 中沿其顶点子预概形（识别为 $Y$）爆破所得的预概形 $Z$；而 $\V(\LL)$ 的零截面（识别为 $X$）是 $C$ 的顶点子预概形 $Y$ 的逆像。\transnote{原版此处把 $\V(\LL)$ 的零截面写作“识别为 $Y$”；按（8.8.2.2）中的 $j:X\to\V(\LL)$，应为 $X$。}

上述结果中的一部分实际上无须诺特性假设即可建立。

\egastatement{8.8.4}{推论} 设 $Y$ 是预概形（分别地，拟紧概形），$f:X\to Y$ 是固有态射，$\LL$ 是相对于 $f$ 丰沛的可逆 $\OO_X$-模。那么（8.8.2.1）是固有态射（分别地，射影态射），其限制（8.8.2.3）是同构。
\begin{proof}
证明 $g$ 固有时可归结于 $Y$ 仿射，故只须考虑 $Y$ 是拟紧概形的情形。与（8.8.3）相同的论证先给出同构 $r:X\simeq P$；因而 $q$ 也是同构。又因 $h|E_P$ 是同构 $E_P\simeq E$，所以（8.8.2.3）已经是同构。余下只须证明 $g$ 射影。

$f$ 按假设是有限型态射，故可对（8.8.1.2）的同态 $\tau$ 应用（3.8.5）：存在整数 $d>0$ 及 $\SSS_d$ 的有限型拟凝聚 $\OO_Y$-子模 $\E$，使得若 $\SSS'$ 是由 $\E$ 生成的 $\SSS$ 的 $\OO_Y$-子代数，$\varphi:\SSS'\to\SSS$ 是典范单射，$\tau'=\tau\circ f^*(\varphi)$，则 $r'=r_{\LL,\tau'}$ 是浸入 $X\to P'=\Proj(\SSS')$。

因为 $\varphi$ 单射，$r'$ 还是占优浸入（3.7.6）。对 $r$ 所作的论证表明，$r'$ 是满射闭浸入。它分解为
\[X\xrightarrow{r}\Proj(\SSS)\xrightarrow{\Phi}\Proj(\SSS'),\qquad\Phi=\Proj(\varphi),\]
所以 $\Phi$ 也是满射闭浸入。但这蕴涵 $\Phi$ 为同构：可归结于 $Y=\Spec(A)$，$\SSS=\wt S$，$\SSS'=\wt{S'}$，其中 $S'$ 是分次 $A$-代数 $S$ 的分次子代数。对每个齐次 $t\in S'$，$S'_{(t)}$ 是 $S_{(t)}$ 的子环。回顾 $\Proj(\varphi)$ 的定义（2.8.1），只须证明：若 $B'$ 是 $B$ 的子环，且由单射 $B'\to B$ 对应的态射 $\Spec(B)\to\Spec(B')$ 是闭浸入，则此态射是同构；这正是（I，4.2.3）的结果。

还有 $\Phi^*(\OO_{P'}(n))=\OO_P(n)$（3.5.2，(ii)及3.5.4），所以 $r'^*(\OO_{P'}(n))\simeq\LL^{\otimes n}$（4.6.3）。置 $\SSS''=\SSS'^{(d)}$；由（3.1.8，(i)），$X$ 典范地识别为 $P''=\Proj(\SSS'')$，而 $\LL''=\LL^{\otimes d}$ 识别为 $\OO_{P''}(1)$（3.2.9，(ii)）。

若 $C''=\Spec(\SSS'')$，则 $\SSS_{P''}^{\geq}=\bigoplus_{n\geq0}\OO_{P''}(n)$ 识别为 $\bigoplus_{n\geq0}\LL''^{\otimes n}$，故 $C_{P''}=\Spec(\SSS_{P''}^{\geq})$ 识别为 $\V(\LL'')$。另一方面，$C_{P''}$ 与 $\Proj(\SSS''^{\natural})$ 在 $C''$ 上同构（8.7.3）。由 $\SSS''$ 的定义，$\SSS''^{\natural}$ 由 $\SSS''^{\natural}_1$ 生成，且为 $\SSS''^{\natural}_0=\SSS''$ 上的有限型代数（8.2.10，(i)及(iii)）；故 $\Proj(\SSS''^{\natural})$ 在 $C''$ 上射影（5.5.1）。考虑图
\begin{equation}\begin{tikzcd}\V(\LL)\arrow[r,"g"]\arrow[d,"u"']&\Spec(\SSS)=C\arrow[d,"v"]\\\V(\LL'')\arrow[r,"g''"']&\Spec(\SSS'')=C''\end{tikzcd}\tag{8.8.4.1}\end{equation}
其中 $g,g''$ 通过（1.5.6）分别对应于典范 $f$-态射
\[\SSS\longrightarrow\bigoplus_{n\geq0}\LL^{\otimes n},\qquad
\SSS''\longrightarrow\bigoplus_{n\geq0}\LL''^{\otimes n}\]
（3.3.2.3；另见下文8.8.5），$v$ 对应于包含态射 $\SSS''\to\SSS$，$u$ 对应于包含态射 $\bigoplus_{n\geq0}\LL^{\otimes nd}\to\bigoplus_{n\geq0}\LL^{\otimes n}$。此图显然交换（3.3.2）。已经证明 $g''$ 是射影态射；另一方面，$u$ 是有限态射。事实上，这一问题在 $X$ 上局部，可假设 $X$ 仿射，环为 $A$，且 $\LL=\OO_X$；只须注意，$A[T]$ 是其子环 $A[T^d]$ 上的有限型模。因为 $Y$ 是拟紧概形，$C''$ 在 $Y$ 上仿射，所以 $C''$ 也是拟紧概形；故 $g''\circ u$ 是射影态射（5.5.5，(ii)）。由（8.8.4.1）交换，$v\circ g$ 也射影。$v$ 仿射，因而分离，最后得到 $g$ 射影（5.5.5，(v)）。
\end{proof}

\egapar{8.8.5} 回到（8.8.1）的情形。态射 $g:\V(\LL)\to C$ 还有另一种定义，对任意可逆 $\OO_X$-模 $\LL$ 都有效，并不要求丰沛。考虑与（8.8.1.2）的 $\tau$ 对应的 $f$-态射
\begin{equation}\tau^\flat:\SSS\longrightarrow\bigoplus_{n\geq0}\LL^{\otimes n}.\tag{8.8.5.1}\end{equation}
它给出态射 $g':V\to C$（1.5.6）。若 $\pi:V\to X$、$\psi:C\to Y$ 是结构态射，则图
\begin{equation}\begin{gathered}
\begin{tikzcd}X\arrow[d,"f"']&V\arrow[l,"\pi"']\arrow[d,"g'"]\\Y&C\arrow[l,"\psi"]\end{tikzcd}
\qquad
\begin{tikzcd}X\arrow[r,"j"]\arrow[d,"f"']&V\arrow[d,"g'"]\\Y\arrow[r,"\varepsilon"']&C\end{tikzcd}
\end{gathered}\tag{8.8.5.2}\end{equation}
交换（8.5.1.2及8.5.1.3）。下面证明，在 $\LL$ 相对于 $f$ 丰沛的假设下，$g=g'$。

问题在 $Y$ 上局部，故可设 $Y=\Spec(A)$，并据（8.8.1.3）将 $X$ 识别为 $P=\Proj(S)$ 的开集，其中 $S=A\oplus\bigoplus_{n\geq1}\Gamma(X,\LL^{\otimes n})$。\transnote{原版本段的求和下限印作 $n\geq0$；为与（8.8.1.1）的零次部分一致，这里写作 $n\geq1$。}由（8.8.1.4），对每个 $n\in\Z$，有 $\Gamma(X,\OO_P(n))=\Gamma(X,\LL^{\otimes n})$。$h=\Spec(\alpha)$ 中的 $\alpha$ 是典范 $p$-态射 $\wt S\to\SSS_P^{\geq}$（8.6.1.2），所以要验证 $\alpha^\#:p^*(\wt S)\to\SSS_P^{\geq}$ 在 $X$ 上的限制等于 $\tau$。由（0，4.4.3），只须验证，对每个 $n>0$，将典范同态 $\alpha_n:S_n\to\Gamma(P,\OO_P(n))$ 与限制同态
\[\Gamma(P,\OO_P(n))\longrightarrow\Gamma(X,\OO_P(n))=\Gamma(X,\LL^{\otimes n})\]
复合，所得为恒等同态。这立即由代数 $S$ 及 $\alpha_n$ 的定义（2.6.2）得到。

\egastatement{8.8.6}{命题} 沿用（8.8.5）的记号，设 $f=(f_0,\lambda)$ 中的同态 $\lambda:\OO_Y\to f_*(\OO_X)$ 为双射。那么：
\begin{enumerate}[label=(\roman*)]
\item 若 $g=(g_0,\mu)$，则 $\mu:\OO_C\to g_*(\OO_V)$ 是同构。
\item 若 $X$ 整（分别地，局部整且正规），则 $C$ 整（分别地，正规）。
\end{enumerate}
\begin{proof}
此时 $f$-态射 $\tau^\flat$ 是同构
\[\tau^\flat:\SSS=\psi_*(\OO_C)\longrightarrow f_*(\pi_*(\OO_V))=\psi_*(g_*(\OO_V)),\]
且 $Y$-态射 $g$ 正是满足 $\A(g)=\tau^\flat$ 的态射（1.1.2）。为证明 $\mu$ 是 $\OO_C$-模同构，只須证明 $\A(\mu):\psi_*(\OO_C)\to\psi_*(g_*(\OO_V))$ 是同构（1.4.2）。按定义（1.1.2），$\A(\mu)=\A(g)$，故(i)成立。

\begin{sloppypar}
证明(ii)时，可归结于 $Y$ 仿射，因而 $\SSS=\wt S$，其中 $S=\bigoplus_{n\geq0}\Gamma(X,\LL^{\otimes n})$。若 $X$ 整，则环 $S$ 整（I，7.4.4），所以 $C$ 整（I，5.1.4）。为证明正规性，使用下面的引理。
\end{sloppypar}
\end{proof}

\egastatement{8.8.6.1}{引理} 若 $Z$ 是整正规预概形，则环 $\Gamma(Z,\OO_Z)$ 整且整闭。
\begin{proof}
由（I，8.2.1.1），$\Gamma(Z,\OO_Z)$ 是有理函数域 $R(Z)$ 中诸整闭环 $\OO_z$（$z\in Z$）的交。
\end{proof}

\noindent\textbf{命题（8.8.6）证明续。} 先证明 $V$ 局部整且正规。可归结于 $X=\Spec(A)$，其中 $A$ 整且整闭（6.3.8），并且 $\LL=\OO_X$。此时 $V=\Spec(A[T])$，而 $A[T]$ 整且整闭（[24]，第99页），故断言成立。另一方面，对 $C$ 的每个仿射开集 $U$，$g^{-1}(U)$ 是拟紧的，因为 $g$ 拟紧。$V$ 局部整，所以 $g^{-1}(U)$ 的连通分支是其中开且整的预概形，因而只有有限个；由于 $V$ 正规，这些预概形都正规（6.3.8）。据(i)，$\Gamma(U,\OO_C)=\Gamma(g^{-1}(U),\OO_V)$；由（8.8.6.1），此环是有限个整且整闭的环的直积，故 $C$ 正规（6.3.4）。\hfill$\square$

\subsection{Grauert 丰沛性判别准则：陈述}
\src{182--184}
我们将证明，（8.8.2）中建立的性质刻画了相对于 $f$ 丰沛的 $\OO_X$-模。准确地说，要证明如下判别准则。

\egastatement{8.9.1}{定理（Grauert 判别准则）} 设 $Y$ 是预概形，$p:X\to Y$ 是分离拟紧态射，$\LL$ 是可逆 $\OO_X$-模。$\LL$ 相对于 $p$ 丰沛，当且仅当存在 $Y$-预概形 $C$、$C$ 的 $Y$-截面 $\varepsilon:Y\to C$ 以及 $Y$-态射 $q:\V(\LL)\to C$，满足以下性质：
\begin{enumerate}[label=(\roman*)]
\item 图
\begin{equation}\begin{tikzcd}X\arrow[r,"j"]\arrow[d,"p"']&\V(\LL)\arrow[d,"q"]\\Y\arrow[r,"\varepsilon"']&C\end{tikzcd}\tag{8.9.1.1}\end{equation}
交换，其中 $j$ 是向量丛 $\V(\LL)$ 的零截面。
\item $q$ 在 $\V(\LL)-j(X)$ 上的限制是拟紧开浸入 $\V(\LL)-j(X)\to C$，其像与 $\varepsilon(Y)$ 不相交。
\end{enumerate}
注意，若 $C$ 在 $Y$ 上分离，则在(ii)中可以删去开浸入拟紧的假设。事实上，为证明拟紧性由其他条件推出，可以归结于 $Y$ 仿射，此时由（I，5.5.1，(i)及5.5.10）即得。若假设 $X$ 诺特，也可删去同一假设，因为此时 $V$ 也诺特，再应用（I，6.3.5）即可。

\egastatement{8.9.2}{推论} 若 $p:X\to Y$ 固有，则在（8.9.1）的陈述中，可以要求 $q$ 固有，并把“开浸入”换为“同构”。

形象地说，当 $p:X\to Y$ 固有时，$\LL$ 相对于 $p$ 丰沛，当且仅当可以将向量丛 $\V(\LL)$ 的零截面“收缩”为底预概形 $Y$。一个重要特例是 $Y$ 为域的谱；此时“收缩”就是将 $\V(\LL)$ 的零截面收缩为一个点。

\egapar{8.9.3} （8.9.1）中条件的必要性以及推论（8.9.2），立即由（8.8.2）及（8.8.4）得到。为证明（8.9.1）的充分性，考虑稍为一般的情形。沿用（8.8.2）的记号，置
\[\SSS'=\bigoplus_{n\geq0}\LL^{\otimes n},\qquad V=\V(\LL)=\Spec(\SSS').\]
闭子预概形 $j(X)$，即 $\V(\LL)$ 的零截面，由 $\OO_V$ 的拟凝聚理想层 $\J=(\SSS'_+)^{\sim}$ 定义（1.4.10）。这个 $\OO_V$-模可逆：问题在 $X$ 上局部，只须注意多项式环 $A[T]$ 的理想 $TA[T]$ 是单生成自由 $A[T]$-模。仍由于问题在 $X$ 上局部，立即得到
\[\LL=j^*(\J),\qquad j_*(\LL)=\J/\J^2.\]
另一方面，若 $\pi:\V(\LL)\to X$ 是结构态射，则 $\pi_*(\J)=\SSS'_+$，$\pi_*(\J/\J^2)=\LL$。因而有两个典范同态 $\LL\to\pi_*(\J)\to\LL$：第一个是典范单射 $\LL\to\SSS'_+$，第二个是 $\SSS'_+$ 到 $\SSS'_1=\LL$ 的典范投影，且二者的复合为恒等同态。又可将 $\pi_*(\J)=\SSS'_+=\bigoplus_{n\geq1}\LL^{\otimes n}$ 典范地嵌入积
\[\prod_{n\geq1}\LL^{\otimes n}=\plim_n\pi_*(\J/\J^{n+1}),\]
因为 $\pi_*(\J/\J^{n+1})=\LL\oplus\LL^{\otimes2}\oplus\cdots\oplus\LL^{\otimes n}$。于是还有两个典范同态
\begin{equation}\LL\longrightarrow\plim_n\pi_*(\J/\J^{n+1})\longrightarrow\LL,\tag{8.9.3.1}\end{equation}
其复合为恒等同态。我们要证明的（8.9.1）的推广如下。

\egastatement{8.9.4}{命题} 设 $Y$ 是预概形，$V$ 是 $Y$-预概形，$X$ 是 $V$ 的闭子预概形，由一个作为 $\OO_V$-模可逆的理想 $\J$ 定义。令 $j:X\to V$ 为典范单射，并置 $\LL=j^*(\J)=\J\otimes_{\OO_V}\OO_X$，因而 $j_*(\LL)=\J/\J^2$。假设结构态射 $p:X\to Y$ 分离且拟紧，并满足下列条件：
\begin{enumerate}[label=(\roman*)]
\item 存在有限型 $Y$-态射 $\pi:V\to X$，使 $\pi\circ j=1_X$，因而 $\pi_*(\J/\J^2)=\LL$。
\item 存在 $\OO_X$-模同态 $\varphi:\LL\to\plim_n\pi_*(\J/\J^{n+1})$，使复合
\[\LL\xrightarrow{\varphi}\plim_n\pi_*(\J/\J^{n+1})\xrightarrow{\alpha}\pi_*(\J/\J^2)=\LL\]
为恒等同态，其中 $\alpha$ 是典范同态。
\item 存在 $Y$-预概形 $C$、$C$ 的 $Y$-截面 $\varepsilon$ 以及 $Y$-态射 $q:V\to C$，使图
\begin{equation}\begin{tikzcd}X\arrow[r,"j"]\arrow[d,"p"']&V\arrow[d,"q"]\\Y\arrow[r,"\varepsilon"']&C\end{tikzcd}\tag{8.9.4.1}\end{equation}
交换。
\item $q$ 在 $W=V-j(X)$ 上的限制是到 $C$ 的拟紧开浸入，其像与 $\varepsilon(Y)$ 不相交。
\end{enumerate}
那么 $\LL$ 相对于 $p$ 丰沛。

\subsection{Grauert 丰沛性判别准则：证明}
\src{184--189}

\egastatement{8.10.1}{引理} 设 $\pi:V\to X$ 是态射，$j:X\to V$ 是作为闭浸入的 $V$ 的 $X$-截面，$\J$ 是定义与 $j$ 相伴的闭子预概形的拟凝聚 $\OO_V$-理想。
\begin{enumerate}[label=(\roman*)]
\item 对每个 $n\geq0$，$\pi_*(\OO_V/\J^{n+1})$ 和 $\pi_*(\J/\J^{n+1})$ 都是拟凝聚 $\OO_X$-模，并且 $\pi_*(\OO_V/\J)=\OO_X$，$\pi_*(\J/\J^2)=j^*(\J)$。
\item 若 $X=\{\xi\}=\Spec(k)$，其中 $k$ 是域，则 $\plim_n\pi_*(\OO_V/\J^{n+1})$ 同构于局部环 $\OO_{j(\xi)}$ 关于 $\mm_{j(\xi)}$-预进拓扑的分离完备化。
\item 假设 $\J$ 是可逆 $\OO_V$-模，因而 $\LL=j^*(\J)=\pi_*(\J/\J^2)$ 是可逆 $\OO_X$-模；又假设存在同态 $\varphi:\LL\to\plim_n\pi_*(\J/\J^{n+1})$，使与典范同态 $\alpha$ 的复合 $\LL\to\plim_n\pi_*(\J/\J^{n+1})\to\pi_*(\J/\J^2)$ 为恒等同态。若 $\SSS=\bigoplus_{n\geq0}\LL^{\otimes n}$，则 $\varphi$ 典范地给出 $\OO_X$-代数同构：将 $\SSS$ 关于典范滤过的完备化 $\widehat\SSS$（同构于积 $\prod_{n\geq0}\LL^{\otimes n}$）映到 $\plim_n\pi_*(\OO_V/\J^{n+1})$。
\end{enumerate}
\begin{proof}
首先注意，$\OO_V/\J^{n+1}$ 的支集是 $j(X)$，而 $\J/\J^{n+1}$ 的支集包含于 $j(X)$。在(ii)中，$j(X)$ 是 $V$ 的闭点 $j(\xi)$；按定义，$\pi_*(\OO_V/\J^{n+1})$ 是 $\OO_V/\J^{n+1}$ 在 $j(\xi)$ 的茎。令 $C=\OO_{j(\xi)}$，$\mm$ 为其极大理想，此茎就是 $C/\mm^{n+1}$，故(ii)显然成立。

为证明(i)，注意问题在 $X$ 上局部，可设 $X$ 仿射。取 $V$ 的仿射开集 $U$；$j(X)\cap U$ 是 $j(X)$ 的仿射开集，因此与它同构的 $U_0=\pi(j(X)\cap U)$ 是 $X$ 的仿射开集。对每个仿射开集 $W_0\subset U_0$，$W=\pi^{-1}(W_0)\cap U$ 是 $V$ 的仿射开集，因为 $X$ 是概形（I，5.5.10）。特别地，$U'=U\cap\pi^{-1}(U_0)$ 仿射，且显然 $\pi(U')=U_0$，$j(U_0)=j(X)\cap U$。按定义，
\[\Gamma(W_0,\pi_*(\OO_V/\J^{n+1}))=\Gamma(\pi^{-1}(W_0),\OO_V/\J^{n+1}).\]
而 $\pi^{-1}(W_0)$ 中不属于 $j(W_0)$ 的每一点，都有不遇 $j(X)$ 的开邻域，在该邻域上 $\OO_V/\J^{n+1}=0$。因此此层在 $\pi^{-1}(W_0)$ 与在 $W$ 上的截面一一对应。若 $\pi'$ 是 $\pi$ 在 $U'$ 上的限制，则 $\pi_*(\OO_V/\J^{n+1})|U_0$ 与 $\pi'_*( (\OO_V/\J^{n+1})|U')$ 相同。$U',U_0$ 都仿射，而诸 $U_0$ 覆盖 $X$，所以 $\pi_*(\OO_V/\J^{n+1})$ 拟凝聚（I，1.6.3）。对 $\pi_*(\J/\J^{n+1})$ 的证明相同。

最后证明(iii)。$\SSS$ 就是 $\Sym_{\OO_X}(\LL)$，所以 $\varphi$ 典范地给出 $\OO_X$-代数同态 $\psi:\SSS\to\plim_n\pi_*(\OO_V/\J^{n+1})$（1.7.4）。它把 $\LL^{\otimes n}$ 映入 $\plim_m\pi_*(\J^n/\J^{m+1})$，因而对于所考虑的拓扑连续，可以延拓为同态 $\widehat\psi:\widehat\SSS\to\plim_n\pi_*(\OO_V/\J^{n+1})$。为证明它是同构，可如(i)的证明归结于 $X=\Spec(A)$、$V=\Spec(B)$，$\J=\wt{\frj}$。$\pi$ 对应的单射 $A\to B$ 将 $A$ 识别为 $B$ 的一个与 $\frj$ 互补的子环；$\LL$ 和 $\pi_*(\OO_V/\J^{n+1})$ 分别与 $A$-模 $L=\frj/\frj^2$ 和 $B/\frj^{n+1}$ 相伴。因为 $\J$ 可逆，还可假设 $\frj=Bt$，其中 $t$ 不是 $B$ 的零因子。由 $B=A\oplus Bt$，对每个 $n>0$ 得到
\[B=A\oplus At\oplus At^2\oplus\cdots\oplus At^n\oplus Bt^{n+1}.\]
因此有典范的 $A$-代数同构 $A[[T]]\simeq C=\plim_n B/\frj^{n+1}$，使 $T$ 对应 $t$。又有 $L=A\bar t$，其中 $\bar t$ 是 $t$ 模 $Bt^2$ 的类；$\varphi$ 按假设将 $\bar t$ 映至一个模 $Ct^2$ 与 $t$ 同余的元素 $t'\in C$。关于 $n$ 归纳得到
\[A\oplus At'\oplus\cdots\oplus At'^n\oplus Ct^{n+1}
=A\oplus At\oplus\cdots\oplus At^n\oplus Ct^{n+1},\]
这说明 $\widehat\psi$ 确实对应于 $\prod_{n\geq0}L^{\otimes n}$ 到 $C$ 的同构。
\end{proof}

\egastatement{8.10.2}{引理} 在（8.10.1）的假设下，设 $g:X'\to X$ 是态射，$V'=V\times_X X'$，$\pi':V'\to X'$、$g':V'\to V$ 是典范投影，从而有交换图
\[\begin{tikzcd}V\arrow[d,"\pi"']&V'\arrow[l,"g'"']\arrow[d,"\pi'"]\\X&X'\arrow[l,"g"]\end{tikzcd}\]
则 $j'=j\times1_{X'}$ 是作为闭浸入的 $V'$ 的 $X'$-截面；$\J'=g'^*(\J)\OO_{V'}$ 是定义与 $j'$ 相伴的闭子预概形的拟凝聚 $\OO_{V'}$-理想。此外，
\[\pi'_*(\OO_{V'}/\J'^{n+1})=g^*(\pi_*(\OO_V/\J^{n+1})).\]
最后，$\J'$ 作为 $\OO_{V'}$-模典范同构于 $g'^*(\J)$；特别地，若 $\J$ 可逆，则 $\J'$ 也可逆。
\begin{proof}
$j'$ 是闭浸入由（I，4.3.1）得到；由底预概形扩张的函子性，它是 $X'$-截面。若 $Z,Z'$ 分别是与 $j,j'$ 相伴的闭子预概形，则 $Z'=g'^{-1}(Z)$（I，4.3.1），故第二项断言由（I，4.4.5）得到。其他断言如（8.10.1）归结于 $X,V,X'$，因而 $V'$，都仿射。沿用该引理证明中的记号，置 $X'=\Spec(A')$，则 $V'=\Spec(B')$，$B'=B\otimes_A A'$，$\J'=\wt{\frj'}$，$\frj'=\Im(\frj\otimes_A A')$。于是 $B'/\frj'^{n+1}=(B/\frj^{n+1})\otimes_A A'$。又因为 $\frj$ 是 $B$ 的 $A$-模直和因子，$\frj\otimes_A A'$ 是 $B'$ 的 $A'$-模直和因子，从而典范地识别为 $\frj'$。
\end{proof}

\egastatement{8.10.3}{推论} 假设（8.10.1）的条件成立，并且 $\pi$ 有限型，$\J$ 是可逆 $\OO_V$-模。那么对每个 $x\in X$，纤维 $\pi^{-1}(x)$ 在点 $j(x)$ 的局部环是维数 $1$ 的正则环，因而是整环；其完备化同构于形式幂级数环 $k(x)[[T]]$。此外，$\pi^{-1}(x)$ 只有一个不可约分支包含 $j(x)$。
\begin{proof}
因为 $\pi^{-1}(x)=V\times_X\Spec(k(x))$，由（8.10.2）可归结于 $X$ 为域 $K$ 的谱。$\pi$ 有限型（I，6.3.4，(iv)），所以 $\OO_{j(x)}$ 是诺特局部环，因而关于 $\mm_{j(x)}$-预进拓扑分离（0，7.3.5）。由（8.10.1，(ii)、(iii)），其完备化同构于 $K[[T]]$，故它正则且维数为 $1$（[1]，第17-01页，定理1）。最后，由于此局部环整，$j(x)$ 只属于 $V$ 的有限个不可约分支中的一个（I，5.1.4）。
\end{proof}

\egastatement{8.10.4}{推论} 假设（8.10.1）的条件成立，且 $\J$ 是可逆 $\OO_V$-模。置 $W=V-j(X)$。对每个拟凝聚 $\OO_X$-理想 $\K$，置 $\K_V=\pi^*(\K)\OO_V$，$\K_W=\K_V|W$。则 $\K_V$ 是限制到 $W$ 后为 $\K_W$ 的最大拟凝聚 $\OO_V$-理想。
\begin{proof}
如（8.10.1），问题在 $X,V$ 上局部，可沿用其证明中的记号，取 $\frj=Bt$，其中 $t$ 不是 $B$ 的零因子。此时 $W=\Spec(B_t)$，$\K=\wt{\frk}$，其中 $\frk$ 是 $A$ 的理想；于是 $\pi^*(\K)\OO_V=(\frk B)^{\sim}$（I，1.6.9），$\K_W=(\frk B_t)^{\sim}$。在 $B_t$ 中的典范像为 $\frk B_t$ 的最大 $B$-理想，就是 $\frk B_t$ 的逆像，即满足某个 $n>0$ 时 $t^ns\in\frk B$ 的元素 $s\in B$ 所成的集合。须证此关系蕴涵 $s\in\frk B$，即 $t$ 在 $B/\frk B=(A/\frk)\otimes_A B$ 中的典范像不是零因子；对 $X'=\Spec(A/\frk)$ 应用（8.10.2）即得。
\end{proof}

\egastatement{8.10.5}{推论} 假设（8.10.3）的条件成立。置 $W=V-j(X)$，设 $x\in X$，$\K$ 是拟凝聚 $\OO_X$-理想，$z$ 是 $\pi^{-1}(x)$ 中包含 $j(x)$ 的不可约分支的泛点（8.10.3）。
\begin{enumerate}[label=(\roman*)]
\item 设 $g$ 是 $\OO_V$ 在 $V$ 上的截面，且 $g|W$ 是 $\K_W$ 在 $W$ 上的截面（8.10.4的记号）。则 $g$ 是 $\K_V$ 的截面。若再有 $g(z)\ne0$，且以 $g_m^x$ 表示 $g$ 在 $\Gamma(X,\pi_*(\OO_V/\J^{m+1}))$ 中的典范像 $g_m$ 在点 $x$ 的茎，则存在整数 $m>0$，使 $g_m^x$ 在
\[(\pi_*(\OO_V/\J^{m+1}))_x\otimes_{\OO_x}k(x)\]
中的像非零。
\item 再假设（8.10.1，(iii)）的条件成立。若存在 $\K_V$ 在 $V$ 上的截面 $g$，使 $g(z)\ne0$，则存在整数 $n\geq0$ 及 $\K\LL^{\otimes n}=\K\otimes\LL^{\otimes n}\subset\LL^{\otimes n}$ 的截面 $f$，使 $f(x)\ne0$。若 $g$ 是 $\J$ 的截面，则可取 $n>0$。
\end{enumerate}
\begin{proof}
(i) 按假设，$g|W$ 生成的 $\OO_W$-理想包含于 $\K_W$，所以由（8.10.4），$g$ 生成的 $\OO_V$-理想包含于 $\K_V$，即 $g$ 是 $\K_V$ 的截面。证明第二项断言时，还可设 $X,V$ 仿射，沿用（8.10.1）的记号。纤维 $\pi^{-1}(x)$ 是环 $B'=B\otimes_A k(x)$ 的谱，且 $B'$ 中存在非零因子 $t'$，使 $B'=k(x)\oplus B't'$。$j(x)$ 是 $z$ 的特殊化，$g(z)\ne0$，所以在纤维的局部环中 $g_{j(x)}\ne0$。此局部环分离（8.10.3），嵌入其完备化，故 $g$ 在完备化中的像非零。而该完备化同构于 $\plim_n B'/B't'^{n+1}$（8.10.3）。令 $g'=g\otimes1\in B'$，则有整数 $m$ 使 $g'\notin B't'^{m+1}$，即其像 $g'_m$ 非零。$g'_m$ 正是 $g_m^x$ 的像，故断言成立。

(ii) 由（8.10.1，(iii)），$\pi_*(\OO_V/\J^{m+1})$ 同构于 $\bigoplus_{0\leq k\leq m}\LL^{\otimes k}$。以 $f_k$ 表示对应于 $g_m$ 的 $\bigoplus_{k=0}^m\Gamma(X,\LL^{\otimes k})$ 中元素的第 $k$ 个分量。按(i)选取 $m$，则某个 $k$ 满足 $f_k(x)\ne0$。为证明 $f_k$ 是 $\K\LL^{\otimes k}$ 的截面，只须如上考虑仿射情形，此时结论立即由 $g\in\frk B$ 得到（8.10.4的记号）。最后，$g\in\Gamma(V,\J)$ 蕴涵 $f_0=0$，故可取 $k>0$。
\end{proof}

\egapar{8.10.6} \textbf{（8.9.4）的证明。} 问题在 $Y$ 上局部（4.6.4）。因为 $\varepsilon$ 是 $Y$-截面，可将 $C$ 换为 $\varepsilon(Y)$ 中一点的仿射开邻域 $U$，使 $\varepsilon(Y)\cap U$ 在 $U$ 中闭。换言之，可设 $C$ 仿射，且 $Y$ 是 $C$ 中由拟凝聚理想层 $\I$ 定义的闭子预概形，因而也仿射。$p$ 分离且拟紧，所以 $X$ 是拟紧概形；须证 $\LL$ 丰沛（4.6.4）。据（4.5.2，a）），只须证明：对每个拟凝聚 $\OO_X$-理想 $\K$ 及每个不属于 $\Supp(\OO_X/\K)$ 的点 $x$，存在整数 $n>0$ 和 $\K\otimes\LL^{\otimes n}$ 在 $X$ 上的截面 $f$，使 $f(x)\ne0$。

置 $\K_V=\pi^*(\K)\OO_V$，$\K_W=\K_V|W$，$W=V-j(X)$。$q|W$ 是到 $C$ 的拟紧浸入，所以由（I，9.4.2），$\K_W$ 是某个拟凝聚 $\OO_V$-理想 $\K'_V$ 在 $W$ 上的限制，其中
\[\K'_V=q^*(\K_C)\OO_V,\]
$\K_C$ 是拟凝聚 $\OO_C$-理想。按假设 $q^{-1}(Y)\subset j(X)$，且 $Y$ 由 $\I$ 定义，所以 $q^*(\I)\OO_V$ 在 $W$ 上的限制等于 $\OO_V|W$。因此 $\K_W$ 也是 $q^*(\I\K_C)\OO_V$ 的限制。于是可设 $\K_C\subset\I$，从而
\begin{equation}\K'_V\subset q^*(\I)\OO_V\subset\J,\tag{8.10.6.1}\end{equation}
这里使用了（I，4.4.6）及（8.9.4.1）的交换性。又由（8.10.4），有
\begin{equation}\K'_V\subset\K_V.\tag{8.10.6.2}\end{equation}
据（8.10.3），$j(x)$ 仅属于 $\pi^{-1}(x)$ 的一个不可约分支。令 $z$ 为该分支的泛点，$z'=q(z)$。由（8.10.5），结合（8.10.6.1）及（8.10.6.2），只须证明存在 $\K'_V$ 在 $V$ 上的截面 $g$，使 $g(z)\ne0$。

按假设，$\K$ 在 $x$ 的某个开邻域上等于 $\OO_X$；另一方面，（8.10.3）给出 $z\ne j(x)$，所以 $z\in W$，进而 $(\K_W)_z=\OO_{V,z}$。由定义，$(\K_C)_{z'}=\OO_{C,z'}$。$C$ 仿射，故有 $\K_C$ 在 $C$ 上的截面 $g'$，使 $g'(z')\ne0$。取与 $g'$ 典范对应的 $\K'_V$ 的截面 $g$，即有 $g(z)\ne0$，证明完成。\hfill$\square$

\egastatement{8.10.7}{注} 我们不知道（8.9.4）的条件(ii)是否多余。\transnote{这句话保留原著撰写时的知识状态。}无论如何，若不假设存在满足 $\pi\circ j=1_X$ 的 $Y$-态射 $\pi:V\to X$，结论就不再成立。简述一个反例的构造，其细节要留待后文才能展开。取 $Y=\Spec(k)$，其中 $k$ 是域；$C=\Spec(A)$，其中 $A=k[T_1,T_2]$，$Y$-截面 $\varepsilon$ 对应于增广同态 $A\to k$。令 $C'$ 为在 $C$ 中爆破闭点 $a=\varepsilon(Y)$ 所得的概形；若 $D$ 是 $a$ 在 $C'$ 中的逆像，在 $D$ 中取一个闭点 $b$，再令 $V$ 为在 $C'$ 中爆破 $b$ 所得的概形。$X$ 是 $V$ 的闭子预概形，等于 $a$ 在结构态射 $q:V\to C$ 下的逆像。可以证明 $X$ 是两个不可约分支 $X_1,X_2$ 的并，其中 $X_1$ 是 $b$ 在 $V$ 中的逆像。定义 $X$ 的 $\OO_V$-理想 $\J$ 显然仍可逆，但 $\LL=j^*(\J)$ 并不丰沛：考虑 $\LL$ 在 $X_1$ 上的逆像的“次数”；若 $\LL$ 丰沛，此次数应 $>0$，而一个初等的交数计算表明它等于 $0$。

\subsection{收缩的唯一性}
\src{189--191}

\egastatement{8.11.1}{引理} 设 $U,V$ 是预概形，$h=(h_0,\lambda):U\to V$ 是满射态射。假设：
\begin{enumerate}[label=\arabic*°]
\item $\lambda:\OO_V\to h_*(\OO_U)=(h_0)_*(\OO_U)$ 是同构。
\item $V$ 的底空间识别为 $U$ 的底空间关于关系 $R:h_0(x)=h_0(y)$ 的商空间。当 $h$ 开或闭，特别是当 $h$ 固有时，此条件总成立。
\end{enumerate}
那么对每个预概形 $W$，将 $v=(v_0,\nu):V\to W$ 映至 $u=v\circ h=(u_0,\mu)$ 的映射
\begin{equation}\Hom(V,W)\longrightarrow\Hom(U,W)\tag{8.11.1.1}\end{equation}
是从 $\Hom(V,W)$ 到满足 $u_0$ 在每条纤维 $h_0^{-1}(z)$ 上恒定的态射 $u$ 所成集合的双射。
\begin{proof}
若 $u=v\circ h$，则 $u_0=v_0\circ h_0$，显然 $u_0$ 在每个 $h_0^{-1}(z)$ 上恒定。反之，若 $u$ 有此性质，证明存在唯一 $v$ 使 $u=v\circ h$。满足 $u_0=v_0\circ h_0$ 的连续映射 $v_0:V\to W$ 的存在及唯一性由假设得到，因为 $h_0$ 识别为 $U\to U/R$ 的典范映射。必要时以同构预概形替换 $V$，还可假设 $\lambda$ 为恒等同态。此时
\[\mu:\OO_W\longrightarrow(u_0)_*(\OO_U)=(v_0)_*((h_0)_*(\OO_U))\]
是环层同态，相应的 $\mu^\#:u_0^{-1}(\OO_W)\to\OO_U$ 在每个茎上局部。由于 $(v_0)_*((h_0)_*(\OO_U))=(v_0)_*(\OO_V)$，必有 $\nu=\mu$；只须验证 $\nu^\#:v_0^{-1}(\OO_W)\to\OO_V$ 在每个茎上局部。

对每个 $y\in V$，可取 $x\in U$ 使 $y=h_0(x)$；置 $z=v_0(y)=u_0(x)$，则（0，3.5.5）有分解
\[\mu_x^\#:\OO_z\xrightarrow{\nu_y^\#}\OO_y\xrightarrow{\lambda_x^\#}\OO_x.\]
按假设，$\lambda_x^\#$ 和 $\mu_x^\#$ 是局部同态；特别地，$\lambda_x^\#$ 将 $\OO_y$ 的可逆元映为可逆元。若 $\nu_y^\#$ 把 $\OO_z$ 的某个不可逆元映为 $\OO_y$ 的可逆元，则 $\mu_x^\#$ 也会将该元映为 $\OO_x$ 的可逆元，与假设矛盾。引理得证。
\end{proof}

\egastatement{8.11.2}{推论} 设 $U$ 是整预概形，$V$ 是正规预概形；则每个普遍闭、双有理且根式的态射 $h:U\to V$ 都是同构。
\begin{proof}
若 $h=(h_0,\lambda)$，则由假设 $h_0$ 单射且闭，$h_0(U)$ 在 $V$ 中稠密，所以 $h_0$ 是 $U$ 到 $V$ 的同胚。只须证明 $\lambda:\OO_V\to(h_0)_*(\OO_U)$ 为同构：此时可应用（8.11.1），由于纤维均为单点，（8.11.1.1）是双射，因而 $h$ 是同构。

问题在 $V$ 上局部，可设 $V=\Spec(A)$，其中 $A$ 整且整闭（8.8.6.1）。$h$ 对应于同态 $\varphi:A\to\Gamma(U,\OO_U)$（I，2.2.4），须证它是同构。若 $K$ 是 $A$ 的分式域，按假设 $\Gamma(U,\OO_U)$ 的分式域也是 $K$，$A$ 是 $\Gamma(U,\OO_U)$ 的子环，且 $\varphi$ 为典范单射（I，8.2.7）。由于 $h$ 满足（7.3.11）的假设，$\Gamma(U,\OO_U)$ 是 $A$ 在 $K$ 中整闭包的子环；按假设，它等于 $A$。
\end{proof}

\egastatement{8.11.3}{注} 第三章（III，4.4.11）将证明，当 $V$ 局部诺特时，每个固有拟有限态射 $h:U\to V$ 必为有限态射。原著在此特别提及满足（8.11.2）假设的态射。此时（8.11.2）的结论也由（6.1.15）得到。

\egapar{8.11.4} 现在说明，在 Grauert 判别准则（8.9.1）中，往往可以断言预概形 $C$ 及“收缩”$q$ 在本质上唯一确定。

\egastatement{8.11.5}{引理} 设 $Y$ 是预概形，$p:X\to Y$ 是固有态射，$\LL$ 是可逆且 $p$-丰沛的 $\OO_X$-模，$C$ 是 $Y$-预概形，$\varepsilon:Y\to C$ 是 $Y$-截面，$q:V=\V(\LL)\to C$ 是使（8.9.1.1）交换的 $Y$-态射。再假设，若 $p=(p_0,\theta)$，则 $\theta:\OO_Y\to p_*(\OO_X)$ 是同构。令
\[\SSS'=\bigoplus_{n\geq0}p_*(\LL^{\otimes n}),\qquad C'=\Spec(\SSS'),\]
并令 $q':\V(\LL)\to C'$ 为典范 $Y$-态射（8.8.5）。则存在唯一 $Y$-态射 $u:C'\to C$，使 $q=u\circ q'$。
\begin{proof}
关于 $\theta$ 的假设特别蕴涵 $p$ 满射。由（8.8.4），$q'$ 在 $\V(\LL)-j(X)$ 上的限制是到 $C'-\varepsilon'(Y)$ 的同构，其中 $\varepsilon'$ 是 $C'$ 的顶点截面；故由（8.8.4），$q'$ 固有且满射。由（8.8.6），若 $q'=(q'_0,\tau)$，则 $\tau:\OO_{C'}\to q'_*(\OO_V)$ 是同构。因此可应用（8.11.1），只须证明 $q$ 在每条纤维 $q'^{-1}(z')$ 上恒定。

对 $z'\notin\varepsilon'(Y)$，此条件显然成立。若 $z'\in\varepsilon'(Y)$，有唯一 $y\in Y$ 使 $z'=\varepsilon'(y)$。由（8.8.5.2）的交换性，以及 $q'$ 将 $\V(\LL)-j(X)$ 映入 $C'-\varepsilon'(Y)$，得到 $q'^{-1}(z')=j(p^{-1}(y))$。此时（8.9.1.1）的交换性证明所需断言。
\end{proof}

\egastatement{8.11.6}{推论} 在（8.11.5）的假设下，再假设 $q$ 固有，且它在 $\V(\LL)-j(X)$ 上的限制是到 $C-\varepsilon(Y)$ 的同构。则 $u$ 普遍闭、满射且根式；它在 $C'-\varepsilon'(Y)$ 上的限制是到 $C-\varepsilon(Y)$ 的同构。
\begin{proof}
$q'$ 是 $\V(\LL)-j(X)$ 到 $C'-\varepsilon'(Y)$ 的同构（8.8.4），所以最后一项断言立即由 $q=u\circ q'$ 得到。由（8.8.5.2）及（8.9.1.1）的交换性，$u$ 在闭子预概形 $\varepsilon'(Y)$ 上的限制是到 $\varepsilon(Y)$ 的同构。于是，对 $z'\in\varepsilon'(Y)$，若 $z=u(z')$，则 $u$ 给出同构 $k(z)\simeq k(z')$。这些事实表明 $u$ 双射且根式。若 $\psi,\psi'$ 是 $C\to Y,C'\to Y$ 的结构态射，则 $\psi'=\psi\circ u$；$\psi'$ 分离（1.2.4），所以 $u$ 分离（I，5.5.1，(v)）。在（8.11.5）的证明中已知 $q'$ 满射，而 $q=u\circ q'$ 固有，故由（5.4.3）及（5.4.9），$u$ 普遍闭。
\end{proof}

\egastatement{8.11.7}{命题} 设 $Y$ 是预概形，$X$ 是整预概形，$p:X\to Y$ 是固有态射，$\LL$ 是可逆且 $p$-丰沛的 $\OO_X$-模，$C$ 是正规的 $Y$-预概形，$\varepsilon:Y\to C$ 是 $Y$-截面，$q:V=\V(\LL)\to C$ 是使（8.9.1.1）交换的 $Y$-态射。再假设 $p=(p_0,\theta)$ 中 $\theta:\OO_Y\to p_*(\OO_X)$ 是同构。令 $\SSS'=\bigoplus_{n\geq0}p_*(\LL^{\otimes n})$，$C'=\Spec(\SSS')$，$q':\V(\LL)\to C'$ 为典范 $Y$-态射（8.8.5）。则满足 $q=u\circ q'$ 的唯一 $Y$-态射 $u:C'\to C$ 是同构。\transnote{本条按原版列出的假设翻译。其证明调用（8.11.6），但原版陈述未重述该推论额外要求的“$q$ 固有且在零截面之外为同构”。仅按本条字面假设，例如取 $X=Y=\Spec(k)$、$\LL=\OO_X$、$C=Y$ 及结构态射 $q:\mathbb A^1_k\to\Spec(k)$，结论不成立。应注意这里缺少承接收缩条件的假设，不宜将原版证明的调用理解为已自动满足这些条件。}
\begin{proof}
由（8.8.6），$C'$ 整。$u$ 是底空间的同胚 $C'\to C$，因为由（8.11.6）它双射且闭，所以 $C$ 不可约，因而整。$u$ 在 $C'$ 的一个非空开集上的限制是到 $C$ 的一个开集的同构，故 $u$ 双有理。$C$ 按假设正规，应用（8.11.2）即得结论。
\end{proof}

\egastatement{8.11.8}{注} (i) 注意，$C$ 正规蕴涵 $X$ 也正规。事实上，$C'=\Spec(\SSS')$ 与 $C$ 同构，故正规，且由（8.8.6）为整；于是 $\Proj(\SSS')$ 正规。为此，问题在 $Y$ 上局部；若 $Y$ 仿射，$\SSS'=\wt{S'}$，则 $S'=\Gamma(C',\OO_{C'})$ 整且整闭（8.8.6.1）。对每个齐次 $f\in S'_+$，分次环 $S'_f$ 整且整闭（[13]，卷I，第257、261页）；其零次部分 $S'_{(f)}$ 也整且整闭，因为 $S'_f$ 与 $S'_{(f)}$ 的分式域之交等于 $S'_{(f)}$。于是所需断言由（6.3.4）得到。最后，$X$ 与 $\Proj(\SSS')$ 的一个开子预概形同构（8.8.1），故 $X$ 正规。因此（8.11.7）可以表述如下：若 $X$ 整且正规，$p=(p_0,\theta):X\to Y$ 固有，$\theta:\OO_Y\to p_*(\OO_X)$ 是同构，则对每个 $p$-丰沛的 $\OO_X$-模 $\LL$，存在唯一的方式收缩 $V=\V(\LL)$ 的零截面，得到正规的 $Y$-概形 $C$ 及固有 $Y$-态射 $q:V\to C$。

(ii) 当 $p$ 固有时，可将假设 $p_*(\OO_X)=\OO_Y$ 看作辅助性假设，它并不实质限制结果的一般性。事实上，若它不成立，只須以 $Y$-概形 $Y'=\Spec(p_*(\OO_X))$ 代替 $Y$，并将 $X$ 视作 $Y'$-概形。第三章第4节将再次讨论这一一般方法。

\subsection{投影锥上的拟凝聚层}
\src{191--195}

\egapar{8.12.1} 回到（8.3.1）的假设与记号。设 $\M$ 是分次拟凝聚 $\SSS$-模。为避免混淆，当 $\M$ 被视作不带分次的 $\SSS$-模时，以 $\wt\M$ 表示相伴的拟凝聚 $\OO_C$-模（1.4.3）；当 $\M$ 被视作分次 $\SSS$-模时，以 $\sProj_0(\M)$ 表示相伴的拟凝聚 $\OO_X$-模，即（3.2.2）中记作 $\wt\M$ 的模。再置
\begin{equation}\M_X=\sProj(\M)=\bigoplus_{n\in\Z}\sProj_0(\M(n)).\tag{8.12.1.1}\end{equation}
$\SSS_X$ 是（8.6.1.1）定义的分次拟凝聚 $\OO_X$-代数；通过典范同态（3.2.6.1）
\begin{equation}\begin{aligned}
\OO_X(m)\otimes_{\OO_X}\sProj_0(\M(n))
&\longrightarrow\sProj_0(\SSS(m)\otimes_{\SSS}\M(n))\\
&\longrightarrow\sProj_0(\M(m+n)),
\end{aligned}\tag{8.12.1.2}\end{equation}
可使 $\sProj(\M)$ 成为分次拟凝聚 $\SSS_X$-模，模公理用交换图（2.5.11.4）验证。

若 $Y=\Spec(A)$ 仿射，$\SSS=\wt S$，$\M=\wt M$，其中 $S$ 是分次 $A$-代数，$M$ 是分次 $S$-模，则由定义及（8.2.9.1），对每个齐次 $f\in S_+$，有
\begin{equation}\Gamma(X_f,\sProj(\wt M))=M_f.\tag{8.12.1.3}\end{equation}
现在考虑分次拟凝聚 $\widehat\SSS$-模
\begin{equation}\widehat\M=\M\otimes_{\SSS}\widehat\SSS,\tag{8.12.1.4}\end{equation}
其中 $\widehat\SSS$ 由（8.3.1.1）定义。它给出拟凝聚 $\OO_{\widehat C}$-模 $\sProj_0(\widehat\M)$，另记为
\begin{equation}\M^\square=\sProj_0(\widehat\M).\tag{8.12.1.5}\end{equation}
显然 $\M^\square$ 是关于 $\M$ 的可加正合函子，与直和及归纳极限交换（3.2.4）。

\egastatement{8.12.2}{命题} 沿用（8.3.2）的记号，有典范函子性同构
\begin{equation}i^*(\M^\square)\simeq\wt\M,\qquad j^*(\M^\square)\simeq\sProj_0(\M).\tag{8.12.2.1}\end{equation}
\begin{proof}
由（3.2.3），$i^*(\M^\square)$ 在 $\Spec(\widehat\SSS/(z-1)\widehat\SSS)$ 上典范地识别为 $(\widehat\M/(z-1)\widehat\M)^{\sim}$。于是第一个同构立即由典范同构 $\widehat\M/(z-1)\widehat\M\simeq\M$ 得到（1.4.1）。另一方面，典范浸入 $j:X\to\widehat C$ 对应于核为 $z\widehat\SSS$ 的典范同态 $\widehat\SSS\to\SSS$（8.3.2）。由于 $\widehat\M\otimes_{\widehat\SSS}\SSS=\M$，第二个同态是（3.5.2，(ii)）中典范同态的特例。为验证它是同构，可设 $Y=\Spec(A)$，$\SSS=\wt S$，$\M=\wt M$；回顾（2.8.8），对每个齐次 $f\in S_+$，此同态在 $X_f$ 上的限制为同构，验证立即可得。
\end{proof}
由于（8.12.2.1）的第一个同构存在，也约定称 $\M^\square$ 为 $\OO_C$-模 $\wt\M$ 的射影闭包；这里默认 $\wt\M$ 的数据中包含 $\SSS$-模 $\M$ 的分次。

\egapar{8.12.3} 沿用（8.3.5）的记号，有典范函子性同态
\begin{equation}p^*(\sProj_0(\M))\longrightarrow\M^\square|\widehat E.\tag{8.12.3.1}\end{equation}
它是（3.5.6）中一般定义的同态 $u^\#$ 的特例。若 $Y=\Spec(A)$，$\SSS=\wt S$，$\M=\wt M$，则回顾（2.8.8），并使用（8.2.3.2）和（8.2.5.2），（8.12.3.1）在 $p^{-1}(X_f)=\widehat C_f$ 上的限制，对应于典范同态
\begin{equation}M_{(f)}\otimes_{S_{(f)}}S_f^{\leq}\longrightarrow M_f^{\leq}.\tag{8.12.3.2}\end{equation}

\egapar{8.12.4} 采用（8.5.1）的假设和记号。由（1.5.6），对每个分次拟凝聚 $\SSS$-模 $\M$，一方面有典范的 $\OO_{C'}$-模同构
\begin{equation}\Phi^*(\wt\M)\simeq(q^*(\M)\otimes_{q^*(\SSS)}\SSS')^{\sim};\tag{8.12.4.1}\end{equation}
另一方面，由（3.5.6）有典范的 $\Proj(\varphi)$-态射
\begin{equation}\sProj_0(\M)\longrightarrow
\sProj_0(q^*(\M)\otimes_{q^*(\SSS)}\SSS')|G(\varphi),\tag{8.12.4.2}\end{equation}
以及典范的 $\widehat\Phi$-态射
\begin{equation}\sProj_0(\widehat\M)\longrightarrow
\sProj_0(q^*(\widehat\M)\otimes_{q^*(\widehat\SSS)}\widehat{\SSS'})|G(\widehat\varphi).\tag{8.12.4.3}\end{equation}

\egapar{8.12.5} 现在考虑（8.6.1）的情形，仍用其记号。即在上面取 $Y'=X$，$q:X\to Y$ 为结构态射，$\varphi$ 为典范 $q$-态射（8.6.1.2）。此时有典范同构
\begin{equation}q^*(\M)\otimes_{q^*(\SSS)}\SSS_X^{\geq}\simeq\M_X^{\geq},\tag{8.12.5.1}\end{equation}
其中 $\M_X^{\geq}=\bigoplus_{n\geq0}\sProj_0(\M(n))$。\transnote{这里保留原版的“同构”断言，但非负截断与该张量积之间的相容性需要额外条件。例如取 $Y=\Spec(k)$、$S=k[T]$、$\deg T=1$、$M=S(-1)$，则 $X$ 为一点，式左端同构于 $k[T](-1)$，其零次部分为零；右端的零次部分却为 $kT^{-1}$。原文随后所用的 $M\otimes_S S_f^{\geq}\simeq M_f^{\geq}$ 在这种一般性下不能直接成立。本译稿不把可能需要的修正静默并入正文。}
事实上，可设 $Y=\Spec(A)$ 仿射，$\SSS=\wt S$、$\M=\wt M$，在各仿射开集 $X_f$ 上定义（8.12.5.1），并检验对 $f$ 的齐次倍数的相容性。由（8.6.2.1），左端在 $X_f$ 上的限制为
\[\wt{M'}=\bigl((M\otimes_A S_{(f)})\otimes_{S\otimes_A S_{(f)}}S_f^{\geq}\bigr)^{\sim}.\]
有典范同构 $M\otimes_A S_{(f)}\simeq M\otimes_S(S\otimes_A S_{(f)})$，从而得到 $\wt{M'}\simeq(M\otimes_S S_f^{\geq})^{\sim}$。原文再由（8.2.9.1）将后者典范地识别为右端在 $X_f$ 上的限制，并指出它满足要求的相容条件。

在上述论证中，以 $\widehat\M$ 代替 $\M$、$\widehat\SSS$ 代替 $\SSS$，并相应使用 $\widehat{\SSS_X^{\geq}}$，同样得到典范同构
\begin{equation}q^*(\widehat\M)\otimes_{q^*(\widehat\SSS)}\widehat{\SSS_X^{\geq}}
\simeq\widehat{\M_X^{\geq}}.\tag{8.12.5.2}\end{equation}
回忆结构态射 $u:\Proj(\SSS_X^{\geq})\to X$ 是同构（8.6.2），先由上面及（8.12.4.2）得到典范 $u$-同构
\begin{equation}\sProj_0(\M)\simeq\sProj_0(\M_X^{\geq}).\tag{8.12.5.3}\end{equation}
沿用（8.6.2）证明中的记号，只须验证当 $f\in S_d$ 时，典范同态 $M_{(f)}\otimes_{S_{(f)}}(S_f^{\geq})^{(d)}\to(M_f^{\geq})^{(d)}$ 是同构，此点显然。其次，将（8.12.4.3）应用于典范态射 $r=\Proj(\widehat\alpha):\widehat C_X\to\widehat C$，由（8.12.5.2）得到典范 $r$-态射
\begin{equation}\M^\square\longrightarrow(\M_X^{\geq})^\square.\tag{8.12.5.4}\end{equation}
回忆，$r$ 在去顶点锥 $\widehat E_X,E_X$ 上的限制，分别是到 $\widehat E,E$ 的同构（8.6.2）。此外：

\egastatement{8.12.6}{命题} 典范 $r$-态射（8.12.5.4）在 $\widehat E_X,E_X$ 上的限制是同构
\begin{align}
\M^\square|\widehat E&\simeq(\M_X^{\geq})^\square|\widehat E_X,\tag{8.12.6.1}\\
\wt\M|E&\simeq(\M_X^{\geq})^{\sim}|E_X.\tag{8.12.6.2}
\end{align}
\begin{proof}
如（8.6.2）的证明，归结于 $Y$ 仿射；沿用该证明的记号，并回到（2.8.8）的定义，须证典范同态
\[\widehat M_{(f)}\otimes_{\widehat S_{(f)}}\widehat{(S_f^{\geq})}_{(f/1)}
\longrightarrow\widehat{(M\otimes_S S_f^{\geq})}_{(f/1)}\]
为同构。由（8.2.3.2）及（8.2.5.2），左端典范地识别为 $M_f^{\leq}\otimes_{S_f^{\leq}}((S_f^{\geq})_{f/1})^{\leq}$，再由（8.2.7.2）识别为 $M_f^{\leq}$；右端识别为 $((M_f^{\geq})_{f/1})^{\leq}$，由（8.2.9.2）也识别为 $M_f^{\leq}$，故（8.12.6.1）成立。（8.12.6.2）再由（8.12.6.1）和（8.12.2.1）得到。
\end{proof}

\egastatement{8.12.7}{推论} 在（8.6.3）的识别下，$(\M_X^{\geq})^\square$ 在 $\widehat E_X$ 上的限制识别为 $(\M_X^{\leq})^{\sim}$，在 $E_X$ 上的限制识别为 $\wt{\M_X}$。
\begin{proof}
归结于仿射情形，这由 $((M_f^{\geq})_{f/1})^{\leq}=M_f^{\leq}$ 及 $(M_f^{\geq})_{f/1}=M_f$ 的典范识别得到（8.2.9.2）。
\end{proof}

\egastatement{8.12.8}{命题} 在（8.6.4）的假设下，典范同态（8.12.3.1）为同构。
\begin{proof}
$\Proj(\SSS_X^{\geq})\to X$ 是同构（8.6.2）；结合（8.12.5.4）及（8.12.6.1），可归结于证明 $p_X^*(\sProj_0(\M_X^{\geq}))\to(\M_X^{\geq})^\square|\widehat E_X$ 的相应命题，也就是归结于 $\SSS_1$ 为可逆 $\OO_Y$-模且 $\SSS$ 由 $\SSS_1$ 生成的情形。沿用（8.12.3）的记号，对 $f\in S_1$，有 $S_f^{\leq}=S_{(f)}[1/f]$，典范同态 $M_{(f)}\otimes_{S_{(f)}}S_f^{\leq}\to M_f^{\leq}$ 由 $M_f^{\leq}$ 的定义为同构。
\end{proof}

\egapar{8.12.9} 现在考虑拟凝聚 $\SSS$-模 $\M_{[n]}=\bigoplus_{m\geq n}\M_m$，以及沿用（8.7.2）记号的分次拟凝聚 $\SSS^{\natural}$-模
\begin{equation}\M^{\natural}=\left(\bigoplus_{n\geq0}\M_{[n]}\right)^{\sim}.\tag{8.12.9.1}\end{equation}
（8.7.3）给出典范 $C$-同构 $h:C_X\simeq\Proj(\SSS^{\natural})$。此外：

\egastatement{8.12.10}{命题} 有典范 $h$-同构
\begin{equation}\sProj_0(\M^{\natural})\simeq\wt{\M_X}.\tag{8.12.10.1}\end{equation}
证明如（8.7.3），但以双同构（8.2.9.3）代替（8.2.7.3）。细节留给读者。

\subsection{子层与闭子概形的射影闭包}
\src{195--197}

\egapar{8.13.1} 沿用（8.12.1）的假设与记号，考虑 $\M$ 的拟凝聚 $\SSS$-子模 $\N$，不要求它带有分次。可以考虑与 $\N$ 相伴的拟凝聚 $\OO_C$-模 $\wt\N$，它是 $\wt\M$ 的 $\OO_C$-子模。另一方面，（8.12.2.1）将 $\wt\M$ 识别为 $\M^\square$ 在 $C$ 上的限制。因为典范单射 $i:C\to\widehat C$ 是仿射态射（8.3.2），因而拟紧，所以典范延拓 $(\wt\N)^-$，即包含于 $\M^\square$ 且在 $C$ 上诱导 $\wt\N$ 的最大 $\OO_{\widehat C}$-子模，是拟凝聚的（I，9.4.2）。下面用分次 $\widehat\SSS$-模显式描述它。

\egapar{8.13.2} 对每个整数 $n\geq0$，考虑同态 $\bigoplus_{i\leq n}\M_i\to\M$，它对 $Y$ 的每个开集 $U$，将族 $(s_i)\in\bigoplus_{i\leq n}\Gamma(U,\M_i)$ 映到截面 $\sum_i s_i\in\Gamma(U,\M)$。以 $\N'_n$ 表示 $\N$ 在此同态下的逆像，它是 $\bigoplus_{i\leq n}\M_i$ 的拟凝聚子模。\transnote{原版在这里称 $\N'_n$ 为“$\SSS$-子模”；当 $\SSS$ 有正次部分时，$\bigoplus_{i\leq n}\M_i$ 一般只自然为 $\OO_Y$-模，不对全部 $\SSS$ 的作用封闭。后面拼成 $\overline\N$ 后才有分次 $\widehat\SSS$-模结构。}
再考虑同态 $\bigoplus_{i\leq n}\M_i\to\widehat\M=\M[z]$，将 $(s_i)$ 映为 $\sum_{i\leq n}s_i z^{n-i}\in\Gamma(U,\widehat\M_n)$；以 $\overline\N_n$ 表示 $\N'_n$ 的像。立即可验证，$\overline\N=\bigoplus_{n\geq0}\overline\N_n$ 是 $\widehat\M$ 的分次拟凝聚 $\widehat\SSS$-子模。称 $\overline\N$ 是借助“齐次化变量”$z$，由 $\N$ 齐次化所得的模。注意，若 $\N$ 已是 $\M$ 的分次 $\SSS$-子模，则 $\overline\N$ 识别为 $\widehat\N=\N[z]$ 中次数 $n\geq0$ 的齐次分量 $\widehat\N_n$ 的直和。

\egastatement{8.13.3}{命题} $\OO_{\widehat C}$-模 $\sProj_0(\overline\N)$ 是 $\wt\N$ 到 $\widehat C$ 的典范延拓 $(\wt\N)^-$。
\begin{proof}
由典范延拓的定义（I，9.4.1），问题在 $Y$ 和 $\widehat C$ 上局部。可先设 $Y=\Spec(A)$ 仿射，$\SSS=\wt S$，$\M=\wt M$，$\N=\wt N$，其中 $N$ 是不一定分次的 $S$-子模。又由（8.3.2.6），$\widehat C$ 是诸仿射开集 $\widehat C_z=C$ 和 $\widehat C_f=\Spec(S_f^{\leq})$（$f\in S_+$ 齐次）的并。只须证明：1° $\sProj_0(\overline\N)$ 在 $C$ 上的限制是 $\wt\N$；2°它在每个 $\widehat C_f$ 上的限制，是 $\wt\N$ 在 $C\cap\widehat C_f=\Spec(S_f)$ 上的限制的典范延拓（8.3.2.6）。

第一点中，$\sProj_0(\overline\N)|C$ 识别为 $(\overline N_{(z)})^{\sim}$（8.3.2.4）。由（2.2.5），$\overline N_{(z)}$ 典范地识别为 $\overline N$ 在 $\widehat M/(z-1)\widehat M$ 中的像；再通过后者到 $M$ 的典范同构（8.2.5），由（8.13.2）的定义，此像识别为 $N$。

第二点中，单射 $i:C\cap\widehat C_f\to\widehat C$ 对应于典范单射 $S_f^{\leq}\to S_f$（8.3.2.6）。另一方面，有
\[\Gamma(\widehat C_f,\M^\square)=M_f^{\leq},\quad
\Gamma(\widehat C_f,i_*(\wt\N))=N_f,\quad
\Gamma(\widehat C_f,i_*(i^*(\M^\square)))=M_f,\]
最后一个等式使用（8.12.2.1）。由（I，9.4.2），只须证明 $\overline N_{(f)}\subset\widehat M_{(f)}=M_f^{\leq}$ 典范地识别为 $N_f$ 在单射 $M_f^{\leq}\to M_f$ 下的逆像。令 $d=\deg f>0$。若 $M_f$ 的元素 $(\sum_{k\leq md}x_k)/f^m$，其中 $x_k\in M_k$，也写为 $y/f^m$，$y\in N$，则将 $y$ 和所有 $x_k$ 同乘适当的 $f^h$，可设 $\sum_{k\leq md}x_k=y$。但在（8.2.5.2）的识别中，该元素对应于
\[\frac{\sum_{k\leq md}x_k z^{md-k}}{f^m},\]
它属于 $\overline N_{(f)}$，因为 $\sum_{k\leq md}x_k\in N$。反之显然。
\end{proof}

\egastatement{8.13.4}{注} (i) （8.13.3）最重要的应用是 $\M=\SSS$。此时 $\wt\N$ 是 $\OO_C$ 的任意拟凝聚理想层 $\J$（1.4.3），与 $C$ 的一个闭子预概形 $Z$ 一一对应。$\J$ 的典范延拓 $\overline\J$ 是定义 $Z$ 在 $\widehat C$ 中的闭包 $\overline Z$ 的拟凝聚 $\OO_{\widehat C}$-理想层（I，9.5.10）。命题（8.13.3）给出了利用 $\widehat\SSS=\SSS[z]$ 中分次理想典范地定义 $\overline Z$ 的方法。

(ii) 为简便起见，设 $Y$ 仿射，沿用（8.13.3）证明的记号。对每个非零 $x\in N$，令 $d(x)$ 为 $x$ 在 $M$ 中的非零齐次分量 $x_i$ 的最大次数。按定义，$\overline N$ 由零元及形如
\[h(x,k)=z^k\sum_{i\leq d(x)}x_i z^{d(x)-i}\qquad(k\geq0)\]
的元素组成；因而它作为 $\widehat S=S[z]$-模，由形如
\[h(x,0)=\sum_{i\leq d(x)}x_i z^{d(x)-i}\]
的元素生成。\transnote{原版的“由……元素组成”应结合分次模结构理解：一般还须取不同次数齐次元素的有限和；若允许 $M$ 有负次数，也须遵守（8.13.2）的非负总次数截断。本段保留原著的齐次化公式。}
称 $h(x,0)$ 是借助“齐次化变量”$z$ 由 $x$ 齐次化所得的元素。但 $h(x,0)$ 对 $x$ 不是可加的，当然更不是 $S$-线性的；所以即使 $M=S$，也不能认为当 $x$ 遍历 $S$-模 $N$ 的一组生成元时，$h(x,0)$ 就遍历分次 $\widehat S$-模 $\overline N$ 的一组生成元。不过，当 $N$ 是单生成自由 $S$-模时，情况确实如此——这是初等代数几何中唯一考虑的情形：若 $t$ 是 $N$ 的一组基，则 $h(t,0)$ 生成 $\widehat S$-模 $\overline N$。

\subsection{与分次模相伴的层的补充}
\src{197--203}

\egapar{8.14.1} 设 $Y$ 是预概形，$\SSS$ 是非负分次拟凝聚 $\OO_Y$-代数，$X=\Proj(\SSS)$，$q:X\to Y$ 是结构态射，由（3.1.3）它分离。沿用（8.12.1）的记号，已经定义函子 $\M_X=\sProj(\M)$，从分次拟凝聚 $\SSS$-模范畴到分次拟凝聚 $\SSS_X$-模范畴；它显然可加、正合且与归纳极限交换（3.2.4）。由（8.12.1.1）的定义，对每个 $n\in\Z$，立即得到
\begin{equation}\sProj(\M(n))=(\sProj(\M))(n).\tag{8.14.1.1}\end{equation}

\egapar{8.14.2} 先将（3.2.6）定义的典范同态 $\lambda,\mu$ 推广到形如 $\sProj(\M)$ 的 $\SSS_X$-模。对任意 $m,n\in\Z$，由（2.1.2.1），有典范 $\OO_X$-模同态
\begin{equation}\begin{aligned}
\lambda_{mn}:\sProj_0(\M(m))\otimes_{\OO_X}\sProj_0(\N(n))
\longrightarrow\sProj_0((\M\otimes_{\SSS}\N)(m+n)).
\end{aligned}\tag{8.14.2.1}\end{equation}
同样，由（2.1.2.2），有典范 $\OO_X$-模同态
\begin{equation}\begin{aligned}
\mu_{mn}:&\sProj_0((\shHom_{\SSS}(\M,\N))(n-m))\\
&\longrightarrow\shHom_{\OO_X}(\sProj_0(\M(m)),\sProj_0(\N(n))),
\end{aligned}\tag{8.14.2.2}\end{equation}
其中 $\M,\N$ 是任意分次拟凝聚 $\SSS$-模。由此得到同态
\[\mu_k:\sProj_0((\shHom_{\SSS}(\M,\N))(k))
\longrightarrow(\shHom_{\SSS_X}(\sProj(\M),\sProj(\N)))_k.\]
其定义是：对开集 $U\subset X$ 及
\[u\in\Gamma(U,\sProj_0((\shHom_{\SSS}(\M,\N))(k))),\]
令 $\mu_k(u)$ 为次数 $k$ 的分次 $\Z$-模同态 $\Gamma(U,\sProj(\M))\to\Gamma(U,\sProj(\N))$，它在每个 $\Gamma(U,\sProj_0(\M(m)))$ 上等于 $\mu_{m,m+k}(u)$。回到 $\mu_{mn}$ 的定义（2.5.12.1），立即看到 $\mu_k(u)$ 实为次数 $k$ 的分次 $\Gamma(U,\SSS_X)$-模同态，而且诸 $\mu_k$ 定义分次 $\SSS_X$-模同态
\begin{equation}\sProj(\shHom_{\SSS}(\M,\N))\longrightarrow
\shHom_{\SSS_X}(\sProj(\M),\sProj(\N)).\tag{8.14.2.3}\end{equation}
同样，利用结合性交换图（2.5.11.4），诸同态（8.14.2.1）给出分次 $\SSS_X$-模同态
\begin{equation}\lambda:\sProj(\M)\otimes_{\SSS_X}\sProj(\N)
\longrightarrow\sProj(\M\otimes_{\SSS}\N).\tag{8.14.2.4}\end{equation}

\egastatement{8.14.3}{命题} 同态（8.14.2.4）为双射；若分次 $\SSS$-模 $\M$ 有有限表示（3.1.1），则（8.14.2.3）也为双射。
\begin{proof}
问题在 $X,Y$ 上局部，可设 $Y=\Spec(A)$ 仿射，$\SSS=\wt S$，$\M=\wt M$，$\N=\wt N$，其中 $S$ 是非负分次 $A$-代数，$M,N$ 是分次 $S$-模。若 $f\in S_+$ 齐次，（8.14.2.1）、（8.14.2.2）在仿射开集 $D_+(f)$ 上的限制，分别对应于典范同态（2.5.11.1）和（2.5.12.1）：
\[M(m)_{(f)}\otimes_{S_{(f)}}N(n)_{(f)}\longrightarrow(M\otimes_S N)(m+n)_{(f)},\]
\[(\Hom_S(M,N))(n-m)_{(f)}\longrightarrow
\Hom_{S_{(f)}}(M(m)_{(f)},N(n)_{(f)}).\]
回到这些同态的定义，并利用（8.2.9.1），（8.14.2.4）在 $D_+(f)$ 上的限制对应于（0，1.3.4）中的典范同态 $M_f\otimes_{S_f}N_f\to(M\otimes_S N)_f$，已知它是同构。同样，（8.14.2.3）的限制对应于典范同态
\[(\Hom_S(M,N))_f\longrightarrow\Hom_{S_f}(M_f,N_f)\]
（0，1.3.5）。这里因为 $M$ 有限型，齐次 $S$-模同态所成子群的直和 $\Hom_S(M,N)$（2.1.2）等于所有 $S$-模同态 $M\to N$ 的集合。由 $M$ 有限表示及（0，1.3.5），此典范同态为同构。
\end{proof}

\egastatement{8.14.4}{命题} 若 $U$ 是 $X$ 的拟紧开集，则存在整数 $d$，使对每个 $d$ 的整数倍 $n$，$\OO_X(n)|U$ 可逆，其逆为 $\OO_X(-n)|U$。
\begin{proof}
$q(U)$ 拟紧，可由有限个仿射开集 $V_i$ 覆盖，所以每个 $x\in U$ 包含于形如 $D_+(f)$ 的仿射开集，其中 $f$ 是某个环 $\Gamma(V_i,\SSS)$ 中的正次齐次元素。$U$ 拟紧，故可取有限个这样的 $D_+(f_j)$ 覆盖。令 $d$ 为诸 $f_j$ 的次数的公倍数，由（2.5.17），$d$ 满足要求。
\end{proof}

\egapar{8.14.5} 沿用（8.14.1）的假设与记号，（3.3.2）定义了典范 $\OO_Y$-模同态
\begin{equation}\alpha_n:\M_n\longrightarrow q_*(\sProj_0(\M(n)))\qquad(n\in\Z).\tag{8.14.5.1}\end{equation}
推广（3.3.1）的记号，对每个分次 $\SSS_X$-模 $\F$，置
\begin{equation}\boldsymbol\Gamma_*(\F)=\bigoplus_{n\in\Z}q_*(\F_n).\tag{8.14.5.2}\end{equation}
特别地，$\boldsymbol\Gamma_*(\SSS_X)=\bigoplus_{n\in\Z}q_*(\OO_X(n))$ 是（3.3.1.2）中记作 $\Gamma_*(\OO_X)$ 的分次 $\OO_Y$-代数；显然 $\boldsymbol\Gamma_*(\F)$ 是分次 $\boldsymbol\Gamma_*(\SSS_X)$-模（0，4.2.2）。\transnote{原版在这里写“代数”；$\F$ 仅假设为模，后文也使用其模结构，故此处注明并采用“模”。}
在（8.14.5.1）中取 $\M=\SSS$，得到（3.3.2）已定义的分次 $\OO_Y$-代数同态
\begin{equation}\alpha:\SSS\longrightarrow\boldsymbol\Gamma_*(\SSS_X),\tag{8.14.5.3}\end{equation}
它使 $\boldsymbol\Gamma_*(\F)$ 成为分次 $\SSS$-模。于是诸同态（8.14.5.1）定义次数 $0$ 的分次 $\SSS$-模同态
\begin{equation}\alpha:\M\longrightarrow\boldsymbol\Gamma_*(\sProj(\M)).\tag{8.14.5.4}\end{equation}

\egapar{8.14.6} 一般地，对分次拟凝聚 $\SSS_X$-模 $\F$，不能保证分次 $\SSS$-模 $\boldsymbol\Gamma_*(\F)$ 必拟凝聚。考虑 $X$ 的开集 $X'$，使 $q$ 在它上的限制 $q':X'\to Y$ 拟紧。$q'$ 还分离，故对每个拟凝聚 $\OO_{X'}$-模 $\F'$，$q'_*(\F')$ 是拟凝聚 $\OO_Y$-模（I，9.2.2，b））。置
\begin{equation}\SSS_{X'}=\SSS_X|X'=\bigoplus_{n\in\Z}\OO_X(n)|X',\tag{8.14.6.1}\end{equation}
并对每个分次 $\SSS_{X'}$-模 $\F'$ 置
\begin{equation}\boldsymbol\Gamma'_*(\F')=\bigoplus_{n\in\Z}q'_*(\F'_n).\tag{8.14.6.2}\end{equation}
由上述观察，若 $\F'$ 拟凝聚，则 $\boldsymbol\Gamma'_*(\F')$ 是分次拟凝聚 $\SSS$-模（I，9.6.1）。

典范单射 $j:X'\to X$ 也是拟紧的，因为 $q'=q\circ j$ 拟紧，$q$ 分离（I，6.6.4，(v)）。所以对每个分次拟凝聚 $\SSS_{X'}$-模 $\F'$，$\F=j_*(\F')$ 是分次拟凝聚 $\SSS_X$-模，且由定义
\begin{equation}\boldsymbol\Gamma'_*(\F')=\boldsymbol\Gamma_*(\F).\tag{8.14.6.3}\end{equation}
对 $X'$ 保持相同假设，对每个分次拟凝聚 $\SSS$-模 $\M$，置
\begin{equation}\sProj'(\M)=\sProj(\M)|X',\tag{8.14.6.4}\end{equation}
它是分次拟凝聚 $\SSS_{X'}$-模。典范同态 $\sProj(\M)\to j_*(\sProj'(\M))$（0，4.4.3）给出典范的分次 $\SSS$-模同态
\[\boldsymbol\Gamma_*(\sProj(\M))\longrightarrow\boldsymbol\Gamma'_*(\sProj'(\M)).\]
与（8.14.5.4）复合，得到次数 $0$ 的典范函子性分次拟凝聚 $\SSS$-模同态
\begin{equation}\alpha':\M\longrightarrow\boldsymbol\Gamma'_*(\sProj'(\M)).\tag{8.14.6.5}\end{equation}

\egapar{8.14.7} 保持（8.14.6）关于 $X'$ 的假设，设 $\F'$ 是分次拟凝聚 $\SSS_{X'}$-模，故 $\sProj'(\boldsymbol\Gamma'_*(\F'))$ 也是分次拟凝聚 $\SSS_{X'}$-模。下面定义次数 $0$ 的典范函子性同态
\begin{equation}\beta':\sProj'(\boldsymbol\Gamma'_*(\F'))\longrightarrow\F'.\tag{8.14.7.1}\end{equation}
先设 $Y=\Spec(A)$ 仿射，$\SSS=\wt S$，其中 $S$ 是非负分次 $A$-代数；则 $\boldsymbol\Gamma'_*(\F')=\wt M$，其中 $M=\bigoplus_{n\in\Z}\Gamma(X',\F'_n)$ 是分次 $S$-模。取 $f\in S_d$，使 $D_+(f)\subset X'$。按定义（2.6.2），$\alpha_d(f)$ 在 $D_+(f)$ 上的限制是 $\OO_X(d)$ 的截面，对应于 $(S(d))_{(f)}$ 中的 $f/1$，因而可逆；对每个 $n>0$，$\alpha_{nd}(f^n)$ 也如此。因此，将 $z/f^n\in M_f$ 映到截面
\[(z|D_+(f))(\alpha_{nd}(f^n)|D_+(f))^{-1}\]
便定义了次数 $0$ 的分次 $S_f$-模同态 $\beta_f:M_f\to\Gamma(D_+(f),\F')$。还有与（2.6.4.1）对应的交换图，故在此情形得到 $\beta'$。一般情形下，须考虑 $A$-代数 $A'$、分次 $A'$-代数 $S'=S\otimes_A A'$，并使用与（2.8.13.2）类似的交换图。细节留给读者。

\egastatement{8.14.8}{命题} 若 $X'$ 是 $X=\Proj(\SSS)$ 的开集，且 $q':X'\to Y$ 拟紧，则（8.14.7）定义的同态 $\beta'$ 为双射。
\begin{proof}
显然可设 $Y$ 仿射，沿用（8.14.7）的记号，只须证明 $\beta_f:M_f\to\Gamma(D_+(f),\F')$ 为同构。以 $f$ 的某个幂代替 $f$，并不改变 $D_+(f)$ 或 $\beta_f$。由于 $X'$ 按假设拟紧，由（8.14.4）总可假设 $\OO_X(d)$ 在 $X'$ 上可逆。$X'$ 是概形，因为 $q'$ 分离；此时命题正是（I，9.3.1）。
\end{proof}

\egastatement{8.14.9}{推论} 在（8.14.8）的假设下，每个分次拟凝聚 $\SSS_{X'}$-模 $\F'$ 都同构于形如 $\sProj'(\M)$ 的分次 $\SSS_{X'}$-模，其中 $\M$ 是分次拟凝聚 $\SSS$-模。若 $\F'$ 还有限型，且 $Y$ 是拟紧概形，或底空间诺特的预概形，则可取 $\M$ 有限型。

由（8.14.8）导出本推论的证明，与由（3.4.4）导出（3.4.5）的证明完全相同。细节留给读者。

\egastatement{8.14.10}{命题} 在（8.14.7）的假设下，设 $\M$ 是分次拟凝聚 $\SSS$-模，$\F'$ 是分次拟凝聚 $\SSS_{X'}$-模。下列复合均为恒等同构：
\begin{equation}\begin{aligned}
\sProj'(\M)&\xrightarrow{\sProj'(\alpha')}
\sProj'(\boldsymbol\Gamma'_*(\sProj'(\M)))
\xrightarrow{\beta'}\sProj'(\M),
\end{aligned}\tag{8.14.10.1}\end{equation}
\begin{equation}\begin{aligned}
\boldsymbol\Gamma'_*(\F')&\xrightarrow{\alpha'}
\boldsymbol\Gamma'_*(\sProj'(\boldsymbol\Gamma'_*(\F')))
\xrightarrow{\boldsymbol\Gamma'_*(\beta')}\boldsymbol\Gamma'_*(\F').
\end{aligned}\tag{8.14.10.2}\end{equation}
问题在 $Y$ 上局部，验证与（2.6.5）相同，细节留给读者。

\egastatement{8.14.11}{注} 第三章（III，2.3.1）将证明，若 $Y$ 局部诺特，$\SSS$ 是有限型分次拟凝聚 $\OO_Y$-代数（此时可取 $X'=X$），则对每个满足(TF)的分次拟凝聚 $\SSS$-模 $\M$，同态 $\alpha$（8.14.5.4）是(TN)-双射。

\egastatement{8.14.12}{注} （8.14.4）的情形是下面情形的特例。设 $X$ 是赋环空间，$\SSS$ 是次数可正可负的分次 $\OO_X$-代数。假设存在整数 $d>0$，使 $\SSS_d,\SSS_{-d}$ 可逆，且典范同态
\begin{equation}\SSS_d\otimes_{\OO_X}\SSS_{-d}\longrightarrow\OO_X\tag{8.14.12.1}\end{equation}
为同构，因而 $\SSS_{-d}$ 识别为 $\SSS_d^{-1}$。称分次 $\OO_X$-代数 $\SSS$ 是以 $d$ 为周期的周期代数。这个名称来自以下性质：在上述假设下，对每个分次 $\SSS$-模 $\F$ 及每个 $n\in\Z$，典范同态
\begin{equation}\SSS_d\otimes\F_n\longrightarrow\F_{n+d}\tag{8.14.12.2}\end{equation}
都是同构。事实上，问题在 $X$ 上局部，可设 $\SSS_d$ 有 $X$ 上的可逆截面 $s$，其逆 $s'$ 是 $\SSS_{-d}$ 的截面。将 $z\in\Gamma(U,\F_{n+d})$ 映为 $(s|U)\otimes(s'|U)z$ 的同态 $\F_{n+d}\to\SSS_d\otimes\F_n$，就是（8.14.12.2）的逆，故断言成立。于是对每个 $k\in\Z$，有典范同构
\[(\SSS_d)^{\otimes k}\otimes\F_n\simeq\F_{n+kd}.\]
因此，只要给出 $\SSS_0$-模 $\F_i$（$0\leq i\leq d-1$）及典范同态
\[\SSS_i\otimes\F_j\longrightarrow\F_{i+j}\qquad(0\leq i,j\leq d-1),\]
分次 $\SSS$-模 $\F$ 就已确定；当 $i+j\geq d$ 时，置 $\F_{i+j}=\SSS_d\otimes_{\SSS_0}\F_{i+j-d}$。当然，为使这些同态在诸 $(\SSS_d)^{\otimes k}\otimes\F_i$（$k\in\Z$，$0\leq i\leq d-1$）的直和上定义 $\SSS$-模结构，须满足相应结合性条件，这里不再写出。

在 $d=1$ 的情形，即（3.3）所考虑的情形，分次 $\SSS$-模的范畴等价于任意 $\SSS_0$-模的范畴；若 $X$ 是预概形且 $\SSS$ 拟凝聚，限制于拟凝聚模后仍然等价。正是在此意义下，本小节的讨论推广了第3节的讨论。此外，在适当有限性条件下，本小节的结果，结合（8.14.11），在一定程度上将预概形上的拟凝聚分次代数，以及这些代数上的分次模“模去(TN)”后的研究，归结于所考虑的代数为周期代数的特例；在此特例中，$\M$ 满足(TN)（3.4.2）便蕴涵 $\M=0$。

\egastatement{8.14.13}{注} 在（8.14.1）的假设下，设 $d>0$。已经定义典范的 $Y$-同构 $h:X\simeq X^{(d)}=\Proj(\SSS^{(d)})$（3.1.8）。对每个分次拟凝聚 $\SSS$-模 $\M$，以及每个 $0\leq k\leq d-1$，还有典范 $h$-同构（采用3.1.1的记号）
\begin{equation}(\sProj(\M))^{(d,k)}\xleftarrow{\sim}\sProj(\M^{(d,k)}).\tag{8.14.13.1}\end{equation}
先设 $Y=\Spec(A)$ 仿射，$\SSS=\wt S$，$\M=\wt M$，其中 $S$ 是非负分次 $A$-代数，$M$ 是分次 $S$-模。对每个 $f\in S_e$（$e>0$），$h$ 将 $D_+(f)$ 映至 $D_+(f^d)$，并对应于典范同构 $S_{(f^d)}\simeq S_{(f)}$（2.2.2）。于是（8.14.13.1）在 $D_+(f^d)$ 上的限制，对应于典范双同构 $M_{f^d}\to M_f$ 在次数模 $d$ 同余于 $k$ 的元素上的限制。读者可以验证这些同构与由 $f$ 换为齐次倍数 $fg$ 相容，并且当 $A'$ 是 $A$-代数时，与由 $S$ 换为分次 $A'$-代数 $S'=S\otimes_A A'$ 也相容。

特别地，有 $h$-同构
\begin{equation}(\SSS^{(d)})_{X^{(d)}}\simeq(\SSS_X)^{(d)},\tag{8.14.13.2}\end{equation}
保持两端的乘法结构。借助它，（8.14.13.1）成为从分次 $(\SSS^{(d)})_{X^{(d)}}$-模到分次 $(\SSS_X)^{(d)}$-模的 $h$-双同构。同样，有 $h$-同构
\begin{equation}\sProj_0(\SSS^{(d,k)}(n))\simeq\OO_X(nd+k),\tag{8.14.13.3}\end{equation}
补充了（3.2.9，(ii)）的结果。

由（8.14.13.1），立即得到分次 $\SSS^{(d)}$-模同构
\begin{equation}\boldsymbol\Gamma_*^{(d)}(\sProj(\M^{(d,k)}))
\simeq\boldsymbol\Gamma_*((\sProj(\M))^{(d,k)}),\tag{8.14.13.4}\end{equation}
其中 $\boldsymbol\Gamma_*^{(d)}$ 对应于结构态射 $q^{(d)}:X^{(d)}\to Y$。立即可验证，典范同态 $\alpha$（8.14.5.4）及 $X^{(d)}$ 的相应同态 $\alpha^{(d)}$ 使图
\begin{equation}\begin{tikzcd}[column sep=small]
&\M^{(d,k)}\arrow[dl,"\alpha^{(d)}"']\arrow[dr,"\alpha"]&\\
\boldsymbol\Gamma_*^{(d)}(\sProj(\M^{(d,k)}))\arrow[rr,"\sim"']&&\boldsymbol\Gamma_*((\sProj(\M))^{(d,k)})
\end{tikzcd}\tag{8.14.13.5}\end{equation}
交换。验证时设 $Y$ 仿射，计算同一元素经 $\alpha^{(d)}$ 与 $\alpha$ 所得像分别在 $D_+(f^d)$ 与 $D_+(f)$ 上的限制即可。

\egastatement{8.14.14}{命题} 设 $Y$ 是拟紧预概形，$\SSS$ 是有限型分次拟凝聚 $\OO_Y$-代数，$\M$ 是满足(TF)的分次拟凝聚 $\SSS$-模，$X=\Proj(\SSS)$。则 $\SSS_X$ 是分次周期 $\OO_X$-代数（8.14.12）；并且存在 $\SSS_X$ 的一个周期 $d$，使各 $(\sProj(\M))^{(d,k)}$（$0\leq k\leq d-1$）都是有限型 $(\SSS_X)^{(d)}$-模。
\begin{proof}
由（3.1.10），存在 $d$，使 $\SSS^{(d)}$ 由 $\SSS_d=(\SSS^{(d)})_1$ 生成，且后者是有限型 $\OO_Y$-模。因此，为证明第一项断言，可由（8.14.13.2）归结于 $d=1$，此时命题由（3.2.7）得到。再利用（8.14.13.1），第二项断言由（2.1.6，(iii)）及（3.4.3）得到。
\end{proof}


% Reconstructed from uploaded EGAII.pdf, PDF pages 202--219 (printed pp.205--222).
\clearpage

% END text/ch2.tex


% BEGIN text/ch3_01_05.tex
\renewcommand{\B}{\mathscr B}
\setcounter{chapter}{2}
\chapter{凝聚层的上同调研究}
\src{81--82}
\begin{center}\textbf{内容提要}\end{center}
\begin{enumerate}[label=\S\,\arabic*.,leftmargin=2.8em]
\item 仿射概形的上同调。
\item 射影态射的上同调研究。
\item 固有态射的有限性定理。
\item 固有态射的基本定理及其应用。
\item 凝聚代数层的一个存在定理。
\item 局部与整体的 Tor 和 Ext 函子；Künneth 公式。
\item 模的协变上同调函子中的基变换研究。
\item 射影丛上的对偶定理。
\item 相对上同调与局部上同调；局部对偶。
\item 射影上同调与局部上同调的关系；沿除子的形式完备化技术。
\item 整体与局部的 Picard 群。\origfootnote{第IV章不依赖于 §§8—11，而且很可能在这些节之前发表。}
\end{enumerate}
\transnote{这是原著对整个第III章所列的内容提要。本次上传的“第一部分”实际收入第III章 §§1—5；提要中 §§6—11 的标题不表示本册包含这些节的正文。}

本章给出凝聚代数层上同调的基本定理，但不包括由留数理论推出的定理（即对偶定理）；后者将在以后的某章中研究。在前一类定理中，主要有六个基本定理，分别构成本章前六节的内容。以后这些结果将成为重要工具，即使所讨论的问题本身并非上同调性质的问题，也同样如此；读者从第4节起就能看到最初的例子。第7节给出一些更具技术性的结果，但在应用中经常需要使用。最后，§§8—11 将展开一些与凝聚层对偶有关、对应用特别重要而又能在一般留数理论之前叙述的结果。

§§1—2 的内容归功于 J.-P. Serre；读者将会看到，我们只需沿用（FAC）的叙述。§§8—9 同样受到（FAC）的启发，不过，由于背景不同而必须作出的转化并不那么显然。最后，正如我们在导言中所说，第4节应看作以现代语言重新表述 Zariski 的“全纯函数理论”中的基本“不变性定理”。

还应指出，3.4 小节的结果以及预备命题（0，13.4—13.7）不会用于第III章的后续部分，初读时可以略去。

\section{仿射概形的上同调}
\subsection{外代数复形的回顾}
\src{82--85}
\egapar{1.1.1}
设 $A$ 为环，$\mathbf f=(f_i)_{1\leq i\leq r}$ 为 $A$ 的 $r$ 个元素组成的系统。与 $\mathbf f$ 对应的\textbf{外代数复形} $K_\bullet(\mathbf f)$ 是如下定义的链复形（G，I，2.2）：其分次 $A$-模为外代数 $\bigwedge(A^r)$，赋予通常的分次；边缘算子为内乘 $i_{\mathbf f}$，这里将 $\mathbf f$ 看成对偶模 $(A^r)^\vee$ 的元素。回顾，$i_{\mathbf f}$ 是 $\bigwedge(A^r)$ 上次数为 $-1$ 的反导子；若 $(e_i)_{1\leq i\leq r}$ 是 $A^r$ 的典范基，则 $i_{\mathbf f}(e_i)=f_i$。条件 $i_{\mathbf f}\circ i_{\mathbf f}=0$ 可立即验证。

一个等价定义如下：对每个 $i$，令链复形 $K_\bullet(f_i)$ 满足 $K_0(f_i)=K_1(f_i)=A$，而 $n\ne0,1$ 时 $K_n(f_i)=0$；其边缘算子 $d_1:A\to A$ 是乘以 $f_i$。令
\[
K_\bullet(\mathbf f)=K_\bullet(f_1)\otimes_A\cdots\otimes_A K_\bullet(f_r)
\]
并赋予总次数（G，I，2.7）。此复形与前一定义的复形同构，验证是直接的。

\egapar{1.1.2}
对每个 $A$-模 $M$，定义链复形
\begin{equation}\tag{1.1.2.1}
K_\bullet(\mathbf f,M)=K_\bullet(\mathbf f)\otimes_AM
\end{equation}
及上链复形（G，I，2.2）
\begin{equation}\tag{1.1.2.2}
K^\bullet(\mathbf f,M)=\Hom_A(K_\bullet(\mathbf f),M).
\end{equation}
若 $g$ 是后一复形中的 $k$-上链，令
\[
g(i_1,\ldots,i_k)=g(e_{i_1}\wedge\cdots\wedge e_{i_k}),
\]
则 $g$ 可识别为从 $[1,r]^k$ 到 $M$ 的交错映射。由上述定义有
\begin{equation}\tag{1.1.2.3}
d^kg(i_1,\ldots,i_{k+1})=
\sum_{h=1}^{k+1}(-1)^{h-1}f_{i_h}g(i_1,\ldots,\widehat{i_h},\ldots,i_{k+1}).
\end{equation}

\egapar{1.1.3}
由这些复形，照例得到同调与上同调 $A$-模（G，I，2.2）
\begin{equation}\tag{1.1.3.1}
\mathrm H_\bullet(\mathbf f,M)=\mathrm H_\bullet(K_\bullet(\mathbf f,M)),
\end{equation}
\begin{equation}\tag{1.1.3.2}
\mathrm H^\bullet(\mathbf f,M)=\mathrm H^\bullet(K^\bullet(\mathbf f,M)).
\end{equation}
还可定义 $A$-同构 $K_\bullet(\mathbf f,M)\simeq K^\bullet(\mathbf f,M)$：将链
\[
z=\sum(e_{i_1}\wedge\cdots\wedge e_{i_k})\otimes m_{i_1\cdots i_k}
\]
对应于上链 $g_z$，令 $g_z(j_1,\ldots,j_{r-k})=\varepsilon m_{i_1\cdots i_k}$；这里严格递增序列 $(j_h)$ 是严格递增序列 $(i_h)$ 在 $[1,r]$ 中的补序列，$\varepsilon=(-1)^\nu$，$\nu$ 是排列 $i_1,\ldots,i_k,j_1,\ldots,j_{r-k}$ 的逆序数。验证可得 $g_{dz}=d(g_z)$，从而有同构
\begin{equation}\tag{1.1.3.3}
\mathrm H^i(\mathbf f,M)\simeq\mathrm H_{r-i}(\mathbf f,M)\qquad(0\leq i\leq r).
\end{equation}
本章主要考察上同调模 $\mathrm H^\bullet(\mathbf f,M)$。

固定 $\mathbf f$ 时，立即可见（G，I，2.1），$M\mapsto\mathrm H^\bullet(\mathbf f,M)$ 是从 $A$-模范畴到分次 $A$-模范畴的上同调函子（T，II，2.1），在次数小于 $0$ 或大于 $r$ 时为零。此外，若 $(\mathbf f)$ 表示 $f_1,\ldots,f_r$ 生成的理想，则
\begin{equation}\tag{1.1.3.4}
\mathrm H^0(\mathbf f,M)=\Hom_A(A/(\mathbf f),M).
\end{equation}
这直接来自（1.1.2.3），且 $\mathrm H^0(\mathbf f,M)$ 显然可识别为 $M$ 中被 $(\mathbf f)$ 零化的子模。同样，由（1.1.2.3）有
\begin{equation}\tag{1.1.3.5}
\mathrm H^r(\mathbf f,M)=M\Big/\sum_{i=1}^rf_iM=(A/(\mathbf f))\otimes_AM.
\end{equation}
以下将使用一个已知结果；为叙述完整，回顾其证明。

\egastatement{1.1.4}{命题}
设 $A$ 为环，$\mathbf f=(f_i)_{1\leq i\leq r}$ 是 $A$ 的有限元素族，$M$ 为 $A$-模。若对 $1\leq i\leq r$，在 $M_{i-1}=M/(f_1M+\cdots+f_{i-1}M)$ 上的数乘 $z\mapsto f_iz$ 是单射，则 $i\ne r$ 时 $\mathrm H^i(\mathbf f,M)=0$。
\eproof 由（1.1.3.3），只需证明 $i>0$ 时 $\mathrm H_i(\mathbf f,M)=0$。对 $r$ 归纳，$r=0$ 显然。令 $\mathbf f'=(f_i)_{1\leq i\leq r-1}$；此族满足命题条件。令 $L_\bullet=K_\bullet(\mathbf f',M)$，则归纳假设给出 $i>0$ 时 $\mathrm H_i(L_\bullet)=0$，而（1.1.3.3）、（1.1.3.5）给出 $\mathrm H_0(L_\bullet)=M_{r-1}$。简记 $K_\bullet=K_\bullet(f_r)=K_0\oplus K_1$，其中 $K_0=K_1=A$，$d_1:K_1\to K_0$ 是乘以 $f_r$。按定义（1.1.1），$K_\bullet(\mathbf f,M)=K_\bullet\otimes_AL_\bullet$。使用以下引理。

\egastatement{1.1.4.1}{引理}
设 $K_\bullet$ 是由自由 $A$-模组成的链复形，除维数 $0,1$ 外各项均为零。对每个 $A$-模链复形 $L_\bullet$ 及每个指标 $p$，有正合列
\[
0\to\mathrm H_0(K_\bullet\otimes_A\mathrm H_p(L_\bullet))
\to\mathrm H_p(K_\bullet\otimes_AL_\bullet)
\to\mathrm H_1(K_\bullet\otimes_A\mathrm H_{p-1}(L_\bullet))\to0.
\]
\begin{proof}
这是 Künneth 谱序列低次项正合列的一个特例（M，XV，5.2，a；G，I，5.5.2），也可直接证明。将 $K_0,K_1$ 看作分别只在维数 $0,1$ 非零的链复形，则有复形正合列
\[
0\to K_0\otimes_AL_\bullet\to K_\bullet\otimes_AL_\bullet\to K_1\otimes_AL_\bullet\to0.
\]
对它应用同调长正合列，得到
\[
\begin{aligned}
\cdots&\to\mathrm H_{p+1}(K_1\otimes_AL_\bullet)\xrightarrow{\partial}\mathrm H_p(K_0\otimes_AL_\bullet)\\
&\to\mathrm H_p(K_\bullet\otimes_AL_\bullet)\to\mathrm H_p(K_1\otimes_AL_\bullet)
\xrightarrow{\partial}\mathrm H_{p-1}(K_0\otimes_AL_\bullet)\to\cdots.
\end{aligned}
\]
显然 $\mathrm H_p(K_0\otimes_AL_\bullet)=K_0\otimes_A\mathrm H_p(L_\bullet)$，且
$\mathrm H_p(K_1\otimes_AL_\bullet)=K_1\otimes_A\mathrm H_{p-1}(L_\bullet)$。又立即可验证，算子 $\partial:K_1\otimes_A\mathrm H_p(L_\bullet)\to K_0\otimes_A\mathrm H_p(L_\bullet)$ 就是 $d_1\otimes1$。引理由此正合列以及 $\mathrm H_0(K_\bullet\otimes_A\mathrm H_p(L_\bullet))$、$\mathrm H_1(K_\bullet\otimes_A\mathrm H_p(L_\bullet))$ 的定义得出。
\end{proof}
引理成立后，（1.1.4）的证明立即完成：归纳假设与引理给出 $p\geq2$ 时 $\mathrm H_p(K_\bullet\otimes_AL_\bullet)=0$。若再证明 $\mathrm H_1(K_\bullet\otimes_A\mathrm H_0(L_\bullet))=0$，引理就给出 $\mathrm H_1(K_\bullet\otimes_AL_\bullet)=0$；但按定义，前者正是 $M_{r-1}$ 上数乘 $z\mapsto f_rz$ 的核，而该核按假设为零。\finishproof

\egapar{1.1.5}
设 $\mathbf g=(g_i)_{1\leq i\leq r}$ 为 $A$ 的另一个 $r$ 元素序列，令 $\mathbf{fg}=(f_ig_i)_{1\leq i\leq r}$。可以定义典范复形同态
\begin{equation}\tag{1.1.5.1}
\varphi_{\mathbf g}:K_\bullet(\mathbf{fg})\to K_\bullet(\mathbf f),
\end{equation}
即将 $A$-线性映射 $(z_1,\ldots,z_r)\mapsto(g_1z_1,\ldots,g_rz_r)$ 典范地延拓到外代数 $\bigwedge(A^r)$。为说明它确为复形同态，只需注意更一般的事实：若 $u:E\to F$ 为 $A$-线性映射，$x\in F^\vee$，$y={}^tu(x)\in E^\vee$，则
\begin{equation}\tag{1.1.5.2}
(\bigwedge u)\circ i_y=i_x\circ(\bigwedge u).
\end{equation}
事实上，两边是 $\bigwedge F$ 的反导子，只需验证它们在 $F$ 上相等，这由定义直接得出。\transnote{原文此处写“两边是 $\bigwedge F$ 的反导子，在 $F$ 上验证”。按所列 $u:E\to F$，两边实际是从 $\bigwedge E$ 到 $\bigwedge F$ 的、相对于 $\bigwedge u$ 的反导子，应在 $E$ 上验证；上述公式本身保留。}
若将 $K_\bullet(\mathbf f)$ 识别为各 $K_\bullet(f_i)$ 的张量积（1.1.1），则 $\varphi_{\mathbf g}$ 是各 $\varphi_{g_i}$ 的张量积；其中 $\varphi_{g_i}$ 在零次为恒等，在一次为乘以 $g_i$。

\egapar{1.1.6}
特别地，对 $0\leq n\leq m$，有复形同态
\begin{equation}\tag{1.1.6.1}
\varphi_{\mathbf f^{m-n}}:K_\bullet(\mathbf f^m)\to K_\bullet(\mathbf f^n),
\end{equation}
从而有同态
\begin{equation}\tag{1.1.6.2}
\varphi_{m-n}:K^\bullet(\mathbf f^n,M)\to K^\bullet(\mathbf f^m,M),
\end{equation}
\begin{equation}\tag{1.1.6.3}
\varphi_{m-n}:\mathrm H^\bullet(\mathbf f^n,M)\to\mathrm H^\bullet(\mathbf f^m,M).
\end{equation}
对 $p\leq n\leq m$，后两者显然满足 $\varphi_{m-p}=\varphi_{m-n}\circ\varphi_{n-p}$，因而定义两个 $A$-模的归纳系统。令
\begin{equation}\tag{1.1.6.4}
C^\bullet((\mathbf f),M)=\ilim_nK^\bullet(\mathbf f^n,M),
\end{equation}
\begin{equation}\tag{1.1.6.5}
\mathrm H^\bullet((\mathbf f),M)=\mathrm H^\bullet(C^\bullet((\mathbf f),M))
=\ilim_n\mathrm H^\bullet(\mathbf f^n,M).
\end{equation}
末一等式来自归纳极限与 $\mathrm H^\bullet$ 可交换（G，I，2.1）。以后将看到（1.4.3），$\mathrm H^\bullet((\mathbf f),M)$ 实际只依赖于理想 $(\mathbf f)$，甚至只依赖于 $A$ 上的 $(\mathbf f)$-预进拓扑，因此上述记号是恰当的。显然 $M\mapsto C^\bullet((\mathbf f),M)$ 是 $A$-线性正合函子，而 $M\mapsto\mathrm H^\bullet((\mathbf f),M)$ 是上同调函子。

\egapar{1.1.7}
设 $\mathbf f,\mathbf g\in A^r$；以 $e_{\mathbf g}$ 表示外代数 $\bigwedge(A^r)$ 中以向量 $\mathbf g$ 左乘的算子。有\textbf{同伦公式}
\begin{equation}\tag{1.1.7.1}
i_{\mathbf f}e_{\mathbf g}+e_{\mathbf g}i_{\mathbf f}=(\mathbf g,\mathbf f)1
\end{equation}
在 $A$-模 $A^r$ 上成立（$1$ 表示 $A^r$ 的恒等自同构）；此关系也表示在复形 $K_\bullet(\mathbf f)$ 中有
\begin{equation}\tag{1.1.7.2}
de_{\mathbf g}+e_{\mathbf g}d=(\mathbf g,\mathbf f)1.
\end{equation}
若 $(\mathbf f)=A$，则存在 $\mathbf g\in A^r$ 使 $(\mathbf g,\mathbf f)=\sum_{i=1}^rg_if_i=1$。因此（G，I，2.4）：

\egastatement{1.1.8}{命题}
若 $f_i$ 生成的理想 $(\mathbf f)$ 等于 $A$，则复形 $K_\bullet(\mathbf f)$ 同伦平凡；因而对每个 $A$-模 $M$，$K_\bullet(\mathbf f,M)$、$K^\bullet(\mathbf f,M)$ 也都同伦平凡。

\egastatement{1.1.9}{推论}
若 $(\mathbf f)=A$，则对每个 $A$-模 $M$，有 $\mathrm H^\bullet(\mathbf f,M)=0$ 及 $\mathrm H^\bullet((\mathbf f),M)=0$。
\begin{proof}此时对每个 $n$ 都有 $(\mathbf f^n)=A$。\end{proof}

\egastatement{1.1.10}{注}
采用上述记号，设 $X=\Spec A$，$Y$ 为理想 $(\mathbf f)$ 定义的闭子概形。将在第9节证明，$\mathrm H^\bullet((\mathbf f),M)$ 同构于相对于 $Y$ 的闭子集所成反滤子的上同调 $\mathrm H_Y^\bullet(X,\wt M)$（T，3.2）。还将证明，命题（1.2.3）用于 $X$ 和 $\F=\wt M$ 时，是上同调正合列
\[
\cdots\to\mathrm H_Y^p(X,\F)\to\mathrm H^p(X,\F)\to\mathrm H^p(X-Y,\F)
\to\mathrm H_Y^{p+1}(X,\F)\to\cdots
\]
的一个特例。

\subsection{开覆盖的 Čech 上同调}
\src{85--88}
\egapar{1.2.1}\textbf{记号。}
本小节以 $X$ 表示预概形，$\F$ 表示拟凝聚 $\OO_X$-模，令 $A=\Gamma(X,\OO_X)$、$M=\Gamma(X,\F)$；以 $\mathbf f=(f_i)_{1\leq i\leq r}$ 表示 $A$ 中的有限元素系统，令 $U_i=X_{f_i}$ 为满足 $f_i(x)\ne0$ 的点组成的开集（$\OI$，5.5.2），$U=\bigcup_iU_i$，$\mathfrak U=(U_i)_{1\leq i\leq r}$ 为 $U$ 的覆盖。

\egapar{1.2.2}
假设 $X$ 的底空间诺特，或者 $X$ 是底空间拟紧的概形。于是已知（I，9.3.3），$\Gamma(U_i,\F)=M_{f_i}$。令
\[
U_{i_0\cdots i_p}=\bigcap_{k=0}^pU_{i_k}=X_{f_{i_0}\cdots f_{i_p}}
\]
（$\OI$，5.5.3），故也有
\begin{equation}\tag{1.2.2.1}
\Gamma(U_{i_0\cdots i_p},\F)=M_{f_{i_0}\cdots f_{i_p}}.
\end{equation}
根据（$\OI$，1.6.1），右边识别为系统 $M^{(n)}_{i_0\cdots i_p}=M$ 的归纳极限；对 $m\leq n$，过渡同态 $\varphi_{nm}$ 为乘以 $(f_{i_0}\cdots f_{i_p})^{n-m}$。对每个 $n$，以 $C_n^p(M)$ 表示从 $[1,r]^{p+1}$ 到 $M$ 的交错映射所成 $A$-模。这些模也相对于 $\varphi_{nm}$ 构成归纳系统。若 $C^p(\mathfrak U,\F)$ 是覆盖 $\mathfrak U$ 的、系数在 $\F$ 中的交错 Čech $p$-上链群（G，II，5.1），则
\begin{equation}\tag{1.2.2.2}
C^p(\mathfrak U,\F)=\ilim_n C_n^p(M).
\end{equation}
采用（1.1.2）的记号，$C_n^p(M)$ 识别为 $K^{p+1}(\mathbf f^n,M)$，$\varphi_{nm}$ 识别为（1.1.6）定义的 $\varphi_{n-m}$。故对每个 $p\geq0$ 有关于 $\F$ 的典范函子同构
\begin{equation}\tag{1.2.2.3}
C^p(\mathfrak U,\F)\simeq C^{p+1}((\mathbf f),M).
\end{equation}
公式（1.1.2.3）及覆盖上同调的定义（G，II，5.1）说明这些同构与上边缘算子相容。

\egastatement{1.2.3}{命题}
若 $X$ 是底空间诺特的预概形，或底空间拟紧的概形，则存在关于 $\F$ 的典范函子同构
\begin{equation}\tag{1.2.3.1}
\mathrm H^p(\mathfrak U,\F)\simeq\mathrm H^{p+1}((\mathbf f),M)\quad(p\geq1),
\end{equation}
以及关于 $\F$ 的函子正合列
\begin{equation}\tag{1.2.3.2}
0\to\mathrm H^0((\mathbf f),M)\to M\to\mathrm H^0(\mathfrak U,\F)
\to\mathrm H^1((\mathbf f),M)\to0.
\end{equation}
\begin{proof}
同构由（1.2.2）立即得出。又 $C^0(\mathfrak U,\F)=C^1((\mathbf f),M)$，故 $\mathrm H^0(\mathfrak U,\F)$ 识别为 $C^1((\mathbf f),M)$ 的1-上闭链子群；因 $M=C^0((\mathbf f),M)$，上述正合列就是 $\mathrm H^0((\mathbf f),M)$、$\mathrm H^1((\mathbf f),M)$ 的定义给出的正合列。
\end{proof}

\egastatement{1.2.4}{推论}
假设各 $X_{f_i}$ 拟紧，且存在 $g_i\in\Gamma(U,\OO_X)$ 使 $\sum_ig_i(f_i|U)=1|U$。\transnote{底本将 $g_i$ 写在 $\Gamma(U,\F)$ 中；但随后的单位分解是在函数环中进行，这里将其记为 $\Gamma(U,\OO_X)$，以使数乘和 $1|U$ 的类型一致。}则对每个拟凝聚 $(\OO_X|U)$-模 $\F$，有 $p>0$ 时 $\mathrm H^p(\mathfrak U,\F)=0$；若再有 $U=X$，则（1.2.3.2）的典范同态 $M\to\mathrm H^0(\mathfrak U,\F)$ 为双射。
\begin{proof}
各 $U_i=X_{f_i}$ 拟紧，故 $U$ 也拟紧；因而只须考虑 $U=X$。假设于是给出各 $p\geq0$ 时 $\mathrm H^p((\mathbf f),M)=0$（1.1.9），结论由（1.2.3.1）、（1.2.3.2）立即得出。
\end{proof}
注意，由于 $\mathrm H^0(\mathfrak U,\F)=\mathrm H^0(U,\F)$（G，II，5.2.2），这也将（I，1.3.7）作为特例重新证明了。

\egastatement{1.2.5}{注}
假设 $X$ 为仿射概形；则 $U_i=X_{f_i}=D(f_i)$ 及各 $U_{i_0\cdots i_p}$ 都仿射，但 $U$ 未必仿射。此时 $\Gamma(X,\F)$ 和 $\Gamma(U_{i_0\cdots i_p},\F)$ 都是关于 $\F$ 的正合函子（I，1.3.11）。若 $0\to\F'\to\F\to\F''\to0$ 是拟凝聚 $\OO_X$-模正合列，则
\[
0\to C^\bullet(\mathfrak U,\F')\to C^\bullet(\mathfrak U,\F)\to C^\bullet(\mathfrak U,\F'')\to0
\]
是复形正合列，并给出上同调长正合列
\[
\cdots\to\mathrm H^p(\mathfrak U,\F')\to\mathrm H^p(\mathfrak U,\F)
\to\mathrm H^p(\mathfrak U,\F'')\to\mathrm H^{p+1}(\mathfrak U,\F')\to\cdots.
\]
令 $M'=\Gamma(X,\F')$、$M''=\Gamma(X,\F'')$，则 $0\to M'\to M\to M''\to0$ 也正合。由于 $C^\bullet((\mathbf f),M)$ 关于 $M$ 正合，也有长正合列
\[
\cdots\to\mathrm H^p((\mathbf f),M')\to\mathrm H^p((\mathbf f),M)
\to\mathrm H^p((\mathbf f),M'')\to\mathrm H^{p+1}((\mathbf f),M')\to\cdots.
\]
由于图
\[
\begin{tikzcd}[column sep=small]
0\arrow[r]&C^\bullet(\mathfrak U,\F')\arrow[r]\arrow[d]&C^\bullet(\mathfrak U,\F)\arrow[r]\arrow[d]&C^\bullet(\mathfrak U,\F'')\arrow[r]\arrow[d]&0\\
0\arrow[r]&C^\bullet((\mathbf f),M')\arrow[r]&C^\bullet((\mathbf f),M)\arrow[r]&C^\bullet((\mathbf f),M'')\arrow[r]&0
\end{tikzcd}
\]
交换，故对每个 $p$，下图交换（G，I，2.1.1）：
\begin{equation}\tag{1.2.5.1}
\begin{tikzcd}
\mathrm H^p(\mathfrak U,\F'')\arrow[r,"\delta_p"]\arrow[d]&\mathrm H^{p+1}(\mathfrak U,\F')\arrow[d]\\
\mathrm H^{p+1}((\mathbf f),M'')\arrow[r]&\mathrm H^{p+2}((\mathbf f),M').
\end{tikzcd}
\end{equation}

\subsection{仿射概形的上同调}
\src{88--89}
\egastatement{1.3.1}{定理}
设 $X$ 是仿射概形。对每个拟凝聚 $\OO_X$-模 $\F$，以及每个 $p>0$，有 $\mathrm H^p(X,\F)=0$。
\begin{proof}
设 $\mathfrak U$ 为由仿射开集 $X_{f_i}=D(f_i)$（$1\leq i\leq r$）组成的有限覆盖；已知此时 $f_i$ 在 $A=\Gamma(X,\OO_X)$ 中生成单位理想。由（1.2.4），$p>0$ 时 $\mathrm H^p(\mathfrak U,\F)=0$。这样的仿射开集有限覆盖可以任意加细（I，1.1.10），故 Čech 上同调的定义（G，II，5.8）给出 $p>0$ 时 $\check{\mathrm H}^p(X,\F)=0$。这也适用于每个 $f\in A$ 所给出的预概形 $X_f$（I，1.3.6），所以 $p>0$ 时 $\check{\mathrm H}^p(X_f,\F)=0$。由于 $X_f\cap X_g=X_{fg}$，根据（G，II，5.9.2），可知对每个 $p>0$ 也有 $\mathrm H^p(X,\F)=0$。
\end{proof}

\egastatement{1.3.2}{推论}
设 $Y$ 为预概形，$f:X\to Y$ 为仿射态射（II，1.6.1）。对每个拟凝聚 $\OO_X$-模 $\F$，有 $q>0$ 时 $\mathrm R^qf_*(\F)=0$。
\begin{proof}
按定义，$\mathrm R^qf_*(\F)$ 与预层 $U\mapsto\mathrm H^q(f^{-1}(U),\F)$ 相伴。仿射开集构成 $Y$ 的一组基；对这样的 $U$，$f^{-1}(U)$ 仿射（II，1.3.2），故由（1.3.1），$q>0$ 时相应上同调为零。
\end{proof}

\egastatement{1.3.3}{推论}
设 $Y$ 为预概形，$f:X\to Y$ 为仿射态射。对每个拟凝聚 $\OO_X$-模 $\F$，典范同态 $\mathrm H^p(Y,f_*(\F))\to\mathrm H^p(X,\F)$（0，12.1.3.1）对每个 $p$ 都是双射。
\begin{proof}
由（0，12.1.7），只需证明复合函子 $\Gamma f_*$ 的第二谱序列的边缘同态 ${}''E_2^{p0}=\mathrm H^p(Y,f_*(\F))\to\mathrm H^p(X,\F)$ 是双射。此谱序列的 $E_2$ 项为 ${}''E_2^{pq}=\mathrm H^p(Y,\mathrm R^qf_*(\F))$（G，II，4.17.1）；由（1.3.2），$q>0$ 时该项为零，故谱序列退化，所求结论由（0，11.1.6）得出。
\end{proof}

\egastatement{1.3.4}{推论}
设 $f:X\to Y$ 为仿射态射，$g:Y\to Z$ 为态射。对每个拟凝聚 $\OO_X$-模 $\F$，典范同态 $\mathrm R^pg_*(f_*(\F))\to\mathrm R^p(g\circ f)_*(\F)$（0，12.2.5.1）对每个 $p$ 都是双射。
\begin{proof}
由（1.3.3），对 $Z$ 的每个仿射开集 $W$，典范同态
\[
\mathrm H^p(g^{-1}(W),f_*(\F))\to\mathrm H^p(f^{-1}(g^{-1}(W)),\F)
\]
都是双射。这说明定义所求典范同态的预层同态为双射（0，12.2.5）。
\end{proof}\par

% ---- END c3_s01a.tex ----
% ---- BEGIN c3_s01b.tex ----
\subsection{应用于任意预概形的上同调}
\src{89--95}
\egastatement{1.4.1}{命题}
设 $X$ 为概形，$\mathfrak U=(U_\alpha)$ 为 $X$ 的仿射开覆盖。对每个拟凝聚 $\OO_X$-模 $\F$，上同调 $\Gamma(X,\OO_X)$-模 $\mathrm H^\bullet(X,\F)$ 与 $\mathrm H^\bullet(\mathfrak U,\F)$ 典范同构。
\begin{proof}
由于 $X$ 为概形，覆盖中的任意有限交 $V$ 都仿射（I，5.5.6），所以 $q\geq1$ 时 $\mathrm H^q(V,\F)=0$（1.3.1）。命题由 Leray 定理（G，II，5.4.1）得出。
\end{proof}

\egastatement{1.4.2}{注}
只要各 $U_\alpha$ 的有限交仿射，（1.4.1）的结论仍然成立，即使不假设 $X$ 为概形。

\egastatement{1.4.3}{推论}
设 $X$ 为底空间拟紧的概形，$A=\Gamma(X,\OO_X)$，$\mathbf f=(f_i)_{1\leq i\leq r}$ 为 $A$ 中的有限元素序列，使各 $X_{f_i}$ 仿射。采用（1.2.1）的记号，对每个拟凝聚 $\OO_X$-模 $\F$，有典范函子同构
\begin{equation}\tag{1.4.3.1}
\mathrm H^q(U,\F)\simeq\mathrm H^{q+1}((\mathbf f),M)\quad(q\geq1),
\end{equation}
及函子正合列
\begin{equation}\tag{1.4.3.2}
0\to\mathrm H^0((\mathbf f),M)\to M\to\mathrm H^0(U,\F)
\to\mathrm H^1((\mathbf f),M)\to0.
\end{equation}
\begin{proof}由（1.4.1）和（1.2.3）立即得出。\end{proof}

\egapar{1.4.4}
若 $X$ 为仿射概形，则由（1.2.5）、（1.4.1），对每个 $q\geq0$，与拟凝聚 $\OO_X$-模正合列 $0\to\F'\to\F\to\F''\to0$ 对应的图
\begin{equation}\tag{1.4.4.1}
\begin{tikzcd}
\mathrm H^q(U,\F'')\arrow[r,"\delta_q"]\arrow[d]&\mathrm H^{q+1}(U,\F')\arrow[d]\\
\mathrm H^{q+1}((\mathbf f),M'')\arrow[r]&\mathrm H^{q+2}((\mathbf f),M')
\end{tikzcd}
\end{equation}
交换，记号如（1.2.5）。

\egastatement{1.4.5}{命题}
设 $X$ 为拟紧概形，$\LL$ 为可逆 $\OO_X$-模，考虑分次环 $A_*=\Gamma_*(\LL)$（$\OI$，5.4.6）。则
\[
\mathrm H^\bullet(\F,\LL)=\bigoplus_n\mathrm H^\bullet(X,\F\otimes_{\OO_X}\LL^{\otimes n})
\]
是分次 $A_*$-模；对每个 $f\in A_d$，有 $(A_*)_{(f)}$-模典范同构
\begin{equation}\tag{1.4.5.1}
\mathrm H^\bullet(X_f,\F)\simeq(\mathrm H^\bullet(\F,\LL))_{(f)}.
\end{equation}
\begin{proof}
由于 $X$ 是拟紧概形，可用同一个有限仿射开覆盖 $\mathfrak U=(U_i)$ 计算各 $\F\otimes\LL^{\otimes n}$ 的上同调，并使 $\LL|U_i\simeq\OO_X|U_i$（1.4.1）。各 $U_i\cap X_f$ 显然也是仿射开集（I，1.3.6），因此也可用 $\mathfrak U|X_f=(U_i\cap X_f)$ 计算 $X_f$ 上的这些上同调。对每个 $f\in A_d$，乘以 $f$ 给出复形同态
\[
C^\bullet(\mathfrak U,\F\otimes\LL^{\otimes n})\to
C^\bullet(\mathfrak U,\F\otimes\LL^{\otimes(n+d)}),
\]
从而给出上同调同态，这证明第一断言。另一方面，对固定 $f\in A_d$，（I，9.3.2）及（I，1.3.9，ii）给出 $(A_*)_{(f)}$-模复形的同构
\[
C^\bullet(\mathfrak U|X_f,\F)
\simeq\left(\bigoplus_nC^\bullet(\mathfrak U,\F\otimes\LL^{\otimes n})\right)_{(f)}.
\]
取两边的上同调，并利用分次 $A_*$-模范畴中 $M\mapsto M_{(f)}$ 的正合性，就得到（1.4.5.1）。
\end{proof}

\egastatement{1.4.6}{推论}
在（1.4.5）的假设下，再设 $\LL=\OO_X$。令 $A=\Gamma(X,\OO_X)$，则对每个 $f\in A$，有 $A_f$-模典范同构 $\mathrm H^\bullet(X_f,\F)\simeq\mathrm H^\bullet(X,\F)_f$。

\egastatement{1.4.7}{推论}
设 $X$ 为拟紧概形，$f\in\Gamma(X,\OO_X)$。
\begin{enumerate}[label=(\roman*)]
\item 若 $X_f$ 仿射，则对每个拟凝聚 $\OO_X$-模 $\F$、每个 $i>0$ 及每个 $\xi\in\mathrm H^i(X,\F)$，存在整数 $n>0$ 使 $f^n\xi=0$。
\item 反过来，假设 $X_f$ 拟紧，而且对 $\OO_X$ 的每个拟凝聚理想层 $\J$ 和每个 $\xi\in\mathrm H^1(X,\J)$，存在 $n>0$ 使 $f^n\xi=0$。则 $X_f$ 仿射。
\end{enumerate}
\begin{proof}
（i）若 $X_f$ 仿射，则由（1.3.1），$i>0$ 时 $\mathrm H^i(X_f,\F)=0$；故断言直接来自（1.4.6）。

（ii）由 Serre 判据（II，5.2.1），只需对 $\OO_X|X_f$ 的每个拟凝聚理想 $\K$ 证明 $\mathrm H^1(X_f,\K)=0$。由于 $X_f$ 是拟紧概形 $X$ 的拟紧开集，存在 $\OO_X$ 的拟凝聚理想 $\J$，使 $\K=\J|X_f$（I，9.4.2）。由（1.4.6），$\mathrm H^1(X_f,\K)=\mathrm H^1(X,\J)_f$；假设说明右边为零，结论得证。
\end{proof}

\egastatement{1.4.8}{注}
（1.4.7，i）给出了关系（II，4.5.13.2）的一个更简单证明。

\egastatement{1.4.9}{引理}
设 $X$ 为拟紧概形，$\mathfrak U=(U_i)_{1\leq i\leq n}$ 为有限仿射开覆盖，$\F$ 为拟凝聚 $\OO_X$-模。由覆盖定义的层复形 $\C^\bullet(\mathfrak U,\F)$（G，II，5.2）逐项为拟凝聚 $\OO_X$-模。
\begin{proof}
按定义，$\C^p(\mathfrak U,\F)$ 是各 $\F|U_{i_0\cdots i_p}$ 经典范嵌入 $U_{i_0\cdots i_p}\to X$ 的直接像的直和。由于 $X$ 为概形，这些嵌入都是仿射态射（I，5.5.6），故各 $\C^p(\mathfrak U,\F)$ 拟凝聚（II，1.2.6）。
\end{proof}

\egastatement{1.4.10}{命题}
设 $u:X\to Y$ 为分离拟紧态射。对每个拟凝聚 $\OO_X$-模 $\F$，各 $\mathrm R^qu_*(\F)$ 都是拟凝聚 $\OO_Y$-模。
\begin{proof}
问题在 $Y$ 上局部，可设 $Y$ 仿射。于是 $X$ 是有限个仿射开集 $U_i$ 的并；令 $\mathfrak U=(U_i)$。又因 $Y$ 为概形，对每个仿射开集 $V\subset Y$，典范嵌入 $u^{-1}(V)\to X$ 为仿射态射（I，5.5.10）。由（1.4.1）及（G，II，5.2），有典范同构
\begin{equation}\tag{1.4.10.1}
\mathrm H^\bullet(u^{-1}(V),\F)\simeq\mathrm H^\bullet(\Gamma(V,\K^\bullet)),
\end{equation}
其中 $\K^\bullet=u_*(\C^\bullet(\mathfrak U,\F))$。由（1.4.9）及（I，9.2.2），$\K^\bullet$ 逐项为拟凝聚 $\OO_Y$-模；它也是层复形。由上同调层的定义（G，II，4.1）及（I，4.1.1），$\HH^\bullet(\K^\bullet)$ 逐项拟凝聚。由于 $V$ 仿射时 $\Gamma(V,-)$ 在拟凝聚模范畴中正合，所以（G，II，4.1）
\begin{equation}\tag{1.4.10.2}
\mathrm H^\bullet(\Gamma(V,\K^\bullet))=\Gamma(V,\HH^\bullet(\K^\bullet)).
\end{equation}
最后，由（G，II，5.2）中典范同态 $\mathrm H^\bullet(\mathfrak U,\F)\to\mathrm H^\bullet(X,\F)$ 的定义，对另一仿射开集 $V'\subset V$，有交换图
\[
\begin{tikzcd}
\mathrm H^\bullet(u^{-1}(V),\F)\arrow[r,"\sim"]\arrow[d]&\mathrm H^\bullet(\Gamma(V,\K^\bullet))\arrow[d]\\
\mathrm H^\bullet(u^{-1}(V'),\F)\arrow[r,"\sim"]&\mathrm H^\bullet(\Gamma(V',\K^\bullet)).
\end{tikzcd}
\]
因此（1.4.10.1）的同构定义 $\OO_Y$-模同构
\begin{equation}\tag{1.4.10.3}
\mathrm R^\bullet u_*(\F)\simeq\HH^\bullet(\K^\bullet),
\end{equation}
从而 $\mathrm R^\bullet u_*(\F)$ 拟凝聚。
\end{proof}
由（1.4.10.3）、（1.4.10.2）、（1.4.10.1），又得：

\egastatement{1.4.11}{推论}
在（1.4.10）的假设下，对 $Y$ 的每个仿射开集 $V$，典范同态
\begin{equation}\tag{1.4.11.1}
\mathrm H^q(u^{-1}(V),\F)\to\Gamma(V,\mathrm R^qu_*(\F))
\end{equation}
对每个 $q\geq0$ 都是同构。

\egastatement{1.4.12}{推论}
在（1.4.10）的假设下，再设 $Y$ 拟紧。则存在整数 $r>0$，使对每个拟凝聚 $\OO_X$-模 $\F$ 及每个 $q>r$，有 $\mathrm R^qu_*(\F)=0$。若 $Y$ 仿射，可取 $r$ 为某个仿射开覆盖 $X$ 所需的开集数。
\begin{proof}
可用有限个仿射开集覆盖 $Y$，故由（1.4.11），归结到第二断言。若 $\mathfrak U$ 是由 $r$ 个仿射开集组成的覆盖，则交错性给出 $q>r$ 时 $C^q(\mathfrak U,\F)=0$，从而 $\mathrm H^q(\mathfrak U,\F)=0$；再用（1.4.1）即可。
\end{proof}

\egastatement{1.4.13}{推论}
在（1.4.10）的假设下，再设 $Y=\Spec A$ 仿射。则对每个拟凝聚 $\OO_X$-模 $\F$ 及每个 $f\in A$，有典范同构
\[
\Gamma(Y_f,\mathrm R^qu_*(\F))\simeq\Gamma(Y,\mathrm R^qu_*(\F))_f.
\]
\begin{proof}来自 $\mathrm R^qu_*(\F)$ 的拟凝聚性（I，1.3.7）。\end{proof}

\egastatement{1.4.14}{命题}
设 $f:X\to Y$ 为分离拟紧态射，$g:Y\to Z$ 为仿射态射。对每个拟凝聚 $\OO_X$-模 $\F$，典范同态 $\mathrm R^p(g\circ f)_*(\F)\to g_*(\mathrm R^pf_*(\F))$（0，12.2.5.2）对每个 $p$ 都是双射。
\begin{proof}
对 $Z$ 的仿射开集 $W$，$g^{-1}(W)$ 是 $Y$ 的仿射开集。由（1.4.11），定义所求典范同态的预层同态（0，12.2.5）为双射。
\end{proof}

\egastatement{1.4.15}{命题}
设 $u:X\to Y$ 为分离有限型态射，$v:Y'\to Y$ 为平坦预概形态射（$\OI$，6.7.1）；令 $u'=u_{(Y')}$，使有交换图
\begin{equation}\tag{1.4.15.1}
\begin{tikzcd}
X&X'=X_{(Y')}\arrow[l,"v'"']\arrow[d,"u'"]\\
Y\arrow[u,leftarrow,"u"]&Y'\arrow[l,"v"]
\end{tikzcd}
\end{equation}
则对每个拟凝聚 $\OO_X$-模 $\F$，令 $\F'=v'^*(\F)=\F\otimes_{\OO_Y}\OO_{Y'}$，有对各 $q\geq0$ 的典范同构
\[
\mathrm R^qu'_*(\F')\simeq\mathrm R^qu_*(\F)\otimes_{\OO_Y}\OO_{Y'}=v^*(\mathrm R^qu_*(\F)).
\]
\eproof 典范同态 $\rho:\F\to v'_*(v'^*\F)$（$\OI$，4.4.3.2）以函子方式给出同态
\begin{equation}\tag{1.4.15.2}
\mathrm R^qu_*(\F)\to\mathrm R^qu_*(v'_*(\F')).
\end{equation}
令 $w=u\circ v'=v\circ u'$。还有典范同态（0，12.2.5.1及12.2.5.2）
\begin{equation}\tag{1.4.15.3}
\mathrm R^qu_*(v'_*(\F'))\to\mathrm R^qw_*(\F')\to v_*(\mathrm R^qu'_*(\F')).
\end{equation}
复合得到 $\psi:\mathrm R^qu_*(\F)\to v_*(\mathrm R^qu'_*(\F'))$，从而得到典范同态
\begin{equation}\tag{1.4.15.4}
\psi^\flat:v^*(\mathrm R^qu_*(\F))\to\mathrm R^qu'_*(\F').
\end{equation}
其定义不要求 $v$ 满足任何附加假设；现在要证明 $v$ 平坦时它为同构。问题在 $Y,Y'$ 上局部，故可设 $Y=\Spec A$、$Y'=\Spec B$。将使用以下引理。

\egastatement{1.4.15.5}{引理}
设 $\varphi:A\to B$ 为环同态，$Y=\Spec A$、$X=\Spec B$，$f:X\to Y$ 为相应态射，$M$ 为 $B$-模。则 $\wt M$ 为 $f$-平坦层，当且仅当 $M$ 为平坦 $A$-模；特别地，$f$ 平坦，当且仅当 $B$ 为平坦 $A$-模。
\begin{proof}由（$\OI$，6.7.1）的定义、（$\OI$，6.3.3），结合（I，1.3.4）得出。\end{proof}
由（1.4.11.1）及（1.4.15.3）中同态的定义，$\psi$ 对应于复合
\[
\begin{aligned}
\mathrm H^q(X,\F)&\xrightarrow{\rho_q}\mathrm H^q(X,v'_*(v'^*\F))\\
&\xrightarrow{\theta_q}\mathrm H^q(X',v'^*v'_*(v'^*\F))
\xrightarrow{\sigma_q}\mathrm H^q(X',v'^*\F),
\end{aligned}
\]
其中 $\rho_q,\sigma_q$ 来自典范同态 $\rho$ 及 $\sigma:v'^*(v'_*(\F'))\to\F'$，$\theta_q$ 是（0，12.1.3.1）相对于 $v'_*(v'^*\F)$ 的 $\varphi$-态射。由 $\theta_q$ 的函子性，有交换图
\[
\begin{tikzcd}[column sep=small]
\mathrm H^q(X,\F)\arrow[r,"\rho_q"]\arrow[d,"\theta_q"']&\mathrm H^q(X,v'_*(v'^*\F))\arrow[d,"\theta_q"]\\
\mathrm H^q(X',v'^*\F)\arrow[r,"(v'^*\rho)_q"']&\mathrm H^q(X',v'^*v'_*(v'^*\F)).
\end{tikzcd}
\]
按定义（$\OI$，4.4.3），$v'^*(\rho)$ 是 $\sigma$ 的逆，故上述复合最终就是 $\theta_q$；于是 $\psi^\flat$ 是与它相伴的 $B$-同态
$\mathrm H^q(X,\F)\otimes_AB\to\mathrm H^q(X',\F')$。
\transnote{原文将 $v'^*(\rho)$ 称为 $\sigma$ 的“逆”。伴随关系在这里直接给出的是 $\sigma\circ v'^*(\rho)=1$，这一个方向的恒等式已足以完成本段推导；这里标明原文表述的强弱差别。}

由于 $u$ 为有限型，$X$ 可由有限个仿射开集 $U_i$ 覆盖；记覆盖为 $\mathfrak U$。$v$ 仿射，故 $v'$ 也仿射（II，1.6.2，iii），因而 $U'_i=v'^{-1}(U_i)$ 给出 $X'$ 的仿射开覆盖 $\mathfrak U'$。由（0，12.1.4.2），图
\[
\begin{tikzcd}
\mathrm H^q(\mathfrak U,\F)\arrow[r,"\theta'_q"]\arrow[d]&\mathrm H^q(\mathfrak U',\F')\arrow[d]\\
\mathrm H^q(X,\F)\arrow[r,"\theta_q"']&\mathrm H^q(X',\F')
\end{tikzcd}
\]
交换；因 $X,X'$ 为概形，竖直箭头为同构（1.4.1）。故只需证明与上方 $\varphi$-态射相伴的 $B$-同态
\[
\mathrm H^q(\mathfrak U,\F)\otimes_AB\to\mathrm H^q(\mathfrak U',\F')
\]
是同构。对指标序列 $s=(i_k)_{0\leq k\leq p}$，令 $U_s=\bigcap_kU_{i_k}$、$U'_s=\bigcap_kU'_{i_k}=v'^{-1}(U_s)$，$M_s=\Gamma(U_s,\F)$、$M'_s=\Gamma(U'_s,\F')$。典范映射 $M_s\otimes_AB\to M'_s$ 为同构（I，1.6.5），故
$C^\bullet(\mathfrak U,\F)\otimes_AB\to C^\bullet(\mathfrak U',\F')$ 为同构，并将 $d\otimes1$ 识别为后者的上边缘算子。因 $B$ 为平坦 $A$-模，取上同调就得到所求同构（$\OI$，6.1.1）。此结果以后将在第6节推广。\finishproof

\egastatement{1.4.16}{推论}
设 $A$ 为环，$X$ 为有限型 $A$-概形，$B$ 为忠实平坦的 $A$-代数。则 $X$ 仿射，当且仅当 $X\otimes_AB$ 仿射。
\begin{proof}
必要性显然（I，3.2.2）。为证充分性，由于 $X$ 在 $A$ 上分离，且 $\Spec B\to\Spec A$ 平坦，故对每个 $i\geq0$ 及拟凝聚 $\OO_X$-模 $\F$，有
\begin{equation}\tag{1.4.16.1}
\mathrm H^i(X\otimes_AB,\F\otimes_AB)=\mathrm H^i(X,\F)\otimes_AB.
\end{equation}
若 $X\otimes_AB$ 仿射，$i=1$ 时左边为零；忠实平坦性说明 $\mathrm H^1(X,\F)=0$。因 $X$ 为拟紧概形，应用 Serre 判据（II，5.2.1）即可。\transnote{原文在本证明中引为（1.4.1）；这里使用的基变换公式是前一命题（1.4.15）。}
\end{proof}

\egastatement{1.4.17}{命题}
设 $X$ 为预概形，$0\to\F\xrightarrow{u}\G\xrightarrow{v}\K\to0$ 为 $\OO_X$-模正合列。若 $\F,\K$ 拟凝聚，则 $\G$ 也拟凝聚。
\begin{proof}
问题在 $X$ 上局部，可设 $X=\Spec A$；只需验证（I，1.4.1）的条件 d1）、d2），其中 $V=X$。先验证消失条件：若 $t\in\Gamma(X,\G)$ 在 $D(f)$ 上限制为零，则 $v(t)$ 亦如此；存在 $m>0$ 使 $f^mv(t)=v(f^mt)=0$。因 $\Gamma$ 左正合，$f^mt=u(s)$，其中 $s\in\Gamma(X,\F)$。$u$ 单射，所以 $s$ 在 $D(f)$ 上为零；再由（I，1.4.1），有 $n>0$ 使 $f^ns=0$，故 $f^{m+n}t=u(f^ns)=0$。

再验证延拓条件。设 $t'\in\Gamma(D(f),\G)$。由 $\K$ 拟凝聚，存在 $m$ 使 $f^mv(t')=v(f^mt')$ 延拓为 $z\in\Gamma(X,\K)$。将（1.3.1），或（I，5.1.9.2），用于 $\F$，可知 $\Gamma(X,\G)\to\Gamma(X,\K)\to0$ 正合，因此 $z=v(t)$，其中 $t\in\Gamma(X,\G)$。若 $t''=t|D(f)$，则 $v(f^mt'-t'')=0$，故 $f^mt'-t''=u(s')$，$s'\in\Gamma(D(f),\F)$。由 $\F$ 拟凝聚，存在 $n>0$ 使 $f^ns'$ 延拓为 $s\in\Gamma(X,\F)$。于是 $f^{m+n}t'-f^nt''=u(f^ns')$，所以 $f^{m+n}t'$ 是全局截面 $f^nt+u(s)$ 在 $D(f)$ 上的限制，证明完成。\transnote{原文连续两次把待验证条件写作 d2），最后又将延拓截面写作 $f^nt+u(f^ns)$。此处分别按“消失”“延拓”说明，并在已定义 $s|D(f)=f^ns'$ 的记号下写作 $f^nt+u(s)$。}
\end{proof}\par

% ---- END c3_s01b.tex ----
% ---- BEGIN c3_s02a.tex ----
\section{射影态射的上同调研究}
\subsection{某些上同调群的显式计算}
\src{95--99}
\egapar{2.1.1}
设 $X$ 为预概形，$\LL$ 为可逆 $\OO_X$-模；考虑分次环（$\OI$，5.4.6）
\begin{equation}\tag{2.1.1.1}
S=\Gamma_*(X,\LL)=\bigoplus_{n\in\Z}\Gamma(X,\LL^{\otimes n}).
\end{equation}
设 $(f_i)_{1\leq i\leq r}$ 为 $S$ 中的有限齐次元素族，$f_i\in S_{d_i}$。令 $U_i=X_{f_i}$、$U=\bigcup_iU_i$，以 $\mathfrak U$ 表示覆盖 $(U_i)$。对每个拟凝聚 $\OO_X$-模 $\F$，令
\begin{equation}\tag{2.1.1.2}
\mathrm H^\bullet(\mathfrak U,\F(*))=\bigoplus_{n\in\Z}
\mathrm H^\bullet(\mathfrak U,\F\otimes_{\OO_X}\LL^{\otimes n}),
\end{equation}
\begin{equation}\tag{2.1.1.3}
\mathrm H^\bullet(U,\F(*))=\bigoplus_{n\in\Z}
\mathrm H^\bullet(U,\F\otimes_{\OO_X}\LL^{\otimes n}).
\end{equation}
这些阿贝尔群是双分次的，令
\[
\mathrm H^\bullet(\mathfrak U,\F(*))^{mn}=
\mathrm H^m(\mathfrak U,\F\otimes\LL^{\otimes n}),
\]
并对（2.1.1.3）作类似定义。就第二个次数而言，它们显然是分次 $S$-模；例如，这可由 $\F\mapsto\mathrm H^\bullet(\mathfrak U,\F)$、$\F\mapsto\mathrm H^\bullet(U,\F)$ 的函子性得出。

\egapar{2.1.2}
考虑分次 $S$-模（$\OI$，5.4.6）
\begin{equation}\tag{2.1.2.1}
M=\Gamma_*(\F)=\mathrm H^0(X,\F(*))=\bigoplus_{n\in\Z}\Gamma(X,\F\otimes\LL^{\otimes n}).
\end{equation}
若 $X$ 为底空间诺特的预概形，或拟紧概形，则由（I，9.3.1），通常记 $U_{i_0\cdots i_p}=\bigcap_{k=0}^pU_{i_k}$ 时，有典范同构
\[
\Gamma(U_{i_0\cdots i_p},\F(*))=\mathrm H^0(U_{i_0\cdots i_p},\F(*))
\simeq M_{f_{i_0}\cdots f_{i_p}}.
\]
采用（1.2.2）的记号，右边还可识别为 $\ilim_nM^{(n)}_{i_0\cdots i_p}$。若 $\xi$ 是此极限中的齐次元素，由次数为 $m$ 的 $x\in M^{(n)}_{i_0\cdots i_p}=M$ 表示，就令其次数为
$m-n(d_{i_0}+\cdots+d_{i_p})$。由过渡同态 $\varphi_{nm}$ 的定义，这不依赖于所选代表 $x$，并使上述识别成为分次 $S$-模同构。

仍以 $C_n^p(M)$ 表示从 $[1,r]^{p+1}$ 到 $M$ 的交错映射所成模，用同样方法在 $\ilim_nC_n^p(M)$ 上定义分次 $S$-模结构。像（1.2.2）那样，有保持次数的同构
\begin{equation}\tag{2.1.2.2}
C^p(\mathfrak U,\F(*))=\ilim_nC_n^p(M),
\end{equation}
从而又有
\begin{equation}\tag{2.1.2.3}
C^p(\mathfrak U,\F(*))=C^{p+1}((\mathbf f),M)
=\ilim_nK^{p+1}(\mathbf f^n,M).
\end{equation}
这里，若极限中的元素由上链 $z\in K^{p+1}(\mathbf f^n,M)$ 表示，其相应值 $z(i_0,\ldots,i_p)$ 在齐次分量 $M_m$ 中，则其次数为 $m-n(d_{i_0}+\cdots+d_{i_p})$；这与该元素的代表选择无关。上述同构与上边缘算子相容，故如（1.2.2）可得：

\egastatement{2.1.3}{命题}
设 $X$ 为底空间诺特的预概形，或拟紧概形。存在关于 $\F$ 的典范函子同构
\begin{equation}\tag{2.1.3.1}
\mathrm H^p(\mathfrak U,\F(*))\simeq\mathrm H^{p+1}((\mathbf f),M)\quad(p\geq1),
\end{equation}
及函子正合列
\begin{equation}\tag{2.1.3.2}
0\to\mathrm H^0((\mathbf f),M)\to M\to\mathrm H^0(\mathfrak U,\F(*))
\to\mathrm H^1((\mathbf f),M)\to0.
\end{equation}
而且，关于（2.1.1.1）中环 $S$ 的分次模结构，以上全部同态的次数都为零。

\egastatement{2.1.4}{推论}
若 $X$ 是拟紧概形，且各 $U_i=X_{f_i}$ 仿射，则存在关于 $\F$ 的零次典范函子同构
\begin{equation}\tag{2.1.4.1}
\mathrm H^p(U,\F(*))\simeq\mathrm H^{p+1}((\mathbf f),M)\quad(p\geq1),
\end{equation}
及所有同态均为零次的函子正合列
\begin{equation}\tag{2.1.4.2}
0\to\mathrm H^0((\mathbf f),M)\to M\to\mathrm H^0(U,\F(*))
\to\mathrm H^1((\mathbf f),M)\to0.
\end{equation}
\begin{proof}将（1.4.1）用于（2.1.3）即可。\end{proof}
与（2.1.3）相应的“局部”命题如下。

\egastatement{2.1.5}{命题}
设 $S$ 为非负分次环，$f_i\in S_+$（$1\leq i\leq r$）为次数 $d_i$ 的齐次元素，$M$ 为分次 $S$-模。设 $X=\Proj S$，令 $U_i=D_+(f_i)$、$U=\bigcup_iU_i$，以及
\[
\mathrm H^\bullet(U,\wt M(*))=\bigoplus_{n\in\Z}\mathrm H^\bullet(U,\widetilde{M(n)}).
\]
则存在关于 $M$ 的、在分次 $S$-模结构下次数为零的典范函子同构
\begin{equation}\tag{2.1.5.1}
\mathrm H^p(U,\wt M(*))\simeq\mathrm H^{p+1}((\mathbf f),M)\quad(p\geq1),
\end{equation}
以及所有同态为零次的函子正合列
\begin{equation}\tag{2.1.5.2}
0\to\mathrm H^0((\mathbf f),M)\to M\to\mathrm H^0(U,\wt M(*))
\to\mathrm H^1((\mathbf f),M)\to0.
\end{equation}
\begin{proof}
\begin{sloppypar}
按定义（II，2.5.2），$\Gamma(U_{i_0\cdots i_p},\widetilde{M(n)})=(M_{f_{i_0}\cdots f_{i_p}})_n$，故
$\Gamma(U_{i_0\cdots i_p},\wt M(*))=M_{f_{i_0}\cdots f_{i_p}}$。其余论证与（2.1.4）的证明相同，并利用 $X$ 为概形。
\end{sloppypar}
\end{proof}

\egastatement{2.1.6}{注}
（i）在（2.1.5）的条件下，$\Gamma(U_{i_0\cdots i_p},\wt M)$ 关于 $M$ 正合（$\OI$，1.3.2）。用（1.2.5）的同一论证，若 $0\to M'\to M\to M''\to0$ 是同态次数为零的分次 $S$-模正合列，则对每个 $p\geq0$，有交换图
\begin{equation}\tag{2.1.6.1}
\begin{tikzcd}
\mathrm H^p(U,\wt M''(*))\arrow[r]\arrow[d]&\mathrm H^{p+1}(U,\wt M'(*))\arrow[d]\\
\mathrm H^{p+1}((\mathbf f),M'')\arrow[r]&\mathrm H^{p+2}((\mathbf f),M').
\end{tikzcd}
\end{equation}
（ii）（2.1.5）主要用于 $S$ 是由有限多个一次元素生成的 $A$-代数、且 $A$ 诺特的情形。因为此时每个拟凝聚 $\OO_X$-模都具有 $\wt M$ 的形式（II，2.7.7）。

\egapar{2.1.7}
现在将（2.1.5）用于 $S=A[T_0,\ldots,T_r]$，其中 $A$ 为任意环，各 $T_i$ 为不定元，取 $M=S$ 及 $f_i=T_i$。于是主要归结为计算 $\mathrm H^\bullet((\mathbf T),S)$，其中 $\mathbf T=(T_i)_{0\leq i\leq r}$。

\egastatement{2.1.8}{引理}
若 $S=A[T_0,\ldots,T_r]$，则
\begin{equation}\tag{2.1.8.1}
\mathrm H^i(\mathbf T^n,S)=0\quad(i\ne r+1),
\end{equation}
\begin{equation}\tag{2.1.8.2}
\mathrm H^{r+1}(\mathbf T^n,S)=S/(\mathbf T^n).
\end{equation}
因此，$A$-模 $\mathrm H^{r+1}(\mathbf T^n,S)$ 的一组基由单项式 $T^{\mathbf p}=T_0^{p_0}\cdots T_r^{p_r}$ 模 $(\mathbf T^n)$ 的类组成，其中每个 $i$ 都满足 $0\leq p_i<n$。
\begin{proof}直接来自（1.1.3.5）及（1.1.4），其假设在此显然成立。\end{proof}

\egapar{2.1.9}
对 $n$ 取归纳极限，（2.1.8.1）给出 $i\ne r+1$ 时 $\mathrm H^i((\mathbf T),S)=0$。对 $i=r+1$，归纳系统为 $S/(\mathbf T^n)$；$0\leq n\leq m$ 时的同态 $\varphi_{nm}:S/(\mathbf T^n)\to S/(\mathbf T^m)$ 为乘以 $(T_0\cdots T_r)^{m-n}$。

对 $n>\sup_i p_i$，以 $\xi_{\mathbf p,n}=\xi_{p_0,\ldots,p_r,n}$ 表示 $T_0^{n-p_0}\cdots T_r^{n-p_r}$ 模 $(\mathbf T^n)$ 的类。则 $\varphi_{nm}(\xi_{\mathbf p,n})=\xi_{\mathbf p,m}$，故它们在 $\mathrm H^{r+1}((\mathbf T),S)$ 中具有相同的典范像 $\xi_{\mathbf p}$。依（2.1.2）的次数定义，$\deg\xi_{\mathbf p}=-|\mathbf p|=-(p_0+\cdots+p_r)$。显然，各 $0<p_i\leq n$ 所对应的 $\xi_{\mathbf p,n}$ 构成 $S/(\mathbf T^n)$ 的一组基。由（2.1.8）立即得：

\egastatement{2.1.10}{推论}
采用（2.1.8）的记号，有
\begin{equation}\tag{2.1.10.1}
\mathrm H^i((\mathbf T),S)=0\quad(i\ne r+1),
\end{equation}
而 $\mathrm H^{r+1}((\mathbf T),S)$ 是自由 $A$-模，一组基由所有满足 $p_i>0$（$0\leq i\leq r$）的 $\xi_{p_0,\ldots,p_r}$ 组成。

\egastatement{2.1.11}{注}
设 $N$ 为任意 $A$-模，$M=S\otimes_AN$。（2.1.8）的论证更一般地给出
\begin{equation}\tag{2.1.11.1}
\mathrm H^i(\mathbf T^n,M)=0\quad(i\ne r+1),
\end{equation}
\begin{equation}\tag{2.1.11.2}
\mathrm H^{r+1}(\mathbf T^n,M)=(S/(\mathbf T^n))\otimes_AN.
\end{equation}
末式直接来自（1.1.3.5）。另一方面，$M/(T_0^nM+\cdots+T_{i-1}^nM)$ 识别为
$\bigl(S/(T_0^nS+\cdots+T_{i-1}^nS)\bigr)\otimes_AN$，因为该理想是 $A$-模 $S$ 的直和因子；这使（1.1.4）可用于 $M$，得到（2.1.11.1）。

结合（2.1.10）与（2.1.5），得：
\egastatement{2.1.12}{命题}
设 $A$ 为环，$r>0$ 为整数，$X=\Pb_A^r$（II，4.1.1）。则：
\begin{enumerate}[label=(\roman*)]
\item $i\ne0,r$ 时，$\mathrm H^i(X,\OO_X(*))=0$。
\item 典范同态 $\alpha:S\to\mathrm H^0(X,\OO_X(*))$（II，2.6.2）为双射。
\item $\mathrm H^r(X,\OO_X(*))$ 是自由 $A$-模，其一组基由 $p_i>0$ 的 $\xi_{p_0,\ldots,p_r}$ 组成；该元素次数为 $-|\mathbf p|$，且 $T_i\xi_{p_0,\ldots,p_r}=\xi_{p_0,\ldots,p_i-1,\ldots,p_r}$。
\end{enumerate}
\begin{proof}
将（2.1.5.2）用于 $M=S=A[T_0,\ldots,T_r]$；由（2.1.10.1），$\mathrm H^0((\mathbf T),S)=\mathrm H^1((\mathbf T),S)=0$。由于 $D_+(T_i)$ 覆盖 $X$（II，2.3.14），（2.1.5）可取 $U=X$。只需再将正合列中的映射 $S\to\mathrm H^0(X,\OO_X(*))$ 与 $\alpha$ 识别；这由 $\mathrm H^0(U,\OO_X(*))$ 与 $\mathrm H^0(\mathfrak U,\OO_X(*))$ 的典范识别得出。
\end{proof}

\egastatement{2.1.13}{推论}
只有以下 $(i,n)$ 才可能使 $\mathrm H^i(X,\OO_X(n))\ne0$：$i=0,n\geq0$；或 $i=r,n\leq-(r+1)$。
若 $A\ne0$，这些指标对的上同调确实非零；这由（2.1.12）及 $n\geq0$ 时 $S_n\ne0$ 得出。本章应用中主要使用较弱的结论：

\egastatement{2.1.14}{推论}
各 $A$-模 $\mathrm H^i(X,\OO_X(n))$ 都自由且有限型；若 $i>0$，则 $n\geq0$ 时为零。

\egastatement{2.1.15}{命题}
设 $Y$ 为预概形，$\E$ 为秩 $r+1$ 的局部自由 $\OO_Y$-模，$X=\Pb(\E)$ 为相应射影丛，$f:X\to Y$ 为结构态射。只有 $i=0,n\geq0$ 或 $i=r,n\leq-(r+1)$ 时，$\mathrm R^if_*(\OO_X(n))$ 才可能非零。此外，典范同态（II，3.3.2）
\[
\alpha:\Sym_{\OO_Y}(\E)\longrightarrow\Gamma_*(\OO_X)
=\mathrm R^0f_*(\OO_X(*))=\bigoplus_{n\in\Z}f_*(\OO_X(n))
\]
为同构。
\begin{proof}
问题在 $Y$ 上局部，可设 $Y=\Spec A$，$\E=\wt E$，其中 $E=A^{r+1}$；结合（1.4.11），立即归结为（2.1.12）。
\end{proof}

\egastatement{2.1.16}{注}
以后将用以下结果补充（2.1.15）。令 $\omega=f^*(\bigwedge^{r+1}\E)(-r-1)$，这是可逆 $\OO_X$-模。则：

（i）有典范同构
\begin{equation}\tag{2.1.16.1}
\rho:\mathrm R^rf_*(\omega)\iso\OO_Y.
\end{equation}
（ii）杯积配对（0，12.2.2）
\begin{equation}\tag{2.1.16.2}
\mathrm R^rf_*(\OO_X(n))\times\mathrm R^0f_*(\omega(-n))\to\mathrm R^rf_*(\omega)
\end{equation}
与同构 $\rho^{-1}$ 复合，给出 $\mathrm R^rf_*(\OO_X(n))$ 到局部自由 $\OO_Y$-模
\[
\mathrm R^0f_*(\omega(-n))=\bigl(\bigwedge^{r+1}\E\bigr)\otimes_{\OO_Y}
(\Sym_{\OO_Y}(\E))_{-n}
\]
的对偶的同构。\transnote{本段保留底本的两处记号：按（2.1.16.1）的方向，配对应与 $\rho$ 复合；按 $\omega=f^*(\bigwedge^{r+1}\E)(-r-1)$ 及（2.1.15），末式对称代数分量的下标应为 $-n-r-1$，而底本印作 $-n$。}\par

% ---- END c3_s02a.tex ----
% ---- BEGIN c3_s02b.tex ----
\subsection{射影态射的基本定理}
\src{100--102}
\egastatement{2.2.1}{定理（Serre）}
设 $Y$ 为诺特预概形，$f:X\to Y$ 为固有态射，$\LL$ 为相对于 $f$ 丰沛的可逆 $\OO_X$-模。对每个 $\OO_X$-模 $\F$ 及 $n\in\Z$，令 $\F(n)=\F\otimes_{\OO_X}\LL^{\otimes n}$。则对每个凝聚 $\OO_X$-模 $\F$：
\begin{enumerate}[label=(\roman*)]
\item 各 $\mathrm R^qf_*(\F)$ 是凝聚 $\OO_Y$-模。
\item 存在整数 $N$，使 $n\geq N$ 时，对每个 $q>0$ 有 $\mathrm R^qf_*(\F(n))=0$。
\item 存在整数 $N$，使 $n\geq N$ 时，典范同态 $f^*(f_*(\F(n)))\to\F(n)$ 满射。
\end{enumerate}
\eproof 首先，若将 $\LL$ 换为 $\LL^{\otimes d}$（$d>0$）后定理成立，则原定理也成立。因为可写
$\F(n)=(\F\otimes\LL^{\otimes r})\otimes\LL^{\otimes hd}$，其中 $h\geq0$、$0\leq r<d$。对每个 $r$，假设给出整数 $N_r$，使 $h\geq N_r$ 时，对 $\F\otimes\LL^{\otimes r}$ 成立（ii）、（iii）；取 $N$ 为各 $dN_r$ 的最大值即可。故可设 $\LL$ 相对于 $f$ 很丰沛（II，4.6.11）。于是存在占优开 $Y$-浸入 $i:X\to P=\Proj\SSS$，其中 $\SSS$ 为非负分次拟凝聚 $\OO_Y$-代数，$\SSS_1$ 有限型并生成 $\SSS$，且 $\LL\simeq i^*\OO_P(1)$（II，4.4.7）。由于 $f$ 固有，$i$ 也固有（II，5.4.4），所以 $i$ 是同构 $X\simeq P$。因此只需考虑 $X=\Proj\SSS$、$\LL=\OO_X(1)$。此时定理由以下命题推出。

\egastatement{2.2.2}{命题}
设 $A$ 为诺特环，$S$ 为非负分次 $A$-代数，$S_1$ 作为 $A$-模具有 $r+1$ 个生成元，并生成代数 $S$。令 $X=\Proj S$。对每个凝聚 $\OO_X$-模 $\F$：
\begin{enumerate}[label=(\roman*)]
\item 各 $A$-模 $\mathrm H^q(X,\F)$ 有限型。
\item $q>r$ 时，$\mathrm H^q(X,\F)=0$。
\item 存在整数 $N$，使 $n\geq N$ 时，对每个 $q>0$ 有 $\mathrm H^q(X,\F(n))=0$。
\item 存在整数 $N$，使 $n\geq N$ 时，$\F(n)$ 由其全局截面生成。
\end{enumerate}
先说明（2.2.2）如何推出（2.2.1）。在已经归结的 $X=\Proj\SSS$ 情形，$Y$ 拟紧，可用有限个仿射开集 $U_\alpha$ 覆盖，其环诺特，且 $\SSS|U_\alpha$ 由 $\SSS_1$ 在 $U_\alpha$ 上的有限多个截面生成。若（2.2.2）成立，只需在（2.2.1，ii）、（iii）中取各 $U_\alpha$ 对应整数的最大值，并使用（1.4.11）、（II，3.4.7）。

为证（2.2.2），注意 $X$ 可识别为 $P=\Pb_A^r$ 的闭子概形（II，3.6.2）。若 $j:X\to P$ 为典范嵌入，则 $j_*(\F)$ 凝聚，且 $j_*(\F(n))=(j_*(\F))(n)$（II，3.4.5及3.5.2）。利用（G，II，定理4.9.1的推论），归结为 $X=\Pb_A^r$、$S=A[T_0,\ldots,T_r]$。$X$ 被 $r+1$ 个仿射开集 $D_+(T_i)$ 覆盖，故（ii）由（1.4.12）得出；（iv）已在（II，2.7.9）证明。

现在同时证明（i）、（iii）。对于 $\F=\OO_X(m)$，它们成立（2.1.13），故对于有限个这种层的直和也成立。由（ii），$q>r$ 时两者平凡成立。对 $q$ 作下降归纳。已知 $\F$ 是有限个 $\OO_X(m_j)$ 的直和 $\E$ 的商（II，2.7.10），所以有正合列 $0\to\R\to\E\to\F\to0$，其中 $\R$ 凝聚（$\OI$，5.3.3），且 $\E$ 满足（i）、（iii）。由于扭转是正合函子，各 $n\in\Z$ 都给出
\[
0\to\R(n)\to\E(n)\to\F(n)\to0,
\]
从而有上同调正合列
\[
\mathrm H^{q-1}(X,\E(n))\to\mathrm H^{q-1}(X,\F(n))\to\mathrm H^q(X,\R(n)).
\]
$\E(n)$ 是各 $\OO_X(n+m_j)$ 的直和，故左项有限型；归纳假设说明右项也有限型。$A$ 诺特，故中项对每个 $n$ 有限型，特别适用于 $n=0$。另一方面，归纳假设给出 $N$，使 $n\geq N$ 时右项为零；在证明正次数的（iii）时，还可将 $N$ 取大，使左项为零，因为 $\E$ 满足（iii）。于是中项也为零，证明完成。\finishproof

\egastatement{2.2.3}{推论}
在（2.2.1）的假设下，设 $\F\to\G\to\HH$ 为凝聚 $\OO_X$-模正合列。则存在整数 $N$，使 $n\geq N$ 时
\[
f_*(\F(n))\to f_*(\G(n))\to f_*(\HH(n))
\]
正合。
\begin{proof}
设 $\F',\G',\G''$ 分别为 $\F\to\G$ 的核、像、余核；则 $\G'$ 也是 $\G\to\HH$ 的核，而 $\G''$ 是其像，再以 $\HH''$ 表示其余核。这些模都凝聚（$\OI$，5.3.4）。因扭转正合，只需对足够大的 $n$ 证明
\[
\begin{aligned}
0&\to f_*(\F'(n))\to f_*(\F(n))\to f_*(\G'(n))\to0,\\
0&\to f_*(\G'(n))\to f_*(\G(n))\to f_*(\G''(n))\to0,\\
0&\to f_*(\G''(n))\to f_*(\HH(n))\to f_*(\HH''(n))\to0
\end{aligned}
\]
各列正合。因此可直接假设 $0\to\F\to\G\to\HH\to0$ 正合。其上同调长正合列的开头为
\[
0\to f_*(\F(n))\to f_*(\G(n))\to f_*(\HH(n))\to\mathrm R^1f_*(\F(n))\to\cdots,
\]
结论由（2.2.1，ii）得出。
\end{proof}

\egastatement{2.2.4}{推论}
设 $Y$ 为诺特预概形，$f:X\to Y$ 为有限型态射，$\LL$ 为相对于 $f$ 丰沛的可逆 $\OO_X$-模，仍令 $\F(n)=\F\otimes\LL^{\otimes n}$。设 $\F\to\G\to\HH$ 是凝聚 $\OO_X$-模正合列，且 $\F,\HH$ 的支撑在 $Y$ 上固有（II，5.4.10）。则存在 $N$，使 $n\geq N$ 时 $f_*(\F(n))\to f_*(\G(n))\to f_*(\HH(n))$ 正合。
\begin{proof}
（2.2.1）开头的同一论证说明，可将 $\LL$ 换为正张量幂，故可设其相对于 $f$ 很丰沛（II，4.6.11）。于是可将 $X$ 识别为 $Y$-概形 $Z=\Proj\SSS$ 的开集，其中 $\SSS$ 为非负分次拟凝聚代数，由有限型 $\SSS_1$ 生成，且 $\LL=i^*\OO_Z(1)$，$i:X\to Z$ 为典范嵌入（II，4.4.7）。由于 $\Supp\G$ 在 $X$ 中闭，且包含于 $\Supp\F\cap\Supp\HH$，它在 $Y$ 上固有；因此三个层的支撑都在 $Z$ 中闭（II，5.4.10）。\transnote{底本此处的交号 $\cap$ 应读作并号 $\cup$：由三项正合性得出的是 $\Supp\G\subset\Supp\F\cup\Supp\HH$。有限并仍足以支持本段的固有性论证。}
于是 $\F'=i_*(\F)$、$\G'=i_*(\G)$、$\HH'=i_*(\HH)$ 是凝聚 $\OO_Z$-模，并构成正合列。若 $g:Z\to Y$ 是结构态射，则 $f=g\circ i$，且 $\F'(n)=i_*(\F(n))$，对另两个层同样成立。将（2.2.3）用于 $\F',\G',\HH'$ 即得结论。
\end{proof}

\egastatement{2.2.5}{注}
（i）（2.2.1，i）在仅假设 $Y$ 局部诺特时仍成立，因为该性质在 $Y$ 上局部。又（2.2.1）的假设保证：对每个开集 $U\subset Y$，$f$ 在 $f^{-1}(U)$ 上的限制为射影态射，且 $\LL|f^{-1}(U)$ 相对于它丰沛（II，5.5.5，iii；4.6.4）。

（ii）已经看到，当只假设 $X$ 为拟紧概形或底空间诺特的预概形、$f$ 拟紧时，（2.2.1，iii）仍成立（II，4.6.8）。但即使 $Y$ 是域 $K$ 的谱、$f$ 拟射影，（2.2.1，ii）也未必成立。例如，令 $X'=\Spec K[T_0,\ldots,T_r]$，$X=\bigcup_{i=0}^rD(T_i)$。开浸入 $X\to X'$ 拟紧，所以结构态射 $f:X\to Y$ 拟仿射（II，5.1.10），因而 $\OO_X$ 相对于 $f$ 很丰沛（II，5.1.6）。但 $\Gamma(X,\OO_X)$ 识别为各局部化环 $K[T_0,\ldots,T_r]_{T_i}$ 的交（I，8.2.1.1），也就是 $K[T_0,\ldots,T_r]$。由（1.4.3.1）及（1.1.3.5），对每个 $n$ 有
\[
\mathrm H^r(X,\OO_X^{\otimes n})=\mathrm H^r(X,\OO_X)=A\ne0.
\]
\transnote{此例应在 $r>0$ 下阅读。底本末式保留了没有在此例中解释的记号 $A$；非零性可由（2.1.10）与（1.4.3.1）核对，这里的最高次上同调不应直接识别为上述多项式函数环。}

\subsection{应用于分次代数层和分次模层}
\src{102--107}
\egastatement{2.3.1}{定理}
设 $Y$ 为诺特预概形，$\SSS$ 为非负分次、拟凝聚、有限型的 $\OO_Y$-代数，$X=\Proj\SSS$，$q:X\to Y$ 为结构态射，$\M$ 为满足（TF）的拟凝聚分次 $\SSS$-模。则存在整数 $N$，使 $n\geq N$ 时典范同态（II，8.14.5.1）
\[
\alpha_n:\M_n\to q_*\bigl(\sProj_0(\M(n))\bigr)
=q_*\bigl((\sProj\M)_n\bigr)
\]
为双射。换言之，典范同态 $\alpha:\M\to\Gamma_*(\sProj\M)$ 是（TN）-同构。
\begin{proof}
可以只考虑 $\M$ 有限型的情形（II，3.4.2）。由于 $Y$ 拟紧，存在整数 $d>0$，使 $\SSS^{(d)}$ 由有限型拟凝聚 $\OO_Y$-模 $\SSS_d$ 生成（II，3.1.10）；$Y$ 诺特，故 $\SSS_d$ 凝聚。$\M$ 是各 $\M^{(d,k)}$（$0\leq k<d$）的直和，它们都是有限型拟凝聚分次 $\SSS^{(d)}$-模；这由（II，2.1.6，iii）局部得出。只需证明各同态
\[
\M^{(d,k)}\to\bigl(\Gamma_*(\sProj\M)\bigr)^{(d,k)}
\]
为（TN）-同构。结合（II，8.14.13），特别是其交换图（8.14.13.4），可归结到 $\SSS$ 由凝聚 $\SSS_1$ 生成的情形。像（2.2.2）开头那样，可再设 $Y=\Spec A$、$\SSS=\wt S$、$\M=\wt M$，其中 $A$ 诺特，$S_1$ 为有限型 $A$-模，$M$ 为有限型分次 $S$-模。

说明只需证 $M=S$。一般情况下有正合列 $L'\to L\to M\to0$，其中 $L,L'$ 是形如 $S(m)$ 的分次模的有限直和。若结论对 $S$ 成立，也对 $S(m)$ 以及 $L,L'$ 成立。考虑交换图
\[
\begin{tikzcd}[column sep=small]
L'_n\arrow[r]\arrow[d]&L_n\arrow[r]\arrow[d]&M_n\arrow[r]\arrow[d]&0\\
q_*(\wt L'(n))\arrow[r]&q_*(\wt L(n))\arrow[r]&q_*(\wt M(n))\arrow[r]&0.
\end{tikzcd}
\]
当 $n$ 足够大时，第二行由（2.2.3）正合；第一行亦正合，而左侧两竖箭头为同构，故第三个也为同构。

为证 $M=S$ 时的结论，先设 $S=A[T_0,\ldots,T_r]$；此时即为（2.1.12，ii）。\transnote{底本此处引作（2.1.11，ii），但具有该“（ii）”结论的是（2.1.12）。}
一般情况下，$S$ 是 $S'=A[T_0,\ldots,T_r]$ 模一个分次理想的商，故 $X$ 是 $X'=\Pb_A^r$ 的闭子概形（II，2.9.2）。若 $j:X\to X'$ 为典范嵌入，则 $j_*(\widetilde{S(n)})=(\sProj S)(n)$，右边把 $S$ 看作分次 $S'$-模（II，2.8.7）。由于该分次模满足（TF），由已证的多项式环情形，对足够大的 $n$，典范同态
$S_n\to\Gamma(X',j_*(\widetilde{S(n)}))$ 为双射。再用
$\Gamma(X',j_*(\widetilde{S(n)}))=\Gamma(X,\widetilde{S(n)})$，证明完成。
\end{proof}

\egastatement{2.3.2}{推论}
在（2.3.1）的假设下，令 $\SSS_X=\bigoplus_{n\in\Z}\OO_X(n)$，设 $\F$ 为有限型拟凝聚分次 $\SSS_X$-模。则 $\Gamma_*(\F)$ 满足（TF）。
\begin{proof}
（2.3.1）的证明已经说明，$X\simeq\Proj\SSS^{(d)}$（II，3.1.8）在 $Y$ 上有限型（II，3.4.1）。由（II，8.14.9），$\F$ 同构于 $\sProj\M$，其中 $\M$ 是有限型拟凝聚分次 $\SSS$-模。由（2.3.1），$\Gamma_*(\F)$ 与 $\M$（TN）-同构，因此满足（TF）。
\end{proof}

\egastatement{2.3.3}{说明}
设 $Y$ 诺特，$\SSS$ 满足（2.3.1）的条件，$X=\Proj\SSS$。设 $K_{\SSS}$ 为满足（TF）的拟凝聚分次 $\SSS$-模所成阿贝尔范畴，$K_{\SSS}^0$ 为其中满足（TN）的模构成的子范畴；又设 $K_X$ 为有限型拟凝聚分次 $\SSS_X$-模 $\E$ 的范畴。由于 $\SSS_X$ 周期性（II，8.14.4及8.14.12），后一条件也等价于各 $\E_n$ 为凝聚 $\OO_X$-模。由（II，8.14.8及8.14.10）和（2.3.2），函子 $\M\mapsto\sProj\M$ 与 $\E\mapsto\Gamma_*(\E)$ 给出商范畴 $K_{\SSS}/K_{\SSS}^0$（T，I，1.11）与 $K_X$ 的等价（T，I，1.2）。若 $\SSS$ 由 $\SSS_1$ 生成，可将 $K_X$ 换成凝聚 $\OO_X$-模范畴（II，8.14.12）。

\egastatement{2.3.4}{命题}
设 $Y$ 为诺特预概形。
\begin{enumerate}[label=(\roman*)]
\item 设 $\SSS$ 为非负分次、拟凝聚、有限型的 $\OO_Y$-代数，$X=\Proj\SSS$，$\SSS_X=\sProj\SSS=\bigoplus_{n\in\Z}\OO_X(n)$。则 $\SSS_X$ 为周期性分次 $\OO_X$-代数，其齐次分量为凝聚模；若 $d>0$ 是一个周期，则 $(\SSS_X)_d=\OO_X(d)$ 可逆且在 $Y$ 上丰沛。此外，典范同态 $\alpha:\SSS\to\Gamma_*(\SSS_X)$ 为（TN）-同构。
\item 反过来，设 $q:X\to Y$ 为射影态射，$\SSS'$ 为分次 $\OO_X$-代数，各齐次分量 $\SSS'_n$ 都凝聚；并且存在周期 $d>0$，使 $\SSS'_d$ 为相对于 $q$ 丰沛的可逆模。则 $\SSS=\bigoplus_{n\geq0}q_*(\SSS'_n)$ 为非负分次、拟凝聚、有限型 $\OO_Y$-代数；存在 $Y$-同构 $r:X\iso\Proj\SSS$，使 $r^*(\sProj\SSS)$ 与 $\SSS'$ 作为分次 $\OO_X$-代数同构。
\end{enumerate}
\eproof（i）实质上各断言均已证明。最后一个为（2.3.2）的特例。周期性见（II，8.14.14）；某个周期 $d$ 对应的 $\OO_X(d)$ 可逆且 $Y$-丰沛，见（II，4.6.18）。又对 $0\leq k<d$，$(\SSS_X)^{(d,k)}$ 为有限型 $(\SSS_X)^{(d)}$-模（II，8.14.14），故每个 $(\SSS_X)_n$ 局部由（II，2.1.6，ii）为有限型拟凝聚模。由于 $\OO_X$ 凝聚，各分量凝聚。

（ii）将周期 $d$ 换为其适当倍数，可设 $\LL=\SSS'_d$ 相对于 $q$ 很丰沛（II，4.6.11）。按假设，$\SSS'^{(d)}=\bigoplus_{n\in\Z}\LL^{\otimes n}$，故 $\SSS^{(d)}=\bigoplus_{n\geq0}q_*(\LL^{\otimes n})$。已知（II，3.1.8及3.2.9），存在 $Y$-同构
$s:X'=\Proj\SSS\iso X''=\Proj\SSS^{(d)}$，且 $s^*\OO_{X''}(n)=\OO_{X'}(nd)$。以下命题将给出 $X\simeq X'$。

\egastatement{2.3.4.1}{命题}
设 $Y$ 为诺特预概形，$q:X\to Y$ 为射影态射，$\LL$ 为相对于 $q$ 很丰沛的可逆模。则 $\SSS=\bigoplus_{n\geq0}q_*(\LL^{\otimes n})$ 是有限型拟凝聚分次 $\OO_Y$-代数，且充分大的 $n$ 满足 $\SSS_n=\SSS_1^n$；存在 $Y$-同构 $r:X\iso P=\Proj\SSS$，使 $\LL=r^*\OO_P(1)$。
\begin{proof}
由射影性以及（II，5.4.4及4.4.7），存在 $Y$-同构 $r':X\iso P'=\Proj\TT$，其中 $\TT$ 由有限型 $\TT_1$ 生成，且 $\LL=r'^*\OO_{P'}(1)$。于是 $\SSS=\bigoplus_{n\geq0}q'_*(\OO_{P'}(n))$，$q':P'\to Y$ 为结构态射。由（2.3.1），充分大的 $n$ 使典范同态 $\alpha_n:\TT_n\to\SSS_n=q_*(\LL^{\otimes n})$ 双射。因 $\TT_n=\TT_1^n$，故此时也有 $\SSS_n=\SSS_1^n$。又 $\alpha:\TT\to\SSS$ 是（TN）-同构，故 $\Phi=\Proj\alpha:\Proj\SSS\to\Proj\TT$ 为同构（II，3.6.1），并有 $\Phi^*(\widetilde{\TT(n)})=\widetilde{\SSS(n)_{[\alpha]}}$（II，3.5.2）。因相应的分次模与 $\SSS(n)$（TN）-同构，（II，3.4.2）给出 $\Phi^*\OO_{P'}(n)=\OO_P(n)$。

只需再证 $\SSS$ 有限型。各 $\SSS_n=q'_*(\OO_{P'}(n))$ 由（2.2.1）凝聚；而充分大的 $n$ 有 $\SSS_n=\SSS_1^n$，所以 $\SSS$ 由有限多个齐次分量的和生成，该和凝聚，因而代数有限型（I，9.6.2）。
\end{proof}
回到（2.3.4）。已得到 $Y$-同构 $r'':X\iso X''$，使 $r''_*(\LL^{\otimes n})=\OO_{X''}(n)$；以 $q''$ 表示 $X''\to Y$。$\SSS'$ 为各 $\SSS'^{(d,k)}$ 的直和；由周期性及各 $\SSS'_n$ 有限型，这些都是有限型拟凝聚分次 $\SSS'^{(d)}$-模（II，8.14.12）。令
$\F^{(k)}=r''_*(\SSS'^{(d,k)})$，它们是有限型拟凝聚分次 $\SSS_{X''}$-模。由（II，8.14.8），典范同态
$\beta:\sProj(\Gamma_*(\F^{(k)}))\to\F^{(k)}$ 为同构。

又 $\Gamma_*(\F^{(k)})_n=q''_*(\F_n^{(k)})=q_*((\SSS'^{(d,k)})_n)$，$n\geq0$ 时按定义等于 $(\SSS^{(d,k)})_n$。因此典范嵌入 $\SSS^{(d,k)}\to\Gamma_*(\F^{(k)})$ 是（TN）-同构，故
\[
\sProj(\SSS^{(d,k)})=\sProj(\Gamma_*(\F^{(k)})),\qquad
r''^*(\sProj(\SSS^{(d,k)}))=\SSS'^{(d,k)}.
\]
结合典范同构
$\sProj(\SSS^{(d,k)})=s_*(\sProj\SSS)^{(d,k)}$（II，8.14.13.1），便得 $r^*(\sProj\SSS)\simeq\SSS'$。最后，由（2.3.2），各 $\Gamma_*(\F^{(k)})$ 满足（TF），所以各 $\SSS^{(d,k)}$ 也满足；又 $\SSS'_n$ 凝聚，故 $\SSS_n=q_*(\SSS'_n)$ 凝聚（2.2.1），从而各 $\SSS^{(d,k)}$ 都为有限型 $\SSS^{(d)}$-模。已经由（2.3.4.1）知道 $\SSS^{(d)}$ 为有限型 $\OO_Y$-代数，因此 $\SSS$ 也有限型。\finishproof

\egastatement{2.3.5}{命题}
设 $Y$ 为诺特整预概形，$X$ 为整预概形，$f:X\to Y$ 为射影双有理态射。则存在凝聚分式理想 $\J\subset\R(Y)$（II，8.1.2），使 $X$ 与沿 $\J$ 爆破得到的预概形 $Y$-同构（II，8.1.3）。并且存在 $Y$ 的开集 $U$，使 $f:f^{-1}(U)\to U$ 为同构，且 $\J|U$ 可逆。
\begin{proof}
存在相对于 $f$ 很丰沛的可逆模 $\LL$（II，4.4.2及5.3.2）。由（2.3.4.1），$X\simeq\Proj\SSS$，其中 $\SSS=\bigoplus_{n\geq0}f_*(\LL^{\otimes n})$。各直接像无挠（I，7.4.5），故 $\SSS$ 无挠，可典范地识别为 $\SSS\otimes_{\OO_Y}\R(Y)$ 的子模（I，7.4.1）。后者是常层（I，7.3.6），在任一非空开集上的限制就能确定它。取非空 $U'\subset U$ 使 $\LL|f^{-1}(U')\simeq\OO_X|f^{-1}(U')$；此时各 $f_*(\LL^{\otimes n})|U'\simeq\OO_Y|U'$，故 $\SSS\otimes\R(Y)\simeq\R(Y)[T]$。

由（2.3.4.1），$\SSS$ 与其中由 $f_*(\LL)$ 的典范像生成的子代数（TN）-同构。在上述识别下，该像为 $\J T$，其中 $\J$ 是 $\R(Y)$ 的凝聚子 $\OO_Y$-模（2.2.1），其在 $U'$ 上限制同构于 $\OO_Y|U'$，所以 $\J|U'$ 可逆。于是 $\SSS$ 与 $\bigoplus_{n\geq0}\J^n$（TN）-同构，命题得证。
\end{proof}

\egastatement{2.3.6}{推论}
在（2.3.5）的假设下，再设对每个非零凝聚子模 $\J\subset\R(Y)$，都存在可逆 $\OO_Y$-模 $\LL$，使
$\Gamma(Y,\LL\otimes\shHom_{\OO_Y}(\J,\OO_Y))\ne0$。
则在（2.3.5）中可取 $\J$ 为 $\OO_Y$ 的理想。只要存在丰沛 $\OO_Y$-模，这一附加条件就成立。
\begin{proof}
由（$\OI$，5.4.2），
\[
\begin{aligned}
\LL\otimes\shHom(\J,\OO_Y)
&=\shHom(\LL^{-1},\shHom(\J,\OO_Y))\\
&=\shHom(\J\otimes\LL^{-1},\OO_Y).
\end{aligned}
\]
所以假设意味着有非零同态 $u:\J\otimes\LL^{-1}\to\OO_Y$。对每个 $y\in Y$，$(\J\otimes\LL^{-1})_y$ 可识别为 $\OO_{Y,y}$ 的分式域中的子模（I，7.1.5）；故 $u_y$ 单射，$u$ 将它同构于 $\OO_Y$ 的理想 $\J'$。但 $\Proj(\bigoplus_n\J^n)$ 与 $\Proj(\bigoplus_n(\J\otimes\LL^{-1})^n)$ 为 $Y$-同构（II，3.1.8），第一断言得证。

为证第二断言，$\F=\shHom_{\OO_Y}(\J,\OO_Y)$ 凝聚且非零，因为存在开集 $U$ 使 $\J|U$ 可逆。若 $\LL$ 丰沛，则某个 $n$ 使 $\F(n)=\F\otimes\LL^{\otimes n}$ 由全局截面生成（II，4.5.5），从而 $\Gamma(Y,\F(n))\ne0$。
\end{proof}

\egastatement{2.3.7}{推论}
设 $X,Y$ 为域 $k$ 上的射影整概形，$f:X\to Y$ 为双有理 $k$-态射。则 $X$ 与沿 $Y$ 的某个闭子概形 $Y'$（不要求约化）爆破所得的 $k$-概形同构。
\begin{proof}
$f$ 射影（II，5.5.5，v），且 $Y$ 在 $k$ 上射影，故（2.3.6）的附加条件成立。取该推论中凝聚理想所定义的闭子概形即可。
\end{proof}

\egastatement{2.3.8}{注}
第IV章研究除子时，将看到：若（2.3.5）中各局部环 $\OO_{Y,y}$ 为唯一分解整环，例如 $Y$ 非奇异，则 $X$ 可由沿 $Y$ 的某个闭子预概形 $Y'$ 爆破得到，且 $Y'$ 的底空间包含于 $Y-U$。\par

% ---- END c3_s02b.tex ----
% ---- BEGIN c3_s02c.tex ----
\subsection{基本定理的一个推广}
\src{107--109}
\egastatement{2.4.1}{定理}
设 $Y$ 为诺特预概形，$\SSS$ 为有限型拟凝聚 $\OO_Y$-代数。设 $f:X\to Y$ 为射影态射，$\SSS'=f^*(\SSS)$，$\M$ 为有限型拟凝聚 $\SSS'$-模。则：
\begin{enumerate}[label=(\roman*)]
\item 对每个 $p\in\Z$，$\mathrm R^pf_*(\M)$ 是有限型 $\SSS$-模。
\item 若 $\LL$ 为相对于 $f$ 丰沛的可逆 $\OO_X$-模，并令 $\M(n)=\M\otimes\LL^{\otimes n}$，则存在整数 $N$，使 $n\geq N$ 时，对每个 $p>0$ 有
\begin{equation}\tag{2.4.1.1}
\mathrm R^pf_*(\M(n))=0,
\end{equation}
且典范同态 $f^*(f_*(\M(n)))\to\M(n)$（$\OI$，4.4.3）满射。
\end{enumerate}
\begin{proof}
令 $Y'=\Spec\SSS$、$X'=\Spec\SSS'$，则 $X'=X\times_YY'$（II，1.5.5）。令 $g:Y'\to Y$、$g':X'\to X$ 为仿射结构态射，$f'=f_{(Y')}$；于是有交换图
\[
\begin{tikzcd}
X\arrow[d,"f"']&X'\arrow[l,"g'"']\arrow[d,"f'"]\arrow[dl,"h",description]\\
Y&Y'\arrow[l,"g"]
\end{tikzcd}
\]
其中 $f'$ 射影（II，5.5.5，iii），$h=f\circ g'=g\circ f'$。

（i）以 $\wt\M$ 表示把 $X'$ 看作 $X$ 上仿射概形时、与拟凝聚 $\SSS'$-模 $\M$ 相伴的 $\OO_{X'}$-模（II，1.4.3），故 $\M=g'_*(\wt\M)$。由（II，1.4.5），$\wt\M$ 有限型；$g,f'$ 均有限型，故 $h$ 有限型（II，1.3.7；I，6.3.4，ii），从而 $X'$ 诺特，$\wt\M$ 凝聚。因 $g'$ 仿射，典范同态
$\mathrm R^pf_*(\M)\to\mathrm R^ph_*(\wt\M)$ 为双射（1.3.4）。它也是 $\SSS$-模同态：定义 $\M$ 的 $\SSS'$-模结构的同态为
\begin{equation}\tag{2.4.1.2}
g'_*(\OO_{X'})\otimes g'_*(\wt\M)\longrightarrow g'_*(\wt\M),
\end{equation}
这里 $\SSS'=g'_*(\OO_{X'})$。由它典范地得到（0，12.2.2）
\[
f_*(g'_*(\OO_{X'}))\otimes\mathrm R^pf_*(g'_*(\wt\M))\to
\mathrm R^pf_*(g'_*(\wt\M)).
\]
而（2.4.1.2）本身由 $\OO_{X'}\otimes\wt\M\to\wt\M$ 经（$\OI$，4.2.2.1）得到，所以图
\[
\begin{tikzcd}[column sep=small]
f_*g'_*(\OO_{X'})\otimes\mathrm R^pf_*(g'_*(\wt\M))\arrow[r]\arrow[d]&\mathrm R^pf_*(g'_*(\wt\M))\arrow[d]\\
h_*(\OO_{X'})\otimes\mathrm R^ph_*(\wt\M)\arrow[r]&\mathrm R^ph_*(\wt\M)
\end{tikzcd}
\]
交换（0，12.2.6）。再将水平箭头与来自典范同态
$\SSS\to f_*(f^*\SSS)=f_*(\SSS')=f_*g'_*(\OO_{X'})=h_*(\OO_{X'})$
的同态复合，就得上述断言。

又因 $g$ 仿射、$f'$ 分离拟紧，典范同态
$\mathrm R^ph_*(\wt\M)\to g_*(\mathrm R^pf'_*(\wt\M))$
为双射（1.4.14）；同样用（0，12.2.6.2）的交换性可知它为 $\SSS$-模同构。$f'$ 射影且 $\wt\M$ 凝聚，故 $\mathrm R^pf'_*(\wt\M)$ 凝聚（2.2.1），于是其直接像为有限型 $\SSS$-模（II，1.4.5）。

（ii）令 $\LL'=g'^*\LL$，这是可逆 $\OO_{X'}$-模。由（$\OI$，5.4.10），对每个 $n$ 有典范同构
\[
g'_*(\wt\M\otimes\LL'^{\otimes n})=g'_*(\wt\M)\otimes\LL^{\otimes n}=\M(n).
\]
将（i）的论证用于 $\wt\M\otimes\LL'^{\otimes n}$，得到
$\mathrm R^pf_*(\M(n))\simeq g_*(\mathrm R^pf'_*(\wt\M\otimes\LL'^{\otimes n}))$。
$\LL'$ 相对于 $f'$ 丰沛（II，4.6.13，iii），所以（2.2.1）给出 $N$，使 $p>0,n\geq N$ 时右边为零，证明（2.4.1.1）。又可将 $N$ 取大，使
$f'^*(f'_*(\wt\M\otimes\LL'^{\otimes n}))\to\wt\M\otimes\LL'^{\otimes n}$
满射。$g'_*$ 在相应拟凝聚模范畴中正合（II，1.4.4），故其直接像也满射到 $\M(n)$。而（II，1.5.2）给出
\[
\begin{aligned}
g'_*(f'^*(f'_*(\wt\M\otimes\LL'^{\otimes n})))
&=f^*(g_*(f'_*(\wt\M\otimes\LL'^{\otimes n})))\\
&=f^*(f_*(g'_*(\wt\M\otimes\LL'^{\otimes n})))
=f^*(f_*(\M(n))),
\end{aligned}
\]
证明完成。
\end{proof}

\egapar{2.4.2}
尤其将（2.4.1）用于 $\SSS$ 非负分次、$\M=\bigoplus_k\M_k$ 为分次 $\SSS'$-模的情形。在相同有限性假设下，对每个 $p$，$\bigoplus_k\mathrm R^pf_*(\M_k)$ 是有限型 $\SSS$-模；在（2.4.1，ii）的附加假设下，还存在 $N$，使 $n\geq N$ 时，对所有 $p>0,k\in\Z$，有 $\mathrm R^pf_*(\M_k(n))=0$，并且 $f^*(f_*(\M_k(n)))\to\M_k(n)$ 满射。

\subsection{Euler--Poincaré 特征与 Hilbert 多项式}
\src{109--111}
\egapar{2.5.1}
设 $A$ 为阿廷环，$X$ 为 $Y=\Spec A$ 上的射影概形。对每个凝聚 $\OO_X$-模 $\F$，各 $\mathrm H^i(X,\F)$ 为有限型 $A$-模（2.2.1），故因 $A$ 阿廷而具有有限长度。此外，只有有限多个上同调模非零；因此整数
\begin{equation}\tag{2.5.1.1}
\chi_A(\F)=\sum_{i\geq0}(-1)^i\len_A\mathrm H^i(X,\F)
\end{equation}
有定义。当 $A$ 是阿廷局部环时，称其为 $\F$（相对于 $A$）的\textbf{Euler--Poincaré 特征}。取 $\F=\OO_X$ 时，称 $\chi_A(\OO_X)$ 为 $X$（相对于 $A$）的\textbf{算术亏格}。

\egastatement{2.5.2}{命题}
若 $0\to\F'\to\F\to\F''\to0$ 是凝聚 $\OO_X$-模正合列，则
\begin{equation}\tag{2.5.2.1}
\chi_A(\F)=\chi_A(\F')+\chi_A(\F'').
\end{equation}
\begin{proof}
三个层的上同调除有限多项外均为零，故某个 $r>0$ 使长正合列写成
\[
\begin{aligned}
0&\to\mathrm H^0(X,\F')\to\mathrm H^0(X,\F)\to\mathrm H^0(X,\F'')\to\mathrm H^1(X,\F')\to\cdots\\
&\to\mathrm H^r(X,\F')\to\mathrm H^r(X,\F)\to\mathrm H^r(X,\F'')\to0.
\end{aligned}
\]
两端为零的有限长度模正合列中，长度的交错和为零（0，11.10.1）。应用此结果即得公式。
\end{proof}
只要已知 $X$ 是拟紧 $A$-概形，而且对每个凝聚模 $\F$，各 $A$-模 $\mathrm H^i(X,\F)$ 有限型，这一结果同样适用（1.4.12）。

\egastatement{2.5.3}{定理}
设 $A$ 为阿廷局部环，$X$ 为 $Y=\Spec A$ 上的射影概形，$\LL$ 为相对于 $Y$ 很丰沛的可逆 $\OO_X$-模，$\F$ 为凝聚模；令 $\F(n)=\F\otimes\LL^{\otimes n}$。则：
\begin{enumerate}[label=(\roman*)]
\item 存在唯一多项式 $P\in\mathbf Q[T]$，使每个 $n\in\Z$ 都满足 $\chi_A(\F(n))=P(n)$。称 $P$ 为 $\F$ 相对于 $A$ 的\textbf{Hilbert 多项式}。
\item 充分大的 $n$ 满足 $\chi_A(\F(n))=\len_A\Gamma(X,\F(n))$。
\item $\chi_A(\F(n))$ 的最高次项系数非负。
\end{enumerate}
还应补充，第IV章关于维数的部分将证明：该多项式的次数等于 $\F$ 支撑的维数。
\eproof 由（2.2.1），充分大的 $n$ 使所有正次数上同调为零，所以 $\chi_A(\F(n))=\len_A\mathrm H^0(X,\F(n))=\len_A\Gamma(X,\F(n))$，这证明（ii）。于是充分大的 $n$ 使其非负，故（iii）由（i）得出；（i）的唯一性显然，只需证明存在性。

先归结到 $\mm\F=0$，其中 $\mm$ 为 $A$ 的极大理想。某个 $s>0$ 满足 $\mm^s=0$，所以 $\F(n)$ 有有限滤过
$\F(n)\supset\mm\F(n)\supset\cdots\supset\mm^{s-1}\F(n)\supset0$。
由（2.5.2.1）归纳得
\[
\chi_A(\F(n))=\sum_{k=1}^s\chi_A(\mm^{k-1}\F(n)/\mm^k\F(n)).
\]
因商层等于 $\F_k(n)$，其中 $\F_k=\mm^{k-1}\F/\mm^k\F$，上述归结成立。

设 $\mm\F=0$。令 $X'$ 为 $X\to\Spec A$ 下唯一闭点的闭子概形逆像，$j:X'\to X$ 为嵌入；则 $\F=j_*(\F')$，其中 $\F'$ 为凝聚 $\OO_{X'}$-模。$X'$ 在 $K=A/\mm$ 上射影，$\LL'=j^*\LL$ 很丰沛（II，4.4.10），且 $\F(n)=j_*(\F'(n))$，这里 $\F'(n)=\F'\otimes\LL'^{\otimes n}$（$\OI$，5.4.10）。因而 $\chi_A(\F(n))=\chi_K(\F'(n))$（G，II，4.9.1），归结为 $A$ 是域的情形。

又可将 $X$ 看作某个 $P=\Pb_A^r$ 的闭子概形（II，5.5.4，ii）。若 $i:X\to P$ 为嵌入，则同样有 $\chi_A(\F(n))=\chi_A((i_*\F)(n))$，故可设 $X=\Pb_A^r=\Proj S$，$S=A[T_0,\ldots,T_r]$，$A$ 为域。此时 $\F=\wt M$，$M$ 为有限型分次 $S$-模（II，2.7.8）。由 Hilbert 协同关系定理（M，VIII，6.5），存在有限型分次自由模的有限解消
\[
0\to L_q\to L_{q-1}\to\cdots\to L_1\to M\to0.
\]
由于相伴层函子正合（II，2.5.4），以及扭转正合，对每个 $n$ 有
\[
0\to\wt L_q(n)\to\wt L_{q-1}(n)\to\cdots\to\wt L_1(n)\to\wt M(n)\to0.
\]
反复应用加性公式得
\[
\chi_A(\wt M(n))=\sum_{i=1}^q(-1)^{i+1}\chi_A(\wt L_i(n)).
\]
因此可设 $M$ 分次自由有限型，再归结为 $M=S(h)$。由（II，2.5.15），$\wt M(n)=\widetilde{S(n+h)}$，所以只需以下引理。

\egastatement{2.5.3.1}{引理}
设 $A$ 为域，$r>0$，$X=\Pb_A^r$。则对每个 $n\in\Z$，有
\[
\chi_A(\OO_X(n))=\binom{n+r}{r}.
\]
\begin{proof}
当 $n\geq0$ 时，特征等于 $\mathrm H^0(X,\OO_X(n))$ 的维数，即总次数为 $n$ 的单项式数 $\binom{n+r}{r}$（2.1.12）。当 $n\leq-r-1$ 时，特征为 $(-1)^r\len_A\mathrm H^r(X,\OO_X(n))$。写 $n=-r-h$，则由（2.1.12），此维数是满足 $p_i>0$、$\sum_{i=0}^rp_i=r+h$ 的序列数，也就是满足 $q_i\geq0$、$\sum_iq_i=h-1$ 的序列数，等于 $\binom{h+r-1}{r}=(-1)^r\binom{n+r}{r}$。最后，$-r\leq n<0$ 时二项式为零，所有上同调也为零（2.1.12），引理得证。
\end{proof}
这完成（2.5.3）的证明。\finishproof

\egastatement{2.5.4}{推论}
设 $A$ 为阿廷局部环，$S$ 为由 $S_1$ 生成的有限型分次 $A$-代数，$M$ 为有限型分次 $S$-模，$X=\Proj S$。则充分大的 $n$ 满足 $\chi_A(\wt M(n))=\len_AM_n$。
\begin{proof}
由（2.3.1），充分大的 $n$ 使 $M_n\simeq\Gamma(X,\wt M(n))$。
\end{proof}

\subsection{应用：丰沛性判据}
\src{111--115}
\egastatement{2.6.1}{命题}
设 $Y$ 为诺特预概形，$f:X\to Y$ 为固有态射，$\LL$ 为可逆 $\OO_X$-模。以下条件等价：
\begin{enumerate}[label=\alph*)]
\item $\LL$ 相对于 $f$ 丰沛。
\item 对每个凝聚 $\OO_X$-模 $\F$，存在整数 $N$，使 $n\geq N$ 时，对每个 $q>0$ 有 $\mathrm R^qf_*(\F\otimes\LL^{\otimes n})=0$。
\item 对每个凝聚理想 $\J\subset\OO_X$，存在 $N$，使 $n\geq N$ 时 $\mathrm R^1f_*(\J\otimes\LL^{\otimes n})=0$。
\end{enumerate}
\begin{proof}
由（2.2.1，ii），a）推出 b）；b）推出 c）显然，只需证明 c）推出 a）。可设 $Y$ 仿射（II，4.6.4），并证明 $\LL$ 丰沛。由（II，4.5.2），只需证明：当 $h$ 遍历各 $\LL^{\otimes n}$（$n>0$）的全局截面时，其中仿射的 $X_h$ 覆盖 $X$。

对每个闭点 $x\in X$ 及其仿射开邻域 $U$，证明存在 $n$ 和 $h\in\Gamma(X,\LL^{\otimes n})$，使 $x\in X_h\subset U$；此时 $X_h$ 仿射（I，1.3.6），这些开集的并包含所有闭点，因 $X$ 诺特便等于 $X$（I，6.3.7及$\OI$，2.1.3）。令 $\J,\J'$ 分别定义底空间为 $X-U$ 及 $(X-U)\cup\{x\}$ 的约化闭子预概形（I，5.2.1）。则它们凝聚，$\J'\subset\J$，而 $\J''=\J/\J'$ 凝聚、支撑为 $\{x\}$，且 $\J''_x=k(x)$。

由于 $\LL$ 局部自由，各 $n$ 都给出正合列
\[
0\to\J'\otimes\LL^{\otimes n}\to\J\otimes\LL^{\otimes n}\to\J''\otimes\LL^{\otimes n}\to0.
\]
取足够大的 $n$，假设使 $\mathrm H^1(X,\J'\otimes\LL^{\otimes n})=0$。于是全局截面映射到 $\Gamma(X,\J''\otimes\LL^{\otimes n})$ 满射。取后者在 $x$ 处非零的截面 $g$，提升为 $h\in\Gamma(X,\J\otimes\LL^{\otimes n})\subset\Gamma(X,\LL^{\otimes n})$（$\OI$，5.4.1）。则 $h(x)\ne0$，而 $z\notin U$ 时 $h(z)=0$，所求结论成立。
\end{proof}

\egastatement{2.6.2}{命题}
设 $Y$ 为诺特预概形，$f:X\to Y$ 为有限型态射，$g:X'\to X$ 为有限满射态射，$\LL$ 为可逆 $\OO_X$-模，$\LL'=g^*\LL$。假设存在 $X$ 的子集 $Z$ 在 $Y$ 上固有（II，5.4.10），使每个 $x\in X-Z$ 满足以下两条件之一：$X$ 在 $x$ 处正规；或 $(g_*\OO_{X'})_x$ 为自由 $\OO_{X,x}$-模。则 $\LL$ 相对于 $f$ 丰沛，当且仅当 $\LL'$ 相对于 $f\circ g$ 丰沛。

\egapar{2.6.2.1}
$g$ 仿射，故必要性由（II，5.1.12）成立。为证充分性，可设 $Y$ 仿射（II，4.6.4），且可归结到 $X$ 约化。事实上，令 $j:X_1=X_{\red}\to X$ 为嵌入，$X'_1=X'\times_XX_1$，得到交换图
\begin{equation}\tag{2.6.2.2}
\begin{tikzcd}
X'\arrow[d,"g"']&X'_1\arrow[l,"j'"']\arrow[d,"g_1"]\\
X&X_1\arrow[l,"j"]
\end{tikzcd}
\end{equation}
$f\circ j$ 有限型，$g_1$ 有限。$j'$ 为闭浸入，所以 $\LL'$ 丰沛推出 $j'^*\LL'$ 丰沛（II，5.1.12及I，4.3.2）。令 $Z_1=j^{-1}(Z)$，仍在 $Y$ 上固有。$X$ 在 $x$ 处正规时，$X_{\red}$ 也如此；若 $(g_*\OO_{X'})_x$ 自由，则由（II，1.5.2），$((g_1)_*\OO_{X'_1})_x$ 也自由。又 $X$ 诺特（I，6.3.7），$j^*\LL$ 丰沛推出 $\LL$ 丰沛（II，4.5.14），而 $j'^*\LL'=g_1^*(j^*\LL)$，归结完成。以下设 $Y$ 仿射、$X$ 约化。

此时（II，6.6.11）的假设成立，故存在约化 $Y$-预概形 $X_2$ 和有限双有理 $Y$-态射 $h:X_2\to X$，使 $h$ 在 $h^{-1}(X-Z)$ 上为到 $X-Z$ 的同构，且 $h^*\LL$ 丰沛。以 $X_2$ 代替 $X'$，只需再假设 $g$ 具有这些性质。仍以 $Z$ 表示底空间为该集合、并在 $Y$ 上固有的子预概形（II，5.4.10）。

\egapar{2.6.2.3}
设 $X_1$ 为 $X$ 的闭子预概形，$j:X_1\to X$ 为嵌入，$X'_1=g^{-1}(X_1)$，仍得到（2.6.2.2）。令 $\LL_1=j^*\LL$、$\LL'_1=j'^*\LL'=g_1^*\LL_1$，则 $\LL'_1$ 丰沛（II，5.1.12），$Z_1=j^{-1}(Z)$ 在 $Y$ 上固有。因此（2.6.2）的假设对 $X_1,\LL_1,g_1,Z_1$ 仍成立。

这使我们可在 $g:g^{-1}(X-Z)\to X-Z$ 为同构的情形，用诺特归纳法（$\OI$，2.2.2）证明命题；由上一步，这已足够。只需证明：若对底空间不等于 $X$ 的每个闭子预概形 $X_1$，结论对 $\LL_1$ 成立，则对 $\LL$ 也成立。

\egapar{2.6.2.4}
令 $\A=\OO_X$、$\B=g_*\OO_{X'}$；则 $\B$ 是 $\R(X)$ 的子 $\A$-代数，也是凝聚 $\A$-模，并在 $X-Z$ 上等于 $\A$。令 $\K$ 为 $\B$ 对 $\A$ 的导子，即满足 $\B\K\subset\A$ 的最大的拟凝聚子 $\A$-模 $\K\subset\A$；也就是 $\B/\A$ 的零化子（$\OI$，5.3.7）。于是 $\B\K=\K$。

在 $g$ 局部为同构的所有点，$\K_x=\A_x$；特别在 $X-Z$ 各点以及各不可约分支泛点的某个邻域如此。令 $Z_1=\Spec(\A/\K)$ 为 $\K$ 定义的闭子预概形，其底空间在 $Z$ 中闭，故仍在 $Y$ 上固有。又 $\B|X-Z_1=\A|X-Z_1$。因此可总是归结到 $Z=\Spec(\A/\K)$，而 $X-Z_1$ 非空，故可假设 $Z\ne X$。

\egapar{2.6.2.5}
把 $X'$ 识别为 $\Spec\B$，令 $\K'=\wt\K$ 为 $\OO_{X'}$ 的凝聚理想，使 $g_*\K'=\K$（II，1.4.1）。则 $Z'=g^{-1}(Z)$ 由 $\K'$ 定义，且等于 $\Spec(\B/\K)$（II，1.4.10）。由于 $h:Z'\to Z$ 有限，$Z'$ 也在 $Y$ 上固有（II，6.1.11及5.4.2，ii）。

要证明，对每个 $x\in X$ 及其开邻域 $U$，存在某个 $\LL^{\otimes n}$（$n>0$）的全局截面 $s$ 使 $x\in X_s\subset U$（II，4.5.2）。分两种情形。

\textbf{1° $x\in X-Z$。}可设 $U\subset X-Z$，所以 $g^{-1}(U)$ 不交 $Z'$。由 $\LL'$ 丰沛，存在 $n>0$ 和全局截面 $s'$，使 $x'=g^{-1}(x)\in X'_{s'}\subset g^{-1}(U)$。还可取 $n$ 使 $\K'\otimes\LL'^{\otimes n}$ 由全局截面生成（II，4.5.5）；因 $\K'_{x'}=\OO_{X',x'}$，存在这样的截面 $s''$ 在 $x'$ 处非零。将它乘以 $s'$，相当于将 $n$ 换为 $2n$，便可使 $x'\in X'_{s''}\subset g^{-1}(U)$。由（$\OI$，5.4.10），
\[
\Gamma(X,\K\otimes\LL^{\otimes n})\simeq\Gamma(X',\K'\otimes\LL'^{\otimes n}).
\]
与 $s''$ 对应的截面 $s$ 显然满足要求。

\textbf{2° $x\in Z$。}令 $\J$ 为定义底空间 $X-U$ 的约化闭子预概形的凝聚理想；在 $\B$ 中考虑凝聚理想 $\J\B$ 与 $\J\K=\J(\B\K)=\K(\J\B)$，有嵌入图
\begin{equation}\tag{2.6.2.6}
\begin{tikzcd}[column sep=large]
\J\B\arrow[r]&\B\\
\J\arrow[u]\arrow[r]&\A\arrow[u]\\
\J\K\B=\J\K\arrow[u]\arrow[r]&\K\arrow[u].
\end{tikzcd}
\end{equation}
令 $\J'=(\J\B)^\sim$，则 $g_*\J'=\J\B$，$\J'\K'=(\J\K\B)^\sim$，故
$\J'/\J'\K'=(\J\B/\J\K\B)^\sim$（II，1.4.4）。在不交 $Z$ 的开集 $V$ 上，$\J|V=\J\K|V$，所以此商层支撑包含于 $Z'$。$Z'$ 在 $Y$ 上固有，应用（2.2.4）可知，充分大的 $n$ 使典范映射
\[
\Gamma(X',\J'\otimes\LL'^{\otimes n})\to
\Gamma(X',(\J'/\J'\K')\otimes\LL'^{\otimes n})
\]
满射。由（$\OI$，5.4.10），因此
\[
\Gamma(X,\J\B\otimes\LL^{\otimes n})\to
\Gamma(X,(\J\B/\J\K\B)\otimes\LL^{\otimes n})
\]
也满射。

令 $i:Z\to X$、$i':Z'\to X'$ 为嵌入，有交换图
\[
\begin{tikzcd}
X'\arrow[d,"g"']&Z'\arrow[l,"i'"']\arrow[d,"h"]\\
X&Z\arrow[l,"i"]
\end{tikzcd}
\]
令 $\M=i^*\LL$、$\M'=i'^*\LL'$；$\M'$ 丰沛且等于 $h^*\M$，由 $Z\ne X$ 及归纳假设，$\M$ 丰沛。因此充分大的 $n$ 使 $i^*(\J/\J\K)\otimes\M^{\otimes n}$ 由 $Z$ 上的全局截面生成。因 $\J/\J\K=i_*i^*(\J/\J\K)$，再由（$\OI$，5.4.10），存在 $(\J/\J\K)\otimes\LL^{\otimes n}$ 的全局截面 $s$，使 $s(x)\ne0$；这里 $\J_x=\A_x$，而 $\K_x\ne\A_x$。

图（2.6.2.6）说明，$s$ 也为 $(\J\B/\J\K\B)\otimes\LL^{\otimes n}$ 的全局截面，因此它是 $\J\B\otimes\LL^{\otimes n}$ 的某个全局截面 $t$ 的典范像。这个像属于 $(\J\otimes\LL^{\otimes n})/(\J\K\otimes\LL^{\otimes n})$，故由同一嵌入图，$t$ 实为 $\J\otimes\LL^{\otimes n}$ 的截面，因而也是 $\LL^{\otimes n}$ 的截面。已经知道 $t(x)\ne0$；又按 $\J$ 的定义，在 $X-U=\Supp(\OO_X/\J)$ 上 $t(y)=0$。于是 $x\in X_t\subset U$，证明完成。\finishproof

\egastatement{2.6.3}{注}
当 $X$ 在 $Y$ 上固有时，可借助（2.6.1）及引理（II，6.7.1.1），仿照 Chevalley 定理（II，6.7.1）的论证，以更简单的方法证明（2.6.2）。\par

% ---- END c3_s02c.tex ----
% ---- BEGIN c3_s03.tex ----
\section{固有态射的有限性定理}
\subsection{分解引理}
\src{115--116}
\egastatement{3.1.1}{定义}
设 $\mathbf K$ 是阿贝尔范畴。若其对象集合的子集 $\mathbf K'$ 满足 $0\in\mathbf K'$，并且对于 $\mathbf K$ 中任意正合列 $0\to A'\to A\to A''\to0$，只要 $A,A',A''$ 中两个对象属于 $\mathbf K'$，第三个对象也属于 $\mathbf K'$，就称 $\mathbf K'$ 是\textbf{正合的}。
\egastatement{3.1.2}{定理}
设 $X$ 为诺特预概形，$\mathbf K$ 为凝聚 $\OO_X$-模构成的阿贝尔范畴。设 $\mathbf K'$ 是 $\mathbf K$ 的正合子集，$X'$ 是 $X$ 底空间的闭子集。假设对于 $X'$ 的每个不可约闭子集 $Y$，若 $y$ 为其泛点，则存在 $\OO_X$-模 $\G\in\mathbf K'$，使 $\G_y$ 是一维 $k(y)$-向量空间。那么，支集包含在 $X'$ 中的每个凝聚 $\OO_X$-模都属于 $\mathbf K'$；特别地，若 $X'=X$，则 $\mathbf K'=\mathbf K$。
\eproof
对 $X'$ 的闭子集 $Y$，考虑性质 $P(Y)$：支集包含在 $Y$ 中的每个凝聚 $\OO_X$-模都属于 $\mathbf K'$。由诺特归纳原理（$\OI$，2.2.2），只须证明：若 $Y$ 是 $X'$ 的闭子集，并且对 $Y$ 的每个真闭子集 $Y'$，$P(Y')$ 都成立，那么 $P(Y)$ 成立。

取支集包含在 $Y$ 中的 $\F\in\mathbf K$，证明 $\F\in\mathbf K'$。仍以 $Y$ 表示以此集合为底空间的约化闭子预概形（I，5.2.1），设其定义凝聚理想为 $\J$。存在 $n>0$ 使 $\J^n\F=0$（I，9.3.4）。对 $1\le k\le n$，有凝聚 $\OO_X$-模的正合列
\[
0\to\J^{k-1}\F/\J^k\F\to\F/\J^k\F\to\F/\J^{k-1}\F\to0
\]
（$\OI$，5.3.6、5.3.3）。由于 $\mathbf K'$ 正合，对 $k$ 归纳，只须证明各个 $\F_k=\J^{k-1}\F/\J^k\F$ 属于 $\mathbf K'$。因此可增加假设 $\J\F=0$；等价地，$\F=j_*(j^*\F)$，其中 $j:Y\to X$ 为典范嵌入。分两种情形。

\textit{a）$Y$ 可约。} 写成 $Y=Y'\cup Y''$，其中 $Y',Y''$ 均为 $Y$ 的真闭子集。仍用 $Y',Y''$ 表示相应约化闭子预概形，定义凝聚理想分别为 $\J',\J''$。令
\[
\F'=\F\otimes_{\OO_X}(\OO_X/\J'),\qquad
\F''=\F\otimes_{\OO_X}(\OO_X/\J'').
\]
典范同态给出 $u:\F\to\F'\oplus\F''$。对 $z\notin Y'\cap Y''$，$u_z$ 为双射。事实上，$\J'\cap\J''=\J$；这是局部问题，由（I，5.2.1、1.1.5）得出。若 $z\notin Y''$，则 $\J'_z=\J_z$，所以 $\F'_z=\F_z$、$\F''_z=0$；若 $z\notin Y'$，论证相同。故 $u$ 的核与余核是支集包含在 $Y'\cap Y''$ 中的凝聚模（$\OI$，5.3.4），由归纳假设属于 $\mathbf K'$。同理，$\F',\F''\in\mathbf K'$，故 $\F'\oplus\F''\in\mathbf K'$。再考虑正合列
\[
0\to\Im u\to\F'\oplus\F''\to\Coker u\to0,
\qquad 0\to\Ker u\to\F\to\Im u\to0,
\]
由正合性得到结论。

\textit{b）$Y$ 不可约。} 此时子预概形 $Y$ 是整的。若 $y$ 是其泛点，则 $(\OO_Y)_y=k(y)$；因 $j^*\F$ 凝聚，$\F_y=(j^*\F)_y$ 是某个有限维数 $m$ 的 $k(y)$-向量空间。由假设，存在凝聚 $\OO_X$-模 $\G$（其支集必为 $Y$），使 $\G_y$ 是一维 $k(y)$-向量空间。\transnote{原文陈述只写了 $\G_y$ 的条件，证明此处却直接使用 $\Supp(\G)=Y$，后文（3.1.3）也使用同一支集条件。泛点茎的条件本身并不推出整个支集恰为 $Y$；这里保留原文陈述与证明的区别，使用此分解引理时须另核对所选 $\G$ 的支集。}
于是有 $k(y)$-同构 $(\G_y)^m\simeq\F_y$，亦为 $\OO_{X,y}$-同构。由凝聚性，存在 $y$ 在 $X$ 中的开邻域 $W$ 和同构 $\G^m|W\simeq\F|W$（$\OI$，5.2.7）。其图像 $\HH$ 是 $(\G^m\oplus\F)|W$ 的凝聚子模，典范地同构于 $\G^m|W$ 及 $\F|W$。因 $\G^m,\F$ 的支集均为 $Y$，存在 $\G^m\oplus\F$ 的凝聚子模 $\HH_0$，在 $W$ 上诱导 $\HH$，在 $X-Y$ 上诱导零（I，9.4.7）。典范投影的限制
\[
v:\HH_0\to\G^m,\qquad w:\HH_0\to\F
\]
在 $W$ 及 $X-Y$ 上均为同构，故其核与余核的支集都包含在真闭子集 $Y-(Y\cap W)$ 中，因而属于 $\mathbf K'$。又因 $\G\in\mathbf K'$、$\mathbf K'$ 正合，有 $\G^m\in\mathbf K'$。由正合性依次得到 $\HH_0\in\mathbf K'$、$\F\in\mathbf K'$。\finishproof
\egastatement{3.1.3}{推论}
再假设 $\mathbf K'$ 具有下述性质：属于 $\mathbf K'$ 的凝聚 $\OO_X$-模的每个凝聚直和因子仍属于 $\mathbf K'$。那么，把（3.1.2）中的条件“$\G_y$ 是一维 $k(y)$-向量空间”换成“$\G_y\ne0$”（亦即 $\Supp(\G)=Y$），其结论仍成立。
\eproof
只需修改（3.1.2）的情形 b）。这时 $\G_y$ 是某个维数 $q>0$ 的 $k(y)$-向量空间，故 $(\G_y)^m\simeq(\F_y)^q$。同样的论证得到 $\F^q\in\mathbf K'$；由附加假设可知 $\F\in\mathbf K'$。\finishproof

\subsection{有限性定理：通常概形的情形}
\src{116--118}
\egastatement{3.2.1}{定理}
设 $Y$ 是局部诺特预概形，$f:X\to Y$ 是固有态射。对任意凝聚 $\OO_X$-模 $\F$ 和任意 $q\ge0$，$R^qf_*(\F)$ 都是凝聚 $\OO_Y$-模。
\eproof
问题在 $Y$ 上是局部的，故可设 $Y$ 诺特，于是 $X$ 亦诺特（I，6.3.7）。使定理结论成立的凝聚 $\OO_X$-模构成凝聚模范畴 $\mathbf K$ 的正合子集 $\mathbf K'$。事实上，若 $0\to\F'\to\F\to\F''\to0$ 正合，例如 $\F',\F''\in\mathbf K'$，则上同调正合列
\[
R^{q-1}f_*\F''\to R^qf_*\F'\to R^qf_*\F\to R^qf_*\F''\to R^{q+1}f_*\F'
\]
两端四项凝聚，故中项凝聚（$\OI$，5.3.4、5.3.3）。另外两种情形同理。若 $\F'$ 是 $\F\in\mathbf K'$ 的凝聚直和因子，则 $R^qf_*\F'$ 是 $R^qf_*\F$ 的直和因子（G，II，4.4.4），故有限型；又由（1.4.10）它拟凝聚，$Y$ 诺特，故它凝聚。

由（3.1.3），只须在 $X$ 不可约、泛点为 $x$ 时，构造 $\F\in\mathbf K'$ 且 $\F_x\ne0$。确实，可把这一结果用于 $X$ 的任意不可约闭子预概形 $Z$：若 $j:Z\to X$ 为嵌入，则 $f\circ j$ 固有（II，5.4.2）；若 $\G$ 是支集为 $Z$ 的凝聚 $\OO_Z$-模，则 $j_*\G$ 凝聚，而且
\[
R^q(f\circ j)_*(\G)=R^qf_*(j_*\G)
\]
（G，II，4.9.1），因此满足（3.1.3）的应用条件。

由 Chow 引理（II，5.6.2），存在不可约预概形 $X'$ 以及射影满态射 $g:X'\to X$，使 $f\circ g:X'\to Y$ 射影。取关于 $g$ 丰沛的可逆模 $\LL$（II，5.3.1）。由射影态射基本定理（2.2.1），可取 $n$，使
\[
\F=g_*(\OO_{X'}(n))
\]
凝聚，并使 $R^qg_*(\OO_{X'}(n))=0$ 对所有 $q>0$ 成立。又因充分大的 $n$ 使 $g^*g_*(\OO_{X'}(n))\to\OO_{X'}(n)$ 满射，可使 $\F_x\ne0$（II，3.4.7）。由于 $f\circ g$ 射影、$Y$ 诺特，$R^n(f\circ g)_*(\OO_{X'}(n))$ 凝聚（2.2.1）。Leray 谱序列
\[
E_2^{pq}=R^pf_*(R^qg_*(\OO_{X'}(n)))
\]
收敛到该复合态射的高阶直接像；由上述消失，它退化，故 $E_2^{p0}=R^pf_*\F$ 同构于 $R^p(f\circ g)_*(\OO_{X'}(n))$（0，11.1.6）。定理得证。\finishproof
\egastatement{3.2.2}{推论}
设 $Y$ 局部诺特。对任意固有态射 $f:X\to Y$，任意凝聚 $\OO_X$-模的直接像均为凝聚 $\OO_Y$-模。
\egastatement{3.2.3}{推论}
设 $A$ 为诺特环，$X$ 为 $A$ 上固有的概形。则每个凝聚 $\OO_X$-模 $\F$ 的 $H^p(X,\F)$ 都是有限型 $A$-模；而且存在 $r>0$，使对每个这样的 $\F$ 及每个 $p>r$ 都有 $H^p(X,\F)=0$。
\eproof
第二断言已由（1.4.12）证明；第一断言由有限性定理（3.2.1）及（1.4.11）得出。特别地，当 $X$ 是域 $k$ 上的固有代数概形时，这些上同调群都是有限维 $k$-向量空间。\finishproof
\egastatement{3.2.4}{推论}
设 $Y$ 局部诺特，$f:X\to Y$ 为有限型态射。若凝聚 $\OO_X$-模 $\F$ 的支集在 $Y$ 上固有（II，5.4.10），则 $R^qf_*\F$ 都凝聚。
\eproof
可设 $Y$ 诺特，因而 $X$ 诺特。以 $\Supp(\F)$ 为底空间的每个闭子预概形 $Z$ 都在 $Y$ 上固有。若 $j:Z\to X$ 为嵌入，可取 $Z$ 使 $\F=j_*\G$，其中 $\G=j^*\F$ 为凝聚 $\OO_Z$-模（I，9.3.5）。由（1.3.4），$R^qf_*\F=R^q(f\circ j)_*\G$，故应用（3.2.1）即可。\finishproof

\subsection{有限性定理的推广（通常概形）}
\src{118}
\egastatement{3.3.1}{命题}
设 $Y$ 为诺特预概形，$\SSS$ 为非负分次、拟凝聚、有限型的 $\OO_Y$-代数，$Y'=\Proj(\SSS)$，$g:Y'\to Y$ 为结构态射。设 $f:X\to Y$ 固有，$\SSS'=f^*\SSS$，$\M=\bigoplus_k\M_k$ 是拟凝聚有限型分次 $\SSS'$-模。那么对每个 $p$，
\[
R^pf_*\M=\bigoplus_kR^pf_*\M_k
\]
是有限型分次 $\SSS$-模。若 $\SSS$ 由 $\SSS_1$ 生成，则对每个 $p$ 存在 $k_p$，使 $k\ge k_p$、$r\ge0$ 时
\begin{equation}\tag{3.3.1.1}
R^pf_*(\M_{k+r})=\SSS_rR^pf_*(\M_k).
\end{equation}
\eproof
第一断言是将（2.4.1，i）中的“射影”换成“固有”。该证明中只用射影性来保证 $R^pf'_*(\wt\M)$ 凝聚；现在 $f'$ 固有（II，5.4.2，iii），故由有限性定理（3.2.1），原证明仍全部适用。

对于第二断言，在 $Y$ 的某个有限仿射开覆盖 $(U_\alpha)$ 上，两边的限制在 $k\ge k_{p,\alpha}$ 时相等（II，2.1.6，ii）。取这些 $k_{p,\alpha}$ 的最大值即可。\finishproof
\egastatement{3.3.2}{推论}
设 $A$ 为诺特环，$\mm$ 为其理想，$X$ 为固有 $A$-概形，$\F$ 为凝聚 $\OO_X$-模。对每个 $p\ge0$，$\bigoplus_{k\ge0}H^p(X,\mm^k\F)$ 是环 $S=\bigoplus_{k\ge0}\mm^k$ 上的有限型模。特别地，存在 $k_p\ge0$，使 $k\ge k_p$、$r\ge0$ 时
\begin{equation}\tag{3.3.3.1}
H^p(X,\mm^{k+r}\F)=\mm^rH^p(X,\mm^k\F).
\end{equation}
\eproof
取 $Y=\Spec(A)$、$\SSS=\wt S$、$\M_k=\mm^k\F$，应用上述命题及（1.4.11）即可。\transnote{此式的原版编号为（3.3.3.1），本稿保留；原证明写“应用（3.3.2）”，所需命题实际上是（3.3.1）。}
须记住，上述直和的 $S$-模结构由下列映射定义：对 $a\in\mm^r$，乘法 $\mm^k\F\to\mm^{k+r}\F$ 诱导上同调映射 $H^p(X,\mm^k\F)\to H^p(X,\mm^{k+r}\F)$（2.4.1）。\finishproof

\subsection{有限性定理：形式概形的情形}
\src{119--122}
本小节的结果（定义（3.4.1）除外）不在本章后文使用。
\egapar{3.4.1}
设 $\mathfrak X,\mathfrak S$ 为局部诺特形式预概形（I，10.4.2），$f:\mathfrak X\to\mathfrak S$ 为形式预概形态射。若满足以下条件，就称 $f$ \textbf{固有}：一是 $f$ 有限型（I，10.13.3）；二是取 $\mathfrak S$ 的定义理想 $\K$，令
\[
\J=f^*(\K)\OO_{\mathfrak X},\qquad
X_0=(\mathfrak X,\OO_{\mathfrak X}/\J),\qquad
S_0=(\mathfrak S,\OO_{\mathfrak S}/\K),
\]
则由 $f$ 诱导的 $f_0:X_0\to S_0$ 固有。

这一定义不依赖 $\K$ 的选择。事实上，若另一取法满足 $\K'\subset\K$，令 $\J'=f^*\K'\OO_{\mathfrak X}$，并令 $X_1,S_1$ 为相应商，则有交换图
\[
\begin{tikzcd}
X_0\arrow[r,"f_0"]\arrow[d,"i"']&S_0\arrow[d,"j"]\\
X_1\arrow[r,"f_1"]&S_1
\end{tikzcd}
\]
其中 $i,j$ 为满浸入，故 $f_0$ 与 $f_1$ 同时固有（II，5.4.5）。此外，若令 $X_n=(\mathfrak X,\OO_{\mathfrak X}/\J^{n+1})$、$S_n=(\mathfrak S,\OO_{\mathfrak S}/\K^{n+1})$，则 $f_0$ 固有即推出每个 $f_n:X_n\to S_n$ 固有（II，5.4.6）。

若 $g:Y\to Z$ 为局部诺特通常预概形的固有态射，$Z'\subset Z$、$Y'\subset Y$ 为闭子集且 $g(Y')\subset Z'$，则其到完备化的延拓 $\widehat g:Y_{/Y'}\to Z_{/Z'}$（I，10.9.1）是固有形式预概形态射；这是定义及（II，5.4.5）的结果。\transnote{原文这里仅写 $g(Y')\subset Z'$。是否得到（3.4.1）所要求的“有限型”形式态射，还须核对定义理想所给出的拓扑；此处未补造额外假设，应用时应与（I，10.13.3）对照。}

若 $f:\mathfrak X\to\mathfrak S$ 有限型，称底空间的子集 $Z$ \textbf{在 $\mathfrak S$ 上固有}，或\textbf{关于 $f$ 固有}，是指把它看成 $X_0$ 的子集时，它在 $S_0$ 上固有（II，5.4.10）。由定义，（II，5.4.10）关于固有子集的全部性质仍适用于此。
\egastatement{3.4.2}{定理}
设 $\mathfrak X,\mathfrak Y$ 为局部诺特形式预概形，$f:\mathfrak X\to\mathfrak Y$ 固有。对每个凝聚 $\OO_{\mathfrak X}$-模 $\F$，$R^qf_*\F$ 对所有 $q\ge0$ 均为凝聚 $\OO_{\mathfrak Y}$-模。

取 $\mathfrak Y$ 的定义理想 $\J$，令 $\K=f^*\J\OO_{\mathfrak X}$，并考虑
\begin{equation}\tag{3.4.2.1}
\F_k=\F\otimes_{\OO_{\mathfrak Y}}(\OO_{\mathfrak Y}/\J^{k+1})
=\F/\K^{k+1}\F\qquad(k\ge0).
\end{equation}
它们构成拓扑模的射影系，且 $\F=\plim_k\F_k$（I，10.11.3）。定理还将表明各 $R^qf_*\F$ 凝聚，故自然带有拓扑模结构（I，10.11.6）；$R^qf_*\F_k$ 亦如此。典范商映射给出连续同态 $R^qf_*\F\to R^qf_*\F_k$，进而给出函子性的连续同态
\begin{equation}\tag{3.4.2.2}
R^qf_*\F\longrightarrow\plim_kR^qf_*\F_k.
\end{equation}
我们将同时证明下述推论。
\egastatement{3.4.3}{推论}
每个（3.4.2.2）都是拓扑同构。若 $\mathfrak Y$ 诺特，则射影系 $(R^qf_*(\F/\K^{k+1}\F))_k$ 满足（ML）条件（0，13.1.1）。

先处理 $\mathfrak Y$ 为诺特仿射形式概形的情形。
\egastatement{3.4.4}{推论}
在（3.4.2）条件下，再设 $\mathfrak Y=\Spf(A)$，其中 $A$ 为诺特进制环。取定义理想 $\frj\subset A$，令 $\F_k=\F/\frj^{k+1}\F$。则 $H^n(\mathfrak X,\F)$ 是有限型 $A$-模，射影系 $(H^n(\mathfrak X,\F_k))_{k\ge0}$ 对每个 $n$ 满足（ML）。令
\begin{equation}\tag{3.4.4.1}
N_k^n=\Ker(H^n(\mathfrak X,\F)\to H^n(\mathfrak X,\F_k));
\end{equation}
由上同调正合列，它亦为 $H^n(\mathfrak X,\frj^{k+1}\F)\to H^n(\mathfrak X,\F)$ 的像。这些子模给出 $\frj$-良滤过（0，13.7.7）。最后，典范同态
\begin{equation}\tag{3.4.4.2}
H^n(\mathfrak X,\F)\longrightarrow\plim_kH^n(\mathfrak X,\F_k)
\end{equation}
是拓扑同构，左端取 $\frj$-进制拓扑，各 $H^n(\mathfrak X,\F_k)$ 取离散拓扑。
\eproof
令
\begin{equation}\tag{3.4.4.3}
S=\gr(A)=\bigoplus_{k\ge0}\frj^k/\frj^{k+1},\qquad
\M=\gr(\F)=\bigoplus_{k\ge0}\frj^k\F/\frj^{k+1}\F.
\end{equation}
$\frj^\triangle$ 是 $\mathfrak Y$ 的定义理想（I，10.3.1）。令 $\K=f^*(\frj^\triangle)\OO_{\mathfrak X}$，$X_0=(\mathfrak X,\OO_{\mathfrak X}/\K)$，$Y_0=(\mathfrak Y,\OO_{\mathfrak Y}/\frj^\triangle)=\Spec(A_0)$，$A_0=A/\frj$。各 $\M_k=\frj^k\F/\frj^{k+1}\F$ 是凝聚 $\OO_{X_0}$-模（I，10.11.3）。再考虑非负分次拟凝聚代数
\begin{equation}\tag{3.4.4.4}
\SSS=\OO_{X_0}\otimes_{A_0}S
=\gr(\OO_{\mathfrak X})=\bigoplus_{k\ge0}\K^k/\K^{k+1}.
\end{equation}
\transnote{本式原文写等号。一般由标量乘法首先得到的是从 $\OO_{X_0}\otimes_{A_0}\gr(A)$ 到 $\gr(\OO_{\mathfrak X})$ 的典范满同态，未经平坦性等附加条件不能直接视为同构。下文局部计算的 $M'=S'\otimes N_0$ 同样须注意：用于有限生成性的典范满同态已足够；本稿不把原文的等号静默换成同构定理。}
$\F$ 有限型推出 $\M$ 是有限型分次 $\SSS$-模。此为局部问题，可设 $\mathfrak X=\Spf(B)$，$B$ 诺特进制，$\F=N^\triangle$，$N$ 为有限型 $B$-模（I，10.10.5）。令 $B_0=B/\frj B$、$N_0=N/\frj N$。原文将 $\SSS,\M$ 分别写成与
\[
S'=\bigoplus_k(\frj^k/\frj^{k+1})\otimes_{A_0}B_0,
\qquad M'=\bigoplus_k(\frj^k/\frj^{k+1})\otimes_{A_0}N_0
\]
相伴的层。于是 $M'=S'\otimes_{B_0}N_0$；$N_0$ 有限型，故 $M'$ 在 $S'$ 上有限型，得到所需结论（I，1.3.13）。

由于 $f_0:X_0\to Y_0$ 固有，由（3.3.1）及（1.4.11），每个 $\bigoplus_kH^n(X_0,\M_k)$ 都是有限型分次 $S$-模。故对于带自然滤过 $A$-模结构的严格射影系 $(\F/\frj^{k+1}\F)_{k\ge0}$，（0，13.7.7）的条件 $(F_n)$ 对所有 $n\ge0$ 成立。由该命题可得：（一）$(H^n(\mathfrak X,\F_k))_k$ 满足（ML）；（二）$H'^n=\plim_kH^n(\mathfrak X,\F_k)$ 为有限型 $A$-模；（三）由 $H'^n\to H^n(\mathfrak X,\F_k)$ 的核给出的滤过是 $\frj$-良的。

令 $X_k=(\mathfrak X,\OO_{\mathfrak X}/\K^{k+1})$。$\F_k$ 为凝聚 $\OO_{X_k}$-模。$X_0$ 中任意仿射开集 $U$ 在每个 $X_k$ 中仍仿射（I，5.1.9），故 $H^n(U,\F_k)=0$ 对 $n>0$ 成立，而 $\Gamma(U,\F_h)\to\Gamma(U,\F_k)$ 在 $h\ge k$ 时满射（I，1.3.9）。于是满足（0，13.3.2）的条件；由（0，13.3.1），
\[
H'^n\simeq H^n(\mathfrak X,\plim_k\F_k)=H^n(\mathfrak X,\F).
\]
（3.4.4）得证。\finishproof
\egapar{3.4.5}
回到（3.4.2）、（3.4.3）的证明。先令 $\mathfrak Y=\Spf(A)$。对 $g\in A$，把（3.4.4）用于开形式概形 $\mathfrak Y_g=D(g)=\Spf(A_{\{g\}})$ 及其逆像。$\mathfrak Y_g$ 也是 $Y_k=(\mathfrak Y,\OO_{\mathfrak Y}/(\frj^\triangle)^{k+1})$ 中的仿射开集，故
\[
H^n(f^{-1}(\mathfrak Y_g),\F_k)=\Gamma(\mathfrak Y_g,R^nf_*\F_k)
\]
（1.4.11）。由（3.4.4），$H^n(f^{-1}(\mathfrak Y_g),\F)\to\plim_k\Gamma(\mathfrak Y_g,R^nf_*\F_k)$ 是同构。又截面与射影极限交换（$\OI$，3.2.6），而 $R^nf_*\F$ 与此上同调预层相伴，故（3.4.2.2）为双射。

进一步证明
\begin{equation}\tag{3.4.5.1}
R^nf_*\F\simeq\bigl(H^n(\mathfrak X,\F)\bigr)^\triangle.
\end{equation}
由 $R^nf_*\F_k$ 在 $Y_k$ 上凝聚及（1.4.13），
\[
\Gamma(\mathfrak Y_g,R^nf_*\F_k)
=\Gamma(\mathfrak Y,R^nf_*\F_k)_g
=H^n(\mathfrak X,\F_k)_g.
\]
该上同调射影系满足（ML），极限 $H^n(\mathfrak X,\F)$ 有限型，故由（0，13.7.8）及（I，10.10.8），
\[
\plim_kH^n(\mathfrak X,\F_k)_g
=H^n(\mathfrak X,\F)\otimes_AA_{\{g\}}
=\Gamma\!\left(\mathfrak Y_g,(H^n(\mathfrak X,\F))^\triangle\right).
\]
这就证明了（3.4.5.1）及凝聚性。因此（3.4.2.2）是凝聚模的同构，必为拓扑同构（I，10.11.6）。由 $R^nf_*\F_k=(H^n(\mathfrak X,\F_k))^\triangle$ 还得到其射影系满足（ML）（I，10.10.2）。

一般情形的（3.4.2）及（3.4.3）第一断言都是 $\mathfrak Y$ 上的局部问题，故立即得出。若 $\mathfrak Y$ 诺特，可取有限诺特仿射开覆盖；在每一片上射影系满足（ML），便得到全局的（ML）。同时已证明下述结果。
\egastatement{3.4.6}{推论}
在（3.4.4）的条件下，典范同态
\begin{equation}\tag{3.4.6.1}
H^n(\mathfrak X,\F)\longrightarrow\Gamma(\mathfrak Y,R^nf_*\F)
\end{equation}
为双射。\par

% ---- END c3_s03.tex ----
% ---- BEGIN c3_s04a.tex ----
\section{固有态射的基本定理及其应用}
\subsection{基本定理}
\src{122--129}
\egapar{4.1.1}
设 $X,Y$ 为诺特通常预概形，$f:X\to Y$ 为固有态射，$Y'\subset Y$ 为闭子集，$X'=f^{-1}(Y')$。以 $\mathfrak X=X_{/X'}$、$\mathfrak Y=Y_{/Y'}$ 表示分别沿 $X',Y'$ 的形式完备化（I，10.8.5），以 $\widehat f:\mathfrak X\to\mathfrak Y$ 表示 $f$ 在这些完备化上的延拓（I，10.9.1）。对于凝聚 $\OO_X$-模 $\F$，以 $\widehat\F=\F_{/X'}$ 表示沿 $X'$ 的完备化（I，10.8.4）；它是凝聚 $\OO_{\mathfrak X}$-模（I，10.8.8）。
\egapar{4.1.2}
取凝聚理想 $\J\subset\OO_Y$，使 $\Supp(\OO_Y/\J)=Y'$（I，5.2.1）。则 $\K=f^*\J\OO_X$ 是凝聚理想，且 $\Supp(\OO_X/\K)=X'$（I，4.4.5）。对 $k\ge0$，考虑凝聚模
\[
\F_k=\F\otimes_{\OO_Y}(\OO_Y/\J^{k+1})=\F/\K^{k+1}\F.
\]
由（3.2.1），$R^nf_*\F$、$R^nf_*\F_k$ 凝聚。商同态通过函子性给出
\begin{equation}\tag{4.1.2.1}R^nf_*\F\longrightarrow R^nf_*\F_k.\end{equation}
因为 $\F_k$ 是 $\OO_Y/\J^{k+1}$-模，$R^nf_*\F_k$ 也是该商环层上的模（0，12.2.1），从而得到
\begin{equation}\tag{4.1.2.2}
R^nf_*\F\otimes_{\OO_Y}(\OO_Y/\J^{k+1})\longrightarrow R^nf_*\F_k.
\end{equation}
两端构成射影系，左端的射影极限是 $(R^nf_*\F)_{/Y'}$，记为 $(R^nf_*\F)^\wedge$。这些同态显然相容，取射影极限得到
\begin{equation}\tag{4.1.2.3}
\varphi_n:(R^nf_*\F)^\wedge\longrightarrow\plim_kR^nf_*\F_k.
\end{equation}
（4.1.2.2）作为 $\OO_Y/\J^{k+1}$-模同态，可视为伪离散拓扑 $\OO_{\mathfrak Y}$-模的连续同态（I，10.8.3；$\OI$，3.8.1）；故 $\varphi_n$ 也是拓扑 $\OO_{\mathfrak Y}$-模的连续同态。
\egapar{4.1.3}
设 $i:\mathfrak X\to X$ 为典范赋环空间态射（I，10.8.7），则有交换图
\begin{equation}\tag{4.1.3.1}
\begin{tikzcd}
X_k\arrow[r,"h_k"]\arrow[dr,"i_k"']&\mathfrak X\arrow[d,"i"]\\
&X
\end{tikzcd}
\end{equation}
其中 $X_k$ 为理想 $\K^{k+1}$ 定义的闭子预概形，$i_k$ 为嵌入，$h_k$ 的底空间映射为恒等，环层同态为 $\OO_{\mathfrak X}\to\OO_X/\K^{k+1}$（I，10.5.2）。有典范同构 $\widehat\F=i^*\F$（I，10.8.8），并有
\begin{equation}\tag{4.1.3.2}H^n(X_k,h_k^*\widehat\F)=H^n(X,\F_k),\end{equation}
因为 $\F_k=(i_k)_*(h_k^*\widehat\F)$（G，II，4.9.1）。因此由（0，12.1.3.5）得到的典范同态也可写成
\begin{equation}\tag{4.1.3.3}H^n(\mathfrak X,\widehat\F)\longrightarrow H^n(X,\F_k).\end{equation}
这些同态组成射影系，从而给出
\begin{equation}\tag{4.1.3.4}
\psi_{n,X}:H^n(\mathfrak X,\widehat\F)\longrightarrow\plim_kH^n(X,\F_k).
\end{equation}
对 $Y$ 的仿射开集 $V$，将 $X$ 换成 $f^{-1}(V)$，并利用（1.4.11），得到
\begin{equation}\tag{4.1.3.5}
\psi_{n,V}:H^n(\mathfrak X\cap f^{-1}(V),\widehat\F)
\longrightarrow\plim_k\Gamma(V,R^nf_*\F_k).
\end{equation}
它们与仿射开集间的限制相容，故定义典范层同态
\begin{equation}\tag{4.1.3.6}
\psi_n:R^n\widehat f_*(\widehat\F)\longrightarrow\plim_kR^nf_*\F_k.
\end{equation}
\egapar{4.1.4}
再设 $j:\mathfrak Y\to Y$ 为典范赋环空间态射。$R^nf_*\F$ 凝聚，故 $j^*(R^nf_*\F)\simeq(R^nf_*\F)^\wedge$（I，10.8.8）。于是得到一般赋环空间情形中定义的同态（参看（1.4.15）的证明）
\begin{equation}\tag{4.1.4.1}
\rho_n:(R^nf_*\F)^\wedge=j^*(R^nf_*\F)
\longrightarrow R^n\widehat f_*(i^*\F)=R^n\widehat f_*\widehat\F.
\end{equation}
以下图交换：
\begin{equation}\tag{4.1.4.2}
\begin{tikzcd}[column sep=large]
(R^nf_*\F)^\wedge\arrow[r,"\rho_n"]\arrow[dr,"\varphi_n"']&R^n\widehat f_*\widehat\F\arrow[d,"\psi_n"]\\
&\plim_kR^nf_*\F_k
\end{tikzcd}
\end{equation}
只须验证相应预层图的交换性，故可设 $Y$ 仿射，并验证
\begin{equation}\tag{4.1.4.3}
\begin{tikzcd}[column sep=large]
(H^n(X,\F))^\wedge\arrow[r,"\rho_n"]\arrow[dr,"\varphi_n"']&H^n(\mathfrak X,\widehat\F)\arrow[d,"\psi_{n,X}"]\\
&\plim_kH^n(X,\F_k).
\end{tikzcd}
\end{equation}
而（4.1.3.1）的交换性及（4.1.3）中的关系给出交换图
\[
\begin{tikzcd}[column sep=large]
H^n(X,\F)\arrow[r]\arrow[dr]&H^n(\mathfrak X,\widehat\F)=H^n(\mathfrak X,i^*\F)\arrow[d]\\
&H^n(X,\F_k)=H^n(X_k,i_k^*\F),
\end{tikzcd}
\]
由此立即推出（4.1.4.3）。
\egastatement{4.1.5}{定理}
设 $f:X\to Y$ 为诺特预概形的固有态射，$Y'\subset Y$ 为闭子集，$X'=f^{-1}(Y')$。则对每个凝聚 $\OO_X$-模 $\F$，$R^n\widehat f_*\widehat\F$ 是凝聚 $\OO_{\mathfrak Y}$-模，而（4.1.4.2）中的 $\varphi_n,\psi_n,\rho_n$ 均为拓扑同构。\transnote{原文此处把直接像的模所在环层印为 $\OO_{\mathfrak X}$，应为目标上的 $\OO_{\mathfrak Y}$。}

只须证明 $\varphi_n,\psi_n$ 为同构。由（3.2.1）、（I，10.8.8），$(R^nf_*\F)^\wedge$ 凝聚；上述各同态的双连续性随后由（I，10.11.6）自动得到。
\egastatement{4.1.6}{注记}
（i）令 $\widehat\F_k=\widehat\F/\widehat\K^{k+1}\widehat\F$。立即有 $\F_k=i_*\widehat\F_k$，且（4.1.3.6）正是已在（3.4.2.2）定义的同态
\begin{equation}\tag{4.1.6.1}
R^n\widehat f_*\widehat\F\longrightarrow\plim_kR^n\widehat f_*\widehat\F_k.
\end{equation}
因此 $\psi_n$ 为同构是（3.4.3）的特殊情形。不过，下面将给出直接证明，避免使用该一般结果所依赖的关于谱序列射影极限的细致讨论（0，13.7）。

（ii）在 $\psi_n$ 为同构的条件下，$\varphi_n$ 为同构等价于 $\rho_n=\psi_n^{-1}\varphi_n$ 为同构。定理特别说明，形成 $R^nf$ 与完备化交换；可称它为“代数”理论与“形式”理论的\textbf{第一比较定理}。

先证明其仿射形式。
\egastatement{4.1.7}{推论}
在（4.1.5）的条件下，再设 $Y=\Spec(A)$，$A$ 诺特，$\J=\wt\frj$，$\frj\subset A$ 为理想，故 $\F_k=\F/\frj^{k+1}\F$。则
\begin{equation}\tag{4.1.7.1}
\varphi_n:(H^n(X,\F))^\wedge\longrightarrow\plim_kH^n(X,\F_k)
\end{equation}
为同构，左端表示 $\frj$-预进制拓扑的分离完备化。射影系 $(H^n(X,\F_k))_{k\ge0}$ 对所有 $n$ 满足（ML），且
\begin{equation}\tag{4.1.7.2}
\psi_n:H^n(\mathfrak X,\widehat\F)\longrightarrow\plim_kH^n(X,\F_k)
\end{equation}
为同构。此外，由典范同态 $H^n(X,\F)\to H^n(X,\F_k)$ 的核定义的滤过是 $\frj$-良的（0，13.7.7），而 $\varphi_n$ 为拓扑同构。\origfootnote{以下较原先证明更简单的证明，以及关于 $H^n(X,\F)$ 的滤过的补充，是 J.-P. Serre 告知我们的。}
\eproof
在本证明中固定 $n\ge0$，简记
\begin{equation}\tag{4.1.7.3}H=H^n(X,\F),\qquad H_k=H^n(X,\F_k),\end{equation}
\begin{equation}\tag{4.1.7.4}R_k=\Ker(H\to H_k)\subset H.\end{equation}
上同调正合列
\[
\begin{gathered}
H^n(X,\frj^{k+1}\F)\to H^n(X,\F)\to H^n(X,\F_k)\\
\to H^{n+1}(X,\frj^{k+1}\F)\to H^{n+1}(X,\F)
\end{gathered}
\]
表明 $R_k$ 也是第一箭头的像。令
\begin{equation}\tag{4.1.7.5}
\begin{aligned}
Q_k&=\Ker(H^{n+1}(X,\frj^{k+1}\F)\to H^{n+1}(X,\F))\\
&=\Im(H^n(X,\F_k)\to H^{n+1}(X,\frj^{k+1}\F)).
\end{aligned}
\end{equation}
于是有正合列
\begin{equation}\tag{4.1.7.6}0\to R_k\to H\to H_k\to Q_k\to0.\end{equation}
\egapar{4.1.7.7}
对 $x\in\frj^m$，乘法 $\frj^k\F\to\frj^{k+m}\F$ 诱导
\begin{equation}\tag{4.1.7.8}
\mu_{x,k}:H^n(X,\frj^k\F)\longrightarrow H^n(X,\frj^{k+m}\F).
\end{equation}
令 $S=\bigoplus_{k\ge0}\frj^k$。由（3.3.2），这些乘法在 $E=\bigoplus_{k\ge0}H^n(X,\frj^k\F)$ 上定义有限型分次 $S$-模结构；$S$ 诺特（II，2.1.5）。
\egastatement{4.1.7.9}{引理}
子模 $R_k$ 在 $H$ 上定义 $\frj$-良滤过。
\eproof
先证明
\begin{equation}\tag{4.1.7.10}\frj^mR_k\subset R_{k+m}.\end{equation}
$H$ 中乘以 $x\in\frj^m$ 即为 $\mu_{x,0}$。下图中横箭头为乘以 $x$，竖箭头为典范嵌入：
\[
\begin{tikzcd}
\frj^{k+1}\F\arrow[r,"x"]\arrow[d]&\frj^{k+m+1}\F\arrow[d]\\
\F\arrow[r,"x"]&\F.
\end{tikzcd}
\]
因此相应上同调图也交换：\transnote{原图（4.1.7.11）上方箭头的第二下标印作 $m$；依（4.1.7.8）的源、靶约定，此处写为 $k+1$。}
\begin{equation}\tag{4.1.7.11}
\begin{tikzcd}[column sep=large]
H^n(X,\frj^{k+1}\F)\arrow[r,"\mu_{x,k+1}"]\arrow[d]&H^n(X,\frj^{k+m+1}\F)\arrow[d]\\
H^n(X,\F)\arrow[r,"\mu_{x,0}"]&H^n(X,\F).
\end{tikzcd}
\end{equation}
由 $R_k$ 的像的解释，得到（4.1.7.10）。它还说明分次 $S$-模 $R=\bigoplus_kR_k$ 是 $M=\bigoplus_kH^n(X,\frj^{k+1}\F)$ 的商，后者是 $E$ 的相应子模（按次数平移理解）。因此 $R$ 有限型，这等价于引理（Bourbaki，\textit{Alg. comm.}，III，§3，n°1，定理1）。\finishproof
\egapar{4.1.7.12}
同样定义 $N=\bigoplus_kH^{n+1}(X,\frj^{k+1}\F)$，由（3.3.2）它有限型。将（4.1.7.11）中的 $n$ 换成 $n+1$，并利用 $Q_k\subset N_k$，得到 $S_mQ_k=\frj^mQ_k\subset Q_{k+m}$。故 $Q=\bigoplus_kQ_k$ 是 $N$ 的分次子模，亦有限型。
\egapar{4.1.7.13}
记 $a_k:\frj^k\hookrightarrow A$ 为嵌入，也写成 $S_k\to S_0$。因为 $\frj^{k+1}\F_k=0$，$H_k$ 被 $\frj^{k+1}$ 零化；$Q_k$ 是 $H_k$ 的 $A$-线性像，故也被零化，即
\begin{equation}\tag{4.1.7.14}a_{k+1}(S_{k+1})Q_k=0.\end{equation}
$Q$ 有限型，故存在 $k_0$、$h>0$，使 $Q_{k+h}=S_hQ_k$ 对 $k\ge k_0$ 成立（II，2.1.6，ii）。结合（4.1.7.14），存在 $r>0$ 使
\begin{equation}\tag{4.1.7.15}a_r(S_r)Q=0.\end{equation}
\egapar{4.1.7.16}
典范嵌入 $\frj^{k+m}\F\hookrightarrow\frj^k\F$ 诱导
\begin{equation}\tag{4.1.7.17}
v_m:H^{n+1}(X,\frj^{k+m}\F)\to H^{n+1}(X,\frj^k\F).
\end{equation}
对 $x\in\frj^m$，乘法有分解
\begin{equation}\tag{4.1.7.18}
H^{n+1}(X,\frj^k\F)\xrightarrow{\mu_{x,k}}
H^{n+1}(X,\frj^{k+m}\F)\xrightarrow{v_m}
H^{n+1}(X,\frj^k\F),
\end{equation}
其复合为 $A$-模中的乘以 $x$。故对 $H^{n+1}(X,\frj^k\F)$ 的每个 $A$-子模 $P$，在分次模 $N$ 中有
\begin{equation}\tag{4.1.7.19}v_m(S_mP)=a_m(S_m)P.\end{equation}
\egastatement{4.1.7.20}{引理}
存在 $m>0$，使 $v_m(Q_{k+m})=0$ 对所有 $k\ge k_0$ 成立。
\eproof
取 $h$ 的倍数 $m\ge r$，则 $Q_{k+m}=S_mQ_k$；由（4.1.7.19）、（4.1.7.15），
\[
v_m(Q_{k+m})=a_m(S_m)Q_k\subset a_r(S_r)Q_k=0.
\]
\finishproof
\egapar{4.1.7.21}
典范映射给出交换图
\[
\begin{tikzcd}[column sep=small]
0\arrow[r]&\frj^{k+1}\F\arrow[r]&\F\arrow[r]&\F_k\arrow[r]&0\\
0\arrow[r]&\frj^{k+m+1}\F\arrow[r]\arrow[u]&\F\arrow[r]\arrow[u,equal]&\F_{k+m}\arrow[r]\arrow[u]&0.
\end{tikzcd}
\]
其上同调图为
\[
\begin{tikzcd}[column sep=small]
H^n(X,\F)\arrow[r]&H^n(X,\F_k)\arrow[r]&H^{n+1}(X,\frj^{k+1}\F)\arrow[r]&H^{n+1}(X,\F)\\
H^n(X,\F)\arrow[r]\arrow[u,equal]&H^n(X,\F_{k+m})\arrow[r]\arrow[u]&H^{n+1}(X,\frj^{k+m+1}\F)\arrow[r]\arrow[u,"v_m"']&H^{n+1}(X,\F)\arrow[u,equal].
\end{tikzcd}
\]
从而得到正合行的交换图
\[
\begin{tikzcd}[column sep=small]
0\arrow[r]&R_k\arrow[r]&H\arrow[r]&H_k\arrow[r]&Q_k\arrow[r]&0\\
0\arrow[r]&R_{k+m}\arrow[r]\arrow[u]&H\arrow[r]\arrow[u,equal]&H_{k+m}\arrow[r]\arrow[u]&Q_{k+m}\arrow[r]\arrow[u,"v_m"']&0.
\end{tikzcd}
\]
当 $k\ge k_0$ 时，最右竖箭头为零，故 $H_{k+m}$ 在 $H_k$ 中的像包含于 $\Ker(H_k\to Q_k)=\Im(H\to H_k)$；交换性又给出反向包含，故二者相等。对 $k'\ge k+m$，$H_{k'}$ 在 $H_k$ 中的像也等于此像，证明了（ML）。

此外，对于 $X$ 的任意仿射开集 $U$，$H^i(U,\F_k)=0$（$i>0$），且 $\Gamma(U,\F_{k+m})\to\Gamma(U,\F_k)$ 满射（I，1.3.9）。由（0，13.3.1），$H^n(\mathfrak X,\widehat\F)\to\plim_kH^n(X,\F_k)$ 是双射。

射影系 $(H/R_k)_k$ 严格，故可在正合列
\begin{equation}\tag{4.1.7.22}0\to H/R_k\to H_k\to Q_k\to0\end{equation}
中取射影极限（0，13.2.2）。由 $v_m(Q_{k+m})=0$，有 $\plim_kQ_k=0$；故 $\plim_k(H/R_k)\simeq\plim_kH_k$ 为拓扑同构。由于 $(R_k)$ 是 $\frj$-良滤过，它给出 $\frj$-预进制拓扑，故左端就是 $H$ 的分离完备化。（4.1.7）全部得证。\finishproof
\egapar{4.1.8}
最后证明（4.1.5）。对 $Y$ 的每个仿射开集 $V$，若 $\J|V=\wt\frj$，则 $\Gamma(V,(R^nf_*\F)^\wedge)$ 是 $\Gamma(V,R^nf_*\F)$ 的 $\frj$-预进制分离完备化，因为该模层凝聚（I，10.8.4）。又
\[
\Gamma(V,\plim_kR^nf_*\F_k)=\plim_k\Gamma(V,R^nf_*\F_k)
\]
（$\OI$，3.2.6）。由（4.1.7）及（1.4.11），$\varphi_n$ 为拓扑同构。同样由这两个结果，$\psi_{n,V}$ 为同构；由高阶直接像的定义，$\psi_n$ 亦为同构。
\egastatement{4.1.9}{推论}
在（4.1.4）的条件下，对于 $Y$ 的每个仿射开集 $V$，典范同态
\[
H^n(\mathfrak X\cap f^{-1}(V),\widehat\F)
\longrightarrow\Gamma(\mathfrak Y\cap V,R^n\widehat f_*\widehat\F)
\]
为双射。
\egastatement{4.1.10}{注记}
设 $f:X\to Y$ 为诺特通常预概形的有限型态射，$\F$ 凝聚且支集在 $Y$ 上固有。由（3.2.4），$R^nf_*\F$ 凝聚。可取以 $\Supp(\F)$ 为底空间的闭子预概形 $Z$，使 $\F=u_*\G$，其中 $u:Z\to X$ 为嵌入、$\G=u^*\F$ 凝聚（I，9.3.5）。令 $\G_k=\G\otimes(\OO_Y/\J^{k+1})$。则
\[
\F_k=u_*\G_k,\quad R^nf_*\F=R^n(f\circ u)_*\G,
\quad R^nf_*\F_k=R^n(f\circ u)_*\G_k
\]
（1.3.4）。再利用（I，10.9.5），形式完备化后也有 $R^n\widehat f_*\widehat\F=R^n\widehat{(f\circ u)}_*\widehat\G$。因此把（4.1.5）用于 $\G$ 及固有态射 $f\circ u$，可知其全部结论仍适用于此处的 $f,\F$。\par

% ---- END c3_s04a.tex ----
% ---- BEGIN c3_s04b.tex ----
\subsection{特殊情形与变式}
\src{129--130}
比较定理（4.1.5）最有用的形式如下。
\egastatement{4.2.1}{命题}
设 $Y$ 局部诺特，$f:X\to Y$ 固有，$\F$ 为凝聚 $\OO_X$-模。对每个 $y\in Y$ 及每个 $p$，$(R^pf_*\F)_y$ 是有限型 $\OO_{Y,y}$-模，因而在 $\mm_y$-预进制拓扑下分离；并有典范拓扑同构
\begin{equation}\tag{4.2.1.1}
\bigl((R^pf_*\F)_y\bigr)^\wedge
\simeq\plim_kH^p\!\left(f^{-1}(y),\F\otimes_{\OO_Y}(\OO_{Y,y}/\mm_y^{k+1})\right).
\end{equation}
左端取 $\mm_y$-预进制完备化；右端的 $f^{-1}(y)$，对每个 $k\ge0$，视为预概形 $X\times_Y\Spec(\OO_{Y,y}/\mm_y^{k+1})$ 的底空间（I，3.6.1）。
\eproof
有限生成性来自（3.2.1），分离性来自（$\OI$，7.3.5）。当 $Y$ 诺特且 $y$ 闭时，换成 $y$ 的一个仿射开邻域并取 $Y'=\{y\}$，其余结论即由（4.1.7）和（G，II，4.9.1）得到。一般地，令
\[
Y_1=\Spec(\OO_{Y,y}),\quad X_1=X\times_YY_1,\quad
\F_1=\F\otimes_{\OO_Y}\OO_{Y,y},\quad f_1=f\times1_{Y_1}.
\]
$Y_1$ 诺特，$f_1$ 固有，$\F_1$ 凝聚（II，5.4.2，iii；$\OI$，5.3.11）。设 $y_1$ 为 $Y_1$ 的唯一闭点，命题已适用于 $f_1,\F_1,y_1$。此时 $\OO_{Y_1,y_1}=\OO_{Y,y}$，$f_1^{-1}(y_1)=f^{-1}(y)$（I，3.6.5），相应的各无穷小增厚由纤维积结合律典范认同（I，3.3.9），其模也由（I，9.1.6）认同。最后，局部态射 $\Spec(\OO_{Y,y})\to Y$ 平坦（$\OI$，6.7.1；I，2.4.2），故（1.4.15）给出 $R^pf_{1*}\F_1\simeq R^pf_*\F\otimes_{\OO_Y}\OO_{Y,y}$。\finishproof

以下推论使用第IV章的维数术语，在该章以前不会应用。
\egastatement{4.2.2}{推论}
设 $Y$ 局部诺特，$f:X\to Y$ 固有，$y\in Y$，$r=\dim f^{-1}(y)$。则对每个凝聚 $\OO_X$-模 $\F$ 和每个 $p>r$，$R^pf_*\F$ 在 $y$ 的邻域内为零。
\eproof
对每个 $k$，右端的上同调 $H^p(f^{-1}(y),\F\otimes(\OO_{Y,y}/\mm_y^{k+1}))=0$（G，II，4.15.2）。由（4.2.1），$(R^pf_*\F)_y$ 的分离完备化为零；因该拓扑分离，茎本身为零。由凝聚性，结论成立（$\OI$，5.2.2）。\finishproof
\egapar{4.2.3}
（4.2.1）主要在 $p=0$ 时使用，得到如下推论。
\egastatement{4.2.4}{推论}
在（4.2.1）的条件下，有典范拓扑同构
\[
\bigl((f_*\F)_y\bigr)^\wedge\simeq
\plim_k\Gamma\!\left(f^{-1}(y),\F\otimes_{\OO_Y}(\OO_{Y,y}/\mm_y^{k+1})\right).
\]

\subsection{Zariski 连通性定理}
\src{130--135}
本小节及下一小节的结果推广 Zariski 的若干著名定理，它们都可从（4.2.4）推出；具体地说，都是下述定理的结果。
\egastatement{4.3.1}{定理（连通性定理）}
设 $Y$ 局部诺特，$f:X\to Y$ 固有。则 $\A(X)=f_*\OO_X$ 是凝聚 $\OO_Y$-代数。设 $Y'$ 为 $Y$ 上的有限概形，满足 $\A(Y')=\A(X)$，它在 $Y$-同构意义下确定（II，1.3.1、6.1.3）。由恒等同构 $e:\A(Y')\simeq\A(X)$ 得到 $Y$-态射 $f'=\rho(e):X\to Y'$（II，1.2.7）。那么 $f'$ 固有，$f'_*\OO_X\simeq\OO_{Y'}$，并且对所有 $y'\in Y'$，纤维 $f'^{-1}(y')$ 连通且非空。
\eproof
设 $g:Y'\to Y$ 为结构态射。由于 $g$ 仿射，要证明定义 $f'$ 的同态 $\theta:\OO_{Y'}\to f'_*\OO_X$ 为双射，只须证明 $g_*\theta:g_*\OO_{Y'}\to g_*f'_*\OO_X=f_*\OO_X$ 为恒等（II，1.4.2）；这正是构造的定义。$\A(X)$ 的凝聚性是有限性定理（3.2.1）的特例。$f$ 固有、$g$ 分离，故 $f'$ 固有（II，5.4.3，i）。余下只须证明以下推论。\finishproof
\egastatement{4.3.2}{推论}
在（4.3.1）的条件下，若 $f_*\OO_X$ 同构于 $\OO_Y$，则 $f$ 的每个纤维非空且连通。
\eproof
该假设先推出 $f$ 占优；因为 $f$ 闭，故 $f$ 满射。可如（4.2.1）归结到 $y$ 为闭点的情形。诺特空间 $f^{-1}(y)$ 有有限多个连通分支 $Z_1,\ldots,Z_n$；它也是沿此纤维的形式完备化 $\mathfrak X$ 的底空间。因此 $\Gamma(\mathfrak X,\OO_{\mathfrak X})$ 是环 $\Gamma(Z_i,\OO_{\mathfrak X})$ 的直积，每个因子非零，因为单位截面在每点都不同于零。把（4.1.5）用于 $\F=\OO_X$，可知该截面环同构于局部环 $\OO_{Y,y}$ 的 $\mm_y$-预进制分离完备化。它仍是局部环，不可能是多个非零环的直积，否则会有多个极大理想。因此 $n=1$。\finishproof
\egastatement{4.3.3}{推论}
在（4.3.1）的条件下，对每个 $y\in Y$，纤维 $f^{-1}(y)$ 的连通分支与有限纤维 $g^{-1}(y)$ 的点一一对应；亦即与环 $(f_*\OO_X)_y$ 的极大理想一一对应。
\eproof
$g$ 有限，故 $g^{-1}(y)$ 有限离散（II，6.1.7），而 $f^{-1}(y)=f'^{-1}(g^{-1}(y))$。应用（4.3.1）即可。\finishproof
这给出了 $Y'$ 的一个重要解释。固有态射的分解 $f=g\circ f'$ 与 K. Stein 对解析空间的全纯映射所作分解相似，今后称为 $f$ 的\textbf{Stein 分解}。
\egastatement{4.3.4}{注记}
若 $k$ 是 $k(y)$ 的扩域，且 $f^{-1}(y)\otimes_{k(y)}k=X\times_Y\Spec(k)$ 连通，则作为其投影像的 $f^{-1}(y)$ 也连通（I，3.4.7）。对于任意预概形态射 $f:X\to Y$，若对 $k(y)$ 的每个扩域 $k$，此基变换后的纤维都连通，就称 $f^{-1}(y)$ \textbf{几何连通}。

在（4.3.2）的假设下，纤维事实上几何连通。对任意 $k/k(y)$，由（0，10.3.1），存在诺特局部环 $A$ 及局部同态 $\OO_{Y,y}\to A$，使 $A$ 在 $\OO_{Y,y}$ 上平坦且剩余域与 $k$ 在 $k(y)$ 上同构。令 $Y_1=\Spec(A)$，$h:Y_1\to Y$ 为对应局部态射，$y_1$ 为其闭点，$X_1=X\times_YY_1$、$f_1=f\times1_{Y_1}$。则 $f_1$ 固有，$f_1^{-1}(y_1)\simeq X\times_Y\Spec(k)$。由（I，2.4.2）和（1.4.15.5），$h$ 平坦；由（1.4.15）在 $q=0$ 的情形，
\[
f_{1*}\OO_{X_1}=h^*(f_*\OO_X)=h^*\OO_Y=\OO_{Y_1}.
\]
故（4.3.2）适用。

一般地，同样的论证表明基变换后 $f_1$ 的 Stein 分解的有限部分是 $Y'_1=Y'\times_YY_1$、$g_1=g\times1_{Y_1}$（II，1.5.2）。由纤维的传递性（I，3.6.4）和（4.3.3），$f_1^{-1}(y_1)$ 的连通分支数等于 $g^{-1}(y)\otimes_{k(y)}k$ 的点数。若 $k$ 为代数闭扩域，这个数不依赖其选择，等于 $g^{-1}(y)$ 的几何点数（I，6.4.7），亦即
\[
\sum_{y'_i\in g^{-1}(y)}[k(y'_i):k(y)]_s.
\]
称它为 $f^{-1}(y)$ 的\textbf{几何连通分支数}。各 $k(y'_i)$ 正是半局部环 $(f_*\OO_X)_y$ 的剩余域。
\egastatement{4.3.5}{命题}
设 $X,Y$ 为局部诺特整预概形，$f:X\to Y$ 为固有占优态射。对每个 $y\in Y$，$f^{-1}(y)$ 的连通分支数不超过 $\OO_{Y,y}$ 在有理函数域 $R(X)$ 中的整闭包 $\OO'_y$ 的极大理想数。
\eproof
对于 $Y$ 的开集 $U$，$\Gamma(U,f_*\OO_X)=\Gamma(f^{-1}(U),\OO_X)$ 是 $x\in f^{-1}(U)$ 的各局部环 $\OO_{X,x}$ 在 $R(X)$ 中的交（I，8.2.1.1）。因此 $(f_*\OO_X)_y$ 是 $R(X)$ 中包含 $\OO_{Y,y}$ 的子环。由凝聚性，它在 $\OO_{Y,y}$ 上有限，故包含在 $\OO'_y$ 中。这样的子环的每个极大理想，都是 $\OO'_y$ 的某个极大理想的交（[13]，I，257、259页）；结合（4.3.3）即得。\finishproof
\egastatement{4.3.6}{定义}
若整局部环的整闭包是局部环，就称它\textbf{单支}。若整预概形 $Y$ 在点 $y$ 的局部环 $\OO_{Y,y}$ 单支，就称 $y$ 为单支点；特别地，正规点是单支点。

设整局部环 $A$ 的分式域为 $K$。$A$ 单支，当且仅当 $K$ 中每个包含 $A$ 且在 $A$ 上有限的子环 $A_\alpha$ 都是局部环。事实上，若 $A'$ 为 $A$ 的整闭包，则由 Cohen--Seidenberg 第一定理，$A_\alpha$ 的每个极大理想都是 $A'$ 中极大理想的交；故 $A'$ 局部推出 $A_\alpha$ 局部。反过来，$A'$ 是这些有限子代数的滤过递增并；各包含 $A_\alpha\subset A_\beta$ 是局部同态，故由（0，10.3.1.3）其极限 $A'$ 为局部环。

若诺特局部环 $A$ 的完备化是整环（称 $A$ \textbf{解析整}），则 $A$ 单支。设 $\mm$ 为极大理想，$K$ 为分式域，$K'$ 为 $\widehat A$ 的分式域。原文写 $K'=K\otimes_A\widehat A$。\transnote{一般首先有典范嵌入 $K\otimes_A\widehat A\hookrightarrow\operatorname{Frac}(\widehat A)$；局部化未必已经是整个分式域。下面使用的是其整环性及有限子代数的完备化，不应把此等式无条件引用。}
若 $A_\beta$ 是 $K$ 中有限 $A$-子代数，则 $K'$ 中由 $\widehat A$ 与 $A_\beta$ 生成的子环 $B_\beta\simeq A_\beta\otimes_A\widehat A$，是有限型 $\widehat A$-模，也是 $A_\beta$ 的 $\mm$-进制完备化（$\OI$，7.3.3、7.3.6）。$A_\beta$ 半局部，而其完备化是整环，故只能有一个极大理想，且该理想与 $A$ 的交为 $\mm$。结论随之成立。
\egastatement{4.3.7}{推论}
在（4.3.5）的条件下，设 $R(Y)$ 在 $R(X)$ 中的代数闭包具有可分次数 $n$，且 $y\in Y$ 为单支点。则 $f^{-1}(y)$ 至多有 $n$ 个连通分支。特别地，若上述扩张在 $R(Y)$ 上纯不可分，则该纤维连通。\transnote{最后一句底本印作“在 $R(X)$ 上纯不可分”；所指扩张的基域应是 $R(Y)$。}
\eproof
设 $\OO''_y$ 为 $\OO_{Y,y}$ 的整闭包，则 $\OO_{Y,y}$ 在 $R(X)$ 中的整闭包，也等于 $\OO''_y$ 在 $R(X)$ 中的整闭包。$\OO''_y$ 局部，故后者半局部，且极大理想数不超过 $n$（[13]，I，289页，定理22）。应用（4.3.5）即得。\finishproof
这基本上就是 Zariski 对代数概形陈述其连通性定理的形式。
\egastatement{4.3.8}{注记}
若在（4.3.7）的假设上，再加上 $Y$ 在 $y$ 正规，则纤维 $f^{-1}(y)$ 几何连通，因为采用（4.3.4）的记号，$g^{-1}(y)$ 只有一个点 $y'$，且 $k(y')$ 在 $k(y)$ 上纯不可分。\transnote{此处的几何连通结论承接（4.3.7）最后一句的纯不可分情形，不能理解为对任意可分次数 $n$ 都成立。}
\egastatement{4.3.9}{定义}
给定局部诺特预概形 $Y$。若有限型态射 $f:X\to Y$ 满足：对于每个不可约局部诺特预概形 $Y'$ 及每个占优态射 $g:Y'\to Y$，$X'=X\times_YY'$ 的每个不可约分支都占优于 $Y'$，则称 $f$ \textbf{普遍开}。

当 $Y$ 不可约时，设 $\eta,\eta'$ 分别为 $Y,Y'$ 的泛点，则 $g(\eta')=\eta$。若 $f'=f_{(Y')}$，上述条件等价于 $X'$ 的每个不可约分支都与 $f'^{-1}(\eta')$ 相交（$\OI$，2.1.8）。由此，对每个开集 $U\subset Y$，限制 $f^{-1}(U)\to U$ 仍普遍开。\transnote{本段保留原书在此处给出的定义；它与后文用任意基变换后的开映射所作表述的关系见（4.3.11，ii），不在译文中以现代定义替换。}
\egastatement{4.3.10}{推论}
设 $X,Y$ 为局部诺特整预概形，$f:X\to Y$ 为固有、占优、普遍开的态射。若 $R(Y)$ 在 $R(X)$ 中的代数闭包在 $R(Y)$ 上纯不可分，则每个纤维 $f^{-1}(y)$ 几何连通。
\eproof
可设 $Y=\Spec(B)$，$B$ 诺特整环。由（II，7.1.7），存在支配 $\OO_{Y,y}$、分式域为 $R(Y)$ 的诺特整闭局部环 $A$。令 $Y'=\Spec(A)$，由 $B\hookrightarrow A$ 得到 $h:Y'\to Y$；它双有理，因而占优。设 $y'$ 为 $Y'$ 的闭点，则 $h(y')=y$。令 $X'=X\times_YY'$、$f'=f\times1_{Y'}$。若 $\eta,\eta',\xi$ 分别为 $Y,Y',X$ 的泛点，则 $f(\xi)=\eta$、$h^{-1}(\eta)=\{\eta'\}$ 且 $k(\eta')=k(\eta)=R(Y)$。因此 $f'^{-1}(\eta')\simeq f^{-1}(\eta)$，只有一个泛点。由假设，$X'$ 的每个不可约分支的泛点都在该纤维中，故 $X'$ 不可约，其泛点 $\xi'$ 满足 $k(\xi')=k(\xi)$。

令 $X''=X'_{\red}$、$f''=f'_{\red}$。则 $X''$ 整且诺特，$f''$ 固有（II，5.4.6），它与 $f'$ 在 $y'$ 上的纤维底空间相同；又 $R(X'')=R(X)$。因此可应用（4.3.8），得 $f''^{-1}(y')$ 几何连通。对任意 $k/k(y)$，可取一个共同扩域 $k_1$，包含 $k$ 及 $k(y')$ 作为 $k(y)$-子扩域。则
\[
f''^{-1}(y')\times_{y'}\Spec(k_1)
\]
连通；它与 $f'^{-1}(y')\times_{y'}\Spec(k_1)$ 的约化预概形相同（I，5.1.8），故后者连通。由纤维的传递性，它同构于 $f^{-1}(y)\times_y\Spec(k_1)$，再投影到 $f^{-1}(y)\times_y\Spec(k)$，由（4.3.4）知后者连通。\finishproof
\egastatement{4.3.11}{注记}
（i）上述论证实质上出自 Zariski [20]，但在经典代数几何的局部环情形中，他可取 $\OO_{Y,y}$ 的整闭包作为 $A$，因为该整闭包诺特。此外，Zariski 证明：若 $Y$ 是域 $k$ 上射影空间 $\Pb_k^r$ 的 Chow 簇，$X\subset\Pb_k^r\times_kY$ 是定义 Chow 对应的闭子集，则投影 $X\to Y$ 普遍开（该文82页的引理）。看来这是“Chow 坐标”在某些应用中唯一用到的形式性质。因此，在这种情形下，以“固有态射的纤维”（可再假设普遍开，或施加类似的局部正则性限制）取代“射影空间中圈的特殊化”的语言，是有益的。

（ii）在第IV章将看到，普遍开态射也可由如下条件定义，这将说明其名称：对任意 $Z\to Y$，态射 $f_{(Z)}:X_{(Z)}\to Z$ 是开映射。还可证明，若 $f$ 满足（4.3.10）的假设，$y$ 是 $y'$ 的特殊化，则 $f^{-1}(y)$ 的几何连通分支数不超过 $f^{-1}(y')$ 的几何连通分支数。
\egastatement{4.3.12}{推论}
在（4.3.5）的假设上，再设 $R(Y)$ 在 $R(X)$ 中代数闭，$y$ 为 $Y$ 的正规点。则 $f^{-1}(y)$ 几何连通，且存在 $y$ 的开邻域 $U$，使 $f_*(\OO_X|f^{-1}(U))\simeq\OO_Y|U$。特别地，若 $Y$ 正规，且上述函数域条件成立，则 $f_*\OO_X\simeq\OO_Y$。
\eproof
几何连通性来自（4.3.8）。在 Stein 分解 $f=g\circ f'$ 中，$g^{-1}(y)$ 只有一个点 $y'$；并且
\[
\OO_{Y,y}\subset\OO_{Y',y'}=(f_*\OO_X)_y\subset R(X).
\]
中间环在 $\OO_{Y,y}$ 上有限，其元素因而代数于 $R(Y)$，由假设它包含在 $R(Y)$ 中；$y$ 正规，故 $\OO_{Y',y'}=\OO_{Y,y}$。由（I，6.5.4），$g$ 在 $y'$ 是局部同构，得到第一部分。若 $Y$ 处处正规，则 $g$ 双射且处处局部同构，故为同构，得到第二部分。\finishproof
在概形框架中建立（4.3.7），还可作如下应用。
\egastatement{4.3.13}{命题}
设 $A$ 为诺特单支局部环，$\mathfrak a$ 为其定义理想，$A_0=A/\mathfrak a$，$S=\gr_{\mathfrak a}(A)$ 为相应关联分次环。$S$ 是由有限型 $A_0$-模 $S_1$ 生成的分次 $A_0$-代数；则 $\Proj(S)$ 是连通的 $A_0$-概形。
\eproof
设 $\mm$ 为极大理想，$Y=\Spec(A)$，$y$ 是唯一闭点。由假设，对某个 $p$ 有 $\mm^p\subset\mathfrak a\subset\mm$，故 $V(\mathfrak a)=\{\mm\}$。令 $S'=\bigoplus_k\mathfrak a^k$、$X=\Proj(S')$，即爆破 $\mathfrak a$ 所得的 $Y$-概形。它整，结构态射 $f:X\to Y$ 双有理（II，8.1.4），而且显然射影。由（4.3.7），$f^{-1}(y)$ 连通。该空间也是 $\Proj(S'\otimes_AA_0)$ 的底空间（I，3.6.1；II，2.8.10），而按定义 $S'\otimes_AA_0=S$，故命题成立。\finishproof\par

% ---- END c3_s04b.tex ----
% ---- BEGIN c3_s04c.tex ----
\subsection{Zariski 的“主定理”}
\src{135--138}
\egastatement{4.4.1}{命题}
设 $Y$ 局部诺特，$f:X\to Y$ 固有。令 $X'$ 为那些在各自纤维 $f^{-1}(f(x))$ 中孤立的点 $x\in X$ 所成集合。则 $X'$ 在 $X$ 中开；若 $f=g\circ f'$ 是 Stein 分解（4.3.3），则 $f'$ 在 $X'$ 上的限制把 $X'$ 同构到 $Y'$ 的某个开集 $U$ 上的诱导子预概形，并且 $X'=f'^{-1}(U)$。
\eproof
$g^{-1}(f(x))$ 有限离散（4.3.3；II，6.1.7），故 $x$ 在 $f^{-1}(f(x))$ 中孤立，当且仅当它在 $f'^{-1}(f'(x))$ 中孤立。因此可设 $f'=f$，即 $f_*\OO_X=\OO_Y$。对 $x\in X'$，连通纤维 $f^{-1}(f(x))$ 只能是单点 $\{x\}$（4.3.2）。$f$ 闭，故对 $x$ 的任意开邻域 $V$，$f(X-V)$ 是不含 $y=f(x)$ 的闭集。取其补集 $U$，有 $f^{-1}(U)\subset V$。所以 $y$ 的开邻域基的逆像构成 $x$ 的开邻域基。由直接像的定义及 $f_*\OO_X=\OO_Y$，若 $f=(\psi,\theta)$，则 $\theta_x^\sharp:\OO_{Y,y}\to\OO_{X,x}$ 是同构。由（I，6.5.4），存在开邻域 $V,U$，使 $f|V:V\to U$ 为同构；又可取 $f^{-1}(U)=V$。依定义 $V\subset X'$，结论随之成立。\finishproof
以下命题在代数概形的情形下由 Chevalley 证明。
\egastatement{4.4.2}{命题}
设 $Y$ 局部诺特，$f:X\to Y$ 为态射。以下条件等价：a）$f$ 有限；b）$f$ 仿射且固有；c）$f$ 固有，且每个 $f^{-1}(y)$ 都是有限集。
\eproof
a）推出 b）已知（II，6.1.2、6.1.11）。若 $f$ 仿射固有，则 $f^{-1}(y)\to\Spec(k(y))$ 亦如此。把有限性定理（3.2.1）用于其结构层，可知该纤维是有限 $k(y)$-代数的谱，故点集有限，得到 b）$\Rightarrow$c）。若 c）成立，因纤维是 $k(y)$ 上的代数预概形，其有限底空间必离散（I，6.4.4）。所以在（4.4.1）的记号下，$X'=X$、$f':X\to Y'$ 为同构，而 $g$ 有限，故 $f$ 有限。\finishproof
\egastatement{4.4.3}{定理（Zariski 主定理）}
设 $Y$ 为诺特预概形，$f:X\to Y$ 拟射影。令 $X'$ 为在各自纤维中孤立的点的集合。则 $X'$ 在 $X$ 中开，而且其诱导子预概形同构于某个有限 $Y$-预概形 $Y'$ 的一个开子预概形。
\eproof
存在射影 $Y$-预概形 $Z$，使 $X$ 在 $Y$ 上同构于 $Z$ 的开子预概形（II，5.3.2、5.5.1）。因此只须对射影态射证明，而射影态射固有（II，5.5.3），可直接应用（4.4.1）。\finishproof
\egastatement{4.4.4}{注记}
若 $X$ 约化（分别地，$X$ 不可约且 $X'$ 非空），则在（4.4.3）中可取 $Y'$ 约化（分别地，不可约）。事实上，可将 $Y'$ 换成 $X'$ 在其中的概形论闭包（I，9.5.11；II，6.1.5，i、ii）。约化开子预概形的闭包仍约化（I，9.5.9，i）；非空不可约集的闭包仍不可约。
\egastatement{4.4.5}{推论}
设 $Y$ 为局部诺特概形，$f:X\to Y$ 有限型，$x$ 在其纤维中孤立。则 $x$ 有一个开邻域，在 $Y$ 上同构于某个有限 $Y$-预概形的开子预概形。
\eproof
取 $y=f(x)$ 的仿射开邻域 $U$，再取 $x$ 的仿射开邻域 $V\subset f^{-1}(U)$。$Y$ 分离，故 $U\to Y$ 仿射（II，1.6.3），而 $V\to U$ 仿射，所以 $V\to Y$ 仿射（II，1.6.2，ii）。它又有限型（I，6.3.5），故拟射影（II，5.3.4，i）；对其应用（4.4.3）即得。\finishproof
用交换代数语言，这一推论写成如下形式。
\egastatement{4.4.6}{推论}
设 $A$ 为诺特环，$B$ 为有限型 $A$-代数，$\qq\subset B$ 为素理想，$\pp$ 为其在 $A$ 中的逆像。假设在所有逆像为 $\pp$ 的 $B$ 的素理想中，$\qq$ 同时为极大元和极小元。则存在 $g\in B-\qq$、有限 $A$-代数 $A'$ 和 $g'\in A'$，使 $B_g$ 与 $A'_{g'}$ 作为 $A$-代数同构。
\eproof
取 $Y=\Spec(A)$、$X=\Spec(B)$，应用（4.4.5）；关于 $\qq$ 的假设恰表示它在纤维中孤立（I，1.1.7）。\finishproof
由此得到一个表面上较不一般的结果。
\egastatement{4.4.7}{推论}
设 $A$ 为诺特局部环，$B$ 为有限型 $A$-代数，$\nn\subset B$ 为素理想，其在 $A$ 中的逆像为极大理想 $\mm$。假设 $\nn$ 在 $B$ 中极大，并且在所有逆像为 $\mm$ 的素理想中极小（原文亦表述为 $B\mm$ 是 $\nn$-准素的）。则存在有限 $A$-代数 $A'$ 及其极大理想 $\mm'$，使 $\mm'$ 在 $A$ 中的逆像为 $\mm$，且 $B_{\nn}\simeq A'_{\mm'}$ 为 $A$-代数同构。\transnote{括号中的全局准素性比所列素理想条件更强；在 $B_{\nn}$ 中的相应理想则具有唯一素理想作为其根。这里保留原文括号的表述，不把它当作无需局部化的等价判据。}
下面的特殊情形也有时称为“主定理”。
\egastatement{4.4.8}{推论}
在（4.4.7）的条件下，再设 $A,B$ 均为整环，具有同一分式域 $K$。若 $A$ 整闭，则 $B=A$。
\eproof
由（4.4.4），可在应用（4.4.7）时取 $A'$ 为分式域为 $K$ 的整环。$A$ 整闭而 $A'$ 在 $A$ 上有限，故 $A'=A$，于是 $B_{\nn}=A$。由 $A\subset B\subset B_{\nn}$，得到 $A=B$。\finishproof
这就是 Zariski 对主定理的陈述形式，现已推广到任意诺特整局部环。

以上是全局结果（4.4.3）的局部变式。另一个全局结果如下。
\egastatement{4.4.9}{推论}
设 $Y$ 为局部诺特整预概形，$f:X\to Y$ 分离、有限型且双有理。若 $Y$ 正规、所有纤维 $f^{-1}(y)$ 有限，则 $f$ 为开浸入；若 $f$ 还闭，特别是若 $f$ 固有，则 $f$ 为同构。
\eproof
令 $y=f(x)$。有限的代数纤维离散（I，6.4.4），而 $\OO_{Y,y}$ 整闭，$\OO_{X,x}$ 与 $\OO_{Y,y}$ 具有同一分式域（I，7.1.5）。由（4.4.8），相应茎同态为双射，故 $f$ 为局部同构（I，6.5.4）。$f$ 分离、$X$ 整，故 $f$ 为开浸入（I，8.2.8）。最后一断言由占优性得出。\finishproof
\egastatement{4.4.10}{命题}
设 $Y$ 局部诺特，$f:X\to Y$ 局部有限型。那些在其纤维中孤立的点组成 $X$ 的开子集。
\eproof
问题在 $X,Y$ 上局部，可设二者仿射诺特，$f$ 有限型。此时 $f$ 仿射有限型，故拟射影，应用（4.4.3）即可。\finishproof
\egastatement{4.4.11}{推论}
设 $Y$ 局部诺特，$f:X\to Y$ 固有。使 $f^{-1}(y)$ 离散的点 $y$ 所成集合 $U$ 在 $Y$ 中开，而且 $f^{-1}(U)\to U$ 有限。特别地，固有拟有限态射有限。
\eproof
$Y-U=f(X-X')$，其中 $X'$ 是（4.4.10）中的开集；$f$ 闭，故 $U$ 开。由（II，6.2.2），$U$ 上各纤维有限；限制态射固有，故由（4.4.2）它有限。\finishproof
\egastatement{4.4.12}{注记}
（i）如（II，6.2.7）所预告，第V章将证明：当 $Y$ 局部诺特时，每个拟有限分离态射 $X\to Y$ 都拟仿射，因而拟射影。于是主定理（4.4.3）的结论在只假设 $f$ 分离且有限型时仍成立：由（4.4.10），$X'$ 开；$X$ 局部诺特，故限制仍有限型（I，6.3.5），依定义拟有限，且分离，故可对限制应用（4.4.3）。

（ii）第IV章将用维数理论给出（4.4.10）的较初等证明。

\subsection{同态模的完备化}
\src{138--139}
\egastatement{4.5.1}{命题}
设 $A$ 为诺特环，$\frj\subset A$ 为理想，$X$ 为有限型 $A$-预概形，$\F,\G$ 为凝聚 $\OO_X$-模，其支集的交在 $Y=\Spec(A)$ 上固有。则对每个 $n\ge0$，$\Ext^n_{\OO_X}(X;\F,\G)$ 是有限型 $A$-模，其 $\frj$-预进制分离完备化典范地同构于 $\Ext^n_{\OO_{\mathfrak X}}(\mathfrak X;\widehat\F,\widehat\G)$，其中完备化记号同（4.1.7）。
\eproof
有双正则谱序列 $E(\F,\G)$（T，4.2），终项为全局 $\Ext$，其第二页为
\[
E_2^{pq}=H^p(X,\shExt^q_{\OO_X}(\F,\G)).
\]
层 $\shExt^q_{\OO_X}(\F,\G)$ 凝聚（0，12.3.3），其支集包含在两支集的交中，故在 $Y$ 上固有。由（3.2.4），各 $E_2^{pq}$ 为有限型 $A$-模；由（0，11.1.8），后续各页及终项亦有限型。

典范态射 $i:\mathfrak X\to X$ 平坦，而且 $\widehat\F=i^*\F$、$\widehat\G=i^*\G$（I，10.8.8、10.8.9）。由（0，12.3.4），有典范 $i$-态射
\[
u_q:\shExt^q_{\OO_X}(\F,\G)\to
\shExt^q_{\OO_{\mathfrak X}}(\widehat\F,\widehat\G),
\]
其对应的 $\OO_{\mathfrak X}$-同态 $u_q^\sharp$ 为同构（0，12.3.5）。即右端典范地等于左端的完备化（I，10.8.8）。由比较定理（4.1.10），
\[
H^p(\mathfrak X,\shExt^q_{\OO_{\mathfrak X}}(\widehat\F,\widehat\G))
\]
是 $H^p(X,\shExt^q_{\OO_X}(\F,\G))$ 的 $\frj$-预进制分离完备化。因此若 $\widehat A$ 为 $A$ 的完备化，
\[
E_2^{pq}(\widehat\F,\widehat\G)
\simeq E_2^{pq}(\F,\G)\otimes_A\widehat A.
\]
平坦态射 $i$ 给出谱序列的典范态射 $\varphi:E(\F,\G)\to E(\widehat\F,\widehat\G)$。在第二页及终项上，它分别是由 $u_q$ 及 $u_0$ 的函子性得到的映射
\[
\begin{aligned}
\varphi_2^{pq}:&\ H^p(X,\shExt^q(\F,\G))\to
H^p(\mathfrak X,\shExt^q(\widehat\F,\widehat\G)),\\
\varphi^n:&\ \Ext^n_{\OO_X}(X;\F,\G)\to
\Ext^n_{\OO_{\mathfrak X}}(\mathfrak X;\widehat\F,\widehat\G).
\end{aligned}
\]
与 $\widehat A$ 张量后得到 $\widehat A$-线性映射 $\psi_r^{pq}$、$\psi^n$。$\widehat A$ 在 $A$ 上平坦（$\OI$，7.3.3），故 $E_r^{pq}(\F,\G)\otimes_A\widehat A$ 仍构成双正则谱序列，终项为 $\Ext^n_{\OO_X}(X;\F,\G)\otimes_A\widehat A$。这些映射构成谱序列态射，而 $\psi_2^{pq}$ 均为同构，故 $\psi^n$ 亦为同构（0，11.1.5）。\finishproof
\egastatement{4.5.2}{推论}
在（4.5.1）的条件下，若 $A$ 本身为诺特 $\frj$-进制环，则
\[
\Ext^n_{\OO_X}(X;\F,\G)\simeq
\Ext^n_{\OO_{\mathfrak X}}(\mathfrak X;\widehat\F,\widehat\G)
\]
对所有 $n\ge0$ 典范成立。
\eproof
左端有限型，故已对 $\frj$-预进制拓扑分离完备（$\OI$，7.3.6）。\finishproof
$n=0$ 的特殊情形写成如下形式。
\egastatement{4.5.3}{推论}
在（4.5.1）的条件下，以 $\widehat u:\widehat\F\to\widehat\G$ 表示同态 $u:\F\to\G$ 的完备化（I，10.8.4）。则
\begin{equation}\tag{4.5.3.1}
\bigl(\Hom_{\OO_X}(\F,\G)\bigr)^\wedge
\simeq\Hom_{\OO_{\mathfrak X}}(\widehat\F,\widehat\G),
\end{equation}
其中左端为 $A$-同态模的 $\frj$-预进制分离完备化，此同构由 $u\mapsto\widehat u$ 经分离完备化得到。\par

% ---- END c3_s04c.tex ----
% ---- BEGIN c3_s04d.tex ----
\subsection{形式态射与通常态射之间的关系}
\src{139--144}
\egastatement{4.6.1}{命题}
设 $Y$ 局部诺特，$f:X\to Y$ 固有，$\F$ 为凝聚且 $f$-平坦的 $\OO_X$-模，$y\in Y$。若对某个整数 $n$ 有
\[
H^n(f^{-1}(y),\F\otimes_{\OO_Y}k(y))=0,
\]
则存在 $y$ 的开邻域 $U$，使 $R^nf_*\F|U=0$；并且对每个 $p\ge0$，典范同态
\[
(R^{n-1}f_*\F)_y\longrightarrow
H^{n-1}\!\left(f^{-1}(y),\F\otimes_{\OO_Y}(\OO_{Y,y}/\mm_y^{p+1})\right)
\]
满射。
\eproof
由于 $R^nf_*\F$ 凝聚，第一断言只须证明其在 $y$ 的茎为零（$\OI$，5.2.2）。由（4.2.1），只须证明所有
\[
H^n(f^{-1}(y),\F\otimes(\OO_{Y,y}/\mm_y^{p+1}))
\]
为零。$p=0$ 即为假设，现对 $p$ 归纳。令 $X_p=X\times_Y\Spec(\OO_{Y,y}/\mm_y^{p+1})$；$X_{p-1}$ 是 $X_p$ 的闭子预概形，底空间相同（I，3.6.1）。由正合列
\[
0\to\mm_y^p\F/\mm_y^{p+1}\F\to\F/\mm_y^{p+1}\F\to\F/\mm_y^p\F\to0
\]
得到上同调正合列
\[
H^n(X_p,\mm_y^p\F/\mm_y^{p+1}\F)\to
H^n(X_p,\F/\mm_y^{p+1}\F)\to H^n(X_p,\F/\mm_y^p\F).
\]
最后一项由归纳假设为零，所以只须证明
\begin{equation}\tag{4.6.1.1}H^n(X_p,\mm_y^p\F/\mm_y^{p+1}\F)=0.\end{equation}
令 $Z=f^{-1}(y)=X\times_Y\Spec(k(y))$；中间层被 $\mm_y$ 零化，可看作 $\OO_Z$-模，其在 $X_p$ 与 $Z$ 上的上同调相同。典范同态
\begin{equation}\tag{4.6.1.2}
(\F/\mm_y\F)\otimes_{k(y)}(\mm_y^p/\mm_y^{p+1})
\longrightarrow\mm_y^p\F/\mm_y^{p+1}\F
\end{equation}
为双射：这是逐茎问题，$\F_x$ 在 $\OO_{Y,y}$ 上平坦，应用（0，10.2.1，c）即可。因 $\mm_y^p/\mm_y^{p+1}$ 是自由 $k(y)$-模，便有
\[
H^n(Z,\mm_y^p\F/\mm_y^{p+1}\F)
=H^n(Z,\F/\mm_y\F)\otimes_{k(y)}(\mm_y^p/\mm_y^{p+1})=0
\]
（0，12.2.3），证明了（4.6.1.1）。

为证明第二断言，可如（4.2.1）归结到 $Y$ 仿射、$y$ 闭。用相同论证，对 $k>0$ 有
\begin{equation}\tag{4.6.1.3}
H^n(X_{p+k-1},\mm_y^p\F/\mm_y^{p+k}\F)=0.
\end{equation}
由（4.2.1）得
\begin{equation}\tag{4.6.1.4}(R^nf_*(\mm_y^p\F))_y=0.\end{equation}
因此上同调正合列给出
\[
(R^{n-1}f_*\F)_y\to(R^{n-1}f_*(\F/\mm_y^p\F))_y
\to(R^nf_*(\mm_y^p\F))_y=0.
\]
因 $Y$ 仿射且 $y$ 闭，（1.4.11）及（G，II，4.9.1）将第二项认同为
\[
H^{n-1}(f^{-1}(y),\F/\mm_y^p\F),
\]
因为它本来就是 $\OO_{Y,y}/\mm_y^p$-模。将 $p$ 换成 $p+1$，得到所需满射。\finishproof
\egastatement{4.6.2}{推论}
设 $Y$ 局部诺特，$f:X\to Y$ 固有平坦，$\F,\G$ 为局部自由 $\OO_X$-模，$y\in Y$。令 $X_y=X\otimes_Yk(y)$，$\F_y=\F\otimes_{\OO_Y}k(y)$、$\G_y=\G\otimes_{\OO_Y}k(y)$。若
\begin{equation}\tag{4.6.2.1}H^1(X_y,\shHom_{\OO_{X_y}}(\F_y,\G_y))=0,\end{equation}
则对任意同态 $u_0:\F_y\to\G_y$，存在 $y$ 的开邻域 $U$ 和同态 $u:\F|f^{-1}(U)\to\G|f^{-1}(U)$，使 $u_0=u\otimes1$。
\eproof
可将（4.6.1）在 $n=1,p=0$ 时用于 $\HH=\shHom_{\OO_X}(\F,\G)$。它凝聚且局部自由，因而 $f$-平坦，而 $\HH\otimes_{\OO_Y}k(y)$ 认同于纤维上的同态层。可设 $Y=\Spec(A)$，则
\[
R^0f_*\HH=(\Hom_{\OO_X}(\F,\G))^\sim,
\quad(R^0f_*\HH)_y=\Hom_{\OO_X}(\F,\G)\otimes_A\OO_{Y,y}.
\]
（4.6.1）说明从此局部化到 $\Hom_{\OO_{X_y}}(\F_y,\G_y)$ 的典范同态满射。局部化中的元素可写为 $u'\otimes(1/s)$，$s\notin\mathfrak j_y$；在 $D(s)$ 上即可实现所需同态。\finishproof
\egastatement{4.6.3}{推论}
在（4.6.2）的条件下，若 $u_0$ 单射（分别地，满射、双射），则可取 $u$ 也具有同样性质。
\eproof
可先设 $U=Y$。只须证明对每个 $x\in f^{-1}(y)$，$\Ker u_x=0$（分别地，$\Coker u_x=0$）。核、余核凝聚，故在纤维的某个邻域上为零；$f$ 闭，故可把该邻域缩成 $f^{-1}(U')$。

由假设，$u_x\otimes1:\F_x\otimes_{\OO_{Y,y}}k(y)\to\G_x\otimes_{\OO_{Y,y}}k(y)$ 单射（分别地，满射）。两茎为有限自由 $\OO_{X,x}$-模，而 $\OO_{X,x}$ 在 $\OO_{Y,y}$ 上平坦。单射情形由（0，10.2.4）得到 $u_x$ 单射。满射情形再对 $\OO_{X,x}$ 的极大理想取商，所得剩余向量空间同态满射；由 Nakayama 引理，$u_x$ 满射。\finishproof
\egastatement{4.6.4}{推论}
设 $Y$ 局部诺特，$f:X\to Y$ 固有平坦，$y\in Y$，$X_y=X\otimes_Yk(y)$。设 $\E_0$ 为局部自由 $\OO_{X_y}$-模，满足
\begin{equation}\tag{4.6.4.1}H^1(X_y,\shHom_{\OO_{X_y}}(\E_0,\E_0))=0.\end{equation}
若局部自由 $\OO_X$-模 $\F,\G$ 在 $X_y$ 上的限制都同构于 $\E_0$，则存在 $y$ 的开邻域 $U$，使 $\F|f^{-1}(U)\simeq\G|f^{-1}(U)$。
\egastatement{4.6.5}{推论}
在（4.6.4）关于 $f,X,Y$ 的条件下，若 $H^1(X_y,\OO_{X_y})=0$，且两可逆 $\OO_X$-模 $\F,\G$ 的纤维限制同构，则在 $y$ 的某个开邻域的逆像上，它们同构。
\eproof
将（4.6.4）用于 $\F^{-1}\otimes\G$ 与 $\OO_X$ 即可。\finishproof
\egastatement{4.6.6}{注记}
（i）第V章将用（4.6.5）建立在（II，4.2.7）中预告的射影丛上可逆层的分类。（ii）在§7中，（4.6.1）将作为更一般命题的结果出现。
\egastatement{4.6.7}{命题}
设 $Z$ 局部诺特，$X,Y$ 是 $Z$-预概形，其结构态射 $g:X\to Z$、$h:Y\to Z$ 固有。设 $f:X\to Y$ 为 $Z$-态射，$z\in Z$，$f_z=f\times_Z1:X\otimes_Zk(z)\to Y\otimes_Zk(z)$。

（i）若 $f_z$ 有限（分别地，为闭浸入），则存在 $z$ 的开邻域 $U$，使限制 $g^{-1}(U)\to h^{-1}(U)$ 有限（分别地，为闭浸入）。

（ii）若再假设 $g$ 平坦，而 $f_z$ 是同构，则存在这样的 $U$，使上述限制为同构。
\eproof
只须对每个 $y\in h^{-1}(z)$ 找到开邻域 $V_y$，使 $f^{-1}(V_y)\to V_y$ 具有对应性质；把这些邻域取并得到 $V$，再利用 $h$ 闭，取 $z$ 的开邻域 $U$ 使 $h^{-1}(U)\subset V$ 即可。

（i）$f$ 固有（II，5.4.3）。若 $f_z$ 有限，则由（4.4.11）即可找到各 $V_y$。处理闭浸入情形时，已可假设 $f$ 有限，写成 $X=\Spec(\B)$，其中 $\B$ 为凝聚 $\OO_Y$-代数，$f$ 对应典范同态 $u:\OO_Y\to\B$。只须证明 $u_y$ 满射，因为 $\Coker u$ 凝聚。$f_z$ 对应 $u\otimes1$；其为闭浸入说明 $u_y\otimes1$ 在 $k(z)$ 上满射。由 $\B_y$ 有限型及 Nakayama 引理，$u_y$ 满射。

（ii）同样，只须从 $u_y\otimes1$ 双射推出 $u_y$ 双射。所需单射性由下述引理得到。\finishproof
\egastatement{4.6.7.1}{引理}
设 $A,B$ 为诺特局部环，$\rho:A\to B$ 为局部同态，$u:N\to M$ 为 $B$-模同态。假设 $M$ 在 $A$ 上平坦，$N$ 为有限型 $B$-模，且 $u\otimes1:N\otimes_Ak\to M\otimes_Ak$ 单射，其中 $k$ 是 $A$ 的剩余域。则 $N$ 在 $A$ 上平坦，且 $u$ 单射。
\eproof
对于有限型 $A$-模的单同态 $v:P\to Q$，考虑交换图
\[
\begin{tikzcd}
N\otimes_AP\arrow[r,"1_N\otimes v"]\arrow[d,"u\otimes1_P"']&N\otimes_AQ\arrow[d,"u\otimes1_Q"]\\
M\otimes_AP\arrow[r,"1_M\otimes v"]&M\otimes_AQ.
\end{tikzcd}
\]
下箭头单射，故若左箭头单射，上箭头即单射，进而证明 $N$ 平坦。令 $\mm$ 为 $A$ 的极大理想。$N\otimes_AP$ 作为 $B$-模有限型，其 $\mm$-进制滤过就是 $\mm B$-进制滤过；因为 $\mm B$ 包含在 $B$ 的根基中，该拓扑分离（$\OI$，7.3.5）。因此只须证明关联分次同态 $\gr(u\otimes1_P)$ 单射（Bourbaki，\textit{Alg. comm.}，III，§2，n°8，定理1推论1）。

$M$ 平坦，故 $M\otimes_A\mm^nP\to\mm^n(M\otimes_AP)$ 双射；于是典范同态
\[
\varphi_M:\gr_0(M)\otimes_k\gr(P)\longrightarrow\gr(M\otimes_AP)
\]
双射。而下图交换：
\[
\begin{tikzcd}[column sep=large]
\gr_0(N)\otimes_k\gr(P)\arrow[r,"\gr_0(u)\otimes1"]\arrow[d,"\varphi_N"']&\gr_0(M)\otimes_k\gr(P)\arrow[d,"\varphi_M"]\\
\gr(N\otimes_AP)\arrow[r,"\gr(u\otimes1)"]&\gr(M\otimes_AP).
\end{tikzcd}
\]
$\varphi_M$ 双射、$\varphi_N$ 满射；$\gr_0(u)$ 由假设单射，与 $k$-向量空间 $\gr(P)$ 张量仍单射。因此下箭头单射，证明第一断言。在上述论证中取 $P=A$，即得 $u$ 单射。\finishproof
\egastatement{4.6.8}{命题}
设 $Z$ 局部诺特，$X,Y$ 为结构态射 $g,h$ 固有的 $Z$-预概形。设 $Z'\subset Z$ 闭，$X'=g^{-1}(Z')$、$Y'=h^{-1}(Z')$，$\mathfrak X=X_{/X'}$、$\mathfrak Y=Y_{/Y'}$、$\mathfrak Z=Z_{/Z'}$。设 $f:X\to Y$ 为 $Z$-态射，$\widehat f:\mathfrak X\to\mathfrak Y$ 为其完备化。则 $\widehat f$ 为同构（分别地，为闭浸入），当且仅当存在 $Z'$ 的开邻域 $U\subset Z$，使 $f$ 的限制 $g^{-1}(U)\to h^{-1}(U)$ 为同构（分别地，为闭浸入）。
\eproof
充分性由（I，10.14.7）立即得到。必要性如（4.6.7），只须在 $Y'$ 中每点附近验证。因此可归结到 $Y=Z=\Spec(A)$ 仿射诺特。由假设，各 $y\in Y'$ 上的纤维至多为一点；\transnote{原文写作纤维“缩成一点”。在闭浸入的情形，纤维也可能为空；这里将这一措辞明确为“至多为一点”，所需的有限性不受影响。}$f$ 固有，由（4.4.11）可在 $Y'$ 的邻域上使 $f$ 有限，写成 $X=\Spec(B)$，其中 $B$ 为有限 $A$-代数。

若 $Y'=V(\frj)$，则 $\mathfrak Y=\Spf(\widehat A)$、$\mathfrak X=\Spf(\widehat B)$，其中 $\widehat A$ 为 $\frj$-预进制分离完备化，$\widehat B$ 为 $\frj B$-预进制分离完备化，也就是把 $B$ 视为 $A$-模时的 $\frj$-预进制分离完备化。$\widehat f$ 对应于环同态 $\varphi:A\to B$ 的连续延拓 $\widehat\varphi:\widehat A\to\widehat B$；假设表示它双射（分别地，满射）。它也是 $A$-模同态 $\varphi$ 的完备化。由（I，10.8.14），存在 $Y'$ 的开邻域 $U$，使层同态 $\wt\varphi:\wt A\to\wt B$ 在 $U$ 上双射（分别地，满射），从而得到结论。\finishproof\par

% ---- END c3_s04d.tex ----
% ---- BEGIN c3_s04e.tex ----
\subsection{一个丰沛性判据}
\src{145--146}
\egastatement{4.7.1}{定理}
设 $Y$ 局部诺特，$f:X\to Y$ 固有，$\LL$ 为可逆 $\OO_X$-模，$y\in Y$，$X_y=X\otimes_Yk(y)$，$g:X_y\to X$ 为投影。若 $\LL_y=g^*\LL=\LL\otimes_{\OO_Y}k(y)$ 在 $X_y$ 上丰沛，则存在 $y$ 的开邻域 $U$，使 $\LL|f^{-1}(U)$ 对于限制态射 $f^{-1}(U)\to U$ 丰沛。
\eproof
\textit{I）} 先令 $Y'=\Spec(\OO_{Y,y})$、$X'=X\times_YY'$、$\LL'=\LL\otimes_{\OO_Y}\OO_{Y,y}$，证明 $\LL'$ 对 $f'=f_{(Y')}$ 丰沛。有交换图
\[
\begin{tikzcd}
X\arrow[d,"f"']&X'\arrow[l]\arrow[d,"f'"]&X_y\arrow[l]\arrow[d]\\
Y&Y'\arrow[l]&\Spec(k(y))\arrow[l].
\end{tikzcd}
\]
$f'$ 固有、$\OO_{Y,y}$ 诺特，故只须在 $Y=\Spec(\OO_{Y,y})$ 时，由纤维上的丰沛性证明 $\LL$ 丰沛（II，4.6.6）。我们使用（2.6.1，b），实际证明对每个凝聚 $\OO_X$-模 $\F$，存在 $N$ 使 $H^1(X,\F(n))=0$ 对 $n\ge N$ 成立，其中 $\F(n)=\F\otimes\LL^{\otimes n}$。

记 $\mm$ 为 $\OO_{Y,y}$ 的极大理想。$X_y$ 为凝聚理想 $\J=\mm\OO_X$ 定义的闭子预概形，$g$ 为其嵌入。令有限型分次 $k(y)$-代数 $S=\gr(\OO_{Y,y})=\bigoplus_{j\ge0}\mm^j/\mm^{j+1}$。其拉回 $\SSS=f^*S$ 是拟凝聚有限型 $\OO_X$-代数，并被 $\J$ 零化；故若 $\SSS_y=g^*\SSS$，则 $\SSS=g_*\SSS_y$。再令
\[
\M_j=\mm^j\F/\mm^{j+1}\F,\qquad \M=\bigoplus_{j\ge0}\M_j=\gr(\F).
\]
它是被 $\J$ 零化的拟凝聚有限型 $\SSS$-模（0，10.1.1）。令 $\M'=g^*\M=\bigoplus_j\M'_j$，则 $\M=g_*\M'$；若 $\M'(n)=\M'\otimes\LL_y^{\otimes n}$，有 $\M(n)=g_*\M'(n)$。

$X_y$ 固有且 $\LL_y$ 丰沛，故 $f_y$ 射影（II，5.5.4、4.6.11）。把（2.4.1，ii）用于 $\Spec(k(y)),f_y,\SSS_y,\LL_y,\M'$，得到一个 $N$，使 $n\ge N$ 时，对所有 $q>0$ 和 $j$ 均有 $H^q(X_y,\M'_j(n))=0$。因此 $H^q(X,\M_j(n))=0$。令 $\F(n)_j=(\F/\mm^{j+1}\F)(n)$，则对 $j\ge1$，
\[
\Ker(\F(n)_j\to\F(n)_{j-1})=\M_j(n),\qquad \F(n)_0=\M_0(n).
\]
由上同调正合列和归纳，$H^1(X,\F(n)_j)=0$ 对所有 $j\ge0$ 成立；由（4.2.1），$H^1(X,\F(n))=0$，完成局部环情形。

\textit{II）} 回到最初记号，可设 $Y=\Spec(A)$ 仿射。因为 $f'$ 有限型、$\LL'$ 丰沛，存在 $m>0$，使 $\LL'^{\otimes m}$ 很丰沛（II，4.6.11）。将 $\LL$ 换成 $\LL^{\otimes m}$，可假设 $\LL'$ 很丰沛，并证明 $\LL|f^{-1}(U)$ 很丰沛。$f'$ 固有，故对某个 $r$ 存在闭浸入 $j:X'\to P=\Pb^r_{Y'}$，使 $\LL'\simeq j^*\OO_P(1)$；它对应满同态 $u:\OO_{X'}^{r+1}\to\LL'$（II，4.2.3），亦即生成 $\LL'$ 的 $r+1$ 个全局截面 $s'_0,\ldots,s'_r$。

这些也是 $f'_*\LL'$ 的全局截面，而 $f'_*\LL'=f_*\LL\otimes_{\OO_Y}\OO_{Y,y}$。各 $s'_i$ 可写成 $s''_i/t_i$，其中 $s''_i\in\Gamma(Y,f_*\LL)$、$t_i\notin\mathfrak j_y$。因此存在 $y$ 的仿射开邻域 $V$，以及 $f_*\LL|V$ 的截面 $s_i$，使其在 $Y'$ 上的限制为 $s'_i$。这些截面定义
\[
v:(\OO_X|f^{-1}(V))^{r+1}\longrightarrow\LL|f^{-1}(V),
\]
它在纤维 $f^{-1}(y)$ 上各点满射。余核凝聚，支集闭，故在纤维的某个开邻域 $W$ 上满射。$f$ 闭，可取 $W=f^{-1}(U)$，其中 $U$ 为 $y$ 的开邻域。原文于是由（II，4.2.3）得出结论。\transnote{满同态首先构造的是到射影空间的态射，尚须验证它在缩小后的邻域上为浸入。这里可用已证（4.6.7，i）的闭浸入情形：纤维上的态射来自 $j$ 的基变换，源与射影空间均在 $U$ 上固有。原文未写出这一步，本注仅标明所需补充。}\finishproof

\subsection{形式预概形的有限态射}
\src{146--149}
\egastatement{4.8.1}{命题}
设 $\mathfrak Y$ 为局部诺特形式预概形，$\K$ 为定义理想，$f:\mathfrak X\to\mathfrak Y$ 为形式预概形态射。以下条件等价：
\begin{enumerate}[label=\alph*)]
\item $\mathfrak X$ 局部诺特，$f$ 进制（I，10.12.1）；令 $\J=f^*\K\OO_{\mathfrak X}$，则 $f_0:(\mathfrak X,\OO_{\mathfrak X}/\J)\to(\mathfrak Y,\OO_{\mathfrak Y}/\K)$ 有限。
\item $\mathfrak X$ 局部诺特，且为某个进制 $(Y_n)$-归纳系 $(X_n)$ 的归纳极限，其中 $X_0\to Y_0$ 有限。
\item $\mathfrak Y$ 的每个点有诺特仿射形式开邻域 $V$，使 $f^{-1}(V)$ 仿射形式，且 $\Gamma(f^{-1}(V),\OO_{\mathfrak X})$ 是有限型 $\Gamma(V,\OO_{\mathfrak Y})$-模。
\end{enumerate}
\eproof
a）$\Rightarrow$b）由（I，10.12.3）立即得出。证明 b）$\Rightarrow$c）时，可设 $\mathfrak Y=\Spf(B)$，$B$ 诺特进制，$\K=\frk^\triangle$。$X_0$ 仿射，其环 $A_0$ 在 $B/\frk$ 上有限。由（I，5.1.9），每个 $X_n$ 仿射，记其环为 $A_n$，则 $m\le n$ 时 $A_m\simeq A_n/\frk^{m+1}A_n$。于是 $\mathfrak X\simeq\Spf(A)$，$A=\plim_nA_n$，由（$\OI$，7.2.9）得 c）。最后，证明 c）$\Rightarrow$a）时，可设 $\mathfrak Y=\Spf(B)$、$\mathfrak X=\Spf(A)$，$A$ 为有限 $B$-代数。$A/\frk A$ 在 $B/\frk$ 上有限，由（I，10.10.9）得 a）。\finishproof
\egastatement{4.8.2}{定义}
满足（4.8.1）的等价条件时，称 $f$ 为\textbf{有限态射}，或称 $\mathfrak X$ 为\textbf{有限的 $\mathfrak Y$-形式预概形}。
\egastatement{4.8.3}{命题}
（i）局部诺特形式预概形的闭浸入有限。

（ii）两个这样的有限态射的复合有限。

（iii）设 $\mathfrak X,\mathfrak Y,\mathfrak S$ 局部诺特，$f:\mathfrak X\to\mathfrak S$ 有限，$g:\mathfrak Y\to\mathfrak S$ 为态射，则 $\mathfrak X\times_{\mathfrak S}\mathfrak Y\to\mathfrak Y$ 有限。

（iv）设 $\mathfrak S,\mathfrak X',\mathfrak Y'$ 局部诺特，且 $\mathfrak X'\times_{\mathfrak S}\mathfrak Y'$ 局部诺特。若 $\mathfrak X,\mathfrak Y$ 是局部诺特的 $\mathfrak S$-形式预概形，$f:\mathfrak X\to\mathfrak X'$、$g:\mathfrak Y\to\mathfrak Y'$ 是有限 $\mathfrak S$-态射，则 $f\times_{\mathfrak S}g$ 有限。

（v）设 $f:\mathfrak Z\to\mathfrak Y$、$g:\mathfrak Y\to\mathfrak X$ 为局部诺特形式预概形态射，$g$ 分离且有限型。若 $g\circ f$ 有限，则 $f$ 有限。
\eproof
（i）显然。其余由（4.8.1，a）归结到通常预概形的对应命题（II，6.1.5）；依照（I，10.13.5）的方式，细节留给读者。\finishproof
\egastatement{4.8.4}{推论}
在（I，10.9.9）的条件下，若 $f$ 有限，则其到完备化的延拓 $\widehat f$ 亦有限。
\egastatement{4.8.5}{推论}
若 $\mathfrak X$ 在 $\mathfrak Y$ 上有限，结构态射为 $f$，则对任意开集 $U\subset\mathfrak Y$，$f^{-1}(U)$ 在 $U$ 上有限。
\egastatement{4.8.6}{命题}
若 $f:\mathfrak X\to\mathfrak Y$ 为局部诺特形式预概形的有限态射，则 $f_*\OO_{\mathfrak X}$ 是凝聚 $\OO_{\mathfrak Y}$-代数。
\eproof
将 $f$ 写成 $f_n:X_n\to Y_n$ 的归纳极限。各 $f_n$ 有限，且 $f_*\OO_{\mathfrak X}$ 是 $(f_n)_*\OO_{X_n}$ 的射影极限，即可由（I，10.10.5）得到结论。这些事实均可在 $\mathfrak Y=\Spf(B)$、$\mathfrak X=\Spf(A)$ 时验证：若 $\frk$ 是 $B$ 的定义理想，$A$ 在 $B$ 上有限，则 $A/\frk^{n+1}A$ 在 $B/\frk^{n+1}$ 上有限，而且 $A=\plim_nA/\frk^{n+1}A$。\finishproof
反过来，有如下结果。
\egastatement{4.8.7}{命题}
设 $\mathfrak Y$ 局部诺特，$\A$ 为凝聚 $\OO_{\mathfrak Y}$-代数。则存在在 $\mathfrak Y$ 上有限的形式预概形 $\mathfrak X$，在 $\mathfrak Y$-同构意义下唯一，使结构态射 $f:\mathfrak X\to\mathfrak Y$ 满足 $f_*\OO_{\mathfrak X}=\A$。
\eproof
取定义理想 $\K$，令 $Y_n=(\mathfrak Y,\OO_{\mathfrak Y}/\K^{n+1})$、$\A_n=\A/\K^{n+1}\A$。$\A_n$ 为有限 $\OO_{Y_n}$-代数，给出有限预概形 $X_n=\Spec(\A_n)$。对 $m\le n$，典范满同态 $h_{mn}:\A_n\to\A_m$ 给出交换图
\[
\begin{tikzcd}
X_m\arrow[r,"u_{nm}"]\arrow[d,"f_m"']&X_n\arrow[d,"f_n"]\\
Y_m\arrow[r]&Y_n,
\end{tikzcd}
\]
它将 $X_m$ 认同于 $X_n\times_{Y_n}Y_m$（II，1.4.6）。归纳极限 $\mathfrak X$ 局部诺特，结构态射 $f$ 有限（4.8.1；I，10.12.3.1）\transnote{此处原著将后一个引文印为 II，10.12.3.1；关于进制归纳系的对应内容在第I章，故改记 I。}，且由（4.8.6）的证明，$f_*\OO_{\mathfrak X}=\plim_n\A_n=\A$（I，10.10.6）。唯一性来自下述更一般命题。\finishproof
\egastatement{4.8.8}{命题}
设 $\mathfrak Y$ 局部诺特，$\mathfrak X,\mathfrak X'$ 在 $\mathfrak Y$ 上有限，结构态射分别为 $f,f'$。则有典范双射
\[
\Hom_{\mathfrak Y}(\mathfrak X,\mathfrak X')\simeq
\Hom_{\OO_{\mathfrak Y}}(f'_*\OO_{\mathfrak X'},f_*\OO_{\mathfrak X}).
\]
\origfootnote{后一式表示 $\OO_{\mathfrak Y}$-代数同态的集合。}
\eproof
映射 $h\mapsto\A(h)$ 与（II，1.1.2）定义相同。局部可设 $\mathfrak Y=\Spf(B)$，则 $\mathfrak X=\Spf(A)$、$\mathfrak X'=\Spf(A')$，其中 $A,A'$ 是有限 $B$-代数，直接像分别为 $A^\triangle,A'^\triangle$。$\mathfrak Y$-态射与必然连续的 $B$-代数同态 $A'\to A$ 一一对应（I，10.2.2）；另一方面，$B$-模同态与相伴凝聚模的同态一一对应（I，10.10.2.3）。兼顾乘法便得到结论。\finishproof
\egastatement{4.8.9}{推论}
在（4.8.8）的典范对应中，闭浸入 $\mathfrak X\to\mathfrak X'$ 对应于满的 $\OO_{\mathfrak Y}$-代数同态 $f'_*\OO_{\mathfrak X'}\to f_*\OO_{\mathfrak X}$。
\eproof
仍为局部问题，归结到形式预概形闭浸入的定义（I，10.14.2）。\finishproof
\egastatement{4.8.10}{推论}
在（4.8.1）的记号及假设下，进制态射 $f$ 为闭浸入，当且仅当 $f_0$ 为通常预概形的闭浸入。
\eproof
由（4.8.9）及凝聚模同态的满射判据（I，10.11.5）立即得出。\finishproof
\egastatement{4.8.11}{命题}
局部诺特形式预概形态射 $f:\mathfrak X\to\mathfrak Y$ 有限，当且仅当它固有，且所有底空间纤维 $f^{-1}(y)$ 有限。
\eproof
由定义（3.4.1、4.8.2），归结到 $f_0$ 的同一命题，即（4.4.2）。\finishproof\par

% ---- END c3_s04e.tex ----
% ---- BEGIN c3_s05a.tex ----
\section{凝聚代数层的一个存在性定理}
\subsection{定理的陈述}
\src{149--151}
\egapar{5.1.1}
设 $A$ 为诺特进制环，$\frj$ 为其定义理想，故 $A$ 对 $\frj$-进制拓扑分离完备。若 $Y=\Spec(A)$，则 $\Spf(A)$ 认同于 $Y$ 沿 $Y'=V(\frj)$ 的完备化 $\mathfrak Y$（I，10.10.1）。设 $X$ 为有限型通常 $Y$-预概形，结构态射为 $f$。以 $\mathfrak X$ 表示沿 $X'=f^{-1}(Y')$ 的完备化，也就是 $\mathfrak Y$-形式预概形 $X\times_Y\mathfrak Y$；以 $\widehat f:\mathfrak X\to\mathfrak Y$ 表示 $f$ 的完备化。对于凝聚 $\OO_X$-模 $\F$，记其完备化为 $\widehat\F=\F_{/X'}$。
\egastatement{5.1.2}{命题}
在（5.1.1）的条件下，若凝聚 $\OO_X$-模 $\F$ 的支集在 $Y$ 上固有，则（4.1.4）的典范同态
\[
\rho_n:H^n(X,\F)\longrightarrow H^n(\mathfrak X,\widehat\F)
\]
都是同构。
\eproof
$H^n(X,\F)$ 为有限型 $A$-模（3.2.4），故已对 $\frj$-预进制拓扑分离完备（$\OI$，7.3.6）；于是应用（4.1.10）即可。\finishproof
对于凝聚模的任意短正合列，这些典范同构与上同调连接同态交换（0，12.1.6；I，10.8.9）。
\egastatement{5.1.3}{推论}
设 $\F,\G$ 为凝聚 $\OO_X$-模，支集的交在 $Y$ 上固有。则典范同态
\begin{equation}\tag{5.1.3.1}
\Hom_{\OO_X}(\F,\G)\longrightarrow
\Hom_{\OO_{\mathfrak X}}(\widehat\F,\widehat\G),\qquad u\longmapsto\widehat u
\end{equation}
为同构。若 $f$ 还闭，则 $u$ 单射（分别地，满射）当且仅当 $\widehat u$ 单射（分别地，满射）。
\eproof
第一断言由（4.5.3）及同态模已完备得到。由（I，10.8.14），$\widehat u$ 单射或满射，当且仅当 $u$ 在 $X'$ 的某个邻域中如此；于是第二断言来自以下引理。\finishproof
\egastatement{5.1.3.1}{引理}
在（5.1.1）的条件下，若 $f$ 闭，则 $X'$ 在 $X$ 中的每个邻域都等于 $X$。
\eproof
可先归结到 $f(X)=Y$。事实上，$f(X)=Z$ 闭；把 $f$ 换成 $f_{\red}$，可设 $X,Y$ 约化，再把 $Y$ 换成以 $Z$ 为底空间的约化闭子预概形（I，5.2.2）。因为 $A$ 中每个理想闭，其商仍为诺特进制环，这种替换不改变所需条件。此时 $f(X')=Y'$。若开邻域 $V$ 包含 $X'$，则闭集 $f(X-V)$ 与 $Y'$ 不交；由于 $\frj$ 包含在 $A$ 的根基中（I，1.1.15；$\OI$，7.1.10），除非 $X-V=\varnothing$，这是不可能的。\finishproof

对于支集在 $Y$ 上固有的凝聚模，（5.1.3）表示完备化函子 $\F\mapsto\widehat\F$ 是到凝聚 $\OO_{\mathfrak X}$-模范畴的全忠实函子，因此把原范畴等价到后者的一个满子范畴（0，8.1.6）。存在性定理将说明，当 $X$ 在 $Y$ 上固有时，这个满子范畴就是全部凝聚模范畴。确切地说：
\egastatement{5.1.4}{定理}
设 $A$ 为诺特进制环，$Y=\Spec(A)$，$\frj$ 为定义理想，$Y'=V(\frj)$，$f:X\to Y$ 分离且有限型，$X'=f^{-1}(Y')$。令 $\mathfrak Y=\Spf(A)$、$\mathfrak X=X_{/X'}$，$\widehat f$ 为 $f$ 的完备化。则函子 $\F\mapsto\widehat\F$ 给出如下范畴等价：一边是支集在 $\Spec(A)$ 上固有的凝聚 $\OO_X$-模，另一边是支集在 $\Spf(A)$ 上固有的凝聚 $\OO_{\mathfrak X}$-模。

结合（5.1.3），也就是：
\egastatement{5.1.5}{推论}
一个 $\OO_{\mathfrak X}$-模同构于某个支集在 $\Spec(A)$ 上固有的凝聚 $\OO_X$-模的完备化，当且仅当它凝聚且支集在 $\Spf(A)$ 上固有。
最重要的特殊情形是：
\egastatement{5.1.6}{推论}
若 $X$ 在 $Y=\Spec(A)$ 上固有，则完备化函子给出凝聚 $\OO_X$-模与凝聚 $\OO_{\mathfrak X}$-模范畴的等价。
\egastatement{5.1.7}{附释}
由形式预概形上凝聚层的刻画（I，10.11.3），令 $Y_n=\Spec(A/\frj^{n+1})$、$X_n=X\times_YY_n$，则上述凝聚代数层的给定等价于一个凝聚 $\OO_{X_n}$-模的射影系 $(\F_n)$，使 $m\le n$ 时
\[
\F_m=\F_n\otimes_{\OO_{Y_n}}\OO_{Y_m}
=\F_n/\frj^{m+1}\F_n,
\]
而 $\Supp(\F_0)$ 在 $Y_0$ 上固有。由（I，10.11.4），同态同样解释为这些射影系的同态。

在当时已知的全部应用中，$A$ 实际上是诺特进制局部环，各 $Y_n$ 是阿廷局部环的谱。因此，本节及前几节的结果，在很大程度上把诺特进制局部环上的代数几何归结到阿廷局部环上的代数几何。\transnote{原文此句描述的是作者所考虑的应用。一般“局部进制环”并不单凭局部性保证每个 $A/\frj^{n+1}$ 阿廷；这里相应的应用取极大理想进制，或至少取极大理想准素的定义理想。}
\egastatement{5.1.8}{推论}
在（5.1.4）的条件下，$Z\mapsto\widehat Z=Z_{/(Z\cap X')}$ 给出以下集合之间的双射：$X$ 中在 $Y$ 上固有的闭子预概形，与 $\mathfrak X$ 中在 $\mathfrak Y$ 上固有的闭形式子预概形。
\eproof
闭形式子预概形写作 $(T,(\OO_{\mathfrak X}/\J)|T)$，其中 $\J$ 凝聚理想（I，10.14.2）。若 $T$ 固有，则由（5.1.4），$\OO_{\mathfrak X}/\J\simeq\widehat\F$，其中 $\F$ 为固有支集的凝聚模。由（5.1.3），典范商同态是某个满同态 $u:\OO_X\to\F$ 的完备化。故 $\F=\OO_X/\I$，$\I$ 为凝聚理想，而且 $\J=\widehat\I$（I，10.8.8）；由（I，10.14.7）得到结论。\finishproof

\subsection{存在性定理的证明：射影与拟射影情形}
\src{151--154}
\egapar{5.2.1}
在（5.1.4）的条件下，暂称凝聚 $\OO_{\mathfrak X}$-模\textbf{可代数化}，是指它同构于某个固有支集的凝聚 $\OO_X$-模的完备化。
\egastatement{5.2.2}{引理}
若 $\mathfrak F,\mathfrak F'$ 可代数化，则任意同态 $u:\mathfrak F\to\mathfrak F'$ 的核、像及余核均可代数化。
\eproof
写 $\mathfrak F=\widehat\F$、$\mathfrak F'=\widehat\G$。由（5.1.3），$u=\widehat v$，其中 $v:\F\to\G$。完备化函子正合，故 $\Ker u\simeq(\Ker v)^\wedge$。$\Ker v$ 的支集包含于 $\F$ 的支集，故固有，得到核的可代数化。像与余核同理。\finishproof
\egastatement{5.2.3}{命题}
设 $A$ 为诺特进制环，$\frj$ 为定义理想，$\mathfrak Y=\Spf(A)$，$f:\mathfrak X\to\mathfrak Y$ 为固有形式预概形态射。令 $Y_k=\Spec(A/\frj^{k+1})$、$X_k=\mathfrak X\times_{\mathfrak Y}Y_k$；对任意 $\OO_{\mathfrak X}$-模 $\F$，令 $\F_k=\F/\frj^{k+1}\F$。设 $\LL$ 可逆，且 $\LL_0=\LL/\frj\LL$ 在 $X_0$ 上丰沛，令 $\F(n)=\F\otimes\LL^{\otimes n}$。则对每个凝聚 $\F$，存在 $n_0$，使所有 $n\ge n_0$ 满足：

（i）$H^0(\mathfrak X,\F(n))\to H^0(\mathfrak X,\F_k(n))$ 对所有 $k\ge0$ 满射；

（ii）$H^q(\mathfrak X,\F(n))=0$ 对所有 $q>0$ 成立。
\eproof
$\mathfrak X$ 与 $X_0$ 底空间相同。$\M_k=\frj^k\F/\frj^{k+1}\F$ 被 $\frj$ 零化，可看作凝聚 $\OO_{X_0}$-模；令 $\M_k(n)=\M_k\otimes\LL_0^{\otimes n}$，则 $\M_k(n)=\frj^k\F(n)/\frj^{k+1}\F(n)$。因为 $f_0:X_0\to Y_0$ 固有、$\LL_0$ 丰沛，$f_0$ 射影（II，5.5.4）。令 $S=\gr_{\frj}(A)$；$A$ 诺特，故 $S$ 有限型。其拉回 $\SSS=f_0^*S$ 拟凝聚有限型，而 $\M=\bigoplus_k\M_k$ 是拟凝聚有限型分次 $\SSS$-模，因为 $\M_0$ 凝聚且生成 $\M$。由（2.4.1，ii），存在 $n_0$，使
\begin{equation}\tag{5.2.3.1}
H^q(X_0,\M_k(n))=0\qquad(q>0,\ n\ge n_0,\ k\ge0).
\end{equation}
把 $\M_k(n)$ 看作 $\mathfrak X$ 上的模，仍有同样消失。对
\[
0\to\frj^{k-1}\F(n)/\frj^k\F(n)\to
\frj^h\F(n)/\frj^k\F(n)\to
\frj^h\F(n)/\frj^{k-1}\F(n)\to0
\]
应用上同调正合列，并对 $k-h$ 归纳，得到
\begin{equation}\tag{5.2.3.2}
H^q(\mathfrak X,\frj^h\F(n)/\frj^k\F(n))=0
\quad(0\le h<k,\ q>0,\ n\ge n_0).
\end{equation}
特别地，$h=0$ 给出
\begin{equation}\tag{5.2.3.3}
H^q(\mathfrak X,\F_k(n))=0\qquad(q>0,\ n\ge n_0,\ k\ge0).
\end{equation}
正合列的另一段为
\begin{equation}\tag{5.2.3.4}
H^0(\mathfrak X,\F_h(n))\to H^0(\mathfrak X,\F_k(n))
\to H^1(\mathfrak X,\frj^{k+1}\F(n)/\frj^{h+1}\F(n))=0
\end{equation}
（$h\ge k$），故
\begin{equation}\tag{5.2.3.5}
H^0(\mathfrak X,\F_h(n))\longrightarrow H^0(\mathfrak X,\F_k(n))
\end{equation}
满射。因此对每个 $q$，$(H^q(\mathfrak X,\F_k(n)))_k$ 满足（ML）。又 $\mathfrak X$ 的仿射形式开集在各 $X_k$ 上仿射，其高阶上同调消失，而零阶截面间的商映射满射。由（0，13.3.1），
\[
H^q(\mathfrak X,\F(n))\simeq\plim_kH^q(\mathfrak X,\F_k(n)).
\]
$q>0$ 时右端为零，得到（ii）；$q=0$ 时，因过渡映射满射，极限到每一项亦满射，得到（i）。\finishproof
\egastatement{5.2.4}{推论}
在（5.2.3）的假设下，对于每个凝聚 $\OO_{\mathfrak X}$-模 $\F$，存在 $N$，使 $n\ge N$ 时 $\F(n)$ 由全局截面生成。换言之，$\F$ 同构于某个 $(\LL^{\otimes(-n)})^k$ 的商。
\eproof
$X_0$ 诺特、$\LL_0$ 丰沛，故充分大的 $n$ 使 $\F_0(n)$ 由全局截面生成（II，4.5.5）。再令 $n$ 充分大，使 $\Gamma(\mathfrak X,\F(n))\to\Gamma(\mathfrak X,\F_0(n))$ 满射（5.2.3）。可取有限个 $\F(n)$ 的截面，其像生成 $\F_0(n)$（$\OI$，5.2.3）。$\frj$ 在每个茎环中都包含于极大理想，故由 Nakayama 引理，这些截面生成 $\F(n)$。\finishproof
\egapar{5.2.5}
\textbf{存在性定理的证明：射影情形。}
采用（5.1.4）的记号，设 $f$ 射影，取丰沛 $\OO_X$-模 $\LL$。各 $X_k=X\times_YY_k$ 也是 $\mathfrak X$ 的相应闭子预概形。若 $\LL'$ 为 $\LL$ 的完备化，则 $\LL'_0=\LL/\frj\LL$ 在 $X_0$ 上丰沛（II，4.6.13，i bis）。故可对任意凝聚 $\OO_{\mathfrak X}$-模 $\mathfrak F$ 应用（5.2.4），把它写成某个 $\mathfrak G=(\LL'^{\otimes(-n)})^k$ 的商。$\mathfrak G$ 是 $(\LL^{\otimes(-n)})^k$ 的完备化（I，10.8.10），故可代数化。取满射 $u:\mathfrak G\to\mathfrak F$，其核凝聚；再次应用同一结果，存在可代数化模 $\mathfrak H$ 及同态 $v:\mathfrak H\to\mathfrak G$，使 $\Im v=\Ker u$。于是 $\mathfrak F=\Coker v$，由（5.2.2）可代数化。
\egapar{5.2.6}
\textbf{存在性定理的证明：拟射影情形。}
若 $f$ 拟射影，可把 $X$ 认同为某个射影 $Y$-预概形 $g:Z\to Y$ 的开子预概形。令 $Z'=g^{-1}(Y')$，则 $X'=X\cap Z'$，而 $\mathfrak X$ 是 $\mathfrak Z=Z_{/Z'}$ 在开集 $X\cap Z'$ 上的诱导形式预概形（I，10.8.5）。

取固有支集 $T'$ 的凝聚 $\OO_{\mathfrak X}$-模 $\mathfrak F$。按定义，$T'$ 作为 $X'$ 的子集在 $Y'$ 上固有，因而在 $Y$ 上固有，故在 $Z'$ 中闭（II，5.4.10）。于是可将 $\mathfrak F$ 与 $\mathfrak Z-T'$ 上的零层粘合成凝聚 $\OO_{\mathfrak Z}$-模 $\mathfrak F'$；二者在交集 $\mathfrak X-T'$ 上相同。由射影情形，$\mathfrak F'=\widehat\G$，其中 $\G$ 为凝聚 $\OO_Z$-模。设 $T=\Supp(\G)$，则 $T'=T\cap Z'$（I，10.8.12）。把 $g$ 限制到以 $T$ 为底空间的约化闭子预概形，记为 $h$；则 $T'=h^{-1}(Y')$，而 $X\cap T$ 是 $T$ 中包含 $T'$ 的开集。$h$ 固有，故闭，由（5.1.3.1）知 $X\cap T=T$，亦即 $T\subset X$。$T$ 在 $Z$ 中闭，因此在 $Y$ 上固有。令 $\F=\G|X$，则其完备化是 $\widehat\G$ 在 $\mathfrak X$ 上的限制，即 $\mathfrak F$，从而完成证明。\par

% ---- END c3_s05a.tex ----
% ---- BEGIN c3_s05b.tex ----
\subsection{存在性定理的证明：一般情形}
\src{154--156}
\egastatement{5.3.1}{引理}
在（5.1.4）的条件下，若凝聚 $\OO_{\mathfrak X}$-模的正合列 $0\to\mathfrak H\to\mathfrak F\to\mathfrak G\to0$ 中，$\mathfrak H,\mathfrak G$ 可代数化，则 $\mathfrak F$ 可代数化。
\eproof
写 $\mathfrak G=\widehat\B$、$\mathfrak H=\widehat\C$，其中 $\B,\C$ 凝聚且支集在 $Y$ 上固有。此扩张定义 $\Ext^1_{\OO_{\mathfrak X}}(\mathfrak X;\mathfrak G,\mathfrak H)$ 中的元素（0，12.3.2），该模由（4.5.2）典范地同构于 $\Ext^1_{\OO_X}(X;\B,\C)$。故存在凝聚模正合列 $0\to\C\to\A\to\B\to0$，使其扩张类映到给定扩张类。依定义并利用（I，10.8.8，ii），这表示 $\mathfrak F\simeq\widehat\A$。又 $\Supp(\A)$ 包含在另外两个支集的并中，故固有，引理成立。\finishproof
\egastatement{5.3.2}{推论}
在（5.1.1）的条件下，设 $u:\mathfrak F\to\mathfrak G$ 为凝聚 $\OO_{\mathfrak X}$-模同态。若 $\mathfrak G,\Ker u,\Coker u$ 可代数化，则 $\mathfrak F$ 也可代数化。
\eproof
由（5.2.2）用于 $\mathfrak G\to\Coker u$，得 $\Im u$ 可代数化；再对 $0\to\Ker u\to\mathfrak F\to\Im u\to0$ 应用（5.3.1）。\finishproof
\egastatement{5.3.3}{引理}
在（5.1.1）的条件下，设 $h:Z\to Y$ 有限型，$\mathfrak Z$ 为沿 $h^{-1}(Y')$ 的完备化，$g:Z\to X$ 为固有 $Y$-态射，$\widehat g:\mathfrak Z\to\mathfrak X$ 为其完备化。则可代数化 $\OO_{\mathfrak Z}$-模 $\mathfrak F$ 的直接像 $\widehat g_*\mathfrak F$ 可代数化。
\eproof
若 $\mathfrak F=\widehat\F$，则第一比较定理（4.1.5）给出 $\widehat g_*\mathfrak F\simeq(g_*\F)^\wedge$，其支集仍固有。\finishproof
\egastatement{5.3.4}{引理}
设 $X$ 为诺特通常概形，$X'\subset X$ 为闭集，$f:Z\to X$ 固有，$Z'=f^{-1}(X')$，$\mathfrak X=X_{/X'}$、$\mathfrak Z=Z_{/Z'}$，$\widehat f$ 为完备化。设 $\M\subset\OO_X$ 为凝聚理想，令 $U=X-\Supp(\OO_X/\M)$，并假设 $f^{-1}(U)\to U$ 为同构。则对每个凝聚 $\OO_{\mathfrak X}$-模 $\mathfrak F$，存在 $n>0$，使典范同态
\[
\rho_{\mathfrak F}:\mathfrak F\longrightarrow
\widehat f_*(\widehat f^{\,*}\mathfrak F)
\]
的核及余核都被 $\widehat\M^{\,n}$ 零化。
\eproof
可设 $X=\Spec(B)$，$B$ 诺特，$X'=V(\frk)$。先归结到 $B$ 为诺特进制环、$\frk$ 为定义理想。令 $B_1$ 为 $B$ 的 $\frk$-预进制分离完备化，$\frk_1=\frk B_1$，$X_1=\Spec(B_1)$，$h:X_1\to X$ 为典范态射，$X'_1=h^{-1}(X')=V(\frk_1)$。令 $Z_1=Z\times_XX_1$、$f_1=f_{(X_1)}$，它固有。把 $X_1,Z_1$ 分别沿 $X'_1$ 及其逆像完备化，得到 $\mathfrak X_1,\mathfrak Z_1$ 和 $\widehat f_1$。

$h$ 的完备化 $\widehat h:\mathfrak X_1\to\mathfrak X$ 是由 $B_1$ 的恒等给出的同构；故相应 $\mathfrak Z_1\to\mathfrak Z$ 也为同构，且它们认同 $\widehat f_1,\widehat f$。令 $\M_1=h^*\M\OO_{X_1}$，则
\[
\Supp(\OO_{X_1}/\M_1)=h^{-1}(\Supp(\OO_X/\M))
\]
（I，9.1.13），故 $U_1=h^{-1}(U)$，而 $f_1^{-1}(U_1)\to U_1$ 为同构。两理想的完备化经 $\widehat h$ 认同（I，10.9.5）。因此全部假设在这一替换后保持，可以设 $B$ 已进制完备。

这时 $\mathfrak X=\Spf(B)$，$\mathfrak F=N^\triangle$，其中 $N$ 为有限型 $B$-模。令 $\F=\wt N$，则 $\mathfrak F=\widehat\F$，且
\[
\widehat f^{\,*}\mathfrak F=(f^*\F)^\wedge,
\qquad \widehat f_*((f^*\F)^\wedge)\simeq(f_*f^*\F)^\wedge
\]
（I，10.9.5；4.1.5）。由（5.1.3），$\rho_{\mathfrak F}$ 就是 $\rho_\F:\F\to f_*f^*\F$ 的完备化。其核 $\B$ 和余核 $\C$ 凝聚，并在 $U$ 上为零，故存在 $n>0$ 使 $\M^n\B=\M^n\C=0$（I，9.3.4）。完备化后得到 $\widehat\M^{\,n}\widehat\B=\widehat\M^{\,n}\widehat\C=0$（I，10.8.8、10.8.10），引理得证。\finishproof
\egapar{5.3.5}
\textbf{存在性定理证明的完成。}
在（5.1.4）的条件下，使用诺特归纳原理，假设对于 $X$ 的每个底空间为真闭子集的闭子预概形 $T$，定理已成立；其完备化取沿 $T\cap X'$ 的完备化。可设 $X$ 非空。$f$ 分离且有限型，由 Chow 引理（II，5.6.1），存在 $Y$-概形 $Z$ 和 $Y$-态射 $g:Z\to X$，使结构态射 $h:Z\to Y$ 拟射影，$g$ 射影满射，而且存在非空开集 $U\subset X$，使 $g^{-1}(U)\to U$ 为同构。

取定义底空间 $X-U$ 的凝聚理想 $\M$。设 $\mathfrak F$ 为凝聚 $\OO_{\mathfrak X}$-模，其支集 $E$ 在 $\mathfrak Y$ 上固有。令 $\mathfrak Z$ 为 $Z$ 沿 $h^{-1}(Y')$ 的完备化，$\widehat g$ 为完备化态射。$\widehat g^{\,*}\mathfrak F$ 凝聚，支集包含在 $\widehat g^{-1}(E)$ 中，因 $\widehat g$ 固有，故该支集在 $\mathfrak Y$ 上固有。$h$ 拟射影，由（5.2.6），$\widehat g^{\,*}\mathfrak F$ 可代数化；由（5.3.3），$\widehat g_*\widehat g^{\,*}\mathfrak F$ 亦可代数化。

现在应用（5.3.4）：$\rho_{\mathfrak F}$ 的核 $\mathfrak B$ 和余核 $\mathfrak C$ 被 $\widehat\M^{\,n}$ 零化。\transnote{原文此处引用（5.2.4），但用于零化核与余核的是紧前的引理（5.3.4）。}
令 $T$ 为 $\M^n$ 定义的闭子预概形，$j:T\to X$ 为嵌入，$\widehat j:\mathfrak T\to\mathfrak X$ 为其完备化。因此
\[
\mathfrak B=\widehat j_*\widehat j^{\,*}\mathfrak B,
\qquad \mathfrak C=\widehat j_*\widehat j^{\,*}\mathfrak C.
\]
$U$ 非空，故 $T$ 的底空间为真闭子集，由归纳假设，$\widehat j^{\,*}\mathfrak B$、$\widehat j^{\,*}\mathfrak C$ 可代数化；再由（5.3.3），$\mathfrak B,\mathfrak C$ 可代数化。最后应用（5.3.2），得到 $\mathfrak F$ 可代数化。证毕。

\subsection{应用：通常概形态射与形式概形态射的比较；可代数化形式概形}
\src{156--159}
\egastatement{5.4.1}{定理}
设 $A$ 为诺特进制环，$\frj$ 为定义理想，$S=\Spec(A)$，$S'=V(\frj)$。设 $u:X\to S$ 固有，$v:Y\to S$ 分离且有限型，并令 $\mathfrak S,\mathfrak X,\mathfrak Y$ 分别为沿 $S',u^{-1}(S'),v^{-1}(S')$ 的完备化。对 $S$-态射 $f:X\to Y$，记完备化为 $\widehat f$。则 $f\mapsto\widehat f$ 是双射
\[
\Hom_S(X,Y)\simeq\Hom_{\mathfrak S}(\mathfrak X,\mathfrak Y).
\]
\eproof
先证单射。若 $\widehat f=\widehat g$，则由（I，10.9.4），$f,g$ 在 $X'=u^{-1}(S')$ 的某个开邻域 $V$ 上相同。$u$ 闭，由（5.1.3.1）知 $V=X$，故 $f=g$。

再证满射。取 $\mathfrak S$-态射 $h:\mathfrak X\to\mathfrak Y$。令 $Z=X\times_SY$，典范投影为 $p:Z\to X$、$q:Z\to Y$。它在 $S$ 上有限型，故诺特。其沿 $Z'=p^{-1}(u^{-1}(S'))$ 的完备化 $\mathfrak Z$ 典范地等于 $\mathfrak X\times_{\mathfrak S}\mathfrak Y$，投影对应 $\widehat p,\widehat q$（I，10.9.7）。$Y$ 在 $S$ 上分离，故 $\mathfrak Y$ 在 $\mathfrak S$ 上分离（I，10.15.7）；图态射 $\Gamma_h=(1_{\mathfrak X},h):\mathfrak X\to\mathfrak Z$ 为闭浸入（I，10.15.4）。设它的像为闭形式子预概形 $\mathfrak T$，嵌入为 $j$，写 $\Gamma_h=j\circ w$，其中 $w:\mathfrak X\to\mathfrak T$ 为同构，其逆为 $\widehat p\circ j$。

$\mathfrak T$ 同构于 $\mathfrak X$，故在 $\mathfrak S$ 上固有。由（5.1.8），存在在 $S$ 上固有的闭子预概形 $T\subset Z$，使 $\mathfrak T=\widehat T$，且 $j=\widehat i$，其中 $i:T\to Z$ 为嵌入。$p\circ i:T\to X$ 的完备化为同构；由（4.6.8），并用 $S$ 是 $S'$ 的唯一邻域，可知 $p\circ i$ 本身为同构。设其逆为 $g:X\to T$，令 $f=q\circ i\circ g$。则 $\Gamma_f=i\circ g$，而
\[
\widehat{\Gamma_f}=\widehat i\circ\widehat g=j\circ w=\Gamma_h.
\]
又 $\widehat{\Gamma_f}=\Gamma_{\widehat f}$（I，10.9.8），故 $h=\widehat f$。\finishproof

因此，对每个诺特进制环 $A$，完备化是从 $\Spec(A)$ 上固有概形范畴到 $\Spf(A)$ 上固有形式概形范畴的全忠实函子（0，8.1.6）。其像中对象称为\textbf{可代数化形式概形}。对于这样的 $\mathfrak X$，存在 $\Spec(A)$ 上固有的通常概形 $X$，其完备化同构于 $\mathfrak X$；连同这一认同，$X$ 在唯一同构意义下确定。
\egastatement{5.4.2}{附释}
令 $S_n=\Spec(A/\frj^{n+1})$、$X_n=X\times_SS_n$、$Y_n=Y\times_SS_n$。由（5.4.1）及（I，10.12.3），给定 $S$-态射 $f:X\to Y$，等价于给定 $S_n$-态射 $f_n:X_n\to Y_n$ 的一个进制 $(S_n)$-归纳系（I，10.12.2）。
\egastatement{5.4.3}{注记}
与存在性定理（5.1.6）可能给人的印象不同，存在 $\Spf(A)$ 上固有但不可代数化的形式概形，正如存在不来自复代数簇的紧解析空间一样。我们以后将在“模理论”中遇到这样的形式概形；在基域为 $\mathbf C$ 时，该理论恰研究完备代数簇的复结构的无穷小变化，而这种变化可以产生非代数的解析簇。
\egastatement{5.4.4}{命题}
设 $A$ 为诺特进制环，$\mathfrak S=\Spf(A)$，$g:\mathfrak X\to\mathfrak S$、$h:\mathfrak Y\to\mathfrak S$ 为固有形式概形态射，$f:\mathfrak X\to\mathfrak Y$ 为 $\mathfrak S$-态射。若 $f$ 有限且 $\mathfrak Y$ 可代数化，则 $\mathfrak X$ 可代数化。
\eproof
首先 $f$ 已固有；若其纤维有限，便足以保证有限性（4.8.11）。由（4.8.6），$\mathfrak B=f_*\OO_{\mathfrak X}$ 为凝聚 $\OO_{\mathfrak Y}$-代数。写 $\mathfrak Y=\widehat Y$，其中 $w:Y\to\Spec(A)$ 为固有通常概形，则存在凝聚 $\OO_Y$-代数 $\B$，使 $\mathfrak B=\widehat\B$。令 $X=\Spec(\B)$、$u:X\to Y$ 为结构态射。由（4.8.7）的构造，$\widehat X\simeq\mathfrak X$，且 $\widehat u$ 对应 $f$；只须在 $Y$ 仿射时核对即可。\finishproof
（5.1.8）是本命题的特殊情形。
\egastatement{5.4.5}{定理}
设 $A$ 为诺特进制环，$\frj$ 为定义理想，$S=\Spec(A)$、$\mathfrak S=\Spf(A)$，$f:\mathfrak X\to\mathfrak S$ 为固有形式概形态射。令 $S_k=\Spec(A/\frj^{k+1})$、$X_k=\mathfrak X\times_{\mathfrak S}S_k$，并对模 $\F$ 写 $\F_k=\F/\frj^{k+1}\F$。若可逆 $\OO_{\mathfrak X}$-模 $\LL$ 的 $\LL_0$ 在 $X_0$ 上丰沛，则 $\mathfrak X$ 可代数化；若固有 $S$-概形 $X$ 的完备化同构于 $\mathfrak X$，则存在丰沛 $\OO_X$-模 $\M$，使 $\LL\simeq\widehat\M$，因此 $X$ 在 $S$ 上射影。
\eproof
对 $\F=\OO_{\mathfrak X}$ 应用（5.2.3），得到充分大的 $n$ 使 $\Gamma(\mathfrak X,\LL^{\otimes n})\to\Gamma(X_0,\LL_0^{\otimes n})$ 满射。再取 $n$ 使 $\LL_0^{\otimes n}$ 对 $S_0$ 很丰沛（II，4.5.10）。因 $f_0$ 固有，右端为有限型 $A$-模（3.2.1），故可取左端的有限型 $A$-子模 $E$，其像为整个右端。

对 $k\ge0$，嵌入 $E\to\Gamma(\mathfrak X,\LL^{\otimes n})$ 给出
\[
u_k:E/\frj^{k+1}E\longrightarrow\Gamma(X_k,\LL_k^{\otimes n}).
\]
因为 $(f_k)_*\LL_k^{\otimes n}$ 拟凝聚，且其在 $S_k$ 上的截面等于上述截面，得到层同态 $\wt u_k$，进而得到
\[
\wt u_k^\sharp:f_k^*((E/\frj^{k+1}E)^\sim)\longrightarrow\LL_k^{\otimes n}.
\]
记左端为 $\G_k$，则 $\G_0=\G_k/\frj\G_k$，且 $\wt u_0^\sharp$ 是 $\wt u_k^\sharp$ 的模 $\frj$ 商。由 $E$ 的选取及很丰沛性，$\wt u_0^\sharp$ 是相应的满评价同态（II，4.4.3）。由 Nakayama 引理，各 $\wt u_k^\sharp$ 都满射，故各自定义 $S_k$-态射
\[
g_k:X_k\longrightarrow P_k=\Pb(E/\frj^{k+1}E)
\]
（II，4.2.2）。对 $h\le k$，$P_h=P_k\times_{S_k}S_h$；各同态间的相容关系使 $(g_k)$ 构成进制归纳系（II，4.2.10）。因此得到形式概形态射 $g:\mathfrak X\to\mathfrak P$，其中 $\mathfrak P=\widehat P$、$P=\Pb(E)$。

因为 $\LL_0^{\otimes n}$ 很丰沛，$g_0$ 为闭浸入，故 $g$ 为闭浸入（4.8.10）；由（5.1.8），$\mathfrak X$ 可代数化。存在性定理（5.1.6）给出 $\LL\simeq\widehat\M$，其中 $\M$ 为可逆 $\OO_X$-模。于是 $\LL^{\otimes n}$ 是 $\M^{\otimes n}$ 的完备化，各 $\wt u_k^\sharp$ 定义唯一同态
\[
v:f_X^*\wt E\longrightarrow\M^{\otimes n}
\]
（5.1.7），其中 $f_X:X\to S$ 为结构态射。由（I，10.11.5），$\widehat v$ 满射；由（5.1.3），$v$ 满射。其定义的 $r:X\to P$ 的完备化是 $g$；$g$ 闭浸入，故由（5.1.8）、（5.4.1），$r$ 也闭浸入。因此 $\M^{\otimes n}$ 很丰沛（II，4.4.6），$\M$ 丰沛（II，4.5.10）。\finishproof
\egastatement{5.4.6}{注记}
设 $A$ 为诺特进制环，$\mathfrak S=\Spf(A)$，$f:\mathfrak X\to\mathfrak S$ 为固有形式概形态射。设 $\N\subset\OO_{\mathfrak X}$ 为凝聚理想，使每个仿射形式开集 $U$ 上，$\Gamma(U,\N)$ 为 $\Gamma(U,\OO_{\mathfrak X})$ 的幂零根。其存在性容易由（I，10.10.2）及环同态保持幂零元素得到。令 $\mathfrak X'$ 为 $\N$ 定义的闭形式子概形。若 $\mathfrak X'$ 可代数化，$\mathfrak X$ 是否也可代数化，是一个值得研究的问题。若能对结构层 $\OO_X$（通常预概形或形式预概形）被平方为零理想扩张的情形作分类，例如用上同调性质的不变量分类，或许就能解决此问题。换言之，原文在此考虑带平方为零理想 $\J$ 的代数层 $\A$，使 $\OO_X\simeq\A/\J$。\transnote{本注记中的研究问题及“以后”等措辞均保留1961年原著的时间语境，不表示对现今研究状况的判断。}

\subsection{某些概形的一个分解}
\src{159}
\egastatement{5.5.1}{命题}
设 $A$ 为诺特进制环，$\frj$ 为定义理想，$Y=\Spec(A)$，$f:X\to Y$ 分离且有限型。令 $Y_0=\Spec(A/\frj)$、$X_0=X\times_YY_0$。若 $Z_0$ 为 $X_0$ 中在 $Y_0$ 上固有的开子集，则存在 $X$ 的开闭子集 $Z$，使 $Z$ 在 $Y$ 上固有，且 $Z\cap X_0=Z_0$。
\eproof
取 $X$ 的开集 $T$，使 $T\cap X_0=Z_0$，以 $\mathfrak T$ 表示 $T$ 沿 $Z_0$ 的完备化。其结构层的支集为 $Z_0$，在 $Y_0$ 上固有，故 $\mathfrak T$ 在 $\mathfrak Y=\Spf(A)$ 上固有（3.4.1）。由（5.1.8），存在 $T$ 中在 $Y$ 上固有的闭子概形 $Z$，其嵌入 $i:Z\to T$ 的完备化是同构 $\widehat Z\to\mathfrak T$。由（4.6.8），存在 $Z_0$ 在 $T$ 中的开邻域 $V$，使 $i^{-1}(V)\to V$ 为同构。而 $i^{-1}(V)$ 是 $Z$ 中包含 $Z_0$ 的邻域，由（5.1.3.1）等于整个 $Z$。故 $Z$ 在 $T$ 中开，从而在 $X$ 中开，证明了结论。\transnote{这里的闭性还可由 $Z$ 在 $Y$ 上固有及 $X$ 在 $Y$ 上分离得出；此句为译者补充说明。}\finishproof
\egastatement{5.5.2}{推论}
若 $X_0$ 在 $Y_0$ 上固有，则 $X$ 为两个不交开集 $Z,Z'$ 的并，其中 $Z$ 在 $Y$ 上固有且包含 $X_0$；此外，$X$ 中每个在 $Y$ 上固有的闭子集 $P$ 都包含在 $Z$ 中。
\eproof
前一断言取 $Z_0=X_0$，应用（5.5.1）。对于后一断言，$P\cap Z'$ 在 $P$ 中闭，故在 $Y$ 上固有；若它非空，则 $f(P\cap Z')$ 是 $Y$ 中的非空闭集，必与 $Y_0$ 相交（5.1.3.1），这与 $Z'$ 的定义矛盾。\finishproof
\begin{flushright}（待续。）\end{flushright}\par

% ---- END c3_s05b.tex ----

% END text/ch3_01_05.tex

\renewcommand{\newscope}{III-2}

% BEGIN text/ch3_06a.tex
\NewNote{以下为第三卷第二部分；其§6直接接续本章§5。第一部分的书末资料已移至全书附录，原稿“待续”仍保留。}
\source{\texttt{EGAIII.2.pdf}，PDF第4--13页，原书第5--14页的相应部分}
\setcounter{section}{5}\section{局部与整体 Tor 函子；Künneth 公式}
\subsection{引言}
\egap{6.1.1}
设 $f:X\to Y$ 为预概形态射，$\F$ 为拟凝聚 $\OO_X$-模。在研究“高阶直接像”$R^nf_*(\F)$ 时，会遇到下列一般问题：给定一个“基变换”态射 $g:Y'\to Y$，令
\[
 X'=X_{(Y')}=X\times_Y Y',\qquad
 \F'=\F\otimes_{\OO_Y}\OO_{Y'},\qquad
 f'=f_{(Y')}:X'\to Y',
\]
并在假定已经知道各 $R^nf_*(\F)$ 的情况下，求得关于高阶直接像 $R^nf'_*(\F')$ 的信息。容易看出（例如参见（7.7.2）），这归结为研究：当拟凝聚 $\OO_Y$-模 $\G$ 变化时，$R^nf_*(\F\otimes_{\OO_Y}\G)$ 如何变化；换言之，研究函子
\[
 \G\longmapsto R^nf_*(\F\otimes_{\OO_Y}\G).
\]
若 $\F$ 在 $Y$ 上平坦，则 $\G\mapsto\F\otimes_{\OO_Y}\G$ 是正合函子（$0_{\mathrm I}$，6.7.4），因而复合函子 $\G\mapsto R^nf_*(\F\otimes_{\OO_Y}\G)$ 仍构成上同调函子。但一般情形不再如此；为了应用上同调方法，便需要用另一些总是构成上同调函子的函子，替代上述函子。这些函子推广模论中的“Tor”函子，将在6.3至6.7中定义。实际上有“局部”和“整体”两种推广，它们由6.7讨论的谱序列相联系。作为这些谱序列的应用，在一定条件下可得到“Künneth公式”，用 $R^pf_{1*}(\F_1)$ 与 $R^qf_{2*}(\F_2)$ 表示
\[
 R^n(f_1\times f_2)_*(\F_1\otimes_Y\F_2).
\]
另外一些谱序列（6.8）推广模的 Tor 函子的结合性谱序列；基变换问题本身也引出谱序列（6.9）。

\egap{6.1.2}
此外，将会看到（6.10），这样定义的上同调函子 $\G\mapsto\mathcal T_*(\G)$，局部地（在 $Y$ 上）具有
\[
 \G\longmapsto H_*(\LL_\bullet\otimes_Y\G)
\]
的形式，其中 $\LL_\bullet$ 是局部自由 $\OO_Y$-模组成的复形（确定到同伦），$H_*$ 表示同调。因而，值得暂且忘掉产生 $\LL_\bullet$ 的特殊情形，而一般地研究上述形式的函子；必要时，再对 $\LL_i$ 或 $H_i(\LL_\bullet)$ 加上适当的有限性假设。这正是§7的对象，该节的主要部分可以独立于§6阅读。

这些函子最重要的性质涉及 $\mathcal T_*$ 的某个分量 $\mathcal T_n$ 的正合性。我们将给出判定此类性质的若干准则。作为应用，将得到保证函子
\[
 \G\longmapsto R^nf_*(\F\otimes_{\OO_Y}\G)
\]
正合的条件（使用（6.1.1）的记号）；这时称 $\F$ 在维数 $n$ 上对 $Y$ 上同调平坦。另一重要性质是函数
\[
 y\longmapsto\dim_{k(y)}\bigl(\mathcal T_n(k(y))\bigr)
\]
的半连续性。若 $\mathcal T_n$ 正合，这一性质成为连续性；根据 Grauert 的结果，当 $Y$ 约化时，逆命题也成立（7.8.4）。

\egap{6.1.3}
在§§6、7中，我们系统地使用超上同调，处处以层的复形而非单个层作为变量，尽管这一观点的必要性要到以后各章才会显现。这里发展的上同调形式体系，在本论著专门建立凝聚层上同调函子演算、并纳入对偶形式体系的章节中，将变得更为清晰。但这需要超出本章范围的展开。

\egap{6.1.4}
为简便起见，设 $K'^\bullet,K''^\bullet$ 是阿贝尔范畴 $\mathbf C$ 中的两个复形。如果复形态射 $f:K'^\bullet\to K''^\bullet$ 存在一个复形态射 $g:K''^\bullet\to K'^\bullet$，使 $f\circ g$、$g\circ f$ 都与恒等态射同伦，则称 $f$ 为\textbf{同伦等价}（原文 homotopisme）。存在这样的态射时，也滥用语言说两个复形“同伦”。当可以定义从 $\mathbf C$ 到阿贝尔范畴 $\mathbf C'$ 的协变加性函子 $T$ 关于 $\mathbf C$ 中复形的超上同调时（0，11.4.3），复形的同伦等价 $K'^\bullet\to K''^\bullet$ 显然典范地诱导超上同调同构
\[
 R^*T(K'^\bullet)\longrightarrow R^*T(K''^\bullet)
\]
（同上）。

\subsection{预概形上模复形的超上同调}
\egap{6.2.1}
设 $X$ 为预概形，$\K^\bullet=(\K^i)_{i\in\mathbf Z}$ 是微分算子次数为 $+1$ 的 $\OO_X$-模复形。回顾：对于每个预概形态射 $f:X\to Y$，已经定义了超上同调 $\OO_Y$-模 $\mathbb H^n(f,\K^\bullet)$（也记为 $\mathbb H^n(\K^\bullet)$ 或 $\mathbf R^nf_*(\K^\bullet)$），其中 $n\in\mathbf Z$（0，12.4.1）。\cmt{记号}{这里用 $\mathbb H$ 排印原文的超上同调记号，以区别于复形的上同调层 $\HH$；数学对象、次数与公式编号均不改变。}
超上同调 $\mathbb H^*(f,\K^\bullet)$ 是两个谱函子 ${}'\mathscr E(f,\K^\bullet)$、${}''\mathscr E(f,\K^\bullet)$ 的终值，其 $E_2$ 项为
\begin{align}
 {}'\mathscr E_2^{pq}&=\HH^p\bigl(\HH^q(f,\K^\bullet)\bigr),\tag{6.2.1.1}\\
 {}''\mathscr E_2^{pq}&=\HH^p\bigl(f,\HH^q(\K^\bullet)\bigr)
 =R^pf_*\bigl(\HH^q(\K^\bullet)\bigr).\tag{6.2.1.2}
\end{align}
在（6.2.1.1）中，$\HH^q(f,\K^\bullet)$ 表示其第 $i$ 次分量为
$\HH^q(f,\K^i)=R^qf_*(\K^i)$ 的复形（同上）。还要回顾：若 $Y$ 只有一个点，相应超上同调记为 $\mathbb H^*(X,\K^\bullet)$，它由 $\Gamma(X,\OO_X)$-模组成，且与所选单点预概形 $Y$ 无关；若 $Y=X$ 且 $f=1_X$，则
\[
 \mathbb H^n(f,\K^\bullet)=\HH^n(\K^\bullet).
\]
若除 $i=i_0$ 外，均有 $\K^i=0$，则
\[
 \mathbb H^n(f,\K^\bullet)\simeq R^{n-i_0}f_*(\K^{i_0}).
\]
谱序列 ${}'\mathscr E(f,\K^\bullet)$ 总是正则的；若 $\K^\bullet$ 下有界，两谱序列均为双正则（0，12.4.1）。

每个同伦等价 $h:\K^\bullet\to\K'^\bullet$（6.1.4）诱导超上同调的同构。若仅假设 $\HH^*(h):\HH^*(\K^\bullet)\to\HH^*(\K'^\bullet)$ 是同构，而 $\K^\bullet,\K'^\bullet$ 都下有界，结论仍成立；这是把（0，11.1.5）应用于（6.2.1.2）及 $\K'^\bullet$ 的相应谱序列所得。最后，对于 $X$ 的每个开覆盖 $\mathfrak U=(U_\alpha)$，也已把 $\mathbb H^*(\mathfrak U,\K^\bullet)$ 定义为双复形 $C^\bullet(\mathfrak U,\K^\bullet)$ 的上同调（0，12.4.5）；其双次数 $(i,j)$ 的分量按定义为 $C^i(\mathfrak U,\K^j)$。各 $\mathbb H^n(\mathfrak U,\K^\bullet)$ 仍是 $\Gamma(X,\OO_X)$-模。

\egas{6.2.2}{命题}
设 $X$ 为概形，$\mathfrak U=(U_\alpha)$ 为 $X$ 的仿射开覆盖。对于每个拟凝聚 $\OO_X$-模复形 $\K^\bullet$，超上同调模 $\mathbb H^*(X,\K^\bullet)$ 与 $\mathbb H^*(\mathfrak U,\K^\bullet)$ 典范同构。
\begin{proof}
覆盖中开集的任意有限交 $V$ 都是仿射的（I，5.5.6），所以对每个 $i$ 与每个 $q>0$，有 $H^q(V,\K^i)=0$（1.3.1）。因此本命题是（0，12.4.7）的特殊情形。
\end{proof}

\egas{6.2.3}{命题}
设 $f:X\to Y$ 为预概形的拟紧分离态射。对于每个拟凝聚 $\OO_X$-模复形 $\K^\bullet$，各 $\OO_Y$-模 $\mathbb H^n(f,\K^\bullet)$ 都拟凝聚。
\begin{proof}
各 $\HH^q(f,\K^i)=R^qf_*(\K^i)$ 都是拟凝聚 $\OO_Y$-模（1.4.10），因而 ${}'\mathscr E_2^{pq}$ 也拟凝聚，因为按（6.2.1.1），它是拟凝聚模之间同态的核模去像（I，4.1.1）。同理，第一谱序列中所有的 ${}'\mathscr E_r^{pq}$、$B_k^{pq}({}'\mathscr E)$、$Z_k^{pq}({}'\mathscr E)$ 都拟凝聚。谱序列的正则性说明 $Z_\infty^{pq}({}'\mathscr E)$ 等于某个 $Z_k^{pq}({}'\mathscr E)$，因而拟凝聚；而
\[
 B_\infty^{pq}({}'\mathscr E)=\ilim_k B_k^{pq}({}'\mathscr E)
\]
也拟凝聚（0，11.2.4及I，4.1.1）。于是各 ${}'\mathscr E_\infty^{pq}$ 都拟凝聚。正则性还说明 $\mathbb H^n(f,\K^\bullet)$ 上的滤过是离散且穷尽的：它是一个递增序列 $(\G_k)_{k\geq0}$ 的并，$\G_0=0$，且每个商 $\G_k/\G_{k-1}$ 都是某个 ${}'\mathscr E_\infty^{pq}$，因而拟凝聚。由（1.4.17）对 $k$ 归纳，各 $\G_k$ 都拟凝聚；再由
$\mathbb H^n(f,\K^\bullet)=\ilim_k\G_k$ 和（I，4.1.1）即得结论。
\end{proof}

\egas{6.2.4}{推论}
在（6.2.3）的假设下，对于 $Y$ 的每个仿射开集 $V$，典范同态
\begin{equation}
 \mathbb H^n(f^{-1}(V),\K^\bullet)
 \longrightarrow\Gamma(V,\mathbb H^n(f,\K^\bullet))\tag{6.2.4.1}
\end{equation}
对每个 $n\in\mathbf Z$ 都双射。
\begin{proof}
证明与（1.4.11）相同：使用（6.2.2），依次将其中的 $\F$ 换成 $\K^\bullet$，$\K^\bullet$ 换成 $\bigl(C^\bullet(\mathfrak U,\K^\bullet)\bigr)^\sim$，$\HH^n(\K^\bullet)$ 换成 $\mathbb H^n(f,\K^\bullet)$；并注意到根据（6.2.3），最后这个模是拟凝聚的。
\end{proof}

\egas{6.2.5}{命题}
设 $Y$ 为局部诺特预概形，$f:X\to Y$ 为本征态射，$\K^\bullet$ 为 $\OO_X$-模复形，并且各 $\OO_X$-模 $\HH^n(\K^\bullet)$ 凝聚。则各 $\OO_Y$-模 $\mathbb H^n(f,\K^\bullet)$ 凝聚。\cmt{P02}{原陈述只假设上同调层凝聚，以下证明却直接使用（6.2.2）、（6.2.4）及各项拟凝聚时的 Čech 正合性。需要核对是否已建立适用的拟凝聚复形替代或更广版本；本批不把此处直接宣判为命题错误。详见核查说明P02。}
\begin{proof}
问题在 $Y$ 上是局部的，可设 $Y$ 为诺特仿射概形。由（6.2.4），只需证明各 $\mathbb H^n(X,\K^\bullet)$ 是有限型 $\Gamma(Y,\OO_Y)$-模。根据（6.2.2），
\[
 \mathbb H^n(X,\K^\bullet)=\mathbb H^n(\mathfrak U,\K^\bullet),
\]
其中，由于 $X$ 拟紧，可取有限仿射开覆盖 $\mathfrak U$。各复形 $C^\bullet(\mathfrak U,\K^j)$ 的余链按定义是交替的，因此存在整数 $r>0$，使其在 $i<0$ 及 $i>r$ 时第 $i$ 项为零。由（0，11.3.3），双复形 $C^\bullet(\mathfrak U,\K^\bullet)$ 的两个谱序列都双正则。

各覆盖开集的交都是仿射开集（I，5.5.6），所以每个函子 $\F\mapsto C^i(\mathfrak U,\F)$ 在拟凝聚 $\OO_X$-模范畴中正合。因此
\[
 H^q(C^i(\mathfrak U,\K^\bullet))=C^i(\mathfrak U,\HH^q(\K^\bullet)),
\]
第二谱序列的 $E_2$ 项（0，11.3.2）为
\begin{align*}
 {}''E_2^{pq}
 &=H^p(C^\bullet(\mathfrak U,\HH^q(\K^\bullet)))\\
 &=H^p(\mathfrak U,\HH^q(\K^\bullet))
 =H^p(X,\HH^q(\K^\bullet)),
\end{align*}
最后一步由（1.4.1）。由于 $f$ 本征，这些都是有限型 $\Gamma(Y,\OO_Y)$-模（3.2.1）。第二谱序列双正则，于是各 $\mathbb H^n(X,\K^\bullet)$ 都有限型（0，11.1.8）。
\end{proof}
特别地，如果 $\K^\bullet$ 是凝聚 $\OO_X$-模复形，则在（6.2.5）关于 $Y,f$ 的假设下，各 $\mathbb H^n(f,\K^\bullet)$ 凝聚（$0_{\mathrm I}$，5.3.4）。

\egap{6.2.6}
超上同调 $\mathbb H^*(f,\K^\bullet)$ 在下有界 $\OO_X$-模复形范畴中构成上同调函子（0，12.4.4）。当 $f$ 拟紧、$X$ 的底空间局部诺特时，它在所有 $\OO_X$-模复形的范畴中也构成上同调函子：由（G，II，3.10.1），$f_*$ 与归纳极限交换（问题在 $Y$ 上局部），故可应用（0，11.5.2）。

最后，如果 $f$ 分离，则 $\mathbb H^*(f,\K^\bullet)$ 在拟凝聚 $\OO_X$-模复形范畴中构成上同调函子。当 $Y$ 仿射时，这是直接的，因为此时 $X$ 是概形，由（6.2.2）的典范同构，归结为验证 $\K^\bullet\mapsto\mathbb H^n(\mathfrak U,\K^\bullet)$ 是上同调函子；而 $\K^\bullet\mapsto C^\bullet(\mathfrak U,\K^\bullet)$ 在该范畴中正合（I，1.3.7），故结论立即成立。

一般情形下，对于 $Y$ 的每个仿射开集 $V$，$f^{-1}(V)$ 是概形。为了应用以上论证，只需验证：对于拟凝聚模复形的短正合列
\[
 0\longrightarrow\K'^\bullet\longrightarrow\K^\bullet
 \longrightarrow\K''^\bullet\longrightarrow0,
\]
连接同态
\[
 \delta:\mathbb H^n(f,\K''^\bullet)|_V
 \longrightarrow\mathbb H^{n+1}(f,\K'^\bullet)|_V
\]
不依赖于定义它时所用的 $f^{-1}(V)$ 的仿射开覆盖 $\mathfrak U$。这是因为，如果 $\mathfrak U'$ 是一个更细的仿射开覆盖，则下列由典范同构组成的图交换：
\[
\begin{tikzcd}[column sep=large]
\mathbb H^n(\mathfrak U,\K^\bullet) \arrow[r,"\sim"] \arrow[d,"\sim"'] & \mathbb H^n(f^{-1}(V),\K^\bullet)\\
\mathbb H^n(\mathfrak U',\K^\bullet) \arrow[ur,"\sim"'] &
\end{tikzcd}
\]
而且下面的图也交换：
\[
\begin{tikzcd}[column sep=large]
\mathbb H^n(\mathfrak U,\K''^\bullet) \arrow[r,"\delta"] \arrow[d] & \mathbb H^{n+1}(\mathfrak U,\K'^\bullet) \arrow[d]\\
\mathbb H^n(\mathfrak U',\K''^\bullet) \arrow[r,"\delta"] & \mathbb H^{n+1}(\mathfrak U',\K'^\bullet).
\end{tikzcd}
\]
当上述条件之一成立且 $\K^\bullet\to\K'^\bullet$ 是同伦等价（6.1.4）时，相应同构是一个 $\delta$-函子同构（0，11.4.4）。

\egap{6.2.7}
以上全部结果，自然也适用于微分次数为 $-1$ 的拟凝聚 $\OO_X$-模复形 $\K_\bullet$，除记号外无需改变；只须考虑 $\K^i=\K_{-i}$ 的上标复形即可。

\subsection{两个模复形的超 Tor}
\egap{6.3.1}
设 $A$ 为交换环，$P_\bullet,Q_\bullet$ 是微分次数为 $-1$ 的两个 $A$-模复形。取 $P_\bullet$、$Q_\bullet$ 的 Cartan--Eilenberg 投射解消 $L_{\bullet\bullet}$、$M_{\bullet\bullet}$（0，11.6.1）。按第一次数相加、第二次数相加，$L_{\bullet\bullet}\otimes_A M_{\bullet\bullet}$ 是一个微分次数为 $-1$ 的双复形。它的同调不依赖所选解消，按定义就是双函子 $P_\bullet\otimes_AQ_\bullet$ 关于两个变量的超同调（0，11.6.5）。定义
\begin{equation}
 \Tor_n^A(P_\bullet,Q_\bullet)
 =H_n(L_{\bullet\bullet}\otimes_A M_{\bullet\bullet}).\tag{6.3.1.1}
\end{equation}
称此 $A$-模为两个复形的第 $n$ 个\textbf{超 Tor}。在下有界 $A$-模复形范畴中，各 $\Tor_n^A$ 在两个变量中构成同调双函子（0，11.6.5）。此外：

\egas{6.3.2}{命题}
双函子 $\Tor_*^A(P_\bullet,Q_\bullet)$ 是两个谱双函子 ${}'E(P_\bullet,Q_\bullet)$、${}''E(P_\bullet,Q_\bullet)$ 的公共终值，其 $E_2$ 项为
\begin{align}
 {}'E^2_{pq}&=H_p\bigl(\Tor_q^A(P_\bullet,Q_\bullet)\bigr),\tag{6.3.2.1}\\
 {}''E^2_{pq}&=\bigoplus_{i+j=q}\Tor_p^A\bigl(H_i(P_\bullet),H_j(Q_\bullet)\bigr).\tag{6.3.2.2}
\end{align}
在（6.3.2.1）中，$\Tor_q^A(P_\bullet,Q_\bullet)$ 指由各普通 Tor 模 $\Tor_q^A(P_i,Q_j)$ 构成的双复形。谱序列（6.3.2.2）总是正则的；如果两个复形都下有界，或者 $A$ 的上同调维数有限，则两个谱序列都双正则。
\begin{proof}
这是（0，11.6.5）的结果；当 $A$ 的上同调维数为有限整数 $n$ 时，每个 $A$-模都存在长度为 $n$ 的投射解消（M，VI，2.1）。
\end{proof}

\egas{6.3.3}{推论}
设 $P'_\bullet,Q'_\bullet$ 为两个 $A$-模复形，$u:P_\bullet\to P'_\bullet$、$v:Q_\bullet\to Q'_\bullet$ 为复形同态。如果诱导的同调同态 $H_*(u)$ 与 $H_*(v)$ 都双射，则所诱导的同态
\[
 \Tor_*^A(P_\bullet,Q_\bullet)\longrightarrow
 \Tor_*^A(P'_\bullet,Q'_\bullet)
\]
也双射。
\begin{proof}
由 $u,v$ 诱导的第二谱序列同态在 $E_2$ 项上是同构；由（6.3.2），这些谱序列正则，结论遂由（0，11.1.5）得出。
\end{proof}

\egas{6.3.4}{命题}
设 $P_\bullet,Q_\bullet$ 为下有界 $A$-模复形。设 $L_{\bullet\bullet}$、$M_{\bullet\bullet}$ 是由平坦 $A$-模构成的双复形，且每个 $L_{i,\bullet}$、$M_{i,\bullet}$ 分别是 $P_i$、$Q_i$ 的解消。则有典范同构
\begin{equation}
\begin{split}
 \Tor_*^A(P_\bullet,Q_\bullet)
 &\simeq H_*(L_{\bullet\bullet}\otimes_AQ_\bullet)\\
 &\simeq H_*(P_\bullet\otimes_AM_{\bullet\bullet})\\
 &\simeq H_*(L_{\bullet\bullet}\otimes_AM_{\bullet\bullet}).
\end{split}\tag{6.3.4.1}
\end{equation}
\begin{proof}
由（0，11.6.5，（ii）、（iii））及平坦模的定义得出。
\end{proof}

\egas{6.3.5}{备注}
（i）使用（6.3.1）的记号，双复形 $L_{\bullet\bullet}\otimes_AM_{\bullet\bullet}$ 与 $M_{\bullet\bullet}\otimes_AL_{\bullet\bullet}$ 典范同构，从而
\[
 \Tor_*^A(P_\bullet,Q_\bullet)\simeq\Tor_*^A(Q_\bullet,P_\bullet).
\]
（ii）若 $F,G$ 为两个 $A$-模，$P_\bullet,Q_\bullet$ 分别只在次数0处为 $F,G$ 而其他次数为零，则它们的普通投射解消 $L_\bullet,M_\bullet$ 通过补零可视为 Cartan--Eilenberg 解消。因此，此时 $\Tor_*^A(P_\bullet,Q_\bullet)=\Tor_*^A(F,G)$。

\egas{6.3.6}{命题}
设 $(P_\bullet^\lambda)$、$(Q_\bullet^\mu)$ 为两个 $A$-模复形的滤过归纳系统。则有典范同构
\begin{equation}
 \ilim_{\lambda,\mu}\Tor_*^A(P_\bullet^\lambda,Q_\bullet^\mu)
 \simeq\Tor_*^A\left(\ilim_\lambda P_\bullet^\lambda,
                         \ilim_\mu Q_\bullet^\mu\right).\tag{6.3.6.1}
\end{equation}
\begin{proof}
令 $P_\bullet=\ilim_\lambda P_\bullet^\lambda$、$Q_\bullet=\ilim_\mu Q_\bullet^\mu$。由函子性，各 $\Tor_*^A(P_\bullet^\lambda,Q_\bullet^\mu)$ 构成归纳系统，典范同态到 $\Tor_*^A(P_\bullet,Q_\bullet)$ 也构成归纳系统。由此得到（6.3.6.1）的同态；更一般地，还得到
\[
 \ilim_{\lambda,\mu}{}''E(P_\bullet^\lambda,Q_\bullet^\mu)
 \longrightarrow{}''E(P_\bullet,Q_\bullet),
\]
（6.3.6.1）就是其终值的同态。第二谱序列正则（6.3.2），其归纳极限也正则，这是定义（0，11.1.7）和（0，11.3.3）的证明的结果。所以由（0，11.1.5），只需证明
\begin{equation}
 \ilim_{\lambda,\mu}{}''E(P_\bullet^\lambda,Q_\bullet^\mu)
 \longrightarrow{}''E(P_\bullet,Q_\bullet)\tag{6.3.6.2}
\end{equation}
在 $E_2$ 项上双射。

由于同调函子 $H_*$ 与模复形的归纳极限交换，最后归结为证明：对于模的两个滤过归纳系统 $(F^\lambda)$、$(G^\mu)$，典范同态
\[
 \ilim_{\lambda,\mu}\Tor_*^A(F^\lambda,G^\mu)
 \longrightarrow\Tor_*^A(\ilim_\lambda F^\lambda,\ilim_\mu G^\mu)
\]
双射。对每个 $F^\lambda$，取其典范自由解消
\[
 \cdots\longrightarrow L_{i+1}^\lambda\longrightarrow L_i^\lambda
 \longrightarrow\cdots\longrightarrow L_0^\lambda\longrightarrow F^\lambda
 \longrightarrow0,
\]
其中 $L_0^\lambda$ 是以 $F^\lambda$ 的元素为基的形式线性组合模，而 $L_{i+1}^\lambda$ 是以相应核 $\Ker(L_i^\lambda\to L_{i-1}^\lambda)$ 的元素为基的形式线性组合模。\cmt{记号}{在 $i=0$ 处，核取增广映射 $L_0^\lambda\to F^\lambda$ 的核；这是原文“典范自由解消”的约定。}
容易验证这些复形组成归纳系统。置 $F=\ilim F^\lambda$、$L_\bullet=\ilim L_\bullet^\lambda$；滤过归纳极限正合，所以 $L_\bullet$ 是 $F$ 的解消。各 $L_i$ 是自由模的滤过归纳极限，因而平坦（$0_{\mathrm I}$，6.1.2）。同样构造 $G^\mu$ 的典范自由解消 $M_\bullet^\mu$，其极限 $M_\bullet$ 是 $G=\ilim G^\mu$ 的平坦解消。由（6.3.5）和（6.3.4），
\[
 \Tor_*^A(F,G)=H_*(L_\bullet\otimes_AM_\bullet).
\]
又因为 $H_*$ 与模复形的归纳极限交换，
\[
 H_*(L_\bullet\otimes_AM_\bullet)
 =\ilim_{\lambda,\mu}H_*(L_\bullet^\lambda\otimes_AM_\bullet^\mu).
\]
右边各项就是 $\Tor_*^A(F^\lambda,G^\mu)$，证明完成。
\end{proof}
若还存在统一整数 $i_0$，使对所有 $\lambda,\mu$ 及 $i<i_0$ 都有 $P_i^\lambda=Q_i^\mu=0$，则同样证明典范同态
\begin{equation}
 \ilim_{\lambda,\mu}{}'E(P_\bullet^\lambda,Q_\bullet^\mu)
 \longrightarrow{}'E(P_\bullet,Q_\bullet)\tag{6.3.6.3}
\end{equation}
双射。

\egas{6.3.7}{命题}
假设 $P_\bullet,Q_\bullet$ 下有界。若 $P_\bullet$ 由平坦 $A$-模组成，则关于变量 $Q_\bullet$ 有典范的 $\delta$-函子 $A$-同构
\begin{equation}
 \Tor_*^A(P_\bullet,Q_\bullet)\simeq H_*(P_\bullet\otimes_AQ_\bullet).\tag{6.3.7.1}
\end{equation}
\begin{proof}
谱序列（6.3.2.1）双正则且退化，故同构存在（0，11.1.6）。从 $P_\bullet$ 的 Cartan--Eilenberg 投射解消计算超 Tor（6.3.4），也立即看出该同构与连接同态相容，因而是 $\delta$-函子同构。
\end{proof}

\egap{6.3.8}
设 $\rho:A\to A'$ 为环同态。我们要定义一个与 $\rho$ 典范相伴、具有函子性的零次 $A$-同态
\begin{equation}
 \rho_{P_\bullet,Q_\bullet}:\Tor_*^A(P_\bullet,Q_\bullet)
 \longrightarrow\Tor_*^{A'}(P_\bullet\otimes_AA',Q_\bullet\otimes_AA').\tag{6.3.8.1}
\end{equation}
为此，取 $P_\bullet$ 的 Cartan--Eilenberg 投射解消 $L_{\bullet\bullet}$，以及 $P_\bullet\otimes_AA'$ 的 Cartan--Eilenberg 投射解消 $L'_{\bullet\bullet}$。我们将说明，可以定义一个确定到同伦的 $A'$-复形同态
\[
 L_{\bullet\bullet}\otimes_AA'\longrightarrow L'_{\bullet\bullet}.
\]
事实上，对于每个 $i$，任意给定 $B_i(P_\bullet)$ 的投射解消 $(X^B_{ij})_{j\geq0}$ 与 $H_i(P_\bullet)$ 的投射解消 $(X^H_{ij})_{j\geq0}$ 后，$L_{\bullet\bullet}$ 的构造便完全确定；它们分别等于 $B_i^{\mathrm I}(L_{\bullet\bullet})$ 与 $H_i^{\mathrm I}(L_{\bullet\bullet})$。依次得到
\[
 Z_i^{\mathrm I}(L_{\bullet\bullet})
 =H_i^{\mathrm I}(L_{\bullet\bullet})\oplus B_i^{\mathrm I}(L_{\bullet\bullet}),
\quad
 L_{i,\bullet}=Z_i^{\mathrm I}(L_{\bullet\bullet})\oplus B_{i-1}^{\mathrm I}(L_{\bullet\bullet}).
\]
此时 $X^B_{i,\bullet}\otimes_AA'$ 一般不再是 $P_i\otimes_AA'$ 的解消，但仍是投射 $A'$-模复形。\cmt{T03}{原文此处确实写 $P_i\otimes_AA'$，而 $X^B_{i,\bullet}$ 原本是 $B_i(P_\bullet)$ 的解消。该说明的对象下标混淆；应区分它增广到 $B_i(P_\bullet)\otimes_AA'$，以及后续到新复形边界的典范映射。这里保留原句，不用其错误对象作为推论依据。}
故存在一个与增广相容且确定到同伦的 $A'$-同态
\[
 X^B_{i,\bullet}\otimes_AA'\longrightarrow B_i^{\mathrm I}(L'_{\bullet\bullet})
\]
（M，V，1.1）。同样存在确定到同伦的同态
$X^H_{i,\bullet}\otimes_AA'\to H_i^{\mathrm I}(L'_{\bullet\bullet})$。由上述构造，对每个 $i$ 得到 $A'$-同态 $L_{i,\bullet}\otimes_AA'\to L'_{i,\bullet}$；它们与微分 $L_{i,\bullet}\to L_{i-1,\bullet}$ 及 $L'$ 的相应微分相容，遂构成所求同态。

为了定义（6.3.8.1），同样取 $Q_\bullet$ 与 $Q_\bullet\otimes_AA'$ 的 Cartan--Eilenberg 投射解消 $M_{\bullet\bullet}$、$M'_{\bullet\bullet}$ 及同态 $M_{\bullet\bullet}\otimes_AA'\to M'_{\bullet\bullet}$。由这些同态得到
\[
 (L_{\bullet\bullet}\otimes_AA')\otimes_{A'}
 (M_{\bullet\bullet}\otimes_AA')
 \longrightarrow L'_{\bullet\bullet}\otimes_{A'}M'_{\bullet\bullet},
\]
再经复合，得到 $A$-双复形同态
\[
 L_{\bullet\bullet}\otimes_AM_{\bullet\bullet}
 \longrightarrow L'_{\bullet\bullet}\otimes_{A'}M'_{\bullet\bullet}.
\]
取同调便得到（6.3.8.1）；由于它来自确定到同伦的复形态射，所以定义良好。

若 $\rho':A'\to A''$ 为另一个环同态，$\rho''=\rho'\circ\rho$，则显然
\[
 \rho''_{P_\bullet,Q_\bullet}
 =\rho'_{P'_\bullet,Q'_\bullet}\circ\rho_{P_\bullet,Q_\bullet},
 \qquad P'_\bullet=P_\bullet\otimes_AA',\quad Q'_\bullet=Q_\bullet\otimes_AA'.
\]
还须注意，上述双复形态射在两个变量中函子性地定义谱序列同态
\begin{align*}
 {}'E^r_{pq}(P_\bullet,Q_\bullet)&\longrightarrow{}'E^r_{pq}(P_\bullet\otimes_AA',Q_\bullet\otimes_AA'),\\
 {}''E^r_{pq}(P_\bullet,Q_\bullet)&\longrightarrow{}''E^r_{pq}(P_\bullet\otimes_AA',Q_\bullet\otimes_AA'),
\end{align*}
它们不依赖所取的 Cartan--Eilenberg 解消，也满足上述传递性。

\egas{6.3.9}{命题}
设 $\rho:A\to A'$ 为环同态，并且 $A'$ 是平坦 $A$-模。则有函子性的典范同构
\begin{equation}
 \Tor_*^{A'}(P_\bullet\otimes_AA',Q_\bullet\otimes_AA')
 \simeq\Tor_*^A(P_\bullet,Q_\bullet)\otimes_AA',\tag{6.3.9.1}
\end{equation}
以及
\begin{equation}
\begin{aligned}
 {}'E(P_\bullet\otimes_AA',Q_\bullet\otimes_AA')&\simeq{}'E(P_\bullet,Q_\bullet)\otimes_AA',\\
 {}''E(P_\bullet\otimes_AA',Q_\bullet\otimes_AA')&\simeq{}''E(P_\bullet,Q_\bullet)\otimes_AA'.
\end{aligned}\tag{6.3.9.2}
\end{equation}
\begin{proof}
由函子 $M\mapsto M\otimes_AA'$ 的正合性，$L_{\bullet\bullet}\otimes_AA'$ 与 $M_{\bullet\bullet}\otimes_AA'$ 分别仍是 $P_\bullet\otimes_AA'$ 与 $Q_\bullet\otimes_AA'$ 的 Cartan--Eilenberg 投射解消，结论即得。
\end{proof}

\egap{6.3.10}
设 $\rho:A\to A'$ 为环同态。对于每个 $A'$-模复形 $P'_\bullet$，$P'_{\bullet[\rho]}$ 是限制标量所得的 $A$-模复形；恒等映射 $P'_{\bullet[\rho]}\to P'_\bullet$ 可看作典范映射的复合
\[
 P'_{\bullet[\rho]}\longrightarrow P'_\bullet\otimes_AA'
 \xrightarrow{\mu}P'_\bullet,
 \qquad\mu(x\otimes a')=a'x.
\]
若 $Q'_\bullet$ 为另一个 $A'$-模复形，则有函子性的零次典范同态
\begin{equation}
\begin{split}
 \Tor_*^A(P'_{\bullet[\rho]},Q'_{\bullet[\rho]})
 &\longrightarrow\Tor_*^{A'}(P'_\bullet\otimes_AA',Q'_\bullet\otimes_AA')\\
 &\longrightarrow\Tor_*^{A'}(P'_\bullet,Q'_\bullet).
\end{split}\tag{6.3.10.1}
\end{equation}
第一箭头是（6.3.8）定义的 $A$-同态；第二箭头由两个乘法同态及函子性得到。对于（6.3.2）的谱序列也有类似同态，并有显然的传递性性质，其表述留给读者。

\egas{6.3.11}{命题}
设 $\rho:A\to A'$ 使 $A'$ 成为平坦 $A$-模。对于下有界的 $A'$-模复形 $P'_\bullet$ 与 $A$-模复形 $Q_\bullet$，有函子性的典范同构
\begin{equation}
 \Tor_*^A(P'_{\bullet[\rho]},Q_\bullet)
 \simeq\Tor_*^{A'}(P'_\bullet,Q_\bullet\otimes_AA').\tag{6.3.11.1}
\end{equation}
\begin{proof}
若 $M_{\bullet\bullet}$ 是 $Q_\bullet$ 的 Cartan--Eilenberg 投射解消，则 $M_{\bullet\bullet}\otimes_AA'$ 是 $Q_\bullet\otimes_AA'$ 的相应解消，且有典范同构
\[
 P'_{\bullet[\rho]}\otimes_AM_{\bullet\bullet}
 \simeq P'_\bullet\otimes_{A'}(M_{\bullet\bullet}\otimes_AA').
\]
由（6.3.4）得出结论。
\end{proof}

\egas{6.3.12}{备注}
设 $(A^\lambda)$ 为环的滤过归纳系统，$(P_\bullet^\lambda)$、$(Q_\bullet^\lambda)$ 是 $A^\lambda$-模复形的两个归纳系统。则有推广（6.3.6.1）的典范同构
\begin{equation}
 \ilim_\lambda\Tor_*^{A^\lambda}(P_\bullet^\lambda,Q_\bullet^\lambda)
 \simeq\Tor_*^A(P_\bullet,Q_\bullet),\tag{6.3.12.1}
\end{equation}
其中 $A=\ilim A^\lambda$、$P_\bullet=\ilim P_\bullet^\lambda$、$Q_\bullet=\ilim Q_\bullet^\lambda$。用（6.3.10）定义了 $\lambda\leq\mu$ 时的过渡同态后，证明就是（6.3.6）的证明。

\egas{6.3.13}{命题}
设 $S$ 为 $A$ 的乘性子集，$P_\bullet,Q_\bullet$ 是两个 $A$-模复形，其中由 $S$ 中每个元素定义的数乘都双射。因此，若 $A'=S^{-1}A$，两个复形都由 $A'$-模组成。则有典范同构
\[
 \Tor_*^A(P_\bullet,Q_\bullet)\simeq\Tor_*^{A'}(P_\bullet,Q_\bullet).
\]
\begin{proof}
假设说明典范同态 $P_\bullet\to P_\bullet\otimes_AA'$、$Q_\bullet\to Q_\bullet\otimes_AA'$ 都双射。另一方面，超 Tor 的函子性说明每个 $s\in S$ 在 $\Tor_*^A(P_\bullet,Q_\bullet)$ 上都定义双射的数乘，因而典范同态
\[
 \Tor_*^A(P_\bullet,Q_\bullet)
 \longrightarrow\Tor_*^A(P_\bullet,Q_\bullet)\otimes_AA'
\]
也双射。由于 $A'$ 是平坦 $A$-模，结论由（6.3.9）得到；相应谱序列也有典范同构。
\end{proof}

% END text/ch3_06a.tex


% BEGIN text/ch3_06b.tex
% 本次续译：EGA III 第二部分，6.4--6.5。严格依上传法文底本。
\subsection{拟凝聚模复形的局部超 Tor 函子：仿射概形情形}
\source{\texttt{EGAIII.2.pdf}，PDF第13--15页，原书第14--16页的相应部分}
\egap{6.4.1}
设 $S$ 是环为 $A$ 的仿射概形，$X,Y$ 是两个仿射 $S$-概形，其环分别为 $B,C$；于是 $B,C$ 都是 $A$-代数。每个拟凝聚 $\OO_X$-模复形 $\PP_\bullet$（分别地，拟凝聚 $\OO_Y$-模复形 $\mathscr Q_\bullet$）都具有 $\widetilde{P_\bullet}$（分别地，$\widetilde{Q_\bullet}$）的形式，其中 $P_\bullet$ 是 $B$-模复形，$Q_\bullet$ 是 $C$-模复形（I，1.3.7、1.3.8）。自然可以把两者看作 $A$-模复形，形成 $\Tor_n^A(P_\bullet,Q_\bullet)$。此外，由超 Tor 的双函子性，$A$-代数 $B$ 和 $C$ 在这个 $A$-模上作用；这些作用使它成为一个 $(B,C)$-双模，等价地，成为
\[
 B\otimes_A C=A(X\times_S Y)
\]
上的模。因此定义了拟凝聚 $\OO_{X\times_S Y}$-模
\begin{equation}
 \lTor_n^S(\PP_\bullet,\mathscr Q_\bullet)
 =\bigl(\Tor_n^A(P_\bullet,Q_\bullet)\bigr)^{\sim}.
 \tag{6.4.1.1}
\end{equation}
称它为复形 $\PP_\bullet,\mathscr Q_\bullet$ 的第 $n$ 个\textbf{局部超 Tor}；也记作 $\lTor_n(\PP_\bullet,\mathscr Q_\bullet)$。

\egas{6.4.2}{引理}
采用（6.4.1）的记号。设 $A=R^{-1}A'$，其中 $A'$ 是环，$R$ 是 $A'$ 的乘法子集。令 $S'=\Spec(A')$，于是 $X,Y$ 都可看作 $S'$-预概形，且
\[
 X\times_{S'}Y=X\times_S Y
\]
（I，1.6.2、3.2.4）。则
\[
 \lTor_n^{S'}(\PP_\bullet,\mathscr Q_\bullet)
 =\lTor_n^S(\PP_\bullet,\mathscr Q_\bullet).
\]
\begin{proof}
由公式（6.4.1.1）及（6.3.13）得出。
\end{proof}

\egap{6.4.3}
采用（6.4.1）的记号和假设，设 $\F=\widetilde F$、$\G=\widetilde G$ 分别是拟凝聚 $\OO_X$-模和 $\OO_Y$-模。把它们看作模复形后，将其第 $n$ 个超 Tor 记为 $\lTor_n^S(\F,\G)$ 或 $\lTor_n(\F,\G)$。由（6.3.5，（ii）），
\begin{equation}
 \lTor_n^S(\F,\G)=\bigl(\Tor_n^A(F,G)\bigr)^{\sim}.
 \tag{6.4.3.1}
\end{equation}
现在回到两个拟凝聚模复形 $\PP_\bullet,\mathscr Q_\bullet$ 的一般情形。由公式（6.4.1.1）、（6.4.3.1）及命题（6.3.2），$\lTor_*^S(\PP_\bullet,\mathscr Q_\bullet)$ 是两个谱序列 ${}'\spE(\PP_\bullet,\mathscr Q_\bullet)$、${}''\spE(\PP_\bullet,\mathscr Q_\bullet)$ 的终值，其 $E_2$ 项为
\begin{align}
 {}'\spE^2_{pq}
 &=\HH_p\bigl(\lTor_q^S(\PP_\bullet,\mathscr Q_\bullet)\bigr),
 \tag{6.4.3.2}\\
 {}''\spE^2_{pq}
 &=\bigoplus_{i+j=q}\lTor_p^S\bigl(\HH_i(\PP_\bullet),
                         \HH_j(\mathscr Q_\bullet)\bigr).
 \tag{6.4.3.3}
\end{align}
其中，（6.4.3.2）中的 $\lTor_q^S(\PP_\bullet,\mathscr Q_\bullet)$ 指由拟凝聚 $\OO_{X\times_S Y}$-模 $\lTor_q^S(\PP_i,\mathscr Q_j)$ 组成的双复形。

\egap{6.4.4}
再考虑两个仿射概形 $X^{(1)}=\Spec(B^{(1)})$、$Y^{(1)}=\Spec(C^{(1)})$，其中 $B^{(1)},C^{(1)}$ 是 $A$-代数。设给定两个 $S$-态射
\[
 u:X^{(1)}\longrightarrow X,\qquad v:Y^{(1)}\longrightarrow Y,
\]
对应于 $A$-同态 $\varphi:B\to B^{(1)}$、$\psi:C\to C^{(1)}$。考虑 $\OO_{X^{(1)}}$-模复形
$u^*(\PP_\bullet)=(P_\bullet\otimes_B B^{(1)})^{\sim}$ 和 $\OO_{Y^{(1)}}$-模复形
$v^*(\mathscr Q_\bullet)=(Q_\bullet\otimes_C C^{(1)})^{\sim}$。
典范 $A$-同态
\[
 P_\bullet\longrightarrow P_\bullet\otimes_B B^{(1)},\qquad
 Q_\bullet\longrightarrow Q_\bullet\otimes_C C^{(1)}
\]
由函子性给出 $A$-同态
\[
 \Tor_n^A(P_\bullet,Q_\bullet)\longrightarrow
 \Tor_n^A(P_\bullet\otimes_B B^{(1)},Q_\bullet\otimes_C C^{(1)}).
\]
仍由函子性，这实际上是 $(B\otimes_A C)$-模同态。因此得到一个 $(u\times_S v)$-态射
\begin{equation}
 \theta:\lTor_n^S(\PP_\bullet,\mathscr Q_\bullet)
 \longrightarrow\lTor_n^S(u^*\PP_\bullet,v^*\mathscr Q_\bullet),
 \tag{6.4.4.1}
\end{equation}
进而得到 $\OO_{X^{(1)}\times_S Y^{(1)}}$-模同态
\begin{equation}
 \theta^\sharp:(u\times_S v)^*\lTor_n^S(\PP_\bullet,\mathscr Q_\bullet)
 \longrightarrow\lTor_n^S(u^*\PP_\bullet,v^*\mathscr Q_\bullet).
 \tag{6.4.4.2}
\end{equation}
在下有界拟凝聚模复形的范畴中，这显然是双 $\delta$-函子的态射。同态（6.4.4.2）未必双射，不过有：

\egas{6.4.5}{引理}
采用（6.4.4）的记号，若 $u,v$ 都是开浸入，则（6.4.4.2）双射。
\begin{proof}
把 $X^{(1)}$、$Y^{(1)}$ 分别与 $X,Y$ 的开集等同。前者是形如 $D(f)$（$f\in B$）的开集的并，后者是形如 $D(g)$（$g\in C$）的开集的并；诱导预概形 $D(\varphi(f))$ 与 $D(f)$、$D(\psi(g))$ 与 $D(g)$ 分别同构。

只须在 $X^{(1)}=D(f)$、$Y^{(1)}=D(g)$ 时证明引理。确实，若这一点已经证明，则回到一般情形，只需证明 $\theta^\sharp$ 在每个 $D(f)\times_S D(g)$ 上的限制是同构。设 $u_1:D(f)\to X^{(1)}$、$v_1:D(g)\to Y^{(1)}$ 是典范嵌入，则这一限制就是 $(u_1\times_S v_1)^*(\theta^\sharp)$。令 $u'=u\circ u_1$、$v'=v\circ v_1$，由定义（6.4.4）及（$0_{\mathrm I}$，4.4.8），将它与典范同态
\begin{equation}
\begin{split}
 (u_1\times_S v_1)^*\lTor_n^S(u^*\PP_\bullet,v^*\mathscr Q_\bullet)
 \longrightarrow\lTor_n^S(u'^*\PP_\bullet,v'^*\mathscr Q_\bullet)
\end{split}\tag{6.4.5.1}
\end{equation}
复合，就得到典范同态
\begin{equation}
 (u'\times_S v')^*\lTor_n^S(\PP_\bullet,\mathscr Q_\bullet)
 \longrightarrow\lTor_n^S(u'^*\PP_\bullet,v'^*\mathscr Q_\bullet).
 \tag{6.4.5.2}
\end{equation}
所以若（6.4.5.1）、（6.4.5.2）都是同构，则 $(u_1\times_S v_1)^*(\theta^\sharp)$ 也是同构。

现在设 $X^{(1)}=D(f)$、$Y^{(1)}=D(g)$，从而 $B^{(1)}=B_f$、$C^{(1)}=C_g$。于是 $u^*\PP_\bullet$、$v^*\mathscr Q_\bullet$ 分别与 $\widetilde{(P_\bullet)_f}$、$\widetilde{(Q_\bullet)_g}$ 等同；而 $X^{(1)}\times_S Y^{(1)}$ 与 $X\times_S Y=\Spec(B\otimes_A C)$ 的开子概形 $D(f\otimes g)$ 等同（II，4.3.2.4）。所要证明的是由典范映射按函子性得到的同态
\begin{equation}
 \bigl(\Tor_n^A(P_\bullet,Q_\bullet)\bigr)_{f\otimes g}
 \longrightarrow\Tor_n^A((P_\bullet)_f,(Q_\bullet)_g)
 \tag{6.4.5.3}
\end{equation}
双射。由（$0_{\mathrm I}$，1.6.1），可以写
$(P_\bullet)_f=\ilim_m P_\bullet^{(m)}$，其中每个 $P_\bullet^{(m)}$ 都是与 $P_\bullet$ 相同的 $B$-模复形，而 $m\leq n$ 时的过渡映射是乘以 $f^{n-m}$。对 $Q_\bullet$ 用 $g$ 代替 $f$，也有同样描述。相应同态
\[
 \Tor_n^A(P_\bullet^{(m)},Q_\bullet^{(m)})
 \longrightarrow\Tor_n^A(P_\bullet^{(r)},Q_\bullet^{(r)})
 \qquad(m\leq r)
\]
按定义就是乘以 $(f\otimes g)^{r-m}$。对（6.4.5.3）左端应用（$0_{\mathrm I}$，1.6.1），再用（6.3.6），即得结论。
\end{proof}

\egap{6.4.6}
采用（6.4.4）的记号，同样定义谱函子的典范同态
\begin{equation}
\begin{aligned}
 (u\times_S v)^*{}'\spE(\PP_\bullet,\mathscr Q_\bullet)
 &\longrightarrow{}'\spE(u^*\PP_\bullet,v^*\mathscr Q_\bullet),\\
 (u\times_S v)^*{}''\spE(\PP_\bullet,\mathscr Q_\bullet)
 &\longrightarrow{}''\spE(u^*\PP_\bullet,v^*\mathscr Q_\bullet).
\end{aligned}\tag{6.4.6.1}
\end{equation}
（6.4.5）的论证表明，当 $u,v$ 是开浸入时，这两个同态双射。事实上，结合（6.3.6.2）、（6.3.6.3），可知它们在 $E_2$ 项上是同构，而（6.4.5）表明它们在终值上是同构；再由（0，11.1.2、11.2.4）得到结论。

\subsection{拟凝聚模复形的局部超 Tor 函子：一般情形}
\source{\texttt{EGAIII.2.pdf}，PDF第15--20页，原书第16--21页的相应部分}
\egap{6.5.1}
现在设 $S$ 为任意预概形，$X,Y$ 为任意两个 $S$-预概形，$\PP_\bullet$、$\mathscr Q_\bullet$ 分别是拟凝聚 $\OO_X$-模复形和 $\OO_Y$-模复形。置 $Z=X\times_S Y$。我们将定义拟凝聚 $\OO_Z$-模 $\lTor_n^S(\PP_\bullet,\mathscr Q_\bullet)$，称为两复形的\textbf{局部超 Tor}；当 $S,X,Y$ 均仿射时，它回到（6.4）中已给出的定义。当两个复形分别只在次数0处为 $\F,\G$，而其他次数为零时，仍写 $\lTor_n^S(\F,\G)$。

\egap{6.5.2}
先设 $S$ 仿射，取 $X,Y$ 的仿射开覆盖 $(X_\lambda)$、$(Y_\mu)$。则 $Z_{\lambda\mu}=X_\lambda\times_S Y_\mu$ 是 $Z$ 的一个仿射开覆盖。置
$\PP_\bullet^\lambda=\PP_\bullet|_{X_\lambda}$、$\mathscr Q_\bullet^\mu=\mathscr Q_\bullet|_{Y_\mu}$。于是对每一对 $(\lambda,\mu)$，有拟凝聚 $\OO_{Z_{\lambda\mu}}$-模
\[
 \F_{\lambda\mu}=\lTor_n^S(\PP_\bullet^\lambda,\mathscr Q_\bullet^\mu).
\]
需证明这些模满足粘合条件（$0_{\mathrm I}$，3.3.1）。只须验证：对任意仿射开集
$U\subset X_\lambda\cap X_{\lambda'}$、$V\subset Y_\mu\cap Y_{\mu'}$，
$\F_{\lambda\mu}$ 和 $\F_{\lambda'\mu'}$ 在 $U\times_S V$ 上的限制典范同构。由（6.4.5），这两个限制都有到 $\lTor_n^S(\PP_\bullet|_U,\mathscr Q_\bullet|_V)$ 的典范同构，故结论成立。

由此定义和（6.4.5），所得 $\OO_Z$-模在同构意义下不依赖所取覆盖，故记为 $\lTor_n^S(\PP_\bullet,\mathscr Q_\bullet)$。最后，（6.4.5）还说明，对 $X,Y$ 的任意开集 $U,V$，它在 $U\times_S V$ 上的限制典范同构于 $\lTor_n^S(\PP_\bullet|_U,\mathscr Q_\bullet|_V)$。

\egap{6.5.3}
再设 $S$ 任意，取它的仿射开覆盖 $(S_\alpha)$，以 $X_\alpha,Y_\alpha$ 表示 $S_\alpha$ 在 $X,Y$ 中的逆像。还需证明各层
\[
 \F_\alpha=\lTor_n^{S_\alpha}
 (\PP_\bullet|_{X_\alpha},\mathscr Q_\bullet|_{Y_\alpha})
\]
满足粘合条件。设 $T$ 是 $S_\alpha\cap S_\beta$ 中的仿射开集，$U,V$ 分别为它在 $X,Y$ 中的逆像；只须构造两层在 $U\times_S V$ 上的限制到 $\lTor_n^T(\PP_\bullet|_U,\mathscr Q_\bullet|_V)$ 的典范同构。还可限于 $T$ 同时写为 $D(f_\alpha)$ 和 $D(f_\beta)$ 的情形，其中 $f_\alpha,f_\beta$ 分别是 $S_\alpha,S_\beta$ 上的结构层截面。由（6.4.2），$\lTor_n^T$ 分别典范同构于以 $S_\alpha,S_\beta$ 为底的相应超 Tor；而（6.5.2）已给出 $\F_\alpha,\F_\beta$ 的限制到后两者的典范同构。这就完成了 $\OO_Z$-模 $\lTor_n^S(\PP_\bullet,\mathscr Q_\bullet)$ 的定义。对任意开集 $U\subset X$、$V\subset Y$，它在 $U\times_S V$ 上的限制仍与 $\lTor_n^S(\PP_\bullet|_U,\mathscr Q_\bullet|_V)$ 典范同构。

在下有界拟凝聚模复形范畴中，这样定义了一个取值于 $\OO_Z$-模范畴的双 $\delta$-函子 $\lTor_*^S$。由（6.4.5）以及（6.4.4.2）是双 $\delta$-函子态射的事实，问题在 $X,Y,S$ 上都是局部的，故这一断言立即成立。若两个复形分别只在次数0处为 $\F,\G$，则由（6.4.1.1），$\lTor_0^S(\F,\G)$ 恰是（I，9.1.2）定义的外张量积 $\F\otimes_S\G$；这由（I，9.1.3）得出。

\egap{6.5.4}
由上述构造和（6.4.6）的说明，局部超 Tor 是两个谱函子 ${}'\spE$、${}''\spE$ 的终值，其 $E_2$ 项为
\begin{align}
 {}'\spE^2_{pq}&=\HH_p\bigl(\lTor_q^S(\PP_\bullet,\mathscr Q_\bullet)\bigr),
 \tag{6.5.4.1}\\
 {}''\spE^2_{pq}&=\bigoplus_{i+j=q}\lTor_p^S
 \bigl(\HH_i(\PP_\bullet),\HH_j(\mathscr Q_\bullet)\bigr).
 \tag{6.5.4.2}
\end{align}
谱序列（6.5.4.2）总是正则的；若两复形都下有界，则两个谱序列都双正则。下面给出另一种双正则情形。

\egap{6.5.5}
对于拓扑空间 $T$ 上的环层 $\A$，若每个 $t\in T$ 处的环 $\A_t$ 的上同调维数都 $\leq n$，则称 $\A$ 的\textbf{上同调维数 $\leq n$}；也称赋环空间 $(T,\A)$ 的上同调维数 $\leq n$。若存在这样的整数 $n$，则称环层（或赋环空间）的上同调维数\textbf{有限}。

若各 $\A_t$ 都是诺特交换局部环，那么它们的上同调维数 $\leq n$，等价于它们都是正则环且 Krull 维数 $\leq n$（$0_{\mathrm{IV}}$，17.3.1）。用第IV章将引入的维数理论术语，局部诺特预概形 $T$ 的上同调维数 $\leq n$，等价于它正则（$0_{\mathrm I}$，4.1.4）且维数 $\leq n$；这又表示对每个仿射开集 $U\subset T$，环 $\Gamma(U,\OO_T)$ 的上同调维数 $\leq n$（$0_{\mathrm{IV}}$，17.2.6）。结合（6.3.2），可知若 $S$ 局部诺特且上同调维数有限，则 ${}'\spE,{}''\spE$ 都双正则。\cmt{B2-H01}{本条“上同调维数”按各茎环的模论同调维数定义，不是结构层上同调消失的最高次数；不能将二者混用。}

通过 $X\times_S Y\simeq Y\times_S X$ 的典范同构，$\lTor_n^S(\PP_\bullet,\mathscr Q_\bullet)$ 显然转化为 $\lTor_n^S(\mathscr Q_\bullet,\PP_\bullet)$（在同构意义下）。

\egas{6.5.6}{命题}
设 $(\PP_{\bullet\alpha})$ 是拟凝聚 $\OO_X$-模复形的一个滤过归纳系统，则有典范同构
\begin{equation}
 \ilim_\alpha\lTor_n^S(\PP_{\bullet\alpha},\mathscr Q_\bullet)
 \simeq\lTor_n^S(\ilim_\alpha\PP_{\bullet\alpha},\mathscr Q_\bullet).
 \tag{6.5.6.1}
\end{equation}
\begin{proof}
问题在 $S,X,Y$ 上局部，故可设三者均仿射，此时命题归结为（6.3.6）。
\end{proof}

\egas{6.5.7}{备注}
（i）特别考虑 $S=X=Y$ 的情形，$\PP_\bullet,\mathscr Q_\bullet$ 都是拟凝聚 $\OO_S$-模复形。此时各局部超 Tor 是拟凝聚 $\OO_S$-模。对每个 $s\in S$，由（6.5.6）得到典范同构
\begin{equation}
 \bigl(\lTor_n^S(\PP_\bullet,\mathscr Q_\bullet)\bigr)_s
 \simeq\Tor_n^{\OO_{S,s}}((\PP_\bullet)_s,(\mathscr Q_\bullet)_s).
 \tag{6.5.7.1}
\end{equation}
因为问题是局部的，利用（6.4.1.1）就归结为模的情形。

（ii）定义可推广到同一个赋环空间 $(X,\OO_X)$ 上的两个 $\OO_X$-模复形 $\PP_\bullet,\mathscr Q_\bullet$。对每个开集 $U$，置
\[
 A(U)=\Gamma(U,\OO_X),\quad P_\bullet(U)=\Gamma(U,\PP_\bullet),\quad
 Q_\bullet(U)=\Gamma(U,\mathscr Q_\bullet).
\]
这些 $A(U)$-模 $\Tor_n^{A(U)}(P_\bullet(U),Q_\bullet(U))$ 构成 $X$ 上的预层；以 $\lTor_n^X(\PP_\bullet,\mathscr Q_\bullet)$ 表示其相伴 $\OO_X$-模。当 $X$ 是预概形时，由（6.3.12），这个模与上面定义的超 Tor 典范同构。这里不再进一步展开这一推广。

\egas{6.5.8}{命题}
设 $X,Y$ 为两个 $S$-预概形，$\F,\G$ 分别是拟凝聚 $\OO_X$-模和 $\OO_Y$-模。如果其中之一是 $S$-平坦的，则对于 $n\ne0$，有 $\lTor_n^S(\F,\G)=0$。
\begin{proof}
问题在 $X,Y$ 上局部，可设 $X,Y,S$ 均仿射，其环分别为 $B,C,A$，且 $\F=\widetilde M$、$\G=\widetilde N$，其中 $M$ 是 $B$-模，$N$ 是 $C$-模。例如设 $\F$ 是 $S$-平坦的；原文以局部化表述说，对于每个 $s\in S$，$M_s$ 是平坦 $A_s$-模（$0_{\mathrm I}$，6.7.1），故 $M$ 是平坦 $A$-模（$0_{\mathrm I}$，6.3.3）。已知对于 $n>0$ 及每个 $C$-模 $N$，有 $\Tor_n^A(M,N)=0$（$0_{\mathrm I}$，6.1.1），由（6.4.1.1）得出结论。\cmt{B2-H02}{这里 $M_s$ 是底环 $A$ 在 $s$ 处局部化后的模，不是 $X$ 上某个未指定点的茎；仿射情形的等价性亦见III（1.4.15.1）。负次数的普通Tor按定义为零。}
\end{proof}

\egas{6.5.9}{推论}
设 $X,Y$ 为两个 $S$-预概形，$\PP_\bullet,\mathscr Q_\bullet$ 分别为下有界拟凝聚 $\OO_X$-模复形和 $\OO_Y$-模复形。若所有 $\PP_i$ 均 $S$-平坦，则关于 $\mathscr Q_\bullet$ 有典范的 $\delta$-函子同构
\begin{equation}
 \lTor_*^S(\PP_\bullet,\mathscr Q_\bullet)
 \simeq\HH_*(\PP_\bullet\otimes_S\mathscr Q_\bullet).
 \tag{6.5.9.1}
\end{equation}
\begin{proof}
当 $S,X,Y$ 仿射时，这就是（6.3.7）；一般情形由（6.5.2）、（6.5.3）的论证得到。
\end{proof}

\egas{6.5.10}{推论}
若 $X$ 对 $S$ 平坦（$0_{\mathrm I}$，6.7.1），两复形下有界，且所有 $\PP_i$ 都是局部自由的 $\OO_X$-模（不必有限型），则（6.5.9.1）双射。
\begin{proof}
平坦性按定义是 $X$ 上逐点的性质，而任意一族平坦模的直和仍平坦（$0_{\mathrm I}$，6.1.2），故（6.5.9）的假设满足。
\end{proof}

\egas{6.5.11}{命题}
设 $X',Y'$ 为两个 $S$-预概形，$f:X\to X'$、$g:Y\to Y'$ 是仿射 $S$-态射。对于拟凝聚模复形 $\PP_\bullet,\mathscr Q_\bullet$，有典范且具有函子性的同构
\begin{equation}
 (f\times_S g)_*\lTor_n^S(\PP_\bullet,\mathscr Q_\bullet)
 \simeq\lTor_n^S(f_*\PP_\bullet,g_*\mathscr Q_\bullet).
 \tag{6.5.11.1}
\end{equation}
\begin{proof}
由于 $f,g$ 仿射，$f_*\PP_\bullet$、$g_*\mathscr Q_\bullet$ 是拟凝聚模复形（II，1.2.6）；置 $Z'=X'\times_S Y'$，两边都是拟凝聚 $\OO_{Z'}$-模（6.5.1）。容易归结到 $S,X',Y'$ 仿射，此时 $X,Y$ 也仿射，结论由（6.4.1.1）和（I，1.6.3）立即得到。
\end{proof}

\egas{6.5.12}{备注}
设 $X',Y'$ 为两个 $S$-预概形，并设 $X'$ 对 $S$ 平坦。令 $X$ 为 $X'$ 的闭子预概形，$i:X\to X'$ 为典范嵌入；设 $\PP_\bullet$、$\mathscr Q_\bullet$ 分别为下有界拟凝聚 $\OO_X$-模复形和 $\OO_{Y'}$-模复形。再设 $\LL_{\bullet\bullet}$ 为 $i_*\PP_\bullet$ 的一个由局部自由 $\OO_{X'}$-模组成的解消，且 $X'$ 的每一点都有仿射开邻域 $U$，使每个 $\LL_{j,\bullet}|_U$ 都是 $(i_*\PP_j)|_U$ 的自由解消。则有典范同构
\begin{equation}
 (i\times_S1_{Y'})_*\lTor_n^S(\PP_\bullet,\mathscr Q_\bullet)
 \simeq\HH_n(\LL_{\bullet\bullet}\otimes_S\mathscr Q_\bullet).
 \tag{6.5.12.1}
\end{equation}
若 $S,X',Y'$ 都仿射且每个 $\LL_{j,\bullet}$ 是 $i_*\PP_j$ 的自由解消，则由（6.5.11）归结到 $X'=X$ 且它对 $S$ 平坦的情形，应用（6.3.4）即可。一般情形下先局部定义此同构，再验证它是整体同构。为此，须回顾（6.3.4.1）中第一个同构的定义：它来自谱序列的同构（0，11.6.5、11.5.3），后者又来自双复形态射 $L_{\bullet\bullet}\to L'_{\bullet\bullet}$，其中 $L$ 为所给解消，$L'$ 为 $P_\bullet$ 在下有界 $A$-模复形范畴中的投射解消（参见0，11.5.2.2）。两个这样的投射解消之间存在同伦等价（M，V，1.2），因而（6.3.4.1）的同构不依赖所选的 $L'$，所需断言由此得出。

\egas{6.5.13}{命题}
设 $X,Y$ 是两个 $S$-预概形，并满足下列条件之一：
\begin{enumerate}[label=(\roman*)]
\item $X$ 和 $Z=X\times_S Y$ 局部诺特，且 $X$ 对 $S$ 平坦；
\item $S,X$ 局部诺特，且 $Y$ 对 $S$ 有限型。
\end{enumerate}
设 $\PP_\bullet,\mathscr Q_\bullet$ 分别是下有界拟凝聚 $\OO_X$-模复形和 $\OO_Y$-模复形，并且对每个 $n$，$\HH_n(\PP_\bullet)$、$\HH_n(\mathscr Q_\bullet)$ 分别为有限型模。那么各 $\lTor_n^S(\PP_\bullet,\mathscr Q_\bullet)$ 都是凝聚 $\OO_Z$-模。
\begin{proof}
两复形下有界，故第二谱序列双正则（6.5.4）。在两种情况下 $Z$ 均局部诺特；由（0，11.1.8），只需证明各 ${}''\spE^2_{pq}$ 凝聚。由同调层有限型的假设及（6.5.4.2），这归结为两个复形均集中在次数0的特殊情形，也就是以下推论。
\end{proof}

\egas{6.5.14}{推论}
满足（6.5.13）的（i）或（ii），且 $\F,\G$ 分别为拟凝聚有限型 $\OO_X$-模和 $\OO_Y$-模时，各 $\lTor_n^S(\F,\G)$ 凝聚。
\begin{proof}
问题在 $X,Y$ 上局部，可设 $S,X,Y$ 均仿射。

（i）原文接着写：“在（i）的假设下，$S,X,Z$ 是诺特的。”\cmt{B2-G02}{其中 $S$ 诺特不能由（i）推出；本证明实际只需要 $X,Z$ 诺特。报告给出非诺特底环的反例，并说明去掉这项多余断言后，证明仍成立。}
于是存在 $\F$ 的一个由有限型 $\OO_X$-模组成的局部自由解消 $\LL_\bullet$（I，1.3.7）。由于 $X$ 对 $S$ 平坦，由（6.3.4），
\[
 \lTor_n^S(\F,\G)=\HH_n(\LL_\bullet\otimes_S\G).
\]
各 $\LL_i\otimes_S\G$ 是拟凝聚有限型 $\OO_Z$-模（I，9.1.1），因而凝聚；于是其同调凝聚（$0_{\mathrm I}$，5.3.4）。

（ii）环 $A(Y)$ 是 $A(S)$ 上有限个未定元的多项式代数的商（I，6.3.3），因此 $Y$ 是某个对 $S$ 平坦且有限型的仿射 $S$-预概形 $Y'$ 的闭子预概形。$Y'$ 诺特（I，6.3.7），故对典范嵌入 $j:Y\to Y'$，存在 $j_*\G$ 的一个由有限型 $\OO_{Y'}$-模组成的局部自由解消 $\LL_\bullet$。由（6.5.12），$\lTor_n^S(\F,\G)$ 是 $\OO_{Z'}$-模 $\HH_n(\F\otimes_S\LL_\bullet)$ 沿 $1_X\times_Sj$ 的逆像，其中 $Z'=X\times_S Y'$。与（i）一样，各 $\F\otimes_S\LL_i$ 凝聚，再用（$0_{\mathrm I}$，5.3.4）即得结论。
\end{proof}

\egap{6.5.15}
以上关于两个 $S$-预概形上的拟凝聚模复形的理论，可直接推广到任意有限个 $S$-预概形 $X_i$（$1\leq i\leq m$），以及各 $X_i$ 上的拟凝聚 $\OO_{X_i}$-模复形 $\PP_\bullet^{(i)}$。置 $Z=X_1\times_S\cdots\times_S X_m$，就定义出拟凝聚 $\OO_Z$-模
\[
 \lTor_n^S(\PP_\bullet^{(1)},\ldots,\PP_\bullet^{(m)}).
\]
一般情形的理论留给读者展开；为以后引用，这里只写出以该超 Tor 为终值的第二个正则谱序列的 $E_2$ 项：
\begin{equation}
 {}''\spE^2_{pq}
 =\bigoplus_{q_1+\cdots+q_m=q}\lTor_p^S
 \bigl(\HH_{q_1}(\PP_\bullet^{(1)}),\ldots,
       \HH_{q_m}(\PP_\bullet^{(m)})\bigr).
 \tag{6.5.15.1}
\end{equation}
（6.8）将研究任意有限多个复形的这些超 Tor 函子所给出的结合性谱序列。
\par

% END text/ch3_06b.tex


% BEGIN text/ch3_06c.tex
\subsection{拟凝聚模复形的整体超 Tor 函子与 Künneth 谱序列：仿射底的情形}
\source{\texttt{EGAIII.2.pdf}，PDF第20--24页，原书第21--25页的相应部分}
\egap{6.6.1}
考虑仿射概形 $S=\Spec A$ 和两个拟紧 $S$-概形 $X^{(i)}$（$i=1,2$）。设 $\mathscr P^{(i)}_\bullet$ 是拟凝聚 $\OO_{X^{(i)}}$-模组成的下有界复形，其微分算子的次数为 $-1$（$i=1,2$）。另取 $X^{(i)}$ 的一个有限仿射开覆盖 $\mathfrak U^{(i)}=(U^{(i)}_\alpha)$。令
\[
 X=X^{(1)}\times_SX^{(2)},\qquad
 \mathfrak U=\mathfrak U^{(1)}\times_S\mathfrak U^{(2)},
\]
其中 $X$ 是拟紧 $S$-概形（I，5.5.1及6.6.4），而 $\mathfrak U$ 由仿射开集 $U^{(1)}_\alpha\times_SU^{(2)}_\beta$ 组成。对于每一对整数 $p\leq0$、$q\in\mathbf Z$，覆盖 $\mathfrak U^{(i)}$ 以层 $\mathscr P^{(i)}_q$ 为系数的交错 $(-p)$-上链群
\[
 C^{-p}(\mathfrak U^{(i)},\mathscr P^{(i)}_q)
\]
是一个 $A$-模（G，II，5.1）；对于 $p>0$，置此模为零。由此定义 $A$-模的双复形 $C^\bullet(\mathfrak U^{(i)},\mathscr P^{(i)}_\bullet)$，其两个微分算子的次数均为 $-1$。由定义（0，12.1.2），把这一双复形按总次数视为单复形时，其同调模
\[
 H_n\bigl(C^\bullet(\mathfrak U^{(i)},\mathscr P^{(i)}_\bullet)\bigr)
\]
就是超上同调模 $\mathbb H^{-n}(\mathfrak U^{(i)},\mathscr P^{(i)\bullet})$，其中微分次数为 $+1$ 的复形 $\mathscr P^{(i)\bullet}$ 以 $\mathscr P^{(i)}_{-q}$ 为第 $q$ 次分量。滥用记号，将此超上同调写为 $\mathbb H^{-n}(\mathfrak U^{(i)},\mathscr P^{(i)}_\bullet)$。由（6.2.2），这个 $A$-模典范同构于 $\mathbb H^{-n}(X^{(i)},\mathscr P^{(i)\bullet})$，后者也写作 $\mathbb H^{-n}(X^{(i)},\mathscr P^{(i)}_\bullet)$；因此它与所选的有限覆盖 $\mathfrak U^{(i)}$ 无关。

\egap{6.6.2}
我们将关于双复形的函子超同调一般理论（0，11.7.4），应用于两个 $A$-模双复形
\[
 L^{(i)}_{\bullet\bullet}=C^\bullet(\mathfrak U^{(i)},\mathscr P^{(i)}_\bullet),\qquad i=1,2,
\]
以及关于这两个双复形的协变双函子 $L^{(1)}_{\bullet\bullet}\otimes_AL^{(2)}_{\bullet\bullet}$。由于上链取成交错上链，且覆盖 $\mathfrak U^{(i)}$ 有限，模 $C^{-p}(\mathfrak U^{(i)},\mathscr P^{(i)}_q)$ 只可能在有限多个 $p$ 值上非零，而且这些 $p$ 值与 $q$ 无关；特别地，每个 $L^{(i)}_{\bullet\bullet}$ 的两个次数都下有界。以
\[
 \bTor^A_n(L^{(1)}_{\bullet\bullet},L^{(2)}_{\bullet\bullet})
 \quad\text{或}\quad
 \bTor^S_n(\mathfrak U^{(1)},\mathfrak U^{(2)};
             \mathscr P^{(1)}_\bullet,\mathscr P^{(2)}_\bullet)
\]
表示 $L^{(1)}_{\bullet\bullet}\otimes_AL^{(2)}_{\bullet\bullet}$ 的第 $n$ 个超同调模，称为两个复形关于覆盖 $\mathfrak U^{(1)},\mathfrak U^{(2)}$ 的第 $n$ 个\textbf{超 Tor}。
\cmt{排印记号}{本小节用粗体 $\bTor$ 区分超 Tor 与逐项普通 $\Tor$；用 $\mathbb H$ 区分超上同调与逐项上同调 $H$。这只是恢复底本所区分的对象，不改变其定义。双复形的上标 $\mathrm I,\mathrm{II}$ 指相应方向，沿用第0章的约定。}
若两个复形各自只有第零项 $\F^{(1)},\F^{(2)}$ 非零，则写作 $\Tor^S_n(\mathfrak U^{(1)},\mathfrak U^{(2)};\F^{(1)},\F^{(2)})$。按照一般约定，
\[
 \Tor^S_n(\mathfrak U^{(1)},\mathfrak U^{(2)};
          \mathscr P^{(1)}_\bullet,\mathscr P^{(2)}_\bullet)
\]
则表示双次数 $(j,k)$ 的分量为
$\Tor^S_n(\mathfrak U^{(1)},\mathfrak U^{(2)};\mathscr P^{(1)}_j,\mathscr P^{(2)}_k)$ 的双复形。

覆盖中开集的交均仿射，故 $L^{(i)}_{\bullet\bullet}$ 关于 $\mathscr P^{(i)}_\bullet$ 是正合函子（I，5.5.6及1.3.11）。因此上述超 Tor 关于两个复形构成取值于 $A$-模范畴的协变双 $\delta$-函子（0，11.7.3）。又由（0，11.7.2），这一双函子是六个双正则谱函子的共同终值；记为
\[
 {}^{(t)}E(\mathfrak U^{(1)},\mathfrak U^{(2)};
             \mathscr P^{(1)}_\bullet,\mathscr P^{(2)}_\bullet),
 \qquad t=a,b,a',b',c,d,
\]
或简记 ${}^{(t)}E(\mathfrak U^{(1)},\mathfrak U^{(2)};\mathscr P^{(1)},\mathscr P^{(2)})$。其 $E^2$ 项如下；为免横向过宽，在这些公式中简写 $L^{(i)}=L^{(i)}_{\bullet\bullet}$：
\begin{align*}
 {}^{(a)}E^2_{pq}
 &=\bigoplus_{q_1+q_2=q}\Tor^A_p
       \bigl(H_{q_1}(L^{(1)}),H_{q_2}(L^{(2)})\bigr),\\
 {}^{(b)}E^2_{pq}
 &=H_p\bigl(\bTor^{A,\mathrm{II}}_q(L^{(1)},L^{(2)})\bigr),\\
 {}^{(a')}E^2_{pq}
 &=\bigoplus_{q_1+q_2=q}\bTor^A_p
       \bigl(H^{\mathrm I}_{q_1}(L^{(1)}),H^{\mathrm I}_{q_2}(L^{(2)})\bigr),\\
 {}^{(b')}E^2_{pq}
 &=H_p\bigl(\Tor^A_q(L^{(1)},L^{(2)})\bigr),\\
 {}^{(c)}E^2_{pq}
 &=\bigoplus_{q_1+q_2=q}\bTor^A_p
       \bigl(H^{\mathrm{II}}_{q_1}(L^{(1)}),H^{\mathrm{II}}_{q_2}(L^{(2)})\bigr),\\
 {}^{(d)}E^2_{pq}
 &=H_p\bigl(\bTor^{A,\mathrm I}_q(L^{(1)},L^{(2)})\bigr).
\end{align*}
记号符合超同调一般理论的约定。下面进一步明确这些初始项。

\egap{6.6.3}
\textbf{谱序列 $(a)$ 与 $(a')$。}由（6.6.1），双复形 $L^{(i)}$ 的同调模 $H_{q_i}(L^{(i)})$ 等于 $\mathbb H^{-q_i}(X^{(i)},\mathscr P^{(i)}_\bullet)$，故
\[
 {}^{(a)}E^2_{pq}=
 \bigoplus_{q_1+q_2=q}\Tor^A_p
 \bigl(\mathbb H^{-q_1}(X^{(1)},\mathscr P^{(1)}_\bullet),
       \mathbb H^{-q_2}(X^{(2)},\mathscr P^{(2)}_\bullet)\bigr).
\]
按定义，复形 $H^{\mathrm I}_{q_i}(L^{(i)})$ 的第 $k$ 项是同调模
$H_{q_i}(C^\bullet(\mathfrak U^{(i)},\mathscr P^{(i)}_k))$，即上同调模
$H^{-q_i}(\mathfrak U^{(i)},\mathscr P^{(i)}_k)$。由（1.4.1），它典范同构于
$H^{-q_i}(X^{(i)},\mathscr P^{(i)}_k)$。所以
\[
 {}^{(a')}E^2_{pq}=
 \bigoplus_{q_1+q_2=q}\bTor^A_p
 \bigl(H^{-q_1}(X^{(1)},\mathscr P^{(1)}_\bullet),
       H^{-q_2}(X^{(2)},\mathscr P^{(2)}_\bullet)\bigr).
\]

\egap{6.6.4}
\textbf{谱序列 $(b)$ 与 $(b')$。}按定义，$\bTor^{A,\mathrm{II}}_q(L^{(1)},L^{(2)})$ 的双次数 $(h,k)$ 分量是
\[
 \bTor^A_q\bigl(C^{-h}(\mathfrak U^{(1)},\mathscr P^{(1)}_\bullet),
                  C^{-k}(\mathfrak U^{(2)},\mathscr P^{(2)}_\bullet)\bigr).
\]
设 $\Phi^{(i)}$ 为覆盖 $\mathfrak U^{(i)}$ 的指标集。按定义，复形
$C^j(\mathfrak U^{(i)},\mathscr P^{(i)}_\bullet)$（$j\geq0$）是各复形
$\Gamma(U^{(i)}_\rho,\mathscr P^{(i)}_\bullet)$ 的直和，其中
$U^{(i)}_\rho$ 为各 $\alpha\in\rho$ 对应的 $U^{(i)}_\alpha$ 之交，而 $\rho$ 遍历相应的有限指标子集。因此上述 $A$-模是各模
\[
 \bTor^A_q\bigl(\Gamma(U^{(1)}_\sigma,\mathscr P^{(1)}_\bullet),
                  \Gamma(U^{(2)}_\tau,\mathscr P^{(2)}_\bullet)\bigr)
\]
的直和。底本在此规定
$\#\sigma=-(h+1)$、$\#\tau=-(k+1)$。
\cmt{B2-T03}{这两个基数公式是底本所印表达；对于通常的 Čech 上链 $C^r$，应有 $r+1$ 个指标，因此这里应是 $\#\sigma=1-h$、$\#\tau=1-k$。后文的基数总和也有对应的符号错误。详见核查报告；本处不把负基数当成有效定义使用。}
因 $X^{(1)},X^{(2)}$ 为概形，各 $U^{(i)}_\rho$ 仿射，由（6.4.1.1）得
\begin{align*}
 &\bTor^A_q\bigl(\Gamma(U^{(1)}_\sigma,\mathscr P^{(1)}_\bullet),
                  \Gamma(U^{(2)}_\tau,\mathscr P^{(2)}_\bullet)\bigr)\\
 &\hspace{1em}=\Gamma\bigl(U^{(1)}_\sigma\times_SU^{(2)}_\tau,
                  \lTor^S_q(\mathscr P^{(1)}_\bullet,\mathscr P^{(2)}_\bullet)\bigr).
\end{align*}
因此 ${}^{(b)}E^2_{pq}$ 是 $\Phi^{(1)},\Phi^{(2)}$ 上以系数系统
\[
 \mathscr V:(\sigma,\tau)\longmapsto
 \Gamma\bigl(U^{(1)}_\sigma\times_SU^{(2)}_\tau,
                  \lTor^S_q(\mathscr P^{(1)}_\bullet,\mathscr P^{(2)}_\bullet)\bigr)
\]
为值的双交错上链复形 $L^\bullet(\Phi^{(1)},\Phi^{(2)};\mathscr V)$ 的第 $(-p)$ 个上同调模（0，11.8.4）。由（0，11.8.5及11.8.6），其上同调与所有上链组成的复形 $C^\bullet(\Phi^{(1)},\Phi^{(2)};\mathscr V)$ 相同，也与复形 $P^\bullet(\Phi^{(1)},\Phi^{(2)};\mathscr V)$ 相同；后者的元素为
\[
 \lambda(\sigma,\tau)\in
 \Gamma\bigl(U^{(1)}_\sigma\times_SU^{(2)}_\tau,
                  \lTor^S_q(\mathscr P^{(1)}_\bullet,\mathscr P^{(2)}_\bullet)\bigr)
\]
的线性组合，$\sigma=(\alpha_0,\ldots,\alpha_m)$ 与
$\tau=(\beta_0,\ldots,\beta_m)$ 是项数相同的序列。此时
\[
 U^{(1)}_\sigma\times_SU^{(2)}_\tau
 =\bigcap_{r=0}^m\bigl(U^{(1)}_{\alpha_r}\times_SU^{(2)}_{\beta_r}\bigr)
\]
（I，3.2.7）。以 $\mathfrak U$ 表示 $Z=X^{(1)}\times_SX^{(2)}$ 的这一仿射开覆盖；考虑到 $Z$ 为概形，最后由（1.3.1）得到
\[
 {}^{(b)}E^2_{pq}=H^{-p}\bigl(X^{(1)}\times_SX^{(2)},
               \lTor^S_q(\mathscr P^{(1)}_\bullet,\mathscr P^{(2)}_\bullet)\bigr).
\]

其次，$\Tor^A_q(L^{(1)},L^{(2)})$ 的双次数 $(h,k)$ 分量，是满足
$h_1+h_2=h$、$k_1+k_2=k$ 的各模
\[
 \Tor^A_q\bigl(C^{-h_1}(\mathfrak U^{(1)},\mathscr P^{(1)}_{k_1}),
                 C^{-h_2}(\mathfrak U^{(2)},\mathscr P^{(2)}_{k_2})\bigr)
\]
之直和。如上明确各上链模后，可见这一分量又是各模
\[
 \Gamma\bigl(U^{(1)}_\sigma\times_SU^{(2)}_\tau,
             \lTor^S_q(\mathscr P^{(1)}_{k_1},\mathscr P^{(2)}_{k_2})\bigr)
\]
之直和，其中 $k_1+k_2=k$；底本写 $\#\sigma+\#\tau=-h-2$。
\cmt{B2-T03续}{沿用同一 Čech 次数约定，总和应为 $2-h$，而不是 $-h-2$。这与前一处是同一类基数符号误植，未重复计为另一项定理错误。}
我们计算的 ${}^{(b')}E^2_{pq}$ 是双复形 $N^{\bullet\bullet}=(N^{hk})$ 的第 $(-p)$ 个上同调模，其中单复形 $N^{\bullet,k}$ 为 $\Phi^{(1)},\Phi^{(2)}$ 上以
\[
 \mathscr V_k:(\sigma,\tau)\longmapsto
 \Gamma\left(U^{(1)}_\sigma\times_SU^{(2)}_\tau,
       \bigoplus_{k_1+k_2=k}\lTor^S_q
                (\mathscr P^{(1)}_{-k_1},\mathscr P^{(2)}_{-k_2})\right)
\]
为系数系统的双交错上链复形。这些系数系统构成复形 $\mathscr V^\bullet$，其微分来自逐项 Tor 双复形所对应单复形的微分。由（0，11.8.9），$N^{\bullet\bullet}$ 的上同调与
$C^\bullet(\Phi^{(1)},\Phi^{(2)};\mathscr V^\bullet)$ 相同；也与
$P^\bullet(\Phi^{(1)},\Phi^{(2)};\mathscr V^\bullet)$ 相同，后者的双次数 $(h,k)$ 元素为
\[
 \lambda(\sigma,\tau)\in
 \Gamma\left(U^{(1)}_\sigma\times_SU^{(2)}_\tau,
       \bigoplus_{k_1+k_2=k}\lTor^S_q
                (\mathscr P^{(1)}_{-k_1},\mathscr P^{(2)}_{-k_2})\right)
\]
的线性组合，$\sigma$、$\tau$ 是项数相同的序列（0，11.8.10）。如上可见，${}^{(b')}E^2_{pq}$ 是双复形 $C^\bullet(\mathfrak U,\mathscr Q^\bullet)$ 的第 $(-p)$ 个上同调模，其中 $\mathscr Q^\bullet$ 为逐项 Tor 双复形对应的单复形。按照（6.6.1）的约定，故有
\[
 {}^{(b')}E^2_{pq}=\mathbb H^{-p}\bigl(X^{(1)}\times_SX^{(2)},
                    \lTor^S_q(\mathscr P^{(1)}_\bullet,\mathscr P^{(2)}_\bullet)\bigr),
\]
这里第二个变量中的 $\lTor$ 指上述逐项构成的双复形。

\egap{6.6.5}
\textbf{谱序列 $(c)$ 与 $(d)$。}由上链函子 $C^{-h}$ 的正合性，复形
$H^{\mathrm{II}}_{q_i}(L^{(i)})$ 的第 $h$ 项是
$C^{-h}(\mathfrak U^{(i)},\HH_{q_i}(\mathscr P^{(i)}_\bullet))$。由两个覆盖对应的超 Tor 定义（6.6.2），得到
\[
 {}^{(c)}E^2_{pq}=\bigoplus_{q_1+q_2=q}
 \Tor^S_p\bigl(\mathfrak U^{(1)},\mathfrak U^{(2)};
       \HH_{q_1}(\mathscr P^{(1)}_\bullet),\HH_{q_2}(\mathscr P^{(2)}_\bullet)\bigr).
\]
最后，按定义，$\bTor^{A,\mathrm I}_q(L^{(1)},L^{(2)})$ 的双次数 $(h,k)$ 分量是
\[
 \Tor^S_q(\mathfrak U^{(1)},\mathfrak U^{(2)};\mathscr P^{(1)}_h,\mathscr P^{(2)}_k),
\]
故
\[
 {}^{(d)}E^2_{pq}=H_p\bigl(\Tor^S_q(\mathfrak U^{(1)},\mathfrak U^{(2)};
                    \mathscr P^{(1)}_\bullet,\mathscr P^{(2)}_\bullet)\bigr).
\]

\egap{6.6.6}
关于双复形函子的超同调理论（0，11.7.3）表明，如同（6.3.4），若
$M^{(i)}_{\bullet\bullet\bullet}$ 是 $L^{(i)}_{\bullet\bullet}$ 在下有界模复形范畴中的一个平坦 Cartan--Eilenberg 解消（$i=1,2$），则有双 $\delta$-函子的典范同构
\begin{equation}\tag{6.6.6.1}
\begin{aligned}
 &\bTor^S_*(\mathfrak U^{(1)},\mathfrak U^{(2)};
                    \mathscr P^{(1)}_\bullet,\mathscr P^{(2)}_\bullet)\\
 &\quad\simeq H_*(M^{(1)}_{\bullet\bullet\bullet}\otimes_AM^{(2)}_{\bullet\bullet\bullet})\\
 &\quad\simeq H_*(M^{(1)}_{\bullet\bullet\bullet}\otimes_AL^{(2)}_{\bullet\bullet})
 \simeq H_*(L^{(1)}_{\bullet\bullet}\otimes_AM^{(2)}_{\bullet\bullet\bullet}).
\end{aligned}
\end{equation}

\egap{6.6.7}
现在证明（6.6.2）定义的整体超 Tor 及相应六个谱序列，在典范同构意义下，与定义时使用的有限仿射开覆盖无关。只须证明：若 $\mathfrak V^{(i)}$ 是另两个同类覆盖，且分别比 $\mathfrak U^{(i)}$ 更细，则有谱函子的典范同构
\begin{equation}\tag{6.6.7,$t$}
\begin{aligned}
 &{}^{(t)}E(\mathfrak U^{(1)},\mathfrak U^{(2)};
                  \mathscr P^{(1)}_\bullet,\mathscr P^{(2)}_\bullet)\\
 &\hspace{2em}\simeq{}^{(t)}E(\mathfrak V^{(1)},\mathfrak V^{(2)};
                  \mathscr P^{(1)}_\bullet,\mathscr P^{(2)}_\bullet),
\end{aligned}
\end{equation}
其中 $t=a,b,a',b',c,d$。

对 $i=1,2$，有在同伦意义下良好定义的双复形同态
\[
 C^\bullet(\mathfrak U^{(i)},\mathscr P^{(i)}_\bullet)
 \longrightarrow C^\bullet(\mathfrak V^{(i)},\mathscr P^{(i)}_\bullet)
\]
（G，II，5.7.1）。所以（6.6.7，$t$）中的同态已有典范定义，并在终值上与边缘算子相容（0，11.3.2）。此外，上面对 $(a),(b),(a'),(b')$ 的 $E^2$ 项的计算表明，对这些谱序列，（6.6.7，$t$）在 $E^2$ 项上是同构；由于这些谱序列双正则，便知它对这“三个谱函子”是同构，因而在共同终值上是双 $\delta$-函子的同构（0，11.1.5）。
\cmt{B2-T04}{底本这里印为“这三个”（ces trois），但前一句明确列举 $(a),(b),(a'),(b')$ 四个。应理解为“四个”；不影响所用的谱序列比较论证。}
特别地，对于拟凝聚 $\OO_{X^{(i)}}$-模 $\F^{(i)}$（$i=1,2$），典范同态
\[
 \Tor^S_n(\mathfrak U^{(1)},\mathfrak U^{(2)};\F^{(1)},\F^{(2)})
 \longrightarrow
 \Tor^S_n(\mathfrak V^{(1)},\mathfrak V^{(2)};\F^{(1)},\F^{(2)})
\]
双射。由（6.6.5）的计算，（6.6.7，$t$）在 $t=c,d$ 时也在 $E^2$ 项上给出同构。如上即可推出它在 $t=c,d$ 时也是谱序列的同构。

可把这些同构视为：在由 $X^{(1)},X^{(2)}$ 的有限仿射开覆盖对组成的滤过指标集上，定义了谱函子的归纳系统。以
\[
 {}^{(t)}E(X^{(1)},X^{(2)};\mathscr P^{(1)}_\bullet,\mathscr P^{(2)}_\bullet)
\]
表示此系统的归纳极限；也可略去 $\bullet$。以
\[
 \bTor^S_n(X^{(1)},X^{(2)};\mathscr P^{(1)}_\bullet,\mathscr P^{(2)}_\bullet)
\]
表示这一谱函子的终值，称为两个复形的\textbf{整体超 Tor}。当两个复形只在第零项非零时，写作
$\Tor^S_n(X^{(1)},X^{(2)};\F^{(1)},\F^{(2)})$。依一般约定，
$\Tor^S_n(X^{(1)},X^{(2)};\mathscr P^{(1)}_\bullet,\mathscr P^{(2)}_\bullet)$ 则表示由逐项的 $\Tor^S_n(X^{(1)},X^{(2)};\mathscr P^{(1)}_j,\mathscr P^{(2)}_k)$ 组成的双复形。

\egap{6.6.8}
保持（6.6.1）的假设，另考虑两个 $S$-态射
$f_i:X^{(i)}\to Y^{(i)}$，其中 $Y^{(i)}=\Spec B_i$ 是仿射 $S$-概形，故 $B_i$ 为 $A$-代数（$i=1,2$）。这给出 $A$-同态
$B_i\to\Gamma(X^{(i)},\OO_{X^{(i)}})$（I，2.2.4），所以（6.6.2）中的每个 $L^{(i)}$ 都是 $B_i$-模双复形。因而 $L^{(1)}\otimes_AL^{(2)}$ 为 $(B_1\otimes_AB_2)$-模的四重复形，其六个超同调谱函子可看成取值于 $(B_1\otimes_AB_2)$-模谱序列范畴的函子。置
\[
 Y=Y^{(1)}\times_SY^{(2)}=\Spec(B_1\otimes_AB_2),
\]
并考虑这些模所对应的拟凝聚 $\OO_Y$-模（I，1.3.4）。记六个 $\OO_Y$-模谱序列
\[
 \bigl({}^{(t)}E(X^{(1)},X^{(2)};
                  \mathscr P^{(1)}_\bullet,\mathscr P^{(2)}_\bullet)\bigr)^\sim
\]
为 ${}^{(t)}\spE(f_1,f_2;\mathscr P^{(1)}_\bullet,\mathscr P^{(2)}_\bullet)$；其中
$t=a,b,a',b',c,d$。其共同终值
\[
 \bigl(\bTor^S_n(X^{(1)},X^{(2)};
                  \mathscr P^{(1)}_\bullet,\mathscr P^{(2)}_\bullet)\bigr)^\sim
\]
记为 $\lTor^S_n(f_1,f_2;\mathscr P^{(1)}_\bullet,\mathscr P^{(2)}_\bullet)$。若 $\mathscr P^{(i)}_\bullet$ 只在第零项 $\F^{(i)}$ 非零，则记为
$\lTor^S_n(f_1,f_2;\F^{(1)},\F^{(2)})$。



% END text/ch3_06c.tex


% BEGIN continuations/iii2_67.tex
% 直接依据用户上传 EGAIII.2.pdf 第24--31页；原书第25--32页。
\subsection{拟凝聚模复形的整体超 Tor 函子与 Künneth 谱序列：一般情形}
\source{\texttt{EGAIII.2.pdf}，PDF第24--31页；原书第25--32页}
\egap{6.7.1}
现在将（6.6.8）的定义推广到如下情形：$S$ 为任意预概形，$X^{(i)},Y^{(i)}$ 是 $S$-预概形，$f_i:X^{(i)}\to Y^{(i)}$ 为分离拟紧态射。对于每一对下有界的拟凝聚 $\OO_{X^{(i)}}$-模复形 $\mathscr P^{(i)}_\bullet$（$i=1,2$），要对每个 $n$ 定义拟凝聚 $\OO_Y$-模
\[
 \lTor_n^S(f_1,f_2;\mathscr P^{(1)}_\bullet,\mathscr P^{(2)}_\bullet),
 \qquad Y=Y^{(1)}\times_SY^{(2)},
\]
以及六个谱函子；当 $S,Y^{(1)},Y^{(2)}$ 仿射时，它们应当还原为（6.6.8）的定义。

先设 $S=\Spec A$ 仿射，而 $Y^{(1)},Y^{(2)}$ 任意。取 $Y^{(i)}$ 的仿射开集 $W^{(i)}$；$f_i^{-1}(W^{(i)})$ 是拟紧 $S$-概形，而 $W=W^{(1)}\times_SW^{(2)}$ 是 $Y$ 的仿射开集。以 $f'_i:f_i^{-1}(W^{(i)})\to W^{(i)}$ 表示 $f_i$ 的限制，以 $\mathscr P'^{(i)}_\bullet$ 表示 $\mathscr P^{(i)}_\bullet|f_i^{-1}(W^{(i)})$（$i=1,2$）。由（6.6.8），有拟凝聚 $(\OO_Y|W)$-模谱序列
\[
 {}^{(t)}\spE(f'_1,f'_2;\mathscr P'^{(1)}_\bullet,\mathscr P'^{(2)}_\bullet).
\]
须验证它们满足粘合条件（$0_{\mathrm I}$，3.3.1）。立即归结为 $Y^{(i)}=\Spec B_i$ 仿射、$W^{(i)}=D(g_i)$（$g_i\in B_i$）的情形；此时 $W=D(g_1\otimes g_2)\subset Y=\Spec(B_1\otimes_AB_2)$（II，4.3.2.1）。若 $X'^{(i)}=f_i^{-1}(W^{(i)})$，所求为谱函子的典范同构
\begin{equation}\tag{6.7.1.1}
\begin{aligned}
 {}^{(t)}E(X'^{(1)},X'^{(2)};\mathscr P'^{(1)}_\bullet,\mathscr P'^{(2)}_\bullet)
 \simeq{}^{(t)}E(X^{(1)},X^{(2)};\mathscr P^{(1)}_\bullet,\mathscr P^{(2)}_\bullet)\otimes_BB_g,
\end{aligned}
\end{equation}
其中 $B=B_1\otimes_AB_2$，$g=g_1\otimes g_2$。

为此，对 $X^{(i)}$ 取有限仿射开覆盖 $\mathfrak U^{(i)}$，并以 $\mathfrak U'^{(i)}$ 表示它在 $X'^{(i)}$ 上的迹；后者仍由仿射开集组成（I，5.5.10）。确切地说，
\[
 C^\bullet(\mathfrak U'^{(i)},\mathscr P'^{(i)}_\bullet)
 =C^\bullet(\mathfrak U^{(i)},\mathscr P^{(i)}_\bullet)\otimes_{B_i}(B_i)_{g_i}.
\]
置 $L'^{(i)}_{\bullet\bullet}=C^\bullet(\mathfrak U'^{(i)},\mathscr P'^{(i)}_\bullet)$，则
\[
 L'^{(1)}_{\bullet\bullet}\otimes_AL'^{(2)}_{\bullet\bullet}
 =\bigl(L^{(1)}_{\bullet\bullet}\otimes_{B_1}(B_1)_{g_1}\bigr)
 \otimes_A\bigl(L^{(2)}_{\bullet\bullet}\otimes_{B_2}(B_2)_{g_2}\bigr).
\]
由于 $B_g\simeq(B_1)_{g_1}\otimes_A(B_2)_{g_2}$，有典范同构
\[
 L'^{(1)}_{\bullet\bullet}\otimes_AL'^{(2)}_{\bullet\bullet}
 \simeq(L^{(1)}_{\bullet\bullet}\otimes_AL^{(2)}_{\bullet\bullet})\otimes_BB_g.
\]
若 $M^{(i)}_{\bullet\bullet\bullet}$ 为 $L^{(i)}_{\bullet\bullet}$ 的一个投射 Cartan--Eilenberg 解消，并可设其由 $B_i$-模组成，则由 $(B_i)_{g_i}$ 在 $B_i$ 上平坦，
\[
 M'^{(i)}_{\bullet\bullet\bullet}
 =M^{(i)}_{\bullet\bullet\bullet}\otimes_{B_i}(B_i)_{g_i}
\]
是双复形 $L'^{(i)}_{\bullet\bullet}$ 的投射 Cartan--Eilenberg 解消；并有
\[
 M'^{(1)}_{\bullet\bullet\bullet}\otimes_AM'^{(2)}_{\bullet\bullet\bullet}
 =(M^{(1)}_{\bullet\bullet\bullet}\otimes_AM^{(2)}_{\bullet\bullet\bullet})\otimes_BB_g.
\]
所求同构（6.7.1.1）遂由双复形超同调的定义，以及关于 $B$-模 $G$ 的函子 $G\otimes_BB_g$ 的正合性得到。
\cmt{C3-P01}{原文此处没有明确“投射”相对于哪个环。一般不能把任意 $A$-投射解消自动赋予 $B_i$-模结构，也不能把 $B_i$-投射自动等同于 $A$-投射；这里保留原文步骤并标为需补明的解消选择问题，不据此判定局部化同构错误。与（6.4.5）中通过滤过归纳极限证明局部化相容的论证应加区分。}

\egap{6.7.2}
现设 $S$ 任意，$u_i:Y^{(i)}\to S$ 为结构态射（$i=1,2$）。取 $S$ 的仿射开覆盖 $(S_\alpha)$，置
\[
 Y^{(i)}_\alpha=u_i^{-1}(S_\alpha),\qquad
 X^{(i)}_\alpha=f_i^{-1}(Y^{(i)}_\alpha),
\]
并以 $f_{i\alpha}:X^{(i)}_\alpha\to Y^{(i)}_\alpha$ 表示 $f_i$ 的限制；它分离且拟紧。各 $Y_\alpha=Y^{(1)}_\alpha\times_{S_\alpha}Y^{(2)}_\alpha$ 构成 $Y$ 的开覆盖，在每个 $Y_\alpha$ 上，由（6.7.1）已定义谱函子
\[
 {}^{(t)}\spE^{S_\alpha}\bigl(f_{1\alpha},f_{2\alpha};
 \mathscr P^{(1)}_\bullet|X^{(1)}_\alpha,
 \mathscr P^{(2)}_\bullet|X^{(2)}_\alpha\bigr).
\]
仍须证明这些函子满足粘合条件。立即归结为如下情形：$S=\Spec A$ 仿射，$S'=D(h)$，$h\in A$，且 $u_i(Y^{(i)})\subset S'$；还可假设 $Y^{(i)}=\Spec B_i$ 仿射。需定义典范同构
\begin{equation}\tag{6.7.2.1}
 {}^{(t)}E^S(X^{(1)},X^{(2)};\mathscr P^{(1)}_\bullet,\mathscr P^{(2)}_\bullet)
 \simeq{}^{(t)}E^{S'}(X^{(1)},X^{(2)};\mathscr P^{(1)}_\bullet,\mathscr P^{(2)}_\bullet).
\end{equation}
采用（6.6.2）的记号，$L^{(i)}_{\bullet\bullet}$ 由 $A_h$-模组成，因此
\[
 L^{(1)}_{\bullet\bullet}\otimes_{A_h}L^{(2)}_{\bullet\bullet}
 \simeq L^{(1)}_{\bullet\bullet}\otimes_AL^{(2)}_{\bullet\bullet}.
\]
由于可以取由 $A_h$-模组成的 $L^{(i)}_{\bullet\bullet}$ 的投射 Cartan--Eilenberg 解消 $M^{(i)}_{\bullet\bullet\bullet}$，立即给出所求的典范同构。
\cmt{C3-P01}{这句话的环上投射性约定同样需要说明；$A_h$ 在 $A$ 上平坦，但一般不投射。保留原文，不把平坦与投射混为一谈。}
综上，得到下述定理。

\egas{6.7.3}{定理}
设 $S$ 为预概形，$f_i:X^{(i)}\to Y^{(i)}$ 为 $S$-预概形间的分离拟紧 $S$-态射，$\mathscr P^{(i)}_\bullet$ 为下有界拟凝聚 $\OO_{X^{(i)}}$-模复形（$i=1,2$）。置 $Y=Y^{(1)}\times_SY^{(2)}$。存在取值于拟凝聚 $\OO_Y$-模范畴的双 $\delta$-函子
$\lTor_*^S(f_1,f_2;\mathscr P^{(1)}_\bullet,\mathscr P^{(2)}_\bullet)$，使得：若 $V^{(i)}$ 是 $Y^{(i)}$ 的仿射开集，$V=V^{(1)}\times_SV^{(2)}$，则
\[
\begin{aligned}
 &\lTor_*^S(f_1,f_2;\mathscr P^{(1)}_\bullet,\mathscr P^{(2)}_\bullet)|V\\
 &\quad=\bigl(\bTor_*^S(f_1^{-1}(V^{(1)}),f_2^{-1}(V^{(2)});\\
 &\hspace{4em}\mathscr P^{(1)}_\bullet|f_1^{-1}(V^{(1)}),
 \mathscr P^{(2)}_\bullet|f_2^{-1}(V^{(2)}))\bigr)^\sim.
\end{aligned}
\]
\cmt{C3-H01}{这一仿射计算式须在（6.7.1）--（6.7.2）的局部化语境中理解：先在 $S$ 上取共同的仿射开片，再取其上方的 $V^{(i)}$。对任意 $S$，两个仿射开片的纤维积未必仿射；不能脱离该约定，将右侧理解为任意非仿射 $V$ 上未经定义的“模的相伴层”。}
此双函子是六个双正则谱函子
\[
 {}^{(t)}\spE(f_1,f_2;\mathscr P^{(1)}_\bullet,\mathscr P^{(2)}_\bullet),
 \qquad t=a,b,a',b',c,d,
\]
的共同终值。其 $E^2$ 项如下：
\begin{align*}
 {}^{(a)}\spE^2_{pq}
 &=\bigoplus_{q_1+q_2=q}\lTor_p^S
 \bigl(\cH^{-q_1}(f_1,\mathscr P^{(1)}_\bullet),
       \cH^{-q_2}(f_2,\mathscr P^{(2)}_\bullet)\bigr),\\
 {}^{(b)}\spE^2_{pq}
 &=\cH^{-p}\bigl(f_1\times_Sf_2,
                  \lTor_q^S(\mathscr P^{(1)}_\bullet,\mathscr P^{(2)}_\bullet)\bigr),\\
 {}^{(a')}\spE^2_{pq}
 &=\bigoplus_{q_1+q_2=q}\lTor_p^S
 \bigl(\cH^{-q_1}(f_1,\mathscr P^{(1)}),
       \cH^{-q_2}(f_2,\mathscr P^{(2)})\bigr),\\
 {}^{(b')}\spE^2_{pq}
 &=\cH^{-p}\bigl(f_1\times_Sf_2,
                  \termTor_q^S(\mathscr P^{(1)}_\bullet,\mathscr P^{(2)}_\bullet)\bigr),\\
 {}^{(c)}\spE^2_{pq}
 &=\bigoplus_{q_1+q_2=q}\lTor_p^S
 \bigl(f_1,f_2;\HH_{q_1}(\mathscr P^{(1)}_\bullet),
                 \HH_{q_2}(\mathscr P^{(2)}_\bullet)\bigr),\\
 {}^{(d)}\spE^2_{pq}
 &=\HH_p\bigl(\termTor_q^S(f_1,f_2;
                   \mathscr P^{(1)}_\bullet,\mathscr P^{(2)}_\bullet)\bigr).
\end{align*}
\cmt{C3-N01}{为避免原书不同“Tor”字体在中文排版中混淆，本式额外用上标 $\mathrm{term}$ 标明逐项应用模的 Tor 后组成的双复形，不是两个复形的超 Tor 终值。$(a')$ 中不带下标点的 $\cH^{-q_i}(f_i,\mathscr P^{(i)})$ 也指逐项的高阶直接像复形；与 $(a)$ 中整个复形的超上同调须区别。此为记号说明，不是增添数学假设。}
称 $(a)$ 与 $(b)$ 为 \textbf{Künneth 谱序列}。当 $\mathscr P^{(1)}_\bullet,\mathscr P^{(2)}_\bullet$ 都仅在次数0处非零时，$(a)$ 与 $(a')$ 相同，$(b)$ 与 $(b')$ 相同；这时 $(c),(d)$ 退化，因而不再提供信息。

\egas{6.7.4}{注记}
上述整体超 Tor 同时包括（6.2.1）的超上同调模以及（6.5.3）的局部超 Tor。对任意分离拟紧态射 $f:X\to Y$ 和下有界拟凝聚 $\OO_X$-模复形 $\PP_\bullet$，有关于 $\PP_\bullet$ 的 $\delta$-函子典范同构
\begin{equation}\tag{6.7.4.1}
 \lTor_n^Y(f,1_Y;\PP_\bullet,\OO_Y)
 \simeq\cH^{-n}(f,\PP_\bullet)\qquad(n\in\mathbf Z).
\end{equation}
事实上，由（6.7.2）的粘合方法立即归结为 $Y$ 仿射。利用（6.2.2），可用同一个有限仿射开覆盖 $\mathfrak U$ 计算两边，并对第一边将 $Y$ 自身作为 $Y$ 的覆盖。
\cmt{C3-T01}{原书第27页把 $\mathfrak U$ 写成“$Y$ 的”有限仿射覆盖；这里所需的是 $X$ 的有限仿射覆盖，否则不能计算 $\cH(f,\PP_\bullet)$。上文略去错置对象，正确对象在本注说明。}
按（6.6.2）的记号，$L^{(2)}_{\bullet\bullet}$ 仅在双次数 $(0,0)$ 处非零，等于 $A$；由（0，11.7.5）即得结论。更一般的结果见（6.7.7）。但若把（6.7.4.1）左边的 $\OO_Y$ 换成任意拟凝聚 $\OO_Y$-模 $\F$，一般不再与 $\cH^{-n}(f,\PP_\bullet\otimes_{\OO_Y}\F)$ 同构，尽管上面计算中的 $L^{(1)}\otimes_AL^{(2)}$ 仍可等同于 $C^\bullet(\mathfrak U,\PP_\bullet\otimes_Y\F)$。

另一方面，有双 $\delta$-函子的典范同构
\begin{equation}\tag{6.7.4.2}
 \lTor_*^S(1_{X^{(1)}},1_{X^{(2)}};
          \mathscr P^{(1)}_\bullet,\mathscr P^{(2)}_\bullet)
 \simeq\lTor_*^S(\mathscr P^{(1)}_\bullet,\mathscr P^{(2)}_\bullet).
\end{equation}
仍由（6.7.1）、（6.7.2）归结为 $S,X^{(1)},X^{(2)}$ 仿射。计算左边时，可取只含 $X^{(i)}$ 自身的覆盖 $\mathfrak U^{(i)}$；此时 $L^{(i)}$ 退化为 $\Gamma(X^{(i)},\mathscr P^{(i)}_\bullet)$，视为第一次数非零项皆为零的双复形。两边相等便由（6.4.1.1）及（6.3.1）得到。

\egas{6.7.5}{命题}
设 $u:\mathscr P^{(1)}_\bullet\to\mathscr Q^{(1)}_\bullet$ 是下有界拟凝聚 $\OO_{X^{(1)}}$-模复形的同态，且诱导同态
$\HH_*(u):\HH_*(\mathscr P^{(1)}_\bullet)\to\HH_*(\mathscr Q^{(1)}_\bullet)$ 为同构。则由 $u$ 诱导的同态
\[
 {}^{(t)}\spE(f_1,f_2;\mathscr P^{(1)}_\bullet,\mathscr P^{(2)}_\bullet)
 \longrightarrow{}^{(t)}\spE(f_1,f_2;\mathscr Q^{(1)}_\bullet,\mathscr P^{(2)}_\bullet)
\]
在 $t=a,b,c$ 时为同构。
\begin{proof}
关于 $(c)$ 的断言，由该谱序列双正则及假设保证 $E^2$ 项上的同态为同构得到（0，11.1.5）。这已说明共同终值之间的同态为同构。应用（6.7.4.1）和（6.7.4.2），先得到由 $u$ 诱导的同态
\[
\begin{gathered}
 \cH^{-n}(f_1,\mathscr P^{(1)}_\bullet)\to
 \cH^{-n}(f_1,\mathscr Q^{(1)}_\bullet),\\
 \lTor_*^S(\mathscr P^{(1)}_\bullet,\mathscr P^{(2)}_\bullet)\to
 \lTor_*^S(\mathscr Q^{(1)}_\bullet,\mathscr P^{(2)}_\bullet)
\end{gathered}
\]
均为同构。再由 $(a),(b)$ 双正则（6.7.3），且相应同态在 $E^2$ 项上双射，得到这两个谱序列的断言（0，11.1.5）。
\end{proof}
此外，若 $u$ 是同伦等价，则六个谱序列全部得到典范同构。事实上，若 $S$ 和 $Y^{(i)}$ 仿射，$u$ 诱导双复形的同伦等价
\[
 C^\bullet(\mathfrak U^{(1)},\mathscr P^{(1)}_\bullet)
 \longrightarrow C^\bullet(\mathfrak U^{(1)},\mathscr Q^{(1)}_\bullet),
\]
由一般超同调理论得到结论（0，11.3.2）。一般情形由粘合得到，并使用如下事实：复形的同伦等价诱导它们的投射 Cartan--Eilenberg 解消之间的同伦等价（M，XVII，1.2）。

\egas{6.7.6}{命题}
设 $\mathscr P^{(1)}_\bullet$ 或 $\mathscr P^{(2)}_\bullet$ 由 $S$-平坦模组成，且两个复形均下有界。则有双 $\delta$-函子的典范同构
\begin{equation}\tag{6.7.6.1}
 \lTor_n^S(f_1,f_2;\mathscr P^{(1)}_\bullet,\mathscr P^{(2)}_\bullet)
 \simeq\cH^{-n}(f_1\times_Sf_2,
                 \mathscr P^{(1)}_\bullet\otimes_S\mathscr P^{(2)}_\bullet).
\end{equation}
\begin{proof}
先设 $S,Y^{(1)},Y^{(2)}$ 仿射，处于（6.6.2）的情形并沿用其记号。例如，设 $\mathscr P^{(1)}_\bullet$ 由 $S$-平坦模组成。按（6.6.6）计算超 Tor：它是 $L^{(1)}_{\bullet\bullet}\otimes_AM^{(2)}_{\bullet\bullet\bullet}$ 的同调，其中 $M^{(2)}$ 是（0，11.7.1）意义下 $L^{(2)}$ 的投射 Cartan--Eilenberg 解消。另一方面，由假设，各 $L^{(1)}_{pq}$ 在 $A$ 上平坦（1.4.15.1）；因此（0，11.7.5）给出典范同构
\begin{equation}\tag{6.7.6.2}
 \bTor_*^S(\mathfrak U^{(1)},\mathfrak U^{(2)};
             \mathscr P^{(1)}_\bullet,\mathscr P^{(2)}_\bullet)
 \simeq H_*(L^{(1)}_{\bullet\bullet}\otimes_AL^{(2)}_{\bullet\bullet}).
\end{equation}
又有从 $L^{(1)}\otimes_AL^{(2)}$ 到 $C^\bullet(\mathfrak U,\mathscr Q_\bullet)$ 的自然双复形同态，其中 $\mathfrak U$ 是 $Z=X^{(1)}\times_SX^{(2)}$ 由 $U^{(1)}_\alpha\times_SU^{(2)}_\beta$ 组成的覆盖，且 $\mathscr Q_\bullet=\mathscr P^{(1)}_\bullet\otimes_S\mathscr P^{(2)}_\bullet$ 按总次数视为单复形。此同态的定义，实际上已在（6.6.4）计算 $(b')$ 谱序列的 $q=0$ 情形时给出：沿用当时记号，先有从 $N^{\bullet\bullet}$ 到 $C^{\bullet\bullet}(\Phi^{(1)},\Phi^{(2)};\mathscr S^\bullet)$ 的自然同态（0，11.8.5），再有从后者到 $P^{\bullet\bullet}(\Phi^{(1)},\Phi^{(2)};\mathscr S^\bullet)$ 的自然同态（0，11.8.6），最后从后者的上链复形映入交错上链子复形（0，11.8.7）。如（6.6.4）所见，这个同态在同调上给出同构；与（6.7.6.2）复合，得到
\begin{equation}\tag{6.7.6.3}
 \bTor_n^S(\mathfrak U^{(1)},\mathfrak U^{(2)};
             \mathscr P^{(1)}_\bullet,\mathscr P^{(2)}_\bullet)
 \simeq\Hh^{-n}(\mathfrak U,
             \mathscr P^{(1)}_\bullet\otimes_S\mathscr P^{(2)}_\bullet).
\end{equation}
随后须证明该同构与所选开覆盖无关；（6.7.6.3）右边由（6.2.2）典范同构于
$\Hh^{-n}(X^{(1)}\times_SX^{(2)},\mathscr P^{(1)}_\bullet\otimes_S\mathscr P^{(2)}_\bullet)$。利用（6.6.7），注意有下图在同伦意义下交换：
\[
\begin{tikzcd}[column sep=small,row sep=large]
 C^\bullet(\mathfrak U^{(1)},\mathscr P^{(1)}_\bullet)\otimes_A
 C^\bullet(\mathfrak U^{(2)},\mathscr P^{(2)}_\bullet)
 \arrow[r]\arrow[d]&C^\bullet(\mathfrak U,\mathscr Q_\bullet)\arrow[d]\\
 C^\bullet(\mathfrak V^{(1)},\mathscr P^{(1)}_\bullet)\otimes_A
 C^\bullet(\mathfrak V^{(2)},\mathscr P^{(2)}_\bullet)
 \arrow[r]&C^\bullet(\mathfrak V,\mathscr Q_\bullet).
\end{tikzcd}
\]
其中横箭头为刚才定义的同态。最后，如（6.7.1）、（6.7.2）一样，通过粘合过渡到一般情形，没有困难；细节留给读者。
\end{proof}

\egas{6.7.7}{命题}
设两个复形下有界，并且所有 $\cH^{-n}(f_1,\mathscr P^{(1)}_\bullet)$，或所有 $\cH^{-n}(f_2,\mathscr P^{(2)}_\bullet)$，均在 $S$ 上平坦。则有双 $\delta$-函子的典范同构（$n\in\mathbf Z$）
\begin{equation}\tag{6.7.7.1}
\begin{aligned}
 &\lTor_n^S(f_1,f_2;\mathscr P^{(1)}_\bullet,\mathscr P^{(2)}_\bullet)\\
 &\qquad\simeq\bigoplus_{q_1+q_2=n}
 \cH^{-q_1}(f_1,\mathscr P^{(1)}_\bullet)\otimes_S
 \cH^{-q_2}(f_2,\mathscr P^{(2)}_\bullet).
\end{aligned}
\end{equation}
由（6.5.8），（6.7.3）的谱序列 $(a)$ 退化；它双正则，所以（0，11.1.6）立即给出结论。

\egas{6.7.8}{定理}
假设：一、$\mathscr P^{(1)}_\bullet,\mathscr P^{(2)}_\bullet$ 下有界；二、其中一个复形由 $S$-平坦模组成；三、所有 $\cH^{-n}(f_1,\mathscr P^{(1)}_\bullet)$ 或所有 $\cH^{-n}(f_2,\mathscr P^{(2)}_\bullet)$ 在 $S$ 上平坦。则有典范的双 $\delta$-函子同构（$n\in\mathbf Z$）
\begin{equation}\tag{6.7.8.1}
\begin{aligned}
 &\cH^n(f_1\times_Sf_2,
               \mathscr P^{(1)}_\bullet\otimes_S\mathscr P^{(2)}_\bullet)\\
 &\qquad\simeq\bigoplus_{n_1+n_2=n}
     \cH^{n_1}(f_1,\mathscr P^{(1)}_\bullet)\otimes_S
     \cH^{n_2}(f_2,\mathscr P^{(2)}_\bullet)
\end{aligned}
\end{equation}
（“Künneth公式”）。

这由（6.7.6）和（6.7.7）得到。当 $S,Y^{(1)},Y^{(2)}$ 仿射时，（6.7.8.1）的逆同构，按（6.7.6）的记号，由双复形同态
\[
 C^\bullet(\mathfrak U^{(1)},\mathscr P^{(1)}_\bullet)\otimes_A
 C^\bullet(\mathfrak U^{(2)},\mathscr P^{(2)}_\bullet)
 \longrightarrow C^\bullet(\mathfrak U,\mathscr Q_\bullet)
\]
经（G，I，2.7）的过程导出；这由（G，I，5.5）可见。

\egas{6.7.9}{命题}
假设以下三个条件成立：一、$S,Y^{(1)},Y^{(2)}$ 局部诺特，$f_1,f_2$ 固有，且 $Y^{(1)}$ 或 $Y^{(2)}$ 在 $S$ 上有限型；二、$\mathscr P^{(1)}_\bullet,\mathscr P^{(2)}_\bullet$ 下有界；三、每个 $n\in\mathbf Z$ 对应的 $\HH_n(\mathscr P^{(i)}_\bullet)$ 均凝聚（$i=1,2$）。则
$\lTor_n^S(f_1,f_2;\mathscr P^{(1)}_\bullet,\mathscr P^{(2)}_\bullet)$ 是凝聚 $\OO_Y$-模，其中 $Y=Y^{(1)}\times_SY^{(2)}$。

由（6.5.13），局部超 Tor 是凝聚 $\OO_X$-模，其中 $X=X^{(1)}\times_SX^{(2)}$ 局部诺特：因为某个 $X^{(i)}$ 由假设在 $S$ 上有限型（I，6.3.4、6.3.8）。$Y$ 局部诺特，$f_1\times_Sf_2$ 固有（II，5.4.2），故由（6.2.5），（6.7.3）的 ${}^{(b)}\spE^2_{pq}$ 为凝聚 $\OO_Y$-模。由条件二，各谱序列都双正则，于是应用（0，11.1.8）即可。

\egap{6.7.10}
现在取两个 $S$-预概形 $Y'^{(i)}$、$S$-态射 $v_i:Y'^{(i)}\to Y^{(i)}$（$i=1,2$），以及其乘积
$v=v_1\times_Sv_2:Y'\to Y$，其中 $Y'=Y'^{(1)}\times_SY'^{(2)}$。对每个 $i$，另取 $S$-预概形 $X'^{(i)}$ 和两个 $S$-态射 $u_i:X'^{(i)}\to X^{(i)}$、$f'_i:X'^{(i)}\to Y'^{(i)}$，使下图交换：
\begin{equation}\tag{6.7.10.1}
\begin{tikzcd}[column sep=large,row sep=large]
 X'^{(i)}\arrow[r,"u_i"]\arrow[d,"f'_i"']&X^{(i)}\arrow[d,"f_i"]\\
 Y'^{(i)}\arrow[r,"v_i"']&Y^{(i)}.
\end{tikzcd}
\end{equation}
设 $f'_i$ 分离且拟紧。于是有谱函子的典范 $\OO_{Y'}$-同态
\begin{equation}\tag{6.7.10.2}
\begin{aligned}
 &v^*\bigl({}^{(t)}\spE(f_1,f_2;
               \mathscr P^{(1)}_\bullet,\mathscr P^{(2)}_\bullet)\bigr)\\
 &\qquad\longrightarrow{}^{(t)}\spE(f'_1,f'_2;
               u_1^*\mathscr P^{(1)}_\bullet,u_2^*\mathscr P^{(2)}_\bullet),
\end{aligned}
\end{equation}
其中 $t=a,a',b,b',c,d$。

为定义它们，先设 $S=\Spec A$、$Y^{(i)}=\Spec B_i$、$Y'^{(i)}=\Spec B'_i$ 仿射；于是 $X^{(i)},X'^{(i)}$ 是拟紧概形。计算未加撇的谱序列时，像（6.6.1）一样取 $X^{(i)}$ 的有限仿射开覆盖 $\mathfrak U^{(i)}$；计算加撇的谱序列时，取 $X'^{(i)}$ 的有限仿射开覆盖 $\mathfrak U'^{(i)}$，分别细于 $u_i^{-1}(\mathfrak U^{(i)})$。原文说，双复形
$L^{(i)}=C^\bullet(\mathfrak U^{(i)},\mathscr P^{(i)}_\bullet)$ 可典范地看作
$C^\bullet(u_i^{-1}(\mathfrak U^{(i)}),u_i^*\mathscr P^{(i)}_\bullet)$ 的“子双复形”（$0_{\mathrm I}$，4.4.3.2）。
\cmt{C3-G01}{“子双复形”一般不成立：取 $\Spec k\to\Spec k[t]$ 为 $t=0$ 的闭浸入，结构层在单开集覆盖上的映射是 $k[t]\to k$，不单射。这里确有典范同态；后续构造只需要同态，不需要单射。正文保留原文说法并在此指出错误，不否定后续函子性结论。}
选取从 $\mathfrak U'^{(i)}$ 到 $u_i^{-1}(\mathfrak U^{(i)})$ 的单纯映射（G，II，5.7），便定义同态
\[
 C^\bullet(u_i^{-1}(\mathfrak U^{(i)}),u_i^*\mathscr P^{(i)}_\bullet)
 \longrightarrow C^\bullet(\mathfrak U'^{(i)},u_i^*\mathscr P^{(i)}_\bullet),
\]
由复合得到 $L^{(i)}\to L'^{(i)}$。改变单纯映射，只将此同态替换为与其同伦的同态（G，II，5.7.1）；所以得到良好定义的谱函子同态
\begin{equation}\tag{6.7.10.3}
\begin{aligned}
 &{}^{(t)}E(\mathfrak U^{(1)},\mathfrak U^{(2)};
              \mathscr P^{(1)}_\bullet,\mathscr P^{(2)}_\bullet)\\
 &\qquad\longrightarrow{}^{(t)}E(\mathfrak U'^{(1)},\mathfrak U'^{(2)};
              u_1^*\mathscr P^{(1)}_\bullet,u_2^*\mathscr P^{(2)}_\bullet).
\end{aligned}
\end{equation}
若 $\mathfrak V^{(i)}$ 是 $X^{(i)}$ 的有限仿射开覆盖，细于 $\mathfrak U^{(i)}$，而 $\mathfrak V'^{(i)}$ 是 $X'^{(i)}$ 的有限仿射开覆盖，同时细于 $u_i^{-1}(\mathfrak V^{(i)})$ 和 $\mathfrak U'^{(i)}$，则容易验证下图交换：
\[
\begin{tikzcd}[column sep=large,row sep=large]
 C^\bullet(\mathfrak U^{(i)},\mathscr P^{(i)}_\bullet)\arrow[r]\arrow[d]&
 C^\bullet(\mathfrak V^{(i)},\mathscr P^{(i)}_\bullet)\arrow[d]\\
 C^\bullet(\mathfrak U'^{(i)},u_i^*\mathscr P^{(i)}_\bullet)\arrow[r]&
 C^\bullet(\mathfrak V'^{(i)},u_i^*\mathscr P^{(i)}_\bullet).
\end{tikzcd}
\]
所以（6.7.10.3）本质上不依赖这些覆盖。实际定义了 $A$-模同态
\begin{equation}\tag{6.7.10.4}
\begin{aligned}
 &{}^{(t)}E(X^{(1)},X^{(2)};
                  \mathscr P^{(1)}_\bullet,\mathscr P^{(2)}_\bullet)\\
 &\qquad\longrightarrow{}^{(t)}E(X'^{(1)},X'^{(2)};
                  u_1^*\mathscr P^{(1)}_\bullet,u_2^*\mathscr P^{(2)}_\bullet).
\end{aligned}
\end{equation}
由拉回的定义和（6.7.10.1）的交换性，这同时是 $(B_1\otimes_AB_2)$-模同态。右边由 $(B'_1\otimes_AB'_2)$-模组成，故它典范地诱导
\begin{equation}\tag{6.7.10.5}
\begin{aligned}
 &{}^{(t)}E(X^{(1)},X^{(2)};
           \mathscr P^{(1)}_\bullet,\mathscr P^{(2)}_\bullet)
           \otimes_{B_1\otimes_AB_2}(B'_1\otimes_AB'_2)\\
 &\qquad\longrightarrow{}^{(t)}E(X'^{(1)},X'^{(2)};
           u_1^*\mathscr P^{(1)}_\bullet,u_2^*\mathscr P^{(2)}_\bullet).
\end{aligned}
\end{equation}
由（I，1.6.5），这正是当前特例中所求的（6.7.10.2）。

最后，按（6.7.1）、（6.7.2）的粘合过渡到一般情形。第二次过渡是直接的；对于第一次，像（6.7.1）那样，取 $g_i\in B_i$、其像 $g'_i\in B'_i$，以及 $B=B_1\otimes_AB_2$ 中的 $g=g_1\otimes g_2$ 和它在 $B'=B'_1\otimes_AB'_2$ 中的像 $g'=g'_1\otimes g'_2$；一切归结为使用典范同构
\[
 (M\otimes_BB')_{g'}\simeq M_g\otimes_{B_g}B'_{g'}
\]
（$0_{\mathrm I}$，1.5.4）。细节留给读者。

\egap{6.7.11}
上述对两个 $S$-态射 $X^{(i)}\to Y^{(i)}$ 及两个下有界拟凝聚模复形所发展的整体超 Tor 理论，立即推广到以下一般情形：给定预概形 $S$，有限个 $S$-预概形 $Y^{(i)}$（$i\in I$），有限个分离拟紧 $S$-态射 $f_i:X^{(i)}\to Y^{(i)}$，以及每个 $X^{(i)}$ 上的下有界拟凝聚 $\OO_{X^{(i)}}$-模复形 $\mathscr P^{(i)}_\bullet$。若 $Y$ 是各 $Y^{(i)}$ 在 $S$ 上的乘积，则对每个 $n\in\mathbf Z$，可定义拟凝聚 $\OO_Y$-模
\[
 \lTor_n^S\bigl((f_i)_{i\in I};(\mathscr P^{(i)}_\bullet)_{i\in I}\bigr).
\]
这些模对每个复形变量都构成协变 $\delta$-函子，并是六个谱函子
\[{}^{(t)}\spE((f_i)_{i\in I};(\mathscr P^{(i)}_\bullet)_{i\in I})\]的共同终值。将 $I=\{1,2\}$ 时的定义与论证重复到此一般情形，留给读者。只须注意，当 $I$ 只含一个元素时，便得到（6.2.7）的超上同调 $\cH^\bullet(f,\PP_\bullet)$，如（6.7.4）已观察到的。若 $I=\{1,\ldots,m\}$，记作
\[
 \lTor_n^S(f_1,\ldots,f_m;
                 \mathscr P^{(1)}_\bullet,\ldots,\mathscr P^{(m)}_\bullet).
\]

\egas{6.7.12}{命题}
采用（6.7.11）的记号。设 $J\subset I$，且对每个 $i\in I\setminus J$，有 $X^{(i)}=Y^{(i)}=S$、$f_i=1_S$，而 $\mathscr P^{(i)}_\bullet$ 仅在第零项等于 $\OO_S$。则有 $\delta$-函子的典范同构
\begin{equation}\tag{6.7.12.1}
 \lTor_n^S\bigl((f_i)_{i\in I};(\mathscr P^{(i)}_\bullet)_{i\in I}\bigr)
 \simeq\lTor_n^S\bigl((f_i)_{i\in J};(\mathscr P^{(i)}_\bullet)_{i\in J}\bigr).
\end{equation}
只须在 $S$ 和各 $Y^{(i)}$ 仿射时定义同构，随后照常粘合。对 $i\in I\setminus J$，取仅含 $S$ 的覆盖 $\mathfrak U^{(i)}$；于是 $L^{(i)}=C^\bullet(\mathfrak U^{(i)},\mathscr P^{(i)}_\bullet)$ 仅在双次数 $(0,0)$ 处非零，等于 $\Gamma(S,\OO_S)=A(S)$。此时（6.7.12.1）显然成立。

\egas{6.7.13}{注记}
采用（6.7.3）的记号，考虑典范 $S$-同构
$Y^{(1)}\times_SY^{(2)}\to Y^{(2)}\times_SY^{(1)}$（I，3.3.5）。它把
$\lTor_*^S(f_1,f_2;\mathscr P^{(1)}_\bullet,\mathscr P^{(2)}_\bullet)$ 变为
$\lTor_*^S(f_2,f_1;\mathscr P^{(2)}_\bullet,\mathscr P^{(1)}_\bullet)$。问题是局部的，故可归结为（6.6.2）的情形。若 $M^{(i)}$ 为 $L^{(i)}$ 的投射 Cartan--Eilenberg 解消，则该同构把 $M^{(1)}\otimes_AM^{(2)}$ 变为 $M^{(2)}\otimes_AM^{(1)}$；取这些三重复形所对应单复形的同调，即得断言。

% END continuations/iii2_67.tex


% BEGIN continuations/iii2_68.tex
\subsection{整体超 Tor 的结合性谱序列}
\source{\texttt{EGAIII.2.pdf}，PDF第31--32页；原书第32--33页}
\egap{6.8.1}
采用（6.7.11）的假设与记号，特别地，各 $\mathscr P^{(i)}_\bullet$ 下有界。设给定指标集 $I$ 的一个分割 $(I_j)_{j\in J}$。要在整体超 Tor
\[
 \lTor_*^S\bigl((f_i)_{i\in I};(\mathscr P^{(i)}_\bullet)_{i\in I}\bigr)
\]
与各“部分”超 Tor
\[
 \mathscr A^{(j)}_\bullet
 =\lTor_*^S\bigl((f_i)_{i\in I_j};(\mathscr P^{(i)}_\bullet)_{i\in I_j}\bigr)
\]
之间给出一种“结合性”关系。为简化书写，只考虑 $I=\{1,\ldots,m\}$，且分割由两个区间 $\{1,\ldots,r\}$ 与 $\{r+1,\ldots,m\}$ 组成的情形。

\egas{6.8.2}{命题}
存在典范的双正则谱函子，称为“结合性谱函子”，记作
\[
 {}^{(r)}\spE(f_1,\ldots,f_m;
                  \mathscr P^{(1)}_\bullet,\ldots,\mathscr P^{(m)}_\bullet)
\]
（也可在复形记号中略去下标点），其终值为
\[\lTor_*^S(f_1,\ldots,f_m;\mathscr P^{(1)}_\bullet,\ldots,\mathscr P^{(m)}_\bullet),\]且 $E^2$ 项为
\[
\begin{aligned}
 {}^{(r)}\spE^2_{pq}
 =\bigoplus_{q_1+q_2=q}\lTor_p^S\bigl(&
   \lTor_{q_1}^S(f_1,\ldots,f_r;
                      \mathscr P^{(1)}_\bullet,\ldots,\mathscr P^{(r)}_\bullet),\\
 &\lTor_{q_2}^S(f_{r+1},\ldots,f_m;
                      \mathscr P^{(r+1)}_\bullet,\ldots,\mathscr P^{(m)}_\bullet)\bigr).
\end{aligned}
\]
在此，已将 $Y$ 典范地等同于 $Z^{(1)}\times_SZ^{(2)}$，其中
\[
 Z^{(1)}=Y^{(1)}\times_S\cdots\times_SY^{(r)},\qquad
 Z^{(2)}=Y^{(r+1)}\times_S\cdots\times_SY^{(m)}.
\]
我们只处理 $S$ 和各 $Y^{(i)}$ 仿射的情形；用（6.7.1）、（6.7.2）的方法从这一特例过渡到一般情形，没有困难，论证细节留给读者。因此证明下面的推论即可。

\egas{6.8.3}{推论}
设 $A$ 为环，$S=\Spec A$，$X^{(i)}$（$1\leq i\leq m$）为拟紧 $S$-概形；对每个 $i$，设 $\mathscr P^{(i)}_\bullet$ 为下有界拟凝聚 $\OO_{X^{(i)}}$-模复形。存在一个典范双正则谱函子，其终值为
\[
 \bTor_*^S(X^{(1)},\ldots,X^{(m)};
                  \mathscr P^{(1)}_\bullet,\ldots,\mathscr P^{(m)}_\bullet),
\]
且 $E^2$ 项为
\[
\begin{aligned}
 {}^{(r)}E^2_{pq}=\bigoplus_{q_1+q_2=q}\Tor_p^A\bigl(&
   \bTor_{q_1}^S(X^{(1)},\ldots,X^{(r)};
                    \mathscr P^{(1)}_\bullet,\ldots,\mathscr P^{(r)}_\bullet),\\
 &\bTor_{q_2}^S(X^{(r+1)},\ldots,X^{(m)};
                    \mathscr P^{(r+1)}_\bullet,\ldots,\mathscr P^{(m)}_\bullet)\bigr).
\end{aligned}
\]
\begin{proof}
按（6.6.2）的定义，计算所讨论的超 Tor 时，对每个 $i$ 取 $X^{(i)}$ 的有限仿射开覆盖 $\mathfrak U^{(i)}$，考虑双复形
$L^{(i)}_{\bullet\bullet}=C^\bullet(\mathfrak U^{(i)},\mathscr P^{(i)}_\bullet)$，并取它的投射 Cartan--Eilenberg 解消 $M^{(i)}_{\bullet\bullet\bullet}$，这里解消按（0，11.7.1）理解。形成这些三重复形的张量积
\[
 M_{\bullet\bullet\bullet}=\bigotimes_{i=1}^mM^{(i)}_{\bullet\bullet\bullet},
\]
然后取其同调。把 $M$ 看作单复形 $N_\bullet$，它是下列两个单复形的张量积：
\[
 N'_\bullet=\bigotimes_{i=1}^rM^{(i)}_{\bullet\bullet\bullet},\qquad
 N''_\bullet=\bigotimes_{i=r+1}^mM^{(i)}_{\bullet\bullet\bullet};
\]
其中各单复形按相应总次数之和分次。$N'_\bullet,N''_\bullet$ 中各 $A$-模均投射，故（6.5.9）给出
$H_*(M)=\Tor_*^A(N'_\bullet,N''_\bullet)$。所求谱序列就是将（6.3.2.2）用于 $N'_\bullet,N''_\bullet$；上述构造用于两个部分乘积时，给出这两个复形的同调模的所需解释。最后，由于各 $\mathscr P^{(i)}_\bullet$ 下有界，各 $M^{(i)}$ 也下有界，因而 $N'_\bullet,N''_\bullet$ 下有界；所需的正则性由（6.3.2）得到。
\end{proof}

% END continuations/iii2_68.tex


% BEGIN continuations/iii2_69.tex
\subsection{整体超 Tor 中的基变换谱序列}
\source{\texttt{EGAIII.2.pdf}，PDF第33--38页；原书第34--39页}
\egap{6.9.1}
仍采用（6.7.11）的假设和记号，特别地，各 $\mathscr P^{(i)}_\bullet$ 下有界。考虑预概形态射 $g:S'\to S$，并置
\[
 Y'^{(i)}=Y^{(i)}_{(S')},\quad X'^{(i)}=X^{(i)}_{(S')},\quad
 \mathscr P'^{(i)}_\bullet=\mathscr P^{(i)}_\bullet\otimes_{\OO_S}\OO_{S'}.
\]
于是 $\mathscr P'^{(i)}_\bullet$ 是拟凝聚 $\OO_{X'^{(i)}}$-模复形。令 $f'_i=(f_i)_{(S')}:X'^{(i)}\to Y'^{(i)}$；它分离且拟紧（I，5.5.1、6.6.4）。要研究拟凝聚 $\OO_{Y'}$-模
\[
 \lTor_n^{S'}\bigl((f'_i)_{i\in I};(\mathscr P'^{(i)}_\bullet)_{i\in I}\bigr)
 \quad\hbox{与}\quad
 \lTor_n^S\bigl((f_i)_{i\in I};(\mathscr P^{(i)}_\bullet)_{i\in I}\bigr)
                 \otimes_{\OO_S}\OO_{S'}
\]
之间的关系；这里 $Y'=Y\times_SS'=Y_{(S')}$。下面是特别简单的一种情形；当 $I$ 只有一个元素且复形只含一个模时，它还原为（1.4.15）。

\egas{6.9.2}{命题}
若 $g:S'\to S$ 平坦，则有关于各 $\mathscr P^{(i)}_\bullet$ 的 $\delta$-函子的典范同构
\begin{equation}\tag{6.9.2.1}
\begin{aligned}
 &\lTor_*^S\bigl((f_i)_{i\in I};(\mathscr P^{(i)}_\bullet)_{i\in I}\bigr)
           \otimes_{\OO_S}\OO_{S'}\\
 &\qquad\simeq\lTor_*^{S'}\bigl((f'_i)_{i\in I};(\mathscr P'^{(i)}_\bullet)_{i\in I}\bigr).
\end{aligned}
\end{equation}
\begin{proof}
仍可只考虑 $S,S'$ 和各 $Y^{(i)}$ 仿射的情形，随后按（6.7.1）、（6.7.2）的方法粘合。设 $S=\Spec A$、$S'=\Spec A'$，并对每个 $i$ 取 $X^{(i)}$ 的仿射开覆盖 $\mathfrak U^{(i)}$。若 $u_i:X'^{(i)}\to X^{(i)}$ 是典范投影，则 $u_i^{-1}(\mathfrak U^{(i)})$ 是 $X'^{(i)}$ 的仿射开覆盖（II，1.5.5），记为 $\mathfrak U'^{(i)}$。显然，
\[
 C^\bullet(\mathfrak U'^{(i)},\mathscr P'^{(i)}_\bullet)
 =C^\bullet(\mathfrak U^{(i)},\mathscr P^{(i)}_\bullet)\otimes_AA'.
\]
若 $M^{(i)}$ 是 $L^{(i)}=C^\bullet(\mathfrak U^{(i)},\mathscr P^{(i)}_\bullet)$ 按（0，11.7.1）取得的、由 $A$-模组成的投射 Cartan--Eilenberg 解消，则由 $A'$ 在 $A$ 上平坦，$M^{(i)}\otimes_AA'$ 是 $L^{(i)}\otimes_AA'$ 的同类投射解消，且
\[
 H_*\left(\bigotimes_{i=1}^m(M^{(i)}\otimes_AA')\right)
 =H_*\left(\bigotimes_{i=1}^mM^{(i)}\right)\otimes_AA'.
\]
这立即给出（6.9.2.1）。
\end{proof}
当 $I=\{1\}$ 时，（6.9.2.1）也可直接从（6.7.7）得到：取 $Y_2=X_2=S'$、$f_2=1_{S'}$，并将 $\mathscr P^{(2)}_\bullet$ 取为仅第零项等于 $\OO_{S'}$ 的复形；此时 $\cH^n(1_{S'},\OO_{S'})$ 在 $n\ne0$ 时为零，在 $n=0$ 时等于 $\OO_{S'}$（6.2.1）。

在一般情形中，以一个下有界拟凝聚 $\OO_{S'}$-模复形 $\mathscr Q'_\bullet$ 代替 $\OO_{S'}$。为简化书写，取 $I=\{1,\ldots,m\}$；于是可以考虑 $\delta$-函子
\[
 \lTor_*^S(f_1,\ldots,f_m,1_{S'};
         \mathscr P^{(1)}_\bullet,\ldots,\mathscr P^{(m)}_\bullet,\mathscr Q'_\bullet).
\]

\egas{6.9.3}{命题}
存在三个典范双正则谱函子，记为
\[
 {}^{(t)}\spE(f_1,\ldots,f_m;
       \mathscr P^{(1)}_\bullet,\ldots,\mathscr P^{(m)}_\bullet,\mathscr Q'_\bullet),
 \qquad t=e,f,f',
\]
其共同终值为上述 $\lTor_*^S$，且 $E^2$ 项分别为
\begin{align*}
 {}^{(e)}\spE^2_{pq}
 &=\bigoplus_{q'+q''=q}\lTor_p^S\bigl(
   \lTor_{q'}^S(f_1,\ldots,f_m;
       \mathscr P^{(1)}_\bullet,\ldots,\mathscr P^{(m)}_\bullet),
       \HH_{q''}(\mathscr Q'_\bullet)\bigr),\\
 {}^{(f)}\spE^2_{pq}
 &=\bigoplus_{q_1+\cdots+q_{m+1}=q}\lTor_p^{S'}\bigl(
   \lTor_{q_1}^S(f_1,1_{S'};\mathscr P^{(1)}_\bullet,\OO_{S'}),\ldots,\\[-.6em]
 &\hspace{7em}\lTor_{q_m}^S(f_m,1_{S'};\mathscr P^{(m)}_\bullet,\OO_{S'}),
      \HH_{q_{m+1}}(\mathscr Q'_\bullet)\bigr),\\
 {}^{(f')}\spE^2_{pq}
 &=\bigoplus_{q_1+\cdots+q_m=q}\lTor_p^{S'}\bigl(
  f'_1,\ldots,f'_m,1_{S'};
  \termTor_{q_1}^{S}(\mathscr P^{(1)}_\bullet,\OO_{S'}),\ldots,\\[-.6em]
 &\hspace{10em}\termTor_{q_m}^{S}(\mathscr P^{(m)}_\bullet,\OO_{S'}),
  \mathscr Q'_\bullet\bigr).
\end{align*}
这里 $\mathrm{term}$ 沿用（6.7.3）译校注中的记号约定，表示逐项构成的复形。

谱序列 $(e)$ 就是（6.8.2）在 $r=m$ 时的结合性谱序列。为定义另外两个谱序列，仍只处理 $S,S'$ 和各 $Y^{(i)}$ 仿射的情形；通过（6.7.1）、（6.7.2）的方法过渡到一般情形，留给读者。因此证明下述推论。

\egas{6.9.4}{推论}
设 $A$ 为环，$A'$ 为 $A$-代数，$S=\Spec A$、$S'=\Spec A'$；$X^{(i)}$（$1\leq i\leq m$）为拟紧 $S$-概形，且置 $X'^{(i)}=X^{(i)}_{(S')}$，它是拟紧 $S'$-概形。对每个 $i$，取拟凝聚 $\OO_{X^{(i)}}$-模复形 $\mathscr P^{(i)}_\bullet$；再取 $A'$-模复形 $Q'_\bullet$，所有这些复形均下有界。存在关于各复形的三个双正则谱函子，共同终值为
\[
 \bTor_*^S(X^{(1)},\ldots,X^{(m)},S';
         \mathscr P^{(1)}_\bullet,\ldots,\mathscr P^{(m)}_\bullet,\widetilde{Q'_\bullet}),
\]
且 $E^2$ 项分别为
\begin{align*}
 {}^{(e)}E^2_{pq}
 &=\bigoplus_{q'+q''=q}\Tor_p^A\bigl(
     \bTor_{q'}^S(X^{(1)},\ldots,X^{(m)};
          \mathscr P^{(1)}_\bullet,\ldots,\mathscr P^{(m)}_\bullet),H_{q''}(Q'_\bullet)\bigr),\\
 {}^{(f)}E^2_{pq}
 &=\bigoplus_{q_1+\cdots+q_{m+1}=q}\Tor_p^{A'}\bigl(
      \bTor_{q_1}^S(X^{(1)},S';\mathscr P^{(1)}_\bullet,\OO_{S'}),\ldots,\\[-.6em]
 &\hspace{8em}\bTor_{q_m}^S(X^{(m)},S';\mathscr P^{(m)}_\bullet,\OO_{S'}),
        H_{q_{m+1}}(Q'_\bullet)\bigr),\\
 {}^{(f')}E^2_{pq}
 &=\bigoplus_{q_1+\cdots+q_m=q}\bTor_p^{S'}\bigl(
   X'^{(1)},\ldots,X'^{(m)},S';
   \termTor_{q_1}^S(\mathscr P^{(1)}_\bullet,\OO_{S'}),\ldots,\\[-.6em]
 &\hspace{10em}\termTor_{q_m}^S(\mathscr P^{(m)}_\bullet,\OO_{S'}),
       \widetilde{Q'_\bullet}\bigr).
\end{align*}
\begin{proof}
第一个谱函子已在（6.8.3）处理，此处仅为备查而重列。为定义另外两个，对每个 $i$ 取 $X^{(i)}$ 的有限仿射开覆盖 $\mathfrak U^{(i)}$；若 $u_i:X'^{(i)}\to X^{(i)}$ 为典范投影，则相应有限仿射开覆盖为 $\mathfrak U'^{(i)}=u_i^{-1}(\mathfrak U^{(i)})$。由（6.6.6），上述超 Tor 可如下计算：对每个 $i$ 取 $L^{(i)}=C^\bullet(\mathfrak U^{(i)},\mathscr P^{(i)}_\bullet)$ 按（0，11.7.1）意义的投射 Cartan--Eilenberg 解消 $M^{(i)}$，形成三重复形
\[
 M_{\bullet\bullet\bullet}
 =M^{(1)}_{\bullet\bullet\bullet}\otimes_A\cdots\otimes_AM^{(m)}_{\bullet\bullet\bullet}\otimes_AQ'_\bullet,
\]
其中 $Q'_\bullet$ 被视为后两个次数均为零的三重复形，然后取同调。置 $M'^{(i)}=M^{(i)}\otimes_AA'$；因为 $Q'_\bullet$ 是 $A'$-模复形，有
\[
 M=M'^{(1)}\otimes_{A'}\cdots\otimes_{A'}M'^{(m)}\otimes_{A'}Q'_\bullet.
\]
将每个 $M'^{(i)}$ 按总次数视为单复形，其各项都是投射 $A'$-模。由（6.3.7）对任意有限个复形的推广，$H_*(M)$ 也等于
$\Tor_*^{A'}(M'^{(1)},\ldots,M'^{(m)},Q'_\bullet)$。因此它是（6.5.15）中某个谱序列的终值；其 $E^2$ 项为
\[
 E^2_{pq}=\bigoplus_{q_1+\cdots+q_{m+1}=q}
 \Tor_p^{A'}\bigl(H_{q_1}(M'^{(1)}),\ldots,H_{q_m}(M'^{(m)}),
                   H_{q_{m+1}}(Q'_\bullet)\bigr).
\]
由整体超 Tor 的定义，
\[
 H_{q_i}(M'^{(i)})=H_{q_i}(M^{(i)}\otimes_AA')
 =\bTor_{q_i}^S(X^{(i)},S';\mathscr P^{(i)}_\bullet,\OO_{S'}),
\]
故得到所求 $(f)$。当 $\mathscr P^{(i)}_\bullet$ 下有界时，$M^{(i)}$ 的三个次数均下有界，这也给出所需的正则性。

另一方面，可把 $M'^{(i)}$ 看作双复形：第一次数为原三重复形的第一、第二次数之和，第二次数为原第三次数。各项为投射 $A'$-模，故由一般超同调理论，其张量积双复形的同调典范同构于超同调（0，11.6.5）。因此它是一个 $E^2$ 项为
\[
 E^2_{pq}=\bigoplus_{q_1+\cdots+q_{m+1}=q}\Tor_p^{A'}
 \bigl(H^{\mathrm{II}}_{q_1}(M'^{(1)}),\ldots,
       H^{\mathrm{II}}_{q_m}(M'^{(m)}),H^{\mathrm{II}}_{q_{m+1}}(Q'_\bullet)\bigr)
\]
的谱序列的终值。因 $Q'_\bullet$ 的第二次数只为零，$H_n^{\mathrm{II}}(Q'_\bullet)=0$（$n\ne0$），而 $H_0^{\mathrm{II}}(Q'_\bullet)=Q'_\bullet$，所以该式成为
\[
 E^2_{pq}=\bigoplus_{q_1+\cdots+q_m=q}\Tor_p^{A'}
 \bigl(H^{\mathrm{II}}_{q_1}(M'^{(1)}),\ldots,
       H^{\mathrm{II}}_{q_m}(M'^{(m)}),Q'_\bullet\bigr).
\]
又由（6.3.4），$H^{\mathrm{II}}_{q_i}(M'^{(i)})=
H^{\mathrm{III}}_{q_i}(M^{(i)}\otimes_AA')=\Tor_{q_i}^A(L^{(i)},A')$。各 $L^{(i)}_{jk}=C^j(\mathfrak U^{(i)},\mathscr P^{(i)}_k)$ 是模 $\Gamma(V,\mathscr P^{(i)}_k)$ 的直和，其中 $V$ 遍历覆盖中 $j+1$ 个开集的仿射交。令 $V'=u_i^{-1}(V)$，它在 $X'^{(i)}$ 中仿射，由（6.4.1.1），
\[
 \Gamma(V',\lTor_{q_i}^S(\mathscr P^{(i)}_k,\OO_{S'}))
 =\Tor_{q_i}^A(\Gamma(V,\mathscr P^{(i)}_k),A').
\]
从而该同调双复形可写成
\[
 C^\bullet\bigl(\mathfrak U'^{(i)},
                \termTor_{q_i}^S(\mathscr P^{(i)}_\bullet,\OO_{S'})\bigr),
\]
最终给出 $(f')$ 的 $E^2$ 项。双正则性照常由各 $M^{(i)}$ 的三个次数下有界验证。
\end{proof}

\egas{6.9.5}{注记}
如（6.7.6）所示，将各 $\mathscr P^{(i)}_\bullet$ 及 $\mathscr Q'_\bullet$ 分别替换为与其同伦等价的复形，不改变 $(e),(f),(f')$，只相差典范同构。此外，对于 $(f)$，若复形同态
$\mathscr P^{(i)}_\bullet\to\mathscr R^{(i)}_\bullet$、
$\mathscr Q'_\bullet\to\mathscr T'_\bullet$ 在同调上给出同构，则它们也给出相应 $(f)$ 谱序列的同构。证明与（6.7.6）相同，并使用该处结果以及 $(f)$ 的正则性。
\cmt{C3-T02}{本段三次引（6.7.6）；对应“同伦等价”和“同调同构”比较的一般命题实际为（6.7.5）。保留底本引号，并在此提示其可能错引；不得把同伦等价与仅同调同构当作同一个条件。}

\egas{6.9.6}{推论}
在（6.9.1）的条件下，设：一、各 $\mathscr P^{(i)}_\bullet$ 的项均在 $S$ 上平坦，且各
$\lTor_n^S(f_1,\ldots,f_m;\mathscr P^{(1)}_\bullet,\ldots,\mathscr P^{(m)}_\bullet)$ 在 $S$ 上平坦；二、各 $\mathscr P^{(i)}_\bullet$ 及 $\mathscr Q'_\bullet$ 下有界。置
$\mathscr P'^{(i)}_\bullet=\mathscr P^{(i)}_\bullet\otimes_{\OO_S}\OO_{S'}$，则有函子性的典范同构
\begin{equation}\tag{6.9.6.1}
\begin{aligned}
 &\lTor_n^{S'}(f'_1,\ldots,f'_m,1_{S'};
      \mathscr P'^{(1)}_\bullet,\ldots,\mathscr P'^{(m)}_\bullet,\mathscr Q'_\bullet)\\
 &\quad\simeq\bigoplus_{n'+n''=n}
 \lTor_{n'}^S(f_1,\ldots,f_m;
       \mathscr P^{(1)}_\bullet,\ldots,\mathscr P^{(m)}_\bullet)
          \otimes_{\OO_S}\HH_{n''}(\mathscr Q'_\bullet).
\end{aligned}
\end{equation}
特别地，当 $\mathscr Q'_\bullet$ 仅在次数0处等于 $\F'$ 时，
\begin{equation}\tag{6.9.6.2}
\begin{aligned}
 &\lTor_n^{S'}(f'_1,\ldots,f'_m,1_{S'};
     \mathscr P'^{(1)}_\bullet,\ldots,\mathscr P'^{(m)}_\bullet,\F')\\
 &\qquad\simeq\lTor_n^S(f_1,\ldots,f_m;
     \mathscr P^{(1)}_\bullet,\ldots,\mathscr P^{(m)}_\bullet)\otimes_{\OO_S}\F'.
\end{aligned}
\end{equation}
再取 $\F'=\OO_{S'}$，则
\begin{equation}\tag{6.9.6.3}
\begin{aligned}
 &\lTor_n^{S'}(f'_1,\ldots,f'_m;
     \mathscr P'^{(1)}_\bullet,\ldots,\mathscr P'^{(m)}_\bullet)\\
 &\qquad\simeq\lTor_n^S(f_1,\ldots,f_m;
     \mathscr P^{(1)}_\bullet,\ldots,\mathscr P^{(m)}_\bullet)\otimes_{\OO_S}\OO_{S'}.
\end{aligned}
\end{equation}
\begin{proof}
各项的平坦性使逐项复形 $\termTor^S_q(\mathscr P^{(i)}_\bullet,\OO_{S'})$ 在 $q\ne0$ 时为零（6.5.8）。因此 $(f')$ 退化；由条件二，它双正则（6.9.3），所以边缘同态
\begin{equation}\tag{6.9.6.4}
\begin{aligned}
 &\lTor_n^S(f_1,\ldots,f_m,1_{S'};
     \mathscr P^{(1)}_\bullet,\ldots,\mathscr P^{(m)}_\bullet,\mathscr Q'_\bullet)\\
 &\quad\longrightarrow{}^{(f')}\spE^2_{n0}
 =\lTor_n^{S'}(f'_1,\ldots,f'_m,1_{S'};
     \mathscr P'^{(1)}_\bullet,\ldots,\mathscr P'^{(m)}_\bullet,\mathscr Q'_\bullet)
\end{aligned}
\end{equation}
双射（0，11.1.6）。整体超 Tor 模的平坦性又使 ${}^{(e)}\spE^2_{pq}=0$（$p\ne0$）（6.5.8），所以 $(e)$ 也退化且双正则，其边缘同态
\begin{equation}\tag{6.9.6.5}
\begin{aligned}
 &\lTor_n^S(f_1,\ldots,f_m,1_{S'};
     \mathscr P^{(1)}_\bullet,\ldots,\mathscr P^{(m)}_\bullet,\mathscr Q'_\bullet)\\
 &\quad\longrightarrow{}^{(e)}\spE^2_{0n}
 =\bigoplus_{n'+n''=n}\lTor_{n'}^S(f_1,\ldots,f_m;
     \mathscr P^{(1)}_\bullet,\ldots,\mathscr P^{(m)}_\bullet)
       \otimes_{\OO_S}\HH_{n''}(\mathscr Q'_\bullet)
\end{aligned}
\end{equation}
也双射（0，11.1.6）。组合这两个同构便得（6.9.6.1）。若只剩第零项 $\F'$，同调只在次数0处非零，故得（6.9.6.2）；再用（6.7.12）得到（6.9.6.3）。
\end{proof}

\egas{6.9.7}{推论}
在（6.9.1）的条件下，设 $S,S'$ 仿射，对每个 $i$ 给定整数 $d_i$（$1\leq i\leq m$）。存在只依赖于 $S$、各 $X^{(i)}$ 和 $d_i$ 的整数 $N$，具有如下性质：对每个整数 $n_0$，以及满足下列条件的任意一组复形，当 $n\leq n_0$ 时有（6.9.6.3）的典范同构：一、$\mathscr P^{(i)}_k=0$（$k<d_i$）；二、$\mathscr P^{(i)}_k$ 在 $S$ 上平坦（$k<n_0+N$）；三、整体超 Tor 模在次数 $q<n_0+N$ 时在 $S$ 上平坦。
\begin{proof}
若 $\mathscr P^{(i)}_k$ 在 $k<r$ 时在 $S$ 上平坦，则
$\lTor^S_{q_i}(\mathscr P^{(i)}_k,\OO_{S'})=0$（$k<r$ 且 $q_i\ne0$）。按（6.6.2）的方法，借助 $X^{(i)}$ 上某个固定有限仿射覆盖 $\mathfrak U^{(i)}$ 的逆像来计算
\begin{equation}\tag{6.9.7.1}
 \lTor_p^{S'}\bigl(f'_1,\ldots,f'_m;
    \termTor^S_{q_1}(\mathscr P^{(1)}_\bullet,\OO_{S'}),\ldots,
    \termTor^S_{q_m}(\mathscr P^{(m)}_\bullet,\OO_{S'})\bigr).
\end{equation}
这些覆盖与 $S'$ 和各复形无关。相应 $L^{(i)}_{jk}$ 在 $j<-N_i$ 时为零，其中 $N_i$ 只依赖于覆盖。若某个 $q_i\ne0$，则用于计算（6.9.7.1）的单复形在所有次数
\[
 <r+\sum_{j\ne i}d_j-\sum_{j=1}^mN_j
\]
处为零。因此（6.9.7.1）在 $p<r-N$ 时为零，其中 $N$ 取各数
$\sum_jN_j-\sum_{j\ne i}d_j$ 的最大值。于是 ${}^{(f')}\spE^2_{pq}=0$，只要 $q\ne0$ 且 $p<r-N$；又它在 $q<0$ 时为零，故边缘同态（6.9.6.4）在 $n<r-N$ 时双射（M，XV，5.6），此处取 $\mathscr Q'_\bullet=\OO_{S'}$。另一方面，若整体超 Tor 模在 $q<r$ 时在 $S$ 上平坦，则 ${}^{(e)}\spE^2_{pq}=0$（$p\ne0$ 且 $q<r$）；它还在 $p<0$ 时为零，所以（6.9.6.5）在 $n<r$ 时双射。这完成证明。
\cmt{C3-H02}{底本陈述使用 $n\leq n_0$，证明的估计使用严格不等式。因结论仅断言某个常数 $N$ 存在，可以先把所得常数增大1，从而把 $n<n_0+1$ 写为 $n\leq n_0$。这里不把端点的记号差异当作反例。}
\end{proof}

应用中（6.9.3）最重要的情形是 $m=1$，且 $\mathscr Q'_\bullet$ 仅含次数0的项 $\F'$。为便于以后引用，重新陈述如下。
\orig{（6.9.8）所处理的情形，特别是谱序列（6.9.8.3）、（6.9.8.4），是J.-P. Serre于1957年向我们指出的。}
\egas{6.9.8}{命题}
设 $S$ 为预概形，$g:S'\to S$ 为态射，$f:X\to Y$ 为 $S$-预概形间的分离拟紧 $S$-态射，$\PP_\bullet$ 为下有界拟凝聚 $\OO_X$-模复形，$\F'$ 为拟凝聚 $\OO_{S'}$-模。存在关于 $\PP_\bullet,\F'$ 的两个双正则谱函子，取值于拟凝聚 $\OO_{Y_{(S')}}$-模范畴，共同终值为 $\lTor_*^S(f,1_{S'};\PP_\bullet,\F')$，其 $E^2$ 项为
\begin{align}
 {}'\spE^2_{pq}&=\lTor_p^S(\cH^{-q}(f,\PP_\bullet),\F'),\tag{6.9.8.1}\\
 {}''\spE^2_{pq}&=\cH^{-p}(f',\termTor_q^S(\PP_\bullet,\F')),\tag{6.9.8.2}
\end{align}
其中 $f'=f_{(S')}:X_{(S')}\to Y_{(S')}$。

这两个谱序列也可不经（6.9.3），而直接从（6.7.3）的 $(a),(b')$ 得到：取 $X^{(1)}=X$、$Y^{(1)}=Y$、$X^{(2)}=Y^{(2)}=S'$、$f_1=f$、$f_2=1_{S'}$。当 $S=S'=Y$ 且 $Y$ 仿射时，得到 $E^2$ 项为
\begin{align}
 {}'\spE^2_{pq}&=\lTor_p^Y(\cH^{-q}(f,\PP_\bullet),\F),\tag{6.9.8.3}\\
 {}''\spE^2_{pq}&=\cH^{-p}(f,\termTor_q^Y(\PP_\bullet,\F))\tag{6.9.8.4}
\end{align}
的两个谱序列。对任意拟凝聚且 $Y$-平坦的 $\OO_Y$-模 $\F$，或当 $\PP_\bullet$ 由 $Y$-平坦 $\OO_X$-模组成时对任意拟凝聚 $\OO_Y$-模 $\F$，它们由（6.7.6）收敛到函子 $f_*$ 关于复形 $\PP_\bullet\otimes_Y\F$ 的超上同调。这些谱序列不同于（6.2.1）的谱序列。

\egas{6.9.9}{推论}
在（6.9.8）的条件下，设 $\PP_\bullet$ 下有界，且由 $S$-平坦模组成，并设各 $\cH^n(f,\PP_\bullet)$ 在 $S$ 上平坦。则有函子性的典范同构
\begin{equation}\tag{6.9.9.1}
 \lTor_n^{S'}(f',1_{S'};\PP'_\bullet,\F')
 \simeq\cH^{-n}(f,\PP_\bullet)\otimes_{\OO_S}\F',
\end{equation}
其中 $\PP'_\bullet=\PP_\bullet\otimes_{\OO_S}\OO_{S'}$；特别地，当 $\F'=\OO_{S'}$ 时，
\begin{equation}\tag{6.9.9.2}
 \cH^n(f',\PP'_\bullet)\simeq\cH^n(f,\PP_\bullet)\otimes_{\OO_S}\OO_{S'}.
\end{equation}
这是（6.9.6）的 $m=1$ 特例。进一步得到：

\egas{6.9.10}{推论}
设 $S$ 为预概形，$f:X\to Y$ 为 $S$-预概形间的分离拟紧 $S$-态射，$\PP_\bullet$ 为由拟凝聚且 $S$-平坦的 $\OO_X$-模组成的下有界复形。又设各 $\cH^n(f,\PP_\bullet)$ 在 $S$ 上平坦。对每个 $s\in S$，以 $X_s,Y_s$ 表示纤维，$f_s:X_s\to Y_s$ 表示 $f\times_S1$，$\PP^s_\bullet=\PP_\bullet\otimes_{\OO_S}k(s)$。则有函子性的典范同构
\begin{equation}\tag{6.9.10.1}
 \cH^n(f_s,\PP^s_\bullet)\simeq\cH^n(f,\PP_\bullet)\otimes_{\OO_S}k(s).
\end{equation}
由此，在适当的平坦性假设下，导出函子 $R^nf_*$ 的形成“与取纤维交换”。第7节还将用另一种方法重新得到这一情形。

% END continuations/iii2_69.tex


% BEGIN continuations/iii2_610.tex
\subsection{某些上同调函子的局部结构}
\source{\texttt{EGAIII.2.pdf}，PDF第38--42页；原书第39--43页}
\egas{6.10.1}{命题}
设 $S=\Spec A$ 为仿射概形，$Y^{(i)}$（$1\leq i\leq n$）为有限个在 $S$ 上平坦的仿射 $S$-概形。对每个 $i$，取分离拟紧 $S$-态射 $f_i:X^{(i)}\to Y^{(i)}$ 及下有界拟凝聚 $\OO_{X^{(i)}}$-模复形 $\mathscr P^{(i)}_\bullet$。以 $Y$ 表示各 $Y^{(i)}$ 在 $S$ 上的乘积。存在由拟凝聚且 $S$-平坦的 $\OO_Y$-模组成的复形 $\mathscr R_\bullet$，具有如下性质：对每个仿射 $S$-概形 $S'$ 和每个下有界拟凝聚 $\OO_{S'}$-模复形 $\mathscr Q'_\bullet$，有同构
\begin{equation}\tag{6.10.1.1}
\begin{aligned}
 &\lTor_*^S(f_1,\ldots,f_n,1_{S'};
        \mathscr P^{(1)}_\bullet,\ldots,\mathscr P^{(n)}_\bullet,\mathscr Q'_\bullet)\\
 &\qquad\simeq\HH_*(\mathscr R_\bullet\otimes_S\mathscr Q'_\bullet),
\end{aligned}
\end{equation}
且这是关于 $\mathscr Q'_\bullet$ 的 $\delta$-函子同构。又对任意仿射 $S$-概形间的 $S$-态射 $u:S''\to S'$，以下图交换：
\begin{equation}\tag{6.10.1.2}
\begin{tikzcd}[column sep=small,row sep=large]
 \lTor_*^S(f_1,\ldots,f_n,1_{S'};\mathscr P^{(1)}_\bullet,\ldots,\mathscr P^{(n)}_\bullet,\mathscr Q'_\bullet)
 \arrow[r,"\sim"]\arrow[d]&
 \HH_*(\mathscr R_\bullet\otimes_S\mathscr Q'_\bullet)\arrow[d]\\
 \lTor_*^S(f_1,\ldots,f_n,1_{S''};\mathscr P^{(1)}_\bullet,\ldots,\mathscr P^{(n)}_\bullet,u^*\mathscr Q'_\bullet)
 \arrow[r,"\sim"]&
 \HH_*(\mathscr R_\bullet\otimes_Su^*\mathscr Q'_\bullet).
\end{tikzcd}
\end{equation}
其中竖箭头为（6.7.10）定义的典范 $(1_Y\times_Su)$-态射。
\begin{proof}
按照（6.6.2）的方法并结合（6.6.6）来计算超 Tor。以 $A_i$ 表示 $Y^{(i)}$ 的环，则每个 $L^{(i)}$ 是 $A_i$-模双复形，故有一个按（0，11.7.1）意义、由投射 $A_i$-模组成的 Cartan--Eilenberg 解消 $M^{(i)}$。由（6.6.6），（6.10.1.1）左边典范同构于 $H_*(M\otimes_AQ'_\bullet)$，其中
\[
 M=M^{(1)}\otimes_A\cdots\otimes_AM^{(n)},\qquad
 \mathscr Q'_\bullet=\widetilde{Q'_\bullet}.
\]
由于各 $A_i$ 在 $A$ 上平坦，$M^{(i)}$ 都是平坦 $A$-模三重复形（$0_{\mathrm I}$，6.2.1），$M$ 也如此。又若 $B$ 为 $Y$ 的环，即各 $A_i$ 的张量积，则 $M$ 是 $B$-模三重复形。将 $M$ 按总次数视为单复形后，$\mathscr R_\bullet=\widetilde M$ 即满足要求；它与（6.7.10）的相容性也容易看出。
\end{proof}

\egas{6.10.2}{推论}
在（6.10.1）中，可令 $\mathscr R_\bullet$ 下有界。若各 $\mathscr P^{(i)}_\bullet$ 上有界，且各 $Y^{(i)}$ 的上同调维数有限，还可令 $\mathscr R_\bullet$ 上有界。

第一断言由各 $M^{(i)}$ 的三个次数均下有界得到。若各环 $A_i$ 的上同调维数有限，则每个 $M^{(i)}$ 的第三次数只取有限多个值；按构造，它的第一次数也只取有限多个值（6.6.2）。若 $\mathscr P^{(i)}_\bullet$ 的次数上有界，则第二次数也上有界（6.6.2），第二断言由此得到。

\egas{6.10.3}{注记}
（i）采用（6.10.1）的记号，因 $\mathscr R_\bullet$ 由 $S$-平坦 $\OO_Y$-模组成，
$\HH_*(\mathscr R_\bullet\otimes_S\mathscr Q'_\bullet)$ 与
$\lTor_*(\mathscr R_\bullet,\mathscr Q'_\bullet)$ 同构（6.5.9）。因此它是一个正则谱序列的终值（6.5.4），其 $E^2$ 项为
\begin{equation}\tag{6.10.3.1}
 E^2_{pq}=\bigoplus_{q'+q''=q}\lTor_p^S
       (\HH_{q'}(\mathscr R_\bullet),\HH_{q''}(\mathscr Q'_\bullet)).
\end{equation}
这正是基变换谱序列 $(e)$（6.9.3）。

（ii）设 $\mathscr R'_\bullet$ 是另一拟凝聚且 $S$-平坦的 $\OO_Y$-模复形，$g:\mathscr R'_\bullet\to\mathscr R_\bullet$ 是复形同态，且
$\HH_*(g):\HH_*(\mathscr R'_\bullet)\to\HH_*(\mathscr R_\bullet)$ 为同构。由（6.3.3）、（6.5.9），$g$ 诱导关于 $\mathscr Q'_\bullet$ 的 $\delta$-函子同构
\[
 \HH_*(\mathscr R'_\bullet\otimes_S\mathscr Q'_\bullet)
 \simeq\HH_*(\mathscr R_\bullet\otimes_S\mathscr Q'_\bullet),
\]
并使下图交换：
\[
\begin{tikzcd}[column sep=large,row sep=large]
 \HH_*(\mathscr R'_\bullet\otimes_S\mathscr Q'_\bullet)
 \arrow[r,"\sim"]\arrow[d]&
 \HH_*(\mathscr R_\bullet\otimes_S\mathscr Q'_\bullet)\arrow[d]\\
 \HH_*(\mathscr R'_\bullet\otimes_Su^*\mathscr Q'_\bullet)
 \arrow[r,"\sim"]&
 \HH_*(\mathscr R_\bullet\otimes_Su^*\mathscr Q'_\bullet).
\end{tikzcd}
\]
这说明，（6.10.1）的性质并不完全决定复形 $\mathscr R_\bullet$。

（iii）在（6.10.1）的证明中，可令各 $M^{(i)}$ 由自由 $A_i$-模组成；这由将（0，11.5.2.1）的证明对偶化容易得到。于是 $M^{(i)}\otimes_{A_i}B$ 由自由 $B$-模组成，而 $M$ 是这些复形在 $B$ 上的张量积，所以还可令 $\mathscr R_\bullet$ 与一个自由 $B$-模复形相伴。又由（M，XVII，1.2），三重复形 $M$ 函子性地依赖于各 $L^{(i)}$，故在固定每个 $X^{(i)}$ 的有限覆盖时，也函子性地依赖于各 $\mathscr P^{(i)}_\bullet$；这里三重复形的“态射”理解为同态按同伦关系所得的等价类。把 $X^{(i)}$ 的覆盖换成更细的覆盖，恰好在各 $L^{(i)}$ 间产生只确定到同伦的同态（6.6.8）。采用上述态射约定，最终可见 $M$ 是关于每个 $\mathscr P^{(i)}_\bullet$ 的函子。在（6.1.3）所提到的、专论上同调函子一般代数的章中，将进一步说明这种函子性，特别是 $\mathscr R_\bullet$ 对复形短正合列
\[
 0\longrightarrow\mathscr P'^{(i)}_\bullet
 \longrightarrow\mathscr P^{(i)}_\bullet
 \longrightarrow\mathscr P''^{(i)}_\bullet\longrightarrow0
\]
的行为。

\egas{6.10.4}{说明}
$\mathscr R_\bullet$ 由 $S$-平坦 $\OO_Y$-模组成，容易推出
$\HH_*(\mathscr R_\bullet\otimes_S\mathscr Q'_\bullet)$ 是关于 $\mathscr Q'_\bullet$ 的同调函子，参见（7.7.1）的论证。正如（6.1.1）所说，这一性质正是引入超 Tor 的动机。置
\[
\begin{gathered}
 X=X^{(1)}\times_S\cdots\times_SX^{(n)},\quad
 f=f_1\times_S\cdots\times_Sf_n,\\
 \PP_\bullet=\mathscr P^{(1)}_\bullet\otimes_S\cdots\otimes_S\mathscr P^{(n)}_\bullet,\\
 X'=X\times_SS',\qquad Y'=Y\times_SS',\qquad f'=f\times_S1_{S'}.
\end{gathered}
\]
基变换问题引导我们研究 $\cH^\bullet(f',\PP_\bullet\otimes_S\N')$ 关于拟凝聚 $\OO_{S'}$-模 $\N'$ 的函子性，或研究
$\cH^\bullet(f,\PP_\bullet\otimes_S\N)$ 关于拟凝聚 $\OO_S$-模 $\N$ 的函子性。若各 $\mathscr P^{(i)}_\bullet$，因而 $\PP_\bullet$，在 $S$ 上平坦，则由上述结果及（6.7.6），它确为关于 $\N$ 的上同调函子。但若不再假设各 $\mathscr P^{(i)}_\bullet$ 平坦，则不再如此，因而无法用通常的同调代数方法研究这个函子。

我们尤其要用 $n=1$、$Y=S$，且 $\PP_\bullet$ 由 $Y$-平坦 $\OO_X$-模组成的情形。此时有：

\egas{6.10.5}{定理}
设 $Y=\Spec A$ 是诺特仿射概形，$f:X\to Y$ 为固有态射，$\PP_\bullet$ 是由凝聚且 $Y$-平坦 $\OO_X$-模组成的下有界复形。存在下有界 $\OO_Y$-模复形 $\LL_\bullet$，其各项 $\LL_j$ 具有 $\OO_Y^{r_j}$ 的形式，以及同构
\begin{equation}\tag{6.10.5.1}
 \cH^\bullet(f,\PP_\bullet\otimes_Y\mathscr Q_\bullet)
 \simeq\HH_*(\LL_\bullet\otimes_Y\mathscr Q_\bullet),
\end{equation}
且这是关于下有界拟凝聚 $\OO_Y$-模复形 $\mathscr Q_\bullet$ 的 $\delta$-函子同构。又对于任意态射 $u:Y'\to Y$，置
\[
 X'=X_{(Y')},\qquad f'=f_{(Y')},\qquad
 \PP'_\bullet=\PP_\bullet\otimes_Y\OO_{Y'},\qquad
 \LL'_\bullet=u^*(\LL_\bullet),
\]
其中 $\LL'_\bullet$ 是有限型局部自由 $\OO_{Y'}$-模复形；则有同构
\begin{equation}\tag{6.10.5.2}
 \cH^\bullet(f',\PP'_\bullet\otimes_{Y'}\mathscr Q'_\bullet)
 \simeq\HH_*(\LL'_\bullet\otimes_{Y'}\mathscr Q'_\bullet),
\end{equation}
为关于下有界拟凝聚 $\OO_{Y'}$-模复形 $\mathscr Q'_\bullet$ 的 $\delta$-函子同构，并使图
\begin{equation}\tag{6.10.5.3}
\begin{tikzcd}[column sep=large,row sep=large]
 \cH^\bullet(f,\PP_\bullet\otimes_Y\mathscr Q_\bullet)
 \arrow[r,"\sim"]\arrow[d]&
 \HH_*(\LL_\bullet\otimes_Y\mathscr Q_\bullet)\arrow[d]\\
 \cH^\bullet(f',\PP'_\bullet\otimes_{Y'}u^*\mathscr Q_\bullet)
 \arrow[r,"\sim"]&
 \HH_*(\LL'_\bullet\otimes_{Y'}u^*\mathscr Q_\bullet)
\end{tikzcd}
\end{equation}
交换。
\cmt{C3-N02}{本节沿用原著将次数为 $-1$ 的复形转成上同调次数的约定；上述带点的同构在具体次数上是 $\cH^{-j}\simeq\HH_j$，不是 $\cH^j\simeq\HH_j$。}
\begin{proof}
应用（6.10.1），先得到由拟凝聚且 $Y$-自由 $\OO_Y$-模组成的下有界复形 $\mathscr R_\bullet$（6.10.2、6.10.3，（ii）），以及关于 $\mathscr Q_\bullet$ 的 $\delta$-函子同构
\begin{equation}\tag{6.10.5.4}
 \cH^\bullet(f,\PP_\bullet\otimes_Y\mathscr Q_\bullet)
 \simeq\HH_*(\mathscr R_\bullet\otimes_Y\mathscr Q_\bullet).
\end{equation}
\cmt{C3-T03}{此处原文引（6.10.3，（ii））说明可取自由项；直接给出这一点的是同条的（iii）。后文用（ii）比较同调同构则有其实际作用，不能将全部（ii）机械改为（iii）。}
但先验地，$\mathscr R_\bullet$ 的各项不一定有限型。将（6.10.5.4）用于仅第零项为 $\OO_Y$ 的复形，可见 $\HH_*(\mathscr R_\bullet)$ 同构于 $\cH^\bullet(f,\PP_\bullet)$，因而由凝聚 $\OO_Y$-模组成（6.2.5）。于是（0，11.9.2）保证存在下有界复形 $\LL_\bullet$，各项与有限型自由 $A$-模相伴，以及同态 $\LL_\bullet\to\mathscr R_\bullet$，使同调上的相应同态双射。由（6.10.3，（ii）），得到（6.10.5.1）。当 $Y'$ 仿射时，其余断言由（6.10.1）及（6.10.3，（ii））得到。一般情形，只须取 $Y'$ 的仿射开覆盖 $(V_\alpha)$，并验证：在任意仿射开集 $W\subset V_\alpha\cap V_\beta$ 上，对应 $V_\alpha,V_\beta$ 的同构（6.10.5.2）的限制，都等于对应 $W$ 的同构。这由把（6.10.1.2）应用于典范嵌入 $W\to V_\alpha$、$W\to V_\beta$ 的交换性得到。
\end{proof}

\egas{6.10.6}{注记}
后续各章主要将（6.10.5）用于 $\PP_\bullet$ 只含一个凝聚且 $Y$-平坦的 $\OO_X$-模 $\F$ 的情形。这时
\[
 \HH_n(\LL_\bullet)\simeq\cH^{-n}(f,\F)=R^{-n}f_*(\F)
\]
（6.2.1），所以 $\HH_n(\LL_\bullet)=0$（$n>0$）。后面（7.7.12，（i））将说明：若把各 $\LL_n$ 与有限型自由 $A$-模相伴的要求，改为有限型局部自由，则还可令 $\LL_\bullet$ 只含次数不大于0的项，因而只含有限项。

与这样的 $\F$ 对应的复形 $\LL_\bullet$，除上述次数限制外，似乎没有特别性质。因此反过来可以问：给定由与有限型投射 $A$-模相伴的 $\OO_Y$-模组成、下有界且次数大于0的项全为零的复形 $\LL_\bullet$，是否存在一个在 $Y$ 上射影且平坦的 $Y$-概形 $X$ 和一个局部自由 $\OO_X$-模 $\F$，使得有关于 $\mathscr Q_\bullet$ 函子性的同构
\[
 \cH^\bullet(f,\F\otimes_Y\mathscr Q_\bullet)
 \simeq\HH_*(\LL_\bullet\otimes_Y\mathscr Q_\bullet)?
\]
这样的结果之所以有意义，是因为它将把固有 $Y$-概形上凝聚且 $Y$-平坦模的上同调理论，完全归结为诺特环 $A$ 上有限型投射模复形的“同伦意义下”的理论。
\cmt{C3-H03}{这里是1963年原著提出的问题，译文保留其历史语境；不据此陈述该问题今天是否已有答案。}

% END continuations/iii2_610.tex


% BEGIN continuations/iii2_71_72.tex
% 依据 EGAIII.2.pdf，第42--47页（原印刷43--48页）逐段翻译。
\section{模的协变同调函子中的基变换研究}
\subsection{\texorpdfstring{$A$-模的函子}{A-模的函子}}
\source{第三卷第二部分，原印刷第43--44页}

\egap{7.1.1}
给定一个环 $A$（不必交换），我们用 $\mathbf{Ab}_A$ 表示左 $A$-模范畴，用 $\mathbf{Ab}$ 表示 $\mathbf Z$-模范畴，即交换群范畴。设 $T:\mathbf{Ab}_A\to\mathbf{Ab}$ 是协变加性函子，$M$ 是一个 $(A,A)$-双模；则 $T(M)$ 自然带有右 $A$-模结构。事实上，对每个 $a\in A$，记 $h_a^M$（或简记 $h_a$）为左 $A$-模 $M$ 的自同态 $x\mapsto xa$。依假设，$T(h_a)$ 是 $\mathbf Z$-模 $T(M)$ 的自同态；又因 $T$ 是协变加性函子，对 $a,b\in A$ 有
\[
\begin{aligned}
 T(h_{ab})&=T(h_b\circ h_a)=T(h_b)\circ T(h_a),\\
 T(h_{a+b})&=T(h_a+h_b)=T(h_a)+T(h_b).
\end{aligned}
\]
因此，$(a,y)\mapsto T(h_a)(y)$ 给出一个右 $A$-模的外部运算。特别地，$T(A_{\mathrm s})$ 是右 $A$-模。
\cmt{C3-N03}{沿用底本的 $A_{\mathrm s}$，表示把 $A$ 看成左 $A$-模；右模结构来自右乘映射，不应与输入函子时的左模结构混淆。}

\egap{7.1.2}
当 $A$ 交换时，由（7.1.1），对任意 $A$-模 $M$，$T(M)$ 自然带有 $A$-模结构。若 $u:M\to N$ 为 $A$-线性同态，则对任意 $a\in A$ 有
\[
 u\circ h_a^M=h_a^N\circ u,
 \qquad T(u)\circ T(h_a^M)=T(h_a^N)\circ T(u),
\]
所以 $T(u):T(M)\to T(N)$ 为 $A$-线性同态。因此可把 $T$ 看成 $\mathbf{Ab}_A$ 到自身的协变加性函子。更准确地说，这样就在协变加性函子 $\mathbf{Ab}_A\to\mathbf{Ab}$ 的范畴与协变 $A$-线性函子 $T:\mathbf{Ab}_A\to\mathbf{Ab}_A$ 的范畴之间定义了典范等价；这里 $A$-线性指对每个 $a\in A$ 有 $T(h_a^M)=h_a^{T(M)}$。由于将 $A$-模送到其底层 $\mathbf Z$-模的遗忘函子 $\mathbf{Ab}_A\to\mathbf{Ab}$ 正合且忠实，上述等价所对应的两个函子具有相同的正合性性质。

\egap{7.1.3}
仍设 $A$ 交换，$B$ 为一个 $A$-代数（不必交换），$\rho:A\to B$ 为给出此代数结构的环同态。它定义一个协变加性函子
\[
 \rho_*:\mathbf{Ab}_B\longrightarrow\mathbf{Ab}_A,
 \qquad M\longmapsto M_{[\rho]}.
\]
复合后得到显然协变且加性的函子
\[
 T_{(B)}:\mathbf{Ab}_B\xrightarrow{\rho_*}\mathbf{Ab}_A
 \xrightarrow{T}\mathbf{Ab},
\]
出于排印便利，也记作 $T^{(B)}$ 或 $T\otimes_A B$，并称它由 $T$ 经从 $A$ 到 $B$ 的标量扩张得到。若 $B$ 交换，当然可把 $T_{(B)}$ 看作 $\mathbf{Ab}_B$ 到自身的函子（7.1.2）。当 $B$ 交换而 $C$ 为 $B$-代数时，立即有 $T_{(C)}=(T_{(B)})_{(C)}$。标量扩张对 $T$ 函子性且加性；若 $T$ 与归纳极限或直和交换（分别若 $T$ 左正合、右正合或正合），则 $T_{(B)}$ 亦然，因为 $\rho_*$ 正合，并与归纳极限和直和交换。

\egap{7.1.4}
仍设 $A$ 交换，并设 $T:\mathbf{Ab}_A\to\mathbf{Ab}_A$ 为与归纳极限交换的协变加性 $A$-线性函子。那么对 $A$ 的任一乘法子集 $S$ 和任一 $A$-模 $M$，有关于 $M$ 函子性的典范 $A$-模同构
\begin{equation}
 T(S^{-1}M)\simeq S^{-1}T(M).\tag{7.1.4.1}
\end{equation}
先设 $S$ 为某个 $f\in A$ 的各次幂 $f^n$（$n\geq0$）组成的集合。已知 $M_f=\varinjlim M_n$，其中每个 $M_n=M$，过渡同态 $M_n\to M_m$ 为乘以 $f^{m-n}$（$m\geq n$）（$0_{\mathrm I}$，1.6.1）；因此，由关于 $T$ 的假设得到（7.1.4.1）。对一般 $S$，$S^{-1}M$ 是各 $M_f$（$f\in S$）的归纳极限（$0_{\mathrm I}$，1.4.5），同理得证。同构（7.1.4.1）的函子性还说明它为 $S^{-1}A$-模同构，因此，在相差一个典范同构的意义下，可写成
\begin{equation}
 T_{(S^{-1}A)}(S^{-1}M)=S^{-1}T(M)=T(M)_{(S^{-1}A)}.
 \tag{7.1.4.2}
\end{equation}
若 $S=A-\mathfrak p$ 为素理想的补集，则以 $T_{\mathfrak p}$ 代替 $T^{(A_{\mathfrak p})}$。

\egas{7.1.5}{命题}
在（7.1.4）的假设下，若对 $A$ 的每个极大理想 $\mathfrak m$，$T_{\mathfrak m}$ 均左正合（分别右正合、正合），则 $T$ 左正合（分别右正合、正合）。

事实上，已知当 $A$-模 $M$ 的两个子模 $N,P$ 对每个极大理想 $\mathfrak m$ 都满足 $N_{\mathfrak m}=P_{\mathfrak m}$ 时，有 $N=P$（Bourbaki，\textit{Alg. comm.}，第II章，§3，第3目，定理1）。

\subsection{张量积函子的刻画}
\source{第三卷第二部分，原印刷第44--48页}

\egap{7.2.1}
设 $A$ 是环（不必交换），$M$ 为左 $A$-模，$N$ 为右 $A$-模，$P$ 为 $\mathbf Z$-模。回顾：给定一个 $\mathbf Z$-线性同态 $v:N\otimes_A M\to P$，等价于给定一个 $\mathbf Z$-双线性映射 $u:N\times M\to P$，满足
\[
 u(ta,x)=u(t,ax)\qquad(a\in A,\ t\in N,\ x\in M),
\]
二者由 $v(t\otimes x)=u(t,x)$ 联系。另一方面，给定 $u$ 又等价于给定 $\mathbf Z$-线性同态 $x\mapsto f_x$，从 $M$ 到 $\Hom_{\mathbf Z}(N,P)$，使得 $f_{ax}(t)=f_x(ta)$；联系式为 $u(t,x)=f_x(t)$。

\egap{7.2.2}
设 $T:\mathbf{Ab}_A\to\mathbf{Ab}$ 为协变加性函子。对每个左 $A$-模 $M$，我们将定义一个关于 $M$ 函子性的典范 $\mathbf Z$-线性同态
\begin{equation}
 t_M:T(A_{\mathrm s})\otimes_A M\longrightarrow T(M).
 \tag{7.2.2.1}
\end{equation}
由（7.2.1），只须定义 $\mathbf Z$-线性同态
\[
 x\longmapsto t_M^\#(x):M\longrightarrow
 \Hom_{\mathbf Z}(T(A_{\mathrm s}),T(M)),
\]
满足 $t_M^\#(ax)(y)=t_M^\#(x)(ya)$。注意右模 $T(A_{\mathrm s})$ 的结构在该 Hom 群上给出左 $A$-模结构，规定
\[
 (a\cdot v)(y)=v(ya).
\]
把 $t_M^\#$ 定义为两个典范同态的复合
\[
 M\longrightarrow\Hom_A(A_{\mathrm s},M)
 \longrightarrow\Hom_{\mathbf Z}(T(A_{\mathrm s}),T(M)),
\]
其中第二个映射是 $u\mapsto T(u)$，第一个是 $x\mapsto\theta_x$，$\theta_x(\xi)=\xi x$，给出典范 $A$-模同构。由 $\theta_{ax}=\theta_x\circ h_a$，有
\[
 T(\theta_{ax})(y)=T(\theta_x)(T(h_a)(y))=T(\theta_x)(ya),
\]
这就证明了 $t_M$ 的存在。直接验证它关于 $M$ 函子性：对左 $A$-模同态 $w:M\to M'$，图
\begin{equation}
 \begin{tikzcd}
 T(A_{\mathrm s})\otimes_A M\arrow[r,"t_M"]\arrow[d,"1\otimes w"']&T(M)\arrow[d,"T(w)"]\\
 T(A_{\mathrm s})\otimes_A M'\arrow[r,"t_{M'}"']&T(M')
 \end{tikzcd}\tag{7.2.2.2}
\end{equation}
交换。当 $A$ 交换时，由此函子性知（7.2.2.1）为 $A$-线性同态（参见7.1.2）。

\egap{7.2.3}
当 $A$ 交换时，更一般地，对每个 $A$-模 $N$ 可定义典范 $A$-线性同态
\begin{equation}
 T(N)\otimes_A M\longrightarrow T(N\otimes_A M).
 \tag{7.2.3.1}
\end{equation}
只须在（7.2.2）的构造中，把 $\theta_x$ 换成 $A$-线性同态 $N\to N\otimes_A M$，$y\mapsto y\otimes x$。该同态对 $M,N$ 均函子性。

特别地，若 $B$ 为 $A$-代数（不必交换），则有关于 $M$ 函子性的同态
\begin{equation}
 \begin{aligned}
 (T(M))_{(B)}=T(M)\otimes_A B
 &\longrightarrow T(M\otimes_A B)\\
 &=T^{(B)}(M_{(B)}),
 \end{aligned}\tag{7.2.3.2}
\end{equation}
由（7.2.3.1）对 $M$ 的函子性，它为 $B$-线性同态。还有交换图
\begin{equation}
 \begin{tikzcd}[column sep=large]
 T(A)\otimes_A M\arrow[r,"t_M"]\arrow[d]&T(M)\arrow[d]\\
 T_{(B)}(B_{\mathrm s})\otimes_B M_{(B)}\arrow[r]&T_{(B)}(M_{(B)})
 \end{tikzcd}\tag{7.2.3.3}
\end{equation}
右侧竖箭头为典范同态与（7.2.3.2）的复合
\[
 T(M)\longrightarrow T(M)\otimes_A B\longrightarrow
 T(M\otimes_A B).
\]
左侧竖箭头为
\[
 \begin{aligned}
 T(A)\otimes_A M&\longrightarrow
 T_{(B)}(B_{\mathrm s})\otimes_B(B\otimes_A M)\\
 &=T_{(B)}(B_{\mathrm s})\otimes_A M,
 \end{aligned}
\]
其中 $T(A)\to T_{(B)}(B_{\mathrm s})=T(B)$ 是 $T(\rho)$，把 $\rho$ 看作 $A$-模同态 $A\to B_{[\rho]}$。

\egas{7.2.4}{引理}
若 $T:\mathbf{Ab}_A\to\mathbf{Ab}$ 是与直和交换的协变加性函子，则对每个自由 $A$-模 $L$，典范同态 $t_L$（7.2.2.1）为同构。

事实上，写 $L=\bigoplus_{\alpha\in I}L_\alpha$，各 $L_\alpha\simeq A_{\mathrm s}$。由（7.2.2）的定义及 $T$ 与直和交换，$t_L=\bigoplus_\alpha t_{L_\alpha}$；其中作用于 Hom 的映射由各个 $\mathbf Z$-线性映射
\[
 T_\alpha:\Hom_A(A_{\mathrm s},L_\alpha)\longrightarrow
 \Hom_{\mathbf Z}(T(A_{\mathrm s}),T(L_\alpha))
\]
给出。因此归结为 $L=A_{\mathrm s}$，此时 $t_L$ 就是对任意右模成立的典范同构 $T(A_{\mathrm s})\otimes_A A\simeq T(A_{\mathrm s})$。

\egas{7.2.5}{命题}
设 $T:\mathbf{Ab}_A\to\mathbf{Ab}$ 为与直和交换的协变加性函子。下列条件等价：
\begin{enumerate}[label=\alph*)]
\item $T$ 右正合。
\item 对每个左 $A$-模 $M$，典范同态 $t_M$（7.2.2.1）为同构。
\item[\textup{b')}] $T$ 半正合，且对每个左 $A$-模 $M$，$t_M$ 满射。
\item[\textup{c)}] $T$ 同构于某个形如 $M\mapsto N\otimes_A M$ 的函子，其中 $N$ 为右 $A$-模。
\end{enumerate}
\cmt{C3-N04}{这里“半正合”指把短正合列送为在中间项正合的列；不能将它译成左正合、右正合或完全正合。}

b) 蕴含 c)，c) 蕴含 a)。证明 a) 蕴含 b)。令 $T'(M)=T(A_{\mathrm s})\otimes_A M$。取自由表示 $L'\to L\to M\to0$。因 $T,T'$ 右正合，得到两行正合的交换图
\[
 \begin{tikzcd}
 T'(L')\arrow[r]\arrow[d,"t_{L'}"']&T'(L)\arrow[r]\arrow[d,"t_L"']&T'(M)\arrow[r]\arrow[d,"t_M"]&0\\
 T(L')\arrow[r]&T(L)\arrow[r]&T(M)\arrow[r]&0.
 \end{tikzcd}
\]
由（7.2.4），$t_{L'},t_L$ 为同构，故五引理给出 $t_M$ 为同构。b) 显然蕴含 b')；为证明 b') 蕴含 a)，只须证明下述引理。

\egas{7.2.5.1}{引理}
设 $\mathbf K,\mathbf K'$ 为阿贝尔范畴，$F,G:\mathbf K\to\mathbf K'$ 为协变加性函子，$f:F\to G$ 为函子态射（T，I，1.2），且对每个对象 $E$，$f_E:F(E)\to G(E)$ 为满态射。若 $F$ 右正合而 $G$ 半正合，则 $G$ 右正合。

只须证明：对 $\mathbf K$ 中每个满态射 $v:E'\to E$，$G(v)$ 也是满态射。在交换图
\[
 \begin{tikzcd}
 F(E')\arrow[r,"F(v)"]\arrow[d,"f_{E'}"']&F(E)\arrow[d,"f_E"]\\
 G(E')\arrow[r,"G(v)"']&G(E)
 \end{tikzcd}
\]
中，$F(v),f_{E'},f_E$ 都是满态射，故 $G(v)$ 也是满态射。

\egas{7.2.6}{注记}
对每个右 $A$-模 $N$，置 $T_N(M)=N\otimes_A M$。它是从 $\mathbf{Ab}_A$ 到 $\mathbf{Ab}$ 的右正合、与直和交换的协变加性函子。将 $T_N(A_{\mathrm s})$ 典范地认作 $N$ 后，直接验证相应的（7.2.2.1）变成恒等同态。因此（7.2.5，c）中的右模 $N$ 由该性质确定至唯一同构，并典范同构于 $T(A_{\mathrm s})$。也可说，$T\mapsto T(A_{\mathrm s})$ 与 $N\mapsto T_N$ 定义右 $A$-模范畴和右正合且与直和交换的协变加性函子 $\mathbf{Ab}_A\to\mathbf{Ab}$ 的范畴之间的等价（T，I，1.2）。

\egas{7.2.7}{命题}
设 $A$ 为左阿廷环，其模根基 $\mathfrak m$ 的商为一个域 $k$。设 $T:\mathbf{Ab}_A\to\mathbf{Ab}$ 为与直和交换的协变加性函子。则（7.2.5）的条件还等价于：

\noindent d) $T$ 半正合，且由典范同态 $\varepsilon:A_{\mathrm s}\to k$ 导出的 $T(\varepsilon):T(A_{\mathrm s})\to T(k)$ 满射。

（7.2.5）的 b') 显然蕴含 d)。证明 d) 蕴含 b')。存在整数 $n$ 使 $\mathfrak m^n=0$。对任意 $A$-模 $M$ 置 $M_h=\mathfrak m^hM$，按 $h$ 作下降归纳证明 $t_{M_h}$ 满射。$h=n$ 时显然；当 $h<n$ 时有短正合列
\[
 0\longrightarrow M_{h+1}\longrightarrow M_h
 \longrightarrow M_h/M_{h+1}\longrightarrow0.
\]
归纳假设给出 $t_{M_{h+1}}$ 满射。商 $M_h/M_{h+1}$ 被 $\mathfrak m$ 消去，故为 $k$-模，即若干个 $k$ 的直和。因 $T$ 与直和交换，只须证明 $t_k$ 满射。考虑交换图
\[
 \begin{tikzcd}
 T(A_{\mathrm s})\otimes_A A_{\mathrm s}\arrow[r,"t_A"]\arrow[d,"1\otimes\varepsilon"']&T(A_{\mathrm s})\arrow[d,"T(\varepsilon)"]\\
 T(A_{\mathrm s})\otimes_A k\arrow[r,"t_k"']&T(k).
 \end{tikzcd}
\]
由（7.2.4）与 d)，$t_k$ 满射。

最后，只须证明：对短正合列 $0\to M'\to M\xrightarrow{v}M''\to0$，若 $t_{M'},t_{M''}$ 满射，则 $t_M$ 满射。令 $T'=T(A_{\mathrm s})\otimes_A-$；有交换图
\[
 \begin{tikzcd}[column sep=small]
 T'(M')\arrow[r]\arrow[d]&T'(M)\arrow[r]\arrow[d]&T'(M'')\arrow[r]\arrow[d]&0\arrow[d]\\
 T(M')\arrow[r]&T(M)\arrow[r]&T(M'')\arrow[r]&\Coker T(v).
 \end{tikzcd}
\]
因 $T$ 半正合，两行均正合。外侧两个 $t$ 为满态射，最右竖箭头为单态射，五引理（M，I，1.1）遂给出 $t_M$ 满射。

% END continuations/iii2_71_72.tex


% BEGIN continuations/iii2_73.tex
\subsection{模的同调函子的正合性判据}
\source{第三卷第二部分，原印刷第48--53页}

\egas{7.3.1}{命题}
设 $A$ 是环（不必交换），$T_\bullet$ 为从 $\mathbf{Ab}_A$ 到 $\mathbf{Ab}$ 的协变同调函子（T，II，2.1），并与直和交换。设整数 $p$ 使得 $T_p,T_{p-1}$ 均已定义。下列条件等价：
\begin{enumerate}[label=\alph*)]
\item $T_p$ 右正合。
\item $T_{p-1}$ 左正合。
\item 对每个左 $A$-模 $M$，典范函子同态
\begin{equation}
 T_p(A_{\mathrm s})\otimes_A M\longrightarrow T_p(M)\tag{7.3.1.1}
\end{equation}
为同构。
\item 对每个左 $A$-模 $M$，（7.3.1.1）为满态射。
\item $T_p$ 同构于某个 $M\mapsto N\otimes_A M$，其中 $N$ 为右 $A$-模。
\end{enumerate}
若 $A,\mathfrak m$ 还满足（7.2.7）的条件，则以上条件还等价于：

\noindent f) 典范同态 $T_p(\varepsilon):T_p(A_{\mathrm s})\to T_p(k)$ 为满态射。

由同调函子的定义，每个已定义的 $T_i$ 都半正合；对短正合列 $0\to M'\xrightarrow{u}M\xrightarrow{v}M''\to0$，有
\[
 \Ker T_{p-1}(u)\simeq\Coker T_p(v).
\]
因此 a)、b) 等价，其余断言由（7.2.5）和（7.2.7）得出。

\egas{7.3.2}{推论}
设 $A$ 交换。沿用（7.3.1）的记号，并设 $T_p$ 右正合。若 $f\in A$ 不属于某个 $A$-模 $M$ 中任何非零元素的消去子，则它也不属于 $T_{p-1}(M)$ 中任何非零元素的消去子。特别地，若 $A$ 为整环，则 $T_{p-1}(A)$ 为无挠 $A$-模。

事实上，若 $h_f$ 表示 $M$ 中的数乘 $x\mapsto fx$，假设就是 $h_f$ 单射；由（7.3.1，b），$T_{p-1}(h_f)$ 也单射。

\egas{7.3.3}{命题}
设 $A$ 为环，$T_\bullet:\mathbf{Ab}_A\to\mathbf{Ab}$ 为与直和交换的协变同调函子。设 $T_{p-1},T_p,T_{p+1}$ 均已定义。下列条件等价：
\begin{enumerate}[label=\alph*)]
\item $T_p$ 正合。
\item $T_{p+1},T_p$ 都右正合。
\item $T_p,T_{p-1}$ 都左正合。
\item $T_{p+1}$ 右正合且 $T_{p-1}$ 左正合。
\item 对每个 $A$-模 $M$，典范同态
\begin{equation}
 T_i(A_{\mathrm s})\otimes_A M\longrightarrow T_i(M)\tag{7.3.3.1}
\end{equation}
在 $i=p,p+1$ 时均为同构。
\item[\textup{e')}] 对每个 $M$，（7.3.3.1）在 $i=p,p+1$ 时均为满态射。
\item[\textup{f)}] 对每个 $M$，（7.3.3.1）在 $i=p$ 时为同构，且右 $A$-模 $T_p(A_{\mathrm s})$ 平坦。
\item[\textup{f')}] 对每个 $M$，（7.3.3.1）在 $i=p$ 时为满态射，且右 $A$-模 $T_p(A_{\mathrm s})$ 平坦。
\end{enumerate}
a)、b)、c)、d) 的等价性来自（7.3.1）中 a)、b) 的等价性。b)、e)、e') 的等价性来自该命题 a)、c)、d) 的等价性。最后，$T_p(A_{\mathrm s})$ 平坦意味着 $M\mapsto T_p(A_{\mathrm s})\otimes_A M$ 左正合；a)、f)、f') 的等价性仍由（7.3.1）的 a)、c)、d) 得出。

\egas{7.3.4}{推论}
设 $A$ 交换，$T_p$ 正合，且 $T_p(A)$ 为有限表示 $A$-模。则函数
\[
 x\longmapsto\dim_{k(x)}T_p(k(x))
\]
在 $X=\Spec A$ 上局部恒定；故当 $\Spec A$ 连通时恒定。

由（7.3.3，f），$T_p(A)$ 平坦，因而为有限型投射模，其相伴层局部自由（Bourbaki，\textit{Alg. comm.}，第II章，§5，第2目，定理1）。又由（7.3.3，e）有 $T_p(k(x))\simeq T_p(A)\otimes_A k(x)$。模 $T_p(A)$ 在 $x$ 处的秩局部恒定（同上），故得结论。

\egas{7.3.5}{命题}
设 $A$ 为左阿廷环，$A/\mathfrak m=k$ 为域，其中 $\mathfrak m$ 是其根基。则（7.3.3）的条件还分别等价于：

\noindent g) 典范同态 $T_i(\varepsilon):T_i(A_{\mathrm s})\to T_i(k)$ 在 $i=p,p+1$ 时均为满态射。

\noindent h) $T_p(\varepsilon)$ 为满态射，且右 $A$-模 $T_p(A_{\mathrm s})$ 平坦；等价地，它自由（Bourbaki，\textit{Alg. comm.}，第II章，§3，第2目，命题5推论2）。

若再设 $A$ 交换且 $A$-模 $T_p(k)$ 的长度为有限数 $d$，则这些条件还分别等价于：

\noindent i) 对每个有限长度 $A$-模 $M$，有
\begin{equation}
 \len T_p(M)=d\,\len M.\tag{7.3.5.1}
\end{equation}
\noindent j) 有
\begin{equation}
 \len T_p(A)=d\,\len A.\tag{7.3.5.2}
\end{equation}
g)、h) 与（7.3.3）的条件等价，直接来自（7.2.7）。为证明其余断言，使用下述引理。

\egas{7.3.5.3}{引理}
设 $\mathbf K,\mathbf K'$ 为阿贝尔范畴，$F:\mathbf K\to\mathbf K'$ 为协变加性、半正合函子，且对 $\mathbf K$ 的每个简单对象 $S$，$F(S)$ 在 $\mathbf K'$ 中具有有限长度。则每个有限长度对象 $E$ 的像 $F(E)$ 都有有限长度；对有限长度对象的短正合列 $0\to E'\to E\to E''\to0$，有
\begin{equation}
 \len F(E)\leq\len F(E')+\len F(E'').\tag{7.3.5.4}
\end{equation}
等号成立当且仅当 $0\to F(E')\to F(E)\to F(E'')\to0$ 正合。

设短正合列的两个同态为 $u,v$。由假设，$F(E')\xrightarrow{F(u)}F(E)\xrightarrow{F(v)}F(E'')$ 正合。若两端有有限长度，则两像也有有限长度，由 $\Ker F(v)=\Im F(u)$ 得
\begin{equation}
 \begin{aligned}
 \len F(E)&=\len\Im F(u)+\len\Im F(v)\\
 &\leq\len F(E')+\len F(E'').
 \end{aligned}\tag{7.3.5.5}
\end{equation}
对 $E$ 的长度归纳，已得第一断言。等号成立当且仅当
\[
 \len\Im F(u)=\len F(E'),\qquad
 \len\Im F(v)=\len F(E''),
\]
也就是 $\Ker F(u)=0$ 且 $\Coker F(v)=0$。

现在 $A$ 交换，有限长度 $M$ 的 Jordan--Hölder 列的商均同构于 $k$。由（7.3.5.4）对长度归纳可得
\begin{equation}
 \len T_p(M)\leq d\,\len M.\tag{7.3.5.6}
\end{equation}
若 $T_p$ 正合，（7.3.5.3）给出（7.3.5.1）；故（7.3.3，a）蕴含 i)。i) 蕴含 j)。最后证明下述引理。

\egas{7.3.5.7}{引理}
等式 $\len T_p(A)=d\,\len A$ 蕴含 $T_p(\varepsilon)$ 为满态射，且 $T_p(A)$ 为平坦 $A$-模。

从 $0\to\mathfrak m\to A\xrightarrow{\varepsilon}k\to0$ 出发，由（7.3.5.4）及（7.3.5.6）得
\[
 \len T_p(A)\leq\len T_p(\mathfrak m)+\len T_p(k)
 \leq d(\len\mathfrak m+\len k)=d\,\len A.
\]
等号成立时，由（7.3.5.3），列
\begin{equation}
 0\longrightarrow T_p(\mathfrak m)\longrightarrow T_p(A)
 \longrightarrow T_p(k)\longrightarrow0\tag{7.3.5.8}
\end{equation}
正合。由（7.2.7）及（7.2.5），$T_p$ 同构于 $M\mapsto N\otimes_A M$。由 Tor 长正合列，（7.3.5.8）的正合性给出 $\Tor_1^A(N,k)=0$；因此 $N=T_p(A)$ 平坦（0，10.1.3）。

\egas{7.3.6}{引理}
设 $A$ 为环，$T_\bullet:\mathbf{Ab}_A\to\mathbf{Ab}$ 为与直和交换的协变同调函子。假设 $T_p,T_{p+1}$ 已定义且 $T_p$ 左正合。则 $T_{p+1}$ 正合当且仅当右 $A$-模 $T_{p+1}(A_{\mathrm s})$ 平坦。

由（7.3.1），典范同态 $T_{p+1}(A_{\mathrm s})\otimes_A M\to T_{p+1}(M)$ 为函子同构；应用平坦模的定义即可。

\egas{7.3.7}{命题}
设 $A$ 为环，$T_\bullet:\mathbf{Ab}_A\to\mathbf{Ab}$ 为与直和交换的协变同调函子，且存在 $i_0$ 使得 $i\leq i_0$ 时 $T_i$ 正合。那么对每个 $p>i_0$，下列条件等价：
\begin{enumerate}[label=\alph*)]
\item $T_q$ 对所有 $q\leq p$ 正合；
\item $T_q(A_{\mathrm s})$ 对所有 $q\leq p$ 为平坦右 $A$-模；
\item 对每个 $A$-模 $M$ 及每个 $q\leq p+1$，典范同态 $T_q(A_{\mathrm s})\otimes_A M\to T_q(M)$ 满射。
\end{enumerate}
由 $T_{i_0}$ 正合及（7.3.6），对 $q$ 归纳可得 a)、b) 等价；a)、c) 的等价性来自（7.3.3）的 a)、e') 等价。

\egap{7.3.8}
若 $A$ 为交换环，$B$ 为 $A$-代数（不必交换），$T_\bullet:\mathbf{Ab}_A\to\mathbf{Ab}$ 为协变同调函子，则由定义（7.1.3），从 $A$ 到 $B$ 标量扩张所得的函子
\[
 T_\bullet^{(B)}=(T_i^{(B)}):\mathbf{Ab}_B\longrightarrow\mathbf{Ab}
\]
仍为同调函子。

\egas{7.3.9}{推论}
设 $T_\bullet$ 满足（7.3.7）的一般条件，并与归纳极限交换；再设 $A$ 为整环，各 $T_i(A)$ 均为有限表示 $A$-模。则对每个整数 $N$，存在 $f\in A-\{0\}$，使得 $p\leq N$ 时 $T_p^{(A_f)}$ 正合。

当 $i\leq i_0$ 时，$T_i$ 依假设正合，故 $T_i(A)$ 平坦。由（7.3.7，b），只须选择 $f$，使 $i_0\leq p\leq N$ 时 $T_p^{(A_f)}(A_f)=T_p(A_f)$ 自由。由（7.1.4），$T_p(A_f)=(T_p(A))_f$。若 $x$ 是 $\Spec A$ 的泛点，则 $(T_p(A))_x$ 为分式域上的有限维向量空间。各 $T_p(A)$ 有限表示，故存在这样的 $f$（Bourbaki，\textit{Alg. comm.}，第II章，§5，第1目，命题2推论）。

注意：若仅有有限个 $T_i$ 非零，则可选择同一个 $f\ne0$，使所有 $T_i^{(A_f)}$ 正合。

\egas{7.3.10}{推论}
设 $T_\bullet$ 满足（7.3.7）的一般条件，且与归纳极限交换；设 $A$ 交换且诺特，各 $T_i(A)$ 均为有限型 $A$-模。那么对每个整数 $N$，存在 $\Spec A$ 的稠密开子集 $U$，使对每个 $p\leq N$，函数 $x\mapsto\dim_{k(x)}T_p(k(x))$ 在 $U$ 中恒定。
\cmt{C3-E01}{“恒定”按字面不成立，一般只能断言“局部恒定”（并可在各不可约分支的稠密开子集上分别恒定）。反例及原证明中的另一处问题见核查报告；此处保留原陈述。}

设 $\mathfrak p$ 为 $A$ 的极小素理想，$B=A/\mathfrak p$ 为整环，$\Spec B$ 认作 $\Spec A$ 的一个不可约分支。我们对 $p\leq N$ 归纳，证明存在 $f_p\in B-\{0\}$，令 $B'=B_{f_p}$ 后，对 $i\leq p$，$T_i^{(B')}$ 正合且 $T_i(B')$ 为有限型 $B'$-模。当 $p\leq i_0$ 时取 $f_p=1$：因 $T_p$ 正合，$T_p(B)$ 是 $T_p(A)$ 的商，故为有限型 $A$-模，也为有限型 $B$-模。

设已得某个 $p$ 的结论。$f_p$ 是某个 $g_p\in A$ 的像。置 $A'=A_{g_p}$，则 $B'=A'/\mathfrak p'$，其中 $\mathfrak p'=\mathfrak p_{g_p}$ 为 $A'$ 的极小素理想。由 $T_i(A_{g_p})=(T_i(A))_{g_p}$，$T_i(A')$ 有限型。因此 $T_\bullet^{(A')}$ 满足相同的假设，而原文在此将 $i_0$ 换成 $p$，并归结为 $A'=A$ 且 $T_p$ 正合的情形。于是短正合列 $0\to\mathfrak p\to A\to A/\mathfrak p\to0$ 给出
\[
 T_{p+1}(A)\longrightarrow T_{p+1}(A/\mathfrak p)
 \longrightarrow T_p(\mathfrak p)\longrightarrow T_p(A).
\]
最后一个同态单射，故 $T_{p+1}(A/\mathfrak p)$ 为 $T_{p+1}(A)$ 的商，因而有限型。注意（7.3.9）的论证只对 $p\leq N$ 使用 $T_p(A)$ 有限型这一事实，所以可将它用于整环 $B$、函子 $T_\bullet^{(B)}$ 以及 $N=p+1$，完成归纳。
\cmt{C3-G02}{这里从商环上的函子正合性提升到原环上的正合性，未由先前归纳假设证明；报告将此与下段局部环等同的错误一起分列为证明缺口，不用一个排印改动掩盖。}

于是可取 $f_{\mathfrak p}\in B-\{0\}$，使 $\Spec B_{f_{\mathfrak p}}$ 为 $\Spec B$ 中处处稠密且不遇 $\Spec A$ 其他不可约分支的开集 $V$。若对 $B_{f_{\mathfrak p}}$ 已证明命题，就会在 $V$ 中得到处处稠密的开集 $W$，使所述函数在 $W$ 恒定；原文在此使用了对 $x\in W$ 的等式 $A_x=(B_{f_{\mathfrak p}})_x$。
\cmt{C3-G03}{上述局部环等式一般错误：取 $A=k[\epsilon]/(\epsilon^2)$、$\mathfrak p=(\epsilon)$，便得到 $A_x=A\ne k=B_x$。余域却可典范等同，不能将这两个说法混同。}
对每个不可约分支重复这一论证，原文便认为推论得证。于是归结为 $A$ 整；此时（7.3.9）的论证可给出 $f\in A-\{0\}$，使 $p\leq N$ 时 $T_p(A_f)$ 为有限型自由 $A_f$-模，再由（7.3.4）得出结论。

\egas{7.3.11}{命题}
设 $A$ 为交换局部环，$k$ 为其剩余域，$T_\bullet:\mathbf{Ab}_A\to\mathbf{Ab}$ 为与直和交换的协变同调函子。假设存在 $i_0$ 使 $i\leq i_0$ 时 $T_i$ 正合，并且各 $T_i(A)$ 均有限表示。则（7.3.7）中等价的 a)、b)、c) 蕴含以下两个条件；若 $A$ 还约化，则也与它们等价：

\noindent d) 对任意 $x\in\Spec A$ 及 $q\leq p$，有 $\dim_{k(x)}T_q(k(x))=\dim_kT_q(k)$。

\noindent d') 对每个不可约分支的泛点 $x_i$ 及 $q\leq p$，有 $\dim_{k(x_i)}T_q(k(x_i))=\dim_kT_q(k)$。

因 $T_q(A)$ 有限表示，（7.3.7，b）等价于 $q\leq p$ 时 $T_q(A)$ 自由（Bourbaki，\textit{Alg. comm.}，第II章，§3，第2目，命题5推论2）；条件 c) 又给出 $T_q(k(x))\simeq T_q(A)\otimes_A k(x)$，故得到 d)。d) 显然蕴含 d')。

当 $A$ 约化时，证明 d') 蕴含 a)。因 $q\leq i_0$ 时 $T_q$ 正合，可作归纳。设对 $j\leq q<p$ 已知 $T_j$ 正合，证明 $T_{q+1}(A)$ 自由。由（7.3.7，c）及（7.3.3），对任意 $M$ 有 $T_{q+1}(A)\otimes_A M\simeq T_{q+1}(M)$。分别取 $M=k(x_i)$ 和 $M=k$，由 d') 得
\[
 \dim_{k(x_i)}\bigl(T_{q+1}(A)\otimes_A k(x_i)\bigr)
 =\dim_kT_{q+1}(k).
\]
这蕴含 $T_{q+1}(A)$ 自由（Bourbaki，\textit{Alg. comm.}，第II章，§3，第2目，命题7），完成证明。

对下一目（7.4）将研究的特殊同调函子，上述结果可大为加强：将得到只涉及某一个 $T_p$ 的正合性判据。

% END continuations/iii2_73.tex


% BEGIN continuations/iii2_74.tex
\subsection{\texorpdfstring{函子 $H_\bullet(P_\bullet\otimes_A M)$ 的正合性判据}{张量复形同调函子的正合性判据}}
\source{第三卷第二部分，原印刷第53--57页}

\egap{7.4.1}
设 $A$ 为环（不必交换），$P_\bullet$ 为平坦右 $A$-模的复形。因每个函子 $M\mapsto P_k\otimes_A M$ 在 $\mathbf{Ab}_A$ 中正合，$\delta$-函子
\begin{equation}
 T_\bullet(M)=H_\bullet(P_\bullet\otimes_A M)\tag{7.4.1.1}
\end{equation}
为从 $\mathbf{Ab}_A$ 到 $\mathbf{Ab}$ 的同调函子；当 $A$ 交换时它显然 $A$-线性（7.1.2），且与归纳极限交换。

当 $A$ 交换时，对任意 $A$-代数 $B$，同调函子 $T_\bullet^{(B)}$（7.3.8）由定义满足
\begin{equation}
 T_\bullet^{(B)}(N)=H_\bullet(P_\bullet\otimes_A N_{[\rho]}),
 \tag{7.4.1.2}
\end{equation}
其中 $\rho:A\to B$ 给出 $B$ 的代数结构。又有
\[
 P_\bullet\otimes_A N_{[\rho]}
 \simeq P_\bullet\otimes_A(B\otimes_B N)
 \simeq(P_\bullet\otimes_A B)\otimes_B N,
\]
故对任意 $B$-模 $N$，
\begin{equation}
 T_\bullet^{(B)}(N)=H_\bullet(P'_\bullet\otimes_B N),
 \qquad P'_\bullet=P_\bullet\otimes_A B,\tag{7.4.1.3}
\end{equation}
且 $P'_\bullet$ 由平坦 $B$-模组成（$0_{\mathrm I}$，6.2.1）。

\egas{7.4.2}{命题}
在（7.4.1）的一般条件下，给定整数 $p\in\mathbf Z$，下列性质等价：
\begin{enumerate}[label=\alph*)]
\item $T_p$ 左正合；等价地，$T_{p+1}$ 右正合。
\item $Z'_p(P_\bullet)=\Coker(P_{p+1}\to P_p)$ 为平坦右 $A$-模。
\item 存在平坦右 $A$-模复形 $P'_\bullet$，其微分 $d'_{p+1}:P'_{p+1}\to P'_p$ 为零，且同调函子 $H_\bullet(P_\bullet\otimes_A M)$ 与 $H_\bullet(P'_\bullet\otimes_A M)$ 同构。
\end{enumerate}
依定义，有关于 $M$ 函子性的正合列
\[
 0\longrightarrow T_p(M)\longrightarrow Z'_p(P_\bullet\otimes M)
 \longrightarrow P_{p-1}\otimes M.
\]
由张量积的右正合性，$Z'_p(P_\bullet\otimes M)=Z'_p(P_\bullet)\otimes M$。对任一同态 $f:M\to N$，得到两行正合的交换图
\begin{equation}
 \begin{tikzcd}[column sep=small]
 0\arrow[r]&T_p(M)\arrow[r]\arrow[d,"u"]&Z'_p(P_\bullet)\otimes M\arrow[r]\arrow[d,"v"]&P_{p-1}\otimes M\arrow[d,"w"]\\
 0\arrow[r]&T_p(N)\arrow[r]&Z'_p(P_\bullet)\otimes N\arrow[r]&P_{p-1}\otimes N.
 \end{tikzcd}\tag{7.4.2.1}
\end{equation}
若 $f$ 单射，则 $w$ 单射，因为 $P_{p-1}$ 平坦；若 $T_p$ 左正合，则 $u$ 也单射，因而 $v$ 单射，说明 $Z'_p(P_\bullet)$ 平坦。反之，若该模平坦，则任一单射 $f$ 都给出单射 $v$，由图知 $u$ 单射；$T_p$ 已经半正合，所以左正合。a)、b) 因而等价。

c) 蕴含 a)：当 $d'_{p+1}=0$ 时，对短正合列 $0\to M'\to M\to M''\to0$，由定义（M，IV，1），正合列
\[
 H_{p+1}(P'_\bullet\otimes M'')\xrightarrow{\delta}
 H_p(P'_\bullet\otimes M')\longrightarrow H_p(P'_\bullet\otimes M)
\]
中的边缘同态 $\delta$ 为零，因此 $T_p$ 左正合。

反过来证明 b) 蕴含 c)。记 $Z_{p+1}(P_\bullet)=\Ker(P_{p+1}\to P_p)$。原文在此写出正合列
\[
 0\longrightarrow Z_{p+1}(P_\bullet)\longrightarrow P_{p+1}
 \longrightarrow Z'_p(P_\bullet)\longrightarrow0,
\]
据 $P_{p+1},Z'_p(P_\bullet)$ 平坦推出 $Z_{p+1}(P_\bullet)$ 平坦（$0_{\mathrm I}$，6.1.2）。
\cmt{C3-G04}{此列最右项应为 $B_p(P_\bullet)$，不是余核 $Z'_p(P_\bullet)$；$B_p$ 的平坦性须先由本证明末段的短正合列得出。下文关于任意次数循环和边界与张量积交换的说法也过强。报告给出同一简单复形的反例，并列明只在所需次数使用平坦性的修补。}
取
\[
 P'_i=P_i\ (i\ne p,p+1),\qquad
 P'_p=Z'_p(P_\bullet),\quad P'_{p+1}=Z_{p+1}(P_\bullet).
\]
微分 $d'_i$ 在 $i\ne p,p+1$ 时取 $P_\bullet$ 的微分（在相应子模或商模间理解）；在 $i=p+1$ 时取零；在 $i=p$ 时取由 $d_p$ 通过商诱导的 $Z'_p(P_\bullet)\to P_{p-1}$。原文接着由 $P_i$ 平坦写出
\[
 \begin{gathered}
 Z'_i(P_\bullet\otimes M)=Z'_i(P_\bullet)\otimes M,\\
 Z_i(P_\bullet\otimes M)=Z_i(P_\bullet)\otimes M,\qquad
 B_i(P_\bullet\otimes M)=B_i(P_\bullet)\otimes M,
 \end{gathered}
\]
其中 $B_i(P_\bullet)=\Im(P_{i+1}\to P_i)$；并由此对每个 $M$ 和每个 $i$ 得到函子性同构 $H_i(P_\bullet\otimes M)\simeq H_i(P'_\bullet\otimes M)$。它与 $\delta$ 相容这一点，可由 $\delta$ 的定义（M，IV，1）核验。

注意，（7.4.2）的条件也蕴含 $B_p(P_\bullet)$ 平坦，因为有短正合列
\[
 0\longrightarrow B_p(P_\bullet)\longrightarrow P_p
 \longrightarrow Z'_p(P_\bullet)\longrightarrow0,
\]
且后两模平坦（$0_{\mathrm I}$，6.1.2）。

\egas{7.4.3}{推论}
设 $A$ 为维数1的正则诺特环（也就是说，为 Dedekind 环的乘积；$0_{\mathrm{IV}}$，17.1.3、17.3.7；例如主理想环）。则 $T_p$ 左正合当且仅当 $T_p(A)$ 为平坦 $A$-模；$T_p$ 正合当且仅当 $T_p(A)$ 与 $T_{p-1}(A)$ 都平坦。

回顾：Dedekind 环上的模平坦当且仅当无挠（$0_{\mathrm I}$，6.3.3、6.3.4），故这里平坦模的任一子模也平坦。第二断言来自第一断言，因为 $T_p$ 正合等价于 $T_p,T_{p-1}$ 均左正合。为证明第一断言，考察短正合列
\[
 0\longrightarrow H_p(P_\bullet)\longrightarrow Z'_p(P_\bullet)
 \longrightarrow B_{p-1}(P_\bullet)\longrightarrow0.
\]
$B_{p-1}$ 为平坦模 $P_{p-1}$ 的子模，故平坦；因此 $H_p(P_\bullet)$ 平坦与 $Z'_p(P_\bullet)$ 平坦等价（$0_{\mathrm I}$，6.1.2）。

（7.4.2）最重要的应用如下。

\egas{7.4.4}{命题}
设 $A$ 为诺特环，$P_\bullet$ 为平坦 $A$-模复形。假设各 $P_i$ 有限型；或者假设各 $H_i(P_\bullet)$ 有限型且存在 $i_0$ 使 $i<i_0$ 时 $H_i(P_\bullet)=0$。令 $T_\bullet$ 为（7.4.1.1）定义的同调函子。那么，使 $(T_p)_y$（7.1.4）右正合（分别左正合、正合）的点 $y\in\Spec A$ 所成集合 $U$ 是开集。

在第二种假设下，可将 $P_\bullet$ 换成有限型自由 $A$-模复形 $P'_\bullet$，使 $H_\bullet(P'_\bullet\otimes_A M)$ 与 $T_\bullet(M)$ 作为 $\delta$-函子同构（0，11.9.3）。所以总可归结为第一种假设；此时各 $Z'_i(P_\bullet)$ 有限型。由（7.4.2，a），只须证明左正合的情形。若 $x\in U$，因局部化正合，有
\[
 (Z'_p(P_\bullet))_x=Z'_p((P_\bullet)_x).
\]
由（7.4.1.3）及（7.4.2，b），该模在 $A_x$ 上平坦，因有限型且 $A_x$ 局部诺特，故自由（Bourbaki，\textit{Alg. comm.}，第II章，§3，第2目，命题5推论2）。由同书第II章§5第1目命题2推论，存在 $f\in A$，$x\in D(f)$，使 $(Z'_p(P_\bullet))_f$ 在 $A_f$ 上自由；于是每个 $y\in D(f)$ 的局部化也自由，再用（7.4.2，b）即可。

\egas{7.4.5}{推论}
在（7.4.4）的假设下，再设 $A$ 为整环。则使 $(T_p)_x$ 正合的 $x\in\Spec A$ 所成集合 $U$ 为非空开集。

只须说明泛点处正合。此时 $A_x$ 是域，域上任意加性函子都是正合函子，故得结论。

\egas{7.4.6}{命题}
在（7.4.4）的一般假设下，（7.4.2）的 a)、b)、c) 还等价于：

\noindent d) 存在 $A$-模 $Q$ 及函子性同构
\begin{equation}
 T_p(M)\simeq\Hom_A(Q,M).\tag{7.4.6.1}
\end{equation}
此外，$Q$ 由该性质确定至唯一同构，并且有限型。

$Q$ 的唯一性是表示函子的对象之唯一性的特殊情形（0，8.1.5）。右边显然左正合。反之，当 $T_p$ 左正合时，如（7.4.4）可先归结为各 $P_i$ 平坦且有限型；因 $A$ 诺特，它们是有限型投射模（Bourbaki，\textit{Alg. comm.}，第II章，§5，第2目，定理1推论）。记 $P_i^\vee$ 为其对偶，仍为有限型投射模；$P_i$ 典范同构于其双对偶，且
\[
 P_i\otimes_A M\longrightarrow\Hom_A(P_i^\vee,M)
\]
为典范同构（Bourbaki，\textit{Alg.}，第II章，第3版，§4，第2目，命题2）。由（7.4.2，c），可设 $d_{p+1}=0$，故有正合列
\[
 0\longrightarrow T_p(M)\longrightarrow P_p\otimes M
 \xrightarrow{d_p\otimes1}P_{p-1}\otimes M.
\]
置 $Q'=\Ker d_p$，有 $0\to Q'\to P_p\xrightarrow{d_p}P_{p-1}$ 正合。原文转置此列，写出
\[
 P_{p-1}^\vee\xrightarrow{{}^td_p}P_p^\vee
 \longrightarrow (Q')^\vee\longrightarrow0,
\]
并令 $Q=(Q')^\vee=\Coker({}^td_p)$。
\cmt{C3-G05}{转置一般不能得到这里所写的右正合列，$(\Ker d_p)^\vee$ 也不一定等于转置映射的余核。有效的定义是直接取 $Q=\Coker({}^td_p)$；随后的 Hom 正合列即可证明命题，不需要错误的对偶核等式。报告以 $d_p$ 为乘以 $t$ 给出反例。}
要证明 $Q$ 满足要求，考察正合列
\[
 0\longrightarrow\Hom_A(Q,M)\longrightarrow\Hom_A(P_p^\vee,M)
 \xrightarrow{\Hom({}^td_p,1)}\Hom_A(P_{p-1}^\vee,M).
\]
在上述典范同构下，最右同态与 $d_p\otimes1$ 对应，因而得到（7.4.6.1）。又 $Q$ 是 $P_p^\vee$ 的商，故有限型。

\egas{7.4.7}{命题}
在（7.4.4）的一般假设下，对每个有限型 $A$-模 $M$：
\begin{enumerate}[label=(\roman*)]
\item 各 $T_p(M)$ 都是有限型 $A$-模；
\item 对 $A$ 的每个理想 $\mathfrak m$，典范同态
\begin{equation}
 \widehat{T_p(M)}\longrightarrow
 \varprojlim_n T_p\bigl(M\otimes_A(A/\mathfrak m^{n+1})\bigr)
 \tag{7.4.7.1}
\end{equation}
为双射；左边为 $T_p(M)$ 关于 $\mathfrak m$-预进制拓扑的分离完备化。
\end{enumerate}
如（7.4.4），先归结为各 $P_i$ 有限型。因 $A$ 诺特，$P_i\otimes M$ 的子模有限型，（i）随即成立。（ii）来自更一般的下述引理。

\egas{7.4.7.2}{引理}
设 $A$ 为诺特环，$u:E\to F$ 为有限型 $A$-模同态。对每个有限型 $A$-模 $M$，置
\[
 K(M)=\Ker(u\otimes1_M),\qquad C(M)=\Coker(u\otimes1_M).
\]
则对 $A$ 的每个理想 $\mathfrak m$，典范同态
\begin{equation}
 \widehat{K(M)}\longrightarrow\varprojlim_nK(M_n),\qquad
 \widehat{C(M)}\longrightarrow\varprojlim_nC(M_n)
 \tag{7.4.7.3}
\end{equation}
为双射，其中 $M_n=M\otimes_A(A/\mathfrak m^{n+1})=M/\mathfrak m^{n+1}M$。

$E\otimes M,F\otimes M$ 有限型，而完备化函子在有限型 $A$-模范畴中正合（$0_{\mathrm I}$，7.3.3），故 $\widehat{K(M)},\widehat{C(M)}$ 分别为 $\widehat{u\otimes1}:\widehat{E\otimes M}\to\widehat{F\otimes M}$ 的核与余核。射影极限函子的左正合性给出 $\widehat{K(M)}=\varprojlim K(M_n)$；张量积的右正合性则给出 $C(M_n)=C(M)\otimes_A(A/\mathfrak m^{n+1})$，从而依定义 $\widehat{C(M)}=\varprojlim C(M_n)$。

\egas{7.4.8}{注记}
结合（6.10.5）和（6.10.6），可见在附加平坦性假设下，（7.4.7）重新给出了（4.1.7.1）是同构这一事实，也就是固有态射“第一比较定理”的主要内容；而且不再只适用于一个凝聚模，还适用于这类模的复形。若能得到同时包括（7.4.7）与（4.1.7.1）的命题，将很有意义。注意，当各 $P_i$ 有限型时，（7.4.7）的证明并未用到它们平坦；应考察：若不再假设各 $P_i$ 平坦或有限型，而只假设各 $H_i(P_\bullet)$ 有限型且 $i<i_0$ 时为零，结论是否仍成立？这时能否用有限型 $A$-模复形 $P'_\bullet$ 代替 $P_\bullet$，使已不再是同调函子的 $H_\bullet(P_\bullet\otimes M)$ 与 $H_\bullet(P'_\bullet\otimes M)$ 仍同构？

% END continuations/iii2_74.tex


% BEGIN continuations/iii2_75_76.tex
\subsection{局部诺特环的情形}
\source{第三卷第二部分，原印刷第58--60页}

\egap{7.5.1}
设 $A$ 是局部诺特环，$\mathfrak m$ 为其极大理想。对每个 $A$-模 $M$，以 $\widehat M$ 表示关于 $\mathfrak m$-预进制拓扑的分离完备化，即
\[
 \widehat M\simeq\varprojlim_n\bigl(M\otimes_A(A/\mathfrak m^{n+1})\bigr)
 =\varprojlim_n M/\mathfrak m^{n+1}M.
\]
设 $T:\mathbf{Ab}_A\to\mathbf{Ab}$ 为协变加性函子。典范同态（7.2.3.1）
\[
 T(M)\otimes_A(A/\mathfrak m^{n+1})\longrightarrow
 T\bigl(M\otimes_A(A/\mathfrak m^{n+1})\bigr)
\]
组成 $A$-线性同态的射影系，取极限便得到关于 $M$ 函子性的 $A$-线性同态
\begin{equation}
 \widehat{T(M)}\longrightarrow\varprojlim_n T(M_n),
 \qquad M_n=M\otimes_A A_n,\quad A_n=A/\mathfrak m^{n+1}.
 \tag{7.5.1.1}
\end{equation}

\egas{7.5.2}{命题}
设 $A$ 为极大理想 $\mathfrak m$、剩余域 $k=A/\mathfrak m$ 的局部诺特环。设 $T:\mathbf{Ab}_A\to\mathbf{Ab}$ 为协变加性、半正合且与归纳极限交换的函子。假设对每个有限型 $A$-模 $M$，$T(M)$ 有限型且（7.5.1.1）为同构。那么下列性质等价：
\begin{enumerate}[label=\alph*)]
\item $T$ 右正合。
\item 对每个 $n$，$N\mapsto T(N)$ 在有限型 $A_n$-模范畴中右正合；也就是说，$T$ 在有限长度 $A$-模范畴中右正合。
\item 典范同态 $T(\varepsilon):T(A)\to T(k)$ 满射。
\item 对一切充分大的 $n$，典范同态 $T(A_n)\to T(k)$ 满射。
\end{enumerate}
a) 蕴含 b)。证明 b) 蕴含 a)：若 $u:M\to N$ 为满射，须证 $T(u)$ 满射。因 $T$ 与归纳极限交换，且模的滤过归纳极限正合，可归结为 $M,N$ 有限型。此时 $T(M),T(N)$ 有限型，且 $A$ 局部诺特，故只须证明其完备化 $\widehat{T(u)}$ 满射（$0_{\mathrm I}$，7.3.5、6.4.1）。依假设，它是同态系 $T(u\otimes1_{A_n}):T(M_n)\to T(N_n)$ 的射影极限。由 b)，各同态满射；又 $T(M_n)$ 为阿廷环 $A_n$ 上的有限型模，所以极限同态满射（0，13.1.2、13.2.2）。

a) 蕴含 c)；因 $T(\varepsilon)$ 分解为 $T(A)\to T(A_n)\to T(k)$，c) 蕴含 d)。最后，由（7.2.7）及 $T$ 的半正合性，b)、d) 等价。

\egas{7.5.3}{推论}
在（7.5.2）的一般假设下，若 $T(k)=0$，则对每个 $A$-模 $M$，$T(M)=0$。

$k$ 是唯一的简单 $A$-模，由（7.3.5.4），每个有限长度模 $E$ 都满足 $T(E)=0$。当 $M$ 有限型时，$\widehat{T(M)}\simeq\varprojlim T(M_n)=0$；因 $T(M)$ 有限型，它同构于自身完备化的子模（$0_{\mathrm I}$，7.3.5），故为零。任意 $M$ 是其有限型子模 $N_\lambda$ 的滤过归纳极限，所以 $T(M)=\varinjlim T(N_\lambda)=0$。

\egas{7.5.4}{命题}
设 $A$ 为极大理想 $\mathfrak m$、剩余域 $k$ 的局部诺特环，$T_\bullet:\mathbf{Ab}_A\to\mathbf{Ab}$ 为与归纳极限交换的同调函子。假设对每个 $i$ 和每个有限型 $A$-模 $M$，$T_i(M)$ 有限型且典范同态 $\widehat{T_i(M)}\to\varprojlim T_i(M_n)$ 为双射。给定整数 $p$，下列条件等价：
\begin{enumerate}[label=\alph*)]
\item $T_p$ 正合。
\item $T_p$ 右正合，且 $T_p(A)$ 自由。
\item $T_{p+1}(A)\to T_{p+1}(k)$ 和 $T_p(A)\to T_p(k)$ 均满射。
\item 对每个 $n$，$T_{p+1}(A_n)\to T_{p+1}(k)$ 和 $T_p(A_n)\to T_p(k)$ 均满射。
\item 对每个 $n$，$N\mapsto T_p(N)$ 在有限型 $A_n$-模范畴中正合。
\end{enumerate}
由（7.3.3），a) 等价于 $T_{p+1},T_p$ 均右正合。限制到 $A_n$-模范畴后仍为同调函子，故如（7.3.1）的论证，e) 等价于这两个函子在有限型 $A_n$-模范畴中均右正合。由（7.5.2）得 a)、e) 等价，也得到 a)、c)、d) 等价。有限型平坦模在局部环上自由（Bourbaki，\textit{Alg. comm.}，第II章，§3，第2目，命题5推论2），因此 a)、b) 的等价性由（7.3.1）和（7.3.3）得到。

\egas{7.5.5}{推论}
在（7.5.4）的一般条件下：
\begin{enumerate}[label=(\roman*)]
\item 若 $T_p(k)=0$，则 $T_p=0$，$T_{p+1}$ 右正合，$T_{p-1}$ 左正合。
\item 若 $T_{p-1}(k)=T_{p+1}(k)=0$，则 $T_p$ 正合，典范同态 $T_p(A)\otimes_A M\to T_p(M)$ 为双射，且 $T_p(A)$ 为自由 $A$-模。
\end{enumerate}
因 $T_p$ 半正合，（i）由（7.5.3）及同调函子定义得到。（ii）的前两断言来自（i）及（7.3.3）；自由性由（7.5.4）得到。

\egas{7.5.6}{推论}
在（7.5.4）的一般假设下，再设 $A$ 为离散赋值环。
\begin{enumerate}[label=(\roman*)]
\item $T_p$ 右正合当且仅当 $T_{p-1}(A)$ 自由。
\item $T_p$ 正合当且仅当 $T_p(A),T_{p-1}(A)$ 均自由。
\end{enumerate}
由（7.3.3），（i）蕴含（ii）。为证（i），取极大理想的生成元 $f$，即一致化元。有限型 $A$-模 $M$ 自由（等价于平坦）当且仅当乘以 $f$ 的同态单射，因为这里也等价于 $M$ 无挠（$0_{\mathrm I}$，6.3.4）。短正合列 $0\to A\xrightarrow{f}A\to k\to0$ 给出
\[
 T_p(A)\longrightarrow T_p(k)\longrightarrow T_{p-1}(A)
 \xrightarrow{f}T_{p-1}(A).
\]
所以 $T_{p-1}(A)$ 自由当且仅当 $T_p(A)\to T_p(k)$ 满射，再用（7.5.2）即可。

\egas{7.5.7}{注记}
由（7.4.7），（7.4.4）的一般假设使（7.4.1.1）定义的同调函子满足（7.5.4）的一般条件。在此情形，（7.5.6）已包含于（7.4.3）。

\subsection{正合性的下降；半连续性定理与 Grauert 正合性判据}
\source{第三卷第二部分，原印刷第60--65页}

\egas{7.6.1}{命题}
在（7.4.1）的条件下，设 $B$ 为交换 $A$-代数。若 $T_p$ 右正合（分别左正合、正合），则 $T_p^{(B)}$ 亦然；当 $B$ 在 $A$ 上忠实平坦时，逆命题成立。

第一断言是（7.1.3）的特殊情形。反之，先假设 $B$ 平坦。对任一 $A$-模 $M$，有
\[
 H_\bullet(P_\bullet\otimes_A(M\otimes_A B))
 \simeq H_\bullet(P_\bullet\otimes_A M)\otimes_A B,
\]
亦即典范同构
\begin{equation}
 T_p(M)\otimes_A B\simeq T_p^{(B)}(M_{(B)}).\tag{7.6.1.1}
\end{equation}
若 $T_p^{(B)}$ 具有所述正合性，因 $M\mapsto M_{(B)}$ 正合，左边关于 $M$ 也具有该正合性；若 $B$ 忠实平坦，即推出 $T_p$ 同样正合（$0_{\mathrm I}$，6.4.1）。

\egas{7.6.2}{命题}
在（7.4.1）的条件下，再设 $A$ 为约化诺特环，各 $P_i$ 有限型。则 $T_p$ 右正合（分别左正合、正合）当且仅当对每个为离散赋值环的 $A$-代数 $B$，$T_p^{(B)}$ 具有同一正合性。

由（7.3.1）、（7.3.3），只须考虑右正合；且只须证明充分性（7.6.1）。由（7.4.2），只须说明 $Z'_{p-1}(P_\bullet)$ 平坦。它为有限型模，由（0，10.2.8）的判据，只须对每个离散赋值 $A$-代数 $B$，证明 $Z'_{p-1}(P_\bullet)\otimes_A B$ 平坦。又
\[
 Z'_{p-1}(P_\bullet)\otimes_A B=Z'_{p-1}(P_\bullet\otimes_A B),
\]
$P_\bullet\otimes_A B$ 由平坦 $B$-模组成（$0_{\mathrm I}$，6.2.1），且对 $B$-模 $N$ 有 $T_p^{(B)}(N)=H_p((P_\bullet\otimes_A B)\otimes_B N)$。再用（7.4.2），所需平坦性正是 $T_p^{(B)}$ 右正合。

下面仔细考察离散赋值环的情形。

\egas{7.6.3}{命题}
在（7.4.1）的条件下，设 $A$ 为维数1的正则诺特环；也就是说，$A$ 诺特且每个局部环 $A_x$ 是域或离散赋值环。那么对任一整数 $p$ 和任一 $A$-模 $M$，有关于 $M$ 函子性的典范短正合列
\begin{equation}
 0\longrightarrow T_p(A)\otimes_A M\xrightarrow{t_M}T_p(M)
 \longrightarrow\Tor_1^A(T_{p-1}(A),M)\longrightarrow0.
 \tag{7.6.3.1}
\end{equation}
以下省略 $H_p(P_\bullet),B_p(P_\bullet),Z_p(P_\bullet),Z'_p(P_\bullet)$ 中的 $P_\bullet$。有三条短正合列
\[
\begin{gathered}
 0\to H_p\to Z'_p\to B_{p-1}\to0,\\
 0\to B_{p-1}\to Z_{p-1}\to H_{p-1}\to0,\\
 0\to Z_{p-1}\to P_{p-1}\to B_{p-2}\to0.
\end{gathered}
\]
因 $P_{p-1},P_{p-2}$ 平坦，它们的子模 $B_{p-1},Z_{p-1},B_{p-2}$ 也平坦：在每个 $A_x$ 上，平坦与无挠等价。张量 $M$ 后得到
\begin{align}
 0=\Tor_1^A(B_{p-1},M)&\longrightarrow H_p\otimes M\longrightarrow Z'_p\otimes M
 \xrightarrow{u}B_{p-1}\otimes M\longrightarrow0,\tag{7.6.3.2}\\
 0=\Tor_1^A(Z_{p-1},M)&\longrightarrow\Tor_1^A(H_{p-1},M)
 \longrightarrow B_{p-1}\otimes M\xrightarrow{v}Z_{p-1}\otimes M,\tag{7.6.3.3}\\
 0=\Tor_1^A(B_{p-2},M)&\longrightarrow Z_{p-1}\otimes M
 \xrightarrow{w}P_{p-1}\otimes M.\tag{7.6.3.4}
\end{align}
依定义，$T_p(M)=\Ker(d_p\otimes1)/\Im(d_{p+1}\otimes1)$，也就是诱导同态 $Z'_p\otimes M\to P_{p-1}\otimes M$ 的核。该同态分解为
\[
 Z'_p\otimes M\xrightarrow{u}B_{p-1}\otimes M
 \xrightarrow{v}Z_{p-1}\otimes M\xrightarrow{w}P_{p-1}\otimes M.
\]
由（7.6.3.4），$w$ 单射，故有 $0\to\Ker u\to T_p(M)\to\Ker v\to0$。结合（7.6.3.2）、（7.6.3.3）以及 $H_p=T_p(A)$，正是（7.6.3.1）。

\egas{7.6.4}{注记}
（i）$H_\bullet(P_\bullet\otimes_A M)$ 是把 $M$ 置于次数0后得到的双复形的同调；因此由（6.3.6）、（6.3.2），它是初始项为
\[
 E^2_{pq}=\Tor_p^A(H_q(P_\bullet),M)=\Tor_p^A(T_q(A),M)
\]
的正则谱序列的终项。关于 $A$ 的假设使 $p\geq2$ 时任意模 $E,F$ 满足 $\Tor_p^A(E,F)=0$（$0_{\mathrm{IV}}$，17.2.2）；因此有短正合列（M，XV）
\[
 0\longrightarrow E^2_{0p}\longrightarrow H_p(P_\bullet\otimes M)
 \longrightarrow E^2_{1,p-1}\longrightarrow0,
\]
这就是（7.6.3.1）。

（ii）结合（7.3.1），（7.6.3.1）以（7.4.3）为特殊情形。

\egas{7.6.5}{推论}
在（7.4.1）的条件下，设 $A$ 为离散赋值环，分式域为 $K$，剩余域为 $k$，各 $T_i(A)$ 有限型。那么
\begin{equation}
 \dim_kT_p(k)\geq\dim_k(T_p(A)\otimes_A k)
 \geq\rk_A T_p(A)=\dim_KT_p(K).\tag{7.6.5.1}
\end{equation}
两个极端项相等，当且仅当 $T_p$ 正合；也等价于 $T_p(A),T_{p-1}(A)$ 均自由。

在（7.6.3.1）中取 $M=k$，得到
\[
 \dim_kT_p(k)=\dim_k(T_p(A)\otimes_A k)
 +\dim_k\Tor_1^A(T_{p-1}(A),k).
\]
又因 $T_p(A)$ 为局部整环上的有限型模，有（Bourbaki，\textit{Alg. comm.}，第II章，§3，第2目，命题4推论1）
\begin{equation}
 \dim_k(T_p(A)\otimes_A k)\geq\rk_A T_p(A)
 =\dim_K(T_p(A)\otimes_A K),\tag{7.6.5.2}
\end{equation}
且等号成立当且仅当 $T_p(A)$ 自由（同处，命题7）。由于 $K$ 平坦，有 $T_p(A)\otimes_A K=H_p(P_\bullet)\otimes_A K=H_p(P_\bullet\otimes_A K)=T_p(K)$。因此得到（7.6.5.1）。极端相等还要求 $\Tor_1^A(T_{p-1}(A),k)=0$，这等价于 $T_{p-1}(A)$ 自由（0，10.1.3）。最后，有限型模的平坦与自由在这里等价，再用（7.4.3）即可。

\egap{7.6.6}
仍沿用（7.4.1）的假设；对每个 $x\in\Spec A$，置
\begin{equation}
 d_p(x)=d_p^T(x)=\dim_{k(x)}T_p(k(x)).\tag{7.6.6.1}
\end{equation}

\egas{7.6.7}{引理}
设 $\varphi:A\to A'$ 为环同态，$f={}^a\varphi:\Spec A'\to\Spec A$ 为相应映射（I，1.2.1）。若 $T'_\bullet=T_\bullet^{(A')}$，则
\begin{equation}
 d_p^{T'}=d_p^T\circ f.\tag{7.6.7.1}
\end{equation}
事实上，取 $x'\in\Spec A'$，置 $x=f(x')$。由于 $k(x')$ 在 $k(x)$ 上平坦，
\[
 \begin{aligned}
 H_\bullet(P_\bullet\otimes_A k(x'))
 &=H_\bullet((P_\bullet\otimes_A k(x))\otimes_{k(x)}k(x'))\\
 &=H_\bullet(P_\bullet\otimes_A k(x))\otimes_{k(x)}k(x'),
 \end{aligned}
\]
故有（7.6.7.1）。

\egas{7.6.8}{引理}
若 $A$ 诺特且 $P_\bullet$ 各项有限型，则 $x\mapsto d_p(x)$ 是 $\Spec A$ 上的可构造函数。

须证：对 $X=\Spec A$ 的每个不可约闭集 $Y$，存在 $Y$ 的非空开集 $U$，使 $d_p$ 在 $U$ 恒定（0，9.2.2）。写 $Y=\Spec(A/\mathfrak a)$，其中 $A/\mathfrak a$ 约化。由（7.6.7），可归结为 $Y=X$ 且 $A$ 诺特整；再用（7.4.5）。

\egas{7.6.9}{定理}
设 $A$ 为诺特环，$P_\bullet$ 为有限型平坦 $A$-模复形，$T_\bullet(M)=H_\bullet(P_\bullet\otimes_A M)$。置 $d_p(x)=\dim_{k(x)}T_p(k(x))$。则：
\begin{enumerate}[label=(\roman*)]
\item $d_p$ 在 $\Spec A$ 上可构造且上半连续。
\item 若 $T_p$ 正合，则 $d_p$ 连续，因而局部恒定；当 $A$ 约化时，逆命题成立。
\end{enumerate}
（i）的可构造性已经由（7.6.8）证明。由（0，9.2.4），为证上半连续，只须对 $x$ 的一般化 $x'\ne x$ 证明 $d_p(x')\leq d_p(x)$。存在离散赋值环 $B$ 及态射 $f:\Spec B\to\Spec A$，将闭点 $a$ 送到 $x$，泛点 $b$ 送到 $x'$（II，7.1.9）。由（7.6.7.1），归结为 $d_p^{T^{(B)}}(a)\geq d_p^{T^{(B)}}(b)$，这正是（7.6.5.1）。
\orig{将（i）的证明归结为离散赋值环情形这一思想，由 Hironaka 告知我们。}

（ii）的第一断言已经在（7.3.4）证明。反之，使用赋值判据（7.6.2）和（7.6.7.1），可归结为 $A$ 是离散赋值环。此时 $\Spec A$ 只有两个点，$d_p$ 恒定便由（7.6.5）给出 $T_p$ 正合。

\egas{7.6.10}{推论}
设 $A$ 为诺特环，$\mathfrak p_i$（$1\leq i\leq r$）为其极小素理想，$k_i=k(\mathfrak p_i)$。
\begin{enumerate}[label=(\roman*)]
\item 对任意 $x\in\Spec A$，至少存在一个 $i$ 使
\begin{equation}
 d_p^T(x)\geq\dim_{k_i}T_p(k_i).\tag{7.6.10.1}
\end{equation}
特别地，若 $A$ 整、$K$ 为其分式域，则
\begin{equation}
 d_p^T(x)\geq\dim_KT_p(K)\qquad(x\in\Spec A).\tag{7.6.10.2}
\end{equation}
\item 再设 $A$ 局部且约化，$k$ 为其剩余域。则 $T_p$ 正合当且仅当
\begin{equation}
 \dim_kT_p(k)=\dim_{k_i}T_p(k_i)\qquad(1\leq i\leq r).
 \tag{7.6.10.3}
\end{equation}
\end{enumerate}
（i）由上半连续性及每个 $x$ 的邻域含有某个极小素点得到。另一方面，若 $A$ 局部，极大理想在 $\Spec A$ 中唯一的开邻域就是全空间，因此 $d_p(x)\leq\dim_kT_p(k)$。结合（i），（7.6.10.3）使 $d_p$ 恒定，于是（7.6.9，（ii））给出正合性；必要性同样由该定理得到。

\egas{7.6.11}{注记}
可以问，（7.6.9，（i））能否加强为对每个 $x$ 都有
\begin{equation}
 \dim_{k(x)}T_p(k(x))\geq
 \dim_{k(x)}(T_p(A)\otimes_A k(x)).\tag{7.6.11.1}
\end{equation}
当 $A$ 为离散赋值环且 $x$ 是其极大理想时，这确由（7.6.5）成立。只考虑局部诺特环 $A$，极大理想为 $\mathfrak m$、剩余域为 $k$。以下条件等价：

\noindent a) 对每个有限型平坦 $A$-模复形 $P_\bullet$ 及每个整数 $p$，有
\begin{equation}
 \dim_kT_p(k)\geq\dim_k(T_p(A)\otimes_A k).\tag{7.6.11.2}
\end{equation}
\noindent b) 对每个有限型 $A$-模 $M$，有
\begin{equation}
 \dim_k(M\otimes_A k)\geq\dim_k(M^\vee\otimes_A k).
 \tag{7.6.11.3}
\end{equation}
\noindent c) 对每个有限型 $A$-模 $N$，有
\begin{equation}
 \dim_k\Tor_1^A(N,k)\geq\dim_k\Tor_2^A(N,k).
 \tag{7.6.11.4}
\end{equation}
由维数移位（M，V，7.2），c) 也等价于对所有 $i\geq1$ 有
\begin{equation}
 \dim_k\Tor_i^A(N,k)\geq\dim_k\Tor_{i+1}^A(N,k).
 \tag{7.6.11.5}
\end{equation}

简述证明。为证 a) 蕴含 b)，取正合列 $L_1\xrightarrow{d}L_0\to M\to0$，其中 $L_0,L_1$ 有限型自由；将 a) 用于双项复形 $P_1\xrightarrow{{}^td}P_0$，其中 $P_0=L_1^\vee$，$P_1=L_0^\vee$，其余项为零。此时 $T_1(A)=M^\vee$，且
\[
 T_1(k)=\Hom_A(M,k)=\Hom_k(M/\mathfrak mM,k),
\]
所以它与 $M\otimes_A k$ 有相同维数。为证 b) 蕴含 c)，先证明下述引理。

\egas{7.6.11.6}{引理}
给定平坦 $A$-模复形 $\cdots\to0\to P_1\xrightarrow{d}P_0\to0\to\cdots$，有正合列
\begin{equation}
 0\to\Tor_2^A(Z'_0,k)\to T_1(A)\otimes_A k\to T_1(k)
 \to\Tor_1^A(Z'_0,k)\to0.\tag{7.6.11.7}
\end{equation}
事实上，从 $0\to Z_1\to P_1\to B_0\to0$ 得到
\[
 0\to\Tor_1^A(B_0,k)\to Z_1\otimes k\to P_1\otimes k
 \to B_0\otimes k\to0.
\]
从 $0\to B_0\to Z_0\to Z'_0\to0$，因 $Z_0=P_0$ 平坦，得到 $\Tor_1^A(B_0,k)\simeq\Tor_2^A(Z'_0,k)$。又 $Z_1=T_1(A)$。同态 $d\otimes1$ 分解为
\[
 P_1\otimes k\xrightarrow{u}B_0\otimes k\xrightarrow{v}Z_0\otimes k=P_0\otimes k.
\]
若 $R=\Ker v$，则 $T_1(k)=u^{-1}(R)$，且由 $u$ 满射，$R=u(T_1(k))$；由 $v$ 的定义，$R=\Tor_1^A(Z'_0,k)$。这就得到（7.6.11.7）。

现在由 b) 推出 c)。取有限型自由表示 $L_1\xrightarrow{d}L_0\to N\to0$，令 $T$ 对应于由 $L_1,L_0$ 构成的双项复形。因它们等于其双对偶，若 $M=\Coker({}^td)$，则 $T_1(A)=\Ker d=M^\vee$；而 $M\otimes_A k=\Coker({}^td\otimes1)$ 与 $\Ker(d\otimes1)=T_1(k)$ 有相同的 $k$-维数。b) 因而给出
\[
 \dim_k(T_1(A)\otimes_A k)\leq\dim_kT_1(k).
\]
此时 $Z'_0=N$，所以（7.6.11.7）给出（7.6.11.4）。

最后，由 c) 推出 a)。将（7.6.11.6）用于双项 $P_p\to P_{p-1}$，对其余核应用 c)，得到 $\dim_kR\geq\dim_kS$，其中
\[
 R=\Ker(P_p\otimes k\to P_{p-1}\otimes k),\qquad S=Z_p\otimes k.
\]
将 $d_{p+1}$ 分解为 $P_{p+1}\xrightarrow{u}Z_p\xrightarrow{j}P_p$；则
\[
 \Im(d_{p+1}\otimes1)=(j\otimes1)(\Im(u\otimes1)),\qquad
 T_p(A)\otimes k=(Z_p\otimes k)/\Im(u\otimes1),
\]
而 $T_p(k)=R/\Im(d_{p+1}\otimes1)$，故得到（7.6.11.2）。

设局部环 $A$ 正则且维数为 $n$。已知 [17] 给出 $\Tor_i^A(k,k)\simeq\bigwedge^i(\mathfrak m/\mathfrak m^2)$；因此当 $n\geq4$ 时，取 $N=k$ 即不满足（7.6.11.4）。另一方面，若局部整环 $A$ 上每个有限型自反模都自由（二维正则局部环具有此性质），则（7.6.11.3）成立：原文使用有限型模的对偶 $M^\vee$ 自反这一事实，从而它自由，于是若 $K$ 为分式域，
\[
 \dim_k(M^\vee\otimes k)=\dim_K(M^\vee\otimes K)=\dim_K(M\otimes K).
\]
又 $M\otimes_A k$ 的每组 $k$-基都来自 $M$ 的一组生成元（Bourbaki，\textit{Alg. comm.}，第II章，§3，第2目，命题4推论2），所以 $\dim_K(M\otimes K)\leq\dim_k(M\otimes k)$，从而得出所述断言。

% END continuations/iii2_75_76.tex


% BEGIN continuations/iii2_77.tex
\subsection{对固有态射的应用：I．交换性质}
\source{第三卷第二部分，原印刷第65--72页}
接下来的三个目，主要是把前面各目的结果译写为预概形态射的语言。

\egap{7.7.1}
设 $f:X\to Y$ 为拟紧且分离的预概形态射，$\PP_\bullet$ 为拟凝聚 $\OO_X$-模复形，微分次数为 $-1$，并设各 $\PP_i$ 在 $Y$ 上平坦（$0_{\mathrm I}$，6.7.1）。考虑拟凝聚 $\OO_Y$-模范畴中的 $\delta$-函子 $\M\mapsto\TT_\bullet(\M)$，也记作 $\TT_\bullet(\PP_\bullet,\M)$，定义为
\begin{equation}
 \TT_n(\PP_\bullet,\M)=\TT_n(\M)
 =\cH^{-n}(f,\PP^\bullet\otimes_{\OO_Y}\M),\qquad n\in\mathbf Z.
 \tag{7.7.1.1}
\end{equation}
这里 $\PP^j=\PP_{-j}$，微分次数因而为 $+1$。由（6.2.3），其值为拟凝聚 $\OO_Y$-模；由（6.2.6），它是关于 $\M$ 的同调函子。

\egap{7.7.2}
设 $g:Y'\to Y$ 为态射，置 $X'=X_{(Y')}=X\times_Y Y'$，$f'=f_{(Y')}:X'\to Y'$，它仍拟紧且分离。置 $\PP'_\bullet=\PP_\bullet\otimes_{\OO_Y}\OO_{Y'}$；由（I，9.1.12）和（$0_{\mathrm I}$，6.2.1），其各项拟凝聚且在 $Y'$ 上平坦。采用同样的次数约定，置
\begin{equation}
 \begin{aligned}
 \TT'_\bullet(\M')
 &=\TT_\bullet(f',\PP'_\bullet\otimes_{\OO_{Y'}}\M')\\
 &=\cH^{-\bullet}(f',\PP'^\bullet\otimes_{\OO_{Y'}}\M').
 \end{aligned}\tag{7.7.2.1}
\end{equation}
这是关于拟凝聚 $\OO_{Y'}$-模 $\M'$ 的同调函子。当 $Y'=\Spec A'$ 时，以 $\TT^{A'}$ 代替 $\TT'$；对每个 $A'$-模 $M'$，置
\[
 T^{A'}(M')=\Gamma(Y',\TT^{A'}(\widetilde{M'})),
 \qquad\TT^{A'}(\widetilde{M'})=\widetilde{T^{A'}(M')}.
\]
$T^{A'}$ 为取值于 $A'$-模范畴的同调函子。

若 $Y=\Spec A$ 也仿射，则 $T^{A'}$ 与从 $T^A$ 标量扩张到 $A'$ 所得函子相同（7.1.3）。事实上，若 $g':X'\to X$ 为相应投影，则它仿射（II，1.6.2）；仿射开覆盖 $\mathfrak U$ 的逆像 $\mathfrak U'=g'^{-1}(\mathfrak U)$ 仍为仿射开覆盖。由（6.2.2），只须验证
\[
 C^\bullet(\mathfrak U,\PP_\bullet\otimes_{\OO_Y}g_*(\M'))
 =C^\bullet(\mathfrak U',\PP'_\bullet\otimes_{\OO_{Y'}}\M').
\]
最终归结为对仿射开集 $U\subseteq X$ 及 $U'=g'^{-1}(U)$，
\[
 \Gamma(U,\PP_\bullet\otimes_{\OO_Y}g_*(\M'))
 =\Gamma(U',\PP'_\bullet\otimes_{\OO_{Y'}}\M'),
\]
这由（I，1.3、3.2）直接得到。特别地，对 $Y$ 的开集 $U$ 及任意拟凝聚 $\M$，
\begin{equation}
 \TT^U_\bullet(\M|U)=(\TT_\bullet(\M))|U.\tag{7.7.2.2}
\end{equation}

\egap{7.7.3}
对每个拟凝聚 $\OO_Y$-模 $\M$，有关于 $\M$ 函子性的典范同态
\begin{equation}
 \TT_p(\OO_Y)\otimes_{\OO_Y}\M\longrightarrow\TT_p(\M).
 \tag{7.7.3.1}
\end{equation}
当 $Y$ 仿射时，这已经在（7.2.2）定义。一般情形通过限制相容性得到：对仿射开集 $V\subseteq U\subseteq Y$，令
\[
 A_W=\TT_p^W(\OO_W)\otimes_{\OO_W}(\M|W),\qquad
 B_W=\TT_p^W(\M|W)\quad(W=U,V),
\]
则在与限制层作典范识别后，图
\[
 \begin{tikzcd}
 A_U\arrow[r]\arrow[d,"\text{限制}"']&B_U\arrow[d,"\text{限制}"]\\
 A_V\arrow[r]&B_V
 \end{tikzcd}
\]
由（7.2.3.3）交换。

对任一态射 $g:Y'\to Y$，还有典范同态
\begin{equation}
 \TT_p(\OO_Y)\otimes_{\OO_Y}\OO_{Y'}\longrightarrow
 \TT'_p(\OO_{Y'}).\tag{7.7.3.2}
\end{equation}
它是（6.7.11.2）关于终项的特殊情形：取 $S=Y$、$u_1=f$、$v_1=1_Y$，并把第二个复形置为次数0的模。若 $Y=\Spec A$、$Y'=\Spec A'$ 均仿射，则（7.7.3.2）与（7.2.2）中典范 $A'$-线性同态
\[
 T_p^A(A)\otimes_A A'\longrightarrow T_p^{A'}(A')=T_p^A(A')
\]
相伴。这可从（6.7.11）看出，因为可在那里选取 $\mathfrak U'^{(i)}=u_i^{-1}(\mathfrak U^{(i)})$。

\egap{7.7.4}
当 $f$ 固有、$Y=\Spec A$ 仿射诺特、$\PP_\bullet$ 下有界且由凝聚、$Y$-平坦的模组成时，由（6.10.5），在同构意义下可写
\[
 \TT_\bullet(\M)=\HH_\bullet(\LL_\bullet\otimes_{\OO_Y}\M),
 \qquad\LL_\bullet=\widetilde{L_\bullet},
\]
其中 $L_\bullet$ 为下有界的有限型自由 $A$-模复形。因此 $T^A$ 正属于（7.4）、（7.6）详细研究的类型。现在译写该研究的结果。

\egas{7.7.5}{定理}
设 $Y$ 为局部诺特预概形，$(U_\alpha)$ 为其仿射开覆盖，$f:X\to Y$ 固有，$\PP_\bullet$ 为下有界的凝聚、$Y$-平坦 $\OO_X$-模复形。由（7.7.1.1）定义的同调函子具有以下性质。

\textbf{I．半连续性。}函数
\begin{equation}
 y\longmapsto d_p(y)=\dim_{k(y)}T_p^{k(y)}(k(y))\tag{7.7.5.1}
\end{equation}
上半连续。
\orig{本定理的一个特殊情形已经出现在 Chow--Igusa 的札记 [3] 中。对解析空间，在较特殊的假设下，Kodaira--Spencer 已证明半连续性（\textit{On the variations of almost-complex structures}，载 \textit{Algebraic Geometry and Topology, A Symposium in honor of S. Lefschetz}，Princeton 数学丛书第12册，第139--150页，Princeton，1957年）。一般版本由 Grauert [31] 证明。}

\textbf{II．交换性质。}给定整数 $p$，下列条件等价：
\begin{enumerate}[label=\alph*)]
\item $\TT_p$ 右正合。
\item[\textup{a')}] $\TT_p(\M)$ 同构于某个形如 $\N\otimes_{\OO_Y}\M$ 的函子；此时 $\N$ 必同构于 $\TT_p(\OO_Y)=\cH^{-p}(f,\PP^\bullet)$。
\item[\textup{a'')}] 典范函子同态（7.7.3.1）为同构。
\item[\textup{b)}] $\TT_{p-1}$ 左正合。
\item[\textup{b')}] 存在 $\OO_Y$-模 $\mathscr Q$（必凝聚，并确定至唯一同构）及函子同构
\begin{equation}
 \TT_{p-1}(\M)\simeq\shHom_{\OO_Y}(\mathscr Q,\M).
 \tag{7.7.5.2}
\end{equation}
\item[\textup{c)}] 若 $A_\alpha$ 为 $U_\alpha$ 的环，则每个 $T_p^{A_\alpha}$ 都右正合。
\item[\textup{d)}] 对每个态射 $g:Y'\to Y$，典范同态
\begin{equation}
 \TT_p(\OO_Y)\otimes_{\OO_Y}\OO_{Y'}\longrightarrow\TT'_p(\OO_{Y'})
 \tag{7.7.5.3}
\end{equation}
为同构。
\end{enumerate}
半连续性是 $Y$ 上的局部性质，故来自（7.7.4）和（7.6.9）。a'') 蕴含 a')，a') 蕴含 a)。$Y$ 仿射时，由（7.7.4），a)、a'')、b)、b') 的等价性已经由（7.3.1）、（7.4.6）证明。

为推广到一般情形，先证明 a) 与 c) 等价，这也说明 a) 在 $Y$ 上局部。相同论证适用于 a'')、b)。c) 蕴含 a)；反之，只须对任一仿射开集 $U\subseteq Y$ 及拟凝聚 $\OO_U$-模短正合列
\[
 0\longrightarrow\F'\longrightarrow\F\longrightarrow\F''\longrightarrow0
\]
找出 $Y$ 上的拟凝聚模短正合列，其限制为该列。由 $Y$ 局部诺特及（I，9.4.2），先把 $\F$ 延拓为 $\G$，把 $\F'$ 延拓为 $\G$ 的拟凝聚子模 $\G'$，再取 $\G''=\G/\G'$ 即可。

对一般 $Y$，证明 b)、b') 等价。仿射情形中的表示模 $\mathscr Q$ 确定至唯一同构；因此对仿射 $Y$ 的仿射开集 $U$，其表示模就是 $\mathscr Q|U$。一般地，在每个仿射开集 $U$ 上取凝聚表示模 $\mathscr Q_U$。若仿射开集 $V\subseteq U$，则 $\mathscr Q_U|V$ 与 $\mathscr Q_V$ 典范同构；故可粘合，得到满足（7.7.5.2）的唯一表示模。

最后证明 a)、d) 等价。二者在 $Y$ 上均局部，且 d) 在 $Y'$ 上也局部。故设 $Y=\Spec A$、$Y'=\Spec A'$。已知 $T^{A'}$ 由 $T^A$ 标量扩张得到，a') 便给出（7.7.5.3）为同构。反之，设 d) 成立。对任一 $A$-模 $M$，取平方零扩张 $A'=A\oplus M$，乘法为
\[
 (a_1,m_1)(a_2,m_2)=(a_1a_2,a_1m_2+a_2m_1).
\]
此时 $T_p(A')=T_p(A)\oplus T_p(M)$；（7.7.5.3）为双射，推出典范同态 $T_p(A)\otimes_A M\to T_p(M)$ 也为双射，即 a'')。

\egas{7.7.6}{定理}
设 $Y$ 局部诺特，$f:X\to Y$ 固有，$\F$ 为凝聚且 $Y$-平坦的 $\OO_X$-模。则存在凝聚 $\OO_Y$-模 $\mathscr Q$（确定至唯一同构），使对拟凝聚 $\OO_Y$-模 $\M$ 有函子同构
\begin{align}
 f_*(\F\otimes_{\OO_Y}\M)&\simeq\shHom_{\OO_Y}(\mathscr Q,\M),\tag{7.7.6.1}\\
 \Gamma(X,\F\otimes_{\OO_Y}\M)&\simeq\Hom_{\OO_Y}(\mathscr Q,\M).
 \tag{7.7.6.2}
\end{align}
因 $\M\mapsto\F\otimes_{\OO_Y}\M$ 正合（$0_{\mathrm I}$，6.7.4），$f_*$ 左正合，复合函子左正合。取 $p=1$，应用（7.7.5，b）与 b') 的等价性即可。

\egas{7.7.7}{推论}
设 $Y$ 局部诺特，$f:X\to Y$ 固有，$\F,\F'$ 凝聚且 $Y$-平坦，$u:\F\to\F'$ 为同态。对拟凝聚 $\OO_Y$-模 $\M$，置
\[
 \begin{aligned}
 \TT(\M)&=\Ker\bigl(f_*(\F\otimes\M)\to f_*(\F'\otimes\M)\bigr),\\
 T(\M)&=\Gamma(Y,\TT(\M))\\
 &=\Ker\bigl(\Gamma(X,\F\otimes\M)\to\Gamma(X,\F'\otimes\M)\bigr).
 \end{aligned}
\]
则存在凝聚 $\OO_Y$-模 $\B$（确定至唯一同构），使
\begin{align}
 \TT(\M)&\simeq\shHom_{\OO_Y}(\B,\M),\tag{7.7.7.1}\\
 T(\M)&\simeq\Hom_{\OO_Y}(\B,\M).\tag{7.7.7.2}
\end{align}
只须证明（7.7.7.2）：它在仿射情形给出（7.7.7.1），再像（7.7.5）的 b)、b') 那样，利用表示对象的唯一性（0，8.1.8）粘合。由（7.7.6），存在凝聚 $\mathscr Q,\mathscr Q'$，分别表示 $\Gamma(X,\F\otimes-)$、$\Gamma(X,\F'\otimes-)$。$u$ 定义这两个函子之间的态射，由（0，8.1.4），唯一对应于 $v:\mathscr Q'\to\mathscr Q$，使图
\[
 \begin{tikzcd}
 \Gamma(X,\F\otimes\M)\arrow[r]\arrow[d,"\sim"']&\Gamma(X,\F'\otimes\M)\arrow[d,"\sim"]\\
 \Hom_{\OO_Y}(\mathscr Q,\M)\arrow[r]&\Hom_{\OO_Y}(\mathscr Q',\M)
 \end{tikzcd}
\]
交换。反变函子 $\N\mapsto\Hom_{\OO_Y}(\N,\M)$ 左正合，故取 $\B=\Coker v$ 即可。

\egas{7.7.8}{推论}
设 $X,Y,f$ 满足（7.7.6）的假设，$\F,\G$ 为凝聚 $\OO_X$-模，满足：（i）$\F$ 在 $Y$ 上平坦；（ii）$\G$ 同构于有限型局部自由模之间某同态 $v:\E_1\to\E_0$ 的余核。置
\[
 \TT(\M)=f_*\shHom_{\OO_X}(\G,\F\otimes_{\OO_Y}\M),\qquad
 T(\M)=\Hom_{\OO_X}(\G,\F\otimes_{\OO_Y}\M).
\]
则存在凝聚 $\OO_Y$-模 $\N$（确定至唯一同构），使
\begin{align}
 \TT(\M)&\simeq\shHom_{\OO_Y}(\N,\M),\tag{7.7.8.1}\\
 T(\M)&\simeq\Hom_{\OO_Y}(\N,\M).\tag{7.7.8.2}
\end{align}
由（$0_{\mathrm I}$，5.4.2.1），对 $i=0,1$ 有函子性同构
\[
 \begin{aligned}
 \shHom_{\OO_X}(\E_i,\F\otimes\M)
 &\simeq\E_i^\vee\otimes_{\OO_X}(\F\otimes\M)\\
 &\simeq(\E_i^\vee\otimes_{\OO_X}\F)\otimes\M
 \simeq\shHom_{\OO_X}(\E_i,\F)\otimes\M.
 \end{aligned}
\]
令 $\F_i=\shHom_{\OO_X}(\E_i,\F)$；它们凝聚（$0_{\mathrm I}$，5.3.5）且 $Y$-平坦（同章5.4.2）。置 $u=\shHom(v,1_\F):\F_0\to\F_1$。由 Hom 函子的左正合性，
\[
 \shHom_{\OO_X}(\G,\F\otimes\M)
 \simeq\Ker(\F_0\otimes\M\to\F_1\otimes\M).
\]
再因 $f_*$ 左正合，$\TT(\M)$ 同构于 $\Ker(f_*(\F_0\otimes\M)\to f_*(\F_1\otimes\M))$，应用（7.7.7）即可。

\egas{7.7.9}{注记}
（i）在（7.7.6）—（7.7.8）中，$\mathscr Q,\B,\N$ 的构造与基变换交换。例如采用（7.7.2）的记号，在（7.7.6）情形，对拟凝聚 $\OO_{Y'}$-模 $\M'$ 有
\[
 f'_*(\F'\otimes_{\OO_{Y'}}\M')
 \simeq\shHom_{\OO_{Y'}}(g^*\mathscr Q,\M').
\]
由（7.7.2）的说明，归结为伴随同构
\[
 \Hom_{\OO_Y}(\mathscr Q,g_*\M')
 =\Hom_{\OO_{Y'}}(g^*\mathscr Q,\M')
\]
（$0_{\mathrm I}$，4.4.3.1）。同样，在（7.7.7）的基变换中用 $g^*\B$ 代替 $\B$；在（7.7.8）中用 $g^*\N$ 代替 $\N$，因为基变换后的 $\E'_1\to\E'_0\to\G'\to0$ 仍正合，并且
\[
 \shHom_{\OO_{X'}}(\E'_i,\F\otimes\OO_{Y'})
 \simeq\E_i^\vee\otimes_{\OO_X}\F\otimes_{\OO_Y}\OO_{Y'}.
\]

（ii）若存在 $Y$-丰沛可逆 $\OO_X$-模，例如 $Y$ 仿射且 $f$ 射影，则任意凝聚 $\G$ 均满足（7.7.8，ii）。结合（II，5.5.1），由（II，2.7.10）可取有限型局部自由 $\E_0$ 满射到 $\G$；其核凝聚，对核再次应用同一结果，即得所需表示。

（iii）我们将在第V章证明，（7.7.8）的限制性假设（ii）是多余的。

\egas{7.7.10}{命题（交换性质的局部判据）}
在（7.7.5）的一般条件下，给定 $y\in Y$ 和整数 $p$，以下性质等价：
\begin{enumerate}[label=\alph*)]
\item $T_p^{\OO_y}$ 右正合。
\item 典范同态 $T_p^{\OO_y}(\OO_y)\to T_p^{\OO_y}(k(y))$ 满射。
\item 对每个整数 $n$，典范同态 $T_p^{\OO_y}(\OO_y/\mathfrak m_y^{n+1})\to T_p^{\OO_y}(k(y))$ 满射。
\end{enumerate}
此外，满足这些条件的点构成使 $\TT_p^U$ 右正合的最大开集 $U\subseteq Y$。

结合（7.7.4），等价性来自（7.4.7）及（7.5.2）；开性来自（7.4.4）。反之，若 $\TT_p^U$ 右正合，则由（7.7.5，c）及（7.6.1），$y\in U$ 时 $T_p^{\OO_y}$ 右正合。

\egas{7.7.11}{推论}
若 $\TT_p$ 右正合（分别左正合），则任意基变换 $g:Y'\to Y$ 后的 $\TT'_p$ 也如此；若 $g$ 忠实平坦，则逆命题成立。

第一断言来自（7.6.1）以及正合性对 $Y,Y'$ 局部（7.7.5，c、b）。反之，由（7.7.10），只须在每个 $y\in Y$ 证明 $T_p^{\OO_y}$ 的相应正合性。因 $g$ 忠实平坦，存在 $y'$ 映到 $y$，且 $\OO_{y'}$ 在 $\OO_y$ 上忠实平坦；结论来自（7.6.1）。

\egas{7.7.12}{注记}
（i）在（7.7.4）的假设下，再设 $\PP_\bullet$ 为有限复形。取 $X$ 的有限仿射开覆盖 $\mathfrak U$，则 $C^\bullet(\mathfrak U,\PP^\bullet\otimes\M)$ 为有限双复形。准确地说，存在与 $\M$ 无关的有限指标集 $E$，使所有 $(h,k)\notin E$ 的项都为零。因此存在 $i_0$，使 $i\geq i_0$ 时对任意拟凝聚 $\M$ 都有 $\TT_i(\M)=0$。这些函子正合，故相应的 $Z'_{i_0}(L_\bullet)$ 为有限型平坦、因而有限型投射 $A$-模。取
\[
 L'_i=L_i\ (i<i_0),\qquad L'_{i_0}=Z'_{i_0}(L_\bullet),
 \qquad L'_i=0\ (i>i_0),
\]
并令 $\LL'_\bullet=\widetilde{L'_\bullet}$。显然 $i<i_0-1$ 时两复形张量 $\M$ 后的同调相同，$i\geq i_0$ 时也相同（均为零）；$i=i_0-1$ 时，由 $\Im(Z'_{i_0}\otimes_A M\to L_{i_0-1}\otimes_A M)=\Im(L_{i_0}\otimes_A M\to L_{i_0-1}\otimes_A M)$，同调亦相同。故在（7.7.4）中，若只要求 $\LL_i$ 有限型局部自由，即对应有限型投射 $A$-模，就可令 $\LL_\bullet$ 为有限复形。

特别地，若 $\PP_\bullet$ 只在次数0有一非零项 $\F$，则 $\TT_p(\M)=R^{-p}f_*(\F\otimes\M)$，可令 $i>0$ 时 $\LL_i=0$；这时宜采用上同调记号，以 $\TT^{-p}$ 代替 $\TT_p$。

（ii）在（7.7.5）中若去掉各 $\PP_i$ 的 $Y$-平坦性，结论在改定义为
\begin{equation}
 \TT_p(\M)=\lTor_p^Y(f,1_Y;\PP_\bullet,\M)\tag{7.7.12.1}
\end{equation}
后仍成立。由（6.7.9），$\lTor_p^Y(f,1_Y;\PP_\bullet,\OO_Y)$ 凝聚；结合（6.10.1），（6.10.5）的证明不变，仍在仿射底空间上得到 $\TT_\bullet(\M)=\HH_\bullet(\LL_\bullet\otimes\M)$，其中 $\LL_\bullet$ 与有限型自由模复形相伴。这就证明了断言。

% END continuations/iii2_77.tex


% BEGIN continuations/iii2_78_79.tex
\subsection{对固有态射的应用：II．上同调平坦性判据}
\source{第三卷第二部分，原印刷第72--76页}

\egas{7.8.1}{定义}
设 $f:X\to Y$ 为拟紧且分离的预概形态射，$\PP_\bullet$ 为拟凝聚且 $Y$-平坦的 $\OO_X$-模复形，$\TT_\bullet$ 为（7.7.1.1）定义的同调函子，$y\in Y$。若 $y$ 在 $Y$ 中存在开邻域 $U$，使 $\M\mapsto\TT_p^U(\M)$ 正合，则称 $\PP_\bullet$ 在 $y$ 处关于 $Y$ 在维数 $p$ 上\emph{同调平坦}，或在维数 $-p$ 上\emph{上同调平坦}。若在每个 $y\in Y$ 都如此，则称它关于 $Y$ 在维数 $p$ 上同调平坦，或在维数 $-p$ 上上同调平坦。

若对所有维数 $p$ 都同调平坦（分别在 $y$ 处同调平坦），则简称关于 $Y$ 同调平坦（分别在 $y$ 处同调平坦），也称关于 $Y$ 上同调平坦（分别在 $y$ 处上同调平坦）。

\egap{7.8.2}
由定义，同调平坦性在 $Y$ 上是局部性质。若 $Y$ 局部诺特，或 $Y$ 是概形，则关于 $Y$ 在维数 $p$ 上同调平坦，当且仅当 $\TT_p$ 正合。局部诺特情形已在（7.7.5）的证明中说明；概形情形用同一论证，将（I，9.4.2）用于拟紧概形中的仿射开集。

\egas{7.8.3}{命题}
沿用（7.8.1）的记号与假设。以下条件等价：
\begin{enumerate}[label=\alph*)]
\item $\PP_\bullet$ 在 $y$ 处关于 $Y$ 在维数 $p$ 上同调平坦。
\item 存在 $y$ 的开邻域 $U$，使 $\TT_p^U,\TT_{p+1}^U$ 均右正合。
\item 存在 $y$ 的开邻域 $U$，使 $\TT_p^U,\TT_{p-1}^U$ 均左正合。
\item 存在 $y$ 的开邻域 $U$，使 $\TT_{p+1}^U$ 右正合且 $\TT_{p-1}^U$ 左正合。
\end{enumerate}
考虑 $Y$ 仿射时 $\TT_\bullet$ 的解释，这只是（7.3.3）的部分结论的译写。

\egas{7.8.4}{命题}
设 $Y$ 局部诺特，$f:X\to Y$ 固有，$\PP_\bullet$ 下有界且各项凝聚、$Y$-平坦，$\TT_\bullet$ 由（7.7.1.1）定义。对每个 $y\in Y$，以下条件等价：
\begin{enumerate}[label=\alph*)]
\item $\PP_\bullet$ 在 $y$ 处关于 $Y$ 在维数 $p$ 上同调平坦。
\item $T_p^{\OO_y}$ 正合。
\item 存在 $n_0$，使 $n\geq n_0$ 时
\begin{equation}
 \begin{aligned}
 \len_{\OO_y}T_p^{\OO_y}(\OO_y/\mathfrak m_y^{n+1})
 &=\dim_{k(y)}T_p^{\OO_y}(k(y))\\
 &\quad\cdot\len_{\OO_y}(\OO_y/\mathfrak m_y^{n+1}).
 \end{aligned}\tag{7.8.4.1}
\end{equation}
\item 存在 $y$ 的开邻域 $U$，使 $\cH^{-p}(f,\PP^\bullet)|U$ 同构于某个 $(\OO_Y|U)^r$，并对任意拟凝聚 $\OO_U$-模 $\M$，典范同态
\begin{equation}
 \cH^{-p}(f,\PP^\bullet)|U\otimes_{\OO_U}\M
 \longrightarrow\cH^{-p}(f_U,\PP^\bullet|f^{-1}(U)\otimes_{\OO_U}\M)
 \tag{7.8.4.2}
\end{equation}
为双射。
\end{enumerate}
这些条件成立时，还具有以下性质：

\noindent e) $y$ 有一个邻域，使其中 $z\mapsto d_p(z)$（7.7.5.1）恒定。

若 $Y$ 在 $y$ 处约化（$0_{\mathrm I}$，4.1.4），则 e) 也与上述条件等价。

由（7.3.3），b) 等价于 $T_p^{\OO_y},T_{p+1}^{\OO_y}$ 均右正合。a)、b) 等价来自（7.7.10）及（7.8.3）。$\OO_y/\mathfrak m_y^{n+1}$ 阿廷，相关模由（7.7.4）有限型，因而有限长度；b)、c) 的等价性来自（7.7.10）及（7.3.5.7）。a) 蕴含 e)，以及局部环约化时二者等价，来自（7.6.9）。a) 又由（7.7.5）给出（7.8.4.2）双射；同时 $\cH^{-p}(f,\PP^\bullet)_y$ 平坦（7.3.3，f），且有限型，故自由（0，10.1.3）。该层凝聚，遂在 $y$ 的某邻域局部自由（$0_{\mathrm I}$，5.2.7）。反之，由（7.7.2.2）中函子的定义，d) 蕴含 a)。

\egas{7.8.5}{命题}
在（7.8.4）的假设下，以下条件等价：
\begin{enumerate}[label=\alph*)]
\item $\PP_\bullet$ 关于 $Y$ 在每个维数 $i\leq p$ 上同调平坦。
\item $i\leq p+1$ 时，各 $\TT_i$ 均右正合。
\item $i\geq-p$ 时，各 $\cH^i(f,\PP^\bullet)$ 均局部自由。
\end{enumerate}
a)、b) 等价由（7.8.3）得到，a) 蕴含 c) 由（7.8.4）得到。反之，设 c) 成立。由（7.7.4），表示复形在充分小的次数为零，同样相应 $\TT_i$ 为零。每个 $y$ 有仿射邻域 $U=\Spec A$，使 $i\leq p$ 时 $T_i^A(A)$ 自由；由（7.3.7），$\TT_i^U$ 在这些次数均正合。

我们主要将这些判据用于 $\PP_\bullet$ 仅在次数0有一个凝聚且 $Y$-平坦模 $\F$ 的情形。此时 $\TT_p(\M)=R^{-p}f_*(\F\otimes_{\OO_Y}\M)$。

\egas{7.8.6}{命题}
设 $Y$ 为局部诺特预概形，$f:X\to Y$ 固有且平坦，$y\in Y$，$X_y=f^{-1}(y)=X\otimes_Y k(y)$。若 $R=\Gamma(X_y,\OO_{X_y})$ 是可分 $k(y)$-代数（Bourbaki，\textit{Alg.}，第VIII章，§7，第5目），即有限个 $k(y)$ 的有限次可分扩张的直积，则 $\OO_X$ 在 $y$ 处关于 $Y$ 在维数0上上同调平坦。

由（7.8.4），可设 $Y=\Spec A$，$A=\OO_y$。因 $f$ 平坦，$\TT_1=0$，所以 $T_0^A$ 左正合；只须证明它右正合。由（7.7.10，c），可进一步设 $A$ 阿廷。取 $k(y)$ 的有限扩张 $k'$，使 $R\otimes_{k(y)}k'$ 为有限个 $k'$ 的直积。存在局部同态 $A\to A'$，使 $A'$ 为有限自由 $A$-代数，且其剩余域同构于 $k'$（0，10.3.2）。由（7.6.1），可归结为 $R$ 已是 $m$ 个 $k(y)$ 的直积。使用以下初等引理。

\egas{7.8.6.1}{引理}
局部赋环空间 $Z$ 连通，当且仅当 $\Gamma(Z,\OO_Z)$ 不是两个非零环的直积。

若 $Z=Z_1\sqcup Z_2$ 为两个非空开集的并，则全局截面环就是两个非零截面环的直积。反之，环有这种直积分解，当且仅当存在异于0、1的幂等元 $s$。对每个 $z\in Z$，其芽 $s_z$ 在局部环中为0或1。$s_z=0$ 的点形成开集；而 $s_z=1$ 时 $s(z)\ne0$，这类点也形成开集（$0_{\mathrm I}$，5.5.2），从而得到断言。

由引理，$X_y$ 恰有 $m$ 个连通分支 $X'_i$，每个满足 $\Gamma(X'_i,\OO_{X'_i})=k(y)$。因为 $A$ 局部阿廷，其谱仅含一点，$X,X_y$ 具有相同底空间；故 $X$ 有 $m$ 个连通分支 $X_i$，且 $X'_i=X_i\otimes_Yk(y)$。于是归结为 $R=k(y)$。由（7.7.10，b），只须证明 $\Gamma(X,\OO_X)\to\Gamma(X_y,\OO_{X_y})$ 满射；复合
\[
 A=\Gamma(Y,\OO_Y)\longrightarrow\Gamma(X,\OO_X)
 \longrightarrow\Gamma(X_y,\OO_{X_y})=k(y)
\]
已满射，故所需结论成立。

\egas{7.8.7}{推论}
在（7.8.6）的假设下，存在 $y$ 的开邻域 $U$，使得：
\begin{enumerate}[label=(\roman*)]
\item $f_*(\OO_X)|U\simeq(\OO_Y|U)^m$；
\item 对任意 $z\in U$，典范同态
\[
 (f_*(\OO_X))_z\otimes_{\OO_z}k(z)
 \longrightarrow\Gamma(X_z,\OO_{X_z})
\]
为双射。
\end{enumerate}
（i）由（7.8.6）、（7.8.4）得到；（ii）由适当 $U$ 上 $\TT_0^U$ 正合及（7.7.5.3）得到。

\egas{7.8.8}{推论}
在（7.8.6）的条件下，若还有 $\Gamma(X_y,\OO_{X_y})=k(y)$，则存在 $y$ 的开邻域 $U$，使典范同态 $\OO_Y|U\to f_*(\OO_X)|U$ 为双射。

由（7.8.7，（ii）），其（i）中的整数 $m$ 必为1。

\egas{7.8.9}{推论}
在（7.8.6）的假设下，存在 $y$ 的开邻域 $U$、凝聚 $\OO_U$-模 $\mathscr Q$（确定至唯一同构），以及关于拟凝聚 $\OO_U$-模 $\M$ 的函子同构
\begin{equation}
 R^1(f_U)_*(f_U^*\M)\simeq\shHom_{\OO_U}(\mathscr Q,\M).
 \tag{7.8.9.1}
\end{equation}
事实上，适当 $U$ 上 $\TT_0^U$ 正合。令 $p=0$、$\PP_\bullet$ 仅在次数0为 $\OO_X$，应用（7.7.5，a）与 b') 的等价性即可。

\egas{7.8.10}{注记}
（i）在（7.8.6）的条件下，考虑 Stein 分解（4.3.3）
\[
 X\xrightarrow{f'}Y'\xrightarrow{g}Y,
 \qquad Y'=\Spec(f_*\OO_X).
\]
有限态射 $g$ 满足 $g_*\OO_{Y'}=f_*\OO_X$ 在 $y$ 附近局部自由，其在 $y$ 处的纤维是一个可分 $k(y)$-代数的谱（II，1.5.1）。将在第IV章推出：存在 $y$ 的开邻域 $U$，使对限制 $g^{-1}(U)\to U$，每个纤维 $g^{-1}(z)$ 都是可分 $k(z)$-代数的谱；这将称为 $U$ 的 étale 覆盖。因此由（7.8.7，（ii）），（7.8.6）对 $y$ 所作假设也在 $y$ 的某邻域所有点成立。

（ii）将在第V章看到：即使 $X$ 在 $Y$ 上射影，并且在 $Y$ 上“简单”（此性质将在第IV章定义），（7.8.9）的 $\mathscr Q$ 也未必局部自由。换言之，$\OO_X$ 未必在维数1上关于 $Y$ 在 $y$ 处上同调平坦。第V章将把 $\mathscr Q$ 解释为 $X/Y$ 的 Picard 概形沿单位截面的1-微分层。

\subsection{对固有态射的应用：III．Euler--Poincaré 特征数与 Hilbert 多项式的不变性}
\source{第三卷第二部分，原印刷第76--80页}

\egap{7.9.1}
设 $A$ 为环，$M$ 为有限型投射 $A$-模。回顾（Bourbaki，\textit{Alg. comm.}，第II章，§5，第2目），这等价于其在 $X=\Spec A$ 上的相伴层 $\widetilde M$ 有限型局部自由。对每个 $\mathfrak p\in\Spec A$，称自由 $A_{\mathfrak p}$-模 $M_{\mathfrak p}$ 的秩为 $M$ 在 $\mathfrak p$ 处的秩，记 $\rk_{\mathfrak p}M$。它也是 $\widetilde M$ 在该点的秩。因此
\begin{equation}
 \rk_{\mathfrak p}M=\rk_{A_{\mathfrak p}}M_{\mathfrak p}
 =\dim_{k(\mathfrak p)}(M\otimes_A k(\mathfrak p)).\tag{7.9.1.1}
\end{equation}

\egas{7.9.2}{命题}
设 $P_\bullet$ 为有限型投射 $A$-模的有限复形，$T_\bullet(M)=H_\bullet(P_\bullet\otimes_A M)$。则对每个 $\mathfrak p\in\Spec A$，有
\begin{equation}
 \sum_i(-1)^i\dim_{k(\mathfrak p)}T_i(k(\mathfrak p))
 =\sum_i(-1)^i\rk_{\mathfrak p}P_i.\tag{7.9.2.1}
\end{equation}
由定义，$T_i(k(\mathfrak p))=H_i(P_\bullet\otimes_A k(\mathfrak p))$。结合（7.9.1.1），该公式就是有限维向量空间的有限复形在取同调时 Euler--Poincaré 特征数不变（0，11.10.2）。
\cmt{C3-T05}{底本证明将本式引为（7.9.1.2）；此处应为（7.9.2.1）。本节其他引文和符号误植集中见核查报告。}

\egas{7.9.3}{推论}
函数 $\mathfrak p\mapsto\sum_i(-1)^i\dim_{k(\mathfrak p)}T_i(k(\mathfrak p))$ 在 $\Spec A$ 上局部恒定。

\egas{7.9.4}{定理}
设 $Y$ 局部诺特，$f:X\to Y$ 固有，$\PP_\bullet$ 为凝聚且 $Y$-平坦 $\OO_X$-模的有限复形。以（7.7.1.1）定义 $\TT_\bullet$，则函数
\begin{equation}
 y\longmapsto\sum_i(-1)^i\dim_{k(y)}T_i(k(y))\tag{7.9.4.1}
\end{equation}
在 $Y$ 上局部恒定。

可设 $Y=\Spec A$ 仿射诺特。因 $\PP_\bullet$ 有限，由（7.7.12，（i））可写 $\TT_\bullet(\M)=\HH_\bullet(\LL_\bullet\otimes\M)$，其中 $\LL_\bullet=\widetilde{L_\bullet}$，$L_\bullet$ 为有限型投射 $A$-模的有限复形。结论来自（7.9.3）。
\cmt{C3-T06}{底本本条引入 $\TT_p$ 时写成 $\cH^p$，而（7.7.1.1）采用 $\cH^{-p}$；译文明确沿用后者。对这里的有限交错和，把指标 $p$ 换为 $-p$ 不改变数值。}

\egap{7.9.5}
在（7.9.4）的条件下，若 $Y$ 连通，则（7.9.4.1）恒定。若 $Y$ 连通且非空，将其唯一整数值记为 $\operatorname{EP}(f,\PP_\bullet)$、$\operatorname{EP}(Y,\PP_\bullet)$，或在不致混淆时记 $\operatorname{EP}(\PP_\bullet)$，并称之为 $\PP_\bullet$ 关于 $f$（或关于 $Y$）的 Euler--Poincaré 特征数。一般情形中，（7.9.4.1）的值记为 $\operatorname{EP}(f,\PP_\bullet;y)$、$\operatorname{EP}(Y,\PP_\bullet;y)$ 或 $\operatorname{EP}(\PP_\bullet;y)$。

\egap{7.9.6}
设 $X,Y,f$ 满足（7.9.4）的假设，并设
\[
 0\longrightarrow\PP'_\bullet\xrightarrow{u}\PP_\bullet
 \xrightarrow{v}\PP''_\bullet\longrightarrow0
\]
为凝聚且 $Y$-平坦模的有限复形正合列，其中 $u,v$ 的次数分别为偶数 $2d,2d'$。因 $\TT_\bullet$ 为同调函子，有仅含有限项的长正合列
\[
\begin{aligned}
 \cdots\to\TT_i(\PP'_\bullet,k(y))
 &\to\TT_{i+2d}(\PP_\bullet,k(y))\\
 &\to\TT_{i+2d+2d'}(\PP''_\bullet,k(y))\\
 &\to\TT_{i-1}(\PP'_\bullet,k(y))\to\cdots.
\end{aligned}
\]
该复形的 Euler--Poincaré 特征数为零（0，11.10.1），故
\begin{equation}
 \operatorname{EP}(\PP_\bullet;y)
 =\operatorname{EP}(\PP'_\bullet;y)+\operatorname{EP}(\PP''_\bullet;y).
 \tag{7.9.6.1}
\end{equation}
例如设 $\PP_i=0$（$i<0$）。底本给出下列逐项正合图：
\[
 \begin{tikzcd}[column sep=small,row sep=small]
 \cdots\arrow[r]&0\arrow[r]\arrow[d]&0\arrow[r]\arrow[d]&0\arrow[r]\arrow[d]&0\arrow[r]\arrow[d]&\cdots\\
 \cdots\arrow[r]&0\arrow[r]\arrow[d]&0\arrow[r]\arrow[d]&\PP_1\arrow[r]\arrow[d,equal]&\PP_2\arrow[r]\arrow[d,equal]&\cdots\\
 \cdots\arrow[r]&0\arrow[r]\arrow[d]&\PP_0\arrow[r]\arrow[d,equal]&\PP_1\arrow[r]\arrow[d]&\PP_2\arrow[r]\arrow[d]&\cdots\\
 \cdots\arrow[r]&0\arrow[r]\arrow[d]&\PP_0\arrow[r]\arrow[d]&0\arrow[r]\arrow[d]&0\arrow[r]\arrow[d]&\cdots\\
 \cdots\arrow[r]&0\arrow[r]&0\arrow[r]&0\arrow[r]&0\arrow[r]&\cdots
 \end{tikzcd}
\]
非零竖箭头均为恒等同态。应用（7.9.6.1），对复形长度归纳，得到
\begin{equation}
 \operatorname{EP}(\PP_\bullet;y)
 =\sum_i(-1)^i\operatorname{EP}(\PP_i;y).\tag{7.9.6.2}
\end{equation}
这里对凝聚且 $Y$-平坦的模 $\F$，$\operatorname{EP}(\F;y)$ 指将其置于次数0、其余项为零所定义的值。因而可归结为研究单个模的 Euler--Poincaré 特征数。
\cmt{C3-T07}{上图忠实保留底本的右向箭头。若严格沿用本章同调次数 $d:P_i\to P_{i-1}$，应使用截断短正合列 $0\to P_0[0]\to P_\bullet\to P_{\geq1}\to0$，而不是把右向图当作下标递减的微分图。交错和公式不受影响；详见报告。}

\egas{7.9.7}{命题}
在（7.9.4）的假设下，设 $Y'$ 局部诺特，$g:Y'\to Y$ 为态射，$X'=X\times_Y Y'$，$f'=f_{(Y')}$，$\PP'_\bullet=\PP_\bullet\otimes_{\OO_X}\OO_{X'}$。则 $\PP'_\bullet$ 的各项凝聚且 $Y'$-平坦，对每个 $y'\in Y'$ 有
\begin{equation}
 \operatorname{EP}(\PP'_\bullet;y')
 =\operatorname{EP}(\PP_\bullet;g(y')).\tag{7.9.7.1}
\end{equation}
各项为原凝聚模的逆像，故凝聚；平坦性由（$0_{\mathrm I}$，6.2.1）及（1.4.14.5）局部验证。由（II，5.4.2），$f'$ 固有，左边遂有定义。归结到 $Y,Y'$ 仿射后，所述公式由底本引用的“（6.10.4.2）”、（7.7.2）和（7.6.7）得到。
\cmt{C3-T08}{本分册（6.10.4）没有编号为（6.10.4.2）的公式，属待订正的引文。这里仅凭（7.7.2）、（7.6.7）逐项保持余域上的维数，已足以推出有限交错和相等。}

\egas{7.9.8}{命题}
在（7.9.4）的假设下，再设存在整数 $i_0$，使每个 $y\in Y$ 及每个 $i\ne i_0$ 均有 $T_i(k(y))=0$。则
\[
 \TT_{i_0}(\OO_Y)=\cH^{-i_0}(f,\PP^\bullet)
\]
局部自由，其在 $y$ 处的秩为 $(-1)^{i_0}\operatorname{EP}(f,\PP_\bullet;y)$。

各 $T^{\OO_y}$ 满足（7.4.4），故可用（7.4.7）；由（7.5.3），$i\ne i_0$ 时 $T_i^{\OO_y}=0$。由（7.3.3），$\TT_{i_0}$ 正合，再由（7.8.4）知相应上同调层局部自由。由特征数定义，其秩为 $\dim_{k(y)}T_{i_0}(k(y))$，也就是 $(-1)^{i_0}\operatorname{EP}(f,\PP_\bullet;y)$。
\cmt{C3-T09}{底本本条陈述包含正确的因子 $(-1)^{i_0}$，但证明最后一行漏印该因子。译文此处按定义写全，并明确记录；未把正确的命题称为错误。}

\egas{7.9.9}{推论}
设 $Y$ 局部诺特，$f:X\to Y$ 固有，$\F$ 凝聚且 $Y$-平坦。若存在整数 $i_0$，使每个 $y\in Y$ 及每个 $i\ne i_0$ 均有
\[
 H^i(f^{-1}(y),\F\otimes_{\OO_Y}k(y))=0,
\]
则 $R^{i_0}f_*(\F)$ 局部自由，其在 $y$ 处的秩为 $(-1)^{i_0}\operatorname{EP}(f,\F;y)$。

\egas{7.9.10}{推论}
设 $X,Y,\F$ 满足（7.9.9）的基本条件，并设 $i>0$ 时 $R^if_*(\F)=0$。则 $f_*(\F)$ 局部自由，其在 $y$ 处的秩为 $\operatorname{EP}(f,\F;y)$。

由（7.9.9），只须证明下述引理。

\egas{7.9.10.1}{引理}
在（7.9.10）的假设下，对每个 $i>0$ 和 $y\in Y$，有
\[
 H^i(f^{-1}(y),\F\otimes_{\OO_Y}k(y))=0.
\]
可设 $Y=\Spec A$。在（7.9.4）中取只在次数0为 $\F$ 的复形，依假设，$p<0$ 时 $\TT_p(\OO_Y)=0$。由（7.3.7）知 $p<0$ 时 $\TT_p$ 正合，再由（7.7.5，a）与 d) 的等价性得到引理。

\egas{7.9.11}{命题}
在（7.9.4）的假设下，设 $\LL$ 为关于 $Y$ 非常丰沛的可逆 $\OO_X$-模，对每个 $n\in\mathbf Z$ 置 $\PP_\bullet(n)=\PP_\bullet\otimes_{\OO_X}\LL^{\otimes n}$。则对每个 $y\in Y$，函数
\begin{equation}
 n\longmapsto\operatorname{EP}(f,\PP_\bullet(n);y)\tag{7.9.11.1}
\end{equation}
是有理系数多项式，且在 $Y$ 的同一连通分支各点为同一个多项式。

$\PP_\bullet(n)$ 仍由 $Y$-平坦模组成。由（7.9.6.2），可归结为单个非零项 $\F$ 在次数0的情形；因问题对 $Y$ 局部，可设 $Y$ 仿射、$f$ 射影（II，5.5.3）。置 $X_y=f^{-1}(y)$、$\LL_y=\LL\otimes_{\OO_Y}k(y)$，它非常丰沛（II，4.4.10）。对 $\PP_\bullet(n)$ 的相应函子，由（7.7.2）有
\[
 T_i(k(y))=H^{-i}(X_y,\F_y\otimes\LL_y^{\otimes n}),
 \qquad\F_y=\F\otimes_{\OO_Y}k(y).
\]
因此所述特征数就是（2.5.1）定义的 $\chi(\F_y\otimes\LL_y^{\otimes n})$；多项式性来自（2.5.3），而每个 $n$ 的值在连通分支上由（7.9.4）恒定，故得命题。

记此有理系数多项式为 $\operatorname{PH}(f,\PP_\bullet;y)$ 或 $\operatorname{PH}(\PP_\bullet;y)$，称为关于 $\PP_\bullet,f,\LL$ 的在 $y$ 处的 Hilbert 多项式；不致混淆时也称 $\PP_\bullet$ 或 $f$ 的 Hilbert 多项式。若 $Y$ 连通且非空，省去 $y$。它将在第V章研究给定凝聚层的凝聚商层之“模空间”时发挥关键作用。

\egap{7.9.12}
沿用（7.9.6）及（7.9.11）的记号，有
\begin{align}
 \operatorname{PH}(\PP_\bullet;y)
 &=\operatorname{PH}(\PP'_\bullet;y)+\operatorname{PH}(\PP''_\bullet;y),\tag{7.9.12.1}\\
 \operatorname{PH}(\PP_\bullet;y)
 &=\sum_i(-1)^i\operatorname{PH}(\PP_i;y).\tag{7.9.12.2}
\end{align}
这直接来自（7.9.6.1）、（7.9.6.2）。在（7.9.7）的记号与假设下，同样有
\begin{equation}
 \operatorname{PH}(\PP'_\bullet;y')
 =\operatorname{PH}(\PP_\bullet;g(y')).\tag{7.9.12.3}
\end{equation}
（7.9.12.2）将复形的 Hilbert 多项式归结为单个 $Y$-平坦模的情形；后者有不依赖同调考虑的解释。

\egas{7.9.13}{推论}
设 $Y$ 诺特，$f:X\to Y$ 固有，$\LL$ 为关于 $Y$ 非常丰沛的可逆模，$\F$ 凝聚且 $Y$-平坦。则存在 $n_0$，使 $n\geq n_0$ 时 $f_*(\F(n))$ 局部自由，其在 $y$ 处的秩等于 $\operatorname{PH}(f,\F;y)(n)$。

由（II，5.5.3），$f$ 射影；故由（2.2.1），存在 $n_0$ 使 $n\geq n_0$、$i>0$ 时 $R^if_*(\F(n))=0$。应用（7.9.10）即得结论。

以下平坦性判据将在第V章凝聚层“模空间”的理论中十分重要。

\egas{7.9.14}{命题}
设 $Y$ 诺特，$f:X\to Y$ 射影，$\LL$ 为 $f$-丰沛可逆模，且对每个 $\OO_X$-模 $\F$ 及整数 $n$ 置 $\F(n)=\F\otimes\LL^{\otimes n}$。则凝聚模 $\F$ 在 $Y$ 上平坦，当且仅当存在 $n_0$，使所有 $n\geq n_0$ 的 $f_*(\F(n))$ 都局部自由。

必要性与（7.9.13）同证，因为当 $f$ 射影时，（2.2.1）也适用于丰沛可逆层。为证充分性，可设 $Y=\Spec A$。由假设及（2.2.2，（i）），各 $\Gamma(X,\F(n))$（$n\geq n_0$）为有限型投射 $A$-模（Bourbaki，\textit{Alg. comm.}，第II章，§5，第2目，定理1）。令
\[
 S=\bigoplus_{n\geq0}\Gamma(X,\LL^{\otimes n}),\qquad
 M=\bigoplus_{n\geq n_0}\Gamma(X,\F(n)).
\]
由（II，4.5.2，b；5.4.4），$X$ 典范同构于 $\Proj S$。必要时以 $\LL$ 的某正次幂代替它，可设 $S$ 由有限个一次元素生成（2.3.5.1）；由（II，2.7.5、2.7.2），$\F$ 与 $M$ 的相伴层同一。对正次齐次元 $g\in S$，因此有 $\Gamma(X_g,\F)=M_{(g)}$。$M$ 为投射 $A$-模的直和，故平坦；由（$0_{\mathrm I}$，6.3.2），$M_g$ 也平坦；$M_{(g)}$ 为其直和因子，故仍平坦（同章6.1.2）。由（1.4.14.5），$\F$ 在 $X_g$ 的每点 $Y$-平坦；这些 $X_g$ 覆盖 $X$，故命题得证。

\begin{flushright}（待续。）\end{flushright}

% END continuations/iii2_78_79.tex


% BEGIN text/remaining.tex
\head{第IV章及后续材料的待译范围}
\noindent\textbf{本页是范围说明，不是正文译文。}按原著内容顺序，第III章之后应接第IV章《概形与概形态射的局部研究》。当前已译的第四卷材料只是其所附第0章§§14--18及§19的19.0--19.6，已移到本书第0章；第IV章主体并未因标题或目录出现而被计作译完。

核对上传分册的原目录，剩余连续译文包括：第四卷第一部分的第0章19.7--19.10及§§20--23、第四章章首及§1；第二部分的第四章§§2--7；第三部分的§§8--15；第四部分的§§16--21及书末资料。各节内部仍须以原段落为单位翻译，不用提纲替代。第0章下一接续条目为19.7小节章首及（19.7.1）。

第V卷预备笔记不补作一个并不存在于当前译稿中的完整“第V章”。下编仅收录已译出的短篇及“超平面截面与锥投影”开头；该篇下一接续点为（7.1）。短篇中原有未完成句子仍按原状标示。

% END text/remaining.tex

\part*{预备笔记：已译部分}
\addcontentsline{toc}{part}{预备笔记：已译部分}
\renewcommand{\newscope}{V-S}

% BEGIN text/v1_2.tex
\head{EGA V 短篇：函数的奇异零点与超奇异零点}
\source{\texttt{EGAV1-2.ps}，转换后共6页：标题页及正文印刷第1--5页}
\noindent\textbf{载体标题。} EGA V：§1及§§2.15、2.16（原编号为EGA IV：§16及§§17.15、17.16）。Grothendieck 的 EGA V“预备笔记”的编辑英译，由 Piotr Blass 与 Joseph Blass 完成。

\egasection{V.1（原IV.16）：函数的奇异集、超奇异集与微分判据}
这一节将在关于超平面截面的V.5（原IV.20）中使用，不过在我看来，它自然的位置是在这里。

\egas{1}{定义}
设 $X$ 为正则预概形，$\varphi$ 为 $\OO_X$ 的截面。若 $\varphi_x\in\mm_x^2$，则称 $x\in X$ 为 $\varphi$ 的\textbf{奇异零点（或根）}。若它是奇异零点，并且由它定义的元素
\[
 \varphi_x\bmod\mm_x^3\in\mm_x^2/\mm_x^3\simeq\Sym(\mm_x/\mm_x^2),
\]
在解释为 $\mm_x/\mm_x^2$ 的对偶 $t_x$ 上、取值于 $k(x)$ 的二次型时为退化二次型，则称它为\textbf{超奇异零点}。
\cmt{T08}{这里的 $\Sym$ 应取二次部分 $\Sym^2$，不是整个对称代数。底本还把点的最大理想记作 $\mm_X$，本译文用 $\mm_x$ 明确其属于点 $x$ 的局部环。}
（不是超奇异零点的奇异零点，有时称为\textbf{普通奇异零点}。）\cmt{P05}{“退化二次型”的约定在特征2时必须明确；后文用极化双线性映射的行列式判断，需要另核其与本定义是否完全一致。这里不擅自加入特征不等于2的假设。}

\egas{2}{备注}
若 $x\in V(\varphi)$，则 $x$ 是 $\varphi$ 的非奇异零点，当且仅当 $\varphi_x\ne0$，且 $x$ 是 $V(\varphi)$ 的非奇异点，即正则点；也就是说，当且仅当 $x$ 是 $V(\varphi)$ 的正则点，且在 $x$ 的某个邻域中 $V(\varphi)\ne X$。

\egas{3}{定义}
设 $X$ 是域 $k$ 上的光滑预概形，$\varphi$ 为 $\OO_X$ 的截面，$x\in V(\varphi)$。若对 $k$ 的每个扩域 $k'$，以及 $X$ 的每个以 $k$ 为值域、定位于 $x$ 的点 $\xi$，其在 $X_{k'}$ 中的对应点 $x'$ 都是 $\varphi_{k'}$ 的奇异零点（分别地，超奇异零点），则称 $x$ 为 $\varphi$ 相对于 $k$ 的\textbf{几何奇异零点}（分别地，\textbf{几何超奇异零点}）。
\cmt{T08}{底本此处明确印作“values in $k$”，但随后的 $X_{k'}$ 点和备注4(a)均要求“以 $k'$ 为值域”。本段保留该误写以便对照；按定义使用时应改为 $k'$，否则非 $k$-有理点可能根本没有这种测试点，造成空泛量化。}

\egas{4}{备注}
\begin{enumerate}[label=\alph*)]
\item 由下文将建立的判据可知，在定义3中，只须用某个 $k'$ 上的一个值点进行检验；例如可取 $k'=k(x)$ 或其代数闭包 $\overline{k(x)}$，以及相应的典范值点。
\item 由备注2可知，$x$ 对 $\varphi$ 几何非奇异，当且仅当 $\varphi_x\ne0$，并且 $V(\varphi)$ 在 $x$ 处对 $k$ 光滑。
\item 设预概形 $X$ 对另一预概形 $Y$ 光滑，$\varphi$ 为 $\OO_X$ 的截面，$x\in V(\varphi)$。若 $x$ 在纤维 $X_s$ 上是相对于 $k(s)$ 的奇异零点（分别地，超奇异零点），则称它为\textbf{相对于 $Y$ 的奇异零点}（分别地，超奇异零点）；这里 $s$ 是 $x$ 在 $Y$ 中的像。
\item 在定义1的条件下，若将 $\varphi$ 换成 $\varphi'=u\varphi$，其中 $u$ 在 $x$ 处可逆，则 $x\in V(\varphi)$ 的奇异性或超奇异性不变。因而定义1以及定义3都可以自然地推广到 $\varphi$ 为某个可逆模 $L$ 的截面的情形，并在 $L=\OO_X$ 时恢复原定义。
\end{enumerate}

设预概形 $X$ 对预概形 $Y$ 光滑，$\varphi$ 为 $\OO_X$ 的截面。它给出 $\PP^2_{X/Y}$ 的截面 $d^2_{X/Y}\varphi$，后者约化为 $\PP^1_{X/Y}$ 的截面 $d^1_{X/Y}\varphi$，再约化为 $\PP^0_{X/Y}=\OO_X$ 的截面 $d^0\varphi=\varphi$。

\egas{5}{命题}
$d^0\varphi$ 的零点集是 $\varphi$ 的零点集 $V(\varphi)$；$d^1\varphi$ 的零点集是 $\varphi$ 相对于 $S$ 的奇异零点集 $V(\varphi)_{\mathrm{sing}}$。
\cmt{T08}{此前的相对底为 $Y$；本句的 $S$ 尚未引入，故此处应为 $Y$。}
\begin{proof}[原文证明]
第一断言是平凡的。第二断言就是 Jacobian 判据；也可由如下典范同构得出：若 $x$ 是 $k$-预概形 $X$ 的一个 $k$-有理点，则
\[
 \mm_x/\mm_x^2\simeq\Omega^1_{X/k}(x).
\]
\end{proof}
注意 $\operatorname{gr}^1(\PP^1_{X/Y})\simeq\Omega^1_{X/Y}$，所以限制 $d^1\varphi|V(\varphi)$ 可看成 $\Omega^1_{X/Y}\otimes\OO_{V(\varphi)}$ 的截面，它就是 $d_{X/Y}\varphi$ 在 $V(\varphi)$ 上的限制。因而可以考虑该截面的零点预概形，记为 $V(\varphi)_{\mathrm{sing}}$；由命题5，其底集正是 $\varphi$ 相对于 $Y$ 的奇异零点集。

（注意：若 $\psi$ 是预概形 $X$ 上有限型局部自由模 $E$ 的截面，则可以自然地定义 $\psi$ 的零点子预概形，例如由转置映射 $E^\vee\to\OO_X$ 的像理想定义。若 $E=\OO_X^n$、$\psi=(\psi_1,\ldots,\psi_n)$，则该理想为 $\sum_i\psi_i\OO_X$，定义的子预概形为 $V(\psi_1,\ldots,\psi_n)$。）

现在限制 $d^2\varphi$ 到 $V(\varphi)_{\mathrm{sing}}$，并注意
\[
 \operatorname{gr}^2(\PP^2_{X/Y})\simeq\Sym^2(\Omega^1_{X/Y}),
\]
便得到 $\Sym^2(\Omega^1_{X/Y})\otimes\OO_{V_{\mathrm{sing}}}$ 的典范截面 $M(\varphi)$。用域值点考察可知，这个截面正是定义1中决定二次型的截面（当 $X_k$ 由 $X/S$ 经 $\Spec k\to S$ 得到时）。\cmt{T08}{这句又改用了 $S$，而本段相对底仍记为 $Y$。它涉及同一类未统一的相对下标，不另计为一个定理反例。}
由此可用该截面描述 $V(\varphi)_{\mathrm{sup\,sing}}$：将 $M(\varphi)$ 解释为同态
\[
 M(\varphi)':G_{X/Y}\otimes_{\OO_X}\OO_{V_{\mathrm{sing}}}
 \longrightarrow\Omega^1_{X/Y}\otimes_{\OO_X}\OO_{V_{\mathrm{sing}}},
\]
取这个同态不是同构的点集。特别地，这说明 $V(\varphi)_{\mathrm{sup\,sing}}$ 为闭集。假定 $X$ 在 $Y$ 上的相对维数恒为 $d$，可以引入
\[
 D(\varphi)=\det M(\varphi)
 \in\Gamma\left(V_{\mathrm{sing}},
       (\Omega^d_{X/Y})^{\otimes2}\otimes\OO_{V_{\mathrm{sing}}}\right)
\]
使描述更精确。于是也可用 $V(\varphi)_{\mathrm{sup\,sing}}$ 表示 $V(\varphi)_{\mathrm{sing}}$（从而也为 $X$）的闭子概形，由这个可逆模截面的消失定义；它有正确的底集。

\noindent\textbf{原稿编写提示：}最好把上述构造归纳为一个“命题6”。\cmt{P01}{所上传文件没有接着给出一个正式编号为6的命题陈述；本译文不补造。}

一般情形下，不能对 $V(\varphi)_{\mathrm{sing}}$ 与 $V(\varphi)_{\mathrm{sup\,sing}}$ 说得更精确。下面考察一个在某些应用中有意义的特殊情形。设 $Y$ 也对预概形 $S$ 光滑，为明确起见，设其相对维数恒为 $m$。为简便起见，将 $V(\varphi)_{\mathrm{sing}}$ 记为 $V'$；它由秩为 $d+1$ 的局部自由模 $\PP^1_{X/Y}$ 的截面 $d^1\varphi$ 的消失定义。再假设 $V'$ 对 $S$ 光滑，相对维数为原稿所写的\cmt{T09}{该等式左边实际为 $m-d$。由 $X/Y$ 相对维数 $d$、$Y/S$ 相对维数 $m$，预期的计算应为 $(m+d)-(d+1)=m-1$。}
\[
 (m+1)-(d+1)=m-1.
\]

（注意，$V(\varphi)_{\mathrm{sing}}$ 与 $V(\varphi)_{\mathrm{sup\,sing}}$ 的记号有歧义，因为没有显示相对的预概形；此处约定相对底为 $Y$。还应注意，由假设可知，$\varphi$ 的每个相对于 $Y$ 的奇异零点，相对于 $S$ 都非奇异。）

在此情形下，可以写出下面的局部自由模层图。为使列图在默认版心内可读，设\cmt{图中文字}{以上 $A,B,C,D,R,W$ 只是本译稿为重排所加的缩写；原稿在图内展开写出各模。本图把原稿的 $V^1$ 统一显示为前文已定义的 $V'$，没有增加任何交换关系。$G$ 的含义沿用英译稿后面的编者疑问。}
\[
 A=(\PP^1_{X/Y})^\vee|_{V'},\quad
 B=\Omega^1_{X/S}|_{V'},\quad C=\Omega^1_{X/Y}|_{V'},
\]
\[
 D=\Omega^1_{Y/S}|_{V'},\quad R=\Omega^1_{V'/S},\quad
 W=\omega_{X/Y}^{-2}|_{V'}.
\]

\[
\begin{tikzcd}[row sep=large,column sep=large]
 &&0\\
 &&C\arrow[u]\\
 0\arrow[r]&A\arrow[r]\arrow[ur,"\mu"]&B\arrow[u]\arrow[r]&R\arrow[r]&0\\
 &\OO_{V'}\arrow[u]\arrow[r]&D\arrow[u]\arrow[ur,"\nu"']\\
 W\arrow[ur]&&0\arrow[u]
\end{tikzcd}
\]
竖直列来自光滑态射 $X\to Y$ 和 $Y\to S$ 的传递性正合列，再与 $\OO_{V'}$ 张量；由于列中各模均局部自由，张量后仍正合。横行是以下正合列的特例：在 $X/S$ 上给定局部自由模 $F$ 的截面 $\psi$，令 $W$ 为其零点概形，则有
\[
 F^\vee\otimes\OO_X\longrightarrow
 \Omega^1_{X/S}\otimes\OO_W\longrightarrow\Omega^1_{W/S}\longrightarrow0.
\]
原稿接着说：若 $X/S$ 光滑，则第一个同态恰在 $W$ 对 $X$ 光滑且具有“正确”相对维数的点处单射（也就是说，在当前情形下处处单射）。\cmt{E06、T10}{若把它理解为 $W$ 上的模列，首项应限制到 $W$；作为 $X$ 上的模列则可经限制映射解释。相对光滑的底应核为 $S$，而层单射不足以推出光滑，须检验纤维满秩。尖点例子见E06。}
这个正合列直接由第16段中的下列正合列得到（可以把上述版本列为推论）：\cmt{T10}{末项在底本中重复写作 $\Omega^1_{X/S}$；余法正合列的末项应为 $\Omega^1_{W/S}$。这里保留错误式，而非无说明地改写。}
\[
 J/J^2\longrightarrow\Omega^1_{X/S}\otimes\OO_W
      \longrightarrow\Omega^1_{X/S}\longrightarrow0.
\]

能够在最左边添上一个零的点集的刻画，包含在 Jacobian 判据中。\cmt{E06}{这一句也必须结合纤维满秩与适当余维假设理解，不能把抽象单射与光滑判据等同。}

注意，有典范同构
\[
 \PP^1_{X/S}=\Omega^1_{X/S}+\OO_X,
 \qquad (\PP^1_{X/S})^\vee=G_{X/S}+\OO_X
\]
（原稿中的 $G$ 是拉长的字形）。\ednote{拉长的 $G$ 是否就是切层？}
\cmt{T08}{这里以及紧接的 $\Omega^1_{X/S}$、$G_{X/S}$ 与前面 Hessian 构造的 $X/Y$ 下标不一致。若解释图中的 $\mu$，应使用相对于 $Y$ 的主部与微分。}
另一方面，可验证图中复合同态 $\mu$ 在因子 $\OO_{V'}$ 上为零，在因子 $G_{X/S}\otimes\OO_{V'}$ 上化为从上述 $\Sym^2(\Omega^1_{X/S})\otimes\OO_{V'}$ 的截面 $M(\varphi)$ 得到的同态 $M(\varphi)$〔此处字迹不清〕。因此，在 $X$ 的点 $x$ 处，$M(\varphi)$ 非退化，即 $M(\varphi)'$ 满射，当且仅当原稿所写的 $M$ 在 $x$ 处满射；图中这也等价于原稿所写的 $V$ 在 $x$ 处满射，因为两者都表示：所提到的 $\Omega^1_{X/S}|_{V'}$ 的两个子模之和到该模的典范同态在 $x$ 处满射。\cmt{T08}{本句的 $M,V$ 依图应读为 $\mu,\nu$；这是印在英译载体中的字母错置。}
由此得到：

\egas{7}{命题}
在以上条件下（需重新列出），$V(\varphi)_{\mathrm{sup\,sing}}$ 的底集恰为 $V(\varphi)_{\mathrm{sing}}$ 中使态射
\[
 V(\varphi)_{\mathrm{sing}}\longrightarrow Y
\]
分歧的点集；这里源与靶都是 $S$ 上光滑的预概形，相对维数分别为 $m-1$ 与 $m$。

用前辈的说法（可把这一点写成备注），点 $x\in V(\varphi)$ 相对于 $Y$ 超奇异，当且仅当“它由至少两个重合的、无穷邻近的奇异点组成……”

可以而且必须把命题7精确化为子预概形之间的相等，而不仅是子集之间的相等。事实上，$V(\varphi)_{\mathrm{sup\,sing}}$ 已定义为 $X$ 的一个闭子预概形；同样也可自然地定义 $V'$ 的一个闭子概形，其底集为相对于 $Y$ 的分歧点集。只须表示局部自由模同态
\[
 Q=\Omega^1_{Y/S}|_{V'}\longrightarrow M=\Omega^1_{V'/S}
\]
不满射的点集。若两者的秩分别为 $q,r$，原英译在此写成
\[
 \Lambda^1Q\longrightarrow\Lambda\sqrt M
\]
不满射的点集，并将截面所属的 Hom 表达式写成\cmt{T01}{这里的根号式字形在 PS 渲染页面上确实可见，并非本次文本抽取造成。它不能构成有效的外幂公式。可有数学意义的表达应为 $\bigwedge^rQ\to\bigwedge^rM$，视为 $\mathscr Hom(\bigwedge^rQ,\det M)$ 的截面。本译文不把破损式冒充可读的完整公式。}
\[
\begin{split}
 \Hom(\Lambda^1Q,\Lambda\sqrt M)
 &\simeq(\approx)(\Lambda\sqrt Q)\\
 &\qquad\otimes(\Lambda^rR)\otimes(\Lambda^1Q)^\vee.
\end{split}
\]

原稿随即把它描述为这个显然截面的零点集，即一个闭零点子预概形的底集，并把它记为 $\operatorname{Ram}(V'/Y)$。我断言，这个子概形等于 $V(\varphi)_{\mathrm{sup\,sing}}$。这是关于上述图的一个简单练习：由 $\mu$ 的描述可知，$V(\varphi)_{\mathrm{sup\,sing}}$ 是对同态
\[
 P=(\PP^1_{X/Y})^\vee|_{V'}\longrightarrow
 S=\Omega^1_{X/Y}|_{V'}
\]
使用与刚才对 $Q\to R$ 所作相同的构造而定义的。

因此归结为以下一般情形。在赋环空间 $W$ 上给定秩 $m$ 的局部自由模 $M$，以及秩分别为 $p,q$ 的两个局部自由子模 $P,Q$，使原稿所写的
\[
 p+1=m+1
\]
成立；对态射 $P\to M/W=S$ 和 $Q\to M/P=R$ 使用上述构造，得到以下两个模的截面：
\begin{align*}
 &P\otimes\det S\otimes(\det P)^{-1}\\
 &\hspace{1em}=P\otimes\det M\otimes(\det P)^{-1}\otimes(\det Q)^{-1},\\
 &Q\otimes\det M\otimes(\det P)^{-1}\otimes(\det Q)^{-1}.
\end{align*}
\cmt{T01}{按上一句的余秩关系，这里应为 $p+q=m+1$；$M/W$ 则应为 $M/Q$，因为 $W$ 是赋环空间而非 $M$ 的子模。这些修正均需明示。}
也可把这两个截面视为从原稿所写的
\[
 L=\det P\otimes\det Q\otimes\det M
\]
分别到 $P,Q$ 的同态。\cmt{T01}{由上面两个截面所属的模计算，$L$ 的末因子应为 $(\det M)^{-1}$。详见核查报告中的类型核对。}
（注意：对局部自由模 $F$，用 $\det F$ 表示它的最高外幂，并使用如下事实：若
\[
 0\longrightarrow F'\longrightarrow F\longrightarrow F''\longrightarrow0
\]
是这种模的短正合列，则有典范同构 $\det F=\det F'\otimes\det F''$。）在这些条件下，有下面这个图的交换性——

\par\bigskip\noindent\textbf{原文件在此结束。}标题提及的其他小节不能仅凭文件名认定已收录；所上传的六页文件没有在最后一句之后提供交换图或证明结尾。本译稿已翻译其全部可读正文，但没有补造所缺内容。

% END text/v1_2.tex

\renewcommand{\newscope}{V-H}

% BEGIN text/v5_01_02.tex
\head{EGA V 预备笔记：超平面截面与锥投影}
\source{\texttt{EGAVaiello.pdf} 第1--12页；本篇与 \texttt{EGAV5.ps} 的开头部分对应}
\noindent\textbf{底本标题署名：}Alexandre Grothendieck 的 EGA V；其“预备笔记”的翻译与编辑：Piotr Blass、Joseph Blass。

\egasection{编写摘要}
这一稿把应收入最终稿的全部结果杂陈在一起，作了详细概述。要达到最终稿的形式，需要对目前这个“第零阶段”彻底重新组织。第一步大概应当是重订计划（毫无疑问，现在的第11、12、14、15节将大幅提前）。我甚至还没有写第16节；原则上，它应当没有困难，也不以任何方式影响前面的各节，因为涉及的只是转换表述。

你会注意到，这里有一条命题10.3，它应当出现在前面的某个段落。

就此，我还想告诉你，我另有若干彼此很不相同、但都涉及双有理映射的结果，很愿意把它们放到某处。我觉得这些还不足以独立构成一节。你对放在哪里有什么建议吗？我打算不久把它们连同这份笔记的第16节寄给你。\ednote{询问 AG：第16节究竟是否写过。}

此外，目前的第20段大概会扩展成两段，其中一段由 Grassmann 概形上属于“初等几何”类型的结果组成。必要时，是否也可以把我跟你提过的关于双有理变换的补充材料放在这里（如果找不到更合适的地方）？

\egasection{拟定目录}
\begin{enumerate}[label=\arabic*)]
\item 记号。
\item 泛超平面截面：局部性质。
\item 泛超平面截面：几何不可约性与连通性。
\item 变动超平面截面：“充分一般”的截面。
\item Seidenberg 类型定理。
\item 任意超平面截面的连通性。
\item 特殊超平面截面与多截面构造的应用。
\item 例外超平面集合的维数。
\item 射影浸入的改变。
\item 超平面截面的铅笔与爆破簇的纤维化。
\item Grassmann 概形。
\item 前述结果向线性截面的推广。
\item 初等态射与 M. Artin 的一个定理。
\item 锥投影。
\item 部分前述结果的公理化。
\item 转换为线性系的语言。
\end{enumerate}
以上是底本列出的拟定目录，不表示全部节均已经写成，也不表示本批已翻译目录中全部内容。

\egasection{§1\quad 预备知识与记号}
设 $S$ 为预概形，$E$ 为 $S$ 上有限型局部自由模，$E^\vee$ 为其对偶。记 $P=\mathbf P(E)$ 为 $E$ 定义的射影丛，$P^\vee$ 为 $E^\vee$ 定义的射影丛。我们称 $P^\vee$ 为 $P$ 的\textbf{超平面概形}。这个术语可作如下解释。

设 $\xi$ 是 $P^\vee$ 在 $S$ 上的截面，因而由 $E^\vee$ 的一个可逆商模 $L$ 确定。由它得到 $(E^\vee)_P=(E_P)^\vee$ 的可逆商模 $L_P$；另一方面，我们有 $E_P$ 的可逆商模 $\OO_P(1)$。取对偶后，可以把 $L_P^{-1}$、$\OO_P(-1)$ 分别看成 $E_P$、$(E_P)^\vee$ 的可逆子模，且局部为直和因子。配对 $E_P\otimes(E_P)^\vee\to\OO_P$ 因而定义自然配对
\begin{equation}
 \OO_P(-1)\otimes L_P^{-1}\longrightarrow\OO_P,\tag{*}
\end{equation}
或者其转置同态
\begin{equation}
 \OO_P\longrightarrow\OO_P(1)\otimes L_P=L_P(1).\tag{**}
\end{equation}
后者就是一个由 $\xi$ 典范确定的 $L_P(1)$ 的截面。此截面的“除子”，即由（*）的像理想定义的 $P$ 的闭子概形 $H_\xi$，称为点 $\xi\in P^\vee(S)$ 定义的\textbf{超平面}。

还可以这样描述它：在 $S$ 上局部地，$\xi$ 由 $E$ 的一个截面 $\phi$ 给出，且对每个 $s$ 都有 $\phi(s)\ne0$；$\xi$ 确定 $\phi$，相差一个 $\OO_S$ 可逆截面的倍数。由于 $E=p_*\OO_P(1)$（$p:P\to S$ 为投影），$\phi$ 也可视为 $\OO_P(1)$ 的截面，而它的除子正是 $H_\xi$。

当然，若把 $L^{-1}$ 看成 $E$ 的可逆子模，且局部为直和因子，则 $\xi$（即 $L$ 或 $L^{-1}\subset E$）与 $\phi$ 的对应，由取 $L^{-1}$ 的一个处处不为零的截面得到，也就是选择 $L^{-1}$ 的一个平凡化；这种平凡化无论如何总在局部存在。注意，$H_\xi$ 也就是 $\mathbf P(E/L^{-1})$（典范同构）。这是第三种描述方式。特别地，$\mathbf P(E/L^{-1})$ 典范嵌入 $P=\mathbf P(E)$，这一描述还证明了 $H_\xi$ 是 $S$ 上的射影丛，因而在 $S$ 上光滑。（还应当在 EGA IV 第17段中说一下，射影丛是光滑的。）无疑，从这一描述开始会更好。

\noindent\textbf{备注。}
$H_\xi$ 的构造与基变换相容；换言之，得到函子范畴 $(\mathrm{Sch}/S)^{\mathrm{op}}\to\mathrm{Ens}$ 中的自然变换
\[
 P^\vee\longrightarrow\operatorname{Div}(P/S),
\]
其中右端表示 $P/S$ 的“相对除子”函子：它在任意 $S$-预概形 $S'$ 上的值，是 $P_{S'}$ 中相对 $S'$ 为横截完全交且余维为1的闭子概形集合（比较第19段）。\ednote{这里使用 EGA I 新版的记号；第19段指 EGA IV。}
容易证明这是函子的单态射，也就是说，$\xi$ 由 $H_\xi$ 确定。这说明“超平面概形”的术语是恰当的。我们将看到，$\operatorname{Div}(P/S)$ 可由各 $\mathbf P(\operatorname{Sym}^m_k(E^\vee))$ 的不交并预概形表示，从而 $P^\vee$ 可视为其一个既开又闭的子概形。\ednote{比较 Mumford 的 \textit{Lectures on Curves on an Algebraic Surface}。}\cmt{P03}{底本此处在一般底预概形 $S$ 上使用了未定义的下标 $k$，对称幂的对偶位置和秩的退化情形也需核对。这段是后文拟证明的宣告，不能仅据此视为已完成证明。本文按所上传英文稿保留表达。}
（严格说，对 $P/S$ 相对除子的确定，用目前已有的工具和第II章的结果就能进行，也可以作为例子补到 EGA IV 第19段中。）

现在作基变换 $S'=P^\vee\to S$，并取 $P^\vee_{S'}=\mathbf P(E^\vee_{S'})$ 在 $S'$ 上的对角截面，或称“泛截面”。得到 $P_{S'}=P\times_SP^\vee$ 的闭子概形 $H$，有时称为 $P$ 与 $P^\vee$ 的\textbf{关联概形}。它由典范同态
\[
 \OO_P(-1)\otimes_S\OO_{P^\vee}(-1)
 \longrightarrow\OO_{P\times_SP^\vee}
\]
的像理想定义。因此它是 $P^\vee$ 上的射影丛，由对称性，也是 $P$ 上的射影丛。

注意，给定 $P^\vee$ 在 $S$ 上的截面 $\xi$ 后，可以从这个“普遍超平面”$H$ 出发，沿基变换 $\xi:S\to P^\vee$ 取逆像，恢复“特殊”超平面 $H_\xi$。对 $P^\vee$ 的每个取值于任意 $S$-预概形 $S'$ 的点也一样：把它看成 $P^\vee_{S'}$ 在 $S'$ 上的截面，便定义 $H_\xi\subset P_{S'}$；它正是沿 $\xi:S'\to P^\vee$ 对 $H$ 作基变换的逆像。

以下假设有一个在 $P$ 上有限型的预概形 $X$，以及一个 $S$-态射 $f:X\to P$。\ednote{或者应当是“在 $S$ 上有限型”，我不能确定。}
这一段的主要目标之一，是对每个超平面 $H_\xi$ 研究其逆像
\[
 Y_\xi=f^{-1}(H_\xi)=X\times_PH_\xi,
\]
特别是联系 $X$ 与 $Y_\xi$ 的性质。像通常一样，也须考虑取值于任意 $S$-概形 $S'$ 的点；此时 $H_\xi$ 为 $P_{S'}$ 中的超平面，并置
\[
 Y_\xi=f_{S'}^{-1}(H_\xi)=X_{S'}\times_{P_{S'}}H_\xi
       =X\times_PH_\xi.
\]
下标 $S'$ 表示基变换 $S'\to S$ 的效果；在最后一式中，通过复合 $H_\xi\to P_{S'}\to P$ 把 $H_\xi$ 视为 $P$-概形。因而仍宜先考虑 $\xi$“普遍”的情况：$S'=P^\vee$，$\xi$ 为对角截面，于是 $H_\xi=H$；此时记 $Y=Y_\xi$（等待 Dieudonné 提出更合适的记号）。对于一般的 $\xi:S'\to P^\vee$，也有 $Y_\xi=Y\times_{P^\vee}S'$。

最后，若 $\F$ 是 $X$ 上的模层，\ednote{询问 A. G.：这里的模是否始终指凝聚或拟凝聚模层。}
以 $\G_\xi$ 表示它在 $Y_\xi$ 上的逆像；原英译稿接着写，“以 $\G$ 表示它在 $H$ 上的逆像”，从而 $\G_\xi=\G\otimes_{\OO_{P^\vee}}\OO_{S'}$。\cmt{T04}{这里 $\G$ 应当定义在 $Y=X\times_PH$ 上，而非一般没有典范态射到 $X$ 的 $H$ 上。下面的图也把 $\G$ 标在 $Y$ 上，提供了同一底本内部的更正依据。}

把主要构造与记号归纳如下（$\F,\G,\G_\eta$ 分别位于 $X,Y,Y_\eta$ 上）：
\[
\begin{tikzcd}[column sep=1.6em,row sep=2.1em]
X \arrow[d,"f"'] & X\times_SP^\vee \arrow[l] \arrow[d] & Y \arrow[l,hook'] \arrow[d] & Y_\eta \arrow[l] \arrow[d]\\
P \arrow[d] & P\times_SP^\vee \arrow[l] \arrow[d] & H \arrow[l,hook'] \arrow[dl] & H_\eta \arrow[l] \arrow[d]\\
S & P^\vee \arrow[l] & & S' \arrow[ll,"\xi"]
\end{tikzcd}
\]
图中的方形与菱形都为笛卡尔图。下一节将系统研究如下情况：$S'=\Spec K$，其中 $K$ 为域，它在 $P^\vee$ 中的像是相应纤维 $P_s^\vee$ 的泛点。作基变换 $\Spec k(s)\to S$ 后，归结到 $S=\Spec k$、$k$ 为域；下一节始终作此假设。

所研究的 $X,Y_\xi$ 的多数性质具有“几何性质”，因此在基变换下不变。这使我们可不失一般性地限制到 $K$ 代数闭，或者 $K=k(\eta)$（$\eta$ 为 $P^\vee$ 的泛点、$\xi:\Spec K\to P^\vee$ 为典范态射）的情形。对于涉及 $X,Y_\xi$ 的几何问题，亦可在基变换后设 $k$ 代数闭。

术语提醒：若 $f$ 为浸入，通常称 $Y_\xi$ 为 $X$ 的\textbf{超平面截面}（相对于射影浸入 $f$ 与超平面 $H_\xi$）。没有理由不把这个术语推广到任意 $f$ 的情形。

\egasection{§2\quad 泛超平面截面的研究：局部性质}
现在 $S=\Spec k$，$k$ 为域。若 $\eta$ 是 $P^\vee$ 的点，$\xi:\Spec k(\eta)\to P^\vee$ 为典范态射，则也写 $H_\eta,Y_\eta,\G_\eta$，代替 $H_\xi,Y_\xi,\G_\xi$。在本节中，$\eta$ 始终表示 $P^\vee$ 的泛点。

\egas{2.1}{命题}
设 $X$ 不可约。则 $Y_\eta$ 不可约或为空；在非空情形下，它占优于 $X$。〔原英译接有“字迹不清，向 AG 核实”以及“$Y$ 不可约”的编者文字。〕\ednote{向 Grothendieck 询问：例如第一节中带下划线的大写字母的含义和作用是什么？}
\begin{proof}
由于 $H\to P$ 是射影丛，$Y\to X$ 也是，所以 $X$ 不可约时 $Y$ 不可约。因此，$Y$ 在 $P^\vee$ 上的泛纤维 $Y_\eta$ 不可约或为空；非空时，它的泛点就是 $Y$ 的泛点，因而位于 $X$ 的泛点之上。
\end{proof}

\egas{2.2}{命题}
设 $Z$ 为 $P$ 的一个子集。则它在 $H_\eta$ 中的逆像 $Z_\eta$ 为空，当且仅当 $Z$ 的每个点都为闭点。特别地，若 $Z$ 可构造，则 $Z_\eta=\varnothing$ 当且仅当 $Z$ 有限。
\begin{proof}
可归结到 $Z$ 只有一个点 $z$；只须证明 $H_\eta$ 在 $P$ 中的像恰好由非闭点构成。令 $X$ 为 $z$ 的闭包，使用（2.1），只须证明 $Z_\eta=\varnothing$ 当且仅当闭子概形 $X\subset P$ 有限。

把 $X$ 换成 $X_{k(\eta)}\hookrightarrow P_{k(\eta)}$ 后，“必要”方向来自下面的事实——这里应给出参考，而且值得作为引理重述：若 $Y$ 是 $X$ 的某个超平面截面，且该截面为空，则 $X$ 有限；事实上，$X\subset P\setminus H$ 同时仿射和射影……〔底本此处在句中又写 $Y_\eta=\varnothing$，记号未统一。〕

“充分”方向是直接的。例如，注意 $Y$ 是 $X$ 上相对维数为 $n-1$ 的射影丛（$n$ 为 $P,P^\vee$ 在 $S$ 上的相对维数）。当 $X$ 在 $k$ 上有限时，$Y$ 的绝对维数为 $n-1$，小于 $\dim P^\vee=n$，所以 $Y\to P^\vee$ 不占优，泛纤维 $Y_\eta$ 为空。
\end{proof}

\egas{2.3}{推论}
设 $f:X\to P$ 为有限型态射，$Z$ 为 $X$ 的可构造子集。它在 $Y_\eta$ 中的逆像为空，当且仅当 $f(Z)$ 有限。特别地，$Y_\eta$ 为空，当且仅当 $f(X)$ 有限。

\egas{2.4}{推论〔第二断言有误〕}
设 $Z,Z'$ 为 $X$ 的两个闭子集，且 $Z$ 不可约；记它们在 $Y_\eta$ 中的逆像为 $Z_\eta,Z'_\eta$。则 $Z_\eta\subset Z'_\eta$ 当且仅当 $f(Z)$ 有限或 $Z\subset Z'$。

原文还声称：$Z_\eta=Z'_\eta$ 当且仅当 $f(Z),f(Z')$ 都有限，或者 $Z=Z'$。\cmt{E03}{后一断言按现有假设不成立。可以给 $Z$ 添加一个与它不相交的闭点而不改变泛截面。详见E03；若再假设 $Z'$ 也不可约，才能用前一断言的对称应用得到所写等号判据。}
\begin{proof}[原文证明]
这是（2.3）的直接结果。因为 $f(Z\setminus(Z\cap Z'))$ 只有在 $Z\subset Z'$ 或 $f(Z)$ 有限时才可能有限：若 $Z\not\subset Z'$，则 $Z\setminus(Z\cap Z')$ 在 $Z$ 中稠密，故其像在 $f(Z)$ 中稠密；若前者有限，它因可构造而闭，因此后者也有限。
\end{proof}

\egas{2.5}{推论}
对 $X$ 的每个满足 $\dim\overline{f(X_i)}>0$ 的不可约分支 $X_i$，取其在 $Y_\eta$ 中的逆像 $Y_{i\eta}$。则 $Y_{i\eta}$ 是 $Y_\eta$ 的一个不可约分支，并由此得到上述 $X_i$ 与 $Y_\eta$ 的不可约分支之间的一一对应。
\begin{proof}
由（2.3），$Y_\eta$ 是这些 $Y_{i\eta}$ 的并，且后者均闭且非空。由（2.1），它们不可约；由（2.4）的包含判据，它们之间互不包含，故结论成立。
\end{proof}
注意，若 $\dim X_i=d_i$，则 $\dim Y_i=d_i-1$。\cmt{记号}{此句底本省略了 $\eta$。依上下文应指 $Y_{i\eta}$，不能误读为普遍关联空间中位于 $X_i$ 上的整族。}
更一般地：

\egas{2.6}{命题}
假设对每个不可约分支 $X_i$ 都有 $\dim\overline{f(X_i)}>0$，亦即 $Y_{i\eta}\ne\varnothing$；或者，$f$ 为浸入且 $\dim\overline{f(X)}>0$〔此处原英译注明字迹稍不清，待向 AG 核实〕。则
\[
 \dim Y_\eta=\dim X-1.
\]
\ednote{法文为“ou que”；我认为应当译为“where”。}
\begin{proof}
由（2.5）归结到 $X$ 不可约。按构造，$Y_\eta$ 在 $X_{k(\eta)}$ 中是某个可逆模截面的除子，这个可逆模是 $\OO_P(1)$ 的逆像。$X_{k(\eta)}$ 不可约，因为 $X$ 不可约且 $k(\eta)/k$ 是纯超越扩张；这一点本应在本节开头指出。又有 $Y_\eta\ne X_{k(\eta)}$，因为 $Y_\eta\to X$ 的像并不是整个 $X$，与忠实平坦的 $X_{k(\eta)}\to X$ 不同：由（2.3），前者的像不含 $X$ 的闭点。所以
\[
 \dim Y_\eta=\dim X_{k(\eta)}-1=\dim X-1
\]
（最后一个等式需给参考）。
\end{proof}

\egas{2.7}{命题}
设 $\F$ 为 $X$ 上拟凝聚模，因而 $\G_\eta$ 在 $Y_\eta$ 上拟凝聚。取 $\F$ 的所有满足 $\dim\overline{f(Z_i)}>0$ 的关联素循环 $Z_i$。记 $Z_{i\eta}$ 为其在 $Y_\eta$ 中的逆像。则这些 $Z_{i\eta}$ 正是 $\G_\eta$ 的全部关联素循环；它们之间的包含关系与相应 $Z_i$ 之间相同。
\begin{proof}
最后一断言包含在（2.4）的包含判据中。由于 $Y\to X$ 是射影丛，所以它平坦且纤维不可约并满足 $(S_1)$。由 EGA IV 第3段，$\F$ 在 $Y$ 上的逆像 $\G$ 的关联素循环，就是 $\F$ 的关联素循环的逆像。再限制到 $Y\to P^\vee$ 的泛纤维，即知 $\G_\eta$ 的关联素循环为上述 $Z_i$ 的非空逆像；由（2.3）得到结论。

严格说，经由 $Y$ 的这一步并非必要。可直接用 $Y_\eta\to X$ 平坦且纤维满足 $(S_1)$、并不可约的事实；其纤维甚至几何正则（即该态射正则）且几何不可约，因为它们是射影概形的局部化。这一点也可用于（2.1）的证明。
\end{proof}

\egas{2.8}{命题}
设 $\F$ 在 $X$ 上凝聚，$y\in Y_\eta$，$x$ 为其在 $X$ 中的像。设 $P(M)$ 是诺特局部环 $A$ 上有限型模 $M$ 的下列性质之一：
\begin{enumerate}[label=(\roman*)]
\item $\coprof M\leq n$〔参考待补〕；
\item $M$ 满足 $(S_k)$〔参考待补〕；
\item $M$ 是 Cohen--Macaulay 模；
\item $M$ 约化〔参考待补〕；
\item $M$ 整〔参考待补〕。
\end{enumerate}
则 $\G_{\eta,y}$ 满足 $P$，当且仅当 $\F_x$ 满足 $P$。
\begin{proof}
这直接来自第6段的结果，\ednote{澄清这一参考：是 EGA IV 吗？}
并使用 $Y_\eta\to Y$ 是正则态射，从而 $\OO_{X,x}\to\OO_{Y_\eta,y}$ 正则这一事实。\cmt{T04}{要推出所写局部环同态正则，此处应援引刚在（2.7）讨论的 $Y_\eta\to X$；底本把靶写为 $Y$，需核对并区分两条态射。}
结合（2.3），得到下列结论。
\end{proof}

\egas{2.9}{推论}
沿用（2.8）的记号，设 $Z$ 为 $X$ 上不满足 $P(\F_x)$ 的点的集合。则 $\G_\eta$ 处处满足 $P$，当且仅当 $f(Z)$ 是 $P$ 的有限子集，或者说 $\dim\overline{f(Z)}=0$。\cmt{T05}{“有限”包括空集，而按第0章约定 $\dim\varnothing=-\infty$；若要包含无坏点的情形，应写维数 $\leq0$，或单独说明 $Z$ 为空的情况。}
\begin{proof}
（2.8）说明 $h^{-1}(Z)$ 是 $\G_\eta$ 的 $P$-奇异点集；由（2.3），其为空当且仅当 $f(Z)$ 有限。（这里 $h$ 表示 $Y_\eta\to X$；我刚刚发现，在（2.8）中字母 $P$ 被使用了两次。）
\end{proof}

\egas{2.10}{推论〔原文仅有提纲〕}
“$Y_\eta$ 约化、分别地局部整的条件。”\cmt{P01}{英文底本在此只有这一句，没有完整陈述或证明。本译稿不补造原文。}

\egas{2.11}{推论}
设 $y\in Y_\eta$，$x$ 为其像。$Y_\eta$ 在 $y$ 处正则，分别地满足 $(R_k)$〔参考待补〕，分别地正规，当且仅当 $X$ 在 $x$ 处满足同一性质。设 $Z$ 为 $X$ 的非正则点集，分别地不满足 $(R_k)$ 的点集，分别地非正规点集；则 $Y_\eta$ 正则、分别地满足 $(R_k)$、分别地正规，当且仅当 $f(Z)$ 有限，即 $\dim\overline{f(Z)}=0$。\cmt{T05}{最后的等号仍需像（2.9）一样处理空集。底本中还将一次 $R_k$ 排成了 $E_k$。}
\begin{proof}
证明与（2.8）、（2.9）相同。必须给出保证 $Z$ 闭的各个参考，因为为应用（2.3），必须知道它可构造。
\end{proof}
注意，（2.10）、（2.11）完全没有讨论相应的几何性质；这里的结果是“绝对”的。现在考察几何性质。（也可以在这里另起一节。）

\subsec{几何性质〔英译编者所加标题〕}
\egas{2.12}{定理〔原证明含错误步骤〕}
设 $f:X\to P$ 非分歧。设 $y\in Y_\eta$，$x$ 为其在 $X$ 中的像。则 $X$ 在 $x$ 处对 $k$ 光滑，当且仅当 $Y_\eta$ 在 $y$ 处对 $k(\eta)$ 光滑。
\begin{proof}[原文证明]
可设 $k$ 代数闭。若 $Y_\eta$ 在 $y$ 处对 $k(\eta)$ 光滑，则它在该点正则；$Y_\eta\to X$ 平坦，所以 $X$ 在 $x$ 处正则〔参考待补〕。由于 $k$ 代数闭，因而完美，$X$ 在 $x$ 处对 $k$ 光滑〔参考待补〕。

反过来，可将 $X$ 换成 $x$ 的开邻域而设 $X$ 光滑；原文进而说，由光滑性的 Jacobian 判据，可以在 $P$ 的某个开子集 $U$ 中用 $p$ 个方程表示为
\[
 X=V(f_1,\ldots,f_p),
\]
其中 $df_1,\ldots,df_p$ 处处线性无关。\cmt{G01}{非分歧态射不自动在 Zariski 局部成为浸入；此处还需适当的 étale 局部化或另作论证。详见G01。}
在 $P^\vee$ 中引入仿射坐标 $S_1,\ldots,S_n$，在 $x$ 的邻域中引入 $P$ 的仿射坐标 $T_1,\ldots,T_n$（选取不包含 $x$ 的无穷远超平面 $H_\infty$），则 $Y_\eta\hookrightarrow U_{k(\eta)}$ 为
\[
 V\left(f_1,\ldots,f_p,\sum_iS_iT_i-1\right).
\]
只须验证 $f_1,\ldots,f_p,\sum_iS_iT_i-1$ 相对于 $k(\eta)$ 的微分线性无关。这些微分是 $\Omega^1_{U/k}\otimes_kk(\eta)$ 的截面
\[
 df_1,\ldots,df_p,\sum_iS_i\,dT_i.
\]
由于 $df_i$ 在 $U$ 的每一点线性无关，而且 $dT_i$ 处处构成 $\Omega^1_{U/S}$ 的基，因而构成生成系，原文断言这些微分在 $U_{k(\eta)}$ 的每一点都线性无关，至少在 $p\leq n$、即
\[
 E=\Omega^1_{U/k}\Big/\sum_{1\leq i\leq n}\OO_Udf_i\ne0
\]
时如此。\cmt{T06}{底本在商的求和范围写到 $n$，而已列出的方程只有 $f_1,\ldots,f_p$；此外，商的秩为 $n-p$，非零对应 $p<n$，不是 $p\leq n$。}
这是关于非零局部自由模 $E$ 的一族生成元 $a_i$（$1\leq i\leq n$）的“小引理”：截面 $\sum_iS_ia_i$ 在 $E\otimes_kk(\eta)$ 中处处不消失。\cmt{E04}{这条辅助断言为假：在 $\mathbf A^1_k$ 的平凡线丛上取生成元 $1,t$，$S_1+S_2t$ 在 $t=-S_1/S_2$ 处为零。此反例否定这一步，不否定定理本身。详见E04及G01。}
另一方面，$p=n$ 时是平凡情形，因为 $Y_\eta=\varnothing$。
\end{proof}

\egas{2.13}{推论}
设 $Z$ 为 $X$ 不在 $k$ 上光滑的点集。则 $Y_\eta$ 对 $k(\eta)$ 光滑，当且仅当 $Z$ 有限。特别地，若 $X$ 光滑，则 $Y_\eta$ 也光滑。
\begin{proof}
由（2.12）、（2.3）得出。更一般地，有以下结果。\cmt{假设继承}{这里继承（2.12）的非分歧假设；非分歧有限型态射拟有限，因此有限的 $f(Z)$ 与有限的 $Z$ 在此可联系起来。不能把这一推论脱离其前提独立引用。}
\end{proof}

\egas{2.14}{定理}
设 $y\in Y_\eta$，$x$ 为其在 $X$ 中的像。设 $P(A,K)$ 为域 $K$ 上诺特局部代数 $A$ 的下列性质之一：
\begin{enumerate}[label=(\roman*)]
\item $A$ 在 $K$ 上几何正则；
\item $A$ 在 $K$ 上几何满足 $(R_k)$；
\item $A$ 在 $K$ 上可分；
\item $A$ 在 $K$ 上几何正规。
\end{enumerate}
原英译稿写出的结论是\cmt{T07}{右端底域应核为 $k(\eta)$，不是 $k(y)$。局部环有到剩余域 $k(y)$ 的商映射，却一般没有相反方向的典范系数域嵌入；因而底本右端甚至没有给出所需的代数结构。}
\[
 P(\OO_{X,x},k)\Longleftrightarrow P(\OO_{Y_\eta,y},k(y)).
\]

\ednote{旁注：“$X$ 非分歧，或 $k$ 的特征为0。”〔这里按与上下文相符的意义，应理解为态射 $f$ 非分歧。〕}
\begin{proof}[原文证明]
考虑第6段的结果后，（iii）、（iv）由（ii）及（2.8，（ii））得出。另一方面，（i）已在（2.12）中证明；只须由（i）推出（ii）。令 $Z$ 为 $X$ 不光滑于 $k$ 的点集，则由（2.12），其逆像 $Z_\eta$ 正是 $Y_\eta$ 不光滑于 $k(\eta)$ 的点集。由平坦性，$Z$ 在 $X$ 中于 $x$ 处的余维数等于 $Z_\eta$ 在 $Y_\eta$ 中于 $y$ 处的余维数〔第6段的参考待补〕。因而其中一个大于 $k$ 当且仅当另一个也大于 $k$，证明完成。
\end{proof}

\egas{2.15}{推论}
设 $Z$ 为 $X$ 不在 $k$ 上光滑的点集，分别地不在 $k$ 上几何满足 $(R_k)$、不可分、非正规的点集。则 $Y_\eta$ 对 $k(\eta)$ 光滑，分别地几何满足 $(R_k)$、可分、几何正规，当且仅当 $Z$ 有限。\cmt{P04}{本句的最后一项左侧只写“非正规”，右侧却写“几何正规”，应核对是否漏掉“几何”；又若采用脚注允许的“特征0而 $f$ 任意”的版本，不能把“$f(Z)$ 有限”不加条件地换成“$Z$ 有限”。详见E05及P04。}

从编写角度看，（2.14）、（2.15）包含（2.12）、（2.13），后两者可以不再独立陈述；另一方面，推论在实际使用中比定理更重要，可以把定理放在一个预备命题或引理里，使推论更为突出。对于性质（iii），还可给出一个变体：

\egas{2.16}{推论}
（仍设 $f$ 为浸入，且 $\F$ 凝聚。）在（2.7）的条件下，$Z_i$ 不是嵌入关联分支，当且仅当 $Z_{i\eta}$ 不是嵌入关联分支。\cmt{用语}{底本写“should not be immersed”；在关联素循环语境中，这里应理解为“不是嵌入分支”，不是态射的浸入性质。由于未核对法文手稿，此处明确标示这一译校理解。}
此时，$\F$ 在 $Z_i$ 处相对于 $k$ 的根重数，等于 $\G_\eta$ 在 $Z_{i\eta}$ 处相对于 $k(\eta)$ 的根重数。
\begin{proof}[原文证明]
第一断言包含在（2.7）的最后一断言中。为证第二断言，因为 $Y_\eta\to X$ 平坦，所以 $\F\mapsto\G_\eta$ 正合；限制到 $Z_i$ 泛点的邻域并使用 dévissage（逐层分解），归结为 $\F=\OO_{Z_i}$，从而可设 $Z_i=X$。接着原英译写道：“也可先假设 $X$ 在 $k$ 上分离，归结到 $k$ 代数闭。”\ednote{无法理解。}\cmt{P01}{此句在所上传英译中语法和数学衔接均未完整；这里不擅自补成一个有效的下降论证。}
于是该断言包含在（2.15，（iii））中。像通常一样，再得到下面的推论。
\end{proof}

\egas{2.17}{推论}
设 $Z$ 为 $X$ 上 $\F$ 不对 $k$ 可分的点集〔参考待补〕。则 $\G_\eta$ 对 $k(\eta)$ 可分，当且仅当 $Z$ 有限。特别地，若 $\F$ 对 $k$ 可分，则 $\G_\eta$ 对 $k(\eta)$ 可分。

\egas{2.18}{备注}
若舍弃 $f$ 是浸入的假设，关键结果（2.12）及其推论（2.13）、（2.17）便不再成立；例如可取 $X$ 为 $P$ 的一个纯不可分覆盖。但是，当 $k$ 的特征为0时，无须假设 $f$ 为浸入，（2.12）、（2.17）也成立。
\begin{proof}[原文论证]
只须对（2.12）验证；它由（2.11）以及特征0的代数预概形上“光滑等于正则”的事实得出。
\end{proof}

% 接续：第一篇3.1；范围已在卷首与进度文件标明。

% END text/v5_01_02.tex


% BEGIN text/v5_03_06.tex
\egasection{§3\quad 泛超平面截面：几何不可约性与连通性}
\source{\texttt{EGAVaiello.pdf} 第12--13页}
\egas{3.1}{定理（Bertini--Zariski）}
假设 $\dim\overline{f(X)}\geq2$，且 $X$ 几何不可约。则泛超平面截面 $Y_\eta$ 也几何不可约。
\begin{proof}
设 $K/k$ 是 $X$ 的函数域，$n=\dim P$。在 $P$ 中选取仿射坐标 $T_1,\ldots,T_n$（选择无穷远超平面 $H^\infty$，使 $f(X)$ 不包含于其中），并在 $P^\vee$ 中取仿射坐标 $S_1,\ldots,S_n$。可见 $Y_\eta$ 的函数域 $L$ 可与整环
\[
 K[S_1,\ldots,S_n]\Big/\Bigl(\sum_{i=1}^n t_iS_i-1\Bigr)
\]
的分式域相等同，其中 $t_i\in K$ 是 $T_i$ 沿 $f:X\to P$ 的像。由于 $\dim\overline{f(X)}>0$，这些 $t_i$ 并非全在 $k$ 上代数，因而更不可能全为零；例如可设 $t_n\ne0$。于是立即得到
\[
 L=K(S_1,\ldots,S_{n-1}),
\]
这是纯超越扩张；元素 $S_n\in L$ 由方程 $\sum_it_iS_i-1=0$ 表成各 $S_i$（$1\leq i\leq n-1$）和 $t_i$（$1\leq i\leq n$）的函数。
\cmt{B2-P01}{定理只假设几何不可约，没有假设约化。这里的“函数域”应先在约化不可约预概形 $X_{\mathrm{red}}$ 上解释，或使用泛点的剩余域。这样补足不改变所研究的几何不可约性；不是把原假设偷偷加强为几何整。}
另一方面，$k'=k(\eta)$ 可与 $k(S_1,\ldots,S_n)$ 等同，而典范嵌入 $k'\to L$ 由把 $S_i$ 映到 $S_i$ 得到〔PB：核对这一点！〕。
\ednote{大概应是把 $S_i$ 映到它在 $L$ 中的等价类 $[S_i]$。}
换言之，作为 $L$ 的子扩张，$k'$ 由这些 $S_i$（$1\leq i\leq n$）生成；等价地，由 $S_1,\ldots,S_{n-1}$ 和
\[
 S_n=a_0+a_1S_1+\cdots+a_{n-1}S_{n-1},
 \qquad a_0=t_n^{-1},\quad a_i=-t_it_n^{-1}\ (1\leq i<n)
\]
生成。注意，各 $a_i$ 生成的域与各 $t_i$ 生成的域显然相同；它们在 $k$ 上的共同超越次数就是 $\overline{f(X)}$ 的维数。（附记：把这个双有理描述至少作为（2.1）的推论收入，应当是合适的。）所以（3.1）的证明归结为下面的引理。
\end{proof}

\egas{3.1.1}{引理（Zariski）}
〔参看本节末的英译编者注。〕设 $k$ 为域，$K/k$ 为有限型域扩张，$m\geq0$ 为整数，$a_i\in K$（$0\leq i\leq m$），且 $k(a_0,\ldots,a_m)$ 在 $k$ 上的超越次数至少为2。令
\[
 L=K(S_1,\ldots,S_m),\qquad
 k'=k\left(S_1,\ldots,S_m,a_0+\sum_{i=1}^m a_iS_i\right)\subset L,
\]
其中 $S_i$ 为不定元。如果 $K$ 是 $k$ 的准素扩张，则 $L$ 是 $k'$ 的准素扩张。
\ednote{primary extension 大概表示较小的域在较大的域中代数闭，或者准代数闭。参见 Jouanolou，定理3.6。}
\cmt{B2-H07}{这里保留英译编者的猜测，但它不能替代定义。所上传 EGA IV 第二部分（4.3.1），PDF第55页，明确说：$L/k$ 为准素扩张，当且仅当 $L$ 内包含的 $k$ 的最大可分代数扩张就是 $k$ 本身。它不要求一切代数元素都已在 $k$ 中；非平凡纯不可分扩张就说明二者不同。}

这一引理，或者同它如同兄弟般相似的引理，在文献中几乎随处可见。因此，我把选取一份证明出处的工作留给你；也就是说，我没有兴致凭自己的办法再找一个证明。
\cmt{B2-P02}{这句话正是上传笔记的内容；此处没有引理的证明。故（3.1）在本载体内只是归结到一个外部引理，不能宣称本节已经给出了自足的完整证明。这里不补造原文未写的证明。}

\egas{3.2}{推论}
假设 $f$ 非分歧，或者 $k$ 的特征为0，并且 $\dim\overline{f(X)}\geq2$。如果 $X$ 几何整，则 $Y_\eta$ 也几何整。
\begin{proof}
因为“几何整”等价于“几何不可约且可分”。
\end{proof}

\egas{3.3}{推论}
假设 $k$ 代数闭，对 $X$ 的每个不可约分支 $X_i$ 都有 $\dim\overline{f(X_i)}\geq2$。再设 $X$ 是 $\sigma$-连通的，其中 $\sigma$ 为满足 $\dim\overline{f(Z)}=0$ 的 $X$ 的闭子集 $Z$ 所组成的集合；这就是说，对于每个这样的 $Z$，$X-Z$ 都连通。在这些条件下，$Y_\eta$ 在 $k(\eta)$ 上几何连通。
\begin{proof}
按照一个应当与 Hartshorne 定理一起出现在第6段中的引理，\ednote{询问 A. G.。}
上述假设表示：任意两个不可约分支 $X',X''$，都能由不可约分支链
$X_0=X',\ldots,X_n=X''$ 联系起来，使相邻两项的交不属于 $\sigma$。因此它们的逆像 $X'_\eta,X''_\eta$ 由一条 $X_{i\eta}$ 链相连；由（3.1），各 $X_{i\eta}$ 在 $k(\eta)$ 上几何连通，由（2.3），相邻两项的交非空。因为 $Y_\eta$ 是所有这些 $X_{i\eta}$ 的并，故 $Y_\eta$ 在 $k(\eta)$ 上几何连通，证毕。
\cmt{B2-T06}{本段英译稿把链的首项写成“$X_0$ and $X'$”，又在末句写 $Y_\eta=X_\eta$。这里分别按同句定义明确排为 $X_0=X'$ 和 $Y_\eta=\bigcup_iX_{i\eta}$，记录这一记号整理，不将其另算为定理反例。}
\end{proof}
\noindent\textbf{英译编者关于（3.1.1）的注。}
应当与 Zariski 的\textit{论文集}（MIT Press）第1卷第174页、第2卷第304页比较；还可参看 Zariski--Samuel，Springer 的 GTM 版第1卷第196页、第2卷第230页，以及 Jouanolou 的 \textit{Théorèmes de Bertini et applications}，定理3.6和第6节。

\egasection{§4\quad 变动超平面截面的研究：“充分一般”的截面}
\source{\texttt{EGAVaiello.pdf} 第13--15页}
我们回到§1的一般情形，$S$ 为任意预概形；另假设 $X$ 在 $S$ 上有限表示。

一般说，若 $P(Z,k)$ 是域 $k$ 上代数预概形 $Z$ 的一种“可构造”性质，则满足 $P(Y_\xi,k(\xi))$ 的 $\xi\in P^\vee$ 组成可构造集合：只须注意 $Y_\xi$ 是有限表示态射 $Y\to P^\vee$ 在 $\xi$ 处的纤维，并应用第9段。\ednote{指 EGA IV。}
对于性质 $P(Z,M,k)$ 也有类似说明，其中 $Z,k$ 同前，$M$ 是 $Z$ 上的凝聚模；底本接着假设“$G$ 在 $X$ 上有限表示”，并写满足 $P(Y_\eta,G_\eta,k(\eta))$ 的 $\xi\in P^\vee$ 所成的集合可构造。
\cmt{B2-T07}{此处前后使用的载体、变动指标不一致。按§1，$F$ 在 $X$ 上，$G$ 在 $Y$ 上；而讨论随 $\xi$ 变化的纤维时，应写 $P(Y_\xi,G_\xi,k(\xi))$。保持底本措辞可看清这一问题，不能据此把一个固定泛纤维性质直接当成各纤维性质。}
另一方面，在前两节中，当 $S=\Spec k$ 时，我们已在若干情形给出了上述集合 $E$ 包含 $P^\vee$ 泛点的判据。由可构造性，$E$ 便包含非空开集：这正是从泛超平面截面的结论，过渡到“充分一般”超平面截面的相应结论。

还应注意，若 $S=\Spec k$，则有逆向表述：泛超平面截面具有性质 $P$，当且仅当 $\eta$ 的某个合适邻域内各 $Y_\xi$ 都具有该性质；只要求其中的闭点 $\xi$ 满足它，也已足够。若 $k$ 代数闭，这归结为考虑 $k$-有理点，也就是 $X$ 本身的超平面截面，而无须先扩张底域。这来自 $E$ 的可构造性以及 $P^\vee$ 是 Jacobson 概形的事实。下面列举几个适用的特殊情形。

\egas{4.2}{命题}
设 $G$ 是 $X$ 上的有限表示模。
\cmt{B2-T07续}{底本此处确实写 $G$，随后的(b)又写 $F$；这里仍标明此处原字母，后文按各自原文保留。原稿没有单独印出“4.1”条目标记，虽然以后用4.1指本节开头的讨论，译文不补造一个新的原编号。}
设 $P$ 是域 $K$ 上代数概形 $Z$ 的模 $M$ 的下列性质之一：
\begin{enumerate}[label=(\roman*)]
\item $\coprof(M)\leq n$；
\item $M$ 满足 $(S_k)$；
\item $M$ 是 Cohen--Macaulay 模；
\item $M$ 没有嵌入的相伴素循环分支；
\item $M$ 在 $K$ 上可分。
\end{enumerate}
在这些记号下，令 $E$ 为使 $G_\xi$ 满足性质 $P$ 的 $\xi\in P$ 的集合。则：(a) $E$ 是 $P^\vee$ 的可构造子集；(b) 若 $S=\Spec k$，$k$ 为域，且 $F$ 满足性质 $P$，再在(v)的情形假设 $k$ 的特征为0或者 $f:X\to P$ 非分歧，则 $E$ 包含 $P^\vee$ 的一个稠密开子集。
\cmt{B2-T07续}{这里的参数应在 $P^\vee$ 而不是 $P$ 中；(a)以及§1的参数空间定义已给出内部依据。}
\begin{proof}
(a)来自我们在 EGA IV 第9段中见到的事实：$P$ 是可构造性质。(b)则来自前两节的相应结果。
\end{proof}
\noindent\textbf{对(b)的补充念头。}更一般地，设 $Z$ 是 $X$ 上 $F$ 不满足 $P$ 的点集，只要求 $f(Z)$ 有限，即 $\dim\overline{f(Z)}\leq0$。

\egas{4.3}{命题}
设 $P$ 是域 $K$ 上代数预概形的下列性质之一：
\begin{enumerate}[label=(\roman*)]
\item $Z$ 在 $K$ 上光滑；
\item $Z$ 在 $K$ 上满足几何 $(R_k)$；
\item $Z$ 在 $K$ 上可分；
\item $Z$ 在 $K$ 上几何正规；
\item $Z$ 在 $K$ 上几何整；
\item $Z$ 在 $K$ 上几何不可约。
\end{enumerate}
令 $E$ 为使 $Y_\xi$ 满足 $P$ 的 $\xi\in P^\vee$ 的集合。则：(a) $E$ 为 $P^\vee$ 的可构造子集；(b) 假设 $S=\Spec k$，$k$ 为域；在(i)至(v)的情形，假设 $k$ 的特征为0，\emph{并且} $f:X\to P$ 非分歧。最后，在(v)、(vi)的情形假设 $\dim\overline{f(X)}\geq2$。若 $X$ 满足 $P$，则 $E$ 包含 $P^\vee$ 的一个稠密开子集。
\cmt{B2-H08}{本段英文使用 and，译文保留“并且”。前面（3.2）使用 or；把两个充分条件同时要求只是更强的假设，不能仅因其较强就断言这里的定理为假，也不应静默把 and 改成 or。}
\begin{proof}
证明与（4.2）相同。还应指出，在(i)至(v)的情形，对(b)只假设 $f(Z)$ 有限便足够，其中 $Z$ 是 $X$ 上性质 $P$ 不成立的点集。
\end{proof}

\egas{4.4}{推论}
考虑性质 $(C_m)$：“$\overline Z$ 不能被一个维数 $\leq m$ 的闭子集断开”，其中
$\overline Z=Z\times_K\overline K$，$\overline K$ 是 $K$ 的代数闭包。令 $E$ 为使 $Y_\xi$ 在 $k(\xi)$ 上满足 $(C_m)$ 的 $\xi\in P^\vee$ 的集合。则：(a) $E$ 可构造；(b) 假设 $S=\Spec k$，$k$ 为域，并且对 $X$ 的每个不可约分支 $X_i$，都有 $\dim\overline{f(X)}\geq2$。若 $X$ 在 $k$ 上满足 $(C_{m+1})$，则 $E$ 包含 $P$ 的一个稠密开子集 $U$。
\cmt{B2-E02}{底本此处还有 $f(X)$ 与 $f(X_i)$、$P$ 与 $P^\vee$ 的记号错置；但即便把这两个记号整理成预期形式，一般态射下的连通性断言仍不成立。报告给出 $X=V(uv)\subset\mathbf A^4$、$f=[1:u+v:(u+v)w]$ 的具体反例：$X$ 满足 $(C_1)$，而泛截面不满足 $(C_0)$。}

可构造性通过 AQT 得到。\ednote{AQT 是什么？询问 AG。}
这是 EGA IV 第9段里遗漏的一个事实，也许还来得及补上。(b)原则上来自（3.3），只是（3.3）的表述过于特殊，因而应作推广；证明基本不变。此外，（4.4）的陈述中有一个错误，它按现在的写法“在 $f$ 拟有限时不成立”；在一般情形中，应当不再考虑 $X$、$Y_\xi$ 中闭子集本身的维数，而考虑它们在 $P$、$H_\xi$ 中的像的维数。\textit{Redactor demerdetur.}
\ednote{拉丁文；待翻译。}
\cmt{B2-P03}{所上传英文确实写“not valid as such if $f$ is quasi-finite”，这一句与随后“在一般情形应看像维数”的说明关系可疑，可能缺少否定或限制词；未见法文手稿，不能替作者裁定。拉丁语短句也保留原字，暂不强作可靠译释。明确反例针对一般非拟有限态射，不宣称它同时推翻拟有限版本。}

\egasection{§5\quad Seidenberg 类型定理}
\source{\texttt{EGAVaiello.pdf} 第16--18页的相应部分}
\egap{5.1}
本节给出（4.1）所定义集合 $E$ 为开集的条件。这里研究的是对 $X$、分别对 $Y_\xi$ 为局部性质的 $P$，从而可以定义 $Y$ 中的点集 $U$：若 $y$ 在 $P^\vee$ 中的像为 $\xi$，则要求 $Y_\xi$ 在 $y$ 处满足 $P$，分别要求 $G_\xi$ 在 $y$ 处满足 $P$。我们先利用第12段给出 $U$ 为开集的判据。\ednote{须查明引用；很可能是 EGA IV。是的。}
如常有 $E=P^\vee-h(Y-U)$〔英译编者补字〕。因此，若 $U$ 为开集且 $X$ 在 $S$ 上固有，则因 $h$ 固有、从而闭，$E$ 也为开集。
\ednote{旁注：因为 $Y$ 在 $X$ 上固有，而 $P^\vee$ 在 $S$ 上分离。}

\egap{5.2}
如§1所见，$Y$ 在 $X\times_SP^\vee=X_{P^\vee}$ 中由
$\OO_X(1)\otimes_S\OO_{P^\vee}(1)$ 的截面 $\phi$ 的“除子”定义。对每个 $\xi\in P^\vee$，它诱导
\[
 \OO_X(1)\otimes_{k(s)}\OO_{P^\vee}(1)(\xi)
\]
的截面 $\phi_\xi$。底本说，这个层非典范地同构于
$\OO_X(1)\otimes_{k(s)}k(\xi)=\OO_{X_{k(s)}}(1)$；$Y_\xi$ 就是该截面的“除子”。
\cmt{B2-T08}{这里应把 $\OO_X(1)$ 先限制到 $X_s$，而扩张到的纤维是 $X_{k(\xi)}$，不是右端所印的 $X_{k(s)}$。这是载体及底域的记号问题。}
（附记：可逆模 $L$ 的截面 $\phi$ 的除子，定义为同态
$\phi^{-1}:L^{-1}\to\OO$ 的像理想所定义的闭子预概形。）若 $F$ 为 $X$ 上的模层，其在 $Y$ 上的逆像，即 $F\otimes_{\OO_S}\OO_{P^\vee}=F_{P^\vee}$ 在闭子概形 $Y\subset X_{P^\vee}$ 上的逆像，正是同态
\[
 \phi^{-1}\otimes\id_{F_{P^\vee}}:
 F_{P^\vee}(-1,-1)\longrightarrow F_{P^\vee}
\]
的余核；正如 Mike 所说，记号 $(-1,-1)$ 不言自明。
\ednote{我猜这里指 Mike Artin。——P. B.}
同样，$G_\xi$ 是相应同态 $F_{k(\xi)}(-1,-1)\to F_{k(\xi)}$ 的余核，其中 $\xi$ 为 $P$ 中的点；若 $\xi$ 不是 $P^\vee$ 中的普通点，而是一个取值于 $S$ 上的 $S'$ 的点，也有相应解释。\cmt{B2-T08续}{参数点所在的 $P$ 应为 $P^\vee$，同段下一句已经再次使用后者。}

一般地，若 $L$ 是某处的可逆模，$\phi$ 为定义子预概形 $V(\phi)$ 的截面，则对每个模 $F$，其在 $V(\phi)$ 上的逆像，按通常的滥用语言，与
$\id_F\otimes{}\,\cdots:F\otimes L^{-1}\to F$ 的余核等同。\ednote{底本此处有“??”，同态的另一张量因子缺失。}
若这个同态单射，则称 $\phi$ 是 $F$-正则的。底本接着说：“若选取 $F$ 与 $\OO_X$ 的一个同构，这在局部是可能的，使 $\phi$ 与 $\OO_X$ 的一个截面等同，那么这一术语与此前引入的术语相容。”
\cmt{B2-T08续}{可在局部平凡化的是可逆模 $L$，不是任意的 $F$。例如 $F=\OO_X^{\oplus2}$ 一般不可能同构于 $\OO_X$。将 $L$ 平凡化后，$\phi$ 给出函数，以此函数在 $F$ 上的乘法是否单射来定义 $F$-正则。$\phi^{-1}$ 在这里是诱导的对偶同态记号，并不表示一个处处存在的倒数函数。}

\egas{5.3}{命题}
使用上述记号，令 $U$ 为 $X_P$ 中满足下列条件的点 $x$ 的集合：若 $x$ 在 $P^\vee$ 中的像为 $\xi$，则 $\phi_\xi$ 在 $x$ 处是 $F_{k(\xi)}$-正则的。于是：
\begin{enumerate}[label=(\alph*)]
\item 若 $F$ 有限表示且相对 $S$ 平坦，则 $U$ 为开集，而 $G|_U$ 相对 $P^\vee$ 平坦。
\item 对每个 $s\in P^\vee$，若 $\eta$ 为 $P^\vee_s$ 的泛点，则 $U$ 包含 $X_{k(\eta)}$。
\end{enumerate}
\cmt{B2-T09}{底本在此写 $X_P$，但参数化的空间为 $X_{P^\vee}$；(b)中要谈 $P^\vee$ 在 $s$ 上的纤维，应有 $s\in S$。此外 $G$ 原定义在 $Y$ 上，写 $G|_U$ 时应理解为限制到 $Y\cap U$，或先沿闭浸入直接像延拓到 $X_{P^\vee}$。}
\begin{proof}
(a) 因为 $F_{P^\vee}$ 有限表示且相对 $P^\vee$ 平坦，结论来自11.3；若 $S$ 局部诺特，也来自 $0_{\mathrm{III}}\ldots$。\ednote{查明这一引文。指 EGA IV。}

(b) 可设 $S=\Spec k$。由第3段，$F_{k(\eta)}$ 的相伴循环是 $F$ 的相伴循环 $Z_i$ 的逆像。\ednote{字迹不清；询问 AG 或自行辨认——大概是 $\phi^{-1}$。}
若 $f(Z_i)$ 有限，则由（2.3），$Z_{i,k(\eta)}\cap Y=\varnothing$；否则，例如由（2.6），也有
$Z_{i,k(\eta)}\cap Y=Z_{i,k(\eta)}$（出于维数的理由；已经在（2.6）中用过的（2.3），无疑是更好的理由）。这说明 $\phi$ 不在任何 $Z_{i,k(\eta)}$ 上恒等于零，因而证明(b)。
\cmt{B2-T10}{上一句等号与下一句结论相反；这里应是 $Z_{i,k(\eta)}\cap Y\ne Z_{i,k(\eta)}$。等号意味着整个相伴循环被包含在零点集内，恰不能证明非零因子性。译文保留所上传英文的等号，并在此明示问题。}
\end{proof}

\egas{5.4}{推论}
令 $V$ 为使 $\phi_\xi$ 是 $F_{k(\xi)}$-正则的 $\xi\in P^\vee$ 的集合。若 $F$ 有限表示，则 $V$ 可构造，并包含 $P^\vee$ 在 $S$ 上各纤维的泛点。若另假设 $X$ 在 $S$ 上固有且 $F$ 在 $S$ 上平坦，则 $V$ 为开集。

\egas{5.5}{备注}
取位于 $s\in S$ 上方的 $\xi\in P^\vee$，并假设 $F_s$ 没有嵌入相伴循环。底本写道：“于是立即看到 $\xi\in V$（使用（5.4）的记号），这也就是说 $\Supp F_{k(\xi)}$ 的每个不可约分支都不位于 $H_\xi$ 上”；还说，在此判据中以 $k(\xi)$ 的任意扩张替换 $k(\xi)$，也成立，不过这一点稍不明显。
\cmt{B2-E03}{按字面，“没有嵌入相伴循环”不推出任意 $\xi\in V$。取 $F=i_*\OO_L$ 为 $\mathbf P^2$ 中直线 $L$ 的结构层，选 $H_\xi=L$，截面在 $F$ 上的乘法为零，因而不正则。合理的读法是把后半句作为 $\xi\in V$ 的充要条件，而不是从无嵌入循环自动推出它；详见报告。}
注意，刚才对 $F_s$ 所加的 $(S_1)$ 假设，在 $F_s$ 为 Cohen--Macaulay 时尤其成立，更不用说 $F$ 相对 $S$ 为 Cohen--Macaulay 的情形。笔记接着说，在这种情况下 $G_s$ 为 Cohen--Macaulay，因为局部上它由同样为 Cohen--Macaulay 的 $F_{k(s)}$ 模去 $\phi\cdot F_{k(s)}$ 得到，其中 $\phi$ 对 $F_{k(s)}$ 正则。上述说明同样应该、并且将会在前面作局部表述，以刻画 $U$ 的点而不是 $V$ 的点。
\cmt{B2-E03续}{这里使用模去正则元素的结论，仍须保留“截面正则”这一条件；不能只靠 Cohen--Macaulay 性就略去它。纤维下标 $s$ 与 $\xi$ 也应按具体基变换区分。}

现在利用12.1.1与12.1.4，得到：\ednote{询问 AG 这些引文——大概是 EGA IV。12.1.4核对不上。}
\egas{5.6}{定理}
假设 $F$ 有限表示且相对 $S$ 平坦。令 $P$ 为12.1.1中(i)至(viii)的性质之一，或者在假设 $F=\OO_X$ 时，为 EGA IV 的12.1.4中(i)至(iv)的性质之一〔英译编者补入卷号〕。以 $U_P$ 表示 $X_P$ 中的点 $x$ 的集合：若 $\xi$ 是其在 $P^\vee$ 中的像，则要求 $G_\xi$（分别地，$Y_\xi$）在 $x$ 处满足 $P$，且 $\phi_\xi$ 在 $x$ 处为 $F_{k(\xi)}$-正则。则 $U_P$ 为开集，且 $G|_{U_P}$ 相对 $S$ 平坦。
\begin{proof}
按定义，$U_P\subset U$（使用（5.3，a）的记号）；把第12段应用于
$U\to P^\vee$ 和 $F_{P^\vee}|_U$ 即得。
\cmt{B2-P04}{英译编者已经明确说引文核对不上，所列性质未在这里实际展开；证明最后写 $F_{P^\vee}|_U$，而结论涉及除子上的余核 $G$，也需与引用的原命题逐项对应。本次不把未确定的八项、四项性质凭常识补进正文。}
\end{proof}

\egas{5.7}{推论}
假设 $F$ 有限表示、相对 $S$ 平坦，并且 $\Supp F$ 在 $S$ 上固有（例如 $X$ 在 $S$ 上固有）。令 $V_P$ 为 $\xi\in P^\vee$ 的集合，使 $G_\xi$（分别地，底本此处写 $Y$）满足性质 $P$，并且“它是 $F_{k(\xi)}$-正则的”。在这些条件下，$V_P$ 为开集；而且在任何情况下，即不加平坦性或固有性假设时，它都是可构造的。
\cmt{B2-P04续}{“正则”的主语应为截面 $\phi_\xi$；谈该参数对应的纤维时应为 $Y_\xi$。此处仍受（5.6）未明确列出性质清单的限制，未据这一句话作超出来源的普遍推广。}

从表述的角度说，我觉得不能把（5.6）保留为只指向另一本中的条件编号，而应明确列出我们要用的性质(i)、(ii)、……。还应指出，也许可以在（5.1）中说，“几何正规”的情形（当然 $S=\Spec k$）是 Seidenberg 的结果。\ednote{或者意为“为稳妥起见”。}

\egasection{§6\quad 任意超平面截面的连通性}
\source{\texttt{EGAVaiello.pdf} 第18页}
这里把已知的泛超平面截面几何连通性判据（3.3），同 Zariski 连通性定理结合，以得到任意超平面截面的连通性结论。

\egas{6.1}{命题}
设 $S=\Spec k$，$k$ 为代数闭域〔假设 $X$ 在 $k$ 上固有〕；\ednote{这一处字迹不清。}
并设对 $X$ 的每个不可约分支 $X_i$，都有 $\dim\overline{f(X_i)}\geq2$。最后，假设 $X$ 不能被满足 $\dim\overline{f(Z)}\leq0$ 的闭子集 $Z$ 断开。在这些条件下，对每个 $\xi\in P^\vee$，$Y_\xi$ 都几何连通。
\cmt{B2-H09}{方括号中的固有性出自上传英译稿的补读，不能省去。报告给出删去它后特殊超平面截面不连通的例子；该例不是对保留固有性版本的反例。}
\begin{proof}
由于没有一个 $f(X_i)$ 为有限集，可见 $Y$ 的每个不可约分支 $Y_i$ 都占优于 $P^\vee$；另一方面，$Y\to P^\vee$ 固有。底本在括号中解释说：“如果 $Y$ 在 $k$ 上固有，因为它在 $X$ 上固有，而 $X$ 在 $k$ 上固有。”由（3.3），$Y\to P^\vee$ 的泛纤维 $Y_\eta$ 几何连通。最后，$P^\vee$ 正则，因而几何单支。应用15.6.3——Zariski 连通性定理的一个变体——就能推出 $Y\to P^\vee$ 的所有纤维几何连通，证毕。\ednote{指 EGA IV。}
\end{proof}
实际上，用同类证明，不难在（4.4）所示的方向上推广（6.1）。如果你不愿为这个练习费心，至少也应把它作为备注提一下。还应说清，（6.1）并不排除维数过大的超平面截面。从编排的角度看，把所有连通性问题，包括（3.3）、（4.4），集中在同一节中，也许更清楚。



% END text/v5_03_06.tex

\backmatter
\part*{各原分册的书末资料}
\addcontentsline{toc}{part}{各原分册的书末资料}

% BEGIN backmatter/volume_I.tex
\clearpage\makeatletter\@mainmatterfalse\makeatother

% BEGIN bibliography.tex
\renewcommand{\bibname}{第一卷原著参考文献}
\begin{thebibliography}{99}
\addcontentsline{toc}{chapter}{第一卷原著参考文献}
\markboth{第一卷原著参考文献}{第一卷原著参考文献}
\small
\bibitem[1]{b1}
H. Cartan 与 C. Chevalley：巴黎高等师范学校讨论班，第 8 年（1955—1956），代数几何。
\textit{Séminaire de l'École Normale Supérieure}, 8\textsuperscript{e} année (1955–56),
\textit{Géométrie algébrique}.

\bibitem[2]{b2}
H. Cartan 与 S. Eilenberg：\textit{Homological Algebra}（《同调代数》），
Princeton Math. Series，Princeton University Press，1956。

\bibitem[3]{b3}
W. L. Chow 与 J. Igusa：Cohomology theory of varieties over rings（环上簇的上同调理论），
\textit{Proc. Nat. Acad. Sci. U.S.A.}，第 XLIV 卷（1958），1244–1248 页。

\bibitem[4]{b4}
R. Godement：\textit{Théorie des faisceaux}（《层论》），
Actual. Scient. et Ind.，第 1252 号，Paris：Hermann，1958。

\bibitem[5]{b5}
H. Grauert：Ein Theorem der analytischen Garbentheorie und die Modulräume komplexer Strukturen
（解析层论的一个定理与复结构的模空间），
\textit{Publ. Math. Inst. Hautes Études Scient.}，第 5 号，1960。

\bibitem[6]{b6}
A. Grothendieck：Sur quelques points d'algèbre homologique（同调代数的若干问题），
\textit{Tôhoku Math. Journ.}，第 IX 卷（1957），119–221 页。

\bibitem[7]{b7}
A. Grothendieck：Cohomology theory of abstract algebraic varieties（抽象代数簇的上同调理论），
\textit{Proc. Intern. Congress of Math.}，103–118 页，Edinburgh（1958）。

\bibitem[8]{b8}
A. Grothendieck：Géométrie formelle et géométrie algébrique（形式几何与代数几何），
\textit{Séminaire Bourbaki}，第 11 年（1958—1959），第 182 讲。

\bibitem[9]{b9}
M. Nagata：A general theory of algebraic geometry over Dedekind domains
（Dedekind 整环上代数几何的一般理论），\textit{Amer. Math. Journ.}：
I，第 LXXVIII 卷，78–116 页（1956）；II，第 LXXX 卷，382–420 页（1958）。

\bibitem[10]{b10}
D. G. Northcott：\textit{Ideal theory}（《理想论》），Cambridge University Press，1953。

\bibitem[11]{b11}
P. Samuel：\textit{Commutative algebra}（《交换代数》；D. Herzig 记录），Cornell University，1953。

\bibitem[12]{b12}
P. Samuel：\textit{Algèbre locale}（《局部代数》），Mém. Sci. Math.，第 123 号，Paris，1953。

\bibitem[13]{b13}
P. Samuel 与 O. Zariski：\textit{Commutative algebra}（《交换代数》），两卷，
New York：Van Nostrand，1958–60。

\bibitem[14]{b14}
J.-P. Serre：Faisceaux algébriques cohérents（凝聚代数层），
\textit{Ann. of Math.}，第 LXI 卷（1955），197–278 页。

\bibitem[15]{b15}
J.-P. Serre：Sur la cohomologie des variétés algébriques（论代数簇的上同调），
\textit{Journ. de Math.} (9)，第 XXXVI 卷（1957），1–16 页。

\bibitem[16]{b16}
J.-P. Serre：Géométrie algébrique et géométrie analytique（代数几何与解析几何），
\textit{Ann. Inst. Fourier}，第 VI 卷（1955–56），1–42 页。

\bibitem[17]{b17}
J.-P. Serre：Sur la dimension homologique des anneaux et des modules noethériens
（论诺特环与诺特模的同调维数），
\textit{Proc. Intern. Symp. on Alg. Number Theory}，176–189 页，Tokyo—Nikko，1955。

\bibitem[18]{b18}
A. Weil：\textit{Foundations of algebraic geometry}（《代数几何基础》），
Amer. Math. Soc. Coll. Publ.，第 29 号，1946。

\bibitem[19]{b19}
A. Weil：Numbers of solutions of equations in finite fields（有限域中方程的解的数目），
\textit{Bull. Amer. Math. Soc.}，第 LV 卷（1949），497–508 页。

\bibitem[20]{b20}
O. Zariski：\textit{Theory and applications of holomorphic functions on algebraic varieties over arbitrary ground fields}
（任意基域上的代数簇的全纯函数理论及应用），Mem. Amer. Math. Soc.，第 5 号（1951）。

\bibitem[21]{b21}
O. Zariski：A new proof of Hilbert's Nullstellensatz（Hilbert 零点定理的一个新证明），
\textit{Bull. Amer. Math. Soc.}，第 LIII 卷（1947），362–368 页。

\bibitem[22]{b22}
E. Kähler：Geometria Arithmetica（算术几何），\textit{Ann. di Mat.} (4)，第 XLV 卷（1958），1–368 页。
\end{thebibliography}


% END bibliography.tex


% BEGIN notation_index.tex
\chapter*{第一卷原著记号索引}
\addcontentsline{toc}{chapter}{第一卷原著记号索引}
\markboth{第一卷原著记号索引}{第一卷原著记号索引}
\src{217—218}
\noindent\textbf{译注。}以下按原著次序译出全部索引条目。右栏是原索引所列的章、条目号，不是本译稿页码；原索引中个别沿用旧编号的引文仍予保留，不静默改写。查阅时须与正文条目核对。
\begingroup\small
\renewcommand{\arraystretch}{1.15}
\begin{longtable}{@{}p{0.78\linewidth}p{0.17\linewidth}@{}}
\toprule 记号及说明&原索引引文\\\midrule\endfirsthead
\toprule 记号及说明（续）&原索引引文\\\midrule\endhead
\bottomrule\endfoot
$M_{[\varphi]}$（$M$ 为 $B$-模，$\varphi:A\to B$ 为环同态）&0，1.0.2\\
$BJ,JB$（$J$ 为环 $A$ 的理想，已给定 $A\to B$ 的同态）&0，1.0.3\\
$r(\mathfrak a)$（$\mathfrak a$ 为理想）&0，1.1.1\\
$\mathfrak R(A)$（$A$ 为环）&0，1.1.2\\
$S_f$（$f$ 为交换环 $A$ 的元素）&0，1.2.1\\
$S^{-1}A$，$S^{-1}M$，$m/s$，$i_A^S$，$i_M^S$，$i^S$（$A$ 为环，$M$ 为 $A$-模，$S$ 为 $A$ 的乘法子集）&0，1.2.2\\
$A_f,M_f,A_{\mathfrak p},M_{\mathfrak p}$（$M$ 为 $A$-模，$f\in A$，$\mathfrak p$ 为 $A$ 的素理想）&0，1.2.3\\
$S^{-1}u$（$u$ 为 $A$-模同态，$S$ 为 $A$ 的乘法子集）&0，1.3.1\\
$\rho_A^{T,S},\rho_M^{T,S},\rho^{T,S}$（$M$ 为 $A$-模，$S,T$ 为 $A$ 的乘法子集）&0，1.4.1\\
$\Supp(M)$（$M$ 为 $A$-模）&0，1.7.1\\
$\F|U,u|U$（$\F$ 为 $X$ 上的层，$u$ 为 $X$ 上的层态射，$U$ 为 $X$ 的开集）&0，3.1.5\\
$\F_x,s_x,\Gamma(U,\F),u(s),\Supp(\F)$（$\F$ 为 $X$ 上的集合层，$x\in X$，$U$ 为开集，$s\in\Gamma(U,\F)$，$u$ 为层同态）&0，3.1.6\\
$\psi_*(\F)$（$\F$ 为 $X$ 上的层，$\psi:X\to Y$ 连续）&0，3.4.1\\
$\psi_*(u)$（$u$ 为 $X$ 上的层同态）&0，3.4.2\\
$\psi_x$&0，3.4.4\\
$\G\to\F$（$\F$ 为 $X$ 上的层，$\G$ 为 $Y$ 上的层）&0，3.5.1\\
$u^\sharp,v^\flat,\rho_{\G}$&0，3.5.3\\
$\psi^*(\G),\psi^*(v),\sigma_{\F}$&0，3.5.5\\
$\OO_X,\OO_{X,x},\OO_x,1,e$（$X$ 为赋环空间）&0，4.1.1\\
$\dual\F,\bigwedge^p\F$（$\F$ 为 $\OO_X$-模层）&0，4.1.4\\
$\J\F$（$\J$ 为 $\OO_X$ 的理想层，$\F$ 为 $\OO_X$-模层）&0，4.1.5\\
$\Psi_*(\F),\Psi_*(u)$（$\F$ 为 $\OO_X$-模层，$u$ 为 $\OO_X$-模层同态）&0，4.2.1\\
$\Psi_*(\C)$（$\C$ 为 $\OO_X$-代数层）&0，4.2.4\\
$\Psi^*(\G),\Psi^*(v)$（$\G$ 为 $\OO_Y$-模层，$v$ 为 $\OO_Y$-模层同态）&0，4.3.1\\
$\Psi^*(\C)$（$\C$ 为 $\OO_Y$-代数层）&0，4.3.4\\
$\Psi^*(\J)\OO_X,\J\OO_X$（$\J$ 为 $\OO_Y$ 的理想）&0，4.3.5\\
$\G\to\F$（$\F$ 为 $\OO_X$-模层，$\G$ 为 $\OO_Y$-模层）&0，4.4.1\\
$u^\sharp_\Psi,u^\sharp,v^\flat_\Psi,v^\flat,\rho_{\G},\sigma_{\F}$&0，4.4.3\\
$u_1\otimes u_2$（$u_1,u_2$ 为从 $\OO_Y$-模层到 $\OO_X$-模层的同态）&0，4.4.4\\
$\LL^{-1}$（$\LL$ 为可逆 $\OO_X$-模层）&0，5.4.3\\
$\LL^{\otimes n}$（$\LL$ 为可逆 $\OO_X$-模层）&0，5.4.4\\
$\Gamma_*(\LL),\Gamma_*(\LL,\F)$（$\LL$ 可逆，$\F$ 为 $\OO_X$-模层）&0，5.4.6\\
$\OO_X^*$&0，5.4.7\\
$\mathfrak m_x,k(x),f(x)$&0，5.5.1\\
$X_f$&0，5.5.2\\
$\Im(M'\otimes_A N')$（$M',N'$ 分别为 $M,N$ 的子 $A$-模）&0，6.0\\
$\widehat A,\widehat M$&0，7.2.3及7.3.1\\
$A\{T_1,\ldots,T_r\}$&0，7.5.1\\
$A\{S^{-1}\}$&0，7.6.1\\
$\mathfrak a\{S^{-1}\}$（$\mathfrak a$ 为 $A$ 的开理想）&0，7.6.9\\
$A_{\{f\}},\mathfrak a_{\{f\}}$&0，7.6.14\\
$A_{\{S\}}$&0，7.6.14\\
$\widehat{M\otimes_A N},M\ctensor_A N$&0，7.7.1\\
$u\ctensor v$&0，7.7.3\\
\midrule
$\Spec(A),\jj_x,\mathfrak m_x,k(x),f(x),M_x$，$r(E),V(E),V(f),D(f)$（$A$ 为环，$M$ 为 $A$-模，$f\in A$，$E\subset A$，$x\in\Spec(A)$）&I，1.1.1\\
$\jj(Y)$（$Y\subset\Spec(A)$）&I，1.1.3\\
${}^{a}\varphi$（$\varphi$ 为环同态）&I，1.2.1\\
$S_f'$（$f$ 为环的元素）&I，1.3.1\\
$\rho_{g,f}$（$f,g$ 为环的元素）&I，1.3.3\\
$\wt A,\wt M,\theta_f$（$A$ 为环，$f\in A$，$M$ 为 $A$-模）&I，1.3.4\\
$\wt u$（$u$ 为 $A$-模同态）&I，1.3.5\\
$\wt\varphi$（$\varphi$ 为环同态）&I，1.6.1\\
$A(X)$（$X$ 为仿射概形）&I，1.7.1\\
$\OO_{X/Y}$（$X$ 为预概形）&I，2.1.6\\
$\Hom(X,Y)$（$X,Y$ 为预概形）&I，2.2.1\\
$\Hom_S(X,Y),1_X$（$X,Y$ 为 $S$-预概形）&I，2.5.2\\
$\Gamma(X/S)$（$S$ 为预概形，$X$ 为 $S$-预概形）&I，2.5.5\\
$X\amalg Y$（$X,Y$ 为预概形）&I，3.1.1\\
$X\times_S Y,X\times Y,(g,h)_S,u\times_Sv,u\times v$（$X,Y$ 为 $S$-预概形，$g,h,u,v$ 为 $S$-态射）&I，3.2.1\\
$X\times_A Y,X\otimes_A B,(g,h)_A,u\times_Av$（$A$ 为环，$X,Y$ 为 $A$-预概形，$B$ 为 $A$-代数，$g,h,u,v$ 为 $A$-态射）&I，3.2.1\\
$X_{(S')}$（$X,S'$ 为 $S$-预概形）&I，3.3.6\\
$f_{(S')}$（$S'$ 为 $S$-预概形，$f$ 为 $S$-态射）&I，3.3.7\\
$\Gamma_f$（$f$ 为 $S$-态射）&I，3.3.14\\
$X(T)$（$X,T$ 为预概形）&I，3.4.1\\
$P\times_R Q$（$P,Q$ 为 $R$ 上的集合）&I，3.4.2\\
$X(T)_S$（$X,T$ 为 $S$-预概形）&I，3.4.3\\
$X(B),X(B)_A$（$X$ 为 $A$-预概形，$B$ 为 $A$-代数）&I，3.4.4\\
$X\otimes_Y B,X\otimes_{\OO_y}B$（$y\in Y$，$B$ 为 $\OO_y$-代数）&I，3.6.3\\
$Z\le Y$（$Y,Z$ 为一个预概形的子预概形）&I，4.1.10\\
$f^{-1}(Y)$（$f$ 为态射，$Y$ 为子预概形）&I，4.4.1\\
$\mathscr N_X$（$X$ 为预概形）&I，5.1.1\\
$X_{\mathrm{red}}$（$X$ 为预概形）&I，5.1.3\\
$f_{\mathrm{red}}$（$f$ 为态射）&I，5.1.5\\
$\Delta_{X/S},\Delta_X,\Delta$（$X$ 为 $S$-预概形）&I，5.3.1\\
$\operatorname{rg}_K(X)$（$K$ 为域，$X$ 为有限 $K$-概形）&I，6.4.5\\
$n(X)$（$X$ 为域上的有限概形）&I，6.4.8\\
$\Gamma_{\mathrm{rat}}(X/Y)$&I，7.1.2\\
$R(X)$（$X$ 为预概形）&I，7.1.3\\
$\mathscr R(X)$（$X$ 为预概形）&I，7.3.1\\
$L(A)$（$A$ 为整环）&I，8.1.2\\
$\delta(f)$（$f$ 为有理映射）&I，8.2.1\\
$\F\otimes_{\OO_S}\G,\F\otimes_S\G$（$\F,\G$ 为 $S$-预概形上的模层）&I，9.1.2\\
$\overline\G$（$\G$ 为 $\OO_X$-模层）&I，9.4.1\\
$\overline Y$（$Y$ 为子预概形）&I，9.5.11\\
$\Spf(A),\OO_{\frX}$（$A$ 为可容许环，$\frX=\Spf(A)$）&I，10.1.2\\
$\frD(f)$（$f$ 为可容许环的元素）&I，10.1.4\\
${}^{a}\varphi,\wt\varphi$（$\varphi$ 为可容许环的连续同态）&I，10.2.1\\
$J^\Delta$（$J$ 为定义理想）&I，10.3.1\\
$\frX\times_{\frS}\frY$（$\frX,\frY$ 为 $\frS$-形式预概形）&I，10.7.3\\
$\F_{/X'},\widehat\F,u_{/X'},\widehat u$（$\F$ 为 $\OO_X$-模层，$u$ 为模层同态，$X'$ 为 $X$ 的闭子集）&I，10.8.4\\
$X_{/X'},\widehat X$（$X$ 为预概形，$X'$ 为闭子集）&I，10.8.5\\
$\widehat f$（$f$ 为预概形态射）&I，10.9.1\\
$M^\Delta,u^\Delta$（$M$ 为进制环 $A$ 上的模，$u$ 为 $A$-模的连续同态）&I，10.10.1\\
$\Delta_{\frX/\frS},\Delta_{\frX}$（$\frX$ 为 $\frS$-形式预概形）&I，10.15.1\\
\end{longtable}
\endgroup


% END notation_index.tex

\par

% BEGIN terminology_index.tex
\chapter*{第一卷原著术语索引}
\addcontentsline{toc}{chapter}{第一卷原著术语索引}
\markboth{第一卷原著术语索引}{第一卷原著术语索引}
\src{219—224}
\noindent\textbf{译注。}本索引依原著法文词条的次序逐条翻译；右栏保留原著所列的章号及条目号，不是中文页码。原索引有若干旧编号或误引，例如“闭点”所引的（0，2.2.6）在本卷正文中并不存在；这里不静默改动原引文。查阅时请同时参照本书目录和正文。
\begingroup\small
\renewcommand{\arraystretch}{1.16}
\begin{longtable}{@{}p{0.76\linewidth}p{0.19\linewidth}@{}}
\toprule 术语&原索引引文\\\midrule\endfirsthead
\toprule 术语（续）&原索引引文\\\midrule\endhead
\bottomrule\endfoot
子预概形的闭包&I，9.5.11\\
$\OO_X$-代数层&0，4.1.3\\
进制环；$J$-进制环&0，7.1.9\\
可容许环&0，7.1.2\\
完备分式环&0，7.6.5\\
分式环&0，1.2.2\\
有理函数环&I，2.1.6；7.1.3\\
仿射概形的环&I，1.7.1\\
整环&0，1.0.6\\
线性拓扑环&0，7.1.1\\
局部环&0，1.0.7\\
支配局部环&I，8.1.1\\
$X$ 沿 $Y$ 的局部环；$Y$ 在 $X$ 中的局部环&I，2.1.6\\
预进制环；$J$-预进制环&0，7.1.9\\
预可容许环&0，7.1.2\\
约化环&0，1.1.1\\
正则环&0，4.1.3\\
相关的局部环&I，8.1.4\\
$\OO_X$-模层的零化理想&0，5.3.7\\
与环同态相伴的环谱映射&I，1.2.1\\
有理映射；$S$-有理映射&I，7.1.2\\
在一点有定义的有理映射&I，7.2.1\\
在开集上诱导的有理映射&I，7.1.2\\
在 $\Spec(\OO_x)$ 上诱导的有理映射&I，7.2.8\\
凝聚的 $\OO_X$-模层&0，5.3.1\\
凝聚的 $\OO_X$-代数层&0，5.3.6\\
$\OO_X$-模层及其同态沿闭子集的完备化&I，10.8.4\\
预概形沿闭子集的完备化&I，10.8.5\\
不可约分支&0，2.1.5\\
$\psi$-态射与 $\psi'$-态射的复合&0，3.5.2\\
$\Psi$-态射与 $\Psi'$-态射的复合&0，4.4.2\\
粘合条件&0，3.3.1；4.1.6\\
代数预概形的基域&I，6.4.1\\
几何点的值域&I，3.4.5\\
$X\times_S X$ 的对角线&I，5.3.9\\
双同态&0，1.0.2\\
有理映射的定义域&I，7.2.1\\
$\OO_X$-模层的对偶&0，4.1.4\\
拓扑幂零元素&0，7.1.1\\
由一族截面生成的 $\OO_X$-模层&0，5.1.2\\
截面为零的点集&0，5.5.1\\
整代数；有限整代数&0，1.0.5\\
赋环空间&0，4.1.1\\
局部环赋环空间&0，5.5.1\\
在开集上诱导的赋环空间&0，4.1.2\\
正规的、约化的、正则的赋环空间&0，4.1.3\\
由粘合得到的赋环空间&0，4.1.6\\
柯尔莫哥洛夫空间&0，2.1.2\\
不可约空间&0，2.1.1\\
诺特空间&0，2.2.1\\
拟紧空间&0，2.2.4\\
赋环空间的底空间&0，4.1.1\\
拓扑赋环空间&0，4.1.1\\
赋环空间上的代数层&0，4.1.3\\
与预层相伴的层&0，3.5.7\\
$\Spec(A)$ 上与一个 $A$-模相伴的层&I，1.3.4\\
取值于一个范畴的层&0，3.1.1\\
凝聚环层&0，5.3.5\\
在一点正规的、正规的、在一点约化的、约化的、在一点正则的、正则的环层&0，4.1.3\\
分次环层&0，4.1.3\\
有理函数层&I，7.3.1\\
挠层&I，7.4.1\\
理想层&0，4.1.3\\
定义理想层&I，10.3.3；10.5.1\\
诱导层&0，3.7.1\\
局部常值层&0，3.6.1\\
由粘合得到的层&0，3.3.1\\
伪离散层&0，3.8.1\\
常值层&0，3.6.1\\
赋环空间的结构层&0，4.1.1\\
仿射概形的结构层&I，1.3.4\\
开集基上的层&0，3.2.2\\
层的茎&0，3.1.6\\
预概形态射的纤维&I，3.6.3\\
有限代数&0，1.0.5\\
有理函数&I，7.1.2\\
一点的一般化&0，2.1.2\\
分次 $\OO_X$-模层&0，4.1.3\\
态射的图&I，5.3.11\\
拓扑环层的连续同态&0，3.1.4\\
由截面定义的同态&0，5.1.1\\
局部环的局部同态&0，1.0.7\\
模的 $\varphi$-同态&0，1.0.2\\
可容许环的定义理想&0，7.1.2\\
$\OO_X$-理想&0，4.1.3\\
形式预概形 $\frX$ 的 $\OO_{\frX}$-定义理想&I，10.3.3；10.5.1\\
素理想&0，1.0.6\\
$\OO_X$-模层的直像&0，4.2.1\\
预层的直像&0，3.4.1\\
$S$-截面的像&I，5.3.11\\
预概形在态射下的闭像&I，9.5.3\\
$\OO_Y$-模层的逆像&0，4.3.1\\
$S$-态射的逆像&I，3.3.7\\
预层的逆像&0，3.5.3\\
$S$-预概形的逆像&I，3.3.6\\
子预概形的逆像&I，4.4.1\\
预概形的浸入、闭浸入、开浸入&I，4.2.1\\
形式预概形的闭浸入&I，10.14.2\\
预概形的局部浸入&I，4.5.1\\
开集上诱导的赋环空间的典范嵌入&0，4.1.2\\
子预概形的典范嵌入&I，4.1.7\\
形式子预概形的典范嵌入&I，10.14.2\\
几何单射&I，3.5.4\\
可逆 $\OO_X$-模层的逆&0，5.4.3\\
可逆 $\OO_X$-模层&0，5.4.1\\
与浸入相伴的同构&I，4.2.1\\
预概形的局部同构&I，4.5.1\\
$\OO_{X_n}$-模层的射影极限&I，10.6.6\\
局部自由 $\OO_X$-模层&0，5.4.1\\
取值于局部环的点的落点&I，3.4.5\\
分式模&0，1.2.2\\
忠实平坦模&0，6.4.1\\
平坦模&0，6.1.1\\
$A$-平坦模&0，6.2\\
拟有限模&0，7.4.1\\
$\OO_X$-模层&0，4.1.3\\
赋环空间、拓扑赋环空间的态射&0，4.1.1\\
忠实平坦态射&0，6.7.8\\
平坦态射&0，6.7.1\\
开集基上所定义预层的态射&0，3.2.3\\
预概形态射&I，2.2.1\\
$S$-预概形、$A$-预概形的态射&I，2.5.2\\
双有理态射&I，6.5.6\\
有限型态射&I，6.3.1\\
对角态射&I，5.3.1\\
支配态射&I，2.2.6\\
闭态射&I，2.2.6\\
一个态射的图态射&I，3.3.14\\
局部有限型态射&I，6.6.2\\
不大于另一态射的态射&I，4.1.8\\
开态射&I，2.2.6\\
拟紧态射&I，6.6.1\\
根式态射&I，3.5.4\\
约化态射&I，5.1.5\\
分离态射&I，5.4.1\\
$S$-预概形的结构态射&I，2.5.1\\
满态射&I，2.2.6\\
普遍单射态射&I，3.5.4\\
等价的态射&I，7.1.1\\
预概形的 $A$-态射、$S$-态射&I，2.5.2\\
形式预概形态射&I，10.4.5\\
形式预概形的进制态射&I，10.12.1\\
形式预概形的对角态射&I，10.15.1\\
形式预概形的有限型态射&I，10.13.3\\
形式预概形的分离态射&I，10.15.1\\
$\frS$-形式预概形的结构态射&I，10.4.7\\
形式预概形的 $A$-态射、$\frS$-态射&I，10.4.7\\
预层的 $\psi$-态射&0，3.5.1\\
从 $\OO_Y$-模层到 $\OO_X$-模层的 $\Psi$-态射&0，4.4.1\\
环的幂零根&0，1.1.1\\
$\OO_X$-代数层的幂零根&I，5.1.1\\
预概形的几何点数&I，6.4.8\\
仿射开集&I，2.1.1\\
仿射形式开集、进制仿射形式开集、诺特仿射形式开集&I，10.4.1\\
环的乘法子集&0，1.2.1\\
饱和乘法子集&0，1.4.3\\
$f$-平坦的 $\OO_X$-模层&0，6.7.1\\
预概形的环值点&I，3.4.4\\
预概形的预概形值点&I，3.4.1\\
$A$-预概形的 $A$-代数值点&I，3.4.4\\
$S$-预概形在 $s\in S$ 上方的点&I，2.5.1\\
$S$-预概形的 $S$-预概形值点&I，3.4.3\\
闭点&0，2.2.6\\
泛点&0，2.1.2\\
预概形的几何点&I，3.4.5\\
$s$ 上方的几何点&I，3.4.5\\
落在 $x$ 的几何点&I，3.4.5\\
$K$ 上的有理点&I，3.4.5\\
常值预层&0，3.6.1\\
开集基上的预层&0，3.2.1\\
预概形&I，2.1.2\\
阿廷预概形&I，6.2.1\\
连通预概形&I，2.1.7\\
基预概形&I，2.5.1\\
模 $J$ 约化所得的预概形&I，3.7.1\\
在开集上诱导的预概形&I，2.1.8\\
整预概形&I，2.1.7\\
不可约预概形&I，2.1.7\\
局部整预概形&I，2.1.7\\
局部诺特预概形&I，6.1.1\\
与预概形相伴的约化预概形&I，5.1.3\\
$A$-预概形；$A$ 上的预概形&I，2.5.1\\
代数 $K$-预概形&I，6.4.5\\
$S$-预概形；$S$ 上的预概形&I，2.5.1\\
支配的 $S$-预概形&I，2.5.1\\
$S$ 上有限型预概形；有限型 $S$-预概形&I，6.3.1\\
由基预概形扩张得到的预概形&I，3.3.6\\
在 $S$ 上分离的预概形&I，5.4.1\\
形式预概形&I，10.4.2\\
进制形式预概形&I，10.4.2\\
局部诺特形式预概形&I，10.4.2\\
诺特形式预概形&I，10.4.2\\
$A$ 上的形式预概形；$A$-形式预概形&I，10.4.7\\
$\frS$ 上的形式预概形；$\frS$-形式预概形&I，10.4.7\\
$\frS$ 上有限型形式预概形；有限型 $\frS$-形式预概形&I，10.13.3\\
在 $S$ 上分离的形式预概形&I，10.15.1\\
具有有限表示的模&0，1.0.5\\
具有有限表示的 $\OO_X$-模层&0，5.2.5\\
诺特归纳原理&0，2.2.2\\
$S$-预概形的积&I，3.2.1\\
$\frS$-形式预概形的积&I，10.7.1\\
集合的纤维积&I，3.4.2\\
不同预概形上层的张量积&I，9.1.2\\
代数的完备张量积&0，7.7.5\\
模的完备张量积&0，7.7.2\\
积的典范投影&I，3.2.1\\
子模层的典范延拓&I，9.4.1\\
态射到完备化的延拓&I，10.9.1\\
$\OO_X$-模层的第 $p$ 外幂&0，4.1.4\\
拟凝聚的 $\OO_X$-模层&0，5.1.3\\
拟凝聚的 $\OO_X$-代数层&0，5.1.3\\
理想的根&0，1.1.1\\
环的根&0，1.1.2\\
局部自由 $\OO_X$-模层的秩&0，5.4.1\\
无挠 $\OO_X$-模层的秩&I，7.4.2\\
有限 $K$-概形的秩&I，6.4.5\\
有限 $K$-概形的可分秩&I，6.4.8\\
赋环空间在开集上的限制&0，4.1.2\\
态射在开集上的限制&0，4.1.2\\
预概形态射在子预概形上的限制&I，4.1.7\\
预概形在开集上的限制&I，2.1.8\\
有理映射在开集上的限制&I，7.1.2\\
概形&I，5.4.1\\
仿射概形&I，1.7.1\\
局部概形；预概形在一点的局部概形&I，2.4.1\\
代数 $K$-概形&I，6.4.1\\
$K$ 上的有限概形；有限 $K$-概形&I，6.4.5\\
$S$-概形&I，5.4.1\\
仿射形式概形&I，10.1.2\\
$\frS$-形式概形&I，10.15.1\\
层的截面&0，3.1.6\\
$S$-预概形的 $S$-截面&I，2.5.5；5.3.11\\
$S$-预概形的有理 $S$-截面&I，7.1.2\\
$\OO_X$ 的单位截面&0，4.1.1\\
限制形式幂级数&0，7.5.1\\
预概形的和&I，3.6.1\\
由子模层生成的 $\OO_X$-子代数层&0，4.1.3\\
子预概形&I，4.1.3\\
与浸入相伴的子预概形&I，4.2.1\\
闭子预概形&I，4.1.3\\
由理想层定义的闭子预概形&I，4.1.2\\
闭形式子预概形&I，10.14.2\\
一点的特化&0，2.1.2\\
环的谱&I，1.1.1\\
可容许环的形式谱&I，10.1.2\\
模的支集&0，1.7.1\\
群层的支集&0，3.1.6\\
定义理想的基本系&I，10.3.6；10.5.1\\
预概形的进制归纳系&I，10.12.2\\
$J$-进制拓扑\textsuperscript{1}；$J$-预进制拓扑&0，7.1.9；7.2.3\\
谱拓扑&I，1.1.2\\
平凡的可逆 $\OO_X$-模层&I，2.4.8\\
有限型 $\OO_X$-模层&0，5.2.1\\
截面在一点的值&0，5.5.1\\
\end{longtable}
\noindent\textsuperscript{1}\textbf{原注：}注意，我们的术语与通常的定义略有不同：只有在所涉及的环或模\emph{分离且完备}时，我们才使用“进制”一词。
\endgroup


% END terminology_index.tex

\par

% BEGIN original_contents.tex
\chapter*{第一卷原著目录}
\addcontentsline{toc}{chapter}{第一卷原著目录}
\markboth{第一卷原著目录}{第一卷原著目录}
\src{225—228}
\noindent\textbf{译注。}本目录保留原书印刷页码，供与法文底本互查；本译稿的实际页码请见卷首目录。第一卷原著目录中的标题依原文译出，个别标题与正文略有差别的情形不另行统一。
\begingroup\small
\renewcommand{\arraystretch}{1.1}
\begin{longtable}{@{}p{0.84\linewidth}r@{}}
\toprule 标题&原书页码\\\midrule\endfirsthead
\toprule 标题（续）&原书页码\\\midrule\endhead
\bottomrule\endfoot
\textbf{导言}&5\\
\textbf{第 0 章\quad 预备知识}&11\\
\textbf{§1\quad 分式环}&11\\
1.0\quad 环与代数&11\\
1.1\quad 理想的根；环的幂零根与根&12\\
1.2\quad 分式模与分式环&13\\
1.3\quad 函子性质&14\\
1.4\quad 乘法子集的变更&15\\
1.5\quad 底环的变更&17\\
1.6\quad 将模 $M_f$ 等同于一个归纳极限&19\\
1.7\quad 模的支集&20\\
\textbf{§2\quad 不可约空间；诺特空间}&21\\
2.1\quad 不可约空间&21\\
2.2\quad 诺特空间&23\\
\textbf{§3\quad 层的补充知识}&23\\
3.1\quad 取值于一个范畴的层&23\\
3.2\quad 开集基上的预层&25\\
3.3\quad 层的粘合&28\\
3.4\quad 预层的直像&29\\
3.5\quad 预层的逆像&30\\
3.6\quad 常值层与局部常值层&33\\
3.7\quad 群预层或环预层的逆像&34\\
3.8\quad 伪离散空间层&35\\
\textbf{§4\quad 赋环空间}&35\\
4.1\quad 赋环空间、$\A$-模层、$\A$-代数层&35\\
4.2\quad $\A$-模层的直像&39\\
4.3\quad $\A$-模层的逆像&40\\
4.4\quad 直像与逆像的关系&42\\
\textbf{§5\quad 拟凝聚层与凝聚层}&44\\
5.1\quad 拟凝聚层&44\\
5.2\quad 有限型层&45\\
5.3\quad 凝聚层&47\\
5.4\quad 局部自由层&48\\
5.5\quad 局部环赋环空间上的层&53\\
\textbf{§6\quad 平坦性}&54\\
6.1\quad 平坦模&55\\
6.2\quad 底环的变更&55\\
6.3\quad 平坦性的局部化&56\\
6.4\quad 忠实平坦模&57\\
6.5\quad 标量限制&58\\
6.6\quad 忠实平坦环&58\\
6.7\quad 赋环空间的平坦态射&59\\
\textbf{§7\quad 进制环}&60\\
7.1\quad 可容许环&60\\
7.2\quad 进制环与射影极限&62\\
7.3\quad 诺特预进制环&66\\
7.4\quad 局部环上的拟有限模&68\\
7.5\quad 限制形式幂级数环&69\\
7.6\quad 完备分式环&72\\
7.7\quad 完备张量积&75\\
7.8\quad 同态模上的拓扑&77\\[0.3em]
\textbf{第 I 章\quad 概形的语言}&79\\
\textbf{§1\quad 仿射概形}&80\\
1.1\quad 环的素谱&80\\
1.2\quad 环素谱的函子性质&83\\
1.3\quad 与模相伴的层&84\\
1.4\quad 素谱上的拟凝聚层&90\\
1.5\quad 素谱上的凝聚层&92\\
1.6\quad 素谱上拟凝聚层的函子性质&93\\
1.7\quad 仿射概形态射的刻画&96\\
\textbf{§2\quad 预概形与预概形态射}&97\\
2.1\quad 预概形的定义&97\\
2.2\quad 预概形态射&98\\
2.3\quad 预概形的粘合&101\\
2.4\quad 局部概形&101\\
2.5\quad 预概形上的预概形&103\\
\textbf{§3\quad 预概形的积}&104\\
3.1\quad 预概形的和&104\\
3.2\quad 预概形的积&104\\
3.3\quad 积的形式性质；基预概形的变更&108\\
3.4\quad 预概形的预概形值点；几何点&111\\
3.5\quad 满射与单射&114\\
3.6\quad 纤维&117\\
3.7\quad 应用：预概形模 $J$ 的约化&118\\
\textbf{§4\quad 子预概形与浸入态射}&119\\
4.1\quad 子预概形&119\\
4.2\quad 浸入态射&122\\
4.3\quad 浸入的积&124\\
4.4\quad 预概形的逆像&125\\
4.5\quad 局部浸入与局部同构&126\\
\textbf{§5\quad 约化预概形；分离条件}&127\\
5.1\quad 约化预概形&127\\
5.2\quad 具有给定底空间的子预概形的存在性&131\\
5.3\quad 对角线；态射的图&132\\
5.4\quad 分离态射与分离预概形&135\\
5.5\quad 分离性判别准则&136\\
\textbf{§6\quad 有限性条件}&140\\
6.1\quad 诺特预概形与局部诺特预概形&140\\
6.2\quad 阿廷预概形&143\\
6.3\quad 有限型态射&144\\
6.4\quad 代数预概形&147\\
6.5\quad 态射的局部确定&150\\
6.6\quad 拟紧态射与局部有限型态射&152\\
\textbf{§7\quad 有理映射}&155\\
7.1\quad 有理映射与有理函数&155\\
7.2\quad 有理映射的定义域&158\\
7.3\quad 有理函数层&161\\
7.4\quad 挠层与无挠层&163\\
\textbf{§8\quad Chevalley 概形}&164\\
8.1\quad 相关的局部环&164\\
8.2\quad 整概形的局部环&165\\
8.3\quad Chevalley 概形&168\\
\textbf{§9\quad 拟凝聚层的补充知识}&169\\
9.1\quad 拟凝聚层的张量积&169\\
9.2\quad 拟凝聚层的直像&171\\
9.3\quad 拟凝聚层截面的延拓&172\\
9.4\quad 拟凝聚层的延拓&174\\
9.5\quad 预概形的闭像；子预概形的闭包&176\\
9.6\quad 拟凝聚代数层；结构层的变更&179\\
\textbf{§10\quad 形式概形}&180\\
10.1\quad 仿射形式概形&180\\
10.2\quad 仿射形式概形的态射&182\\
10.3\quad 仿射形式概形的定义理想&183\\
10.4\quad 形式预概形与形式预概形态射&185\\
10.5\quad 形式预概形的定义理想&186\\
10.6\quad 作为预概形归纳极限的形式预概形&188\\
10.7\quad 形式预概形的积&193\\
10.8\quad 预概形沿闭子集的形式完备化&194\\
10.9\quad 态射到完备化的延拓&198\\
10.10\quad 对仿射形式概形上凝聚层的应用&201\\
10.11\quad 形式预概形上的凝聚层&204\\
10.12\quad 形式预概形的进制态射&206\\
10.13\quad 有限型态射&207\\
10.14\quad 形式预概形的闭子预概形&209\\
10.15\quad 分离形式预概形&212\\
\textbf{参考文献}&215\\
\textbf{记号索引}&217\\
\textbf{术语索引}&219\\
\end{longtable}
\endgroup


% END original_contents.tex

\par

% BEGIN colophon.tex
\chapter*{第一卷书末出版记载}
\addcontentsline{toc}{chapter}{第一卷书末出版记载}
\markboth{第一卷书末出版记载}{第一卷书末出版记载}
\src{228}
收稿日期：1959 年 10 月 17 日。

\medskip
1960 年。法国大学出版社印刷厂（Imprimerie des Presses Universitaires de France）。法国旺多姆（Vendôme）。

\medskip
版号：25937。法国印制。印刷编号：16186。

\bigskip
\noindent\textbf{译注。}以上为法文原著的出版记载，不是本中文译稿的出版或审定信息。

\bigskip
\noindent\textbf{原扉页与版权页的出版信息。}
原扉页列明：高等科学研究所；《代数几何基础》，A. Grothendieck 著，
与 J. Dieudonné 合作撰写；《概形的语言》；1960 年；《数学出版物》第 4 号。
地址为法国巴黎第十六区 Bugeaud 环岛 5 号。

\medskip
原版权页载：法定送存；第一版，1960 年第三季度；在所有国家保留一切权利；
\copyright\ 1960，高等科学研究所（Institut des Hautes Études Scientifiques）。


% END colophon.tex

\par

% BEGIN copyright.tex
\chapter{第一卷底本的档案与版权说明}
\markboth{档案与版权说明}{档案与版权说明}
以下中文为上传文件首页所附档案说明的译文；随后照录其法文版权声明。

\medskip
\noindent\textbf{文献信息}

ALEXANDER GROTHENDIECK\par
《代数几何基础》（与 Jean Dieudonné 合作撰写）：I．概形的语言。\par
\textit{Publications mathématiques de l'I.H.É.S.}，第 4 卷（1960），5–228 页。\par
档案条目标识：\texttt{PMIHES\_1960\_\_4\_\_5\_0}。

\medskip
\noindent\textbf{档案说明译文}

\copyright\ \textit{Publications mathématiques de l'I.H.É.S.}，1960，保留所有权利。
访问本刊档案，即表示同意其一般使用条件。
原档案声明：任何商业使用或系统性印刷均构成刑事侵权；本文件的任何复制或打印都必须包含这一版权声明。
本文由“古旧数学文献数字化”（Numérisation de documents anciens mathématiques）计划数字化。

\medskip
\noindent\textbf{法文版权声明照录}

\begingroup\small
\noindent
\copyright\ Publications mathématiques de l'I.H.É.S., 1960, tous droits réservés.

L'accès aux archives de la revue « Publications mathématiques de l'I.H.É.S. »
(\url{http://www.ihes.fr/IHES/Publications/Publications.html}), implique l'accord avec les conditions générales
 d'utilisation (\url{http://www.numdam.org/legal.php}).
Toute utilisation commerciale ou impression systématique est constitutive d'une infraction pénale.
Toute copie ou impression de ce fichier doit contenir la présente mention de copyright.

Article numérisé dans le cadre du programme\par
Numérisation de documents anciens mathématiques\par
\url{http://www.numdam.org/}
\endgroup

\medskip
\noindent\textbf{版本说明}

本书为根据用户所上传原文制作的全卷中文译稿，并非 IHÉS、NUMDAM 或原权利人发布的中文版本。
原版中的历史陈述、书目年代与引文编号均按源文件保留。


% END copyright.tex


% END backmatter/volume_I.tex


% BEGIN backmatter/volume_II.tex
\unnumberedchapter{第二卷参考文献（续）}
\src{205}
\noindent\textbf{合订编者说明。}第二卷此前最终版的书末源码未能恢复。以下书末资料依据本对话上传的同一法文底本重新译排；它们是原著资料的译文，不是对正文的新增数学审校。原著索引的编号与原著目录的页码均保持原义。
\begin{description}[leftmargin=3em,style=nextline]
\item[{[23]}] C. CHEVALLEY，\textit{Introduction to the theory of algebraic functions of one variable}（《单变量代数函数理论引论》），Math. Surveys，第6号，Amer. Math. Soc.，纽约，1951。
\item[{[24]}] P. JAFFARD，\textit{Les systèmes d'idéaux}（《理想系统》），巴黎，Dunod，1960。
\item[{[25]}] M. NAGATA，\textit{On the derived normal rings of Noetherian integral domains}（《论诺特整环的导出正规环》），\textit{Mem. Coll. Sci. Kyoto}，A辑，第XXIX卷（1955），293--303页。
\item[{[26]}] M. NAGATA，\textit{Existence theorems for non projective complete algebraic varieties}（《非射影完全代数簇的存在性定理》），\textit{Ill. J. Math.}，第II卷（1958），490--498页。
\end{description}

\unnumberedchapter{第二卷原著记号索引}
\src{207--208}
\begingroup\small
\begin{longtable}{@{}p{0.77\textwidth}@{\hspace{1em}}p{0.18\textwidth}@{}}
\toprule 记号与说明 & 原著条目\\\midrule\endfirsthead
\toprule 记号与说明 & 原著条目\\\midrule\endhead
\bottomrule\endfoot
$\A(X),\A(\F),\A(\B)$（$X$ 为 $S$-预概形，$\F$ 为 $\OO_X$-模，$\B$ 为 $\OO_X$-代数）&1.1.1\\
$\A(h)$（$h$ 为 $S$-态射）&1.1.2\\
$\Spec(\B)$（$\B$ 为 $\OO_S$-代数）&1.3.1\\
$\wt\M$（$\M$ 为 $\B$-模，$\B$ 为 $\OO_S$-代数）&1.4.3\\
$\mathrm T(E),\Sym(E),\Sym_A(E)$（$E$ 为 $A$-模）&1.7.1\\
$\Sym(\E),\Sym_{\A}(\E)$（$\E$ 为 $\A$-模）&1.7.4\\
$\V(\E)$（$\E$ 为拟凝聚 $\OO_S$-模）&1.7.8\\
$S_n,S_+,M_n,S^{(d)},M^{(d,k)},M^{(d)}$（$S$ 为分次环，$M$ 为分次 $S$-模）&2.1.1\\
$M(n),S(n)$（$S$ 为分次环，$M$ 为分次 $S$-模）&2.1.1\\
$\Hom_S(M,N)$（$S$ 为分次环，$M,N$ 为分次 $S$-模）&2.1.2\\
$\nrad_+,\mathfrak r_+(\frj)$（$\frj$ 为 $S_+$ 的理想）&2.1.10\\
$S_{(f)},M_{(f)}$（$f$ 为 $S_+$ 的齐次元素）&2.2.1\\
$S_{(T)},M_{(T)},S_{(\pp)},M_{(\pp)}$（$T$ 为 $S_+$ 中由齐次元素组成的乘法部分，$\pp$ 为 $S_+$ 的分次素理想）&2.2.7\\
$\Proj(S)$（$S$ 为分次环）&2.3.1\\
$V_+(E)$（$E$ 为 $S_+$ 的子集）&2.3.2\\
$D_+(f)$（$f$ 为 $S$ 的元素）&2.3.3\\
$\mathfrak j_+(Y)$（$Y$ 为 $\Proj(S)$ 的子集）&2.3.10\\
$\wt M$（$M$ 为分次 $S$-模）&2.5.2\\
$\wt u$（$u$ 为零次同态）&2.5.4\\
$\OO_X(n),\F(n)$（$X=\Proj(S)$，$\F$ 为 $\OO_X$-模）&2.5.10\\
$\lambda$（典范同态）&2.5.11\\
$\mu$（典范同态）&2.5.12\\
$X^{(d)}$（$X=\Proj(S)$，$d$ 为正整数）&2.5.16\\
$\Gamma_*(\F)$（$\F$ 为 $\OO_X$-模，$X=\Proj(S)$）&2.6.1\\
$\alpha_n,\alpha,\alpha_M$（典范同态）&2.6.2\\
$\beta,\beta_{\F}$（典范同态）&2.6.4\\
$G(\varphi),{}^a\varphi$（后一记号为滥用记号；$\varphi$ 为分次环同态）&2.8.1\\
$\Proj(\varphi)$（$\varphi$ 为分次环同态）&2.8.2\\
$\SSS^{(d)},\M^{(d,k)},\M^{(d)},\M(n)$（$\SSS$ 为分次 $\OO_Y$-代数，$\M$ 为分次 $\SSS$-模）&3.1.1\\
$\Proj(\SSS)$（$\SSS$ 为拟凝聚分次 $\OO_Y$-代数）&3.1.3\\
$X_f$（$X=\Proj(\SSS)$，$f$ 为 $\SSS_d$ 在 $Y$ 上的截面）&3.1.4\\
$\wt\M$（$\M$ 为分次 $\SSS$-模）&3.2.2\\
$\OO_X(n),\F(n)$（$X=\Proj(\SSS)$，$\F$ 为 $\OO_X$-模）&3.2.5\\
$\lambda,\mu$（典范同态）&3.2.6\\
$\Gamma_*(\F)$（$\F$ 为 $\OO_X$-模，$X=\Proj(\SSS)$）&3.3.1\\
$\alpha_n,\alpha,\alpha_{\M}$（典范同态）&3.3.2\\
$\beta,\beta_{\F}$（典范同态）&3.3.4\\
$G(\varphi),\Proj(\varphi)$（$\varphi$ 为分次 $\OO_Y$-代数同态）&3.5.1\\
$\Proj(u)$（$u$ 为从分次 $\OO_Y$-代数到分次 $\OO_{Y'}$-代数的 $g$-态射）&3.5.6\\
$r_{\LL,\psi}$（$\LL$ 为可逆 $\OO_X$-模）&3.7.1\\
$\Pb(\E),\Pb(E),\Pb_Y^n,\Pb_A^n$&4.1.1\\
$\Pb(u)$（$u$ 为 $\OO_Y$-模的满同态）&4.1.2\\
$\operatorname{Hyp}_Y(X,\E)$（$\E$ 为拟凝聚 $\OO_Y$-模）&4.2.5\\
$\varsigma$（Segre 态射）&4.3.1\\
$\F(n)$（$\F$ 为 $\OO_X$-模，$X$ 配备一个可逆 $\OO_X$-模）&4.5.1\\
$\varepsilon$（恒等自同构）&4.5.1\\
$P(u,T),\sigma_i(u)$（$u$ 为有限基自由模的自同态）&6.4.1\\
$P(u,T),\sigma_i(u),\det u$（$u$ 为整环或约化诺特环上有限型模的自同态）&6.4.2；6.4.7\\
$\sigma_i(u),\det u$（$u$ 为局部自由 $\A$-模的自同态）&6.4.8\\
$\sigma_i,\det$（层的同态）&6.4.8--6.4.10\\
$\sigma_i(u),\det u$（$u$ 为局部整预概形或局部诺特约化预概形上有限型 $\OO_X$-模的自同态）&6.4.9；6.4.10\\
$N_{\B/\A}(f)$（$f$ 为 $\A$-代数 $\B$ 的截面）&6.5.1\\
$N_{\B/\A}(\LL')$（$\LL'$ 为可逆 $\B$-模）&6.5.2\\
$N_{\B/\A}(h')$（$h'$ 为可逆 $\B$-模的同态）&6.5.3\\
$N_{X'/X}(\LL')$（$X'$ 在 $X$ 上有限，$\LL'$ 为可逆 $\OO_{X'}$-模）&6.5.5\\
$N_{X'/X}(h')$（$h'$ 为可逆 $\OO_{X'}$-模的同态）&6.5.5\\
$\SSS_X$（$X=\Proj(\SSS)$，$\SSS$ 为 $\OO_Y$ 的理想层之直和）&8.1.5\\
$S^{\geq},S^{\leq},S_f^{\geq},S_f^{\leq},M^{\geq},M^{\leq},M_f^{\geq},M_f^{\leq}$（$S$ 为分次环，$M$ 为分次 $S$-模，$f$ 为 $S_+$ 的齐次元素）&8.2.1\\
$\widehat S$（$S$ 为分次环）&8.2.2\\
$\widehat M$（$M$ 为分次 $S$-模）&8.2.4\\
$S_{[n]},S^{\natural},f^{\natural}$（$S$ 为分次环，$f$ 为其中的齐次元素）&8.2.6\\
$M_{[n]},M^{\natural}$（$M$ 为分次 $S$-模）&8.2.8\\
$\widehat\SSS$（$\SSS$ 为分次 $\OO_Y$-代数）&3.8.1\textsuperscript{*}\\
$\mathbf G_m$（群概形）&8.3.9\\
$\SSS_X$（$X=\Proj(\SSS)$，$\SSS$ 为分次 $\OO_Y$-代数）&8.6.1\\
$\SSS_X^{\geq},C_X,\widehat C_X,E_X,\widehat E_X$（$X=\Proj(\SSS)$，$\SSS_0=\OO_Y$）&8.6.1\\
$\SSS_{[n]},\SSS^{\natural}$（$\SSS$ 为分次 $\OO_Y$-代数）&8.7.2\\
$\sProj_0(\M),\sProj(\M),\M_X,\widehat\M,\M^\square$（$\M$ 为分次 $\SSS$-模）&8.12.1\\
$\M_X^{\geq}$（$X=\Proj(\SSS)$，$\M$ 为分次 $\SSS$-模）&8.12.5\\
$\M_{[n]},\M^{\natural}$（$\M$ 为分次 $\SSS$-模）&8.12.9\\
$(\wt\N)^-$（$\N$ 为分次 $\SSS$-模的 $\SSS$-子模）&8.13.1\\
$\overline\N$（$\N$ 为分次 $\SSS$-模的 $\SSS$-子模）&8.13.2\\
$\boldsymbol\Gamma_*(\F)$（$X=\Proj(\SSS)$，$\F$ 为分次 $\SSS_X$-模）&8.14.5\\
$\SSS_{X'},\boldsymbol\Gamma'_*(\F'),\sProj'(\M),\alpha',\beta'$（$X'$ 为 $X=\Proj(\SSS)$ 的开集，$\F'$ 为分次 $\SSS_{X'}$-模，$\M$ 为分次 $\SSS$-模）&8.14.6\\
\end{longtable}
\endgroup

\noindent\textbf{译注。}\textsuperscript{*}原著记号索引此处印为3.8.1；相应记号见正文8.3.1。索引原编号予以保留。

\unnumberedchapter{第二卷原著术语索引}
\src{209--211}
\begingroup\small
\begin{longtable}{@{}p{0.77\textwidth}@{\hspace{1em}}p{0.18\textwidth}@{}}
\toprule 术语 & 原著条目\\\midrule\endfirsthead
\toprule 术语 & 原著条目\\\midrule\endhead\bottomrule\endfoot
仿射态射 &1.6.1\\
一个预概形之上的仿射预概形 &1.2.1\\
分次环上的分次代数 &2.1.2\\
$A$-模的对称代数 &1.7.1\\
$\A$-模的对称 $\A$-代数 &1.7.4\\
本质约化的分次 $\OO_Y$-代数 &3.1.12\\
整分次 $\OO_Y$-代数 &3.1.12\\
丰沛可逆 $\OO_X$-模 &4.5.3\\
相对于 $f$ 丰沛、$f$-丰沛、$Y$-丰沛的可逆 $\OO_X$-模 &4.6.1\\
本质整分次环 &2.1.11\\
本质约化分次环 &2.1.10\\
与 $\B$-模相伴的层 &1.4.3\\
与分次 $S$-模相伴的层 &2.5.3\\
与分次 $\SSS$-模相伴的层 &3.2.2\\
与 $\OO_S$-代数相伴的 $S$-概形 &1.3.1\\
（TF）条件、（TN）条件 &2.7.2\\
（TF）条件、（TN）条件 &3.4.2\\
由分次 $\OO_Y$-代数定义的仿射锥、射影锥 &8.3.1\\
由分次 $\OO_Y$-代数定义的去顶点仿射锥、去顶点射影锥 &8.3.4\\
预概形 $\Proj(\SSS)$ 的仿射投影锥、射影投影锥 &8.3.1\\
域上的代数曲线 &7.4.2\\
完全代数曲线 &7.4.7\\
Serre 判别准则 &5.2.1\\
整环上有限型模自同态的行列式 &6.4.2\\
约化诺特环上有限型模自同态的行列式 &6.4.7\\
局部自由 $\A$-模自同态的行列式 &6.4.8\\
局部整预概形 $X$ 上有限型 $\OO_X$-模自同态的行列式 &6.4.9\\
局部诺特约化预概形 $X$ 上有限型 $\OO_X$-模自同态的行列式 &6.4.10\\
爆破预概形 &8.1.3\\
整态射 &6.1.1\\
一个预概形之上的整预概形 &6.1.1\\
整拟凝聚 $\OO_S$-代数 &6.1.2\\
$\A$-代数中在 $\A$ 上整的截面 &6.3.1\\
射影丛的基本层 &4.1.1\\
环层 $\A$ 在一个 $\A$-代数中的整闭包 &6.3.2\\
预概形 $X$ 关于一个 $\OO_X$-代数的整闭包 &6.3.4\\
仿射锥的射影闭包 &8.3.1\\
射影丛在 $k(y)$ 的扩张 $K$ 上有理的几何纤维 &4.2.6\\
由 $\OO_Y$-模定义的射影丛 &4.1.1\\
由 $\OO_S$-模定义的向量丛 &1.7.8\\
有限态射 &6.1.1\\
一个预概形之上的有限预概形 &6.1.1\\
有限拟凝聚 $\OO_S$-代数 &6.1.2\\
局部自由模自同态的初等对称函数 &6.4.1\\
整环上有限型模自同态的初等对称函数 &6.4.2\\
约化诺特环上有限型模自同态的初等对称函数 &6.4.7\\
常秩局部自由 $\A$-模自同态的初等对称函数 &6.4.8\\
局部整预概形上有限型 $\OO_X$-模自同态的初等对称函数 &6.4.9\\
局部诺特约化预概形上有限型 $\OO_X$-模自同态的初等对称函数 &6.4.10\\
一个 $\A$-代数的截面的初等对称函数 &6.5.1\\
分次 $\SSS$-模的 $\SSS$-子模的齐次化 &8.13.2\\
分次环同态 &2.1.2\\
分次模之间的 $k$ 次同态 &2.1.2\\
$S_+$ 的理想、$S_+$ 的分次素理想（$S$ 为分次环）&2.1.10\\
Chow 引理 &5.6.1\\
射影锥的无穷远处 &8.3.3\\
典范态射 $G(\LL)\to\Proj(S)$，其中 $S=\bigoplus_{n\geq0}\Gamma(X,\LL^{\otimes n})$ &4.5.1\\
典范态射 $X\to\Spec(\Gamma(X,\OO_X))$ &5.1.1\\
Segre 态射 &4.3.1\\
$S_+$ 的幂零根（$S$ 为分次环）&2.1.10\\
$\SSS_+$ 的幂零根（$\SSS$ 为拟凝聚分次 $\OO_Y$-代数）&3.1.12\\
一个 $\A$-代数的截面的范数 &6.5.1\\
可逆 $\B$-模的范数 &6.5.2.3\\
可逆 $\B$-模的截面的范数 &6.5.3\\
可逆 $\OO_{X'}$-模的范数 &6.5.5\\
周期分次代数 &8.14.12\\
整环上有限型模自同态的特征多项式 &6.4.2\\
约化诺特环上有限型模自同态的特征多项式 &6.4.7\\
仿射锥（射影锥）的顶点预概形 &8.3.3\\
有限表示的分次模 &2.1.1\\
有限表示的分次 $\SSS$-模 &3.1.1\\
两个分次模的分次张量积 &2.1.2\\
射影态射 &5.5.2\\
一个预概形之上的射影预概形 &5.5.2\\
固有态射 &5.4.1\\
固有子集 &5.4.10\\
一个预概形之上的固有预概形 &5.4.1\\
拟仿射态射 &5.1.1\\
一个预概形之上的拟仿射预概形 &5.1.1\\
拟仿射概形 &5.1.1\\
拟有限态射 &6.2.3\\
一个预概形之上的拟有限预概形 &6.2.3\\
拟射影态射 &5.3.1\\
一个预概形之上的拟射影预概形 &5.3.1\\
$S_+$ 中理想的根（$S$ 为分次环）&2.1.10\\
去顶点射影锥到其无穷远处的典范收缩 &8.3.5\\
$\OO_X(-1)$ 的典范截面（$X$ 为爆破预概形）&8.1.9\\
仿射锥的零截面 &8.3.3\\
向量丛的零截面 &1.7.9\\
仿射锥（射影锥）的顶点截面 &8.3.3\\
一个 $\OO_S$-代数的谱 &1.3.1\\
分次环的齐次素谱 &2.3.1\\
拟凝聚分次 $\OO_S$-代数的齐次谱 &3.1.3\\
分次模同态的（TN）-单射、（TN）-满射、（TN）-双射 &2.7.2\\
分次模的（TN）-同构 &2.7.2\\
分次环同态的（TN）-单射、（TN）-满射、（TN）-双射 &2.9.1\\
分次环的（TN）-同构 &2.9.1\\
分次 $\SSS$-模同态的（TN）-单射、（TN）-满射、（TN）-双射 &3.4.2\\
分次 $\SSS$-模的（TN）-同构 &3.4.2\\
分次 $\OO_Y$-代数同态的（TN）-单射、（TN）-满射、（TN）-双射 &3.6.1\\
分次 $\OO_Y$-代数的（TN）-同构 &3.6.1\\
$\Proj(S)$ 上的谱拓扑 &2.3.3\\
相对于 $q$ 甚丰沛、相对于 $Y$ 甚丰沛的可逆 $\OO_X$-模 &4.4.2\\
有限型分次 $\SSS$-模 &3.1.1\\
普遍闭态射 &5.4.9\\
\end{longtable}
\endgroup

\unnumberedchapter{第二卷原著目录}
\src{213--215}
\noindent 下表保留法文原著的印刷页码，不是本合订本的页码。
\begingroup\small
\begin{longtable}{@{}p{0.85\textwidth}@{\hspace{1em}}r@{}}
\toprule 内容 & 原书页码\\\midrule\endfirsthead
\toprule 内容 & 原书页码\\\midrule\endhead\bottomrule\endfoot
\textbf{第II章　若干类态射的初等整体研究} &5\\
\textbf{§1　仿射态射}&5\\
1.1　$S$-预概形与 $\OO_S$-代数&5\\
1.2　一个预概形之上的仿射预概形&6\\
1.3　与一个 $\OO_S$-代数相伴的、在 $S$ 之上仿射的预概形&8\\
1.4　在 $S$ 之上仿射的预概形上的拟凝聚层&9\\
1.5　底预概形的变换&12\\
1.6　仿射态射&14\\
1.7　与一个模层相伴的向量丛&14\\
\textbf{§2　齐次素谱}&19\\
2.1　分次环与分次模的一般性质&19\\
2.2　分次环的分式环&23\\
2.3　分次环的齐次素谱&25\\
2.4　$\Proj(S)$ 上的概形结构&28\\
2.5　与一个分次模相伴的层&30\\
2.6　与 $\Proj(S)$ 上的层相伴的分次 $S$-模&36\\
2.7　有限性条件&38\\
2.8　函子性质&41\\
2.9　概形 $\Proj(S)$ 的闭子预概形&48\\
\textbf{§3　分次代数层的齐次谱}&49\\
3.1　拟凝聚分次 $\OO_Y$-代数的齐次谱&49\\
3.2　与分次 $\SSS$-模相伴的 $\Proj(\SSS)$ 上的层&54\\
3.3　与 $\Proj(\SSS)$ 上的层相伴的分次 $\SSS$-模&56\\
3.4　有限性条件&59\\
3.5　函子性质&61\\
3.6　预概形 $\Proj(\SSS)$ 的闭子预概形&64\\
3.7　从预概形到齐次谱的态射&65\\
3.8　浸入齐次谱的判别准则&69\\
\textbf{§4　射影丛。丰沛层}&71\\
4.1　射影丛的定义&71\\
4.2　从预概形到射影丛的态射&72\\
4.3　Segre 态射&76\\
4.4　到射影丛的浸入。甚丰沛层&78\\
4.5　丰沛层&83\\
4.6　相对丰沛层&89\\
\textbf{§5　拟仿射态射；拟射影态射；固有态射；射影态射}&94\\
5.1　拟仿射态射&94\\
5.2　Serre 判别准则&97\\
5.3　拟射影态射&99\\
5.4　固有态射与普遍闭态射&100\\
5.5　射影态射&103\\
5.6　Chow 引理&106\\
\textbf{§6　整态射与有限态射}&110\\
6.1　在另一预概形上为整的预概形&110\\
6.2　拟有限态射&114\\
6.3　预概形的整闭包&116\\
6.4　模层自同态的行列式&120\\
6.5　可逆层的范数&125\\
6.6　应用：丰沛性的判别准则&130\\
6.7　Chevalley 定理&135\\
\textbf{§7　赋值判别准则}&138\\
7.1　赋值环的回顾&138\\
7.2　分离性的赋值判别准则&141\\
7.3　固有性的赋值判别准则&143\\
7.4　代数曲线与一维函数域&148\\
\textbf{§8　爆破概形；投影锥；射影闭包}&152\\
8.1　爆破预概形&152\\
8.2　分次环局部化的初步结果&157\\
8.3　投影锥&162\\
8.4　向量丛的射影闭包&168\\
8.5　函子性质&169\\
8.6　去顶点锥的一个典范同构&171\\
8.7　投影锥的爆破&173\\
8.8　丰沛层与收缩&177\\
8.9　Grauert 丰沛性判别准则：陈述&182\\
8.10　Grauert 丰沛性判别准则：证明&184\\
8.11　收缩的唯一性&189\\
8.12　投影锥上的拟凝聚层&191\\
8.13　子层与闭子概形的射影闭包&195\\
8.14　与分次 $\SSS$-模相伴的层的补充&197\\
参考文献（续）&205\\
记号索引&207\\
术语索引&209\\
勘误与增补（清单一）&217\\
\end{longtable}\endgroup
\begin{flushright}1960年7月22日收到。\end{flushright}

\unnumberedchapter{勘误与增补（清单一）}
\src{217--222}
\noindent\textbf{合订编者说明。}本清单是原作者随第二卷发表的、针对第0章和第I章的勘误与增补；其中的页码均指第一卷原版。保留其独立位置，不将这些改动无说明地并入第一卷正文。以下条目号及公式号按清单原文保留。

\section*{第0章}
\addcontentsline{toc}{section}{勘误与增补：第0章}
\begin{description}[style=nextline,leftmargin=2em]
\item[（0，2.1.3）] 第21页倒数第15行，将 \textit{inter-section} 改为 \textit{intersection}。
\item[（0，2.1.6）] 第22页第13行，将 $Z_i\cap\complement(\bigcup_{j\ne i}Z_j)$ 改为 $\bigcup_{j\ne i}Z_j$。\transnote{此处按勘误清单的可见字面转录。原段落中涉及开集与补集的定义，不能据此字面条目无说明地自动替换正文。}
\item[（0，3.1.6）] 第25页倒数第17行，将 $\Gamma(U,\F)$ 改为 $\Gamma(V,\F)$。
\item[（0，4.1.1）] 第36页第12行，将 $\G\to\F$ 改为 $\B\to\A$。
\item[（0，4.2.1）] 第39页倒数第11行，将 \textit{faisceaux}（层，复数）改为 \textit{faisceau}（层，单数）。
\item[（0，5.1.2）] 在第45页第14行之后加入：若 $\G$ 是由其在 $Y$ 上的截面生成的 $\OO_Y$-模，则 $f^*(\G)$ 由其在 $X$ 上的截面生成，因为 $f^*$ 是右正合函子。
\item[（0，5.2.4）] 第46页倒数第18行，将 $f_U$ 改为 $f_U^*$。
\item[（0，5.3.4）] 在第47页倒数第16行之前加入：若 $\A\to\B\to\C\to\mathscr D\to\E$ 是 $\OO_X$-模正合列，且 $\A,\B,\mathscr D,\E$ 都凝聚，则 $\C$ 凝聚。
\item[（0，5.4.1）] 在第49页第2行之后加入：假设 $\OO_X$ 凝聚，且 $\F$ 为凝聚 $\OO_X$-模。若在点 $x\in X$，$\F_x$ 是秩为 $n$ 的自由 $\OO_x$-模，则存在 $x$ 的邻域 $U$，使 $\F|U$ 局部自由且秩为 $n$。事实上，此时 $\F_x\simeq\OO_x^n$，由（5.2.7）即得此命题。
\item[（0，6.6.2）] 将第59页第10、11行替换为：“事实上，这由（6.4.1，d））得到。”
\item[（0，6.7.1）] 在第59页倒数第10行之前加入：若 $f$ 为平坦态射，也称 $X$ 在 $Y$ 上平坦，或称 $X$ 是 $Y$-平坦的。
\item[（0，7.2.9）] 第65页倒数第11行，将 $M_n$ 改为 $M_{n-1}$。
\end{description}

\section*{第I章}
\addcontentsline{toc}{section}{勘误与增补：第I章}
\begin{description}[style=nextline,leftmargin=2em]
\item[（I，1.2.1）] 第83页第18行，两处 $K$ 均改为 $k$。
\item[（I，1.2.7）] 第84页倒数第16行，将 $A$（两处）改为 $A'$，将 $\nrad$ 改为 $\nrad'$；倒数第17行，将 $X$ 改为 $X'$。
\item[（I，1.7.3）及（I，2.2.4）] 根据 J. Tate 的一个观察，这些条目中的陈述可推广如下。
\end{description}
\subsection*{1.8　从局部赋环空间到仿射概形的态射}
\addcontentsline{toc}{subsection}{增补 I，1.8　从局部赋环空间到仿射概形的态射}
\erritem{1.8.1}{命题}设 $(S,\OO_S)$ 为仿射概形，$(X,\OO_X)$ 为局部赋环空间。则环同态
$\Gamma(S,\OO_S)\to\Gamma(X,\OO_X)$ 的集合，与满足下述条件的赋环空间态射
$(\psi,\theta):(X,\OO_X)\to(S,\OO_S)$ 的集合之间有典范双射：对每个 $x\in X$，
$\theta_x^\sharp:\OO_{\psi(x)}\to\OO_x$ 是局部同态。
\begin{proof}
首先，若 $(X,\OO_X)$、$(S,\OO_S)$ 为任意两个赋环空间，则从前者到后者的赋环空间态射 $(\psi,\theta)$ 典范地确定一个环同态
$\Gamma(\theta):\Gamma(S,\OO_S)\to\Gamma(X,\OO_X)$，因而得到第一个映射
\begin{equation}
\rho:\Hom((X,\OO_X),(S,\OO_S))\longrightarrow
\Hom(\Gamma(S,\OO_S),\Gamma(X,\OO_X)).\tag{1.8.1.1}
\end{equation}
反之，在命题的假设下，令 $A=\Gamma(S,\OO_S)$，考虑环同态 $\varphi:A\to\Gamma(X,\OO_X)$。对每个 $x\in X$，满足 $\varphi(f)(x)=0$ 的 $f\in A$ 所成集合显然是 $A$ 的素理想，因为 $\OO_x/\mm_x=k(x)$ 是域。因此它是 $S=\Spec(A)$ 的一个元素，仍记为 ${}^a\varphi(x)$。此外，对每个 $f\in A$，依定义（0，5.5.2），
$({}^a\varphi)^{-1}(D(f))=X_f$，这证明 ${}^a\varphi:X\to S$ 连续。

再定义 $\OO_S$-模同态 $\wt\varphi:\OO_S\to{}^a\varphi_*(\OO_X)$。对每个 $f\in A$，有 $\Gamma(D(f),\OO_S)=A_f$（1.3.6）；对每个 $s\in A$，将 $s/f\in A_f$ 对应于
\[
(\varphi(s)|X_f)(\varphi(f)|X_f)^{-1}\in
\Gamma(X_f,\OO_X)=\Gamma(D(f),{}^a\varphi_*(\OO_X)).
\]
通过从 $D(f)$ 换为 $D(fg)$，立即可验证这确实定义一个 $\OO_S$-模同态，从而得到赋环空间态射 $({}^a\varphi,\wt\varphi)$。仍沿用这些记号，为简便令 $y={}^a\varphi(x)$，立即有（0，3.7.1）
\[
\wt\varphi_x^\sharp(s_y/f_y)=(\varphi(s)_x)(\varphi(f)_x)^{-1}.
\]
按定义，$s_y\in\mm_y$ 等价于 $\varphi(s)_x\in\mm_x$，故 $\wt\varphi_x^\sharp:\OO_y\to\OO_x$ 是局部同态。由此定义第二个映射
$\sigma:\Hom(\Gamma(S,\OO_S),\Gamma(X,\OO_X))\to\Omega$，其中 $\Omega$ 为满足每个 $\theta_x^\sharp$ 都局部的态射 $(\psi,\theta):(X,\OO_X)\to(S,\OO_S)$ 所成的集合。

余下须证明 $\sigma$ 与 $\rho$ 在 $\Omega$ 上的限制互逆。由 $\wt\varphi$ 的定义立即有 $\Gamma(\wt\varphi)=\varphi$，故 $\rho\circ\sigma$ 为恒等映射。为证明 $\sigma\circ\rho$ 为恒等映射，取 $(\psi,\theta)\in\Omega$ 并令 $\varphi=\Gamma(\theta)$。由 $\theta_x^\sharp$ 的假设，通过取商得到域的单同态
$\theta_x^\natural:k(\psi(x))\to k(x)$，满足对每个 $f\in A=\Gamma(S,\OO_S)$，
$\theta_x^\natural(f(\psi(x)))=\varphi(f)(x)$。因此 $f(\psi(x))=0$ 等价于 $\varphi(f)(x)=0$，已证明 ${}^a\varphi=\psi$。另一方面，由定义，图
\[
\begin{tikzcd}
A\arrow[r,"\varphi"]\arrow[d]&\Gamma(X,\OO_X)\arrow[d]\\
A_{\psi(x)}\arrow[r,"\theta_x^\sharp"']&\OO_x
\end{tikzcd}
\]
交换；以 $\wt\varphi_x^\sharp$ 代替 $\theta_x^\sharp$ 所得的类似图也交换，故（0，1.2.4）$\wt\varphi_x^\sharp=\theta_x^\sharp$，从而 $\wt\varphi=\theta$。
\end{proof}

\erritem{1.8.2}{}当 $(X,\OO_X)$、$(Y,\OO_Y)$ 为局部赋环空间时，我们将考虑满足对每个 $x\in X$，$\theta_x^\sharp:\OO_{\psi(x)}\to\OO_x$ 为局部同态的态射 $(\psi,\theta):(X,\OO_X)\to(Y,\OO_Y)$。今后说到局部赋环空间的态射时，均指这种态射。用这一态射定义，局部赋环空间显然组成一个范畴；对此范畴的两个对象 $X,Y$，$\Hom(X,Y)$ 表示从 $X$ 到 $Y$ 的局部赋环空间态射集合，即（1.8.1）中记作 $\Omega$ 的集合。需要考虑从 $X$ 到 $Y$ 的全部赋环空间态射时，为免混淆，记为 $\Hom_{\mathrm{an}}(X,Y)$。于是（1.8.1.1）写作
\begin{equation}
\rho:\Hom_{\mathrm{an}}(X,Y)\longrightarrow
\Hom(\Gamma(Y,\OO_Y),\Gamma(X,\OO_X)),\tag{1.8.2.1}
\end{equation}
其限制
\begin{equation}
\rho':\Hom(X,Y)\longrightarrow
\Hom(\Gamma(Y,\OO_Y),\Gamma(X,\OO_X))\tag{1.8.2.2}
\end{equation}
是局部赋环空间范畴中关于 $X,Y$ 的函子性映射。

\erritem{1.8.3}{推论}设 $(Y,\OO_Y)$ 为局部赋环空间。$Y$ 是仿射概形，当且仅当对于每个局部赋环空间 $(X,\OO_X)$，映射（1.8.2.2）为双射。
\begin{proof}
命题（1.8.1）说明必要性。反之，假设此条件成立，令 $A=\Gamma(Y,\OO_Y)$。由假设和（1.8.1），从局部赋环空间范畴到集合范畴的函子 $X\mapsto\Hom(X,Y)$ 和 $X\mapsto\Hom(X,\Spec(A))$ 同构。已知这蕴含一个典范同构 $Y\simeq\Spec(A)$ 的存在（参见0，8）。\transnote{原版末句写作 $X\simeq\Spec(A)$；其所指的固定对象是 $Y$，故在此说明该字母误植。}
\end{proof}

\erritem{1.8.4}{}设 $S=\Spec(A)$ 为仿射概形。以 $(S',A')$ 表示底空间只有一个点、结构层 $A'$ 为环 $A$ 在 $S'$ 上定义的层（必为常值层）的赋环空间。令 $\pi:S\to S'$ 为唯一的映射。另一方面，对 $S$ 的每个开集 $U$，有典范映射
$\Gamma(S',A')=\Gamma(S,\OO_S)\to\Gamma(U,\OO_S)$；它定义环层的 $\pi$-态射 $\iota:A'\to\OO_S$。这样典范地定义赋环空间态射 $i=(\pi,\iota):(S,\OO_S)\to(S',A')$。对每个 $A$-模 $M$，以 $M'$ 表示由 $M$ 在 $S'$ 上定义的常值层，它显然是一个 $A'$-模。显然 $i_*(\wt M)=M'$（1.3.7）。

\erritem{1.8.5}{引理}采用（1.8.4）的记号，对每个 $A$-模 $M$，典范函子性 $\OO_S$-同态（0，4.4.3.3）
\begin{equation}i^*(i_*(\wt M))\longrightarrow\wt M\tag{1.8.5.1}\end{equation}
为同构。
\begin{proof}
（1.8.5.1）的两端都右正合（函子 $M\mapsto i_*(\wt M)$ 显然正合），并与直和交换。将 $M$ 看成同态 $A^{(I)}\to A^{(J)}$ 的余核，可归结到 $M=A$ 的情形；在此情形结论显然。
\end{proof}

\erritem{1.8.6}{推论}设 $(X,\OO_X)$ 为赋环空间，$u:X\to S$ 为赋环空间态射。对每个 $A$-模 $M$，采用（1.8.4）的记号，有 $\OO_X$-模的典范函子性同构
\begin{equation}u^*(\wt M)\xrightarrow{\sim}u^*(i^*(M')).\tag{1.8.6.1}\end{equation}

\erritem{1.8.7}{推论}在（1.8.6）的假设下，对每个 $A$-模 $M$ 和每个 $\OO_X$-模 $\F$，有关于 $M,\F$ 的典范函子性同构
\begin{equation}
\Hom_{\OO_S}(\wt M,u_*(\F))\xrightarrow{\sim}
\Hom_A(M,\Gamma(X,\F)).\tag{1.8.7.1}
\end{equation}
事实上，由（0，4.4.3）及引理（1.8.5），有双函子的典范同构
\[
\Hom_{\OO_S}(\wt M,u_*(\F))\simeq
\Hom_{A'}(M',i_*(u_*(\F))).
\]
右端显然就是 $\Hom_A(M,\Gamma(X,\F))$。注意，（1.8.7.1）将每个 $\OO_S$-同态 $h:\wt M\to u_*(\F)$，也就是每个 $u$-态射 $\wt M\to\F$，映到 $A$-同态
$\Gamma(h):M\to\Gamma(S,u_*(\F))=\Gamma(X,\F)$。

\erritem{1.8.8}{}采用（1.8.4）的记号，由（0，4.1.1），每个赋环空间态射 $(\psi,\theta):X\to S'$ 显然等价于给定环同态 $A\to\Gamma(X,\OO_X)$。所以也可将（1.8.1）解释为定义一个典范双射
$\Hom(X,S)\simeq\Hom(X,S')$，其中右端当然指赋环空间态射，因为一般 $A$ 并非局部环。更一般地，若 $X,Y$ 为局部赋环空间，$(Y',A')$ 为底空间只有一点且环层 $A'$ 是由 $\Gamma(Y,\OO_Y)$ 定义的常值层的赋环空间，则可将（1.8.2.1）解释为映射
\begin{equation}\rho:\Hom_{\mathrm{an}}(X,Y)\to\Hom(X,Y').\tag{1.8.8.1}\end{equation}
因而（1.8.3）的结果说明，局部赋环空间中的仿射概形可刻画为满足如下条件的对象：对每个局部赋环空间 $X$，$\rho$ 在 $\Hom(X,Y)$ 上的限制
\begin{equation}\rho':\Hom(X,Y)\to\Hom(X,Y')\tag{1.8.8.2}\end{equation}
为双射。在以后的一个章节中，我们将推广这一定义，使之能够将一个局部赋环空间与任意赋环空间 $Z$ 相联系，而不再只限于底空间只有一点的赋环空间；我们仍将它称为 $\Spec(Z)$。这将成为任意赋环空间之上预概形的“相对”理论的出发点，并推广第I章的结果。

\erritem{1.8.9}{}可以把局部赋环空间 $X$ 与 $\OO_X$-模 $\F$ 构成的二元组 $(X,\F)$ 看作一个范畴的对象。其中态射 $(X,\F)\to(Y,\G)$ 是一个二元组 $(u,h)$，由局部赋环空间态射 $u:X\to Y$ 和模的 $u$-态射 $h:\G\to\F$ 构成。固定 $(X,\F),(Y,\G)$ 时，这些态射所成集合记为 $\Hom((X,\F),(Y,\G))$。映射 $(u,h)\mapsto(\rho'(u),\Gamma(h))$ 是典范映射
\begin{equation}
\begin{aligned}
\Hom((X,\F),(Y,\G))\longrightarrow
\Hom\bigl(& (\Gamma(Y,\OO_Y),\Gamma(Y,\G)),\\
& (\Gamma(X,\OO_X),\Gamma(X,\F))\bigr),
\end{aligned}\tag{1.8.9.1}
\end{equation}
它关于 $(X,\F),(Y,\G)$ 是函子性的；右端为与所考虑的环和模相应的双同态集合（0，1.0.2）。

\erritem{1.8.10}{推论}设 $Y$ 为局部赋环空间，$\G$ 为 $\OO_Y$-模。$Y$ 是仿射概形且 $\G$ 是拟凝聚 $\OO_Y$-模，当且仅当对于每个由局部赋环空间 $X$ 和 $\OO_X$-模 $\F$ 构成的二元组 $(X,\F)$，典范映射（1.8.8.1）为双射。\transnote{此处保留勘误清单中的引文编号；按所述二元组及映射的定义，其应指（1.8.9.1）。}

将这一论证留给读者；它仿照（1.8.3）的论证，并使用（1.8.1）与（1.8.7）。

\erritem{1.8.11}{注}（1.7.3）、（1.7.4）、（2.2.4）的陈述及定义（1.6.1）都是（1.8.1）的特殊情形；同样，（2.2.5）由（1.8.7）得到。命题（1.8.7）也推出（1.6.3），因而还推出（1.6.4），作为其特殊情形。事实上，若 $X$ 仿射且 $\Gamma(X,\F)=N$，则函子
\[
M\longmapsto\Hom_{\OO_S}(\wt M,u_*(\wt N)),\qquad
M\longmapsto\Hom_{\OO_S}(\wt M,(N_{[\varphi]})^{\sim})
\]
由（1.8.7）及（1.3.8）同构，其中 $\varphi:A\to\Gamma(X,\OO_X)$ 对应于 $u$。最后，（1.6.5），从而（1.6.6），由（1.8.6）以及以下事实得到：对每个 $f\in A$，两个 $A_f$-模
$N'\otimes_{A'}A_f$ 和 $(N'\otimes_{A'}A)_f$ 典范同构，记号如（1.6.5）。

\medskip\noindent 在（I，3.2.8）之后插入：
\erritem{3.2.9}{}由（1.8.1），可将（3.2.2）精确化如下：$Z=\Spec(B\otimes_A C)$ 不仅是 $X=\Spec(B)$ 与 $Y=\Spec(C)$ 在 $S$-预概形范畴中的积，也是它们在 $S$ 上的局部赋环空间范畴中的积，其中 $S$-态射仿照（2.5.2）定义。（3.2.6）的证明因而实际上表明，对任意两个 $S$-预概形 $X,Y$，预概形 $X\times_S Y$ 不仅是它们在 $S$-预概形范畴中的积，也是它们在预概形 $S$ 上的局部赋环空间范畴中的积。

\begin{description}[style=nextline,leftmargin=2em]
\item[（I，4.4.5）] 第126页倒数第14行，将 $B$ 改为 $A$；倒数第13行，将“$A$-代数”改为“$B$-代数”。
\item[（I，5.3.5）] 第132页倒数第2行，将 $Y\to S$ 改为 $g:Y\to S$。
\item[（I，6.3.2.1）] 第144页倒数第6行，应读作 $D(g_i)\subset W$。
\item[（I，7.3.8）] 以如下文字替换原有的错误文字：
\end{description}
设 $X,Y$ 为两个整预概形，因此 $\R(X)$（分别地，$\R(Y)$）是拟凝聚 $\OO_X$-模（分别地，$\OO_Y$-模）（7.3.3）。设 $f:X\to Y$ 为占优态射；则有 $\OO_X$-模的典范同态
\begin{equation}\tau:f^*(\R(Y))\longrightarrow\R(X).\tag{7.3.8.1}\end{equation}
先设 $X=\Spec(A)$、$Y=\Spec(B)$ 仿射，$A,B$ 为整环，$f$ 对应于单同态 $B\to A$；它延拓为从 $B$ 的分式域 $L$ 到 $A$ 的分式域 $K$ 的单同态 $L\to K$。于是（7.3.8.1）对应于典范同态 $L\otimes_B A\to K$（1.6.5）。一般地，对每对满足 $f(U)\subset V$ 的非空仿射开集 $U\subset X,V\subset Y$，按上述方式定义同态 $\tau_{U,V}$。立即可见，若 $U'\subset U,V'\subset V,f(U')\subset V'$，则 $\tau_{U,V}$ 延拓 $\tau_{U',V'}$，由此得到所需断言。

若 $x,y$ 分别为 $X,Y$ 的泛点，则 $f(x)=y$，并且
\[
(f^*(\R(Y)))_x=\OO_y\otimes_{\OO_y}\OO_x=\OO_x
\]
（0，4.3.1），因此 $\tau_x$ 是同构。
\begin{description}[style=nextline,leftmargin=2em]
\item[（I，9.5.2）] 在第176页最后3行和第177页最前5行中，将所有 $X'$ 改为 $Y'$，所有 $X''$ 改为 $Y''$。
\item[（I，10.5.4）] 第187页倒数第11行，将（10.3.4）改为（10.3.5）。
\item[（I，10.6.3）] 第189页倒数第3行，将 $X$ 改为 $\frX$。
\item[（I，10.14.2）] 第210页第5行，将“凝聚 $\OO_{\frX}$-模”改为“$\OO_{\frX}$ 的凝聚理想”；第7行，将“凝聚 $\OO_{\frX}$-模”改为“$\OO_{\frX}$ 的凝聚理想”。
\end{description}

% END backmatter/volume_II.tex


% BEGIN backmatter/volume_III_1.tex
\clearpage\makeatletter\@mainmatterfalse\makeatother
% ---- BEGIN backmatter.tex ----
\chapter{第三卷第一部分记号索引}
\noindent\textbf{译者说明。} 本索引依底本次序收录；记号中的点号统一排作 $\bullet$，层仍用花体字母。条目后所列为原著章、节、条号，而不是本译稿页码。底本 PDF 第158页是记号索引，第159页是参考文献；其页脚总页码分别为506、505，原著目录所列本卷页码分别为162、161。本译稿保留上传文件的排列次序及原目录的页码。索引中的引文按原著保留，未将疑似错引静默改成别的条号。
\par\medskip
\begingroup
\small
\newcommand{\notationitem}[2]{\par\addvspace{0.45\baselineskip}\noindent #1\quad：#2。\par}
\notationitem{$\Ens$}{0，8.1.1}
\notationitem{$h_X(Y),\ h_X(u)$（$X,Y$ 为范畴 $C$ 的对象，$u$ 为 $C$ 的态射）}{0，8.1.1}
\notationitem{$h_w(u)$（$u,w$ 为 $C$ 的态射）}{0，8.1.2}
\notationitem{$Z_k(E_r^{pq}),\ B_k(E_r^{pq})$（$k\geq r$）}{0，11.1.1}
\notationitem{$Z_\infty(E_2^{pq}),\ B_\infty(E_2^{pq})$}{0，11.1.1}
\notationitem{$K^\bullet,\ K_\bullet$（复形）}{0，11.2.1}
\notationitem{$K^{\bullet\bullet},\ K_{\bullet\bullet}$（双复形）}{0，11.2.1}
\notationitem{$B^i(K^\bullet),\ Z^i(K^\bullet),\ H^i(K^\bullet),\ H^\bullet(K^\bullet)$}{0，11.2.1}
\notationitem{$B_i(K_\bullet),\ Z_i(K_\bullet),\ H_i(K_\bullet),\ H_\bullet(K_\bullet)$}{0，11.2.1}
\notationitem{$T(K^\bullet),\ T(K_\bullet)$（$T$ 为函子）}{0，11.2.1}
\notationitem{$E(K^\bullet)$（$K^\bullet$ 为滤过复形）}{0，11.2.2}
\notationitem{$K^{i,\bullet},\ K^{\bullet,j},\ Z_{\mathrm{II}}^p(K^{i,\bullet}),\ B_{\mathrm{II}}^p(K^{i,\bullet}),\ H_{\mathrm{II}}^p(K^{i,\bullet}),\ Z_{\mathrm I}^p(K^{\bullet,j}),\ B_{\mathrm I}^p(K^{\bullet,j}),\ H_{\mathrm I}^p(K^{\bullet,j})$}{0，11.3.1}
\notationitem{$H_{\mathrm{II}}^p(K^{\bullet\bullet}),\ Z_{\mathrm I}^q(H_{\mathrm{II}}^p(K^{\bullet\bullet})),\ B_{\mathrm I}^q(H_{\mathrm{II}}^p(K^{\bullet\bullet})),\ H_{\mathrm I}^q(H_{\mathrm{II}}^p(K^{\bullet\bullet})),\ H^n(K^{\bullet\bullet})$}{0，11.3.1}
\notationitem{$F_{\mathrm I}^p(K^{\bullet\bullet}),\ F_{\mathrm{II}}^p(K^{\bullet\bullet}),\ {}'E(K^{\bullet\bullet}),\ {}''E(K^{\bullet\bullet})$}{0，11.3.2}
\notationitem{$K_{i,\bullet},\allowbreak\ K_{\bullet,j},\allowbreak\ H_p^{\mathrm{II}}(K_{i,\bullet}),\allowbreak\ H_p^{\mathrm I}(K_{\bullet,j}),\allowbreak\ H_p^{\mathrm{II}}(K_{\bullet\bullet}),\allowbreak\ H_p^{\mathrm I}(K_{\bullet\bullet}),\allowbreak\ H_q^{\mathrm I}(H_p^{\mathrm{II}}(K_{\bullet\bullet})),\allowbreak\ H_q^{\mathrm{II}}(H_p^{\mathrm I}(K_{\bullet\bullet})),\allowbreak\ H_n(K_{\bullet\bullet})$}{0，11.3.5}
\notationitem{$\mathrm R^\bullet T(K^\bullet)$（$T$ 为协变加性函子）}{0，11.4.3}
\notationitem{$\mathrm R^\bullet T(K^\bullet,K'^{\bullet})$（$T$ 为协变加性双函子）}{0，11.4.6}
\notationitem{$\mathrm L_\bullet T(K_\bullet)$（$T$ 为协变加性函子）}{0，11.6.2}
\notationitem{$\mathrm L_\bullet T(K_\bullet,K'_{\bullet})$（$T$ 为协变加性双函子）}{0，11.6.6}
\notationitem{$\mathrm L_\bullet T(K_{\bullet\bullet})$（$T$ 为协变加性函子）}{0，11.7.2}
\notationitem{$|\sigma|,\ \Sigma(A),\ C_\bullet(A),\ L_\bullet(A)$（$A$ 为有限集，$\sigma$ 为 $A$ 的单形）}{0，11.8.1}
\notationitem{$C^\bullet(A,B;\F),\ L^\bullet(A,B;\F)$（$A,B$ 为有限集，$\F$ 为系数系）}{0，11.8.4}
\notationitem{$P^\bullet(A,B;\F)$（$A,B$ 为有限集，$\F$ 为系数系）}{0，11.8.6}
\notationitem{$C^\bullet(A,B;\F^\bullet),\ L^\bullet(A,B;\F^\bullet)$（$A,B$ 为有限集，$\F^\bullet$ 为系数系的复形）}{0，11.8.8}
\notationitem{$P^\bullet(A,B;\F^\bullet)$（$A,B$ 为有限集，$\F^\bullet$ 为系数系的复形）}{0，11.8.10}
\notationitem{$\chi(M^\bullet)$（$M^\bullet$ 为有限长度模的有限复形）}{0，11.10.1}
\notationitem{$C^\bullet(\mathfrak U,\F)$（$\F$ 为 $X$ 上的阿贝尔群层，$\mathfrak U$ 为 $X$ 的开覆盖）}{0，12.1.2}
\notationitem{$\mathrm R^pf_*(\F)$（$f:X\to Y$ 为赋环空间态射，$\F$ 为 $\OO_X$-模层）}{0，12.2.1}
\notationitem{$\mathscr H^p(f,\K^\bullet),\ \mathscr H_f^p(\K^\bullet),\ {}'\mathscr E(f,\K^\bullet),\ {}''\mathscr E(f,\K^\bullet)$（$f:X\to Y$ 为赋环空间态射，$\K^\bullet$ 为 $\OO_X$-模的复形）}{0，12.4.1}
\notationitem{$H^\bullet(X,\K^\bullet),\ H^\bullet(V,\K^\bullet)$（$X$ 为赋环空间，$\K^\bullet$ 为 $\OO_X$-模的复形，$V$ 为 $X$ 的开集）}{0，12.4.2}
\notationitem{$\gr^\bullet(\mathbf A)$（$\mathbf A$ 为满足（ML）的射影系统）}{0，13.4.2}
\notationitem{$E(A)$（$A$ 为滤过对象）}{0，13.6.4}
\notationitem{$K_\bullet(\mathbf f),\ i_{\mathbf f}$（$\mathbf f$ 为一个环的元素组成的有限族）}{III，1.1.1}
\notationitem{$K^\bullet(\mathbf f,M),\ K_\bullet(\mathbf f,M)$（$\mathbf f$ 为环 $A$ 的元素组成的有限族，$M$ 为 $A$-模）}{III，1.1.2}
\notationitem{$H^\bullet(\mathbf f,M),\ H_\bullet(\mathbf f,M)$（$\mathbf f$ 为环 $A$ 的元素组成的有限族，$M$ 为 $A$-模）}{III，1.1.3}
\notationitem{$C^\bullet((\mathbf f),M),\ H^\bullet((\mathbf f),M)$（$\mathbf f$ 为环 $A$ 的元素组成的有限族，$M$ 为 $A$-模）}{III，1.1.6}
\notationitem{$H^\bullet(U,\F(*)),\ H^\bullet(\mathfrak U,\F(*))$（$\F$ 为拟凝聚 $\OO_X$-模，$U$ 为 $X$ 的开集，$\mathfrak U$ 为 $X$ 的开覆盖）}{III，2.1.1}
\endgroup

\chapter{第三卷第一部分参考文献（续）}
\noindent\textbf{译者说明。} 本部分原著新增的参考文献为 [27]—[29]，编号承接此前各卷；以下保留原题名与出版资料，并附中文题名以便阅读。正文的（M）、（G）、（T）、（FAC）等缩写沿用原著，不能理解成本册另有同名章节。
\begin{enumerate}[label={[\arabic*]},start=27,leftmargin=3em,itemsep=1.2em]
\item P. Gabriel，\textit{Des catégories abéliennes}（《阿贝尔范畴》），博士论文，巴黎，1961年。
\item J. E. Roos，\textit{Sur les foncteurs dérivés de $\plim$. Applications}（《关于射影极限的导出函子及其应用》），\textit{C. R. Acad. Sci. Paris}，第252卷，1961年，第3702—3704页。
\item H. Cartan，\textit{Séminaire de l'École Normale Supérieure}（巴黎高等师范学校讨论班），第13学年（1960—1961），第11讲。
\end{enumerate}

\chapter{第三卷第一部分术语索引}
\noindent\textbf{译者说明。} 按原著法文词目次序排列，附法文词目用于对照；条号仍为原著条号。部分词目在原著中合并列出，本译稿亦保持合并。
\par\medskip
\begingroup\small
\newcommand{\termitem}[3]{\par\addvspace{0.5\baselineskip}\noindent\textbf{#1}（#2）\quad：#3。\par}
\termitem{谱序列的终项}{Aboutissement d'une suite spectrale}{0，11.1.1}
\termitem{可代数化的模层}{Algébrisable, $\OO_{\mathfrak X}$-Module}{III，5.2.1}
\termitem{可代数化的形式概形}{Algébrisable, schéma formel}{III，5.4.2}
\termitem{解析整环}{Analytiquement intègre, anneau}{III，4.3.6}
\termitem{拟紧映射}{Application quasi-compacte}{0，9.1.1}
\termitem{解消的增广}{Augmentation d'une résolution}{0，11.4.1}
\termitem{复形的 Euler--Poincaré 示性数}{Caractéristique d'Euler--Poincaré d'un complexe}{0，11.10.1}
\termitem{凝聚模层的 Euler--Poincaré 示性数；$X$ 在 $\Spec(A)$ 上射影、$A$ 为阿廷环}{Caractéristique d'Euler--Poincaré d'un $\OO_X$-Module cohérent}{III，2.5.1}
\termitem{双交错上链}{Cochaîne bi-alternée}{0，11.8.4}
\termitem{由双复形定义的复形}{Complexe défini par un bicomplexe}{0，11.3.1}
\termitem{外代数复形}{Complexe de l'algèbre extérieure}{III，1.1.1}
\termitem{（ML）条件}{Condition (ML)}{0，13.1.2}
\termitem{可构造子集、可构造集}{Constructible, partie ou ensemble}{0，9.1.2}
\termitem{可构造函数}{Constructible, fonction}{0，9.3.1}
\termitem{杯积}{Cup-produit}{0，12.1.2}
\termitem{范畴中代数结构的双同态}{Dihomomorphisme pour les structures algébriques sur les catégories}{0，8.2.1}
\termitem{阿贝尔范畴中的正合子集}{Exact, sous-ensemble dans une catégorie abélienne}{III，3.1.1}
\termitem{最终常值的射影系统}{Essentiellement constant, système projectif}{0，13.4.2}
\termitem{Stein 分解}{Factorisation de Stein}{III，4.3.3}
\termitem{余离散、余分离、离散、穷尽、有限、分离滤过}{Filtration co-discrète, co-séparée, discrète, exhaustive, finie, séparée}{0，11.1.3}
\termitem{范畴的终对象}{Final, objet d'une catégorie}{0，8.1.10}
\termitem{形式预概形的有限态射}{Fini, morphisme de préschémas formels}{III，4.8.2}
\termitem{在 $\mathfrak Y$ 上有限的形式预概形；$\mathfrak Y$-有限形式预概形}{Fini, préschéma formel au-dessus de $\mathfrak Y$}{III，4.8.2}
\termitem{典范协变函子 $C\to\catHom(C^\circ,\Ens)$}{Foncteur covariant canonique}{0，8.1.2}
\termitem{算术亏格}{Genre arithmétique}{III，2.5.1}
\termitem{几何连通的纤维}{Géométriquement connexe, fibre}{III，4.3.4}
\termitem{范畴中代数结构的同态}{Homomorphisme de structures algébriques sur les catégories}{0，8.2.1}
\termitem{函子关于复形 $K^\bullet$ 的超上同调}{Hypercohomologie d'un foncteur par rapport à un complexe}{0，11.4.3}
\termitem{双函子关于两个复形 $K^\bullet,K'^{\bullet}$ 的超上同调}{Hypercohomologie d'un bifoncteur par rapport à deux complexes}{0，11.4.6}
\termitem{函子关于复形 $K_\bullet$ 的超同调}{Hyperhomologie d'un foncteur par rapport à un complexe}{0，11.6.2}
\termitem{双函子关于两个复形 $K_\bullet,K'_{\bullet}$ 的超同调}{Hyperhomologie d'un bifoncteur par rapport à deux complexes}{0，11.6.6}
\termitem{函子关于双复形的超同调}{Hyperhomologie d'un foncteur par rapport à un bicomplexe}{0，11.7.2}
\termitem{局部可构造集}{Localement constructible, ensemble}{0，9.1.11}
\termitem{范畴对象上的外合成律、内合成律}{Loi de composition externe ou interne sur les objets d'une catégorie}{0，8.2.1}
\termitem{滤过 $S$-$C$-模；分次、双分次 $\gr^\bullet(S)$-$C$-模}{Module filtré; module gradué, bigradué}{0，13.6.6}
\termitem{谱序列态射}{Morphisme de suites spectrales}{0，11.1.2}
\termitem{纤维的连通分支的几何个数}{Nombre géométrique de composantes connexes d'une fibre}{III，4.3.4}
\termitem{$C$-群对象、$C$-群、$C$-环、$C$-模}{Objet en groupes, groupe, anneau, module dans $C$}{0，8.2.3}
\termitem{边界、上边界、上闭链、闭链对象}{Objet des bords, des cobords, des cocycles, des cycles}{0，11.2.1}
\termitem{普遍像对象}{Objet des images universelles}{0，13.1.1}
\termitem{满足（ML）的射影系统的相伴分次对象}{Objet gradué associé à un système projectif satisfaisant à (ML)}{0，13.4.2}
\termitem{有下界、有上界的分次对象}{Objet gradué limité inférieurement ou supérieurement}{0，11.2.1}
\termitem{形式预概形中在 $\mathfrak Y$ 上固有的子集}{Partie propre d'un préschéma formel}{III，3.4.1}
\termitem{满子范畴}{Pleine, sous-catégorie}{0，8.1.5}
\termitem{全忠实函子}{Pleinement fidèle, foncteur}{0，8.1.5}
\termitem{Hilbert 多项式}{Polynôme de Hilbert}{III，2.5.3}
\termitem{形式预概形的固有态射}{Propre, morphisme de préschémas formels}{III，3.4.1}
\termitem{可表函子}{Représentable, foncteur}{0，8.1.8}
\termitem{上同调、右、左、同调、内射、自由、平坦、投射解消}{Résolution cohomologique, droite, gauche, homologique, injective, libre, plate, projective}{0，11.4}
\termitem{右、内射 Cartan--Eilenberg 解消}{Résolution de Cartan--Eilenberg droite, injective}{0，11.4.2}
\termitem{左、投射 Cartan--Eilenberg 解消}{Résolution de Cartan--Eilenberg gauche, projective}{0，11.6.1}
\termitem{逆紧子集}{Rétrocompact, ensemble}{0，9.1.1}
\termitem{严格射影系统}{Strict, système projectif}{0，13.4.2}
\termitem{阿贝尔范畴中的谱序列}{Suite spectrale dans une catégorie abélienne}{0，11.1.1}
\termitem{弱收敛、正则、余正则、双正则谱序列}{Suite spectrale faiblement convergente, régulière, corégulière, birégulière}{0，11.1.3}
\termitem{退化谱序列}{Suite spectrale dégénérée}{0，11.1.6}
\termitem{函子相对于滤过对象的谱序列}{Suite spectrale d'un foncteur relativement à un objet filtré}{0，13.6.4}
\termitem{双复形的谱序列}{Suites spectrales d'un bicomplexe}{0，11.3.2及11.3.5}
\termitem{函子关于复形的超上同调谱序列}{Suites spectrales d'hypercohomologie d'un foncteur par rapport à un complexe}{0，11.4.3}
\termitem{函子关于复形的超同调谱序列}{Suites spectrales d'hyperhomologie d'un foncteur par rapport à un complexe}{0，11.6.2}
\termitem{函子关于双复形的超同调谱序列}{Suites spectrales d'hyperhomologie d'un foncteur par rapport à un bicomplexe}{0，11.7.2}
\termitem{系数系}{Système de coefficients}{0，11.8.4}
\termitem{单分支环、单分支点}{Unibranche, anneau ou point}{III，4.3.6}
\termitem{普遍开态射}{Universellement ouvert, morphisme}{III，4.3.9}
\endgroup\par

% ---- END backmatter.tex ----
% ---- BEGIN original_toc.tex ----

\chapter{第三卷第一部分原著目录}
\noindent 以下为原著书末目录的中文译文。右栏保留原著所印本卷页码，不是本译稿页码；本译稿页码见卷首的统一目录。原书目录中参考文献与记号索引的先后，与上传 PDF 的实际物理页序不同，详见本译稿“记号索引”前的说明。
\par\medskip
\begingroup\small
\setlength{\LTpre}{0pt}\setlength{\LTpost}{0pt}
\begin{longtable}{@{}p{0.86\textwidth}>{\raggedleft\arraybackslash}p{0.09\textwidth}@{}}
\toprule
内容 & 原书页码\\\midrule
\endfirsthead
\toprule 内容 & 原书页码\\\midrule
\endhead
\bottomrule\endfoot
\textbf{第0章　预备知识（续）} & 5\\
\textbf{\S\,8　可表函子} & 5\\
\hspace*{1em}8.1　可表函子 & 5\\
\hspace*{1em}8.2　范畴中的代数结构 & 9\\
\textbf{\S\,9　可构造集} & 12\\
\hspace*{1em}9.1　可构造集 & 12\\
\hspace*{1em}9.2　诺特空间中的可构造集 & 14\\
\hspace*{1em}9.3　可构造函数 & 16\\
\textbf{\S\,10　平坦模的补充知识} & 17\\
\hspace*{1em}10.1　平坦模与自由模之间的关系 & 17\\
\hspace*{1em}10.2　平坦性的局部判别准则 & 18\\
\hspace*{1em}10.3　局部环的平坦扩张的存在性 & 20\\
\textbf{\S\,11　同调代数的补充知识} & 23\\
\hspace*{1em}11.1　谱序列回顾 & 23\\
\hspace*{1em}11.2　滤过复形的谱序列 & 27\\
\hspace*{1em}11.3　双复形的谱序列 & 29\\
\hspace*{1em}11.4　函子相对于复形的超上同调 & 32\\
\hspace*{1em}11.5　超上同调中的归纳极限 & 35\\
\hspace*{1em}11.6　函子相对于复形的超同调 & 39\\
\hspace*{1em}11.7　函子相对于双复形的超同调 & 41\\
\hspace*{1em}11.8　单纯复形上同调的补充知识 & 43\\
\hspace*{1em}11.9　关于有限型复形的一个引理 & 46\\
\hspace*{1em}11.10　有限长度模复形的 Euler--Poincaré 示性数 & 48\\
\textbf{\S\,12　层上同调的补充知识} & 49\\
\hspace*{1em}12.1　赋环空间上模层的上同调 & 49\\
\hspace*{1em}12.2　高阶直像 & 57\\
\hspace*{1em}12.3　层的 \texorpdfstring{$\Ext$}{Ext} 函子的补充知识 & 60\\
\hspace*{1em}12.4　直像函子的超上同调 & 62\\
\textbf{\S\,13　同调代数中的射影极限} & 64\\
\hspace*{1em}13.1　Mittag--Leffler 条件 & 64\\
\hspace*{1em}13.2　阿贝尔群的 Mittag--Leffler 条件 & 65\\
\hspace*{1em}13.3　应用：层的射影极限的上同调 & 68\\
\hspace*{1em}13.4　Mittag--Leffler 条件与射影系统的相伴分次对象 & 69\\
\hspace*{1em}13.5　滤过复形谱序列的射影极限 & 71\\
\hspace*{1em}13.6　函子相对于有限滤过对象的谱序列 & 73\\
\hspace*{1em}13.7　变元的射影极限与导出函子 & 75\\
\textbf{第III章　凝聚层的上同调研究} & 81\\
\textbf{\S\,1　仿射概形的上同调} & 82\\
\hspace*{1em}1.1　外代数复形的回顾 & 82\\
\hspace*{1em}1.2　开覆盖的 Čech 上同调 & 85\\
\hspace*{1em}1.3　仿射概形的上同调 & 88\\
\hspace*{1em}1.4　应用于任意预概形的上同调 & 89\\
\textbf{\S\,2　射影态射的上同调研究} & 95\\
\hspace*{1em}2.1　某些上同调群的显式计算 & 95\\
\hspace*{1em}2.2　射影态射的基本定理 & 100\\
\hspace*{1em}2.3　应用于分次代数层和分次模层 & 102\\
\hspace*{1em}2.4　基本定理的一个推广 & 107\\
\hspace*{1em}2.5　Euler--Poincaré 特征与 Hilbert 多项式 & 109\\
\hspace*{1em}2.6　应用：丰沛性判据 & 111\\
\textbf{\S\,3　固有态射的有限性定理} & 115\\
\hspace*{1em}3.1　分解引理 & 115\\
\hspace*{1em}3.2　有限性定理：通常概形的情形 & 116\\
\hspace*{1em}3.3　有限性定理的推广（通常概形） & 118\\
\hspace*{1em}3.4　有限性定理：形式概形的情形 & 119\\
\textbf{\S\,4　固有态射的基本定理及其应用} & 122\\
\hspace*{1em}4.1　基本定理 & 122\\
\hspace*{1em}4.2　特殊情形与变式 & 129\\
\hspace*{1em}4.3　Zariski 连通性定理 & 130\\
\hspace*{1em}4.4　Zariski 的“主定理” & 135\\
\hspace*{1em}4.5　同态模的完备化 & 138\\
\hspace*{1em}4.6　形式态射与通常态射之间的关系 & 139\\
\hspace*{1em}4.7　一个丰沛性判据 & 145\\
\hspace*{1em}4.8　形式预概形的有限态射 & 146\\
\textbf{\S\,5　凝聚代数层的一个存在性定理} & 149\\
\hspace*{1em}5.1　定理的陈述 & 149\\
\hspace*{1em}5.2　存在性定理的证明：射影与拟射影情形 & 151\\
\hspace*{1em}5.3　存在性定理的证明：一般情形 & 154\\
\hspace*{1em}5.4　应用：通常概形态射与形式概形态射的比较；可代数化形式概形 & 156\\
\hspace*{1em}5.5　某些概形的一个分解 & 159\\
参考文献（续） & 161\\
记号索引 & 162\\
术语索引 & 163\\
\end{longtable}
\endgroup
\begin{flushright}1961年6月28日收到。\end{flushright}\par

% ---- END original_toc.tex ----

% END backmatter/volume_III_1.tex


% BEGIN backmatter/volume_III_2.tex
\head{第三卷第二部分：出版信息与书末索引}
\source{EGAIII.2.pdf，封面及书末；原印刷第81--83页（通卷页码213--215）}
\noindent\textbf{书名。}《代数几何基础》第III章：《凝聚层的上同调研究》，第二部分。A. Grothendieck 著，与 J. Dieudonné 合作撰写。法国高等科学研究所，1963年；\textit{Publications mathématiques de l'I.H.É.S.}，第17卷，第5--91页。初版法定送存：1963年第四季度。各国版权均保留。©1963，Institut des Hautes Études Scientifiques。

\noindent\textbf{底本的数字化版权说明（译文）。}©\textit{Publications mathématiques de l'I.H.É.S.}，1963，保留所有权利。访问该刊档案意味着接受其一般使用条件。任何商业使用或系统性印刷均构成违法行为。本文件的任何复制件或印刷件均须包含此版权声明。本文由“古数学文献数字化”计划数字化。

\NewNote{以下保留本分册原有索引的组织和指向；号码为原著条目号，不是本合订本页码。上传底本在正文之后依次为记号索引、术语索引、原目录、空白页和《勘误与增补（清单二）》，没有另列参考文献续表，故不补造一份续表。正文末的“待续”是原著历史标记。}

\section*{记号索引}
\addcontentsline{toc}{section}{第三卷第二部分记号索引}
\begin{longtable}{@{}p{.77\linewidth}p{.18\linewidth}@{}}
\toprule 记号 & 原著条目\\\midrule\endhead
$\operatorname{Tor}_n^A(P_\bullet,Q_\bullet)$ & 6.3.1\\
$\mathfrak{Tor}_n^S(\PP_\bullet,\mathscr Q_\bullet)$、$\mathfrak{Tor}_n(\PP_\bullet,\mathscr Q_\bullet)$ & 6.4.1；6.5.1\\
$\lTor_n^S(\F,\G)$ & 6.5.1\\
$\mathfrak{Tor}_q^S(\PP_\bullet^{(1)},\PP_\bullet^{(2)},\ldots,\PP_\bullet^{(m)})$ & 6.5.15\\
$H^{-n}(\mathfrak U^{(i)},\PP_\bullet^{(i)})$、$H^{-n}(X^{(i)},\PP_\bullet^{(i)})$ & 6.6.1\\
$\operatorname{Tor}_n^A(L_{\bullet\bullet}^{(1)},L_{\bullet\bullet}^{(2)})$；$\operatorname{Tor}_n^S(\mathfrak U^{(1)},\mathfrak U^{(2)};\PP_\bullet^{(1)},\PP_\bullet^{(2)})$ & 6.6.2\\
$\lTor_n^S(\mathfrak U^{(1)},\mathfrak U^{(2)};\F^{(1)},\F^{(2)})$；$\lTor_n^S(\mathfrak U^{(1)},\mathfrak U^{(2)};\PP_\bullet^{(1)},\PP_\bullet^{(2)})$ & 6.6.2\\
$\mathfrak{Tor}_n^S(f_1,f_2;\PP_\bullet^{(1)},\PP_\bullet^{(2)})$ & 6.6.8；6.7.1；6.7.3\\
$\mathfrak{Tor}_n^S((f_i)_{i\in I};(\PP_\bullet^{(i)})_{i\in I})$；$\mathfrak{Tor}_n^S(f_1,\ldots,f_m;\PP_\bullet^{(1)},\ldots,\PP_\bullet^{(m)})$ & 6.7.11\\
$\mathbf{Ab}$、$\mathbf{Ab}_A$ & 7.1.1\\
$T_{(B)}$、$T^{(B)}$、$T\otimes_A B$（$T$ 为从 $\mathbf{Ab}_A$ 到 $\mathbf{Ab}$ 的函子） & 7.1.3\\
$t_M$ & 7.2.2\\
$T_{\mathfrak p}$ & 7.1.4\\
$\TT_\bullet^{Y'}(\M')$、$\TT_\bullet^U(\M|U)$ & 7.7.2\\
$\operatorname{EP}(f,\PP_\bullet;y)$、$\operatorname{EP}(Y,\PP_\bullet;y)$、$\operatorname{EP}(\PP_\bullet;y)$ & 7.9.5\\
$\operatorname{PH}(f,\PP_\bullet;y)$、$\operatorname{PH}(\PP_\bullet;y)$ & 7.9.12\\
\bottomrule
\end{longtable}
\NewNote{原记号索引以不同字形区分超 Tor 和逐项 Tor。本译稿正文采用 $\bTor$ 与带 term 标识的记号来避免字体难辨，并在定义处说明。此表保留原索引的区分及其条目指向；Hilbert 多项式的命名在（7.9.11）末出现，但原索引指向（7.9.12），此处不静默改号。}

\section*{术语索引}
\addcontentsline{toc}{section}{第三卷第二部分术语索引}
\begin{longtable}{@{}p{.77\linewidth}p{.18\linewidth}@{}}
\toprule 术语 & 原著条目\\\midrule\endhead
模复形的 Euler--Poincaré 特征数 & 7.9.5\\
在 $y\in Y$ 处关于 $Y$ 在维数 $p$ 上上同调平坦；关于 $Y$ 在维数 $p$ 上上同调平坦；在 $y$ 处关于 $Y$ 上同调平坦；关于 $Y$ 上同调平坦（$Y$-平坦的 $\OO_X$-模复形） & 7.8.1\\
同伦等价的复形 & 6.1.4\\
上同调维数 $\leq n$ 的赋环空间 & 6.5.5\\
协变加性函子的标量扩张 & 7.1.3\\
上同调维数 $\leq n$ 的环层 & 6.5.5\\
Künneth 公式 & 6.7.8\\
在 $y\in Y$ 处关于 $Y$ 在维数 $p$ 上同调平坦；关于 $Y$ 在维数 $p$ 上同调平坦；在 $y$ 处关于 $Y$ 同调平坦；关于 $Y$ 同调平坦（$Y$-平坦的 $\OO_X$-模复形） & 7.8.1\\
同伦等价态射 & 6.1.4\\
两个 $A$-模复形的超 Tor & 6.3.1\\
两个模层复形的局部超 Tor & 6.4.1\\
两个模层复形的整体超 Tor & 6.6.2；6.6.8；6.7.3\\
模复形的 Hilbert 多项式 & 7.9.12\\
Künneth 谱序列 & 6.7.3\\
\bottomrule
\end{longtable}

\section*{原著目录}
\addcontentsline{toc}{section}{第三卷第二部分原著目录}
\NewNote{右栏页码是该分册原印刷页码，专供与底本核对；本合订本的实际页码见总目录。}
\begin{longtable}{@{}p{.86\linewidth}r@{}}
\toprule 内容 & 原页\\\midrule\endhead
第III章：凝聚层的上同调研究（续）&\\
§6．局部与整体 Tor 函子；Künneth 公式 &5\\
6.1．引言 &5\\
6.2．预概形上模复形的超上同调 &6\\
6.3．两个模复形的超 Tor &9\\
6.4．拟凝聚模层复形的局部超 Tor 函子：仿射概形情形 &14\\
6.5．拟凝聚模层复形的局部超 Tor 函子：一般情形 &16\\
6.6．拟凝聚模层复形的整体超 Tor 函子与 Künneth 谱序列：仿射底情形 &21\\
6.7．拟凝聚模层复形的整体超 Tor 函子与 Künneth 谱序列：一般情形 &25\\
6.8．整体超 Tor 的结合性谱序列 &32\\
6.9．整体超 Tor 的基变换谱序列 &34\\
6.10．某些上同调函子的局部结构 &39\\
§7．模的协变同调函子中的基变换研究 &43\\
7.1．$A$-模函子 &43\\
7.2．张量积函子的刻画 &44\\
7.3．模的同调函子的正合性判据 &48\\
7.4．函子 $H_\bullet(P_\bullet\otimes_A M)$ 的正合性判据 &53\\
7.5．诺特局部环的情形 &58\\
7.6．正合性的下降；半连续性定理与 Grauert 正合性判据 &60\\
7.7．对固有态射的应用：I．交换性质 &65\\
7.8．对固有态射的应用：II．上同调平坦性判据 &72\\
7.9．对固有态射的应用：III．Euler--Poincaré 特征数与 Hilbert 多项式的不变性 &76\\
\bottomrule
\end{longtable}
\begin{flushright}收稿日期：1962年12月15日。\end{flushright}

% END backmatter/volume_III_2.tex


% BEGIN backmatter/errata_III_2.tex
\head{原著勘误与增补（清单二）}
\source{第三卷第二部分，上传PDF第84--90页；原分册第85--91页，通卷页码217--223}
\NewNote{本章是原作者刊载的勘误全文译文，不是本次核查报告。所有页、行定位均指原版；本合订本页码已改变。原版 A 部分没有连续条号，这里仅按原有条目分段。B 部分保留 $\operatorname{Err}_{\mathrm{III}}$ 的1--30号。涉及难辨字形，均以底本图像而非抽取文本为准。}

\section*{A．排印错误}
\addcontentsline{toc}{section}{原勘误清单二：A．排印错误}
\newcommand{\typo}[2]{\par\smallskip\noindent\textbf{（#1）}\quad #2\par}
\typo{$0_{\mathrm I}$，3.2.1}{第27页第2行，将 $X$ 改为 $U$。}
\typo{$0_{\mathrm I}$，3.2.6}{第27页倒数第1行及第28页第1行，将 $\mathfrak H_\lambda$ 改为 $\mathscr H_\lambda$（共三处）。}
\typo{$0_{\mathrm I}$，3.4.5}{在“3.4.5)”之前加一个左括号。}
\typo{$0_{\mathrm I}$，4.2.1}{第39页第10行，将 $\psi_*(A)$ 改为 $\psi_*(\A)$。}
\typo{$0_{\mathrm I}$，5.1.3}{第45页第16行，将 $\F|V$ 改为 $\F|U$。}
\typo{$0_{\mathrm I}$，5.3.9}{第48页第12行，在“使得”之前加上“在 $x$ 的一个开邻域 $U$ 上”。}
\typo{$0_{\mathrm I}$，7.6.9}{第73页倒数第3行，将 $J_\lambda$ 改为 $\mathfrak J_\lambda$。}
\typo{I，1.1.15}{第83页第3行，将 $\mathfrak a$ 改为 $\mathfrak a\ne A$。}
\typo{I，1.4.1}{第90页倒数第2行，将 $\widetilde N$ 改为 $N$。}
\typo{I，2.3.1}{第101页第3行，将（0，4.1.6）改为（0，4.1.7）。}
\typo{I，2.3.2}{第101页第11行，将 $B$ 改为 $B_s$，将 $C$ 改为 $C_t$。}
\typo{I，2.5.5}{第104页第19行，将“$S$-态射”改为“$S$-预概形”。}
\typo{I，3.3.7}{第109页第17行，将 $Y_{(S)}$ 改为 $Y_{(S')}$。}
\typo{I，3.4.5}{第113页第9行，将 $\underline s$ 改为 $s$。}
\typo{I，3.4.8}{第114页第4行，将 $p(x)$ 改为 $f(x)$。}
\typo{I，3.5.1}{第115页第3行，将 $1_X\times g$ 改为 $1_{X'}\times g$。}
\typo{I，3.7.2}{第119页第8行，将第一个 $X$ 改为 $X'$。}
\typo{I，4.2.2}{第123页第8行，在交换图的左侧竖箭头处，将 $\alpha_{\psi(y)}$ 改为 $\rho_{\psi(y)}$。}
\typo{I，4.4.3}{第126页第16行，将法文误拼“isomorphee”改为“isomorphes”。}
\typo{I，4.5.5}{第127页倒数第17行，将（4.2.4）改为（4.2.5）。}
\typo{I，5.1.4}{第128页倒数第14行，将（2.1.7）改为（2.1.8）。}
\typo{I，5.3.13}{第134页第20行，将（4.2.4）改为（4.2.5）。}
\typo{I，6.4.2}{第147页倒数第1行，将“覆盖”改为“有限覆盖”。}
\typo{I，9.1.13}{第171页第15行，将 $p'^{-1}(\F)$ 改为 $p'^*(\F)$。}
\typo{I，9.5.11}{第179页倒数第7行，将 $Y$ 改为 $Y'$。}
\typo{I，9.6.5}{第180页倒数第17行，将（0，4.1.4）改为（0，4.1.3）。}
\typo{I，10.12.2}{第206页倒数第14行，将 $(X_n)_{(S_n)}$ 改为 $(X_n)_{(S_m)}$。}
\typo{第一卷术语索引}{第219页倒数第6行，将0，4.1.4改为0，4.1.5。第222页倒数第16、17、18行，将I，2.1.7改为I，2.1.8；第222页倒数第19行，将I，2.1.8改为I，2.1.7。}
\typo{II，1.5.2}{第12页第17行，在交换图左侧竖箭头旁加上 $f$。}
\typo{II，4.2.7}{第75页倒数第16行，将（III，2.1.14）改为（III，2.1.13）。}
\typo{II，5.2.1}{第97页倒数第16行，将 $X-U$ 改为 $U$。同页倒数第8行，将 $X^f$ 改为 $X_f$。}
\typo{II，5.3.6}{第100页第16行，将（4.6.18）改为（4.6.17）。}
\typo{II，6.2.7}{第116页倒数第9行，将“第V章”改为“第IV章”。}
\typo{II，6.6.5}{第132页倒数第15行，将“第V章”改为“第IV章”。}
\typo{II，7.4.12}{第151页第1行，将“第V章”改为“第III章”。}
\typo{II，8.1.4}{第154页第15行，将（III，2.3.8）改为（III，2.3.7）。}
\typo{II，8.10.5}{第187页倒数第17行，将 $g_{(j)x}$ 改为 $g_{j(x)}$。}
\typo{$0_{\mathrm{III}}$，10.2.7}{第20页第1行，将“$\mathfrak u$-进制”改为“$\mathfrak n$-进制”；同页倒数第13行，将 $k$ 改为 $k_i$。}
\typo{$0_{\mathrm{III}}$，11.1.1}{第23页第18行，将 $\Im(d_r^{p+r,q-r+1})$ 改为 $\Im(d_r^{p-r,q+r-1})$。}
\typo{$0_{\mathrm{III}}$，11.8.6}{第45页倒数第14行，将箭头 $\to$ 改为 $\rightsquigarrow$（四处）。}
\typo{$0_{\mathrm{III}}$，12.3.4}{第61页第4行，将 $\xrightarrow{\sim}$ 改为 $\to$。}
\typo{$0_{\mathrm{III}}$，13.2.4}{第67页倒数第2行，将 $\varinjlim$ 改为 $\varprojlim$（两处）。}
\typo{III，1.1.7}{第85页第15行，将 $A^r$ 改为 $\Lambda(A^r)$（两处）。}
\typo{III，1.2.4}{第87页第7行，将 $\Gamma(U,\F)$ 改为 $\Gamma(U,\OO_X)$，并将 $\G$ 改为 $\F$；同页第8行，将 $H^p(\mathfrak U,\G)$ 改为 $H^p(\mathfrak U,\F)$。}
\typo{III，2.5.3.1}{第111页第3行，将 $n>0$ 改为 $n\geq0$；同页第8行，将 $-r\leq n\leq0$ 改为 $-r\leq n<0$。}
\typo{III，3.1.2}{第116页第13行，将 $\G$ 改为 $\G\in\mathbf K'$。}
\typo{III，4.3.6}{第133页第7行，将 $K'=K\otimes_A\widehat A$ 改为 $K'\supset K\otimes_A\widehat A$；同页第16行，将 $R(X)$ 改为 $R(Y)$。}
\typo{III，4.4.12}{第138页第10行，将“第V章”改为“第IV章”。}
\typo{III，4.6.6}{第142页第19行，将“第V章”改为“在随后的一节中”。}

\section*{B．正文修改}
\addcontentsline{toc}{section}{原勘误清单二：B．正文修改}
\orig{为便于查找引文，今后将插入第N章的勘误清单中的正文修改，记为 $\operatorname{Err}_{N}$ 后接一个号码。}
\newcommand{\errIII}[1]{\par\addvspace{.6\baselineskip}\Needspace{3\baselineskip}\noindent\hypertarget{errIII:#1}{}\textbf{（$\operatorname{Err}_{\mathrm{III}}$，#1）}\quad\ignorespaces}

\errIII{1}在（$0_{\mathrm I}$，5.1.3）第45页第18行，将“直和”改为“有限直和”。

\errIII{2}在（$0_{\mathrm I}$，5.4.3）第49页倒数第8至第4行，自“在一般情形……”起，替换为下文：

在一般情形，设 $X$ 是赋环空间，每个 $x\in X$ 的 $\OO_x$ 都是局部环；若 $\LL$ 是有限型 $\OO_X$-模，且存在 $\OO_X$-模 $\F$，使 $\LL\otimes_{\OO_X}\F$ 同构于 $\OO_X$，则前面的论证说明，对每个 $x$，$\LL_x$ 都同构于 $\OO_x$。由此可知 $\LL$ 可逆。事实上，取 $x$ 的开邻域 $U$，使 $\LL|U$ 由 $U$ 上的 $n$ 个截面 $s_i$ 生成（5.2.1）。例如可设 $s=s_1$ 满足 $s_x\ne0$。\cmt{C3-A03}{底本这里写的是芽 $s_x\ne0$。仅非零并不足以推出它生成秩一自由模；所需的是选择一个在剩余模中非零的生成元，即 $s(x)\ne0$。这种选择由有限生成元在剩余模中仍生成而得到，因此可补足论证，不是否定本条结论。}由于 $\LL_x\simeq\OO_x$，对每个 $i$，在 $x$ 的某开邻域 $V_i\subset U$ 上存在 $\OO_X$ 的截面 $t_i$，使 $(t_is)_x=(s_i)_x$。因此存在开邻域 $V\subset\bigcap_{i=1}^nV_i$，使 $V$ 上 $t_is=s_i$，即 $\LL|V$ 由单个截面 $s|V$ 生成。再设 $z$ 为由 $s$ 定义的同态 $\OO_X|V\to\LL|V$ 的核在某开集 $W\subset V$ 上的截面（5.1.1）。对每个 $y\in W$，$z_y$ 零化 $\LL_y$，因而由假设 $z_y=0$，于是 $z=0$。断言得证。再考察 $\LL^{-1}\otimes\LL\otimes\F$，立即得到 $\F\simeq\LL^{-1}$。

\errIII{3}在（$0_{\mathrm I}$，7.2.4）中，必须增加条件：各幂 $\mathfrak J^n$ 是可容许环 $A$ 中的闭理想，该命题才正确。已给出的证明不正确；修正后的命题可由 Bourbaki，\textit{Top. gén.}，第III章，第3版，§3，第5目，命题9的推论1推出。在（$0_{\mathrm I}$，7.2.5）与（$0_{\mathrm I}$，7.2.6）的陈述中，同样须假设各 $\mathfrak J^n$ 在 $A$ 中闭。

\errIII{4}在命题（I，2.2.5）之后增加：沿用（I，2.2.5）的记号，若 $\mathfrak J$ 是 $A$ 的理想，则用 $\mathfrak J\OO_X$ 表示（0，4.3.5）中定义的 $\OO_X$-子模。

\errIII{5}在（I，3.4.5）第112页倒数第1行，将“一个域 $K$”改为“一个代数闭域 $K$”。第113页第1行，将“扩张”改为“代数闭扩张”。第113页第1至第4行，从“$K$ 称为……”起，替换为：

对于 $X$ 的任意一个取值于域 $K$ 的点，称 $K$ 为该点的值域；若 $x$ 是这个点的所在地，也说该点位于 $x$。由此定义映射 $X(K)\to X$，将一个取值于 $K$ 的点对应到它的所在地。

第113页第7与第16行删去“几何”；同页倒数第10与第13行也删去“几何”。第115页倒数第10与第11行删去“几何”。

\errIII{6}在（I，3.6.1）证明末尾，第117页倒数第12行，增加：若置 $X'=X\times_Y\Spec(\OO_y/\mathfrak a_y)$，则对每个经 $p$ 与 $X$ 的点相识别的 $x\in X'$，有
\[
 \OO_{X',x}=\OO_{X,x}/\mathfrak a_y\OO_{X,x}.
\]
事实上，问题在 $X,Y$ 上局部，故可设 $X=\Spec B$、$Y=\Spec A$；由平坦性，有 $(B/\mathfrak a B)_x=B_x/\mathfrak a B_x$（0，1.3.2）。

\errIII{7}在（I，5.1.1）第127页倒数第3行，将“拟凝聚 $\OO_X$-模”改为“$\B$ 的拟凝聚理想”。

\errIII{8}在（I，5.1.10）之后，第131页倒数第16行，增加：

\noindent\textbf{注记（5.1.11）。}对（5.1.9）的论证略作调整即可证明：不必预先假设 $X$ 是预概形，结论仍成立；只须假设 $(X,\OO_X)$ 是局部赋环空间，$\I$ 为 $\OO_X$ 的理想，$\I^n=0$，赋环空间 $X_0=(X,\OO_X/\I)$ 为仿射概形，且各 $\I^j/\I^{j+1}$ 为拟凝聚 $\OO_{X_0}$-模。在论证中，只须用（1.8.1）（见 $\operatorname{Err}_{\mathrm{II}}$）代替（2.2.4）。

\errIII{9}在（I，5.3.1）第132页倒数第11行，于“或 $\Delta_X$”之后插入“或 $\Delta_\varphi$，其中 $\varphi:X\to S$ 为结构态射”。

\errIII{10}在（I，5.3.9）中，所给证明不足，因为没有证明 $\Delta_X(X)$ 在 $X\times_SX$ 中局部闭。要得到正确证明，只须使用（4.2.4，a）：对每个 $x\in X$ 及 $x$ 在 $X$ 中的任意仿射邻域 $U$，$U\times_SU$ 是 $\Delta_X(x)$ 的仿射邻域；考虑（5.3.16）（其证明仅使用定义（5.3.1.1）），即可将（5.3.9）归结为 $S=\Spec B$、$X=\Spec A$ 均仿射的情形。此时由（5.3.1.1），$\Delta_X$ 对应于典范同态 $A\otimes_BA\to A$，$x\otimes y\mapsto xy$；这个同态满射，故 $\Delta_X$ 是闭浸入（4.2.3），证明完毕。
\cmt{C3-A02}{上句关于 $U\times_SU$ 仿射，须配合先将底限制到包含 $f(x)$ 的仿射开集，再取映入该开集的仿射 $U$；仅由 $U$ 仿射而不限制 $S$，并不能自动推出任意 $U\times_SU$ 仿射。此处保留作者勘误原文，在核查报告中写出完整局部化步骤。}

\errIII{11}在（I，7.1.15）及（I，7.1.16）第158页倒数第7及第15行，删去“几何”。

\errIII{12}将（I，7.4.7）替换为：将（7.4.1）的定义推广（略有滥用术语）到以下情形：$X$ 是约化预概形，且每个点都有一个仅含有限个不可约分支的开邻域。由（7.3.4）、（7.4.6）可知，对有限型拟凝聚 $\OO_X$-模 $\F$，称其为挠层，等价于 $\Supp(\F)$ 不包含 $X$ 的任何不可约分支。

\errIII{13}在（I，10.11.7）第205页第17至第23行，从“尚须证明……”起的证明结尾无须保留：由（0，5.3.5），所考虑的层都是凝聚层，故可直接应用（10.11.6）。

\errIII{14}在（II，1.7.8）第16页倒数第10行之后增加：当 $\E=\OO_S^n$ 时，也把 $\mathbf V(\E)$ 写成 $\mathbf V_S^n$；若 $S=\Spec A$ 仿射，则把 $\mathbf V_S^n$ 写成 $\mathbf V_A^n$；有 $\mathbf V_A^n=\Spec(A[T_1,\ldots,T_n])$，其中 $T_i$ 为不定元。

\errIII{15}在（II，1.7.10）第17页第10行及第12、16--17行，删去“几何”。

\errIII{16}在（II，1.7.12）中增加：同样约定
\[
 \mathbf V(\OO_S^n)=\mathbf V_S^n=S[T_1,\ldots,T_n].
\]

\errIII{17}在（II，4.2.6）第75页第7、8行，删去“几何”。

\errIII{18}在（II，4.6.13）第92页第7行，须将论证说得更精确：沿用（4.4.10）的记号，这里须证明浸入 $\Gamma_f$ 拟紧，才能应用（i bis）。由（4.6.4），为证明 $\LL$ 关于 $f$ 丰沛，可设 $Y$ 仿射。又在（v）的每一种假设下，$f$ 都拟紧（I，6.6.4）。若 $g$ 分离，则 $\Gamma_f$ 为闭浸入（I，5.4.3），因而拟紧（I，6.6.4）。若 $X$ 局部诺特，则因 $Y$ 仿射且 $f$ 拟紧，$X$ 的底空间拟紧，从而诺特；再次应用（I，6.6.4），可知 $\Gamma_f$ 拟紧。

\errIII{19}（II，5.2.2）的陈述不正确：若不知道对 $Y$ 的开集 $U$，每个拟凝聚 $(\OO_X|f^{-1}(U))$-模都是某个拟凝聚 $\OO_X$-模在 $f^{-1}(U)$ 上的限制，就不能说条件 b)、c)、c') 在 $Y$ 上是局部的。因此，应将第98页第16行的“设 $f:X\to Y$ 为分离拟紧态射”替换为：“设 $X,Y$ 为两个预概形，且 $X$ 是概形，或 $X$ 的底空间局部诺特；设 $f:X\to Y$ 为拟紧态射。”证明中应注意：在 $X$ 为概形的情形，由（I，5.5.5、5.5.8），这些假设蕴含 $f$ 分离，因此两种情形下都可应用（I，9.2.2）。

\errIII{20}在（II，6.2.3）第115页倒数第18行之后增加：若存在 $y=f(x)$ 的仿射开邻域 $V$ 及 $x$ 的仿射开邻域 $U$，使 $f(U)\subset V$，且限制态射 $U\to V$ 拟有限，则称态射 $f:X\to Y$ 在点 $x\in X$ 处拟有限。若它在每个 $x\in X$ 处都拟有限，则称 $f$ 局部拟有限。

\errIII{21}（II，6.4.3）的证明不正确：$K$ 的一个有限型 $A$-子模的元素，不一定在 $A$ 上整。将第121页最后九行替换为下文：因为可设 $E$ 无挠，因而忠实，$E$ 也是忠实 $A[u]$-模（这里 $A$ 嵌入 $E$ 的自同态环）。又 $E$ 是有限型 $A$-模，故 $u$ 在 $A$ 上整（Bourbaki，\textit{Alg. comm.}，第V章，§1，第1目，引理1）。立即推出 $u\otimes1$ 的特征值（在 $K$ 的某个代数闭包中）在 $A$ 上整，因而各 $a_i(u)$ 也在 $A$ 上整。

\errIII{22}在（$0_{\mathrm{III}}$，9.1.1）第12页倒数第18行，删去“开”。

\errIII{23}在（$0_{\mathrm{III}}$，11.7.3）第43页倒数第10行，于“$C''$”之后增加：使得对 $C$ 的每个投射对象 $P$（分别对 $C'$ 的每个投射对象 $P'$），函子 $A'\mapsto T(P,A')$（分别 $A\mapsto T(A,P')$）在 $C'$（分别 $C$）中正合。

\errIII{24}（$0_{\mathrm{III}}$，13.7.7）的陈述不准确，须作如下修改：第78页第6行，删去“以及 $(R^{n+1}u(A_i))_{i\in I}$”；同页第7行，将“$R^nu(A)$ 是有限型 $S$-模”改为“$R^nu(A)$ 与 $R^{n+1}u(A)$ 均为有限型 $S$-模”；同页第12--13行，删去“以及 $p+q=n+1$”。证明中，第78页第23行，删去“以及对于 $n+1$”。

\errIII{25}在（III，1.4.15）第92页倒数第6行，将“有限型”改为“拟紧”。

\errIII{26}在（III，2.2.4）第102页第1行，将“$\F$ 与 $\G$ 的支集在 $Y$ 上固有”改为“$\Im(\F\to\G)$ 的支集在 $Y$ 上固有”。证明中，从“既如此……”开始的第10至第16行，替换为下文：

置 $\G_1=\Im(\F\to\G)=\Ker(\G\to\HH)$，它凝聚。由正合列 $0\to\G_1(n)\to\G(n)\to\HH(n)$ 及 $f_*$ 左正合，对每个 $n$，序列 $0\to f_*(\G_1(n))\to f_*(\G(n))\to f_*(\HH(n))$ 正合。因此，只须证明 $n$ 充分大时，$f_*(\F(n))\to f_*(\G_1(n))\to0$ 正合。换言之，可设 $\HH=0$，且 $\G$ 的支集在 $Y$ 上固有，因而在 $Z$ 中闭（II，5.4.10）。于是 $\G'=i_*(\G)$ 为凝聚 $\OO_Z$-模，且 $\G'|X=\G$。由（I，9.4.3），存在凝聚 $\OO_Z$-模 $\F'$，使 $\F'|X=\F$。又因 $X$ 在 $Z$ 中开，且 $\Supp(\G)\subset X$ 在 $Z$ 中闭，可定义满同态 $u':\F'\to\G'$：当开集 $U\subset Z$ 不交 $\Supp(\G)$ 时取 $u'|U=0$，当 $U\subset X$ 时取 $u'|U=u|U$。所得限制 $u'|X$ 正是给定的满同态 $u:\F\to\G$。

既如此，如（2.2.2）证明开头，可设 $Y$ 仿射。只须证明 $n$ 充分大时，$\Gamma(u):\Gamma(X,\F(n))\to\Gamma(X,\G(n))$ 满射。有交换图
\[
\begin{tikzcd}[column sep=large]
\Gamma(Z,\F'(n))\arrow[r,"\Gamma(u')"]\arrow[d,"v"']&\Gamma(Z,\G'(n))\arrow[d,"w"]\\
\Gamma(X,\F(n))\arrow[r,"\Gamma(u)"']&\Gamma(X,\G(n)).
\end{tikzcd}
\]
其中 $v,w$ 是限制同态。由 $\G'$ 的定义，$w$ 双射；由（2.2.3），$n$ 充分大时 $\Gamma(u')$ 满射，故 $\Gamma(u)$ 也满射。

\errIII{27}在（III，2.2.5）第102页倒数第7行之后增加：

（iii）在（2.2.4）关于 $X,Y,f,\LL$ 的假设下，若 $\F$ 是其支集在 $Y$ 上固有的凝聚 $\OO_X$-模，类似（2.2.4）的论证给出整数 $N$，使 $n\geq N$ 时，对所有 $i>0$ 均有 $R^if_*(\F(n))=0$。由上同调正合列可知：若 $u:\F\to\G$ 为凝聚 $\OO_X$-模的同态，且 $\Ker u$ 与 $\Coker u$ 的支集均在 $Y$ 上固有，则存在 $N$，使 $n\geq N$ 时，对所有 $i>0$，相应同态 $R^if_*(\F(n))\to R^if_*(\G(n))$ 均双射。

\errIII{28}在（III，4.4.9）第137页倒数第17行，将“设 $Y$ 为整的局部诺特预概形”改为“设 $X,Y$ 为两个整的局部诺特预概形”；同页倒数第16行，将“有限型”改为“局部有限型”。证明实质上不变：$f^{-1}(y)$ 在 $k(y)$ 上局部有限型，因而仍是离散的。

\errIII{29}在（I，9.3.4）第173页倒数第19行，将“诺特”改为“拟紧”；同页倒数第18、19行，将“凝聚”改为“有限型拟凝聚”。证明中，归结为 $X=\Spec A$、$\F=\widetilde M$，其中 $M$ 是有限型 $A$-模，$\J=\widetilde{\mathfrak J}$，其中 $\mathfrak J$ 是 $A$ 的有限型理想；其余论证不变。

\errIII{30}将（I，9.3.5）第173页最后七行替换为下文：

\noindent\textbf{命题（9.3.5）。}设 $X$ 为预概形，$\F$ 为有限型拟凝聚 $\OO_X$-模。则存在 $X$ 的闭子预概形 $Y$，其底空间等于 $\Supp(\F)$；还存在有限型拟凝聚 $\OO_Y$-模 $\G$，使得若 $j:Y\to X$ 为典范包含，则 $\F\simeq j_*(\G)$。

只须证明 $\F$ 的零化理想层 $\J$ 拟凝聚。取 $Y$ 为由 $\J$ 定义的闭子预概形（I，4.1.2）；因 $\J\F=0$，$\F$ 是 $\OO_X/\J$-模，取 $\G=j^*(\F)$ 即可。为证明 $\J$ 拟凝聚，可局部化为 $X=\Spec A$、$\F=\widetilde M$，其中 $M$ 由有限个元素 $x_i$（$1\leq i\leq r$）生成。于是 $\J$ 是各 $x_i$ 的零化理想层的交。$x_i$ 的零化理想层是同态 $\OO_X\to\F$ 的核，该同态对应于 $A\to M$、$s\mapsto sx_i$；故它拟凝聚（I，4.1.1），这类理想层的有限交也拟凝聚（I，1.3.10）。

% END backmatter/errata_III_2.tex


% BEGIN backmatter/bibliography_IV_1.tex
\head{第四卷第一部分：参考文献续表}
\source{EGA IV第一部分，印刷第251页（期刊连续页码347）；上传PDF第250页}
\NewNote{本页为该分册的三条参考文献续表全文。书目题名译成中文并保留原题名以便核对；这一页的收入不表示该分册其余正文和索引已译完。}
\begin{description}[style=nextline,leftmargin=2.8em,labelwidth=2.4em]
\item[{[30]}] M. Nagata，\textit{局部环}（\textit{Local rings}），\textit{Interscience Tracts}，13（1962），Interscience，纽约。
\item[{[31]}] M. Nagata，\textit{单纯点的一个雅可比判据}（\textit{A Jacobian criterion of simple points}），\textit{Ill. Journ. of Math.}，1（1957），427--432页。
\item[{[32]}] A. Néron，\textit{局部域和整体域上阿贝尔簇的极小模型}（\textit{Modèles minimaux des variétés abéliennes sur les corps locaux et globaux}），\textit{Publ. Inst. Hautes Études Sci.}，第21号（1964）。
\end{description}

% END backmatter/bibliography_IV_1.tex


% BEGIN backmatter/errata_IV_selected.tex
\head{第四卷勘误清单三：与第0章§§16--19已译部分有关的选译}
\source{EGA IV第四部分，勘误清单三，印刷第345、347页；上传PDF第344、346页}
\NewNote{本节只收入与已译范围直接相关的七项作者勘误，保留其原分类和原编号。它不是整份勘误清单三的全译；未列出的条目不因此被删除或判定无关。下文的“页”均指作者所列原出版分册的印刷页码。}

\egasection{A）排印错误：三项选译}
\hypertarget{ErrIV:A16222}{}\noindent\textbf{（$0_{\mathrm{IV}}$，16.2.2.2）。}
第26页第13行，将 $M'\cap\mathfrak q^nM'$ 改为 $M'\cap\mathfrak q^nM$。

\medskip
\hypertarget{ErrIV:A1737}{}\noindent\textbf{（$0_{\mathrm{IV}}$，17.3.7）。}
第50页倒数第16行，将（17.1.10）改为（17.3.3）。

\medskip
\hypertarget{ErrIV:A1813}{}\noindent\textbf{（$0_{\mathrm{IV}}$，18.1.3）。}
在图（18.1.3.1）中，$p_2$ 应标在箭头 $G\to F$ 的旁边。

\egasection{B）正文修订：四项选译}
\hypertarget{ErrIV:5}{}
\noindent\textbf{（Err$_{\mathrm{IV}}$，5）。}
在（$0_{\mathrm{IV}}$，16.1.4）中，第23页倒数第8行，将“说 $A$ 是悬链的”改为：“说一个整环 $A$ 是悬链的”。

\medskip
\hypertarget{ErrIV:6}{}
\noindent\textbf{（Err$_{\mathrm{IV}}$，6）。}
在（$0_{\mathrm{IV}}$，16.1.9）中，可删去 $A$ 为诺特环的假设。事实上，对于每个 $u\in B_M$，$A[u]$ 是 $C$ 的一个子 $A$-代数，故 $u$ 整于 $A$（Bourbaki，\textit{交换代数}，第五章，§1，第1目，命题1的推论1），因而也整于 $A_M$；仍由（16.1.5）得出结论。

\medskip
\hypertarget{ErrIV:7}{}
\noindent\textbf{（Err$_{\mathrm{IV}}$，7）。}
在（$0_{\mathrm{IV}}$，16.5.12）中，应注意：关系（16.5.11.1）等价于将 $A$ 换成各商环 $A/\mathfrak r$ 后得到的相同关系，其中 $\mathfrak r$ 遍历 $A$ 的极小素理想。因此，在（16.5.12）的证明中可限于 $A$ 为整环的情形，并应用（16.1.4.2）。


\medskip
\hypertarget{ErrIV:8}{}
\noindent\textbf{（Err$_{\mathrm{IV}}$，8）。}
删去（$0_{\mathrm{IV}}$，19.3.12），因为其证明错误（集合 $W_\lambda,U_\lambda,V_\lambda$ 不是线性流形）。

% END backmatter/errata_IV_selected.tex

\part*{译稿来源与版本记录}
\addcontentsline{toc}{part}{译稿来源与版本记录}
\markboth{译稿来源与版本记录}{译稿来源与版本记录}
\NewNote{下面三份说明来自既有分卷，保留其中关于当时稿件范围和制作过程的历史记录。旧页数、旧接续点不代表本合订本的当前页数或当前完成度；本版范围以卷首说明及工程包PROGRESS.md为准。}

% BEGIN backmatter/old_note_I.tex
\chapter*{第一卷既有译稿说明}
\addcontentsline{toc}{chapter}{第一卷既有译稿说明}
\markboth{译稿说明}{译稿说明}
本译稿以用户上传的 \texttt{EGAI.pdf} 为底本。原著为 A. Grothendieck 与 J. Dieudonné
合作撰写的 \textit{Éléments de géométrie algébrique：I. Le langage des schémas}，
刊于 \textit{Publications mathématiques de l'I.H.É.S.} 第 4 卷（1960），5—228 页。
本书中的未来计划、历史叙述及跨卷引文，均保留原著写作时的语境。

\medskip
\noindent\textbf{译法与记号。}
沿用“典范”“归纳极限”“射影极限”“诺特”“有限型”“拟凝聚层”“凝聚层”等译法。
层的一点处的 \textit{fibre} 译为“茎”，截面的 \textit{germe} 译为“芽”；
态射的 \textit{fibre} 译为“纤维”。
\textit{préschéma} 与 \textit{schéma} 分别译为“预概形”与“概形”，
不将本版的定义替换为其他版本的定义。
\textit{préadique} 与 \textit{adique} 分别译为“预进制”与“进制”。
条目及公式编号依照原文；不带章号的交叉引用指其所在章，“（0，…）”指第 0 章。
中文分页不与原文逐页对应；各小节注明原文印刷页码以便对照。

\medskip
\noindent\textbf{原文与译注。}
译文保留原著的组织、证明详略与引文。原著只引用、略去或留给读者的证明，不自行补写。
原作者的脚注标为“原注”；译者对术语、记号和可辨疑点的说明标为“译注”。
扫描文本的识别错误按页面图像恢复；涉及原版自身疑点的处理须另作明确说明，
不将译者的解释冒充原著内容。

\medskip
\noindent\textbf{版本说明。}
使用实际的 \texttt{ctexbook} 文档类及其默认字体配置，以 XeLaTeX 编译。
中文翻译与排版由 ChatGPT 生成，供研读及原文对照使用，未经原作者或出版机构审定。
\par\medskip\noindent\textbf{收录范围。}本合订本收入本卷导言、第 0 章 §§1—7、第 I 章 §§1—10，以及原著参考文献、记号索引、术语索引和书末资料。


\medskip
\noindent\textbf{V2 校订范围（2026 年 10 月 6 日）。}
本次在第 0 章（4.1.5）、（6.7.8），以及第 I 章（5.3.9）、（7.4.4）
修订或扩充译注：分别给出外幂、忠实平坦性、任意无挠层情形的反例，并补足对角像局部闭性的论证。
反例与补证均属译者说明，原著正文没有因此被静默改写。
另为正文中可按章号和条目号确定目标的交叉引文增加点击跳转；
跳转表示该编号的查阅位置，不代表已独立验证每个引文在论证中的适用性。
原著索引中的旧编号保留原样，不为其自动建立可能误导的链接。

本版附独立的《第一轮校订报告》，列出反例的完整计算、保留的待核问题及实际编译检查结果。
全卷结构和链接目标的核对，不等同于全卷逐句翻译复核，也不等同于所有数学证明已获独立审定。

% END frontmatter_full.tex




% END backmatter/old_note_I.tex


% BEGIN backmatter/old_note_II.tex
\chapter*{第二卷既有译稿说明}
\addcontentsline{toc}{chapter}{第二卷既有译稿说明}
本译稿依据用户提供的 \texttt{EGAII.pdf}：A. Grothendieck 与 J. Dieudonné 合作撰写的 \textit{Éléments de géométrie algébrique, II: Étude globale élémentaire de quelques classes de morphismes}，载于 \textit{Publications mathématiques de l'I.H.É.S.} 第8卷（1961），原书印刷页码5--222，上传扫描文件共219页。

\textbf{收录范围。} 本合订本使用已恢复的连续正文，最后编号条目为（8.14.14）及其证明；书末资料依据同一上传底本重新译排，详见合订说明。

\textbf{组织与记号。} 沿用原著章节、条目、公式和交叉引文编号。各小节标注原文印刷页码，区别于扫描页序和译稿页码。“0”“I”“III”等章号保持原义；第二章中不带章号的引文仍指本章。保留“预概形／概形”“拟凝聚／凝聚”“双同态”等既有译法，不把原著历史定义直接替换为现代定义。分次 Hom 仍依（2.1.2）解释；保留 $S_f$ 与 $S_{(f)}$、$M_f$ 与 $M_{(f)}$ 的区别。

\textbf{译注与审校。} 公式的扫描识别错误依据原页还原。原著脚注标作“原注”，另加的疑点说明标作“译注”。不在正文中静默补入原文没有的假设；原著论证中的省略、原版疑点与译者说明分开处理。编译成功与版面检查，不等于全部数学内容经过独立审定。本稿为研读用工作译稿，尚未经独立出版级逐句审校。

\textbf{排版与源码。} 采用实际的 \texttt{ctexbook} 文档类，未指定其他中文或数学字体。公式与交换图使用 TeX 排版。合订本的单文件源码可独立编译，不依赖原版PDF或外部截图。请使用 XeLaTeX 连续编译至少两次；完整工程另外保留分节源文件和进度记录。

\subsection*{底本版权说明}
下面保留底本版权文字；不表示原作者或出版者审定或认可本中文译稿。
\begin{quote}\small
© Publications mathématiques de l'I.H.É.S., 1961, tous droits réservés.

L'accès aux archives de la revue « Publications mathématiques de l'I.H.É.S. » (\url{http://www.ihes.fr/IHES/Publications/Publications.html}), implique l'accord avec les conditions générales d'utilisation (\url{http://www.numdam.org/legal.php}). Toute utilisation commerciale ou impression systématique est constitutive d'une infraction pénale. Toute copie ou impression de ce fichier doit contenir la présente mention de copyright.

Article numérisé dans le cadre du programme Numérisation de documents anciens mathématiques.
\end{quote}



% END backmatter/old_note_II.tex


% BEGIN backmatter/old_note_III_1.tex
\chapter*{第三卷第一部分既有译稿说明}
\addcontentsline{toc}{chapter}{第三卷第一部分既有译稿说明}
本译稿以用户上传的 \texttt{EGAIII.1.pdf} 为底本。该文件共164个 PDF 页面，所载作品为《代数几何基础》第III卷《凝聚层的上同调研究》第一部分。译文包括第0章补篇 §§8—13、第III章内容提要及 §§1—5、记号索引、参考文献续表、术语索引与原著书末目录。正文止于第III章（5.5.2）及其证明。原书的空白页未在译稿中机械复制。

原著第III章的内容提要预告了 §§1—11，但本次底本只收入 §§1—5。本译稿是\textbf{这份“第一部分”文件的完整工作译稿}，不是整个第III章的完整译本；卷末“待续”为原著结束语。本卷现与前两卷统一编排。

条目、命题及公式采用原著编号；第0章中与第一章同时发表的内容仍按原著记作 $0_{\mathrm I}$，不带这一角标的第0章引文根据上下文指本补篇或原著约定的对应部分。原著的预概形／概形、拟凝聚／凝聚、双同态、丰沛等术语区别均保留，不以现代习惯改写历史定义。法文 \textit{propre} 在本稿中译作“固有”，\textit{très ample} 译作“很丰沛”。

正文保持原著的陈述与证明层次；原著省略的证明仍按原文省略。脚注分别冠以“原注”和“译注”。识别出的原版符号、复合次序、引文或假设疑点，以译注说明；译注是翻译过程中作出的文本分析，不属于原著，也不代替独立数学审校。原著关于“以后”“将证明”及研究问题的措辞均处于1961年的写作语境，不表示现今研究状况。

小节开头标出原著印刷页码范围，便于对照。原稿同时带有本卷页码与另一套累计页码；除索引前特别说明外，本译稿的“原文印刷页码”指本卷页码。记号索引与参考文献在上传文件中存在页序与印刷页码不一致的情况，已注明，原著目录中的页码不作静默调换。

本稿采用 \texttt{ctexbook} 文档类及其默认中文字体配置，使用 XeLaTeX 编译；公式与交换图均由 TeX 排版，不使用底本截图替代正文。卷首目录使用译稿页码，卷末另附保留原页码的原著目录。本稿供研读与对照使用，尚未经出版机构或独立数学审校者审定。

\chapter*{底本出版与版权说明}
\addcontentsline{toc}{chapter}{底本出版与版权说明}
A. Grothendieck（与 J. Dieudonné 合作撰写）：\textit{Éléments de géométrie algébrique : III. Étude cohomologique des faisceaux cohérents, Première partie}，\textit{Publications mathématiques de l'I.H.É.S.}，第11卷，1961年，第5—167页。

底本版权声明的中文译文如下：

\begin{quote}
© Publications mathématiques de l'I.H.É.S.，1961年，保留所有权利。

使用本刊档案，意味着同意其一般使用条款。任何商业使用或系统性印刷均构成刑事违法。本文件的任何复制件或印刷件均须包含本版权声明。

本文在“古数学文献数字化”计划框架内数字化。
\end{quote}
以上为底本所载声明的翻译；另保留法文原声明及其所列地址，以供核对：
\begin{quote}\small
© Publications mathématiques de l'I.H.É.S., 1961, tous droits réservés.
L'accès aux archives de la revue « Publications mathématiques de l'I.H.É.S. »
(\url{http://www.ihes.fr/IHES/Publications/Publications.html}), implique l'accord avec les conditions générales d'utilisation
(\url{http://www.numdam.org/legal.php}). Toute utilisation commerciale ou impression systématique est constitutive d'une infraction pénale. Toute copie ou impression de ce fichier doit contenir la présente mention de copyright.

Article numérisé dans le cadre du programme Numérisation de documents anciens mathématiques.
\url{http://www.numdam.org/}
\end{quote}\par

% ---- END frontmatter.tex ----


% ---- BEGIN c0_s08.tex ----

% END backmatter/old_note_III_1.tex


% BEGIN backmatter/additional_sources.tex
\head{来源、署名与版权说明}
本册翻译依据为用户提供的文件本身；数学核查与原文翻译分别标示。本文没有访问未提供的法文手稿，也没有把网络上其他译本的补写移入原文。

\egasection{已译部分使用的底本}
A. Grothendieck，与J. Dieudonné合作撰写：\textit{Éléments de géométrie algébrique III. Étude cohomologique des faisceaux cohérents, Seconde partie}，\textit{Publications mathématiques de l'I.H.É.S.} 17（1963），5--91。用户文件：\texttt{EGAIII.2.pdf}。

A. Grothendieck，与J. Dieudonné合作撰写：\textit{Éléments de géométrie algébrique IV. Étude locale des schémas et des morphismes de schémas, Première partie}，\textit{Publications mathématiques de l'I.H.É.S.} 20（1964），5--259。用户文件：\texttt{EGAIV.1.pdf}。

\textit{Alexandre Grothendieck's EGA V. Translation and Editing of his `prenotes'}，英译及编辑Piotr Blass、Joseph Blass。用户文件：\texttt{EGAVaiello.pdf}。书内自称预备笔记，保留多处编写提示及待核说明；不附会一个本载体没有给出的正式出版年份。

\textit{EGA V: §1 and §§2.15, 2.16 (formerly numbered as EGA IV: §16 and §§17.15, 17.16)}，Grothendieck预备笔记的编辑英译，Piotr Blass、Joseph Blass。用户文件：\texttt{EGAV1-2.ps}，此前转换为六页PDF后对照。\texttt{EGAV6.ps} 的（1.1）只用于核查摘译；其余内容未计入连续译文。

第四版还直接核对A. Grothendieck与J. Dieudonné合作撰写的\textit{Éléments de géométrie algébrique IV, Quatrième partie}，\textit{Publications mathématiques de l'I.H.É.S.} 32（1967），5--361，用户文件\texttt{EGAIV.4.pdf}中的相关作者勘误；仅选译了与本轮§§16--17对应的五项，不表示全分册已译。该底本标示\copyright\ 1967 Publications mathématiques de l'I.H.É.S.／Institut des Hautes Études Scientifiques，保留所有权利。

\egasection{外部核查资料}
Katharina Heinrich，\textit{Some remarks on biequidimensionality of topological spaces and Noetherian schemes}，\href{https://arxiv.org/abs/1403.5814}{arXiv:1403.5814}（2014）。此前用于核对EGA $0_{\mathrm{IV}}$（14.3.3）的错误性质；本文的六点反例另行逐步说明。

The Stacks Project，\href{https://stacks.math.columbia.edu/tag/0FD4}{Tag 0FD4：Bertini theorems}，特别是Lemma 33.47.3。仅用于核对明确浸入假设下的标准结论与证明路线，不据此把所有草稿变体宣称为已证明。

\egasection{保留的版权信息}
法文底本所附说明标示相应年度 \copyright\ Publications mathématiques de l'I.H.É.S.，保留所有权利，并要求复制件保留版权提示。III第二部分另有 \copyright\ 1963 Institut des Hautes Études Scientifiques；IV第一部分另有 \copyright\ 1964 Institut des Hautes Études Scientifiques。底本由Numdam古数学文献数字化项目提供；原文件中的使用条件及原权利标示没有因本次翻译消失。

EGA V载体的原著署名、英译编辑署名及编者疑问照原件保留。本册仅为用户提供材料的中文工作译稿及核查，未获原作者、译编者或出版机构审定，不表示拥有原著版权或正式出版许可。

% END backmatter/additional_sources.tex

\end{document}
